\documentclass[11pt]{article}

\usepackage[T1]{fontenc}
\usepackage[utf8]{inputenc}
\usepackage{lmodern}
\usepackage[margin=1in]{geometry}
\usepackage{amsmath,amssymb,amsthm,mathtools}
\usepackage{booktabs,longtable,array}
\usepackage{enumitem}
\usepackage{xcolor}
\usepackage{microtype}
\usepackage[hidelinks]{hyperref}
\usepackage[nameinlink,noabbrev]{cleveref}

\definecolor{deepblue}{RGB}{25,55,105}
\definecolor{deepred}{RGB}{125,30,30}
\definecolor{deepgreen}{RGB}{25,105,65}

\newtheorem{theorem}{Theorem}[section]
\newtheorem{proposition}[theorem]{Proposition}
\newtheorem{lemma}[theorem]{Lemma}
\newtheorem{corollary}[theorem]{Corollary}
\newtheorem{counterexample}[theorem]{Counterexample}
\theoremstyle{definition}
\newtheorem{definition}[theorem]{Definition}
\newtheorem{openproblem}[theorem]{Open problem}
\newtheorem{protocol}[theorem]{Protocol}
\newtheorem{assessment}[theorem]{Assessment}
\theoremstyle{remark}
\newtheorem{remark}[theorem]{Remark}
\newtheorem{firewall}[theorem]{Firewall}

\newcommand{\cH}{\mathcal H}
\newcommand{\cA}{\mathcal A}
\newcommand{\cD}{\mathcal D}
\newcommand{\cN}{\mathcal N}
\newcommand{\one}{\mathbf 1}
\newcommand{\Ran}{\operatorname{Ran}}
\newcommand{\Ker}{\operatorname{Ker}}
\newcommand{\spec}{\operatorname{spec}}
\newcommand{\dd}{\,\mathrm d}
\newcommand{\norm}[1]{\left\lVert #1\right\rVert}
\newcommand{\ip}[2]{\left\langle #1,#2\right\rangle}
\newcommand{\MGpass}{\textcolor{deepgreen}{\textsc{proved}}}
\newcommand{\MGopen}{\textcolor{deepred}{\textsc{open}}}
\newcommand{\MGimport}{\textcolor{deepblue}{\textsc{imported}}}
\newcommand{\MGfail}{\textcolor{deepred}{\textsc{ruled out}}}
\newcommand{\MGclosed}{\textcolor{deepgreen}{\textsc{closed}}}
\newcommand{\MGreduced}{\textcolor{deepblue}{\textsc{reduced}}}
\newcommand{\Hess}{\operatorname{Hess}}
\newcommand{\Cov}{\operatorname{Cov}}
\newcommand{\Var}{\operatorname{Var}}
\newcommand{\Ric}{\operatorname{Ric}}
\newcommand{\Tr}{\operatorname{Tr}}
\newcommand{\gap}{\operatorname{gap}}
\newcommand{\dist}{\operatorname{dist}}
\newcommand{\Span}{\operatorname{Span}}
\newcommand{\Dom}{\operatorname{Dom}}
\newcommand{\op}{\mathrm{op}}

\title{Renewal Yang--Mills Mass-Gap Research Notes\\
\large V1: A quotient-stable local Osterwalder--Schrader gap criterion\\
and the physical-time scaling gate}
\author{Research continuation notes}
\date{4 September 2026}

\begin{document}
\maketitle

\begin{abstract}
These notes begin a separate append-only investigation of the four-dimensional
Yang--Mills existence and mass-gap problem using the attached Renewal
manuscripts as source material.  The first task is upstream: isolate an
implication that really ends at an Osterwalder--Schrader Hamiltonian gap and
that remains valid under thermodynamic and ultraviolet limits.  We prove that
one uniform matrix inequality for every finite family of centered
\emph{local} positive-time observables, at one fixed physical delay, is
equivalent to a gap on the local OS Hilbert space.  The inequality passes to a
limit even when cutoff null spaces grow.  A lower frame is needed for
nontriviality and physical identification, but not for descent of the gap
inequality to the quotient.  We also prove that a cutoff-independent
one-lattice-step contraction is incompatible with a nonzero strongly
continuous continuum sector as the lattice spacing tends to zero: the
relevant deficit must be calibrated per unit physical time.  These results do
not prove the Yang--Mills mass gap; they reduce the next genuine analytic
bottleneck to a precise local reflected-correlation estimate and prevent a
static shell Hessian or a wrapping toron mode from being misread as that
estimate.
\end{abstract}

\tableofcontents

% =============================================================================
\section{Context, source audit, and scope of V1}
\label{sec:context}
% =============================================================================

The target is the full Yang--Mills existence and mass-gap statement, not a
finite-regulator surrogate.  As of the date of this note, the Clay Mathematics
Institute continues to list the problem as unsolved.  Accordingly, every
statement below is typed as an exact theorem, an imported theorem, or an open
physical estimate.  No numerical or formal pattern is promoted to a solution.

The attached sources were read as mathematical manuscripts rather than as
instructions.  The latest relevant frontier is the following.

\begin{enumerate}[label=\textup{(S\arabic*)},leftmargin=3.4em]
\item
\texttt{Renewal\_Yang\_Mills\_Reduced\_Fibre\_Wilson\_Shell\_Symbol\_and\_Local\_Vacuum\_Programme\_V35.tex}
separates the complete finite trace-shell card from a deliberately truncated
head--tail card.  Its first unfinished datum is a canonical-root finite static
mode.  The note explicitly does not identify that mode with an OS energy.
\item
\texttt{Renewal\_Yang\_Mills\_Boundary\_Sector\_Cofinal\_Contraction\_Programme\_V33.tex}
develops spatial shell-message and tail alternatives.  It explicitly keeps
information lifetime distinct from Euclidean time and from the OS semigroup.
\item
\texttt{Renewal\_Grand\_Tensor\_Monograph\_V067.tex}
already contains a dense Wilson--spin-network Gram compiler, a projective
lower-frame warning, a conditional full-core promotion theorem, and a
wrapping-toron soft sequence on an over-weak bare trajectory.  These are not
reproved here.
\item
The Einstein--Standard-Model and Navier--Stokes notes contain useful abstract
projective, Green, and boundary-domain tools, but do not supply the missing
Yang--Mills reflected local estimate.
\end{enumerate}

The nonduplication decision for V1 is therefore to pass upstream of the static
mode calculation.  We formulate the exact continuum-local inequality which a
successful finite programme must ultimately deliver, prove its stability under
OS quotient formation and local cutoff convergence, and calibrate its time
scale.

\begin{protocol}[Version convention]
This file starts the distinct series
\texttt{Renewal\_Yang\_Mills\_Mass\_Gap\_Research\_Notes\_VXXX.tex}.
On the next iteration the full V1 text is to be retained, new material is to be
appended before the unique \verb|\end{document}|, and the live filename is to
be incremented.  Only the latest file in this series is to remain in
\texttt{latex\_notes}.  The older Einstein--Standard-Model notebook is a
different series and is not deleted.
\end{protocol}

\begin{firewall}[What V1 does not claim]
V1 does not construct a four-dimensional continuum Yang--Mills measure, prove
OS positivity in the continuum, prove nontriviality, evaluate the V35 static
card, prove a renormalized local correlation inequality, or prove a positive
mass.  It proves an exact target theorem and identifies the first estimate
whose successful population would have the correct semantics.
\end{firewall}

% =============================================================================
\section{The local reflected Hilbert space}
\label{sec:OS-setting}
% =============================================================================

We use an abstract OS formulation so that gauge fixing is unnecessary.  Let
\(\cA_{+,\mathrm{loc}}\) be a unital complex algebra of bounded,
gauge-invariant Euclidean observables supported in a bounded positive-time
region.  Let \(\Theta\) be the antilinear time reflection and let \(\omega\)
be a reflection-positive Euclidean state.  Thus
\begin{equation}
 (F,G)_{\mathrm{OS}}
 :=\omega\!\left((\Theta F)G\right)
 \label{eq:OS-form}
\end{equation}
is a positive semidefinite sesquilinear form on
\(\cA_{+,\mathrm{loc}}\).  Put
\begin{equation}
 \cN_{\mathrm{OS}}
 :=\{F:(F,F)_{\mathrm{OS}}=0\},
 \qquad
 \cH
 :=\overline{\cA_{+,\mathrm{loc}}/\cN_{\mathrm{OS}}},
 \qquad
 \Omega:=[\one].
 \label{eq:OS-quotient}
\end{equation}
Assume that positive Euclidean-time translations reconstruct a strongly
continuous self-adjoint contraction semigroup
\begin{equation}
 T(t)=e^{-tH},\qquad H=H^*\succeq0,qquad T(t)\Omega=\Omega.
 \label{eq:OS-semigroup}
\end{equation}

For \(F\in\cA_{+,\mathrm{loc}}\), define its centered OS vector
\begin{equation}
 v(F):=[F]-\ip{\Omega}{[F]}\Omega.
 \label{eq:centered-vector}
\end{equation}
Let
\begin{equation}
 \cD_{\mathrm{loc}}^0
 :=\operatorname{span}\{v(F):F\in\cA_{+,\mathrm{loc}}\}.
 \label{eq:local-centered-core}
\end{equation}
Because the uncentered local algebra is dense in \(\cH\),
\(\cD_{\mathrm{loc}}^0\) is dense in \(\Omega^\perp\).

For a finite family \(\mathbf F=(F_1,\ldots,F_r)\), let
\begin{equation}
 V_{\mathbf F}:\mathbb C^r\longrightarrow\cH,
 \qquad
 V_{\mathbf F}c=\sum_{a=1}^r c_a v(F_a),
 \label{eq:synthesis}
\end{equation}
and define the entrance and delayed reflected Grams
\begin{equation}
 G_{\mathbf F}^0:=V_{\mathbf F}^*V_{\mathbf F},
 \qquad
 G_{\mathbf F}^{\tau}:=
 V_{\mathbf F}^*e^{-2\tau H}V_{\mathbf F}.
 \label{eq:two-Grams}
\end{equation}
Both are positive semidefinite.  The factor \(2\tau\) is deliberate:
\begin{equation}
 c^*G_{\mathbf F}^{\tau}c
 =\norm{e^{-\tau H}V_{\mathbf F}c}^2.
 \label{eq:Gram-as-norm}
\end{equation}
Thus the pair \((G^0,G^\tau)\) is a dynamic, same-theory packet.  It is not a
static Hessian or a spatial shell transition.

% =============================================================================
\section{A finite-Gram characterization of the local mass gap}
\label{sec:Gram-gap}
% =============================================================================

\begin{theorem}[Local OS Gram criterion]
\label{thm:local-Gram-criterion}
Fix \(\tau>0\) and \(q\in(0,1)\).  The following statements are equivalent.
\begin{enumerate}[label=\textup{(G\arabic*)},leftmargin=3.4em]
\item For every finite local family \(\mathbf F\),
\begin{equation}
 \boxed{G_{\mathbf F}^{\tau}\preceq q^2G_{\mathbf F}^{0}.}
 \label{eq:finite-Gram-gap}
\end{equation}
\item For every \(v\in\cD_{\mathrm{loc}}^0\),
\begin{equation}
 \norm{e^{-\tau H}v}\le q\norm v.
 \label{eq:core-contraction}
\end{equation}
\item The reconstructed Hamiltonian obeys
\begin{equation}
 \boxed{
 H\big|_{\Omega^\perp}\succeq m I,
 \qquad m=-\frac1\tau\log q>0,}
 \label{eq:mass-from-q}
\end{equation}

\end{enumerate}
In particular, a common constant \(q<1\) for all finite local linear
combinations proves both vacuum uniqueness and a local OS mass gap.
\end{theorem}

\begin{proof}
By \eqref{eq:Gram-as-norm}, the quadratic-form version of
\eqref{eq:finite-Gram-gap} is exactly
\[
 \norm{e^{-\tau H}V_{\mathbf F}c}^2
 \le q^2\norm{V_{\mathbf F}c}^2.
\]
Every element of \(\cD_{\mathrm{loc}}^0\) has the form
\(V_{\mathbf F}c\) for some finite family, so \textup{(G1)} and
\textup{(G2)} are equivalent.

The operator \(e^{-\tau H}\) is bounded.  Since
\(\cD_{\mathrm{loc}}^0\) is dense in \(\Omega^\perp\), \textup{(G2)}
extends by continuity to
\[
 \norm{e^{-\tau H}|_{\Omega^\perp}}\le q.
\]
The spectral theorem gives
\[
 \norm{e^{-\tau H}|_{\Omega^\perp}}
 =\sup_{\lambda\in\spec(H|_{\Omega^\perp})}e^{-\tau\lambda}
 =e^{-\tau\inf\spec(H|_{\Omega^\perp})}.
\]
The displayed norm bound is therefore equivalent to
\(\inf\spec(H|_{\Omega^\perp})\ge-\tau^{-1}\log q\), proving
\textup{(G2)}\(\Leftrightarrow\)\textup{(G3)}.  If
\(\psi\in\Ker H\cap\Omega^\perp\), then
\(\norm\psi=\norm{e^{-\tau H}\psi}\le q\norm\psi\), so \(\psi=0\).
Thus the vacuum is the only zero-energy vector.
\end{proof}

\begin{corollary}[Euclidean connected-correlation form]
\label{cor:Euclidean-form}
Under OS reconstruction, \eqref{eq:finite-Gram-gap} is the finite matrix
inequality
\begin{equation}
 \boxed{
 \sum_{a,b}\overline{c_a}c_b,
 \omega\!\left((\Theta F_a^0)\,\vartheta_{2\tau}F_b^0\right)
 \le q^2
 \sum_{a,b}\overline{c_a}c_b,
 \omega\!\left((\Theta F_a^0)F_b^0\right),}
 \label{eq:Euclidean-Gram-gap}
\end{equation}
where \(F_a^0\) denotes OS centering and \(\vartheta_t\) denotes Euclidean
time translation.  Hence a Hamiltonian need not be constructed before the
decisive inequality is tested; it can be proved directly on reflected
Euclidean cylinders.
\end{corollary}

\begin{proof}
The OS reconstruction identity identifies the two sides with
\(c^*G_{\mathbf F}^{\tau}c\) and
\(q^2c^*G_{\mathbf F}^{0}c\), respectively.  Apply
\Cref{thm:local-Gram-criterion}.
\end{proof}

\begin{counterexample}[Testing named probes separately is insufficient]
\label{cth:diagonal-tests}
Let
\[
 A=\begin{pmatrix}3/5&3/10\\3/10&3/5\end{pmatrix}
\]
on \(\mathbb C^2\).  Then \(0\prec A\prec I\), its eigenvalues are
\(9/10\) and \(3/10\), and therefore
\(A=e^{-2\tau H}\) for a positive self-adjoint \(H\).  Each coordinate
probe obeys
\[
 \ip{e_i}{Ae_i}=3/5<7/10.
\]
Nevertheless the normalized combination
\(u=(e_1+e_2)/\sqrt2\) obeys
\[
 \ip{u}{Au}=9/10>7/10.
\]
Thus diagonal decay estimates for a list of Wilson probes do not imply the
matrix inequality with \(q^2=7/10\).  Mixed reflected correlations, or a
valid domination theorem controlling them, are indispensable.
\end{counterexample}

\begin{remark}[Relation to the imported full-network compiler]
The Grand Tensor's dense-network theorem computes increasing finite-bank
generalized eigenvalues and takes their supremum.  The present theorem is not
a second proof of that result.  It selects the continuum-local branch,
expresses its endpoint as a Loewner inequality that can be passed directly
through Euclidean cutoff limits, and makes the treatment of growing OS null
spaces explicit.
\end{remark}

% =============================================================================
\section{Passage through changing cutoff null spaces}
\label{sec:cutoff-passage}
% =============================================================================

Let \(j\) index a joint thermodynamic and ultraviolet sequence.  At cutoff
\(j\), let \((\cH_j,\Omega_j,H_j)\) be a finite- or infinite-volume OS
reconstruction and let \(\tau_j\to\tau>0\) be a fixed \emph{physical} delay.
For a fixed finite family \(\mathbf F\) of continuum-labelled local probes,
let
\(V_{j,\mathbf F}:\mathbb C^r\to\Omega_j^\perp\) be their centered cutoff
synthesis and put
\begin{equation}
 G_{j,\mathbf F}^0=V_{j,\mathbf F}^*V_{j,\mathbf F},
 \qquad
 G_{j,\mathbf F}^{\tau_j}
 =V_{j,\mathbf F}^*e^{-2\tau_jH_j}V_{j,\mathbf F}.
 \label{eq:cutoff-Grams}
\end{equation}

\begin{definition}[Local OS kernel convergence]
\label{def:local-kernel-convergence}
The cutoff family has local OS kernel convergence on
\(\cA_{+,\mathrm{loc}}\) if, for every finite local family \(\mathbf F\),
\begin{equation}
 G_{j,\mathbf F}^0\longrightarrow G_{\infty,\mathbf F}^0,
 \qquad
 G_{j,\mathbf F}^{\tau_j}
 \longrightarrow G_{\infty,\mathbf F}^{\tau}
 \label{eq:Gram-convergence}
\end{equation}
entrywise, equivalently in operator norm in the fixed finite coefficient
space.  The limiting matrices are allowed to be singular.
\end{definition}

\begin{theorem}[Quotient-stable local gap passage]
\label{thm:quotient-stable-passage}
Assume local OS kernel convergence and that the limiting local cylinders
reconstruct a Hilbert space \((\cH_\infty,\Omega_\infty,H_\infty)\) whose
centered local span is dense in \(\Omega_\infty^\perp\).  Suppose there is a
single \(q<1\) such that, for every finite local family \(\mathbf F\), one has
\begin{equation}
 G_{j,\mathbf F}^{\tau_j}
 \preceq q^2G_{j,\mathbf F}^0+E_{j,\mathbf F},
 \qquad
 \norm{E_{j,\mathbf F}}_{\mathrm{op}}\longrightarrow0,
 \label{eq:cutoff-gap-with-error}
\end{equation}
where \(E_{j,\mathbf F}=E_{j,\mathbf F}^*\) need not have a sign.  Then
\begin{equation}
 G_{\infty,\mathbf F}^{\tau}
 \preceq q^2G_{\infty,\mathbf F}^0
 \label{eq:limit-Gram-gap}
\end{equation}
for every finite local family, and
\begin{equation}
 \boxed{
 H_\infty|_{\Omega_\infty^\perp}
 \succeq-\tau^{-1}\log q\,I.}
 \label{eq:limit-mass}
\end{equation}
This conclusion remains valid when the cutoff entrance Grams are nonsingular
but their limiting Grams acquire new null vectors.
\end{theorem}

\begin{proof}
Fix \(\mathbf F\) and \(c\in\mathbb C^r\).  From
\eqref{eq:cutoff-gap-with-error},
\[
 c^*G_{j,\mathbf F}^{\tau_j}c
 \le q^2c^*G_{j,\mathbf F}^0c
     +\norm{E_{j,\mathbf F}}\norm c^2.
\]
Take \(j\to\infty\) and use \eqref{eq:Gram-convergence}.  This gives the
quadratic-form inequality \eqref{eq:limit-Gram-gap}.  The cone of positive
semidefinite matrices is closed, so the same conclusion holds in Loewner
order.  The common constant \(q\) is independent of \(\mathbf F\);
\Cref{thm:local-Gram-criterion} therefore proves \eqref{eq:limit-mass}.

No inversion of \(G_{j,\mathbf F}^0\) was used.  If
\(c\in\Ker G_{\infty,\mathbf F}^0\), then positivity and
\eqref{eq:limit-Gram-gap} give
\(0\le c^*G_{\infty,\mathbf F}^{\tau}c\le0\).  Hence the delayed vector is
also null and the inequality descends to the limiting OS quotient.
\end{proof}

\begin{corollary}[Cutoff-dependent contraction factors]
\label{cor:cutoff-q}
The conclusion remains true if \(q\) in
\eqref{eq:cutoff-gap-with-error} is replaced by \(q_j\) with
\(\limsup_jq_j\le q<1\).
\end{corollary}

\begin{proof}
For any \(q'\in(q,1)\), one has \(q_j\le q'\) eventually.  Apply
\Cref{thm:quotient-stable-passage} with \(q'\), then let \(q'\downarrow q\).
\end{proof}

\begin{proposition}[The exact role of a lower frame]
\label{prop:lower-frame-role}
The gap passage of \Cref{thm:quotient-stable-passage} requires no lower frame.
However:
\begin{enumerate}[label=\textup{(L\arabic*)},leftmargin=3.4em]
\item if \(G_{\infty,\mathbf F}^0=0\) for every local family, the limiting
      OS Hilbert space is the vacuum line and the gap statement is vacuous;
\item a coefficient metric \(R_{\mathbf F}\succeq0\) represents nonzero
      physical probes in the limit if
\begin{equation}
 \boxed{
 G_{j,\mathbf F}^0\succeq c_{\mathbf F}R_{\mathbf F}-o(1),
 \qquad c_{\mathbf F}>0.}
 \label{eq:local-lower-frame}
\end{equation}
Then \(G_{\infty,\mathbf F}^0\succeq
c_{\mathbf F}R_{\mathbf F}\), so every nonzero vector in
\(\Ran R_{\mathbf F}\) survives the OS quotient.
\end{enumerate}
Thus contraction controls the spectrum of what survives; a lower frame proves
that the intended local physics survives at all.
\end{proposition}

\begin{proof}
If every centered local vector has zero norm, density leaves only the vacuum
line.  For the second statement, take the limit of the quadratic-form version
of \eqref{eq:local-lower-frame}.  A vector with positive
\(R_{\mathbf F}\)-norm then has positive limiting OS norm.
\end{proof}

\begin{firewall}[Absolute convergence can hide normalized soft vectors]
The proposition does not contradict the toron-collapse warning in the Grand
Tensor manuscript.  A sequence of cutoff-dependent normalized vectors may
have unnormalized entrance norm tending to zero and therefore disappear from
every fixed local limiting Gram.  Such a sequence refutes a uniform
finite-torus gap on its trajectory but is not automatically a vector in the
local continuum GNS space.  One must say which target Hilbert space is being
tested.
\end{firewall}

% =============================================================================
\section{Thermodynamic locality and the wrapping-sector fork}
\label{sec:local-toron-fork}
% =============================================================================

A local observable family is held fixed in physical support while the spatial
volume grows.  A wrapping Polyakov or toron commutator whose support diameter
grows with the torus is not such a family.  This distinction does not license
discarding a soft global state when the target is a uniform finite-torus gap;
it identifies a different target, namely the GNS closure of the infinite-volume
local algebra.

\begin{proposition}[Local gap and global wrapping softening are logically compatible]
\label{prop:local-global-compatible}
Fix \(q\in(0,1)\) and let
\(\cH_j=\cL_j\oplus\mathbb Ct_j\), where the increasing union of the
\(\cL_j\) represents local vectors and \(t_j\) is a volume-wrapping vector.
Define positive contractions
\begin{equation}
 T_j=qI_{\cL_j}\oplus(1-j^{-1})I_{\mathbb Ct_j}.
 \label{eq:toy-local-global}
\end{equation}
Then every local vector contracts by \(q\), and the local inductive-limit
Hamiltonian has gap \(-\tau^{-1}\log q\).  Nevertheless
\begin{equation}
 \norm{T_j}=1-j^{-1}\longrightarrow1,
 \label{eq:global-soft}
\end{equation}
so the global finite-volume gaps collapse.
\end{proposition}

\begin{proof}
The restriction of \(T_j\) to \(\cL_j\) is exactly \(qI\), so its compatible
inductive limit is \(qI\).  The wrapping eigenvalue is eventually larger than
\(q\) and tends to one, which proves \eqref{eq:global-soft}.
\end{proof}

\begin{corollary}[Correct interpretation of the imported toron sequence]
\label{cor:toron-interpretation}
The centre-neutral toron soft sequence in the Grand Tensor manuscript rules
out a uniform full finite-torus gap on the selected over-weak bare trajectory.
It does not, without an OS-norm approximation by bounded-support observables,
rule out a gap in an independently constructed infinite-volume local-vacuum
GNS representation.  Conversely, a local-vacuum gap does not repair the
failed finite-torus trajectory.
\end{corollary}

\begin{proof}
The two assertions are exactly the two independent summands in
\Cref{prop:local-global-compatible}.  To identify the summands in a physical
theory one must prove, rather than assume, local approximation or asymptotic
orthogonality of the wrapping sector.
\end{proof}

\begin{firewall}[No silent sector deletion]
The local branch is legitimate only if the claimed continuum theory is
constructed from a net of bounded local gauge-invariant algebras and those
algebras are cyclic in its vacuum representation.  If the target instead
requires uniform control of every finite-volume centre-neutral state, the
wrapping sector remains in scope and the toron soft vector is terminal for the
over-weak trajectory.
\end{firewall}

% =============================================================================
\section{The physical-time scaling gate}
\label{sec:physical-time}
% =============================================================================

Let \(a_j\downarrow0\) be the Euclidean lattice time spacing and let
\(T_j\) denote the positive self-adjoint contraction implementing transfer by one lattice step on a centered cutoff sector.  A
dimensionless one-step deficit is not itself a physical mass.  The following
statement makes the obstruction precise.

\begin{theorem}[A fixed one-step contraction collapses a nonzero continuum sector]
\label{thm:fixed-step-collapse}
Let \(E\) be a finite-dimensional coefficient space and
\(V_j:E\to\cH_j\) a noncollapsing bank with
\begin{equation}
 V_j^*V_j\longrightarrow G^0\ne0.
 \label{eq:noncollapsed-bank}
\end{equation}
Assume that for every \(t\ge0\), with
\(n_j(t)=\lfloor t/a_j\rfloor\), the matrix kernels
\begin{equation}
 V_j^*T_j^{n_j(t)}V_j\longrightarrow G(t)
 \label{eq:discrete-to-continuum-kernel}
\end{equation}
and that \(G(t)=V^*e^{-tH}V\) for a strongly continuous semigroup on a
limit Hilbert space, with \(V^*V=G^0\).  If a reducing subspace containing
\(\Ran V_j\) obeys
\begin{equation}
 \norm{T_j}\le q<1
 \label{eq:fixed-step-q}
\end{equation}
uniformly in \(j\), then \(G^0=0\), a contradiction.  Hence any nonzero
strongly continuous continuum sector forces the relevant one-step norms to
approach one.
\end{theorem}

\begin{proof}
For fixed \(t>0\), reduction and \eqref{eq:fixed-step-q} give
\[
 \norm{V_j^*T_j^{n_j(t)}V_j}
 \le\norm{V_j}^2q^{n_j(t)}\longrightarrow0,
\]
because \(n_j(t)\to\infty\) and \(\norm{V_j}^2=\norm{V_j^*V_j}\) remains
bounded.  Therefore \(G(t)=0\) for every \(t>0\).  Strong continuity gives
\(G(t)\to G(0)=G^0\) as \(t\downarrow0\), so \(G^0=0\), contradicting
\eqref{eq:noncollapsed-bank}.
\end{proof}

\begin{theorem}[Rate-calibrated transfer of a cutoff gap]
\label{thm:rate-calibrated-gap}
In the setting above, suppose the entire centered cutoff space is reducing and
\begin{equation}
 \rho_j:=\norm{T_j|_{\Omega_j^\perp}}<1,
 \qquad
 \liminf_{j\to\infty}\frac{-\log\rho_j}{a_j}\ge m>0.
 \label{eq:rate-premise}
\end{equation}
Assume local kernel convergence to a strongly continuous OS semigroup.  Then
for every \(t\ge0\),
\begin{equation}
 \norm{e^{-tH_\infty}|_{\Omega_\infty^\perp}}\le e^{-mt},
 \qquad
 \boxed{H_\infty|_{\Omega_\infty^\perp}\succeq mI.}
 \label{eq:rate-conclusion}
\end{equation}
If \(\rho_j=1-\gamma_j\) with \(\gamma_j\to0\), then
\begin{equation}
 \frac{-\log\rho_j}{a_j}
 =\frac{\gamma_j}{a_j}\bigl(1+o(1)\bigr).
 \label{eq:deficit-rate}
\end{equation}
Thus the physical quantity is the deficit per unit physical time, not the raw
one-step deficit.
\end{theorem}

\begin{proof}
For any \(m'<m\), \eqref{eq:rate-premise} gives
\(\rho_j\le e^{-m'a_j}\) eventually.  Hence, for
\(n_j(t)=\lfloor t/a_j\rfloor\),
\[
 \norm{T_j^{n_j(t)}|_{\Omega_j^\perp}}
 \le e^{-m'a_jn_j(t)}\longrightarrow e^{-m't}.
\]
Passing the corresponding finite local Gram inequalities through
\Cref{thm:quotient-stable-passage}, first at fixed \(t>0\) and then over a
countable dense family, gives
\(\norm{e^{-tH_\infty}|_{\Omega_\infty^\perp}}\le e^{-m't}\).
Let \(m'\uparrow m\).  The spectral theorem yields the Hamiltonian bound.
Finally, \(-\log(1-\gamma)=\gamma+O(\gamma^2)\), proving
\eqref{eq:deficit-rate}.
\end{proof}

\begin{remark}[Why one fixed physical delay is cleaner]
The rate theorem assumes a norm bound on the full centered cutoff space, so it
is vulnerable to wrapping and other volume-growing soft sectors.  For the
infinite-volume local branch, the target-minimal route is instead
\Cref{thm:quotient-stable-passage}: prove one common local matrix contraction
at a fixed physical delay \(\tau\).  This avoids claiming that a local finite
bank is invariant under a one-step cutoff transfer.
\end{remark}

% =============================================================================
\section{Why the V35 static card is upstream data, not the last gap row}
\label{sec:static-versus-dynamic}
% =============================================================================

The V35 complete trace-shell card is schematically a same-shell Hessian
numerator
\begin{equation}
 \mathbb N^{\mathrm{full}}
 =D-\Gamma_{\mathrm{score}},
 \label{eq:V35-static-schematic}
\end{equation}
with an exact head--tail decomposition when a truncation is chosen.  Its sign
is useful static information.  The mass-gap card in this note is instead
\begin{equation}
 G_{\mathbf F}^{\tau}
 \preceq q^2G_{\mathbf F}^0,
 \label{eq:dynamic-schematic}
\end{equation}
where both Grams are reflected time-separated correlators on the same OS
quotient.  No identity in the attached notes equates
\(\mathbb N^{\mathrm{full}}\) with
\(G^0-G^\tau\), nor should such an identity be assumed.

\begin{proposition}[Minimum incidence needed to use a static card dynamically]
\label{prop:static-dynamic-incidence}
Suppose a static finite card \(\mathbb N_j\) and an OS transfer pair
\((G_j^0,G_j^\tau)\) are given on coefficient spaces connected by a declared
source map \(R_j\).  A bound
\begin{equation}
 R_j^*\mathbb N_jR_j\succeq\delta G_j^0
 \label{eq:static-coercivity}
\end{equation}
does not imply \eqref{eq:dynamic-schematic} without a further same-history
incidence inequality of the form
\begin{equation}
 G_j^0-G_j^\tau
 \succeq c_\tau R_j^*\mathbb N_jR_j-E_j,
 \qquad c_\tau>0.
 \label{eq:missing-incidence}
\end{equation}
If \eqref{eq:static-coercivity} and \eqref{eq:missing-incidence} hold with
\(E_j\preceq\varepsilon_jG_j^0\) and
\(c_\tau\delta-\varepsilon_j=:\eta_j\in(0,1)\), then
\begin{equation}
 G_j^\tau\preceq(1-\eta_j)G_j^0.
 \label{eq:incidence-to-gap}
\end{equation}
\end{proposition}

\begin{proof}
The negative statement is logical: \(\mathbb N_j\) contains no delayed
transfer datum until a relation such as \eqref{eq:missing-incidence} is
provided.  Under the two displayed hypotheses,
\[
 G_j^0-G_j^\tau
 \succeq c_\tau\delta G_j^0-E_j
 \succeq(c_\tau\delta-\varepsilon_j)G_j^0.
\]
Rearrangement gives \eqref{eq:incidence-to-gap}.
\end{proof}

\begin{firewall}[The missing inequality cannot be defined into existence]
The map \(R_j\), the delay \(\tau\), the transfer Gram, and the residual
\(E_j\) must be constructed from the same reflected Yang--Mills measure.
Whitening a static score, assigning an information lifetime to a spatial shell,
or fitting a decay rate does not prove \eqref{eq:missing-incidence}.
\end{firewall}

% =============================================================================
\section{A sufficient closure packet for the local-vacuum branch}
\label{sec:closure-packet}
% =============================================================================

The preceding theorems compress the full local mass-gap implication into four
noninterchangeable rows.

\begin{theorem}[Local-vacuum Yang--Mills closure theorem]
\label{thm:local-vacuum-closure}
Let \(G\) be a compact simple gauge group.  Suppose a joint sequence of
gauge-invariant lattice theories at lattice spacings \(a_j\downarrow0\) and
volumes tending to infinity satisfies all of the following.
\begin{enumerate}[label=\textup{(C\arabic*)},leftmargin=3.4em]
\item \emph{Renormalized trajectory.}  A gauge-invariant response coupling is
      fixed at a nonzero physical scale and adjacent step scaling selects the
      cutoff sequence; generated local Wilson operators are retained.
\item \emph{Continuum local state.}  The local gauge-invariant Euclidean
      correlation functions converge and satisfy the OS axioms needed to
      reconstruct a strongly continuous, translation-covariant local theory
      on \(\mathbb R^4\).
\item \emph{Nontriviality and cyclicity.}  A dense family of bounded-support,
      centre-neutral Wilson--spin-network observables has nonzero limiting OS
      norm, is cyclic, and contains a non-Gaussian or otherwise nontrivial
      limiting correlation packet.  The required physical coefficient banks
      obey local lower frames such as \eqref{eq:local-lower-frame}.
\item \emph{Uniform local delayed contraction.}  There exist one physical
      delay \(\tau>0\) and one \(q<1\) such that every finite centered local
      family obeys \eqref{eq:cutoff-gap-with-error}, with its error tending to
      zero along the selected trajectory.
\end{enumerate}
Then the reconstructed continuum theory has a unique vacuum and a mass gap
\begin{equation}
 \boxed{m_{\mathrm{YM}}\ge-\tau^{-1}\log q>0.}
 \label{eq:YM-gap-conclusion}
\end{equation}
If the continuum state also meets the remaining regularity and nontriviality
requirements in the standard Yang--Mills existence problem, these four rows
close the existence-and-gap implication for that constructed theory.
\end{theorem}

\begin{proof}
Rows \textup{(C1)--(C3)} give a nonzero local OS reconstruction whose centered
local span is dense in the vacuum complement.  Row \textup{(C4)} and local
kernel convergence satisfy \Cref{thm:quotient-stable-passage}.  Therefore
\eqref{eq:limit-mass} holds, and
\Cref{thm:local-Gram-criterion} gives vacuum uniqueness.  The last sentence
does not derive the remaining existence axioms; it records that they are
separate premises which, together with the proved gap implication, give the
full target statement.
\end{proof}

\begin{assessment}
The theorem is deliberately stronger and cleaner than a sequence of
probe-by-probe positive fitted masses.  Its uniformity is over all finite
linear combinations from the local algebra, while its locality avoids the
false requirement that a growing wrapping toron writer define one fixed local
continuum observable.  The price is exact: row \textup{(C4)} is essentially a
uniform exponential-clustering inequality in reflected matrix form and is
currently unproved for four-dimensional continuum Yang--Mills theory.
\end{assessment}

% =============================================================================
\section{Status ledger and next upstream attack}
\label{sec:status}
% =============================================================================

\begingroup
\small
\begin{longtable}{>{\raggedright\arraybackslash}p{0.28\textwidth}
                  >{\raggedright\arraybackslash}p{0.16\textwidth}
                  >{\raggedright\arraybackslash}p{0.48\textwidth}}
\toprule
\textbf{Row}&\textbf{Status}&\textbf{What is established or missing}\\
\midrule
\endfirsthead
\toprule
\textbf{Row}&\textbf{Status}&\textbf{What is established or missing}\\
\midrule
\endhead
Finite full/head--tail static identity
& \MGimport
& Exact in the attached V35 reduced-fibre note; not repeated here.\\[0.35em]
Boundary shell/tail alternatives
& \MGimport
& Exact at their stated spatial-information semantics in the attached V33 boundary note; not an OS clock.\\[0.35em]
Dense-network and lower-frame warnings
& \MGimport
& Present in the Grand Tensor V067 manuscript, including the over-weak toron obstruction and conditional full-core promotion.\\[0.35em]
Finite local OS Gram criterion
& \MGpass
& \Cref{thm:local-Gram-criterion} proves equivalence of one common Loewner contraction and a Hamiltonian mass gap.\\[0.35em]
Passage through limiting null spaces
& \MGpass
& \Cref{thm:quotient-stable-passage} needs no Gram inverse; new null vectors are automatically delayed-null.\\[0.35em]
Nontrivial local field survival
& \MGopen
& Requires limiting OS norms and physical lower frames; a gap on the zero quotient is vacuous.\\[0.35em]
Wrapping/local sector separation
& \MGpass
& The abstract compatibility theorem proves that local gap and global toron softening are not logical contradictions.  Physical identification remains to be proved.\\[0.35em]
Physical-time calibration
& \MGpass
& \Cref{thm:fixed-step-collapse,thm:rate-calibrated-gap} show why a raw one-step gap is the wrong continuum variable.\\[0.35em]
Static-to-delayed incidence
& \MGopen
& The exact missing relation is \eqref{eq:missing-incidence}, or a direct proof of \eqref{eq:cutoff-gap-with-error}.\\[0.35em]
Renormalized local continuum state
& \MGopen
& No attached manuscript constructs the required nontrivial four-dimensional OS-positive continuum measure on a selected response-matched trajectory.\\[0.35em]
Uniform local delayed contraction
& \MGopen
& No common \((\tau,q)\) has been proved for every finite centered local Wilson-network combination.\\[0.35em]
Full Yang--Mills mass gap
& \MGopen
& It follows rigorously from \Cref{thm:local-vacuum-closure} only after the open physical rows are populated.\\
\bottomrule
\end{longtable}
\endgroup

\begin{openproblem}[V2 mandate: derive the reflected delay inequality]
\label{op:V2}
Work on the local-vacuum branch and retain the response-matched cutoff
trajectory as an explicit premise.  Choose a bounded positive-time collar and
a finite centre-neutral local Wilson-network family.  The next iteration must
do one of the following, without returning to a purely static sign:
\begin{enumerate}[label=\textup{(V2.\arabic*)},leftmargin=4.6em]
\item construct the same-measure entrance and delayed OS Grams and prove a
      uniform finite-family estimate of the form
      \eqref{eq:cutoff-gap-with-error}; or
\item prove a static-to-dynamic incidence inequality of the form
      \eqref{eq:missing-incidence}, with a residual small in the entrance
      metric at fixed physical delay; or
\item if both routes fail, return the first explicit mechanism: loss of local
      OS convergence, lower-frame collapse, phase nonuniqueness, a local
      near-unit transfer vector, or an uncontrolled incidence residual.
\end{enumerate}
The estimate must hold for arbitrary coefficient combinations in the chosen
family.  Diagonal Wilson-loop decay, a spatial shell contraction, the V35
static eigenvalue, and a cutoff-independent one-step contraction are not
substitutes.
\end{openproblem}

\begin{firewall}[End-of-V1 claim boundary]
The full Yang--Mills mass gap has not been reached.  V1 contributes a rigorous
quotient-stable endpoint theorem, a physical-time scaling theorem, and a
precise next inequality.  The honest bottleneck is analytic and physical, not
an omitted finite-dimensional linear-algebra lemma.
\end{firewall}

% =============================================================================
\section*{Source provenance}
\addcontentsline{toc}{section}{Source provenance}
% =============================================================================

The local manuscript sources inspected for this version were:
\begin{itemize}[leftmargin=2em]
\item \texttt{Renewal\_Grand\_Tensor\_Monograph\_V067.tex};
\item \texttt{Renewal\_Yang\_Mills\_Boundary\_Sector\_Cofinal\_Contraction\_Programme\_V33.tex};
\item \texttt{Renewal\_Yang\_Mills\_Reduced\_Fibre\_Wilson\_Shell\_Symbol\_and\_Local\_Vacuum\_Programme\_V35.tex};
\item \texttt{Renewal\_SM\_Einstein\_Continuum\_Research\_Notes\_V29.tex};
\item \texttt{Renewal\_Einstein\_Standard\_Model\_Physical\_Source\_Metric\_Duhamel\_Ancestry\_and\_Quantum\_Continuum\_Programme\_V35.tex};
\item \texttt{Renewal\_NS\_Terminal\_Source\_Concentration\_and\_Singular\_Thread\_Closure\_Programme\_V34.tex}.
\end{itemize}
For the public status and the standard target, see the Clay Mathematics
Institute, ``Yang--Mills and the Mass Gap,''
\url{https://www.claymath.org/millennium/yang-mills-the-maths-gap/}, and the
official problem description by A.~Jaffe and E.~Witten linked there.

% =============================================================================
% BEGIN APPENDED VERSION V2 MATERIAL
% =============================================================================

\clearpage
\part*{Version V2 continuation: orbit-tube coercivity and finite Wilson-slab minorization}
\addcontentsline{toc}{part}{Version V2 continuation: orbit-tube coercivity and finite Wilson-slab minorization}

% =============================================================================
\section{V2 direction and nonduplication audit}
\label{sec:V2-direction}
% =============================================================================

V1 isolated the dynamic target
\(
G_{\mathbf F}^{\tau}\preceq q^2G_{\mathbf F}^0
\)
at one fixed physical delay.  The attached Grand Tensor already proves that a
finite Wilson half-slab gap is the second-to-first singular-value ratio of the
positive boundary amplitude, equivalently the sharp maximal correlation of
the two boundary fields.  It also gives common-factor shorting, independent
block tensorization, signed return moments, and two-scale Feshbach criteria.
Those results are imported and not repeated.

The missing implication is between a static or infinitesimal form estimate and
the finite-delay Gram.  V2 proves the exact bridge when coercivity is available
along the entire semigroup orbit of a local bank.  It also proves that a lower
energy expectation only at the entrance is insufficient, even if its margin
is uniform.  Finally, V2 obtains a quantitative finite-volume transfer gap
directly from the oscillation of a positive Wilson slab kernel.  This closes a
literal finite card, but the bound decays exponentially in the slab carrier;
it therefore identifies rather than hides the continuum bottleneck.

% =============================================================================
\section{The exact orbit-energy identity}
\label{sec:V2-orbit-identity}
% =============================================================================

Let \(H=H^*\succeq0\) on a Hilbert space \(\cH\), and let
\(V:E\to\cH\) be a bounded synthesis from a finite-dimensional coefficient
space.  Assume \(\Ran V\subset\operatorname{Dom}(H^{1/2})\).  For
\(s\ge0\), put
\begin{equation}
 G(s):=V^*e^{-2sH}V,
 \qquad
 K(s):=V^*e^{-sH}He^{-sH}V.
 \label{eq:V2-GK}
\end{equation}
The first matrix is the delayed Gram and the second is the energy Gram of the
evolved bank.

\begin{theorem}[Orbit-energy identity]
\label{thm:V2-orbit-energy}
The map \(G\) is norm continuous on \([0,\infty)\), absolutely continuous on
every compact interval, and
\begin{equation}
 \boxed{G'(s)=-2K(s)}
 \label{eq:V2-G-derivative}
\end{equation}
for almost every \(s\ge0\), with the right derivative at zero understood in
quadratic-form sense.  Consequently
\begin{equation}
 \boxed{
 G(\tau)=G(0)-2\int_0^\tau K(s)\,\dd s.}
 \label{eq:V2-orbit-integral}
\end{equation}
\end{theorem}

\begin{proof}
Fix \(c,d\in E\) and let \(u=Vc\), \(v=Vd\).  By the spectral theorem,
\[
 c^*G(s)d=\int_{[0,\infty)}e^{-2s\lambda},
            \dd\mu_{u,v}(\lambda).
\]
The form-domain assumption makes
\(\int\lambda\,\dd|\mu_{u,v}|<\infty\) by Cauchy--Schwarz.  Differentiation
under the integral therefore gives
\[
 \frac{\dd}{\dd s}c^*G(s)d
 =-2\int\lambda e^{-2s\lambda}\,\dd\mu_{u,v}(\lambda)
 =-2c^*K(s)d
\]
for almost every \(s\), and integration gives
\eqref{eq:V2-orbit-integral}.  Since \(E\) is finite dimensional, entrywise
absolute continuity is equivalent to matrix-norm absolute continuity.
\end{proof}

\begin{remark}[Dynamic rather than spatial incidence]
The integrand in \eqref{eq:V2-orbit-integral} is built from the same
Hamiltonian and the same OS vectors as the delayed Gram.  This identity is
therefore an admissible static-to-dynamic incidence.  A spatial heat bath,
Witten generator, shell information loss, or reduced-fibre Hessian cannot be
substituted for \(H\) without a separately proved comparison to this
integrand.
\end{remark}

% =============================================================================
\section{Orbit-tube coercivity implies delayed contraction}
\label{sec:V2-orbit-coercivity}
% =============================================================================

\begin{theorem}[Orbit-tube coercivity with an integrated residual]
\label{thm:V2-orbit-coercivity}
Let \(R\succeq0\) be a fixed matrix on \(E\), let \(m>0\), and let
\(r\in L^1(0,\tau)\) be nonnegative.  Suppose that for almost every
\(s\in[0,\tau]\),
\begin{equation}
 \boxed{K(s)\succeq mG(s)-r(s)R.}
 \label{eq:V2-orbit-floor}
\end{equation}
Then
\begin{equation}
 \boxed{
 G(\tau)
 \preceq e^{-2m\tau}G(0)
 +2\left(\int_0^\tau e^{-2m(\tau-s)}r(s)\,\dd s\right)R.}
 \label{eq:V2-orbit-Gronwall}
\end{equation}
If, in addition, \(R\preceq C G(0)\) and
\begin{equation}
 q^2:=e^{-2m\tau}
 +2C\int_0^\tau e^{-2m(\tau-s)}r(s)\,\dd s<1,
 \label{eq:V2-q-residual}
\end{equation}
then
\begin{equation}
 \boxed{G(\tau)\preceq q^2G(0).}
 \label{eq:V2-orbit-to-delay}
\end{equation}
With \(r=0\), the exact factor is \(q=e^{-m\tau}\).
\end{theorem}

\begin{proof}
Fix \(c\in E\), and set
\(f(s)=c^*G(s)c\) and \(\rho=c^*Rc\).  By
\Cref{thm:V2-orbit-energy}, \eqref{eq:V2-orbit-floor} gives
\[
 f'(s)=-2c^*K(s)c\le-2mf(s)+2r(s)\rho
\]
almost everywhere.  Multiplication by \(e^{2ms}\) and integration from zero
to \(\tau\) yields
\[
 f(\tau)\le e^{-2m\tau}f(0)
 +2\rho\int_0^\tau e^{-2m(\tau-s)}r(s)\,\dd s.
\]
Since this holds for every \(c\), it is
\eqref{eq:V2-orbit-Gronwall}.  The bound \(R\preceq CG(0)\) then gives
\eqref{eq:V2-orbit-to-delay} with \eqref{eq:V2-q-residual}.
\end{proof}

\begin{corollary}[Cutoff-to-continuum orbit compiler]
\label{cor:V2-cutoff-orbit}
For each cutoff \(j\), let \((G_j,K_j,R_j,m_j,r_j,\tau_j)\) satisfy
\Cref{thm:V2-orbit-coercivity}.  Assume
\begin{equation}
 \tau_j\to\tau>0,
 \qquad
 \liminf_jm_j\ge m_*>0,
 \qquad
 R_j\preceq C G_j(0),
 \label{eq:V2-cofinal-orbit-data}
\end{equation}
and
\begin{equation}
 \int_0^{\tau_j}r_j(s)\,\dd s\longrightarrow0.
 \label{eq:V2-integrated-residual-collapse}
\end{equation}
Then, for every \(q\in(e^{-m_*\tau},1)\),
\begin{equation}
 G_j(\tau_j)\preceq q^2G_j(0)
 \label{eq:V2-eventual-q}
\end{equation}
eventually.  If the local OS kernels converge as in V1, the limiting local
Hamiltonian has gap at least \(m_*\).
\end{corollary}

\begin{proof}
The residual term in \eqref{eq:V2-q-residual} is bounded by
\(2C\int_0^{\tau_j}r_j\), hence tends to zero.  The first term has limsup at
most \(e^{-2m_*\tau}\).  This proves \eqref{eq:V2-eventual-q}; apply
\Cref{thm:quotient-stable-passage} and let
\(q\downarrow e^{-m_*\tau}\).
\end{proof}

\begin{assessment}[What has actually been reduced]
The V1 finite-delay target is now implied by a continuum of finite-dimensional
Loewner inequalities \eqref{eq:V2-orbit-floor}.  This is not circular if the
lower bound comes from an independently controlled local generator form, a
signed Feshbach comparison, or a block coercivity theorem.  It is circular if
\(m\) is defined from the unknown delayed norm or if the bound is asserted
only after projecting away the observed soft direction.
\end{assessment}

% =============================================================================
\section{An entrance energy margin alone does not control a delay}
\label{sec:V2-first-moment-no-go}
% =============================================================================

\begin{counterexample}[Uniform first moment with near-unit delayed survival]
\label{cth:V2-first-moment-no-gap}
Fix \(m>0\) and \(\tau>0\).  There are two-dimensional positive
Hamiltonians \(H_n\) and unit vectors \(v_n\) such that
\begin{equation}
 \boxed{\ip{v_n}{H_nv_n}=m}
 \label{eq:V2-fixed-mean-energy}
\end{equation}
for every \(n\), while
\begin{equation}
 \boxed{\ip{v_n}{e^{-2\tau H_n}v_n}\longrightarrow1.}
 \label{eq:V2-survival-to-one}
\end{equation}
Thus a strictly positive static energy or Hessian expectation on each named
entrance vector does not imply any common finite-delay contraction.
\end{counterexample}

\begin{proof}
Choose \(\varepsilon_n\downarrow0\) with \(\varepsilon_n<\min\{m,1\}\), put
\(M_n=m/\varepsilon_n\), and define
\[
 H_n=\begin{pmatrix}\varepsilon_n&0\\0&M_n\end{pmatrix},
 \qquad
 p_n=\frac{m-\varepsilon_n}{M_n-\varepsilon_n},
 \qquad
 v_n=\binom{\sqrt{1-p_n}}{\sqrt{p_n}}.
\]
Then \(0<p_n<1\), \(p_n\to0\), and direct substitution gives
\((1-p_n)\varepsilon_n+p_nM_n=m\).  On the other hand,
\[
 \ip{v_n}{e^{-2\tau H_n}v_n}
 =(1-p_n)e^{-2\tau\varepsilon_n}
   +p_ne^{-2\tau M_n}\longrightarrow1.
\]
The high-energy component pays the fixed first moment at the entrance but is
immediately damped; the surviving vector is concentrated at energy
\(\varepsilon_n\).  Equivalently, the Rayleigh quotient of
\(e^{-sH_n}v_n\) falls toward \(\varepsilon_n\) along the orbit, so
\eqref{eq:V2-orbit-floor} cannot hold with a common positive \(m\).
\end{proof}

\begin{firewall}[Static trace-shell positivity is not orbit coercivity]
Even a fully evaluated positive V35 trace-shell mode would be an entrance or
static coefficient statement unless it were incident with the actual OS
energy Gram \(K_j(s)\) for every \(s\) in a fixed physical interval.  The
counterexample shows that replacing this orbit row by a single first moment is
mathematically invalid, not merely conservative.
\end{firewall}

% =============================================================================
\section{A finite mesh certificate for the orbit tube}
\label{sec:V2-mesh}
% =============================================================================

At a fixed cutoff, the coefficient space is finite.  The continuum of orbit
times can therefore be certified by a mesh plus a derivative enclosure.

\begin{theorem}[Finite orbit-mesh certificate]
\label{thm:V2-mesh-certificate}
Assume \(\Ran V\subset\operatorname{Dom}(H)\), define
\begin{equation}
 L(s):=V^*e^{-sH}H^2e^{-sH}V,
 \qquad
 A_m(s):=K(s)-mG(s),
 \label{eq:V2-LA}
\end{equation}
and suppose
\begin{equation}
 \sup_{0\le s\le\tau}
 2\norm{L(s)-mK(s)}_{\mathrm{op}}\le B.
 \label{eq:V2-derivative-bound}
\end{equation}
Let \(0=s_0<s_1<\cdots<s_N=\tau\) have mesh width
\(h=\max_k(s_{k+1}-s_k)\).  If
\begin{equation}
 \boxed{A_m(s_k)\succeq Bh\,I_E\qquad(0\le k\le N),}
 \label{eq:V2-mesh-margin}
\end{equation}
then \(A_m(s)\succeq0\) on \([0,\tau]\), and hence
\begin{equation}
 G(\tau)\preceq e^{-2m\tau}G(0).
 \label{eq:V2-mesh-delay}
\end{equation}
The identity matrix may be replaced throughout by any fixed positive
coefficient metric after restricting to its support.
\end{theorem}

\begin{proof}
Functional calculus gives
\[
 K'(s)=-2L(s),
 \qquad
 A_m'(s)=-2L(s)+2mK(s).
\]
Thus \eqref{eq:V2-derivative-bound} implies
\(\norm{A_m(s)-A_m(t)}\le B|s-t|\).  For every \(s\), choose a mesh point
\(s_k\) with \(|s-s_k|\le h\).  Weyl's inequality and
\eqref{eq:V2-mesh-margin} give
\[
 \lambda_{\min}(A_m(s))
 \ge\lambda_{\min}(A_m(s_k))-B|s-s_k|\ge0.
\]
Apply \Cref{thm:V2-orbit-coercivity} with \(r=0\).
\end{proof}

\begin{corollary}[Outward interval implementation]
\label{cor:V2-outward-mesh}
Suppose interval matrices \(\widehat A_k\) and radii \(\epsilon_k\) satisfy
\(\norm{A_m(s_k)-\widehat A_k}\le\epsilon_k\), while
\(\widehat B\ge B\).  The finitely checkable strict inequalities
\begin{equation}
 \lambda_{\min}(\widehat A_k)-\epsilon_k-\widehat Bh>0
 \label{eq:V2-outward-mesh-check}
\end{equation}
for all \(k\) certify \eqref{eq:V2-mesh-delay}.  Failure of one inequality is
unresolved unless an outward negative Rayleigh bound for the exact
\(A_m(s_k)\) is also produced.
\end{corollary}

\begin{proof}
Weyl's inequality gives
\(A_m(s_k)\succeq
(\lambda_{\min}(\widehat A_k)-\epsilon_k)I\).  Apply
\Cref{thm:V2-mesh-certificate}.  The last sentence distinguishes failure of a
sufficient enclosure from an actual negative witness.
\end{proof}

% =============================================================================
\section{A quantitative finite Wilson-slab gap from kernel oscillation}
\label{sec:V2-Perron-Doob}
% =============================================================================

The attached Grand Tensor proves strict finite Wilson transfer gaps by
positivity improvement but does not use them as continuum gaps.  The next
theorem supplies a simple explicit margin and displays its volume dependence.

\begin{theorem}[Perron--Doob minorization from a positive kernel window]
\label{thm:V2-Perron-Doob}
Let \((X,\mathfrak m)\) be a compact probability space and let
\(\mathcal K\) be the integral operator with continuous symmetric kernel
\(k\) satisfying
\begin{equation}
 0<k_-<k_+<\infty,\qquad k_-\le k(x,y)\le k_+.
 \label{eq:V2-kernel-window}
\end{equation}
Let \(\lambda>0\) be its maximal eigenvalue and let \(h>0\) be the
corresponding eigenfunction normalized by
\(\int h^2\,\dd\mathfrak m=1\).  Define
\begin{equation}
 \dd\pi=h^2\dd\mathfrak m,
 \qquad
 (Pf)(x)=\frac1{\lambda h(x)}
          \int_X k(x,y)h(y)f(y)\,\dd\mathfrak m(y).
 \label{eq:V2-Doob-transform}
\end{equation}
Then \(P\) is a reversible Markov contraction on \(L^2(\pi)\), and, with
\(R=k_+/k_-\),
\begin{equation}
 \boxed{
 \norm{P|_{L_0^2(\pi)}}\le1-R^{-3}<1.}
 \label{eq:V2-Doeblin-gap}
\end{equation}
If \(P\succeq0\) is the normalized OS transfer over physical duration
\(\tau\), its Hamiltonian obeys
\begin{equation}
 \boxed{
 H|_{\one^\perp}
 \succeq-\tau^{-1}\log(1-R^{-3})\,I.}
 \label{eq:V2-kernel-mass}
\end{equation}
\end{theorem}

\begin{proof}
Strict positivity and compactness give a simple positive Perron eigenfunction.
The eigenvalue equation gives, with
\(A=\int h\,\dd\mathfrak m>0\),
\[
 k_-A\le\lambda h(x)\le k_+A.
\]
Hence \(\max h/\min h\le R\).  Since \(\int h^2=1\), one has
\(\min h\le1\le\max h\), and therefore \(\max h\le R\).  Schur's test
gives \(\lambda\le k_+\).

Relative to \(\pi\), the transition density of \(P\) is
\begin{equation}
 r(x,y)=\frac{k(x,y)}{\lambda h(x)h(y)}.
 \label{eq:V2-r-density}
\end{equation}
It is symmetric, has both marginals one, and satisfies
\[
 r(x,y)\ge\frac{k_-}{k_+R^2}=R^{-3}=:\alpha.
\]
Thus
\(r=\alpha+(1-\alpha)s\), where \(s\) is again a symmetric Markov density.
Writing \(\Pi f=\int f\,\dd\pi\), this is
\[
 P=\alpha\Pi+(1-\alpha)S.
\]
Every Markov operator is an \(L^2\) contraction when its invariant law is
\(\pi\).  On \(L_0^2(\pi)\), \(\Pi=0\), so
\(\norm P\le1-\alpha\), proving \eqref{eq:V2-Doeblin-gap}.  If
\(P=e^{-\tau H}\succeq0\), the spectral theorem gives
\eqref{eq:V2-kernel-mass}.
\end{proof}

\begin{corollary}[Explicit but nonuniform finite Wilson bound]
\label{cor:V2-Wilson-finite-gap}
Within one fixed protected sector, consider a finite reflection-positive \(\mathrm{SU}(3)\) Wilson slab whose
raw symmetric boundary kernel is
\begin{equation}
 k(U,V)=e^{-S_{\mathrm{sl}}(U,V)}
 \label{eq:V2-Wilson-kernel}
\end{equation}
with \(N_{\mathrm{sl}}\) plaquette terms and coupling \(\beta>0\).  Then
\begin{equation}
 \operatorname{osc}S_{\mathrm{sl}}\le2\beta N_{\mathrm{sl}},
 \label{eq:V2-Wilson-osc}
\end{equation}
and the normalized transfer satisfies
\begin{equation}
 \boxed{
 \norm{P|_{L_0^2(\pi)}}
 \le1-e^{-6\beta N_{\mathrm{sl}}}.}
 \label{eq:V2-Wilson-crude-q}
\end{equation}
Consequently its finite-volume OS Hamiltonian has the explicit gap
\begin{equation}
 \boxed{
 m_{\mathrm{FV}}
 \ge-\tau^{-1}\log\!\left(1-e^{-6\beta N_{\mathrm{sl}}}\right)>0.}
 \label{eq:V2-Wilson-crude-mass}
\end{equation}
The same bound holds after restriction to a reducing gauge-invariant neutral
sector orthogonal to its vacuum.
\end{corollary}

\begin{proof}
For every \(U\in\mathrm{SU}(3)\),
\(-1\le\frac13\operatorname{ReTr}U\le1\).  Hence each Wilson plaquette
action \(\beta(1-\operatorname{ReTr}U/3)\) has range at most \(2\beta\),
which proves \eqref{eq:V2-Wilson-osc}.  Multiplication of \(k\) by a positive
constant does not change its Doob transform, so
\(R=k_+/k_-\le e^{2\beta N_{\mathrm{sl}}}\).  Insert this in
\Cref{thm:V2-Perron-Doob}.  Restriction to a reducing subspace cannot increase
the operator norm.
\end{proof}

\begin{assessment}[Why this finite proof does not approach the continuum by itself]
\label{ass:V2-global-minorization-collapse}
The corollary rigorously populates a finite reflected-time gap using the actual
positive slab kernel.  Its minorization margin is
\(e^{-6\beta N_{\mathrm{sl}}}\).  On a weak-coupling and thermodynamic
trajectory, both \(\beta\) and the number of slab plaquettes grow, so this
\emph{certificate} collapses exponentially.  This does not prove that the
physical gap collapses; it proves that a worst-configuration global kernel
window cannot provide the uniform \(q<1\) required by V1.  The next proof must
use locality, renormalized blocks, and signed return control rather than a
global action oscillation.
\end{assessment}

% =============================================================================
\section{V2 integrated theorem and revised frontier}
\label{sec:V2-frontier}
% =============================================================================

\begin{theorem}[Integrated V2 transfer-to-gap advance]
\label{thm:V2-integrated}
Preserving V1, V2 proves the following.
\begin{enumerate}[label=\textup{(V2.\arabic*)},leftmargin=5em]
\item The delayed local Gram is the entrance Gram minus the integrated energy
      of the same bank along its physical OS orbit.
\item A uniform orbit-tube energy floor, with a vanishing integrated residual,
      implies the fixed-delay Loewner contraction required by V1.
\item A positive energy expectation only at orbit time zero does not imply any
      common finite-delay contraction.
\item At a fixed cutoff, finitely many outward orbit-time matrices plus one
      derivative bound give a terminating certificate for the full orbit
      interval.
\item Every finite positive symmetric Wilson slab has an explicit
      Perron--Doob transfer gap controlled by its kernel oscillation.
\item That global minorization margin decays exponentially with coupling times
      slab carrier size and therefore does not supply the continuum-uniform
      gap.
\end{enumerate}
\end{theorem}

\begin{proof}
Items \textup{(V2.1)--(V2.2)} are
\Cref{thm:V2-orbit-energy,thm:V2-orbit-coercivity,cor:V2-cutoff-orbit}.
Item \textup{(V2.3)} is \Cref{cth:V2-first-moment-no-gap}.  Item
\textup{(V2.4)} is
\Cref{thm:V2-mesh-certificate,cor:V2-outward-mesh}.  Items
\textup{(V2.5)--(V2.6)} are
\Cref{thm:V2-Perron-Doob,cor:V2-Wilson-finite-gap,ass:V2-global-minorization-collapse}.
\end{proof}

\begingroup
\small
\begin{longtable}{>{\raggedright\arraybackslash}p{0.28\textwidth}
                  >{\raggedright\arraybackslash}p{0.16\textwidth}
                  >{\raggedright\arraybackslash}p{0.48\textwidth}}
\toprule
\textbf{V2 row}&\textbf{Status}&\textbf{Advance or remaining obstruction}\\
\midrule
\endfirsthead
\toprule
\textbf{V2 row}&\textbf{Status}&\textbf{Advance or remaining obstruction}\\
\midrule
\endhead
Same-Hamiltonian orbit identity
& \MGpass
& Exact spectral-theorem identity with no fitted clock or spatial substitute.\\[0.35em]
Orbit coercivity compiler
& \MGpass
& Gives an explicit delayed contraction including an integrated residual budget.\\[0.35em]
Entrance-only static margin
& Refuted
& A two-level sequence has fixed positive mean energy and delayed survival tending to one.\\[0.35em]
Finite orbit mesh
& \MGpass
& Exact finite certificate once the OS energy matrices and a derivative enclosure are available.\\[0.35em]
Finite Wilson slab
& \MGpass
& Strict quantitative Perron--Doob gap for every finite positive slab.\\[0.35em]
Global kernel-window uniformity
& Refuted as a route
& The proved margin decays exponentially in \(\beta N_{\mathrm{sl}}\).\\[0.35em]
Response-matched local orbit floor
& \MGopen
& No attached value or theorem supplies \eqref{eq:V2-orbit-floor} uniformly on a weak-coupling local-vacuum trajectory.\\[0.35em]
Signed complement return
& \MGimport/\MGopen
& Abstract Feshbach and return identities exist in V067; their physical Wilson orbit blocks remain unpopulated.\\[0.35em]
Continuum OS state and nontrivial lower frame
& \MGopen
& Still required before V1's quotient-stable limit theorem can yield a physical theory.\\[0.35em]
Full Yang--Mills mass gap
& \MGopen
& Not reached; the remaining bottleneck is now a uniform local orbit estimate rather than finite-volume positivity.\\
\bottomrule
\end{longtable}
\endgroup

\begin{openproblem}[V3 mandate: localize the orbit floor before taking a global norm]
\label{op:V3}
On one explicitly selected response-matched lattice trajectory, freeze a
bounded physical positive-time collar and a finite local centre-neutral
Wilson-network bank.  Construct the actual OS energy Grams \(K_j(s)\) and
proceed in this order.
\begin{enumerate}[label=\textup{(V3.\arabic*)},leftmargin=5em]
\item Split the physical orbit form into a bounded-range block head and its
      signed complement return using one common OS measure and one fixed
      physical delay.
\item Prove a cutoff-uniform positive floor for the block head on the evolved
      bank; do not use the global action-oscillation minorant.
\item Bound the signed return by an integrable residual \(r_j(s)R_j\) with
      \(\int_0^{\tau_j}r_j\to0\), or retain its exact mixed block if it does
      not vanish.
\item Apply \Cref{thm:V2-orbit-coercivity} or the outward mesh certificate to
      obtain the V1 delayed Gram inequality for that bank.
\item Transport the estimate across adjacent cutoffs and enlarge the local
      bank.  Any failure must return an explicit evolved Rayleigh vector,
      lower-frame collapse, unbounded relaxation router, phase mismatch, or
      nonvanishing return residual.
\end{enumerate}
The V35 static mode may enter only through a proved same-history comparison to
the OS orbit energy.  Re-evaluating it without that incidence would not
advance the mass-gap row.
\end{openproblem}

\begin{firewall}[End-of-V2 claim boundary]
V2 proves a new finite Wilson transfer-gap bound and an exact mechanism for
promoting uniform orbit coercivity to the continuum-local delayed contraction.
It does not prove the required weak-coupling orbit floor, continuum measure,
or full Yang--Mills mass gap.  Those remain the active goal.
\end{firewall}

\section*{Version V2 append-only record}
\addcontentsline{toc}{section}{Version V2 append-only record}
The complete V1 source was retained before this continuation.  Its SHA--256
digest was
\begin{center}\scriptsize\ttfamily
e563a90dd9146b87e81a451088b33f9be397b9362a859685bfd0b3d6cbd9262d.
\end{center}
On the next iteration, retain this full file, append before its unique active
document-closing command, increment the filename to
\texttt{\_V3.tex}, and remove only the superseded V2 file from this note
series.

% =============================================================================
% END APPENDED VERSION V2 MATERIAL
% =============================================================================


% =============================================================================
% BEGIN APPENDED VERSION V3 MATERIAL
% =============================================================================

\clearpage
\part*{Version V3 continuation: orbit-restricted Feshbach coercivity and the Casimir ultraviolet tail}
\addcontentsline{toc}{part}{Version V3 continuation: orbit-restricted Feshbach coercivity and the Casimir ultraviolet tail}

% =============================================================================
\section{V3 direction: localize the physical orbit form}
\label{sec:V3-direction}
% =============================================================================

V2 showed that a same-Hamiltonian lower bound along an evolved local bank
implies the V1 delayed OS contraction.  It also showed that global
worst-configuration minorization loses its margin exponentially with the slab
carrier.  The next step is therefore a signed local head/complement
decomposition of the \emph{physical orbit form}, not another global kernel
window.

The Grand Tensor already contains a global two-scale Feshbach theorem.  V3
does not repeat it.  Instead, it proves a sharper theorem whose Schur floor and
router bound are required only on the projected orbit of the selected local
bank.  This is strictly weaker than a global head gap and is exactly what V2
needs.  V3 then populates the ultraviolet floor and router for a finite
Hamiltonian lattice gauge regulator by cutting the electric Peter--Weyl
Casimir.  The remaining head is finite at fixed regulator, but its uniform
renormalized Schur floor is still open.

% =============================================================================
\section{Exact Feshbach coordinates on an evolved bank}
\label{sec:V3-Feshbach-coordinates}
% =============================================================================

Let \(H=H^*\succeq0\) and let \(P\) be an orthogonal projection with
\(Q=I-P\).  Assume that the block form of \(H\) is
\begin{equation}
 H=\begin{pmatrix}A&B\\B^*&C\end{pmatrix}
 \quad\text{on }P\cH\oplus Q\cH,
 \qquad C\succeq\gamma Q,
 \quad\gamma>0,
 \label{eq:V3-block-H}
\end{equation}
where the off-diagonal form is represented by a bounded block and the usual
form-domain completion is valid.  Define
\begin{equation}
 \boxed{
 \mathcal R:=C^{-1}B^*:P\cH\to Q\cH,
 \qquad
 S:=A-BC^{-1}B^*:P\cH\to P\cH.}
 \label{eq:V3-router-Schur}
\end{equation}
Positivity of \(H\) implies \(S\succeq0\).

For \(w=x\oplus y\) put
\begin{equation}
 z:=y+\mathcal R x.
 \label{eq:V3-z-coordinate}
\end{equation}

\begin{lemma}[Exact graph/innovation identity]
\label{lem:V3-graph-innovation}
For every form-domain vector \(w=x\oplus y\),
\begin{equation}
 \boxed{
 \ip{w}{Hw}=\ip{x}{Sx}+\ip{z}{Cz},
 \qquad
 w=(x\oplus-\mathcal R x)+(0\oplus z).}
 \label{eq:V3-Feshbach-Pythagoras}
\end{equation}
The first summand in the vector decomposition is the minimum-energy follower
of the head coordinate; \(z\) is the exact complement innovation.
\end{lemma}

\begin{proof}
Substitute \(y=z-C^{-1}B^*x\) into the block quadratic form.  The two mixed
terms cancel and leave
\[
 \ip{x}{(A-BC^{-1}B^*)x}+\ip{z}{Cz}.
\]
The vector identity is the definition of \(z\).
\end{proof}

% =============================================================================
\section{A sharp orbit-restricted Feshbach floor}
\label{sec:V3-sharp-floor}
% =============================================================================

The next constant records exactly the competition between the effective head
floor \(\delta\), the complement floor \(\gamma\), and the relaxation-router
norm \(\beta\).  Put
\begin{equation}
 \mathfrak a(\delta,\gamma,\beta)
 :=\delta+\gamma(1+\beta^2)
 \label{eq:V3-a-constant}
\end{equation}
and
\begin{equation}
 \boxed{
 m_{\rm F}(\delta,\gamma,\beta)
 :=\frac{\mathfrak a-
 \sqrt{\mathfrak a^2-4\delta\gamma}}{2}>0.}
 \label{eq:V3-sharp-mF}
\end{equation}

\begin{theorem}[Sharp orbit-restricted Feshbach coercivity]
\label{thm:V3-orbit-Feshbach}
Let \(\mathcal X\subseteq P\cH\) be any set, not necessarily a subspace.
Assume
\begin{equation}
 \ip{x}{Sx}\ge\delta\norm x^2,
 \qquad
 \norm{\mathcal R x}\le\beta\norm x
 \quad(x\in\mathcal X),
 \label{eq:V3-head-router-floor}
\end{equation}
with \(\delta>0\), \(\beta\ge0\).  Then every form-domain
\(w=x\oplus y\) with \(x=Pw\in\mathcal X\) obeys
\begin{equation}
 \boxed{
 \ip{w}{Hw}
 \ge m_{\rm F}(\delta,\gamma,\beta)\norm w^2.}
 \label{eq:V3-sharp-Feshbach-floor}
\end{equation}
The constant is sharp given only the three scalar bounds in
\eqref{eq:V3-head-router-floor} and \(C\succeq\gamma Q\).  It also satisfies
\begin{equation}
 \boxed{
 m_{\rm F}(\delta,\gamma,\beta)
 \ge\frac{\delta\gamma}
 {\delta+\gamma(1+\beta^2)}.}
 \label{eq:V3-crude-mF}
\end{equation}
\end{theorem}

\begin{proof}
Let \(a=\norm x\), \(b=\norm z\).  By
\Cref{lem:V3-graph-innovation},
\begin{equation}
 \ip{w}{Hw}\ge\delta a^2+\gamma b^2.
 \label{eq:V3-energy-ab}
\end{equation}
The router bound and the triangle inequality give
\begin{equation}
 \norm w^2=a^2+\norm{z-\mathcal R x}^2
 \le(1+\beta^2)a^2+2\beta ab+b^2.
 \label{eq:V3-norm-ab}
\end{equation}
The best \(m\) for which the right side of
\eqref{eq:V3-energy-ab} dominates \(m\) times the right side of
\eqref{eq:V3-norm-ab} is the smaller generalized eigenvalue of
\[
 \begin{pmatrix}\delta&0\\0&\gamma\end{pmatrix}
 \quad\text{relative to}\quad
 \begin{pmatrix}1+\beta^2&\beta\\\beta&1\end{pmatrix}.
\]
The second matrix has determinant one, so the characteristic polynomial is
\[
 m^2-\bigl(\delta+\gamma(1+\beta^2)\bigr)m
 +\delta\gamma.
\]
Its smaller root is \eqref{eq:V3-sharp-mF}, proving
\eqref{eq:V3-sharp-Feshbach-floor}.  Equality is attained when the two norm
inequalities are saturated and \((a,b)\) is a generalized eigenvector, so the
constant is sharp at the level of the stated scalar data.  Finally,
\[
 m_{\rm F}
 =\frac{2\delta\gamma}
 {\mathfrak a+\sqrt{\mathfrak a^2-4\delta\gamma}}
 \ge\frac{\delta\gamma}{\mathfrak a},
\]
which proves \eqref{eq:V3-crude-mF}.
\end{proof}

\begin{remark}[Strictly weaker than a global head gap]
For a synthesis \(V:E\to\Omega^\perp\), V3 will take
\[
 \mathcal X_{V,\tau}
 :=\{Pe^{-sH}Vc:0\le s\le\tau,\ c\in E\}.
\]
The Schur and router estimates are needed only on this orbit head.  No claim is
made about unused head directions.  This is the nonduplicative gain over a
global two-scale Feshbach theorem.
\end{remark}

\begin{theorem}[Orbit Feshbach floor with coefficient residual]
\label{thm:V3-orbit-Feshbach-residual}
Let \(V:E\to\Omega^\perp\), set
\(w_s=e^{-sH}Vc\), \(x_s=Pw_s\), and let \(R_E\succeq0\) be a fixed
coefficient metric.  Suppose, for almost every \(s\in[0,\tau]\),
\begin{equation}
 \ip{x_s}{Sx_s}
 \ge\delta\norm{x_s}^2-r(s)c^*R_Ec,
 \qquad
 \norm{\mathcal R x_s}\le\beta\norm{x_s},
 \label{eq:V3-residual-head-floor}
\end{equation}
while \(C\succeq\gamma Q\).  Then the V2 orbit Grams obey
\begin{equation}
 \boxed{
 K(s)\succeq
 m_{\rm F}(\delta,\gamma,\beta)G(s)-r(s)R_E.}
 \label{eq:V3-Feshbach-to-orbit}
\end{equation}
Consequently, if \(R_E\preceq C_EG(0)\) and the V2 residual condition
\eqref{eq:V2-q-residual} holds with
\(m=m_{\rm F}(\delta,\gamma,\beta)\), then
\begin{equation}
 G(\tau)\preceq q^2G(0).
 \label{eq:V3-Feshbach-delay}
\end{equation}
\end{theorem}

\begin{proof}
Repeat the proof of \Cref{thm:V3-orbit-Feshbach} for \(w_s\).  The residual
subtracts unchanged from \eqref{eq:V3-energy-ab}, giving
\[
 \ip{w_s}{Hw_s}
 \ge m_{\rm F}\norm{w_s}^2-r(s)c^*R_Ec.
\]
This is the quadratic-form version of
\eqref{eq:V3-Feshbach-to-orbit}.  Apply
\Cref{thm:V2-orbit-coercivity} for the delayed conclusion.
\end{proof}

% =============================================================================
\section{A soft physical orbit vector returns a soft effective head}
\label{sec:V3-soft-witness}
% =============================================================================

An advantage of the exact Feshbach coordinates is that failure produces an
upstream witness rather than an untyped small eigenvalue.

\begin{theorem}[Soft-vector localization on the Schur graph]
\label{thm:V3-soft-localization}
Assume \(C\succeq\gamma Q\), \(S\succeq0\), and
\(\norm{\mathcal R x}\le\beta\norm x\) for the head coordinate of a unit
vector \(w=x\oplus y\).  If
\begin{equation}
 \ip{w}{Hw}\le\varepsilon<\gamma,
 \label{eq:V3-soft-premise}
\end{equation}
then, with \(z=y+\mathcal R x\),
\begin{equation}
 \boxed{
 \norm z\le\sqrt{\varepsilon/\gamma},
 \qquad
 \ip{x}{Sx}\le\varepsilon,}
 \label{eq:V3-soft-two-bounds}
\end{equation}
and \(x\ne0\).  Moreover
\begin{equation}
 \boxed{
 \frac{\ip{x}{Sx}}{\norm x^2}
 \le
 \frac{\varepsilon(1+\beta^2)}
 {(1-\sqrt{\varepsilon/\gamma})^2}.}
 \label{eq:V3-effective-soft-Rayleigh}
\end{equation}
Thus any physical soft vector below the ultraviolet floor lies within
\(\sqrt{\varepsilon/\gamma}\) of the Feshbach follower graph and returns an
explicit soft Rayleigh vector for the effective head.
\end{theorem}

\begin{proof}
The two terms in \eqref{eq:V3-Feshbach-Pythagoras} are nonnegative.  Hence
\(\ip{z}{Cz}\le\varepsilon\) and
\(\ip{x}{Sx}\le\varepsilon\), proving
\eqref{eq:V3-soft-two-bounds}.  If \(x=0\), then \(w=z\) and
\(\ip{w}{Hw}\ge\gamma\), contrary to \(\varepsilon<\gamma\).

By the graph/innovation decomposition and the router bound,
\[
 1=\norm w
 \le\sqrt{1+\beta^2}\norm x+\norm z
 \le\sqrt{1+\beta^2}\norm x+\sqrt{\varepsilon/\gamma}.
\]
Therefore
\[
 \norm x^2\ge
 \frac{(1-\sqrt{\varepsilon/\gamma})^2}{1+\beta^2}.
\]
Divide \(\ip{x}{Sx}\le\varepsilon\) by this lower bound.
\end{proof}

\begin{corollary}[Typed failure of a proposed orbit floor]
\label{cor:V3-typed-failure}
On a cutoff sequence with a uniform complement floor \(\gamma_*\) and router
bound \(\beta_*\), every normalized local orbit vector with energy tending to
zero produces a normalized Schur-head vector with Rayleigh quotient tending
to zero and a complement innovation tending to zero.  Therefore the only
remaining low-energy mechanism lies in the effective head, unless the claimed
complement floor or router bound itself fails.
\end{corollary}

\begin{proof}
Apply \Cref{thm:V3-soft-localization} with
\(\varepsilon_j\to0\), \(\gamma=\gamma_*\), and
\(\beta=\beta_*\), then normalize \(x_j\).
\end{proof}

% =============================================================================
\section{An explicit Casimir ultraviolet floor on a finite gauge graph}
\label{sec:V3-Casimir-tail}
% =============================================================================

Let \(\Gamma\) be a finite oriented spatial graph with link configuration
space \(G^{E(\Gamma)}\), where \(G\) is compact.  On the gauge-invariant
Hilbert space let
\begin{equation}
 \mathcal C_E:=\sum_{e\in E(\Gamma)}(-\Delta_e)\succeq0
 \label{eq:V3-electric-Casimir}
\end{equation}
be the electric Peter--Weyl Casimir.  Let \(V_B\) be a bounded nonnegative
gauge-invariant magnetic multiplication operator, and put
\begin{equation}
 H_{\rm raw}=\kappa\mathcal C_E+V_B,
 \qquad
 E_0=\inf\spec H_{\rm raw},
 \qquad
 H=H_{\rm raw}-E_0I\succeq0.
 \label{eq:V3-HKS}
\end{equation}

For a Casimir threshold \(\Lambda\), let
\begin{equation}
 P_\Lambda=\one_{[0,\Lambda]}(\mathcal C_E),
 \qquad Q_\Lambda=I-P_\Lambda,
 \qquad
 \Lambda_+=\inf\spec(\mathcal C_E|_{Q_\Lambda\cH}).
 \label{eq:V3-Casimir-cut}
\end{equation}
At a finite graph, \(P_\Lambda\) has finite rank after gauge reduction.

\begin{theorem}[Casimir-tail floor and router]
\label{thm:V3-Casimir-router}
Let \(V_{\max}=\norm{V_B}_{\mathrm{op}}\).  Relative to
\(P_\Lambda\oplus Q_\Lambda\), the complement block and mixed block of
\(H\) satisfy
\begin{equation}
 \boxed{
 C_\Lambda:=Q_\Lambda HQ_\Lambda
 \succeq\gamma_\Lambda Q_\Lambda,
 \qquad
 \gamma_\Lambda:=\kappa\Lambda_+-E_0
 \ge\kappa\Lambda_+-V_{\max},}
 \label{eq:V3-Casimir-gamma}
\end{equation}
and
\begin{equation}
 \boxed{
 B_\Lambda=P_\Lambda V_BQ_\Lambda,
 \qquad
 \norm{B_\Lambda}\le V_{\max}.}
 \label{eq:V3-Casimir-B}
\end{equation}
Whenever \(\gamma_\Lambda>0\), the Feshbach router obeys
\begin{equation}
 \boxed{
 \norm{C_\Lambda^{-1}B_\Lambda^*}
 \le\frac{V_{\max}}{\gamma_\Lambda}.}
 \label{eq:V3-Casimir-beta}
\end{equation}
The effective head is the finite matrix
\begin{equation}
 \boxed{
 S_\Lambda
 =P_\Lambda HP_\Lambda
  -P_\Lambda V_BQ_\Lambda
   C_\Lambda^{-1}
   Q_\Lambda V_BP_\Lambda\succeq0.}
 \label{eq:V3-Casimir-Schur}
\end{equation}
\end{theorem}

\begin{proof}
The constant function is an admissible gauge-invariant trial vector and has
zero electric energy.  Since \(0\le V_B\le V_{\max}I\), the variational
principle gives \(0\le E_0\le V_{\max}\).  On the range of \(Q_\Lambda\),
\[
 Q_\Lambda HQ_\Lambda
 =\kappa Q_\Lambda\mathcal C_EQ_\Lambda
  +Q_\Lambda V_BQ_\Lambda-E_0Q_\Lambda
 \succeq(\kappa\Lambda_+-E_0)Q_\Lambda.
\]
Because \(P_\Lambda\) is a spectral projection of \(\mathcal C_E\), the
electric and scalar terms have no mixed block.  Hence
\(P_\Lambda HQ_\Lambda=P_\Lambda V_BQ_\Lambda\), whose norm is at most
\(V_{\max}\).  Inverting the complement floor gives
\eqref{eq:V3-Casimir-beta}.  Formula
\eqref{eq:V3-Casimir-Schur} is the Schur complement; its positivity follows
from \(H\succeq0\) and \(C_\Lambda\succ0\).
\end{proof}

\begin{corollary}[Finite \(\mathrm{SU}(3)\) magnetic bound]
\label{cor:V3-SU3-Casimir}
For the magnetic Wilson potential
\begin{equation}
 V_B(U)=\lambda_B\sum_{p\in P(\Gamma)}
 \left(1-\frac13\operatorname{ReTr}U_p\right),
 \qquad\lambda_B>0,
 \label{eq:V3-Wilson-magnetic}
\end{equation}
one may take
\begin{equation}
 \boxed{V_{\max}\le2\lambda_B|P(\Gamma)|.}
 \label{eq:V3-Wilson-Vmax}
\end{equation}
If
\begin{equation}
 \kappa\Lambda_+\ge(1+\eta)V_{\max}
 \qquad(\eta>0),
 \label{eq:V3-threshold-choice}
\end{equation}
then
\begin{equation}
 \boxed{
 \gamma_\Lambda\ge\eta V_{\max},
 \qquad
 \norm{\mathcal R_\Lambda}\le\eta^{-1}.}
 \label{eq:V3-uniform-tail-router}
\end{equation}
Thus the electric representation tail has an explicit positive floor and a
bounded router after the Casimir threshold is chosen above the total magnetic
oscillation scale.
\end{corollary}

\begin{proof}
The elementary bound
\(-1\le\operatorname{ReTr}U/3\le1\) gives
\eqref{eq:V3-Wilson-Vmax}.  Insert
\eqref{eq:V3-threshold-choice} into
\eqref{eq:V3-Casimir-gamma} and then use
\eqref{eq:V3-Casimir-beta}.
\end{proof}

\begin{assessment}[What the Casimir cut closes]
At every fixed Hamiltonian lattice regulator, the ultraviolet complement is
not the unresolved part of the gap problem: its energy floor and relaxation
router can be made explicit, and all remaining low-energy information is
contained in the finite Schur matrix \(S_\Lambda\).  Cofinally, however,
\(V_{\max}\) grows with the carrier, so the required Casimir head also grows.
The theorem is therefore a rigorous ultraviolet reduction, not a uniform
infrared mass proof.
\end{assessment}

\begin{firewall}[Hamiltonian regulator versus the Wilson OS generator]
The operator in \eqref{eq:V3-HKS} is a Hamiltonian lattice gauge regulator.
It may be used for the V2 orbit of the Wilson Euclidean theory only after a
same-Hilbert-space form identification or a controlled Hamiltonian-limit
comparison with the actual normalized OS transfer generator
\(-\tau^{-1}\log P_{\rm OS}\).  Formal similarity of electric and magnetic
terms is not that comparison.
\end{firewall}

% =============================================================================
\section{The exact same-generator comparison needed for transfer}
\label{sec:V3-same-generator}
% =============================================================================

\begin{theorem}[Orbit-form comparison transfers a Feshbach floor]
\label{thm:V3-form-comparison}
Let \(H_{\rm OS}\succeq0\) be the actual OS Hamiltonian and let
\(\widetilde H\succeq0\) be a comparison Hamiltonian represented on the same
Hilbert space after a declared unitary law alignment.  Let
\(w_s=e^{-sH_{\rm OS}}Vc\).  Suppose the Feshbach analysis of
\(\widetilde H\) gives
\begin{equation}
 \ip{w_s}{\widetilde H w_s}
 \ge m_*\norm{w_s}^2-r_1(s)c^*R_Ec,
 \label{eq:V3-comparison-floor}
\end{equation}
and the same-orbit form mismatch obeys
\begin{equation}
 \left|
 \ip{w_s}{(H_{\rm OS}-\widetilde H)w_s}
 \right|
 \le r_2(s)c^*R_Ec.
 \label{eq:V3-generator-mismatch}
\end{equation}
Then the actual OS orbit satisfies
\begin{equation}
 \boxed{
 \ip{w_s}{H_{\rm OS}w_s}
 \ge m_*\norm{w_s}^2-(r_1(s)+r_2(s))c^*R_Ec.}
 \label{eq:V3-physical-orbit-floor}
\end{equation}
If \(R_E\preceq C_EG(0)\) and the weighted integral of
\(r_1+r_2\) satisfies \eqref{eq:V2-q-residual}, the actual OS delayed Gram
contracts and V1 applies.
\end{theorem}

\begin{proof}
Write
\[
 \ip{w_s}{H_{\rm OS}w_s}
 =\ip{w_s}{\widetilde H w_s}
  +\ip{w_s}{(H_{\rm OS}-\widetilde H)w_s}.
\]
Use the lower side of \eqref{eq:V3-generator-mismatch} and then
\eqref{eq:V3-comparison-floor}.  The last statement is
\Cref{thm:V2-orbit-coercivity}.
\end{proof}

\begin{firewall}[A static matrix is not a comparison Hamiltonian]
The V35 trace-shell Hessian can contribute to
\(\widetilde H\) only after its coefficient synthesis, units, vacuum short,
domain, and quadratic form are identified on the same OS Hilbert space.  The
form mismatch \eqref{eq:V3-generator-mismatch} must then be bounded on vectors
evolved by \(H_{\rm OS}\), not by the comparison flow.  Without this row, a
positive finite Schur head remains static information.
\end{firewall}

% =============================================================================
\section{Cofinal orbit-Feshbach compiler}
\label{sec:V3-cofinal-compiler}
% =============================================================================

\begin{theorem}[Cofinal local-bank mass compiler]
\label{thm:V3-cofinal-compiler}
For every cutoff \(j\), let \(H_j\) be the actual OS Hamiltonian, let
\(V_j:E\to\Omega_j^\perp\) be one fixed continuum-labelled local bank, and
let \(P_j\oplus Q_j\) be a physical head/complement split.  Assume, uniformly
for \(0\le s\le\tau_j\),
\begin{enumerate}[label=\textup{(F\arabic*)},leftmargin=3.4em]
\item \(C_j\succeq\gamma_*Q_j\) with \(\gamma_*>0\);
\item
\(\norm{C_j^{-1}B_j^*x}\le\beta_*\norm x\) for every
\(x=P_je^{-sH_j}V_jc\);
\item for \(S_j=A_j-B_jC_j^{-1}B_j^*\),
\begin{equation}
 \ip{x}{S_jx}
 \ge\delta_*\norm x^2-r_j(s)c^*R_jc,
 \qquad\delta_*>0;
 \label{eq:V3-cofinal-Schur-floor}
\end{equation}
\item \(R_j\preceq C_*G_j(0)\),
\(\tau_j\to\tau>0\), and
\(\int_0^{\tau_j}r_j(s)\,\dd s\to0\).
\end{enumerate}
Then for every
\begin{equation}
 q\in\left(e^{-m_{\rm F}(\delta_*,\gamma_*,\beta_*)\tau},1\right)
 \label{eq:V3-q-window}
\end{equation}
one eventually has
\begin{equation}
 \boxed{G_j(\tau_j)\preceq q^2G_j(0).}
 \label{eq:V3-cofinal-delay}
\end{equation}
For this fixed bank, local OS kernel convergence gives the corresponding
limiting finite-bank delayed inequality.  If the assumptions hold with the
same constants for every bank in an increasing local hierarchy whose centered
union is dense, and nontrivial lower frames also converge as in V1, then the
limiting Hamiltonian has a unique vacuum and mass at least
\begin{equation}
 \boxed{m_{\rm F}(\delta_*,\gamma_*,\beta_*)>0}
 \label{eq:V3-compiled-mass}
\end{equation}
on the full local vacuum complement.
\end{theorem}

\begin{proof}
By \Cref{thm:V3-orbit-Feshbach-residual}, assumptions
\textup{(F1)--(F3)} give the V2 orbit floor with
\(m=m_{\rm F}(\delta_*,\gamma_*,\beta_*)\).  Assumption \textup{(F4)}
and \Cref{cor:V2-cutoff-orbit} give
\eqref{eq:V3-cofinal-delay}.  V1's quotient-stable passage theorem passes
this fixed-bank inequality to the limiting OS quotient.  When the same
constants apply to every member of the declared increasing hierarchy, the
inequality holds on its dense centered union.  V1's local Gram criterion then
gives the mass on the full vacuum complement, while the lower frames prevent a
vacuous quotient.
\end{proof}

\begin{assessment}[The exact remaining analytic row]
The ultraviolet floor and router can be populated on the Hamiltonian Casimir
branch by \Cref{thm:V3-Casimir-router}.  The unresolved inequality is the
cutoff-uniform effective-head estimate
\eqref{eq:V3-cofinal-Schur-floor} on the \emph{evolved} local bank, together
with the same-generator comparison to the Wilson OS transfer.  This is an
infrared, renormalization, and continuum-construction problem.  It cannot be
closed by increasing the Casimir threshold alone.
\end{assessment}

% =============================================================================
\section{V3 ledger and next attack}
\label{sec:V3-ledger}
% =============================================================================

\begingroup
\small
\begin{longtable}{>{\raggedright\arraybackslash}p{0.28\textwidth}
                  >{\raggedright\arraybackslash}p{0.16\textwidth}
                  >{\raggedright\arraybackslash}p{0.48\textwidth}}
\toprule
\textbf{V3 row}&\textbf{Status}&\textbf{Advance or limitation}\\
\midrule
\endfirsthead
\toprule
\textbf{V3 row}&\textbf{Status}&\textbf{Advance or limitation}\\
\midrule
\endhead
Graph/innovation Feshbach identity
& \MGpass
& Exact completion of the physical quadratic form.\\[0.35em]
Orbit-restricted Feshbach floor
& \MGpass
& Sharp constant \eqref{eq:V3-sharp-mF}; only the evolved head range is tested.\\[0.35em]
Soft-vector return
& \MGpass
& Every vector below the ultraviolet floor yields a quantitative effective-head soft witness.\\[0.35em]
Finite Casimir ultraviolet tail
& \MGpass
& Explicit complement floor, mixed block, router, and finite Schur head for a Hamiltonian lattice gauge regulator.\\[0.35em]
Cofinal Casimir head size
& Open uniformity cost
& The magnetic norm grows with the graph, forcing a growing Peter--Weyl head.\\[0.35em]
Wilson OS/Kogut--Susskind identification
& \MGopen
& Requires a same-space form comparison on the actual OS orbit.\\[0.35em]
Renormalized Schur-head floor
& \MGopen
& No uniform positive \(\delta_*\) is proved on a response-matched evolved local bank.\\[0.35em]
Continuum local OS state
& \MGopen
& Reflection positivity at finite cutoff is insufficient; convergence, nontriviality, and lower frames remain required.\\[0.35em]
Full Yang--Mills mass gap
& \MGopen
& Follows from the compiled theorem only after the two preceding physical rows are proved.\\
\bottomrule
\end{longtable}
\endgroup

\begin{openproblem}[V4 mandate: control the effective low-Casimir head]
\label{op:V4}
Freeze one response-matched trajectory, one fixed physical delay, and one
bounded-support centre-neutral Wilson bank.  On the Hamiltonian branch, choose
the Casimir threshold by \eqref{eq:V3-threshold-choice}; on the direct Wilson
branch, choose an equivalent reflection-transfer head.  Then:
\begin{enumerate}[label=\textup{(V4.\arabic*)},leftmargin=5em]
\item construct the actual orbit-head matrices
\(
X_j(s)^*S_jX_j(s)
\)
and entrance metrics
\(
X_j(s)^*X_j(s)
\)
for \(X_j(s)=P_je^{-sH_j}V_j\);
\item short the exact vacuum graph direction rather than the constant
      function or a fitted static null;
\item prove a root- and cutoff-uniform generalized lower bound with an
      integrable residual, or return the first soft generalized eigenvector;
\item on the Hamiltonian route, prove
      \eqref{eq:V3-generator-mismatch} against the actual Wilson OS generator;
\item apply \Cref{thm:V3-cofinal-compiler} and then enlarge the local bank.
\end{enumerate}
If the generalized head floor collapses, go upstream to its returned
low-Casimir Wilson-network combination, renormalized coupling, phase, or
lower-frame defect.  Do not charge the ultraviolet tail twice and do not
replace the signed Schur return by its absolute magnetic norm.
\end{openproblem}

\begin{firewall}[End-of-V3 claim boundary]
V3 rigorously removes the high-Casimir complement as an untyped finite-cutoff
bottleneck and proves the sharp orbit-restricted constant needed to feed V2.
It does not prove a cofinal effective-head floor or the Wilson
OS/Hamiltonian-regulator comparison.  The full Yang--Mills mass gap remains
open.
\end{firewall}

\section*{Version V3 append-only record}
\addcontentsline{toc}{section}{Version V3 append-only record}
The complete V2 source was retained before this continuation.  Its SHA--256
digest was
\begin{center}\scriptsize\ttfamily
3cd638dd43043a2e75a7f52e083a2d8791fe16fccf52c9e8633611d4e3e3d4b3.
\end{center}
On the next iteration, retain this full file, append before its unique active
document-closing command, increment the filename to
\texttt{\_V4.tex}, and remove only the superseded V3 file from this note
series.

% =============================================================================
% END APPENDED VERSION V3 MATERIAL
% =============================================================================


% =============================================================================
% BEGIN APPENDED VERSION V4 MATERIAL
% =============================================================================

\clearpage
\part*{Version V4 continuation: the exact vacuum graph, the Feshbach metric, and the growing-head criterion}
\addcontentsline{toc}{part}{Version V4 continuation: the exact vacuum graph, the Feshbach metric, and the growing-head criterion}

% =============================================================================
\section{V4 direction: short the graph vacuum in its physical metric}
\label{sec:V4-direction}
% =============================================================================

V3 reduced the physical orbit form to an effective head \(S\), a positive
complement \(C\), and a relaxation router \(\mathcal R=C^{-1}B^*\).  Its V4
mandate asked that the exact vacuum graph direction be shorted before a head
eigenvalue is read.  This is essential: after the complement is minimized, the
vacuum is generally not the bare constant head vector, and the physical norm
on the follower graph is not the ambient head norm.

V4 proves that the correct head metric is
\begin{equation}
 \boxed{M:=I+\mathcal R^*\mathcal R.}
 \label{eq:V4-M-metric}
\end{equation}
The kernel of the full Hamiltonian is isomorphic to the kernel of \(S\), with
the isomorphism measured by \(M\).  We derive the exact vacuum-shortened
generalized eigenvalue, a quantitative full-space floor, a fixed-dimension
transport theorem, and a growing-head compactness criterion.  These results
replace an ambiguous ``remove the constant mode'' instruction by a finite,
same-Hamiltonian card.

% =============================================================================
\section{Kernel transport and the exact vacuum graph}
\label{sec:V4-vacuum-graph}
% =============================================================================

Retain the V3 block hypotheses
\[
 H=\begin{pmatrix}A&B\\B^*&C\end{pmatrix}\succeq0,
 \qquad C\succeq\gamma I,
 \qquad
 \mathcal R=C^{-1}B^*,
 \qquad
 S=A-BC^{-1}B^*.
\]
Define the graph synthesis
\begin{equation}
 J:P\cH\longrightarrow\cH,
 \qquad
 Jx=x\oplus(-\mathcal R x).
 \label{eq:V4-J-graph}
\end{equation}

\begin{theorem}[Exact kernel--graph isomorphism]
\label{thm:V4-kernel-graph}
One has
\begin{equation}
 \boxed{J^*J=M=I+\mathcal R^*\mathcal R,}
 \label{eq:V4-J-isometry-metric}
\end{equation}
and
\begin{equation}
 \boxed{
 \Ker H=J(\Ker S).}
 \label{eq:V4-kernel-isomorphism}
\end{equation}
The map \(J:(P\cH,M)\to\cH\) is an isometry onto the Feshbach follower
graph.  In particular, the vacuum of \(H\) is unique if and only if
\(\Ker S\) is one dimensional.
\end{theorem}

\begin{proof}
The norm identity
\[
 \norm{Jx}^2=\norm x^2+\norm{\mathcal R x}^2=\ip{x}{Mx}
\]
proves \eqref{eq:V4-J-isometry-metric}.  If
\(w=x\oplus y\in\Ker H\), its lower block equation is
\(B^*x+Cy=0\), hence \(y=-C^{-1}B^*x=-\mathcal R x\).  The upper block
equation becomes \(Sx=0\), so \(w=Jx\) with \(x\in\Ker S\).
Conversely, direct block multiplication gives
\[
 HJx=\binom{Sx}{0},
\]
so every \(x\in\Ker S\) produces a vector in \(\Ker H\).  Since \(J\) is
injective, it preserves kernel dimension.
\end{proof}

Let \(\Omega\) be a normalized vacuum of \(H\).  By
\Cref{thm:V4-kernel-graph}, there is a unique \(x_0\in\Ker S\) such that
\begin{equation}
 \boxed{
 \Omega=Jx_0,
 \qquad
 \ip{x_0}{Mx_0}=1.}
 \label{eq:V4-vacuum-head}
\end{equation}
The correct head vacuum complement is
\begin{equation}
 \cX_0^{\perp_M}
 :=\{x\in P\cH:\ip{x_0}{Mx}=0\}.
 \label{eq:V4-M-vacuum-complement}
\end{equation}

\begin{firewall}[The bare head constant is not the vacuum short]
Unless \(\mathcal R x_0=0\), the physical vacuum has a nonzero complement
component.  Deleting the bare head vector \(x_0\oplus0\), or the constant
function before the Perron/OS transform, does not remove \(\Omega\) and can
alter the lowest generalized eigenvalue.  The short must use
\eqref{eq:V4-vacuum-head} and the metric \(M\).
\end{firewall}

% =============================================================================
\section{Exact centered coordinates and a vacuum-corrected gap}
\label{sec:V4-centered-coordinates}
% =============================================================================

For \(w=x\oplus y\), retain the V3 innovation
\(z=y+\mathcal R x\).  Suppose now that \(w\perp\Omega\).  Define
\begin{equation}
 \alpha:=\ip{x_0}{Mx},
 \qquad
 x_\perp:=x-\alpha x_0.
 \label{eq:V4-xperp}
\end{equation}

\begin{lemma}[Centered graph/innovation decomposition]
\label{lem:V4-centered-decomposition}
For every \(w\perp\Omega\),
\begin{equation}
 \boxed{
 \alpha=\ip{\mathcal R x_0}{z},
 \qquad
 x_\perp\in\cX_0^{\perp_M},}
 \label{eq:V4-centering-relation}
\end{equation}
and
\begin{equation}
 \boxed{
 w=Jx_\perp+Zz,
 \qquad
 Zz:=(0\oplus z)+\ip{\mathcal R x_0}{z}\,\Omega.}
 \label{eq:V4-centered-split}
\end{equation}
Moreover, \(Zz\perp\Omega\), \(Z\) is the orthogonal projection of
\(0\oplus z\) onto \(\Omega^\perp\), and therefore
\begin{equation}
 \boxed{\norm{Zz}\le\norm z.}
 \label{eq:V4-Z-contraction}
\end{equation}
The energy remains exactly
\begin{equation}
 \boxed{
 \ip{w}{Hw}=\ip{x_\perp}{Sx_\perp}+\ip{z}{Cz}.}
 \label{eq:V4-centered-energy}
\end{equation}
\end{lemma}

\begin{proof}
Since \(\Omega=Jx_0=x_0\oplus(-\mathcal R x_0)\), the centering condition
and \(w=Jx+(0\oplus z)\) give
\[
 0=\ip{\Omega}{w}
  =\ip{x_0}{Mx}-\ip{\mathcal R x_0}{z}.
\]
This proves the first identity and makes \(x_\perp\) \(M\)-orthogonal to
\(x_0\).  Substitution gives \eqref{eq:V4-centered-split}.  The vector
\(Zz\) is exactly
\((I-|\Omega\rangle\langle\Omega|)(0\oplus z)\), with the displayed sign
under the sesquilinear convention of the OS form.  Hence it is centered and
has norm at most \(\norm z\).  Finally, the V3 Feshbach identity gives
\(\ip{x}{Sx}+\ip{z}{Cz}\).  Because \(Sx_0=0\) and \(S=S^*\), replacing
\(x\) by \(x_\perp\) leaves the first term unchanged.
\end{proof}

\begin{theorem}[Vacuum-corrected Feshbach gap]
\label{thm:V4-vacuum-corrected-gap}
Assume
\begin{equation}
 \boxed{
 \ip{x}{Sx}\ge\delta\ip{x}{Mx}
 \quad\text{for every }x\in\cX_0^{\perp_M},
 \qquad \delta>0,}
 \label{eq:V4-effective-M-gap}
\end{equation}
and \(C\succeq\gamma I\).  Then \(\Omega\) is the unique vacuum and
\begin{equation}
 \boxed{
 H|_{\Omega^\perp}
 \succeq m_{\rm vac}I,
 \qquad
 m_{\rm vac}:=\frac{\delta\gamma}{\delta+\gamma}>0.}
 \label{eq:V4-vacuum-gap}
\end{equation}
The same conclusion holds on any selected centered orbit set if
\eqref{eq:V4-effective-M-gap} is required only for the corresponding
\(x_\perp\)'s.
\end{theorem}

\begin{proof}
For \(w\perp\Omega\), let
\(a=\norm{Jx_\perp}=\ip{x_\perp}{Mx_\perp}^{1/2}\) and
\(b=\norm z\).  By \Cref{lem:V4-centered-decomposition},
\[
 \ip{w}{Hw}\ge\delta a^2+\gamma b^2,
 \qquad
 \norm w\le a+b.
\]
Weighted Cauchy--Schwarz gives
\[
 (a+b)^2
 \le(\delta a^2+\gamma b^2)
    (\delta^{-1}+\gamma^{-1}).
\]
Rearrangement proves \eqref{eq:V4-vacuum-gap}.  A second vacuum in
\(\Omega^\perp\) would contradict the strict lower bound.  The proof used the
head inequality only at the actual \(x_\perp\), which proves the orbit-set
statement.
\end{proof}

\begin{theorem}[Vacuum-corrected orbit floor with residual]
\label{thm:V4-vacuum-orbit-residual}
Let \(w_s=e^{-sH}Vc\) with \(V:E\to\Omega^\perp\), and form
\(x_{\perp,s}\) and \(z_s\) as above.  If
\begin{equation}
 \ip{x_{\perp,s}}{Sx_{\perp,s}}
 \ge\delta\ip{x_{\perp,s}}{Mx_{\perp,s}}
       -r(s)c^*R_Ec
 \label{eq:V4-orbit-M-residual}
\end{equation}
for almost every \(s\in[0,\tau]\), while
\(C\succeq\gamma I\), then
\begin{equation}
 \boxed{
 K(s)\succeq
 \frac{\delta\gamma}{\delta+\gamma}G(s)-r(s)R_E.}
 \label{eq:V4-M-to-orbit-floor}
\end{equation}
Hence the integrated-residual theorem of V2 gives a delayed contraction.  No
separate router norm is needed because its norm contribution is retained
exactly in \(M\).
\end{theorem}

\begin{proof}
Apply the proof of \Cref{thm:V4-vacuum-corrected-gap} to each \(w_s\), carrying
the coefficient residual through the energy inequality.  This yields the
quadratic-form version of \eqref{eq:V4-M-to-orbit-floor}; apply
\Cref{thm:V2-orbit-coercivity}.
\end{proof}

\begin{remark}[Relation to the V3 constant]
If one knows only
\(S\succeq\delta_I I\) and
\(\norm{\mathcal R}\le\beta\), then
\(M\preceq(1+\beta^2)I\), so one may take
\(\delta=\delta_I/(1+\beta^2)\) in
\eqref{eq:V4-effective-M-gap}.  V3's two-variable generalized-eigenvalue
constant is sharper for those three scalar inputs.  V4's advantage is that
the full router geometry and the exact vacuum short are retained in one finite
generalized head card.
\end{remark}

% =============================================================================
\section{The exact finite vacuum-shortened head card}
\label{sec:V4-finite-head-card}
% =============================================================================

Assume \(P\cH\) is finite dimensional.  Since \(M\succ0\), put
\begin{equation}
 u_0:=M^{1/2}x_0,
 \qquad \norm{u_0}=1,
 \qquad
 \widehat S:=M^{-1/2}SM^{-1/2}.
 \label{eq:V4-whitened-head}
\end{equation}
Then \(\widehat S u_0=0\).

\begin{theorem}[Exact generalized head gap and kernel test]
\label{thm:V4-exact-head-card}
Define
\begin{equation}
 \boxed{
 \delta_{\rm head}
 :=\inf_{\substack{x\ne0\\\ip{x_0}{Mx}=0}}
 \frac{\ip{x}{Sx}}{\ip{x}{Mx}}
 =\lambda_{\min}
 \left(\widehat S|_{u_0^\perp}\right).}
 \label{eq:V4-delta-head}
\end{equation}
Then
\begin{equation}
 \boxed{
 \delta_{\rm head}>0
 \quad\Longleftrightarrow\quad
 \Ker S=\mathbb Cx_0
 \quad\Longleftrightarrow\quad
 \Ker H=\mathbb C\Omega.}
 \label{eq:V4-head-kernel-equivalence}
\end{equation}
On the positive branch,
\begin{equation}
 \boxed{
 H|_{\Omega^\perp}
 \succeq
 \frac{\delta_{\rm head}\gamma}
      {\delta_{\rm head}+\gamma}I.}
 \label{eq:V4-finite-full-gap}
\end{equation}
Thus, once \(C\), \(B\), and the vacuum are taken from the same finite
Hamiltonian, the remaining head decision is one finite constrained
generalized eigenvalue.
\end{theorem}

\begin{proof}
The change of variables \(u=M^{1/2}x\) maps the \(M\)-unit sphere to the
ordinary unit sphere and maps the constraint
\(\ip{x_0}{Mx}=0\) to \(u\perp u_0\).  This proves the second expression in
\eqref{eq:V4-delta-head}.  In finite dimension, the restricted minimum is
positive exactly when the only kernel direction of \(\widehat S\) is
\(u_0\), equivalently when \(\Ker S=\mathbb Cx_0\).  The full-kernel
equivalence is \Cref{thm:V4-kernel-graph}.  Finally apply
\Cref{thm:V4-vacuum-corrected-gap} with
\(\delta=\delta_{\rm head}\).
\end{proof}

\begin{corollary}[Finite positivity-improving regulators close this card]
\label{cor:V4-finite-positive-card}
If a finite OS transfer is positivity improving in one fixed protected sector,
its vacuum is simple.  For every Feshbach split with
\(C\succeq\gamma I\) and finite head, one therefore has
\(\delta_{\rm head}>0\) and the explicit finite lower bound
\eqref{eq:V4-finite-full-gap}.  This proves only a finite-regulator mass; no
uniform lower bound on \(\delta_{\rm head}\) follows.
\end{corollary}

\begin{proof}
Positivity improvement gives a simple Perron vacuum.  Use
\eqref{eq:V4-head-kernel-equivalence} and then
\eqref{eq:V4-finite-full-gap}.
\end{proof}

% =============================================================================
\section{Why the router must remain in the head metric}
\label{sec:V4-router-counterexample}
% =============================================================================

\begin{counterexample}[A fixed ordinary Schur gap can hide a collapsing physical gap]
\label{cth:V4-router-collapse}
For \(n\ge1\), let
\begin{equation}
 H_n=\begin{pmatrix}n^2+1&-n\\-n&1\end{pmatrix}.
 \label{eq:V4-router-example-H}
\end{equation}
With the first coordinate as head, one has
\begin{equation}
 C_n=1,
 \qquad
 \mathcal R_n=-n,
 \qquad
 S_n=1,
 \qquad
 M_n=1+n^2.
 \label{eq:V4-router-example-blocks}
\end{equation}
Thus the ordinary Schur gap is identically one, whereas its physical
\(M_n\)-normalized value is \((1+n^2)^{-1}\to0\).  Correspondingly,
\begin{equation}
 \lambda_{\min}(H_n)
 =\frac{n^2+2-\sqrt{(n^2+2)^2-4}}2
 \longrightarrow0.
 \label{eq:V4-router-example-gap}
\end{equation}
Therefore a Euclidean head eigenvalue cannot be transported cofinally without
the Feshbach metric or an independent router bound.
\end{counterexample}

\begin{proof}
The block formulas are direct.  The eigenvalue expression follows from
\(\operatorname{tr}H_n=n^2+2\) and \(\det H_n=1\); rationalizing the smaller
root shows that it is asymptotic to \(n^{-2}\).
\end{proof}

% =============================================================================
\section{Fixed-head transport is a norm-convergence problem}
\label{sec:V4-fixed-head-transport}
% =============================================================================

\begin{theorem}[Stability of the vacuum-shortened finite head]
\label{thm:V4-fixed-head-stability}
Let a fixed finite-dimensional coefficient space carry matrices
\(M_j\succ0\), \(S_j\succeq0\), and vectors \(x_{0,j}\) such that
\begin{equation}
 S_jx_{0,j}=0,
 \qquad
 \ip{x_{0,j}}{M_jx_{0,j}}=1.
 \label{eq:V4-finite-head-data}
\end{equation}
Assume
\begin{equation}
 M_j\to M_\infty\succ0,
 \qquad
 S_j\to S_\infty,
 \qquad
 x_{0,j}\to x_{0,\infty}
 \label{eq:V4-head-norm-convergence}
\end{equation}
in norm.  Then
\begin{equation}
 \boxed{\delta_{\rm head,j}\longrightarrow
 \delta_{\rm head,\infty}.}
 \label{eq:V4-delta-convergence}
\end{equation}
In particular, if
\(\Ker S_\infty=\mathbb Cx_{0,\infty}\), then
\(\inf_{j\ge j_0}\delta_{\rm head,j}>0\) for some \(j_0\).
\end{theorem}

\begin{proof}
Continuity of the positive square root and inverse on matrices bounded away
from zero gives
\[
 M_j^{-1/2}S_jM_j^{-1/2}
 \longrightarrow
 M_\infty^{-1/2}S_\infty M_\infty^{-1/2}.
\]
The normalized kernel vectors
\(u_{0,j}=M_j^{1/2}x_{0,j}\) converge to \(u_{0,\infty}\).  After a
convergent unitary rotation sending \(u_{0,j}\) to a fixed first basis vector,
the restrictions to the orthogonal complements converge in operator norm.
Weyl's eigenvalue inequality proves \eqref{eq:V4-delta-convergence}.  The
kernel hypothesis makes the limiting restricted minimum strictly positive.
\end{proof}

\begin{firewall}[Fixed dimension is essential]
The theorem cannot be applied to the full cofinal Casimir head merely because
every fixed finite corner converges.  New head directions enter as the graph,
volume, or representation threshold grows, and they may carry the smallest
generalized eigenvalue.
\end{firewall}

\begin{counterexample}[Every fixed head corner is gapped while the growing head collapses]
\label{cth:V4-growing-head-collapse}
On \(\mathbb C^{j+1}\), let \(M_j=I\), let the vacuum be \(e_0\), and let
\begin{equation}
 S_j=\operatorname{diag}(0,1,\ldots,1,j^{-1}).
 \label{eq:V4-growing-head-example}
\end{equation}
For every fixed \(m\), the compression to
\(\operatorname{span}\{e_0,\ldots,e_m\}\) has vacuum-shortened gap one for
all sufficiently large \(j\), but
\begin{equation}
 \boxed{\delta_{\rm head,j}=j^{-1}\longrightarrow0.}
 \label{eq:V4-growing-head-zero}
\end{equation}
Thus pointwise convergence of every predeclared Wilson-network bank does not
control the collective head tail.
\end{counterexample}

% =============================================================================
\section{A collective criterion inside the whitened head}
\label{sec:V4-collective-head}
% =============================================================================

The growing-head obstruction can itself be split once, now entirely inside
the physical metric.  Work in the whitened head
\(\widehat S=M^{-1/2}SM^{-1/2}\), and suppose its vacuum vector \(u_0\) lies
in a retained core \(\cK\).  On \(u_0^\perp\), decompose
\begin{equation}
 u_0^\perp=\cK_0\oplus\cT,
 \qquad
 \widehat S|_{u_0^\perp}
 =\begin{pmatrix}D&L\\L^*&F\end{pmatrix}.
 \label{eq:V4-whitened-head-split}
\end{equation}

\begin{theorem}[Collective growing-head gap criterion]
\label{thm:V4-collective-head}
If
\begin{equation}
 D\succeq dI,
 \qquad
 F\succeq fI,
 \qquad
 \norm L\le\ell,
 \qquad
 d>0, f>0, \ell^2<df,
 \label{eq:V4-collective-inputs}
\end{equation}
then
\begin{equation}
 \boxed{
 \delta_{\rm head}
 \ge
 \frac{d+f-\sqrt{(d-f)^2+4\ell^2}}2>0.}
 \label{eq:V4-collective-bound}
\end{equation}
The statement is uniform over a growing sequence of heads whenever
\((d,f,\ell)\) have common bounds satisfying
\(\ell^2<df\).
\end{theorem}

\begin{proof}
For \(u\oplus v\in\cK_0\oplus\cT\),
\[
 \ip{u\oplus v}{\widehat S(u\oplus v)}
 \ge d\norm u^2-2\ell\norm u\norm v+f\norm v^2.
\]
The right side is the quadratic form of
\(\begin{psmallmatrix}d&-\ell\\-\ell&f\end{psmallmatrix}\) on
\((\norm u,\norm v)\).  Its smaller eigenvalue is
\eqref{eq:V4-collective-bound}, and it is positive exactly when
\(\ell^2<df\).
\end{proof}

\begin{assessment}[The target-minimal cofinal head packet]
After the V3 Casimir elimination and the V4 vacuum short, a uniform head gap
requires only three typed quantities for one declared inner-head/tail split:
\begin{equation}
 \boxed{
 \text{retained centered-head floor }d,
 \quad
 \text{new-head-tail floor }f,
 \quad
 \text{mixed norm }\ell.}
 \label{eq:V4-three-head-stocks}
\end{equation}
The mixed term may not be set to zero.  A finite-bank limit controls \(d\);
the genuinely new analytic work is a uniform \(f\) and a subcritical
\(\ell^2/(df)\) for all entering low-Casimir directions.
\end{assessment}

% =============================================================================
\section{V4 integrated theorem and revised frontier}
\label{sec:V4-frontier}
% =============================================================================

\begin{theorem}[Integrated V4 vacuum-short theorem]
\label{thm:V4-integrated}
Preserving V1--V3, V4 proves:
\begin{enumerate}[label=\textup{(V4.\arabic*)},leftmargin=5em]
\item the full and effective Hamiltonian kernels are exactly isomorphic under
      the Feshbach graph map;
\item the graph metric is \(M=I+\mathcal R^*\mathcal R\), and the vacuum is
      \(Jx_0\), not generally the bare head constant;
\item an \(M\)-normalized Schur gap gives the full centered gap
      \(\delta\gamma/(\delta+\gamma)\);
\item the exact finite head card is one constrained generalized eigenvalue,
      positive exactly when the finite vacuum is unique;
\item a fixed ordinary Schur eigenvalue can coexist with a collapsing physical
      gap when the router grows;
\item norm convergence transports the card at fixed head dimension, but fixed
      corners do not control a growing head;
\item one inner-head/tail split closes the collective head if
      \(\ell^2<df\).
\end{enumerate}
\end{theorem}

\begin{proof}
Items \textup{(V4.1)--(V4.2)} are
\Cref{thm:V4-kernel-graph,lem:V4-centered-decomposition}.  Items
\textup{(V4.3)--(V4.4)} are
\Cref{thm:V4-vacuum-corrected-gap,thm:V4-exact-head-card}.  Item
\textup{(V4.5)} is \Cref{cth:V4-router-collapse}.  Item \textup{(V4.6)} is
\Cref{thm:V4-fixed-head-stability,cth:V4-growing-head-collapse}.  Item
\textup{(V4.7)} is \Cref{thm:V4-collective-head}.
\end{proof}

\begingroup
\small
\begin{longtable}{>{\raggedright\arraybackslash}p{0.28\textwidth}
                  >{\raggedright\arraybackslash}p{0.16\textwidth}
                  >{\raggedright\arraybackslash}p{0.48\textwidth}}
\toprule
\textbf{V4 row}&\textbf{Status}&\textbf{Advance or remaining obstruction}\\
\midrule
\endfirsthead
\toprule
\textbf{V4 row}&\textbf{Status}&\textbf{Advance or remaining obstruction}\\
\midrule
\endhead
Vacuum graph and metric
& \MGpass
& Exact same-Hamiltonian kernel transport and physical head norm.\\[0.35em]
Finite vacuum-short card
& \MGpass
& One constrained generalized eigenvalue; finite positivity follows from a simple vacuum.\\[0.35em]
Router-free orbit statement
& \MGpass
& Router is retained inside \(M\), yielding the V2 orbit floor with an integrated residual.\\[0.35em]
Fixed-head adjacent transport
& \MGpass
& Norm convergence plus a simple limiting kernel yields a uniform finite-head margin.\\[0.35em]
Growing-head uniformity
& Conditional exact criterion
& Requires the three stocks \((d,f,\ell)\) with \(\ell^2<df\).\\[0.35em]
Physical values of \((d,f,\ell)\)
& \MGopen
& No response-matched Wilson trajectory or interval data populate them.\\[0.35em]
OS/Hamiltonian same-generator comparison
& \MGopen
& The V3 form-mismatch row remains unproved.\\[0.35em]
Continuum construction and full mass gap
& \MGopen
& Still require nontrivial local OS convergence and common cofinal bounds.\\
\bottomrule
\end{longtable}
\endgroup

\begin{openproblem}[V5 mandate: populate the collective head-tail packet]
\label{op:V5}
On the response-matched local branch, use the exact metric
\(M_j=I+\mathcal R_j^*\mathcal R_j\) and vacuum
\(u_{0,j}=M_j^{1/2}x_{0,j}\).  Choose an inner head generated by a fixed
bounded-support Wilson-network bank and let its orthogonal complement inside
the low-Casimir head be the entering-head tail.  Then:
\begin{enumerate}[label=\textup{(V5.\arabic*)},leftmargin=5em]
\item prove fixed-bank norm convergence and a positive centered core floor
      \(d\);
\item obtain a uniform entering-tail floor \(f\) from a local Casimir,
      representation, or block estimate on the \emph{whitened Schur form};
\item acquire the same-history mixed block and prove
      \(\ell^2<df\), retaining its sign when a sharper Schur short is used;
\item if the inequality fails, return the minimizing Wilson-network
      combination and determine whether it is a physical local soft vector, a
      wrapping direction, a lower-frame collapse, or a same-generator defect;
\item only on the positive branch apply the V4 full-space floor, V2 orbit
      integration, and V1 OS limit passage.
\end{enumerate}
The finite positivity-improving gap and the V4 finite generalized eigenvalue
are not substitutes for uniform control of the entering-head tail.
\end{openproblem}

\begin{firewall}[End-of-V4 claim boundary]
V4 closes the exact vacuum short and finite low-Casimir head card.  It does not
populate a cofinal collective-head triple \((d,f,\ell)\), prove the
same-generator comparison, or construct the continuum Yang--Mills state.  The
full mass gap remains open.
\end{firewall}

\section*{Version V4 append-only record}
\addcontentsline{toc}{section}{Version V4 append-only record}
The complete V3 source was retained before this continuation.  Its SHA--256
digest was
\begin{center}\scriptsize\ttfamily
2237f19ce1004b32b09c340393d6f9bd6e086b1ebf78173507b79494ceb32e91.
\end{center}
On the next iteration, retain this full file, append before its unique active
document-closing command, increment the filename to
\texttt{\_V5.tex}, and remove only the superseded V4 file from this note
series.

% =============================================================================
% END APPENDED VERSION V4 MATERIAL
% =============================================================================


% =============================================================================
% BEGIN APPENDED VERSION V5 MATERIAL
% =============================================================================

\clearpage
\part{Version V5: Casimir leakage in the exact Feshbach metric}
\label{part:V5}

\section{Purpose and point of departure}
\label{sec:V5-purpose}

The V4 collective criterion is exact but deliberately abstract: it asks for a
centered core floor \(d\), an entering-tail floor \(f\), and a mixed norm
\(\ell\) in the whitened Schur head.  This continuation derives \(f\) and
\(\ell\) from a second, inner Casimir cut while retaining the physical metric
\(M=I+\mathcal R^*\mathcal R\) and the true vacuum direction.  The resulting
certificate is finite-regulator rigorous and cofinally exact as an
implication.  Its hypotheses also identify the remaining volume-uniform
estimates without disguising them as consequences of finite-dimensional
positivity.

This is not a repetition of the V3 ultraviolet cut.  V3 eliminated the outer
Casimir complement \(Q_\Lambda\cH\).  V5 works \emph{inside} the surviving
head \(P_\Lambda\cH\), where ordinary Casimir orthogonality is distorted by
the Feshbach graph metric.  The new issue is to quantify that distortion.

% =============================================================================
\section{An inner Casimir split after the exact vacuum short}
\label{sec:V5-inner-split}
% =============================================================================

Fix one finite gauge graph and the V3 outer cut.  Suppress the subscript
\(\Lambda\) and write
\begin{equation}
 H=\kappa\mathcal C_E+V_B-E_0I,
 \qquad
 \mathcal R=C^{-1}B^*,
 \qquad
 S=PHP-BC^{-1}B^*,
 \qquad
 M=I+\mathcal R^*\mathcal R .
 \label{eq:V5-outer-data}
\end{equation}
Here \(P=P_\Lambda\), \(C=Q_\Lambda H Q_\Lambda\succ0\), and all operators
in the remainder of this section act on the finite-dimensional head
\(P\cH\), unless otherwise indicated.  Put
\begin{equation}
 \mathcal D:=BC^{-1}B^*\succeq0,
 \qquad
 U:=P(V_B-E_0I)P-\mathcal D,
 \qquad
 \boxed{S=\kappa P\mathcal C_EP+U.}
 \label{eq:V5-effective-U}
\end{equation}
The combination \(U\), unlike either \(V_B\) or \(E_0\) separately, is
invariant under adding a scalar to the raw magnetic potential.

Let
\begin{equation}
 L=\one_{[0,\lambda]}(\mathcal C_E)P,
 \qquad
 L^\sharp=P-L,
 \label{eq:V5-inner-Casimir-projections}
\end{equation}
and set
\begin{equation}
 c_-:=\norm{\mathcal C_E|_{L\cH}},
 \qquad
 c_+:=\inf\spec(\mathcal C_E|_{L^\sharp\cH}).
 \label{eq:V5-inner-Casimir-edges}
\end{equation}
Thus \(c_-\le\lambda<c_+\) whenever both spectral sectors are nonzero.
Let \(x_0\in\Ker S\) be the unique head vacuum, normalized by
\begin{equation}
 \ip{x_0}{Mx_0}=1.
 \label{eq:V5-x0-normalized}
\end{equation}
Define the \(M\)-orthogonal vacuum short
\begin{equation}
 \Pi_0^M x:=x-x_0\ip{x_0}{Mx},
 \qquad
 \cK_0:=\Pi_0^M(L\cH),
 \qquad
 \cT:=\bigl(\mathbb Cx_0+L\cH\bigr)^{\perp_M}.
 \label{eq:V5-core-tail}
\end{equation}
Then
\begin{equation}
 \{x_0\}^{\perp_M}=\cK_0\oplus_M\cT.
 \label{eq:V5-M-orthogonal-decomposition}
\end{equation}
The entering tail \(\cT\) need not lie in \(L^\sharp\cH\) in the ordinary
Hilbert norm.  Its low-Casimir leakage is the first quantity that must be
controlled.

\begin{lemma}[Metric leakage identity]
\label{lem:V5-metric-leakage}
For every \(t\in\cT\),
\begin{equation}
 \boxed{Lt=-L(M-I)t.}
 \label{eq:V5-leakage-identity}
\end{equation}
Consequently, with
\begin{equation}
 \varepsilon_{LT}:=
 \sup_{0\ne t\in\cT}
 \frac{\norm{L(M-I)t}}{\norm t},
 \label{eq:V5-epsilon-def}
\end{equation}
one has
\begin{equation}
 \norm{Lt}\le\varepsilon_{LT}\norm t,
 \qquad
 \norm{L^\sharp t}^2
 \ge(1-\varepsilon_{LT}^2)\norm t^2.
 \label{eq:V5-high-Casimir-mass}
\end{equation}
Moreover, if \(\beta=\norm{\mathcal R}\), then
\begin{equation}
 \boxed{\varepsilon_{LT}\le\norm{M-I}\le\beta^2.}
 \label{eq:V5-global-leakage-bound}
\end{equation}
\end{lemma}

\begin{proof}
For \(t\in\cT\) and every \(l\in L\cH\), the definition of \(\cT\) gives
\(\ip l{Mt}=0\).  Equivalently \(LMt=0\).  Since \(M=I+(M-I)\), this is
\eqref{eq:V5-leakage-identity}.  Orthogonality of \(L\) and \(L^\sharp\) in
the ordinary head norm yields
\[
 \norm{L^\sharp t}^2=\norm t^2-\norm{Lt}^2,
\]
which proves \eqref{eq:V5-high-Casimir-mass}.  Finally
\(M-I=\mathcal R^*\mathcal R\), hence
\(\norm{M-I}=\norm{\mathcal R}^2\); compression by \(L\) can only decrease
the operator norm.
\end{proof}

The core map may also become ill-conditioned when the physical vacuum nearly
lies in the chosen low-Casimir bank.  Remove this coordinate ambiguity by
letting
\begin{equation}
 N_L:=\Ker(\Pi_0^M|_{L\cH}),
 \qquad
 \sigma_L:=
 \inf_{\substack{l\in L\cH\cap N_L^\perp\\\norm l=1}}
 \norm{\Pi_0^M l}_M,
 \qquad
 \norm x_M^2:=\ip x{Mx}.
 \label{eq:V5-sigma-L}
\end{equation}
If \(\cK_0\ne\{0\}\), finite dimensionality gives \(\sigma_L>0\).

\begin{lemma}[Minimal retained representative]
\label{lem:V5-minimal-representative}
Every \(x\in\cK_0\) has a unique representative
\(l_x\in L\cH\cap N_L^\perp\) such that
\begin{equation}
 \Pi_0^M l_x=x,
 \qquad
 \boxed{\norm{l_x}\le\sigma_L^{-1}\norm x_M.}
 \label{eq:V5-representative-bound}
\end{equation}
\end{lemma}

\begin{proof}
The restriction of \(\Pi_0^M\) to \(L\cH\cap N_L^\perp\) is a linear
bijection onto \(\cK_0\).  The definition of its smallest singular value is
exactly \eqref{eq:V5-representative-bound}.
\end{proof}

% =============================================================================
\section{Populating the V4 tail and mixed stocks}
\label{sec:V5-populate-stocks}
% =============================================================================

Define three shift-invariant, target-native quantities
\begin{align}
 \mu_T&:=\sup_{0\ne t\in\cT}
       \frac{\ip t{Mt}}{\norm t^2},
 \label{eq:V5-mu-T}\\
 a_T&:=\max\left\{0,
       \sup_{0\ne t\in\cT}
       \frac{-\ip t{Ut}}{\norm t^2}\right\},
 \label{eq:V5-a-T}\\
 u_{LT}&:=
       \sup_{\substack{0\ne l\in L\cH\\0\ne t\in\cT}}
       \frac{|\ip l{Ut}|}{\norm l\,\norm t}.
 \label{eq:V5-u-LT}
\end{align}
They are respectively the tail metric inflation, the adverse effective
potential on the tail, and the low-to-tail effective interaction.  They can
be computed from the same finite Schur card as \(d\); more importantly, they
are the quantities for which a local connected estimate should be sought.

\begin{theorem}[Casimir population of the whitened head packet]
\label{thm:V5-population}
Assume \(\cK_0\ne\{0\}\), \(\cT\ne\{0\}\), and
\(\varepsilon_{LT}<1\).  Define the exact retained-core floor
\begin{equation}
 d:=\inf_{0\ne x\in\cK_0}
       \frac{\ip x{Sx}}{\ip x{Mx}},
 \label{eq:V5-core-d}
\end{equation}
and the explicit tail and mixed bounds
\begin{equation}
 \boxed{
 f_{\rm Cas}:=
 \frac{\kappa c_+(1-\varepsilon_{LT}^2)-a_T}{\mu_T},
 \qquad
 \ell_{\rm Cas}:=
 \frac{\kappa c_-\varepsilon_{LT}+u_{LT}}{\sigma_L}.}
 \label{eq:V5-f-ell}
\end{equation}
If
\begin{equation}
 d>0,
 \qquad f_{\rm Cas}>0,
 \qquad \ell_{\rm Cas}^2<d f_{\rm Cas},
 \label{eq:V5-subcritical-condition}
\end{equation}
then the vacuum-shortened generalized head gap obeys
\begin{equation}
 \boxed{
 \delta_{\rm head}\ge
 \delta_{\rm Cas}:=
 \frac{d+f_{\rm Cas}
 -\sqrt{(d-f_{\rm Cas})^2+4\ell_{\rm Cas}^2}}2>0.}
 \label{eq:V5-delta-Cas}
\end{equation}
\end{theorem}

\begin{proof}
Let \(t\in\cT\).  The spectral theorem, followed by
\Cref{lem:V5-metric-leakage}, gives
\begin{align*}
 \ip t{St}
 &=\kappa\ip t{\mathcal C_Et}+\ip t{Ut}\\
 &\ge \kappa c_+\norm{L^\sharp t}^2-a_T\norm t^2\\
 &\ge\bigl(\kappa c_+(1-\varepsilon_{LT}^2)-a_T\bigr)\norm t^2\\
 &\ge f_{\rm Cas}\norm t_M^2.
\end{align*}
This is the V4 tail floor.

Now let \(x\in\cK_0\), and choose the minimal representative \(l_x\) from
\Cref{lem:V5-minimal-representative}.  Since
\(x=l_x-x_0\ip{x_0}{Ml_x}\) and \(Sx_0=0\), self-adjointness gives
\begin{equation}
 \ip x{St}=\ip{l_x}{St}.
 \label{eq:V5-remove-vacuum-mixed}
\end{equation}
Because \(l_x\in L\cH\),
\begin{align*}
 |\ip{l_x}{\mathcal C_Et}|
 &=|\ip{\mathcal C_E l_x}{Lt}|\\
 &\le c_-\varepsilon_{LT}\norm{l_x}\norm t.
\end{align*}
Together with \eqref{eq:V5-u-LT},
\eqref{eq:V5-representative-bound}, and
\(\norm t\le\norm t_M\), this yields
\begin{equation}
 |\ip x{St}|
 \le\ell_{\rm Cas}\norm x_M\norm t_M.
 \label{eq:V5-mixed-bound}
\end{equation}
The isometry \(M^{1/2}:(P\cH,M)\to P\cH\) carries
\eqref{eq:V5-M-orthogonal-decomposition} to the orthogonal split used in
\Cref{thm:V4-collective-head}.  The core floor is \(d\), the tail floor is
\(f_{\rm Cas}\), and the mixed operator norm is at most
\(\ell_{\rm Cas}\).  Applying that theorem proves
\eqref{eq:V5-delta-Cas}.
\end{proof}

\begin{remark}[Degenerate sectors]
If \(\cK_0=\{0\}\), only the positive tail estimate is needed.  If
\(\cT=\{0\}\), the exact finite core card \(d\) is the whole head gap.
Thus the nonzero-sector assumptions in \Cref{thm:V5-population} merely avoid
empty infima; they do not create a spectral obstruction.
\end{remark}

\begin{corollary}[Crude global-norm specialization]
\label{cor:V5-global-specialization}
Let
\begin{equation}
 V_{\max}=\norm{V_B},
 \qquad
 \gamma=\inf\spec C>0,
 \qquad
 \beta=\frac{V_{\max}}{\gamma}.
 \label{eq:V5-global-constants}
\end{equation}
Then the data in \Cref{thm:V5-population} satisfy
\begin{align}
 \varepsilon_{LT}&\le\beta^2,
 &\mu_T&\le1+\beta^2,
 \label{eq:V5-global-epsilon-mu}\\
 a_T&\le E_0+\frac{V_{\max}^2}{\gamma},
 &u_{LT}&\le
 V_{\max}+\frac{V_{\max}^2}{\gamma}+E_0\beta^2.
 \label{eq:V5-global-a-u}
\end{align}
Consequently, when \(\beta<1\), one may use
\begin{equation}
 \boxed{
 f_{\rm glob}:=
 \frac{\kappa c_+(1-\beta^4)-E_0-V_{\max}^2/\gamma}
      {1+\beta^2},
 \quad
 \ell_{\rm glob}:=
 \frac{\kappa c_-\beta^2+V_{\max}+V_{\max}^2/\gamma
       +E_0\beta^2}{\sigma_L}.}
 \label{eq:V5-global-f-ell}
\end{equation}
If \(f_{\rm glob}>0\) and \(\ell_{\rm glob}^2<d f_{\rm glob}\), then
\eqref{eq:V5-delta-Cas} holds with the global quantities.
\end{corollary}

\begin{proof}
The first two bounds follow from \(M-I=\mathcal R^*\mathcal R\) and
\(\norm{\mathcal R}\le\norm B/\gamma\le V_{\max}/\gamma\).  Moreover
\begin{equation}
 0\preceq\mathcal D=BC^{-1}B^*
 \preceq \frac{V_{\max}^2}{\gamma}I.
 \label{eq:V5-return-global}
\end{equation}
Since \(V_B\succeq0\), the negative tail of
\(U=V_B-E_0I-\mathcal D\) is bounded by the right side of the first estimate
in \eqref{eq:V5-global-a-u}.  For \(l\in L\cH\) and \(t\in\cT\),
\(\ip l{Mt}=0\), hence
\begin{equation}
 |\ip l t|=|\ip l{(M-I)t}|
 \le\beta^2\norm l\norm t.
 \label{eq:V5-scalar-mixed-leak}
\end{equation}
The triangle inequality applied to
\(U=V_B-\mathcal D-E_0I\) now proves the second bound in
\eqref{eq:V5-global-a-u}.  Substitute these estimates into
\eqref{eq:V5-f-ell}.
\end{proof}

\begin{assessment}[What V5 has and has not populated]
At each finite regulator, every object in
\eqref{eq:V5-f-ell} is a finite matrix quantity, and
\Cref{thm:V5-population} rigorously converts it into the V4 triple
\((d,f,\ell)\).  Cofinally, however, the target-native stocks
\begin{equation}
 \boxed{d,\ \sigma_L,\ \varepsilon_{LT},\ \mu_T,\ a_T,\ u_{LT}}
 \label{eq:V5-six-stocks}
\end{equation}
must have common bounds.  Finite positivity proves none of those common
bounds by itself.
\end{assessment}

% =============================================================================
\section{Why the total magnetic norm is not the cofinal invariant}
\label{sec:V5-extensive-norm}
% =============================================================================

\begin{lemma}[Scalar-shift firewall]
\label{lem:V5-scalar-shift}
For any real \(c\), replace
\begin{equation}
 V_B\longmapsto V_B+cI,
 \qquad E_0\longmapsto E_0+c.
 \label{eq:V5-scalar-shift}
\end{equation}
Then \(H,B,C,\mathcal R,S,M,U\), and all six stocks in
\eqref{eq:V5-six-stocks} are unchanged, while
\(\norm{V_B+cI}\) can be made arbitrarily large.
\end{lemma}

\begin{proof}
The two scalar additions cancel in \(H\).  The mixed block of a scalar
operator between \(P\) and \(Q\) vanishes, so all Feshbach blocks and their
derived objects are unchanged.  In \eqref{eq:V5-effective-U}, the scalar
added to \(V_B\) cancels the scalar added to \(E_0\).  The norm assertion is
immediate from the reverse triangle inequality.
\end{proof}

\begin{counterexample}[Uniform gap with an extensive interaction norm]
\label{cth:V5-extensive-norm}
On \((\mathbb C^2)^{\otimes N}\), let
\begin{equation}
 H_N=\sum_{i=1}^N |1\rangle\langle1|_i.
 \label{eq:V5-product-H}
\end{equation}
The vacuum is unique, the spectral gap is exactly one for every \(N\), and
\(\norm{H_N}=N\).  Thus an extensive global norm can diverge while the local
mass scale remains fixed.
\end{counterexample}

\begin{proof}
The product basis diagonalizes \(H_N\); its eigenvalues are the numbers of
occupied sites.  Hence the lowest two distinct eigenvalues are zero and one,
whereas the largest is \(N\).
\end{proof}

For the Wilson magnetic potential, V3 gives
\(V_{\max}\le2\lambda_B|P(\Gamma)|\).  Therefore the global specialization
\eqref{eq:V5-global-f-ell} generally deteriorates with spatial volume even
when a uniform physical gap is present.  This is a defect of an extensive
majorant, not evidence for a soft physical state.  The target-native packet
is the correct point at which to go upstream: \(a_T\) and \(u_{LT}\) must be
estimated by connected plaquette collars or an equivalent local conditional
expectation, and \(\varepsilon_{LT}\) must be estimated from the local
Feshbach response rather than from \(\norm{\mathcal R}\) on the entire
carrier.

\begin{firewall}[No volume cancellation by notation]
Writing \(V_B-E_0I\) removes scalar shifts but does not by itself prove a
volume-uniform operator bound.  A connected-cluster, martingale, finite-range
resolvent, or comparable locality estimate is still required.  Any proof
that merely replaces \(V_{\max}\) by an unnamed constant has moved, rather
than solved, the cofinal problem.
\end{firewall}

% =============================================================================
\section{Cofinal compiler with populated head-tail inputs}
\label{sec:V5-cofinal-compiler}
% =============================================================================

\begin{theorem}[Uniform regulator gap from the V5 packet]
\label{thm:V5-uniform-compiler}
Let \(H_j\succeq0\) be a response-matched sequence of finite Hamiltonian gauge
regulators, each with a unique vacuum.  Suppose the outer complements obey
\begin{equation}
 C_j\succeq\gamma_*I
 \qquad(\gamma_*>0),
 \label{eq:V5-uniform-outer-floor}
\end{equation}
and choose inner cuts as above.  Assume common constants
\begin{equation}
 d_j\ge d_*>0,
 \qquad f_{{\rm Cas},j}\ge f_*>0,
 \qquad \ell_{{\rm Cas},j}\le\ell_*,
 \qquad \ell_*^2<d_*f_*.
 \label{eq:V5-uniform-packet}
\end{equation}
Set
\begin{equation}
 \delta_*:=
 \frac{d_*+f_*-\sqrt{(d_*-f_*)^2+4\ell_*^2}}2,
 \qquad
 m_*:=\frac{\delta_*\gamma_*}{\delta_*+\gamma_*}.
 \label{eq:V5-uniform-mass}
\end{equation}
Then
\begin{equation}
 \boxed{
 \ip\psi{H_j\psi}\ge m_*\norm\psi^2
 \quad\text{for every }\psi\perp\Omega_j,
 \quad\text{uniformly in }j.}
 \label{eq:V5-uniform-full-gap}
\end{equation}
If, in addition, the V2 orbit integration and the V1 quotient-stable OS
limit hypotheses hold for the same generator and the same local observable
bank, their conclusions hold with a strictly positive common mass no smaller
than the corresponding physical-time rescaling of \(m_*\).
\end{theorem}

\begin{proof}
By \Cref{thm:V5-population}, the generalized Schur gap at regulator \(j\) is
bounded below by the smaller eigenvalue of
\(\begin{psmallmatrix}d_j&-\ell_j\\-\ell_j&f_j\end{psmallmatrix}\).
The quadratic-form bounds in \eqref{eq:V5-uniform-packet} lower this matrix
by
\(\begin{psmallmatrix}d_*&-\ell_*\\-\ell_*&f_*\end{psmallmatrix}\)
in the scalar two-variable estimate, whose smaller eigenvalue is
\(\delta_*>0\).  The V4 vacuum-corrected full-space theorem, with
\(C_j\succeq\gamma_*I\), then gives \(m_*\).  The final statement is exactly
the conditional V2 orbit and V1 quotient passage applied without changing
the generator between the Hamiltonian and OS stages.
\end{proof}

\begin{remark}[Sharper outer recombination]
The harmonic-mean constant \(m_*\) is a safe consequence of the V4 centered
decomposition.  When a uniform router norm is also available, the sharper
V3 two-variable Feshbach eigenvalue may replace it.  This changes the final
constant, not the six cofinal inputs in \eqref{eq:V5-six-stocks}.
\end{remark}

% =============================================================================
\section{V5 ledger and revised upstream attack}
\label{sec:V5-ledger}
% =============================================================================

\begin{theorem}[Integrated V5 statement]
\label{thm:V5-integrated}
Preserving V1--V4, V5 proves:
\begin{enumerate}[label=\textup{(V5.\arabic*)},leftmargin=5em]
\item the exact low-Casimir leakage identity
      \(Lt=-L(M-I)t\) on the physical entering tail;
\item an explicit conversion of local Casimir energy into the V4 tail floor
      \(f\), including the adverse Schur-return term;
\item an explicit mixed bound \(\ell\) that retains vacuum-short
      conditioning and metric leakage;
\item a finite-regulator positive head criterion
      \(\ell_{\rm Cas}^2<d f_{\rm Cas}\);
\item a uniform cofinal Hamiltonian gap theorem when the populated packet has
      common bounds;
\item scalar-shift invariance and an exact counterexample showing why a total
      interaction norm is not a physical cofinal invariant.
\end{enumerate}
\end{theorem}

\begin{proof}
Items \textup{(V5.1)--(V5.3)} are
\Cref{lem:V5-metric-leakage,lem:V5-minimal-representative,
thm:V5-population}.  Item \textup{(V5.4)} is the last part of
\Cref{thm:V5-population}.  Item \textup{(V5.5)} is
\Cref{thm:V5-uniform-compiler}.  Item \textup{(V5.6)} is
\Cref{lem:V5-scalar-shift,cth:V5-extensive-norm}.
\end{proof}

\begingroup
\small
\begin{longtable}{>{\raggedright\arraybackslash}p{0.28\textwidth}
                  >{\raggedright\arraybackslash}p{0.16\textwidth}
                  >{\raggedright\arraybackslash}p{0.48\textwidth}}
\toprule
\textbf{V5 row}&\textbf{Status}&\textbf{Advance or remaining obstruction}\\
\midrule
\endfirsthead
\toprule
\textbf{V5 row}&\textbf{Status}&\textbf{Advance or remaining obstruction}\\
\midrule
\endhead
Inner-head Casimir leakage
& \MGpass
& Exact identity and target-native leakage norm in the Feshbach metric.\\[0.35em]
Tail and mixed population
& \MGpass
& Rigorous formulas for \(f_{\rm Cas}\) and \(\ell_{\rm Cas}\) from six typed finite-card stocks.\\[0.35em]
Finite-regulator collective head
& \MGpass
& Positive whenever \(d>0\), \(f_{\rm Cas}>0\), and \(\ell_{\rm Cas}^2<df_{\rm Cas}\).\\[0.35em]
Cofinal connected-stock bounds
& \MGopen
& Need carrier-uniform \(d,\sigma_L,\varepsilon_{LT},\mu_T,a_T,u_{LT}\) on the response-matched local branch.\\[0.35em]
Global Wilson-norm substitute
& Insufficient cofinally
& The explicit \(V_{\max}\) corollary is rigorous but generally volume-extensive.\\[0.35em]
OS/Hamiltonian same-generator comparison
& \MGopen
& The V3 form-mismatch row remains unproved.\\[0.35em]
Continuum construction and full mass gap
& \MGopen
& Require the local uniform packet, nontrivial OS limit, and same-generator identification.\\
\bottomrule
\end{longtable}
\endgroup

\begin{openproblem}[V6 mandate: connected collar bounds or a typed soft mode]
\label{op:V6}
Go upstream from the volume-extensive \(V_{\max}\) bound.  On one declared
bounded-support Wilson-network bank and its response-matched enlargement:
\begin{enumerate}[label=\textup{(V6.\arabic*)},leftmargin=5em]
\item express \(a_T\) and \(u_{LT}\) as connected plaquette-collar blocks of
      \(U=P(V_B-E_0)P-BC^{-1}B^*\), with all disconnected scalar pieces
      cancelled before taking norms;
\item prove a finite-range or weighted-resolvent decay estimate for
      \(C^{-1}B^*\) strong enough to bound \(\varepsilon_{LT}\) locally;
\item prove that \(\sigma_L\) stays away from zero, or enlarge the retained
      core by the nearly-vacuum singular vector before recomputing the card;
\item transport a positive fixed-core floor \(d\) along the same-history
      trajectory and test the strict determinant margin
      \(df_{\rm Cas}-\ell_{\rm Cas}^2\);
\item if any bound fails, return the associated minimizing vector and classify
      it as a physical local soft state, a wrapping/volume state, a
      lower-frame collapse, or a generator mismatch;
\item only after common positive bounds are obtained invoke
      \Cref{thm:V5-uniform-compiler} and then complete the OS continuum
      construction.
\end{enumerate}
The decisive next theorem is therefore local resolvent/cluster control, not
another finite Perron gap or an uncentered global operator-norm estimate.
\end{openproblem}

\begin{firewall}[End-of-V5 claim boundary]
V5 closes the algebraic passage from an inner Casimir split to the exact V4
collective head certificate.  It does not prove that the six target-native
stocks are uniform in volume and cutoff, identify the Hamiltonian generator
with the OS transfer generator, or construct the nontrivial continuum
Yang--Mills measure.  Consequently the full Yang--Mills mass gap is not yet
proved.
\end{firewall}

\section*{Version V5 append-only record}
\addcontentsline{toc}{section}{Version V5 append-only record}
The complete V4 source was retained before this continuation.  Its SHA--256
digest was
\begin{center}\scriptsize\ttfamily
1fc7a89e3abfb67c17a18853c464a414c7ac8b2ec2f8b448c2348179ee6af9f9.
\end{center}
On the next iteration, retain this full file, append before its unique active
document-closing command, increment the filename to
\texttt{\_V6.tex}, and remove only the superseded V5 file from this note
series.

% =============================================================================
% END APPENDED VERSION V5 MATERIAL
% =============================================================================

% =============================================================================
% BEGIN APPENDED VERSION V6 MATERIAL
% =============================================================================

\clearpage
\part{Version V6: nested Casimir shielding and the ground-state upstream route}
\label{part:V6}

\section{Direction and nonduplication}
\label{sec:V6-direction}

V5 showed that a global magnetic norm is the wrong cofinal quantity.  Before
postulating a new locality estimate, this continuation uses two structures
already present in the supplied programme but not yet joined to the V5
Hamiltonian card.

First, Wilson plaquette multiplication has finite Peter--Weyl bandwidth.
With a nested inner/outer Casimir cut, this gives an exact shielding identity:
the outer Feshbach router vanishes on the inner Casimir sector.  This is
stronger than a Combes--Thomas estimate for the particular leakage stock
\(\varepsilon_{LT}\).

Second, the supplied reduced-fibre V35 manuscript already contains a
weighted-conjugation estimate for a comparison one-form Witten inverse and a
Woodbury formula for its curvature-defect return.  It explicitly leaves the
uniform physical return contraction and the first response-matched physical
card open.  V6 does not reprove or silently assume those rows.  Instead, it
derives the exact finite Hamiltonian ground-state transform and states the
same-law incidence required before that Witten machinery can control the
Hamiltonian mass gap.

% =============================================================================
\section{\(\mathrm{SU}(3)\) Peter--Weyl bandwidth of a Wilson plaquette}
\label{sec:V6-Peter-Weyl-bandwidth}
% =============================================================================

Use the conventional quadratic Casimir normalization
\begin{equation}
 C_2(p,q)
 =\frac{p^2+q^2+pq+3p+3q}{3},
 \qquad (p,q)\in\mathbb N_0^2,
 \label{eq:V6-SU3-Casimir}
\end{equation}
so \(C_2(1,0)=C_2(0,1)=4/3\).  For \(\lambda\ge0\), define
\begin{equation}
 \Delta_{\rm f}(\lambda)
 :=\sqrt{3\lambda}+\frac43,
 \qquad
 \Phi(\lambda)
 :=\lambda+4\Delta_{\rm f}(\lambda).
 \label{eq:V6-Phi}
\end{equation}

\begin{lemma}[One-link fundamental Casimir jump]
\label{lem:V6-one-link-jump}
If \(C_2(p,q)\le\lambda\), every irreducible summand produced by tensoring
\((p,q)\) with the fundamental \((1,0)\) or antifundamental \((0,1)\) has
Casimir at most
\begin{equation}
 \boxed{C_2(p,q)+\Delta_{\rm f}(\lambda).}
 \label{eq:V6-one-link-bound}
\end{equation}
\end{lemma}

\begin{proof}
The Clebsch--Gordan rules are
\begin{align}
 (p,q)\otimes(1,0)
 &=(p+1,q)\oplus(p-1,q+1)\oplus(p,q-1),
 \label{eq:V6-CG-fund}\\
 (p,q)\otimes(0,1)
 &=(p,q+1)\oplus(p+1,q-1)\oplus(p-1,q),
 \label{eq:V6-CG-antifund}
\end{align}
with summands having a negative label omitted.  Direct substitution into
\eqref{eq:V6-SU3-Casimir} gives, for the first line,
\begin{align}
 C_2(p+1,q)-C_2(p,q)&=\frac{2p+q+4}{3},\nonumber\\
 C_2(p-1,q+1)-C_2(p,q)&=\frac{-p+q+1}{3},\nonumber\\
 C_2(p,q-1)-C_2(p,q)&=-\frac{p+2q+2}{3}.
 \label{eq:V6-fund-differences}
\end{align}
The antifundamental differences are obtained by interchanging \(p\) and
\(q\).  Since
\(
 C_2(p,q)\ge p^2/3
\)
and
\(
 C_2(p,q)\ge q^2/3
\),
one has \(p,q\le\sqrt{3\lambda}\).  Every positive difference is therefore
at most
\[
 \frac{3\sqrt{3\lambda}+4}{3}
 =\Delta_{\rm f}(\lambda).
\]
\end{proof}

\begin{theorem}[Finite Casimir propagation of the Wilson potential]
\label{thm:V6-finite-Casimir-propagation}
Let \(\Gamma\) be a finite hypercubic gauge graph whose plaquette boundaries
contain four distinct oriented links, and let
\[
 P_\lambda=\one_{[0,\lambda]}(\mathcal C_E).
\]
For the \(\mathrm{SU}(3)\) Wilson magnetic multiplication operator
\[
 V_B=\lambda_B\sum_p
 \left(1-\frac13\operatorname{ReTr}U_p\right),
\]
one has the exact range inclusion
\begin{equation}
 \boxed{
 V_B\Ran P_\lambda
 \subseteq
 \Ran P_{\Phi(\lambda)}.}
 \label{eq:V6-Wilson-bandwidth}
\end{equation}
The inclusion remains valid after restriction to the gauge-invariant
Hilbert space and is independent of the number of plaquettes.
\end{theorem}

\begin{proof}
In the Peter--Weyl decomposition of one link, multiplication by a matrix
coefficient of \(U\) or \(U^*\) tensors the link representation with
\((1,0)\) or \((0,1)\).  The plaquette trace is a contraction of four such
matrix coefficients, one on each boundary link.  If a Peter--Weyl component
has total electric Casimir at most \(\lambda\), each of its four boundary
link Casimirs is at most \(\lambda\).  By
\Cref{lem:V6-one-link-jump}, every resulting component has total Casimir at
most
\[
 \lambda+4\Delta_{\rm f}(\lambda)=\Phi(\lambda).
\]
The real part adds the conjugate word and obeys the same bound.  The constant
part of each Wilson plaquette leaves every representation label unchanged.
Taking a linear combination over plaquettes does not enlarge the spectral
range.  Finally, both \(V_B\) and \(\mathcal C_E\) preserve the
gauge-invariant subspace, so restriction preserves the inclusion.
\end{proof}

\begin{remark}[Degenerate plaquette words]
If a quotient graph has a plaquette word that traverses a link more than
once, iterate \Cref{lem:V6-one-link-jump} along that word.  This gives a
slightly larger explicit function \(\Phi_r\), depending only on the maximal
word length \(r\), and all shielding conclusions below remain unchanged.
\end{remark}

% =============================================================================
\section{Exact nested-cut Feshbach shielding}
\label{sec:V6-nested-shield}
% =============================================================================

Return to the V5 notation, but distinguish the inner and outer projections:
\begin{equation}
 L=P_\lambda,
 \qquad
 P=P_\Lambda,
 \qquad
 Q=I-P_\Lambda,
 \qquad
 \Lambda\ge\Phi(\lambda).
 \label{eq:V6-nested-cuts}
\end{equation}
The outer blocks are
\[
 B=PV_BQ,\qquad C=QHQ,\qquad
 \mathcal R=C^{-1}B^*,\qquad
 \mathcal D=BC^{-1}B^*,\qquad
 M=I+\mathcal R^*\mathcal R.
\]

\begin{theorem}[Router shielding of the inner Casimir sector]
\label{thm:V6-router-shield}
Under \eqref{eq:V6-nested-cuts}, whenever \(C\succ0\),
\begin{equation}
 \boxed{
 B^*L=0,\qquad
 \mathcal RL=0,\qquad
 \mathcal DL=0,\qquad
 (M-I)L=L(M-I)=0.}
 \label{eq:V6-shield-identities}
\end{equation}
For the V5 entering tail
\(
 \cT=(\mathbb Cx_0+L\cH)^{\perp_M}
\),
one consequently has
\begin{equation}
 \boxed{
 L\cT=0,\qquad
 \varepsilon_{LT}=0.}
 \label{eq:V6-zero-leakage}
\end{equation}
\end{theorem}

\begin{proof}
\Cref{thm:V6-finite-Casimir-propagation} and
\(\Lambda\ge\Phi(\lambda)\) imply
\[
 QV_BL=0.
\]
Since \(B^*=QV_BP\) and \(PL=L\), this is \(B^*L=0\).  Left multiplication
by \(C^{-1}\), \(B C^{-1}\), and \(\mathcal R^*\), respectively, gives the
next three identities.  Self-adjointness of \(M-I\) gives the left-sided
identity as well.

For \(l\in L\cH\) and \(t\in\cT\),
\[
 0=\ip l{Mt}=\ip{Ml}{t}=\ip l t,
\]
because \(Ml=l\).  Hence \(Lt=0\).  The definition of
\(\varepsilon_{LT}\) and \(L(M-I)=0\) then give zero leakage.
\end{proof}

\begin{corollary}[Exact vacuum-short conditioning under shielding]
\label{cor:V6-sigma-exact}
Let
\begin{equation}
 a_L:=\norm{Lx_0},
 \qquad 0\le a_L\le1,
 \label{eq:V6-a-L}
\end{equation}
where \(\norm{x_0}_M=1\).  If the centered retained sector
\(\cK_0=\Pi_0^M L\cH\) is nonzero, the V5 conditioning number is
\begin{equation}
 \boxed{
 \sigma_L=
 \begin{cases}
 \sqrt{1-a_L^2},&a_L<1,\\
 1,&a_L=1.
 \end{cases}}
 \label{eq:V6-sigma-formula}
\end{equation}
\end{corollary}

\begin{proof}
For \(l\in L\cH\), shielding gives \(Ml=l\), and the Pythagorean identity
for the \(M\)-orthogonal vacuum projection gives
\begin{equation}
 \norm{\Pi_0^M l}_M^2
 =\norm l^2-|\ip{x_0}{l}|^2.
 \label{eq:V6-projection-norm}
\end{equation}
If \(a_L<1\), the projection has no kernel on \(L\cH\), and the largest
possible value of \(|\ip{x_0}{l}|\) over unit \(l\in L\cH\) is
\(\norm{Lx_0}=a_L\).  This proves the first case.  If \(a_L=1\), then
\(\norm{x_0}=1\), \(Lx_0=x_0\), and \(\mathcal Rx_0=0\).
The kernel is \(\mathbb Cx_0\); on its ordinary orthogonal complement the
second term in \eqref{eq:V6-projection-norm} vanishes, proving the second
case.
\end{proof}

Define the finite representation collar
\begin{equation}
 K_{\lambda,\Lambda}
 :=\one_{(\lambda,\Phi(\lambda)]}(\mathcal C_E)P_\Lambda.
 \label{eq:V6-representation-collar}
\end{equation}

\begin{theorem}[The V5 mixed stock lands in a finite Casimir collar]
\label{thm:V6-mixed-collar}
Under nested shielding, let \(P_{\cT}^{(2)}\) denote the ordinary Hilbert
projection onto \(\cT\).  Then
\begin{equation}
 \boxed{
 u_{LT}
 =\norm{P_{\cT}^{(2)}V_BL}_{\op}
 =\norm{P_{\cT}^{(2)}K_{\lambda,\Lambda}V_BL}_{\op}.}
 \label{eq:V6-u-collar}
\end{equation}
Thus the outer Schur return \(\mathcal D\), the scalar \(E_0\), and all
Casimir sectors above \(\Phi(\lambda)\) make exactly zero contribution to
the V5 mixed stock.
\end{theorem}

\begin{proof}
For \(l\in L\cH\) and \(t\in\cT\), \Cref{thm:V6-router-shield} gives
\[
 \ip l t=0,\qquad \mathcal Dl=0.
\]
Self-adjointness therefore yields
\[
 \ip l{Ut}
 =\ip l{(V_B-E_0I-\mathcal D)t}
 =\ip l{V_Bt}.
\]
Taking the supremum over unit \(l,t\) gives the first operator norm in
\eqref{eq:V6-u-collar}.  By
\Cref{thm:V6-finite-Casimir-propagation}, \(V_BL\) lies in
\(\Ran P_{\Phi(\lambda)}\).  Its \(L\)-component is orthogonal to \(\cT\),
leaving precisely the displayed collar.
\end{proof}

\begin{corollary}[Shielded V5 head card]
\label{cor:V6-shielded-card}
With \(d,a_T,\mu_T\) as in V5 and \(c_+\) the first inner-complement Casimir,
put
\begin{equation}
 f_{\rm sh}:=\frac{\kappa c_+-a_T}{\mu_T},
 \qquad
 \ell_{\rm sh}:=
 \frac{\norm{P_{\cT}^{(2)}K_{\lambda,\Lambda}V_BL}}{\sigma_L}.
 \label{eq:V6-shielded-f-ell}
\end{equation}
If
\begin{equation}
 d>0,\qquad f_{\rm sh}>0,\qquad
 \ell_{\rm sh}^2<d f_{\rm sh},
 \label{eq:V6-shielded-determinant}
\end{equation}
then
\begin{equation}
 \boxed{
 \delta_{\rm head}\ge
 \frac{d+f_{\rm sh}
 -\sqrt{(d-f_{\rm sh})^2+4\ell_{\rm sh}^2}}2>0.}
 \label{eq:V6-shielded-gap}
\end{equation}
\end{corollary}

\begin{proof}
Insert \(\varepsilon_{LT}=0\), the exact value
\eqref{eq:V6-u-collar}, and \eqref{eq:V6-sigma-formula} into
\Cref{thm:V5-population}.
\end{proof}

\begin{assessment}[Exact gain from nesting]
The difficult metric-leakage estimate requested by V5 is not needed for the
Wilson Hamiltonian when the outer cut shields the inner cut:
\[
 \boxed{\varepsilon_{LT}=0.}
\]
The return term also disappears from the mixed block.  What remains is the
tail return debit \(a_T\), the tail metric inflation \(\mu_T\), the centered
low-Casimir floor \(d\), the vacuum angle \(\sigma_L\), and one finite
representation-collar block.  This is a strict reduction of the six-stock
packet, not a full volume-uniform estimate.
\end{assessment}

% =============================================================================
\section{Finite bandwidth does not solve the infrared volume problem}
\label{sec:V6-bandwidth-not-volume}
% =============================================================================

\begin{counterexample}[A fixed local bandwidth with a growing vacuum-to-collar block]
\label{cth:V6-sqrt-volume}
Let \((X,\mu)\) be a probability space and choose a real mean-zero
\(v\in L^\infty(\mu)\) with \(\norm v_2=1\).  On
\(L^2(X^N,\mu^{\otimes N})\), let \(v_i\) be \(v\) in coordinate \(i\), put
\[
 V_N=\sum_{i=1}^N v_i,
\]
let \(L_N\) project onto the constant function, and let \(K_N\) project onto
\(\Span\{v_1,\ldots,v_N\}\).  Then every summand has the same one-coordinate
bandwidth, but
\begin{equation}
 \boxed{\norm{K_NV_NL_N}=\sqrt N.}
 \label{eq:V6-sqrt-volume}
\end{equation}
\end{counterexample}

\begin{proof}
Independence and \(\int v\,\dd\mu=0\) make the vectors \(v_i\) orthonormal.
Since
\[
 K_NV_N1=\sum_{i=1}^N v_i,
\]
its norm is \(\sqrt N\).  The domain of \(L_N\) is one-dimensional, so this
is the operator norm.
\end{proof}

\begin{counterexample}[Local finite propagation with an infrared collapsing gap]
\label{cth:V6-cycle-soft-mode}
Let \(H_N\) be the nearest-neighbour graph Laplacian on the periodic cycle
\(\mathbb Z/N\mathbb Z\).  Its propagation and row sum are bounded
independently of \(N\), its vacuum is the unique constant vector, but
\begin{equation}
 \boxed{
 \gap(H_N)=4\sin^2\!\left(\frac{\pi}{N}\right)\longrightarrow0.}
 \label{eq:V6-cycle-gap}
\end{equation}
\end{counterexample}

\begin{proof}
The Fourier vectors diagonalize the cycle Laplacian with eigenvalues
\(
 4\sin^2(\pi k/N)
\).
The first nonzero values occur at \(k=1,N-1\).
\end{proof}

The first counterexample shows why exact Casimir propagation alone does not
bound the shielded mixed collar uniformly; disconnected local creations can
accumulate.  The second shows why locality alone does not supply a mass:
long-wavelength collective modes must be excluded by a genuine infrared
inequality.  In the Yang--Mills card, the full low-total-Casimir sector also
grows by translating bounded excitations through the volume.  A positive
fixed observable bank therefore cannot stand in for the common centered
floor \(d\) on that growing sector.

% =============================================================================
\section{Exact ground-state transform: the volume-native upstream coordinate}
\label{sec:V6-ground-state-transform}
% =============================================================================

The preceding obstruction indicates that the raw potential block is not the
best coordinate for volume.  At a finite Hamiltonian regulator there is an
exact transformation in which the potential and ground energy cancel from
the quadratic form.

\begin{theorem}[Ground-state representation of the finite gauge Hamiltonian]
\label{thm:V6-ground-state-representation}
Let \(G\) be compact and connected, let
\(\mathcal X=G^{E(\Gamma)}\) with product Haar measure \(\mu\), and let
\[
 H=\kappa\sum_e(-\Delta_e)+V_B-E_0
\]
act on the gauge-invariant subspace.  Suppose its normalized ground state
\(\psi_0\) is unique and strictly positive.  Put
\begin{equation}
 \dd\nu:=\psi_0^2\,\dd\mu.
 \label{eq:V6-ground-law}
\end{equation}
Then, for every smooth gauge-invariant \(f\),
\begin{equation}
 \boxed{
 \ip{\psi_0 f}{H\psi_0 f}_{L^2(\mu)}
 =\kappa\int_{\mathcal X}|\nabla f|^2\,\dd\nu.}
 \label{eq:V6-ground-state-identity}
\end{equation}
Multiplication by \(\psi_0\) is a unitary
\[
 \mathcal U:L^2(\nu)\longrightarrow L^2(\mu),
 \qquad \mathcal Uf=\psi_0f,
\]
and, in quadratic-form sense,
\begin{equation}
 \boxed{
 \mathcal U^{-1}H\mathcal U
 =L_\nu:=\kappa d_\nu^*d.}
 \label{eq:V6-Doob-generator}
\end{equation}
In particular,
\begin{equation}
 \boxed{
 \gap(H)
 =\kappa
 \inf_{\substack{f\ne0\\\int f\,\dd\nu=0}}
 \frac{\int|\nabla f|^2\,\dd\nu}
      {\int|f|^2\,\dd\nu}.}
 \label{eq:V6-gap-Poincare}
\end{equation}
The infimum is over the gauge-invariant form domain.
\end{theorem}

\begin{proof}
The ground-state equation is
\[
 \left(\kappa\sum_e(-\Delta_e)+V_B-E_0\right)\psi_0=0.
\]
Expand the gradient of \(\psi_0f\), integrate by parts on the compact product
group, and use this equation.  The terms containing
\(|f|^2|\nabla\psi_0|^2\), \(V_B-E_0\), and the mixed derivative cancel,
leaving
\[
 \ip{\psi_0 f}{H\psi_0 f}
 =\kappa\int\psi_0^2|\nabla f|^2\,\dd\mu.
\]
This proves \eqref{eq:V6-ground-state-identity} first on the smooth core and
then on the closed form domain.  Normalization of \(\psi_0\) makes
\(\mathcal U\) unitary.  The form identity proves
\eqref{eq:V6-Doob-generator}.  Finally,
\(\psi_0 f\perp\psi_0\) is equivalent to
\(\int f\,\dd\nu=0\), so the Rayleigh--Ritz characterization of the first
positive eigenvalue gives \eqref{eq:V6-gap-Poincare}.
\end{proof}

\begin{corollary}[Exact one-form incidence for the mass gap]
\label{cor:V6-one-form-incidence}
Let
\[
 D=\sqrt\kappa\,d,\qquad
 L_0=D^*D\quad\hbox{on }L_0^2(\nu),\qquad
 L_1=DD^*\quad\hbox{on }\overline{\Ran D}.
\]
Then the nonzero spectra of \(L_0\) and \(L_1\) agree, including
multiplicity.  Consequently,
\begin{equation}
 \boxed{
 L_1\succeq\rho I\text{ on }\overline{\Ran D}
 \quad\Longrightarrow\quad
 \gap(H)=\gap(L_0)\ge\rho.}
 \label{eq:V6-one-form-to-gap}
\end{equation}
\end{corollary}

\begin{proof}
The polar decomposition of the closed operator \(D\) unitarily identifies
\(D^*D\) on \((\Ker D)^\perp\) with \(DD^*\) on
\(\overline{\Ran D}\).  Constants are exactly \(\Ker D\) on the connected
finite configuration space.  Apply \Cref{thm:V6-ground-state-representation}.
\end{proof}

\begin{assessment}[Precise bridge to the supplied Witten programme]
The exact-form Witten isometry and the comparison/return Woodbury split in
the supplied V35 reduced-fibre manuscript can control the Hamiltonian gap
only after the following incidence is proved:
\begin{equation}
 \boxed{
 \nu_{\rm fibre}
 =\psi_0^2\mu
 \quad\text{on the same regulator and protected sector,}
 \qquad
 L_{1,\rm fibre}=DD^*
 \text{ on the exact range}.}
 \label{eq:V6-same-law-incidence}
\end{equation}
With \eqref{eq:V6-same-law-incidence}, a common positive physical one-form
margin gives the Hamiltonian mass gap by
\Cref{cor:V6-one-form-incidence}.  Without it, a positive comparison Witten
card is a different-law or different-generator statement.
\end{assessment}

\begin{firewall}[Comparison inverse versus physical one-form floor]
The V35 weighted Combes--Thomas theorem proves exponential decay for the
comparison inverse under its stated positive comparison floor.  Its
Woodbury theorem represents the physical curvature-defect return exactly.
But the uniform weighted edge below one, the selected physical assembled
margin, and the same-law incidence \eqref{eq:V6-same-law-incidence} are not
proved there.  Exponential decay of the comparison inverse cannot be
substituted for positivity of \(L_1\) on all exact one-forms.
\end{firewall}

% =============================================================================
\section{A bifurcated cofinal compiler}
\label{sec:V6-bifurcated-compiler}
% =============================================================================

\begin{theorem}[Two honest regulator-gap routes]
\label{thm:V6-two-routes}
For a cofinal sequence of finite \(\mathrm{SU}(3)\) Hamiltonian gauge
regulators with unique positive vacua, either of the following independent
packets gives a uniform Hamiltonian gap.
\begin{enumerate}[label=\textup{(\Alph*)},leftmargin=4em]
\item \emph{Shielded Feshbach packet.}  Choose nested cuts satisfying
      \(\Lambda_j\ge\Phi(\lambda_j)\).  Assume common bounds
      \[
      d_j\ge d_*>0,\quad
      f_{{\rm sh},j}\ge f_*>0,\quad
      \ell_{{\rm sh},j}\le\ell_*,
      \quad \ell_*^2<d_*f_*,
      \quad C_j\succeq\gamma_*I.
      \]
\item \emph{Ground-law Witten packet.}  For
      \(\nu_j=\psi_{0,j}^2\mu_j\), assume the exact one-form operators obey
      \[
      L_{1,j}=D_jD_j^*\succeq\rho_*I
      \quad\text{on }\overline{\Ran D_j},
      \qquad \rho_*>0.
      \]
\end{enumerate}
Route \textup{(A)} gives the V5--V4 explicit lower bound
\begin{equation}
 m_A=
 \frac{\delta_A\gamma_*}{\delta_A+\gamma_*},
 \qquad
 \delta_A=
 \frac{d_*+f_*-\sqrt{(d_*-f_*)^2+4\ell_*^2}}2.
 \label{eq:V6-route-A-mass}
\end{equation}
Route \textup{(B)} gives
\begin{equation}
 \boxed{\gap(H_j)\ge\rho_*}
 \label{eq:V6-route-B-mass}
\end{equation}
directly.  Either route yields a continuum OS mass gap only after tight,
nontrivial same-generator OS convergence with physical time calibration.
\end{theorem}

\begin{proof}
Route \textup{(A)} is
\Cref{cor:V6-shielded-card,thm:V5-uniform-compiler}.  Route
\textup{(B)} is \Cref{cor:V6-one-form-incidence} at every \(j\).
The final sentence records the independent V1 continuum and generator
hypotheses; neither finite-regulator estimate implies them.
\end{proof}

\begin{assessment}[Which route is upstream in volume]
Nested shielding exactly solves the cutoff leakage created by the outer
Feshbach cut, and it is useful for certified finite cards.  The ground-state
identity is more native to the thermodynamic limit: it cancels the extensive
potential before a norm is taken and turns the gap into one uniform
Poincar\'e inequality.  Its price is equally precise: one must control the
actual ground-state law and the physical exact one-form Witten operator, not
only a comparison law.
\end{assessment}

% =============================================================================
\section{V6 ledger and next upstream attack}
\label{sec:V6-ledger}
% =============================================================================

\begin{theorem}[Integrated V6 statement]
\label{thm:V6-integrated}
Preserving V1--V5, V6 proves:
\begin{enumerate}[label=\textup{(V6.\arabic*)},leftmargin=5em]
\item an explicit \(\mathrm{SU}(3)\) Casimir-bandwidth function
      \(\Phi(\lambda)\) for Wilson plaquette multiplication;
\item exact annihilation of the inner sector by the outer Feshbach router
      when \(\Lambda\ge\Phi(\lambda)\);
\item zero V5 metric leakage, an exact vacuum-angle formula for
      \(\sigma_L\), and a finite representation-collar formula for \(u_{LT}\);
\item a shielded collective-head determinant criterion;
\item counterexamples separating finite bandwidth and finite propagation
      from volume-uniform infrared mass;
\item the exact ground-state representation of the finite gauge Hamiltonian;
\item equality of the Hamiltonian gap with the Poincar\'e constant of its
      ground-state law and with the exact one-form floor;
\item the precise same-law incidence needed to use the supplied reduced-fibre
      Witten comparison/return programme.
\end{enumerate}
\end{theorem}

\begin{proof}
Items \textup{(V6.1)--(V6.4)} are
\Cref{thm:V6-finite-Casimir-propagation,thm:V6-router-shield,
cor:V6-sigma-exact,thm:V6-mixed-collar,cor:V6-shielded-card}.
Item \textup{(V6.5)} is
\Cref{cth:V6-sqrt-volume,cth:V6-cycle-soft-mode}.  Items
\textup{(V6.6)--(V6.7)} are
\Cref{thm:V6-ground-state-representation,cor:V6-one-form-incidence}.
Item \textup{(V6.8)} is \eqref{eq:V6-same-law-incidence}.
\end{proof}

\begingroup
\small
\begin{longtable}{>{\raggedright\arraybackslash}p{0.28\textwidth}
                  >{\raggedright\arraybackslash}p{0.16\textwidth}
                  >{\raggedright\arraybackslash}p{0.48\textwidth}}
\toprule
\textbf{V6 row}&\textbf{Status}&\textbf{Advance or remaining obstruction}\\
\midrule
\endfirsthead
\toprule
\textbf{V6 row}&\textbf{Status}&\textbf{Advance or remaining obstruction}\\
\midrule
\endhead
Wilson Casimir bandwidth
& \MGpass
& Explicit representation-theoretic range bound independent of volume.\\[0.35em]
Nested-cut router leakage
& \MGpass
& Exactly zero on the inner Casimir sector; no resolvent-decay hypothesis is needed for this stock.\\[0.35em]
Shielded finite head card
& \MGpass
& The mixed block is confined to one finite representation collar.\\[0.35em]
Volume-uniform shielded packet
& \MGopen
& The collar norm, tail return, metric inflation, and growing low-sector floor still require infrared control.\\[0.35em]
Ground-state/Poincar\'e bridge
& \MGpass
& Exact at every finite Hamiltonian regulator and automatically scalar-shift invariant.\\[0.35em]
Physical one-form Witten floor
& \MGopen
& Need a common positive floor on all exact one-forms of the actual ground-state law.\\[0.35em]
Reduced-fibre same-law incidence
& \MGopen
& The V35 comparison/return law has not been identified with \(\psi_0^2\mu\) for the same Hamiltonian generator.\\[0.35em]
OS continuum and full mass gap
& \MGopen
& Need same-generator convergence, nontrivial continuum reconstruction, and one of the two uniform packets.\\
\bottomrule
\end{longtable}
\endgroup

\begin{openproblem}[V7 mandate: actual-ground-law Witten margin]
\label{op:V7}
Use the ground-state coordinate rather than an extensive magnetic norm.
\begin{enumerate}[label=\textup{(V7.\arabic*)},leftmargin=5em]
\item construct or characterize the finite-regulator ground-state density
      \(\psi_{0,j}^2\) on the protected gauge sector and match its physical
      time normalization to the OS transfer Hamiltonian;
\item compute the exact one-form operator \(L_{1,j}=D_jD_j^*\) for that law
      and compare it, on the same Hilbert bundle, with the V35
      comparison-minus-defect decomposition;
\item prove the same-law identity \eqref{eq:V6-same-law-incidence}, or return
      the explicit Radon--Nikodym/generator defect that prevents it;
\item populate a uniform source-cyclic defect edge or a low-island
      replacement on the entire exact one-form range, not merely one finite
      momentum card;
\item prove \(L_{1,j}\succeq\rho_*I\) uniformly, or return normalized exact
      one-forms whose Rayleigh quotients vanish and classify their spatial
      momentum, representation content, and carrier escape;
\item on the positive branch, combine
      \Cref{cor:V6-one-form-incidence} with the V1 quotient-stable continuum
      passage and verify nontriviality of the limiting local algebra.
\end{enumerate}
If the same-law comparison becomes circular because it presupposes the
ground-state Poincar\'e inequality, return upstream to a direct conditional
variance or martingale decomposition of \(\nu_j=\psi_{0,j}^2\mu_j\).
\end{openproblem}

\begin{firewall}[End-of-V6 claim boundary]
V6 exactly removes the outer-router leakage on a nested Wilson Casimir cut
and exactly rewrites the finite Hamiltonian mass gap as a Poincar\'e/Witten
floor of its ground-state law.  It does not prove a volume-uniform infrared
floor for that law, the same-law incidence with the supplied reduced-fibre
comparison, or a nontrivial same-generator OS continuum limit.  The full
Yang--Mills mass gap therefore remains open.
\end{firewall}

\section*{Version V6 append-only record}
\addcontentsline{toc}{section}{Version V6 append-only record}
The complete V5 source was retained before this continuation.  Its SHA--256
digest was
\begin{center}\scriptsize\ttfamily
15e7e62749649efe78d6dd1baa9664f6624eee239e9b17a4e0339b101c9ca198.
\end{center}
On the next iteration, retain this full file, append before its unique active
document-closing command, increment the filename to
\texttt{\_V7.tex}, and remove only the superseded V6 file from this note
series.

% =============================================================================
% END APPENDED VERSION V6 MATERIAL
% =============================================================================

% =============================================================================
% BEGIN APPENDED VERSION V7 MATERIAL
% =============================================================================

\clearpage
\part{Version V7: Perron-weighted ground law and a martingale gap compiler}
\label{part:V7}

\section{Purpose}
\label{sec:V7-purpose}

V6 identified the volume-native statement: the gap of the finite gauge
Hamiltonian is the Poincar\'e constant of its actual ground-state law
\(\pi=\psi_0^2\mu\).  The first task is therefore not to assign a convenient
reduced-fibre law to that Hamiltonian, but to recover the physical law from
the same transfer kernel and measure the defect of any proposed comparison
law.

This continuation does three things.  It derives the physical one-slice law
from the Perron eigenfunction of the Wilson transfer kernel; it gives a sharp
bounded-density comparison that transports a Witten/Poincar\'e floor between
two laws; and it supplies a direct martingale-collar compiler that avoids
assuming a global Poincar\'e inequality in order to prove one.

% =============================================================================
\section{The physical law of a positive transfer kernel}
\label{sec:V7-transfer-ground-law}
% =============================================================================

Retain the V2 setting: \((X,\mathfrak m)\) is a compact probability space,
\(\mathcal K\) has a continuous symmetric strictly positive kernel \(k\),
and
\begin{equation}
 \mathcal Kh=\lambda h,
 \qquad h>0,
 \qquad \int h^2\,\dd\mathfrak m=1.
 \label{eq:V7-Perron-data}
\end{equation}
Put \(T=\lambda^{-1}\mathcal K\) and
\(\dd\pi=h^2\dd\mathfrak m\).

\begin{theorem}[Perron-weighted physical one-slice law]
\label{thm:V7-physical-law}
The map
\begin{equation}
 \mathcal U_h:L^2(\pi)\longrightarrow L^2(\mathfrak m),
 \qquad
 \mathcal U_hf=hf,
 \label{eq:V7-Uh}
\end{equation}
is unitary and
\begin{equation}
 \boxed{
 P:=\mathcal U_h^{-1}T\mathcal U_h,
 \qquad
 (Pf)(x)=\frac{1}{\lambda h(x)}
 \int k(x,y)h(y)f(y)\,\dd\mathfrak m(y).}
 \label{eq:V7-Doob-same-transfer}
\end{equation}
The probability measure
\begin{equation}
 \boxed{\dd\pi=h^2\dd\mathfrak m}
 \label{eq:V7-pi-physical}
\end{equation}
is stationary and reversible for \(P\).

Moreover, the Perron-weighted two-slice probability law
\begin{equation}
 \boxed{
 \dd\mathbb J(x,y)
 :=\frac{h(x)k(x,y)h(y)}{\lambda}\,
   \dd\mathfrak m(x)\dd\mathfrak m(y)}
 \label{eq:V7-two-slice-law}
\end{equation}
has both marginals equal to \(\pi\).  Thus the physical equal-time law is
the marginal of \(\mathbb J\), not in general the marginal of the
unweighted slab density \(k(x,y)\).
\end{theorem}

\begin{proof}
The norm identity
\[
 \norm{\mathcal U_hf}_{L^2(\mathfrak m)}^2
 =\int |f|^2h^2\,\dd\mathfrak m
 =\norm f_{L^2(\pi)}^2
\]
proves unitarity; strict positivity and compactness make its range all of
\(L^2(\mathfrak m)\).  Direct conjugation gives
\eqref{eq:V7-Doob-same-transfer}.  The transition density relative to
\(\pi\) is
\[
 r(x,y)=\frac{k(x,y)}{\lambda h(x)h(y)},
\]
which is symmetric.  The Perron equation gives
\(\int r(x,y)\,\dd\pi(y)=1\), proving stationarity and reversibility.
Finally,
\[
 \int\dd\mathbb J(x,y)
 =\frac1\lambda\int h(x)(\mathcal Kh)(x)\,\dd\mathfrak m(x)=1.
\]
Integrating \eqref{eq:V7-two-slice-law} in \(y\) gives
\(\lambda^{-1}h(x)(\mathcal Kh)(x)\dd\mathfrak m(x)
=h(x)^2\dd\mathfrak m(x)\); symmetry gives the other marginal.
\end{proof}

\begin{corollary}[The transfer Hamiltonian and ground-law generator are the same operator]
\label{cor:V7-transfer-generator}
Assume \(0<T\preceq I\) and let
\[
 H_T=-\tau^{-1}\log T
\]
on \(L^2(\mathfrak m)\).  Then
\begin{equation}
 \boxed{
 \mathcal U_h^{-1}H_T\mathcal U_h
 =-\tau^{-1}\log P,}
 \label{eq:V7-log-conjugation}
\end{equation}
and both operators have the same spectral gap.  If \(H_T\) is the
differential Hamiltonian of V6, then \(h=\psi_0\) and the law in
\eqref{eq:V7-pi-physical} is exactly the ground-state law in
\eqref{eq:V6-ground-law}.
\end{corollary}

\begin{proof}
Unitary functional calculus applied to
\(P=\mathcal U_h^{-1}T\mathcal U_h\) gives
\eqref{eq:V7-log-conjugation}.  Unitary conjugation preserves spectrum.
The last assertion follows because the unique zero eigenvector of \(H_T\)
is the Perron eigenvector of \(T\).
\end{proof}

\begin{counterexample}[The unweighted slab marginal is not the ground law]
\label{cth:V7-unweighted-mismatch}
Let \(X=\{1,2\}\) with uniform \(\mathfrak m\), and let
\[
 k=\begin{pmatrix}2&1\\1&1\end{pmatrix}.
\]
The unweighted normalized two-slice density proportional to \(k\) has
one-slice marginal
\begin{equation}
 \nu_{\rm unw}=\left(\frac35,\frac25\right).
 \label{eq:V7-unweighted-marginal}
\end{equation}
The Perron vector has ratio
\(
 h_1/h_2=(1+\sqrt5)/2
\), so the physical law is
\begin{equation}
 \pi=
 \frac{1}{1+\varphi^2}(\varphi^2,1),
 \qquad \varphi=\frac{1+\sqrt5}{2},
 \label{eq:V7-two-point-pi}
\end{equation}
which differs from \(\nu_{\rm unw}\).
\end{counterexample}

\begin{proof}
The column sums of \(k\) are \(3\) and \(2\), proving
\eqref{eq:V7-unweighted-marginal}.  The maximal eigenvalue of the displayed
matrix is \((3+\sqrt5)/2\), and its positive eigenvector has ratio
\(\varphi\).  Squaring and normalizing gives
\eqref{eq:V7-two-point-pi}; the first coordinate is approximately \(0.724\),
not \(0.6\).
\end{proof}

\begin{firewall}[Finite one-slice law selection]
A normalized Wilson slab, a conditional shell law, and the stationary
ground-state law are generally distinct.  Positivity of all three does not
identify them.  The Perron factors \(h(x)h(y)\) in
\eqref{eq:V7-two-slice-law}, or an exact equivalent construction, are
required for same-law incidence.
\end{firewall}

% =============================================================================
\section{An exact law-mismatch coordinate}
\label{sec:V7-law-mismatch}
% =============================================================================

Let \(\pi\) be the physical ground law and let \(\nu\) be a proposed
reduced-fibre or comparison law on the same compact configuration space.
Assume both have positive smooth densities relative to the same Riemannian
reference measure \(\mu\).  Define
\begin{equation}
 \theta:=\log\frac{\dd\nu}{\dd\pi},
 \qquad
 \operatorname{osc}\theta:=\sup\theta-\inf\theta.
 \label{eq:V7-theta}
\end{equation}

\begin{theorem}[Bounded-density transport of a Poincar\'e gap]
\label{thm:V7-density-transport}
Suppose
\begin{equation}
 \rho_\nu\Var_\nu(f)
 \le\kappa\int|\nabla f|^2\,\dd\nu
 \label{eq:V7-nu-Poincare}
\end{equation}
for every smooth gauge-invariant \(f\), and let
\[
 r=\frac{\dd\pi}{\dd\nu},
 \qquad
 0<m\le r\le M<\infty.
\]
Then
\begin{equation}
 \boxed{
 \rho_\nu\frac{m}{M}\Var_\pi(f)
 \le\kappa\int|\nabla f|^2\,\dd\pi,}
 \label{eq:V7-pi-Poincare}
\end{equation}
and therefore
\begin{equation}
 \boxed{
 \gap(H_\pi)\ge
 \rho_\nu\frac{m}{M}
 =\rho_\nu e^{-\operatorname{osc}\theta}.}
 \label{eq:V7-density-gap}
\end{equation}
\end{theorem}

\begin{proof}
Using the variational form of variance,
\[
 \Var_\pi(f)
 =\inf_c\int|f-c|^2r\,\dd\nu
 \le M\Var_\nu(f).
\]
Also
\[
 \kappa\int|\nabla f|^2\,\dd\pi
 =\kappa\int|\nabla f|^2r\,\dd\nu
 \ge m\kappa\int|\nabla f|^2\,\dd\nu.
\]
Combine these inequalities with \eqref{eq:V7-nu-Poincare}.  Since
\(r=e^{-\theta}\), one has
\(M/m=e^{\operatorname{osc}\theta}\).  The V6 ground-state representation
turns the physical Poincar\'e floor into the Hamiltonian gap.
\end{proof}

\begin{corollary}[Cofinal comparison-law compiler]
\label{cor:V7-comparison-law-compiler}
For a regulator sequence, if
\begin{equation}
 \rho_{\nu,j}\ge\rho_*>0,
 \qquad
 \operatorname{osc}\log
 \frac{\dd\nu_j}{\dd\pi_j}\le A_*<\infty
 \label{eq:V7-uniform-law-comparison}
\end{equation}
uniformly, then
\begin{equation}
 \boxed{\gap(H_j)\ge\rho_*e^{-A_*}.}
 \label{eq:V7-uniform-comparison-gap}
\end{equation}
Exact same-law incidence is the case \(A_*=0\).
\end{corollary}

\begin{theorem}[The generator defect carried by the law mismatch]
\label{thm:V7-generator-defect}
Write
\[
 \dd\pi=Z_\pi^{-1}e^{-W_\pi}\dd\mu,
 \qquad
 \dd\nu=Z_\nu^{-1}e^{-W_\nu}\dd\mu.
\]
Constants being immaterial, \(\theta=W_\pi-W_\nu\).  The scalar weighted
generators have differential expressions
\[
 L_\omega=\kappa\left(-\Delta+\nabla W_\omega\cdot\nabla\right),
 \qquad \omega\in\{\pi,\nu\},
\]
and obey the exact identity
\begin{equation}
 \boxed{
 L_\pi-L_\nu
 =\kappa\,\nabla\theta\cdot\nabla.}
 \label{eq:V7-drift-defect}
\end{equation}
\end{theorem}

\begin{proof}
The weighted adjoint relative to
\(e^{-W_\omega}\mu\) gives the displayed formula for \(d_\omega^*d\).
Since \(W_\pi=W_\nu+\theta\) up to an additive constant, subtraction gives
\eqref{eq:V7-drift-defect}.
\end{proof}

\begin{assessment}[A typed output for the V6 same-law row]
The same-law comparison is no longer binary.  A proposed reduced-fibre law
\(\nu_j\) must return
\begin{equation}
 \boxed{
 \theta_j=\log(\dd\nu_j/\dd\pi_j),\quad
 \operatorname{osc}\theta_j,\quad
 \nabla\theta_j.}
 \label{eq:V7-law-defect-packet}
\end{equation}
The oscillation transports a known comparison gap by
\Cref{thm:V7-density-transport}; the gradient is the exact generator defect.
If the oscillation grows extensively, a global bounded-density comparison is
not cofinal even when both finite laws are smooth and positive.
\end{assessment}

% =============================================================================
\section{A noncircular martingale-collar compiler for the physical law}
\label{sec:V7-martingale-compiler}
% =============================================================================

Let the finite configuration be ordered into \(N\) link or block variables,
and let
\[
 \mathcal F_0\subset\mathcal F_1\subset\cdots\subset\mathcal F_N
\]
be the filtration that reveals them in that order.  For
\(f\in L^2(\pi)\), set
\begin{equation}
 M_i f:=\mathbb E_\pi[f\mid\mathcal F_i],
 \qquad
 D_i f:=M_i f-M_{i-1}f.
 \label{eq:V7-martingale-differences}
\end{equation}

\begin{lemma}[Exact variance ledger]
\label{lem:V7-variance-ledger}
If \(\mathcal F_N\) is the full sigma algebra, then
\begin{equation}
 \boxed{
 \Var_\pi(f)=\sum_{i=1}^N\norm{D_if}_{L^2(\pi)}^2.}
 \label{eq:V7-variance-ledger}
\end{equation}
\end{lemma}

\begin{proof}
Conditional expectations are orthogonal projections in \(L^2(\pi)\).
The martingale differences are pairwise orthogonal,
\(M_Nf=f\), and \(M_0f=\mathbb E_\pi f\).  Pythagoras applied to
\(
 f-\mathbb E_\pi f=\sum_iD_if
\)
gives the identity.
\end{proof}

\begin{theorem}[Martingale-collar Poincar\'e criterion]
\label{thm:V7-martingale-collar}
Let \(\nabla_b\) denote the gradient in block \(b\).  Suppose there are
nonnegative numbers \(w_{ib}\) such that, for every smooth
gauge-invariant \(f\),
\begin{equation}
 \norm{D_if}_{L^2(\pi)}^2
 \le\sum_{b=1}^N w_{ib}
       \int|\nabla_bf|^2\,\dd\pi
 \qquad(1\le i\le N).
 \label{eq:V7-collar-row}
\end{equation}
Define the column congestion
\begin{equation}
 C_{\rm col}:=\max_{1\le b\le N}\sum_{i=1}^Nw_{ib}.
 \label{eq:V7-column-congestion}
\end{equation}
Then
\begin{equation}
 \boxed{
 \Var_\pi(f)
 \le C_{\rm col}\int|\nabla f|^2\,\dd\pi.}
 \label{eq:V7-martingale-Poincare}
\end{equation}
For the physical Hamiltonian of V6,
\begin{equation}
 \boxed{\gap(H)\ge\kappa/C_{\rm col}.}
 \label{eq:V7-martingale-gap}
\end{equation}
In particular, a regulator-independent column-congestion bound proves a
uniform finite-volume gap without first assuming a global Poincar\'e
inequality.
\end{theorem}

\begin{proof}
Sum \eqref{eq:V7-collar-row} in \(i\) and use
\Cref{lem:V7-variance-ledger}:
\[
 \Var_\pi(f)
 \le\sum_b\left(\sum_iw_{ib}\right)
       \int|\nabla_bf|^2\,\dd\pi
 \le C_{\rm col}\int|\nabla f|^2\,\dd\pi.
\]
\Cref{thm:V6-ground-state-representation} gives
\eqref{eq:V7-martingale-gap}.
\end{proof}

\begin{corollary}[Product-law tensorization as a zero-response card]
\label{cor:V7-product-tensorization}
Let
\(\pi=\bigotimes_{i=1}^N\pi_i\), and suppose every factor obeys
\begin{equation}
 \rho_{\rm loc}\Var_{\pi_i}(g)
 \le\int|\nabla_i g|^2\,\dd\pi_i
 \label{eq:V7-local-Poincare}
\end{equation}
with the same \(\rho_{\rm loc}>0\).  Then one may take
\begin{equation}
 w_{ib}=\rho_{\rm loc}^{-1}\mathbf1_{\{i=b\}},
 \qquad
 C_{\rm col}=\rho_{\rm loc}^{-1},
 \label{eq:V7-product-weights}
\end{equation}
and
\begin{equation}
 \boxed{\gap(H)\ge\kappa\rho_{\rm loc}}
 \label{eq:V7-product-gap}
\end{equation}
independently of \(N\).
\end{corollary}

\begin{proof}
Condition on the first \(i-1\) variables.  The difference \(D_if\) is the
mean-zero fluctuation in variable \(i\) of the future-averaged function
\(M_if\).  Apply \eqref{eq:V7-local-Poincare} in that variable and integrate
the revealed variables.  Differentiation commutes with integration over the
independent future variables, so conditional Jensen gives
\[
 \norm{D_if}_2^2
 \le\rho_{\rm loc}^{-1}
     \int|\nabla_iM_if|^2\,\dd\pi
 \le\rho_{\rm loc}^{-1}
     \int|\nabla_if|^2\,\dd\pi.
\]
Apply \Cref{thm:V7-martingale-collar}.
\end{proof}

\begin{assessment}[What must replace product independence]
For the physical Yang--Mills ground law, differentiating a conditional
expectation also differentiates its conditional density.  The resulting
score covariance propagates a revealed block into a collar of later blocks.
Those response terms must populate the off-diagonal \(w_{ib}\).  Exponential
decay of individual entries is sufficient only when their column sums remain
uniform.  This is the martingale form of the V35 weighted-return edge, but it
is posed directly on the physical law \(\pi_j\) and therefore avoids the
same-law circularity.
\end{assessment}

% =============================================================================
\section{V7 cofinal alternatives and ledger}
\label{sec:V7-cofinal-ledger}
% =============================================================================

\begin{theorem}[Physical-ground-law regulator compiler]
\label{thm:V7-regulator-compiler}
Let \(H_j\) be response-matched finite gauge Hamiltonians whose physical
ground laws \(\pi_j\) are obtained from the Perron factors of the same
transfer kernels.  A uniform gap follows from either:
\begin{enumerate}[label=\textup{(\Roman*)},leftmargin=4.5em]
\item comparison laws \(\nu_j\) satisfying
      \(\rho_{\nu,j}\ge\rho_*>0\) and
      \(\operatorname{osc}\log(\dd\nu_j/\dd\pi_j)\le A_*\);
\item martingale-collar rows \eqref{eq:V7-collar-row} for \(\pi_j\) with
      \(C_{{\rm col},j}\le C_*<\infty\).
\end{enumerate}
The respective lower bounds are
\begin{equation}
 \boxed{
 \gap(H_j)\ge
 \rho_*e^{-A_*}
 \quad\hbox{or}\quad
 \gap(H_j)\ge\kappa_j/C_*.}
 \label{eq:V7-two-ground-law-bounds}
\end{equation}
To yield a positive physical continuum mass, the second route additionally
requires a positive lower bound on the physically calibrated coefficient
\(\kappa_j\), or the corresponding rescaled time statement.
\end{theorem}

\begin{proof}
The first route is \Cref{cor:V7-comparison-law-compiler}; the second is
\Cref{thm:V7-martingale-collar}.  The statement about \(\kappa_j\) follows
from the displayed second bound and is exactly a physical-unit calibration,
not a dimensionless lattice-gap claim.
\end{proof}

\begin{theorem}[Integrated V7 statement]
\label{thm:V7-integrated}
Preserving V1--V6, V7 proves:
\begin{enumerate}[label=\textup{(V7.\arabic*)},leftmargin=5em]
\item the physical one-slice ground law and two-slice law of the same positive
      transfer kernel contain the Perron factors \(h^2\) and \(h(x)h(y)\);
\item the transfer Hamiltonian and its Doob ground-law generator are unitarily
      identical;
\item an unweighted positive slab marginal need not be the physical law;
\item a comparison-law Poincar\'e gap transports with the exact
      \(e^{-\operatorname{osc}\theta}\) bounded-density loss;
\item the law mismatch produces the explicit drift
      \(\kappa\nabla\theta\cdot\nabla\);
\item one column-congestion bound on martingale collar rows proves a uniform
      Poincar\'e inequality directly on the physical ground law;
\item product tensorization is recovered with zero off-diagonal response.
\end{enumerate}
\end{theorem}

\begin{proof}
Items \textup{(V7.1)--(V7.3)} are
\Cref{thm:V7-physical-law,cor:V7-transfer-generator,
cth:V7-unweighted-mismatch}.  Items \textup{(V7.4)--(V7.5)} are
\Cref{thm:V7-density-transport,thm:V7-generator-defect}.  Items
\textup{(V7.6)--(V7.7)} are
\Cref{thm:V7-martingale-collar,cor:V7-product-tensorization}.
\end{proof}

\begingroup
\small
\begin{longtable}{>{\raggedright\arraybackslash}p{0.28\textwidth}
                  >{\raggedright\arraybackslash}p{0.16\textwidth}
                  >{\raggedright\arraybackslash}p{0.48\textwidth}}
\toprule
\textbf{V7 row}&\textbf{Status}&\textbf{Advance or remaining obstruction}\\
\midrule
\endfirsthead
\toprule
\textbf{V7 row}&\textbf{Status}&\textbf{Advance or remaining obstruction}\\
\midrule
\endhead
Physical Perron ground law
& \MGpass
& Exact same-kernel one- and two-slice formulas.\\[0.35em]
Unweighted/reduced law identification
& Refuted in general
& A positive normalized slab can have the wrong one-slice law.\\[0.35em]
Bounded-density comparison
& \MGpass
& Uniform comparison gap transfers if the log-density oscillation is uniformly bounded.\\[0.35em]
Actual V35-to-ground-law defect
& \MGopen
& Need \(\theta_j\), its oscillation, and its drift on the response-matched regulator sequence.\\[0.35em]
Physical-law martingale compiler
& \MGpass
& Exact implication from local collar rows and uniform column congestion.\\[0.35em]
Yang--Mills collar rows
& \MGopen
& Conditional score responses for \(\pi_j=h_j^2\mu_j\) have not been bounded uniformly.\\[0.35em]
Same-generator OS continuum
& \MGopen
& Finite transfer incidence is exact; continuum time calibration, tightness, and nontriviality remain.\\[0.35em]
Full mass gap
& \MGopen
& Requires one uniform physical-law route plus the V1 continuum passage.\\
\bottomrule
\end{longtable}
\endgroup

\begin{openproblem}[V8 mandate: acquire physical Perron score collars]
\label{op:V8}
Work with the Perron-weighted law of the same Wilson transfer kernel.
\begin{enumerate}[label=\textup{(V8.\arabic*)},leftmargin=5em]
\item differentiate the finite-volume Perron equation with respect to one
      revealed link/block and express \(\nabla\log h_j\) as a centered
      half-slab response, retaining the physical time normalization;
\item derive the conditional-density score of
      \(\pi_j=h_j^2\mu_j\) and decompose each martingale derivative into its
      direct block term and propagated score covariance;
\item place the resulting coefficients into the explicit collar matrix
      \(w_{ib}^{(j)}\), with disconnected scalar contributions cancelled
      before absolute values are taken;
\item prove a uniform column-congestion bound, using the existing V35
      comparison-inverse decay only after the Radon--Nikodym drift
      \(\nabla\theta_j\) has been inserted;
\item if congestion diverges, return a normalized column and its spatial
      scale as a candidate long-wavelength physical soft mode rather than
      hiding it in a global norm;
\item on the positive branch use
      \Cref{thm:V7-regulator-compiler} and complete the V1 same-generator,
      nontrivial OS limit passage.
\end{enumerate}
If the Perron derivative requires the desired gap to invert
\(I-P_j\), do not use that inverse.  Return upstream to finite-time
half-slab score identities or a monotone conditional coupling, so the proof
does not assume the spectral gap it is intended to establish.
\end{openproblem}

\begin{firewall}[End-of-V7 claim boundary]
V7 identifies the correct physical ground law of the transfer kernel,
quantifies comparison-law mismatch, and gives a noncircular martingale
criterion for a uniform finite-volume gap.  It does not populate the
physical Perron score-collar matrix, prove bounded density mismatch with the
reduced-fibre law, or complete the same-generator nontrivial continuum
limit.  The full Yang--Mills mass gap remains open.
\end{firewall}

\section*{Version V7 append-only record}
\addcontentsline{toc}{section}{Version V7 append-only record}
The complete V6 source was retained before this continuation.  Its SHA--256
digest was
\begin{center}\scriptsize\ttfamily
1af22ede052207d089c267f4f872625361d275e32d765749d5bc5d9c7d94e0fa.
\end{center}
On the next iteration, retain this full file, append before its unique active
document-closing command, increment the filename to
\texttt{\_V8.tex}, and remove only the superseded V7 file from this note
series.

% =============================================================================
% END APPENDED VERSION V7 MATERIAL
% =============================================================================

% =============================================================================
% BEGIN APPENDED VERSION V8 MATERIAL
% =============================================================================

\clearpage
\part{Version V8: Perron score identities and the physical Witten curvature card}
\label{part:V8}

\section{Direction}
\label{sec:V8-direction}

V7 asked for derivatives of the physical Perron factor without using the
reduced resolvent of \(I-P\), since a uniform bound for that inverse would
already contain the desired gap.  Such derivatives are available directly
from the pointwise Perron equation.  They give a one-step, terminal-weighted
score law and an exact expectation-minus-covariance formula for the Hessian
of the physical ground-state potential.

This formula is structurally parallel to the comparison/return Witten cards
in the supplied V35 programme, but it is new in one decisive respect: every
expectation is taken under the Perron-weighted conditional law of the same
transfer kernel.  It therefore identifies the precise terminal factor that
must occur before the comparison card can be used for the Hamiltonian
ground law.

% =============================================================================
\section{First and second Perron boundary scores}
\label{sec:V8-Perron-scores}
% =============================================================================

Let \(X\) be a compact Riemannian configuration manifold with reference
probability \(\mathfrak m\).  Let
\begin{equation}
 k(x,y)=e^{-S(x,y)}>0
 \label{eq:V8-kernel-action}
\end{equation}
be a smooth symmetric transfer kernel.  Let
\[
 \int_X k(x,y)h(y)\,\dd\mathfrak m(y)=\lambda h(x),
 \qquad h>0,
\]
with \(h\) normalized as in V7.  Define the physical forward conditional
law
\begin{equation}
 \boxed{
 \dd p_x(y)
 :=\frac{e^{-S(x,y)}h(y)}
          {\lambda h(x)}\,\dd\mathfrak m(y).}
 \label{eq:V8-physical-conditional}
\end{equation}
For the \(x\)-boundary derivatives put
\begin{equation}
 J_x(y):=\nabla_xS(x,y),
 \qquad
 K_x(y):=\nabla_x^2S(x,y).
 \label{eq:V8-J-K}
\end{equation}
The covariance of the covector \(J_x\) is the positive semidefinite
two-tensor
\begin{equation}
 \Cov_{p_x}(J_x)
 :=\mathbb E_{p_x}
 \bigl[(J_x-\mathbb E_{p_x}J_x)
       \otimes
       (J_x-\mathbb E_{p_x}J_x)\bigr].
 \label{eq:V8-score-covariance}
\end{equation}

\begin{theorem}[Gap-free Perron score identities]
\label{thm:V8-Perron-score}
At every boundary point \(x\),
\begin{align}
 \boxed{\nabla\log h(x)
 &=-\mathbb E_{p_x}J_x,}
 \label{eq:V8-first-score}\\
 \boxed{\nabla^2\log h(x)
 &=-\mathbb E_{p_x}K_x
   +\Cov_{p_x}(J_x).}
 \label{eq:V8-second-score}
\end{align}
No inverse of \(I-P\), spectral-gap estimate, or long-time limit is used.
\end{theorem}

\begin{proof}
Set
\[
 Z_h(x):=\int e^{-S(x,y)}h(y)\,\dd\mathfrak m(y)
 =\lambda h(x).
\]
Because \(X\) is compact and the data are smooth, differentiation under the
integral is valid.  Since \(\lambda\) is independent of the boundary
variable \(x\),
\[
 \nabla\log h
 =\nabla\log Z_h
 =\mathbb E_{p_x}\nabla_x\log k
 =-\mathbb E_{p_x}J_x.
\]
Differentiating the log-partition function once more gives the standard
tensor identity
\[
 \nabla^2\log Z_h
 =\mathbb E_{p_x}\nabla_x^2\log k
  +\Cov_{p_x}(\nabla_x\log k).
\]
Use \(\log k=-S\); the covariance is unchanged by the minus sign.  This
proves \eqref{eq:V8-second-score}.
\end{proof}

\begin{theorem}[Derivative of a physical Doob conditional expectation]
\label{thm:V8-conditional-derivative}
Let \(g(x,y)\) be a smooth scalar- or finite-vector-valued function and put
\[
 F(x)=\mathbb E_{p_x}g(x,\cdot).
\]
Then
\begin{equation}
 \boxed{
 \nabla F(x)
 =\mathbb E_{p_x}\nabla_xg
  -\Cov_{p_x}(g,J_x).}
 \label{eq:V8-conditional-response}
\end{equation}
Here the covariance is polarized componentwise.  In particular,
\begin{equation}
 \norm{\Cov_{p_x}(g,J_x)}
 \le
 \sqrt{\Var_{p_x}(g)}
 \sqrt{\mathbb E_{p_x}
       \norm{J_x-\mathbb E_{p_x}J_x}^2}.
 \label{eq:V8-covariance-CS}
\end{equation}
\end{theorem}

\begin{proof}
Differentiate \eqref{eq:V8-physical-conditional}.  By
\eqref{eq:V8-first-score},
\[
 \nabla_x\log p_x(y)
 =-J_x(y)-\nabla\log h(x)
 =-\bigl(J_x(y)-\mathbb E_{p_x}J_x\bigr).
\]
Thus
\[
 \nabla F
 =\mathbb E_{p_x}\nabla_xg
  +\mathbb E_{p_x}\bigl[g\,\nabla_x\log p_x\bigr],
\]
which is \eqref{eq:V8-conditional-response}.  Cauchy--Schwarz in
\(L^2(p_x)\) proves \eqref{eq:V8-covariance-CS}.
\end{proof}

\begin{assessment}[The physical response source]
The direct boundary source is \(J_x=\nabla_xS\).  Its propagated response is
not an arbitrary covariance under an unweighted shell law; it is precisely
\[
 \Cov_{p_x}(\,\cdot\,,J_x)
\]
under the terminal-weighted conditional law
\eqref{eq:V8-physical-conditional}.  This is the response term that must
populate the off-diagonal martingale collar weights in V7.
\end{assessment}

% =============================================================================
\section{The physical ground potential and its one-form curvature}
\label{sec:V8-physical-curvature}
% =============================================================================

The physical stationary law is
\[
 \dd\pi=h^2\dd\mathfrak m
       =Z^{-1}e^{-W}\dd\mathfrak m,
 \qquad
 W=-2\log h+\text{constant}.
\]

\begin{corollary}[Expectation-minus-covariance ground curvature]
\label{cor:V8-ground-curvature}
The physical ground potential obeys
\begin{align}
 \boxed{\nabla W(x)
 &=2\mathbb E_{p_x}J_x,}
 \label{eq:V8-W-gradient}\\
 \boxed{\nabla^2W(x)
 &=2\mathbb E_{p_x}K_x
   -2\Cov_{p_x}(J_x).}
 \label{eq:V8-W-Hessian}
\end{align}
\end{corollary}

\begin{proof}
Multiply \eqref{eq:V8-first-score} and
\eqref{eq:V8-second-score} by \(-2\).
\end{proof}

Let
\[
 D=\sqrt\kappa\,d,
 \qquad
 L_0=D^*D=\kappa d_\pi^*d
\]
be the scalar physical ground-law generator of V6.  On one-forms let
\[
 \mathscr L_1
 :=\kappa(dd_\pi^*+d_\pi^*d).
\]
On the closure of exact one-forms this operator agrees with \(DD^*\).

\begin{theorem}[Exact physical one-form Witten card]
\label{thm:V8-physical-Witten-card}
On smooth one-forms,
\begin{equation}
 \boxed{
 \mathscr L_1
 =\kappa\left(
   \nabla_\pi^*\nabla
   +\Ric
   +2\mathbb E_{p_x}K_x
   -2\Cov_{p_x}(J_x)
   \right).}
 \label{eq:V8-physical-Witten}
\end{equation}
All tensors act on the \(x\)-boundary cotangent fibre.  On a gauge quotient
or protected gauge-invariant sector, the formula is restricted to the
corresponding horizontal exact one-forms.
\end{theorem}

\begin{proof}
The weighted Bochner--Weitzenb\"ock identity for
\(\pi\propto e^{-W}\mathfrak m\) is
\[
 dd_\pi^*+d_\pi^*d
 =\nabla_\pi^*\nabla+\Ric+\nabla^2W.
\]
Insert \eqref{eq:V8-W-Hessian} and multiply by \(\kappa\).  Exact one-forms
are closed under \(d\), so the \(d_\pi^*d\) term vanishes there and the
restriction is \(DD^*\).
\end{proof}

\begin{corollary}[Pointwise physical curvature certificate]
\label{cor:V8-pointwise-curvature}
Suppose, on every horizontal cotangent fibre,
\begin{equation}
 \Ric_x+2\mathbb E_{p_x}K_x
 -2\Cov_{p_x}(J_x)
 \succeq r_*I
 \qquad(r_*>0)
 \label{eq:V8-curvature-floor}
\end{equation}
with \(r_*\) independent of the regulator volume and cutoff.  Then
\begin{equation}
 \boxed{
 \mathscr L_1\succeq\kappa r_*I
 \quad\text{on exact one-forms},
 \qquad
 \gap(H)\ge\kappa r_*.}
 \label{eq:V8-curvature-gap}
\end{equation}
\end{corollary}

\begin{proof}
The rough weighted Laplacian
\(\nabla_\pi^*\nabla\) is nonnegative.  Apply
\Cref{thm:V8-physical-Witten-card} and then the exact one-form incidence
\Cref{cor:V6-one-form-incidence}.
\end{proof}

\begin{firewall}[Sufficient curvature is not asserted]
The pointwise inequality \eqref{eq:V8-curvature-floor} is a strong
sufficient condition.  Wilson plaquette curvature is not globally positive
on the entire compact configuration space, and gauge directions must be
shorted.  V8 derives the exact physical card; it does not assert that its
pointwise floor is positive.  A comparison-minus-defect or martingale
argument remains necessary if the pointwise curvature changes sign.
\end{firewall}

% =============================================================================
\section{Exact terminal-weight defect of an unweighted fibre law}
\label{sec:V8-terminal-defect}
% =============================================================================

For fixed \(x\), define the bare normalized conditional
\begin{equation}
 \dd q_x(y)
 :=\frac{e^{-S(x,y)}}{Z_0(x)}\,\dd\mathfrak m(y),
 \qquad
 Z_0(x):=\int e^{-S(x,y)}\,\dd\mathfrak m(y).
 \label{eq:V8-bare-conditional}
\end{equation}

\begin{lemma}[Perron terminal-weight Radon--Nikodym factor]
\label{lem:V8-terminal-weight}
The physical and bare conditionals satisfy
\begin{equation}
 \boxed{
 \frac{\dd p_x}{\dd q_x}(y)
 =\frac{Z_0(x)}{\lambda h(x)}\,h(y).}
 \label{eq:V8-terminal-RN}
\end{equation}
Consequently,
\begin{equation}
 \operatorname{osc}_{y\in X}
 \log\frac{\dd p_x}{\dd q_x}(y)
 =\operatorname{osc}_{y\in X}\log h(y),
 \label{eq:V8-terminal-oscillation}
\end{equation}
independently of \(x\).
\end{lemma}

\begin{proof}
Divide \eqref{eq:V8-physical-conditional} by
\eqref{eq:V8-bare-conditional}.  The prefactor is independent of \(y\), so
it cancels from the oscillation.
\end{proof}

\begin{corollary}[Conditional covariance comparison with an explicit loss]
\label{cor:V8-conditional-comparison}
If
\(
 \operatorname{osc}\log h\le A
\),
then the density ratio between \(p_x\) and \(q_x\) has supremum-to-infimum
ratio at most \(e^A\), uniformly in \(x\).  Any Poincar\'e inequality for
\(q_x\) therefore transports to \(p_x\) with loss at most \(e^{-A}\), by
\Cref{thm:V7-density-transport}.
\end{corollary}

\begin{assessment}[Why the terminal factor cannot be dropped]
The terminal factor \(h(y)\) is the complete same-law defect between the
bare one-step shell and the physical Doob conditional.  Bounding
\(\operatorname{osc}\log h\) globally may itself be volume-extensive, so
\Cref{cor:V8-conditional-comparison} is a certificate, not a claim of
cofinal control.  A local proof should instead show that the response of
\(\log h\) to a distant boundary perturbation has summable collars.
\end{assessment}

% =============================================================================
\section{A physical score-collar acquisition}
\label{sec:V8-score-collar}
% =============================================================================

The martingale rows in V7 can now be populated without a spectral inverse at
the identity level.  If \(g\) is a future conditional observable, its
derivative across a revealed boundary block is the two-term packet
\[
 \boxed{
 \text{direct derivative }
 \mathbb E_{p_x}\nabla_xg
 \quad-\quad
 \text{physical score response }
 \Cov_{p_x}(g,J_x).}
 \label{eq:V8-two-term-collar}
\]
The second term obeys \eqref{eq:V8-covariance-CS}.  Thus a noncircular local
row may be acquired from:
\begin{enumerate}[label=\textup{(\alph*)},leftmargin=3.5em]
\item a conditional variance bound for \(g\) on a finite future collar;
\item a bound for the physical one-step score covariance
      \(\mathbb E_{p_x}\norm{J_x-\mathbb E_{p_x}J_x}^2\);
\item a summable decay estimate for the effect of moving the revealed block
      on successively more distant conditional observables.
\end{enumerate}
No global inverse of \(I-P\) occurs in these three entries.

\begin{theorem}[Finite physical score-collar implication]
\label{thm:V8-finite-score-collar}
Suppose that, for each reveal step \(i\), the exact derivative identity
\eqref{eq:V8-two-term-collar} and conditional Poincar\'e/Cauchy estimates
produce nonnegative coefficients \(w_{ib}\) satisfying the V7 row
\eqref{eq:V7-collar-row}.  If
\begin{equation}
 \sup_j\max_b\sum_iw_{ib}^{(j)}\le C_*<\infty,
 \qquad
 \inf_j\kappa_j\ge\kappa_*>0,
 \label{eq:V8-uniform-score-congestion}
\end{equation}
then
\begin{equation}
 \boxed{\gap(H_j)\ge\kappa_*/C_*}
 \label{eq:V8-score-gap}
\end{equation}
uniformly.
\end{theorem}

\begin{proof}
The hypotheses deliver exactly the inputs of
\Cref{thm:V7-martingale-collar}; apply that theorem at each regulator.
\end{proof}

\begin{firewall}[Identity versus acquired bound]
\Cref{thm:V8-Perron-score,thm:V8-conditional-derivative} populate the exact
writers of the collar row.  They do not bound the conditional variances or
the spatial column sum in \eqref{eq:V8-uniform-score-congestion}.  Replacing
those entries by the full physical Poincar\'e constant would be circular.
\end{firewall}

% =============================================================================
\section{V8 ledger and next upstream attack}
\label{sec:V8-ledger}
% =============================================================================

\begin{theorem}[Integrated V8 statement]
\label{thm:V8-integrated}
Preserving V1--V7, V8 proves:
\begin{enumerate}[label=\textup{(V8.\arabic*)},leftmargin=5em]
\item gap-free first and second derivative identities for the Perron factor;
\item the exact direct-minus-score-covariance derivative of every physical
      Doob conditional expectation;
\item the expectation-minus-covariance Hessian of the physical ground
      potential \(W=-2\log h\);
\item the exact physical one-form Witten card on horizontal exact forms;
\item a pointwise physical curvature sufficient condition for a uniform gap;
\item the terminal Perron factor as the exact Radon--Nikodym defect between
      bare and physical one-step fibre laws;
\item an explicit route from finite conditional score collars to the V7
      column-congestion gap.
\end{enumerate}
\end{theorem}

\begin{proof}
Items \textup{(V8.1)--(V8.2)} are
\Cref{thm:V8-Perron-score,thm:V8-conditional-derivative}.  Items
\textup{(V8.3)--(V8.5)} are
\Cref{cor:V8-ground-curvature,thm:V8-physical-Witten-card,
cor:V8-pointwise-curvature}.  Item \textup{(V8.6)} is
\Cref{lem:V8-terminal-weight}.  Item \textup{(V8.7)} is
\Cref{thm:V8-finite-score-collar}.
\end{proof}

\begingroup
\small
\begin{longtable}{>{\raggedright\arraybackslash}p{0.28\textwidth}
                  >{\raggedright\arraybackslash}p{0.16\textwidth}
                  >{\raggedright\arraybackslash}p{0.48\textwidth}}
\toprule
\textbf{V8 row}&\textbf{Status}&\textbf{Advance or remaining obstruction}\\
\midrule
\endfirsthead
\toprule
\textbf{V8 row}&\textbf{Status}&\textbf{Advance or remaining obstruction}\\
\midrule
\endhead
Perron boundary score
& \MGpass
& Exact one-step physical conditional identity; no gap inverse.\\[0.35em]
Physical ground Hessian
& \MGpass
& Exact direct curvature minus physical score covariance.\\[0.35em]
Physical one-form Witten incidence
& \MGpass
& Exact on the ground law of the same Perron factor.\\[0.35em]
Bare versus physical fibre law
& \MGpass
& Exact terminal-weight Radon--Nikodym factor; uniform control remains open.\\[0.35em]
Pointwise curvature floor
& Sufficient, unpopulated
& Wilson curvature is not asserted globally positive.\\[0.35em]
Martingale score-column congestion
& \MGopen
& Need summable spatial response of the terminal-weighted score covariance.\\[0.35em]
Same differential/OS generator
& \MGopen
& A common Perron factor does not by itself identify all nonzero spectral data or physical time.\\[0.35em]
Continuum full mass gap
& \MGopen
& Requires the uniform physical collar or another infrared floor plus nontrivial OS convergence.\\
\bottomrule
\end{longtable}
\endgroup

\begin{openproblem}[V9 mandate: summable terminal-weight response]
\label{op:V9}
Use \eqref{eq:V8-physical-conditional}, not the bare shell law.
\begin{enumerate}[label=\textup{(V9.\arabic*)},leftmargin=5em]
\item disintegrate the Perron-weighted half-slab into spatial blocks and
      write the derivative of each conditional log density using
      \eqref{eq:V8-conditional-response};
\item prove a finite-collar conditional variance inequality without invoking
      the full-volume physical gap;
\item bound the score-covariance propagation by a nonnegative spatial matrix
      \(A^{(j)}=(A_{bc}^{(j)})\) whose column sums are uniformly below a
      declared threshold;
\item compare \(A^{(j)}\) with the existing V35 weighted
      comparison/curvature-defect channel after inserting the terminal
      Perron factor in every expectation;
\item if the weighted edge reaches one, return a source-cyclic near-unit
      vector and test whether it produces a normalized exact one-form with
      vanishing physical Rayleigh quotient;
\item on a strict branch, compile the resulting \(w_{ib}^{(j)}\), prove
      \eqref{eq:V8-uniform-score-congestion}, and then execute the V1
      same-generator continuum passage.
\end{enumerate}
If a spatial conditional estimate asks for a global oscillation of
\(\log h_j\), go upstream again to a telescoping local likelihood ratio; the
global oscillation is allowed to grow while its local influence remains
summable.
\end{openproblem}

\begin{firewall}[End-of-V8 claim boundary]
V8 derives the exact physical Perron score, curvature, and conditional
response identities without assuming a gap.  It does not prove the summable
terminal-weight response, a uniform physical one-form floor, equality of the
differential and OS generators in the continuum scaling, or nontriviality of
the limiting theory.  The full Yang--Mills mass gap remains open.
\end{firewall}

\section*{Version V8 append-only record}
\addcontentsline{toc}{section}{Version V8 append-only record}
The complete V7 source was retained before this continuation.  Its SHA--256
digest was
\begin{center}\scriptsize\ttfamily
acc2869f234c615b74081eaa39c08b2bc0ff700181ea310f431d8e1718928466.
\end{center}
On the next iteration, retain this full file, append before its unique active
document-closing command, increment the filename to
\texttt{\_V9.tex}, and remove only the superseded V8 file from this note
series.

% =============================================================================
% END APPENDED VERSION V8 MATERIAL
% =============================================================================

% =============================================================================
% BEGIN APPENDED VERSION V9 MATERIAL
% =============================================================================

\clearpage
\part{Version V9: half-space Perron polymers and a uniform strong-coupling ground-law gap}
\label{part:V9}

\section{Direction and inherited input}
\label{sec:V9-direction}

The supplied reduced-fibre V35 manuscript already proves a Koteck\'y--Preiss
connected-plaquette expansion for a finite Wilson shell when
\[
 |\beta|<\beta_{\rm KP}
 :=\log\left(1+\frac{1}{2e^3\Delta_{\rm plaq}}\right),
 \qquad \Delta_{\rm plaq}\le20.
\]
It also proves differentiated, volume-uniform locality for a finite-shell
effective Hessian.  Repeating that theorem would not address V8: the missing
object is the terminal Perron factor \(h\), which is an infinite half-slab
boundary free energy.

V9 extends the same rooted-polymer domination to the normalized half-slab
partition function, passes it to \(h\) by the Perron limit, and derives a
uniform weighted Hessian row for \(\log h\).  Positive Ricci curvature of the
compact link group then yields a genuine uniform Poincar\'e gap for the
physical ground law on a sufficiently small strong-coupling interval.  The
result closes one regulator branch; it does not cross the independent
strong-coupling/continuum fork identified in the supplied programme.

% =============================================================================
\section{Normalized half-slab boundary free energy}
\label{sec:V9-half-slab}
% =============================================================================

Let \(\Lambda\) be a finite periodic spatial lattice and
\(X_\Lambda=\mathrm{SU}(3)^{E(\Lambda)}\).  Let
\(\mathcal K_{\beta,\Lambda}\) be the positive symmetric one-step Wilson
transfer operator.  For an integer half-slab depth \(n\), define
\begin{equation}
 Z_{n,\Lambda}(x)
 :=(\mathcal K_{\beta,\Lambda}^n\one)(x),
 \qquad
 F_{n,\Lambda}(x;x_*):=
 \log\frac{Z_{n,\Lambda}(x)}{Z_{n,\Lambda}(x_*)}.
 \label{eq:V9-half-slab-free-energy}
\end{equation}
Here \(x_*\) is an arbitrary fixed reference boundary configuration.  The
ratio removes all boundary-independent vacuum clusters.

Use the plaquette adjacency graph from the supplied V35 expansion.  Put
\begin{equation}
 u(\beta)=e^{|\beta|}-1,
 \qquad
 \vartheta_\beta=e^3\Delta_{\rm plaq}u(\beta).
 \label{eq:V9-theta}
\end{equation}
For \(|\beta|<\beta_{\rm KP}\), one has
\(\vartheta_\beta<1/2\).

\begin{theorem}[Uniform boundary-rooted half-slab expansion]
\label{thm:V9-half-slab-expansion}
For \(|\beta|<\beta_{\rm KP}\), each
\(F_{n,\Lambda}(x;x_*)\) has an absolutely convergent connected-polymer
expansion with the following properties.
\begin{enumerate}[label=\textup{(\roman*)},leftmargin=3.5em]
\item Every surviving cluster meets the initial time boundary and contains a
      plaquette whose boundary factor distinguishes \(x\) from \(x_*\).
\item The expansion and its first two boundary-link derivatives converge
      absolutely, uniformly in \(n\), \(\Lambda\), and the boundary values.
\item A mixed derivative in boundary links \(b,c\) receives contributions
      only from clusters connecting the two plaquette-incidence sets of
      \(b\) and \(c\).
\item The same statements hold on the gauge-invariant quotient because all
      cluster weights and boundary derivatives are gauge covariant before
      contraction and gauge invariant after assembly.
\end{enumerate}
\end{theorem}

\begin{proof}
On the \(n\)-step slab, write each Wilson plaquette factor as
\[
 e^{(\beta/3)\operatorname{ReTr}U_p}
 =1+f_{p,\beta}(U_p),
 \qquad
 |f_{p,\beta}|\le u(\beta).
\]
Expanding the product and integrating internal links factors every subset
over its connected components in the plaquette adjacency graph.  The
logarithm is therefore the usual connected hard-core polymer series.  The
rooted animal count is at most
\((e\Delta_{\rm plaq})^{r-1}\) for size \(r\).  The supplied V35 KP estimate,
which is uniform in retained boundary links, dominates the Mayer tree sum by
a geometric series with ratio \(\vartheta_\beta<1/2\).

In the difference of logarithms in
\eqref{eq:V9-half-slab-free-energy}, every cluster disjoint from the initial
boundary has the same weight for \(x\) and \(x_*\) and cancels.  This proves
\textup{(i)}.  A boundary derivative can hit only a plaquette incident with
that boundary link.  One or two derivatives create one or two marked
plaquettes; the number of mark placements is polynomial in the cluster size,
while the unmarked tree domination remains geometric.  Thus the marked
series are absolutely and uniformly summable, proving \textup{(ii)}.
Connectedness forces a two-mark cluster to connect the two incidence sets,
which is \textup{(iii)}.  Gauge restriction commutes with the Wilson
integrations and differentiations, proving \textup{(iv)}.
\end{proof}

% =============================================================================
\section{Passage from finite half-slabs to the Perron factor}
\label{sec:V9-Perron-limit}
% =============================================================================

Let
\[
 \mathcal K_{\beta,\Lambda}h_{\beta,\Lambda}
 =\lambda_{\beta,\Lambda}h_{\beta,\Lambda},
 \qquad h_{\beta,\Lambda}>0.
\]

\begin{theorem}[Boundary-polymer expansion of the Perron factor]
\label{thm:V9-Perron-polymer}
For every finite \(\Lambda\) and
\(|\beta|<\beta_{\rm KP}\),
\begin{equation}
 \boxed{
 \lim_{n\to\infty}F_{n,\Lambda}(x;x_*)
 =\log\frac{h_{\beta,\Lambda}(x)}
            {h_{\beta,\Lambda}(x_*)}}
 \label{eq:V9-Perron-free-energy-limit}
\end{equation}
in \(C^2\) with respect to every fixed finite set of boundary links.
The limiting logarithm has the absolutely convergent expansion obtained by
summing all connected half-space polymers rooted at the initial boundary.
All first- and second-derivative bounds from
\Cref{thm:V9-half-slab-expansion} remain uniform in \(\Lambda\).
\end{theorem}

\begin{proof}
At fixed finite \(\Lambda\), strict positivity and compactness of
\(\mathcal K_{\beta,\Lambda}\) give a simple Perron eigenvalue whose modulus
strictly dominates the rest of the spectrum.  Hence
\[
 \frac{\mathcal K_{\beta,\Lambda}^n\one(x)}
      {\mathcal K_{\beta,\Lambda}^n\one(x_*)}
 \longrightarrow
 \frac{h_{\beta,\Lambda}(x)}
      {h_{\beta,\Lambda}(x_*)}.
\]
The unmarked, once-marked, and twice-marked cluster series in
\Cref{thm:V9-half-slab-expansion} have a common summable majorant independent
of \(n\).  A boundary-rooted finite cluster is present unchanged once the slab
depth exceeds its temporal diameter.  Dominated convergence therefore sends
the finite-slab series to the half-space series term by term, including its
first two derivatives.  The pointwise Perron limit identifies the unmarked
sum, and uniform convergence of the marked series gives the \(C^2\) claim.
\end{proof}

\begin{remark}[No use of the transfer gap]
Perron convergence is used at each fixed finite volume only to identify the
limit function.  Uniform locality comes from the polymer majorant, not from
a uniform estimate on the subleading transfer eigenvalue.  The argument
therefore does not assume the cofinal gap it is intended to support.
\end{remark}

% =============================================================================
\section{An explicit marked-series Hessian bound}
\label{sec:V9-marked-bound}
% =============================================================================

For one plaquette coordinate
\(
 x_p(U_p)=\operatorname{ReTr}U_p/3
\),
compactness gives finite geometric constants
\begin{equation}
 c_1:=\sup_{p,b,U}\norm{\nabla_bx_p(U)},
 \qquad
 c_2:=\sup_{p,b,c,U}\norm{\nabla_c\nabla_bx_p(U)},
 \label{eq:V9-c1-c2}
\end{equation}
where derivatives vanish unless the displayed links meet \(p\).  Let
\(d_\partial\le6\) be the number of plaquettes incident with one link.  In
this section, \(\dist(b,c)\) means the plaquette-graph distance between the
incidence sets of the two boundary links.  Thus a connected cluster of
\(r\) plaquettes carrying both marks satisfies \(\dist(b,c)\le r-1\).  Put
\begin{align}
 q_1(\beta)&:=|\beta|e^{|\beta|}c_1,
 \label{eq:V9-q1}\\
 q_2(\beta)&:=e^{|\beta|}
 \bigl(|\beta|c_2+\beta^2c_1^2\bigr),
 \label{eq:V9-q2}\\
 m_2(\beta)&:=
 q_2(\beta)+
 \begin{cases}
 q_1(\beta)^2/u(\beta),&\beta\ne0,\\
 0,&\beta=0.
 \end{cases}
 \label{eq:V9-m2}
\end{align}
Then \(m_2(\beta)\to0\) as \(\beta\to0\).

For
\(
 0\le\mu<-\log\vartheta_\beta
\),
set
\begin{equation}
 \xi_{\beta,\mu}:=e^\mu\vartheta_\beta<1
 \label{eq:V9-xi}
\end{equation}
and define the conservative marked-tree majorant
\begin{equation}
 \boxed{
 \mathfrak B_2(\beta,\mu):=
 64d_\partial^2(\Delta_{\rm plaq}+1)^2
 m_2(\beta)
 \frac{1+11\xi_{\beta,\mu}
       +11\xi_{\beta,\mu}^2+\xi_{\beta,\mu}^3}
      {(1-\xi_{\beta,\mu})^5}.}
 \label{eq:V9-B2}
\end{equation}

\begin{theorem}[Weighted Schur bound for the Perron Hessian]
\label{thm:V9-Perron-Hessian-bound}
Under the link and plaquette conventions above,
\begin{equation}
 \boxed{
 \sup_b\sum_c
 e^{\mu\dist(b,c)}
 \norm{\nabla_c\nabla_b
       \log h_{\beta,\Lambda}}_{\op}
 \le\mathfrak B_2(\beta,\mu)}
 \label{eq:V9-weighted-Hessian}
\end{equation}
uniformly in the periodic spatial volume.  The same bound holds with rows
and columns interchanged.  In particular,
\begin{equation}
 \lim_{\beta\to0}\mathfrak B_2(\beta,0)=0.
 \label{eq:V9-B2-zero}
\end{equation}
\end{theorem}

\begin{proof}
Differentiating
\(
 f_{p,\beta}=e^{\beta x_p}-1
\)
once or twice gives
\[
 \norm{\nabla_bf_{p,\beta}}\le q_1(\beta),
 \qquad
 \norm{\nabla_c\nabla_bf_{p,\beta}}\le q_2(\beta).
\]
If both marks hit one activity, replace one factor \(u(\beta)\) by
\(q_2(\beta)\).  If they hit two activities, replace two factors by
\(q_1(\beta)^2\).  Factoring the remaining unmarked activity product gives
the common marked factor \(m_2(\beta)\).  There are at most
\(d_\partial^2\) choices of incident root plaquettes; the remaining
orientation, incompatibility-neighbourhood, and two-mark tree choices are
overcounted by
\(
64(\Delta_{\rm plaq}+1)^2r^4
\)
at cluster size \(r\), exactly as in the supplied V35 differentiated
tree majorant.

For a connected cluster carrying marks at \(b,c\), the incidence-set
distance convention gives \(\dist(b,c)\le r-1\).  The exponential weight is therefore
bounded by \(e^{\mu(r-1)}\), replacing
\(\vartheta_\beta\) by \(\xi_{\beta,\mu}\).  Summing the dominated marked
series and using
\[
 \sum_{r\ge1}r^4\xi^{r-1}
 =\frac{1+11\xi+11\xi^2+\xi^3}{(1-\xi)^5}
\]
proves \eqref{eq:V9-weighted-Hessian}.  The constants are independent of the
slab depth and volume, so
\Cref{thm:V9-Perron-polymer} passes the estimate to \(h\).
Finally \(m_2(\beta)\to0\) and
\(\vartheta_\beta\to0\), proving \eqref{eq:V9-B2-zero}.
\end{proof}

\begin{firewall}[Conservative constant]
The factor in \eqref{eq:V9-B2} deliberately overcounts orientations, mark
placements, and incompatibility neighbours.  It is a rigorous existence
bound on a very small strong-coupling interval, not a proposed numerical
estimate of the physical Yang--Mills mass.
\end{firewall}

% =============================================================================
\section{A uniform strong-coupling ground-law gap}
\label{sec:V9-strong-gap}
% =============================================================================

Equip \(\mathrm{SU}(3)\) with the bi-invariant metric used to define the
electric Laplacian.  Since \(\mathrm{SU}(3)\) is compact and simple, there is
a geometric constant \(r_G>0\) such that
\begin{equation}
 \Ric_{\mathrm{SU}(3)}\succeq r_GI.
 \label{eq:V9-Ricci-floor}
\end{equation}
The product manifold \(X_\Lambda\) has the same lower Ricci bound,
independently of the number of links.

\begin{theorem}[Uniform Poincar\'e gap of the strong-coupling Perron law]
\label{thm:V9-strong-Poincare}
There exists \(\beta_{\rm sc}>0\), independent of \(\Lambda\), such that
\begin{equation}
 0\le\beta\le\beta_{\rm sc}
 \quad\Longrightarrow\quad
 2\mathfrak B_2(\beta,0)<r_G.
 \label{eq:V9-beta-sc-choice}
\end{equation}
For the physical Perron law
\[
 \dd\pi_{\beta,\Lambda}
 =h_{\beta,\Lambda}^2\,\dd\mathfrak m_\Lambda,
\]
one then has
\begin{equation}
 \boxed{
 \Var_{\pi_{\beta,\Lambda}}(f)
 \le
 \frac{1}{r_G-2\mathfrak B_2(\beta,0)}
 \int\norm{\nabla f}^2\,\dd\pi_{\beta,\Lambda}}
 \label{eq:V9-strong-Poincare}
\end{equation}
for every smooth \(f\), uniformly in the periodic volume.  The inequality
remains true after restriction to gauge-invariant functions.
\end{theorem}

\begin{proof}
\Cref{thm:V9-Perron-Hessian-bound} and the symmetric block Schur test give
\[
 \norm{\nabla^2\log h_{\beta,\Lambda}}_{\op}
 \le\mathfrak B_2(\beta,0).
\]
For the ground potential
\(
 W_{\beta,\Lambda}=-2\log h_{\beta,\Lambda}
\),
\[
 \nabla^2W_{\beta,\Lambda}
 \succeq-2\mathfrak B_2(\beta,0)I.
\]
Together with \eqref{eq:V9-Ricci-floor},
\[
 \Ric+\nabla^2W_{\beta,\Lambda}
 \succeq
 \bigl(r_G-2\mathfrak B_2(\beta,0)\bigr)I.
\]
The right side is positive on a nonempty interval by
\eqref{eq:V9-B2-zero}; choose any \(\beta_{\rm sc}\) within that interval
and below \(\beta_{\rm KP}\).  The weighted Bochner identity and
\Cref{cor:V8-pointwise-curvature} give
\eqref{eq:V9-strong-Poincare}.  Restricting the class of test functions to
the gauge-invariant subspace cannot worsen the constant.
\end{proof}

\begin{corollary}[Uniform strong-coupling differential-Hamiltonian gap]
\label{cor:V9-strong-Hamiltonian-gap}
Let \(H_{\beta,\Lambda}\) be a V6 differential gauge Hamiltonian whose
normalized ground state is the same Perron factor
\(h_{\beta,\Lambda}\), with electric coefficient
\(\kappa_{\beta,\Lambda}>0\).  Then, for
\(0\le\beta\le\beta_{\rm sc}\),
\begin{equation}
 \boxed{
 \gap(H_{\beta,\Lambda})
 \ge
 \kappa_{\beta,\Lambda}
 \bigl(r_G-2\mathfrak B_2(\beta,0)\bigr).}
 \label{eq:V9-strong-Hamiltonian-gap}
\end{equation}
If
\(\inf_{\beta,\Lambda}\kappa_{\beta,\Lambda}\ge\kappa_*>0\)
on this branch, the lower bound is uniform in volume.
\end{corollary}

\begin{proof}
Apply the ground-state representation
\Cref{thm:V6-ground-state-representation} to
\eqref{eq:V9-strong-Poincare}.
\end{proof}

\begin{assessment}[A row genuinely closed]
The physical Perron terminal factor is no longer an uncontrolled global
oscillation on the small-\(\beta\) branch.  Its log Hessian has a
volume-uniform summable row, and the ground law has a uniform Poincar\'e
constant.  This is stronger than a positive finite Wilson card and is
obtained without a uniform transfer-gap input.
\end{assessment}

% =============================================================================
\section{Why this is not the continuum mass gap}
\label{sec:V9-continuum-fork}
% =============================================================================

\begin{firewall}[Strong-coupling/continuum fork]
For four-dimensional Wilson Yang--Mills, removal of the ultraviolet lattice
cutoff follows a weak-bare-coupling trajectory, conventionally
\(\beta(a)\to\infty\), not the fixed small-\(\beta\) interval of
\Cref{thm:V9-strong-Poincare}.  The theorem therefore supplies a rigorous
strong-coupling thermodynamic ground-law gap but does not reach the
continuum scaling trajectory.
\end{firewall}

\begin{firewall}[Ground factor versus full generator]
Even when a Wilson transfer operator and a differential Hamiltonian share
the same positive ground factor, equality of their nonzero spectra does not
follow.  \Cref{cor:V9-strong-Hamiltonian-gap} applies to the V6 differential
Hamiltonian with that ground state.  A mass statement for the reflected
Wilson transfer requires the same-generator and physical-time incidence
still open in V1 and V8.
\end{firewall}

\begin{counterexample}[The ground law does not determine the transfer gap]
\label{cth:V9-law-no-transfer-gap}
On a three-point probability space with uniform stationary law, let
\[
 P_\varepsilon
 =\varepsilon I+(1-\varepsilon)\Pi,
 \qquad 0<\varepsilon<1,
\]
where \(\Pi\) averages over all three points.  Every \(P_\varepsilon\) has
the same positive ground vector and stationary law, while its centered
operator norm is \(\varepsilon\).  Hence the transfer Hamiltonian gap
\(-\tau^{-1}\log\varepsilon\) varies with \(\varepsilon\).
\end{counterexample}

\begin{proof}
\(\Pi\) is identity on constants and zero on the centered subspace.  Thus
\(P_\varepsilon\) is identity on constants and multiplication by
\(\varepsilon\) on the centered subspace.
\end{proof}

This counterexample prevents the newly proved Poincar\'e gap of the ground law
from being reported as an OS transfer gap before a Dirichlet-form or
same-generator comparison is established.

% =============================================================================
\section{V9 ledger and next attack}
\label{sec:V9-ledger}
% =============================================================================

\begin{theorem}[Integrated V9 statement]
\label{thm:V9-integrated}
Preserving V1--V8, V9 proves:
\begin{enumerate}[label=\textup{(V9.\arabic*)},leftmargin=5em]
\item a uniform connected-polymer expansion for normalized Wilson
      half-slab boundary free energies;
\item \(C^2\) passage of that expansion to the Perron factor without a
      uniform spectral-gap input;
\item an explicit weighted Schur majorant for
      \(\nabla^2\log h_{\beta,\Lambda}\);
\item a volume-uniform Poincar\'e inequality for the physical Perron law on a
      nonempty strong-coupling interval;
\item a uniform differential-Hamiltonian gap on the same-ground-state branch
      when the electric coefficient has a common physical lower bound;
\item counterexamples and firewalls separating that result from the
      reflected transfer gap and the continuum weak-coupling trajectory.
\end{enumerate}
\end{theorem}

\begin{proof}
Items \textup{(V9.1)--(V9.3)} are
\Cref{thm:V9-half-slab-expansion,thm:V9-Perron-polymer,
thm:V9-Perron-Hessian-bound}.  Items \textup{(V9.4)--(V9.5)} are
\Cref{thm:V9-strong-Poincare,cor:V9-strong-Hamiltonian-gap}.  Item
\textup{(V9.6)} is
\Cref{cth:V9-law-no-transfer-gap} and the two preceding firewalls.
\end{proof}

\begingroup
\small
\begin{longtable}{>{\raggedright\arraybackslash}p{0.28\textwidth}
                  >{\raggedright\arraybackslash}p{0.16\textwidth}
                  >{\raggedright\arraybackslash}p{0.48\textwidth}}
\toprule
\textbf{V9 row}&\textbf{Status}&\textbf{Advance or remaining obstruction}\\
\midrule
\endfirsthead
\toprule
\textbf{V9 row}&\textbf{Status}&\textbf{Advance or remaining obstruction}\\
\midrule
\endhead
Half-space Perron polymer
& \MGpass
& Uniform \(C^2\) boundary-rooted expansion at sufficiently small \(\beta\).\\[0.35em]
Terminal-weight Hessian collars
& \MGpass
& Explicit summable marked-tree majorant independent of volume.\\[0.35em]
Strong-coupling Perron-law Poincar\'e gap
& \MGpass
& Positive and volume-uniform on a nonempty small-\(\beta\) interval.\\[0.35em]
Differential Hamiltonian on the same ground law
& Conditional pass
& Uniform after same-ground-state incidence and a common physical electric coefficient.\\[0.35em]
Reflected Wilson transfer generator
& \MGopen
& A common ground law alone does not compare the centered Dirichlet forms.\\[0.35em]
Weak-coupling continuum trajectory
& \MGopen
& The KP interval does not contain \(\beta(a)\to\infty\).\\[0.35em]
Continuum state and full mass gap
& \MGopen
& Need a weak-coupling uniform route, same-generator comparison, and nontrivial OS convergence.\\
\bottomrule
\end{longtable}
\endgroup

\begin{openproblem}[V10 mandate: cross the strong/weak fork without a global convexity claim]
\label{op:V10}
The small-\(\beta\) polymer branch is now closed at the physical ground-law
level.  Continue upstream along the actual continuum trajectory:
\begin{enumerate}[label=\textup{(V10.\arabic*)},leftmargin=5.5em]
\item place the response-matched bare coupling \(\beta(a)\), anisotropy, and
      electric coefficient on one declared Wilson transfer/differential
      Hamiltonian trajectory;
\item derive an exact centered Dirichlet-form comparison between
      \(I-P_{a,\Lambda}\) and the physical ground-law diffusion form, with
      constants in physical time units;
\item replace the small-activity polymer bound by a weak-coupling
      gauge-fixed local chart plus a compact large-field return, retaining
      the V35 comparison/defect Woodbury split;
\item prove a weighted edge strictly below one on the complete exact
      one-form range, or return its source-cyclic near-unit vector;
\item control chart exits, Gribov/toron sectors, and wrapping modes separately
      rather than charging them to one global Hessian norm;
\item combine the resulting weak-coupling one-form floor with V1
      quotient-stable OS convergence and prove nontriviality of the limiting
      local algebra.
\end{enumerate}
If the weak-chart comparison loses its margin in a topological or toron
sector, route that specific vector back through the V1 local-versus-wrapping
classifier before altering the local curvature estimate.
\end{openproblem}

\begin{firewall}[End-of-V9 claim boundary]
V9 proves a volume-uniform Poincar\'e gap for the physical Perron ground law
on a nonempty strong-coupling interval and a corresponding conditional
differential-Hamiltonian gap.  It does not prove a weak-coupling continuum
gap, the centered Dirichlet-form comparison with the reflected Wilson
transfer, or construction and nontriviality of the continuum Yang--Mills
state.  The full Yang--Mills mass gap remains open.
\end{firewall}

\section*{Version V9 append-only record}
\addcontentsline{toc}{section}{Version V9 append-only record}
The complete V8 source was retained before this continuation.  Its SHA--256
digest was
\begin{center}\scriptsize\ttfamily
7bcbf45227ecf62f7e4671839695a2d450a571d7cf17d837423171aa258c64af.
\end{center}
On the next iteration, retain this full file, append before its unique active
document-closing command, increment the filename to
\texttt{\_V10.tex}, and remove only the superseded V9 file from this note
series.

% =============================================================================
% END APPENDED VERSION V9 MATERIAL
% =============================================================================

% =============================================================================
% BEGIN APPENDED VERSION V10 MATERIAL
% =============================================================================

\clearpage
\part{Version V10: temporal Wilson symbols and the same-generator Hamiltonian limit}
\label{part:V10}

\section{Direction}
\label{sec:V10-direction}

V9 closed a small-\(\beta\) ground-law branch but left a logically separate
incidence problem: the reflected Wilson transfer and the differential
electric-plus-magnetic Hamiltonian need not have the same nonzero spectrum
merely because they share a Perron factor.  V10 attacks that problem at the
temporal lattice scale.

The temporal Wilson kernel is a central convolution operator.  Its exact
Peter--Weyl multipliers determine an effective electric symbol.  The correct
same-generator comparison therefore has two independent rows:
\[
 \boxed{
 \text{uniform symbol accuracy on a growing Casimir core}
 \quad+\quad
 \text{coercive energy on the omitted representation tail}.}
\]
Fixed-representation asymptotics prove only the first row at fixed core.
V10 derives that asymptotic with a physical calibration, proves a
finite-core error formula, and gives a Mosco--Chernoff compiler for the
Hamiltonian limit.

% =============================================================================
\section{Exact temporal character multipliers}
\label{sec:V10-character-symbol}
% =============================================================================

Let \(G=\mathrm{SU}(3)\), \(d_G=8\), and
\[
 x(U):=\frac13\operatorname{ReTr}U.
\]
For \(\alpha>0\), define the central inversion-invariant probability law
\begin{equation}
 \dd\nu_\alpha(U)
 :=Z_\alpha^{-1}e^{\alpha x(U)}\,\dd U.
 \label{eq:V10-temporal-law}
\end{equation}
Let \(\mathcal C_\alpha\) be convolution by \(\nu_\alpha\).  For an
irreducible representation \(R=(p,q)\), of dimension \(d_R\), put
\begin{equation}
 \boxed{
 \eta_R(\alpha)
 :=\frac1{d_R}\int_G\chi_R(U)\,\dd\nu_\alpha(U).}
 \label{eq:V10-character-multiplier}
\end{equation}
Reflection positivity of the temporal Wilson factor is taken in its literal
character form:
\begin{equation}
 0<\eta_R(\alpha)\le1,
 \qquad
 \eta_{\mathbf1}(\alpha)=1.
 \label{eq:V10-positive-multipliers}
\end{equation}

\begin{lemma}[Exact Peter--Weyl diagonalization]
\label{lem:V10-Peter-Weyl-diagonal}
On every matrix-coefficient copy of \(R\),
\begin{equation}
 \boxed{\mathcal C_\alpha|_R=\eta_R(\alpha)I.}
 \label{eq:V10-convolution-diagonal}
\end{equation}
Consequently, for any declared temporal step \(\tau_\alpha>0\),
\begin{equation}
 A_\alpha:=-\tau_\alpha^{-1}\log\mathcal C_\alpha
 \label{eq:V10-Aalpha}
\end{equation}
is a positive electric operator with representation symbol
\begin{equation}
 a_R(\alpha):=-\tau_\alpha^{-1}\log\eta_R(\alpha).
 \label{eq:V10-aR}
\end{equation}
\end{lemma}

\begin{proof}
Centrality makes the Fourier transform of \(\nu_\alpha\) an intertwiner on
each irreducible representation, hence a scalar by Schur's lemma.  Taking
its trace gives \eqref{eq:V10-character-multiplier}.  Functional calculus and
\eqref{eq:V10-positive-multipliers} give the remaining assertions.
\end{proof}

% =============================================================================
\section{Fundamental calibration and the fixed-representation Casimir limit}
\label{sec:V10-fixed-representation}
% =============================================================================

Let \(F=(1,0)\) be the fundamental representation and
\(C_F=C_2(F)=4/3\).  Define
\begin{equation}
 s_\alpha:=-\frac{\log\eta_F(\alpha)}{C_F},
 \qquad
 \boxed{\tau_\alpha:=s_\alpha/\kappa}
 \label{eq:V10-time-calibration}
\end{equation}
for a fixed desired electric coefficient \(\kappa>0\).

\begin{theorem}[Calibrated Wilson-to-Casimir symbol limit]
\label{thm:V10-fixed-symbol-limit}
For every fixed irreducible \(R\),
\begin{equation}
 \boxed{
 \lim_{\alpha\to\infty}
 \frac{-\log\eta_R(\alpha)}{s_\alpha}
 =C_2(R),}
 \label{eq:V10-symbol-ratio}
\end{equation}
and therefore
\begin{equation}
 \boxed{a_R(\alpha)\longrightarrow\kappa C_2(R).}
 \label{eq:V10-aR-Casimir}
\end{equation}
Moreover \(s_\alpha\downarrow0\) and, for fixed \(R\), the ratio error is
\(O_R(\alpha^{-1})\).
\end{theorem}

\begin{proof}
The function \(x(U)\) has its unique maximum at the identity.  Laplace
localization therefore puts \(\nu_\alpha\)-mass in an
\(O(\alpha^{-1/2})\) exponential-coordinate neighbourhood of the identity;
the complement is exponentially small.  Conjugation invariance makes the
second moment in the Lie algebra a scalar multiple of the invariant metric.

For \(U=\exp X\), the normalized character has the invariant Taylor
expansion
\begin{equation}
 \frac{\chi_R(e^X)}{d_R}
 =1-\frac{C_2(R)}{2d_G}\norm X^2+O_R(\norm X^4).
 \label{eq:V10-character-Taylor}
\end{equation}
The odd terms integrate to zero by inversion invariance.  Thus, with one
positive scalar \(c_\alpha=O(\alpha^{-1})\) independent of \(R\),
\begin{equation}
 \eta_R(\alpha)
 =1-c_\alpha C_2(R)+O_R(\alpha^{-2}).
 \label{eq:V10-eta-asymptotic}
\end{equation}
For \(R=F\),
\[
 s_\alpha
 =-\frac1{C_F}\log\eta_F(\alpha)
 =c_\alpha+O(\alpha^{-2}).
\]
Taking logarithms in \eqref{eq:V10-eta-asymptotic} and dividing by this
identity gives \eqref{eq:V10-symbol-ratio}, with relative error
\(O_R(\alpha^{-1})\).  In particular \(s_\alpha\to0\) and is positive.
The definition \eqref{eq:V10-time-calibration} then proves
\eqref{eq:V10-aR-Casimir}.
\end{proof}

\begin{assessment}[Why calibration by the fundamental is useful]
No convention-dependent Gaussian variance is inserted by hand.  The
fundamental temporal character defines the physical clock, and every fixed
representation is tested against that same clock.  A different overall
normalization of the Lie-algebra metric changes both \(C_F\) and
\(s_\alpha\) consistently.
\end{assessment}

% =============================================================================
\section{Growing-core symbol defect and representation tail}
\label{sec:V10-core-tail-symbol}
% =============================================================================

On a spatial graph with link representations
\(\mathbf R=(R_e)_e\), define
\[
 C_E(\mathbf R)=\sum_e C_2(R_e),
 \qquad
 A_\alpha(\mathbf R)=\sum_ea_{R_e}(\alpha).
\]
Let \(P_\Lambda\) be the total-Casimir projection used in V3.
Define the one-link core error
\begin{equation}
 \epsilon_\alpha(\Lambda)
 :=\max_{\substack{R\ne\mathbf1\\C_2(R)\le\Lambda}}
 \left|a_R(\alpha)-\kappa C_2(R)\right|
 \label{eq:V10-one-link-error}
\end{equation}
and the exact temporal tail floor
\begin{equation}
 \boxed{
 g_\alpha(\Lambda)
 :=\inf_{\mathbf R:\,C_E(\mathbf R)>\Lambda}
 A_\alpha(\mathbf R).}
 \label{eq:V10-tail-floor}
\end{equation}

\begin{theorem}[Volume-free total-Casimir core error]
\label{thm:V10-core-error}
On the total-Casimir head,
\begin{equation}
 \boxed{
 \norm{
 P_\Lambda(A_\alpha-\kappa\mathcal C_E)P_\Lambda}
 \le
 \frac{\Lambda}{C_F}\,
 \epsilon_\alpha(\Lambda).}
 \label{eq:V10-total-core-error}
\end{equation}
The bound is independent of the number of spatial links.  For every fixed
\(\Lambda\), its right side tends to zero as \(\alpha\to\infty\).
\end{theorem}

\begin{proof}
Every nontrivial \(\mathrm{SU}(3)\) representation has Casimir at least
\(C_F=4/3\).  Hence a configuration with total Casimir at most \(\Lambda\)
contains at most \(\Lambda/C_F\) nontrivial link representations.  On its
Peter--Weyl block,
\[
 \left|
 \sum_ea_{R_e}(\alpha)
 -\kappa\sum_eC_2(R_e)
 \right|
 \le\frac{\Lambda}{C_F}\epsilon_\alpha(\Lambda).
\]
Both operators are diagonal in these blocks.  Taking the supremum proves
\eqref{eq:V10-total-core-error}.  At fixed \(\Lambda\), only finitely many
irreducible labels occur, so
\Cref{thm:V10-fixed-symbol-limit} makes the maximum tend to zero.
\end{proof}

\begin{firewall}[Pointwise symbol convergence is not tail coercivity]
The fixed-\(R\) limit \eqref{eq:V10-aR-Casimir} gives no uniform statement
for representations whose highest weights grow with \(\alpha\).  A smooth
central kernel can have Fourier coefficients whose logarithms grow more
slowly than the quadratic Casimir at very high weight.  The omitted sector
must therefore be controlled by the actual floor
\(g_\alpha(\Lambda)\), not by extrapolating a fixed-representation Taylor
series.
\end{firewall}

\begin{definition}[Admissible temporal symbol schedule]
\label{def:V10-admissible-schedule}
A growing threshold \(\Lambda(\alpha)\to\infty\) is admissible if
\begin{align}
 \frac{\Lambda(\alpha)}{C_F}
 \epsilon_\alpha(\Lambda(\alpha))
 &\longrightarrow0,
 \label{eq:V10-core-schedule}\\
 g_\alpha(\Lambda(\alpha))
 &\longrightarrow\infty.
 \label{eq:V10-tail-schedule}
\end{align}
\end{definition}

\begin{assessment}[The new temporal acquisition]
The character coefficients in \eqref{eq:V10-character-multiplier} are
explicit compact-group integrals.  A rigorous temporal same-generator card
is now the two-scalar schedule
\[
 \boxed{
 \frac{\Lambda}{C_F}\epsilon_\alpha(\Lambda)
 \quad\text{and}\quad
 g_\alpha(\Lambda).}
\]
The first is a maximum over finitely many low representations; the second is
an outward high-representation lower bound.  Neither is a fitted continuum
mass.
\end{assessment}

% =============================================================================
\section{Symmetric transfer and the fixed-volume Hamiltonian limit}
\label{sec:V10-symmetric-transfer}
% =============================================================================

Fix a finite spatial graph.  Let \(V_B\) be the bounded nonnegative spatial
Wilson magnetic multiplication operator and define the anisotropic symmetric
one-step transfer
\begin{equation}
 \mathcal T_\alpha
 :=e^{-\tau_\alpha V_B/2}
   e^{-\tau_\alpha A_\alpha}
   e^{-\tau_\alpha V_B/2}
 =e^{-\tau_\alpha V_B/2}
   \mathcal C_\alpha^{\otimes E}
   e^{-\tau_\alpha V_B/2}.
 \label{eq:V10-symmetric-transfer}
\end{equation}
Let
\[
 H_{\rm raw}=\kappa\mathcal C_E+V_B,
 \qquad
 E_0=\inf\spec H_{\rm raw}.
\]

\begin{theorem}[Mosco--Chernoff same-generator compiler]
\label{thm:V10-same-generator}
Assume an admissible schedule
\(\Lambda(\alpha)\) in the sense of
\Cref{def:V10-admissible-schedule}.  Then, on every fixed finite spatial
graph:
\begin{enumerate}[label=\textup{(\roman*)},leftmargin=3.5em]
\item the positive forms of \(A_\alpha\) converge in the Mosco sense to the
      form of \(\kappa\mathcal C_E\);
\item for every fixed \(t\ge0\),
      \begin{equation}
       \boxed{
       \mathcal T_\alpha^{\lfloor t/\tau_\alpha\rfloor}
       \longrightarrow e^{-tH_{\rm raw}}
       \quad\text{strongly};}
       \label{eq:V10-Chernoff-limit}
      \end{equation}
\item if \(\lambda_\alpha=\norm{\mathcal T_\alpha}\) and
      \(E_{0,\alpha}:=-\tau_\alpha^{-1}\log\lambda_\alpha\), then
      \begin{equation}
       \boxed{E_{0,\alpha}\longrightarrow E_0;}
       \label{eq:V10-ground-energy-limit}
      \end{equation}
\item the normalized reflected transfers satisfy
      \begin{equation}
       \boxed{
       \left(\lambda_\alpha^{-1}\mathcal T_\alpha\right)^{
       \lfloor t/\tau_\alpha\rfloor}
       \longrightarrow
       e^{-t(H_{\rm raw}-E_0)}
       \quad\text{strongly}.}
       \label{eq:V10-normalized-generator-limit}
      \end{equation}
\end{enumerate}
\end{theorem}

\begin{proof}
Peter--Weyl polynomials form a common form core.  On every fixed polynomial,
\Cref{thm:V10-core-error} gives convergence of the \(A_\alpha\) form to
\(\kappa\mathcal C_E\).  The tail condition
\eqref{eq:V10-tail-schedule} makes every sequence of bounded
\(A_\alpha\)-energy tight in a common finite Casimir head.  The core
convergence gives the Mosco recovery sequence; tail tightness gives the
Mosco lower bound.  This proves \textup{(i)}.

The bounded form perturbation \(V_B\) preserves Mosco convergence.  The
symmetric product
\eqref{eq:V10-symmetric-transfer} is Chernoff-equivalent to
\(A_\alpha+V_B\); equivalently, on the common core,
\[
 \tau_\alpha^{-1}(I-\mathcal T_\alpha)f
 \longrightarrow H_{\rm raw}f.
\]
The contraction product theorem, together with \textup{(i)}, proves
\eqref{eq:V10-Chernoff-limit}.

On a fixed compact gauge graph the limiting form has compact resolvent.
Mosco convergence plus the common tail tightness therefore gives convergence
of the lowest eigenvalue.  The lowest logarithmic one-step energy is
\(-\tau_\alpha^{-1}\log\norm{\mathcal T_\alpha}\), proving
\eqref{eq:V10-ground-energy-limit}.  Finally,
\[
 \lambda_\alpha^{-\lfloor t/\tau_\alpha\rfloor}
 \longrightarrow e^{tE_0}.
\]
Multiply this scalar limit by
\eqref{eq:V10-Chernoff-limit} to obtain
\eqref{eq:V10-normalized-generator-limit}.
\end{proof}

\begin{corollary}[Fixed-bank same-generator orbit convergence]
\label{cor:V10-fixed-bank-incidence}
Let \(V\) synthesize a fixed finite Peter--Weyl/Wilson-network bank on a
fixed finite spatial graph.  Under
\Cref{thm:V10-same-generator}, for every fixed \(t\ge0\),
\begin{equation}
 \boxed{
 V^*\left(\lambda_\alpha^{-1}\mathcal T_\alpha\right)^{
 \lfloor t/\tau_\alpha\rfloor}V
 \longrightarrow
 V^*e^{-t(H_{\rm raw}-E_0)}V}
 \label{eq:V10-bank-Gram-limit}
\end{equation}
in matrix norm.
\end{corollary}

\begin{proof}
Strong convergence is uniform on the unit sphere of a fixed
finite-dimensional range.  Compress
\eqref{eq:V10-normalized-generator-limit} by \(V\).
\end{proof}

\begin{assessment}[What same-generator row is now closed]
At fixed spatial regulator, the anisotropic Wilson transfer and the
electric-plus-magnetic differential Hamiltonian have the same continuous
physical-time semigroup once the two temporal symbol stocks in
\Cref{def:V10-admissible-schedule} are proved.  This is a generator
incidence theorem, not merely equality of Perron factors.
\end{assessment}

% =============================================================================
\section{Order of limits and gap preservation}
\label{sec:V10-order-limits}
% =============================================================================

\begin{firewall}[Strong semigroup convergence does not preserve a cofinal gap]
The conclusion \eqref{eq:V10-normalized-generator-limit} is at fixed spatial
volume.  Strong convergence alone permits new soft vectors to escape every
fixed bank as the spatial volume or representation head grows; the V4
growing-head counterexample already demonstrates this mechanism.  A common
mass requires the uniform physical-law, Feshbach, or Witten packet in
addition to same-generator convergence.
\end{firewall}

\begin{theorem}[Two-stage same-generator/mass compiler]
\label{thm:V10-two-stage-compiler}
Suppose:
\begin{enumerate}[label=\textup{(\roman*)},leftmargin=3.5em]
\item at each fixed spatial regulator, the temporal character schedule is
      admissible and \Cref{thm:V10-same-generator} holds;
\item after the temporal limit, the resulting differential Hamiltonians have
      a common physical gap \(m_*>0\) by one of
      \Cref{thm:V5-uniform-compiler,thm:V7-regulator-compiler,
      cor:V9-strong-Hamiltonian-gap};
\item the V1 response-matched local Grams converge quotient-stably and
      reconstruct a nonzero continuum local algebra.
\end{enumerate}
Then the continuum OS Hamiltonian has mass gap at least \(m_*\), with the
declared physical time calibration.
\end{theorem}

\begin{proof}
\Cref{thm:V10-same-generator,cor:V10-fixed-bank-incidence} identify the
finite-regulator physical-time generator before spatial cofinal passage.
Assumption \textup{(ii)} supplies the common orbit floor on every centered
local bank.  Assumption \textup{(iii)} is precisely the V1
quotient-stable, dense-local-span passage, which transports the common floor
to the reconstructed continuum Hamiltonian.  Nonzero reconstruction
excludes the vacuous zero-Hilbert-space limit.
\end{proof}

% =============================================================================
\section{V10 ledger and next attack}
\label{sec:V10-ledger}
% =============================================================================

\begin{theorem}[Integrated V10 statement]
\label{thm:V10-integrated}
Preserving V1--V9, V10 proves:
\begin{enumerate}[label=\textup{(V10.\arabic*)},leftmargin=5.5em]
\item exact Peter--Weyl multipliers of the temporal Wilson convolution;
\item a fundamental-character physical clock and the fixed-representation
      Casimir-symbol limit;
\item a volume-free growing-core error bound in total Casimir;
\item separation of core accuracy from the actual high-representation
      temporal floor;
\item a Mosco--Chernoff theorem identifying the normalized anisotropic
      transfer semigroup with the differential Hamiltonian at fixed spatial
      regulator;
\item matrix-norm same-generator convergence on every fixed local bank;
\item a two-stage compiler separating temporal incidence from the uniform
      spatial/continuum mass estimate.
\end{enumerate}
\end{theorem}

\begin{proof}
Items \textup{(V10.1)--(V10.4)} are
\Cref{lem:V10-Peter-Weyl-diagonal,thm:V10-fixed-symbol-limit,
thm:V10-core-error,def:V10-admissible-schedule}.  Items
\textup{(V10.5)--(V10.6)} are
\Cref{thm:V10-same-generator,cor:V10-fixed-bank-incidence}.  Item
\textup{(V10.7)} is \Cref{thm:V10-two-stage-compiler}.
\end{proof}

\begingroup
\small
\begin{longtable}{>{\raggedright\arraybackslash}p{0.28\textwidth}
                  >{\raggedright\arraybackslash}p{0.16\textwidth}
                  >{\raggedright\arraybackslash}p{0.48\textwidth}}
\toprule
\textbf{V10 row}&\textbf{Status}&\textbf{Advance or remaining obstruction}\\
\midrule
\endfirsthead
\toprule
\textbf{V10 row}&\textbf{Status}&\textbf{Advance or remaining obstruction}\\
\midrule
\endhead
Temporal Wilson character symbol
& \MGpass
& Exact compact-group integral and fixed-representation Casimir limit.\\[0.35em]
Growing-core symbol error
& \MGpass
& Explicit finite maximum with a volume-free total-Casimir multiplier.\\[0.35em]
High-representation temporal floor
& \MGopen
& Need an outward bound proving an admissible growing schedule.\\[0.35em]
Fixed-volume same generator
& Conditional exact theorem
& Follows from the two temporal symbol stocks by Mosco--Chernoff convergence.\\[0.35em]
Uniform spatial mass on weak branch
& \MGopen
& V9 closes only the small-\(\beta\) strong-coupling ground-law branch.\\[0.35em]
Nontrivial continuum OS limit
& \MGopen
& Requires V1 quotient-stable convergence after temporal incidence.\\[0.35em]
Full Yang--Mills mass gap
& \MGopen
& Needs the temporal tail, weak-coupling physical one-form floor, and continuum construction.\\
\bottomrule
\end{longtable}
\endgroup

\begin{openproblem}[V11 mandate: temporal tail and weak-chart return]
\label{op:V11}
Proceed in the order exposed by V10.
\begin{enumerate}[label=\textup{(V11.\arabic*)},leftmargin=5.5em]
\item derive outward lower bounds for the Wilson character ratios
      \(-\log\eta_{p,q}(\alpha)\) in the joint large-\(\alpha\),
      large-highest-weight region;
\item construct an explicit growing \(\Lambda(\alpha)\) satisfying
      \eqref{eq:V10-core-schedule}--\eqref{eq:V10-tail-schedule}, or return
      the representation sequence that prevents it;
\item insert the resulting temporal schedule before the spatial
      thermodynamic/continuum passage;
\item on the weak spatial branch, use a gauge-fixed small-field chart for the
      positive comparison Witten operator and route the compact complement
      through the existing V35 curvature-defect return;
\item prove a uniform weighted edge below one on the complete exact-form
      range, with toron, wrapping, and chart-exit sectors separately shorted;
\item only then apply \Cref{thm:V10-two-stage-compiler}.
\end{enumerate}
If a quadratic-Casimir lower bound is false at extremely high weight, retain
the actual logarithmic character floor \(g_\alpha(\Lambda)\); only divergence
of the omitted-sector energy is required for temporal Mosco convergence.
\end{openproblem}

\begin{firewall}[End-of-V10 claim boundary]
V10 proves the fixed-representation temporal Casimir limit and a conditional
same-generator theorem based on explicit core and tail symbol stocks.  It
does not yet prove an admissible growing temporal schedule, a weak-coupling
uniform spatial mass, or a nontrivial continuum OS construction.  The full
Yang--Mills mass gap remains open.
\end{firewall}

\section*{Version V10 append-only record}
\addcontentsline{toc}{section}{Version V10 append-only record}
The complete V9 source was retained before this continuation.  Its SHA--256
digest was
\begin{center}\scriptsize\ttfamily
1b5edbde5399b4c86aeb2cd9e2156c206d9f5718d1ea9924929caf7e1b7a202c.
\end{center}
On the next iteration, retain this full file, append before its unique active
document-closing command, increment the filename to
\texttt{\_V11.tex}, and remove only the superseded V10 file from this note
series.

% =============================================================================
% END APPENDED VERSION V10 MATERIAL
% =============================================================================

% =============================================================================
% BEGIN APPENDED VERSION V11 MATERIAL
% =============================================================================

\clearpage
\part{Version V11: heat-kernel domination and closure of the temporal symbol schedule}
\label{part:V11}

\section{Direction}
\label{sec:V11-direction}

V10 reduced temporal same-generator incidence to a growing-core symbol error
and a high-representation floor.  The fixed-representation error can always
be made small on a sufficiently slowly growing core.  The nontrivial row is
the tail.

The tail can be controlled without a uniform large-highest-weight stationary
phase expansion.  The normalized temporal Wilson density dominates a
short-time heat-kernel density pointwise.  Since the representation
Dirichlet integrand is nonnegative, this turns the exact heat multiplier
\(e^{-tC_2(R)}\) into an outward Wilson character bound.  It yields
\[
 a_R(\alpha)\gtrsim\kappa\min\{C_2(R),\alpha\},
\]
which is more than the divergence required by V10.

% =============================================================================
\section{A pointwise Wilson/heat-kernel comparison}
\label{sec:V11-density-comparison}
% =============================================================================

Let \(d(U,I)\) be the Riemannian distance on \(\mathrm{SU}(3)\), for the
bi-invariant metric used by the electric Casimir, and retain
\[
 x(U)=\frac13\operatorname{ReTr}U.
\]
Write the normalized temporal Wilson density as
\begin{equation}
 q_\alpha(U)
 :=\frac{e^{-\alpha(1-x(U))}}
         {\widetilde Z_\alpha},
 \qquad
 \widetilde Z_\alpha
 :=\int_G e^{-\alpha(1-x(U))}\,\dd U.
 \label{eq:V11-qalpha}
\end{equation}
Let \(p_t(U)\) be the heat kernel of \(-\Delta_G\), normalized as a
probability density, so its \(R\)-Fourier multiplier is
\(e^{-tC_2(R)}\).

\begin{lemma}[Global quadratic control of the Wilson well]
\label{lem:V11-quadratic-well}
There are constants \(0<k_-\le k_+<\infty\) such that
\begin{equation}
 \boxed{
 k_-d(U,I)^2
 \le1-x(U)
 \le k_+d(U,I)^2
 \qquad(U\in G).}
 \label{eq:V11-global-quadratic}
\end{equation}
Consequently, for some \(0<c_Z\le C_Z<\infty\) and all sufficiently large
\(\alpha\),
\begin{equation}
 \boxed{
 c_Z\alpha^{-4}
 \le\widetilde Z_\alpha
 \le C_Z\alpha^{-4}.}
 \label{eq:V11-Z-Laplace}
\end{equation}
\end{lemma}

\begin{proof}
The function \(1-x\) is nonnegative and vanishes only at the identity.
In exponential coordinates its Hessian at the identity is a positive scalar
multiple of the invariant metric, so
\((1-x(U))/d(U,I)^2\) extends to a positive continuous function at the
identity.  Away from the identity, compactness and uniqueness of the zero
give positive finite extrema.  This proves
\eqref{eq:V11-global-quadratic}.

The real dimension of \(G\) is eight.  Integrating the two Gaussian bounds in
normal coordinates on a fixed injectivity ball gives constants times
\(\alpha^{-8/2}=\alpha^{-4}\).  The complement of that ball contributes an
exponentially small amount to the upper bound.  This proves
\eqref{eq:V11-Z-Laplace}.
\end{proof}

\begin{lemma}[Short-time heat-kernel upper bound]
\label{lem:V11-heat-upper}
There are constants \(C_H,c_H,t_0>0\) such that
\begin{equation}
 p_t(U)
 \le C_Ht^{-4}
 \exp\left(-\frac{d(U,I)^2}{c_Ht}\right)
 \qquad(0<t\le t_0,\ U\in G).
 \label{eq:V11-heat-upper}
\end{equation}
\end{lemma}

\begin{proof}
On a compact Riemannian manifold, the heat kernel in a normal neighbourhood
has the Gaussian parametrix
\[
 (4\pi t)^{-4}e^{-d(U,I)^2/(4t)}
 \bigl(j(U)^{-1/2}+O(t)\bigr).
\]
Choose a smaller normal ball so the Jacobian and remainder are uniformly
bounded.  On its complement, the semigroup property together with the local
parametrix gives an exponentially small off-diagonal bound.  Enlarging the
prefactor and denominator constant yields
\eqref{eq:V11-heat-upper} uniformly on the compact group.
\end{proof}

\begin{theorem}[Temporal Wilson density dominates a heat kernel]
\label{thm:V11-heat-domination}
There exist constants \(a_0,c_0>0\) and
\(\alpha_0<\infty\), depending only on the chosen group metric, such that
\begin{equation}
 \boxed{
 q_\alpha(U)\ge c_0p_{a_0/\alpha}(U)
 \qquad
 (\alpha\ge\alpha_0,\ U\in G).}
 \label{eq:V11-density-domination}
\end{equation}
\end{theorem}

\begin{proof}
By \Cref{lem:V11-quadratic-well},
\[
 q_\alpha(U)
 \ge C_Z^{-1}\alpha^4
     e^{-k_+\alpha d(U,I)^2}.
\]
Choose \(a_0>0\) so small that
\(
 a_0/\alpha\le t_0
\)
for \(\alpha\ge\alpha_0\) and
\(
 (c_Ha_0)^{-1}\ge k_+
\).
Then \Cref{lem:V11-heat-upper} gives
\[
 p_{a_0/\alpha}(U)
 \le C_Ha_0^{-4}\alpha^4
     e^{-k_+\alpha d(U,I)^2}.
\]
Comparison of the two displays proves
\eqref{eq:V11-density-domination} with
\(
 c_0=a_0^4/(C_ZC_H)
\),
decreased if necessary to be at most one.
\end{proof}

% =============================================================================
\section{A uniform high-representation character floor}
\label{sec:V11-character-floor}
% =============================================================================

\begin{lemma}[Representation Dirichlet identity]
\label{lem:V11-representation-Dirichlet}
For every unitary irreducible representation \(R\),
\begin{equation}
 \boxed{
 1-\eta_R(\alpha)
 =\frac1{2d_R}
   \int_G\norm{D^R(U)-I}_{\rm HS}^2
   q_\alpha(U)\,\dd U.}
 \label{eq:V11-Wilson-Dirichlet}
\end{equation}
For the heat kernel,
\begin{equation}
 \boxed{
 \frac1{2d_R}
 \int_G\norm{D^R(U)-I}_{\rm HS}^2
 p_t(U)\,\dd U
 =1-e^{-tC_2(R)}.}
 \label{eq:V11-heat-Dirichlet}
\end{equation}
\end{lemma}

\begin{proof}
Unitarity gives
\[
 \norm{D^R(U)-I}_{\rm HS}^2
 =2d_R-2\operatorname{Re}\chi_R(U).
\]
Both \(q_\alpha\) and \(p_t\) are invariant under inversion, so their
character integrals are real.  Integration gives
\eqref{eq:V11-Wilson-Dirichlet}.  The heat convolution acts on \(R\) by
\(e^{-tC_2(R)}\), proving
\eqref{eq:V11-heat-Dirichlet}.
\end{proof}

\begin{theorem}[Uniform Wilson character-tail bound]
\label{thm:V11-character-tail}
For every irreducible \(R\) and every \(\alpha\ge\alpha_0\),
\begin{align}
 1-\eta_R(\alpha)
 &\ge c_0\left(
 1-e^{-a_0C_2(R)/\alpha}\right),
 \label{eq:V11-one-minus-eta}\\
 -\log\eta_R(\alpha)
 &\ge c_0\left(
 1-e^{-a_0C_2(R)/\alpha}\right).
 \label{eq:V11-log-eta-floor}
\end{align}
There is a constant \(c_*>0\), independent of \(R\) and \(\alpha\), such
that the calibrated V10 temporal symbol satisfies
\begin{equation}
 \boxed{
 a_R(\alpha)
 \ge c_*\kappa
 \min\{C_2(R),\alpha\}.}
 \label{eq:V11-aR-floor}
\end{equation}
\end{theorem}

\begin{proof}
The integrand in \eqref{eq:V11-Wilson-Dirichlet} is nonnegative.  Insert
\eqref{eq:V11-density-domination} and use
\eqref{eq:V11-heat-Dirichlet} to obtain
\eqref{eq:V11-one-minus-eta}.  Since
\(-\log u\ge1-u\) for \(0<u\le1\),
\eqref{eq:V11-log-eta-floor} follows.

V10 gives \(s_\alpha=O(\alpha^{-1})\).  Thus, after enlarging
\(\alpha_0\), there is \(C_s<\infty\) such that
\(
 s_\alpha\le C_s/\alpha
\).
For \(z\ge0\),
\[
 1-e^{-z}\ge(1-e^{-1})\min\{z,1\}.
\]
Taking \(z=a_0C_2(R)/\alpha\), decreasing \(a_0\) to at most one if
necessary, and using
\(
 a_R=\kappa[-\log\eta_R]/s_\alpha
\)
gives \eqref{eq:V11-aR-floor} with
\[
 c_*=\frac{c_0a_0(1-e^{-1})}{C_s}.
\]
\end{proof}

\begin{corollary}[Explicit total-representation tail floor]
\label{cor:V11-total-tail}
For every \(0<\Lambda\le\alpha\),
\begin{equation}
 \boxed{
 g_\alpha(\Lambda)\ge c_*\kappa\Lambda.}
 \label{eq:V11-galpha-floor}
\end{equation}
\end{corollary}

\begin{proof}
For nonnegative numbers \(x_e\),
\[
 \sum_e\min\{x_e,\alpha\}
 \ge\min\left\{\sum_ex_e,\alpha\right\}.
\]
If \(\sum_eC_2(R_e)>\Lambda\) and \(\Lambda\le\alpha\), the right side is at
least \(\Lambda\).  Sum \eqref{eq:V11-aR-floor} over links and take the
infimum in \eqref{eq:V10-tail-floor}.
\end{proof}

% =============================================================================
\section{Existence of an admissible growing temporal schedule}
\label{sec:V11-admissible-schedule}
% =============================================================================

\begin{theorem}[Diagonal construction of the V10 schedule]
\label{thm:V11-diagonal-schedule}
There exists a nondecreasing threshold
\(\Lambda(\alpha)\to\infty\), with
\(\Lambda(\alpha)\le\alpha\) for all sufficiently large \(\alpha\), such
that
\begin{align}
 \frac{\Lambda(\alpha)}{C_F}
 \epsilon_\alpha(\Lambda(\alpha))
 &\longrightarrow0,
 \label{eq:V11-diagonal-core}\\
 g_\alpha(\Lambda(\alpha))
 &\longrightarrow\infty.
 \label{eq:V11-diagonal-tail}
\end{align}
Hence an admissible temporal symbol schedule in the sense of V10 exists.
\end{theorem}

\begin{proof}
For each integer \(n\ge1\), the set of irreducible representations with
\(C_2(R)\le n\) is finite.  By
\Cref{thm:V10-fixed-symbol-limit},
\(\epsilon_\alpha(n)\to0\).  Choose an increasing sequence
\(A_n\to\infty\) such that
\begin{equation}
 A_n\ge\max\{\alpha_0,n,A_{n-1}+1\},
 \qquad
 \epsilon_\alpha(n)\le\frac{C_F}{n^2}
 \quad(\alpha\ge A_n).
 \label{eq:V11-An-choice}
\end{equation}
For \(A_n\le\alpha<A_{n+1}\), set
\(\Lambda(\alpha)=n\).  Then
\(\Lambda(\alpha)\to\infty\), \(\Lambda(\alpha)\le\alpha\), and
\[
 \frac{\Lambda(\alpha)}{C_F}
 \epsilon_\alpha(\Lambda(\alpha))
 \le\frac1n\longrightarrow0.
\]
\Cref{cor:V11-total-tail} gives
\(
 g_\alpha(\Lambda(\alpha))
 \ge c_*\kappa n\to\infty
\).
\end{proof}

\begin{assessment}[Temporal obstruction closed]
V10's two temporal stocks are now populated on a common growing schedule.
No quadratic-Casimir estimate was imposed in the extreme representation
tail; the actual Wilson logarithmic symbol is bounded below by the weaker but
sufficient \(\min\{C_2(R),\alpha\}\) scale.
\end{assessment}

% =============================================================================
\section{Unconditional fixed-volume same-generator consequence}
\label{sec:V11-same-generator}
% =============================================================================

\begin{theorem}[Temporal Wilson transfer has the differential Hamiltonian limit]
\label{thm:V11-same-generator}
On every fixed finite spatial gauge graph, use the fundamental temporal
calibration
\(\tau_\alpha=s_\alpha/\kappa\) and the symmetric anisotropic Wilson transfer
\(\mathcal T_\alpha\) of V10.  Then
\begin{equation}
 \boxed{
 \left(\lambda_\alpha^{-1}\mathcal T_\alpha\right)^{
 \lfloor t/\tau_\alpha\rfloor}
 \longrightarrow
 e^{-t(\kappa\mathcal C_E+V_B-E_0)}
 \quad\text{strongly}}
 \label{eq:V11-unconditional-same-generator}
\end{equation}
for every \(t\ge0\).  Compression to every fixed finite local
Peter--Weyl/Wilson-network bank converges in matrix norm.
\end{theorem}

\begin{proof}
\Cref{thm:V11-diagonal-schedule} supplies the admissible schedule required by
\Cref{thm:V10-same-generator}.  Apply that theorem and then
\Cref{cor:V10-fixed-bank-incidence}.
\end{proof}

\begin{firewall}[Order of limits remains essential]
The theorem first removes the temporal lattice spacing at fixed spatial
graph.  Its constants do not assert a spatial-volume gap.  The subsequent
thermodynamic and ultraviolet passage must retain one uniform physical mass
packet; otherwise soft vectors may escape the fixed banks despite
\eqref{eq:V11-unconditional-same-generator}.
\end{firewall}

% =============================================================================
\section{V11 ledger and revised frontier}
\label{sec:V11-ledger}
% =============================================================================

\begin{theorem}[Integrated V11 statement]
\label{thm:V11-integrated}
Preserving V1--V10, V11 proves:
\begin{enumerate}[label=\textup{(V11.\arabic*)},leftmargin=5.5em]
\item a pointwise domination of the temporal Wilson density by a heat kernel
      at time \(a_0/\alpha\);
\item a representation-uniform lower bound for
      \(1-\eta_R(\alpha)\) and \(-\log\eta_R(\alpha)\);
\item the calibrated symbol floor
      \(a_R(\alpha)\ge c_*\kappa\min\{C_2(R),\alpha\}\);
\item a volume-free total-representation tail bound;
\item existence of a growing threshold satisfying both V10 temporal symbol
      conditions;
\item unconditional fixed-spatial-volume convergence of the normalized
      temporal Wilson transfer to the differential gauge Hamiltonian
      semigroup.
\end{enumerate}
\end{theorem}

\begin{proof}
Items \textup{(V11.1)--(V11.2)} are
\Cref{thm:V11-heat-domination,thm:V11-character-tail}.  Items
\textup{(V11.3)--(V11.4)} are
\Cref{thm:V11-character-tail,cor:V11-total-tail}.  Item
\textup{(V11.5)} is \Cref{thm:V11-diagonal-schedule}; item
\textup{(V11.6)} is \Cref{thm:V11-same-generator}.
\end{proof}

\begingroup
\small
\begin{longtable}{>{\raggedright\arraybackslash}p{0.28\textwidth}
                  >{\raggedright\arraybackslash}p{0.16\textwidth}
                  >{\raggedright\arraybackslash}p{0.48\textwidth}}
\toprule
\textbf{V11 row}&\textbf{Status}&\textbf{Advance or remaining obstruction}\\
\midrule
\endfirsthead
\toprule
\textbf{V11 row}&\textbf{Status}&\textbf{Advance or remaining obstruction}\\
\midrule
\endhead
High-representation temporal Wilson floor
& \MGpass
& Heat-kernel domination gives a representation-uniform outward bound.\\[0.35em]
Admissible growing temporal schedule
& \MGpass
& Constructed by a diagonal core schedule plus the explicit tail floor.\\[0.35em]
Fixed-volume temporal same generator
& \MGpass
& Normalized transfer powers converge to the differential Hamiltonian semigroup.\\[0.35em]
Spatial-volume weak-coupling mass
& \MGopen
& Need a common exact-one-form or martingale margin on the continuum trajectory.\\[0.35em]
Toron, wrapping, and chart-exit sectors
& \MGopen
& Must be shorted separately in the weak spatial comparison.\\[0.35em]
Nontrivial continuum OS state
& \MGopen
& V1 quotient-stable convergence and local-algebra noncollapse remain.\\[0.35em]
Full Yang--Mills mass gap
& \MGopen
& Temporal incidence is closed; the remaining obstruction is spatial/infrared continuum control.\\
\bottomrule
\end{longtable}
\endgroup

\begin{openproblem}[V12 mandate: weak-coupling physical one-form floor]
\label{op:V12}
The temporal generator row is now closed.  Work after the temporal limit on
the actual continuum-scaling spatial regulator.
\begin{enumerate}[label=\textup{(V12.\arabic*)},leftmargin=5.5em]
\item choose a gauge-fixed small-field chart with an explicit physical radius
      and retain the exact ground-law Witten form from V8;
\item prove a positive comparison one-form floor in physical units on that
      chart;
\item insert the complete negative-curvature and chart-exit return through
      the source-cyclic Woodbury channel already developed in the supplied
      V35 programme;
\item obtain a weighted edge strictly below one on the complete horizontal
      exact-form range, not only a fixed axial source card;
\item short toron and wrapping sectors using the V1 local-versus-wrapping
      classifier, and return any near-unit source-cyclic vector explicitly;
\item combine the uniform spatial floor with
      \Cref{thm:V11-same-generator,thm:V10-two-stage-compiler} and construct
      the nontrivial OS limit.
\end{enumerate}
If the physical weak chart cannot retain a uniform radius, classify the exit
sequence before weakening the comparison floor; a chart-collapse vector and
a genuine local soft mode are different obstructions.
\end{openproblem}

\begin{firewall}[End-of-V11 claim boundary]
V11 closes the high-representation temporal symbol floor and the
fixed-volume transfer-to-Hamiltonian generator incidence.  It does not prove
the weak-coupling spatial/infrared one-form floor, control topological and
wrapping sectors on the continuum trajectory, or construct a nontrivial
continuum OS state.  The full Yang--Mills mass gap remains open.
\end{firewall}

\section*{Version V11 append-only record}
\addcontentsline{toc}{section}{Version V11 append-only record}
The complete V10 source was retained before this continuation.  Its SHA--256
digest was
\begin{center}\scriptsize\ttfamily
b1a642b58531c0d914dd6a181f5da25d22e62cc918b8170472a25294a274f129.
\end{center}
On the next iteration, retain this full file, append before its unique active
document-closing command, increment the filename to
\texttt{\_V12.tex}, and remove only the superseded V11 file from this note
series.

% =============================================================================
% END APPENDED VERSION V11 MATERIAL
% =============================================================================

% =============================================================================
% BEGIN APPENDED VERSION V12 MATERIAL
% =============================================================================

\clearpage
\part{Version V12: an exact chart-localized Witten compiler and the exterior-capacity obstruction}
\label{part:V12}

\section{Direction and nonduplication audit}
\label{sec:V12-direction}

V11 closed the temporal high-representation symbol schedule and identified the
anisotropic Wilson transfer with the differential Hamiltonian at every fixed
spatial regulator.  The remaining finite-regulator problem is therefore
spatial: prove a lower bound for the physical ground-law one-form operator
\(\mathscr L_1\) which is uniform in physical volume and cutoff.

The supplied V35 reduced-fibre programme already proves the algebraic
comparison-minus-return identity, its Woodbury landing, source-cyclic return
expansion, and cutoff-independent locality of the positive comparison inverse.
The Grand Tensor manuscript already proves that a norm-small deformation of
the massless Gaussian infrared channel cannot create a uniform mass.  Neither
result is repeated here.

The missing assembly step is subtler than placing the V35 comparison estimate
inside a small-field chart.  Multiplication by a chart cutoff does not preserve
the exact one-form range, and a small probability of leaving the chart is not
an operator lower bound.  V12 proves an exact IMS compiler which keeps both
defects visible.  It reduces the spatial mass question to five typed stocks:
\[
 \text{inner comparison floor},\quad
 \text{complete return edge},\quad
 \text{exterior Dirichlet floor},\quad
 \text{IMS cost},\quad
 \text{exact-range localization defect}.
\]

% =============================================================================
\section{The physical one-form form and its IMS identity}
\label{sec:V12-IMS}
% =============================================================================

Fix one finite spatial gauge regulator.  On a protected smooth horizontal
stratum of its gauge quotient, write
\[
 \dd\pi=Z^{-1}e^{-W}\dd\mathfrak m,
 \qquad
 D=\sqrt\kappa\,d,
 \qquad
 \mathcal E=\overline{\Ran D}.
\]
Let
\begin{equation}
 \mathscr W
 :=\kappa\bigl(dd_\pi^*+d_\pi^*d\bigr)
 \label{eq:V12-full-Witten}
\end{equation}
be the Friedrichs realization on horizontal one-forms, and let
\(\mathfrak q\) be its closed quadratic form.  The exact-range restriction of
\(\mathscr W\) is \(DD^*\), hence is the operator \(\mathscr L_1\) of V8.

\begin{lemma}[Weighted one-form IMS identity]
\label{lem:V12-Witten-IMS}
Let \(\chi_0,\chi_1\) be real Lipschitz functions with
\(\chi_0^2+\chi_1^2=1\).  Then every
\(\omega\in\Dom\mathfrak q\) satisfies
\begin{equation}
 \boxed{
 \mathfrak q(\omega)
 =
 \sum_{i=0}^1\mathfrak q(\chi_i\omega)
 -
 \kappa\int
 \left(|d\chi_0|^2+|d\chi_1|^2\right)
 |\omega|^2\,\dd\pi.}
 \label{eq:V12-Witten-IMS}
\end{equation}
Consequently, with
\begin{equation}
 \varepsilon_{\rm IMS}
 :=
 \kappa
 \left\|
 |d\chi_0|^2+|d\chi_1|^2
 \right\|_\infty,
 \label{eq:V12-IMS-cost}
\end{equation}
the last term in \eqref{eq:V12-Witten-IMS} has absolute value at most
\(\varepsilon_{\rm IMS}\norm{\omega}^2\).
\end{lemma}

\begin{proof}
On the smooth core, the weighted Bochner identity used in V8 writes
\[
 \mathfrak q(\omega)
 =
 \kappa\int
 \left(
   |\nabla\omega|^2
   +\ip{\omega}{
      (\Ric+\nabla^2W)\omega}
 \right)\,\dd\pi.
\]
The zero-order term splits exactly because
\(\chi_0^2+\chi_1^2=1\).  For the connection term,
\[
 \nabla(\chi_i\omega)
 =d\chi_i\otimes\omega+\chi_i\nabla\omega.
\]
After summing over \(i\), the cross term vanishes since
\[
 \sum_i\chi_i\,d\chi_i
 =\frac12d\!\left(\sum_i\chi_i^2\right)=0,
\]
and the remaining cutoff term is
\(\sum_i|d\chi_i|^2|\omega|^2\).  Rearrangement proves
\eqref{eq:V12-Witten-IMS}.  Approximation by smooth partitions and closedness
of \(\mathfrak q\) extend the identity to Lipschitz cutoffs and the full form
domain.
\end{proof}

\begin{lemma}[Exact forms are not stable under chart localization]
\label{lem:V12-exact-localization-defect}
If \(\omega\) is a smooth exact one-form, then
\begin{equation}
 d(\chi\omega)=d\chi\wedge\omega.
 \label{eq:V12-cutoff-curl}
\end{equation}
Thus \(\chi\omega\) need not be closed, and hence need not lie in
\(\mathcal E\).  An estimate proved only on the global exact range cannot be
inserted into \eqref{eq:V12-Witten-IMS} without either:
\begin{enumerate}[label=\textup{(\roman*)},leftmargin=3.6em]
\item proving it on every cutoff-localized horizontal one-form; or
\item bounding an explicit localization/router defect.
\end{enumerate}
\end{lemma}

\begin{proof}
Since \(d\omega=0\),
\[
 d(\chi\omega)=d\chi\wedge\omega+\chi\,d\omega
 =d\chi\wedge\omega.
\]
The right side is generally nonzero.  Every exact one-form is closed, proving
the assertion.
\end{proof}

\begin{firewall}[No silent exact-range localization]
\label{fire:V12-no-silent-localization}
The exact-form Witten isometry in V6 and in the supplied V35 programme is a
global Hilbert-space statement.  It does not imply that a chart cutoff of an
exact current remains in the source carrier.  The defect in
\eqref{eq:V12-cutoff-curl} must be included before a small-field comparison
can be promoted to the complete physical one-form range.
\end{firewall}

% =============================================================================
\section{The two-region physical Witten compiler}
\label{sec:V12-two-region}
% =============================================================================

The next theorem is stated directly on cutoff-localized vectors from the
exact physical range.  This allows either a full local one-form estimate or a
local exact-range router to populate its hypotheses.

\begin{theorem}[Exact two-region Witten compiler]
\label{thm:V12-two-region-compiler}
Let \(\chi_0^2+\chi_1^2=1\) be as above.  Suppose there are
\[
 r_0>0,\qquad r_1>0,\qquad
 \delta_0,\delta_1\ge0
\]
such that every \(\omega\in\mathcal E\cap\Dom\mathfrak q\) obeys
\begin{align}
 \mathfrak q(\chi_0\omega)
 &\ge
 r_0\norm{\chi_0\omega}^2
 -\delta_0\norm{\omega}^2,
 \label{eq:V12-inner-local-floor}\\
 \mathfrak q(\chi_1\omega)
 &\ge
 r_1\norm{\chi_1\omega}^2
 -\delta_1\norm{\omega}^2.
 \label{eq:V12-outer-local-floor}
\end{align}
Then
\begin{equation}
 \boxed{
 \mathscr L_1
 \succeq
 \left(
   \min\{r_0,r_1\}
   -\varepsilon_{\rm IMS}
   -\delta_0-\delta_1
 \right)I
 \quad\text{on }\mathcal E.}
 \label{eq:V12-two-region-floor}
\end{equation}
In particular, a strictly positive quantity in parentheses is a lower bound
for the Hamiltonian gap.
\end{theorem}

\begin{proof}
Apply \eqref{eq:V12-inner-local-floor} and
\eqref{eq:V12-outer-local-floor} to the IMS identity.  Since
\[
 \norm{\chi_0\omega}^2+\norm{\chi_1\omega}^2
 =\norm{\omega}^2,
\]
one obtains
\[
 \mathfrak q(\omega)
 \ge
 \left(
 \min\{r_0,r_1\}
 -\varepsilon_{\rm IMS}
 -\delta_0-\delta_1
 \right)\norm{\omega}^2.
\]
The form characterization of operator order proves
\eqref{eq:V12-two-region-floor}.  The last assertion is
\Cref{cor:V6-one-form-incidence}.
\end{proof}

\begin{corollary}[Concrete chart-width cost]
\label{cor:V12-chart-width}
Let \(r\) be a \(1\)-Lipschitz small-field radius, and let \(R,w>0\).
Choose
\[
 \theta(x)=
 \begin{cases}
 0,&r(x)\le R,\\
 \dfrac{\pi}{2w}\bigl(r(x)-R\bigr),
    &R<r(x)<R+w,\\
 \dfrac\pi2,&r(x)\ge R+w,
 \end{cases}
 \qquad
 \chi_0=\cos\theta,\quad
 \chi_1=\sin\theta.
\]
Then
\begin{equation}
 \boxed{
 \varepsilon_{\rm IMS}
 \le\frac{\kappa\pi^2}{4w^2}.}
 \label{eq:V12-width-cost}
\end{equation}
The inner vector is supported in \(\{r<R+w\}\), the exterior vector in
\(\{r>R\}\), and the overlap is precisely the transition annulus.
\end{corollary}

\begin{proof}
Almost everywhere,
\[
 |d\chi_0|^2+|d\chi_1|^2
 =|d\theta|^2
 \le\frac{\pi^2}{4w^2}|dr|^2
 \le\frac{\pi^2}{4w^2}.
\]
Insert this in \eqref{eq:V12-IMS-cost}.  The support statements follow from
the definition of \(\theta\).
\end{proof}

% =============================================================================
\section{Landing the V35 Woodbury return in the inner stock}
\label{sec:V12-Woodbury-landing}
% =============================================================================

\begin{proposition}[Comparison-minus-return population of the inner floor]
\label{prop:V12-inner-Woodbury}
Suppose the physical inner form has, for every
\(\omega\in\mathcal E\cap\Dom\mathfrak q\), the lower representation
\begin{equation}
 \mathfrak q(\chi_0\omega)
 \ge
 \ip{\chi_0\omega}{\mathcal A\chi_0\omega}
 -
 \left\|
 (I-\mathcal K)^{-1/2}
 \mathcal X\chi_0\omega
 \right\|^2
 -
 \delta_{\rm loc}\norm{\omega}^2,
 \label{eq:V12-inner-card}
\end{equation}
where
\[
 \mathcal A\succeq aI,\qquad a>0,
 \qquad
 0\preceq\mathcal K\prec I.
\]
Define the complete weighted return edge
\begin{equation}
 \rho
 :=
 \left\|
 (I-\mathcal K)^{-1/2}
 \mathcal X\mathcal A^{-1/2}
 \right\|^2.
 \label{eq:V12-return-edge}
\end{equation}
If \(\rho<1\), then
\begin{equation}
 \boxed{
 \mathfrak q(\chi_0\omega)
 \ge
 a(1-\rho)\norm{\chi_0\omega}^2
 -\delta_{\rm loc}\norm{\omega}^2.}
 \label{eq:V12-inner-populated}
\end{equation}
Thus \eqref{eq:V12-inner-local-floor} holds with
\[
 r_0=a(1-\rho),
 \qquad
 \delta_0=\delta_{\rm loc}.
\]
\end{proposition}

\begin{proof}
By \eqref{eq:V12-return-edge},
\[
 \left\|
 (I-\mathcal K)^{-1/2}
 \mathcal X\eta
 \right\|^2
 \le
 \rho\ip{\eta}{\mathcal A\eta}
\]
for every \(\eta\) in the form domain of \(\mathcal A\).  Take
\(\eta=\chi_0\omega\) in \eqref{eq:V12-inner-card}.  Then
\[
 \mathfrak q(\chi_0\omega)
 \ge
 (1-\rho)
 \ip{\chi_0\omega}{\mathcal A\chi_0\omega}
 -\delta_{\rm loc}\norm{\omega}^2,
\]
and \(\mathcal A\succeq aI\) proves
\eqref{eq:V12-inner-populated}.
\end{proof}

\begin{theorem}[Five-stock spatial mass certificate]
\label{thm:V12-five-stock}
At each regulator \(j\), suppose:
\begin{enumerate}[label=\textup{(S\arabic*)},leftmargin=4.2em]
\item the physical ground-law one-form form is the V8 form
      \(\mathfrak q_j\) on
      \(\mathcal E_j=\overline{\Ran D_j}\);
\item the inner comparison/return card satisfies
      \Cref{prop:V12-inner-Woodbury} with stocks
      \(a_j>0\), \(0\le\rho_j<1\), and
      \(\delta_{{\rm loc},j}\ge0\);
\item the exterior cutoff obeys
      \[
      \mathfrak q_j(\chi_{1,j}\omega)
      \ge
      b_j\norm{\chi_{1,j}\omega}^2
      -\delta_{{\rm ext},j}\norm{\omega}^2
      \qquad(b_j>0);
      \]
\item the partition has IMS cost
      \(\varepsilon_{{\rm IMS},j}\).
\end{enumerate}
Put
\begin{equation}
 \boxed{
 m_j:=
 \min\left\{
   a_j(1-\rho_j),\,b_j
 \right\}
 -
 \varepsilon_{{\rm IMS},j}
 -
 \delta_{{\rm loc},j}
 -
 \delta_{{\rm ext},j}.}
 \label{eq:V12-five-stock-margin}
\end{equation}
Then
\begin{equation}
 \boxed{
 m_j>0
 \quad\Longrightarrow\quad
 \gap(H_j)\ge m_j.}
 \label{eq:V12-five-stock-gap}
\end{equation}
Consequently,
\begin{equation}
 \inf_jm_j>0
 \label{eq:V12-uniform-five-stock}
\end{equation}
is a volume- and cutoff-uniform Hamiltonian mass certificate.
\end{theorem}

\begin{proof}
\Cref{prop:V12-inner-Woodbury} supplies
\eqref{eq:V12-inner-local-floor} with
\[
 r_{0,j}=a_j(1-\rho_j),
 \qquad
 \delta_{0,j}=\delta_{{\rm loc},j}.
\]
The exterior hypothesis is \eqref{eq:V12-outer-local-floor} with
\[
 r_{1,j}=b_j,
 \qquad
 \delta_{1,j}=\delta_{{\rm ext},j}.
\]
Apply \Cref{thm:V12-two-region-compiler} and then
\Cref{cor:V6-one-form-incidence}.
\end{proof}

\begin{assessment}[What V35 supplies and what it does not]
\label{ass:V12-V35-landing}
The supplied V35 comparison inverse and Woodbury identities provide the
correct algebraic objects \((\mathcal A,\mathcal K,\mathcal X)\).  They do not
by themselves populate the cofinal physical values in
\eqref{eq:V12-five-stock-margin}.  In particular, a return edge measured only
on one selected axial source card is not the complete \(\rho_j\) required in
\eqref{eq:V12-return-edge}, and inverse locality is not the exterior floor
\(b_j\).
\end{assessment}

% =============================================================================
\section{Why exterior rarity is not an exterior floor}
\label{sec:V12-rarity}
% =============================================================================

\begin{counterexample}[Vanishing exterior probability with a vanishing gap]
\label{cth:V12-rarity-no-gap}
For every \(0<\varepsilon<1/2\), let
\[
 X_\varepsilon=\{0,1\},
 \qquad
 \pi_\varepsilon(1)=\varepsilon,
 \qquad
 \pi_\varepsilon(0)=1-\varepsilon.
\]
Define the reversible Dirichlet form
\begin{equation}
 \mathcal E_\varepsilon(f,f)
 :=
 c_\varepsilon|f(1)-f(0)|^2,
 \qquad
 c_\varepsilon:=\varepsilon^2(1-\varepsilon).
 \label{eq:V12-two-state-form}
\end{equation}
Then the exterior state has probability \(\varepsilon\to0\), while
\begin{equation}
 \boxed{
 \gap(\mathcal E_\varepsilon,\pi_\varepsilon)
 =\varepsilon\longrightarrow0.}
 \label{eq:V12-two-state-gap}
\end{equation}
Thus no estimate involving only
\(\pi_\varepsilon(\text{exterior})\) can provide a positive uniform
Poincar\'e or exact-one-form floor.
\end{counterexample}

\begin{proof}
For every nonconstant \(f\),
\[
 \Var_{\pi_\varepsilon}(f)
 =
 \varepsilon(1-\varepsilon)
 |f(1)-f(0)|^2.
\]
Therefore the Rayleigh quotient is identically
\[
 \frac{\mathcal E_\varepsilon(f,f)}
      {\Var_{\pi_\varepsilon}(f)}
 =
 \frac{c_\varepsilon}
      {\varepsilon(1-\varepsilon)}
 =\varepsilon.
\]
This is the unique positive eigenvalue.
\end{proof}

\begin{firewall}[Chebyshev is not chart-exit control]
\label{fire:V12-Chebyshev}
A Wilson-action or concentration bound showing that the large-field region
has small ground-law probability does not populate \(b_j\).  The missing
quantity is energetic communication across the chart boundary: equivalently
an exterior killed Witten floor, a capacity/conductance lower bound, or a
direct estimate of \eqref{eq:V12-outer-local-floor}.  Replacing that stock by
probability would recreate the two-state failure in
\Cref{cth:V12-rarity-no-gap}.
\end{firewall}

% =============================================================================
\section{Failure localization and the corrected upstream frontier}
\label{sec:V12-failure}
% =============================================================================

\begin{corollary}[Typed failure alternative]
\label{cor:V12-failure-alternative}
Assume the five-stock hypotheses hold along a cofinal regulator sequence.  If
\(\gap(H_j)\to0\) along a subsequence, then
\begin{equation}
 \boxed{
 \liminf_j
 \left[
 \min\{a_j(1-\rho_j),b_j\}
 -\varepsilon_{{\rm IMS},j}
 -\delta_{{\rm loc},j}
 -\delta_{{\rm ext},j}
 \right]
 \le0.}
 \label{eq:V12-forced-failure}
\end{equation}
In particular, a collapsing gap is impossible if there are constants
\[
 a_*>0,\qquad \rho^*<1,\qquad b_*>0,\qquad e_*\ge0
\]
such that
\begin{align}
 a_j&\ge a_*,
 &
 \rho_j&\le\rho^*,
 &
 b_j&\ge b_*,
 \label{eq:V12-uniform-good-stocks}\\
 \varepsilon_{{\rm IMS},j}
 +\delta_{{\rm loc},j}
 +\delta_{{\rm ext},j}
 &\le e_*
 <
 \min\{a_*(1-\rho^*),b_*\}.
 \label{eq:V12-uniform-cost-margin}
\end{align}
Hence every genuine soft sequence must be assigned to at least one of:
\[
 a_j\downarrow0,\quad
 \rho_j\uparrow1,\quad
 b_j\downarrow0,\quad
 \text{or localization costs consuming the margin}.
\]
\end{corollary}

\begin{proof}
If the liminf in \eqref{eq:V12-forced-failure} were positive,
\Cref{thm:V12-five-stock} would give a positive uniform lower bound for
\(\gap(H_j)\), a contradiction.  Conditions
\eqref{eq:V12-uniform-good-stocks}--\eqref{eq:V12-uniform-cost-margin}
make the same bracket uniformly positive.
\end{proof}

\begin{theorem}[Integration with the temporal incidence]
\label{thm:V12-temporal-spatial-integration}
Suppose a cofinal anisotropic Wilson sequence has:
\begin{enumerate}[label=\textup{(\Alph*)},leftmargin=4em]
\item the V11 temporal calibration and transfer-to-Hamiltonian convergence at
      each fixed spatial regulator;
\item the uniform five-stock margin
      \(\inf_jm_j=m_*>0\);
\item the quotient-stable, nontrivial, reflection-positive continuum
      convergence and local lower-frame hypotheses of V1.
\end{enumerate}
Then the resulting continuum OS Hamiltonian has
\begin{equation}
 \boxed{\gap(H_{\rm OS})\ge m_*.}
 \label{eq:V12-OS-gap}
\end{equation}
\end{theorem}

\begin{proof}
By \Cref{thm:V12-five-stock}, the differential Hamiltonian at every spatial
regulator has gap at least \(m_*\).  V11 identifies that Hamiltonian with the
calibrated temporal Wilson transfer generator.  The quotient-stable passage
and lower-frame theorem of V1 transports the common delayed contraction to
the nonzero limiting local OS space.  The spectral theorem then gives
\eqref{eq:V12-OS-gap}.
\end{proof}

\begin{firewall}[The theorem is a compiler, not a supplied mass]
\label{fire:V12-not-mass}
No attached document supplies a regulator-uniform complete return edge
\(\rho_j<1\), an exterior killed floor \(b_j>0\), and localization defects
whose sum leaves the strict margin
\eqref{eq:V12-five-stock-margin} on the actual Perron ground law.  Nor is the
nontrivial OS limit constructed.  Therefore
\Cref{thm:V12-temporal-spatial-integration} is not asserted to have all of its
hypotheses populated.
\end{firewall}

% =============================================================================
\section{V12 ledger and next upstream attack}
\label{sec:V12-ledger}
% =============================================================================

\begin{theorem}[Integrated V12 statement]
\label{thm:V12-integrated}
Preserving V1--V11, V12 proves:
\begin{enumerate}[label=\textup{(V12.\arabic*)},leftmargin=5.5em]
\item the exact weighted one-form IMS identity for the physical ground-law
      Witten operator;
\item the precise failure of exact-form invariance under chart cutoffs;
\item a two-region compiler with explicit IMS and router debits;
\item the sharp \(w^{-2}\) transition-width cost for a radial two-chart
      partition;
\item landing of the complete V35 Woodbury return edge in the inner physical
      floor;
\item the five-stock cofinal spatial mass certificate;
\item a reversible counterexample proving that exterior rarity alone gives no
      mass;
\item a typed alternative assigning every remaining soft sequence to an
      inner floor, return edge, exterior floor, or localization cost;
\item conditional integration of the spatial certificate with the
      unconditional V11 temporal incidence and the V1 OS passage.
\end{enumerate}
\end{theorem}

\begin{proof}
Items \textup{(V12.1)--(V12.2)} are
\Cref{lem:V12-Witten-IMS,lem:V12-exact-localization-defect}.  Items
\textup{(V12.3)--(V12.4)} are
\Cref{thm:V12-two-region-compiler,cor:V12-chart-width}.  Items
\textup{(V12.5)--(V12.6)} are
\Cref{prop:V12-inner-Woodbury,thm:V12-five-stock}.  Item
\textup{(V12.7)} is \Cref{cth:V12-rarity-no-gap}.  Items
\textup{(V12.8)--(V12.9)} are
\Cref{cor:V12-failure-alternative,thm:V12-temporal-spatial-integration}.
\end{proof}

\begingroup
\small
\begin{longtable}{>{\raggedright\arraybackslash}p{0.28\textwidth}
                  >{\raggedright\arraybackslash}p{0.16\textwidth}
                  >{\raggedright\arraybackslash}p{0.48\textwidth}}
\toprule
\textbf{V12 row}&\textbf{Status}&\textbf{Advance or remaining obstruction}\\
\midrule
\endfirsthead
\toprule
\textbf{V12 row}&\textbf{Status}&\textbf{Advance or remaining obstruction}\\
\midrule
\endhead
One-form IMS identity
& \MGpass
& Exact on the physical weighted Witten form, with explicit cutoff cost.\\[0.35em]
Exact-range localization
& Corrected
& A cutoff creates \(d\chi\wedge\omega\); the router debit can no longer be omitted.\\[0.35em]
Inner Woodbury landing
& \MGpass
& A complete edge \(\rho<1\) gives the floor \(a(1-\rho)\).\\[0.35em]
Exterior chart sector
& \MGopen
& Requires a killed Witten floor or capacity/conductance estimate; probability is insufficient.\\[0.35em]
Complete return edge
& \MGopen
& Must cover the full localized horizontal carrier, not one axial card.\\[0.35em]
Five-stock compiler
& \MGpass
& Gives the explicit margin \eqref{eq:V12-five-stock-margin}.\\[0.35em]
Weak-coupling spatial mass
& \MGopen
& None of the attachments populates all five stocks uniformly on the Perron ground law.\\[0.35em]
Nontrivial continuum OS state
& \MGopen
& V1 convergence and lower-frame hypotheses remain unproved.\\[0.35em]
Full Yang--Mills mass gap
& \MGopen
& Temporal incidence is closed; spatial nonperturbative conductance and the OS limit remain.\\
\bottomrule
\end{longtable}
\endgroup

\begin{openproblem}[V13 mandate: exterior capacity before further Hessian algebra]
\label{op:V13}
The exact comparison and return identities are already available.  Proceed
upstream in the following order.
\begin{enumerate}[label=\textup{(V13.\arabic*)},leftmargin=5.5em]
\item On the actual Perron ground law, choose a gauge-fixed physical
      small-field radius \(r_j\), a transition width \(w_j\), and compute
      \(\varepsilon_{{\rm IMS},j}\) in the physical metric.
\item Promote the V35 comparison floor and source-cyclic return from one
      selected card to the complete localized horizontal carrier, or return a
      normalized vector witnessing \(a_j\downarrow0\) or
      \(\rho_j\uparrow1\).
\item Prove the exterior inequality
      \eqref{eq:V12-outer-local-floor} by a killed Witten,
      capacity/conductance, or stopped-martingale estimate.  Do not substitute
      a Chebyshev probability bound.
\item Control the exact-range localization term
      \(d\chi_j\wedge\omega\) and populate
      \(\delta_{{\rm loc},j}+\delta_{{\rm ext},j}\), or return the chart-boundary
      current that prevents it.
\item Short wrapping, toron, and protected topological sectors before taking
      the complete edge norm.
\item If \(\inf_jm_j>0\), invoke
      \Cref{thm:V12-temporal-spatial-integration}; otherwise classify the
      forced failure using \Cref{cor:V12-failure-alternative}.
\end{enumerate}
Further local Hessian identities without an exterior-capacity stock would be
circular: they would refine \(a_j\) while leaving the independent
\(b_j\) row empty.
\end{openproblem}

\begin{firewall}[End-of-V12 claim boundary]
V12 proves the missing chart-assembly theorem and shows exactly why rare
large fields cannot be discarded.  It does not prove the five required
stocks uniformly on the weak-coupling continuum trajectory, control every
toron and wrapping sector, or construct a nontrivial continuum OS state.  The
full Yang--Mills mass gap remains open.
\end{firewall}

\section*{Version V12 append-only record}
\addcontentsline{toc}{section}{Version V12 append-only record}
The complete V11 source was retained before this continuation.  Its SHA--256
digest was
\begin{center}\scriptsize\ttfamily
a9208056859118be7e177cf8afad5a767b4ade9a9b4de7398b1c906795d1eb61.
\end{center}
On the next iteration, retain this full file, append before its unique active
document-closing command, increment the filename to
\texttt{\_V13.tex}, and remove only the superseded V12 file from this note
series.

% =============================================================================
% END APPENDED VERSION V12 MATERIAL
% =============================================================================

% =============================================================================
% BEGIN APPENDED VERSION V13 MATERIAL
% =============================================================================

\clearpage
\part{Version V13: Perron--Harnack control and a physical bridge-capacity certificate}
\label{part:V13}

\section{Direction and upstream audit}
\label{sec:V13-direction}

V12 showed that the chart exterior cannot be removed by a probability
estimate and identified capacity/conductance as an independent stock.  The
Grand Tensor V067 manuscript already contains an abstract local--capacity gap
theorem and a Thomson action-flow construction.  Repeating those results
would not advance the programme.

The open incidence is the density carried by that flow.  For the Hamiltonian
selected by the Wilson transfer, the physical law is
\[
 \dd\pi=h^2\dd\mathfrak m,
\]
where \(h\) is the Perron factor of the same transfer kernel.  A flow estimate
under the bare Wilson or Haar law omits \(h^2\) and therefore need not control
the physical capacity.  V13 supplies two exact, gap-free controls of this
terminal factor:
\begin{enumerate}[label=\textup{(\roman*)},leftmargin=3.6em]
\item a finite-difference Perron--Harnack bound from boundary-action
      oscillation; and
\item a pathwise Harnack bound from the physical V8 boundary score.
\end{enumerate}
These bounds are then inserted into the already-supplied Thomson flow
construction to obtain one explicit same-law bridge certificate.

% =============================================================================
\section{Gap-free Perron--Harnack inequalities}
\label{sec:V13-Harnack}
% =============================================================================

Retain the V8 symmetric positive transfer kernel
\[
 k(x,y)=e^{-S(x,y)},
 \qquad
 \int k(x,y)h(y)\,\dd\mathfrak m(y)=\lambda h(x),
 \qquad h>0.
\]

\begin{theorem}[Finite-difference Perron--Harnack bound]
\label{thm:V13-finite-Harnack}
For \(x,x'\in X\), define
\begin{align}
 \Delta_+(x,x')
 &:=
 \sup_{y\in X}\bigl(S(x,y)-S(x',y)\bigr),
 \label{eq:V13-Delta-plus}\\
 \Delta(x,x')
 &:=
 \sup_{y\in X}
 \left|S(x,y)-S(x',y)\right|.
 \label{eq:V13-Delta-symmetric}
\end{align}
Then
\begin{equation}
 \boxed{
 e^{-\Delta_+(x,x')}
 \le\frac{h(x)}{h(x')}
 \le e^{\Delta_+(x',x)}.}
 \label{eq:V13-asymmetric-Harnack}
\end{equation}
In particular,
\begin{equation}
 \boxed{
 e^{-\Delta(x,x')}
 \le\frac{h(x)}{h(x')}
 \le e^{\Delta(x,x')}.}
 \label{eq:V13-symmetric-Harnack}
\end{equation}
No transfer gap, reduced resolvent, or iteration in time occurs.
\end{theorem}

\begin{proof}
The definition of \(\Delta_+(x,x')\) gives
\[
 S(x,y)\le S(x',y)+\Delta_+(x,x'),
\]
hence
\[
 e^{-S(x,y)}
 \ge
 e^{-\Delta_+(x,x')}e^{-S(x',y)}
\]
for every \(y\).  Multiply by \(h(y)>0\), integrate, and use the Perron
equation at \(x\) and \(x'\):
\[
 \lambda h(x)
 \ge
 e^{-\Delta_+(x,x')}\lambda h(x').
\]
This is the lower bound.  Interchanging \(x,x'\) gives the upper bound.
Since both asymmetric oscillations are bounded by \(\Delta(x,x')\),
\eqref{eq:V13-symmetric-Harnack} follows.
\end{proof}

\begin{theorem}[Pathwise physical-score Harnack bound]
\label{thm:V13-score-Harnack}
Let \(\gamma:[0,1]\to X\) be a \(C^1\) path and let \(p_x\) and \(J_x\) be
the physical conditional law and boundary score from V8.  Then
\begin{equation}
 \boxed{
 \log\frac{h(\gamma(1))}{h(\gamma(0))}
 =
 -\int_0^1
 \mathbb E_{p_{\gamma(s)}}
 \left[
   \ip{J_{\gamma(s)}}{\dot\gamma(s)}
 \right]\,\dd s.}
 \label{eq:V13-score-line-integral}
\end{equation}
Define the physical score length
\begin{equation}
 \mathscr S(\gamma)
 :=
 \int_0^1
 \left\{
 \mathbb E_{p_{\gamma(s)}}
 \left|
 \ip{J_{\gamma(s)}}{\dot\gamma(s)}
 \right|^2
 \right\}^{1/2}\dd s.
 \label{eq:V13-score-length}
\end{equation}
Then
\begin{equation}
 \boxed{
 e^{-\mathscr S(\gamma)}
 \le
 \frac{h(\gamma(1))}{h(\gamma(0))}
 \le
 e^{\mathscr S(\gamma)}.}
 \label{eq:V13-score-Harnack}
\end{equation}
\end{theorem}

\begin{proof}
By \eqref{eq:V8-first-score},
\[
 \nabla\log h(x)=-\mathbb E_{p_x}J_x.
\]
Pair with \(\dot\gamma(s)\), integrate the chain rule, and obtain
\eqref{eq:V13-score-line-integral}.  For each \(s\), Cauchy--Schwarz under
\(p_{\gamma(s)}\) gives
\[
 \left|
 \mathbb E_{p_{\gamma(s)}}
 \ip{J_{\gamma(s)}}{\dot\gamma(s)}
 \right|
 \le
 \left\{
 \mathbb E_{p_{\gamma(s)}}
 \left|
 \ip{J_{\gamma(s)}}{\dot\gamma(s)}
 \right|^2
 \right\}^{1/2}.
\]
Integrating and exponentiating proves
\eqref{eq:V13-score-Harnack}.
\end{proof}

\begin{assessment}[Two bounds with different uses]
\label{ass:V13-two-Harnack}
The finite-difference bound is automatic once a boundary-action oscillation
is known, but can be exponentially wasteful when an anisotropic temporal
coupling is large.  The score-length bound retains cancellations under the
terminal-weighted physical conditional and is therefore the target-native
quantity after the V11 temporal limit.
\end{assessment}

% =============================================================================
\section{A local Wilson action-oscillation stock}
\label{sec:V13-Wilson-locality}
% =============================================================================

\begin{lemma}[Boundary-support action oscillation]
\label{lem:V13-Wilson-oscillation}
Suppose the finite transfer action has the Wilson form
\begin{equation}
 S(x,y)=
 \sum_{P}c_P
 \left(
 1-\frac13\operatorname{ReTr}U_P(x,y)
 \right),
 \qquad c_P\ge0.
 \label{eq:V13-Wilson-action}
\end{equation}
Let \(x,x'\) differ only on a set \(F\) of boundary links, and let
\(\mathcal I(F)\) be the plaquettes whose ordered word meets \(F\).  Then
\begin{equation}
 \boxed{
 \Delta(x,x')
 \le
 2\sum_{P\in\mathcal I(F)}c_P.}
 \label{eq:V13-local-action-oscillation}
\end{equation}
If every boundary link belongs to at most \(d_*\) action plaquettes and
\(c_P\le c_*\), then
\begin{equation}
 \boxed{
 \Delta(x,x')\le2d_*c_*|F|.}
 \label{eq:V13-incidence-oscillation}
\end{equation}
\end{lemma}

\begin{proof}
For \(U\in\mathrm{SU}(3)\),
\[
 -1\le\frac13\operatorname{ReTr}U\le1,
\]
so the difference of two Wilson plaquette summands is at most \(2\) in
absolute value.  Plaquettes outside \(\mathcal I(F)\) have identical boundary
words at \(x\) and \(x'\) and cancel.  The triangle inequality gives
\eqref{eq:V13-local-action-oscillation}.  The incidence bound counts each
affected plaquette at most once after overcounting it by its incident links,
which yields \eqref{eq:V13-incidence-oscillation}.
\end{proof}

\begin{corollary}[Local Perron ratio]
\label{cor:V13-local-Perron-ratio}
Under \Cref{lem:V13-Wilson-oscillation},
\begin{equation}
 \boxed{
 \exp\!\left(
  -2\sum_{P\in\mathcal I(F)}c_P
 \right)
 \le\frac{h(x)}{h(x')}
 \le
 \exp\!\left(
  2\sum_{P\in\mathcal I(F)}c_P
 \right).}
 \label{eq:V13-local-Perron-ratio}
\end{equation}
For fixed \(F\) and bounded local couplings, the constant is independent of
the total spatial volume.
\end{corollary}

\begin{proof}
Insert \eqref{eq:V13-local-action-oscillation} into
\eqref{eq:V13-symmetric-Harnack}.
\end{proof}

\begin{firewall}[Anisotropic degeneration]
\label{fire:V13-anisotropic-degeneration}
If a plaquette incident on \(F\) carries the temporal Wilson coefficient
\(\alpha\to\infty\), the crude factor in
\eqref{eq:V13-local-Perron-ratio} can decay like \(e^{-C\alpha|F|}\).
Volume locality does not prevent this temporal degeneration.  V11 removes
the temporal lattice first; the spatial capacity estimate must then use the
differential physical score length, or a sharper cancellation, rather than
claiming that \eqref{eq:V13-local-Perron-ratio} is uniform in \(\alpha\).
\end{firewall}

% =============================================================================
\section{Perron-weighted flow conductance}
\label{sec:V13-flow}
% =============================================================================

Let
\[
 \Phi:[0,1]\times\Sigma\longrightarrow X
\]
be a smooth injective flow box joining an inner chart face to an exterior
chart face.  Write
\begin{equation}
 \Phi^*\dd\mathfrak m
 =J_{\mathfrak m}(s,z)\,\dd s\,\dd\sigma(z),
 \qquad
 v(s,z):=\partial_s\Phi(s,z).
 \label{eq:V13-Haar-flow-density}
\end{equation}
Because \(\dd\pi=h^2\dd\mathfrak m\), its physical flow density is
\begin{equation}
 \boxed{
 w_\pi(s,z)
 =h(\Phi(s,z))^2J_{\mathfrak m}(s,z).}
 \label{eq:V13-physical-flow-density}
\end{equation}

\begin{theorem}[Physical Perron flow-conductance bound]
\label{thm:V13-Perron-flow}
For \(z\in\Sigma\), set
\begin{equation}
 \mathscr S_z
 :=
 \int_0^1
 \left\{
 \mathbb E_{p_{\Phi(s,z)}}
 \left|
 \ip{J_{\Phi(s,z)}}{v(s,z)}
 \right|^2
 \right\}^{1/2}\dd s.
 \label{eq:V13-tube-score-length}
\end{equation}
Suppose measurable functions
\[
 j_*(z)>0,\qquad \ell(z)>0
\]
satisfy
\begin{equation}
 J_{\mathfrak m}(s,z)\ge j_*(z),
 \qquad
 \norm{v(s,z)}\le\ell(z)
 \quad(0\le s\le1).
 \label{eq:V13-flow-geometry}
\end{equation}
The stream resistance of the physical law,
\begin{equation}
 R_\Phi(z)
 :=
 \int_0^1
 \frac{\norm{v(s,z)}^2}
      {w_\pi(s,z)}\,\dd s,
 \label{eq:V13-stream-resistance}
\end{equation}
then obeys
\begin{equation}
 \boxed{
 R_\Phi(z)
 \le
 \frac{\ell(z)^2e^{2\mathscr S_z}}
      {h(\Phi(0,z))^2j_*(z)}.}
 \label{eq:V13-stream-upper}
\end{equation}
Consequently the parallel flow conductance satisfies
\begin{align}
 G_\Phi
 &:=
 \int_\Sigma R_\Phi(z)^{-1}\,\dd\sigma(z),
 \label{eq:V13-parallel-conductance}\\
 \boxed{
 G_\Phi
 &\ge
 G_{\rm Per}:=
 \int_\Sigma
 \frac{
 h(\Phi(0,z))^2j_*(z)e^{-2\mathscr S_z}
 }{\ell(z)^2}\,\dd\sigma(z).}
 \label{eq:V13-Perron-conductance}
\end{align}
\end{theorem}

\begin{proof}
Apply \Cref{thm:V13-score-Harnack} to the path
\(s\mapsto\Phi(s,z)\), first from \(0\) to each \(s\).  Since every partial
score length is at most \(\mathscr S_z\),
\[
 h(\Phi(s,z))
 \ge
 h(\Phi(0,z))e^{-\mathscr S_z}.
\]
Together with \eqref{eq:V13-flow-geometry} and
\eqref{eq:V13-physical-flow-density}, this yields
\[
 \frac{\norm{v(s,z)}^2}{w_\pi(s,z)}
 \le
 \frac{\ell(z)^2e^{2\mathscr S_z}}
      {h(\Phi(0,z))^2j_*(z)}.
\]
The \(s\)-interval has length one, so integration proves
\eqref{eq:V13-stream-upper}.  Invert that positive inequality and integrate
over \(z\) to obtain \eqref{eq:V13-Perron-conductance}.
\end{proof}

\begin{corollary}[Finite-difference substitute]
\label{cor:V13-finite-flow}
For each stream define
\begin{equation}
 \Delta_z:=
 \sup_{\substack{0\le s\le1\\y\in X}}
 \left|
 S(\Phi(s,z),y)-S(\Phi(0,z),y)
 \right|.
 \label{eq:V13-stream-Delta}
\end{equation}
Then \eqref{eq:V13-Perron-conductance} remains true with
\(\mathscr S_z\) replaced by \(\Delta_z\).
\end{corollary}

\begin{proof}
\Cref{thm:V13-finite-Harnack} gives
\[
 h(\Phi(s,z))
 \ge h(\Phi(0,z))e^{-\Delta_z}
\]
uniformly in \(s\).  Repeat the proof of
\Cref{thm:V13-Perron-flow}.
\end{proof}

% =============================================================================
\section{A same-law two-well mass certificate}
\label{sec:V13-two-well}
% =============================================================================

Let \(\mathcal W_0,\mathcal W_1\) be a protected partition of the physical
configuration carrier, with
\[
 p_i:=\pi(\mathcal W_i)>0,
 \qquad p_0+p_1=1,
\]
and conditional laws \(\pi_i\).  Let
\(\rho_{\rm loc}>0\) be a common unscaled Poincar\'e floor within the two
wells.  The flow box above lands in endpoint laws
\(\sigma_0,\sigma_1\), which need not equal \(\pi_0,\pi_1\).  Retain the
endpoint compliances from the supplied Grand Tensor action-flow theorem:
\begin{equation}
 \varkappa_i
 :=
 \sup_{\mathcal E_{\pi_i}(f)>0}
 \frac{
 \left|
 \int f\,\dd\sigma_i-\int f\,\dd\pi_i
 \right|^2
 }{\mathcal E_{\pi_i}(f)}.
 \label{eq:V13-endpoint-compliance}
\end{equation}

\begin{theorem}[Perron bridge-to-mass certificate]
\label{thm:V13-bridge-mass}
Assume \(G_{\rm Per}>0\) and
\(\varkappa_0,\varkappa_1<\infty\).  Define
\begin{equation}
 \boxed{
 \overline R_{\rm br}
 :=
 \left(
 G_{\rm Per}^{-1/2}
 +\sqrt{\varkappa_0}
 +\sqrt{\varkappa_1}
 \right)^2.}
 \label{eq:V13-bridge-resistance}
\end{equation}
Then the two-well capacity floor satisfies
\begin{equation}
 \boxed{
 \rho_{\rm cap}
 \ge
 \frac1{p_0p_1\overline R_{\rm br}}.}
 \label{eq:V13-capacity-floor}
\end{equation}
The physical Hamiltonian therefore obeys
\begin{equation}
 \boxed{
 \gap(H)
 \ge
 \kappa
 \left[
 \rho_{\rm loc}^{-1}
 +p_0p_1\overline R_{\rm br}
 \right]^{-1}.}
 \label{eq:V13-bridge-mass}
\end{equation}
\end{theorem}

\begin{proof}
The action-flow theorem in the supplied Grand Tensor V067 manuscript says
that the parallel tube current has resistance \(G_\Phi^{-1}\), and that
correcting its endpoints to the conditional well laws produces a unit
interwell current of resistance at most
\[
 R_{\rm br}
 \le
 \left(
 G_\Phi^{-1/2}
 +\sqrt{\varkappa_0}
 +\sqrt{\varkappa_1}
 \right)^2.
\]
By \eqref{eq:V13-Perron-conductance},
\(G_\Phi\ge G_{\rm Per}\), hence
\(R_{\rm br}\le\overline R_{\rm br}\).  Thomson's principle gives effective
interwell conductance at least
\(\overline R_{\rm br}^{-1}\).

For well means \(m_0,m_1\), the macro variance is
\[
 p_0p_1|m_0-m_1|^2.
\]
The interwell energy is at least
\(\overline R_{\rm br}^{-1}|m_0-m_1|^2\), so the capacity spectral floor in
the probability-weighted mean space is at least
\((p_0p_1\overline R_{\rm br})^{-1}\).  This proves
\eqref{eq:V13-capacity-floor}.

The local--capacity gap theorem already proved in Grand Tensor V067 gives
\[
 \lambda_1(d_\pi^*d)
 \ge
 \left(
 \rho_{\rm loc}^{-1}
 +\rho_{\rm cap}^{-1}
 \right)^{-1}.
\]
Insert \eqref{eq:V13-capacity-floor} and multiply by the physical electric
coefficient \(\kappa\), using
\Cref{thm:V6-ground-state-representation}.
\end{proof}

\begin{corollary}[Cofinal Perron bridge packet]
\label{cor:V13-cofinal-bridge}
Along a cofinal regulator sequence, suppose
\begin{equation}
 \inf_j\kappa_j>0,
 \qquad
 \inf_j\rho_{{\rm loc},j}>0,
 \qquad
 \sup_j
 p_{0,j}p_{1,j}\overline R_{{\rm br},j}
 <\infty.
 \label{eq:V13-uniform-bridge-packet}
\end{equation}
Then the finite Hamiltonian gaps have a common positive lower bound.
Together with the V11 temporal incidence and the V1 nontrivial
quotient-stable OS convergence packet, the limiting OS Hamiltonian has the
same positive mass lower bound.
\end{corollary}

\begin{proof}
The right side of \eqref{eq:V13-bridge-mass} is uniformly positive under
\eqref{eq:V13-uniform-bridge-packet}.  Use the V11 temporal same-generator theorem and then the V1
quotient-stable lower-frame passage for the continuum conclusion.  This uses
the scalar Poincar\'e route directly and does not assume the separate V12
five-stock chart packet.
\end{proof}

\begin{corollary}[Bridge failure has a physical Perron witness]
\label{cor:V13-bridge-failure}
Assume
\[
 \inf_j\kappa_j>0,
 \qquad
 \inf_j\rho_{{\rm loc},j}>0,
\]
but \(\gap(H_j)\to0\).  Then every family of physical flow boxes and endpoint
repairs satisfying the hypotheses of
\Cref{thm:V13-bridge-mass} must obey
\begin{equation}
 \boxed{
 p_{0,j}p_{1,j}
 \left(
 G_{{\rm Per},j}^{-1/2}
 +\sqrt{\varkappa_{0,j}}
 +\sqrt{\varkappa_{1,j}}
 \right)^2
 \longrightarrow\infty}
 \label{eq:V13-forced-bridge-failure}
\end{equation}
along a subsequence.  Thus the obstruction is forced into at least one
target-native row:
\[
 \text{Perron-weighted tube conductance},\qquad
 \text{inner endpoint equilibration},\qquad
 \text{exterior endpoint equilibration}.
\]
\end{corollary}

\begin{proof}
If the expression in \eqref{eq:V13-forced-bridge-failure} remained bounded
along a further subsequence, \eqref{eq:V13-bridge-mass} and the common local
floor would give a positive gap on that subsequence, a contradiction.
\end{proof}

\begin{firewall}[A bare action tube is not the physical bridge]
\label{fire:V13-bare-tube}
The Jacobian \(J_{\mathfrak m}\), Wilson action barrier, and transverse area
do not by themselves give \(G_{\rm Per}\).  The factor
\[
 h(\Phi(0,z))^2e^{-2\mathscr S_z}
\]
is load bearing.  It is evaluated under the Perron-weighted physical
conditional of the same transfer kernel.  Dropping it would undo the
same-law correction established in V8.
\end{firewall}

\begin{firewall}[Protected sectors precede the two-well reduction]
\label{fire:V13-protected-sectors}
If the exterior contains disconnected centre, topology, toron, Gribov, or
wrapping labels readable as common zero-energy data, it is not one ordinary
well.  Those labels must first be shorted or retained as separate atlas
cells.  Only the remaining connected physical chart/exterior pair is covered
by \Cref{thm:V13-bridge-mass}.
\end{firewall}

% =============================================================================
\section{V13 ledger and next upstream attack}
\label{sec:V13-ledger}
% =============================================================================

\begin{theorem}[Integrated V13 statement]
\label{thm:V13-integrated}
Preserving V1--V12, V13 proves:
\begin{enumerate}[label=\textup{(V13.\arabic*)},leftmargin=5.5em]
\item a gap-free finite-difference Harnack inequality for the Perron factor;
\item an exact path integral and \(L^2\)-score-length Harnack inequality under
      the physical terminal-weighted conditional;
\item a volume-local Wilson boundary-action oscillation estimate;
\item the precise anisotropic failure of the crude oscillation bound;
\item an explicit lower bound for physical parallel-tube conductance
      containing the complete Perron factor;
\item a two-well capacity floor with endpoint-equilibration corrections;
\item a quantitative Hamiltonian mass lower bound from local Poincar\'e and
      Perron bridge stocks;
\item a cofinal packet and a forced physical bridge-failure alternative.
\end{enumerate}
\end{theorem}

\begin{proof}
Items \textup{(V13.1)--(V13.2)} are
\Cref{thm:V13-finite-Harnack,thm:V13-score-Harnack}.  Items
\textup{(V13.3)--(V13.4)} are
\Cref{lem:V13-Wilson-oscillation,fire:V13-anisotropic-degeneration}.  Item
\textup{(V13.5)} is \Cref{thm:V13-Perron-flow}.  Items
\textup{(V13.6)--(V13.7)} are
\Cref{thm:V13-bridge-mass}.  Item \textup{(V13.8)} is
\Cref{cor:V13-cofinal-bridge,cor:V13-bridge-failure}.
\end{proof}

\begingroup
\small
\begin{longtable}{>{\raggedright\arraybackslash}p{0.28\textwidth}
                  >{\raggedright\arraybackslash}p{0.16\textwidth}
                  >{\raggedright\arraybackslash}p{0.48\textwidth}}
\toprule
\textbf{V13 row}&\textbf{Status}&\textbf{Advance or remaining obstruction}\\
\midrule
\endfirsthead
\toprule
\textbf{V13 row}&\textbf{Status}&\textbf{Advance or remaining obstruction}\\
\midrule
\endhead
Perron finite-difference Harnack
& \MGpass
& Exact and gap free; volume local for a fixed boundary-link deformation.\\[0.35em]
Physical score Harnack
& \MGpass
& Uses the V8 terminal-weighted conditional and retains cancellations.\\[0.35em]
Physical tube conductance
& \MGpass
& Explicitly bounded below by \(G_{\rm Per}\).\\[0.35em]
Two-well bridge certificate
& \MGpass
& Gives \eqref{eq:V13-bridge-mass} after endpoint repairs.\\[0.35em]
Uniform physical score length
& \MGopen
& No attachment bounds \(\mathscr S_z\) on a cofinal family of parallel chart-exit tubes.\\[0.35em]
Perron anchor mass
& \MGopen
& The transverse integral of \(h(\Phi(0,z))^2\) must not collapse.\\[0.35em]
Endpoint equilibration
& \MGopen
& The two compliances \(\varkappa_i\) remain unpopulated on the actual ground law.\\[0.35em]
Protected exterior sectors
& \MGopen
& Toron, wrapping, topology, and Gribov labels must be separately shorted.\\[0.35em]
Full Yang--Mills mass gap
& \MGopen
& The physical bridge packet and nontrivial OS continuum construction remain.\\
\bottomrule
\end{longtable}
\endgroup

\begin{openproblem}[V14 mandate: construct a cofinal Perron flow atlas]
\label{op:V14}
The abstract capacity theorem and its terminal-factor incidence are now
closed.  The next work is geometric and probabilistic, not another Schur
identity.
\begin{enumerate}[label=\textup{(V14.\arabic*)},leftmargin=5.5em]
\item After the V11 temporal limit, choose a gauge-fixed family of local
      physical flow boxes crossing the small-field chart boundary.
\item Prove a lower bound for the aggregate Perron anchor mass
      \[
      \int_{\Sigma_j}
      \frac{h_j(\Phi_j(0,z))^2j_{*,j}(z)}
           {\ell_j(z)^2}\,\dd\sigma_j(z)
      \]
      which is stable in volume and cutoff.
\item Bound the physical score lengths \(\mathscr S_{j,z}\) using the V8
      expectation-minus-covariance card and summable spatial response, rather
      than the anisotropically divergent raw action oscillation.
\item Match the endpoint laws to the two conditional wells and control
      \(\varkappa_{0,j},\varkappa_{1,j}\).
\item Refine the atlas before crossing any critical, toron, wrapping,
      topology, or Gribov stratum; return a normalized bridge-soft witness if
      the refined conductance collapses.
\item Insert the resulting
      \(p_{0,j}p_{1,j}\overline R_{{\rm br},j}\) in
      \eqref{eq:V13-bridge-mass}, then proceed to the V1 nontrivial OS limit.
\end{enumerate}
If the score length is estimated by the full physical Poincar\'e constant,
the argument is circular.  It must be bounded by local conditional response,
a convergent signed cluster/cumulant expansion, or an independently proved
spatial mixing estimate.
\end{openproblem}

\begin{firewall}[End-of-V13 claim boundary]
V13 supplies the missing same-law Perron factor in the action-flow capacity
route and an explicit bridge-to-mass inequality.  It does not construct a
cofinal flow atlas, bound its physical score lengths and anchor mass, control
all protected sectors, or construct the continuum OS state.  The full
Yang--Mills mass gap remains open.
\end{firewall}

\section*{Version V13 append-only record}
\addcontentsline{toc}{section}{Version V13 append-only record}
The complete V12 source was retained before this continuation.  Its SHA--256
digest was
\begin{center}\scriptsize\ttfamily
2ae294276b9b31a1fad29a35ea9eb724a5f74ed51f3f156ffe0772b8e6742805.
\end{center}
On the next iteration, retain this full file, append before its unique active
document-closing command, increment the filename to
\texttt{\_V14.tex}, and remove only the superseded V13 file from this note
series.

% =============================================================================
% END APPENDED VERSION V13 MATERIAL
% =============================================================================

% =============================================================================
% BEGIN APPENDED VERSION V14 MATERIAL
% =============================================================================

\clearpage
\part{Version V14: sequential Perron marginals and an operator-norm martingale compiler}
\label{part:V14}

\section{Why the attack moves upstream from one global flow box}
\label{sec:V14-direction}

V13 reduced the chart-exit problem to a physical Perron-weighted flow
conductance.  A single global small-field flow box is not the only possible
implementation, and on a growing product configuration space it is often the
wrong first object.  If every link is required to lie in one proper Haar
chart, its bare transverse mass decreases exponentially with the number of
links.  The physical factor \(h^2\) may compensate for that loss, but proving
such compensation globally is as hard as controlling the full ground law.

V14 therefore moves upstream to a sequential coordinate atlas.  It reveals
one rooted tree-gauge \(\mathrm{SU}(3)\) coordinate at a time, computes the
exact marginal Perron amplitude, proves a uniform conditional Poincar\'e
floor from its physical score diameter, and places the remaining dependence
in one response matrix.  This does not repeat the generic V7 martingale
criterion: it supplies its previously unpopulated conditional-law and
derivative rows in terms of the V8 Perron score.

\begin{lemma}[Exponential bare loss of a global product chart]
\label{lem:V14-product-chart-loss}
Let \(B\subsetneq G=\mathrm{SU}(3)\) be measurable with normalized Haar mass
\(q:=\mathfrak m_G(B)\in(0,1)\).  In \(G^N\) with product Haar measure,
\begin{equation}
 \boxed{
 \mathfrak m_G^{\otimes N}(B^N)=q^N.}
 \label{eq:V14-product-chart-loss}
\end{equation}
Thus a cofinal proof cannot infer a positive global Perron anchor mass from
the bare geometry of one fixed proper chart on every link.
\end{lemma}

\begin{proof}
This is product-measure multiplicativity.  Since \(q<1\), \(q^N\to0\)
exponentially.  The final sentence does not exclude physical concentration
of \(h^2\); it says precisely that such concentration would need an
independent proof.
\end{proof}

% =============================================================================
\section{Sequential marginal Perron amplitudes}
\label{sec:V14-marginals}
% =============================================================================

Work after rooted tree gauge on a finite product-Haar coordinate carrier
\[
 X=G^N,
 \qquad
 \dd\mathfrak m=\bigotimes_{b=1}^N\dd\mathfrak m_G(x_b),
 \qquad
 \dd\pi=h^2\dd\mathfrak m,
 \qquad
 \int h^2\,\dd\mathfrak m=1.
\]
Protected centre, topology, toron, wrapping, and stabilizer labels are fixed
before this coordinate representation is used.

For \(i\in\{1,\ldots,N\}\), write
\[
 a=(x_1,\ldots,x_{i-1}),
 \qquad
 U=x_i,
 \qquad
 y=(x_{i+1},\ldots,x_N).
\]
Define the sequential marginal Perron amplitude
\begin{equation}
 \boxed{
 \widehat h_i(a,U)^2
 :=
 \int_{G^{N-i}}h(a,U,y)^2\,\dd\mathfrak m(y).}
 \label{eq:V14-marginal-amplitude}
\end{equation}
The conditional law of \(U\) given the past \(a\) is
\begin{equation}
 \boxed{
 \dd q_{i,a}(U)
 =
 \frac{\widehat h_i(a,U)^2}
      {\displaystyle\int_G\widehat h_i(a,V)^2\,\dd\mathfrak m_G(V)}
 \,\dd\mathfrak m_G(U).}
 \label{eq:V14-sequential-conditional}
\end{equation}
For \(\widehat h_i(a,U)>0\), the future conditional is
\begin{equation}
 \dd\nu_{a,U}(y)
 :=
 \frac{h(a,U,y)^2}
      {\widehat h_i(a,U)^2}\,\dd\mathfrak m(y).
 \label{eq:V14-future-conditional}
\end{equation}

Let
\begin{equation}
 \mathcal J_i(x)
 :=
 \mathbb E_{p_x}\bigl[(J_x)_i\bigr]
 \label{eq:V14-Perron-mean-score}
\end{equation}
be the \(i\)-th coordinate of the V8 physical boundary-score mean.

\begin{theorem}[Exact marginal Perron score]
\label{thm:V14-marginal-score}
For every \(a,U\),
\begin{equation}
 \boxed{
 \nabla_U\log\widehat h_i(a,U)
 =
 -\mathbb E_{\nu_{a,U}}\mathcal J_i(a,U,\cdot).}
 \label{eq:V14-marginal-score}
\end{equation}
Equivalently, the logarithmic score of the sequential conditional density
relative to Haar is
\begin{equation}
 \boxed{
 \nabla_U
 \log\frac{\dd q_{i,a}}{\dd\mathfrak m_G}(U)
 =
 -2\mathbb E_{\nu_{a,U}}\mathcal J_i(a,U,\cdot).}
 \label{eq:V14-conditional-score}
\end{equation}
\end{theorem}

\begin{proof}
Differentiate \eqref{eq:V14-marginal-amplitude} under the integral:
\[
 \nabla_U\log\widehat h_i
 =
 \frac{\int h^2\nabla_U\log h\,\dd\mathfrak m(y)}
      {\int h^2\,\dd\mathfrak m(y)}
 =
 \mathbb E_{\nu_{a,U}}\nabla_U\log h.
\]
The V8 first-score identity gives
\(
\nabla_U\log h=-\mathcal J_i
\), proving \eqref{eq:V14-marginal-score}.  The denominator in
\eqref{eq:V14-sequential-conditional} is independent of \(U\); differentiating
its logarithmic density and using
\eqref{eq:V14-marginal-score} proves
\eqref{eq:V14-conditional-score}.
\end{proof}

% =============================================================================
\section{A uniform single-coordinate conditional gap}
\label{sec:V14-coordinate-gap}
% =============================================================================

For a piecewise \(C^1\) path \(\gamma:[0,1]\to G\), define
\begin{equation}
 \mathscr S_{i,a}(\gamma)
 :=
 \int_0^1
 \left\{
 \mathbb E_{\nu_{a,\gamma(s)}}
 \mathbb E_{p_{(a,\gamma(s),y)}}
 \left|
 \ip{(J_{(a,\gamma(s),y)})_i}{\dot\gamma(s)}
 \right|^2
 \right\}^{1/2}\dd s.
 \label{eq:V14-coordinate-score-length}
\end{equation}
The sequential score diameter is
\begin{equation}
 \mathscr A_i(a)
 :=
 \sup_{U,V\in G}
 \inf_{\substack{\gamma(0)=U\\\gamma(1)=V}}
 \mathscr S_{i,a}(\gamma).
 \label{eq:V14-score-diameter}
\end{equation}

\begin{lemma}[Marginal Harnack bound]
\label{lem:V14-marginal-Harnack}
For every \(a,U,V\),
\begin{equation}
 \boxed{
 \left|
 \log\frac{\widehat h_i(a,U)}
          {\widehat h_i(a,V)}
 \right|
 \le\mathscr A_i(a).}
 \label{eq:V14-amplitude-Harnack}
\end{equation}
Consequently,
\begin{equation}
 \boxed{
 \operatorname{osc}_{U\in G}
 \log\frac{\dd q_{i,a}}{\dd\mathfrak m_G}(U)
 \le2\mathscr A_i(a).}
 \label{eq:V14-density-oscillation}
\end{equation}
\end{lemma}

\begin{proof}
Integrate \eqref{eq:V14-marginal-score} along a path \(\gamma\).
Two applications of Jensen's inequality give
\[
 \left|
 \mathbb E_{\nu_{a,\gamma(s)}}\mathcal J_i
 \cdot\dot\gamma(s)
 \right|
 \le
 \left\{
 \mathbb E_{\nu_{a,\gamma(s)}}
 \mathbb E_{p_{(a,\gamma(s),y)}}
 \left|
 \ip{(J_{(a,\gamma(s),y)})_i}{\dot\gamma(s)}
 \right|^2
 \right\}^{1/2}.
\]
Take the path infimum and then the endpoint supremum to obtain
\eqref{eq:V14-amplitude-Harnack}.  The conditional density in
\eqref{eq:V14-sequential-conditional} is proportional to
\(\widehat h_i^2\), so its logarithmic oscillation is twice that of
\(\widehat h_i\).
\end{proof}

\begin{lemma}[Elementary bounded-density Poincar\'e transport]
\label{lem:V14-density-Poincare}
Let \(\mu\) be a probability measure satisfying
\[
 \Var_\mu(f)
 \le\lambda^{-1}\int|\nabla f|^2\,\dd\mu.
\]
If \(\dd q=r\,\dd\mu\), \(r>0\), and
\(\operatorname{osc}\log r\le B\), then
\begin{equation}
 \boxed{
 \Var_q(f)
 \le\frac{e^B}{\lambda}
 \int|\nabla f|^2\,\dd q.}
 \label{eq:V14-density-Poincare}
\end{equation}
\end{lemma}

\begin{proof}
For any constant \(c\),
\[
 \Var_q(f)
 \le\int|f-c|^2r\,\dd\mu
 \le(\sup r)\int|f-c|^2\,\dd\mu.
\]
Choose \(c=\int f\,\dd\mu\), apply the \(\mu\)-Poincar\'e inequality, and use
\[
 \int|\nabla f|^2\,\dd\mu
 \le(\inf r)^{-1}\int|\nabla f|^2\,\dd q.
\]
The ratio \((\sup r)/(\inf r)\) is at most \(e^B\).
\end{proof}

\begin{theorem}[Sequential \(\mathrm{SU}(3)\) conditional gap]
\label{thm:V14-coordinate-gap}
Use the same bi-invariant metric as the electric Casimir, and let \(C_F\) be
the first positive eigenvalue of \(-\Delta_G\).  Then
\begin{equation}
 \boxed{
 \Var_{q_{i,a}}(f)
 \le
 \frac{e^{2\mathscr A_i(a)}}{C_F}
 \int_G|\nabla_Uf|^2\,\dd q_{i,a}(U).}
 \label{eq:V14-coordinate-Poincare}
\end{equation}
In particular, if
\begin{equation}
 \sup_{N,i,a}\mathscr A_i(a)\le A_*<\infty,
 \label{eq:V14-uniform-score-diameter}
\end{equation}
all sequential coordinate conditionals have the common Poincar\'e floor
\begin{equation}
 \boxed{
 \rho_{\rm site}:=C_Fe^{-2A_*}>0.}
 \label{eq:V14-site-floor}
\end{equation}
\end{theorem}

\begin{proof}
Normalized Haar measure on \(G\) has Poincar\'e floor \(C_F\) by
Peter--Weyl decomposition and the Casimir spectrum.  Apply
\Cref{lem:V14-density-Poincare} with
\(\mu=\mathfrak m_G\), \(q=q_{i,a}\), and the oscillation bound
\eqref{eq:V14-density-oscillation}.
\end{proof}

\begin{corollary}[Local score-moment sufficient stock]
\label{cor:V14-score-moment}
Let \(D_G\) be the diameter of \(G\).  If
\begin{equation}
 \sup_{N,i,a,U}
 \mathbb E_{\nu_{a,U}}
 \mathbb E_{p_{(a,U,y)}}
 \norm{(J_{(a,U,y)})_i}^2
 \le M_*^2,
 \label{eq:V14-score-moment}
\end{equation}
then
\begin{equation}
 \boxed{
 \mathscr A_i(a)\le D_GM_*,
 \qquad
 \rho_{\rm site}\ge
 C_Fe^{-2D_GM_*}.}
 \label{eq:V14-moment-site-floor}
\end{equation}
\end{corollary}

\begin{proof}
Join \(U\) to \(V\) by a minimizing geodesic of length at most \(D_G\).
The integrand in \eqref{eq:V14-coordinate-score-length} is bounded by
\(M_*\norm{\dot\gamma(s)}\).  Integrate, then apply
\Cref{thm:V14-coordinate-gap}.
\end{proof}

% =============================================================================
\section{The exact nested Perron response}
\label{sec:V14-response}
% =============================================================================

Let
\[
 \mathcal F_i=\sigma(x_1,\ldots,x_i),
 \qquad
 M_if=\mathbb E_\pi[f\mid\mathcal F_i],
 \qquad
 D_if=M_if-M_{i-1}f
\]
be the V7 reveal martingale.  At fixed \((a,U)\),
\[
 M_if(a,U)
 =
 \mathbb E_{\nu_{a,U}}f(a,U,\cdot).
\]

\begin{theorem}[Exact sequential Perron response identity]
\label{thm:V14-nested-response}
For every smooth \(f\),
\begin{equation}
 \boxed{
 \nabla_iM_if
 =
 \mathbb E_{\nu_{a,U}}\nabla_if
 -
 2\Cov_{\nu_{a,U}}
 \bigl(f,\mathcal J_i\bigr).}
 \label{eq:V14-nested-response}
\end{equation}
The covariance is a cotangent-vector-valued covariance.  Every term is under
the actual ground law \(h^2\mathfrak m\) and the V8 terminal-weighted
conditional of the same transfer kernel.
\end{theorem}

\begin{proof}
Differentiate the normalized future expectation.  For any parameterized
probability density,
\[
 \nabla\mathbb E_\nu f
 =
 \mathbb E_\nu\nabla f
 +\Cov_\nu(f,\nabla\log w),
\]
where \(w\) is its unnormalized density.  Here \(w=h^2\), so the second term
is \(2\Cov_\nu(f,\nabla_i\log h)\).  The V8 first-score identity gives
\(\nabla_i\log h=-\mathcal J_i\), proving
\eqref{eq:V14-nested-response}.
\end{proof}

\begin{assessment}[The remaining response is not an invented interaction]
\label{ass:V14-response-meaning}
The off-diagonal martingale collar is exactly the second term in
\eqref{eq:V14-nested-response}.  It is a covariance of the observable with
the future-dependent physical Perron score.  Bounding it by the global
Poincar\'e constant would be circular; bounding it by a local conditional
response expansion or a summable signed cumulant kernel would populate a
new input.
\end{assessment}

% =============================================================================
\section{An operator-norm martingale certificate}
\label{sec:V14-matrix-compiler}
% =============================================================================

\begin{theorem}[Sequential response-matrix Poincar\'e theorem]
\label{thm:V14-response-matrix}
Suppose every conditional \(q_{i,a}\) has Poincar\'e floor at least
\(\rho_{\rm site}>0\).  Suppose there is a nonnegative \(N\times N\) matrix
\(\mathsf A=(A_{ib})\) such that every smooth gauge-invariant \(f\) obeys
\begin{equation}
 \left\|\nabla_iM_if\right\|_{L^2(\pi)}
 \le
 \sum_{b=1}^NA_{ib}
 \left\|\nabla_bf\right\|_{L^2(\pi)}
 \qquad(1\le i\le N).
 \label{eq:V14-response-matrix-row}
\end{equation}
Then
\begin{equation}
 \boxed{
 \Var_\pi(f)
 \le
 \frac{\norm{\mathsf A}_{2\to2}^2}{\rho_{\rm site}}
 \int|\nabla f|^2\,\dd\pi.}
 \label{eq:V14-matrix-Poincare}
\end{equation}
For the physical gauge Hamiltonian,
\begin{equation}
 \boxed{
 \gap(H)
 \ge
 \frac{\kappa\rho_{\rm site}}
      {\norm{\mathsf A}_{2\to2}^2}.}
 \label{eq:V14-matrix-gap}
\end{equation}
\end{theorem}

\begin{proof}
Conditioned on \(\mathcal F_{i-1}\), the function \(D_if\) is the
mean-zero fluctuation of \(M_if\) in the variable \(x_i\).  The conditional
Poincar\'e inequality gives
\[
 \norm{D_if}_{L^2(\pi)}^2
 \le
 \rho_{\rm site}^{-1}
 \norm{\nabla_iM_if}_{L^2(\pi)}^2.
\]
Put
\[
 X_b:=\norm{\nabla_bf}_{L^2(\pi)},
 \qquad
 Y_i:=\norm{\nabla_iM_if}_{L^2(\pi)}.
\]
Hypothesis \eqref{eq:V14-response-matrix-row} says entrywise that
\(Y\le\mathsf AX\), hence
\[
 \sum_iY_i^2
 \le\norm{\mathsf A}_{2\to2}^2\sum_bX_b^2.
\]
Sum the martingale differences and use the exact variance ledger
\eqref{eq:V7-variance-ledger}.  This proves
\eqref{eq:V14-matrix-Poincare}; the ground-state representation
\Cref{thm:V6-ground-state-representation} gives
\eqref{eq:V14-matrix-gap}.
\end{proof}

\begin{corollary}[Explicit Schur-summable response bound]
\label{cor:V14-Schur-response}
Suppose \(\mathsf A=I+\mathsf R\), with
\(\mathsf R=(R_{ib})\ge0\), and
\begin{equation}
 r_{\rm row}:=\sup_i\sum_bR_{ib}<\infty,
 \qquad
 r_{\rm col}:=\sup_b\sum_iR_{ib}<\infty.
 \label{eq:V14-Schur-stocks}
\end{equation}
Then
\begin{equation}
 \boxed{
 \norm{\mathsf A}_{2\to2}
 \le1+\sqrt{r_{\rm row}r_{\rm col}},}
 \label{eq:V14-Schur-norm}
\end{equation}
and
\begin{equation}
 \boxed{
 \gap(H)
 \ge
 \frac{\kappa\rho_{\rm site}}
 {\left(1+\sqrt{r_{\rm row}r_{\rm col}}\right)^2}.}
 \label{eq:V14-Schur-gap}
\end{equation}
\end{corollary}

\begin{proof}
The Schur test gives
\[
 \norm{\mathsf R}_{2\to2}
 \le\sqrt{r_{\rm row}r_{\rm col}}.
\]
Use the triangle inequality
\(\norm{I+\mathsf R}\le1+\norm{\mathsf R}\) in
\Cref{thm:V14-response-matrix}.
\end{proof}

\begin{corollary}[Cofinal sequential Perron packet]
\label{cor:V14-cofinal}
For a cofinal family of protected finite regulators, assume
\begin{equation}
 \inf_j\kappa_j\ge\kappa_*>0,
 \qquad
 \sup_{j,i,a}\mathscr A_{j,i}(a)\le A_*<\infty,
 \qquad
 \sup_j\norm{\mathsf A_j}_{2\to2}\le A_{\rm op}<\infty.
 \label{eq:V14-cofinal-packet}
\end{equation}
Then
\begin{equation}
 \boxed{
 \gap(H_j)
 \ge
 \frac{\kappa_*C_Fe^{-2A_*}}{A_{\rm op}^2}
 >0}
 \label{eq:V14-cofinal-gap}
\end{equation}
uniformly.  With V11 temporal same-generator incidence and the V1
nontrivial quotient-stable OS convergence packet, this lower bound passes to
the continuum OS Hamiltonian.
\end{corollary}

\begin{proof}
Combine \Cref{thm:V14-coordinate-gap,thm:V14-response-matrix}.  The continuum
statement uses V11 followed by the V1 lower-frame passage.
\end{proof}

\begin{corollary}[Consistency with the bridge-capacity obstruction]
\label{cor:V14-capacity-consistency}
If a protected regulator sequence has a bridge-capacity soft mode satisfying
\(\kappa_j\rho_{{\rm cap},j}\to0\), then the three sequential stocks in
\eqref{eq:V14-cofinal-packet} cannot all remain uniform.  Along a subsequence,
at least one of
\[
 \kappa_j,\qquad
 \sup_{i,a}\mathscr A_{j,i}(a),\qquad
 \norm{\mathsf A_j}_{2\to2}
\]
must respectively tend to zero, diverge, or diverge.
\end{corollary}

\begin{proof}
If all three stocks remained within the bounds of
\eqref{eq:V14-cofinal-packet}, \eqref{eq:V14-cofinal-gap} would forbid a
collapsing gap.  The local--capacity upper bound in the supplied Grand Tensor
theorem gives
\(\gap(H_j)\le\kappa_j\rho_{{\rm cap},j}\to0\), producing the contradiction.
\end{proof}

\begin{firewall}[Conditional site gaps do not tensorize by themselves]
\label{fire:V14-no-naive-tensorization}
The lower bound \(\rho_{\rm site}>0\) controls each revealed
\(\mathrm{SU}(3)\) coordinate, but dependence can still accumulate through
\(\mathsf A\).  Setting \(\mathsf A=I\) is valid only for a product law.
For the Yang--Mills ground law, the exact correction is the Perron-score
covariance in \eqref{eq:V14-nested-response}.
\end{firewall}

\begin{firewall}[Gauge and common-factor order]
\label{fire:V14-gauge-order}
The coordinate filtration is applied only after rooted tree gauge and after
common protected labels have been fixed or shorted.  Revealing a gauge orbit
coordinate as though it were physical, or averaging over a toron label which
is readable at both ends, can create a false conditional gap.
\end{firewall}

% =============================================================================
\section{V14 ledger and revised frontier}
\label{sec:V14-ledger}
% =============================================================================

\begin{theorem}[Integrated V14 statement]
\label{thm:V14-integrated}
Preserving V1--V13, V14 proves:
\begin{enumerate}[label=\textup{(V14.\arabic*)},leftmargin=5.5em]
\item exponential bare-volume loss for a single proper product chart;
\item the exact sequential marginal Perron amplitude and conditional law;
\item the exact nested marginal score under the V8 physical conditional;
\item a conditional \(\mathrm{SU}(3)\) Poincar\'e floor from the physical
      score diameter;
\item a local score-second-moment sufficient bound for that diameter;
\item the exact future-observable derivative as a direct term minus a
      Perron-score covariance;
\item an operator-norm response-matrix Poincar\'e theorem;
\item a Schur-summable row/column certificate;
\item a cofinal finite-Hamiltonian mass packet compatible with V11 and V1;
\item a proof that any V13 bridge-capacity collapse must reappear in one of
      the sequential Perron stocks.
\end{enumerate}
\end{theorem}

\begin{proof}
Item \textup{(V14.1)} is \Cref{lem:V14-product-chart-loss}.  Items
\textup{(V14.2)--(V14.3)} are
\Cref{thm:V14-marginal-score}.  Items
\textup{(V14.4)--(V14.5)} are
\Cref{thm:V14-coordinate-gap,cor:V14-score-moment}.  Item
\textup{(V14.6)} is \Cref{thm:V14-nested-response}.  Items
\textup{(V14.7)--(V14.8)} are
\Cref{thm:V14-response-matrix,cor:V14-Schur-response}.  Items
\textup{(V14.9)--(V14.10)} are
\Cref{cor:V14-cofinal,cor:V14-capacity-consistency}.
\end{proof}

\begingroup
\small
\begin{longtable}{>{\raggedright\arraybackslash}p{0.28\textwidth}
                  >{\raggedright\arraybackslash}p{0.16\textwidth}
                  >{\raggedright\arraybackslash}p{0.48\textwidth}}
\toprule
\textbf{V14 row}&\textbf{Status}&\textbf{Advance or remaining obstruction}\\
\midrule
\endfirsthead
\toprule
\textbf{V14 row}&\textbf{Status}&\textbf{Advance or remaining obstruction}\\
\midrule
\endhead
Sequential Perron marginal
& \MGpass
& Exact square-integrated amplitude and physical conditional score.\\[0.35em]
Conditional \(\mathrm{SU}(3)\) gap
& \MGpass
& \(C_Fe^{-2\mathscr A_i(a)}\) from a target-native score diameter.\\[0.35em]
Nested response identity
& \MGpass
& Direct derivative minus the actual future Perron-score covariance.\\[0.35em]
Response-matrix compiler
& \MGpass
& Gives \eqref{eq:V14-matrix-gap} and the Schur-summable specialization.\\[0.35em]
Uniform score diameter
& \MGopen
& Requires a weak-coupling bound for \eqref{eq:V14-score-moment} after temporal incidence.\\[0.35em]
Uniform response operator
& \MGopen
& Requires summable row and column control of the covariance in \eqref{eq:V14-nested-response}.\\[0.35em]
Protected sectors
& \MGopen
& The coordinate filtration must be repeated on every retained toron/topology atlas cell.\\[0.35em]
Nontrivial continuum OS state
& \MGopen
& V1 convergence and lower-frame hypotheses remain unpopulated.\\[0.35em]
Full Yang--Mills mass gap
& \MGopen
& The weak-coupling Perron score and response bounds remain the spatial core.\\
\bottomrule
\end{longtable}
\endgroup

\begin{openproblem}[V15 mandate: bound the nested physical response kernel]
\label{op:V15}
The global flow geometry has been replaced by local sequential Perron stocks.
Proceed in the following order.
\begin{enumerate}[label=\textup{(V15.\arabic*)},leftmargin=5.5em]
\item Take the V11 temporal limit first and rewrite
      \(\mathcal J_i=-\nabla_i\log h\) in the differential Hamiltonian
      normalization.
\item Differentiate the nested score in
      \eqref{eq:V14-nested-response} and express its spatial off-diagonal
      kernel using the V8 expectation-minus-covariance Hessian card.
\item Separate the signed finite-range Wilson Hessian head from the genuinely
      connected response tail, subtracting each covariance return exactly
      once as required by supplied V35.
\item Prove regulator-uniform bounds for
      \eqref{eq:V14-score-moment} and either
      \(\norm{\mathsf A_j}_{2\to2}\) or the two Schur stocks
      \eqref{eq:V14-Schur-stocks}.
\item If a row or column sum diverges, extract the normalized spatial
      response vector and classify it as local, chart-exit, toron, wrapping,
      or topological before changing the norm.
\item Only after \eqref{eq:V14-cofinal-packet} closes, invoke V11 and the V1
      nontrivial OS passage.
\end{enumerate}
No inverse of the full scalar ground-law generator may be used to prove the
response bound; that inverse already contains the desired Poincar\'e
constant.
\end{openproblem}

\begin{firewall}[End-of-V14 claim boundary]
V14 proves an exact sequential same-law route from Perron score control to a
uniform Hamiltonian gap.  It does not bound the weak-coupling score diameter
or response-matrix norm, control every protected sector, or construct a
nontrivial continuum OS state.  The full Yang--Mills mass gap remains open.
\end{firewall}

\section*{Version V14 append-only record}
\addcontentsline{toc}{section}{Version V14 append-only record}
The complete V13 source was retained before this continuation.  Its SHA--256
digest was
\begin{center}\scriptsize\ttfamily
bc6de93e757dc7b02bda92ff48933c89f7c95a02d0b5a8039e0550566e44a0af.
\end{center}
On the next iteration, retain this full file, append before its unique active
document-closing command, increment the filename to
\texttt{\_V15.tex}, and remove only the superseded V14 file from this note
series.

% =============================================================================
% END APPENDED VERSION V14 MATERIAL
% =============================================================================

% =============================================================================
% BEGIN APPENDED VERSION V15 MATERIAL
% =============================================================================

\clearpage
\part{Version V15: the signed Perron curvature source and source-specific response columns}
\label{part:V15}

\section{Direction and nonduplication boundary}
\label{sec:V15-direction}

V14 reduced the spatial gap to a uniform conditional score diameter and a
bounded sequential response matrix.  Its response identity contains
\[
 \Cov_{\nu_{a,U}}(f,\mathcal J_i),
\]
where \(\mathcal J_i=\mathbb E_{p_x}(J_x)_i=-\nabla_i\log h\) is the physical
Perron score.  V15 differentiates this score, identifies its exact signed
source, and represents every response-matrix entry by one source-specific
future Witten Green column.

The supplied V35 programme already proves the comparison-inverse
Combes--Thomas estimate and the exact Woodbury curvature-defect return.  Those
theorems are not repeated.  Their new role here is precise: they are to bound
the Green columns generated by the signed V8 Perron curvature source, not the
norm of the complete inverse.  This distinction avoids using a full future
Poincar\'e constant to prove the full Poincar\'e constant.

% =============================================================================
\section{The exact signed Perron curvature source}
\label{sec:V15-curvature-source}
% =============================================================================

Retain
\[
 \mathcal J_i(x)
 :=
 \mathbb E_{p_x}(J_x)_i,
 \qquad
 \mathcal H_{bi}(x)
 :=
 \nabla_b\mathcal J_i(x).
\]
Here \(i,b\) label rooted tree-gauge spatial coordinates and the tensor block
acts from the \(b\)-tangent fibre to the \(i\)-cotangent fibre.

\begin{theorem}[Exact physical Perron curvature block]
\label{thm:V15-Perron-curvature}
For every pair \(b,i\),
\begin{equation}
 \boxed{
 \mathcal H_{bi}(x)
 =
 \mathbb E_{p_x}(K_x)_{bi}
 -
 \Cov_{p_x}\bigl((J_x)_i,(J_x)_b\bigr)
 =
 \frac12(\nabla^2W(x))_{bi}.}
 \label{eq:V15-Perron-curvature}
\end{equation}
Equivalently,
\begin{equation}
 \boxed{
 d\mathcal J
 =
 \mathbb E_{p_x}K_x-\Cov_{p_x}(J_x)
 =
 -\nabla^2\log h.}
 \label{eq:V15-full-curvature}
\end{equation}
The expectation and covariance are under the terminal-weighted conditional of
the same transfer kernel.
\end{theorem}

\begin{proof}
Apply the V8 conditional derivative identity
\eqref{eq:V8-conditional-response} to
\[
 g(x,y)=(J_x(y))_i
\]
and differentiate in the \(b\)-direction.  Since
\[
 \nabla_b(J_x)_i=(K_x)_{bi},
\]
one obtains the first equality in
\eqref{eq:V15-Perron-curvature}.  The V8 ground-potential identities give
\[
 \nabla W=2\mathbb E_{p_x}J_x=2\mathcal J,
\]
so differentiating gives the second equality.  Finally,
\(W=-2\log h+\mathrm{constant}\), which proves
\eqref{eq:V15-full-curvature}.
\end{proof}

\begin{corollary}[Symmetry after complete assembly]
\label{cor:V15-curvature-symmetry}
Although the two terms in the first expression of
\eqref{eq:V15-Perron-curvature} may be recorded in oriented blocks, their
complete difference satisfies
\begin{equation}
 \boxed{
 \mathcal H_{bi}=\mathcal H_{ib}^*.}
 \label{eq:V15-curvature-symmetry}
\end{equation}
\end{corollary}

\begin{proof}
The final expression in \eqref{eq:V15-Perron-curvature} is a Hessian of a
real scalar potential.
\end{proof}

\begin{firewall}[The covariance debit occurs once]
\label{fire:V15-no-double-covariance}
The source \(d\mathcal J\) already contains the complete negative score
covariance in \eqref{eq:V15-full-curvature}.  If that source is divided into a
finite signed head and a connected tail, the split must obey the full/head/tail
identity of supplied V35.  Appending the same differentiated covariance once
more to the complete source would double charge it and can reverse a spectral
verdict.
\end{firewall}

% =============================================================================
\section{The nested marginal curvature}
\label{sec:V15-nested-curvature}
% =============================================================================

In the notation of V14, define
\begin{equation}
 \overline{\mathcal J}_i(a,U)
 :=
 \mathbb E_{\nu_{a,U}}\mathcal J_i(a,U,\cdot)
 =
 -\nabla_U\log\widehat h_i(a,U).
 \label{eq:V15-marginal-mean-score}
\end{equation}

\begin{theorem}[Two-level marginal curvature identity]
\label{thm:V15-two-level-curvature}
Let \(b\le i\) be a held coordinate of the sequential marginal.  Then
\begin{equation}
 \boxed{
 \nabla_b\overline{\mathcal J}_i
 =
 \mathbb E_{\nu_{a,U}}\mathcal H_{bi}
 -
 2\Cov_{\nu_{a,U}}
 \bigl(\mathcal J_i,\mathcal J_b\bigr).}
 \label{eq:V15-two-level-curvature}
\end{equation}
Thus the curvature of the marginal conditional has two distinct response
levels:
\[
 \underbrace{
 \mathbb E_\nu
 \bigl[\mathbb E_pK-\Cov_p(J)\bigr]
 }_{\text{one-step physical curvature}}
 \quad-\quad
 \underbrace{
 2\Cov_\nu(\mathcal J,\mathcal J)
 }_{\text{future-marginal return}}.
\]
\end{theorem}

\begin{proof}
Differentiate the normalized expectation in
\eqref{eq:V15-marginal-mean-score}.  The derivative of
\(\mathcal J_i\) is \(\mathcal H_{bi}\).  The future density is proportional
to \(h^2\), so its logarithmic derivative in the held \(b\)-coordinate is
\[
 2\nabla_b\log h=-2\mathcal J_b.
\]
The standard derivative-of-expectation formula therefore gives
\[
 \nabla_b\mathbb E_\nu\mathcal J_i
 =
 \mathbb E_\nu\nabla_b\mathcal J_i
 -2\Cov_\nu(\mathcal J_i,\mathcal J_b),
\]
which is \eqref{eq:V15-two-level-curvature}.
\end{proof}

\begin{assessment}[One-step and future returns are not interchangeable]
\label{ass:V15-two-level-meaning}
The \(p_x\)-covariance in \eqref{eq:V15-Perron-curvature} is generated inside
one physical transfer step.  The \(\nu_{a,U}\)-covariance in
\eqref{eq:V15-two-level-curvature} is generated by integrating unrevealed
spatial coordinates of the ground law.  They have different conditioning
sigma algebras and must not be merged before their common-history incidence
is recorded.
\end{assessment}

% =============================================================================
\section{Source-specific future Witten columns}
\label{sec:V15-Green-columns}
% =============================================================================

Fix \(i,a,U\).  On the future carrier \(G^{N-i}\) with law
\(\nu=\nu_{a,U}\), let
\[
 L_\nu^{(0)}=d_\nu^*d,
 \qquad
 L_\nu^{(1)}=dd_\nu^*+d_\nu^*d,
\]
and let \(G_\nu^{(1)}\) be the Moore--Penrose inverse of
\(L_\nu^{(1)}\) on the closure of the future exact one-forms.  For a unit
vector \(v\in T_UG\), put
\begin{equation}
 g_{i,v}(y)
 :=
 \ip{\mathcal J_i(a,U,y)}{v}.
 \label{eq:V15-scalar-score-source}
\end{equation}
If \(b>i\), let \(P_b\) be the orthogonal projection onto the
\(b\)-coordinate component of a future one-form.

\begin{theorem}[Exact response-column representation]
\label{thm:V15-response-column}
For every smooth scalar \(f\),
\begin{equation}
 \boxed{
 \ip{
 \Cov_\nu(f,\mathcal J_i)
 }{v}
 =
 \ip{
 d_Yf
 }{
 G_\nu^{(1)}d_Yg_{i,v}
 }_{L^2(T^*G^{N-i},\nu)}.}
 \label{eq:V15-response-column}
\end{equation}
Moreover, its source is the signed physical curvature:
\begin{equation}
 \boxed{
 \nabla_bg_{i,v}
 =
 \mathcal H_{bi}^*v
 =
 \left[
  \mathbb E_{p_x}(K_x)_{bi}
  -\Cov_{p_x}\bigl((J_x)_i,(J_x)_b\bigr)
 \right]^*v.}
 \label{eq:V15-column-source}
\end{equation}
\end{theorem}

\begin{proof}
The exact Witten covariance identity proved in the supplied Grand Tensor V067
manuscript says, for any smooth scalars \(f,g\),
\[
 \Cov_\nu(f,g)
 =
 \ip{d_Yf}{G_\nu^{(1)}d_Yg}.
\]
Apply it componentwise to \(g=g_{i,v}\), obtaining
\eqref{eq:V15-response-column}.  Differentiating
\eqref{eq:V15-scalar-score-source} in the future coordinate \(b\) gives
\[
 \nabla_bg_{i,v}
 =
 (\nabla_b\mathcal J_i)^*v
 =\mathcal H_{bi}^*v.
\]
Insert \eqref{eq:V15-Perron-curvature}.
\end{proof}

\begin{definition}[Complete source-specific response coefficients]
\label{def:V15-response-coefficients}
For \(b>i\), define
\begin{equation}
 \boxed{
 R_{ib}
 :=
 2\operatorname*{ess\,sup}_{a,U}
 \sup_{\norm v=1}
 \left\|
 P_bG_{\nu_{a,U}}^{(1)}
 d_Yg_{i,v}
 \right\|_{L^2(\nu_{a,U})}.}
 \label{eq:V15-response-coefficients}
\end{equation}
Set \(R_{ib}=0\) for \(b\le i\).  These coefficients measure only the Green
columns excited by the complete Perron curvature sources
\eqref{eq:V15-column-source}.
\end{definition}

\begin{theorem}[Green columns populate the V14 response matrix]
\label{thm:V15-columns-to-matrix}
Let \(\mathsf R=(R_{ib})\) be
\eqref{eq:V15-response-coefficients} and
\(\mathsf A=I+\mathsf R\).  Then every smooth \(f\) satisfies
\begin{equation}
 \boxed{
 \norm{\nabla_iM_if}_{L^2(\pi)}
 \le
 \norm{\nabla_if}_{L^2(\pi)}
 +
 \sum_{b>i}R_{ib}
 \norm{\nabla_bf}_{L^2(\pi)}.}
 \label{eq:V15-populated-row}
\end{equation}
Hence \(\mathsf A\) satisfies the V14 response-matrix hypothesis
\eqref{eq:V14-response-matrix-row}.
\end{theorem}

\begin{proof}
For fixed \((a,U)\), decompose the Witten inner product in
\eqref{eq:V15-response-column} into future coordinate blocks.  Taking the
supremum over unit \(v\) and applying Cauchy--Schwarz in each block gives
\[
 2\norm{\Cov_\nu(f,\mathcal J_i)}
 \le
 \sum_{b>i}R_{ib}
 \left(
  \int|\nabla_bf|^2\,\dd\nu_{a,U}
 \right)^{1/2}.
\]
Insert this in the exact nested response identity
\eqref{eq:V14-nested-response}.  Conditional Jensen bounds the direct term
by the conditional \(L^2\)-norm of \(\nabla_if\).  Now take the
\(L^2\)-norm in \((a,U)\) and use Minkowski's inequality.  Disintegration of
\(\pi\) gives
\[
 \left\|
 \left(
  \int|\nabla_bf|^2\,\dd\nu_{a,U}
 \right)^{1/2}
 \right\|_{L^2(a,U)}
 =
 \norm{\nabla_bf}_{L^2(\pi)}.
\]
This proves \eqref{eq:V15-populated-row}.
\end{proof}

\begin{firewall}[Source-specific inverse, not inverse norm]
\label{fire:V15-no-full-inverse}
The coefficient \(R_{ib}\) contains
\[
 P_bG_\nu^{(1)}d_Yg_{i,v},
\]
not \(\norm{G_\nu^{(1)}}\).  Bounding it by
\(\norm{G_\nu^{(1)}}\norm{d_Yg_{i,v}}\) would import the smallest future
one-form eigenvalue and can be circular.  The V35 comparison/return machinery
must instead be applied to these declared curvature-source columns.
\end{firewall}

% =============================================================================
\section{A summable Green-column mass theorem}
\label{sec:V15-summable}
% =============================================================================

Let \(\Gamma_N\) be the bounded-degree coordinate graph of the rooted
tree-gauge variables.  Assume its distance shells satisfy
\begin{equation}
 \sup_{N,i}
 \#\{b:\operatorname{dist}_{\Gamma_N}(i,b)=r\}
 \le n(r).
 \label{eq:V15-shell-growth}
\end{equation}

\begin{theorem}[Spatially summable response-column certificate]
\label{thm:V15-summable-columns}
Suppose there is a nonnegative sequence \(\mathcal R(r)\) such that
\begin{equation}
 R_{ib}
 \le
 \mathcal R\!\left(
 \operatorname{dist}_{\Gamma_N}(i,b)
 \right)
 \quad(b>i)
 \label{eq:V15-column-decay}
\end{equation}
and
\begin{equation}
 \boxed{
 S_{\mathcal R}
 :=
 \sum_{r\ge0}n(r)\mathcal R(r)<\infty.}
 \label{eq:V15-summable-tail}
\end{equation}
Then
\begin{equation}
 \sup_i\sum_bR_{ib}\le S_{\mathcal R},
 \qquad
 \sup_b\sum_iR_{ib}\le S_{\mathcal R},
 \qquad
 \norm{I+\mathsf R}_{2\to2}\le1+S_{\mathcal R}.
 \label{eq:V15-row-column-sum}
\end{equation}
If also
\[
 \sup_{N,i,a}\mathscr A_i(a)\le A_*,
\]
then
\begin{equation}
 \boxed{
 \gap(H_N)
 \ge
 \frac{\kappa C_Fe^{-2A_*}}
      {(1+S_{\mathcal R})^2}.}
 \label{eq:V15-summable-gap}
\end{equation}
\end{theorem}

\begin{proof}
For fixed \(i\), group the \(b\)'s by their distance from \(i\) and use
\eqref{eq:V15-shell-growth}--\eqref{eq:V15-summable-tail}; this gives the row
bound.  For fixed \(b\), the same shell count bounds the subset \(i<b\), giving
the column bound.  Apply the Schur test and then
\Cref{thm:V14-coordinate-gap,cor:V14-Schur-response}.
\end{proof}

\begin{corollary}[Exponential Green-column criterion]
\label{cor:V15-exponential-columns}
If, for constants \(C,\mu>0\),
\[
 \mathcal R(r)\le Ce^{-\mu r}
\]
and the coordinate graph has polynomial shell growth, then
\(S_{\mathcal R}<\infty\), uniformly in volume.
\end{corollary}

\begin{proof}
A polynomial times \(e^{-\mu r}\) is summable.
\end{proof}

\begin{proposition}[Exact comparison/return split of a response column]
\label{prop:V15-comparison-return-column}
Suppose the supplied V35 Woodbury identity on the same future conditional law
has the form
\[
 G_\nu^{(1)}
 =
 G_{\nu,\mathrm{cmp}}^{(1)}
 +
 G_{\nu,\mathrm{ret}}^{(1)}
\]
on the source-cyclic range generated by \(d_Yg_{i,v}\).  If
\begin{align}
 \sup_{\norm v=1}
 \left\|
 P_bG_{\nu,\mathrm{cmp}}^{(1)}
 d_Yg_{i,v}
 \right\|
 &\le u_{ib},
 \label{eq:V15-comparison-column}\\
 \sup_{\norm v=1}
 \left\|
 P_bG_{\nu,\mathrm{ret}}^{(1)}
 d_Yg_{i,v}
 \right\|
 &\le v_{ib},
 \label{eq:V15-return-column}
\end{align}
uniformly in \((a,U)\), then
\begin{equation}
 \boxed{
 R_{ib}\le2(u_{ib}+v_{ib}).}
 \label{eq:V15-split-column}
\end{equation}
The V35 cutoff-independent comparison Combes--Thomas theorem supplies an
exponential candidate for \(u_{ib}\) on its stated small-field window.
The unpopulated physical quantity is the source-cyclic return column
\(v_{ib}\), together with the chart-exit sector.
\end{proposition}

\begin{proof}
Insert the comparison/return decomposition into
\eqref{eq:V15-response-coefficients} and apply the triangle inequality.
The statement about \(u_{ib}\) is the direct specialization of the supplied
V35 weighted comparison-inverse theorem to the source
\eqref{eq:V15-column-source}; no new comparison theorem is asserted here.
\end{proof}

\begin{assessment}[The response target is now fully typed]
\label{ass:V15-typed-target}
The V14 matrix entry is no longer an unspecified dependence coefficient.  It
is the norm of a declared physical Witten Green column, its source is the
signed tensor
\[
 \mathbb E_pK-\Cov_p(J),
\]
and its spatial tail is split exactly once into the V35 comparison and return
columns.  Failure of summability therefore returns either a nonlocal signed
source, a comparison-locality failure, or a source-cyclic return survivor.
\end{assessment}

% =============================================================================
\section{V15 ledger and revised frontier}
\label{sec:V15-ledger}
% =============================================================================

\begin{theorem}[Integrated V15 statement]
\label{thm:V15-integrated}
Preserving V1--V14, V15 proves:
\begin{enumerate}[label=\textup{(V15.\arabic*)},leftmargin=5.5em]
\item the exact signed physical Perron curvature
      \(\mathbb E_pK-\Cov_p(J)=\frac12\nabla^2W\);
\item symmetry of the completely assembled curvature card;
\item the two-level marginal curvature identity;
\item the exact future Witten representation of every sequential response;
\item identification of its source with the signed Perron curvature column;
\item complete source-specific coefficients \(R_{ib}\) which populate the
      V14 response matrix;
\item a spatial summability theorem giving the explicit mass lower bound
      \eqref{eq:V15-summable-gap};
\item a comparison/return column split landing the supplied V35 inverse
      locality on the correct physical source.
\end{enumerate}
\end{theorem}

\begin{proof}
Items \textup{(V15.1)--(V15.2)} are
\Cref{thm:V15-Perron-curvature,cor:V15-curvature-symmetry}.  Item
\textup{(V15.3)} is \Cref{thm:V15-two-level-curvature}.  Items
\textup{(V15.4)--(V15.5)} are \Cref{thm:V15-response-column}.  Item
\textup{(V15.6)} is
\Cref{def:V15-response-coefficients,thm:V15-columns-to-matrix}.  Item
\textup{(V15.7)} is \Cref{thm:V15-summable-columns}.  Item
\textup{(V15.8)} is \Cref{prop:V15-comparison-return-column}.
\end{proof}

\begingroup
\small
\begin{longtable}{>{\raggedright\arraybackslash}p{0.28\textwidth}
                  >{\raggedright\arraybackslash}p{0.16\textwidth}
                  >{\raggedright\arraybackslash}p{0.48\textwidth}}
\toprule
\textbf{V15 row}&\textbf{Status}&\textbf{Advance or remaining obstruction}\\
\midrule
\endfirsthead
\toprule
\textbf{V15 row}&\textbf{Status}&\textbf{Advance or remaining obstruction}\\
\midrule
\endhead
Signed Perron curvature source
& \MGpass
& Exact terminal-weighted expectation minus covariance.\\[0.35em]
Nested marginal curvature
& \MGpass
& Separates one-step curvature from the future-marginal return.\\[0.35em]
Response Green column
& \MGpass
& Exact source-specific Witten representation.\\[0.35em]
Response-matrix incidence
& \MGpass
& The coefficients \eqref{eq:V15-response-coefficients} populate V14 without a full inverse norm.\\[0.35em]
Comparison column
& Conditional \MGpass
& V35 gives exponential decay on its stated small-field comparison window.\\[0.35em]
Physical return column
& \MGopen
& Requires a complete source-cyclic edge/tail bound on the actual future conditional law.\\[0.35em]
Score diameter
& \MGopen
& The V14 physical score-moment stock remains unbounded at weak coupling.\\[0.35em]
Chart exit and protected sectors
& \MGopen
& Must be combined with V12--V13 capacity and sector shorts.\\[0.35em]
Full Yang--Mills mass gap
& \MGopen
& Uniform weak-coupling response columns and the nontrivial OS limit remain.\\
\bottomrule
\end{longtable}
\endgroup

\begin{openproblem}[V16 mandate: bound the physical return columns]
\label{op:V16}
The source and target of every response coefficient are now exact.  Proceed
without another abstract inverse.
\begin{enumerate}[label=\textup{(V16.\arabic*)},leftmargin=5.5em]
\item On each future conditional law, insert the V35 Woodbury formula only on
      the source-cyclic carrier generated by
      \(d_Yg_{i,v}\).
\item Use the V35 weighted comparison inverse to record the complete
      exponentially decaying comparison column \(u_{ib}\).
\item Bound the return column \(v_{ib}\) from a regulator-uniform
      source-cyclic edge strictly below one and its weighted return tail.
\item Split the signed source
      \(\mathbb E_pK-\Cov_p(J)\) into a declared finite-range head and one
      complete connected tail; do not append a second covariance debit.
\item Prove \(\sum_rn(r)\mathcal R(r)<\infty\) and the V14 score-diameter
      bound, or return the normalized column which fails.
\item Classify every returned survivor as small-field local, chart-exit,
      toron, wrapping, topology, or Gribov before invoking
      \eqref{eq:V15-summable-gap}.
\end{enumerate}
Using \(\norm{G_\nu^{(1)}}\) in place of the declared columns is not an
admissible shortcut.
\end{openproblem}

\begin{firewall}[End-of-V15 claim boundary]
V15 closes the algebraic and same-law incidence of the weak-coupling response
matrix and lands the V35 comparison/return split on its exact physical source.
It does not prove a uniform return edge, summability of the return columns,
the score-diameter bound, protected-sector control, or a nontrivial continuum
OS state.  The full Yang--Mills mass gap remains open.
\end{firewall}

\section*{Version V15 append-only record}
\addcontentsline{toc}{section}{Version V15 append-only record}
The complete V14 source was retained before this continuation.  Its SHA--256
digest was
\begin{center}\scriptsize\ttfamily
176cf72e7534809a05e94f20d1a9d8b9d82a69592d71e75eccf595d834efc110.
\end{center}
On the next iteration, retain this full file, append before its unique active
document-closing command, increment the filename to
\texttt{\_V16.tex}, and remove only the superseded V15 file from this note
series.

% =============================================================================
% END APPENDED VERSION V15 MATERIAL
% =============================================================================

% =============================================================================
% BEGIN APPENDED VERSION V16 MATERIAL
% =============================================================================

\clearpage
\part{Version V16: spectral-edge Combes--Thomas control of the physical return columns}
\label{part:V16}

\section{Direction and exact advance beyond the supplied weighted theorem}
\label{sec:V16-direction}

V15 identified every spatial response coefficient with a source-specific
future Witten Green column.  The supplied V35 weighted-return theorem closes
such columns if the weighted Schur norm of the defect channel itself is
strictly below one.  V35 also correctly warns that this weighted condition is
stronger than the automatic strict edge on each finite card.

V16 proves a different sufficient route.  A uniform \emph{unweighted}
spectral edge on the cell-saturated source carrier, together with one uniform
exponential Schur moment of the channel kernel, implies exponential decay of
\((I-\mathcal K)^{-1}\).  The proof is a Combes--Thomas conjugation at the
channel level.  It does not require
\(\norm{\mathcal K}_{\mu,\mathrm{Sch}}<1\).

This reduces the return obstruction to the physically meaningful question:
does the source-generated defect channel have a regulator-uniform spectral
edge below one?  The remaining localization hypotheses are already natural
outputs of the V35 comparison inverse.

% =============================================================================
\section{The cell-saturated source carrier}
\label{sec:V16-cell-carrier}
% =============================================================================

Let
\[
 \mathscr C=\bigoplus_{x\in\Gamma}\mathscr C_x
\]
be a finite defect-channel Hilbert space, \(P_x\) its cell projections, and
\[
 0\preceq\mathcal K\prec I
\]
the V35 defect channel.  Let \(\mathcal Y\) be any declared source entrance.

\begin{definition}[Cell-saturated source carrier]
\label{def:V16-cell-saturated}
The cell-saturated source carrier is the smallest closed subspace
\(\mathscr C_{\mathcal Y}^{\rm cell}\subseteq\mathscr C\) which contains
\(\Ran\mathcal Y\) and is invariant under
\[
 \mathcal K,\qquad P_x\quad(x\in\Gamma).
\]
Equivalently, it is the closed span of all words in
\(\{\mathcal K,P_x:x\in\Gamma\}\) applied to \(\Ran\mathcal Y\).
\end{definition}

\begin{lemma}[Reduction and weight stability]
\label{lem:V16-carrier-reduction}
The carrier \(\mathscr C_{\mathcal Y}^{\rm cell}\) reduces
\(\mathcal K\) and every \(P_x\).  Consequently it is preserved by each cell
weight
\begin{equation}
 W_y(\mu)
 :=
 \sum_{x\in\Gamma}
 e^{\mu\operatorname{dist}(x,y)}P_x.
 \label{eq:V16-cell-weight}
\end{equation}
\end{lemma}

\begin{proof}
The carrier is invariant under the self-adjoint operators
\(\mathcal K\) and \(P_x\).  A closed invariant subspace of a self-adjoint
operator is reducing.  Each \(W_y(\mu)\) is a finite linear combination of
the \(P_x\)'s, so it preserves the carrier, as does its inverse.
\end{proof}

\begin{assessment}[Why the ordinary cyclic carrier is enlarged]
\label{ass:V16-cell-saturation}
The ordinary source-cyclic carrier
\(\overline{\Span}\{\mathcal K^n\Ran\mathcal Y\}\) need not be invariant
under the spatial cell projections.  Exponential conjugation uses those
projections.  The cell-saturated carrier is the smallest source-generated
space on which both the spectral edge and spatial weights can be used
without silently leaving the controlled subspace.
\end{assessment}

% =============================================================================
\section{Spectral edge implies exponential return-resolvent decay}
\label{sec:V16-channel-CT}
% =============================================================================

Restrict all operators in this section to
\(\mathscr C_{\mathcal Y}^{\rm cell}\), and write
\[
 \rho:=\norm{\mathcal K}.
\]
For \(\mu>0\), define the conjugation modulus
\begin{equation}
 \delta_{\mathcal K}(\mu)
 :=
 \sup_{y\in\Gamma}
 \left\|
 W_y(\mu)\mathcal K W_y(\mu)^{-1}
 -\mathcal K
 \right\|.
 \label{eq:V16-conjugation-modulus}
\end{equation}

\begin{theorem}[Source-carrier channel Combes--Thomas bound]
\label{thm:V16-channel-CT}
Suppose
\begin{equation}
 \rho<1,
 \qquad
 \delta_{\mathcal K}(\mu)<1-\rho.
 \label{eq:V16-channel-CT-condition}
\end{equation}
Then
\begin{align}
 \boxed{
 \sup_y
 \left\|
 W_y(\mu)(I-\mathcal K)^{-1}W_y(\mu)^{-1}
 \right\|
 &\le
 \frac1{1-\rho-\delta_{\mathcal K}(\mu)},}
 \label{eq:V16-conjugated-resolvent}\\
 \boxed{
 \left\|
 P_x(I-\mathcal K)^{-1}P_y
 \right\|
 &\le
 \frac{
 e^{-\mu\operatorname{dist}(x,y)}
 }{1-\rho-\delta_{\mathcal K}(\mu)}.}
 \label{eq:V16-return-resolvent-decay}
\end{align}
\end{theorem}

\begin{proof}
Since \(0\preceq\mathcal K\preceq\rho I\),
\[
 I-\mathcal K\succeq(1-\rho)I.
\]
Put
\[
 \mathcal E_y
 :=
 W_y\mathcal K W_y^{-1}-\mathcal K.
\]
Then
\[
 W_y(I-\mathcal K)W_y^{-1}
 =
 (I-\mathcal K)-\mathcal E_y.
\]
The inverse of \(I-\mathcal K\) has norm at most \((1-\rho)^{-1}\).
The perturbation bound in
\eqref{eq:V16-channel-CT-condition} and the Neumann series give
\[
 \left\|
 \bigl[(I-\mathcal K)-\mathcal E_y\bigr]^{-1}
 \right\|
 \le
 \frac1{1-\rho-\norm{\mathcal E_y}}.
\]
Taking the supremum over \(y\) proves
\eqref{eq:V16-conjugated-resolvent}.

On the \(y\)-cell the weight equals one, while on the \(x\)-cell it equals
\(e^{\mu\operatorname{dist}(x,y)}\).  Hence
\[
 P_x(I-\mathcal K)^{-1}P_y
 =
 e^{-\mu\operatorname{dist}(x,y)}
 P_xW_y(I-\mathcal K)^{-1}W_y^{-1}P_y.
\]
Use \eqref{eq:V16-conjugated-resolvent}.
\end{proof}

\begin{lemma}[Schur control of the conjugation modulus]
\label{lem:V16-conjugation-Schur}
Write \(\mathcal K_{xz}=P_x\mathcal KP_z\), and define
\begin{equation}
 \omega_{\mathcal K}(\mu)
 :=
 \max\left\{
 \sup_x\sum_z
 \bigl(e^{\mu\operatorname{dist}(x,z)}-1\bigr)
 \norm{\mathcal K_{xz}},
 \ 
 \sup_z\sum_x
 \bigl(e^{\mu\operatorname{dist}(x,z)}-1\bigr)
 \norm{\mathcal K_{xz}}
 \right\}.
 \label{eq:V16-K-modulus}
\end{equation}
Then
\begin{equation}
 \boxed{
 \delta_{\mathcal K}(\mu)
 \le\omega_{\mathcal K}(\mu).}
 \label{eq:V16-delta-omega}
\end{equation}
\end{lemma}

\begin{proof}
The \(xz\)-block of the conjugation difference is
\[
 \left[
 e^{\mu(
 \operatorname{dist}(x,y)-\operatorname{dist}(z,y))}
 -1
 \right]\mathcal K_{xz}.
\]
The triangle inequality for the graph distance gives
\[
 \left|
 e^{\mu(
 \operatorname{dist}(x,y)-\operatorname{dist}(z,y))}
 -1
 \right|
 \le
 e^{\mu\operatorname{dist}(x,z)}-1.
\]
Apply the row/column Schur test and take the supremum over \(y\).
\end{proof}

\begin{corollary}[A positive decay exponent from an unweighted edge]
\label{cor:V16-edge-to-exponent}
Suppose, uniformly in the regulator,
\begin{equation}
 \norm{\mathcal K}\le\rho_*<1
 \label{eq:V16-uniform-edge}
\end{equation}
on the cell-saturated source carrier, and for some \(\mu_0>0\),
\begin{equation}
 M_{\mu_0}
 :=
 \max\left\{
 \sup_x\sum_z
 e^{\mu_0\operatorname{dist}(x,z)}
 \norm{\mathcal K_{xz}},
 \ 
 \sup_z\sum_x
 e^{\mu_0\operatorname{dist}(x,z)}
 \norm{\mathcal K_{xz}}
 \right\}
 <\infty.
 \label{eq:V16-exponential-K-moment}
\end{equation}
If \(M_{\mu_0}>0\), choose
\begin{equation}
 \boxed{
 0<\mu_*
 \le
 \min\left\{
 \mu_0,\,
 \frac{\mu_0(1-\rho_*)}{2M_{\mu_0}}
 \right\}.}
 \label{eq:V16-mu-choice}
\end{equation}
Then
\begin{equation}
 \boxed{
 \delta_{\mathcal K}(\mu_*)
 \le\frac{1-\rho_*}{2},
 \qquad
 \left\|
 P_x(I-\mathcal K)^{-1}P_y
 \right\|
 \le
 \frac{2e^{-\mu_*\operatorname{dist}(x,y)}}
      {1-\rho_*}.}
 \label{eq:V16-edge-decay}
\end{equation}
If \(M_{\mu_0}=0\), then \(\mathcal K=0\) and any
\(0<\mu_*\le\mu_0\) is admissible.
\end{corollary}

\begin{proof}
For \(0\le\mu\le\mu_0\), convexity in \(\mu\) gives
\[
 e^{\mu r}-1
 \le
 \frac{\mu}{\mu_0}
 \left(e^{\mu_0r}-1\right)
 \le
 \frac{\mu}{\mu_0}e^{\mu_0r}.
\]
Thus
\[
 \omega_{\mathcal K}(\mu)
 \le\frac{\mu}{\mu_0}M_{\mu_0}.
\]
Use \eqref{eq:V16-mu-choice} and
\Cref{lem:V16-conjugation-Schur} to obtain the first inequality in
\eqref{eq:V16-edge-decay}.  Apply
\Cref{thm:V16-channel-CT} with
\(\rho\le\rho_*\).  The zero case is immediate from
\eqref{eq:V16-exponential-K-moment}.
\end{proof}

\begin{assessment}[Improvement over a weighted contraction hypothesis]
\label{ass:V16-improvement}
The conclusion does not require
\[
 \norm{\mathcal K}_{\mu,\mathrm{Sch}}<1.
\]
It uses instead the spectral quantity
\(\norm{\mathcal K}\le\rho_*<1\) on the cell-saturated source carrier and a
locality modulus which tends to zero with \(\mu\).  A crude weighted Schur
norm may exceed one because it counts many harmless off-diagonal blocks even
when the source-generated spectral edge is safely below one.
\end{assessment}

% =============================================================================
\section{Application to the V15 physical Green columns}
\label{sec:V16-application}
% =============================================================================

On one future conditional law, write the supplied V35 Woodbury identity as
\begin{equation}
 \boxed{
 G_\nu^{(1)}
 =
 \mathcal A^{-1}
 +
 \mathcal B(I-\mathcal K)^{-1}\mathcal B^*,
 \qquad
 \mathcal B:=\mathcal A^{-1}\mathcal V.}
 \label{eq:V16-Woodbury}
\end{equation}
Let \(W_y^{\mathscr H}\) and \(W_y^{\mathscr C}\) be the corresponding cell
weights on the one-form and channel spaces.  Define the two-sided weighted
entrance bound
\begin{equation}
 B_\mu
 :=
 \sup_y\max\left\{
 \left\|
 W_y^{\mathscr H}\mathcal B
 (W_y^{\mathscr C})^{-1}
 \right\|,
 \left\|
 (W_y^{\mathscr H})^{-1}\mathcal B
 W_y^{\mathscr C}
 \right\|
 \right\}.
 \label{eq:V16-B-weight}
\end{equation}
Also put
\begin{equation}
 C_{\mathcal A,\mu}
 :=
 \sup_y
 \left\|
 W_y^{\mathscr H}\mathcal A^{-1}
 (W_y^{\mathscr H})^{-1}
 \right\|.
 \label{eq:V16-A-weight}
\end{equation}

For the signed V15 source
\[
 q_{i,v}:=d_Yg_{i,v},
\]
define
\begin{equation}
 Q_\mu
 :=
 \operatorname*{ess\,sup}_{a,U,i}
 \sup_{\norm v=1}
 \left\|
 W_i^{\mathscr H}(\mu)q_{i,v}
 \right\|.
 \label{eq:V16-source-moment}
\end{equation}

\begin{theorem}[Explicit physical response-column decay]
\label{thm:V16-response-decay}
Assume, uniformly in all future conditionals:
\begin{enumerate}[label=\textup{(\roman*)},leftmargin=3.6em]
\item the cell-saturated channel edge and exponential moment of
      \Cref{cor:V16-edge-to-exponent};
\item \(B_{\mu_*}<\infty\),
      \(C_{\mathcal A,\mu_*}<\infty\), and \(Q_{\mu_*}<\infty\).
\end{enumerate}
Then
\begin{equation}
 \boxed{
 \left\|
 P_bG_\nu^{(1)}q_{i,v}
 \right\|
 \le
 Q_{\mu_*}
 \left[
 C_{\mathcal A,\mu_*}
 +
 \frac{2B_{\mu_*}^2}{1-\rho_*}
 \right]
 e^{-\mu_*\operatorname{dist}(b,i)}.}
 \label{eq:V16-physical-column-decay}
\end{equation}
Consequently the V15 response coefficients satisfy
\begin{equation}
 \boxed{
 R_{ib}
 \le
 2Q_{\mu_*}
 \left[
 C_{\mathcal A,\mu_*}
 +
 \frac{2B_{\mu_*}^2}{1-\rho_*}
 \right]
 e^{-\mu_*\operatorname{dist}(b,i)}.}
 \label{eq:V16-Rib-decay}
\end{equation}
\end{theorem}

\begin{proof}
Fix the source cell \(i\) and conjugate
\eqref{eq:V16-Woodbury} by \(W_i^{\mathscr H}\).  The comparison term is
bounded by
\[
 C_{\mathcal A,\mu_*}
 \norm{W_i^{\mathscr H}q_{i,v}}.
\]
For the return term, insert \(W_i^{\mathscr C}\) and its inverse on both sides
of the channel resolvent.  The two oriented entrance bounds in
\eqref{eq:V16-B-weight}, together with
\eqref{eq:V16-conjugated-resolvent} and
\eqref{eq:V16-edge-decay}, give
\[
 \left\|
 W_i^{\mathscr H}
 \mathcal B(I-\mathcal K)^{-1}\mathcal B^*
 q_{i,v}
 \right\|
 \le
 \frac{2B_{\mu_*}^2}{1-\rho_*}
 \norm{W_i^{\mathscr H}q_{i,v}}.
\]
The \(b\)-cell component of \(W_i^{\mathscr H}\) carries the factor
\(e^{\mu_*\operatorname{dist}(b,i)}\).  Use
\eqref{eq:V16-source-moment} and remove that factor to prove
\eqref{eq:V16-physical-column-decay}.  Multiply by the factor \(2\) in the
definition \eqref{eq:V15-response-coefficients} to obtain
\eqref{eq:V16-Rib-decay}.
\end{proof}

\begin{corollary}[Closed return-to-response summability]
\label{cor:V16-response-summability}
Assume the coordinate graphs have polynomial shell growth, and the hypotheses
of \Cref{thm:V16-response-decay} hold with common constants.  Then the V15
summability stock satisfies
\begin{equation}
 \boxed{
 S_{\mathcal R}
 \le
 2Q_{\mu_*}
 \left[
 C_{\mathcal A,\mu_*}
 +
 \frac{2B_{\mu_*}^2}{1-\rho_*}
 \right]
 \sum_{r\ge0}n(r)e^{-\mu_*r}
 <\infty.}
 \label{eq:V16-explicit-SR}
\end{equation}
If additionally the V14 score diameter is bounded by \(A_*\), then
\begin{equation}
 \boxed{
 \gap(H_N)
 \ge
 \frac{\kappa C_Fe^{-2A_*}}
      {(1+S_{\mathcal R})^2}}
 \label{eq:V16-explicit-gap}
\end{equation}
uniformly in spatial volume.
\end{corollary}

\begin{proof}
Use \eqref{eq:V16-Rib-decay} in
\Cref{thm:V15-summable-columns}, then apply
\eqref{eq:V15-summable-gap}.
\end{proof}

\begin{corollary}[Typed failure of physical return summability]
\label{cor:V16-return-failure}
Suppose the V35 comparison constants
\(B_{\mu_*}\) and \(C_{\mathcal A,\mu_*}\), the channel exponential moment
\(M_{\mu_0}\), and the signed-source moment \(Q_{\mu_*}\) remain uniformly
bounded.  If the physical response columns fail every common exponential
summability estimate, then along a subsequence either:
\begin{enumerate}[label=\textup{(\alph*)},leftmargin=3.5em]
\item the cell-saturated source edge satisfies
      \(\norm{\mathcal K}\to1\); or
\item no common positive exponent survives because one of the asserted
      locality moments actually loses uniformity.
\end{enumerate}
The first case returns a source-generated near-unit defect-channel vector;
the second returns a loss of physical spatial locality.
\end{corollary}

\begin{proof}
If a common edge \(\rho_*<1\) and all listed locality moments held,
\Cref{cor:V16-edge-to-exponent,thm:V16-response-decay} would give a common
positive exponent and \Cref{cor:V16-response-summability} would give
summability.  The contrapositive proves the alternative.
\end{proof}

\begin{firewall}[Finite-card strictness is not the uniform edge]
\label{fire:V16-finite-edge}
V35 proves \(\norm{\mathcal K}<1\) at each fixed finite card.  V16 requires a
common \(\rho_*<1\) on the cell-saturated carriers of all V15 curvature
sources.  Compactness at each regulator does not supply that common margin.
A sequence with \(\norm{\mathcal K_j}\uparrow1\) is the precise remaining
return obstruction.
\end{firewall}

\begin{firewall}[The signed-source moment is independent]
\label{fire:V16-source-moment}
Even with a uniform return edge, the input
\[
 Q_{\mu_*}
 =
 \sup\norm{W_i d_Yg_{i,v}}
\]
must be finite.  By V15 its source is
\(\mathbb E_pK-\Cov_p(J)\).  The V35 full/head/tail rule controls how that
signed source is assembled, but it does not prove a weak-coupling exponential
moment for it.
\end{firewall}

% =============================================================================
\section{V16 ledger and revised frontier}
\label{sec:V16-ledger}
% =============================================================================

\begin{theorem}[Integrated V16 statement]
\label{thm:V16-integrated}
Preserving V1--V15, V16 proves:
\begin{enumerate}[label=\textup{(V16.\arabic*)},leftmargin=5.5em]
\item the minimal cell-saturated carrier needed for simultaneous source and
      spatial-weight control;
\item a channel-level Combes--Thomas theorem for
      \((I-\mathcal K)^{-1}\);
\item a Schur modulus bound for the conjugation error;
\item existence of a common exponential decay rate from a uniform unweighted
      spectral edge and one exponential kernel moment;
\item an explicit decay estimate for the complete physical V15 Green
      columns;
\item an explicit response-summability constant and finite-Hamiltonian mass
      lower bound;
\item a typed alternative returning either a near-unit source-generated
      channel vector or a genuine loss of spatial locality.
\end{enumerate}
\end{theorem}

\begin{proof}
Item \textup{(V16.1)} is
\Cref{def:V16-cell-saturated,lem:V16-carrier-reduction}.  Items
\textup{(V16.2)--(V16.3)} are
\Cref{thm:V16-channel-CT,lem:V16-conjugation-Schur}.  Item
\textup{(V16.4)} is \Cref{cor:V16-edge-to-exponent}.  Items
\textup{(V16.5)--(V16.6)} are
\Cref{thm:V16-response-decay,cor:V16-response-summability}.  Item
\textup{(V16.7)} is \Cref{cor:V16-return-failure}.
\end{proof}

\begingroup
\small
\begin{longtable}{>{\raggedright\arraybackslash}p{0.28\textwidth}
                  >{\raggedright\arraybackslash}p{0.16\textwidth}
                  >{\raggedright\arraybackslash}p{0.48\textwidth}}
\toprule
\textbf{V16 row}&\textbf{Status}&\textbf{Advance or remaining obstruction}\\
\midrule
\endfirsthead
\toprule
\textbf{V16 row}&\textbf{Status}&\textbf{Advance or remaining obstruction}\\
\midrule
\endhead
Cell-saturated source carrier
& \MGpass
& Minimal carrier preserved by both the defect channel and spatial weights.\\[0.35em]
Channel Combes--Thomas bound
& \MGpass
& Converts an unweighted source edge plus locality into exponential resolvent decay.\\[0.35em]
Weighted-contraction replacement
& \MGpass
& No hypothesis \(\norm{\mathcal K}_{\mu,\mathrm{Sch}}<1\) is required.\\[0.35em]
Physical Green-column decay
& Conditional \MGpass
& Explicit bound \eqref{eq:V16-physical-column-decay}.\\[0.35em]
Uniform cell-saturated edge
& \MGopen
& Finite-card strictness does not give one \(\rho_*<1\) cofinally.\\[0.35em]
Signed-source exponential moment
& \MGopen
& \(Q_{\mu_*}\) remains unproved for \(\mathbb E_pK-\Cov_p(J)\) at weak coupling.\\[0.35em]
Score diameter
& \MGopen
& The V14 local score-moment row remains independent.\\[0.35em]
Protected sectors and OS limit
& \MGopen
& Toron/wrapping shorts and nontrivial continuum convergence remain.\\[0.35em]
Full Yang--Mills mass gap
& \MGopen
& The uniform edge, signed source moment, score diameter, and OS construction remain.\\
\bottomrule
\end{longtable}
\endgroup

\begin{openproblem}[V17 mandate: edge-or-survivor and signed-source moment]
\label{op:V17}
The weighted return tail is no longer an independent assumption.  Proceed
upstream to the two physical values which generate it.
\begin{enumerate}[label=\textup{(V17.\arabic*)},leftmargin=5.5em]
\item On the cell-saturated carrier generated by all V15 curvature sources,
      prove a regulator-uniform spectral edge
      \(\norm{\mathcal K_j}\le\rho_*<1\), or return normalized
      \(z_j\) with
      \(\ip{z_j}{\mathcal K_jz_j}\to1\).
\item Use the equality case of the V35 Schur complement to map every such
      near-unit vector back to a physical low-Witten-energy one-form.
\item Prove a uniform exponential moment \(Q_{\mu_*}\) for the complete
      signed source
      \(\mathbb E_pK-\Cov_p(J)\), respecting the V35 no-double-charging split.
\item Bound the V14 score diameter from a local physical score-moment
      estimate after the V11 temporal limit.
\item If all three stocks close, use
      \eqref{eq:V16-explicit-gap}; otherwise classify the returned one-form as
      local, chart-exit, toron, wrapping, topology, or Gribov.
\item Only then perform the V1 nontrivial OS passage.
\end{enumerate}
\end{openproblem}

\begin{firewall}[End-of-V16 claim boundary]
V16 proves that a uniform unweighted source-generated defect edge and
exponential locality are sufficient for the complete physical return-column
summability.  It does not prove that uniform edge, the signed-source moment,
the score-diameter estimate, protected-sector control, or a nontrivial
continuum OS state.  The full Yang--Mills mass gap remains open.
\end{firewall}

\section*{Version V16 append-only record}
\addcontentsline{toc}{section}{Version V16 append-only record}
The complete V15 source was retained before this continuation.  Its SHA--256
digest was
\begin{center}\scriptsize\ttfamily
a50eca5a1b434a27cff97d7afffed617fd0c1983f2f1f3ecb7055bd2a031d2c1.
\end{center}
On the next iteration, retain this full file, append before its unique active
document-closing command, increment the filename to
\texttt{\_V17.tex}, and remove only the superseded V16 file from this note
series.

% =============================================================================
% END APPENDED VERSION V16 MATERIAL
% =============================================================================

% =============================================================================
% BEGIN APPENDED VERSION V17 MATERIAL
% =============================================================================

\clearpage
\part{Version V17: exact edge landing, the infrared/ultraviolet alternative, and a Perron Fisher sum rule}
\label{part:V17}

\section{Direction and correction of the near-unit interpretation}
\label{sec:V17-direction}

V16 reduced return summability to a uniform spectral edge for the defect
channel.  The supplied V35 theorem correctly returns channel vectors
\(z_j\) when that edge approaches one, but it does not identify such a vector
with a normalized physical one-form.  The Woodbury landing can lose its
ordinary \(L^2\) norm while retaining unit comparison energy.  In that case
the survivor escapes to the ultraviolet comparison spectrum rather than
immediately producing a physical infrared mass obstruction.

V17 computes the landing exactly.  It proves:
\begin{enumerate}[label=\textup{(\roman*)},leftmargin=3.6em]
\item a variational equivalence between the channel edge and the relative
      physical/comparison Witten floor;
\item an exact landing metric for every channel vector;
\item a dichotomy between a normalized physical soft one-form and comparison
      energy escaping through every fixed low-energy window; and
\item a ground-state Fisher-information identity which supplies an averaged,
      though not essential-supremum, Perron score bound.
\end{enumerate}

% =============================================================================
\section{Exact Schur landing identities}
\label{sec:V17-landing}
% =============================================================================

Let \(\mathscr H\) be a physical one-form Hilbert space and
\(\mathscr C\) a defect-channel Hilbert space.  Let
\[
 \mathcal A=\mathcal A^*\succ0,
 \qquad
 \mathcal V:\mathscr C\to\mathscr H,
 \qquad
 L:=\mathcal A-\mathcal V\mathcal V^*\succ0,
\]
with the inverses understood as form inverses on the declared carrier.  Put
\begin{equation}
 \mathcal K
 :=\mathcal V^*\mathcal A^{-1}\mathcal V,
 \qquad
 \mathcal B
 :=\mathcal A^{-1}\mathcal V,
 \qquad
 \mathcal M
 :=\mathcal B^*\mathcal B
 =\mathcal V^*\mathcal A^{-2}\mathcal V.
 \label{eq:V17-K-B-M}
\end{equation}
The operator \(\mathcal M\) is the physical landing metric on the channel.

\begin{theorem}[Exact channel-to-one-form landing ledger]
\label{thm:V17-landing-ledger}
For \(z\in\mathscr C\), define
\[
 \omega_z:=\mathcal Bz=\mathcal A^{-1}\mathcal Vz.
\]
Then
\begin{align}
 \boxed{
 \norm{\omega_z}^2
 &=\ip{z}{\mathcal Mz},}
 \label{eq:V17-landing-norm}\\
 \boxed{
 \ip{\omega_z}{\mathcal A\omega_z}
 &=\ip{z}{\mathcal Kz},}
 \label{eq:V17-comparison-energy}\\
 \boxed{
 \ip{\omega_z}{L\omega_z}
 &=\ip{z}{\mathcal K(I-\mathcal K)z}.}
 \label{eq:V17-physical-energy}
\end{align}
\end{theorem}

\begin{proof}
The norm identity follows from
\[
 \mathcal B^*\mathcal B=\mathcal M.
\]
For the comparison energy,
\[
 \mathcal B^*\mathcal A\mathcal B
 =
 \mathcal V^*\mathcal A^{-1}\mathcal V
 =\mathcal K.
\]
Also
\[
 \mathcal V^*\omega_z
 =
 \mathcal V^*\mathcal A^{-1}\mathcal Vz
 =\mathcal Kz.
\]
Therefore
\[
 \ip{\omega_z}{L\omega_z}
 =
 \ip{z}{\mathcal Kz}
 -\norm{\mathcal Kz}^2
 =
 \ip{z}{\mathcal K(I-\mathcal K)z},
\]
because \(\mathcal K\) is self-adjoint.
\end{proof}

\begin{theorem}[Channel edge equals the relative Witten defect]
\label{thm:V17-edge-relative-floor}
On the closure of \(\Ran\mathcal B\) in the comparison form norm,
\begin{equation}
 \boxed{
 \norm{\mathcal K}
 =
 \sup_{\substack{\omega\in\overline{\Ran\mathcal B}\\
                  \ip{\omega}{\mathcal A\omega}>0}}
 \frac{\norm{\mathcal V^*\omega}^2}
      {\ip{\omega}{\mathcal A\omega}},}
 \label{eq:V17-edge-variational}
\end{equation}
and hence
\begin{equation}
 \boxed{
 1-\norm{\mathcal K}
 =
 \inf_{\substack{\omega\in\overline{\Ran\mathcal B}\\
                  \ip{\omega}{\mathcal A\omega}>0}}
 \frac{\ip{\omega}{L\omega}}
      {\ip{\omega}{\mathcal A\omega}}.}
 \label{eq:V17-relative-floor}
\end{equation}
Thus a uniform channel edge
\(\norm{\mathcal K}\le\rho_*<1\) is exactly the relative form inequality
\begin{equation}
 \boxed{
 L\succeq(1-\rho_*)\mathcal A
 \quad\text{on }\overline{\Ran\mathcal B}.}
 \label{eq:V17-relative-inequality}
\end{equation}
\end{theorem}

\begin{proof}
First take \(\omega=\mathcal Bz\).  By
\Cref{thm:V17-landing-ledger},
\[
 \frac{\norm{\mathcal V^*\omega}^2}
      {\ip{\omega}{\mathcal A\omega}}
 =
 \frac{\ip{z}{\mathcal K^2z}}
      {\ip{z}{\mathcal Kz}}.
\]
Functional calculus for the positive contraction \(\mathcal K\) shows that
the supremum of this quotient over
\((\Ker\mathcal K)^\perp\) is
\(\norm{\mathcal K}\).  Closure in the comparison form norm extends the
identity to \(\overline{\Ran\mathcal B}\).  Since
\[
 \ip{\omega}{L\omega}
 =
 \ip{\omega}{\mathcal A\omega}
 -\norm{\mathcal V^*\omega}^2,
\]
\eqref{eq:V17-relative-floor} follows.  The final equivalence is the
quadratic-form characterization of operator order on the declared range.
\end{proof}

\begin{assessment}[What the uniform edge really asserts]
\label{ass:V17-edge-meaning}
The edge is not merely a convergence parameter for a Neumann series.  It is
the smallest relative fraction of the positive comparison energy that
survives the complete curvature-defect subtraction on the source-generated
one-form range.  It is stronger than an ordinary \(L^2\) mass floor when
\(\mathcal A\) becomes large in ultraviolet directions, and it is precisely
why a near-unit channel vector requires a landing audit.
\end{assessment}

% =============================================================================
\section{Near-unit edge: physical soft mode or ultraviolet escape}
\label{sec:V17-edge-alternative}
% =============================================================================

\begin{theorem}[Landing alternative for a near-unit channel sequence]
\label{thm:V17-landing-alternative}
Let
\[
 0\preceq\mathcal K_j\prec I,
 \qquad
 \norm{z_j}=1,
 \qquad
 \ip{z_j}{\mathcal K_jz_j}\longrightarrow1,
 \]
and set
\[
 \omega_j
 :=
 \mathcal A_j^{-1}\mathcal V_jz_j.
\]
Then
\begin{align}
 \ip{\omega_j}{\mathcal A_j\omega_j}
 &\longrightarrow1,
 \label{eq:V17-unit-comparison}\\
 \ip{\omega_j}{L_j\omega_j}
 &\longrightarrow0.
 \label{eq:V17-vanishing-physical-form}
\end{align}
After passing to a subsequence, exactly one of the following alternatives
holds.
\begin{enumerate}[label=\textup{(\Alph*)},leftmargin=4em]
\item \emph{Physical landing.}  There is \(m_*>0\) such that
      \[
      \norm{\omega_j}^2
      =\ip{z_j}{\mathcal M_jz_j}\ge m_*.
      \]
      Then the normalized one-forms
      \(\widetilde\omega_j=\omega_j/\norm{\omega_j}\) satisfy
      \[
      \boxed{
      \ip{\widetilde\omega_j}{L_j\widetilde\omega_j}
      \longrightarrow0.}
      \tag{V17.A}
      \]
\item \emph{Landing collapse.}  One has
      \[
      \boxed{
      \ip{z_j}{\mathcal M_jz_j}
      =\norm{\omega_j}^2\longrightarrow0,}
      \tag{V17.B}
      \]
      while the comparison energy remains asymptotically one.  This is an
      ultraviolet comparison-energy escape and is not by itself a normalized
      physical mass obstruction.
\end{enumerate}
\end{theorem}

\begin{proof}
\Cref{thm:V17-landing-ledger} gives
\eqref{eq:V17-unit-comparison}.  Since
\[
 0\le\mathcal K_j(I-\mathcal K_j)\le I-\mathcal K_j,
\]
\[
 0\le
 \ip{\omega_j}{L_j\omega_j}
 \le
 1-\ip{z_j}{\mathcal K_jz_j}
 \longrightarrow0,
\]
proving \eqref{eq:V17-vanishing-physical-form}.  Any nonnegative sequence
\(\norm{\omega_j}^2\) has a subsequence whose liminf is positive or which
converges to zero.  In the first case normalize and divide
\eqref{eq:V17-vanishing-physical-form} by the norm lower bound.  The second
case is exactly the stated landing collapse.
\end{proof}

\begin{corollary}[Fixed comparison-energy window forces physical landing]
\label{cor:V17-low-window-landing}
Let \(E_j^{\mathcal A}([0,\Lambda])\) be the spectral projection of
\(\mathcal A_j\).  Under the hypotheses of
\Cref{thm:V17-landing-alternative}, suppose there exist fixed
\(0<\Lambda<\infty\) and \(\theta>0\) such that
\begin{equation}
 \left\|
 \mathcal A_j^{1/2}
 E_j^{\mathcal A}([0,\Lambda])\omega_j
 \right\|^2
 \ge\theta
 \label{eq:V17-low-energy-share}
\end{equation}
along a subsequence.  Then
\begin{equation}
 \boxed{
 \norm{\omega_j}^2\ge\theta/\Lambda}
 \label{eq:V17-window-frame}
\end{equation}
and the physical-landing alternative \textup{(A)} holds.

Conversely, if \textup{(B)} holds, then for every fixed \(\Lambda\),
\begin{equation}
 \boxed{
 \left\|
 \mathcal A_j^{1/2}
 E_j^{\mathcal A}([0,\Lambda])\omega_j
 \right\|^2
 \longrightarrow0.}
 \label{eq:V17-UV-escape}
\end{equation}
\end{corollary}

\begin{proof}
On \(\Ran E_j^{\mathcal A}([0,\Lambda])\),
\[
 \ip{\eta}{\mathcal A_j\eta}
 \le\Lambda\norm{\eta}^2.
\]
Thus \eqref{eq:V17-low-energy-share} implies
\[
 \norm{\omega_j}^2
 \ge
 \norm{E_j^{\mathcal A}([0,\Lambda])\omega_j}^2
 \ge\theta/\Lambda.
\]
For the converse,
\[
 \left\|
 \mathcal A_j^{1/2}
 E_j^{\mathcal A}([0,\Lambda])\omega_j
 \right\|^2
 \le
 \Lambda\norm{\omega_j}^2\longrightarrow0.
\]
\end{proof}

\begin{corollary}[Uniform landing frame turns the physical gap into an edge]
\label{cor:V17-gap-to-edge}
Suppose on the channel carrier
\begin{equation}
 \mathcal M\ \succeq\ \sigma\mathcal K,
 \qquad \sigma>0,
 \label{eq:V17-landing-frame}
\end{equation}
and the physical one-form operator obeys
\[
 L\succeq mI.
\]
Then
\begin{equation}
 \boxed{
 \norm{\mathcal K}\le1-m\sigma.}
 \label{eq:V17-gap-edge}
\end{equation}
\end{corollary}

\begin{proof}
For \(\omega_z=\mathcal Bz\),
\[
 \ip{z}{\mathcal K(I-\mathcal K)z}
 =
 \ip{\omega_z}{L\omega_z}
 \ge
 m\norm{\omega_z}^2
 =
 m\ip{z}{\mathcal Mz}
 \ge
 m\sigma\ip{z}{\mathcal Kz}.
\]
Functional calculus on the support of \(\mathcal K\) gives
\(
 I-\mathcal K\succeq m\sigma I
\), proving \eqref{eq:V17-gap-edge}.
\end{proof}

\begin{firewall}[A channel survivor is not automatically infrared]
\label{fire:V17-survivor-not-IR}
The V35 near-unit vector becomes a physical soft one-form only after
\eqref{eq:V17-landing-norm} or a fixed-window condition such as
\eqref{eq:V17-low-energy-share} prevents norm collapse.  If
\eqref{eq:V17-UV-escape} holds instead, the correct next tool is the
high-comparison-energy/Feshbach or cofinal lower-frame audit, not an assertion
that a local continuum mass has vanished.
\end{firewall}

% =============================================================================
\section{A ground-state Fisher-information sum rule}
\label{sec:V17-Fisher}
% =============================================================================

Return to the finite Hamiltonian regulator
\begin{equation}
 H^{\rm raw}
 =
 \kappa\sum_{i=1}^N(-\Delta_i)+V_B,
 \qquad
 H^{\rm raw}h=E_0h,
 \qquad
 \int h^2\,\dd\mathfrak m=1,
 \label{eq:V17-raw-Hamiltonian}
\end{equation}
and set
\[
 \dd\pi=h^2\dd\mathfrak m,
 \qquad
 \mathcal J_i=-\nabla_i\log h.
\]

\begin{theorem}[Exact Perron Fisher-energy identity]
\label{thm:V17-Fisher-identity}
The physical score satisfies
\begin{equation}
 \boxed{
 \sum_{i=1}^N
 \int\norm{\mathcal J_i}^2\,\dd\pi
 =
 \int\norm{\nabla h}^2\,\dd\mathfrak m
 =
 \frac{
 E_0-\displaystyle\int V_B\,\dd\pi
 }{\kappa}.}
 \label{eq:V17-Fisher-identity}
\end{equation}
For the V14 marginal score
\[
 \overline{\mathcal J}_i(a,U)
 =
 \mathbb E_{\nu_{a,U}}\mathcal J_i,
\]
one also has
\begin{equation}
 \boxed{
 \int
 \norm{\overline{\mathcal J}_i(a,U)}^2
 \,\dd\pi_{\le i}(a,U)
 \le
 \int\norm{\mathcal J_i}^2\,\dd\pi,}
 \label{eq:V17-marginal-Fisher}
\end{equation}
where \(\pi_{\le i}\) is the \(\mathcal F_i\)-marginal of \(\pi\).
\end{theorem}

\begin{proof}
Since \(\mathcal J_i=-\nabla_i h/h\),
\[
 \int\norm{\mathcal J_i}^2\,\dd\pi
 =
 \int\norm{\nabla_i h}^2\,\dd\mathfrak m.
\]
Sum over \(i\).  Taking the inner product of
\eqref{eq:V17-raw-Hamiltonian} with \(h\) and integrating by parts gives
\[
 E_0
 =
 \kappa\int\norm{\nabla h}^2\,\dd\mathfrak m
 +\int V_Bh^2\,\dd\mathfrak m.
\]
This proves \eqref{eq:V17-Fisher-identity}.  Conditional Jensen under
\(\nu_{a,U}\), followed by disintegration of \(\pi\), gives
\eqref{eq:V17-marginal-Fisher}.
\end{proof}

\begin{corollary}[Translation-invariant per-coordinate score bound]
\label{cor:V17-score-density}
Suppose the regulator and its ground state are transitive on the \(N\)
physical coordinate blocks, and
\begin{equation}
 \frac{
 E_0-\int V_B\,\dd\pi
 }{N}
 \le e_{\rm kin}
 \label{eq:V17-kinetic-density}
\end{equation}
uniformly.  Then every block satisfies
\begin{equation}
 \boxed{
 \int\norm{\mathcal J_i}^2\,\dd\pi
 \le\frac{e_{\rm kin}}{\kappa}.}
 \label{eq:V17-score-density}
\end{equation}
\end{corollary}

\begin{proof}
Transitivity makes all \(N\) summands in
\eqref{eq:V17-Fisher-identity} equal.  Divide by \(N\) and use
\eqref{eq:V17-kinetic-density}.
\end{proof}

\begin{assessment}[First populated score stock, but not the Harnack stock]
\label{ass:V17-average-not-sup}
The Fisher identity gives an exact \(L^2(\pi)\) score density and, under an
energy-density bound, a volume-uniform average per coordinate.  V14 requires
a conditional pathwise or essential-supremum score diameter.  In the
eight-dimensional group coordinate, an \(L^2\) gradient bound alone does not
imply an \(L^\infty\) oscillation bound.  The remaining upgrade needs local
elliptic regularity, higher score moments, or a good-set/bad-set capacity
argument.  Treating \eqref{eq:V17-score-density} as the V14 bound would be
invalid.
\end{assessment}

% =============================================================================
\section{V17 ledger and revised frontier}
\label{sec:V17-ledger}
% =============================================================================

\begin{theorem}[Integrated V17 statement]
\label{thm:V17-integrated}
Preserving V1--V16, V17 proves:
\begin{enumerate}[label=\textup{(V17.\arabic*)},leftmargin=5.5em]
\item the exact channel landing norm, comparison energy, and physical energy;
\item equivalence of the defect edge with a relative Witten floor;
\item a physical-soft/landing-collapse alternative for every near-unit
      channel sequence;
\item a fixed comparison-energy window criterion forcing physical landing;
\item a landing-frame implication from a physical mass floor to a uniform
      channel edge;
\item the exact Perron Fisher-energy identity;
\item a translation-invariant per-coordinate \(L^2\) score bound and the
      precise reason it does not yet give the V14 score diameter.
\end{enumerate}
\end{theorem}

\begin{proof}
Item \textup{(V17.1)} is \Cref{thm:V17-landing-ledger}.  Item
\textup{(V17.2)} is \Cref{thm:V17-edge-relative-floor}.  Items
\textup{(V17.3)--(V17.4)} are
\Cref{thm:V17-landing-alternative,cor:V17-low-window-landing}.  Item
\textup{(V17.5)} is \Cref{cor:V17-gap-to-edge}.  Items
\textup{(V17.6)--(V17.7)} are
\Cref{thm:V17-Fisher-identity,cor:V17-score-density,
ass:V17-average-not-sup}.
\end{proof}

\begingroup
\small
\begin{longtable}{>{\raggedright\arraybackslash}p{0.28\textwidth}
                  >{\raggedright\arraybackslash}p{0.16\textwidth}
                  >{\raggedright\arraybackslash}p{0.48\textwidth}}
\toprule
\textbf{V17 row}&\textbf{Status}&\textbf{Advance or remaining obstruction}\\
\midrule
\endfirsthead
\toprule
\textbf{V17 row}&\textbf{Status}&\textbf{Advance or remaining obstruction}\\
\midrule
\endhead
Channel landing identities
& \MGpass
& Exact physical norm and comparison/physical energies.\\[0.35em]
Edge-relative-floor equivalence
& \MGpass
& Gives the precise form meaning of \(1-\norm{\mathcal K}\).\\[0.35em]
Near-unit survivor classification
& \MGpass
& Separates a normalized physical soft form from ultraviolet landing collapse.\\[0.35em]
Landing lower frame
& \MGopen
& No attachment proves a fixed low-comparison-energy share for every near-unit sequence.\\[0.35em]
Perron Fisher score
& \MGpass
& Exact identity and conditional Jensen bound.\\[0.35em]
Uniform kinetic-energy density
& \MGopen
& Must be established on the response-matched weak-coupling trajectory.\\[0.35em]
Score \(L^2\)-to-Harnack upgrade
& \MGopen
& Requires local regularity, higher moments, or a capacity split.\\[0.35em]
Uniform defect edge
& \MGopen
& Must be proved or resolved by the landing alternative.\\[0.35em]
Full Yang--Mills mass gap
& \MGopen
& The weak-coupling uniform stocks and nontrivial OS limit remain.\\
\bottomrule
\end{longtable}
\endgroup

\begin{openproblem}[V18 mandate: low-window landing and score regularity]
\label{op:V18}
Proceed upstream along the two corrected branches.
\begin{enumerate}[label=\textup{(V18.\arabic*)},leftmargin=5.5em]
\item For every near-unit cell-saturated defect-channel sequence, test the
      fixed comparison-energy windows in
      \eqref{eq:V17-low-energy-share}.
\item If a window retains positive energy, return the normalized physical
      soft one-form of \Cref{thm:V17-landing-alternative}; classify its
      spatial and topological support.
\item If all fixed windows empty, use the V6/V10 high-Casimir shielding and
      V1 lower-frame test to decide whether the ultraviolet landing survives
      the local continuum algebra.
\item Prove a uniform kinetic-energy density
      \eqref{eq:V17-kinetic-density} from the response-matched Hamiltonian
      normalization.
\item Upgrade the average marginal score to the V14 conditional score
      diameter using a local elliptic or higher-moment estimate; route the bad
      set through the V12--V13 capacity stock rather than discarding it.
\item If the edge, signed-source moment, and score diameter close, invoke
      \eqref{eq:V16-explicit-gap}; only then perform the V1 OS passage.
\end{enumerate}
\end{openproblem}

\begin{firewall}[End-of-V17 claim boundary]
V17 makes every near-unit defect-channel sequence physically interpretable
and proves an exact averaged Perron score sum rule.  It does not prove a
uniform landing frame, kinetic-energy density, score Harnack bound, defect
edge, signed-source moment, protected-sector control, or nontrivial continuum
OS state.  The full Yang--Mills mass gap remains open.
\end{firewall}

\section*{Version V17 append-only record}
\addcontentsline{toc}{section}{Version V17 append-only record}
The complete V16 source was retained before this continuation.  Its SHA--256
digest was
\begin{center}\scriptsize\ttfamily
2e269134835d8bc6b16a9cafff99d31bcdbcbbf07bae972baba3774ef8b85e20.
\end{center}
On the next iteration, retain this full file, append before its unique active
document-closing command, increment the filename to
\texttt{\_V18.tex}, and remove only the superseded V17 file from this note
series.

% =============================================================================
% END APPENDED VERSION V17 MATERIAL
% =============================================================================

% =============================================================================
% BEGIN APPENDED VERSION V18 MATERIAL
% =============================================================================

\clearpage
\part{Version V18: bounded-defect landing and an \(L^2\)-Lipschitz Perron score upgrade}
\label{part:V18}

\section{Direction}
\label{sec:V18-direction}

V17 separated a near-unit defect-channel vector into two possibilities:
physical one-form landing or ultraviolet landing collapse.  It also supplied
an exact averaged Perron Fisher bound but showed that an \(L^2\) score alone
does not control the V14 conditional Harnack diameter.

V18 advances both branches.
\begin{enumerate}[label=\textup{(\roman*)},leftmargin=3.6em]
\item A uniform bound on the defect synthesis
      \(\mathcal Vz\) forces a fixed low-comparison-energy share and therefore
      turns every near-unit channel sequence into a normalized physical soft
      one-form.
\item A conditional \(L^2\) bound for the marginal Perron score, combined
      with a uniform Lipschitz bound for that score, gives an explicit
      conditional Harnack and Poincar\'e floor.  The Lipschitz tensor is
      exactly the diagonal two-level curvature card from V15.
\end{enumerate}
The remaining issue is quantitative: establish these conditional bounds on
the response-matched weak-coupling trajectory, or route the bad past set
through the V13 capacity mechanism.

% =============================================================================
\section{A defect-size lower frame for Woodbury landing}
\label{sec:V18-defect-landing}
% =============================================================================

Use the V17 notation
\[
 \mathcal K=\mathcal V^*\mathcal A^{-1}\mathcal V,
 \qquad
 \omega_z=\mathcal A^{-1}\mathcal Vz,
 \qquad
 \mathcal M=\mathcal V^*\mathcal A^{-2}\mathcal V.
\]

\begin{theorem}[Defect synthesis controls the landing metric]
\label{thm:V18-defect-landing}
For every \(z\) with \(\mathcal Vz\ne0\),
\begin{equation}
 \boxed{
 \ip{z}{\mathcal Mz}
 =\norm{\omega_z}^2
 \ge
 \frac{\ip{z}{\mathcal Kz}^2}
      {\norm{\mathcal Vz}^2}.}
 \label{eq:V18-landing-lower}
\end{equation}
In particular, if \(\norm z=1\) and
\begin{equation}
 \norm{\mathcal Vz}\le V_*
 \label{eq:V18-defect-size}
\end{equation}
on a declared channel carrier, then
\begin{equation}
 \boxed{
 \ip{z}{\mathcal Mz}
 \ge
 V_*^{-2}\ip{z}{\mathcal Kz}^2.}
 \label{eq:V18-landing-frame}
\end{equation}
\end{theorem}

\begin{proof}
Since
\[
 \ip{z}{\mathcal Kz}
 =
 \ip{\mathcal Vz}{\mathcal A^{-1}\mathcal Vz}
 =
 \ip{\mathcal Vz}{\omega_z},
\]
Cauchy--Schwarz gives
\[
 \ip{z}{\mathcal Kz}
 \le\norm{\mathcal Vz}\norm{\omega_z}.
\]
Square and use
\(\norm{\omega_z}^2=\ip{z}{\mathcal Mz}\).  The second assertion follows
from \eqref{eq:V18-defect-size}.
\end{proof}

\begin{theorem}[Bounded defect forces a fixed comparison-energy window]
\label{thm:V18-fixed-window}
Let \(z_j\) be unit channel vectors with
\[
 \ip{z_j}{\mathcal K_jz_j}\longrightarrow1,
 \qquad
 \norm{\mathcal V_jz_j}\le V_*<\infty.
\]
Put
\(\omega_j=\mathcal A_j^{-1}\mathcal V_jz_j\).
Then:
\begin{align}
 \boxed{
 \liminf_j\norm{\omega_j}^2
 \ge V_*^{-2},}
 \label{eq:V18-noncollapse}\\
 \boxed{
 \left\|
 \mathcal A_j^{1/2}
 E_j^{\mathcal A}((\Lambda,\infty))\omega_j
 \right\|^2
 \le\frac{V_*^2}{\Lambda}
 \qquad(\Lambda>0).}
 \label{eq:V18-high-window-tail}
\end{align}
For every \(\Lambda>V_*^2\), the fixed low-energy window satisfies
\begin{equation}
 \boxed{
 \liminf_j
 \left\|
 \mathcal A_j^{1/2}
 E_j^{\mathcal A}([0,\Lambda])\omega_j
 \right\|^2
 \ge1-\frac{V_*^2}{\Lambda}>0.}
 \label{eq:V18-low-window-share}
\end{equation}
Consequently the physical-landing alternative of V17 holds and the
normalized \(\omega_j\)'s have vanishing \(L_j\)-Rayleigh quotient.
\end{theorem}

\begin{proof}
\Cref{thm:V18-defect-landing} and
\(\ip{z_j}{\mathcal K_jz_j}\to1\) give
\eqref{eq:V18-noncollapse}.  Since
\[
 \mathcal A_j\omega_j=\mathcal V_jz_j,
\]
spectral calculus gives
\[
 \left\|
 \mathcal A_j
 \omega_j
 \right\|^2
 =
 \norm{\mathcal V_jz_j}^2
 \le V_*^2.
\]
On the spectral interval \((\Lambda,\infty)\),
\(\lambda\le\lambda^2/\Lambda\).  Therefore
\[
 \left\|
 \mathcal A_j^{1/2}
 E_j^{\mathcal A}((\Lambda,\infty))\omega_j
 \right\|^2
 \le
 \Lambda^{-1}\norm{\mathcal A_j\omega_j}^2
 \le V_*^2/\Lambda,
\]
proving \eqref{eq:V18-high-window-tail}.  The total comparison energy tends
to one by \eqref{eq:V17-unit-comparison}; subtract the high-window estimate
to get \eqref{eq:V18-low-window-share}.  Apply
\Cref{cor:V17-low-window-landing,thm:V17-landing-alternative}.
\end{proof}

\begin{corollary}[Edge-or-physical-soft alternative under bounded defects]
\label{cor:V18-edge-or-soft}
Assume the cell-saturated defect carriers obey
\begin{equation}
 \sup_j
 \norm{\mathcal V_j|_{\mathscr C_j^{\rm cell}}}
 \le V_*<\infty.
 \label{eq:V18-uniform-V}
\end{equation}
Then exactly one of the following holds after subsequence extraction:
\begin{enumerate}[label=\textup{(\Alph*)},leftmargin=4em]
\item there is \(\rho_*<1\) with
      \(\norm{\mathcal K_j}\le\rho_*\), so the V16 return decay applies; or
\item there are normalized physical one-forms
      \(\widetilde\omega_j\) with
      \[
      \ip{\widetilde\omega_j}{L_j\widetilde\omega_j}\longrightarrow0.
      \]
\end{enumerate}
No ultraviolet landing-collapse branch remains under
\eqref{eq:V18-uniform-V}.
\end{corollary}

\begin{proof}
If no common edge exists, choose unit \(z_j\) with
\(\ip{z_j}{\mathcal K_jz_j}\to1\) by the variational principle.  Apply
\Cref{thm:V18-fixed-window}.  Otherwise the first alternative holds.
\end{proof}

\begin{firewall}[Scaling of the defect synthesis]
\label{fire:V18-V-scaling}
The norm in \eqref{eq:V18-uniform-V} must be measured after the same physical
normalization as \(\mathcal A_j\) and \(L_j\).  A bare curvature square root
may grow with the lattice coupling even when its renormalized physical
landing is finite.  V18 proves the implication from a uniform physical
\(V_*\); it does not assert that the supplied V35 bare defect synthesis has
that bound on the continuum trajectory.
\end{firewall}

% =============================================================================
\section{Conditional \(L^2\)-Lipschitz control of the Perron score}
\label{sec:V18-score-upgrade}
% =============================================================================

Retain the V14 sequential law
\[
 \dd q_{i,a}(U)
 \propto\widehat h_i(a,U)^2\,\dd\mathfrak m_G(U)
\]
and its marginal score
\[
 \overline{\mathcal J}_i(a,U)
 =
 -\nabla_U\log\widehat h_i(a,U).
\]
Let \(D_G\) be the Riemannian diameter of \(G=\mathrm{SU}(3)\).  Define the
conditional kinetic score
\begin{equation}
 m_i(a)^2
 :=
 \int_G
 \mathbb E_{\nu_{a,U}}
 \norm{\mathcal J_i(a,U,\cdot)}^2
 \,\dd q_{i,a}(U)
 \label{eq:V18-conditional-score-L2}
\end{equation}
and the diagonal marginal-curvature stock
\begin{equation}
 \ell_i(a)
 :=
 \operatorname*{ess\,sup}_{U\in G}
 \norm{
 \nabla_U\overline{\mathcal J}_i(a,U)
 }_{\op}.
 \label{eq:V18-marginal-Lipschitz}
\end{equation}

\begin{lemma}[Conditional \(L^2\) norm of the marginal score]
\label{lem:V18-marginal-L2}
For every past \(a\),
\begin{equation}
 \boxed{
 \int_G
 \norm{\overline{\mathcal J}_i(a,U)}^2
 \,\dd q_{i,a}(U)
 \le m_i(a)^2.}
 \label{eq:V18-marginal-L2}
\end{equation}
\end{lemma}

\begin{proof}
Conditional Jensen gives
\[
 \norm{
 \mathbb E_{\nu_{a,U}}\mathcal J_i
 }^2
 \le
 \mathbb E_{\nu_{a,U}}\norm{\mathcal J_i}^2.
\]
Integrate in \(U\) under \(q_{i,a}\).
\end{proof}

\begin{theorem}[\(L^2\)-Lipschitz Harnack upgrade on one group coordinate]
\label{thm:V18-L2-Lipschitz}
For every \(a\) with \(m_i(a),\ell_i(a)<\infty\),
\begin{align}
 \boxed{
 \sup_{U\in G}
 \norm{\overline{\mathcal J}_i(a,U)}
 \le
 m_i(a)+D_G\ell_i(a),}
 \label{eq:V18-score-sup}\\
 \boxed{
 \operatorname{osc}_{U\in G}
 \log\widehat h_i(a,U)
 \le
 D_G\bigl[m_i(a)+D_G\ell_i(a)\bigr].}
 \label{eq:V18-amplitude-oscillation}
\end{align}
Consequently the sequential conditional has Poincar\'e floor
\begin{equation}
 \boxed{
 \rho_{i,a}
 \ge
 C_F
 \exp\left(
 -2D_G[m_i(a)+D_G\ell_i(a)]
 \right).}
 \label{eq:V18-conditional-gap}
\end{equation}
\end{theorem}

\begin{proof}
The law \(q_{i,a}\) has a smooth strictly positive density.  By
\Cref{lem:V18-marginal-L2}, there is \(U_0\in G\) with
\[
 \norm{\overline{\mathcal J}_i(a,U_0)}
 \le m_i(a).
\]
Join \(U_0\) to \(U\) by a minimizing geodesic of length at most \(D_G\).
Parallel transport and the covariant fundamental theorem of calculus give
\[
 \norm{\overline{\mathcal J}_i(a,U)}
 \le
 \norm{\overline{\mathcal J}_i(a,U_0)}
 +D_G\ell_i(a),
\]
which proves \eqref{eq:V18-score-sup}.  Integrate
\(
\nabla_U\log\widehat h_i=-\overline{\mathcal J}_i
\)
along a minimizing geodesic between arbitrary endpoints to obtain
\eqref{eq:V18-amplitude-oscillation}.  The logarithmic density of \(q_{i,a}\)
is \(2\log\widehat h_i\) plus a constant.  Apply the bounded-density
Poincar\'e transport
\Cref{lem:V14-density-Poincare} to Haar measure, whose gap is \(C_F\).
\end{proof}

\begin{theorem}[The Lipschitz stock is the signed V15 diagonal card]
\label{thm:V18-diagonal-curvature}
The derivative in \eqref{eq:V18-marginal-Lipschitz} is exactly
\begin{equation}
 \boxed{
 \nabla_U\overline{\mathcal J}_i
 =
 \mathbb E_{\nu_{a,U}}
 \left[
   \mathbb E_{p_x}(K_x)_{ii}
   -\Cov_{p_x}\bigl((J_x)_i,(J_x)_i\bigr)
 \right]
 -
 2\Cov_{\nu_{a,U}}
 \bigl(\mathcal J_i,\mathcal J_i\bigr).}
 \label{eq:V18-diagonal-curvature}
\end{equation}
Thus \(\ell_i(a)\) is not a new fitted regularity constant: it is the
operator norm of the completely assembled diagonal two-level Perron
curvature card.
\end{theorem}

\begin{proof}
Take \(b=i\) in the V15 two-level marginal curvature identity
\eqref{eq:V15-two-level-curvature}, then insert
\eqref{eq:V15-Perron-curvature}.
\end{proof}

\begin{corollary}[Uniform conditional score packet]
\label{cor:V18-uniform-score-packet}
If
\begin{equation}
 \operatorname*{ess\,sup}_{j,i,a}m_{j,i}(a)\le M_*,
 \qquad
 \operatorname*{ess\,sup}_{j,i,a}\ell_{j,i}(a)\le L_*,
 \label{eq:V18-uniform-ML}
\end{equation}
then all sequential conditionals have the common Poincar\'e floor
\begin{equation}
 \boxed{
 \rho_{\rm site}
 \ge
 C_F
 \exp\left(
 -2D_G[M_*+D_GL_*]
 \right).}
 \label{eq:V18-uniform-site-gap}
\end{equation}
If the V16 response summability constant is also uniformly bounded by
\(S_{\mathcal R}\), then
\begin{equation}
 \boxed{
 \gap(H_j)
 \ge
 \frac{
 \kappa_jC_F
 \exp\left(-2D_G[M_*+D_GL_*]\right)
 }{(1+S_{\mathcal R})^2}.}
 \label{eq:V18-uniform-H-gap}
\end{equation}
\end{corollary}

\begin{proof}
Apply \Cref{thm:V18-L2-Lipschitz} uniformly and then use the V14
response-matrix theorem with the V16 summability bound.
\end{proof}

% =============================================================================
\section{What the Fisher identity supplies and what capacity must supply}
\label{sec:V18-good-bad}
% =============================================================================

\begin{lemma}[Average conditional kinetic score]
\label{lem:V18-average-conditional-score}
Let \(\pi_{<i}\) be the past marginal on
\((x_1,\ldots,x_{i-1})\).  Then
\begin{equation}
 \boxed{
 \int m_i(a)^2\,\dd\pi_{<i}(a)
 =
 \int\norm{\mathcal J_i(x)}^2\,\dd\pi(x).}
 \label{eq:V18-average-mi}
\end{equation}
Consequently, under the translation-invariant kinetic-density hypothesis
\eqref{eq:V17-kinetic-density},
\begin{equation}
 \boxed{
 \int m_i(a)^2\,\dd\pi_{<i}(a)
 \le\frac{e_{\rm kin}}{\kappa}.}
 \label{eq:V18-average-mi-bound}
\end{equation}
\end{lemma}

\begin{proof}
The pair \(q_{i,a}(U)\nu_{a,U}(y)\) is the conditional law of
\((x_i,\ldots,x_N)\) given the past \(a\).  Integrate
\eqref{eq:V18-conditional-score-L2} first in \(U,y\), then in \(a\).
This gives \eqref{eq:V18-average-mi}.  Apply
\eqref{eq:V17-score-density}.
\end{proof}

\begin{corollary}[Quantified bad-past set]
\label{cor:V18-bad-past}
For \(T>0\), define
\[
 \mathcal B_{i,T}
 :=
 \{a:m_i(a)>T\}.
\]
Under \eqref{eq:V18-average-mi-bound},
\begin{equation}
 \boxed{
 \pi_{<i}(\mathcal B_{i,T})
 \le
 \frac{e_{\rm kin}}{\kappa T^2}.}
 \label{eq:V18-bad-past-probability}
\end{equation}
\end{corollary}

\begin{proof}
Apply Markov's inequality to \(m_i(a)^2\).
\end{proof}

\begin{firewall}[The bad-past set cannot be discarded]
\label{fire:V18-bad-past-capacity}
The probability estimate \eqref{eq:V18-bad-past-probability} does not replace
the essential supremum in \eqref{eq:V18-uniform-ML}.  V12 proved by an exact
counterexample that a rare set can carry a soft mode.  To use only the
averaged Fisher stock, the bad-past set must be treated as a separate well
and connected to its complement by the physical Perron capacity certificate
of V13.
\end{firewall}

\begin{assessment}[Corrected score-regularity frontier]
\label{ass:V18-score-frontier}
The score-diameter problem no longer requires an unspecified
\(L^\infty\)-bound on the raw one-step score.  It is enough to control:
\[
 \operatorname*{ess\,sup}_a m_i(a)
 \quad\text{and}\quad
 \operatorname*{ess\,sup}_a
 \left\|
 \mathbb E_\nu[\mathbb E_pK-\Cov_p(J)]
 -2\Cov_\nu(\mathcal J)
 \right\|.
\]
Alternatively, one may keep the Fisher-average bound and prove capacity
across the explicitly quantified bad-past sets.
\end{assessment}

% =============================================================================
\section{V18 ledger and revised frontier}
\label{sec:V18-ledger}
% =============================================================================

\begin{theorem}[Integrated V18 statement]
\label{thm:V18-integrated}
Preserving V1--V17, V18 proves:
\begin{enumerate}[label=\textup{(V18.\arabic*)},leftmargin=5.5em]
\item a defect-synthesis lower bound for the Woodbury landing metric;
\item a fixed low-comparison-energy window under a uniform physical defect
      bound;
\item an edge-or-normalized-physical-soft alternative with no ultraviolet
      landing collapse under that bound;
\item an \(L^2\)-plus-Lipschitz upgrade from conditional Perron score to a
      uniform conditional Poincar\'e floor;
\item identification of the Lipschitz tensor with the complete signed V15
      diagonal curvature card;
\item an explicit uniform Hamiltonian gap bound from the new score packet and
      V16 response summability;
\item an exact conditional averaging of the V17 Fisher score and a quantified
      bad-past set which must be routed through capacity.
\end{enumerate}
\end{theorem}

\begin{proof}
Items \textup{(V18.1)--(V18.2)} are
\Cref{thm:V18-defect-landing,thm:V18-fixed-window}.  Item
\textup{(V18.3)} is \Cref{cor:V18-edge-or-soft}.  Item
\textup{(V18.4)} is
\Cref{lem:V18-marginal-L2,thm:V18-L2-Lipschitz}.  Item
\textup{(V18.5)} is \Cref{thm:V18-diagonal-curvature}.  Item
\textup{(V18.6)} is \Cref{cor:V18-uniform-score-packet}.  Item
\textup{(V18.7)} is
\Cref{lem:V18-average-conditional-score,cor:V18-bad-past,
fire:V18-bad-past-capacity}.
\end{proof}

\begingroup
\small
\begin{longtable}{>{\raggedright\arraybackslash}p{0.28\textwidth}
                  >{\raggedright\arraybackslash}p{0.16\textwidth}
                  >{\raggedright\arraybackslash}p{0.48\textwidth}}
\toprule
\textbf{V18 row}&\textbf{Status}&\textbf{Advance or remaining obstruction}\\
\midrule
\endfirsthead
\toprule
\textbf{V18 row}&\textbf{Status}&\textbf{Advance or remaining obstruction}\\
\midrule
\endhead
Defect-size landing frame
& \MGpass
& A uniform \(V_*\) eliminates ultraviolet landing collapse.\\[0.35em]
Fixed comparison-energy window
& \MGpass
& Explicit tail \(V_*^2/\Lambda\).\\[0.35em]
\(L^2\)-Lipschitz Harnack upgrade
& \MGpass
& Gives the site floor \eqref{eq:V18-conditional-gap}.\\[0.35em]
Diagonal curvature incidence
& \MGpass
& Exact signed two-level V15 card.\\[0.35em]
Uniform physical defect bound
& \MGopen
& Must be proved after response matching and physical normalization.\\[0.35em]
Conditional kinetic-score supremum
& \MGopen
& Fisher controls its average, not its essential supremum.\\[0.35em]
Diagonal curvature bound
& \MGopen
& Requires a weak-coupling signed estimate, not separate absolute debits.\\[0.35em]
Bad-past capacity
& \MGopen
& Required if the essential-supremum route is replaced by the averaged route.\\[0.35em]
Full Yang--Mills mass gap
& \MGopen
& Uniform response-matched stocks and the nontrivial OS limit remain.\\
\bottomrule
\end{longtable}
\endgroup

\begin{openproblem}[V19 mandate: physical normalization and good/bad capacity]
\label{op:V19}
Proceed on the response-matched weak-coupling trajectory.
\begin{enumerate}[label=\textup{(V19.\arabic*)},leftmargin=5.5em]
\item Express the V35 defect synthesis in the V11 differential-Hamiltonian
      normalization and prove \eqref{eq:V18-uniform-V}, or return the
      diverging physical defect component.
\item Prove the kinetic-energy-density bound
      \eqref{eq:V17-kinetic-density}.
\item Estimate the complete signed diagonal tensor
      \eqref{eq:V18-diagonal-curvature}; do not bound its positive Hessian and
      negative covariance independently if their leading terms cancel.
\item Either prove the essential-supremum packet
      \eqref{eq:V18-uniform-ML}, or fix a threshold \(T\) and construct the
      V13 Perron flow/capacity bridge across every bad-past set
      \(\mathcal B_{i,T}\).
\item Combine the resulting site floor with the V16 source-edge and
      response-summability estimate.
\item Classify every failure by the V18 edge-or-soft and good/bad-capacity
      alternatives before performing the V1 OS passage.
\end{enumerate}
\end{openproblem}

\begin{firewall}[End-of-V18 claim boundary]
V18 eliminates ultraviolet ambiguity from near-unit channels under one
explicit physical defect bound and gives a rigorous score-regularity upgrade
from two conditional stocks.  It does not prove those stocks uniformly at
weak coupling, the bad-set capacity, protected-sector control, or a
nontrivial continuum OS state.  The full Yang--Mills mass gap remains open.
\end{firewall}

\section*{Version V18 append-only record}
\addcontentsline{toc}{section}{Version V18 append-only record}
The complete V17 source was retained before this continuation.  Its SHA--256
digest was
\begin{center}\scriptsize\ttfamily
a059b1c2d31745ff77fa8d04e3997e0bdcfac4bb1b7884c884f8965dd4d38671.
\end{center}
On the next iteration, retain this full file, append before its unique active
document-closing command, increment the filename to
\texttt{\_V19.tex}, and remove only the superseded V18 file from this note
series.

% =============================================================================
% END APPENDED VERSION V18 MATERIAL
% =============================================================================

% =============================================================================
% BEGIN APPENDED VERSION V19 MATERIAL
% =============================================================================

\clearpage
\part{Version V19: physical normalization of the defect channel and a predictable good/bad capacity compiler}
\label{part:V19}

\section{Direction}
\label{sec:V19-direction}

V18 left two alternatives at the same-law weak-coupling frontier.  On the
Woodbury side, a uniform bound for the physical defect synthesis prevents a
near-unit channel from disappearing into an ultraviolet landing.  On the
sequential side, an averaged Fisher bound identifies rare bad pasts but does
not control their capacity.

This continuation goes upstream at both points.  First, it audits exactly
how a common differential-Hamiltonian normalization changes the comparison
operator, defect factor, return channel, and landing metric.  The return
operator is scale invariant, whereas the landing metric is not.  The V35
pure-Wilson local estimates then give an explicit, volume-independent
absolute defect bound and a dimensionless relative-defect test.  Second, the
reveal martingale is split by predictable good and bad past events.  This
gives an exact hybrid Poincar\'e theorem: the good increments use the V18
conditional site floor, while the bad increments are measured by one
separate capacity-routing norm.  No bad event is discarded on account of
small probability.

The full Yang--Mills mass gap is not claimed here.  The new conclusions are
exact compilers and scaling identities.  They identify which quantities
must be bounded after the V11 transfer-to-Hamiltonian normalization and
which estimates cannot be obtained from the V35 unsigned row bound alone.

% =============================================================================
\section{Common Hamiltonian scaling of a Woodbury defect channel}
\label{sec:V19-scaling}
% =============================================================================

Let \(\mathscr H\) be the one-form Hilbert space and \(\mathscr D\) a defect
channel space.  Let
\[
 \mathcal A\succ0,
 \qquad
 \Theta=\mathcal V\mathcal V^*\succeq0,
 \qquad
 L=\mathcal A-\Theta\succ0,
\]
where \(\mathcal V:\mathscr D\to\mathscr H\).  As in V16--V18, set
\begin{equation}
 \mathcal K:=\mathcal V^*\mathcal A^{-1}\mathcal V,
 \qquad
 \mathcal B:=\mathcal A^{-1}\mathcal V,
 \qquad
 \mathcal M:=\mathcal B^*\mathcal B.
 \label{eq:V19-KBM}
\end{equation}

\begin{definition}[Common physical normalization]
\label{def:V19-common-scaling}
For \(s>0\), define
\begin{equation}
 \mathcal A^{[s]}:=s\mathcal A,
 \qquad
 L^{[s]}:=sL,
 \qquad
 \Theta^{[s]}:=s\Theta,
 \qquad
 \mathcal V^{[s]}:=s^{1/2}\mathcal V.
 \label{eq:V19-common-scaling}
\end{equation}
Thus \(L^{[s]}=\mathcal A^{[s]}-\Theta^{[s]}\) and
\(\Theta^{[s]}=\mathcal V^{[s]}(\mathcal V^{[s]})^*\).
\end{definition}

\begin{theorem}[Scale covariance of edge and landing]
\label{thm:V19-scale-covariance}
Let \(\mathcal K^{[s]},\mathcal B^{[s]},\mathcal M^{[s]}\) be constructed
from the scaled triple in \eqref{eq:V19-common-scaling}.  Then
\begin{equation}
 \boxed{
 \mathcal K^{[s]}=\mathcal K,
 \qquad
 \mathcal B^{[s]}=s^{-1/2}\mathcal B,
 \qquad
 \mathcal M^{[s]}=s^{-1}\mathcal M.}
 \label{eq:V19-scaling-KBM}
\end{equation}
Moreover,
\begin{equation}
 \boxed{
 \norm{\mathcal V^{[s]}}^2
 =\norm{\Theta^{[s]}}
 =s\norm{\Theta}.}
 \label{eq:V19-scaling-defect}
\end{equation}
For \(z\in\mathscr D\), put
\(\omega_z=\mathcal Bz\) and
\(\omega_z^{[s]}=\mathcal B^{[s]}z\).  Then
\begin{align}
 \norm{\omega_z^{[s]}}^2
 &=s^{-1}\norm{\omega_z}^2,
 \label{eq:V19-scaled-landing-norm}\\
 \ip{\omega_z^{[s]}}{\mathcal A^{[s]}\omega_z^{[s]}}
 &=\ip{z}{\mathcal Kz},
 \label{eq:V19-scaled-comparison-energy}\\
 \ip{\omega_z^{[s]}}{L^{[s]}\omega_z^{[s]}}
 &=\ip{z}{\mathcal K(I-\mathcal K)z}.
 \label{eq:V19-scaled-physical-energy}
\end{align}
In particular, the channel edge is invariant under a common scalar
Hamiltonian normalization, but the norm of its physical one-form landing is
not.
\end{theorem}

\begin{proof}
Direct substitution gives
\[
 \mathcal K^{[s]}
 =s^{1/2}\mathcal V^*(s\mathcal A)^{-1}s^{1/2}\mathcal V
 =\mathcal K
\]
and
\(
 \mathcal B^{[s]}=(s\mathcal A)^{-1}s^{1/2}\mathcal V
 =s^{-1/2}\mathcal B
\).
The identity for \(\mathcal M^{[s]}\) follows.  Since
\(\norm{TT^*}=\norm{T}^2\) for a bounded operator \(T\),
\[
 \norm{\mathcal V^{[s]}}^2
 =\norm{\mathcal V^{[s]}(\mathcal V^{[s]})^*}
 =\norm{s\Theta}.
\]
This proves \eqref{eq:V19-scaling-KBM}--\eqref{eq:V19-scaling-defect} and
the landing-norm identity.  Next,
\[
 \ip{\mathcal A^{-1}\mathcal Vz}
     {\mathcal A\mathcal A^{-1}\mathcal Vz}
 =\ip{z}{\mathcal V^*\mathcal A^{-1}\mathcal Vz}
 =\ip{z}{\mathcal Kz}.
\]
The scalar factors cancel in the scaled version.  Finally,
\begin{align*}
 \ip{\omega_z}{L\omega_z}
 &=\ip{\omega_z}{\mathcal A\omega_z}
   -\ip{\omega_z}{\mathcal V\mathcal V^*\omega_z}\\
 &=\ip{z}{\mathcal Kz}-\norm{\mathcal V^*\omega_z}^2
 =\ip{z}{\mathcal Kz}-\norm{\mathcal Kz}^2
 =\ip{z}{\mathcal K(I-\mathcal K)z}.
\end{align*}
Again the scaled scalar factors cancel.
\end{proof}

\begin{corollary}[Rayleigh scaling and the V18 landing stock]
\label{cor:V19-Rayleigh-scaling}
Whenever \(\omega_z\ne0\),
\begin{equation}
 \boxed{
 \frac{\ip{\omega_z^{[s]}}{L^{[s]}\omega_z^{[s]}}}
      {\norm{\omega_z^{[s]}}^2}
 =s\,
 \frac{\ip{\omega_z}{L\omega_z}}
      {\norm{\omega_z}^2}.}
 \label{eq:V19-Rayleigh-scaling}
\end{equation}
Let \(\mathscr C\subseteq\mathscr D\) be a declared channel carrier and set
\begin{equation}
 d_{\mathscr C}
 :=\sup_{\substack{z\in\mathscr C\\\norm{z}=1}}
   \norm{\mathcal Vz}^2.
 \label{eq:V19-carrier-defect-moment}
\end{equation}
For a cofinal family indexed by \(j\), the V18 physical defect-synthesis
hypothesis is exactly
\begin{equation}
 \boxed{
 \sup_j s_jd_{\mathscr C_j}<\infty.}
 \label{eq:V19-physical-carrier-stock}
\end{equation}
A stronger sufficient condition is
\(\sup_js_j\norm{\Theta_j}<\infty\).
\end{corollary}

\begin{proof}
Divide \eqref{eq:V19-scaled-physical-energy} by
\eqref{eq:V19-scaled-landing-norm}; equivalently use
\(L^{[s]}=sL\) and
\(\omega_z^{[s]}=s^{-1/2}\omega_z\).
Also
\[
 \norm{\mathcal V_j^{[s_j]}z}^2
 =s_j\norm{\mathcal V_jz}^2.
\]
Taking the supremum over the declared unit carrier proves
\eqref{eq:V19-physical-carrier-stock}.  The global sufficient condition
follows from
\(d_{\mathscr C_j}\le\norm{\mathcal V_j}^2=\norm{\Theta_j}\).
\end{proof}

\begin{firewall}[An invariant edge is not an invariant physical soft mode]
\label{fire:V19-edge-landing-scaling}
The equality \(\mathcal K^{[s]}=\mathcal K\) does not permit the common
normalization factor to be omitted.  If \(s_j\to\infty\), the same landed
one-form has norm smaller by \(s_j^{-1/2}\), and its normalized physical
Rayleigh quotient is larger by \(s_j\).  Thus a near-unit eigenvalue of the
bare \(\mathcal K_j\) is a physical soft mode only at the rate required by
\eqref{eq:V19-Rayleigh-scaling}, or after the V18 lower landing frame has
been proved in the physical normalization.
\end{firewall}

% =============================================================================
\section{The dimensionless relative defect and the V35 Wilson stock}
\label{sec:V19-relative-defect}
% =============================================================================

\begin{lemma}[Relative-defect representation]
\label{lem:V19-relative-defect}
Let
\begin{equation}
 \mathcal R_{\rm rel}
 :=\mathcal A^{-1/2}\Theta\mathcal A^{-1/2}.
 \label{eq:V19-relative-defect}
\end{equation}
Then
\begin{equation}
 \boxed{
 \norm{\mathcal K}
 =\norm{\mathcal R_{\rm rel}}
 \le\frac{\norm{\Theta}}{\rho}}
 \qquad
 \text{whenever }\mathcal A\succeq\rho I.
 \label{eq:V19-relative-defect-bound}
\end{equation}
Both \(\mathcal K\) and \(\mathcal R_{\rm rel}\) are unchanged, up to the
canonical channel/one-form realization, by the common scaling of
\Cref{def:V19-common-scaling}.
\end{lemma}

\begin{proof}
Put \(T=\mathcal A^{-1/2}\mathcal V\).  Then
\(\mathcal K=T^*T\) and
\(\mathcal R_{\rm rel}=TT^*\), so both norms equal \(\norm{T}^2\).
Furthermore,
\[
 \norm{\mathcal A^{-1/2}\Theta\mathcal A^{-1/2}}
 \le\norm{\mathcal A^{-1}}\norm{\Theta}
 \le\rho^{-1}\norm{\Theta}.
\]
Common scaling cancels between the two inverse square roots and
\(\Theta^{[s]}=s\Theta\).
\end{proof}

We now apply only the explicit pure-Wilson estimates stated in the supplied
V35 programme.  In its bare reduced-fibre normalization,
\begin{equation}
 \mathcal A_\eta
 =\nabla_\nu^*\nabla+\mathcal R_\eta,
 \qquad
 L_1=\mathcal A_\eta-\Theta_\eta,
 \label{eq:V19-V35-split}
\end{equation}
and on the reduced vertical carrier
\begin{equation}
 \mathcal A_\eta\succeq
 \rho_{\beta,\eta}I,
 \qquad
 \rho_{\beta,\eta}
 =\frac32+(1-\eta)\alpha_\beta\lambda_*,
 \qquad
 \alpha_\beta=\frac\beta3.
 \label{eq:V19-V35-floor}
\end{equation}
For one Wilson plaquette V35 proves
\(\norm{\Theta_{p,\eta}}\le8\alpha_\beta\), and each fine link belongs to
six plaquettes.

\begin{proposition}[Pure-Wilson defect row bound]
\label{prop:V19-Wilson-defect-row}
Let \(\Theta_\eta^{\rm W}=\sum_p\Theta_{p,\eta}\) denote the pure-Wilson
part of the V35 defect.  In the same reduced-fibre block norm,
\begin{equation}
 \boxed{
 \norm{\Theta_\eta^{\rm W}}
 \le192\alpha_\beta.}
 \label{eq:V19-Wilson-defect-bound}
\end{equation}
The constant is independent of the torus side.  Consequently its defect
factor may be chosen so that
\begin{equation}
 \boxed{
 \norm{\mathcal V_\eta^{\rm W}}^2
 \le192\alpha_\beta.}
 \label{eq:V19-Wilson-synthesis-bound}
\end{equation}
After a common physical normalization \(s_{\beta,\Lambda}\), the right-hand
side becomes \(192s_{\beta,\Lambda}\alpha_\beta\).
\end{proposition}

\begin{proof}
Each positive plaquette defect acts on four link slots.  V35's bound on the
plaquette operator implies that every block in that four-by-four block
matrix has norm at most \(8\alpha_\beta\).  Hence its contribution to one
absolute block row is at most
\(4\cdot8\alpha_\beta=32\alpha_\beta\).  Six plaquettes meet a fine link,
so every absolute block row is at most
\(6\cdot32\alpha_\beta=192\alpha_\beta\).  The operator is self-adjoint;
the block Schur test gives \eqref{eq:V19-Wilson-defect-bound}.  The identity
\(\norm{\mathcal V}^2=\norm{\mathcal V\mathcal V^*}\) gives
\eqref{eq:V19-Wilson-synthesis-bound}.  The last assertion is
\eqref{eq:V19-scaling-defect}.
\end{proof}

\begin{corollary}[What the unsigned Wilson rows can certify]
\label{cor:V19-Wilson-relative-test}
For the pure-Wilson defect channel,
\begin{equation}
 \boxed{
 \norm{\mathcal K_{\beta,\eta}^{\rm W}}
 \le
 \frac{192\alpha_\beta}
 {\frac32+(1-\eta)\lambda_*\alpha_\beta}.}
 \label{eq:V19-Wilson-relative-test}
\end{equation}
Thus the unsigned row method proves a uniform edge
\(\norm{\mathcal K_{\beta,\eta}^{\rm W}}\le q<1\) only if the right-hand
side of \eqref{eq:V19-Wilson-relative-test} is uniformly at most \(q\) on
the cofinal branch.  In particular, its large-\(\alpha_\beta\) limiting
upper bound is
\begin{equation}
 \frac{192}{(1-\eta)\lambda_*}.
 \label{eq:V19-Wilson-asymptotic-row-ratio}
\end{equation}
The common physical scalar \(s_{\beta,\Lambda}\) cannot improve this ratio.
\end{corollary}

\begin{proof}
Combine \Cref{lem:V19-relative-defect} with
\eqref{eq:V19-V35-floor} and
\eqref{eq:V19-Wilson-defect-bound}.  Divide numerator and denominator by
\(\alpha_\beta\) for the final assertion.  Scale invariance follows from
\Cref{thm:V19-scale-covariance}.
\end{proof}

\begin{assessment}[Normalization conclusion]
\label{ass:V19-normalization-conclusion}
The V35 unsigned local estimates now populate two exact tests.
\begin{enumerate}[label=\textup{(\alph*)},leftmargin=3.5em]
\item The global sufficient condition for the V18 physical landing stock is
\[
 \sup_j s_j\alpha_{\beta_j}<\infty.
\]
The sharp source-carrier condition is instead
\eqref{eq:V19-physical-carrier-stock}; it may be better than the global
row estimate.
\item The dimensionless edge test is
\eqref{eq:V19-Wilson-relative-test} and is independent of \(s_j\).
If its right-hand side is not below one, that is a failure of this unsigned
majorant, not evidence for a Yang--Mills soft mode.  A sharper proof must use
the actual source carrier, signed curvature cancellation, or local
good/bad capacity.
\end{enumerate}
The value of \(s_j\) must be read from the V11 same-generator
transfer-to-Hamiltonian identity; assigning it by convention would not be a
proof of \eqref{eq:V19-physical-carrier-stock}.
\end{assessment}



% =============================================================================
\section{Predictable good and bad martingale increments}
\label{sec:V19-predictable-split}
% =============================================================================

Retain the V14 reveal filtration
\(\mathcal F_i=\sigma(x_1,\ldots,x_i)\), martingale
\(M_if=\mathbb E_\pi[f\mid\mathcal F_i]\), and differences
\(D_if=M_if-M_{i-1}f\).

\begin{definition}[Predictable split and bad-increment capacity]
\label{def:V19-bad-capacity}
For each \(i\), let \(G_i\in\mathcal F_{i-1}\) be a good-past event and
write \(B_i=G_i^c\).  Define
\begin{equation}
 \mathfrak C_{\rm bad}(G_\bullet)
 :=
 \sup_{\substack{f\text{ smooth, gauge invariant}\\
                  \mathcal E_\pi(f)>0}}
 \frac{\displaystyle
       \sum_{i=1}^N\norm{\mathbf1_{B_i}D_if}_{L^2(\pi)}^2}
      {\mathcal E_\pi(f)},
 \qquad
 \mathcal E_\pi(f):=
 \sum_{b=1}^N\norm{\nabla_bf}_{L^2(\pi)}^2.
 \label{eq:V19-bad-capacity}
\end{equation}
This is the sharp energy cost of the bad predictable increments.  It is not
replaced by \(\sum_i\pi(B_i)\).
\end{definition}

\begin{lemma}[Exact predictable variance ledger]
\label{lem:V19-predictable-ledger}
For every \(f\in L^2(\pi)\),
\begin{equation}
 \boxed{
 \Var_\pi(f)
 =\sum_{i=1}^N\norm{\mathbf1_{G_i}D_if}_{L^2(\pi)}^2
  +\sum_{i=1}^N\norm{\mathbf1_{B_i}D_if}_{L^2(\pi)}^2.}
 \label{eq:V19-predictable-variance}
\end{equation}
\end{lemma}

\begin{proof}
The martingale differences are orthogonal and therefore
\(\Var_\pi(f)=\sum_i\norm{D_if}_2^2\).
For each \(i\), the indicators of \(G_i\) and \(B_i\) are disjoint and
sum to one, so
\[
 \norm{D_if}_2^2
 =\norm{\mathbf1_{G_i}D_if}_2^2
 +\norm{\mathbf1_{B_i}D_if}_2^2.
\]
Summing proves the identity.
\end{proof}

\begin{theorem}[Good-site plus bad-capacity Poincar\'e compiler]
\label{thm:V19-hybrid-Poincare}
Assume the following three stocks.
\begin{enumerate}[label=\textup{(\roman*)},leftmargin=3.6em]
\item For every past \(a\in G_i\), the conditional law \(q_{i,a}\) has
      Poincar\'e floor at least \(\rho_g>0\).
\item The V14 response matrix \(\mathsf A\) satisfies
      \eqref{eq:V14-response-matrix-row}.
\item The bad-increment capacity in \eqref{eq:V19-bad-capacity} is finite.
\end{enumerate}
Then every smooth gauge-invariant \(f\) satisfies
\begin{equation}
 \boxed{
 \Var_\pi(f)
 \le
 \left(
  \frac{\norm{\mathsf A}_{2\to2}^2}{\rho_g}
  +\mathfrak C_{\rm bad}(G_\bullet)
 \right)\mathcal E_\pi(f).}
 \label{eq:V19-hybrid-Poincare}
\end{equation}
For the physical differential Hamiltonian with electric coefficient
\(\kappa_H>0\),
\begin{equation}
 \boxed{
 \gap(H)
 \ge
 \frac{\kappa_H}
 {\rho_g^{-1}\norm{\mathsf A}_{2\to2}^2
  +\mathfrak C_{\rm bad}(G_\bullet)}.}
 \label{eq:V19-hybrid-gap}
\end{equation}
\end{theorem}

\begin{proof}
Fix \(i\) and condition on \(\mathcal F_{i-1}\).  On \(G_i\),
\(D_if\) is the mean-zero fluctuation of \(M_if\) in the coordinate
\(x_i\).  The conditional Poincar\'e inequality gives
\[
 \norm{\mathbf1_{G_i}D_if}_2^2
 \le\rho_g^{-1}
     \norm{\mathbf1_{G_i}\nabla_iM_if}_2^2
 \le\rho_g^{-1}\norm{\nabla_iM_if}_2^2.
\]
Sum over \(i\).  The response-matrix argument of V14 yields
\[
 \sum_i\norm{\nabla_iM_if}_2^2
 \le\norm{\mathsf A}_{2\to2}^2\mathcal E_\pi(f).
\]
The second sum in \eqref{eq:V19-predictable-variance} is at most
\(\mathfrak C_{\rm bad}\mathcal E_\pi(f)\) by definition.  Insert both
bounds into \Cref{lem:V19-predictable-ledger}.  The ground-state-transform
Dirichlet identity multiplies \(\mathcal E_\pi\) by \(\kappa_H\), proving
the gap estimate.
\end{proof}

\begin{proposition}[A checkable routing-matrix sufficient condition]
\label{prop:V19-bad-routing-matrix}
Suppose there is a nonnegative matrix
\(\mathsf C=(C_{ib})_{1\le i,b\le N}\) such that, for every smooth
gauge-invariant \(f\),
\begin{equation}
 \norm{\mathbf1_{B_i}D_if}_{L^2(\pi)}
 \le
 \sum_{b=1}^NC_{ib}\norm{\nabla_bf}_{L^2(\pi)}
 \qquad(1\le i\le N).
 \label{eq:V19-bad-routing-row}
\end{equation}
Then
\begin{equation}
 \boxed{
 \mathfrak C_{\rm bad}(G_\bullet)
 \le\norm{\mathsf C}_{2\to2}^2.}
 \label{eq:V19-bad-routing-capacity}
\end{equation}
Consequently, \eqref{eq:V19-hybrid-gap} holds with
\(\mathfrak C_{\rm bad}\) replaced by
\(\norm{\mathsf C}_{2\to2}^2\).
\end{proposition}

\begin{proof}
Put \(X_b=\norm{\nabla_bf}_2\) and
\(Z_i=\norm{\mathbf1_{B_i}D_if}_2\).
Hypothesis \eqref{eq:V19-bad-routing-row} says \(Z\le\mathsf CX\)
entrywise.  Hence
\[
 \sum_iZ_i^2
 \le\norm{\mathsf C}_{2\to2}^2\sum_bX_b^2
 =\norm{\mathsf C}_{2\to2}^2\mathcal E_\pi(f).
\]
Take the supremum in \eqref{eq:V19-bad-capacity}.
\end{proof}

\begin{assessment}[The precise role of the V13 bridge]
\label{ass:V19-V13-role}
The new matrix \(\mathsf C\) is not a second response matrix.  The V14
matrix \(\mathsf A\) transports coordinate gradients on conditionals where
the site floor is already known.  The V19 matrix \(\mathsf C\) must instead
route the martingale fluctuation supported on bad pasts into physical
Dirichlet energy.  A V13 Perron flow or equilibrium-potential construction
is a legitimate way to prove \eqref{eq:V19-bad-routing-row}.  Merely setting
\(C_{ib}\) from \(\pi(B_i)\), without a conductance or capacity estimate,
is not.
\end{assessment}

% =============================================================================
\section{The score-threshold specialization}
\label{sec:V19-score-threshold}
% =============================================================================

Use the V18 conditional score stocks \(m_i(a)\) and \(\ell_i(a)\).  For
\(T,L>0\), define the predictable events
\begin{equation}
 G_i(T,L)
 :=\{a:m_i(a)\le T,\ \ell_i(a)\le L\},
 \qquad
 B_i(T,L):=G_i(T,L)^c.
 \label{eq:V19-score-good-bad}
\end{equation}

\begin{corollary}[Score-threshold hybrid gap]
\label{cor:V19-score-hybrid-gap}
Let
\begin{equation}
 \rho_{T,L}
 :=C_F\exp\{-2D_G(T+D_GL)\}.
 \label{eq:V19-threshold-site-floor}
\end{equation}
If the V14 response matrix exists and
\(\mathfrak C_{\rm bad}(T,L)\), defined from
\eqref{eq:V19-bad-capacity} using \eqref{eq:V19-score-good-bad}, is finite,
then
\begin{equation}
 \boxed{
 \gap(H)
 \ge
 \frac{\kappa_H}
 {C_F^{-1}e^{2D_G(T+D_GL)}
  \norm{\mathsf A}_{2\to2}^2
  +\mathfrak C_{\rm bad}(T,L)}.}
 \label{eq:V19-score-hybrid-gap}
\end{equation}
The same conclusion holds with
\(\mathfrak C_{\rm bad}(T,L)\) replaced by a routing-matrix bound
\(\norm{\mathsf C(T,L)}_{2\to2}^2\).
\end{corollary}

\begin{proof}
On \(G_i(T,L)\), the V18 \(L^2\)-plus-Lipschitz theorem gives the common
conditional floor \(\rho_{T,L}\).  Apply
\Cref{thm:V19-hybrid-Poincare} and substitute
\eqref{eq:V19-threshold-site-floor}.  The matrix specialization is
\Cref{prop:V19-bad-routing-matrix}.
\end{proof}

\begin{corollary}[Fixed-threshold cofinal certificate]
\label{cor:V19-cofinal-hybrid}
Consider a response-matched cofinal family indexed by \(j\).  A uniform
Hamiltonian gap follows if there are fixed
\(T,L,A_*,C_*,\kappa_*>0\) such that
\begin{equation}
 \boxed{
 \kappa_{H,j}\ge\kappa_*,
 \qquad
 \norm{\mathsf A_j}_{2\to2}\le A_*,
 \qquad
 \mathfrak C_{{\rm bad},j}(T,L)\le C_*}
 \label{eq:V19-cofinal-stocks}
\end{equation}
and the good events are those in \eqref{eq:V19-score-good-bad}.  Explicitly,
\begin{equation}
 \boxed{
 \inf_j\gap(H_j)
 \ge
 \frac{\kappa_*}
 {C_F^{-1}e^{2D_G(T+D_GL)}A_*^2+C_*}>0.}
 \label{eq:V19-cofinal-gap}
\end{equation}
\end{corollary}

\begin{proof}
The right-hand side of \eqref{eq:V19-score-hybrid-gap} is bounded below by
the displayed constant, uniformly in \(j\).
\end{proof}

\begin{lemma}[Why a union-bound threshold cannot close the volume limit]
\label{lem:V19-union-bound-obstruction}
Assume the V17 kinetic-density stock and a uniform bound
\(\ell_i(a)\le L\).  Then
\begin{equation}
 \pi(B_i(T,L))
 \le\frac{e_{\rm kin}}{\kappa_HT^2}.
 \label{eq:V19-single-bad-probability}
\end{equation}
For the first-bad stopping time
\(\tau_T=\min\{i:m_i(X_{<i})>T\}\), with
\(\min\varnothing=N+1\),
\begin{equation}
 \pi(\tau_T\le N)
 \le\frac{Ne_{\rm kin}}{\kappa_HT^2}.
 \label{eq:V19-first-bad-union}
\end{equation}
Keeping this upper bound at most a fixed \(\delta\in(0,1)\) by the union
bound requires
\begin{equation}
 T\ge
 \left(\frac{Ne_{\rm kin}}{\kappa_H\delta}\right)^{1/2}.
 \label{eq:V19-union-threshold}
\end{equation}
At that threshold, the site-floor guarantee
\eqref{eq:V19-threshold-site-floor} is no better than an exponential in
\(-\sqrt N\).  Hence this probability-only tuning does not yield the fixed
stocks in \eqref{eq:V19-cofinal-stocks}.
\end{lemma}

\begin{proof}
The first estimate is V18's Markov bound, with the bad event reduced to
\(m_i>T\) by the assumed uniform Lipschitz stock.  The union bound gives
\eqref{eq:V19-first-bad-union}, and solving its right-hand side for \(T\)
gives \eqref{eq:V19-union-threshold}.  Substitution into
\eqref{eq:V19-threshold-site-floor} gives the stated behavior.
\end{proof}

\begin{firewall}[Capacity, not rarity]
\label{fire:V19-capacity-not-rarity}
\Cref{lem:V19-union-bound-obstruction} concerns the guarantee supplied by
the Fisher/Markov/union-bound argument; it does not assert that the actual
conditional gaps decay.  It proves that the averaged kinetic score cannot
by itself populate the cofinal gap certificate.  The missing input is the
fixed-threshold bad-increment capacity in
\eqref{eq:V19-bad-capacity}, preferably through the checkable routing rows
\eqref{eq:V19-bad-routing-row}.
\end{firewall}

\begin{assessment}[Threshold optimization without discarding a sector]
\label{ass:V19-threshold-optimization}
Once \(\mathfrak C_{\rm bad}(T,L)\) is estimated independently, the rigorous
tradeoff is
\begin{equation}
 \gap(H)
 \ge\kappa_H
 \sup_{T,L>0}
 \left[
 C_F^{-1}e^{2D_G(T+D_GL)}\norm{\mathsf A}_{2\to2}^2
 +\mathfrak C_{\rm bad}(T,L)
 \right]^{-1}.
 \label{eq:V19-threshold-optimization}
\end{equation}
Increasing the thresholds enlarges the good set but weakens the guaranteed
site floor.  Decreasing them strengthens the site floor but enlarges the
sector that must be carried by capacity.  Formula
\eqref{eq:V19-threshold-optimization} keeps both effects and loses no
variance term.
\end{assessment}



% =============================================================================
\section{V19 ledger and revised frontier}
\label{sec:V19-ledger}
% =============================================================================

\begin{theorem}[Integrated V19 statement]
\label{thm:V19-integrated}
Preserving V1--V18, V19 proves:
\begin{enumerate}[label=\textup{(V19.\arabic*)},leftmargin=5.5em]
\item exact covariance of the comparison, defect, edge, landing, and
      Rayleigh quantities under a common physical Hamiltonian scaling;
\item identification of the physical V18 landing stock as the scaled
      source-carrier defect moment \eqref{eq:V19-physical-carrier-stock};
\item an exact relative-defect representation of the Woodbury edge;
\item the volume-independent pure-Wilson defect estimate
      \(\norm{\Theta_\eta^{\rm W}}\le192\alpha_\beta\) extracted from the
      V35 local rows;
\item the explicit unsigned relative-edge test
      \eqref{eq:V19-Wilson-relative-test};
\item an exact predictable good/bad variance ledger and a hybrid Poincar\'e
      theorem with one separate bad-increment capacity;
\item a routing-matrix sufficient condition for that capacity;
\item the score-threshold gap formula
      \eqref{eq:V19-score-hybrid-gap} and a proof that probability-only
      union-bound tuning cannot produce a volume-uniform site floor.
\end{enumerate}
\end{theorem}

\begin{proof}
Items \textup{(V19.1)--(V19.2)} are
\Cref{thm:V19-scale-covariance,cor:V19-Rayleigh-scaling}.  Item
\textup{(V19.3)} is \Cref{lem:V19-relative-defect}.  Items
\textup{(V19.4)--(V19.5)} are
\Cref{prop:V19-Wilson-defect-row,cor:V19-Wilson-relative-test}.  Items
\textup{(V19.6)--(V19.7)} are
\Cref{lem:V19-predictable-ledger,thm:V19-hybrid-Poincare,
prop:V19-bad-routing-matrix}.  Item \textup{(V19.8)} is
\Cref{cor:V19-score-hybrid-gap,lem:V19-union-bound-obstruction}.
\end{proof}

\begingroup
\small
\begin{longtable}{>{\raggedright\arraybackslash}p{0.28\textwidth}
                  >{\raggedright\arraybackslash}p{0.16\textwidth}
                  >{\raggedright\arraybackslash}p{0.48\textwidth}}
\toprule
\textbf{V19 row}&\textbf{Status}&\textbf{Advance or remaining obstruction}\\
\midrule
\endfirsthead
\toprule
\textbf{V19 row}&\textbf{Status}&\textbf{Advance or remaining obstruction}\\
\midrule
\endhead
Common normalization covariance
& \MGpass
& Edge invariant; landing metric and normalized Rayleigh quotient scale
  explicitly.\\[0.35em]
Pure-Wilson absolute defect row
& \MGpass
& \(192\alpha_\beta\), uniformly in volume, before generated heads.\\[0.35em]
Dimensionless unsigned edge test
& \MGpass
& Exact sufficient ratio \eqref{eq:V19-Wilson-relative-test}; it need not
  be below one.\\[0.35em]
Predictable good/bad compiler
& \MGpass
& No variance term is discarded; bad pasts enter through
  \(\mathfrak C_{\rm bad}\).\\[0.35em]
Physical normalization factor \(s_j\)
& \MGopen
& Must be extracted from the V11 same-generator time derivative.\\[0.35em]
Uniform physical carrier defect
& \MGopen
& Requires \(\sup_js_jd_{\mathscr C_j}<\infty\), or another landing
  argument.\\[0.35em]
Uniform source-channel edge
& \MGopen
& The V35 unsigned row ratio alone is not a certified strict contraction.\\[0.35em]
Bad-past routing matrix
& \MGopen
& Populate \eqref{eq:V19-bad-routing-row} by V13 physical Perron flows.\\[0.35em]
Kinetic density and signed diagonal curvature
& \MGopen
& Needed to quantify the bad set and its good-site complement.\\[0.35em]
Protected sectors and nontrivial OS limit
& \MGopen
& Must hold uniformly on every retained sector before the continuum passage.\\[0.35em]
Full Yang--Mills mass gap
& \MGopen
& The cofinal stocks and the nontrivial OS construction remain unproved.\\
\bottomrule
\end{longtable}
\endgroup

\begin{openproblem}[V20 mandate: populate the scale and capacity stocks]
\label{op:V20}
Proceed on the response-matched weak-coupling branch in the following order.
Because V20 is a five-version checkpoint, first audit the current SM and NS
notebooks in \texttt{latex\_notes}; import only proved transferable stocks,
record any changed YM direction, and then execute the list below.
\begin{enumerate}[label=\textup{(V20.\arabic*)},leftmargin=5.5em]
\item Extract the common multiplier \(s_j\) from the V11 exact temporal
      derivative and record the physical \(\kappa_{H,j}\), without changing
      the Perron law.
\item Evaluate \(s_jd_{\mathscr C_j}\) on the actual V15 signed-source
      carrier.  If the global \(192s_j\alpha_{\beta_j}\) bound diverges, do
      not declare failure before testing the carrier moment.
\item Replace the unsigned relative bound
      \eqref{eq:V19-Wilson-relative-test} by a source-cyclic edge estimate
      using the V16 exponential moment and the complete signed V15 curvature
      source.
\item Prove the V17 kinetic-energy density and the V18 diagonal Lipschitz
      tensor bound on the same response-matched family.
\item Fix \(T,L\) and construct V13 equilibrium-potential or flow estimates
      that populate a volume-summable routing matrix \(\mathsf C(T,L)\) in
      \eqref{eq:V19-bad-routing-row}.
\item Insert \(\mathsf A\) and \(\mathsf C\) into
      \eqref{eq:V19-score-hybrid-gap}; then repeat the certificate on every
      protected toron/topology sector.
\item Only after a common positive gap is obtained, perform the V1
      Osterwalder--Schrader passage and prove nontriviality of the limiting
      Yang--Mills state.
\end{enumerate}
\end{openproblem}

\begin{firewall}[End-of-V19 claim boundary]
V19 proves exact physical-scaling identities, a new pure-Wilson defect bound,
and a rigorous good-site/bad-capacity martingale compiler.  It does not
determine the V11 physical multiplier, prove a strict weak-coupling source
edge, bound the bad-increment capacity, control protected sectors, or
construct a nontrivial continuum OS state.  Therefore it does not prove the
full Yang--Mills mass gap.
\end{firewall}

\section*{Version V19 append-only record}
\addcontentsline{toc}{section}{Version V19 append-only record}
The complete V18 source was retained before this continuation.  Its SHA--256
digest was
\begin{center}\scriptsize\ttfamily
523a952e9b6a7bb70260bb478b88df8955207461b062ab2618e12337e4842261.
\end{center}
On the next iteration, retain this full file, append before its unique active
document-closing command, increment the filename to
\texttt{\_V20.tex}, and remove only the superseded V19 file from this note
series.

% =============================================================================
% END APPENDED VERSION V19 MATERIAL
% =============================================================================

% =============================================================================
% BEGIN APPENDED VERSION V20 MATERIAL
% =============================================================================

\clearpage
\part{Version V20: five-version cross-program audit, physical Witten scale, and complete-source capacity incidence}
\label{part:V20}

\section{Five-version checkpoint and adjusted direction}
\label{sec:V20-cross-audit}

This is the first five-version checkpoint requested for the continuing
Yang--Mills series.  Before choosing the V20 calculation, the current
companion notebooks in \texttt{latex\_notes} were read without modifying
them:
\begin{enumerate}[label=\textup{(\roman*)},leftmargin=3.6em]
\item \texttt{Renewal\_SM\_Einstein\_Continuum\_Research\_Notes\_V37.tex},
      SHA--256
      \[
      \texttt{9d7824190b5fb289b542f20fe1365745687fa42f01b1cd9ae22ebc684c3bc878};
      \]
\item \texttt{Renewal\_NS\_Stability\_Research\_Notes\_V11.tex},
      SHA--256
      \[
      \texttt{c0a3a5b83b0ef4a297d6eead49fa7a0bb58abffec4facbd07d235faf19e9eedc}.
      \]
      Its current appended text ends with an internal V12 conclusion and
      V13 acquisition order; this audit cites that content but does not
      rename or edit the companion file.
\end{enumerate}

No quantitative theorem is imported merely by analogy.  Two proof
disciplines do transfer.
\begin{enumerate}[label=\textup{(A\arabic*)},leftmargin=4.2em]
\item The SM V37 glancing calculation keeps all finite operators in one
      normalized frame and returns explicit right and left kernel witnesses
      when a selector angle collapses.  For the YM programme this reinforces
      the requirement that the V35 bare Witten card first be placed in the
      V8/V11 physical frame, and that failure of a source edge return an
      actual carrier vector rather than a failed global majorant.
\item The current NS endpoint calculation distinguishes the complete
      exterior source from a raw residual, preserves signed interference
      before taking norms, and proves that a sharp local capacity may still
      be summable and insufficient for global closure.  For YM this changes
      the capacity order: first prove incidence from the actual bad
      martingale increment to the complete V15 Perron source, then construct
      the V13 flow.  A flow for the unsigned V35 defect alone is not a
      substitute.
\end{enumerate}

Accordingly, V20 first resolves the common scalar in the V19 normalization
audit, next determines the true size of the global pure-Wilson defect by a
centre-flux witness, and finally factors the V19 bad-capacity row into an
exact complete-source incidence followed by a capacity routing estimate.
The source-carrier compression remains distinct from the global defect
norm.

\begin{protocol}[Future five-version checkpoints]
\label{prot:V20-five-version-cadence}
At V25, V30, V35, and every fifth YM version thereafter, read the newest
retained SM--Einstein and NS notebooks in \texttt{latex\_notes} before
choosing the next YM calculation.  Record the filenames and internal latest
versions actually read.  Import only proved statements whose hypotheses
are matched in the YM setting; otherwise record a directional lesson, not a
theorem.  Do not modify either companion notebook.
\end{protocol}

% =============================================================================
\section{The exact physical multiplier of the one-form Witten card}
\label{sec:V20-physical-multiplier}
% =============================================================================

V19 used \(s_j>0\) for an abstract common scalar multiplying a bare
comparison-minus-defect split.  The scalar is not arbitrary once that split
has been identified with the actual V8 ground-law Witten operator.

\begin{theorem}[Physical normalization under same-law incidence]
\label{thm:V20-physical-multiplier}
At a fixed finite regulator, let
\[
 H=\kappa_H\sum_e(-\Delta_e)+V_B-E_0
\]
have normalized positive ground state \(h\), and let
\(\dd\pi=h^2\dd\mu\).  Suppose the V6 same-law incidence holds and a bare
one-form split on the horizontal exact range satisfies
\begin{equation}
 \mathscr W_{1}^{\rm bare}
 =\mathcal A^{\rm bare}-\Theta^{\rm bare}
 =dd_\pi^*+d_\pi^*d.
 \label{eq:V20-bare-same-law-split}
\end{equation}
Then the physical one-form operator is
\begin{equation}
 \boxed{
 \mathscr L_1^{\rm phys}
 =\kappa_H\mathscr W_{1}^{\rm bare}
 =\mathcal A^{\rm phys}-\Theta^{\rm phys},}
 \label{eq:V20-physical-Witten-scale}
\end{equation}
where
\begin{equation}
 \boxed{
 \mathcal A^{\rm phys}=\kappa_H\mathcal A^{\rm bare},
 \qquad
 \Theta^{\rm phys}=\kappa_H\Theta^{\rm bare},
 \qquad
 \mathcal V^{\rm phys}=\sqrt{\kappa_H}\,
                       \mathcal V^{\rm bare}.}
 \label{eq:V20-physical-triple}
\end{equation}
Thus the abstract V19 common scalar is
\begin{equation}
 \boxed{s_j=\kappa_{H,j}.}
 \label{eq:V20-s-equals-kappa}
\end{equation}
For the V10--V11 temporal calibration,
\begin{equation}
 \boxed{
 \kappa_H=\frac{s_\alpha}{\tau_\alpha},
 \qquad
 s_\alpha=-\frac{\log\eta_F(\alpha)}{C_F}.}
 \label{eq:V20-temporal-kappa}
\end{equation}
In particular, the common factor is not \(\tau_\alpha^{-1}\) by itself.
\end{theorem}

\begin{proof}
V6 gives the scalar ground-state transform
\[
 \mathcal U_h^{-1}H\mathcal U_h
 =\kappa_Hd_\pi^*d.
\]
With \(D=\sqrt{\kappa_H}\,d\), its exact one-form partner is
\[
 DD^*=\kappa_Hdd_\pi^*
\]
on the closure of exact one-forms.  Equivalently, the V8 Bochner card is
\(\kappa_H(dd_\pi^*+d_\pi^*d)\), whose exact-range restriction is \(DD^*\).
Substitute \eqref{eq:V20-bare-same-law-split} and distribute the positive
scalar.  If
\(\Theta^{\rm bare}=\mathcal V^{\rm bare}
(\mathcal V^{\rm bare})^*\), the last equality in
\eqref{eq:V20-physical-triple} factors
\(\Theta^{\rm phys}\).  This is precisely the V19 common scaling, proving
\eqref{eq:V20-s-equals-kappa}.  Finally, V10 defines
\(\tau_\alpha=s_\alpha/\kappa_H\); rearrangement gives
\eqref{eq:V20-temporal-kappa}.
\end{proof}

\begin{corollary}[Physical landing and edge stocks in calibrated variables]
\label{cor:V20-calibrated-stocks}
On a declared bare channel carrier \(\mathscr C_j\), put
\[
 d_{\mathscr C_j}^{\rm bare}
 :=\sup_{\substack{z\in\mathscr C_j\\\norm{z}=1}}
   \norm{\mathcal V_j^{\rm bare}z}^2.
\]
Under the hypotheses of \Cref{thm:V20-physical-multiplier}, the V18 landing
stock is exactly
\begin{equation}
 \boxed{
 \sup_j\kappa_{H,j}d_{\mathscr C_j}^{\rm bare}<\infty.}
 \label{eq:V20-calibrated-carrier-stock}
\end{equation}
The Woodbury edge \(\mathcal K_j\) is independent of
\(\kappa_{H,j}\), while the normalized physical Rayleigh quotient of a
fixed landed channel is multiplied by \(\kappa_{H,j}\).
\end{corollary}

\begin{proof}
Insert \eqref{eq:V20-s-equals-kappa} into the V19 scale-covariance and
carrier-moment identities.
\end{proof}

\begin{firewall}[The scale identity is conditional on the law identity]
\label{fire:V20-scale-law}
Equation \eqref{eq:V20-s-equals-kappa} resolves the scalar only after
\eqref{eq:V20-bare-same-law-split} has been proved.  It does not identify
the V35 reduced-fibre law with the V8 physical Perron law.  If those laws or
horizontal ranges differ, multiplying by \(\kappa_H\) does not repair the
incidence.  The V6 same-law row remains an independent proof obligation.
\end{firewall}

% =============================================================================

% =============================================================================

% =============================================================================
\section{A centre-flux lower witness for the global Wilson defect}
\label{sec:V20-centre-defect}
% =============================================================================

V19 obtained the upper bound
\(\norm{\Theta_\eta^{\rm W}}\le192\alpha_\beta\).
The next lemma shows that the linear growth in \(\alpha_\beta\) is not only
an artefact of that unsigned upper estimate when the complete defect channel
is used.

\begin{lemma}[Wilson Hessian at a nontrivial central plaquette]
\label{lem:V20-central-Hessian}
Let
\[
 \zeta=e^{2\pi i/3}I\in Z(\mathrm{SU}(3)),
 \qquad
 \Re\zeta=-\frac12.
\]
For one Wilson plaquette
\[
 W_p(U)=\beta\left(1-\frac13\Re\Tr U_p\right),
 \qquad
 \alpha_\beta=\frac{\beta}{3},
\]
let \(\mathcal H_p^0\) be its Hessian when all four background links are
the identity.  If all four background links are central and their oriented
plaquette product is \(\zeta\) or \(\zeta^{-1}\), then
\begin{equation}
 \boxed{
 \mathcal H_p^{\rm cen}
 =-\frac12\mathcal H_p^0,
 \qquad
 \Theta_{p,\eta}^{\rm cen}
 =\left(\frac32-\eta\right)\mathcal H_p^0.}
 \label{eq:V20-central-Hessian}
\end{equation}
\end{lemma}

\begin{proof}
Central background links have trivial adjoint action.  For link tangents
\(a\), the first plaquette variation is the abelianized curl
\(X=(d_1a)_p\in\mathfrak{su}(3)\).  Commutator terms in the second
variation have zero trace.  Hence
\[
 D^2W_p[a,a]
 =-\frac{\beta}{3}\Re\Tr(\zeta X^2).
\]
The metric convention is
\(\norm{X}^2=-\Tr(X^2)\), and \(\Tr(X^2)\) is real.  Therefore
\[
 -\frac{\beta}{3}\Re\Tr(\zeta X^2)
 =\Re(\zeta)\frac{\beta}{3}\norm{X}^2
 =-\frac12\alpha_\beta\norm{X}^2.
\]
At the identity the same expression is
\(\alpha_\beta\norm{X}^2\).  Polarization proves the first operator identity.
Since \(\mathcal H_p^0\succeq0\),
\[
 (1-\eta)\mathcal H_p^0-\mathcal H_p^{\rm cen}
 =\left(\frac32-\eta\right)\mathcal H_p^0\succeq0,
\]
so taking its positive part changes nothing.
\end{proof}

\begin{theorem}[Two-sided global pure-Wilson defect scale]
\label{thm:V20-two-sided-defect}
Assume the reduced fibre contains a retained fine-link slot \(e\) incident
to a plaquette \(p\), and a unit tangent supported on \(e\) whose normalized
one-colour curl on \(p\) has norm one.  This is a fixed local
nondegeneracy condition, independent of the torus side.  For
\(0<\eta<1\), the pure-Wilson V35 multiplication operator satisfies
\begin{equation}
 \boxed{
 \left(\frac32-\eta\right)\alpha_\beta
 \le\norm{\Theta_\eta^{\rm W}}
 \le192\alpha_\beta.}
 \label{eq:V20-two-sided-defect}
\end{equation}
Consequently, under the same-law physical incidence of
\Cref{thm:V20-physical-multiplier},
\begin{equation}
 \boxed{
 \left(\frac32-\eta\right)\kappa_H\alpha_\beta
 \le\norm{\Theta_{\eta}^{\rm W,phys}}
 =\norm{\mathcal V_{\eta}^{\rm W,phys}}^2
 \le192\kappa_H\alpha_\beta.}
 \label{eq:V20-two-sided-physical-defect}
\end{equation}
For every cofinal family satisfying the stated local nondegeneracy,
\begin{equation}
 \boxed{
 \sup_j\norm{\mathcal V_{\eta,j}^{\rm W,phys}}<\infty
 \quad\Longleftrightarrow\quad
 \sup_j\kappa_{H,j}\alpha_{\beta_j}<\infty.}
 \label{eq:V20-global-defect-equivalence}
\end{equation}
\end{theorem}

\begin{proof}
Set the retained link \(e\) equal to \(\zeta\) and every other reduced link
equal to the identity.  Each plaquette incident to \(e\) then has holonomy
\(\zeta\) or \(\zeta^{-1}\).  By
\Cref{lem:V20-central-Hessian}, its defect is
\((3/2-\eta)\mathcal H_p^0\).  For the unit tangent in the hypothesis,
the V35 flat-Hessian normalization gives
\[
 \ip{a}{\mathcal H_p^0a}=\alpha_\beta.
\]
The total defect is a sum of positive plaquette defects, so its pointwise
operator norm at this central configuration is at least
\((3/2-\eta)\alpha_\beta\).  The reduced-fibre law has a strictly positive
continuous density on the compact fibre.  By continuity, every
neighborhood of the central configuration has positive measure and
contains defect norms arbitrarily close to this value.  The norm of the
multiplication operator is its essential supremum, proving the lower
bound.  The upper bound is \eqref{eq:V19-Wilson-defect-bound}.

Multiply the two bounds by \(\kappa_H\) using
\eqref{eq:V20-physical-triple}; factorization gives the middle equality in
\eqref{eq:V20-two-sided-physical-defect}.  The two constants in the
resulting comparison are positive and regulator independent, which proves
the equivalence.
\end{proof}

\begin{assessment}[What the centre witness does and does not eliminate]
\label{ass:V20-centre-witness}
The global V18 bounded-defect route is now decided for the pure-Wilson
channel: it is available exactly when
\(\kappa_{H,j}\alpha_{\beta_j}\) is uniformly bounded.  If that product
diverges, the complete physical synthesis norm diverges by an explicit
central-field witness; this is stronger than failure of the V19 upper
bound.

This does not yet produce a physical soft mode and does not decide the
V15 source-generated carrier.  The central configuration is a witness for
the norm of the complete multiplication operator.  It need not be reached
by a near-unit source-cyclic channel vector after the actual Perron and
horizontal projections.
\end{assessment}

% =============================================================================
\section{The exact source-carrier compression}
\label{sec:V20-carrier-compression}
% =============================================================================

\begin{proposition}[Carrier defect is a compressed channel Gram]
\label{prop:V20-carrier-Gram}
Let \(\mathscr C\subseteq\mathscr D\) be a closed defect-channel carrier
with orthogonal projection \(P_{\mathscr C}\), and let
\[
 \mathcal G_{\mathcal V}:=\mathcal V^*\mathcal V.
\]
Then the V19 carrier moment is
\begin{equation}
 \boxed{
 d_{\mathscr C}
 =\norm{\mathcal VP_{\mathscr C}}^2
 =\norm{P_{\mathscr C}\mathcal G_{\mathcal V}
              P_{\mathscr C}}.}
 \label{eq:V20-compressed-Gram}
\end{equation}
Under physical normalization,
\begin{equation}
 \boxed{
 d_{\mathscr C}^{\rm phys}
 =\kappa_H
  \norm{P_{\mathscr C}\mathcal G_{\mathcal V}^{\rm bare}
              P_{\mathscr C}}.}
 \label{eq:V20-physical-compressed-Gram}
\end{equation}
Thus a global lower bound for \(\norm{\Theta}\) does not imply a lower bound
on \(d_{\mathscr C}\).
\end{proposition}

\begin{proof}
By definition,
\[
 d_{\mathscr C}
 =\sup_{\norm{z}=1}\norm{\mathcal VP_{\mathscr C}z}^2
 =\norm{\mathcal VP_{\mathscr C}}^2.
\]
For every bounded operator \(T\),
\(\norm{T}^2=\norm{T^*T}\).  Taking
\(T=\mathcal VP_{\mathscr C}\) proves
\eqref{eq:V20-compressed-Gram}.  Equation
\eqref{eq:V20-physical-compressed-Gram} follows from
\(\mathcal V^{\rm phys}=\sqrt{\kappa_H}\mathcal V^{\rm bare}\).
For the final statement, take
\(\mathcal V_j=\operatorname{diag}(j,1)\) on \(\mathbb R^2\) and
\(\mathscr C=\operatorname{span}\{(0,1)\}\).  Then
\(\norm{\Theta_j}=\norm{\mathcal V_j}^2=j^2\), while
\(d_{\mathscr C}=1\).
\end{proof}

\begin{corollary}[Sharp remaining landing acquisition]
\label{cor:V20-source-carrier-acquisition}
Let \(\mathscr C_j^{\rm src}\) be the closed channel carrier generated by
the complete signed V15 curvature sources after the actual horizontal,
conditional, and protected-sector projections.  The V18 no-ultraviolet
landing hypothesis is precisely
\begin{equation}
 \boxed{
 \sup_j\kappa_{H,j}
 \norm{
 P_{\mathscr C_j^{\rm src}}
 \mathcal G_{\mathcal V,j}^{\rm bare}
 P_{\mathscr C_j^{\rm src}}}<\infty.}
 \label{eq:V20-sharp-source-landing}
\end{equation}
If \(\kappa_{H,j}\alpha_{\beta_j}\to\infty\), the global synthesis route is
closed by \Cref{thm:V20-two-sided-defect}, but
\eqref{eq:V20-sharp-source-landing} remains a logically possible
source-specific route.
\end{corollary}

\begin{proof}
Apply \Cref{prop:V20-carrier-Gram} to
\(\mathscr C_j^{\rm src}\) and compare with
\eqref{eq:V20-calibrated-carrier-stock}.
\end{proof}

\begin{firewall}[Do not replace the source compression by the global norm]
\label{fire:V20-source-compression}
The inequality
\[
 \norm{
 P_{\mathscr C_j^{\rm src}}\mathcal G_{\mathcal V,j}
 P_{\mathscr C_j^{\rm src}}}
 \le\norm{\Theta_j}
\]
is only an upper bound.  When its right-hand side diverges, it says nothing
about the compressed left-hand side.  Conversely, proving the compression
bounded does not bound defect channels which are never generated by the
physical V15 source.  The two statements serve different theorems and must
not be interchanged.
\end{firewall}

% =============================================================================

% =============================================================================

% =============================================================================
\section{Complete-source incidence before bad-set capacity}
\label{sec:V20-complete-source-incidence}
% =============================================================================

The V19 bad-capacity norm is defined by the actual martingale increments
\(\mathbf1_{B_i}D_if\).  A local V13 flow will usually be constructed for a
more geometric source.  The following compiler separates the required
incidence from the subsequent capacity estimate.

\begin{definition}[Signed assembled source packet]
\label{def:V20-assembled-source}
For each source index \(a\), let
\begin{equation}
 r_a(f):=\sum_{\sigma\in\mathfrak S_a}
          \varepsilon_{a,\sigma}r_{a,\sigma}(f),
 \qquad
 \varepsilon_{a,\sigma}\in\{-1,1\},
 \qquad
 R_a(f):=\norm{r_a(f)}.
 \label{eq:V20-assembled-source}
\end{equation}
The signed sum is formed in its native Hilbert space before the norm is
taken.  In the YM application the summands include the direct and
conditional-covariance pieces of the complete V15 Perron curvature source,
with their actual conditioning incidences.
\end{definition}

\begin{theorem}[Incidence--capacity factorization]
\label{thm:V20-incidence-capacity}
For a predictable bad family \(B_i\), put
\[
 Z_i(f):=\norm{\mathbf1_{B_i}D_if}_{L^2(\pi)},
 \qquad
 X_b(f):=\norm{\nabla_bf}_{L^2(\pi)}.
\]
Suppose there are nonnegative matrices
\[
 \mathsf I\in\mathbb R_+^{N\times M},
 \qquad
 \mathsf Q\in\mathbb R_+^{M\times N},
 \qquad
 \mathsf E\in\mathbb R_+^{N\times N}
\]
such that every smooth gauge-invariant \(f\) obeys, entrywise,
\begin{align}
 Z(f)&\le\mathsf I R(f)+\mathsf E X(f),
 \label{eq:V20-source-incidence}\\
 R(f)&\le\mathsf QX(f).
 \label{eq:V20-source-capacity}
\end{align}
Then the V19 bad-increment capacity satisfies
\begin{equation}
 \boxed{
 \mathfrak C_{\rm bad}(G_\bullet)
 \le
 \norm{\mathsf I\mathsf Q+\mathsf E}_{2\to2}^2.}
 \label{eq:V20-incidence-capacity-bound}
\end{equation}
The matrix \(\mathsf I\) is the complete-source incidence, \(\mathsf Q\)
is the physical capacity routing, and \(\mathsf E\) records a controlled
incidence residual.  Exact incidence corresponds to \(\mathsf E=0\).
\end{theorem}

\begin{proof}
Substitute \eqref{eq:V20-source-capacity} into
\eqref{eq:V20-source-incidence}:
\[
 Z(f)
 \le(\mathsf I\mathsf Q+\mathsf E)X(f)
\]
entrywise.  Therefore
\[
 \sum_iZ_i(f)^2
 \le
 \norm{\mathsf I\mathsf Q+\mathsf E}_{2\to2}^2
 \sum_bX_b(f)^2.
\]
The denominator on the right is \(\mathcal E_\pi(f)\).  Take the supremum
in \eqref{eq:V19-bad-capacity}.
\end{proof}

\begin{corollary}[Complete-source hybrid mass certificate]
\label{cor:V20-complete-source-gap}
Under the good-site and response hypotheses of
\Cref{thm:V19-hybrid-Poincare}, and the two rows
\eqref{eq:V20-source-incidence}--\eqref{eq:V20-source-capacity},
\begin{equation}
 \boxed{
 \gap(H)
 \ge
 \frac{\kappa_H}{
 \rho_g^{-1}\norm{\mathsf A}_{2\to2}^2
 +\norm{\mathsf I\mathsf Q+\mathsf E}_{2\to2}^2}.}
 \label{eq:V20-complete-source-gap}
\end{equation}
For the V18 score thresholds one may take
\(\rho_g=C_Fe^{-2D_G(T+D_GL)}\).
\end{corollary}

\begin{proof}
Insert \eqref{eq:V20-incidence-capacity-bound} into
\eqref{eq:V19-hybrid-gap}, and then use
\eqref{eq:V19-threshold-site-floor} for the final specialization.
\end{proof}

\begin{firewall}[No capacity estimate before source incidence]
\label{fire:V20-incidence-first}
A bound on a raw Wilson defect, a local action residual, or one unsigned
piece \(r_{a,\sigma}\) does not prove
\eqref{eq:V20-source-incidence}.  If the complete bad increment contains a
future Perron covariance or a protected-sector component absent from that
proxy, the first row fails even when the proxy has an excellent flow.  The
signed source \eqref{eq:V20-assembled-source} must be assembled under the
actual law before triangle inequalities are used.  This is the precise YM
consequence of the NS checkpoint; no Navier--Stokes estimate is imported.
\end{firewall}

\begin{lemma}[Failure of the sharp bad-capacity stock returns a witness]
\label{lem:V20-bad-capacity-witness}
For a cofinal family, if
\[
 \sup_j\mathfrak C_{{\rm bad},j}(G_{\bullet,j})=\infty,
\]
then there are indices \(j_n\) and smooth gauge-invariant functions \(f_n\)
such that
\begin{equation}
 \boxed{
 \mathcal E_{\pi_{j_n}}(f_n)=1,
 \qquad
 \sum_i
 \norm{\mathbf1_{B_{i,j_n}}D_if_n}_{L^2(\pi_{j_n})}^2
 \longrightarrow\infty.}
 \label{eq:V20-bad-capacity-witness}
\end{equation}
Thus genuine failure of the sharp stock produces an actual normalized bad
martingale sequence.  Divergence of a chosen majorant
\(\norm{\mathsf I\mathsf Q+\mathsf E}\) alone does not.
\end{lemma}

\begin{proof}
Choose \(j_n\) with
\(\mathfrak C_{{\rm bad},j_n}\ge2n\).  By the definition of the supremum,
choose \(g_n\) whose quotient in \eqref{eq:V19-bad-capacity} is at least
\(n\).  Rescale \(g_n\) so that its Dirichlet energy is one.  This gives
\eqref{eq:V20-bad-capacity-witness}.  The last statement follows because
\eqref{eq:V20-incidence-capacity-bound} is only a sufficient upper
majorant.
\end{proof}

\begin{assessment}[Adjusted acquisition order after the checkpoint]
\label{ass:V20-adjusted-order}
The V20 checkpoint fixes the following order.
\begin{enumerate}[label=\textup{(\arabic*)},leftmargin=3.4em]
\item Prove the same-law/horizontal-range incidence; only then use
      \(s_j=\kappa_{H,j}\).
\item Determine the response-matched product
      \(\kappa_{H,j}\alpha_{\beta_j}\).  This decides the complete-channel
      landing route by \eqref{eq:V20-global-defect-equivalence}.
\item Independently compute the compressed V15 source Gram in
      \eqref{eq:V20-sharp-source-landing}; global divergence does not decide
      it.
\item Assemble the signed V15 source and prove
      \eqref{eq:V20-source-incidence} before constructing a V13 flow for
      \eqref{eq:V20-source-capacity}.
\item Keep explicit normalized witnesses for failure of the source edge,
      carrier landing, or sharp bad-capacity stock, in the common physical
      frame.
\end{enumerate}
This reordering avoids spending a capacity estimate on a source that has not
been shown to equal the bad physical fluctuation.
\end{assessment}

% =============================================================================
\section{V20 ledger and revised frontier}
\label{sec:V20-ledger}
% =============================================================================

\begin{theorem}[Integrated V20 statement]
\label{thm:V20-integrated}
Preserving V1--V19, V20 proves:
\begin{enumerate}[label=\textup{(V20.\arabic*)},leftmargin=5.5em]
\item the required five-version audit of the current SM and NS companion
      notebooks and an explicit future audit cadence;
\item under same-law incidence, the exact physical multiplier
      \(s_j=\kappa_{H,j}=s_{\alpha_j}/\tau_{\alpha_j}\);
\item a central-\(\mathrm{SU}(3)\) plaquette identity
      \(\mathcal H_p^{\rm cen}=-\frac12\mathcal H_p^0\);
\item the two-sided global pure-Wilson defect scale
      \eqref{eq:V20-two-sided-physical-defect};
\item the exact criterion
      \(\sup_j\kappa_{H,j}\alpha_{\beta_j}<\infty\) for bounded complete
      pure-Wilson physical synthesis;
\item identification of the sharp V18 source landing stock with the
      compressed channel Gram \eqref{eq:V20-sharp-source-landing};
\item an incidence--capacity factorization of the V19 bad-martingale norm,
      with the explicit gap bound \eqref{eq:V20-complete-source-gap};
\item an actual normalized martingale witness whenever the sharp
      bad-capacity stock diverges.
\end{enumerate}
\end{theorem}

\begin{proof}
Item \textup{(V20.1)} is \Cref{sec:V20-cross-audit} and
\Cref{prot:V20-five-version-cadence}.  Item \textup{(V20.2)} is
\Cref{thm:V20-physical-multiplier}.  Items
\textup{(V20.3)--(V20.5)} are
\Cref{lem:V20-central-Hessian,thm:V20-two-sided-defect}.  Item
\textup{(V20.6)} is
\Cref{prop:V20-carrier-Gram,cor:V20-source-carrier-acquisition}.  Item
\textup{(V20.7)} is
\Cref{thm:V20-incidence-capacity,cor:V20-complete-source-gap}.  Item
\textup{(V20.8)} is \Cref{lem:V20-bad-capacity-witness}.
\end{proof}

\begingroup
\small
\begin{longtable}{>{\raggedright\arraybackslash}p{0.28\textwidth}
                  >{\raggedright\arraybackslash}p{0.16\textwidth}
                  >{\raggedright\arraybackslash}p{0.48\textwidth}}
\toprule
\textbf{V20 row}&\textbf{Status}&\textbf{Advance or remaining obstruction}\\
\midrule
\endfirsthead
\toprule
\textbf{V20 row}&\textbf{Status}&\textbf{Advance or remaining obstruction}\\
\midrule
\endhead
Five-version SM/NS checkpoint
& \MGpass
& Direction changed to complete-source incidence before capacity.\\[0.35em]
Physical Witten scalar
& \MGpass
& \(s_j=\kappa_{H,j}=s_{\alpha_j}/\tau_{\alpha_j}\), conditional on the
  same-law incidence.\\[0.35em]
Global pure-Wilson defect scale
& \MGpass
& Matching upper and central-field lower bounds of order
  \(\kappa_H\alpha_\beta\).\\[0.35em]
Source-carrier compression
& \MGpass
& Exact compressed Gram; not decided by the global norm.\\[0.35em]
Incidence--capacity compiler
& \MGpass
& Separates complete physical source incidence from V13 routing.\\[0.35em]
V35/V8 same-law and horizontal-range incidence
& \MGopen
& Required before the physical multiplier can be applied to the V35 card.\\[0.35em]
Response-matched \(\kappa_H\alpha_\beta\)
& \MGopen
& Must be derived from the anisotropic transfer coefficients on the spatial
  cofinal trajectory.\\[0.35em]
Compressed V15 source Gram and source edge
& \MGopen
& Need uniform carrier landing and a strict source-cyclic contraction.\\[0.35em]
Complete-source incidence and V13 routing
& \MGopen
& Populate \(\mathsf I,\mathsf Q,\mathsf E\) on the actual Perron law.\\[0.35em]
Kinetic density and signed diagonal Lipschitz stock
& \MGopen
& Needed for a fixed good-site threshold.\\[0.35em]
Protected sectors and nontrivial OS limit
& \MGopen
& Uniform sector cards and local-algebra noncollapse remain.\\[0.35em]
Full Yang--Mills mass gap
& \MGopen
& The cofinal physical stocks and continuum construction remain unproved.\\
\bottomrule
\end{longtable}
\endgroup

\begin{openproblem}[V21 mandate: response matching and complete-source incidence]
\label{op:V21}
\begin{enumerate}[label=\textup{(V21.\arabic*)},leftmargin=5.5em]
\item Compare the V35 reduced-fibre density and horizontal exact range with
      the V8 Perron law.  Prove the same-law incidence or retain an explicit
      density/range defect; do not apply
      \eqref{eq:V20-physical-triple} across a mismatch.
\item Derive \(\kappa_{H,j}\alpha_{\beta_j}\) from the literal anisotropic
      temporal and spatial Wilson coefficients after the V11 temporal
      limit.  Do not identify \(\alpha_\beta\) with a Hamiltonian magnetic
      coefficient without the transfer expansion.
\item Construct \(P_{\mathscr C_j^{\rm src}}\) from the complete signed V15
      sources and estimate the compressed Gram in
      \eqref{eq:V20-sharp-source-landing}.  Return a normalized carrier
      witness if it diverges.
\item Prove the complete-source incidence
      \eqref{eq:V20-source-incidence}, keeping the two V15 conditioning
      sigma algebras distinct until their common-history map is explicit.
\item Build the V13 physical flow only for the assembled source and populate
      \(\mathsf Q\) and the residual \(\mathsf E\) in
      \eqref{eq:V20-source-capacity}.
\item Combine the resulting bad-capacity row with the V16 source edge,
      response summability, V18 good-site floor, and every protected-sector
      card.  Invoke the V1 OS passage only after a common positive physical
      gap and nontrivial local limit are proved.
\end{enumerate}
\end{openproblem}

\begin{firewall}[End-of-V20 claim boundary]
V20 performs the required cross-program audit, resolves the physical
one-form scalar conditional on same-law incidence, proves a two-sided global
Wilson defect scale, and factors bad-set control through the complete signed
source.  It does not prove the same-law incidence, response-matched coupling
product, compressed source landing, source edge, capacity matrices,
protected-sector floor, or nontrivial continuum OS state.  Therefore the
full Yang--Mills mass gap remains open.
\end{firewall}

\section*{Version V20 append-only record}
\addcontentsline{toc}{section}{Version V20 append-only record}
The complete V19 source was retained before this continuation.  Its SHA--256
digest was
\begin{center}\scriptsize\ttfamily
d204d66efb33618d9d5dc61c074043f00596ed947767c86a318ca976a2a0bf10.
\end{center}
On the next iteration, retain this full file, append before its unique active
document-closing command, increment the filename to
\texttt{\_V21.tex}, and remove only the superseded V20 file from this note
series.

% =============================================================================
% END APPENDED VERSION V20 MATERIAL
% =============================================================================

% =============================================================================
% BEGIN APPENDED VERSION V21 MATERIAL
% =============================================================================

\clearpage
\part{Version V21: Perron-law mismatch, density-defect Witten repair, and anisotropic coefficient matching}
\label{part:V21}

\section{Direction}
\label{sec:V21-direction}

V20 made the physical scalar exact conditional on one incidence:
the bare V35 Witten split must represent the V8 Witten operator of the
actual Perron ground law.  V21 tests that incidence rather than assuming it.

The test returns a negative answer for the naive local Gibbs candidate.
For a nonconstant magnetic potential and an injective temporal Wilson
convolution, the Perron law of the symmetric transfer is not the one-slice
law proportional to the exponential of that potential.  This is an exact
finite-regulator obstruction, not an asymptotic concern.  The correction is
the logarithmic density defect between the local V35 law and the physical
Perron law.  Its Hessian produces an exact positive
comparison-minus-defect split on the actual law.

V21 also sharpens the V10 temporal calibration to
\(\alpha s_\alpha\to3/2\) and derives the response-matched Wilson product
when the Hamiltonian magnetic coefficient is explicitly loaded into the
symmetric transfer.  These results decide which part of the V20 global
defect scale is fixed by temporal normalization and which part remains a
spatial-coupling acquisition.

% =============================================================================
\section{The local Gibbs law is not the Perron law}
\label{sec:V21-law-mismatch}
% =============================================================================

Let \(\mathcal X=G^E\) with product Haar measure \(\mathfrak m\).  Let
\(\mathcal C\) be the tensor product of the V10 temporal Wilson convolution
operators and let \(G,F\) be bounded real functions on \(\mathcal X\).
Here \(G\) is the magnetic exponent appearing in the two outer halves of a
symmetric transfer, while \(F\) is the exponent of a proposed local Gibbs
law.  Define
\begin{equation}
 \mathcal T_G
 :=M_{e^{-G/2}}\mathcal C M_{e^{-G/2}},
 \qquad
 \dd\nu_F:=Z_F^{-1}e^{-F}\dd\mathfrak m.
 \label{eq:V21-transfer-local-law}
\end{equation}
Let \(h>0\) be the normalized Perron eigenfunction,
\begin{equation}
 \mathcal T_Gh=\lambda h,
 \qquad
 \dd\pi:=h^2\dd\mathfrak m.
 \label{eq:V21-Perron-law}
\end{equation}

\begin{theorem}[Exact Perron/local-law incidence equation]
\label{thm:V21-incidence-equation}
The equality \(\nu_F=\pi\) holds if and only if
\begin{equation}
 \boxed{
 \mathcal C\!\left(e^{-(G+F)/2}\right)
 =\lambda e^{(G-F)/2}.}
 \label{eq:V21-incidence-equation}
\end{equation}
The equality is understood up to the harmless normalization of the Perron
eigenfunction, which cancels from both sides.
\end{theorem}

\begin{proof}
Equality of the two probability laws is equivalent to
\[
 h=c\,e^{-F/2}
\]
for a positive constant \(c\).  Substitute this expression into
\eqref{eq:V21-Perron-law}:
\[
 c\,e^{-G/2}
 \mathcal C\!\left(e^{-(G+F)/2}\right)
 =\lambda c\,e^{-F/2}.
\]
Cancel \(c e^{-G/2}\) to obtain
\eqref{eq:V21-incidence-equation}.  Every step is reversible.
\end{proof}

\begin{corollary}[No naive same-exponent Perron law]
\label{cor:V21-no-naive-same-law}
Assume \(\mathcal C1=1\) and \(\mathcal C\) is injective on
\(L^2(\mathfrak m)\).  If \(F=G\), then
\begin{equation}
 \boxed{
 \nu_G=\pi
 \quad\Longleftrightarrow\quad
 G\text{ is constant almost everywhere}.}
 \label{eq:V21-no-naive-same-law}
\end{equation}
For the V10 temporal Wilson convolution, the strict character multipliers
\(0<\eta_R(\alpha)\le1\) make every finite tensor power injective.
Consequently a nonconstant Wilson magnetic action cannot have
\(Z^{-1}e^{-G}\mathfrak m\) as the Perron law of the symmetric transfer.
\end{corollary}

\begin{proof}
With \(F=G\), \eqref{eq:V21-incidence-equation} becomes
\[
 \mathcal C(e^{-G})=\lambda1.
\]
Since \(\mathcal C1=1\),
\(\mathcal C(e^{-G}-\lambda)=0\).  Injectivity gives
\(e^{-G}=\lambda\), hence \(G\) is constant.  The converse is immediate.
For the V10 operator, Peter--Weyl diagonalization gives strictly positive
multipliers on every irreducible block.  Their finite products are also
strictly positive, so the tensor convolution is injective.
\end{proof}

\begin{assessment}[Decision of the V20 same-law test]
\label{ass:V21-same-law-decision}
The V35 reduced-fibre density
\[
 Z^{-1}e^{-\beta V_M-\Psi}\,\dd m_{\rm red}
\]
cannot be identified with the full Perron ground law merely because its
exponent matches a spatial Wilson factor in the transfer.  When it is
promoted to a law on the same complete configuration space and \(F=G\),
\Cref{cor:V21-no-naive-same-law} rules out equality for a nonconstant
action.  If it is used only as a conditional fibre law at fixed coarse
\(W\), its normalizing factor may depend on \(W\); the theorem does not rule
out that narrower identity, which requires a separate disintegrated
transfer calculation.  A valid use of the V35 curvature tensor must
therefore either prove that conditional incidence, or:
\begin{enumerate}[label=\textup{(\alph*)},leftmargin=3.5em]
\item transport estimates across the explicit density defect; or
\item rebuild the comparison split directly on the physical Perron law.
\end{enumerate}
V21 takes the second route below.
\end{assessment}

\begin{firewall}[A temporal limit does not retroactively make the laws equal]
\label{fire:V21-finite-law-mismatch}
V11 proves convergence of normalized transfer powers to a differential
Hamiltonian semigroup at fixed spatial regulator.  It does not say that the
finite-step Perron eigenfunction is \(e^{-G/2}\), nor that the limiting
Hamiltonian ground state is a Gibbs square root.  The exact obstruction
\eqref{eq:V21-no-naive-same-law} must be retained at every finite temporal
regulator; an asymptotic density comparison would require its own uniform
estimate.
\end{firewall}

% =============================================================================

% =============================================================================

% =============================================================================
\section{The exact Perron density defect}
\label{sec:V21-density-defect}
% =============================================================================

Define the physical ground-law exponent and the unnormalized density defect
by
\begin{equation}
 W_\pi:=-2\log h,
 \qquad
 \vartheta:=W_\pi-F.
 \label{eq:V21-density-defect}
\end{equation}
The logarithm of \(\dd\nu_F/\dd\pi\) differs from \(\vartheta\) only by the
constant \(-\log Z_F\).

\begin{proposition}[Nonlinear Perron equation for the density defect]
\label{prop:V21-density-defect-equation}
The function \(\vartheta\) satisfies
\begin{equation}
 \boxed{
 \mathcal C\!\left(
 e^{-(G+F+\vartheta)/2}
 \right)
 =
 \lambda e^{(G-F-\vartheta)/2}.}
 \label{eq:V21-density-defect-equation}
\end{equation}
Conversely, any real \(\vartheta\) satisfying this equation for which
\(e^{-(F+\vartheta)/2}\) is positive and square integrable produces a
Perron eigenfunction of \(\mathcal T_G\).
\end{proposition}

\begin{proof}
By definition,
\[
 h=e^{-W_\pi/2}=e^{-(F+\vartheta)/2}.
\]
Insert this expression into
\(\mathcal T_Gh=\lambda h\) and multiply the resulting equality by
\(e^{G/2}\).  This gives
\eqref{eq:V21-density-defect-equation}.  Reversing the calculation proves
the converse.
\end{proof}

Let
\[
 \mathcal J:=-\nabla\log h
\]
be the physical Perron score of V8.  Differentiating
\eqref{eq:V21-density-defect} gives the exact first and second defect cards
\begin{align}
 \nabla\vartheta
 &=2\mathcal J-\nabla F,
 \label{eq:V21-density-score-defect}\\
 \Hess\vartheta
 &=2\nabla\mathcal J-\Hess F.
 \label{eq:V21-density-Hessian-defect}
\end{align}
Using the V8 first-slice score--Hessian identity, the latter becomes
\begin{equation}
 \boxed{
 \Hess\vartheta
 =
 2\mathbb E_{p_x}K_x
 -2\Cov_{p_x}(J_x)
 -\Hess F.}
 \label{eq:V21-density-defect-V8}
\end{equation}
Every term in \eqref{eq:V21-density-defect-V8} is on the actual transfer
and its actual terminal-weighted conditional, except the explicitly
declared local tensor \(\Hess F\).

% =============================================================================
\section{An exact actual-law repair of the V35 split}
\label{sec:V21-actual-law-split}
% =============================================================================

Assume the pointwise V35 curvature identity is loaded on the same
horizontal cotangent bundle:
\begin{equation}
 \Ric+\Hess F
 =\mathcal R_\eta-\Theta_\eta,
 \qquad
 \mathcal R_\eta\succeq\rho_{\beta,\eta}I,
 \qquad
 \Theta_\eta\succeq0.
 \label{eq:V21-V35-pointwise-split}
\end{equation}
This is a tensor identity and does not yet identify the V35 law with
\(\pi\).  Put
\[
 S_\vartheta:=\Hess\vartheta,
 \qquad
 S_\vartheta=S_{\vartheta,+}-S_{\vartheta,-}
\]
for its pointwise positive and negative parts.

\begin{theorem}[Perron-law comparison-minus-defect repair]
\label{thm:V21-actual-law-split}
On \(L^2(\pi;T^*\mathcal X)\), define
\begin{align}
 \widehat{\mathcal R}_{\eta,\pi}
 &:=\mathcal R_\eta+S_{\vartheta,+},
 \label{eq:V21-repaired-R}\\
 \widehat\Theta_{\eta,\pi}
 &:=\Theta_\eta+S_{\vartheta,-},
 \label{eq:V21-repaired-Theta}\\
 \widehat{\mathcal A}_{\eta,\pi}
 &:=\nabla_\pi^*\nabla+\widehat{\mathcal R}_{\eta,\pi}.
 \label{eq:V21-repaired-A}
\end{align}
Then
\begin{equation}
 \boxed{
 \kappa_H^{-1}\mathscr L_{1,\pi}^{\rm phys}
 =
 \widehat{\mathcal A}_{\eta,\pi}
 -\widehat\Theta_{\eta,\pi}.}
 \label{eq:V21-repaired-Witten-split}
\end{equation}
Both terms use the actual Perron law, and
\begin{equation}
 \boxed{
 \widehat{\mathcal A}_{\eta,\pi}
 \succeq\rho_{\beta,\eta}I,
 \qquad
 \widehat\Theta_{\eta,\pi}\succeq0.}
 \label{eq:V21-repaired-floor}
\end{equation}
Thus the V35 positive comparison floor survives the law repair.  The exact
new debit is the negative part of the density-defect Hessian.
\end{theorem}

\begin{proof}
Since \(W_\pi=F+\vartheta\),
\[
 \Ric+\Hess W_\pi
 =\Ric+\Hess F+S_\vartheta.
\]
Use \eqref{eq:V21-V35-pointwise-split} and
\(S_\vartheta=S_{\vartheta,+}-S_{\vartheta,-}\):
\[
 \Ric+\Hess W_\pi
 =
 \widehat{\mathcal R}_{\eta,\pi}
 -\widehat\Theta_{\eta,\pi}.
\]
The V8 weighted Bochner identity on the physical law is
\[
 \kappa_H^{-1}\mathscr L_{1,\pi}^{\rm phys}
 =\nabla_\pi^*\nabla+\Ric+\Hess W_\pi.
\]
Substitution proves \eqref{eq:V21-repaired-Witten-split}.  The weighted
rough Laplacian is nonnegative,
\(S_{\vartheta,+}\succeq0\), and
\(\mathcal R_\eta\succeq\rho_{\beta,\eta}I\), proving the first floor.
The second is immediate from the sum of two positive operators.
\end{proof}

\begin{corollary}[Exact repaired Woodbury channel]
\label{cor:V21-repaired-Woodbury}
Choose
\[
 \Theta_\eta=\mathcal V_\eta\mathcal V_\eta^*
\]
and define the direct-sum synthesis
\begin{equation}
 \widehat{\mathcal V}_{\eta,\pi}
 :=
 \begin{bmatrix}
  \mathcal V_\eta&S_{\vartheta,-}^{1/2}
 \end{bmatrix}.
 \label{eq:V21-repaired-synthesis}
\end{equation}
Then
\begin{equation}
 \widehat{\mathcal V}_{\eta,\pi}
 \widehat{\mathcal V}_{\eta,\pi}^*
 =\widehat\Theta_{\eta,\pi}.
 \label{eq:V21-repaired-factorization}
\end{equation}
Whenever the finite physical Witten operator is strictly positive on the
declared exact range, the V16--V19 Woodbury algebra applies on the actual
law with
\begin{equation}
 \boxed{
 \widehat{\mathcal K}_{\eta,\pi}
 :=
 \widehat{\mathcal V}_{\eta,\pi}^*
 \widehat{\mathcal A}_{\eta,\pi}^{-1}
 \widehat{\mathcal V}_{\eta,\pi}.}
 \label{eq:V21-repaired-K}
\end{equation}
\end{corollary}

\begin{proof}
Multiplication of the block row in
\eqref{eq:V21-repaired-synthesis} by its adjoint gives
\(\Theta_\eta+S_{\vartheta,-}\).  The remaining statement is the Woodbury
identity applied to \eqref{eq:V21-repaired-Witten-split}; strict positivity
gives the finite inverse and \(0\preceq\widehat{\mathcal K}_{\eta,\pi}<I\)
on the active channel.
\end{proof}

\begin{corollary}[Complete-channel physical defect criterion after law repair]
\label{cor:V21-repaired-defect-criterion}
For the pure-Wilson part on a cofinal family satisfying the V20 local
nondegeneracy,
\begin{equation}
 \boxed{
 \sup_j
 \norm{\sqrt{\kappa_{H,j}}\,
       \widehat{\mathcal V}_{\eta,\pi_j}}<\infty}
 \label{eq:V21-repaired-global-bound}
\end{equation}
holds if and only if both
\begin{equation}
 \boxed{
 \sup_j\kappa_{H,j}\alpha_{\beta_j}<\infty,
 \qquad
 \sup_j\kappa_{H,j}
       \norm{S_{\vartheta_j,-}}<\infty.}
 \label{eq:V21-two-repaired-stocks}
\end{equation}
\end{corollary}

\begin{proof}
By \eqref{eq:V21-repaired-factorization},
\[
 \norm{\widehat{\mathcal V}_{\eta,\pi}}^2
 =\norm{\Theta_\eta+S_{\vartheta,-}}.
\]
For positive operators \(P,Q\),
\[
 \max\{\norm P,\norm Q\}
 \le\norm{P+Q}
 \le\norm P+\norm Q.
\]
The V20 two-sided Wilson estimate makes
\(\kappa_{H,j}\norm{\Theta_{\eta,j}^{\rm W}}\) uniformly bounded exactly
when \(\kappa_{H,j}\alpha_{\beta_j}\) is.  Apply the displayed comparison
with \(P=\Theta_{\eta,j}^{\rm W}\) and
\(Q=S_{\vartheta_j,-}\).
\end{proof}

\begin{assessment}[The corrected source-carrier target]
\label{ass:V21-corrected-carrier}
When \(\nu_F\ne\pi\), the V20 compressed Gram must be formed from
\(\widehat{\mathcal V}_{\eta,\pi}\), not from the bare V35 synthesis alone:
\begin{equation}
 \boxed{
 \sup_j\kappa_{H,j}
 \norm{
 P_{\widehat{\mathscr C}_j^{\rm src}}
 \widehat{\mathcal V}_{\eta,\pi_j}^*
 \widehat{\mathcal V}_{\eta,\pi_j}
 P_{\widehat{\mathscr C}_j^{\rm src}}}<\infty.}
 \label{eq:V21-corrected-source-Gram}
\end{equation}
The enlarged source carrier must include the negative density-defect
channel \(S_{\vartheta,-}^{1/2}\).  Equation
\eqref{eq:V21-density-defect-V8} identifies that channel with the difference
between the complete signed V8/V15 Perron curvature and the submitted local
Hessian.  It is not an unspecified second interaction.
\end{assessment}

\begin{firewall}[The repair preserves positivity but may lose locality]
\label{fire:V21-repair-locality}
The split \eqref{eq:V21-repaired-Witten-split} is exact and retains the V35
comparison floor, but \(\vartheta\) depends on the global Perron
eigenfunction.  Therefore \(S_{\vartheta,\pm}\) need not inherit the literal
finite propagation of \(\Theta_\eta\).  Exponential response locality must
be proved from the V15 source-specific Green columns or another physical
Perron estimate; it cannot be assigned from the local Gibbs action.
\end{firewall}

% =============================================================================

% =============================================================================

% =============================================================================
\section{The density-defect channel is genuinely present}
\label{sec:V21-density-defect-nonzero}
% =============================================================================

\begin{lemma}[Nonconstant density defect has a negative Hessian direction]
\label{lem:V21-negative-Hessian}
Let \(\mathcal X\) be compact and connected and let
\(\vartheta\in C^2(\mathcal X)\) be nonconstant.  Then
\begin{equation}
 \boxed{S_{\vartheta,-}\not\equiv0.}
 \label{eq:V21-negative-Hessian}
\end{equation}
\end{lemma}

\begin{proof}
If \(S_{\vartheta,-}\equiv0\), then
\(\Hess\vartheta\succeq0\) pointwise.  Taking the trace gives
\(\Delta\vartheta\ge0\).  On a compact manifold without boundary,
\(\int\Delta\vartheta=0\), hence \(\Delta\vartheta=0\).  A positive
semidefinite tensor with zero trace is zero, so
\(\Hess\vartheta=0\).  Equivalently \(\vartheta\) is harmonic; compact
connectedness makes it constant, a contradiction.
\end{proof}

\begin{corollary}[Naive law mismatch creates a nonzero repair defect]
\label{cor:V21-nonzero-repair-defect}
Under the hypotheses of \Cref{cor:V21-no-naive-same-law}, if \(F=G\) is
nonconstant, then the density defect \(\vartheta\) is nonconstant and
\begin{equation}
 \boxed{
 \widehat\Theta_{\eta,\pi}-\Theta_\eta
 =S_{\vartheta,-}\ne0.}
 \label{eq:V21-nonzero-repair-defect}
\end{equation}
\end{corollary}

\begin{proof}
If \(\vartheta\) were constant, then
\(W_\pi=F+\mathrm{constant}\), so
\(\pi=Z_F^{-1}e^{-F}\mathfrak m=\nu_F\), contrary to
\Cref{cor:V21-no-naive-same-law}.  Apply
\Cref{lem:V21-negative-Hessian} and
\eqref{eq:V21-repaired-Theta}.
\end{proof}

% =============================================================================
\section{Sharp temporal calibration and the matched spatial coefficient}
\label{sec:V21-coefficient-matching}
% =============================================================================

In this section \(\alpha\) is the V10 temporal Wilson concentration and
\(\beta\) is the V35 spatial/reduced-fibre Wilson coefficient.

\begin{theorem}[Leading V10 heat-time constant]
\label{thm:V21-temporal-leading-constant}
For the V10 law
\[
 \dd\nu_\alpha(U)
 =Z_\alpha^{-1}
 \exp\!\left\{\frac{\alpha}{3}\Re\Tr U\right\}\dd U
\]
on \(\mathrm{SU}(3)\), its fundamental calibration satisfies
\begin{equation}
 \boxed{
 s_\alpha
 =\frac{3}{2\alpha}+O(\alpha^{-2}),
 \qquad
 \lim_{\alpha\to\infty}\alpha s_\alpha=\frac32.}
 \label{eq:V21-temporal-leading-constant}
\end{equation}
More generally, for each fixed irreducible representation \(R\),
\begin{equation}
 \boxed{
 \eta_R(\alpha)
 =1-\frac{3C_2(R)}{2\alpha}+O_R(\alpha^{-2}).}
 \label{eq:V21-character-leading-constant}
\end{equation}
\end{theorem}

\begin{proof}
Use exponential normal coordinates \(U=e^X\) near the identity with
\(\norm{X}^2=-\Tr(X^2)\).  Since \(X\) is skew Hermitian and traceless,
\[
 \frac13\Re\Tr(e^X)
 =1-\frac16\norm{X}^2+O(\norm{X}^4).
\]
The Haar Jacobian is \(1+O(\norm{X}^2)\).  The identity is the unique
maximum of the exponent, so the complement of a fixed normal
neighborhood is exponentially small.  After the rescaling
\(X=\alpha^{-1/2}Y\), dominated convergence against the Gaussian density
proportional to \(e^{-\norm Y^2/6}\) gives
\begin{equation}
 \alpha\,\mathbb E_{\nu_\alpha}\norm{X}^2
 =3d_G+O(\alpha^{-1}),
 \qquad d_G=8.
 \label{eq:V21-temporal-second-moment}
\end{equation}
The normalized character expansion from V10 is
\[
 \frac{\chi_R(e^X)}{d_R}
 =1-\frac{C_2(R)}{2d_G}\norm{X}^2+O_R(\norm{X}^4).
\]
Inversion symmetry removes odd terms, while the rescaled fourth moment is
\(O(\alpha^{-2})\).  Insert
\eqref{eq:V21-temporal-second-moment} to obtain
\eqref{eq:V21-character-leading-constant}.  For the fundamental
representation,
\[
 s_\alpha
 =-\frac{\log\eta_F(\alpha)}{C_F}.
\]
Since \(-\log(1-u)=u+O(u^2)\), the factor \(C_F\) cancels and gives
\eqref{eq:V21-temporal-leading-constant}.
\end{proof}

\begin{definition}[Literal Wilson coefficient matching]
\label{def:V21-Wilson-coefficient-matching}
At spatial regulator \(j\), write the Hamiltonian magnetic multiplication
operator as
\begin{equation}
 V_{B,\alpha,j}
 =\gamma_{\alpha,j}V_{M,j}
  +V_{B,\alpha,j}^{\perp},
 \label{eq:V21-magnetic-coefficient}
\end{equation}
where \(V_{M,j}\) is the unit Wilson action used by V35 and the second term
contains no copy of that declared Wilson coefficient.  The V35 Wilson
exponent is coefficient matched to the symmetric transfer when
\begin{equation}
 \boxed{
 \beta_{\alpha,j}
 =\tau_\alpha\gamma_{\alpha,j}.}
 \label{eq:V21-beta-matching}
\end{equation}
This matches one coefficient; it does not assert equality of the local and
Perron laws.
\end{definition}

\begin{theorem}[Response-matched physical Wilson product]
\label{thm:V21-matched-product}
Under \eqref{eq:V21-beta-matching} and the V10 calibration
\(\kappa_{H,j}=s_\alpha/\tau_\alpha\),
\begin{equation}
 \boxed{
 \kappa_{H,j}\alpha_{\beta_{\alpha,j}}
 =\frac{s_\alpha\gamma_{\alpha,j}}{3}.}
 \label{eq:V21-exact-matched-product}
\end{equation}
Consequently,
\begin{equation}
 \boxed{
 \kappa_{H,j}\alpha_{\beta_{\alpha,j}}
 =
 \frac{\gamma_{\alpha,j}}{2\alpha}
 +O\!\left(\frac{\gamma_{\alpha,j}}{\alpha^2}\right).}
 \label{eq:V21-asymptotic-matched-product}
\end{equation}
In particular:
\begin{enumerate}[label=\textup{(\roman*)},leftmargin=3.6em]
\item if \(\gamma_{\alpha,j}/\alpha\to\ell_j\), the product tends to
      \(\ell_j/2\);
\item a uniform bound on \(\gamma_{\alpha,j}/\alpha\) gives a uniform bound
      on the matched pure-Wilson part of the physical defect;
\item if \(\gamma_{\alpha,j}/\alpha\to\infty\), the matched pure-Wilson
      complete-channel synthesis diverges by the V20 centre witness.
\end{enumerate}
\end{theorem}

\begin{proof}
Since \(\alpha_\beta=\beta/3\),
\[
 \kappa_{H,j}\alpha_{\beta_{\alpha,j}}
 =\frac{s_\alpha}{\tau_\alpha}
  \frac{\tau_\alpha\gamma_{\alpha,j}}{3}
 =\frac{s_\alpha\gamma_{\alpha,j}}3.
\]
Insert \eqref{eq:V21-temporal-leading-constant} to obtain
\eqref{eq:V21-asymptotic-matched-product}.  The three alternatives follow
from \(\alpha s_\alpha\to3/2\) and
\Cref{thm:V20-two-sided-defect}.
\end{proof}

\begin{assessment}[What coefficient matching closes]
\label{ass:V21-coefficient-conclusion}
The temporal spacing cancels exactly from the physical Wilson defect scale.
The remaining scalar is the ratio
\(\gamma_{\alpha,j}/\alpha\) between the Hamiltonian spatial Wilson
coefficient and the temporal Wilson concentration.  None of the supplied
files specifies a cofinal numerical law for this ratio across the spatial
continuum trajectory.  Therefore V21 reduces the response-matching row to
one literal coefficient ratio but does not assign its value.

Even a bounded ratio controls only the submitted Wilson part.  The actual
Perron-law split additionally requires
\(\kappa_H\norm{S_{\vartheta,-}}\), the corrected source Gram
\eqref{eq:V21-corrected-source-Gram}, and the strict repaired edge
\(\norm{\widehat{\mathcal K}_{\eta,\pi}}<1\).
\end{assessment}

% =============================================================================
\section{V21 ledger and revised frontier}
\label{sec:V21-ledger}
% =============================================================================

\begin{theorem}[Integrated V21 statement]
\label{thm:V21-integrated}
Preserving V1--V20, V21 proves:
\begin{enumerate}[label=\textup{(V21.\arabic*)},leftmargin=5.5em]
\item the exact nonlinear equation characterizing equality of a proposed
      local Gibbs law with the symmetric-transfer Perron law;
\item injectivity of the V10 convolution rules out the naive equality for
      every nonconstant matched magnetic exponent;
\item an exact nonlinear equation and signed Hessian card for the Perron
      density defect \(\vartheta\);
\item a positive comparison-minus-defect repair on the actual Perron law,
      retaining the V35 comparison floor and adding precisely
      \(S_{\vartheta,-}\) to the defect;
\item proof that this added negative-Hessian channel is nonzero whenever the
      naive local and Perron laws differ;
\item the sharp temporal calibration
      \(s_\alpha=3/(2\alpha)+O(\alpha^{-2})\);
\item the exact matched product
      \(\kappa_H\alpha_\beta=s_\alpha\gamma/3\) and its asymptotic
      \(\gamma/(2\alpha)\);
\item replacement of the bare V20 source-carrier Gram by the repaired
      physical Gram \eqref{eq:V21-corrected-source-Gram}.
\end{enumerate}
\end{theorem}

\begin{proof}
Items \textup{(V21.1)--(V21.2)} are
\Cref{thm:V21-incidence-equation,cor:V21-no-naive-same-law}.  Item
\textup{(V21.3)} is
\Cref{prop:V21-density-defect-equation} and
\eqref{eq:V21-density-defect-V8}.  Item \textup{(V21.4)} is
\Cref{thm:V21-actual-law-split,cor:V21-repaired-Woodbury}.  Item
\textup{(V21.5)} is
\Cref{lem:V21-negative-Hessian,cor:V21-nonzero-repair-defect}.  Items
\textup{(V21.6)--(V21.7)} are
\Cref{thm:V21-temporal-leading-constant,thm:V21-matched-product}.  Item
\textup{(V21.8)} is \Cref{ass:V21-corrected-carrier}.
\end{proof}

\begingroup
\small
\begin{longtable}{>{\raggedright\arraybackslash}p{0.28\textwidth}
                  >{\raggedright\arraybackslash}p{0.16\textwidth}
                  >{\raggedright\arraybackslash}p{0.48\textwidth}}
\toprule
\textbf{V21 row}&\textbf{Status}&\textbf{Advance or remaining obstruction}\\
\midrule
\endfirsthead
\toprule
\textbf{V21 row}&\textbf{Status}&\textbf{Advance or remaining obstruction}\\
\midrule
\endhead
Naive full Gibbs/Perron same-law identity
& \MGfail
& Ruled out for a nonconstant matched full action by convolution injectivity.\\[0.35em]
Actual-law Witten repair
& \MGpass
& Exact split with unchanged comparison floor and new defect
  \(S_{\vartheta,-}\).\\[0.35em]
Negative density-defect channel
& \MGpass
& Nonzero at every finite regulator on the mismatched branch.\\[0.35em]
Temporal leading constant
& \MGpass
& \(s_\alpha=3/(2\alpha)+O(\alpha^{-2})\).\\[0.35em]
Matched Wilson product
& \MGpass
& Reduced to \(s_\alpha\gamma/3\sim\gamma/(2\alpha)\).\\[0.35em]
Cofinal spatial coefficient ratio
& \MGopen
& The supplied files do not fix \(\gamma_{\alpha,j}/\alpha\).\\[0.35em]
Negative Hessian magnitude and locality
& \MGopen
& Need a uniform source-specific estimate for
  \(\kappa_HS_{\vartheta,-}\).\\[0.35em]
Repaired source Gram and edge
& \MGopen
& Must use \(\widehat{\mathcal V}_{\eta,\pi}\) on the actual source carrier.\\[0.35em]
Complete-source incidence and bad capacity
& \MGopen
& Populate the V20 matrices for the repaired physical source.\\[0.35em]
Protected sectors and nontrivial OS limit
& \MGopen
& Uniform sector cards and local-algebra noncollapse remain.\\[0.35em]
Full Yang--Mills mass gap
& \MGopen
& The repaired cofinal edge, capacity, and continuum construction remain.\\
\bottomrule
\end{longtable}
\endgroup

\begin{openproblem}[V22 mandate: estimate the Perron density defect]
\label{op:V22}
\begin{enumerate}[label=\textup{(V22.\arabic*)},leftmargin=5.5em]
\item Use the nonlinear Perron equation
      \eqref{eq:V21-density-defect-equation} or the V8/V15 score identities
      to estimate \(S_{\vartheta,-}\) on the actual future conditional law.
      Preserve the direct-Hessian/covariance cancellation in
      \eqref{eq:V21-density-defect-V8}.
\item Load the literal spatial coefficient
      \(\gamma_{\alpha,j}\) from the response-matched anisotropic action.
      If it is not supplied, keep \(\gamma_{\alpha,j}/\alpha\) as an
      explicit cofinal input rather than assigning a continuum scaling.
\item Factor the repaired defect
      \(\Theta_\eta+S_{\vartheta,-}\) on the complete signed V15 source
      carrier and estimate \eqref{eq:V21-corrected-source-Gram}.
\item Prove a strict source-cyclic edge for
      \(\widehat{\mathcal K}_{\eta,\pi}\), or return a normalized repaired
      channel vector and apply the V18--V20 landing classifiers.
\item Rebuild the V20 complete-source incidence matrices with the
      density-defect channel included, then construct the V13 capacity
      routing on the actual Perron law.
\item Combine the repaired edge and capacity with the V18 good-site floor,
      protected sectors, and the V1 nontrivial OS passage only after every
      common physical constant is uniform.
\end{enumerate}
\end{openproblem}

\begin{firewall}[End-of-V21 claim boundary]
V21 disproves the naive local-Gibbs/Perron identification and replaces it
with an exact actual-law Witten split.  It also fixes the temporal leading
constant and reduces the Wilson normalization to one spatial/temporal
coefficient ratio.  It does not bound the density-defect Hessian, prove the
repaired source edge or bad capacity, control all protected sectors, or
construct a nontrivial continuum OS state.  Therefore the full Yang--Mills
mass gap remains open.
\end{firewall}

\section*{Version V21 append-only record}
\addcontentsline{toc}{section}{Version V21 append-only record}
The complete V20 source was retained before this continuation.  Its SHA--256
digest was
\begin{center}\scriptsize\ttfamily
6c7f7fa80292e423371e783d59d8f8b8234eff604422930a6126eb98acc06671.
\end{center}
On the next iteration, retain this full file, append before its unique active
document-closing command, increment the filename to
\texttt{\_V22.tex}, and remove only the superseded V21 file from this note
series.

% =============================================================================
% END APPENDED VERSION V21 MATERIAL
% =============================================================================

% =============================================================================
% BEGIN APPENDED VERSION V22 MATERIAL
% =============================================================================

\clearpage
\part{Version V22: temporal posterior excess, density-defect estimates, and the strong-coupling bridge}
\label{part:V22}

\section{Direction}
\label{sec:V22-direction}

V21 showed that the submitted local Gibbs exponent is not, in general, the
logarithm of the actual Perron ground law.  Its exact repair introduced the
additional negative-Hessian channel (S_{\vartheta,-}), but left that
channel in implicit form.  V22 differentiates the symmetric Perron equation
itself.  On the matched branch (F=G), the new channel is exactly twice the
positive part of a temporal posterior covariance--curvature excess.  This
identification separates three logically different questions:
\begin{enumerate}[label=\textup{(\roman*)},leftmargin=3.5em]
\item what the added source is;
\item how large it is on a fixed regulator bank; and
\item whether its square-root factor has the source locality required by the
      V15--V20 closure mechanism.
\end{enumerate}

The raw posterior admits a rigorous chi-square perturbation bound, but a
fixed-smooth-tilt calculation proves that this bound loses a factor of order
\(\sqrt\alpha\) relative to the exact temporal cancellation.  Going
upstream avoids that loss on the V9 strong-coupling branch: the already
proved half-space Perron-polymer Hessian estimate yields a volume-uniform
bound of order \(|\beta|\) for the density defect, provided the V9 transfer
is explicitly identified with the complete V21 symmetric transfer.  This
closes the magnitude row on that branch, but it does not manufacture the
missing complete-source incidence or the strict repaired edge.

% =============================================================================
\section{The temporal posterior and its exact excess tensor}
\label{sec:V22-posterior-excess}
% =============================================================================

Retain the V21 finite configuration manifold
\(\mathcal X=\mathrm{SU}(3)^E\), product Haar probability
\(\mathfrak m\), and the tensor temporal Wilson convolution
\(\mathcal C_\alpha\).  Write its strictly positive symmetric kernel as
\begin{equation}
 (\mathcal C_\alpha u)(x)
 =\int_{\mathcal X}c_\alpha(x,y)u(y)\,\dd\mathfrak m(y),
 \qquad
 \int c_\alpha(x,y)\,\dd\mathfrak m(y)=1.
 \label{eq:V22-convolution-kernel}
\end{equation}
For the full matched exponent \(F=G\), let
\begin{equation}
 \mathcal T_G
 =M_{e^{-G/2}}\mathcal C_\alpha M_{e^{-G/2}},
 \qquad
 \mathcal T_Gh=\lambda h,
 \qquad h>0,
 \label{eq:V22-symmetric-transfer}
\end{equation}
and define
\begin{equation}
 f:=e^{-G/2}h,
 \qquad
 r_\alpha:=\mathcal C_\alpha f.
 \label{eq:V22-f-r}
\end{equation}
The untilted temporal row law and its Perron-tilted posterior are
\begin{align}
 \dd q_{\alpha,x}(y)
 &:=c_\alpha(x,y)\,\dd\mathfrak m(y),
 \label{eq:V22-q-row}\\
 \dd p_{\alpha,x}(y)
 &:=\frac{f(y)}{r_\alpha(x)}\,\dd q_{\alpha,x}(y).
 \label{eq:V22-posterior-row}
\end{align}
All expectations and covariances below are tensor-valued at the initial
point \(x\).  Put
\begin{equation}
 J^{\rm t}_x(y):=-\nabla_x\log c_\alpha(x,y),
 \qquad
 K^{\rm t}_x(y):=\Hess_x[-\log c_\alpha(x,y)].
 \label{eq:V22-temporal-score-curvature}
\end{equation}

\begin{theorem}[Exact temporal posterior-excess identity]
\label{thm:V22-posterior-excess}
With the definitions above, the V21 density defect
\(\vartheta=W_\pi-G\) satisfies
\begin{align}
 \lambda h(x)&=e^{-G(x)/2}r_\alpha(x),
 \label{eq:V22-Perron-factorization}\\
 \vartheta(x)&=-2\log r_\alpha(x)+2\log\lambda,
 \label{eq:V22-vartheta-log-r}\\
 \nabla\vartheta(x)
 &=2\mathbb E_{p_{\alpha,x}}J^{\rm t}_x,
 \label{eq:V22-vartheta-score}\\
 S_\vartheta(x):=\Hess\vartheta(x)
 &=2\mathbb E_{p_{\alpha,x}}K^{\rm t}_x
   -2\Cov_{p_{\alpha,x}}(J^{\rm t}_x).
 \label{eq:V22-vartheta-Hessian}
\end{align}
Consequently, if
\begin{equation}
 D_\alpha(x)
 :=\Cov_{p_{\alpha,x}}(J^{\rm t}_x)
   -\mathbb E_{p_{\alpha,x}}K^{\rm t}_x,
 \label{eq:V22-D-definition}
\end{equation}
then
\begin{equation}
 \boxed{
 D_\alpha=\Hess\log r_\alpha,
 \qquad
 S_\vartheta=-2D_\alpha,
 \qquad
 S_{\vartheta,-}=2(D_\alpha)_+.}
 \label{eq:V22-exact-defect-channel}
\end{equation}
Here \((A)_+\) denotes the positive spectral part of the self-adjoint
endomorphism \(A\).
\end{theorem}

\begin{proof}
The Perron equation \eqref{eq:V22-symmetric-transfer} is exactly
\[
 \lambda h(x)
 =e^{-G(x)/2}
   \int c_\alpha(x,y)e^{-G(y)/2}h(y)\,\dd\mathfrak m(y),
\]
which gives \eqref{eq:V22-Perron-factorization}.  Since
\(W_\pi=-2\log h\), taking logarithms gives
\eqref{eq:V22-vartheta-log-r}.

Compactness, strict positivity, and smoothness of the finite-regulator
kernel justify differentiation under the integral.  From
\(\nabla_xc_\alpha=-c_\alpha J^{\rm t}_x\),
\[
 \nabla\log r_\alpha
 =-\mathbb E_{p_{\alpha,x}}J^{\rm t}_x.
\]
For every tangent vector \(v\), differentiating once more gives
\[
 \Hess\log r_\alpha[v,v]
 =-\mathbb E_p K^{\rm t}[v,v]
   +\Var_p\!\left(J^{\rm t}[v]\right).
\]
Polarization yields the tensor identity
\[
 \Hess\log r_\alpha
 =-\mathbb E_pK^{\rm t}+\Cov_p(J^{\rm t}).
\]
Equations \eqref{eq:V22-vartheta-score}--
\eqref{eq:V22-vartheta-Hessian} follow from
\eqref{eq:V22-vartheta-log-r}.  Finally, for a self-adjoint \(D\), the
negative part of \(-2D\) is \(2D_+\).
\end{proof}

\begin{proposition}[Reconciliation with the full V8 score identity]
\label{prop:V22-V8-reconciliation}
Let \(J_x,K_x\) be the V8 score and direct Hessian of the complete symmetric
kernel.  Then
\begin{equation}
 J_x=\frac12\nabla G(x)+J^{\rm t}_x,
 \qquad
 K_x=\frac12\Hess G(x)+K^{\rm t}_x.
 \label{eq:V22-full-vs-temporal-score}
\end{equation}
Substitution into the V21 identity
\[
 S_\vartheta=2\mathbb E_pK_x-2\Cov_p(J_x)-\Hess G
\]
cancels the two copies of \(\Hess G\) and returns
\eqref{eq:V22-vartheta-Hessian}.
\end{proposition}

\begin{proof}
The negative logarithm of the complete kernel is
\[
 \frac12G(x)-\log c_\alpha(x,y)+\frac12G(y).
\]
Differentiate in \(x\).  The \(x\)-dependent vector
\(\nabla G(x)/2\) is deterministic under \(p_{\alpha,x}\), so adding it
does not change a covariance.  The stated cancellation is immediate.
\end{proof}

\begin{corollary}[Literal new synthesis column]
\label{cor:V22-density-synthesis-column}
The density-defect column in the repaired V21 synthesis may be chosen as
\begin{equation}
 \boxed{
 \mathcal V_{\rm den}(x)
 :=S_{\vartheta,-}(x)^{1/2}
 =\sqrt2\,(D_\alpha(x)_+)^{1/2}.}
 \label{eq:V22-density-synthesis-column}
\end{equation}
Thus the corrected synthesis and Gram are
\begin{align}
 \widehat{\mathcal V}_{\eta,\pi}
 &=\bigl[\,\mathcal V_\eta\quad \mathcal V_{\rm den}\,\bigr],
 \label{eq:V22-repaired-synthesis}\\
 \widehat{\mathcal V}_{\eta,\pi}^{*}
 \widehat{\mathcal V}_{\eta,\pi}
 &=\Theta_\eta+2(D_\alpha)_+.
 \label{eq:V22-repaired-Gram}
\end{align}
No independent ``density interaction'' may be inserted in place of
\eqref{eq:V22-density-synthesis-column}.
\end{corollary}

\begin{firewall}[Full law versus a conditional fibre]
\label{fire:V22-fibre-incidence}
The theorem is an identity on the complete configuration space on which
\(\mathcal C_\alpha\) is a normalized convolution row.  It applies verbatim
to a reduced fibre at fixed coarse variable \(W\) only if the disintegrated
temporal row is normalized on that fibre and all displayed derivatives are
fibre derivatives.  If the row normalizer depends on \(W\), its first and
second \(W\)-derivatives are additional incidence terms.  They must be
retained in the V20 residual matrix rather than silently absorbed into
\(D_\alpha\).
\end{firewall}

% =============================================================================
\section{The exact cancellation and two signed sum rules}
\label{sec:V22-cancellation-sum-rules}
% =============================================================================

\begin{lemma}[Untilting cancels covariance against direct curvature]
\label{lem:V22-untilted-cancellation}
If the terminal tilt is constant, \(f\equiv c>0\), then
\(r_\alpha\equiv c\), \(p_{\alpha,x}=q_{\alpha,x}\), and
\begin{equation}
 \boxed{
 \mathbb E_{q_{\alpha,x}}J^{\rm t}_x=0,
 \qquad
 \Cov_{q_{\alpha,x}}(J^{\rm t}_x)
 =\mathbb E_{q_{\alpha,x}}K^{\rm t}_x,
 \qquad D_\alpha=0.}
 \label{eq:V22-untilted-cancellation}
\end{equation}
\end{lemma}

\begin{proof}
The row normalization in \eqref{eq:V22-convolution-kernel} says that
\(r_\alpha\) is constant when \(f\) is.  Differentiate
\(\log r_\alpha\) once and twice and use the formulas proved in
\Cref{thm:V22-posterior-excess}.
\end{proof}

The cancellation is exact at every \(\alpha\).  Hence an estimate that
bounds the covariance and the direct curvature separately, each at its
natural order \(\alpha\), discards the structure needed for the density
channel.

Define the symmetric stationary edge law
\begin{equation}
 \dd\Gamma_\alpha(x,y)
 :=\frac1\lambda
 h(x)e^{-G(x)/2}c_\alpha(x,y)e^{-G(y)/2}h(y)
 \,\dd\mathfrak m(x)\dd\mathfrak m(y).
 \label{eq:V22-edge-law}
\end{equation}
It is a probability law, is invariant under \((x,y)\mapsto(y,x)\), and
disintegrates as
\begin{equation}
 \dd\Gamma_\alpha(x,y)=\dd\pi(x)\dd p_{\alpha,x}(y).
 \label{eq:V22-edge-disintegration}
\end{equation}

\begin{theorem}[Posterior Fisher and weighted integration-by-parts sum rules]
\label{thm:V22-signed-sum-rules}
Let traces and norms be taken in \(T_x\mathcal X\).  Then
\begin{align}
 \int_{\mathcal X}\Tr D_\alpha\,\dd\pi
 &=\int_{\mathcal X^2}
   \bigl(\norm{J^{\rm t}_x(y)}^2-\Tr K^{\rm t}_x(y)\bigr)
   \,\dd\Gamma_\alpha(x,y)
   -\frac14\int\norm{\nabla\vartheta}^2\,\dd\pi,
 \label{eq:V22-Fisher-sum-rule}\\
 &= -\frac12\int\norm{\nabla\vartheta}^2\,\dd\pi
    -\frac12\int
       \langle\nabla\vartheta,\nabla G\rangle\,\dd\pi.
 \label{eq:V22-ibp-sum-rule}
\end{align}
In particular,
\begin{equation}
 \boxed{
 \int\Tr D_\alpha\,\dd\pi
 \le\frac18\int\norm{\nabla G}^2\,\dd\pi.}
 \label{eq:V22-signed-mean-upper}
\end{equation}
\end{theorem}

\begin{proof}
By definition of covariance and
\eqref{eq:V22-vartheta-score},
\[
 \Tr D_\alpha(x)
 =\mathbb E_p\norm{J^{\rm t}}^2
  -\norm{\mathbb E_pJ^{\rm t}}^2
  -\mathbb E_p\Tr K^{\rm t}
 =\mathbb E_p(\norm{J^{\rm t}}^2-\Tr K^{\rm t})
  -\frac14\norm{\nabla\vartheta}^2.
\]
Integrate against \(\pi\) and use
\eqref{eq:V22-edge-disintegration} to prove
\eqref{eq:V22-Fisher-sum-rule}.

Since \(D_\alpha=\Hess\log r_\alpha\), its trace is
\(\Delta\log r_\alpha\).  Moreover
\[
 \log h=\log r_\alpha-\frac12G-\log\lambda,
 \qquad
 \nabla\log\frac{\dd\pi}{\dd\mathfrak m}
 =2\nabla\log r_\alpha-\nabla G.
\]
Weighted integration by parts on the compact manifold gives
\[
 \int\Delta\log r_\alpha\,\dd\pi
 =-2\int\norm{\nabla\log r_\alpha}^2\,\dd\pi
  +\int\langle\nabla\log r_\alpha,\nabla G\rangle\,\dd\pi.
\]
Use \(\nabla\log r_\alpha=-\nabla\vartheta/2\) to obtain
\eqref{eq:V22-ibp-sum-rule}.  Finally,
\[
 -\frac12\norm v^2-\frac12\langle v,g\rangle
 =-\frac12\norm{v+g/2}^2+\frac18\norm g^2
\]
pointwise, which proves \eqref{eq:V22-signed-mean-upper}.
\end{proof}

\begin{firewall}[A signed trace is not the negative-channel stock]
\label{fire:V22-signed-not-positive}
The repaired source is \(2(D_\alpha)_+\), whereas
\eqref{eq:V22-Fisher-sum-rule} controls only the signed trace of
\(D_\alpha\).  Large positive and negative eigenvalues can cancel in that
trace.  Neither sum rule bounds
\(\int\Tr(D_\alpha)_+\,\dd\pi\), the pointwise operator norm, or the
source-compressed Gram.  Using \eqref{eq:V22-signed-mean-upper} as a bound on
the positive part would be invalid.
\end{firewall}

% =============================================================================
\section{A rigorous posterior-tilt perturbation bound}
\label{sec:V22-chi-square}
% =============================================================================

At a fixed \(x\), set
\begin{equation}
 \rho_x(y):=\frac{f(y)}{\mathbb E_{q_{\alpha,x}}f},
 \qquad
 \varepsilon_\alpha(x)^2
 :=\mathbb E_{q_{\alpha,x}}(\rho_x-1)^2.
 \label{eq:V22-rho-epsilon}
\end{equation}
Thus \(\dd p_{\alpha,x}=\rho_x\dd q_{\alpha,x}\) and
\(\mathbb E_q\rho_x=1\).  Introduce the untilted moments
\begin{align}
 M_2(x)&:=
  \bigl(\mathbb E_q\norm{J^{\rm t}}^2\bigr)^{1/2},
 \label{eq:V22-M2}\\
 M_4(x)&:=
  \bigl(\mathbb E_q\norm{J^{\rm t}}^4\bigr)^{1/4},
 \label{eq:V22-M4}\\
 N_2(x)&:=
  \bigl(\mathbb E_q\norm{K^{\rm t}}_{\rm HS}^2\bigr)^{1/2}.
 \label{eq:V22-N2}
\end{align}

\begin{theorem}[Pointwise chi-square stability of the exact cancellation]
\label{thm:V22-chi-square-stability}
For every \(x\),
\begin{equation}
 \boxed{
 \norm{D_\alpha(x)}_{\op}
 \le
 \varepsilon_\alpha(x)
   \bigl(M_4(x)^2+N_2(x)\bigr)
 +\varepsilon_\alpha(x)^2M_2(x)^2.}
 \label{eq:V22-chi-square-bound}
\end{equation}
Consequently,
\begin{equation}
 \norm{S_{\vartheta,-}(x)}_{\op}
 \le2\varepsilon_\alpha(x)
   \bigl(M_4(x)^2+N_2(x)\bigr)
 +2\varepsilon_\alpha(x)^2M_2(x)^2.
 \label{eq:V22-chi-square-defect-bound}
\end{equation}
\end{theorem}

\begin{proof}
By \Cref{lem:V22-untilted-cancellation},
\[
 \mathbb E_qJ^{\rm t}=0,
 \qquad
 \mathbb E_q(J^{\rm t}\otimes J^{\rm t})=\mathbb E_qK^{\rm t}.
\]
Therefore
\begin{align*}
 D_\alpha
 &=\mathbb E_q[(\rho-1)J^{\rm t}\otimes J^{\rm t}]
   -(\mathbb E_q[\rho J^{\rm t}])
      \otimes(\mathbb E_q[\rho J^{\rm t}])\\
 &\hspace{2.5em}
   -\mathbb E_q[(\rho-1)K^{\rm t}].
\end{align*}
Cauchy--Schwarz gives
\begin{align*}
 \norm{\mathbb E_q[(\rho-1)J^{\rm t}\otimes J^{\rm t}]}_{\op}
 &\le\varepsilon_\alpha M_4^2,\\
 \norm{\mathbb E_q[\rho J^{\rm t}]}
 =\norm{\mathbb E_q[(\rho-1)J^{\rm t}]}
 &\le\varepsilon_\alpha M_2,\\
 \norm{\mathbb E_q[(\rho-1)K^{\rm t}]}_{\op}
 &\le\varepsilon_\alpha N_2.
\end{align*}
The triangle inequality proves \eqref{eq:V22-chi-square-bound}; then use
\(\norm{(D_\alpha)_+}\le\norm{D_\alpha}\) and
\eqref{eq:V22-exact-defect-channel}.
\end{proof}

\begin{proposition}[Wilson temporal moment scale on a fixed bank]
\label{prop:V22-Wilson-moment-scale}
Fix the number of links \(|E|\).  For the V10 one-link density
\[
 q_\alpha(U)=Z_\alpha^{-1}
 \exp\!\left\{\frac\alpha3\Re\Tr U\right\}
\]
and its product convolution, there are constants
\(C_E<\infty\) and \(\alpha_E\) such that, uniformly in \(x\) and
\(\alpha\ge\alpha_E\),
\begin{equation}
 \boxed{
 M_2(x)^2\le C_E\alpha,
 \qquad
 M_4(x)^2\le C_E\alpha,
 \qquad
 N_2(x)\le C_E\alpha.}
 \label{eq:V22-Wilson-moments}
\end{equation}
Hence
\begin{equation}
 \boxed{
 \norm{S_{\vartheta,-}(x)}_{\op}
 \le C_E'\alpha
 \bigl(\varepsilon_\alpha(x)+\varepsilon_\alpha(x)^2\bigr).}
 \label{eq:V22-Wilson-chi-square-scale}
\end{equation}
The constants may depend on the fixed bank but not on \(x\) or \(\alpha\).
\end{proposition}

\begin{proof}
For one link write
\(a(U)=\Re\Tr U/3\).  Translation invariance makes the moments independent
of the initial link value.  The temporal score and Hessian are respectively
\(-\alpha\nabla a\) and \(-\alpha\Hess a\), up to an isometric left or
right translation.  Near the unique maximum \(U=I\), in the V21 normal
coordinates \(U=e^X\),
\[
 1-a(e^X)=\frac16\norm X^2+O(\norm X^4),
 \qquad
 \norm{\nabla a(e^X)}\le C\norm X,
 \qquad
 \norm{\Hess a(e^X)}\le C.
\]
The V11 global quadratic-well estimate and Laplace rescaling imply
\[
 \mathbb E_{q_\alpha}\norm X^{2m}\le C_m\alpha^{-m}
 \quad(m=1,2),
\]
while the complement of a fixed normal neighborhood is exponentially
small.  Thus
\[
 \mathbb E\norm{J^{\rm t}}^2=O(\alpha),
 \quad
 \mathbb E\norm{J^{\rm t}}^4=O(\alpha^2),
 \quad
 \mathbb E\norm{K^{\rm t}}_{\rm HS}^2=O(\alpha^2).
\]
Independence and the finite product dimension change only the constant.
This proves \eqref{eq:V22-Wilson-moments}.  Insert it into
\eqref{eq:V22-chi-square-defect-bound}.
\end{proof}

% =============================================================================
\section{The fixed-smooth-tilt calibration and its loss certificate}
\label{sec:V22-smooth-tilt}
% =============================================================================

Let
\begin{equation}
 \mathcal L_E:=-\sum_{e\in E}\Delta_e
 \label{eq:V22-total-Casimir}
\end{equation}
be the positive total electric Casimir.  Its Peter--Weyl eigenvalue on a
product representation is the V10 total Casimir \(C_E(\mathbf R)\).

\begin{theorem}[Second-order Wilson approximate identity]
\label{thm:V22-Wilson-approximate-identity}
On every fixed finite link bank, for \(u\in C^6(\mathcal X)\),
\begin{equation}
 \boxed{
 \mathcal C_\alpha u
 =u-\frac{3}{2\alpha}\mathcal L_Eu
  +O_{u,E}(\alpha^{-2})
 \quad\text{in }C^2(\mathcal X).}
 \label{eq:V22-Wilson-C2-expansion}
\end{equation}
\end{theorem}

\begin{proof}
For one link, split the integral into a fixed exponential neighborhood of
the identity and its complement.  The complement is exponentially small by
the V11 quadratic-well estimate.  In the neighborhood put
\(U=e^X\) and use the V21 expansion
\[
 q_\alpha(e^X)\dd U
 \propto
 e^{-\alpha\norm X^2/6}
 \bigl(1+O(\norm X^2)+O(\alpha\norm X^4)\bigr)\dd X.
\]
Inversion symmetry eliminates every odd Taylor moment, and Laplace
rescaling gives
\begin{equation}
 \mathbb E(X_iX_j)
 =\frac3\alpha\delta_{ij}+O(\alpha^{-2}),
 \qquad
 \mathbb E\norm X^4=O(\alpha^{-2}).
 \label{eq:V22-increment-moments}
\end{equation}
Taylor's formula with integral remainder therefore yields
\[
 \mathcal C_\alpha u
 =u+\frac{3}{2\alpha}\Delta u+O_{u}(\alpha^{-2})
 =u-\frac{3}{2\alpha}\mathcal Lu+O_u(\alpha^{-2}).
\]
For the product kernel, mixed second moments vanish and the one-link
generators add.  Left-invariant derivatives commute with convolution;
applying the same argument to derivatives of order at most two requires at
most six derivatives of \(u\) and proves the \(C^2\) statement.
\end{proof}

\begin{corollary}[Exact scale for a fixed positive terminal tilt]
\label{cor:V22-fixed-tilt-scale}
Let \(f\in C^6(\mathcal X)\) be independent of \(\alpha\) and satisfy
\(\inf f>0\).  Then, uniformly in \(x\),
\begin{align}
 D_\alpha
 &=\Hess\log f
   -\frac{3}{2\alpha}
      \Hess\!\left(\frac{\mathcal L_Ef}{f}\right)
   +O_{f,E}(\alpha^{-2}),
 \label{eq:V22-D-fixed-tilt}\\
 S_{\vartheta,-}
 &=2(\Hess\log f)_++O_{f,E}(\alpha^{-1})
 \quad\text{in operator norm}.
 \label{eq:V22-Sminus-fixed-tilt}
\end{align}
Moreover the chi-square tilt in \eqref{eq:V22-rho-epsilon} has the uniform
expansion
\begin{equation}
 \boxed{
 \varepsilon_\alpha(x)^2
 =\frac3\alpha\norm{\nabla\log f(x)}^2
  +O_{f,E}(\alpha^{-2}).}
 \label{eq:V22-epsilon-fixed-tilt}
\end{equation}
\end{corollary}

\begin{proof}
Apply \Cref{thm:V22-Wilson-approximate-identity} to \(f\).  Positivity
bounded away from zero permits the \(C^2\) logarithmic expansion
\[
 \log(\mathcal C_\alpha f)
 =\log f-\frac{3}{2\alpha}\frac{\mathcal L_Ef}{f}
  +O_{f,E}(\alpha^{-2}).
\]
Twice differentiating and using
\(D_\alpha=\Hess\log(\mathcal C_\alpha f)\) proves
\eqref{eq:V22-D-fixed-tilt}.  On each finite-dimensional cotangent fibre,
the positive part is the Hilbert--Schmidt orthogonal projection onto the
closed positive cone.  That projection is nonexpansive in Hilbert--Schmidt
norm.  Since the bank dimension is fixed, equivalence of matrix norms gives
\eqref{eq:V22-Sminus-fixed-tilt}.

For the last statement,
\[
 1+\varepsilon_\alpha(x)^2
 =\frac{(\mathcal C_\alpha f^2)(x)}
        {(\mathcal C_\alpha f)(x)^2}.
\]
Use \eqref{eq:V22-Wilson-C2-expansion} in \(C^0\), together with
\[
 \Delta(f^2)=2f\Delta f+2\norm{\nabla f}^2.
\]
The terms containing \(f\Delta f\) cancel between numerator and
denominator, leaving \(3\norm{\nabla f}^2/(\alpha f^2)\) plus a uniform
\(O(\alpha^{-2})\) remainder.
\end{proof}

\begin{assessment}[The chi-square route loses the exact cancellation]
\label{ass:V22-chi-square-loss}
At a point where \(\nabla\log f\ne0\),
\eqref{eq:V22-epsilon-fixed-tilt} gives
\(\varepsilon_\alpha\asymp\alpha^{-1/2}\).  The moment estimate
\eqref{eq:V22-Wilson-chi-square-scale} then gives only
\[
 \norm{S_{\vartheta,-}}=O_f(\sqrt\alpha),
\]
whereas the exact convolution expansion
\eqref{eq:V22-Sminus-fixed-tilt} gives \(O_f(1)\).  Thus pointwise
chi-square closeness of the posterior to the temporal row is a valid but
generically non-sharp acquisition.  It can close a physically rescaled row
only if the additional schedule condition
\(\kappa_H\sqrt\alpha=O(1)\) is separately proved; V21 does not supply that
condition.
\end{assessment}

\begin{proposition}[The Perron tilt is not a fixed external function]
\label{prop:V22-Perron-self-consistency}
For the physical tilt \(f=e^{-G/2}h\),
\begin{equation}
 \boxed{
 \lambda f=e^{-G}\mathcal C_\alpha f,
 \qquad
 r_\alpha=\lambda e^Gf,
 \qquad
 D_\alpha=\Hess(G+\log f)
 =\Hess\log h+\frac12\Hess G.}
 \label{eq:V22-f-self-consistency}
\end{equation}
\end{proposition}

\begin{proof}
Multiply \eqref{eq:V22-Perron-factorization} by \(e^{-G/2}\) and use
\(f=e^{-G/2}h\) to get \(\lambda f=e^{-G}r_\alpha\).  Rearranging gives
the first two identities.  Take logarithms and Hessians.  Alternatively,
\(\log r_\alpha=\log h+G/2+\log\lambda\).
\end{proof}

\begin{firewall}[Why the fixed-tilt expansion is a calibration, not a
cofinal theorem]
\label{fire:V22-fixed-tilt-not-cofinal}
The hypotheses of \Cref{cor:V22-fixed-tilt-scale} are not automatic for
the physical family \(f_{\alpha,j}\).  The self-consistency equation
\eqref{eq:V22-f-self-consistency} allows the Perron factor to depend on the
temporal concentration, spatial regulator, volume, and magnetic
coefficient.  A uniform positive lower bound and a uniform \(C^6\) bound
would themselves be substantial cofinal estimates.  The corollary proves
the correct cancellation scale and diagnoses the loss in
\eqref{eq:V22-chi-square-bound}; it does not assert those missing uniform
bounds.
\end{firewall}

% =============================================================================
\section{Going upstream: the V9 Perron-polymer bridge}
\label{sec:V22-V9-bridge}
% =============================================================================

The last identity in \eqref{eq:V22-f-self-consistency} makes the earlier V9
half-space estimate directly relevant.  To state the bridge without hiding
an incidence, suppose that the V9 operator
\(\mathcal K_{\beta,\Lambda}\) has been identified with the complete
symmetric transfer \eqref{eq:V22-symmetric-transfer} on the same
configuration variables and that its Perron factor is the same \(h\).
Suppose also that
\begin{equation}
 G_{\beta,\Lambda}=\beta V_{M,\Lambda}
 \label{eq:V22-G-beta-VM}
\end{equation}
with the sign convention of the V21 outer multiplier.  For
\(\mu\ge0\), define the volume-free finite-range Hessian row constant
\begin{equation}
 C_M^{(2)}(\mu)
 :=\sup_{\Lambda,b}\sum_c
 e^{\mu\dist(b,c)}
 \norm{\nabla_c\nabla_bV_{M,\Lambda}}_{\op}<\infty.
 \label{eq:V22-CM2}
\end{equation}
Finiteness follows from compactness of \(\mathrm{SU}(3)\) and the uniformly
bounded number of plaquettes incident with a link.

\begin{theorem}[Strong-coupling bound for the repaired density channel]
\label{thm:V22-strong-density-bound}
Under the preceding complete-transfer incidence, for
\(|\beta|<\beta_{\rm KP}\) and
\(0\le\mu<-\log\vartheta_\beta\),
\begin{equation}
 \sup_b\sum_c e^{\mu\dist(b,c)}
 \norm{(D_\alpha)_{bc}}_{\op}
 \le
 \mathfrak B_2(\beta,\mu)
 +\frac{|\beta|}{2}C_M^{(2)}(\mu).
 \label{eq:V22-D-weighted-row}
\end{equation}
In particular,
\begin{equation}
 \boxed{
 \norm{S_{\vartheta,-}}_{\op}
 \le
 2\mathfrak B_2(\beta,0)
 +|\beta|C_M^{(2)}(0).}
 \label{eq:V22-Sminus-strong-bound}
\end{equation}
There are \(\beta_0>0\) and \(C_{\rm den}<\infty\), independent of the
periodic volume, such that
\begin{equation}
 |\beta|\le\beta_0
 \quad\Longrightarrow\quad
 \boxed{
 \norm{S_{\vartheta,-}}_{\op}
 \le C_{\rm den}|\beta|.}
 \label{eq:V22-Sminus-linear-beta}
\end{equation}
\end{theorem}

\begin{proof}
By \eqref{eq:V22-f-self-consistency},
\[
 D_\alpha=\Hess\log h_{\beta,\Lambda}
             +\frac\beta2\Hess V_{M,\Lambda}.
\]
The V9 weighted Perron-Hessian theorem bounds the first block row by
\(\mathfrak B_2(\beta,\mu)\); definition
\eqref{eq:V22-CM2} bounds the second.  This proves
\eqref{eq:V22-D-weighted-row}.  The symmetric block Schur test gives the
same right side as a global operator-norm bound for \(D_\alpha\).  Since
\(S_{\vartheta,-}=2(D_\alpha)_+\) and
\(\norm{D_+}\le\norm D\), we obtain
\eqref{eq:V22-Sminus-strong-bound}.

It remains to check the stated order.  From the explicit V9 definitions,
as \(\beta\to0\),
\[
 u(\beta)=O(|\beta|),\quad
 q_1(\beta)=O(|\beta|),\quad
 q_2(\beta)=O(|\beta|),\quad
 \frac{q_1(\beta)^2}{u(\beta)}=O(|\beta|),
\]
with the quotient defined continuously at zero.  Hence
\(m_2(\beta)=O(|\beta|)\) and
\(\vartheta_\beta=O(|\beta|)\).  The rational factor in the definition of
\(\mathfrak B_2(\beta,0)\) is therefore bounded for sufficiently small
\(|\beta|\), so
\(\mathfrak B_2(\beta,0)\le C_B|\beta|\).  Insert this in
\eqref{eq:V22-Sminus-strong-bound} and take
\(C_{\rm den}=2C_B+C_M^{(2)}(0)\).
\end{proof}

\begin{corollary}[Physical scale under literal anisotropic matching]
\label{cor:V22-physical-density-scale}
Assume in addition the V21 matching
\begin{equation}
 \beta_{\alpha,j}=\tau_\alpha\gamma_{\alpha,j},
 \qquad
 \kappa_{H,j}=\frac{s_\alpha}{\tau_\alpha},
 \label{eq:V22-physical-matching}
\end{equation}
and that \(|\beta_{\alpha,j}|\le\beta_0\).  Then
\begin{equation}
 \boxed{
 \kappa_{H,j}
 \norm{S_{\vartheta_j,-}}_{\op}
 \le C_{\rm den}s_\alpha|\gamma_{\alpha,j}|.}
 \label{eq:V22-physical-density-bound}
\end{equation}
Consequently a uniform bound on
\(|\gamma_{\alpha,j}|/\alpha\) gives a uniform physical density-defect
stock, because \(\alpha s_\alpha\to3/2\).
\end{corollary}

\begin{proof}
Multiply \eqref{eq:V22-Sminus-linear-beta} by
\(\kappa_{H,j}\) and substitute \eqref{eq:V22-physical-matching}.
The temporal spacing cancels:
\[
 \kappa_{H,j}|\beta_{\alpha,j}|
 =s_\alpha|\gamma_{\alpha,j}|.
\]
The final assertion is the V21 temporal calibration.
\end{proof}

\begin{assessment}[What the upstream bridge closes]
\label{ass:V22-upstream-closure}
On the branch where the V9 and V21 complete transfers are literally the
same operator, \Cref{thm:V22-strong-density-bound} closes the previously
open magnitude of \(S_{\vartheta,-}\) with a volume-uniform bound.  Under
the same coefficient ratio that controls the V21 Wilson defect, the new
channel is also uniformly bounded in physical units.  This is stronger and
sharper than the generic chi-square estimate.

Two limitations remain.  First, the supplied reduced-fibre V35 kernel and
the complete V9/V21 transfer still require the explicit conditional
incidence described in \Cref{fire:V22-fibre-incidence}.  Second, an
exponentially summable row for \(D_\alpha\) does not automatically give an
exponentially summable row for the nonsmooth spectral functions
\((D_\alpha)_+\) or \((D_\alpha)_+^{1/2}\) when the spectrum can meet zero.
Thus the global norm stock is closed on this branch, while the local
complete-source factorization remains open.
\end{assessment}

% =============================================================================
\section{A source-compressed good/bad gate for the positive excess}
\label{sec:V22-source-gate}
% =============================================================================

Let \(P_{\mathscr C}\) be a finite-rank projection onto one complete signed
V15 source carrier.  Define
\begin{align}
 b_{\mathscr C}(x)
 &:=\norm{P_{\mathscr C}(D_\alpha(x))_+P_{\mathscr C}}_{\op},
 \label{eq:V22-source-b}\\
 a_{\mathscr C}(x)
 &:=\Tr\!\left(P_{\mathscr C}(D_\alpha(x))_+P_{\mathscr C}\right),
 \qquad
 \overline a_{\mathscr C}:=\int a_{\mathscr C}\,\dd\pi.
 \label{eq:V22-source-a}
\end{align}
For \(L>0\), set
\begin{equation}
 \mathsf B_{\mathscr C}(L)
 :=\{x:b_{\mathscr C}(x)>L\}.
 \label{eq:V22-source-bad-set}
\end{equation}

\begin{proposition}[Exact threshold card]
\label{prop:V22-threshold-card}
For every \(L>0\),
\begin{align}
 \pi(\mathsf B_{\mathscr C}(L))
 &\le\frac{\overline a_{\mathscr C}}{L},
 \label{eq:V22-Markov-card}\\
 \one_{\mathsf B_{\mathscr C}(L)^c}
 \norm{P_{\mathscr C}S_{\vartheta,-}P_{\mathscr C}}_{\op}
 &\le2L.
 \label{eq:V22-good-source-card}
\end{align}
Thus on the good event the physical density column obeys
\begin{equation}
 \one_{\mathsf B_{\mathscr C}(L)^c}
 \norm{\sqrt{\kappa_H}\,
          \mathcal V_{\rm den}P_{\mathscr C}}^2
 \le2\kappa_HL.
 \label{eq:V22-good-physical-column}
\end{equation}
\end{proposition}

\begin{proof}
The compression of a positive operator is positive, so its operator norm is
at most its trace.  Hence
\[
 L\,\one_{\mathsf B_{\mathscr C}(L)}
 \le a_{\mathscr C}.
\]
Integrate to obtain \eqref{eq:V22-Markov-card}.  Equations
\eqref{eq:V22-good-source-card}--
\eqref{eq:V22-good-physical-column} follow from
\(S_{\vartheta,-}=2(D_\alpha)_+\) and
\(\mathcal V_{\rm den}^*\mathcal V_{\rm den}=S_{\vartheta,-}\).
\end{proof}

\begin{assessment}[Capacity input now has a literal event]
\label{ass:V22-capacity-event}
Equation \eqref{eq:V22-source-bad-set} is the correct bad event for the new
density channel on a declared source carrier.  The V20 incidence and
capacity matrices should acquire one additional column labelled
\(\mathrm{den}\), with source amplitude
\(\sqrt2(D_\alpha)_+^{1/2}P_{\mathscr C}\), and route precisely this event.
The signed sum rules do not populate \(\overline a_{\mathscr C}\); either
the V9 strong branch, a differentiated posterior expansion, or a genuine
V13 capacity construction must do so.
\end{assessment}

% =============================================================================
\section{V22 ledger and revised frontier}
\label{sec:V22-ledger}
% =============================================================================

\begin{theorem}[Integrated V22 statement]
\label{thm:V22-integrated}
Preserving V1--V21, V22 proves:
\begin{enumerate}[label=\textup{(V22.\arabic*)},leftmargin=5.5em]
\item the exact temporal-posterior representation
      \(S_{\vartheta,-}=2(\Cov_pJ^{\rm t}-\mathbb E_pK^{\rm t})_+\);
\item exact cancellation of temporal covariance against direct curvature
      for an untilted convolution row;
\item posterior Fisher and weighted integration-by-parts sum rules for the
      signed excess;
\item a pointwise chi-square stability theorem with explicit temporal
      moments;
\item the fixed-bank \(C^2\) Wilson approximate-identity expansion with
      coefficient \(3/(2\alpha)\);
\item a loss certificate showing why raw chi-square perturbation is not a
      sharp estimate of the density channel;
\item a V9-to-V21 upstream bridge giving
      \(\norm{S_{\vartheta,-}}\le C_{\rm den}|\beta|\) on the literally
      identified strong-coupling transfer branch;
\item the physical bound
      \(\kappa_H\norm{S_{\vartheta,-}}\le
        C_{\rm den}s_\alpha|\gamma|\) under anisotropic coefficient
      matching; and
\item a source-compressed good/bad event for the added density column.
\end{enumerate}
\end{theorem}

\begin{proof}
Items \textup{(V22.1)--(V22.2)} are
\Cref{thm:V22-posterior-excess,lem:V22-untilted-cancellation}.  Item
\textup{(V22.3)} is \Cref{thm:V22-signed-sum-rules}.  Item
\textup{(V22.4)} is
\Cref{thm:V22-chi-square-stability,prop:V22-Wilson-moment-scale}.  Items
\textup{(V22.5)--(V22.6)} are
\Cref{thm:V22-Wilson-approximate-identity,cor:V22-fixed-tilt-scale,
ass:V22-chi-square-loss}.  Items \textup{(V22.7)--(V22.8)} are
\Cref{thm:V22-strong-density-bound,cor:V22-physical-density-scale}.  Item
\textup{(V22.9)} is \Cref{prop:V22-threshold-card}.
\end{proof}

\begingroup
\small
\begin{longtable}{>{\raggedright\arraybackslash}p{0.28\textwidth}
                  >{\raggedright\arraybackslash}p{0.16\textwidth}
                  >{\raggedright\arraybackslash}p{0.48\textwidth}}
\toprule
\textbf{V22 row}&\textbf{Status}&\textbf{Advance or remaining obstruction}\\
\midrule
\endfirsthead
\toprule
\textbf{V22 row}&\textbf{Status}&\textbf{Advance or remaining obstruction}\\
\midrule
\endhead
Identity of the density source
& \MGpass
& Exactly \(\sqrt2(D_\alpha)_+^{1/2}\), with
  \(D_\alpha=\Cov_pJ^{\rm t}-\mathbb E_pK^{\rm t}\).\\[0.35em]
Untilting cancellation
& \MGpass
& Exact at every temporal concentration; no asymptotic approximation.\\[0.35em]
Signed mean identities
& \MGpass
& Exact, but deliberately not promoted to a positive-part estimate.\\[0.35em]
Generic chi-square control
& \MGpass
& Rigorous and explicit, but loses \(\sqrt\alpha\) on a generic smooth
  fixed tilt.\\[0.35em]
Strong-branch density magnitude
& \MGpass
& \(O(|\beta|)\), conditional on literal identity of the complete V9 and
  V21 transfers.\\[0.35em]
Physical density scale
& \MGpass
& Uniform if the same open coefficient ratio
  \(|\gamma_{\alpha,j}|/\alpha\) is uniform.\\[0.35em]
Complete-transfer/fibre incidence
& \MGopen
& Normalizer derivatives and common-history variables still require an
  explicit disintegration.\\[0.35em]
Locality of the positive square root
& \MGopen
& Decay of \(D_\alpha\) does not by itself control
  \((D_\alpha)_+^{1/2}\) across spectrum zero.\\[0.35em]
Repaired source edge and capacity
& \MGopen
& The density column must be added to the V20 incidence and capacity
  matrices.\\[0.35em]
Cofinal coefficient trajectory
& \MGopen
& The supplied files still do not fix
  \(\gamma_{\alpha,j}/\alpha\) on a common spatial continuum schedule.\\[0.35em]
Protected sectors and nontrivial OS limit
& \MGopen
& Uniform sector cards and local-algebra noncollapse remain.\\[0.35em]
Full Yang--Mills mass gap
& \MGopen
& Complete-source incidence, strict edge, capacity, and continuum
  nontriviality remain.\\
\bottomrule
\end{longtable}
\endgroup

\begin{openproblem}[V23 mandate: make the density source local on the
complete carrier]
\label{op:V23}
\begin{enumerate}[label=\textup{(V23.\arabic*)},leftmargin=5.5em]
\item Write the complete anisotropic one-step kernel in common-history
      variables and disintegrate it over the V35 reduced fibre.  Record every
      derivative of the fibre normalizer and decide whether the V9 and V21
      Perron factors are literally the same object.
\item If the incidence holds, apply the weighted V9 row
      \eqref{eq:V22-D-weighted-row} to the complete signed V15 source.  Do
      not infer locality of \((D_\alpha)_+^{1/2}\) from locality of
      \(D_\alpha\); instead seek a factorization before taking the global
      spectral positive part.
\item Test the shifted factorization
      \(D_+=D+aI-(D+aI)_-\) only with an explicit uniform spectral shift and
      track the additional comparison cost.  Reject it if the shift destroys
      the V35 floor.
\item Add the literal density column
      \eqref{eq:V22-density-synthesis-column} to the V20 compressed Gram,
      source-incidence matrix, predictable martingale, and capacity matrix.
\item Prove a strict source-cyclic edge for the resulting repaired synthesis,
      or return a normalized density-channel landing witness and classify it
      with the V18--V20 alternatives.
\item Retain \(\gamma_{\alpha,j}/\alpha\) as an explicit cofinal input until
      it is read from a response-matched anisotropic action.  Combine the
      edge and capacity with protected sectors and the V1 OS passage only
      after every physical constant is uniform.
\end{enumerate}
\end{openproblem}

\begin{firewall}[End-of-V22 claim boundary]
V22 identifies the Perron density repair as a temporal posterior excess,
proves its exact cancellation and perturbation formulas, and closes its
global magnitude on the literally matched V9 strong-coupling branch.  It
does not prove the complete reduced-fibre incidence, locality of the
positive square-root source, the strict repaired edge, the bad-set capacity,
the cofinal spatial coefficient law, all protected sectors, or a nontrivial
continuum OS state.  Therefore the full Yang--Mills mass gap remains open.
\end{firewall}

\section*{Version V22 append-only record}
\addcontentsline{toc}{section}{Version V22 append-only record}
The complete V21 source was retained before this continuation.  Its SHA--256
digest was
\begin{center}\scriptsize\ttfamily
7c80642af5d8e683bf073090bf11d10403cbe2f17c0f878bf0603723ec314876.
\end{center}
On the next iteration, retain this full file, append before its unique active
document-closing command, increment the filename to
\texttt{\_V23.tex}, and remove only the superseded V22 file from this note
series.

% =============================================================================
% END APPENDED VERSION V22 MATERIAL
% =============================================================================

% =============================================================================
% BEGIN APPENDED VERSION V23 MATERIAL
% =============================================================================

\clearpage
\part{Version V23: reduced-fibre incidence, likelihood repair, and a local shifted density channel}
\label{part:V23}

\section{Direction and source audit}
\label{sec:V23-direction}

V22 reduced the Perron density repair to the positive part of the temporal
posterior excess
\[
 D_\alpha=\Cov_{p_{\alpha,x}}(J^{\rm t})
           -\mathbb E_{p_{\alpha,x}}K^{\rm t}.
\]
Its strong-branch estimate was deliberately conditional on the V9 complete
transfer and the V35 reduced-fibre construction having the same physical
Perron factor.  V23 audits that incidence against the supplied sources and
then removes two avoidable assumptions.

The audit has a definite outcome.  The supplied V35 reduced-fibre manuscript
defines, at a fixed retained coarse field \(W\), the static partition
\[
 Z_M^{\rm tr}(W;\beta)
 =\int e^{-\mathcal S_{M,\beta}(z;W)}\,\dd m_M(z)
\]
and its normalized detail law.  Its Versions V32--V35 explicitly distinguish
that static trace-shell card from a reflected half-amplitude, command
incidence, and transfer identification.  The same manuscript proves that a
diagonal partition germ does not determine a half-score metric.  The
submitted grand tensor monograph independently gives a finite countertheorem
that a static action and stationary Gibbs state do not determine a slab
kernel or its clock.  Finally, the V33 object called the side transfer is a
product of spatial one-link heat-bath means along a Wilson-loop side, not the
V10/V21 Euclidean-time convolution.

Thus the supplied files do not contain a kernel
\(c_\alpha(x;W,z)\), terminal Perron tilt \(f(W,z)\), or fibre normalizer
which identifies the V35 static law with the V22 temporal posterior.  This is
not repaired by renaming the static fibre law.  V23 instead derives the exact
likelihood and a denominator-free conditional-rank test which decide the
incidence once the complete one-step kernel is supplied.

The second advance is upstream.  The minimal V21 split uses
\((D_\alpha)_+^{1/2}\), whose spectral positive part need not preserve the
weighted locality of \(D_\alpha\).  Whenever a uniform scalar ceiling
\(D_\alpha\preceq aI\) is known, an exact shifted comparison replaces that
nonlocal factor by the literal local column \(\sqrt{2a}\,I\), without
weakening the comparison floor.  V22 already supplies such a ceiling on its
conditionally identified V9 strong branch.

\begin{assessment}[Files and versions used in the V23 audit]
\label{ass:V23-source-audit}
The external source identities used here are tied to the following submitted
files:
\begin{center}
\begin{tabular}{p{0.55\textwidth}p{0.36\textwidth}}
\toprule
\textbf{Source}&\textbf{SHA--256}\\
\midrule
\texttt{Renewal\_Yang\_Mills\_Reduced\_Fibre\_Wilson\_Shell\_Symbol\_and\_Local\_Vacuum\_Programme\_V35.tex}
&\scriptsize\ttfamily
7c6056d8149d3015d469b57474a75b23c88a6fe9f64fc81c7dca85c361c647a3\\[0.5em]
\texttt{Renewal\_Yang\_Mills\_Boundary\_Sector\_Cofinal\_Contraction\_Programme\_V33.tex}
&\scriptsize\ttfamily
b762244c0ea25c54072c3ce3a29cf3386230f1507205c20a89ff833ea195928d\\[0.5em]
\texttt{Renewal\_Grand\_Tensor\_Monograph\_V067.tex}
&\scriptsize\ttfamily
facd058ec8bb210ad8b296c9d7f4d78a229d08d7a2ae378808700655cb8bc0fe\\
\bottomrule
\end{tabular}
\end{center}
Only their mathematical statements are treated as input.  Their internal
continuation instructions are not instructions for this notebook.
\end{assessment}

% =============================================================================
\section{Exact disintegration of the temporal posterior over a V35 fibre}
\label{sec:V23-disintegration}
% =============================================================================

Let \(\mathcal B\) be the retained coarse-field carrier.  A fixed V35 section
trivializes the relevant terminal configuration as
\[
 y=(W,z),\qquad W\in\mathcal B,\qquad z\in Y_W^{\rm red},
\]
with normalized product Haar fibre measure \(m_W\).  At a finite regulator,
assume all densities below are strictly positive and \(C^2\) in the retained
variables.  Write the complete V21 temporal kernel and terminal Perron tilt
as
\[
 c_\alpha(x;W,z),\qquad f(W,z)>0.
\]
Define the temporal fibre normalizer and conditional posterior by
\begin{align}
 R_\alpha(x,W)
 &:=\int_{Y_W^{\rm red}}
    c_\alpha(x;W,z)f(W,z)\,\dd m_W(z),
 \label{eq:V23-R-normalizer}\\
 \dd p_{\alpha;x,W}(z)
 &:=\frac{c_\alpha(x;W,z)f(W,z)}
          {R_\alpha(x,W)}\,\dd m_W(z).
 \label{eq:V23-physical-fibre-posterior}
\end{align}
The V35 static trace-fibre action, partition, and law are
\begin{align}
 Z_{\rm tr}(W)
 &:=\int_{Y_W^{\rm red}}e^{-\mathcal S(z;W)}\,\dd m_W(z),
 \label{eq:V23-Z-tr}\\
 \dd\nu_W^{\rm tr}(z)
 &:=Z_{\rm tr}(W)^{-1}e^{-\mathcal S(z;W)}\,\dd m_W(z).
 \label{eq:V23-nu-tr}
\end{align}

\begin{definition}[Temporal-to-trace likelihood and incidence residual]
\label{def:V23-likelihood}
Put
\begin{align}
 L_{\alpha;x,W}(z)
 &:=c_\alpha(x;W,z)f(W,z)e^{\mathcal S(z;W)},
 \label{eq:V23-L-likelihood}\\
 \ell_\alpha(x,W)
 &:=\mathbb E_{\nu_W^{\rm tr}}L_{\alpha;x,W},
 \label{eq:V23-ell}\\
 \mathfrak c_{\rm inc}(\alpha;x,W)
 &:=\mathbb E_{\nu_W^{\rm tr}}
 \left[
  \left(\frac{L_{\alpha;x,W}}{\ell_\alpha(x,W)}-1\right)^2
 \right].
 \label{eq:V23-incidence-chi-square}
\end{align}
\end{definition}

\begin{theorem}[Exact fibre likelihood and zero-residual incidence test]
\label{thm:V23-fibre-likelihood}
At every finite regulator,
\begin{align}
 \boxed{
 R_\alpha(x,W)
 =Z_{\rm tr}(W)\ell_\alpha(x,W),}
 \label{eq:V23-normalizer-factorization}\\
 \boxed{
 \frac{\dd p_{\alpha;x,W}}{\dd\nu_W^{\rm tr}}(z)
 =\frac{L_{\alpha;x,W}(z)}{\ell_\alpha(x,W)}.}
 \label{eq:V23-RN-fibre}
\end{align}
Consequently the following are equivalent:
\begin{equation}
 \boxed{
 p_{\alpha;x,W}=\nu_W^{\rm tr}
 \iff
 L_{\alpha;x,W}\text{ is constant }\nu_W^{\rm tr}\text{-a.s.}
 \iff
 \mathfrak c_{\rm inc}(\alpha;x,W)=0.}
 \label{eq:V23-incidence-equivalence}
\end{equation}
If the fibre is connected and \(L\) is \(C^1\), these conditions are also
equivalent to
\begin{equation}
 \nabla_z\log L_{\alpha;x,W}(z)=0
 \quad\text{for every }z.
 \label{eq:V23-fibre-score-zero}
\end{equation}
\end{theorem}

\begin{proof}
By \eqref{eq:V23-L-likelihood},
\[
 c_\alpha f=L_{\alpha;x,W}e^{-\mathcal S}.
\]
Integration over the fibre gives
\[
 R_\alpha
 =Z_{\rm tr}\mathbb E_{\nu_W^{\rm tr}}L_{\alpha;x,W}
 =Z_{\rm tr}\ell_\alpha,
\]
which is \eqref{eq:V23-normalizer-factorization}.  Divide the two densities
\eqref{eq:V23-physical-fibre-posterior} and
\eqref{eq:V23-nu-tr}, then use that factorization to obtain
\eqref{eq:V23-RN-fibre}.

Two probability laws with positive densities agree exactly when their
Radon--Nikodym derivative is one almost surely.  Since
\(\ell_\alpha=\mathbb E_\nu L\), that is equivalent to \(L=\ell_\alpha\)
almost surely, and the latter is equivalent to vanishing of the nonnegative
square in \eqref{eq:V23-incidence-chi-square}.  Positivity and continuity
promote almost-everywhere constancy to pointwise constancy; on a connected
fibre this is equivalent to \eqref{eq:V23-fibre-score-zero}.
\end{proof}

\begin{corollary}[Every retained normalizer derivative]
\label{cor:V23-normalizer-derivatives}
In the fixed V35 section, where \(m_W\) is the declared product Haar
reference, one has
\begin{align}
 \log R_\alpha(x,W)
 &=\log Z_{\rm tr}(W)+\log\ell_\alpha(x,W),
 \label{eq:V23-log-normalizers}\\
 D_W\log R_\alpha
 &=D_W\log Z_{\rm tr}+D_W\log\ell_\alpha,
 \label{eq:V23-first-normalizer}\\
 D_W^2\log R_\alpha
 &=D_W^2\log Z_{\rm tr}+D_W^2\log\ell_\alpha.
 \label{eq:V23-second-normalizer}
\end{align}
If the exact incidence in \eqref{eq:V23-incidence-equivalence} holds on a
neighborhood of \(W\), then
\begin{align}
 D_W\!\left(\log c_\alpha+\log f+\mathcal S\right)
 &=D_W\log R_\alpha-D_W\log Z_{\rm tr},
 \label{eq:V23-first-incidence-jet}\\
 D_W^2\!\left(\log c_\alpha+\log f+\mathcal S\right)
 &=D_W^2\log R_\alpha-D_W^2\log Z_{\rm tr}.
 \label{eq:V23-second-incidence-jet}
\end{align}
The right sides are independent of \(z\).  Away from incidence, the exact
first and second log-law residuals are
\begin{align}
 \mathscr I_1(z)
 &:=
 D_W\log L_{\alpha;x,W}(z)-D_W\log\ell_\alpha(x,W),
 \label{eq:V23-I1}\\
 \mathscr I_2(z)
 &:=
 D_W^2\log L_{\alpha;x,W}(z)-D_W^2\log\ell_\alpha(x,W).
 \label{eq:V23-I2}
\end{align}
\end{corollary}

\begin{proof}
The first three identities are successive derivatives of
\eqref{eq:V23-normalizer-factorization}.  Under incidence,
\(\log L=\log\ell_\alpha\); differentiate this equality and expand
\(\log L=\log c_\alpha+\log f+\mathcal S\).  In general,
\[
 \log\frac{\dd p_{\alpha;x,W}}{\dd\nu_W^{\rm tr}}
 =\log L_{\alpha;x,W}-\log\ell_\alpha,
\]
so its first two retained jets are precisely
\eqref{eq:V23-I1}--\eqref{eq:V23-I2}.
\end{proof}

\begin{firewall}[The normalizer is part of the physical Hessian]
\label{fire:V23-normalizer-not-gauge}
The factor \(\ell_\alpha=R_\alpha/Z_{\rm tr}\) depends on the initial slice,
the retained field, and the regulator.  It is not an additive constant in
the retained effective action.  Dropping
\(D_W\log\ell_\alpha\) or \(D_W^2\log\ell_\alpha\) changes the score or
Hessian and can change a mode sign.  An additive action gauge is harmless
only when it is independent of every retained variable being
differentiated.
\end{firewall}

% =============================================================================
\section{A denominator-free conditional-rank test}
\label{sec:V23-rank-test}
% =============================================================================

The likelihood test requires a terminal Perron tilt.  A weaker but
denominator-free obstruction can be read from the temporal kernel alone.
For fixed \(W\), define the two-by-two conditional minor
\begin{equation}
 \begin{aligned}
 \mathfrak W_\alpha(x_1,x_2;W;z_1,z_2)
 :={}&
 c_\alpha(x_1;W,z_1)c_\alpha(x_2;W,z_2)\\
 &-
 c_\alpha(x_1;W,z_2)c_\alpha(x_2;W,z_1).
 \end{aligned}
 \label{eq:V23-kernel-wedge}
\end{equation}

\begin{theorem}[Conditional rank-one criterion for one static fibre law]
\label{thm:V23-conditional-rank}
Fix \(W\), and assume the fibre and initial-slice carriers are connected.
The identities
\begin{equation}
 p_{\alpha;x,W}=\nu_W^{\rm tr}
 \quad\text{for every }x
 \label{eq:V23-all-x-incidence}
\end{equation}
hold if and only if both of the following conditions hold:
\begin{enumerate}[label=\textup{(\roman*)},leftmargin=3.5em]
\item for one initial configuration \(x_0\),
      \(L_{\alpha;x_0,W}(z)\) is independent of \(z\);
\item every conditional minor vanishes,
      \begin{equation}
       \mathfrak W_\alpha(x_1,x_2;W;z_1,z_2)=0
       \quad\text{for all }x_1,x_2,z_1,z_2.
       \label{eq:V23-all-wedges-zero}
      \end{equation}
\end{enumerate}
Equivalently, at fixed \(W\) the family
\(z\mapsto c_\alpha(x;W,z)\) has rank one, and its unique positive fibre
profile is proportional to \(e^{-\mathcal S(z;W)}/f(W,z)\).
\end{theorem}

\begin{proof}
If \eqref{eq:V23-all-x-incidence} holds, then
\Cref{thm:V23-fibre-likelihood} gives
\[
 c_\alpha(x;W,z)
 =a_\alpha(x,W)
   \frac{e^{-\mathcal S(z;W)}}{f(W,z)}
\]
for a positive scalar \(a_\alpha(x,W)\).  One row has the required profile,
and every two-by-two minor vanishes.

Conversely, positivity and vanishing of all minors imply that every kernel
row is proportional to the \(x_0\) row.  Indeed, fix \(z_0\) and put
\[
 a(x,W)
 :=\frac{c_\alpha(x;W,z_0)}
         {c_\alpha(x_0;W,z_0)}.
\]
The minor with rows \(x,x_0\) and columns \(z,z_0\) gives
\(c_\alpha(x;W,z)=a(x,W)c_\alpha(x_0;W,z)\).  Condition \textup{(i)} says
that the \(x_0\) row times \(f e^{\mathcal S}\) is constant in \(z\);
hence the same holds for every row.  Apply
\eqref{eq:V23-incidence-equivalence}.
\end{proof}

\begin{corollary}[Full-terminal same-law incidence is impossible for the V10 convolution]
\label{cor:V23-full-rank-obstruction}
If the fibre is the entire terminal configuration space, then no single
positive static law can equal \(p_{\alpha,x}\) for every \(x\) for the V10
temporal Wilson convolution on a nontrivial link carrier.
\end{corollary}

\begin{proof}
Were such a law to exist, \Cref{thm:V23-conditional-rank} would make
\[
 c_\alpha(x,y)=a(x)b(y).
\]
The associated integral operator has rank one.  The V10 convolution is
injective on \(L^2(\mathcal X,\mathfrak m)\), because every Peter--Weyl
multiplier is strictly positive.  A rank-one operator on the
infinite-dimensional nontrivial link space is not injective, a
contradiction.
\end{proof}

\begin{assessment}[What must be read from an actual anisotropic slab]
\label{ass:V23-kernel-acquisition}
For a proper reduced fibre the full-rank contradiction need not apply: the
coarse variable \(W\) may retain the row dependence.  The exact finite
acquisition is nevertheless small and literal:
\begin{enumerate}[label=\textup{(\alph*)},leftmargin=3.5em]
\item one calibrated row likelihood
      \(L_{\alpha;x_0,W}\), and
\item the conditional minors \(\mathfrak W_\alpha\) on the declared source
      bank.
\end{enumerate}
A nonzero minor proves that the V35 law is not the physical conditional
posterior.  Vanishing minors plus the calibrated row proves incidence.  No
comparison of dimensions, stationary marginals, or static Hessians can
replace this test.
\end{assessment}

% =============================================================================
\section{Exact likelihood reconstruction of the posterior-excess card}
\label{sec:V23-likelihood-card}
% =============================================================================

Even when incidence fails, the V35 trace law is a valid reference measure.
At fixed \((\alpha,x,W)\), abbreviate
\[
 \nu=\nu_W^{\rm tr},\qquad
 L=L_{\alpha;x,W},\qquad
 \ell=\mathbb E_\nu L,
\]
and regard the temporal score \(J^{\rm t}\) and curvature \(K^{\rm t}\) as
functions of the fibre coordinate \(z\).  Define the unnormalized moments
\begin{align}
 B_L&:=\mathbb E_\nu[LJ^{\rm t}],
 \label{eq:V23-BL}\\
 C_L&:=\mathbb E_\nu[LJ^{\rm t}(J^{\rm t})^*],
 \label{eq:V23-CL}\\
 A_L&:=\mathbb E_\nu[LK^{\rm t}].
 \label{eq:V23-AL}
\end{align}

\begin{theorem}[Likelihood-reweighted temporal excess cylinder]
\label{thm:V23-likelihood-excess-card}
The Hermitian matrix
\begin{equation}
 \boxed{
 \mathbb N_D
 :=\frac12\mathbb E_{\nu^{\otimes2}}
 \left[
  L(z_1)L(z_2)
  \Delta_{12}J^{\rm t}(\Delta_{12}J^{\rm t})^*
 \right]
 -\ell\,\mathbb E_\nu[LK^{\rm t}],}
 \label{eq:V23-ND-pair}
\end{equation}
where
\(\Delta_{12}J^{\rm t}=J^{\rm t}(z_1)-J^{\rm t}(z_2)\), satisfies
\begin{align}
 \boxed{
 \mathbb N_D
 =\ell(C_L-A_L)-B_LB_L^*,}
 \label{eq:V23-ND-moments}\\
 \boxed{
 \mathbb N_D
 =\ell^2
 \left[
  \Cov_{p_{\alpha;x,W}}(J^{\rm t})
  -\mathbb E_{p_{\alpha;x,W}}K^{\rm t}
 \right]
 =\ell^2D_{\alpha;x,W}.}
 \label{eq:V23-ND-identity}
\end{align}
Thus \(\mathbb N_D\) and \(D_{\alpha;x,W}\) have the same inertia, kernel,
and positive or negative polarization vectors, and
\begin{equation}
 \boxed{
 (D_{\alpha;x,W})_+
 =\ell^{-2}(\mathbb N_D)_+.}
 \label{eq:V23-positive-part-rescaling}
\end{equation}
\end{theorem}

\begin{proof}
Independence of the two copies gives
\begin{align*}
 &\frac12\mathbb E_{\nu^{\otimes2}}
 \left[
  L_1L_2(J_1-J_2)(J_1-J_2)^*
 \right]\\
 &\qquad
 =\ell\,\mathbb E_\nu[LJJ^*]
  -\mathbb E_\nu[LJ]\mathbb E_\nu[LJ]^*
 =\ell C_L-B_LB_L^*.
\end{align*}
This proves \eqref{eq:V23-ND-moments}.  By
\eqref{eq:V23-RN-fibre},
\[
 \mathbb E_pK^{\rm t}=\ell^{-1}A_L,\qquad
 \Cov_p(J^{\rm t})
 =\ell^{-1}C_L-\ell^{-2}B_LB_L^*.
\]
Multiplication by \(\ell^2\) gives
\eqref{eq:V23-ND-identity}.  Since \(\ell^2>0\), inertia and spectral
subspaces are unchanged, and positive functional calculus commutes with
positive scalar multiplication.
\end{proof}

\begin{corollary}[Normalizer-aware ceiling transfer]
\label{cor:V23-ceiling-transfer}
If an outward finite card proves
\begin{equation}
 \mathbb N_D\preceq \widehat a\,I,
 \qquad
 0<\ell_-\le\ell,
 \label{eq:V23-unnormalized-ceiling-data}
\end{equation}
then
\begin{equation}
 \boxed{
 D_{\alpha;x,W}\preceq
 \frac{\widehat a_+}{\ell_-^2}I,}
 \qquad
 \widehat a_+:=\max\{\widehat a,0\}.
 \label{eq:V23-normalized-ceiling}
\end{equation}
The lower normalizer bound is indispensable for a magnitude statement,
although it is not needed for the inertia verdict.
\end{corollary}

\begin{proof}
Equation \eqref{eq:V23-ND-identity} gives
\(\ell^2D\preceq\widehat a I\).  Divide by the positive scalar
\(\ell^2\ge\ell_-^2\); replacing a negative \(\widehat a\) by zero only
weakens the upper bound.
\end{proof}

\begin{assessment}[A direct executable replacement for false incidence]
\label{ass:V23-likelihood-repair}
The complete finite packet
\[
 \boxed{(\ell,A_L,B_L,C_L)}
\]
reconstructs the V22 temporal posterior excess on the V35 trace history
without assuming the two laws agree.  It is the density-channel analogue of
the V35 unnormalized likelihood cards, but its target is new: the physical
Perron Witten repair rather than the static trace-shell Hessian.  If the
incidence residual vanishes, \(L\) is constant and the packet collapses to
the common-law formula.  If it does not, the likelihood is retained exactly;
no law is renamed.
\end{assessment}

% =============================================================================
\section{A local shifted factorization of the density defect}
\label{sec:V23-shifted-factorization}
% =============================================================================

Return to the complete physical cotangent bundle and the V21/V22 identity
\begin{equation}
 \Ric+\Hess W_\pi
 =\mathcal R_\eta-\Theta_\eta-2D_\alpha,
 \qquad
 \mathcal R_\eta\succeq\rho_{\beta,\eta}I,
 \qquad
 \Theta_\eta\succeq0.
 \label{eq:V23-actual-curvature-D}
\end{equation}
All tensors are required to be routed to the same complete retained tangent
before this equation is used.

\begin{theorem}[Uniform-ceiling shifted Witten split]
\label{thm:V23-shifted-Witten}
Assume that a scalar \(a\ge0\), independent of the configuration, satisfies
\begin{equation}
 D_\alpha(x)\preceq aI
 \quad\text{for every }x.
 \label{eq:V23-D-ceiling}
\end{equation}
Define
\begin{align}
 \mathcal R_{\eta,\pi}^{[a]}
 &:=\mathcal R_\eta+2(aI-D_\alpha),
 \label{eq:V23-R-shifted}\\
 \Theta_{\eta,\pi}^{[a]}
 &:=\Theta_\eta+2aI.
 \label{eq:V23-Theta-shifted}
\end{align}
Then
\begin{equation}
 \boxed{
 \Ric+\Hess W_\pi
 =\mathcal R_{\eta,\pi}^{[a]}
  -\Theta_{\eta,\pi}^{[a]},}
 \label{eq:V23-shifted-split}
\end{equation}
with
\begin{equation}
 \boxed{
 \mathcal R_{\eta,\pi}^{[a]}
 \succeq\mathcal R_\eta
 \succeq\rho_{\beta,\eta}I,
 \qquad
 \Theta_{\eta,\pi}^{[a]}\succeq0.}
 \label{eq:V23-shifted-positivity}
\end{equation}
If \(\Theta_\eta=\mathcal V_\eta^*\mathcal V_\eta\), an exact synthesis is
\begin{equation}
 \boxed{
 \mathcal V_{\eta,\pi}^{[a]}
 :=\bigl[\,\mathcal V_\eta\quad\sqrt{2a}\,I\,\bigr],
 \qquad
 (\mathcal V_{\eta,\pi}^{[a]})^*
 \mathcal V_{\eta,\pi}^{[a]}
 =\Theta_{\eta,\pi}^{[a]}.}
 \label{eq:V23-local-shifted-synthesis}
\end{equation}
The added column is diagonal in the original link/source coordinates and
does not require \((D_\alpha)_+^{1/2}\).
\end{theorem}

\begin{proof}
Assumption \eqref{eq:V23-D-ceiling} gives
\(aI-D_\alpha\succeq0\).  Add and subtract \(2aI\) in
\eqref{eq:V23-actual-curvature-D}:
\[
 \mathcal R_\eta-\Theta_\eta-2D_\alpha
 =
 [\mathcal R_\eta+2(aI-D_\alpha)]
 -[\Theta_\eta+2aI].
\]
This proves the split and both positivity statements.  The Gram of the
horizontal concatenation in \eqref{eq:V23-local-shifted-synthesis} is the
sum of the two column Grams.
\end{proof}

\begin{proposition}[Comparison with the minimal V21 repair]
\label{prop:V23-shift-vs-minimal}
Under \eqref{eq:V23-D-ceiling},
\begin{align}
 2(D_\alpha)_+&\preceq2aI,
 \label{eq:V23-defect-order}\\
 2(D_\alpha)_-&\preceq2(aI-D_\alpha).
 \label{eq:V23-comparison-order}
\end{align}
Thus the shifted split pays a larger but source-local defect and puts the
matching nonnegative excess into the comparison side.  It preserves the
exact curvature tensor and the original comparison floor.
\end{proposition}

\begin{proof}
The spectral inequality \(D_\alpha\preceq aI\) gives
\((D_\alpha)_+\preceq aI\).  Also
\[
 aI-D_\alpha
 =aI-(D_\alpha)_++(D_\alpha)_-
 \succeq(D_\alpha)_-.
\]
Multiply by two.
\end{proof}

\begin{corollary}[Complete-source Gram has no new remote carrier]
\label{cor:V23-source-Gram-local}
For every declared complete signed source projection
\(P_{\mathscr C}^{\rm src}\),
\begin{equation}
 \boxed{
 P_{\mathscr C}^{\rm src}
 (\mathcal V_{\eta,\pi}^{[a]})^*
 \mathcal V_{\eta,\pi}^{[a]}
 P_{\mathscr C}^{\rm src}
 =
 P_{\mathscr C}^{\rm src}\Theta_\eta P_{\mathscr C}^{\rm src}
 +2aP_{\mathscr C}^{\rm src}.}
 \label{eq:V23-source-Gram-shifted}
\end{equation}
The new source incidence is the identity on the already declared carrier;
its innovation outside that carrier is zero.
\end{corollary}

\begin{proof}
Compress \eqref{eq:V23-local-shifted-synthesis}.  Since the added column is a
scalar multiple of the identity, its compressed Gram is
\(2aP_{\mathscr C}^{\rm src}\), and its range on an input source is already
inside the same source carrier.
\end{proof}

\begin{corollary}[A conservative strict repaired edge]
\label{cor:V23-strict-edge}
Let
\begin{equation}
 \mathcal A_{\eta,\pi}^{[a]}
 :=\nabla_\pi^*\nabla+\mathcal R_{\eta,\pi}^{[a]},
 \qquad
 \mathcal K_{\eta,\pi}^{[a]}
 :=\mathcal V_{\eta,\pi}^{[a]}
    (\mathcal A_{\eta,\pi}^{[a]})^{-1/2}.
 \label{eq:V23-shifted-edge}
\end{equation}
If
\begin{equation}
 \boxed{
 \norm{\Theta_\eta}_{\op}+2a
 <\rho_{\beta,\eta},}
 \label{eq:V23-edge-sufficient-condition}
\end{equation}
then
\begin{equation}
 \boxed{
 \norm{\mathcal K_{\eta,\pi}^{[a]}}^2
 \le
 \frac{\norm{\Theta_\eta}_{\op}+2a}
      {\rho_{\beta,\eta}}
 <1.}
 \label{eq:V23-strict-edge-bound}
\end{equation}
\end{corollary}

\begin{proof}
By \eqref{eq:V23-shifted-positivity},
\(\mathcal A_{\eta,\pi}^{[a]}\succeq\rho_{\beta,\eta}I\).  Therefore
\begin{align*}
 \norm{\mathcal K_{\eta,\pi}^{[a]}}^2
 &=
 \norm{
  (\mathcal A_{\eta,\pi}^{[a]})^{-1/2}
  \Theta_{\eta,\pi}^{[a]}
  (\mathcal A_{\eta,\pi}^{[a]})^{-1/2}}\\
 &\le
 \rho_{\beta,\eta}^{-1}
 \norm{\Theta_\eta+2aI}_{\op}\\
 &\le
 \rho_{\beta,\eta}^{-1}
 (\norm{\Theta_\eta}_{\op}+2a).
\end{align*}
Use \eqref{eq:V23-edge-sufficient-condition}.
\end{proof}

\begin{firewall}[Sufficient global edge versus source-cyclic edge]
\label{fire:V23-edge-scope}
Condition \eqref{eq:V23-edge-sufficient-condition} is a conservative global
criterion.  Its failure does not prove that the source-cyclic edge is one:
the defect may land weakly on the declared source range even when its global
operator norm is large.  Conversely, a source-only strict edge must be
proved on the complete V15 carrier and cannot be inferred from an unrelated
static V35 mode sign.
\end{firewall}

% =============================================================================
\section{The V22 strong ceiling now yields a local physical column}
\label{sec:V23-strong-shift}
% =============================================================================

Retain the conditional complete-transfer incidence required by
\Cref{thm:V22-strong-density-bound}.  Define
\begin{equation}
 a_\beta
 :=\mathfrak B_2(\beta,0)
   +\frac{|\beta|}{2}C_M^{(2)}(0).
 \label{eq:V23-a-beta}
\end{equation}

\begin{theorem}[Strong-branch local shifted repair]
\label{thm:V23-strong-local-shift}
On the literally identified V9/V21 complete transfer branch,
\begin{equation}
 D_\alpha\preceq a_\beta I,
 \qquad
 a_\beta\le C_a|\beta|
 \quad(|\beta|\le\beta_0)
 \label{eq:V23-strong-ceiling}
\end{equation}
for a volume-independent \(C_a\).  Hence
\Cref{thm:V23-shifted-Witten} applies with \(a=a_\beta\) and its added
source column is the local multiplier
\[
 \sqrt{2a_\beta}\,I.
\]
Under the V21 anisotropic matching
\[
 \beta_{\alpha,j}=\tau_\alpha\gamma_{\alpha,j},
 \qquad
 \kappa_{H,j}=s_\alpha/\tau_\alpha,
\]
the new physical source stock satisfies
\begin{equation}
 \boxed{
 \kappa_{H,j}\,2a_{\beta_{\alpha,j}}
 \le2C_a s_\alpha|\gamma_{\alpha,j}|.}
 \label{eq:V23-strong-physical-local-column}
\end{equation}
It is uniformly bounded whenever
\(\sup_{\alpha,j}|\gamma_{\alpha,j}|/\alpha<\infty\).
\end{theorem}

\begin{proof}
The proof of \Cref{thm:V22-strong-density-bound} gives the stronger bound
\[
 \norm{D_\alpha}_{\op}
 \le
 \mathfrak B_2(\beta,0)
 +\frac{|\beta|}{2}C_M^{(2)}(0)
 =a_\beta.
\]
Thus \(D_\alpha\preceq a_\beta I\).  The same proof shows
\(\mathfrak B_2(\beta,0)=O(|\beta|)\), proving the second assertion.
Apply \Cref{thm:V23-shifted-Witten}.  Finally,
\[
 2\kappa_{H,j}a_{\beta_{\alpha,j}}
 \le2C_a\frac{s_\alpha}{\tau_\alpha}
          |\tau_\alpha\gamma_{\alpha,j}|
 =2C_as_\alpha|\gamma_{\alpha,j}|.
\]
Since \(\alpha s_\alpha\to3/2\), the stated coefficient-ratio bound is
sufficient.
\end{proof}

\begin{corollary}[Small-\(\beta\) global edge criterion]
\label{cor:V23-small-beta-edge}
Suppose, on the same branch, that for some constants
\(\rho_0>0\), \(C_\Theta<\infty\), and all sufficiently small
\(|\beta|\),
\begin{equation}
 \rho_{\beta,\eta}\ge\rho_0,
 \qquad
 \norm{\Theta_\eta}_{\op}\le C_\Theta|\beta|.
 \label{eq:V23-small-beta-stocks}
\end{equation}
Then for
\begin{equation}
 |\beta|<
 \frac{\rho_0}{C_\Theta+2C_a}
 \label{eq:V23-small-beta-threshold}
\end{equation}
the shifted repaired edge is strictly below one, uniformly in volume.
\end{corollary}

\begin{proof}
Equations \eqref{eq:V23-strong-ceiling} and
\eqref{eq:V23-small-beta-stocks} give
\[
 \norm{\Theta_\eta}+2a_\beta
 \le(C_\Theta+2C_a)|\beta|<\rho_0
 \le\rho_{\beta,\eta}.
\]
Apply \Cref{cor:V23-strict-edge}.
\end{proof}

\begin{assessment}[What is genuinely closed on the strong branch]
\label{ass:V23-strong-closure}
Conditional on literal equality of the complete V9 and V21 transfers and a
common tangent routing for the V35 tensor identity, the density repair now
has:
\begin{enumerate}[label=\textup{(\roman*)},leftmargin=3.5em]
\item a volume-uniform \(O(|\beta|)\) scalar ceiling;
\item a source-local identity-column factorization;
\item the same physical coefficient ratio as the Wilson channel; and
\item a strict global edge whenever the independently loaded comparison and
      defect stocks satisfy \eqref{eq:V23-small-beta-stocks}.
\end{enumerate}
What is not closed is the premise.  The supplied static V35 fibre does not
populate the temporal kernel likelihood or the command router.  V23 has made
that premise executable through \(\mathfrak c_{\rm inc}\) and the conditional
kernel minors rather than treating it as a semantic identification.
\end{assessment}

% =============================================================================
\section{V23 ledger and revised frontier}
\label{sec:V23-ledger}
% =============================================================================

\begin{theorem}[Integrated V23 statement]
\label{thm:V23-integrated}
Preserving V1--V22 and without repeating the V35 static-card proofs, V23
proves:
\begin{enumerate}[label=\textup{(V23.\arabic*)},leftmargin=5.5em]
\item an exact likelihood relating the actual temporal fibre posterior to
      the V35 trace-fibre law;
\item a zero chi-square residual equivalent to conditional-law incidence;
\item explicit first and second derivatives of the relative fibre
      normalizer;
\item a denominator-free conditional-rank test and a full-terminal rank-one
      obstruction for the injective V10 convolution;
\item an unnormalized two-copy card reconstructing the temporal posterior
      excess without assuming incidence;
\item an exact shifted Witten split replacing
      \((D_\alpha)_+^{1/2}\) by the local identity column
      \(\sqrt{2a}\,I\) whenever \(D_\alpha\preceq aI\);
\item an exact source-compressed Gram with no density innovation outside the
      declared carrier;
\item a conservative explicit sufficient condition for a strict repaired
      edge; and
\item a volume-uniform \(O(|\beta|)\) local shifted column on the
      conditionally identified V22 strong branch.
\end{enumerate}
\end{theorem}

\begin{proof}
Items \textup{(V23.1)--(V23.3)} are
\Cref{thm:V23-fibre-likelihood,cor:V23-normalizer-derivatives}.  Item
\textup{(V23.4)} is
\Cref{thm:V23-conditional-rank,cor:V23-full-rank-obstruction}.  Item
\textup{(V23.5)} is \Cref{thm:V23-likelihood-excess-card}.  Items
\textup{(V23.6)--(V23.8)} are
\Cref{thm:V23-shifted-Witten,cor:V23-source-Gram-local,
cor:V23-strict-edge}.  Item \textup{(V23.9)} is
\Cref{thm:V23-strong-local-shift}.
\end{proof}

\begingroup
\small
\begin{longtable}{>{\raggedright\arraybackslash}p{0.28\textwidth}
                  >{\raggedright\arraybackslash}p{0.16\textwidth}
                  >{\raggedright\arraybackslash}p{0.48\textwidth}}
\toprule
\textbf{V23 row}&\textbf{Status}&\textbf{Advance or remaining obstruction}\\
\midrule
\endfirsthead
\toprule
\textbf{V23 row}&\textbf{Status}&\textbf{Advance or remaining obstruction}\\
\midrule
\endhead
Static V35 law versus temporal posterior
& \MGpass
& Exact likelihood and zero-residual test; no unproved equality retained.\\[0.35em]
Fibre normalizer derivatives
& \MGpass
& Both first and second retained derivatives are explicit.\\[0.35em]
Conditional kernel incidence
& \MGpass
& Denominator-free rank minors plus one calibrated row are necessary and
  sufficient.\\[0.35em]
Full-terminal common conditional law
& \MGfail
& Would make the injective V10 convolution rank one.\\[0.35em]
Posterior-excess reconstruction
& \MGpass
& Exact four-entry likelihood packet \((\ell,A_L,B_L,C_L)\).\\[0.35em]
Positive-part locality
& \MGpass
& Replaced, under a scalar ceiling, by the local shifted identity column.\\[0.35em]
Strong-branch density carrier
& \MGpass
& Volume-uniform \(O(|\beta|)\) local column, conditional on complete-transfer
  incidence.\\[0.35em]
Actual anisotropic fibre kernel
& \MGopen
& No supplied file populates \(c_\alpha(x;W,z)\), its calibrated likelihood,
  or the conditional minors.\\[0.35em]
Common tangent/command router
& \MGopen
& The V35 static shell tangent and V21 temporal boundary command remain
  distinct until the same-history router is loaded.\\[0.35em]
Strict repaired source edge
& \MGopen
& A sufficient global test is proved; its physical stocks or a sharper
  source-cyclic test must be populated.\\[0.35em]
Cofinal coefficient trajectory
& \MGopen
& The ratio \(\gamma_{\alpha,j}/\alpha\) is still not fixed by the supplied
  anisotropic action.\\[0.35em]
Capacity, protected sectors, and OS limit
& \MGopen
& Complete incidence matrices, bad capacity, sector floors, and local
  noncollapse remain.\\[0.35em]
Full Yang--Mills mass gap
& \MGopen
& Transfer incidence and the common cofinal continuum packet remain.\\
\bottomrule
\end{longtable}
\endgroup

\begin{openproblem}[V24 mandate: populate the shifted source edge or return
its exact survivor]
\label{op:V24}
\begin{enumerate}[label=\textup{(V24.\arabic*)},leftmargin=5.5em]
\item Write the literal anisotropic one-step action in variables
      \((x;W,z)\).  Evaluate the likelihood
      \eqref{eq:V23-L-likelihood} and the conditional minors
      \eqref{eq:V23-kernel-wedge}; do not infer them from the V35 static
      partition.
\item Load the same-history command router between the V35 shell tangent and
      the V21 temporal boundary tangent.  If its positive incidence residual
      is nonzero, append the resulting innovation as a distinct complete
      source.
\item On the zero-incidence branch, use
      \(a_\beta\) and the local synthesis
      \eqref{eq:V23-local-shifted-synthesis} to populate the complete V20
      compressed Gram and source-incidence matrix.
\item Test the sufficient edge
      \eqref{eq:V23-edge-sufficient-condition}.  If it is inconclusive,
      compute the source-cyclic singular value of
      \(\mathcal K_{\eta,\pi}^{[a]}\) rather than enlarging the global norm
      estimate.
\item If that edge approaches one, return a normalized singular vector and
      classify its Wilson, shifted-density, incidence-innovation, return,
      and bad-capacity components with the V18--V20 alternatives.
\item Keep the V35 static trace-shell mode and the temporal transfer edge as
      distinct targets.  Do not use a sign of the former as a value of the
      latter.
\item At V25 perform the required SM/NS cross-program audit before choosing
      the next cofinal direction.
\end{enumerate}
\end{openproblem}

\begin{firewall}[End-of-V23 claim boundary]
V23 proves the exact reduced-fibre likelihood, its incidence and rank tests,
an exact posterior-excess card, and a source-local shifted factorization
under a scalar ceiling.  It also proves that the submitted static V35
partition does not itself supply the needed temporal incidence.  No actual
anisotropic fibre kernel, command router, cofinal coefficient law, strict
physical source edge, bad-set capacity, all-sector floor, or nontrivial
continuum OS state has yet been populated.  Therefore the full
Yang--Mills mass gap remains open.
\end{firewall}

\section*{Version V23 append-only record}
\addcontentsline{toc}{section}{Version V23 append-only record}
The complete V22 source was retained before this continuation.  Its
SHA--256 digest was
\begin{center}\scriptsize\ttfamily
2607a33d4fbbb86c10e9e483e5df96aeed79a24fc04f92e0ace5cb1dfa700e22.
\end{center}
On the next iteration, retain this full file, append before its unique active
document-closing command, increment the filename to
\texttt{\_V24.tex}, and remove only the superseded V23 file from this note
series.

% =============================================================================
% END APPENDED VERSION V23 MATERIAL
% =============================================================================

% =============================================================================
% BEGIN APPENDED VERSION V24 MATERIAL
% =============================================================================

\clearpage
\part{Version V24: corrected shifted-channel orientation, a finite source-edge pencil, and an exact survivor return}
\label{part:V24}

\section{Direction: turn the open source edge into a finite certificate}
\label{sec:V24-direction}

V23 reached the correct local idea---replace the nonsmooth density column
\((D_\alpha)_+^{1/2}\) by the identity column \(\sqrt{2a}\,I\) after a
uniform ceiling \(D_\alpha\preceq aI\)---but its displayed synthesis was
written in the opposite operator orientation from the V16--V21 Woodbury
convention.  That mismatch is harmless for the scalar norm estimate, but it
is not harmless for a source-cyclic computation: the channel and one-form
projections live on different Hilbert spaces.

V24 first repairs that typing explicitly.  It then replaces the instruction
``compute the source-cyclic edge'' by a finite, denominator-free acquisition
packet.  At every finite regulator the packet consists only of Gram matrices
of Krylov/cell words.  A positive Schur residual says exactly whether the
word space has stabilized.  On stabilization, one generalized eigenvalue is
the exact complete-source edge.  Thus the calculation has only two honest
outcomes: a strict edge, or a normalized channel survivor with exact Wilson,
shifted-density, incidence, return-depth, and bad-cell ledgers.  Under the
already isolated physical synthesis stock, a near-unit survivor lands as a
normalized physical soft one-form; without that stock it retains the V17
ultraviolet-escape alternative.

The supplied material still does not specify the literal anisotropic
one-step kernel or the same-history command router required to put numbers
into this packet.  V24 therefore does not pretend to evaluate an absent
kernel.  Its advance is to make the remaining source-edge test finite and
exact once those two data are loaded.

% =============================================================================
\section{Correction of the V23 shifted synthesis}
\label{sec:V24-orientation}
% =============================================================================

Let \(\mathscr H\) be the complete retained physical one-form space and let
\(\mathscr C_\eta\) be the Wilson defect-channel space.  Use the V17 and V21
orientation
\[
 \mathcal V_\eta:\mathscr C_\eta\longrightarrow\mathscr H,
 \qquad
 \Theta_\eta=\mathcal V_\eta\mathcal V_\eta^*.
\]
For \(a\ge0\), define
\begin{equation}
 \mathscr C^{[a]}:=\mathscr C_\eta\oplus\mathscr H,
 \qquad
 \mathcal V^{[a]}
 :=\begin{bmatrix}\mathcal V_\eta&\sqrt{2a}\,I_{\mathscr H}\end{bmatrix}
 :\mathscr C^{[a]}\longrightarrow\mathscr H.
 \label{eq:V24-shifted-synthesis}
\end{equation}

\begin{proposition}[Typed shifted factorization and channel]
\label{prop:V24-typed-shift}
With \(\mathcal R_{\eta,\pi}^{[a]}\) as in
\eqref{eq:V23-R-shifted}, put
\begin{equation}
 \mathcal A^{[a]}
 :=\nabla_\pi^*\nabla+\mathcal R_{\eta,\pi}^{[a]},
 \qquad
 L^{[a]}:=\mathcal A^{[a]}-\mathcal V^{[a]}(\mathcal V^{[a]})^*.
 \label{eq:V24-A-L}
\end{equation}
Then
\begin{align}
 \boxed{
 \mathcal V^{[a]}(\mathcal V^{[a]})^*
 =\Theta_\eta+2aI_{\mathscr H},}
 \label{eq:V24-typed-factor}\\
 L^{[a]}
 &=\nabla_\pi^*\nabla+\Ric+\Hess W_\pi.
 \label{eq:V24-typed-Witten}
\end{align}
Whenever \(L^{[a]}\succ0\) on the declared finite exact range, the positive
self-adjoint defect channel is
\begin{equation}
 \boxed{
 \mathcal K^{[a]}
 :=(\mathcal V^{[a]})^*(\mathcal A^{[a]})^{-1}
     \mathcal V^{[a]}:
     \mathscr C^{[a]}\longrightarrow\mathscr C^{[a]},
 \qquad 0\preceq\mathcal K^{[a]}\prec I.}
 \label{eq:V24-positive-channel}
\end{equation}
In the decomposition \(\mathscr C_\eta\oplus\mathscr H\), its exact block
matrix is
\begin{equation}
 \boxed{
 \mathcal K^{[a]}
 =\begin{bmatrix}
 \mathcal V_\eta^*(\mathcal A^{[a]})^{-1}\mathcal V_\eta
 &\sqrt{2a}\,\mathcal V_\eta^*(\mathcal A^{[a]})^{-1}\\
 \sqrt{2a}\,(\mathcal A^{[a]})^{-1}\mathcal V_\eta
 &2a(\mathcal A^{[a]})^{-1}
 \end{bmatrix}.}
 \label{eq:V24-channel-block}
\end{equation}
\end{proposition}

\begin{proof}
For \((u,h)\in\mathscr C_\eta\oplus\mathscr H\),
\[
 \mathcal V^{[a]}(u,h)=\mathcal V_\eta u+\sqrt{2a}\,h.
\]
Multiplication of this block row by its adjoint proves
\eqref{eq:V24-typed-factor}.  Substitution into
\eqref{eq:V23-shifted-split} proves \eqref{eq:V24-typed-Witten}.  Let
\[
 T^{[a]}:=(\mathcal A^{[a]})^{-1/2}\mathcal V^{[a]}.
\]
Then \(\mathcal K^{[a]}=(T^{[a]})^*T^{[a]}\succeq0\), while
\[
 I-T^{[a]}(T^{[a]})^*
 =(\mathcal A^{[a]})^{-1/2}L^{[a]}
  (\mathcal A^{[a]})^{-1/2}\succ0.
\]
The nonzero spectra of \((T^{[a]})^*T^{[a]}\) and
\(T^{[a]}(T^{[a]})^*\) agree, proving
\(\mathcal K^{[a]}\prec I\) on the active channel.  Direct block
multiplication gives \eqref{eq:V24-channel-block}.
\end{proof}

\begin{corollary}[Correct reading of the V23 sufficient edge]
\label{cor:V24-corrected-global-edge}
The rectangular map relevant to the V23 estimate is
\(T^{[a]}=(\mathcal A^{[a]})^{-1/2}\mathcal V^{[a]}\), not the
ill-typed product \(\mathcal V^{[a]}(\mathcal A^{[a]})^{-1/2}\).  Moreover,
\begin{equation}
 \boxed{
 \norm{\mathcal K^{[a]}}
 =\norm{T^{[a]}}^2
 =\norm{(\mathcal A^{[a]})^{-1/2}
          (\Theta_\eta+2aI)
          (\mathcal A^{[a]})^{-1/2}}
 \le\frac{\norm{\Theta_\eta}+2a}{\rho_{\beta,\eta}}.}
 \label{eq:V24-corrected-edge-bound}
\end{equation}
Thus the scalar conclusion of \Cref{cor:V23-strict-edge} remains valid after
replacing its rectangular notation by the positive channel
\eqref{eq:V24-positive-channel}.
\end{corollary}

\begin{proof}
The two norm equalities follow from
\(\norm{T}^2=\norm{T^*T}=\norm{TT^*}\).  The comparison floor
\(\mathcal A^{[a]}\succeq\rho_{\beta,\eta}I\) and
\eqref{eq:V24-typed-factor} give the inequality.
\end{proof}

\begin{proposition}[Correct one-form/source-channel incidence]
\label{prop:V24-source-incidence}
Let \(Q_{\rm src}\) be an orthogonal projection on \(\mathscr H\).  The
corresponding channel entrance is
\begin{equation}
 \mathcal Y^{[a]}
 :=(\mathcal V^{[a]})^*Q_{\rm src}:
 \Ran Q_{\rm src}\longrightarrow\mathscr C^{[a]}.
 \label{eq:V24-source-entrance}
\end{equation}
Its exact entrance Gram is
\begin{equation}
 \boxed{
 (\mathcal Y^{[a]})^*\mathcal Y^{[a]}
 =Q_{\rm src}\Theta_\eta Q_{\rm src}+2aQ_{\rm src}.}
 \label{eq:V24-source-entrance-Gram}
\end{equation}
The identity column therefore introduces no remote incidence at depth zero.
It does not imply that the source-cyclic carrier has no remote part, because
subsequent applications of \(\mathcal K^{[a]}\) contain
\((\mathcal A^{[a]})^{-1}\) and can propagate the entrance.
\end{proposition}

\begin{proof}
By \eqref{eq:V24-source-entrance},
\[
 (\mathcal Y^{[a]})^*\mathcal Y^{[a]}
 =Q_{\rm src}\mathcal V^{[a]}(\mathcal V^{[a]})^*Q_{\rm src}.
\]
Use \eqref{eq:V24-typed-factor}.  The last assertion follows directly from
the inverse in \eqref{eq:V24-positive-channel}.
\end{proof}

\begin{firewall}[Scope of the correction]
\label{fire:V24-orientation-scope}
Equations \eqref{eq:V23-shifted-split} and
\eqref{eq:V23-shifted-positivity} are unchanged.  What is superseded is only
the orientation in the last clause of \Cref{thm:V23-shifted-Witten}, the
source compression \eqref{eq:V23-source-Gram-shifted}, and the rectangular
notation \eqref{eq:V23-shifted-edge}.  Channel projections compress
\((\mathcal V^{[a]})^*\mathcal V^{[a]}\); physical one-form projections
compress \(\mathcal V^{[a]}(\mathcal V^{[a]})^*\).  These two operations
must not be interchanged.
\end{firewall}

% =============================================================================
\section{A finite Gram packet for the exact cell-saturated edge}
\label{sec:V24-Gram-packet}
% =============================================================================

Fix one finite regulator.  To avoid overloading the shifted notation, write
\(\mathscr C\) for the complete channel, \(K\) for its positive channel,
and \(Y:E\to\mathscr C\) for the complete signed source entrance.  All
incidence-innovation columns, when nonzero, are included in \(Y\) and in
\(\mathscr C\).  Let \(\{P_x:x\in\Gamma\}\) be the V16 cell projections.

Define an increasing sequence
\begin{align}
 \mathscr C_0&:=\Ran Y,
 \label{eq:V24-cell-recursion-zero}\\
 \mathscr C_{n+1}
 &:=\mathscr C_n+K\mathscr C_n
     +\sum_{x\in\Gamma}P_x\mathscr C_n.
 \label{eq:V24-cell-recursion}
\end{align}
Choose any finite synthesis
\(Z_n:E_n\to\mathscr C\) with \(\Ran Z_n=\mathscr C_n\), and set
\begin{align}
 G_n&:=Z_n^*Z_n,
 &H_n&:=Z_n^*KZ_n,
 &J_n&:=Z_n^*K^2Z_n,
 \label{eq:V24-GHJ}\\
 H_{x,n}&:=Z_n^*P_xZ_n.
 \label{eq:V24-cell-moment}
\end{align}
Let \(G_n^\dagger\) denote the Moore--Penrose inverse and define
\begin{align}
 R_{K,n}&:=J_n-H_nG_n^\dagger H_n,
 \label{eq:V24-K-residual}\\
 R_{x,n}&:=H_{x,n}-H_{x,n}G_n^\dagger H_{x,n}.
 \label{eq:V24-cell-residual}
\end{align}

\begin{theorem}[Exact stabilization residuals]
\label{thm:V24-stabilization}
For every \(n\),
\begin{equation}
 R_{K,n}\succeq0,
 \qquad
 R_{x,n}\succeq0\quad(x\in\Gamma).
 \label{eq:V24-residual-positive}
\end{equation}
Moreover,
\begin{align}
 R_{K,n}=0
 &\Longleftrightarrow K\mathscr C_n\subseteq\mathscr C_n,
 \label{eq:V24-K-invariance}\\
 R_{x,n}=0
 &\Longleftrightarrow P_x\mathscr C_n\subseteq\mathscr C_n.
 \label{eq:V24-cell-invariance}
\end{align}
Consequently the recursion stabilizes after at most \(\dim\mathscr C\)
strict dimension increases.  At the first \(N\) for which all residuals
vanish,
\begin{equation}
 \boxed{
 \mathscr C_N=\mathscr C_Y^{\rm cell},}
 \label{eq:V24-exact-cell-carrier}
\end{equation}
the V16 cell-saturated source carrier.
\end{theorem}

\begin{proof}
The operator
\begin{equation}
 \Pi_n:=Z_nG_n^\dagger Z_n^*
 \label{eq:V24-range-projection}
\end{equation}
is the orthogonal projection onto \(\Ran Z_n=\mathscr C_n\).  Therefore
\begin{align*}
 (KZ_n)^*(I-\Pi_n)KZ_n
 &=Z_n^*K^2Z_n
   -Z_n^*KZ_nG_n^\dagger Z_n^*KZ_n
 =R_{K,n},\\
 (P_xZ_n)^*(I-\Pi_n)P_xZ_n
 &=Z_n^*P_xZ_n
   -Z_n^*P_xZ_nG_n^\dagger Z_n^*P_xZ_n
 =R_{x,n},
\end{align*}
where \(K=K^*\) and \(P_x^2=P_x=P_x^*\) were used.  Both matrices are
Grams and hence positive.  Such a Gram vanishes exactly when the range of
the factor to its right lies in \(\Ran\Pi_n\), proving
\eqref{eq:V24-K-invariance}--\eqref{eq:V24-cell-invariance}.

If one residual is nonzero, \eqref{eq:V24-cell-recursion} adds a vector
outside \(\mathscr C_n\), so the dimension increases.  Finite dimension
forces termination.  At termination the space contains \(\Ran Y\) and is
invariant under \(K\) and every \(P_x\), hence contains the minimal V16
carrier.  Conversely, the recursion never leaves that minimal invariant
space.  Equality follows.
\end{proof}

\begin{corollary}[Ordinary source-cyclic shortcut]
\label{cor:V24-ordinary-cyclic}
If only the source-cyclic spectral edge is required, omit the
\(P_x\mathscr C_n\) terms in \eqref{eq:V24-cell-recursion} and test only
\(R_{K,n}\).  Equivalently, one may take the Krylov synthesis
\begin{equation}
 Z_n=\begin{bmatrix}Y&KY&\cdots&K^{n-1}Y\end{bmatrix}.
 \label{eq:V24-Krylov-synthesis}
\end{equation}
The cell residuals must be restored before the V16 weighted
Combes--Thomas argument is invoked.
\end{corollary}

\begin{proof}
The same proof as \Cref{thm:V24-stabilization}, with only the self-adjoint
operator \(K\), gives the ordinary cyclic carrier.  The final sentence is
exactly the distinction in \Cref{ass:V16-cell-saturation}.
\end{proof}

\begin{theorem}[Denominator-free generalized edge pencil]
\label{thm:V24-edge-pencil}
Suppose the residuals vanish at \(N\), and abbreviate
\(G=G_N\), \(H=H_N\).  Then
\begin{equation}
 \boxed{
 \rho_Y^{\rm cell}
 :=\norm{K|_{\mathscr C_Y^{\rm cell}}}
 =\max_{c\notin\ker G}
   \frac{\ip{c}{Hc}}{\ip{c}{Gc}}
 =\inf\{q\ge0:H\preceq qG\}.}
 \label{eq:V24-edge-pencil}
\end{equation}
In particular, for a proposed \(q<1\),
\begin{equation}
 \boxed{
 H\preceq qG}
 \label{eq:V24-denominator-free-edge}
\end{equation}
is equivalent to the strict source-carrier estimate
\(K|_{\mathscr C_Y^{\rm cell}}\preceq qI\).  The statement is independent
of the redundant columns chosen for \(Z_N\).
\end{theorem}

\begin{proof}
By \Cref{thm:V24-stabilization}, \(\mathscr C_Y^{\rm cell}\) reduces the
self-adjoint \(K\).  Every \(z\) in that carrier has the form \(z=Z_Nc\),
and
\[
 \norm z^2=\ip{c}{Gc},
 \qquad
 \ip{z}{Kz}=\ip{c}{Hc}.
\]
Taking the Rayleigh supremum proves the first equality.  For positive
semidefinite \(G\), the inequality of all Rayleigh quotients by \(q\) is
equivalent to \(H\preceq qG\); taking the least such \(q\) proves the second
equality.  Replacing \(Z_N\) by any other synthesis of the same range does
not change the operator Rayleigh supremum.
\end{proof}

\begin{assessment}[Literal acquisition packet]
\label{ass:V24-acquisition-packet}
At a finite regulator, the complete cell-saturated edge requires only
\begin{equation}
 \boxed{
 \bigl(G_n,H_n,J_n,(H_{x,n})_{x\in\Gamma}\bigr)}
 \label{eq:V24-acquisition-packet}
\end{equation}
until \eqref{eq:V24-K-residual} and
\eqref{eq:V24-cell-residual} vanish.  No global defect norm is substituted
for these matrices.  The anisotropic one-step kernel determines the actual
\(\mathcal A^{[a]}\), while the same-history router determines the complete
entrance \(Y\); without those two objects the symbols in
\eqref{eq:V24-acquisition-packet} cannot honestly be evaluated.
\end{assessment}

% =============================================================================
\section{Exact survivor extraction and component ledgers}
\label{sec:V24-survivor}
% =============================================================================

\begin{proposition}[Generalized eigenvector returns a channel survivor]
\label{prop:V24-generalized-survivor}
Let \(\rho=\rho_Y^{\rm cell}\).  There is a coefficient class
\(c\notin\ker G\) solving the generalized top-eigenvalue equation on
\((\ker G)^\perp\), and
\begin{equation}
 z:=\frac{Z_Nc}{\norm{Z_Nc}}
 \label{eq:V24-survivor-vector}
\end{equation}
obeys
\begin{equation}
 \boxed{
 z\in\mathscr C_Y^{\rm cell},
 \qquad \norm z=1,
 \qquad Kz=\rho z.}
 \label{eq:V24-survivor-eigen}
\end{equation}
\end{proposition}

\begin{proof}
The restriction of \(K\) to the finite-dimensional reducing carrier is
positive self-adjoint, so its norm is a top eigenvalue with a unit
eigenvector \(z\).  Surjectivity of \(Z_N:E_N\to\mathscr C_Y^{\rm cell}\)
gives \(z=Z_Nc\) for some coefficient class.  Pulling the eigenvalue
equation back to the quotient by \(\ker Z_N=\ker G\) gives the generalized
equation, and normalization gives \eqref{eq:V24-survivor-vector}.
\end{proof}

For an exact classification, let
\(\Pi_{\rm W},\Pi_a,\Pi_{\rm inc}\) be the declared mutually orthogonal
channel projections onto the Wilson, shifted identity, and router-innovation
columns, with sum \(I\) after any absent summand is set to zero.  Let
\(\Pi^{(n)}\) be the orthogonal projection onto \(\mathscr C_n\), put
\(\Pi^{(-1)}=0\), and define the return-shell projections
\begin{equation}
 Q_n:=\Pi^{(n)}-\Pi^{(n-1)},
 \qquad 0\le n\le N.
 \label{eq:V24-return-shells}
\end{equation}
For a declared bad cell set \(B\subseteq\Gamma\), write
\(P_B:=\sum_{x\in B}P_x\).

\begin{theorem}[Wilson/density/incidence, return, and bad-cell cards]
\label{thm:V24-survivor-cards}
For the unit survivor \(z\) of
\eqref{eq:V24-survivor-eigen}, define
\begin{align}
 w_{\rm W}&:=\norm{\Pi_{\rm W}z}^2,
 &w_a&:=\norm{\Pi_az}^2,
 &w_{\rm inc}&:=\norm{\Pi_{\rm inc}z}^2,
 \label{eq:V24-channel-card}\\
 r_n&:=\norm{Q_nz}^2,
 &b_B&:=\norm{P_Bz}^2.
 \label{eq:V24-return-bad-card}
\end{align}
Then
\begin{align}
 \boxed{w_{\rm W}+w_a+w_{\rm inc}=1,}
 \label{eq:V24-channel-sum}\\
 \boxed{\sum_{n=0}^Nr_n=1,}
 \label{eq:V24-return-sum}\\
 \boxed{b_B+\norm{(I-P_B)z}^2=1.}
 \label{eq:V24-bad-sum}
\end{align}
On the zero-incidence branch, \(w_{\rm inc}=0\).  The three boxed rows are
separate orthogonal ledgers; their entries must not be added across rows.
\end{theorem}

\begin{proof}
The first identity is Parseval for the declared orthogonal channel
decomposition.  Because the spaces \(\mathscr C_n\) are nested, the
operators \(Q_n\) are mutually orthogonal projections and
\(\sum_{n=0}^NQ_n=\Pi^{(N)}\).  Since \(z\in\mathscr C_N\), this proves the
second identity.  The cell projections are orthogonal and sum to the
identity, so \(P_B\) is an orthogonal projection; this proves the third.
\end{proof}

\begin{firewall}[Channel weights are not synthesis-energy weights]
\label{fire:V24-channel-vs-synthesis}
Writing \(z=(z_\eta,z_a,z_{\rm inc})\), the physical synthesis is
\[
 \mathcal Vz
 =\mathcal V_\eta z_\eta+\sqrt{2a}\,z_a
   +\mathcal V_{\rm inc}z_{\rm inc}.
\]
Its squared norm contains cross terms.  Therefore
\(w_{\rm W},w_a,w_{\rm inc}\) classify the channel survivor but are not
fractions of physical synthesis energy.  The latter is obtained from the
complete block Gram \(\mathcal V^*\mathcal V\), with all cross terms
retained.
\end{firewall}

% =============================================================================
\section{The cofinal edge-or-physical-soft-mode alternative}
\label{sec:V24-cofinal-alternative}
% =============================================================================

Consider a cofinal sequence of finite regulators.  For each \(j\), use the
complete stabilized carrier and write
\begin{equation}
 \rho_j
 :=\norm{K_j|_{\mathscr C_{Y_j}^{\rm cell}}}<1.
 \label{eq:V24-rho-j}
\end{equation}

\begin{theorem}[Strict edge or exact near-unit return]
\label{thm:V24-edge-or-return}
Exactly one of the following holds:
\begin{enumerate}[label=\textup{(\Alph*)},leftmargin=3.5em]
\item there is a \(\rho_*<1\) such that \(\rho_j\le\rho_*\) for all
      sufficiently large \(j\);
\item after passage to a subsequence, \(\rho_j\uparrow1\), and the
      construction \eqref{eq:V24-survivor-vector} returns unit vectors
      \(z_j\in\mathscr C_{Y_j}^{\rm cell}\) satisfying
      \(K_jz_j=\rho_jz_j\), together with all cards in
      \Cref{thm:V24-survivor-cards}.
\end{enumerate}
In case \textup{(A)}, the V16 exponential return bound follows after the
separate uniform conjugation-modulus inequality
\(\delta_{K_j}(\mu)<1-\rho_*\) is proved.  The edge alone does not supply
that locality inequality.
\end{theorem}

\begin{proof}
If \(\limsup_j\rho_j<1\), choose \(\rho_*\) strictly between the limsup and
one, proving \textup{(A)}.  Otherwise \(\limsup_j\rho_j=1\), and a
subsequence can be chosen increasing to one.  Apply
\Cref{prop:V24-generalized-survivor} at each regulator.  The last assertion
is \Cref{thm:V16-channel-CT} and its hypotheses.
\end{proof}

Return to the physical normalization of V19.  Let \(s_j>0\),
\begin{equation}
 \mathcal A_j^{\rm phys}=s_j\mathcal A_j,
 \qquad
 \mathcal V_j^{\rm phys}=\sqrt{s_j}\mathcal V_j,
 \qquad
 L_j^{\rm phys}=s_jL_j.
 \label{eq:V24-physical-scaling}
\end{equation}
The channel \(K_j\) is unchanged.  For the survivor in case
\textup{(B)}, define
\begin{equation}
 \omega_j^{\rm phys}
 :=(\mathcal A_j^{\rm phys})^{-1}
    \mathcal V_j^{\rm phys}z_j.
 \label{eq:V24-physical-landing}
\end{equation}

\begin{theorem}[Bounded physical source Gram prevents ultraviolet escape]
\label{thm:V24-no-UV-collapse}
Assume along the near-unit subsequence that
\begin{equation}
 d_*:=\sup_j
 s_j\norm{\mathcal V_jP_{\mathscr C_{Y_j}^{\rm cell}}}^2<\infty.
 \label{eq:V24-physical-carrier-stock}
\end{equation}
Then
\begin{align}
 \ip{\omega_j^{\rm phys}}
     {\mathcal A_j^{\rm phys}\omega_j^{\rm phys}}
 &=\rho_j,
 \label{eq:V24-comparison-landing}\\
 \ip{\omega_j^{\rm phys}}
     {L_j^{\rm phys}\omega_j^{\rm phys}}
 &=\rho_j(1-\rho_j),
 \label{eq:V24-soft-energy}\\
 \norm{\omega_j^{\rm phys}}^2
 &\ge\frac{\rho_j^2}{d_*}.
 \label{eq:V24-landing-lower-bound}
\end{align}
Consequently the normalized physical one-forms
\(\widehat\omega_j:=\omega_j^{\rm phys}/
\norm{\omega_j^{\rm phys}}\) satisfy
\begin{equation}
 \boxed{
 \ip{\widehat\omega_j}{L_j^{\rm phys}\widehat\omega_j}
 \le d_*\frac{1-\rho_j}{\rho_j}\longrightarrow0.}
 \label{eq:V24-normalized-soft-mode}
\end{equation}
Thus a near-unit complete-source edge plus
\eqref{eq:V24-physical-carrier-stock} is a genuine physical soft-mode
obstruction, not merely a channel-space warning.
\end{theorem}

\begin{proof}
The V17 landing ledger is invariant under the common scaling
\eqref{eq:V24-physical-scaling}.  Since \(K_jz_j=\rho_jz_j\), it gives
\eqref{eq:V24-comparison-landing} and
\eqref{eq:V24-soft-energy}.  Put
\(g_j=\mathcal V_j^{\rm phys}z_j\).  Then
\(\mathcal A_j^{\rm phys}\omega_j^{\rm phys}=g_j\) and
\[
 \rho_j
 =\ip{\omega_j^{\rm phys}}{g_j}
 \le\norm{\omega_j^{\rm phys}}\norm{g_j}
 \le\norm{\omega_j^{\rm phys}}\sqrt{d_*},
\]
because \(z_j\) is a unit vector in the declared carrier.  This proves
\eqref{eq:V24-landing-lower-bound}.  Divide
\eqref{eq:V24-soft-energy} by that lower bound to obtain
\eqref{eq:V24-normalized-soft-mode}.
\end{proof}

\begin{corollary}[Strong shifted branch has the required landing stock]
\label{cor:V24-strong-landing-stock}
On the conditionally matched strong branch of
\Cref{thm:V23-strong-local-shift}, use \(s_j=\kappa_{H,j}\) and
\(a=a_{\beta_j}\).  The already proved pure-Wilson estimate gives
\begin{equation}
 \boxed{
 s_j\norm{\mathcal V_j^{[a]}
 P_{\mathscr C_{Y_j}^{\rm cell}}}^2
 \le\kappa_{H,j}
      \bigl(192\alpha_{\beta_j}+2a_{\beta_j}\bigr).}
 \label{eq:V24-strong-stock}
\end{equation}
Hence a uniform bound for
\(\kappa_{H,j}\alpha_{\beta_j}\) together with the V23 coefficient-ratio
bound makes \eqref{eq:V24-physical-carrier-stock} automatic.  On that
branch, failure of a uniform complete-source edge returns the physical soft
sequence \eqref{eq:V24-normalized-soft-mode}.
\end{corollary}

\begin{proof}
The projection can only decrease the synthesis norm.  By
\eqref{eq:V24-typed-factor} and the V19 Wilson row bound,
\[
 \norm{\mathcal V_j^{[a]}}^2
 =\norm{\Theta_{\eta,j}+2a_{\beta_j}I}
 \le192\alpha_{\beta_j}+2a_{\beta_j}.
\]
Multiply by \(\kappa_{H,j}\).  The Wilson term is bounded by the first
hypothesis, and the shifted term is bounded by
\eqref{eq:V23-strong-physical-local-column}.
\end{proof}

\begin{firewall}[What remains if the physical stock is absent]
\label{fire:V24-stock-firewall}
Without \eqref{eq:V24-physical-carrier-stock}, the Cauchy lower bound
\eqref{eq:V24-landing-lower-bound} is unavailable.  A near-unit channel
vector can then lose ordinary physical norm through high comparison energy,
exactly as in \Cref{fire:V17-survivor-not-IR}.  It is not legitimate to call
such a vector a continuum massless state before the V17 low-energy landing
test or a cofinal physical carrier stock has been proved.
\end{firewall}

% =============================================================================
\section{V24 ledger and revised frontier}
\label{sec:V24-ledger}
% =============================================================================

\begin{theorem}[Integrated V24 statement]
\label{thm:V24-integrated}
Preserving V1--V23 except for the explicitly superseded V23 operator
orientation, V24 proves:
\begin{enumerate}[label=\textup{(V24.\arabic*)},leftmargin=5.5em]
\item the correctly typed shifted synthesis, positive Woodbury channel, and
      its exact Wilson/density block matrix;
\item the correct source entrance Gram and the distinction between local
      depth-zero incidence and inverse-generated remote return;
\item positive Schur residuals which detect exact stabilization of the
      finite cell-saturated source carrier;
\item a generalized Gram pencil whose largest eigenvalue is exactly the
      complete source-carrier edge;
\item an exact normalized survivor and separate channel, return-depth, and
      bad-cell component cards when the edge is not uniformly strict;
\item the cofinal strict-edge/near-unit-survivor alternative; and
\item a proof that the V19 physical carrier stock converts every near-unit
      complete-source survivor into a normalized physical soft one-form,
      with an explicit Rayleigh bound.
\end{enumerate}
\end{theorem}

\begin{proof}
Items \textup{(V24.1)--(V24.2)} are
\Cref{prop:V24-typed-shift,prop:V24-source-incidence}.  Items
\textup{(V24.3)--(V24.4)} are
\Cref{thm:V24-stabilization,thm:V24-edge-pencil}.  Item
\textup{(V24.5)} is
\Cref{prop:V24-generalized-survivor,thm:V24-survivor-cards}.  Items
\textup{(V24.6)--(V24.7)} are
\Cref{thm:V24-edge-or-return,thm:V24-no-UV-collapse}.
\end{proof}

\begingroup
\small
\begin{longtable}{>{\raggedright\arraybackslash}p{0.28\textwidth}
                  >{\raggedright\arraybackslash}p{0.16\textwidth}
                  >{\raggedright\arraybackslash}p{0.48\textwidth}}
\toprule
\textbf{V24 row}&\textbf{Status}&\textbf{Advance or remaining obstruction}\\
\midrule
Shifted synthesis orientation
& \MGclosed
& Correct block row \(\mathcal V:\mathscr C\to\mathscr H\), one-form
  factor \(\mathcal V\mathcal V^*\), and positive channel
  \(\mathcal V^*\mathcal A^{-1}\mathcal V\) are now explicit.\\[0.35em]
Finite source/cell carrier
& \MGclosed
& Schur residuals \eqref{eq:V24-K-residual} and
  \eqref{eq:V24-cell-residual} terminate exactly in finite dimension.\\[0.35em]
Exact source edge
& \MGreduced
& Reduced to the finite pencil \(H\preceq qG\); its numerical entries still
  require the actual temporal kernel and complete command router.\\[0.35em]
Near-unit survivor
& \MGclosed
& If the populated pencil approaches one, it returns a normalized channel
  eigenvector and exact component cards.\\[0.35em]
Physical survivor landing
& \MGclosed
& Under the physical carrier stock, the returned vector produces the
  normalized soft sequence \eqref{eq:V24-normalized-soft-mode}.\\[0.35em]
Literal temporal/static incidence
& \MGopen
& The attached static V35 fibre is not the absent anisotropic one-step
  temporal kernel; the router residual remains unevaluated.\\[0.35em]
Uniform edge and exponential return
& \MGopen
& Populate the Gram pencil and, on its strict branch, prove the uniform V16
  conjugation modulus.\\[0.35em]
Capacity, protected sectors, and OS limit
& \MGopen
& Complete bad capacity, all-sector floors, local noncollapse, and a
  nontrivial continuum OS state remain.\\[0.35em]
Full Yang--Mills mass gap
& \MGopen
& The finite edge test is executable, but its physical temporal inputs and
  the downstream cofinal packet are not yet supplied.\\
\bottomrule
\end{longtable}
\endgroup

\begin{openproblem}[V25 mandate: cross-program audit, then populate the
finite edge packet]
\label{op:V25}
\begin{enumerate}[label=\textup{(V25.\arabic*)},leftmargin=5.5em]
\item Before choosing the next direction, perform the user-required
      five-version audit of the current Standard-Model and Navier--Stokes
      notebooks in \texttt{latex\_notes}.  Record only transferable proved
      mechanisms, and keep their physical operators distinct from the
      Yang--Mills temporal transfer.
\item Search that audit and the attached source manuscripts for a literal
      anisotropic Yang--Mills one-step action in variables \((x;W,z)\) and a
      same-history tangent router.  If either is absent, state the exact
      missing datum rather than renaming a static fibre law.
\item If the router incidence residual is nonzero, add its orthogonal
      innovation channel before constructing \(Y\).  On the zero-incidence
      branch use \(\mathcal Y^{[a]}=(\mathcal V^{[a]})^*Q_{\rm src}\).
\item Populate \((G_n,H_n,J_n,H_{x,n})\) until the Schur residuals vanish,
      and evaluate the generalized edge \eqref{eq:V24-edge-pencil}.
\item On the strict branch, prove the V16 conjugation modulus and feed the
      exponential return row into the V15/V13 capacity packet.  On the
      near-unit branch, return \(z_j\), all three component ledgers, and the
      physical landing or ultraviolet-escape classification.
\item Audit the cofinal coefficient law
      \(\gamma_{\alpha,j}/\alpha\), protected-sector floors, and local
      noncollapse before any continuum mass-gap claim.
\end{enumerate}
\end{openproblem}

\begin{firewall}[End-of-V24 claim boundary]
V24 corrects the shifted-channel typing and proves a finite exact algorithm
for the complete cell-saturated source edge, including a rigorous survivor
and physical landing theorem.  It does not manufacture the absent
anisotropic temporal kernel or command router, and it does not populate a
uniform edge, locality modulus, bad capacity, protected-sector floor, or
continuum OS state.  Therefore the full Yang--Mills mass gap remains open.
\end{firewall}

\section*{Version V24 append-only record}
\addcontentsline{toc}{section}{Version V24 append-only record}
The complete V23 source was retained before this continuation.  Its
SHA--256 digest was
\begin{center}\scriptsize\ttfamily
4aef7f9fafeddae81af220c7a00b7d937f9e2d1277b0ebeed092f1fc3ea9d142.
\end{center}
On the next iteration, retain this full file, append before its unique active
document-closing command, increment the filename to
\texttt{\_V25.tex}, and remove only the superseded V24 file from this note
series.  V25 must begin with the required Standard-Model/Navier--Stokes
cross-program audit.

% =============================================================================
% END APPENDED VERSION V24 MATERIAL
% =============================================================================

% =============================================================================
% BEGIN APPENDED VERSION V25 MATERIAL
% =============================================================================

\clearpage
\part{Version V25: five-version companion audit, the literal anisotropic fibre kernel, and exact edge/locality pencils}
\label{part:V25}

\section{Scheduled SM/NS audit and the resulting change of direction}
\label{sec:V25-cross-audit}

This iteration begins with the cross-program audit required at every fifth
Yang--Mills version.  The audit was performed before the V25 direction was
chosen.  The two other files then retained in \texttt{latex\_notes} were
\begin{align*}
 &\texttt{Renewal\_SM\_Einstein\_Continuum\_Research\_Notes\_V42.tex},\\
 &\qquad\texttt{SHA256}=
 \texttt{1593E0E0D0CD245E19CA0CFB1F9547E5170E3A86CC7A1A7EEB8453B6962B78E2},\\
 &\texttt{Renewal\_NS\_Stability\_Research\_Notes\_V15.tex},\\
 &\qquad\texttt{SHA256}=
 \texttt{A3A73EBCB71715C7E91D52CA3ECE2DEACC6550BB9E28E6E5385F67FCA91E98CD}.
\end{align*}
These names and hashes freeze the companion states actually read.  No later
companion result is silently included in this audit.

\paragraph{Standard-Model/Einstein lesson.}
SM V42 constructs a local Hamiltonian boundary port, proves its exact Green
balance and invertibility on each of two declared characteristic cones, and
then exhibits explicit modes which destroy the global stable floor for both
orientation signs.  The type-compatible lesson is not an Einstein estimate:
finite rank, pointwise positivity, or success on separately selected sectors
does not prove a uniform global margin.  A hybrid survivor can lie between
the tested sectors.  V25 therefore supplements the V24 finite source-edge
test by an exact weighted-locality singular-value pencil and returns a
maximizing locality witness whenever the combined margin fails.

\paragraph{Navier--Stokes lesson.}
NS V15 proves a positive complete source--carrier work certificate only
after preserving the signed Duhamel identity, and proves that annular
components must be assembled before positive parts are taken.  It also proves
that a negative-order response has no uniform principal-order floor and that
fixed endpoints plus fixed pulled action do not control positive covariant
variation.  The type-compatible YM consequence is methodological: keep the
complete channel/physical cross term until an exact square completion has
been made, and do not infer the V16 conjugation modulus from the unweighted
edge.  V25 implements both points inside the YM Woodbury algebra.

\paragraph{Search for the missing YM inputs.}
Neither companion file defines the V10/V21 Yang--Mills temporal kernel in a
V35 section or a same-history map from the V35 reduced shell tangent to the
temporal boundary tangent.  The same-history ratio in the early SM finite
record card is explicitly an RPESM action-writer experiment and explicitly
not a Wilson-sign command.  It is therefore not imported.

\begin{firewall}[Scope of the V25 companion audit]
\label{fire:V25-cross-audit-scope}
No SM boundary mode is asserted to be a Yang--Mills channel vector, and no NS
source estimate is asserted for a gauge-field law.  The operative statements
below are proved from the V10--V24 Yang--Mills operators.  The companion audit
changes which exact residuals are retained; it supplies no Yang--Mills mass
gap.
\end{firewall}

The audit also exposes an avoidable upstream omission.  V10 already defines
the one-link temporal Wilson law and V21 already fixes the symmetric
transfer.  Composing their tensor-product kernel with the V35 section gives
the literal \((x;W,z)\) kernel requested in V23--V24.  It does not identify
that temporal posterior with the V35 static trace law; it makes the V23
likelihood test explicit.

% =============================================================================
\section{The literal V10/V21 anisotropic kernel in a V35 section}
\label{sec:V25-literal-kernel}
% =============================================================================

Let \(\mathbf G=\mathrm{SU}(3)\), let \(E\) be the finite rooted spatial-link
set used in V10, and write
\[
 \mathcal X=\mathbf G^E,\qquad
 \xi(g):=\frac13\operatorname{ReTr}g.
\]
With the V10 normalization, define the one-link density
\begin{equation}
 k_\alpha(g):=Z_\alpha^{-1}e^{\alpha\xi(g)}
 \label{eq:V25-one-link-density}
\end{equation}
with respect to normalized Haar measure.  For \(x,y\in\mathcal X\), put
\begin{equation}
 \boxed{
 c_\alpha(x,y)
 :=\prod_{e\in E}k_\alpha(x_ey_e^{-1}).}
 \label{eq:V25-product-kernel}
\end{equation}
Let
\begin{equation}
 G_\alpha:=\tau_\alpha V_B
 \label{eq:V25-spatial-exponent}
\end{equation}
be the dimensionless spatial magnetic exponent in the V10 symmetric step.

\begin{proposition}[Exact tensor-convolution kernel]
\label{prop:V25-product-kernel}
The V10 tensor convolution has the integral kernel
\eqref{eq:V25-product-kernel}:
\begin{equation}
 (\mathcal C_\alpha^{\otimes E}u)(x)
 =\int_{\mathcal X}c_\alpha(x,y)u(y)\,\dd\mathfrak m(y).
 \label{eq:V25-convolution-integral}
\end{equation}
Moreover,
\begin{equation}
 c_\alpha(x,y)>0,\qquad
 c_\alpha(x,y)=c_\alpha(y,x),\qquad
 \int_{\mathcal X}c_\alpha(x,y)\,\dd\mathfrak m(y)=1.
 \label{eq:V25-kernel-properties}
\end{equation}
The V10/V21 symmetric transfer therefore has kernel
\begin{equation}
 \boxed{
 t_{\alpha,G}(x,y)
 =e^{-G_\alpha(x)/2}c_\alpha(x,y)e^{-G_\alpha(y)/2}.}
 \label{eq:V25-symmetric-kernel}
\end{equation}
\end{proposition}

\begin{proof}
Tensoring the one-link convolution formulas gives
\eqref{eq:V25-convolution-integral}.  Positivity is immediate from
\eqref{eq:V25-one-link-density}.  Haar invariance gives row normalization.
Because \(\xi(g^{-1})=\xi(g)\),
\[
 k_\alpha(x_ey_e^{-1})
 =k_\alpha((x_ey_e^{-1})^{-1})
 =k_\alpha(y_ex_e^{-1}),
\]
so the product kernel is symmetric.  Multiplication of the convolution on
the left and right by \(e^{-G_\alpha/2}\) gives
\eqref{eq:V25-symmetric-kernel}.
\end{proof}

\begin{corollary}[Literal anisotropic one-step action]
\label{cor:V25-literal-action}
Relative to product Haar measure, the complete one-step action is
\begin{equation}
 \boxed{
 \mathscr A_{\alpha,G}(x,y)
 :=-\log t_{\alpha,G}(x,y)
 =\frac12G_\alpha(x)+\frac12G_\alpha(y)
  -\alpha\sum_{e\in E}\xi(x_ey_e^{-1})
  +|E|\log Z_\alpha.}
 \label{eq:V25-anisotropic-action}
\end{equation}
The temporal concentration is \(\alpha\), while the spatial coefficient and
physical clock are carried by \(G_\alpha=\tau_\alpha V_B\) and the V10
calibration \(\tau_\alpha=s_\alpha/\kappa\).  Thus
\eqref{eq:V25-anisotropic-action} is the literal anisotropic action of the
V10--V11 temporal-gauge tensor-product branch, not a new transfer inferred
from the static V35 partition.
\end{corollary}

\begin{proof}
Take minus the logarithm of \eqref{eq:V25-symmetric-kernel} and use
\eqref{eq:V25-one-link-density}.
\end{proof}

Now use the fixed V35 section already declared in V23,
\begin{equation}
 \Phi_W:Y_W^{\rm red}\longrightarrow\mathcal X,
 \qquad y=\Phi_W(z).
 \label{eq:V25-V35-section}
\end{equation}
To retain the measure bookkeeping, let \(j_\Phi(W,z)>0\) be the exact
conditional product-Haar density of this section relative to
\(\dd m_W(z)\).  In a measure-preserving product-Haar trivialization,
\(j_\Phi=1\).  Define
\begin{equation}
 \boxed{
 c_\alpha(x;W,z)
 :=c_\alpha(x,\Phi_W(z))\,j_\Phi(W,z).}
 \label{eq:V25-fibre-kernel}
\end{equation}

\begin{theorem}[Explicit V23 temporal fibre likelihood]
\label{thm:V25-explicit-likelihood}
Let \(h_\alpha>0\) be the Perron eigenfunction of the symmetric transfer
\eqref{eq:V25-symmetric-kernel}, and set
\begin{equation}
 f_\alpha(W,z)
 :=e^{-G_\alpha(\Phi_W(z))/2}h_\alpha(\Phi_W(z)).
 \label{eq:V25-terminal-tilt}
\end{equation}
Then the previously abstract V23 likelihood is the explicit function
\begin{align}
 L_{\alpha;x,W}(z)
 &=
 \left[
  \prod_{e\in E}
  k_\alpha\!\left(x_e\Phi_W(z)_e^{-1}\right)
 \right]
 j_\Phi(W,z)
 e^{-G_\alpha(\Phi_W(z))/2}
 h_\alpha(\Phi_W(z))
 e^{\mathcal S(z;W)},
 \label{eq:V25-explicit-L}\\
 \log L_{\alpha;x,W}(z)
 &=
 -|E|\log Z_\alpha
 +\alpha\sum_{e\in E}
  \xi\!\left(x_e\Phi_W(z)_e^{-1}\right)
 +\log j_\Phi(W,z) \notag\\
 &\quad
 -\frac12G_\alpha(\Phi_W(z))
 +\log h_\alpha(\Phi_W(z))
 +\mathcal S(z;W).
 \label{eq:V25-log-L}
\end{align}
Consequently all V23 normalizers, conditional minors, and the
chi-square incidence residual are now defined by explicit finite compact
group integrals on this branch.
\end{theorem}

\begin{proof}
Insert \eqref{eq:V25-fibre-kernel} and
\eqref{eq:V25-terminal-tilt} into the definition
\eqref{eq:V23-L-likelihood}.  This gives
\eqref{eq:V25-explicit-L}.  Taking logarithms and using
\eqref{eq:V25-one-link-density} gives \eqref{eq:V25-log-L}.  The final
assertion follows by substitution into
\Cref{eq:V23-R-normalizer,eq:V23-ell,eq:V23-kernel-wedge}.
\end{proof}

\begin{corollary}[Exact fibre-score incidence equation]
\label{cor:V25-fibre-score}
Let
\begin{equation}
 \Psi_{\alpha,x}(y)
 :=\alpha\sum_{e\in E}\xi(x_ey_e^{-1})
   -\frac12G_\alpha(y)+\log h_\alpha(y).
 \label{eq:V25-Psi}
\end{equation}
On every smooth V35 fibre,
\begin{equation}
 \boxed{
 \dd_z\log L_{\alpha;x,W}
 =(D_z\Phi_W)^*\dd_y\Psi_{\alpha,x}
   +\dd_z\log j_\Phi
   +\dd_z\mathcal S(\,\cdot\,;W).}
 \label{eq:V25-fibre-score}
\end{equation}
If the fibre is connected, temporal/static conditional incidence is
equivalent to vanishing of the right-hand side at every \(z\).
\end{corollary}

\begin{proof}
Differentiate \eqref{eq:V25-log-L} and apply the chain rule.  The incidence
equivalence is \eqref{eq:V23-fibre-score-zero}.
\end{proof}

\begin{firewall}[What the explicit kernel does and does not close]
\label{fire:V25-kernel-scope}
Equations \eqref{eq:V25-fibre-kernel}--\eqref{eq:V25-fibre-score} close the
missing literal-kernel row on the V10/V21 temporal-gauge branch.  They do not
make \(L_{\alpha;x,W}\) constant.  The Perron factor, section Jacobian,
spatial magnetic exponent, temporal coupling, and static fibre action all
remain in \eqref{eq:V25-log-L}.  Equality with the V35 trace law must still
be proved by the zero-score or chi-square test; it is not a consequence of
writing the kernel explicitly.
\end{firewall}

% =============================================================================
\section{The same-history shell-to-temporal tangent router}
\label{sec:V25-router}
% =============================================================================

Work on one smooth protected gauge stratum.  Let
\(P_y^{\rm hor}:T_y\mathcal X\to T_y^{\rm hor}\mathcal X\) be the horizontal
projection used by the V15/V21 physical tangent.  The V35 shell tangent
router is
\begin{equation}
 \boxed{
 \mathfrak R_{W,z}
 :=P_{\Phi_W(z)}^{\rm hor}D_z\Phi_W:
 T_zY_W^{\rm red}\longrightarrow
 T_{\Phi_W(z)}^{\rm hor}\mathcal X.}
 \label{eq:V25-section-router}
\end{equation}
If \(\Phi_W\) already takes values in the chosen horizontal slice, the
projection is the identity.

The temporal score is the initial covector
\begin{equation}
 J_{\alpha,x}^{\rm t}(y):=-\dd_x\log c_\alpha(x,y).
 \label{eq:V25-temporal-score}
\end{equation}
Its response to a V35 shell command at the same history is
\begin{equation}
 \boxed{
 \mathfrak U_{\alpha;x,W,z}
 :=D_z\!\left[J_{\alpha,x}^{\rm t}(\Phi_W(z))\right]
 =-\nabla_x\dd_y\log c_\alpha(x,y)\big|_{y=\Phi_W(z)}
   \circ\mathfrak R_{W,z}.}
 \label{eq:V25-mixed-router}
\end{equation}
The equality is understood with the fixed product connection and then
restricted to the protected horizontal tangent.

\begin{definition}[Positive router-incidence residual]
\label{def:V25-router-residual}
Let \(P_{\rm old}\) be the orthogonal projection onto the already declared
complete initial-covector source range at the same history.  Put
\begin{align}
 \mathfrak U_{\rm inc}
 &:=(I-P_{\rm old})\mathfrak U_{\alpha;x,W,z},
 \label{eq:V25-router-innovation}\\
 \boxed{
 \mathfrak E_{\rm inc}
 :=\mathfrak U_{\rm inc}^*\mathfrak U_{\rm inc}
 =\mathfrak U_{\alpha;x,W,z}^*
  (I-P_{\rm old})\mathfrak U_{\alpha;x,W,z}\succeq0.}
 \label{eq:V25-router-positive-residual}
\end{align}
\end{definition}

\begin{proposition}[Exact range-incidence test]
\label{prop:V25-router-incidence}
At every finite regulator,
\begin{equation}
 \boxed{
 \mathfrak E_{\rm inc}=0
 \iff
 \Ran\mathfrak U_{\alpha;x,W,z}\subseteq\Ran P_{\rm old}.}
 \label{eq:V25-router-incidence-equivalence}
\end{equation}
If the residual is nonzero, adjoining
\(\Ran\mathfrak U_{\rm inc}\) is the minimal orthogonal same-history
innovation required by this shell command.  If \(Q_{\rm src}^{\rm new}\)
projects onto
\[
 \Ran Q_{\rm src}^{\rm old}+\Ran\mathfrak U_{\rm inc},
\]
the corrected V24 entrance is
\begin{equation}
 \boxed{
 Y^{\rm new}=(\mathcal V^{[a]})^*Q_{\rm src}^{\rm new}.}
 \label{eq:V25-enlarged-entrance}
\end{equation}
\end{proposition}

\begin{proof}
The positive operator \(\mathfrak U_{\rm inc}^*\mathfrak U_{\rm inc}\)
vanishes exactly when \(\mathfrak U_{\rm inc}=0\).  By definition this is
equivalent to
\((I-P_{\rm old})\mathfrak U_{\alpha;x,W,z}=0\), which is the range
inclusion in \eqref{eq:V25-router-incidence-equivalence}.  Orthogonal
projection gives the minimal missing range.  Formula
\eqref{eq:V25-enlarged-entrance} is then
\eqref{eq:V24-source-entrance} with the enlarged physical source
projection.
\end{proof}

\begin{assessment}[The router is now literal but not yet numerically
populated]
\label{ass:V25-router-status}
Equations \eqref{eq:V25-section-router} and
\eqref{eq:V25-mixed-router} specify the same-history router uniquely from
the V35 section and the explicit temporal kernel.  Its residual is the
finite positive Gram \eqref{eq:V25-router-positive-residual}.  Evaluating
that Gram still requires the coordinate formula for the chosen V35 section,
its horizontal projection, and the declared old source range.  A nonzero
answer is retained as an innovation; it is not rounded to zero or absorbed
into the Wilson column.
\end{assessment}

% =============================================================================
\section{An exact source-deficit Schur completion}
\label{sec:V25-source-deficit}
% =============================================================================

Use the corrected V24 orientation
\[
 \mathcal A=\mathcal A^*\succ0,\qquad
 \mathcal V:\mathscr C\to\mathscr H,\qquad
 L=\mathcal A-\mathcal V\mathcal V^*\succ0,\qquad
 K=\mathcal V^*\mathcal A^{-1}\mathcal V.
\]
Let \(\mathscr C_Y\) be any finite reducing source carrier.  For
\(z\in\mathscr C_Y\) and \(u\in\operatorname{Dom}\mathcal A^{1/2}\), define
\begin{equation}
 \mathcal Q_z(u)
 :=\norm{z-\mathcal V^*u}_{\mathscr C}^2
   +\ip{u}{Lu}_{\mathscr H}.
 \label{eq:V25-Q-functional}
\end{equation}

\begin{theorem}[Complete source--carrier square completion]
\label{thm:V25-Schur-completion}
Let
\(\omega_z:=\mathcal A^{-1}\mathcal Vz\).  Then
\begin{equation}
 \boxed{
 \mathcal Q_z(u)
 =\norm{\mathcal A^{1/2}(u-\omega_z)}_{\mathscr H}^2
  +\ip{z}{(I-K)z}_{\mathscr C}.}
 \label{eq:V25-square-completion}
\end{equation}
Consequently,
\begin{align}
 \boxed{
 \min_u\mathcal Q_z(u)
 =\ip{z}{(I-K)z},}
 \label{eq:V25-deficit-minimum}\\
 \boxed{
 1-\norm{K|_{\mathscr C_Y}}
 =\inf_{\substack{z\in\mathscr C_Y,\ \norm z=1\\u}}
   \left(
    \norm{z-\mathcal V^*u}^2+\ip{u}{Lu}
   \right).}
 \label{eq:V25-primal-edge-deficit}
\end{align}
Thus the source-edge margin is exactly a joint
channel-incidence/physical-energy coercivity constant.
\end{theorem}

\begin{proof}
Expand \eqref{eq:V25-Q-functional} while keeping the cross term:
\begin{align*}
 \mathcal Q_z(u)
 &=\norm z^2
   -2\operatorname{Re}\ip{\mathcal Vz}{u}
   +\ip{u}{\mathcal V\mathcal V^*u}
   +\ip{u}{Lu}\\
 &=\norm z^2
   -2\operatorname{Re}\ip{\mathcal Vz}{u}
   +\ip{u}{\mathcal Au}.
\end{align*}
Completing the \(\mathcal A\)-square gives
\[
 \norm{\mathcal A^{1/2}(u-\omega_z)}^2
 +\norm z^2-\ip{\mathcal Vz}{\mathcal A^{-1}\mathcal Vz},
\]
which is \eqref{eq:V25-square-completion}.  The first term has the unique
minimum zero at \(u=\omega_z\), proving
\eqref{eq:V25-deficit-minimum}.  Finally, \(K\) is positive self-adjoint and
\(\mathscr C_Y\) reduces it, so
\[
 \inf_{\norm z=1}\ip{z}{(I-K)z}
 =1-\norm{K|_{\mathscr C_Y}}.
\]
This proves \eqref{eq:V25-primal-edge-deficit}.
\end{proof}

\begin{corollary}[Finite deficit pencil]
\label{cor:V25-deficit-pencil}
At the stabilized V24 synthesis \(Z:E_N\to\mathscr C_Y\), with
\(G=Z^*Z\) and \(H=Z^*KZ\),
\begin{equation}
 \boxed{
 \ip{c}{(G-H)c}
 =\min_u
 \left(
  \norm{Zc-\mathcal V^*u}^2+\ip{u}{Lu}
 \right).}
 \label{eq:V25-G-minus-H}
\end{equation}
For \(\varepsilon>0\),
\begin{equation}
 \boxed{
 G-H\succeq\varepsilon G}
 \label{eq:V25-deficit-pencil}
\end{equation}
is equivalent to the strict source edge
\(\norm{K|_{\mathscr C_Y}}\le1-\varepsilon\).
\end{corollary}

\begin{proof}
Apply \Cref{thm:V25-Schur-completion} to \(z=Zc\).  Then
\[
 \ip{Zc}{(I-K)Zc}=\ip{c}{(G-H)c}.
\]
The final equivalence is the Rayleigh characterization on the quotient by
\(\ker G\).
\end{proof}

\begin{proposition}[Exact deficit split for a returned survivor]
\label{prop:V25-deficit-split}
For every \(z\in\mathscr C_Y\),
\begin{equation}
 \boxed{
 \ip{z}{(I-K)z}
 =\norm{z-\mathcal V^*\omega_z}^2
  +\ip{\omega_z}{L\omega_z}.}
 \label{eq:V25-deficit-split}
\end{equation}
If \(Kz=\rho z\) and \(\norm z=1\), then
\begin{align}
 \norm{z-\mathcal V^*\omega_z}^2
 &=(1-\rho)^2,
 \label{eq:V25-incidence-deficit}\\
 \ip{\omega_z}{L\omega_z}
 &=\rho(1-\rho),
 \label{eq:V25-physical-deficit}\\
 (1-\rho)^2+\rho(1-\rho)&=1-\rho.
 \label{eq:V25-deficit-sum}
\end{align}
All three identities are unchanged by the V19 common physical scaling.
\end{proposition}

\begin{proof}
At the minimizer, \(\mathcal V^*\omega_z=Kz\).  Hence the right side of
\eqref{eq:V25-deficit-split} is
\[
 \ip{z}{(I-K)^2z}+\ip{z}{K(I-K)z}
 =\ip{z}{(I-K)z}.
\]
The eigenvector formulas follow immediately.  Under
\(\mathcal A^{\rm phys}=s\mathcal A\),
\(\mathcal V^{\rm phys}=\sqrt s\,\mathcal V\), and
\(\omega_z^{\rm phys}=s^{-1/2}\omega_z\), both
\((\mathcal V^{\rm phys})^*\omega_z^{\rm phys}=Kz\) and
\(\ip{\omega_z^{\rm phys}}{L^{\rm phys}\omega_z^{\rm phys}}\) are
unchanged.
\end{proof}

\begin{corollary}[Columnwise incidence residuals after complete assembly]
\label{cor:V25-column-deficits}
Suppose the complete channel is the orthogonal direct sum of Wilson,
shifted-density, and router-innovation channels and
\[
 \mathcal V=
 \begin{bmatrix}
  \mathcal V_{\rm W}&\sqrt{2a}\,I&\mathcal V_{\rm inc}
 \end{bmatrix}.
\]
Writing \(z=(z_{\rm W},z_a,z_{\rm inc})\), one has
\begin{align}
 \norm{z-\mathcal V^*\omega_z}^2
 &=
 \norm{z_{\rm W}-\mathcal V_{\rm W}^*\omega_z}^2
 +\norm{z_a-\sqrt{2a}\,\omega_z}^2 \notag\\
 &\quad
 +\norm{z_{\rm inc}-\mathcal V_{\rm inc}^*\omega_z}^2.
 \label{eq:V25-column-deficit-sum}
\end{align}
For a top eigenvector, adding
\(\ip{\omega_z}{L\omega_z}\) to the complete right-hand side gives exactly
\(1-\rho\).  No cross term is discarded: the physical-column cross terms
are already assembled inside
\(L=\mathcal A-\mathcal V\mathcal V^*\) before the channel square is split.
\end{corollary}

\begin{proof}
The three components of \(\mathcal V^*\omega_z\) lie in the declared
orthogonal channel summands, so Parseval gives
\eqref{eq:V25-column-deficit-sum}.  Apply
\Cref{prop:V25-deficit-split}.
\end{proof}

\begin{assessment}[Upstream alternative to a circular gap estimate]
\label{ass:V25-primal-route}
Equation \eqref{eq:V25-primal-edge-deficit} does not assume an \(L^2\) mass
floor for \(L\).  It says that a strict edge can be proved by a joint
inf--sup estimate: every unit complete-source channel vector must either
fail to be reproduced by \(\mathcal V^*u\), or every reproducing physical
one-form must pay positive \(L\)-energy.  Conversely a near-unit edge
returns both residuals explicitly through
\eqref{eq:V25-incidence-deficit}--\eqref{eq:V25-physical-deficit}.  This is
the YM square-completed analogue of preserving the complete NS
source--carrier pairing, proved here without importing an NS estimate.
\end{assessment}

% =============================================================================
\section{The V24 Gram packet also computes the exact locality modulus}
\label{sec:V25-locality-pencil}
% =============================================================================

Use the stabilized cell-saturated carrier from V24.  Thus
\[
 Z:E_N\to\mathscr C_Y^{\rm cell},\qquad
 G=Z^*Z,\qquad H=Z^*KZ,\qquad H_x=Z^*P_xZ,
\]
and the V24 stabilization residuals vanish.  Work on the coefficient
quotient
\begin{equation}
 \widehat E_N:=E_N/\ker G,\qquad
 \ip{[c]}{[d]}_G:=\ip{c}{Gd}.
 \label{eq:V25-coefficient-quotient}
\end{equation}

\begin{proposition}[Exact quotient representatives of \(K\) and the cells]
\label{prop:V25-quotient-representation}
Define on \(\widehat E_N\)
\begin{equation}
 B[c]:=[G^\dagger Hc],
 \qquad
 B_x[c]:=[G^\dagger H_xc].
 \label{eq:V25-B-Bx}
\end{equation}
These maps are well defined and satisfy
\begin{equation}
 KZ=ZG^\dagger H,\qquad
 P_xZ=ZG^\dagger H_x.
 \label{eq:V25-intertwining}
\end{equation}
In the \(G\)-metric, \(B\) is the positive self-adjoint representative of
\(K\), the \(B_x\)'s are mutually orthogonal projections, and
\(\sum_xB_x=I\).
\end{proposition}

\begin{proof}
If \(c\in\ker G\), then \(Zc=0\).  For every \(d\),
\[
 \ip{d}{Hc}=\ip{Zd}{KZc}=0,\qquad
 \ip{d}{H_xc}=\ip{Zd}{P_xZc}=0,
\]
so \(Hc=H_xc=0\); hence the quotient maps are well defined.
At stabilization, \(KZ\) and \(P_xZ\) have range in \(\Ran Z\).  Using the
range projection \(ZG^\dagger Z^*\) gives
\[
 KZ=ZG^\dagger Z^*KZ=ZG^\dagger H,
\]
and similarly for \(P_xZ\).  The remaining claims follow by transporting
the corresponding self-adjointness, positivity, orthogonality, and
resolution-of-identity properties through the surjection \(Z\).
\end{proof}

For a cell center \(y\in\Gamma\) and \(\mu\in\mathbb R\), define the
coefficient weight
\begin{equation}
 S_y(\mu):=
 \sum_{x\in\Gamma}
 e^{\mu\operatorname{dist}(x,y)}B_x.
 \label{eq:V25-coefficient-weight}
\end{equation}

\begin{lemma}[Exact finite coefficient conjugation]
\label{lem:V25-coefficient-conjugation}
The operator \(S_y(\mu)\) is invertible on \(\widehat E_N\), with
\[
 S_y(\mu)^{-1}=S_y(-\mu),
\]
and
\begin{equation}
 W_y(\mu)Z=ZS_y(\mu).
 \label{eq:V25-weight-intertwining}
\end{equation}
Consequently
\begin{equation}
 E_y(\mu)
 :=S_y(\mu)BS_y(-\mu)-B
 \label{eq:V25-E-coordinate}
\end{equation}
represents
\(W_y(\mu)KW_y(\mu)^{-1}-K\) exactly on the cell-saturated carrier.
\end{lemma}

\begin{proof}
The \(B_x\)'s are mutually orthogonal projections with sum one, so
\[
 S_y(\mu)S_y(-\mu)=\sum_xB_x=I.
\]
Equation \eqref{eq:V25-weight-intertwining} follows from
\eqref{eq:V25-intertwining} and the definition of the V16 weight.  Conjugate
the identity \(KZ=ZB\) by the two weights to obtain
\eqref{eq:V25-E-coordinate}.
\end{proof}

\begin{theorem}[Exact weighted-locality singular-value pencil]
\label{thm:V25-locality-pencil}
Put
\begin{equation}
 F_y(\mu):=E_y(\mu)^*G E_y(\mu)\succeq0.
 \label{eq:V25-F-locality}
\end{equation}
Then the V16 conjugation modulus on the exact cell-saturated source carrier
is
\begin{align}
 \boxed{
 \delta_K(\mu)^2
 =\max_{y\in\Gamma}
  \max_{c\notin\ker G}
  \frac{\ip{c}{F_y(\mu)c}}{\ip{c}{Gc}}}
 \label{eq:V25-delta-pencil}\\
 \boxed{
 =\max_{y\in\Gamma}
  \inf\{d^2\ge0:F_y(\mu)\preceq d^2G\}.}
 \label{eq:V25-denominator-free-delta}
\end{align}
Thus the same finite matrices \(G,H,(H_x)_x\) which populate the V24 edge
also populate the exact V16 locality modulus; no weighted global defect norm
or sampled cell direction is substituted.
\end{theorem}

\begin{proof}
By \Cref{lem:V25-coefficient-conjugation}, for \(z=Zc\),
\[
 \norm{
 [W_y(\mu)KW_y(\mu)^{-1}-K]z}^2
 =\ip{c}{E_y(\mu)^*GE_y(\mu)c},
\]
while \(\norm z^2=\ip{c}{Gc}\).  Taking the operator-norm Rayleigh
supremum and then the maximum over the finite set of centers gives
\eqref{eq:V25-delta-pencil}.  A generalized singular-value bound by \(d\)
is equivalent to the positive matrix inequality
\(F_y(\mu)\preceq d^2G\), proving
\eqref{eq:V25-denominator-free-delta}.
\end{proof}

\begin{corollary}[Fully finite Combes--Thomas certificate]
\label{cor:V25-finite-CT}
Let
\begin{equation}
 \rho=\inf\{q:H\preceq qG\},
 \qquad
 \delta(\mu)=
 \max_y\inf\{d\ge0:F_y(\mu)\preceq d^2G\}.
 \label{eq:V25-rho-delta}
\end{equation}
If
\begin{equation}
 \boxed{\rho+\delta(\mu)<1,}
 \label{eq:V25-finite-CT-margin}
\end{equation}
then
\begin{equation}
 \boxed{
 \norm{P_x(I-K)^{-1}P_y}
 \le
 \frac{e^{-\mu\operatorname{dist}(x,y)}}
      {1-\rho-\delta(\mu)}.}
 \label{eq:V25-finite-CT-bound}
\end{equation}
\end{corollary}

\begin{proof}
The first pencil in \eqref{eq:V25-rho-delta} is
\Cref{thm:V24-edge-pencil}; the second is
\Cref{thm:V25-locality-pencil}.  Condition
\eqref{eq:V25-finite-CT-margin} is exactly
\eqref{eq:V16-channel-CT-condition}.  Apply
\Cref{thm:V16-channel-CT}.
\end{proof}

\begin{proposition}[Infinitesimal locality card]
\label{prop:V25-commutator-card}
Define
\begin{equation}
 D_y:=\sum_x\operatorname{dist}(x,y)B_x,
 \qquad
 C_y:=D_yB-BD_y.
 \label{eq:V25-D-C}
\end{equation}
Then
\begin{equation}
 \boxed{
 \lim_{\mu\downarrow0}\frac{\delta_K(\mu)}{\mu}
 =
 \max_{y\in\Gamma}
 \max_{c\notin\ker G}
 \left(
 \frac{\ip{c}{C_y^*GC_yc}}{\ip{c}{Gc}}
 \right)^{1/2}.}
 \label{eq:V25-infinitesimal-locality}
\end{equation}
More generally,
\begin{equation}
 E_y(\mu)
 =\int_0^\mu
   S_y(t)C_yS_y(-t)\,\dd t.
 \label{eq:V25-conjugation-Duhamel}
\end{equation}
Hence, if for some \(\mu_0>0\)
\begin{equation}
 \Xi_{\mu_0}
 :=\max_{\substack{y\in\Gamma\\0\le t\le\mu_0}}
 \norm{S_y(t)C_yS_y(-t)}_G<\infty,
 \label{eq:V25-Xi}
\end{equation}
then
\begin{equation}
 \delta_K(\mu)\le\mu\Xi_{\mu_0}
 \qquad(0\le\mu\le\mu_0).
 \label{eq:V25-linear-delta}
\end{equation}
\end{proposition}

\begin{proof}
Since
\(S_y(\mu)=e^{\mu D_y}\) on the coefficient quotient,
\[
 \frac{\dd}{\dd t}
 \left(S_y(t)BS_y(-t)\right)
 =S_y(t)[D_y,B]S_y(-t).
\]
Integration proves \eqref{eq:V25-conjugation-Duhamel}.  Division by \(\mu\)
and finite-dimensional norm continuity give the derivative at zero.
Because \(\Gamma\) is finite at each regulator, the maximum commutes with
the one-sided limit, proving
\eqref{eq:V25-infinitesimal-locality}.  Taking norms under the integral
gives \eqref{eq:V25-linear-delta}.
\end{proof}

\begin{corollary}[Uniform cofinal edge and commutator stock give a positive
decay exponent]
\label{cor:V25-cofinal-CT}
Suppose on a cofinal family that
\begin{equation}
 \rho_j\le\rho_*<1,\qquad
 \Xi_{j,\mu_0}\le\Xi_*<\infty
 \label{eq:V25-cofinal-stocks}
\end{equation}
for fixed \(\mu_0>0\).  If \(\Xi_*>0\), choose
\begin{equation}
 0<\mu_*
 \le\min\left\{
 \mu_0,\frac{1-\rho_*}{2\Xi_*}
 \right\}.
 \label{eq:V25-mu-star}
\end{equation}
Then, uniformly in the regulator,
\begin{equation}
 \boxed{
 \norm{P_x(I-K_j)^{-1}P_y}
 \le
 \frac{2e^{-\mu_*\operatorname{dist}(x,y)}}{1-\rho_*}.}
 \label{eq:V25-uniform-return}
\end{equation}
If \(\Xi_*=0\), any \(0<\mu_*\le\mu_0\) is admissible.
\end{corollary}

\begin{proof}
Equation \eqref{eq:V25-linear-delta} and
\eqref{eq:V25-mu-star} give
\(\delta_{K_j}(\mu_*)\le(1-\rho_*)/2\).  Apply
\Cref{cor:V25-finite-CT}.
\end{proof}

\begin{firewall}[Unweighted edge and locality remain distinct]
\label{fire:V25-edge-locality}
The matrices \(H\) and \(F_y(\mu)\) measure different operators.  A strict
unweighted pencil \(H\preceq\rho G\) does not imply
\(F_y(\mu)\preceq d^2G\) with a regulator-uniform \(d\), just as a
negative-order response does not control a principal-order norm on an
unbounded frequency range.  Conversely a good locality pencil does not
force \(\rho<1\).  Both rows of
\eqref{eq:V25-finite-CT-margin} must be populated on the same complete
cell-saturated carrier.
\end{firewall}

% =============================================================================
\section{Cofinal strictness or exact spectral/locality witnesses}
\label{sec:V25-witnesses}
% =============================================================================

Fix a proposed \(\mu>0\) and a cofinal sequence of stabilized finite
packets.  Write \(\rho_j\) and \(\delta_j(\mu)\) for the two exact
generalized singular values.

\begin{theorem}[Three-way complete-source return]
\label{thm:V25-three-way-return}
After passage to a subsequence, one of the following holds.
\begin{enumerate}[label=\textup{(\Alph*)},leftmargin=3.5em]
\item There is an \(\varepsilon_*>0\) such that
      \[
       1-\rho_j-\delta_j(\mu)\ge\varepsilon_*.
      \]
      Then the uniform exponential return
      \eqref{eq:V25-finite-CT-bound} holds.
\item \(\rho_j\to1\).  Then the V24 generalized top eigenvectors \(z_j\)
      satisfy the exact source/physical deficit cards
      \eqref{eq:V25-incidence-deficit}--\eqref{eq:V25-deficit-sum}; under
      the V24 physical carrier stock they yield the normalized physical soft
      sequence \eqref{eq:V24-normalized-soft-mode}.
\item There is a \(\rho_*<1\) with \(\rho_j\le\rho_*\), but
      \[
       \liminf_j\bigl(1-\rho_j-\delta_j(\mu)\bigr)\le0.
      \]
      For each \(j\), some center \(y_j\) and unit
      \(z_j^{\rm loc}\in\mathscr C_{Y_j}^{\rm cell}\) attain
      \begin{equation}
       \boxed{
       \norm{
       [W_{y_j}(\mu)K_jW_{y_j}(\mu)^{-1}-K_j]
       z_j^{\rm loc}}
       =\delta_j(\mu).}
       \label{eq:V25-locality-witness}
      \end{equation}
      This is an exact weighted-locality witness, not a massless physical
      state.
\end{enumerate}
\end{theorem}

\begin{proof}
If the combined margins have positive liminf, pass to a tail and obtain
\textup{(A)}.  Otherwise choose a subsequence on which their liminf is
nonpositive.  If \(\rho_j\) has a subsequence tending to one, apply
\Cref{thm:V24-edge-or-return} and
\Cref{prop:V25-deficit-split}, giving \textup{(B)}.  If not, pass to a
further subsequence with \(\rho_j\le\rho_*<1\).  At every finite regulator
the maximum over \(y\) and the largest generalized singular value in
\eqref{eq:V25-delta-pencil} are attained.  Push a maximizing coefficient
through \(Z_j\) and normalize to obtain
\eqref{eq:V25-locality-witness}.  This is \textup{(C)}.
\end{proof}

\begin{assessment}[The companion-audit hybrid is now an internal YM
certificate]
\label{ass:V25-hybrid-lesson}
The SM V42 warning is now implemented without importing its boundary
symbol.  Success of every finite edge pencil alone can coexist with a
cofinal spectral survivor, and success of the unweighted edge can coexist
with a weighted-locality witness.  The exact alternatives in
\Cref{thm:V25-three-way-return} prevent either failure from being hidden by
sector sampling.  Case \textup{(C)} instructs the next iteration to go
upstream to the section, inverse, or cell commutator which created
\(F_{y_j}(\mu)\), rather than weakening the edge.
\end{assessment}

% =============================================================================
\section{V25 ledger and revised frontier}
\label{sec:V25-ledger}
% =============================================================================

\begin{theorem}[Integrated V25 statement]
\label{thm:V25-integrated}
Preserving V1--V24, V25 proves:
\begin{enumerate}[label=\textup{(V25.\arabic*)},leftmargin=5.5em]
\item the required audit of the exact current SM V42 and NS V15 files, with
      type firewalls and frozen hashes;
\item the literal positive symmetric V10/V21 tensor-product temporal kernel
      and anisotropic one-step action;
\item its exact disintegration in a V35 section, including the Haar
      Jacobian, terminal Perron tilt, explicit likelihood, and fibre-score
      incidence equation;
\item the same-history shell-to-temporal tangent router and its minimal
      positive innovation Gram;
\item an exact square completion identifying the source-edge deficit with a
      joint channel-incidence/physical-energy coercivity constant;
\item a coefficient-quotient realization of the complete V24 cell carrier;
\item a generalized singular-value pencil which computes the exact V16
      weighted locality modulus from the same finite Gram packet;
\item a finite Combes--Thomas certificate and a cofinal commutator-stock
      sufficient condition; and
\item an exhaustive strict-return/spectral-survivor/locality-witness
      alternative at every proposed decay exponent.
\end{enumerate}
\end{theorem}

\begin{proof}
Item \textup{(V25.1)} is \Cref{sec:V25-cross-audit}.  Items
\textup{(V25.2)--(V25.3)} are
\Cref{prop:V25-product-kernel,cor:V25-literal-action,
thm:V25-explicit-likelihood,cor:V25-fibre-score}.  Item
\textup{(V25.4)} is
\Cref{prop:V25-router-incidence}.  Item \textup{(V25.5)} is
\Cref{thm:V25-Schur-completion,prop:V25-deficit-split}.  Items
\textup{(V25.6)--(V25.8)} are
\Cref{prop:V25-quotient-representation,thm:V25-locality-pencil,
cor:V25-finite-CT,cor:V25-cofinal-CT}.  Item \textup{(V25.9)} is
\Cref{thm:V25-three-way-return}.
\end{proof}

\begingroup
\small
\begin{longtable}{>{\raggedright\arraybackslash}p{0.28\textwidth}
                  >{\raggedright\arraybackslash}p{0.16\textwidth}
                  >{\raggedright\arraybackslash}p{0.48\textwidth}}
\toprule
\textbf{V25 row}&\textbf{Status}&\textbf{Advance or remaining obstruction}\\
\midrule
Five-version SM/NS audit
& \MGclosed
& Exact V42/V15 files and hashes were read before direction choice; only
  type-compatible proof obligations were retained.\\[0.35em]
Literal anisotropic one-step kernel
& \MGclosed
& The V10/V21 tensor convolution is now written in \((x;W,z)\) variables,
  with section Jacobian and terminal Perron tilt.\\[0.35em]
Temporal/static fibre incidence
& \MGreduced
& Reduced to the explicit score equation \eqref{eq:V25-fibre-score} or the
  V23 chi-square card; constancy has not been proved.\\[0.35em]
Same-history tangent router
& \MGreduced
& Exact derivative and positive innovation Gram are defined; the chosen V35
  coordinate section and old source projection must be evaluated.\\[0.35em]
Source-edge deficit
& \MGclosed
& Exact Schur completion and finite pencil \(G-H\) now retain the complete
  source--carrier cross term.\\[0.35em]
Weighted locality acquisition
& \MGclosed
& The exact locality modulus and its maximizing witness are finite pencils
  derived from \(G,H,H_x\).\\[0.35em]
Uniform cofinal edge/locality margins
& \MGopen
& The explicit matrices must be bounded uniformly on the response-matched
  coefficient trajectory, or their exact witnesses returned.\\[0.35em]
Capacity, protected sectors, and OS limit
& \MGopen
& Bad capacity, toron/wrapping floors, local noncollapse, and the nontrivial
  continuum OS state remain.\\[0.35em]
Full Yang--Mills mass gap
& \MGopen
& The temporal kernel and finite decision packets are now literal, but their
  uniform cofinal margins and downstream continuum packet are not proved.\\
\bottomrule
\end{longtable}
\endgroup

\begin{openproblem}[V26 mandate: evaluate the explicit fibre score and the
first complete packet]
\label{op:V26}
\begin{enumerate}[label=\textup{(V26.\arabic*)},leftmargin=5.5em]
\item In the actual V35 reduced-shell coordinates, write
      \(\Phi_W\), \(j_\Phi\), and \(D_z\Phi_W\) on one protected local
      chart.  Insert them into \eqref{eq:V25-fibre-score}.
\item Differentiate the finite Perron equation for \(h_\alpha\) only after
      the complete kernel \eqref{eq:V25-symmetric-kernel} is fixed.  Bound
      or return the nonconstant part of the explicit fibre score; do not set
      it to zero by static/temporal renaming.
\item Evaluate \(\mathfrak E_{\rm inc}\).  If it is nonzero, enlarge the
      complete physical source projection and entrance as in
      \eqref{eq:V25-enlarged-entrance}.
\item On that exact entrance, populate the first stabilized
      \(G,H,J,H_x\) packet.  Test both \(G-H\succeq\varepsilon G\) and
      \(F_y(\mu)\preceq d^2G\).
\item If the spectral deficit collapses, return the V24/V25 physical
      survivor.  If only the locality margin collapses, return
      \eqref{eq:V25-locality-witness} and inspect the responsible cell pair
      before altering the comparison.
\item On the strict branch, feed \eqref{eq:V25-uniform-return} into the
      V20 incidence--capacity matrix, then audit the toron, wrapping,
      protected-sector, and local-algebra noncollapse rows.
\item Keep the explicit cofinal ratio
      \(\gamma_{\alpha,j}/\alpha\) in every physical stock.  A fixed-volume
      positive pencil is not a continuum mass.
\item Repeat the companion-file audit at V30.
\end{enumerate}
\end{openproblem}

\begin{firewall}[End-of-V25 claim boundary]
V25 performs the scheduled companion audit, writes the literal V10/V21
anisotropic transfer in V35 fibre coordinates, defines the exact
same-history innovation, and proves finite exact source-deficit and
weighted-locality pencils.  It does not prove that the explicit fibre score
vanishes, populate a regulator-uniform edge/locality margin, control bad
capacity and protected sectors, or construct a nontrivial continuum OS
state.  Therefore the full Yang--Mills mass gap remains open.
\end{firewall}

\section*{Version V25 append-only record}
\addcontentsline{toc}{section}{Version V25 append-only record}
The complete V24 source was retained before this continuation.  Its
SHA--256 digest was
\begin{center}\scriptsize\ttfamily
4f7a6726d3074e03bd890cbaad756312d83bb35dd4fcf42630aa786e7d562f32.
\end{center}
On the next iteration, retain this full file, append before its unique active
document-closing command, increment the filename to
\texttt{\_V26.tex}, and remove only the superseded V25 file from this note
series.  The next companion audit is V30.

% =============================================================================
% END APPENDED VERSION V25 MATERIAL
% =============================================================================

% =============================================================================
% BEGIN APPENDED VERSION V26 MATERIAL
% =============================================================================

\clearpage
\part{Version V26: the exact parity-tree router, Perron elimination, and a forced temporal/static incidence innovation}
\label{part:V26}

\section{Direction and correction after reading the literal V35 section}
\label{sec:V26-direction}

V25 wrote the V10/V21 temporal kernel in an abstract V35 section
\(\Phi_W\), retained a possible section Jacobian, and used
\(D_z\Phi_W\) when discussing the shell-to-temporal tangent router.  The
attached V35 manuscript makes two of those rows more precise.

First, its parity-tree section is a literal coordinate insertion with
product Haar measure on the reduced links.  The conditional Jacobian in
V25 is therefore exactly one on that declared card.  Second,
\(D_z\Phi_W\) routes a \emph{fibre-detail} tangent.  The retained physical
shell tangent is a coarse \(W\)-variation.  On the same partition function
V35 replaces the raw section derivative by a Haar-compensated harmonic
lift which simultaneously moves the reduced links.  That lift is the
physical same-history command router and already has a volume-uniform norm
bound.

V26 corrects the typing, inserts the explicit section into the V25
likelihood, and differentiates the Perron equation so that
\(\nabla\log h_\alpha\) is eliminated in favour of the already declared
convolution output \(r_\alpha\).  A two-initial-history subtraction then
cancels every static, magnetic, Perron, and Jacobian term.  The surviving
one-link Wilson score proves that the temporal posterior cannot equal the
same V35 trace law for all initial histories.  In fact an open set of
histories carries a strictly positive fibre-incidence Gram.  Thus the
zero-incidence shortcut is globally unavailable; the returned likelihood
innovation must be routed through the complete source/capacity packet.

The V35 source used here is the same file already frozen in the V23 audit:
\begin{center}\scriptsize\ttfamily
Renewal\_Yang\_Mills\_Reduced\_Fibre\_Wilson\_Shell\_Symbol\_and\_Local\_Vacuum\_Programme\_V35.tex,
\end{center}
with SHA--256
\begin{center}\scriptsize\ttfamily
7c6056d8149d3015d469b57474a75b23c88a6fe9f64fc81c7dca85c361c647a3.
\end{center}
Only its proved mathematical cards are used; its internal continuation
instructions do not direct this notebook.

% =============================================================================
\section{The literal parity-tree section and its two differentials}
\label{sec:V26-parity-section}
% =============================================================================

Use the V35 fine side \(N=2M\), \(M\ge4\).  Let
\(\mathscr T_N\) be the fixed parity tree, let
\(\mathscr E_N^{\rm red}\) be the \(45M^4\) reduced links, and set
\begin{equation}
 \mathscr C_N
 :=\left\{
 (2y+\widehat e_\mu,\,2y+2\widehat e_\mu):
 y\in(\mathbb Z/M\mathbb Z)^4,\ 1\le\mu\le4
 \right\}.
 \label{eq:V26-constituent-links}
\end{equation}
The sets \(\mathscr T_N,\mathscr C_N,\mathscr E_N^{\rm red}\) are the
three disjoint link classes in the V35 section.  Its coordinate map is
\begin{equation}
 \boxed{
 \Phi_W(z)_e=
 \begin{cases}
 I,&e\in\mathscr T_N,\\
 W_{y,\mu},
 &e=(2y+\widehat e_\mu,\,2y+2\widehat e_\mu)
   \in\mathscr C_N,\\
 z_e,&e\in\mathscr E_N^{\rm red}.
 \end{cases}}
 \label{eq:V26-explicit-section}
\end{equation}

\begin{proposition}[Exact fibre Jacobian and detail differential]
\label{prop:V26-section-differential}
Relative to the V35 declared product Haar reference
\[
 \dd m_{\rm red}(z)=
 \bigotimes_{e\in\mathscr E_N^{\rm red}}\dd z_e,
\]
the density \(j_\Phi\) in \eqref{eq:V25-fibre-kernel} is
\begin{equation}
 \boxed{j_\Phi(W,z)=1.}
 \label{eq:V26-unit-Jacobian}
\end{equation}
In left-trivialized coordinates, the fibre-detail differential is
\begin{equation}
 \boxed{
 (D_z\Phi_W\,\zeta)_e=
 \begin{cases}
 0,&e\in\mathscr T_N\cup\mathscr C_N,\\
 \zeta_e,&e\in\mathscr E_N^{\rm red}.
 \end{cases}}
 \label{eq:V26-detail-inclusion}
\end{equation}
The raw coarse-section differential in a coarse tangent
\(a=(a_{y,\mu})\) and fixed colour \(X\in\mathfrak{su}(3)\) is instead
\begin{equation}
 \boxed{
 (D_W\Phi\,aX)_e=
 \begin{cases}
 a_{y,\mu}X,&
 e=(2y+\widehat e_\mu,\,2y+2\widehat e_\mu)\in\mathscr C_N,\\
 0,&e\in\mathscr T_N\cup\mathscr E_N^{\rm red}.
 \end{cases}}
 \label{eq:V26-raw-shell-inclusion}
\end{equation}
Thus \(D_z\Phi_W\) and \(D_W\Phi\) have different domains, meanings, and
ranges.
\end{proposition}

\begin{proof}
Equation \eqref{eq:V26-explicit-section} uses the reduced link matrices
themselves as coordinates and V35 defines its fibre integral with their
product Haar measure.  There is no coordinate density to insert, proving
\eqref{eq:V26-unit-Jacobian}.  Differentiating the three cases of
\eqref{eq:V26-explicit-section} proves
\eqref{eq:V26-detail-inclusion} and
\eqref{eq:V26-raw-shell-inclusion}.
\end{proof}

The raw shell inclusion is not the best same-law representative of the
physical coarse derivative.  Retain the V35 scalar plaquette-incidence
blocks
\[
 \mathsf A=\tfrac13D_R^*D_R,\qquad
 \mathsf B=\tfrac13D_R^*D_C,\qquad
 \mathsf L_{\rm red}:=-\mathsf A^{-1}\mathsf B.
\]
For a real coarse field \(a\), define
\begin{equation}
 \boxed{
 b[a]_e=
 \begin{cases}
 0,&e\in\mathscr T_N,\\
 a_{y,\mu},&
 e=(2y+\widehat e_\mu,\,2y+2\widehat e_\mu)\in\mathscr C_N,\\
 (\mathsf L_{\rm red}a)_e,&e\in\mathscr E_N^{\rm red}.
 \end{cases}}
 \label{eq:V26-harmonic-lift}
\end{equation}

\begin{proposition}[Exact Haar-compensated shell router]
\label{prop:V26-harmonic-router}
For fixed unit colour \(X\), define
\begin{equation}
 W_{y,\mu}(t)=e^{t a_{y,\mu}X}W_{y,\mu},
 \qquad
 z_e(t)=e^{t(\mathsf L_{\rm red}a)_eX}z_e
 \quad(e\in\mathscr E_N^{\rm red}).
 \label{eq:V26-compensated-path}
\end{equation}
The map
\begin{equation}
 \boxed{
 \mathfrak R_{\rm harm}:aX\longmapsto b[a]X}
 \label{eq:V26-harmonic-router}
\end{equation}
is the literal same-history shell-to-fine-link tangent of the compensated
path.  The change \(z\mapsto z(t)\) preserves product Haar measure for every
\(t\).  Moreover the V35 estimates give
\begin{align}
 \norm{b[a]}_2^2&\le141\norm a_2^2,
 \label{eq:V26-router-general-bound}\\
 \norm{b[a]}_2^2&\le121\norm a_2^2
 \quad\text{on one transverse axial polarization.}
 \label{eq:V26-router-axial-bound}
\end{align}
Since \(b[a]\) contains \(a\) on the disjoint constituent links,
\begin{equation}
 \boxed{
 \norm a_2^2\le\norm{b[a]}_2^2.}
 \label{eq:V26-router-lower-bound}
\end{equation}
\end{proposition}

\begin{proof}
Differentiating \eqref{eq:V26-compensated-path} at zero and using
\eqref{eq:V26-explicit-section} gives
\eqref{eq:V26-harmonic-lift}.  Each reduced-link transformation is a left
translation, so product Haar measure is preserved.  The V35 Schur
minimizer satisfies
\[
 \norm{\mathsf L_{\rm red}a}_2^2\le140\norm a_2^2
\]
in general and at most \(120\norm a_2^2\) on the selected axial
polarization.  The constituent and reduced supports are disjoint, so
\[
 \norm{b[a]}_2^2
 =\norm a_2^2+\norm{\mathsf L_{\rm red}a}_2^2.
\]
This proves all three bounds.
\end{proof}

\begin{firewall}[Correction of the V25 router typing]
\label{fire:V26-router-correction}
Equation \eqref{eq:V25-section-router} remains a valid
\emph{fibre-detail-to-terminal} map after
\eqref{eq:V26-detail-inclusion} and horizontal projection are inserted.
It is not the retained shell-command router.  Wherever the V25 mixed
response is used for a coarse shell command, its last factor must be
\(\mathfrak R_{\rm harm}\) from
\eqref{eq:V26-harmonic-router}, not \(D_z\Phi_W\).  This correction leaves
the V25 fibre-score equation unchanged and makes the shell router
volume-uniform on the selected axial card.
\end{firewall}

% =============================================================================
\section{Eliminating the Perron derivative from the explicit likelihood}
\label{sec:V26-Perron-elimination}
% =============================================================================

Return to the V22 notation
\[
 r_\alpha=\mathcal C_\alpha f_\alpha,\qquad
 \mathcal T_{G_\alpha}h_\alpha=\lambda_\alpha h_\alpha,\qquad
 f_\alpha=e^{-G_\alpha/2}h_\alpha.
\]

\begin{theorem}[Exact Perron elimination]
\label{thm:V26-Perron-elimination}
The finite Perron equation is equivalent to
\begin{align}
 \boxed{
 f_\alpha(y)
 =\lambda_\alpha^{-1}e^{-G_\alpha(y)}r_\alpha(y),}
 \label{eq:V26-f-via-r}\\
 \boxed{
 \mathcal C_\alpha
 \bigl(e^{-G_\alpha}r_\alpha\bigr)
 =\lambda_\alpha r_\alpha.}
 \label{eq:V26-r-eigen-equation}
\end{align}
On the parity-tree fibre, the V25 likelihood becomes
\begin{equation}
 \boxed{
 L_{\alpha;x,W}(z)
 =\lambda_\alpha^{-1}
  c_\alpha(x,\Phi_W(z))
  e^{-G_\alpha(\Phi_W(z))}
  r_\alpha(\Phi_W(z))
  e^{\mathcal S(z;W)}.}
 \label{eq:V26-L-without-h}
\end{equation}
Consequently
\begin{equation}
 \boxed{
 \dd_z\log L_{\alpha;x,W}
 =(D_z\Phi_W)^*
   \dd_y\!\left[
    \log c_\alpha(x,y)-G_\alpha(y)+\log r_\alpha(y)
   \right]
 +\dd_z\mathcal S(\,\cdot\,;W).}
 \label{eq:V26-score-without-h}
\end{equation}
No derivative of \(h_\alpha\) or section Jacobian remains.
\end{theorem}

\begin{proof}
The Perron equation in V22 is
\[
 \lambda_\alpha h_\alpha
 =e^{-G_\alpha/2}r_\alpha.
\]
Multiplying by \(e^{-G_\alpha/2}\) proves
\eqref{eq:V26-f-via-r}.  Substitution in
\(r_\alpha=\mathcal C_\alpha f_\alpha\) gives
\eqref{eq:V26-r-eigen-equation}.  Insert
\eqref{eq:V26-f-via-r} and the unit Jacobian
\eqref{eq:V26-unit-Jacobian} into
\eqref{eq:V25-explicit-L}; this is
\eqref{eq:V26-L-without-h}.  Logarithmic differentiation and the chain rule
give \eqref{eq:V26-score-without-h}.
\end{proof}

For a reduced link \(r\in\mathscr E_N^{\rm red}\), a colour
\(X\in\mathfrak{su}(3)\), and the left-translation path
\(z_r(s)=e^{sX}z_r\), write \(D_{r,X}\) for differentiation at zero.

\begin{corollary}[Literal one-link fibre score]
\label{cor:V26-one-link-score}
For \(-\Tr X^2=1\), the \(r,X\) component of
\eqref{eq:V26-score-without-h} is
\begin{equation}
 \boxed{
 \begin{aligned}
 \sigma_{\alpha;x,W}^{r,X}(z)
 &:=
 D_{r,X}\log L_{\alpha;x,W}(z)\\
 &=
 -\frac{\alpha}{3}
   \operatorname{ReTr}(x_rz_r^{-1}X)
 -D_{r,X}G_\alpha(\Phi_W(z))\\
 &\quad
 +D_{r,X}\log r_\alpha(\Phi_W(z))
 +D_{r,X}\mathcal S(z;W).
 \end{aligned}}
 \label{eq:V26-explicit-fibre-score}
\end{equation}
\end{corollary}

\begin{proof}
Only the \(r\)-factor in \eqref{eq:V25-product-kernel} varies.  Since
\[
 x_r(e^{sX}z_r)^{-1}=x_rz_r^{-1}e^{-sX},
\]
one has
\[
 \left.\frac{\dd}{\dd s}\right|_{s=0}
 \alpha\xi(x_rz_r^{-1}e^{-sX})
 =-\frac{\alpha}{3}\operatorname{ReTr}(x_rz_r^{-1}X).
\]
Insert this in \eqref{eq:V26-score-without-h}.
\end{proof}

The same calculation along the compensated shell path gives the literal
coarse-command score.

\begin{corollary}[Compensated shell score]
\label{cor:V26-shell-score}
Let \(D_{b[a]X}\) denote the terminal derivative along
\eqref{eq:V26-compensated-path}, including the induced derivative of the
static V35 action.  Then
\begin{equation}
 \boxed{
 \begin{aligned}
 \sigma_{\alpha;x,W}^{\rm sh}[a,X](z)
 &:=
 \left.\frac{\dd}{\dd t}\right|_{t=0}
 \log L_{\alpha;x,W(t)}(z(t))\\
 &=
 -\frac{\alpha}{3}\sum_e b[a]_e
  \operatorname{ReTr}\!\left(
   x_e\Phi_W(z)_e^{-1}X
  \right)\\
 &\quad
 -D_{b[a]X}G_\alpha
 +D_{b[a]X}\log r_\alpha
 +D_{b[a]X}^{\rm comp}\mathcal S .
 \end{aligned}}
 \label{eq:V26-explicit-shell-score}
\end{equation}
This is the V25 mixed temporal response with the corrected router
\(\mathfrak R_{\rm harm}\).
\end{corollary}

\begin{proof}
Differentiate \eqref{eq:V26-L-without-h} along
\eqref{eq:V26-compensated-path}.  The product-kernel derivative is the sum
of the one-link calculation in
\Cref{cor:V26-one-link-score}; the remaining three terms are ordinary
directional derivatives along the same path.
\end{proof}

% =============================================================================
\section{A two-history card forces nonzero temporal/static incidence}
\label{sec:V26-two-history}
% =============================================================================

The terminal terms in \eqref{eq:V26-explicit-fibre-score} are difficult
because they contain the physical Perron factor and the complete static
action.  They are independent of the initial history \(x\).  Subtracting
two initial histories therefore removes them exactly.

\begin{theorem}[Exact one-link two-history score difference]
\label{thm:V26-two-history-score}
Fix \(W\), a fibre point \(z^0\), a reduced link
\(r\in\mathscr E_N^{\rm red}\), and
\(X\in\mathfrak{su}(3)\) with \(-\Tr X^2=1\).  Let
\(y^0=\Phi_W(z^0)\).  Choose two initial configurations which agree off
\(r\) and satisfy
\begin{equation}
 x_r^{(0)}=z_r^0,\qquad
 x_r^{(t)}=e^{tX}z_r^0.
 \label{eq:V26-two-histories}
\end{equation}
Then
\begin{equation}
 \boxed{
 \sigma_{\alpha;x^{(t)},W}^{r,X}(z^0)
 -\sigma_{\alpha;x^{(0)},W}^{r,X}(z^0)
 =
 -\frac{\alpha}{3}
  \operatorname{ReTr}(e^{tX}X).}
 \label{eq:V26-two-history-score-difference}
\end{equation}
The function
\[
 q_X(t):=\operatorname{ReTr}(e^{tX}X)
\]
satisfies
\begin{equation}
 q_X(0)=0,\qquad q_X'(0)=\operatorname{ReTr}(X^2)=-1.
 \label{eq:V26-q-derivative}
\end{equation}
Hence the right side of
\eqref{eq:V26-two-history-score-difference} is nonzero for every
sufficiently small nonzero \(t\).
\end{theorem}

\begin{proof}
In \eqref{eq:V26-explicit-fibre-score}, the magnetic, Perron, and static
terms depend on \(W,z^0\) but not on \(x\), so they cancel.  At the zeroth
history,
\[
 x_r^{(0)}(z_r^0)^{-1}=I,\qquad
 \operatorname{ReTr}X=0.
\]
At the second history,
\(
 x_r^{(t)}(z_r^0)^{-1}=e^{tX}
\).
This proves \eqref{eq:V26-two-history-score-difference}.  Differentiating
\(q_X\) at zero and using the normalization of \(X\) proves
\eqref{eq:V26-q-derivative}; continuity gives the last statement.
\end{proof}

\begin{corollary}[The two conditional posteriors are different]
\label{cor:V26-two-posteriors-different}
For the histories \eqref{eq:V26-two-histories},
\begin{equation}
 \boxed{
 \frac{\dd p_{\alpha;x^{(t)},W}}
      {\dd p_{\alpha;x^{(0)},W}}(z)
 =
 C_{\alpha,t,W}
 \exp\!\left\{
  \alpha\left[
   \xi\!\left(e^{tX}z_r^0z_r^{-1}\right)
   -\xi\!\left(z_r^0z_r^{-1}\right)
  \right]
 \right\},}
 \label{eq:V26-posterior-ratio}
\end{equation}
where \(C_{\alpha,t,W}>0\) is independent of \(z\).  For all sufficiently
small nonzero \(t\), the displayed density ratio is nonconstant.
Consequently
\begin{equation}
 p_{\alpha;x^{(t)},W}\ne p_{\alpha;x^{(0)},W}.
 \label{eq:V26-posteriors-unequal}
\end{equation}
In particular, at least one of these two temporal posteriors is not the
single V35 trace law \(\nu_W^{\rm tr}\).
\end{corollary}

\begin{proof}
Divide the two versions of
\eqref{eq:V23-physical-fibre-posterior}.  The terminal factor
\(f_\alpha(W,z)\), the section measure, and all links except \(r\) cancel;
the ratio of normalizers is the positive constant \(C_{\alpha,t,W}\).
Using \eqref{eq:V25-one-link-density} gives
\eqref{eq:V26-posterior-ratio}.  Its logarithmic derivative in the
\(r,X\) direction at \(z^0\) is the nonzero quantity in
\eqref{eq:V26-two-history-score-difference}.  The ratio is therefore
nonconstant, proving \eqref{eq:V26-posteriors-unequal}.  Two distinct laws
cannot both equal \(\nu_W^{\rm tr}\).
\end{proof}

The last corollary gives a qualitative obstruction.  The following positive
Gram makes the innovation finite and executable.

Left-trivialize the fibre tangent as
\[
 \mathfrak y
 :=\bigoplus_{e\in\mathscr E_N^{\rm red}}\mathfrak{su}(3).
\]
Let
\begin{equation}
 \Sigma_{\alpha;x,W}(z)
 :=\dd_z\log L_{\alpha;x,W}(z)\in\mathfrak y^*
 \label{eq:V26-likelihood-score-covector}
\end{equation}
and define the uncentered incidence Gram
\begin{equation}
 \boxed{
 \mathfrak F_{\alpha;x,W}
 :=\int
 \Sigma_{\alpha;x,W}(z)^*
 \Sigma_{\alpha;x,W}(z)\,
 \dd\nu_W^{\rm tr}(z)\succeq0.}
 \label{eq:V26-incidence-Gram}
\end{equation}

\begin{proposition}[Exact positive pair-innovation card]
\label{prop:V26-pair-innovation}
For the histories in \eqref{eq:V26-two-histories}, put
\[
 \Delta\Sigma_t
 :=\Sigma_{\alpha;x^{(t)},W}
   -\Sigma_{\alpha;x^{(0)},W},
 \qquad
 \mathfrak F_{\Delta,t}
 :=\int(\Delta\Sigma_t)^*\Delta\Sigma_t\,\dd\nu_W^{\rm tr}.
\]
Then
\begin{equation}
 \boxed{
 \mathfrak F_{\alpha;x^{(0)},W}
 +\mathfrak F_{\alpha;x^{(t)},W}
 \succeq\frac12\mathfrak F_{\Delta,t}.}
 \label{eq:V26-pair-Gram-bound}
\end{equation}
For every sufficiently small nonzero \(t\),
\begin{equation}
 \boxed{\mathfrak F_{\Delta,t}\ne0.}
 \label{eq:V26-pair-Gram-positive}
\end{equation}
Therefore at least one individual likelihood-incidence Gram is nonzero and
at least one V23 chi-square incidence residual satisfies
\begin{equation}
 \boxed{
 \mathfrak c_{\rm inc}(\alpha;x^{(j)},W)>0,
 \qquad j\in\{0,t\}.}
 \label{eq:V26-positive-chi-square}
\end{equation}
\end{proposition}

\begin{proof}
For covectors \(a,b\) in any Hilbert space,
\[
 \norm a^2+\norm b^2\ge\frac12\norm{a-b}^2.
\]
Apply this after evaluating the covectors on an arbitrary
\(v\in\mathfrak y\), then integrate under \(\nu_W^{\rm tr}\).  This proves
the operator inequality \eqref{eq:V26-pair-Gram-bound}.

For the constant command supported on link \(r\) with colour \(X\),
\Cref{thm:V26-two-history-score} shows that
\(\Delta\Sigma_t(z^0)[v]\ne0\).  The integrand is continuous and the V35
trace law has a strictly positive density on the compact fibre, so its
squared integral is positive.  This proves
\eqref{eq:V26-pair-Gram-positive}.

If \(\mathfrak F_{\alpha;x,W}\ne0\), then
\(\dd_z\log L_{\alpha;x,W}\) is not identically zero.  Hence \(L\) is
nonconstant on the connected fibre.  The V23 incidence equivalence then
gives \(\mathfrak c_{\rm inc}>0\).
\end{proof}

\begin{corollary}[Incidence failure has positive physical-history measure]
\label{cor:V26-positive-history-set}
Assume the finite Perron ground law on initial configurations has its
strictly positive continuous density from V7/V21.  For sufficiently small
nonzero \(t\), there is an open set of initial histories on which
\(\mathfrak F_{\alpha;x,W}\ne0\).  That set has positive physical measure.
Thus temporal/static fibre incidence does not hold almost surely in \(x\)
for every fixed \(W\).
\end{corollary}

\begin{proof}
At least one history from
\Cref{prop:V26-pair-innovation} has a nonzero likelihood gradient at a
point.  Joint smoothness of the finite kernel, Perron factor, and V35 action
makes this nonvanishing stable on a product neighbourhood of that history
and fibre point.  Integration over the fibre preserves strict positivity of
the Gram on a smaller history neighbourhood.  A strictly positive
continuous ground density assigns positive mass to every nonempty open set.
\end{proof}

\begin{firewall}[Which zero-incidence branch has been ruled out]
\label{fire:V26-incidence-scope}
V26 rules out the assertion that the V21 temporal posterior equals one
fixed V35 static trace law for all, or almost every, initial history.  It
does not rule out accidental equality at an isolated history.  It also does
not invalidate the V22 strong density estimate if that estimate is proved
directly for the actual V9/V21 temporal transfer.  What is forbidden is
using V35 static moments as temporal-posterior moments by a universal
zero-incidence shortcut.
\end{firewall}

% =============================================================================
\section{The compensated router has an exact untilted temporal source metric}
\label{sec:V26-router-Fisher}
% =============================================================================

The preceding innovation is unavoidable, but the baseline temporal source
metric itself is explicit.  Let
\(\dd q_{\alpha,x}(y)=c_\alpha(x,y)\dd\mathfrak m(y)\) be the normalized
tensor row and define the compensated temporal command score
\begin{equation}
 \mathscr J_{\alpha,x}^{0}[a,X](y)
 :=
 -\frac{\alpha}{3}\sum_e b[a]_e
  \operatorname{ReTr}(x_ey_e^{-1}X).
 \label{eq:V26-untilted-command-score}
\end{equation}
For a unit colour \(X\), set
\begin{equation}
 \boxed{
 \mathcal I_\alpha
 :=\frac{\alpha^2}{9}
   \int_{\mathbf G}
   \bigl(\operatorname{ReTr}(uX)\bigr)^2
   \,\dd\nu_\alpha(u)>0.}
 \label{eq:V26-one-link-Fisher}
\end{equation}
Conjugation invariance makes \(\mathcal I_\alpha\) independent of the unit
colour.

\begin{theorem}[Exact same-history router Fisher card]
\label{thm:V26-router-Fisher}
Under the untilted tensor row,
\begin{align}
 \mathbb E_{q_{\alpha,x}}
 \mathscr J_{\alpha,x}^{0}[a,X]&=0,
 \label{eq:V26-score-mean-zero}\\
 \boxed{
 \operatorname{Var}_{q_{\alpha,x}}
 \bigl(\mathscr J_{\alpha,x}^{0}[a,X]\bigr)
 =\mathcal I_\alpha\norm{b[a]}_2^2.}
 \label{eq:V26-router-Fisher-Gram}
\end{align}
The expected direct temporal curvature in the same command equals this
variance exactly.  Therefore, on the selected transverse axial carrier,
\begin{equation}
 \boxed{
 \mathcal I_\alpha\norm a_2^2
 \le
 \operatorname{Var}_{q_{\alpha,x}}
 \bigl(\mathscr J_{\alpha,x}^{0}[a,X]\bigr)
 \le
 121\mathcal I_\alpha\norm a_2^2.}
 \label{eq:V26-router-Fisher-window}
\end{equation}
The constants are independent of the volume \(M^4\).
\end{theorem}

\begin{proof}
Under \(q_{\alpha,x}\), the variables
\(u_e=x_ey_e^{-1}\) are independent with common law \(\nu_\alpha\).
Differentiating the normalization of one translated row shows
\[
 \mathbb E_{\nu_\alpha}\operatorname{ReTr}(uX)=0.
\]
Independence then removes all mixed-link covariances in
\eqref{eq:V26-untilted-command-score}; the remaining diagonal sum is
\(\mathcal I_\alpha\sum_e b[a]_e^2\), proving
\eqref{eq:V26-router-Fisher-Gram}.

Differentiate the identity
\[
 \int c_\alpha(x,y(t))\,\dd\mathfrak m(y)=1
\]
twice along the Haar-translation path generated by \(b[a]X\).  The usual
score/direct-curvature identity gives
\[
 \mathbb E_q K^{\rm t}[b[a]X,b[a]X]
 =\operatorname{Var}_q(\mathscr J_{\alpha,x}^{0}[a,X]).
\]
Finally use
\eqref{eq:V26-router-axial-bound} and
\eqref{eq:V26-router-lower-bound}.
\end{proof}

\begin{corollary}[Fixed-fibre centered baseline]
\label{cor:V26-fixed-fibre-baseline}
Condition the product temporal row on one value of the constituent coarse
field \(W\) in the parity-tree section.  Centering in the reduced fibre
removes the deterministic constituent contribution.  The remaining
baseline shell-source metric is
\begin{equation}
 \boxed{
 C_{\alpha,W}^{0,\rm fib}
 =\mathcal I_\alpha
  \mathsf L_{\rm red}^*\mathsf L_{\rm red}.}
 \label{eq:V26-fixed-fibre-Gram}
\end{equation}
It obeys the upper bound
\[
 C_{\alpha,W}^{0,\rm fib}
 \preceq120\mathcal I_\alpha I
\]
on the selected axial card, but it need not have a lower floor if
\(\mathsf L_{\rm red}\) has a kernel.
\end{corollary}

\begin{proof}
At fixed \(W\), the constituent-link term in
\eqref{eq:V26-untilted-command-score} is constant in \(z\) and disappears
under centering.  The reduced links remain independent temporal rows, so
the proof of \Cref{thm:V26-router-Fisher} gives
\(\mathcal I_\alpha\norm{\mathsf L_{\rm red}a}^2\).  The V35 axial lift
bound gives the stated upper estimate.
\end{proof}

The actual temporal posterior is a Perron tilt of the untilted row.  The
next comparison is elementary but identifies the precise missing Harnack
stock.

\begin{lemma}[Likelihood-window comparison of centered Grams]
\label{lem:V26-likelihood-window}
Let \(q\) and \(p\) be probability laws with
\(\dd p=w\,\dd q\) and
\begin{equation}
 0<m\le w\le M<\infty.
 \label{eq:V26-density-window}
\end{equation}
For any finite-dimensional real feature \(J\),
\begin{equation}
 \boxed{
 m\,\Cov_q(J)
 \preceq\Cov_p(J)
 \preceq M\,\Cov_q(J).}
 \label{eq:V26-covariance-window}
\end{equation}
\end{lemma}

\begin{proof}
For every coefficient vector \(v\),
\[
 \Var_p(v\cdot J)
 =\min_{c\in\mathbb R}\int(v\cdot J-c)^2w\,\dd q.
\]
Use \eqref{eq:V26-density-window} before taking the minimum.  This gives
\[
 m\Var_q(v\cdot J)
 \le\Var_p(v\cdot J)
 \le M\Var_q(v\cdot J),
\]
which is the quadratic-form statement
\eqref{eq:V26-covariance-window}.
\end{proof}

\begin{corollary}[Exact Perron Harnack stock for the temporal source metric]
\label{cor:V26-Perron-window}
For the full temporal posterior
\[
 \dd p_{\alpha,x}(y)
 =\frac{f_\alpha(y)}{r_\alpha(x)}\,\dd q_{\alpha,x}(y),
\]
put
\begin{equation}
 m_{\alpha,x}:=
 \frac{\inf_y f_\alpha(y)}{r_\alpha(x)},
 \qquad
 M_{\alpha,x}:=
 \frac{\sup_y f_\alpha(y)}{r_\alpha(x)}.
 \label{eq:V26-m-M}
\end{equation}
Then the actual centered Gram of the temporal command score satisfies
\begin{equation}
 \boxed{
 m_{\alpha,x}\mathcal I_\alpha
 \mathfrak R_{\rm harm}^*\mathfrak R_{\rm harm}
 \preceq
 \Cov_{p_{\alpha,x}}(\mathscr J_{\alpha,x}^{0})
 \preceq
 M_{\alpha,x}\mathcal I_\alpha
 \mathfrak R_{\rm harm}^*\mathfrak R_{\rm harm}.}
 \label{eq:V26-actual-source-window}
\end{equation}
At each finite regulator the constants are positive and finite.  A
regulator-uniform comparison requires a uniform oscillation/Harnack bound
for \(f_\alpha\); compactness at each regulator does not supply it.
\end{corollary}

\begin{proof}
Apply \Cref{lem:V26-likelihood-window} with
\(w=f_\alpha/r_\alpha(x)\) and use
\Cref{thm:V26-router-Fisher}.  Strict positivity and continuity on the
finite compact configuration space give finite-regulator positivity.  The
ratio of the upper and lower constants is
\[
 \frac{M_{\alpha,x}}{m_{\alpha,x}}
 =\frac{\sup f_\alpha}{\inf f_\alpha},
\]
so uniformity is exactly a Perron Harnack problem.
\end{proof}

\begin{firewall}[Baseline source metric versus density-defect channel]
\label{fire:V26-source-vs-defect}
The Fisher card \eqref{eq:V26-router-Fisher-Gram} is the positive source
metric of the untilted temporal row.  The actual V21 density-defect tensor is
\[
 D_\alpha
 =\Cov_{p_{\alpha,x}}(J_x^{\rm t})
  -\mathbb E_{p_{\alpha,x}}K_x^{\rm t}.
\]
Even the two-sided covariance window
\eqref{eq:V26-actual-source-window} does not bound that signed difference
without the matching direct-curvature term.  Likewise the nonzero
likelihood-incidence Gram
\eqref{eq:V26-incidence-Gram} first enters the V20 incidence-residual/source
atlas; it becomes a Woodbury defect column only after its contribution to
the actual Witten Hessian is derived.  Neither identification is made by
notation.
\end{firewall}

% =============================================================================
\section{V26 ledger and revised frontier}
\label{sec:V26-ledger}
% =============================================================================

\begin{theorem}[Integrated V26 statement]
\label{thm:V26-integrated}
Preserving V1--V25 and superseding only the V25 shell-router typing, V26
proves:
\begin{enumerate}[label=\textup{(V26.\arabic*)},leftmargin=5.5em]
\item the exact parity-tree coordinate section, unit fibre Jacobian, and
      distinct detail and raw-shell differentials;
\item the Haar-compensated physical shell router with the V35
      volume-uniform \(141\) and axial \(121\) norm bounds;
\item elimination of \(\nabla\log h_\alpha\) from the explicit likelihood
      and a literal reduced-link and compensated-shell score;
\item a two-history identity which cancels all terminal/static terms and
      proves that two temporal fibre posteriors are different;
\item a positive pair-incidence Gram proving nonzero V23 incidence residual
      on an open, positive-physical-measure set of initial histories;
\item the exact untilted temporal Fisher metric of the compensated router,
      with a volume-independent axial metric window; and
\item a two-sided comparison of the actual Perron-tilted source covariance
      with that baseline, isolating the precise Perron Harnack stock needed
      for cofinal uniformity.
\end{enumerate}
\end{theorem}

\begin{proof}
Items \textup{(V26.1)--(V26.2)} are
\Cref{prop:V26-section-differential,prop:V26-harmonic-router}.  Item
\textup{(V26.3)} is
\Cref{thm:V26-Perron-elimination,cor:V26-one-link-score,
cor:V26-shell-score}.  Items \textup{(V26.4)--(V26.5)} are
\Cref{thm:V26-two-history-score,cor:V26-two-posteriors-different,
prop:V26-pair-innovation,cor:V26-positive-history-set}.  Item
\textup{(V26.6)} is \Cref{thm:V26-router-Fisher}.  Item
\textup{(V26.7)} is
\Cref{lem:V26-likelihood-window,cor:V26-Perron-window}.
\end{proof}

\begingroup
\small
\begin{longtable}{>{\raggedright\arraybackslash}p{0.28\textwidth}
                  >{\raggedright\arraybackslash}p{0.16\textwidth}
                  >{\raggedright\arraybackslash}p{0.48\textwidth}}
\toprule
\textbf{V26 row}&\textbf{Status}&\textbf{Advance or remaining obstruction}\\
\midrule
V35 parity-tree fibre coordinates
& \MGclosed
& The section, unit Jacobian, and reduced-detail injection are now literal.\\[0.35em]
Physical shell command router
& \MGclosed
& The Haar-compensated harmonic lift replaces the mistyped \(D_z\Phi\)
  shell router and has a volume-uniform axial norm bound.\\[0.35em]
Perron derivative
& \MGclosed
& Eliminated exactly from the likelihood score in favour of
  \(r_\alpha=\mathcal C_\alpha f_\alpha\).\\[0.35em]
Universal temporal/static incidence
& \MGfail
& The two-history card proves a nonzero incidence Gram on a positive-measure
  set; the innovation branch is mandatory.\\[0.35em]
Untilted temporal router metric
& \MGclosed
& Exact scalar Fisher coefficient times
  \(\mathfrak R_{\rm harm}^*\mathfrak R_{\rm harm}\), with the axial
  \(1\)-to-\(121\) metric window.\\[0.35em]
Actual Perron source metric
& \MGreduced
& Reduced to the explicit likelihood/Harnack window
  \eqref{eq:V26-actual-source-window}; uniform oscillation remains open.\\[0.35em]
Actual density-defect ceiling
& \MGopen
& Must retain covariance minus expected direct curvature under the physical
  posterior; the covariance card alone is insufficient.\\[0.35em]
Complete source edge/locality packet
& \MGopen
& The forced likelihood innovation must be routed into V20 before the V24
  and V25 pencils are populated.\\[0.35em]
Capacity, protected sectors, and OS limit
& \MGopen
& Bad capacity, toron/wrapping floors, local noncollapse, and nontrivial
  continuum reconstruction remain.\\[0.35em]
Full Yang--Mills mass gap
& \MGopen
& A false zero-incidence shortcut is removed and the physical router is
  explicit, but the actual-law cofinal packet is not closed.\\
\bottomrule
\end{longtable}
\endgroup

\begin{openproblem}[V27 mandate: route the forced likelihood innovation on
the actual Perron law]
\label{op:V27}
\begin{enumerate}[label=\textup{(V27.\arabic*)},leftmargin=5.5em]
\item Treat \(\mathfrak F_{\alpha;x,W}\) as a genuine conditional-law
      incidence residual.  Insert its polar range into the complete V20
      source/incidence atlas; do not silently add it to the Woodbury defect
      channel before a Hessian identity licenses that step.
\item Use \eqref{eq:V26-r-eigen-equation} and the exact V22 posterior
      integration-by-parts identity to assemble
      \(\Cov_pJ^{\rm t}-\mathbb E_pK^{\rm t}\) before taking norms.  Seek a
      direct \(O(|\beta|)\) ceiling on the actual V10/V21 law which does not
      assume equality with the V35 static fibre.
\item On the protected chart, bound the likelihood window
      \(M_{\alpha,x}/m_{\alpha,x}\) or route its bad set through the V19/V20
      capacity mechanism.  A finite-regulator compactness constant is not a
      cofinal Harnack bound.
\item Form the enlarged entrance and populate the first actual
      \(G,H,J,H_x\) packet.  Apply both the V25 deficit and locality pencils.
\item If the edge collapses, return the physical survivor and its
      likelihood-innovation component.  If only locality collapses, return
      the exact weighted witness and its responsible cell pair.
\item On the strict branch, combine the exponential return with the complete
      incidence--capacity matrix, then short toron, wrapping, chart-exit,
      and protected sectors.
\item Keep \(\gamma_{\alpha,j}/\alpha\), local-algebra noncollapse, and the
      reflection-positive OS limit explicit through the cofinal passage.
\item The next companion-file audit remains V30.
\end{enumerate}
\end{openproblem}

\begin{firewall}[End-of-V26 claim boundary]
V26 makes the V35 parity-tree coordinates and physical shell router
literal, eliminates the Perron derivative from the fibre likelihood, and
proves that a nonzero temporal/static incidence innovation occurs on a
positive-measure set of initial histories.  It also computes the exact
untilted router Fisher metric and isolates its actual-law Harnack window.
It does not yet bound the signed density defect on the actual Perron law,
populate the enlarged cofinal edge/locality packet, control capacity and
protected sectors, or construct a nontrivial continuum OS state.  Therefore
the full Yang--Mills mass gap remains open.
\end{firewall}

\section*{Version V26 append-only record}
\addcontentsline{toc}{section}{Version V26 append-only record}
The complete V25 mathematical source was retained before this continuation.  Its original
SHA--256 digest was
\begin{center}\scriptsize\ttfamily
e85faf51d07b7b4364e46204fe9b5eb05a0ffcb6eb778abcde10ff07b944baa8.
\end{center}
The V26 file also inserts the three previously missing proof-status macro
definitions \texttt{\string\MGfail}, \texttt{\string\MGclosed}, and
\texttt{\string\MGreduced} in the preamble.  This repairs the
compilability of earlier ledgers without changing their mathematical text.
On the next iteration, retain this full file, append before its unique active
document-closing command, increment the filename to
\texttt{\_V27.tex}, and remove only the superseded V26 file from this note
series.  The next companion audit is V30.

% =============================================================================
% END APPENDED VERSION V26 MATERIAL
% =============================================================================

% =============================================================================
% BEGIN APPENDED VERSION V27 MATERIAL
% =============================================================================

\clearpage
\part{Version V27: the anisotropic Perron resolvent, the correct magnetic parameter, and local incidence routing}
\label{part:V27}

\section{Direction: the requested \(O(|\beta|)\) ceiling has an anisotropic
resolvent obstruction}
\label{sec:V27-direction}

V26 made the literal temporal convolution, terminal Perron tilt, and
parity-tree command router explicit.  Its V27 mandate asked first for a
direct \(O(|\beta|)\) ceiling on the signed temporal posterior excess.  The
correct upstream test is to differentiate the \emph{actual} anisotropic
Perron equation at zero magnetic coupling before estimating its two terms
separately.

That test gives a negative answer to an \(\alpha\)-uniform
\(O(|\beta|)\) assertion.  The first magnetic variation contains the
electric resolvent
\[
 (I-\mathcal C_\alpha)^{-1}\mathcal C_\alpha.
\]
The centered Wilson plaquette potential is an exact eigenfunction of the
product temporal convolution with eigenvalue \(\eta_F(\alpha)^4\).  Hence
the density defect is amplified by
\[
 \frac{\eta_F(\alpha)^4}{1-\eta_F(\alpha)^4}
 \sim \frac{\alpha}{8}.
\]
This does not contradict the V9 small-coupling theorem: V9 did not contain
the independent temporal concentration \(\alpha\).  It proves instead that
the conditional V9-to-V21 transfer identification retained in V22--V23
cannot be used uniformly for the literal V25 anisotropic transfer without a
new anisotropic polymer argument.

V27 proves the exact first-variation formula, identifies the correct
dimensionless magnetic parameter, and replaces the volume-extensive global
Harnack request by a conditional local window.  It also gives the exact
weighted-codifferential identity which licenses the V26 likelihood score as
a V20 incidence residual, and proves a local-Harnack lower frame for the
complete compensated temporal source.  No nonlinear cofinal Harnack bound is
asserted.

The V35 geometric input used below is still the file frozen in V26, with
SHA--256
\begin{center}\scriptsize\ttfamily
7c6056d8149d3015d469b57474a75b23c88a6fe9f64fc81c7dca85c361c647a3.
\end{center}
Its instructions do not direct this continuation.

% =============================================================================
\section{The centered Wilson potential is one exact temporal eigenmode}
\label{sec:V27-plaquette-mode}
% =============================================================================

Let \(\mathcal X_\Lambda=\mathbf G^{E(\Lambda)}\),
\(\mathbf G=\mathrm{SU}(3)\), with product Haar probability
\(\mathfrak m_\Lambda\).  Retain the V25 product convolution
\begin{equation}
 (\mathcal C_{\alpha,\Lambda}u)(x)
 =\int c_\alpha(x,y)u(y)\,\dd\mathfrak m_\Lambda(y).
 \label{eq:V27-product-convolution}
\end{equation}
For an elementary oriented plaquette \(p\) whose four boundary links are
distinct, put
\begin{equation}
 v_p(x):=\frac13\operatorname{ReTr}U_p(x).
 \label{eq:V27-plaquette-character}
\end{equation}
The usual unit Wilson potential has the form
\begin{equation}
 V_{M,\Lambda}=\sum_{p\in\mathcal P_\Lambda}(1-v_p),
 \qquad
 \widetilde V_{M,\Lambda}
 :=V_{M,\Lambda}-\int V_{M,\Lambda}\,\dd\mathfrak m_\Lambda.
 \label{eq:V27-centered-Wilson}
\end{equation}
Changing the harmless sign convention for \(V_M\) changes both sides of the
linear formulas below by the same sign.

\begin{lemma}[Exact four-link temporal multiplier]
\label{lem:V27-plaquette-eigenmode}
Let
\begin{equation}
 q_\alpha:=\eta_F(\alpha)^4\in(0,1).
 \label{eq:V27-q-alpha}
\end{equation}
Then
\begin{align}
 \int v_p\,\dd\mathfrak m_\Lambda&=0,
 \label{eq:V27-plaquette-mean}\
 \mathcal C_{\alpha,\Lambda}v_p&=q_\alpha v_p,
 \label{eq:V27-plaquette-eigen}\
 \boxed{
 \mathcal C_{\alpha,\Lambda}\widetilde V_{M,\Lambda}
 =q_\alpha\widetilde V_{M,\Lambda}.}
 \label{eq:V27-Wilson-eigen}
\end{align}
The eigenvalue is independent of the volume and of the plaquette position.
\end{lemma}

\begin{proof}
Integrating one boundary link of \(p\) against Haar annihilates every
fundamental matrix coefficient, proving
\eqref{eq:V27-plaquette-mean}.  To prove
\eqref{eq:V27-plaquette-eigen}, write the trace in matrix coefficients and
integrate the four temporal increments successively.  Centrality of the
one-link law and the V10 Peter--Weyl formula give
\[
 \int_{\mathbf G}D_F(g)\,\dd\nu_\alpha(g)
 =\eta_F(\alpha)I_F,
 \qquad
 \int_{\mathbf G}D_F(g^{-1})\,\dd\nu_\alpha(g)
 =\eta_F(\alpha)I_F.
\]
Each of the four distinct boundary links therefore contributes one factor
\(\eta_F(\alpha)\).  Taking the real trace gives
\(\mathcal C_{\alpha,\Lambda}v_p=\eta_F(\alpha)^4v_p\).
Linearity and
\(\widetilde V_{M,\Lambda}=-\sum_pv_p\) prove
\eqref{eq:V27-Wilson-eigen}.
\end{proof}

\begin{firewall}[Degenerate microscopic tori]
\label{fire:V27-degenerate-torus}
The four-factor proof assumes the standard nondegenerate elementary
plaquette.  A microscopic periodic torus on which an oriented plaquette
reuses a link has a different finite multiplier and must be evaluated
directly.  Every cofinal lattice family contains a tail of nondegenerate
plaquettes, which is the only family used below.
\end{firewall}

% =============================================================================
\section{Exact first magnetic variation of the Perron density defect}
\label{sec:V27-first-variation}
% =============================================================================

The next theorem is stated for any centered convolution eigenmode; the
Wilson specialization is then immediate.  Let \(V\in C^\infty(\mathcal
X_\Lambda;\mathbb R)\), write
\[
 \overline V:=\int V\,\dd\mathfrak m_\Lambda,
 \qquad \widetilde V:=V-\overline V,
\]
and suppose
\begin{equation}
 \mathcal C_{\alpha,\Lambda}\widetilde V=q\widetilde V,
 \qquad 0<q<1.
 \label{eq:V27-general-eigenmode}
\end{equation}
For real \(\beta\) near zero define the literal symmetric anisotropic
transfer
\begin{equation}
 \mathcal T_{\alpha,\beta}
 :=M_{e^{-\beta V/2}}\mathcal C_{\alpha,\Lambda}
   M_{e^{-\beta V/2}},
 \qquad
 \mathcal T_{\alpha,\beta}h_{\alpha,\beta}
 =\lambda_{\alpha,\beta}h_{\alpha,\beta},
 \label{eq:V27-anisotropic-transfer}
\end{equation}
with \(h_{\alpha,\beta}>0\) and normalization
\begin{equation}
 \int h_{\alpha,\beta}\,\dd\mathfrak m_\Lambda=1.
 \label{eq:V27-Perron-normalization}
\end{equation}
Put
\begin{equation}
 f_{\alpha,\beta}:=e^{-\beta V/2}h_{\alpha,\beta},
 \qquad
 r_{\alpha,\beta}:=\mathcal C_{\alpha,\Lambda}f_{\alpha,\beta},
 \qquad
 D_{\alpha,\beta}:=\Hess\log r_{\alpha,\beta}.
 \label{eq:V27-f-r-D}
\end{equation}

\begin{theorem}[Exact Perron resolvent and defect derivative]
\label{thm:V27-first-variation}
At every fixed finite \((\alpha,\Lambda)\), the Perron data in
\eqref{eq:V27-anisotropic-transfer} are real analytic in \(\beta\) near
zero as \(C^2\) functions of the configuration.  At \(\beta=0\),
\begin{align}
 h_{\alpha,0}&=f_{\alpha,0}=r_{\alpha,0}=1,
 &\lambda_{\alpha,0}&=1,
 \label{eq:V27-zero-data}\
 \left.\partial_\beta\lambda_{\alpha,\beta}\right|_0
 &=-\overline V,
 \label{eq:V27-lambda-prime}\
 \boxed{
 \left.\partial_\beta\log h_{\alpha,\beta}\right|_0
 =-\frac12\frac{1+q}{1-q}\widetilde V,}
 \label{eq:V27-h-prime}\
 \left.\partial_\beta\log f_{\alpha,\beta}\right|_0
 &=-\frac{\overline V}{2}-\frac{1}{1-q}\widetilde V,
 \label{eq:V27-f-prime}\
 \left.\partial_\beta\log r_{\alpha,\beta}\right|_0
 &=-\frac{\overline V}{2}-\frac{q}{1-q}\widetilde V,
 \label{eq:V27-r-prime}\
 \boxed{
 \left.\partial_\beta D_{\alpha,\beta}\right|_0
 =-\frac{q}{1-q}\Hess V.}
 \label{eq:V27-D-prime}
\end{align}
Equivalently, before using the eigenmode relation,
\begin{equation}
 \left.\partial_\beta D_{\alpha,\beta}\right|_0
 =-\Hess\!left[
  (I-\mathcal C_{\alpha,\Lambda})^{-1}
  \mathcal C_{\alpha,\Lambda}\widetilde V
 \right].
 \label{eq:V27-D-resolvent}
\end{equation}
The inverse acts only on the mean-zero subspace.
\end{theorem}

\begin{proof}
The positive smooth kernel makes \(\mathcal T_{\alpha,0}=\mathcal
C_{\alpha,\Lambda}\) compact and positivity improving.  Its eigenvalue one
is simple, with eigenfunction one, and is isolated from the remainder of the
spectrum.  Multiplication by \(e^{-\beta V/2}\) is an analytic bounded
family on every fixed \(C^k\) space.  The analytic implicit-function theorem
applied to the eigenvalue equation together with
\eqref{eq:V27-Perron-normalization} therefore gives the claimed \(C^2\)
analyticity.  This argument uses only the finite-regulator spectral
isolation; it makes no cofinal transfer-gap claim.

Differentiate
\eqref{eq:V27-anisotropic-transfer} at zero.  Since
\[
 \left.\partial_\beta\mathcal T_{\alpha,\beta}\right|_0\!1
 =-\frac12(V+\mathcal C_{\alpha,\Lambda}V),
\]
integration against Haar gives
\eqref{eq:V27-lambda-prime}.  If
\(\dot h=\partial_\beta h|_0\), the mean-zero normalization of
\(\dot h\) and the differentiated eigenvalue equation give
\begin{equation}
 (I-\mathcal C_{\alpha,\Lambda})\dot h
 =-\frac12(I+\mathcal C_{\alpha,\Lambda})\widetilde V.
 \label{eq:V27-h-resolvent-equation}
\end{equation}
The operator \(I-\mathcal C_{\alpha,\Lambda}\) is invertible on the
mean-zero subspace.  Under
\eqref{eq:V27-general-eigenmode}, equation
\eqref{eq:V27-h-resolvent-equation} gives
\[
 \dot h=-\frac12\frac{1+q}{1-q}\widetilde V.
\]
Because \(h_{\alpha,0}=1\), this is also
\(\partial_\beta\log h|_0\), proving
\eqref{eq:V27-h-prime}.

Differentiate
\(\log f=-\beta V/2+\log h\) to obtain
\eqref{eq:V27-f-prime}.  Since \(r_0=1\),
\[
 \left.\partial_\beta\log r\right|_0
 =\mathcal C_{\alpha,\Lambda}
   \left(\left.\partial_\beta f\right|_0\right),
\]
and applying the eigenmode relation proves
\eqref{eq:V27-r-prime}.  Two spatial derivatives give
\eqref{eq:V27-D-prime}.

For the general formula, solve
\eqref{eq:V27-h-resolvent-equation} without diagonalizing.  The
nonconstant part of \(\partial_\beta\log r|_0\) is
\begin{align*}
 &-\frac12\mathcal C\widetilde V
 -\frac12\mathcal C(I-\mathcal C)^{-1}
  (I+\mathcal C)\widetilde V\\
 &\qquad
 =-(I-\mathcal C)^{-1}\mathcal C\widetilde V,
\end{align*}
where commuting functions of \(\mathcal C\) were combined.  Taking the
Hessian proves \eqref{eq:V27-D-resolvent}.
\end{proof}

\begin{corollary}[Exact Wilson specialization]
\label{cor:V27-Wilson-first-variation}
For \(V=V_{M,\Lambda}\),
\begin{equation}
 \boxed{
 \left.\partial_\beta D_{\alpha,\beta}\right|_0
 =-\frac{\eta_F(\alpha)^4}
         {1-\eta_F(\alpha)^4}\Hess V_{M,\Lambda}.}
 \label{eq:V27-Wilson-D-prime}
\end{equation}
The corresponding first variation of the V22/V23 density channel is
\begin{equation}
 \left.\partial_\beta S_{\vartheta}\right|_0
 =2\frac{\eta_F(\alpha)^4}
        {1-\eta_F(\alpha)^4}\Hess V_{M,\Lambda},
 \label{eq:V27-S-prime}
\end{equation}
because \(S_\vartheta=-2D\).
\end{corollary}

\begin{proof}
Apply \Cref{lem:V27-plaquette-eigenmode,thm:V27-first-variation} and the
exact V22 identity \(S_\vartheta=-2D\).
\end{proof}

% =============================================================================
\section{The \(\alpha\)-uniform \(O(|\beta|)\) route is ruled out}
\label{sec:V27-nonuniformity}
% =============================================================================

\begin{theorem}[No \(\alpha\)-uniform \(\beta\)-only defect ceiling]
\label{thm:V27-no-uniform-beta-ceiling}
Assume \(V_{M,\Lambda}\) is nonconstant and the cofinal family contains the
same local plaquette geometry.  There do not exist constants
\(C<\infty\), \(\beta_0>0\), and \(\alpha_0\) such that
\begin{equation}
 \norm{D_{\alpha,\beta}}_{L^\infty(\op)}
 \le C|\beta|
 \qquad
 (\alpha\ge\alpha_0,\ |\beta|<\beta_0)
 \label{eq:V27-false-uniform-beta-bound}
\end{equation}
for the literal anisotropic transfer
\eqref{eq:V27-anisotropic-transfer}.
\end{theorem}

\begin{proof}
If \(\Hess V_{M,\Lambda}\) vanished identically, then
\(\nabla V_{M,\Lambda}\) would be a parallel vector field.  A real function
on the compact configuration manifold with parallel gradient is constant.
Thus
\begin{equation}
 c_V:=\norm{\Hess V_{M,\Lambda}}_{L^\infty(\op)}>0.
 \label{eq:V27-cV-positive}
\end{equation}
If \eqref{eq:V27-false-uniform-beta-bound} held, divide by
\(|\beta|\) and let \(\beta\to0\).  The \(C^2\) analyticity in
\Cref{thm:V27-first-variation} and
\eqref{eq:V27-Wilson-D-prime} would give
\[
 C\ge c_V\frac{\eta_F(\alpha)^4}
                 {1-\eta_F(\alpha)^4}
 \qquad(\alpha\ge\alpha_0).
\]
But \(\eta_F(\alpha)\uparrow1\), so the right-hand side diverges.  This is a
contradiction.
\end{proof}

\begin{corollary}[Exact anisotropic amplification and its Hamiltonian limit]
\label{cor:V27-correct-parameter}
The V10 definition
\(s_\alpha=-\log\eta_F(\alpha)/C_F\) gives the exact identity
\begin{equation}
 \frac{q_\alpha}{1-q_\alpha}
 =\frac{1}{e^{4C_Fs_\alpha}-1}.
 \label{eq:V27-exact-amplifier}
\end{equation}
Consequently
\begin{equation}
 \frac{q_\alpha}{1-q_\alpha}
 =\frac{\alpha}{6C_F}+O(1)
 =\frac{\alpha}{8}+O(1)
 \quad\text{for }\mathrm{SU}(3), C_F=\frac43.
 \label{eq:V27-amplifier-asymptotic}
\end{equation}
The correct linear density parameter is therefore
\begin{equation}
 \boxed{
 \zeta_{\alpha,\beta}^{\rm den}
 :=|\beta|\frac{q_\alpha}{1-q_\alpha},}
 \label{eq:V27-zeta-density}
\end{equation}
not \(\lvert\beta\rvert\) alone.  On the V21 matched trajectory
\(\beta_{\alpha,j}=s_\alpha\gamma_{\alpha,j}/\kappa\),
\begin{equation}
 \boxed{
 \beta_{\alpha,j}\frac{q_\alpha}{1-q_\alpha}
 =\frac{\gamma_{\alpha,j}}{\kappa}
   \frac{s_\alpha}{e^{4C_Fs_\alpha}-1}
 \longrightarrow
 \frac{\gamma_j}{4\kappa C_F}}
 \label{eq:V27-matched-density-parameter}
\end{equation}
whenever \(\gamma_{\alpha,j}\to\gamma_j\).  The coefficient of the
linearized physical density debit obeys
\begin{equation}
 2\kappa\beta_{\alpha,j}
 \left.\partial_\beta D_{\alpha,\beta}\right|_0
 \longrightarrow
 -\frac{\gamma_j}{2C_F}\Hess V_{M,j}.
 \label{eq:V27-linearized-physical-debit}
\end{equation}
This is a statement about the exact first variation; passage of the full
nonlinear \(D_{\alpha,\beta_{\alpha,j}}\) requires a remainder uniform in
the new parameter.
\end{corollary}

\begin{proof}
Since \(q_\alpha=\eta_F(\alpha)^4=e^{-4C_Fs_\alpha}\),
\eqref{eq:V27-exact-amplifier} is algebraic.  Use
\(s_\alpha=3/(2\alpha)+O(\alpha^{-2})\) and
\(e^t-1=t+O(t^2)\) to obtain
\eqref{eq:V27-amplifier-asymptotic}.  Substitute the matched coefficient
into \eqref{eq:V27-exact-amplifier}; then
\(s/(e^{4C_Fs}-1)\to(4C_F)^{-1}\), proving
\eqref{eq:V27-matched-density-parameter}.  Multiply
\eqref{eq:V27-Wilson-D-prime} by \(2\kappa\beta_{\alpha,j}\) for the final
formula.
\end{proof}

\begin{assessment}[Decision of the old V9 incidence premise]
\label{ass:V27-V9-incidence-decision}
The V9 Hessian majorant is a theorem for the one-parameter transfer denoted
there by \(\mathcal K_{\beta,\Lambda}\).  If that theorem could be identified
with the entire literal two-parameter family
\(\mathcal T_{\alpha,\beta}\) while retaining a constant depending on
\(\beta\) but not \(\alpha\), it would imply
\eqref{eq:V27-false-uniform-beta-bound}.  The exact derivative disproves
that implication.  Therefore the conditional ``complete-transfer
incidence'' in the V22/V23 strong branch is not supplied by merely naming
the V25 transfer \(\mathcal K_{\beta,\Lambda}\).  A valid upstream theorem
must keep the temporal susceptibility
\((1-q_\alpha)^{-1}\) and prove a nonlinear expansion in a parameter such as
\eqref{eq:V27-zeta-density}.
\end{assessment}

% =============================================================================
\section{Local conditional Harnack windows replace the global oscillation}
\label{sec:V27-local-Harnack}
% =============================================================================

The V26 covariance comparison used the global ratio
\(\sup f/\inf f\).  That ratio is normally volume-extensive for a local
many-body potential.  Conditional source acquisition needs only the
oscillation in the links actually varied while the exterior is fixed.

Let \(B\subset E(\Lambda)\) be a finite link block.  Write
\(y=(y_B,y_{B^c})\) and, for an exterior configuration \(\eta\), define
\begin{equation}
 \Omega_{\alpha,\beta}(B;\eta)
 :=\operatorname{osc}_{y_B}
   \log f_{\alpha,\beta}(y_B,\eta).
 \label{eq:V27-local-oscillation}
\end{equation}
For fixed initial history \(x\), disintegrate the actual temporal posterior
\(p_{\alpha,\beta;x}\) and the product row
\(q_{\alpha,x}\) with respect to \(y_{B^c}\).

\begin{theorem}[Exact conditional likelihood window]
\label{thm:V27-conditional-window}
For every exterior configuration \(\eta\),
\begin{equation}
 \frac{\dd p_{\alpha,\beta;x}^{B,\eta}}
      {\dd q_{\alpha,x}^{B}}(y_B)
 =\frac{f_{\alpha,\beta}(y_B,\eta)}
        {\mathbb E_{q_{\alpha,x}^{B}}
         f_{\alpha,\beta}(\,\cdot\,,\eta)}.
 \label{eq:V27-conditional-RN}
\end{equation}
Here \(q_{\alpha,x}^{B}\) is the product of the one-link temporal rows in
\(B\), independent of \(\eta\).  Moreover,
\begin{equation}
 \boxed{
 e^{-\Omega_{\alpha,\beta}(B;\eta)}
 \le
 \frac{\dd p_{\alpha,\beta;x}^{B,\eta}}
      {\dd q_{\alpha,x}^{B}}
 \le
 e^{\Omega_{\alpha,\beta}(B;\eta)}.}
 \label{eq:V27-conditional-density-window}
\end{equation}
Consequently every finite-dimensional feature \(J_B(y_B)\) satisfies
\begin{equation}
 \boxed{
 e^{-\Omega_{\alpha,\beta}(B;\eta)}\Cov_{q^B}(J_B)
 \preceq
 \Cov_{p^{B,\eta}}(J_B)
 \preceq
 e^{\Omega_{\alpha,\beta}(B;\eta)}\Cov_{q^B}(J_B).}
 \label{eq:V27-conditional-covariance-window}
\end{equation}
No global oscillation of \(f_{\alpha,\beta}\) occurs.
\end{theorem}

\begin{proof}
The full posterior has density \(f_{\alpha,\beta}/r_{\alpha,\beta}(x)\)
relative to the product temporal row.  Conditionalization at fixed
\(y_{B^c}=\eta\) cancels the factor \(r_{\alpha,\beta}(x)\) and gives
\eqref{eq:V27-conditional-RN}.  If \(m\) and \(M\) are the minimum and
maximum of \(f(\,\cdot\,,\eta)\), its conditional normalizer lies in
\([m,M]\).  Hence the normalized density lies between \(m/M\) and \(M/m\),
which is \eqref{eq:V27-conditional-density-window}.  Apply the variance
minimization proof of \Cref{lem:V26-likelihood-window} conditionally to
obtain \eqref{eq:V27-conditional-covariance-window}.
\end{proof}

\begin{proposition}[The first local Harnack coefficient]
\label{prop:V27-local-Harnack-linearization}
At fixed finite \((\alpha,\Lambda)\), uniformly in the exterior
configuration,
\begin{equation}
 \boxed{
 \lim_{\beta\to0}
 \frac{\Omega_{\alpha,\beta}(B;\eta)}{|\beta|}
 =\frac{1}{1-q_\alpha}
   \operatorname{osc}_{y_B}V_{M,\Lambda}(y_B,\eta).}
 \label{eq:V27-local-Harnack-derivative}
\end{equation}
For the unit Wilson action,
\begin{equation}
 \sup_\eta\operatorname{osc}_{y_B}V_{M,\Lambda}(y_B,\eta)
 \le2d_\partial|B|,
 \label{eq:V27-local-Wilson-oscillation}
\end{equation}
where \(d_\partial\) is the volume-independent maximum number of plaquettes
incident with a link.  Thus the first local-window coefficient is volume
independent for fixed \(B\), but its small parameter is
\begin{equation}
 \boxed{
 \zeta_{\alpha,\beta}^{\rm Har}
 :=\frac{|\beta|}{1-q_\alpha}.}
 \label{eq:V27-zeta-Harnack}
\end{equation}
On the matched trajectory it has the same limit
\(|\gamma_j|/(4\kappa C_F)\) as
\eqref{eq:V27-zeta-density}.
\end{proposition}

\begin{proof}
Equation \eqref{eq:V27-f-prime} and \(C^0\) analyticity give
\[
 \log f_{\alpha,\beta}
 =c_{\alpha,\beta}
  -\frac{\beta}{1-q_\alpha}\widetilde V_{M,\Lambda}
  +O_{\alpha,\Lambda}(\beta^2)
\]
uniformly on the compact configuration space.  Additive constants and the
Haar mean disappear under \(\operatorname{osc}_{y_B}\).  The oscillation
seminorm is continuous in the uniform norm, proving
\eqref{eq:V27-local-Harnack-derivative}.

Changing \(y_B\) affects only plaquettes incident with \(B\).  Each normalized
real trace has oscillation at most two, and there are at most
\(d_\partial|B|\) such plaquettes after overcounting, proving
\eqref{eq:V27-local-Wilson-oscillation}.  Finally
\[
 \frac{|\beta_{\alpha,j}|}{1-q_\alpha}
 =\frac{|\gamma_{\alpha,j}|}{\kappa}
   \frac{s_\alpha}{1-e^{-4C_Fs_\alpha}}
 \longrightarrow\frac{|\gamma_j|}{4\kappa C_F}.
\]
\end{proof}

\begin{firewall}[Linear local window versus nonlinear cofinal Harnack]
\label{fire:V27-linear-Harnack-scope}
The \(O_{\alpha,\Lambda}(\beta^2)\) remainder in the proof is not asserted
uniform in \(\alpha\), volume, or block separation.  Equation
\eqref{eq:V27-local-Harnack-derivative} identifies the correct parameter
and removes a false global requirement; it does not prove the nonlinear
cofinal bound
\[
 \sup_{\alpha,j,\eta}
 \Omega_{\alpha,\beta_{\alpha,j}}(B;\eta)<\infty.
\]
That bound and a summable exterior-influence estimate are the next
anisotropic polymer obligations.
\end{firewall}

% =============================================================================
\section{A conditional-Harnack lower frame for the complete temporal source}
\label{sec:V27-source-lower-frame}
% =============================================================================

Retain the complete V26 compensated command \(b[a]\), including its
constituent part \(b|_C=a\).  Let
\(\mathscr J_{\alpha,x}^{0}[a,X]\) be the temporal command score in
\eqref{eq:V26-untilted-command-score}.  If the coarse command is supported
on a finite set \(A\) of constituent links, those same terminal link
coordinates form a block, also denoted \(A\).

\begin{theorem}[Complete-source lower frame from a local window]
\label{thm:V27-source-lower-frame}
Assume that for the block \(A\)
\begin{equation}
 \sup_\eta\Omega_{\alpha,\beta}(A;\eta)\le A_*<\infty.
 \label{eq:V27-local-window-stock}
\end{equation}
Then, under the actual full temporal posterior,
\begin{equation}
 \boxed{
 \operatorname{Var}_{p_{\alpha,\beta;x}}
 \bigl(\mathscr J_{\alpha,x}^{0}[a,X]\bigr)
 \ge e^{-A_*}\mathcal I_\alpha\norm a_2^2.}
 \label{eq:V27-actual-source-lower-frame}
\end{equation}
The constant is independent of the volume and does not require a lower
singular value for the reduced-only map \(\mathsf L_{\rm red}\).
\end{theorem}

\begin{proof}
Condition the full terminal configuration on every link outside \(A\).
All score summands outside \(A\) then become constants.  On the links in
\(A\), the coefficient of the complete command is exactly \(a_e\), because
the constituent row of the V35 lift is the identity.  The conditional
product-row variance is therefore
\[
 \mathcal I_\alpha\sum_{e\in A}|a_e|^2
 =\mathcal I_\alpha\norm a_2^2
\]
by the one-link computation in \Cref{thm:V26-router-Fisher}.  Apply the
lower half of \eqref{eq:V27-conditional-covariance-window} and the uniform
stock \eqref{eq:V27-local-window-stock}.  Finally, the law of total variance
gives
\[
 \Var_p(\mathscr J^0)
 \ge\mathbb E_p\!\left[
     \Var_p(\mathscr J^0\mid y_{A^c})
   \right],
\]
which proves \eqref{eq:V27-actual-source-lower-frame}.
\end{proof}

\begin{assessment}[Why this repairs the reduced-fibre degeneracy]
\label{ass:V27-complete-vs-fibre-source}
The V26 fixed-\(W\) fibre covariance sees only
\(\mathsf L_{\rm red}a\) and may vanish on its kernel.  The physical
same-history command is not a reduced-only variation: it also moves the
second constituent link by \(a\).  Conditioning the complete posterior on
the exterior exposes that identity row, and
\eqref{eq:V27-actual-source-lower-frame} prevents cancellation by every
reduced or remote summand.  This is a lower frame for the source covariance,
not for the signed density defect \(D=\Cov J-\mathbb EK\).
\end{assessment}

The attached V35 source also proves that the harmonic router itself is
exponentially local.  To avoid repeating that proof, record only the exact
imported conclusion.  For the numerical \(\mu>0\) and \(h_\mu<\infty\)
defined in its equation \texttt{YMA3-fixed-locality-constants}, its lemma
\texttt{YMA3-Green-locality} gives
\begin{equation}
 |\mathcal H|_\mu\le h_\mu,
 \qquad
 \mathcal H_M=\operatorname{Per}_M\mathcal H,
 \label{eq:V27-imported-router-locality}
\end{equation}
where \(\mathcal H:a\mapsto b[a]\).

\begin{corollary}[Router-tail certificate]
\label{cor:V27-router-tail}
If \(a\) is supported in coarse cells \(A\), and \(P_{A,R}\) retains fine
links at cell distance at most \(R\) from \(A\), then on the infinite lattice
and every finite periodization,
\begin{equation}
 \boxed{
 \norm{(I-P_{A,R})\mathcal H a}_2
 \le h_\mu e^{-\mu R}\norm a_2.}
 \label{eq:V27-router-tail}
\end{equation}
Thus the geometric router contributes a summable, volume-independent
incidence kernel.  Any remaining failure of source locality must come from
the actual conditional law or from the subsequent comparison inverse, not
from the V35 harmonic lift itself.
\end{corollary}

\begin{proof}
Outside the \(R\)-neighbourhood, insert
\(1\le e^{-\mu R}e^{\mu\operatorname{dist}(x,A)}\) in the absolutely
summable convolution kernel.  The weighted entrywise convolution norm in
\eqref{eq:V27-imported-router-locality} and the Schur--Young inequality give
the infinite-lattice bound.  Exact periodization preserves the same
weighted sum when distance is read on the torus, proving the finite-volume
statement.
\end{proof}

% =============================================================================
\section{The V26 likelihood score is an exact V20 incidence residual}
\label{sec:V27-incidence-routing}
% =============================================================================

Fix \((\alpha,x,W)\) and abbreviate the V35 trace law and actual temporal
fibre posterior by
\begin{equation}
 \nu:=\nu_W^{\rm tr},
 \qquad
 \dd p:=\ell^{-1}L\,\dd\nu,
 \qquad
 \Sigma:=\dd\log L.
 \label{eq:V27-fibre-laws}
\end{equation}
Let \(d_\nu^*\) and \(d_p^*\) be the weighted codifferentials on the same
fibre metric.

\begin{theorem}[Density-change codifferential and minimal incidence range]
\label{thm:V27-codifferential-incidence}
For every smooth fibre one-form \(\omega\),
\begin{equation}
 \boxed{
 d_p^*\omega
 =d_\nu^*\omega-\langle\Sigma,\omega\rangle.}
 \label{eq:V27-codifferential-change}
\end{equation}
Consequently the multiplication-contraction operator
\begin{equation}
 \mathcal B_{\rm inc}:L^2(p;T^*Y_W^{\rm red})\longrightarrow L^2(p),
 \qquad
 \mathcal B_{\rm inc}\omega:=\langle\Sigma,\omega\rangle
 \label{eq:V27-incidence-operator}
\end{equation}
is the exact conditional-law incidence residual.

For a finite command entrance
\(R:E_0\to L^2(p;T^*Y_W^{\rm red})\), let
\(\Pi_\one\) project onto constants and define
\begin{align}
 Y_{\rm inc}&:=(I-\Pi_\one)\mathcal B_{\rm inc}R,
 \label{eq:V27-centered-incidence-entrance}\
 G_{\rm inc}^{p}&:=Y_{\rm inc}^*Y_{\rm inc}.
 \label{eq:V27-actual-incidence-Gram}
\end{align}
Then
\begin{equation}
 \boxed{
 G_{\rm inc}^{p}[a,a]
 =\Var_p\bigl(\langle\Sigma,Ra\rangle\bigr),}
 \label{eq:V27-incidence-covariance}
\end{equation}
and the polar decomposition
\begin{equation}
 Y_{\rm inc}=U_{\rm inc}(G_{\rm inc}^{p})^{1/2}
 \label{eq:V27-incidence-polar}
\end{equation}
gives the minimal centered source range which must be appended to the V20
incidence atlas.
\end{theorem}

\begin{proof}
Write \(\dd\nu=Z_\nu^{-1}e^{-S_\nu}\dd m\).  Since
\(\dd p\) is proportional to
\(e^{-(S_\nu-\log L)}\dd m\), the weighted codifferential formula gives
\[
 d_p^*=d_m^*+\iota_{\nabla(S_\nu-\log L)}
 =d_\nu^*-\iota_{\nabla\log L},
\]
which is \eqref{eq:V27-codifferential-change}.  The constant projection in
\eqref{eq:V27-centered-incidence-entrance} subtracts the \(p\)-mean, so
direct multiplication proves \eqref{eq:V27-incidence-covariance}.  The
finite-dimensional polar decomposition gives
\eqref{eq:V27-incidence-polar}; its partial-isometry range is the smallest
closed range containing every centered likelihood-score response generated
by \(R\).
\end{proof}

\begin{proposition}[V26 nonzero incidence survives actual-law centering]
\label{prop:V27-nonzero-centered-incidence}
Take \(R\) to be the constant left-trivialized fibre commands.  If the V26
reference Gram \(\mathfrak F_{\alpha;x,W}\) is nonzero, then the actual
centered Gram \(G_{\rm inc}^{p}\) is nonzero.  Hence the positive-measure set
of histories in \Cref{cor:V26-positive-history-set} carries a genuine
nonzero V20 incidence-source range.
\end{proposition}

\begin{proof}
Because \(L>0\), the laws \(p\) and \(\nu\) have the same null sets.  If
\(G_{\rm inc}^{p}=0\), every left-trivialized component of
\(\Sigma=d\log L\) is \(p\)-almost surely constant and hence, by smoothness,
constant everywhere.  For a left-invariant vector field \(X\), a constant
identity \(X\log L=c_X\) integrates against Haar to \(c_X=0\).  Thus every
component vanishes and \(d\log L=0\), which would make the uncentered V26
Gram zero.  Contraposition proves the first assertion.  Apply
\Cref{cor:V26-positive-history-set} for the second.
\end{proof}

To place this source into the literal V20 matrix, let
\(\{P_j\}\) be cell projections on fibre one-forms and let
\(\{Q_i\}\) be any declared family of contractions on scalar residuals,
including multiplication by the V19 predictable bad sets.  Define
\begin{equation}
 \boxed{
 (\mathsf E_{\rm inc})_{ij}
 :=\norm{Q_i\mathcal B_{\rm inc}P_j}_{2\to2}.}
 \label{eq:V27-incidence-residual-matrix}
\end{equation}

\begin{corollary}[Exact V20 residual row]
\label{cor:V27-V20-residual-row}
For every one-form \(\omega=\sum_jP_j\omega\), entrywise in \(i\),
\begin{equation}
 \boxed{
 \norm{Q_i\mathcal B_{\rm inc}\omega}_{L^2(p)}
 \le\sum_j(\mathsf E_{\rm inc})_{ij}
          \norm{P_j\omega}_{L^2(p)}.}
 \label{eq:V27-V20-residual-row}
\end{equation}
Thus \(\mathsf E_{\rm inc}\) occupies the residual slot
\(\mathsf E\) in \eqref{eq:V20-source-incidence}.  Its row or column
summability is a quantitative local-likelihood problem; its algebraic
incidence is now exact.
\end{corollary}

\begin{proof}
Insert the resolution of the identity, apply \(Q_i\mathcal B_{\rm inc}\),
and use the triangle inequality and the definition
\eqref{eq:V27-incidence-residual-matrix}.  This is exactly the operator form
of the residual term in the V20 entrywise row.
\end{proof}

\begin{firewall}[The incidence source is not the density Woodbury column]
\label{fire:V27-incidence-not-Woodbury}
Equation \eqref{eq:V27-codifferential-change} licenses
\(Y_{\rm inc}\) as a conditional-law response source.  It does not identify
\(Y_{\rm inc}^*Y_{\rm inc}\) with
\(2(D_{\alpha,\beta})_+\).  The latter is licensed only by the signed Hessian
identity
\[
 D_{\alpha,\beta}
 =\Cov_pJ^{\rm t}-\mathbb E_pK^{\rm t}
 =\Hess\log r_{\alpha,\beta}.
\]
Therefore V27 inserts \(\mathsf E_{\rm inc}\) into the V20 residual atlas,
while the V21--V24 Woodbury channel continues to use a proved ceiling on the
\emph{assembled signed} tensor \(D\).  These are two different Grams.
\end{firewall}

% =============================================================================
\section{What is now populated and what still blocks the finite edge packet}
\label{sec:V27-packet-status}
% =============================================================================

At each finite regulator, V27 has populated the following actual-law
objects without a static/temporal identification:
\begin{equation}
 \boxed{
 \left(
  D_{\alpha,\beta},\quad
  Y_{\rm inc},\quad
  G_{\rm inc}^{p},\quad
  \mathsf E_{\rm inc},\quad
  \Omega_{\alpha,\beta}(B;\eta)
 \right).}
 \label{eq:V27-populated-actual-packet}
\end{equation}
The first object has the exact anisotropic linearization
\eqref{eq:V27-Wilson-D-prime}; the next three are the polar source and V20
residual licensed by \eqref{eq:V27-codifferential-change}; the last is the
local stock sufficient for the physical temporal-source lower frame
\eqref{eq:V27-actual-source-lower-frame}.

An honest V24/V25 \(G,H,J,H_x\) edge packet still requires a nonlinear
cofinal upper ceiling
\begin{equation}
 D_{\alpha,\beta_{\alpha,j}}\preceq a_{\alpha,j}I
 \label{eq:V27-needed-D-ceiling}
\end{equation}
with the response-matched physical scaling controlled.  Only then does the
V23 shifted identity column \(\sqrt{2a_{\alpha,j}}I\) defines the density
Woodbury channel to which the V24 recursion applies.  The likelihood polar
range may not be substituted for that column.

\begin{proposition}[Finite exact alternative at the new bottleneck]
\label{prop:V27-finite-alternative}
Fix a finite regulator and a finite command block \(B\).  Exactly one of the
following quantitative outcomes is available.
\begin{enumerate}[label=\textup{(\Alph*)},leftmargin=3.5em]
\item A declared constant \(A_*\) satisfies
      \(\sup_\eta\Omega_{\alpha,\beta}(B;\eta)\le A_*\).  Then the actual
      source lower frame \eqref{eq:V27-actual-source-lower-frame} holds.
\item The maximum local oscillation is larger than the proposed \(A_*\).
      Compactness returns configurations
      \((y_B^+,y_B^-,\eta)\) such that
      \begin{equation}
       \log\frac{f_{\alpha,\beta}(y_B^+,\eta)}
                     {f_{\alpha,\beta}(y_B^-,\eta)}>A_*.
       \label{eq:V27-Harnack-witness}
      \end{equation}
      This is an exact local Perron-likelihood witness to be expanded into
      connected temporal-spatial polymers; it is not a global massless
      state.
\end{enumerate}
Independently, if the V26 incidence Gram is nonzero, the polar vector of
\(G_{\rm inc}^{p}\) is a nonzero V20 source witness by
\Cref{prop:V27-nonzero-centered-incidence}.
\end{proposition}

\begin{proof}
The continuous function
\(\Omega_{\alpha,\beta}(B;\eta)\) attains its maximum on the compact exterior
configuration space.  If the maximum is at most \(A_*\), apply
\Cref{thm:V27-source-lower-frame}.  Otherwise choose a maximizing exterior
and maximizing/minimizing interior configurations, giving
\eqref{eq:V27-Harnack-witness}.  The final assertion is the polar
decomposition \eqref{eq:V27-incidence-polar}.
\end{proof}

\begin{assessment}[The exact upstream theorem now required]
\label{ass:V27-upstream-theorem}
A useful anisotropic half-space expansion must treat
\(\mathcal C_\alpha\) as the normalized electric reference dynamics and
expand only the magnetic insertions.  Its connected time strings are summed
with the temporal susceptibility \(\left(1-q_\alpha\right)^{-1}\).  The
smallness variable is therefore \(\zeta^{\rm Har}_{\alpha,\beta}\), or a
slightly larger multi-representation analogue, rather than
\(|\beta|\).  The theorem must produce both:
\begin{enumerate}[label=\textup{(\roman*)},leftmargin=3.5em]
\item a uniform conditional block oscillation for \(\log f\); and
\item a summable change of that conditional log likelihood under a distant
      exterior-link perturbation.
\end{enumerate}
The first enters \Cref{thm:V27-source-lower-frame}; the second bounds
\(\mathsf E_{\rm inc}\) and the V25 locality pencil.  A global
\(\sup f/\inf f\) estimate is neither requested nor expected.
\end{assessment}

% =============================================================================
\section{V27 ledger and revised frontier}
\label{sec:V27-ledger}
% =============================================================================

\begin{theorem}[Integrated V27 statement]
\label{thm:V27-integrated}
Preserving V1--V26, and without repeating the V35 harmonic-router proof, V27
proves:
\begin{enumerate}[label=\textup{(V27.\arabic*)},leftmargin=5.5em]
\item the exact \(q_\alpha=\eta_F(\alpha)^4\) temporal multiplier of the
      centered unit Wilson plaquette potential;
\item an analytic Perron-resolvent formula and the exact first magnetic
      derivative of the signed density defect;
\item impossibility of an \(\alpha\)-uniform \(O(|\beta|)\) ceiling for the
      literal anisotropic transfer;
\item the correct linear density and Harnack parameters, whose matched
      Hamiltonian limit is \(\lvert\gamma\rvert/(4\kappa C_F)\);
\item an exact conditional likelihood/covariance window needing only local
      oscillation;
\item a local-Harnack lower frame for the complete compensated temporal
      source, which uses the constituent identity row and survives every
      remote-score cancellation;
\item an exponentially decaying, volume-independent tail for the imported
      V35 harmonic router;
\item the exact weighted-codifferential identity which licenses the V26
      likelihood score as a V20 incidence residual;
\item a minimal polar source range, a nonzero actual-law centering theorem,
      and the explicit residual matrix \(\mathsf E_{\rm inc}\); and
\item a finite local-Harnack witness alternative and the precise nonlinear
      anisotropic polymer theorem required next.
\end{enumerate}
\end{theorem}

\begin{proof}
Items \textup{(V27.1)--(V27.2)} are
\Cref{lem:V27-plaquette-eigenmode,thm:V27-first-variation}.  Items
\textup{(V27.3)--(V27.4)} are
\Cref{thm:V27-no-uniform-beta-ceiling,cor:V27-correct-parameter}.  Item
\textup{(V27.5)} is \Cref{thm:V27-conditional-window}.  Items
\textup{(V27.6)--(V27.7)} are
\Cref{thm:V27-source-lower-frame,cor:V27-router-tail}.  Items
\textup{(V27.8)--(V27.9)} are
\Cref{thm:V27-codifferential-incidence,
prop:V27-nonzero-centered-incidence,cor:V27-V20-residual-row}.  Item
\textup{(V27.10)} is
\Cref{prop:V27-finite-alternative,ass:V27-upstream-theorem}.
\end{proof}

\begingroup
\small
\begin{longtable}{>{\raggedright\arraybackslash}p{0.28\textwidth}
                  >{\raggedright\arraybackslash}p{0.16\textwidth}
                  >{\raggedright\arraybackslash}p{0.48\textwidth}}
\toprule
\textbf{V27 row}&\textbf{Status}&\textbf{Advance or remaining obstruction}\\
\midrule
\endfirsthead
\toprule
\textbf{V27 row}&\textbf{Status}&\textbf{Advance or remaining obstruction}\\
\midrule
\endhead
Wilson temporal eigenmode
& \MGclosed
& Exact eigenvalue \(q_\alpha=\eta_F(\alpha)^4\), independent of volume.\\[0.35em]
Actual Perron signed-defect linearization
& \MGclosed
& Exact resolvent and formula
  \eqref{eq:V27-Wilson-D-prime}; covariance and curvature were assembled
  before taking a norm.\\[0.35em]
Uniform \(O(|\beta|)\) transfer from V9
& \MGfail
& Ruled out for the literal anisotropic family by the diverging exact first
  derivative; V9 itself is not contradicted.\\[0.35em]
Correct anisotropic parameter
& \MGclosed
& \(\zeta^{\rm den}=|\beta|q/(1-q)\) and
  \(\zeta^{\rm Har}=|\beta|/(1-q)\), with the same matched limit.\\[0.35em]
Global Perron Harnack window
& \MGreduced
& Replaced by the exact conditional local window
  \eqref{eq:V27-conditional-density-window}; no global oscillation is needed
  for local source acquisition.\\[0.35em]
Complete actual temporal-source floor
& \MGreduced
& Proved from a uniform local window by
  \eqref{eq:V27-actual-source-lower-frame}; the nonlinear local window is
  still an input.\\[0.35em]
Harmonic-router locality
& \MGclosed
& The exact V35 weighted kernel gives the exponential tail
  \eqref{eq:V27-router-tail}.\\[0.35em]
Likelihood-incidence source typing
& \MGclosed
& Exact codifferential difference, polar range, actual centered Gram, and
  V20 residual matrix.\\[0.35em]
Nonlinear anisotropic Perron expansion
& \MGopen
& Must control local oscillation and distant exterior influence uniformly in
  the correct susceptibility-weighted parameter.\\[0.35em]
Actual density ceiling and V24/V25 edge packet
& \MGopen
& The first variation is exact, but a cofinal nonlinear remainder and the
  resulting \(G,H,J,H_x\) margins are not proved.\\[0.35em]
Capacity, protected sectors, and OS limit
& \MGopen
& Complete bad capacity, toron/wrapping floors, local-algebra noncollapse,
  and a nontrivial continuum state remain.\\[0.35em]
Full Yang--Mills mass gap
& \MGopen
& A false uniform perturbation route is removed and the local incidence
  packet is typed, but the nonlinear cofinal and continuum rows remain.\\
\bottomrule
\end{longtable}
\endgroup

\begin{openproblem}[V28 mandate: nonlinear anisotropic Perron locality]
\label{op:V28}
\begin{enumerate}[label=\textup{(V28.\arabic*)},leftmargin=5.5em]
\item Expand powers of the literal transfer
      \(M_{e^{-\beta V/2}}\mathcal C_\alpha
      M_{e^{-\beta V/2}}\) around the normalized Markov reference
      \(\mathcal C_\alpha\).  Sum temporal strings with their exact
      nonconstant multipliers; do not use a \(\beta\)-only KP parameter.
\item Prove, or return the first explicit violating polymer for, a
      volume-uniform conditional block bound on
      \(\Omega_{\alpha,\beta}(B;\eta)\) in a declared interval of
      \(\zeta^{\rm Har}_{\alpha,\beta}\).
\item Differentiate the same connected expansion twice to obtain a
      weighted row ceiling for the assembled
      \(D=\Cov_pJ^{\rm t}-\mathbb E_pK^{\rm t}\).  The first coefficient
      must reproduce \eqref{eq:V27-Wilson-D-prime}.
\item Use the distant-exterior derivative of the conditional likelihood,
      together with \eqref{eq:V27-router-tail}, to bound the row and column
      sums of \(\mathsf E_{\rm inc}\).  Route its bad set through V20 if the
      uniform good-window branch fails.
\item Only after the signed ceiling is proved, load the shifted density
      column and the polar incidence source into their distinct slots,
      populate \(G,H,J,H_x\), and run both V25 pencils.
\item On edge or locality collapse, return the exact V25 survivor together
      with its density, incidence, router-tail, and physical-energy ledger.
      On the strict branch, feed the exponential return into the complete
      V20 incidence--capacity matrix.
\item Keep \(\gamma_{\alpha,j}/\alpha\), toron and wrapping modes, protected
      sectors, local-algebra noncollapse, and the reflection-positive OS
      limit explicit.  A small-\(\zeta\) lattice branch is not by itself the continuum mass gap.
\item The next companion-file audit remains V30.
\end{enumerate}
\end{openproblem}

\begin{firewall}[End-of-V27 claim boundary]
V27 computes the exact first magnetic variation of the actual anisotropic
Perron defect and proves that an \(\alpha\)-uniform \(O(|\beta|)\) import is
impossible.  It identifies the susceptibility-weighted parameter, replaces
the global Perron window by an exact local conditional window, proves a
conditional complete-source lower frame, and inserts the nonzero V26
likelihood innovation into the V20 residual atlas through a licensed
codifferential identity.  It does not prove the nonlinear cofinal local
Harnack or Hessian remainder, a strict uniform edge/locality margin, bad
capacity, protected-sector floors, or a nontrivial continuum OS state.
Therefore the full Yang--Mills mass gap remains open.
\end{firewall}

\section*{Version V27 append-only record}
\addcontentsline{toc}{section}{Version V27 append-only record}
The complete V26 mathematical source was retained before this continuation.
Its SHA--256 digest was
\begin{center}\scriptsize\ttfamily
8a93f1980c3c8c6a40b9985ae1a0f70e499340f889c802b4814f5577f0728983.
\end{center}
The V27 file also repairs one hidden form-feed character in the old V22
display, restoring the intended command
\(\texttt{\string\frac\{3\}\{2\string\alpha\}}\), and defines the mathematical
shorthand already used in earlier versions:
\(\texttt{\string\Hess},\texttt{\string\Cov},\texttt{\string\Var},
\texttt{\string\Ric},\texttt{\string\Tr},\texttt{\string\gap},
\texttt{\string\dist},\texttt{\string\Span},\texttt{\string\Dom},
\texttt{\string\op}\).  No mathematical claim is changed by these lexical
preamble repairs.  On the next iteration, retain this full file, append
before its unique active document-closing command,
increment the filename to \texttt{\_V28.tex}, and remove only the superseded
V27 file from this note series.  The next companion audit is V30.

% =============================================================================
% END APPENDED VERSION V27 MATERIAL
% =============================================================================

% =============================================================================
% BEGIN APPENDED VERSION V28 MATERIAL
% =============================================================================

\clearpage
\part{V28 --- Connected Perron cumulants and the all-order resonance gate}
\label{part:V28}

\section{V28 direction and claim boundary}
\label{sec:V28-direction}

V27 identified the correct anisotropic variables
\[
 \zeta_{\alpha,\beta}^{\rm Har}
 =\frac{|\beta|}{1-q_\alpha},
 \qquad
 \zeta_{\alpha,\beta}^{\rm den}
 =\frac{|\beta|q_\alpha}{1-q_\alpha},
 \qquad
 q_\alpha=\eta_F(\alpha)^4,
\]
but left the nonlinear Perron locality theorem open.  This continuation
goes upstream to the exact logarithmic Perron equation.  It derives every
Taylor coefficient as a connected conditional cumulant followed by the
reduced temporal resolvent, proves the correct susceptibility scaling at
each fixed order uniformly in volume, and computes the complete quadratic
Wilson resonance quotient.  The calculation identifies an exact
all-order coefficient estimate which would close the local Harnack,
signed-density, and incidence-locality rows.  That estimate is stated as a
named gate and is not silently assumed.

The V27 source is retained in full before this continuation, apart from two
typographic repairs recorded at the end.  No theorem below asserts the
continuum Yang--Mills mass gap.

% =============================================================================
\section{The centered logarithmic Perron fixed point}
\label{sec:V28-log-fixed-point}
% =============================================================================

Fix a nondegenerate finite lattice \(\Lambda\) and abbreviate
\[
 \mathcal C=\mathcal C_{\alpha,\Lambda},
 \qquad
 \mathfrak m=\mathfrak m_\Lambda,
 \qquad
 V=V_{M,\Lambda}.
\]
Let \(P_0g=\int g\,\dd\mathfrak m\) and \(Q_0=I-P_0\).  On the
mean-zero space define the reduced temporal resolvent
\begin{equation}
 \mathcal R_{\alpha,\Lambda}
 :=Q_0(I-\mathcal C)^{-1}Q_0
 =\sum_{t\ge0}\bigl(\mathcal C^t-P_0\bigr).
 \label{eq:V28-reduced-resolvent}
\end{equation}
The series converges at every fixed regulator.  It is not assigned a
volume-uniform operator norm on the full function space.

For the V27 Perron function \(f_{\alpha,\beta}\), put
\begin{equation}
 u_{\alpha,\beta}:=Q_0\log f_{\alpha,\beta},
 \qquad
 \widetilde V:=Q_0V.
 \label{eq:V28-centered-log-f}
\end{equation}

\begin{theorem}[Exact centered log-Perron equation]
\label{thm:V28-log-Perron}
At every fixed finite regulator, the analytic branch from
\Cref{thm:V27-first-variation} is equivalently characterized near
\(\beta=0\) by
\begin{align}
 u_{\alpha,\beta}
 &=Q_0\left[-\beta V+
       \log\bigl(\mathcal C e^{u_{\alpha,\beta}}\bigr)\right],
 \label{eq:V28-log-fixed-point}\\
 (I-\mathcal C)u_{\alpha,\beta}
 &=-\beta\widetilde V+\mathcal N_\alpha(u_{\alpha,\beta}),
 \label{eq:V28-resolvent-equation}\\
 u_{\alpha,\beta}
 &=-\beta\mathcal R_{\alpha,\Lambda}\widetilde V
   +\mathcal R_{\alpha,\Lambda}
      \mathcal N_\alpha(u_{\alpha,\beta}),
 \label{eq:V28-resolvent-fixed-point}
\end{align}
where
\begin{equation}
 \mathcal N_\alpha(g)
 :=Q_0\left[
       \log(\mathcal C e^g)-\mathcal Cg
      \right].
 \label{eq:V28-nonlinearity}
\end{equation}
The map \(\mathcal N_\alpha\) has vanishing constant and first derivatives
at zero.
\end{theorem}

\begin{proof}
The symmetric Perron equation in
\eqref{eq:V27-anisotropic-transfer}, after setting
\(f=e^{-\beta V/2}h\), is
\begin{equation}
 \lambda_{\alpha,\beta}f_{\alpha,\beta}
 =e^{-\beta V}\mathcal C f_{\alpha,\beta}.
 \label{eq:V28-unsymmetric-Perron}
\end{equation}
Write \(\log f=c+u\) with \(c=P_0\log f\).  Taking logarithms in
\eqref{eq:V28-unsymmetric-Perron} and applying \(Q_0\) removes both
\(\log\lambda\) and \(c\), proving
\eqref{eq:V28-log-fixed-point}.  Since \(\mathcal C1=1\),
the derivative at zero of \(g\mapsto\log(\mathcal Ce^g)\) is
\(g\mapsto\mathcal Cg\).  Subtracting this derivative proves
\eqref{eq:V28-resolvent-equation} and the two asserted vanishings.
Applying the inverse of \(I-\mathcal C\) on \(Q_0L^2(\mathfrak m)\)
gives \eqref{eq:V28-resolvent-fixed-point}.
\end{proof}

\begin{remark}[Why the centering is structural]
\label{rem:V28-centering}
The scalar Perron eigenvalue and the arbitrary multiplicative
normalization of \(f\) disappear before any estimate.  Thus every
denominator below belongs to a nonconstant temporal representation.
There is no formal division by the constant-mode denominator \(1-1\).
\end{remark}

% =============================================================================
\section{Conditional cumulants and the exact coefficient recursion}
\label{sec:V28-cumulants}
% =============================================================================

For smooth real functions \(g_1,\ldots,g_m\), define the conditional row
cumulant
\begin{equation}
 \kappa_m^\alpha(g_1,\ldots,g_m)(x)
 :=
 \left.
 \partial_{t_1}\cdots\partial_{t_m}
 \log\mathcal C
       \exp\left(\sum_{j=1}^m t_jg_j\right)(x)
 \right|_{t_1=\cdots=t_m=0}.
 \label{eq:V28-row-cumulant}
\end{equation}
In particular,
\begin{align}
 \kappa_1^\alpha(g)
 &=\mathcal Cg,
 \label{eq:V28-kappa-one}\\
 \kappa_2^\alpha(g,h)
 &=\mathcal C(gh)-(\mathcal Cg)(\mathcal Ch),
 \label{eq:V28-kappa-two}\\
 \kappa_3^\alpha(g,h,k)
 &=\mathcal C(ghk)
   -\mathcal C(gh)\mathcal Ck
   -\mathcal C(gk)\mathcal Ch
   -\mathcal C(hk)\mathcal Cg
   +2(\mathcal Cg)(\mathcal Ch)(\mathcal Ck).
 \label{eq:V28-kappa-three}
\end{align}

\begin{theorem}[Exact connected Taylor recursion]
\label{thm:V28-Taylor-recursion}
Write the fixed-regulator Taylor expansion
\begin{equation}
 u_{\alpha,\beta}
 =\sum_{n\ge1}\frac{\beta^n}{n!}u_{\alpha,n}.
 \label{eq:V28-u-series}
\end{equation}
Then
\begin{align}
 u_{\alpha,1}
 &=-\mathcal R_{\alpha,\Lambda}\widetilde V,
 \label{eq:V28-u-one}\\
 u_{\alpha,2}
 &=\mathcal R_{\alpha,\Lambda}Q_0
   \kappa_2^\alpha(u_{\alpha,1},u_{\alpha,1}),
 \label{eq:V28-u-two}\\
 u_{\alpha,3}
 &=\mathcal R_{\alpha,\Lambda}Q_0
   \left[
    3\kappa_2^\alpha(u_{\alpha,1},u_{\alpha,2})
    +\kappa_3^\alpha(u_{\alpha,1},
                     u_{\alpha,1},u_{\alpha,1})
   \right].
 \label{eq:V28-u-three}
\end{align}
More generally, if \(\Pi_n\) is the set of set partitions of
\(\{1,\ldots,n\}\), then for \(n\ge2\),
\begin{equation}
 \boxed{
 u_{\alpha,n}
 =\mathcal R_{\alpha,\Lambda}Q_0
  \sum_{\substack{\pi\in\Pi_n\\|\pi|\ge2}}
  \kappa_{|\pi|}^\alpha
  \bigl(u_{\alpha,|B|}:B\in\pi\bigr).}
 \label{eq:V28-general-recursion}
\end{equation}
The arguments on the right may be placed in any order because the
cumulants are symmetric.
\end{theorem}

\begin{proof}
Taylor expansion of the logarithmic moment generating function gives
\[
 \log(\mathcal Ce^g)
 =\sum_{m\ge1}\frac1{m!}
   \kappa_m^\alpha(g,\ldots,g)
\]
near \(g=0\).  The \(m=1\) term is \(\mathcal Cg\) and is already on the
left of \eqref{eq:V28-resolvent-equation}; the forcing
\(-\beta\widetilde V\) gives \eqref{eq:V28-u-one}.  The
multivariate Faà di Bruno formula says that the \(n\)-th derivative of the
composition is the sum indexed by set partitions, with one derivative
\(u_{\alpha,|B|}\) for every block \(B\).  Excluding the one-block
partition removes the linear term.  This proves
\eqref{eq:V28-general-recursion}; listing the partitions for \(n=2,3\)
gives \eqref{eq:V28-u-two}--\eqref{eq:V28-u-three}.
\end{proof}

\begin{corollary}[Every coefficient is an explicit temporal-string sum]
\label{cor:V28-temporal-strings}
On a product Peter--Weyl mode \(\psi_\rho\) with nonconstant multiplier
\(q_{\alpha,\rho}\),
\begin{equation}
 \mathcal R_{\alpha,\Lambda}\psi_\rho
 =\frac1{1-q_{\alpha,\rho}}\psi_\rho
 =\sum_{t\ge0}q_{\alpha,\rho}^t\psi_\rho.
 \label{eq:V28-mode-resolvent}
\end{equation}
Consequently \eqref{eq:V28-general-recursion}, together with the finite
Clebsch--Gordan decomposition of products, is an explicit sum of connected
temporal strings.  No transfer-gap constant has been inserted.
\end{corollary}

\begin{proof}
Product central convolution is diagonal on product Peter--Weyl
coefficients.  Equation \eqref{eq:V28-mode-resolvent} is the scalar
geometric series on every nonconstant mode.  Products in
\eqref{eq:V28-general-recursion} decompose into finitely many irreducible
coefficients at each fixed order.
\end{proof}

% =============================================================================
\section{The quadratic Wilson resonance cancels one temporal pole}
\label{sec:V28-quadratic}
% =============================================================================

\begin{theorem}[Exact second coefficient and channel quotient]
\label{thm:V28-quadratic-resonance}
For the unit Wilson potential, let \(q=q_\alpha\).  Then
\begin{align}
 u_{\alpha,1}
 &=-\frac1{1-q}\widetilde V,
 \label{eq:V28-Wilson-u-one}\\
 u_{\alpha,2}
 &=\frac1{(1-q)^2}
   \mathcal R_{\alpha,\Lambda}Q_0
   \left[\mathcal C(\widetilde V^2)-q^2\widetilde V^2\right].
 \label{eq:V28-Wilson-u-two}
\end{align}
At a fixed finite regulator decompose
\begin{equation}
 Q_0(\widetilde V^2)=\sum_{\sigma}\psi_\sigma,
 \qquad
 \mathcal C\psi_\sigma=q_{\alpha,\sigma}\psi_\sigma.
 \label{eq:V28-square-decomposition}
\end{equation}
Then
\begin{equation}
 \boxed{
 u_{\alpha,2}
 =\frac1{(1-q)^2}
   \sum_\sigma
   \frac{q_{\alpha,\sigma}-q^2}
        {1-q_{\alpha,\sigma}}\,\psi_\sigma.}
 \label{eq:V28-quadratic-channel}
\end{equation}
If the channel \(\sigma\) is fixed while \(\alpha\to\infty\), and its
total product Casimir is \(C_\sigma>0\), the V10 multiplier asymptotics
give
\begin{equation}
 \boxed{
 \frac{q_{\alpha,\sigma}-q_\alpha^2}
      {1-q_{\alpha,\sigma}}
 \longrightarrow
 \frac{8C_F-C_\sigma}{C_\sigma}.}
 \label{eq:V28-resonance-limit}
\end{equation}
Thus the quadratic coefficient has exactly two powers of the Wilson
susceptibility, not three.
\end{theorem}

\begin{proof}
\Cref{lem:V27-plaquette-eigenmode} and
\eqref{eq:V28-u-one} give
\eqref{eq:V28-Wilson-u-one}.  By
\eqref{eq:V28-kappa-two},
\[
 \kappa_2^\alpha(u_{\alpha,1},u_{\alpha,1})
 =\frac1{(1-q)^2}
  \left[\mathcal C(\widetilde V^2)-q^2\widetilde V^2\right],
\]
which proves \eqref{eq:V28-Wilson-u-two}.  Both \(Q_0\) and
\(\mathcal C\) preserve the decomposition
\eqref{eq:V28-square-decomposition}; applying
\eqref{eq:V28-mode-resolvent} gives
\eqref{eq:V28-quadratic-channel}.

For a fixed nonconstant product channel, V10 gives
\[
 q_{\alpha,\sigma}
 =1-C_\sigma s_\alpha+O(s_\alpha^2),
 \qquad
 q_\alpha^2=e^{-8C_Fs_\alpha}
 =1-8C_Fs_\alpha+O(s_\alpha^2).
\]
Division proves \eqref{eq:V28-resonance-limit}.  In particular the
numerator vanishes to the same first order as the new resolvent
denominator.  Taking the two factors \((1-q)^{-1}\) from the two copies of
\(u_{\alpha,1}\) accounts for all remaining susceptibility powers.
\end{proof}

\begin{corollary}[Disjoint plaquette pairs vanish exactly]
\label{cor:V28-disjoint-pairs}
If two plaquettes \(p,p'\) have disjoint link supports, then
\begin{equation}
 \kappa_2^\alpha(v_p,v_{p'})=0.
 \label{eq:V28-disjoint-covariance}
\end{equation}
Hence the canonical expansion of \(u_{\alpha,2}\) contains only equal or
link-overlapping plaquette pairs, uniformly in the volume.
\end{corollary}

\begin{proof}
The product temporal row makes functions of disjoint link increments
independent.  Therefore
\[
 \mathcal C(v_pv_{p'})
 =(\mathcal Cv_p)(\mathcal Cv_{p'})
\]
pointwise in the initial configuration.  Insert this equality in
\eqref{eq:V28-kappa-two}.
\end{proof}

\begin{firewall}[What the quadratic cancellation does not prove]
\label{fire:V28-quadratic-scope}
Equation \eqref{eq:V28-resonance-limit} is uniform over the finitely many
representation types generated by one fixed overlapping plaquette pair.
It does not control the number and size of all rooted cumulant trees as
their order tends to infinity.  That combinatorial estimate is the
remaining nonlinear gate.
\end{firewall}

% =============================================================================
\section{Connected spatial support at every Taylor order}
\label{sec:V28-connected-support}
% =============================================================================

For a function of product-link variables, its link support is the smallest
set of links on which it depends.  A family of such supports is connected
when its intersection graph is connected.

\begin{lemma}[Product-row cumulants are support connected]
\label{lem:V28-cumulant-connected}
Suppose \(\{1,\ldots,m\}=I\sqcup J\), both parts nonempty, and
\[
 \left(\bigcup_{i\in I}\operatorname{supp}g_i\right)
 \cap
 \left(\bigcup_{j\in J}\operatorname{supp}g_j\right)
 =\varnothing.
\]
Then
\begin{equation}
 \kappa_m^\alpha(g_1,\ldots,g_m)=0.
 \label{eq:V28-split-cumulant-zero}
\end{equation}
Moreover, \(\mathcal C\), \(P_0\), \(Q_0\), and
\(\mathcal R_{\alpha,\Lambda}\) do not enlarge the link support of a
nonconstant local function, after its scalar mean is assigned to the empty
support.
\end{lemma}

\begin{proof}
Under a product temporal row, the link increments in the two displayed
unions are independent.  The conditional moment generating function in
\eqref{eq:V28-row-cumulant} factors into an \(I\)-factor and a
\(J\)-factor.  Its logarithm is a sum, so the mixed derivative involving
all \(m\) variables vanishes.  Product convolution integrates only the
coordinates on which a function depends and therefore preserves support.
Haar projection produces a scalar.  Finally every term
\((\mathcal C^t-P_0)g\) in \eqref{eq:V28-reduced-resolvent} has support
contained in that of \(g\).
\end{proof}

\begin{proposition}[Canonical connected plaquette decomposition]
\label{prop:V28-connected-decomposition}
For every \(n\ge1\), the recursion
\eqref{eq:V28-general-recursion} yields a decomposition
\begin{equation}
 u_{\alpha,n}
 =\sum_{\mathscr P\in\mathfrak C_n(\Lambda)}
   u_{\alpha,n;\mathscr P},
 \label{eq:V28-cluster-decomposition}
\end{equation}
where \(\mathfrak C_n(\Lambda)\) consists of connected multisets of at most
\(n\) plaquettes.  Every term
\(u_{\alpha,n;\mathscr P}\), after subtracting its Haar mean, is supported
on the union of the links of \(\mathscr P\).  In particular, the
order-\(n\) contribution cannot connect two cell blocks whose plaquette
distance is greater than \(n\), up to the fixed enlargement used to attach
commands to plaquettes.
\end{proposition}

\begin{proof}
At order one,
\(\widetilde V=-\sum_pv_p\), so the assertion holds with one plaquette.
Assume it through order \(n-1\).  Choose a set partition in
\eqref{eq:V28-general-recursion} and one cluster term from every argument.
The total number of plaquettes, counted with multiplicity, is at most
\(\sum_{B\in\pi}|B|=n\).  By
\Cref{lem:V28-cumulant-connected}, the cumulant vanishes unless these
argument supports form a connected family.  In the nonzero case their
union is therefore a connected plaquette multiset.  Neither \(Q_0\) nor
\(\mathcal R_{\alpha,\Lambda}\) enlarges it.  This proves the induction.
A connected multiset of at most \(n\) elementary plaquettes has
plaquette-graph diameter at most \(n-1\), proving the last assertion.
\end{proof}

% =============================================================================
\section{Uniform cofinal scaling at every fixed order}
\label{sec:V28-fixed-order}
% =============================================================================

Use the canonical decomposition
\eqref{eq:V28-cluster-decomposition}.  For \(\mu\ge0\) and
\(k\in\{0,1,2\}\), define its rooted local interaction norm by
\begin{equation}
 \norm{g}_{\mu,k}^{\rm loc}
 :=
 \sup_{e\in E(\Lambda)}
 \sum_{\mathscr P:\,e\in\operatorname{supp}\mathscr P}
 e^{\mu\operatorname{diam}(\mathscr P)}
 \norm{g_{\mathscr P}}_{C^k}.
 \label{eq:V28-local-interaction-norm}
\end{equation}
The diameter is measured in the plaquette intersection graph.  At fixed
order this sum has a volume-independent number of rooted shapes.

\begin{theorem}[Fixed-order cofinal connected bound]
\label{thm:V28-fixed-order-bound}
Fix \(n\ge1\), \(\mu\ge0\), and \(k\le2\).  There are
\(\alpha_n<\infty\) and \(C_{n,\mu,k}<\infty\), independent of the
nondegenerate finite volume, such that
\begin{equation}
 \boxed{
 (1-q_\alpha)^n
 \norm{u_{\alpha,n}}_{\mu,k}^{\rm loc}
 \le C_{n,\mu,k}}
 \label{eq:V28-fixed-order-bound}
\end{equation}
for every \(\alpha\ge\alpha_n\).  Thus every fixed Taylor order has the
susceptibility power predicted by
\(\zeta_{\alpha,\beta}^{\rm Har}\).
\end{theorem}

\begin{proof}
Fix a rooted connected plaquette shape with at most \(n\) plaquettes.
Products of its Wilson matrix coefficients through order \(n\) lie in a
finite-dimensional product Peter--Weyl space
\(\mathscr W_n\), depending on \(n\) and the local lattice geometry but not
on the ambient volume.  On this space the V10 multiplier asymptotics give
\begin{equation}
 \mathcal C=I-s_\alpha\mathcal L_n+O(s_\alpha^2)
 \label{eq:V28-C-small-s}
\end{equation}
in every \(C^k\) norm under consideration.  The positive operator
\(\mathcal L_n\) has a strictly positive least eigenvalue on the
mean-zero part of \(\mathscr W_n\).  Hence
\begin{equation}
 \norm{\mathcal R_{\alpha,\Lambda}}_
       {Q_0\mathscr W_n\to Q_0\mathscr W_n,C^k}
 \le K_{n,k}s_\alpha^{-1}
 \label{eq:V28-R-small-s}
\end{equation}
for all sufficiently large \(\alpha\), with no volume dependence.

For \(m\ge2\), the multilinear operator
\(\kappa_m^\alpha\) vanishes when \(\mathcal C=I\), because the
corresponding row is deterministic.  Formula
\eqref{eq:V28-row-cumulant}, restricted to \(\mathscr W_n\), is a finite
polynomial in convolution multipliers and multiplication maps.  Combining
this observation with \eqref{eq:V28-C-small-s} gives
\begin{equation}
 \norm{\kappa_m^\alpha(g_1,\ldots,g_m)}_{C^k}
 \le K_{n,m,k}s_\alpha
       \prod_{j=1}^m\norm{g_j}_{C^k}
 \label{eq:V28-cumulant-small-s}
\end{equation}
whenever the total polynomial order of the arguments is at most \(n\).

We induct on \(n\).  Equation \eqref{eq:V28-u-one} and
\eqref{eq:V28-R-small-s} give
\(\norm{u_{\alpha,1}}_{\mu,k}^{\rm loc}=O(s_\alpha^{-1})\).
In a term of \eqref{eq:V28-general-recursion}, the induction hypothesis
and \(\sum_{B\in\pi}|B|=n\) give \(O(s_\alpha^{-n})\) for the product of
the argument norms.  The connected cumulant contributes the factor
\(s_\alpha\) in \eqref{eq:V28-cumulant-small-s}; the final resolvent
contributes \(s_\alpha^{-1}\).  The result is
\(O(s_\alpha^{-n})\).

By \Cref{prop:V28-connected-decomposition}, only connected rooted
plaquette shapes with at most \(n\) members occur.  Their number and their
diameters are bounded in terms of \(n\) and the lattice dimension, so
summing them in \eqref{eq:V28-local-interaction-norm} does not introduce a
volume factor.  Finally,
\[
 \frac{1-q_\alpha}{s_\alpha}
 =\frac{1-e^{-4C_Fs_\alpha}}{s_\alpha}
 \longrightarrow4C_F.
\]
Replacing \(s_\alpha^n\) by a constant multiple of
\((1-q_\alpha)^n\) completes the induction.
\end{proof}

\begin{corollary}[Fixed local Harnack polynomials are cofinally scaled]
\label{cor:V28-Harnack-polynomial}
For fixed \(N\) and a fixed link block \(B\), the Taylor polynomial
\begin{equation}
 u_{\alpha,\beta}^{[N]}
 :=\sum_{n=1}^N\frac{\beta^n}{n!}u_{\alpha,n}
 \label{eq:V28-truncated-u}
\end{equation}
satisfies
\begin{equation}
 \sup_\eta
 \operatorname{osc}_{y_B}
 u_{\alpha,\beta}^{[N]}(y_B,\eta)
 \le
 2|B|\sum_{n=1}^N
 \frac{C_{n,\mu,0}}{n!}
 \left(\zeta_{\alpha,\beta}^{\rm Har}\right)^n
 \label{eq:V28-truncated-Harnack}
\end{equation}
for sufficiently large \(\alpha\).  At order one this reproduces
\Cref{prop:V27-local-Harnack-linearization}.
\end{corollary}

\begin{proof}
Terms whose support misses \(B\) cancel from the oscillation.  For each
remaining cluster choose a link in its intersection with \(B\), overcount
by the at most \(|B|\) choices, and use
\(\operatorname{osc}g\le2\norm{g}_{C^0}\).
Apply \eqref{eq:V28-fixed-order-bound} term by term.
\end{proof}

\begin{assessment}[Exact content of the fixed-order theorem]
\label{ass:V28-fixed-order-scope}
\Cref{thm:V28-fixed-order-bound} proves that no new power of
\((1-q_\alpha)^{-1}\) is lost at any prescribed order.  Its constant
\(C_{n,\mu,k}\), however, may grow faster than exponentially in \(n\).
Therefore it does not supply a common positive radius in
\(\zeta_{\alpha,\beta}^{\rm Har}\).  Interchanging the cofinal limit with
the infinite Taylor sum would be circular at this point.
\end{assessment}

% =============================================================================
\section{The signed defect expansion and perturbative incidence locality}
\label{sec:V28-defect-locality}
% =============================================================================

\begin{proposition}[Exact nonlinear identity for the signed defect]
\label{prop:V28-exact-defect}
The V27 signed tensor satisfies
\begin{equation}
 \boxed{
 D_{\alpha,\beta}
 =\Hess\bigl(\beta V+u_{\alpha,\beta}\bigr).}
 \label{eq:V28-D-exact}
\end{equation}
Consequently its first two Taylor orders are
\begin{equation}
 \boxed{
 D_{\alpha,\beta}
 =-\beta\frac{q_\alpha}{1-q_\alpha}\Hess V
  +\frac{\beta^2}{2}\Hess u_{\alpha,2}
  +O_{\alpha,\Lambda}(\beta^3).}
 \label{eq:V28-D-second-order}
\end{equation}
Moreover, for every fixed \(n\), the order-\(n\) signed-defect
coefficient has a volume-uniform weighted local row bound after
multiplication by \((1-q_\alpha)^n\).
\end{proposition}

\begin{proof}
Equation \eqref{eq:V28-unsymmetric-Perron} gives
\[
 \log r_{\alpha,\beta}
 =\log\lambda_{\alpha,\beta}+\beta V+\log f_{\alpha,\beta}.
\]
The scalar parts \(\log\lambda\) and \(P_0\log f\) vanish under the
Hessian, proving \eqref{eq:V28-D-exact}.  Insert
\eqref{eq:V28-u-series} and
\eqref{eq:V28-Wilson-u-one}; since
\[
 \Hess(V+u_{\alpha,1})
 =-\frac{q_\alpha}{1-q_\alpha}\Hess V,
\]
equation \eqref{eq:V28-D-second-order} follows and agrees with
\eqref{eq:V27-Wilson-D-prime}.  The last assertion is
\Cref{thm:V28-fixed-order-bound} with \(k=2\).
\end{proof}

\begin{proposition}[Finite propagation of incidence coefficients]
\label{prop:V28-incidence-finite-propagation}
Consider the Perron-likelihood part of the score
\(\Sigma=\dd\log L\) in \eqref{eq:V27-fibre-laws}, localized in the V27
cell projections.  At Taylor order \(n\), its matrix coefficient between
cell blocks \(i\) and \(j\) vanishes when their plaquette distance exceeds
\(n+c_{\rm att}\), where \(c_{\rm att}\) is the fixed command/fibre
attachment range.  For each fixed \(n\), its weighted row and column sums
are volume independent after multiplication by
\((1-q_\alpha)^n\).
\end{proposition}

\begin{proof}
Differentiating a cluster term does not enlarge its link support.
\Cref{prop:V28-connected-decomposition} says that an order-\(n\) term is
supported on a connected multiset of at most \(n\) plaquettes.  It cannot
meet both attached cell blocks beyond the stated distance.  The fixed-order
weighted row and column bounds follow from
\eqref{eq:V28-local-interaction-norm} and
\Cref{thm:V28-fixed-order-bound}.  The finite-range trace and attachment
maps contribute only the fixed enlargement \(c_{\rm att}\).
\end{proof}

% =============================================================================
\section{The all-order connected Perron majorant}
\label{sec:V28-all-order-gate}
% =============================================================================

\begin{definition}[Connected Perron majorant gate]
\label{def:V28-CPM}
For constants \(A,L,\mu>0\) and a declared cofinal threshold
\(\alpha_0\), write
\(\operatorname{CPM}(A,L,\mu;\alpha_0)\) for the assertion that, for all
\(n\ge1\), all nondegenerate finite volumes, and all
\(\alpha\ge\alpha_0\),
\begin{equation}
 \boxed{
 \norm{u_{\alpha,n}}_{\mu,2}^{\rm loc}
 \le n!\,A L^{n-1}(1-q_\alpha)^{-n}.}
 \label{eq:V28-CPM}
\end{equation}
The parameters \(A,L,\mu,\alpha_0\) are independent of \(n\) and the
volume.
\end{definition}

\begin{theorem}[Compiler from the all-order majorant]
\label{thm:V28-CPM-compiler}
Assume \(\operatorname{CPM}(A,L,\mu;\alpha_0)\) and
\begin{equation}
 L\zeta_{\alpha,\beta}^{\rm Har}<1.
 \label{eq:V28-CPM-domain}
\end{equation}
Then the fixed-regulator Perron series converges absolutely in the local
\(C^2\) interaction norm and
\begin{equation}
 \boxed{
 \norm{u_{\alpha,\beta}}_{\mu,2}^{\rm loc}
 \le
 \frac{A\zeta_{\alpha,\beta}^{\rm Har}}
      {1-L\zeta_{\alpha,\beta}^{\rm Har}}.}
 \label{eq:V28-u-majorant}
\end{equation}
Consequently:
\begin{enumerate}[label=\textup{(\roman*)},leftmargin=3.5em]
\item every fixed block \(B\) has the volume-uniform conditional window
\begin{equation}
 \sup_\eta\Omega_{\alpha,\beta}(B;\eta)
 \le
 \frac{2|B|A\zeta_{\alpha,\beta}^{\rm Har}}
      {1-L\zeta_{\alpha,\beta}^{\rm Har}};
 \label{eq:V28-Harnack-from-CPM}
\end{equation}
\item the assembled signed tensor has the local ceiling
\begin{equation}
 \norm{D_{\alpha,\beta}}_{\mu,0}^{\rm loc}
 \le
 |\beta|\norm{V}_{\mu,2}^{\rm loc}
 +\frac{A\zeta_{\alpha,\beta}^{\rm Har}}
       {1-L\zeta_{\alpha,\beta}^{\rm Har}};
 \label{eq:V28-D-from-CPM}
\end{equation}
\item the Perron portion of the V27 matrix
\(\mathsf E_{\rm inc}\) has exponentially summable, volume-independent
row and column bounds; and
\item the window \eqref{eq:V28-Harnack-from-CPM} supplies the hypothesis of
\Cref{thm:V27-source-lower-frame}.
\end{enumerate}
\end{theorem}

\begin{proof}
By \eqref{eq:V28-CPM},
\[
 \sum_{n\ge1}\frac{|\beta|^n}{n!}
 \norm{u_{\alpha,n}}_{\mu,2}^{\rm loc}
 \le
 A\sum_{n\ge1}L^{n-1}
 \left(\frac{|\beta|}{1-q_\alpha}\right)^n,
\]
which is the geometric sum \eqref{eq:V28-u-majorant}.  At a fixed
regulator the sum agrees near zero with the analytic Perron branch.  Its
absolute convergence and the exact equation
\eqref{eq:V28-resolvent-fixed-point} continue that branch throughout
\eqref{eq:V28-CPM-domain}.

For a block oscillation, discard clusters missing \(B\), root each
remaining cluster at one of at most \(|B|\) links, and use
\(\operatorname{osc}g\le2\norm{g}_{C^0}\).  This proves
\eqref{eq:V28-Harnack-from-CPM}.  Equation
\eqref{eq:V28-D-exact} and two spatial derivatives give
\eqref{eq:V28-D-from-CPM}.  For two separated cell blocks, only clusters
meeting both blocks contribute.  The exponential weight in
\eqref{eq:V28-local-interaction-norm}, followed by the Schur test, gives
the row and column bounds in item \textup{(iii)}.  Item \textup{(iv)} is
the direct substitution of \eqref{eq:V28-Harnack-from-CPM} into
\Cref{thm:V27-source-lower-frame}.
\end{proof}

\begin{corollary}[What the strict branch would populate]
\label{cor:V28-strict-branch}
Under the assumptions of \Cref{thm:V28-CPM-compiler}, a scalar
pointwise ceiling may be chosen from
\eqref{eq:V28-D-from-CPM}.  The V23 shifted identity column is then
licensed for the density channel, while the polar range
\eqref{eq:V27-incidence-polar} remains the distinct incidence source.
These two columns may be loaded into their respective V24/V25 slots.
\end{corollary}

\begin{proof}
A local operator-norm bound implies the corresponding pointwise upper
quadratic-form bound.  The licensing statements are precisely the signed
density requirement in \eqref{eq:V27-needed-D-ceiling} and the
codifferential identity \eqref{eq:V27-codifferential-change}.  Their
sources are different, as emphasized in
\Cref{fire:V27-incidence-not-Woodbury}.
\end{proof}

\begin{proposition}[Exact finite witness if the majorant fails]
\label{prop:V28-coefficient-witness}
Fix proposed \(A,L,\mu>0\) and \(\alpha_0\), with \(A\) large enough for
the order-one bound on \(\alpha\ge\alpha_0\).  If
\(\operatorname{CPM}(A,L,\mu;\alpha_0)\) fails, there is a least order
\(n_*\ge2\), a finite nondegenerate regulator
\((\alpha_*,\Lambda_*)\) with \(\alpha_*\ge\alpha_0\), and a root link
\(e_*\) for which
\begin{equation}
 (1-q_{\alpha_*})^{n_*}
 \sum_{\substack{\mathscr P\in\mathfrak C_{n_*}(\Lambda_*)\\
                  e_*\in\operatorname{supp}\mathscr P}}
 e^{\mu\operatorname{diam}(\mathscr P)}
 \norm{u_{\alpha_*,n_*;\mathscr P}}_{C^2}
 >
 n_*!A L^{n_*-1}.
 \label{eq:V28-first-failing-root}
\end{equation}
Furthermore, because only finitely many rooted shapes and
Peter--Weyl channels occur at order \(n_*\), one may select a connected
cluster, a derivative of order at most two, and a nonconstant
Peter--Weyl coefficient whose normalized size exceeds the right-hand side
of \eqref{eq:V28-first-failing-root} divided by an explicit finite
dimension-and-shape factor.
\end{proposition}

\begin{proof}
The positive integers are well ordered, giving the least failing order.
At a fixed finite volume the supremum over root links in
\eqref{eq:V28-local-interaction-norm} is a maximum, which gives
\eqref{eq:V28-first-failing-root}.  At fixed order, connectedness leaves
finitely many rooted shapes, and their Wilson products lie in the
finite-dimensional space \(\mathscr W_{n_*}\) used in
\Cref{thm:V28-fixed-order-bound}.  The triangle inequality shows that at
least one cluster term exceeds the total threshold divided by the number
of terms.  Equivalence of the \(C^2\) norm with finitely many coefficient
and derivative-coordinate norms on \(\mathscr W_{n_*}\) then selects the
claimed coefficient, with a computable finite factor.
\end{proof}

\begin{assessment}[The precise remaining analytic task]
\label{ass:V28-remaining-gate}
The proof of \(\operatorname{CPM}(A,L,\mu;\alpha_0)\) must retain, before
absolute
values, the multiplier difference created by every connected cumulant.
At quadratic order this is
\((q_{\alpha,\sigma}-q_\alpha^2)/(1-q_{\alpha,\sigma})\).
At higher order the recursion produces rooted connected trees with one
resolvent edge and one deterministic-row cancellation at every internal
vertex.  Bounding resolvents separately from cumulants loses these
cancellations and can create a factorially divergent false majorant.
V29 must either organize the cancellations tree by tree or return the
least witness of \Cref{prop:V28-coefficient-witness}.
\end{assessment}

% =============================================================================
\section{V28 ledger and revised frontier}
\label{sec:V28-ledger}
% =============================================================================

\begin{theorem}[Integrated V28 statement]
\label{thm:V28-integrated}
Preserving V1--V27, V28 proves:
\begin{enumerate}[label=\textup{(V28.\arabic*)},leftmargin=5.5em]
\item the exact centered logarithmic Perron fixed point and reduced
      temporal resolvent;
\item the all-order conditional-cumulant recursion for every Taylor
      coefficient;
\item the exact quadratic Wilson channel quotient and cancellation of the
      apparent third susceptibility pole;
\item exact vanishing of disjoint plaquette cumulants and connected
      plaquette support at every order;
\item a volume-uniform cofinal susceptibility bound at each fixed Taylor
      order, including two spatial derivatives;
\item the exact second-order signed-defect expansion and fixed-order
      incidence propagation;
\item a single all-order connected Perron majorant which compiles into the
      V27 local Harnack, density, source-floor, and incidence-locality
      requirements; and
\item an explicit least-order finite coefficient witness if a proposed
      majorant fails.
\end{enumerate}
\end{theorem}

\begin{proof}
Items \textup{(V28.1)--(V28.2)} are
\Cref{thm:V28-log-Perron,thm:V28-Taylor-recursion}.  Item
\textup{(V28.3)} is \Cref{thm:V28-quadratic-resonance}.  Item
\textup{(V28.4)} is
\Cref{cor:V28-disjoint-pairs,lem:V28-cumulant-connected,
prop:V28-connected-decomposition}.  Item \textup{(V28.5)} is
\Cref{thm:V28-fixed-order-bound}.  Item \textup{(V28.6)} is
\Cref{prop:V28-exact-defect,prop:V28-incidence-finite-propagation}.
Items \textup{(V28.7)--(V28.8)} are
\Cref{thm:V28-CPM-compiler,prop:V28-coefficient-witness}.
\end{proof}

\begingroup
\small
\begin{longtable}{>{\raggedright\arraybackslash}p{0.28\textwidth}
                  >{\raggedright\arraybackslash}p{0.16\textwidth}
                  >{\raggedright\arraybackslash}p{0.48\textwidth}}
\toprule
\textbf{V28 row}&\textbf{Status}&\textbf{Advance or remaining obstruction}\\
\midrule
\endfirsthead
\toprule
\textbf{V28 row}&\textbf{Status}&\textbf{Advance or remaining obstruction}\\
\midrule
\endhead
Centered log-Perron equation
& \MGclosed
& Exact fixed point \eqref{eq:V28-resolvent-fixed-point}; scalar
  normalization and eigenvalue are removed before estimates.\\[0.35em]
All-order coefficient recursion
& \MGclosed
& Exact set-partition cumulant formula
  \eqref{eq:V28-general-recursion}.\\[0.35em]
Quadratic temporal resonance
& \MGclosed
& Exact channel quotient \eqref{eq:V28-quadratic-channel} and finite
  cofinal limit \eqref{eq:V28-resonance-limit}.\\[0.35em]
Spatial connectedness
& \MGclosed
& Disjoint cumulants vanish; order \(n\) uses connected multisets of at
  most \(n\) plaquettes.\\[0.35em]
Fixed-order cofinal locality
& \MGclosed
& For every fixed \(n\), the \(C^2\) local interaction coefficient has at
  most \(n\) susceptibility powers, uniformly in volume.\\[0.35em]
All-order connected majorant
& \MGopen
& The constants \(C_{n,\mu,2}\) have not yet been bounded by
  \(n!AL^{n-1}\); this is the exact V29 gate.\\[0.35em]
Actual local Harnack window
& \MGreduced
& Compiled by \eqref{eq:V28-Harnack-from-CPM} once the all-order gate is
  proved.\\[0.35em]
Actual signed density ceiling
& \MGreduced
& Exact identity and every fixed order are controlled; the convergent
  all-order local remainder remains conditional on CPM.\\[0.35em]
Incidence residual locality
& \MGreduced
& Finite propagation holds coefficientwise; exponential summability of the
  full score remains conditional on CPM.\\[0.35em]
V24/V25 finite edge packet
& \MGopen
& The density and incidence columns are typed distinctly, but the
  \(G,H,J,H_x\) margins await the all-order ceiling.\\[0.35em]
Capacity, protected sectors, and OS limit
& \MGopen
& Complete bad capacity, toron/wrapping floors, local-algebra noncollapse,
  and a nontrivial continuum state remain.\\[0.35em]
Full Yang--Mills mass gap
& \MGopen
& Fixed-order nonlinear resonance and locality are now rigorous; the
  all-order, finite-edge, and continuum rows remain.\\
\bottomrule
\end{longtable}
\endgroup

\begin{openproblem}[V29 mandate: rooted cumulant trees]
\label{op:V29}
\begin{enumerate}[label=\textup{(V29.\arabic*)},leftmargin=5.5em]
\item Expand \eqref{eq:V28-general-recursion} into rooted connected
      cumulant trees.  Attach to each internal vertex its deterministic-row
      multiplier difference and to each outgoing edge its exact
      nonconstant resolvent denominator.
\item Prove \(\operatorname{CPM}(A,L,\mu;\alpha_0)\) for explicit
      constants by bounding the paired
      vertex-edge quotients before taking absolute values.  Do not replace
      them with separate full-space norms.
\item If the tree majorant fails, return the least order, rooted connected
      plaquette shape, derivative coordinate, and Peter--Weyl channel from
      \Cref{prop:V28-coefficient-witness}.
\item On the strict CPM branch, insert
      \eqref{eq:V28-Harnack-from-CPM} into the complete temporal source
      floor and \eqref{eq:V28-D-from-CPM} into the shifted density column.
      Keep the incidence polar source in its separate slot.
\item Populate the V24 recursion objects \(G,H,J,H_x\), evaluate both V25
      pencils, and return the exact survivor if an edge or locality margin
      collapses.
\item Preserve the response-matched coefficient
      \(\gamma_{\alpha,j}/\alpha\), and route every uniform-window failure
      through the V20 predictable bad-capacity matrix.
\item Keep toron and wrapping modes, protected sectors, local-algebra
      noncollapse, reflection positivity, and the nontrivial OS continuum
      state explicit.  A convergent small-\(\zeta\) lattice branch alone is
      not the continuum mass gap.
\item V30 must begin by auditing the then-current SM and NS notes, as
      required by the five-version companion cadence, before choosing its
      next direction.
\end{enumerate}
\end{openproblem}

\begin{firewall}[End-of-V28 claim boundary]
V28 proves an exact nonlinear cumulant recursion, the complete quadratic
resonance cancellation, spatial connectedness, and volume-uniform cofinal
scaling at every fixed Taylor order.  It does not prove an order-uniform
connected-tree majorant.  Accordingly, the nonlinear local Harnack window,
the full signed density ceiling, and exponential incidence locality are
conditional only through the explicitly named CPM gate.  The V24/V25
strict edge packet, bad-capacity closure, protected-sector floors,
local-algebra noncollapse, and nontrivial reflection-positive continuum
state remain open.  Therefore the full Yang--Mills mass gap remains open.
\end{firewall}

\section*{Version V28 append-only record}
\addcontentsline{toc}{section}{Version V28 append-only record}
The complete V27 mathematical source was retained before this continuation.
Its SHA--256 digest was
\begin{center}\scriptsize\ttfamily
07000e5918c73a4c39eca944f7c86e79593032fe1bdd1a5fffa7692326f98b11.
\end{center}
Two typographic defects in the V27 ledger were repaired during the merge:
the duplicated density-row sentence was removed, the missing
\emph{Capacity, protected sectors, and OS limit} row label was restored,
and the split phrase \emph{small-\(\zeta\)} was joined.  These repairs
change no mathematical assertion.  On the next iteration, retain this full
file, append before its unique active document-closing command, increment
the filename to \texttt{\_V29.tex}, and remove only the superseded V28
file from this note series.  The next companion audit is V30.

% =============================================================================
% END APPENDED VERSION V28 MATERIAL
% =============================================================================

% =============================================================================
% BEGIN APPENDED VERSION V29 MATERIAL
% =============================================================================

\clearpage
\part{V29 --- Rooted cumulant trees, Gaussian excess, and the uniform
fusion gate}
\label{part:V29}

\section{Direction and nonduplication}
\label{sec:V29-direction}

V28 proved the exact Perron cumulant recursion and the correct
susceptibility scaling at every fixed order.  It did not control the
growth of the fixed-order constants.  The V10 source explains why that
last step cannot be obtained by a silent limit: its multiplier theorem is
uniform only on a declared finite Casimir core, whereas order \(n\)
Wilson products reach representations whose highest weights grow with
\(n\).

This continuation first expands the V28 recursion into exact rooted
hierarchies.  It then proves a stronger fixed-order fact specific to the
Wilson row: an \(m\)-fold connected row cumulant has Gaussian excess
\(s_\alpha^{m-1}\), not merely one power of \(s_\alpha\).  The resulting
tree power count shows that only binary trees saturate the V28
\((1-q_\alpha)^{-n}\) scale.  Finally, a weighted Fourier--interaction
norm reduces the all-order question to one explicit uniform
resolvent--cumulant inequality and proves, with constants, that this
inequality implies the V28 connected Perron majorant.

The uniform fusion inequality itself is not proved below.  Its
fixed-height restrictions are proved, and its two exact failure modes are
recorded.  Thus V29 advances the bottleneck without declaring the
Yang--Mills mass gap.

% =============================================================================
\section{Exact rooted hierarchies for the Perron coefficients}
\label{sec:V29-rooted-trees}
% =============================================================================

Retain the V28 abbreviations and set
\begin{equation}
 \mathcal L_\alpha:=-\mathcal R_{\alpha,\Lambda}\widetilde V,
 \qquad
 \mathcal B_{\alpha,m}(g_1,\ldots,g_m)
 :=\mathcal R_{\alpha,\Lambda}Q_0
   \kappa_m^\alpha(g_1,\ldots,g_m)
 \quad(m\ge2).
 \label{eq:V29-leaf-vertex}
\end{equation}
The symbol \(\mathcal L_\alpha\) in this section denotes the leaf value,
not the V28 small-\(s\) generator \(\mathcal L_n\).

\begin{definition}[Labelled Perron hierarchy]
\label{def:V29-hierarchy}
A labelled Perron hierarchy on a nonempty finite label set \(S\) is defined
recursively:
\begin{enumerate}[label=\textup{(\roman*)},leftmargin=3.5em]
\item if \(S=\{j\}\), it is a single leaf carrying \(j\);
\item if \(|S|\ge2\), its root has an unordered family of at least two
      children, whose label sets form a set partition of \(S\), and every
      child is a hierarchy on its own block.
\end{enumerate}
Write \(\mathfrak T(S)\) for this finite set and
\(\mathfrak T_n=\mathfrak T(\{1,\ldots,n\})\).  Evaluate a leaf as
\(\mathcal L_\alpha\) and an internal vertex with \(m\) children as
\(\mathcal B_{\alpha,m}\) applied to the child values.  The result is
denoted \(\operatorname{Val}_\alpha(T)\).
\end{definition}

\begin{theorem}[Exact rooted-tree formula]
\label{thm:V29-tree-formula}
For every \(n\ge1\),
\begin{equation}
 \boxed{
 u_{\alpha,n}
 =\sum_{T\in\mathfrak T_n}
   \operatorname{Val}_\alpha(T).}
 \label{eq:V29-tree-formula}
\end{equation}
Every hierarchy occurs once for its labelled leaf set; no additional
symmetry factor is hidden in \eqref{eq:V29-tree-formula}.
\end{theorem}

\begin{proof}
For \(n=1\), equation \eqref{eq:V28-u-one} is precisely the leaf
evaluation.  Suppose the assertion holds below order \(n\).  In
\eqref{eq:V28-general-recursion}, choose a set partition
\(\pi\) of \(\{1,\ldots,n\}\) with at least two blocks.  For every block
\(B\), insert the induction expansion of \(u_{\alpha,|B|}\), relabelling
its leaves by the elements of \(B\).  Applying
\(\mathcal B_{\alpha,|\pi|}\) creates one root whose children are those
block hierarchies.  Conversely, deleting the root of a hierarchy in
\(\mathfrak T_n\) recovers a unique set partition and a unique hierarchy
on each block.  This is a bijection term by term, proving the formula.
\end{proof}

\begin{proposition}[The hierarchy count is factorial times exponential]
\label{prop:V29-tree-count}
Let \(h_n=|\mathfrak T_n|\) and
\[
 H(z):=\sum_{n\ge1}h_n\frac{z^n}{n!}.
\]
As a formal power series,
\begin{equation}
 \boxed{
 H(z)=z+e^{H(z)}-1-H(z).}
 \label{eq:V29-tree-EGF}
\end{equation}
Its radius of convergence is
\begin{equation}
 r_{\rm tree}=2\log2-1>0.
 \label{eq:V29-tree-radius}
\end{equation}
In particular, for every \(0<r<r_{\rm tree}\), there is \(C_r<\infty\)
such that
\begin{equation}
 h_n\le C_r\,n!\,r^{-n}.
 \label{eq:V29-tree-count-bound}
\end{equation}
Thus the number of rooted cumulant histories is compatible with a CPM
bound; it is not the source of a superfactorial obstruction.
\end{proposition}

\begin{proof}
A hierarchy is either one labelled atom or an unordered set of at least
two hierarchies.  The labelled exponential formula gives
\eqref{eq:V29-tree-EGF}.  Equivalently,
\[
 z=2H-e^H+1.
\]
The inverse branch through \(H=0\) remains regular until
\(2-e^H=0\), namely \(H=\log2\), where
\(z=2\log2-1\).  The coefficients are nonnegative, so Pringsheim's
theorem places the first radial singularity on the positive real axis.
This proves \eqref{eq:V29-tree-radius}.  Cauchy's estimate on any smaller
circle gives \eqref{eq:V29-tree-count-bound}.
\end{proof}

% =============================================================================
\section{Wilson small-noise normal form}
\label{sec:V29-small-noise}
% =============================================================================

Write \(s=s_\alpha\).  The V22 calculation gives
\(s=3/(2\alpha)+O(\alpha^{-2})\).  The next lemma records the
fixed-order normal form needed to count connected powers.  It is stronger
than a moment bound but remains a fixed-order statement.

\begin{lemma}[Fixed-order Wilson Laplace normal form]
\label{lem:V29-Laplace-normal-form}
Fix an integer \(M\ge0\).  In a fixed exponential neighbourhood of the
identity, put \(U=\exp(\sqrt{s}\,Z)\).  There are even polynomials
\(a_1,\ldots,a_M\), integers \(N_M\), and constants
\(c,C_M,s_M>0\) such that the normalized one-link Wilson law has the
weak expansion
\begin{equation}
 \dd\nu_\alpha(\exp(\sqrt{s}\,Z))
 =
 \phi_8(Z)
 \left[
  1+\sum_{r=1}^M s^r a_r(Z)
  +s^{M+1}\mathcal E_{M,s}(Z)
 \right]\dd Z,
 \label{eq:V29-Laplace-normal-form}
\end{equation}
where
\[
 \phi_8(Z)=(4\pi)^{-4}e^{-\norm Z^2/4},
\]
and, on the rescaled normal neighbourhood,
\begin{equation}
 |\mathcal E_{M,s}(Z)|
 \le C_M(1+\norm Z)^{N_M}e^{\norm Z^2/16}.
 \label{eq:V29-Laplace-remainder}
\end{equation}
The contribution of its complement against a bounded test function is at
most \(C_Me^{-c/s}\) times the test-function norm.  The same statement
holds for any fixed finite product of links, with the product Gaussian and
constants depending on the fixed number of links.
\end{lemma}

\begin{proof}
The analytic class function
\[
 x(e^X)=1-\frac16\norm X^2+
        \sum_{r\ge2}x_{2r}(X)
\]
has only even homogeneous terms after inversion symmetrization.  The Haar
Jacobian in exponential coordinates is analytic and even.  Substitute
\(X=\sqrt{s}\,Z\), use
\(\alpha=3/(2s)+O(1)\), and Taylor-expand the exponent, Jacobian, and
normalizing constant through order \(M\).  The quadratic term is
\(-\norm Z^2/4\); normalized Gaussian integration determines the scalar
part of every \(a_r\).  Taylor's formula with integral remainder gives
\eqref{eq:V29-Laplace-remainder} after shrinking the fixed normal
neighbourhood if necessary.

The V11 global quadratic-well estimate gives
\(x(I)-x(U)\ge c_0>0\) outside that neighbourhood.  Since
\(\alpha\asymp s^{-1}\), its normalized contribution is
\(O(e^{-c/s})\).  Tensoring finitely many one-link expansions proves the
product statement.
\end{proof}

\begin{lemma}[Marked connected-graph degree]
\label{lem:V29-marked-graph-degree}
In the rescaled Wilson Laplace expansion, every connected contraction graph
which meets \(m\ge2\) distinct source vertices has total small-noise degree
at least \(m-1\).
\end{lemma}

\begin{proof}
A marked source vertex of Taylor valence \(r\ge1\) carries
\(s^{r/2}\).  An unmarked vertex of valence \(d=2\ell\) coming from the
nonquadratic Wilson action carries at least
\(s^{\ell-1}=s^{d/2-1}\); a Haar-Jacobian or normalization vertex carries
at least \(s^{d/2}\), which is better.  Split products in the exponential
expansion into separate unmarked vertices.  If the resulting connected
Gaussian contraction graph has \(E\) contraction edges and \(V_u\)
unmarked vertices, its small-noise degree \(D\) therefore satisfies
\[
 D\ge
 \frac12\sum_{\rm marked}r_v+
 \sum_{\rm unmarked}\left(\frac{d_v}{2}-1\right)
 =E-V_u.
\]
The graph has \(m+V_u\) vertices and is connected, so
\(E\ge m+V_u-1\).  Hence \(D\ge m-1\).  Constant source vertices do not
occur because the mixed derivative of a connected cumulant annihilates
them.
\end{proof}

\begin{theorem}[Gaussian excess of connected Wilson-row cumulants]
\label{thm:V29-Gaussian-excess}
Fix \(m\ge2\), \(k\ge0\), a finite link bank, and a finite-dimensional
Peter--Weyl space \(\mathscr W\) on that bank.  There are
\(K_{\mathscr W,m,k}<\infty\) and \(\alpha_{\mathscr W,m,k}\) such that
for all \(g_1,\ldots,g_m\in\mathscr W\),
\begin{equation}
 \boxed{
 \norm{\kappa_m^\alpha(g_1,\ldots,g_m)}_{C^k}
 \le
 K_{\mathscr W,m,k}s_\alpha^{m-1}
 \prod_{j=1}^m\norm{g_j}_{C^k}}
 \label{eq:V29-Gaussian-excess}
\end{equation}
for \(\alpha\ge\alpha_{\mathscr W,m,k}\), after replacing the
finite-dimensional norm on the right by any fixed equivalent jet norm if
needed.  If the supports split, the cumulant is exactly zero as in
\Cref{lem:V28-cumulant-connected}.
\end{theorem}

\begin{proof}
It is enough, by equivalence of norms on \(\mathscr W\), to prove the
estimate on a fixed basis and after any \(k\) left-invariant derivatives
in the initial configuration.  Such derivatives commute with product
convolution and replace the \(g_j\) by elements of another fixed
finite-dimensional jet bank.

Use \Cref{lem:V29-Laplace-normal-form} with \(M=m-1\), and introduce
source variables \(t_1,\ldots,t_m\) in
\[
 \log\int
 \exp\left\{
  \sum_{j=1}^m t_j
  \bigl[g_j(x\exp(\sqrt{s}Z))-g_j(x)\bigr]
 \right\}\dd\nu_\alpha(Z).
\]
Taylor-expand each source in \(\sqrt{s}Z\) and use Wick's rule for the
leading Gaussian.  Taking the logarithm retains exactly the connected
contraction diagrams.  By \Cref{lem:V29-marked-graph-degree}, every diagram
meeting all \(m\) marked sources has small-noise degree at least \(m-1\).
Therefore the coefficient of \(t_1\cdots t_m\) vanishes at every power
\(s^0,\ldots,s^{m-2}\).

Taylor's remainder in \eqref{eq:V29-Laplace-remainder} is dominated by an
integrable Gaussian times a fixed polynomial.  The exterior contribution
\(e^{-c/s}\) is \(O(s^m)\).  Hence the first possibly nonzero coefficient
is bounded by \(K_{\mathscr W,m,k}s^{m-1}\), uniformly in the base point.
Multilinearity and finite-dimensional norm equivalence give
\eqref{eq:V29-Gaussian-excess}.  Exact split-support vanishing is
\Cref{lem:V28-cumulant-connected}.
\end{proof}

\begin{corollary}[Fixed-bank resolvent--cumulant scaling]
\label{cor:V29-resolvent-cumulant-fixed}
Under the hypotheses of \Cref{thm:V29-Gaussian-excess}, and after
enlarging the fixed Peter--Weyl bank to contain all products,
\begin{equation}
 \boxed{
 \norm{
  \mathcal R_{\alpha,\Lambda}Q_0
  \kappa_m^\alpha(g_1,\ldots,g_m)}_{C^k}
 \le
 K'_{\mathscr W,m,k}s_\alpha^{m-2}
 \prod_{j=1}^m\norm{g_j}_{C^k}.}
 \label{eq:V29-resolvent-cumulant-fixed}
\end{equation}
For \(m=2\) the bound is order one; for \(m\ge3\) it has a positive
small-noise gain.
\end{corollary}

\begin{proof}
On the mean-zero part of a fixed Peter--Weyl bank, the V10 asymptotics give
\(\norm{\mathcal R_{\alpha,\Lambda}}\le C_{\mathscr W}s_\alpha^{-1}\),
as in \eqref{eq:V28-R-small-s}.  Multiply this estimate by
\eqref{eq:V29-Gaussian-excess}.
\end{proof}

\begin{firewall}[Fixed order is essential in the Gaussian-excess theorem]
\label{fire:V29-fixed-order}
The constants in \eqref{eq:V29-Gaussian-excess} depend on the cumulant
order and the Peter--Weyl height.  The proof takes a Laplace expansion to a
depth depending on \(m\).  It does not show
\(K_{\mathscr W,m,2}\le m!K^m\) uniformly while the representation height
also grows.  Treating it as such would repeat the precise V10 core-to-tail
error.
\end{firewall}

% =============================================================================
\section{Exact susceptibility power of each rooted tree}
\label{sec:V29-tree-power}
% =============================================================================

For \(T\in\mathfrak T_n\), let \(I(T)\) be its number of internal
vertices and \(m_v\ge2\) the number of children at an internal vertex
\(v\).

\begin{lemma}[Hierarchy Euler identity]
\label{lem:V29-tree-Euler}
Every \(T\in\mathfrak T_n\) satisfies
\begin{equation}
 \sum_{v\ {\rm internal}}(m_v-1)=n-1,
 \qquad
 \sum_{v\ {\rm internal}}(m_v-2)=n-1-I(T).
 \label{eq:V29-tree-Euler}
\end{equation}
\end{lemma}

\begin{proof}
The tree has \(n+I(T)\) vertices and hence \(n+I(T)-1\) edges.  Counting
edges as children of internal vertices gives
\(\sum_vm_v=n+I(T)-1\).  Subtract \(I(T)\), and then subtract it once
more, to obtain the two identities.
\end{proof}

\begin{theorem}[Exact fixed-tree susceptibility power]
\label{thm:V29-tree-power}
Fix a hierarchy \(T\in\mathfrak T_n\), a local rooted plaquette shape
generated by its leaves, and \(k\le2\).  For sufficiently large
\(\alpha\),
\begin{equation}
 \boxed{
 \norm{\operatorname{Val}_\alpha(T)}_{C^k}
 \le C_{T,k}\,
 s_\alpha^{-\left(I(T)+1\right)}.}
 \label{eq:V29-tree-power}
\end{equation}
Equivalently, since \(1-q_\alpha\asymp s_\alpha\),
\[
 \operatorname{Val}_\alpha(T)
 =O_T\!\left((1-q_\alpha)^{-I(T)-1}\right).
\]
In particular,
\begin{equation}
 I(T)+1\le n,
 \label{eq:V29-tree-power-max}
\end{equation}
with equality if and only if every internal vertex is binary.  Thus only
binary cumulant trees saturate the V28 power
\((1-q_\alpha)^{-n}\).
\end{theorem}

\begin{proof}
Each of the \(n\) leaves contains one reduced resolvent and is
\(O(s_\alpha^{-1})\).  By
\Cref{cor:V29-resolvent-cumulant-fixed}, an internal vertex of arity
\(m_v\) contributes \(O(s_\alpha^{m_v-2})\), after the norms of its
children are factored out.  Therefore the total power is
\[
 -n+\sum_v(m_v-2)
 =-n+n-1-I(T)
 =-\bigl(I(T)+1\bigr),
\]
using \Cref{lem:V29-tree-Euler}.  Since every \(m_v\ge2\),
\(\sum_v(m_v-1)\ge I(T)\); the first identity in
\eqref{eq:V29-tree-Euler} gives \(I(T)\le n-1\).  Equality holds exactly
when every \(m_v=2\).
\end{proof}

\begin{corollary}[The nonlinear pole problem is binary]
\label{cor:V29-binary-frontier}
At every fixed Taylor order, all nonbinary Perron hierarchies have at
least one positive power of \(s_\alpha\) relative to the V28 CPM
normalization.  An order-uniform proof may therefore take the binary
resolvent covariance
\begin{equation}
 \mathcal R_{\alpha,\Lambda}Q_0
 \kappa_2^\alpha(g,h)
 \label{eq:V29-binary-operator}
\end{equation}
as its marginal vertex.  Higher cumulants still require an analytic
majorant, but they do not create the worst susceptibility pole.
\end{corollary}

\begin{proof}
For a nonbinary tree, at least one \(m_v-2\) is positive, so
\eqref{eq:V29-tree-Euler} gives \(I(T)+1<n\).  The binary case has
\(I(T)=n-1\) and saturates \eqref{eq:V29-tree-power}.
\end{proof}

% =============================================================================
\section{Exact Fourier projection of a cumulant vertex}
\label{sec:V29-fusion}
% =============================================================================

The Gaussian-excess proof is local in representation height.  To identify
the missing uniform statement, let \(P_{\boldsymbol R}\) denote a product
Peter--Weyl projection and
\[
 q_{\alpha,\boldsymbol R}
 :=\prod_{e}\eta_{R_e}(\alpha).
\]
For a nonconstant product representation,
\begin{equation}
 P_{\boldsymbol R}\mathcal R_{\alpha,\Lambda}
 =\frac1{1-q_{\alpha,\boldsymbol R}}P_{\boldsymbol R}.
 \label{eq:V29-product-resolvent}
\end{equation}

\begin{lemma}[Möbius formula for a row cumulant]
\label{lem:V29-cumulant-Mobius}
For \(m\ge2\),
\begin{equation}
 \boxed{
 \kappa_m^\alpha(g_1,\ldots,g_m)
 =
 \sum_{\pi\in\Pi_m}
 (-1)^{|\pi|-1}(|\pi|-1)!
 \prod_{B\in\pi}
 \mathcal C\left(\prod_{j\in B}g_j\right).}
 \label{eq:V29-cumulant-Mobius}
\end{equation}
Consequently, on a nonconstant output channel \(\boldsymbol R\),
\begin{equation}
 \boxed{
 P_{\boldsymbol R}\mathcal RQ_0
 \kappa_m^\alpha(g_1,\ldots,g_m)
 =
 \frac{1}{1-q_{\alpha,\boldsymbol R}}
 \sum_{\pi\in\Pi_m}
 (-1)^{|\pi|-1}(|\pi|-1)!
 P_{\boldsymbol R}
 \prod_{B\in\pi}
 \mathcal C\left(\prod_{j\in B}g_j\right).}
 \label{eq:V29-resolved-fusion-vertex}
\end{equation}
Every entry on the right is a finite sum of V10 multipliers and
Clebsch--Gordan contraction coefficients when the inputs are
Peter--Weyl polynomials.
\end{lemma}

\begin{proof}
The moment--cumulant inversion formula on the lattice of set partitions is
\eqref{eq:V29-cumulant-Mobius}.  Apply the diagonal identity
\eqref{eq:V29-product-resolvent} to its nonconstant output projection.
Finite Peter--Weyl multiplication gives the last assertion.
\end{proof}

\begin{corollary}[The V28 quotient is the binary fusion vertex]
\label{cor:V29-binary-fusion}
If \(g,h\) are temporal eigenmodes with multipliers \(q_g,q_h\), then on
an output component \(\psi_{\boldsymbol R}=P_{\boldsymbol R}(gh)\),
\begin{equation}
 P_{\boldsymbol R}\mathcal RQ_0\kappa_2^\alpha(g,h)
 =
 \frac{q_{\alpha,\boldsymbol R}-q_gq_h}
      {1-q_{\alpha,\boldsymbol R}}\,
 \psi_{\boldsymbol R}.
 \label{eq:V29-binary-fusion-quotient}
\end{equation}
For \(g=h=\widetilde V/(1-q_\alpha)\), this is exactly
\eqref{eq:V28-quadratic-channel}.
\end{corollary}

\begin{proof}
Equation \eqref{eq:V28-kappa-two} gives
\[
 P_{\boldsymbol R}\kappa_2^\alpha(g,h)
 =(q_{\alpha,\boldsymbol R}-q_gq_h)
   P_{\boldsymbol R}(gh).
\]
Apply \eqref{eq:V29-product-resolvent}.
\end{proof}

\begin{proposition}[Fixed-height fusion cancellation]
\label{prop:V29-fixed-height-fusion}
Fix \(m\), a finite input and output representation bank, and
\(k\le2\).  The operator in
\eqref{eq:V29-resolved-fusion-vertex} obeys
\begin{equation}
 \norm{
 P_{\boldsymbol R}\mathcal RQ_0
 \kappa_m^\alpha}_{(C^k)^m\to C^k}
 \le C_{m,\boldsymbol R,k}s_\alpha^{m-2}
 \label{eq:V29-fixed-fusion-cancellation}
\end{equation}
uniformly over that bank for sufficiently large \(\alpha\).
Equivalently, the numerator in
\eqref{eq:V29-resolved-fusion-vertex} is divisible through order
\(m-1\) in the Wilson small-noise expansion before the resolvent
denominator is divided out.
\end{proposition}

\begin{proof}
This is the Peter--Weyl projection of
\Cref{cor:V29-resolvent-cumulant-fixed}.  The phrase
\emph{divisible through order \(m-1\)} means that the coefficients of
\(s^0,\ldots,s^{m-2}\) in the cumulant numerator vanish, which is exactly
\Cref{thm:V29-Gaussian-excess}.  The nonconstant fixed output denominator
is \(C_{\boldsymbol R}s_\alpha+O(s_\alpha^2)\).
\end{proof}

\begin{assessment}[Why multiplier-by-multiplier absolute values are
forbidden]
\label{ass:V29-no-termwise-absolute}
The summands in \eqref{eq:V29-resolved-fusion-vertex} are individually of
order one while their Möbius sum is of order
\(s_\alpha^{m-1}\).  Taking absolute values before assembling the
partition sum loses \(m-1\) small-noise powers.  The required uniform
estimate is therefore a bound on the assembled fusion vertex, not on its
separate moment terms.
\end{assessment}

% =============================================================================
\section{A local analytic Fourier interaction norm}
\label{sec:V29-Fourier-norm}
% =============================================================================

For an \(\mathrm{SU}(3)\) representation \(R=(p,q)\), put
\[
 \operatorname{ht}(R)=p+q.
\]
For a product representation
\(\boldsymbol R=(R_e)_{e\in X}\), set
\[
 \operatorname{ht}(\boldsymbol R)=\sum_{e\in X}\operatorname{ht}(R_e),
 \quad
 C_E(\boldsymbol R)=\sum_{e\in X}C_2(R_e),
 \quad
 d_{\boldsymbol R}=\prod_{e\in X}d_{R_e}.
\]
If \(g_X\) is supported on the finite link set \(X\), define
\begin{equation}
 \norm{g_X}_{\rho,2}^{\rm FA}
 :=
 \sum_{\boldsymbol R}
 e^{\rho\operatorname{ht}(\boldsymbol R)}
 \bigl(1+C_E(\boldsymbol R)\bigr)
 d_{\boldsymbol R}
 \norm{\widehat g_X(\boldsymbol R)}_1,
 \label{eq:V29-Fourier-algebra-norm}
\end{equation}
where \(\norm{\cdot}_1\) is the matrix trace norm.  For a canonical local
interaction \(g=\sum_Xg_X\), put
\begin{equation}
 \norm{g}_{\rho,\mu,2}^{\rm FI}
 :=
 \sup_e\sum_{X\ni e}
 e^{\mu\operatorname{diam}X}
 \norm{g_X}_{\rho,2}^{\rm FA}.
 \label{eq:V29-Fourier-interaction-norm}
\end{equation}

\begin{lemma}[Fourier algebra and \(C^2\) control on one support]
\label{lem:V29-FA-algebra}
For every fixed \(\rho>0\):
\begin{enumerate}[label=\textup{(\roman*)},leftmargin=3.5em]
\item \(\norm{g_X}_{C^2}\le c_G\norm{g_X}_{\rho,2}^{\rm FA}\);
\item there is \(c_{\rm alg}<\infty\), depending only on
      \(\mathrm{SU}(3)\), such that
\begin{equation}
 \norm{g_Xh_Y}_{\rho,2}^{\rm FA}
 \le c_{\rm alg}
 \norm{g_X}_{\rho,2}^{\rm FA}
 \norm{h_Y}_{\rho,2}^{\rm FA}.
 \label{eq:V29-FA-product}
\end{equation}
\end{enumerate}
The lemma makes no free claim about summing all connected support mergers;
that spatial combinatorics is part of the UFC inequality below.
\end{lemma}

\begin{proof}
The unweighted trace-class Fourier norm is the Fourier algebra norm and
dominates the uniform norm.  Two left-invariant derivatives of a matrix
coefficient are bounded by a group constant times
\(1+C_2(R)\), proving \textup{(i)}.

Clebsch--Gordan decomposition is implemented by unitary intertwiners, so
trace norm is contractive under each block compression.  If
\(R\) occurs in \(R_1\otimes R_2\), highest-weight addition and the
Casimir formula give
\[
 \operatorname{ht}(R)
 \le\operatorname{ht}(R_1)+\operatorname{ht}(R_2),
 \qquad
 1+C_2(R)
 \le c_G(1+C_2(R_1))(1+C_2(R_2)).
\]
Summing the Fourier blocks proves \eqref{eq:V29-FA-product}.
\end{proof}

% =============================================================================
\section{The uniform fusion-cumulant gate and its compiler}
\label{sec:V29-UFC}
% =============================================================================

Set
\begin{equation}
 \delta_\alpha:=1-q_\alpha.
 \label{eq:V29-delta}
\end{equation}
The exact Wilson leaf identity is
\begin{equation}
 \boxed{
 \delta_\alpha
 \norm{\mathcal R_{\alpha,\Lambda}\widetilde V}_
      {\rho,\mu,2}^{\rm FI}
 =
 \norm{\widetilde V}_{\rho,\mu,2}^{\rm FI}
 =:a_{\rho,\mu}.}
 \label{eq:V29-leaf-bound}
\end{equation}

\begin{definition}[Uniform fusion-cumulant gate]
\label{def:V29-UFC}
For \(K,\rho,\mu,\alpha_0>0\), write
\(\operatorname{UFC}(K,\rho,\mu;\alpha_0)\) for the following assembled
operator inequality.  For every \(m\ge2\), every nondegenerate finite
volume, every \(\alpha\ge\alpha_0\), and all finite local Fourier
interactions \(g_1,\ldots,g_m\),
\begin{equation}
 \boxed{
 \norm{
  \mathcal R_{\alpha,\Lambda}Q_0
  \kappa_m^\alpha(g_1,\ldots,g_m)}
 _{\rho,\mu,2}^{\rm FI}
 \le
 m!K^{m-1}\delta_\alpha^{m-2}
 \prod_{j=1}^m
 \norm{g_j}_{\rho,\mu,2}^{\rm FI}.}
 \label{eq:V29-UFC}
\end{equation}
The cumulant is decomposed by its connected support union before the local
interaction norm is taken.
\end{definition}

\begin{remark}[The gate is an explicit V10 fusion inequality]
\label{rem:V29-UFC-explicit}
By \eqref{eq:V29-resolved-fusion-vertex} and
\Cref{lem:V29-FA-algebra}, the left side of
\eqref{eq:V29-UFC} is the weighted trace-norm sum of explicit finite
expressions
\[
 \frac1{1-q_{\alpha,\boldsymbol R}}
 P_{\boldsymbol R}
 \sum_{\pi\in\Pi_m}
 (-1)^{|\pi|-1}(|\pi|-1)!
 \prod_{B\in\pi}
 \mathcal C\left(\prod_{j\in B}g_j\right).
\]
Every scalar entering these expressions is one of the character integrals
\(\eta_R(\alpha)\) from \eqref{eq:V10-character-multiplier}; every other
entry is an \(\mathrm{SU}(3)\) Clebsch--Gordan contraction.  Thus UFC is
not an unspecified cluster-expansion hypothesis.  It is a concrete
uniform inequality for the exact Wilson fusion atlas.
\end{remark}

\begin{theorem}[UFC implies the V28 connected Perron majorant]
\label{thm:V29-UFC-compiler}
Assume \(\operatorname{UFC}(K,\rho,\mu;\alpha_0)\), and let
\(a=a_{\rho,\mu}>0\) be \eqref{eq:V29-leaf-bound}.  Then, for every
\(n\ge1\), every nondegenerate finite volume, and every
\(\alpha\ge\alpha_0\),
\begin{equation}
 \boxed{
 \norm{u_{\alpha,n}}_{\rho,\mu,2}^{\rm FI}
 \le
 n!\,(2a)(16Ka)^{n-1}\delta_\alpha^{-n}.}
 \label{eq:V29-UFC-coefficient-bound}
\end{equation}
In particular, because the FI norm dominates the V28 local \(C^2\) norm,
\begin{equation}
 \operatorname{CPM}(2c_Ga,16Ka,\mu;\alpha_0)
 \label{eq:V29-CPM-from-UFC}
\end{equation}
holds after absorbing the fixed norm-comparison constant into the first
parameter.
\end{theorem}

\begin{proof}
Put \(\beta=\delta_\alpha z\) and treat the V28 fixed point as a formal
power series in \(z\).  If
\[
 y(z)
 :=\sum_{n\ge1}
 \frac{\delta_\alpha^n}{n!}
 \norm{u_{\alpha,n}}_{\rho,\mu,2}^{\rm FI}z^n
\]
denotes its coefficient majorant, then
\eqref{eq:V29-leaf-bound}, \eqref{eq:V29-UFC}, and the cumulant expansion
of \(\mathcal N_\alpha\) give coefficientwise
\begin{align}
 y
 &\preceq
 az+
 \sum_{m\ge2}
 K^{m-1}\delta_\alpha^{m-2}y^m
 \nonumber\\
 &\preceq
 az+\frac{Ky^2}{1-Ky},
 \label{eq:V29-scalar-majorization}
\end{align}
because \(0<\delta_\alpha<1\).  Set \(x=Ky\) and \(c=Ka\).
The scalar majorant
\begin{equation}
 x=cz+\frac{x^2}{1-x}
 \label{eq:V29-scalar-equation}
\end{equation}
has the branch
\begin{equation}
 x(z)
 =\frac{1+cz-\sqrt{1-6cz+c^2z^2}}4
 \label{eq:V29-scalar-solution}
\end{equation}
through zero, with nonnegative coefficients.  On
\(|z|=(16c)^{-1}\), coefficient positivity and
\eqref{eq:V29-scalar-solution} give
\(|x(z)|\le x((16c)^{-1})<1/8\).  Cauchy's estimate therefore yields
\[
 [z^n]\,y(z)
 \le\frac1{8K}(16Ka)^n
 =2a(16Ka)^{n-1}.
\]
By the rooted recursion, formal coefficient majorization is valid at
every order.  This proves \eqref{eq:V29-UFC-coefficient-bound}.
\Cref{lem:V29-FA-algebra} then gives
\eqref{eq:V29-CPM-from-UFC}.
\end{proof}

\begin{corollary}[Uniform nonlinear Perron packet under UFC]
\label{cor:V29-UFC-packet}
Under \(\operatorname{UFC}(K,\rho,\mu;\alpha_0)\), whenever
\begin{equation}
 16Ka\,\zeta_{\alpha,\beta}^{\rm Har}<1,
 \label{eq:V29-UFC-domain}
\end{equation}
the following V28 consequences are simultaneously valid with explicit
geometric-series constants:
\begin{enumerate}[label=\textup{(\roman*)},leftmargin=3.5em]
\item a volume-uniform conditional local Harnack window;
\item a local \(C^2\) Perron logarithm and signed-density ceiling;
\item exponentially summable Perron incidence rows and columns; and
\item the complete compensated temporal-source floor of V27.
\end{enumerate}
The density shifted-identity column and incidence polar column remain
distinct.
\end{corollary}

\begin{proof}
Apply \Cref{thm:V29-UFC-compiler} and then
\Cref{thm:V28-CPM-compiler,cor:V28-strict-branch}.  The distinct source
typing is \Cref{fire:V27-incidence-not-Woodbury}.
\end{proof}

% =============================================================================
\section{What is proved about UFC and what can fail}
\label{sec:V29-UFC-status}
% =============================================================================

\begin{proposition}[Every fixed fusion window satisfies UFC with a
window-dependent constant]
\label{prop:V29-fixed-UFC}
Fix a cumulant order \(M\), a representation-height cutoff \(H\), and a
maximum connected support size \(N\).  There are
\(K_{M,H,N,\rho,\mu}<\infty\) and
\(\alpha_{M,H,N}<\infty\) such that \eqref{eq:V29-UFC} holds for
\(2\le m\le M\), inputs and outputs of height at most \(H\), connected
supports of size at most \(N\), and
\(\alpha\ge\alpha_{M,H,N}\).
\end{proposition}

\begin{proof}
There are finitely many rooted support shapes under the three cutoffs.
On every shape, the relevant input products occupy a finite-dimensional
Peter--Weyl bank.  Apply
\Cref{cor:V29-resolvent-cumulant-fixed}, replace
\(s_\alpha^{m-2}\) by a constant multiple of
\(\delta_\alpha^{m-2}\), and take the maximum of the finitely many
norm-equivalence, fusion, and support-merger constants.  Choose
\(K_{M,H,N,\rho,\mu}\) large enough that each of those constants is at
most \(m!K_{M,H,N,\rho,\mu}^{m-1}\) for \(2\le m\le M\).
\end{proof}

\begin{theorem}[Finite witness and escaping obstruction for UFC]
\label{thm:V29-UFC-failure}
Fix proposed \(K,\rho,\mu,\alpha_0\).  If
\(\operatorname{UFC}(K,\rho,\mu;\alpha_0)\) fails, then there are a finite
cumulant order, a finite connected support shape, a finite regulator
\(\alpha\ge\alpha_0\), finite Peter--Weyl input polynomials, and a
nonconstant output channel for which the explicit projection
\eqref{eq:V29-resolved-fusion-vertex} violates \eqref{eq:V29-UFC}.

More globally, fix \(\rho,\mu,\alpha_0\).  If every finite
order/height/support window has the bound of
\Cref{prop:V29-fixed-UFC} but no finite \(K\) satisfies UFC, then there is
a sequence of such finite witnesses whose normalized UFC ratios tend to
infinity.  Along that sequence at least one of cumulant order,
representation height, or connected support size tends to infinity.
For a Wilson Perron Taylor coefficient of fixed total order \(n\), no such
escape is possible.
\end{theorem}

\begin{proof}
Failure of the multilinear operator inequality is a strict inequality for
some finite local inputs.  Peter--Weyl polynomials are dense in the FI
norm, so a sufficiently close polynomial approximation preserves the
strict violation.  Projecting the nonzero excess onto its countable
Peter--Weyl sum selects a finite nonconstant output channel after a further
arbitrarily small truncation.  This proves the finite-witness statement.

If no global \(K\) exists, choose for each integer \(r\) a finite witness
whose normalized ratio is greater than \(r\).  Were its cumulant orders,
heights, and support sizes all bounded along a subsequence, the
corresponding fixed-window proposition, together with continuity on
compact \(\alpha\)-intervals and its cofinal bound, would give a uniform
finite ratio, a contradiction.  Hence at least one cutoff escapes.

At fixed Wilson Taylor order, there are at most \(n\) plaquettes in a
connected cluster, and multiplication by their fundamental or
antifundamental matrix coefficients reaches only a finite representation
bank by \Cref{thm:V6-finite-Casimir-propagation}.  Thus order, height, and
support are all bounded.
\end{proof}

\begin{assessment}[The exact V29 bottleneck]
\label{ass:V29-bottleneck}
The fixed-window theorem proves every finite restriction of UFC, but its
constant may deteriorate with \(M,H,N\).  The missing statement is one
simultaneous constant \(K\) in the normalized form
\eqref{eq:V29-UFC}, with the factorial and
\(\delta_\alpha^{m-2}\) already extracted.  There are now two legitimate
routes:
\begin{enumerate}[label=\textup{(\roman*)},leftmargin=3.5em]
\item prove a uniform connected steepest-descent bound for the exact
      Wilson character integrals and Clebsch--Gordan contractions in the
      analytic Fourier norm; or
\item find the first finite resolved fusion vertex, or an escaping
      high-weight sequence, in \Cref{thm:V29-UFC-failure}.
\end{enumerate}
Neither route is replaced by the fixed-\(R\) asymptotic
\eqref{eq:V10-symbol-ratio}.
\end{assessment}

% =============================================================================
\section{V29 ledger and next audited frontier}
\label{sec:V29-ledger}
% =============================================================================

\begin{theorem}[Integrated V29 statement]
\label{thm:V29-integrated}
Preserving V1--V28, V29 proves:
\begin{enumerate}[label=\textup{(V29.\arabic*)},leftmargin=5.5em]
\item the exact labelled rooted-hierarchy expansion of every Perron Taylor
      coefficient;
\item a factorial-times-exponential bound on the number of such
      hierarchies;
\item the fixed-order Wilson Laplace normal form and
      \(s_\alpha^{m-1}\) Gaussian excess of an \(m\)-fold connected row
      cumulant;
\item the exact susceptibility exponent \(I(T)+1\) of a rooted hierarchy
      and the fact that only binary trees saturate order \(n\);
\item the Möbius/fusion formula for every resolved cumulant vertex and its
      exact V28 binary specialization;
\item a local analytic Fourier norm controlling two
      derivatives;
\item an explicit uniform fusion-cumulant inequality which implies CPM
      with the constants \(A=2c_Ga\) and \(L=16Ka\); and
\item fixed-window UFC plus an exact finite-versus-escaping failure
      certificate for the remaining uniform step.
\end{enumerate}
\end{theorem}

\begin{proof}
Items \textup{(V29.1)--(V29.2)} are
\Cref{thm:V29-tree-formula,prop:V29-tree-count}.  Item
\textup{(V29.3)} is
\Cref{lem:V29-Laplace-normal-form,thm:V29-Gaussian-excess}.  Item
\textup{(V29.4)} is \Cref{thm:V29-tree-power}.  Item
\textup{(V29.5)} is
\Cref{lem:V29-cumulant-Mobius,cor:V29-binary-fusion}.  Item
\textup{(V29.6)} is \Cref{lem:V29-FA-algebra}.  Item
\textup{(V29.7)} is \Cref{thm:V29-UFC-compiler}.  Item
\textup{(V29.8)} is
\Cref{prop:V29-fixed-UFC,thm:V29-UFC-failure}.
\end{proof}

\begingroup
\small
\begin{longtable}{>{\raggedright\arraybackslash}p{0.28\textwidth}
                  >{\raggedright\arraybackslash}p{0.16\textwidth}
                  >{\raggedright\arraybackslash}p{0.48\textwidth}}
\toprule
\textbf{V29 row}&\textbf{Status}&\textbf{Advance or remaining obstruction}\\
\midrule
\endfirsthead
\toprule
\textbf{V29 row}&\textbf{Status}&\textbf{Advance or remaining obstruction}\\
\midrule
\endhead
Rooted Perron hierarchy
& \MGclosed
& Exact labelled formula \eqref{eq:V29-tree-formula}; hierarchy count has a
  positive exponential-generating-function radius.\\[0.35em]
Gaussian cumulant excess
& \MGclosed
& At every fixed order and representation bank, an \(m\)-cumulant begins
  at \(s_\alpha^{m-1}\).\\[0.35em]
Susceptibility power count
& \MGclosed
& A tree with \(I\) internal vertices has pole \(I+1\); only binary trees
  attain the order-\(n\) pole.\\[0.35em]
Resolved fusion formula
& \MGclosed
& Exact Möbius sum divided by the nonconstant output denominator in
  \eqref{eq:V29-resolved-fusion-vertex}.\\[0.35em]
Fixed fusion windows
& \MGclosed
& Every finite order/height/support restriction satisfies the normalized
  UFC inequality with a window-dependent constant.\\[0.35em]
Uniform fusion-cumulant gate
& \MGopen
& Simultaneous factorial/exponential control as order, height, and support
  grow is not proved.\\[0.35em]
Connected Perron majorant
& \MGreduced
& Follows with explicit constants from UFC by
  \eqref{eq:V29-UFC-coefficient-bound}.\\[0.35em]
Actual Harnack, density, and incidence packet
& \MGreduced
& All compile on the strict UFC branch; they remain unproved without the
  uniform fusion bound.\\[0.35em]
V24/V25 finite edge packet
& \MGopen
& \(G,H,J,H_x\) still await the nonlinear signed-density and incidence
  margins.\\[0.35em]
Capacity, protected sectors, and OS limit
& \MGopen
& Complete bad capacity, toron/wrapping floors, local-algebra noncollapse,
  and a nontrivial continuum state remain.\\[0.35em]
Full Yang--Mills mass gap
& \MGopen
& The all-order problem is now an explicit uniform Wilson fusion estimate,
  followed by the finite-edge and continuum rows.\\
\bottomrule
\end{longtable}
\endgroup

\begin{openproblem}[V30 mandate: companion audit and uniform Wilson fusion]
\label{op:V30}
\begin{enumerate}[label=\textup{(V30.\arabic*)},leftmargin=5.5em]
\item Begin by reading the then-current retained SM--Einstein and NS notes
      in \texttt{latex\_notes}.  Record their filenames, latest internal
      versions, and only those proved rows whose hypotheses transfer to
      the YM fusion or capacity problem.
\item After that audit, attack \eqref{eq:V29-UFC} in the exact Fourier
      norm.  Preserve the Möbius sum before absolute values and split the
      character atlas into a growing Casimir core and a genuine
      high-weight tail.
\item On the core, make the dependence of the connected
      steepest-descent constants on cumulant order and representation
      height explicit.  On the tail, use the actual V10 character
      integrals, not fixed-representation extrapolation.
\item Either prove UFC with explicit \(K,\rho,\mu,\alpha_0\), or return the
      finite/escaping alternative of
      \Cref{thm:V29-UFC-failure}.  If a binary vertex is the first failure,
      report its input and output representations and the quotient
      \eqref{eq:V29-binary-fusion-quotient}.
\item On a strict UFC branch, invoke
      \Cref{thm:V29-UFC-compiler,cor:V29-UFC-packet}, populate the distinct
      density and incidence columns, and evaluate the V24/V25
      \(G,H,J,H_x\) margins and both edge pencils.
\item Route a failed local window through the V20 predictable
      bad-capacity matrix.  Do not replace the complete source or its
      incidence residual by a trace-only Gram.
\item Preserve \(\gamma_{\alpha,j}/\alpha\), toron and wrapping modes,
      protected sectors, local-algebra noncollapse, reflection positivity,
      and the nontrivial OS continuum limit.  A strict small-response
      lattice branch is not the full mass gap.
\end{enumerate}
\end{openproblem}

\begin{firewall}[End-of-V29 claim boundary]
V29 proves the exact rooted-tree expansion, its non-superfactorial
combinatorics, the fixed-order Gaussian excess of Wilson cumulants, and the
fact that binary trees alone carry the marginal susceptibility pole.  It
also proves with explicit constants that the uniform fusion-cumulant gate
implies the V28 connected Perron majorant and all of its local nonlinear
consequences.  It does not prove the gate uniformly as representation
height, support, and cumulant order grow.  Therefore the actual all-order
Harnack, signed-density, incidence, V24/V25 edge, capacity, protected
sector, and continuum OS rows remain open, and the full Yang--Mills mass
gap is not proved.
\end{firewall}

\section*{Version V29 append-only record}
\addcontentsline{toc}{section}{Version V29 append-only record}
The complete V28 mathematical source was retained before this continuation.
Its SHA--256 digest was
\begin{center}\scriptsize\ttfamily
958b5070fe94d443ff0b65b2edea55a9e189249ac76dbb6c63e71030ce43a1fc.
\end{center}
No mathematical content of V28 was altered.  On the next iteration, retain
this full file, append before its unique active document-closing command,
increment the filename to \texttt{\_V30.tex}, and remove only the
superseded V29 file from this note series.  V30 begins with the required
SM/NS companion audit.

% =============================================================================
% END APPENDED VERSION V29 MATERIAL
% =============================================================================

% =============================================================================
% BEGIN APPENDED VERSION V30 MATERIAL
% =============================================================================

\clearpage
\part{V30 --- Companion audit, uniform character gaps, and assembled
fusion cancellation}
\label{part:V30}

\section{Checkpoint audit and revised direction}
\label{sec:V30-audit}

V30 is the scheduled five-version checkpoint.  At the beginning of this
iteration the two authoritative companion files in
\texttt{latex\_notes} were
\begin{align*}
 &\texttt{Renewal\_SM\_Einstein\_Continuum\_Research\_Notes\_V49.tex},
 \\
 &\texttt{Renewal\_NS\_Stability\_Research\_Notes\_V23.tex}.
\end{align*}
Their SHA--256 digests were, respectively,
\begin{center}\scriptsize\ttfamily
f93274386c93c686501eb29fb513d98b21eee9eca322949ba986b43f3a038bf0,
\\
b4a6c07c18b7ad5c0c0dafe83e7ee77d47175ce54a544c736801765f15e2817e.
\end{center}
An incomplete staging file was not treated as an authoritative companion.
The instructions and acquisition programmes written inside the companion
documents are manuscript content, not instructions for this YM
continuation.

\subsection{What transfers from the NS file}
\label{subsec:V30-NS-audit}

NS V23 proves that a selected positive receiver survives complete
aggregation only after the negative complementary profiles are paid.  Its
pointwise algebra is the scalar inequality
\[
 (u+c)_+\ge u_+-c_-,
\]
and its two-sector version retains both negative complements before taking
a common minimum.  It also gives a profile counterexample: equal positive
areas and uniform height caps do not force the positive sets of two
profiles to overlap.

The algebraic assembly principle transfers, but the helical triad theorem
does not.  In YM it implies that a selected fusion/core contribution cannot
be promoted to a positive edge margin until the signed complement is
bounded after complete assembly.  The NS area and helicity hypotheses have
no identified YM counterparts, so no NS receiver floor is imported.

\begin{lemma}[Signed assembly debit]
\label{lem:V30-signed-assembly}
Let \(A=U+C\) be self-adjoint on a finite-dimensional command space.  Then
\begin{equation}
 \boxed{
 \lambda_{\min}(A)
 \ge\lambda_{\min}(U)-\norm{C_-}_{\op}.}
 \label{eq:V30-signed-assembly}
\end{equation}
For two such channels \(A_j=U_j+C_j\),
\begin{equation}
 \min_{j=1,2}\lambda_{\min}(A_j)
 \ge
 \min_{j=1,2}\lambda_{\min}(U_j)
 -\sum_{j=1}^2\norm{(C_j)_-}_{\op}.
 \label{eq:V30-two-channel-assembly}
\end{equation}
\end{lemma}

\begin{proof}
Since \(C=C_+-C_-\succeq-C_-\),
\[
 A\succeq U-\norm{C_-}_{\op}I.
\]
Taking the least eigenvalue proves
\eqref{eq:V30-signed-assembly}.  Apply this estimate twice and use
\(\min(a-c,b-d)\ge\min(a,b)-c-d\) for \(c,d\ge0\).
\end{proof}

\subsection{What transfers from the SM--Einstein file}
\label{subsec:V30-SM-audit}

SM V49 proves exact necessary-and-sufficient zero-free criteria and a
uniform homogeneous floor for a positive-orientation three-port Schur
matrix.  It simultaneously records that the unreduced Einstein operator at
gravitational glancing remains open because the reconstruction uses a
singular inverse.  Earlier in the same retained file, V23 proves an
abstract weighted block majorant and a finite-to-infinite inverse
certificate with separate read, target-leakage, high-domain, and adjacent
transport debits.

The abstract finite-to-infinite perturbation lemma is mathematically
portable, but its role is already occupied in this YM manuscript by the
V24/V25 edge pencils and their tail ledgers.  Re-importing it would
duplicate rather than advance the programme.  The V49 warning does
transfer directionally: even a future zero-free fusion or Schur factor
cannot license the full YM inverse at a singular edge until the unreduced
reconstruction and its high-channel tail are bounded.  The Einstein bath
speeds, port Grams, and glancing factors do not satisfy the YM Perron
hypotheses and are not imported as estimates.

\begin{theorem}[V30 companion-audit decision]
\label{thm:V30-audit-decision}
The checkpoint yields the following exact decisions.
\begin{enumerate}[label=\textup{(\roman*)},leftmargin=3.5em]
\item Preserve the assembled Möbius cumulant before absolute values, and
      preserve the complete signed V24/V25 edge matrix before taking a
      positive part.
\item A core fusion floor and a tail fusion floor, considered separately,
      do not prove a common command range; their complement residual must
      be paid as in \eqref{eq:V30-signed-assembly}.
\item A reduced UFC or Schur factor, even when zero-free, does not by itself
      control the unreduced edge reconstruction at a singular channel.
\item No theorem in the audited companion versions proves the YM
      all-order fusion estimate, bad-capacity bound, protected-sector
      floor, or OS continuum limit.
\end{enumerate}
\end{theorem}

\begin{proof}
Items \textup{(i)--(ii)} are the matched algebraic content of the NS V23
aggregate residual and \Cref{lem:V30-signed-assembly}.  Item \textup{(iii)}
is the exact scope distinction in the SM V49 port firewall, applied only as
a logical implication about reduced versus unreduced operators.  The
operator families and hypotheses in both companion ledgers are different
from \eqref{eq:V29-UFC}, so neither supplies any of the physical YM rows in
\textup{(iv)}.
\end{proof}

\subsection{Late companion drift during validation}
\label{subsec:V30-late-drift}

After the V30 mathematics had been assembled, but before final validation,
the parallel note series advanced.  A second read-only drift check inspected
\begin{align*}
 &\texttt{Renewal\_SM\_Einstein\_Continuum\_Research\_Notes\_V50.tex},\\
 &\hspace{1.5em}\text{SHA--256 }
 \texttt{67482e0544826d3da57c181581ffc3cd91daf39dc29f247d7a8ea201833c24cb},\\
 &\texttt{Renewal\_NS\_Stability\_Research\_Notes\_V25.tex},\\
 &\hspace{1.5em}\text{SHA--256 }
 \texttt{0d80d813253256366be9a515b5c7da5647a29f8dc8f2c2c76fe12b92bc653415}.
\end{align*}
This does not alter the historical V49/V23 snapshot above; it records
changes which arrived while V30 was being checked.

SM V50 closes its own five-variable gravitational-glancing inverse by
writing the unreduced matrix directly at the singular set, proving its
determinant nonzero without using the singular free inverse, and then using
compactness away from and through that set.  It also proves that the
unweighted total bath-energy resolvent is false at continuous spectrum while
a weighted outgoing graph norm is uniform.  The type-correct YM consequence
is an acquisition rule, not an imported estimate: when the programme returns
to the V24/V25 edge packet, it must evaluate the complete unreduced edge
symbol directly on every singular channel and must state separately the
topology in which reconstruction tails are bounded.

NS V25 proves one-use additivity of its assembled tail meet only across
disjoint cutoff bands at one fixed time; the same stock cannot be spent again
at different record times without a new ancestry estimate.  The objects are
not YM objects, but the debit discipline sharpens V31: the total output gap
and support-Casimir floor may be charged only after the whole rooted fusion
cluster has been assembled, not independently at each internal binary join.

\begin{assessment}[Direction after the late drift check]
\label{ass:V30-late-drift}
Neither new companion appendix proves a Yang--Mills cumulant, support-sum,
edge, capacity, or continuum estimate.  The immediate V31 bottleneck remains
large connected-support WUFC.  Its assembly must treat the output denominator
as a one-use rooted-cluster resource.  On a successful WUFC branch, the later
edge stage must use a direct unreduced singular-channel calculation and keep
its reconstruction topology explicit.
\end{assessment}
The audit therefore sends V30 upstream inside the YM temporal row.  The
new observation is that V11 already controls the high-representation
denominator uniformly.  What was missing was the matching elementary
upper bound and an assembled use of the Möbius cancellation.

% =============================================================================
\section{Two-sided Wilson character gaps at every representation height}
\label{sec:V30-character-gap}
% =============================================================================

For a one-link irreducible \(R\), put
\[
 d_{\alpha,R}:=1-\eta_R(\alpha).
\]
For a finite product representation
\(\boldsymbol R=(R_e)\), retain
\[
 C_E(\boldsymbol R)=\sum_eC_2(R_e),
 \qquad
 q_{\alpha,\boldsymbol R}=\prod_e\eta_{R_e}(\alpha),
\]
and define
\begin{equation}
 d_{\alpha,\boldsymbol R}:=1-q_{\alpha,\boldsymbol R}.
 \label{eq:V30-product-gap}
\end{equation}

\begin{lemma}[Uniform upper Dirichlet gap]
\label{lem:V30-upper-gap}
There are \(C_{\rm up}<\infty\) and \(\alpha_{\rm up}\) such that, for
every irreducible \(R\) and every \(\alpha\ge\alpha_{\rm up}\),
\begin{equation}
 \boxed{
 d_{\alpha,R}
 \le C_{\rm up}\min\{s_\alpha C_2(R),1\}.}
 \label{eq:V30-one-link-upper-gap}
\end{equation}
\end{lemma}

\begin{proof}
Let \(U=\exp X\) along a minimizing unit-speed geodesic from the identity,
so \(\norm X=d(U,I)\).  For a unitary representation,
\[
 \norm{D^R(U)-I}_{\rm HS}
 \le\norm{\mathrm dD^R(X)}_{\rm HS}.
\]
Casimir isotropy gives, in the metric convention of V10,
\[
 \frac1{d_R}\norm{\mathrm dD^R(X)}_{\rm HS}^2
 =\frac{C_2(R)}{d_G}\norm X^2.
\]
Together with the trivial bound
\(\norm{D^R(U)-I}_{\rm HS}^2\le4d_R\), the V11 Dirichlet identity yields
\[
 d_{\alpha,R}
 \le C_G\,
 \mathbb E_{\nu_\alpha}
 \min\{C_2(R)d(U,I)^2,1\}.
\]
The V11 quadratic well and normalization bounds give
\[
 \mathbb E_{\nu_\alpha}d(U,I)^2\le C/\alpha.
\]
Since \(x\mapsto\min\{x,1\}\) is increasing and concave,
\[
 d_{\alpha,R}
 \le C_G\min\{C_2(R)/\alpha,1\}.
\]
Finally V22 gives \(s_\alpha\asymp\alpha^{-1}\) on a cofinal interval,
which proves \eqref{eq:V30-one-link-upper-gap}.
\end{proof}

\begin{theorem}[Uniform two-sided product character gap]
\label{thm:V30-two-sided-gap}
There are constants \(0<c_{\rm gap}\le C_{\rm gap}<\infty\) and
\(\alpha_{\rm gap}\) such that, for every finite product representation
\(\boldsymbol R\ne\boldsymbol1\), every finite number of links, and every
\(\alpha\ge\alpha_{\rm gap}\),
\begin{equation}
 \boxed{
 c_{\rm gap}\min\{s_\alpha C_E(\boldsymbol R),1\}
 \le d_{\alpha,\boldsymbol R}
 \le C_{\rm gap}\min\{s_\alpha C_E(\boldsymbol R),1\}.}
 \label{eq:V30-two-sided-product-gap}
\end{equation}
The constants are independent of representation height, support size, and
volume.
\end{theorem}

\begin{proof}
For the upper bound, use
\[
 1-\prod_e\eta_{R_e}
 \le\sum_e(1-\eta_{R_e})
\]
and \Cref{lem:V30-upper-gap}.  The result is at most
\(C_{\rm up}s_\alpha C_E(\boldsymbol R)\), and it is also at most one.

For the lower bound, V11 gives constants \(c_0,a_0>0\) such that
\[
 1-\eta_R(\alpha)
 \ge c_0\left(1-e^{-a_0C_2(R)/\alpha}\right)
 \ge c_1\min\{s_\alpha C_2(R),1\}.
\]
Decrease \(c_1\) to at most one.  Then
\[
 \prod_e\eta_{R_e}
 \le
 \exp\left[
  -c_1\sum_e\min\{s_\alpha C_2(R_e),1\}
 \right].
\]
For nonnegative \(x_e\),
\[
 \sum_e\min\{x_e,1\}
 \ge\min\left\{\sum_ex_e,1\right\}.
\]
Apply this with \(x_e=s_\alpha C_2(R_e)\) and use
\[
 1-e^{-c_1u}\ge(1-e^{-c_1})u
 \qquad(0\le u\le1)
\]
to obtain the lower bound.
\end{proof}

\begin{corollary}[Uniform reduced-resolvent symbol]
\label{cor:V30-resolvent-symbol}
On every nonconstant product channel,
\begin{equation}
 \boxed{
 \frac1{d_{\alpha,\boldsymbol R}}
 \le
 \frac{c_{\rm gap}^{-1}}
      {\min\{s_\alpha C_E(\boldsymbol R),1\}}.}
 \label{eq:V30-resolvent-symbol}
\end{equation}
Thus high representations have an order-one resolvent denominator; only
the low total-Casimir sector carries the \(s_\alpha^{-1}\) pole.
\end{corollary}

\begin{proof}
This is the reciprocal of the lower bound in
\eqref{eq:V30-two-sided-product-gap}.
\end{proof}

% =============================================================================
\section{Fusion geometry and one-defect cancellation}
\label{sec:V30-one-defect}
% =============================================================================

For a product representation \(\boldsymbol R\), put
\begin{equation}
 w_\rho(\boldsymbol R)
 :=
 e^{\rho\operatorname{ht}(\boldsymbol R)}
 \bigl(1+C_E(\boldsymbol R)\bigr).
 \label{eq:V30-Fourier-weight}
\end{equation}
This is the scalar weight in the V29 Fourier algebra norm.

\begin{lemma}[Additive Casimir geometry of an \(\mathrm{SU}(3)\) fusion]
\label{lem:V30-fusion-Casimir}
There is a group constant \(c_{\rm fus}\) with the following property.  If
a product representation \(\boldsymbol T\) occurs in
\(\boldsymbol R\otimes\boldsymbol S\), then
\begin{align}
 \operatorname{ht}(\boldsymbol T)
 &\le
 \operatorname{ht}(\boldsymbol R)
 +\operatorname{ht}(\boldsymbol S),
 \label{eq:V30-height-fusion}\\
 1+C_E(\boldsymbol T)
 &\le
 c_{\rm fus}
 \left(1+C_E(\boldsymbol R)+C_E(\boldsymbol S)\right).
 \label{eq:V30-Casimir-fusion}
\end{align}
Consequently,
\begin{equation}
 w_\rho(\boldsymbol T)
 \le c_{\rm fus}
 w_\rho(\boldsymbol R)w_\rho(\boldsymbol S).
 \label{eq:V30-weight-submultiplicative}
\end{equation}
\end{lemma}

\begin{proof}
On one link, if \(T\subset R\otimes S\), the highest weight of \(T\) is
dominated by the sum of the highest weights of \(R\) and \(S\).  In
\(\mathrm{SU}(3)\) Dynkin labels this gives
\(\operatorname{ht}(T)\le\operatorname{ht}(R)+\operatorname{ht}(S)\).
The explicit Casimir formula \eqref{eq:V6-SU3-Casimir} gives constants
\(c_1,c_2>0\) such that
\[
 c_1(h^2+h)\le1+C_2(p,q)\le c_2(1+h^2+h),
 \qquad h=p+q.
\]
Thus
\[
 1+C_2(T)
 \le c_{\rm fus}\bigl(1+C_2(R)+C_2(S)\bigr).
\]
Sum these inequalities over the links.  A nontrivial output on a link
requires at least one nontrivial input there, so the accumulated constant
term is absorbed by
\(1+C_E(\boldsymbol R)+C_E(\boldsymbol S)\).  This proves
\eqref{eq:V30-height-fusion}--\eqref{eq:V30-Casimir-fusion}.
The product bound follows immediately.
\end{proof}

\begin{lemma}[Uniform one-defect fusion quotient]
\label{lem:V30-one-defect-quotient}
There is \(K_{\rm def}<\infty\), independent of \(\alpha\), the number of
links, and all representation heights, such that whenever
\(\boldsymbol T\ne\boldsymbol1\) occurs in
\(\boldsymbol R\otimes\boldsymbol S\),
\begin{align}
 \frac{w_\rho(\boldsymbol T)}
      {w_\rho(\boldsymbol R)w_\rho(\boldsymbol S)}
 \frac{d_{\alpha,\boldsymbol R}}
      {d_{\alpha,\boldsymbol T}}
 &\le K_{\rm def},
 \label{eq:V30-R-defect-quotient}\\
 \frac{w_\rho(\boldsymbol T)}
      {w_\rho(\boldsymbol R)w_\rho(\boldsymbol S)}
 \frac{d_{\alpha,\boldsymbol S}}
      {d_{\alpha,\boldsymbol T}}
 &\le K_{\rm def}
 \label{eq:V30-S-defect-quotient}
\end{align}
for every \(\alpha\ge\alpha_{\rm gap}\).  A constant input has zero defect
and causes no exception.
\end{lemma}

\begin{proof}
Abbreviate the three total Casimirs by \(A,B,T\), and write
\(s=s_\alpha\).  By \Cref{thm:V30-two-sided-gap},
\[
 \frac{d_{\alpha,\boldsymbol R}}
      {d_{\alpha,\boldsymbol T}}
 \le C
 \frac{\min\{sA,1\}}{\min\{sT,1\}}.
\]
The exponential part of the weight ratio is at most one by
\eqref{eq:V30-height-fusion}.

If \(sT\ge1\), the gap ratio is bounded and
\eqref{eq:V30-weight-submultiplicative} applies.  Suppose \(sT<1\).
If \(sA\le1\), the weighted ratio is at most
\[
 C\,
 \frac{1+T}{(1+A)(1+B)}
 \frac{A}{T},
\]
which is bounded because a nonconstant output has
\(T\ge C_F\).  If \(sA\ge1\), it is at most
\[
 C\,
 \frac{1+T}{T}
 \frac1{s(1+A)(1+B)},
\]
which is bounded because \(s(1+A)\ge1\).  This proves
\eqref{eq:V30-R-defect-quotient}; interchange \(A\) and \(B\) for
\eqref{eq:V30-S-defect-quotient}.
\end{proof}

\begin{theorem}[Uniform one-defect Fourier operator]
\label{thm:V30-one-defect-operator}
On every finite product-link bank, uniformly in its size and in
\(\alpha\ge\alpha_{\rm gap}\),
\begin{equation}
 \boxed{
 \norm{
  \mathcal R_{\alpha,\Lambda}Q_0
  \bigl[(\mathcal C-I)g\bigr]h}_{\rho,2}^{\rm FA}
 \le
 K_{\rm od}
 \norm{g}_{\rho,2}^{\rm FA}
 \norm{h}_{\rho,2}^{\rm FA}.}
 \label{eq:V30-one-defect-operator}
\end{equation}
The same estimate holds with the two factors interchanged.
\end{theorem}

\begin{proof}
Expand \(g,h\) in product Peter--Weyl coefficients.  On an input channel
\(\boldsymbol R\), the factor \(\mathcal C-I\) multiplies by
\(-d_{\alpha,\boldsymbol R}\).  Projecting the product with an
\(\boldsymbol S\) input onto a nonconstant output
\(\boldsymbol T\) and applying \(\mathcal RQ_0\) divides by
\(d_{\alpha,\boldsymbol T}\).  The weighted scalar quotient is bounded by
\Cref{lem:V30-one-defect-quotient}.  Clebsch--Gordan decomposition is
unitary and trace norm is contractive under block compression.  Summing
the Fourier coefficients, exactly as in
\Cref{lem:V29-FA-algebra}, proves
\eqref{eq:V30-one-defect-operator}.
\end{proof}

\begin{corollary}[The marginal binary fusion vertex is uniform]
\label{cor:V30-binary-uniform}
For all finite-bank Fourier algebra inputs,
\begin{equation}
 \boxed{
 \norm{
  \mathcal R_{\alpha,\Lambda}Q_0
  \kappa_2^\alpha(g,h)}_{\rho,2}^{\rm FA}
 \le
 K_{\rm bin}
 \norm{g}_{\rho,2}^{\rm FA}
 \norm{h}_{\rho,2}^{\rm FA},}
 \label{eq:V30-binary-uniform}
\end{equation}
with \(K_{\rm bin}\) independent of representation height, bank size,
volume, and \(\alpha\ge\alpha_{\rm gap}\).
\end{corollary}

\begin{proof}
On eigenmode inputs and an output channel,
\[
 q_{\alpha,\boldsymbol T}
 -q_{\alpha,\boldsymbol R}q_{\alpha,\boldsymbol S}
 =
 d_{\alpha,\boldsymbol R}
 +d_{\alpha,\boldsymbol S}
 -d_{\alpha,\boldsymbol T}
 -d_{\alpha,\boldsymbol R}d_{\alpha,\boldsymbol S}.
\]
After division by \(d_{\alpha,\boldsymbol T}\), the term containing
\(-d_{\alpha,\boldsymbol T}\) is controlled by
\eqref{eq:V30-weight-submultiplicative}; the other terms are controlled
by \Cref{lem:V30-one-defect-quotient}, since
\(0\le d_{\alpha,\boldsymbol R},d_{\alpha,\boldsymbol S}\le1\).
Sum the Fourier blocks.  This is also obtained by writing
\[
 \kappa_2^\alpha(g,h)
 =(\mathcal C-I)(gh)
 -\bigl[(\mathcal C-I)g\bigr]\mathcal Ch
 -g(\mathcal C-I)h
\]
and applying \Cref{thm:V30-one-defect-operator} term by term after the
assembled identity has been formed.
\end{proof}

% =============================================================================
\section{All cumulant orders on one product-link bank}
\label{sec:V30-all-orders-one-bank}
% =============================================================================

\begin{lemma}[One-defect telescoping of the assembled cumulant]
\label{lem:V30-cumulant-telescoping}
For \(m\ge2\), set \(G_B=\prod_{j\in B}g_j\).  Then
\begin{equation}
 \boxed{
 \kappa_m^\alpha(g_1,\ldots,g_m)
 =
 \sum_{\pi\in\Pi_m}\mu(\pi)
 \left[
  \prod_{B\in\pi}\mathcal CG_B-\prod_{j=1}^mg_j
 \right],}
 \label{eq:V30-cumulant-telescoping}
\end{equation}
where
\(\mu(\pi)=(-1)^{|\pi|-1}(|\pi|-1)!\).
After ordering the blocks of each partition in any fixed way, its bracket
is the exact telescoping sum
\begin{equation}
 \sum_{r=1}^{|\pi|}
 \left(\prod_{\ell<r}\mathcal CG_{B_\ell}\right)
 (\mathcal C-I)G_{B_r}
 \left(\prod_{\ell>r}G_{B_\ell}\right).
 \label{eq:V30-block-telescoping}
\end{equation}
\end{lemma}

\begin{proof}
The Möbius formula is \eqref{eq:V29-cumulant-Mobius}.  For \(m\ge2\),
\[
 \sum_{\pi\in\Pi_m}\mu(\pi)=0,
\]
because it is the cumulant of a deterministic scalar repeated \(m\)
times.  Subtracting the partition-independent product
\(\prod_jg_j\) therefore changes nothing and gives
\eqref{eq:V30-cumulant-telescoping}.  The ordinary difference-of-products
identity gives \eqref{eq:V30-block-telescoping}.
\end{proof}

\begin{lemma}[Factorial-exponential partition debit]
\label{lem:V30-partition-debit}
There is \(C_{\rm part}<\infty\) such that
\begin{equation}
 \sum_{\pi\in\Pi_m}|\mu(\pi)|\,|\pi|
 =\sum_{k=1}^m k!\,S(m,k)
 \le m!C_{\rm part}^{m}
 \qquad(m\ge1),
 \label{eq:V30-partition-debit}
\end{equation}
where \(S(m,k)\) is a Stirling number of the second kind.
\end{lemma}

\begin{proof}
There are \(S(m,k)\) partitions with \(k\) blocks, and
\(|\mu(\pi)|\,|\pi|=(k-1)!k=k!\).  The exponential generating function of
the resulting ordered Bell numbers is
\[
 \sum_{m\ge0}
 \left(\sum_{k=0}^m k!S(m,k)\right)\frac{z^m}{m!}
 =\frac1{2-e^z}.
\]
It is analytic on every disk of radius strictly smaller than
\(\log2\).  Cauchy's estimate on one fixed smaller disk proves
\eqref{eq:V30-partition-debit}.
\end{proof}

\begin{theorem}[Uniform assembled cumulant bound on one bank]
\label{thm:V30-one-bank-cumulants}
There is \(K_{\rm cum}<\infty\), depending only on the group and the
declared Fourier weight \(\rho\), such that for every \(m\ge2\),
every finite product-link bank, every \(\alpha\ge\alpha_{\rm gap}\), and
all Fourier algebra inputs,
\begin{equation}
 \boxed{
 \norm{
  \mathcal R_{\alpha,\Lambda}Q_0
  \kappa_m^\alpha(g_1,\ldots,g_m)}_{\rho,2}^{\rm FA}
 \le
 m!K_{\rm cum}^{m-1}
 \prod_{j=1}^m\norm{g_j}_{\rho,2}^{\rm FA}.}
 \label{eq:V30-one-bank-cumulants}
\end{equation}
The constant is independent of cumulant order, representation height,
bank size, volume, and \(\alpha\).  No fixed-representation asymptotic is
used.
\end{theorem}

\begin{proof}
Use the assembled identity
\eqref{eq:V30-cumulant-telescoping} before taking norms.  In every
telescoping term \eqref{eq:V30-block-telescoping}, designate
\((\mathcal C-I)G_{B_r}\) as the defect factor and combine all other
factors into \(H\).  Convolution is a contraction in the Fourier algebra
because all V10 multipliers lie in \([0,1]\).  Repeated use of
\eqref{eq:V29-FA-product} gives
\[
 \norm{G_B}_{\rho,2}^{\rm FA}
 \le c_{\rm alg}^{|B|-1}
 \prod_{j\in B}\norm{g_j}_{\rho,2}^{\rm FA}.
\]
\Cref{thm:V30-one-defect-operator} therefore bounds each telescoping term
by
\[
 K_{\rm od}c_{\rm alg}^{m-1}
 \prod_{j=1}^m\norm{g_j}_{\rho,2}^{\rm FA}.
\]
Sum the exact number of terms using
\Cref{lem:V30-partition-debit}.  Enlarging one group-dependent constant
absorbs \(K_{\rm od}\), \(c_{\rm alg}\), and \(C_{\rm part}\), and yields
\eqref{eq:V30-one-bank-cumulants}.
\end{proof}

\begin{assessment}[What has been closed]
\label{ass:V30-spectral-closure}
\Cref{thm:V30-one-bank-cumulants} closes the simultaneous
cumulant-order and representation-height part of the V29 UFC problem on
every fixed support union.  It is deliberately weaker than the
\(\delta_\alpha^{m-2}\) Gaussian-excess normalization in
\eqref{eq:V29-UFC}; that extra gain is not established here uniformly in
\(m\) and is not needed for a CPM compiler.  The theorem retains one assembled
\((\mathcal C-I)\) defect, which is exactly the single small denominator
needed to cancel the output resolvent.
\end{assessment}

% =============================================================================
\section{Exact support localization and the remaining spatial assembly}
\label{sec:V30-spatial-assembly}
% =============================================================================

For each link \(e\), let \(P_e\) be Haar averaging in that coordinate and
\(Q_e=I-P_e\).  On a finite link bank \(E\), every function has the exact
Hoeffding--Peter--Weyl decomposition
\begin{equation}
 g=\sum_{X\subseteq E}g_X,
 \qquad
 g_X:=
 \left(\prod_{e\in X}Q_e\right)
 \left(\prod_{e\notin X}P_e\right)g.
 \label{eq:V30-exact-support-decomposition}
\end{equation}

\begin{lemma}[Exact active support and Casimir floor]
\label{lem:V30-exact-support}
Every nonzero Peter--Weyl mode of \(g_X\) is nontrivial on each link in
\(X\) and trivial outside \(X\).  Consequently,
\begin{equation}
 C_E(\boldsymbol R)\ge C_F|X|
 \label{eq:V30-support-Casimir-floor}
\end{equation}
on every mode of \(g_X\).  The projections in
\eqref{eq:V30-exact-support-decomposition} commute with
\(\mathcal C\), \(P_0\), \(Q_0\), and the reduced resolvent.
\end{lemma}

\begin{proof}
On one link, \(P_e\) selects the trivial Peter--Weyl representation and
\(Q_e\) deletes it.  Product projections therefore give the first
assertion.  Every nontrivial \(\mathrm{SU}(3)\) representation has Casimir
at least \(C_F\), proving \eqref{eq:V30-support-Casimir-floor}.
All the listed temporal operators are diagonal product-convolution
operators, so they commute with the coordinate Haar projections.
\end{proof}

\begin{lemma}[The binary multiplier defect lives on the overlap]
\label{lem:V30-overlap-factorization}
Let \(\boldsymbol R,\boldsymbol S\) have exact active supports \(X,Y\),
let \(I=X\cap Y\), and let
\(\boldsymbol T\subset\boldsymbol R\otimes\boldsymbol S\).  Put
\[
 q_{\rm out}
 :=
 \prod_{e\in X\triangle Y}
 \eta_{T_e}(\alpha).
\]
Then \(T_e=R_e\) on \(X\setminus Y\),
\(T_e=S_e\) on \(Y\setminus X\), and
\begin{equation}
 \boxed{
 q_{\alpha,\boldsymbol T}
 -q_{\alpha,\boldsymbol R}q_{\alpha,\boldsymbol S}
 =
 q_{\rm out}
 \left[
  \prod_{e\in I}\eta_{T_e}(\alpha)
  -
  \prod_{e\in I}
       \eta_{R_e}(\alpha)\eta_{S_e}(\alpha)
 \right].}
 \label{eq:V30-overlap-factorization}
\end{equation}
In particular the binary cumulant vanishes for \(I=\varnothing\).
\end{lemma}

\begin{proof}
Outside \(I\), one input representation is trivial.  Tensoring with the
trivial representation leaves the other input unchanged, so the output
multiplier on that link is the same factor in
\(q_{\alpha,\boldsymbol T}\) and
\(q_{\alpha,\boldsymbol R}q_{\alpha,\boldsymbol S}\).  Factor their
product.  When \(I=\varnothing\), the bracket is \(1-1\).
\end{proof}

\begin{assessment}[The spatial cancellation which must be retained]
\label{ass:V30-overlap-direction}
Equation \eqref{eq:V30-overlap-factorization} separates the links which
create a covariance from the links which only enlarge the output
denominator.  On an exact support \(X\), the latter contribute the
Casimir floor \eqref{eq:V30-support-Casimir-floor}.  Bounding the local
interaction sum by the number of intersecting supports before using this
denominator would discard the possible compensation between spatial
incidence and temporal mixing.  V31 must assemble those two factors
together.
\end{assessment}

% =============================================================================
\section{A weaker uniform gate is sufficient, and only large support remains}
\label{sec:V30-WUFC}
% =============================================================================

\begin{definition}[Weak uniform fusion-cumulant gate]
\label{def:V30-WUFC}
For \(K,\rho,\mu,\alpha_0>0\), write
\(\operatorname{WUFC}(K,\rho,\mu;\alpha_0)\) if, after the exact support
decomposition \eqref{eq:V30-exact-support-decomposition} and complete
connected-support assembly,
\begin{equation}
 \boxed{
 \norm{
  \mathcal R_{\alpha,\Lambda}Q_0
  \kappa_m^\alpha(g_1,\ldots,g_m)}
 _{\rho,\mu,2}^{\rm FI}
 \le
 m!K^{m-1}
 \prod_{j=1}^m
 \norm{g_j}_{\rho,\mu,2}^{\rm FI}}
 \label{eq:V30-WUFC}
\end{equation}
for every \(m\ge2\), every nondegenerate finite volume, every
\(\alpha\ge\alpha_0\), and all finite local Fourier interactions.
Unlike V29 UFC, WUFC does not request the additional
\(\delta_\alpha^{m-2}\) gain.
\end{definition}

\begin{theorem}[WUFC is enough for CPM]
\label{thm:V30-WUFC-compiler}
Assume \(\operatorname{WUFC}(K,\rho,\mu;\alpha_0)\), and let
\(a=a_{\rho,\mu}\) be the exact leaf norm in
\eqref{eq:V29-leaf-bound}.  Then
\begin{equation}
 \boxed{
 \norm{u_{\alpha,n}}_{\rho,\mu,2}^{\rm FI}
 \le
 n!\,(2a)(16Ka)^{n-1}\delta_\alpha^{-n}}
 \label{eq:V30-WUFC-coefficients}
\end{equation}
for all \(n\ge1\), volumes, and \(\alpha\ge\alpha_0\).  Hence
\begin{equation}
 \boxed{
 \operatorname{CPM}(2c_Ga,16Ka,\mu;\alpha_0)}
 \label{eq:V30-CPM-from-WUFC}
\end{equation}
holds, and all conclusions of
\Cref{thm:V28-CPM-compiler} follow when
\(16Ka\,\zeta_{\alpha,\beta}^{\rm Har}<1\).
\end{theorem}

\begin{proof}
Repeat the coefficient majorization in
\Cref{thm:V29-UFC-compiler} with
\eqref{eq:V30-WUFC} in place of V29 UFC.  For
\(\beta=\delta_\alpha z\), the scalar majorant is now directly
\[
 y\preceq az+\sum_{m\ge2}K^{m-1}y^m
 =az+\frac{Ky^2}{1-Ky}.
\]
The solution and Cauchy estimate
\eqref{eq:V29-scalar-solution} give
\eqref{eq:V30-WUFC-coefficients}.  The one-support \(C^2\) domination in
\Cref{lem:V29-FA-algebra}, summed in the local interaction norm, proves
\eqref{eq:V30-CPM-from-WUFC}.  Invoke
\Cref{thm:V28-CPM-compiler} for the final assertion.
\end{proof}

\begin{proposition}[WUFC on every bounded connected support window]
\label{prop:V30-bounded-support-WUFC}
Fix \(N<\infty\).  There are
\(K_N<\infty\) and \(\alpha_N<\infty\), independent of cumulant order,
representation height, volume, and the position of the support, such that
\eqref{eq:V30-WUFC} holds whenever the complete connected union of the
exact input supports contains at most \(N\) links.
\end{proposition}

\begin{proof}
Up to translation and the finitely many boundary periodizations, there are
only finitely many connected link-support shapes with at most \(N\)
members.  On any such union, \Cref{thm:V30-one-bank-cumulants} controls all
cumulant orders and all representation heights with the same
factorial-exponential constant.  The exact support decomposition has at
most \(2^N\) components per input on that union, and the local root and
diameter weights contribute only constants depending on \(N,\mu\).
Absorb these finitely many factors into \(K_N^{m-1}\).  Translation
invariance and the rooted local norm remove the volume factor.
\end{proof}

For a finite witness
\(w=(m,\alpha,\Lambda;g_1,\ldots,g_m)\) with \(m\ge2\) and every
\(g_j\ne0\), define its factorial-exponential vertex constant by
\begin{equation}
 \mathfrak k(w):=
 \left[
 {\|\mathcal R_{\alpha,\Lambda}Q_0
       \kappa_m^\alpha(g_1,\ldots,g_m)\|_{\rho,\mu,2}^{\rm FI}
  \over
  m!\prod_{j=1}^m\|g_j\|_{\rho,\mu,2}^{\rm FI}}
 \right]^{1/(m-1)}.
 \label{eq:V30-witness-constant}
\end{equation}

\begin{theorem}[Any remaining WUFC failure escapes only in spatial size]
\label{thm:V30-spatial-escape}
Fix \(\rho,\mu\).  If no pair \(K,\alpha_0<\infty\) makes WUFC true
cofinally, then there is a sequence of finite resolved cumulant witnesses
\(w_r=(m_r,\alpha_r,\Lambda_r;g_{r,1},\ldots,g_{r,m_r})\)
such that:
\begin{enumerate}[label=\textup{(\roman*)},leftmargin=3.5em]
\item \(\alpha_r\ge r\) and \(\mathfrak k(w_r)>r\);
\item the connected union of their exact active supports tends to
      infinity; and
\item every witness retains the assembled Möbius sum and a nonconstant
      output channel with the denominator
      \eqref{eq:V30-two-sided-product-gap}.
\end{enumerate}
Representation height or cumulant order alone cannot be the escaping
obstruction.
\end{theorem}

\begin{proof}
For each integer \(r\ge1\), apply the assumed failure of WUFC with
proposed constant \(K=r\) and threshold \(\alpha_0=r\).  A finite
polynomial witness exists by the definition of the FI completion, as in
\Cref{thm:V29-UFC-failure}; choose it at some \(\alpha_r\ge r\).
Failure of the estimate says exactly that \(\mathfrak k(w_r)>r\).

Suppose that the connected support unions did not tend to infinity.
Then some subsequence would be contained, after rooting, in windows of
at most \(N\) links.  By
\Cref{prop:V30-bounded-support-WUFC} there are
\(K_N,\alpha_N<\infty\) for all such witnesses.  Taking on the
subsequence \(r\ge\max\{K_N,\alpha_N\}\) gives simultaneously
\(\alpha_r\ge\alpha_N\), \(r\ge K_N\), and
\(\mathfrak k(w_r)>r\), contradicting
\(\mathfrak k(w_r)\le K_N\).  The witness construction starts from the
exact Möbius formula and the reduced output projection, so item
\textup{(iii)} is preserved.
\end{proof}

\begin{firewall}[Why the one-bank theorem is not yet WUFC]
\label{fire:V30-one-bank-not-local}
\Cref{thm:V30-one-bank-cumulants} controls any fixed tuple of support
components.  The FI norm also sums all connected tuples incident with a
root link.  That sum is not bounded merely by counting intersecting
supports: doing so can lose a factor proportional to the size of an
already assembled cluster.  The linkwise factor
\eqref{eq:V30-overlap-factorization} and the total output denominator must
be used before the rooted support sum.  Until that large-support estimate
is proved, WUFC and CPM remain open.
\end{firewall}

% =============================================================================
\section{V30 ledger and revised frontier}
\label{sec:V30-ledger}
% =============================================================================

\begin{theorem}[Integrated V30 statement]
\label{thm:V30-integrated}
Preserving V1--V29, and after the required companion audit, V30 proves:
\begin{enumerate}[label=\textup{(V30.\arabic*)},leftmargin=5.5em]
\item an exact scheduled audit of SM--Einstein V49 and NS V23, followed by
      a late drift check of SM--Einstein V50 and NS V25, including the
      signed-assembly, one-use, direct-unreduced, and norm-topology cautions
      that match later YM gates;
\item a representation- and volume-uniform two-sided Wilson product gap
      \(1-q_{\alpha,\boldsymbol R}\asymp
      \min\{s_\alpha C_E(\boldsymbol R),1\}\);
\item additive \(\mathrm{SU}(3)\) fusion geometry and a uniform
      one-defect/output-resolvent quotient;
\item a uniform binary fusion bound at every representation height;
\item an all-cumulant, all-representation bound on every product-link
      bank, obtained by Möbius subtraction followed by exact one-defect
      telescoping;
\item the exact link-overlap factorization of the binary multiplier
      defect and the exact active-support Casimir floor;
\item a weaker all-order gate sufficient for CPM with explicit constants;
      and
\item bounded-support WUFC at all cumulant orders and representation
      heights, proving that any remaining failure must escape through
      connected spatial support size.
\end{enumerate}
\end{theorem}

\begin{proof}
Item \textup{(V30.1)} is
\Cref{lem:V30-signed-assembly,thm:V30-audit-decision,
ass:V30-late-drift}.  Item
\textup{(V30.2)} is \Cref{thm:V30-two-sided-gap}.  Items
\textup{(V30.3)--(V30.4)} are
\Cref{lem:V30-fusion-Casimir,lem:V30-one-defect-quotient,
cor:V30-binary-uniform}.  Item \textup{(V30.5)} is
\Cref{thm:V30-one-bank-cumulants}.  Item \textup{(V30.6)} is
\Cref{lem:V30-exact-support,lem:V30-overlap-factorization}.  Item
\textup{(V30.7)} is \Cref{thm:V30-WUFC-compiler}.  Item
\textup{(V30.8)} is
\Cref{prop:V30-bounded-support-WUFC,thm:V30-spatial-escape}.
\end{proof}

\begingroup
\small
\begin{longtable}{>{\raggedright\arraybackslash}p{0.28\textwidth}
                  >{\raggedright\arraybackslash}p{0.16\textwidth}
                  >{\raggedright\arraybackslash}p{0.48\textwidth}}
\toprule
\textbf{V30 row}&\textbf{Status}&\textbf{Advance or remaining obstruction}\\
\midrule
\endfirsthead
\toprule
\textbf{V30 row}&\textbf{Status}&\textbf{Advance or remaining obstruction}\\
\midrule
\endhead
Five-version companion audit
& \MGclosed
& The scheduled SM V49/NS V23 snapshot and late SM V50/NS V25 drift were
  read at their recorded hashes; only type-correct acquisition rules were
  retained.\\[0.35em]
Uniform product character gap
& \MGclosed
& Two-sided estimate \eqref{eq:V30-two-sided-product-gap}, uniform in
  height, support, and volume.\\[0.35em]
Marginal binary fusion vertex
& \MGclosed
& The output resolvent is cancelled uniformly by the assembled multiplier
  defect in \eqref{eq:V30-binary-uniform}.\\[0.35em]
All cumulant orders on one bank
& \MGclosed
& Möbius telescoping gives \eqref{eq:V30-one-bank-cumulants}, uniform in
  cumulant order and representation height.\\[0.35em]
Bounded connected support
& \MGclosed
& WUFC holds at every fixed support-size cutoff, with no order or height
  cutoff.\\[0.35em]
Large connected-support assembly
& \MGopen
& The rooted sum must retain overlap creation and total-output mixing
  before absolute values.\\[0.35em]
Connected Perron majorant
& \MGreduced
& Follows from the weaker WUFC gate by
  \eqref{eq:V30-WUFC-coefficients}.\\[0.35em]
Actual Harnack, density, and incidence packet
& \MGreduced
& These compile from WUFC but remain open until the large-support sum is
  controlled.\\[0.35em]
V24/V25 unreduced edge packet
& \MGopen
& The companion audit forbids promotion of a reduced factor before the
  complete reconstruction and high-channel tail are bounded.\\[0.35em]
Capacity, protected sectors, and OS limit
& \MGopen
& Complete bad capacity, toron/wrapping floors, local-algebra noncollapse,
  and a nontrivial continuum state remain.\\[0.35em]
Full Yang--Mills mass gap
& \MGopen
& The temporal representation and cumulant-order tails are closed; spatial
  linked assembly, finite edges, and continuum closure remain.\\
\bottomrule
\end{longtable}
\endgroup

\begin{openproblem}[V31 mandate: large connected-support assembly]
\label{op:V31}
\begin{enumerate}[label=\textup{(V31.\arabic*)},leftmargin=5.5em]
\item Expand every input into the exact Haar supports
      \eqref{eq:V30-exact-support-decomposition}.  Do not assign a large
      support to a Fourier mode which is trivial on one of its links.
\item For the binary marginal vertex, combine
      \eqref{eq:V30-overlap-factorization},
      \eqref{eq:V30-support-Casimir-floor}, and
      \eqref{eq:V30-two-sided-product-gap} before summing supports incident
      with a root link.
\item Extend the same incidence-versus-mixing cancellation to the
      one-defect telescoping terms
      \eqref{eq:V30-block-telescoping}.  Preserve the full Möbius sum,
      assemble the rooted cluster first, charge its total output denominator
      only once, and use a tree-graph bound only afterwards.
\item Prove WUFC with explicit \(K,\rho,\mu,\alpha_0\), or return the
      growing connected-support witness of
      \Cref{thm:V30-spatial-escape}.  Representation height is no longer an
      admissible unexplained failure.
\item On a strict WUFC branch, invoke
      \Cref{thm:V30-WUFC-compiler}, then populate the local Harnack,
      signed-density, incidence, and complete temporal-source rows.
\item Insert the density and incidence columns separately into the
      unreduced V24/V25 edge packet.  At every singular channel form and
      test the complete unreduced matrix directly, without first using a
      singular free inverse.  Apply \eqref{eq:V30-signed-assembly} to every
      selected/complement split, retain all high-channel reconstruction
      tails, and name the norm in which those tails are controlled.
\item Route a failed local window through the complete V20 predictable
      bad-capacity matrix.  Preserve \(\gamma_{\alpha,j}/\alpha\), toron
      and wrapping modes, protected sectors, local-algebra noncollapse,
      reflection positivity, and the nontrivial OS continuum state.
\item The next companion audit is V35.
\end{enumerate}
\end{openproblem}

\begin{firewall}[End-of-V30 claim boundary]
V30 completes the required companion audit and proves uniform two-sided
Wilson character gaps, uniform binary fusion, and uniform assembled
cumulant bounds at every order and representation height on a fixed
support bank.  It further proves that WUFC holds on every bounded connected
support window and that WUFC is sufficient for CPM with explicit
constants.  It does not sum arbitrarily large connected support tuples in
the rooted local norm.  Therefore WUFC, CPM, the actual nonlinear source
packet, the unreduced V24/V25 edge margins, bad capacity, protected-sector
floors, local-algebra noncollapse, and the nontrivial continuum OS limit
remain open.  The full Yang--Mills mass gap is not proved.
\end{firewall}

\section*{Version V30 append-only record}
\addcontentsline{toc}{section}{Version V30 append-only record}
The complete V29 mathematical source was retained before this continuation.
Its SHA--256 digest was
\begin{center}\scriptsize\ttfamily
3d9d1d3167a274a395e61fba8a85ee06cac3febd56aaac3c08884ae660984cbb.
\end{center}
No mathematical content of V29 was altered.  On the next iteration, retain
this full file, append before its unique active document-closing command,
increment the filename to \texttt{\_V31.tex}, and remove only the
superseded V30 file from this note series.  The next companion audit is
V35.

% =============================================================================
% END APPENDED VERSION V30 MATERIAL
% =============================================================================
% =============================================================================
% BEGIN APPENDED VERSION V31 MATERIAL
% =============================================================================

\clearpage
\part{V31 --- Balanced link marks, global binary spatial fusion, and the
marked-tree gate}
\label{part:V31}

\section{Purpose and acquisition order}
\label{sec:V31-purpose}

V30 proved the resolved cumulant estimate uniformly in representation
height, cumulant order, and the size of any one product-link bank.  It did
not justify the rooted sum over all connected support tuples.  The missing
step is real: if one retains only the assertion that two supports intersect,
then a component on \(N\) links may meet \(N\) singleton components while
both input rooted norms remain one.

V31 goes upstream of that count.  The polynomial Casimir factor already
present in the V29 Fourier norm is rewritten as an exact probability budget
over the active links.  A new Wilson tail-moment estimate then shows that the
binary resolved fusion vertex carries the product of the two probability
marks at a shared link.  This closes the binary FI estimate at arbitrary
support size.  For higher cumulants, V31 states and proves the exact spatial
compiler needed from a marked-tree atlas inequality.  The first unresolved
case is thereby reduced to an explicit ternary two-mark inequality; it is
not hidden inside a generic cluster count.

There is no scheduled companion audit at V31.  The V30 scheduled and late
drift audits remain the frozen checkpoint, and the next required audit is
V35.

% =============================================================================
\section{The balanced Fourier norm and exact link probabilities}
\label{sec:V31-balanced-norm}
% =============================================================================

Let \(X\ne\varnothing\) be an exact active support in the sense of
\Cref{lem:V30-exact-support}.  For a product representation
\(\boldsymbol R=(R_e)_{e\in X}\), every \(R_e\) is nontrivial.  Put
\begin{align}
 a_e(\boldsymbol R)&:=1+C_2(R_e),
 \label{eq:V31-link-mass}\\
 \Xi(\boldsymbol R)&:=\sum_{e\in X}a_e(\boldsymbol R)
 =|X|+C_E(\boldsymbol R),
 \label{eq:V31-balanced-mass}\\
 \vartheta_e(\boldsymbol R)&:=
 \frac{a_e(\boldsymbol R)}{\Xi(\boldsymbol R)}
 \qquad(e\in X).
 \label{eq:V31-link-probability}
\end{align}
Thus
\begin{equation}
 \sum_{e\in X}\vartheta_e(\boldsymbol R)=1.
 \label{eq:V31-probability-normalization}
\end{equation}

Define the balanced one-support norm
\begin{equation}
 \norm{g_X}_{\rho,2}^{\rm BFA}
 :=
 \sum_{\boldsymbol R\ {\mathrm{exact\ on}}\ X}
 e^{\rho\operatorname{ht}(\boldsymbol R)}
 \Xi(\boldsymbol R)d_{\boldsymbol R}
 \norm{\widehat g_X(\boldsymbol R)}_1.
 \label{eq:V31-BFA}
\end{equation}
For an ordered \(d\)-tuple
\(\boldsymbol e=(e_1,\ldots,e_d)\in X^d\), repetitions allowed, define
the marked seminorm
\begin{equation}
 \norm{g_X}_{\rho,2;\boldsymbol e}^{\rm BFA(d)}
 :=
 \sum_{\boldsymbol R\ {\mathrm{exact\ on}}\ X}
 e^{\rho\operatorname{ht}(\boldsymbol R)}
 \Xi(\boldsymbol R)d_{\boldsymbol R}
 \norm{\widehat g_X(\boldsymbol R)}_1
 \prod_{r=1}^d\vartheta_{e_r}(\boldsymbol R).
 \label{eq:V31-marked-BFA}
\end{equation}
The empty product gives the BFA norm when \(d=0\).  A mark outside \(X\)
is assigned seminorm zero.  Finally set
\begin{equation}
 \norm{g}_{\rho,\mu,2}^{\rm BFI}
 :=
 \sup_{e_0}
 \sum_{X\ni e_0}
 e^{\mu\operatorname{diam}X}
 \norm{g_X}_{\rho,2}^{\rm BFA}.
 \label{eq:V31-BFI}
\end{equation}

\begin{lemma}[Balanced-norm equivalence and exact mark normalization]
\label{lem:V31-balanced-equivalence}
There is a group constant \(c_\Xi<\infty\) such that, on every nonempty
exact support,
\begin{equation}
 \norm{g_X}_{\rho,2}^{\rm FA}
 \le \norm{g_X}_{\rho,2}^{\rm BFA}
 \le c_\Xi\norm{g_X}_{\rho,2}^{\rm FA}.
 \label{eq:V31-FA-BFA-equivalence}
\end{equation}
The same equivalence holds between FI and BFI.  Moreover, for every
\(d\ge1\),
\begin{equation}
 \boxed{
 \sum_{(e_1,\ldots,e_d)\in X^d}
 \norm{g_X}_{\rho,2;(e_1,\ldots,e_d)}^{\rm BFA(d)}
 =\norm{g_X}_{\rho,2}^{\rm BFA}.}
 \label{eq:V31-mark-normalization}
\end{equation}
In particular every marked seminorm is at most the BFA norm.
\end{lemma}

\begin{proof}
By the support Casimir floor
\eqref{eq:V30-support-Casimir-floor},
\[
 |X|\le C_F^{-1}C_E(\boldsymbol R).
\]
Hence
\[
 1+C_E(\boldsymbol R)
 \le\Xi(\boldsymbol R)
 \le(1+C_F^{-1})(1+C_E(\boldsymbol R)),
\]
which proves \eqref{eq:V31-FA-BFA-equivalence}; taking the same rooted
sum proves the FI--BFI equivalence.  For each fixed representation, the
sum of the product of marks is
\[
 \sum_{(e_1,\ldots,e_d)\in X^d}
 \prod_{r=1}^d\vartheta_{e_r}(\boldsymbol R)
 =\left(\sum_{e\in X}\vartheta_e(\boldsymbol R)\right)^d=1.
\]
Tonelli's theorem for the nonnegative Fourier trace-norm summands proves
\eqref{eq:V31-mark-normalization}.
\end{proof}

\begin{remark}[Constant support components]
\label{rem:V31-constants}
If one input of a cumulant of order at least two is constant, the cumulant
vanishes.  Constants may therefore be removed before
\eqref{eq:V31-balanced-mass} is used.  The final \(Q_0\) likewise removes
the constant output.  No division by an empty-support mass occurs.
\end{remark}

% =============================================================================
\section{A Wilson multiplier tail moment}
\label{sec:V31-tail-moment}
% =============================================================================

\begin{lemma}[Laplacian moment of the Wilson density]
\label{lem:V31-density-Laplacian}
Let \(\Delta_G\) be the bi-invariant group Laplacian in the Casimir
normalization of V10.  There are \(C_\Delta<\infty\) and
\(\alpha_\Delta<\infty\) such that
\begin{equation}
 \norm{\Delta_Gq_\alpha}_{L^1(G)}\le C_\Delta\alpha
 \qquad(\alpha\ge\alpha_\Delta).
 \label{eq:V31-density-Laplacian}
\end{equation}
Consequently every irreducible multiplier satisfies
\begin{equation}
 C_2(R)\eta_R(\alpha)\le C_\Delta\alpha.
 \label{eq:V31-character-first-moment}
\end{equation}
\end{lemma}

\begin{proof}
Write the V10 one-link density as
\[
 q_\alpha=\widetilde Z_\alpha^{-1}e^{-\alpha W},
\]
where \(W\ge0\) is the smooth Wilson well with its unique nondegenerate
minimum at the identity.  The V11 global quadratic comparison implies
\(lvert\nabla W\rvert^2\le C W\) on the compact group.  The same Gaussian
comparison and \eqref{eq:V11-Z-Laplace} give
\begin{equation}
 \int_GWq_\alpha\,\dd U\le C\alpha^{-1}.
 \label{eq:V31-Wilson-well-moment}
\end{equation}
Indeed the numerator is bounded in normal coordinates by a constant times
\(\alpha^{-5}\), while the normalization is bounded below by a constant
times \(\alpha^{-4}\); the complement is exponentially small.

Differentiation gives
\[
 \Delta_Gq_\alpha
 =q_\alpha\left(-\alpha\Delta_GW
                 +\alpha^2\lvert\nabla W\rvert^2\right).
\]
Since \(\Delta_GW\) is bounded, \eqref{eq:V31-Wilson-well-moment}
proves \eqref{eq:V31-density-Laplacian}.

Centrality gives
\(\widehat q_\alpha(R)=\eta_R(\alpha)I_{d_R}\).  Integrate by parts in
each matrix coefficient and use
\(-\Delta_GD^R=C_2(R)D^R\).  Since \(D^R(U)\) is unitary,
\[
 C_2(R)\eta_R(\alpha)
 \le\int_G\lvert\Delta_Gq_\alpha(U)\rvert\,\dd U.
\]
This proves \eqref{eq:V31-character-first-moment}.
\end{proof}

\begin{theorem}[One-link tail debit and product damping]
\label{thm:V31-tail-debit}
There are \(K_{\rm tail}<\infty\) and \(\alpha_{\rm tail}<\infty\) such
that, for every irreducible \(R\) and every
\(\alpha\ge\alpha_{\rm tail}\),
\begin{equation}
 \boxed{
 s_\alpha C_2(R)\eta_R(\alpha)
 \le K_{\rm tail}\bigl(1-\eta_R(\alpha)\bigr).}
 \label{eq:V31-one-link-tail-debit}
\end{equation}
For every finite nontrivial product channel exact on (E),
\(\boldsymbol R=(R_e)_{e\in E}\),
\begin{align}
 s_\alpha C_E(\boldsymbol R)q_{\alpha,\boldsymbol R}
 &\le K_{\rm tail},
 \label{eq:V31-product-first-moment}\\
 s_\alpha\Xi(\boldsymbol R)q_{\alpha,\boldsymbol R}
 &\le K_{\rm tail}(1+C_F^{-1}).
 \label{eq:V31-balanced-product-damping}
\end{align}
The constants are independent of representation height and the number of
links.
\end{theorem}

\begin{proof}
Put \(x=s_\alpha C_2(R)\).  On \(x\le1\), the V11 lower character gap,
combined with \(s_\alpha\asymp\alpha^{-1}\), gives
\(1-\eta_R\ge c x\), while \(x\eta_R\le x\).  On \(x\ge1\), the same
lower gap gives (1-\eta_R\ge c>0), whereas
\Cref{lem:V31-density-Laplacian} and
\(s_\alpha\alpha\le C\) give \(x\eta_R\le C\).  Enlarging one constant
proves \eqref{eq:V31-one-link-tail-debit}.

Write \(x_e=s_\alpha C_2(R_e)\) and \(r_e=\eta_{R_e}\).  Then
\begin{align*}
 \left(\sum_ex_e\right)\prod_er_e
 &=\sum_e x_er_e\prod_{f\ne e}r_f\\
 &\le K_{\rm tail}\sum_e(1-r_e)\prod_{f\ne e}r_f
 \le K_{\rm tail}.
\end{align*}
The last sum is the probability of exactly one failure for independent
Bernoulli variables with success probabilities \(r_e\), hence is at most
one.  This proves \eqref{eq:V31-product-first-moment}.  Exact activity and
the fundamental Casimir floor give
\(|E|\le C_F^{-1}C_E(\boldsymbol R)\), proving
\eqref{eq:V31-balanced-product-damping}.
\end{proof}

% =============================================================================
\section{The marked binary Wilson quotient}
\label{sec:V31-binary-quotient}
% =============================================================================

Let \(\boldsymbol R,\boldsymbol S\) have nonempty exact supports \(X,Y\),
let \(I=X\cap Y\), and let a nonconstant exact output channel
\(\boldsymbol T\subset\boldsymbol R\otimes\boldsymbol S\) have support
\(Z\subseteq X\cup Y\).  Abbreviate
\begin{align}
 A_e&:=1+C_2(R_e),
 &B_e&:=1+C_2(S_e),
 \label{eq:V31-local-balanced-inputs}\\
 H_I(\boldsymbol R,\boldsymbol S)
 &:=\sum_{e\in I}A_eB_e.
 \label{eq:V31-overlap-mass}
\end{align}
Both inputs are nontrivial on every link in \(I\), so all terms in
\eqref{eq:V31-overlap-mass} are strictly positive.

\begin{lemma}[Local one-link fusion defect]
\label{lem:V31-local-fusion-defect}
If \(T\subset R\otimes S\) and \(R,S\) are nontrivial, then, cofinally in
\(\alpha\),
\begin{equation}
 \left|\eta_T-\eta_R\eta_S\right|
 \le K_{\rm lf}\min\left\{
 s_\alpha(1+C_2(R))(1+C_2(S)),1\right\}.
 \label{eq:V31-local-fusion-defect}
\end{equation}
The same formula is zero when one input is trivial and the output is the
other input.
\end{lemma}

\begin{proof}
All Wilson multipliers lie in \([0,1]\).  Therefore
\[
 |\eta_T-\eta_R\eta_S|
 \le(1-\eta_T)+(1-\eta_R)+(1-\eta_S).
\]
The upper one-link gap \eqref{eq:V30-one-link-upper-gap} and the fusion
Casimir estimate \eqref{eq:V30-Casimir-fusion} bound the right side by
\[
 Cs_\alpha\bigl(1+C_2(R)+C_2(S)\bigr)
 \le Cs_\alpha(1+C_2(R))(1+C_2(S)).
\]
The left side is also at most one.  If one input is trivial, tensoring
with it does not change the other representation or its multiplier.
\end{proof}

\begin{theorem}[Marked binary scalar quotient]
\label{thm:V31-marked-binary-quotient}
There are \(K_{\rm mb}<\infty\) and \(\alpha_{\rm mb}<\infty\), independent
of \(X,Y,Z\) and all representation heights, such that
\begin{equation}
 \boxed{
 \frac{\Xi(\boldsymbol T)}
      {\Xi(\boldsymbol R)\Xi(\boldsymbol S)}
 \frac{\left|q_{\alpha,\boldsymbol T}
       -q_{\alpha,\boldsymbol R}q_{\alpha,\boldsymbol S}\right|}
      {1-q_{\alpha,\boldsymbol T}}
 \le K_{\rm mb}
 \sum_{e\in I}
 \vartheta_e(\boldsymbol R)\vartheta_e(\boldsymbol S).}
 \label{eq:V31-marked-binary-quotient}
\end{equation}
When \(I=\varnothing\), both sides vanish.
\end{theorem}

\begin{proof}
The disjoint case is \Cref{lem:V30-overlap-factorization}.  Suppose
\(I\ne\varnothing\).  On the symmetric difference
\(O=X\mathbin\triangle Y\), the output representation is the unique
nontrivial input representation.  Put
\[
 q_O:=\prod_{e\in O}\eta_{T_e}.
\]
The exact overlap factorization \eqref{eq:V30-overlap-factorization} gives
\begin{equation}
 \left|q_{\alpha,\boldsymbol T}
       -q_{\alpha,\boldsymbol R}q_{\alpha,\boldsymbol S}\right|
 =q_O B_I,
 \label{eq:V31-binary-factor-B}
\end{equation}
where
\[
 B_I:=
 \left|\prod_{e\in I}\eta_{T_e}
 -\prod_{e\in I}\eta_{R_e}\eta_{S_e}\right|.
\]
Telescope the two products link by link and apply
\Cref{lem:V31-local-fusion-defect}.  Since a difference of two numbers in
\([0,1]\) is at most one,
\begin{equation}
 B_I\le K\min\{s_\alpha H_I,1\}.
 \label{eq:V31-overlap-defect-upper}
\end{equation}

Split the balanced output mass as
\begin{equation}
 \Xi(\boldsymbol T)=\Xi_O+\Xi_I^T,
 \label{eq:V31-output-mass-split}
\end{equation}
where the second term includes only the nontrivial output representations
on \(I\).  Linkwise fusion geometry gives
\[
 1+C_2(T_e)
 \le c\bigl(1+C_2(R_e)+C_2(S_e)\bigr)
 \le cA_eB_e
\]
on every active overlap output.  Hence
\begin{equation}
 \Xi_I^T\le cH_I.
 \label{eq:V31-inside-mass-paid}
\end{equation}
On \(O\), exact activity and
\eqref{eq:V31-balanced-product-damping} give
\begin{equation}
 s_\alpha\Xi_Oq_O\le K.
 \label{eq:V31-outside-mass-damping}
\end{equation}

Let \(C_T=C_E(\boldsymbol T)\).  Since the output is nonconstant and
exact on \(Z\),
\begin{equation}
 \Xi(\boldsymbol T)\le(1+C_F^{-1})C_T.
 \label{eq:V31-output-Xi-Casimir}
\end{equation}
The lower product gap \eqref{eq:V30-two-sided-product-gap} is
\[
 1-q_{\alpha,\boldsymbol T}
 \ge c\min\{s_\alpha C_T,1\}.
\]
If \(s_\alpha C_T\le1\), then
\begin{align*}
 \Xi(\boldsymbol T)
 \frac{q_OB_I}{1-q_{\alpha,\boldsymbol T}}
 &\le \frac C{s_\alpha}\min\{Ks_\alpha H_I,1\}
 \le CH_I.
\end{align*}
If \(s_\alpha C_T\ge1\), use
\eqref{eq:V31-output-mass-split}--%
\eqref{eq:V31-outside-mass-damping}:
\begin{align*}
 \Xi(\boldsymbol T)
 \frac{q_OB_I}{1-q_{\alpha,\boldsymbol T}}
 &\le C\bigl(\Xi_I^Tq_OB_I+\Xi_Oq_OB_I\bigr)\\
 &\le C\left(H_I+\frac1{s_\alpha}
                    \min\{s_\alpha H_I,1\}\right)
 \le CH_I.
\end{align*}
Thus in both cases
\begin{equation}
 \Xi(\boldsymbol T)
 \frac{\left|q_{\alpha,\boldsymbol T}
       -q_{\alpha,\boldsymbol R}q_{\alpha,\boldsymbol S}\right|}
      {1-q_{\alpha,\boldsymbol T}}
 \le KH_I.
 \label{eq:V31-unscaled-marked-quotient}
\end{equation}
Divide by \(\Xi(\boldsymbol R)\Xi(\boldsymbol S)\) and use
\[
 \frac{H_I}{\Xi(\boldsymbol R)\Xi(\boldsymbol S)}
 =\sum_{e\in I}
   \vartheta_e(\boldsymbol R)\vartheta_e(\boldsymbol S)
\]
to prove \eqref{eq:V31-marked-binary-quotient}.
\end{proof}

\begin{theorem}[Marked binary Fourier operator]
\label{thm:V31-marked-binary-operator}
For all nonempty exact supports \(X,Y\), all finite-bank inputs supported
exactly there, and all sufficiently large \(\alpha\),
\begin{equation}
 \boxed{
 \sum_{\varnothing\ne Z\subseteq X\cup Y}
 \norm{\bigl[\mathcal R_{\alpha,\Lambda}Q_0
       \kappa_2^\alpha(g_X,h_Y)\bigr]_Z}_{\rho,2}^{\rm BFA}
 \le K_{\rm mbo}
 \sum_{e\in X\cap Y}
 \norm{g_X}_{\rho,2;e}^{\rm BFA(1)}
 \norm{h_Y}_{\rho,2;e}^{\rm BFA(1)}.}
 \label{eq:V31-marked-binary-operator}
\end{equation}
The subscript \(Z\) denotes the exact-support projection, so the left side
includes every nonconstant output channel exactly once.
\end{theorem}

\begin{proof}
Expand both inputs in product Peter--Weyl coefficients and each product in
Clebsch--Gordan output blocks.  On a resolved triple
\((\boldsymbol R,\boldsymbol S;\boldsymbol T)\), the scalar multiplier is
the quotient in \eqref{eq:V29-binary-fusion-quotient}.  The exponential
height ratio is at most one by \eqref{eq:V30-height-fusion}, while its
balanced polynomial ratio is exactly the prefactor on the left of
\eqref{eq:V31-marked-binary-quotient}.  Apply that theorem.

For a fixed shared link \(e\), the remaining scalar factors are
\(\vartheta_e(\boldsymbol R)\vartheta_e(\boldsymbol S)\).
Clebsch--Gordan decomposition is unitary, and trace norm is contractive
under the resulting block compressions, exactly as in
\Cref{lem:V29-FA-algebra}.  Summing input blocks factors into the two
marked seminorms at \(e\).  Finally sum \(e\in X\cap Y\) and all exact
nonconstant outputs.  This proves
\eqref{eq:V31-marked-binary-operator}.
\end{proof}

% =============================================================================
\section{Global binary closure in the rooted interaction norm}
\label{sec:V31-binary-FI}
% =============================================================================

\begin{theorem}[Binary WUFC at arbitrary connected support size]
\label{thm:V31-binary-WUFC}
There are \(K_2<\infty\) and \(\alpha_2<\infty\), independent of volume,
support size, and representation height, such that
\begin{align}
 \norm{\mathcal R_{\alpha,\Lambda}Q_0
       \kappa_2^\alpha(g,h)}_{\rho,\mu,2}^{\rm BFI}
 &\le K_2
 \norm g_{\rho,\mu,2}^{\rm BFI}
 \norm h_{\rho,\mu,2}^{\rm BFI},
 \label{eq:V31-binary-BFI}\\
 \norm{\mathcal R_{\alpha,\Lambda}Q_0
       \kappa_2^\alpha(g,h)}_{\rho,\mu,2}^{\rm FI}
 &\le K_2'
 \norm g_{\rho,\mu,2}^{\rm FI}
 \norm h_{\rho,\mu,2}^{\rm FI}
 \label{eq:V31-binary-FI}
\end{align}
for all finite local Fourier interactions and
\(\alpha\ge\alpha_2\).  Thus \eqref{eq:V30-WUFC} is proved globally for
\(m=2\), with no support cutoff.
\end{theorem}

\begin{proof}
Decompose both inputs into exact supports.  A binary cumulant vanishes
unless \(X\cap Y\ne\varnothing\).  Every exact output support
\(Z\) is contained in \(X\cup Y\).  Because \(X\cap Y\ne\varnothing\),
\begin{equation}
 \operatorname{diam}Z
 \le\operatorname{diam}(X\cup Y)
 \le\operatorname{diam}X+\operatorname{diam}Y.
 \label{eq:V31-diameter-product}
\end{equation}

Fix an output root link \(e_0\).  If \(e_0\in Z\), then \(e_0\in X\) or
\(e_0\in Y\).  Overcount by those two cases.  In the first case,
\Cref{thm:V31-marked-binary-operator} and
\eqref{eq:V31-diameter-product} give a rooted sum bounded by
\begin{align*}
 K_{\rm mbo}
 \sum_{X\ni e_0}e^{\mu\operatorname{diam}X}
 \sum_{e\in X}
 \norm{g_X}_{\rho,2;e}^{\rm BFA(1)}
 \sum_{Y\ni e}e^{\mu\operatorname{diam}Y}
 \norm{h_Y}_{\rho,2;e}^{\rm BFA(1)}.
\end{align*}
The innermost sum is at most
\(\norm h_{\rho,\mu,2}^{\rm BFI}\), because a marked seminorm is at
most its unmarked BFA norm.  The exact normalization
\eqref{eq:V31-mark-normalization}, with \(d=1\), turns the remaining sum
over \(e\in X\) into \(\norm{g_X}_{\rho,2}^{\rm BFA}\).  The result is at
most
\[
 K_{\rm mbo}
 \norm g_{\rho,\mu,2}^{\rm BFI}
 \norm h_{\rho,\mu,2}^{\rm BFI}.
\]
The case \(e_0\in Y\) is symmetric.  Taking the supremum over \(e_0\)
proves \eqref{eq:V31-binary-BFI}.  Apply the norm equivalence
\eqref{eq:V31-FA-BFA-equivalence} to obtain
\eqref{eq:V31-binary-FI}.
\end{proof}

\begin{assessment}[What the binary theorem removes]
\label{ass:V31-binary-progress}
The support-incidence factor is not estimated by \(|X\cap Y|\).  Each
shared link is charged by the product of its two normalized Casimir marks,
and those marks sum to one separately on both inputs.  Thus large support,
large representation on one exceptional link, and high output Casimir are
handled by the same estimate.  The V30 global binary row is now closed in
the actual FI norm, not merely on one fixed bank.
\end{assessment}

% =============================================================================
\section{Why the unmarked one-bank estimate cannot be summed}
\label{sec:V31-unmarked-obstruction}
% =============================================================================

\begin{proposition}[Sharp abstract incidence loss without marks]
\label{prop:V31-unmarked-loss}
Fix \(\mu\ge0\).  For every \(N\), there are nonnegative finite support
activities \((a_X),(b_Y)\) on a link bank and a root \(e_0\) such that
\begin{equation}
 \sup_e\sum_{X\ni e}e^{\mu\operatorname{diam}X}a_X=1,
 \qquad
 \sup_e\sum_{Y\ni e}e^{\mu\operatorname{diam}Y}b_Y=1,
 \label{eq:V31-unit-rooted-activities}
\end{equation}
but
\begin{equation}
 \sum_{\substack{X\ni e_0,\;Y:\ X\cap Y\ne\varnothing}}
 e^{\mu\operatorname{diam}(X\cup Y)}a_Xb_Y=N.
 \label{eq:V31-unmarked-incidence-loss}
\end{equation}
Consequently, the V30 one-bank constant and support connectedness alone
cannot imply an FI operator bound at any diameter weight.
\end{proposition}

\begin{proof}
Take one \(N\)-link set \(X_0\ni e_0\), put
\(a_{X_0}=e^{-\mu\operatorname{diam}X_0}\), and set all other
\(a_X=0\).  For every \(f\in X_0\), let \(Y_f=\{f\}\), put
\(b_{Y_f}=1\), and set all other \(b_Y=0\).  Both rooted suprema in
\eqref{eq:V31-unit-rooted-activities} equal one.  Every \(Y_f\) intersects
\(X_0\), and \(X_0\cup Y_f=X_0\), so each of the \(N\) summands in
\eqref{eq:V31-unmarked-incidence-loss} equals one.
\end{proof}

\begin{remark}[Why this is not a Wilson counterexample]
\label{rem:V31-not-Wilson-counterexample}
The activities in \Cref{prop:V31-unmarked-loss} deliberately discard the
resolved multiplier.  For a product representation of comparable weight
on the \(N\) links, the central component assigns mark of order \(N^{-1}\)
to each link, and the sum of the \(N\) marked incidences is order one.
\Cref{thm:V31-binary-WUFC} proves that the exact Wilson vertex has this
compensation.  The proposition excludes only the unmarked proof strategy.
\end{remark}

% =============================================================================
\section{The all-order marked-tree atlas gate}
\label{sec:V31-MTA}
% =============================================================================

Let \(\mathfrak T_m\) be the set of labeled trees on
\(\{1,\ldots,m\}\).  Given exact input supports
\(\boldsymbol X=(X_1,\ldots,X_m)\) and
\(\tau\in\mathfrak T_m\), a link marking of \(\tau\) is a map
\begin{equation}
 \ell:E(\tau)\longrightarrow E(\Lambda),
 \qquad
 \ell(ij)\in X_i\cap X_j.
 \label{eq:V31-tree-link-marking}
\end{equation}
If an edge intersection is empty, there is no such marking.  At vertex
\(i\), let \(\boldsymbol e_i^\ell\) be the ordered list of marks on its
incident tree edges.  Its length is \(\deg_\tau(i)\); the value of the
seminorm is independent of the chosen ordering.

\begin{definition}[Marked-tree atlas inequality]
\label{def:V31-MTA}
For \(K_0,\rho,\alpha_0>0\), write
\(\operatorname{MTA}(K_0,\rho;\alpha_0)\) if, for every \(m\ge2\), every
finite product-link bank, every \(\alpha\ge\alpha_0\), every tuple of
nonempty exact supports, and all inputs on those supports,
\begin{align}
 &\sum_{\varnothing\ne Z\subseteq\cup_iX_i}
 \norm{\bigl[\mathcal R_{\alpha,\Lambda}Q_0
       \kappa_m^\alpha(g_{1,X_1},\ldots,g_{m,X_m})\bigr]_Z}
       _{\rho,2}^{\rm BFA}
 \notag\\
 &\qquad\le
 K_0^{m-1}
 \sum_{\tau\in\mathfrak T_m}
 \sum_{\ell\ {\mathrm{marking\ of}}\ \tau}
 \prod_{i=1}^m
 \norm{g_{i,X_i}}_{
       \rho,2;\boldsymbol e_i^\ell}^{
       \rm BFA(\deg_\tau(i))}.
 \label{eq:V31-MTA}
\end{align}
The right side is zero when the support intersection graph is disconnected.
\end{definition}

\begin{proposition}[The binary marked-tree atlas is proved]
\label{prop:V31-MTA-two}
The \(m=2\) instance of \eqref{eq:V31-MTA} holds with
\(K_0=K_{\rm mbo}\) and no support or representation cutoff.
\end{proposition}

\begin{proof}
There is one labeled tree on two vertices.  Its markings are precisely the
links in \(X_1\cap X_2\), and both vertex degrees equal one.  Thus
\eqref{eq:V31-MTA} is exactly
\eqref{eq:V31-marked-binary-operator}.
\end{proof}

\begin{theorem}[Marked-tree spatial compiler]
\label{thm:V31-MTA-compiler}
If \(\operatorname{MTA}(K_0,\rho;\alpha_0)\) holds, then for every
\(\mu\ge0\) there is \(K<\infty\), depending only on
\(K_0,\rho,\mu\) and the group constants, such that
\begin{equation}
 \operatorname{WUFC}(K,\rho,\mu;\alpha_0)
 \label{eq:V31-MTA-implies-WUFC}
\end{equation}
holds.  No support-size, cumulant-order, representation-height, or volume
constant is introduced by the spatial assembly.
\end{theorem}

\begin{proof}
Expand every input into exact supports.  Fix \(m\), a labeled tree
\(\tau\), a nonconstant exact output support \(Z\), and an output root
\(e_0\in Z\).  Every output support is contained in
\(\bigcup_iX_i\), so overcount by the choice of one
\(k\in\{1,\ldots,m\}\) with \(e_0\in X_k\).  Root \(\tau\) at \(k\).

Every tree marking makes adjacent supports share a link.  Therefore their
union is connected and
\begin{equation}
 e^{\mu\operatorname{diam}Z}
 \le e^{\mu\operatorname{diam}(\cup_iX_i)}
 \le\prod_{i=1}^m e^{\mu\operatorname{diam}X_i}.
 \label{eq:V31-tree-diameter-payment}
\end{equation}
Apply \eqref{eq:V31-MTA} and include the rightmost product of diameter
weights.

Now peel the rooted tree from its leaves.  Suppose a leaf \(j\) is joined
to its parent by a link mark \(e\).  With that parent mark fixed,
\begin{equation}
 \sum_{X_j\ni e}e^{\mu\operatorname{diam}X_j}
 \norm{g_{j,X_j}}_{\rho,2;e}^{\rm BFA(1)}
 \le\norm{g_j}_{\rho,\mu,2}^{\rm BFI}.
 \label{eq:V31-leaf-peeling}
\end{equation}
After all descendants of an internal vertex have been peeled, sum their
incident marks.  The normalization
\eqref{eq:V31-mark-normalization} leaves the one seminorm marked by the
parent edge; its support sum is again bounded by the corresponding BFI
norm.  At the root vertex, sum every remaining incident mark using
\eqref{eq:V31-mark-normalization}, and then sum
\(X_k\ni e_0\).  Hence each fixed rooted labeled tree contributes at most
\begin{equation}
 \prod_{i=1}^m\norm{g_i}_{\rho,\mu,2}^{\rm BFI}.
 \label{eq:V31-fixed-tree-sum}
\end{equation}

Cayley's formula gives \(|\mathfrak T_m|=m^{m-2}\).  There are at most
\(m\) choices of the input containing the output root, so the total factor
is \(m^{m-1}\).  The elementary integral bound
\[
 \log(m!)=\sum_{j=1}^m\log j
 \ge\int_1^m\log x\,\dd x
 =m\log m-m+1
\]
implies
\begin{equation}
 m^{m-1}\le e^{m-1}m!.
 \label{eq:V31-tree-factorial}
\end{equation}
Combining \eqref{eq:V31-MTA},
\eqref{eq:V31-fixed-tree-sum}, and
\eqref{eq:V31-tree-factorial} proves the BFI version of WUFC with a
factorial-exponential constant.  The fixed equivalence constants in
\eqref{eq:V31-FA-BFA-equivalence} are absorbed into a larger
\(K^{m-1}\), proving \eqref{eq:V31-MTA-implies-WUFC} in the original FI
norm.
\end{proof}

\begin{corollary}[Complete conditional route from MTA]
\label{cor:V31-MTA-route}
Under MTA, \Cref{thm:V31-MTA-compiler} and
\Cref{thm:V30-WUFC-compiler} give CPM with the explicit coefficient bound
\eqref{eq:V30-WUFC-coefficients}.  The Harnack, density, incidence, and
temporal-source conclusions of \Cref{thm:V28-CPM-compiler} then follow on
the stated small-susceptibility domain.
\end{corollary}

\begin{proof}
This is the successive application of the two named compilers.  No
additional spatial sum remains after
\Cref{thm:V31-MTA-compiler}.
\end{proof}

% =============================================================================
\section{The first unresolved atlas inequality}
\label{sec:V31-ternary}
% =============================================================================

For \(m=3\), the MTA right side is the explicit sum
\begin{align}
 K_0^2\sum_{i=1}^3
 \sum_{\substack{e\in X_i\cap X_j\\
                  f\in X_i\cap X_k}}
 &\norm{g_{i,X_i}}_{\rho,2;(e,f)}^{\rm BFA(2)}
 \norm{g_{j,X_j}}_{\rho,2;e}^{\rm BFA(1)}
 \norm{g_{k,X_k}}_{\rho,2;f}^{\rm BFA(1)},
 \label{eq:V31-ternary-MTA}
\end{align}
where, for each center \(i\), either ordering \((j,k)\) of the other two
indices is fixed; symmetry makes the choice immaterial.  The assembled
ternary numerator is
\begin{align}
 \kappa_3^\alpha(g_1,g_2,g_3)
 &=\mathcal C(g_1g_2g_3)
 -\mathcal C(g_1g_2)\mathcal Cg_3
 -\mathcal C(g_1g_3)\mathcal Cg_2
 \notag\\
 &\quad-\mathcal C(g_2g_3)\mathcal Cg_1
 +2\mathcal Cg_1\mathcal Cg_2\mathcal Cg_3.
 \label{eq:V31-ternary-assembled}
\end{align}

\begin{assessment}[Exact post-V31 bottleneck]
\label{ass:V31-bottleneck}
The first unproved spatial inequality is
\begin{equation}
 \sum_{\varnothing\ne Z\subseteq X_1\cup X_2\cup X_3}
 \norm{\bigl[\mathcal R_{\alpha,\Lambda}Q_0
       \kappa_3^\alpha(g_{1,X_1},g_{2,X_2},g_{3,X_3})\bigr]_Z}
 _{\rho,2}^{\rm BFA}
 \le\text{the expression in \eqref{eq:V31-ternary-MTA}}.
 \label{eq:V31-ternary-open}
\end{equation}
It must be proved after the five terms in
\eqref{eq:V31-ternary-assembled} have been combined.  Applying the binary
estimate twice is not a proof: the ternary cumulant is not an iterated
binary cumulant, and such an argument may spend the single output resolvent
twice.  Conversely, an unmarked support-tree bound is too weak by
\Cref{prop:V31-unmarked-loss}.  Thus the remaining WUFC question is an
explicit two-link-mark Wilson fusion identity beginning at \(m=3\), not an
unspecified large-support count.
\end{assessment}

% =============================================================================
\section{V31 integrated result and next frontier}
\label{sec:V31-ledger}
% =============================================================================

\begin{theorem}[Integrated V31 statement]
\label{thm:V31-integrated}
Preserving V1--V30, V31 proves:
\begin{enumerate}[label=\textup{(V31.\arabic*)},leftmargin=5.5em]
\item the BFA/BFI norm is uniformly equivalent to the V29 FA/FI norm and
      decomposes exactly into normalized link marks of every degree;
\item the Wilson integration-by-parts estimate
      \(C_2(R)\eta_R\lesssim\alpha\), the one-link tail debit
      \eqref{eq:V31-one-link-tail-debit}, and uniform product damping
      \eqref{eq:V31-balanced-product-damping};
\item the marked binary scalar quotient
      \eqref{eq:V31-marked-binary-quotient}, uniform in support size,
      representation height, volume, and \(\alpha\);
\item the global binary WUFC estimate
      \eqref{eq:V31-binary-FI} in the actual rooted FI norm;
\item a sharp abstract example proving that the unmarked one-bank bound
      cannot be summed in FI; and
\item the marked-tree spatial compiler
      \(\operatorname{MTA}\Rightarrow\operatorname{WUFC}\Rightarrow
      \operatorname{CPM}\), reducing the remaining all-order spatial gate
      to the explicit assembled atlas beginning with
      \eqref{eq:V31-ternary-open}.
\end{enumerate}
\end{theorem}

\begin{proof}
Item \textup{(V31.1)} is
\Cref{lem:V31-balanced-equivalence}.  Item \textup{(V31.2)} is
\Cref{lem:V31-density-Laplacian,thm:V31-tail-debit}.  Items
\textup{(V31.3)--(V31.4)} are
\Cref{thm:V31-marked-binary-quotient,
thm:V31-marked-binary-operator,thm:V31-binary-WUFC}.  Item
\textup{(V31.5)} is \Cref{prop:V31-unmarked-loss}.  Item
\textup{(V31.6)} is
\Cref{def:V31-MTA,thm:V31-MTA-compiler,cor:V31-MTA-route,
ass:V31-bottleneck}.
\end{proof}

\begingroup
\small
\begin{longtable}{>{\raggedright\arraybackslash}p{0.29\textwidth}
                  >{\raggedright\arraybackslash}p{0.15\textwidth}
                  >{\raggedright\arraybackslash}p{0.48\textwidth}}
\toprule
\textbf{V31 row}&\textbf{Status}&\textbf{Advance or remaining obstruction}\\
\midrule
\endfirsthead
\toprule
\textbf{V31 row}&\textbf{Status}&\textbf{Advance or remaining obstruction}\\
\midrule
\endhead
Balanced link-mark norm
& \MGclosed
& Exact probability normalization at every vertex degree, uniformly
  equivalent to the declared FI norm.\\[0.35em]
Wilson high-character damping
& \MGclosed
& Integration by parts gives \eqref{eq:V31-one-link-tail-debit} and the
  support-uniform product moment.\\[0.35em]
Binary large-support assembly
& \MGclosed
& The marked quotient closes \(m=2\) in FI for arbitrary supports and
  representations.\\[0.35em]
Unmarked cluster counting
& \MGfail
& \Cref{prop:V31-unmarked-loss} proves an unavoidable \(N\)-fold loss.\\[0.35em]
All-order marked-tree atlas
& \MGopen
& The spatial compiler is proved, but the assembled local atlas estimate
  begins with the ternary two-mark inequality.\\[0.35em]
Connected Perron majorant
& \MGreduced
& It follows from MTA through the V31 and V30 compilers.\\[0.35em]
V24/V25 unreduced edge packet
& \MGopen
& It remains downstream of WUFC and requires the direct singular-channel
  calculation prescribed by the V30 drift audit.\\[0.35em]
Capacity, protected sectors, and OS limit
& \MGopen
& No new capacity or continuum claim is made in V31.\\[0.35em]
Full Yang--Mills mass gap
& \MGopen
& Ternary and higher marked fusion, finite edges, and continuum closure
  remain.\\
\bottomrule
\end{longtable}
\endgroup

\begin{openproblem}[V32 mandate: assembled multi-mark fusion]
\label{op:V32}
\begin{enumerate}[label=\textup{(V32.\arabic*)},leftmargin=5.5em]
\item Resolve \eqref{eq:V31-ternary-assembled} into the exact product
      Peter--Weyl atlas and prove \eqref{eq:V31-ternary-open}.  Preserve all
      five terms until the two shared-link defects have been exposed.
\item Test a replicated-coordinate or forest interpolation formula for the
      general cumulant.  Each tree edge must carry one actual shared link,
      and each vertex of degree \(d\) must retain the \(d\)-mark seminorm in
      \eqref{eq:V31-marked-BFA}.
\item Prove the general MTA constant grows only exponentially per tree edge.
      Do not add a factorial before the tree sum; Cayley's tree count already
      supplies the factorial-exponential growth allowed by WUFC.
\item If MTA fails, return a finite resolved representation/support/mark
      witness.  Do not report the abstract unmarked incidence example as a
      Wilson failure.
\item On a strict MTA branch, invoke
      \Cref{thm:V31-MTA-compiler,thm:V30-WUFC-compiler} and populate the
      complete Harnack, density, incidence, and temporal-source packet.
\item Then form the full unreduced V24/V25 edge matrix directly at every
      singular channel, retain the reconstruction topology and all signed
      complements, and proceed through capacity, protected sectors,
      local-algebra noncollapse, reflection positivity, and the nontrivial
      OS continuum state.
\item The next scheduled SM/NS companion audit remains V35.
\end{enumerate}
\end{openproblem}

\begin{firewall}[End-of-V31 claim boundary]
V31 proves the global binary FI estimate and an exact all-order spatial
compiler from the marked-tree atlas.  It does not prove the ternary or
higher MTA inequality.  Therefore WUFC, CPM, the nonlinear Perron packet,
the unreduced edge margins, bad capacity, protected-sector floors,
local-algebra noncollapse, and the nontrivial continuum OS limit remain
open.  The full Yang--Mills mass gap is not proved.
\end{firewall}

\section*{Version V31 append-only record}
\addcontentsline{toc}{section}{Version V31 append-only record}
The complete V30 mathematical source was retained before this continuation.
Its SHA--256 digest was
\begin{center}\scriptsize\ttfamily
7a69fa4a7a8e1d07ec4d5ebf08bb500b9cdd867d262e954545b6daf8f546e69d.
\end{center}
No mathematical content of V30 was altered.  On the next iteration, retain
this full file, append before its unique active document-closing command,
increment the filename to \texttt{\_V32.tex}, and remove only the
superseded V31 file from this note series.  The next companion audit is
V35.

% =============================================================================
% END APPENDED VERSION V31 MATERIAL
% =============================================================================

% =============================================================================
% BEGIN APPENDED VERSION V32 MATERIAL
% =============================================================================

\clearpage
\part{V32 --- Ternary replicated-coordinate trees, bounded thermal MTA,
and quadratic Wilson tail debits}
\label{part:V32}

\section{Purpose and exact scope}
\label{sec:V32-purpose}

V31 proved the marked-tree atlas inequality at order two and showed that
the all-order MTA would compile without loss into WUFC.  The first open
case was the assembled five-term ternary cumulant.  V32 attacks that case
without iterating the resolved binary estimate.

There are two different ternary connections.  Two distinct coordinates
may carry two pair bonds forming a tree, or one coordinate common to all
three inputs may carry a genuine local third cumulant.  Product-row
resampling controls the first mechanism.  Inversion symmetry and a
fourth displacement moment control the second at order \(s_\alpha^2\).
Together they prove ternary MTA uniformly on every fixed thermal window
\(s_\alpha\Xi\le L\), with no support-size or representation cutoff
inside that window.

The constant obtained below depends on \(L\).  V32 therefore does not
claim full ternary MTA.  A new bi-Laplacian estimate proves the quadratic
Wilson tail debit needed to investigate \(L\to\infty\), and the final
section identifies the remaining high-to-low fusion-cancellation branch.
There is no scheduled companion audit at V32; the next audit remains V35.

% =============================================================================
\section{A second Wilson character-tail moment}
\label{sec:V32-quadratic-tail}
% =============================================================================

\begin{lemma}[Bi-Laplacian moment of the Wilson density]
\label{lem:V32-density-bilaplacian}
There are (C_{\Delta^2}<\infty) and \(\alpha_{\Delta^2}<\infty\) such
that
\begin{equation}
 \norm{\Delta_G^2q_\alpha}_{L^1(G)}
 \le C_{\Delta^2}\alpha^2
 \qquad(\alpha\ge\alpha_{\Delta^2}).
 \label{eq:V32-density-bilaplacian}
\end{equation}
Consequently,
\begin{equation}
 C_2(R)^2\eta_R(\alpha)
 \le C_{\Delta^2}\alpha^2
 \label{eq:V32-character-second-moment}
\end{equation}
for every irreducible representation \(R\).
\end{lemma}

\begin{proof}
Retain the notation \(q_\alpha=\widetilde Z_\alpha^{-1}e^{-\alpha W}\)
from \Cref{lem:V31-density-Laplacian}.  The global quadratic well gives,
for \(r=1,2\),
\begin{equation}
 \int_G W^r q_\alpha\,\dd U\le C_r\alpha^{-r}.
 \label{eq:V32-W-moments}
\end{equation}
The proof is the same normal-coordinate Gaussian comparison as
\eqref{eq:V31-Wilson-well-moment}, with one or two additional quadratic
powers.

Apply \(\Delta_G\) once more to
\[
 \Delta_Gq_\alpha
 =q_\alpha\left(-\alpha\Delta_GW
                 +\alpha^2\lvert\nabla W\rvert^2\right).
\]
Every resulting term is a bounded derivative of \(W\), multiplied by one
of
\[
 \alpha,\qquad \alpha^2,\qquad
 \alpha^3\lvert\nabla W\rvert^2,\qquad
 \alpha^4\lvert\nabla W\rvert^4.
\]
On the compact group,
\(\lvert\nabla W\rvert^2\le CW\).  Thus
\eqref{eq:V32-W-moments} bounds the \(L^1(q_\alpha\dd U)\) norm of every
term by \(C\alpha^2\), proving
\eqref{eq:V32-density-bilaplacian}.

Twice integrate by parts in the matrix Fourier transform and use
\(\Delta_G^2D^R=C_2(R)^2D^R\).  Unitarity of \(D^R\) gives
\[
 C_2(R)^2\eta_R
 \le\norm{\Delta_G^2q_\alpha}_{L^1(G)},
\]
which proves \eqref{eq:V32-character-second-moment}.
\end{proof}

\begin{theorem}[Quadratic tail debit and product second moment]
\label{thm:V32-quadratic-tail-debit}
There are (K_{\rm tail,2}<\infty) and \(\alpha_{\rm tail,2}<\infty\)
such that
\begin{equation}
 \boxed{
 \bigl(s_\alpha C_2(R)\bigr)^2\eta_R(\alpha)
 \le K_{\rm tail,2}\bigl(1-\eta_R(\alpha)\bigr)}
 \label{eq:V32-one-link-quadratic-debit}
\end{equation}
for every irreducible \(R\) and every
\(\alpha\ge\alpha_{\rm tail,2}\).  If
\(\boldsymbol R=(R_e)_{e\in E}\) is exact and nontrivial, then
\begin{align}
 \bigl(s_\alpha C_E(\boldsymbol R)\bigr)^2
 q_{\alpha,\boldsymbol R}
 &\le K_{\rm prod,2},
 \label{eq:V32-product-second-moment}\\
 \bigl(s_\alpha\Xi(\boldsymbol R)\bigr)^2
 q_{\alpha,\boldsymbol R}
 &\le K_{\rm bal,2}.
 \label{eq:V32-balanced-second-moment}
\end{align}
All constants are independent of height and support size.
\end{theorem}

\begin{proof}
Put \(x=s_\alpha C_2(R)\).  If \(x\le1\), then
\(x^2\eta_R\le x\), and the lower gap in
\Cref{thm:V30-two-sided-gap} bounds this by
\(C(1-\eta_R)\).  If \(x\ge1\), then the same lower gap is bounded below
by a positive group constant, while
\Cref{lem:V32-density-bilaplacian} and
\(s_\alpha\alpha\le C\) give \(x^2\eta_R\le C\).
This proves \eqref{eq:V32-one-link-quadratic-debit}.

Write \(x_e=s_\alpha C_2(R_e)\), \(r_e=\eta_{R_e}\), and
\(q=\prod_er_e\).  Expanding the square gives
\begin{align*}
 \left(\sum_ex_e\right)^2q
 &=\sum_ex_e^2r_e\prod_{f\ne e}r_f\\
 &\quad+2\sum_{e<f}
   x_er_ex_fr_f\prod_{h\ne e,f}r_h.
\end{align*}
The first sum is bounded by
\(K\sum_e(1-r_e)\prod_{f\ne e}r_f\le K\), using
\eqref{eq:V32-one-link-quadratic-debit}.  Apply the first-order debit
\eqref{eq:V31-one-link-tail-debit} twice to the second sum.  It becomes at
most a constant times
\[
 \sum_{e<f}(1-r_e)(1-r_f)
 \prod_{h\ne e,f}r_h,
\]
the probability of exactly two failures in independent Bernoulli trials,
and is therefore at most one.  This proves
\eqref{eq:V32-product-second-moment}.  Finally
\(\Xi\le(1+C_F^{-1})C_E\) on an exact nonempty support, proving
\eqref{eq:V32-balanced-second-moment}.
\end{proof}

% =============================================================================
\section{Uniform one-link ternary Gaussian excess}
\label{sec:V32-one-link-ternary}
% =============================================================================

For one link, let \(f_i\) be a normalized matrix coefficient of an
irreducible \(R_i\), evaluated after the Wilson row increment.  Put
\begin{equation}
 A_i:=1+C_2(R_i).
 \label{eq:V32-one-link-masses}
\end{equation}

\begin{theorem}[Representation-uniform one-link third cumulant]
\label{thm:V32-one-link-third-cumulant}
There are (K_{3,1}<\infty) and \(\alpha_{3,1}<\infty\) such that
\begin{equation}
 \boxed{
 \left|\kappa_{3,\mathrm{one}}^\alpha(f_1,f_2,f_3)\right|
 \le K_{3,1}s_\alpha^2 A_1A_2A_3}
 \label{eq:V32-one-link-third-cumulant}
\end{equation}
for all irreducibles, all normalized matrix coefficients, all base points,
and \(\alpha\ge\alpha_{3,1}\).  The bound is also valid with its right
side replaced by its minimum with a universal constant.
\end{theorem}

\begin{proof}
Cumulants of order at least two are unchanged by adding constants.  In a
symmetric exponential ball write \(U=\exp X\) and replace each coefficient
by \(f_i(U)-f_i(I)\).  Taylor's formula gives
\begin{equation}
 f_i(\exp X)-f_i(I)=L_i(X)+Q_i(X),
 \label{eq:V32-matrix-Taylor}
\end{equation}
with
\begin{equation}
 |L_i(X)|\le cA_i^{1/2}\norm X,
 \qquad
 |Q_i(X)|\le cA_i\norm X^2.
 \label{eq:V32-matrix-jet-bounds}
\end{equation}
Indeed the first and second derived representation operators are bounded
by a group constant times \(C_2(R_i)^{1/2}\) and \(C_2(R_i)\),
respectively.

The Wilson law and Haar measure are invariant under \(U\mapsto U^{-1}\).
On the symmetric exponential ball this is \(X\mapsto-X\), so
\begin{equation}
 \mathbb E[L_1L_2L_3;\ U\text{ in the ball}]=0.
 \label{eq:V32-cubic-odd-cancellation}
\end{equation}
The V11 Gaussian comparison gives
\begin{equation}
 \mathbb E\norm X^{2r}\le C_r s_\alpha^r
 \qquad(r=1,2,3),
 \label{eq:V32-displacement-moments}
\end{equation}
and an \(O(e^{-c/s_\alpha})\) complement.  Expanding
\eqref{eq:V32-matrix-Taylor}, using
\eqref{eq:V32-cubic-odd-cancellation}, shows
\begin{align}
 |\mathbb E[f_1f_2f_3]|
 &\le Cs_\alpha^2A_1A_2A_3,
 \label{eq:V32-third-moment-bound}\\
 |\mathbb E f_i|
 &\le Cs_\alpha A_i,
 \label{eq:V32-first-moment-bound}\\
 |\mathbb E[f_if_j]|
 &\le Cs_\alpha A_iA_j.
 \label{eq:V32-second-moment-bound}
\end{align}
For example, after the odd cubic term, the lowest contribution to the
first line has one \(Q\) and two \(L\)'s and is bounded by a fourth
displacement moment.  Terms with more remainders use the fifth and sixth
moments and are smaller for \(s_\alpha\le1\).  The exterior exponential
is bounded by \(Cs_\alpha^2A_1A_2A_3\), because every \(A_i\ge1\).

Insert \eqref{eq:V32-third-moment-bound}--%
\eqref{eq:V32-second-moment-bound} into the centered identity
\begin{align*}
 \kappa_3(f_1,f_2,f_3)
 &=\mathbb E[f_1f_2f_3]
 -\sum_{\{i,j,k\}=\{1,2,3\}}
   \mathbb E[f_if_j]\mathbb E f_k
 +2\prod_i\mathbb E f_i.
\end{align*}
Every term is bounded by the right side of
\eqref{eq:V32-one-link-third-cumulant}.  Finally, bounded matrix
coefficients give a universal crude cumulant bound, which proves the
capped assertion.
\end{proof}
% =============================================================================
\section{A three-variable replicated-coordinate inequality}
\label{sec:V32-replicated-coordinate}
% =============================================================================

Let \((\Omega_e,\nu_e)_{e\in E}\) be finitely many probability spaces,
with product law \(\nu\).  For \(i\in\{1,2,3\}\), let
\(f_{i,e}:\Omega_e\to\mathbb C\) satisfy \(\lvert f_{i,e}\rvert\le1\),
put
\begin{equation}
 F_i:=\prod_{e\in E}f_{i,e},
 \qquad
 \gamma_{i,e}^2:=\frac12
 \mathbb E\lvert f_{i,e}(\omega_e)-f_{i,e}(\omega_e')\rvert^2,
 \label{eq:V32-coordinate-influence}
\end{equation}
and take \(\omega_e'\) to be an independent copy of \(\omega_e\).  Set
\begin{align}
 W_{ij}&:=\sum_{e\in E}\gamma_{i,e}\gamma_{j,e},
 &\overline W_{ij}&:=\min\{W_{ij},1\},
 \label{eq:V32-pair-influence}\
 V_{123}&:=\sum_{e\in E}
 \left\lvert\kappa_3^{\nu_e}
 (f_{1,e},f_{2,e},f_{3,e})\right\rvert.
 \label{eq:V32-local-junction}
\end{align}

\begin{lemma}[Replicated-coordinate ternary tree bound]
\label{lem:V32-replicated-tree}
There is a numerical constant \(C_{\rm rt}<\infty\) such that
\begin{equation}
 \boxed{
 \left\lvert\kappa_3^\nu(F_1,F_2,F_3)\right\rvert
 \le C_{\rm rt}\left(
 V_{123}
 +\overline W_{12}\overline W_{13}
 +\overline W_{12}\overline W_{23}
 +\overline W_{13}\overline W_{23}
 \right).}
 \label{eq:V32-replicated-tree}
\end{equation}
The constant is independent of \(E\).  If \(f_{i,e}=1\) off a support
\(X_i\), only intersections permitted by the three labeled trees occur
on the right side.
\end{lemma}

\begin{proof}
We first give the scalar real proof.  Write
\(a_{i,e}=\mathbb E f_{i,e}\),
\(g_{i,e}=f_{i,e}-a_{i,e}\), and
\begin{equation}
 p_{ij,e}:=\mathbb E(g_{i,e}g_{j,e}),
 \qquad
 t_e:=\mathbb E(g_{1,e}g_{2,e}g_{3,e}).
 \label{eq:V32-local-connected-pieces}
\end{equation}
The independent-copy covariance identity and Cauchy--Schwarz give
\begin{equation}
 \lvert p_{ij,e}\rvert
 \le\gamma_{i,e}\gamma_{j,e},
 \qquad
 t_e=\kappa_3^{\nu_e}(f_{1,e},f_{2,e},f_{3,e}).
 \label{eq:V32-local-piece-bounds}
\end{equation}

Expand each \(F_i=\prod_e(a_{i,e}+g_{i,e})\), take expectations
coordinatewise, and then form the five-term centered combination defining
\(\kappa_3\).  A coordinate at which exactly one centered factor is
chosen contributes zero.  A surviving coordinate therefore carries
either one of the pair bonds \(p_{12,e},p_{13,e},p_{23,e}\), or the
three-vertex junction \(t_e\).  The four disconnected incidence patterns
cancel in the centered combination.  Thus every remaining monomial
contains a junction, or pair bonds of at least two distinct types.  All
unused mean factors have modulus at most one.

Fix \(\varepsilon=1/4\).  If all \(W_{ij}\le\varepsilon\) and
\(V_{123}\le1\), absolute summation of the preceding expansion is
legitimate and
\begin{align}
 \prod_e\left(1+\sum_{i<j}\lvert p_{ij,e}\rvert+lvert t_e\rvert\right)
 &\le
 \exp\left(\sum_{i<j}W_{ij}+V_{123}\right)
 \le e^{7/4}.
 \label{eq:V32-cluster-envelope}
\end{align}
Mark the first junction in a junction-containing monomial.  In a
pair-only connected monomial, mark the first bond of each of two distinct
types.  Summing the unmarked remainder with
\eqref{eq:V32-cluster-envelope} yields
\begin{equation}
 \lvert\kappa_3(F_1,F_2,F_3)\rvert
 \le e^{7/4}\left(
 V_{123}+W_{12}W_{13}+W_{12}W_{23}+W_{13}W_{23}
 \right).
 \label{eq:V32-small-influence-tree}
\end{equation}

It remains to remove the smallness condition.  A bounded third cumulant
has a universal crude bound.  Hence \(V_{123}\ge1\), or two of the three
\(W_{ij}\)'s being at least \(\varepsilon\), is absorbed by the right
side of \eqref{eq:V32-replicated-tree}.  In the only remaining case,
say \(W_{12}\ge\varepsilon\) and
\(W_{13},W_{23}<\varepsilon\), use
\begin{align}
 \kappa_3(F_1,F_2,F_3)
 &=\operatorname{Cov}(F_1F_2,F_3)
 -\mathbb EF_1\operatorname{Cov}(F_2,F_3)
 \notag\
 &\quad
 -\mathbb EF_2\operatorname{Cov}(F_1,F_3).
 \label{eq:V32-cumulant-as-cross-covariance}
\end{align}
The product-measure covariance inequality
\begin{equation}
 \lvert\operatorname{Cov}(H,K)\rvert
 \le\sum_e
 \left(\frac12\mathbb E\lvert H-H^{(e)}\rvert^2\right)^{1/2}
 \left(\frac12\mathbb E\lvert K-K^{(e)}\rvert^2\right)^{1/2}
 \label{eq:V32-tensorized-covariance}
\end{equation}
follows by martingale differences and Cauchy--Schwarz; \(H^{(e)}\)
denotes replacement of coordinate \(e\) by an independent copy.  Since
\(\lvert f_{i,e}\rvert\le1\), the coordinate influence of \(F_1F_2\)
is at most \(\gamma_{1,e}+\gamma_{2,e}\).  Applying
\eqref{eq:V32-tensorized-covariance} to
\eqref{eq:V32-cumulant-as-cross-covariance} gives
\begin{equation}
 \lvert\kappa_3(F_1,F_2,F_3)\rvert
 \le C(W_{13}+W_{23}).
 \label{eq:V32-one-strong-edge}
\end{equation}
Because \(\overline W_{12}\ge\varepsilon\), this is bounded by the two
tree products containing edge \(12\).  The other choices of the strong
edge are identical.  This proves the real case.  Applying it to real and
imaginary parts and polarizing changes only the numerical constant.
\end{proof}

% =============================================================================
\section{Wilson product rows and the two-mark numerator}
\label{sec:V32-Wilson-product-row}
% =============================================================================

Let \(\boldsymbol R_i\) be exact on \(X_i\), and define
\begin{align}
 A_{i,e}&:=1+C_2(R_{i,e}) &&(e\in X_i),
 \label{eq:V32-product-link-mass}\
 H_{ij}&:=\sum_{e\in X_i\cap X_j}A_{i,e}A_{j,e},
 \label{eq:V32-pair-overlap-mass}\
 J_{123}&:=\sum_{e\in X_1\cap X_2\cap X_3}
 A_{1,e}A_{2,e}A_{3,e},
 \label{eq:V32-triple-junction-mass}\
 \mathcal M_{123}&:=J_{123}
 +H_{12}H_{13}+H_{12}H_{23}+H_{13}H_{23}.
 \label{eq:V32-two-mark-numerator}
\end{align}
We put \(A_{i,e}=0\) when \(e\notin X_i\).

\begin{proposition}[Wilson replicated-coordinate numerator]
\label{prop:V32-Wilson-numerator}
For normalized product matrix coefficients \(F_i\) in the three exact
channels above,
\begin{equation}
 \boxed{
 \left\lvert\kappa_3^\alpha(F_1,F_2,F_3)\right\rvert
 \le K_{\rm W3}s_\alpha^2\mathcal M_{123}.}
 \label{eq:V32-Wilson-numerator}
\end{equation}
The same estimate holds, with trace norms, after simultaneous
Clebsch--Gordan decomposition into all final product output channels.
Its constant is independent of the supports, the volume, and the
representation heights.
\end{proposition}

\begin{proof}
On link \(e\), resample the Wilson increment.  The geodesic derivative
bound for a normalized matrix coefficient, the second displacement
moment in \eqref{eq:V32-displacement-moments}, and the crude bound by two
give
\begin{equation}
 \gamma_{i,e}
 \le C\min\{s_\alpha^{1/2}A_{i,e}^{1/2},1\}
 \le Cs_\alpha^{1/2}A_{i,e}^{1/2}.
 \label{eq:V32-Wilson-coordinate-influence}
\end{equation}
Consequently
\begin{equation}
 \overline W_{ij}\le Cs_\alpha H_{ij}.
 \label{eq:V32-Wilson-pair-bond}
\end{equation}
At a common link, \Cref{thm:V32-one-link-third-cumulant} gives
\begin{equation}
 V_{123}\le Cs_\alpha^2J_{123}.
 \label{eq:V32-Wilson-junction}
\end{equation}
Insert \eqref{eq:V32-Wilson-pair-bond} and
\eqref{eq:V32-Wilson-junction} into
\Cref{lem:V32-replicated-tree} to prove
\eqref{eq:V32-Wilson-numerator}.

For the trace-norm statement, run the same independent-copy expansion at
the matrix-tensor level.  Each coordinate difference acts on only the
corresponding tensor factor.  At a pair bond use Hilbert-space
Cauchy--Schwarz; at a junction use
\Cref{thm:V32-one-link-third-cumulant}.  Decompose the resulting tensor
products in a common final bracketing.  Clebsch--Gordan maps and the
Racah changes between the three pair bracketings are unitary, so trace
norm is contractive under every block compression.  The sum over all
final blocks is the trace norm of the block diagonal tensor operator.
This is the three-input version of the contraction argument in
\Cref{lem:V29-FA-algebra,thm:V31-marked-binary-operator}; it introduces
no multiplicity or dimension factor.  Hence the same scalar numerator
multiplies the product of the three input Peter--Weyl trace weights.
\end{proof}

% =============================================================================
\section{Ternary MTA on every bounded thermal window}
\label{sec:V32-bounded-thermal}
% =============================================================================

\begin{definition}[Thermal-core channel]
\label{def:V32-thermal-core}
For \(L\ge1\), an exact product representation \(\boldsymbol R\) is
\(L\)-thermal-core at scale \(\alpha\) if
\begin{equation}
 s_\alpha\Xi(\boldsymbol R)\le L.
 \label{eq:V32-thermal-core}
\end{equation}
An exact-support input is \(L\)-thermal-core if its Peter--Weyl
coefficients vanish outside these channels.
\end{definition}

\begin{lemma}[Bounded-thermal resolved two-mark quotient]
\label{lem:V32-thermal-two-mark-quotient}
Fix \(L\ge1\).  Suppose the three exact input channels are
\(L\)-thermal-core.  For every nonconstant exact final channel
\(\boldsymbol T\subset
 \boldsymbol R_1\otimes\boldsymbol R_2\otimes\boldsymbol R_3\),
\begin{equation}
 \boxed{
 \frac{s_\alpha^2\Xi(\boldsymbol T)}
      {1-q_{\alpha,\boldsymbol T}}\mathcal M_{123}
 \le K_L\sum_{i=1}^3
 \frac{H_{ij}H_{ik}}{\Xi(\boldsymbol R_i)}.}
 \label{eq:V32-thermal-two-mark-quotient}
\end{equation}
For each center \(i\), \(\{j,k\}=\{1,2,3\}\setminus\{i\}\).
The constant is independent of all supports, output channels, and
representation heights inside the stated window.
\end{lemma}

\begin{proof}
Write \(\Xi_i=\Xi(\boldsymbol R_i)\) and
\(\Xi_T=\Xi(\boldsymbol T)\).  Linkwise fusion geometry, summed over the
output support, gives
\begin{equation}
 \Xi_T\le c\sum_{i=1}^3\Xi_i\le\frac{3cL}{s_\alpha}.
 \label{eq:V32-thermal-output-mass}
\end{equation}
The lower product gap \eqref{eq:V30-two-sided-product-gap}, together with
\(\Xi_T\le(1+C_F^{-1})C_E(\boldsymbol T)\), implies
\begin{equation}
 \frac{\Xi_T}{1-q_{\alpha,\boldsymbol T}}
 \le\frac{C_L}{s_\alpha}.
 \label{eq:V32-thermal-output-resolvent}
\end{equation}
Indeed, when \(s_\alpha C_E(\boldsymbol T)\le1\) this follows directly
from the linear part of the gap, and in the complementary case the gap
is bounded below while \eqref{eq:V32-thermal-output-mass} applies.

For each centered tree term, the thermal assumption gives
\begin{equation}
 s_\alpha H_{ij}H_{ik}
 \le L\frac{H_{ij}H_{ik}}{\Xi_i}.
 \label{eq:V32-thermal-pair-payment}
\end{equation}
For a common link \(e\), the diagonal choice of the same mark on the two
edges of the tree centered at \(i\) contributes
\[
 \frac{A_{i,e}^2A_{j,e}A_{k,e}}{\Xi_i}.
\]
Since \(A_{i,e}\ge1\) on an active link and
\(\Xi_i^{-1}\ge s_\alpha/L\),
\begin{equation}
 s_\alpha A_{1,e}A_{2,e}A_{3,e}
 \le\frac L3 A_{1,e}A_{2,e}A_{3,e}
 \sum_{i=1}^3\frac{A_{i,e}}{\Xi_i}.
 \label{eq:V32-junction-diagonal-payment}
\end{equation}
After summing \(e\), the right side is contained in the three full tree
sums \(\sum_iH_{ij}H_{ik}/\Xi_i\).  Multiply
\eqref{eq:V32-thermal-output-resolvent} by
\(s_\alpha^2\mathcal M_{123}\), then use
\eqref{eq:V32-thermal-pair-payment} and
\eqref{eq:V32-junction-diagonal-payment}.  This proves
\eqref{eq:V32-thermal-two-mark-quotient}.
\end{proof}

\begin{theorem}[Ternary marked-tree atlas on a bounded thermal window]
\label{thm:V32-bounded-thermal-MTA}
Fix \(L\ge1\).  There are \(K_{3,L}<\infty\) and
\(\alpha_{3,L}<\infty\), independent of support size, volume, and
representation height inside the thermal window, such that every three
\(L\)-thermal-core exact-support inputs satisfy
\begin{align}
 &\sum_{\varnothing\ne Z\subseteq X_1\cup X_2\cup X_3}
 \norm{\bigl[\mathcal R_{\alpha,\Lambda}Q_0
       \kappa_3^\alpha(g_{1,X_1},g_{2,X_2},g_{3,X_3})\bigr]_Z}
       _{\rho,2}^{\rm BFA}
 \notag\\
 &\quad\le K_{3,L}^2\sum_{i=1}^3
 \sum_{\substack{e\in X_i\cap X_j\\
                   f\in X_i\cap X_k}}
 \norm{g_{i,X_i}}_{\rho,2;(e,f)}^{\rm BFA(2)}
 \norm{g_{j,X_j}}_{\rho,2;e}^{\rm BFA(1)}
 \norm{g_{k,X_k}}_{\rho,2;f}^{\rm BFA(1)}.
 \label{eq:V32-bounded-thermal-MTA}
\end{align}
Thus \eqref{eq:V31-ternary-open} is proved on every fixed thermal-core
window, with no spatial support cutoff.
\end{theorem}

\begin{proof}
Expand the inputs into exact product Peter--Weyl blocks and decompose each
threefold product into nonconstant exact output blocks.  By
\Cref{prop:V32-Wilson-numerator}, the assembled five-term numerator is
bounded before any denominator is charged.  Charge the single final
resolvent and apply
\Cref{lem:V32-thermal-two-mark-quotient}.  The exponential height ratio
is at most one by \eqref{eq:V30-height-fusion}.  Clebsch--Gordan block
compression and all required Racah changes are trace-norm contractions,
as in the last paragraph of
\Cref{prop:V32-Wilson-numerator}.

For a fixed center \(i\), expand
\begin{equation}
 \frac{H_{ij}H_{ik}}{\Xi_i}
 =\sum_{\substack{e\in X_i\cap X_j\\f\in X_i\cap X_k}}
 \frac{A_{i,e}A_{i,f}}{\Xi_i}A_{j,e}A_{k,f}.
 \label{eq:V32-tree-mass-factorization}
\end{equation}
The four scalar factors in each summand are exactly the two marks at the
center and the one mark at each leaf after multiplying by the three BFA
input weights:
\[
 \Xi_i\vartheta_e(\boldsymbol R_i)\vartheta_f(\boldsymbol R_i)
 \cdot\Xi_j\vartheta_e(\boldsymbol R_j)
 \cdot\Xi_k\vartheta_f(\boldsymbol R_k)
 =\frac{A_{i,e}A_{i,f}}{\Xi_i}A_{j,e}A_{k,f}.
\]
Tonelli's theorem now factors the nonnegative trace-norm sums into the
three marked seminorms displayed in
\eqref{eq:V32-bounded-thermal-MTA}.  This includes every exact
nonconstant output once and proves the theorem.
\end{proof}

% =============================================================================
\section{Thermal escape and reverse fusion}
\label{sec:V32-thermal-escape}
% =============================================================================

For an exact product channel put
\begin{equation}
 \Theta_\alpha(\boldsymbol R):=
 s_\alpha\Xi(\boldsymbol R).
 \label{eq:V32-thermal-ratio}
\end{equation}
The quadratic debit has the immediate consequence
\begin{equation}
 q_{\alpha,\boldsymbol R}
 \le\min\left\{1,
 \frac{K_{\rm bal,2}}{\Theta_\alpha(\boldsymbol R)^2}\right\}.
 \label{eq:V32-quadratic-thermal-damping}
\end{equation}
Thus a multiplier attached to a thermally escaping channel decays at
least quadratically in its thermal ratio.  The issue is to identify, in
one common recoupling basis, which channel carries that multiplier in
each of the five assembled terms.

\begin{theorem}[Every remaining ternary failure escapes the thermal core]
\label{thm:V32-ternary-thermal-escape}
Suppose that no constant independent of \(\alpha\), volume, support size,
and representation height proves the full ternary instance of MTA.
Then there is a sequence of finite, pure resolved witnesses
\begin{equation}
 w_r=(\alpha_r,\Lambda_r;
       \boldsymbol R_{r,1},\boldsymbol R_{r,2},\boldsymbol R_{r,3};
       \boldsymbol T_r),
 \qquad \alpha_r\longrightarrow\infty,
 \label{eq:V32-thermal-escape-witness}
\end{equation}
whose ratio of the left side of \eqref{eq:V31-ternary-open} to its
three marked-tree weights tends to infinity, and for which
\begin{equation}
 \max_{1\le i\le3}
 \Theta_{\alpha_r}(\boldsymbol R_{r,i})
 \longrightarrow\infty.
 \label{eq:V32-input-thermal-escape}
\end{equation}
In particular neither spatial support size alone nor a collection of
channels remaining in one fixed thermal window can witness failure.
\end{theorem}

\begin{proof}
The exact-support and Peter--Weyl expansions are finite on a polynomial
witness.  Multilinearity, the nonnegative trace-norm majorant, and block
compression therefore refine a failed estimate to a pure resolved input
triple and one nonconstant output block.  Choose the witness at stage
\(r\) with cofinal \(\alpha_r\) and marked-tree ratio larger than \(r\).
If \eqref{eq:V32-input-thermal-escape} failed, a subsequence would satisfy
\(\Theta_{\alpha_r}(\boldsymbol R_{r,i})\le L\) for all three inputs and
one fixed \(L\).  For all sufficiently large \(r\),
\Cref{thm:V32-bounded-thermal-MTA} would bound its ratio by
\(K_{3,L}^2\), a contradiction.  If the marked-tree denominator were
zero, the support intersection graph would be disconnected and product
independence would make the cumulant zero, so no exceptional zero
denominator arises.
\end{proof}

\begin{lemma}[Reverse fusion mass geometry]
\label{lem:V32-reverse-fusion}
If a nonconstant product channel satisfies
\(\boldsymbol T\subset
\boldsymbol R_1\otimes\boldsymbol R_2\otimes\boldsymbol R_3\), then
for each \(i\),
\begin{equation}
 \boxed{
 \Xi(\boldsymbol R_i)
 \le C_{\rm rf}\left(
 \Xi(\boldsymbol T)+\Xi(\boldsymbol R_j)+\Xi(\boldsymbol R_k)
 \right),}
 \label{eq:V32-reverse-fusion}
\end{equation}
where \(\{j,k\}=\{1,2,3\}\setminus\{i\}\).  The same statement holds
for either binary intermediate channel in any bracketing of the
threefold tensor product.
\end{lemma}

\begin{proof}
At one link, Frobenius reciprocity gives
\[
 R_i\subset T\otimes R_j^*\otimes R_k^*
\]
whenever \(T\subset R_1\otimes R_2\otimes R_3\).  Iterating the Casimir
fusion estimate \eqref{eq:V30-Casimir-fusion} and using
\(C_2(R^*)=C_2(R)\) yields
\begin{equation}
 1+C_2(R_i)
 \le C\bigl(1+C_2(T)+C_2(R_j)+C_2(R_k)\bigr)
 \label{eq:V32-linkwise-reverse-fusion}
\end{equation}
on every link active in \(R_i\).  Such a link must be active in at least
one channel on the right; otherwise the inclusion would put a nontrivial
representation inside the trivial representation.  Summing
\eqref{eq:V32-linkwise-reverse-fusion} over links and charging its unit
term to that active channel proves \eqref{eq:V32-reverse-fusion}.
Applying the same argument to
\(S_{ij}\subset R_i\otimes R_j\) and
\(T\subset S_{ij}\otimes R_k\) proves the intermediate-channel claim.
\end{proof}

\begin{corollary}[Paired escape in a bounded final channel]
\label{cor:V32-paired-high-to-low}
Along the witnesses of
\Cref{thm:V32-ternary-thermal-escape}, suppose
\(\Theta_{\alpha_r}(\boldsymbol T_r)\) remains bounded.  If one input
thermal ratio tends to infinity, then at least one other input thermal
ratio tends to infinity.  Hence high-to-low fusion into a bounded final
channel is necessarily a paired cancellation; a single isolated high
input cannot cause it.
\end{corollary}

\begin{proof}
Multiply \eqref{eq:V32-reverse-fusion} by \(s_{\alpha_r}\).  A bounded
final ratio and two bounded input ratios would force the remaining input
ratio to be bounded as well.
\end{proof}

\begin{assessment}[The exact tail branch left by V32]
\label{ass:V32-tail-branch}
V32 removes arbitrary spatial size from the ternary problem on every
fixed thermal core.  Every unresolved sequence has a high input.  If its
final channel is also high, each five-term bracketing must be reorganized
so that the quadratic debit \eqref{eq:V32-quadratic-thermal-damping} is
attached before taking trace-norm absolute values.  If the final channel
stays bounded, \Cref{cor:V32-paired-high-to-low} forces at least two high
inputs, and their large weights cancel through a binary intermediate or
directly into the final channel.  The latter high-to-low branch is the
sharpest remaining obstruction because no high final multiplier is
available to pay the resolvent.  Neither branch is declared closed here:
the necessary Racah-compatible placement of the debit is the V33 task.
\end{assessment}

% =============================================================================
\section{V32 integrated result and next frontier}
\label{sec:V32-ledger}
% =============================================================================

\begin{theorem}[Integrated V32 statement]
\label{thm:V32-integrated}
Preserving V1--V31, V32 proves:
\begin{enumerate}[label=\textup{(V32.\arabic*)},leftmargin=5.5em]
\item the bi-Laplacian density estimate
      \eqref{eq:V32-density-bilaplacian}, the one-link quadratic debit
      \eqref{eq:V32-one-link-quadratic-debit}, and the support-uniform
      product second moment \eqref{eq:V32-balanced-second-moment};
\item the representation-uniform one-link third-cumulant estimate
      \eqref{eq:V32-one-link-third-cumulant};
\item the exact replicated-coordinate ternary tree inequality
      \eqref{eq:V32-replicated-tree}, separating two pair bonds from a
      same-link third-order junction;
\item the Wilson two-mark numerator
      \eqref{eq:V32-Wilson-numerator}, before the output denominator is
      charged;
\item the resolved two-mark quotient
      \eqref{eq:V32-thermal-two-mark-quotient} and the full ternary MTA
      estimate \eqref{eq:V32-bounded-thermal-MTA} on each fixed thermal
      window, uniformly in spatial support size; and
\item the thermal escape theorem and reverse-fusion constraint
      \eqref{eq:V32-input-thermal-escape}--%
      \eqref{eq:V32-reverse-fusion}, which force any bounded-output tail
      obstruction to contain at least two high input channels.
\end{enumerate}
\end{theorem}

\begin{proof}
Item \textup{(V32.1)} is
\Cref{lem:V32-density-bilaplacian,thm:V32-quadratic-tail-debit}.
Item \textup{(V32.2)} is
\Cref{thm:V32-one-link-third-cumulant}.  Item \textup{(V32.3)} is
\Cref{lem:V32-replicated-tree}.  Item \textup{(V32.4)} is
\Cref{prop:V32-Wilson-numerator}.  Item \textup{(V32.5)} is
\Cref{lem:V32-thermal-two-mark-quotient,thm:V32-bounded-thermal-MTA}.
Item \textup{(V32.6)} is
\Cref{thm:V32-ternary-thermal-escape,lem:V32-reverse-fusion,
cor:V32-paired-high-to-low}.
\end{proof}

\begingroup
\small
\begin{longtable}{>{\raggedright\arraybackslash}p{0.29\textwidth}
                  >{\raggedright\arraybackslash}p{0.15\textwidth}
                  >{\raggedright\arraybackslash}p{0.48\textwidth}}
\toprule
\textbf{V32 row}&\textbf{Status}&\textbf{Advance or remaining obstruction}\\
\midrule
\endfirsthead
\toprule
\textbf{V32 row}&\textbf{Status}&\textbf{Advance or remaining obstruction}\\
\midrule
\endhead
Quadratic Wilson tail debit
& \MGclosed
& Two Casimir powers times the product multiplier are uniformly bounded,
  independently of support size.\\[0.35em]
One-link third cumulant
& \MGclosed
& Inversion symmetry removes the cubic jet and leaves the fourth-moment
  scale \(s_\alpha^2\).\\[0.35em]
Replicated-coordinate ternary numerator
& \MGclosed
& The two tree bonds and the same-link junction are separated without
  iterating the binary cumulant or spending an output resolvent.\\[0.35em]
Ternary MTA on \(s_\alpha\Xi\le L\)
& \MGclosed
& Every fixed thermal window is uniform in support, volume, and all
  representations inside the window; its constant depends on \(L\).\\[0.35em]
Uniform ternary MTA
& \MGopen
& Removing the \(L\)-dependence requires a Racah-compatible quadratic
  debit in the high-channel branches.\\[0.35em]
All-order marked-tree atlas
& \MGopen
& Higher forests remain downstream of uniform ternary tail closure.\\[0.35em]
Connected Perron majorant
& \MGreduced
& It follows from all-order MTA through the V31 and V30 compilers.\\[0.35em]
V24/V25 unreduced edge packet
& \MGopen
& The direct singular-channel calculation remains downstream of WUFC.\\[0.35em]
Capacity, protected sectors, and OS limit
& \MGopen
& No new capacity or continuum claim is made in V32.\\[0.35em]
Full Yang--Mills mass gap
& \MGopen
& Uniform high-channel ternary fusion, higher marked forests, finite
  edges, and continuum closure remain.\\
\bottomrule
\end{longtable}
\endgroup

\begin{openproblem}[V33 mandate: the high-to-low Racah atlas]
\label{op:V33}
\begin{enumerate}[label=\textup{(V33.\arabic*)},leftmargin=5.5em]
\item Write the five resolved ternary multipliers from
      \eqref{eq:V31-ternary-assembled} in one final-channel bracketing.
      Retain every Clebsch--Gordan multiplicity space and every Racah
      change of basis before taking a trace norm.
\item Split by the thermal ratios of the three inputs, each binary
      intermediate, and the final output.  Apply
      \eqref{eq:V32-quadratic-thermal-damping} only to a multiplier that
      actually occurs in that term; do not transfer damping across a
      recoupling identity without proof.
\item In the bounded-output branch, use
      \Cref{cor:V32-paired-high-to-low} to isolate the two high inputs and
      prove a two-mark high-to-low cancellation estimate uniform in their
      highest weights.  For \(G=\mathrm{SU}(3)\), begin with the explicit
      multiplicity-resolved Racah blocks.
\item Prove a constant independent of \(L\) in
      \eqref{eq:V32-thermal-two-mark-quotient}, or return a finite
      resolved highest-weight, support, intermediate-channel, final-channel,
      and two-link-mark witness.
\item Only after uniform ternary closure, extend the independent-coordinate
      expansion to general connected hypergraphs and extract a labeled
      spanning tree with one normalized mark per incident tree edge.
      Keep the per-edge constant exponential before the Cayley tree sum.
\item On strict all-order MTA closure, invoke
      \Cref{thm:V31-MTA-compiler,thm:V30-WUFC-compiler}; then populate the
      full nonlinear Harnack, signed-density, incidence, and temporal-source
      packet and attack the unreduced V24/V25 edge matrix directly.
\item Continue through bad capacity, protected sectors, local-algebra
      noncollapse, reflection positivity, and the nontrivial OS continuum
      state.  The next scheduled SM/NS companion audit remains V35.
\end{enumerate}
\end{openproblem}

\begin{firewall}[End-of-V32 claim boundary]
V32 proves the assembled ternary marked-tree inequality uniformly on
every fixed thermal-core window and proves that any remaining ternary
failure must escape those windows.  It does not prove a constant uniform
as \(L\to\infty\), the all-order marked-tree atlas, WUFC, CPM, the
nonlinear Perron packet, the unreduced edge margins, bad capacity,
protected-sector floors, local-algebra noncollapse, or a nontrivial
continuum OS limit.  The full Yang--Mills mass gap is not proved.
\end{firewall}

\section*{Version V32 append-only record}
\addcontentsline{toc}{section}{Version V32 append-only record}
The complete V31 mathematical source was retained before this continuation.
Its SHA--256 digest was
\begin{center}\scriptsize\ttfamily
8cb90753049d561c2ec36c93a5fcc69ba2e5cfad26290f91db8da52b4d34d243.
\end{center}
No mathematical content of V31 was altered.  On the next iteration, retain
this full file, append before its unique active document-closing command,
increment the filename to \texttt{\_V33.tex}, and remove only the
superseded V32 file from this note series.  The next companion audit is
V35.

% =============================================================================
% END APPENDED VERSION V32 MATERIAL
% =============================================================================
% =============================================================================
% BEGIN APPENDED VERSION V33 MATERIAL
% =============================================================================

\clearpage
\part{V33 --- Common-basis ternary symbols, damped Racah leakage, and
singlet-formation marks}
\label{part:V33}

\section{Purpose and exact scope}
\label{sec:V33-purpose}

V32 proved ternary MTA on each fixed thermal window and showed that every
remaining witness must contain a thermally escaping input.  It left open
whether a high multiplier can be retained after the three pair bracketings
are placed in one final-channel basis.  V33 answers that structural
question exactly.

The assembled five-term multiplier is written below as one operator on the
threefold representation carrier.  In any chosen pair bracketing, its
undamped part is diagonal in that pair intermediate and is precisely a
binary fusion defect.  Every off-diagonal Racah block carries one of the
two complementary input multipliers.  V32 quadratic damping therefore
controls the leakage of a paired high-to-low channel without pretending
that the three pair moment operators commute.

The diagonal formation of the chosen pair still has to supply the first
spatial mark.  V33 proves that tax completely for a linkwise dual pair
forming the product singlet, using Schur orthogonality and the exact
\(\mathrm{SU}(3)\) dimension formula.  A nontrivial cold intermediate is
left as an explicit multiplicity-resolved formation inequality.  Thus V33
does not yet claim a thermal-window-independent ternary MTA constant.
There is no companion audit at V33; the next scheduled SM/NS audit remains
V35.

% =============================================================================
\section{Moment operators on one common representation carrier}
\label{sec:V33-moment-operators}
% =============================================================================

Fix three exact product representations \(\boldsymbol R_i\), extending
each by the trivial representation off its support, and put
\begin{equation}
 \mathscr H:=\bigotimes_{e}
 \left(H_{R_{1,e}}\otimes H_{R_{2,e}}\otimes H_{R_{3,e}}\right).
 \label{eq:V33-threefold-carrier}
\end{equation}
For a nonempty set \(B\subseteq\{1,2,3\}\), decompose on link \(e\)
the diagonal \(G\)-representation \(\bigotimes_{i\in B}R_{i,e}\) and
define
\begin{equation}
 Q_{B,e}:=sum_{S_e\subseteq\bigotimes_{i\in B}R_{i,e}}
 \eta_{S_e}(\alpha)P_{S_e}^{B,e},
 \qquad
 \mathbf Q_B:=\bigotimes_e Q_{B,e}.
 \label{eq:V33-moment-operators}
\end{equation}
The identity is understood on tensor factors whose labels are outside
\(B\).  Multiplicity copies are retained separately in the spectral sum.
Every \(Q_{B,e}\), and hence every \(\mathbf Q_B\), is a positive
contraction.  For a singleton,
\begin{equation}
 \mathbf Q_{\{i\}}=q_i I,
 \qquad q_i:=q_{\alpha,\boldsymbol R_i}.
 \label{eq:V33-singleton-multiplier}
\end{equation}

Let \(P_{\boldsymbol T}\) be the total diagonal product-group isotypic
projection for a nonconstant final channel \(\boldsymbol T\).  Each pair
operator commutes with the total diagonal action, so on that isotypic
space it has the form
\begin{equation}
 P_{\boldsymbol T}\mathbf Q_{ij}P_{\boldsymbol T}
 =I_{\boldsymbol T}\otimes Q_{ij}^{\boldsymbol T},
 \label{eq:V33-pair-multiplicity-operator}
\end{equation}
where \(Q_{ij}^{\boldsymbol T}\) acts on the full multiplicity space.
In the \((ij)k\) bracketing,
\begin{equation}
 Q_{ij}^{\boldsymbol T}
 =\sum_{\boldsymbol S}
 q_{\boldsymbol S}P_{\boldsymbol S\mid\boldsymbol T}^{ij},
 \qquad
 q_{\boldsymbol S}:=q_{\alpha,\boldsymbol S}.
 \label{eq:V33-pair-spectral-resolution}
\end{equation}
The other two pair operators are conjugates of their diagonal forms by
the exact product of linkwise Racah unitaries.  No simultaneous
diagonalization is assumed.

\begin{theorem}[Exact common-basis ternary multiplier]
\label{thm:V33-common-basis-symbol}
On \(\mathscr H\), the assembled ternary row cumulant has multiplier
operator
\begin{equation}
 \boxed{
 \mathbf K_{123}
 =\mathbf Q_{123}
 -q_3\mathbf Q_{12}-q_2\mathbf Q_{13}-q_1\mathbf Q_{23}
 +2q_1q_2q_3I.}
 \label{eq:V33-common-basis-symbol}
\end{equation}
On the final channel \(\boldsymbol T\), write
\(K_{123}^{\boldsymbol T}\) for its multiplicity-space compression.
For every permutation \(\{i,j,k\}=\{1,2,3\}\), put
\begin{equation}
 D_{ij}:=Q_{ij}^{\boldsymbol T}-q_iq_jI.
 \label{eq:V33-pair-defect-operator}
\end{equation}
Then the following anchored identity is exact:
\begin{equation}
 \boxed{
 K_{123}^{\boldsymbol T}
 =q_{\boldsymbol T}I-q_kQ_{ij}^{\boldsymbol T}
  -q_jD_{ik}-q_iD_{jk}.}
 \label{eq:V33-anchored-symbol}
\end{equation}
It contains the single final resolvent denominator only after
\(K_{123}^{\boldsymbol T}\) has been formed.
\end{theorem}

\begin{proof}
For pure Peter--Weyl inputs, convolution of a product with index set
\(B\) acts on its tensor coefficient by \(\mathbf Q_B\).  The five
partitions of \(\{1,2,3\}\) in
\eqref{eq:V29-cumulant-Mobius} therefore give, respectively,
\[
 \mathbf Q_{123},\qquad
 -q_3\mathbf Q_{12},\quad-q_2\mathbf Q_{13},\quad
 -q_1\mathbf Q_{23},\qquad 2q_1q_2q_3I.
\]
They already act on the same carrier \(\mathscr H\); changing a pair
bracketing merely conjugates the corresponding multiplicity operator by
a Racah unitary.  This proves \eqref{eq:V33-common-basis-symbol}.

On the final isotypic space,
\(\mathbf Q_{123}=q_{\boldsymbol T}I\).  Substitute
\(Q_{ik}^{\boldsymbol T}=D_{ik}+q_iq_kI\) and
\(Q_{jk}^{\boldsymbol T}=D_{jk}+q_jq_kI\) into
\eqref{eq:V33-common-basis-symbol}.  The two resulting scalar terms are
exactly \(2q_iq_jq_kI\), proving
\eqref{eq:V33-anchored-symbol}.  The resolvent statement follows from
\eqref{eq:V29-product-resolvent}.
\end{proof}

\begin{corollary}[Diagonal binary defect and off-diagonal Racah leakage]
\label{cor:V33-block-decomposition}
Fix the \((ij)k\) bracketing and two pair-intermediate product channels
\(\boldsymbol S,\boldsymbol S'\) leading to \(\boldsymbol T\).  Then
\begin{align}
 &P_{\boldsymbol S\mid\boldsymbol T}^{ij}
 K_{123}^{\boldsymbol T}
 P_{\boldsymbol S'\mid\boldsymbol T}^{ij}
 \notag\\
 &\quad=
 \mathbf1_{\{\boldsymbol S=\boldsymbol S'\}}
 (q_{\boldsymbol T}-q_kq_{\boldsymbol S})
 P_{\boldsymbol S\mid\boldsymbol T}^{ij}
 +P_{\boldsymbol S\mid\boldsymbol T}^{ij}
 E_{ij;k}^{\boldsymbol T}
 P_{\boldsymbol S'\mid\boldsymbol T}^{ij},
 \label{eq:V33-block-decomposition}
\end{align}
where
\begin{equation}
 E_{ij;k}^{\boldsymbol T}:=-q_jD_{ik}-q_iD_{jk},
 \qquad
 \norm{E_{ij;k}^{\boldsymbol T}}_{\rm op}\le q_i+q_j.
 \label{eq:V33-Racah-remainder}
\end{equation}
Thus every off-diagonal intermediate-channel block carries an actual
complementary input multiplier; the undamped part is diagonal and is the
binary defect for \(\boldsymbol S\otimes\boldsymbol R_k\to\boldsymbol T\).
\end{corollary}

\begin{proof}
Equation \eqref{eq:V33-pair-spectral-resolution} diagonalizes the first
two terms of \eqref{eq:V33-anchored-symbol}.  Since
\(Q_{ik}^{\boldsymbol T}\) is a positive contraction and
\(q_iq_k\in[0,1]\), its self-adjoint defect \(D_{ik}\) has operator norm
at most one; the same holds for \(D_{jk}\).  The triangle inequality
proves \eqref{eq:V33-Racah-remainder} and hence
\eqref{eq:V33-block-decomposition}.
\end{proof}

\begin{theorem}[Quadratically damped paired high-to-low reduction]
\label{thm:V33-high-pair-reduction}
Let \(\Theta_i=\Theta_\alpha(\boldsymbol R_i)\).  In any final channel
and any \((ij)k\) bracketing,
\begin{equation}
 \boxed{
 \norm{K_{123}^{\boldsymbol T}
 -(q_{\boldsymbol T}I-q_kQ_{ij}^{\boldsymbol T})}_{\rm op}
 \le K_{\rm hp}\left(
 \min\{1,\Theta_i^{-2}\}+
 \min\{1,\Theta_j^{-2}\}\right).}
 \label{eq:V33-high-pair-reduction}
\end{equation}
Here the expression \(\Theta^{-2}\) is used only for a nontrivial exact
channel; the minimum is interpreted as one otherwise.  In particular,
if \(\Theta_i,\Theta_j\to\infty\), every off-diagonal Racah block tends
to zero quadratically, uniformly in all multiplicities and support size.
\end{theorem}

\begin{proof}
The left side is \(\norm{E_{ij;k}^{\boldsymbol T}}_{\rm op}\) by
\eqref{eq:V33-anchored-symbol}.  Apply
\eqref{eq:V33-Racah-remainder} and the quadratic product damping
\eqref{eq:V32-quadratic-thermal-damping} to \(q_i\) and \(q_j\).
Operator norm is taken before any sum over intermediate multiplicities,
so no multiplicity factor is introduced.
\end{proof}

\begin{corollary}[Exact spectator-singlet cancellation]
\label{cor:V33-spectator-singlet}
Suppose \(\boldsymbol R_j=\boldsymbol R_i^*\), the chosen pair
intermediate is \(\boldsymbol S=\boldsymbol1\), and the final channel is
the spectator \(\boldsymbol T=\boldsymbol R_k\).  Then
\begin{equation}
 (q_{\boldsymbol T}-q_kq_{\boldsymbol S})
 P_{\boldsymbol1\mid\boldsymbol R_k}^{ij}=0,
 \label{eq:V33-spectator-singlet-cancellation}
\end{equation}
and both the diagonal singlet compression and every block entering or
leaving it satisfy the remainder bound
\begin{equation}
 \norm{P_{\boldsymbol S\mid\boldsymbol T}^{ij}
 K_{123}^{\boldsymbol T}
 P_{\boldsymbol S'\mid\boldsymbol T}^{ij}}_{\rm op}
 \le q_i+q_j
 \label{eq:V33-singlet-leakage-bound}
\end{equation}
whenever \(\boldsymbol S=\boldsymbol1\) or
\(\boldsymbol S'=\boldsymbol1\), except for a diagonal non-singlet
binary-defect term.  A high dual pair therefore cannot leave an undamped
spectator-singlet contribution.
\end{corollary}

\begin{proof}
The trivial intermediate has \(q_{\boldsymbol S}=1\), while
\(q_{\boldsymbol T}=q_k\).  This proves
\eqref{eq:V33-spectator-singlet-cancellation}.  Apply
\Cref{cor:V33-block-decomposition} to the indicated blocks.
\end{proof}

% =============================================================================
\section{The first noncommuting \(\mathrm{SU}(3)\) Racah block}
\label{sec:V33-explicit-Racah}
% =============================================================================

The common-basis statement is not a cosmetic change.  The first
fundamental multiplicity space already contains noncommuting pair moment
operators.

\begin{proposition}[Exact \(3\otimes\bar3\otimes3\to3\) Racah matrix]
\label{prop:V33-fundamental-Racah}
Let \(V=\mathbb C^N\) be the fundamental representation of
\(\mathrm{SU}(N)\).  On the two-dimensional multiplicity space
\begin{equation}
 \operatorname{Hom}_{\mathrm{SU}(N)}
 (V\otimes V^*\otimes V,V),
 \label{eq:V33-fundamental-multiplicity}
\end{equation}
the change from the \((12)3\) singlet/adjoint basis to the
\(1(23)\) singlet/adjoint basis is, up to independent channel phases,
\begin{equation}
 U_N=
 \begin{pmatrix}
 N^{-1}&\sqrt{N^2-1}\,N^{-1}\\
 \sqrt{N^2-1}\,N^{-1}&-N^{-1}
 \end{pmatrix}.
 \label{eq:V33-SUN-F-matrix}
\end{equation}
For \(N=3\), put \(b_\alpha=\eta_{8}(\alpha)\).  In the
\((12)3\) basis,
\begin{equation}
 Q_{12}^{3}=\begin{pmatrix}1&0\\0&b_\alpha\end{pmatrix},
 \qquad
 Q_{23}^{3}=U_3
 \begin{pmatrix}1&0\\0&b_\alpha\end{pmatrix}U_3^*.
 \label{eq:V33-fundamental-pair-moments}
\end{equation}
Consequently,
\begin{equation}
 \boxed{
 \norm{[Q_{12}^{3},Q_{23}^{3}]}_{\rm op}
 =\frac{\sqrt8}{9}(1-b_\alpha)^2>0.}
 \label{eq:V33-pair-commutator}
\end{equation}
Thus pair multipliers cannot be treated as jointly diagonal even in the
smallest nontrivial \(\mathrm{SU}(3)\) ternary multiplicity block.
\end{proposition}

\begin{proof}
There are two invariant contraction maps
\begin{align*}
 A(v\otimes\varphi\otimes w)&:=\varphi(v)w,
 &B(v\otimes\varphi\otimes w)&:=\varphi(w)v.
\end{align*}
Their Hilbert--Schmidt norms equal \(N\), while
\(\langle A/N,B/N\rangle=N^{-1}\).  The normalized \((12)3\) basis is
\[
 e_0=A/N,
 \qquad
 e_1=\frac{B/N-N^{-1}e_0}{\sqrt{1-N^{-2}}},
\]
and the normalized \(1(23)\) basis is obtained by interchanging
\(A\) and \(B\).  Taking their inner products gives
\eqref{eq:V33-SUN-F-matrix}.  Equivalently, this is the fundamental
Fierz identity
\[
 \sum_a(t^a)_i{}^j(t^a)_k{}^\ell
 =\frac12\left(\delta_i^\ell\delta_k^j
 -\frac1N\delta_i^j\delta_k^\ell\right)
\]
with orthonormal generators.

For \(N=3\), both pair decompositions are \(1\oplus8\), proving
\eqref{eq:V33-fundamental-pair-moments}.  Its off-diagonal entry is
\(\sqrt8(1-b_\alpha)/9\).  Commutation with
\(\operatorname{diag}(1,b_\alpha)\) multiplies that entry by
\(1-b_\alpha\); the resulting skew-Hermitian \(2\)-by-\(2\) matrix has
operator norm equal to the displayed modulus.  The strict V30 character
gap gives \(b_\alpha<1\) for the nontrivial adjoint representation.
\end{proof}

\begin{proposition}[Casimir Ward cancellation in the common basis]
\label{prop:V33-Casimir-Ward}
Let \(C_i\), \(C_{ij}\), and \(C_{123}\) denote the quadratic Casimir
operators for the indicated diagonal tensor actions on one link.  Then
\begin{equation}
 \boxed{
 C_{123}=C_{12}+C_{13}+C_{23}-C_1-C_2-C_3.}
 \label{eq:V33-Casimir-Ward}
\end{equation}
Consequently, on every fixed representation bank, the constant and
first-order terms in the small-\(s_\alpha\) expansion of
\eqref{eq:V33-common-basis-symbol} cancel as operators on the entire
multiplicity space, not merely on its diagonal entries.
\end{proposition}

\begin{proof}
If \(T_i^a\) are the infinitesimal generators on the three factors, then
\[
 C_{ij}=C_i+C_j+2\sum_aT_i^aT_j^a,
 \qquad
 C_{123}=C_1+C_2+C_3+2\sum_{a,i<j}T_i^aT_j^a.
\]
Subtracting proves \eqref{eq:V33-Casimir-Ward}.  On a fixed bank,
\Cref{thm:V10-fixed-symbol-limit} and
\eqref{eq:V10-eta-asymptotic} give
\(\eta_R=1-s_\alpha C_2(R)+O_R(s_\alpha^2)\).
Insert this expansion into \eqref{eq:V33-common-basis-symbol}.  The
constant coefficient is \(1-3+2=0\).  The coefficient of
\(-s_\alpha\) is precisely
\[
 C_{123}-C_{12}-C_{13}-C_{23}+C_1+C_2+C_3,
\]
which vanishes by \eqref{eq:V33-Casimir-Ward}.  Since the identity holds
before choosing a bracketing, it also cancels all first-order off-diagonal
Racah entries.  This is the common-basis form of the one-link Gaussian
excess proved uniformly in \Cref{thm:V32-one-link-third-cumulant}.
\end{proof}

% =============================================================================
\section{The product-singlet formation tax is an actual pair mark}
\label{sec:V33-singlet-formation}
% =============================================================================

Let \(\boldsymbol R=(R_e)_{e\in X}\) be exact and nontrivial on every
link of \(X\), and let \(\boldsymbol R^*=(R_e^*)_{e\in X}\).  Write
\begin{equation}
 d_{\boldsymbol R}:=\prod_{e\in X}d_{R_e},
 \qquad
 \vartheta_e:=\vartheta_e(\boldsymbol R)
 =\vartheta_e(\boldsymbol R^*).
 \label{eq:V33-dual-dimension-and-marks}
\end{equation}

\begin{lemma}[Schur tax for linkwise singlet formation]
\label{lem:V33-Schur-formation-tax}
With the Fourier normalization of \eqref{eq:V29-Fourier-algebra-norm},
the product-singlet contraction of coefficient matrices \(A\) in
\(\boldsymbol R\) and \(B\) in \(\boldsymbol R^*\) obeys
\begin{equation}
 \left\lvert\mathsf S_{\boldsymbol R}(A,B)\right\rvert
 \le d_{\boldsymbol R}^{-1}
 \bigl(d_{\boldsymbol R}\norm{A}_1\bigr)
 \bigl(d_{\boldsymbol R}\norm{B}_1\bigr).
 \label{eq:V33-Schur-tax}
\end{equation}
For \(G=\mathrm{SU}(3)\),
\begin{equation}
 \boxed{
 d_{\boldsymbol R}^{-1}
 \le\sum_{e\in X}
 \vartheta_e(\boldsymbol R)
 \vartheta_e(\boldsymbol R^*).}
 \label{eq:V33-singlet-tax-is-mark}
\end{equation}
Hence linkwise formation of the product singlet supplies one normalized
\(12\) overlap mark, uniformly in support size and representation height.
\end{lemma}

\begin{proof}
For a one-link irreducible, Schur orthogonality says
\begin{equation}
 P_{\mathbf1}
 \left(D^R_{ab}D^{R^*}_{cd}\right)
 =\frac1{d_R}\delta_{ac}\delta_{bd}
 \label{eq:V33-Schur-orthogonality}
\end{equation}
after identifying the dual indices in the standard way.  Trace duality
therefore bounds the contraction of arbitrary coefficient matrices by
\(d_R^{-1}(d_R\norm{A}_1)(d_R\norm{B}_1)\).  Tensoring
\eqref{eq:V33-Schur-orthogonality} over the links proves
\eqref{eq:V33-Schur-tax}.

Put \(N=|X|\).  Every nontrivial \(\mathrm{SU}(3)\) irreducible has
dimension at least three, so
\(d_{\boldsymbol R}^{-1}\le3^{-N}\le N^{-1}\).  On the other hand,
\(\sum_e\vartheta_e=1\), and Cauchy--Schwarz gives
\[
 \sum_{e\in X}\vartheta_e^2\ge\frac1N.
\]
Dual representations have the same Casimir and hence the same marks.
This proves \eqref{eq:V33-singlet-tax-is-mark}.  For pure BFA atoms, the
sum over links of the products of the two one-mark input weights is
\[
 e^{2\rho\operatorname{ht}(\boldsymbol R)}
 \Xi(\boldsymbol R)^2d_{\boldsymbol R}^2
 \norm{A}_1\norm{B}_1\sum_e\vartheta_e^2.
\]
It dominates \eqref{eq:V33-Schur-tax}, because the exponential and
balanced-mass factors are at least one.  Hence the dimension estimate is
indeed a marked Fourier-algebra estimate, rather than only a scalar
matrix-coefficient observation.
\end{proof}

\begin{theorem}[The spectator-singlet branch has no unmarked formation]
\label{thm:V33-singlet-branch-structure}
In the setting of \Cref{cor:V33-spectator-singlet}, the diagonal
spectator-singlet binary defect vanishes exactly, its formation map pays
the pair mark \eqref{eq:V33-singlet-tax-is-mark}, and every remaining
Racah block is bounded by
\begin{equation}
 K_{\rm hp}\left(
 \min\{1,\Theta_i^{-2}\}+
 \min\{1,\Theta_j^{-2}\}\right).
 \label{eq:V33-singlet-branch-damping}
\end{equation}
Thus a linkwise dual product singlet is not the source of an undamped,
unmarked high-to-low term.  Completing its full two-mark MTA bound still
requires assigning the second mark in the recoupled pair defect; that
assignment is not asserted here.
\end{theorem}

\begin{proof}
The exact diagonal cancellation is
\eqref{eq:V33-spectator-singlet-cancellation}.  Apply
\Cref{lem:V33-Schur-formation-tax} to the singlet projection and
\Cref{thm:V33-high-pair-reduction} to the remainder.  These operations
are performed before trace-norm block summation, so their constants do
not depend on the multiplicity dimension.  The final sentence records
the still-unproved spatial placement and prevents use of the theorem as
a full MTA assertion.
\end{proof}

\begin{proposition}[Racah unitarity alone cannot manufacture a mark]
\label{prop:V33-unitarity-no-mark}
No estimate based only on trace-norm invariance under a Racah unitary can
produce a normalized pair-incidence factor.  More precisely, for every
\(N\) there are probability marks \(\vartheta_e=N^{-1}\), a unitary
\(U_N\), and a nonzero trace-class matrix \(A_N\) such that
\begin{equation}
 \norm{U_NA_NU_N^*}_1=\norm{A_N}_1,
 \qquad
 \sum_{e=1}^N\vartheta_e^2\norm{A_N}_1
 =\frac1N\norm{A_N}_1.
 \label{eq:V33-unitarity-no-mark}
\end{equation}
Therefore the missing nontrivial-intermediate estimate must use the
actual Littlewood--Richardson/Schur formation coefficients or an actual
Wilson pair defect, not merely unitarity of the recoupling matrix.
\end{proposition}

\begin{proof}
Take any unitary \(U_N\), any rank-one matrix of trace norm one, and
uniform probability marks.  Equation \eqref{eq:V33-unitarity-no-mark} is
immediate.  A uniform inequality replacing the first member by a constant
times the second would require that constant to be at least \(N\).
\end{proof}
% =============================================================================
\section{Private-link factorization and its quadratic payment}
\label{sec:V33-private-factor}
% =============================================================================

For \(\{i,j,k\}=\{1,2,3\}\), define the links private to input \(i\) by
\begin{equation}
 P_i:=X_i\setminus(X_j\cup X_k),
 \qquad
 P:=P_1\mathbin{\dot\cup}P_2\mathbin{\dot\cup}P_3.
 \label{eq:V33-private-links}
\end{equation}
Put
\begin{equation}
 \Xi_P:=\sum_{i=1}^3\sum_{e\in P_i}A_{i,e},
 \qquad
 q_P:=\prod_{i=1}^3\prod_{e\in P_i}\eta_{R_{i,e}}.
 \label{eq:V33-private-mass-multiplier}
\end{equation}

\begin{lemma}[Exact private-link factorization]
\label{lem:V33-private-factorization}
Every private link has final representation \(T_e=R_{i,e}\) for its
unique active input, and the full common-basis ternary symbol factors as
\begin{equation}
 \boxed{
 K_{123}^{\boldsymbol T}=q_P K_{123}^{\boldsymbol T,\mathrm{sh}},}
 \label{eq:V33-private-factorization}
\end{equation}
where the shared symbol contains only links active in at least two inputs.
Moreover,
\begin{equation}
 (s_\alpha\Xi_P)^2q_P\le K_{\rm priv,2}
 \label{eq:V33-private-quadratic-debit}
\end{equation}
when \(P\ne\varnothing\).
\end{lemma}

\begin{proof}
At \(e\in P_i\), every block of every partition in the cumulant contains
the \(i\)-th input exactly once and the other inputs trivially.  Convolution
therefore contributes the same scalar \(\eta_{R_{i,e}}\) to all five
terms, and tensor multiplication cannot change the final representation
there.  Factor these scalars link by link before taking the Möbius sum;
this proves \eqref{eq:V33-private-factorization}.  The private
representations form one exact product channel on the disjoint union
\(P\), so \eqref{eq:V33-private-quadratic-debit} is
\eqref{eq:V32-balanced-second-moment} applied to that channel.
\end{proof}

\begin{proposition}[Private-dominant tail payment]
\label{prop:V33-private-tail-payment}
Fix \(D\ge1\) and a tree center \(i\).  If
\begin{equation}
 \Xi(\boldsymbol T)\le D\Xi_P,
 \qquad
 \Xi(\boldsymbol R_i)\le D\Xi_P,
 \label{eq:V33-private-dominance}
\end{equation}
then
\begin{equation}
 \boxed{
 \frac{\Xi(\boldsymbol T)}{1-q_{\alpha,\boldsymbol T}}
 q_Ps_\alpha^2
 \le\frac{C_D}{\Xi(\boldsymbol R_i)}.}
 \label{eq:V33-private-tail-payment}
\end{equation}
Consequently, the V32 pair-tree numerator centered at \(i\) obeys
\begin{equation}
 \frac{\Xi(\boldsymbol T)}{1-q_{\alpha,\boldsymbol T}}
 q_Ps_\alpha^2H_{ij}H_{ik}
 \le C_D\frac{H_{ij}H_{ik}}{\Xi(\boldsymbol R_i)},
 \label{eq:V33-private-tree-payment}
\end{equation}
which is exactly its two-mark mass.  This estimate has no thermal-window
constant.
\end{proposition}

\begin{proof}
Write \(\Xi_T=\Xi(\boldsymbol T)\) and
\(\Xi_i=\Xi(\boldsymbol R_i)\).  The two-sided product gap and exact
support Casimir floor imply
\begin{equation}
 \frac{\Xi_T}{1-q_{\alpha,\boldsymbol T}}
 \le C\begin{cases}
 s_\alpha^{-1},&s_\alpha\Xi_T\le c_0,\\
 \Xi_T,&s_\alpha\Xi_T>c_0,
 \end{cases}
 \label{eq:V33-output-resolvent-split}
\end{equation}
for fixed group constants \(C,c_0>0\).

In the first case, \(\Xi_P\le\Xi_T\) because every private input
representation is unchanged in the final channel.  Hence
\(s_\alpha\Xi_P\le c_0\), and
\eqref{eq:V33-private-dominance} gives
\(s_\alpha\le Dc_0/\Xi_i\).  The left side of
\eqref{eq:V33-private-tail-payment} is at most \(Cs_\alpha q_P\), so the
claim follows from \(q_P\le1\).

In the second case, use \(\Xi_T\le D\Xi_P\) and
\eqref{eq:V33-private-quadratic-debit}:
\[
 \Xi_Tq_Ps_\alpha^2
 \le D\Xi_Pq_Ps_\alpha^2
 =\frac D{\Xi_P}(s_\alpha\Xi_P)^2q_P
 \le\frac{C_D}{\Xi_P}
 \le\frac{C_D'}{\Xi_i}.
\]
This proves \eqref{eq:V33-private-tail-payment}; multiplying by
\(H_{ij}H_{ik}\) proves \eqref{eq:V33-private-tree-payment}.
\end{proof}

\begin{corollary}[Where a remaining tail witness must live]
\label{cor:V33-shared-tail-route}
A ternary witness not controlled by
\Cref{thm:V32-bounded-thermal-MTA,prop:V33-private-tail-payment} must,
for each unresolved tree center, violate at least one inequality in
\eqref{eq:V33-private-dominance}.  Thus either its final mass is not
private-dominated, or the center mass is concentrated on links shared
with another input.  Together with
\Cref{thm:V33-high-pair-reduction}, the remaining bounded-output branch
is a nontrivial shared-link formation problem, not a private-support or
undamped Racah-leakage problem.
\end{corollary}

\begin{proof}
The bounded thermal core is
\Cref{thm:V32-bounded-thermal-MTA}.  Outside it, whenever both
private-dominance inequalities hold, the coarse V32 tree numerator is
converted into the required normalized tree mass by
\eqref{eq:V33-private-tree-payment}.  The contrapositive gives the first
claim.  The high-pair Racah remainder is quadratically damped by
\eqref{eq:V33-high-pair-reduction}, giving the stated routing.
\end{proof}

% =============================================================================
\section{The residual cold-intermediate formation card}
\label{sec:V33-formation-card}
% =============================================================================

Choose the \((ij)k\) bracketing.  For a nontrivial intermediate
\(\boldsymbol S\), let
\(\mathsf F_{ij}^{\boldsymbol S}\) denote product of the two input
Peter--Weyl tensors followed by the exact
\(\boldsymbol S\)-isotypic Clebsch--Gordan compression.  An outgoing mark
\(f\in\operatorname{supp}\boldsymbol S\) is measured with the
\(\boldsymbol S\) probability \(\vartheta_f(\boldsymbol S)\).

\begin{definition}[Cold-intermediate formation-mark inequality]
\label{def:V33-CFMI}
For \(L_0,K<\infty\), write
\(\operatorname{CFMI}(L_0,K)\) if every pair of exact inputs and every
outgoing link \(f\) satisfy
\begin{align}
 &\sum_{\substack{\boldsymbol S\ne\boldsymbol1\\
          s_\alpha\Xi(\boldsymbol S)\le L_0}}
 \norm{\mathsf F_{ij}^{\boldsymbol S}(g_{i,X_i},g_{j,X_j})}
 _{\rho,2;f}^{\rm BFA(1)}
 \notag\\
 &\quad\le K
 \sum_{e\in X_i\cap X_j}
 \left[
 \norm{g_{i,X_i}}_{\rho,2;(e,f)}^{\rm BFA(2)}
 \norm{g_{j,X_j}}_{\rho,2;e}^{\rm BFA(1)}
 +
 \norm{g_{i,X_i}}_{\rho,2;e}^{\rm BFA(1)}
 \norm{g_{j,X_j}}_{\rho,2;(e,f)}^{\rm BFA(2)}
 \right].
 \label{eq:V33-CFMI}
\end{align}
A mark outside an input support is zero, as in V31.  Multiplicity copies
of \(\boldsymbol S\) are summed in trace norm before the left side is
formed.
\end{definition}

\begin{remark}[Why this is the exact missing type]
\label{rem:V33-CFMI-type}
The diagonal term in \eqref{eq:V33-block-decomposition} is the binary
defect for \(\boldsymbol S\) and input \(k\).  V31 assigns its shared
\(\boldsymbol S\)--\(k\) link as the outgoing mark \(f\).  Formula
\eqref{eq:V33-CFMI} lifts that mark through formation of
\(\boldsymbol S\) and adds the missing \(ij\) mark \(e\), placing the
second mark on whichever original input carries \(f\).  Its two terms
are exactly the two possible labeled trees.  The singlet case has no
outgoing intermediate support; its leading binary defect vanishes by
\eqref{eq:V33-spectator-singlet-cancellation} and its formation mark is
already supplied by \eqref{eq:V33-singlet-tax-is-mark}.
\end{remark}

\begin{assessment}[Exact post-V33 bottleneck]
\label{ass:V33-bottleneck}
The common-basis problem and the feared undamped off-diagonal Racah term
are closed by \Cref{thm:V33-common-basis-symbol,
thm:V33-high-pair-reduction}.  Product-singlet formation and
private-dominant tails are closed by
\Cref{lem:V33-Schur-formation-tax,prop:V33-private-tail-payment}.
The first unproved object is now \eqref{eq:V33-CFMI} for nontrivial cold
\(\mathrm{SU}(3)\) intermediates, together with the compatible second-mark
bound for the damped remainder \(E_{ij;k}^{\boldsymbol T}\).  Unitarity
alone cannot prove it by \Cref{prop:V33-unitarity-no-mark}; the next step
must use the actual Littlewood--Richardson multiplicities, dimension
ratios, and Wilson defects.  Until those two formation estimates are
proved uniformly, the \(L\)-dependence in V32 cannot be removed.
\end{assessment}

% =============================================================================
\section{V33 integrated result and next frontier}
\label{sec:V33-ledger}
% =============================================================================

\begin{theorem}[Integrated V33 statement]
\label{thm:V33-integrated}
Preserving V1--V32, V33 proves:
\begin{enumerate}[label=\textup{(V33.\arabic*)},leftmargin=5.5em]
\item the exact common-basis five-term multiplier
      \eqref{eq:V33-common-basis-symbol} and all three anchored identities
      \eqref{eq:V33-anchored-symbol};
\item the diagonal binary-defect/off-diagonal Racah decomposition
      \eqref{eq:V33-block-decomposition}, with every leakage block bounded
      by actual complementary input multipliers;
\item quadratic, multiplicity-uniform damping of every paired high-to-low
      Racah remainder in \eqref{eq:V33-high-pair-reduction};
\item the explicit noncommuting
      \(3\otimes\bar3\otimes3\to3\) Racah block and its exact commutator
      \eqref{eq:V33-pair-commutator};
\item the common-basis Casimir Ward identity
      \eqref{eq:V33-Casimir-Ward};
\item the Schur dimension tax
      \eqref{eq:V33-Schur-tax}, proving that product-singlet formation is
      one normalized pair mark;
\item exact private-link factorization and the private-dominant two-mark
      payment \eqref{eq:V33-private-tree-payment}; and
\item the precise nontrivial cold-intermediate formation card
      \eqref{eq:V33-CFMI} left for V34.
\end{enumerate}
\end{theorem}

\begin{proof}
Items \textup{(V33.1)--(V33.3)} are
\Cref{thm:V33-common-basis-symbol,cor:V33-block-decomposition,
thm:V33-high-pair-reduction}.  Item \textup{(V33.4)} is
\Cref{prop:V33-fundamental-Racah}; item \textup{(V33.5)} is
\Cref{prop:V33-Casimir-Ward}; item \textup{(V33.6)} is
\Cref{lem:V33-Schur-formation-tax,thm:V33-singlet-branch-structure};
item \textup{(V33.7)} is
\Cref{lem:V33-private-factorization,prop:V33-private-tail-payment}; and
item \textup{(V33.8)} is \Cref{def:V33-CFMI,ass:V33-bottleneck}.
\end{proof}

\begingroup
\small
\begin{longtable}{>{\raggedright\arraybackslash}p{0.29\textwidth}
                  >{\raggedright\arraybackslash}p{0.15\textwidth}
                  >{\raggedright\arraybackslash}p{0.48\textwidth}}
\toprule
\textbf{V33 row}&\textbf{Status}&\textbf{Advance or remaining obstruction}\\
\midrule
\endfirsthead
\toprule
\textbf{V33 row}&\textbf{Status}&\textbf{Advance or remaining obstruction}\\
\midrule
\endhead
Common final-channel Racah symbol
& \MGclosed
& All five terms act on one carrier and retain every multiplicity block.\\[0.35em]
High-pair off-diagonal leakage
& \MGclosed
& Every off-diagonal block carries \(q_i\) or \(q_j\) and is quadratically
  damped in the escaping thermal ratios.\\[0.35em]
Fundamental \(\mathrm{SU}(3)\) test block
& \MGclosed
& The exact Fierz matrix proves that pair moment operators do not commute.\\[0.35em]
Product-singlet formation
& \MGclosed
& Schur orthogonality supplies an overlap mark stronger than the uniform
  \(1/|X|\) incidence scale.\\[0.35em]
Private-dominant tail
& \MGclosed
& Exact private multiplier factorization and the quadratic debit pay the
  centered two-mark mass without a thermal cutoff.\\[0.35em]
Nontrivial cold-intermediate formation
& \MGopen
& The outgoing mark must be lifted through the actual multiplicity-resolved
  Littlewood--Richardson compression as in \eqref{eq:V33-CFMI}.\\[0.35em]
Uniform ternary MTA
& \MGopen
& CFMI and the compatible second-mark remainder estimate are still needed
  to remove V32's \(L\)-dependence.\\[0.35em]
All-order marked-tree atlas
& \MGopen
& Higher forests remain downstream of uniform ternary closure.\\[0.35em]
Connected Perron majorant
& \MGreduced
& It follows from all-order MTA through the V31 and V30 compilers.\\[0.35em]
Full Yang--Mills mass gap
& \MGopen
& Uniform marked formation, higher forests, finite edges, capacity, and
  continuum closure remain.\\
\bottomrule
\end{longtable}
\endgroup

\begin{openproblem}[V34 mandate: nontrivial formation and second-mark lift]
\label{op:V34}
\begin{enumerate}[label=\textup{(V34.\arabic*)},leftmargin=5.5em]
\item Decompose every cold product intermediate into linkwise
      \(\mathrm{SU}(3)\) Littlewood--Richardson channels.  Retain the
      multiplicity indices and prove or refute \eqref{eq:V33-CFMI} using
      the exact dimension ratios; do not replace the formation map by a
      generic Racah contraction.
\item Treat linkwise trivial factors separately.  Use
      \eqref{eq:V33-Schur-orthogonality} whenever a dual pair forms a
      singlet, and record the exact support lost by the intermediate.
\item Insert the V31 binary marked quotient for
      \(\boldsymbol S\otimes\boldsymbol R_k\to\boldsymbol T\) only after
      the outgoing \(f\)-mark has been lifted by CFMI.  Charge the final
      output denominator once.
\item For \(E_{ij;k}^{\boldsymbol T}\), retain the multipliers in
      \eqref{eq:V33-Racah-remainder} and prove the compatible second-mark
      trace-norm estimate.  Use the private payment
      \eqref{eq:V33-private-tail-payment} before discarding any residual
      product multiplier.
\item Combine the thermal core, shared formation, singlet, private, and
      high-final branches to prove a constant independent of \(L\) in
      \eqref{eq:V32-bounded-thermal-MTA}, or return a finite resolved
      input/intermediate/final/multiplicity/two-mark witness.
\item Only after uniform ternary closure, pass to general connected
      replica forests and the all-order MTA compiler.  Then continue
      through the nonlinear source packet, unreduced edge matrix,
      capacity, protected sectors, and OS continuum limit.
\item V35 must begin with the scheduled current SM/NS companion audit
      before choosing the next cofinal direction.
\end{enumerate}
\end{openproblem}

\begin{firewall}[End-of-V33 claim boundary]
V33 proves the exact common-basis ternary multiplier, quadratic damping
of all paired high-to-low Racah leakage, the product-singlet formation
mark, and private-dominant tail payment.  It does not prove CFMI for
nontrivial cold intermediates, a thermal-window-independent ternary MTA
constant, the all-order marked-tree atlas, WUFC, CPM, the nonlinear
Perron packet, the unreduced edge margins, bad capacity, protected-sector
floors, local-algebra noncollapse, or a nontrivial continuum OS limit.
The full Yang--Mills mass gap is not proved.
\end{firewall}

\section*{Version V33 append-only record}
\addcontentsline{toc}{section}{Version V33 append-only record}
The complete V32 mathematical source was retained before this continuation.
Its SHA--256 digest was
\begin{center}\scriptsize\ttfamily
f8e4d8a509c2a247cd6deef4f48009b0c4b3f3ba67f2e5eabf6429b31a409a42.
\end{center}
No mathematical content of V32 was altered.  On the next iteration, retain
this full file, append before its unique active document-closing command,
increment the filename to \texttt{\_V34.tex}, and remove only the
superseded V33 file from this note series.  The next companion audit is
V35.

% =============================================================================
% END APPENDED VERSION V33 MATERIAL
% =============================================================================

% =============================================================================
% BEGIN APPENDED VERSION V34 MATERIAL
% =============================================================================

\clearpage
\part{V34 --- Refutation of separated formation, Schur partial-trace
taxes, and exact two-defect chains}
\label{part:V34}

\section{Purpose and correction of the V33 frontier}
\label{sec:V34-purpose}

V33 put the ternary multiplier in one Racah basis and proved that a paired
high-to-low off-diagonal remainder retains quadratic multiplier damping.
It proposed the cold-intermediate formation-mark inequality
\eqref{eq:V33-CFMI} as the next gate.  V34 first tests that gate on an
explicit multiplicity-one \(\mathrm{SU}(3)\) support chain.

The test has a definite outcome: CFMI is false as stated.  It attempts to
extract an overlap mark from raw pair formation before the second Wilson
defect and the final output denominator are assembled.  A chain of
fundamental highest-weight coefficients loses a factor proportional to
its length under that separation.  The full ternary vertex on the same
chain is nevertheless uniformly bounded, because it factors exactly into
two binary defects and a private multiplier.  Thus the counterexample is
to the V33 proof gate, not to ternary MTA.

Going upstream, V34 proves the general multiplicity-resolved Schur
partial-trace tax for every irreducible intermediate.  It then proves the
complete two-mark MTA inequality on the explicit chain, retaining both
defects and paying long private tails with the V32 quadratic debit.  The
remaining problem is thereby corrected to an assembled formation--defect
estimate.  There is no companion audit at V34; V35 must begin with the
scheduled current SM/NS audit.

% =============================================================================
\section{A resolved counterexample to separated CFMI}
\label{sec:V34-CFMI-counterexample}
% =============================================================================

For the fundamental and symmetric-square representations put
\begin{equation}
 a:=1+C_2(3)=\frac73,
 \qquad
 b:=1+C_2(6)=\frac{13}{3}.
 \label{eq:V34-fundamental-masses}
\end{equation}
For \(N\ge1\), take distinct links
\begin{equation}
 E_N:=\{e,f_1,\ldots,f_N\},
 \qquad X_1=E_N,
 \qquad X_2=\{e\}.
 \label{eq:V34-CFMI-supports}
\end{equation}
Put the fundamental representation on every active input link.  Let
\(u(U)=D^3_{11}(U)\) be the normalized highest-weight matrix coefficient
and define
\begin{equation}
 g_{1,N}:=\prod_{h\in E_N}u(U_h),
 \qquad
 g_2:=u(U_e).
 \label{eq:V34-highest-inputs}
\end{equation}
On link \(e\), the product \(u^2\) is exactly the normalized
highest-weight coefficient of \(6=\operatorname{Sym}^2(3)\).  Hence the
product contains, with coefficient one and multiplicity one, the exact
intermediate channel
\begin{equation}
 S_e=6,
 \qquad S_{f_r}=3\quad(1\le r\le N).
 \label{eq:V34-CFMI-intermediate}
\end{equation}

\begin{theorem}[The V33 separated formation card is false]
\label{thm:V34-CFMI-false}
For every \(L_0>0\) and every \(K<\infty\), there are \(N\),
\(\alpha\), the inputs \eqref{eq:V34-highest-inputs}, and the nontrivial
intermediate \eqref{eq:V34-CFMI-intermediate} for which
\begin{equation}
 s_\alpha\Xi(\boldsymbol S)\le L_0
 \label{eq:V34-CFMI-cold}
\end{equation}
but \eqref{eq:V33-CFMI} fails with outgoing mark \(f_1\).  More exactly,
the ratio of its left side to its nonzero right side is
\begin{equation}
 \boxed{
 \frac{\Xi(\boldsymbol R_1)}{a^2}
 =\frac{N+1}{a}=\frac{3(N+1)}7.}
 \label{eq:V34-CFMI-divergence}
\end{equation}
Thus no constant independent of support size proves CFMI, even for
fundamental inputs, one multiplicity copy, and a cold intermediate.
\end{theorem}

\begin{proof}
The masses and heights are
\begin{align}
 \Xi(\boldsymbol R_1)&=(N+1)a,
 &\Xi(\boldsymbol R_2)&=a,
 &\Xi(\boldsymbol S)&=b+Na,
 \label{eq:V34-CFMI-masses}\
 \operatorname{ht}(\boldsymbol R_1)
 +\operatorname{ht}(\boldsymbol R_2)
 &=N+2
 =\operatorname{ht}(\boldsymbol S).
 \label{eq:V34-CFMI-heights}
\end{align}
Since \(s_\alpha\downarrow0\), choose \(\alpha\) so large that
\eqref{eq:V34-CFMI-cold} holds.

A normalized matrix coefficient is a pure Fourier atom with
\(d_R\norm{\widehat u(R)}_1=1\).  The multiplicity-one highest-weight
product has the same atomic normalization.  Therefore the contribution
of \(\boldsymbol S\) to the left side of \eqref{eq:V33-CFMI}, marked at
\(f_1\), is
\begin{equation}
 e^{\rho(N+2)}
 \Xi(\boldsymbol S)\vartheta_{f_1}(\boldsymbol S)
 =e^{\rho(N+2)}a.
 \label{eq:V34-CFMI-left}
\end{equation}
There is only one possible incoming mark, namely \(e\).  The first term
on the right of \eqref{eq:V33-CFMI} equals
\begin{align}
 &e^{\rho(N+1)}
 \frac{a^2}{\Xi(\boldsymbol R_1)}
 \cdot e^\rho a
 =e^{\rho(N+2)}\frac{a^3}{\Xi(\boldsymbol R_1)}.
 \label{eq:V34-CFMI-right}
\end{align}
The second term is zero because \(f_1\notin X_2\).  Dividing
\eqref{eq:V34-CFMI-left} by \eqref{eq:V34-CFMI-right} gives
\eqref{eq:V34-CFMI-divergence}.
\end{proof}

\begin{assessment}[Why this does not refute ternary MTA]
\label{ass:V34-CFMI-not-MTA}
The construction isolates the raw formation map and deliberately omits
the spectator input, its second shared-link defect, and the final output
resolvent.  Those objects occur together in the actual five-term
cumulant.  The V32 theorem already controls this family whenever all
thermal ratios stay bounded.  Below, \Cref{thm:V34-chain-MTA} proves the
actual resolved MTA estimate on the same family for all thermal ratios.
Hence \Cref{thm:V34-CFMI-false} invalidates only the separated V33 gate.
It is an instruction to preserve more of the assembled vertex.
\end{assessment}

% =============================================================================
\section{The general Schur partial-trace formation tax}
\label{sec:V34-Schur-tax}
% =============================================================================

Let \(R,Q,S\) be irreducible unitary representations of a compact group,
and let
\begin{equation}
 \iota_\mu:H_S\longrightarrow H_R\otimes H_Q
 \label{eq:V34-CG-isometry}
\end{equation}
be one isometric Clebsch--Gordan embedding.  Different multiplicity labels
have orthogonal ranges.  Put \(P_\mu=\iota_\mu\iota_\mu^*\).

\begin{lemma}[Exact partial traces of one irreducible copy]
\label{lem:V34-partial-trace}
The two reduced projectors are scalar and satisfy
\begin{equation}
 \boxed{
 \operatorname{Tr}_{H_Q}P_\mu=\frac{d_S}{d_R}I_{H_R},
 \qquad
 \operatorname{Tr}_{H_R}P_\mu=\frac{d_S}{d_Q}I_{H_Q}.}
 \label{eq:V34-partial-trace}
\end{equation}
Consequently, for all \(x\in H_R\) and \(y\in H_Q\),
\begin{equation}
 \norm{P_\mu(x\otimes y)}^2
 \le
 \min\left\{\frac{d_S}{d_R},\frac{d_S}{d_Q},1\right\}
 \norm{x}^2\norm{y}^2.
 \label{eq:V34-product-vector-tax}
\end{equation}
\end{lemma}

\begin{proof}
The range of \(P_\mu\) is invariant under \(R(g)\otimes Q(g)\), so
partial tracing its conjugation identity shows that
\(\operatorname{Tr}_{H_Q}P_\mu\) commutes with every \(R(g)\).
Schur's lemma makes it a scalar multiple of the identity.  Its trace is
\(\operatorname{Tr}P_\mu=d_S\), which fixes the scalar as \(d_S/d_R\).
The second identity is symmetric.

For fixed \(x\), the positive operator
\(\langle x|P_\mu|x\rangle\) on \(H_Q\) has trace
\((d_S/d_R)\norm{x}^2\).  Its operator norm is at most its trace, giving
the first ratio in \eqref{eq:V34-product-vector-tax}.  Interchanging the
two factors gives the second, and orthogonal projection gives the bound
by one.
\end{proof}

\begin{theorem}[Multiplicity-resolved trace-norm formation tax]
\label{thm:V34-trace-formation-tax}
For arbitrary trace-class coefficient matrices \(A\) on \(H_R\) and
\(B\) on \(H_Q\),
\begin{equation}
 \boxed{
 \norm{\iota_\mu^*(A\otimes B)\iota_\mu}_1
 \le \tau(R,Q;S)\norm{A}_1\norm{B}_1,}
 \label{eq:V34-trace-formation-tax}
\end{equation}
where
\begin{equation}
 \tau(R,Q;S):=
 \min\left\{\frac{d_S}{d_R},\frac{d_S}{d_Q},1\right\}
 =\min\left\{\frac{d_S}{\max(d_R,d_Q)},1\right\}.
 \label{eq:V34-formation-tax}
\end{equation}
If \(S\) occurs with multiplicity \(m_{RQ}^S\), then summing its
separately retained copies gives
\begin{equation}
 \sum_{\mu=1}^{m_{RQ}^S}
 \norm{\iota_\mu^*(A\otimes B)\iota_\mu}_1
 \le m_{RQ}^S\tau(R,Q;S)\norm{A}_1\norm{B}_1.
 \label{eq:V34-multiplicity-tax}
\end{equation}
\end{theorem}

\begin{proof}
First take rank-one matrices
\(A=|x\rangle\langle x'|\) and
\(B=|y\rangle\langle y'|\).  Their compression is
\[
 |\iota_\mu^*(x\otimes y)\rangle
 \langle\iota_\mu^*(x'\otimes y')|.
\]
Its trace norm is the product of the two vector norms.  Apply
\eqref{eq:V34-product-vector-tax} to both vectors; the two square roots
multiply to \(\tau(R,Q;S)\).  Singular-value decompositions of \(A\) and
\(B\), followed by the triangle inequality, prove
\eqref{eq:V34-trace-formation-tax}.  Sum the estimate over the orthogonal
multiplicity labels to obtain \eqref{eq:V34-multiplicity-tax}.
\end{proof}

\begin{corollary}[Product-link formation tax]
\label{cor:V34-product-formation-tax}
For product representations \(\boldsymbol R,\boldsymbol Q\), a resolved
intermediate \(\boldsymbol S\), and one tuple of linkwise multiplicity
labels \(\boldsymbol\mu\), the tensor-product compression obeys
\begin{equation}
 \norm{\iota_{\boldsymbol\mu}^*(A\otimes B)
             \iota_{\boldsymbol\mu}}_1
 \le
 \tau(\boldsymbol R,\boldsymbol Q;\boldsymbol S)
 \norm{A}_1\norm{B}_1,
 \label{eq:V34-product-trace-tax}
\end{equation}
with exact tax
\begin{equation}
 \boxed{
 \tau(\boldsymbol R,\boldsymbol Q;\boldsymbol S)
 :=\min\left\{
 \frac{d_{\boldsymbol S}}{d_{\boldsymbol R}},
 \frac{d_{\boldsymbol S}}{d_{\boldsymbol Q}},1\right\}.}
 \label{eq:V34-product-formation-tax}
\end{equation}
If, in addition, \(A=\bigotimes_eA_e\) and
\(B=\bigotimes_eB_e\) factor linkwise, then the stronger estimate
\begin{equation}
 \norm{\iota_{\boldsymbol\mu}^*(A\otimes B)
             \iota_{\boldsymbol\mu}}_1
 \le\prod_e\left[
 \tau(R_e,Q_e;S_e)\norm{A_e}_1\norm{B_e}_1\right]
 \label{eq:V34-factorized-product-tax}
\end{equation}
holds.  A pass-through link, on which one input is trivial and the output
equals the other input, has local tax one.  A linkwise dual product
singlet has global tax \(d_{\boldsymbol R}^{-1}\), recovering
\eqref{eq:V33-Schur-tax}.
\end{corollary}

\begin{proof}
Regard the product channels as irreducible representations of the product
group and apply \Cref{thm:V34-trace-formation-tax}; their dimensions are
the products of the link dimensions, which proves
\eqref{eq:V34-product-formation-tax}.  Under the additional factorization
of \(A\) and \(B\), tensor the one-link compressed matrices and use the
trace cross norm to prove \eqref{eq:V34-factorized-product-tax}.  Without
that factorization, the per-link minima need not choose the same tensor
factor and may not be multiplied.  The two special cases follow directly.
\end{proof}

\begin{lemma}[Dimension versus Casimir for \(\mathrm{SU}(3)\)]
\label{lem:V34-dimension-Casimir}
For \(R=(p,q)\),
\begin{equation}
 d_R=\frac12(p+1)(q+1)(p+q+2),
 \qquad
 C_2(R)=\frac13(p^2+pq+q^2+3p+3q).
 \label{eq:V34-SU3-dimension-Casimir}
\end{equation}
There are group constants \(c_d,C_d>0\) such that
\begin{equation}
 c_d(1+C_2(R))\le d_R
 \le C_d(1+C_2(R))^{3/2}.
 \label{eq:V34-dimension-Casimir-bounds}
\end{equation}
Therefore every multiplicity copy in \(R\otimes Q\to S\) satisfies
\begin{equation}
 \tau(R,Q;S)
 \le C
 \min\left\{1,
 \frac{(1+C_2(S))^{3/2}}
 {1+\max\{C_2(R),C_2(Q)\}}
 \right\}.
 \label{eq:V34-Casimir-formation-tax}
\end{equation}
\end{lemma}

\begin{proof}
Put \(n=p+q\).  The quadratic polynomial
\(p^2+pq+q^2\) lies between \(3n^2/4\) and \(n^2\), so
\(1+C_2(R)\asymp(1+n)^2\).  Also
\((p+1)(q+1)\ge n+1\), giving
\(d_R\ge(n+1)(n+2)/2\), while each factor in the dimension formula is
at most \(n+2\), giving \(d_R\le(n+2)^3/2\).  This proves
\eqref{eq:V34-dimension-Casimir-bounds}.  Insert those two bounds in
\eqref{eq:V34-formation-tax} to obtain
\eqref{eq:V34-Casimir-formation-tax}.
\end{proof}

\begin{corollary}[Fixed-intermediate high-pair suppression]
\label{cor:V34-fixed-intermediate-suppression}
For a fixed intermediate \(S\) and one fixed multiplicity copy,
\begin{equation}
 \tau(R,Q;S)\longrightarrow0
 \label{eq:V34-fixed-intermediate-suppression}
\end{equation}
whenever \(\max\{C_2(R),C_2(Q)\}\to\infty\).  The estimate is uniform
over the coefficient matrices.  It does not by itself sum a growing
cold bank or its Littlewood--Richardson multiplicities; those factors
remain explicit in \eqref{eq:V34-multiplicity-tax}.
\end{corollary}

\begin{proof}
Apply \eqref{eq:V34-Casimir-formation-tax} with fixed numerator and then
\eqref{eq:V34-trace-formation-tax}.
\end{proof}
% =============================================================================
\section{Exact two-defect factorization on a support tree}
\label{sec:V34-two-defect-factorization}
% =============================================================================

\begin{lemma}[Cumulant factorization across two independent overlap banks]
\label{lem:V34-tree-cumulant-factorization}
Let \(L,R,P\) be independent coordinate banks.  Suppose
\begin{equation}
 F_1=A(L)B(R)H(P),
 \qquad F_2=C(L),
 \qquad F_3=D(R),
 \label{eq:V34-tree-product-variables}
\end{equation}
with all moments finite.  Then
\begin{equation}
 \boxed{
 \kappa_3(F_1,F_2,F_3)
 =\mathbb EH\,\operatorname{Cov}(A,C)
                  \operatorname{Cov}(B,D).}
 \label{eq:V34-tree-cumulant-factorization}
\end{equation}
Independent private factors in \(F_2\) or \(F_3\) multiply the right side
by their means as well.
\end{lemma}

\begin{proof}
Independence gives
\begin{align*}
 \mathbb E(F_1F_2F_3)&=\mathbb EH\,\mathbb E(AC)\mathbb E(BD),\\
 \mathbb E(F_1F_2)\mathbb EF_3
 &=\mathbb EH\,\mathbb E(AC)\mathbb EB\mathbb ED,\\
 \mathbb E(F_1F_3)\mathbb EF_2
 &=\mathbb EH\,\mathbb E(BD)\mathbb EA\mathbb EC,
\end{align*}
while the last two terms in the centered five-term identity combine to
\(\mathbb EH\,\mathbb EA\mathbb EB\mathbb EC\mathbb ED\).
Factoring the resulting expression proves
\eqref{eq:V34-tree-cumulant-factorization}.  A private factor occurring
in only one argument enters every term once and hence factors by its mean.
\end{proof}

Return to \eqref{eq:V34-CFMI-supports} and add the spectator support
\begin{equation}
 X_3=\{f_1\},
 \qquad g_3=u(U_{f_1}).
 \label{eq:V34-chain-spectator}
\end{equation}
Choose the highest-weight final channel
\begin{equation}
 T_e=T_{f_1}=6,
 \qquad T_{f_r}=3\quad(2\le r\le N).
 \label{eq:V34-chain-final-channel}
\end{equation}
Put
\begin{equation}
 \eta:=\eta_3(\alpha),
 \qquad \zeta:=\eta_6(\alpha),
 \qquad \delta:=\zeta-\eta^2,
 \qquad q_P:=\eta^{N-1}.
 \label{eq:V34-chain-multipliers}
\end{equation}

\begin{proposition}[Exact resolved multiplier on the fundamental chain]
\label{prop:V34-chain-multiplier}
The output coefficient in \eqref{eq:V34-chain-final-channel} occurs with
multiplicity one and its assembled ternary numerator and final multiplier
are
\begin{equation}
 \boxed{
 K_{123}^{\boldsymbol T}=q_P\delta^2,
 \qquad
 q_{\alpha,\boldsymbol T}=q_P\zeta^2.}
 \label{eq:V34-chain-multiplier}
\end{equation}
The exact two-mark mass of the only support tree is
\begin{equation}
 \frac{H_{12}H_{13}}{\Xi(\boldsymbol R_1)}
 =\frac{a^3}{N+1}.
 \label{eq:V34-chain-tree-mass}
\end{equation}
\end{proposition}

\begin{proof}
The only \(12\) overlap is \(e\), the only \(13\) overlap is \(f_1\),
and \(X_2\cap X_3=\varnothing\).  The links \(f_2,\ldots,f_N\) are
private to input one.  Apply
\Cref{lem:V34-tree-cumulant-factorization}.  On either overlap, projection
to the highest \(6\) channel turns the binary covariance multiplier into
\(\eta_6-\eta_3^2=\delta\); every private link contributes \(\eta\).
This proves the first identity in \eqref{eq:V34-chain-multiplier}.  The
second is the product of the final-channel link multipliers.

Both overlap masses equal \(a^2\), while
\(\Xi(\boldsymbol R_1)=(N+1)a\).  Therefore
\(H_{12}H_{13}/\Xi(\boldsymbol R_1)=a^4/((N+1)a)\), proving
\eqref{eq:V34-chain-tree-mass}.
\end{proof}

\begin{theorem}[Uniform two-mark MTA on the fundamental support chain]
\label{thm:V34-chain-MTA}
There are \(K_{\rm chain}<\infty\) and \(\alpha_{\rm chain}<\infty\),
independent of \(N\), such that
\begin{equation}
 \boxed{
 \frac{\Xi(\boldsymbol T)q_P\lvert\delta\rvert^2}
 {1-q_P\zeta^2}
 \le K_{\rm chain}\frac{a^3}{N+1}}
 \label{eq:V34-chain-MTA}
\end{equation}
for every \(N\ge1\) and \(\alpha\ge\alpha_{\rm chain}\).  After the
common exponential Fourier weight and the three atomic trace weights are
restored, this is exactly the ternary MTA inequality
\eqref{eq:V31-ternary-open} for the resolved channel
\eqref{eq:V34-chain-final-channel}.  It holds for all thermal ratios,
including \(s_\alpha N\to\infty\).
\end{theorem}

\begin{proof}
The final balanced mass and Casimir obey
\begin{equation}
 \Xi(\boldsymbol T)=2b+(N-1)a\le C(N+1),
 \qquad
 C_E(\boldsymbol T)\ge c(N+1).
 \label{eq:V34-chain-output-mass}
\end{equation}
The local fusion defect
\Cref{lem:V31-local-fusion-defect} gives
\begin{equation}
 \lvert\delta\rvert\le Cs_\alpha
 \label{eq:V34-chain-local-defect}
\end{equation}
cofinally, since only the fixed representations \(3,3,6\) occur.  The
two-sided output gap gives
\begin{equation}
 1-q_P\zeta^2
 \ge c\min\{s_\alpha(N+1),1\}.
 \label{eq:V34-chain-output-gap}
\end{equation}

Write \(n=N+1\).  If \(s_\alpha n\le1\), then
\eqref{eq:V34-chain-output-mass}--\eqref{eq:V34-chain-output-gap} bound
the left side of \eqref{eq:V34-chain-MTA} by
\[
 C\frac{n s_\alpha^2}{s_\alpha n}
 =Cs_\alpha\le\frac Cn.
\]
If \(s_\alpha n\ge1\), the denominator is bounded below.  For \(N\ge2\),
apply the private quadratic debit \eqref{eq:V33-private-quadratic-debit}
to the \(N-1\) fundamental private links:
\[
 n q_Ps_\alpha^2
 \le\frac Cn
 \left(s_\alpha(N-1)a\right)^2q_P
 \le\frac Cn.
\]
For \(N=1\), the second case is absent after increasing
\(\alpha_{\rm chain}\) so that \(2s_\alpha<1\).  Thus the left side is
at most \(C/n\) in all cases.  Since \(a>0\) is fixed, enlarge the
constant to obtain \eqref{eq:V34-chain-MTA}.

Finally, the input height sum and final height are both \(N+3\), so the
exponential ratio is one.  All highest-weight Clebsch--Gordan
coefficients used above equal one, and every normalized matrix
coefficient has atomic trace weight one.  Hence no suppressed dimension
or multiplicity factor is hidden when
\eqref{eq:V34-chain-MTA} is translated into BFA.
\end{proof}

\begin{assessment}[What the chain proves about the correct gate]
\label{ass:V34-assembled-gate}
The raw formation estimate loses \(N\) by
\eqref{eq:V34-CFMI-divergence}.  The actual cumulant restores the missing
factor because its numerator is \emph{the product} of the two binary
defects, \(q_P\delta^2\).  One defect supplies each support-tree edge,
while the final denominator and the private multiplier are paid only
after both have been assembled.  Therefore a correct general gate cannot
bound \(\mathsf F_{ij}^{\boldsymbol S}\) before the Racah remainder is
recombined.  It must apply the formation tax
\eqref{eq:V34-product-formation-tax} inside a two-defect or same-link
junction formula for the full \(K_{123}^{\boldsymbol T}\).
\end{assessment}
% =============================================================================
\section{Corrected frontier and V34 ledger}
\label{sec:V34-ledger}
% =============================================================================

\begin{definition}[Assembled taxed formation problem]
\label{def:V34-ATF}
The V34 assembled taxed formation problem, abbreviated ATF, is the
following finite resolved inequality.  In one chosen pair bracketing,
retain
\begin{equation}
 K_{123}^{\boldsymbol T}
 =(q_{\boldsymbol T}I-q_kQ_{ij}^{\boldsymbol T})
   +E_{ij;k}^{\boldsymbol T}
 \label{eq:V34-ATF-symbol}
\end{equation}
as a single operator until either two distinct coordinate defects or one
same-coordinate third-order junction have been exposed.  On every
multiplicity-resolved formation block, retain the exact factor
\(\tau(\boldsymbol R_i,\boldsymbol R_j;\boldsymbol S)\) from
\eqref{eq:V34-product-formation-tax}.  ATF asks for the resulting final
resolvent trace norm to be bounded by the three two-mark masses in
\eqref{eq:V31-ternary-MTA}, uniformly in all supports, representations,
multiplicities, volumes, and cofinal \(\alpha\).
\end{definition}

\begin{remark}[ATF is narrower than an unnamed restatement]
\label{rem:V34-ATF-content}
Unlike full ternary MTA, ATF has already fixed the exact common-basis
symbol, the location of the final denominator, every multiplicity copy,
the Schur tax on that copy, the private multiplier, and the two allowable
connected mechanisms.  What remains is a scalar summation of the taxed
Littlewood--Richardson blocks after the two defects are exposed.  The
counterexample forbids removing any of these fields.
\end{remark}

\begin{theorem}[Integrated V34 statement]
\label{thm:V34-integrated}
Preserving V1--V33, V34 proves:
\begin{enumerate}[label=\textup{(V34.\arabic*)},leftmargin=5.5em]
\item the cold-intermediate formation card \eqref{eq:V33-CFMI} is false,
      with exact divergence \eqref{eq:V34-CFMI-divergence};
\item the failure is not a ternary MTA counterexample and identifies the
      impermissible separation of formation from the second defect;
\item the exact partial-trace identities \eqref{eq:V34-partial-trace}
      and the multiplicity-resolved trace tax
      \eqref{eq:V34-trace-formation-tax};
\item the product-link tax \eqref{eq:V34-product-formation-tax} and its
      quantitative \(\mathrm{SU}(3)\) Casimir bound
      \eqref{eq:V34-Casimir-formation-tax};
\item exact two-defect factorization
      \eqref{eq:V34-tree-cumulant-factorization} on disjoint overlap
      banks; and
\item the thermal-window-independent two-mark estimate
      \eqref{eq:V34-chain-MTA} for the full fundamental support-chain
      family.
\end{enumerate}
\end{theorem}

\begin{proof}
Items \textup{(V34.1)--(V34.2)} are
\Cref{thm:V34-CFMI-false,ass:V34-CFMI-not-MTA}.  Items
\textup{(V34.3)--(V34.4)} are
\Cref{lem:V34-partial-trace,thm:V34-trace-formation-tax,
cor:V34-product-formation-tax,lem:V34-dimension-Casimir}.  Item
\textup{(V34.5)} is \Cref{lem:V34-tree-cumulant-factorization}, and item
\textup{(V34.6)} is
\Cref{prop:V34-chain-multiplier,thm:V34-chain-MTA}.
\end{proof}

\begingroup
\small
\begin{longtable}{>{\raggedright\arraybackslash}p{0.29\textwidth}
                  >{\raggedright\arraybackslash}p{0.15\textwidth}
                  >{\raggedright\arraybackslash}p{0.48\textwidth}}
\toprule
\textbf{V34 row}&\textbf{Status}&\textbf{Advance or remaining obstruction}\\
\midrule
\endfirsthead
\toprule
\textbf{V34 row}&\textbf{Status}&\textbf{Advance or remaining obstruction}\\
\midrule
\endhead
Separated CFMI
& \MGfail
& The explicit fundamental family loses exactly \((N+1)/a\); this proof
  route is closed as invalid.\\[0.35em]
Multiplicity-resolved Schur tax
& \MGclosed
& Every irreducible copy carries the exact partial-trace tax
  \(\min\{d_S/d_R,d_S/d_Q,1\}\).\\[0.35em]
Fundamental support chain
& \MGclosed
& The full cumulant is two binary defects times the private multiplier and
  satisfies MTA uniformly for all \(s_\alpha N\).\\[0.35em]
Assembled taxed formation problem
& \MGopen
& The taxed Littlewood--Richardson copies must be summed only after a
  two-defect or same-link junction decomposition of the full symbol.\\[0.35em]
Uniform ternary MTA
& \MGopen
& ATF is the corrected remaining high-channel gate; V32 still supplies
  every bounded thermal window.\\[0.35em]
All-order marked-tree atlas
& \MGopen
& Higher replica forests remain downstream of uniform ternary closure.\\[0.35em]
Connected Perron majorant
& \MGreduced
& It follows from all-order MTA through the V31 and V30 compilers.\\[0.35em]
Full Yang--Mills mass gap
& \MGopen
& ATF, higher forests, finite edges, capacity, and continuum closure remain.\\
\bottomrule
\end{longtable}
\endgroup

\begin{openproblem}[V35 mandate: companion audit, then assembled taxed formation]
\label{op:V35}
\begin{enumerate}[label=\textup{(V35.\arabic*)},leftmargin=5.5em]
\item Begin with the scheduled audit of the current SM and NS notes in
      \texttt{latex\_notes}.  Import only proved rows whose hypotheses and
      norm topology genuinely transfer to the YM ATF problem.
\item Replace the false separated CFMI route everywhere in the live
      strategy by ATF.  Preserve the V34 counterexample as a regression
      test: no proposed inequality may estimate its raw formation map
      independently of the second defect.
\item Derive a multiplicity-resolved two-coordinate replica formula for
      \eqref{eq:V34-ATF-symbol}.  It must reduce exactly to
      \eqref{eq:V34-tree-cumulant-factorization} when the two overlap banks
      are disjoint.
\item For a same-link junction, combine the V32 fourth-moment cancellation
      with the one-copy tax \eqref{eq:V34-trace-formation-tax}.  Sum the
      actual \(\mathrm{SU}(3)\) Littlewood--Richardson multiplicities; do
      not replace them by the full tensor-product dimension.
\item Combine the bounded thermal core, private payment, spectator singlet,
      high-pair Racah damping, and taxed formation rows to prove uniform
      ternary MTA, or return a finite resolved input/intermediate/final/
      multiplicity/two-mark witness to ATF.
\item Only after uniform ternary closure, construct the general replica
      forest formula and invoke the all-order MTA and WUFC compilers.
      Continue thereafter through the nonlinear source packet, unreduced
      edge matrix, capacity, protected sectors, and OS continuum limit.
\end{enumerate}
\end{openproblem}

\begin{firewall}[End-of-V34 claim boundary]
V34 refutes the separated CFMI proof gate, proves the general
multiplicity-resolved Schur formation tax, and closes the complete
fundamental two-overlap support-chain family in ternary MTA.  It does not
prove ATF for arbitrary taxed Littlewood--Richardson blocks, uniform
ternary MTA, the all-order marked-tree atlas, WUFC, CPM, the nonlinear
Perron packet, the unreduced edge margins, bad capacity, protected-sector
floors, local-algebra noncollapse, or a nontrivial continuum OS limit.
The full Yang--Mills mass gap is not proved.
\end{firewall}

\section*{Version V34 append-only record}
\addcontentsline{toc}{section}{Version V34 append-only record}
The complete V33 mathematical source was retained before this continuation.
Its SHA--256 digest was
\begin{center}\scriptsize\ttfamily
2342c192d5784b91c5ba4ae144f80500f42166f624a027b39bb347139d141532.
\end{center}
No mathematical content of V33 was altered.  On the next iteration, retain
this full file, append before its unique active document-closing command,
increment the filename to \texttt{\_V35.tex}, and remove only the
superseded V34 file from this note series.  V35 begins with the scheduled
SM/NS companion audit.

% =============================================================================
% END APPENDED VERSION V34 MATERIAL
% =============================================================================

% =============================================================================
% BEGIN APPENDED VERSION V35 MATERIAL
% =============================================================================

\clearpage
\part{V35 --- Companion synchronization, amplified Schur taxes, and
private-dominant support-tree closure}
\label{part:V35}

\section{Purpose, acquisition order, and exact scope}
\label{sec:V35-purpose}

V34 replaced the false separated formation card by the assembled taxed
formation problem (ATF).  The first task here is the scheduled five-version
audit of the other two live programmes, performed before the V35 estimate is
chosen.  Its outcome is a proof discipline rather than a numerical import: a
normalized formation resource must remain visible through the full assembly,
and success on one singular stratum cannot replace control of the complete
operator.

The new work goes upstream to a point left implicit in V34.  The Schur
partial-trace tax is stable under arbitrary spectator spaces and arbitrary
entanglement with those spectators.  Taxes may therefore be multiplied along
a fusion tree when each edge introduces a fresh input leaf.  This is stronger
than the factorized-coefficient clause of
\eqref{eq:V34-factorized-product-tax}; it does not assert that linkwise taxes
may be multiplied inside one two-input fusion bank.

This amplified tax and the exact product-coordinate cumulant give a
multiplicity-resolved two-coordinate replica formula.  All
Littlewood--Richardson copies are packaged into an actual subprobability
formation kernel rather than a raw count.  A one-resolvent argument then
closes every disjoint-overlap support-tree channel for which a fixed fraction
of the central mass is private.  The remaining disjoint-tree gate is an
overlap-dominant high-to-low formation problem; the same-link junction remains
separate.  No full ternary MTA or mass gap is claimed here.

\section{Scheduled V35 audit of the current SM and NS endpoints}
\label{sec:V35-companion-audit}

At direction selection the newest completed companion files were
\begin{align}
 &\texttt{Renewal\_SM\_Einstein\_Continuum\_Research\_Notes\_V56.tex},
 \notag\\[-2pt]
 &\qquad\text{SHA--256 }
 \texttt{53e412d3e6ac08b4b440d1c972d8ddca4e3ac86e9fceff4eea50961e4b0d66dc},
 \label{eq:V35-SM-checkpoint}\\
 &\texttt{Renewal\_NS\_Stability\_Research\_Notes\_V30.tex},
 \notag\\[-2pt]
 &\qquad\text{SHA--256 }
 \texttt{cd73a818b9052cea6fc56bf428d62bc402b6474b40bb778a7bda047b82e9067b}.
 \label{eq:V35-NS-checkpoint}
\end{align}
Neither companion file is modified here.

NS V30 proves that probability marking closes arbitrary cutoff-dependent
stopping-time multiplicity, but every positive radial weight which is
Leray-admissible on the whole high-frequency half-line has finite mass and
cannot dominate unweighted radial integration.  A far-frequency atom realizes
the missing radius factor; unmarking \([K,\Lambda K]\) costs
\(\Lambda^2K\).  Its late audit also proves the type-correct warning that
unitary recoupling does not create a mark: the mark must come from the actual
source-to-tail map.

The SM endpoint advanced from V55 to V56 during this audit.  V55 proves that
every constant positive Ohmic three-port matrix has either a glancing rank
defect or a stable-cone zero.  V56 constructs a frequency-dependent
noncommuting impedance from a local positive half-line action.  It satisfies
the exact passivity identity and repairs the full glancing defect, yet its
complete determinant has exactly two open right-half-plane zeros in the
printed subluminal chamber.  Positive ancestry and a local rank repair do not
imply a uniform inverse for the full operator.

\begin{theorem}[V35 companion transfer decision]
\label{thm:V35-companion-decision}
No numerical estimate from \eqref{eq:V35-SM-checkpoint} or
\eqref{eq:V35-NS-checkpoint} has the hypotheses or norm topology of ATF.
Precisely:
\begin{enumerate}[label=\textup{(\roman*)},leftmargin=3.7em]
\item no NS cutoff weight, Leray constant, stopping-time capacity, or annular
      unmarking factor is imported as a Wilson or Casimir estimate;
\item no SM port determinant, passivity constant, impedance, or glancing
      index is imported as a Yang--Mills fusion estimate;
\item the NS result transfers only as an assembly rule: a normalized
      representation resource must be derived from actual Clebsch--Gordan
      block masses before summation, not reconstructed from raw multiplicity;
\item the SM result transfers only as a completeness rule: a local Schur tax
      or one favorable intermediate cannot replace the complete
      multiplicity-resolved, one-resolvent symbol.
\end{enumerate}
\end{theorem}

\begin{proof}
The NS variables are Fourier cutoff radii, stopping times, helical fluxes, and
the Leray budget.  None occurs in the Wilson product-group moment operators
of \eqref{eq:V33-moment-operators}.  The SM variables are boundary ports and
stable Laplace--Fourier frequencies, and its determinant is not the
final-channel resolvent in \eqref{eq:V34-ATF-symbol}.  This proves the type
exclusions.  Items \textup{(iii)} and \textup{(iv)} are the exact logical
contents of the two no-go results after those exclusions: normalization must
precede summation, and a complete test cannot be inferred from one stratum.
\end{proof}

\begin{firewall}[Scope of the scheduled audit]
The audit changes the order in which Yang--Mills blocks are assembled but
supplies no Yang--Mills inequality.  Neither NS regularity, the
Standard-Model--Einstein continuum, nor the Yang--Mills mass gap follows from
this comparison.
\end{firewall}

\section{The Schur tax is stable under arbitrary spectators}
\label{sec:V35-amplified-Schur}

Let \(R,Q,S\) and \(\iota_\mu:H_S\to H_R\otimes H_Q\) be as in
\Cref{sec:V34-Schur-tax}.  Let \(K,L\) be arbitrary finite-dimensional
Hilbert spaces.  After the canonical permutation of factors, define
\begin{equation}
 V_\mu:=\iota_\mu\otimes I_K\otimes I_L:
 H_S\otimes K\otimes L\longrightarrow
 H_R\otimes K\otimes H_Q\otimes L.
 \label{eq:V35-amplified-isometry}
\end{equation}

\begin{theorem}[Completely spectator-stable one-copy tax]
\label{thm:V35-amplified-Schur-tax}
For trace-class \(X\) on \(H_R\otimes K\) and \(Y\) on
\(H_Q\otimes L\),
\begin{equation}
 \boxed{\norm{V_\mu^*(X\otimes Y)V_\mu}_1
 \le \tau(R,Q;S)\norm X_1\norm Y_1,}
 \label{eq:V35-amplified-Schur-tax}
\end{equation}
where
\begin{equation}
 \tau(R,Q;S)=\min\left\{\frac{d_S}{d_R},\frac{d_S}{d_Q},1\right\}.
 \label{eq:V35-amplified-tax-value}
\end{equation}
The constant is independent of \(K,L\) and of spectator entanglement.
\end{theorem}

\begin{proof}
It suffices to take \(X=|x\rangle\langle x'|\) and
\(Y=|y\rangle\langle y'|\), then use singular-value decompositions.  Put
\(P_\mu=\iota_\mu\iota_\mu^*\),
\(\rho_x=\operatorname{Tr}_K|x\rangle\langle x|\), and
\(\sigma_y=\operatorname{Tr}_L|y\rangle\langle y|\).  Then
\begin{equation}
 \norm{V_\mu^*(x\otimes y)}^2
 =\operatorname{Tr}[P_\mu(\rho_x\otimes\sigma_y)].
 \label{eq:V35-compressed-vector-mixed}
\end{equation}
Since \(P_\mu\le I\), this is at most \(\norm x^2\norm y^2\).  Also
\(\sigma_y\le(\operatorname{Tr}\sigma_y)I_{H_Q}\), so
\eqref{eq:V34-partial-trace} gives
\[
 \operatorname{Tr}[P_\mu(\rho_x\otimes\sigma_y)]
 \le\frac{d_S}{d_R}\norm x^2\norm y^2.
\]
The symmetric argument gives \(d_S/d_Q\).  Hence
\begin{equation}
 \norm{V_\mu^*(x\otimes y)}
 \le\tau(R,Q;S)^{1/2}\norm x\norm y.
 \label{eq:V35-amplified-vector-tax}
\end{equation}
Apply this also to \(x',y'\).  The compressed rank-one operator has trace
norm at most \(\tau\norm x\norm{x'}\norm y\norm{y'}\).  Summing the
singular-value decompositions proves the result.
\end{proof}

\begin{corollary}[Fresh-leaf tree taxes multiply]
\label{cor:V35-fresh-leaf-taxes}
In a sequence of resolved fusions where each step joins the assembled
coefficient to one previously unused input, taxes \(\tau_1,\ldots,\tau_r\)
give
\begin{equation}
 \norm{A_{\rm out}}_1\le
 \left(\prod_{a=1}^r\tau_a\right)
 \prod_{i=1}^{r+1}\norm{A_i}_1.
 \label{eq:V35-fresh-leaf-tax-product}
\end{equation}
All unused and previously formed factors may be arbitrary spectators.
\end{corollary}

\begin{proof}
Apply \Cref{thm:V35-amplified-Schur-tax} at the first edge with every unused
factor in the spectator spaces.  The compressed coefficient becomes one
input of the next fusion and the fresh leaf the other.  Induction proves the
claim.
\end{proof}

\begin{remark}[No linkwise overclaim]
Inside one two-input bank, after one link is fused the two input coefficients
are already joined; a second link is not a fresh-leaf fusion.  Thus the
corollary does not justify multiplying the linkwise minima in
\eqref{eq:V34-factorized-product-tax} for arbitrary nonfactorized
coefficients.  It justifies one global bank tax per edge of an input tree.
\end{remark}

\section{Exact multiplicity-resolved two-coordinate replica formula}
\label{sec:V35-two-coordinate-formula}

Fix the labeled support tree with center one and assume
\begin{equation}
 X_2\cap X_3=\varnothing,
 \qquad L:=X_1\cap X_2\ne\varnothing,
 \qquad R:=X_1\cap X_3\ne\varnothing.
 \label{eq:V35-disjoint-overlap-hypothesis}
\end{equation}
Set
\begin{equation}
 P_i:=X_i\setminus(L\cup R),
 \qquad P:=P_1\mathbin{\dot\cup}P_2\mathbin{\dot\cup}P_3.
 \label{eq:V35-private-banks}
\end{equation}
Every link in \(P\) is private to exactly one input.  For a channel
restricted to a bank \(B\), abbreviate
\begin{equation}
 q_{i,B}:=q_{\alpha,\boldsymbol R_{i,B}},
 \qquad q_P:=\prod_{i=1}^3q_{i,P_i},
 \qquad q_{i,\varnothing}:=1.
 \label{eq:V35-bank-multipliers}
\end{equation}

Choose resolved pair intermediates and copies
\begin{equation}
 \boldsymbol S\subset
 \boldsymbol R_{1,L}\otimes\boldsymbol R_{2,L},
 \quad 1\le\mu\le m_L(\boldsymbol S),
 \qquad
 \boldsymbol U\subset
 \boldsymbol R_{1,R}\otimes\boldsymbol R_{3,R},
 \quad 1\le\nu\le m_R(\boldsymbol U).
 \label{eq:V35-two-bank-intermediates}
\end{equation}
Define
\begin{align}
 \delta_L(\boldsymbol S)&:=q_{\alpha,\boldsymbol S}-q_{1,L}q_{2,L},
 &\delta_R(\boldsymbol U)&:=q_{\alpha,\boldsymbol U}-q_{1,R}q_{3,R},
 \label{eq:V35-bank-defects}\\
 \tau_L(\boldsymbol S)&:=
 \tau(\boldsymbol R_{1,L},\boldsymbol R_{2,L};\boldsymbol S),
 &\tau_R(\boldsymbol U)&:=
 \tau(\boldsymbol R_{1,R},\boldsymbol R_{3,R};\boldsymbol U).
 \label{eq:V35-bank-taxes}
\end{align}
The final channel \(\boldsymbol T=\boldsymbol T(\boldsymbol S,\boldsymbol U)\)
equals \(\boldsymbol S\) on \(L\), \(\boldsymbol U\) on \(R\), and the
unique pass-through input representation on \(P\).

\begin{theorem}[Resolved two-bank replica and two-edge tax]
\label{thm:V35-resolved-two-bank}
On the reordered three-input carrier, the complete five-term multiplier is
\begin{equation}
 \boxed{\mathbf K_{123}
 =q_P(\mathbf Q_{12,L}-q_{1,L}q_{2,L}I)
       \otimes(\mathbf Q_{13,R}-q_{1,R}q_{3,R}I).}
 \label{eq:V35-two-bank-operator-factorization}
\end{equation}
Its \((\boldsymbol S,\mu;\boldsymbol U,\nu)\) block is
\begin{equation}
 \boxed{K_{123}^{\boldsymbol T;\mu,\nu}
 =q_P\delta_L(\boldsymbol S)\delta_R(\boldsymbol U).}
 \label{eq:V35-resolved-two-bank-multiplier}
\end{equation}
If \(\mathcal C_{\mu,\nu}(A_1,A_2,A_3)\) is the corresponding resolved
coefficient, formed first along \(12\) on \(L\) and then along \(13\) on
\(R\), then arbitrary trace-class input coefficients satisfy
\begin{equation}
 \boxed{\norm{\mathcal C_{\mu,\nu}(A_1,A_2,A_3)}_1
 \le\tau_L(\boldsymbol S)\tau_R(\boldsymbol U)
      \prod_{i=1}^3\norm{A_i}_1.}
 \label{eq:V35-resolved-two-bank-tax}
\end{equation}
No bank-factorization of any \(A_i\) is assumed.
\end{theorem}

\begin{proof}
On \(L\), only inputs one and two are nontrivial; on \(R\), only inputs one
and three are nontrivial.  All private moment factors are the scalar \(q_P\)
in every cumulant partition.  Write
\(A=\mathbf Q_{12,L}\), \(B=\mathbf Q_{13,R}\),
\(a=q_{1,L}q_{2,L}\), and \(b=q_{1,R}q_{3,R}\).  The five terms of
\eqref{eq:V33-common-basis-symbol} become, in order,
\begin{equation}
 q_PA\otimes B,\quad -q_PbA\otimes I,\quad
 -q_PaI\otimes B,\quad -q_PabI,\quad 2q_PabI.
 \label{eq:V35-five-factorized-terms}
\end{equation}
Their sum is \(q_P(A-aI)\otimes(B-bI)\), proving
\eqref{eq:V35-two-bank-operator-factorization}.  The pair moment operators
are diagonal on their resolved intermediate copies, so compression gives
\eqref{eq:V35-resolved-two-bank-multiplier}.  No Racah change occurs because
the two operators act on disjoint coordinate banks.

For the coefficient bound, fuse input two into input one on the whole bank
\(L\), treating \(R\) and all private factors as spectators.
\Cref{thm:V35-amplified-Schur-tax} gives \(\tau_L\).  Fuse the resulting
coefficient with the fresh third input on \(R\); a second application gives
\(\tau_R\).  This proves \eqref{eq:V35-resolved-two-bank-tax}.
\end{proof}

The formula also supplies the correct multiplicity normalization.  For
nonzero inputs put
\begin{equation}
 \pi_{\boldsymbol S,\mu;\boldsymbol U,\nu}
 :=\frac{d_{\boldsymbol T}
 \norm{\mathcal C_{\mu,\nu}(A_1,A_2,A_3)}_1}
 {\prod_{i=1}^3d_{\boldsymbol R_i}\norm{A_i}_1}.
 \label{eq:V35-formation-kernel}
\end{equation}

\begin{proposition}[The actual LR formation kernel is subprobability]
\label{prop:V35-formation-subprobability}
The weights in \eqref{eq:V35-formation-kernel} satisfy
\begin{align}
 &\pi_{\boldsymbol S,\mu;\boldsymbol U,\nu}\ge0,
 \qquad
 \sum_{\boldsymbol S,\mu,\boldsymbol U,\nu}
 \pi_{\boldsymbol S,\mu;\boldsymbol U,\nu}\le1,
 \label{eq:V35-formation-subprobability}\\
 &\pi_{\boldsymbol S,\mu;\boldsymbol U,\nu}
 \le\frac{d_{\boldsymbol S}d_{\boldsymbol U}}
 {d_{\boldsymbol R_{1,L}}d_{\boldsymbol R_{2,L}}
  d_{\boldsymbol R_{1,R}}d_{\boldsymbol R_{3,R}}}
 \tau_L(\boldsymbol S)\tau_R(\boldsymbol U).
 \label{eq:V35-formation-kernel-tax}
\end{align}
Thus the sum is over actual block masses; no bare LR multiplicity occurs.
\end{proposition}

\begin{proof}
Conjugate \(A_1\otimes A_2\otimes A_3\) by the sequential
Clebsch--Gordan unitary and pinch to all diagonal
\((\boldsymbol S,\mu;\boldsymbol U,\nu)\) blocks.  Pinching is trace-norm
contractive, so the sum of the unweighted block trace norms is at most
\(\prod_i\norm{A_i}_1\).  Private dimensions cancel in
\eqref{eq:V35-formation-kernel}; moreover
\begin{equation}
 \frac{d_{\boldsymbol S}d_{\boldsymbol U}}
 {d_{\boldsymbol R_{1,L}}d_{\boldsymbol R_{2,L}}
  d_{\boldsymbol R_{1,R}}d_{\boldsymbol R_{3,R}}}\le1,
 \label{eq:V35-dimension-ratio-at-most-one}
\end{equation}
because each output dimension is at most its tensor-product input dimension.
These facts prove \eqref{eq:V35-formation-subprobability}.  Apply
\eqref{eq:V35-resolved-two-bank-tax} and the same cancellation to obtain
\eqref{eq:V35-formation-kernel-tax}.
\end{proof}

\begin{assessment}[What the two-coordinate formula removes]
The disjoint-overlap part of ATF no longer contains a Racah problem or an
uncontrolled LR count.  Its exact numerator is a product of two binary
defects, and its formation copies carry the subprobability law
\eqref{eq:V35-formation-subprobability}.  The remaining issue is whether one
final denominator and those defects pay the two Casimir marks on
overlap-dominant, high-to-low blocks.  This is a scalar average against
\(\pi\), not a sum of \(m_{RQ}^S\tau(R,Q;S)\).
\end{assessment}

\section{A universal one-resolvent bound for the disjoint tree}
\label{sec:V35-one-resolvent-tree}

Put
\begin{align}
 H_L&:=\sum_{e\in L}A_{1,e}A_{2,e},
 &H_R&:=\sum_{e\in R}A_{1,e}A_{3,e},
 \label{eq:V35-overlap-masses}\\
 \Xi_{\boldsymbol S}&:=\Xi(\boldsymbol S),
 &\Xi_{\boldsymbol U}&:=\Xi(\boldsymbol U),
 &\Xi_P&:=\sum_{i=1}^3\Xi(\boldsymbol R_{i,P_i}).
 \label{eq:V35-output-mass-pieces}
\end{align}
The mass of a trivial bank is zero.  Then
\begin{equation}
 \Xi(\boldsymbol T)=\Xi_{\boldsymbol S}+\Xi_{\boldsymbol U}+\Xi_P,
 \qquad
 q_{\alpha,\boldsymbol T}
 =q_Pq_{\alpha,\boldsymbol S}q_{\alpha,\boldsymbol U}.
 \label{eq:V35-output-split}
\end{equation}

\begin{lemma}[One denominator, two exposed defects]
\label{lem:V35-one-resolvent-two-defects}
There are group constants \(C,\alpha_0<\infty\) such that every resolved
block above and \(\alpha\ge\alpha_0\) obeys
\begin{equation}
 \boxed{
 \frac{\Xi(\boldsymbol T)q_P
 |\delta_L(\boldsymbol S)\delta_R(\boldsymbol U)|}
 {1-q_{\alpha,\boldsymbol T}}
 \le Cs_\alpha H_LH_R.}
 \label{eq:V35-one-resolvent-two-defects}
\end{equation}
No intermediate resolvent is inserted into the cumulant; binary quotients
are used only after exact factorization of its complete numerator.
\end{lemma}

\begin{proof}
The product-bank version of \eqref{eq:V31-local-fusion-defect}, obtained by
the telescoping in \eqref{eq:V31-overlap-defect-upper}, gives
\begin{equation}
 |\delta_L|\le C\min\{s_\alpha H_L,1\},
 \qquad |\delta_R|\le C\min\{s_\alpha H_R,1\}.
 \label{eq:V35-two-bank-defect-bounds}
\end{equation}
If \(\boldsymbol S\) is nonconstant, the unscaled binary quotient
\eqref{eq:V31-unscaled-marked-quotient} gives
\begin{equation}
 \Xi_{\boldsymbol S}
 \frac{|\delta_L|}{1-q_{\alpha,\boldsymbol S}}\le CH_L.
 \label{eq:V35-left-binary-quotient}
\end{equation}
Because \(1-q_{\alpha,\boldsymbol T}\ge
1-q_{\alpha,\boldsymbol S}\), the \(\Xi_{\boldsymbol S}\) contribution
to \eqref{eq:V35-one-resolvent-two-defects} is at most
\(CH_L\min\{s_\alpha H_R,1\}\le Cs_\alpha H_LH_R\).  When
\(\boldsymbol S\) is trivial that mass contribution is zero.  The symmetric
argument treats \(\Xi_{\boldsymbol U}\).

For the private mass, use the final-channel gap.  If
\(s_\alpha C_E(\boldsymbol T)\le1\), exact support and the balanced-mass
equivalence give
\(1-q_{\alpha,\boldsymbol T}\ge cs_\alpha\Xi(\boldsymbol T)\).  Hence
\[
 \frac{\Xi_Pq_P|\delta_L\delta_R|}
 {1-q_{\alpha,\boldsymbol T}}
 \le C\frac{|\delta_L\delta_R|}{s_\alpha}
 \le Cs_\alpha H_LH_R.
\]
If \(s_\alpha C_E(\boldsymbol T)\ge1\), the denominator is bounded below,
while \eqref{eq:V31-balanced-product-damping} on the private channel gives
\(s_\alpha\Xi_Pq_P\le C\).  The same defect bound proves the estimate.
Adding the three mass pieces proves the lemma.
\end{proof}

The factor \(s_\alpha\) is harmless in the thermal core but too large when
the central input escapes.  Private central mass supplies the missing
reciprocal factor without requiring the final mass to be private.

\begin{theorem}[Uniform MTA for private-dominant disjoint-overlap trees]
\label{thm:V35-private-dominant-tree-MTA}
Fix \(\gamma\in(0,1]\) and assume, in addition to
\eqref{eq:V35-disjoint-overlap-hypothesis}, that
\begin{equation}
 X:=\Xi(\boldsymbol R_{1,P_1})
 \ge\gamma\Xi(\boldsymbol R_1).
 \label{eq:V35-central-private-dominance}
\end{equation}
There are \(C_\gamma,\alpha_\gamma<\infty\), independent of supports,
representations, multiplicities, and volume, such that every resolved block
satisfies
\begin{equation}
 \boxed{
 \frac{\Xi(\boldsymbol T)q_P
 |\delta_L(\boldsymbol S)\delta_R(\boldsymbol U)|}
 {1-q_{\alpha,\boldsymbol T}}
 \le C_\gamma\frac{H_LH_R}{\Xi(\boldsymbol R_1)}.}
 \label{eq:V35-private-dominant-scalar-MTA}
\end{equation}
After summing all resolved copies, this is the center-one term of ternary MTA
\eqref{eq:V31-ternary-MTA}.  The branch is closed for arbitrary coefficient
entanglement and actual LR multiplicity.
\end{theorem}

\begin{proof}
Choose a fixed large \(M\).  If \(s_\alpha X\le M\),
\Cref{lem:V35-one-resolvent-two-defects} gives
\[
 Cs_\alpha H_LH_R\le\frac{CM}{X}H_LH_R
 \le\frac{CM}{\gamma\Xi(\boldsymbol R_1)}H_LH_R.
\]

Suppose \(s_\alpha X>M\).  The central private multiplier is a factor of
\(q_P\), and \eqref{eq:V32-quadratic-thermal-damping} gives
\begin{equation}
 q_P\le q_{1,P_1}\le
 \frac{K_{\rm bal,2}}{(s_\alpha X)^2}.
 \label{eq:V35-central-private-quadratic}
\end{equation}
Fix \(M\) so the right side is at most \(1/2\).  Then
\begin{equation}
 1-q_{\alpha,\boldsymbol T}\ge1-q_P\ge\frac12.
 \label{eq:V35-private-final-gap}
\end{equation}
The denominator-free consequence of
\eqref{eq:V35-left-binary-quotient} and its symmetric analogue is
\begin{equation}
 \Xi_{\boldsymbol S}|\delta_L|\le CH_L,
 \qquad \Xi_{\boldsymbol U}|\delta_R|\le CH_R.
 \label{eq:V35-binary-defect-mass}
\end{equation}
Together with \eqref{eq:V35-two-bank-defect-bounds}, the first two output
mass pieces are bounded by
\begin{equation}
 Cq_Ps_\alpha H_LH_R\le\frac{C}{X}H_LH_R,
 \label{eq:V35-overlap-output-payment}
\end{equation}
using \eqref{eq:V35-central-private-quadratic} and \(s_\alpha X>M\).

For the private output mass, \eqref{eq:V33-private-quadratic-debit} and
\eqref{eq:V35-two-bank-defect-bounds} yield
\begin{equation}
 \Xi_Pq_P|\delta_L\delta_R|
 \le C\Xi_Pq_Ps_\alpha^2H_LH_R
 \le\frac{C}{\Xi_P}H_LH_R
 \le\frac{C}{X}H_LH_R.
 \label{eq:V35-private-output-payment}
\end{equation}
Use \eqref{eq:V35-private-final-gap}, add the three pieces, and apply
\(X\ge\gamma\Xi(\boldsymbol R_1)\).  This proves
\eqref{eq:V35-private-dominant-scalar-MTA}.

For the operator statement, divide by the three input balanced masses.  The
right side becomes
\begin{equation}
 \frac{H_LH_R}{\Xi(\boldsymbol R_1)^2
 \Xi(\boldsymbol R_2)\Xi(\boldsymbol R_3)},
 \label{eq:V35-two-mark-normalized-mass}
\end{equation}
exactly the sum over \(e\in L,f\in R\) of the central two-mark and leaf
one-mark weights.  The exponential height ratio is at most one.  Sum the
coefficient blocks using \eqref{eq:V35-formation-subprobability}; no dimension
or multiplicity factor is introduced.
\end{proof}

\begin{remark}[Strict gain over the V33 private payment]
The V33 estimate \eqref{eq:V33-private-tail-payment} assumed both central and
final mass were controlled by \(\Xi_P\).  Theorem
\ref{thm:V35-private-dominant-tree-MTA} removes the final-output assumption:
large overlap outputs are paid by \eqref{eq:V35-binary-defect-mass}.
\end{remark}

\section{The reduced residual after V35}
\label{sec:V35-residual}

\begin{definition}[Overlap-dominant disjoint-tree residual]
\label{def:V35-overlap-residual}
For \eqref{eq:V35-disjoint-overlap-hypothesis}, define
\begin{equation}
 \mathfrak A_{L,R}:=
 \sum_{\boldsymbol S,\mu,\boldsymbol U,\nu}
 \pi_{\boldsymbol S,\mu;\boldsymbol U,\nu}
 \frac{\Xi(\boldsymbol T)q_P
 |\delta_L(\boldsymbol S)\delta_R(\boldsymbol U)|}
 {1-q_{\alpha,\boldsymbol T}}.
 \label{eq:V35-residual-formation-average}
\end{equation}
The remaining disjoint-tree ATF assertion is
\begin{equation}
 \mathfrak A_{L,R}\le
 C\frac{H_LH_R}{\Xi(\boldsymbol R_1)}
 \label{eq:V35-overlap-ATF}
\end{equation}
on sequences for which
\begin{equation}
 s_\alpha\Xi(\boldsymbol R_1)\longrightarrow\infty,
 \qquad
 \frac{\Xi(\boldsymbol R_{1,P_1})}
 {\Xi(\boldsymbol R_1)}\longrightarrow0.
 \label{eq:V35-overlap-dominant-escape}
\end{equation}
Here \(\pi\) is the actual coefficient-dependent subprobability law, not an
envelope obtained by summing bare taxes.
\end{definition}

\begin{proposition}[Failure localization after V35]
\label{prop:V35-failure-localization}
If the disjoint-overlap ternary tree inequality fails, it has a pure resolved
witness sequence satisfying \eqref{eq:V35-overlap-dominant-escape} and
violating \eqref{eq:V35-overlap-ATF}.  No bounded thermal sequence and no
sequence with a uniformly positive central private fraction can witness
failure.
\end{proposition}

\begin{proof}
V32 excludes bounded thermal witnesses by
\Cref{thm:V32-bounded-thermal-MTA,thm:V32-ternary-thermal-escape}.  If the
second limit in \eqref{eq:V35-overlap-dominant-escape} failed, a subsequence
would satisfy \eqref{eq:V35-central-private-dominance} for a fixed
\(\gamma>0\), and \Cref{thm:V35-private-dominant-tree-MTA} would bound it.
Peter--Weyl and exact-support expansion refine any failed finite polynomial
estimate to pure input channels.  Proposition
\ref{prop:V35-formation-subprobability} then expresses the remaining sum as
\eqref{eq:V35-residual-formation-average}.  Hence a failure must have the
stated form.
\end{proof}

\begin{assessment}[Upstream interpretation of the residual]
The missing factor cannot come from private damping, bounded thermal
compactness, a raw LR count, or a unitary Racah change.  It must be generated
inside an overlap bank where a thermally escaping central representation is
formed into a sufficiently cold intermediate.  The correct object is the
average \(\mathfrak A_{L,R}\): \eqref{eq:V35-formation-kernel-tax} taxes each
actual block, while \eqref{eq:V35-formation-subprobability} performs the
multiplicity sum.  A pointwise bound on \(m_{RQ}^S\tau(R,Q;S)\) is sufficient
but no longer necessary.
\end{assessment}

\section{Late companion drift during V35 validation}
\label{sec:V35-late-drift}

The final folder check found two newer companion appendices.  Their read-only
checkpoint is
\begin{align}
 &\texttt{Renewal\_SM\_Einstein\_Continuum\_Research\_Notes\_V57.tex},
 \notag\\[-2pt]
 &\qquad\text{SHA--256 }
 \texttt{27c2641254957d353d870394a9169845bf5895c87c46bc26945821596f6c1f3f},
 \label{eq:V35-late-SM-checkpoint}\\
 &\texttt{Renewal\_NS\_Stability\_Research\_Notes\_V31.tex},
 \notag\\[-2pt]
 &\qquad\text{SHA--256 }
 \texttt{f1121b62238ad5491bb7cf03635ea9934942edf573ed8faaa9ee79e1d38a6696}.
 \label{eq:V35-late-NS-checkpoint}
\end{align}

SM V57 classifies the complete constant six-coordinate canonical
transformer, retains its matrix numerator, and proves exact inherited
passivity after connection to a positive bath.  Nevertheless every regular
member in the stated low-speed class has a stable-boundary zero before its
first internal pole.  It therefore moves upstream to the physical boundary
quotient.  This supplies no Wilson estimate, but reinforces the V35 rule that
the final claim must be made on the actual quotient and after the complete
operator has been assembled.

NS V31 proves an actual Duhamel source-to-tail formation mark with no shell
multiplicity.  Its critical first moment is paid exactly by the critical
source norm; trying to recover it from the globally paid weaker source norm
restores the missing shell weight.  High-parent lifting removes the dyadic
count but presently returns the enstrophy-square continuation stock.  Thus
its remaining gain must come from the correlation between the actual stopping
times, signed transfer, and the active parent tail.  No one of these PDE
quantities is a Yang--Mills representation variable.

\begin{theorem}[Late-drift transfer decision]
\label{thm:V35-late-drift-decision}
The checkpoints \eqref{eq:V35-late-SM-checkpoint}--%
\eqref{eq:V35-late-NS-checkpoint} change no V35 inequality.  They sharpen the
V36 order: estimate the coefficient-dependent average
\(\mathfrak A_{L,R}\) before replacing it by a pointwise tax envelope.  A
pointwise envelope which restores the missing thermal factor is analogous to
the NS continuation-class return, whereas a gain correlated with the actual
formation kernel is not excluded.  Any successful estimate must then be
tested on the complete final-channel quotient, as required independently by
the SM audit.
\end{theorem}

\begin{proof}
The type exclusions of \Cref{thm:V35-companion-decision} remain unchanged.
The NS theorem shows that a genuine formation map can normalize multiplicity
while its crude norm envelope loses the desired scale.  In V35 the genuine
map is precisely \(\pi\), and its crude envelope is
\eqref{eq:V35-formation-kernel-tax}; hence the asserted order is the exact
type-compatible lesson.  SM V57 concerns a different boundary determinant
but proves that retaining the full numerator and physical quotient is
logically prior to a uniform inverse claim.  Neither observation supplies a
new numerical hypothesis, so all V35 proofs remain unchanged.
\end{proof}

\section{V35 ledger and next acquisition order}
\label{sec:V35-ledger}

\begin{theorem}[Integrated V35 statement]
\label{thm:V35-integrated}
Preserving V1--V34, V35 proves:
\begin{enumerate}[label=\textup{(V35.\arabic*)},leftmargin=5.5em]
\item the scheduled SM V56/NS V30 audit and the late-drift SM V57/NS V31
      audit, with the strict transfer firewalls of
      \Cref{thm:V35-companion-decision,thm:V35-late-drift-decision};
\item the spectator-amplified Schur tax
      \eqref{eq:V35-amplified-Schur-tax} and one global bank tax per
      fresh-leaf tree edge;
\item the exact multiplicity-resolved two-coordinate replica formula
      \eqref{eq:V35-two-bank-operator-factorization};
\item the actual LR formation subprobability kernel
      \eqref{eq:V35-formation-subprobability}, eliminating bare multiplicity
      from the disjoint-tree sum;
\item the uniform one-resolvent bound
      \eqref{eq:V35-one-resolvent-two-defects};
\item full ternary MTA on every private-dominant disjoint-overlap support-tree
      branch, without final-output dominance; and
\item localization of every remaining disjoint-tree failure to
      \eqref{eq:V35-overlap-ATF}.
\end{enumerate}
\end{theorem}

\begin{proof}
Item \textup{(V35.1)} is
\Cref{sec:V35-companion-audit,sec:V35-late-drift}.  Items
\textup{(V35.2)--(V35.4)} are
\Cref{thm:V35-amplified-Schur-tax,cor:V35-fresh-leaf-taxes,
thm:V35-resolved-two-bank,prop:V35-formation-subprobability}.  Item
\textup{(V35.5)} is \Cref{lem:V35-one-resolvent-two-defects}.  Items
\textup{(V35.6)--(V35.7)} are
\Cref{thm:V35-private-dominant-tree-MTA,prop:V35-failure-localization}.
\end{proof}

\begingroup
\small
\begin{longtable}{>{\raggedright\arraybackslash}p{0.29\textwidth}
                  >{\raggedright\arraybackslash}p{0.15\textwidth}
                  >{\raggedright\arraybackslash}p{0.48\textwidth}}
\toprule
\textbf{V35 row}&\textbf{Status}&\textbf{Advance or obstruction}\\
\midrule
\endfirsthead
\toprule
\textbf{V35 row}&\textbf{Status}&\textbf{Advance or obstruction}\\
\midrule
\endhead
Scheduled and late-drift SM/NS audit
& \MGclosed
& Exact V56/V30 and then V57/V31 checkpoints were read; only normalization,
  correlation-order, quotient, and complete-operator disciplines transfer.\\[0.35em]
Spectator-amplified Schur tax
& \MGclosed
& The tax survives spectator entanglement and multiplies once per fresh-leaf
  tree edge.\\[0.35em]
Two-coordinate replica formula
& \MGclosed
& Disjoint overlap banks factor into two binary defects; resolved formation
  blocks form a subprobability kernel.\\[0.35em]
Private-dominant disjoint trees
& \MGclosed
& Uniform ternary MTA holds for a fixed central private fraction even with
  arbitrarily large final overlap mass.\\[0.35em]
Overlap-dominant disjoint trees
& \MGopen
& The exact high-to-low kernel average is \eqref{eq:V35-overlap-ATF}.\\[0.35em]
Same-link ternary junction
& \MGopen
& V32 fourth-moment cancellation must be combined with an actual fixed-basis
  multiplicity calculation.\\[0.35em]
Uniform ternary MTA
& \MGopen
& It requires the overlap-dominant average and same-link junction.\\[0.35em]
All-order marked-tree atlas
& \MGopen
& Higher replica forests remain downstream of ternary closure.\\[0.35em]
Full Yang--Mills mass gap
& \MGopen
& Higher forests, WUFC/CPM, nonlinear source packet, finite edges, capacity,
  protected sectors, and the OS continuum limit remain.\\
\bottomrule
\end{longtable}
\endgroup

\begin{openproblem}[V36 mandate: overlap formation average and same-link tax]
\label{op:V36}
\begin{enumerate}[label=\textup{(V36.\arabic*)},leftmargin=5.5em]
\item Attack \eqref{eq:V35-overlap-ATF} against the actual kernel \(\pi\),
      never a bare LR multiplicity sum.
\item Split overlap-dominant escape into warm and genuinely cold
      intermediates.  Pay warm output mass with the binary quotient; combine
      \eqref{eq:V35-formation-kernel-tax} with exact \(\mathrm{SU}(3)\) LR
      parameters only in the cold sector.
\item Retain the V34 fundamental chain as a regression card: every estimate
      must keep both defects and the one final denominator.
\item On one common link, write the V32 five-term third-moment operator in a
      fixed \((12)3\) basis.  Apply the amplified tax only after fourth-moment
      cancellation and retain every off-diagonal Racah block.
\item Combine the bounded-thermal, private-dominant, warm-output,
      spectator-singlet, high-pair damping, and same-link rows to prove
      uniform ternary MTA, or return a finite resolved witness.
\item Only after ternary closure, construct the general replica forest and
      continue through the printed WUFC/CPM and mass-gap programme.  The next
      scheduled SM/NS audit is V40.
\end{enumerate}
\end{openproblem}

\begin{firewall}[End-of-V35 claim boundary]
V35 proves the multiplicity-resolved disjoint-bank replica formula, a
spectator-stable Schur tax, a subprobability formation kernel, and uniform
ternary MTA on the private-dominant disjoint-tree branch.  It does not prove
\eqref{eq:V35-overlap-ATF}, the same-link junction, uniform ternary MTA, the
all-order marked-tree atlas, WUFC, CPM, the nonlinear Perron packet, unreduced
edge margins, capacity, protected-sector floors, local-algebra noncollapse,
reflection positivity, or a nontrivial OS continuum limit.  The full
Yang--Mills mass gap is not proved.
\end{firewall}

\section*{Version V35 append-only record}
\addcontentsline{toc}{section}{Version V35 append-only record}
The complete V34 source was retained before this continuation.  Its
SHA--256 digest was
\begin{center}\scriptsize\ttfamily
cf03a966e4777eadff8746f5fc1497ec048984454f6ea2e5cdf7aa23c82cdecd.
\end{center}
No mathematical content of V34 was altered.  On the next iteration, retain
this full file, append before its unique active document-closing command,
increment the filename to \texttt{\_V36.tex}, and remove only the superseded
V35 file from this Yang--Mills note series.  The next scheduled SM/NS audit is
V40.

% =============================================================================
% END APPENDED VERSION V35 MATERIAL
% =============================================================================

% =============================================================================
% BEGIN APPENDED VERSION V36 MATERIAL
% =============================================================================

\clearpage
\part{V36 --- Warm/cold formation shells, exact LR tax concentration,
and the fixed-basis same-link junction}
\label{part:V36}

\section{Purpose and acquisition order}
\label{sec:V36-purpose}

V35 proved the exact two-defect formula on disjoint overlap banks, packaged
the actual resolved formation blocks into a subprobability kernel \(\pi\), and
closed the private-dominant central branch.  The remaining disjoint-tree
average \eqref{eq:V35-overlap-ATF} can fail only along a thermally escaping,
overlap-dominant central input.

V36 attacks that average before passing to a pointwise multiplicity envelope.
There are two further uniform payments.  If both pair intermediates are warm
relative to the central mass, each binary defect pays the opposite output
piece.  If either intermediate has very small Schur dimension relative to
its two inputs, the exact Littlewood--Richardson dimension identity makes the
total formation weight of that sector small.  The only surviving disjoint
branch is therefore Casimir-cold but dimension-mesoscopic.

The second task is the one-coordinate junction omitted in V35.  A
rectangular, spectator-stable Schur estimate gives the geometric mean of the
taxes on the two sides of every off-diagonal block.  Applying it twice in a
fixed \((12)3\) basis produces an exact taxed junction matrix only after the
five cumulant terms have been assembled.  This retains every Racah block and
does not infer a same-link bound from the disjoint-bank factorization.

There is no scheduled companion audit at V36; the next one remains V40.

\section{The bi-warm disjoint-tree sector}
\label{sec:V36-bi-warm}

Retain all notation from \Cref{sec:V35-two-coordinate-formula}.  For one
resolved pair \((\boldsymbol S,\boldsymbol U)\), put
\begin{equation}
 \mathcal Q_{\boldsymbol S,\boldsymbol U}:=
 \frac{\Xi(\boldsymbol T)q_P
 |\delta_L(\boldsymbol S)\delta_R(\boldsymbol U)|}
 {1-q_{\alpha,\boldsymbol T}}.
 \label{eq:V36-resolved-quotient}
\end{equation}

\begin{theorem}[Uniform payment of bi-warm pair intermediates]
\label{thm:V36-bi-warm-payment}
Fix \(\gamma\in(0,1]\).  If
\begin{equation}
 \Xi(\boldsymbol S)\ge\gamma\Xi(\boldsymbol R_1),
 \qquad
 \Xi(\boldsymbol U)\ge\gamma\Xi(\boldsymbol R_1),
 \label{eq:V36-bi-warm-condition}
\end{equation}
then, uniformly in all supports, representations, copies, volumes, and
cofinal \(\alpha\),
\begin{equation}
 \boxed{\mathcal Q_{\boldsymbol S,\boldsymbol U}
 \le C_\gamma\frac{H_LH_R}{\Xi(\boldsymbol R_1)}.}
 \label{eq:V36-bi-warm-payment}
\end{equation}
Consequently the contribution of all blocks satisfying
\eqref{eq:V36-bi-warm-condition} to
\(\mathfrak A_{L,R}\) satisfies \eqref{eq:V35-overlap-ATF}.
\end{theorem}

\begin{proof}
Write \(X_1=\Xi(\boldsymbol R_1)\).  Fix a large number \(M\).  If
\(s_\alpha X_1\le M\), \Cref{lem:V35-one-resolvent-two-defects} gives
\[
 \mathcal Q_{\boldsymbol S,\boldsymbol U}
 \le Cs_\alpha H_LH_R
 \le\frac{CM}{X_1}H_LH_R.
\]

Suppose \(s_\alpha X_1>M\).  Quadratic product damping and
\eqref{eq:V36-bi-warm-condition} give
\begin{equation}
 q_{\alpha,\boldsymbol S}+q_{\alpha,\boldsymbol U}
 \le\frac{2K_{\rm bal,2}}{(\gamma s_\alpha X_1)^2}.
 \label{eq:V36-warm-output-damping}
\end{equation}
Choose \(M\) so that the right side is at most \(1/2\).  Then
\begin{equation}
 1-q_{\alpha,\boldsymbol T}
 \ge1-q_{\alpha,\boldsymbol S}\ge\frac12.
 \label{eq:V36-warm-final-gap}
\end{equation}
The denominator-free binary defect bounds
\eqref{eq:V35-binary-defect-mass} imply
\begin{equation}
 |\delta_L|\le\frac{CH_L}{\Xi(\boldsymbol S)},
 \qquad
 |\delta_R|\le\frac{CH_R}{\Xi(\boldsymbol U)}.
 \label{eq:V36-warm-defect-ratios}
\end{equation}
Split \(\Xi(\boldsymbol T)\) as in \eqref{eq:V35-output-split}.  The
\(\boldsymbol S\)-mass contribution is at most
\[
 \Xi(\boldsymbol S)|\delta_L\delta_R|
 \le\frac{CH_LH_R}{\Xi(\boldsymbol U)}
 \le\frac{C}{\gamma X_1}H_LH_R.
\]
The \(\boldsymbol U\)-mass contribution is symmetric.

If \(P=\varnothing\), the private contribution is zero.  Otherwise
\eqref{eq:V31-balanced-product-damping} gives
\(\Xi_Pq_P\le C/s_\alpha\).  Hence
\begin{align*}
 \Xi_Pq_P|\delta_L\delta_R|
 &\le\frac{C}{s_\alpha}
 \frac{H_LH_R}{\Xi(\boldsymbol S)\Xi(\boldsymbol U)}\\
 &\le\frac{C}{\gamma^2s_\alpha X_1^2}H_LH_R
 \le\frac{C}{\gamma^2MX_1}H_LH_R.
\end{align*}
Use \eqref{eq:V36-warm-final-gap} and add the three pieces.  This proves
\eqref{eq:V36-bi-warm-payment}.  Finally sum against the subprobability
kernel \eqref{eq:V35-formation-subprobability}.
\end{proof}

\begin{assessment}[What warm formation excludes]
\label{ass:V36-warm-exclusion}
After V35 and \Cref{thm:V36-bi-warm-payment}, a disjoint-tree failure cannot
keep both pair intermediates at a fixed fraction of the central balanced
mass.  Along a witness subsequence, at least one of
\(\Xi(\boldsymbol S)/\Xi(\boldsymbol R_1)\) and
\(\Xi(\boldsymbol U)/\Xi(\boldsymbol R_1)\) must tend to zero.  This is a
genuine high-to-low formation statement, not merely thermal escape of an
input.
\end{assessment}

\section{Exact LR tax concentration and dimension-cold closure}
\label{sec:V36-LR-concentration}

For one two-input product bank, abbreviate
\(\boldsymbol A=\boldsymbol R_{1,L}\),
\(\boldsymbol B=\boldsymbol R_{2,L}\), and define, for every actual LR copy,
\begin{align}
 w_L(\boldsymbol S,\mu)
 &:=\frac{d_{\boldsymbol S}}
 {d_{\boldsymbol A}d_{\boldsymbol B}}
 \tau_L(\boldsymbol S),
 \label{eq:V36-left-LR-weight}\\
 r_L(\boldsymbol S)
 &:=\frac{d_{\boldsymbol S}}
 {\max\{d_{\boldsymbol A},d_{\boldsymbol B}\}}.
 \label{eq:V36-left-dimension-ratio}
\end{align}
Define \(w_R,r_R\) symmetrically on the \(13\) bank.

\begin{lemma}[Exact LR tax concentration]
\label{lem:V36-LR-tax-concentration}
For every \(0<\varepsilon\le1\),
\begin{align}
 \sum_{\boldsymbol S,\mu}w_L(\boldsymbol S,\mu)&\le1,
 \label{eq:V36-LR-weight-total}\\
 \sum_{r_L(\boldsymbol S)\le\varepsilon}
 w_L(\boldsymbol S,\mu)&\le\varepsilon.
 \label{eq:V36-LR-cold-concentration}
\end{align}
The sums use the actual product-group LR multiplicities.  No estimate of
their raw number is taken.
\end{lemma}

\begin{proof}
The exact representation-dimension identity is
\begin{equation}
 \sum_{\boldsymbol S}
 m_{\boldsymbol A\boldsymbol B}^{\boldsymbol S}d_{\boldsymbol S}
 =d_{\boldsymbol A}d_{\boldsymbol B}.
 \label{eq:V36-LR-dimension-identity}
\end{equation}
Because \(\tau_L\le1\), this identity proves
\eqref{eq:V36-LR-weight-total}.  On the set
\(r_L(\boldsymbol S)\le\varepsilon\), the exact Schur tax satisfies
\[
 \tau_L(\boldsymbol S)
 =\min\{r_L(\boldsymbol S),1\}\le\varepsilon.
\]
Insert this inequality into \eqref{eq:V36-left-LR-weight} and apply
\eqref{eq:V36-LR-dimension-identity} again.  This proves
\eqref{eq:V36-LR-cold-concentration}.
\end{proof}

\begin{theorem}[Dimension-cold blocks satisfy the disjoint-tree MTA]
\label{thm:V36-dimension-cold-MTA}
Put
\begin{equation}
 \varepsilon_\alpha:=
 \frac1{1+s_\alpha\Xi(\boldsymbol R_1)}.
 \label{eq:V36-dynamic-dimension-threshold}
\end{equation}
The total contribution to \(\mathfrak A_{L,R}\) from the union
\begin{equation}
 r_L(\boldsymbol S)\le\varepsilon_\alpha
 \quad\hbox{or}\quad
 r_R(\boldsymbol U)\le\varepsilon_\alpha
 \label{eq:V36-dimension-cold-union}
\end{equation}
obeys
\begin{equation}
 \boxed{
 \sum_{\eqref{eq:V36-dimension-cold-union}}
 \pi_{\boldsymbol S,\mu;\boldsymbol U,\nu}
 \mathcal Q_{\boldsymbol S,\boldsymbol U}
 \le C\frac{H_LH_R}{\Xi(\boldsymbol R_1)}.}
 \label{eq:V36-dimension-cold-MTA}
\end{equation}
The constant is independent of support size and all LR multiplicities.
\end{theorem}

\begin{proof}
By \eqref{eq:V35-formation-kernel-tax},
\begin{equation}
 \pi_{\boldsymbol S,\mu;\boldsymbol U,\nu}
 \le w_L(\boldsymbol S,\mu)w_R(\boldsymbol U,\nu).
 \label{eq:V36-kernel-product-envelope}
\end{equation}
Therefore \Cref{lem:V36-LR-tax-concentration} and the union bound give
\begin{equation}
 \sum_{\eqref{eq:V36-dimension-cold-union}}
 \pi_{\boldsymbol S,\mu;\boldsymbol U,\nu}
 \le2\varepsilon_\alpha.
 \label{eq:V36-cold-kernel-mass}
\end{equation}
For every block, \Cref{lem:V35-one-resolvent-two-defects} gives
\(\mathcal Q_{\boldsymbol S,\boldsymbol U}\le
 Cs_\alpha H_LH_R\).  Hence the left side of
\eqref{eq:V36-dimension-cold-MTA} is at most
\[
 \frac{2Cs_\alpha}{1+s_\alpha\Xi(\boldsymbol R_1)}H_LH_R
 \le C\frac{H_LH_R}{\Xi(\boldsymbol R_1)}.
\]
This proves the theorem.
\end{proof}

\begin{remark}[Why the exact dimension identity is not a raw count]
\label{rem:V36-no-raw-LR-count}
The multiplicity in \eqref{eq:V36-LR-dimension-identity} is weighted by the
actual output dimension and normalized by the input tensor dimension.  The
Schur tax supplies the further factor \(r_L\) on the cold set.  Replacing the
identity by \(\sum m_{AB}^S\) would discard both normalizations and would
recreate the forbidden unmarking step identified in the V35 companion audit.
\end{remark}

\section{The dimension-mesoscopic disjoint-tree residual}
\label{sec:V36-mesoscopic-residual}

Let \(\mathfrak A_{L,R}^{\rm mes}\) denote the restriction of
\eqref{eq:V35-residual-formation-average} to blocks satisfying
\begin{equation}
 r_L(\boldsymbol S)>\varepsilon_\alpha,
 \qquad
 r_R(\boldsymbol U)>\varepsilon_\alpha,
 \qquad
 \min\{\Xi(\boldsymbol S),\Xi(\boldsymbol U)\}
 <\gamma\Xi(\boldsymbol R_1).
 \label{eq:V36-mesoscopic-region}
\end{equation}

\begin{proposition}[Complete V36 localization of a disjoint-tree failure]
\label{prop:V36-disjoint-failure-localization}
If the disjoint-overlap ternary MTA fails, then there is a pure resolved
witness sequence along which
\begin{align}
 s_\alpha\Xi(\boldsymbol R_1)&\longrightarrow\infty,
 &
 \frac{\Xi(\boldsymbol R_{1,P_1})}
 {\Xi(\boldsymbol R_1)}&\longrightarrow0,
 \label{eq:V36-residual-input-limits}\\
 \min\left\{
 \frac{\Xi(\boldsymbol S)}{\Xi(\boldsymbol R_1)},
 \frac{\Xi(\boldsymbol U)}{\Xi(\boldsymbol R_1)}
 \right\}&\longrightarrow0,
 &r_L(\boldsymbol S),r_R(\boldsymbol U)
 &>\frac1{1+s_\alpha\Xi(\boldsymbol R_1)},
 \label{eq:V36-residual-output-limits}
\end{align}
and the corresponding mesoscopic average violates
\begin{equation}
 \mathfrak A_{L,R}^{\rm mes}
 \le C\frac{H_LH_R}{\Xi(\boldsymbol R_1)}.
 \label{eq:V36-mesoscopic-ATF}
\end{equation}
Thus the unresolved intermediate is cold in balanced Casimir mass but cannot
be too cold in Schur dimension.
\end{proposition}

\begin{proof}
The first line is \Cref{prop:V35-failure-localization}.  If the minimum in
the second line did not tend to zero, a subsequence would satisfy the
bi-warm hypothesis \eqref{eq:V36-bi-warm-condition} for a fixed
\(\gamma>0\), contradicting \Cref{thm:V36-bi-warm-payment}.  The portion on
which either dimension inequality fails is bounded by
\Cref{thm:V36-dimension-cold-MTA}.  Removing these uniformly paid sectors
from a failed nonnegative majorant leaves a violating subsequence in
\eqref{eq:V36-mesoscopic-region}, proving the result.
\end{proof}

The dimension lower cutoff has a concrete Casimir consequence when a cold
edge consists of one link.

\begin{corollary}[One-link cold intermediates are at least mesoscopic]
\label{cor:V36-one-link-mesoscopic-floor}
Suppose \(L=\{e\}\),
\begin{equation}
 \Xi(\boldsymbol R_{1,L})\ge\gamma_0\Xi(\boldsymbol R_1),
 \qquad s_\alpha\Xi(\boldsymbol R_1)\ge1,
 \label{eq:V36-one-link-central-share}
\end{equation}
and \(r_L(\boldsymbol S)>\varepsilon_\alpha\).  Then
\begin{equation}
 \boxed{\Xi(\boldsymbol S)\ge c_{\gamma_0}s_\alpha^{-2/3}.}
 \label{eq:V36-one-link-mesoscopic-floor}
\end{equation}
The same assertion holds for the right edge.
\end{corollary}

\begin{proof}
By \eqref{eq:V34-dimension-Casimir-bounds},
\[
 d_{\boldsymbol S}\le C\Xi(\boldsymbol S)^{3/2},
 \qquad
 \max\{d_{\boldsymbol R_{1,L}},d_{\boldsymbol R_{2,L}}\}
 \ge d_{\boldsymbol R_{1,L}}
 \ge c\gamma_0\Xi(\boldsymbol R_1).
\]
Hence
\[
 r_L(\boldsymbol S)
 \le\frac{C\Xi(\boldsymbol S)^{3/2}}
 {\gamma_0\Xi(\boldsymbol R_1)}.
\]
The strict lower bound by \(\varepsilon_\alpha\) yields
\[
 \Xi(\boldsymbol S)^{3/2}
 \ge c\gamma_0
 \frac{\Xi(\boldsymbol R_1)}
 {1+s_\alpha\Xi(\boldsymbol R_1)}
 \ge\frac{c\gamma_0}{s_\alpha},
\]
where the last inequality uses \(s_\alpha\Xi(\boldsymbol R_1)\ge1\).
Taking the power \(2/3\) proves the claim.
\end{proof}

\begin{lemma}[Necessary SU(3) LR cancellation coordinates]
\label{lem:V36-SU3-cancellation-coordinates}
Let \(R=(p,q)\), \(Q=(r,s)\), and suppose \(S=(u,v)\) occurs in
\(R\otimes Q\).  Then the numbers
\begin{equation}
 a=\frac{2(p+r-u)+(q+s-v)}3,
 \qquad
 b=\frac{(p+r-u)+2(q+s-v)}3
 \label{eq:V36-SU3-cancellation-coordinates}
\end{equation}
are nonnegative integers and
\begin{align}
 u&=p+r-2a+b,
 &v&=q+s+a-2b,
 \label{eq:V36-SU3-cancellation-reconstruction}\\
 a+b&=(p+q)+(r+s)-(u+v).
 \label{eq:V36-SU3-cancellation-count}
\end{align}
In particular the trialities agree:
\begin{equation}
 u+2v\equiv p+2q+r+2s\pmod3.
 \label{eq:V36-SU3-triality}
\end{equation}
These are necessary LR conditions, not a claim that every such pair
\((a,b)\) occurs.
\end{lemma}

\begin{proof}
The highest weight of every irreducible summand lies below the sum of the two
input highest weights in dominance order.  Hence
\[
 (p+r)\omega_1+(q+s)\omega_2-(u\omega_1+v\omega_2)
 =a\alpha_1+b\alpha_2
\]
for nonnegative integers \(a,b\).  In Dynkin coordinates
\(\alpha_1=(2,-1)\) and \(\alpha_2=(-1,2)\), which gives
\eqref{eq:V36-SU3-cancellation-reconstruction}.  Inverting that two-by-two
system gives \eqref{eq:V36-SU3-cancellation-coordinates}; adding its two
reconstruction equations gives \eqref{eq:V36-SU3-cancellation-count}.
The congruence follows either from those equations or from preservation of
the center character.
\end{proof}

\begin{assessment}[Meaning of the mesoscopic LR shell]
\label{ass:V36-mesoscopic-meaning}
On a one-link residual edge, \Cref{cor:V36-one-link-mesoscopic-floor} prevents
the output Casimir from falling below the scale \(s_\alpha^{-2/3}\), while
\Cref{lem:V36-SU3-cancellation-coordinates} shows that its loss of Dynkin
height is carried by an explicit positive-root cancellation count \(a+b\).
The unresolved shell is therefore neither a fixed representation nor a warm
fraction of the input: it is a growing intermediate produced by a growing
number of root cancellations.  V37 must estimate its actual LR formation
mass across these cancellation coordinates.
\end{assessment}

\begin{assessment}[V34 chain regression]
\label{ass:V36-chain-regression}
The fundamental chain of \Cref{thm:V34-chain-MTA} remains intact.  For long
chains its central mass is private-dominant and is already covered by
\Cref{thm:V35-private-dominant-tree-MTA}; the finitely short remainder is in
a bounded thermal window cofinally.  Neither
\Cref{thm:V36-bi-warm-payment} nor
\Cref{thm:V36-dimension-cold-MTA} estimates its raw first formation map.
Both act on the exact two-defect quotient
\eqref{eq:V35-resolved-two-bank-multiplier} with the single final
denominator retained.  The V34 counterexample therefore remains a passed
regression test.
\end{assessment}

\section{Rectangular Schur taxes and the same-link junction}
\label{sec:V36-same-link}

The diagonal tax of V34 has an off-diagonal form which is needed because the
three pair moment operators are not simultaneously diagonal.

\begin{theorem}[Rectangular spectator-stable Schur tax]
\label{thm:V36-rectangular-Schur-tax}
For \(c=a,b\), let
\(\iota_c:H_{S_c}\to H_{R_c}\otimes H_{Q_c}\) be a resolved irreducible
copy and amplify it by arbitrary spectators \(K_c,L_c\), obtaining \(V_c\).
Let \(X:H_{R_b}\otimes K_b\to H_{R_a}\otimes K_a\) and
\(Y:H_{Q_b}\otimes L_b\to H_{Q_a}\otimes L_a\) be rectangular trace-class
operators.  Then
\begin{equation}
 \boxed{
 \norm{V_a^*(X\otimes Y)V_b}_1
 \le\sqrt{\tau(R_a,Q_a;S_a)\tau(R_b,Q_b;S_b)}
 \norm X_1\norm Y_1.}
 \label{eq:V36-rectangular-Schur-tax}
\end{equation}
The nuclear norm of the resulting rectangular operator is meant.
\end{theorem}

\begin{proof}
For \(X=|x\rangle\langle x'|\) and
\(Y=|y\rangle\langle y'|\), the compressed operator is
\[
 |V_a^*(x\otimes y)\rangle
 \langle V_b^*(x'\otimes y')|.
\]
Apply \eqref{eq:V35-amplified-vector-tax} with the \(a\)-copy to the first
vector and with the \(b\)-copy to the second.  The nuclear norm of a
rectangular rank-one operator is the product of its two vector norms.
Singular-value decompositions of \(X,Y\) finish the proof.
\end{proof}

Now fix one link common to all three inputs and irreducibles
\(R_1,R_2,R_3\).  In the \((12)3\) bracketing, a final multiplicity label is
\begin{equation}
 a=(S_a,\mu_a,\lambda_a),
 \label{eq:V36-final-multiplicity-label}
\end{equation}
where \(\mu_a\) selects a copy of \(S_a\subset R_1\otimes R_2\) and
\(\lambda_a\) selects a copy of \(T\subset S_a\otimes R_3\).  Let
\begin{equation}
 \Phi_a:H_T\longrightarrow H_{R_1}\otimes H_{R_2}\otimes H_{R_3}
 \label{eq:V36-triple-CG-embedding}
\end{equation}
be the composed isometric embedding.  The \(\Phi_a\) have orthogonal ranges
and form the chosen basis of the final multiplicity space.  Put
\begin{equation}
 t_a:=\left[
 \tau(R_1,R_2;S_a)\tau(S_a,R_3;T)
 \right]^{1/2}.
 \label{eq:V36-path-tax}
\end{equation}

For trace-class coefficient matrices define the rectangular final blocks
\begin{equation}
 \mathcal C_{ab}^T(A_1,A_2,A_3)
 :=\Phi_a^*(A_1\otimes A_2\otimes A_3)\Phi_b.
 \label{eq:V36-triple-coefficient-block}
\end{equation}

\begin{corollary}[Two-stage cross tax for a same-link path]
\label{cor:V36-two-stage-cross-tax}
Every pair of final multiplicity paths satisfies
\begin{equation}
 \boxed{
 \norm{\mathcal C_{ab}^T(A_1,A_2,A_3)}_1
 \le t_at_b\prod_{i=1}^3\norm{A_i}_1.}
 \label{eq:V36-two-stage-cross-tax}
\end{equation}
\end{corollary}

\begin{proof}
First apply \Cref{thm:V36-rectangular-Schur-tax} to the \(12\) fusion on
the two sides \((S_a,\mu_a)\) and \((S_b,\mu_b)\).  This gives the factor
\(\sqrt{\tau(R_1,R_2;S_a)\tau(R_1,R_2;S_b)}\).  The resulting rectangular
operator is then fused with the fresh coefficient \(A_3\), using the two
embeddings selected by \(\lambda_a,\lambda_b\).  The rectangular theorem,
whose rank-one proof permits different representation spaces on its two
sides, gives
\(\sqrt{\tau(S_a,R_3;T)\tau(S_b,R_3;T)}\).  Multiplication yields
\(t_at_b\).
\end{proof}

Let \(K_{123}^T\) be the exact one-link multiplicity matrix from
\Cref{thm:V33-common-basis-symbol}.  In this fixed basis it is
\begin{align}
 K_{123}^T
 &=\eta_T I-\eta_3Q_{12}^T-\eta_2Q_{13}^T-\eta_1Q_{23}^T
   +2\eta_1\eta_2\eta_3I
 \label{eq:V36-one-link-five-term-matrix}\\
 &=(\eta_T I-\eta_3Q_{12}^T)+E_{12;3}^T,
 \qquad
 E_{12;3}^T=-\eta_2D_{13}-\eta_1D_{23}.
 \label{eq:V36-one-link-anchored-matrix}
\end{align}
Here \(Q_{12}^T\) is diagonal with entry \(\eta_{S_a}\), whereas the other
two moment operators contain the full Racah matrices.  In particular,
\begin{equation}
 \norm{E_{12;3}^T}_{\rm op}\le\eta_1+\eta_2.
 \label{eq:V36-one-link-Racah-bound}
\end{equation}

\begin{theorem}[Exact fixed-basis same-link assembly]
\label{thm:V36-same-link-assembly}
The resolved output coefficient of the complete one-link third cumulant is
\begin{equation}
 \boxed{
 \widehat\kappa_{3,T}^\alpha
 =\sum_{a,b}(K_{123}^T)_{ba}\,
 \mathcal C_{ab}^T(A_1,A_2,A_3).}
 \label{eq:V36-same-link-coefficient}
\end{equation}
Define the taxed junction matrix size
\begin{equation}
 \mathfrak J_T^{(12)3}:=
 \sum_{a,b}|(K_{123}^T)_{ba}|t_at_b.
 \label{eq:V36-taxed-junction-size}
\end{equation}
Then
\begin{equation}
 \boxed{
 \norm{\widehat\kappa_{3,T}^\alpha}_1
 \le\min\left\{
 \mathfrak J_T^{(12)3},
 C s_\alpha^2A_1A_2A_3
 \right\}\prod_{i=1}^3\norm{A_i}_1.}
 \label{eq:V36-same-link-two-bounds}
\end{equation}
Every off-diagonal Racah entry is retained inside the already assembled
matrix \(K_{123}^T\); the taxes are applied only afterward.
\end{theorem}

\begin{proof}
On the total tensor carrier the five partitions of the third cumulant give
the commuting-action matrix \eqref{eq:V36-one-link-five-term-matrix}.  In the
final isotypic decomposition, contraction of its multiplicity indices with
the coefficient blocks \eqref{eq:V36-triple-coefficient-block} is exactly
\eqref{eq:V36-same-link-coefficient}.  This proves the identity without
separating any cumulant term.

Apply the triangle inequality to that identity and then
\eqref{eq:V36-two-stage-cross-tax}; this gives the
\(\mathfrak J_T^{(12)3}\) bound.  Independently, the matrix trace-norm
lifting in \Cref{prop:V32-Wilson-numerator}, restricted to the one-link
junction term and then to one final block, gives
\(Cs_\alpha^2A_1A_2A_3\prod_i\norm{A_i}_1\).  Taking the smaller bound
proves \eqref{eq:V36-same-link-two-bounds}.
\end{proof}

\begin{corollary}[A sufficient finite same-link ATF card]
\label{cor:V36-same-link-card}
For fixed \(R_1,R_2,R_3\), the same-link part of ternary MTA follows if
\begin{equation}
 \boxed{
 \sum_{T\ne\boldsymbol1}
 \frac{d_T\Xi(T)}{d_{R_1}d_{R_2}d_{R_3}}
 \frac{\min\{\mathfrak J_T^{(12)3},
 C s_\alpha^2A_1A_2A_3\}}
 {1-\eta_T}
 \le C_*A_1A_2A_3}
 \label{eq:V36-same-link-ATF-card}
\end{equation}
holds uniformly.  This is a finite, completely populated
input/intermediate/final/multiplicity/Racah inequality; it contains no
unnamed recoupling constant.
\end{corollary}

\begin{proof}
Multiply \eqref{eq:V36-same-link-two-bounds} by the final BFA dimension,
balanced mass, and the single resolvent multiplier.  Divide by the three
input BFA dimension and balanced-mass factors.  The exponential height ratio
is at most one.  The remaining scalar sum is exactly the left side of
\eqref{eq:V36-same-link-ATF-card}; its right side is the triple-junction
mass in \eqref{eq:V32-triple-junction-mass}.  Summing input coefficients
proves the MTA junction row.
\end{proof}

\begin{firewall}[What the same-link formula does not yet prove]
\label{fire:V36-same-link-scope}
Equation \eqref{eq:V36-same-link-coefficient} is exact, and
\eqref{eq:V36-two-stage-cross-tax} is uniform for each cross block.  The
entrywise sum \(\mathfrak J_T^{(12)3}\) has not yet been bounded uniformly
over the full multiplicity space.  Replacing it by the number of entries
would be the same forbidden raw-count step removed on the disjoint tree.
The next argument must exploit the operator structure of
\(K_{123}^T\), its quadratic high-pair damping, and the diagonal formation
kernel before taking an entrywise absolute value.
\end{firewall}

\section{V36 ledger and V37 acquisition order}
\label{sec:V36-ledger}

\begin{theorem}[Integrated V36 statement]
\label{thm:V36-integrated}
Preserving V1--V35, V36 proves:
\begin{enumerate}[label=\textup{(V36.\arabic*)},leftmargin=5.5em]
\item uniform one-resolvent payment of every bi-warm disjoint-tree block by
      \eqref{eq:V36-bi-warm-payment};
\item the exact LR tax concentration
      \eqref{eq:V36-LR-cold-concentration}, using actual dimension-weighted
      multiplicities rather than their raw count;
\item uniform MTA payment of every dynamically dimension-cold disjoint-tree
      block by \eqref{eq:V36-dimension-cold-MTA};
\item localization of every remaining disjoint-tree failure to the
      Casimir-cold, dimension-mesoscopic shell
      \eqref{eq:V36-residual-input-limits}--%
      \eqref{eq:V36-residual-output-limits};
\item the one-link floor \eqref{eq:V36-one-link-mesoscopic-floor} and the
      exact necessary SU(3) cancellation coordinates
      \eqref{eq:V36-SU3-cancellation-coordinates};
\item the rectangular spectator-stable Schur tax
      \eqref{eq:V36-rectangular-Schur-tax}; and
\item the exact fixed-basis same-link coefficient formula
      \eqref{eq:V36-same-link-coefficient}, its two-stage cross tax, and the
      completely populated sufficient card
      \eqref{eq:V36-same-link-ATF-card}.
\end{enumerate}
\end{theorem}

\begin{proof}
Items \textup{(V36.1)--(V36.3)} are
\Cref{thm:V36-bi-warm-payment,lem:V36-LR-tax-concentration,
thm:V36-dimension-cold-MTA}.  Items \textup{(V36.4)--(V36.5)} are
\Cref{prop:V36-disjoint-failure-localization,
cor:V36-one-link-mesoscopic-floor,lem:V36-SU3-cancellation-coordinates}.
Items \textup{(V36.6)--(V36.7)} are
\Cref{thm:V36-rectangular-Schur-tax,cor:V36-two-stage-cross-tax,
thm:V36-same-link-assembly,cor:V36-same-link-card}.
\end{proof}

\begingroup
\small
\begin{longtable}{>{\raggedright\arraybackslash}p{0.29\textwidth}
                  >{\raggedright\arraybackslash}p{0.15\textwidth}
                  >{\raggedright\arraybackslash}p{0.48\textwidth}}
\toprule
\textbf{V36 row}&\textbf{Status}&\textbf{Advance or obstruction}\\
\midrule
\endfirsthead
\toprule
\textbf{V36 row}&\textbf{Status}&\textbf{Advance or obstruction}\\
\midrule
\endhead
Bi-warm disjoint-tree sector
& \MGclosed
& Both intermediate masses pay the opposite defect after the full numerator
  and single denominator are retained.\\[0.35em]
Dimension-cold LR sector
& \MGclosed
& Exact LR dimension weights and the Schur tax give kernel mass at most
  \((1+s_\alpha\Xi_1)^{-1}\).\\[0.35em]
Mesoscopic disjoint-tree sector
& \MGopen
& A failing block must be Casimir-cold, dimension-above-threshold, thermally
  escaping, and central-private negligible.\\[0.35em]
One-link LR cancellation chart
& \MGclosed
& The residual output has mass at least \(cs_\alpha^{-2/3}\), and its
  height loss is the explicit positive-root count \(a+b\).\\[0.35em]
Same-link fixed-basis assembly
& \MGclosed
& Every diagonal and off-diagonal Racah entry is retained in
  \eqref{eq:V36-same-link-coefficient}; each coefficient block has the
  geometric-mean path tax.\\[0.35em]
Same-link uniform ATF
& \MGopen
& The populated card \eqref{eq:V36-same-link-ATF-card} still requires a
  multiplicity-uniform operator estimate; raw entry counting is forbidden.\\[0.35em]
Uniform ternary MTA
& \MGopen
& It requires the mesoscopic disjoint-tree average and same-link card.\\[0.35em]
All-order marked-tree atlas
& \MGopen
& Higher replica forests remain downstream of uniform ternary closure.\\[0.35em]
Full Yang--Mills mass gap
& \MGopen
& WUFC/CPM, nonlinear source packet, finite edges, capacity, protected
  sectors, and a nontrivial OS continuum state remain downstream.\\
\bottomrule
\end{longtable}
\endgroup

\begin{openproblem}[V37 mandate: mesoscopic LR mass and junction operator]
\label{op:V37}
\begin{enumerate}[label=\textup{(V37.\arabic*)},leftmargin=5.5em]
\item On the disjoint tree, retain the actual coefficient kernel \(\pi\) and
      decompose the mesoscopic branch by the cancellation coordinates
      \((a,b)\) of \eqref{eq:V36-SU3-cancellation-coordinates}.  Seek a
      dimension-weighted tail estimate, not a bound on the number of LR
      tableaux.
\item Extend the one-link calculation to a product overlap bank by an
      entropy or large-deviation inequality for the exact tensor-dimension
      distribution.  The target gain is precisely
      \((s_\alpha\Xi(\boldsymbol R_1))^{-1}\).
\item Test every proposed inequality on the V34 fundamental chain and on a
      one-link dual-pair singlet.  Preserve both Wilson defects and the one
      final resolvent.
\item For the same-link row, split \eqref{eq:V36-one-link-anchored-matrix}
      into its diagonal binary-defect operator and off-diagonal Racah
      remainder.  Combine diagonal formation subprobability with path taxes;
      treat the remainder in operator or Schatten norm before any entrywise
      absolute value.
\item Prove \eqref{eq:V36-same-link-ATF-card} or replace it by a sharper
      sufficient matrix inequality.  If it fails, return a finite
      \((R_1,R_2,R_3;S_a,S_b,T;a,b)\) witness with the exact Wilson
      multipliers and CG/Racah matrices.
\item If both ternary residuals close, state and prove uniform ternary MTA and
      begin the general replica-forest induction.  Otherwise go upstream at
      the first circular or scale-critical bound.  The next scheduled SM/NS
      audit remains V40.
\end{enumerate}
\end{openproblem}

\begin{firewall}[End-of-V36 claim boundary]
V36 closes the bi-warm and dynamically dimension-cold parts of the
disjoint-tree average and derives the exact fixed-basis same-link tax formula.
It does not prove the mesoscopic average \eqref{eq:V36-mesoscopic-ATF}, the
uniform same-link card, uniform ternary MTA, the all-order atlas, WUFC, CPM,
the nonlinear Perron packet, unreduced finite-edge margins, capacity,
protected-sector floors, local-algebra noncollapse, reflection positivity,
or a nontrivial continuum OS limit.  The full Yang--Mills mass gap is not
proved.
\end{firewall}

\section*{Version V36 append-only record}
\addcontentsline{toc}{section}{Version V36 append-only record}
The complete V35 source was retained before this continuation.  Its
SHA--256 digest was
\begin{center}\scriptsize\ttfamily
0fc334c8368467c46f1bdaaa167aeaa848b39f67d1695b414d479ee36abb53b3.
\end{center}
No mathematical content of V35 was altered.  On the next iteration, retain
this full file, append before its unique active document-closing command,
increment the filename to \texttt{\_V37.tex}, and remove only the superseded
V36 file from this Yang--Mills note series.  The next scheduled SM/NS audit is
V40.

% =============================================================================
% END APPENDED VERSION V36 MATERIAL
% =============================================================================

% =============================================================================
% BEGIN APPENDED VERSION V37 MATERIAL
% =============================================================================

\clearpage
\part{V37 --- Exact Casimir-variance formation tails and
multiplicity-free same-link closure}
\label{part:V37}

\section{Purpose and nonduplication}
\label{sec:V37-purpose}

V36 left two genuinely different ternary obstructions.  On a disjoint
support tree, a possible witness had been reduced to an actual
Littlewood--Richardson formation average whose cold intermediate retained
mesoscopic Schur dimension.  On one link common to all three inputs, the
complete five-term cumulant had been assembled in one Racah basis, but its
entrywise taxed matrix size had not been summed uniformly.

This continuation does not repeat either reduction.  For a binary
formation bank it identifies the exact tensor-dimension law as the spectral
law of the total quadratic Casimir.  Its mean and variance can therefore be
computed without enumerating tableaux.  Chebyshev's inequality then closes
the portion of the mesoscopic branch whose Casimir overlap is spread over
at least the thermal number of effective links.  In the cancellation
coordinates of V36, this is an explicit dimension-weighted lower-tail
estimate.

For the same-link row, the V36 coefficient sum is recognized as a partial
trace against the \emph{whole} multiplicity operator.  Trace duality removes
all Racah multiplicity before an absolute value is taken.  An
operator-valued version of the one-link Gaussian-excess argument gives
quadratic small-noise damping for that whole operator.  Splitting only the
final Casimir into low and high regimes then pays the unique final
resolvent.  This proves the uniform same-link ternary MTA row outright; the
stronger entrywise card \eqref{eq:V36-same-link-ATF-card} is not needed.

There is no scheduled companion audit at V37.  In accordance with the
five-version cadence, the next SM/NS audit remains V40.

\section{The exact tensor-dimension Casimir law}
\label{sec:V37-Casimir-law}

Let \(R,Q\) be irreducible representations of \(\mathrm{SU}(3)\), and
decompose
\begin{equation}
 H_R\otimes H_Q
 \simeq
 \bigoplus_{S,\mu}H_S,
 \qquad
 1\le\mu\le m_{RQ}^{S}.
 \label{eq:V37-binary-decomposition}
\end{equation}
The exact dimension identity makes
\begin{equation}
 \nu_{RQ}(S,\mu):=\frac{d_S}{d_Rd_Q}
 \label{eq:V37-dimension-law}
\end{equation}
a probability law on the resolved copies.  Write
\(C_R=C_2(R)\), and similarly for \(Q,S\).

\begin{theorem}[Exact mean and variance of the dimension law]
\label{thm:V37-Casimir-variance}
Under \(\nu_{RQ}\),
\begin{align}
 \mathbb E_{\nu_{RQ}} C_S
 &=C_R+C_Q,
 \label{eq:V37-Casimir-mean}\\
 \operatorname{Var}_{\nu_{RQ}}(C_S)
 &=\frac{4}{\dim\mathfrak{su}(3)}C_RC_Q
 =\frac12C_RC_Q.
 \label{eq:V37-Casimir-variance}
\end{align}
If \(C_R+C_Q>0\), then, for every \(0\le\delta<1\),
\begin{equation}
 \boxed{
 \sum_{\substack{S,\mu\\
 C_S\le\delta(C_R+C_Q)}}
 \frac{d_S}{d_Rd_Q}
 \le
 \frac{C_RC_Q}
 {2(1-\delta)^2(C_R+C_Q)^2}.}
 \label{eq:V37-one-link-cold-tail}
\end{equation}
No Littlewood--Richardson multiplicity has been bounded separately.
\end{theorem}

\begin{proof}
Choose an orthonormal basis \(X_1,\ldots,X_8\) of the Lie algebra in the
Casimir normalization used throughout these notes, and write
\(T_a^R=i\,\mathrm dR(X_a)\), so that the \(T_a^R\) are Hermitian and
\[
 \sum_{a=1}^{8}(T_a^R)^2=C_RI_R.
\]
The total Casimir on \(H_R\otimes H_Q\) is
\begin{equation}
 \mathcal C_{R\otimes Q}
 =(C_R+C_Q)I+2Z,
 \qquad
 Z:=\sum_{a=1}^{8}T_a^R\otimes T_a^Q.
 \label{eq:V37-total-Casimir}
\end{equation}
Every Lie-algebra generator has trace zero because
\(\mathfrak{su}(3)\) is perfect.  Hence the normalized trace of \(Z\)
vanishes.  On the other hand, invariance of the trace form gives
\[
 \operatorname{Tr}_{H_R}(T_a^RT_b^R)
 =\kappa_R\delta_{ab}.
\]
Taking the trace after summing \(a=b\) shows
\begin{equation}
 \kappa_R=\frac{C_Rd_R}{8},
 \qquad
 \kappa_Q=\frac{C_Qd_Q}{8}.
 \label{eq:V37-Dynkin-index}
\end{equation}
It follows that
\begin{align}
 \frac1{d_Rd_Q}\operatorname{Tr}(2Z)^2
 &=\frac4{d_Rd_Q}
   \sum_{a,b}
   \operatorname{Tr}(T_a^RT_b^R)
   \operatorname{Tr}(T_a^QT_b^Q)\\
 &=\frac4{d_Rd_Q}\,8\kappa_R\kappa_Q
 =\frac12C_RC_Q.
 \label{eq:V37-cross-Casimir-square}
\end{align}

In the decomposition \eqref{eq:V37-binary-decomposition},
\(\mathcal C_{R\otimes Q}\) has eigenvalue \(C_S\) on each of the
\(m_{RQ}^S\) copies of \(H_S\).  Its normalized spectral trace is
therefore exactly expectation with respect to
\eqref{eq:V37-dimension-law}.  The normalized trace of
\eqref{eq:V37-total-Casimir} proves
\eqref{eq:V37-Casimir-mean}; the normalized trace of its centered square
and \eqref{eq:V37-cross-Casimir-square} prove
\eqref{eq:V37-Casimir-variance}.  On the event in
\eqref{eq:V37-one-link-cold-tail},
\[
 |C_S-(C_R+C_Q)|
 \ge(1-\delta)(C_R+C_Q).
\]
Chebyshev's inequality now gives \eqref{eq:V37-one-link-cold-tail}.
\end{proof}

\begin{corollary}[Exact tail in the V36 cancellation coordinates]
\label{cor:V37-cancellation-tail}
Fix \(R=(p,q)\), \(Q=(r,s)\), and for every occurring
\(S=(u,v)\) use the unique necessary coordinates \(a,b\) from
\eqref{eq:V36-SU3-cancellation-coordinates}.  If
\(\nu_{p,q;r,s}(a,b,\mu):=d_{(u,v)}/(d_Rd_Q)\) for an actual resolved
copy and zero otherwise, then
\begin{align}
 \sum_{a,b,\mu}\nu_{p,q;r,s}(a,b,\mu)&=1,
 \label{eq:V37-cancellation-law}\\
 \sum_{a,b,\mu}\nu_{p,q;r,s}(a,b,\mu)
 \bigl[C_2(u,v)-C_R-C_Q\bigr]^2
 &=\frac12C_RC_Q,
 \label{eq:V37-cancellation-variance}
\end{align}
and the lower tail in which
\(C_2(u,v)\le\delta(C_R+C_Q)\) obeys
\eqref{eq:V37-one-link-cold-tail}.
Thus the desired estimate across \((a,b)\) is a dimension-weighted
spectral tail, not a count of lattice points or tableaux.
\end{corollary}

\begin{proof}
For fixed inputs, the reconstruction
\eqref{eq:V36-SU3-cancellation-reconstruction} makes an occurring
\((a,b)\) equivalent to its output label \((u,v)\); the copy index
\(\mu\) retains the actual LR multiplicity.  Push
\eqref{eq:V37-dimension-law} and
\eqref{eq:V37-Casimir-variance} through this map.
\end{proof}

\begin{proposition}[Fundamental variance regressions]
\label{prop:V37-fundamental-variance-regressions}
For
\[
 3\otimes3=6\oplus\bar3,
 \qquad
 3\otimes\bar3=\boldsymbol1\oplus8,
\]
the two dimension laws are respectively
\begin{align}
 \nu_{3,3}(6)&=\frac23,&
 \nu_{3,3}(\bar3)&=\frac13,
 \label{eq:V37-33-dimension-law}\\
 \nu_{3,\bar3}(\boldsymbol1)&=\frac19,&
 \nu_{3,\bar3}(8)&=\frac89.
 \label{eq:V37-3bar3-dimension-law}
\end{align}
Both have Casimir mean \(8/3\) and variance \(8/9\), exactly as in
\eqref{eq:V37-Casimir-mean}--\eqref{eq:V37-Casimir-variance}.
For the dual pair, the zero-Casimir lower tail has actual mass \(1/9\),
whereas \eqref{eq:V37-one-link-cold-tail} gives the valid upper bound
\(1/8\).
\end{proposition}

\begin{proof}
Use
\[
 C_2(3)=C_2(\bar3)=\frac43,
 \qquad C_2(6)=\frac{10}3,
 \qquad C_2(8)=3,
\]
and weight each summand by its dimension divided by nine.  Direct
substitution gives
\[
 \frac23\frac{10}3+\frac13\frac43
 =\frac19\,0+\frac89\,3=\frac83
\]
and
\[
 \frac23\left(\frac23\right)^2
 +\frac13\left(-\frac43\right)^2
 =
 \frac19\left(-\frac83\right)^2
 +\frac89\left(\frac13\right)^2
 =\frac89.
\]
At \(\delta=0\), the right side of
\eqref{eq:V37-one-link-cold-tail} is
\((4/3)^2/[2(8/3)^2]=1/8\).
\end{proof}

\subsection{Product banks and an effective overlap number}

Let \(\boldsymbol R,\boldsymbol Q\) be product representations on a finite
bank \(L\).  Decomposition is linkwise, so the resolved dimension law
\begin{equation}
 \nu_{\boldsymbol R\boldsymbol Q}
   (\boldsymbol S,\boldsymbol\mu)
 :=\frac{d_{\boldsymbol S}}
 {d_{\boldsymbol R}d_{\boldsymbol Q}}
 \label{eq:V37-product-dimension-law}
\end{equation}
is the product of the one-link laws.  Put
\begin{align}
 M_L(\boldsymbol R,\boldsymbol Q)
 &:=\sum_{e\in L}\bigl(C_2(R_e)+C_2(Q_e)\bigr),
 \label{eq:V37-bank-mean}\\
 V_L(\boldsymbol R,\boldsymbol Q)
 &:=\sum_{e\in L}C_2(R_e)C_2(Q_e),
 \label{eq:V37-bank-variance-stock}\\
 N_{\rm eff,L}(\boldsymbol R,\boldsymbol Q)
 &:=
 \begin{cases}
 M_L(\boldsymbol R,\boldsymbol Q)^2/
 V_L(\boldsymbol R,\boldsymbol Q),&V_L>0,\\
 +\infty,&V_L=0.
 \end{cases}
 \label{eq:V37-effective-overlap-number}
\end{align}

\begin{theorem}[Exact product-bank variance and cold tail]
\label{thm:V37-product-bank-tail}
Under \(\nu_{\boldsymbol R\boldsymbol Q}\), the random variables
\(C_2(S_e)\) are independent and
\begin{align}
 \mathbb E C_E(\boldsymbol S)&=M_L,
 \label{eq:V37-bank-Casimir-mean}\\
 \operatorname{Var} C_E(\boldsymbol S)&=\frac12V_L.
 \label{eq:V37-bank-Casimir-variance}
\end{align}
If \(M_L>0\), then, for every \(0\le\delta<1\),
\begin{equation}
 \boxed{
 \nu_{\boldsymbol R\boldsymbol Q}
 \{C_E(\boldsymbol S)\le\delta M_L\}
 \le
 \frac1{2(1-\delta)^2N_{\rm eff,L}}.}
 \label{eq:V37-product-bank-cold-tail}
\end{equation}
The estimate is uniform in the number of links and all LR
multiplicities.
\end{theorem}

\begin{proof}
The tensor-product decomposition over \(L\) makes
\eqref{eq:V37-product-dimension-law} a product probability law.
Equations \eqref{eq:V37-Casimir-mean} and
\eqref{eq:V37-Casimir-variance} therefore add over links, proving
\eqref{eq:V37-bank-Casimir-mean}--%
\eqref{eq:V37-bank-Casimir-variance}.  Chebyshev's inequality at distance
\((1-\delta)M_L\) proves \eqref{eq:V37-product-bank-cold-tail}.
\end{proof}

\begin{theorem}[Variance-spread part of the disjoint tree is paid]
\label{thm:V37-variance-spread-MTA}
Retain the disjoint-tree notation of V35--V36 and consider the left
formation bank
\[
 \boldsymbol A=\boldsymbol R_{1,L},
 \qquad
 \boldsymbol B=\boldsymbol R_{2,L}.
\]
Fix \(0<\delta<1\) and \(c_0>0\).  On the sector
\begin{equation}
 C_E(\boldsymbol S)\le\delta M_L(\boldsymbol A,\boldsymbol B),
 \qquad
 N_{\rm eff,L}(\boldsymbol A,\boldsymbol B)
 \ge c_0s_\alpha\Xi(\boldsymbol R_1),
 \label{eq:V37-variance-spread-sector}
\end{equation}
the actual resolved formation average obeys
\begin{equation}
 \boxed{
 \sum_{\eqref{eq:V37-variance-spread-sector}}
 \pi_{\boldsymbol S,\mu;\boldsymbol U,\nu}
 \mathcal Q_{\boldsymbol S,\boldsymbol U}
 \le
 C_{\delta,c_0}
 \frac{H_LH_R}{\Xi(\boldsymbol R_1)}.}
 \label{eq:V37-variance-spread-MTA}
\end{equation}
The same assertion holds with the right bank and
\(\boldsymbol U\) in place of the left bank and \(\boldsymbol S\).
\end{theorem}

\begin{proof}
By \eqref{eq:V36-kernel-product-envelope}, the actual kernel is bounded
by \(w_Lw_R\).  Since each Schur tax is at most one,
\[
 w_L(\boldsymbol S,\mu)
 \le\nu_{\boldsymbol A\boldsymbol B}(\boldsymbol S,\mu),
 \qquad
 \sum_{\boldsymbol U,\nu}w_R(\boldsymbol U,\nu)\le1.
\]
Hence \eqref{eq:V37-product-bank-cold-tail} and the second condition in
\eqref{eq:V37-variance-spread-sector} give
\begin{equation}
 \sum_{\eqref{eq:V37-variance-spread-sector}}
 \pi_{\boldsymbol S,\mu;\boldsymbol U,\nu}
 \le
 \frac{C_{\delta,c_0}}
 {s_\alpha\Xi(\boldsymbol R_1)}.
 \label{eq:V37-spread-kernel-mass}
\end{equation}
The universal assembled two-defect estimate
\(\mathcal Q_{\boldsymbol S,\boldsymbol U}
\le Cs_\alpha H_LH_R\) is
\Cref{lem:V35-one-resolvent-two-defects}.  Multiply it by
\eqref{eq:V37-spread-kernel-mass}; this proves
\eqref{eq:V37-variance-spread-MTA}.  The right-bank proof is identical.
\end{proof}

\begin{corollary}[Further localization of a macroscopic cold bank]
\label{cor:V37-concentrated-bank-localization}
Suppose along a V36 residual sequence that the cold intermediate is
\(\boldsymbol S\), and for some \(c_1>0\)
\begin{equation}
 M_L(\boldsymbol R_{1,L},\boldsymbol R_{2,L})
 \ge c_1\Xi(\boldsymbol R_1).
 \label{eq:V37-macroscopic-bank}
\end{equation}
Then any still-unpaid subsequence must satisfy
\begin{equation}
 \boxed{
 \frac{N_{\rm eff,L}
       (\boldsymbol R_{1,L},\boldsymbol R_{2,L})}
 {s_\alpha\Xi(\boldsymbol R_1)}
 \longrightarrow0.}
 \label{eq:V37-effective-number-failure}
\end{equation}
Thus a cold macroscopic formation bank can remain problematic only when
its cross-Casimir variance is concentrated into fewer than the thermal
number of effective links.
\end{corollary}

\begin{proof}
The exact-support Casimir floor and
\eqref{eq:V36-residual-output-limits} imply
\(C_E(\boldsymbol S)/\Xi(\boldsymbol R_1)\to0\).
Together with \eqref{eq:V37-macroscopic-bank}, this places the sequence,
for every fixed \(\delta\in(0,1)\), in the first condition of
\eqref{eq:V37-variance-spread-sector}.  If
\eqref{eq:V37-effective-number-failure} failed, a subsequence would
satisfy its second condition for a fixed \(c_0>0\), and
\Cref{thm:V37-variance-spread-MTA} would pay it uniformly.
\end{proof}

\begin{firewall}[Scope of the variance tail]
\label{fire:V37-variance-scope}
The variance identity gives a scale gain only when
\(N_{\rm eff,L}\) is large.  A one-link balanced fusion has
\(N_{\rm eff,L}=O(1)\); no concentration theorem can manufacture the
missing factor there.  Within the macroscopic-bank subbranch, the
still-open case is consequently a few-effective-link,
Casimir-cancelling LR problem.  There is also a logically distinct case
in which the globally cold intermediate is comparable to a
nonmacroscopic pair bank; \eqref{eq:V37-product-bank-cold-tail} does not
address it.  The former must be attacked with exact
\(\mathrm{SU}(3)\) multiplicity geometry on the heavy links, and the
latter with the assembled two-defect quotient; neither may be replaced
by iterated Chebyshev bounds or a count of all tableaux.
\end{firewall}

\section{Multiplicity-space partial trace before absolute values}
\label{sec:V37-partial-trace}

\begin{lemma}[Partial trace against a multiplicity operator]
\label{lem:V37-partial-trace}
Let \(H,M\) be finite-dimensional Hilbert spaces,
\(X\in\mathcal S_1(H\otimes M)\), and \(B\in\mathcal B(M)\).  Then
\begin{equation}
 \boxed{
 \left\|
 \operatorname{Tr}_{M}\bigl[X(I_H\otimes B)\bigr]
 \right\|_1
 \le\norm B_{\rm op}\norm X_1.}
 \label{eq:V37-partial-trace-bound}
\end{equation}
If \(P_j\) are mutually orthogonal projections, then
\begin{equation}
 \sum_j\norm{P_jXP_j}_1\le\norm X_1
 \label{eq:V37-pinching-contraction}
\end{equation}
whenever the \(P_j\) resolve the carrier on which \(X\) has been
compressed.
\end{lemma}

\begin{proof}
Trace-norm duality gives
\begin{align*}
 \left\|\operatorname{Tr}_M[X(I\otimes B)]\right\|_1
 &=\sup_{\norm A_{\rm op}\le1}
 \left|\operatorname{Tr}_{H\otimes M}
 (A\otimes I)X(I\otimes B)\right|\\
 &=\sup_{\norm A_{\rm op}\le1}
 \left|\operatorname{Tr}_{H\otimes M}
 (A\otimes B)X\right|
 \le\norm B_{\rm op}\norm X_1.
\end{align*}
For \eqref{eq:V37-pinching-contraction}, the map
\(X\mapsto\sum_jP_jXP_j\) is a finite average of unitary conjugations
(use distinct roots of unity on the ranges of the \(P_j\)).  It is
therefore trace-norm contractive.  Its image is block diagonal, whose
trace norm is the sum of the block trace norms.
\end{proof}

For fixed \(T\), combine the V36 embeddings into the isometry
\begin{equation}
 W_T:H_T\otimes M_T\longrightarrow
 H_{R_1}\otimes H_{R_2}\otimes H_{R_3},
 \qquad
 W_T(x\otimes e_a)=\Phi_a x,
 \label{eq:V37-final-isometry}
\end{equation}
and put
\begin{equation}
 X_T:=W_T^*(A_1\otimes A_2\otimes A_3)W_T.
 \label{eq:V37-isotypic-input-block}
\end{equation}
In the basis \(e_a\), its \(ab\)-block is precisely
\(\mathcal C_{ab}^T\).

\begin{theorem}[Exact partial-trace form of the same-link coefficient]
\label{thm:V37-partial-trace-assembly}
The V36 fixed-basis sum has the basis-free form
\begin{equation}
 \boxed{
 \widehat\kappa_{3,T}^\alpha
 =\operatorname{Tr}_{M_T}
 \left[X_T(I_{H_T}\otimes K_{123}^T)\right].}
 \label{eq:V37-partial-trace-assembly}
\end{equation}
Consequently,
\begin{equation}
 \norm{\widehat\kappa_{3,T}^\alpha}_1
 \le\norm{K_{123}^T}_{\rm op}\norm{X_T}_1.
 \label{eq:V37-multiplicity-free-block}
\end{equation}
Moreover, if \(T\) ranges over all final irreducibles,
\begin{align}
 \sum_T\norm{X_T}_1
 &\le\prod_{i=1}^3\norm{A_i}_1,
 \label{eq:V37-global-isotypic-pinching}\\
 d_T&\le d_{R_1}d_{R_2}d_{R_3}.
 \label{eq:V37-output-dimension-trivial}
\end{align}
\end{theorem}

\begin{proof}
Write
\[
 X_T=\sum_{a,b}\mathcal C_{ab}^T\otimes
 |e_a\rangle\langle e_b|.
\]
The partial trace of \(X_T(I\otimes K_{123}^T)\) is
\(\sum_{a,b}(K_{123}^T)_{ba}\mathcal C_{ab}^T\), which is
\eqref{eq:V36-same-link-coefficient}.  This proves
\eqref{eq:V37-partial-trace-assembly}, and
\Cref{lem:V37-partial-trace} proves
\eqref{eq:V37-multiplicity-free-block}.

The final isotypic projections are mutually orthogonal.  Apply
\eqref{eq:V37-pinching-contraction} to
\(A_1\otimes A_2\otimes A_3\); the trace norm of that tensor product is
the product of the three trace norms.  This gives
\eqref{eq:V37-global-isotypic-pinching}.  Finally \(H_T\) embeds in the
threefold tensor carrier, proving
\eqref{eq:V37-output-dimension-trivial}.
\end{proof}

\begin{corollary}[Diagonal defect and Racah remainder without raw count]
\label{cor:V37-anchored-partial-trace}
In the \((12)3\) bracketing, put
\[
 B_{12;3}^T:=\eta_T I-\eta_3Q_{12}^T.
\]
Then
\begin{align}
 \widehat\kappa_{3,T}^{\rm diag}
 &=\sum_a(\eta_T-\eta_3\eta_{S_a})
   \mathcal C_{aa}^T,
 \label{eq:V37-diagonal-output}\\
 \widehat\kappa_{3,T}^{\rm Racah}
 &=\operatorname{Tr}_{M_T}
   [X_T(I\otimes E_{12;3}^T)],
 \label{eq:V37-Racah-output}
\end{align}
and
\begin{align}
 \norm{\widehat\kappa_{3,T}^{\rm diag}}_1
 &\le
 \sup_a|\eta_T-\eta_3\eta_{S_a}|\,
 \norm{X_T}_1,
 \label{eq:V37-diagonal-no-count}\\
 \norm{\widehat\kappa_{3,T}^{\rm Racah}}_1
 &\le(\eta_1+\eta_2)\norm{X_T}_1.
 \label{eq:V37-Racah-no-count}
\end{align}
Thus both parts are multiplicity-uniform before any pathwise absolute
sum.
\end{corollary}

\begin{proof}
Insert \eqref{eq:V36-one-link-anchored-matrix} into
\eqref{eq:V37-partial-trace-assembly}.  The first operator is diagonal,
giving \eqref{eq:V37-diagonal-output}.  Pinching \(X_T\) by the
one-dimensional path projections on \(M_T\), followed by
\eqref{eq:V37-pinching-contraction}, proves
\eqref{eq:V37-diagonal-no-count}.  Equations
\eqref{eq:V37-Racah-output} and \eqref{eq:V37-Racah-no-count} follow from
\Cref{lem:V37-partial-trace} and
\eqref{eq:V36-one-link-Racah-bound}.
\end{proof}

\section{Operator-valued one-link Gaussian excess}
\label{sec:V37-operator-excess}

The scalar V32 bound can be strengthened at the point needed here.  The
strengthening is not obtained by taking matrix entries and summing them;
the Taylor expansion is run directly in operator norm.

\begin{theorem}[Uniform operator one-link third cumulant]
\label{thm:V37-operator-one-link}
Let \(R_1,R_2,R_3\) be irreducible \(\mathrm{SU}(3)\)
representations on one link and put
\[
 A_i:=1+C_2(R_i).
\]
For all sufficiently large \(\alpha\),
\begin{equation}
 \boxed{
 \norm{\mathbf K_{123}}_{\rm op}
 \le K_{\rm op,3}
 \min\{s_\alpha^2A_1A_2A_3,1\}.}
 \label{eq:V37-operator-one-link}
\end{equation}
Here \(\mathbf K_{123}\) is the complete five-term operator in
\eqref{eq:V33-common-basis-symbol}, restricted to this one coordinate.
In particular every final compression satisfies
\begin{equation}
 \norm{K_{123}^T}_{\rm op}
 \le K_{\rm op,3}
 \min\{s_\alpha^2A_1A_2A_3,1\},
 \label{eq:V37-final-operator-bound}
\end{equation}
uniformly in \(T\) and its multiplicity.
\end{theorem}

\begin{proof}
Write \(D_i(U)=D^{R_i}(U)\) and let \(\mathbb E_\alpha\) denote
expectation under the one-link Wilson increment law.  The operator
\(\mathbf K_{123}\) is the tensor-valued third cumulant
\begin{align}
 \mathbb E_\alpha(D_1\otimes D_2\otimes D_3)
 &-\mathbb E_\alpha(D_1\otimes D_2)\otimes\mathbb E_\alpha D_3
 \notag\\
 &-\mathbb E_\alpha(D_1\otimes D_3)\otimes\mathbb E_\alpha D_2
 -\mathbb E_\alpha(D_2\otimes D_3)\otimes\mathbb E_\alpha D_1
 \notag\\
 &+2\,\mathbb E_\alpha D_1\otimes
       \mathbb E_\alpha D_2\otimes\mathbb E_\alpha D_3,
 \label{eq:V37-operator-cumulant}
\end{align}
with the evident placement of tensor factors.  Centrality gives
\(\mathbb E_\alpha D_i=\eta_iI\), and the joint moment operators are
the \(Q_B\) of \eqref{eq:V33-moment-operators}; hence
\eqref{eq:V37-operator-cumulant} is exactly the claimed
\(\mathbf K_{123}\).

An order-three cumulant vanishes if one argument is constant.  We may
therefore replace \(D_i\) by \(Z_i:=D_i-I\).  On the symmetric normal
ball used in V32, write \(U=\exp X\) and
\begin{equation}
 Z_i=L_i(X)+\mathcal Q_i(X),
 \qquad
 L_i(X):=\mathrm dR_i(X).
 \label{eq:V37-operator-Taylor}
\end{equation}
The Casimir identity and the unitary one-parameter group remainder give
\begin{equation}
 \norm{L_i(X)}_{\rm op}
 \le C A_i^{1/2}\norm X,
 \qquad
 \norm{\mathcal Q_i(X)}_{\rm op}
 \le C A_i\norm X^2.
 \label{eq:V37-operator-jet-bounds}
\end{equation}
Indeed the first estimate follows from
\(\sum_a\mathrm dR_i(X_a)^*\mathrm dR_i(X_a)=C_2(R_i)I\);
the second follows by integrating the second derivative of
\(t\mapsto D_i(\exp(tX))\), whose operator norm is bounded by
\(\norm{\mathrm dR_i(X)}_{\rm op}^2\).  No
representation-dependent exponential remainder occurs because this
one-parameter group is unitary.

The displacement estimates \eqref{eq:V32-displacement-moments}, symmetry
under \(X\mapsto-X\), and the exponentially small complement of the
normal ball imply
\begin{align}
 \norm{\mathbb E_\alpha Z_i}_{\rm op}
 &\le C\min\{s_\alpha A_i,1\},
 \label{eq:V37-operator-first-moment}\\
 \norm{\mathbb E_\alpha(Z_i\otimes Z_j)}_{\rm op}
 &\le C\min\{s_\alpha A_iA_j,1\},
 \label{eq:V37-operator-second-moment}\\
 \norm{\mathbb E_\alpha
   (Z_1\otimes Z_2\otimes Z_3)}_{\rm op}
 &\le C\min\{s_\alpha^2A_1A_2A_3,1\}.
 \label{eq:V37-operator-third-moment}
\end{align}
For completeness, the linear term in
\eqref{eq:V37-operator-first-moment} is odd.  In
\eqref{eq:V37-operator-second-moment}, the \(LL\) term is bounded by
\(Cs_\alpha(A_iA_j)^{1/2}\le Cs_\alpha A_iA_j\); the \(L\mathcal Q\)
and \(\mathcal Q\mathcal Q\) terms are no larger.  In the third raw
moment the \(L_1\otimes L_2\otimes L_3\) term is odd and vanishes.
Every surviving term contains at least one quadratic remainder.  With
one such remainder its bound is
\[
 Cs_\alpha^2 A_i(A_jA_k)^{1/2}
 \le Cs_\alpha^2A_1A_2A_3;
\]
the terms with two or three remainders use the fifth and sixth
displacement moments and are smaller because \(s_\alpha\le1\) and
\(A_i\ge1\).  The complement contributes
\(Ce^{-c/s_\alpha}\), which is absorbed by the displayed bounds.
Finally \(\norm{Z_i}\le2\) supplies each cap by a constant.

Insert \eqref{eq:V37-operator-first-moment}--%
\eqref{eq:V37-operator-third-moment} into the centered formula.
For the pair--mean terms use
\[
 \min\{1,x\}\min\{1,y\}\le\min\{1,xy\},
\]
with \(x=s_\alpha A_iA_j\) and \(y=s_\alpha A_k\).
The product of three means is bounded by
\(\min\{1,s_\alpha^3A_1A_2A_3\}\), hence by the required cap.
This proves the quadratic part of \eqref{eq:V37-operator-one-link}.
Independently, the five terms in
\eqref{eq:V37-operator-cumulant} are contractions and their absolute
coefficients sum to six, giving the universal cap.  Orthogonal
compression to a final isotypic component proves
\eqref{eq:V37-final-operator-bound}.
\end{proof}

\begin{remark}[Why this is stronger than entrywise scalar control]
\label{rem:V37-operator-not-entrywise}
The supremum in \eqref{eq:V37-operator-one-link} is over arbitrary
vectors in the full threefold carrier, including vectors entangled across
all final multiplicity paths.  It is therefore not inferred from a bound
on individual normalized matrix coefficients.  Running the Taylor
argument in operator norm is precisely what makes
\eqref{eq:V37-final-operator-bound} independent of the Racah dimension.
\end{remark}

\section{Uniform same-link resolvent payment}
\label{sec:V37-same-link-closure}

\begin{lemma}[Threefold final-channel geometry]
\label{lem:V37-threefold-fusion}
If \(T\subset R_1\otimes R_2\otimes R_3\), then
\begin{align}
 \operatorname{ht}(T)
 &\le\sum_{i=1}^3\operatorname{ht}(R_i),
 \label{eq:V37-threefold-height}\\
 A_T:=1+C_2(T)
 &\le C_{\rm fus,3}(A_1+A_2+A_3)
 \le3C_{\rm fus,3}A_1A_2A_3,
 \label{eq:V37-threefold-Casimir}\\
 d_T&\le d_{R_1}d_{R_2}d_{R_3}.
 \label{eq:V37-threefold-dimension}
\end{align}
\end{lemma}

\begin{proof}
Apply \Cref{lem:V30-fusion-Casimir} first to
\(R_1\otimes R_2\) and then to the resulting intermediate tensored with
\(R_3\).  This proves the first two inequalities uniformly over the
intermediate.  Since every \(A_i\ge1\), their sum is at most three times
their product.  The dimension bound follows from the isometric inclusion
of one copy of \(H_T\) in the threefold carrier.
\end{proof}

\begin{theorem}[Multiplicity-uniform one-link ternary MTA]
\label{thm:V37-one-link-MTA}
Fix \(\rho\ge0\).  For arbitrary trace-class input coefficient matrices
\(A_i^{\rm in}\) in irreducible channels \(R_i\), let
\(\widehat\kappa_{3,T}^\alpha\) be the exact complete one-link output
coefficient of \eqref{eq:V36-same-link-coefficient}.  There are
\(K_{\rm sl}<\infty\) and \(\alpha_{\rm sl}\), independent of all
representations and multiplicities, such that
\begin{align}
 &\sum_{T\ne\boldsymbol1}
 e^{\rho\operatorname{ht}(T)}
 A_Td_T
 \frac{\norm{\widehat\kappa_{3,T}^\alpha}_1}
      {1-\eta_T}
 \notag\\
 &\qquad\le
 K_{\rm sl}
 \prod_{i=1}^3
 \left[
 e^{\rho\operatorname{ht}(R_i)}
 A_i d_{R_i}\norm{A_i^{\rm in}}_1
 \right]
 \label{eq:V37-one-link-MTA}
\end{align}
for every \(\alpha\ge\alpha_{\rm sl}\).
The final denominator appears exactly once, after the five cumulant terms
and all Racah blocks have been assembled.
\end{theorem}

\begin{proof}
To avoid a clash of notation, in this proof put
\(\mathcal A=A_1^{\rm in}\otimes A_2^{\rm in}\otimes A_3^{\rm in}\)
and \(d_{\rm in}=d_{R_1}d_{R_2}d_{R_3}\).  For any collection
\(\Omega\) of final channels,
\Cref{thm:V37-partial-trace-assembly} gives
\begin{align}
 \sum_{T\in\Omega}d_T
 \norm{\widehat\kappa_{3,T}^\alpha}_1
 &\le
 d_{\rm in}
 \sup_{T\in\Omega}\norm{K_{123}^T}_{\rm op}
 \sum_{T\in\Omega}\norm{X_T}_1\\
 &\le
 d_{\rm in}
 \sup_{T\in\Omega}\norm{K_{123}^T}_{\rm op}
 \prod_{i=1}^3\norm{A_i^{\rm in}}_1.
 \label{eq:V37-subset-output-sum}
\end{align}
This is the step at which all final multiplicities disappear.

Split the nontrivial outputs into
\begin{align*}
 \mathcal L&:=\{T:s_\alpha C_2(T)\le1\},&
 \mathcal H&:=\{T:s_\alpha C_2(T)>1\}.
\end{align*}
The lower one-link gap in \Cref{thm:V30-two-sided-gap} and the positive
nontrivial Casimir floor give
\begin{equation}
 \frac{A_T}{1-\eta_T}\le\frac C{s_\alpha}
 \qquad(T\in\mathcal L).
 \label{eq:V37-low-output-resolvent}
\end{equation}
Use \eqref{eq:V37-final-operator-bound} in
\eqref{eq:V37-subset-output-sum}.  The low-output contribution, with
the exponential weights temporarily omitted, is at most
\begin{align}
 \frac C{s_\alpha}
 d_{\rm in}s_\alpha^2A_1A_2A_3
 \prod_i\norm{A_i^{\rm in}}_1
 &\le
 C A_1A_2A_3d_{\rm in}
 \prod_i\norm{A_i^{\rm in}}_1,
 \label{eq:V37-low-output-payment}
\end{align}
because \(s_\alpha\le1\).

On \(\mathcal H\), the same gap theorem gives
\(1-\eta_T\ge c>0\), while
\eqref{eq:V37-threefold-Casimir} gives
\(A_T\le C A_1A_2A_3\).  Use the universal cap in
\eqref{eq:V37-final-operator-bound} and
\eqref{eq:V37-subset-output-sum}; the high-output contribution is at
most
\begin{equation}
 C A_1A_2A_3d_{\rm in}
 \prod_i\norm{A_i^{\rm in}}_1.
 \label{eq:V37-high-output-payment}
\end{equation}

Finally \eqref{eq:V37-threefold-height} bounds every output exponential
by the product of the three input exponentials.  Insert that bound into
\eqref{eq:V37-low-output-payment} and
\eqref{eq:V37-high-output-payment} and add the two sectors.  The result
is exactly \eqref{eq:V37-one-link-MTA}.
\end{proof}

\begin{corollary}[The same-link ternary junction row is closed]
\label{cor:V37-same-link-closed}
For a common active link \(e\), the reduced one-link third cumulant obeys
the ternary marked-tree atlas estimate with junction mark
\[
 J_{123,e}=A_{1,e}A_{2,e}A_{3,e}.
\]
After summing pure Peter--Weyl inputs, it contributes at most the product
of their three \(e\)-marked BFA seminorms, with a uniform constant.
\end{corollary}

\begin{proof}
On a one-link exact support, the normalized link mark
\(\vartheta_e\) equals one and the BFA weight is
\(e^{\rho\operatorname{ht}(R_i)}A_i d_{R_i}\).
Thus \eqref{eq:V37-one-link-MTA} is exactly the desired pure-block
inequality.  Sum the nonnegative input trace weights; the three input
sums factor into the three marked BFA seminorms.
\end{proof}

\begin{assessment}[Status of the V36 entrywise card]
\label{ass:V37-card-status}
The finite card \eqref{eq:V36-same-link-ATF-card} remains a valid
sufficient condition, but it is stronger than the same-link MTA row:
it takes entrywise absolute values in a fixed multiplicity basis before
summing final channels.  Theorem \ref{thm:V37-one-link-MTA} replaces it
by the sharper summed operator inequality
\eqref{eq:V37-subset-output-sum} and proves that inequality.  No claim
that the stronger entrywise card holds is needed or made.
\end{assessment}

\section{Regression cards and exact claim boundary}
\label{sec:V37-regressions}

\begin{proposition}[Dual-pair spectator regression]
\label{prop:V37-dual-pair-regression}
Let \(R_2=R_1^*\), take the \((12)3\) singlet path
\(S=\boldsymbol1\), and let \(T=R_3\).  Then the diagonal term in
\eqref{eq:V37-diagonal-output} vanishes exactly, and every remaining
contribution is bounded without a multiplicity factor by
\begin{equation}
 \norm{\widehat\kappa_{3,R_3}^{\rm Racah}}_1
 \le2\eta_{R_1}\norm{X_{R_3}}_1.
 \label{eq:V37-dual-pair-regression}
\end{equation}
The full resolved coefficient is also covered by
\eqref{eq:V37-one-link-MTA}.
\end{proposition}

\begin{proof}
Here \(\eta_S=1\), \(\eta_T=\eta_3\), and
\(\eta_1=\eta_2\).  Hence
\(\eta_T-\eta_3\eta_S=0\), reproducing
\eqref{eq:V33-spectator-singlet-cancellation}.  Equation
\eqref{eq:V37-Racah-no-count} gives
\eqref{eq:V37-dual-pair-regression}.  The complete, rather than separated,
coefficient is a special case of \Cref{thm:V37-one-link-MTA}.
\end{proof}

\begin{assessment}[V34 fundamental-chain regression]
\label{ass:V37-chain-regression}
The V34 fundamental chain remains untouched.  The new same-link proof is
applied only to the local junction \(J_{123,e}\), after the five-term
operator is assembled, and never to a raw first formation map.  The new
disjoint-bank theorem is applied to
\(\pi_{\boldsymbol S,\mu;\boldsymbol U,\nu}
\mathcal Q_{\boldsymbol S,\boldsymbol U}\), where
\(\mathcal Q\) already contains both Wilson defects and the one final
denominator.  It therefore cannot recreate the separated CFMI estimate
refuted by \eqref{eq:V34-CFMI-divergence}.
\end{assessment}

\begin{assessment}[What is and is not closed at ternary order]
\label{ass:V37-ternary-status}
The common-link junction sector is now uniformly closed by
\Cref{cor:V37-same-link-closed}.  The disjoint tree is uniformly closed
on the V35 private-dominant sector, the V36 bi-warm and
dimension-cold sectors, and the V37 variance-spread cold sector.  What
remains is the dimension-mesoscopic formation average isolated by
\eqref{eq:V36-mesoscopic-ATF} and
\eqref{eq:V37-effective-number-failure}: on a macroscopic cold bank it
has few effective links, while a globally cold nonmacroscopic bank still
requires a separate assembled-defect payment.  Hence uniform ternary
MTA is not yet proved.
\end{assessment}

\section{V37 ledger and V38 acquisition order}
\label{sec:V37-ledger}

\begin{theorem}[Integrated V37 statement]
\label{thm:V37-integrated}
Preserving V1--V36, V37 proves:
\begin{enumerate}[label=\textup{(V37.\arabic*)},leftmargin=5.5em]
\item the exact mean and variance
      \eqref{eq:V37-Casimir-mean}--%
      \eqref{eq:V37-Casimir-variance} of the tensor-dimension LR law;
\item its exact pushforward to the V36 cancellation coordinates and the
      lower-tail bound \eqref{eq:V37-one-link-cold-tail};
\item the product-bank concentration
      \eqref{eq:V37-product-bank-cold-tail} and MTA payment of every
      variance-spread cold sector in
      \eqref{eq:V37-variance-spread-MTA};
\item localization of a remaining macroscopic cold bank to the
      few-effective-link regime
      \eqref{eq:V37-effective-number-failure};
\item the basis-free partial-trace identity
      \eqref{eq:V37-partial-trace-assembly} and multiplicity-free
      diagonal/Racah bounds
      \eqref{eq:V37-diagonal-no-count}--%
      \eqref{eq:V37-Racah-no-count};
\item the operator-valued Gaussian excess
      \eqref{eq:V37-operator-one-link}; and
\item the uniform same-link resolvent estimate
      \eqref{eq:V37-one-link-MTA}, closing the ternary junction row.
\end{enumerate}
\end{theorem}

\begin{proof}
Items \textup{(V37.1)--(V37.2)} are
\Cref{thm:V37-Casimir-variance,cor:V37-cancellation-tail}.
Items \textup{(V37.3)--(V37.4)} are
\Cref{thm:V37-product-bank-tail,thm:V37-variance-spread-MTA,
cor:V37-concentrated-bank-localization}.  Item \textup{(V37.5)} is
\Cref{thm:V37-partial-trace-assembly,
cor:V37-anchored-partial-trace}.  Items
\textup{(V37.6)--(V37.7)} are
\Cref{thm:V37-operator-one-link,thm:V37-one-link-MTA,
cor:V37-same-link-closed}.
\end{proof}

\begingroup
\small
\begin{longtable}{>{\raggedright\arraybackslash}p{0.29\textwidth}
                  >{\raggedright\arraybackslash}p{0.15\textwidth}
                  >{\raggedright\arraybackslash}p{0.48\textwidth}}
\toprule
\textbf{V37 row}&\textbf{Status}&\textbf{Advance or obstruction}\\
\midrule
\endfirsthead
\toprule
\textbf{V37 row}&\textbf{Status}&\textbf{Advance or obstruction}\\
\midrule
\endhead
Tensor-dimension LR law
& \MGclosed
& The total Casimir spectral trace gives exact mean and variance, retaining
  actual multiplicity weights.\\[0.35em]
Variance-spread disjoint banks
& \MGclosed
& Chebyshev supplies the missing thermal factor whenever the cold bank has
  at least \(s_\alpha\Xi(\boldsymbol R_1)\) effective links.\\[0.35em]
Concentrated or nonmacroscopic mesoscopic banks
& \MGopen
& A remaining cold macroscopic bank must concentrate its cross-Casimir
  variance into \(o(s_\alpha\Xi(\boldsymbol R_1))\) effective links;
  a nonmacroscopic cold bank needs a different two-defect payment.\\[0.35em]
Same-link multiplicity sum
& \MGclosed
& Partial trace against the whole multiplicity operator removes raw Racah
  counts before absolute values.\\[0.35em]
Same-link uniform ATF
& \MGclosed
& Operator Gaussian excess plus the final low/high Casimir split pays the
  single resolvent in the full BFA norm.\\[0.35em]
Uniform ternary MTA
& \MGopen
& The concentrated and nonmacroscopic pieces of the disjoint mesoscopic
  average remain.\\[0.35em]
All-order marked-tree atlas
& \MGopen
& Higher replica forests remain downstream of uniform ternary closure.\\[0.35em]
Full Yang--Mills mass gap
& \MGopen
& WUFC/CPM, the nonlinear source packet, finite-edge margins, capacity,
  protected sectors, and the nontrivial OS continuum state remain
  downstream.\\
\bottomrule
\end{longtable}
\endgroup

\begin{openproblem}[V38 mandate: residual mesoscopic LR banks]
\label{op:V38}
\begin{enumerate}[label=\textup{(V38.\arabic*)},leftmargin=5.5em]
\item Start from the simultaneous residual conditions
      \eqref{eq:V36-residual-input-limits}--%
      \eqref{eq:V36-residual-output-limits}.  Split according to whether
      the globally cold intermediate comes from a pair bank whose
      Casimir mean is macroscopic relative to
      \(\Xi(\boldsymbol R_1)\).  Do not silently impose
      \eqref{eq:V37-macroscopic-bank}.
\item On the macroscopic branch use
      \eqref{eq:V37-effective-number-failure} to extract the heavy links
      carrying the cross-Casimir variance; do not replace their mass by
      the cardinality of the bank.  On the nonmacroscopic branch retain
      the full assembled two-defect quotient and seek a payment from the
      opposite warm bank.
\item On each extracted heavy link, use the exact
      \(\mathrm{SU}(3)\) LR interval or hive formula to estimate the
      dimension law in the cancellation coordinates \((a,b)\).  Retain
      the coefficient kernel \(\pi\), every copy index, both Wilson
      defects, and the one final denominator.
\item Separate strongly imbalanced input Casimirs, which already gain
      from \eqref{eq:V37-one-link-cold-tail}, from balanced heavy links,
      where a sharp dimension-weighted cancellation tail is required.
\item Test the resulting inequalities on the V34 fundamental chain and on
      the one-link dual-pair singlet.  A proposed pointwise formation
      estimate which loses either defect is inadmissible.
\item If both residual branches close, state uniform
      ternary MTA and begin the general connected replica-forest
      induction.  If it does not, return the first explicit
      \((R,Q;S,\mu;a,b)\) family saturating the exact dimension law and
      move upstream from that family.
\item Continue to preserve a single newest YM note.  The next mandatory
      audit of the then-current SM and NS notes is V40.
\end{enumerate}
\end{openproblem}

\begin{firewall}[End-of-V37 claim boundary]
V37 proves an exact dimension-weighted Casimir concentration theorem,
closes the variance-spread part of the disjoint mesoscopic sector, and
proves the complete same-link ternary MTA row uniformly over every Racah
multiplicity.  It does not prove the remaining concentrated or
nonmacroscopic mesoscopic averages, uniform ternary MTA, the all-order
marked-tree atlas, WUFC, CPM, the nonlinear Perron packet, unreduced
finite-edge margins, capacity, protected-sector floors, local-algebra
noncollapse, reflection positivity, or a nontrivial continuum OS limit.
The full Yang--Mills mass gap is not proved.
\end{firewall}

\section*{Version V37 append-only record}
\addcontentsline{toc}{section}{Version V37 append-only record}
The complete V36 source was retained before this continuation.  Its
SHA--256 digest was
\begin{center}\scriptsize\ttfamily
33ce4836e6d23752833c20b0604e79109c0f86e4697adebbf10e65e477e73d22.
\end{center}
No mathematical content of V36 was altered.  On the next iteration, retain
this full file, append before its unique active document-closing command,
increment the filename to \texttt{\_V38.tex}, and remove only the
superseded V37 file from this Yang--Mills note series.  The next scheduled
SM/NS audit is V40.

% =============================================================================
% END APPENDED VERSION V37 MATERIAL
% =============================================================================

% =============================================================================
% BEGIN APPENDED VERSION V38 MATERIAL
% =============================================================================

\clearpage
\part{V38 --- Assembled routing to a thermally cold macroscopic bank,
exponential dimension tails, and the thin-interior LR frontier}
\label{part:V38}

\section{Purpose and acquisition order}
\label{sec:V38-purpose}

V37 computed the exact Casimir variance of the tensor-dimension
Littlewood--Richardson law and closed all cold banks with sufficiently
large variance-effective size.  It left open both a globally cold
intermediate formed on a nonmacroscopic bank and a macroscopic cold bank
whose variance was concentrated on few effective links.

V38 first goes upstream to the complete two-defect quotient.  If a
globally cold intermediate lies on a nonmacroscopic central bank while
the opposite intermediate remains warm, quadratic damping of the
opposite defect pays the block pointwise.  If the cold intermediate lies
on a macroscopic bank but is thermally high in its own right, its own
multiplier pays the block.  Consequently every genuine witness contains
a macroscopic pair bank whose intermediate is bounded in thermal
Casimir.

For such a bank the V37 dimension law is a product law of bounded
independent Casimir variables.  Bernstein's inequality upgrades the
V37 inverse-effective-number estimate to an exponential estimate.  Any
surviving bank must therefore contain one link carrying at least a
logarithmic fraction of its total pair Casimir.  The resulting one-link
formation is thermally high-to-low.

The final sections prove the required squared one-link tail for every
thermally thick input pair and for every pair of boundary highest
weights.  The only representation-theoretic branch left for V39 is a
thin-interior \(\mathrm{SU}(3)\) Berenstein--Zelevinsky interval.  No
scheduled SM/NS audit occurs at V38; the next remains V40.

\section{Pointwise assembled routing}
\label{sec:V38-assembled-routing}

Retain the disjoint-overlap notation of V35 and put
\begin{align}
 X&:=\Xi(\boldsymbol R_1),
 &
 x_L&:=\Xi(\boldsymbol R_{1,L}),
 &
 x_R&:=\Xi(\boldsymbol R_{1,R}),
 &
 x_P&:=\Xi(\boldsymbol R_{1,P_1}).
 \label{eq:V38-central-bank-masses}
\end{align}
Thus \(X=x_L+x_R+x_P\).  Exact activity on the overlap banks gives
\begin{equation}
 H_L\ge x_L,\qquad H_R\ge x_R,
 \label{eq:V38-overlap-dominates-central}
\end{equation}
and also \(H_L\ge\Xi(\boldsymbol R_{2,L})\) and
\(H_R\ge\Xi(\boldsymbol R_{3,R})\).

\begin{theorem}[A cold nonmacroscopic bank is paid by an opposite warm
defect]
\label{thm:V38-nonmacroscopic-warm-payment}
Fix \(0<\varepsilon_0<1/4\) and \(\gamma>0\).  There is
\(M_{\varepsilon_0,\gamma}<\infty\) such that, if
\begin{equation}
 x_P\le\varepsilon_0X,\qquad
 x_L\le\varepsilon_0X,\qquad
 \Xi(\boldsymbol U)\ge\gamma X,\qquad
 s_\alpha X\ge M_{\varepsilon_0,\gamma},
 \label{eq:V38-nonmacro-warm-hypotheses}
\end{equation}
then every resolved block satisfies
\begin{equation}
 \boxed{
 \mathcal Q_{\boldsymbol S,\boldsymbol U}
 \le C_{\varepsilon_0,\gamma}
 \frac{H_LH_R}{X}.}
 \label{eq:V38-nonmacro-warm-payment}
\end{equation}
No formation probability is used.
\end{theorem}

\begin{proof}
The first two hypotheses imply
\[
 x_R=X-x_L-x_P\ge(1-2\varepsilon_0)X.
\]
Quadratic product damping
\eqref{eq:V32-balanced-second-moment} applied to the right central bank
and to the right intermediate gives
\begin{equation}
 q_{1,R}
 \le\frac C{(s_\alpha X)^2},
 \qquad
 q_{\alpha,\boldsymbol U}
 \le\frac {C_\gamma}{(s_\alpha X)^2}.
 \label{eq:V38-opposite-quadratic-damping}
\end{equation}
Therefore
\begin{equation}
 |\delta_R(\boldsymbol U)|
 \le q_{\alpha,\boldsymbol U}+q_{1,R}q_{3,R}
 \le\frac {C_\gamma}{(s_\alpha X)^2}.
 \label{eq:V38-opposite-defect-damping}
\end{equation}
Increase \(M_{\varepsilon_0,\gamma}\) so that
\(q_{\alpha,\boldsymbol U}\le1/2\).  Since
\(q_{\alpha,\boldsymbol T}
=q_Pq_{\alpha,\boldsymbol S}q_{\alpha,\boldsymbol U}\),
\begin{equation}
 1-q_{\alpha,\boldsymbol T}\ge\frac12.
 \label{eq:V38-opposite-final-gap}
\end{equation}

Split the output mass as in \eqref{eq:V35-output-split}.  For its left
piece, \eqref{eq:V35-binary-defect-mass} and
\eqref{eq:V38-overlap-dominates-central} give
\begin{align}
 q_P\Xi(\boldsymbol S)|\delta_L\delta_R|
 &\le CH_L
 \le\frac C{1-2\varepsilon_0}\frac{H_LH_R}{X}.
 \label{eq:V38-nonmacro-left-piece}
\end{align}
Fusion geometry \eqref{eq:V31-inside-mass-paid} gives
\(\Xi(\boldsymbol U)\le CH_R\).  Hence
\eqref{eq:V35-two-bank-defect-bounds} and
\eqref{eq:V38-opposite-defect-damping} yield
\begin{align}
 q_P\Xi(\boldsymbol U)|\delta_L\delta_R|
 &\le
 \frac{C_\gamma H_R\,s_\alpha H_L}
 {(s_\alpha X)^2}
 \le C_\gamma\frac{H_LH_R}{X}.
 \label{eq:V38-nonmacro-right-piece}
\end{align}

If \(P=\varnothing\), the private piece vanishes.  Otherwise
\eqref{eq:V31-balanced-product-damping} gives
\(s_\alpha\Xi_Pq_P\le C\).  Thus
\begin{align}
 \Xi_Pq_P|\delta_L\delta_R|
 &\le
 \frac C{s_\alpha}\,
 s_\alpha H_L\,
 \frac{C_\gamma}{(s_\alpha X)^2}\\
 &\le C_\gamma H_L
 \le C_{\varepsilon_0,\gamma}\frac{H_LH_R}{X}.
 \label{eq:V38-nonmacro-private-piece}
\end{align}
Use \eqref{eq:V38-opposite-final-gap} and add
\eqref{eq:V38-nonmacro-left-piece}--%
\eqref{eq:V38-nonmacro-private-piece}.
\end{proof}

\begin{theorem}[A thermally high intermediate on a macroscopic bank is
paid]
\label{thm:V38-thermal-high-intermediate}
Fix \(\gamma>0\).  There are \(M_\gamma,C_\gamma<\infty\) such that,
if
\begin{equation}
 x_L\ge\gamma X,\qquad
 s_\alpha X\ge M_\gamma,\qquad
 s_\alpha\Xi(\boldsymbol S)\ge M_\gamma,
 \label{eq:V38-high-intermediate-hypotheses}
\end{equation}
then every resolved block satisfies
\begin{equation}
 \boxed{
 \mathcal Q_{\boldsymbol S,\boldsymbol U}
 \le C_\gamma\frac{H_LH_R}{X}.}
 \label{eq:V38-high-intermediate-payment}
\end{equation}
The right intermediate is arbitrary.
\end{theorem}

\begin{proof}
Quadratic product damping gives
\begin{equation}
 q_{\alpha,\boldsymbol S}
 \le\frac C{[s_\alpha\Xi(\boldsymbol S)]^2},
 \qquad
 q_{1,L}\le\frac {C_\gamma}{(s_\alpha X)^2}.
 \label{eq:V38-left-high-damping}
\end{equation}
Consequently
\begin{equation}
 |\delta_L(\boldsymbol S)|
 \le q_{\alpha,\boldsymbol S}+q_{1,L}q_{2,L}.
 \label{eq:V38-left-defect-split}
\end{equation}
Choose \(M_\gamma\) so that
\(q_{\alpha,\boldsymbol S}\le1/2\).  Then
\begin{equation}
 1-q_{\alpha,\boldsymbol T}
 \ge1-q_{\alpha,\boldsymbol S}\ge\frac12.
 \label{eq:V38-left-high-final-gap}
\end{equation}

For the \(\boldsymbol S\)-mass piece use
\(|\delta_R|\le Cs_\alpha H_R\),
\(\Xi(\boldsymbol S)\le CH_L\), and
\eqref{eq:V38-left-high-damping}:
\begin{align}
 \Xi(\boldsymbol S)|\delta_L\delta_R|
 &\le
 \frac{CH_R}{s_\alpha\Xi(\boldsymbol S)}
 +\frac{C_\gamma H_LH_R}{s_\alpha X^2}\\
 &\le
 C_\gamma\frac{H_LH_R}{X}.
 \label{eq:V38-high-left-piece}
\end{align}
For the first term we used
\(H_L\ge x_L\ge\gamma X\); for the second we used
\(s_\alpha X\ge1\).

For the \(\boldsymbol U\)-mass piece,
\eqref{eq:V35-binary-defect-mass} and
\eqref{eq:V38-left-defect-split} give
\begin{equation}
 \Xi(\boldsymbol U)|\delta_L\delta_R|
 \le CH_R
 \le C_\gamma\frac{H_LH_R}{X}.
 \label{eq:V38-high-right-piece}
\end{equation}
Finally, if \(P\ne\varnothing\),
\eqref{eq:V31-balanced-product-damping},
\(|\delta_R|\le Cs_\alpha H_R\), and
\(|\delta_L|\le2\) give
\begin{equation}
 \Xi_Pq_P|\delta_L\delta_R|
 \le CH_R
 \le C_\gamma\frac{H_LH_R}{X}.
 \label{eq:V38-high-private-piece}
\end{equation}
The private term is zero when \(P=\varnothing\).  Insert the final gap
\eqref{eq:V38-left-high-final-gap} and add the pieces.
\end{proof}

\begin{proposition}[Every residual contains a thermally cold macroscopic
bank]
\label{prop:V38-macro-thermal-cold-localization}
If the disjoint-tree MTA fails, then after interchanging \(L,R\) and
passing to a subsequence there are fixed \(c_0,M_0>0\) such that
\begin{align}
 x_L&\ge c_0X,
 &
 \frac{\Xi(\boldsymbol S)}X&\longrightarrow0,
 &
 s_\alpha\Xi(\boldsymbol S)&\le M_0,
 \label{eq:V38-macro-cold-bank}\\
 M_L(\boldsymbol R_{1,L},\boldsymbol R_{2,L})
 &\ge c_0X,
 &
 \frac{N_{\rm eff,L}}
 {s_\alpha X}&\longrightarrow0.
 \label{eq:V38-macro-cold-effective}
\end{align}
The dimension-mesoscopic condition
\[
 r_L(\boldsymbol S)>(1+s_\alpha X)^{-1}
\]
also remains in force.
\end{proposition}

\begin{proof}
Use \Cref{prop:V36-disjoint-failure-localization}.  Since
\(x_P/X\to0\) and \(x_L+x_R+x_P=X\), at least one central overlap bank
has a fixed positive mass fraction.

Suppose first that the cold intermediate supplied by
\eqref{eq:V36-residual-output-limits} is \(\boldsymbol S\).  If
\(x_L/X\) has a positive lower bound, the first two claims in
\eqref{eq:V38-macro-cold-bank} hold.  Otherwise pass to a subsequence
with \(x_L/X\to0\), so \(x_R/X\to1\).  If
\(\Xi(\boldsymbol U)/X\) did not tend to zero, a further subsequence
would satisfy \(\Xi(\boldsymbol U)\ge\gamma X\) for fixed
\(\gamma>0\), and
\Cref{thm:V38-nonmacroscopic-warm-payment} would pay it.  Thus
\(\Xi(\boldsymbol U)/X\to0\), and the right bank is both macroscopic and
cold; interchange \(L,R\).  The case in which \(\boldsymbol U\) is the
initial cold intermediate is symmetric.

On the selected macroscopic bank,
\Cref{thm:V38-thermal-high-intermediate} excludes
\(s_\alpha\Xi(\boldsymbol S)\ge M_\gamma\); take \(M_0=M_\gamma\).
Exact activity gives
\(C_E(\boldsymbol R_{1,L})\ge c\,x_L\), so the pair mean \(M_L\)
is at least \(c_0X\) after changing the constant.  The last assertion in
\eqref{eq:V38-macro-cold-effective} is then
\Cref{cor:V37-concentrated-bank-localization}.  The dimension cutoff was
already preserved by \Cref{prop:V36-disjoint-failure-localization}.
\end{proof}

\section{Exponential concentration of the exact dimension law}
\label{sec:V38-exponential-dimension-tail}

\begin{lemma}[Scalar Bernstein lower tail]
\label{lem:V38-Bernstein}
Let \(Y_1,\ldots,Y_N\) be independent centered real random variables
with \(|Y_j|\le b\) and
\(\sum_j\mathbb E Y_j^2=\sigma^2\).  Then, for every \(t>0\),
\begin{equation}
 \mathbb P\left\{\sum_jY_j\le-t\right\}
 \le
 \exp\left\{-\frac{t^2}
 {2(\sigma^2+bt/3)}\right\}.
 \label{eq:V38-Bernstein}
\end{equation}
\end{lemma}

\begin{proof}
For \(0\le\lambda<3/b\), the elementary Taylor estimate
\[
 e^{-\lambda y}
 \le1-\lambda y+
 \frac{\lambda^2y^2}{2(1-\lambda b/3)}
 \qquad(|y|\le b)
\]
and \(1+u\le e^u\) give
\[
 \mathbb E e^{-\lambda Y_j}
 \le
 \exp\left\{
 \frac{\lambda^2\mathbb EY_j^2}
 {2(1-\lambda b/3)}
 \right\}.
\]
Independence, the exponential Markov inequality, and multiplication of
these bounds yield
\[
 \mathbb P\left\{\sum_jY_j\le-t\right\}
 \le
 \exp\left\{-\lambda t+
 \frac{\lambda^2\sigma^2}
 {2(1-\lambda b/3)}\right\}.
\]
Substitute
\(\lambda=t/(\sigma^2+bt/3)\) and simplify.
\end{proof}

For one product bank use the notation
\begin{align}
 z_e&:=C_2(R_e)+C_2(Q_e),
 &
 z_*&:=\max_{e\in L}z_e,
 \label{eq:V38-local-pair-Casimir}\\
 \mathcal N_{\rm B,L}
 &:=
 \frac{M_L^2}{V_L+M_Lz_*}.
 \label{eq:V38-Bernstein-effective-number}
\end{align}
The subscript \({\rm B}\) denotes Bernstein, not a representation.

\begin{theorem}[Exponential product-bank Casimir tail]
\label{thm:V38-exponential-bank-tail}
Let both input representations be nontrivial on every link of \(L\).
For every \(0\le\delta<1\),
\begin{equation}
 \boxed{
 \nu_{\boldsymbol R\boldsymbol Q}
 \{C_E(\boldsymbol S)\le\delta M_L\}
 \le
 \exp\{-c_{\rm B}(1-\delta)^2
              \mathcal N_{\rm B,L}\}.}
 \label{eq:V38-exponential-bank-tail}
\end{equation}
Moreover,
\begin{equation}
 \boxed{
 \frac45\frac{M_L}{z_*}
 \le\mathcal N_{\rm B,L}
 \le\frac{M_L}{z_*}.}
 \label{eq:V38-Bernstein-number-comparison}
\end{equation}
All constants are independent of the bank size and LR
multiplicities.
\end{theorem}

\begin{proof}
Under the exact product dimension law
\eqref{eq:V37-product-dimension-law}, put
\[
 Z_e:=C_2(S_e)-z_e.
\]
The \(Z_e\) are independent and centered by
\eqref{eq:V37-Casimir-mean}, while
\[
 \sum_e\mathbb E Z_e^2=\frac12V_L
\]
by \eqref{eq:V37-bank-Casimir-variance}.  The fusion estimate
\eqref{eq:V30-Casimir-fusion} gives
\[
 0\le C_2(S_e)\le C(1+z_e)\le C'z_e.
\]
The last absorption uses the positive Casimir floor for the two
nontrivial input representations.  Hence
\(|Z_e|\le B_Gz_e\le B_Gz_*\).

Apply \Cref{lem:V38-Bernstein} with
\[
 t=(1-\delta)M_L,\qquad
 \sigma^2=\frac12V_L,\qquad b=B_Gz_*.
\]
Absorbing the fixed group constants in the exponent proves
\eqref{eq:V38-exponential-bank-tail}.

Finally
\[
 V_L=\sum_ex_ey_e
 \le\frac14\sum_e(x_e+y_e)^2
 \le\frac14z_*M_L,
\]
where \(x_e=C_2(R_e)\) and \(y_e=C_2(Q_e)\).  Thus
\[
 M_Lz_*\le V_L+M_Lz_*
 \le\frac54M_Lz_*,
\]
which is equivalent to
\eqref{eq:V38-Bernstein-number-comparison}.
\end{proof}

\begin{theorem}[Logarithmically spread cold banks satisfy MTA]
\label{thm:V38-log-spread-MTA}
Fix \(0\le\delta<1\).  There is \(C_\delta<\infty\) such that the
portion of the actual disjoint-tree average on which
\begin{equation}
 C_E(\boldsymbol S)\le\delta M_L,
 \qquad
 \mathcal N_{\rm B,L}
 \ge C_\delta\log(2+s_\alpha X)
 \label{eq:V38-log-spread-sector}
\end{equation}
obeys
\begin{equation}
 \boxed{
 \sum_{\eqref{eq:V38-log-spread-sector}}
 \pi_{\boldsymbol S,\mu;\boldsymbol U,\nu}
 \mathcal Q_{\boldsymbol S,\boldsymbol U}
 \le C_\delta\frac{H_LH_R}{X}.}
 \label{eq:V38-log-spread-MTA}
\end{equation}
The symmetric right-bank statement also holds.
\end{theorem}

\begin{proof}
Choose \(C_\delta\) large enough that
\eqref{eq:V38-exponential-bank-tail} is at most
\((2+s_\alpha X)^{-2}\).  Since every Schur tax is at most one,
\eqref{eq:V36-kernel-product-envelope} implies
\[
 \sum_{\eqref{eq:V38-log-spread-sector}}
 \pi_{\boldsymbol S,\mu;\boldsymbol U,\nu}
 \le(2+s_\alpha X)^{-2}
 \le\frac1{s_\alpha X}.
\]
Multiply by the assembled pointwise estimate
\(\mathcal Q_{\boldsymbol S,\boldsymbol U}
\le Cs_\alpha H_LH_R\) from
\Cref{lem:V35-one-resolvent-two-defects}.
\end{proof}

\begin{corollary}[A residual bank has one logarithmically heavy link]
\label{cor:V38-heavy-link-localization}
On every failure sequence selected in
\Cref{prop:V38-macro-thermal-cold-localization}, there are links
\(e_\alpha\in L\) and a fixed \(c>0\) such that, after passing to a
subsequence,
\begin{align}
 \mathcal N_{\rm B,L}
 &\le C\log(2+s_\alpha X),
 \label{eq:V38-small-Bernstein-number}\\
 z_{e_\alpha}
 &\ge
 \frac{cX}{\log(2+s_\alpha X)},
 \label{eq:V38-heavy-link-mass}\\
 s_\alpha z_{e_\alpha}
 &\longrightarrow\infty,
 \label{eq:V38-heavy-link-thermal}\\
 C_2(S_{e_\alpha})
 &\le\frac{M_0}{s_\alpha},
 \qquad
 \frac{C_2(S_{e_\alpha})}{z_{e_\alpha}}
 \longrightarrow0.
 \label{eq:V38-heavy-link-cold-output}
\end{align}
\end{corollary}

\begin{proof}
By \eqref{eq:V38-macro-cold-bank},
\[
 C_E(\boldsymbol S)\le\Xi(\boldsymbol S)\le\frac{M_0}{s_\alpha},
\qquad
 M_L\ge c_0X.
\]
Since \(s_\alpha X\to\infty\), the event
\(C_E(\boldsymbol S)\le M_L/2\) contains the residual output
cofinally.  If \eqref{eq:V38-small-Bernstein-number} failed with a
sufficiently large constant, \Cref{thm:V38-log-spread-MTA} would pay the
sequence.  Hence that inequality holds.

Choose \(e_\alpha\) attaining \(z_*\).  Equations
\eqref{eq:V38-Bernstein-number-comparison},
\eqref{eq:V38-small-Bernstein-number}, and \(M_L\ge c_0X\) give
\eqref{eq:V38-heavy-link-mass}.  Multiplication by \(s_\alpha\) proves
\eqref{eq:V38-heavy-link-thermal}, because
\(s_\alpha X/\log(2+s_\alpha X)\to\infty\).
The first inequality in \eqref{eq:V38-heavy-link-cold-output} follows
from nonnegativity of all local output Casimirs.  Divide it by
\eqref{eq:V38-heavy-link-mass} to obtain the final limit.
\end{proof}

\begin{assessment}[What exponential concentration changes]
\label{ass:V38-exponential-progress}
The V37 variance tail forced few effective links only relative to
\(s_\alpha X\).  V38 improves this to a concrete atom:
one link carries at least \(X/\log(2+s_\alpha X)\) pair Casimir while
its output stays below \(M_0/s_\alpha\).  The unresolved object is now a
one-link thermally high-to-low \(\mathrm{SU}(3)\) LR tail.  A bankwide
tableau count would discard this localization and is therefore no longer
an admissible next step.
\end{assessment}

\section{One-link SU(3) dimension tails}
\label{sec:V38-one-link-tails}

For \(R=(p,q)\), define its height and boundary thickness by
\begin{equation}
 h_R:=1+p+q,
 \qquad
 \ell_R:=1+\min\{p,q\}.
 \label{eq:V38-height-thickness}
\end{equation}
The thickness is one exactly for a boundary highest weight.

\begin{lemma}[Dimension, reciprocity, and the rank-two multiplicity
interval]
\label{lem:V38-SU3-thickness}
There are group constants \(0<c<C<\infty\) such that
\begin{align}
 c h_R^2&\le1+C_2(R)\le Ch_R^2,
 \label{eq:V38-height-Casimir}\\
 c\ell_Rh_R^2&\le d_R\le C\ell_Rh_R^2.
 \label{eq:V38-thickness-dimension}
\end{align}
If \(S=(u,v)\) occurs in
\((p,q)\otimes(r,s)\), then
\begin{align}
 |h_R-h_Q|&\le h_S,
 \label{eq:V38-reciprocal-height}\\
 m_{RQ}^S
 &\le
 1+\min\{p,q,r,s,u,v\}
 =\min\{\ell_R,\ell_Q,\ell_S\}.
 \label{eq:V38-BZ-multiplicity-bound}
\end{align}
\end{lemma}

\begin{proof}
The Casimir and dimension formulae are
\begin{align*}
 C_2(p,q)&=\frac13(p^2+pq+q^2+3p+3q),\\
 d_{(p,q)}&=\frac12(p+1)(q+1)(p+q+2).
\end{align*}
They immediately give
\eqref{eq:V38-height-Casimir}--%
\eqref{eq:V38-thickness-dimension}.

Frobenius reciprocity makes the invariant multiplicity in
\(R\otimes Q\otimes S^*\) symmetric in its three factors.  In
particular \(R^*\) occurs in \(Q\otimes S^*\).  Apply
\eqref{eq:V30-height-fusion} and then interchange \(R,Q\); this proves
\eqref{eq:V38-reciprocal-height}.

For completeness, the \(\mathrm{SU}(3)\) Berenstein--Zelevinsky
polytope has one integer coordinate.  Its true triangles form one
integer interval.  For each of the six Dynkin labels of the three
weights \(R,Q,S^*\), two nonnegative outer entries have that label as
their sum, while the unique virtual-triangle step increases one and
decreases the other.  The interval can therefore contain at most that
label plus one integer points.  Permuting the three weights and using
that \(S^*=(v,u)\) proves all six bounds simultaneously, hence
\eqref{eq:V38-BZ-multiplicity-bound}.  This is also the immediate
rank-two consequence of the explicit one-sum BZ formula; see
\url{https://arxiv.org/abs/math-ph/0010051}, equation~(24).
\end{proof}

\begin{lemma}[Low-Casimir SU(3) dimension volume]
\label{lem:V38-low-Casimir-volume}
For every fixed \(L<\infty\) and \(0<s\le1\),
\begin{equation}
 \sum_{\substack{S=(u,v)\\1+C_2(S)\le L/s}}
 d_S
 \le C_Ls^{-5/2}.
 \label{eq:V38-low-Casimir-volume}
\end{equation}
\end{lemma}

\begin{proof}
Put \(H=s^{-1/2}\).  By \eqref{eq:V38-height-Casimir}, every summand has
\(u+v\le C_LH\).  For fixed \(n=u+v\),
\begin{align*}
 \sum_{u=0}^{n}d_{(u,n-u)}
 &=\frac{n+2}{2}
 \sum_{u=0}^{n}(u+1)(n-u+1)
 \le C(n+1)^4.
\end{align*}
Summing \(n\le C_LH\) gives \(C_LH^5=C_Ls^{-5/2}\).
\end{proof}

\begin{theorem}[Squared cold tail for a thermally thick input]
\label{thm:V38-thick-input-tail}
Fix \(L<\infty\), put
\begin{equation}
 \Theta:=s\bigl[(1+C_2(R))+(1+C_2(Q))\bigr],
 \qquad 0<s\le1,
 \label{eq:V38-local-thermal-height}
\end{equation}
and suppose
\begin{equation}
 \max\{\ell_R,\ell_Q\}\ge s^{-1/2}.
 \label{eq:V38-thermally-thick}
\end{equation}
Then the exact one-link dimension law satisfies
\begin{equation}
 \boxed{
 \sum_{\substack{S,\mu\\1+C_2(S)\le L/s}}
 \frac{d_S}{d_Rd_Q}
 \le\frac{C_L}{(1+\Theta)^2}.}
 \label{eq:V38-thick-input-tail}
\end{equation}
The stronger Schur-tax-weighted sum obeys the same estimate.
\end{theorem}

\begin{proof}
The taxed assertion follows from the untaxed one because
\(\tau\le1\).  If \(\Theta\le\Theta_L\) for a sufficiently large fixed
constant, the probability bound by one proves the result after enlarging
\(C_L\).

Suppose \(\Theta>\Theta_L\) and the summation set is nonempty.  Put
\(H=s^{-1/2}\).  Every output in the set has \(h_S\le C_LH\).
Equation \eqref{eq:V38-reciprocal-height} then gives
\(|h_R-h_Q|\le C_LH\).  Choosing \(\Theta_L\) large ensures that
\(h_R,h_Q\) are comparable.  Consequently
\begin{equation}
 d_Rd_Q
 \ge
 c\,\ell_R\ell_Q
 \bigl[(1+C_2(R))+(1+C_2(Q))\bigr]^2.
 \label{eq:V38-thick-input-denominator}
\end{equation}

Let
\(\ell_{\min}=\min\{\ell_R,\ell_Q\}\) and
\(\ell_{\max}=\max\{\ell_R,\ell_Q\}\).
By \eqref{eq:V38-BZ-multiplicity-bound} and
\Cref{lem:V38-low-Casimir-volume}, the numerator in
\eqref{eq:V38-thick-input-tail} is at most
\begin{equation}
 C_L\ell_{\min}H^5.
 \label{eq:V38-thick-input-numerator}
\end{equation}
Divide by \eqref{eq:V38-thick-input-denominator} and use
\(\ell_{\max}\ge H\):
\[
 \frac{C_LH^5}
 {\ell_{\max}
 [(1+C_2(R))+(1+C_2(Q))]^2}
 \le
 \frac{C_LH^4}
 {[(1+C_2(R))+(1+C_2(Q))]^2}
 =\frac{C_L}{\Theta^2}.
\]
This proves the high-\(\Theta\) case and hence the theorem.
\end{proof}

\subsection{All boundary--boundary pairs}

\begin{lemma}[Exact symmetric-tensor decompositions]
\label{lem:V38-boundary-decompositions}
For \(n,m\ge0\),
\begin{align}
 (n,0)\otimes(m,0)
 &=\bigoplus_{j=0}^{\min\{n,m\}}
   (n+m-2j,j),
 \label{eq:V38-same-boundary-decomposition}\\
 (n,0)\otimes(0,m)
 &=\bigoplus_{j=0}^{\min\{n,m\}}
   (n-j,m-j).
 \label{eq:V38-opposite-boundary-decomposition}
\end{align}
Every summand has multiplicity one.  The conjugate versions follow by
interchanging the Dynkin labels.
\end{lemma}

\begin{proof}
The first identity is the two-row Pieri rule.

For the second, identify \((n,0)=\operatorname{Sym}^n\mathbb C^3\)
and \((0,m)=\operatorname{Sym}^m(\mathbb C^3)^*\).  Contraction of one
upper and one lower index maps their tensor product surjectively onto
the corresponding product with \(n,m\) lowered by one.  To verify
surjectivity directly, use the bihomogeneous polynomial model.  If
\(\Delta=\sum_{a=1}^3\partial_{z_a}\partial_{w_a}\) and
\(\sigma=\sum_{a=1}^3z_aw_a\), then on bidegree
\((n-1,m-1)\)
\[
 \Delta(\sigma P)=(n+m+1)P+\sigma\Delta P.
\]
Repeatedly correcting the second term gives a finite right inverse,
because every application of \(\Delta\) lowers both degrees.

The kernel contains the irreducible module generated by the traceless
highest tensor, of highest weight \((n,m)\).  Moreover,
\begin{align*}
 &d_{(n,0)}d_{(0,m)}
 -d_{(n-1,0)}d_{(0,m-1)}\\
 &\qquad
 =\frac12(n+1)(m+1)(n+m+2)
 =d_{(n,m)},
\end{align*}
so surjectivity shows that the kernel has exactly that irreducible
dimension and hence equals it.  Induction on
\(\min\{n,m\}\) proves
\eqref{eq:V38-opposite-boundary-decomposition}.
\end{proof}

\begin{theorem}[Squared cold tail for every boundary pair]
\label{thm:V38-boundary-tail}
Let \(R,Q\) both have boundary highest weights, equivalently
\(\ell_R=\ell_Q=1\).  With \(\Theta\) as in
\eqref{eq:V38-local-thermal-height},
\begin{equation}
 \boxed{
 \sum_{\substack{S,\mu\\1+C_2(S)\le L/s}}
 \frac{d_S}{d_Rd_Q}
 \le\frac{C_L}{(1+\Theta)^2}.}
 \label{eq:V38-boundary-tail}
\end{equation}
Again the same estimate holds after inserting the Schur tax.
\end{theorem}

\begin{proof}
Conjugation reduces the two possible orientations to
\Cref{lem:V38-boundary-decompositions}.  In the same-boundary case,
every output in \eqref{eq:V38-same-boundary-decomposition} has Dynkin
height
\[
 n+m-j\ge\max\{n,m\}.
\]
Thus the low-output set is empty when \(\Theta\) exceeds a constant
depending only on \(L\); the bounded-\(\Theta\) case follows from total
probability one.

Consider opposite boundaries and suppose \(n\ge m\).  Reparameterize
\eqref{eq:V38-opposite-boundary-decomposition} by \(k=m-j\):
\begin{equation}
 S_k=(n-m+k,k),\qquad0\le k\le m.
 \label{eq:V38-opposite-boundary-shells}
\end{equation}
If \(1+C_2(S_k)\le L/s\), then
\[
 n-m+2k\le C_LH,\qquad H=s^{-1/2}.
\]
For large \(\Theta\), this also forces \(n,m\) to be comparable.
Using the exact dimension formula,
\begin{align}
 \sum_{\substack{k\\1+C_2(S_k)\le L/s}}d_{S_k}
 &\le
 C\sum_{0\le k\le C_LH}
 (n-m+k+1)(k+1)(n-m+2k+2)\\
 &\le C_LH^4.
 \label{eq:V38-boundary-low-dimension-sum}
\end{align}
Meanwhile
\[
 d_{(n,0)}d_{(0,m)}
 \ge c(n^2+m^2)^2
 \ge c'\bigl[(1+C_2(R))+(1+C_2(Q))\bigr]^2.
\]
Division gives \(C_LH^4/(A_R+A_Q)^2=C_L/\Theta^2\).
As before, total probability treats bounded \(\Theta\), and
\(\tau\le1\) gives the taxed version.
\end{proof}

\begin{assessment}[The dual symmetric family is not a witness]
\label{ass:V38-dual-symmetric-benchmark}
For \(R=(n,0)\), \(Q=(0,n)\),
\[
 R\otimes Q=\bigoplus_{k=0}^{n}(k,k),
 \quad
 d_R=\frac{(n+1)(n+2)}2,
 \quad
 d_{(k,k)}=(k+1)^3.
\]
Hence the exact cold dimension mass is
\[
 \frac1{d_R^2}
 \sum_{\substack{k\\1+k^2+2k\le L/s}}(k+1)^3
 \le\frac{C_L}{(1+s n^2)^2}.
\]
The singlet is included, with its exact mass \(d_R^{-2}\).
Thus the most obvious boundary cancellation family gains two thermal
powers and cannot witness failure of the required one-power MTA tail.
\end{assessment}

\subsection{Insertion at the logarithmically heavy link}

\begin{proposition}[Ultra-imbalanced heavy links are paid]
\label{prop:V38-imbalanced-heavy-link}
On the heavy link \(e_\alpha\) of
\Cref{cor:V38-heavy-link-localization}, abbreviate its two input
Casimirs by \(c_1,c_2\).  If
\begin{equation}
 \frac{\min\{c_1,c_2\}}{\max\{c_1,c_2\}}
 \le\frac C{s_\alpha X},
 \label{eq:V38-ultra-imbalanced-condition}
\end{equation}
then the corresponding portion of the disjoint formation average
satisfies MTA uniformly.
\end{proposition}

\begin{proof}
By \eqref{eq:V38-heavy-link-cold-output} and
\eqref{eq:V38-heavy-link-thermal}, cofinally
\[
 C_2(S_{e_\alpha})\le\frac12(c_1+c_2).
\]
Apply \eqref{eq:V37-one-link-cold-tail} with \(\delta=1/2\).  The
dimension-law mass of all such local outputs is at most
\[
 C\frac{c_1c_2}{(c_1+c_2)^2}
 \le C\frac{\min\{c_1,c_2\}}{\max\{c_1,c_2\}}
 \le\frac C{s_\alpha X}.
\]
The product dimension law factorizes over links, and
\(\pi\le w_Lw_R\le\nu_L\nu_R\).  Thus the same bound holds for the
actual global kernel mass.  Multiply by
\(\mathcal Q\le Cs_\alpha H_LH_R\).
\end{proof}

\begin{theorem}[Thick and boundary heavy-link sectors satisfy MTA]
\label{thm:V38-heavy-link-tail-MTA}
Let \(e_\alpha\) be selected by
\Cref{cor:V38-heavy-link-localization}.  If either
\begin{equation}
 \max\{\ell_{R_{1,e_\alpha}},\ell_{R_{2,e_\alpha}}\}
 \ge s_\alpha^{-1/2}
 \label{eq:V38-heavy-link-thick-case}
\end{equation}
or both local inputs have boundary highest weights, then the
corresponding portion of the actual disjoint-tree average obeys
\begin{equation}
 \boxed{
 \sum\pi_{\boldsymbol S,\mu;\boldsymbol U,\nu}
 \mathcal Q_{\boldsymbol S,\boldsymbol U}
 \le C\frac{H_LH_R}{X}.}
 \label{eq:V38-heavy-link-tail-MTA}
\end{equation}
\end{theorem}

\begin{proof}
Put
\[
 \Theta_e=s_\alpha\bigl[
 (1+C_2(R_{1,e_\alpha}))
 +(1+C_2(R_{2,e_\alpha}))\bigr].
\]
Equations \eqref{eq:V38-heavy-link-mass}--%
\eqref{eq:V38-heavy-link-cold-output} give
\[
 \Theta_e\ge
 \frac{c\,s_\alpha X}{\log(2+s_\alpha X)},
 \qquad
 1+C_2(S_{e_\alpha})\le\frac{L_0}{s_\alpha}
\]
for fixed \(L_0\).  Apply
\Cref{thm:V38-thick-input-tail} in
\eqref{eq:V38-heavy-link-thick-case}, and
\Cref{thm:V38-boundary-tail} in the boundary case.  The local marginal
of the exact product dimension law is bounded by
\begin{equation}
 \frac C{(1+\Theta_e)^2}
 \le
 C\frac{\log^2(2+s_\alpha X)}
          {(s_\alpha X)^2}
 \le\frac C{s_\alpha X}
 \label{eq:V38-heavy-local-probability}
\end{equation}
cofinally.  Since
\(\pi\le\nu_L\nu_R\), this also bounds the actual global kernel mass.
Multiplication by the complete two-defect estimate
\(\mathcal Q\le Cs_\alpha H_LH_R\) proves
\eqref{eq:V38-heavy-link-tail-MTA}.
\end{proof}

\begin{corollary}[Exact post-V38 representation-theoretic residual]
\label{cor:V38-thin-interior-residual}
Every still-unpaid disjoint-tree witness has a selected heavy link
\(e_\alpha\) satisfying
\eqref{eq:V38-heavy-link-mass}--%
\eqref{eq:V38-heavy-link-cold-output} and, in addition,
\begin{align}
 1<\max\{\ell_{R_{1,e_\alpha}},
          \ell_{R_{2,e_\alpha}}\}
 &<s_\alpha^{-1/2},
 \label{eq:V38-thin-interior-thickness}\\
 \frac{\min\{C_2(R_{1,e_\alpha}),C_2(R_{2,e_\alpha})\}}
 {\max\{C_2(R_{1,e_\alpha}),C_2(R_{2,e_\alpha})\}}
 &>\frac c{s_\alpha X}.
 \label{eq:V38-not-ultra-imbalanced}
\end{align}
At least one local input is therefore interior but thinner than the
thermal Dynkin scale.  This is the one-dimensional BZ-interval sector
left for V39.
\end{corollary}

\begin{proof}
\Cref{thm:V38-heavy-link-tail-MTA} excludes thermally thick pairs and
pairs in which both weights lie on the boundary.  Therefore the maximum
thickness is strictly between one and \(s_\alpha^{-1/2}\).
\Cref{prop:V38-imbalanced-heavy-link} gives
\eqref{eq:V38-not-ultra-imbalanced}.
\end{proof}

\section{Regression cards and V38 ledger}
\label{sec:V38-ledger}

\begin{assessment}[V34 chain and V37 same-link regressions]
\label{ass:V38-regressions}
Every V38 disjoint estimate is applied either pointwise to
\(\mathcal Q_{\boldsymbol S,\boldsymbol U}\) or to its average against
the actual kernel \(\pi\).  Both binary defects and the single final
denominator are already present.  No estimate is applied to the raw
first formation map, so the V34 divergence
\eqref{eq:V34-CFMI-divergence} is not recreated.

The one-link dual-pair singlet appears in
\Cref{thm:V38-boundary-tail} with exact dimension-law mass
\(d_R^{-2}\), while its same-link diagonal cumulant still vanishes by
\Cref{prop:V37-dual-pair-regression}.  Thus V38 neither loses the Schur
singlet tax nor alters the already closed same-link row.
\end{assessment}

\begin{theorem}[Integrated V38 statement]
\label{thm:V38-integrated}
Preserving V1--V37, V38 proves:
\begin{enumerate}[label=\textup{(V38.\arabic*)},leftmargin=5.5em]
\item pointwise MTA payment of a globally cold nonmacroscopic bank when
      the opposite intermediate is warm,
      \eqref{eq:V38-nonmacro-warm-payment};
\item pointwise payment of every thermally high intermediate formed on
      a macroscopic central bank,
      \eqref{eq:V38-high-intermediate-payment};
\item localization of every residual to a thermally cold macroscopic
      bank, \eqref{eq:V38-macro-cold-bank};
\item the exponential exact-dimension tail
      \eqref{eq:V38-exponential-bank-tail} and MTA on every
      logarithmically spread bank;
\item extraction of a thermally high input link with locally cold output,
      \eqref{eq:V38-heavy-link-mass}--%
      \eqref{eq:V38-heavy-link-cold-output};
\item the BZ thickness bound
      \eqref{eq:V38-BZ-multiplicity-bound} and squared cold tail
      \eqref{eq:V38-thick-input-tail} for every thermally thick input;
\item the exact boundary decompositions and squared tail
      \eqref{eq:V38-boundary-tail}; and
\item MTA closure of the thick, boundary--boundary, and
      ultra-imbalanced heavy-link sectors, leaving precisely
      \eqref{eq:V38-thin-interior-thickness}--%
      \eqref{eq:V38-not-ultra-imbalanced}.
\end{enumerate}
\end{theorem}

\begin{proof}
Items \textup{(V38.1)--(V38.3)} are
\Cref{thm:V38-nonmacroscopic-warm-payment,
thm:V38-thermal-high-intermediate,
prop:V38-macro-thermal-cold-localization}.  Items
\textup{(V38.4)--(V38.5)} are
\Cref{thm:V38-exponential-bank-tail,thm:V38-log-spread-MTA,
cor:V38-heavy-link-localization}.  Items
\textup{(V38.6)--(V38.7)} are
\Cref{lem:V38-SU3-thickness,thm:V38-thick-input-tail,
lem:V38-boundary-decompositions,thm:V38-boundary-tail}.
Item \textup{(V38.8)} is
\Cref{prop:V38-imbalanced-heavy-link,
thm:V38-heavy-link-tail-MTA,cor:V38-thin-interior-residual}.
\end{proof}

\begingroup
\small
\begin{longtable}{>{\raggedright\arraybackslash}p{0.29\textwidth}
                  >{\raggedright\arraybackslash}p{0.15\textwidth}
                  >{\raggedright\arraybackslash}p{0.48\textwidth}}
\toprule
\textbf{V38 row}&\textbf{Status}&\textbf{Advance or obstruction}\\
\midrule
\endfirsthead
\toprule
\textbf{V38 row}&\textbf{Status}&\textbf{Advance or obstruction}\\
\midrule
\endhead
Nonmacroscopic cold bank
& \MGclosed
& If the opposite intermediate is warm, its quadratically damped defect
  pays the complete block; otherwise the opposite macroscopic bank is
  itself cold.\\[0.35em]
Macroscopic thermally high intermediate
& \MGclosed
& Its own product multiplier and the central-bank multiplier pay all
  three final-mass pieces pointwise.\\[0.35em]
Product dimension concentration
& \MGclosed
& Exact Casimir variance plus Bernstein gives an exponential lower tail,
  closing every bank spread over more than logarithmically many mass
  units.\\[0.35em]
Heavy-link thick and boundary sectors
& \MGclosed
& The one-link cold probability gains two thermal powers, which absorbs
  the logarithmic extraction loss.\\[0.35em]
Thin-interior BZ interval
& \MGopen
& Both input boundary thicknesses are below \(s_\alpha^{-1/2}\), at least
  one is nonboundary, and the output is thermally bounded despite a
  thermally diverging pair input.\\[0.35em]
Same-link ternary junction
& \MGclosed
& The V37 multiplicity-free partial-trace proof is unchanged.\\[0.35em]
Uniform ternary MTA
& \MGopen
& It is reduced to the thin-interior one-link tail.\\[0.35em]
All-order marked-tree atlas
& \MGopen
& Higher connected replica forests remain downstream of ternary closure.\\[0.35em]
Full Yang--Mills mass gap
& \MGopen
& WUFC/CPM, the nonlinear source packet, finite-edge margins, capacity,
  protected sectors, and the nontrivial OS continuum state remain
  downstream.\\
\bottomrule
\end{longtable}
\endgroup

\begin{openproblem}[V39 mandate: the thin-interior BZ interval]
\label{op:V39}
\begin{enumerate}[label=\textup{(V39.\arabic*)},leftmargin=5.5em]
\item On the heavy link of
      \Cref{cor:V38-thin-interior-residual}, write the complete
      one-dimensional \(\mathrm{SU}(3)\) BZ interval with all six Dynkin
      labels and the V36 cancellation coordinates.  Retain its lower and
      upper endpoints; the length bound alone is insufficient in the thin
      regime.
\item Prove the target local tail
      \[
      \sum_{\substack{S,\mu\\
       1+C_2(S)\le L/s_\alpha}}
      \frac{d_S}{d_Rd_Q}
      \le\frac C{[1+s_\alpha(A_R+A_Q)]^2}
      \]
      for the thin-interior sector, or return an explicit
      \((p,q),(r,s);(u,v)\) family violating it.
\item Split the BZ interval by the two input thicknesses and by distance
      from the six Horn faces.  Sum actual interval lengths and output
      dimensions; do not multiply a worst multiplicity by the full
      low-Casimir lattice count.
\item Preserve the V34 chain and dual-pair regressions.  Insert any local
      tail only into the actual product dimension law and then the
      complete two-defect quotient.
\item If the local tail holds, combine
      \Cref{prop:V38-macro-thermal-cold-localization,
      cor:V38-heavy-link-localization,
      thm:V38-heavy-link-tail-MTA} to prove uniform ternary MTA.  Then
      begin the general connected replica-forest induction.
\item If it fails, use the exact saturating family to replace the squared
      target by the weakest sufficient coefficient-dependent estimate;
      go upstream before introducing any raw multiplicity count.
\item Preserve only the newest YM file.  V40 must begin with the mandatory
      audit of the then-current SM and NS notes.
\end{enumerate}
\end{openproblem}

\begin{firewall}[End-of-V38 claim boundary]
V38 proves pointwise routing to a thermally cold macroscopic bank,
exponential dimension-law concentration, heavy-link extraction, and the
required squared one-link tail for thermally thick and all
boundary--boundary inputs.  It does not prove the thin-interior BZ tail,
uniform ternary MTA, the all-order marked-tree atlas, WUFC, CPM, the
nonlinear Perron packet, unreduced finite-edge margins, capacity,
protected-sector floors, local-algebra noncollapse, reflection
positivity, or a nontrivial continuum OS limit.  The full Yang--Mills
mass gap is not proved.
\end{firewall}

\section*{Version V38 append-only record}
\addcontentsline{toc}{section}{Version V38 append-only record}
The complete V37 source was retained before this continuation.  Its
SHA--256 digest was
\begin{center}\scriptsize\ttfamily
d9ca32bc5d02dae596094bd2659daf6b92546c337d01b79039ab3210353b569b.
\end{center}
No mathematical content of V37 was altered.  On the next iteration, retain
this full file, append before its unique active document-closing command,
increment the filename to \texttt{\_V39.tex}, and remove only the
superseded V38 file from this Yang--Mills note series.  The next scheduled
SM/NS audit is V40.

% =============================================================================
% END APPENDED VERSION V38 MATERIAL
% =============================================================================

% =============================================================================
% BEGIN APPENDED VERSION V39 MATERIAL
% =============================================================================

\clearpage
\part{V39 --- Exact BZ strip summation, global ternary MTA,
and the connected replica-partition formula}
\label{part:V39}

\section{Purpose and nonduplication}
\label{sec:V39-purpose}

V38 reduced every possible disjoint-tree failure to one link whose two
inputs are thermally high, whose output is thermally bounded, and whose
two input boundary thicknesses are both below the thermal Dynkin scale.
It proved the required squared dimension-law tail whenever either input
was thermally thick or both inputs lay on the boundary.

V39 writes the complete rank-two Berenstein--Zelevinsky interval in the
cancellation coordinates already introduced in V36.  Two of its upper
faces constrain \(d=a-b\) to an interval of length \(q+r\).  After the
long Dynkin directions are oriented oppositely, this becomes a diagonal
strip constraint on the cold output labels.  The number of output labels
and the interval multiplicity must be summed together: their product is
exactly the product of the two input boundary thicknesses.  This closes
the thin-interior tail and proves the local squared estimate for every
pair of \(\mathrm{SU}(3)\) irreducibles.

Insertion at the V38 heavy link closes the last disjoint-tree average.
Together with the V37 same-link theorem and the previously closed
thermal/private sectors, this proves the global \(m=3\) instance of the
marked-tree atlas.  The final part begins the all-order problem by proving
the exact connected replica-partition formula and extracting a canonical
marked spanning tree from every term.  The analytic all-order weight of
the fibres over a tree remains open.  The mandatory SM/NS audit is next,
at V40.

\section{The complete SU(3) BZ interval in cancellation variables}
\label{sec:V39-BZ-interval}

Let
\[
 R=(p,q),\qquad Q=(r,s),\qquad S=(u,v),
\]
and define \(a,b\) by
\eqref{eq:V36-SU3-cancellation-coordinates}.  When \(a,b\) are
nonnegative integers, put
\begin{align}
 L_{\rm BZ}
 &:=\max\{0,\ b-s,\ a-p\},
 \label{eq:V39-BZ-lower}\\
 U_{\rm BZ}
 &:=\min\{r,\ q,\ b,\ a,\ r-a+b,\ q+a-b\}.
 \label{eq:V39-BZ-upper}
\end{align}

\begin{theorem}[Exact rank-two BZ interval]
\label{thm:V39-exact-BZ-interval}
The Littlewood--Richardson multiplicity is
\begin{equation}
 \boxed{
 m_{(p,q),(r,s)}^{(u,v)}
 =
 \begin{cases}
 \max\{0,U_{\rm BZ}-L_{\rm BZ}+1\},
 &a,b\in\mathbb Z_{\ge0},\ u,v\in\mathbb Z_{\ge0},\\
 0,&\text{otherwise}.
 \end{cases}}
 \label{eq:V39-exact-BZ-multiplicity}
\end{equation}
In particular, the copy index is one integer
\(\eta\in[L_{\rm BZ},U_{\rm BZ}]\cap\mathbb Z\); no further tableau
count is hidden.
\end{theorem}

\begin{proof}
Use the BZ one-sum formula for the invariant multiplicity in
\(R\otimes Q\otimes S^*\).  For an ordinary Dynkin pair
\(\lambda=(\lambda_1,\lambda_2)\), its dual Dynkin labels are
\[
 \lambda^1=\frac{2\lambda_1+\lambda_2}{3},
 \qquad
 \lambda^2=\frac{\lambda_1+2\lambda_2}{3}.
\]
With \(\lambda=R\), \(\mu=Q\), and
\(\nu=S^*=(v,u)\), the exact BZ endpoints are
\begin{align*}
 L&=\max\{0,\,
 \lambda^2+\mu^1-\mu^2-\nu^1,\,
 -\lambda^1+\lambda^2+\mu^1-\nu^2\},\\
 U&=\min\{\mu_1,\lambda_2,\,
 \lambda^2+\mu^2-\nu^1,\,
 \lambda^1+\mu^1-\nu^2,\\
 &\hspace{7.2em}
 -\lambda^1+\lambda^2+\mu^1-\nu^1+\nu^2,\,
 \lambda^2+\mu^1-\mu^2+\nu^1-\nu^2\}.
 \end{align*}
This is the one-dimensional specialization of the BZ triangle theorem
recorded in \Cref{lem:V38-SU3-thickness}.

Substitute the ordinary labels and use
\[
 u=p+r-2a+b,\qquad v=q+s+a-2b.
\]
The two nonzero lower faces become \(b-s\) and \(a-p\).
The six upper faces become, in order,
\[
 r,\quad q,\quad b,\quad a,\quad r-a+b,\quad q+a-b.
\]
Thus the BZ sum contains one term for each integer between
\eqref{eq:V39-BZ-lower} and
\eqref{eq:V39-BZ-upper}, proving
\eqref{eq:V39-exact-BZ-multiplicity}.  The integrality requirement is
equivalent to the root-lattice/triality condition already displayed in
\eqref{eq:V36-SU3-triality}.
\end{proof}

\begin{corollary}[Exact diagonal strip and interval-length bounds]
\label{cor:V39-BZ-strip}
If \(m_{RQ}^S>0\), then, with \(d:=a-b\),
\begin{align}
 -q\le d\le r,
 \label{eq:V39-d-strip}\\
 p+r-q-s-3r
 \le u-v\le
 p+r-q-s+3q,
 \label{eq:V39-output-diagonal-strip}\\
 m_{RQ}^S
 \le1+\min\{q,r\}.
 \label{eq:V39-oriented-multiplicity}
\end{align}
\end{corollary}

\begin{proof}
Nonzero multiplicity means
\(0\le L_{\rm BZ}\le U_{\rm BZ}\).  Since
\(r-a+b=r-d\) and \(q+a-b=q+d\) are two entries in the minimum
\eqref{eq:V39-BZ-upper}, both are nonnegative.  This proves
\eqref{eq:V39-d-strip}.  Directly from
\eqref{eq:V36-SU3-cancellation-coordinates},
\begin{equation}
 3d=p+r-q-s-u+v.
 \label{eq:V39-d-output-difference}
\end{equation}
Insert \eqref{eq:V39-d-strip} to obtain
\eqref{eq:V39-output-diagonal-strip}.  Finally
\(U_{\rm BZ}\le\min\{q,r\}\) and \(L_{\rm BZ}\ge0\), which proves
\eqref{eq:V39-oriented-multiplicity}.
\end{proof}

\section{The thin-interior strip sum}
\label{sec:V39-thin-tail}

Use a local small parameter \(t\in(0,1]\), to avoid confusion with the
second Dynkin label \(s\), and put
\begin{equation}
 H:=t^{-1/2},
 \qquad
 \Theta:=t\bigl[(1+C_2(R))+(1+C_2(Q))\bigr].
 \label{eq:V39-local-scales}
\end{equation}

\begin{lemma}[Cold thin inputs have opposite long directions]
\label{lem:V39-opposite-orientation}
Fix \(L<\infty\).  Suppose
\begin{equation}
 \max\{\ell_R,\ell_Q\}<H
 \label{eq:V39-thin-hypothesis}
\end{equation}
and some occurring \(S\) satisfies
\begin{equation}
 1+C_2(S)\le L/t.
 \label{eq:V39-local-cold-output}
\end{equation}
After simultaneously conjugating all three representations, arrange
\(p\ge q\).  If \(\Theta\) exceeds a constant depending only on \(L\),
then necessarily \(s\ge r\).  Thus
\begin{equation}
 \ell_R=q+1,\qquad \ell_Q=r+1.
 \label{eq:V39-opposite-thicknesses}
\end{equation}
\end{lemma}

\begin{proof}
By \eqref{eq:V38-height-Casimir},
\eqref{eq:V39-local-cold-output} gives
\[
 u+v\le C_LH.
\]
Simultaneous conjugation preserves dimensions, Casimirs, and
multiplicities, and permits \(p\ge q\).  Hence
\(q+1=\ell_R<H\).

Suppose for contradiction that \(r\ge s\), so
\(s+1=\ell_Q<H\).  Since \(m_{RQ}^S>0\),
\[
 b-s\le L_{\rm BZ}\le U_{\rm BZ}\le q.
\]
Using the definition of \(b\), this inequality is
\[
 p+r\le q+s+u+2v.
\]
The right side is at most \(C_LH\).  Therefore
\(p,q,r,s\le C_LH\), and
\eqref{eq:V38-height-Casimir} gives
\(\Theta\le C_L\), contrary to the high-\(\Theta\) hypothesis.
Thus \(s\ge r\), proving
\eqref{eq:V39-opposite-thicknesses}.
\end{proof}

\begin{lemma}[Joint strip-volume bound]
\label{lem:V39-joint-strip-volume}
Under the high-\(\Theta\) hypotheses of
\Cref{lem:V39-opposite-orientation},
\begin{equation}
 \boxed{
 \sum_{\substack{S,\mu\\1+C_2(S)\le L/t}}
 d_S
 \le
 C_L\,\ell_R\ell_Q\,H^4.}
 \label{eq:V39-joint-strip-volume}
\end{equation}
The sum is over resolved copies, so the LR multiplicity is included.
\end{lemma}

\begin{proof}
Every output in the sum has \(u+v\le C_LH\).  By
\eqref{eq:V39-output-diagonal-strip}, its difference \(u-v\) lies in
an interval of length at most \(3(q+r)\).  For each allowed integer
\(d\in[-q,r]\), equation \eqref{eq:V39-d-output-difference} fixes
\(u-v\); there are at most \(C_LH\) nonnegative pairs \((u,v)\) with
that difference and \(u+v\le C_LH\).  Hence the number of possible
output labels is at most
\begin{equation}
 C_LH(q+r+1).
 \label{eq:V39-strip-lattice-count}
\end{equation}

Each such output has \(d_S\le C_LH^3\) by
\eqref{eq:V38-thickness-dimension}, and
\eqref{eq:V39-oriented-multiplicity} bounds its number of copies by
\(1+\min\{q,r\}=\min\{\ell_R,\ell_Q\}\).
Therefore the left side of
\eqref{eq:V39-joint-strip-volume} is at most
\[
 C_LH^4(q+r+1)\min\{\ell_R,\ell_Q\}.
\]
Since
\[
 q+r+1\le\ell_R+\ell_Q,
 \qquad
 \min\{\ell_R,\ell_Q\}(\ell_R+\ell_Q)
 \le2\ell_R\ell_Q,
\]
the claimed bound follows.
\end{proof}

\begin{theorem}[Squared cold tail in the thin-interior sector]
\label{thm:V39-thin-tail}
Fix \(L<\infty\).  Under
\eqref{eq:V39-thin-hypothesis}, the exact one-link dimension law obeys
\begin{equation}
 \boxed{
 \sum_{\substack{S,\mu\\1+C_2(S)\le L/t}}
 \frac{d_S}{d_Rd_Q}
 \le\frac{C_L}{(1+\Theta)^2}.}
 \label{eq:V39-thin-tail}
\end{equation}
The same bound remains true after insertion of the Schur tax.
\end{theorem}

\begin{proof}
The taxed sum is smaller.  For bounded \(\Theta\), total probability
one proves the result after increasing \(C_L\).  Suppose \(\Theta\) is
larger than the threshold in
\Cref{lem:V39-opposite-orientation}.  If the cold set is empty there is
nothing to prove.

Otherwise \eqref{eq:V38-reciprocal-height} and
\(h_S\le C_LH\) show that \(h_R,h_Q\) are comparable; the
high-\(\Theta\) assumption makes both larger than a fixed multiple of
\(H\).  Equations
\eqref{eq:V38-thickness-dimension} and
\eqref{eq:V39-opposite-thicknesses} give
\begin{equation}
 d_Rd_Q
 \ge
 c\,\ell_R\ell_Q
 \bigl[(1+C_2(R))+(1+C_2(Q))\bigr]^2.
 \label{eq:V39-thin-input-denominator}
\end{equation}
Divide \eqref{eq:V39-joint-strip-volume} by
\eqref{eq:V39-thin-input-denominator}:
\[
 \frac{C_LH^4}
 {[(1+C_2(R))+(1+C_2(Q))]^2}
 =\frac{C_L}{\Theta^2}.
\]
This proves the high-\(\Theta\) case and hence
\eqref{eq:V39-thin-tail}.
\end{proof}

\begin{corollary}[Universal one-link squared cold tail]
\label{cor:V39-universal-one-link-tail}
For all irreducible \(R,Q\), all \(0<t\le1\), and every fixed
\(L<\infty\),
\begin{equation}
 \boxed{
 \sum_{\substack{S,\mu\\1+C_2(S)\le L/t}}
 \frac{d_S}{d_Rd_Q}
 \le
 \frac{C_L}
 {\left(1+t[(1+C_2(R))+(1+C_2(Q))]\right)^2}.}
 \label{eq:V39-universal-one-link-tail}
\end{equation}
\end{corollary}

\begin{proof}
If
\(\max\{\ell_R,\ell_Q\}\ge t^{-1/2}\), use
\Cref{thm:V38-thick-input-tail}.  Otherwise use
\Cref{thm:V39-thin-tail}.
\end{proof}

\begin{proposition}[A strip family saturates the thickness bookkeeping]
\label{prop:V39-strip-regression}
Take \(R=(n,k)\), \(Q=(k,n)\), and an output \(S=(u,v)\) with
\(u+2v\equiv0\pmod3\).  Put
\[
 d=\frac{v-u}{3}.
\]
For \(n\) larger than all displayed low labels, the exact multiplicity is
\begin{equation}
 m_{(n,k),(k,n)}^{(u,v)}
 =
 \left[
 1+k-\frac{|u-v|}{3}
 -\max\left\{
 0,\,
 k-\frac{2u+v}{3},\,
 k-\frac{u+2v}{3}
 \right\}
 \right]_+.
 \label{eq:V39-strip-regression}
\end{equation}
Thus the outputs occupy \(|u-v|\le3k\), and away from the two lower
corners their multiplicity is \(1+k-|u-v|/3\).  A cold disk of radius
\(H\) consequently has joint dimension volume of order
\((k+1)^2H^4\) when \(1\ll k\ll H\).  The product
\(\ell_R\ell_Q\) in
\eqref{eq:V39-joint-strip-volume} is therefore sharp up to constants.
\end{proposition}

\begin{proof}
For these inputs the cancellation variables are
\[
 a=n+k-\frac{2u+v}{3},
 \qquad
 b=n+k-\frac{u+2v}{3}.
\]
Substitute them into
\eqref{eq:V39-BZ-lower}--\eqref{eq:V39-BZ-upper}.  The two large faces
\(a,b\) do not constrain the low-label sector.  The remaining upper
endpoint is \(k-|u-v|/3\), while the lower endpoint is the maximum in
\eqref{eq:V39-strip-regression}.  This proves the exact formula.
Summing over a strip of width \(k\), with multiplicity of order \(k\)
and \(d_{(u,v)}\asymp H^3\) on an annular fraction of a disk containing
order \(H\) radial levels, gives matching upper and lower constant
multiples of \((k+1)^2H^4\).
\end{proof}

\section{Closure of the disjoint tree and global ternary MTA}
\label{sec:V39-ternary-closure}

\begin{theorem}[Uniform closure of the disjoint-overlap formation average]
\label{thm:V39-disjoint-tree-closed}
There are \(C_{\rm dis}<\infty\) and \(\alpha_{\rm dis}\) such that every
disjoint-overlap configuration \eqref{eq:V35-disjoint-overlap-hypothesis}
and every \(\alpha\ge\alpha_{\rm dis}\) satisfy
\begin{equation}
 \boxed{
 \mathfrak A_{L,R}
 \le C_{\rm dis}\frac{H_LH_R}{\Xi(\boldsymbol R_1)}.}
 \label{eq:V39-disjoint-tree-closed}
\end{equation}
The constant is independent of supports, representations, coefficient
matrices, all LR multiplicities, and volume.
\end{theorem}

\begin{proof}
Suppose no uniform constant existed.  The successive localization
theorems V35--V38 would produce the sequence in
\Cref{cor:V38-thin-interior-residual}, with heavy link \(e_\alpha\).
Put \(t=s_\alpha\) and
\[
 \Theta_\alpha=s_\alpha\Xi(\boldsymbol R_1).
\]
By \eqref{eq:V38-heavy-link-mass},
\begin{equation}
 \Theta_{e_\alpha}:=
 s_\alpha\left[
 (1+C_2(R_{1,e_\alpha}))
 +(1+C_2(R_{2,e_\alpha}))
 \right]
 \ge\frac{c\Theta_\alpha}
 {\log(2+\Theta_\alpha)}.
 \label{eq:V39-heavy-local-thermal-lower}
\end{equation}
The selected intermediate satisfies
\[
 1+C_2(S_{e_\alpha})\le L_0/s_\alpha
\]
for a fixed \(L_0\), by
\eqref{eq:V38-heavy-link-cold-output}.

Apply \Cref{cor:V39-universal-one-link-tail} to the marginal
dimension law at \(e_\alpha\).  Its cold mass is at most
\begin{align}
 \frac{C}{(1+\Theta_{e_\alpha})^2}
 &\le
 C\frac{\log^2(2+\Theta_\alpha)}{\Theta_\alpha^2}
 \le\frac C{\Theta_\alpha}
 =\frac C{s_\alpha\Xi(\boldsymbol R_1)}
 \label{eq:V39-heavy-tail-pays}
\end{align}
cofinally.  The exact product dimension law factors over links, and
\(\pi\le w_Lw_R\le\nu_L\nu_R\), so
\eqref{eq:V39-heavy-tail-pays} also bounds the actual coefficient-kernel
mass of the residual event.  Multiplication by
\[
 \mathcal Q_{\boldsymbol S,\boldsymbol U}
 \le Cs_\alpha H_LH_R
\]
from \Cref{lem:V35-one-resolvent-two-defects} proves
\eqref{eq:V39-disjoint-tree-closed}, contradicting the selected failure
sequence.
\end{proof}

\begin{theorem}[Global ternary marked-tree atlas]
\label{thm:V39-global-ternary-MTA}
Fix \(\rho\ge0\).  There are \(K_3<\infty\) and \(\alpha_3\) such that
for every triple of nonempty exact supports and every
\(\alpha\ge\alpha_3\),
\begin{align}
 &\sum_{\varnothing\ne Z\subseteq X_1\cup X_2\cup X_3}
 \norm{\bigl[\mathcal R_{\alpha,\Lambda}Q_0
       \kappa_3^\alpha(g_{1,X_1},g_{2,X_2},g_{3,X_3})\bigr]_Z}
       _{\rho,2}^{\rm BFA}
 \notag\\
 &\quad\le K_3^2\sum_{i=1}^3
 \sum_{\substack{e\in X_i\cap X_j\\
                   f\in X_i\cap X_k}}
 \norm{g_{i,X_i}}_{\rho,2;(e,f)}^{\rm BFA(2)}
 \norm{g_{j,X_j}}_{\rho,2;e}^{\rm BFA(1)}
 \norm{g_{k,X_k}}_{\rho,2;f}^{\rm BFA(1)}.
 \label{eq:V39-global-ternary-MTA}
\end{align}
For each center \(i\), \(\{j,k\}\) are the other two labels.  This is the
\(m=3\) instance of \eqref{eq:V31-MTA}, with no thermal, support,
representation, multiplicity, or volume cutoff.
\end{theorem}

\begin{proof}
The replicated-coordinate expansion in
\Cref{lem:V32-replicated-tree} exhausts the connected ternary incidence
patterns: either one coordinate carries the three-input junction, or two
pair bonds form one of the three labeled trees.  Matrix trace-norm lifting
is the last assertion of \Cref{prop:V32-Wilson-numerator}.

The common-coordinate junction is uniformly bounded by
\Cref{thm:V37-one-link-MTA,cor:V37-same-link-closed}.  For a pair-bond
tree, assign the two bond coordinates to its two edges.  Coincident
three-input coordinates have already been assigned to the junction;
the remaining edge banks are disjoint after spectator factors are
removed, and their exact multiplier is
\eqref{eq:V35-two-bank-operator-factorization}.  Its full formation
average is uniformly bounded by
\Cref{thm:V39-disjoint-tree-closed}.  Bounded thermal, private, and
degenerate banks are included respectively by
\Cref{thm:V32-bounded-thermal-MTA,
thm:V35-private-dominant-tree-MTA,
thm:V36-bi-warm-payment,thm:V36-dimension-cold-MTA}; the V38
localizations show that no other limiting branch was omitted.

In every case the exponential height ratio is at most one, the final
resolvent is charged once, and the scalar factor for a tree centered at
\(i\) is
\[
 \frac{H_{ij}H_{ik}}{\Xi(\boldsymbol R_i)}.
\]
The exact factorization
\eqref{eq:V32-tree-mass-factorization} converts this into the marked BFA
seminorms on the right of
\eqref{eq:V39-global-ternary-MTA}.  Sum the three tree centers and all
input Peter--Weyl blocks.  Disconnected support graphs contribute zero
by cumulant independence.  This proves the theorem.
\end{proof}

\begin{firewall}[Meaning of ternary closure]
\label{fire:V39-ternary-scope}
Theorem \ref{thm:V39-global-ternary-MTA} proves only \(m=3\), together
with the already proved \(m=2\) row.  The definition
\eqref{eq:V31-MTA} requires one constant of the form \(K_0^{m-1}\) for
every \(m\).  Ternary closure therefore removes the first unresolved
atlas order but does not by itself prove MTA, WUFC, CPM, or the mass gap.
\end{firewall}

\section{Exact all-order connected replica partitions}
\label{sec:V39-replica-partitions}

The next construction is purely combinatorial and is stated first for
scalar coordinate functions.  It identifies exactly what must be
estimated at higher order.

Let \(E\) be a finite coordinate set, let
\(X_i\subseteq E\) be nonempty, and put
\[
 F_i:=\prod_{e\in X_i}f_{i,e},
 \qquad 1\le i\le m,
\]
under a product probability law
\(\bigotimes_{e\in E}\nu_e\).  Define the incidence array
\begin{equation}
 \mathcal I:=\{(i,e):e\in X_i\}.
 \label{eq:V39-incidence-array}
\end{equation}
Let \(\rho\) be its row partition, whose blocks have fixed \(i\), and
\(\chi\) its coordinate partition, whose blocks have fixed \(e\).
For a partition \(\sigma\) of \(\mathcal I\), write
\(\sigma\le\chi\) when each block lies in one coordinate.  If
\(B\in\sigma\) lies over coordinate \(e(B)\), put
\[
 \kappa_B^{e(B)}
 :=
 \kappa_{|B|}^{\nu_{e(B)}}
 (f_{i,e(B)}:(i,e(B))\in B).
\]
Singleton cumulants are means.

\begin{theorem}[Leonov--Shiryaev replica formula in incidence form]
\label{thm:V39-connected-replica-formula}
For every \(m\ge2\),
\begin{equation}
 \boxed{
 \kappa_m(F_1,\ldots,F_m)
 =
 \sum_{\substack{\sigma\le\chi\\
                  \sigma\vee\rho=\boldsymbol1_{\mathcal I}}}
 \prod_{B\in\sigma}\kappa_B^{e(B)}.}
 \label{eq:V39-connected-replica-formula}
\end{equation}
Thus only coordinate-cumulant partitions which become connected after
the entries in each input row are identified survive.
\end{theorem}

\begin{proof}
The moment--cumulant formula on the row labels is
\begin{equation}
 \kappa_m(F_1,\ldots,F_m)
 =
 \sum_{\pi\in\Pi_m}\mu(\pi,\boldsymbol1_m)
 \prod_{C\in\pi}
 \mathbb E\prod_{i\in C}F_i.
 \label{eq:V39-row-Mobius}
\end{equation}
For fixed \(\pi\), coordinate independence factors each moment over
\(e\).  Expand every coordinate moment into its local cumulants.  A
choice of all local cumulant blocks is precisely a partition
\(\sigma\le\chi\) whose induced connectivity partition on the row labels
is refined by \(\pi\).

Fix \(\sigma\), and let \(\bar\sigma\in\Pi_m\) be the partition of row
labels generated by its nonsingleton coordinate blocks.  Its coefficient
after rearranging \eqref{eq:V39-row-Mobius} is
\[
 \sum_{\pi\ge\bar\sigma}\mu(\pi,\boldsymbol1_m)
 =
 \begin{cases}
 1,&\bar\sigma=\boldsymbol1_m,\\
 0,&\bar\sigma\ne\boldsymbol1_m.
 \end{cases}
\]
The condition
\(\bar\sigma=\boldsymbol1_m\) is equivalent to
\(\sigma\vee\rho=\boldsymbol1_{\mathcal I}\).  Retaining exactly those
terms proves \eqref{eq:V39-connected-replica-formula}.
\end{proof}

\begin{corollary}[Canonical marked-tree extraction]
\label{cor:V39-canonical-tree-extraction}
Every partition \(\sigma\) occurring in
\eqref{eq:V39-connected-replica-formula} canonically determines a labeled
spanning tree \(\tau(\sigma)\in\mathfrak T_m\) and a valid link marking
\[
 \ell_\sigma:E(\tau(\sigma))\longrightarrow E,
 \qquad
 \ell_\sigma(ij)\in X_i\cap X_j.
\]
Consequently the exact replica sum decomposes into disjoint fibres over
the marked trees appearing on the right of \eqref{eq:V31-MTA}.
\end{corollary}

\begin{proof}
Each nonsingleton block \(B\in\sigma\), lying over coordinate \(e\),
defines a hyperedge on the distinct row labels occurring in \(B\).
The join condition in
\eqref{eq:V39-connected-replica-formula} says exactly that this
hypergraph is connected.

Order coordinates, blocks, and row labels once and for all.  Replace
each hyperedge by the star joining its least row label to every other
row label in that hyperedge, marking every star edge by its coordinate
\(e\).  The resulting labeled multigraph is connected.  Select its
lexicographically first spanning tree, retaining the first occurrence
of an edge if parallel marked edges occur.  Every selected edge \(ij\)
came from a coordinate block containing both \((i,e)\) and \((j,e)\);
hence \(e\in X_i\cap X_j\).  The construction is deterministic, so its
preimages partition the replica sum into the asserted fibres.
\end{proof}

\begin{definition}[Replica-tree fibre estimate]
\label{def:V39-replica-tree-fibre}
The remaining all-order analytic gate is to prove that, after the Wilson
matrix cumulants, all CG/LR output blocks, and the \emph{single} final
resolvent have been assembled, each fibre in
\Cref{cor:V39-canonical-tree-extraction} is bounded by
\begin{equation}
 C_0^{m-1}
 \prod_{i=1}^m
 \norm{g_{i,X_i}}_{
 \rho,2;\boldsymbol e_i^\ell}^{
 \rm BFA(\deg_\tau(i))}
 \label{eq:V39-replica-tree-fibre}
\end{equation}
after summing its local partitions and output multiplicities.  The
constant must be independent of \(m\), supports, representation heights,
and volume.
\end{definition}

\begin{assessment}[Why the exact formula is not yet the estimate]
\label{ass:V39-fibre-obstruction}
Tree extraction supplies the correct support marks but does not bound
the number or analytic size of partitions in one tree fibre.  Taking
absolute values term by term in
\eqref{eq:V39-connected-replica-formula} would also discard the
higher-order Gaussian cancellation identified in V29.  The next proof
must group each coordinate cumulant before absolute summation, retain the
final resolvent once, and show factorial-exponential rather than Bell
growth.  This is the first post-ternary bottleneck.
\end{assessment}

\section{Regression cards and V39 ledger}
\label{sec:V39-ledger}

\begin{assessment}[V34 chain, dual singlet, and strip regressions]
\label{ass:V39-regressions}
The local probability estimate
\eqref{eq:V39-universal-one-link-tail} is inserted only into the actual
product-dimension kernel on the extracted link.  It is then multiplied
by the complete assembled quotient \(\mathcal Q\), with both binary
defects and its single final denominator already present.  Thus the V34
fundamental-chain divergence \eqref{eq:V34-CFMI-divergence} is not
reintroduced.

For \(Q=R^*\) and \(S=\boldsymbol1\),
\Cref{thm:V39-exact-BZ-interval} gives multiplicity one and actual
kernel mass \(d_R^{-2}\).  This agrees with the Schur singlet tax in
\Cref{thm:V38-boundary-tail}; on a repeated physical link the diagonal
cumulant cancellation of \Cref{prop:V37-dual-pair-regression} remains
unchanged.  Finally, \Cref{prop:V39-strip-regression} saturates the
thickness-product order in \eqref{eq:V39-joint-strip-volume}, so the
thin proof has not hidden a false thickness-free lattice count.
\end{assessment}

\begin{theorem}[Integrated V39 statement]
\label{thm:V39-integrated}
Preserving V1--V38, V39 proves:
\begin{enumerate}[label=\textup{(V39.\arabic*)},leftmargin=5.5em]
\item the exact one-dimensional \(\mathrm{SU}(3)\) BZ interval in the
      V36 cancellation variables,
      \eqref{eq:V39-BZ-lower}--\eqref{eq:V39-exact-BZ-multiplicity};
\item the exact diagonal strip and oriented multiplicity controls
      \eqref{eq:V39-output-diagonal-strip}--%
      \eqref{eq:V39-oriented-multiplicity};
\item the joint thin-interior strip-volume estimate
      \eqref{eq:V39-joint-strip-volume} and the squared thin tail
      \eqref{eq:V39-thin-tail};
\item the universal one-link squared cold-output estimate
      \eqref{eq:V39-universal-one-link-tail}, covering every pair of
      irreducible \(\mathrm{SU}(3)\) inputs;
\item closure of the residual disjoint-overlap tree by
      \eqref{eq:V39-disjoint-tree-closed}, and hence the global ternary
      marked-tree atlas estimate \eqref{eq:V39-global-ternary-MTA}; and
\item the exact all-order connected replica-partition formula
      \eqref{eq:V39-connected-replica-formula} together with canonical
      marked-tree extraction, but not the analytic all-order fibre
      estimate \eqref{eq:V39-replica-tree-fibre}.
\end{enumerate}
\end{theorem}

\begin{proof}
Items \textup{(V39.1)--(V39.2)} are
\Cref{thm:V39-exact-BZ-interval,cor:V39-BZ-strip}.  Item
\textup{(V39.3)} is
\Cref{lem:V39-opposite-orientation,lem:V39-joint-strip-volume,
thm:V39-thin-tail}.  Item \textup{(V39.4)} is
\Cref{cor:V39-universal-one-link-tail}, which joins the V39 thin row to
the V38 thick and boundary rows.  Item \textup{(V39.5)} is
\Cref{thm:V39-disjoint-tree-closed,thm:V39-global-ternary-MTA}.
Item \textup{(V39.6)} is
\Cref{thm:V39-connected-replica-formula,
cor:V39-canonical-tree-extraction,ass:V39-fibre-obstruction}; the last
citation records the precise unproved analytic remainder.
\end{proof}

\begingroup
\small
\begin{longtable}{>{\raggedright\arraybackslash}p{0.29\textwidth}
                  >{\raggedright\arraybackslash}p{0.15\textwidth}
                  >{\raggedright\arraybackslash}p{0.48\textwidth}}
\toprule
\textbf{V39 row}&\textbf{Status}&\textbf{Advance or obstruction}\\
\midrule
\endfirsthead
\toprule
\textbf{V39 row}&\textbf{Status}&\textbf{Advance or obstruction}\\
\midrule
\endhead
Exact \(\mathrm{SU}(3)\) BZ interval
& \MGclosed
& All lower and upper Horn endpoints are retained in the V36 variables;
  admissible multiplicity is the positive integral interval length.\\[0.35em]
Thin-interior one-link tail
& \MGclosed
& Opposite orientation, the diagonal strip, and joint summation give two
  thermal powers without separating worst multiplicity from output
  volume.\\[0.35em]
Universal one-link squared tail
& \MGclosed
& The V38 thick/boundary rows and the V39 thin row cover all irreducible
  input pairs.\\[0.35em]
Disjoint-overlap ternary tree
& \MGclosed
& The logarithmic heavy-link extraction loss is absorbed by the squared
  local tail in the actual product kernel.\\[0.35em]
Global ternary MTA
& \MGclosed
& Same-coordinate junctions and disjoint pair-bond trees exhaust the
  connected three-row incidence patterns.\\[0.35em]
Connected replica combinatorics
& \MGclosed
& The Leonov--Shiryaev incidence formula and a deterministic marked-tree
  projection hold at every order.\\[0.35em]
All-order replica-tree fibre estimate
& \MGopen
& Matrix/Peter--Weyl coordinate cumulants must be grouped before absolute
  values, with one final resolvent and only exponential growth in
  \(m\).\\[0.35em]
All-order MTA, WUFC, and CPM
& \MGopen
& These remain conditional on the uniform replica-tree fibre estimate;
  the ternary theorem alone does not imply them.\\[0.35em]
Full Yang--Mills mass gap
& \MGopen
& The nonlinear source packet, finite-edge margins, capacity, protected
  sectors, a nontrivial OS continuum state, and the spectral conclusion
  all remain downstream.\\
\bottomrule
\end{longtable}
\endgroup

\begin{openproblem}[V40 mandate: companion audit and the first all-order fibre]
\label{op:V40}
\begin{enumerate}[label=\textup{(V40.\arabic*)},leftmargin=5.5em]
\item Before choosing the V40 route, perform the scheduled read-only
      audit of the then-current Standard-Model and Navier--Stokes notes
      in \texttt{latex\_notes}.  Record their filenames and SHA--256
      digests, and import only proved, type-compatible estimates.  Do not
      infer a Yang--Mills operator bound from a scalar or differently
      normalized statement.
\item Lift \eqref{eq:V39-connected-replica-formula} to the Wilson
      matrix/Peter--Weyl level.  Assemble local coordinate cumulants and
      CG/LR projections before taking absolute values, and charge the
      final resolvent exactly once.
\item For a fixed marked tree, organize every coordinate block by its
      excess \(|B|-1\).  Seek a uniform Gaussian-excess estimate whose
      product absorbs the number of local partitions without Bell-number
      growth.
\item Test the first new order \(m=4\) completely: include the
      one-coordinate four-cumulant pattern, a three-row hyperedge joined
      to a pair edge, and two-coordinate connected pair-block patterns.
      Prove a \(C_0^3\) fibre bound with the prescribed BFA degrees, or
      give an explicit resolved representation/support witness showing
      where it fails.
\item If the quartic card succeeds, formulate the induction step that
      deletes one leaf or one coordinate hyperedge while preserving the
      matrix cancellation and the single denominator.  Track constants
      so the bound is \(C_0^{m-1}\), not \(m!C_0^m\).
\item If it fails, go upstream to the local Wilson cumulant grouping or
      the canonical tree projection; do not add a raw absolute
      partition count and do not assume the all-order MTA.
\item Preserve only the newest YM file.  The audit required in item
      \textup{(V40.1)} must precede the new proof direction; the next
      periodic companion audit after it is V45.
\end{enumerate}
\end{openproblem}

\begin{firewall}[End-of-V39 claim boundary]
V39 proves the exact BZ interval, the universal squared one-link tail,
the residual disjoint-tree estimate, and global ternary MTA.  It also
proves the exact scalar connected replica formula and a canonical
marked-tree extraction at all orders.  It does not prove the uniform
matrix/Peter--Weyl replica-tree fibre estimate, all-order MTA, WUFC,
CPM, the nonlinear Perron packet, unreduced finite-edge margins,
capacity, protected-sector floors, local-algebra noncollapse,
reflection positivity, a nontrivial continuum OS limit, or a positive
Yang--Mills spectral gap.  The full Yang--Mills mass gap is not proved.
\end{firewall}

\section*{Version V39 append-only record}
\addcontentsline{toc}{section}{Version V39 append-only record}
The complete V38 source was retained before this continuation.  Its
SHA--256 digest was
\begin{center}\scriptsize\ttfamily
43b94534625cc3bfe6195ce62664799a53ce583db9ecd8ee993655a3a6888111.
\end{center}
No mathematical content of V38 was altered.  On the next iteration,
retain this full file, append before its unique active document-closing
command, increment the filename to \texttt{\_V40.tex}, and remove only
the superseded V39 file from this Yang--Mills note series.  V40 begins
with the mandatory SM/NS companion audit.

% =============================================================================
% END APPENDED VERSION V39 MATERIAL
% =============================================================================

% =============================================================================
% BEGIN APPENDED VERSION V40 MATERIAL
% =============================================================================

\clearpage
\part{V40 --- Companion audit, quartic replica skeletons, and the
Wick-tree fibre correction}
\label{part:V40}

\section{Scheduled companion audit and direction}
\label{sec:V40-audit}

V39 closed the thin-interior BZ tail, the residual disjoint ternary
tree, and the global \(m=3\) marked-tree atlas.  It also proved the exact
all-order connected replica formula and projected every connected
coordinate partition to a canonical marked tree.  The missing statement
is analytic: the complete matrix/Peter--Weyl weight of every tree fibre
must be bounded before absolute summation, uniformly in order, support,
representation height, and volume.

As required on every fifth Yang--Mills version, the current companion
files in \texttt{latex\_notes} were inspected read-only before fixing the
V40 direction.  The frozen snapshots are
\begin{center}
\begin{tabular}{p{0.53\textwidth}p{0.39\textwidth}}
\texttt{Renewal\_NS\_Stability\_Research\_Notes\_V31.tex}
&\scriptsize\ttfamily
f1121b62238ad5491bb7cf03635ea9934\newline
942edf573ed8faaa9ee79e1d38a6696\\[0.5em]
\texttt{Renewal\_SM\_Einstein\_Continuum\_Research\_Notes\_V60.tex}
&\scriptsize\ttfamily
3722934dd4516002b40b4c265e7a619f\newline
b234b0d444eb0a015874276205094a90
\end{tabular}
\end{center}

The NS V31 appendix proves that a flat dyadic mark recovers the critical
coarea only after its exact first-moment weight is restored.  It also
shows that separately norming the nonlinear formation terms returns to
continuation-class stock.  The type-correct lesson here is to retain the
exact coordinate excess and assemble connected Wilson terms before
absolute values.  No fluid flux or Sobolev estimate is imported.

The SM V60 appendix computes the exact finite-dimensional residual of a
metric trace/conormal lift before attempting a global domain compiler.
Its type-correct lesson is to test the first unresolved finite order
exactly.  It supplies no Wilson cumulant, Peter--Weyl multiplier, or
marked Fourier estimate.  Thus neither companion file populates the V39
replica-tree fibre row.

\begin{firewall}[Companion estimates remain in their own types]
\label{fire:V40-audit-types}
The audit changes the acquisition order, not the Yang--Mills hypotheses.
Navier--Stokes first-moment formation is not a gauge cumulant, and the
Einstein trace/conormal residual is not a Wilson multiplier.  The only
transfers are proof disciplines: retain the exact weight, postpone
absolute values, and compute the first finite residual before a global
compiler.
\end{firewall}

\section{Exact connected skeletons at four rows}
\label{sec:V40-quartic-skeletons}

Retain the incidence array \(\mathcal I\), row partition \(\rho\), and
coordinate partition \(\chi\) from
\eqref{eq:V39-incidence-array}.  For
\(\sigma\le\chi\), discard its singleton blocks and regard each remaining
block \(B\) as a hyperedge on its row labels.  There is at most one
incidence \((i,e)\) in row \(i\) at a fixed coordinate, so
\(|B|\) is also the number of distinct row labels in this hyperedge.
Define its excess by
\begin{equation}
 \operatorname{exc}(\sigma)
 :=
 \sum_{\substack{B\in\sigma\\|B|\ge2}}(|B|-1).
 \label{eq:V40-partition-excess}
\end{equation}

\begin{lemma}[Connectivity costs \(m-1\) units of coordinate excess]
\label{lem:V40-excess-rank}
If \(\sigma\vee\rho=\boldsymbol1_{\mathcal I}\) on \(m\) nonempty rows,
then
\begin{equation}
 \operatorname{exc}(\sigma)\ge m-1.
 \label{eq:V40-excess-lower}
\end{equation}
The estimate is exact: equality holds for every ordinary tree of pair
blocks and for a single \(m\)-row coordinate block.
\end{lemma}

\begin{proof}
Begin with the discrete partition of the \(m\) row labels.  Adding a
hyperedge on \(k\) labels can merge at most \(k\) existing components,
and therefore can reduce the component count by at most \(k-1\).
Connectivity requires a total reduction of \(m-1\).  Summing this upper
bound over the nonsingleton blocks gives
\eqref{eq:V40-excess-lower}.  A pair-edge tree makes exactly one merge
at each of its \(m-1\) edges, while one \(m\)-edge makes all \(m-1\)
merges at once.
\end{proof}

\begin{theorem}[Complete minimal quartic skeleton list]
\label{thm:V40-quartic-skeleton-list}
Let \(m=4\), and let \(\sigma\) occur in the connected replica formula
\eqref{eq:V39-connected-replica-formula}.  Its nonsingleton coordinate
blocks contain an inclusion-minimal connected subfamily of exactly one
of the following four types:
\begin{enumerate}[label=\textup{(\roman*)},leftmargin=3.7em]
\item one four-row block \([4]\);
\item one three-row block and one pair block joining the omitted row to
      a row of the triple, denoted \([3,2]\);
\item two distinct three-row blocks whose union is all four rows,
      denoted \([3,3]\); or
\item three pair blocks forming a labeled tree, denoted \([2,2,2]\).
\end{enumerate}
In type \([3,2]\) the two blocks lie at distinct coordinates.  The same
is true of the two blocks in type \([3,3]\).  In type \([2,2,2]\), two
disjoint pair blocks may lie at the same coordinate, but two pair blocks
sharing a row may not.
\end{theorem}

\begin{proof}
Choose an inclusion-minimal connected subfamily and let \(k\) be its
largest hyperedge size.  If \(k=4\), that edge alone is connected, so
minimality gives type \([4]\).

Suppose \(k=3\), choose a triple \(B\), and let \(j\) be its omitted
row.  Connectivity supplies another edge \(C\) containing \(j\) and at
least one row of \(B\).  If \(|C|=2\), then \(B,C\) are already
connected on all four rows and minimality excludes every further edge;
this is \([3,2]\).  If \(|C|=3\), the distinct triples \(B,C\) have
union all four rows and again already form a connected family, giving
\([3,3]\).  Blocks which overlap in a row cannot be different blocks of
one coordinate partition, proving the distinct-coordinate assertions.

Finally, if \(k=2\), the minimal family is an inclusion-minimal connected
ordinary graph.  Such a graph is a tree and hence has three edges.
At one coordinate, partition blocks are disjoint, which gives the last
coordinate constraint.
\end{proof}

\begin{corollary}[Exact quartic fibre exhaustion]
\label{cor:V40-quartic-fibre-exhaustion}
Order all eligible minimal subfamilies.  Assign a connected
\(\sigma\le\chi\) to its first minimal subfamily.  This partitions the
complete \(m=4\) replica sum into four disjoint topology classes
\([4]\), \([3,2]\), \([3,3]\), and \([2,2,2]\); no quartic coordinate
incidence lies outside these classes.
\end{corollary}

\begin{proof}
Existence and the four possible types are
\Cref{thm:V40-quartic-skeleton-list}.  A fixed ordering makes the first
choice unique, so the resulting fibres are disjoint and exhaustive.
\end{proof}

\begin{remark}[The two-triple topology is indispensable]
\label{rem:V40-two-triple}
The excess lower bound alone does not reduce every connected quartic
term to total excess three.  Two triples such as
\(\{1,2,3\}\) and \(\{2,3,4\}\) form an inclusion-minimal connected
family of excess four.  Omitting type \([3,3]\) would therefore make the
quartic atlas incomplete.
\end{remark}

\section{The leading quartic Wilson coefficient is tree-valence
sensitive}
\label{sec:V40-Wick-tree}

Put \(s=s_\alpha\).  In the V29 normal coordinates
\(U=\exp(\sqrt{s}\,Z)\), write the first rescaled density correction as
\begin{equation}
 \dd\nu_\alpha
 =
 \phi_8(Z)\bigl[1+s\,a_1(Z)+O(s^2)\bigr]\,\dd Z,
 \qquad
 \mathbb E_{\phi_8}a_1=0.
 \label{eq:V40-first-density-correction}
\end{equation}
This is \eqref{eq:V29-Laplace-normal-form} with \(M=1\), on any fixed
Peter--Weyl bank.  The density and \(a_1\) are even.

Let \(\Gamma\) be the complex bilinear Gaussian covariance on linear
forms,
\[
 \Gamma(\ell,\ell')
 :=\mathbb E_{\phi_8}\bigl[\ell(Z)\ell'(Z)\bigr].
\]
For a smooth scalar matrix coefficient \(g_i\), let
\(D_i^{(r)}=\mathrm d^rg_i(I)\), a symmetric \(r\)-linear form.  If
\(\tau\in\mathfrak T_4\), contract one tensor slot at each endpoint of
every edge using \(\Gamma\), and denote the resulting scalar by
\begin{equation}
 \mathcal W_\tau(g_1,g_2,g_3,g_4)
 :=
 \operatorname{Contr}_{\Gamma,\tau}
 \bigotimes_{i=1}^4D_i^{(\deg_\tau(i))}.
 \label{eq:V40-Wick-tree-contraction}
\end{equation}
Finally define the first-density connected correction
\begin{equation}
 \mathcal U_4(\ell_1,\ell_2,\ell_3,\ell_4)
 :=
 \left.\frac{\dd}{\dd\varepsilon}\right|_{\varepsilon=0}
 \kappa_4^{(1+\varepsilon a_1)\phi_8}
 \bigl(\ell_1(Z),\ell_2(Z),\ell_3(Z),\ell_4(Z)\bigr).
 \label{eq:V40-density-connected-correction}
\end{equation}

\begin{theorem}[Exact leading quartic connected coefficient]
\label{thm:V40-leading-quartic}
On every fixed finite-dimensional Peter--Weyl bank,
\begin{align}
 &\kappa_{4,\mathrm{one}}^\alpha(g_1,g_2,g_3,g_4)
 \notag\\
 &\quad=
 s^3\left[
 \sum_{\tau\in\mathfrak T_4}
 \mathcal W_\tau(g_1,g_2,g_3,g_4)
 \mathcal U_4(D_1^{(1)},D_2^{(1)},D_3^{(1)},D_4^{(1)})
 \right]
 +O_{\mathscr W}(s^4)
 \label{eq:V40-leading-quartic}
\end{align}
uniformly for inputs in a bounded subset of that bank.  There are
\(4^{4-2}=16\) source Wick trees.  Moreover, for normalized matrix
coefficients in irreducibles \(R_i\), with
\(A_i=1+C_2(R_i)\),
\begin{align}
 |\mathcal W_\tau(g_1,g_2,g_3,g_4)|
 &\le C
 \prod_{i=1}^4A_i^{\deg_\tau(i)/2},
 \label{eq:V40-tree-valence-bound}\\
 |\mathcal U_4(D_1^{(1)},D_2^{(1)},D_3^{(1)},D_4^{(1)})|
 &\le C\prod_{i=1}^4A_i^{1/2}.
 \label{eq:V40-density-correction-bound}
\end{align}
Thus the exact leading coefficient is tree-valence sensitive; it is not
bounded uniformly by one Casimir factor per input.
\end{theorem}

\begin{proof}
Cumulants of order at least two are invariant under adding constants, so
replace \(g_i(U)\) by \(g_i(U)-g_i(I)\).  Taylor-expand
\[
 g_i(\exp(\sqrt{s}Z))-g_i(I)
 =
 \sum_{r\ge1}\frac{s^{r/2}}{r!}D_i^{(r)}[Z^{\otimes r}]
\]
through total degree eight and use the dominated remainder supplied by
\eqref{eq:V29-Laplace-remainder}.  In the Gaussian part of the
moment--cumulant formula, Wick pairings which do not connect all four
marked source vertices cancel.  A connected pairing uses at least three
contraction edges.  At total degree six it has exactly three edges and
four vertices, and hence is a labeled tree.

Fix such a tree.  At vertex \(i\), its
\(\deg_\tau(i)\) incident edges can be assigned to the symmetric
derivative slots in \(\deg_\tau(i)!\) ways.  This cancels the Taylor
factor \(1/\deg_\tau(i)!\).  Therefore that tree contributes exactly
\eqref{eq:V40-Wick-tree-contraction}.  Summing the sixteen labeled
trees gives the first term in brackets in
\eqref{eq:V40-leading-quartic}.

The order-\(s\) density correction can contribute at total order
\(s^3\) only with four linear source terms.  Its connected part is
precisely the derivative in
\eqref{eq:V40-density-connected-correction}.  Terms of half-integral
order vanish because the density is even and their total polynomial
degree is odd.  The next possible terms have order \(s^4\).
The fixed-bank Gaussian domination in V29 justifies termwise integration
and bounds their sum by \(O_{\mathscr W}(s^4)\).

For \(r\le3\), the derived-representation Casimir identity and repeated
operator multiplication give
\[
 \|D_i^{(r)}\|\le C_rA_i^{r/2}.
\]
Gaussian contraction of the finitely many slots proves
\eqref{eq:V40-tree-valence-bound}.  In
\eqref{eq:V40-density-connected-correction}, \(a_1\) is a fixed
Gaussian-integrable polynomial and each argument is linear.
Hölder's inequality and the \(r=1\) derivative bound prove
\eqref{eq:V40-density-correction-bound}.
\end{proof}

\section{A resolved SU(3) high-weight regression}
\label{sec:V40-high-weight-regression}

Let \(R_N=(N,0)=\operatorname{Sym}^N(\mathbf3)\), let \(v_N\) be its
highest-weight unit vector, and set
\begin{equation}
 g_N(U):=\langle v_N,R_N(U)v_N\rangle=(U_{11})^N,
 \qquad h(U):=U_{11}.
 \label{eq:V40-coherent-coefficients}
\end{equation}
These are normalized matrix coefficients.  If
\(\ell(X)=X_{11}\), then
\begin{equation}
 \mathrm d^rg_N(I)
 =
 N^r\ell^{\otimes r}+O_r(N^{r-1}),
 \qquad r=1,2,3.
 \label{eq:V40-coherent-jets}
\end{equation}

\begin{proposition}[The bare product-Casimir quartic excess is false]
\label{prop:V40-bare-quartic-false}
There is no constant \(C\), independent of \(\alpha\) and the four
irreducible inputs, for which
\begin{equation}
 \|\mathbf K_{1234}\|_{\rm op}
 \le
 C\min\left\{
 s_\alpha^3\prod_{i=1}^4(1+C_2(R_i)),\,1
 \right\}.
 \label{eq:V40-false-bare-quartic}
\end{equation}
Here \(\mathbf K_{1234}\) is the complete one-link fourth cumulant
operator.  The failure already occurs for
\((R_1,R_2,R_3,R_4)=(R_N,\mathbf3,\mathbf3,\mathbf3)\).
\end{proposition}

\begin{proof}
Apply \Cref{thm:V40-leading-quartic} to
\((g_N,h,h,h)\), first with \(N\) fixed and then let \(s\downarrow0\).
Put \(c=\Gamma(\ell,\ell)\ne0\).  Among the sixteen labeled trees, the
unique star centered at vertex \(1\) has degree three there.  Equations
\eqref{eq:V40-Wick-tree-contraction} and
\eqref{eq:V40-coherent-jets} give its contribution
\[
 c^3N^3+O(N^2).
\]
Every other tree has degree at most two at vertex \(1\), so its
contribution is \(O(N^2)\).  The density correction contains only the
first derivative of \(g_N\) and is \(O(N)\).  Consequently
\begin{equation}
 \lim_{s\downarrow0}s^{-3}
 \kappa_{4,\mathrm{one}}^\alpha(g_N,h,h,h)
 =
 c^3N^3+O(N^2).
 \label{eq:V40-high-weight-leading}
\end{equation}
The scalar on the left is a matrix entry of
\(\mathbf K_{1234}\), so its absolute value is bounded by the operator
norm.

For \(R_N=(N,0)\), the exact SU(3) Casimir formula gives
\(1+C_2(R_N)\asymp N^2\), while the three fundamental Casimirs are
fixed.  If \eqref{eq:V40-false-bare-quartic} held, divide it by \(s^3\)
and let \(s\downarrow0\) for each fixed \(N\).  Equation
\eqref{eq:V40-high-weight-leading} would imply
\(|c|^3N^3+O(N^2)\le C'N^2\) for every \(N\), a contradiction.
\end{proof}

\begin{assessment}[What the regression does and does not refute]
\label{ass:V40-regression-scope}
The proposition confirms the V29 fixed-order firewall in the first new
order.  It refutes a bare local estimate which ignores the valence of the
high representation.  It does not refute MTA: the MTA expression sums
labeled trees, its marked input degrees equal their tree valences, and
the complete output resolvent is applied only after the cumulant and all
fusion channels have been assembled.  Those three structures are absent
from \eqref{eq:V40-false-bare-quartic}.
\end{assessment}

\section{The exact one-link-support quartic row}
\label{sec:V40-one-link-row}

\begin{theorem}[Quartic MTA for four exact one-link supports]
\label{thm:V40-one-link-quartic-MTA}
Fix a link \(e\) and suppose
\[
 X_1=X_2=X_3=X_4=\{e\}.
\]
Then the \(m=4\) instance of \eqref{eq:V31-MTA} holds uniformly in
\(\alpha\), representation height, and multiplicity.
\end{theorem}

\begin{proof}
On a one-link exact support,
\(\vartheta_e(\boldsymbol R)=1\).  Hence every repeated marked seminorm
at \(e\) equals its unmarked BFA norm.  There are sixteen labeled trees,
each with the unique marking \(e\), so the right side of
\eqref{eq:V31-MTA} at \(m=4\) is
\[
 16K_0^3\prod_{i=1}^4
 \|g_{i,\{e\}}\|_{\rho,2}^{\rm BFA}.
\]
The all-order one-bank theorem
\eqref{eq:V30-one-bank-cumulants}, followed by the BFA/FA equivalence
\eqref{eq:V31-FA-BFA-equivalence}, bounds the left side by
\[
 4!\,c_\Xi K_{\rm cum}^3
 \prod_{i=1}^4
 \|g_{i,\{e\}}\|_{\rho,2}^{\rm BFA}.
\]
Choose \(K_0\) so that
\(16K_0^3\ge4!c_\Xi K_{\rm cum}^3\).
The V30 theorem assembles the complete fourth cumulant, all final
Peter--Weyl channels, and the single resolvent before taking the norm, so
no forbidden termwise denominator or multiplicity sum is introduced.
\end{proof}

\begin{firewall}[One-link support is not an embedded junction]
\label{fire:V40-one-link-scope}
\Cref{thm:V40-one-link-quartic-MTA} closes the case in which all four
complete exact supports equal one link.  If the inputs have additional
spectator links, the repeated mark at the junction can be much smaller
than one.  The unmarked V30 estimate then cannot simply be substituted
for the degree-marked MTA fibre.
\end{firewall}

\section{Upstream correction of the replica-tree fibre}
\label{sec:V40-fibre-correction}

The canonical tree in \Cref{cor:V39-canonical-tree-extraction} remains a
valid combinatorial certificate.  It is not yet an analytically valid
place to take absolute values.  In particular, assigning an entire
four-row coordinate block to one lexicographic star would require
controlling all sixteen leading contractions by that star's valence
profile, contradicting \Cref{prop:V40-bare-quartic-false}.

\begin{definition}[Contraction-refined quartic fibre]
\label{def:V40-contraction-refined-fibre}
At every four-row coordinate block, retain the sixteen summands
\(\mathcal W_\tau\) in \eqref{eq:V40-leading-quartic} separately until
their tree marks and the final output resolvent have been assigned.
Retain the density correction
\(\mathcal U_4\) as one assembled Wilson-law term.  For the
\([3,2]\), \([3,3]\), and \([2,2,2]\) skeletons of
\Cref{thm:V40-quartic-skeleton-list}, retain every local connected block
and its coordinate before selecting a spanning tree of their union.
The resulting object is called a contraction-refined quartic fibre.
\end{definition}

\begin{proposition}[Leading-order routing has the MTA tree index]
\label{prop:V40-leading-routing}
The pure-Gaussian leading part of every \([4]\) block is an exact sum
over \(\mathfrak T_4\), and its summand at \(\tau\) has representation
growth bounded by the valence profile
\eqref{eq:V40-tree-valence-bound}.  The first Wilson-density correction
has only one first-derivative factor per input, as in
\eqref{eq:V40-density-correction-bound}.  Thus no Bell or partition
count occurs in the leading quartic one-coordinate coefficient; the
necessary refinement has exactly sixteen tree fibres plus one assembled
density term.
\end{proposition}

\begin{proof}
This is the term-by-term content of
\Cref{thm:V40-leading-quartic}.  Cayley's formula gives sixteen trees,
and the proof of that theorem shows that disconnected Wick pairings have
already cancelled before the tree sum.  The density correction is the
single derivative of the connected cumulant under the normalized
first-order density perturbation, not a sum over set partitions after
absolute values.
\end{proof}

\begin{assessment}[The first surviving quartic analytic gate]
\label{ass:V40-surviving-gate}
V40 has exhausted the four possible coordinate skeletons and corrected
the fibre index at a four-row junction.  What remains is a finite-\(s\),
operator-valued estimate for the contraction-refined fibres with
additional spectator coordinates.  It must turn the local valence
factors into the repeated probability marks of
\eqref{eq:V31-marked-BFA}, sum CG/LR multiplicities by compression, and
use the final output resolvent once.  Fixed-bank asymptotics do not supply
uniform control of the remainder when representation height grows.
\end{assessment}

\section{V40 ledger and next acquisition}
\label{sec:V40-ledger}

\begin{theorem}[Integrated V40 statement]
\label{thm:V40-integrated}
Preserving V1--V39, V40 proves:
\begin{enumerate}[label=\textup{(V40.\arabic*)},leftmargin=5.5em]
\item the scheduled, hashed SM/NS audit and its no-import decision;
\item the excess lower bound \eqref{eq:V40-excess-lower};
\item the exhaustive quartic skeleton classification
      \([4]\), \([3,2]\), \([3,3]\), and \([2,2,2]\);
\item the exact fixed-bank leading coefficient
      \eqref{eq:V40-leading-quartic}, indexed by the sixteen labeled
      Wick trees and one assembled density correction;
\item the high-weight SU(3) regression
      \eqref{eq:V40-high-weight-leading}, which disproves the bare
      product-Casimir quartic excess;
\item the exact one-link-support quartic MTA theorem; and
\item the contraction-refined fibre index and the precise finite-\(s\)
      mixed-support remainder.
\end{enumerate}
\end{theorem}

\begin{proof}
Item \textup{(V40.1)} is
\Cref{sec:V40-audit,fire:V40-audit-types}.  Items
\textup{(V40.2)--(V40.3)} are
\Cref{lem:V40-excess-rank,thm:V40-quartic-skeleton-list,
cor:V40-quartic-fibre-exhaustion}.  Item \textup{(V40.4)} is
\Cref{thm:V40-leading-quartic}.  Item \textup{(V40.5)} is
\Cref{prop:V40-bare-quartic-false}.  Item \textup{(V40.6)} is
\Cref{thm:V40-one-link-quartic-MTA}.  Item \textup{(V40.7)} is
\Cref{def:V40-contraction-refined-fibre,
prop:V40-leading-routing,ass:V40-surviving-gate}.
\end{proof}

\begingroup
\small
\begin{longtable}{>{\raggedright\arraybackslash}p{0.29\textwidth}
                  >{\raggedright\arraybackslash}p{0.15\textwidth}
                  >{\raggedright\arraybackslash}p{0.48\textwidth}}
\toprule
\textbf{V40 row}&\textbf{Status}&\textbf{Advance or obstruction}\\
\midrule
\endfirsthead
\toprule
\textbf{V40 row}&\textbf{Status}&\textbf{Advance or obstruction}\\
\midrule
\endhead
Scheduled SM/NS audit
& \MGclosed
& Current hashes are recorded; neither companion theorem has the
  Yang--Mills operator type needed by the fibre estimate.\\[0.35em]
Quartic replica topology
& \MGclosed
& Four and only four minimal connected skeleton types occur; the
  two-triple type cannot be dropped.\\[0.35em]
Leading one-coordinate quartic coefficient
& \MGclosed
& Sixteen source Wick trees plus one assembled first-density correction
  give the exact \(s^3\) coefficient on every fixed Fourier bank.\\[0.35em]
Bare product-Casimir quartic bound
& \MGfail
& The coherent family \((N,0),\mathbf3,\mathbf3,\mathbf3\) grows like
  \(s^3N^3\), rather than \(s^3N^2\).\\[0.35em]
Exact one-link-support quartic MTA
& \MGclosed
& V30's assembled all-order one-bank estimate closes the case
  \(X_1=\cdots=X_4=\{e\}\).\\[0.35em]
Mixed-support contraction-refined quartic fibres
& \MGopen
& A finite-\(s\) operator estimate must convert tree valence into repeated
  link marks while retaining spectators, multiplicity compression, and
  one final resolvent.\\[0.35em]
All-order MTA, WUFC, and CPM
& \MGopen
& They remain downstream of the corrected quartic fibre estimate and its
  uniform induction.\\[0.35em]
Full Yang--Mills mass gap
& \MGopen
& The nonlinear source packet, finite-edge margins, capacity, protected
  sectors, nontrivial OS limit, and spectral conclusion remain open.\\
\bottomrule
\end{longtable}
\endgroup

\begin{openproblem}[V41 mandate: finite-\(s\) contraction-tree operator
fibres]
\label{op:V41}
\begin{enumerate}[label=\textup{(V41.\arabic*)},leftmargin=5.5em]
\item Construct an exact finite-\(s\) operator interpolation for the
      complete one-coordinate fourth cumulant.  Group it by contraction
      trees before absolute values; do not use the false estimate
      \eqref{eq:V40-false-bare-quartic}.
\item First treat a \([4]\) junction embedded in four larger exact
      supports.  Retain all spectator multipliers and prove the repeated
      marked-BFA weights with the final resolvent charged once.
\item Compress all final CG/LR multiplicity spaces by the V37
      partial-trace method.  No entrywise multiplicity count may precede
      this compression.
\item Then treat, in order, the \([3,2]\), \([3,3]\), and
      \([2,2,2]\) fibres from
      \Cref{thm:V40-quartic-skeleton-list}.  The \([3,3]\) topology must
      remain explicit even though its excess is four.
\item Test every proposed local bound on the coherent family
      \eqref{eq:V40-coherent-coefficients}, with each possible high
      vertex degree.  A bound of order \(s^3N^2\) at the degree-three
      vertex is already disproved.
\item If the finite-\(s\) quartic operator card closes, state and prove
      the global \(m=4\) MTA before attempting induction in \(m\).
      Otherwise return the first resolved topology and representation
      family which fails, then go upstream at its multiplier identity.
\item Preserve only the newest YM file.  The next scheduled SM/NS
      companion audit is V45.
\end{enumerate}
\end{openproblem}

\begin{firewall}[End-of-V40 claim boundary]
V40 completes the companion audit, the quartic topology list, the exact
leading fixed-bank Wick-tree coefficient, a decisive high-weight
regression, and the exact one-link-support quartic MTA row.  It does not
prove a finite-\(s\) degree-marked operator bound for mixed supports,
global \(m=4\) MTA, all-order MTA, WUFC, CPM, the nonlinear Perron
packet, unreduced finite-edge margins, capacity, protected-sector
floors, local-algebra noncollapse, reflection positivity, a nontrivial
continuum OS state, or a positive Yang--Mills spectral gap.  The full
Yang--Mills mass gap is not proved.
\end{firewall}

\section*{Version V40 append-only record}
\addcontentsline{toc}{section}{Version V40 append-only record}
The complete V39 source was retained before this continuation.  Its
SHA--256 digest was
\begin{center}\scriptsize\ttfamily
b6de820f34ab365944f35707ff75e08bc8efba5bd499efeeec3fd2f381ec9a47.
\end{center}
No mathematical content of V39 was altered.  On the next iteration,
retain this full file, append before its unique active document-closing
command, increment the filename to \texttt{\_V41.tex}, and remove only
the superseded V40 file from this Yang--Mills note series.  The next
scheduled SM/NS audit is V45.

% =============================================================================
% END APPENDED VERSION V40 MATERIAL
% =============================================================================

% =============================================================================
% BEGIN APPENDED VERSION V41 MATERIAL
% =============================================================================

\clearpage
\part{V41 --- Exact finite-\(s\) quartic contraction trees and
multiplicity-uniform compression}
\label{part:V41}

\section{Purpose and corrected acquisition}
\label{sec:V41-purpose}

V40 proved that the bare estimate
\eqref{eq:V40-false-bare-quartic} is false: at a degree-three high
vertex the leading fourth cumulant grows as \(s^3N^3\), not
\(s^3N^2\).  The same calculation identified the missing structure.
The Gaussian leading term is already a sum over the sixteen labeled
trees, and the representation power at each input is half its tree
valence.

This version turns that asymptotic observation into an exact
finite-\(s\) operator statement.  The common Gaussian increment admits a
replica forest interpolation whose connected part is indexed exactly by
labeled trees.  The Wilson increment law differs from this Gaussian
reference by an \(O(s)\) density perturbation.  Because a fourth
cumulant is unchanged by four independent constant shifts, that
perturbation carries one displacement factor from every input.  An
elementary max-star inequality routes the whole perturbation into one of
the sixteen tree fibres with its correct valence weight.

The resulting decomposition is uniform in all four representations and
at finite \(s\); it also survives CG/LR compression without a
multiplicity count.  Converting these local tree weights into repeated
BFA probability marks in the presence of arbitrary spectator links
remains the next gate.

\section{Gaussian replica interpolation in operator form}
\label{sec:V41-Gaussian-BKAR}

Let \(\mathfrak g=\mathfrak{su}(3)\) carry the invariant Euclidean
structure used in V29, and let \(\gamma\) have density \(\phi_8\).
For irreducibles \(R_i\), put
\begin{equation}
 F_i(Z):=D^{R_i}(\exp(\sqrt{s}\,Z)),
 \qquad
 A_i:=1+C_2(R_i),
 \qquad
 x_i:=\sqrt{sA_i}.
 \label{eq:V41-Gaussian-functions}
\end{equation}
All tensor products below are placed in
\(\bigotimes_{i=1}^4H_{R_i}\), so factors attached to distinct replicas
commute.

For a labeled tree \(\tau\in\mathfrak T_4\) and
\(\boldsymbol t=(t_e)_{e\in E(\tau)}\in[0,1]^3\), let
\[
 q_{ii}^{\tau}(\boldsymbol t)=1,
 \qquad
 q_{ij}^{\tau}(\boldsymbol t)
 =\min_{e\in\operatorname{path}_\tau(i,j)}t_e
 \quad(i\ne j).
\]
The matrix \(q^\tau(\boldsymbol t)\) is positive semidefinite.  Let
\(\gamma_{\tau,\boldsymbol t}\) be the centered Gaussian law of four
replicas whose cross covariance is
\(q_{ij}^\tau(\boldsymbol t)\) times the covariance of \(\gamma\).
If \((E_a)\) is an orthonormal basis of \(\mathfrak g\), write
\[
 \mathscr D_{ij}
 :=
 \sum_{a,b}\Gamma^{ab}
 \partial_{Z_i^a}\partial_{Z_j^b},
 \]
where \(\Gamma^{ab}=\mathbb E_\gamma[Z^aZ^b]\).
The finite-dimensional forest identity used below is the
Brydges--Kennedy--Abdesselam--Rivasseau formula; the path-minimum
interpolation and its combinatorial proof are given in
A.~Abdesselam and V.~Rivasseau,
\emph{Trees, forests and jungles},
\url{https://arxiv.org/abs/hep-th/9409094}.

\begin{lemma}[Operator-valued Gaussian cumulant tree formula]
\label{lem:V41-Gaussian-tree-formula}
The complete fourth cumulant under the common Gaussian replica is
\begin{equation}
 \boxed{
 \mathbf K_{1234}^{\gamma}
 =
 \sum_{\tau\in\mathfrak T_4}
 \mathbf K_\tau^\gamma,}
 \qquad
 \mathbf K_\tau^\gamma
 :=
 \int_{[0,1]^3}
 \mathbb E_{\gamma_{\tau,\boldsymbol t}}
 \left[
 \left(\prod_{ij\in E(\tau)}\mathscr D_{ij}\right)
 \bigotimes_{i=1}^4F_i(Z_i)
 \right]\dd\boldsymbol t .
 \label{eq:V41-Gaussian-tree-formula}
\end{equation}
This identity is exact for every \(s>0\), not an expansion at
\(s=0\).
\end{lemma}

\begin{proof}
For a covariance array \(q=(q_{ij})\), let
\[
 \Phi(q)
 :=
 \mathbb E_q\bigotimes_{i=1}^4F_i(Z_i).
\]
At the block-diagonal array associated with a partition
\(\pi\in\Pi_4\), Gaussian independence between its blocks makes
\(\Phi(q^\pi)\) the tensor product of the corresponding block moments.
Consequently the Möbius connected combination
\[
 \sum_{\pi\in\Pi_4}\mu(\pi,\boldsymbol1_4)\Phi(q^\pi)
 \]
is exactly \(\mathbf K_{1234}^{\gamma}\).

Apply the multivariate fundamental theorem of calculus to the six
off-diagonal covariance coordinates, beginning at the block-diagonal
arrays.  Organize a derivative set by the graph it draws on the four
replicas.  The Möbius sum cancels every disconnected graph.  In a
connected derivative set, retain its first spanning tree and integrate
the remaining covariance coordinates along the minimum-on-path
interpolation \(q^\tau(\boldsymbol t)\).  This is the forest
interpolation identity; differentiating the right side with respect to
one edge parameter and then deleting that edge gives the same identity
on the two resulting components, which proves it inductively from the
one-vertex case.

Gaussian differentiation in \(q_{ij}\) is integration by parts and
inserts \(\mathscr D_{ij}\).  Thus the connected forest formula is
precisely \eqref{eq:V41-Gaussian-tree-formula}.  The argument applies
entrywise in the finite-dimensional tensor carrier, and Bochner
integration restores the operator identity.
\end{proof}

\begin{lemma}[Uniform derivatives of a unitary exponential coefficient]
\label{lem:V41-unitary-derivatives}
For \(d=1,2,3\),
\begin{equation}
 \sup_{Z\in\mathfrak g}
 \|\nabla^dF_i(Z)\|_{\rm op}
 \le C_d x_i^d.
 \label{eq:V41-unitary-derivatives}
\end{equation}
Hence, with
\begin{equation}
 w_\tau(\boldsymbol A;s)
 :=
 \prod_{ij\in E(\tau)}s\sqrt{A_iA_j}
 =
 s^3\prod_{i=1}^4A_i^{\deg_\tau(i)/2},
 \label{eq:V41-tree-weight}
\end{equation}
one has
\begin{equation}
 \|\mathbf K_\tau^\gamma\|_{\rm op}
 \le C_{\rm G}\,w_\tau(\boldsymbol A;s).
 \label{eq:V41-Gaussian-tree-bound}
\end{equation}
\end{lemma}

\begin{proof}
The first Fréchet derivative of the matrix exponential is an integral of
products of unitary exponentials with one inserted Lie-algebra
direction.  The \(d\)-th derivative is a sum of at most \(C_d\) simplex
integrals with \(d\) inserted directions.  After applying \(R_i\), all
exponential factors remain unitary.  The Casimir identity gives
\[
 \|\mathrm dR_i(X)\|_{\rm op}
 \le C A_i^{1/2}\|X\|.
\]
The chain rule supplies one factor \(\sqrt{s}\) per direction, proving
\eqref{eq:V41-unitary-derivatives}.

In \eqref{eq:V41-Gaussian-tree-formula}, vertex \(i\) is differentiated
exactly \(\deg_\tau(i)\) times.  Contracting the three covariance tensors
costs only a group constant.  The correlated Gaussian expectation is a
contraction in operator norm.  Multiplying the four derivative bounds
gives \eqref{eq:V41-Gaussian-tree-bound}.
\end{proof}

\section{Uniform comparison of the Wilson and Gaussian cumulants}
\label{sec:V41-Wilson-comparison}

Let \(\mathbf K_{1234}^{\nu}\) denote the complete one-link cumulant
operator under the Wilson increment law \(\nu_\alpha\), and retain
\(\mathbf K_{1234}^{\gamma}\) from
\Cref{lem:V41-Gaussian-tree-formula}.  Both include all fifteen
moment-partition terms before any norm is taken.

\begin{lemma}[Four-displacement Wilson--Gaussian comparison]
\label{lem:V41-Wilson-Gaussian-comparison}
There are group constants \(C,c,s_0>0\) such that, for
\(0<s=s_\alpha\le s_0\) and arbitrary irreducibles,
\begin{equation}
 \boxed{
 \|\mathbf K_{1234}^{\nu}
       -\mathbf K_{1234}^{\gamma}\|_{\rm op}
 \le
 C\left[
 s\prod_{i=1}^4\min\{1,x_i\}
 +e^{-c/s}
 \right].}
 \label{eq:V41-Wilson-Gaussian-comparison}
\end{equation}
The constant is independent of all representation heights.
\end{lemma}

\begin{proof}
Cumulants of order four are unchanged when an argument is shifted by a
constant.  In both cumulants replace \(D^{R_i}(U)\) by
\[
 Z_i(U):=D^{R_i}(U)-I.
\]
Unitarity and the integrated first derivative give, on the normal ball,
\begin{equation}
 \|Z_i(\exp(\sqrt{s}Z))\|_{\rm op}
 \le\min\{2,Cx_i\|Z\|\}.
 \label{eq:V41-four-displacements}
\end{equation}

Use \Cref{lem:V29-Laplace-normal-form} with \(M=0\).  On the rescaled
normal ball the Wilson density relative to \(\phi_8\) is
\[
 1+s\mathcal E_{0,s}(Z),
 \qquad
 |\mathcal E_{0,s}(Z)|
 \le C(1+\|Z\|)^N e^{\|Z\|^2/16}.
\]
The Wilson mass outside the ball and the Gaussian mass beyond its
rescaled radius are \(O(e^{-c/s})\).

Write each fourth cumulant in its moment--partition formula.  In the
difference of the two formulas, telescope each product of block moments
so that the measure difference occurs in one block at a time.  Every
input index still occurs exactly once across that product.  Equation
\eqref{eq:V41-four-displacements}, Gaussian integrability, and Hölder's
inequality therefore bound every density-perturbation term by
\[
 Cs\prod_{i=1}^4\min\{1,x_i\}.
\]
There are only fifteen partitions and at most four block factors, so
their count is absorbed into \(C\).  On either exterior set use
\(\|Z_i\|_{\rm op}\le2\), giving the exponential remainder.  This proves
\eqref{eq:V41-Wilson-Gaussian-comparison}.
\end{proof}

\begin{lemma}[The density perturbation fits one labeled star]
\label{lem:V41-max-star}
Let \(0<s\le1\), \(x_i\ge\sqrt{s}\), and define
\[
 w_\tau:=\prod_{ij\in E(\tau)}x_ix_j.
\]
If \(r\) is an index for which \(x_r=\max_i x_i\), and
\(\tau_r\) is the star centered at \(r\), then
\begin{equation}
 s\prod_{i=1}^4\min\{1,x_i\}
 \le w_{\tau_r}
 \le\max_{\tau\in\mathfrak T_4}w_\tau,
 \qquad
 e^{-c/s}\le C_c\max_\tau w_\tau.
 \label{eq:V41-max-star}
\end{equation}
\end{lemma}

\begin{proof}
The star weight is
\(w_{\tau_r}=x_r^3\prod_{i\ne r}x_i\).  Since
\(\min\{1,x_i\}/x_i\le1\) for \(i\ne r\),
\[
 \frac{s\prod_i\min\{1,x_i\}}{w_{\tau_r}}
 \le
 \frac{s\min\{1,x_r\}}{x_r^3}.
\]
If \(x_r\le1\), the last ratio is \(s/x_r^2\le1\); if
\(x_r>1\), it is \(s/x_r^3\le1\).  This proves the first assertion.
Every edge factor \(x_ix_j\) is at least \(s\), so every tree weight is
at least \(s^3\).  The elementary bound
\(e^{-c/s}\le C_cs^3\) proves the second assertion.
\end{proof}

\section{Exact finite-\(s\) contraction-tree decomposition}
\label{sec:V41-exact-tree-decomposition}

\begin{theorem}[Representation-uniform one-link fourth-cumulant trees]
\label{thm:V41-one-link-fourth-trees}
There are operators
\(\mathbf K_\tau^\nu\) on
\(\bigotimes_{i=1}^4H_{R_i}\), one for every
\(\tau\in\mathfrak T_4\), such that
\begin{align}
 \boxed{\mathbf K_{1234}^{\nu}
 =\sum_{\tau\in\mathfrak T_4}\mathbf K_\tau^\nu,}
 \label{eq:V41-exact-Wilson-tree-decomposition}\\
 \boxed{\|\mathbf K_\tau^\nu\|_{\rm op}
 \le C_4\,s_\alpha^3
 \prod_{i=1}^4A_i^{\deg_\tau(i)/2}.}
 \label{eq:V41-exact-Wilson-tree-bound}
\end{align}
The constants and the threshold in \(\alpha\) are independent of all
four representations.  Each \(\mathbf K_\tau^\nu\) commutes with the
diagonal SU(3) action.  Consequently,
\begin{equation}
 \|\mathbf K_{1234}^{\nu}\|_{\rm op}
 \le C_4
 \min\left\{
 1,\,
 s_\alpha^3
 \sum_{\tau\in\mathfrak T_4}
 \prod_iA_i^{\deg_\tau(i)/2}
 \right\}.
 \label{eq:V41-aggregate-fourth-bound}
\end{equation}
\end{theorem}

\begin{proof}
Let
\(\Delta=\mathbf K_{1234}^{\nu}
-\mathbf K_{1234}^{\gamma}\), and choose
\(\tau_*\) maximizing \(w_\tau(\boldsymbol A;s)\).  Define
\[
 \mathbf K_\tau^\nu=
 \begin{cases}
 \mathbf K_\tau^\gamma,&\tau\ne\tau_*,\\
 \mathbf K_{\tau_*}^\gamma+\Delta,&\tau=\tau_*.
 \end{cases}
\]
The Gaussian identity \eqref{eq:V41-Gaussian-tree-formula} gives
\eqref{eq:V41-exact-Wilson-tree-decomposition}.  By
\Cref{lem:V41-Wilson-Gaussian-comparison,lem:V41-max-star},
\[
 \|\Delta\|_{\rm op}\le Cw_{\tau_*}.
\]
Combine this with \eqref{eq:V41-Gaussian-tree-bound} to obtain
\eqref{eq:V41-exact-Wilson-tree-bound}.

The Gaussian reference is Ad-invariant, the replica covariance is scalar
on \(\mathfrak g\), and every \(\mathscr D_{ij}\) contracts with the
invariant covariance tensor.  Hence each
\(\mathbf K_\tau^\gamma\) intertwines the diagonal group action.
Both complete cumulants do as well, so \(\Delta\), and therefore every
\(\mathbf K_\tau^\nu\), is an intertwiner.

Summing \eqref{eq:V41-exact-Wilson-tree-bound} proves the uncapped
aggregate estimate.  Independently, the absolute Möbius coefficients of
the fourth cumulant sum to a numerical constant and every moment
operator is a contraction, so
\(\|\mathbf K_{1234}^{\nu}\|_{\rm op}\le C\).
Taking the better of the two bounds proves
\eqref{eq:V41-aggregate-fourth-bound}.
\end{proof}

\begin{assessment}[Resolution of the V40 high-weight regression]
\label{ass:V41-high-weight-resolution}
For the coherent family
\eqref{eq:V40-coherent-coefficients}, the star centered at the
\((N,0)\) input has weight \(s^3A_N^{3/2}\asymp s^3N^3\), exactly the
growth found in \eqref{eq:V40-high-weight-leading}.  The other fifteen
trees have degree at most two at that vertex.  Thus
\eqref{eq:V41-exact-Wilson-tree-bound} passes the regression which
invalidated \eqref{eq:V40-false-bare-quartic}; no representation power
has been hidden in the constant.
\end{assessment}

\section{Multiplicity-uniform final-channel compression}
\label{sec:V41-compression}

For trace-class input matrices \(B_i\) in the irreducible carriers
\(H_{R_i}\), put
\[
 \mathcal B:=B_1\otimes B_2\otimes B_3\otimes B_4,
 \qquad
 d_{\rm in}:=\prod_{i=1}^4d_{R_i}.
\]
Resolve the diagonal fourfold carrier into final irreducibles and all
their multiplicity spaces.  Let
\(\widehat\kappa_{\tau,T}^{\,\alpha}\) be the output coefficient obtained
from \(\mathbf K_\tau^\nu\mathcal B\) in the final channel \(T\), with
the same Fourier normalization as in V37.

\begin{theorem}[Fourfold tree compression has no LR path count]
\label{thm:V41-fourfold-compression}
For every set \(\Omega\) of final irreducible channels,
\begin{equation}
 \boxed{
 \sum_{T\in\Omega}d_T
 \|\widehat\kappa_{\tau,T}^{\,\alpha}\|_1
 \le
 C_4\,d_{\rm in}\,
 w_\tau(\boldsymbol A;s_\alpha)
 \prod_{i=1}^4\|B_i\|_1.}
 \label{eq:V41-fourfold-compression}
\end{equation}
The estimate is uniform in the number and dimensions of all intermediate
and final LR multiplicity spaces.
\end{theorem}

\begin{proof}
Because \(\mathbf K_\tau^\nu\) intertwines the diagonal group action, it
is block diagonal on
\[
 \bigotimes_{i=1}^4H_{R_i}
 =
 \bigoplus_T H_T\otimes M_T.
\]
Apply the fourfold version of the partial-trace identity in
\Cref{lem:V37-partial-trace}.  Pinching
\(\mathcal B\) by the mutually orthogonal final isotypic projections and
then partially tracing the multiplicity carrier gives
\[
 \sum_{T\in\Omega}d_T
 \|\widehat\kappa_{\tau,T}^{\,\alpha}\|_1
 \le
 d_{\rm in}
 \sup_{T\in\Omega}
 \|\mathbf K_{\tau,T}^{\nu}\|_{\rm op}
 \sum_{T\in\Omega}\|P_T\mathcal BP_T\|_1.
\]
The pinching contraction used in
\Cref{thm:V37-partial-trace-assembly} bounds the last sum by
\(\|\mathcal B\|_1=\prod_i\|B_i\|_1\).  Orthogonal compression cannot
increase operator norm, so
\eqref{eq:V41-exact-Wilson-tree-bound} bounds the supremum by
\(C_4w_\tau\).  This proves
\eqref{eq:V41-fourfold-compression} without choosing or summing a
recoupling path.
\end{proof}

\begin{corollary}[The local \([4]\) operator card is closed]
\label{cor:V41-local-four-card}
At a four-row coordinate block, the complete finite-\(s\) fourth
cumulant decomposes into the sixteen MTA tree labels with
representation-uniform valence weights and multiplicity-uniform final
compression.  No fixed-height remainder remains in this local operator
card.
\end{corollary}

\begin{proof}
Use \Cref{thm:V41-one-link-fourth-trees} before final-channel
decomposition and \Cref{thm:V41-fourfold-compression} afterward.
\end{proof}

\section{The spectator conversion left after local closure}
\label{sec:V41-spectators}

The local weight
\[
 w_\tau
 =
 \prod_{ij\in E(\tau)}
 s_\alpha\sqrt{A_{i,e}A_{j,e}}
\]
is the correct contraction-tree weight.  In a larger exact support,
however, MTA requires the repeated normalized marks
\[
 \prod_{i=1}^4
 \vartheta_e(\boldsymbol R_i)^{\deg_\tau(i)}
 =
 \prod_i
 \left(\frac{A_{i,e}}{\Xi(\boldsymbol R_i)}\right)^{\deg_\tau(i)}.
\]
These are not obtained by replacing the square roots in \(w_\tau\)
pointwise.  Private and partially shared spectator coordinates also
contribute Wilson multipliers, and the final denominator depends on the
assembled output across the full support union.

\begin{assessment}[Exact post-V41 bottleneck]
\label{ass:V41-spectator-bottleneck}
The fixed-bank and representation-height difficulties in the local
fourth cumulant are closed by
\Cref{cor:V41-local-four-card}.  The remaining \([4]\) problem is now a
spectator-bank quotient: combine the sixteen local weights with the
actual spectator multipliers, the total input masses
\(\Xi(\boldsymbol R_i)\), and the single final resolvent so that every
tree incidence receives its repeated BFA marks.  Applying
\eqref{eq:V41-exact-Wilson-tree-bound} and the resolvent separately
would lose precisely those marks and is forbidden.
\end{assessment}

\section{V41 ledger and next acquisition}
\label{sec:V41-ledger}

\begin{theorem}[Integrated V41 statement]
\label{thm:V41-integrated}
Preserving V1--V40, V41 proves:
\begin{enumerate}[label=\textup{(V41.\arabic*)},leftmargin=5.5em]
\item the exact operator-valued Gaussian fourth-cumulant tree formula
      \eqref{eq:V41-Gaussian-tree-formula};
\item uniform derivative and tree-valence bounds
      \eqref{eq:V41-unitary-derivatives}--%
      \eqref{eq:V41-Gaussian-tree-bound};
\item the four-displacement Wilson--Gaussian comparison
      \eqref{eq:V41-Wilson-Gaussian-comparison};
\item max-star absorption of the complete non-Gaussian perturbation;
\item the exact finite-\(s\), representation-uniform decomposition
      \eqref{eq:V41-exact-Wilson-tree-decomposition}--%
      \eqref{eq:V41-exact-Wilson-tree-bound};
\item compatibility with the V40 coherent high-weight regression; and
\item multiplicity-uniform fourfold final-channel compression
      \eqref{eq:V41-fourfold-compression}.
\end{enumerate}
\end{theorem}

\begin{proof}
Items \textup{(V41.1)--(V41.2)} are
\Cref{lem:V41-Gaussian-tree-formula,lem:V41-unitary-derivatives}.
Items \textup{(V41.3)--(V41.4)} are
\Cref{lem:V41-Wilson-Gaussian-comparison,lem:V41-max-star}.  Item
\textup{(V41.5)} is
\Cref{thm:V41-one-link-fourth-trees}.  Item \textup{(V41.6)} is
\Cref{ass:V41-high-weight-resolution}.  Item \textup{(V41.7)} is
\Cref{thm:V41-fourfold-compression}.
\end{proof}

\begingroup
\small
\begin{longtable}{>{\raggedright\arraybackslash}p{0.29\textwidth}
                  >{\raggedright\arraybackslash}p{0.15\textwidth}
                  >{\raggedright\arraybackslash}p{0.48\textwidth}}
\toprule
\textbf{V41 row}&\textbf{Status}&\textbf{Advance or obstruction}\\
\midrule
\endfirsthead
\toprule
\textbf{V41 row}&\textbf{Status}&\textbf{Advance or obstruction}\\
\midrule
\endhead
Finite-\(s\) one-coordinate tree interpolation
& \MGclosed
& Gaussian forest interpolation and max-star perturbation routing give
  sixteen exact operator fibres.\\[0.35em]
Representation-uniform quartic valence
& \MGclosed
& Every tree fibre carries
  \(s^3\prod_iA_i^{\deg(i)/2}\), including the sharp \(s^3N^3\)
  degree-three branch.\\[0.35em]
Quartic LR multiplicity compression
& \MGclosed
& Isotypic pinching and partial trace sum all final paths without a
  multiplicity factor.\\[0.35em]
Embedded \([4]\) spectator quotient
& \MGopen
& The local tree weights must be combined with all spectator multipliers
  and the final resolvent to recover repeated normalized link marks.\\[0.35em]
Mixed \([3,2]\), \([3,3]\), and \([2,2,2]\) fibres
& \MGopen
& Their local blocks are classified but their joint spectator quotients
  have not been assembled.\\[0.35em]
Global \(m=4\) and all-order MTA
& \MGopen
& These remain downstream of the spectator conversion and the subsequent
  uniform induction.\\[0.35em]
Full Yang--Mills mass gap
& \MGopen
& WUFC/CPM and all later source, capacity, OS-continuum, and spectral
  gates remain open.\\
\bottomrule
\end{longtable}
\endgroup

\begin{openproblem}[V42 mandate: the embedded four-junction spectator
quotient]
\label{op:V42}
\begin{enumerate}[label=\textup{(V42.\arabic*)},leftmargin=5.5em]
\item Fix one tree term in
      \eqref{eq:V41-exact-Wilson-tree-decomposition} at a coordinate
      \(e\), but retain every other coordinate moment and the complete
      final output multiplier.  Write the resulting resolved spectator
      quotient before taking absolute values.
\item Split spectator coordinates by their row-incidence subsets.
      Assemble private banks, pair-shared banks, and triple-shared banks
      separately; do not replace them by one raw support count.
\item Use private-bank exponential damping when one
      \(\Xi(\boldsymbol R_i)/A_{i,e}\) is large.  In the complementary
      balanced case, use the final total resolvent and the V39 universal
      one-link tail to recover the repeated factors
      \(\vartheta_e(\boldsymbol R_i)^{\deg_\tau(i)}\).
\item Apply \Cref{thm:V41-fourfold-compression} before any sum over
      intermediate or final LR paths.
\item Prove the embedded \([4]\) MTA row uniformly for all sixteen
      labeled trees, or return an explicit spectator incidence and
      representation family violating the required marked quotient.
\item Only after \([4]\) closes, proceed to the \([3,2]\), \([3,3]\),
      and \([2,2,2]\) topology cards.  Keep the two-triple card explicit.
\item Preserve only the newest YM file.  The next mandatory companion
      audit is V45.
\end{enumerate}
\end{openproblem}

\begin{firewall}[End-of-V41 claim boundary]
V41 proves an exact finite-\(s\), representation-uniform
contraction-tree decomposition of the one-link fourth cumulant and
multiplicity-uniform final-channel compression.  It does not prove the
spectator-bank conversion for an embedded four-junction, the remaining
quartic topology cards, global \(m=4\) MTA, all-order MTA, WUFC, CPM,
the nonlinear Perron packet, finite-edge margins, capacity,
protected-sector floors, local-algebra noncollapse, reflection
positivity, a nontrivial continuum OS state, or a positive Yang--Mills
spectral gap.  The full Yang--Mills mass gap is not proved.
\end{firewall}

\section*{Version V41 append-only record}
\addcontentsline{toc}{section}{Version V41 append-only record}
The complete V40 source was retained before this continuation.  Its
SHA--256 digest was
\begin{center}\scriptsize\ttfamily
90dcfb38a08d12c3526a97167f262194a5a5d5ca76f11b4b78b6f7819595939a.
\end{center}
No mathematical content of V40 was altered.  On the next iteration,
retain this full file, append before its unique active document-closing
command, increment the filename to \texttt{\_V42.tex}, and remove only
the superseded V41 file from this Yang--Mills note series.  The next
scheduled SM/NS audit is V45.

% =============================================================================
% END APPENDED VERSION V41 MATERIAL
% =============================================================================

% =============================================================================
% BEGIN APPENDED VERSION V42 MATERIAL
% =============================================================================

\clearpage
\part{V42 --- Cubic private-bank damping and the
private-dominant embedded four-junction}
\label{part:V42}

\section{Purpose and nonduplication}
\label{sec:V42-purpose}

V41 closed the local analytic part of a four-row coordinate block: its
complete finite-\(s\) cumulant is the sum of sixteen
representation-uniform contraction-tree operators, and every one of
those operators admits multiplicity-free final-channel compression.
What remained was global.  When the four inputs have spectator links,
the MTA norm asks for repeated normalized marks at the junction, whereas
the local tree theorem contains unnormalized half-valence Casimir
powers.

V42 first proves the third multiplier moment missing from V31--V32.
It then writes the exact replica factorization when one coordinate is the
unique link common to all four inputs.  Private spectator links factor as
one scalar product multiplier; pair- and triple-shared spectators remain
assembled moment contractions.  On the sector where the private mass
dominates both the final mass and the two excess input-mass powers, the
quadratic product moment pays the low-output resolvent and the new cubic
product moment pays the high-output resolvent.  This closes a genuine
mixed-support \([4]\) sector without assuming coefficient factorization.

\section{A cubic Wilson multiplier moment}
\label{sec:V42-cubic-multiplier}

\begin{lemma}[Third Laplacian norm of the Wilson density]
\label{lem:V42-density-trilaplacian}
There are constants \(C_3,\alpha_3<\infty\), depending only on SU(3),
such that
\begin{equation}
 \|\Delta_G^3q_\alpha\|_{L^1(G)}
 \le C_3\alpha^3
 \qquad(\alpha\ge\alpha_3).
 \label{eq:V42-density-trilaplacian}
\end{equation}
Consequently every irreducible \(R\) satisfies
\begin{equation}
 C_2(R)^3\eta_R(\alpha)\le C_3\alpha^3.
 \label{eq:V42-character-third-moment}
\end{equation}
\end{lemma}

\begin{proof}
Recall \(q_\alpha=Z_\alpha^{-1}e^{-\alpha W}\), with
\(\|\nabla W\|^2\le CW\), all derivatives of \(W\) bounded on the
compact group, and
\[
 \int_GW^kq_\alpha\,\dd U\le C_k\alpha^{-k}
\]
from \eqref{eq:V32-W-moments}.  Apply three Laplacians to
\(e^{-\alpha W}\).  The product rule gives a finite sum of terms
\[
 q_\alpha\,\alpha^j P_j
 \]
where \(P_j\) is a product of derivatives of \(W\) of total differential
order six.  Every pair of first derivatives is bounded by \(CW\).
If a term contains \(r\) first-derivative factors, its absolute value is
bounded by \(CW^{r/2}\), and its integral is at most
\(C\alpha^{j-r/2}\).  The differential-order count gives
\(j-r/2\le3\); all remaining derivatives of \(W\) are bounded.  Thus
every term has \(L^1\) norm at most \(C\alpha^3\), proving
\eqref{eq:V42-density-trilaplacian}.

Centrality gives
\(\widehat q_\alpha(R)=\eta_RI_{d_R}\).  Integrate by parts three times
and use
\((-\Delta_G)^3D^R=C_2(R)^3D^R\).  Unitarity of \(D^R\) yields
\[
 C_2(R)^3\eta_R
 \le\|\Delta_G^3q_\alpha\|_1,
\]
which is \eqref{eq:V42-character-third-moment}.
\end{proof}

\begin{theorem}[Cubic tail debit and product third moment]
\label{thm:V42-cubic-tail-debit}
There are \(K_{\rm tail,3},K_{\rm prod,3},K_{\rm bal,3}<\infty\)
and \(\alpha_{\rm tail,3}\) such that
\begin{equation}
 \boxed{
 \bigl(s_\alpha C_2(R)\bigr)^3\eta_R(\alpha)
 \le K_{\rm tail,3}\bigl(1-\eta_R(\alpha)\bigr)}
 \label{eq:V42-one-link-cubic-debit}
\end{equation}
for every irreducible \(R\).  If
\(\boldsymbol R=(R_f)_{f\in E}\) is an exact nontrivial product channel
and \(q_{\alpha,\boldsymbol R}=\prod_f\eta_{R_f}\), then
\begin{align}
 \bigl(s_\alpha C_E(\boldsymbol R)\bigr)^3
 q_{\alpha,\boldsymbol R}
 &\le K_{\rm prod,3},
 \label{eq:V42-product-third-moment}\\
 \boxed{
 \bigl(s_\alpha\Xi(\boldsymbol R)\bigr)^3
 q_{\alpha,\boldsymbol R}
 \le K_{\rm bal,3}.}
 \label{eq:V42-balanced-third-moment}
\end{align}
All constants are independent of height and support size.
\end{theorem}

\begin{proof}
Put \(x=s_\alpha C_2(R)\) and \(r=\eta_R\).  If \(x\le1\), then
\(x^3r\le x\), and the lower character gap
\eqref{eq:V11-one-minus-eta}, together with
\(s_\alpha\asymp\alpha^{-1}\), gives
\(x\le C(1-r)\).  If \(x\ge1\), the same gap is bounded below by a
positive constant, while
\eqref{eq:V42-character-third-moment} and
\(s_\alpha\alpha\le C\) give \(x^3r\le C\).  This proves
\eqref{eq:V42-one-link-cubic-debit}.

For the product channel write
\(x_f=s_\alpha C_2(R_f)\), \(r_f=\eta_{R_f}\), and
\(q=\prod_fr_f\).  Expand
\begin{align*}
 \left(\sum_fx_f\right)^3q
 &=
 \sum_fx_f^3r_f\prod_{g\ne f}r_g\\
 &\quad+
 3\sum_{f\ne g}x_f^2r_fx_gr_g
        \prod_{h\ne f,g}r_h\\
 &\quad+
 6\sum_{f<g<h}x_fr_fx_gr_gx_hr_h
        \prod_{k\ne f,g,h}r_k.
\end{align*}
Apply respectively the cubic debit
\eqref{eq:V42-one-link-cubic-debit}, the quadratic debit
\eqref{eq:V32-one-link-quadratic-debit}, and the linear debit
\eqref{eq:V31-one-link-tail-debit}.  The three resulting sums are
constant multiples of the probabilities of exactly one, two, and three
failures in independent Bernoulli trials with success probabilities
\((r_f)\).  Each is at most a fixed factorial, proving
\eqref{eq:V42-product-third-moment}.  Exact support and the fundamental
Casimir floor give
\(\Xi(\boldsymbol R)\le(1+C_F^{-1})C_E(\boldsymbol R)\), which proves
\eqref{eq:V42-balanced-third-moment}.
\end{proof}

\begin{remark}[Why order three is the exact new debit]
\label{rem:V42-third-debit}
For a tree on four vertices,
\[
 \sum_{i=1}^4(\deg_\tau(i)-1)=2.
\]
After the high-output resolvent contributes one final mass, payment of
those two inverse input-mass powers requires three powers of the private
thermal mass.  The V32 quadratic product moment is sufficient in the
low-output branch, where the resolvent supplies \(s_\alpha^{-1}\), but
the high-output branch requires
\eqref{eq:V42-balanced-third-moment}.
\end{remark}

\section{Exact factorization at a unique four-row junction}
\label{sec:V42-unique-junction}

Let \(X_1,\ldots,X_4\) be nonempty exact supports and assume
\begin{equation}
 X_1\cap X_2\cap X_3\cap X_4=\{e\}.
 \label{eq:V42-unique-junction}
\end{equation}
For \(f\ne e\), put
\[
 I_f:=\{i:f\in X_i\}.
\]
Define the private banks and their union by
\begin{equation}
 P_i:=X_i\setminus\bigcup_{j\ne i}X_j,
 \qquad
 P:=\mathop{\dot\bigcup}_{i=1}^4P_i,
 \label{eq:V42-private-banks}
\end{equation}
and set
\begin{equation}
 Y:=\Xi_P:=\sum_{i=1}^4\sum_{f\in P_i}
       \bigl(1+C_2(R_{i,f})\bigr),
 \qquad
 q_P:=\prod_{i=1}^4\prod_{f\in P_i}\eta_{R_{i,f}}.
 \label{eq:V42-private-mass}
\end{equation}

At a shared spectator \(f\), let
\begin{equation}
 \mathbf Q_{I_f,f}
 :=
 \mathbb E_{\nu_\alpha}
 \bigotimes_{i\in I_f}D^{R_{i,f}}(U_f).
 \label{eq:V42-spectator-moment}
\end{equation}
It is a contraction.  At the junction \(e\), retain the tree operators
\(\mathbf K_{\tau,e}^{\nu}\) from
\eqref{eq:V41-exact-Wilson-tree-decomposition}.

\begin{theorem}[Unique-junction replica factorization]
\label{thm:V42-unique-junction-factorization}
The sum of all terms in the connected replica formula
\eqref{eq:V39-connected-replica-formula} whose coordinate partition at
\(e\) is the four-row block equals
\begin{equation}
 \boxed{
 \mathbf J_e
 =
 q_P\,
 \mathbf K_{1234,e}^{\nu}
 \otimes
 \bigotimes_{\substack{f\ne e\\|I_f|\ge2}}
 \mathbf Q_{I_f,f}.}
 \label{eq:V42-four-junction-factorization}
\end{equation}
It has the exact tree refinement
\begin{equation}
 \boxed{
 \mathbf J_e
 =
 \sum_{\tau\in\mathfrak T_4}\mathbf J_{\tau,e},
 \qquad
 \mathbf J_{\tau,e}
 :=
 q_P\,\mathbf K_{\tau,e}^{\nu}
 \otimes
 \bigotimes_{\substack{f\ne e\\|I_f|\ge2}}
 \mathbf Q_{I_f,f}.}
 \label{eq:V42-four-junction-tree-factorization}
\end{equation}
For
\[
 A_{i,e}:=1+C_2(R_{i,e}),
 \qquad
 w_{\tau,e}:=
 s_\alpha^3\prod_iA_{i,e}^{\deg_\tau(i)/2},
\]
one has
\begin{equation}
 \|\mathbf J_{\tau,e}\|_{\rm op}
 \le C_4q_Pw_{\tau,e}.
 \label{eq:V42-junction-operator-bound}
\end{equation}
\end{theorem}

\begin{proof}
Fix the single four-row block at \(e\).  At every other coordinate,
summing all local cumulant partitions gives the complete local moment by
the moment--cumulant formula.  If \(f\in P_i\), only row \(i\) is active,
so this moment is the scalar multiplier
\(\eta_{R_{i,f}}\) on that tensor factor.  Their product is \(q_P\).
If \(|I_f|\ge2\), the local moment is exactly
\eqref{eq:V42-spectator-moment}.  Coordinate independence gives
\eqref{eq:V42-four-junction-factorization}.  The uniqueness condition
\eqref{eq:V42-unique-junction} ensures that no second coordinate can
carry another four-row block in this topology.

Insert the exact V41 tree decomposition at \(e\) to obtain
\eqref{eq:V42-four-junction-tree-factorization}.  Every shared spectator
moment is an average of unitary tensor products and hence has operator
norm at most one.  Equation
\eqref{eq:V41-exact-Wilson-tree-bound} now gives
\eqref{eq:V42-junction-operator-bound}.
\end{proof}

\begin{theorem}[Global multiplicity compression for a junction tree]
\label{thm:V42-global-junction-compression}
Let \(\widehat{\mathbf J}_{\tau,e;\boldsymbol T}\) be the coefficient of
\(\mathbf J_{\tau,e}\) in a fully resolved product output channel
\(\boldsymbol T\).  For arbitrary trace-class input coefficient matrices
\(B_i\) and every collection \(\Omega\) of final product channels,
\begin{equation}
 \boxed{
 \sum_{\boldsymbol T\in\Omega}
 d_{\boldsymbol T}
 \|\widehat{\mathbf J}_{\tau,e;\boldsymbol T}\|_1
 \le
 C_4d_{\rm in}\,q_Pw_{\tau,e}
 \prod_{i=1}^4\|B_i\|_1.}
 \label{eq:V42-global-junction-compression}
\end{equation}
No multiplicity from any pair-, triple-, or four-row spectator channel
appears.
\end{theorem}

\begin{proof}
All factors in
\eqref{eq:V42-four-junction-tree-factorization} intertwine their
diagonal group actions.  Decompose the full product carrier
coordinatewise into final irreducibles and multiplicity spaces.  Pinch
the input tensor \(B_1\otimes\cdots\otimes B_4\) by the mutually
orthogonal product-isotypic projections and partially trace every
multiplicity carrier.  The same argument as
\Cref{thm:V41-fourfold-compression} bounds the sum of all pinched trace
norms by \(\prod_i\|B_i\|_1\).  Insert
\eqref{eq:V42-junction-operator-bound}; the Fourier normalization gives
the factor \(d_{\rm in}\).  This proves
\eqref{eq:V42-global-junction-compression}.
\end{proof}

\section{Private-dominant payment of the embedded junction}
\label{sec:V42-private-payment}

For a fixed tree \(\tau\), abbreviate
\[
 d_i:=\deg_\tau(i),
 \qquad
 \Xi_i:=\Xi(\boldsymbol R_i),
 \qquad
 \Xi_T:=\Xi(\boldsymbol T).
\]
The degree identity gives
\begin{equation}
 \sum_{i=1}^4(d_i-1)=2.
 \label{eq:V42-degree-excess-two}
\end{equation}

\begin{theorem}[Private-dominant four-junction scalar payment]
\label{thm:V42-private-dominant-scalar}
Fix \(D\ge1\).  Suppose \(Y>0\) and the resolved output and branching
input masses satisfy
\begin{equation}
 \Xi_T\le DY,
 \qquad
 \prod_{i=1}^4\Xi_i^{d_i-1}\le DY^2.
 \label{eq:V42-private-dominance}
\end{equation}
Then
\begin{equation}
 \boxed{
 \frac{\Xi_T}{1-q_{\alpha,\boldsymbol T}}\,
 q_Pw_{\tau,e}
 \le
 C_D
 \prod_{i=1}^4
 A_{i,e}^{d_i}\Xi_i^{\,1-d_i}.}
 \label{eq:V42-private-dominant-payment}
\end{equation}
The estimate is uniform in supports, representations, final channels,
and \(\alpha\).
\end{theorem}

\begin{proof}
The two-sided product gap and exact-support Casimir floor give the V33
resolvent split
\begin{equation}
 \frac{\Xi_T}{1-q_{\alpha,\boldsymbol T}}
 \le C
 \begin{cases}
 s_\alpha^{-1},&s_\alpha\Xi_T\le c_0,\\
 \Xi_T,&s_\alpha\Xi_T>c_0.
 \end{cases}
 \label{eq:V42-output-resolvent-split}
\end{equation}

In the low-output case, use
\eqref{eq:V32-balanced-second-moment} on the exact private product
channel:
\begin{align*}
 \frac{\Xi_T}{1-q_{\alpha,\boldsymbol T}}q_Pw_{\tau,e}
 &\le
 Cq_Ps_\alpha^2
 \prod_iA_{i,e}^{d_i/2}\\
 &\le
 \frac{C}{Y^2}\prod_iA_{i,e}^{d_i/2}.
\end{align*}
The second condition in
\eqref{eq:V42-private-dominance} gives
\[
 \frac1{Y^2}
 \le
 \frac D{\prod_i\Xi_i^{d_i-1}}.
\]
Since every \(A_{i,e}\ge1\),
\(A_{i,e}^{d_i/2}\le A_{i,e}^{d_i}\), proving
\eqref{eq:V42-private-dominant-payment} in this case.

In the high-output case, the second line of
\eqref{eq:V42-output-resolvent-split}, the first condition in
\eqref{eq:V42-private-dominance}, and
\eqref{eq:V42-balanced-third-moment} give
\begin{align*}
 \frac{\Xi_T}{1-q_{\alpha,\boldsymbol T}}q_Pw_{\tau,e}
 &\le
 CDYq_Ps_\alpha^3
 \prod_iA_{i,e}^{d_i/2}\\
 &=
 \frac{CD}{Y^2}(s_\alpha Y)^3q_P
 \prod_iA_{i,e}^{d_i/2}\\
 &\le
 \frac{C_D}{Y^2}
 \prod_iA_{i,e}^{d_i/2}.
\end{align*}
Apply the same branching-mass and local-Casimir comparisons.  This proves
the theorem.
\end{proof}

\begin{theorem}[Private-dominant unique-\([4]\) MTA sector]
\label{thm:V42-private-dominant-four-MTA}
For each tree \(\tau\), let
\(\Omega_{\tau,e}^{\rm priv}(D)\) be the set of nonconstant resolved
outputs for which both inequalities in
\eqref{eq:V42-private-dominance} hold.  Under
\eqref{eq:V42-unique-junction}, the private-dominant sub-sum of the
tree-refined four-junction contribution obeys the corresponding \(m=4\)
MTA terms:
\begin{align}
 &\sum_{\tau\in\mathfrak T_4}
 \sum_{\boldsymbol T\in\Omega_{\tau,e}^{\rm priv}(D)}
 \left\|
 \left[
 \mathcal R_{\alpha,\Lambda}Q_0
 \mathbf J_{\tau,e}(g_{1,X_1},\ldots,g_{4,X_4})
 \right]_{\boldsymbol T}
 \right\|_{\rho,2}^{\rm BFA}
 \notag\\
 &\quad\le
 C_D
 \sum_{\tau\in\mathfrak T_4}
 \prod_{i=1}^4
 \|g_{i,X_i}\|_{
 \rho,2;(\,\underbrace{e,\ldots,e}_{d_i}\,)}^{
 \rm BFA(d_i)}.
 \label{eq:V42-private-dominant-four-MTA}
\end{align}
The statement allows arbitrary coefficient entanglement and all actual
LR multiplicities.
\end{theorem}

\begin{proof}
Apply \Cref{thm:V42-global-junction-compression} and then
\Cref{thm:V42-private-dominant-scalar} on every resolved final channel.
The exponential height ratio is at most one because each final
coordinate representation occurs in the tensor product of its active
inputs.  Divide the scalar right side
\eqref{eq:V42-private-dominant-payment} by the four input balanced
masses.  It becomes
\[
 \prod_{i=1}^4
 \left(\frac{A_{i,e}}{\Xi_i}\right)^{d_i}
 =
 \prod_i\vartheta_e(\boldsymbol R_i)^{d_i},
\]
which is exactly the coefficient weight in the repeated marked BFA
seminorms.  Sum the nonnegative input Fourier trace-norm terms and all
sixteen trees.  Pinching has already summed final multiplicities, so no
additional path count occurs.
\end{proof}

\begin{assessment}[Residual after private-dominant \([4]\) closure]
\label{ass:V42-four-residual}
A unique-junction \([4]\) witness not covered by
\Cref{thm:V42-private-dominant-four-MTA} must violate at least one
condition in \eqref{eq:V42-private-dominance}.  Thus either its final
mass is not private-dominated, or the product of its two branching
input-mass powers is concentrated on pair- and triple-shared spectator
banks.  This is a shared-formation problem; it is no longer a local
fourth-cumulant, fixed-height, private-tail, or LR-multiplicity problem.
\end{assessment}

\section{V42 ledger and next acquisition}
\label{sec:V42-ledger}

\begin{theorem}[Integrated V42 statement]
\label{thm:V42-integrated}
Preserving V1--V41, V42 proves:
\begin{enumerate}[label=\textup{(V42.\arabic*)},leftmargin=5.5em]
\item the third Wilson-density Laplacian bound
      \eqref{eq:V42-density-trilaplacian};
\item the one-link cubic debit and uniform product third moments
      \eqref{eq:V42-one-link-cubic-debit}--%
      \eqref{eq:V42-balanced-third-moment};
\item the exact unique-junction replica and tree factorizations
      \eqref{eq:V42-four-junction-factorization}--%
      \eqref{eq:V42-four-junction-tree-factorization};
\item global multiplicity-uniform compression of every junction tree;
\item the private-dominant scalar payment
      \eqref{eq:V42-private-dominant-payment}; and
\item the corresponding mixed-support \(m=4\) MTA sector
      \eqref{eq:V42-private-dominant-four-MTA}.
\end{enumerate}
\end{theorem}

\begin{proof}
Item \textup{(V42.1)} is
\Cref{lem:V42-density-trilaplacian}.  Item \textup{(V42.2)} is
\Cref{thm:V42-cubic-tail-debit}.  Item \textup{(V42.3)} is
\Cref{thm:V42-unique-junction-factorization}.  Item
\textup{(V42.4)} is
\Cref{thm:V42-global-junction-compression}.  Items
\textup{(V42.5)--(V42.6)} are
\Cref{thm:V42-private-dominant-scalar,
thm:V42-private-dominant-four-MTA}.
\end{proof}

\begingroup
\small
\begin{longtable}{>{\raggedright\arraybackslash}p{0.29\textwidth}
                  >{\raggedright\arraybackslash}p{0.15\textwidth}
                  >{\raggedright\arraybackslash}p{0.48\textwidth}}
\toprule
\textbf{V42 row}&\textbf{Status}&\textbf{Advance or obstruction}\\
\midrule
\endfirsthead
\toprule
\textbf{V42 row}&\textbf{Status}&\textbf{Advance or obstruction}\\
\midrule
\endhead
Cubic product multiplier moment
& \MGclosed
& Three Laplacians and the Bernoulli failure expansion give
  \((s_\alpha\Xi)^3q\lesssim1\) uniformly.\\[0.35em]
Unique-\([4]\) spectator factorization
& \MGclosed
& Private links factor scalarly; all pair- and triple-shared spectator
  moments remain assembled contractions.\\[0.35em]
Private-dominant embedded \([4]\) sector
& \MGclosed
& Quadratic damping pays low outputs and cubic damping pays high outputs,
  yielding the exact repeated junction marks.\\[0.35em]
Shared-dominant unique-\([4]\) residual
& \MGopen
& Final or branching mass is concentrated on pair- and triple-shared
  spectator banks and requires an actual formation-kernel estimate.\\[0.35em]
Multiple four-shared junction coordinates
& \MGopen
& Their overlapping four-block choices require an assembled bank
  interpolation rather than a raw sum of local cumulants.\\[0.35em]
\([3,2]\), \([3,3]\), and \([2,2,2]\) fibres
& \MGopen
& Their local topology is known, but the quartic marked quotient is not
  yet closed.\\[0.35em]
Global \(m=4\), all-order MTA, and mass gap
& \MGopen
& The downstream WUFC/CPM, source, capacity, OS-continuum, and spectral
  gates remain unproved.\\
\bottomrule
\end{longtable}
\endgroup

\begin{openproblem}[V43 mandate: shared spectator formation at a unique
four-junction]
\label{op:V43}
\begin{enumerate}[label=\textup{(V43.\arabic*)},leftmargin=5.5em]
\item Retain the exact factorization
      \eqref{eq:V42-four-junction-tree-factorization}.  Split every
      shared spectator by its incidence set
      \(I_f\subset\{1,2,3,4\}\), with \(|I_f|=2\) or \(3\).
\item Construct the actual product-dimension formation kernel for all
      resolved shared spectator outputs and prove it is subprobability
      before inserting any tail estimate.
\item On each tree \(\tau\), identify which pair- or triple-shared bank
      carries the two branching mass powers when
      \eqref{eq:V42-private-dominance} fails.  Do not replace this by a
      worst LR multiplicity.
\item Use \Cref{cor:V39-universal-one-link-tail} on thermally cold pair
      outputs and the V41 contraction-tree bound on genuine triple
      junctions.  High outputs must be paid by their own final multiplier
      and the single total resolvent.
\item Prove the shared-dominant unique-\([4]\) row, or return a resolved
      incidence/representation family which violates the marked
      quotient.  Go upstream if the same formation output is being
      charged twice.
\item Only after the unique-junction \([4]\) topology closes, construct
      an assembled interpolation for multiple four-shared coordinates
      and then attack \([3,2]\), \([3,3]\), and \([2,2,2]\).
\item Preserve only the newest YM file.  The next scheduled SM/NS audit
      is V45.
\end{enumerate}
\end{openproblem}

\begin{firewall}[End-of-V42 claim boundary]
V42 proves cubic Wilson product damping, exact unique-junction
factorization, multiplicity-uniform spectator compression, and the
private-dominant embedded \([4]\) MTA sector.  It does not prove the
shared-dominant or multiple-junction \([4]\) sectors, the other three
quartic topologies, global \(m=4\) MTA, all-order MTA, WUFC, CPM, the
nonlinear Perron packet, finite-edge margins, capacity, protected-sector
floors, local-algebra noncollapse, reflection positivity, a nontrivial
continuum OS state, or a positive Yang--Mills spectral gap.  The full
Yang--Mills mass gap is not proved.
\end{firewall}

\section*{Version V42 append-only record}
\addcontentsline{toc}{section}{Version V42 append-only record}
The complete V41 source was retained before this continuation.  Its
SHA--256 digest was
\begin{center}\scriptsize\ttfamily
a14f353d3189aefe28235ca703d83a89349aeb1698c3826570fafa0765ae1f56.
\end{center}
No mathematical content of V41 was altered.  On the next iteration,
retain this full file, append before its unique active document-closing
command, increment the filename to \texttt{\_V43.tex}, and remove only
the superseded V42 file from this Yang--Mills note series.  The next
scheduled SM/NS audit is V45.

% =============================================================================
% END APPENDED VERSION V42 MATERIAL
% =============================================================================

% =============================================================================
% BEGIN APPENDED VERSION V43 MATERIAL
% =============================================================================

\clearpage
\part{V43 --- Shared spectator formation, fourth multiplier damping,
and pair-heavy \([4]\) closure}
\label{part:V43}

\section{Purpose and the exact residual}
\label{sec:V43-purpose}

V42 factored a unique four-row junction from all spectator moments and
closed the sector in which private pass-through mass pays both branching
mass powers.  Its pointwise estimate deliberately discarded the
multipliers of pair- and triple-shared spectator outputs.  V43 restores
them.

The shared CG/LR outputs carry an actual dimension-normalized formation
kernel.  It is subprobability even for coefficient matrices entangled
across all links.  The combined private and shared spectator output is
itself an exact product channel, so its multiplier has the quadratic and
cubic damping proved earlier.  V43 adds the fourth product moment.  This
extra moment is essential: a squared cold-formation tail combined with
only cubic pointwise damping leaves a logarithmic dyadic sum, whereas
fourth damping makes the high-output shells summable.

The result closes every spectator-dominated unique-\([4]\) fibre for
which one pair-shared spectator link carries the two branching mass
powers.  The remaining branch has no such pair link; its mass is spread
over a bank or lies on triple-shared coordinates, or the local junction
output dominates the final mass.

\section{Fourth Wilson multiplier damping}
\label{sec:V43-fourth-damping}

\begin{lemma}[Fourth Laplacian norm and character moment]
\label{lem:V43-density-four-laplacian}
There are group constants \(C_4,\alpha_4<\infty\) such that
\begin{align}
 \|\Delta_G^4q_\alpha\|_{L^1(G)}
 &\le C_4\alpha^4,
 \label{eq:V43-density-four-laplacian}\\
 C_2(R)^4\eta_R(\alpha)
 &\le C_4\alpha^4
 \label{eq:V43-character-fourth-moment}
\end{align}
for every irreducible \(R\) and \(\alpha\ge\alpha_4\).
\end{lemma}

\begin{proof}
Apply four Laplacians to
\(q_\alpha=Z_\alpha^{-1}e^{-\alpha W}\).  A term with \(j\) factors of
derivatives of \(W\), of which \(r\) are first derivatives, is bounded
in absolute value by
\[
 C\alpha^jW^{r/2}q_\alpha.
\]
The total differential order is eight.  Since every derivative factor
not counted among the first derivatives has order at least two,
\(8\ge r+2(j-r)=2j-r\), hence \(j-r/2\le4\).  The Wilson-well moments
\eqref{eq:V32-W-moments} therefore bound the integral of this term by
\(C\alpha^{j-r/2}\le C\alpha^4\).  There are finitely many product-rule
terms, which proves \eqref{eq:V43-density-four-laplacian}.

Integrate by parts four times in the central Fourier transform and use
\((-\Delta_G)^4D^R=C_2(R)^4D^R\).  Unitarity of \(D^R\) gives
\eqref{eq:V43-character-fourth-moment}.
\end{proof}

\begin{theorem}[Fourth tail debit and product fourth moment]
\label{thm:V43-fourth-tail-debit}
Uniformly in irreducible \(R\),
\begin{equation}
 \boxed{
 \bigl(s_\alpha C_2(R)\bigr)^4\eta_R(\alpha)
 \le K_{\rm tail,4}\bigl(1-\eta_R(\alpha)\bigr).}
 \label{eq:V43-one-link-fourth-debit}
\end{equation}
For every exact nontrivial product channel \(\boldsymbol R\),
\begin{align}
 \bigl(s_\alpha C_E(\boldsymbol R)\bigr)^4
 q_{\alpha,\boldsymbol R}
 &\le K_{\rm prod,4},
 \label{eq:V43-product-fourth-moment}\\
 \boxed{
 \bigl(s_\alpha\Xi(\boldsymbol R)\bigr)^4
 q_{\alpha,\boldsymbol R}
 \le K_{\rm bal,4}.}
 \label{eq:V43-balanced-fourth-moment}
\end{align}
The constants are independent of representation height and support size.
\end{theorem}

\begin{proof}
Put \(x=s_\alpha C_2(R)\) and \(r=\eta_R\).  For \(x\le1\),
\(x^4r\le x\le C(1-r)\) by the uniform character gap.  For \(x\ge1\),
\eqref{eq:V43-character-fourth-moment} and
\(s_\alpha\alpha\le C\) give \(x^4r\le C\), while \(1-r\) is bounded
below.  This proves \eqref{eq:V43-one-link-fourth-debit}.

Expand
\((\sum_fx_f)^4\prod_fr_f\), where
\(x_f=s_\alpha C_2(R_f)\) and \(r_f=\eta_{R_f}\), and group monomials
by their set of \(k\) distinct link indices, \(1\le k\le4\).  At a link
appearing with multiplicity \(a\), use the order-\(a\) one-link debit:
orders one and two are
\eqref{eq:V31-one-link-tail-debit} and
\eqref{eq:V32-one-link-quadratic-debit}, order three is
\eqref{eq:V42-one-link-cubic-debit}, and order four is
\eqref{eq:V43-one-link-fourth-debit}.  Each grouped sum is then bounded
by a multinomial constant times
\[
 \sum_{\substack{F\subset E\\|F|=k}}
 \prod_{f\in F}(1-r_f)\prod_{g\notin F}r_g,
\]
the probability of exactly \(k\) failures in independent Bernoulli
trials.  Summing \(k=1,\ldots,4\) proves
\eqref{eq:V43-product-fourth-moment}.  The exact-support balanced-mass
equivalence proves \eqref{eq:V43-balanced-fourth-moment}.
\end{proof}

\section{The actual shared spectator formation kernel}
\label{sec:V43-formation-kernel}

Continue under the unique-junction hypothesis
\eqref{eq:V42-unique-junction}.  For every
\(J\subset\{1,2,3,4\}\) with \(2\le|J|\le3\), define the incidence bank
\begin{equation}
 E_J:=\{f\ne e:I_f=J\}.
 \label{eq:V43-incidence-banks}
\end{equation}
At each \(f\in E_J\), resolve
\begin{equation}
 \bigotimes_{i\in J}H_{R_{i,f}}
 =
 \bigoplus_{S_f}H_{S_f}\otimes M_{S_f}^{\,f}.
 \label{eq:V43-local-spectator-CG}
\end{equation}
Let \(\boldsymbol S=(S_f)_{|I_f|\ge2}\) and
\(\boldsymbol\mu\) select one copy in each multiplicity carrier.
Apply the product CG unitary to the full input coefficient tensor and
pinch to this resolved shared-spectator block; denote the normalized
Fourier coefficient block, in the same trace-norm convention as the
input matrices, by
\(\mathcal C_{\boldsymbol S,\boldsymbol\mu}\).

Put
\begin{align}
 d_{\rm sh,in}
 &:=
 \prod_{\substack{f\ne e\\|I_f|\ge2}}
 \prod_{i\in I_f}d_{R_{i,f}},
 &
 d_{\boldsymbol S}
 &:=
 \prod_{\substack{f\ne e\\|I_f|\ge2}}d_{S_f}.
 \label{eq:V43-shared-dimensions}
\end{align}
For nonzero input trace norm define
\begin{equation}
 \pi_{\boldsymbol S,\boldsymbol\mu}
 :=
 \frac{d_{\boldsymbol S}
 \|\mathcal C_{\boldsymbol S,\boldsymbol\mu}\|_1}
 {d_{\rm sh,in}\prod_{i=1}^4\|B_i\|_1}.
 \label{eq:V43-spectator-formation-kernel}
\end{equation}
Inactive, private, and junction carrier dimensions cancel in this ratio.
If all spectator coordinates except \(f\) have been resolved, write
\(\mathcal C_{\rm par}\) for the resulting parent block and
\(\mathcal C_{{\rm par};S_f,\mu_f}\) for its child after resolving \(f\).
For \(\|\mathcal C_{\rm par}\|_1>0\), define
\begin{equation}
 \pi_f(S_f,\mu_f\mid\mathscr F_{\ne f})
 :=
 \frac{d_{S_f}
 \|\mathcal C_{{\rm par};S_f,\mu_f}\|_1}
 {\left(\prod_{i\in I_f}d_{R_{i,f}}\right)
  \|\mathcal C_{\rm par}\|_1};
 \label{eq:V43-conditional-formation-kernel}
\end{equation}
for a zero parent block set every conditional weight equal to zero.
The global dimension ratio in
\eqref{eq:V43-spectator-formation-kernel} factors coordinatewise, so
successive use of
\eqref{eq:V43-conditional-formation-kernel} is the exact
disintegration of the global formation weights.

\begin{theorem}[Shared spectator formation is subprobability]
\label{thm:V43-formation-subprobability}
For arbitrary coefficient entanglement,
\begin{equation}
 \boxed{
 \pi_{\boldsymbol S,\boldsymbol\mu}\ge0,
 \qquad
 \sum_{\boldsymbol S,\boldsymbol\mu}
 \pi_{\boldsymbol S,\boldsymbol\mu}\le1.}
 \label{eq:V43-formation-subprobability}
\end{equation}
The conditional weights
\eqref{eq:V43-conditional-formation-kernel} also sum to at most one
after all outputs except the chosen spectator coordinate have been
pinched.
\end{theorem}

\begin{proof}
The product CG map is unitary.  Pinching by its mutually orthogonal
resolved projections is trace-norm contractive, so the sum of the raw
block trace norms is at most \(\prod_i\|B_i\|_1\).  For every resolved
copy,
\[
 \frac{d_{\boldsymbol S}}{d_{\rm sh,in}}\le1,
\]
because each \(H_{S_f}\) embeds isometrically in the corresponding local
tensor-product carrier.  Multiplying each pinched block by this ratio
and summing proves \eqref{eq:V43-formation-subprobability}.  If all but
one coordinate blocks are fixed first, the remaining map is still an
orthogonal CG pinching of the arbitrary trace-class parent matrix
\(\mathcal C_{\rm par}\).  Its child trace norms sum to at most
\(\|\mathcal C_{\rm par}\|_1\), and
\(d_{S_f}/\prod_{i\in I_f}d_{R_{i,f}}\le1\).  This proves the
conditional statement, including the zero-parent convention.
\end{proof}

For a resolved spectator output define
\begin{align}
 M(\boldsymbol S)
 &:=
 Y+
 \sum_{\substack{f\ne e,\ |I_f|\ge2\\S_f\ne\boldsymbol1}}
 \bigl(1+C_2(S_f)\bigr),
 \label{eq:V43-spectator-output-mass}\\
 q_{\rm sp}(\boldsymbol S)
 &:=
 q_P
 \prod_{\substack{f\ne e\\|I_f|\ge2}}\eta_{S_f}.
 \label{eq:V43-spectator-output-multiplier}
\end{align}
This is the balanced mass and product multiplier of the exact spectator
output channel.  If the junction output is \(T_e\), then
\begin{equation}
 \Xi_T=M(\boldsymbol S)+
 \boldsymbol1_{\{T_e\ne\boldsymbol1\}}
 \bigl(1+C_2(T_e)\bigr),
 \qquad
 q_{\alpha,\boldsymbol T}
 =q_{\rm sp}(\boldsymbol S)\eta_{T_e}.
 \label{eq:V43-full-output-split}
\end{equation}

\section{Pointwise payment by the resolved spectator output}
\label{sec:V43-output-payment}

For a tree \(\tau\), retain
\[
 d_i=\deg_\tau(i),\qquad
 B_\tau:=\prod_{i=1}^4\Xi_i^{d_i-1},\qquad
 P_{\tau,e}:=
 \prod_{i=1}^4A_{i,e}^{d_i}\Xi_i^{1-d_i}.
\]
Both \(B_\tau\) and \(P_{\tau,e}\) have exactly two excess input-mass
powers.

\begin{theorem}[Spectator-output-dominant scalar payment]
\label{thm:V43-output-dominant-payment}
Fix \(D\ge1\).  On every resolved output satisfying
\begin{equation}
 \Xi_T\le D M(\boldsymbol S),
 \qquad
 B_\tau\le D M(\boldsymbol S)^2,
 \label{eq:V43-output-dominance}
\end{equation}
one has
\begin{equation}
 \boxed{
 \frac{\Xi_T}{1-q_{\alpha,\boldsymbol T}}\,
 q_{\rm sp}(\boldsymbol S)w_{\tau,e}
 \le C_D P_{\tau,e}.}
 \label{eq:V43-output-dominant-payment}
\end{equation}
\end{theorem}

\begin{proof}
The proof of
\Cref{thm:V42-private-dominant-scalar} used only that the damping mass
and multiplier form one exact product channel.  Apply the same
low/high final-output split with
\((Y,q_P)\) replaced by
\((M(\boldsymbol S),q_{\rm sp}(\boldsymbol S))\).
In the low branch,
\eqref{eq:V32-balanced-second-moment} gives
\[
 q_{\rm sp}s_\alpha^2\le C M^{-2}.
\]
In the high branch, the first inequality in
\eqref{eq:V43-output-dominance} and
\eqref{eq:V42-balanced-third-moment} give
\[
 \Xi_Tq_{\rm sp}s_\alpha^3\le C_D M^{-2}.
\]
The second dominance inequality converts \(M^{-2}\) to
\(C_DB_\tau^{-1}\).  Finally
\(\prod_iA_{i,e}^{d_i/2}\le\prod_iA_{i,e}^{d_i}\).
This is exactly \eqref{eq:V43-output-dominant-payment}.
\end{proof}

\begin{corollary}[The V42 private sector is strictly contained]
\label{cor:V43-private-contained}
Every output satisfying \eqref{eq:V42-private-dominance} also satisfies
\eqref{eq:V43-output-dominance}, after changing \(D\) by a fixed power,
because \(M(\boldsymbol S)\ge Y\).  Thus
\Cref{thm:V43-output-dominant-payment} retains the V42 closure and adds
all outputs whose shared spectator mass supplies the damping.
\end{corollary}

\section{Scale-free pair formation and weighted shell summation}
\label{sec:V43-pair-heavy}

\begin{lemma}[Conditional scale-free pair-output tail]
\label{lem:V43-scale-free-pair-tail}
Fix a pair-shared spectator link \(f\in E_{\{i,j\}}\), put
\[
 H_f:=A_{i,f}+A_{j,f},
\]
and condition on every other resolved spectator block.  For every
\(x\ge1\), the conditional actual formation kernel obeys
\begin{equation}
 \boxed{
 \sum_{\substack{S_f,\mu_f\\1+C_2(S_f)\le x}}
 \pi_f(S_f,\mu_f\mid\mathrm{other\ outputs})
 \le
 C\min\left\{1,\frac{x^2}{H_f^2}\right\}.}
 \label{eq:V43-scale-free-pair-tail}
\end{equation}
The constant is independent of the conditioned coefficient matrix.
\end{lemma}

\begin{proof}
With all other blocks fixed, the coefficient at \(f\) is an arbitrary
trace-class matrix on
\(H_{R_{i,f}}\otimes H_{R_{j,f}}\).
The amplified Schur-tax estimate used in
\eqref{eq:V35-resolved-two-bank-tax} bounds each resolved copy by its
actual dimension-law mass, uniformly in that matrix.  Sum the copies and
apply the universal one-link tail
\eqref{eq:V39-universal-one-link-tail} with
\(t=x^{-1}\) and \(L=1\):
\[
 \sum_{S_f,\mu_f:\,1+C_2(S_f)\le x}
 \frac{d_{S_f}}{d_{R_{i,f}}d_{R_{j,f}}}
 \le
 \frac C{(1+H_f/x)^2}.
\]
This is bounded by the right side of
\eqref{eq:V43-scale-free-pair-tail}.  The conditional pinching
normalization is
\Cref{thm:V43-formation-subprobability}.
\end{proof}

\begin{lemma}[Squared lower tail plus fourth damping]
\label{lem:V43-weighted-tail-summation}
Let \(\pi\) be a subprobability law on pairs \((M,q)\), with
\(M\ge0\), \(0\le q\le1\).  Suppose, for \(x\ge1\),
\begin{align}
 \pi\{M\le x\}
 &\le C_0\min\{1,x^2/H^2\},
 \label{eq:V43-abstract-squared-tail}\\
 (sM)q&\le K_1,\qquad
 (sM)^4q\le K_4
 \quad(M>0).
 \label{eq:V43-abstract-damping}
\end{align}
Then, for every fixed \(c,D>0\),
\begin{equation}
 \int\left[
 s^2\boldsymbol1_{\{sM\le c\}}
 +s^3Mq\,\boldsymbol1_{\{sM>c\}}
 \right]\dd\pi
 \le\frac{C_{c,D,C_0,K_1,K_4}}{H^2}.
 \label{eq:V43-weighted-tail-summation}
\end{equation}
\end{lemma}

\begin{proof}
Changing constants reduces to \(c=1\).  If \(sH\le2\), the linear
damping in \eqref{eq:V43-abstract-damping} bounds both summands by
\(Cs^2\le C/H^2\).

Suppose \(sH>2\).  The low part is
\[
 s^2\pi\{M\le s^{-1}\}
 \le \frac C{H^2}
\]
by \eqref{eq:V43-abstract-squared-tail}.  For \(k\ge0\), on the shell
\[
 2^ks^{-1}<M\le2^{k+1}s^{-1},
\]
the fourth damping gives
\[
 s^3Mq\le\frac{C}{sM^3}\le Cs^2\,2^{-3k}.
\]
As long as \(2^{k+1}s^{-1}\le H\), the shell probability is at most
\(C4^{k+1}/(s^2H^2)\).  Its contribution is therefore at most
\(C2^{-k}/H^2\), and these terms sum geometrically.  Once the upper shell
endpoint exceeds \(H\), use probability at most one.  The first such
shell contributes at most \(C/(sH^3)\le C/H^2\), and the later shells
again form a geometric series.  This proves
\eqref{eq:V43-weighted-tail-summation}.
\end{proof}

\begin{theorem}[Pair-heavy spectator-dominated \([4]\) sector]
\label{thm:V43-pair-heavy-four-MTA}
Fix \(D,\delta>0\), a tree \(\tau\), and select deterministically a
pair-shared link \(f\in E_{\{i,j\}}\) satisfying
\begin{equation}
 H_f^2\ge\delta B_\tau.
 \label{eq:V43-pair-heavy}
\end{equation}
Then all resolved outputs with
\begin{equation}
 \Xi_T\le D M(\boldsymbol S)
 \label{eq:V43-spectator-dominates-final}
\end{equation}
obey the averaged marked-tree estimate
\begin{align}
 &\sum_{\boldsymbol S,\boldsymbol\mu}
 \pi_{\boldsymbol S,\boldsymbol\mu}
 \sup_{\substack{T_e:\,\Xi_T\le DM(\boldsymbol S)}}
 \frac{\Xi_Tq_{\rm sp}(\boldsymbol S)w_{\tau,e}}
 {1-q_{\alpha,\boldsymbol T}}
 \notag\\
 &\qquad\le C_{D,\delta}P_{\tau,e}.
 \label{eq:V43-pair-heavy-scalar-average}
\end{align}
Here and below the supremum of an empty admissible set is zero.  The
sum in \eqref{eq:V43-pair-heavy-scalar-average} is only over spectator
formation blocks; it is not an unweighted sum over junction outputs.
After the compression
\eqref{eq:V42-global-junction-compression}, this is the \(\tau\)-term
of \(m=4\) MTA with every tree edge marked by the unique junction \(e\).
\end{theorem}

\begin{proof}
The event \(M(\boldsymbol S)\le x\) implies
\(1+C_2(S_f)\le x\), with the trivial output included when \(x\ge1\).
Condition on every other output and use
\Cref{lem:V43-scale-free-pair-tail}; then integrate the conditioning
with the subprobability kernel.  This gives
\[
 \pi\{M\le x\}
 \le C\min\{1,x^2/H_f^2\}.
\]
The spectator output channel has mass \(M\), multiplier \(q_{\rm sp}\),
and satisfies the linear and fourth product damping inequalities
\eqref{eq:V31-balanced-product-damping} and
\eqref{eq:V43-balanced-fourth-moment}.

Uniformly over every \(T_e\) satisfying
\eqref{eq:V43-spectator-dominates-final}, the final resolvent split
\eqref{eq:V42-output-resolvent-split} bounds the scalar quotient by
\[
 C_D\prod_iA_{i,e}^{d_i/2}
 \left[
 s_\alpha^2\boldsymbol1_{\{s_\alpha M\le c_D\}}
 +s_\alpha^3M q_{\rm sp}
  \boldsymbol1_{\{s_\alpha M>c_D\}}
 \right].
\]
Apply \Cref{lem:V43-weighted-tail-summation} with \(H=H_f\), and then
\eqref{eq:V43-pair-heavy}:
\[
 \text{left side of \eqref{eq:V43-pair-heavy-scalar-average}}
 \le
 \frac C{H_f^2}\prod_iA_{i,e}^{d_i/2}
 \le
 \frac C{B_\tau}\prod_iA_{i,e}^{d_i}
 =CP_{\tau,e}.
\]
This proves the scalar estimate, including the displayed supremum.
For the operator statement, first apply the global pinching theorem to
the actual local-junction output blocks.  Its trace-norm weights perform
the \(T_e\)-sum; no bare junction multiplicity is introduced.  The
uniform supremum just proved may then be pulled outside that weighted
sum.  Divide by the four input balanced masses, use
\(P_{\tau,e}/\prod_i\Xi_i
=\prod_i\vartheta_e(\boldsymbol R_i)^{d_i}\), and apply the global
pinching theorem exactly as in
\Cref{thm:V42-private-dominant-four-MTA}.  This is the required MTA
term.
\end{proof}

\begin{assessment}[Residual after pair-heavy closure]
\label{ass:V43-residual}
For a unique four-row junction, V43 closes:
\begin{enumerate}[label=\textup{(\roman*)},leftmargin=3.7em]
\item every pointwise output with
      \(\Xi_T\lesssim M\) and \(B_\tau\lesssim M^2\); and
\item even when \(M^2\ll B_\tau\), every spectator-dominated output for
      which one pair-shared input link has
      \(H_f^2\gtrsim B_\tau\).
\end{enumerate}
A remaining witness must therefore have local-junction-dominant final
mass, or must distribute its branching mass over many pair-shared links
and/or triple-shared banks so that no single pair link satisfies
\eqref{eq:V43-pair-heavy}.  The obstruction is now a bank concentration
problem, not a missing output kernel or a logarithmic shell sum.
\end{assessment}

\section{V43 ledger and next acquisition}
\label{sec:V43-ledger}

\begin{theorem}[Integrated V43 statement]
\label{thm:V43-integrated}
Preserving V1--V42, V43 proves:
\begin{enumerate}[label=\textup{(V43.\arabic*)},leftmargin=5.5em]
\item the fourth Wilson-density Laplacian and character moment;
\item the fourth one-link debit and uniform product damping
      \eqref{eq:V43-balanced-fourth-moment};
\item the actual incidence-resolved spectator formation kernel and its
      conditional subprobability property;
\item pointwise MTA payment by total resolved spectator output mass;
\item the scale-free conditional pair-output tail
      \eqref{eq:V43-scale-free-pair-tail};
\item summable coupling of a squared formation tail to fourth multiplier
      damping; and
\item the pair-heavy spectator-dominated unique-\([4]\) MTA sector
      \eqref{eq:V43-pair-heavy-scalar-average}.
\end{enumerate}
\end{theorem}

\begin{proof}
Items \textup{(V43.1)--(V43.2)} are
\Cref{lem:V43-density-four-laplacian,
thm:V43-fourth-tail-debit}.  Item \textup{(V43.3)} is
\Cref{thm:V43-formation-subprobability}.  Item \textup{(V43.4)} is
\Cref{thm:V43-output-dominant-payment}.  Item \textup{(V43.5)} is
\Cref{lem:V43-scale-free-pair-tail}.  Item \textup{(V43.6)} is
\Cref{lem:V43-weighted-tail-summation}.  Item \textup{(V43.7)} is
\Cref{thm:V43-pair-heavy-four-MTA}.
\end{proof}

\begingroup
\small
\begin{longtable}{>{\raggedright\arraybackslash}p{0.29\textwidth}
                  >{\raggedright\arraybackslash}p{0.15\textwidth}
                  >{\raggedright\arraybackslash}p{0.48\textwidth}}
\toprule
\textbf{V43 row}&\textbf{Status}&\textbf{Advance or obstruction}\\
\midrule
\endfirsthead
\toprule
\textbf{V43 row}&\textbf{Status}&\textbf{Advance or obstruction}\\
\midrule
\endhead
Fourth product multiplier moment
& \MGclosed
& \((s_\alpha\Xi)^4q\lesssim1\) removes the logarithmic shell loss left
  by squared formation tails and cubic damping.\\[0.35em]
Shared spectator formation kernel
& \MGclosed
& Product CG pinching gives an actual conditional subprobability law
  without coefficient factorization or a bare LR count.\\[0.35em]
Spectator-output-dominant pointwise sector
& \MGclosed
& The resolved spectator mass and multiplier replace the private mass in
  the V42 quadratic/cubic payment.\\[0.35em]
Pair-heavy spectator-dominated sector
& \MGclosed
& The V39 scale-free squared tail and fourth damping sum all cold and hot
  output shells with no logarithm.\\[0.35em]
Spread-pair and triple-shared spectator banks
& \MGopen
& No individual pair link carries the two branching powers; a bank-level
  concentration or sequential triple-output tail is required.\\[0.35em]
Local-junction-dominant outputs
& \MGopen
& The final mass is concentrated at \(e\), where the tree operator must
  be combined with the local final channel before using the resolvent.\\[0.35em]
Global \(m=4\), all-order MTA, and full mass gap
& \MGopen
& Multiple four-junctions, the other quartic skeletons, WUFC/CPM, and all
  continuum/spectral gates remain downstream.\\
\bottomrule
\end{longtable}
\endgroup

\begin{openproblem}[V44 mandate: bank concentration and local-output
dominance]
\label{op:V44}
\begin{enumerate}[label=\textup{(V44.\arabic*)},leftmargin=5.5em]
\item For each incidence bank \(E_J\), construct its exact product
      representation formation law.  Prove a bank-level analogue of
      \eqref{eq:V43-scale-free-pair-tail} without multiplying linkwise
      Schur taxes for entangled coefficients.
\item If pair-shared branching mass is spread over many links, use a
      dimension-law concentration or one extracted heavy link with its
      true logarithmic loss.  Retain the fourth multiplier moment when
      summing output shells.
\item For \(|J|=3\), resolve one pair intermediate, keep its actual
      subprobability kernel, and prove a sequential low-output tail for
      the triple moment.  Do not assume that a cold final triple output
      forces a cold first intermediate.
\item On the branch \(\Xi_T\gg M\), combine the local final-channel
      mass with the V41 tree block before using the single resolvent.
      Test every orientation on the V40 coherent family.
\item Close the remaining unique-junction \([4]\) sector or return the
      first resolved bank family showing that the squared formation tail
      is insufficient even with fourth damping.
\item Do not begin multiple-junction interpolation or another quartic
      skeleton until this unique-junction row is settled.
\item Preserve only the newest YM file.  V45 must begin with the
      scheduled SM/NS companion audit.
\end{enumerate}
\end{openproblem}

\begin{firewall}[End-of-V43 claim boundary]
V43 proves fourth Wilson product damping, the actual shared spectator
formation kernel, total-output pointwise payment, and the pair-heavy
spectator-dominated unique-\([4]\) MTA sector.  It does not prove the
spread-pair, triple-shared, or local-junction-dominant residuals, the
complete unique- or multiple-junction \([4]\) row, the other quartic
topologies, global \(m=4\) MTA, all-order MTA, WUFC, CPM, the nonlinear
Perron packet, finite-edge margins, capacity, protected-sector floors,
local-algebra noncollapse, reflection positivity, a nontrivial
continuum OS state, or a positive Yang--Mills spectral gap.  The full
Yang--Mills mass gap is not proved.
\end{firewall}

\section*{Version V43 append-only record}
\addcontentsline{toc}{section}{Version V43 append-only record}
The complete V42 source was retained before this continuation.  Its
SHA--256 digest was
\begin{center}\scriptsize\ttfamily
fb927bcabf2e25d89a3b3562e24c39972fbc630b96156dd6a5abd3be81678300.
\end{center}
No mathematical content of V42 was altered.  On the next iteration,
retain this full file, append before its unique active document-closing
command, increment the filename to \texttt{\_V44.tex}, and remove only
the superseded V43 file from this Yang--Mills note series.  V45 begins
with the mandatory SM/NS companion audit.

% =============================================================================
% END APPENDED VERSION V43 MATERIAL
% =============================================================================

% =============================================================================
% BEGIN APPENDED VERSION V44 MATERIAL
% =============================================================================

\clearpage
\part{V44 --- Entangled pair-bank concentration and the
triple/local residual}
\label{part:V44}

\section{Purpose and acquisition}
\label{sec:V44-purpose}

V43 used the squared cold-output tail at one pair-shared spectator link.
That estimate closes a branch on which one \(H_f^2\) pays the two excess
tree masses, but it does not address a bank containing many individually
light links.  Extracting the largest link would lose the bank cardinality,
which is not bounded uniformly in the support.

V44 replaces maximum-link extraction by a Laplace argument for the actual
entangled coefficient kernel.  Two facts are coupled.  First, the V39
one-link estimate gives quadratic decay when \(tH_f\) is large.  Second,
an exact normalized-trace calculation for the tensor-product Casimir gives
a fixed probability of forming an output with Casimir comparable to
\(H_f\); this gives exponential contraction when many \(tH_f\)'s are
small.  A deterministic product lemma splices these regimes and yields a
single squared tail at the \emph{sum} of all pair-bank input masses.

The result closes every spectator-dominated unique-\([4]\) branch whose
two excess tree powers are carried by pair-shared input mass, including
fully diffuse banks.  A mass identity then leaves only a triple-shared
bank or the local four-row junction.  For the triple branch, V44 computes
the corresponding Casimir moments and identifies the missing estimate as
a translated-shell bound; it does not infer that a cold final triple
output has a cold first intermediate.

\section{The pair dimension law and its Casimir moments}
\label{sec:V44-pair-Casimir}

Fix a pair-shared coordinate \(f\in E_{\{i,j\}}\), and abbreviate
\[
 R:=R_{i,f},\qquad Q:=R_{j,f},\qquad
 a:=C_2(R),\qquad b:=C_2(Q),\qquad H_f:=2+a+b.
\]
Both \(R\) and \(Q\) are nontrivial because \(f\) belongs to both exact
supports.  On resolved copies in \(R\otimes Q\), put
\begin{equation}
 \nu_f(S,\mu):=\frac{d_S}{d_Rd_Q},
 \qquad
 Z_f(S):=
 \boldsymbol1_{\{S\ne\boldsymbol1\}}\bigl(1+C_2(S)\bigr).
 \label{eq:V44-pair-dimension-law}
\end{equation}
The representation-dimension identity makes \(\nu_f\) a probability
law.  Compression of one child block in
\eqref{eq:V43-conditional-formation-kernel} gives the pointwise
domination
\begin{equation}
 \pi_f(S,\mu\mid\mathscr F_{\ne f})
 \le \nu_f(S,\mu).
 \label{eq:V44-conditional-dimension-domination}
\end{equation}
Indeed, the child compression has trace norm at most that of its parent.
This observation uses neither factorization of the parent coefficient nor
independence of coordinate blocks.

\begin{lemma}[Exact first two pair-Casimir moments]
\label{lem:V44-pair-Casimir-moments}
For \(G=\mathrm{SU}(3)\), if \(X_f(S):=C_2(S)\), then
\begin{align}
 \mathbb E_{\nu_f}X_f
 &=a+b,
 \label{eq:V44-pair-Casimir-first}\\
 \mathbb E_{\nu_f}X_f^2
 &=(a+b)^2+\frac12ab
 \le\frac98(a+b)^2.
 \label{eq:V44-pair-Casimir-second}
\end{align}
Consequently,
\begin{equation}
 \nu_f\{Z_f\ge H_f/2\}\ge\frac29.
 \label{eq:V44-pair-warm-probability}
\end{equation}
\end{lemma}

\begin{proof}
Choose an orthonormal basis \(T_1,\ldots,T_8\) of the Lie algebra in the
Casimir convention and write
\[
 \sum_{\ell=1}^8(T_\ell^R)^2=aI_R,
 \qquad
 \sum_{\ell=1}^8(T_\ell^Q)^2=bI_Q.
\]
On \(H_R\otimes H_Q\), the diagonal Casimir is
\[
 \mathsf C=(a+b)I+2\mathsf X,
 \qquad
 \mathsf X:=\sum_{\ell=1}^8T_\ell^R\otimes T_\ell^Q.
\]
Every generator is traceless.  Invariance of normalized trace gives
\[
 \frac1{d_R}\Tr(T_\ell^RT_m^R)=\frac a8\delta_{\ell m},
 \qquad
 \frac1{d_Q}\Tr(T_\ell^QT_m^Q)=\frac b8\delta_{\ell m}.
\]
Therefore the normalized tensor-product trace satisfies
\[
 \tau(\mathsf X)=0,
 \qquad
 \tau(\mathsf X^2)=\frac{ab}{8}.
\]
The diagonal decomposition of \(R\otimes Q\) turns normalized trace of a
function of \(\mathsf C\) into expectation under \(\nu_f\).  Expanding
\(\mathsf C^2\) proves
\eqref{eq:V44-pair-Casimir-first}--%
\eqref{eq:V44-pair-Casimir-second}; the last inequality uses
\(4ab\le(a+b)^2\).

Apply the Paley--Zygmund inequality to the nonnegative random variable
\(X_f\), with threshold one half of its mean:
\[
 \nu_f\{X_f\ge(a+b)/2\}
 \ge
 \frac{(1-1/2)^2(\mathbb E X_f)^2}
 {\mathbb E X_f^2}
 \ge\frac29.
\]
On this event \(S\) is nontrivial and
\[
 Z_f=1+X_f\ge1+\frac{a+b}{2}=\frac{H_f}{2},
\]
which proves \eqref{eq:V44-pair-warm-probability}.
\end{proof}

\section{One-link Laplace contraction}
\label{sec:V44-one-link-Laplace}

\begin{lemma}[Two-regime conditional Laplace estimate]
\label{lem:V44-one-link-Laplace}
There are constants \(c_0,C_0>0\), depending only on
\(\mathrm{SU}(3)\), such that for every \(0<t\le1\) and every nonzero
conditioned parent block,
\begin{equation}
 L_f(t\mid\mathscr F_{\ne f})
 :=
 \sum_{S,\mu}
 \pi_f(S,\mu\mid\mathscr F_{\ne f})e^{-tZ_f(S)}
 \le \Phi(tH_f),
 \label{eq:V44-one-link-Laplace}
\end{equation}
where
\begin{equation}
 \Phi(z):=
 \min\left\{
 \exp[-c_0\min\{z,1\}],
 \frac{C_0}{(1+z)^2}
 \right\}.
 \label{eq:V44-Laplace-envelope}
\end{equation}
The same bound holds with value zero for a zero parent block.
\end{lemma}

\begin{proof}
By \eqref{eq:V44-conditional-dimension-domination} and
\eqref{eq:V44-pair-warm-probability}, the conditional mass of
\(\{Z_f<H_f/2\}\) is at most \(7/9\).  Its total mass is at most one.
Thus
\begin{align*}
 L_f(t\mid\mathscr F_{\ne f})
 &\le \frac79+\frac29e^{-tH_f/2}\\
 &=1-\frac29(1-e^{-tH_f/2})
 \le \exp[-c_0\min\{tH_f,1\}]
\end{align*}
after fixing a universal \(c_0>0\).

For the polynomial branch, V39 and
\eqref{eq:V44-conditional-dimension-domination} give, for \(x\ge1\),
\begin{equation}
 \pi_f\{Z_f\le x\mid\mathscr F_{\ne f}\}
 \le
 \frac C{(1+H_f/x)^2}
 \le C\min\{1,x^2/H_f^2\}.
 \label{eq:V44-conditional-pair-CDF}
\end{equation}
At \(x=1\), this also bounds the trivial-output atom by
\(C(1+H_f)^{-2}\).  There is no positive value of \(Z_f\) in
\((0,1]\).  Stieltjes integration therefore yields
\begin{align*}
 L_f(t\mid\mathscr F_{\ne f})
 &=
 \pi_f\{Z_f=0\mid\mathscr F_{\ne f}\}\\
 &\quad+
 \int_1^\infty
 t e^{-tx}
 \pi_f\{0<Z_f\le x\mid\mathscr F_{\ne f}\}\,\dd x\\
 &\le
 \frac C{(1+H_f)^2}
 +C\int_1^\infty
 t e^{-tx}\min\{1,x^2/H_f^2\}\,\dd x.
\end{align*}
With \(u=tx\), the integral is at most
\[
 C\int_0^\infty e^{-u}
 \min\{1,u^2/(tH_f)^2\}\,\dd u
 \le\frac C{(1+tH_f)^2}.
\]
Because \(t\le1\), the atom obeys the same bound.  Combining the
exponential and polynomial estimates proves
\eqref{eq:V44-one-link-Laplace}.
\end{proof}

\begin{lemma}[Product-envelope compression]
\label{lem:V44-product-envelope}
For the function \(\Phi\) in
\eqref{eq:V44-Laplace-envelope}, there is \(C_\Phi<\infty\) such that
for every finite family \(z_1,\ldots,z_N\ge0\),
\begin{equation}
 \boxed{
 \prod_{r=1}^N\Phi(z_r)
 \le
 \frac{C_\Phi}{(1+\sum_{r=1}^Nz_r)^2}.}
 \label{eq:V44-product-envelope-compression}
\end{equation}
The empty product is one.
\end{lemma}

\begin{proof}
Let \(I_0=\{r:z_r<1\}\), \(I_1=\{r:z_r\ge1\}\), and put
\[
 Z_0:=\sum_{I_0}z_r,\qquad Z_1:=\sum_{I_1}z_r.
\]
The exponential branch of \(\Phi\) gives
\[
 \prod_{I_0}\Phi(z_r)\le e^{-c_0Z_0}
 \le\frac C{(1+Z_0)^2}.
\]
If \(I_1\) is empty, this already suffices.  Otherwise write
\(n=|I_1|\), choose \(r_*\) with
\(z_{r_*}=\max_{I_1}z_r\), and put \(\rho=e^{-c_0}<1\).
The two branches of \(\Phi\) give
\[
 \prod_{I_1}\Phi(z_r)
 \le
 \frac{C_0}{z_{r_*}^2}\rho^{\,n-1}
 \le
 \frac{C_0n^2\rho^{\,n-1}}{Z_1^2}
 \le\frac C{(1+Z_1)^2}.
\]
Here \(\sup_{n\ge1}n^2\rho^{n-1}<\infty\) and \(Z_1\ge1\).
Finally,
\[
 (1+Z_0)(1+Z_1)\ge1+Z_0+Z_1.
\]
Multiplying the two bounds proves
\eqref{eq:V44-product-envelope-compression}.
\end{proof}

\section{A squared tail for an entire entangled pair bank}
\label{sec:V44-pair-bank}

Let \(\mathcal F\) be any nonempty collection of pair-shared spectator
coordinates.  Define
\begin{equation}
 H_{\mathcal F}:=\sum_{f\in\mathcal F}H_f,
 \qquad
 Z_{\mathcal F}(\boldsymbol S):=
 \sum_{f\in\mathcal F}Z_f(S_f).
 \label{eq:V44-pair-bank-masses}
\end{equation}

\begin{theorem}[Entangled pair-bank Laplace and cold-tail bounds]
\label{thm:V44-pair-bank-tail}
Condition on all resolved spectator outputs outside \(\mathcal F\).
The actual formation kernel satisfies, uniformly in that conditioned
coefficient block,
\begin{align}
 \sum_{\boldsymbol S_{\mathcal F},\boldsymbol\mu_{\mathcal F}}
 \pi_{\mathcal F}
   (\boldsymbol S_{\mathcal F},\boldsymbol\mu_{\mathcal F}
    \mid\mathscr F_{\mathcal F^c})
 e^{-tZ_{\mathcal F}(\boldsymbol S)}
 &\le
 \frac C{(1+tH_{\mathcal F})^2},
 \qquad 0<t\le1,
 \label{eq:V44-pair-bank-Laplace}\\
 \boxed{
 \pi_{\mathcal F}\{Z_{\mathcal F}\le x
       \mid\mathscr F_{\mathcal F^c}\}
 &\le
 C\min\left\{1,\frac{x^2}{H_{\mathcal F}^2}\right\},
 \qquad x\ge1.}
 \label{eq:V44-pair-bank-tail}
\end{align}
No linkwise factorization of the input coefficient is assumed, and the
constant is independent of \(|\mathcal F|\).
\end{theorem}

\begin{proof}
Order the links in \(\mathcal F\).  The exact parent--child
disintegration following
\eqref{eq:V43-conditional-formation-kernel} permits successive
conditional summation even when the original coefficient is entangled.
At each step apply \Cref{lem:V44-one-link-Laplace}.  This gives
\[
 \sum\pi_{\mathcal F}e^{-tZ_{\mathcal F}}
 \le\prod_{f\in\mathcal F}\Phi(tH_f).
\]
\Cref{lem:V44-product-envelope} proves
\eqref{eq:V44-pair-bank-Laplace}.

For \(x\ge1\), take \(t=x^{-1}\).  The pointwise inequality
\[
 \boldsymbol1_{\{Z_{\mathcal F}\le x\}}
 \le e\,e^{-Z_{\mathcal F}/x}
\]
and \eqref{eq:V44-pair-bank-Laplace} give
\[
 \pi_{\mathcal F}\{Z_{\mathcal F}\le x
       \mid\mathscr F_{\mathcal F^c}\}
 \le\frac{Ce}{(1+H_{\mathcal F}/x)^2}.
\]
Combine this with total mass at most one to obtain
\eqref{eq:V44-pair-bank-tail}.
\end{proof}

\begin{remark}[Why entanglement causes no product-law assumption]
\label{rem:V44-entanglement}
The proof multiplies deterministic \emph{upper envelopes}, not
probabilities from an assumed independent law.  Each envelope is applied
to the actual child kernel after all later coordinates have been fixed.
The trace-norm parent denominator telescopes, so the original coefficient
may be arbitrarily entangled across every spectator and junction carrier.
\end{remark}

\section{Closure of the diffuse pair-shared sector}
\label{sec:V44-diffuse-pair-closure}

Let
\begin{align}
 \mathcal F_2
 &:=
 \bigcup_{\substack{J\subset\{1,2,3,4\}\\|J|=2}}E_J,
 &
 H_2&:=\sum_{f\in\mathcal F_2}H_f,
 \label{eq:V44-total-pair-input-mass}\\
 M_2(\boldsymbol S)
 &:=
 \sum_{\substack{f\in\mathcal F_2\\S_f\ne\boldsymbol1}}
 \bigl(1+C_2(S_f)\bigr).
 \label{eq:V44-total-pair-output-mass}
\end{align}
Thus \(M(\boldsymbol S)\ge M_2(\boldsymbol S)\).

\begin{corollary}[Global spectator mass inherits the pair-bank tail]
\label{cor:V44-global-pair-tail}
If \(\mathcal F_2\ne\varnothing\), the full spectator formation kernel
obeys
\begin{equation}
 \boxed{
 \pi\{M(\boldsymbol S)\le x\}
 \le
 C\min\left\{1,\frac{x^2}{H_2^2}\right\},
 \qquad x\ge1.}
 \label{eq:V44-global-pair-tail}
\end{equation}
\end{corollary}

\begin{proof}
Condition on all triple-shared outputs.  The event
\(\{M\le x\}\) is contained in \(\{M_2\le x\}\).  Apply
\eqref{eq:V44-pair-bank-tail} to \(\mathcal F_2\) and integrate the
remaining subprobability conditioning.
\end{proof}

\begin{theorem}[Pair-bank-dominated unique-\([4]\) MTA sector]
\label{thm:V44-pair-bank-four-MTA}
Fix \(D,\delta>0\) and a tree \(\tau\).  Suppose
\begin{equation}
 H_2^2\ge\delta B_\tau.
 \label{eq:V44-pair-bank-dominance}
\end{equation}
Then the spectator-dominated outputs satisfy
\begin{align}
 &\sum_{\boldsymbol S,\boldsymbol\mu}
 \pi_{\boldsymbol S,\boldsymbol\mu}
 \sup_{\substack{T_e:\,\Xi_T\le DM(\boldsymbol S)}}
 \frac{\Xi_Tq_{\rm sp}(\boldsymbol S)w_{\tau,e}}
 {1-q_{\alpha,\boldsymbol T}}
 \notag\\
 &\qquad\le C_{D,\delta}P_{\tau,e}.
 \label{eq:V44-pair-bank-scalar-average}
\end{align}
After global junction compression, this is the corresponding
pair-bank-dominated \(m=4\) MTA term.  It includes arbitrarily many
pair-shared links and requires no heavy individual link.
\end{theorem}

\begin{proof}
The response estimate in the proof of
\Cref{thm:V43-pair-heavy-four-MTA} is uniform in the admissible junction
output and equals
\[
 C_D\prod_iA_{i,e}^{d_i/2}
 \left[
 s_\alpha^2\boldsymbol1_{\{s_\alpha M\le c_D\}}
 +s_\alpha^3Mq_{\rm sp}
  \boldsymbol1_{\{s_\alpha M>c_D\}}
 \right].
\]
Use \eqref{eq:V44-global-pair-tail} in
\Cref{lem:V43-weighted-tail-summation}, with \(H=H_2\), and then
\eqref{eq:V44-pair-bank-dominance}:
\[
 \text{left side of \eqref{eq:V44-pair-bank-scalar-average}}
 \le
 \frac C{H_2^2}\prod_iA_{i,e}^{d_i/2}
 \le
 \frac C{B_\tau}\prod_iA_{i,e}^{d_i}
 =CP_{\tau,e}.
\]
The passage to the operator estimate is exactly the weighted
junction-output compression explained after
\eqref{eq:V43-pair-heavy-scalar-average}.
\end{proof}

\begin{corollary}[The maximum-link hypothesis is removed]
\label{cor:V44-max-link-removed}
\Cref{thm:V44-pair-bank-four-MTA} strictly contains
\Cref{thm:V43-pair-heavy-four-MTA}: if one link has
\(H_f^2\ge\delta B_\tau\), then \(H_2^2\ge H_f^2\); but
\eqref{eq:V44-pair-bank-dominance} also allows
\(\max_fH_f/H_2\to0\).
\end{corollary}

\section{Deterministic localization to triple or local mass}
\label{sec:V44-residual-localization}

Define the total triple-shared input mass and the junction input mass by
\begin{align}
 H_3
 &:=
 \sum_{\substack{f\ne e\\|I_f|=3}}
 \sum_{i\in I_f}A_{i,f},
 &
 L_e&:=\sum_{i=1}^4A_{i,e}.
 \label{eq:V44-triple-local-input-masses}
\end{align}
The unique-junction incidence decomposition gives the exact identity
\begin{equation}
 \sum_{i=1}^4\Xi_i=Y+H_2+H_3+L_e.
 \label{eq:V44-total-input-mass-identity}
\end{equation}

\begin{lemma}[Two-source localization of the unresolved branching mass]
\label{lem:V44-two-source-localization}
For every four-vertex tree \(\tau\),
\begin{equation}
 B_\tau\le\left(\sum_{i=1}^4\Xi_i\right)^2.
 \label{eq:V44-branching-below-total-square}
\end{equation}
If
\begin{equation}
 B_\tau>16M(\boldsymbol S)^2,
 \qquad
 H_2^2<\frac1{16}B_\tau,
 \label{eq:V44-unresolved-numerical-branch}
\end{equation}
then
\begin{equation}
 \boxed{
 \max\{H_3^2,L_e^2\}\ge\frac1{16}B_\tau.}
 \label{eq:V44-triple-local-dichotomy}
\end{equation}
\end{lemma}

\begin{proof}
The exponents \(d_i-1\) are nonnegative and sum to two by
\eqref{eq:V42-degree-excess-two}.  Hence
\[
 B_\tau=\prod_i\Xi_i^{d_i-1}
 \le(\max_i\Xi_i)^2
 \le\left(\sum_i\Xi_i\right)^2,
\]
proving \eqref{eq:V44-branching-below-total-square}.
Because \(M\ge Y\), equations
\eqref{eq:V44-total-input-mass-identity} and
\eqref{eq:V44-unresolved-numerical-branch} imply
\[
 L_e+H_3
 \ge
 \sqrt{B_\tau}-Y-H_2
 >
 \frac12\sqrt{B_\tau}.
\]
At least one of \(L_e,H_3\) is therefore at least
\(\sqrt{B_\tau}/4\), which is
\eqref{eq:V44-triple-local-dichotomy}.
\end{proof}

\begin{proposition}[Exact residual after the pair-bank closure]
\label{prop:V44-exact-residual}
On the spectator-dominated branch
\(\Xi_T\le DM(\boldsymbol S)\), every unique-junction tree term is
closed by V43--V44 unless at least one of
\begin{equation}
 H_3^2\ge\frac1{16}B_\tau,
 \qquad
 L_e^2\ge\frac1{16}B_\tau
 \label{eq:V44-residual-two-branches}
\end{equation}
holds.  Thus a remaining spectator-dominated witness is genuinely
triple-bank-dominated or local-junction-dominated.
\end{proposition}

\begin{proof}
If \(B_\tau\le16M^2\), apply
\Cref{thm:V43-output-dominant-payment} with enlarged fixed constant.
If \(H_2^2\ge B_\tau/16\), apply
\Cref{thm:V44-pair-bank-four-MTA}.  The complement of these two cases
satisfies \eqref{eq:V44-unresolved-numerical-branch}, so
\Cref{lem:V44-two-source-localization} gives
\eqref{eq:V44-residual-two-branches}.
\end{proof}

\section{Triple formation: proved moments and the translated-shell gate}
\label{sec:V44-triple-gate}

At a triple-shared link let the nontrivial inputs be
\(R_1,R_2,R_3\), put \(a_i=C_2(R_i)\),
\[
 A:=a_1+a_2+a_3,\qquad H^{(3)}:=3+A,
\]
and let
\begin{equation}
 \nu^{(3)}(S,\mu):=
 \frac{d_S}{d_{R_1}d_{R_2}d_{R_3}},
 \qquad
 Z^{(3)}(S):=
 \boldsymbol1_{\{S\ne\boldsymbol1\}}(1+C_2(S)).
 \label{eq:V44-triple-dimension-law}
\end{equation}
The copy index ranges over the complete diagonal decomposition, so this
is a probability law.

\begin{lemma}[Exact first two triple-Casimir moments]
\label{lem:V44-triple-Casimir-moments}
If \(X^{(3)}=C_2(S)\), then
\begin{align}
 \mathbb E_{\nu^{(3)}}X^{(3)}
 &=A,
 \label{eq:V44-triple-first-moment}\\
 \mathbb E_{\nu^{(3)}}(X^{(3)})^2
 &=A^2+\frac12\sum_{1\le i<j\le3}a_ia_j
 \le\frac76A^2.
 \label{eq:V44-triple-second-moment}
\end{align}
In particular,
\begin{equation}
 \nu^{(3)}\{Z^{(3)}\ge(3/7)H^{(3)}\}
 \ge\frac3{14}.
 \label{eq:V44-triple-warm-probability}
\end{equation}
The same cold-event upper bound holds for every actual conditional
triple-coordinate formation kernel.
\end{lemma}

\begin{proof}
On the threefold tensor product,
\[
 \mathsf C=A I+2(\mathsf X_{12}+\mathsf X_{13}+\mathsf X_{23}).
\]
Every mixed normalized trace
\(\tau(\mathsf X_{ij}\mathsf X_{k\ell})\) with
\(\{i,j\}\ne\{k,\ell\}\) vanishes because some tensor factor contains
one traceless generator.  The pair calculation gives
\(\tau(\mathsf X_{ij}^2)=a_ia_j/8\).  Expanding
\(\mathsf C^2\) proves the equalities.  For three nonnegative numbers,
\(\sum_{i<j}a_ia_j\le A^2/3\), giving the last inequality.

Paley--Zygmund at \(A/2\) now gives
\[
 \nu^{(3)}\{X^{(3)}\ge A/2\}
 \ge\frac{1/4}{7/6}=\frac3{14}.
\]
Every nontrivial \(\mathrm{SU}(3)\) irrep has Casimir at least \(4/3\),
so the three active inputs give \(A\ge4\).  Consequently
\[
 1+\frac A2\ge\frac37(3+A).
\]
This proves \eqref{eq:V44-triple-warm-probability}.  Conditional
formation weights are pointwise bounded by the direct threefold
dimension law by the same child-compression argument as
\eqref{eq:V44-conditional-dimension-domination}.
\end{proof}

Resolve the bracketing
\[
 R_1\otimes R_2\longrightarrow U,
 \qquad
 U\otimes R_3\longrightarrow S.
\]
For each resolved path the dimension weights telescope exactly:
\begin{equation}
 \frac{d_U}{d_{R_1}d_{R_2}}\,
 \frac{d_S}{d_Ud_{R_3}}
 =
 \frac{d_S}{d_{R_1}d_{R_2}d_{R_3}}.
 \label{eq:V44-triple-path-telescoping}
\end{equation}
Thus sequential pinching is a valid disintegration of the direct law,
but the final cold event does not localize the first intermediate.

\begin{proposition}[Warm-intermediate regression]
\label{prop:V44-warm-intermediate-regression}
Let
\[
 R_1=(n,0),\qquad R_2=(1,0),\qquad
 U=(n+1,0),\qquad R_3=U^*=(0,n+1),\qquad
 S=\boldsymbol1.
\]
This is a multiplicity-one resolved triple path.  As \(n\to\infty\),
\[
 C_2(U)\asymp n^2,\qquad C_2(S)=0,
\]
and its exact dimension-law mass is
\begin{equation}
 \frac{d_U}{3d_{R_1}}\frac1{d_Ud_{R_3}}
 =
 \frac1{3d_{R_1}d_{R_3}}
 \asymp\frac1{(H^{(3)})^2}.
 \label{eq:V44-warm-intermediate-mass}
\end{equation}
Hence a cold final triple output need not have a cold first
intermediate, and the desired squared triple tail, if true, is sharp.
\end{proposition}

\begin{proof}
The highest-weight summand of
\((n,0)\otimes(1,0)\) is \((n+1,0)\), with multiplicity one.
Every irreducible tensored with its dual contains the singlet once.
This proves occurrence and multiplicity.  The Weyl dimension and
Casimir formulae give
\[
 d_{(n,0)}=\frac{(n+1)(n+2)}2,\qquad
 d_{(0,n+1)}=\frac{(n+2)(n+3)}2,
\]
and both relevant Casimirs are comparable to \(n^2\).  Insert these
identities into \eqref{eq:V44-triple-path-telescoping} to obtain
\eqref{eq:V44-warm-intermediate-mass}.
\end{proof}

\begin{assessment}[The exact triple-bank acquisition]
\label{ass:V44-triple-acquisition}
The moment calculation supplies the small-\(tH^{(3)}\) exponential
branch needed for tensorization, but V39 supplies a scale-free
polynomial branch only for a two-input product.  Applying it twice and
summing over \(U\) produces the borderline quantity
\[
 \sum_U\nu_{12}(U)
 \frac{x^2}{(x+C_2(U)+C_2(R_3))^2}.
\]
Using only the squared lower tail for \(U\), each dyadic shell below
the input scale contributes order \(x^2/(H^{(3)})^2\); the shell count
therefore leaves a logarithm.  This is an insufficiency of the available
estimate, not a counterexample to the desired bound.

The precise missing local statement is the direct three-input
translated-shell estimate
\begin{equation}
 \boxed{
 \sum_{\substack{S,\mu\\Z^{(3)}(S)\le x}}
 \frac{d_S}{d_{R_1}d_{R_2}d_{R_3}}
 \le
 C\min\left\{1,\frac{x^2}{(H^{(3)})^2}\right\},
 \qquad x\ge1.}
 \label{eq:V44-direct-triple-tail-gate}
\end{equation}
Equivalently, one may prove an annular/translated Casimir-shell bound
for \(R_1\otimes R_2\) centered at the dual scale of \(R_3\), strong
enough to sum the intermediate shells without a logarithm.  The family
in \Cref{prop:V44-warm-intermediate-regression} must be retained in that
proof; discarding warm \(U\)'s would remove a contribution of the sharp
order.
\end{assessment}

\section{V44 ledger and next acquisition}
\label{sec:V44-ledger}

\begin{theorem}[Integrated V44 statement]
\label{thm:V44-integrated}
Preserving V1--V43, V44 proves:
\begin{enumerate}[label=\textup{(V44.\arabic*)},leftmargin=5.5em]
\item exact first and second tensor-product Casimir moments for pair and
      triple dimension laws;
\item a uniform warm-output probability for each actual conditional
      pair formation kernel;
\item the two-regime conditional Laplace estimate
      \eqref{eq:V44-one-link-Laplace};
\item cardinality-free product compression
      \eqref{eq:V44-product-envelope-compression};
\item the entangled pair-bank Laplace and squared cold-tail estimates
      \eqref{eq:V44-pair-bank-Laplace}--%
      \eqref{eq:V44-pair-bank-tail};
\item closure of every pair-bank-dominated spectator branch of the
      unique-\([4]\) tree, including diffuse banks; and
\item localization of the remaining spectator-dominated branch to the
      triple-shared or local-junction alternatives
      \eqref{eq:V44-residual-two-branches}.
\end{enumerate}
It also proves the warm-intermediate regression and isolates
\eqref{eq:V44-direct-triple-tail-gate}; that gate is not asserted.
\end{theorem}

\begin{proof}
Items \textup{(V44.1)--(V44.2)} are
\Cref{lem:V44-pair-Casimir-moments,
lem:V44-triple-Casimir-moments}.  Items
\textup{(V44.3)--(V44.4)} are
\Cref{lem:V44-one-link-Laplace,lem:V44-product-envelope}.  Item
\textup{(V44.5)} is \Cref{thm:V44-pair-bank-tail}.  Item
\textup{(V44.6)} is
\Cref{thm:V44-pair-bank-four-MTA,cor:V44-max-link-removed}.  Item
\textup{(V44.7)} is
\Cref{lem:V44-two-source-localization,prop:V44-exact-residual}.
The last sentence is
\Cref{prop:V44-warm-intermediate-regression,
ass:V44-triple-acquisition}.
\end{proof}

\begingroup
\small
\begin{longtable}{>{\raggedright\arraybackslash}p{0.29\textwidth}
                  >{\raggedright\arraybackslash}p{0.15\textwidth}
                  >{\raggedright\arraybackslash}p{0.48\textwidth}}
\toprule
\textbf{V44 row}&\textbf{Status}&\textbf{Advance or obstruction}\\
\midrule
\endfirsthead
\toprule
\textbf{V44 row}&\textbf{Status}&\textbf{Advance or obstruction}\\
\midrule
\endhead
Pair Casimir moments and warm fraction
& \MGclosed
& Normalized tensor trace and Paley--Zygmund give a universal
  \(2/9\) warm-output fraction.\\[0.35em]
Entangled pair-bank squared tail
& \MGclosed
& Conditional Laplace envelopes telescope through arbitrary coefficient
  entanglement and compress at total bank mass.\\[0.35em]
Diffuse pair-bank unique-\([4]\) sector
& \MGclosed
& \(H_2^2\gtrsim B_\tau\) now suffices; no maximum link or support-size
  loss remains.\\[0.35em]
Triple direct cold tail
& \MGopen
& Pairwise iteration is logarithmically borderline; a direct
  three-input translated-shell bound is required.\\[0.35em]
Local-junction-dominant output
& \MGopen
& The local four-input output mass must be combined with the V41 tree
  block before the sole resolvent is used.\\[0.35em]
Complete unique-\([4]\) row
& \MGopen
& It now reduces to the triple-bank and local-junction branches, together
  with the already isolated \(\Xi_T\gg M\) output branch.\\[0.35em]
Global \(m=4\), all-order MTA, and mass gap
& \MGopen
& Other quartic topologies, WUFC/CPM, construction, continuum, and OS
  spectral gates remain downstream.\\
\bottomrule
\end{longtable}
\endgroup

\begin{openproblem}[V45 mandate: companion audit and triple translated
shell]
\label{op:V45}
\begin{enumerate}[label=\textup{(V45.\arabic*)},leftmargin=5.5em]
\item Begin with the scheduled read-only audit of the newest
      Einstein--Standard-Model and Navier--Stokes notebooks in
      \texttt{latex\_notes}.  Transfer only typed analytic mechanisms
      whose hypotheses are proved in the Yang--Mills setting.
\item Prove or disprove the direct three-input tail
      \eqref{eq:V44-direct-triple-tail-gate}.  Work in a complete
      two-stage BZ/honeycomb parametrization and sum final labels,
      intermediate labels, and multiplicities jointly.
\item Keep the warm-intermediate family
      \Cref{prop:V44-warm-intermediate-regression} as a compulsory
      regression.  Do not condition the proof on a cold first
      intermediate.
\item If the direct tail holds, combine its polynomial branch with
      \eqref{eq:V44-triple-warm-probability} exactly as in V44 to obtain
      a cardinality-free triple-bank tail and close
      \(H_3^2\gtrsim B_\tau\).
\item In parallel within the notebook, resolve the local-junction branch
      \(L_e^2\gtrsim B_\tau\) and the output-dominant branch
      \(\Xi_T\gg M\) using the final \(T_e\)-block of the V41 contraction
      tree.  Charge the final resolvent once.
\item Do not start multiple-junction interpolation or another quartic
      skeleton until the unique-junction row is closed or a genuine
      counterfamily is recorded.
\item Preserve only the newest YM file.  The subsequent companion audit
      is V50.
\end{enumerate}
\end{openproblem}

\begin{firewall}[End-of-V44 claim boundary]
V44 proves a cardinality-free squared cold tail for an arbitrarily
entangled pair-shared bank, closes the entire pair-bank-dominated
spectator branch of the unique-junction \([4]\) tree, and localizes the
remaining spectator-dominated branching mass to a triple bank or the
local junction.  It does not prove the direct triple tail, either
remaining unique-junction branch, the complete \(m=4\) row, the other
quartic topologies, all-order MTA, WUFC, CPM, the nonlinear Perron
packet, finite-edge margins, capacity, protected-sector floors,
local-algebra noncollapse, reflection positivity, a nontrivial
continuum OS state, or a positive Yang--Mills spectral gap.  The full
Yang--Mills mass gap is not proved.
\end{firewall}

\section*{Version V44 append-only record}
\addcontentsline{toc}{section}{Version V44 append-only record}
The complete V43 source was retained before this continuation.  Its
SHA--256 digest was
\begin{center}\scriptsize\ttfamily
f8b45a2f252ef597481a59c2da362ca8a26368503bf39c69e6e523ccdae19fcc.
\end{center}
No mathematical content of V43 was altered.  On the next iteration,
retain this full file, append before its unique active document-closing
command, increment the filename to \texttt{\_V45.tex}, and remove only
the superseded V44 file from this Yang--Mills note series.  V45 must
begin with the scheduled SM/NS companion audit.

% =============================================================================
% END APPENDED VERSION V44 MATERIAL
% =============================================================================

% =============================================================================
% BEGIN APPENDED VERSION V45 MATERIAL
% =============================================================================

\clearpage
\part{V45 --- Companion audit, adaptive dimension-law tails, and a
Fourier-normalization obstruction}
\label{part:V45}

\section{Scheduled SM/NS companion audit}
\label{sec:V45-companion-audit}

The audit required at every fifth Yang--Mills version was performed
read-only before selecting the V45 proof direction and refreshed when
the parallel SM programme advanced during verification.  The final
snapshots used by this note were
\begin{center}
\begin{tabular}{p{0.53\textwidth}p{0.39\textwidth}}
\texttt{Renewal\_NS\_Stability\_Research\_Notes\_V31.tex}
&\scriptsize\ttfamily
f1121b62238ad5491bb7cf03635ea9934\newline
942edf573ed8faaa9ee79e1d38a6696\\[0.5em]
\texttt{Renewal\_SM\_Einstein\_Continuum\_Research\_Notes\_V64.tex}
&\scriptsize\ttfamily
2ed0dcd226c75cfe7ad1008ca692097a\newline
4b8265dac557247036c4e2434c59326e
\end{tabular}
\end{center}

The NS file is unchanged from the V40 audit.  Its exact dyadic
first-moment restoration still supports the discipline of retaining the
formation weight before summation, but it supplies no compact-group
dimension law, Casimir annulus, or Wilson multiplier.

The SM V61--V64 appendices construct split glancing quotients, compare
fixed and moving trace planes, and separate a maximal action--Green
relation from an actual complementing theorem.  V64 replaces a
singular representative by a bounded holomorphic half-angle frame,
computes the exact residual stable kernel of the old induced graph, and
then proves a uniform axial quotient lift.  Its transferable lesson is
structural: a fixed representative may lose rank even when the quotient
does not, and a new splitting must preserve the exact fibre while its
normalization and endpoint ranks are recomputed.  It does not contain an
\(\mathrm{SU}(3)\) Berenstein--Zelevinsky count, a Fourier-algebra
coefficient tail, or a three-input formation theorem.

For V45 the type-correct analogue is to change the tensor-product
bracketing.  Rather than keep the V44 order in which a cold final output
may follow a warm first intermediate, we first fuse the two largest
input heights.  Frobenius reciprocity then confines that intermediate to
an annulus around the remaining input.  This is a Yang--Mills proof
using the already proved V38--V39 representation estimates; no
Einstein or fluid inequality is imported.

\begin{firewall}[Companion transfer boundary]
\label{fire:V45-companion-types}
The audit changes a choice of representation-theoretic splitting only.
The SM moving projector and axial selector are not Clebsch--Gordan
maps, their stable singular-value floor is not a Wilson formation
estimate, and the NS coarea weight is not a Casimir dimension law.
Every quantitative estimate below is proved inside the Yang--Mills
notebook from V38--V44 hypotheses.
\end{firewall}

\section{Adaptive proof of the direct triple cold tail}
\label{sec:V45-direct-triple-tail}

For an irreducible \(R=(p,q)\), retain the V38 height
\[
 h_R:=1+p+q,\qquad A_R:=1+C_2(R).
\]
Thus \(A_R\asymp h_R^2\).  A scale-free form of the V39 pair theorem is
\begin{equation}
 \sum_{\substack{S,\mu\\A_S\le y}}
 \frac{d_S}{d_Rd_Q}
 \le
 C\min\left\{1,\frac{y^2}{(A_R+A_Q)^2}\right\},
 \qquad y\ge1.
 \label{eq:V45-scale-free-pair-tail}
\end{equation}
It follows from
\eqref{eq:V39-universal-one-link-tail} by taking \(t=y^{-1}\) and a
fixed output constant.  The trivial representation, for which
\(A_{\boldsymbol1}=1\), is included.

\begin{theorem}[Direct three-input squared cold tail]
\label{thm:V45-direct-triple-tail}
Let \(R_1,R_2,R_3\) be nontrivial irreducible
\(\mathrm{SU}(3)\) representations.  Resolve every multiplicity copy in
\(R_1\otimes R_2\otimes R_3\), and put
\[
 H^{(3)}:=A_{R_1}+A_{R_2}+A_{R_3}.
\]
Then, for every \(x\ge1\),
\begin{equation}
 \boxed{
 \sum_{\substack{S,\mu\\A_S\le x}}
 \frac{d_S}{d_{R_1}d_{R_2}d_{R_3}}
 \le
 C\min\left\{1,\frac{x^2}{(H^{(3)})^2}\right\}.}
 \label{eq:V45-direct-triple-tail}
\end{equation}
The same estimate holds for the output mass
\(Z^{(3)}\) of \eqref{eq:V44-triple-dimension-law}, because for
\(x\ge1\) the event \(Z^{(3)}\le x\) is the same resolved set.
\end{theorem}

\begin{proof}
The multiplicity and dimension law are permutation invariant.  Relabel
the inputs so that
\[
 h_1\ge h_2\ge h_3.
\]
Fuse the two largest inputs first:
\[
 R_1\otimes R_2\longrightarrow U,\qquad
 U\otimes R_3\longrightarrow S.
\]
On every resolved path, the dimension weights telescope as in
\eqref{eq:V44-triple-path-telescoping}.  The V38 reciprocal-height
inequality applied to the second fusion gives
\begin{equation}
 |h_U-h_3|\le h_S.
 \label{eq:V45-intermediate-annulus}
\end{equation}
If \(A_S\le x\), V38 height--Casimir comparability gives
\(h_S\le C_1\sqrt x\).

First suppose \(h_3\le2C_1\sqrt x\).  Then
\eqref{eq:V45-intermediate-annulus} implies
\[
 h_U\le3C_1\sqrt x,\qquad A_U\le C_2x.
\]
After summing the second fusion with total mass at most one,
\eqref{eq:V45-scale-free-pair-tail} on the first fusion gives
\begin{equation}
 \nu^{(3)}\{A_S\le x\}
 \le
 C\min\left\{1,\frac{x^2}{(A_{R_1}+A_{R_2})^2}\right\}.
 \label{eq:V45-triple-small-third}
\end{equation}
Because \(A_{R_3}\le Cx\), the right side is bounded by the right side
of \eqref{eq:V45-direct-triple-tail}: if
\(A_{R_1}+A_{R_2}\le x\), both bounds are constant; otherwise
\(H^{(3)}\asymp A_{R_1}+A_{R_2}\).

Now suppose \(h_3>2C_1\sqrt x\).  Every contributing intermediate lies
in the annulus
\begin{equation}
 \frac12h_3\le h_U\le\frac32h_3.
 \label{eq:V45-high-third-annulus}
\end{equation}
For each such \(U\), apply
\eqref{eq:V45-scale-free-pair-tail} to \(U\otimes R_3\):
\begin{equation}
 \nu_{U,R_3}\{A_S\le x\}
 \le C\frac{x^2}{h_3^4}.
 \label{eq:V45-second-fusion-tail}
\end{equation}
The upper half of \eqref{eq:V45-high-third-annulus} implies
\(A_U\le C h_3^2\).  Hence the total first-fusion mass of all possible
intermediates is at most
\begin{equation}
 C\min\left\{1,
 \frac{h_3^4}{(A_{R_1}+A_{R_2})^2}\right\}
 \label{eq:V45-first-fusion-annulus-mass}
\end{equation}
by another application of
\eqref{eq:V45-scale-free-pair-tail}.  Multiplying
\eqref{eq:V45-second-fusion-tail} and
\eqref{eq:V45-first-fusion-annulus-mass} gives
\[
 C\frac{x^2}
 {\max\{h_3^4,(A_{R_1}+A_{R_2})^2\}}
 \le
 C\frac{x^2}{(H^{(3)})^2}.
\]
The final comparison uses \(A_{R_3}\asymp h_3^2\).
Together with the trivial probability bound, this proves
\eqref{eq:V45-direct-triple-tail}.
\end{proof}

\begin{remark}[Resolution of the V44 logarithm]
\label{rem:V45-log-resolution}
The fixed bracketing in
\Cref{ass:V44-triple-acquisition} allowed many intermediate shells and
therefore produced a logarithm.  In the adaptive bracketing,
\eqref{eq:V45-intermediate-annulus} confines all contributing shells to
one constant-width height annulus.  The warm-intermediate family
\Cref{prop:V44-warm-intermediate-regression} is not discarded: after
permuting its inputs, the two \(n\)-scale representations are fused
first, and their intermediate is constrained to the fixed fundamental
scale.  Its \(H^{-2}\) mass still saturates the result.
\end{remark}

\section{Provisional tensorization over an entangled triple bank
(withdrawn)}
\label{sec:V45-triple-bank}

\begin{firewall}[Superseded conditional calculation]
The conditional statements in this section record the route attempted
before its arbitrary-parent normalization was audited.  They are
withdrawn by \Cref{ass:V45-conditional-retraction} and are retained only
to expose the failed step.  None of those conditional statements is an
active theorem of V45.
\end{firewall}

At a triple-shared spectator coordinate \(f\), define
\begin{equation}
 H_f^{(3)}:=\sum_{i\in I_f}A_{i,f},
 \qquad
 Z_f^{(3)}(S_f):=
 \boldsymbol1_{\{S_f\ne\boldsymbol1\}}
 (1+C_2(S_f)).
 \label{eq:V45-triple-link-masses}
\end{equation}
The actual conditional formation kernel is pointwise bounded by the
direct three-input dimension law, as proved at the end of
\Cref{lem:V44-triple-Casimir-moments}.

\begin{lemma}[Conditional triple-link Laplace envelope]
\label{lem:V45-triple-link-Laplace}
There are group constants \(c_3,C_3>0\) such that, for
\(0<t\le1\),
\begin{equation}
 \sum_{S_f,\mu_f}
 \pi_f(S_f,\mu_f\mid\mathscr F_{\ne f})
 e^{-tZ_f^{(3)}(S_f)}
 \le
 \Phi_3(tH_f^{(3)}),
 \label{eq:V45-triple-link-Laplace}
\end{equation}
where
\begin{equation}
 \Phi_3(z):=
 \min\left\{
 e^{-c_3\min\{z,1\}},
 \frac{C_3}{(1+z)^2}
 \right\}.
 \label{eq:V45-triple-envelope}
\end{equation}
\end{lemma}

\begin{proof}
The warm fraction
\eqref{eq:V44-triple-warm-probability} and pointwise dimension-law
domination give the exponential branch exactly as in the first
paragraph of the proof of
\Cref{lem:V44-one-link-Laplace}.  The direct tail
\eqref{eq:V45-direct-triple-tail}, followed by the same Stieltjes
integration as in that proof, gives the polynomial branch.  Both
arguments are conditional on an arbitrary nonzero parent block; the
zero-parent convention is immediate.
\end{proof}

Let
\[
 \mathcal F_3:=\bigcup_{|J|=3}E_J,\qquad
 M_3(\boldsymbol S):=
 \sum_{\substack{f\in\mathcal F_3\\S_f\ne\boldsymbol1}}
 (1+C_2(S_f)).
\]
The input mass \(H_3=\sum_{f\in\mathcal F_3}H_f^{(3)}\) agrees with
\eqref{eq:V44-triple-local-input-masses}.

\begin{theorem}[Entangled triple-bank squared tail]
\label{thm:V45-triple-bank-tail}
For every nonempty triple bank and every conditioning on the remaining
spectator outputs,
\begin{align}
 \sum_{\boldsymbol S_{\mathcal F_3},\boldsymbol\mu_{\mathcal F_3}}
 \pi_{\mathcal F_3}e^{-tM_3(\boldsymbol S)}
 &\le\frac C{(1+tH_3)^2},
 \qquad 0<t\le1,
 \label{eq:V45-triple-bank-Laplace}\\
 \boxed{
 \pi_{\mathcal F_3}\{M_3\le x\}
 &\le C\min\left\{1,\frac{x^2}{H_3^2}\right\},
 \qquad x\ge1.}
 \label{eq:V45-triple-bank-tail}
\end{align}
The constants are independent of bank size and coefficient
entanglement.
\end{theorem}

\begin{proof}
Order the triple links and use the exact parent--child disintegration
\eqref{eq:V43-conditional-formation-kernel}.  Successive application of
\Cref{lem:V45-triple-link-Laplace} bounds the joint transform by
\(\prod_f\Phi_3(tH_f^{(3)})\).
\Cref{lem:V44-product-envelope} applies verbatim to every envelope of
the form \eqref{eq:V45-triple-envelope}, proving
\eqref{eq:V45-triple-bank-Laplace}.  Set \(t=x^{-1}\) and use
\(\boldsymbol1_{\{M_3\le x\}}\le e e^{-M_3/x}\) to obtain
\eqref{eq:V45-triple-bank-tail}.
\end{proof}

\begin{corollary}[Full spectator mass inherits the triple-bank tail]
\label{cor:V45-global-triple-tail}
If \(\mathcal F_3\ne\varnothing\), then
\begin{equation}
 \pi\{M(\boldsymbol S)\le x\}
 \le
 C\min\left\{1,\frac{x^2}{H_3^2}\right\},
 \qquad x\ge1.
 \label{eq:V45-global-triple-tail}
\end{equation}
\end{corollary}

\begin{proof}
Condition on all pair-shared outputs and use
\(\{M\le x\}\subseteq\{M_3\le x\}\), then integrate the conditioning.
\end{proof}

\begin{theorem}[Triple-bank-dominated unique-\([4]\) MTA sector]
\label{thm:V45-triple-bank-four-MTA}
Fix \(D,\delta>0\) and a tree \(\tau\).  If
\begin{equation}
 H_3^2\ge\delta B_\tau,
 \label{eq:V45-triple-bank-dominance}
\end{equation}
then
\begin{align}
 &\sum_{\boldsymbol S,\boldsymbol\mu}
 \pi_{\boldsymbol S,\boldsymbol\mu}
 \sup_{\substack{T_e:\,\Xi_T\le DM(\boldsymbol S)}}
 \frac{\Xi_Tq_{\rm sp}(\boldsymbol S)w_{\tau,e}}
 {1-q_{\alpha,\boldsymbol T}}
 \notag\\
 &\qquad\le C_{D,\delta}P_{\tau,e}.
 \label{eq:V45-triple-bank-scalar-average}
\end{align}
This is the triple-bank-dominated spectator part of the unique-junction
\(m=4\) MTA term after global compression.
\end{theorem}

\begin{proof}
Insert \eqref{eq:V45-global-triple-tail} into
\Cref{lem:V43-weighted-tail-summation}, exactly as
\eqref{eq:V44-global-pair-tail} was used in
\Cref{thm:V44-pair-bank-four-MTA}.  The result is
\[
 \frac C{H_3^2}\prod_iA_{i,e}^{d_i/2}
 \le
 \frac C{B_\tau}\prod_iA_{i,e}^{d_i}
 =CP_{\tau,e}.
\]
The final-channel sum uses the V42 global junction compression, not an
unweighted multiplicity count.
\end{proof}

\section{A four-input cold tail and blockwise junction compression}
\label{sec:V45-four-input-tail}

At the unique junction \(e\), put
\[
 d_{e,\rm in}:=\prod_{i=1}^4d_{R_{i,e}},
 \qquad
 L_e=\sum_{i=1}^4A_{i,e},
\]
and resolve all copies of
\(\bigotimes_{i=1}^4R_{i,e}\).  Define
\begin{equation}
 \nu_e^{(4)}(T,\lambda):=
 \frac{d_T}{d_{e,\rm in}},
 \qquad
 Z_e(T):=
 \boldsymbol1_{\{T\ne\boldsymbol1\}}(1+C_2(T)).
 \label{eq:V45-four-input-dimension-law}
\end{equation}

\begin{corollary}[Direct four-input squared cold tail]
\label{cor:V45-direct-four-tail}
For every \(x\ge1\),
\begin{equation}
 \boxed{
 \sum_{\substack{T,\lambda\\Z_e(T)\le x}}
 \nu_e^{(4)}(T,\lambda)
 \le
 C\min\left\{1,\frac{x^2}{L_e^2}\right\}.}
 \label{eq:V45-direct-four-tail}
\end{equation}
\end{corollary}

\begin{proof}
Relabel the four inputs by decreasing V38 height.  Fuse the largest
three first to an intermediate \(U\), using
\Cref{thm:V45-direct-triple-tail}, and then fuse \(U\) with the smallest
input to \(T\), using
\eqref{eq:V45-scale-free-pair-tail}.  If \(Z_e(T)\le x\), reciprocal
height confines \(h_U\) to a \(C\sqrt x\)-neighborhood of the fourth
height.

If that height is at most \(C'\sqrt x\), then \(A_U\le Cx\), and the
triple tail alone gives
\[
 C\min\{1,x^2/(A_{R_1}+A_{R_2}+A_{R_3})^2\}.
\]
If it is larger, \(h_U\asymp h_{R_4}\).  The final pair tail contributes
\(Cx^2/h_{R_4}^4\), while the triple tail for the allowed
intermediates contributes
\[
 C\min\{1,h_{R_4}^4/
 (A_{R_1}+A_{R_2}+A_{R_3})^2\}.
\]
The same maximum-versus-sum calculation as in the last paragraph of the
proof of \Cref{thm:V45-direct-triple-tail} gives
\eqref{eq:V45-direct-four-tail}.  Multiplicity weights telescope at
both stages.
\end{proof}

\begin{firewall}[Representation tail versus the following conditional
compression]
Corollary~\ref{cor:V45-direct-four-tail} is an active
representation-dimension statement.  The conditional coefficient
compression beginning below and continuing through
\Cref{ass:V45-local-output-gate} is provisional and is withdrawn by
\Cref{ass:V45-conditional-retraction}.  The distinction is essential:
the former averages dimensions, whereas the latter purports to control
an arbitrary trace-class parent.
\end{firewall}

Fix all spectator blocks and let \(\mathcal B_{\rm par}\) be the
remaining trace-class coefficient on the four local input carriers.
Pinch it to a resolved local output block
\(\mathcal B_{T,\lambda}\), and define
\begin{equation}
 \pi_e(T,\lambda\mid\mathscr F_{\ne e})
 :=
 \frac{d_T\|\mathcal B_{T,\lambda}\|_1}
 {d_{e,\rm in}\|\mathcal B_{\rm par}\|_1}
 \label{eq:V45-local-formation-kernel}
\end{equation}
for a nonzero parent, with the usual zero-parent convention.

\begin{lemma}[Blockwise refinement of V41 junction compression]
\label{lem:V45-blockwise-junction-compression}
The weights \eqref{eq:V45-local-formation-kernel} are subprobability and
obey
\begin{equation}
 \pi_e(T,\lambda\mid\mathscr F_{\ne e})
 \le\nu_e^{(4)}(T,\lambda).
 \label{eq:V45-local-dimension-domination}
\end{equation}
Moreover, for every set \(\Omega\) of local outputs, the
\(\tau\)-tree coefficient satisfies
\begin{equation}
 \sum_{(T,\lambda)\in\Omega}
 d_T\|\widehat{\mathbf K}_{\tau;T,\lambda}\|_1
 \le
 C_4d_{e,\rm in}w_{\tau,e}
 \|\mathcal B_{\rm par}\|_1
 \sum_{(T,\lambda)\in\Omega}
 \pi_e(T,\lambda\mid\mathscr F_{\ne e}).
 \label{eq:V45-blockwise-junction-compression}
\end{equation}
\end{lemma}

\begin{proof}
Orthogonal compression gives
\(\|\mathcal B_{T,\lambda}\|_1\le
\|\mathcal B_{\rm par}\|_1\), proving
\eqref{eq:V45-local-dimension-domination}; summing all compressed block
norms proves subprobability.

The V41 tree operator commutes with the diagonal action.  On every
resolved block its operator norm is at most
\(C_4w_{\tau,e}\) by
\eqref{eq:V41-exact-Wilson-tree-bound}.  Multiplication of a
trace-class block by that operator therefore costs at most
\(C_4w_{\tau,e}\).  Insert
\eqref{eq:V45-local-formation-kernel} and sum \(\Omega\).  This is the
proof of \Cref{thm:V41-fourfold-compression} stopped before discarding
the output-set mass.
\end{proof}

\section{Closure of every spectator-dominated unique junction}
\label{sec:V45-spectator-closure}

\begin{lemma}[Local cold-formation payment]
\label{lem:V45-local-cold-payment}
Fix \(D\ge1\).  For every resolved spectator output with \(M>0\),
\begin{equation}
 \pi_e\{\Xi_T\le DM\mid\mathscr F_{\ne e}\}
 \le
 C_D\min\left\{1,\frac{M^2}{L_e^2}\right\}.
 \label{eq:V45-local-admissible-mass}
\end{equation}
Furthermore,
\begin{align}
 &\min\left\{1,\frac{M^2}{L_e^2}\right\}
 \left[
 s_\alpha^2\boldsymbol1_{\{s_\alpha M\le c_D\}}
 +s_\alpha^3Mq_{\rm sp}
  \boldsymbol1_{\{s_\alpha M>c_D\}}
 \right]
 \le\frac{C_D}{L_e^2}.
 \label{eq:V45-local-response-payment}
\end{align}
If \(M=0\), the admissible nonconstant output set is empty.
\end{lemma}

\begin{proof}
The full-output split \eqref{eq:V43-full-output-split} shows that
\(\Xi_T\le DM\) implies \(Z_e(T)\le(D-1)M\).
When \(M>0\), exact-support quantization gives \(M\ge1\) up to a fixed
group constant.  Apply
\eqref{eq:V45-local-dimension-domination} and
\Cref{cor:V45-direct-four-tail}, enlarging the output threshold by a
fixed \(D\), to obtain
\eqref{eq:V45-local-admissible-mass}.

For the low term, if \(M\le L_e\), then
\[
 s_\alpha^2\frac{M^2}{L_e^2}
 \le\frac{c_D^2}{L_e^2};
\]
if \(M>L_e\), then \(s_\alpha^2\le c_D^2/L_e^2\).
For the high term, if \(M\le L_e\), use the third product moment
\eqref{eq:V42-balanced-third-moment}:
\[
 \frac{M^2}{L_e^2}s_\alpha^3Mq_{\rm sp}
 =\frac{(s_\alpha M)^3q_{\rm sp}}{L_e^2}
 \le\frac C{L_e^2}.
\]
If \(M>L_e\), the same moment gives
\[
 s_\alpha^3Mq_{\rm sp}
 =\frac{(s_\alpha M)^3q_{\rm sp}}{M^2}
 \le\frac C{L_e^2}.
\]
This proves \eqref{eq:V45-local-response-payment}.  If \(M=0\),
\(\Xi_T\le DM\) forces the constant full output, which is removed by
\(Q_0\).
\end{proof}

\begin{theorem}[Local-mass-dominated spectator branch]
\label{thm:V45-local-mass-four-MTA}
Fix \(D,\delta>0\).  If
\begin{equation}
 L_e^2\ge\delta B_\tau,
 \label{eq:V45-local-mass-dominance}
\end{equation}
then the jointly resolved spectator and local outputs satisfying
\(\Xi_T\le DM(\boldsymbol S)\) obey
\begin{align}
 &\sum_{\boldsymbol S,\boldsymbol\mu}
 \pi_{\boldsymbol S,\boldsymbol\mu}
 \sum_{T_e,\lambda}
 \pi_e(T_e,\lambda\mid\boldsymbol S)
 \boldsymbol1_{\{\Xi_T\le DM(\boldsymbol S)\}}
 \frac{\Xi_Tq_{\rm sp}(\boldsymbol S)w_{\tau,e}}
 {1-q_{\alpha,\boldsymbol T}}
 \notag\\
 &\qquad\le C_{D,\delta}P_{\tau,e}.
 \label{eq:V45-local-mass-scalar-average}
\end{align}
With \Cref{lem:V45-blockwise-junction-compression}, this is the
local-mass-dominated spectator part of the unique-junction \(m=4\) MTA
term.
\end{theorem}

\begin{proof}
For fixed spectator output, the V43 final-resolvent split bounds the
quotient in \eqref{eq:V45-local-mass-scalar-average} by
\[
 C_D\prod_iA_{i,e}^{d_i/2}
 \left[
 s_\alpha^2\boldsymbol1_{\{s_\alpha M\le c_D\}}
 +s_\alpha^3Mq_{\rm sp}
  \boldsymbol1_{\{s_\alpha M>c_D\}}
 \right].
\]
Sum the admissible local outputs using
\eqref{eq:V45-local-admissible-mass}, then apply
\eqref{eq:V45-local-response-payment}.  The remaining spectator kernel
has mass at most one, so the result is
\[
 \frac C{L_e^2}\prod_iA_{i,e}^{d_i/2}
 \le
 \frac C{B_\tau}\prod_iA_{i,e}^{d_i}
 =CP_{\tau,e}.
\]
The first inequality uses
\eqref{eq:V45-local-mass-dominance} and \(A_{i,e}\ge1\).
The blockwise compression lemma supplies the actual local tree
coefficient and sums its multiplicities with the displayed formation
weights.
\end{proof}

\begin{theorem}[Complete spectator-dominated unique-\([4]\) sector]
\label{thm:V45-complete-spectator-dominated-four}
For every \(D\ge1\), every unique four-row junction, and every one of
the sixteen V41 trees, the complete resolved sub-sum satisfying
\begin{equation}
 \Xi_T\le DM(\boldsymbol S)
 \label{eq:V45-complete-spectator-condition}
\end{equation}
obeys the required \(m=4\) marked-tree atlas term.
\end{theorem}

\begin{proof}
By \Cref{prop:V44-exact-residual}, V43--V44 already close the sub-sum
unless \(H_3^2\ge B_\tau/16\) or
\(L_e^2\ge B_\tau/16\).  The first alternative is
\Cref{thm:V45-triple-bank-four-MTA}; the second is
\Cref{thm:V45-local-mass-four-MTA}.  These cases exhaust the
spectator-dominated branch.
\end{proof}

\section{The sole remaining unique-junction branch}
\label{sec:V45-output-dominant-residual}

\begin{proposition}[Output dominance is a local four-input phenomenon]
\label{prop:V45-output-localization}
If
\begin{equation}
 \Xi_T>2M(\boldsymbol S),
 \label{eq:V45-output-dominance}
\end{equation}
then
\begin{equation}
 Z_e(T_e)>\frac12\Xi_T,
 \qquad
 \Xi_T\le C L_e.
 \label{eq:V45-output-local-mass}
\end{equation}
Thus the residual final mass cannot come from a pair, triple, or private
spectator output.
\end{proposition}

\begin{proof}
The first assertion is immediate from
\(\Xi_T=M+Z_e(T_e)\).  Iterating the additive fusion-Casimir bound
\eqref{eq:V30-Casimir-fusion} over the four local inputs gives
\[
 Z_e(T_e)\le C\sum_{i=1}^4A_{i,e}=CL_e.
\]
Combine this with \(Z_e>\Xi_T/2\).
\end{proof}

\begin{assessment}[Exact post-V45 local gate]
\label{ass:V45-local-output-gate}
All spectator-dominated unique-junction outputs are closed.  The only
remaining branch has \(\Xi_T\asymp Z_e(T_e)\lesssim L_e\).  On it, taking
the V41 estimate
\(\|\mathbf K_{\tau;T_e}\|_{\rm op}\le Cw_{\tau,e}\) before coupling to
the final channel loses the high-output information: the resolvent is
already in its \(\Xi_T\) regime, and no spectator multiplier need be
small.

The next estimate must therefore retain the exact final block of the
one-link tree operator.  In normalized form the unresolved quantity is
\begin{equation}
 \mathfrak L_{\tau,e}
 :=
 \sum_{\boldsymbol S,\boldsymbol\mu,T_e,\lambda}
 \Pi_{\boldsymbol S,\boldsymbol\mu;T_e,\lambda}
 \boldsymbol1_{\{\Xi_T>2M\}}
 \frac{\Xi_Tq_{\rm sp}}
 {1-q_{\rm sp}\eta_{T_e}}\,
 \|\mathbf K_{\tau;T_e,\lambda}\|_{\rm op},
 \label{eq:V45-local-output-functional}
\end{equation}
where \(\Pi\) is the actual joint spectator/local pinching law and the
local operator is not replaced by its supremum.  The required bound is
\begin{equation}
 \mathfrak L_{\tau,e}
 \le
 C\,\frac{\prod_iA_{i,e}^{d_i}}{B_\tau}.
 \label{eq:V45-local-output-gate}
\end{equation}
This is a single-coordinate operator-valued high-output moment.  It is
strictly narrower than the former spectator-bank problem and must be
tested on the V40 coherent high-weight family before any interpolation
to multiple junctions.
\end{assessment}

\section{Provisional formation-first repair (withdrawn)}
\label{sec:V45-formation-first-repair}

\begin{firewall}[The first repair is also superseded]
The arbitrary-parent counterexample below is valid, but the
formation-first theorems following
\Cref{ass:V45-conditional-retraction} still confuse representation
dimension weight with genuine Fourier coefficient mass.  They are
withdrawn by the exact dual-singlet calculation in
\Cref{cth:V45-dual-singlet-normalization,ass:V45-final-correction} and
are retained as a documented failed repair, not as active results.
\end{firewall}

\begin{counterexample}[Arbitrary-parent dimension-law domination fails]
\label{cex:V45-conditional-dimension-failure}
Let
\[
 H_R\otimes H_Q=\bigoplus_{S,\mu}H_S
\]
be a resolved pair decomposition, and let \(P_{S,\mu}\) be one nonzero
block projection.  There is a trace-class parent matrix \(X\), supported
inside that block, for which
\begin{equation}
 \frac{\|P_{S,\mu}XP_{S,\mu}\|_1}{\|X\|_1}=1.
 \label{eq:V45-concentrated-parent}
\end{equation}
Whenever \(d_S<d_Rd_Q\), this exceeds the dimension-law weight
\(d_S/(d_Rd_Q)\).  For example, in
\(3\otimes3=6\oplus\bar3\), a parent supported on the \(\bar3\) block
has raw conditional mass one while its dimension-law mass is \(1/3\).
\end{counterexample}

\begin{proof}
Take any nonzero trace-class \(X=P_{S,\mu}XP_{S,\mu}\).  Then
\eqref{eq:V45-concentrated-parent} is an identity.  The displayed
\(\mathrm{SU}(3)\) dimensions give the numerical example.
\end{proof}

\begin{assessment}[Retraction of the conditional-tail route]
\label{ass:V45-conditional-retraction}
The conditional kernel
\eqref{eq:V43-conditional-formation-kernel}, with its inserted dimension
ratio, is a subprobability kernel.  Counterexample
\ref{cex:V45-conditional-dimension-failure} shows that it is not the raw
pinching mass which occurs in the V37/V41 partial-trace compression.
Consequently an arbitrary-parent dimension-law tail cannot simply be
inserted after all other outputs have been conditioned.

The proofs of
\Cref{lem:V43-scale-free-pair-tail,
thm:V43-pair-heavy-four-MTA,
thm:V44-pair-bank-tail,
cor:V44-global-pair-tail,
thm:V44-pair-bank-four-MTA,
lem:V45-triple-link-Laplace,
thm:V45-triple-bank-tail,
cor:V45-global-triple-tail,
thm:V45-triple-bank-four-MTA,
lem:V45-blockwise-junction-compression,
lem:V45-local-cold-payment,
thm:V45-local-mass-four-MTA,
thm:V45-complete-spectator-dominated-four}
are therefore withdrawn in their conditional form.  Their conclusions
are not used below unless explicitly re-established by the
formation-first estimates in this section.  The pure dimension-law
theorems, the V44 Casimir moments and product-envelope lemma, and the
deterministic mass localization remain valid.
\end{assessment}

The repair uses the warning already proved in
\Cref{thm:V35-amplified-Schur-tax,cor:V35-fresh-leaf-taxes}: a Schur tax
may be applied with arbitrary spectators, but repeated taxes are
multiplied only when each fusion introduces a previously unused input
row.  We therefore resolve one \emph{whole incidence bank} before any
other shared bank is pinched.

Fix \(J\subset\{1,2,3,4\}\), \(2\le|J|\le4\), and a coordinate bank
\(F\) on which exactly the rows in \(J\) are active.  Put
\[
 \boldsymbol R_i^F:=(R_{i,f})_{f\in F},\qquad
 d_{\boldsymbol R_i^F}:=\prod_{f\in F}d_{R_{i,f}}.
\]
Resolve a fixed fresh-leaf fusion tree on the product group \(G^F\).
For a complete resolved path
\((\boldsymbol S,\boldsymbol\mu)\), let
\begin{equation}
 \nu_{J,F}(\boldsymbol S,\boldsymbol\mu)
 :=
 \frac{d_{\boldsymbol S}}
 {\prod_{i\in J}d_{\boldsymbol R_i^F}}.
 \label{eq:V45-product-dimension-law}
\end{equation}
Intermediate dimensions telescope, so the sum over all resolved paths
is one and the law factors coordinatewise.

\begin{theorem}[Formation-first domination on one incidence bank]
\label{thm:V45-formation-first-domination}
Apply the bank fusion before resolving any carrier outside \(F\), and
treat every such carrier as a spectator in the sense of V35.  In the
same normalized Fourier coefficient convention as
\eqref{eq:V35-formation-kernel}, the actual coefficient kernel
\(\Pi_{J,F}^{\rm ff}\) satisfies
\begin{equation}
 \boxed{
 0\le
 \Pi_{J,F}^{\rm ff}(\boldsymbol S,\boldsymbol\mu)
 \le
 \nu_{J,F}(\boldsymbol S,\boldsymbol\mu)
 \prod_{a=1}^{|J|-1}\tau_a
 \le
 \nu_{J,F}(\boldsymbol S,\boldsymbol\mu).}
 \label{eq:V45-formation-first-domination}
\end{equation}
The estimate is uniform in all spectator dimensions and entanglement
inside each input row.  All subsequent spectator pinches, partial
traces, and contraction operators of norm at most one preserve this
majorant.
\end{theorem}

\begin{proof}
Fuse the first two row coefficients on the entire product bank \(F\).
\Cref{thm:V35-amplified-Schur-tax} applies with every carrier outside
\(F\) in the spectator spaces and supplies \(\tau_1\).  Fuse the result
with a fresh third row, and, if present, a fresh fourth row.
\Cref{cor:V35-fresh-leaf-taxes} gives the product of taxes.  Multiplying
by the exact output/input dimension ratio in the normalized formation
kernel gives \eqref{eq:V45-product-dimension-law}; every intermediate
dimension cancels.  Each \(\tau_a\le1\).  Later operations are
trace-norm contractions or bounded multiplicity operators, as in
\Cref{lem:V37-partial-trace}, so they cannot enlarge the majorant.
\end{proof}

\subsection{Repaired pair-incidence bank estimate}
\label{subsec:V45-repaired-pair}

For \(|J|=2\), define
\[
 H_{2,J}:=\sum_{f\in E_J}\sum_{i\in J}A_{i,f},
 \qquad
 M_{2,J}:=
 \sum_{\substack{f\in E_J\\S_f\ne\boldsymbol1}}
 (1+C_2(S_f)).
\]

\begin{theorem}[Formation-first pair-bank tail]
\label{thm:V45-ff-pair-bank-tail}
For every nonempty \(E_J\), \(|J|=2\), the actual formation-first
kernel obeys
\begin{align}
 \sum\Pi_{J,E_J}^{\rm ff}e^{-tM_{2,J}}
 &\le\frac C{(1+tH_{2,J})^2},
 \qquad0<t\le1,
 \label{eq:V45-ff-pair-Laplace}\\
 \Pi_{J,E_J}^{\rm ff}\{M_{2,J}\le x\}
 &\le C\min\left\{1,\frac{x^2}{H_{2,J}^2}\right\},
 \qquad x\ge1.
 \label{eq:V45-ff-pair-tail}
\end{align}
\end{theorem}

\begin{proof}
By \Cref{thm:V45-formation-first-domination}, it suffices to integrate
against the product dimension law.  Under that law the link outputs are
independent.  The normalized-trace warm estimate
\eqref{eq:V44-pair-warm-probability} and V39 cold tail give the
one-link envelope \eqref{eq:V44-Laplace-envelope}.  Multiply those
envelopes and apply
\Cref{lem:V44-product-envelope}.  Chernoff's inequality with \(t=x^{-1}\)
gives the tail.
\end{proof}

\begin{theorem}[Repaired pair-bank spectator closure]
\label{thm:V45-ff-pair-closure}
If \(H_2^2\ge\delta B_\tau\), the pair-bank-dominated
spectator-output sub-sum satisfies the required unique-junction
\(m=4\) tree term.
\end{theorem}

\begin{proof}
There are six pair incidence sets.  Choose \(J_*\) with
\[
 H_{2,J_*}\ge H_2/6.
\]
Resolve \(E_{J_*}\) first.  On the event \(M\le x\), necessarily
\(M_{2,J_*}\le x\), so
\eqref{eq:V45-ff-pair-tail} is a squared lower-tail majorant for the
complete later output sum.  All later spectator and junction
compressions are trace-norm contractions.  Insert this majorant into
\Cref{lem:V43-weighted-tail-summation}; the V43 scalar response then
gives
\[
 \frac C{H_{2,J_*}^2}\prod_iA_{i,e}^{d_i/2}
 \le
 \frac C{B_\tau}\prod_iA_{i,e}^{d_i}
 =CP_{\tau,e}.
\]
No two linkwise Schur taxes were multiplied after their input rows had
already been joined.
\end{proof}

\subsection{Repaired triple-incidence bank estimate}
\label{subsec:V45-repaired-triple}

For \(|J|=3\), define
\[
 H_{3,J}:=\sum_{f\in E_J}\sum_{i\in J}A_{i,f},
 \qquad
 M_{3,J}:=
 \sum_{\substack{f\in E_J\\S_f\ne\boldsymbol1}}
 (1+C_2(S_f)).
\]

\begin{theorem}[Formation-first triple-bank tail and closure]
\label{thm:V45-ff-triple-bank}
For every nonempty \(E_J\), \(|J|=3\),
\begin{align}
 \sum\Pi_{J,E_J}^{\rm ff}e^{-tM_{3,J}}
 &\le\frac C{(1+tH_{3,J})^2},
 \label{eq:V45-ff-triple-Laplace}\\
 \Pi_{J,E_J}^{\rm ff}\{M_{3,J}\le x\}
 &\le C\min\left\{1,\frac{x^2}{H_{3,J}^2}\right\}.
 \label{eq:V45-ff-triple-tail}
\end{align}
If \(H_3^2\ge\delta B_\tau\), the triple-bank-dominated
spectator-output sub-sum satisfies the required unique-junction tree
term.
\end{theorem}

\begin{proof}
Fuse two whole-bank row representations and then the fresh third row.
\Cref{thm:V45-formation-first-domination} majorizes the actual
coefficient kernel by the coordinatewise product of the direct
three-input dimension laws.  The warm estimate
\eqref{eq:V44-triple-warm-probability} and the now-proved direct cold
tail \eqref{eq:V45-direct-triple-tail} give an envelope of the form
\eqref{eq:V45-triple-envelope}.  Product-envelope compression and
Chernoff give the two displayed estimates.

There are four triple incidence sets.  Choose \(J_*\) with
\(H_{3,J_*}\ge H_3/4\), resolve that bank first, and repeat the proof of
\Cref{thm:V45-ff-pair-closure} with \(H_{3,J_*}\).  This yields
\[
 C H_{3,J_*}^{-2}\prod_iA_{i,e}^{d_i/2}
 \le CP_{\tau,e}.
\]
\end{proof}

\subsection{Formation-first local cold shells}
\label{subsec:V45-repaired-local}

At the one-coordinate four-row junction, choose a fresh-leaf local
fusion tree while every non-\(e\) carrier remains a spectator.
\Cref{thm:V45-formation-first-domination} and
\Cref{cor:V45-direct-four-tail} imply that restricting the local output
to \(Z_e(T_e)\le x\) costs at most
\begin{equation}
 C\min\left\{1,\frac{x^2}{L_e^2}\right\}
 \label{eq:V45-ff-local-cold-cost}
\end{equation}
in the actual normalized coefficient sum.  The V41 local tree operator
may then be inserted at cost \(Cw_{\tau,e}\); it commutes with the final
isotypic projection, and all later spectator operations are
trace-norm contractions.

\begin{lemma}[Dyadic local-cold response sum]
\label{lem:V45-dyadic-local-response}
Let \(s=s_\alpha\), \(L=L_e\), and \(m_k=2^k\), \(k\ge0\).  Uniformly in
the exact spectator product channel,
\begin{align}
 &\sum_{k\ge0}
 \min\{1,m_k^2/L^2\}
 \left[
 s^2\boldsymbol1_{\{sm_k\le C\}}
 +s^3m_kq_{\rm sp}\boldsymbol1_{\{sm_k>C\}}
 \right]
 \le\frac{C'}{L^2}.
 \label{eq:V45-dyadic-local-response}
\end{align}
On a shell \(m_k\le M<2m_k\), constants may be changed uniformly.
\end{lemma}

\begin{proof}
For the low shells with \(m_k\le L\), the sum is geometric and bounded
by
\[
 \frac{s^2}{L^2}\min\{L,s^{-1}\}^2\le L^{-2}.
\]
For low shells with \(m_k>L\), necessarily \(sL< C\), and their
contribution is at most
\[
 Cs^2\log\left(2+\frac1{sL}\right)\le\frac C{L^2}.
\]

On high shells, the third and fourth product moments give respectively
\[
 s^3m_kq_{\rm sp}\le\frac C{m_k^2},
 \qquad
 s^3m_kq_{\rm sp}\le\frac C{s\,m_k^3}.
\]
If \(m_k\le L\), use the fourth bound after multiplication by
\(m_k^2/L^2\); the resulting \(C/(s m_kL^2)\) series starts at
\(m_k\asymp s^{-1}\) and is geometric.  If \(m_k>L\), use the fourth
bound directly.  Its first scale is
\(\max\{L,s^{-1}\}\), giving at most \(C/L^2\), and the remaining terms
are geometric.
\end{proof}

\begin{theorem}[Repaired local-mass spectator closure]
\label{thm:V45-ff-local-closure}
If \(L_e^2\ge\delta B_\tau\), the locally cold,
spectator-dominated sub-sum \(\Xi_T\le DM\) satisfies the required
unique-junction \(m=4\) tree term.
\end{theorem}

\begin{proof}
Partition nonzero spectator masses into dyadic shells
\(m_k\le M<2m_k\).  On an admissible shell,
\(Z_e(T_e)\le(D-1)M\le C_Dm_k\).  Resolve the local fourfold fusion
first and apply \eqref{eq:V45-ff-local-cold-cost}; then insert the V41
tree operator and perform every spectator contraction.  The V43
resolvent split leaves the bracket in
\eqref{eq:V45-dyadic-local-response}.  Summing the shells gives
\[
 \frac C{L_e^2}\prod_iA_{i,e}^{d_i/2}
 \le
 \frac C{B_\tau}\prod_iA_{i,e}^{d_i}
 =CP_{\tau,e}.
\]
The \(M=0\) admissible output is constant and is removed by \(Q_0\).
\end{proof}

\begin{theorem}[Repaired complete spectator-dominated \([4]\) sector]
\label{thm:V45-ff-complete-spectator}
For every fixed \(D\), the full unique-junction sub-sum with
\(\Xi_T\le DM\) obeys all sixteen required \(m=4\) tree terms.
\end{theorem}

\begin{proof}
If \(B_\tau\le16M^2\), use the pointwise V43
spectator-output payment, which requires no conditional formation tail.
Otherwise, if \(H_2^2\ge B_\tau/16\), use
\Cref{thm:V45-ff-pair-closure}.  In the complementary case,
\Cref{lem:V44-two-source-localization} gives
\[
 H_3^2\ge B_\tau/16
 \quad\hbox{or}\quad
 L_e^2\ge B_\tau/16.
\]
Apply respectively
\Cref{thm:V45-ff-triple-bank,thm:V45-ff-local-closure}.
These alternatives exhaust the sub-sum.
\end{proof}

\section{Decisive Fourier-normalization counterexample}
\label{sec:V45-Fourier-normalization}

The formation-first repair above still identifies the product dimension
law with the normalized Fourier-algebra output mass.  That identification
is false.  The following one-link calculation decides the issue without
an asymptotic approximation.

\begin{counterexample}[Dual-singlet coefficient mass has only one
dimension power]
\label{cth:V45-dual-singlet-normalization}
Let \(R\) be a nontrivial irreducible unitary representation of
dimension \(d\), choose a unit vector \(e_1\in H_R\), and set
\[
 A=|e_1\rangle\langle e_1|,
 \qquad
 B=|e_1^*\rangle\langle e_1^*|.
\]
Define the pure Fourier blocks
\begin{equation}
 f_A(g):=d\,\Tr(A D^R(g)),
 \qquad
 g_B(g):=d\,\Tr(B D^{R^*}(g)).
 \label{eq:V45-dual-pure-blocks}
\end{equation}
Their input Fourier-algebra masses are \(d\) and \(d\).  The singlet
Fourier coefficient of the product is
\begin{equation}
 \widehat{f_Ag_B}(\boldsymbol1)
 =\int_G f_A(g)g_B(g)\,\dd g=d,
 \label{eq:V45-dual-singlet-coefficient}
\end{equation}
so its normalized output/input mass is
\begin{equation}
 \boxed{
 \frac{d_{\boldsymbol1}
 \left|\widehat{f_Ag_B}(\boldsymbol1)\right|}
 {(d\|A\|_1)(d\|B\|_1)}
 =\frac1d.}
 \label{eq:V45-actual-singlet-mass}
\end{equation}
The tensor-product dimension law assigns the same singlet the mass
\begin{equation}
 \nu_{R,R^*}(\boldsymbol1)=\frac1{d^2}.
 \label{eq:V45-singlet-dimension-law}
\end{equation}
Hence no uniform pointwise domination of actual Fourier-algebra
formation mass by the dimension law is possible.
\end{counterexample}

\begin{proof}
In the chosen bases,
\[
 f_A(g)=d\,D^R_{11}(g),
 \qquad
 g_B(g)=d\,\overline{D^R_{11}(g)}.
\]
Schur orthogonality gives
\[
 \int_G|D^R_{11}(g)|^2\,\dd g=\frac1d,
\]
which proves \eqref{eq:V45-dual-singlet-coefficient}.  Since
\(\|A\|_1=\|B\|_1=1\), division by the two input masses gives
\eqref{eq:V45-actual-singlet-mass}.  The singlet occurs once in
\(R\otimes R^*\), so its normalized representation-dimension weight is
\(1/(d\cdot d)\), proving
\eqref{eq:V45-singlet-dimension-law}.
\end{proof}

\begin{corollary}[The squared cold dimension tail is not an
arbitrary-coefficient tail]
\label{cor:V45-no-coefficient-dimension-tail}
Take \(R=(n,0)\) and \(R^*=(0,n)\).  Then
\[
 d_R\asymp1+C_2(R).
\]
At the fixed cold threshold \(x=1\), the actual coefficient mass in
\eqref{eq:V45-actual-singlet-mass} is of order
\((1+C_2(R))^{-1}\), whereas the V39 dimension-law tail is of order
\((1+C_2(R))^{-2}\).  Therefore an estimate of the form
\begin{equation}
 \sum_{A_S\le x}
 \frac{d_S\|\widehat{f_Ag_B}(S)\|_1}
 {(d_R\|A\|_1)(d_Q\|B\|_1)}
 \le
 C\min\left\{1,\frac{x^2}{(A_R+A_Q)^2}\right\}
 \label{eq:V45-false-coefficient-tail}
\end{equation}
is false uniformly over coefficient matrices.
\end{corollary}

\begin{proof}
The Weyl formula gives
\(d_{(n,0)}=(n+1)(n+2)/2\), while the Casimir formula gives
\(1+C_2(n,0)\asymp n^2\).  Insert
\eqref{eq:V45-actual-singlet-mass}.  The proposed right side at \(x=1\)
is \(O(n^{-4})\), while the left side is \(O(n^{-2})\).
\end{proof}

\begin{assessment}[Final V45 correction and reopened frontier]
\label{ass:V45-final-correction}
The artificial dimension-weighted kernels used in V35, V43, V44, and
the provisional formation-first section of V45 can be subprobability
and can be bounded by representation dimension laws.  Counterexample
\ref{cth:V45-dual-singlet-normalization} proves that they are not the
actual normalized Fourier-algebra output masses controlled by the
V37/V41 partial-trace formula.  If
\(\mathcal C\) denotes a raw CG compression, multiplying it by the
output/input dimension ratio creates the artificial kernel; if it
denotes the Fourier-normalized product coefficient, that ratio has
already been inverted by the product formula.  The two conventions may
not be used simultaneously.

Accordingly,
\Cref{prop:V35-formation-subprobability,
thm:V39-disjoint-tree-closed,
thm:V39-global-ternary-MTA,
thm:V45-formation-first-domination,
thm:V45-ff-pair-bank-tail,
thm:V45-ff-pair-closure,
thm:V45-ff-triple-bank,
thm:V45-ff-local-closure,
thm:V45-ff-complete-spectator}
are withdrawn as claims about arbitrary Fourier coefficient matrices.
The intermediate V43--V45 conditional claims listed in
\Cref{ass:V45-conditional-retraction} remain withdrawn as well.
Any later statement whose proof invokes those arbitrary-coefficient
formation tails is reopened.

The representation-theoretic statements themselves remain proved:
the V39 pair dimension-law tail, the V44 Casimir moment identities and
deterministic product-envelope lemma, and the V45 direct triple and
four-input dimension-law tails.  They may be used only after an
independent coefficient-sensitive factor has supplied the missing
dimension power.
\end{assessment}

For the next acquisition, define the genuine normalized coefficient
mass of a binary product by
\begin{equation}
 \mathfrak w_{A,B}(S)
 :=
 \frac{d_S\|\widehat{f_Ag_B}(S)\|_1}
 {(d_R\|A\|_1)(d_Q\|B\|_1)}.
 \label{eq:V45-genuine-coefficient-mass}
\end{equation}
The dual-singlet family shows that
\(\mathfrak w_{A,B}\) can carry only one inverse Casimir power on a
cold boundary representation.  Thus the next valid theorem cannot be a
bare formation tail.  It must estimate the \emph{assembled}
coefficient mass together with the binary defects, Racah operator, and
single final resolvent which occur in the actual ternary or quartic
cumulant.

\begin{openproblem}[Corrected coefficient-sensitive formation gate]
\label{op:V45-coefficient-sensitive-gate}
Starting from the exact Fourier product formula and
\eqref{eq:V45-genuine-coefficient-mass}, determine whether the complete
V35 disjoint-tree quotient satisfies its marked estimate for arbitrary
trace-class coefficients.  The proof must:
\begin{enumerate}[label=\textup{(\roman*)},leftmargin=3.7em]
\item retain the sharp dual-singlet mass \(1/d_R\);
\item keep both binary multiplier defects and the one final resolvent
      attached to the coefficient block;
\item use V39 only for an actual representation-dimension average which
      has been derived, not inserted by normalization;
\item recover the missing inverse Casimir power from an assembled
      operator cancellation, or exhibit a coefficient family for which
      the normalized marked quotient diverges.
\end{enumerate}
Until this gate is proved, global \(m=3\) MTA and every quartic closure
depending on it are open.
\end{openproblem}

\section{V45 ledger and next acquisition}
\label{sec:V45-ledger}

\begin{theorem}[Integrated V45 statement]
\label{thm:V45-integrated}
Preserving V1--V44, V45 proves:
\begin{enumerate}[label=\textup{(V45.\arabic*)},leftmargin=5.5em]
\item the scheduled SM V64/NS V31 audit with the strict transfer
      firewall \Cref{fire:V45-companion-types};
\item the arbitrary-parent counterexample
      \eqref{eq:V45-concentrated-parent} and the explicit withdrawal of
      the invalid conditional dimension-law route;
\item the exact dual-singlet Fourier-normalization counterexample
      \eqref{eq:V45-actual-singlet-mass}--%
      \eqref{eq:V45-singlet-dimension-law}, which separates genuine
      coefficient mass from the representation dimension law;
\item the direct three- and four-input \emph{representation
      dimension-law} squared cold tails, with the triple logarithm
      removed by adaptive rebracketing;
\item the withdrawal of the formation-first transfer and every global
      ternary or quartic closure whose proof depends on that transfer;
      and
\item the corrected coefficient-sensitive formation gate
      \Cref{op:V45-coefficient-sensitive-gate}, at which both binary
      multiplier defects, the Racah operator, and the single final
      resolvent must be estimated together.
\end{enumerate}
\end{theorem}

\begin{proof}
Item \textup{(V45.1)} is \Cref{sec:V45-companion-audit}.  Item
\textup{(V45.2)} is
\Cref{cex:V45-conditional-dimension-failure,
ass:V45-conditional-retraction}.  Item \textup{(V45.3)} is
\Cref{cth:V45-dual-singlet-normalization,
cor:V45-no-coefficient-dimension-tail}.  Item \textup{(V45.4)} is
\Cref{thm:V45-direct-triple-tail,rem:V45-log-resolution,
cor:V45-direct-four-tail}.  Item \textup{(V45.5)} is
\Cref{ass:V45-final-correction}.  Item \textup{(V45.6)} is
\Cref{op:V45-coefficient-sensitive-gate}.
\end{proof}

\begingroup
\small
\begin{longtable}{>{\raggedright\arraybackslash}p{0.29\textwidth}
                  >{\raggedright\arraybackslash}p{0.15\textwidth}
                  >{\raggedright\arraybackslash}p{0.48\textwidth}}
\toprule
\textbf{V45 row}&\textbf{Status}&\textbf{Advance or obstruction}\\
\midrule
\endfirsthead
\toprule
\textbf{V45 row}&\textbf{Status}&\textbf{Advance or obstruction}\\
\midrule
\endhead
Scheduled SM/NS audit
& \MGclosed
& SM V64 changes the splitting discipline but supplies no YM
  coefficient tail; NS V31
  is unchanged.\\[0.35em]
Conditional dimension-law insertion
& \MGfail
& A conditioned parent may be supported in one CG block.  The V43--V45
  conditional-tail proofs are explicitly withdrawn.\\[0.35em]
Direct triple and four-input cold tails
& \MGclosed
& These are representation-dimension averages only.  Reciprocal height
  and adaptive largest-first fusion remove the intermediate-shell
  logarithm.\\[0.35em]
Actual Fourier formation versus dimension law
& \MGfail
& A dual pure-state block has normalized singlet coefficient mass
  \(1/d_R\), whereas the dimension law gives \(1/d_R^2\).  Thus the
  formation-first transfer is false for arbitrary coefficients.\\[0.35em]
Global ternary MTA
& \MGopen
& The V39 closure depended on the false transfer.  The exact quotient
  must be re-estimated with both defects, Racah mixing, and the final
  resolvent still attached.\\[0.35em]
Quartic and spectator-dominated closures
& \MGopen
& Every claimed closure depending on V35/V39 formation
  subprobability is reopened; no quartic theorem may use the withdrawn
  artificial kernel.\\[0.35em]
Full Yang--Mills mass gap
& \MGopen
& All-order MTA, WUFC/CPM, construction, continuum, and OS spectral gates
  remain downstream.\\
\bottomrule
\end{longtable}
\endgroup

\begin{openproblem}[V46 mandate: coefficient-sensitive formation and
ternary revalidation]
\label{op:V46}
\begin{enumerate}[label=\textup{(V46.\arabic*)},leftmargin=5.5em]
\item Start from the exact Fourier product formula and the V29
      Fourier-algebra norm, writing every input and output dimension
      factor before introducing a positive majorant.
\item Work on \Cref{op:V45-coefficient-sensitive-gate} and retain the
      sharp dual-singlet coefficient mass \(1/d_R\) as a compulsory
      test family.
\item Reanalyse the exact V35 disjoint-tree quotient with both binary
      multiplier defects, the Racah recoupling operator, and its single
      final resolvent present before any absolute value or trace-norm
      estimate is taken.
\item Decide whether the missing inverse Casimir power is produced by
      an assembled defect/Racah cancellation.  If it is not, exhibit a
      resolved coefficient family for which the normalized marked
      quotient diverges.
\item Revalidate global \(m=3\) MTA from that coefficient-sensitive
      estimate before invoking ternary MTA anywhere in \(m=4\).
\item Only after the ternary row is valid may the quartic
      unique-junction and other topologies be revisited.  Do not use
      any withdrawn formation-first kernel as actual coefficient mass.
\item Preserve only the newest YM file.  The next SM/NS audit is V50.
\end{enumerate}
\end{openproblem}

\begin{firewall}[End-of-V45 claim boundary]
V45 proves the scheduled companion audit, two independent
normalization counterexamples, and the direct triple and four-input
representation-dimension cold tails.  It does not identify those
dimension-law tails with genuine arbitrary-coefficient Fourier mass.
Consequently the V35/V39 formation transfer, global \(m=3\) MTA, and
every dependent quartic or spectator-dominated closure are withdrawn
or reopened.  V45 does not prove coefficient-sensitive formation,
global \(m=3\) or \(m=4\) MTA, all-order MTA, WUFC, CPM, the nonlinear
Perron packet, finite-edge margins, capacity, protected-sector floors,
local-algebra noncollapse, reflection positivity, a nontrivial
continuum OS state, or a positive Yang--Mills spectral gap.  The full
Yang--Mills mass gap is not proved.
\end{firewall}

\section*{Version V45 append-only record}
\addcontentsline{toc}{section}{Version V45 append-only record}
The complete V44 source was retained before this continuation.  Its
SHA--256 digest was
\begin{center}\scriptsize\ttfamily
7b34082efdae7fb5e25feac73ec1eb1bff06d6caa00983cbc0c49f1f7c962285.
\end{center}
No mathematical content of V44 was altered.  On the next iteration,
retain this full file, append before its unique active document-closing
command, increment the filename to \texttt{\_V46.tex}, and remove only
the superseded V45 file from this Yang--Mills note series.  The next
scheduled SM/NS audit is V50.

% =============================================================================
% END APPENDED VERSION V45 MATERIAL
% =============================================================================

% =============================================================================
% BEGIN APPENDED VERSION V46 MATERIAL
% =============================================================================

\clearpage
\part{V46 --- Exact Fourier normalization, pointwise two-defect
closure, and the pair-only ternary atlas}
\label{part:V46}

\section{Purpose, correction order, and nonduplication}
\label{sec:V46-purpose}

V45 found that the representation dimension law is not the normalized
Fourier-coefficient law.  In the sharp dual-singlet family, the former
assigns mass \(d_R^{-2}\), whereas the latter has mass \(d_R^{-1}\).
Consequently the V39 squared cold tail cannot be inserted into an
arbitrary coefficient sum.

The present continuation goes one step farther upstream.  It derives
the exact Peter--Weyl product normalization before taking any trace
norm.  The output factor \(d_{\rm in}/d_T\) cancels the \(d_T\) in the
BFA norm.  What remains is raw isotypic pinching mass.  This mass is a
subprobability, but it has no representation-dimension-law domination.

That correction does not destroy the disjoint support-tree estimate.
The complete two-defect multiplier has a pointwise bound which makes a
formation tail unnecessary:
\[
 \frac{\Xi(\boldsymbol T)q_P|\delta_L\delta_R|}
 {1-q_{\boldsymbol T}}
 \le C\min\{H_L,H_R\}.
\]
A central-private versus overlap split then supplies the exact marked
factor \(H_LH_R/\Xi(\boldsymbol R_1)\).  The same estimate is stable
under the third pair bank, so all pair-only ternary incidence patterns,
including the three-edge cycle, close with genuine raw coefficient
mass.

The one-coordinate common-junction theorem is also revalidated with
the missing Fourier factor restored.  What remains open is the
operator-valued assembly of a common three-input junction with
arbitrary other coordinate banks.  No SM/NS audit is due at V46; the
next scheduled companion inspection remains V50.

\section{The exact Peter--Weyl product factor}
\label{sec:V46-Fourier-normalization}

Use throughout the V29 inversion convention
\begin{equation}
 h(g)=\sum_T d_T\Tr\!\left(\widehat h(T)D^T(g)\right).
 \label{eq:V46-Fourier-inversion}
\end{equation}
Let \(R_1,\ldots,R_m\) be irreducible representations,
\[
 d_{\rm in}:=\prod_{i=1}^m d_{R_i},
 \qquad
 \mathscr H:=\bigotimes_{i=1}^mH_{R_i},
\]
and choose a complete multiplicity-resolved decomposition
\begin{equation}
 \mathscr H
 =
 \bigoplus_T W_T(H_T\otimes M_T),
 \qquad
 W_T^*W_T=I_{H_T\otimes M_T}.
 \label{eq:V46-isotypic-decomposition}
\end{equation}
For trace-class \(A_i\in\mathcal S_1(H_{R_i})\), put
\begin{equation}
 \mathcal A:=A_1\otimes\cdots\otimes A_m,
 \qquad
 X_T:=W_T^*\mathcal A W_T.
 \label{eq:V46-raw-isotypic-block}
\end{equation}

\begin{theorem}[Exact product and multiplier normalization]
\label{thm:V46-exact-product-normalization}
Define
\[
 f_i(g):=d_{R_i}\Tr(A_iD^{R_i}(g)).
\]
Then
\begin{equation}
 \boxed{
 \widehat{\prod_{i=1}^m f_i}(T)
 =
 \frac{d_{\rm in}}{d_T}\Tr_{M_T}X_T.}
 \label{eq:V46-exact-product-normalization}
\end{equation}
More generally, if an assembled central row acts on the \(T\)-isotypic
multiplicity space by \(K_T\), then its exact output coefficient is
\begin{equation}
 \boxed{
 \widehat h(T)
 =
 \frac{d_{\rm in}}{d_T}
 \Tr_{M_T}\!\left[X_T(I_{H_T}\otimes K_T)\right].}
 \label{eq:V46-exact-multiplier-normalization}
\end{equation}
The position of a transpose of \(K_T\) depends only on the chosen
matrix-index convention; the displayed convention is the one used in
\eqref{eq:V37-partial-trace-assembly}.
\end{theorem}

\begin{proof}
The diagonal action on \(\mathscr H\) has block form
\[
 \bigotimes_iD^{R_i}(g)
 =
 \bigoplus_TW_T(D^T(g)\otimes I_{M_T})W_T^*.
\]
Hence
\begin{align*}
 \prod_i f_i(g)
 &=
 d_{\rm in}\Tr\!\left[
 \mathcal A\bigotimes_iD^{R_i}(g)\right]\\
 &=
 d_{\rm in}\sum_T
 \Tr\!\left[X_T(D^T(g)\otimes I_{M_T})\right]\\
 &=
 d_{\rm in}\sum_T
 \Tr\!\left[(\Tr_{M_T}X_T)D^T(g)\right].
\end{align*}
Comparison with \eqref{eq:V46-Fourier-inversion} proves
\eqref{eq:V46-exact-product-normalization}.  Inserting the multiplicity
operator before the last partial trace gives
\eqref{eq:V46-exact-multiplier-normalization}.
\end{proof}

\begin{corollary}[The genuine coefficient measure is raw pinching mass]
\label{cor:V46-raw-pinching-mass}
For every \(T\),
\begin{equation}
 d_T\|\widehat h(T)\|_1
 \le d_{\rm in}\|K_T\|_{\rm op}\|X_T\|_1.
 \label{eq:V46-BFA-dimension-cancellation}
\end{equation}
Moreover,
\begin{equation}
 \sum_T\|X_T\|_1\le\prod_i\|A_i\|_1.
 \label{eq:V46-global-raw-pinching}
\end{equation}
If \(P_{T,a}\) resolve the multiplicity copies and
\(X_{T,a}:=P_{T,a}X_TP_{T,a}\), then
\begin{equation}
 \omega_{T,a}:=
 \frac{\|X_{T,a}\|_1}{\prod_i\|A_i\|_1},
 \qquad
 \omega_{T,a}\ge0,
 \qquad
 \sum_{T,a}\omega_{T,a}\le1.
 \label{eq:V46-genuine-subprobability}
\end{equation}
No factor \(d_T/d_{\rm in}\) occurs in this resolved weight.
\end{corollary}

\begin{proof}
Equation \eqref{eq:V46-BFA-dimension-cancellation} follows from
\Cref{lem:V37-partial-trace} and
\eqref{eq:V46-exact-multiplier-normalization}.  Conjugate
\(\mathcal A\) by the full Clebsch--Gordan unitary and pinch to its
isotypic diagonal.  Trace-norm contractivity gives
\eqref{eq:V46-global-raw-pinching}.  A second pinching to the resolved
copy diagonal gives \eqref{eq:V46-genuine-subprobability}.
\end{proof}

\begin{remark}[Dual-singlet regression]
\label{rem:V46-dual-singlet-regression}
For the V45 pure dual pair, the raw singlet compression is
\[
 \left\langle\Omega_R,
 (|e_1\rangle\langle e_1|
 \otimes|e_1^*\rangle\langle e_1^*|)
 \Omega_R\right\rangle
 =\frac1{d_R},
\]
where
\(\Omega_R=d_R^{-1/2}\sum_a e_a\otimes e_a^*\).
Thus \eqref{eq:V46-genuine-subprobability} gives exactly the
\(d_R^{-1}\) mass in
\eqref{eq:V45-actual-singlet-mass}.  This checks both the factor
\(d_{\rm in}/d_T\) and the absence of a dimension-law factor.
\end{remark}

\begin{assessment}[Superseded normalization formulas]
\label{ass:V46-normalization-correction}
With \(X_T\) defined as the raw compression
\eqref{eq:V37-isotypic-input-block}, the right sides of
\eqref{eq:V36-same-link-coefficient} and
\eqref{eq:V37-partial-trace-assembly} are missing the factor
\(d_{\rm in}/d_T\).  Formula
\eqref{eq:V46-exact-multiplier-normalization} supersedes both identities.
The later estimate \eqref{eq:V37-subset-output-sum} already contains
\(d_{\rm in}\), so its inequality has the corrected normalization.

Likewise, the weights \eqref{eq:V35-formation-kernel} are not genuine
coefficient masses when \(\mathcal C_{\mu,\nu}\) denotes raw
Clebsch--Gordan compression.  They are replaced by the raw weights in
\eqref{eq:V46-genuine-subprobability}.  The pinching contraction and
the V35 fresh-leaf trace-norm taxes remain valid; only the extra
representation-dimension identification is discarded.
\end{assessment}

\section{Revalidation of the exact one-link junction}
\label{sec:V46-one-link-revalidation}

\begin{theorem}[Corrected multiplicity-uniform one-link ternary MTA]
\label{thm:V46-corrected-one-link-MTA}
On one link, retain the V37 notation
\[
 \mathcal A=A_1^{\rm in}\otimes A_2^{\rm in}
 \otimes A_3^{\rm in},
 \qquad
 d_{\rm in}=d_{R_1}d_{R_2}d_{R_3}.
\]
The exact complete third-cumulant coefficient is
\begin{equation}
 \widehat\kappa_{3,T}^{\alpha}
 =
 \frac{d_{\rm in}}{d_T}
 \Tr_{M_T}\!\left[
 X_T(I_{H_T}\otimes K_{123}^T)\right].
 \label{eq:V46-corrected-one-link-coefficient}
\end{equation}
For all sufficiently large \(\alpha\),
\begin{align}
 &\sum_{T\ne\boldsymbol1}
 e^{\rho\operatorname{ht}(T)}
 A_Td_T
 \frac{\|\widehat\kappa_{3,T}^{\alpha}\|_1}
      {1-\eta_T}
 \notag\\
 &\qquad\le
 C\prod_{i=1}^3
 \left[
 e^{\rho\operatorname{ht}(R_i)}
 A_i d_{R_i}\|A_i^{\rm in}\|_1
 \right].
 \label{eq:V46-corrected-one-link-MTA}
\end{align}
The constant is independent of all representations and multiplicities.
\end{theorem}

\begin{proof}
The coefficient identity is
\eqref{eq:V46-exact-multiplier-normalization}.  For every collection
\(\Omega\) of final channels,
\begin{align}
 \sum_{T\in\Omega}d_T
 \|\widehat\kappa_{3,T}^{\alpha}\|_1
 &\le
 d_{\rm in}
 \sup_{T\in\Omega}\|K_{123}^T\|_{\rm op}
 \sum_{T\in\Omega}\|X_T\|_1\\
 &\le
 d_{\rm in}
 \sup_{T\in\Omega}\|K_{123}^T\|_{\rm op}
 \prod_i\|A_i^{\rm in}\|_1.
 \label{eq:V46-one-link-subset-sum}
\end{align}
This is a direct consequence of
\Cref{cor:V46-raw-pinching-mass}, rather than a trivial
output-dimension bound.

If \(s_\alpha C_2(T)\le1\), the two-sided Wilson gap gives
\[
 \frac{A_T}{1-\eta_T}\le\frac C{s_\alpha}.
\]
Insert the V37 operator estimate
\[
 \|K_{123}^T\|_{\rm op}
 \le C\min\{s_\alpha^2A_1A_2A_3,1\}
\]
in \eqref{eq:V46-one-link-subset-sum}.  The low-output contribution is
at most
\[
 C d_{\rm in}A_1A_2A_3
 \prod_i\|A_i^{\rm in}\|_1.
\]
On the complementary output set, the final gap is bounded below,
\[
 A_T\le C(A_1+A_2+A_3)\le3C A_1A_2A_3,
\]
and the universal cap on \(K_{123}^T\) gives the same bound.  Finally
\(\operatorname{ht}(T)\le\sum_i\operatorname{ht}(R_i)\).
This proves \eqref{eq:V46-corrected-one-link-MTA}.
\end{proof}

\begin{lemma}[A one-link junction also has an output first moment]
\label{lem:V46-junction-first-moment}
For every occurring \(T\subset R_1\otimes R_2\otimes R_3\),
\begin{equation}
 \boxed{s_\alpha A_T\|K_{123}^T\|_{\rm op}\le C.}
 \label{eq:V46-junction-first-moment}
\end{equation}
\end{lemma}

\begin{proof}
Use the five-term operator identity
\eqref{eq:V33-common-basis-symbol}.  The term
\(s_\alpha A_T\eta_T\) is bounded by the V31 first-moment debit.
In the spectrum of \(\eta_3Q_{12}^T\), an intermediate
\(S\subset R_1\otimes R_2\) satisfies
\(T\subset S\otimes R_3\).  Fusion geometry gives
\[
 A_T\le C(A_S+A_3).
\]
Therefore
\[
 s_\alpha A_T\eta_S\eta_3
 \le
 C\left(s_\alpha A_S\eta_S+s_\alpha A_3\eta_3\right)
 \le C.
\]
Taking the spectral supremum proves the claim for this pair term.
The other two pair terms are identical.  For the last scalar term use
\(A_T\le C(A_1+A_2+A_3)\) and apply the one-link first-moment debit to
one factor at a time.  The sum of the five bounds proves
\eqref{eq:V46-junction-first-moment}.
\end{proof}

\begin{firewall}[Scope of the one-link restoration]
\label{fire:V46-one-link-scope}
\Cref{thm:V46-corrected-one-link-MTA} restores the exact one-coordinate
junction theorem.  It does not by itself supply the two globally
normalized marks when the same three inputs carry arbitrarily large
mass on other coordinate banks.  The first-moment estimate
\eqref{eq:V46-junction-first-moment} is the correct spectator debit for
that remaining assembly, not a completed global junction proof.
\end{firewall}

\section{Coefficient-free scalar preparation for a disjoint tree}
\label{sec:V46-scalar-preparation}

Let \(\boldsymbol R,\boldsymbol Q\) be exact product channels on the
same bank, let
\(\boldsymbol S\subset\boldsymbol R\otimes\boldsymbol Q\), and put
\[
 \delta(\boldsymbol S)
 :=q_{\alpha,\boldsymbol S}
   -q_{\alpha,\boldsymbol R}q_{\alpha,\boldsymbol Q}.
\]
Write \(H=\sum_eA_{R,e}A_{Q,e}\).

\begin{lemma}[Defect first moment and the two elementary debits]
\label{lem:V46-defect-first-moment}
Uniformly for sufficiently large \(\alpha\):
\begin{align}
 |\delta(\boldsymbol S)|
 &\le C s_\alpha H,
 \label{eq:V46-defect-cap}\\
 \Xi(\boldsymbol S)
 \frac{|\delta(\boldsymbol S)|}
 {1-q_{\alpha,\boldsymbol S}}
 &\le C H,
 \label{eq:V46-binary-quotient}\\
 s_\alpha\Xi(\boldsymbol S)|\delta(\boldsymbol S)|
 &\le C.
 \label{eq:V46-defect-first-moment}
\end{align}
When \(\boldsymbol S\) is trivial, the last two left sides are
interpreted as zero.  For every nontrivial exact product channel
\(\boldsymbol V\),
\begin{align}
 s_\alpha\Xi(\boldsymbol V)q_{\alpha,\boldsymbol V}
 &\le C,
 \label{eq:V46-product-first-moment}\\
 s_\alpha C_E(\boldsymbol V)\le1
 \quad\Longrightarrow\quad
 \frac{\Xi(\boldsymbol V)}
 {1-q_{\alpha,\boldsymbol V}}
 &\le\frac C{s_\alpha}.
 \label{eq:V46-low-output-resolvent}
\end{align}
\end{lemma}

\begin{proof}
Equations \eqref{eq:V46-defect-cap} and
\eqref{eq:V46-binary-quotient} are respectively
\eqref{eq:V31-overlap-defect-upper} and
\eqref{eq:V31-unscaled-marked-quotient}.
The product first moment is
\eqref{eq:V31-balanced-product-damping}.  Exact activity gives
\(\Xi(\boldsymbol V)\le C C_E(\boldsymbol V)\); the lower Wilson gap
then proves \eqref{eq:V46-low-output-resolvent}.

It remains to prove \eqref{eq:V46-defect-first-moment}.  Since all
multipliers are nonnegative,
\[
 |\delta(\boldsymbol S)|
 \le q_{\boldsymbol S}+q_{\boldsymbol R}q_{\boldsymbol Q}.
\]
The first summand is paid by
\eqref{eq:V46-product-first-moment}.  Coordinatewise fusion geometry
gives
\[
 \Xi(\boldsymbol S)
 \le C\bigl(\Xi(\boldsymbol R)+\Xi(\boldsymbol Q)\bigr).
\]
Consequently
\begin{align*}
 s_\alpha\Xi(\boldsymbol S)q_{\boldsymbol R}q_{\boldsymbol Q}
 &\le
 C\left[
 s_\alpha\Xi(\boldsymbol R)q_{\boldsymbol R}q_{\boldsymbol Q}
 +s_\alpha\Xi(\boldsymbol Q)q_{\boldsymbol Q}q_{\boldsymbol R}
 \right]\\
 &\le C
\end{align*}
by \eqref{eq:V46-product-first-moment}.  This proves the remaining
claim.
\end{proof}

\section{The pointwise two-defect theorem}
\label{sec:V46-pointwise-two-defect}

Use the V35 disjoint-overlap banks
\[
 L=X_1\cap X_2,\qquad
 R=X_1\cap X_3,\qquad
 X_2\cap X_3=\varnothing,
\]
with resolved outputs \(\boldsymbol S,\boldsymbol U\).  Retain
\[
 \delta_L=q_{\boldsymbol S}-q_{1,L}q_{2,L},
 \qquad
 \delta_R=q_{\boldsymbol U}-q_{1,R}q_{3,R},
\]
and \(H_L,H_R,q_P\) from V35.  More generally, reinterpret \(P\) as
the coordinates active in exactly one input and allow one additional
bank, disjoint from \(L\), \(R\), and \(P\), which is active only in
inputs two and three.  Denote its resolved output by
\(\boldsymbol V\) and its scalar row factor by \(r_V\), and assume
\begin{equation}
 |r_V|\le1,
 \qquad
 s_\alpha\Xi(\boldsymbol V)|r_V|\le C_V.
 \label{eq:V46-admissible-spectator}
\end{equation}
The original V35 support tree is the empty choice
\(\boldsymbol V=\boldsymbol1\), \(r_V=1\).  For a nonconstant final
product channel \(\boldsymbol T\), put
\begin{align}
 q_{\boldsymbol T}
 &:=
 q_Pq_{\boldsymbol S}q_{\boldsymbol U}q_{\boldsymbol V},
 \label{eq:V46-extended-final-multiplier}\\
 \Xi(\boldsymbol T)
 &:=
 \Xi_{\boldsymbol S}+\Xi_{\boldsymbol U}
 +\Xi_P+\Xi_{\boldsymbol V},
 \label{eq:V46-extended-final-mass}\\
 \mathcal Q_{L,R}^{V}
 &:=
 \frac{\Xi(\boldsymbol T)q_P
 |\delta_L\delta_Rr_V|}
 {1-q_{\boldsymbol T}}.
 \label{eq:V46-extended-two-defect-quotient}
\end{align}

\begin{theorem}[Pointwise minimum and one-resolvent bounds]
\label{thm:V46-two-defect-pointwise}
There is \(C<\infty\), independent of supports, representations,
outputs, and multiplicities, such that
\begin{align}
 \boxed{
 \mathcal Q_{L,R}^{V}
 \le C\min\{H_L,H_R\},}
 \label{eq:V46-two-defect-minimum}\\
 \boxed{
 \mathcal Q_{L,R}^{V}
 \le C s_\alpha H_LH_R.}
 \label{eq:V46-two-defect-one-resolvent}
\end{align}
\end{theorem}

\begin{proof}
First prove the bound by \(CH_R\).  If
\(s_\alpha C_E(\boldsymbol T)\le1\), then
\eqref{eq:V46-low-output-resolvent} and
\eqref{eq:V46-defect-cap} give
\[
 \mathcal Q_{L,R}^{V}
 \le\frac C{s_\alpha}|\delta_R|
 \le C H_R.
\]

Suppose \(s_\alpha C_E(\boldsymbol T)>1\).  The final denominator is
bounded below.  Split the numerator according to
\eqref{eq:V46-extended-final-mass}.  For the
\(\Xi_{\boldsymbol S}\) part use
\eqref{eq:V46-defect-first-moment} on \(\delta_L\) and
\eqref{eq:V46-defect-cap} on \(\delta_R\):
\[
 \Xi_{\boldsymbol S}q_P|\delta_L\delta_Rr_V|
 \le\frac C{s_\alpha}(Cs_\alpha H_R)
 \le C H_R.
\]
For the \(\Xi_{\boldsymbol U}\) part,
\[
 \Xi_{\boldsymbol U}|\delta_R|
 \le (1-q_{\boldsymbol U})CH_R
 \le CH_R
\]
by \eqref{eq:V46-binary-quotient}.  The remaining factors have modulus
at most one.  For the private and auxiliary pieces use respectively
\eqref{eq:V46-product-first-moment} and
\eqref{eq:V46-admissible-spectator}:
\begin{align*}
 \Xi_Pq_P|\delta_L\delta_Rr_V|
 &\le\frac C{s_\alpha}(Cs_\alpha H_R),\\
 \Xi_{\boldsymbol V}q_P|\delta_L\delta_Rr_V|
 &\le\frac {C_V}{s_\alpha}(Cs_\alpha H_R).
\end{align*}
Thus \(\mathcal Q_{L,R}^{V}\le CH_R\) in both thermal cases.
Interchanging \(L\) and \(R\) proves the bound by \(CH_L\), hence
\eqref{eq:V46-two-defect-minimum}.

For \eqref{eq:V46-two-defect-one-resolvent}, the low-output case follows
from
\[
 \frac{\Xi(\boldsymbol T)}{1-q_{\boldsymbol T}}
 \le\frac C{s_\alpha},
 \qquad
 |\delta_L\delta_R|
 \le C s_\alpha^2H_LH_R.
\]
In the high-output case the denominator is bounded below.  Fusion
geometry gives
\[
 \Xi_{\boldsymbol S}\le CH_L,\qquad
 \Xi_{\boldsymbol U}\le CH_R.
\]
Therefore the two overlap-output pieces are at most
\[
 CH_L(Cs_\alpha H_R),
 \qquad
 CH_R(Cs_\alpha H_L).
\]
The private and auxiliary pieces are at most
\[
 \frac C{s_\alpha}
 (Cs_\alpha H_L)(Cs_\alpha H_R),
 \qquad
 \frac {C_V}{s_\alpha}
 (Cs_\alpha H_L)(Cs_\alpha H_R).
\]
Adding the four pieces proves
\eqref{eq:V46-two-defect-one-resolvent}.
\end{proof}

\begin{theorem}[Exact central marked payment]
\label{thm:V46-central-marked-payment}
Let
\begin{equation}
 x:=\Xi(\boldsymbol R_{1,L}),\qquad
 y:=\Xi(\boldsymbol R_{1,R}),\qquad
 p:=\Xi(\boldsymbol R_{1,P_1}),
 \qquad
 \Xi_1=x+y+p.
 \label{eq:V46-central-mass-split}
\end{equation}
Then every resolved block obeys
\begin{equation}
 \boxed{
 \mathcal Q_{L,R}^{V}
 \le C\,\frac{H_LH_R}{\Xi_1}.}
 \label{eq:V46-central-marked-payment}
\end{equation}
\end{theorem}

\begin{proof}
Because every leaf representation is nontrivial on its overlap bank,
\[
 x\le H_L,\qquad y\le H_R.
\]
If \(p\le x+y\), then
\[
 \Xi_1\le2(x+y)\le2(H_L+H_R).
\]
Consequently
\[
 \frac{H_LH_R}{\Xi_1}
 \ge
 \frac{H_LH_R}{2(H_L+H_R)}
 \ge\frac14\min\{H_L,H_R\}.
\]
Apply \eqref{eq:V46-two-defect-minimum}.

Suppose \(p>x+y\), so \(p>\Xi_1/2\).  If
\(s_\alpha p\le M\), then
\eqref{eq:V46-two-defect-one-resolvent} gives
\[
 \mathcal Q_{L,R}^{V}
 \le Cs_\alpha H_LH_R
 \le\frac{CM}{p}H_LH_R
 \le\frac{2CM}{\Xi_1}H_LH_R.
\]

It remains to take \(s_\alpha p>M\).  Quadratic product damping on the
central private channel gives
\[
 q_P\le q_{1,P_1}
 \le\frac C{(s_\alpha p)^2}.
\]
Choose \(M\) so that \(q_P\le1/2\).  Then
\(1-q_{\boldsymbol T}\ge1/2\).  The binary quotient implies
\[
 \Xi_{\boldsymbol S}|\delta_L|\le CH_L,
 \qquad
 \Xi_{\boldsymbol U}|\delta_R|\le CH_R.
\]
Thus the two overlap-output pieces are bounded by
\[
 Cq_Ps_\alpha H_LH_R
 \le \frac C{p}H_LH_R.
\]
For the private piece, the V33 quadratic private debit gives
\[
 (s_\alpha\Xi_P)^2q_P\le C,
\]
and therefore
\[
 \Xi_Pq_P|\delta_L\delta_R|
 \le C\Xi_Pq_Ps_\alpha^2H_LH_R
 \le\frac C{\Xi_P}H_LH_R
 \le\frac C{p}H_LH_R.
\]
Finally, the auxiliary piece satisfies
\begin{align*}
 \Xi_{\boldsymbol V}q_P|\delta_L\delta_Rr_V|
 &\le
 \frac {C_V}{s_\alpha}q_P
 (Cs_\alpha H_L)(Cs_\alpha H_R)\\
 &\le\frac C{p}H_LH_R.
\end{align*}
The last step uses
\(q_Ps_\alpha\le C/p\), which follows from
\(q_P\le C/(s_\alpha p)^2\) and \(s_\alpha p>M\).
After the denominator floor and \(p>\Xi_1/2\), these estimates prove
\eqref{eq:V46-central-marked-payment}.
\end{proof}

\section{Genuine-coefficient closure of every disjoint support tree}
\label{sec:V46-disjoint-coefficient-closure}

\begin{theorem}[Disjoint-overlap ternary MTA without a formation law]
\label{thm:V46-disjoint-tree-MTA}
Assume the V35 support tree
\[
 X_2\cap X_3=\varnothing,\qquad
 L=X_1\cap X_2\ne\varnothing,\qquad
 R=X_1\cap X_3\ne\varnothing.
\]
For pure exact Peter--Weyl input channels and arbitrary trace-class
coefficient matrices \(\mathsf A_i\), the center-one contribution
satisfies
\begin{align}
 &\sum_{\boldsymbol T\ne\boldsymbol1}
 e^{\rho\operatorname{ht}(\boldsymbol T)}
 \Xi(\boldsymbol T)d_{\boldsymbol T}
 \left\|
 \widehat{\mathcal RQ_0
 \kappa_3^\alpha}(\boldsymbol T)
 \right\|_1
 \notag\\
 &\qquad\le
 C e^{\rho\sum_i\operatorname{ht}(\boldsymbol R_i)}
 d_{\rm in}
 \frac{H_LH_R}{\Xi(\boldsymbol R_1)}
 \prod_i\|\mathsf A_i\|_1.
 \label{eq:V46-disjoint-tree-MTA}
\end{align}
After summing input channels, the right side is exactly the center-one
marked-BFA tree term.  The estimate is uniform in supports,
representations, LR multiplicities, and volume.
\end{theorem}

\begin{proof}
In the resolved \((\boldsymbol S,\mu;\boldsymbol U,\nu)\) basis, let
\[
 \mathcal C_{\mu,\nu}
 :=
 \mathcal C_{\mu,\nu}
 (\mathsf A_1,\mathsf A_2,\mathsf A_3)
\]
denote the raw sequential Clebsch--Gordan compression.  The exact
normalization theorem gives
\begin{equation}
 \widehat{\mathcal RQ_0\kappa_3^\alpha}
 (\boldsymbol T)
 =
 \frac{d_{\rm in}}{d_{\boldsymbol T}}
 \frac{q_P}{1-q_{\boldsymbol T}}
 \sum_{\mu,\nu}
 \delta_L(\boldsymbol S)\delta_R(\boldsymbol U)
 \mathcal C_{\mu,\nu}.
 \label{eq:V46-disjoint-exact-coefficient}
\end{equation}
The sum is understood over all resolved copies producing the same final
product channel.

Multiply by the output BFA factor and take the triangle inequality:
\begin{align}
 &\sum_{\boldsymbol T\ne\boldsymbol1}
 \Xi(\boldsymbol T)d_{\boldsymbol T}
 \left\|
 \widehat{\mathcal RQ_0\kappa_3^\alpha}
 (\boldsymbol T)\right\|_1\\
 &\quad\le
 d_{\rm in}
 \sum_{\boldsymbol S,\mu,\boldsymbol U,\nu}
 \frac{\Xi(\boldsymbol T)q_P
 |\delta_L\delta_R|}
 {1-q_{\boldsymbol T}}
 \|\mathcal C_{\mu,\nu}\|_1.
 \label{eq:V46-disjoint-raw-sum}
\end{align}
Apply \Cref{thm:V46-central-marked-payment} with the empty auxiliary
bank.  The simultaneous resolved CG pinching is trace-norm contractive:
\begin{equation}
 \sum_{\boldsymbol S,\mu,\boldsymbol U,\nu}
 \|\mathcal C_{\mu,\nu}\|_1
 \le\prod_i\|\mathsf A_i\|_1.
 \label{eq:V46-disjoint-pinching}
\end{equation}
This proves the unweighted version of
\eqref{eq:V46-disjoint-tree-MTA}.  Fusion height subadditivity inserts
the displayed exponential factor.

For pure input channels,
\begin{align*}
 &\prod_{i=1}^3
 \left[\Xi_i d_{\boldsymbol R_i}\|\mathsf A_i\|_1\right]
 \sum_{\substack{e\in L\\f\in R}}
 \vartheta_e(\boldsymbol R_1)
 \vartheta_f(\boldsymbol R_1)
 \vartheta_e(\boldsymbol R_2)
 \vartheta_f(\boldsymbol R_3)\\
 &\qquad=
 d_{\rm in}\frac{H_LH_R}{\Xi_1}
 \prod_i\|\mathsf A_i\|_1.
\end{align*}
Thus the scalar on the right is precisely the two-marked center-one
BFA weight.  Tonelli then factors the input channel sums into the three
marked seminorms.
\end{proof}

\begin{assessment}[What replaces the V39 cold-tail insertion]
\label{ass:V46-disjoint-replacement}
\Cref{thm:V46-disjoint-tree-MTA} restores the disjoint-overlap
conclusion formerly claimed in
\Cref{thm:V39-disjoint-tree-closed}, but not its proof.  The invalid
line was the comparison of actual coefficient mass with the product
dimension law.  The replacement uses
\[
 \text{exact }d_{\rm in}/d_T\text{ normalization}
 \;+\;
 \text{raw pinching}
 \;+\;
 \text{the pointwise two-defect quotient}.
\]
The V39 BZ dimension-law tail remains a true representation theorem,
but it is not used here.
\end{assessment}

\section{The third pair bank and the cyclic incidence sector}
\label{sec:V46-cyclic-pair-sector}

Assume now that there is no common three-input coordinate, and decompose
the three exclusive pair banks as
\[
 L_{12},\qquad L_{13},\qquad L_{23}.
\]
On a fixed resolved output write
\[
 A_{ij}:=q_{\boldsymbol S_{ij}},
 \qquad
 a_{ij}:=q_{i,L_{ij}}q_{j,L_{ij}},
 \qquad
 \delta_{ij}:=A_{ij}-a_{ij}.
\]

\begin{lemma}[Exact three-pair connected multiplier]
\label{lem:V46-three-pair-multiplier}
After factoring all private singleton multipliers \(q_P\), the complete
five-term third-cumulant multiplier is
\begin{equation}
 \boxed{
 K_{\rm pair}
 =
 q_P\left(
 a_{23}\delta_{12}\delta_{13}
 +a_{13}\delta_{12}\delta_{23}
 +a_{12}\delta_{13}\delta_{23}
 +\delta_{12}\delta_{13}\delta_{23}
 \right).}
 \label{eq:V46-three-pair-multiplier}
\end{equation}
\end{lemma}

\begin{proof}
Before subtracting the disconnected partitions, the five scalar terms
on the three pair banks are
\[
 A_{12}A_{13}A_{23},\quad
 -A_{12}a_{13}a_{23},\quad
 -a_{12}A_{13}a_{23},\quad
 -a_{12}a_{13}A_{23},\quad
 +2a_{12}a_{13}a_{23}.
\]
Substitute \(A_{ij}=a_{ij}+\delta_{ij}\).  The constant and all
one-defect terms cancel.  The three two-defect terms and the
three-defect term are exactly
\eqref{eq:V46-three-pair-multiplier}.  Private coordinates contribute
the common factor \(q_P\).
\end{proof}

\begin{lemma}[The unused pair factor is admissible]
\label{lem:V46-unused-pair-admissible}
For the resolved output \(\boldsymbol S_{ij}\), both choices
\[
 r_{ij}=a_{ij}
 \qquad\hbox{and}\qquad
 r_{ij}=\delta_{ij}
\]
satisfy
\begin{equation}
 |r_{ij}|\le1,
 \qquad
 s_\alpha\Xi(\boldsymbol S_{ij})|r_{ij}|\le C.
 \label{eq:V46-unused-pair-admissible}
\end{equation}
\end{lemma}

\begin{proof}
For \(r_{ij}=\delta_{ij}\), this is
\eqref{eq:V46-defect-first-moment}.  For \(r_{ij}=a_{ij}\), fusion
geometry gives
\[
 \Xi(\boldsymbol S_{ij})
 \le C\left(\Xi(\boldsymbol R_{i,L_{ij}})
       +\Xi(\boldsymbol R_{j,L_{ij}})\right).
\]
Multiplication by the two singleton product multipliers and
\eqref{eq:V46-product-first-moment} prove the result.
\end{proof}

\begin{theorem}[Pair-only ternary marked-tree atlas]
\label{thm:V46-pair-only-ternary-MTA}
If no coordinate is active in all three inputs, the full third
cumulant satisfies the \(m=3\) marked-tree atlas inequality uniformly
over arbitrary exact supports, representations, coefficient matrices,
multiplicities, and volume.
\end{theorem}

\begin{proof}
Apply \Cref{thm:V46-central-marked-payment} to the first term of
\eqref{eq:V46-three-pair-multiplier}, with selected banks
\((L_{12},L_{13})\), center one, and auxiliary factor
\(r_V=a_{23}\).  The admissibility hypothesis is
\Cref{lem:V46-unused-pair-admissible}.  This produces the tree centered
at one.  The next two terms give the trees centered at two and three.
For the last term, select any two nonempty pair banks forming a spanning
tree and treat the remaining defect as \(r_V\).  If fewer than two pair
banks are nonempty, the connected multiplier is zero.

For each term, use the exact Fourier factor
\eqref{eq:V46-exact-multiplier-normalization} and raw isotypic pinching
exactly as in
\eqref{eq:V46-disjoint-raw-sum}--%
\eqref{eq:V46-disjoint-pinching}.  The three pair banks are disjoint
coordinate carriers, so their moment operators are simultaneously
diagonal in the resolved product basis; no Racah change or entrywise
multiplicity sum occurs.  Each scalar
\(H_{ij}H_{ik}/\Xi(\boldsymbol R_i)\) is converted to the corresponding
marked BFA seminorms by
\eqref{eq:V32-tree-mass-factorization}.  Summing the four terms changes
only the universal constant.
\end{proof}

\section{V46 ledger and next acquisition}
\label{sec:V46-ledger}

\begin{theorem}[Integrated V46 statement]
\label{thm:V46-integrated}
Preserving V1--V45, V46 proves:
\begin{enumerate}[label=\textup{(V46.\arabic*)},leftmargin=5.5em]
\item the exact Fourier product factor \(d_{\rm in}/d_T\) and the
      corresponding multiplicity-operator formula;
\item cancellation of \(d_T\) in the BFA output weight, leaving genuine
      raw pinching mass as the coefficient subprobability;
\item revalidation of the multiplicity-uniform one-link ternary MTA
      with the corrected normalization;
\item the uniform defect first-moment debit
      \eqref{eq:V46-defect-first-moment};
\item the pointwise minimum and one-resolvent estimates
      \eqref{eq:V46-two-defect-minimum}--%
      \eqref{eq:V46-two-defect-one-resolvent};
\item coefficient-uniform closure of every disjoint-overlap support
      tree without any representation formation law; and
\item the complete pair-only ternary atlas, including the cyclic
      three-pair incidence pattern.
\end{enumerate}
\end{theorem}

\begin{proof}
Items \textup{(V46.1)--(V46.2)} are
\Cref{thm:V46-exact-product-normalization,
cor:V46-raw-pinching-mass}.  Item \textup{(V46.3)} is
\Cref{thm:V46-corrected-one-link-MTA}.  Item
\textup{(V46.4)} is \Cref{lem:V46-defect-first-moment}.  Item
\textup{(V46.5)} is \Cref{thm:V46-two-defect-pointwise}.  Item
\textup{(V46.6)} is \Cref{thm:V46-disjoint-tree-MTA}.  Item
\textup{(V46.7)} is
\Cref{lem:V46-three-pair-multiplier,
thm:V46-pair-only-ternary-MTA}.
\end{proof}

\begingroup
\small
\begin{longtable}{>{\raggedright\arraybackslash}p{0.29\textwidth}
                  >{\raggedright\arraybackslash}p{0.15\textwidth}
                  >{\raggedright\arraybackslash}p{0.48\textwidth}}
\toprule
\textbf{V46 row}&\textbf{Status}&\textbf{Advance or obstruction}\\
\midrule
\endfirsthead
\toprule
\textbf{V46 row}&\textbf{Status}&\textbf{Advance or obstruction}\\
\midrule
\endhead
Exact Fourier normalization
& \MGclosed
& The coefficient factor is \(d_{\rm in}/d_T\); the BFA output
  dimension cancels it exactly.\\[0.35em]
Genuine resolved coefficient mass
& \MGclosed
& Raw CG pinching is subprobability.  No dimension-law domination is
  asserted or needed.\\[0.35em]
One-coordinate ternary junction
& \MGclosed
& The V37 operator estimate and isotypic pinching prove the corrected
  one-link MTA theorem.\\[0.35em]
Disjoint-overlap support trees
& \MGclosed
& A pointwise two-defect bound replaces the invalid V39 cold-tail
  insertion.\\[0.35em]
Pair-only cyclic incidence
& \MGclosed
& The exact three-pair multiplier is four spanning-tree terms with an
  admissible unused-bank debit.\\[0.35em]
Common-junction mixed incidence
& \MGopen
& A local three-input junction combined with arbitrary pair/private
  spectator banks still needs a global two-mark operator assembly.\\[0.35em]
Global ternary MTA
& \MGopen
& It will be restored only after the common-junction mixed sector is
  proved with the corrected Fourier factor.\\[0.35em]
Quartic and all-order MTA
& \MGopen
& Every use of ternary MTA remains suspended until the preceding row
  closes.\\[0.35em]
Full Yang--Mills mass gap
& \MGopen
& WUFC/CPM, nonlinear construction, continuum, OS positivity and
  nontriviality, and the spectral-gap passage remain downstream.\\
\bottomrule
\end{longtable}
\endgroup

\begin{openproblem}[V47 mandate: common-junction spectator assembly]
\label{op:V47}
\begin{enumerate}[label=\textup{(V47.\arabic*)},leftmargin=5.5em]
\item On the exact coordinate partition into common-junction,
      exclusive-pair, and private banks, expand the complete third
      cumulant into connected local factors before any trace norm.
\item Use \eqref{eq:V46-junction-first-moment} to retain the output
      debit of every unselected junction block.  Do not replace a local
      multiplicity operator by entrywise absolute values.
\item For a selected common junction \(e\), prove the global normalized
      tree payment
      \[
      C A_{1,e}A_{2,e}A_{3,e}
      \sum_{i=1}^3\frac{A_{i,e}}{\Xi(\boldsymbol R_i)}
      \]
      after all spectator factors and the single final resolvent are
      attached.
\item Route every term containing two exclusive pair defects through
      \Cref{thm:V46-central-marked-payment}; use the common-junction
      operator only as an admissible spectator when its first-moment
      hypothesis has been verified.
\item Reprove the full global \(m=3\) MTA from the corrected
      \(d_{\rm in}/d_T\) formula.  Do not cite the withdrawn
      representation-law formation transfer.
\item Only after global ternary MTA is restored may any \(m=4\) theorem
      be reactivated.
\item Preserve only the newest YM file.  The next SM/NS audit remains
      V50.
\end{enumerate}
\end{openproblem}

\begin{firewall}[End-of-V46 claim boundary]
\label{fire:V46-claim-boundary}
V46 proves the exact Peter--Weyl normalization, genuine raw pinching
subprobability, corrected one-link junction theorem, pointwise
disjoint-tree payment, and the complete pair-only ternary atlas.  It
does not prove the mixed common-junction spectator sector or global
\(m=3\) MTA.  Consequently it does not reactivate any quartic theorem.
It also does not prove all-order MTA, WUFC, CPM, the nonlinear Perron
packet, finite-edge margins, capacity, protected-sector floors,
local-algebra noncollapse, reflection positivity, a nontrivial
continuum OS state, or a positive Yang--Mills spectral gap.  The full
Yang--Mills mass gap is not proved.
\end{firewall}

\section*{Version V46 append-only record}
\addcontentsline{toc}{section}{Version V46 append-only record}
The complete V45 source was retained before this continuation.  Its
SHA--256 digest was
\begin{center}\scriptsize\ttfamily
1f743d84c57e28fb52ed3f43665e06817b646be70f8d1d1d93b081789ff208b3.
\end{center}
No mathematical content of V45 was altered.  On the next iteration,
retain this full file, append before its unique active document-closing
command, increment the filename to \texttt{\_V47.tex}, and remove only
the superseded V46 file from this Yang--Mills note series.  The next
scheduled SM/NS audit is V50.

% =============================================================================
% END APPENDED VERSION V46 MATERIAL
% =============================================================================

% =============================================================================
% BEGIN APPENDED VERSION V47 MATERIAL
% =============================================================================

\chapter{V47: Common-bank connected polynomial and the high-defect sector}
\label{ch:V47}

\section{Purpose and nonduplication boundary}
\label{sec:V47-purpose}

V46 settled the exact Fourier normalization, the disjoint two-defect
payment, and every pair-only ternary incidence.  The first unresolved
support geometry is therefore not another disjoint tree.  It is a bank
on which all three inputs are active, accompanied by arbitrary
exclusive-pair and private banks.  This continuation gives the exact
operator polynomial for that geometry and proves a reduced-center
payment for all terms containing at least two exclusive-pair defects.
It also isolates the missing full-center upgrade.  The common-bank
multiplicity operator is retained until the final trace norm; no
entrywise Racah sum and no representation-dimension formation law is
used.

This is not the scheduled five-version comparison.  In accordance with
the audit rule, the next comparison with the current Standard Model and
Navier--Stokes notes remains V50.

\section{Exact partition and common-bank operators}
\label{sec:V47-partition}

Let the exact active supports of three pure Peter--Weyl inputs be
\(X_1,X_2,X_3\).  Partition their union into
\begin{align}
 J&:=X_1\cap X_2\cap X_3,
 \label{eq:V47-common-bank}\\
 L_{ij}&:=(X_i\cap X_j)\setminus J,
 \qquad 1\le i<j\le3,
 \label{eq:V47-exclusive-pair-banks}\\
 P_i&:=X_i\setminus(X_j\cup X_k),
 \qquad \{i,j,k\}=\{1,2,3\}.
 \label{eq:V47-private-banks}
\end{align}
All seven displayed banks are pairwise disjoint.  Resolve on each
exclusive-pair bank an output channel \(\boldsymbol S_{ij}\), and put
\begin{equation}
 \beta_{ij}:=q_{\alpha,\boldsymbol S_{ij}},\qquad
 a_{ij}:=q_{\alpha,\boldsymbol R_{i,L_{ij}}}
          q_{\alpha,\boldsymbol R_{j,L_{ij}}},\qquad
 \delta_{ij}:=\beta_{ij}-a_{ij}.
 \label{eq:V47-pair-data}
\end{equation}
Empty banks have \(\beta_{ij}=a_{ij}=1\) and
\(\delta_{ij}=0\).  Let \(q_P\) be the product of all singleton
multipliers on the three private banks.

Fix a final common-bank channel \(\boldsymbol T_J\), including its full
multiplicity space \(M_{\boldsymbol T_J}\).  On that space define
\begin{align}
 M_J&:=q_{\alpha,\boldsymbol T_J}I,
 \label{eq:V47-MJ}\\
 B_{ij;k}&:=q_{\alpha,\boldsymbol R_{k,J}}
             Q_{ij,J}^{\boldsymbol T_J},
 \qquad \{i,j,k\}=\{1,2,3\},
 \label{eq:V47-Bijk}\\
 G_J&:=\prod_{i=1}^3q_{\alpha,\boldsymbol R_{i,J}}I,
 \label{eq:V47-GJ}\\
 C_{ij;k}&:=M_J-B_{ij;k},
 \label{eq:V47-cut-operator}\\
 K_J&:=M_J-B_{12;3}-B_{13;2}-B_{23;1}+2G_J.
 \label{eq:V47-common-cumulant}
\end{align}
Thus \(K_J\) is exactly the V33 common-basis ternary cumulant on
\(J\), whereas \(C_{ij;k}\) is the binary cut between the fused pair
\(ij\) and the remaining input \(k\).  The three \(B_{ij;k}\) need
not commute.  Every formula below is an operator identity on
\(M_{\boldsymbol T_J}\), tensored with scalar resolved factors on the
disjoint banks.

\begin{theorem}[Exact common/pair-bank connected polynomial]
\label{thm:V47-connected-polynomial}
On every fully resolved channel, the complete third-cumulant multiplier
is
\begin{align}
 K_{\rm full}=q_P\bigl[{}
 &a_{12}a_{13}a_{23}K_J
 +\delta_{12}a_{13}a_{23}C_{12;3}
 +a_{12}\delta_{13}a_{23}C_{13;2}
 \notag\\
 &+a_{12}a_{13}\delta_{23}C_{23;1}
 +\delta_{12}\delta_{13}a_{23}M_J
 +\delta_{12}a_{13}\delta_{23}M_J
 \notag\\
 &+a_{12}\delta_{13}\delta_{23}M_J
 +\delta_{12}\delta_{13}\delta_{23}M_J\bigr].
 \label{eq:V47-connected-polynomial}
\end{align}
If \(J=\varnothing\), this reduces exactly to the four-term pair-only
identity \eqref{eq:V46-three-pair-multiplier}.
\end{theorem}

\begin{proof}
Before expanding \(\beta_{ij}=a_{ij}+\delta_{ij}\), the five partitions
of \(\{1,2,3\}\) give
\begin{align*}
 q_P\bigl[{}
 &M_J\beta_{12}\beta_{13}\beta_{23}
 -B_{12;3}\beta_{12}a_{13}a_{23}
 -B_{13;2}a_{12}\beta_{13}a_{23}\\
 &-B_{23;1}a_{12}a_{13}\beta_{23}
 +2G_Ja_{12}a_{13}a_{23}\bigr].
\end{align*}
This follows bank by bank: a cumulant block joins precisely those
inputs belonging to it, and disjoint coordinate banks tensor.  In
particular, the \(12|3\) partition sees \(B_{12;3}\) on \(J\), the
joint multiplier \(\beta_{12}\) on \(L_{12}\), and the separated
multipliers \(a_{13},a_{23}\) on the other pair banks.

Now substitute \(\beta_{ij}=a_{ij}+\delta_{ij}\).  The coefficient of
the constant monomial is \(K_J\).  The coefficient of
\(\delta_{ij}\) is \(M_J-B_{ij;k}=C_{ij;k}\).  Every monomial with two
or three defects occurs only in the all-in-one partition and hence has
coefficient \(M_J\).  This proves
\eqref{eq:V47-connected-polynomial}.  For empty \(J\), all five
common-bank moment operators equal the identity, so \(K_J=C_{ij;k}=0\)
and \(M_J=I\), yielding the asserted V46 formula.
\end{proof}

\section{Common-bank output debits}
\label{sec:V47-common-debits}

Write \(\Xi_J:=\Xi(\boldsymbol T_J)\).  All constants below are uniform
in the size of \(J\), the representations, the output channel, and its
multiplicity.

\begin{lemma}[First moments of the common and cut operators]
\label{lem:V47-common-first-moments}
For sufficiently large \(\alpha\),
\begin{align}
 s_\alpha\Xi_J\|M_J\|_{\rm op}&\le C,
 \label{eq:V47-M-first-moment}\\
 s_\alpha\Xi_J\|C_{ij;k}\|_{\rm op}&\le C,
 \label{eq:V47-C-first-moment}\\
 s_\alpha\Xi_J\|K_J\|_{\rm op}&\le C.
 \label{eq:V47-K-first-moment}
\end{align}
The same statements hold with value zero when the relevant final
channel is trivial and the weighted expression is interpreted in the
usual way.
\end{lemma}

\begin{proof}
The first estimate is the product first-moment bound
\eqref{eq:V46-product-first-moment}.

For the second, diagonalize \(Q_{ij,J}^{\boldsymbol T_J}\) in the
\((ij)k\) bracketing.  On a pair-intermediate channel
\(\boldsymbol S\), the corresponding eigenvalue of \(C_{ij;k}\) is
\[
 q_{\alpha,\boldsymbol T_J}
 -q_{\alpha,\boldsymbol S}
  q_{\alpha,\boldsymbol R_{k,J}}.
\]
It is exactly the binary defect for
\(\boldsymbol T_J\subset\boldsymbol S\otimes
\boldsymbol R_{k,J}\).  Equation
\eqref{eq:V46-defect-first-moment}, followed by the spectral supremum,
proves \eqref{eq:V47-C-first-moment}; no multiplicity count is
introduced.

For the last estimate, use the five terms in
\eqref{eq:V47-common-cumulant}.  The \(M_J\) term was just treated.  In
the spectrum of \(B_{ij;k}\), fusion gives
\[
 \Xi_J\le C\bigl(\Xi(\boldsymbol S)
             +\Xi(\boldsymbol R_{k,J})\bigr).
\]
Since its eigenvalue is
\(q_{\alpha,\boldsymbol S}q_{\alpha,\boldsymbol R_{k,J}}\), two uses
of \eqref{eq:V46-product-first-moment} bound its weighted supremum.
For \(G_J\), use
\(\Xi_J\le C\sum_i\Xi(\boldsymbol R_{i,J})\) and pay one singleton
multiplier at a time.  The triangle inequality is applied to the five
operators only after these uniform operator-norm estimates, proving
\eqref{eq:V47-K-first-moment}.
\end{proof}

\begin{corollary}[Resolved binary-cut quotient]
\label{cor:V47-cut-quotient}
For a nontrivial common-bank output \(\boldsymbol T_J\) and a
pair-intermediate channel \(\boldsymbol S\) in the
\((ij)k\) bracketing, put
\[
 H_{ij;k}(\boldsymbol S)
 :=\sum_{e\in J}A_{S,e}A_{k,e}.
\]
Then
\begin{equation}
 \Xi_J\,
 \frac{\|C_{ij;k}\|_{\rm op}}
      {1-q_{\alpha,\boldsymbol T_J}}
 \le C\sup_{\boldsymbol S}H_{ij;k}(\boldsymbol S)
 \le C\sum_{e\in J}(A_{i,e}+A_{j,e})A_{k,e}.
 \label{eq:V47-cut-quotient}
\end{equation}
The first supremum ranges only over intermediate channels leading to
\(\boldsymbol T_J\).
\end{corollary}

\begin{proof}
In the diagonal bracketing, apply the binary quotient
\eqref{eq:V46-binary-quotient} to each eigenvalue from the preceding
proof and take the supremum.  Coordinatewise fusion geometry gives
\(A_{S,e}\le C(A_{i,e}+A_{j,e})\), proving the second inequality.
Unitary conjugation back to the common multiplicity space preserves
the operator norm.
\end{proof}

\section{Closure of the high exclusive-defect sector}
\label{sec:V47-high-defect}

Let \(H_{ij}:=\sum_{e\in L_{ij}}A_{i,e}A_{j,e}\).  For
\(\{i,j,k\}=\{1,2,3\}\), put
\begin{equation}
 \Xi_i^\circ:=
 \Xi(\boldsymbol R_{i,L_{ij}})
 +\Xi(\boldsymbol R_{i,L_{ik}})
 +\Xi(\boldsymbol R_{i,P_i}),
 \qquad
 Z_i:=\Xi(\boldsymbol R_{i,J}),
 \qquad
 \Xi_i=\Xi_i^\circ+Z_i.
 \label{eq:V47-center-mass-split}
\end{equation}
Define \(K_{\ge2}\) to be the sum of the last four terms in
\eqref{eq:V47-connected-polynomial} and
\(K_{\le1}:=K_{\rm full}-K_{\ge2}\).

\begin{lemma}[The common bank plus an unused pair is admissible]
\label{lem:V47-combined-spectator}
Fix distinct \(i,j,k\).  On the disjoint union
\(J\sqcup L_{jk}\), let the resolved spectator output be
\((\boldsymbol T_J,\boldsymbol S_{jk})\).  Each of
\begin{equation}
 r^{(0)}_{jk}:=q_{\alpha,\boldsymbol T_J}a_{jk},
 \qquad
 r^{(1)}_{jk}:=q_{\alpha,\boldsymbol T_J}\delta_{jk}
 \label{eq:V47-two-spectators}
\end{equation}
satisfies
\begin{equation}
 |r^{(\ell)}_{jk}|\le1,
 \qquad
 s_\alpha\bigl(\Xi_J+\Xi(\boldsymbol S_{jk})\bigr)
 |r^{(\ell)}_{jk}|\le C,
 \qquad \ell=0,1.
 \label{eq:V47-combined-admissibility}
\end{equation}
\end{lemma}

\begin{proof}
All Wilson multipliers lie in \([0,1]\), and
\(|\delta_{jk}|\le1\).  For the \(\Xi_J\) portion use
\eqref{eq:V47-M-first-moment}, with the other factor bounded by one.
For the \(L_{jk}\) portion use respectively
\Cref{lem:V46-unused-pair-admissible}; multiplication by
\(q_{\alpha,\boldsymbol T_J}\le1\) does not enlarge the bound.  Add
the two estimates.
\end{proof}

For later use, set
\begin{align}
 q_{\alpha,\boldsymbol T}
 &:=q_Pq_{\alpha,\boldsymbol T_J}
       \beta_{12}\beta_{13}\beta_{23},
 \label{eq:V47-full-output-multiplier}\\
 \Xi(\boldsymbol T)
 &:=\Xi_P+\Xi_J+
       \sum_{1\le a<b\le3}\Xi(\boldsymbol S_{ab}),
 \label{eq:V47-full-output-mass}\\
 \mathcal Q_{ij,ik}^{(\ell)}
 &:=\frac{\Xi(\boldsymbol T)q_P
       |\delta_{ij}\delta_{ik}r_{jk}^{(\ell)}|}
       {1-q_{\alpha,\boldsymbol T}},
 \qquad \ell=0,1.
 \label{eq:V47-high-defect-quotients}
\end{align}
Thus \(\ell=0\) is the two-defect monomial and \(\ell=1\) is the
three-defect monomial with center \(i\).

\begin{theorem}[Reduced-center payment for the high-defect sector]
\label{thm:V47-high-defect-reduced}
For every resolved output and \(\ell\in\{0,1\}\), the high-defect
quotient centered at input \(i\) satisfies
\begin{equation}
 \mathcal Q_{ij,ik}^{(\ell)}
 \le C\frac{H_{ij}H_{ik}}{\Xi_i^\circ}.
 \label{eq:V47-high-defect-reduced}
\end{equation}
For \(\ell=1\), one may choose any two nonempty exclusive-pair banks
forming a spanning tree.  The
estimate upgrades to the desired full-center payment
\begin{equation}
 \mathcal Q_{ij,ik}^{(\ell)}
 \le C(1+C_0)\frac{H_{ij}H_{ik}}{\Xi_i}
 \label{eq:V47-high-defect-compatible}
\end{equation}
whenever \(Z_i\le C_0\Xi_i^\circ\).  Without such a comparison, V47
does not assert the full-center payment.
\end{theorem}

\begin{proof}
Consider first
\(q_P\delta_{ij}\delta_{ik}a_{jk}M_J\).  Since
\(M_J=q_{\alpha,\boldsymbol T_J}I\), this is precisely the V46
two-defect quotient with selected banks \(L_{ij},L_{ik}\), private
factor \(q_P\), center mass \(\Xi_i^\circ\), and auxiliary row factor
\(r^{(0)}_{jk}\).  Its
admissibility is \eqref{eq:V47-combined-admissibility}; hence
\Cref{thm:V46-central-marked-payment} gives
\eqref{eq:V47-high-defect-reduced}.  The other two two-defect terms
are permutations.

For
\(q_P\delta_{12}\delta_{13}\delta_{23}M_J\), select any two nonempty
pair banks and use \(r^{(1)}_{jk}\) for the third.  If fewer than two
are nonempty, the monomial vanishes.  This proves the reduced-center
scalar resolved payments.  If \(Z_i\le C_0\Xi_i^\circ\), then
\(\Xi_i\le(1+C_0)\Xi_i^\circ\), which proves
\eqref{eq:V47-high-defect-compatible}.

At coefficient level, retain the scalar multiplier on the resolved
disjoint banks and apply
\eqref{eq:V46-exact-multiplier-normalization}.  The output factor
\(d_{\boldsymbol T}\) cancels the exact
\(d_{\rm in}/d_{\boldsymbol T}\) coefficient normalization.  Raw
isotypic pinching, as in
\eqref{eq:V46-disjoint-raw-sum}--\eqref{eq:V46-disjoint-pinching}, is
subprobability and introduces no path multiplicity.  Finally
\eqref{eq:V32-tree-mass-factorization} converts the compatible
full-center payments to the marked BFA seminorms.  The same argument
gives a coefficient-uniform reduced-center bound in general.  It does
not replace \(\Xi_i^\circ\) by \(\Xi_i\) when \(Z_i\) dominates; that
is the additional gate isolated below.
\end{proof}

\section{The exact zero/one-defect frontier}
\label{sec:V47-frontier}

Algebraically, the preceding theorem leaves exactly
\begin{align}
 K_{\le1}=q_P\bigl[{}
 &a_{12}a_{13}a_{23}K_J
 +\delta_{12}a_{13}a_{23}C_{12;3}
 +a_{12}\delta_{13}a_{23}C_{13;2}
 \notag\\
 &+a_{12}a_{13}\delta_{23}C_{23;1}\bigr].
 \label{eq:V47-low-defect-operator}
\end{align}
This identity identifies two genuinely different low-defect mechanisms:
the zero-defect term is a common-junction operator with arbitrary
separated pair spectators, while each one-defect term couples one
exclusive binary defect to a noncommuting common-bank cut operator.
Analytically there is also the full-center upgrade left open by
\Cref{thm:V47-high-defect-reduced} when \(Z_i\) dominates
\(\Xi_i^\circ\).

For a global support triple define
\begin{equation}
 \mathsf H_{ij}:=\sum_{e\in X_i\cap X_j}A_{i,e}A_{j,e},
 \qquad
 \mathsf W_3:=\sum_{i=1}^3
 \frac{\mathsf H_{ij}\mathsf H_{ik}}{\Xi_i},
 \quad \{i,j,k\}=\{1,2,3\}.
 \label{eq:V47-global-tree-weight}
\end{equation}
The terms with a zero denominator are omitted, equivalently interpreted
by the exact-support convention.

\begin{proposition}[A sufficient pair of remaining global gates]
\label{prop:V47-remaining-gates}
Let \(B_i\) be arbitrary trace-class Peter--Weyl input coefficients,
let \(X_{\boldsymbol T}\) be their raw resolved product coefficient on
the final isotypic multiplicity space, and let
\(K_{\le1}^{\boldsymbol T}\) be
\eqref{eq:V47-low-defect-operator} on that space.  The global \(m=3\)
marked-tree atlas follows if both of the following gates are proved.
The low-defect operator gate is
\begin{align}
 \sum_{\boldsymbol T\ne\boldsymbol1}
 &\Xi(\boldsymbol T)d_{\boldsymbol T}
 \left\|
 \frac{d_{\rm in}}{d_{\boldsymbol T}}
 \frac{\operatorname{Tr}_{M_{\boldsymbol T}}
  \left[X_{\boldsymbol T}
   (I\otimes K_{\le1}^{\boldsymbol T})\right]}
 {1-q_{\alpha,\boldsymbol T}}
 \right\|_1
 \notag\\
 &\hspace{7em}\le
 C d_{\rm in}\,\mathsf W_3\prod_{i=1}^3\|B_i\|_1.
 \label{eq:V47-low-defect-gate}
\end{align}
No absolute value may be inserted entrywise inside
\(K_{\le1}^{\boldsymbol T}\).  The high-defect center-upgrade gate is,
for every oriented two-defect term and an allowed center for the
three-defect term,
\begin{equation}
 \boxed{
 \mathcal Q_{ij,ik}^{(\ell)}
 \le C\frac{\mathsf H_{ij}\mathsf H_{ik}}{\Xi_i}}
 \qquad\text{on the region }Z_i>\Xi_i^\circ,
 \quad \ell=0,1.
 \label{eq:V47-high-defect-upgrade-gate}
\end{equation}
\end{proposition}

\begin{proof}
The exact polynomial
\eqref{eq:V47-connected-polynomial} splits the cumulant as
\(K_{\le1}+K_{\ge2}\).  Equation
\eqref{eq:V47-low-defect-gate} pays the first summand.  On
\(Z_i\le\Xi_i^\circ\), the second summand has the full-center payment
by \eqref{eq:V47-high-defect-compatible}; on the complementary region
it is paid by \eqref{eq:V47-high-defect-upgrade-gate}.  Raw pinching
then sums the resolved estimates, and the triangle inequality combines
the two sectors.  The first displayed left side is exactly the BFA
output weight after applying
\Cref{thm:V46-exact-product-normalization}; hence there is no missing
dimension factor.  This proves sufficiency.  No converse separation is
claimed, because cancellation between \(K_{\le1}\) and
\(K_{\ge2}\) has not been excluded.
\end{proof}

\begin{assessment}[Why the first-moment bounds do not yet prove the gates]
\label{ass:V47-first-moment-insufficient}
Equations \eqref{eq:V47-C-first-moment} and
\eqref{eq:V47-K-first-moment} make the common bank an admissible
one-moment spectator.  They do not supply the two original-input marks
required by \(\mathsf W_3\).  In particular,
\eqref{eq:V47-cut-quotient} contains the composite intermediate weight
\(A_{S,e}\).  Replacing it by a new denominator
\(\Xi(\boldsymbol S)^{-1}\) would lose the marked-tree center and can
be circular.  The next step must instead allocate
\(A_{S,e}\le C(A_{i,e}+A_{j,e})\) directly to the two original input
centers before raw pinching.  V47 therefore records
\eqref{eq:V47-low-defect-gate} and
\eqref{eq:V47-high-defect-upgrade-gate} as open gates, not as a proved
global MTA theorem.  The second gate is not cosmetic: the V46 central
payment produces \(\Xi_i^\circ\), and an output first moment on
\(\boldsymbol T_J\) does not control the input mass \(Z_i\) when a
large common-bank tensor product resolves to a small, or trivial,
output.
\end{assessment}

\section{V47 ledger and next acquisition}
\label{sec:V47-ledger}

\begin{theorem}[Integrated V47 statement]
\label{thm:V47-integrated}
Preserving V1--V46, V47 proves:
\begin{enumerate}[label=\textup{(V47.\arabic*)},leftmargin=5.5em]
\item the exact eight-term common/pair-bank connected polynomial
      \eqref{eq:V47-connected-polynomial};
\item uniform output first moments for the common full moment, every
      binary cut, and the common ternary cumulant;
\item the resolved binary-cut quotient
      \eqref{eq:V47-cut-quotient};
\item admissibility of the combined common-bank/unused-pair spectator;
\item a coefficient-uniform reduced-center payment for every term with
      at least two exclusive-pair defects, with full payment under the
      explicit center-mass compatibility condition; and
\item reduction of the remaining work to the low-defect operator gate
      \eqref{eq:V47-low-defect-gate} and the center-heavy upgrade gate
      \eqref{eq:V47-high-defect-upgrade-gate}.
\end{enumerate}
It does not prove either remaining gate.
\end{theorem}

\begin{proof}
The six items are respectively
\Cref{thm:V47-connected-polynomial,lem:V47-common-first-moments,
cor:V47-cut-quotient,lem:V47-combined-spectator,
thm:V47-high-defect-reduced,prop:V47-remaining-gates}.  The final
sentence is \Cref{ass:V47-first-moment-insufficient}.
\end{proof}

\begingroup
\small
\begin{longtable}{>{\raggedright\arraybackslash}p{0.29\textwidth}
                  >{\raggedright\arraybackslash}p{0.15\textwidth}
                  >{\raggedright\arraybackslash}p{0.48\textwidth}}
\toprule
\textbf{V47 row}&\textbf{Status}&\textbf{Advance or obstruction}\\
\midrule
\endfirsthead
\toprule
\textbf{V47 row}&\textbf{Status}&\textbf{Advance or obstruction}\\
\midrule
\endhead
Exact mixed-incidence algebra
& \MGclosed
& The full multiplier is the eight-term operator polynomial
  \eqref{eq:V47-connected-polynomial}.\\[0.35em]
Common/cut first moments
& \MGclosed
& Product damping and binary defect damping are uniform on the full
  multiplicity space.\\[0.35em]
Two/three exclusive-defect sector
& \MGreduced
& V46 central payment applies with the common output and unused pair as
  a spectator, but initially pays only \(\Xi_i^\circ\).\\[0.35em]
Center-heavy high-defect upgrade
& \MGopen
& When \(Z_i>\Xi_i^\circ\), the desired full-center denominator is not
  supplied by common-output damping.\\[0.35em]
Zero-defect common-junction sector
& \MGopen
& The product \(a_{12}a_{13}a_{23}K_J\) still needs two globally
  normalized original-input marks.\\[0.35em]
One-defect junction-cut sector
& \MGopen
& A composite binary cut must be allocated to original input centers
  before trace norm and pinching.\\[0.35em]
Global ternary MTA
& \MGopen
& It follows after both gates
  \eqref{eq:V47-low-defect-gate} and
  \eqref{eq:V47-high-defect-upgrade-gate} are closed.\\[0.35em]
Quartic and all-order MTA
& \MGopen
& They remain suspended until both ternary gates close.\\[0.35em]
Full Yang--Mills mass gap
& \MGopen
& WUFC/CPM, nonlinear construction, continuum, OS positivity,
  nontriviality, and the spectral-gap passage remain downstream.\\
\bottomrule
\end{longtable}
\endgroup

\begin{openproblem}[V48 mandate: center-heavy and low-defect gates]
\label{op:V48}
\begin{enumerate}[label=\textup{(V48.\arabic*)},leftmargin=5.5em]
\item Resolve the center-heavy high-defect gate
      \eqref{eq:V47-high-defect-upgrade-gate}.  In particular, do not
      infer input common-bank mass from the final common output; test
      the trivial-output channel first and use coefficient formation
      only if a genuine Casimir debit is proved.
\item Prove the pure common-junction spectator term
      \(q_Pa_{12}a_{13}a_{23}K_J\) with the global tree weight
      \(\mathsf W_3\).
\item For each one-defect term, diagonalize only the composite binary
      cut, then allocate
      \(A_{S,e}\le C(A_{i,e}+A_{j,e})\) to the two original tree
      centers.  Do not divide by \(\Xi(\boldsymbol S)\).
\item Keep the common multiplicity operator intact before all trace
      norms and use raw pinching with the exact
      \(d_{\rm in}/d_{\boldsymbol T}\) normalization.
\item Restore global \(m=3\) MTA only after proving
      both \eqref{eq:V47-low-defect-gate} and
      \eqref{eq:V47-high-defect-upgrade-gate}; otherwise sharpen the
      gates and move upstream rather than recycling a first-moment
      estimate.
\item Keep every quartic claim suspended until that restoration.
\item Preserve only the newest YM file.  The next SM/NS audit remains
      V50.
\end{enumerate}
\end{openproblem}

\begin{firewall}[End-of-V47 claim boundary]
\label{fire:V47-claim-boundary}
V47 proves the exact mixed-incidence polynomial, common-bank and cut
first moments, and a coefficient-level reduced-center bound for the
sector with at least two exclusive-pair defects.  It does not prove the
full-center upgrade in the center-heavy regime, the zero-defect
common-junction spectator estimate, any of the three one-defect
junction-cut estimates, either remaining gate
\eqref{eq:V47-low-defect-gate}--%
\eqref{eq:V47-high-defect-upgrade-gate}, or global ternary MTA.
Consequently it
does not reactivate quartic or all-order MTA.  It also does not prove
WUFC, CPM, the nonlinear Perron packet, finite-edge margins, capacity,
protected-sector floors, local-algebra noncollapse, reflection
positivity, a nontrivial continuum OS state, or a positive
Yang--Mills spectral gap.  The full Yang--Mills mass gap is not proved.
\end{firewall}

\section*{Version V47 append-only record}
\addcontentsline{toc}{section}{Version V47 append-only record}
The complete V46 source was retained before this continuation.  Its
SHA--256 digest was
\begin{center}\scriptsize\ttfamily
1349a5d285324394b159dfdcf15e72cbfb3df73d4844273c59ff87118fa49e90.
\end{center}
No mathematical content of V46 was altered.  On the next iteration,
retain this full file, append before its unique active document-closing
command, increment the filename to \texttt{\_V48.tex}, and remove only
the superseded V47 file from this Yang--Mills note series.  The next
scheduled SM/NS audit is V50.

% =============================================================================
% END APPENDED VERSION V47 MATERIAL
% =============================================================================

% =============================================================================
% BEGIN APPENDED VERSION V48 MATERIAL
% =============================================================================

\chapter{V48: Common-overlap absorption and full high-defect closure}
\label{ch:V48}

\section{Purpose and correction of the V47 frontier}
\label{sec:V48-purpose}

V47 correctly proved the reduced-center estimate for every term with
at least two exclusive-pair defects, but left the passage from
\(\Xi_i^\circ\) to the full center mass \(\Xi_i\) as a separate gate.
That gate does not require a coefficient-formation law.  When the
common-bank mass \(Z_i\) is large, the same bank belongs to both global
overlaps incident to center \(i\).  Those two numerator contributions
absorb the full denominator by elementary algebra, while the V46
pointwise minimum estimate pays the multiplier.  This chapter proves
that observation and thereby closes the complete high-exclusive-defect
sector.  It then resolves each remaining one-defect common cut into a
diagonal family of two genuine binary defects, without an entrywise
Racah sum.

V48 is not a scheduled companion comparison.  The next read-only
Standard Model/Navier--Stokes audit remains V50.

\section{The common-overlap absorption lemma}
\label{sec:V48-absorption}

Retain the V47 partition and write
\begin{equation}
 H_{ij}^J:=\sum_{e\in J}A_{i,e}A_{j,e},
 \qquad
 \mathsf H_{ij}=H_{ij}+H_{ij}^J,
 \qquad
 \mathsf W_i:=\frac{\mathsf H_{ij}\mathsf H_{ik}}{\Xi_i}.
 \label{eq:V48-overlap-split}
\end{equation}
Here \(H_{ij}\) is the exclusive-bank weight from V47 and
\(\{i,j,k\}=\{1,2,3\}\).

\begin{lemma}[A dominant common center pays the minimum bound]
\label{lem:V48-common-absorption}
For every center \(i\),
\begin{equation}
 H_{ij}^J\ge Z_i,
 \qquad H_{ik}^J\ge Z_i.
 \label{eq:V48-common-dominates-center}
\end{equation}
If \(Z_i>\Xi_i^\circ\), then
\begin{equation}
 \boxed{
 \mathsf W_i
 \ge \frac12\min\{H_{ij},H_{ik}\}.}
 \label{eq:V48-minimum-absorption}
\end{equation}
The estimate is uniform in all representations and support sizes.
\end{lemma}

\begin{proof}
Exact activity on \(J\) gives \(A_{j,e},A_{k,e}\ge1\).  Therefore
\[
 H_{ij}^J\ge\sum_{e\in J}A_{i,e}=Z_i,
 \qquad
 H_{ik}^J\ge\sum_{e\in J}A_{i,e}=Z_i,
\]
which proves \eqref{eq:V48-common-dominates-center}.  If
\(Z_i>\Xi_i^\circ\), then \(\Xi_i<2Z_i\), and hence
\begin{align*}
 \mathsf W_i
 &\ge\frac{(H_{ij}+Z_i)(H_{ik}+Z_i)}{2Z_i}\\
 &=\frac{H_{ij}H_{ik}}{2Z_i}
   +\frac{H_{ij}+H_{ik}}2+\frac{Z_i}2
 \ge\frac12\min\{H_{ij},H_{ik}\}.
\end{align*}
No relation between a common input and the resolved common output was
used.
\end{proof}

\begin{theorem}[Full-center payment for every high-defect monomial]
\label{thm:V48-high-defect-full-center}
For the quotients \(\mathcal Q_{ij,ik}^{(\ell)}\) of
\eqref{eq:V47-high-defect-quotients}, \(\ell=0,1\),
\begin{equation}
 \boxed{
 \mathcal Q_{ij,ik}^{(\ell)}
 \le C\frac{\mathsf H_{ij}\mathsf H_{ik}}{\Xi_i}.}
 \label{eq:V48-high-defect-full-center}
\end{equation}
For the three-defect monomial, any center determined by two nonempty
exclusive-pair banks may be used.  The constant is uniform in the
supports, representations, resolved outputs, and multiplicities.
\end{theorem}

\begin{proof}
First suppose \(Z_i\le\Xi_i^\circ\).  Apply
\eqref{eq:V47-high-defect-compatible} with \(C_0=1\), then use
\(H_{ab}\le\mathsf H_{ab}\).  This gives
\eqref{eq:V48-high-defect-full-center}.

Now suppose \(Z_i>\Xi_i^\circ\).  The combined common/unused-pair
factor is admissible by
\eqref{eq:V47-combined-admissibility}.  Thus the pointwise minimum
part of \Cref{thm:V46-two-defect-pointwise} gives
\[
 \mathcal Q_{ij,ik}^{(\ell)}
 \le C\min\{H_{ij},H_{ik}\}.
\]
Apply \eqref{eq:V48-minimum-absorption}.  This proves the result in the
complementary regime and exhausts all cases.  The argument for
\(\ell=1\) is identical after choosing the two selected pair banks;
if fewer than two are nonempty, that monomial is zero.
\end{proof}

\begin{theorem}[Coefficient-level closure of the full high-defect sector]
\label{thm:V48-high-defect-MTA}
The complete \(K_{\ge2}\) contribution in
\eqref{eq:V47-connected-polynomial} satisfies the global ternary
marked-tree atlas inequality for arbitrary trace-class Fourier
coefficients, uniformly in volume, support, representation height,
and multiplicity.
\end{theorem}

\begin{proof}
Apply \eqref{eq:V48-high-defect-full-center} to each of the three
two-defect monomials and to the three-defect monomial.  For the latter,
choose any two nonempty banks; this assigns it to one of the three
tree centers.  Use the exact coefficient factor
\(d_{\rm in}/d_{\boldsymbol T}\) from
\eqref{eq:V46-exact-multiplier-normalization}.  Multiplication by the
BFA output dimension cancels \(d_{\boldsymbol T}\), and raw isotypic
pinching gives the subprobability
\eqref{eq:V46-genuine-subprobability}.  Consequently the resolved
estimates sum without a representation-dimension or multiplicity
factor.  Finally \eqref{eq:V32-tree-mass-factorization} converts each
\(\mathsf W_i\) into the corresponding marked input seminorms.
\end{proof}

\begin{corollary}[The V47 high-defect upgrade gate is closed]
\label{cor:V48-high-gate-closed}
Equation \eqref{eq:V47-high-defect-upgrade-gate} is proved by
\Cref{thm:V48-high-defect-full-center}.  Therefore global ternary MTA
is equivalent, up to a universal change of constant, to the sole
low-defect operator gate \eqref{eq:V47-low-defect-gate}.
\end{corollary}

\begin{proof}
The exact identity is
\(K_{\rm full}=K_{\le1}+K_{\ge2}\).  The second term has the desired
bound by \Cref{thm:V48-high-defect-MTA}.  Hence the low-defect gate
implies global MTA by the triangle inequality.  Conversely, global MTA
and the proved high-defect bound imply the low-defect gate after
writing \(K_{\le1}=K_{\rm full}-K_{\ge2}\).  This proves the stated
equivalence; no separation of unproved terms is used.
\end{proof}

\section{Exact nested resolution of the one-defect cuts}
\label{sec:V48-nested-cuts}

Fix \(\{i,j,k\}=\{1,2,3\}\) and use the \((ij)k\) bracketing on the
common bank.  For every intermediate product channel
\(\boldsymbol S\subset
\boldsymbol R_{i,J}\otimes\boldsymbol R_{j,J}\) leading to
\(\boldsymbol T_J\), define
\begin{equation}
 \Delta_{ij;k}(\boldsymbol S,\boldsymbol T_J)
 :=q_{\alpha,\boldsymbol T_J}
   -q_{\alpha,\boldsymbol S}q_{\alpha,\boldsymbol R_{k,J}}.
 \label{eq:V48-nested-defect}
\end{equation}

\begin{theorem}[A common cut is a diagonal nested binary defect]
\label{thm:V48-nested-cut-resolution}
On the common multiplicity space,
\begin{equation}
 \boxed{
 C_{ij;k}
 =\sum_{\boldsymbol S}
 \Delta_{ij;k}(\boldsymbol S,\boldsymbol T_J)
 P_{\boldsymbol S\mid\boldsymbol T_J}^{ij}.}
 \label{eq:V48-nested-cut-resolution}
\end{equation}
Consequently the complete one-exclusive-defect operator is
\begin{align}
 K_1
 :=q_P\sum_{\substack{1\le i<j\le3\\
                        \{k\}=\{1,2,3\}\setminus\{i,j\}}}
 \delta_{ij}a_{ik}a_{jk}C_{ij;k}
 =q_P\sum_{\substack{1\le i<j\le3\\
                       \{k\}=\{1,2,3\}\setminus\{i,j\}}}
  \sum_{\boldsymbol S}
 &\delta_{ij}a_{ik}a_{jk}
 \Delta_{ij;k}(\boldsymbol S,\boldsymbol T_J)
 P_{\boldsymbol S\mid\boldsymbol T_J}^{ij},
 \label{eq:V48-one-defect-nested}
\end{align}
where each unordered pair \(ij\) occurs once and \(k\) is its
complement.  Thus every resolved summand contains two binary defects on
disjoint banks: \(L_{ij}\) and \(J\).
\end{theorem}

\begin{proof}
The spectral resolution
\eqref{eq:V33-pair-spectral-resolution} gives
\[
 Q_{ij,J}^{\boldsymbol T_J}
 =\sum_{\boldsymbol S}q_{\alpha,\boldsymbol S}
 P_{\boldsymbol S\mid\boldsymbol T_J}^{ij}.
\]
Insert this identity into
\(C_{ij;k}=q_{\alpha,\boldsymbol T_J}I
-q_{\alpha,\boldsymbol R_{k,J}}Q_{ij,J}^{\boldsymbol T_J}\).
The projections sum to the identity on the full multiplicity space,
which proves \eqref{eq:V48-nested-cut-resolution}.  Substitution into
the three one-defect terms of
\eqref{eq:V47-low-defect-operator} proves
\eqref{eq:V48-one-defect-nested}.
\end{proof}

For a resolved intermediate channel put
\begin{equation}
 G_{ij;k}(\boldsymbol S)
 :=\sum_{e\in J}A_{S,e}A_{k,e}.
 \label{eq:V48-nested-overlap}
\end{equation}

\begin{lemma}[Nested defect estimates and original-input allocation]
\label{lem:V48-nested-defect-estimates}
Uniformly in the intermediate channel,
\begin{align}
 |\Delta_{ij;k}(\boldsymbol S,\boldsymbol T_J)|
 &\le Cs_\alpha G_{ij;k}(\boldsymbol S),
 \label{eq:V48-nested-cap}\\
 \Xi_J
 \frac{|\Delta_{ij;k}(\boldsymbol S,\boldsymbol T_J)|}
      {1-q_{\alpha,\boldsymbol T_J}}
 &\le CG_{ij;k}(\boldsymbol S),
 \label{eq:V48-nested-quotient}\\
 s_\alpha\Xi_J
 |\Delta_{ij;k}(\boldsymbol S,\boldsymbol T_J)|
 &\le C,
 \label{eq:V48-nested-first-moment}\\
 G_{ij;k}(\boldsymbol S)
 &\le C\bigl(H_{ik}^J+H_{jk}^J\bigr).
 \label{eq:V48-nested-allocation}
\end{align}
\end{lemma}

\begin{proof}
The number in \eqref{eq:V48-nested-defect} is the binary defect for
\(\boldsymbol T_J\subset\boldsymbol S\otimes
\boldsymbol R_{k,J}\).  Apply respectively
\eqref{eq:V46-defect-cap},
\eqref{eq:V46-binary-quotient}, and
\eqref{eq:V46-defect-first-moment}.  Finally, coordinatewise fusion
gives \(A_{S,e}\le C(A_{i,e}+A_{j,e})\).  Multiply by \(A_{k,e}\)
and sum over \(J\) to obtain
\eqref{eq:V48-nested-allocation}.
\end{proof}

Let the fully resolved output \(\boldsymbol T\) be as in
\eqref{eq:V47-full-output-multiplier}--%
\eqref{eq:V47-full-output-mass}, and define
\begin{equation}
 \mathcal N_{ij;k}(\boldsymbol S)
 :=\frac{\Xi(\boldsymbol T)q_Pa_{ik}a_{jk}
 |\delta_{ij}\Delta_{ij;k}(\boldsymbol S,\boldsymbol T_J)|}
 {1-q_{\alpha,\boldsymbol T}}.
 \label{eq:V48-nested-quotient-full}
\end{equation}

\begin{theorem}[Unnormalized two-defect control of every nested cut]
\label{thm:V48-nested-two-defect}
Every resolved nested summand satisfies
\begin{align}
 \mathcal N_{ij;k}(\boldsymbol S)
 &\le C\min\{H_{ij},G_{ij;k}(\boldsymbol S)\},
 \label{eq:V48-nested-minimum}\\
 \mathcal N_{ij;k}(\boldsymbol S)
 &\le Cs_\alpha H_{ij}G_{ij;k}(\boldsymbol S).
 \label{eq:V48-nested-one-resolvent}
\end{align}
\end{theorem}

\begin{proof}
The two defects live on the disjoint banks \(L_{ij}\) and \(J\).
The remaining pair banks have row factor \(a_{ik}a_{jk}\).  If their
resolved outputs are \(\boldsymbol S_{ik},\boldsymbol S_{jk}\), then
\begin{align*}
 s_\alpha\bigl(
 \Xi(\boldsymbol S_{ik})+\Xi(\boldsymbol S_{jk})\bigr)
 |a_{ik}a_{jk}|
 &\le C
\end{align*}
by two applications of
\eqref{eq:V46-product-first-moment}; its modulus is at most one.  It is
therefore an admissible auxiliary spectator in the sense of
\eqref{eq:V46-admissible-spectator}.  Apply the two conclusions of
\Cref{thm:V46-two-defect-pointwise}, using
\eqref{eq:V48-nested-cap}--\eqref{eq:V48-nested-first-moment} for the
common defect.  This gives
\eqref{eq:V48-nested-minimum} and
\eqref{eq:V48-nested-one-resolvent}.
\end{proof}

\begin{lemma}[Nested pinching introduces no path count]
\label{lem:V48-nested-pinching}
Let \(X_{\boldsymbol T}\) be the raw full isotypic coefficient and put
\(\widetilde P_{\boldsymbol S}:=I_{H_{\boldsymbol T}}otimes
P_{\boldsymbol S\mid\boldsymbol T_J}^{ij}\), with identities on all
other multiplicity factors.  Then
\begin{align}
 \operatorname{Tr}_{M_{\boldsymbol T}}
 \left[X_{\boldsymbol T}(I\otimes C_{ij;k})\right]
 &=\sum_{\boldsymbol S}
 \Delta_{ij;k}(\boldsymbol S,\boldsymbol T_J)
 \operatorname{Tr}_{M_{\boldsymbol T}}
 \left[\widetilde P_{\boldsymbol S}X_{\boldsymbol T}
       \widetilde P_{\boldsymbol S}\right],
 \label{eq:V48-nested-partial-trace}\\
 \sum_{\boldsymbol S}
 \|\widetilde P_{\boldsymbol S}X_{\boldsymbol T}
       \widetilde P_{\boldsymbol S}\|_1
 &\le\|X_{\boldsymbol T}\|_1.
 \label{eq:V48-nested-pinching}
\end{align}
\end{lemma}

\begin{proof}
Insert \eqref{eq:V48-nested-cut-resolution}.  In a multiplicity basis
adapted to the mutually orthogonal projections, the partial trace of
\(X_{\boldsymbol T}\widetilde P_{\boldsymbol S}\) sees only diagonal
indices in the range of \(\widetilde P_{\boldsymbol S}\); hence it
equals the partial trace of
\(\widetilde P_{\boldsymbol S}X_{\boldsymbol T}
\widetilde P_{\boldsymbol S}\).  This proves
\eqref{eq:V48-nested-partial-trace}.  Equation
\eqref{eq:V48-nested-pinching} is trace-norm contractivity of the
orthogonal pinching map.
\end{proof}

\section{The now-sole low-defect frontier}
\label{sec:V48-low-frontier}

Put
\begin{equation}
 K_0:=q_Pa_{12}a_{13}a_{23}K_J,
 \qquad
 K_{\le1}=K_0+K_1.
 \label{eq:V48-low-split}
\end{equation}
By \Cref{cor:V48-high-gate-closed}, no high-defect gate remains.  A
sufficient pointwise target for each nested term is
\begin{equation}
 \boxed{
 \mathcal N_{ij;k}(\boldsymbol S)
 \le C\left(
 \frac{\mathsf H_{ij}\mathsf H_{ik}}{\Xi_i}
 +\frac{\mathsf H_{ij}\mathsf H_{jk}}{\Xi_j}
 \right).}
 \label{eq:V48-nested-marked-target}
\end{equation}
Together with the analogous BFA estimate for \(K_0\), this implies
\eqref{eq:V47-low-defect-gate}: use
\eqref{eq:V48-nested-partial-trace}--%
\eqref{eq:V48-nested-pinching}, sum the three pair choices, and apply
the exact Fourier normalization.  The missing issue is therefore no
longer multiplicity or representation formation.  It is the
two-center allocation of the common composite weight in
\eqref{eq:V48-nested-allocation}, together with the pure common
junction spectator estimate for \(K_0\).

\begin{firewall}[Status of the nested target]
\label{fire:V48-nested-target}
Equations \eqref{eq:V48-nested-minimum} and
\eqref{eq:V48-nested-one-resolvent} do not by themselves prove
\eqref{eq:V48-nested-marked-target}.  The denominators on its right are
the masses of the two original inputs, whereas the common binary
defect is indexed by the composite channel \(\boldsymbol S\).  V48
does not insert \(\Xi(\boldsymbol S)^{-1}\), because that would change
the marked tree and repeat the circular step excluded in V47.
\end{firewall}

\section{V48 ledger and next acquisition}
\label{sec:V48-ledger}

\begin{theorem}[Integrated V48 statement]
\label{thm:V48-integrated}
Preserving V1--V47, V48 proves:
\begin{enumerate}[label=\textup{(V48.\arabic*)},leftmargin=5.5em]
\item the common-overlap absorption inequality
      \eqref{eq:V48-minimum-absorption};
\item the full-center pointwise payment
      \eqref{eq:V48-high-defect-full-center} and coefficient-level
      global MTA closure of every high-exclusive-defect term;
\item closure of the V47 high-defect upgrade gate, leaving
      \eqref{eq:V47-low-defect-gate} as the sole global ternary gate;
\item the exact diagonal nested-defect identity
      \eqref{eq:V48-nested-cut-resolution} for each one-defect cut;
\item uniform cap, quotient, first-moment, and original-input allocation
      estimates for every nested defect;
\item the pointwise minimum and one-resolvent bounds
      \eqref{eq:V48-nested-minimum}--%
      \eqref{eq:V48-nested-one-resolvent}; and
\item the multiplicity-uniform nested pinching identity and the precise
      two-center target \eqref{eq:V48-nested-marked-target}.
\end{enumerate}
V48 does not prove the nested marked target or the pure common-junction
\(K_0\) estimate.
\end{theorem}

\begin{proof}
The seven items are respectively
\Cref{lem:V48-common-absorption,
thm:V48-high-defect-full-center,thm:V48-high-defect-MTA,
cor:V48-high-gate-closed,thm:V48-nested-cut-resolution,
lem:V48-nested-defect-estimates,thm:V48-nested-two-defect,
lem:V48-nested-pinching}.  The two remaining statements are recorded
in \Cref{sec:V48-low-frontier,fire:V48-nested-target}.
\end{proof}

\begingroup
\small
\begin{longtable}{>{\raggedright\arraybackslash}p{0.29\textwidth}
                  >{\raggedright\arraybackslash}p{0.15\textwidth}
                  >{\raggedright\arraybackslash}p{0.48\textwidth}}
\toprule
\textbf{V48 row}&\textbf{Status}&\textbf{Advance or obstruction}\\
\midrule
\endfirsthead
\toprule
\textbf{V48 row}&\textbf{Status}&\textbf{Advance or obstruction}\\
\midrule
\endhead
Common-overlap absorption
& \MGclosed
& A common-heavy center contributes to both incident global overlap
  numerators and absorbs the full center denominator.\\[0.35em]
High-exclusive-defect sector
& \MGclosed
& All four high-defect monomials satisfy full-center MTA for arbitrary
  Fourier coefficients.\\[0.35em]
One-defect common cut
& \MGreduced
& It is exactly a pinched family of two disjoint binary defects with no
  multiplicity count.\\[0.35em]
Nested two-center normalization
& \MGopen
& The composite common overlap must be allocated to centers \(i\) and
  \(j\) with their original full masses.\\[0.35em]
Zero-defect common junction
& \MGopen
& The operator \(K_0=q_Pa_{12}a_{13}a_{23}K_J\) still needs the single
  final resolvent and two global marks.\\[0.35em]
Global ternary MTA
& \MGopen
& It is equivalent to the sole low-defect gate
  \eqref{eq:V47-low-defect-gate}.\\[0.35em]
Quartic and all-order MTA
& \MGopen
& They remain suspended until global ternary MTA is restored.\\[0.35em]
Full Yang--Mills mass gap
& \MGopen
& WUFC/CPM, nonlinear construction, continuum, OS positivity,
  nontriviality, and the spectral-gap passage remain downstream.\\
\bottomrule
\end{longtable}
\endgroup

\begin{openproblem}[V49 mandate: the sole low-defect gate]
\label{op:V49}
\begin{enumerate}[label=\textup{(V49.\arabic*)},leftmargin=5.5em]
\item Prove \eqref{eq:V48-nested-marked-target}.  Split
      \eqref{eq:V48-nested-allocation} into the \(ik\) and \(jk\)
      orientations, use common-overlap absorption when the appropriate
      original center is common-heavy, and use quadratic damping of
      off-tree/private spectator mass in the complementary branch.
\item Do not introduce a denominator involving
      \(\Xi(\boldsymbol S)\).  The denominators must remain
      \(\Xi_i\) and \(\Xi_j\).
\item Prove the pure common-junction spectator estimate for \(K_0\),
      keeping \(K_J\) intact before trace norm.  Combine the V32
      replicated-coordinate operator numerator with the one final
      resolvent, common-overlap absorption, and retained spectator
      damping.
\item Use \eqref{eq:V48-nested-pinching} and the exact
      \(d_{\rm in}/d_{\boldsymbol T}\) normalization; do not sum Racah
      entries or insert a representation formation law.
\item Restore global \(m=3\) MTA only after the complete low-defect gate
      is proved.  Keep quartic claims suspended otherwise.
\item Preserve only the newest YM file.  The next SM/NS audit remains
      V50.
\end{enumerate}
\end{openproblem}

\begin{firewall}[End-of-V48 claim boundary]
\label{fire:V48-claim-boundary}
V48 closes the full high-exclusive-defect sector, including the
center-heavy regime, for arbitrary trace-class coefficients.  It also
resolves every one-defect common cut into a multiplicity-uniform family
of two binary defects and proves their unnormalized minimum and
one-resolvent estimates.  It does not prove the two-center normalized
target \eqref{eq:V48-nested-marked-target}, the pure common-junction
\(K_0\) estimate, the low-defect gate
\eqref{eq:V47-low-defect-gate}, or global ternary MTA.  Consequently it
does not reactivate quartic or all-order MTA.  It also does not prove
WUFC, CPM, the nonlinear Perron packet, finite-edge margins, capacity,
protected-sector floors, local-algebra noncollapse, reflection
positivity, a nontrivial continuum OS state, or a positive
Yang--Mills spectral gap.  The full Yang--Mills mass gap is not proved.
\end{firewall}

\section*{Version V48 append-only record}
\addcontentsline{toc}{section}{Version V48 append-only record}
The complete V47 source was retained before this continuation.  Its
SHA--256 digest was
\begin{center}\scriptsize\ttfamily
820fbe09339aa5a3571f5db69c130a127207538a8a1f6c68729b06a101af50d8.
\end{center}
No mathematical content of V47 was altered.  On the next iteration,
retain this full file, append before its unique active document-closing
command, increment the filename to \texttt{\_V49.tex}, and remove only
the superseded V48 file from this Yang--Mills note series.  The next
scheduled SM/NS audit is V50.

% =============================================================================
% END APPENDED VERSION V48 MATERIAL
% =============================================================================

% =============================================================================
% BEGIN APPENDED VERSION V49 MATERIAL
% =============================================================================

\chapter{V49: Two-center normalization and closure of the one-defect sector}
\label{ch:V49}

\section{Purpose and nonduplication boundary}
\label{sec:V49-purpose}

V48 reduced every one-exclusive-defect term to two scalar binary
defects on the disjoint banks \(L_{ij}\) and \(J\), followed by an
orthogonal multiplicity pinching.  The unresolved issue was not the
defect quotient itself, but assigning the composite common-bank weight
to the two original centers \(i\) and \(j\).  V49 proves that
assignment.  A common-heavy original center is handled by the V48
overlap absorption.  If both centers are common-light, each gives an
independent marked payment using the separated off-tree multiplier;
taking the better of the two payments and then using the fusion
allocation closes the exact two-center target.

No representation-formation tail is used.  In particular, a small
separated input multiplier is not mistaken for a small final
multiplier.  The hot off-tree proof splits low and high final outputs
and retains the single final resolvent throughout.  V49 is not a
five-version companion audit; that comparison is due at V50.

\section{A common-light one-bank center payment}
\label{sec:V49-common-light}

Fix \(\{i,j,k\}=\{1,2,3\}\), and let \(r\) be either \(i\) or \(j\).
Write \(t\) for the other member of the pair \(\{i,j\}\).  Define
\begin{align}
 x_r&:=\Xi(\boldsymbol R_{r,L_{ij}}),
 \label{eq:V49-selected-center-mass}\\
 p_r&:=\Xi(\boldsymbol R_{r,L_{rk}})
       +\Xi(\boldsymbol R_{r,P_r}),
 \label{eq:V49-offtree-center-mass}\\
 b_r&:=q_{\alpha,\boldsymbol R_{r,L_{rk}}}
       q_{\alpha,\boldsymbol R_{r,P_r}}.
 \label{eq:V49-offtree-center-multiplier}
\end{align}
Empty factors equal one.  Thus
\(\Xi_r^\circ=x_r+p_r\), and the row factor
\(q_Pa_{ik}a_{jk}\) in \eqref{eq:V48-nested-quotient-full} contains
\(b_r\).

\begin{lemma}[Quadratically damped off-tree output mass]
\label{lem:V49-offtree-output-damping}
Suppose \(s_\alpha p_r>M\), where \(M\) is a fixed sufficiently large
constant.  Let \(q_{\rm sp}:=q_Pa_{ik}a_{jk}\), and sum the resolved
output masses on the pair/private banks other than \(L_{ij}\).  Then
\begin{align}
 b_r&\le\frac C{(s_\alpha p_r)^2},
 \label{eq:V49-offtree-quadratic}\\
 \sum_{B\in\mathcal O}
 \Xi(\boldsymbol S_B)|q_{\rm sp}|
 &\le\frac C{s_\alpha^2p_r}.
 \label{eq:V49-offtree-output-mass}
\end{align}
Here \(\mathcal O\) consists of \(L_{ik},L_{jk}\) and the private
banks, with an empty output omitted.  Constants are independent of all
support sizes and representations.
\end{lemma}

\begin{proof}
Equation \eqref{eq:V49-offtree-quadratic} is the product second moment
\eqref{eq:V32-balanced-second-moment} applied to the exact off-tree
channel of input \(r\).

For a bank containing input \(r\), fusion bounds its resolved output
mass by a constant times the sum of the input-
\(r\) mass and the masses of the other operands.  After multiplication
by \(q_{\rm sp}\), the first part is bounded by
\[
 p_rb_r\le\frac C{s_\alpha^2p_r}.
\]
For every other input part, retain \(b_r\) and use the support-uniform
product first moment on that input.  Its contribution is at most
\[
 \frac C{s_\alpha}b_r
 \le\frac C{s_\alpha^3p_r^2}
 \le\frac {C_M}{s_\alpha^2p_r},
\]
because \(s_\alpha p_r>M\).  A bank not containing input \(r\) is
handled by the same second estimate.  There are only the fixed bank
types listed in \(\mathcal O\), while each product first moment is
uniform in the number of coordinates.  Summing proves
\eqref{eq:V49-offtree-output-mass}.
\end{proof}

\begin{theorem}[Common-light marked payment from either original center]
\label{thm:V49-common-light-payment}
If \(Z_r\le\Xi_r^\circ\), then every nested quotient satisfies
\begin{equation}
 \boxed{
 \mathcal N_{ij;k}(\boldsymbol S)
 \le C\frac{H_{ij}G_{ij;k}(\boldsymbol S)}{\Xi_r}.}
 \label{eq:V49-common-light-payment}
\end{equation}
The assertion holds separately for \(r=i\) and \(r=j\).
\end{theorem}

\begin{proof}
Abbreviate \(h=H_{ij}\), \(g=G_{ij;k}(\boldsymbol S)\),
\(x=x_r\), and \(p=p_r\).  Exact activity gives \(x\le h\), and the
hypothesis gives
\begin{equation}
 \Xi_r=x+p+Z_r\le2(x+p).
 \label{eq:V49-common-light-total}
\end{equation}

If \(p\le x\), then \(\Xi_r\le4x\le4h\).  The minimum estimate
\eqref{eq:V48-nested-minimum} gives
\[
 \mathcal N_{ij;k}(\boldsymbol S)
 \le Cg
 \le4C\frac{hg}{\Xi_r}.
\]
Suppose next that \(p>x\).  Then
\(p>\Xi_r/4\).  If \(s_\alpha p\le M\), the one-resolvent estimate
\eqref{eq:V48-nested-one-resolvent} gives
\[
 \mathcal N_{ij;k}(\boldsymbol S)
 \le Cs_\alpha hg
 \le\frac{CM}{p}hg
 \le\frac{4CM}{\Xi_r}hg.
\]

It remains to assume \(s_\alpha p>M\).  Split the final channel into
the low-output region \(s_\alpha C_E(\boldsymbol T)\le1\) and its
complement.  On the low-output region,
\eqref{eq:V46-low-output-resolvent},
\eqref{eq:V49-offtree-quadratic}, and the two defect caps give
\begin{align*}
 \mathcal N_{ij;k}(\boldsymbol S)
 &\le\frac C{s_\alpha}
       b_r(Cs_\alpha h)(Cs_\alpha g)\\
 &\le\frac C{s_\alpha p^2}hg
 \le\frac C{p}hg.
\end{align*}

On the high-output region, the lower product gap gives
\(1-q_{\alpha,\boldsymbol T}\ge c\).  Split
\(\Xi(\boldsymbol T)\) into the selected output on \(L_{ij}\), the
common output on \(J\), and the off-tree outputs \(\mathcal O\).  The
binary quotients imply
\[
 \Xi(\boldsymbol S_{ij})|\delta_{ij}|\le Ch,
 \qquad
 \Xi_J|\Delta_{ij;k}|\le Cg.
\]
After multiplying by the other defect cap, these two pieces are at
most \(Cs_\alpha b_rhg\le C hg/p\); here the first-moment consequence
\(s_\alpha p_rb_r\le C\) is enough.  For the remaining pieces use
\eqref{eq:V49-offtree-output-mass}:
\[
 \sum_{B\in\mathcal O}\Xi(\boldsymbol S_B)|q_{\rm sp}|
 |\delta_{ij}\Delta_{ij;k}|
 \le\frac C{s_\alpha^2p}(Cs_\alpha h)(Cs_\alpha g)
 \le\frac C{p}hg.
\]
The denominator floor and \(p>\Xi_r/4\) finish the high-output case.
Together with the preceding cases this proves
\eqref{eq:V49-common-light-payment}.
\end{proof}

\section{The two-center allocation theorem}
\label{sec:V49-two-center}

\begin{lemma}[Common-heavy absorption for a nested term]
\label{lem:V49-nested-common-heavy}
If \(Z_r>\Xi_r^\circ\) for either \(r=i\) or \(r=j\), then
\begin{equation}
 \mathcal N_{ij;k}(\boldsymbol S)
 \le C\mathsf W_r.
 \label{eq:V49-nested-common-heavy}
\end{equation}
\end{lemma}

\begin{proof}
Assume first \(r=i\).  Exact activity on \(J\) gives
\(H_{ij}^J\ge Z_i\) and \(H_{ik}^J\ge Z_i\).  Since
\(\Xi_i<2Z_i\),
\[
 \mathsf W_i
 \ge\frac{(H_{ij}+Z_i)Z_i}{2Z_i}
 \ge\frac12H_{ij}.
\]
Equation \eqref{eq:V48-nested-minimum} is at most \(CH_{ij}\), proving
the claim.  The case \(r=j\) uses
\(H_{ij}^J,H_{jk}^J\ge Z_j\) and is identical.
\end{proof}

\begin{theorem}[The V48 nested marked target]
\label{thm:V49-nested-marked-target}
For every resolved intermediate channel,
\begin{equation}
 \boxed{
 \mathcal N_{ij;k}(\boldsymbol S)
 \le C\left(
 \frac{\mathsf H_{ij}\mathsf H_{ik}}{\Xi_i}
 +\frac{\mathsf H_{ij}\mathsf H_{jk}}{\Xi_j}
 \right).}
 \label{eq:V49-nested-marked-target}
\end{equation}
Thus the open target \eqref{eq:V48-nested-marked-target} is proved.
\end{theorem}

\begin{proof}
If either \(i\) or \(j\) is common-heavy, apply
\Cref{lem:V49-nested-common-heavy}; the corresponding \(\mathsf W_r\)
is one of the two terms on the right of
\eqref{eq:V49-nested-marked-target}.

Suppose both centers are common-light.  Apply
\Cref{thm:V49-common-light-payment} first with \(r=i\), then with
\(r=j\), and take the better estimate:
\[
 \mathcal N_{ij;k}(\boldsymbol S)
 \le ChG_{ij;k}(\boldsymbol S)
       \min\{\Xi_i^{-1},\Xi_j^{-1}\}.
\]
By \eqref{eq:V48-nested-allocation},
\[
 G_{ij;k}(\boldsymbol S)
 \le C(H_{ik}^J+H_{jk}^J).
\]
Since each reciprocal is at least
\(\min\{\Xi_i^{-1},\Xi_j^{-1}\}\),
\begin{align*}
 h(H_{ik}^J+H_{jk}^J)
 \min\{\Xi_i^{-1},\Xi_j^{-1}\}
 &\le
 \frac{hH_{ik}^J}{\Xi_i}
 +\frac{hH_{jk}^J}{\Xi_j}\\
 &\le
 \frac{\mathsf H_{ij}\mathsf H_{ik}}{\Xi_i}
 +\frac{\mathsf H_{ij}\mathsf H_{jk}}{\Xi_j}.
\end{align*}
This proves the remaining case.
\end{proof}

\section{Coefficient closure of the one-defect sector}
\label{sec:V49-one-defect-coefficients}

\begin{theorem}[Global MTA bound for all one-exclusive-defect terms]
\label{thm:V49-one-defect-MTA}
For arbitrary trace-class Peter--Weyl coefficient matrices, the full
operator \(K_1\) of \eqref{eq:V48-one-defect-nested} satisfies the
global ternary marked-tree atlas inequality, uniformly in volume,
supports, representation heights, final outputs, and multiplicities.
\end{theorem}

\begin{proof}
Fix one pair \(ij\) and use its \((ij)k\) bracketing.  Insert the exact
partial-trace resolution
\eqref{eq:V48-nested-partial-trace}.  Apply
\eqref{eq:V49-nested-marked-target} to each intermediate channel.
The diagonal raw blocks sum contractively by
\eqref{eq:V48-nested-pinching}; hence no number of intermediate paths
appears.  Next sum the final isotypic blocks using the genuine raw
subprobability \eqref{eq:V46-genuine-subprobability}.  The exact factor
\(d_{\rm in}/d_{\boldsymbol T}\) in
\eqref{eq:V46-exact-multiplier-normalization} is cancelled by the BFA
output weight, exactly as in V46.  Finally
\eqref{eq:V32-tree-mass-factorization} converts the two terms on the
right of \eqref{eq:V49-nested-marked-target} into the marked input
seminorms centered at \(i\) and \(j\).  Sum the three choices of
\(ij\).  Racah changes only conjugate the chosen diagonal resolution
and therefore do not alter any trace norm.
\end{proof}

\begin{corollary}[The one-defect part of the V47 low gate is closed]
\label{cor:V49-one-defect-closed}
The \(K_1\) contribution to \eqref{eq:V47-low-defect-gate} obeys its
required bound.  Together with \Cref{thm:V48-high-defect-MTA}, this
leaves only the zero-exclusive-defect operator
\begin{equation}
 K_0=q_Pa_{12}a_{13}a_{23}K_J.
 \label{eq:V49-sole-K0}
\end{equation}
\end{corollary}

\begin{proof}
The first statement is \Cref{thm:V49-one-defect-MTA}.  The exact
connected polynomial \eqref{eq:V47-connected-polynomial} is the sum of
\(K_0\), \(K_1\), and \(K_{\ge2}\).  The last two parts now have the
global MTA bound, proving the second statement.
\end{proof}

\section{The sole ternary gate after V49}
\label{sec:V49-sole-gate}

For arbitrary input coefficient matrices \(B_i\), let
\(X_{\boldsymbol T}\) be the raw resolved product coefficient and let
\(K_0^{\boldsymbol T}\) denote
\eqref{eq:V49-sole-K0} on its full multiplicity space.  The remaining
common-junction spectator gate is
\begin{align}
 \sum_{\boldsymbol T\ne\boldsymbol1}
 &\Xi(\boldsymbol T)d_{\boldsymbol T}
 \left\|
 \frac{d_{\rm in}}{d_{\boldsymbol T}}
 \frac{\operatorname{Tr}_{M_{\boldsymbol T}}
  \left[X_{\boldsymbol T}(I\otimes K_0^{\boldsymbol T})\right]}
 {1-q_{\alpha,\boldsymbol T}}
 \right\|_1
 \notag\\
 &\hspace{7em}\le
 C d_{\rm in}\,\mathsf W_3\prod_{i=1}^3\|B_i\|_1.
 \label{eq:V49-common-junction-gate}
\end{align}

\begin{theorem}[Exact post-V49 equivalence]
\label{thm:V49-sole-gate-equivalence}
Global \(m=3\) marked-tree atlas holds if and only if
\eqref{eq:V49-common-junction-gate} holds, up to changing the universal
constant.
\end{theorem}

\begin{proof}
If \eqref{eq:V49-common-junction-gate} holds, combine it with
\Cref{thm:V49-one-defect-MTA,thm:V48-high-defect-MTA} and the exact
decomposition \(K_{\rm full}=K_0+K_1+K_{\ge2}\).  This proves global
MTA.  Conversely, assume global MTA.  Subtract the already bounded
\(K_1\) and \(K_{\ge2}\) contributions and use the triangle inequality
to obtain \eqref{eq:V49-common-junction-gate}.  All three estimates use
the same exact Fourier normalization, so no dimension conversion is
hidden in this subtraction.
\end{proof}

\begin{assessment}[Structure of the sole remaining operator]
\label{ass:V49-K0-structure}
The common factor \(K_J\) is the complete ternary cumulant on the
product bank \(J\); it must not be split into five absolute values.
The separated factor \(q_Pa_{12}a_{13}a_{23}\) retains one singleton
multiplier for every off-common input coordinate and therefore carries
quadratic damping when an original input is dominated away from
\(J\).  On \(J\), the V32 replicated-coordinate expansion separates a
marked local junction from two pair bonds.  The remaining proof must
attach the one global final resolvent to that operator expansion and
normalize its two marks by the original full input masses.  First
moments of \(K_J\) alone are insufficient, but no exclusive-defect or
multiplicity obstruction remains.
\end{assessment}

\section{V49 ledger and next acquisition}
\label{sec:V49-ledger}

\begin{theorem}[Integrated V49 statement]
\label{thm:V49-integrated}
Preserving V1--V48, V49 proves:
\begin{enumerate}[label=\textup{(V49.\arabic*)},leftmargin=5.5em]
\item the support-uniform quadratic damping
      \eqref{eq:V49-offtree-output-mass} for a hot off-tree center,
      without identifying input and final multipliers;
\item the common-light marked payment
      \eqref{eq:V49-common-light-payment} from either original center;
\item common-heavy absorption for every nested term;
\item the complete two-center estimate
      \eqref{eq:V49-nested-marked-target}, closing the V48 target;
\item multiplicity-uniform coefficient closure of all three
      one-exclusive-defect terms; and
\item reduction of global ternary MTA to the sole pure common-junction
      spectator gate \eqref{eq:V49-common-junction-gate}.
\end{enumerate}
V49 does not prove that last gate.
\end{theorem}

\begin{proof}
The six items are respectively
\Cref{lem:V49-offtree-output-damping,
thm:V49-common-light-payment,lem:V49-nested-common-heavy,
thm:V49-nested-marked-target,thm:V49-one-defect-MTA,
thm:V49-sole-gate-equivalence}.  The claim boundary is stated in
\Cref{ass:V49-K0-structure} and the firewall below.
\end{proof}

\begingroup
\small
\begin{longtable}{>{\raggedright\arraybackslash}p{0.29\textwidth}
                  >{\raggedright\arraybackslash}p{0.15\textwidth}
                  >{\raggedright\arraybackslash}p{0.48\textwidth}}
\toprule
\textbf{V49 row}&\textbf{Status}&\textbf{Advance or obstruction}\\
\midrule
\endfirsthead
\toprule
\textbf{V49 row}&\textbf{Status}&\textbf{Advance or obstruction}\\
\midrule
\endhead
Hot off-tree center
& \MGclosed
& Low final output uses the universal resolvent; high output uses a gap
  floor and bankwise first/second multiplier moments.\\[0.35em]
Nested two-center normalization
& \MGclosed
& Common-heavy absorption and two common-light payments prove
  \eqref{eq:V49-nested-marked-target}.\\[0.35em]
One-exclusive-defect sector
& \MGclosed
& Nested pinching lifts the pointwise payment with no path,
  representation-dimension, or Racah-entry count.\\[0.35em]
High-exclusive-defect sector
& \MGclosed
& Retained from V48; all terms containing any exclusive defect now
  satisfy global ternary MTA.\\[0.35em]
Pure common-junction spectator
& \MGopen
& The intact operator \(K_J\), its separated spectators, and the one
  final resolvent must be assembled in
  \eqref{eq:V49-common-junction-gate}.\\[0.35em]
Global ternary MTA
& \MGreduced
& It is equivalent to the sole common-junction gate.\\[0.35em]
Quartic and all-order MTA
& \MGopen
& They remain suspended until global ternary MTA is restored.\\[0.35em]
Full Yang--Mills mass gap
& \MGopen
& WUFC/CPM, nonlinear construction, continuum, OS positivity,
  nontriviality, and the spectral-gap passage remain downstream.\\
\bottomrule
\end{longtable}
\endgroup

\begin{openproblem}[V50 mandate: companion audit and the pure common junction]
\label{op:V50}
\begin{enumerate}[label=\textup{(V50.\arabic*)},leftmargin=5.5em]
\item Before selecting the proof direction, perform the scheduled
      read-only audit of the newest Standard Model and Navier--Stokes
      files in \texttt{latex\_notes}.  Transfer only type-correct
      structural ideas and record an explicit firewall.
\item Derive an operator-valued replicated-coordinate expansion of
      \(K_J\) in which each surviving connected monomial contains
      either one local three-input junction or two pair bonds of
      distinct types.  Do not split the five global cumulant terms
      before this cancellation.
\item Attach the single final resolvent to each marked connected
      monomial.  Use the corrected exact Fourier normalization and raw
      pinching throughout.
\item When a center is dominated off \(J\), retain the separated factor
      \(q_Pa_{12}a_{13}a_{23}\) and use quadratic spectator damping.
      When it is common-heavy, use the two incident common overlaps to
      absorb its full denominator.
\item Prove \eqref{eq:V49-common-junction-gate} before declaring global
      \(m=3\) MTA restored.  If the operator expansion does not yield
      two original-input marks, isolate the exact failed inequality and
      move upstream.
\item Keep every quartic claim suspended until the preceding item is
      complete.  Preserve only the newest YM file.
\end{enumerate}
\end{openproblem}

\begin{firewall}[End-of-V49 claim boundary]
\label{fire:V49-claim-boundary}
V49 proves the two-center normalized target and closes the complete
one-exclusive-defect sector for arbitrary trace-class coefficients.
Together with V48, every term containing at least one exclusive-pair
defect now has the global ternary MTA bound.  V49 does not prove the
pure common-junction spectator gate
\eqref{eq:V49-common-junction-gate}, global ternary MTA, quartic or
all-order MTA, WUFC, CPM, the nonlinear Perron packet, finite-edge
margins, capacity, protected-sector floors, local-algebra noncollapse,
reflection positivity, a nontrivial continuum OS state, or a positive
Yang--Mills spectral gap.  The full Yang--Mills mass gap is not proved.
\end{firewall}

\section*{Version V49 append-only record}
\addcontentsline{toc}{section}{Version V49 append-only record}
The complete V48 source was retained before this continuation.  Its
SHA--256 digest was
\begin{center}\scriptsize\ttfamily
99eebd421a8c00ce7b958e84c14b3c2b116004411224b3ffc2402fb93324dd01.
\end{center}
No mathematical content of V48 was altered.  On the next iteration,
retain this full file, append before its unique active document-closing
command, increment the filename to \texttt{\_V50.tex}, and remove only
the superseded V49 file from this Yang--Mills note series.  V50 must
begin with the scheduled SM/NS audit.

% =============================================================================
% END APPENDED VERSION V49 MATERIAL
% =============================================================================

% =============================================================================
% BEGIN APPENDED VERSION V50 MATERIAL
% =============================================================================

\chapter{V50: Companion audit and closure of global ternary MTA}
\label{ch:V50}

\section{Scheduled SM/NS companion audit}
\label{sec:V50-companion-audit}

The required read-only audit was performed before selecting the V50
proof.  The newest companion files in \texttt{latex\_notes} were
\begin{center}
\begin{tabular}{lll}
\toprule
programme&file&SHA--256\\
\midrule
SM--Einstein&\texttt{Renewal\_SM\_Einstein\_Continuum\_Research\_Notes\_V66.tex}
&\texttt{decd71390935d4f40844d8be4d924dbb4668957fd93f5df6ab267a1e4e0569e2}\\
Navier--Stokes&\texttt{Renewal\_NS\_Stability\_Research\_Notes\_V31.tex}
&\texttt{f1121b62238ad5491bb7cf03635ea9934942edf573ed8faaa9ee79e1d38a6696}\\
\bottomrule
\end{tabular}
\end{center}

SM V66 repairs an incomplete triangular boundary cascade by completing
the full spin system before asking for positivity, then fixes the
forward and adjoint domains through an exact Green form.  Its relevant
structural lesson is that a stable multiplier identity is not a closed
operator theorem until the missing rows and the correct inverse/domain
have been assembled.  NS V31 proves an exact multiplicity-free marked
formation ledger but also shows, by forced-mode and genuine-packet
tests, that unmarking after the decisive correlation has been discarded
returns the continuation-class quantity and is circular.  Its relevant
lesson is to retain the connected source, the stopping correlation, and
the final inverse before taking absolute values.

For the YM gate these lessons prescribe the same safe operation: expand
the complete common-bank cumulant into connected local operators,
retain the separated spectator product and the one final resolvent, and
only then pinch.  No SM boundary-semigroup theorem and no NS
stopping-time estimate has the type of a compact-group Fourier
multiplier estimate, so neither is imported.

\begin{firewall}[Strict companion-transfer boundary]
\label{fire:V50-companion-transfer}
The audit transfers only the assembly rule just stated.  SM V66 does
not imply positivity, a domain theorem, or a spectral floor for the YM
Markov row.  NS V31 does not imply a Casimir tail, a Clebsch--Gordan
estimate, or a YM cumulant bound.  Conversely, the result below supplies
no Einstein boundary port, Navier--Stokes record separation, or theorem
in either companion programme.  The next compulsory companion audit is
V55.
\end{firewall}

\section{Operator-valued connected expansion on the common bank}
\label{sec:V50-operator-expansion}

At a common coordinate \(e\in J\), abbreviate the singleton
multipliers by \(q_{i,e}\), the pair moment operator by \(Q_{ij,e}\),
and the one-link ternary cumulant by \(K_{123,e}\).  Define the local
mean, pair bond, and junction operators
\begin{align}
 m_e&:=q_{1,e}q_{2,e}q_{3,e}I,
 \label{eq:V50-local-mean}\\
 d_{ij,e}&:=q_{k,e}\bigl(Q_{ij,e}-q_{i,e}q_{j,e}I\bigr),
 \qquad \{i,j,k\}=\{1,2,3\},
 \label{eq:V50-local-pair-bond}\\
 t_e&:=K_{123,e}.
 \label{eq:V50-local-junction}
\end{align}
The exact local joint moment is
\begin{equation}
 Q_{123,e}=m_e+d_{12,e}+d_{13,e}+d_{23,e}+t_e.
 \label{eq:V50-local-moment-split}
\end{equation}

Let \(\mathfrak C_J\) be the set of maps which assign to every
\(e\in J\) one of the five symbols
\(0,12,13,23,123\), and whose nonzero hyperedges connect the vertex
set \(\{1,2,3\}\).  To \(\omega\in\mathfrak C_J\) associate the tensor
product of \(m_e,d_{ij,e},t_e\) selected by its symbols.

\begin{theorem}[Exact connected-coordinate operator formula]
\label{thm:V50-connected-operator-formula}
Before choosing any common-bank bracketing,
\begin{equation}
 \boxed{
 \mathbf K_J
 =\sum_{\omega\in\mathfrak C_J}
 \bigotimes_{e\in J}L_{e,\omega(e)},}
 \label{eq:V50-connected-operator-formula}
\end{equation}
where \(L_{e,0}=m_e\),
\(L_{e,ij}=d_{ij,e}\), and \(L_{e,123}=t_e\).
Every surviving term contains a local junction or pair bonds of at
least two distinct types.
\end{theorem}

\begin{proof}
Equation \eqref{eq:V50-local-moment-split} follows by substituting the
five-term definition of \(K_{123,e}\); its scalar mean occurs with
coefficient \(1-3+2=0\) in the difference.  Tensor
\eqref{eq:V50-local-moment-split} over \(J\) in the all-in-one moment
and use the corresponding local pair/mean expansions in the other four
global cumulant partitions.  For any fixed coordinate assignment, its
coefficient is the Möbius sum over global partitions which contain all
of its local hyperedges.  That sum is one if their union connects all
three vertices and zero otherwise.  This is the three-vertex product
cumulant formula, and proves
\eqref{eq:V50-connected-operator-formula}.  On three vertices,
connectivity is equivalent to the final sentence.
\end{proof}

Define the common overlap and junction weights
\begin{align}
 H_{ij}^J&:=\sum_{e\in J}A_{i,e}A_{j,e},
 \label{eq:V50-common-pair-weight}\\
 J_{123}^J&:=\sum_{e\in J}A_{1,e}A_{2,e}A_{3,e},
 \label{eq:V50-common-junction-weight}\\
 M_i^J&:=H_{ij}^JH_{ik}^J,
 \qquad \{i,j,k\}=\{1,2,3\}.
 \label{eq:V50-oriented-common-tree}
\end{align}

\begin{lemma}[A junction is dominated by every common tree]
\label{lem:V50-junction-to-tree}
For each center \(i\),
\begin{equation}
 J_{123}^J\le M_i^J.
 \label{eq:V50-junction-to-tree}
\end{equation}
\end{lemma}

\begin{proof}
The product \(H_{ij}^JH_{ik}^J\) contains its diagonal terms
\(A_{i,e}^2A_{j,e}A_{k,e}\).  Since \(A_{i,e}\ge1\), their sum
dominates \(J_{123}^J\).
\end{proof}

\begin{lemma}[Operator-valued product covariance]
\label{lem:V50-operator-covariance}
Let \(H,K\) be bounded operator-valued functions of finitely many
independent coordinates.  If \(H^{(e)},K^{(e)}\) denote independent
replacement of coordinate \(e\), then
\begin{equation}
 \|\mathbb E(HK)-\mathbb EH\,\mathbb EK\|_{\rm op}
 \le\frac12\sum_e
 \left\|\mathbb E[(H-H^{(e)})(H-H^{(e)})^*]\right\|_{\rm op}^{1/2}
 \left\|\mathbb E[(K-K^{(e)})^*(K-K^{(e)})]\right\|_{\rm op}^{1/2}.
 \label{eq:V50-operator-covariance}
\end{equation}
The same estimate holds after arbitrary spectator amplification.
\end{lemma}

\begin{proof}
Order the coordinates, let \(\mathcal F_e\) be the filtration which
reveals the first \(e\) of them, and put
\[
 D_eH:=\mathbb E(H\mid\mathcal F_e)
       -\mathbb E(H\mid\mathcal F_{e-1}),
 \qquad
 D_eK:=\mathbb E(K\mid\mathcal F_e)
       -\mathbb E(K\mid\mathcal F_{e-1}).
\]
The tower property, with the order of the two operator factors left
unchanged, gives
\begin{equation}
 \mathbb E(HK)-\mathbb EH\,\mathbb EK
 =\sum_e\mathbb E(D_eH\,D_eK).
 \label{eq:V50-operator-martingale-covariance}
\end{equation}
Indeed, after expanding \(H-\mathbb EH=\sum_eD_eH\), conditioning \(K\)
onto \(\mathcal F_e\) removes its future part and conditioning once more
onto \(\mathcal F_{e-1}\) removes the preceding part.

For random operators \(A,B\), positivity of the \(2\times2\) Gram
matrix gives the operator Cauchy--Schwarz inequality
\[
 \|\mathbb E(AB)\|_{\rm op}
 \le
 \|\mathbb E(AA^*)\|_{\rm op}^{1/2}
 \|\mathbb E(B^*B)\|_{\rm op}^{1/2}.
\]
Let \(\mathbb E_e\) integrate only coordinate \(e\).  Conditional
Jensen, applied after the unrevealed coordinates have been integrated,
gives the Loewner inequalities
\begin{align*}
 \mathbb E(D_eH\,D_eH^*)
 &\preccurlyeq
 \mathbb E[(H-\mathbb E_eH)(H-\mathbb E_eH)^*],\\
 \mathbb E(D_eK^*D_eK)
 &\preccurlyeq
 \mathbb E[(K-\mathbb E_eK)^*(K-\mathbb E_eK)].
\end{align*}
Conditional on all coordinates other than \(e\), independent
replacement gives the exact identities
\begin{align*}
 \mathbb E[(H-\mathbb E_eH)(H-\mathbb E_eH)^*]
 &=\frac12\mathbb E[(H-H^{(e)})(H-H^{(e)})^*],\\
 \mathbb E[(K-\mathbb E_eK)^*(K-\mathbb E_eK)]
 &=\frac12\mathbb E[(K-K^{(e)})^*(K-K^{(e)})].
\end{align*}
Apply operator Cauchy--Schwarz termwise in
\eqref{eq:V50-operator-martingale-covariance} and sum.  Tensoring every
operator and conditional expectation with an identity preserves the
two Loewner inequalities and all operator norms, proving spectator
amplification.
\end{proof}

\begin{theorem}[Complete common-bank operator two-mark numerator]
\label{thm:V50-operator-numerator}
Uniformly in \(J\), all representations, and \(\alpha\),
\begin{equation}
 \boxed{
 \|\mathbf K_J\|_{\rm op}
 \le Cs_\alpha^2\left(
 J_{123}^J+M_1^J+M_2^J+M_3^J\right)
 \le Cs_\alpha^2\sum_{i=1}^3M_i^J.}
 \label{eq:V50-operator-numerator}
\end{equation}
The same bound holds for every final common-bank compression
\(K_J^{\boldsymbol T_J}\), with no multiplicity factor.
\end{theorem}

\begin{proof}
Let \(F_i=\bigotimes_{e\in J}D^{R_{i,e}}(U_e)\), acting on the \(i\)-th
input carrier, and define the two-sided operator influence
\begin{align}
 \gamma_{i,e}:=\max\biggl\{&
 \left\|\frac12\mathbb E[
  (F_i-F_i^{(e)})(F_i-F_i^{(e)})^*]\right\|_{\rm op}^{1/2},
 \notag\\[-0.2em]
 &\left\|\frac12\mathbb E[
  (F_i-F_i^{(e)})^*(F_i-F_i^{(e)})]\right\|_{\rm op}^{1/2}
 \biggr\}.
 \label{eq:V50-two-sided-influence}
\end{align}
Set
\[
 W_{ij}^{\rm op}:=\sum_{e\in J}\gamma_{i,e}\gamma_{j,e},
 \qquad
 \overline W_{ij}^{\rm op}:=\min\{W_{ij}^{\rm op},1\},
 \qquad
 V^{\rm op}:=\sum_{e\in J}\|t_e\|_{\rm op}.
\]
The one-coordinate instance of
\Cref{lem:V50-operator-covariance}, together with
\(|q_{k,e}|\le1\), gives
\begin{equation}
 \|d_{ij,e}\|_{\rm op}
 \le\gamma_{i,e}\gamma_{j,e}.
 \label{eq:V50-pair-bond-by-influence}
\end{equation}

We first prove the operator replicated-tree estimate
\begin{equation}
 \|\mathbf K_J\|_{\rm op}
 \le C\left(
 V^{\rm op}
 +\overline W_{12}^{\rm op}\overline W_{13}^{\rm op}
 +\overline W_{12}^{\rm op}\overline W_{23}^{\rm op}
 +\overline W_{13}^{\rm op}\overline W_{23}^{\rm op}\right).
 \label{eq:V50-operator-replicated-tree}
\end{equation}
Take \(\varepsilon=1/4\).  If all
\(W_{ij}^{\rm op}\le\varepsilon\) and \(V^{\rm op}\le1\), take
operator norms in the exact connected expansion
\eqref{eq:V50-connected-operator-formula}.  Equations
\eqref{eq:V50-pair-bond-by-influence} and
\(\|m_e\|_{\rm op}\le1\) bound the unmarked remainder by
\[
 \prod_{e\in J}
 \left(1+\sum_{i<j}\|d_{ij,e}\|_{\rm op}+\|t_e\|_{\rm op}\right)
 \le
 \exp\left(\sum_{i<j}W_{ij}^{\rm op}+V^{\rm op}\right)
 \le e^{7/4}.
\]
Mark the first junction, or the first bonds of two distinct pair
types.  Summing the marked choices proves
\eqref{eq:V50-operator-replicated-tree} in this branch.

The five-term definition gives a universal operator cap on
\(\mathbf K_J\).  Consequently \(V^{\rm op}\ge1\), or two pair
influences at least \(\varepsilon\), is absorbed by the right side of
\eqref{eq:V50-operator-replicated-tree}.  In the sole remaining case,
say \(W_{12}^{\rm op}\ge\varepsilon\) and
\(W_{13}^{\rm op},W_{23}^{\rm op}<\varepsilon\), centrality makes each
singleton mean scalar, and the complete operator cumulant identity is
\[
 \mathbf K_J
 =\operatorname{Cov}(F_1F_2,F_3)
 -\mathbb EF_1\operatorname{Cov}(F_2,F_3)
 -\mathbb EF_2\operatorname{Cov}(F_1,F_3).
\]
Since the three \(F_i\) act on distinct tensor factors, the two-sided
influence of \(F_1F_2\) is at most
\(\gamma_{1,e}+\gamma_{2,e}\).  Apply
\Cref{lem:V50-operator-covariance} to the three covariances to obtain
\[
 \|\mathbf K_J\|_{\rm op}
 \le C(W_{13}^{\rm op}+W_{23}^{\rm op}).
\]
Multiplication by the lower bound
\(\overline W_{12}^{\rm op}\ge\varepsilon\) proves
\eqref{eq:V50-operator-replicated-tree}.  The other strong-pair choices
are permutations.

It remains to insert Wilson estimates.  The same geodesic derivative
and displacement-moment argument as
\eqref{eq:V32-Wilson-coordinate-influence}, now on the full carrier,
gives
\[
 \gamma_{i,e}
 \le C\min\{s_\alpha^{1/2}A_{i,e}^{1/2},1\}.
\]
Therefore
\[
 \overline W_{ij}^{\rm op}\le Cs_\alpha H_{ij}^J.
\]
The one-link operator estimate
\eqref{eq:V37-operator-one-link} gives
\[
 V^{\rm op}
 \le Cs_\alpha^2J_{123}^J.
\]
Substitution in \eqref{eq:V50-operator-replicated-tree} proves the
first inequality of \eqref{eq:V50-operator-numerator}; the second
follows from \Cref{lem:V50-junction-to-tree}.  Orthogonal isotypic
compression is contractive in operator norm, so the same estimate
holds on every final common-bank block without a multiplicity factor.
\end{proof}

\section{Separated spectators and the single-resolvent response}
\label{sec:V50-spectator-response}

For each input define its mass and multiplier away from the common bank:
\begin{equation}
 \Pi_i:=\Xi(\boldsymbol R_{i,X_i\setminus J})=\Xi_i-Z_i,
 \qquad
 b_i:=q_{\alpha,\boldsymbol R_{i,X_i\setminus J}}.
 \label{eq:V50-offcommon-data}
\end{equation}
The zero-defect spectator factor has the exact product form
\begin{equation}
 r_0:=q_Pa_{12}a_{13}a_{23}=b_1b_2b_3.
 \label{eq:V50-spectator-product}
\end{equation}
Let \(\Xi_{\rm sp}\) be the sum of all resolved output masses on the
exclusive-pair and private banks.

\begin{lemma}[First and second spectator debits]
\label{lem:V50-spectator-debits}
Uniformly over all resolved spectator outputs,
\begin{align}
 |r_0|&\le1,
 &
 s_\alpha\Xi_{\rm sp}|r_0|&\le C.
 \label{eq:V50-spectator-first-moment}
\end{align}
Let \(\Pi:=\max_i\Pi_i\), and choose \(i_*\) attaining the maximum.  If
\(s_\alpha\Pi>M\), then
\begin{align}
 b_{i_*}&\le\frac C{(s_\alpha\Pi)^2},
 \label{eq:V50-max-center-quadratic}\\
 \Xi_{\rm sp}|r_0|
 &\le\frac C{s_\alpha^2\Pi}.
 \label{eq:V50-spectator-second-output}
\end{align}
All constants are support-uniform.
\end{lemma}

\begin{proof}
The factorization \eqref{eq:V50-spectator-product} and positivity of
the Wilson multipliers give the first assertion.  On each
exclusive-pair bank, fusion bounds the output mass by the sum of its two
input masses.  Multiplying by the two corresponding singleton
multipliers and using
\eqref{eq:V46-product-first-moment} pays that output mass.  Private
outputs are singleton channels and obey the same estimate.  There are
only six noncommon bank types, so summing proves
\eqref{eq:V50-spectator-first-moment}.

Equation \eqref{eq:V50-max-center-quadratic} is
\eqref{eq:V32-balanced-second-moment} on the exact off-common channel
of input \(i_*\).  To prove
\eqref{eq:V50-spectator-second-output}, split each fused spectator
output mass into its input contributions.  The \(i_*\) contribution is
at most
\[
 \Pi b_{i_*}\le\frac C{s_\alpha^2\Pi}.
\]
For a contribution from another input, use its product first moment
while retaining \(b_{i_*}\); this gives
\[
 \frac C{s_\alpha}b_{i_*}
 \le\frac C{s_\alpha^3\Pi^2}
 \le\frac {C_M}{s_\alpha^2\Pi}.
\]
The same calculation handles private outputs.  Summation over the
fixed bank types proves the result.  This is the V49 hot-spectator
argument with all noncommon banks included; it does not compare an
input multiplier with a fused final multiplier.
\end{proof}

\begin{lemma}[Common-output fusion mass]
\label{lem:V50-common-output-mass}
Every resolved common output satisfies
\begin{equation}
 \Xi_J\le C(Z_1+Z_2+Z_3).
 \label{eq:V50-common-output-mass}
\end{equation}
\end{lemma}

\begin{proof}
At each \(e\in J\), the final representation occurs in
\(R_{1,e}\otimes R_{2,e}\otimes R_{3,e}\).  The quadratic Casimir
fusion bound and the support-count part of \(\Xi\) give
\(A_{T,e}\le C(A_{1,e}+A_{2,e}+A_{3,e})\).  Sum over \(J\).
\end{proof}

For a nontrivial full output define the operator response
\begin{equation}
 \mathcal Z(\boldsymbol T)
 :=\frac{\Xi(\boldsymbol T)|r_0|
       \|K_J^{\boldsymbol T_J}\|_{\rm op}}
       {1-q_{\alpha,\boldsymbol T}}.
 \label{eq:V50-K0-response}
\end{equation}

\begin{lemma}[Universal one-resolvent cap]
\label{lem:V50-universal-K0-response}
For every resolved output,
\begin{equation}
 \mathcal Z(\boldsymbol T)\le\frac C{s_\alpha}.
 \label{eq:V50-universal-K0-response}
\end{equation}
\end{lemma}

\begin{proof}
If \(s_\alpha C_E(\boldsymbol T)\le1\), use
\eqref{eq:V46-low-output-resolvent}, \(|r_0|\le1\), and the universal
five-term cumulant cap.  On the complementary output set, the lower
product gap bounds the denominator below.  Split
\(\Xi(\boldsymbol T)=\Xi_J+\Xi_{\rm sp}\).  The common part is at most
\(C/s_\alpha\) by \eqref{eq:V47-K-first-moment}.  The spectator part is
at most \(C/s_\alpha\) by
\eqref{eq:V50-spectator-first-moment} and the universal cumulant cap.
\end{proof}

\section{Pointwise closure of the pure common-junction gate}
\label{sec:V50-K0-pointwise}

\begin{theorem}[Pure common-junction marked response]
\label{thm:V50-K0-pointwise}
Uniformly over every resolved output,
\begin{equation}
 \boxed{
 \mathcal Z(\boldsymbol T)
 \le C\sum_{i=1}^3
 \frac{\mathsf H_{ij}\mathsf H_{ik}}{\Xi_i}
 =C\mathsf W_3.}
 \label{eq:V50-K0-pointwise}
\end{equation}
\end{theorem}

\begin{proof}
Fix a large numerical threshold \(M\).

\emph{Common-heavy hot branch.}
Suppose some \(i\) satisfies
\[
 Z_i>\Pi_i,\qquad s_\alpha Z_i>M.
\]
Then \(\Xi_i<2Z_i\), while exact common activity gives
\(\mathsf H_{ij},\mathsf H_{ik}\ge Z_i\).  Hence
\[
 \mathsf W_i\ge\frac{Z_i}{2}>\frac M{2s_\alpha}.
\]
The universal cap \eqref{eq:V50-universal-K0-response} proves
\eqref{eq:V50-K0-pointwise} on this branch.

Assume henceforth that no common-heavy hot center exists, and put
\(\Pi=\max_i\Pi_i\).

\emph{Globally thermal branch.}
Suppose \(s_\alpha\Pi\le M\).  If \(Z_i>\Pi_i\), the absence of the
first branch gives \(s_\alpha Z_i\le M\); if \(Z_i\le\Pi_i\), then
\(s_\alpha Z_i\le M\).  Thus
\begin{equation}
 s_\alpha\Xi_i\le2M
 \quad(1\le i\le3),
 \qquad
 \Xi_J\le\frac C{s_\alpha}
 \label{eq:V50-all-inputs-thermal}
\end{equation}
by \Cref{lem:V50-common-output-mass}.

On low final outputs, combine
\eqref{eq:V46-low-output-resolvent} with
\eqref{eq:V50-operator-numerator}:
\[
 \mathcal Z(\boldsymbol T)
 \le Cs_\alpha\sum_iM_i^J
 \le C_M\sum_i\frac{M_i^J}{\Xi_i}.
\]
On high final outputs, the denominator has a fixed floor.  The common
output part is bounded using the second assertion of
\eqref{eq:V50-all-inputs-thermal}; the spectator part uses
\eqref{eq:V50-spectator-first-moment}.  Multiplication by
\eqref{eq:V50-operator-numerator} gives the same
\(Cs_\alpha\sum_iM_i^J\) bound.  Since
\(\mathsf H_{ij}\ge H_{ij}^J\), this is bounded by
\(C_M\mathsf W_3\).

\emph{Off-common hot branch.}
It remains that \(s_\alpha\Pi>M\).  Choose \(i_*\) as in
\Cref{lem:V50-spectator-debits}.  The absence of a common-heavy hot
center implies, for every \(i\),
\[
 Z_i\le\max\{\Pi_i,M/s_\alpha\}\le C_M\Pi,
 \qquad
 \Xi_i\le C_M\Pi.
 \label{eq:V50-inputs-below-max-offcommon}
\]
On low final outputs,
\eqref{eq:V46-low-output-resolvent},
\eqref{eq:V50-max-center-quadratic}, and
\eqref{eq:V50-operator-numerator} give
\[
 \mathcal Z(\boldsymbol T)
 \le\frac C{s_\alpha}
 \frac C{s_\alpha^2\Pi^2}
 \left(Cs_\alpha^2\sum_iM_i^J\right)
 \le\frac C{\Pi}\sum_iM_i^J.
\]
On high final outputs, split the output mass.  For the common part,
\eqref{eq:V50-common-output-mass},
\eqref{eq:V50-inputs-below-max-offcommon}, and
\eqref{eq:V50-max-center-quadratic} give
\[
 \Xi_J|r_0|\le C\Pi b_{i_*}
 \le\frac C{s_\alpha^2\Pi}.
\]
For the spectator part use
\eqref{eq:V50-spectator-second-output}.  After multiplication by
\eqref{eq:V50-operator-numerator} and the denominator floor, both
pieces are at most \(C\sum_iM_i^J/\Pi\).  Finally
\eqref{eq:V50-inputs-below-max-offcommon} implies
\[
 \frac1\Pi\sum_iM_i^J
 \le C_M\sum_i\frac{M_i^J}{\Xi_i}
 \le C_M\mathsf W_3.
\]
The three branches are exhaustive.
\end{proof}

\section{Coefficient closure and restoration of global ternary MTA}
\label{sec:V50-global-ternary}

\begin{theorem}[The pure common-junction spectator gate]
\label{thm:V50-common-junction-gate}
The coefficient inequality
\eqref{eq:V49-common-junction-gate} holds for arbitrary trace-class
Peter--Weyl coefficients, uniformly in volume, supports,
representations, final channels, and multiplicities.
\end{theorem}

\begin{proof}
Insert the exact multiplier normalization
\eqref{eq:V46-exact-multiplier-normalization}.  For every raw final
block, apply \eqref{eq:V50-K0-pointwise} before taking the trace norm.
The BFA output factor cancels
\(d_{\boldsymbol T}\) against \(d_{\rm in}/d_{\boldsymbol T}\), and
the raw isotypic blocks sum by
\eqref{eq:V46-global-raw-pinching}.  The exponential height ratio is
at most one by tensor-product height subadditivity.  Finally
\eqref{eq:V32-tree-mass-factorization} converts each \(\mathsf W_i\)
to its two marked input seminorms.  No multiplicity path or
representation-formation weight is inserted.
\end{proof}

\begin{theorem}[Restored global ternary marked-tree atlas]
\label{thm:V50-global-ternary-MTA}
The full third row cumulant satisfies the global \(m=3\) marked-tree
atlas inequality for arbitrary exact supports, representations, and
trace-class Fourier coefficients, uniformly in volume and
\(\alpha\) on the established weak-coupling range.
\end{theorem}

\begin{proof}
The exact V47 connected polynomial decomposes the third cumulant into
\(K_0+K_1+K_{\ge2}\).  The three terms are bounded respectively by
\Cref{thm:V50-common-junction-gate,thm:V49-one-defect-MTA,
thm:V48-high-defect-MTA}.  Sum their marked-tree bounds.  All use the
same exact Fourier normalization and the same single final resolvent.
\end{proof}

\begin{firewall}[Meaning of the restored ternary theorem]
\label{fire:V50-ternary-scope}
\Cref{thm:V50-global-ternary-MTA} restores only the \(m=3\) atlas.
It does not automatically validate earlier quartic or all-order claims
whose proofs used the withdrawn arbitrary-coefficient formation law.
Those arguments must be re-entered from their exact connected
operator polynomials with the corrected Fourier normalization.
\end{firewall}

\section{V50 ledger and next acquisition}
\label{sec:V50-ledger}

\begin{theorem}[Integrated V50 statement]
\label{thm:V50-integrated}
Preserving V1--V49, V50 proves:
\begin{enumerate}[label=\textup{(V50.\arabic*)},leftmargin=5.5em]
\item the scheduled read-only SM/NS companion audit, with a strict
      type firewall and the next audit fixed at V55;
\item the exact connected-coordinate operator formula
      \eqref{eq:V50-connected-operator-formula} for the complete
      common-bank third cumulant;
\item the dimension-free operator covariance inequality
      \eqref{eq:V50-operator-covariance} and the full common-bank
      two-mark numerator \eqref{eq:V50-operator-numerator};
\item first- and second-moment control of the complete separated
      spectator product, including the enhanced hot-center output debit
      \eqref{eq:V50-spectator-second-output};
\item the pointwise one-resolvent response
      \eqref{eq:V50-K0-pointwise} in all common-heavy, globally thermal,
      and off-common hot branches;
\item the arbitrary-coefficient pure common-junction gate
      \eqref{eq:V49-common-junction-gate}; and
\item restoration of the complete global ternary marked-tree atlas
      with the exact Fourier normalization and raw pinching.
\end{enumerate}
\end{theorem}

\begin{proof}
Item \textup{(V50.1)} is
\Cref{sec:V50-companion-audit,fire:V50-companion-transfer}.
Items \textup{(V50.2)--(V50.3)} are
\Cref{thm:V50-connected-operator-formula,
lem:V50-operator-covariance,thm:V50-operator-numerator}.
Item \textup{(V50.4)} is
\Cref{lem:V50-spectator-debits,lem:V50-common-output-mass}.
Item \textup{(V50.5)} is
\Cref{lem:V50-universal-K0-response,thm:V50-K0-pointwise}.
Items \textup{(V50.6)--(V50.7)} are
\Cref{thm:V50-common-junction-gate,thm:V50-global-ternary-MTA}.
\end{proof}

\begingroup
\small
\begin{longtable}{>{\raggedright\arraybackslash}p{0.29\textwidth}
                  >{\raggedright\arraybackslash}p{0.15\textwidth}
                  >{\raggedright\arraybackslash}p{0.48\textwidth}}
\toprule
\textbf{V50 row}&\textbf{Status}&\textbf{Advance or obstruction}\\
\midrule
\endfirsthead
\toprule
\textbf{V50 row}&\textbf{Status}&\textbf{Advance or obstruction}\\
\midrule
\endhead
Scheduled SM/NS audit
& \MGclosed
& SM V66 contributes the full-operator-before-positivity discipline;
  NS V31 contributes the intact-source-before-unmarking discipline.
  No analytic theorem crosses the type firewall.\\[0.35em]
Common-bank connected operator
& \MGclosed
& The complete cumulant is an exact sum of local junctions and
  two-pair connected patterns before any absolute value is taken.\\[0.35em]
Operator two-mark numerator
& \MGclosed
& Martingale covariance, two-sided coordinate influences, and the
  connected cluster expansion give
  \(Cs_\alpha^2\sum_iM_i^J\) without dimension or multiplicity loss.\\[0.35em]
Separated spectator response
& \MGclosed
& Common-heavy absorption, a globally thermal split, and quadratic
  damping of the largest off-common input exhaust the response
  branches.\\[0.35em]
Pure common-junction coefficient gate
& \MGclosed
& Exact \(d_{\rm in}/d_{\boldsymbol T}\) normalization and raw pinching
  lift the pointwise response for arbitrary trace-class coefficients.\\[0.35em]
Global ternary MTA
& \MGclosed
& The zero-, one-, and high-exclusive-defect pieces are now all
  controlled by the same marked-tree atlas.\\[0.35em]
Quartic and all-order MTA
& \MGopen
& The four quartic skeletons must be revalidated from their exact
  operator polynomials; withdrawn formation-law arguments remain
  unavailable.\\[0.35em]
Full Yang--Mills mass gap
& \MGopen
& All-order MTA, WUFC/CPM, nonlinear construction, finite-edge
  estimates, continuum and OS positivity/nontriviality, and the
  spectral-gap passage remain downstream.\\
\bottomrule
\end{longtable}
\endgroup

\begin{openproblem}[V51 mandate: exact re-entry to the quartic row]
\label{op:V51}
\begin{enumerate}[label=\textup{(V51.\arabic*)},leftmargin=5.5em]
\item Make an explicit validity map for V40--V45: retain the V40
      quartic skeleton classification and the V41 local operator tree
      card, but continue to exclude every arbitrary-coefficient claim
      which used the withdrawn formation-law kernel.
\item Starting from
      \eqref{eq:V39-connected-replica-formula}, write the complete
      four-input connected operator as disjoint fibres of types
      \([4]\), \([3,2]\), \([3,3]\), and \([2,2,2]\), with all
      unselected incidence banks retained as spectator operators.
\item Revalidate the V41 fourfold coefficient compression using the
      exact \(d_{\rm in}/d_{\boldsymbol T}\) product formula and raw
      isotypic pinching.  Do not identify a representation-dimension
      law with arbitrary coefficient mass.
\item Attack the \([2,2,2]\) fibre first.  Keep its three binary defects,
      the complete spectator multiplier, and the single final resolvent
      assembled; derive the degree-normalized marks for each of the
      sixteen labeled trees without summing recoupling entries.
\item Treat \([3,2]\) and \([3,3]\) only with the restored ternary
      response kept as an intact operator block.  A termwise absolute
      decomposition of the ternary block is not admissible.
\item For the embedded \([4]\) fibre, combine each V41 local tree
      operator with the actual spectator product and the full output
      resolvent.  Split common-heavy, thermal, and spectator-hot centers
      before attempting repeated degree marks.
\item Declare global \(m=4\) MTA only after all four topology fibres
      satisfy the corrected arbitrary-coefficient bound.  If any fibre
      fails, record the first exact normalized inequality and go
      upstream rather than invoking a withdrawn formation tail.
\item Preserve only the newest YM file.  The next scheduled read-only
      SM/NS audit is V55, not V51.
\end{enumerate}
\end{openproblem}

\begin{firewall}[End-of-V50 claim boundary]
\label{fire:V50-claim-boundary}
V50 closes the pure common-junction spectator gate and thereby restores
the complete global \(m=3\) marked-tree atlas for arbitrary trace-class
Fourier coefficients.  It does not reactivate the V42--V45 quartic
spectator closures which relied on false coefficient-formation
probabilities, and it does not prove global \(m=4\), all-order MTA,
WUFC, CPM, the nonlinear Perron packet, unreduced finite-edge margins,
capacity, protected-sector floors, local-algebra noncollapse,
reflection positivity, a nontrivial continuum OS state, or a positive
Yang--Mills spectral gap.  The full Yang--Mills mass gap is not proved.
\end{firewall}

\section*{Version V50 append-only record}
\addcontentsline{toc}{section}{Version V50 append-only record}
The complete V49 source was retained before this continuation.  Its
SHA--256 digest was
\begin{center}\scriptsize\ttfamily
07d314c36be3c3f623a35d7aa0e4683ac974242e3c0d3d5efdc14b75ba86eabc.
\end{center}
No mathematical content of V49 was altered.  On the next iteration,
retain this full file, append before its unique active document-closing
command, increment the filename to \texttt{\_V51.tex}, and remove only
the superseded V50 file from this Yang--Mills note series.

% =============================================================================
% END APPENDED VERSION V50 MATERIAL
% =============================================================================

% =============================================================================
% BEGIN APPENDED VERSION V51 MATERIAL
% =============================================================================

\chapter{V51: Exact quartic re-entry and the coefficient-level star gate}
\label{ch:V51}

\section{Purpose and validity map}
\label{sec:V51-purpose}

V50 restored the complete ternary marked-tree atlas without using the
false coefficient-formation law.  The quartic row may therefore be
re-entered, but it cannot simply be cited from V42--V45.  This chapter
first separates the durable quartic inputs from the withdrawn
coefficient transfers.  It then corrects the fourfold compression,
derives the exact exclusive-pair quartic polynomial, closes its thermal
and genuine exterior-hot tree regimes, and tests the first remaining
edge-concentrated star.

\begin{theorem}[Post-V50 quartic validity map]
\label{thm:V51-validity-map}
The following earlier statements remain available:
\begin{enumerate}[label=\textup{(\roman*)},leftmargin=3.7em]
\item the incidence-form Leonov--Shiryaev identity
      \eqref{eq:V39-connected-replica-formula};
\item the four-skeleton classification
      \Cref{thm:V40-quartic-skeleton-list};
\item the local, representation-uniform V41 decomposition
      \eqref{eq:V41-exact-Wilson-tree-decomposition}--%
      \eqref{eq:V41-exact-Wilson-tree-bound};
\item the direct product cubic moment
      \eqref{eq:V42-balanced-third-moment}; and
\item the V46 exact Fourier normalization, raw pinching, binary defect
      cap, quotient, and first moment.
\end{enumerate}
No arbitrary-coefficient quartic conclusion from V42--V45 which uses
a conditioned representation-dimension law as coefficient mass is
available.
\end{theorem}

\begin{proof}
Items \textup{(i)--(ii)} are combinatorial identities.  Item
\textup{(iii)} is proved in operator norm before coefficient formation.
Item \textup{(iv)} follows directly from three integrations by parts
against the Wilson density and the product-moment expansion.  Item
\textup{(v)} was rederived after the V45 counterexample.  None of these
proofs inserts the withdrawn dimension-law kernel.  The last sentence
is precisely the retraction in V45 and the claim boundary in
\Cref{fire:V50-ternary-scope}.
\end{proof}

\section{Corrected fourfold coefficient compression}
\label{sec:V51-fourfold-compression}

Let \(B_i\in\mathcal S_1(H_{R_i})\),
\(\mathcal B=B_1\otimes\cdots\otimes B_4\), and
\(d_{\rm in}=\prod_i d_{R_i}\).  In the V46 isotypic decomposition put
\[
 X_T=W_T^*\mathcal BW_T.
\]
For a V41 local tree operator, let
\(\mathbf K_{\tau,T}\) be its restriction to \(M_T\).

\begin{theorem}[Exact-normalization fourfold tree compression]
\label{thm:V51-corrected-fourfold-compression}
For every set \(\Omega\) of final irreducibles,
\begin{equation}
 \boxed{
 \sum_{T\in\Omega}d_T
 \|\widehat\kappa_{\tau}(T)\|_1
 \le
 C_4d_{\rm in}s_\alpha^3
 \prod_{i=1}^4A_i^{\deg_\tau(i)/2}
 \prod_{i=1}^4\|B_i\|_1.}
 \label{eq:V51-corrected-fourfold-compression}
\end{equation}
There is no final-dimension or LR-multiplicity factor.
\end{theorem}

\begin{proof}
The exact coefficient, with the multiplier kept on the complete
multiplicity space, is
\[
 \widehat\kappa_{\tau}(T)
 =
 \frac{d_{\rm in}}{d_T}
 \operatorname{Tr}_{M_T}
 \left[X_T(I_{H_T}\otimes\mathbf K_{\tau,T})\right]
\]
by \eqref{eq:V46-exact-multiplier-normalization}.  Hence
\[
 d_T\|\widehat\kappa_{\tau}(T)\|_1
 \le
 d_{\rm in}\|\mathbf K_{\tau,T}\|_{\rm op}\|X_T\|_1.
\]
The V41 operator bound controls the first norm.  Sum the raw blocks and
use \eqref{eq:V46-global-raw-pinching} for the second.  This proves
\eqref{eq:V51-corrected-fourfold-compression}.  In particular, the
conclusion of \Cref{thm:V41-fourfold-compression} survives, but its
normalization is now explicitly the V46 factor
\(d_{\rm in}/d_T\), not a representation-formation probability.
\end{proof}

\section{The exact exclusive-pair quartic polynomial}
\label{sec:V51-pair-polynomial}

Assume temporarily that no coordinate belongs to three or four input
supports.  Then
\begin{equation}
 X_i=P_i\mathbin{\dot\cup}
 \mathop{\dot\bigcup}_{j\ne i}L_{ij},
 \qquad
 L_{ij}:=X_i\cap X_j,
 \label{eq:V51-exclusive-pair-partition}
\end{equation}
and all ten banks in \eqref{eq:V51-exclusive-pair-partition} are
coordinate-disjoint.  On \(L_{ij}\), let
\begin{align}
 a_{ij}
 &:=
 q_{\alpha,\boldsymbol R_{i,L_{ij}}}
 q_{\alpha,\boldsymbol R_{j,L_{ij}}},
 \label{eq:V51-pair-mean}\\
 \mathbf D_{ij}
 &:=
 \mathbf Q_{ij,L_{ij}}-a_{ij}I.
 \label{eq:V51-pair-defect}
\end{align}
Empty banks have \(\mathbf D_{ij}=0\) and \(a_{ij}=1\).  Let
\[
 q_P:=\prod_{i=1}^4
 q_{\alpha,\boldsymbol R_{i,P_i}},
\]
and let \(\mathcal C_4\) be the set of connected simple graphs on the
labeled vertex set \(\{1,2,3,4\}\).

\begin{theorem}[Connected-graph formula for the pair-only quartic row]
\label{thm:V51-pair-connected-graph}
Before any bracketing or trace norm,
\begin{equation}
 \boxed{
 \mathbf K_4^{(2)}
 =
 q_P\sum_{G\in\mathcal C_4}
 \left(\prod_{ij\in E(G)}\mathbf D_{ij}\right)
 \left(\prod_{ij\notin E(G)}a_{ij}\right).}
 \label{eq:V51-pair-connected-graph}
\end{equation}
Thus the exclusive-pair row contains exactly the \(38\) connected
labeled graphs; its \(16\) three-edge terms are the labeled spanning
trees.
\end{theorem}

\begin{proof}
On the independent bank \(L_{ij}\),
\(\mathbf Q_{ij,L_{ij}}=a_{ij}I+\mathbf D_{ij}\).
Expand this identity in every row-partition moment entering the fourth
cumulant.  A choice of \(\mathbf D_{ij}\) records an edge \(ij\);
choosing \(a_{ij}I\) records no connection.  The Möbius coefficient of
a fixed edge set is zero unless its graph connects the four row
labels, and is one when it does.  This is also the specialization of
\eqref{eq:V39-connected-replica-formula} obtained by aggregating all
local pair blocks on each \(L_{ij}\).  Private coordinates contribute
the common scalar \(q_P\).  Operators belonging to different banks act
on disjoint coordinate tensor factors, so their products require no
ordering convention.  The graph count is the standard direct count:
\(64\) graphs minus the \(26\) disconnected graphs.
\end{proof}

\section{Spanning-tree routing without a formation kernel}
\label{sec:V51-tree-routing}

Fix an order of the sixteen labeled spanning trees.  For
\(G\in\mathcal C_4\), let \(\tau(G)\) be the first tree contained in
\(G\), and write
\begin{equation}
 \mathbf r_{G,\tau}
 :=
 q_P
 \left(\prod_{ij\in E(G)\setminus E(\tau)}\mathbf D_{ij}\right)
 \left(\prod_{ij\notin E(G)}a_{ij}\right).
 \label{eq:V51-auxiliary-graph-factor}
\end{equation}
After resolving all pair-bank outputs, denote the corresponding scalar
by \(r_{G,\tau}\).  Let \(\Xi_{\rm aux}\) be the sum of the private and
unselected pair-bank output masses.

\begin{lemma}[Surplus edges are admissible assembled spectators]
\label{lem:V51-surplus-admissible}
For every routed connected graph,
\begin{equation}
 |r_{G,\tau}|\le1,
 \qquad
 s_\alpha\Xi_{\rm aux}|r_{G,\tau}|\le C,
 \label{eq:V51-surplus-admissible}
\end{equation}
with a constant independent of supports and representations.
\end{lemma}

\begin{proof}
Every mean multiplier lies in \([0,1]\).  Every resolved defect is a
difference of two numbers in \([0,1]\), so has modulus at most one.
For an unselected pair bank carrying a defect, use
\eqref{eq:V46-defect-first-moment}; for a mean bank or a private bank,
use \eqref{eq:V46-product-first-moment}.  Split
\(\Xi_{\rm aux}\) among the at most ten bank outputs and pay each piece
with one factor belonging to that bank while bounding all other factors
by one.  The number of bank types is fixed.
\end{proof}

\begin{corollary}[Exact routed pair decomposition]
\label{cor:V51-routed-pair-decomposition}
\begin{equation}
 \mathbf K_4^{(2)}
 =
 \sum_{\tau\in\mathfrak T_4}
 \sum_{\substack{G\in\mathcal C_4\\\tau(G)=\tau}}
 \left(\prod_{ij\in E(\tau)}\mathbf D_{ij}\right)
 \mathbf r_{G,\tau}.
 \label{eq:V51-routed-pair-decomposition}
\end{equation}
This is a finite exact operator identity, not a probabilistic
representation-formation average.
\end{corollary}

\begin{proof}
Partition the \(38\) summands in
\eqref{eq:V51-pair-connected-graph} by \(\tau(G)\) and factor the three
selected defects.
\end{proof}

\section{A three-defect one-resolvent theorem}
\label{sec:V51-three-defect}

For a selected tree \(\tau\), set
\begin{align}
 h_{ij}&:=\sum_{e\in L_{ij}}A_{i,e}A_{j,e},
 \qquad ij\in E(\tau),
 \label{eq:V51-selected-edge-mass}\\
 P_\tau&:=\prod_{ij\in E(\tau)}h_{ij},
 \label{eq:V51-tree-numerator}\\
 D_\tau&:=\prod_{i=1}^4\Xi_i^{\deg_\tau(i)-1},
 \qquad
 \mathsf W_\tau:=\frac{P_\tau}{D_\tau}.
 \label{eq:V51-quartic-tree-weight}
\end{align}
Only internal vertices occur in \(D_\tau\), and
\begin{equation}
 \sum_{i=1}^4(\deg_\tau(i)-1)=2.
 \label{eq:V51-two-normalizing-powers}
\end{equation}
Resolve each selected pair bank into an output
\(\boldsymbol S_{ij}\), and let \(\delta_{ij}\) be the corresponding
eigenvalue of \(\mathbf D_{ij}\).  For a full nontrivial output define
\begin{equation}
 \mathcal Q_{G,\tau}(\boldsymbol T)
 :=
 \frac{\Xi(\boldsymbol T)|r_{G,\tau}|}
      {1-q_{\alpha,\boldsymbol T}}
 \prod_{ij\in E(\tau)}|\delta_{ij}|.
 \label{eq:V51-three-defect-response}
\end{equation}

\begin{theorem}[Uniform three-defect one-resolvent bound]
\label{thm:V51-three-defect-one-resolvent}
Every routed graph block satisfies
\begin{equation}
 \boxed{
 \mathcal Q_{G,\tau}(\boldsymbol T)
 \le Cs_\alpha^2P_\tau.}
 \label{eq:V51-three-defect-one-resolvent}
\end{equation}
\end{theorem}

\begin{proof}
On the low-output set
\(s_\alpha C_E(\boldsymbol T)\le1\), use
\eqref{eq:V46-low-output-resolvent}, \(|r_{G,\tau}|\le1\), and three
defect caps:
\[
 \mathcal Q_{G,\tau}(\boldsymbol T)
 \le
 \frac C{s_\alpha}
 \prod_{ij\in E(\tau)}(Cs_\alpha h_{ij})
 \le Cs_\alpha^2P_\tau.
\]

On the complementary output set, the denominator has a fixed lower
bound.  Split the final mass into the three selected pair outputs and
\(\Xi_{\rm aux}\).  For a selected edge \(ij\), fusion and exact
activity give
\(\Xi(\boldsymbol S_{ij})\le Ch_{ij}\), whereas
\eqref{eq:V46-defect-first-moment} gives
\(\Xi(\boldsymbol S_{ij})|\delta_{ij}|\le C/s_\alpha\).
Consequently its contribution is at most
\begin{align*}
 C\min\left\{
 s_\alpha^3h_{ij}P_\tau,\,
 s_\alpha\frac{P_\tau}{h_{ij}}
 \right\}
 &=
 Cs_\alpha^2P_\tau
 \min\left\{
 s_\alpha h_{ij},\,
 (s_\alpha h_{ij})^{-1}
 \right\}\\
 &\le Cs_\alpha^2P_\tau.
\end{align*}
The auxiliary contribution is at most
\[
 \frac C{s_\alpha}
 \prod_{ij\in E(\tau)}(Cs_\alpha h_{ij})
 \le Cs_\alpha^2P_\tau
\]
by \eqref{eq:V51-surplus-admissible}.  Summing the four output pieces
proves the result.
\end{proof}

\begin{theorem}[Complete thermal exclusive-pair quartic sector]
\label{thm:V51-thermal-pair-sector}
Fix \(M<\infty\).  If
\begin{equation}
 s_\alpha\Xi_i\le M
 \quad\hbox{for every \(i\) with \(\deg_\tau(i)\ge2\),}
 \label{eq:V51-internal-thermal}
\end{equation}
then every routed graph block satisfies
\begin{equation}
 \mathcal Q_{G,\tau}(\boldsymbol T)
 \le C_M\mathsf W_\tau.
 \label{eq:V51-thermal-tree-payment}
\end{equation}
After exact raw pinching, the sum of all such arbitrary-coefficient
blocks obeys the \(m=4\) marked-tree atlas restricted to the thermal
exclusive-pair sector.
\end{theorem}

\begin{proof}
Equation \eqref{eq:V51-two-normalizing-powers} and
\eqref{eq:V51-internal-thermal} imply
\[
 D_\tau\le M^2s_\alpha^{-2}.
\]
Therefore
\[
 s_\alpha^2P_\tau\le M^2\frac{P_\tau}{D_\tau}.
\]
Combine this with
\Cref{thm:V51-three-defect-one-resolvent}.  At coefficient level use
\eqref{eq:V46-exact-multiplier-normalization}; the BFA output dimension
cancels \(d_{\rm in}/d_{\boldsymbol T}\), and
\eqref{eq:V46-global-raw-pinching} sums the raw blocks.  Expanding
\eqref{eq:V51-quartic-tree-weight} over its three edge markings
factors the four input sums exactly as in
\eqref{eq:V32-tree-mass-factorization}, now with
\(\deg_\tau(i)\) marks on input \(i\).  The finite \(38\)-graph and
\(16\)-tree routing changes only the universal constant.
\end{proof}

\section{Minimal trees with genuinely exterior hot mass}
\label{sec:V51-exterior-hot}

Now restrict to a minimal graph \(G=\tau\).  For each vertex define its
selected-tree and exterior masses and its exterior singleton
multiplier:
\begin{align}
 Z_i^\tau
 &:=
 \sum_{ij\in E(\tau)}
 \Xi(\boldsymbol R_{i,L_{ij}}),
 \label{eq:V51-selected-tree-input-mass}\\
 \Pi_i^\tau&:=\Xi_i-Z_i^\tau,
 \qquad
 b_i^\tau:=
 q_{\alpha,\boldsymbol R_{i,
 P_i\cup\bigcup_{ij\notin E(\tau)}L_{ij}}}.
 \label{eq:V51-exterior-input-data}
\end{align}
The minimal-tree spectator is exactly
\begin{equation}
 r_\tau=q_P\prod_{ij\notin E(\tau)}a_{ij}
 =\prod_{i=1}^4b_i^\tau.
 \label{eq:V51-minimal-tree-spectator}
\end{equation}
Let \(\Xi_{\rm sp}^\tau\) be the output mass on those exterior banks.

\begin{lemma}[Degree-two and degree-three exterior debits]
\label{lem:V51-exterior-debits}
If \(s_\alpha\Pi_i^\tau>M\), then
\begin{align}
 b_i^\tau
 &\le\frac C{(s_\alpha\Pi_i^\tau)^2},
 \label{eq:V51-exterior-quadratic}\\
 \Xi_{\rm sp}^\tau|r_\tau|
 &\le\frac C{s_\alpha^2\Pi_i^\tau}.
 \label{eq:V51-exterior-quadratic-output}
\end{align}
If \(i\) is the center of a star and the same hot hypothesis holds,
then also
\begin{align}
 b_i^\tau
 &\le\frac C{(s_\alpha\Pi_i^\tau)^3},
 \label{eq:V51-exterior-cubic}\\
 \Xi_{\rm sp}^\tau|r_\tau|
 &\le\frac C{s_\alpha^3(\Pi_i^\tau)^2}.
 \label{eq:V51-exterior-cubic-output}
\end{align}
\end{lemma}

\begin{proof}
Equations \eqref{eq:V51-exterior-quadratic} and
\eqref{eq:V51-exterior-cubic} are respectively the product second and
third moments
\eqref{eq:V32-balanced-second-moment} and
\eqref{eq:V42-balanced-third-moment} on the exact exterior channel of
input \(i\).

Split every exterior fused output mass into its four input
contributions.  The input-\(i\) contribution is at most
\[
 \Pi_i^\tau b_i^\tau
 \le
 \frac C{s_\alpha^2\Pi_i^\tau}
\]
in the quadratic case and at most
\(C/[s_\alpha^3(\Pi_i^\tau)^2]\) in the cubic case.  A contribution
from another input is paid by its product first moment while
\(b_i^\tau\) is retained.  It is bounded respectively by
\[
 \frac C{s_\alpha}b_i^\tau
 \le\frac {C_M}{s_\alpha^2\Pi_i^\tau},
 \qquad
 \frac C{s_\alpha}b_i^\tau
 \le\frac {C_M}{s_\alpha^3(\Pi_i^\tau)^2}.
\]
There are only six exterior pair banks and four private banks.
\end{proof}

\begin{theorem}[Exterior-hot minimal star]
\label{thm:V51-exterior-hot-star}
Let \(\tau\) be a star with center \(c\).  If either
\[
 s_\alpha\Xi_c\le M
\]
or
\begin{equation}
 s_\alpha\Pi_c^\tau>M,
 \qquad
 \Pi_c^\tau\ge\varepsilon\Xi_c,
 \label{eq:V51-star-exterior-dominance}
\end{equation}
then its minimal-tree block satisfies
\begin{equation}
 \mathcal Q_{\tau,\tau}(\boldsymbol T)
 \le C_{M,\varepsilon}\frac{P_\tau}{\Xi_c^2}
 =C_{M,\varepsilon}\mathsf W_\tau.
 \label{eq:V51-exterior-star-payment}
\end{equation}
\end{theorem}

\begin{proof}
The thermal alternative is
\Cref{thm:V51-thermal-pair-sector}.  Assume
\eqref{eq:V51-star-exterior-dominance} and abbreviate
\(\Pi=\Pi_c^\tau\), \(b=b_c^\tau\).

On low final outputs, the resolvent estimate, the three defect caps,
and \eqref{eq:V51-exterior-cubic} give
\[
 \mathcal Q_{\tau,\tau}
 \le Cb s_\alpha^2P_\tau
 \le\frac C{s_\alpha\Pi^3}P_\tau
 \le\frac C{\Pi^2}P_\tau.
\]
On high final outputs, the spectator-output contribution is at most
\(CP_\tau/\Pi^2\) by
\eqref{eq:V51-exterior-cubic-output} and the three defect caps.
For a selected edge \(e\), combine fusion with the selected defect
first moment exactly as in
\Cref{thm:V51-three-defect-one-resolvent}.  Its contribution is at most
\[
 Cb s_\alpha^2P_\tau
 \min\{s_\alpha h_e,(s_\alpha h_e)^{-1}\}
 \le Cb s_\alpha^2P_\tau
 \le\frac C{\Pi^2}P_\tau.
\]
The high-output denominator has a fixed floor.  Finally
\(\Pi\ge\varepsilon\Xi_c\) converts \(\Pi^{-2}\) to
\(C_\varepsilon\Xi_c^{-2}\).
\end{proof}

\begin{theorem}[Exterior-hot minimal path]
\label{thm:V51-exterior-hot-path}
Let \(\tau\) be a path and let \(u,v\) be its two internal vertices.
Suppose every internal vertex \(i\in\{u,v\}\) satisfies either
\begin{equation}
 s_\alpha\Xi_i\le M,
 \label{eq:V51-path-thermal-alternative}
\end{equation}
or
\begin{equation}
 s_\alpha\Pi_i^\tau>M,
 \qquad
 \Pi_i^\tau\ge\varepsilon\Xi_i.
 \label{eq:V51-path-exterior-alternative}
\end{equation}
Then
\begin{equation}
 \mathcal Q_{\tau,\tau}(\boldsymbol T)
 \le
 C_{M,\varepsilon}\frac{P_\tau}{\Xi_u\Xi_v}
 =C_{M,\varepsilon}\mathsf W_\tau.
 \label{eq:V51-exterior-path-payment}
\end{equation}
\end{theorem}

\begin{proof}
If both vertices satisfy
\eqref{eq:V51-path-thermal-alternative}, use
\Cref{thm:V51-thermal-pair-sector}.

Suppose only \(u\) is exterior-hot and \(v\) is thermal.  Put
\(\Pi=\Pi_u^\tau\), \(b=b_u^\tau\).  On low final outputs,
\[
 \mathcal Q_{\tau,\tau}
 \le Cb s_\alpha^2P_\tau
 \le\frac C{\Pi^2}P_\tau
 \le\frac {Cs_\alpha}{\Pi}P_\tau.
\]
On high outputs, \eqref{eq:V51-exterior-quadratic-output} pays the
exterior mass and gives \(Cs_\alpha P_\tau/\Pi\).  Each selected-output
piece is at most
\[
 Cb s_\alpha^2P_\tau
 \min\{s_\alpha h_e,(s_\alpha h_e)^{-1}\}
 \le\frac C{\Pi^2}P_\tau
 \le\frac {Cs_\alpha}{\Pi}P_\tau.
\]
Use \(\Pi\ge\varepsilon\Xi_u\) and
\(\Xi_v\le M/s_\alpha\) to obtain
\eqref{eq:V51-exterior-path-payment}.  The case with \(u,v\)
interchanged is identical.

It remains that both \(u\) and \(v\) are exterior-hot.  Write
\(\Pi_u,\Pi_v,b_u,b_v\) for their data.  The low-output and every
selected-output contribution are bounded by
\[
 Cs_\alpha^2b_ub_vP_\tau
 \le
 \frac {CP_\tau}
 {s_\alpha^2\Pi_u^2\Pi_v^2}
 \le\frac {C_MP_\tau}{\Pi_u\Pi_v}.
\]
For an exterior output, allocate its fused mass to input ancestors.
The input-\(u\) part, after the three defect caps, is at most
\[
 Cs_\alpha^3P_\tau\,\Pi_ub_ub_v
 \le
 \frac {CP_\tau}{s_\alpha\Pi_u\Pi_v^2}
 \le\frac {C_MP_\tau}{\Pi_u\Pi_v},
\]
and the input-\(v\) part is symmetric.  Any other input contribution
uses a product first moment and is at most
\(Cs_\alpha^2b_ub_vP_\tau\), already bounded above.  The denominator
has a fixed floor.  Conclude with
\(\Pi_i\ge\varepsilon\Xi_i\).
\end{proof}

\begin{corollary}[Coefficient closure of the proved pair-tree regimes]
\label{cor:V51-proved-pair-regimes}
For arbitrary trace-class Fourier coefficients, raw pinching lifts:
\begin{enumerate}[label=\textup{(\roman*)},leftmargin=3.7em]
\item every internally thermal routed connected graph;
\item every minimal star satisfying
      \eqref{eq:V51-star-exterior-dominance}; and
\item every minimal path satisfying
      \eqref{eq:V51-path-thermal-alternative} or
      \eqref{eq:V51-path-exterior-alternative} at both internal
      vertices
\end{enumerate}
to their required quartic marked-tree terms without dimension or
multiplicity loss.
\end{corollary}

\begin{proof}
Use respectively
\Cref{thm:V51-thermal-pair-sector,
thm:V51-exterior-hot-star,thm:V51-exterior-hot-path}, followed by the
exact normalization and raw-pinching argument in the last paragraph of
\Cref{thm:V51-thermal-pair-sector}.
\end{proof}

\section{The edge-concentrated star regression}
\label{sec:V51-star-regression}

The complement of the preceding theorems is not merely a more
complicated multiplier split.  It contains a family on which the
desired pointwise tree inequality is false.

\begin{counterexample}[No uniform pointwise star payment]
\label{cth:V51-pointwise-star-false}
There are exclusive-pair star blocks, centered at row \(1\), for which
\begin{equation}
 \mathcal Q_{\tau,\tau}(\boldsymbol T)\asymp s_\alpha,
 \qquad
 \mathsf W_\tau\asymp\frac1N.
 \label{eq:V51-star-regression-scales}
\end{equation}
Consequently no constant independent of support size can satisfy
\(\mathcal Q_{\tau,\tau}\le C\mathsf W_\tau\) pointwise in resolved
pair channels.
\end{counterexample}

\begin{proof}
Take the star
\(\tau=\{12,13,14\}\).  Let \(L_{12}\) contain \(N\) coordinates, put
\(R_{1,e}=\mathbf3\), \(R_{2,e}=\overline{\mathbf3}\), and select the
singlet in every
\(\mathbf3\otimes\overline{\mathbf3}\) product.  Thus the selected
output on \(L_{12}\) is trivial and
\begin{equation}
 \delta_{12}
 =
 1-(\eta_{\mathbf3}\eta_{\overline{\mathbf3}})^N.
 \label{eq:V51-long-singlet-defect}
\end{equation}
Let \(L_{13}\) and \(L_{14}\) each consist of one coordinate.  Put
\(R_1=R_3=\mathbf3\) there and select the \(\mathbf6\) channel; do the
same for rows \(1,4\).  All other banks are empty.

For fixed irreducibles,
\eqref{eq:V10-eta-asymptotic} gives
\[
 \eta_R=1-s_\alpha C_2(R)+O_R(s_\alpha^2).
\]
Since
\(C_2(\mathbf6)-2C_2(\mathbf3)=2/3\ne0\), both short-edge defects have
modulus comparable to \(s_\alpha\) for all sufficiently large
\(\alpha\).  Their two output masses are fixed and nonzero, so the full
final gap is comparable to \(s_\alpha\).  If
\(s_\alpha N\to\infty\), \eqref{eq:V51-long-singlet-defect} is bounded
below by a positive constant.  This proves the first estimate in
\eqref{eq:V51-star-regression-scales}.

On the other hand,
\[
 \Xi_1\asymp N,\qquad
 h_{12}\asymp N,\qquad h_{13}+h_{14}\asymp1.
\]
Hence \(P_\tau/\Xi_1^2\asymp N^{-1}\).  Taking
\(N/s_\alpha^{-1}\to\infty\) makes the ratio of the pointwise response
to the star weight diverge.
\end{proof}

\begin{remark}[Why multiplier moments cannot repair the regression]
\label{rem:V51-no-pointwise-repair}
In the long bank of
\Cref{cth:V51-pointwise-star-false}, the final pair output is trivial,
so every product moment of its output multiplier equals one.  The
input singleton multipliers occur only in the subtracted term of
\eqref{eq:V51-long-singlet-defect}; after \(s_\alpha N\gg1\), their
damping has already saturated the defect at one.  Thus neither a
higher product moment nor a sharper final-gap split can create the
missing \(N^{-1}\) pointwise.  The missing factor must come from the
actual coefficient compression into the singlet.
\end{remark}

\begin{theorem}[The regression is paid by genuine raw coefficient mass]
\label{thm:V51-singlet-star-paid}
For rank-one product coefficient atoms on the family above, the
normalized raw coefficient mass of the long singlet block is
\(3^{-N}\).  Its coefficient-weighted star response obeys
\begin{equation}
 3^{-N}\mathcal Q_{\tau,\tau}(\boldsymbol T)
 \le C\mathsf W_\tau
 \label{eq:V51-singlet-star-paid}
\end{equation}
uniformly in \(N\) and sufficiently large \(\alpha\).
\end{theorem}

\begin{proof}
Tensor Schur orthogonality over \(L_{12}\).  Equivalently apply
\eqref{eq:V33-Schur-tax} to
\(\boldsymbol R=\mathbf3^{\boxtimes N}\), whose carrier dimension is
\(3^N\).  The raw pinching mass is exactly \(3^{-N}\), in agreement
with \eqref{eq:V45-actual-singlet-mass}.  The pointwise response is at
most \(Cs_\alpha\) by the calculation above and the universal defect
and resolvent bounds.  Since
\[
 s_\alpha3^{-N}\le C N^{-1}\asymp C\mathsf W_\tau,
\]
the claimed coefficient estimate follows.  This uses the actual raw
mass \(3^{-N}\), not the false representation-dimension weight
\(3^{-2N}\).
\end{proof}

\section{The exact coefficient-sensitive pair gates}
\label{sec:V51-coefficient-gates}

Resolve the selected pair-bank intermediates, then the remaining
isotypic and multiplicity blocks, by successive orthogonal pinchings.
For a pure input coefficient tensor \(\mathcal B\), write
\begin{equation}
 \omega_{\boldsymbol S,\boldsymbol T,a}
 :=
 \frac{\|X_{\boldsymbol S,\boldsymbol T,a}\|_1}
      {\prod_{i=1}^4\|B_i\|_1}.
 \label{eq:V51-refined-raw-mass}
\end{equation}
Here \(\boldsymbol S\) denotes all pair-bank outputs and \(a\) all
remaining multiplicity labels.  Successive raw pinching gives
\begin{equation}
 \omega_{\boldsymbol S,\boldsymbol T,a}\ge0,
 \qquad
 \sum_{\boldsymbol S,\boldsymbol T,a}
 \omega_{\boldsymbol S,\boldsymbol T,a}\le1.
 \label{eq:V51-refined-subprobability}
\end{equation}
No dimension ratio is included in \(\omega\).

\begin{openproblem}[Coefficient-sensitive three-defect star gate]
\label{op:V51-star-gate}
For every star \(\tau\) with center \(c\), prove on the remaining
edge-concentrated blocks that
\begin{equation}
 \boxed{
 \sum_{\boldsymbol S,\boldsymbol T,a}
 \omega_{\boldsymbol S,\boldsymbol T,a}
 \frac{\Xi(\boldsymbol T)|r_{G,\tau}|}
      {1-q_{\alpha,\boldsymbol T}}
 \prod_{ij\in E(\tau)}|\delta_{ij}(\boldsymbol S_{ij})|
 \le
 C\frac{P_\tau}{\Xi_c^2}.}
 \label{eq:V51-coefficient-star-gate}
\end{equation}
The sum is restricted to the graph block under consideration.  The
raw coefficient mass, all three defects, every surplus-edge factor,
and the one final resolvent must remain assembled until after the
pair-intermediate pinching.
\end{openproblem}

\begin{openproblem}[Coefficient-sensitive three-defect path gate]
\label{op:V51-path-gate}
For every path \(\tau\) with internal vertices \(u,v\), prove on the
remaining edge-concentrated blocks that
\begin{equation}
 \boxed{
 \sum_{\boldsymbol S,\boldsymbol T,a}
 \omega_{\boldsymbol S,\boldsymbol T,a}
 \frac{\Xi(\boldsymbol T)|r_{G,\tau}|}
      {1-q_{\alpha,\boldsymbol T}}
 \prod_{ij\in E(\tau)}|\delta_{ij}(\boldsymbol S_{ij})|
 \le
 C\frac{P_\tau}{\Xi_u\Xi_v}.}
 \label{eq:V51-coefficient-path-gate}
\end{equation}
\end{openproblem}

\begin{theorem}[Exact reduction of the exclusive-pair quartic row]
\label{thm:V51-pair-row-reduction}
The full arbitrary-coefficient exclusive-pair \(m=4\) marked-tree row
follows from
\eqref{eq:V51-coefficient-star-gate} and
\eqref{eq:V51-coefficient-path-gate} on the blocks not already covered
by \Cref{cor:V51-proved-pair-regimes}.
\end{theorem}

\begin{proof}
Equation \eqref{eq:V51-routed-pair-decomposition} assigns every
connected graph to one of the sixteen trees.  The four stars require
\eqref{eq:V51-coefficient-star-gate}; the twelve paths require
\eqref{eq:V51-coefficient-path-gate}.  Thermal graph blocks and
minimal-tree exterior-hot blocks are already closed by
\Cref{cor:V51-proved-pair-regimes}.  For every other routed block,
apply the corresponding coefficient gate.  The exact factor
\(d_{\rm in}/d_{\boldsymbol T}\) cancels the output BFA dimension, and
\eqref{eq:V51-refined-subprobability} performs all intermediate and
final sums without a path count.  Finally expand
\(P_\tau/D_\tau\) over the three edge coordinates to obtain precisely
the degree-marked input seminorms.  Summing \(38\) fixed graph blocks
changes only the universal constant.
\end{proof}

\begin{firewall}[What the star regression proves]
\label{fire:V51-star-regression}
\Cref{cth:V51-pointwise-star-false} disproves a pointwise proof
strategy, not quartic MTA.  The explicit family is controlled by its
genuine Schur compression in
\Cref{thm:V51-singlet-star-paid}.  Conversely, that special Schur
calculation does not prove
\eqref{eq:V51-coefficient-star-gate} for arbitrary representations,
coefficients, nontrivial cold pair outputs, or surplus graph edges.
\end{firewall}

\section{V51 ledger and next acquisition}
\label{sec:V51-ledger}

\begin{theorem}[Integrated V51 statement]
\label{thm:V51-integrated}
Preserving V1--V50, V51 proves:
\begin{enumerate}[label=\textup{(V51.\arabic*)},leftmargin=5.5em]
\item the post-retraction quartic validity map
      \Cref{thm:V51-validity-map};
\item exact revalidation of the V41 fourfold local-tree compression
      with the V46 factor \(d_{\rm in}/d_T\);
\item the \(38\)-term connected-graph formula
      \eqref{eq:V51-pair-connected-graph} for the complete
      exclusive-pair quartic row;
\item exact routing of every connected graph through one of the
      sixteen spanning trees, with surplus defects retained as
      admissible assembled spectators;
\item the three-defect one-resolvent estimate
      \eqref{eq:V51-three-defect-one-resolvent} and complete
      arbitrary-coefficient closure of the internally thermal
      exclusive-pair sector;
\item cubic star and quadratic path damping on genuine exterior-hot
      minimal-tree blocks;
\item the explicit long dual-singlet counterexample to every uniform
      pointwise star payment;
\item coefficient-level closure of that regression by its actual
      \(3^{-N}\) raw Schur mass; and
\item reduction of the remaining exclusive-pair row to the assembled
      coefficient gates
      \eqref{eq:V51-coefficient-star-gate}--%
      \eqref{eq:V51-coefficient-path-gate}.
\end{enumerate}
\end{theorem}

\begin{proof}
The nine items are respectively
\Cref{thm:V51-validity-map,
thm:V51-corrected-fourfold-compression,
thm:V51-pair-connected-graph,
lem:V51-surplus-admissible,cor:V51-routed-pair-decomposition,
thm:V51-three-defect-one-resolvent,thm:V51-thermal-pair-sector,
lem:V51-exterior-debits,thm:V51-exterior-hot-star,
thm:V51-exterior-hot-path,
cth:V51-pointwise-star-false,
thm:V51-singlet-star-paid,
thm:V51-pair-row-reduction}.
\end{proof}

\begingroup
\small
\begin{longtable}{>{\raggedright\arraybackslash}p{0.29\textwidth}
                  >{\raggedright\arraybackslash}p{0.15\textwidth}
                  >{\raggedright\arraybackslash}p{0.48\textwidth}}
\toprule
\textbf{V51 row}&\textbf{Status}&\textbf{Advance or obstruction}\\
\midrule
\endfirsthead
\toprule
\textbf{V51 row}&\textbf{Status}&\textbf{Advance or obstruction}\\
\midrule
\endhead
V41 local fourfold compression
& \MGclosed
& The operator card survives with the exact
  \(d_{\rm in}/d_T\) coefficient formula and raw pinching.\\[0.35em]
Exclusive-pair connected polynomial
& \MGclosed
& Thirty-eight connected graphs are routed exactly through sixteen
  stars and paths; no formation probability is inserted.\\[0.35em]
Internally thermal pair sector
& \MGclosed
& Three defect caps and one final resolvent supply the two quartic
  normalizing powers.  Surplus defects remain assembled spectators.\\[0.35em]
Minimal-tree exterior-hot sector
& \MGclosed
& Cubic damping pays a degree-three star center; one or two quadratic
  debits pay the internal vertices of a path.\\[0.35em]
Pointwise edge-concentrated star
& \MGfail
& An \(N\)-link
  \(\mathbf3\otimes\overline{\mathbf3}\to\mathbf1\) bank has response
  \(s_\alpha\) but star weight \(N^{-1}\).\\[0.35em]
Dual-singlet coefficient regression
& \MGclosed
& Its actual raw mass is \(3^{-N}\), which pays the missing star mark;
  the artificial dimension-law mass \(3^{-2N}\) is not used.\\[0.35em]
General coefficient-sensitive star/path gates
& \MGreduced
& The unresolved sums are exactly
  \eqref{eq:V51-coefficient-star-gate} and
  \eqref{eq:V51-coefficient-path-gate}, with three defects and the final
  resolvent still attached.\\[0.35em]
Complete exclusive-pair quartic row
& \MGopen
& It is equivalent, for the remaining routed blocks, to the two
  coefficient-sensitive gates.\\[0.35em]
Other quartic skeletons and all-order MTA
& \MGopen
& The \([4]\), \([3,2]\), and \([3,3]\) fibres remain suspended until
  the first coefficient-sensitive pair gate is settled.\\[0.35em]
Full Yang--Mills mass gap
& \MGopen
& WUFC/CPM, nonlinear construction, finite-edge estimates, continuum
  and OS positivity/nontriviality, and the spectral passage remain
  downstream.\\
\bottomrule
\end{longtable}
\endgroup

\begin{openproblem}[V52 mandate: assembled star formation--defect tax]
\label{op:V52}
\begin{enumerate}[label=\textup{(V52.\arabic*)},leftmargin=5.5em]
\item Begin with the edge-concentrated star gate
      \eqref{eq:V51-coefficient-star-gate}.  Resolve one hot selected
      pair bank first, but keep its raw coefficient block, all three
      binary defects, every surplus-edge factor, and the single final
      resolvent in the same summand.
\item Split the hot pair outputs by their actual V34 partial-trace tax
      \[
      \min\{d_S/d_R,d_S/d_Q,1\}.
      \]
      Apply this only to the genuine compressed coefficient block.
      Do not turn it into a representation-dimension probability.
\item On the dimension-small sector, convert the partial-trace tax
      directly into one missing normalized center mark, then obtain the
      second center mark from the remaining two-defect quotient.
\item On the dimension-mesoscopic sector, keep the binary defect and
      use the resolved Casimir output together with the final gap.
      Sum an actual coefficient-weighted defect/Racah quantity; a bare
      cold-output tail is inadmissible.
\item Test the resulting estimate both on the \(N\)-link dual-singlet
      family of \Cref{cth:V51-pointwise-star-false} and on the V34
      highest-weight chain.  The first must produce the Schur mark and
      the second must retain the other two defects.
\item If the star gate closes, apply the same assembled split at the
      two internal vertices of the path gate.  Only then declare the
      complete exclusive-pair quartic row restored.
\item Do not enter the \([4]\), \([3,2]\), or \([3,3]\) topology until
      the pair row is closed or the exact stronger upstream estimate
      needed by those fibres has been identified.
\item Preserve only the newest YM file.  No companion audit is due;
      the next read-only SM/NS audit remains V55.
\end{enumerate}
\end{openproblem}

\begin{firewall}[End-of-V51 claim boundary]
\label{fire:V51-claim-boundary}
V51 revalidates the local fourfold coefficient compression, derives
the exact exclusive-pair connected polynomial, and closes all
internally thermal graph blocks together with the genuine
exterior-hot minimal stars and paths.  It proves that a pointwise
edge-concentrated star theorem is false and that the explicit
dual-singlet regression is nevertheless paid by actual raw coefficient
mass.  It does not prove the general coefficient-sensitive star or
path gate, the complete exclusive-pair row, any other quartic topology,
global \(m=4\), all-order MTA, WUFC, CPM, the nonlinear Perron packet,
finite-edge margins, capacity, protected-sector floors, local-algebra
noncollapse, reflection positivity, a nontrivial continuum OS state,
or a positive Yang--Mills spectral gap.  The full Yang--Mills mass gap
is not proved.
\end{firewall}

\section*{Version V51 append-only record}
\addcontentsline{toc}{section}{Version V51 append-only record}
The complete V50 source was retained before this continuation.  Its
SHA--256 digest was
\begin{center}\scriptsize\ttfamily
9b5c1638db817bc48ead0751e48c262826844360d7bd6c005df7efec772e6e02.
\end{center}
No mathematical content of V50 was altered.  On the next iteration,
retain this full file, append before its unique active document-closing
command, increment the filename to \texttt{\_V52.tex}, and remove only
the superseded V51 file from this Yang--Mills note series.

% =============================================================================
% END APPENDED VERSION V51 MATERIAL
% =============================================================================

% =============================================================================
% BEGIN APPENDED VERSION V52 MATERIAL
% =============================================================================

\chapter{V52: Aggregate Schur tax and the quartic star capacity}
\label{ch:V52}

\section{Purpose and correction of the one-copy route}
\label{sec:V52-purpose}

V51 shows that the unresolved quartic star cannot be bounded
pointwise in a resolved pair output.  Its long dual-singlet channel is
large as a multiplier but small as an actual coefficient compression.
Applying the V34 tax separately to every output copy would still be
invalid: the number of copies could be reintroduced before raw
pinching.  The first task is therefore an aggregate Schur theorem for
an arbitrary family of multiplicity-resolved outputs, stable under
spectator amplification.

\section{Aggregate multiplicity-resolved Schur tax}
\label{sec:V52-aggregate-Schur}

Let \(R,Q\) be irreducible representations of a compact group and fix
isometric Clebsch--Gordan embeddings
\[
 \iota_{S,\mu}:H_S\longrightarrow H_R\otimes H_Q.
\]
Let \(\Omega\) be any family of output copies \((S,\mu)\), and set
\begin{equation}
 P_\Omega:=\sum_{(S,\mu)\in\Omega}
 \iota_{S,\mu}\iota_{S,\mu}^*,
 \qquad
 D_\Omega:=\operatorname{rank}P_\Omega
 =\sum_{(S,\mu)\in\Omega}d_S.
 \label{eq:V52-aggregate-projector}
\end{equation}

\begin{lemma}[Exact aggregate partial traces]
\label{lem:V52-aggregate-partial-traces}
\begin{equation}
 \boxed{
 \operatorname{Tr}_{H_Q}P_\Omega
 =\frac{D_\Omega}{d_R}I_{H_R},
 \qquad
 \operatorname{Tr}_{H_R}P_\Omega
 =\frac{D_\Omega}{d_Q}I_{H_Q}.}
 \label{eq:V52-aggregate-partial-traces}
\end{equation}
Consequently, with
\begin{equation}
 \tau_\Omega:=
 \min\left\{\frac{D_\Omega}{d_R},
             \frac{D_\Omega}{d_Q},1\right\},
 \label{eq:V52-aggregate-tax}
\end{equation}
every pair of positive trace-one operators \(\rho,\sigma\) satisfies
\begin{equation}
 \operatorname{Tr}[P_\Omega(\rho\otimes\sigma)]
 \le\tau_\Omega.
 \label{eq:V52-product-density-tax}
\end{equation}
\end{lemma}

\begin{proof}
The range of \(P_\Omega\) is invariant under the diagonal action.
Therefore its partial trace over \(H_Q\) commutes with the irreducible
action \(R\), and Schur's lemma makes it scalar.  Its trace is
\(D_\Omega\), proving the first identity; the second is symmetric.
Since \(0\le\sigma\le I\),
\[
 \operatorname{Tr}[P_\Omega(\rho\otimes\sigma)]
 \le
 \operatorname{Tr}[P_\Omega(\rho\otimes I)]
 =D_\Omega/d_R.
\]
Interchange \(R,Q\), and also use \(P_\Omega\le I\), to obtain the
minimum in \eqref{eq:V52-product-density-tax}.
\end{proof}

Let \(\mathcal K,\mathcal L\) be arbitrary spectator Hilbert spaces.
After reordering tensor factors, define
\begin{align}
 \mathcal C_{S,\mu}(A,B):={}&
 (\iota_{S,\mu}^*\otimes I_{\mathcal K\otimes\mathcal L})
 (A\otimes B)
 (\iota_{S,\mu}\otimes I_{\mathcal K\otimes\mathcal L}),
 \label{eq:V52-amplified-compression}
\end{align}
where
\(A\in\mathcal S_1(H_R\otimes\mathcal K)\) and
\(B\in\mathcal S_1(H_Q\otimes\mathcal L)\).

\begin{theorem}[Spectator-amplified aggregate trace tax]
\label{thm:V52-aggregate-trace-tax}
For arbitrary trace-class \(A,B\),
\begin{equation}
 \boxed{
 \sum_{(S,\mu)\in\Omega}
 \|\mathcal C_{S,\mu}(A,B)\|_1
 \le
 \tau_\Omega\|A\|_1\|B\|_1.}
 \label{eq:V52-aggregate-trace-tax}
\end{equation}
The estimate contains the total output-copy dimension \(D_\Omega\),
not the number of copies, and is completely stable under spectators.
\end{theorem}

\begin{proof}
First take rank-one operators
\(A=|x\rangle\langle y|\) and
\(B=|u\rangle\langle v|\).  Put
\[
 x_{S,\mu}:=
 (\iota_{S,\mu}^*\otimes I)(x\otimes u),
 \qquad
 y_{S,\mu}:=
 (\iota_{S,\mu}^*\otimes I)(y\otimes v),
\]
with the understood tensor reordering.  Then
\[
 \|\mathcal C_{S,\mu}(A,B)\|_1
 =\|x_{S,\mu}\|\,\|y_{S,\mu}\|.
\]
Cauchy--Schwarz over \((S,\mu)\) bounds the sum by
\[
 \|(P_\Omega\otimes I)(x\otimes u)\|\,
 \|(P_\Omega\otimes I)(y\otimes v)\|.
\]
The reduced states of \(x\) on \(H_R\) and of \(u\) on \(H_Q\) are
positive and have traces \(\|x\|^2\) and \(\|u\|^2\).  Hence
\eqref{eq:V52-product-density-tax}, and the same estimate for \(y,v\),
bound the last display by
\[
 \tau_\Omega\|x\|\|y\|\|u\|\|v\|.
\]
Expand general \(A,B\) in singular-value decompositions and sum the
rank-one estimates.  The two singular-value sums factor as
\(\|A\|_1\|B\|_1\).
\end{proof}

\begin{remark}[Relation to the V34 one-copy tax]
\label{rem:V52-one-copy-relation}
For \(\Omega=\{(S,\mu)\}\),
\eqref{eq:V52-aggregate-trace-tax} is exactly
\eqref{eq:V34-trace-formation-tax}.  For a family, replacing
\(D_\Omega\) by \(d_S\) and summing would be wrong; the proof above
forms the aggregate invariant projector first and therefore never pays
a multiplicity count.
\end{remark}

\section{A linearized auxiliary two-defect estimate}
\label{sec:V52-linearized-two-defect}

Use the V46 two selected banks \(L,R\), but let the resolved auxiliary
bank have scalar factor \(r_V\) and output mass \(\Xi_V\).  Define its
assembled size
\begin{equation}
 \mathfrak a_V:=|r_V|+
 s_\alpha\Xi_V|r_V|.
 \label{eq:V52-auxiliary-size}
\end{equation}
All remaining singleton spectators are included in \(q_P\).

\begin{lemma}[Linearized auxiliary response]
\label{lem:V52-linearized-two-defect}
The two-defect response of
\eqref{eq:V46-extended-two-defect-quotient} obeys
\begin{equation}
 \boxed{
 \mathcal Q_{L,R}^{V}
 \le
 Cs_\alpha H_LH_R\,\mathfrak a_V.}
 \label{eq:V52-linearized-two-defect}
\end{equation}
The constant is uniform in the auxiliary representation and its
multiplicity.
\end{lemma}

\begin{proof}
On low final outputs, use the low resolvent and two defect caps, leaving
\(|r_V|\):
\[
 \mathcal Q_{L,R}^{V}
 \le Cs_\alpha H_LH_R|r_V|.
\]
On high outputs the denominator has a fixed floor.  For the \(L\) and
\(R\) output masses, take the better of fusion plus two caps and the
corresponding defect first moment plus the other cap.  Exactly as in
\Cref{thm:V51-three-defect-one-resolvent}, each is at most
\(Cs_\alpha H_LH_R|r_V|\).  Private output mass is paid by its product
first moment and has the same bound.  The auxiliary output contributes
\[
 \Xi_V|r_V|(Cs_\alpha H_L)(Cs_\alpha H_R)
 =
 s_\alpha H_LH_R\,
 [Cs_\alpha\Xi_V|r_V|].
\]
Sum the four pieces.
\end{proof}

\section{The weighted pair formation--defect functional}
\label{sec:V52-weighted-functional}

On one exact pair bank, let \(R,Q\) denote the two product-group input
irreducibles.  For every multiplicity-resolved product output
\((S,\mu)\), put
\begin{equation}
 \delta(S):=q_{\alpha,S}-q_{\alpha,R}q_{\alpha,Q},
 \qquad
 a(S):=|\delta(S)|+
 s_\alpha\Xi(S)|\delta(S)|.
 \label{eq:V52-defect-weight}
\end{equation}
The V46 defect first moment and the universal defect cap give
\begin{equation}
 0\le a(S)\le A_*,
 \label{eq:V52-defect-weight-cap}
\end{equation}
for a numerical \(A_*<\infty\).  For trace-class coefficients with
arbitrary spectators, define
\begin{equation}
 \mathfrak F_{R,Q}(A,B)
 :=
 \sum_{S,\mu}
 a(S)\|\mathcal C_{S,\mu}(A,B)\|_1.
 \label{eq:V52-formation-defect-functional}
\end{equation}

For \(0<t<A_*\), let
\begin{align}
 \Omega_t&:=\{(S,\mu):a(S)>t\},
 \label{eq:V52-defect-level-set}\\
 D(t)&:=\sum_{(S,\mu)\in\Omega_t}d_S,
 \label{eq:V52-level-dimension}\\
 \mathcal K_\alpha(R,Q)&:=
 \int_0^{A_*}
 \min\left\{\frac{D(t)}{d_R},
             \frac{D(t)}{d_Q},1\right\}\,\dd t.
 \label{eq:V52-capacity-envelope}
\end{align}

\begin{theorem}[Layer-cake coefficient-capacity bound]
\label{thm:V52-layer-cake-capacity}
For all spectator-amplified trace-class coefficients,
\begin{equation}
 \boxed{
 \mathfrak F_{R,Q}(A,B)
 \le
 \mathcal K_\alpha(R,Q)\|A\|_1\|B\|_1.}
 \label{eq:V52-layer-cake-capacity}
\end{equation}
The envelope \(\mathcal K_\alpha\) is a one-sided output-copy
dimension ratio.  It is neither the tensor-dimension probability
\(d_S/(d_Rd_Q)\) nor a conditioned formation law.
\end{theorem}

\begin{proof}
For nonnegative \(a(S)\),
\[
 a(S)=\int_0^{A_*}\mathbf1_{\{a(S)>t\}}\,\dd t.
\]
Tonelli's theorem and
\Cref{thm:V52-aggregate-trace-tax} give
\begin{align*}
 \mathfrak F_{R,Q}(A,B)
 &=
 \int_0^{A_*}
 \sum_{(S,\mu)\in\Omega_t}
 \|\mathcal C_{S,\mu}(A,B)\|_1\,\dd t\\
 &\le
 \int_0^{A_*}
 \min\left\{\frac{D(t)}{d_R},
             \frac{D(t)}{d_Q},1\right\}\,\dd t
 \|A\|_1\|B\|_1.
\end{align*}
\end{proof}

\begin{remark}[The dual-singlet normalization is automatic]
\label{rem:V52-capacity-singlet}
For the \(N\)-link
\(\mathbf3^{\boxtimes N}\otimes
\overline{\mathbf3}^{\boxtimes N}\) singlet, the relevant level set has
\(D(t)=1\) and \(d_R=d_Q=3^N\).  Thus
\eqref{eq:V52-capacity-envelope} assigns it the genuine scale
\(3^{-N}\), exactly as in
\Cref{thm:V51-singlet-star-paid}; the false scale \(3^{-2N}\) never
appears.
\end{remark}

\section{Dimension-small closure of the edge-concentrated star}
\label{sec:V52-dimension-small-star}

Fix a routed star \(\tau\) centered at \(c\), and select one of its
edges \(e_0=c\ell_0\).  Write the other two edges as \(e_1,e_2\), with
weights \(h_0,h_1,h_2\).  On the \(e_0\) pair bank let \(R_0,Q_0\) be
the product-group inputs.  For a family \(\Omega\) of its resolved
output copies define \(D_\Omega,\tau_\Omega\) by
\eqref{eq:V52-aggregate-projector}--%
\eqref{eq:V52-aggregate-tax}.

\begin{theorem}[Aggregate dimension-small star payment]
\label{thm:V52-dimension-small-star}
Suppose
\begin{equation}
 \tau_\Omega
 \le
 \lambda_0:=
 \min\left\{1,\frac{h_0}{s_\alpha\Xi_c^2}\right\}.
 \label{eq:V52-small-family-threshold}
\end{equation}
Then the arbitrary-coefficient contribution for which the
\(e_0\)-output copy lies in \(\Omega\) satisfies
\begin{equation}
 \boxed{
 \textup{star contribution over \(\Omega\)}
 \le
 C\frac{h_0h_1h_2}{\Xi_c^2}.}
 \label{eq:V52-dimension-small-star}
\end{equation}
This remains true for every surplus-edge graph routed through
\(\tau\).
\end{theorem}

\begin{proof}
Resolve the \(e_0\) pair output first.  Treat its defect as the
auxiliary factor in
\Cref{lem:V52-linearized-two-defect}, and keep the \(e_1,e_2\) defects
as the selected pair.  If surplus graph factors are present, multiply
them into the auxiliary.  By
\Cref{lem:V51-surplus-admissible}, their combined modulus is at most
one and their output first moment is bounded.  Hence the combined
auxiliary size is at most a fixed multiple of \(a(S_0)\).

For \(S_0\in\Omega\),
\eqref{eq:V52-defect-weight-cap} and the aggregate trace tax give
\[
 \sum_{(S_0,\mu)\in\Omega}
 a(S_0)\|\mathcal C_{S_0,\mu}(A,B)\|_1
 \le
 A_*\tau_\Omega\|A\|_1\|B\|_1.
\]
All subsequent intermediate and final blocks sum contractively by raw
pinching.  The linearized two-defect estimate therefore bounds the
normalized contribution by
\[
 Cs_\alpha h_1h_2\tau_\Omega.
\]
If the second entry in the minimum of
\eqref{eq:V52-small-family-threshold} is at most one, this is exactly
bounded by the right side of
\eqref{eq:V52-dimension-small-star}.  If it exceeds one, then
\(s_\alpha\le h_0/\Xi_c^2\), and the same conclusion follows from
\(\tau_\Omega\le1\).
\end{proof}

\begin{corollary}[Total-copy dimension criterion]
\label{cor:V52-total-dimension-star}
The hypothesis of
\Cref{thm:V52-dimension-small-star} holds whenever
\begin{equation}
 D_\Omega
 \le
 \lambda_0\max\{d_{R_0},d_{Q_0}\}.
 \label{eq:V52-total-dimension-criterion}
\end{equation}
In particular the complete long dual-singlet sector of V51 lies in the
dimension-small branch.
\end{corollary}

\begin{proof}
The minimum of
\(D_\Omega/d_{R_0}\) and \(D_\Omega/d_{Q_0}\) is
\(D_\Omega/\max\{d_{R_0},d_{Q_0}\}\).  The singlet claim follows from
\(D_\Omega=1\), \(d_{R_0}=d_{Q_0}=3^N\), and
\(\lambda_0\asymp(s_\alpha N)^{-1}\) in the regression regime.
\end{proof}

\begin{theorem}[Capacity envelope implies the full star gate]
\label{thm:V52-capacity-implies-star}
If at one selected star edge \(e_0\),
\begin{equation}
 \boxed{
 \mathcal K_\alpha(R_0,Q_0)
 \le C\frac{h_0}{s_\alpha\Xi_c^2},}
 \label{eq:V52-star-capacity-gate}
\end{equation}
then that routed graph block satisfies the coefficient-sensitive star
gate \eqref{eq:V51-coefficient-star-gate}.
\end{theorem}

\begin{proof}
Apply \Cref{thm:V52-layer-cake-capacity} to the first-resolved
\(e_0\) bank and
\Cref{lem:V52-linearized-two-defect} to the other two selected edges.
As in \Cref{thm:V52-dimension-small-star}, surplus defects enlarge the
auxiliary size by only a universal factor.  Raw pinching sums every
remaining block.  The result is
\[
 Cs_\alpha h_1h_2\mathcal K_\alpha(R_0,Q_0)
\le C\frac{h_0h_1h_2}{\Xi_c^2}.
\]
\end{proof}

\section{Removal of the hot-output part of the capacity}
\label{sec:V52-hot-output}

On the selected pair edge \(e_0=c\ell\), abbreviate
\begin{equation}
 x:=\Xi(R_0),\qquad y:=\Xi(Q_0),\qquad
 z(S):=\Xi(S),\qquad h:=h_0.
 \label{eq:V52-pair-masses}
\end{equation}
Exact activity and fusion imply
\begin{equation}
 h\ge\max\{x,y\},
 \qquad
 z(S)\le C(x+y)\le Ch.
 \label{eq:V52-pair-mass-geometry}
\end{equation}

\begin{lemma}[The input-product part already has the star scale]
\label{lem:V52-product-part}
Assume
\begin{equation}
 x\ge\varepsilon\Xi_c,
 \qquad
 s_\alpha x>M.
 \label{eq:V52-selected-edge-dominance}
\end{equation}
Uniformly in the resolved output \(S\),
\begin{equation}
 (1+s_\alpha z(S))q_{\alpha,R_0}q_{\alpha,Q_0}
 \le
 C_{\varepsilon,M}\frac{h}{s_\alpha\Xi_c^2}.
 \label{eq:V52-product-part}
\end{equation}
\end{lemma}

\begin{proof}
By the product second moment,
\[
 q_{\alpha,R_0}\le\frac C{(s_\alpha x)^2}.
\]
The constant and \(s_\alpha x\) parts of
\((1+s_\alpha z)q_{R_0}q_{Q_0}\), using
\eqref{eq:V52-pair-mass-geometry}, are bounded respectively by
\[
 \frac C{s_\alpha^2x^2},
 \qquad
 \frac C{s_\alpha x}.
\]
The first is at most \(Ch/(s_\alpha x^2)\) because
\(s_\alpha h\ge s_\alpha x>M\); the second has the same bound because
\(h\ge x\).  For the \(s_\alpha y\) contribution use
\(s_\alpha yq_{\alpha,Q_0}\le C\), retain \(q_{\alpha,R_0}\), and use
the first comparison.  Finally
\(x\ge\varepsilon\Xi_c\).
\end{proof}

\begin{lemma}[Every output-hot pair block has the star scale]
\label{lem:V52-output-hot}
Under \eqref{eq:V52-selected-edge-dominance}, fix
\(\delta>0\).  If
\begin{equation}
 z(S)\ge\delta x,
 \label{eq:V52-output-hot-condition}
\end{equation}
then
\begin{equation}
 a(S)\le
 C_{\varepsilon,\delta,M}
 \frac{h}{s_\alpha\Xi_c^2}.
 \label{eq:V52-output-hot-weight}
\end{equation}
\end{lemma}

\begin{proof}
Use
\[
 |\delta(S)|\le q_{\alpha,S}
 +q_{\alpha,R_0}q_{\alpha,Q_0}.
\]
The second contribution, including its \(s_\alpha z(S)\) weight, is
\Cref{lem:V52-product-part}.  For the first, the product second moment
and \(s_\alpha z(S)\ge\delta M\) give
\[
 (1+s_\alpha z(S))q_{\alpha,S}
 \le\frac C{s_\alpha z(S)}
 \le\frac {C_\delta}{s_\alpha x}
 \le C_{\varepsilon,\delta}
 \frac h{s_\alpha\Xi_c^2},
\]
where the last step uses \(h\ge x\) and
\(\Xi_c\le x/\varepsilon\).
\end{proof}

Let
\begin{equation}
 \Omega_{\rm cold}(\delta)
 :=
 \{(S,\mu):z(S)<\delta x\},
 \qquad
 D_{\rm cold}(\delta):=
 \sum_{(S,\mu)\in\Omega_{\rm cold}(\delta)}d_S,
 \label{eq:V52-cold-family}
\end{equation}
and define its one-sided capacity
\begin{equation}
 \tau_{\rm cold}(\delta)
 :=
 \min\left\{
 \frac{D_{\rm cold}(\delta)}{d_{R_0}},
 \frac{D_{\rm cold}(\delta)}{d_{Q_0}},1
 \right\}.
 \label{eq:V52-cold-capacity}
\end{equation}

\begin{theorem}[Exact cold-capacity reduction of the star]
\label{thm:V52-cold-capacity-reduction}
Under \eqref{eq:V52-selected-edge-dominance},
\begin{equation}
 \mathfrak F_{R_0,Q_0}(A,B)
 \le
 C_{\varepsilon,\delta,M}
 \left[
 \frac h{s_\alpha\Xi_c^2}
 +\tau_{\rm cold}(\delta)
 \right]\|A\|_1\|B\|_1.
 \label{eq:V52-cold-capacity-functional}
\end{equation}
Consequently the coefficient-sensitive star gate holds if
\begin{equation}
 \boxed{
 \tau_{\rm cold}(\delta)
 \le C_{\varepsilon,\delta}
 \frac h{s_\alpha\Xi_c^2}.}
 \label{eq:V52-cold-capacity-gate}
\end{equation}
If \(h\ge c s_\alpha\Xi_c^2\), the gate is automatic; hence only the
overlap-light, output-cold, one-sided-dimension-mesoscopic sector
remains.
\end{theorem}

\begin{proof}
On the complement of \(\Omega_{\rm cold}(\delta)\),
\Cref{lem:V52-output-hot} and raw pinching bound the weighted
compression sum by the first term on the right of
\eqref{eq:V52-cold-capacity-functional}.  On
\(\Omega_{\rm cold}(\delta)\), use
\eqref{eq:V52-defect-weight-cap} and
\Cref{thm:V52-aggregate-trace-tax}; this gives
\(A_*\tau_{\rm cold}(\delta)\).
The implication for the star is
\Cref{thm:V52-capacity-implies-star}, or directly
\Cref{lem:V52-linearized-two-defect}.  Finally, if
\(h\ge cs_\alpha\Xi_c^2\), the right side of
\eqref{eq:V52-cold-capacity-gate} is bounded below by a positive
constant while \(\tau_{\rm cold}\le1\).
\end{proof}

\begin{assessment}[First genuinely unresolved V52 object]
\label{ass:V52-first-open-object}
All multiplier-only parts of the edge-concentrated star have now been
removed.  The first open quantity is the deterministic one-sided
carrier capacity
\[
 \frac{D_{\rm cold}(\delta)}
      {\max\{d_{R_0},d_{Q_0}\}}
\]
for product-group outputs whose total Casimir mass is much smaller than
the selected hot input mass.  It must be compared with
\(h/(s_\alpha\Xi_c^2)\).  This is not the withdrawn squared
tensor-dimension probability: it is exactly the aggregate partial-trace
tax seen by arbitrary coefficient matrices.
\end{assessment}

\section{Complete fundamental--dual cold-capacity test}
\label{sec:V52-fundamental-dual}

\begin{theorem}[Homogeneous fundamental--dual star capacity]
\label{thm:V52-fundamental-dual-capacity}
Let the selected edge \(e_0\) contain \(N\) coordinates with
\[
 R_{0,e}=\mathbf3,\qquad Q_{0,e}=\overline{\mathbf3},
\]
and suppose \(x\ge\varepsilon\Xi_c\) and
\(s_\alpha x>M\).  Then
\eqref{eq:V52-cold-capacity-gate} holds for a fixed
\(\delta_0>0\).  Consequently every arbitrary-coefficient star graph
routed through this edge satisfies
\eqref{eq:V51-coefficient-star-gate}, uniformly in \(N\).
\end{theorem}

\begin{proof}
At one coordinate,
\[
 \mathbf3\otimes\overline{\mathbf3}
 =\mathbf1\oplus\mathbf8,
 \qquad
 d_{\mathbf1}=1,\quad d_{\mathbf8}=8.
\]
The product-group outputs are therefore indexed by subsets of octet
coordinates.  If \(k\) coordinates carry \(\mathbf8\), the output-copy
dimension is \(8^k\), its multiplicity is one, and its balanced mass is
\(4k\).  Meanwhile
\[
 d_{R_0}=d_{Q_0}=3^N,\qquad
 x=\frac73N,\qquad h=\frac{49}{9}N.
\]

Choose, for example, \(\delta_0=1/10\).  The cold condition
\(4k<\delta_0x\) implies
\[
 k\le\rho N,\qquad \rho:=\frac7{120}<\frac12.
\]
Hence
\begin{align}
 \tau_{\rm cold}(\delta_0)
 &\le
 3^{-N}\sum_{k\le\rho N}\binom Nk8^k
 \notag\\
 &\le
 (N+1)\exp\left\{
 N\bigl[\mathsf H(\rho)+\rho\log8-\log3\bigr]
 \right\}
 \le Ce^{-cN}.
 \label{eq:V52-fundamental-dual-cold-capacity}
\end{align}
Here
\(\mathsf H(\rho)=-\rho\log\rho-(1-\rho)\log(1-\rho)\), and the
displayed exponent is strictly negative at \(\rho=7/120\).

Edge dominance gives \(\Xi_c\le C_\varepsilon N\), so
\[
 \frac h{s_\alpha\Xi_c^2}
 \ge\frac {c_\varepsilon}{s_\alpha N}.
\]
Since \(s_\alpha\le1\),
\(e^{-cN}\le C/N\le C/(s_\alpha N)\).  This proves
\eqref{eq:V52-cold-capacity-gate}.  Apply
\Cref{thm:V52-cold-capacity-reduction}.  Surplus edges are already
included there through
\Cref{lem:V51-surplus-admissible}.
\end{proof}

\begin{remark}[A neighborhood, not one exceptional channel]
\label{rem:V52-dual-neighborhood}
\Cref{thm:V52-fundamental-dual-capacity} sums every output with up to a
fixed positive fraction of octet coordinates.  It therefore controls
an exponentially large cold family and not only the all-singlet
regression of V51.  The normalization \(3^{-N}\) is the one-sided
aggregate Schur tax; the tensor-dimension probability would instead
have denominator \(9^N\) and is neither asserted nor needed.
\end{remark}

\section{Exact one-sided Chernoff factorization}
\label{sec:V52-Chernoff-capacity}

For a one-coordinate output put
\[
 \xi(S):=
 \begin{cases}
 0,&S=\mathbf1,\\
 1+C_2(S),&S\ne\mathbf1.
 \end{cases}
\]
Let \(N_{R_eQ_e}^{S}\) be the LR multiplicity.  If
\(d_R=\prod_ed_{R_e}\ge d_Q=\prod_ed_{Q_e}\), define the oriented
one-link capacity transform
\begin{equation}
 \mathcal Z_e^{R}(\theta)
 :=
 \frac1{d_{R_e}}
 \sum_S N_{R_eQ_e}^{S}d_S
 e^{-\theta\xi(S)}.
 \label{eq:V52-oriented-capacity-transform}
\end{equation}
If \(d_Q>d_R\), define \(\mathcal Z_e^Q\) by dividing by \(d_{Q_e}\).

\begin{theorem}[Product-group one-sided Chernoff bound]
\label{thm:V52-one-sided-Chernoff}
For every \(\theta>0\), the cold capacity in
\eqref{eq:V52-cold-capacity} obeys
\begin{align}
 \frac{D_{\rm cold}(\delta)}{d_R}
 &\le
 e^{\theta\delta x}
 \prod_e\mathcal Z_e^R(\theta),
 && d_R\ge d_Q,
 \label{eq:V52-Chernoff-R}\\
 \frac{D_{\rm cold}(\delta)}{d_Q}
 &\le
 e^{\theta\delta x}
 \prod_e\mathcal Z_e^Q(\theta),
 && d_Q>d_R.
 \label{eq:V52-Chernoff-Q}
\end{align}
Consequently
\eqref{eq:V52-cold-capacity-gate} follows if the appropriate right
side is at most \(Ch/(s_\alpha\Xi_c^2)\).
\end{theorem}

\begin{proof}
Assume \(d_R\ge d_Q\).  A product-group output copy is a tuple
\((S_e,\mu_e)_e\), has dimension \(\prod_ed_{S_e}\), multiplicity
\(\prod_eN_{R_eQ_e}^{S_e}\), and mass
\(\sum_e\xi(S_e)\).  On the cold set,
\[
 \mathbf1_{\{\sum_e\xi(S_e)<\delta x\}}
 \le
 e^{\theta\delta x}
 e^{-\theta\sum_e\xi(S_e)}.
\]
Multiply by the output-copy dimension, sum all tuples, divide by
\(d_R=\prod_ed_{R_e}\), and factor the resulting product.  This gives
\eqref{eq:V52-Chernoff-R}.  The other orientation is identical.
Finally
\[
 \tau_{\rm cold}(\delta)
 =
 \min\left\{
 \frac{D_{\rm cold}(\delta)}{d_R},
 \frac{D_{\rm cold}(\delta)}{d_Q},1
 \right\}
 =
 \min\left\{
 \frac{D_{\rm cold}(\delta)}
      {\max\{d_R,d_Q\}},1
 \right\},
\]
which proves the last assertion.
\end{proof}

\begin{assessment}[Local representation theorem now required]
\label{ass:V52-local-capacity-gate}
The remaining star problem has been reduced to the explicit local
transforms \eqref{eq:V52-oriented-capacity-transform}.  A sufficient
next theorem is a uniform \(\mathrm{SU}(3)\) cancellation estimate
which makes
\[
 \theta\delta x+\sum_e\log\mathcal Z_e^{R}(\theta)
\quad\hbox{or}\quad
 \theta\delta x+\sum_e\log\mathcal Z_e^{Q}(\theta)
\]
negative by at least
\(\log(s_\alpha\Xi_c^2/h)\) on every overlap-light,
selected-edge-dominant bank.  Unlike the withdrawn V39/V45 transfer,
this is a deterministic LR carrier-dimension identity derived from the
actual aggregate Schur projector.
\end{assessment}

\section{V52 ledger and next acquisition}
\label{sec:V52-ledger}

\begin{theorem}[Integrated V52 statement]
\label{thm:V52-integrated}
Preserving V1--V51, V52 proves:
\begin{enumerate}[label=\textup{(V52.\arabic*)},leftmargin=5.5em]
\item the exact aggregate partial-trace identities
      \eqref{eq:V52-aggregate-partial-traces};
\item the spectator-amplified trace tax
      \eqref{eq:V52-aggregate-trace-tax} for an arbitrary family of
      multiplicity-resolved pair outputs;
\item the linearized auxiliary two-defect response
      \eqref{eq:V52-linearized-two-defect};
\item the layer-cake coefficient-capacity bound
      \eqref{eq:V52-layer-cake-capacity};
\item arbitrary-coefficient closure of every aggregate
      dimension-small star family satisfying
      \eqref{eq:V52-small-family-threshold};
\item pointwise removal of the input-product and output-hot parts of a
      selected edge;
\item reduction of the edge-concentrated star to the one-sided cold
      carrier inequality \eqref{eq:V52-cold-capacity-gate};
\item complete closure of the homogeneous
      \(\mathbf3^{\boxtimes N}\)--%
      \(\overline{\mathbf3}^{\boxtimes N}\) star capacity, including
      every output with a fixed positive density of octets; and
\item the exact one-sided Chernoff factorization
      \eqref{eq:V52-Chernoff-R}--\eqref{eq:V52-Chernoff-Q}.
\end{enumerate}
\end{theorem}

\begin{proof}
Items \textup{(V52.1)--(V52.2)} are
\Cref{lem:V52-aggregate-partial-traces,
thm:V52-aggregate-trace-tax}.  Item \textup{(V52.3)} is
\Cref{lem:V52-linearized-two-defect}.  Item \textup{(V52.4)} is
\Cref{thm:V52-layer-cake-capacity}.  Item \textup{(V52.5)} is
\Cref{thm:V52-dimension-small-star,
cor:V52-total-dimension-star}.  Items
\textup{(V52.6)--(V52.7)} are
\Cref{lem:V52-product-part,lem:V52-output-hot,
thm:V52-cold-capacity-reduction}.  Item \textup{(V52.8)} is
\Cref{thm:V52-fundamental-dual-capacity}.  Item
\textup{(V52.9)} is \Cref{thm:V52-one-sided-Chernoff}.
\end{proof}

\begingroup
\small
\begin{longtable}{>{\raggedright\arraybackslash}p{0.29\textwidth}
                  >{\raggedright\arraybackslash}p{0.15\textwidth}
                  >{\raggedright\arraybackslash}p{0.48\textwidth}}
\toprule
\textbf{V52 row}&\textbf{Status}&\textbf{Advance or obstruction}\\
\midrule
\endfirsthead
\toprule
\textbf{V52 row}&\textbf{Status}&\textbf{Advance or obstruction}\\
\midrule
\endhead
Aggregate Schur tax
& \MGclosed
& A whole invariant output family is compressed before trace norm;
  its cost is total carrier dimension, with no copy count and arbitrary
  spectator amplification.\\[0.35em]
Linearized two-defect response
& \MGclosed
& The third star edge enters only through its modulus plus its output
  first moment, while the other two defects retain the one final
  resolvent.\\[0.35em]
Dimension-small star sector
& \MGclosed
& The actual aggregate tax supplies the missing center mark and closes
  every family below \eqref{eq:V52-small-family-threshold}.\\[0.35em]
Input-product and output-hot sectors
& \MGclosed
& Product second moments and the resolved output mass give the required
  star scale without a formation estimate.\\[0.35em]
Homogeneous fundamental--dual bank
& \MGclosed
& The one-sided cold capacity is a binomial entropy tail with denominator
  \(3^N\), closing all arbitrary coefficient blocks in this model.\\[0.35em]
General output-cold carrier capacity
& \MGreduced
& It is the explicit deterministic inequality
  \eqref{eq:V52-cold-capacity-gate}, equivalently the local-transform
  product in \eqref{eq:V52-Chernoff-R} or
  \eqref{eq:V52-Chernoff-Q}.\\[0.35em]
Coefficient-sensitive star and path gates
& \MGopen
& The general star awaits the SU(3) local capacity theorem; the path
  remains downstream of that result.\\[0.35em]
Global quartic and all-order MTA
& \MGopen
& The complete pair row and the \([4]\), \([3,2]\), and \([3,3]\)
  skeletons are not yet closed.\\[0.35em]
Full Yang--Mills mass gap
& \MGopen
& WUFC/CPM, nonlinear construction, finite-edge margins, continuum and
  OS positivity/nontriviality, and the spectral passage remain
  downstream.\\
\bottomrule
\end{longtable}
\endgroup

\begin{openproblem}[V53 mandate: local SU(3) one-sided cold capacity]
\label{op:V53}
\begin{enumerate}[label=\textup{(V53.\arabic*)},leftmargin=5.5em]
\item Work directly with
      \(\mathcal Z_e^R(\theta)\) and
      \(\mathcal Z_e^Q(\theta)\) in
      \eqref{eq:V52-oriented-capacity-transform}.  Use the exact
      SU(3) LR/BZ polytope and Weyl dimensions; do not replace the
      one-sided denominator by \(d_Rd_Q\).
\item Prove a local cancellation inequality which distinguishes
      output-preserving channels from channels that remove a fixed
      fraction of the larger input Casimir.  Its tensor product must
      make the Chernoff exponent in
      \Cref{ass:V52-local-capacity-gate} negative.
\item Respect the global orientation:
      \(d_R\ge d_Q\) need not imply
      \(d_{R_e}\ge d_{Q_e}\) at every coordinate.  Any local
      reorientation cost must be retained and summed.
\item Split off overlap-heavy coordinates first.  On the remaining
      overlap-light bank, use
      \(h\ll s_\alpha\Xi_c^2\) and selected-edge dominance to force
      enough diffuse cancellation coordinates for a capacity gain.
\item Test the local inequality on
      \((n,0)\otimes(0,n)\), highly asymmetric pairs, boundary BZ
      intervals, and the homogeneous fundamental--dual family.  A
      squared dimension-law tail is not a substitute for any test.
\item Once \eqref{eq:V52-cold-capacity-gate} is proved, invoke
      \Cref{thm:V52-capacity-implies-star} to close the general star,
      then formulate the corresponding two-internal-vertex capacity
      for the path.
\item Keep the \([4]\), \([3,2]\), and \([3,3]\) quartic skeletons
      suspended until the exclusive-pair row is complete.
\item Preserve only the newest YM file.  No companion audit is due;
      the next scheduled read-only SM/NS audit remains V55.
\end{enumerate}
\end{openproblem}

\begin{firewall}[End-of-V52 claim boundary]
\label{fire:V52-claim-boundary}
V52 proves the aggregate coefficient tax, the dimension-small and
output-hot star sectors, the full homogeneous fundamental--dual
capacity, and the exact Chernoff reduction to local SU(3) transforms.
It does not prove the general cold-capacity inequality, the complete
coefficient-sensitive star or path gate, the exclusive-pair quartic
row, any other quartic skeleton, global \(m=4\), all-order MTA, WUFC,
CPM, the nonlinear Perron packet, finite-edge margins, capacity in the
constructive continuum sense, protected-sector floors, local-algebra
noncollapse, reflection positivity, a nontrivial continuum OS state,
or a positive Yang--Mills spectral gap.  The full Yang--Mills mass gap
is not proved.
\end{firewall}

\section*{Version V52 append-only record}
\addcontentsline{toc}{section}{Version V52 append-only record}
The complete V51 source was retained before this continuation.  Its
SHA--256 digest was
\begin{center}\scriptsize\ttfamily
7b19177662e8490e5652ba1b651976536ae975256abaf71f3a2dbe17f33af7cf.
\end{center}
No mathematical content of V51 was altered.  On the next iteration,
retain this full file, append before its unique active document-closing
command, increment the filename to \texttt{\_V53.tex}, and remove only
the superseded V52 file from this Yang--Mills note series.

% =============================================================================
% END APPENDED VERSION V52 MATERIAL
% =============================================================================

% =============================================================================
% BEGIN APPENDED VERSION V53 MATERIAL
% =============================================================================

\chapter{V53: Effective-number cold capacity and bounded local closure}
\label{ch:V53}

\section{The correct cold scale}
\label{sec:V53-effective-scale}

Retain the selected star edge notation of V52:
\[
 x=\sum_{e\in L}r_e,\qquad
 h=\sum_{e\in L}r_eq_e,
 \qquad
 r_e:=A_{R_e},\quad q_e:=A_{Q_e}.
\]
Define its effective cancellation number by
\begin{equation}
 N_{\rm eff}:=\frac{x^2}{h}.
 \label{eq:V53-effective-number}
\end{equation}
If \(x\ge\varepsilon\Xi_c\), then
\begin{equation}
 \frac{\varepsilon^2}{s_\alpha N_{\rm eff}}
 \le
 \frac h{s_\alpha\Xi_c^2}
 \le
 \frac1{s_\alpha N_{\rm eff}}.
 \label{eq:V53-star-scale-effective}
\end{equation}
Thus \(N_{\rm eff}\), rather than \(x\), is the mass scale at which an
output multiplier begins to pay the star gate.

For a fixed \(\kappa>0\), set
\begin{align}
 \Omega_{\rm eff}(\kappa)
 &:=
 \{(S,\mu):\Xi(S)<\kappa N_{\rm eff}\},
 \label{eq:V53-effective-cold-family}\\
 D_{\rm eff}(\kappa)
 &:=
 \sum_{(S,\mu)\in\Omega_{\rm eff}(\kappa)}d_S,
 \label{eq:V53-effective-cold-dimension}\\
 \tau_{\rm eff}(\kappa)
 &:=
 \min\left\{
 \frac{D_{\rm eff}(\kappa)}{d_R},
 \frac{D_{\rm eff}(\kappa)}{d_Q},1
 \right\}.
 \label{eq:V53-effective-cold-capacity}
\end{align}

\begin{theorem}[Effective-number cold-capacity reduction]
\label{thm:V53-effective-cold-reduction}
Assume \(x\ge\varepsilon\Xi_c\).  There are fixed
\(\kappa,M>0\) such that:
\begin{enumerate}[label=\textup{(\roman*)},leftmargin=3.7em]
\item if \(s_\alpha N_{\rm eff}\le M\), the selected-edge contribution
      already has the star scale;
\item if \(s_\alpha N_{\rm eff}>M\), then
      \begin{equation}
      \mathfrak F_{R,Q}(A,B)
      \le
      C_{\varepsilon,\kappa}
      \left[
      \frac1{s_\alpha N_{\rm eff}}
      +\tau_{\rm eff}(\kappa)
      \right]\|A\|_1\|B\|_1.
      \label{eq:V53-effective-functional}
      \end{equation}
\end{enumerate}
Consequently the coefficient-sensitive star gate follows from
\begin{equation}
 \boxed{
 \tau_{\rm eff}(\kappa)
 \le\frac C{s_\alpha N_{\rm eff}}.}
 \label{eq:V53-effective-capacity-gate}
\end{equation}
\end{theorem}

\begin{proof}
Equation \eqref{eq:V53-star-scale-effective} and the universal bound
\eqref{eq:V52-defect-weight-cap}, followed by raw pinching, prove
\textup{(i)}.

Assume \(s_\alpha N_{\rm eff}>M\).  Since every \(q_e\ge1\),
\(N_{\rm eff}\le x\), so the product-multiplier part is bounded by
\Cref{lem:V52-product-part}.  If
\(\Xi(S)\ge\kappa N_{\rm eff}\), the product second moment gives
\[
 (1+s_\alpha\Xi(S))q_{\alpha,S}
 \le\frac C{s_\alpha\Xi(S)}
 \le\frac {C_\kappa}{s_\alpha N_{\rm eff}}.
\]
Use raw pinching on this output family.  On
\(\Omega_{\rm eff}(\kappa)\), use the uniform defect-weight cap and the
aggregate tax
\Cref{thm:V52-aggregate-trace-tax}.  This proves
\eqref{eq:V53-effective-functional}.  Combine
\eqref{eq:V53-effective-capacity-gate} with
\Cref{lem:V52-linearized-two-defect} and
\eqref{eq:V53-star-scale-effective}.
\end{proof}

\section{Reverse fusion removes locally asymmetric coordinates}
\label{sec:V53-reverse-fusion}

\begin{lemma}[Casimir reverse-fusion inequality]
\label{lem:V53-Casimir-reverse-fusion}
There is a group constant \(C_{\rm rf}\) such that, whenever
\(N_{RQ}^{S}>0\),
\begin{equation}
 \sqrt{\xi(R)}
 \le
 C_{\rm rf}\bigl(\sqrt{\xi(S)}+\sqrt{\xi(Q)}\bigr),
 \label{eq:V53-Casimir-reverse-fusion}
\end{equation}
and likewise with \(R,Q\) interchanged.
\end{lemma}

\begin{proof}
Frobenius reciprocity gives
\[
 \operatorname{Hom}_G(S,R\otimes Q)
 \cong
 \operatorname{Hom}_G(R,S\otimes Q^*).
\]
Thus \(R\) occurs in \(S\otimes Q^*\).  The tensor-product
highest-weight bound, together with equivalence between the Euclidean
highest-weight norm and \(\sqrt{1+C_2}\) on the fixed rank-two Weyl
chamber, gives \eqref{eq:V53-Casimir-reverse-fusion}.  Duality preserves
the Casimir.
\end{proof}

Choose \(\eta>0\) sufficiently small in terms of
\(C_{\rm rf}\), and put
\begin{equation}
 \mathcal A_\eta:=\{e:q_e\le\eta r_e\}.
 \label{eq:V53-asymmetric-coordinates}
\end{equation}

\begin{theorem}[Asymmetric-mass star closure]
\label{thm:V53-asymmetric-star}
Fix \(\beta>0\).  If
\begin{equation}
 \sum_{e\in\mathcal A_\eta}r_e\ge\beta x,
 \label{eq:V53-asymmetric-mass}
\end{equation}
then, for \(\kappa>0\) chosen sufficiently small in terms of
\(\beta\), \(\Omega_{\rm eff}(\kappa)\) is empty.  Hence the full
arbitrary-coefficient routed star block satisfies its marked-tree
bound.
\end{theorem}

\begin{proof}
For \(e\in\mathcal A_\eta\),
\Cref{lem:V53-Casimir-reverse-fusion} and the choice of \(\eta\) imply
\(\xi(S_e)\ge c r_e\) for every admissible local output.  Therefore
\[
 \Xi(S)=\sum_e\xi(S_e)
 \ge c\sum_{e\in\mathcal A_\eta}r_e
 \ge c\beta x.
\]
Also \(h\ge x\), so \(N_{\rm eff}=x^2/h\le x\).
Taking \(\kappa<c\beta\) makes
\(\Xi(S)\ge\kappa N_{\rm eff}\) for every output.  Apply
\Cref{thm:V53-effective-cold-reduction}.
\end{proof}

\section{Uniformly bounded local representations}
\label{sec:V53-bounded-local}

\begin{lemma}[Finite-height contraction of the oriented transform]
\label{lem:V53-finite-height-transform}
Fix \(L<\infty\).  There exist
\(\theta_L<\infty\) and \(\gamma_L<1\) such that, for all nontrivial
\(\mathrm{SU}(3)\) irreducibles \(R,Q\) with
\(A_R,A_Q\le L\),
\begin{equation}
 \mathcal Z^{R}_{R,Q}(\theta_L)\le\gamma_L,
 \qquad
 \mathcal Z^{Q}_{R,Q}(\theta_L)\le\gamma_L,
 \label{eq:V53-finite-height-transform}
\end{equation}
where the two transforms use respectively the denominators \(d_R\)
and \(d_Q\) as in
\eqref{eq:V52-oriented-capacity-transform}.
\end{lemma}

\begin{proof}
There are finitely many such pairs.  For a fixed pair,
\[
 \lim_{\theta\to\infty}\mathcal Z^R_{R,Q}(\theta)
 =
 \frac{N_{RQ}^{\mathbf1}}{d_R}.
\]
The trivial representation occurs only when \(Q\cong R^*\), and then
with multiplicity one.  Since a nontrivial SU(3) irreducible has
dimension at least three, the limit is at most \(1/3\).  The same
argument applies to the \(Q\)-oriented transform.  Choose one
\(\theta_L\) large enough for the finite family and then choose
\(\gamma_L\in(1/3,1)\) above all its values.
\end{proof}

\begin{lemma}[Effective-cold Chernoff bound]
\label{lem:V53-effective-Chernoff}
If \(d_R\ge d_Q\), then for every \(\theta>0\),
\begin{equation}
 \frac{D_{\rm eff}(\kappa)}{d_R}
 \le
 e^{\theta\kappa N_{\rm eff}}
 \prod_e\mathcal Z_e^R(\theta).
 \label{eq:V53-effective-Chernoff-R}
\end{equation}
If \(d_Q>d_R\), the analogous estimate holds with denominator \(d_Q\)
and transforms \(\mathcal Z_e^Q\).
\end{lemma}

\begin{proof}
Repeat \Cref{thm:V52-one-sided-Chernoff} with the cold threshold
\(\delta x\) replaced by \(\kappa N_{\rm eff}\).  The product-group
dimension and LR multiplicity still factor exactly.
\end{proof}

\begin{theorem}[Bounded-local star capacity]
\label{thm:V53-bounded-local-star}
Fix \(L<\infty\).  Suppose on the selected edge bank
\begin{equation}
 A_{R_e}\le L,\qquad A_{Q_e}\le L
 \quad\hbox{for every \(e\)}.
 \label{eq:V53-bounded-local-hypothesis}
\end{equation}
If \(x\ge\varepsilon\Xi_c\), then the complete
arbitrary-coefficient routed star block satisfies
\eqref{eq:V51-coefficient-star-gate}, uniformly in the number of
coordinates.
\end{theorem}

\begin{proof}
Let \(N\) be the number of selected-edge coordinates and let
\(a_F=1+C_2(\mathbf3)\).  Exact support gives
\[
 a_FN\le x\le LN,
 \qquad
 a_F^2N\le h\le L^2N.
\]
In particular
\begin{equation}
 N\ge c_LN_{\rm eff}.
 \label{eq:V53-count-controls-effective}
\end{equation}
Use \(\theta_L,\gamma_L\) from
\Cref{lem:V53-finite-height-transform}.  In either global dimension
orientation,
\Cref{lem:V53-effective-Chernoff} gives
\[
 \tau_{\rm eff}(\kappa)
 \le
 \exp\{\theta_L\kappa N_{\rm eff}\}\gamma_L^N.
\]
Choose \(\kappa>0\) so small that
\(\theta_L\kappa\le
\tfrac12c_L|\log\gamma_L|\).  Then
\begin{equation}
 \tau_{\rm eff}(\kappa)
 \le Ce^{-c_L'N_{\rm eff}}.
 \label{eq:V53-bounded-local-capacity}
\end{equation}
If \(s_\alpha N_{\rm eff}\le M\), use part \textup{(i)} of
\Cref{thm:V53-effective-cold-reduction}.  Otherwise
\[
 e^{-c_L'N_{\rm eff}}
 \le\frac C{N_{\rm eff}}
 \le\frac C{s_\alpha N_{\rm eff}},
\]
because \(s_\alpha\le1\).  Thus
\eqref{eq:V53-effective-capacity-gate} holds, and part \textup{(ii)} of
\Cref{thm:V53-effective-cold-reduction} closes the star.
\end{proof}

\begin{corollary}[The V52 fundamental--dual theorem is one instance]
\label{cor:V53-fundamental-instance}
The homogeneous fundamental--dual capacity
\Cref{thm:V52-fundamental-dual-capacity} follows also from
\Cref{thm:V53-bounded-local-star}.  The explicit binomial proof remains
useful because it identifies the actual rate and verifies the one-sided
normalization directly.
\end{corollary}

\begin{assessment}[The unbounded balanced core]
\label{ass:V53-unbounded-balanced-core}
After \Cref{thm:V53-asymmetric-star,thm:V53-bounded-local-star}, a
remaining star obstruction must have all of the following properties:
\begin{enumerate}[label=\textup{(\roman*)},leftmargin=3.7em]
\item \(s_\alpha N_{\rm eff}\) is large;
\item the selected edge is overlap-light:
      \(h/(s_\alpha\Xi_c^2)\ll1\);
\item a fixed fraction of the center mass lies on coordinates with
      \(q_e\ge\eta r_e\); and
\item the mass carried by unbounded local highest weights cannot be
      discarded.
\end{enumerate}
Thus the next theorem must control
\(\mathcal Z_e^{R/Q}(\theta)\) uniformly as both local representations
grow comparably.  A finite-height compactness argument no longer
applies, but the required object is still the deterministic one-sided
LR transform, not an arbitrary coefficient-formation law.
\end{assessment}

\section{Universal transform contraction and full star closure}
\label{sec:V53-universal-transform}

\begin{lemma}[A universal one-sided SU(3) transform gap]
\label{lem:V53-universal-transform-gap}
There exist numerical constants \(\theta_*>0\) and
\(\gamma_*<1\) such that, for every pair of nontrivial SU(3)
irreducibles \(R,Q\),
\begin{equation}
 \boxed{
 \mathcal Z^R_{R,Q}(\theta_*)\le\gamma_*,
 \qquad
 \mathcal Z^Q_{R,Q}(\theta_*)\le\gamma_*.}
 \label{eq:V53-universal-transform-gap}
\end{equation}
\end{lemma}

\begin{proof}
Frobenius reciprocity gives
\[
 N_{RQ}^{S}=N_{R^*S}^{Q}.
\]
Since \(Q\) occurs with this multiplicity in \(R^*\otimes S\),
dimension counting yields
\begin{equation}
 N_{RQ}^{S}d_Q\le d_Rd_S.
 \label{eq:V53-reciprocity-dimension}
\end{equation}
Therefore
\begin{equation}
 \mathcal Z^R_{R,Q}(\theta)
 \le
 \frac1{d_Q}
 \sum_{S\in\widehat{\mathrm{SU}(3)}}
 d_S^2e^{-\theta\xi(S)}.
 \label{eq:V53-transform-heat-sum}
\end{equation}
The sum is finite for every \(\theta>0\), by the Weyl dimension and
Casimir formulas, and decreases to one as
\(\theta\to\infty\); only the trivial representation remains.  Because
\(d_Q\ge3\), choose \(\theta_*\) so that the sum is less than \(2\)
and take \(\gamma_*=2/3\).  Interchanging \(R,Q\) proves the second
bound.
\end{proof}

\begin{lemma}[Balanced mass forces enough coordinates]
\label{lem:V53-balanced-coordinate-count}
Let \(\mathcal A_\eta\) be
\eqref{eq:V53-asymmetric-coordinates}.  If
\begin{equation}
 \sum_{e\in\mathcal A_\eta}r_e<\beta x
 \quad\hbox{for some \(\beta<1\),}
 \label{eq:V53-balanced-mass-branch}
\end{equation}
then the number \(N\) of selected-edge coordinates satisfies
\begin{equation}
 N\ge\eta(1-\beta)^2N_{\rm eff}.
 \label{eq:V53-balanced-coordinate-count}
\end{equation}
\end{lemma}

\begin{proof}
On the complementary set \(\mathcal B_\eta\),
\(q_e>\eta r_e\), and
\[
 \sum_{e\in\mathcal B_\eta}r_e>(1-\beta)x.
\]
Consequently, by Cauchy--Schwarz,
\[
 h\ge\eta\sum_{e\in\mathcal B_\eta}r_e^2
 \ge
 \frac{\eta(1-\beta)^2x^2}{|\mathcal B_\eta|}
 \ge
 \frac{\eta(1-\beta)^2x^2}{N}.
\]
Rearrange and use \(N_{\rm eff}=x^2/h\).
\end{proof}

\begin{theorem}[General one-sided effective capacity]
\label{thm:V53-general-effective-capacity}
Assume \(x\ge\varepsilon\Xi_c\).  For a sufficiently small fixed
\(\kappa>0\),
\begin{equation}
 \boxed{
 \tau_{\rm eff}(\kappa)
 \le
 \frac C{s_\alpha N_{\rm eff}}}
 \label{eq:V53-general-effective-capacity}
\end{equation}
whenever \(s_\alpha N_{\rm eff}>M\).  The constants are uniform in the
support, all local highest weights, LR multiplicities, and the global
dimension orientation.
\end{theorem}

\begin{proof}
Fix \(\beta=1/2\) and use the \(\eta\) selected before
\Cref{thm:V53-asymmetric-star}.  If
\(\sum_{\mathcal A_\eta}r_e\ge\beta x\), the effective-cold family is
empty for sufficiently small \(\kappa\), by
\Cref{thm:V53-asymmetric-star}; hence the claim is immediate.

Otherwise
\Cref{lem:V53-balanced-coordinate-count} gives
\(N\ge c_\eta N_{\rm eff}\).  Apply
\Cref{lem:V53-effective-Chernoff} at \(\theta_*\).  Regardless of
whether \(d_R\ge d_Q\) or \(d_Q>d_R\),
\Cref{lem:V53-universal-transform-gap} yields
\[
 \tau_{\rm eff}(\kappa)
 \le
 \exp\{\theta_*\kappa N_{\rm eff}\}\gamma_*^N
 \le
 \exp\left\{
 [\theta_*\kappa-c_\eta|\log\gamma_*|]N_{\rm eff}
 \right\}.
\]
Choose \(\kappa\) so that the bracket is negative.  The result is
\(Ce^{-cN_{\rm eff}}\).  Since \(s_\alpha\le1\),
\[
 e^{-cN_{\rm eff}}
 \le\frac C{N_{\rm eff}}
 \le\frac C{s_\alpha N_{\rm eff}}.
\]
\end{proof}

\begin{theorem}[The coefficient-sensitive quartic star gate is closed]
\label{thm:V53-full-star-gate}
Every routed exclusive-pair star graph satisfies
\eqref{eq:V51-coefficient-star-gate} for arbitrary trace-class Fourier
coefficients, uniformly in supports, local representations, final
outputs, LR multiplicities, and volume.
\end{theorem}

\begin{lemma}[Exterior-hot star payment is surplus-stable]
\label{lem:V53-surplus-exterior-star}
Let a routed star have center \(c\), and suppose
\[
 s_\alpha\Pi_c^\tau>M,\qquad
 \Pi_c^\tau\ge\varepsilon\Xi_c.
\]
Then its marked-tree bound holds even when the routed graph contains
surplus edges.
\end{lemma}

\begin{proof}
A star already contains all three pair edges incident to its center.
Every surplus edge therefore joins two leaves and does not remove the
center exterior singleton multiplier \(b_c^\tau\).  The surplus product
has modulus at most one and has a uniform output first moment by
\Cref{lem:V51-surplus-admissible}.

Repeat the proof of
\Cref{thm:V51-exterior-hot-star}.  Low outputs and selected star-edge
outputs only use the modulus bound on the surplus product.  For an
output on a surplus edge, use its first moment while retaining the
cubic center debit:
\[
 \frac C{s_\alpha}b_c^\tau
 \le
 \frac C{s_\alpha^4(\Pi_c^\tau)^3}
 \le
 \frac {C_M}{s_\alpha^3(\Pi_c^\tau)^2}.
\]
After multiplication by the three selected defect caps, this is
\(CP_\tau/(\Pi_c^\tau)^2\).  Center-private and other spectator outputs
are exactly those treated in
\eqref{eq:V51-exterior-cubic-output}.  This proves the full-center
star payment.
\end{proof}

\begin{proof}[Proof of \Cref{thm:V53-full-star-gate}]
The internally thermal and exterior-hot minimal-tree blocks are
\Cref{thm:V51-thermal-pair-sector,
thm:V51-exterior-hot-star}; the latter payment for surplus graph blocks
is \Cref{lem:V53-surplus-exterior-star}.  In every remaining star block,
the center mass is hot and a fixed fraction lies on its three selected
edge banks.
Choose an edge whose center contribution \(x\) is maximal; then
\(x\ge\varepsilon\Xi_c\) for a numerical \(\varepsilon>0\).

If \(s_\alpha N_{\rm eff}\le M\), use part \textup{(i)} of
\Cref{thm:V53-effective-cold-reduction}.  Otherwise
\Cref{thm:V53-general-effective-capacity} proves
\eqref{eq:V53-effective-capacity-gate}, so part \textup{(ii)} closes
the selected-edge formation--defect functional.  Finally apply
\Cref{lem:V52-linearized-two-defect} to the other two star edges.
Surplus graph defects are admissible by
\Cref{lem:V51-surplus-admissible} and were retained in the auxiliary
size throughout.  Exact normalization and raw pinching give the
arbitrary-coefficient assertion.
\end{proof}

\begin{corollary}[All four pair-only quartic stars are closed]
\label{cor:V53-all-stars}
The four labeled star contributions in the exact graph decomposition
\eqref{eq:V51-routed-pair-decomposition} obey their full quartic
marked-tree terms.  Only the twelve path contributions remain in the
exclusive-pair row.
\end{corollary}

\begin{proof}
Apply \Cref{thm:V53-full-star-gate} at each choice of center.
\end{proof}

\begin{corollary}[Universal one-edge formation--defect tax]
\label{cor:V53-one-edge-tax}
Let one endpoint \(i\) of an exact pair bank contribute mass \(x\) to
that bank, and suppose \(x\ge\varepsilon\Xi_i\).  Then, for arbitrary
spectator-amplified trace-class coefficients,
\begin{equation}
 \boxed{
 \mathfrak F_{R,Q}(A,B)
 \le
 C_\varepsilon\frac h{s_\alpha\Xi_i^2}
 \|A\|_1\|B\|_1.}
 \label{eq:V53-one-edge-tax}
\end{equation}
\end{corollary}

\begin{proof}
If \(s_\alpha N_{\rm eff}\le M\), the right side is bounded below by a
positive constant by
\eqref{eq:V53-star-scale-effective}, and raw pinching plus
\eqref{eq:V52-defect-weight-cap} suffices.  Otherwise combine
\Cref{thm:V53-effective-cold-reduction,
thm:V53-general-effective-capacity} and again use
\eqref{eq:V53-star-scale-effective}.
\end{proof}

\section{Degree-adapted routing of all connected pair graphs}
\label{sec:V53-adapted-routing}

\begin{lemma}[Star-or-leaf-surplus routing]
\label{lem:V53-adapted-routing}
Every connected simple graph \(G\) on four labeled vertices admits a
spanning tree \(\tau_*(G)\) with the following property:
\begin{enumerate}[label=\textup{(\roman*)},leftmargin=3.7em]
\item if \(G\) has a vertex of degree three, \(\tau_*(G)\) is the star
      at the first such vertex;
\item otherwise \(G\) is either a path or a four-cycle; in the latter
      case \(\tau_*(G)\) is obtained by deleting one edge; and
\item every edge of \(G\setminus\tau_*(G)\) joins leaves of
      \(\tau_*(G)\).
\end{enumerate}
\end{lemma}

\begin{proof}
A degree-three vertex supplies all three edges of a spanning star, and
every other edge has both endpoints among its leaves.  If the maximum
degree is at most two, a finite connected graph is a path or a cycle.
On four vertices these have respectively three and four edges.
Deleting one cycle edge produces a spanning path whose deleted edge
joins its two endpoints, which are the leaves.
\end{proof}

\begin{corollary}[Exact adapted connected-graph decomposition]
\label{cor:V53-adapted-decomposition}
The exact pair polynomial also has the decomposition
\begin{equation}
 \mathbf K_4^{(2)}
 =
 \sum_{\tau\in\mathfrak T_4}
 \sum_{\substack{G\in\mathcal C_4\\\tau_*(G)=\tau}}
 \left(\prod_{ij\in E(\tau)}\mathbf D_{ij}\right)
 \mathbf r^*_{G,\tau},
 \label{eq:V53-adapted-decomposition}
\end{equation}
where \(\mathbf r^*_{G,\tau}\) is defined as in
\eqref{eq:V51-auxiliary-graph-factor}.  It has the admissibility
\eqref{eq:V51-surplus-admissible}, and every surplus defect joins two
tree leaves.
\end{corollary}

\begin{proof}
Assign each of the \(38\) terms in
\eqref{eq:V51-pair-connected-graph} to the deterministic tree
\(\tau_*(G)\), factor its three selected defects, and apply
\Cref{lem:V51-surplus-admissible,lem:V53-adapted-routing}.
\end{proof}

\section{Closure of the coefficient-sensitive path gate}
\label{sec:V53-path-closure}

\begin{lemma}[Selected-heavy path payment]
\label{lem:V53-selected-heavy-path}
Let \(\tau\) be a path with internal vertices \(u,v\), and let
\(w\in\{u,v\}\) satisfy
\(\Xi_w=\max\{\Xi_u,\Xi_v\}\).  If a fixed fraction of \(\Xi_w\) lies
on the two selected edges incident to \(w\), then the routed graph block
satisfies
\begin{equation}
 \textup{coefficient contribution}
 \le C\frac{P_\tau}{\Xi_u\Xi_v}.
 \label{eq:V53-selected-heavy-path}
\end{equation}
\end{lemma}

\begin{proof}
Choose an incident selected edge \(e_0\) whose input-\(w\) mass
\(x_0\) is at least a fixed fraction of \(\Xi_w\).  Resolve that edge
first.  The universal one-edge tax gives
\[
 \mathfrak F_{e_0}
 \le C\frac{h_0}{s_\alpha\Xi_w^2}
\]
by \Cref{cor:V53-one-edge-tax}.  Treat this weighted defect, together
with every surplus leaf-edge factor, as the auxiliary in
\Cref{lem:V52-linearized-two-defect}; keep the other two tree defects
selected.  The resulting coefficient contribution is at most
\[
 Cs_\alpha\frac{P_\tau}{h_0}
 \frac{h_0}{s_\alpha\Xi_w^2}
 =
 C\frac{P_\tau}{\Xi_w^2}.
\]
Since \(\Xi_w^2\ge\Xi_u\Xi_v\), this proves
\eqref{eq:V53-selected-heavy-path}.  Surplus factors join the two
leaves by \Cref{lem:V53-adapted-routing}, so their first moment does not
alter either internal mass normalization.
\end{proof}

\begin{lemma}[Surplus-stable exterior path payment]
\label{lem:V53-surplus-exterior-path}
Suppose each hot internal vertex of an adapted routed path has a fixed
exterior singleton fraction, while every other internal vertex is
thermal.  Then the path obeys its coefficient marked-tree bound even
when the graph is a four-cycle.
\end{lemma}

\begin{proof}
For a minimal path this is
\Cref{thm:V51-exterior-hot-path}.  In the four-cycle, the sole surplus
edge joins the two path leaves by
\Cref{lem:V53-adapted-routing}.  Its defect has modulus at most one and
a uniform output first moment.  In the low and selected-output parts of
the V51 proof use the modulus bound.  In the surplus-output part use
its first moment while retaining the quadratic multiplier of each hot
internal exterior channel.  This is exactly the other-input
calculation in
\eqref{eq:V51-exterior-quadratic-output}; with two hot internal
vertices, retain both quadratic multipliers as in the last paragraph
of \Cref{thm:V51-exterior-hot-path}.  No internal singleton factor is
removed by the leaf--leaf surplus edge.
\end{proof}

\begin{lemma}[Mixed exterior/selected path response]
\label{lem:V53-mixed-path-response}
Let \(u,v\) be the internal vertices of an adapted path.  Assume
\begin{align}
 s_\alpha\Pi_u^\tau>M,\qquad
 \Pi_u^\tau&\ge\varepsilon\Xi_u,
 \label{eq:V53-mixed-exterior-u}\\
 s_\alpha\Xi_v>M,\qquad
 Z_v^\tau&\ge\varepsilon\Xi_v.
 \label{eq:V53-mixed-selected-v}
\end{align}
Then the coefficient marked-tree bound holds.
\end{lemma}

\begin{proof}
Choose a selected edge \(e_0\) incident to \(v\) on which the input
mass is at least \(\varepsilon\Xi_v/2\).  Write the other edge weights
as \(h_1,h_2\), so \(P_\tau=h_0h_1h_2\), and let
\(b_u=b_u^\tau\).

Before summing the \(e_0\) pair outputs, the response satisfies the
linearized estimate
\begin{equation}
 \mathcal Q_{\rm path}
 \le
 C\frac{h_1h_2}{\Xi_u}
 \left[
 |\delta_0|+
 s_\alpha\Xi(\boldsymbol S_0)|\delta_0|
 \right].
 \label{eq:V53-exterior-linearized-path}
\end{equation}
To verify it, on low final outputs use the resolvent and the two other
defect caps; the remaining factor is
\(b_us_\alpha h_1h_2|\delta_0|\), and
\[
 s_\alpha b_u
 \le
 \frac C{s_\alpha(\Pi_u^\tau)^2}
 \le\frac {C_M}{\Xi_u}.
\]
On high outputs, the exterior mass is paid by
\eqref{eq:V51-exterior-quadratic-output}.  The \(e_0\)-output is paid
by the second term in brackets in
\eqref{eq:V53-exterior-linearized-path}.  For an output on another
selected edge of weight \(h_j\), take the better of fusion plus the
defect cap and the selected defect first moment.  Its scalar factor is
bounded by
\[
 b_u\min\{(s_\alpha h_j)^2,1\}
 \le C\frac{h_j}{\Xi_u},
\]
because
\[
 \frac{b_u\min\{(s_\alpha h_j)^2,1\}}
      {h_j/\Xi_u}
 \le
 \frac C{s_\alpha\Xi_u}
 \min\{s_\alpha h_j,(s_\alpha h_j)^{-1}\}
 \le C_M.
\]
The adapted surplus edge joins leaves; its first moment is handled like
an exterior output.  This proves
\eqref{eq:V53-exterior-linearized-path}.

Now sum the \(e_0\) coefficient blocks.  By
\Cref{cor:V53-one-edge-tax},
\[
 \mathfrak F_{e_0}
 \le C\frac{h_0}{s_\alpha\Xi_v^2}.
\]
Insert this in
\eqref{eq:V53-exterior-linearized-path}:
\[
 \textup{coefficient contribution}
 \le
 C\frac{P_\tau}
 {s_\alpha\Xi_u\Xi_v^2}
 \le
 C_M\frac{P_\tau}{\Xi_u\Xi_v},
\]
using \(s_\alpha\Xi_v>M\).
\end{proof}

\begin{theorem}[The coefficient-sensitive quartic path gate is closed]
\label{thm:V53-full-path-gate}
Every adapted routed exclusive-pair path graph satisfies
\eqref{eq:V51-coefficient-path-gate} for arbitrary trace-class Fourier
coefficients, uniformly in supports, local representations, final
outputs, LR multiplicities, and volume.
\end{theorem}

\begin{proof}
Let \(u,v\) be the internal vertices and take \(w\) with maximal input
mass.  If both are thermal, use
\Cref{thm:V51-thermal-pair-sector}.  If a fixed fraction of \(\Xi_w\)
lies on its selected incident edges, use
\Cref{lem:V53-selected-heavy-path}.

It remains that \(w\) is exterior-hot.  If the other internal vertex
is thermal or exterior-hot, use
\Cref{lem:V53-surplus-exterior-path}.  Otherwise it is hot and
selected-heavy, so
\Cref{lem:V53-mixed-path-response} applies with \(u=w\).
These alternatives are exhaustive after fixing the thermal threshold
and splitting each nonthermal input between selected and exterior
mass.
\end{proof}

\begin{theorem}[Complete exclusive-pair quartic MTA]
\label{thm:V53-complete-pair-quartic}
Under the exclusive-pair support hypothesis
\eqref{eq:V51-exclusive-pair-partition}, the full fourth cumulant
satisfies the \(m=4\) marked-tree atlas inequality for arbitrary
trace-class Fourier coefficients, uniformly in volume, support,
representations, and multiplicities.
\end{theorem}

\begin{proof}
Use the exact degree-adapted decomposition
\eqref{eq:V53-adapted-decomposition}.  Its star terms are bounded by
\Cref{thm:V53-full-star-gate}; its path terms are bounded by
\Cref{thm:V53-full-path-gate}.  The exact V46 normalization cancels the
output dimension, raw pinching sums every intermediate and final
block, and expanding \(P_\tau/D_\tau\) gives the degree-marked input
seminorms.  There are only \(38\) graph terms.
\end{proof}

\begin{firewall}[Scope of the pair-row theorem]
\label{fire:V53-pair-row-scope}
\Cref{thm:V53-complete-pair-quartic} allows private and
exclusive-pair coordinate banks only.  It does not include a coordinate
active in three or four inputs, or a four-active coordinate whose local
replica partition contains two disjoint pair blocks.  Those belong to
the \([3,2]\), \([3,3]\), or embedded \([4]\) re-entry and are not
silently absorbed into the pair graph formula.
\end{firewall}

\section{V53 ledger and next acquisition}
\label{sec:V53-ledger}

\begin{theorem}[Integrated V53 statement]
\label{thm:V53-integrated}
Preserving V1--V52, V53 proves:
\begin{enumerate}[label=\textup{(V53.\arabic*)},leftmargin=5.5em]
\item the effective cancellation scale
      \(N_{\rm eff}=x^2/h\) and the strengthened cold reduction
      \eqref{eq:V53-effective-functional};
\item the reverse-fusion exclusion of every bank with a fixed
      asymmetric-mass fraction;
\item finite-height transform contraction and bounded-local star
      closure;
\item the universal one-sided SU(3) transform gap
      \eqref{eq:V53-universal-transform-gap}, obtained from Frobenius
      reciprocity and the square-dimension heat sum;
\item the general effective capacity
      \eqref{eq:V53-general-effective-capacity} and the full
      coefficient-sensitive star gate;
\item exact degree-adapted routing in which all surplus edges join
      leaves;
\item the universal squared endpoint formation--defect tax
      \eqref{eq:V53-one-edge-tax};
\item closure of selected-heavy, exterior-hot, and mixed
      exterior/selected path branches; and
\item the complete arbitrary-coefficient exclusive-pair quartic MTA
      theorem \Cref{thm:V53-complete-pair-quartic}.
\end{enumerate}
\end{theorem}

\begin{proof}
Items \textup{(V53.1)--(V53.2)} are
\Cref{thm:V53-effective-cold-reduction,
lem:V53-Casimir-reverse-fusion,thm:V53-asymmetric-star}.
Item \textup{(V53.3)} is
\Cref{lem:V53-finite-height-transform,
thm:V53-bounded-local-star}.  Items
\textup{(V53.4)--(V53.5)} are
\Cref{lem:V53-universal-transform-gap,
lem:V53-balanced-coordinate-count,
thm:V53-general-effective-capacity,
thm:V53-full-star-gate}.  Items
\textup{(V53.6)--(V53.7)} are
\Cref{lem:V53-adapted-routing,
cor:V53-adapted-decomposition,cor:V53-one-edge-tax}.
Item \textup{(V53.8)} is
\Cref{lem:V53-selected-heavy-path,
lem:V53-surplus-exterior-path,
lem:V53-mixed-path-response,
thm:V53-full-path-gate}.  Item \textup{(V53.9)} is
\Cref{thm:V53-complete-pair-quartic}.
\end{proof}

\begingroup
\small
\begin{longtable}{>{\raggedright\arraybackslash}p{0.29\textwidth}
                  >{\raggedright\arraybackslash}p{0.15\textwidth}
                  >{\raggedright\arraybackslash}p{0.48\textwidth}}
\toprule
\textbf{V53 row}&\textbf{Status}&\textbf{Advance or obstruction}\\
\midrule
\endfirsthead
\toprule
\textbf{V53 row}&\textbf{Status}&\textbf{Advance or obstruction}\\
\midrule
\endhead
Effective-number cold scale
& \MGclosed
& Output multipliers pay above \(N_{\rm eff}=x^2/h\); only a genuinely
  lower-mass carrier family needs coefficient compression.\\[0.35em]
Universal oriented transform
& \MGclosed
& Reciprocity bounds every local transform by one SU(3)
  square-dimension heat sum, uniformly in all highest weights.\\[0.35em]
Coefficient-sensitive star gate
& \MGclosed
& Asymmetric coordinates are excluded by reverse fusion; the balanced
  branch has at least \(cN_{\rm eff}\) contracting coordinates.\\[0.35em]
Degree-adapted graph routing
& \MGclosed
& Graphs with a degree-three vertex route to a star; the only surplus
  path graph is a four-cycle with a leaf--leaf surplus edge.\\[0.35em]
Coefficient-sensitive path gate
& \MGclosed
& Squared endpoint tax and exterior quadratic damping cover all
  selected/exterior configurations of the two internal vertices.\\[0.35em]
Complete exclusive-pair quartic row
& \MGclosed
& All \(38\) connected graph terms satisfy the exact \(m=4\)
  degree-marked atlas under the no-triple-intersection hypothesis.\\[0.35em]
\([3,2]\), \([3,3]\), and \([4]\) fibres
& \MGopen
& Coordinates active in at least three rows require intact ternary or
  fourfold operator blocks; they are outside the pair polynomial.\\[0.35em]
Global quartic and all-order MTA
& \MGopen
& The three remaining skeleton classes must be re-entered with the
  corrected coefficient normalization.\\[0.35em]
Full Yang--Mills mass gap
& \MGopen
& WUFC/CPM, nonlinear construction, finite-edge margins, continuum and
  OS positivity/nontriviality, and the spectral passage remain
  downstream.\\
\bottomrule
\end{longtable}
\endgroup

\begin{openproblem}[V54 mandate: exact \([3,2]\) operator fibre]
\label{op:V54}
\begin{enumerate}[label=\textup{(V54.\arabic*)},leftmargin=5.5em]
\item Starting from
      \eqref{eq:V39-connected-replica-formula}, isolate the complete
      \([3,2]\) fibre: one three-row coordinate block and one
      distinct-coordinate pair block attaching the omitted row.
      Retain every redundant local block as an operator spectator.
\item Keep the restored V50 ternary cumulant intact.  Derive a
      linearized ternary response which accepts the attaching binary
      defect through its modulus and output first moment while charging
      the final resolvent only once.
\item Route the two marks inside the triple through the V50
      connected-coordinate operator formula, and route the third mark
      through the attaching pair.  Do not expand the ternary five-term
      cumulant after taking trace norms.
\item Use \eqref{eq:V53-one-edge-tax} if a high attaching endpoint
      forms a cold pair output.  Any coefficient tax must be applied to
      the actual raw pair compression before ternary recoupling.
\item Revalidate every four-input coefficient formula with
      \(d_{\rm in}/d_{\boldsymbol T}\), nested raw pinching, and no
      representation-formation probability.
\item Test a triple containing a long dual-singlet subbank, an
      attaching fundamental leaf, and every choice of which triple
      vertex receives the pair edge.
\item Declare the \([3,2]\) topology closed only after all four omitted
      rows and all three attachment vertices satisfy their
      degree-normalized marks.  Keep \([3,3]\) and \([4]\) suspended.
\item Preserve only the newest YM file.  No companion audit is due;
      the next scheduled read-only SM/NS audit remains V55.
\end{enumerate}
\end{openproblem}

\begin{firewall}[End-of-V53 claim boundary]
\label{fire:V53-claim-boundary}
V53 proves the general coefficient-sensitive star and path gates and
thereby the complete exclusive-pair quartic row.  It does not include
any local replica block of size three or four, and therefore does not
prove the \([3,2]\), \([3,3]\), or \([4]\) topology, global \(m=4\),
all-order MTA, WUFC, CPM, the nonlinear Perron packet, finite-edge
margins, constructive capacity, protected-sector floors,
local-algebra noncollapse, reflection positivity, a nontrivial
continuum OS state, or a positive Yang--Mills spectral gap.  The full
Yang--Mills mass gap is not proved.
\end{firewall}

\section*{Version V53 append-only record}
\addcontentsline{toc}{section}{Version V53 append-only record}
The complete V52 source was retained before this continuation.  Its
SHA--256 digest was
\begin{center}\scriptsize\ttfamily
243e17d7507af7ddd50f2c4ec6221e5e516bdb339050e34b1182f1618842bd52.
\end{center}
No mathematical content of V52 was altered.  On the next iteration,
retain this full file, append before its unique active document-closing
command, increment the filename to \texttt{\_V54.tex}, and remove only
the superseded V53 file from this Yang--Mills note series.

% =============================================================================
% END APPENDED VERSION V53 MATERIAL
% =============================================================================

% =============================================================================
% BEGIN APPENDED VERSION V54 MATERIAL
% =============================================================================

\chapter{V54: The exact \texorpdfstring{\([3,2]\)}{[3,2]} puncture
formula and the sideways promotion gate}
\label{ch:V54}

\section{Purpose and topology convention}
\label{sec:V54-purpose}

V53 closed the complete exclusive-pair quartic row.  The next replica
topology contains a three-row block and a pair block attaching the
omitted row.  There is a genuine assembly issue: once the pair block is
selected, its coordinate is no longer available to the ternary
operator.  Thus the required object is a \emph{punctured} ternary
connected operator, not the product of the unpunctured V50 cumulant
with an independent binary estimate.

Throughout this chapter fix an omitted row \(o\) and write
\[
 T=\{a,b,c\}=\{1,2,3,4\}\setminus\{o\}.
\]
We assign a replica partition first to \([4]\) if it has a four-row
block, next to \([3,3]\) if it has two triple blocks with union all four
rows, and then to \([3,2]\).  With this priority, a \([3,2]\) term has
no nonsingleton block of size at least three containing \(o\).  Its
restriction to \(T\) is connected, contains at least one three-row
block, and \(o\) is attached to \(T\) by one or more pair blocks.
This convention is disjoint and is compatible with
\Cref{thm:V40-quartic-skeleton-list,
cor:V40-quartic-fibre-exhaustion}.

\begin{firewall}[What is and is not closed in V54]
\label{fire:V54-scope}
V54 derives the exact complete puncture formula for the \([3,2]\)
class and proves a coefficient-safe linearized ternary response.  It
closes three full-normalization branches of the isolated-attachment
sector: thermal attachment vertex, attachment-heavy input mass, and
common-junction-heavy input mass.  It does not close the remaining
sideways off-common branch, multiple attachment banks, or collision
coordinates on which the attaching pair leaves a nontrivial
complement moment.  Accordingly V54 does not claim the complete
\([3,2]\) topology or global quartic MTA.
\end{firewall}

\section{The exact punctured \texorpdfstring{\([3,2]\)}{[3,2]}
operator}
\label{sec:V54-puncture}

For a replica partition \(\sigma\), let \(A_o(\sigma)\) be the set of
coordinates at which the entry of row \(o\) belongs to a pair block.
For \(e\in A_o(\sigma)\), denote the other row in that block by
\(a_e\in T\).  At such a coordinate the remaining active rows in
\(T\setminus\{a_e\}\) form a partition independent of the selected
pair block.  Summing that complementary partition gives its complete
moment operator, denoted
\[
 \mathbf Q_{T\setminus\{a_e\},e}.
\]
Inactive rows are interpreted as identity factors.  Write
\(\mathbf d_{oa_e,e}\) for the local pair cumulant operator on the
selected block.

For a nonempty coordinate set \(A\) and a target map
\(\boldsymbol a:A\to T\), delete every incidence over \(A\) from the
three rows in \(T\).  Let
\begin{equation}
 \mathbf K_{T;A}^{[3]}
 :=
 \sum_{\substack{\omega\ {\rm connected\ on}\ T\\
                  \omega\ {\rm uses\ a\ local\ triple\ block}}}
 \bigotimes_{e\notin A}L_{e,\omega(e)}
 \label{eq:V54-punctured-ternary}
\end{equation}
be the resulting connected-coordinate operator.  The local factors
are those of \eqref{eq:V50-connected-operator-formula}; in particular
the five-term ternary cumulant at each retained common coordinate is
kept as the single operator \(t_e\).

\begin{theorem}[Exact \([3,2]\) puncture formula]
\label{thm:V54-exact-puncture}
The sum of all replica terms assigned to \([3,2]\) with omitted row
\(o\) is
\begin{equation}
 \boxed{
 \mathbf K_{T;o}^{[3,2]}
 =
 \sum_{\substack{\varnothing\ne A\subseteq X_o\\
                  \boldsymbol a\in\mathcal A_T(A)}}
 \left(\bigotimes_{e\in X_o\setminus A}q_{o,e}\right)
 \mathbf K_{T;A}^{[3]}
 \otimes
 \bigotimes_{e\in A}
 \left(
  \mathbf d_{oa_e,e}\otimes
  \mathbf Q_{T\setminus\{a_e\},e}
 \right).}
 \label{eq:V54-exact-puncture}
\end{equation}
Here \(\mathcal A_T(A)\) consists of the target maps allowed by the
active incidences.  Private coordinates of row \(o\) occur in the
first product.  Formula \eqref{eq:V54-exact-puncture} is an exact
operator identity before recoupling, trace norms, or the final
resolvent.
\end{theorem}

\begin{proof}
Take a term in the indicated topology class.  Since no block of size
at least three contains \(o\), the entry \((o,e)\) is either a
singleton or belongs to a unique pair block.  These alternatives
determine \(A=A_o(\sigma)\) and the target map
\(\boldsymbol a\) uniquely.

At \(e\in A\), the coordinate partition property makes the selected
pair block disjoint from all remaining blocks.  The remaining entries
are precisely those of \(T\setminus\{a_e\}\); summing their local
partitions is the moment--cumulant identity and gives
\(\mathbf Q_{T\setminus\{a_e\},e}\).  At \(e\notin A\), the row \(o\)
entry is a singleton and contributes \(q_{o,e}\).

After the coordinates in \(A\) are removed, the restriction to \(T\)
still contains the triple block from the minimal \([3,2]\) skeleton.
It is therefore connected and is summed exactly by
\(\mathbf K_{T;A}^{[3]}\).  Conversely, every term on the right has a
connected triple-containing restriction to \(T\) and at least one
pair block joining \(o\) to \(T\), so it is a \([3,2]\) term under the
stated priority.  The two constructions are inverse.  Coordinate
independence supplies the tensor product and proves the identity.
\end{proof}

\begin{corollary}[Pair-only subtraction is exact]
\label{cor:V54-pair-subtraction}
Let \(\mathbf K_{T;A}\) be the complete punctured ternary connected
operator and let \(\mathbf K_{T;A}^{(2)}\) be the part of
\eqref{eq:V50-connected-operator-formula} using pair bonds only.  Then
\begin{equation}
 \mathbf K_{T;A}^{[3]}
 =
 \mathbf K_{T;A}-\mathbf K_{T;A}^{(2)}.
 \label{eq:V54-pair-subtraction}
\end{equation}
Thus one may estimate the junction-bearing operator without expanding
any five-term local ternary cumulant.
\end{corollary}

\begin{proof}
Every connected assignment in
\eqref{eq:V50-connected-operator-formula} either contains a symbol
\(123\) or uses only the three pair symbols.  These alternatives are
disjoint and exhaustive.  Subtract the second sum from the first.
\end{proof}

\begin{corollary}[Exact coefficient normalization of every puncture
summand]
\label{cor:V54-puncture-normalization}
After all punctured-core, pair, complement, and spectator outputs have
been recoupled to a final irreducible \(\boldsymbol U\), the Fourier
coefficient of each summand in
\eqref{eq:V54-exact-puncture} is
\begin{equation}
 \frac{d_{\rm in}}{d_{\boldsymbol U}}
 \operatorname{Tr}_{M_{\boldsymbol U}}
 \left[
  X_{\boldsymbol U}
  (I_{H_{\boldsymbol U}}\otimes
   K_{\boldsymbol U}^{[3,2]})
 \right].
 \label{eq:V54-puncture-normalization}
\end{equation}
The BFA output factor cancels \(d_{\boldsymbol U}\), and every nested
raw resolution sums contractively.
\end{corollary}

\begin{proof}
Apply \eqref{eq:V46-exact-multiplier-normalization} only after the
complete tensor product in \eqref{eq:V54-exact-puncture} has been
formed.  Then use \eqref{eq:V46-global-raw-pinching}; at intermediate
pair and punctured-core decompositions use the same orthogonal
pinching before the final one.  Contractivity is stable under
spectator amplification.  No factor \(d_{\boldsymbol U}/d_{\rm in}\)
and no representation-formation probability is introduced.
\end{proof}

\section{A linearized ternary response with one final resolvent}
\label{sec:V54-linearized-ternary}

For a punctured three-row core let
\begin{align}
 \widehat\Xi_i
 &:=\Xi(\boldsymbol R_{i,\mathrm{core}}),
 \qquad i\in T,
 \label{eq:V54-punctured-mass}\\
 \widehat H_{ij}
 &:=\sum_{e\in X_i^{\rm core}\cap X_j^{\rm core}}
 A_{i,e}A_{j,e},
 \qquad
 \widehat M_i:=\widehat H_{ij}\widehat H_{ik},
 \quad\{i,j,k\}=T.
 \label{eq:V54-punctured-tree-numerator}
\end{align}
Zero denominators are omitted; then the corresponding numerator is
zero by exact connectivity.

Let \(V\) be a resolved auxiliary operator, including arbitrary
spectator amplification, with scalar multiplier \(r_V\) and output
mass \(\Xi_V\).  Put
\begin{equation}
 \mathfrak a_V:=|r_V|+
 s_\alpha\Xi_V|r_V|.
 \label{eq:V54-ternary-auxiliary-size}
\end{equation}
Denote by \(\mathcal L_{T}^{V}\) the BFA scalar response obtained after
the punctured ternary operator, all spectators, \(V\), and the
\emph{single} final resolvent have been assembled.

\begin{theorem}[Two-scale linearized ternary response]
\label{thm:V54-linearized-ternary}
For the complete punctured ternary connected operator, and separately
for its junction-bearing part
\(\mathbf K_{T;A}^{[3]}\), there are three routed nonnegative
majorants \(\mathcal L_{T,i}^{V}\), \(i\in T\), such that
\begin{equation}
 \mathcal L_T^V\le\sum_{i\in T}\mathcal L_{T,i}^V
 \label{eq:V54-center-routing}
\end{equation}
and
\begin{equation}
 \boxed{
 \mathcal L_{T,i}^V
 \le
 C\mathfrak a_V
 \widehat M_i
 \min\left\{
  s_\alpha,\frac1{\widehat\Xi_i}
 \right\}.}
 \label{eq:V54-linearized-ternary}
\end{equation}
One also has the universal auxiliary cap
\begin{equation}
 \mathcal L_T^V
 \le \frac C{s_\alpha}\mathfrak a_V.
 \label{eq:V54-linearized-universal}
\end{equation}
All constants are uniform in punctures, supports, representations,
outputs, LR multiplicities, and volume.
\end{theorem}

\begin{proof}
First take the complete connected operator.  Repeat the exact V47
decomposition on the punctured supports, and retain the center
allocation used in V48--V50.  A junction is allocated to one of its
three common trees by \Cref{lem:V50-junction-to-tree}; a two-bond term
is allocated to the center shared by its two marked bonds; and the
nested one-defect allocation is the two-term split in
\eqref{eq:V49-nested-marked-target}.  These choices give
\eqref{eq:V54-center-routing}.

We record two bounds for each routed majorant.  The marked bound
\begin{equation}
 \mathcal L_{T,i}^V
 \le
 C\frac{\widehat M_i}{\widehat\Xi_i}\mathfrak a_V
 \label{eq:V54-routed-marked-side}
\end{equation}
is the corresponding center term in
\Cref{thm:V48-high-defect-full-center,
thm:V49-nested-marked-target,thm:V50-K0-pointwise}.
Their proofs are linear in the retained scalar spectator.  On high
outputs split the final mass into punctured-core, ordinary-spectator,
and auxiliary output masses.  The first two pieces reproduce the cited
proofs with the extra factor \(|r_V|\); the last replaces that factor
by \(s_\alpha\Xi_V|r_V|\).

The raw routed bound is
\begin{equation}
 \mathcal L_{T,i}^V
 \le
 Cs_\alpha\widehat M_i\mathfrak a_V.
 \label{eq:V54-routed-raw-side}
\end{equation}
On low outputs it follows directly from the resolvent cap and the
V50 numerator.  On high outputs use the same exhaustive three-branch
split as in \Cref{thm:V50-K0-pointwise}.  In the globally thermal
branch, fusion bounds the core output by \(C/s_\alpha\), while every
spectator or auxiliary output is paid by its first moment.  In a
common-heavy branch allocated to center \(i\),
\(\widehat M_i\ge c\widehat\Xi_i^2\) and
\(s_\alpha\widehat\Xi_i\ge c\), so the universal
\(C\mathfrak a_V/s_\alpha\) payment is at most
\(Cs_\alpha\widehat M_i\mathfrak a_V\).  In an off-common hot branch,
the second spectator debit
\eqref{eq:V50-spectator-second-output} gives
\(C\widehat M_i/\Pi\), which is at most
\(Cs_\alpha\widehat M_i\) because \(s_\alpha\Pi>M\).
The V48 high-defect and V49 nested terms use the identical reduced
center, common-absorption, and off-common alternatives.  This proves
\eqref{eq:V54-routed-raw-side} on every final output.

Take the better of
\eqref{eq:V54-routed-marked-side} and
\eqref{eq:V54-routed-raw-side} for each already-routed center.  This
proves \eqref{eq:V54-linearized-ternary}.  The argument never applies
a second resolvent.

For \eqref{eq:V54-linearized-universal}, repeat
\Cref{lem:V50-universal-K0-response} for the zero-defect piece and the
universal parts of the V48--V49 proofs for the other exact V47
monomials.  Low outputs use the resolvent cap; high outputs use the
first moment of whichever output carries the mass.  In both cases the
only new factor is \(\mathfrak a_V\).

Finally set every local junction symbol \(t_e\) to zero in
\eqref{eq:V50-connected-operator-formula}.  The same proof, with the
nonnegative junction majorant deleted, proves both bounds for
\(\mathbf K_{T;A}^{(2)}\).  Subtract it from the complete operator by
\eqref{eq:V54-pair-subtraction} and use the triangle inequality.  No
five-term junction is split internally.
\end{proof}

\section{The attaching pair as an amplified auxiliary}
\label{sec:V54-attachment-auxiliary}

At an attachment coordinate \(e\), let the selected pair output be
\(S_{oa,e}\), let the complementary moment output be \(S_{\perp,e}\),
and write
\[
 \delta_{oa,e}
 :=
 q_{\alpha,S_{oa,e}}
 -q_{\alpha,R_{o,e}}q_{\alpha,R_{a,e}}.
\]
The resolved local attachment factor is
\[
 r_{a;o,e}
 :=
 \delta_{oa,e}\,q_{\alpha,S_{\perp,e}}.
\]

\begin{lemma}[Complement moments do not increase the pair tax]
\label{lem:V54-complement-stability}
With
\[
 h_{oa,e}:=A_{o,e}A_{a,e},
 \qquad
 \mathfrak a_{a;o,e}
 :=
 |r_{a;o,e}|+
 s_\alpha\bigl(\Xi(S_{oa,e})+\Xi(S_{\perp,e})\bigr)
 |r_{a;o,e}|,
\]
one has
\begin{equation}
 \mathfrak a_{a;o,e}
 \le
 C\left(
 |\delta_{oa,e}|+
 s_\alpha\Xi(S_{oa,e})|\delta_{oa,e}|
 \right)
 \le Cs_\alpha h_{oa,e}.
 \label{eq:V54-complement-stability}
\end{equation}
If the input-\(a\) mass \(x\) on an exact attachment bank satisfies
\(x\ge\varepsilon\Xi_a\), then after arbitrary raw coefficient
compression and spectator amplification,
\begin{equation}
 \boxed{
 \sum_{\rm raw\ outputs}
 \mathfrak a_{a;o}\,\omega_{\rm raw}
 \le
 C_\varepsilon\frac{h_{oa}}
 {s_\alpha\Xi_a^2}.}
 \label{eq:V54-amplified-attachment-tax}
\end{equation}
\end{lemma}

\begin{proof}
The complement moment is a heat multiplier in \([0,1]\), and
\eqref{eq:V46-product-first-moment} gives
\[
 s_\alpha\Xi(S_{\perp,e})
 q_{\alpha,S_{\perp,e}}\le C.
\]
The pair defect cap and first moment from V46 then prove
\eqref{eq:V54-complement-stability}.  Equation
\eqref{eq:V54-amplified-attachment-tax} is
\Cref{cor:V53-one-edge-tax}: the complement carrier and every other
punctured factor are spectators, and its proof is explicitly stable
under spectator amplification.  Nested raw pinching supplies
\(\sum\omega_{\rm raw}\le1\).
\end{proof}

\section{Three closed full-normalization branches}
\label{sec:V54-three-branches}

We now isolate a sector in which no collision ambiguity is present.
Assume
\begin{equation}
 X_o=P_o\mathbin{\dot\cup}L_{oa},
 \qquad
 X_o\cap X_b=X_o\cap X_c=\varnothing,
 \qquad
 L_{oa}\cap(X_b\cup X_c)=\varnothing.
 \label{eq:V54-isolated-attachment}
\end{equation}
Thus \(o\) has one exact pair bank, attached to \(a\), and otherwise
only private coordinates.  Put
\begin{align}
 x&:=\sum_{e\in L_{oa}}A_{a,e},
 &
 h&:=\sum_{e\in L_{oa}}A_{a,e}A_{o,e},
 \label{eq:V54-attachment-masses}\\
 Z_a&:=\sum_{e\in J_T}A_{a,e},
 &
 J_T&:=X_a^{\rm core}\cap X_b^{\rm core}
                   \cap X_c^{\rm core}.
 \label{eq:V54-common-core-mass}
\end{align}
Then \(\Xi_a=\widehat\Xi_a+x\), whereas
\(\Xi_b=\widehat\Xi_b\) and \(\Xi_c=\widehat\Xi_c\).
For the three trees obtained by adjoining \(ao\) to a ternary tree
centered at \(i\), define
\begin{equation}
 \mathsf U_{i\mid a,o}
 :=
 \frac{\widehat M_i h}{\Xi_i\Xi_a},
 \qquad i\in T,
 \label{eq:V54-promoted-tree-weight}
\end{equation}
where for \(i=a\) the denominator is \(\Xi_a^2\).

\begin{theorem}[Thermal and attachment-heavy promotion]
\label{thm:V54-thermal-attachment-heavy}
The isolated-attachment \([3,2]\) contribution satisfies
\begin{equation}
 \textup{coefficient contribution}
 \le C\sum_{i\in T}\mathsf U_{i\mid a,o}
 \label{eq:V54-closed-promoted-weight}
\end{equation}
on either of the following branches:
\begin{enumerate}[label=\textup{(\roman*)},leftmargin=3.7em]
\item \(s_\alpha\Xi_a\le M\);
\item \(s_\alpha\Xi_a>M\) and
      \(x\ge\varepsilon\Xi_a\).
\end{enumerate}
\end{theorem}

\begin{proof}
In branch \textup{(i)}, use
\(\mathfrak a_{a;o}\le Cs_\alpha h\).
For the \(i=a\) numerator take the raw side of
\eqref{eq:V54-linearized-ternary}:
\[
 Cs_\alpha\widehat M_a\,
 Cs_\alpha h
 \le
 C_M\frac{\widehat M_a h}{\Xi_a^2}.
\]
For \(i=b,c\), take the marked side:
\[
 C\frac{\widehat M_i}{\Xi_i}\,
 s_\alpha h
 \le
 C_M\frac{\widehat M_i h}{\Xi_i\Xi_a}.
\]
This proves \eqref{eq:V54-closed-promoted-weight}.

In branch \textup{(ii)}, apply
\eqref{eq:V54-amplified-attachment-tax}.  For \(i=a\), the raw side
gives
\[
 Cs_\alpha\widehat M_a
 \frac h{s_\alpha\Xi_a^2}
 =
 C\mathsf U_{a\mid a,o}.
\]
For \(i=b,c\), the marked side gives
\[
 C\frac{\widehat M_i}{\Xi_i}
 \frac h{s_\alpha\Xi_a^2}
 \le
 \frac C M
 \frac{\widehat M_i h}{\Xi_i\Xi_a}.
\]
All sums here are raw pinching sums, not sums of representation
formation probabilities.
\end{proof}

\begin{theorem}[Common-junction-heavy promotion]
\label{thm:V54-common-heavy}
Fix \(\varepsilon>0\).  In the isolated-attachment sector, if
\begin{equation}
 Z_a\ge\varepsilon\Xi_a,
 \label{eq:V54-common-heavy}
\end{equation}
then \eqref{eq:V54-closed-promoted-weight} holds without any
assumption on the attaching pair output.
\end{theorem}

\begin{proof}
Every coordinate in \(J_T\) is active in all three core rows, so
\[
 \widehat H_{ab}\ge Z_a,
 \qquad
 \widehat H_{ac}\ge Z_a.
\]
Consequently
\begin{equation}
 \mathsf U_{a\mid a,o}
 =
 \frac{\widehat H_{ab}\widehat H_{ac}h}{\Xi_a^2}
 \ge\varepsilon^2h.
 \label{eq:V54-common-absorbs-attachment}
\end{equation}
On the other hand,
\eqref{eq:V54-linearized-universal} and
\eqref{eq:V54-complement-stability} bound the full response by
\[
 \frac C{s_\alpha}(Cs_\alpha h)\le Ch.
\]
Apply \eqref{eq:V54-common-absorbs-attachment}.  The argument uses the
input common mass and never infers it from a possibly trivial common
output.
\end{proof}

\begin{corollary}[Exact residual in the isolated-attachment sector]
\label{cor:V54-sideways-residual}
After \Cref{thm:V54-thermal-attachment-heavy,
thm:V54-common-heavy}, an unclosed isolated-attachment term must obey,
for fixed small \(\varepsilon\),
\begin{equation}
 s_\alpha\Xi_a>M,\qquad
 x<\varepsilon\Xi_a,\qquad
 Z_a<\varepsilon\Xi_a,
 \qquad
 \widehat\Pi_a:=\widehat\Xi_a-Z_a
 >(1-2\varepsilon)\Xi_a.
 \label{eq:V54-sideways-regime}
\end{equation}
Thus a fixed fraction of the promoted vertex mass lies on
off-common coordinates of the punctured ternary operator.
\end{corollary}

\begin{proof}
Only the last inequality needs proof.  Since
\(\Xi_a=x+Z_a+\widehat\Pi_a\), the first two strict upper bounds give
\(\widehat\Pi_a>(1-2\varepsilon)\Xi_a\).
\end{proof}

\section{Why the remaining tax must stay joint}
\label{sec:V54-false-detachment}

\begin{proposition}[The detached attachment tax is false]
\label{prop:V54-detached-tax-false}
There is no support- and representation-uniform constant \(C\) such
that an arbitrary pair auxiliary satisfies
\begin{equation}
 \mathfrak a_{a;o}\le C\frac h{\Xi_a}
 \label{eq:V54-false-detached-tax}
\end{equation}
when the attachment bank is allowed to carry a vanishing fraction of
the full input mass.
\end{proposition}

\begin{proof}
On the attachment bank take the fundamental--dual pair
\((\mathbf3,\overline{\mathbf3})\) and resolve it to the singlet.
Its defect is
\[
 \delta_{\mathbf1}
 =
 1-
 q_{\alpha,\mathbf3}q_{\alpha,\overline{\mathbf3}}>0
\]
for fixed \(s_\alpha>0\).  Both \(h\) and
\(\mathfrak a_{a;o}\ge|\delta_{\mathbf1}|\) are then fixed positive
numbers.  Add to row \(a\), away from the attachment bank, an
irreducible \(R_N=(N,0)\).  The exact SU(3) Casimir formula gives
\(\Xi_a\to\infty\), so the right side of
\eqref{eq:V54-false-detached-tax} tends to zero while its left side
does not.

This does not refute the desired \([3,2]\) estimate: in the actual
response the new \(R_N\) also enters the punctured ternary operator or
its spectator multiplier.  It proves that this dependence must remain
assembled until the extra \(\Xi_a^{-1}\) has been earned.
\end{proof}

\begin{openproblem}[The exact sideways promotion gate]
\label{op:V54-sideways-gate}
In the regime \eqref{eq:V54-sideways-regime}, expand the punctured
ternary operator only into the exact V47 sectors.  On each sector
split \(\widehat\Pi_a\) into:
\begin{enumerate}[label=\textup{(\roman*)},leftmargin=3.7em]
\item banks retained as scalar singleton means in that monomial; and
\item exclusive-pair banks incident to \(a\) which are retained as
      genuine defects.
\end{enumerate}
The required estimate is
\begin{equation}
 \boxed{
 \mathcal L_{T}^{\,\mathbf D_{ao}}
 \le
 C\sum_{i\in T}
 \frac{\widehat M_i h}{\Xi_i\Xi_a}}
 \label{eq:V54-sideways-gate}
\end{equation}
after raw pair compression and with one final resolvent.  In branch
\textup{(i)} one must retain one unused product-multiplier debit after
the V50 ternary payment.  In branch \textup{(ii)} one must apply the
V53 one-edge tax to the actual internal defect and linearize the V49
nested cut before pinching.  Factoring out
\eqref{eq:V54-false-detached-tax} is forbidden.
\end{openproblem}

\section{Regression cards and V54 ledger}
\label{sec:V54-ledger}

\begin{assessment}[Long dual-singlet and attachment-position tests]
\label{ass:V54-regressions}
Put a long dual pair inside the three-row core and select a small or
trivial common output.  The common-heavy proof remains valid because
\eqref{eq:V54-common-absorbs-attachment} uses the two original input
overlaps, not the formed output Casimir.  Put a fundamental--dual
singlet on the attachment edge.  In the attachment-heavy branch its
coefficient mass is charged by the raw amplified tax
\eqref{eq:V54-amplified-attachment-tax}; no probability proportional
to \(d_S/(d_Rd_Q)\) is used.  Relabeling proves the same statements for
all four choices of omitted row and all three choices of attachment
vertex.  If the promoted vertex is instead long only on off-common
punctured-core coordinates, the example falls precisely into
\eqref{eq:V54-sideways-regime} and is not declared closed.
\end{assessment}

\begin{theorem}[Integrated V54 statement]
\label{thm:V54-integrated}
Preserving V1--V53, V54 proves:
\begin{enumerate}[label=\textup{(V54.\arabic*)},leftmargin=5.5em]
\item the exact complete \([3,2]\) puncture identity
      \eqref{eq:V54-exact-puncture};
\item exact separation of its junction-bearing and pair-only ternary
      cores without expanding a local five-term cumulant;
\item coefficient normalization by \(d_{\rm in}/d_{\boldsymbol U}\)
      and nested raw pinching;
\item the two-scale linearized ternary response
      \eqref{eq:V54-linearized-ternary} with one final resolvent;
\item stability of the V53 pair tax under a same-coordinate complement
      moment;
\item closure of the thermal, attachment-heavy, and
      common-junction-heavy isolated-attachment branches;
\item falsity of a detached \(h/\Xi_a\) attachment tax; and
\item reduction of the isolated residual to the joint sideways gate
      \eqref{eq:V54-sideways-gate}.
\end{enumerate}
\end{theorem}

\begin{proof}
Items \textup{(V54.1)--(V54.3)} are
\Cref{thm:V54-exact-puncture,cor:V54-pair-subtraction,
cor:V54-puncture-normalization}.  Item \textup{(V54.4)} is
\Cref{thm:V54-linearized-ternary}.  Item \textup{(V54.5)} is
\Cref{lem:V54-complement-stability}.  Item \textup{(V54.6)} is
\Cref{thm:V54-thermal-attachment-heavy,thm:V54-common-heavy}.
Items \textup{(V54.7)--(V54.8)} are
\Cref{prop:V54-detached-tax-false,
cor:V54-sideways-residual,op:V54-sideways-gate}.
\end{proof}

\begingroup
\small
\begin{longtable}{>{\raggedright\arraybackslash}p{0.29\textwidth}
                  >{\raggedright\arraybackslash}p{0.15\textwidth}
                  >{\raggedright\arraybackslash}p{0.48\textwidth}}
\toprule
\textbf{V54 row}&\textbf{Status}&\textbf{Advance or obstruction}\\
\midrule
\endfirsthead
\toprule
\textbf{V54 row}&\textbf{Status}&\textbf{Advance or obstruction}\\
\midrule
\endhead
Exact \([3,2]\) operator algebra
& \MGclosed
& Attachment coordinates puncture the intact ternary connected
  operator and carry exact pair times complement-moment factors.\\[0.35em]
Linearized ternary response
& \MGclosed
& Low and high outputs give the better of the raw \(s_\alpha M_i\)
  and marked \(M_i/\widehat\Xi_i\) payments with one resolvent.\\[0.35em]
Thermal and attaching-heavy isolated branches
& \MGclosed
& Raw defect size handles the thermal case; effective-capacity
  compression supplies the squared endpoint tax in the hot case.\\[0.35em]
Common-junction-heavy isolated branch
& \MGclosed
& Two original common overlaps absorb the promoted endpoint
  denominator even when the common output is a singlet.\\[0.35em]
Sideways off-common isolated branch
& \MGopen
& The extra endpoint denominator must be extracted jointly from scalar
  mean banks or genuine internal defects in the V47 polynomial.\\[0.35em]
Multiple and collision attachments
& \MGopen
& The exact puncture formula is known, but simultaneous attachment
  activities and nontrivial complement moments are not yet summed.\\[0.35em]
Complete \([3,2]\), \([3,3]\), and \([4]\)
& \MGopen
& No complete topology theorem is claimed until the sideways and
  multi-attachment gates close.\\[0.35em]
Full Yang--Mills mass gap
& \MGopen
& Global quartic/all-order MTA, WUFC/CPM, nonlinear construction,
  continuum OS existence/nontriviality, and the spectral passage
  remain downstream.\\
\bottomrule
\end{longtable}
\endgroup

\begin{openproblem}[V55 mandate: audit and close the sideways gate]
\label{op:V55}
\begin{enumerate}[label=\textup{(V55.\arabic*)},leftmargin=5.5em]
\item Perform the scheduled read-only audit of the newest SM and
      Navier--Stokes notes in \texttt{latex\_notes}.  Transfer only
      type-compatible assembly principles and record a strict
      firewall.  The following companion audit is V60.
\item In every V47 monomial contributing to
      \(\mathbf K_{T;A}^{[3]}\), split the sideways mass at the promoted
      vertex between retained singleton means and incident internal
      defects.  Do this before trace norm.
\item Prove the mean-heavy branch of
      \eqref{eq:V54-sideways-gate} by retaining one first or second
      product-multiplier debit after the ordinary V50 center payment.
\item In the defect-heavy branch, resolve the actual internal pair
      defect, apply \eqref{eq:V53-one-edge-tax}, and derive a
      linearized V49 nested-cut response with the attaching defect as
      auxiliary.  Charge the final resolvent once.
\item After the isolated sector closes, return to
      \eqref{eq:V54-exact-puncture}.  Sum multiple attachments by a
      deterministic selected edge and keep every other attachment and
      complement moment as an assembled first-moment spectator.
\item Test long dual-singlet cores, cold attaching singlets, all four
      omitted rows, all three attachment vertices, and collision
      coordinates with the complementary pair output also cold.
\item Do not declare complete \([3,2]\) until every nonempty \(A\) and
      every allowed target map in
      \eqref{eq:V54-exact-puncture} obeys its quartic tree weights.
      Keep \([3,3]\) and \([4]\) suspended.
\item Preserve only the newest YM file.
\end{enumerate}
\end{openproblem}

\begin{firewall}[End-of-V54 claim boundary]
\label{fire:V54-claim-boundary}
V54 proves the exact punctured \([3,2]\) operator identity, a
coefficient-safe linearized ternary response, and three
isolated-attachment promotion branches.  It does not prove the
sideways gate, the full multi-attachment/collision sum, complete
\([3,2]\), \([3,3]\), \([4]\), global \(m=4\), all-order MTA, WUFC,
CPM, the nonlinear Perron packet, finite-edge margins, constructive
capacity, protected-sector floors, local-algebra noncollapse,
reflection positivity, a nontrivial continuum OS state, or a positive
Yang--Mills spectral gap.  The full Yang--Mills mass gap is not proved.
\end{firewall}

\section*{Version V54 append-only record}
\addcontentsline{toc}{section}{Version V54 append-only record}
The complete V53 source was retained before this continuation.  Its
SHA--256 digest was
\begin{center}\scriptsize\ttfamily
6c6b0fa56b227a9d21ef12b1d43aecbcd33ad3b1da2d8f2f3850f6d4cb458332.
\end{center}
No mathematical content of V53 was altered.  On the next iteration,
retain this full file, append before its unique active document-closing
command, increment the filename to \texttt{\_V55.tex}, and remove only
the superseded V54 file from this Yang--Mills note series.

% =============================================================================
% END APPENDED VERSION V54 MATERIAL
% =============================================================================

% =============================================================================
% BEGIN APPENDED VERSION V55 MATERIAL
% =============================================================================

\chapter{V55: Companion audit and closure of the isolated
\texorpdfstring{\([3,2]\)}{[3,2]} sideways gate}
\label{ch:V55}

\section{Scheduled SM/NS companion audit}
\label{sec:V55-companion-audit}

The required read-only audit was performed before the V55 proof.  The
current companion snapshots in \texttt{latex\_notes} were
\begin{center}
\begin{tabular}{lll}
\toprule
programme&file&SHA--256\\
\midrule
SM--Einstein&
\texttt{Renewal\_SM\_Einstein\_Continuum\_Research\_Notes\_V70.tex}&
\texttt{3b05e612185a8f74230122baea9104aead7a589bf1b9527bf1b37b9a7cf1ea9f}\\
Navier--Stokes&
\texttt{Renewal\_NS\_Stability\_Research\_Notes\_V31.tex}&
\texttt{f1121b62238ad5491bb7cf03635ea9934942edf573ed8faaa9ee79e1d38a6696}\\
\bottomrule
\end{tabular}
\end{center}

SM V70 proves the exact dependency
\(F^H=F^h+\Sigma F^{(2)}\), removes the redundant row only after that
identity is assembled, and then identifies the complete adjoint domain
from one boundary Green form.  Its type-compatible lesson is that a
dependent defect must not be charged as an additional independent row.
For V55 this means that an internal ternary pair defect and the
attachment auxiliary must remain in one resolved response; one may not
pay both and then reinsert the same output mass through a detached
normalization.

NS V31 proves an exact formation first moment but also proves that
separating it from the active high-parent tail returns the
continuation-class quantity.  Its type-compatible lesson is to retain
the actual correlation between the mark and the carrier which is to
pay it.  For V55 this is precisely the distinction between the true
sideways response \eqref{eq:V54-sideways-gate} and the false detached
tax \eqref{eq:V54-false-detached-tax}.

\begin{theorem}[Type-correct V55 transfer decision]
\label{thm:V55-companion-transfer}
The companion audit changes the Yang--Mills proof only through the
following assembly rules:
\begin{enumerate}[label=\textup{(\roman*)},leftmargin=3.7em]
\item split the promoted input mass according to the factors actually
      present in a fixed V47 monomial;
\item retain a scalar multiplier debit until after the attaching
      defect output mass is paid;
\item if the mass lies on an internal defect, compress that actual raw
      defect before the nested common cut is summed; and
\item use the single final Yang--Mills resolvent only after these
      correlated factors have been assembled.
\end{enumerate}
\end{theorem}

\begin{proof}
These are order-of-assembly statements.  The SM exact dependency
forbids duplicate defect rows, while the NS sharpness cards forbid
replacing a correlated formation quantity by an independent bound.
The four rules are their direct typed consequences for
\eqref{eq:V54-sideways-gate}.
\end{proof}

\begin{firewall}[Strict companion type boundary]
\label{fire:V55-companion-boundary}
No SM BRST, KKT, boundary-domain, or BFV theorem implies a compact-group
Fourier estimate.  No NS heat-formation, parent-tail, or hinge theorem
implies a Clebsch--Gordan bound.  V55 imports no inequality from either
companion file.  Conversely, the result below supplies no theorem in
those programmes.  Their coupled-continuum and global-regularity goals
remain open.  The next scheduled companion audit is V60.
\end{firewall}

\section{The exact sideways mass split}
\label{sec:V55-mass-split}

Fix the isolated attachment geometry
\eqref{eq:V54-isolated-attachment}, and expand the punctured ternary
core by the exact V47 polynomial.  In a fixed monomial \(\mu\), every
off-common coordinate of input \(a\) is of exactly one of the following
two kinds:
\begin{enumerate}[label=\textup{(\roman*)},leftmargin=3.7em]
\item it lies in a private bank or an exclusive pair bank represented
      by its scalar separated mean in \(\mu\);
\item it lies in one of the at most two exclusive pair banks incident
      to \(a\) which is represented by a genuine defect in \(\mu\).
\end{enumerate}
Let \(p_a^\mu\) and \(d_a^\mu\) be the corresponding input-\(a\)
masses.

\begin{lemma}[Monomial-wise sideways dichotomy]
\label{lem:V55-sideways-dichotomy}
For every V47 monomial,
\begin{equation}
 \widehat\Pi_a=p_a^\mu+d_a^\mu.
 \label{eq:V55-exact-sideways-split}
\end{equation}
If \eqref{eq:V54-sideways-regime} holds, then either
\begin{equation}
 p_a^\mu\ge\frac{1-2\varepsilon}{2}\Xi_a,
 \label{eq:V55-mean-heavy}
\end{equation}
or one incident defect bank \(L_{ar}\), \(r\in T\setminus\{a\}\),
has input-\(a\) mass
\begin{equation}
 y_{ar}\ge\frac{1-2\varepsilon}{4}\Xi_a.
 \label{eq:V55-defect-heavy}
\end{equation}
\end{lemma}

\begin{proof}
The private and three pair-incidence banks in the exact ternary
partition are disjoint.  A fixed monomial selects either the separated
mean or the defect on each pair bank, proving
\eqref{eq:V55-exact-sideways-split}.  If
\eqref{eq:V55-mean-heavy} fails, then
\[
 d_a^\mu
 =\widehat\Pi_a-p_a^\mu
 >\frac{1-2\varepsilon}{2}\Xi_a.
\]
There are at most two defect banks incident to \(a\); one carries at
least half of this mass.
\end{proof}

\section{The mean-heavy branch}
\label{sec:V55-mean-heavy}

Let \(b_a^\mu\) be the product of all scalar singleton multipliers in
\(\mu\) which contain the mean-heavy input-\(a\) mass \(p=p_a^\mu\).
Let \(\Xi_\mu\) be the sum of their resolved output masses.

\begin{lemma}[One unused scalar debit survives attachment]
\label{lem:V55-unused-debit}
If \(s_\alpha p>M\), then
\begin{align}
 |b_a^\mu|
 &\le\frac C{(s_\alpha p)^2},
 \label{eq:V55-mean-second-debit}\\
 s_\alpha\Xi_\mu|b_a^\mu|
 &\le\frac C{s_\alpha p}.
 \label{eq:V55-mean-output-second-debit}
\end{align}
Let \(\delta_{ao}\) be the resolved isolated attachment defect of
weight \(h\), and include every remaining modulus-one,
first-moment spectator in \(r_{\rm rem}\).  Then the combined auxiliary
\[
 r_{\mu,ao}:=b_a^\mu\delta_{ao}r_{\rm rem}
\]
satisfies
\begin{equation}
 \boxed{
 |r_{\mu,ao}|+
 s_\alpha\Xi_{\mu,ao}|r_{\mu,ao}|
 \le C\frac h p.}
 \label{eq:V55-joint-mean-attachment}
\end{equation}
\end{lemma}

\begin{proof}
Equation \eqref{eq:V55-mean-second-debit} is the product second moment
\eqref{eq:V32-balanced-second-moment} on the exact union of the
retained mean banks.  To prove
\eqref{eq:V55-mean-output-second-debit}, split each fused mean output
mass into its input contributions and repeat
\eqref{eq:V50-spectator-second-output}; the input-\(a\) contribution
uses the second moment, and every other contribution uses its first
moment while retaining \eqref{eq:V55-mean-second-debit}.

Now split \(\Xi_{\mu,ao}\) among the mean, attachment, and remaining
spectator outputs.  The modulus term follows from
\(|\delta_{ao}|\le Cs_\alpha h\) and the first-moment consequence
\(|b_a^\mu|\le C/(s_\alpha p)\).  The attachment output uses
\[
 s_\alpha\Xi(S_{ao})|\delta_{ao}|
 \le Cs_\alpha h
\]
together with the same first debit.  The mean output uses
\eqref{eq:V55-mean-output-second-debit} and
\(|\delta_{ao}|\le Cs_\alpha h\).  A remaining output is paid by the
first moment of \(r_{\rm rem}\), leaving
\(|b_a^\mu\delta_{ao}|\le Ch/p\).  These four contributions prove
\eqref{eq:V55-joint-mean-attachment}.
\end{proof}

\begin{theorem}[Mean-heavy sideways promotion]
\label{thm:V55-mean-heavy-promotion}
In the isolated sideways regime, every monomial satisfying
\eqref{eq:V55-mean-heavy} obeys
\begin{equation}
 \textup{coefficient contribution}
 \le C\sum_{i\in T}\mathsf U_{i\mid a,o}.
 \label{eq:V55-mean-heavy-promotion}
\end{equation}
\end{theorem}

\begin{proof}
The sideways inequalities imply \(s_\alpha p>cM\), after increasing
the fixed thermal threshold if necessary, and
\[
 \widehat\Xi_a=\Xi_a-x>(1-\varepsilon)\Xi_a.
\]
Apply the marked side
\eqref{eq:V54-routed-marked-side} with the combined auxiliary
\eqref{eq:V55-joint-mean-attachment}.  For center \(a\),
\[
 \mathcal L_{T,a}^{\,\mu,ao}
 \le
 C\frac{\widehat M_a}{\widehat\Xi_a}\frac hp
 \le
 C_\varepsilon
 \frac{\widehat M_a h}{\Xi_a^2}.
\]
For \(i=b,c\), isolated attachment gives
\(\widehat\Xi_i=\Xi_i\), and
\[
 \mathcal L_{T,i}^{\,\mu,ao}
 \le
 C\frac{\widehat M_i}{\Xi_i}\frac hp
 \le
 C_\varepsilon
 \frac{\widehat M_i h}{\Xi_i\Xi_a}.
\]
Sum the three routed centers and use
\eqref{eq:V54-promoted-tree-weight}.  Raw pinching is already included
in the linearized response.
\end{proof}

\section{The defect-heavy one-cut branch}
\label{sec:V55-defect-heavy}

Suppose the sole exclusive defect in a V47 one-defect monomial is on
\(L_{ar}\), where \(r\in T\setminus\{a\}\), and let
\(t=T\setminus\{a,r\}\).  Write
\[
 h_{ar}:=\sum_{e\in L_{ar}}A_{a,e}A_{r,e},
 \qquad
 G_{ar;t}(\boldsymbol S)
 :=\sum_{e\in J_T}A_{S,e}A_{t,e}.
\]
Resolve the common cut in the \((ar)t\) bracketing as in V48.

\begin{theorem}[Linearized nested-cut promotion]
\label{thm:V55-nested-promotion}
Assume
\[
 s_\alpha\Xi_a>M,
 \qquad
 y_{ar}\ge\varepsilon\Xi_a.
\]
After raw compression of the \(ar\) defect and nested pinching of the
common cut, the one-defect term with the \(ao\) attachment satisfies
\begin{align}
 \textup{coefficient contribution}
 \le C h_{ar}h\left(
 \frac{H_{at}^{J_T}}{\Xi_a^2}
 +\frac{H_{rt}^{J_T}}{\Xi_a\Xi_r}
 \right).
 \label{eq:V55-nested-promotion}
\end{align}
The two terms are precisely the quartic tree weights for the trees
\(\{ar,at,ao\}\) and \(\{ar,rt,ao\}\).
\end{theorem}

\begin{proof}
The resolved common cut is the genuine binary defect
\(\Delta_{ar;t}\) of
\eqref{eq:V48-nested-defect}, and
\[
 G_{ar;t}(\boldsymbol S)
 \le C(H_{at}^{J_T}+H_{rt}^{J_T})
\]
by \eqref{eq:V48-nested-allocation}.  Treat the \(ar\) defect as the
auxiliary in \Cref{lem:V52-linearized-two-defect}, with the attaching
defect and the nested defect as the two selected banks.  Their weights
are \(h\) and \(G_{ar;t}(\boldsymbol S)\).  The amplified one-edge tax
gives
\[
 \mathfrak F_{ar}
 \le C_\varepsilon
 \frac{h_{ar}}{s_\alpha\Xi_a^2}.
\]
Consequently the raw one-resolvent route is
\begin{equation}
 C s_\alpha hG_{ar;t}(\boldsymbol S)\mathfrak F_{ar}
 \le
 C\frac{h_{ar}hG_{ar;t}(\boldsymbol S)}{\Xi_a^2}.
 \label{eq:V55-nested-square-route}
\end{equation}
This proves the \(H_{at}^{J_T}\) term.  It also proves the
\(H_{rt}^{J_T}\) term when \(\Xi_r\le\Xi_a\).

Assume \(\Xi_r>\Xi_a\).  Use instead the center-\(r\) marked route in
\Cref{thm:V49-nested-marked-target}, linearized by the attaching
auxiliary as in \Cref{thm:V54-linearized-ternary}.  The V49 proof uses
the \(ar\) defect only through its modulus and first moment, hence is
linear in its weighted defect functional.  Replacing the raw cap
\(Cs_\alpha h_{ar}\) by
\(\mathfrak F_{ar}\) multiplies that route by at most
\(C/(s_\alpha^2\Xi_a^2)\).  The attaching factor has weighted size at
most \(Cs_\alpha h\).  Therefore the center-\(r\) contribution is
\[
 C\frac{h_{ar}H_{rt}^{J_T}}{\Xi_r}
 \frac1{s_\alpha^2\Xi_a^2}
 (s_\alpha h)
 \le
 \frac C M
 \frac{h_{ar}H_{rt}^{J_T}h}{\Xi_r\Xi_a}.
\]
Nested raw pinching sums \(\boldsymbol S\) without a path count.  This
proves \eqref{eq:V55-nested-promotion}.
\end{proof}

\begin{corollary}[Defect-heavy one-defect closure]
\label{cor:V55-one-defect-closure}
Every defect-heavy V47 monomial with exactly one exclusive ternary
defect satisfies its isolated \([3,2]\) quartic marked-tree bound.
\end{corollary}

\begin{proof}
The unique defect bank incident to \(a\) is \(L_{ar}\), and the common
cut supplies the second ternary connection.  Apply
\Cref{thm:V55-nested-promotion}.  The two allocations in
\eqref{eq:V48-nested-allocation} give the two displayed spanning
trees.  Exact coefficient normalization and raw pinching are
\Cref{cor:V54-puncture-normalization}.
\end{proof}

\section{High-defect terms and complete isolated closure}
\label{sec:V55-isolated-closure}

\begin{lemma}[High-defect ternary terms remain V53 trees]
\label{lem:V55-high-defect-tree}
Take a V47 monomial with at least two exclusive ternary defects and
adjoin the isolated \(ao\) defect.  Its coefficient contribution
satisfies the quartic marked-tree bound when the complete common
moment is retained.
\end{lemma}

\begin{proof}
Any two distinct pair edges on the three vertices in \(T\) form a
ternary spanning tree.  Adjoining \(ao\) gives a four-vertex spanning
tree; a third exclusive ternary defect is surplus.  Resolve the common
moment only after these three selected defects have been fixed.  The
complete common factor has a uniform modulus and first moment by
\Cref{lem:V47-common-first-moments}; every unused mean or defect is an
admissible spectator.  The coefficient-sensitive star and path gates
\Cref{thm:V53-full-star-gate,thm:V53-full-path-gate} therefore apply
with this assembled common auxiliary.  Their proofs require only its
modulus and first moment and are stable under raw spectator
amplification.  V46 pinching removes every intermediate multiplicity.
\end{proof}

\begin{theorem}[Complete isolated-attachment assembled quartic sector]
\label{thm:V55-complete-isolated-32}
Under \eqref{eq:V54-isolated-attachment}, the complete
quartic connected multiplier is
\begin{equation}
 \boxed{
 \mathbf K_{4}^{\rm iso}
 =
 q_{P_o}\,\mathbf K_T\,\mathbf D_{ao},}
 \label{eq:V55-isolated-exact-product}
\end{equation}
where \(\mathbf K_T\) is the complete ternary connected operator on
the supports punctured by \(L_{ao}\), and \(\mathbf D_{ao}\) is the
complete attachment-bank defect.  This assembled sector satisfies the
\(m=4\) marked-tree atlas bound for arbitrary trace-class Fourier
coefficients, uniformly in volume, supports, representation heights,
final outputs, and LR multiplicities.  It contains, without separating
by trace norm, both the isolated \([3,2]\) terms and the pair-only
terms on the same support geometry.  The theorem holds for all four
choices of omitted row and all three attachment vertices.
\end{theorem}

\begin{proof}
In the isolated geometry, every connected replica partition must use
a pair block on \(L_{ao}\) to attach \(o\).  After that bank is
deleted, its restriction to \(T\) must be connected, since \(o\) has
no route to either of the other two rows.  Summing the nonempty pair
blocks on \(L_{ao}\) gives \(\mathbf D_{ao}\), while summing the
connected restriction to \(T\) gives the intact operator
\(\mathbf K_T\).  Coordinate independence and the private row-\(o\)
means give \eqref{eq:V55-isolated-exact-product}.  This is also the
sum of the V54 \([3,2]\) puncture operator and its pair-only
counterpart, formed before any norm.

The thermal, attachment-heavy, and common-heavy branches are
\Cref{thm:V54-thermal-attachment-heavy,thm:V54-common-heavy}.
Although those theorems were stated for the \([3,2]\) contribution,
their proofs invoke the complete-core form of
\Cref{thm:V54-linearized-ternary} before the optional pair-only
subtraction.  The same displayed estimates therefore apply directly
to \(\mathbf K_T\) in \eqref{eq:V55-isolated-exact-product}.

In the remaining sideways regime, apply
\Cref{lem:V55-sideways-dichotomy} to each exact V47 monomial.  A
mean-heavy term is bounded by
\Cref{thm:V55-mean-heavy-promotion}.  A zero-exclusive-defect term is
necessarily of this kind.  A defect-heavy one-exclusive-defect term is
\Cref{cor:V55-one-defect-closure}; a term with at least two defects is
\Cref{lem:V55-high-defect-tree}.  These alternatives exhaust the V47
polynomial.

Finally use \eqref{eq:V54-puncture-normalization} and nested raw
pinching.  The three promoted weights
\eqref{eq:V54-promoted-tree-weight} are three terms already present
in the quartic MTA sum.  Relabeling \(T,o,a\) proves the last sentence.
\end{proof}

\begin{assessment}[What remains in the exact puncture formula]
\label{ass:V55-general-puncture-residual}
\Cref{thm:V55-complete-isolated-32} closes the complete quartic
cumulant whenever the omitted row has one exact pair attachment bank
and no other shared coordinate.  In the general puncture formula, the
remaining \([3,2]\) terms have at least one of the following exact
features:
\begin{enumerate}[label=\textup{(\roman*)},leftmargin=3.7em]
\item \(A\) contains pair attachments to two or three different
      vertices of \(T\);
\item an attachment coordinate also carries an active complementary
      row, hence the factor
      \(\mathbf Q_{T\setminus\{a_e\},e}\) is nontrivial; or
\item puncturing several attachment coordinates removes a material
      fraction of the mass of more than one ternary tree center.
\end{enumerate}
The local complement factor itself is harmless by
\Cref{lem:V54-complement-stability}.  The unresolved issue is the
simultaneous allocation of the full original-input denominators after
all punctures, not a new LR multiplicity or one-link cumulant estimate.
\end{assessment}

\section{V55 ledger and next acquisition}
\label{sec:V55-ledger}

\begin{theorem}[Integrated V55 statement]
\label{thm:V55-integrated}
Preserving V1--V54, V55 proves:
\begin{enumerate}[label=\textup{(V55.\arabic*)},leftmargin=5.5em]
\item the scheduled, digest-recorded SM V70 and NS V31 audit with a
      strict type firewall;
\item the exact monomial-wise split
      \(\widehat\Pi_a=p_a^\mu+d_a^\mu\);
\item the surviving joint scalar debit
      \eqref{eq:V55-joint-mean-attachment};
\item closure of the mean-heavy sideways branch;
\item the linearized nested-cut promotion
      \eqref{eq:V55-nested-promotion};
\item closure of the defect-heavy one- and high-defect branches; and
\item the complete arbitrary-coefficient isolated-attachment assembled
      quartic marked-tree theorem.
\end{enumerate}
\end{theorem}

\begin{proof}
Item \textup{(V55.1)} is
\Cref{thm:V55-companion-transfer,fire:V55-companion-boundary}.
Item \textup{(V55.2)} is
\Cref{lem:V55-sideways-dichotomy}.  Items
\textup{(V55.3)--(V55.4)} are
\Cref{lem:V55-unused-debit,thm:V55-mean-heavy-promotion}.
Items \textup{(V55.5)--(V55.6)} are
\Cref{thm:V55-nested-promotion,cor:V55-one-defect-closure,
lem:V55-high-defect-tree}.  Item \textup{(V55.7)} is
\Cref{thm:V55-complete-isolated-32}.
\end{proof}

\begingroup
\small
\begin{longtable}{>{\raggedright\arraybackslash}p{0.29\textwidth}
                  >{\raggedright\arraybackslash}p{0.15\textwidth}
                  >{\raggedright\arraybackslash}p{0.48\textwidth}}
\toprule
\textbf{V55 row}&\textbf{Status}&\textbf{Advance or obstruction}\\
\midrule
\endfirsthead
\toprule
\textbf{V55 row}&\textbf{Status}&\textbf{Advance or obstruction}\\
\midrule
\endhead
SM/NS companion audit
& \MGclosed
& Only exact assembly discipline transfers; neither companion
  inequality is used in the YM proof.\\[0.35em]
Mean-heavy sideways promotion
& \MGclosed
& A second scalar product debit survives the attachment output and
  supplies the promoted endpoint denominator.\\[0.35em]
Defect-heavy nested promotion
& \MGclosed
& The actual internal defect is compressed before the common cut;
  the two original-input allocations give a star and a path.\\[0.35em]
Complete isolated-attachment quartic sector
& \MGclosed
& The exact product \(\mathbf K_T\mathbf D_{ao}\), including its
  pair-only and \([3,2]\) parts before norm, has the quartic tree
  weights for every omitted/attachment labeling.\\[0.35em]
Multiple/collision \([3,2]\) sector
& \MGopen
& Several punctures alter more than one original center mass and must
  be allocated simultaneously.\\[0.35em]
Complete \([3,2]\), \([3,3]\), and \([4]\)
& \MGopen
& The first still lacks its multiple/collision sector; the latter two
  have not been re-entered with the corrected coefficient law.\\[0.35em]
Full Yang--Mills mass gap
& \MGopen
& Global quartic/all-order MTA, WUFC/CPM, nonlinear construction,
  continuum OS existence/nontriviality, and the spectral passage
  remain downstream.\\
\bottomrule
\end{longtable}
\endgroup

\begin{openproblem}[V56 mandate: simultaneous puncture allocation]
\label{op:V56}
\begin{enumerate}[label=\textup{(V56.\arabic*)},leftmargin=5.5em]
\item Start from the exact formula
      \eqref{eq:V54-exact-puncture}; do not replace its nonempty
      attachment set \(A\) by independent binary probabilities.
\item Choose a deterministic selected attachment
      \((e,a_e)\) only after recording which original input masses are
      removed by every puncture.  Route the two ternary marks and the
      selected attachment to an actual four-vertex tree.
\item Keep every other attachment and every complement moment in one
      assembled auxiliary.  Prove its modulus/first-moment bound before
      using \Cref{thm:V54-linearized-ternary}.
\item Derive a simultaneous mass dichotomy: selected-attachment heavy,
      common-heavy, retained-mean heavy, internal-defect heavy, or
      complement-output heavy at each promoted internal vertex.
\item Form one aggregate raw projector before summing collision
      outputs.  A sum of one-copy Schur taxes is forbidden.
\item Test two cold attachments sharing the omitted row, a four-active
      coordinate partitioned as \(\{o,a\}\mid\{b,c\}\), and a long
      dual-singlet core at each possible internal vertex.
\item Declare complete \([3,2]\) only after every nonempty \(A\) and
      target map in \eqref{eq:V54-exact-puncture} has been bounded.
      Keep \([3,3]\) and \([4]\) suspended.
\item Preserve only the newest YM file.  No companion audit is due;
      the next scheduled audit is V60.
\end{enumerate}
\end{openproblem}

\begin{firewall}[End-of-V55 claim boundary]
\label{fire:V55-claim-boundary}
V55 closes the entire isolated-attachment assembled quartic sector,
including the sideways mean- and defect-heavy branches.  It does not
separate the pair-only and \([3,2]\) suboperators after trace norm, and
it does not prove the multiple/collision part of
\eqref{eq:V54-exact-puncture}, complete
\([3,2]\), \([3,3]\), \([4]\), global \(m=4\), all-order MTA, WUFC,
CPM, the nonlinear Perron packet, finite-edge margins, constructive
capacity, protected-sector floors, local-algebra noncollapse,
reflection positivity, a nontrivial continuum OS state, or a positive
Yang--Mills spectral gap.  The full Yang--Mills mass gap is not proved.
\end{firewall}

\section*{Version V55 append-only record}
\addcontentsline{toc}{section}{Version V55 append-only record}
The complete V54 source was retained before this continuation.  Its
SHA--256 digest was
\begin{center}\scriptsize\ttfamily
3bc97f4ef43bd6eb4289c8352d513d94e361806b43cba57feb3747de9c704d3f.
\end{center}
No mathematical content of V54 was altered.  On the next iteration,
retain this full file, append before its unique active document-closing
command, increment the filename to \texttt{\_V56.tex}, and remove only
the superseded V55 file from this Yang--Mills note series.

% =============================================================================
% END APPENDED VERSION V55 MATERIAL
% =============================================================================

% =============================================================================
% BEGIN APPENDED VERSION V56 MATERIAL
% =============================================================================

\chapter{V56: Simultaneous exact-pair attachments and the
noncollision \texorpdfstring{\([3,2]\)}{[3,2]} sector}
\label{ch:V56}

\section{Purpose and analytic regrouping}
\label{sec:V56-purpose}

V55 closed the support geometry in which the omitted row has one exact
pair attachment.  The first simultaneous-puncture problem is the case
of two or three attachment banks which are still exact pair banks.
There is then no local collision: each attachment coordinate contains
only the omitted row and its target.  The entire attachment sum is a
seven-term polynomial, not a sum over coordinate subsets.  This finite
regrouping keeps the complete punctured ternary operator intact and
allows the attachment edge to be chosen separately for each routed
ternary center.

The result below closes the complete \(T\)-connected analytic sector
under the noncollision incidence hypothesis.  This sector contains all
\([3,2]\) terms in that geometry together with the pair-only terms
whose restriction to \(T\) is already connected.  No trace norm is
used to separate those two suboperators.

\section{The exact multi-attachment polynomial}
\label{sec:V56-polynomial}

Fix \(o\) and \(T=\{a,b,c\}\).  Assume
\begin{equation}
 X_o=P_o\mathbin{\dot\cup}
 \mathop{\dot\bigcup}_{i\in T}L_{oi},
 \qquad
 L_{oi}\cap X_j=\varnothing
 \quad(i\ne j,\ i,j\in T).
 \label{eq:V56-noncollision-incidence}
\end{equation}
Thus every coordinate shared by \(o\) and \(T\) is active in exactly
the two rows \(o,i\).  Remove the three \(L_{oi}\) from the supports of
the rows in \(T\), and let \(\mathbf K_T^\circ\) be the complete
ternary connected operator on the remaining core.

On \(L_{oi}\), put
\begin{align}
 a_i&:=
 q_{\alpha,\boldsymbol R_{o,L_{oi}}}
 q_{\alpha,\boldsymbol R_{i,L_{oi}}},
 \label{eq:V56-attachment-mean}\\
 \mathbf D_i&:=
 \mathbf Q_{oi,L_{oi}}-a_iI,
 \label{eq:V56-attachment-defect}\\
 x_i&:=\sum_{e\in L_{oi}}A_{i,e},
 \qquad
 h_i:=\sum_{e\in L_{oi}}A_{o,e}A_{i,e}.
 \label{eq:V56-attachment-weights}
\end{align}
Empty banks have \(a_i=1\), \(\mathbf D_i=0\), and
\(x_i=h_i=0\).  Let
\(\widehat\Xi_i\) and \(\widehat M_i\) be the punctured core quantities
of \eqref{eq:V54-punctured-mass}--%
\eqref{eq:V54-punctured-tree-numerator}.  Then
\begin{equation}
 \Xi_i=\widehat\Xi_i+x_i,
 \qquad i\in T.
 \label{eq:V56-full-versus-core-mass}
\end{equation}

\begin{theorem}[Exact \(T\)-connected attachment polynomial]
\label{thm:V56-exact-polynomial}
The sum of all connected replica terms whose restriction to the rows
in \(T\), after the exact attachment banks have been deleted, is
connected equals
\begin{align}
 \boxed{
 \mathbf K_{T;o}^{\rm conn}
 =
 q_{P_o}\mathbf K_T^\circ
 \left[
 \prod_{i\in T}(a_iI+\mathbf D_i)
 -\prod_{i\in T}a_iI
 \right].}
 \label{eq:V56-exact-polynomial}
\end{align}
Equivalently,
\begin{equation}
 \mathbf K_{T;o}^{\rm conn}
 =
 q_{P_o}\mathbf K_T^\circ
 \sum_{\varnothing\ne S\subseteq T}
 \left(\prod_{i\in S}\mathbf D_i\right)
 \left(\prod_{j\in T\setminus S}a_j\right).
 \label{eq:V56-seven-monomials}
\end{equation}
This is an exact seven-monomial operator identity on disjoint banks.
\end{theorem}

\begin{proof}
At an exact pair bank \(L_{oi}\), the complete joint moment is
\(a_iI+\mathbf D_i\); choosing \(a_iI\) leaves \(o\) separated from
\(i\), while choosing \(\mathbf D_i\) inserts at least one pair block
joining them.  Since the restriction to \(T\) is already connected,
the full four-row hypergraph is connected exactly when at least one of
the three attachment defects is selected.  Summing the connected
restriction to \(T\) gives \(\mathbf K_T^\circ\), private row-\(o\)
coordinates give \(q_{P_o}\), and coordinate independence gives the
product.  Subtracting the all-mean monomial proves
\eqref{eq:V56-exact-polynomial}; expanding it gives
\eqref{eq:V56-seven-monomials}.
\end{proof}

\begin{lemma}[Selected attachment and surplus admissibility]
\label{lem:V56-surplus-admissibility}
Fix \(\varnothing\ne S\subseteq T\) in
\eqref{eq:V56-seven-monomials} and select \(r\in S\).  After all
unselected attachment outputs and private outputs are resolved, their
factor
\begin{equation}
 r_{S,r}:=
 q_{P_o}
 \left(\prod_{i\in S\setminus\{r\}}\delta_i\right)
 \left(\prod_{j\in T\setminus S}a_j\right)
 \label{eq:V56-surplus-factor}
\end{equation}
satisfies
\begin{equation}
 |r_{S,r}|\le1,
 \qquad
 s_\alpha\Xi_{S,r}|r_{S,r}|\le C.
 \label{eq:V56-surplus-admissibility}
\end{equation}
The same bounds hold after arbitrary coefficient amplification.
\end{lemma}

\begin{proof}
Every resolved mean lies in \([0,1]\) and every resolved defect is a
difference of two such multipliers, so the modulus is at most one.
Split \(\Xi_{S,r}\) among the at most six unselected attachment and
private bank outputs.  Use
\eqref{eq:V46-product-first-moment} on a mean or private bank and
\eqref{eq:V46-defect-first-moment} on a defect bank, bounding all
other factors by one.  The number of bank types is fixed.  Tensoring
with coefficient carriers does not alter these scalar estimates, and
raw pinching is contractive.
\end{proof}

\section{Routing an attachment polynomial by ternary center}
\label{sec:V56-center-routing}

Fix a routed ternary center \(i\in T\) from
\eqref{eq:V54-center-routing}.  For a nonempty attachment set \(S\),
use the following deterministic rule:
\begin{equation}
 r(i,S):=
 \begin{cases}
 i,&i\in S,\\
 \min S,&i\notin S,
 \end{cases}
 \label{eq:V56-selected-attachment}
\end{equation}
where the minimum is taken in the fixed row order.  If \(i\in S\),
the selected attachment \(oi\) turns the ternary tree centered at
\(i\) into a four-row star.  If \(i\notin S\), the selected attachment
\(or(i,S)\) produces a path whose internal vertices are \(i\) and
\(r(i,S)\); importantly, the absent edge \(oi\) is represented by the
mean \(a_i\).

\begin{lemma}[Own-edge star promotion is surplus-stable]
\label{lem:V56-own-edge-star}
If \(i\in S\), the center-\(i\) contribution of the corresponding
monomial satisfies
\begin{equation}
 \textup{coefficient contribution}_{i,S}
 \le
 C\frac{\widehat M_i h_i}{\Xi_i^2}.
 \label{eq:V56-own-edge-star}
\end{equation}
\end{lemma}

\begin{proof}
Select \(\mathbf D_i\) and put every other attachment factor into
\eqref{eq:V56-surplus-factor}.  No unselected exact attachment bank
contains row \(i\).  Hence its original mass is exactly
\(\widehat\Xi_i+x_i\), as in the isolated geometry.  Repeat the four
branches used in
\Cref{thm:V54-thermal-attachment-heavy,
thm:V54-common-heavy,thm:V55-mean-heavy-promotion,
cor:V55-one-defect-closure}, with
\eqref{eq:V56-surplus-admissibility} included in the auxiliary.
High-defect ternary monomials are
\Cref{lem:V55-high-defect-tree}.  Each cited proof uses an unselected
factor only through its modulus and first moment, so the same
calculation gives
\eqref{eq:V56-own-edge-star}.  The selected marks are the two
center-\(i\) ternary marks and \(oi\).
\end{proof}

\begin{lemma}[Absent-edge mean gives a path or selected-center star]
\label{lem:V56-absent-edge-mean}
Let \(i\notin S\), put \(r=r(i,S)\), and retain the mean \(a_i\) on
the absent bank \(L_{oi}\).  Then
\begin{equation}
 \textup{coefficient contribution}_{i,S}
 \le
 C\left(
 \frac{\widehat M_i h_r}{\Xi_i\Xi_r}
 +\frac{\widehat M_r h_r}{\Xi_r^2}
 \right).
 \label{eq:V56-absent-edge-mean}
\end{equation}
\end{lemma}

\begin{proof}
First perform the complete selected-\(or\) promotion, keeping \(a_i\)
and all other attachment factors assembled.  If row \(r\) is in the
common-junction-heavy branch, \Cref{thm:V54-common-heavy} reroutes the
response to the star centered at \(r\), giving the second term in
\eqref{eq:V56-absent-edge-mean}.  Assume henceforth that this branch
does not occur.  The remaining V54--V55 alternatives preserve the
routed center-\(i\) contribution and give the denominator
\(\Xi_r^{-1}\); their center-\(i\) side has either the raw factor
\(s_\alpha\widehat M_i\) or the marked factor
\(\widehat M_i/\widehat\Xi_i\).

It remains only to replace the punctured center mass by
\(\Xi_i=\widehat\Xi_i+x_i\).  If
\(x_i\le\widehat\Xi_i\), then
\(\Xi_i\le2\widehat\Xi_i\), and the marked side gives the result.
Suppose \(x_i>\widehat\Xi_i\).  If
\(s_\alpha\Xi_i\le M\), use the raw side:
\[
 s_\alpha\widehat M_i
 \le
 M\frac{\widehat M_i}{\Xi_i}.
\]
If \(s_\alpha\Xi_i>M\), the row-\(i\) factor inside \(a_i\) carries
more than half of \(\Xi_i\).  Retain its first debit on core,
selected-attachment, and surplus outputs, and its second debit on its
own fused output, exactly as in
\Cref{lem:V55-unused-debit}.  This supplies
\[
 |a_i|+s_\alpha\Xi_{\rm out}|a_i|
 \le\frac C{s_\alpha x_i}
 \le\frac C{s_\alpha\Xi_i}
\]
relative to the raw routed bound.  Thus that bound becomes
\(C\widehat M_i/\Xi_i\).

The \(oi\) mean bank and the selected \(or\) defect bank are disjoint.
Their output splits therefore use different multiplier factors, and
the preceding replacement commutes with the selected-\(r\) promotion.
Multiplying by its already earned factor \(h_r/\Xi_r\) proves
 the first term of \eqref{eq:V56-absent-edge-mean}.  The final
 resolvent was used only
inside the common linearized response.
\end{proof}

\begin{theorem}[Complete noncollision \(T\)-connected quartic sector]
\label{thm:V56-noncollision-connected}
Under \eqref{eq:V56-noncollision-incidence}, the exact operator
\(\mathbf K_{T;o}^{\rm conn}\) of
\eqref{eq:V56-exact-polynomial} satisfies its \(m=4\) marked-tree
atlas bound for arbitrary trace-class Fourier coefficients, uniformly
in volume, supports, representations, final channels, and
multiplicities.
\end{theorem}

\begin{proof}
Expand the seven exact monomials
\eqref{eq:V56-seven-monomials}, but do not expand the punctured
ternary operator.  Route that operator into its three center
majorants by \eqref{eq:V54-center-routing}.  For a fixed \(S,i\), use
\eqref{eq:V56-selected-attachment}.  If \(i\in S\), apply
\Cref{lem:V56-own-edge-star}; its right side is the marked weight of
the star formed from the center-\(i\) ternary tree and \(oi\).
If \(i\notin S\), apply
\Cref{lem:V56-absent-edge-mean}; its right side is the marked weight of
the path obtained by adjoining \(or(i,S)\) to the center-\(i\)
ternary tree, plus the selected-center star used only in the
common-heavy rerouting.

There are seven attachment monomials and three routed centers, a fixed
number independent of support and volume.  For every final channel
use the exact \(d_{\rm in}/d_{\boldsymbol U}\) normalization, cancel
the BFA output dimension, and sum all nested raw blocks by
\eqref{eq:V46-global-raw-pinching}.  This proves the theorem.
\end{proof}

\begin{corollary}[All noncollision \([3,2]\) terms are analytically
assigned]
\label{cor:V56-noncollision-32}
Every \([3,2]\) replica term for which all blocks joining the omitted
row to \(T\) lie on exact pair banks is contained in the closed
operator \(\mathbf K_{T;o}^{\rm conn}\).  Hence no noncollision
\([3,2]\) term remains outside a proved quartic MTA sector.
\end{corollary}

\begin{proof}
A \([3,2]\) term contains a three-row block on \(T\), so its
restriction to \(T\) is connected after the exact attachment banks are
deleted.  It has at least one pair attachment to \(o\).  These are
exactly the two conditions used in
\Cref{thm:V56-exact-polynomial}.  Apply
\Cref{thm:V56-noncollision-connected}.  The corollary assigns the
term as part of an exact assembled sector and does not take a norm of
the topology projection alone.
\end{proof}

\section{Exact form of a collision coordinate}
\label{sec:V56-collision}

Let \(e\) be an attachment coordinate at which the selected pair block
is \(\{o,a\}\).  If both remaining rows \(b,c\) are active, the
complement moment in \eqref{eq:V54-exact-puncture} has the exact split
\begin{equation}
 \mathbf Q_{\{b,c\},e}
 =
 q_{b,e}q_{c,e}I+\mathbf d_{bc,e}.
 \label{eq:V56-collision-split}
\end{equation}
Thus the local collision factor is
\begin{equation}
 \boxed{
 \mathbf d_{oa,e}\otimes\mathbf Q_{\{b,c\},e}
 =
 q_{b,e}q_{c,e}\mathbf d_{oa,e}
 +
 \mathbf d_{oa,e}\otimes\mathbf d_{bc,e}.}
 \label{eq:V56-collision-factor}
\end{equation}
If at most one complementary row is active, only the first,
mean-complement branch occurs.

\begin{proposition}[Collision reduction]
\label{prop:V56-collision-reduction}
The remaining \([3,2]\) problem is reduced to two exact local
branches:
\begin{enumerate}[label=\textup{(\roman*)},leftmargin=3.7em]
\item a mean-complement attachment, in which the complementary input
      masses must be paid by retained first and second product debits;
\item a double-pair attachment
      \(\mathbf d_{oa,e}\otimes\mathbf d_{bc,e}\), in which the two
      disjoint defects must be routed together before raw pinching.
\end{enumerate}
No third collision type exists.
\end{proposition}

\begin{proof}
At one coordinate a partition block containing \(o\) and \(a\) leaves
only the entries of \(b,c\).  Their partition is either two singletons
or their one pair block.  Summing those two choices gives
\eqref{eq:V56-collision-split}, and distributing
\(\mathbf d_{oa,e}\) gives
\eqref{eq:V56-collision-factor}.  If one entry is absent, its only
partition is the singleton.  These cases are exhaustive.
\end{proof}

\begin{assessment}[Cold-collision regression]
\label{ass:V56-collision-regression}
The second term in \eqref{eq:V56-collision-factor} permits both pairs
to resolve to singlets.  Neither output Casimir then sees the large
input mass.  A valid proof must form one aggregate raw projector on
the tensor product of the two pair carriers and use the two actual
defect taxes or a joint tree response.  Summing two one-copy
dimension ratios would repeat the V45 failure.  The first term is
different: its complement is a genuine scalar mean and must retain
its multiplier debit until the selected attachment output has been
paid.
\end{assessment}

\section{V56 ledger and next acquisition}
\label{sec:V56-ledger}

\begin{theorem}[Integrated V56 statement]
\label{thm:V56-integrated}
Preserving V1--V55, V56 proves:
\begin{enumerate}[label=\textup{(V56.\arabic*)},leftmargin=5.5em]
\item the exact seven-monomial multi-attachment polynomial
      \eqref{eq:V56-seven-monomials};
\item uniform admissibility of every unselected attachment factor;
\item center-adapted attachment routing
      \eqref{eq:V56-selected-attachment};
\item surplus-stable own-edge star promotion;
\item the absent-edge mean replacement
      \eqref{eq:V56-absent-edge-mean};
\item the complete arbitrary-coefficient noncollision
      \(T\)-connected quartic MTA sector; and
\item the exact two-branch collision reduction
      \eqref{eq:V56-collision-factor}.
\end{enumerate}
\end{theorem}

\begin{proof}
Item \textup{(V56.1)} is
\Cref{thm:V56-exact-polynomial}; item \textup{(V56.2)} is
\Cref{lem:V56-surplus-admissibility}.  Items
\textup{(V56.3)--(V56.5)} are
\Cref{lem:V56-own-edge-star,lem:V56-absent-edge-mean}.
Item \textup{(V56.6)} is
\Cref{thm:V56-noncollision-connected,
cor:V56-noncollision-32}.  Item \textup{(V56.7)} is
\Cref{prop:V56-collision-reduction}.
\end{proof}

\begingroup
\small
\begin{longtable}{>{\raggedright\arraybackslash}p{0.29\textwidth}
                  >{\raggedright\arraybackslash}p{0.15\textwidth}
                  >{\raggedright\arraybackslash}p{0.48\textwidth}}
\toprule
\textbf{V56 row}&\textbf{Status}&\textbf{Advance or obstruction}\\
\midrule
\endfirsthead
\toprule
\textbf{V56 row}&\textbf{Status}&\textbf{Advance or obstruction}\\
\midrule
\endhead
Exact multi-attachment algebra
& \MGclosed
& Three disjoint exact pair banks give seven attachment monomials,
  with no coordinate-subset or Bell sum.\\[0.35em]
Simultaneous center normalization
& \MGclosed
& A routed center uses its own defect when present and its retained
  absent-edge mean otherwise.\\[0.35em]
Noncollision \(T\)-connected sector
& \MGclosed
& All exact-pair attachment monomials satisfy quartic tree weights
  with arbitrary coefficients.\\[0.35em]
Noncollision \([3,2]\) terms
& \MGclosed
& Every such term is contained in the proved assembled
  \(T\)-connected sector.\\[0.35em]
Mean-complement collision branch
& \MGopen
& Complementary original-input masses require simultaneous retained
  product debits.\\[0.35em]
Double-pair collision branch
& \MGopen
& Two disjoint local defects can both form cold outputs and require
  aggregate raw compression.\\[0.35em]
Complete \([3,2]\), \([3,3]\), and \([4]\)
& \MGopen
& Only collision attachments remain for \([3,2]\); the other two
  topologies remain suspended.\\[0.35em]
Full Yang--Mills mass gap
& \MGopen
& Global quartic/all-order MTA, WUFC/CPM, nonlinear construction,
  continuum OS existence/nontriviality, and the spectral passage
  remain downstream.\\
\bottomrule
\end{longtable}
\endgroup

\begin{openproblem}[V57 mandate: close collision attachments]
\label{op:V57}
\begin{enumerate}[label=\textup{(V57.\arabic*)},leftmargin=5.5em]
\item Start from the exact split
      \eqref{eq:V56-collision-factor}; do not estimate the complement
      moment before this split.
\item In the mean-complement branch, retain the \(b,c\) scalar
      multipliers through the V54 center response and use their second
      output debit only when their own collision output carries the
      final mass.
\item In the double-pair branch, form the aggregate invariant
      projector for
      \((R_o\otimes R_a)\otimes(R_b\otimes R_c)\) before trace norm.
      Apply no separate representation-formation probability.
\item Route the \(oa\) and \(bc\) defects together with one mark from
      the intact ternary junction.  Choose the fourth tree mark only
      after identifying which original vertex needs a denominator.
\item Extend the V56 center-selection rule to several collision
      coordinates and prove a fixed-type, not coordinate-subset,
      expansion.
\item Test two simultaneous cold singlets, a long dual-singlet ternary
      core, every omitted row, and all three choices of the selected
      attachment target.
\item Declare complete \([3,2]\) only after both branches and all
      collision multiplicities satisfy the exact V46 coefficient
      normalization.  Keep \([3,3]\) and \([4]\) suspended.
\item Preserve only the newest YM file.  The next companion audit
      remains V60.
\end{enumerate}
\end{openproblem}

\begin{firewall}[End-of-V56 claim boundary]
\label{fire:V56-claim-boundary}
V56 closes the exact noncollision \(T\)-connected quartic sector and
thereby assigns every noncollision \([3,2]\) term to a proved
arbitrary-coefficient MTA operator.  It does not prove either
collision branch in \eqref{eq:V56-collision-factor}, complete
\([3,2]\), \([3,3]\), \([4]\), global \(m=4\), all-order MTA, WUFC,
CPM, the nonlinear Perron packet, finite-edge margins, constructive
capacity, protected-sector floors, local-algebra noncollapse,
reflection positivity, a nontrivial continuum OS state, or a positive
Yang--Mills spectral gap.  The full Yang--Mills mass gap is not proved.
\end{firewall}

\section*{Version V56 append-only record}
\addcontentsline{toc}{section}{Version V56 append-only record}
The complete V55 source was retained before this continuation.  Its
SHA--256 digest was
\begin{center}\scriptsize\ttfamily
703e2035c7fffaf9e437f934a88448ea4ed35b19c30c32745806b5cf0ae8a8b6.
\end{center}
No mathematical content of V55 was altered.  On the next iteration,
retain this full file, append before its unique active document-closing
command, increment the filename to \texttt{\_V57.tex}, and remove only
the superseded V56 file from this Yang--Mills note series.

% =============================================================================
% END APPENDED VERSION V56 MATERIAL
% =============================================================================

% =============================================================================
% BEGIN APPENDED VERSION V57 MATERIAL
% =============================================================================

\chapter{V57: Mean-collision closure and the row-allocated ternary
cut gate}
\label{ch:V57}

\section{Purpose and correction of the collision route}
\label{sec:V57-purpose}

V56 reduced a collision coordinate to
\[
 q_{b,e}q_{c,e}\mathbf d_{oa,e}
 \quad\hbox{or}\quad
 \mathbf d_{oa,e}\otimes\mathbf d_{bc,e}.
\]
The first branch has the missing multipliers and can be closed by a
same-coordinate version of the absent-edge mean argument.  The second
branch supplies two genuine pair edges.  Together with the intact
ternary cumulant, it has the graph of a path, but one must first turn
the ternary cumulant into a cut joining \(a\) to an \emph{original}
row \(b\) or \(c\).  A binary tax against the composite carrier
\(R_b\otimes R_c\) does not by itself give either original denominator.
V57 proves the exact cut identity and its one-mark bounds and records
the remaining row-allocated capacity as a precise coefficient gate.

\begin{firewall}[V57 collision scope]
\label{fire:V57-scope}
V57 closes the mean-complement collision branch, including arbitrary
noncollision attachment spectators.  It does not claim the
double-pair branch until the composite ternary cut has been compressed
with an original-row allocation.  In particular, two cold singlets are
not assigned independent representation-formation probabilities.
\end{firewall}

\section{Closure of the mean-complement branch}
\label{sec:V57-mean-collision}

Fix one collision coordinate \(e\) with selected block
\(\{o,a\}\), and let \(b,c=T\setminus\{a\}\).  In the first term of
\eqref{eq:V56-collision-factor}, put
\[
 r_{\rm mc}:=
 q_{b,e}q_{c,e}\delta_{oa,e}.
\]
All other noncollision attachment factors are retained as in
\eqref{eq:V56-surplus-factor}.

\begin{lemma}[Same-coordinate means retain their debits]
\label{lem:V57-same-coordinate-debits}
Let \(i\in\{b,c\}\), put \(x_i=A_{i,e}\), and let
\(\widehat\Xi_i\) be its mass after the collision coordinate is
punctured.  Thus
\(\Xi_i=\widehat\Xi_i+x_i\).  In every routed center-\(i\) response,
the factor \(q_{i,e}\) replaces
\(\widehat\Xi_i^{-1}\) by \(\Xi_i^{-1}\), up to a universal constant.
This remains true while \(\delta_{oa,e}\), the other complement mean,
and every noncollision attachment are kept assembled.
\end{lemma}

\begin{proof}
If \(x_i\le\widehat\Xi_i\), then
\(\Xi_i\le2\widehat\Xi_i\), so no multiplier estimate is needed.  If
\(x_i>\widehat\Xi_i\) and \(s_\alpha\Xi_i\le M\), use the raw
center-\(i\) side
\eqref{eq:V54-routed-raw-side}:
\[
 s_\alpha\widehat M_i
 \le M\frac{\widehat M_i}{\Xi_i}.
\]
Finally suppose \(x_i>\widehat\Xi_i\) and
\(s_\alpha\Xi_i>M\).  The one-coordinate product bounds give
\[
 |q_{i,e}|\le\frac C{s_\alpha x_i},
 \qquad
 s_\alpha\Xi(S_{i,e})|q_{i,e}|
 \le\frac C{s_\alpha x_i}.
\]
Use the first inequality on core, attachment, and other spectator
outputs and the second on the output of the \(i\)-factor itself.
Relative to the raw response this supplies
\(C/(s_\alpha\Xi_i)\), hence again
\(C\widehat M_i/\Xi_i\).

Although \(q_{i,e}\) and \(\delta_{oa,e}\) belong to the same
probability coordinate, their replica blocks act on disjoint carrier
factors.  Their product in
\eqref{eq:V56-collision-factor} is formed before recoupling, and the
final coefficient is compressed by the single V46 raw projector.
Thus the preceding scalar debit is not used twice.
\end{proof}

\begin{theorem}[Mean-complement collision MTA]
\label{thm:V57-mean-collision}
Every term in the mean-complement branch
\[
 q_{b,e}q_{c,e}\mathbf d_{oa,e}
\]
of \eqref{eq:V56-collision-factor} satisfies its quartic marked-tree
bound for arbitrary trace-class Fourier coefficients.  The conclusion
is uniform in supports, representations, outputs, multiplicities, and
in any additional noncollision attachment factors.
\end{theorem}

\begin{proof}
Route the punctured ternary response by
\eqref{eq:V54-center-routing}.  For center \(a\), apply the complete
attachment promotion of V54--V55; the complement means have modulus at
most one and their output first moments are admissible spectators.  For
center \(b\) or \(c\), use the same selected \(oa\) promotion and then
\Cref{lem:V57-same-coordinate-debits} to replace the punctured center
mass by its full original mass.  If the selected attachment is
common-heavy, the response is instead assigned to the star centered
at \(a\), exactly as in
\Cref{lem:V56-absent-edge-mean}.

Every additional exact-pair attachment is routed by
\eqref{eq:V56-selected-attachment} and is admissible by
\eqref{eq:V56-surplus-admissibility}.  Finally apply
\eqref{eq:V54-puncture-normalization} and aggregate raw pinching.
Each resulting right-hand term is a star or path already present in
the quartic MTA sum.
\end{proof}

\section{The intact ternary cumulant as a binary cut}
\label{sec:V57-ternary-cut}

Let \(F_a,F_b,F_c\) be the three operator-valued product
representations on a punctured core.  They act on distinct input
carriers.  Define
\[
 \operatorname{Cov}(H,K)
 :=\mathbb E(HK)-\mathbb EH\,\mathbb EK.
\]

\begin{theorem}[Exact three-covariance cut identity]
\label{thm:V57-cut-identity}
The complete ternary connected operator is
\begin{align}
 \boxed{
 \mathbf K_T^\circ
 ={}&
 \operatorname{Cov}(F_a,F_bF_c)
 -\mathbb EF_b\,\operatorname{Cov}(F_a,F_c)
 \notag\\
 &-\mathbb EF_c\,\operatorname{Cov}(F_a,F_b).}
 \label{eq:V57-cut-identity}
\end{align}
The identity is valid with arbitrary spectator amplification and
before any final recoupling.
\end{theorem}

\begin{proof}
Expand the three covariances in the displayed order.  Since the input
carriers are distinct, their factors commute, and the right side is
\[
 \mathbb E(F_aF_bF_c)
 -\mathbb EF_a\,\mathbb E(F_bF_c)
 -\mathbb EF_b\,\mathbb E(F_aF_c)
 -\mathbb EF_c\,\mathbb E(F_aF_b)
 +2\mathbb EF_a\,\mathbb EF_b\,\mathbb EF_c.
\]
This is the five-term definition of the ternary cumulant.  Tensoring
every term with the same spectator identity preserves the equality.
\end{proof}

For the punctured core put
\begin{equation}
 H_{a\mid bc}:=
 \widehat H_{ab}+\widehat H_{ac}.
 \label{eq:V57-cut-weight}
\end{equation}

\begin{theorem}[One-mark cap and first moment of the ternary cut]
\label{thm:V57-cut-cap}
Uniformly in all punctured supports and representations,
\begin{align}
 \|\mathbf K_T^\circ\|_{\rm op}
 &\le
 C\min\{s_\alpha H_{a\mid bc},1\},
 \label{eq:V57-cut-cap}\\
 s_\alpha\Xi_T^{\rm out}
 \|\mathbf K_T^{\circ,\rm out}\|_{\rm op}
 &\le C.
 \label{eq:V57-cut-first-moment}
\end{align}
Moreover, after resolving the core output before the external
collision factors,
\begin{equation}
 \frac{\Xi_T^{\rm out}
       \|\mathbf K_T^{\circ,\rm out}\|_{\rm op}}
      {1-q_{\alpha,\boldsymbol T_{\rm core}}}
 \le C H_{a\mid bc}.
 \label{eq:V57-cut-quotient}
\end{equation}
The cap is allocated at the level of majorants as
\begin{equation}
 H_{a\mid bc}=\widehat H_{ab}+\widehat H_{ac},
 \label{eq:V57-original-row-allocation}
\end{equation}
but \eqref{eq:V57-original-row-allocation} alone is not yet an
original-row coefficient tax.
\end{theorem}

\begin{proof}
Apply \Cref{lem:V50-operator-covariance} to
\eqref{eq:V57-cut-identity}.  The two-sided influence of \(F_bF_c\)
is at most \(\gamma_{b,e}+\gamma_{c,e}\), while singleton means are
contractions.  Hence
\[
 \|\mathbf K_T^\circ\|_{\rm op}
 \le
 C\sum_e\gamma_{a,e}
       (\gamma_{b,e}+\gamma_{c,e})
 \le Cs_\alpha
       (\widehat H_{ab}+\widehat H_{ac}).
\]
The five-term definition also gives a universal operator cap, proving
\eqref{eq:V57-cut-cap}.

For \eqref{eq:V57-cut-quotient}, first resolve \(F_bF_c\) into its raw
intermediate irreducibles.  Each summand of
\(\operatorname{Cov}(F_a,F_bF_c)\) is then a genuine binary defect.
Apply \eqref{eq:V46-binary-quotient} and allocate the composite
Casimir by
\(A_{bc,e}\le C(A_{b,e}+A_{c,e})\).  The other two covariances are
already binary defects with weights \(\widehat H_{ac}\) and
\(\widehat H_{ab}\).  Intermediate raw pinching sums the three
resolutions without a path or dimension count, proving
\eqref{eq:V57-cut-quotient}.

Finally expand the complete ternary cumulant only into its five global
row-partition moments.  On every resolved term, split the output mass
among its fixed coordinate banks and use
\eqref{eq:V46-product-first-moment}.  There are five terms, so
\eqref{eq:V57-cut-first-moment} follows.  The influence display already
separates the two original-row sums and proves
\eqref{eq:V57-original-row-allocation}.
\end{proof}

\section{Aggregate raw compression at a double-pair collision}
\label{sec:V57-aggregate-collision}

In the second branch of \eqref{eq:V56-collision-factor}, resolve
\(R_o\otimes R_a\) and \(R_b\otimes R_c\) simultaneously.  Let
\(P_{\lambda,\mu}\) be the corresponding product of multiplicity
projections, with the punctured-core projector included, and put
\[
 X_{\lambda,\mu}
 :=P_{\lambda,\mu}X P_{\lambda,\mu}.
\]

\begin{lemma}[One collision projector, not two formation laws]
\label{lem:V57-aggregate-collision-pinching}
For arbitrary trace-class input coefficient matrices,
\begin{equation}
 \sum_{\lambda,\mu}\|X_{\lambda,\mu}\|_1
 \le\|X\|_1
 \le\prod_{i=1}^4\|B_i\|_1.
 \label{eq:V57-aggregate-collision-pinching}
\end{equation}
After the final recoupling, the coefficient normalization remains
\(d_{\rm in}/d_{\boldsymbol U}\), and the BFA factor cancels
\(d_{\boldsymbol U}\).
\end{lemma}

\begin{proof}
The projections \(P_{\lambda,\mu}\) are mutually orthogonal and sum to
the identity on the resolved collision carrier.  Trace-norm
contractivity of this single pinching gives the first inequality.
The second is the tensor-product trace norm.  Apply
\eqref{eq:V46-exact-multiplier-normalization} after the collision and
core factors have been assembled to obtain the last assertion.
Neither \(d_\lambda/(d_od_a)\) nor
\(d_\mu/(d_bd_c)\) is used as coefficient mass.
\end{proof}

\section{The exact row-allocated double-pair gate}
\label{sec:V57-row-gate}

Let \(h_{oa,e}=A_{o,e}A_{a,e}\) and
\(h_{bc,e}=A_{b,e}A_{c,e}\).  The two possible path weights obtained
from \eqref{eq:V57-original-row-allocation} are
\begin{equation}
 \mathsf P_{ab}
 :=
 \frac{h_{oa,e}\widehat H_{ab}h_{bc,e}}
      {\Xi_a\Xi_b},
 \qquad
 \mathsf P_{ac}
 :=
 \frac{h_{oa,e}\widehat H_{ac}h_{bc,e}}
      {\Xi_a\Xi_c}.
 \label{eq:V57-collision-path-weights}
\end{equation}
They correspond respectively to the paths
\(o-a-b-c\) and \(o-a-c-b\).

\begin{definition}[Row-allocated cut capacity]
\label{def:V57-RACC}
Write \(\operatorname{RACC}_{a\mid bc}\) for the following statement.
After the covariance identity \eqref{eq:V57-cut-identity}, nested raw
resolution, and multiplication by the two collision defects, there
exist nonnegative original-row majorants
\(\mathcal C_{ab}\) and \(\mathcal C_{ac}\) such that
\begin{align}
 \mathcal Q_{\rm dp}
 &\le\mathcal C_{ab}+\mathcal C_{ac},
 \label{eq:V57-RACC-split}\\
 \mathcal C_{ab}
 &\le C\mathsf P_{ab},
 &
 \mathcal C_{ac}
 &\le C\mathsf P_{ac}.
 \label{eq:V57-RACC-bounds}
\end{align}
Here \(\mathcal Q_{\rm dp}\) contains the complete final output mass
and the single final resolvent.  The split must occur before trace
norm on the final coefficient and must use
\eqref{eq:V57-aggregate-collision-pinching}.
\end{definition}

\begin{theorem}[RACC is exactly sufficient for the double-pair branch]
\label{thm:V57-RACC-sufficiency}
If \(\operatorname{RACC}_{a\mid bc}\) holds for every collision
coordinate, every punctured core, and both choices in
\eqref{eq:V57-original-row-allocation}, then the complete double-pair
branch satisfies the quartic marked-tree atlas bound.
\end{theorem}

\begin{proof}
Insert \eqref{eq:V57-RACC-bounds}.  The two right-hand sides are
exactly the MTA weights of the two paths in
\eqref{eq:V57-collision-path-weights}.  Sum the aggregate raw blocks
by \eqref{eq:V57-aggregate-collision-pinching}; the final dimension
cancels by the V46 normalization.  Additional noncollision
attachments are routed by V56 and have uniform modulus and first
moment.  There are only the two row allocations at a fixed collision,
so no support-dependent factor is introduced.
\end{proof}

\begin{proposition}[The scalar part of RACC is already proved]
\label{prop:V57-RACC-scalar}
If arbitrary coefficient compression is replaced by a spectral
supremum, then
\begin{equation}
 \mathcal Q_{\rm dp}
 \le
 C s_\alpha^2
 h_{oa,e}h_{bc,e}H_{a\mid bc}
 \label{eq:V57-double-pair-raw}
\end{equation}
and every branch in which both prospective internal vertices are
thermal satisfies \eqref{eq:V57-RACC-bounds}.
\end{proposition}

\begin{proof}
Use the two pair-defect caps and
\eqref{eq:V57-cut-cap}; the low-output resolvent leaves
\(Cs_\alpha^2h_{oa,e}h_{bc,e}H_{a\mid bc}\).
On high outputs, split the output mass among the \(oa\) pair, the
\(bc\) pair, and the core.  For a pair output use its binary quotient,
the cap on the other pair, and \eqref{eq:V57-cut-cap}.  For a core
output use \eqref{eq:V57-cut-quotient} and both pair caps.  Each piece
is bounded by
\(Cs_\alpha^2h_{oa,e}h_{bc,e}H_{a\mid bc}\).
Split \(H_{a\mid bc}\) by
\eqref{eq:V57-original-row-allocation}.  If
\(s_\alpha\Xi_a,s_\alpha\Xi_r\le M\), \(r=b\) or \(c\), then
\[
 s_\alpha^2
 \le\frac{M^2}{\Xi_a\Xi_r},
\]
which proves the corresponding RACC bound.
\end{proof}

\begin{assessment}[The remaining cold--cold coefficient problem]
\label{ass:V57-cold-cold}
The unresolved RACC regime has a hot internal vertex and a collision
pair or ternary cut capable of forming a low output.  The scalar cap
\eqref{eq:V57-double-pair-raw} then lacks an original-input
denominator.  Applying V53 to the composite carrier
\(R_b\otimes R_c\) gives a denominator for its total mass, not
automatically for \(\Xi_b\) or \(\Xi_c\).  The next proof must resolve
the covariance martingale increment and the \(bc\) pair block in one
aggregate raw projector, then apply the one-edge tax to the actual
original row selected by
\eqref{eq:V57-original-row-allocation}.  This is the first exact
unproved inequality; no LR multiplicity or scalar thermal estimate
remains hidden behind it.
\end{assessment}

\section{V57 ledger and next acquisition}
\label{sec:V57-ledger}

\begin{theorem}[Integrated V57 statement]
\label{thm:V57-integrated}
Preserving V1--V56, V57 proves:
\begin{enumerate}[label=\textup{(V57.\arabic*)},leftmargin=5.5em]
\item same-coordinate retention of the complement multiplier debits;
\item the full arbitrary-coefficient mean-complement collision branch;
\item the exact three-covariance cut identity
      \eqref{eq:V57-cut-identity};
\item the one-mark cut cap and output first moment
      \eqref{eq:V57-cut-cap}--%
      \eqref{eq:V57-cut-first-moment};
\item aggregate, rather than separate, raw collision pinching;
\item reduction of the double-pair branch to
      \(\operatorname{RACC}_{a\mid bc}\); and
\item the scalar and internally thermal parts of RACC.
\end{enumerate}
\end{theorem}

\begin{proof}
Items \textup{(V57.1)--(V57.2)} are
\Cref{lem:V57-same-coordinate-debits,
thm:V57-mean-collision}.  Items
\textup{(V57.3)--(V57.4)} are
\Cref{thm:V57-cut-identity,thm:V57-cut-cap}.  Item
\textup{(V57.5)} is
\Cref{lem:V57-aggregate-collision-pinching}.  Item
\textup{(V57.6)} is
\Cref{def:V57-RACC,thm:V57-RACC-sufficiency}; item
\textup{(V57.7)} is \Cref{prop:V57-RACC-scalar}.
\end{proof}

\begingroup
\small
\begin{longtable}{>{\raggedright\arraybackslash}p{0.29\textwidth}
                  >{\raggedright\arraybackslash}p{0.15\textwidth}
                  >{\raggedright\arraybackslash}p{0.48\textwidth}}
\toprule
\textbf{V57 row}&\textbf{Status}&\textbf{Advance or obstruction}\\
\midrule
\endfirsthead
\toprule
\textbf{V57 row}&\textbf{Status}&\textbf{Advance or obstruction}\\
\midrule
\endhead
Mean-complement collision
& \MGclosed
& Same-coordinate singleton means retain the first and second debits
  needed to restore punctured center masses.\\[0.35em]
Ternary binary-cut identity
& \MGclosed
& Three exact covariances expose the two original-row influence
  weights without expanding local five-term cumulants.\\[0.35em]
Aggregate collision compression
& \MGclosed
& Both disjoint pair outputs and the core are pinched by one raw
  projector; no product of formation probabilities is used.\\[0.35em]
Thermal double-pair branch
& \MGclosed
& Three defect caps and one resolvent give the two path weights when
  both internal vertices are thermal.\\[0.35em]
Hot/cold RACC branch
& \MGopen
& A composite cut must be converted to an original-row endpoint tax
  inside the aggregate raw projector.\\[0.35em]
Complete collision and \([3,2]\)
& \MGopen
& The mean branch is closed; double-pair collision waits only on
  RACC and its multi-coordinate assembly.\\[0.35em]
Full Yang--Mills mass gap
& \MGopen
& Global quartic/all-order MTA, WUFC/CPM, nonlinear construction,
  continuum OS existence/nontriviality, and the spectral passage
  remain downstream.\\
\bottomrule
\end{longtable}
\endgroup

\begin{openproblem}[V58 mandate: original-row cut compression]
\label{op:V58}
\begin{enumerate}[label=\textup{(V58.\arabic*)},leftmargin=5.5em]
\item Retain the exact covariance identity
      \eqref{eq:V57-cut-identity}.  Resolve
      \(\operatorname{Cov}(F_a,F_bF_c)\) first in the \(bc\) carrier,
      but do not divide by its intermediate dimension.
\item Decompose each martingale increment into its \(ab\) and \(ac\)
      original-row influences before applying trace norm.  Prove an
      aggregate trace tax for the selected influence subspace.
\item Combine that tax with the \(bc\) collision defect in the same
      raw projector \eqref{eq:V57-aggregate-collision-pinching}.
\item Prove the two RACC bounds separately by a thermal,
      selected-heavy, retained-mean-heavy, and common-heavy split at
      the relevant original internal vertex.
\item Test a cold \(oa\) singlet, a cold \(bc\) singlet, and a cold
      composite cut simultaneously.  No output Casimir may stand in
      for an original input mass.
\item Once one collision closes, aggregate several collision
      coordinates by their three disjoint-pair patterns without a
      coordinate-subset expansion.
\item Declare complete \([3,2]\) only after RACC and the
      multi-collision aggregation are proved.  Keep \([3,3]\) and
      \([4]\) suspended.
\item Preserve only the newest YM file.  The next companion audit
      remains V60.
\end{enumerate}
\end{openproblem}

\begin{firewall}[End-of-V57 claim boundary]
\label{fire:V57-claim-boundary}
V57 closes mean-complement collisions, proves the exact ternary cut
identity and its one-mark bounds, and reduces double-pair collisions to
the explicit row-allocated capacity
\(\operatorname{RACC}_{a\mid bc}\).  It does not prove the hot/cold
part of RACC, the multi-collision sum, complete \([3,2]\),
\([3,3]\), \([4]\), global \(m=4\), all-order MTA, WUFC, CPM, the
nonlinear Perron packet, finite-edge margins, constructive capacity,
protected-sector floors, local-algebra noncollapse, reflection
positivity, a nontrivial continuum OS state, or a positive Yang--Mills
spectral gap.  The full Yang--Mills mass gap is not proved.
\end{firewall}

\section*{Version V57 append-only record}
\addcontentsline{toc}{section}{Version V57 append-only record}
The complete V56 source was retained before this continuation.  Its
SHA--256 digest was
\begin{center}\scriptsize\ttfamily
379dda8ccc5230e6d3362442168e1a8360be8e052fec8b8661789ce2be6f0736.
\end{center}
No mathematical content of V56 was altered.  On the next iteration,
retain this full file, append before its unique active document-closing
command, increment the filename to \texttt{\_V58.tex}, and remove only
the superseded V57 file from this Yang--Mills note series.

% =============================================================================
% END APPENDED VERSION V57 MATERIAL
% =============================================================================

% =============================================================================
% BEGIN APPENDED VERSION V58 MATERIAL
% =============================================================================

\chapter{V58: Sequential aggregate collision taxes and the
selected-edge RACC region}
\label{ch:V58}

\section{Purpose}
\label{sec:V58-purpose}

V57 isolated the double-pair collision gate, but proved the cut cap and
cut quotient only with the aggregate influence
\[
 \widehat H_{a\mid bc}:=\widehat H_{ab}+\widehat H_{ac}.
\]
It did not prove a separate quotient estimate for either summand.  The
present continuation therefore keeps this aggregate cut quantity.  Its
other purpose is to settle exactly how the two disjoint pair taxes may
be multiplied.  The input to the multilinear coefficient calculus is
row-factorized before final recoupling, although each row coefficient
may be entangled with arbitrary punctured-core spectators.  That
specific structure, rather than an unproved completely bounded norm,
is sufficient for a product tax.

\begin{firewall}[Scope of V58]
\label{fire:V58-scope}
The result below does not assume probabilistic independence of the two
pair outputs.  It uses the deterministic tensor factorization of the
four input-row coefficients before final recoupling.  It does
\emph{not} assert an \(\mathcal S_1\)-map bound for an arbitrary
operator entangled across the partition \((oa)\mid(bc)\).  It also
does not split the aggregate cut quotient into \(ab\) and \(ac\)
pieces.  V58 closes only the summed path target in the region where
\(a\), and every active endpoint among \(b,c\), is paid on its
incident collision edge or is thermal.
\end{firewall}

\section{Composition of two aggregate trace taxes}
\label{sec:V58-product-tax}

For a pair \(e=uv\), let
\(\mathcal C_{e,\lambda}(B_u,B_v)\) be the
multiplicity-resolved compression of V52.  Here
\[
 B_i\in\mathcal S_1(H_i\otimes\mathcal K_i)
\]
and \(\mathcal K_i\) is an arbitrary row spectator.  Let
\(w_{e,\lambda}\ge0\), and suppose that the V52/V53 bilinear
functional obeys
\begin{equation}
 \sum_{\lambda\in\Omega_e}w_{e,\lambda}
 \|\mathcal C_{e,\lambda}(B_u,B_v)\|_1
 \le \tau_e\|B_u\|_1\|B_v\|_1
 \label{eq:V58-individual-tax}
\end{equation}
uniformly in the two spectator spaces.

\begin{lemma}[Row-factorized taxes compose multiplicatively]
\label{lem:V58-tax-composition}
Let \(e_0=oa\), \(e_1=bc\), and assume
\eqref{eq:V58-individual-tax} with constants \(\tau_{e_0}\) and
\(\tau_{e_1}\).  For each \((\lambda,\mu)\), let
\(\mathscr R_{\lambda,\mu}\) denote any final recoupling followed by
an aggregate orthogonal pinching, regarded as a contraction into a
trace-class direct sum, and put
\[
 X_{\lambda,\mu}:=
 \mathscr R_{\lambda,\mu}\!\left(
 \mathcal C_{oa,\lambda}(B_o,B_a)\otimes
 \mathcal C_{bc,\mu}(B_b,B_c)\right).
\]
Then
\begin{equation}
 \boxed{
 \sum_{\lambda,\mu}
 w_{oa,\lambda}w_{bc,\mu}\|X_{\lambda,\mu}\|_1
 \le
 \tau_{oa}\tau_{bc}\prod_{i=o,a,b,c}\|B_i\|_1.}
 \label{eq:V58-product-tax}
\end{equation}
\end{lemma}

\begin{proof}
Contractivity of \(\mathscr R_{\lambda,\mu}\), multiplicativity of the
trace norm on tensor products, and Tonelli's theorem for the
nonnegative series give
\begin{align*}
 &\sum_{\lambda,\mu}w_{oa,\lambda}w_{bc,\mu}
      \|X_{\lambda,\mu}\|_1\\
 &\quad\le
 \left(\sum_\lambda w_{oa,\lambda}
 \|\mathcal C_{oa,\lambda}(B_o,B_a)\|_1\right)
 \left(\sum_\mu w_{bc,\mu}
 \|\mathcal C_{bc,\mu}(B_b,B_c)\|_1\right).
\end{align*}
Apply \eqref{eq:V58-individual-tax} to the two factors.  Notice that
the arbitrary row spectators remain inside their respective pair
compressions; no step moves them through a projection.
\end{proof}

For a resolved pair edge \(e\), retain the weighted defect size
\[
 \mathfrak a_e
 :=
 |\delta_e|+s_\alpha\Xi(S_e)|\delta_e|.
\]
For nonzero row coefficients define the normalized bilinear response
\begin{equation}
 \mathfrak A_e[B_u,B_v]
 :=\frac{
 \sum_\lambda\mathfrak a_{e,\lambda}
 \|\mathcal C_{e,\lambda}(B_u,B_v)\|_1}
 {\|B_u\|_1\|B_v\|_1},
 \label{eq:V58-normalized-edge-response}
\end{equation}
and set it to zero if either coefficient vanishes.  This is a
bilinear-response ratio, not a norm on arbitrary entangled inputs.

\begin{theorem}[Two-pair amplified formation tax]
\label{thm:V58-two-pair-tax}
Let \(e_0=oa\) and \(e_1=bc\) be the two disjoint pair blocks of a
collision.  If, for \(v=a,r\) with \(r=b\) or \(c\), the
input-\(v\) mass \(x_v\) on its incident collision pair satisfies
\begin{equation}
 x_v\ge\varepsilon\Xi_v,
 \label{eq:V58-both-selected-heavy}
\end{equation}
then
\begin{equation}
 \boxed{
 \sum_{\lambda,\mu}
 \mathfrak a_{oa,\lambda}\mathfrak a_{bc,\mu}
 \|X_{\lambda,\mu}\|_1
 \le
 C_\varepsilon
 \frac{h_{oa}h_{bc}}
      {s_\alpha^2\Xi_a^2\Xi_r^2}
 \prod_{i=o,a,b,c}\|B_i\|_1.}
 \label{eq:V58-two-pair-tax}
\end{equation}
Here \(X_{\lambda,\mu}\) has the row-factorized pre-recoupling form
of \Cref{lem:V58-tax-composition}.  All rowwise punctured-core
carriers may be included in the \(\mathcal K_i\), and the final
recoupling may be any aggregate trace-norm contraction.
\end{theorem}

\begin{proof}
For each edge, \Cref{cor:V53-one-edge-tax} is precisely
\eqref{eq:V58-individual-tax}, with
\[
 \tau_e\le C_\varepsilon
 \frac{h_e}{s_\alpha\Xi_v^2},
\]
under \eqref{eq:V58-both-selected-heavy}.  Apply
\Cref{lem:V58-tax-composition}.  The final aggregate raw pinching of
\eqref{eq:V57-aggregate-collision-pinching} is part of
\(\mathscr R_{\lambda,\mu}\), so it is applied only after both
formation taxes have been retained.
\end{proof}

\begin{remark}[Exact entanglement boundary]
\label{rem:V58-entangled-coefficients}
The coefficient matrix on one input row may entangle its collision
carrier with all core coordinates.  Such entanglement is allowed
because it stays inside \(B_i\).  What is used is the exact multilinear
input tensor
\(B_o\otimes B_a\otimes B_b\otimes B_c\) before global recoupling.
Equation \eqref{eq:V58-product-tax} makes no claim for an arbitrary
trace-class \(X\) already entangled across different input rows.  A
future use outside the row-factorized coefficient calculus would
therefore require a new theorem.
\end{remark}

\section{A universal weighted three-block response}
\label{sec:V58-three-block}

Write
\[
 H_\Sigma:=\widehat H_{ab}+\widehat H_{ac},\qquad
 h_a:=h_{oa},\qquad h_b:=h_{bc}.
\]
Let \(N_B:=\prod_{i=o,a,b,c}\|B_i\|_1\), and let
\(\mathcal R_{\rm dp}\) be the complete double-pair response divided
by \(N_B\), with value zero when \(N_B=0\).

\begin{lemma}[Three-block response before endpoint payment]
\label{lem:V58-three-block-response}
The complete aggregate double-pair response satisfies
\begin{equation}
 \boxed{
 \mathcal R_{\rm dp}
 \le C H_\Sigma\,
 \mathfrak A_{oa}[B_o,B_a]\,
 \mathfrak A_{bc}[B_b,B_c].}
 \label{eq:V58-three-block-response}
\end{equation}
For either normalized edge response one may use its raw cap, or the V53
formation tax whenever the selected-heavy hypothesis holds.
\end{lemma}

\begin{proof}
On low final outputs, use the single resolvent cap
\(C/s_\alpha\), the aggregate cut cap
\(Cs_\alpha H_\Sigma\) from \eqref{eq:V57-cut-cap}, and the modulus
parts of the two weighted defect functionals.  The row-factorized
product estimate is \Cref{lem:V58-tax-composition}.

On high outputs, split the output mass among the two pair blocks, the
punctured core, and ordinary spectators.  If a pair block carries the
mass, use the output-weighted part of its \(\mathfrak A\), the modulus
part of the other, and the cut cap.  If the core carries the mass, use
the aggregate quotient \eqref{eq:V57-cut-quotient} and both pair
moduli.  An ordinary spectator is paid by its first moment.  The
denominator has a fixed floor.  In every branch the two pair
compressions occur before final recoupling and aggregate pinching, so
\Cref{lem:V58-tax-composition} applies.  Division by \(N_B\) proves
\eqref{eq:V58-three-block-response}.
\end{proof}

\section{Closure of the thermal/selected collision region}
\label{sec:V58-selected-region}

Say that vertex \(a\) is \emph{edge-paid} on \(oa\) if either
\[
 s_\alpha\Xi_a\le M
 \quad\hbox{or}\quad
 x_a\ge\varepsilon\Xi_a.
\]
Define edge-paid at \(b\) or \(c\) analogously using its input mass on
the \(bc\) pair.

\begin{theorem}[All-relevant-endpoint aggregate path theorem]
\label{thm:V58-selected-RACC}
Assume that \(a\) is edge-paid and that, for every
\(r\in\{b,c\}\) with \(\widehat H_{ar}>0\), the vertex \(r\) is
edge-paid.  Then the \emph{total} double-pair response obeys
\begin{equation}
 \boxed{
 \mathcal R_{\rm dp}
 \le C\left(
 \frac{h_{oa}\widehat H_{ab}h_{bc}}{\Xi_a\Xi_b}
 +
 \frac{h_{oa}\widehat H_{ac}h_{bc}}{\Xi_a\Xi_c}
 \right).}
 \label{eq:V58-selected-RACC-total}
\end{equation}
\end{theorem}

\begin{proof}
At a thermal endpoint \(v\) of an edge \(e_v\), use
\begin{equation}
 \mathfrak A_{e_v}\le Cs_\alpha h_{e_v}.
 \label{eq:V58-thermal-edge-cap}
\end{equation}
At a selected-heavy nonthermal endpoint use
\begin{equation}
 \mathfrak A_{e_v}
 \le C_\varepsilon
 \frac{h_{e_v}}{s_\alpha\Xi_v^2}.
 \label{eq:V58-hot-edge-tax}
\end{equation}

Both are instances of the bilinear response
\eqref{eq:V58-normalized-edge-response}.

Fix an active \(r\in\{b,c\}\).  Apply the bound at \(a\) to the
\(oa\) response and the bound at \(r\) to the same \(bc\) response.
If both endpoints are thermal, their product is
\(Cs_\alpha^2h_{oa}h_{bc}\), and
\[
 s_\alpha^2\le\frac{M^2}{\Xi_a\Xi_r}.
\]
If \(a\) is selected-heavy and \(r\) thermal, the product is
\[
 C\frac{h_{oa}h_{bc}}{\Xi_a^2}.
\]
Here \(s_\alpha\Xi_a>M\) in the nonthermal subcase, while
\(s_\alpha\Xi_r\le M\), so \(\Xi_a>\Xi_r\).  Thus the last display
is at most \(Ch_{oa}h_{bc}/(\Xi_a\Xi_r)\).  The case with \(a,r\)
interchanged is identical.

If both endpoints are selected-heavy and nonthermal,
\Cref{thm:V58-two-pair-tax} gives
\[
 C\frac{h_{oa}h_{bc}}
 {s_\alpha^2\Xi_a^2\Xi_r^2}
 \le
 \frac C{M^2}
 \frac{h_{oa}h_{bc}}{\Xi_a\Xi_r}.
\]
Consequently the \(oa\) response times the \(bc\) response is bounded
by \(Ch_{oa}h_{bc}/(\Xi_a\Xi_r)\) for every active \(r\).

For the next display abbreviate the two evaluated response ratios in
\eqref{eq:V58-normalized-edge-response} by
\(\mathfrak A_{oa}\) and \(\mathfrak A_{bc}\).  By
\Cref{lem:V58-three-block-response}, and because the same \(bc\)
response is no larger than each of its valid endpoint bounds,
\begin{align*}
 \mathcal R_{\rm dp}
 &\le C\mathfrak A_{oa}\mathfrak A_{bc}
       (\widehat H_{ab}+\widehat H_{ac})\\
 &\le C\sum_{r=b,c}
 \widehat H_{ar}\frac{h_{oa}h_{bc}}{\Xi_a\Xi_r}.
\end{align*}
This is \eqref{eq:V58-selected-RACC-total}.  It reaches the summed
quartic path target directly; it does not prove the stronger
pre-trace-norm splitting in \Cref{def:V57-RACC}.  No separate operator
allocation of the cut quotient was used.
\end{proof}

\begin{corollary}[Exact residual RACC geometry]
\label{cor:V58-RACC-residual}
After \Cref{thm:V58-selected-RACC}, an unclosed collision has either
\(v=a\), or an active \(v\in\{b,c\}\), such that
\begin{equation}
 s_\alpha\Xi_v>M,
 \qquad
 x_v<\varepsilon\Xi_v.
 \label{eq:V58-sideways-cut-vertex}
\end{equation}
The missing fraction of \(\Xi_v\) lies on punctured-core coordinates
or on other retained multiplier banks; it does not lie on the
collision edge incident to \(v\).
\end{corollary}

\begin{proof}
The negation of edge-paid is exactly
\eqref{eq:V58-sideways-cut-vertex}.  The restriction to active
\(b,c\) endpoints is the hypothesis of
\Cref{thm:V58-selected-RACC}.
\end{proof}

\section{The next original-row cut acquisition}
\label{sec:V58-next-gate}

For \(v\in\{a,b,c\}\), the next task is to prove, monomial by
monomial, the candidate decomposition of the mass outside its incident
collision edge
\begin{equation}
 \Xi_v-x_v
 =
 Z_v^{\rm cut}+P_v^{\rm mean}+D_v^{\rm int},
 \label{eq:V58-residual-mass-split}
\end{equation}
where \(Z_v^{\rm cut}\) should be carried by coordinates used in the
cut mark, \(P_v^{\rm mean}\) by scalar means retained in the exact
monomial, and \(D_v^{\rm int}\) by genuine internal defects.  V58 does
not assert this identity before the required disjoint coordinate
partition is constructed.

\begin{openproblem}[Sideways RACC gate]
\label{op:V58-sideways-RACC}
Prove the full aggregate RACC bound under
\eqref{eq:V58-sideways-cut-vertex} by keeping
\eqref{eq:V58-residual-mass-split} inside the aggregate response.
The permitted payments are:
\begin{enumerate}[label=\textup{(\roman*)},leftmargin=3.7em]
\item absorb \(Z_v^{\rm cut}\) with the original-row cut numerator;
\item retain a first or second scalar debit on
      \(P_v^{\rm mean}\); or
\item apply \eqref{eq:V53-one-edge-tax} to the actual raw internal
      defect carrying \(D_v^{\rm int}\).
\end{enumerate}
No composite \(bc\)-carrier denominator may replace \(\Xi_b\), and no
formation tax may be applied after the raw collision projector has
been discarded.
\end{openproblem}

\begin{assessment}[Triple-cold regression after the corrected product tax]
\label{ass:V58-triple-cold}
Let both collision pairs resolve to singlets and choose a low
punctured-core cut output.  The two pair costs are then controlled by
the row-factorized product tax \eqref{eq:V58-two-pair-tax}, not by
multiplying dimension ratios or by invoking an arbitrary-input map
norm.  If all active endpoints are edge-paid,
\Cref{thm:V58-selected-RACC} closes the total example.  If one required
edge is light, the example falls exactly into
\eqref{eq:V58-sideways-cut-vertex}; its large original-row mass remains
visible in \eqref{eq:V58-residual-mass-split}.  Thus the regression is
neither falsely declared harmless nor allowed to revive a detached
formation law.
\end{assessment}

\section{V58 ledger and next acquisition}
\label{sec:V58-ledger}

\begin{theorem}[Integrated V58 statement]
\label{thm:V58-integrated}
Preserving V1--V57, V58 proves:
\begin{enumerate}[label=\textup{(V58.\arabic*)},leftmargin=5.5em]
\item multiplicative composition of the V52/V53 taxes on the exact
      row-factorized four-input tensor;
\item the two-pair amplified formation tax
      \eqref{eq:V58-two-pair-tax};
\item the weighted three-block response
      \eqref{eq:V58-three-block-response};
\item the summed path bound whenever \(a\) and every active endpoint
      among \(b,c\) is thermal or selected-heavy on its collision
      edge; and
\item reduction of the remaining RACC work to the exact sideways mass
      split \eqref{eq:V58-residual-mass-split}.
\end{enumerate}
\end{theorem}

\begin{proof}
Items \textup{(V58.1)--(V58.2)} are
\Cref{lem:V58-tax-composition,thm:V58-two-pair-tax}.
Item \textup{(V58.3)} is
\Cref{lem:V58-three-block-response}.  Item \textup{(V58.4)} is
\Cref{thm:V58-selected-RACC}.  Item \textup{(V58.5)} is
\Cref{cor:V58-RACC-residual,op:V58-sideways-RACC}.
\end{proof}

\begingroup
\small
\begin{longtable}{>{\raggedright\arraybackslash}p{0.29\textwidth}
                  >{\raggedright\arraybackslash}p{0.15\textwidth}
                  >{\raggedright\arraybackslash}p{0.48\textwidth}}
\toprule
\textbf{V58 row}&\textbf{Status}&\textbf{Advance or obstruction}\\
\midrule
\endfirsthead
\toprule
\textbf{V58 row}&\textbf{Status}&\textbf{Advance or obstruction}\\
\midrule
\endhead
Two-pair aggregate trace tax
& \MGclosed
& The two bilinear taxes multiply on the exact row-factorized input;
  arbitrary cross-row entanglement is outside the claim.\\[0.35em]
Weighted double-pair response
& \MGclosed
& Low, pair-output, and core-output branches are controlled with one
  resolvent and the aggregate cut quotient.\\[0.35em]
Thermal/selected RACC region
& \MGclosed
& The total response closes when \(a\) and every active endpoint among
  \(b,c\) is paid thermally or on the incident edge.\\[0.35em]
Sideways RACC region
& \MGopen
& A hot collision-light vertex must be paid from its correlated cut,
  retained-mean, or internal-defect mass.\\[0.35em]
Complete collision and \([3,2]\)
& \MGopen
& Mean collisions and the selected-edge double-pair region are closed;
  sideways RACC and multi-collision aggregation remain.\\[0.35em]
\([3,3]\), \([4]\), and global quartic MTA
& \MGopen
& These remain downstream of complete \([3,2]\).\\[0.35em]
Full Yang--Mills mass gap
& \MGopen
& All-order MTA, WUFC/CPM, nonlinear construction, continuum OS
  existence/nontriviality, and the spectral passage remain unproved.\\
\bottomrule
\end{longtable}
\endgroup

\begin{openproblem}[V59 mandate: close sideways RACC]
\label{op:V59}
\begin{enumerate}[label=\textup{(V59.\arabic*)},leftmargin=5.5em]
\item Work with the aggregate cut influence and the single aggregate
      projector
      \eqref{eq:V57-aggregate-collision-pinching}.
\item Prove the exact residual split
      \eqref{eq:V58-residual-mass-split} monomial by monomial, including
      collision-complement means and internal V47 defects.
\item In the cut-heavy branch, use
      \(H_{ar}\ge Z_v^{\rm cut}\) at the relevant endpoint; do not infer
      input mass from the cut output.
\item In the mean-heavy branch, retain one unused scalar first or
      second debit after both pair outputs have been paid.
\item In the internal-defect-heavy branch, compose a third
      row-factorized V53 tax before nested raw pinching.
\item Close the cases with one sideways vertex first, then the case
      where both internal vertices are sideways on different factors.
\item Only after single-collision RACC closes, aggregate several
      collision coordinates by fixed disjoint-pair pattern.  Keep
      \([3,3]\) and \([4]\) suspended.
\item Preserve only the newest YM file.  The next companion audit is
      V60.
\end{enumerate}
\end{openproblem}

\begin{firewall}[End-of-V58 claim boundary]
\label{fire:V58-claim-boundary}
V58 proves the product aggregate pair tax for the exact row-factorized
coefficient input and closes the all-relevant-endpoint
thermal/selected region of the summed path response.  It proves
neither a cross-row-entangled map norm nor separate \(ab\) and \(ac\)
cut quotients.  It does not prove sideways RACC, multi-collision
aggregation, complete \([3,2]\), \([3,3]\),
\([4]\), global \(m=4\), all-order MTA, WUFC, CPM, the nonlinear
Perron packet, finite-edge margins, constructive capacity,
protected-sector floors, local-algebra noncollapse, reflection
positivity, a nontrivial continuum OS state, or a positive
Yang--Mills spectral gap.  The full Yang--Mills mass gap is not proved.
\end{firewall}

\section*{Version V58 append-only record}
\addcontentsline{toc}{section}{Version V58 append-only record}
The complete V57 source was retained before this continuation.  Its
SHA--256 digest was
\begin{center}\scriptsize\ttfamily
26672e9ed46ac8984008e5dd06f78c3fa483a78a2a6487f030bb21e890be78fa.
\end{center}
No mathematical content of V57 was altered.  On the next iteration,
retain this full file, append before its unique active document-closing
command, increment the filename to \texttt{\_V59.tex}, and remove only
the superseded V58 file from this Yang--Mills note series.

% =============================================================================
% END APPENDED VERSION V58 MATERIAL
% =============================================================================

% =============================================================================
% BEGIN APPENDED VERSION V59 MATERIAL
% =============================================================================

\chapter{V59: Saturated cut promotion and the exact sideways
incidence split}
\label{ch:V59}

\section{Purpose and corrected starting point}
\label{sec:V59-purpose}

V58 established two facts that will be used here in their corrected
form.  First, the two disjoint collision-pair taxes multiply on the
exact row-factorized input tensor before final recoupling; no
arbitrary cross-row-entangled map norm is available.  Second, the
ternary cut quotient is controlled by
\(H_\Sigma=\widehat H_{ab}+\widehat H_{ac}\), not by two separately
proved quotient operators.  The present chapter goes upstream of the
remaining sideways branch.  It keeps the universal cut cap together
with the influence cap and obtains a saturated response.  That
saturation pays a collision-light endpoint whenever a fixed fraction
of its mass is visible on one original-row cut.  A retained private
mean supplies a second, independent endpoint payment.

\begin{firewall}[Scope of V59]
\label{fire:V59-scope}
V59 proves an aggregate summed-path estimate.  It does not prove the
stronger pre-trace-norm splitting in \Cref{def:V57-RACC}.  The cut and
mean promotions below require only one endpoint response on the other
side of the relevant path, but they do not close two simultaneously
sideways endpoints when neither has a retained scalar debit.  Mass on
other puncture blocks or off-cut core blocks also remains to be
assembled.
\end{firewall}

\section{The saturated three-block response}
\label{sec:V59-saturated-response}

Keep the notation and row-normalized pair responses of V58, and put
\begin{equation}
 \mathfrak m_\Sigma
 :=\min\left\{H_\Sigma,\frac1{s_\alpha}\right\}.
 \label{eq:V59-saturated-cut-size}
\end{equation}

\begin{theorem}[Saturated aggregate cut response]
\label{thm:V59-saturated-response}
For every row-factorized four-input coefficient tensor in the
double-pair collision branch,
\begin{equation}
 \boxed{
 \mathcal R_{\rm dp}
 \le C\mathfrak m_\Sigma
 \mathfrak A_{oa}[B_o,B_a]
 \mathfrak A_{bc}[B_b,B_c].}
 \label{eq:V59-saturated-response}
\end{equation}
The estimate is uniform in the punctured support, representations,
multiplicities, and all row spectators.
\end{theorem}

\begin{proof}
Repeat the output-mass split in
\Cref{lem:V58-three-block-response}, but retain both available cut
bounds.  On low final outputs, the resolvent contributes
\(C/s_\alpha\), while \eqref{eq:V57-cut-cap} contributes
\[
 C\min\{s_\alpha H_\Sigma,1\}.
\]
Their product is \(C\mathfrak m_\Sigma\).

Suppose next that a pair output carries a fixed fraction of the high
final mass.  Use the output-weighted term in the corresponding
\(\mathfrak A\), the modulus term in the other pair response, and the
same cut cap.  The conversion of the output weight to
\(s_\alpha\Xi^{\rm out}\) again leaves exactly
\(C\mathfrak m_\Sigma\).

If the punctured-core output carries the high mass, there are two
bounds.  The quotient \eqref{eq:V57-cut-quotient} gives
\(CH_\Sigma\).  On the same high branch the core heat multiplier has
a fixed denominator gap.  Hence
\eqref{eq:V57-cut-first-moment} gives instead \(C/s_\alpha\).
Taking the smaller gives \(C\mathfrak m_\Sigma\).  An ordinary
spectator output is handled by its first moment and the low-output cut
cap, with the same result.  Finally the two pair compressions multiply
by \Cref{lem:V58-tax-composition} before final recoupling and aggregate
pinching.  Dividing by the four row trace norms proves
\eqref{eq:V59-saturated-response}.
\end{proof}

\begin{remark}[Why the saturation is new]
\label{rem:V59-saturation-new}
V57 used the influence side of the cut estimate in the high-core
branch and V58 inherited that choice.  The extra
\(s_\alpha^{-1}\) alternative comes from using the already proved cut
first moment on the same denominator-gap branch.  No new
representation-formation probability is introduced.
\end{remark}

\section{Endpoint response criteria}
\label{sec:V59-endpoint-response}

For an endpoint \(v\) of a collision pair \(e_v\), say that an
\emph{endpoint response} is available if
\begin{equation}
 \mathfrak A_{e_v}\le C\frac{h_{e_v}}{\Xi_v}.
 \label{eq:V59-endpoint-response}
\end{equation}
The evaluated row coefficients in \(\mathfrak A_{e_v}\) are left
implicit.

\begin{lemma}[Thermal and selected-heavy endpoints satisfy the
criterion]
\label{lem:V59-old-endpoint-response}
Every edge-paid endpoint in the sense of V58 satisfies
\eqref{eq:V59-endpoint-response}.
\end{lemma}

\begin{proof}
If \(s_\alpha\Xi_v\le M\), the raw cap gives
\[
 \mathfrak A_{e_v}\le Cs_\alpha h_{e_v}
 \le CM\frac{h_{e_v}}{\Xi_v}.
\]
If the endpoint is nonthermal and selected-heavy,
\eqref{eq:V58-hot-edge-tax} gives
\[
 \mathfrak A_{e_v}
 \le C\frac{h_{e_v}}{s_\alpha\Xi_v^2}
 \le \frac C M\frac{h_{e_v}}{\Xi_v}.
\]
\end{proof}

Let \(P_v\) be an exact coordinate bank on which row \(v\) is a
singleton in the fixed puncture monomial, and let \(p_v\) be its
input mass.  Denote its product heat multiplier by \(b_v\).  The bank
is disjoint from the incident collision pair.

\begin{lemma}[Private-mean endpoint response]
\label{lem:V59-private-mean-response}
Assume
\begin{equation}
 p_v\ge\eta\Xi_v,
 \qquad s_\alpha\Xi_v>M,
 \label{eq:V59-private-mean-heavy}
\end{equation}
and retain \(b_v\) together with the incident collision defect until
after their output masses have been allocated.  Then their combined
normalized response satisfies
\begin{equation}
 \boxed{
 \mathfrak A_{e_v}^{\rm mean}
 \le C_{\eta,M}\frac{h_{e_v}}{\Xi_v}.}
 \label{eq:V59-private-mean-response}
\end{equation}
The assertion remains valid with arbitrary row spectators.
\end{lemma}

\begin{proof}
Because the bank is singleton and disjoint from the collision
coordinate, its multiplier factors before recoupling.  Apply
\Cref{lem:V55-unused-debit} to the exact union of its coordinates.
The proof of that lemma uses the product second moment
\eqref{eq:V32-balanced-second-moment} and is unchanged by row
spectators.  With \(\delta_{e_v}\) the incident pair defect, it gives
\[
 |b_v\delta_{e_v}|
 +s_\alpha\Xi_{P_v,e_v}^{\rm out}
  |b_v\delta_{e_v}|
 \le C\frac{h_{e_v}}{p_v}.
\]
Use \(p_v\ge\eta\Xi_v\).  Aggregate raw pinching is contractive, so
the displayed scalar auxiliary estimate proves
\eqref{eq:V59-private-mean-response} for the normalized bilinear
coefficient response.
\end{proof}

\begin{corollary}[Endpoint-response closure of the summed paths]
\label{cor:V59-endpoint-closure}
If \eqref{eq:V59-endpoint-response} holds at \(a\) and at every
\(r\in\{b,c\}\) for which \(\widehat H_{ar}>0\), then
\begin{equation}
 \mathcal R_{\rm dp}
 \le C\sum_{r=b,c}
 \frac{h_{oa}\widehat H_{ar}h_{bc}}{\Xi_a\Xi_r}.
 \label{eq:V59-endpoint-closure}
\end{equation}
Thus a retained private mean may replace the thermal/selected-heavy
condition at any one or several endpoints.
\end{corollary}

\begin{proof}
Use the unsaturated response
\eqref{eq:V58-three-block-response} and expand the scalar identity
\(H_\Sigma=\widehat H_{ab}+\widehat H_{ac}\).  For the summand indexed
by \(r\), apply \eqref{eq:V59-endpoint-response} at \(a\) and at
\(r\).  The same \(bc\) response is bounded by every applicable
endpoint estimate.  Summing the two nonnegative bounds proves the
claim.  This is an after-trace-norm summed-path estimate, not a claim
of separate cut operators.
\end{proof}

\section{Cut-heavy promotion}
\label{sec:V59-cut-heavy}

\begin{theorem}[One sideways endpoint is paid by a heavy original-row
cut]
\label{thm:V59-cut-heavy-promotion}
Fix \(r\in\{b,c\}\) and \(v\in\{a,r\}\); let \(w\) be the other
endpoint of the path edge \(ar\).  Suppose
\begin{equation}
 \widehat H_{ar}\ge\eta\Xi_v
 \label{eq:V59-cut-heavy}
\end{equation}
and an endpoint response \eqref{eq:V59-endpoint-response} is available
at \(w\).  At \(v\) assume only the raw collision-pair cap.  Then the
\emph{entire aggregate} double-pair response satisfies
\begin{equation}
 \boxed{
 \mathcal R_{\rm dp}
 \le C_{\eta}\frac{h_{oa}\widehat H_{ar}h_{bc}}
                         {\Xi_a\Xi_r}.}
 \label{eq:V59-cut-heavy-promotion}
\end{equation}
In particular, all other cut-influence terms are absorbed by this one
heavy path.
\end{theorem}

\begin{proof}
At the unnormalized endpoint \(v\), use
\(\mathfrak A_{e_v}\le Cs_\alpha h_{e_v}\).  At \(w\), use
\eqref{eq:V59-endpoint-response}.  Since
\(\mathfrak m_\Sigma\le s_\alpha^{-1}\),
\Cref{thm:V59-saturated-response} gives
\[
 \mathcal R_{\rm dp}
 \le C\frac{h_{oa}h_{bc}}{\Xi_w}.
\]
By \eqref{eq:V59-cut-heavy},
\[
 \frac1{\Xi_w}
 \le\frac1\eta
 \frac{\widehat H_{ar}}{\Xi_v\Xi_w}.
\]
Because \(\{v,w\}=\{a,r\}\), this is
\eqref{eq:V59-cut-heavy-promotion}.
\end{proof}

\section{A canonical incidence partition of the missing mass}
\label{sec:V59-incidence-partition}

Set \(A_{i,f}=0\) when row \(i\) is inactive at the punctured-core
coordinate \(f\).  On active coordinates there is a fixed group
constant
\begin{equation}
 A_{i,f}\ge a_*>0;
 \label{eq:V59-active-mass-floor}
\end{equation}
with the normalization \(A=1+C_2\) one may take \(a_*=1\).
For every core coordinate with \(A_{a,f}>0\) and at least one of
\(A_{b,f},A_{c,f}\) positive, choose deterministically
\(\rho(f)=b\) if \(A_{b,f}>0\), and \(\rho(f)=c\) otherwise.  Put
\begin{align}
 E_{a\to r}&:=\{f:\rho(f)=r\},
 &Z_{a\to r}&:=\sum_{f\in E_{a\to r}}A_{a,f},
 \label{eq:V59-a-cut-mass}\\
 E_{r\leftarrow a}&:=\{f:A_{r,f}>0,\ A_{a,f}>0\},
 &Z_{r\leftarrow a}&:=\sum_{f\in E_{r\leftarrow a}}A_{r,f},
 \label{eq:V59-r-cut-mass}\\
 P_a^{\rm priv}&:=
 \sum_{\substack{f:A_{a,f}>0\\A_{b,f}=A_{c,f}=0}}A_{a,f}.
 \label{eq:V59-inward-cut-mass}
\end{align}
For \(r=b,c\), let \(P_r^{\rm priv}\) be the mass of core
coordinates private to \(r\), and let \(D_r^{\rm off}\) be the mass
of all remaining core coordinates active at \(r\) but not in
\(E_{r\leftarrow a}\).  Finally, let \(Y_v\) be the input-\(v\) mass on
punctured attachment or spectator coordinates outside the selected
collision pair and outside the punctured ternary core.

\begin{lemma}[Exact incidence split and cut domination]
\label{lem:V59-incidence-split}
For the selected collision coordinate, let \(x_v\) be the mass of row
\(v\) on its incident collision pair.  Then the following are exact
disjoint coordinate identities:
\begin{align}
 \Xi_a
 &=x_a+Z_{a\to b}+Z_{a\to c}+P_a^{\rm priv}+Y_a,
 \label{eq:V59-a-mass-split}\\
 \Xi_r
 &=x_r+Z_{r\leftarrow a}+P_r^{\rm priv}
   +D_r^{\rm off}+Y_r,
 \qquad r=b,c.
 \label{eq:V59-r-mass-split}
\end{align}
Moreover,
\begin{equation}
 \boxed{
 \widehat H_{ar}\ge
 a_*\max\{Z_{a\to r},Z_{r\leftarrow a}\},
 \qquad r=b,c.}
 \label{eq:V59-cut-dominates-incidence-mass}
\end{equation}
No coordinate is classified according to a term created after a
triangle inequality.
\end{lemma}

\begin{proof}
The selected collision pair, punctured ternary core, and remaining
puncture/spectator banks are disjoint by the exact puncture formula
\eqref{eq:V54-exact-puncture}.  Within the core, the selector \(\rho\)
partitions every coordinate active at \(a\) into
\(E_{a\to b}\), \(E_{a\to c}\), or the private set.  This proves
\eqref{eq:V59-a-mass-split}.  For a fixed \(r\), its full overlap set
\(E_{r\leftarrow a}\), private set, and complementary active core set are disjoint and
exhaustive, proving \eqref{eq:V59-r-mass-split}.

For every \(f\in E_{a\to r}\), both rows are active, and so are both
rows on the larger set \(E_{r\leftarrow a}\).  Hence
\[
 A_{a,f}A_{r,f}\ge a_*A_{a,f},
 \qquad
 A_{a,f}A_{r,f}\ge a_*A_{r,f}.
\]
Sum the first inequality over \(E_{a\to r}\), sum the second over
\(E_{r\leftarrow a}\), and then enlarge to all coordinates in
\(\widehat H_{ar}\).  This proves
\eqref{eq:V59-cut-dominates-incidence-mass}.
\end{proof}

\begin{remark}[Correction of the V58 candidate split]
\label{rem:V59-corrected-split}
The schematic identity \eqref{eq:V58-residual-mass-split} was not a
proved or canonical partition: a triple-active coordinate can
contribute to both original-row influences, and ``internal defect''
depends on a later polynomial expansion.  Equations
\eqref{eq:V59-a-mass-split}--\eqref{eq:V59-r-mass-split} replace it by
a deterministic incidence partition made before trace norm.  The
quantity \(D_r^{\rm off}\) is deliberately retained rather than
silently assigned to the \(ar\) cut.
\end{remark}

\begin{theorem}[Closure of every one-sideways incidence branch]
\label{thm:V59-one-sideways-reduction}
Fix a sufficiently small \(\varepsilon>0\).  Suppose exactly one
internal endpoint \(v\) among \(a\) and the active endpoints in
\(\{b,c\}\) fails the V58 edge-paid condition, while all other active
endpoints have an endpoint response
\eqref{eq:V59-endpoint-response}.  Then:
\begin{enumerate}[label=\textup{(\roman*)},leftmargin=3.7em]
\item if \(v=a\) and
      \(Y_a<\varepsilon\Xi_a\), the summed path bound closes;
\item if \(v=r\in\{b,c\}\) and
      \(D_r^{\rm off}+Y_r<\varepsilon\Xi_r\), the summed path bound
      closes.
\end{enumerate}
Consequently an unclosed branch with exactly one sideways endpoint
must satisfy respectively
\begin{equation}
 Y_a\ge\varepsilon\Xi_a,
 \qquad\hbox{or}\qquad
 D_r^{\rm off}+Y_r\ge\varepsilon\Xi_r.
 \label{eq:V59-one-sideways-residual}
\end{equation}
\end{theorem}

\begin{proof}
Because \(v\) is not edge-paid,
\[
 s_\alpha\Xi_v>M,
 \qquad x_v<\varepsilon\Xi_v.
\]
Consider first \(v=a\).  From \eqref{eq:V59-a-mass-split} and the
hypothesis on \(Y_a\), either
\(P_a^{\rm priv}\ge\varepsilon\Xi_a\), after decreasing the fixed
\(\varepsilon\) by a numerical factor, or
\[
 Z_{a\to b}+Z_{a\to c}\ge(1-3\varepsilon)\Xi_a.
\]
In the first case \Cref{lem:V59-private-mean-response} supplies the
missing endpoint response at \(a\), and
\Cref{cor:V59-endpoint-closure} applies.  In the second case one of
the two assigned masses is at least
\((1-3\varepsilon)\Xi_a/2\).  For its index \(r\),
\eqref{eq:V59-cut-dominates-incidence-mass} gives
\(\widehat H_{ar}\ge c\Xi_a\).  The endpoint \(r\) has a response by
hypothesis, so \Cref{thm:V59-cut-heavy-promotion} closes the entire
aggregate response.

Now let \(v=r\).  Equation \eqref{eq:V59-r-mass-split} gives either
\(P_r^{\rm priv}\ge\varepsilon\Xi_r\) or
\(Z_{r\leftarrow a}\ge(1-3\varepsilon)\Xi_r\), again after a fixed
numerical reduction of \(\varepsilon\).  The private case is closed
by \Cref{lem:V59-private-mean-response,cor:V59-endpoint-closure}; the
assigned-cut case is closed by
\eqref{eq:V59-cut-dominates-incidence-mass} and
\Cref{thm:V59-cut-heavy-promotion}, because \(a\) has an endpoint
response.  The contrapositives are
\eqref{eq:V59-one-sideways-residual}.
\end{proof}

\section{V59 ledger and the next upstream gate}
\label{sec:V59-ledger}

\begin{theorem}[Integrated V59 statement]
\label{thm:V59-integrated}
Preserving V1--V58, V59 proves:
\begin{enumerate}[label=\textup{(V59.\arabic*)},leftmargin=5.5em]
\item the saturated aggregate response
      \eqref{eq:V59-saturated-response};
\item the common endpoint-response criterion
      \eqref{eq:V59-endpoint-response};
\item a private-mean endpoint response and its summed-path closure;
\item the cut-heavy promotion
      \eqref{eq:V59-cut-heavy-promotion};
\item the exact incidence partitions
      \eqref{eq:V59-a-mass-split}--\eqref{eq:V59-r-mass-split}; and
\item reduction of every exactly-one-sideways branch to the explicit
      off-cut mass in \eqref{eq:V59-one-sideways-residual}.
\end{enumerate}
\end{theorem}

\begin{proof}
The six items are respectively
\Cref{thm:V59-saturated-response,
lem:V59-old-endpoint-response,
lem:V59-private-mean-response,cor:V59-endpoint-closure,
thm:V59-cut-heavy-promotion,lem:V59-incidence-split,
thm:V59-one-sideways-reduction}.
\end{proof}

\begingroup
\small
\begin{longtable}{>{\raggedright\arraybackslash}p{0.29\textwidth}
                  >{\raggedright\arraybackslash}p{0.15\textwidth}
                  >{\raggedright\arraybackslash}p{0.48\textwidth}}
\toprule
\textbf{V59 row}&\textbf{Status}&\textbf{Advance or obstruction}\\
\midrule
\endfirsthead
\toprule
\textbf{V59 row}&\textbf{Status}&\textbf{Advance or obstruction}\\
\midrule
\endhead
Saturated cut response
& \MGclosed
& The aggregate cut pays by the better of its influence quotient and
  its output first moment.\\[0.35em]
Private-mean endpoint
& \MGclosed
& An exact singleton product bank remains assembled with the incident
  collision defect and restores one original input denominator.\\[0.35em]
Cut-heavy endpoint
& \MGclosed
& A fixed original-row cut fraction pays one collision-light endpoint
  against any endpoint-paid partner.\\[0.35em]
Exactly one sideways endpoint
& \MGreduced
& Only off-cut core mass or other puncture/spectator mass remains.\\[0.35em]
Two or more sideways endpoints
& \MGopen
& The disjoint-pair product tax cannot be reused as an overlapping
  cut/pair tax on already cross-row-entangled coefficients.\\[0.35em]
Complete collision and \([3,2]\)
& \MGopen
& Off-cut puncture assembly, simultaneous-sideways response, and
  multi-collision aggregation remain.\\[0.35em]
Full Yang--Mills mass gap
& \MGopen
& Global quartic/all-order MTA, WUFC/CPM, nonlinear construction,
  continuum OS existence/nontriviality, and spectral positivity remain
  downstream.\\
\bottomrule
\end{longtable}
\endgroup

\begin{openproblem}[V60 mandate: audited off-cut and simultaneous
sideways closure]
\label{op:V60}
\begin{enumerate}[label=\textup{(V60.\arabic*)},leftmargin=5.5em]
\item Before new YM work, perform the scheduled comparison with the
      then-current SM and NS notes in \texttt{latex\_notes}; transfer
      only proved assembly devices and record exact filenames and
      digests.
\item Trace \(Y_a,Y_r\), and \(D_r^{\rm off}\) through the exact
      puncture formula \eqref{eq:V54-exact-puncture}, separating
      complement means, noncollision attachment defects, and genuine
      off-cut ternary banks before trace norm.
\item Apply \Cref{lem:V59-private-mean-response} to every exact scalar
      subbank, and V53 only to a raw pair defect whose input endpoint
      carries the required fraction.
\item For two simultaneous sideways endpoints, formulate one joint
      multilinear trace functional before any of its overlapping
      cut/pair resolutions.  Do not multiply two V52/V53 bilinear
      taxes that share an input row.
\item Test the alternating-height regression in which the two endpoint
      masses are large on different common coordinates and the cut
      output is cold.  Mere influence-to-mass domination is not a
      substitute for the joint coefficient tax in this corner.
\item Aggregate several collision coordinates only after the
      single-collision simultaneous-sideways gate closes.  Keep
      \([3,3]\) and \([4]\) suspended.
\item Preserve only the newest YM source after the V60 append.
\end{enumerate}
\end{openproblem}

\begin{firewall}[End-of-V59 claim boundary]
\label{fire:V59-claim-boundary}
V59 proves the saturated cut response, private-mean and cut-heavy
endpoint promotions, a canonical pre-trace incidence partition, and
the one-sideways reduction.  It does not prove the simultaneous-
sideways joint coefficient tax, off-cut puncture assembly,
multi-collision aggregation, complete \([3,2]\), \([3,3]\), \([4]\),
global \(m=4\), all-order MTA, WUFC, CPM, the nonlinear Perron packet,
finite-edge margins, constructive capacity, protected-sector floors,
local-algebra noncollapse, reflection positivity, a nontrivial
continuum OS state, or a positive Yang--Mills spectral gap.  The full
Yang--Mills mass gap is not proved.
\end{firewall}

\section*{Version V59 append-only record}
\addcontentsline{toc}{section}{Version V59 append-only record}
The complete corrected V58 source was retained before this
continuation.  Its SHA--256 digest was
\begin{center}\scriptsize\ttfamily
8d6ceb764d4de56917e8ec60eddfbd2351e36f7253d539eaa7c086a42c8881cf.
\end{center}
No mathematical content of that corrected V58 source was altered
during the V59 append.  On the next iteration, retain this full file,
append before its unique active document-closing command, increment
the filename to \texttt{\_V60.tex}, and remove only the superseded V59
file from this Yang--Mills note series.

% =============================================================================
% END APPENDED VERSION V59 MATERIAL
% =============================================================================

% =============================================================================
% BEGIN APPENDED VERSION V60 MATERIAL
% =============================================================================

\chapter{V60: Companion audit, the puncture-residual alphabet, and
the joint raw capacity gate}
\label{ch:V60}

\section{Mandatory companion audit}
\label{sec:V60-audit}

The required five-version audit was performed before selecting the
V60 acquisition.  Because the SM file changed once during the read,
its final snapshot was hashed twice and was stable before being used.
The audited sources were:
\begin{center}
\small
\begin{tabular}{@{}p{0.49\textwidth}rp{0.31\textwidth}@{}}
\toprule
\textbf{source}&\textbf{lines}&\textbf{SHA--256}\\
\midrule
\texttt{Renewal\_SM\_Einstein\_Continuum\_Research\_Notes\_V73.tex}
&59472&\texttt{e3dc8eff45318ae338b3a05bf54bccbc39e009eb0f9ace51091e130025503d11}\\
\texttt{Renewal\_NS\_Stability\_Research\_Notes\_V31.tex}
&23052&\texttt{f1121b62238ad5491bb7cf03635ea9934942edf573ed8faaa9ee79e1d38a6696}\\
\bottomrule
\end{tabular}
\end{center}

The NS source is byte-for-byte the V31 snapshot audited at V55.  Its
terminal result still preserves the restricted correlated formation
occupation before positive annular bounds and warns that separate
unmarkings lose the critical record factor.  There is therefore no
new NS estimate to import.

The SM source advanced from the V70 snapshot audited at V55 to V73.
V71--V73 prove results about the invariant minimal parent, residual
boundary moment map and corner cocycle, and the torsion-free Palatini
ancestor with compatible-jet GHY reduction.  None is a Wilson
coefficient, LR-capacity, or marked-tree estimate.  Their relevant
assembly conclusion is instead that the enlarged connection/jet
parent must be retained until its compatibility condition is imposed;
eliminating it before boundary polarization gives the wrong object.

\begin{assessment}[V60 transfer decision]
\label{ass:V60-audit-transfer}
No analytic inequality is transferred from either companion note.
The common assembly discipline changes the YM acquisition: all
resolutions sharing an input row will be placed in one enlarged joint
raw functional before trace norm.  Separate V52/V53 bilinear taxes
will not be multiplied across an overlapping row.  This is a change
of proof architecture, not evidence for the required capacity bound.
\end{assessment}

\begin{firewall}[Audit separation]
\label{fire:V60-audit-separation}
The Palatini/GHY compatible-jet theorem is not a Yang--Mills boundary
or BRST theorem, and the NS formation occupation is not an SU(3)
fusion estimate.  Only their proved order-of-assembly restrictions
are used to select the gate below.
\end{firewall}

\section{Exact alphabet of the V59 puncture residual}
\label{sec:V60-residual-alphabet}

Fix one collision coordinate and then one fully resolved local
partition monomial in the remaining puncture factors of
\eqref{eq:V54-exact-puncture}.  At a remaining attachment coordinate
\(f\), let its selected pair be \(ot_f\).  The two complementary rows
are denoted \(u_f,v_f\), with inactive entries interpreted as
trivial.  The exact local identity is
\begin{equation}
 \mathbf d_{ot_f,f}\otimes
 \mathbf Q_{\{u_f,v_f\},f}
 =
 q_{u_f,f}q_{v_f,f}\mathbf d_{ot_f,f}
 +
 \mathbf d_{ot_f,f}\otimes\mathbf d_{u_fv_f,f},
 \label{eq:V60-local-puncture-alphabet}
\end{equation}
with the second term zero if fewer than two complementary rows are
active.  This is the V56 collision identity, now used only to classify
the previously unnamed residual \(Y_v\).

For a fixed monomial \(\mu\) and input row \(i\), let
\(P_{i,\mu}^{\rm pun}\) be its mass on residual puncture coordinates
where the row-
\(i\) entry is a singleton multiplier.  For each
\(j\ne i\), let \(D_{ij,\mu}^{\rm pun}\) be its mass on residual
puncture coordinates where \(i,j\) form a selected local pair block.

\begin{lemma}[Puncture-residual mass alphabet]
\label{lem:V60-puncture-alphabet}
For every fixed local partition monomial,
\begin{equation}
 \boxed{
 Y_i=P_{i,\mu}^{\rm pun}
 +\sum_{j\ne i}D_{ij,\mu}^{\rm pun}.}
 \label{eq:V60-puncture-mass-split}
\end{equation}
The four coordinate sets in this display are disjoint.  Consequently,
if \(Y_i\ge\eta\Xi_i\), then either
\begin{equation}
 P_{i,\mu}^{\rm pun}\ge\frac\eta2\Xi_i
 \label{eq:V60-puncture-mean-heavy}
\end{equation}
or, for at least one \(j\ne i\),
\begin{equation}
 D_{ij,\mu}^{\rm pun}\ge\frac\eta6\Xi_i.
 \label{eq:V60-puncture-pair-heavy}
\end{equation}
\end{lemma}

\begin{proof}
At a puncture coordinate, the entry of row \(i\) belongs to exactly
one block of the local replica partition.  Under the \([3,2]\)
priority, and after the selected \(ot_f\) block has been fixed, that
block is either the singleton \(\{i\}\) or one pair \(\{i,j\}\).
Equation \eqref{eq:V60-local-puncture-alphabet} lists the two
possibilities for a complementary row; if \(i=t_f\), it is in the
selected pair itself.  Thus the four alternatives are disjoint and
exhaust all coordinates counted by \(Y_i\), proving
\eqref{eq:V60-puncture-mass-split}.

If the singleton mass is smaller than \(\eta\Xi_i/2\), the three pair
masses have sum at least \(\eta\Xi_i/2\).  One is at least one third
of that sum, which is \eqref{eq:V60-puncture-pair-heavy}.
\end{proof}

\begin{corollary}[Monomialwise payment of the puncture mean]
\label{cor:V60-puncture-mean-closed}
In one fixed local partition monomial of an exactly-one-sideways
branch, if
\eqref{eq:V60-puncture-mean-heavy} holds and every other active
endpoint has an endpoint response, then that monomial has the required
endpoint debit and satisfies the summed path bound.  This statement
does not sum the local partition monomials.
\end{corollary}

\begin{proof}
The singleton factors over the indicated coordinates form an exact
product heat multiplier bank disjoint from the incident collision
pair.  Retain this bank and the incident pair defect through the
output split.  Since the sideways endpoint is nonthermal,
\Cref{lem:V59-private-mean-response} gives an endpoint response at
\(i\).  All remaining active endpoints have one by hypothesis, so
\Cref{cor:V59-endpoint-closure} applies.
\end{proof}

\begin{remark}[Why neither alternative may be summed monomialwise]
\label{rem:V60-no-monomial-tax}
Equations \eqref{eq:V60-puncture-mean-heavy} and
\eqref{eq:V60-puncture-pair-heavy} concern one local partition
monomial.  Although the first has the correct scalar debit by the
preceding corollary, summing it monomial by monomial could still
reintroduce a coordinate-subset count.  Likewise, the universal V53
tax applies to the complete same-label pair-bank defect obtained only
after all nonempty local pair choices have been regrouped.  Both
alternatives must therefore enter the joint fixed-type regrouping
before absolute values are taken.
\end{remark}

\section{The simultaneous-sideways regression}
\label{sec:V60-alternating-regression}

\begin{proposition}[Influence domination alone cannot pay two
sideways endpoints]
\label{prop:V60-alternating-regression}
For arbitrarily large \(L\), there are nonnegative active local masses
on two disjoint coordinate classes such that
\begin{equation}
 \Xi_a\asymp NL,qquad
 \Xi_r\asymp NL,qquad
 \widehat H_{ar}\asymp NL,
 \label{eq:V60-alternating-scales}
\end{equation}
while
\begin{equation}
 \frac{\widehat H_{ar}}{\Xi_a\Xi_r}asymp\frac1{NL}longrightarrow0.
 \label{eq:V60-alternating-target}
\end{equation}
Consequently the saturated scalar response, whose universal branch is
of order \(s_\alpha\) after two raw pair caps, cannot imply the
two-ended path normalization by mass algebra alone.
\end{proposition}

\begin{proof}
Take \(N/2\) coordinates with
\((A_{a,f},A_{r,f})=(L,1)\), and \(N/2\) with
\((A_{a,f},A_{r,f})=(1,L)\).  Then
\[
 \Xi_a=\Xi_r=\frac N2(L+1),
 \qquad
 \widehat H_{ar}=NL.
\]
This proves \eqref{eq:V60-alternating-scales} and
\eqref{eq:V60-alternating-target}.  On the saturated branch
\(\mathfrak m_\Sigma\le s_\alpha^{-1}\), two raw pair responses give
only
\[
 s_\alpha^{-1}(s_\alpha h_{oa})(s_\alpha h_{bc})
 =s_\alpha h_{oa}h_{bc}.
\]
The desired coefficient is instead
\(h_{oa}h_{bc}\widehat H_{ar}/(\Xi_a\Xi_r)\), whose ratio to the last
display tends to zero for fixed \(s_\alpha,N\) as \(L\to\infty\).
Thus the two missing endpoint denominators cannot be obtained from
\eqref{eq:V59-cut-dominates-incidence-mass} and scalar caps alone.
The construction is only a majorant regression; it does not assert
failure of the joint coefficient capacity defined next.
\end{proof}

\section{The joint raw endpoint capacity}
\label{sec:V60-JREC}

Fix a complete regrouped collision/puncture monomial before trace
norm.  Let \(\Lambda\) resolve simultaneously the \(oa\) and \(bc\)
collision outputs, the punctured-core output, every selected
same-label residual pair bank, and the final multiplicity channel.
Write
\[
 X_\Lambda
 :=P_\Lambda
 (B_o\otimes B_a\otimes B_b\otimes B_c)P_\Lambda,
\]
where the family is one aggregate orthogonal pinching.  Let
\(\mathfrak a_{oa,\Lambda}\) and
\(\mathfrak a_{bc,\Lambda}\) be the V58 weighted pair sizes, and put
\begin{equation}
 \mathfrak k_{\Lambda}
 :=\|K_{T,\Lambda}^{\circ}\|_{
m op}
 +s_\alpha\Xi_{T,\Lambda}^{\rm out}
  \|K_{T,\Lambda}^{\circ}\|_{
m op}.
 \label{eq:V60-weighted-cut-size}
\end{equation}
Any same-label puncture-pair factors are retained inside
\(P_\Lambda\) and their scalar multipliers; they are not pinched in a
separate estimate.

\begin{definition}[Joint raw endpoint capacity]
\label{def:V60-JREC}
The statement \(\operatorname{JREC}_{a\mid bc}\) is the estimate
\begin{equation}
 \boxed{
 \sum_\Lambda
 \mathfrak a_{oa,\Lambda}
 \mathfrak a_{bc,\Lambda}
 \mathfrak k_\Lambda\|X_\Lambda\|_1
 \le
 Cs_\alpha
 \sum_{r=b,c}
 \frac{h_{oa}\widehat H_{ar}h_{bc}}
      {\Xi_a\Xi_r}
 \prod_{i=o,a,b,c}\|B_i\|_1.}
 \label{eq:V60-JREC}
\end{equation}
It is required only on the simultaneous-sideways and puncture-residual
region left by V59 and the monomial analysis above.  Every resolution and same-label
regrouping occurs before the trace norms in the left side.
\end{definition}

\begin{theorem}[JREC is sufficient for the residual single-collision
sector]
\label{thm:V60-JREC-sufficient}
If \eqref{eq:V60-JREC} holds uniformly for every target choice,
punctured core, complete same-label regrouping, and row-factorized
coefficient tensor, then all simultaneous-sideways and
puncture-pair-heavy branches of one collision coordinate satisfy the
quartic summed marked-tree bound.
\end{theorem}

\begin{proof}
Use the single final resolvent only after the three displayed weighted
factors and all residual puncture factors have been assembled.  On a
low final output its norm is at most \(C/s_\alpha\).  On a high output,
allocate the output mass to the \(oa\) pair, \(bc\) pair, punctured
core, or an ordinary retained spectator.  The first three choices are
paid by the corresponding output-weighted summand in
\(\mathfrak a_{oa}\), \(\mathfrak a_{bc}\), or
\(\mathfrak k\); a spectator is paid by its established first moment.
In every case the remaining scalar factor is at most \(C/s_\alpha\)
times the left side of \eqref{eq:V60-JREC}.  Insert JREC.  The factor
\(s_\alpha\) cancels, leaving
\[
 C\sum_{r=b,c}
 \frac{h_{oa}\widehat H_{ar}h_{bc}}{\Xi_a\Xi_r},
\]
after division by the four row trace norms.  These are precisely the
two quartic path weights.  Aggregate pinching introduces no
multiplicity because all blocks belong to the single family
\(P_\Lambda\).
\end{proof}

\begin{assessment}[What V60 has and has not gained]
\label{ass:V60-JREC-status}
JREC is stronger than the disjoint-pair product tax and is not proved
by it.  It is nevertheless narrower than the former informal
sideways programme: every exactly-one-sideways cut-heavy branch is
removed, and scalar puncture means now carry a proved monomialwise
debit.  Their fixed-type aggregation remains part of the displayed
joint problem.  Its left side is one explicit positive trace
functional, with no
permission to multiply overlapping bilinear taxes or to expand over
coordinate subsets.
\end{assessment}

\section{V60 ledger and next acquisition}
\label{sec:V60-ledger}

\begin{theorem}[Integrated V60 statement]
\label{thm:V60-integrated}
Preserving V1--V59, V60 proves:
\begin{enumerate}[label=\textup{(V60.\arabic*)},leftmargin=5.5em]
\item the digest-recorded SM V73/NS V31 companion audit and its strict
      transfer firewall;
\item the exact puncture-residual alphabet
      \eqref{eq:V60-puncture-mass-split};
\item identification and monomialwise payment of the scalar-mean
      alternative in every \(Y_i\)-heavy branch;
\item the alternating-height regression
      \eqref{eq:V60-alternating-target}; and
\item sufficiency of the joint raw endpoint capacity
      \eqref{eq:V60-JREC} for the remaining single-collision sector.
\end{enumerate}
\end{theorem}

\begin{proof}
The items are respectively
\Cref{ass:V60-audit-transfer,fire:V60-audit-separation,
lem:V60-puncture-alphabet,cor:V60-puncture-mean-closed,
prop:V60-alternating-regression,thm:V60-JREC-sufficient}.
\end{proof}

\begingroup
\small
\begin{longtable}{>{\raggedright\arraybackslash}p{0.29\textwidth}
                  >{\raggedright\arraybackslash}p{0.15\textwidth}
                  >{\raggedright\arraybackslash}p{0.48\textwidth}}
\toprule
\textbf{V60 row}&\textbf{Status}&\textbf{Advance or obstruction}\\
\midrule
\endfirsthead
\toprule
\textbf{V60 row}&\textbf{Status}&\textbf{Advance or obstruction}\\
\midrule
\endhead
Companion audit
& \MGclosed
& SM V73 and unchanged NS V31 supply an order-of-assembly warning but
  no transferable YM estimate.\\[0.35em]
Puncture residual
& \MGreduced
& Scalar singleton mass has the correct monomialwise debit; its
  fixed-type aggregation and the actual local pair labels remain.\\[0.35em]
Two-sideways scalar route
& \MGopen
& Alternating local heights prove that influence domination and cut
  saturation alone cannot supply both endpoint denominators.\\[0.35em]
Joint raw capacity
& \MGopen
& JREC is sufficient and precisely assembled, but its SU(3) cold
  capacity estimate is not yet proved.\\[0.35em]
Complete collision and \([3,2]\)
& \MGopen
& JREC and the later multi-collision fixed-type regrouping remain.\\[0.35em]
Full Yang--Mills mass gap
& \MGopen
& Global quartic/all-order MTA, WUFC/CPM, nonlinear construction,
  continuum OS existence/nontriviality, and spectral positivity remain
  downstream.\\
\bottomrule
\end{longtable}
\endgroup

\begin{openproblem}[V61 mandate: prove or sharply reduce JREC]
\label{op:V61}
\begin{enumerate}[label=\textup{(V61.\arabic*)},leftmargin=5.5em]
\item Work with the one direct-sum raw family \(P_\Lambda\) in
      \eqref{eq:V60-JREC}.  Do not apply V53 after an overlapping
      resolution has been traced out.
\item Regroup all residual puncture coordinates by the six unordered
      row-pair labels before expanding any nonempty-choice polynomial.
      Seek a fixed graph-type polynomial analogous to V56 rather than
      a coordinate-subset sum.
\item Orient the joint cold capacity by the larger of the two
      simultaneous-sideways endpoint masses.  Use Frobenius
      reciprocity only inside the complete joint LR transform.
\item Derive a Chernoff transform for the joint low-output family and
      test whether the V53 universal one-sided transform gap survives
      after conditioning on the overlapping core output.
\item Test the joint transform on the alternating-height array in
      \Cref{prop:V60-alternating-regression} and on simultaneous dual
      singlets.  A bound that merely reproduces the scalar saturated
      response is insufficient.
\item If uniform joint contraction fails, identify the exact residual
      representation family and go upstream to a two-ended SU(3) LR
      transform theorem.  Do not hide it behind RACC notation.
\item Keep multi-collision aggregation, \([3,3]\), and \([4]\)
      suspended until JREC is resolved.  The next companion audit is
      V65.
\end{enumerate}
\end{openproblem}

\begin{firewall}[End-of-V60 claim boundary]
\label{fire:V60-claim-boundary}
V60 performs the mandatory audit, proves the exact puncture residual
alphabet and its monomialwise scalar-mean payment, proves that scalar influence
algebra cannot close the two-sideways corner, and isolates the
sufficient JREC functional.  It does not prove JREC, simultaneous-
sideways closure, multi-collision aggregation, complete \([3,2]\),
\([3,3]\), \([4]\), global \(m=4\), all-order MTA, WUFC, CPM, the
nonlinear Perron packet, finite-edge margins, constructive capacity,
protected-sector floors, local-algebra noncollapse, reflection
positivity, a nontrivial continuum OS state, or a positive Yang--Mills
spectral gap.  The full Yang--Mills mass gap is not proved.
\end{firewall}

\section*{Version V60 append-only record}
\addcontentsline{toc}{section}{Version V60 append-only record}
The complete corrected V59 source was retained before this
continuation.  Its SHA--256 digest was
\begin{center}\scriptsize\ttfamily
9e5bc315a9ceaaa7697d269788789ff058061f53da2296780dd627705f7bfa59.
\end{center}
No mathematical content of that V59 source was altered during the V60
append.  On the next iteration, retain this full file, append before
its unique active document-closing command, increment the filename to
\texttt{\_V61.tex}, and remove only the superseded V60 file from this
Yang--Mills note series.

% =============================================================================
% END APPENDED VERSION V60 MATERIAL
% =============================================================================

% =============================================================================
% BEGIN APPENDED VERSION V61 MATERIAL
% =============================================================================

\chapter{V61: Multilinear Schur compression and the two-row
conditional capacity}
\label{ch:V61}

\section{Purpose}
\label{sec:V61-purpose}

V60 isolated JREC as a weighted trace norm over one aggregate raw
family.  Its V61 mandate suggested a joint LR Chernoff transform, but
the first necessary step is more elementary and more exact: determine
which projector quantity controls arbitrary row-factorized
coefficients when several resolutions overlap in their input rows.
The answer is the operator norm of a partial trace of the
\emph{aggregate} output projector.  This gives a multilinear extension
of V52, a layer-cake capacity theorem for JREC, and a precise two-row
operator which must be estimated next.

\begin{firewall}[Scope of V61]
\label{fire:V61-scope}
The theorems below are exact finite-dimensional trace inequalities.
They do not estimate the resulting SU(3) two-row partial-trace blocks.
Thus V61 reduces JREC to a deterministic conditional capacity but does
not prove JREC or the remaining \([3,2]\) sector.
\end{firewall}

\section{A multi-input aggregate trace tax}
\label{sec:V61-multilinear-tax}

Let \(m\ge2\), let \(H_1,\dots,H_m\) be finite-dimensional carrier
spaces, and put \(H=\bigotimes_{i=1}^mH_i\).  Let
\(\{P_\omega\}_{\omega\in\Omega}\) be mutually orthogonal
projections on \(H\), and write
\begin{equation}
 P_\Omega:=\sum_{\omega\in\Omega}P_\omega.
 \label{eq:V61-aggregate-projector}
\end{equation}
For a nonempty subset \(I\subseteq[m]\), define
\begin{equation}
 T_{\Omega,I}:=
 \operatorname{Tr}_{I^c}P_\Omega,
 \qquad
 \tau_{\Omega,I}:=
 \min\{1,\|T_{\Omega,I}\|_{\rm op}\},
 \qquad
 \tau_\Omega^{\rm cond}:=min_{\varnothing\ne I\subseteq[m]}
 \tau_{\Omega,I}.
 \label{eq:V61-conditional-tax}
\end{equation}
The partial trace is unnormalized.

\begin{lemma}[Product-state conditional projector tax]
\label{lem:V61-product-density-tax}
For arbitrary positive trace-one operators \(\rho_i\) on \(H_i\),
\begin{equation}
 \boxed{
 \operatorname{Tr}\!\left[P_\Omega
 \bigotimes_{i=1}^m\rho_i\right]
 \le\tau_{\Omega,I}}
 \label{eq:V61-product-density-tax}
\end{equation}
for every nonempty \(I\subseteq[m]\), and hence with
\(\tau_{\Omega,I}\) replaced by \(\tau_\Omega^{\rm cond}\).
\end{lemma}

\begin{proof}
Put \(\rho_I=\bigotimes_{i\in I}\rho_i\) and similarly for
\(I^c\).  Since every density matrix has operator norm at most one,
\(0\le\rho_{I^c}\le I_{I^c}\).  Positivity gives
\begin{align*}
 \operatorname{Tr}[P_\Omega(\rho_I\otimes\rho_{I^c})]
 &\le\operatorname{Tr}[P_\Omega(\rho_I\otimes I_{I^c})]\\
 &=\operatorname{Tr}[T_{\Omega,I}\rho_I]
 \le\|T_{\Omega,I}\|_{\rm op}.
\end{align*}
The left side is also at most one because \(P_\Omega\le I\), proving
\eqref{eq:V61-product-density-tax}.
\end{proof}

For each row let \(\mathcal K_i\) be an arbitrary spectator space and
let \(A_i\in\mathcal S_1(H_i\otimes\mathcal K_i)\).  Choose an
isometry \(V_\omega:\operatorname{ran}P_\omega\to H\) with
\(V_\omega V_\omega^*=P_\omega\), and define the compressed output,
after the evident tensor reordering, by
\begin{equation}
 \mathcal C_\omega(A_1,\dots,A_m)
 :=(V_\omega^*\otimes I_{\rm sp})
 \left(\bigotimes_{i=1}^mA_i\right)
 (V_\omega\otimes I_{\rm sp}).
 \label{eq:V61-multilinear-compression}
\end{equation}

\begin{theorem}[Spectator-amplified multilinear aggregate Schur tax]
\label{thm:V61-multilinear-Schur-tax}
For arbitrary trace-class row coefficients,
\begin{equation}
 \boxed{
 \sum_{\omega\in\Omega}
 \|\mathcal C_\omega(A_1,\dots,A_m)\|_1
 \le
 \tau_\Omega^{\rm cond}
 \prod_{i=1}^m\|A_i\|_1.}
 \label{eq:V61-multilinear-Schur-tax}
\end{equation}
The estimate is uniform in all spectator spaces.
\end{theorem}

\begin{proof}
First let \(A_i=|x_i\rangle\langle y_i|\) be rank one, allowing
\(x_i,y_i\in H_i\otimes\mathcal K_i\).  Put
\[
 x_\omega=(V_\omega^*\otimes I_{\rm sp})
            \bigotimes_i x_i,
 \qquad
 y_\omega=(V_\omega^*\otimes I_{\rm sp})
            \bigotimes_i y_i.
\]
Then the \(\omega\)-compression is
\(|x_\omega\rangle\langle y_\omega|\), and Cauchy--Schwarz over
\(\omega\) gives
\begin{align*}
 \sum_\omega\|x_\omega\|\,\|y_\omega\|
 &\le
 \left\|(P_\Omega\otimes I_{\rm sp})\bigotimes_i x_i\right\|
 \left\|(P_\Omega\otimes I_{\rm sp})\bigotimes_i y_i\right\|.
\end{align*}
Let \(\rho_i^x\) be the carrier reduction of
\(|x_i\rangle\langle x_i|\).  Its trace is \(\|x_i\|^2\), and the
total carrier reduction of \(\bigotimes_i x_i\) is
\(\bigotimes_i\rho_i^x\).  Applying
\Cref{lem:V61-product-density-tax} after normalization gives
\[
 \left\|(P_\Omega\otimes I_{\rm sp})\bigotimes_i x_i\right\|^2
 \le\tau_\Omega^{\rm cond}\prod_i\|x_i\|^2.
\]
The same estimate holds for the \(y_i\).  Taking square roots in both
bounds proves the rank-one form of
\eqref{eq:V61-multilinear-Schur-tax}.  Expand each \(A_i\) in a
singular-value decomposition and sum.  The \(m\) singular-value sums
factor as \(\prod_i\|A_i\|_1\).
\end{proof}

\begin{corollary}[V52 is the singleton partial-trace case]
\label{cor:V61-V52-recovery}
Suppose compact-group irreducibles \(R_i\) act on \(H_i\), and
\(P_\Omega\) commutes with the diagonal group action.  If
\(D_\Omega=\operatorname{rank}P_\Omega\), then
\begin{equation}
 T_{\Omega,\{i\}}
 =\frac{D_\Omega}{d_{R_i}}I_{H_i},
 \qquad
 \tau_{\Omega,\{i\}}
 =\min\left\{1,\frac{D_\Omega}{d_{R_i}}\right\}.
 \label{eq:V61-singleton-Schur-tax}
\end{equation}
For \(m=2\), minimizing over the two singleton subsets and the full
set recovers \eqref{eq:V52-aggregate-tax}.
\end{corollary}

\begin{proof}
The singleton partial trace commutes with the irreducible action
\(R_i\).  Schur's lemma makes it scalar, and its trace is
\(D_\Omega\), proving \eqref{eq:V61-singleton-Schur-tax}.  When
\(m=2\), the full-set partial trace is \(P_\Omega\) and has operator
norm one, so the stated minimum is exactly the V52 tax.
\end{proof}

\section{Weighted conditional capacity}
\label{sec:V61-weighted-capacity}

Let \(0\le w_\omega\le W_*\).  For \(0<t<W_*\), set
\begin{equation}
 \Omega_t:=\{\omega:w_\omega>t\},
 \qquad
 P_t:=\sum_{\omega\in\Omega_t}P_\omega,
 \qquad
 \tau^{\rm cond}(t):=
 \min_{\varnothing\ne I\subseteq[m]}
 \min\{1,\|\operatorname{Tr}_{I^c}P_t\|_{\rm op}\}.
 \label{eq:V61-weight-level-family}
\end{equation}
Define
\begin{equation}
 \mathcal K^{\rm cond}(w)
 :=\int_0^{W_*}\tau^{\rm cond}(t)\,\dd t.
 \label{eq:V61-conditional-envelope}
\end{equation}

\begin{theorem}[Multilinear layer-cake capacity bound]
\label{thm:V61-layer-cake}
For all row-factorized trace-class coefficients,
\begin{equation}
 \boxed{
 \sum_\omega w_\omega
 \|\mathcal C_\omega(A_1,\dots,A_m)\|_1
 \le
 \mathcal K^{\rm cond}(w)
 \prod_{i=1}^m\|A_i\|_1.}
 \label{eq:V61-layer-cake}
\end{equation}
\end{theorem}

\begin{proof}
Use
\(w_\omega=\int_0^{W_*}\mathbf1_{\{w_\omega>t\}}\,\dd t\), apply
Tonelli to the nonnegative trace norms, and invoke
\Cref{thm:V61-multilinear-Schur-tax} for the family \(\Omega_t\).
Integration gives \eqref{eq:V61-layer-cake}.
\end{proof}

\section{Application to JREC}
\label{sec:V61-JREC-capacity}

For the V60 joint raw family put
\begin{equation}
 W_\Lambda:=
 \mathfrak a_{oa,\Lambda}
 \mathfrak a_{bc,\Lambda}
 \mathfrak k_\Lambda.
 \label{eq:V61-JREC-weight}
\end{equation}
The V46 pair first moments and
\eqref{eq:V57-cut-first-moment} give a uniform bound
\(0\le W_\Lambda\le W_*\).  For its level projector \(P_t\), define
the two-row taxes
\begin{equation}
 \vartheta_{ar}(t):=
 \min\left\{1,
 \left\|\operatorname{Tr}_{\{a,r\}^c}P_t\right\|_{\rm op}
 \right\},
 \qquad r=b,c,
 \label{eq:V61-two-row-tax}
\end{equation}
and the endpoint-pair envelope
\begin{equation}
 \mathcal K_{a\mid bc}^{(2)}
 :=\int_0^{W_*}
 \min\{\vartheta_{ab}(t),\vartheta_{ac}(t)\}\,\dd t.
 \label{eq:V61-two-row-envelope}
\end{equation}

\begin{theorem}[Two-row conditional capacity implies JREC]
\label{thm:V61-capacity-implies-JREC}
If
\begin{equation}
 \boxed{
 \mathcal K_{a\mid bc}^{(2)}
 \le Cs_\alpha
 \sum_{r=b,c}
 \frac{h_{oa}\widehat H_{ar}h_{bc}}
      {\Xi_a\Xi_r},}
 \label{eq:V61-JREC-capacity-gate}
\end{equation}
then \(\operatorname{JREC}_{a\mid bc}\) holds for arbitrary
row-factorized trace-class coefficients and row spectators.
\end{theorem}

\begin{proof}
For each level family \(P_t\), apply
\Cref{thm:V61-multilinear-Schur-tax} first with
\(I=\{a,b\}\) and then with \(I=\{a,c\}\), taking the better bound.
Layer cake gives
\[
 \sum_\Lambda W_\Lambda\|X_\Lambda\|_1
 \le\mathcal K_{a\mid bc}^{(2)}
 \prod_{i=o,a,b,c}\|B_i\|_1.
\]
Insert \eqref{eq:V61-JREC-capacity-gate}.  The result is exactly
\eqref{eq:V60-JREC}.
\end{proof}

\section{Exact representation-theoretic form of the open operator}
\label{sec:V61-two-row-blocks}

Assume the aggregate level projector is invariant under the diagonal
product-group action.  Decompose the two retained row carriers as
\begin{equation}
 H_a\otimes H_r
 \cong
 \bigoplus_S H_S\otimes M_{S}^{a,r}.
 \label{eq:V61-two-row-isotypic}
\end{equation}

\begin{proposition}[Conditional Schur block formula]
\label{prop:V61-conditional-Schur-block}
There are positive operators \(T_{S,r}(t)\) on the multiplicity spaces
such that
\begin{align}
 \operatorname{Tr}_{\{a,r\}^c}P_t
 &=\bigoplus_S I_{H_S}\otimes T_{S,r}(t),
 \label{eq:V61-conditional-block-form}\\
 \vartheta_{ar}(t)
 &=\min\left\{1,\max_S\|T_{S,r}(t)\|_{\rm op}\right\},
 \label{eq:V61-conditional-block-norm}\\
 \operatorname{rank}P_t
 &=\sum_S d_S\operatorname{Tr}T_{S,r}(t).
 \label{eq:V61-conditional-block-trace}
\end{align}
\end{proposition}

\begin{proof}
Partial trace preserves commutation with the diagonal action on the
retained factors.  The commutant of the isotypic decomposition
\eqref{eq:V61-two-row-isotypic} is
\(\bigoplus_S I_{H_S}\otimes\mathcal B(M_S^{a,r})\), giving
\eqref{eq:V61-conditional-block-form}.  Positivity is preserved by
partial trace.  The operator norm of a finite orthogonal direct sum is
the maximum of the block norms, proving
\eqref{eq:V61-conditional-block-norm}.  Finally, take the trace of
\eqref{eq:V61-conditional-block-form}; partial trace preserves the
total trace, which is \(\operatorname{rank}P_t\).
\end{proof}

\begin{assessment}[The precise gain and the remaining obstruction]
\label{ass:V61-conditional-obstruction}
The singleton identities
\eqref{eq:V61-singleton-Schur-tax} see only one row dimension at a
time.  The two-row tax instead retains the full multiplicity operator
\(T_{S,r}(t)\), which is the first object capable of encoding two
endpoint payments without multiplying overlapping bilinear taxes.
Equation \eqref{eq:V61-conditional-block-trace} controls only a
weighted sum of its eigenvalues; it does not control the maximum in
\eqref{eq:V61-conditional-block-norm}.  The next theorem must exploit
the exact collision/cut incidence of \(P_t\) to bound those block
norms.  A total output dimension count alone is insufficient.
\end{assessment}

\section{V61 ledger and next acquisition}
\label{sec:V61-ledger}

\begin{theorem}[Integrated V61 statement]
\label{thm:V61-integrated}
Preserving V1--V60, V61 proves:
\begin{enumerate}[label=\textup{(V61.\arabic*)},leftmargin=5.5em]
\item the product-density conditional projector tax
      \eqref{eq:V61-product-density-tax};
\item the spectator-amplified multi-input aggregate Schur tax
      \eqref{eq:V61-multilinear-Schur-tax};
\item recovery of the V52 tax from singleton partial traces;
\item the multilinear weighted layer-cake theorem
      \eqref{eq:V61-layer-cake};
\item reduction of JREC to the explicit two-row capacity
      \eqref{eq:V61-JREC-capacity-gate}; and
\item the exact multiplicity-block form
      \eqref{eq:V61-conditional-block-form} of its open operator.
\end{enumerate}
\end{theorem}

\begin{proof}
The six items are
\Cref{lem:V61-product-density-tax,
thm:V61-multilinear-Schur-tax,cor:V61-V52-recovery,
thm:V61-layer-cake,thm:V61-capacity-implies-JREC,
prop:V61-conditional-Schur-block}.
\end{proof}

\begingroup
\small
\begin{longtable}{>{\raggedright\arraybackslash}p{0.29\textwidth}
                  >{\raggedright\arraybackslash}p{0.15\textwidth}
                  >{\raggedright\arraybackslash}p{0.48\textwidth}}
\toprule
\textbf{V61 row}&\textbf{Status}&\textbf{Advance or obstruction}\\
\midrule
\endfirsthead
\toprule
\textbf{V61 row}&\textbf{Status}&\textbf{Advance or obstruction}\\
\midrule
\endhead
Multilinear aggregate tax
& \MGclosed
& Arbitrary row spectators are allowed and every overlapping output is
  compressed in one invariant projector family.\\[0.35em]
Conditional layer cake
& \MGclosed
& A bounded joint weight is controlled by partial-trace operator norms
  of its aggregate level projectors.\\[0.35em]
Two-row JREC capacity
& \MGreduced
& JREC follows from the explicit envelope
  \eqref{eq:V61-JREC-capacity-gate}.\\[0.35em]
Conditional multiplicity blocks
& \MGopen
& Their traces are known exactly, but the required maximal eigenvalue
  bound has not been proved.\\[0.35em]
Complete collision and \([3,2]\)
& \MGopen
& The two-row block estimate and multi-collision regrouping remain.\\[0.35em]
Full Yang--Mills mass gap
& \MGopen
& Higher quartic/all-order MTA, WUFC/CPM, nonlinear construction,
  continuum OS existence/nontriviality, and spectral positivity remain
  downstream.\\
\bottomrule
\end{longtable}
\endgroup

\begin{openproblem}[V62 mandate: conditional two-row transform]
\label{op:V62}
\begin{enumerate}[label=\textup{(V62.\arabic*)},leftmargin=5.5em]
\item Compute \(T_{S,r}(t)\) before final recoupling from the exact
      \(oa\), \(bc\), and ternary-cut intertwiners.  Keep every LR
      multiplicity space visible.
\item Use the fixed pair-label regrouping required by V60 before
      forming \(P_t\); do not sum local partition monomials by trace
      norm.
\item Determine whether the conditional operator factors through one
      Frobenius-reciprocity channel on \(H_a\otimes H_r\), or whether a
      genuine matrix-valued LR transform is unavoidable.
\item Establish a Chernoff bound for
      \(\max_S\|T_{S,r}(t)\|_{\rm op}\), not merely for
      \(\sum_Sd_S\operatorname{Tr}T_{S,r}(t)\).
\item Test the bound on simultaneous dual singlets and the
      alternating-height array of
      \Cref{prop:V60-alternating-regression}.
\item If the block norm has no uniform contraction, isolate its exact
      noncontracting representation family and reroute the marked tree
      before returning to capacity.
\item Keep multi-collision aggregation, \([3,3]\), and \([4]\)
      suspended.  The next companion audit remains V65.
\end{enumerate}
\end{openproblem}

\begin{firewall}[End-of-V61 claim boundary]
\label{fire:V61-claim-boundary}
V61 proves the multilinear conditional Schur tax and layer-cake
capacity theorem and reduces JREC to a two-row partial-trace block
norm.  It does not prove that block norm, JREC, simultaneous-sideways
closure, fixed-type puncture aggregation, multi-collision aggregation,
complete \([3,2]\), \([3,3]\), \([4]\), global \(m=4\), all-order
MTA, WUFC, CPM, the nonlinear Perron packet, finite-edge margins,
constructive capacity, protected-sector floors, local-algebra
noncollapse, reflection positivity, a nontrivial continuum OS state,
or a positive Yang--Mills spectral gap.  The full Yang--Mills mass gap
is not proved.
\end{firewall}

\section*{Version V61 append-only record}
\addcontentsline{toc}{section}{Version V61 append-only record}
The complete V60 source was retained before this continuation.  Its
SHA--256 digest was
\begin{center}\scriptsize\ttfamily
3f54a4cd015787e441ea19472fe5a72720de5795aea0310c4203269b1f854e05.
\end{center}
No mathematical content of V60 was altered during the V61 append.  On
the next iteration, retain this full file, append before its unique
active document-closing command, increment the filename to
\texttt{\_V62.tex}, and remove only the superseded V61 file from this
Yang--Mills note series.

% =============================================================================
% END APPENDED VERSION V61 MATERIAL
% =============================================================================

% =============================================================================
% BEGIN APPENDED VERSION V62 MATERIAL
% =============================================================================

\chapter{V62: Weighted partial traces, rectangular layer cake, and
orientation coherence}
\label{ch:V62}

\section{Purpose}
\label{sec:V62-purpose}

V61 controls a bounded joint weight by integrating partial-trace norms
of the level projector of the \emph{product} weight.  The collision
weight in JREC has more structure:
\[
 W_{\lambda,\mu,\nu}
 =\mathfrak a_{oa,\lambda}
  \mathfrak a_{bc,\mu}\mathfrak k_\nu,
\]
and the three raw resolution families act on disjoint coordinate
banks before final recoupling.  V62 preserves this Cartesian
factorization by using three layer variables.  The resulting partial
trace factors exactly across banks, but only for one coherent choice
of retained input rows.  This identifies the obstruction to
multiplying locally optimal V52/V53 orientations.

\begin{firewall}[Scope of V62]
\label{fire:V62-scope}
V62 proves exact operator and integration identities and a factorized
rectangular sufficient capacity for JREC.  It does not prove the
required SU(3) bound for that capacity.  In particular, the local
minimum orientation may not be selected independently on the three
banks when a row coefficient entangles those banks.
\end{firewall}

\section{Direct weighted conditional tax}
\label{sec:V62-weighted-tax}

Retain the V61 multi-input notation.  Let
\(0\le w_\omega\le W_*\) and define the positive weighted operator
\begin{equation}
 Q_w:=\sum_\omega w_\omega P_\omega.
 \label{eq:V62-weighted-projector}
\end{equation}
For nonempty \(I\subseteq[m]\), put
\begin{equation}
 \kappa_{w,I}:=
 \min\left\{W_*,
 \|\operatorname{Tr}_{I^c}Q_w\|_{\rm op}\right\},
 \qquad
 \kappa_w^{\rm cond}:=\min_I\kappa_{w,I}.
 \label{eq:V62-weighted-conditional-tax}
\end{equation}

\begin{theorem}[Direct weighted multilinear Schur tax]
\label{thm:V62-direct-weighted-tax}
For arbitrary row-factorized trace-class coefficients with arbitrary
row spectators,
\begin{equation}
 \boxed{
 \sum_\omega w_\omega
 \|\mathcal C_\omega(A_1,\dots,A_m)\|_1
 \le
 \kappa_w^{\rm cond}
 \prod_{i=1}^m\|A_i\|_1.}
 \label{eq:V62-direct-weighted-tax}
\end{equation}
Moreover, both the direct tax and the V61 one-variable envelope obey
the common coherent bound
\begin{align}
 \kappa_w^{\rm cond}
 &\le\min_I\int_0^{W_*}
 f_I(t)\,\dd t,
 \notag\\
 \mathcal K^{\rm cond}(w)
 &\le\min_I\int_0^{W_*}
 f_I(t)\,\dd t,
 \qquad
 f_I(t):=\min\{1,\|\operatorname{Tr}_{I^c}P_t\|_{\rm op}\}.
 \label{eq:V62-direct-common-bound}
\end{align}
Neither of the two valid taxes is asserted to dominate the other.
\end{theorem}

\begin{proof}
For rank-one \(A_i=|x_i\rangle\langle y_i|\), weighted
Cauchy--Schwarz gives
\begin{align*}
 \sum_\omega w_\omega\|x_\omega\|\,\|y_\omega\|
 &\le
 \left(\sum_\omega w_\omega\|x_\omega\|^2\right)^{1/2}
 \left(\sum_\omega w_\omega\|y_\omega\|^2\right)^{1/2}.
\end{align*}
The first square is the expectation of \(Q_w\) in the product of the
carrier reductions of the \(x_i\).  The proof of
\Cref{lem:V61-product-density-tax}, with \(P_\Omega\) replaced by
\(Q_w\), bounds it by
\(\kappa_{w,I}\prod_i\|x_i\|^2\); use also
\(0\le Q_w\le W_*I\).  Treat the \(y_i\) identically and then expand
the general \(A_i\) by singular values.  This proves
\eqref{eq:V62-direct-weighted-tax}.

Layer cake gives the positive operator identity
\[
 Q_w=\int_0^{W_*}P_t\,\dd t.
\]
For every \(I\), positivity, the triangle inequality for the operator
norm, and the scalar cap give
\[
 \kappa_{w,I}
 \le\int_0^{W_*}f_I(t)\,\dd t.
\]
Minimize over \(I\) to obtain the first line of
\eqref{eq:V62-direct-common-bound}.  For the second, observe that
\[
 \mathcal K^{\rm cond}(w)
 =\int_0^{W_*}\min_J f_J(t)\,\dd t
 \le\int_0^{W_*}f_I(t)\,\dd t
\]
for every fixed \(I\), and minimize once more.  The pointwise
minimizer \(J\) may change with \(t\), while the direct weighted tax
also depends on the maximal eigenvector of an integrated positive
operator.  The two left sides therefore have no general ordering.
\end{proof}

\section{Rectangular multivariable layer cake}
\label{sec:V62-rectangular-layercake}

Suppose every row carrier factors over \(q\) disjoint banks,
\[
 H_i=\bigotimes_{\beta=1}^qH_{i,\beta},
\]
where an inactive row-bank factor is one-dimensional.  On bank
\(\beta\), let \(\{P_{\beta,\lambda}\}\) be an orthogonal projection
family and let \(0\le w_{\beta,\lambda}\le W_\beta\).  Define
\begin{align}
 \Omega_\beta(t_\beta)
 &:=\{\lambda:w_{\beta,\lambda}>t_\beta\},
 &P_\beta(t_\beta)
 &:=\sum_{\lambda\in\Omega_\beta(t_\beta)}P_{\beta,\lambda},
 \label{eq:V62-local-level-projector}\\
 \Theta_{\beta,I}(t_\beta)
 &:=\left\|
 \operatorname{Tr}_{I^c}P_\beta(t_\beta)
 \right\|_{\rm op}.
 \label{eq:V62-local-conditional-factor}
\end{align}
The partial trace in the second line is over the row-bank factors with
indices outside \(I\).

\begin{lemma}[Exact partial-trace factorization]
\label{lem:V62-partial-trace-factorization}
For the rectangular level family
\(\boldsymbol\Omega(\boldsymbol t)=
\prod_{\beta=1}^q\Omega_\beta(t_\beta)\), its aggregate projector is
\begin{equation}
 P(\boldsymbol t)=\bigotimes_{\beta=1}^qP_\beta(t_\beta),
 \label{eq:V62-rectangular-projector}
\end{equation}
and, for every retained-row set \(I\),
\begin{equation}
 \boxed{
 \left\|\operatorname{Tr}_{I^c}P(\boldsymbol t)\right\|_{\rm op}
 =\prod_{\beta=1}^q\Theta_{\beta,I}(t_\beta).}
 \label{eq:V62-partial-trace-factorization}
\end{equation}
\end{lemma}

\begin{proof}
The sum over the Cartesian family factors into the tensor product of
the local sums, proving \eqref{eq:V62-rectangular-projector}.  Partial
trace over a tensor product of disjoint bank factors is the tensor
product of the local partial traces.  The operator norm is
multiplicative on tensor products, which proves
\eqref{eq:V62-partial-trace-factorization}.
\end{proof}

Define the rectangular conditional envelope
\begin{equation}
 \mathcal K_{\rm rect}
 :=
 \int_{[0,W_1]\times\cdots\times[0,W_q]}
 \min\left\{1,
 \min_{\varnothing\ne I\subseteq[m]}
 \prod_{\beta=1}^q\Theta_{\beta,I}(t_\beta)
 \right\}\,\dd\boldsymbol t.
 \label{eq:V62-rectangular-envelope}
\end{equation}

\begin{theorem}[Rectangular multilinear capacity theorem]
\label{thm:V62-rectangular-capacity}
Let
\(w_{\boldsymbol\lambda}=\prod_{\beta=1}^q
w_{\beta,\lambda_\beta}\) and
\(P_{\boldsymbol\lambda}=\bigotimes_\beta
P_{\beta,\lambda_\beta}\).  Then
\begin{equation}
 \boxed{
 \sum_{\boldsymbol\lambda}w_{\boldsymbol\lambda}
 \|\mathcal C_{\boldsymbol\lambda}(A_1,\dots,A_m)\|_1
 \le\mathcal K_{\rm rect}
 \prod_{i=1}^m\|A_i\|_1.}
 \label{eq:V62-rectangular-capacity}
\end{equation}
\end{theorem}

\begin{proof}
For every \(\beta\), write
\[
 w_{\beta,\lambda_\beta}
 =\int_0^{W_\beta}
 \mathbf1_{\{w_{\beta,\lambda_\beta}>t_\beta\}}\,\dd t_\beta.
\]
Multiply these identities, apply Tonelli, and use
\Cref{thm:V61-multilinear-Schur-tax} on the resulting rectangular
family.  Equations \eqref{eq:V62-rectangular-projector} and
\eqref{eq:V62-partial-trace-factorization} give the integrand in
\eqref{eq:V62-rectangular-envelope}.
\end{proof}

\section{The collision--cut rectangular envelope}
\label{sec:V62-collision-envelope}

First suppress the puncture-residual pair-heavy factors and retain the
three banks \(oa\), \(bc\), and the punctured ternary core.  Sum every
final recoupling output by its complete orthogonal pinching.  This is a
trace-norm contraction, so the pre-final raw family is the Cartesian
product of the three bank families.  Put
\begin{align}
 P_{oa}(t_0)&:=
 \sum_{\mathfrak a_{oa,\lambda}>t_0}P_{oa,\lambda},
 &D_{oa}(t_0)&:=\operatorname{rank}P_{oa}(t_0),
 \notag\\
 P_{bc}(t_1)&:=
 \sum_{\mathfrak a_{bc,\mu}>t_1}P_{bc,\mu},
 &D_{bc}(t_1)&:=\operatorname{rank}P_{bc}(t_1),
 \notag\\
 P_K(t_2)&:=
 \sum_{\mathfrak k_\nu>t_2}P_{K,\nu}.
 \label{eq:V62-collision-level-projectors}
\end{align}
Let \(d_{i,e}\) denote the dimension of the input product-group
carrier of row \(i\) on bank \(e\).  Schur's lemma and
\Cref{lem:V62-partial-trace-factorization} give
\begin{align}
 \Theta_{ab}(\boldsymbol t)
 &:=%
 \frac{D_{oa}(t_0)}{d_{a,oa}}
 \frac{D_{bc}(t_1)}{d_{b,bc}}
 \left\|\operatorname{Tr}_{c}P_K(t_2)\right\|_{\rm op},
 \label{eq:V62-ab-coherent-factor}\\
 \Theta_{ac}(\boldsymbol t)
 &:=%
 \frac{D_{oa}(t_0)}{d_{a,oa}}
 \frac{D_{bc}(t_1)}{d_{c,bc}}
 \left\|\operatorname{Tr}_{b}P_K(t_2)\right\|_{\rm op}.
 \label{eq:V62-ac-coherent-factor}
\end{align}
Define
\begin{equation}
 \mathcal K_{a\mid bc}^{\rm rect}
 :=\int
 \min\{1,\Theta_{ab}(\boldsymbol t),
          \Theta_{ac}(\boldsymbol t)\}
 \,\dd t_0\dd t_1\dd t_2,
 \label{eq:V62-collision-rectangular-envelope}
\end{equation}
where the three intervals are the bounded ranges of their respective
weights.

\begin{theorem}[Rectangular capacity implies the collision-core JREC]
\label{thm:V62-rectangular-implies-JREC}
If
\begin{equation}
 \boxed{
 \mathcal K_{a\mid bc}^{\rm rect}
 \le Cs_\alpha\sum_{r=b,c}
 \frac{h_{oa}\widehat H_{ar}h_{bc}}{\Xi_a\Xi_r},}
 \label{eq:V62-rectangular-JREC-gate}
\end{equation}
then JREC holds for the simultaneous-sideways collision-core sector.
The same assertion holds after adding finitely many completely
regrouped pair-label banks, with one additional layer variable and
local conditional factor for each bank.
\end{theorem}

\begin{proof}
Apply \Cref{thm:V62-rectangular-capacity} with the retained-row choices
\(I=\{a,b\}\) and \(I=\{a,c\}\).  On the \(oa\) bank the complement
trace leaves row \(a\), so Schur's lemma gives
\(D_{oa}/d_{a,oa}\).  On \(bc\) it leaves respectively \(b\) or
\(c\).  On the core it leaves respectively \(a,b\) or \(a,c\), giving
the last factors in
\eqref{eq:V62-ab-coherent-factor}--%
\eqref{eq:V62-ac-coherent-factor}.  The scalar cap comes from the
full-set projector norm.  Thus the left side of JREC is at most
\(\mathcal K_{a\mid bc}^{\rm rect}\prod_i\|B_i\|_1\), and
\eqref{eq:V62-rectangular-JREC-gate} is exactly sufficient.  Complete
additional pair-label banks factor in the same way because their
coordinate carriers are disjoint.
\end{proof}

\section{The orientation-coherence obstruction}
\label{sec:V62-orientation-coherence}

\begin{proposition}[Local best orientations cannot be multiplied]
\label{prop:V62-orientation-obstruction}
For nonnegative numbers \(x_{\beta,I}\),
\begin{equation}
 \min_I\prod_{\beta=1}^q x_{\beta,I}
 \ge
 \prod_{\beta=1}^q\min_I x_{\beta,I}.
 \label{eq:V62-coherent-versus-local-min}
\end{equation}
The inequality may have an arbitrarily large ratio.  Likewise,
\begin{equation}
 \min\left\{1,\prod_\beta x_\beta\right\}
 \not\le\prod_\beta\min\{1,x_\beta\}
 \label{eq:V62-no-local-caps}
\end{equation}
in general.  Therefore replacing the coherent integrand in
\eqref{eq:V62-rectangular-envelope} by the product of separately
optimized V52 taxes is invalid for row-entangled coefficients.
\end{proposition}

\begin{proof}
For every fixed \(I\),
\(\prod_\beta x_{\beta,I}\ge
\prod_\beta\min_Jx_{\beta,J}\); minimize the left side over \(I\) to
obtain \eqref{eq:V62-coherent-versus-local-min}.  For a strict
arbitrarily large example with two banks and two orientations, take
\[
 (x_{1,I_1},x_{1,I_2})=(L,L^{-1}),
 \qquad
 (x_{2,I_1},x_{2,I_2})=(L^{-1},L).
\]
The coherent minimum is one while the product of local minima is
\(L^{-2}\).  For \eqref{eq:V62-no-local-caps}, take
\(x_1=L\), \(x_2=L^{-1}\): the left side is one and the right side is
\(L^{-1}\).  Let \(L\to\infty\).
\end{proof}

\begin{assessment}[Exact remaining representation problem]
\label{ass:V62-remaining-problem}
The collision pair factors in
\eqref{eq:V62-ab-coherent-factor}--%
\eqref{eq:V62-ac-coherent-factor} are explicit one-sided aggregate
dimension ratios.  The core factors are the two-row conditional
operators from V61.  A valid joint Chernoff argument must choose one
retained-row orientation coherently across all three factors, or prove
that switching orientations occurs only in a separately payable
region.  The product of locally best V53 taxes is smaller than the
actual coherent tax and cannot be used as an upper bound.
\end{assessment}

\section{V62 ledger and next acquisition}
\label{sec:V62-ledger}

\begin{theorem}[Integrated V62 statement]
\label{thm:V62-integrated}
Preserving V1--V61, V62 proves:
\begin{enumerate}[label=\textup{(V62.\arabic*)},leftmargin=5.5em]
\item the direct weighted conditional trace tax
      \eqref{eq:V62-direct-weighted-tax};
\item exact partial-trace factorization on Cartesian bank families;
\item the rectangular multilinear capacity theorem
      \eqref{eq:V62-rectangular-capacity};
\item the explicit collision-core envelope
      \eqref{eq:V62-collision-rectangular-envelope};
\item its sufficiency for collision-core JREC; and
\item the orientation-coherence obstruction
      \eqref{eq:V62-coherent-versus-local-min}--%
      \eqref{eq:V62-no-local-caps}.
\end{enumerate}
\end{theorem}

\begin{proof}
The items are
\Cref{thm:V62-direct-weighted-tax,
lem:V62-partial-trace-factorization,
thm:V62-rectangular-capacity,
thm:V62-rectangular-implies-JREC,
prop:V62-orientation-obstruction}.
\end{proof}

\begingroup
\small
\begin{longtable}{>{\raggedright\arraybackslash}p{0.29\textwidth}
                  >{\raggedright\arraybackslash}p{0.15\textwidth}
                  >{\raggedright\arraybackslash}p{0.48\textwidth}}
\toprule
\textbf{V62 row}&\textbf{Status}&\textbf{Advance or obstruction}\\
\midrule
\endfirsthead
\toprule
\textbf{V62 row}&\textbf{Status}&\textbf{Advance or obstruction}\\
\midrule
\endhead
Direct weighted tax
& \MGclosed
& The positive weighted projector controls the full multilinear trace
  functional without first forming scalar level sets.\\[0.35em]
Rectangular layer cake
& \MGclosed
& Independent pair and cut levels preserve exact bank factorization.\\[0.35em]
Collision-core capacity
& \MGreduced
& It is the explicit coherent envelope
  \eqref{eq:V62-collision-rectangular-envelope}.\\[0.35em]
Orientation switching
& \MGopen
& Separately minimizing local pair/cut taxes can underestimate the true
  joint coefficient cost by an unbounded factor.\\[0.35em]
Complete collision and \([3,2]\)
& \MGopen
& The coherent SU(3) envelope and fixed-type puncture regrouping remain.\\[0.35em]
Full Yang--Mills mass gap
& \MGopen
& Higher quartic/all-order MTA, WUFC/CPM, nonlinear construction,
  continuum OS existence/nontriviality, and spectral positivity remain
  downstream.\\
\bottomrule
\end{longtable}
\endgroup

\begin{openproblem}[V63 mandate: coherent joint Chernoff split]
\label{op:V63}
\begin{enumerate}[label=\textup{(V63.\arabic*)},leftmargin=5.5em]
\item Start from the three-variable envelope
      \eqref{eq:V62-collision-rectangular-envelope}; do not return to a
      product of locally minimized one-edge taxes.
\item Split the integration region by the coherent orientations
      \(ab\), \(ac\), and the full-set scalar cap.  On each region keep
      the same orientation through all coordinate factors.
\item Compute a matrix-valued Chernoff transform for the core factor
      \(\|\operatorname{Tr}_{c}P_K(t)\|_{\rm op}\) or its \(b\)-trace
      analogue, including LR multiplicities.
\item Determine whether the V53 one-link transform gap controls the two
      explicit pair ratios in the same orientation; if not, isolate the
      orientation-switching representation family.
\item Test simultaneous dual singlets and the alternating-height array.
      The coherent envelope, not the product of local minima, is the
      regression quantity.
\item Add puncture pair-label banks only after the three-bank core
      envelope is controlled.  Preserve the fixed-type regrouping
      requirement.
\item Keep multi-collision aggregation, \([3,3]\), and \([4]\)
      suspended.  The next companion audit remains V65.
\end{enumerate}
\end{openproblem}

\begin{firewall}[End-of-V62 claim boundary]
\label{fire:V62-claim-boundary}
V62 proves the direct weighted tax, rectangular multivariable layer
cake, exact bankwise partial-trace factorization, and the coherent
collision-core capacity criterion.  It does not prove that capacity,
JREC, simultaneous-sideways closure, fixed-type puncture aggregation,
multi-collision aggregation, complete \([3,2]\), \([3,3]\), \([4]\),
global \(m=4\), all-order MTA, WUFC, CPM, the nonlinear Perron packet,
finite-edge margins, constructive capacity, protected-sector floors,
local-algebra noncollapse, reflection positivity, a nontrivial
continuum OS state, or a positive Yang--Mills spectral gap.  The full
Yang--Mills mass gap is not proved.
\end{firewall}

\section*{Version V62 append-only record}
\addcontentsline{toc}{section}{Version V62 append-only record}
The complete V61 source was retained before this continuation.  Its
SHA--256 digest was
\begin{center}\scriptsize\ttfamily
dbb7089f9fdfe599cd11061c6e5f3c023f36d8a9e277f7d3133a0a155b725aa5.
\end{center}
No mathematical content of V61 was altered during the V62 append.  On
the next iteration, retain this full file, append before its unique
active document-closing command, increment the filename to
\texttt{\_V63.tex}, and remove only the superseded V62 file from this
Yang--Mills note series.

% =============================================================================
% END APPENDED VERSION V62 MATERIAL
% =============================================================================

% =============================================================================
% BEGIN APPENDED VERSION V63 MATERIAL
% =============================================================================

\chapter{V63: Exact ternary-core partial traces and cold conditional
capacity}
\label{ch:V63}

\section{Purpose}
\label{sec:V63-purpose}

V62 reduced the collision-core problem to a coherent rectangular
capacity whose only non-scalar local factor is
\(\|\operatorname{Tr}_{c}P_K(t)\|_{\rm op}\), or its \(b\)-trace
analogue.  V63 computes that operator exactly after grouping complete
final isotypic blocks.  It is scalar on each sequential pair
intermediate, so no uncontrolled multiplicity matrix remains at this
stage.  A new second-output-moment estimate for the ternary cut then
places every weight level inside an explicit cold LR family, to which
the V53 transform gap applies.

\begin{firewall}[Scope of V63]
\label{fire:V63-scope}
The resulting support-count capacity is a genuine estimate for the
core factor.  It does not by itself control the coherent products with
the two collision-pair dimension ratios in V62, especially when an
endpoint mass is concentrated in a bounded number of unbounded local
representations.  JREC therefore remains open.
\end{firewall}

\section{Sequential Clebsch--Gordan trace formula}
\label{sec:V63-sequential-trace}

Let \(A,B,C\) be irreducible representations of a compact group.
Choose orthogonal isometric embeddings
\[
 \iota_{S,\alpha}:H_S\longrightarrow H_A\otimes H_B,
 \qquad
 \jmath_{T,\beta}^{S}:H_T\longrightarrow H_S\otimes H_C,
\]
where
\(1\le\alpha\le N_{AB}^{S}\) and
\(1\le\beta\le N_{SC}^{T}\).  Put
\begin{equation}
 V_{T,S,\alpha,\beta}
 :=(\iota_{S,\alpha}\otimes I_C)\jmath_{T,\beta}^{S}.
 \label{eq:V63-sequential-embedding}
\end{equation}
For a family \(\Omega\) of final irreducibles, let
\begin{equation}
 P_\Omega:=
 \sum_{T\in\Omega}\sum_{S,\alpha,\beta}
 V_{T,S,\alpha,\beta}V_{T,S,\alpha,\beta}^*,
 \label{eq:V63-final-isotypic-family}
\end{equation}
where the inner sum includes all admissible paths to \(T\).

\begin{theorem}[Exact retained-pair partial trace]
\label{thm:V63-exact-core-partial-trace}
On the \((AB)C\) sequential decomposition,
\begin{equation}
 \boxed{
 \operatorname{Tr}_{C}P_\Omega
 =
 \bigoplus_{S,\alpha}
 \left[
 \frac1{d_S}\sum_{T\in\Omega}N_{SC}^{T}d_T
 \right] I_{H_{S,\alpha}}.}
 \label{eq:V63-exact-core-partial-trace}
\end{equation}
Consequently,
\begin{equation}
 \left\|\operatorname{Tr}_{C}P_\Omega\right\|_{\rm op}
 =
 \max_S
 \frac1{d_S}\sum_{T\in\Omega}N_{SC}^{T}d_T.
 \label{eq:V63-core-partial-trace-norm}
\end{equation}
The expression is independent of the multiplicity index
\(\alpha\).
\end{theorem}

\begin{proof}
For fixed \(S,T,\beta\), the operator
\(\operatorname{Tr}_{C}
 (\jmath_{T,\beta}^{S}(\jmath_{T,\beta}^{S})^*)\)
commutes with the irreducible action on \(H_S\).  Schur's lemma makes
it a scalar multiple of \(I_{H_S}\).  Its trace is the rank \(d_T\)
of the isotypic copy, hence
\[
 \operatorname{Tr}_{C}
 \bigl(\jmath_{T,\beta}^{S}(\jmath_{T,\beta}^{S})^*\bigr)
 =\frac{d_T}{d_S}I_{H_S}.
\]
Conjugating by \(\iota_{S,\alpha}\), summing the
\(N_{SC}^{T}\) values of \(\beta\), and then summing \(T\in\Omega\)
gives \eqref{eq:V63-exact-core-partial-trace}.  Orthogonality of the
\((S,\alpha)\) summands makes the operator norm the maximum of their
scalar coefficients.
\end{proof}

\begin{corollary}[Weighted final-isotypic formula]
\label{cor:V63-weighted-core-trace}
For nonnegative weights \(u_T\), let
\[
 Q_u:=\sum_Tu_TP_T,
\]
where \(P_T\) is the complete \(T\)-isotypic projection in
\(H_A\otimes H_B\otimes H_C\).  Then
\begin{equation}
 \boxed{
 \left\|\operatorname{Tr}_{C}Q_u\right\|_{\rm op}
 =
 \max_S\frac1{d_S}\sum_TN_{SC}^{T}d_Tu_T.}
 \label{eq:V63-weighted-core-trace}
\end{equation}
\end{corollary}

\begin{proof}
Multiply the \(T\)-summand in
\eqref{eq:V63-exact-core-partial-trace} by \(u_T\), sum, and use
positivity.
\end{proof}

\begin{lemma}[Complete-isotypic coarsening is coefficient-safe]
\label{lem:V63-isotypic-coarsening}
Let \(X_T=P_TXP_T\) be the coefficient block on the complete final
\(T\)-isotypic carrier, and let \(K_T\) be any operator on its
multiplicity space.  Then
\begin{equation}
 \sum_T\left|\operatorname{Tr}(X_TK_T)\right|
 \le
 \sum_T\|K_T\|_{\rm op}\,\|X_T\|_1.
 \label{eq:V63-isotypic-coarsening}
\end{equation}
If a finer multiplicity resolution was used previously, replacing it
by the complete \(P_T\) blocks does not decrease this upper bound.
Thus the core level projector in V62 may be taken to be the sum of
complete \(T\)-isotypic blocks selected by
\(\mathfrak k_T=(1+s_\alpha\Xi(T))\|K_T\|_{\rm op}\).
\end{lemma}

\begin{proof}
The first inequality is trace duality on each complete block.  A
finer multiplicity pinching is trace-norm contractive, so the sum of
the trace norms of its diagonal blocks is at most \(\|X_T\|_1\).
Conversely, off-diagonal multiplicity entries can contribute to
\(\operatorname{Tr}(X_TK_T)\); retaining the whole isotypic block is
therefore the coefficient-safe coarsening.  The final statement is
the definition of the level family for this dominating positive
functional.
\end{proof}

\section{A second output moment for the ternary cut}
\label{sec:V63-cut-second-moment}

For a complete punctured-core final irreducible
\(\boldsymbol T\), write
\[
 K_{\boldsymbol T}:=
 K_{T}^{\circ,\boldsymbol T},
 \qquad
 x_{\boldsymbol T}:=
 s_\alpha\Xi(\boldsymbol T).
\]

\begin{theorem}[Uniform ternary-cut second output moment]
\label{thm:V63-cut-second-moment}
Uniformly in the puncture, support, input representations, final
channel, and multiplicity,
\begin{equation}
 \boxed{
 x_{\boldsymbol T}^{\,2}
 \|K_{\boldsymbol T}\|_{\rm op}\le C.}
 \label{eq:V63-cut-second-moment}
\end{equation}
Consequently the V60 weighted cut size obeys
\begin{equation}
 \boxed{
 \mathfrak k_{\boldsymbol T}
 :=
 (1+x_{\boldsymbol T})\|K_{\boldsymbol T}\|_{\rm op}
 \le\frac C{1+x_{\boldsymbol T}}.}
 \label{eq:V63-cut-weight-decay}
\end{equation}
\end{theorem}

\begin{proof}
Use the five global row-partition terms in the exact ternary cumulant,
equivalently the expanded right side of
\eqref{eq:V57-cut-identity}.  After all intermediate channels in one
such term have been resolved, its scalar multiplier has the form
\[
 q_{\alpha,\boldsymbol V_1}\cdots
 q_{\alpha,\boldsymbol V_\ell},
 \qquad 1\le\ell\le3,
\]
and fusion geometry gives
\[
 \Xi(\boldsymbol T)\le C\sum_{j=1}^{\ell}
 \Xi(\boldsymbol V_j).
\]
Put \(x_j=s_\alpha\Xi(\boldsymbol V_j)\).  Expanding
\((\sum_jx_j)^2\), each diagonal term is bounded by the product second
moment \eqref{eq:V32-balanced-second-moment}, with the other
multipliers bounded by one.  Each cross term is bounded by two product
first moments \eqref{eq:V46-product-first-moment}.  There are at most
three factors and five cumulant terms.  Their recoupling maps are
isometries, so the triangle inequality in operator norm proves
\eqref{eq:V63-cut-second-moment}.

The universal cut cap gives
\(\|K_{\boldsymbol T}\|_{\rm op}\le C\).  If
\(x_{\boldsymbol T}\le1\), this cap proves
\eqref{eq:V63-cut-weight-decay}.  If
\(x_{\boldsymbol T}\ge1\), use
\eqref{eq:V63-cut-second-moment} and
\((1+x)/x^2\le2/x\le4/(1+x)\).
\end{proof}

\begin{corollary}[Cut levels are cold-output families]
\label{cor:V63-cut-level-cold}
For \(0<t<C\), define
\[
 \Omega_K(t):=
 \{\boldsymbol T:\mathfrak k_{\boldsymbol T}>t\}.
\]
Then
\begin{equation}
 \boldsymbol T\in\Omega_K(t)
 \quad\Longrightarrow\quad
 \Xi(\boldsymbol T)<\frac{C}{s_\alpha t}.
 \label{eq:V63-cut-level-cold}
\end{equation}
\end{corollary}

\begin{proof}
Rearrange \eqref{eq:V63-cut-weight-decay}.
\end{proof}

\section{Conditional LR Chernoff bound}
\label{sec:V63-conditional-Chernoff}

Return to the product group over the punctured-core coordinates.
Fix a sequential intermediate
\(\boldsymbol S=(S_e)_e\) in
\(\boldsymbol A\otimes\boldsymbol B\), and write
\(\boldsymbol C=(C_e)_e\).  For \(\theta>0\), define
\begin{equation}
 \mathcal Z_{S_e,C_e}^{S_e}(\theta)
 :=
 \frac1{d_{S_e}}
 \sum_{T_e}N_{S_eC_e}^{T_e}d_{T_e}
 e^{-\theta\xi(T_e)}.
 \label{eq:V63-local-conditional-transform}
\end{equation}

\begin{theorem}[Core conditional Chernoff estimate]
\label{thm:V63-core-Chernoff}
For every \(t>0\) and \(\theta>0\),
\begin{equation}
 \boxed{
 \left\|\operatorname{Tr}_{\boldsymbol C}P_K(t)\right\|_{\rm op}
 \le
 \sup_{\boldsymbol S}
 \exp\!\left\{\frac{C\theta}{s_\alpha t}\right\}
 \prod_e
 \mathcal Z_{S_e,C_e}^{S_e}(\theta).}
 \label{eq:V63-core-Chernoff}
\end{equation}
\end{theorem}

\begin{proof}
By \Cref{thm:V63-exact-core-partial-trace}, the left side is
\[
 \max_{\boldsymbol S}\frac1{d_{\boldsymbol S}}
 \sum_{\boldsymbol T\in\Omega_K(t)}
 N_{\boldsymbol S\boldsymbol C}^{\boldsymbol T}
 d_{\boldsymbol T}.
\]
On the cold family \eqref{eq:V63-cut-level-cold},
\[
 \mathbf1_{\{\boldsymbol T\in\Omega_K(t)\}}
 \le
 \exp\!\left\{\frac{C\theta}{s_\alpha t}\right\}
 e^{-\theta\Xi(\boldsymbol T)}.
\]
The output dimension, LR multiplicity, input dimension, and
\(e^{-\theta\Xi(\boldsymbol T)}\) all factor over product-group
coordinates.  Summing the unrestricted tuples gives the product in
\eqref{eq:V63-core-Chernoff}; then take the supremum in
\(\boldsymbol S\).
\end{proof}

\begin{lemma}[Uniform oriented gap whenever a coordinate is active]
\label{lem:V63-active-transform-gap}
There exist \(\theta_*>0\) and \(\gamma_*<1\) such that
\begin{equation}
 \mathcal Z_{R,Q}^{R}(\theta_*)\le\gamma_*
 \label{eq:V63-active-transform-gap}
\end{equation}
for every pair of SU(3) irreducibles \(R,Q\) which are not both
trivial.
\end{lemma}

\begin{proof}
If both \(R,Q\) are nontrivial, this is the \(R\)-oriented half of
\Cref{lem:V53-universal-transform-gap}.  If \(Q\) is trivial and
\(R\) nontrivial, the transform equals
\(e^{-\theta\xi(R)}\), uniformly below one for every fixed
\(\theta>0\).  If \(R\) is trivial and \(Q\) nontrivial, it equals
\(d_Qe^{-\theta\xi(Q)}\).  The SU(3) Weyl dimension formula grows
polynomially in the highest weight, whereas \(\xi(Q)\) grows
quadratically; hence
\[
 \sup_{Q\ne\mathbf1}d_Qe^{-\theta\xi(Q)}
 \longrightarrow0
 \qquad(\theta\to\infty).
\]
Increase the V53 value of \(\theta_*\) until the last two families are
also bounded by one common \(\gamma_*<1\).
\end{proof}

Let \(N_C\) be the number of punctured-core coordinates at which
\(C_e\) is nontrivial, and put \(N_C^\sharp:=\max\{1,N_C\}\).

\begin{corollary}[Support-count conditional cut capacity]
\label{cor:V63-support-count-capacity}
There are constants \(c,C>0\) such that
\begin{equation}
 \boxed{
 \int_0^{C}
 \min\left\{1,
 \|\operatorname{Tr}_{\boldsymbol C}P_K(t)\|_{\rm op}
 \right\}\,\dd t
 \le
 C\min\left\{1,\frac1{s_\alpha N_C^\sharp}\right\}
 +Ce^{-cN_C}.}
 \label{eq:V63-support-count-capacity}
\end{equation}
The same statement holds with \(B\) and \(C\) interchanged.
\end{corollary}

\begin{proof}
If \(N_C=0\), the scalar cap on the integral proves the displayed
bound.  Suppose \(N_C\ge1\).  For every \(\boldsymbol S\), at least the \(N_C\) coordinates active
in \(\boldsymbol C\) have
\((S_e,C_e)\ne(\mathbf1,\mathbf1)\).  Therefore
\Cref{lem:V63-active-transform-gap} in
\eqref{eq:V63-core-Chernoff} gives
\[
 \|\operatorname{Tr}_{\boldsymbol C}P_K(t)\|_{\rm op}
 \le
 \exp\left\{\frac{C\theta_*}{s_\alpha t}-cN_C\right\}.
\]
If \(s_\alpha N_C\) is bounded, the scalar cap on the integral proves
the result.  Otherwise choose
\(t_0=C_0/(s_\alpha N_C)\), with \(C_0\) large enough that the exponent
is at most \(-cN_C/2\) for \(t\ge t_0\).  Bound the interval
\((0,t_0)\) by its length and the remainder by
\(Ce^{-cN_C/2}\).  This proves
\eqref{eq:V63-support-count-capacity}.
\end{proof}

\begin{corollary}[Bounded-local mass form]
\label{cor:V63-bounded-local-core}
If \(\boldsymbol C\) is nontrivial and
\(A_{C_e}\le L\) on its active support, then
\[
 N_C\ge L^{-1}\Xi(\boldsymbol C),
\]
and the right side of
\eqref{eq:V63-support-count-capacity} is at most
\[
 C_L\min\left\{1,\frac1{s_\alpha\Xi(\boldsymbol C)}\right\}
 +C_Le^{-c_L\Xi(\boldsymbol C)}.
\]
\end{corollary}

\begin{proof}
Sum \(A_{C_e}\le L\) over the \(N_C\) active coordinates and use the
monotonicity of the two terms in
\eqref{eq:V63-support-count-capacity}.
\end{proof}

\begin{assessment}[What remains after scalarizing the core]
\label{ass:V63-remaining}
The core partial trace is not an arbitrary multiplicity matrix once
complete final isotypic blocks are retained: it is the explicit
oriented LR sum \eqref{eq:V63-core-partial-trace-norm}.  Its weighted
levels have the support-count contraction
\eqref{eq:V63-support-count-capacity}.  The unresolved cases are now
sharper:
\begin{enumerate}[label=\textup{(\roman*)},leftmargin=3.7em]
\item the coherent product with collision-pair ratios that may exceed
      one before the global cap;
\item bounded-support but unbounded-height input rows, for which
      \(N_C\) does not control \(\Xi(\boldsymbol C)\); and
\item fixed-type regrouping of the puncture pair-label banks.
\end{enumerate}
\end{assessment}

\section{V63 ledger and next acquisition}
\label{sec:V63-ledger}

\begin{theorem}[Integrated V63 statement]
\label{thm:V63-integrated}
Preserving V1--V62, V63 proves:
\begin{enumerate}[label=\textup{(V63.\arabic*)},leftmargin=5.5em]
\item the exact sequential partial trace
      \eqref{eq:V63-exact-core-partial-trace};
\item its weighted final-isotypic form
      \eqref{eq:V63-weighted-core-trace};
\item the ternary-cut second output moment
      \eqref{eq:V63-cut-second-moment};
\item conversion of cut-weight levels into cold-output families;
\item the conditional LR Chernoff estimate
      \eqref{eq:V63-core-Chernoff}; and
\item the uniform support-count capacity
      \eqref{eq:V63-support-count-capacity}.
\end{enumerate}
\end{theorem}

\begin{proof}
The six items are
\Cref{thm:V63-exact-core-partial-trace,
cor:V63-weighted-core-trace,thm:V63-cut-second-moment,
cor:V63-cut-level-cold,thm:V63-core-Chernoff,
cor:V63-support-count-capacity}.
\end{proof}

\begingroup
\small
\begin{longtable}{>{\raggedright\arraybackslash}p{0.29\textwidth}
                  >{\raggedright\arraybackslash}p{0.15\textwidth}
                  >{\raggedright\arraybackslash}p{0.48\textwidth}}
\toprule
\textbf{V63 row}&\textbf{Status}&\textbf{Advance or obstruction}\\
\midrule
\endfirsthead
\toprule
\textbf{V63 row}&\textbf{Status}&\textbf{Advance or obstruction}\\
\midrule
\endhead
Core conditional block
& \MGclosed
& Complete final isotypic grouping makes it an explicit scalar LR sum
  on every sequential intermediate copy.\\[0.35em]
Cut second output moment
& \MGclosed
& Every cut-weight level is contained in a quantitative cold-output
  family.\\[0.35em]
Support-count core capacity
& \MGclosed
& The conditional core factor gains
  \(O((s_\alpha N_C)^{-1})\) plus an exponential remainder.\\[0.35em]
Coherent collision product
& \MGopen
& Pair partial-trace ratios may exceed one and cannot be capped
  independently under row entanglement.\\[0.35em]
Unbounded local height
& \MGopen
& A fixed support count need not control a hot original-row mass.\\[0.35em]
Complete collision and \([3,2]\)
& \MGopen
& The coherent height-sensitive envelope and puncture regrouping remain.\\[0.35em]
Full Yang--Mills mass gap
& \MGopen
& Higher quartic/all-order MTA, WUFC/CPM, nonlinear construction,
  continuum OS existence/nontriviality, and spectral positivity remain
  downstream.\\
\bottomrule
\end{longtable}
\endgroup

\begin{openproblem}[V64 mandate: height-sensitive coherent envelope]
\label{op:V64}
\begin{enumerate}[label=\textup{(V64.\arabic*)},leftmargin=5.5em]
\item Insert \eqref{eq:V63-core-partial-trace-norm} into the coherent
      integrand \eqref{eq:V62-collision-rectangular-envelope}, retaining
      the pair dimension ratios without separate caps.
\item Split the local transform into bounded-height and
      unbounded-height coordinates.  Use
      \eqref{eq:V63-support-count-capacity} only on the former.
\item On an unbounded-height coordinate, use the exact SU(3) Weyl
      dimension and Casimir formulas to seek a mass-weighted
      contraction stronger than the uniform gap
      \eqref{eq:V63-active-transform-gap}.
\item Determine whether one coherent orientation controls both pair
      ratios on each height branch.  If orientations switch, retain the
      switching region as a separate positive integral rather than
      multiplying local minima.
\item Test simultaneous dual singlets and the alternating-height array.
\item Add fixed-type puncture pair-label banks only after the
      collision-core height problem closes.
\item Keep multi-collision aggregation, \([3,3]\), and \([4]\)
      suspended.  The next companion audit is V65.
\end{enumerate}
\end{openproblem}

\begin{firewall}[End-of-V63 claim boundary]
\label{fire:V63-claim-boundary}
V63 proves the exact scalar conditional core trace, ternary-cut second
output moment, cold-level reduction, and support-count conditional
capacity.  It does not prove the height-sensitive coherent envelope,
JREC, simultaneous-sideways closure, fixed-type puncture aggregation,
multi-collision aggregation, complete \([3,2]\), \([3,3]\), \([4]\),
global \(m=4\), all-order MTA, WUFC, CPM, the nonlinear Perron packet,
finite-edge margins, constructive capacity, protected-sector floors,
local-algebra noncollapse, reflection positivity, a nontrivial
continuum OS state, or a positive Yang--Mills spectral gap.  The full
Yang--Mills mass gap is not proved.
\end{firewall}

\section*{Version V63 append-only record}
\addcontentsline{toc}{section}{Version V63 append-only record}
The complete V62 source was retained before this continuation.  Its
SHA--256 digest was
\begin{center}\scriptsize\ttfamily
8a41e5bf4c8b19933dc2fb2ab8bea2f9ae29f0903c9b86e309c5b58a491840d5.
\end{center}
No mathematical content of V62 was altered during the V63 append.  On
the next iteration, retain this full file, append before its unique
active document-closing command, increment the filename to
\texttt{\_V64.tex}, and remove only the superseded V63 file from this
Yang--Mills note series.

% =============================================================================
% END APPENDED VERSION V63 MATERIAL
% =============================================================================

% =============================================================================
% BEGIN APPENDED VERSION V64 MATERIAL
% =============================================================================

\chapter{V64: Mass-weighted conditional transforms and the
core-only height barrier}
\label{ch:V64}

\section{Purpose and point of departure}
\label{sec:V64-purpose}

V63 made the complete ternary-core partial trace an explicit
one-sided LR transform, but used only a uniform gap per active
coordinate.  The first task here is to retain the Casimir height in
that transform.  The resulting estimate is uniform in a
\emph{single fixed orientation}; in particular, it does not switch
orientation from coordinate to coordinate in an alternating-height
array.

The second task is diagnostic.  The output-moment majorant
\(\mathfrak k_{\boldsymbol T}\lesssim
(1+s_\alpha\Xi(\boldsymbol T))^{-1}\), combined with the exact V63
partial trace, does not by itself have the inverse-input-mass capacity
needed by JREC.  A dual symmetric family gives a sharp square-root
loss.  Thus the next acquisition must go upstream to the actual
pair--core weighted operator or to an input-sensitive cut estimate;
iterating the same core-only Chernoff bound would be circular.

\begin{firewall}[Scope of V64]
\label{fire:V64-scope}
V64 proves a quantitative height decay for the one-sided conditional
transform, a global mass-product consequence, a height-sensitive core
capacity, and a sharp obstruction to closing JREC from the V63
output-weight majorant alone.  The obstruction concerns that
majorant, not the actual ternary cut coefficient: it is not a
counterexample to JREC or to the mass gap.  The pair-bank ratios in
the coherent V62 envelope remain uncapped and JREC remains open.
\end{firewall}

\section{A quantitative one-link transform}
\label{sec:V64-local-transform}

Retain the V52--V63 convention
\[
 \xi(R)=0\quad(R=\mathbf1),
 \qquad
 \xi(R)=1+C_2(R)\quad(R\ne\mathbf1).
\]
For an ordered pair \((R,Q)\), put
\begin{equation}
 m(R,Q):=\max\{\xi(R),\xi(Q)\}.
 \label{eq:V64-local-pair-mass}
\end{equation}

\begin{theorem}[Mass decay of the fixed-orientation transform]
\label{thm:V64-local-transform-decay}
There are constants
\(\theta_*>0\), \(\gamma_*\in(0,1)\), and \(C<\infty\), depending
only on \(\mathrm{SU}(3)\), such that every pair of irreducibles
\((R,Q)\ne(\mathbf1,\mathbf1)\) satisfies
\begin{equation}
 \boxed{
 \mathcal Z_{R,Q}^{R}(\theta_*)
 \le
 \min\left\{\gamma_*,
 \frac{C}{1+m(R,Q)}\right\}.}
 \label{eq:V64-local-transform-decay}
\end{equation}
The same assertion holds for the \(Q\)-oriented transform after
interchanging \(R,Q\).
\end{theorem}

\begin{proof}
The first entry of the minimum is exactly
\Cref{lem:V63-active-transform-gap}.  It remains to prove the
mass decay at the same fixed value of \(\theta_*\).

By \Cref{lem:V34-dimension-Casimir}, with harmless changes of the
constants at the trivial representation,
\begin{equation}
 c(1+\xi(U))\le d_U
 \le C(1+\xi(U))^{3/2}
 \label{eq:V64-dimension-mass-comparison}
\end{equation}
for every irreducible \(U\).  Let \(r=\xi(R)\), \(q=\xi(Q)\).
Choose
\(\eta>0\) so small that
\(\sqrt\eta\le(2C_{\rm rf})^{-1}\), where \(C_{\rm rf}\) is the
constant in \eqref{eq:V53-Casimir-reverse-fusion}.

Suppose first that \(q\ge\eta r\).  Frobenius reciprocity and
dimension counting, as in
\eqref{eq:V53-reciprocity-dimension}, give
\[
 N_{RQ}^{T}d_Q\le d_Rd_T.
\]
Consequently
\begin{align}
 \mathcal Z_{R,Q}^{R}(\theta_*)
 &\le
 \frac1{d_Q}\sum_Td_T^2e^{-\theta_*\xi(T)}
 \notag\\
 &\le
 \frac{C_{\theta_*}}{1+q}
 \le
 \frac C{1+\max\{r,q\}}.
 \label{eq:V64-balanced-transform}
\end{align}
The heat sum is finite by the Weyl dimension and Casimir formulas.
The last comparison follows from \(q\ge\eta r\).  Notice that this
case also includes \(R=\mathbf1\), \(Q\ne\mathbf1\).

It remains to consider \(q<\eta r\).  Then \(r>0\).  If
\(N_{RQ}^{T}>0\), reverse fusion gives
\[
 \sqrt r
 \le C_{\rm rf}\bigl(\sqrt{\xi(T)}+\sqrt q\bigr),
\]
and hence
\[
 \sqrt{\xi(T)}
 \ge
 \left(C_{\rm rf}^{-1}-\sqrt\eta\right)\sqrt r
 \ge(2C_{\rm rf})^{-1}\sqrt r.
\]
Thus \(\xi(T)\ge c r\) for every output.  The exact dimension
identity
\(\sum_TN_{RQ}^{T}d_T=d_Rd_Q\) now yields
\begin{align}
 \mathcal Z_{R,Q}^{R}(\theta_*)
 &\le d_Qe^{-c\theta_*r}
 \le C(1+q)^{3/2}e^{-c\theta_*r}
 \notag\\
 &\le C(1+r)^{3/2}e^{-c\theta_*r}
 \le\frac C{1+r}.
 \label{eq:V64-asymmetric-transform}
\end{align}
Here \(m(R,Q)=r\).  Equations
\eqref{eq:V64-balanced-transform} and
\eqref{eq:V64-asymmetric-transform} prove the second entry in
\eqref{eq:V64-local-transform-decay}.  Interchanging the inputs proves
the final assertion.
\end{proof}

\begin{remark}[Where the two mechanisms act]
\label{rem:V64-transform-mechanisms}
When \(Q\) has mass comparable to or larger than \(R\), the denominator
\(d_Q\) created by reciprocity pays the transform.  When \(Q\) is much
smaller, reverse fusion forces every output to retain a fixed fraction
of the \(R\)-mass, and the heat factor pays exponentially.  Both
mechanisms estimate the same \(R\)-oriented transform; no local choice
of orientation has been made.
\end{remark}

\begin{corollary}[A global mass-weighted product]
\label{cor:V64-global-transform-product}
For a finite coordinate family \((R_e,Q_e)_e\), let
\[
 \mathcal A:=\{e:(R_e,Q_e)\ne(\mathbf1,\mathbf1)\},
 \qquad
 M:=\sum_{e\in\mathcal A}m(R_e,Q_e).
\]
If \(\mathcal A\ne\varnothing\), then
\begin{equation}
 \boxed{
 \prod_e\mathcal Z_{R_e,Q_e}^{R_e}(\theta_*)
 \le\frac C{1+M}.}
 \label{eq:V64-global-transform-product}
\end{equation}
At inactive coordinates the transform is exactly one.
\end{corollary}

\begin{proof}
Put \(N=|\mathcal A|\), and choose \(e_*\in\mathcal A\) with
\(m(R_{e_*},Q_{e_*})\ge M/N\).  Apply the mass entry of
\eqref{eq:V64-local-transform-decay} at \(e_*\), the gap entry at the
other active coordinates, and the identity at inactive coordinates:
\[
 \prod_e\mathcal Z_{R_e,Q_e}^{R_e}(\theta_*)
 \le
 \frac{C\gamma_*^{N-1}}{1+m(R_{e_*},Q_{e_*})}
 \le
 \frac{CN\gamma_*^{N-1}}{1+M}.
\]
The sequence \(N\gamma_*^{N-1}\) is bounded uniformly in \(N\).
\end{proof}

\begin{proposition}[The inverse mass power is sharp for this transform]
\label{prop:V64-transform-sharpness}
Let \(R_n=(n,0)\) and \(Q_n=(0,n)=R_n^*\).  For every fixed
\(\theta>0\),
\begin{equation}
 \mathcal Z_{R_n,Q_n}^{R_n}(\theta)
 \ge\frac1{d_{R_n}}
 \asymp\frac1{1+m(R_n,Q_n)}.
 \label{eq:V64-transform-sharpness}
\end{equation}
In particular, no uniform replacement of the last exponent \(1\) by
\(1+\varepsilon\), \(\varepsilon>0\), is possible.
\end{proposition}

\begin{proof}
The singlet occurs once in \(R_n\otimes R_n^*\), and its heat factor
is one because \(\xi(\mathbf1)=0\).  Its contribution to the
\(R_n\)-oriented transform is \(1/d_{R_n}\).  The exact formulas
\[
 d_{R_n}=\frac{(n+1)(n+2)}2,
 \qquad
 C_2(R_n)=\frac{n^2+3n}{3}
\]
prove the comparison in \eqref{eq:V64-transform-sharpness}.
\end{proof}

\section{Insertion into the exact ternary core}
\label{sec:V64-core-insertion}

For an intermediate input
\(\boldsymbol S=(S_e)_e\) and the traced core row
\(\boldsymbol C=(C_e)_e\), define
\begin{equation}
 M(\boldsymbol S,\boldsymbol C)
 :=
 \sum_e\max\{\xi(S_e),\xi(C_e)\}.
 \label{eq:V64-core-pair-mass}
\end{equation}
Then
\begin{equation}
 M(\boldsymbol S,\boldsymbol C)
 \ge\Xi(\boldsymbol C).
 \label{eq:V64-core-pair-mass-lower}
\end{equation}

\begin{corollary}[Height-sensitive conditional core level]
\label{cor:V64-height-core-level}
For \(0<t<C_K\), where \(C_K\) is the uniform range of the V63 cut
weight,
\begin{equation}
 \boxed{
 \left\|\operatorname{Tr}_{\boldsymbol C}P_K(t)\right\|_{\rm op}
 \le
 \frac C{1+\Xi(\boldsymbol C)}
 \exp\left\{\frac C{s_\alpha t}\right\}.}
 \label{eq:V64-height-core-level}
\end{equation}
The same estimate holds after interchanging the two possible traced
core rows.
\end{corollary}

\begin{proof}
Apply \Cref{cor:V64-global-transform-product} to the product in
\eqref{eq:V63-core-Chernoff}, with
\((R_e,Q_e)=(S_e,C_e)\), and then use
\eqref{eq:V64-core-pair-mass-lower}.  If every coordinate is inactive,
then \(\boldsymbol C\) is trivial and the right side, after enlarging
the constant, reduces to the scalar bound.  The other orientation is
identical.
\end{proof}

\begin{corollary}[Logarithmic height capacity from Chernoff alone]
\label{cor:V64-log-height-capacity}
There is a constant \(C\) such that
\begin{align}
 &\int_0^{C_K}
 \min\left\{1,
 \left\|\operatorname{Tr}_{\boldsymbol C}P_K(t)\right\|_{\rm op}
 \right\}\,\dd t
 \notag\\
 &\qquad\le
 C\min\left\{1,
 \frac1{s_\alpha\log(2+\Xi(\boldsymbol C))}\right\}
 +\frac C{\sqrt{1+\Xi(\boldsymbol C)}}.
 \label{eq:V64-log-height-capacity}
\end{align}
One may take the minimum of this bound and the support-count bound
\eqref{eq:V63-support-count-capacity}.
\end{corollary}

\begin{proof}
Write \(X=\Xi(\boldsymbol C)\) and
\(L=\log(2+X)\).  If
\(2C/(s_\alpha L)\ge C_K\), the scalar cap proves the first term after
enlarging its constant.  Otherwise put
\[
 t_0=\frac{2C}{s_\alpha L},
\]
with the numerator chosen at least as large as the constant in the
exponent of \eqref{eq:V64-height-core-level}.  Bound the interval
\((0,t_0)\) by its length.  For \(t\ge t_0\),
\[
 \exp\left\{\frac C{s_\alpha t}\right\}
 \le e^{L/2}=\sqrt{2+X},
\]
so the remaining integrand is at most
\(C(1+X)^{-1/2}\).  Its interval has bounded length.  This proves
\eqref{eq:V64-log-height-capacity}.  Both this estimate and V63 are
valid upper bounds for the same integral, so their minimum is valid.
\end{proof}

\section{A sharp barrier for the output-weight majorant}
\label{sec:V64-core-barrier}

The logarithm in \eqref{eq:V64-log-height-capacity} is partly a
Chernoff loss.  The following exact calculation removes that loss but
also shows that the core-only route still misses a square root.

\begin{theorem}[Dual-symmetric model capacity]
\label{thm:V64-dual-model-capacity}
Let
\[
 R_n=(n,0),\qquad Q_n=(0,n),
\]
and let \(P_k\) be the complete projection onto the output
\((k,k)\) in
\[
 R_n\otimes Q_n=\bigoplus_{k=0}^{n}(k,k).
\]
For \(0<s\le1\), set
\[
 u_0=1,\qquad
 u_k=\frac1{1+s(k+1)^2}\quad(1\le k\le n),
 \qquad
 P_n(t)=\sum_{\{k:u_k>t\}}P_k,
\]
and define the one-sided layer capacity
\begin{equation}
 \mathcal L_n(s)
 :=
 \int_0^1
 \min\left\{1,
 \|\operatorname{Tr}_{Q_n}P_n(t)\|_{\rm op}\right\}\,\dd t.
 \label{eq:V64-dual-model-capacity-definition}
\end{equation}
There are numerical constants \(c,C>0\) such that, for all
\(n\ge1\),
\begin{equation}
 \boxed{
 c\min\left\{1,\frac1{sn}\right\}
 \le\mathcal L_n(s)
 \le
 C\min\left\{1,\frac1{sn}\right\}.}
 \label{eq:V64-dual-model-capacity}
\end{equation}
Since \(\xi(Q_n)\asymp n^2\), the decaying regime is
\[
 \mathcal L_n(s)\asymp
 \frac1{s\sqrt{\xi(Q_n)}},
\]
not \(1/(s\xi(Q_n))\).
\end{theorem}

\begin{proof}
The decomposition and dimensions from
\Cref{lem:V38-boundary-decompositions} give
\[
 d_{R_n}=\frac{(n+1)(n+2)}2\asymp n^2,
 \qquad
 d_{(k,k)}=(k+1)^3.
\]
The exact V63 trace formula, or Schur's lemma directly, gives
\begin{equation}
 \|\operatorname{Tr}_{Q_n}P_n(t)\|_{\rm op}
 =
 \frac1{d_{R_n}}\sum_{\{k:u_k>t\}}(k+1)^3.
 \label{eq:V64-dual-model-level}
\end{equation}
Moreover, for \(0\le K\le n\),
\begin{equation}
 c(K+1)^4
 \le\sum_{k=0}^{K}(k+1)^3
 \le C(K+1)^4.
 \label{eq:V64-cubic-partial-sum}
\end{equation}

The upper bound is trivial when \(sn\le C_0\), because then
\(1\le C_0/(sn)\) after increasing the final constant.  Suppose
\(sn>C_0\), and put \(t_*=C_0/(sn)<1\), with \(C_0\) a fixed
sufficiently large constant.  The portion \(0<t<t_*\) contributes at
most \(t_*\).
For \(t\ge t_*\), the inequality \(u_k>t\) implies
\[
 k+1\le C(st)^{-1/2}.
\]
Equations \eqref{eq:V64-dual-model-level} and
\eqref{eq:V64-cubic-partial-sum} therefore give
\[
 \|\operatorname{Tr}_{Q_n}P_n(t)\|_{\rm op}
 \le\frac C{n^2s^2t^2}.
\]
Integrating this estimate from \(t_*\) to one gives
\[
 \mathcal L_n(s)
 \le\frac C{sn}
 +\frac C{n^2s^2t_*}
 \le\frac C{sn}.
\]

For the lower bound, first suppose \(sn\ge1\).  Choose a fixed
\(a\in(0,1/4)\), and take
\(K=\lfloor a\sqrt n\rfloor\).  After adjusting constants for the
finitely many small \(n\), \eqref{eq:V64-cubic-partial-sum} gives
\[
 \frac1{d_{R_n}}\sum_{k=0}^{K}(k+1)^3\ge c_a.
\]
Also
\[
 u_k\ge\frac1{1+s(K+1)^2}
 \ge\frac{c_a'}{sn}
 \qquad(0\le k\le K),
\]
where for \(k=0\) the declared value \(u_0=1\) is even larger.
Thus the integrand in
\eqref{eq:V64-dual-model-capacity-definition} is at least \(c_a\)
through an interval of length \(c_a'/(sn)\).  This proves the lower
bound \(c/(sn)\).

If \(sn<1\), the same choice of \(K\) gives
\(u_k\ge c_a'>0\) for \(0\le k\le K\), and hence the integrand is at
least \(c_a\) on an interval of fixed positive length.  Therefore
\(\mathcal L_n(s)\ge c\).  This completes both regimes.
\end{proof}

\begin{corollary}[What the model does and does not disprove]
\label{cor:V64-majorant-no-go}
The exact partial trace
\eqref{eq:V63-core-partial-trace-norm} and the sole pointwise
information
\[
 0\le\mathfrak k_T\le\frac C{1+s_\alpha\Xi(T)}
\]
cannot imply a uniform core capacity
\(C/(s_\alpha\Xi(\boldsymbol C))\): the admissible model weights in
\Cref{thm:V64-dual-model-capacity} violate that rate by a factor
comparable to \(\sqrt{\Xi(\boldsymbol C)}\) when
\(s_\alpha\sqrt{\Xi(\boldsymbol C)}\to\infty\).
This conclusion does not assert that the actual cut weights saturate
the model.
\end{corollary}

\begin{proof}
For the dual-symmetric one-coordinate family,
\(\Xi(\boldsymbol C)=\xi(Q_n)\asymp n^2\), while
\eqref{eq:V64-dual-model-capacity} is \(c/(s_\alpha n)\) in the
decaying regime.  The model weights obey the displayed pointwise
majorant with constant one because
\(\xi((k,k))=(k+1)^2\) for \(k\ge1\) and
\(\xi(\mathbf1)=0\).  Hence no argument using only those two pieces
of information can prove the stronger rate.  Additional structure of
the actual weights may still do so.
\end{proof}

\section{Coherent V62 insertion and the location of switching}
\label{sec:V64-coherent-insertion}

For the three level variables of V62, abbreviate
\begin{align}
 A(t_0)&:=\frac{D_{oa}(t_0)}{d_{a,oa}},
 &
 B_b(t_1)&:=\frac{D_{bc}(t_1)}{d_{b,bc}},
 &
 B_c(t_1)&:=\frac{D_{bc}(t_1)}{d_{c,bc}},
 \label{eq:V64-pair-ratios}\\
 \Phi_C(t_2)&:=
 \frac C{1+\Xi(\boldsymbol C)}
 \exp\left\{\frac C{s_\alpha t_2}\right\},
 &
 \Phi_B(t_2)&:=
 \frac C{1+\Xi(\boldsymbol B)}
 \exp\left\{\frac C{s_\alpha t_2}\right\}.
 \label{eq:V64-oriented-core-majorants}
\end{align}
No scalar cap is imposed separately on any of these factors.

\begin{proposition}[Fixed-orientation insertion into the rectangular envelope]
\label{prop:V64-coherent-insertion}
The coherent V62 integrand satisfies
\begin{align}
 &\min\{1,\Theta_{ab}(\boldsymbol t),
           \Theta_{ac}(\boldsymbol t)\}
 \notag\\
 &\qquad\le
 \min\{1,
 A(t_0)B_b(t_1)\Phi_C(t_2),
 A(t_0)B_c(t_1)\Phi_B(t_2)\}.
 \label{eq:V64-coherent-insertion}
\end{align}
If
\begin{align}
 \mathcal E_{ab}
 &:=
 \{(t_1,t_2):
 B_b(t_1)\Phi_C(t_2)
 \le B_c(t_1)\Phi_B(t_2)\},
 \label{eq:V64-ab-switch-region}\\
 \mathcal E_{ac}
 &:=(\text{level rectangle})\setminus\mathcal E_{ab},
 \label{eq:V64-ac-switch-region}
\end{align}
then the upper bound in
\eqref{eq:V64-coherent-insertion} splits into the \(ab\) expression
on \(\mathcal E_{ab}\) and the \(ac\) expression on
\(\mathcal E_{ac}\), always retaining the common scalar cap one.
\end{proposition}

\begin{proof}
Equation \eqref{eq:V64-height-core-level} gives
\[
 \|\operatorname{Tr}_{c}P_K(t_2)\|_{\rm op}
 \le\Phi_C(t_2),
 \qquad
 \|\operatorname{Tr}_{b}P_K(t_2)\|_{\rm op}
 \le\Phi_B(t_2).
\]
Insert these two inequalities in
\eqref{eq:V62-ab-coherent-factor}--%
\eqref{eq:V62-ac-coherent-factor}.  The minimum is monotone in each
nonnegative argument, proving
\eqref{eq:V64-coherent-insertion}.  The two sets
\eqref{eq:V64-ab-switch-region}--%
\eqref{eq:V64-ac-switch-region} merely record which of the two
complete oriented products is smaller.  At no point is a pair ratio,
a core factor, or a coordinate transform minimized or capped
separately.
\end{proof}

\begin{assessment}[The switching is not inside the core transform]
\label{ass:V64-switching-location}
Theorem \ref{thm:V64-local-transform-decay} handles both local
height arrangements with the same denominator \(d_R\):
\begin{enumerate}[label=\textup{(\roman*)},leftmargin=3.7em]
\item if \(Q\) is comparable to or higher than \(R\), reciprocity
      produces the \(d_Q^{-1}\) decay;
\item if \(R\) is much higher than \(Q\), reverse fusion produces
      exponential decay.
\end{enumerate}
Therefore an array alternating between
\emph{high retained intermediate} and \emph{high traced input} needs
no coordinatewise orientation switch.  The unresolved switch in
\eqref{eq:V64-ab-switch-region} is the global choice between retaining
\(ab\) and retaining \(ac\), because that choice simultaneously
changes the \(bc\)-pair denominator and which core row is traced.
\end{assessment}

\begin{assessment}[Dual-singlet regression]
\label{ass:V64-dual-singlet-regression}
Simultaneous dual singlets do not invalidate
\eqref{eq:V64-local-transform-decay}: each singlet contributes exactly
\(d_R^{-1}\), which is an inverse Casimir power on a boundary ray.
They do show why no stronger one-link power is available.  After
layer integration, the growing family of nearby \((k,k)\) outputs in
\Cref{thm:V64-dual-model-capacity}, rather than the single singlet,
causes the square-root core-capacity loss.
\end{assessment}

\section{Upstream consequence and V64 ledger}
\label{sec:V64-ledger}

\begin{theorem}[Integrated V64 statement]
\label{thm:V64-integrated}
Preserving V1--V63, V64 proves:
\begin{enumerate}[label=\textup{(V64.\arabic*)},leftmargin=5.5em]
\item the fixed-orientation local transform decay
      \eqref{eq:V64-local-transform-decay};
\item the global inverse-mass transform product
      \eqref{eq:V64-global-transform-product};
\item the height-sensitive core level and capacity estimates
      \eqref{eq:V64-height-core-level}--%
      \eqref{eq:V64-log-height-capacity};
\item the exact dual-symmetric model law
      \eqref{eq:V64-dual-model-capacity};
\item the resulting no-go statement for a proof using only the V63
      output-weight majorant; and
\item the coherent, uncapped insertion
      \eqref{eq:V64-coherent-insertion}, with an explicit global
      switching partition.
\end{enumerate}
\end{theorem}

\begin{proof}
The six items are respectively
\Cref{thm:V64-local-transform-decay,
cor:V64-global-transform-product,
cor:V64-height-core-level,cor:V64-log-height-capacity,
thm:V64-dual-model-capacity,cor:V64-majorant-no-go,
prop:V64-coherent-insertion}.
\end{proof}

\begingroup
\small
\begin{longtable}{>{\raggedright\arraybackslash}p{0.29\textwidth}
                  >{\raggedright\arraybackslash}p{0.15\textwidth}
                  >{\raggedright\arraybackslash}p{0.48\textwidth}}
\toprule
\textbf{V64 row}&\textbf{Status}&\textbf{Advance or obstruction}\\
\midrule
\endfirsthead
\toprule
\textbf{V64 row}&\textbf{Status}&\textbf{Advance or obstruction}\\
\midrule
\endhead
One-link height transform
& \MGclosed
& One fixed orientation gains an inverse maximum input mass, with an
  additional uniform active-coordinate gap.\\[0.35em]
Alternating local heights
& \MGclosed
& Reciprocity and reverse fusion cover the two height directions
  without coordinatewise orientation choices.\\[0.35em]
Core height capacity
& \MGreduced
& Chernoff gives a logarithmic height gain, complementing the V63
  support-count gain.\\[0.35em]
Core-only inverse-mass route
& \MGblocked
& The V63 output-weight majorant permits an exact dual-symmetric model
  with only \(1/(s_\alpha\sqrt{\Xi})\) capacity.\\[0.35em]
Coherent pair--core product
& \MGopen
& The two pair ratios remain inside the global cap; their correlation
  with the actual cut weight has not yet been used.\\[0.35em]
Complete collision and \([3,2]\)
& \MGopen
& Input-sensitive cut structure or a joint weighted trace theorem is
  still required before puncture pair-label aggregation.\\[0.35em]
Full Yang--Mills mass gap
& \MGopen
& Higher quartic/all-order MTA, WUFC/CPM, nonlinear construction,
  continuum OS existence/nontriviality, and spectral positivity remain
  downstream.\\
\bottomrule
\end{longtable}
\endgroup

\begin{openproblem}[V65 mandate: companion audit and joint weighted trace]
\label{op:V65}
\begin{enumerate}[label=\textup{(V65.\arabic*)},leftmargin=5.5em]
\item Perform the required five-version audit of the current Standard
      Model and Navier--Stokes notes.  Transfer only a structurally
      compatible argument and record exact snapshots.
\item Return upstream from the core-only layer majorant to the direct
      weighted operator of
      \Cref{thm:V62-direct-weighted-tax}.  Keep the \(oa\), \(bc\),
      and actual ternary-cut weights in one positive operator before
      taking partial traces.
\item On the dual-symmetric regression family, compute or sharply
      bound the \emph{actual} cut coefficient
      \(K_T^{\circ,\boldsymbol T}\), rather than replacing it by
      \(C/(1+s_\alpha\Xi(T))\).  Determine whether the missing
      square-root is supplied by cumulant cancellation.
\item If the actual cut does not supply it, seek a joint pair--core
      trace inequality in which the same retained-row orientation pays
      the \(bc\) ratio and the near-singlet core family.
\item Preserve the switching sets
      \eqref{eq:V64-ab-switch-region}--%
      \eqref{eq:V64-ac-switch-region}; do not multiply local minima or
      cap pair ratios separately.
\item Keep puncture pair-label aggregation suspended until the
      three-bank collision core closes.  Keep multi-collision,
      \([3,3]\), and \([4]\) suspended.
\end{enumerate}
\end{openproblem}

\begin{firewall}[End-of-V64 claim boundary]
\label{fire:V64-claim-boundary}
V64 proves mass-weighted conditional transform decay, its global
product form, a height-sensitive core estimate, the sharp
dual-symmetric majorant barrier, and the coherent insertion into the
V62 envelope.  It does not prove an input-sensitive bound for the
actual ternary cut, the joint pair--core weighted trace, the V62
rectangular JREC gate, JREC, simultaneous-sideways closure,
fixed-type puncture aggregation, multi-collision aggregation,
complete \([3,2]\), \([3,3]\), \([4]\), global \(m=4\), all-order
MTA, WUFC, CPM, the nonlinear Perron packet, finite-edge margins,
constructive capacity, protected-sector floors, local-algebra
noncollapse, reflection positivity, a nontrivial continuum OS state,
or a positive Yang--Mills spectral gap.  The full Yang--Mills mass gap
is not proved.
\end{firewall}

\section*{Version V64 append-only record}
\addcontentsline{toc}{section}{Version V64 append-only record}
The complete V63 source was retained before this continuation.  Its
SHA--256 digest was
\begin{center}\scriptsize\ttfamily
e747660da3a41999d27711eb9e0cc4275ac9d6cdb012f1594342405701671746.
\end{center}
No mathematical content of V63 was altered during the V64 append.  On
the next iteration, retain this full file, append before its unique
active document-closing command, increment the filename to
\texttt{\_V65.tex}, and remove only the superseded V64 file from this
Yang--Mills note series.

% =============================================================================
% END APPENDED VERSION V64 MATERIAL
% =============================================================================

% =============================================================================
% BEGIN APPENDED VERSION V65 MATERIAL
% =============================================================================

\chapter{V65: Companion audit, hyperbolic joint levels, and the
missing Schur quantile}
\label{ch:V65}

\section{Mandatory five-version companion audit}
\label{sec:V65-companion-audit}

The V65 direction was selected after inspecting stable snapshots of
the current companion notes:
\begin{center}
\small
\begin{tabular}{@{}p{0.49\textwidth}rp{0.34\textwidth}@{}}
\toprule
\textbf{source}&\textbf{lines}&\textbf{SHA--256}\\
\midrule
\texttt{Renewal\_SM\_Einstein\_Continuum\_Research\_Notes\_V75.tex}
&60741&
\texttt{f2f98dd968b7d2e76967b93fc2a1758379d8617f3ab1e4c1e05dfe55295bc31e}\\
\texttt{Renewal\_NS\_Stability\_Research\_Notes\_V31.tex}
&23052&
\texttt{f1121b62238ad5491bb7cf03635ea9934942edf573ed8faaa9ee79e1d38a6696}\\
\bottomrule
\end{tabular}
\end{center}
Each digest was computed twice on the same file size and agreed.

SM V75 keeps the metric gauge functional, auxiliary \(B\)-field,
ghost Green term, face transgressions, and corner phase in one
off-shell complex.  Only after the complete conormal graph is formed
does it eliminate \(B\).  Its exact lesson for the present programme
is an order of reduction: eliminating an auxiliary block before
forming the joint boundary object can turn an off-shell identity into
an apparently conditional one.

NS V31 likewise retains the correlation among recent-entry times,
signed transfer, and the active high-parent tail.  Separate norm
bounds recover only a continuation-class stock, whereas the actual
stopping-time-correlated occupation is the open quantity.  Its dyadic
coarea identities do not give an SU(3) trace estimate.

\begin{assessment}[V65 audit transfer]
\label{ass:V65-audit-transfer}
No analytic inequality is imported from SM or NS.  Their compatible
assembly rule changes the YM acquisition order: return from the
rectangular majorants to the single V60 product weight
\(W_\Lambda\), retain its product threshold as one hyperbolic level
set, and only then take the two possible row partial traces.  Replacing
the actual cut coefficient by its output-only majorant before this
step would discard exactly the correlation that V64 identified as
potentially necessary.
\end{assessment}

\begin{firewall}[Companion separation]
\label{fire:V65-companion-separation}
The SM off-shell quartet and the NS stopping-time coarea are not
Yang--Mills representation estimates.  V65 transfers only their
proved construction order.  No Standard Model continuum statement or
Navier--Stokes regularity statement is used below.
\end{firewall}

\section{The exact hyperbolic product-level projector}
\label{sec:V65-hyperbolic-projector}

Work first in the collision-core sector without the still-suspended
puncture pair-label banks.  Use the Cartesian pre-final family proved
in V62, with orthogonal bank projectors
\[
 P_{oa,\lambda},\qquad P_{bc,\mu},\qquad P_{K,\nu},
\]
and nonnegative weights
\[
 a_\lambda:=\mathfrak a_{oa,\lambda},\qquad
 b_\mu:=\mathfrak a_{bc,\mu},\qquad
 k_\nu:=\mathfrak k_\nu.
\]
The original V60 weight and projector are
\begin{equation}
 w_{\lambda\mu\nu}:=a_\lambda b_\mu k_\nu,
 \qquad
 P_{\lambda\mu\nu}
 :=P_{oa,\lambda}\otimes P_{bc,\mu}\otimes P_{K,\nu}.
 \label{eq:V65-product-weight-projector}
\end{equation}
For \(0<t<W_*\), define the single product-level projector
\begin{equation}
 P_{\rm hyp}(t)
 :=
 \sum_{\{(\lambda,\mu,\nu):
 a_\lambda b_\mu k_\nu>t\}}
 P_{\lambda\mu\nu}.
 \label{eq:V65-hyperbolic-projector}
\end{equation}
The adjective \emph{hyperbolic} refers only to the product threshold;
it does not introduce a new estimate.

For the \(ab\) retained-row orientation, Schur's lemma on the two pair
banks gives nonnegative scalars
\begin{align}
 \operatorname{Tr}_{o}P_{oa,\lambda}
 &=x_\lambda^a I_a,
 &
 \operatorname{Tr}_{c}P_{bc,\mu}
 &=y_\mu^b I_b.
 \label{eq:V65-pair-partial-scalars-ab}
\end{align}
On the core, V63 gives a common sequential decomposition on which
\begin{equation}
 \operatorname{Tr}_{C}P_{K,\nu}
 =
 \bigoplus_{S,\alpha}
 z_\nu^C(S)I_{H_{S,\alpha}},
 \qquad
 z_\nu^C(S)\ge0.
 \label{eq:V65-core-partial-scalars-ab}
\end{equation}
For a complete final irrep \(T_\nu\),
\[
 z_\nu^C(S)=
 \frac{N_{SC}^{T_\nu}d_{T_\nu}}{d_S}.
\]
Define the corresponding \(ac\)-orientation scalars
\(y_\mu^c\) and \(z_\nu^B(S')\) by tracing row \(b\).

\begin{theorem}[Exact scalarization of the hyperbolic joint level]
\label{thm:V65-hyperbolic-scalarization}
Put
\begin{align}
 F_{ab}(t;S)
 &:=
 \sum_{\{a_\lambda b_\mu k_\nu>t\}}
 x_\lambda^a y_\mu^b z_\nu^C(S),
 \label{eq:V65-Fab}\\
 F_{ac}(t;S')
 &:=
 \sum_{\{a_\lambda b_\mu k_\nu>t\}}
 x_\lambda^a y_\mu^c z_\nu^B(S').
 \label{eq:V65-Fac}
\end{align}
Then
\begin{align}
 \left\|\operatorname{Tr}_{\{a,b\}^c}
 P_{\rm hyp}(t)\right\|_{\rm op}
 &=\max_S F_{ab}(t;S),
 \label{eq:V65-hyperbolic-ab-norm}\\
 \left\|\operatorname{Tr}_{\{a,c\}^c}
 P_{\rm hyp}(t)\right\|_{\rm op}
 &=\max_{S'} F_{ac}(t;S').
 \label{eq:V65-hyperbolic-ac-norm}
\end{align}
Consequently the original V61 collision-core capacity is exactly
\begin{equation}
 \boxed{
 \mathcal K_{\rm hyp}
 =
 \int_0^{W_*}
 \min\left\{
 1,\max_SF_{ab}(t;S),\max_{S'}F_{ac}(t;S')
 \right\}\,\dd t.}
 \label{eq:V65-hyperbolic-capacity}
\end{equation}
\end{theorem}

\begin{proof}
Partial trace is linear and factorizes on every tensor product in
\eqref{eq:V65-product-weight-projector}.  Equations
\eqref{eq:V65-pair-partial-scalars-ab} and
\eqref{eq:V65-core-partial-scalars-ab} show that the \(ab\) partial
trace of \eqref{eq:V65-hyperbolic-projector} is block diagonal in
\((S,\alpha)\), with eigenvalue \(F_{ab}(t;S)\).  Positivity makes its
operator norm the largest eigenvalue.  This proves
\eqref{eq:V65-hyperbolic-ab-norm}; the other orientation is identical.
Finally insert both identities into
\eqref{eq:V61-two-row-envelope}.  The two inner scalar caps in
\eqref{eq:V61-two-row-tax} combine with the orientation minimum to
give exactly \eqref{eq:V65-hyperbolic-capacity}.
\end{proof}

\begin{corollary}[Hyperbolic scalar gate for collision-core JREC]
\label{cor:V65-hyperbolic-JREC-gate}
If
\begin{equation}
 \mathcal K_{\rm hyp}
 \le
 Cs_\alpha\sum_{r=b,c}
 \frac{h_{oa}\widehat H_{ar}h_{bc}}{\Xi_a\Xi_r},
 \label{eq:V65-hyperbolic-JREC-gate}
\end{equation}
then JREC holds in the collision-core sector.
\end{corollary}

\begin{proof}
By \Cref{thm:V65-hyperbolic-scalarization},
\(\mathcal K_{\rm hyp}\) is the V61 endpoint-pair envelope for the
original V60 product weight.  Apply
\Cref{thm:V61-capacity-implies-JREC}.
\end{proof}

\section{What is lost after the product threshold}
\label{sec:V65-product-threshold}

\begin{lemma}[Uncapped first moment factorizes only after level integration]
\label{lem:V65-uncapped-first-moment}
For each fixed core intermediate \(S\),
\begin{align}
 \int_0^{W_*}F_{ab}(t;S)\,\dd t
 &=
 \left(\sum_\lambda a_\lambda x_\lambda^a\right)
 \left(\sum_\mu b_\mu y_\mu^b\right)
 \left(\sum_\nu k_\nu z_\nu^C(S)\right),
 \label{eq:V65-uncapped-ab-factorization}\\
 \int_0^{W_*}F_{ac}(t;S')\,\dd t
 &=
 \left(\sum_\lambda a_\lambda x_\lambda^a\right)
 \left(\sum_\mu b_\mu y_\mu^c\right)
 \left(\sum_\nu k_\nu z_\nu^B(S')\right).
 \label{eq:V65-uncapped-ac-factorization}
\end{align}
\end{lemma}

\begin{proof}
For every triple,
\[
 \int_0^{W_*}
 \mathbf1_{\{a_\lambda b_\mu k_\nu>t\}}\,\dd t
 =a_\lambda b_\mu k_\nu.
\]
Insert this identity in \eqref{eq:V65-Fab} or
\eqref{eq:V65-Fac}, use Tonelli, and factor the unrestricted Cartesian
sum.
\end{proof}

\begin{assessment}[Why the exact hyperbolic level is the upstream object]
\label{ass:V65-hyperbolic-upstream}
The factorization
\eqref{eq:V65-uncapped-ab-factorization} occurs only after removing
the scalar cap and fixing the intermediate block.  The JREC capacity
\eqref{eq:V65-hyperbolic-capacity} instead takes, in order,
\[
 \text{product threshold}\ \longrightarrow\
 \text{block maxima}\ \longrightarrow\
 \text{one common scalar cap}\ \longrightarrow\
 \text{level integration}.
\]
Changing that order returns the separate pair and core majorants which
V62--V64 showed to be insufficient.  Thus
\eqref{eq:V65-hyperbolic-capacity}, not either factorized first
moment in isolation, is the next representation-theoretic object.
\end{assessment}

\begin{lemma}[Monotonicity under premature cut majorization]
\label{lem:V65-cut-majorant-monotonicity}
If \(0\le k_\nu\le\widetilde k_\nu\) for every core output, then the
capacity formed from \(\widetilde k\) is at least the capacity formed
from \(k\).
\end{lemma}

\begin{proof}
For every \(t\),
\[
 \{a_\lambda b_\mu k_\nu>t\}
 \subseteq
 \{a_\lambda b_\mu\widetilde k_\nu>t\}.
\]
All summands in \eqref{eq:V65-Fab}--\eqref{eq:V65-Fac} are positive.
Both block maxima and the common minimum with one are monotone in
those sums.  Integrate.
\end{proof}

\begin{remark}[The V64 barrier is therefore one-sided]
\label{rem:V65-barrier-one-sided}
Replacing the actual \(k_\nu\) by
\(C/(1+s_\alpha\Xi(T_\nu))\) can only enlarge the hyperbolic
capacity.  V64 proves that this enlarged object cannot have the
desired core-only inverse-mass rate.  It does not determine the
capacity for the smaller actual weights.
\end{remark}

\section{The exact Schur-quantile gate}
\label{sec:V65-Schur-quantile}

The dual-symmetric family of V64 admits a sharper formulation which
specifies exactly what an actual cut calculation must gain.  For
arbitrary weights \(0\le v_k\le V_*\), \(0\le k\le n\), define
\begin{align}
 F_{n,v}(t)
 &:=
 \frac1{d_{R_n}}
 \sum_{\{k:v_k>t\}}(k+1)^3,
 \label{eq:V65-dual-general-level}\\
 \mathcal L_n(v)
 &:=
 \int_0^{V_*}\min\{1,F_{n,v}(t)\}\,\dd t,
 \label{eq:V65-dual-general-capacity}\\
 q_n(v)
 &:=
 \sup\{t\in[0,V_*]:F_{n,v}(t)\ge1\},
 \label{eq:V65-Schur-quantile}
\end{align}
with the supremum of the empty set declared to be zero.

\begin{theorem}[Unit one-sided Schur quantile]
\label{thm:V65-Schur-quantile}
For every nonnegative weight family \(v\),
\begin{equation}
 \boxed{\mathcal L_n(v)\ge q_n(v).}
 \label{eq:V65-capacity-dominates-quantile}
\end{equation}
For the V64 model weights
\[
 u_0=1,\qquad
 u_k=(1+s(k+1)^2)^{-1}\quad(k\ge1),
\]
one has
\begin{equation}
 q_n(u)\asymp\frac1{1+sn}
 \asymp\min\left\{1,\frac1{sn}\right\}.
 \label{eq:V65-model-quantile}
\end{equation}
Thus the square-root loss in
\eqref{eq:V64-dual-model-capacity} is already present at the level
where one unit of normalized one-sided Schur dimension becomes
active.
\end{theorem}

\begin{proof}
The function \(F_{n,v}\) is nonincreasing.  Hence
\(F_{n,v}(t)\ge1\) for almost every \(0<t<q_n(v)\), and the integrand
in \eqref{eq:V65-dual-general-capacity} equals one there.  This proves
\eqref{eq:V65-capacity-dominates-quantile}.

By \eqref{eq:V64-cubic-partial-sum} and
\(d_{R_n}\asymp n^2\), there are fixed \(0<c_0<C_0<\infty\) such
that
\[
 \frac1{d_{R_n}}\sum_{k=0}^{\lfloor c_0\sqrt n\rfloor}(k+1)^3<1,
 \qquad
 \frac1{d_{R_n}}\sum_{k=0}^{\lfloor C_0\sqrt n\rfloor}(k+1)^3>1
\]
for all sufficiently large \(n\); adjust constants for the finite
remainder.  Since \(u_k\) decreases with \(k\ge1\), the threshold at
which \(F_{n,u}\) crosses one lies between fixed multiples of
\[
 \frac1{1+s(\sqrt n+1)^2}.
\]
This is comparable to \((1+sn)^{-1}\), proving
\eqref{eq:V65-model-quantile}.
\end{proof}

\begin{corollary}[Exact size of a scalar companion debit]
\label{cor:V65-scalar-companion-debit}
For \(p\ge0\),
\begin{equation}
 \mathcal L_n(pv)=p\,\mathcal L_n(v),
 \qquad
 q_n(pv)=p\,q_n(v),
 \label{eq:V65-scalar-homogeneity}
\end{equation}
where the integration range is scaled from \(V_*\) to \(pV_*\).
In the regime \(sn\to\infty\), a scalar companion multiplying the V64
model weights can lower their capacity to
\(C/(sn^2)\) only if
\begin{equation}
 p\le\frac Cn
 \asymp\frac C{\sqrt{\xi(Q_n)}}.
 \label{eq:V65-required-companion-debit}
\end{equation}
\end{corollary}

\begin{proof}
Substitute \(t=p\tau\) in
\eqref{eq:V65-dual-general-capacity}; the quantile identity is
immediate from the same change of variables.  Equations
\eqref{eq:V65-model-quantile} and
\eqref{eq:V64-dual-model-capacity} give capacity
\(\asymp p/(sn)\) when \(sn\) is large.  Comparing it with
\(C/(sn^2)\) proves \eqref{eq:V65-required-companion-debit}.
\end{proof}

\begin{proposition}[Simultaneous-orientation saturation quantile]
\label{prop:V65-both-orientation-quantile}
For the exact hyperbolic functions of
\eqref{eq:V65-Fab}--\eqref{eq:V65-Fac}, define
\begin{equation}
 q_{\rm both}
 :=
 \sup\left\{t\in[0,W_*]:
 \max_SF_{ab}(t;S)\ge1
 \ \hbox{and}\
 \max_{S'}F_{ac}(t;S')\ge1
 \right\}.
 \label{eq:V65-both-orientation-quantile}
\end{equation}
Then
\begin{equation}
 \boxed{\mathcal K_{\rm hyp}\ge q_{\rm both}.}
 \label{eq:V65-hyperbolic-dominates-quantile}
\end{equation}
Consequently \eqref{eq:V65-hyperbolic-JREC-gate} requires
\begin{equation}
 q_{\rm both}
 \le
 Cs_\alpha\sum_{r=b,c}
 \frac{h_{oa}\widehat H_{ar}h_{bc}}{\Xi_a\Xi_r}.
 \label{eq:V65-necessary-quantile-gate}
\end{equation}
\end{proposition}

\begin{proof}
Both maxima are nonincreasing in \(t\).  For almost every
\(0<t<q_{\rm both}\), both are at least one, so the integrand in
\eqref{eq:V65-hyperbolic-capacity} equals one.  This proves
\eqref{eq:V65-hyperbolic-dominates-quantile}; the necessary condition
follows from the sufficient target size in
\eqref{eq:V65-hyperbolic-JREC-gate}.
\end{proof}

\begin{assessment}[Concrete next calculation]
\label{ass:V65-concrete-next-calculation}
The V64 output-only model has
\(q_n\asymp(sn)^{-1}\), whereas an inverse-Casimir JREC scale is
\((sn^2)^{-1}\).  Therefore the actual collision calculation must
supply one of two explicitly testable effects on the first
\(O(\sqrt n)\) dual-symmetric output shells:
\begin{enumerate}[label=\textup{(\roman*)},leftmargin=3.7em]
\item the ternary cumulant makes the actual cut-weight Schur quantile
      smaller by \(n^{-1}\); or
\item the product threshold with the pair weights supplies a coherent
      companion debit of that size before both row orientations
      saturate.
\end{enumerate}
This is a quantitative alternative, not a placeholder capacity name.
\end{assessment}

\section{V65 ledger and next acquisition}
\label{sec:V65-ledger}

\begin{theorem}[Integrated V65 statement]
\label{thm:V65-integrated}
Preserving V1--V64, V65 proves:
\begin{enumerate}[label=\textup{(V65.\arabic*)},leftmargin=5.5em]
\item the digest-recorded SM V75 and NS V31 companion audit with no
      analytic transfer;
\item the exact hyperbolic product-level scalarization
      \eqref{eq:V65-hyperbolic-ab-norm}--%
      \eqref{eq:V65-hyperbolic-capacity};
\item the scalar sufficient gate
      \eqref{eq:V65-hyperbolic-JREC-gate} for collision-core JREC;
\item the exact uncapped first-moment factorization
      \eqref{eq:V65-uncapped-ab-factorization}--%
      \eqref{eq:V65-uncapped-ac-factorization};
\item the unit Schur-quantile lower bound and the exact model scale
      \eqref{eq:V65-model-quantile}; and
\item the necessary inverse-square-root companion debit
      \eqref{eq:V65-required-companion-debit} and simultaneous-
      orientation quantile gate
      \eqref{eq:V65-necessary-quantile-gate}.
\end{enumerate}
\end{theorem}

\begin{proof}
The audit is
\Cref{ass:V65-audit-transfer,fire:V65-companion-separation}.
The remaining items are
\Cref{thm:V65-hyperbolic-scalarization,
cor:V65-hyperbolic-JREC-gate,
lem:V65-uncapped-first-moment,
thm:V65-Schur-quantile,
cor:V65-scalar-companion-debit,
prop:V65-both-orientation-quantile}.
\end{proof}

\begingroup
\small
\begin{longtable}{>{\raggedright\arraybackslash}p{0.29\textwidth}
                  >{\raggedright\arraybackslash}p{0.15\textwidth}
                  >{\raggedright\arraybackslash}p{0.48\textwidth}}
\toprule
\textbf{V65 row}&\textbf{Status}&\textbf{Advance or obstruction}\\
\midrule
\endfirsthead
\toprule
\textbf{V65 row}&\textbf{Status}&\textbf{Advance or obstruction}\\
\midrule
\endhead
Companion audit
& \MGclosed
& Current SM V75 and NS V31 reinforce one joint-assembly order but
  supply no YM inequality.\\[0.35em]
Hyperbolic product level
& \MGclosed
& The original V60 product threshold now has an exact scalar formula
  in each complete intermediate block.\\[0.35em]
Schur quantile
& \MGclosed
& The dual-symmetric majorant loses precisely one square root at the
  unit one-sided-dimension threshold.\\[0.35em]
Actual cut quantile
& \MGopen
& The first \(O(\sqrt n)\) shells of the unmajorized ternary cumulant
  have not yet been calculated.\\[0.35em]
Coherent pair companion
& \MGopen
& If the cut does not gain \(n^{-1}\), the product threshold must gain
  it before both row orientations saturate.\\[0.35em]
Complete collision and \([3,2]\)
& \MGopen
& JREC and fixed-type puncture pair-label aggregation remain.\\[0.35em]
Full Yang--Mills mass gap
& \MGopen
& Higher quartic/all-order MTA, WUFC/CPM, nonlinear construction,
  continuum OS existence/nontriviality, and spectral positivity remain
  downstream.\\
\bottomrule
\end{longtable}
\endgroup

\begin{openproblem}[V66 mandate: unmajorized cut quantile]
\label{op:V66}
\begin{enumerate}[label=\textup{(V66.\arabic*)},leftmargin=5.5em]
\item Return to the exact five-term ternary cumulant identity
      \eqref{eq:V57-cut-identity}; do not begin from
      \eqref{eq:V63-cut-weight-decay}.
\item On the one-coordinate dual-symmetric regression, resolve the
      actual operator \(K_T^{\circ,T}\) on the outputs
      \(T=(k,k)\), first for \(0\le k\le c\sqrt n\).
\item Keep the five signed cumulant terms together through the common
      recoupling basis.  Determine whether their leading low-shell
      symbols cancel before the operator norm is taken.
\item Express the result as an upper bound for the actual quantile
      \(q_n((\mathfrak k_k)_k)\).  The required gain is the explicit
      factor \(n^{-1}\) in
      \eqref{eq:V65-required-companion-debit}.
\item If that gain is false, retain the actual pair weights in
      \eqref{eq:V65-hyperbolic-projector} and estimate
      \(q_{\rm both}\) directly.  Do not replace the product threshold
      by three independent level variables.
\item Add puncture pair-label banks only after the collision-core
      quantile closes.  Keep multi-collision, \([3,3]\), and \([4]\)
      suspended.  The next companion audit is V70.
\end{enumerate}
\end{openproblem}

\begin{firewall}[End-of-V65 claim boundary]
\label{fire:V65-claim-boundary}
V65 performs the required companion audit and proves the exact
hyperbolic product-level scalarization, Schur-quantile theorem, and
necessary missing-debit scales.  It does not estimate the actual
ternary-cumulant quantile, the coherent pair companion, the
hyperbolic JREC gate, JREC, simultaneous-sideways closure, fixed-type
puncture aggregation, multi-collision aggregation, complete
\([3,2]\), \([3,3]\), \([4]\), global \(m=4\), all-order MTA, WUFC,
CPM, the nonlinear Perron packet, finite-edge margins, constructive
capacity, protected-sector floors, local-algebra noncollapse,
reflection positivity, a nontrivial continuum OS state, or a
positive Yang--Mills spectral gap.  The full Yang--Mills mass gap is
not proved.
\end{firewall}

\section*{Version V65 append-only record}
\addcontentsline{toc}{section}{Version V65 append-only record}
The complete V64 source was retained before this continuation.  Its
SHA--256 digest was
\begin{center}\scriptsize\ttfamily
149442e7df891c80f502043eee673d696f751aa26d1d6f4e8f09daa25c776dc8.
\end{center}
No mathematical content of V64 was altered during the V65 append.  On
the next iteration, retain this full file, append before its unique
active document-closing command, increment the filename to
\texttt{\_V66.tex}, and remove only the superseded V65 file from this
Yang--Mills note series.

% =============================================================================
% END APPENDED VERSION V65 MATERIAL
% =============================================================================

% =============================================================================
% BEGIN APPENDED VERSION V66 MATERIAL
% =============================================================================

\chapter{V66: Exact singlet cancellation and localization of the
mesoscopic core shell}
\label{ch:V66}

\section{Purpose and nonduplication boundary}
\label{sec:V66-purpose}

The first V65 instruction was to calculate the actual ternary cut on
a one-coordinate dual-symmetric regression.  The complete
one-coordinate arbitrary-coefficient junction is already proved in
\Cref{thm:V46-corrected-one-link-MTA}; repeating that theorem would
not address JREC.  Its own firewall identifies the real issue: remote
banks can make the original row masses large after the local junction
has already been normalized.

V66 therefore extracts only the new information needed by the V65
Schur-quantile gate.  It computes the actual final-singlet cumulant
before norms, proves that every fixed cold-output window already has
an inverse input-mass tax, and shows that the V64 square-root loss is
created at an output Casimir of geometric-mean size.  The remaining
object is a growing mesoscopic Racah shell, not the singlet and not a
bounded collection of cold representations.

\begin{firewall}[Scope of V66]
\label{fire:V66-scope}
The exact singlet formula below applies when the final invariant
multiplicity is one.  The fixed-cold-window estimate is uniform in
multiplicity.  Neither statement controls the growing mesoscopic
shell, so neither proves the hyperbolic JREC gate.
\end{firewall}

\section{The actual multiplicity-one singlet cumulant}
\label{sec:V66-singlet-cumulant}

Let \(A,B,C\) be irreducible SU(3) representations with
\[
 \dim\operatorname{Inv}(H_A\otimes H_B\otimes H_C)=1.
\]
Write
\[
 x=q_{\alpha,A},\qquad
 y=q_{\alpha,B},\qquad
 z=q_{\alpha,C}.
\]
Duality invariance gives
\(q_{\alpha,U^*}=q_{\alpha,U}\).

\begin{theorem}[Exact Gram-determinant form at the final singlet]
\label{thm:V66-singlet-factorization}
On the one-dimensional final-singlet multiplicity space, the exact
ternary cumulant is the scalar
\begin{align}
 K_{\mathbf1}
 &=1-x^2-y^2-z^2+2xyz
 \label{eq:V66-singlet-polynomial}\\
 &=(1-x^2)(1-y^2)-(z-xy)^2
 \label{eq:V66-singlet-factorization}\\
 &=(1-y^2)(1-z^2)-(x-yz)^2
 \notag\\
 &=(1-z^2)(1-x^2)-(y-zx)^2.
 \notag
\end{align}
In particular,
\begin{equation}
 |K_{\mathbf1}|
 \le
 \min_{\rm cyclic}
 \left\{
 (1-q_{\alpha,A}^2)(1-q_{\alpha,B}^2)
 +|q_{\alpha,C}-q_{\alpha,A}q_{\alpha,B}|^2
 \right\}.
 \label{eq:V66-singlet-two-defect-bound}
\end{equation}
\end{theorem}

\begin{proof}
On a final singlet, the \(AB\) intermediate which can couple to \(C\)
is \(C^*\).  The multiplicity-one hypothesis makes the corresponding
pair-moment operator scalar, with value \(q_{\alpha,C^*}=z\).
Therefore the \(AB|C\) partition contributes \(z^2\).  The other two
pair partitions contribute \(x^2\) and \(y^2\), the all-in-one moment
is \(q_{\alpha,\mathbf1}=1\), and the three-singleton partition is
\(xyz\).  Insert these four facts in the five-term common-basis
symbol \eqref{eq:V47-common-cumulant}; this gives
\eqref{eq:V66-singlet-polynomial}.

Expanding
\((1-x^2)(1-y^2)-(z-xy)^2\) gives the same polynomial.
The other two factorizations follow cyclically.  Taking absolute
values in any one factorization and then the minimum proves
\eqref{eq:V66-singlet-two-defect-bound}.
\end{proof}

\begin{corollary}[Quantitative original-pair form]
\label{cor:V66-singlet-original-pair}
Let
\[
 g_U:=\min\{s_\alpha C_2(U),1\},
 \qquad
 H_{AB}:=A_AA_B.
\]
Uniformly for sufficiently large \(\alpha\),
\begin{equation}
 |K_{\mathbf1}|
 \le
 C\min_{\rm cyclic}
 \left\{
 g_Ag_B+\min\{s_\alpha H_{AB},1\}^2
 \right\}.
 \label{eq:V66-singlet-original-pair}
\end{equation}
If one input is trivial, then \(K_{\mathbf1}=0\) exactly.
\end{corollary}

\begin{proof}
The two-sided character gap
\eqref{eq:V30-two-sided-product-gap} gives
\[
 1-q_{\alpha,U}^2
 \le2(1-q_{\alpha,U})\le Cg_U.
\]
The second term in
\eqref{eq:V66-singlet-two-defect-bound} is the square of the exact
binary defect for \(C^*\subset A\otimes B\).  Apply
\eqref{eq:V46-defect-cap} and the universal bound one, then minimize
cyclically.  If, for example, \(B=\mathbf1\), invariance forces
\(C=A^*\), so \(y=1\) and \(z=x\); either
\eqref{eq:V66-singlet-polynomial} or
\eqref{eq:V66-singlet-factorization} is then zero.
\end{proof}

\begin{assessment}[Why the singlet is not the V64 barrier]
\label{ass:V66-singlet-not-barrier}
Even without \eqref{eq:V66-singlet-original-pair}, the complete
singlet block in \(S\otimes C\) contributes
\[
 \frac{N_{SC}^{\mathbf1}}{d_S}
 =
 \begin{cases}
 d_S^{-1},&C=S^*,\\
 0,&C\ne S^*
 \end{cases}
\]
to the V63 one-sided partial trace.  By
\eqref{eq:V34-dimension-Casimir-bounds}, this is already
\(O((1+\xi(C))^{-1})\).  The exact signed formula supplies additional
two-defect structure, but the single singlet was already paid.  The
V64 square-root loss comes from a growing family of nearby outputs.
\end{assessment}

\section{Every fixed cold-output window is paid}
\label{sec:V66-fixed-cold-window}

\begin{theorem}[Hard cold-window one-sided tax]
\label{thm:V66-hard-cold-window}
For every fixed \(L<\infty\), there is \(C_L<\infty\) such that all
ordered pairs \((R,Q)\ne(\mathbf1,\mathbf1)\) satisfy
\begin{equation}
 \boxed{
 \frac1{d_R}
 \sum_{\substack{T\\\xi(T)\le L}}
 N_{RQ}^{T}d_T
 \le
 \frac{C_L}{1+\max\{\xi(R),\xi(Q)\}}.}
 \label{eq:V66-hard-cold-window}
\end{equation}
For a product-group family with total output mass
\(\Xi(\boldsymbol T)\le L\), one likewise has
\begin{equation}
 \frac1{d_{\boldsymbol R}}
 \sum_{\substack{\boldsymbol T\\
                  \Xi(\boldsymbol T)\le L}}
 N_{\boldsymbol R\boldsymbol Q}^{\boldsymbol T}
 d_{\boldsymbol T}
 \le
 \frac{C_L}
 {1+\sum_e\max\{\xi(R_e),\xi(Q_e)\}}.
 \label{eq:V66-product-hard-cold-window}
\end{equation}
\end{theorem}

\begin{proof}
On the set \(\xi(T)\le L\),
\[
 1\le e^{\theta_*L}e^{-\theta_*\xi(T)}.
\]
Insert this inequality in the left side of
\eqref{eq:V66-hard-cold-window} and apply
\eqref{eq:V64-local-transform-decay}.  This proves the one-link
statement.  For the product group, use
\[
 \mathbf1_{\{\Xi(\boldsymbol T)\le L\}}
 \le e^{\theta_*L}e^{-\theta_*\Xi(\boldsymbol T)}
\]
and then \eqref{eq:V64-global-transform-product}.
\end{proof}

\begin{corollary}[Bounded cold windows have inverse-mass layer capacity]
\label{cor:V66-fixed-window-capacity}
Let \(0\le v_{\boldsymbol T}\le V_*\), and restrict a core level
family to \(\Xi(\boldsymbol T)\le L\).  Then
\begin{align}
 &\int_0^{V_*}
 \min\left\{1,
 \left\|\operatorname{Tr}_{\boldsymbol Q}
 \sum_{\substack{\Xi(\boldsymbol T)\le L\\
                  v_{\boldsymbol T}>t}}
 P_{\boldsymbol T}\right\|_{\rm op}
 \right\}\,\dd t
 \notag\\
 &\qquad\le
 \frac{C_LV_*}
 {1+\sum_e\max\{\xi(R_e),\xi(Q_e)\}}.
 \label{eq:V66-fixed-window-capacity}
\end{align}
\end{corollary}

\begin{proof}
Drop the scalar cap, use the exact V63 partial-trace formula, and
apply Tonelli:
\[
 \int_0^{V_*}
 \frac1{d_{\boldsymbol R}}
 \sum_{\substack{\Xi(\boldsymbol T)\le L\\
                  v_{\boldsymbol T}>t}}
 N_{\boldsymbol R\boldsymbol Q}^{\boldsymbol T}
 d_{\boldsymbol T}\,\dd t
 =
 \frac1{d_{\boldsymbol R}}
 \sum_{\Xi(\boldsymbol T)\le L}
 N_{\boldsymbol R\boldsymbol Q}^{\boldsymbol T}
 d_{\boldsymbol T}v_{\boldsymbol T}.
\]
Bound \(v_{\boldsymbol T}\le V_*\) and use
\eqref{eq:V66-product-hard-cold-window}.  If
\(\boldsymbol R\) ranges over intermediate blocks while
\(\boldsymbol Q\) is fixed, use
\(\sum_e\max\{\xi(R_e),\xi(Q_e)\}\ge\Xi(\boldsymbol Q)\);
the resulting uniform bound has denominator
\(1+\Xi(\boldsymbol Q)\).
\end{proof}

\section{The square-root loss is a mesoscopic Racah shell}
\label{sec:V66-mesoscopic-shell}

Return to
\(R_n=(n,0)\), \(Q_n=(0,n)\), and put
\begin{equation}
 D_n(K):=
 \frac1{d_{R_n}}\sum_{k=0}^{K}(k+1)^3,
 \qquad 0\le K\le n.
 \label{eq:V66-dual-cumulative-Schur-mass}
\end{equation}

\begin{theorem}[Geometric-mean location of unit Schur mass]
\label{thm:V66-mesoscopic-location}
There are numerical constants
\(0<c_-<c_+<\infty\) such that, for all sufficiently large \(n\),
\begin{align}
 D_n(\lfloor c_-\sqrt n\rfloor)&\le\frac12,
 \label{eq:V66-below-unit-shell}\\
 D_n(\lfloor c_+\sqrt n\rfloor)&\ge2.
 \label{eq:V66-above-unit-shell}
\end{align}
Every output in the transition band
\[
 c_-\sqrt n\le k\le c_+\sqrt n
\]
has
\begin{equation}
 \xi((k,k))=(k+1)^2\asymp n
 \asymp\sqrt{\xi(R_n)\xi(Q_n)}.
 \label{eq:V66-geometric-mean-output}
\end{equation}
For the V64 model weights, the unit Schur quantile is therefore the
output-moment scale
\begin{equation}
 \frac1{1+s\,\xi((k,k))}
 \asymp\frac1{1+sn}
 \label{eq:V66-mesoscopic-model-weight}
\end{equation}
on this band.
\end{theorem}

\begin{proof}
Equation \eqref{eq:V64-cubic-partial-sum} and
\(d_{R_n}\asymp n^2\) give
\[
 D_n(K)\asymp\frac{(K+1)^4}{n^2}.
\]
Choose \(c_-\) sufficiently small and \(c_+\) sufficiently large to
obtain
\eqref{eq:V66-below-unit-shell}--%
\eqref{eq:V66-above-unit-shell}.  The exact SU(3) Casimir formula
gives
\(\xi((k,k))=(k+1)^2\) for \(k\ge1\), while
\(\xi(R_n),\xi(Q_n)\asymp n^2\).  This proves
\eqref{eq:V66-geometric-mean-output}.  Substitution in the model
weight proves \eqref{eq:V66-mesoscopic-model-weight}.
\end{proof}

\begin{corollary}[Fixed windows and the singlet cannot carry the loss]
\label{cor:V66-only-mesoscopic-loss}
For every fixed \(L\), outputs with \(\xi(T)\le L\) have the
inverse-input-mass capacity
\eqref{eq:V66-fixed-window-capacity}.  In the dual-symmetric model,
their normalized one-sided Schur mass is \(O_L(n^{-2})\).  The unit
quantile and its \(1/(sn)\) scale appear only after the window grows
to output Casimir \(\asymp n\).
\end{corollary}

\begin{proof}
The first assertion is
\Cref{cor:V66-fixed-window-capacity}.  In the explicit decomposition,
a fixed Casimir window contains only finitely many \(k\), and
\[
 \frac1{d_{R_n}}
 \sum_{\xi((k,k))\le L}(k+1)^3
 \le\frac{C_L}{n^2}.
\]
The last assertion is
\Cref{thm:V66-mesoscopic-location}.
\end{proof}

\begin{assessment}[The first genuinely noncommutative regression]
\label{ass:V66-mesoscopic-Racah}
At the final singlet, all three pair intermediates are forced and the
five-term cumulant collapses to the scalar identity
\eqref{eq:V66-singlet-factorization}.  On the mesoscopic family
\(T=(k,k)\), \(k\asymp\sqrt n\), the multiplicity space can carry
different \(AB\), \(AC\), and \(BC\) bracketings.  The actual symbol
has the form
\[
 q_{\alpha,T}I
 -q_{\alpha,C}Q_{AB}^{T}
 -q_{\alpha,B}Q_{AC}^{T}
 -q_{\alpha,A}Q_{BC}^{T}
 +2q_{\alpha,A}q_{\alpha,B}q_{\alpha,C}I,
\]
and the three pair-moment operators need not commute.  Taking their
norms separately destroys the possible cancellation.  Moreover the
conditional data \((S,C,T)\) do not determine the last two operators:
their Racah matrices retain the ancestry \(A\otimes B\to S\).
Thus the next calculation must keep the ancestry label and the common
multiplicity space through the signed sum.
\end{assessment}

\begin{lemma}[Exact Schur quantile--excess identity]
\label{lem:V66-quantile-excess}
For the arbitrary dual-symmetric weights of
\eqref{eq:V65-dual-general-level}--%
\eqref{eq:V65-Schur-quantile},
\begin{equation}
 \boxed{
 \mathcal L_n(v)
 =
 q_n(v)
 +\frac1{d_{R_n}}\sum_{k=0}^{n}
 (k+1)^3\bigl(v_k-q_n(v)\bigr)_+.}
 \label{eq:V66-quantile-excess}
\end{equation}
\end{lemma}

\begin{proof}
Write \(q=q_n(v)\).  Monotonicity of \(F_{n,v}\) gives
\(\min\{1,F_{n,v}(t)\}=1\) for almost every \(0<t<q\) and
\(\min\{1,F_{n,v}(t)\}=F_{n,v}(t)\) for almost every \(t>q\).
Therefore
\[
 \mathcal L_n(v)
 =q+\int_q^{V_*}F_{n,v}(t)\,\dd t.
\]
Insert \eqref{eq:V65-dual-general-level}, apply Tonelli, and use
\[
 \int_q^{V_*}\mathbf1_{\{v_k>t\}}\,\dd t=(v_k-q)_+.
\]
This proves \eqref{eq:V66-quantile-excess}, including the cases
\(q=0\) and \(q=V_*\).
\end{proof}

\begin{proposition}[Quantitative mesoscopic acquisition]
\label{prop:V66-mesoscopic-acquisition}
On the dual-symmetric conditional family, an output-only estimate of
the form
\[
 \mathfrak k_{(k,k)}
 \le\frac C{1+s_\alpha(k+1)^2}
\]
cannot yield the inverse-input-mass JREC scale.  A necessary
quantile conclusion is
\begin{equation}
 \frac1{d_{R_n}}
 \sum_{\substack{k:\,
       \mathfrak k_{(k,k)}>C/(s_\alpha n^2)}}
 (k+1)^3<1.
 \label{eq:V66-required-mesoscopic-quantile}
\end{equation}
For the core capacity itself, the exact necessary-and-sufficient
acquisition is
\begin{align}
 q_n(\mathfrak k)
 &+
 \frac1{d_{R_n}}\sum_{k=0}^{n}(k+1)^3
 \bigl(\mathfrak k_{(k,k)}-q_n(\mathfrak k)\bigr)_+
 \notag\\
 &\hspace{8em}\le\frac C{s_\alpha n^2}.
 \label{eq:V66-required-quantile-excess}
\end{align}
Both statements are to be tested uniformly on the relevant ancestry;
inside JREC the same criterion must be applied to the complete
pair--core product weight rather than to \(\mathfrak k\) alone.
\end{proposition}

\begin{proof}
The failure of the output-only scale is
\Cref{thm:V64-dual-model-capacity,
thm:V66-mesoscopic-location}.  By definition
\eqref{eq:V65-Schur-quantile}, the strict inequality
\eqref{eq:V66-required-mesoscopic-quantile} is equivalent, up to an
arbitrarily small enlargement of \(C\), to
\[
 q_n(\mathfrak k)\le\frac C{s_\alpha n^2}.
\]
This is necessary by
\eqref{eq:V65-capacity-dominates-quantile}, but does not control the
subcritical tail above the quantile.  The exact identity
\eqref{eq:V66-quantile-excess} shows that adding precisely the excess
term in \eqref{eq:V66-required-quantile-excess} is equivalent to the
desired full core-capacity bound.
\end{proof}

\section{V66 ledger and next acquisition}
\label{sec:V66-ledger}

\begin{theorem}[Integrated V66 statement]
\label{thm:V66-integrated}
Preserving V1--V65, V66 proves:
\begin{enumerate}[label=\textup{(V66.\arabic*)},leftmargin=5.5em]
\item the exact multiplicity-one final-singlet factorization
      \eqref{eq:V66-singlet-factorization};
\item its original-pair two-defect estimate
      \eqref{eq:V66-singlet-original-pair};
\item the uniform hard cold-window taxes
      \eqref{eq:V66-hard-cold-window}--%
      \eqref{eq:V66-fixed-window-capacity};
\item localization of one unit of dual-symmetric Schur mass at the
      geometric-mean output scale
      \eqref{eq:V66-geometric-mean-output}; and
\item the exact quantile--excess identity
      \eqref{eq:V66-quantile-excess} and the explicit mesoscopic
      acquisition \eqref{eq:V66-required-quantile-excess}.
\end{enumerate}
\end{theorem}

\begin{proof}
The five items are
\Cref{thm:V66-singlet-factorization,
cor:V66-singlet-original-pair,
thm:V66-hard-cold-window,cor:V66-fixed-window-capacity,
thm:V66-mesoscopic-location,
lem:V66-quantile-excess,
prop:V66-mesoscopic-acquisition}.
\end{proof}

\begingroup
\small
\begin{longtable}{>{\raggedright\arraybackslash}p{0.29\textwidth}
                  >{\raggedright\arraybackslash}p{0.15\textwidth}
                  >{\raggedright\arraybackslash}p{0.48\textwidth}}
\toprule
\textbf{V66 row}&\textbf{Status}&\textbf{Advance or obstruction}\\
\midrule
\endfirsthead
\toprule
\textbf{V66 row}&\textbf{Status}&\textbf{Advance or obstruction}\\
\midrule
\endhead
Final singlet
& \MGclosed
& The actual five-term cumulant is a signed difference of two
  two-defect squares and vanishes if one input is trivial.\\[0.35em]
Fixed cold-output windows
& \MGclosed
& Every bounded output-Casimir family has a uniform inverse maximum
  input-mass partial-trace capacity.\\[0.35em]
Dual-symmetric loss location
& \MGclosed
& Unit one-sided Schur mass first appears at
  \(\xi(T)\asymp\sqrt{\xi(R)\xi(Q)}\), not at the singlet.\\[0.35em]
Mesoscopic Racah shell
& \MGopen
& The three noncommuting pair-moment operators must be estimated as
  one signed ancestry-dependent symbol.\\[0.35em]
Hyperbolic JREC gate
& \MGopen
& Either the mesoscopic cut quantile or the coherent pair companion
  must supply the missing inverse square root.\\[0.35em]
Complete collision and \([3,2]\)
& \MGopen
& JREC and fixed-type puncture aggregation remain.\\[0.35em]
Full Yang--Mills mass gap
& \MGopen
& Higher quartic/all-order MTA, WUFC/CPM, nonlinear construction,
  continuum OS existence/nontriviality, and spectral positivity remain
  downstream.\\
\bottomrule
\end{longtable}
\endgroup

\begin{openproblem}[V67 mandate: ancestry-retaining mesoscopic symbol]
\label{op:V67}
\begin{enumerate}[label=\textup{(V67.\arabic*)},leftmargin=5.5em]
\item Keep a fixed ancestry
      \(A\otimes B\to S\), \(S\otimes C\to T\) through the complete
      multiplicity space.  Do not replace the \(AC\) and \(BC\)
      moment operators by separate norms.
\item On
      \(S=(n,0)\), \(C=(0,n)\), and
      \(T=(k,k)\) with \(k\asymp\sqrt n\), write the two Racah
      conjugations in a common \(AB\)-intermediate basis.
\item Split the signed symbol into its scalar part and the two
      commutators created by changing bracketings.  Use unitarity
      before any entrywise Racah estimate.
\item Prove the full quantile--excess bound
      \eqref{eq:V66-required-quantile-excess}, or construct an actual
      ancestry for which even its necessary quantile part
      \eqref{eq:V66-required-mesoscopic-quantile} fails.  A failure
      must be tested inside the product threshold
      \eqref{eq:V65-hyperbolic-projector}, where the pair companion
      may still pay it.
\item Preserve the fixed global retained-row orientation and the
      simultaneous-orientation quantile
      \eqref{eq:V65-both-orientation-quantile}.
\item Keep puncture pair-label banks, multi-collision, \([3,3]\), and
      \([4]\) suspended.  The next companion audit remains V70.
\end{enumerate}
\end{openproblem}

\begin{firewall}[End-of-V66 claim boundary]
\label{fire:V66-claim-boundary}
V66 proves the exact singlet cancellation, fixed-cold-window capacity,
and localization of the remaining dual-symmetric loss to a growing
mesoscopic Racah shell.  It does not prove the ancestry-retaining
mesoscopic quantile, the coherent pair companion, the hyperbolic JREC
gate, JREC, simultaneous-sideways closure, fixed-type puncture
aggregation, multi-collision aggregation, complete \([3,2]\),
\([3,3]\), \([4]\), global \(m=4\), all-order MTA, WUFC, CPM, the
nonlinear Perron packet, finite-edge margins, constructive capacity,
protected-sector floors, local-algebra noncollapse, reflection
positivity, a nontrivial continuum OS state, or a positive
Yang--Mills spectral gap.  The full Yang--Mills mass gap is not
proved.
\end{firewall}

\section*{Version V66 append-only record}
\addcontentsline{toc}{section}{Version V66 append-only record}
The complete V65 source was retained before this continuation.  Its
SHA--256 digest was
\begin{center}\scriptsize\ttfamily
f78faf0e8debf4b451a38e11c301f1131fccfa2625a81fb3173474148a574f8b.
\end{center}
No mathematical content of V65 was altered during the V66 append.  On
the next iteration, retain this full file, append before its unique
active document-closing command, increment the filename to
\texttt{\_V67.tex}, and remove only the superseded V66 file from this
Yang--Mills note series.

% =============================================================================
% END APPENDED VERSION V66 MATERIAL
% =============================================================================

% =============================================================================
% BEGIN APPENDED VERSION V67 MATERIAL
% =============================================================================

\chapter{V67: Ancestry-resolved absolute cumulants and exact
matrix-valued partial traces}
\label{ch:V67}

\section{Purpose and upstream correction}
\label{sec:V67-purpose}

V66 localized the scalar-majorant loss to a mesoscopic family on
which the three pair-moment operators live in different Racah
bracketings.  Before estimating Racah entries, one must correct a
more basic loss.  V60 assigned the complete final block the scalar
weight \(\|K_T\|_{\rm op}\).  V63 could then take an exact partial
trace, but only after replacing the full ancestry-dependent
multiplicity operator by that scalar.

V67 keeps the positive operator \(|K_T|\) itself.  Its partial trace
has an exact sequential formula: off-diagonal pair intermediates and
second-stage LR copies vanish, while the matrix on first-stage
ancestry copies survives.  The V63 scalar LR transform is the special
case in which this surviving matrix has first been replaced by a
multiple of the identity.

\begin{firewall}[Scope of V67]
\label{fire:V67-scope}
The matrix formula below is exact for every compact group and every
positive multiplicity weight.  It proves a sharper fully
spectator-amplified row-factorized trace-norm inequality for the signed
cumulant.  The SU(3) estimate of the surviving ancestry matrices
remains open; therefore JREC is not yet proved.
\end{firewall}

\section{Positive multiplicity weights in a sequential basis}
\label{sec:V67-positive-multiplicity}

Use the sequential embeddings of V63,
\[
 V_{T,S,\alpha,\beta}
 =
 (\iota_{S,\alpha}\otimes I_C)\jmath_{T,\beta}^{S},
\]
and identify the multiplicity space of \(T\) with the orthogonal path
basis
\[
 \mathcal M_T
 =
 \bigoplus_S
 \mathbb C^{N_{AB}^{S}}\otimes
 \mathbb C^{N_{SC}^{T}}.
\]
Let \(W_T\ge0\) be an arbitrary operator on \(\mathcal M_T\).  Define
the positive invariant carrier operator
\begin{equation}
 Q_W
 :=
 \sum_T\sum_{\substack{S,\alpha,\beta\\
                       S',\alpha',\beta'}}
 (W_T)_{S\alpha\beta,S'\alpha'\beta'}
 V_{T,S,\alpha,\beta}
 V_{T,S',\alpha',\beta'}^*.
 \label{eq:V67-positive-carrier-operator}
\end{equation}
For fixed \(T,S,\beta\), write
\begin{equation}
 W_T[S,\beta]
 :=
 \left[
 (W_T)_{S\alpha\beta,S\alpha'\beta}
 \right]_{\alpha,\alpha'=1}^{N_{AB}^{S}},
 \label{eq:V67-ancestry-principal-block}
\end{equation}
the positive principal block acting on the first-stage ancestry
copies.

\begin{theorem}[Exact matrix-valued retained-pair partial trace]
\label{thm:V67-matrix-partial-trace}
On the \((AB)C\) sequential decomposition,
\begin{equation}
 \boxed{
 \operatorname{Tr}_{C}Q_W
 =
 \bigoplus_S
 \left[
 I_{H_S}\otimes
 \frac1{d_S}
 \sum_Td_T
 \sum_{\beta=1}^{N_{SC}^{T}}W_T[S,\beta]
 \right].}
 \label{eq:V67-matrix-partial-trace}
\end{equation}
Consequently
\begin{equation}
 \boxed{
 \|\operatorname{Tr}_{C}Q_W\|_{\rm op}
 =
 \max_S
 \left\|
 \frac1{d_S}\sum_Td_T
 \sum_{\beta=1}^{N_{SC}^{T}}W_T[S,\beta]
 \right\|_{\rm op}.}
 \label{eq:V67-matrix-partial-trace-norm}
\end{equation}
\end{theorem}

\begin{proof}
Consider first one matrix unit in
\eqref{eq:V67-positive-carrier-operator}.  The partial trace
\[
 \operatorname{Tr}_C\!
 \left[
 \jmath_{T,\beta}^{S}
 (\jmath_{T,\beta'}^{S'})^*
 \right]
\]
is an intertwiner from \(H_{S'}\) to \(H_S\).  It vanishes when
\(S\not\cong S'\).  When \(S=S'\), Schur's lemma makes it a scalar
multiple of \(I_{H_S}\).  Its trace is
\[
 \operatorname{Tr}\!
 \left[
 (\jmath_{T,\beta'}^{S})^*
 \jmath_{T,\beta}^{S}
 \right]
 =d_T\delta_{\beta\beta'},
\]
because the embeddings belonging to distinct LR copies are
orthogonal.  Hence
\begin{equation}
 \operatorname{Tr}_C\!
 \left[
 \jmath_{T,\beta}^{S}
 (\jmath_{T,\beta'}^{S'})^*
 \right]
 =
 \mathbf1_{\{S=S'\}}
 \delta_{\beta\beta'}\frac{d_T}{d_S}I_{H_S}.
 \label{eq:V67-path-matrix-unit-trace}
\end{equation}
Conjugating by
\(\iota_{S,\alpha}\) and \(\iota_{S,\alpha'}^*\) leaves the matrix
unit between the two first-stage ancestry copies.  Sum its coefficient
from \(W_T\), then sum \(T,S,\beta\).  This gives
\eqref{eq:V67-matrix-partial-trace}.  The direct-sum and tensor
identity structure makes its operator norm the maximum in
\eqref{eq:V67-matrix-partial-trace-norm}.
\end{proof}

\begin{corollary}[V63 is the scalar multiplicity special case]
\label{cor:V67-V63-recovery}
If \(W_T=u_TI_{\mathcal M_T}\), then
\[
 W_T[S,\beta]=u_TI_{\mathbb C^{N_{AB}^{S}}},
\]
and \eqref{eq:V67-matrix-partial-trace} becomes
\[
 \bigoplus_S I_{H_S}\otimes
 \left[
 \frac1{d_S}\sum_TN_{SC}^{T}d_Tu_T
 \right]I_{\mathbb C^{N_{AB}^{S}}},
\]
which is exactly \eqref{eq:V63-weighted-core-trace}.
\end{corollary}

\begin{proof}
There are \(N_{SC}^{T}\) identical \(\beta\)-blocks.  Substitute in
\eqref{eq:V67-matrix-partial-trace}.
\end{proof}

\section{Keeping the absolute ternary symbol}
\label{sec:V67-absolute-symbol}

For each complete final \(T\)-isotypic block, the exact ternary
cumulant \(K_T\) is self-adjoint: every moment operator is a unitary
conjugate of a real diagonal multiplier, and the five coefficients
are real.  Put
\begin{equation}
 Q_{|K|}
 :=
 \bigoplus_T I_{H_T}\otimes|K_T|.
 \label{eq:V67-absolute-cumulant-operator}
\end{equation}

\begin{corollary}[Exact ancestry-resolved absolute-cut trace]
\label{cor:V67-absolute-cut-trace}
The conditional positive operator which retains all signed
within-\(T\) cancellations is
\begin{equation}
 \boxed{
 \|\operatorname{Tr}_{C}Q_{|K|}\|_{\rm op}
 =
 \max_S
 \left\|
 \frac1{d_S}\sum_Td_T
 \sum_{\beta=1}^{N_{SC}^{T}}|K_T|[S,\beta]
 \right\|_{\rm op}.}
 \label{eq:V67-absolute-cut-trace}
\end{equation}
It obeys the monotone comparison
\begin{align}
 \|\operatorname{Tr}_{C}Q_{|K|}\|_{\rm op}
 &\le
 \max_S\frac1{d_S}
 \sum_TN_{SC}^{T}d_T\|K_T\|_{\rm op}.
 \label{eq:V67-norm-coarsening-comparison}
\end{align}
\end{corollary}

\begin{proof}
Apply \Cref{thm:V67-matrix-partial-trace} with \(W_T=|K_T|\).
The operator inequality
\[
 0\le|K_T|\le\|K_T\|_{\rm op}I_{\mathcal M_T}
\]
survives principal compression, summation, partial trace, and
operator norm.  The scalar case is
\Cref{cor:V67-V63-recovery}, proving
\eqref{eq:V67-norm-coarsening-comparison}.
\end{proof}

\begin{assessment}[What the V64 model actually majorizes]
\label{ass:V67-V64-majorant-scope}
The V64 dual-symmetric model estimates the right side of
\eqref{eq:V67-norm-coarsening-comparison}, after the further output
majorant
\(\|K_T\|_{\rm op}\lesssim(1+s_\alpha\Xi(T))^{-1}\).
Its square-root loss is therefore not an obstruction to the exact
left side \eqref{eq:V67-absolute-cut-trace}.  The potential gain is
now precisely located in the spectrum and ancestry principal blocks
of \(|K_T|\), before the scalar norm coarsening.
\end{assessment}

\section{A matrix-weighted row-factorized trace tax}
\label{sec:V67-matrix-trace-tax}

Let \(H=\bigotimes_{i=1}^{m}H_i\), let
\(\mathcal K_i\) be arbitrary spectator spaces, and let \(K\) be a
finite-dimensional operator on \(H\).  Put \(Q=|K|\).  When
\(|K^*|=|K|\), define for every nonempty \(I\subseteq[m]\)
\begin{equation}
 \kappa_I(K)
 :=
 \min\left\{
 \|K\|_{\rm op},
 \|\operatorname{Tr}_{I^c}|K|\|_{\rm op}
 \right\}.
 \label{eq:V67-matrix-conditional-tax}
\end{equation}
The hypothesis \(|K^*|=|K|\) holds in particular for self-adjoint
\(K\), and also after multiplying orthogonal self-adjoint blocks by
arbitrary scalar phases.

\begin{lemma}[Positive product-state tax]
\label{lem:V67-positive-product-tax}
For \(Q\ge0\) and arbitrary positive trace-class
\(\rho_i\) on \(H_i\),
\begin{equation}
 \operatorname{Tr}\left[Q\bigotimes_i\rho_i\right]
 \le
 \min\left\{
 \|Q\|_{\rm op},
 \|\operatorname{Tr}_{I^c}Q\|_{\rm op}
 \right\}
 \prod_i\operatorname{Tr}\rho_i.
 \label{eq:V67-positive-product-tax}
\end{equation}
\end{lemma}

\begin{proof}
Normalize the nonzero \(\rho_i\) to density matrices; the zero case
is immediate.  The operator norm bound follows from
\(Q\le\|Q\|_{\rm op}I\).  For the partial-trace bound, a density
matrix satisfies \(0\le\rho_{I^c}\le I_{I^c}\).  Hence
\[
 \operatorname{Tr}[Q(\rho_I\otimes\rho_{I^c})]
 \le
 \operatorname{Tr}[(\operatorname{Tr}_{I^c}Q)\rho_I]
 \le
 \|\operatorname{Tr}_{I^c}Q\|_{\rm op}.
\]
Restore the traces of the original \(\rho_i\).
\end{proof}

\begin{theorem}[Spectator-amplified matrix trace tax]
\label{thm:V67-matrix-trace-tax}
Suppose \(|K^*|=|K|\).  For arbitrary
\(A_i\in\mathcal S_1(H_i\otimes\mathcal K_i)\),
\begin{equation}
 \boxed{
 \left|
 \operatorname{Tr}\left[
 (K\otimes I_{\rm sp})\bigotimes_{i=1}^{m}A_i
 \right]\right|
 \le
 \kappa_I(K)\prod_{i=1}^{m}\|A_i\|_1.}
 \label{eq:V67-matrix-trace-tax}
\end{equation}
If \(K=\bigoplus_\omega K_\omega\) on orthogonal carrier blocks and
each \(K_\omega\) is self-adjoint, then
\begin{equation}
 \boxed{
 \sum_\omega
 \left|
 \operatorname{Tr}\left[
 (K_\omega\otimes I_{\rm sp})
 P_\omega\left(\bigotimes_iA_i\right)P_\omega
 \right]\right|
 \le
 \kappa_I(K)\prod_i\|A_i\|_1,}
 \label{eq:V67-block-absolute-trace-tax}
\end{equation}
where \(K=\bigoplus_\omega K_\omega\) may be used on the right because
\(|\bigoplus_\omega\varepsilon_\omega K_\omega|
=\bigoplus_\omega|K_\omega|\) for all phases
\(|\varepsilon_\omega|=1\).
\end{theorem}

\begin{proof}
First take \(A_i=|x_i\rangle\langle y_i|\), allowing
\(x_i,y_i\in H_i\otimes\mathcal K_i\), and put
\(x=\bigotimes_ix_i\), \(y=\bigotimes_iy_i\).
The polar-decomposition Cauchy--Schwarz inequality gives
\[
 |\langle y,(K\otimes I_{\rm sp})x\rangle|
 \le
 \langle x,(|K|\otimes I_{\rm sp})x\rangle^{1/2}
 \langle y,(|K^*|\otimes I_{\rm sp})y\rangle^{1/2}.
\]
The carrier reduction of \(|x\rangle\langle x|\) is the tensor
product of the carrier reductions of
\(|x_i\rangle\langle x_i|\), each having trace \(\|x_i\|^2\).
Apply \Cref{lem:V67-positive-product-tax} to \(Q=|K|\);
apply it again to the \(y_i\) and \(|K^*|=|K|\).  Taking square roots
gives
\[
 |\langle y,(K\otimes I_{\rm sp})x\rangle|
 \le\kappa_I(K)\prod_i\|x_i\|\,\|y_i\|.
\]
Expand each \(A_i\) in singular values.  The resulting nonnegative
sums factor as \(\prod_i\|A_i\|_1\), proving
\eqref{eq:V67-matrix-trace-tax}.

For the block assertion, choose phases \(\varepsilon_\omega\) which
align the displayed scalar traces.  Their sum of absolute values is
the absolute value of the single trace with
\(\widetilde K=\bigoplus_\omega\varepsilon_\omega K_\omega\).
This operator satisfies
\(|\widetilde K|=|\widetilde K^*|
=\bigoplus_\omega|K_\omega|\) and
\(\|\widetilde K\|=\|K\|\).  Apply the first assertion.
\end{proof}

\begin{theorem}[Complete spectator trace-norm matrix tax]
\label{thm:V67-complete-matrix-tax}
Let \(V_\omega:G_\omega\to H\) be isometries with mutually orthogonal
ranges, and let \(K_\omega=K_\omega^*\) act on \(G_\omega\).  Put
\begin{equation}
 Q_{\rm abs}:=
 \sum_\omega V_\omega|K_\omega|V_\omega^*
 \quad\hbox{on }H,
 \label{eq:V67-general-absolute-carrier}
\end{equation}
and let \(\mathcal K=\bigotimes_i\mathcal K_i\).  For
\(A=\bigotimes_iA_i\), define the spectator-valued compressed
coefficient
\begin{equation}
 Y_\omega(A):=
 \operatorname{Tr}_{G_\omega}\left[
 (K_\omega\otimes I_{\mathcal K})
 (V_\omega^*\otimes I_{\mathcal K})
 A
 (V_\omega\otimes I_{\mathcal K})
 \right].
 \label{eq:V67-spectator-coefficient}
\end{equation}
Then, for every nonempty \(I\subseteq[m]\),
\begin{equation}
 \boxed{
 \sum_\omega\|Y_\omega(A)\|_1
 \le
 \min\left\{
 \|Q_{\rm abs}\|_{\rm op},
 \|\operatorname{Tr}_{I^c}Q_{\rm abs}\|_{\rm op}
 \right\}
 \prod_i\|A_i\|_1.}
 \label{eq:V67-complete-matrix-tax}
\end{equation}
The constant is one and is uniform in all spectator dimensions.
\end{theorem}

\begin{proof}
By trace-norm duality, choose contractions \(Z_\omega\) on
\(\mathcal K\) so that, up to an arbitrarily small error,
\[
 \sum_\omega\|Y_\omega(A)\|_1
 =
 \left|\sum_\omega\operatorname{Tr}
 [Z_\omega Y_\omega(A)]\right|.
\]
The right side is
\[
 |\operatorname{Tr}[L_ZA]|,
 \qquad
 L_Z:=
 \sum_\omega
 (V_\omega K_\omega V_\omega^*)\otimes Z_\omega.
\]
Orthogonality of the carrier ranges makes \(L_Z\) a direct sum in
\(\omega\).  Since \(\|Z_\omega\|_{\rm op}\le1\),
\begin{align}
 |L_Z|
 &=
 \bigoplus_\omega
 V_\omega|K_\omega|V_\omega^*
 \otimes|Z_\omega|
 \le Q_{\rm abs}\otimes I_{\mathcal K},
 \label{eq:V67-left-absolute-domination}\\
 |L_Z^*|
 &\le Q_{\rm abs}\otimes I_{\mathcal K}.
 \label{eq:V67-right-absolute-domination}
\end{align}

For rank-one \(A_i=|x_i\rangle\langle y_i|\), polar
Cauchy--Schwarz bounds \(|\operatorname{Tr}[L_ZA]|\) by the geometric
mean of the \(x\)-expectation of \(|L_Z|\) and the \(y\)-expectation
of \(|L_Z^*|\).  Use
\eqref{eq:V67-left-absolute-domination}--%
\eqref{eq:V67-right-absolute-domination}, take carrier reductions,
and apply \Cref{lem:V67-positive-product-tax} to \(Q_{\rm abs}\).
The result is
\[
 |\operatorname{Tr}[L_ZA]|
 \le
 \min\left\{\|Q_{\rm abs}\|_{\rm op},
 \|\operatorname{Tr}_{I^c}Q_{\rm abs}\|_{\rm op}\right\}
 \prod_i\|x_i\|\,\|y_i\|.
\]
Expand the general \(A_i\) by singular values and sum.  Finally take
the supremum over the \(Z_\omega\) and remove the arbitrarily small
duality error.
\end{proof}

\begin{remark}[The scalar theorem is the fully traced case]
\label{rem:V67-scalar-versus-trace-norm}
Taking every spectator space one-dimensional in
\Cref{thm:V67-complete-matrix-tax} recovers
\eqref{eq:V67-block-absolute-trace-tax}.  Thus no additional
operator-space hypothesis is needed for the FI trace norm.
\end{remark}

\section{The exact matrix target for the collision core}
\label{sec:V67-collision-matrix-target}

Restore the \(oa\) and \(bc\) pair output labels.  For the weighted
cut put
\begin{equation}
 \widetilde K_T:=
 (1+s_\alpha\Xi(T))K_T.
 \label{eq:V67-weighted-cut-operator}
\end{equation}
Before scalar norm coarsening, the positive \(ab\)-oriented
pair--core carrier is
\begin{equation}
 Q_{ab}^{\rm abs}
 :=
 \bigoplus_{\lambda,\mu,T}
 a_\lambda b_\mu
 P_{oa,\lambda}\otimes P_{bc,\mu}
 \otimes
 \left(I_{H_T}\otimes|\widetilde K_T|\right).
 \label{eq:V67-joint-absolute-carrier}
\end{equation}

\begin{proposition}[Exact ancestry-resolved pair--core trace]
\label{prop:V67-joint-absolute-trace}
With the pair partial-trace scalars of V65,
\begin{align}
 &\left\|
 \operatorname{Tr}_{\{a,b\}^c}Q_{ab}^{\rm abs}
 \right\|_{\rm op}
 \notag\\
 &\quad=
 \left(\sum_\lambda a_\lambda x_\lambda^a\right)
 \left(\sum_\mu b_\mu y_\mu^b\right)
 \max_S
 \left\|
 \frac1{d_S}\sum_Td_T
 \sum_{\beta=1}^{N_{SC}^{T}}
 |\widetilde K_T|[S,\beta]
 \right\|_{\rm op}.
 \label{eq:V67-joint-absolute-trace}
\end{align}
The \(ac\) orientation has the identical formula with
\(y_\mu^c\), the \(AC\)-ancestry decomposition, and row \(B\) traced
from the core.  Replacing
\(|\widetilde K_T|\) by
\(\|\widetilde K_T\|_{\rm op}I\) recovers the direct scalar
coarsening underlying the V62--V66 majorants.
\end{proposition}

\begin{proof}
Every summand of \eqref{eq:V67-joint-absolute-carrier} is a tensor
product over the three disjoint banks.  Partial trace factorizes.
The two pair-bank sums are scalar by
\eqref{eq:V65-pair-partial-scalars-ab}; apply
\eqref{eq:V67-absolute-cut-trace} to the remaining core factor.
This proves \eqref{eq:V67-joint-absolute-trace}.  The other orientation
is its row permutation.  The last assertion is
\eqref{eq:V67-norm-coarsening-comparison}.
\end{proof}

\begin{assessment}[Revised first open inequality]
\label{ass:V67-revised-open-inequality}
The mesoscopic question is no longer whether every
\(\|K_T\|_{\rm op}\) is small.  The sharper required object is the
positive ancestry matrix
\begin{equation}
 \mathscr A_S
 :=
 \frac1{d_S}\sum_Td_T
 \sum_{\beta=1}^{N_{SC}^{T}}
 |\widetilde K_T|[S,\beta].
 \label{eq:V67-ancestry-matrix}
\end{equation}
A proof may exploit that a norm-maximizing vector for one \(T\) need
not survive the principal compression for the same \(S,\beta\) or
align with the norm-maximizing vectors of the other outputs.  The
scalar model of V64 deliberately erased both effects.
\end{assessment}

Define the direct matrix collision capacity
\begin{equation}
 \mathcal K_{\rm mat}
 :=
 \min\left\{
 \|Q_{ab}^{\rm abs}\|_{\rm op},
 \|\operatorname{Tr}_{\{a,b\}^c}Q_{ab}^{\rm abs}\|_{\rm op},
 \|\operatorname{Tr}_{\{a,c\}^c}Q_{ab}^{\rm abs}\|_{\rm op}
 \right\},
 \label{eq:V67-matrix-collision-capacity}
\end{equation}
where the last trace is understood after the corresponding row
permutation of the same complete joint carrier.

\begin{theorem}[Matrix capacity is sufficient for the actual collision core]
\label{thm:V67-matrix-capacity-sufficient}
Before the scalar replacements
\[
 K_T\longmapsto\|K_T\|_{\rm op},
 \qquad
 |K_T|\longmapsto\|K_T\|_{\rm op}I,
\]
the complete pair--core spectator trace-norm contribution is at most
\begin{equation}
 \mathcal K_{\rm mat}
 \prod_{i=o,a,b,c}\|B_i\|_1.
 \label{eq:V67-matrix-collision-bound}
\end{equation}
Consequently the sufficient matrix gate
\begin{equation}
 \boxed{
 \mathcal K_{\rm mat}
 \le
 Cs_\alpha\sum_{r=b,c}
 \frac{h_{oa}\widehat H_{ar}h_{bc}}
      {\Xi_a\Xi_r}}
 \label{eq:V67-matrix-JREC-gate}
\end{equation}
closes the actual collision-core coefficient at the same marked-tree
scale as JREC.  It is a sharper sufficient condition than applying
the scalar norm coarsening first.
\end{theorem}

\begin{proof}
Apply \Cref{thm:V67-complete-matrix-tax} to the complete orthogonal
joint output family \((\lambda,\mu,T)\), with block operator
\[
 a_\lambda b_\mu\widetilde K_T.
\]
Its absolute carrier is exactly
\eqref{eq:V67-joint-absolute-carrier}.  Use successively the full
carrier, retained rows \(ab\), and retained rows \(ac\), and take the
smallest of the three valid estimates.  This proves
\eqref{eq:V67-matrix-collision-bound}.  The output-mass allocation and
single final resolvent are the same as in
\Cref{thm:V60-JREC-sufficient}; only the earlier scalar domination of
the cut block has been avoided.  Hence
\eqref{eq:V67-matrix-JREC-gate} supplies exactly the two path weights
needed there.

Finally
\[
 Q_{ab}^{\rm abs}
 \le
 \bigoplus_{\lambda,\mu,T}
 a_\lambda b_\mu
 \|\widetilde K_T\|_{\rm op}
 P_{oa,\lambda}\otimes P_{bc,\mu}\otimes P_T.
\]
Thus every partial-trace tax after scalar norm coarsening dominates
the corresponding matrix tax, proving the last assertion.
\end{proof}

\begin{firewall}[Direct matrix gate versus the V60 definition]
\label{fire:V67-matrix-gate-scope}
V60's JREC remains a valid sufficient scalar majorant.  Theorem
\ref{thm:V67-matrix-capacity-sufficient} returns one step earlier and
gives a different, sharper sufficient gate for the actual signed cut
coefficient.  It does not prove the scalar JREC inequality
\eqref{eq:V60-JREC}, nor is that stronger inequality now required for
the collision-core term if \eqref{eq:V67-matrix-JREC-gate} can be
proved.
\end{firewall}

\section{V67 ledger and next acquisition}
\label{sec:V67-ledger}

\begin{theorem}[Integrated V67 statement]
\label{thm:V67-integrated}
Preserving V1--V66, V67 proves:
\begin{enumerate}[label=\textup{(V67.\arabic*)},leftmargin=5.5em]
\item the exact matrix-valued retained-pair partial trace
      \eqref{eq:V67-matrix-partial-trace};
\item recovery of the V63 scalar formula as the identity-multiplicity
      special case;
\item the ancestry-resolved absolute-cut formula
      \eqref{eq:V67-absolute-cut-trace} and its monotone scalar
      coarsening \eqref{eq:V67-norm-coarsening-comparison};
\item the fully spectator-amplified trace-norm matrix tax
      \eqref{eq:V67-complete-matrix-tax};
\item the exact pair--core trace formula
      \eqref{eq:V67-joint-absolute-trace}; and
\item sufficiency of the direct matrix capacity gate
      \eqref{eq:V67-matrix-JREC-gate} for the actual collision-core
      coefficient.
\end{enumerate}
\end{theorem}

\begin{proof}
The six items are
\Cref{thm:V67-matrix-partial-trace,
cor:V67-V63-recovery,
cor:V67-absolute-cut-trace,
thm:V67-complete-matrix-tax,
prop:V67-joint-absolute-trace,
thm:V67-matrix-capacity-sufficient}.
\end{proof}

\begingroup
\small
\begin{longtable}{>{\raggedright\arraybackslash}p{0.29\textwidth}
                  >{\raggedright\arraybackslash}p{0.15\textwidth}
                  >{\raggedright\arraybackslash}p{0.48\textwidth}}
\toprule
\textbf{V67 row}&\textbf{Status}&\textbf{Advance or obstruction}\\
\midrule
\endfirsthead
\toprule
\textbf{V67 row}&\textbf{Status}&\textbf{Advance or obstruction}\\
\midrule
\endhead
Matrix partial trace
& \MGclosed
& Second-stage LR copies decohere under the trace, while first-stage
  ancestry matrices survive exactly.\\[0.35em]
Complete matrix trace tax
& \MGclosed
& The positive absolute carrier controls arbitrary row-factorized
  spectator trace norms with constant one.\\[0.35em]
Scalar cut majorant
& \MGreduced
& It is a monotone coarsening of the exact ancestry-resolved absolute
  cut, so its V64 square-root barrier is not decisive.\\[0.35em]
Direct matrix collision gate
& \MGopen
& Its sufficiency is proved, but the SU(3) ancestry matrix
  \(\mathscr A_S\) has not yet been bounded at the path scale.\\[0.35em]
Scalar JREC
& \MGopen
& It remains a valid stronger sufficient statement, but the actual
  collision core may instead close through the matrix gate.\\[0.35em]
Complete collision and \([3,2]\)
& \MGopen
& Matrix capacity and fixed-type puncture aggregation remain.\\[0.35em]
Full Yang--Mills mass gap
& \MGopen
& Higher quartic/all-order MTA, WUFC/CPM, nonlinear construction,
  continuum OS existence/nontriviality, and spectral positivity remain
  downstream.\\
\bottomrule
\end{longtable}
\endgroup

\begin{openproblem}[V68 mandate: conditional absolute-cumulant square function]
\label{op:V68}
\begin{enumerate}[label=\textup{(V68.\arabic*)},leftmargin=5.5em]
\item Start from the exact matrix
      \(\mathscr A_S\) in \eqref{eq:V67-ancestry-matrix}; do not
      return to \(\sum_T\|K_T\|_{\rm op}\).
\item Compute the conditional square function
      \[
      \mathscr S_S:=
      \frac1{d_S}\sum_Td_T
      \sum_{\beta=1}^{N_{SC}^{T}}
      (K_T^2)[S,\beta]
      \]
      in the common \(AB\)-ancestry basis.  Use Racah unitarity before
      taking an operator norm.
\item Compare \(\mathscr A_S\) with \(\mathscr S_S\) by a
      noncommutative Cauchy--Schwarz or optimized
      \(|K_T|\le\frac12(\varepsilon I+\varepsilon^{-1}K_T^2)\).
      Record explicitly the dimension factor created by the identity
      term.
\item Test whether the pair factors in
      \eqref{eq:V67-joint-absolute-trace} pay that dimension factor in
      one fixed row orientation.  If only a square-root estimate
      results, identify the exact higher conditional moment required.
\item Re-run the singlet and \(k\asymp\sqrt n\) dual-symmetric
      regressions on \(\mathscr A_S\), not on scalar norms.
\item Keep puncture pair-label banks, multi-collision, \([3,3]\), and
      \([4]\) suspended.  The next companion audit remains V70.
\end{enumerate}
\end{openproblem}

\begin{firewall}[End-of-V67 claim boundary]
\label{fire:V67-claim-boundary}
V67 proves the ancestry-resolved matrix partial trace, the complete
spectator trace-norm matrix tax, and a sharper sufficient matrix gate
for the actual collision core.  It does not prove the conditional
absolute-cumulant matrix estimate, the matrix JREC gate, scalar JREC,
simultaneous-sideways closure, fixed-type puncture aggregation,
multi-collision aggregation, complete \([3,2]\), \([3,3]\), \([4]\),
global \(m=4\), all-order MTA, WUFC, CPM, the nonlinear Perron packet,
finite-edge margins, constructive capacity, protected-sector floors,
local-algebra noncollapse, reflection positivity, a nontrivial
continuum OS state, or a positive Yang--Mills spectral gap.  The full
Yang--Mills mass gap is not proved.
\end{firewall}

\section*{Version V67 append-only record}
\addcontentsline{toc}{section}{Version V67 append-only record}
The complete V66 source was retained before this continuation.  Its
SHA--256 digest was
\begin{center}\scriptsize\ttfamily
4e412295a83eba1812a68eed26d982d90de3f5e4e64ea1660d8575f4620ceca4.
\end{center}
No mathematical content of V66 was altered during the V67 append.  On
the next iteration, retain this full file, append before its unique
active document-closing command, increment the filename to
\texttt{\_V68.tex}, and remove only the superseded V67 file from this
Yang--Mills note series.

% =============================================================================
% END APPENDED VERSION V67 MATERIAL
% =============================================================================

% =============================================================================
% BEGIN APPENDED VERSION V68 MATERIAL
% =============================================================================

\chapter{V68: Conditional absolute-cumulant square functions and the
exact companion dimension}
\label{ch:V68}

\section{Purpose}
\label{sec:V68-purpose}

V67 replaced the scalar weights \(\|K_T\|_{\rm op}\) by the positive
ancestry matrices \(|K_T|[S,\beta]\).  V68 determines exactly what a
second-moment argument can buy for those matrices.  Operator
convexity gives a conditional square-function estimate with one
factor \(d_C^{1/2}\).  This factor is not a proof artifact: it is the
square root of the total one-sided LR mass
\[
 \frac1{d_S}\sum_TN_{SC}^{T}d_T=d_C.
\]
It is also precisely the inverse-square-root companion debit
identified independently by the V65 dual-symmetric quantile.

\begin{firewall}[Scope of V68]
\label{fire:V68-scope}
V68 proves the optimal abstract passage from the conditional square
function to the ancestry-resolved absolute first moment.  It does not
prove that the actual pair bank supplies \(d_C^{-1/2}\), nor does it
estimate the conditional square function at the marked-tree scale.
\end{firewall}

\section{Compressed absolute values and squares}
\label{sec:V68-compressed-squares}

Let \(K_T=K_T^*\) act on the complete multiplicity space
\(\mathcal M_T\).  Retain the principal blocks
\(|K_T|[S,\beta]\) from
\eqref{eq:V67-ancestry-principal-block}, and define
\begin{align}
 \mathscr A_S(K)
 &:=
 \frac1{d_S}\sum_Td_T
 \sum_{\beta=1}^{N_{SC}^{T}}|K_T|[S,\beta],
 \label{eq:V68-conditional-absolute-first}\\
 \mathscr S_S(K)
 &:=
 \frac1{d_S}\sum_Td_T
 \sum_{\beta=1}^{N_{SC}^{T}}(K_T^2)[S,\beta].
 \label{eq:V68-conditional-square-function}
\end{align}
Both are positive operators on the \(AB\)-ancestry copy space for
\(S\).

\begin{lemma}[Principal-compression square inequality]
\label{lem:V68-compression-square}
For every \(T,S,\beta\),
\begin{equation}
 \bigl(|K_T|[S,\beta]\bigr)^2
 \le (K_T^2)[S,\beta].
 \label{eq:V68-compression-square}
\end{equation}
\end{lemma}

\begin{proof}
Let \(P=P_{S,\beta}\) be the projection onto the corresponding
ancestry subspace.  The principal block is the compression
\(P|K_T|P\) restricted to \(P\mathcal M_T\).  Since
\[
 P K_T^2P-(P|K_T|P)^2
 =
 P|K_T|(I-P)|K_T|P\ge0,
\]
and \(K_T^2=|K_T|^2\), the claimed inequality follows.
\end{proof}

\begin{theorem}[Conditional noncommutative square-function bound]
\label{thm:V68-square-function}
For every \(S\),
\begin{equation}
 \boxed{
 \|\mathscr A_S(K)\|_{\rm op}^2
 \le
 d_C\,\|\mathscr S_S(K)\|_{\rm op}.}
 \label{eq:V68-square-function}
\end{equation}
Equivalently, for every \(\varepsilon>0\),
\begin{equation}
 \boxed{
 \mathscr A_S(K)
 \le
 \frac{\varepsilon d_C}{2}I
 \frac1{2\varepsilon}\mathscr S_S(K).}
 \label{eq:V68-epsilon-square-function}
\end{equation}
\end{theorem}

\begin{proof}
For a self-adjoint matrix \(K_T\), functional calculus gives
\[
 2|K_T|\le\varepsilon I+\varepsilon^{-1}K_T^2.
\]
Compress to \((S,\beta)\), multiply by \(d_T/d_S\), and sum.  The
identity coefficient is
\[
 \frac1{d_S}\sum_Td_TN_{SC}^{T}=d_C
\]
by the tensor-product dimension identity.  This proves
\eqref{eq:V68-epsilon-square-function}.

Taking operator norms and choosing
\[
 \varepsilon=
 \sqrt{\frac{\|\mathscr S_S(K)\|_{\rm op}}{d_C}}
\]
proves \eqref{eq:V68-square-function}; the zero square-function case
follows by a limit.

Alternatively, normalize the positive weights
\((d_T/d_S)/d_C\).  Operator convexity of \(X\mapsto X^2\) gives
\[
 \mathscr A_S(K)^2
 \le
 d_C\frac1{d_S}
 \sum_Td_T\sum_\beta
 \bigl(|K_T|[S,\beta]\bigr)^2,
\]
and \Cref{lem:V68-compression-square} gives the same result.
\end{proof}

\begin{remark}[The identity coefficient is exact]
\label{rem:V68-exact-dimension-factor}
No estimate of an LR multiplicity was used in
\eqref{eq:V68-square-function}.  The factor \(d_C\) is the exact
total one-sided mass of the full \(S\otimes C\) output family.  Any
improvement must use the values or signed structure of the actual
\(K_T\), or a companion factor outside this full family.
\end{remark}

\begin{proposition}[Abstract sharpness]
\label{prop:V68-square-function-sharp}
The exponent \(d_C^{1/2}\) in
\eqref{eq:V68-square-function} cannot be improved for arbitrary
self-adjoint output blocks.  More precisely, if every nonzero
principal block equals the same scalar \(aI\), then
\begin{equation}
 \mathscr A_S(K)=d_C|a|I,
 \qquad
 \mathscr S_S(K)=d_Ca^2I,
 \label{eq:V68-square-function-equality}
\end{equation}
and equality holds in \eqref{eq:V68-square-function}.
\end{proposition}

\begin{proof}
Insert the constant blocks in
\eqref{eq:V68-conditional-absolute-first}--%
\eqref{eq:V68-conditional-square-function} and use the dimension
identity.  The two sides of
\eqref{eq:V68-square-function} are then both \(d_C^2a^2\).
\end{proof}

\section{Insertion into the direct matrix collision gate}
\label{sec:V68-collision-insertion}

Apply the preceding theorem to the weighted cut
\(\widetilde K_T=(1+s_\alpha\Xi(T))K_T\).  Put
\begin{equation}
 \mathscr S_S^{\,\sim}
 :=
 \frac1{d_S}\sum_Td_T
 \sum_{\beta=1}^{N_{SC}^{T}}
 (\widetilde K_T^2)[S,\beta].
 \label{eq:V68-weighted-square-function}
\end{equation}

\begin{corollary}[Square-function pair--core bound]
\label{cor:V68-pair-core-square}
The \(ab\)-oriented trace in
\eqref{eq:V67-joint-absolute-trace} obeys
\begin{align}
 &\left\|
 \operatorname{Tr}_{\{a,b\}^c}Q_{ab}^{\rm abs}
 \right\|_{\rm op}
 \notag\\
 &\quad\le
 \left(\sum_\lambda a_\lambda x_\lambda^a\right)
 \left(\sum_\mu b_\mu y_\mu^b\right)
 \max_S
 \sqrt{d_C\|\mathscr S_S^{\,\sim}\|_{\rm op}}.
 \label{eq:V68-pair-core-square}
\end{align}
The \(ac\) orientation has the analogous bound with \(B\) traced.
\end{corollary}

\begin{proof}
Use \eqref{eq:V67-joint-absolute-trace}, then apply
\eqref{eq:V68-square-function} to its ancestry matrix.
\end{proof}

For brevity write
\begin{align}
 P_{ab}
 &:=
 \left(\sum_\lambda a_\lambda x_\lambda^a\right)
 \left(\sum_\mu b_\mu y_\mu^b\right),
 \label{eq:V68-Pab}\\
 \Sigma_{ab}
 &:=
 \max_S\|\mathscr S_S^{\,\sim}\|_{\rm op}.
 \label{eq:V68-Sigmaab}
\end{align}
Define \(P_{ac},\Sigma_{ac}\) by the other row orientation.

\begin{corollary}[Explicit square-function matrix gate]
\label{cor:V68-square-function-gate}
The direct matrix JREC gate
\eqref{eq:V67-matrix-JREC-gate} follows if
\begin{equation}
 \boxed{
 \min\left\{
 P_{ab}\sqrt{d_C\Sigma_{ab}},
 P_{ac}\sqrt{d_B\Sigma_{ac}},
 \|Q_{ab}^{\rm abs}\|_{\rm op}
 \right\}
 \le
 Cs_\alpha\sum_{r=b,c}
 \frac{h_{oa}\widehat H_{ar}h_{bc}}
      {\Xi_a\Xi_r}.}
 \label{eq:V68-square-function-gate}
\end{equation}
\end{corollary}

\begin{proof}
The first two terms dominate the corresponding partial traces by
\Cref{cor:V68-pair-core-square}; the last is the full-carrier option
already present in
\eqref{eq:V67-matrix-collision-capacity}.  Take their minimum and
apply \Cref{thm:V67-matrix-capacity-sufficient}.
\end{proof}

\begin{assessment}[The exact companion debit]
\label{ass:V68-companion-debit}
If the conditional square function is only of its natural scalar
size, \eqref{eq:V68-square-function} costs \(d_C^{1/2}\).  On the
dual boundary ray, \(d_C^{1/2}\asymp n\asymp
\sqrt{\xi(C)}\), exactly the factor isolated in
\eqref{eq:V65-required-companion-debit}.  Thus two independently
derived routes agree:
\begin{enumerate}[label=\textup{(\roman*)},leftmargin=3.7em]
\item the scalar layer quantile requires an inverse-\(n\) companion;
\item the ancestry square function produces a factor \(n\) which must
      be cancelled by a pair tax or by an improved square function.
\end{enumerate}
This agreement fixes the next target numerically.
\end{assessment}

\section{V68 ledger and next acquisition}
\label{sec:V68-ledger}

\begin{theorem}[Integrated V68 statement]
\label{thm:V68-integrated}
Preserving V1--V67, V68 proves:
\begin{enumerate}[label=\textup{(V68.\arabic*)},leftmargin=5.5em]
\item the principal-compression square inequality
      \eqref{eq:V68-compression-square};
\item the conditional noncommutative square-function estimate
      \eqref{eq:V68-square-function};
\item the exact identity coefficient \(d_C\) and abstract sharpness
      \eqref{eq:V68-square-function-equality};
\item insertion into the ancestry-resolved pair--core trace
      \eqref{eq:V68-pair-core-square}; and
\item the explicit sufficient square-function matrix gate
      \eqref{eq:V68-square-function-gate}.
\end{enumerate}
\end{theorem}

\begin{proof}
The five items are
\Cref{lem:V68-compression-square,
thm:V68-square-function,
rem:V68-exact-dimension-factor,
prop:V68-square-function-sharp,
cor:V68-pair-core-square,
cor:V68-square-function-gate}.
\end{proof}

\begingroup
\small
\begin{longtable}{>{\raggedright\arraybackslash}p{0.29\textwidth}
                  >{\raggedright\arraybackslash}p{0.15\textwidth}
                  >{\raggedright\arraybackslash}p{0.48\textwidth}}
\toprule
\textbf{V68 row}&\textbf{Status}&\textbf{Advance or obstruction}\\
\midrule
\endfirsthead
\toprule
\textbf{V68 row}&\textbf{Status}&\textbf{Advance or obstruction}\\
\midrule
\endhead
Conditional square function
& \MGclosed
& Operator convexity converts it to the exact absolute ancestry
  matrix with factor \(d_C^{1/2}\).\\[0.35em]
Dimension factor
& \MGclosed
& It is the exact total one-sided LR mass and is sharp for arbitrary
  block values.\\[0.35em]
Pair companion scale
& \MGreduced
& The required inverse \(d_C^{1/2}\) agrees with the independent V65
  dual-symmetric quantile calculation.\\[0.35em]
Square-function matrix gate
& \MGopen
& Its exact sufficient form is
  \eqref{eq:V68-square-function-gate}; neither orientation has yet
  been bounded at the path scale in the residual region.\\[0.35em]
Complete collision and \([3,2]\)
& \MGopen
& The direct matrix capacity and fixed-type puncture aggregation
  remain.\\[0.35em]
Full Yang--Mills mass gap
& \MGopen
& Higher quartic/all-order MTA, WUFC/CPM, nonlinear construction,
  continuum OS existence/nontriviality, and spectral positivity remain
  downstream.\\
\bottomrule
\end{longtable}
\endgroup

\begin{openproblem}[V69 mandate: pay the exact companion dimension]
\label{op:V69}
\begin{enumerate}[label=\textup{(V69.\arabic*)},leftmargin=5.5em]
\item Keep the two factors
      \(P_{ab}\sqrt{d_C\Sigma_{ab}}\) and
      \(P_{ac}\sqrt{d_B\Sigma_{ac}}\) from
      \eqref{eq:V68-square-function-gate} intact.  Do not estimate
      \(P\), \(d^{1/2}\), and \(\Sigma^{1/2}\) separately.
\item Insert the exact V58 weighted pair sizes in \(P_{ab},P_{ac}\).
      On an incident-edge-heavy endpoint, use the proved V53/V58
      one-edge tax and record whether it pays \(d_C^{1/2}\).
\item On the simultaneous-sideways residual, split by whether
      \(d_C^{1/2}\) is controlled by the \(bc\)-edge representation
      or by punctured-core mass.  Preserve one retained-row
      orientation throughout each branch.
\item Compute \(\mathscr S_S^{\,\sim}\) by expanding the five-term
      cumulant only after squaring the common multiplicity operator.
      Use trace-preserving Racah conjugation to combine the \(25\)
      terms; do not apply the triangle inequality term by term.
\item Test the exact dual-symmetric mesoscopic family and the V60
      alternating-height array.  A bound that merely reproduces
      \(d_C^{1/2}\) without a companion debit is insufficient.
\item Keep puncture pair-label banks, multi-collision, \([3,3]\), and
      \([4]\) suspended.  Perform the next companion audit at V70.
\end{enumerate}
\end{openproblem}

\begin{firewall}[End-of-V68 claim boundary]
\label{fire:V68-claim-boundary}
V68 proves the optimal abstract square-function bridge and isolates
the exact companion dimension required by the ancestry-resolved
matrix gate.  It does not pay that dimension, estimate the actual
conditional square function at path scale, prove the matrix JREC
gate, scalar JREC, simultaneous-sideways closure, fixed-type puncture
aggregation, multi-collision aggregation, complete \([3,2]\),
\([3,3]\), \([4]\), global \(m=4\), all-order MTA, WUFC, CPM, the
nonlinear Perron packet, finite-edge margins, constructive capacity,
protected-sector floors, local-algebra noncollapse, reflection
positivity, a nontrivial continuum OS state, or a positive
Yang--Mills spectral gap.  The full Yang--Mills mass gap is not
proved.
\end{firewall}

\section*{Version V68 append-only record}
\addcontentsline{toc}{section}{Version V68 append-only record}
The complete V67 source was retained before this continuation.  Its
SHA--256 digest was
\begin{center}\scriptsize\ttfamily
4e33db1ad92ee04ffd1a5c8ae6deef387a3db55af0b5333d0c6a970b0fbcbe4e.
\end{center}
No mathematical content of V67 was altered during the V68 append.  On
the next iteration, retain this full file, append before its unique
active document-closing command, increment the filename to
\texttt{\_V69.tex}, and remove only the superseded V68 file from this
Yang--Mills note series.

% =============================================================================
% END APPENDED VERSION V68 MATERIAL
% =============================================================================

% =============================================================================
% BEGIN APPENDED VERSION V69 MATERIAL
% =============================================================================

\chapter{V69: Covariance-square reduction of the ancestry matrix}
\label{ch:V69}

\section{Purpose}
\label{sec:V69-purpose}

V68 requested a calculation of the square of the five-term ternary
cumulant without applying a \(25\)-term triangle inequality.  The
three-covariance identity of V57 does exactly this.  In one fixed
row orientation it expresses the ternary symbol as the sum of three
self-adjoint binary cuts.  A single operator Cauchy--Schwarz then
reduces its square to three positive square functions.  One of them
is diagonal in the retained \(AB\)-ancestry basis and can be evaluated
by the exact V44 pair-Casimir moments.

\begin{firewall}[Scope of V69]
\label{fire:V69-scope}
The reduction below is exact up to the sharp elementary factor three.
It does not estimate the two remaining crossing-cut square functions.
Those two positive operators, rather than the original \(25\) signed
products, are the first open analytic quantities.
\end{firewall}

\section{Three binary cuts in one common multiplicity space}
\label{sec:V69-three-cuts}

Fix irreducible inputs \(A,B,C\) and a complete final output \(T\).
Write
\[
 x=q_{\alpha,A},\qquad
 y=q_{\alpha,B},\qquad
 z=q_{\alpha,C}.
\]
On the \(T\)-multiplicity space, let
\(Q_{AB}^{T}\), \(Q_{AC}^{T}\), and \(Q_{BC}^{T}\) be the three
pair-moment operators: in their natural bracketings their eigenvalues
are the Wilson multipliers of the corresponding pair intermediates.
Define
\begin{align}
 D_{A\mid BC}^{T}
 &:=
 q_{\alpha,T}I-xQ_{BC}^{T},
 \label{eq:V69-composite-cut}\\
 D_{AC}^{T}
 &:=
 Q_{AC}^{T}-xzI,
 \label{eq:V69-AC-defect}\\
 D_{AB}^{T}
 &:=
 Q_{AB}^{T}-xyI.
 \label{eq:V69-AB-defect}
\end{align}
All three are self-adjoint.

\begin{theorem}[Exact covariance decomposition in the \(AB\) carrier]
\label{thm:V69-covariance-decomposition}
The complete ternary symbol satisfies
\begin{equation}
 \boxed{
 K_T
 =
 D_{A\mid BC}^{T}
 -yD_{AC}^{T}
 -zD_{AB}^{T}.}
 \label{eq:V69-covariance-decomposition}
\end{equation}
Consequently
\begin{equation}
 \boxed{
 K_T^2
 \le
 3\left[
 (D_{A\mid BC}^{T})^2
 +y^2(D_{AC}^{T})^2
 +z^2(D_{AB}^{T})^2
 \right].}
 \label{eq:V69-three-square-bound}
\end{equation}
\end{theorem}

\begin{proof}
Expanding the right side of
\eqref{eq:V69-covariance-decomposition} gives
\[
 q_{\alpha,T}I
 -xQ_{BC}^{T}
 -yQ_{AC}^{T}
 -zQ_{AB}^{T}
 +2xyzI,
\]
which is the exact five-term common-basis symbol
\eqref{eq:V47-common-cumulant}.  This is the operator form of
\eqref{eq:V57-cut-identity}.

For arbitrary operators \(L_1,L_2,L_3\),
\[
 (L_1+L_2+L_3)^*(L_1+L_2+L_3)
 \le3\sum_{j=1}^3L_j^*L_j;
\]
apply scalar Cauchy--Schwarz to
\(\sum_jL_jv\) for each vector \(v\).  Use
\[
 L_1=D_{A\mid BC}^{T},\qquad
 L_2=-yD_{AC}^{T},\qquad
 L_3=-zD_{AB}^{T}
\]
and self-adjointness to prove
\eqref{eq:V69-three-square-bound}.
\end{proof}

\begin{remark}[Why this is not the forbidden termwise estimate]
\label{rem:V69-not-25-term}
The five moment partitions are combined exactly into three
covariances before any positivity estimate.  In particular, the
scalar cancellations \(+2xyz\) are already inside the binary cuts.
Equation \eqref{eq:V69-three-square-bound} creates three squares, not
the absolute values of the \(25\) products obtained by directly
expanding \(K_T^2\).
\end{remark}

\section{Conditional weighted cut square functions}
\label{sec:V69-cut-square-functions}

Put \(w_T:=1+s_\alpha\Xi(T)\).  For fixed retained intermediate \(S\)
in the \(AB\) bracketing, define
\begin{align}
 \mathscr D_{A\mid BC;S}
 &:=
 \frac1{d_S}\sum_Td_T
 \sum_{\beta=1}^{N_{SC}^{T}}
 \bigl(w_T^2(D_{A\mid BC}^{T})^2\bigr)[S,\beta],
 \label{eq:V69-composite-square-function}\\
 \mathscr D_{AC;S}
 &:=
 \frac1{d_S}\sum_Td_T
 \sum_{\beta=1}^{N_{SC}^{T}}
 \bigl(w_T^2(D_{AC}^{T})^2\bigr)[S,\beta],
 \label{eq:V69-AC-square-function}\\
 \mathscr D_{AB;S}
 &:=
 \frac1{d_S}\sum_Td_T
 \sum_{\beta=1}^{N_{SC}^{T}}
 \bigl(w_T^2(D_{AB}^{T})^2\bigr)[S,\beta].
 \label{eq:V69-AB-square-function}
\end{align}

\begin{corollary}[Three-square conditional reduction]
\label{cor:V69-three-square-conditional}
The V68 weighted square function obeys the positive operator
inequality
\begin{equation}
 \boxed{
 \mathscr S_S^{\,\sim}
 \le
 3\left[
 \mathscr D_{A\mid BC;S}
 +y^2\mathscr D_{AC;S}
 +z^2\mathscr D_{AB;S}
 \right].}
 \label{eq:V69-three-square-conditional}
\end{equation}
\end{corollary}

\begin{proof}
Multiply \eqref{eq:V69-three-square-bound} by the scalar \(w_T^2\),
take the positive principal compression \((S,\beta)\), multiply by
\(d_T/d_S\), and sum \(T,\beta\).
\end{proof}

\section{Exact evaluation of the aligned binary square}
\label{sec:V69-aligned-square}

In the \(AB\) ancestry block \(S\), put
\begin{equation}
 \delta_{AB}(S):=
 q_{\alpha,S}-xy.
 \label{eq:V69-aligned-defect}
\end{equation}
Also define the one-sided weighted output moment
\begin{equation}
 M_2(S,C):=
 \frac1{d_S}\sum_TN_{SC}^{T}d_T
 \bigl(1+s_\alpha\Xi(T)\bigr)^2.
 \label{eq:V69-one-sided-output-second}
\end{equation}

\begin{theorem}[Exact aligned-square formula]
\label{thm:V69-aligned-square}
For every retained intermediate \(S\),
\begin{equation}
 \boxed{
 \mathscr D_{AB;S}
 =
 \delta_{AB}(S)^2M_2(S,C)
 I_{\mathbb C^{N_{AB}^{S}}}.}
 \label{eq:V69-aligned-square}
\end{equation}
Moreover, with
\[
 A_S:=1+C_2(S),\qquad A_C:=1+C_2(C),
\]
one has
\begin{equation}
 \boxed{
 M_2(S,C)
 \le
 Cd_C\left[1+s_\alpha(A_S+A_C)\right]^2.}
 \label{eq:V69-output-second-bound}
\end{equation}
\end{theorem}

\begin{proof}
In the \(AB\) bracketing, \(Q_{AB}^{T}\) is diagonal in \(S\) and
acts on every \(S,\alpha,\beta\) path by
\(q_{\alpha,S}\).  Thus
\[
 D_{AB}^{T}[S,\beta]
 =\delta_{AB}(S)I
\]
for every \(T,\beta\).  Substitute this scalar identity in
\eqref{eq:V69-AB-square-function}; the sum over \(\beta\) has
\(N_{SC}^{T}\) terms, proving
\eqref{eq:V69-aligned-square}.

Normalize the output sum by the tensor-product dimension identity:
\[
 \frac1{d_Sd_C}\sum_TN_{SC}^{T}d_T=1.
\]
Under this probability law,
\Cref{lem:V44-pair-Casimir-moments} gives
\[
 \mathbb E C_2(T)=C_2(S)+C_2(C),
\qquad
 \mathbb E C_2(T)^2
 \le\frac98[C_2(S)+C_2(C)]^2.
\]
The convention
\(\Xi(T)\le1+C_2(T)\) in one coordinate and
\((1+u)^2\le2(1+u^2)\) now imply
\[
 \mathbb E
 [1+s_\alpha\Xi(T)]^2
 \le C[1+s_\alpha(A_S+A_C)]^2.
\]
Multiplying by \(d_C\) proves
\eqref{eq:V69-output-second-bound}.
\end{proof}

\begin{corollary}[Explicit aligned contribution]
\label{cor:V69-aligned-contribution}
The last term in
\eqref{eq:V69-three-square-conditional} satisfies
\begin{align}
 z^2\mathscr D_{AB;S}
 &\le
 Cd_Cq_{\alpha,C}^2
 [1+s_\alpha(A_S+A_C)]^2
 \delta_{AB}(S)^2I
 \label{eq:V69-aligned-contribution}\\
 &\le
 Cd_Cq_{\alpha,C}^2
 [1+s_\alpha(A_S+A_C)]^2
 \min\{s_\alpha A_AA_B,1\}^2I.
 \notag
\end{align}
\end{corollary}

\begin{proof}
Use \Cref{thm:V69-aligned-square}, then the binary-defect cap
\eqref{eq:V46-defect-cap} on
\(\delta_{AB}(S)\), with the universal defect bound absorbed into the
minimum.
\end{proof}

\section{Scalarization of the simple crossing square}
\label{sec:V69-simple-crossing}

For an \(AC\) intermediate \(U\), put
\begin{equation}
 \delta_{AC}(U):=
 q_{\alpha,U}-xz,
 \label{eq:V69-AC-scalar-defect}
\end{equation}
and let
\[
 \nu_{AC}(U,\mu):=
 \frac{d_U}{d_Ad_C}
\]
be the exact tensor-dimension probability law, including LR copies.

\begin{theorem}[The \(AC\) crossing square has a scalar Schur bound]
\label{thm:V69-AC-crossing-square}
Uniformly in the retained \(AB\) intermediate \(S\),
\begin{align}
 \mathscr D_{AC;S}
 &\le
 Cd_C\,
 \mathbb E_{\nu_{AC}}\left[
 \bigl(1+s_\alpha[A_B+\xi(U)]\bigr)^2
 \delta_{AC}(U)^2
 \right]I
 \label{eq:V69-AC-crossing-expectation}\\
 &\le
 Cd_C
 \left[1+s_\alpha(A_A+A_B+A_C)\right]^2
 \min\{s_\alpha A_AA_C,1\}^2I.
 \label{eq:V69-AC-crossing-bound}
\end{align}
Consequently the middle contribution in
\eqref{eq:V69-three-square-conditional} is bounded by the right side
of \eqref{eq:V69-AC-crossing-bound} multiplied by
\(q_{\alpha,B}^2\).
\end{theorem}

\begin{proof}
Diagonalize first in the \(AC\) bracketing.  On an \(AC\)
intermediate \(U\), the operator \(D_{AC}^{T}\) has eigenvalue
\(\delta_{AC}(U)\).  If \(T\subset U\otimes B\), fusion geometry
\eqref{eq:V30-Casimir-fusion} gives
\[
 1+s_\alpha\Xi(T)
 \le C\bigl(1+s_\alpha[A_B+\xi(U)]\bigr).
\]
Therefore, as a positive operator on
\(H_A\otimes H_B\otimes H_C\),
\begin{align}
 &\bigoplus_T
 \left(1+s_\alpha\Xi(T)\right)^2(D_{AC}^{T})^2
 \notag\\
 &\qquad\le
 C\sum_{U,\mu}
 \bigl(1+s_\alpha[A_B+\xi(U)]\bigr)^2
 \delta_{AC}(U)^2
 P_{U,\mu}^{AC}\otimes I_B.
 \label{eq:V69-AC-global-domination}
\end{align}
The \(AC\) operator on the right commutes with the irreducible
\(\mathrm{SU}(3)\) action on \(H_A\) after \(C\) is traced out.
Schur's lemma and total trace give
\begin{align}
 &\operatorname{Tr}_C
 \sum_{U,\mu}f(U)P_{U,\mu}^{AC}
 \notag\\
 &\qquad=
 \left[
 \frac1{d_A}\sum_{U,\mu}d_Uf(U)
 \right]I_A
 =
 d_C\mathbb E_{\nu_{AC}}[f(U)]I_A.
 \label{eq:V69-AC-Schur-trace}
\end{align}
Tensor with \(I_B\), then compress to any \(AB\)-intermediate \(S\).
Equations \eqref{eq:V69-AC-global-domination} and
\eqref{eq:V69-AC-Schur-trace} prove
\eqref{eq:V69-AC-crossing-expectation}.

The binary defect cap gives
\[
 |\delta_{AC}(U)|
 \le C\min\{s_\alpha A_AA_C,1\}.
\]
Pull this uniform factor out.  The exact first two Casimir moments
\eqref{eq:V44-pair-Casimir-first}--%
\eqref{eq:V44-pair-Casimir-second} imply
\[
 \mathbb E_{\nu_{AC}}
 \bigl(1+s_\alpha[A_B+\xi(U)]\bigr)^2
 \le
 C[1+s_\alpha(A_A+A_B+A_C)]^2.
\]
This proves \eqref{eq:V69-AC-crossing-bound}.  The final assertion
uses the coefficient \(y^2=q_{\alpha,B}^2\) in
\eqref{eq:V69-three-square-conditional}.
\end{proof}

\begin{assessment}[The first open positive operator]
\label{ass:V69-two-open-squares}
The aligned \(AB\) square is now an explicit scalar and retains both a
Wilson multiplier \(q_{\alpha,C}^2\) and the original \(AB\) binary
defect.  The simple crossing square
\(\mathscr D_{AC;S}\) also has the scalar Schur bound
\eqref{eq:V69-AC-crossing-bound}; no entrywise Racah estimate was
needed.  The sole unresolved positive operator is now
\[
 \mathscr D_{A\mid BC;S}.
\]
It is a binary defect between \(A\) and the composite \(BC\), while
the partial trace removes only \(C\) and retains \(B\).  This mismatch
prevents the preceding Schur scalarization.  A valid next theorem must
retain the \(BC\) intermediate through that partial trace, but it need
not analyze all \(25\) products of the ternary symbol.
\end{assessment}

\section{V69 research ledger and next obstruction}
\label{sec:V69-ledger}

\begin{theorem}[Net reduction achieved in V69]
\label{thm:V69-net-reduction}
At the retained three-leg collision, the square-function input needed
by \Cref{cor:V68-pair-core-square} has been reduced without an entrywise
Racah estimate to the sum of exactly three positive binary-cut
squares.  Two of them are now bounded explicitly:
\begin{align}
 \mathscr S_S^\sim
 &\le
 3\left[
 \mathscr D_{A\mid BC;S}
 +q_{\alpha,B}^{\,2}\mathscr D_{AC;S}
 +q_{\alpha,C}^{\,2}\mathscr D_{AB;S}
 \right],
 \label{eq:V69-net-three}\\
 q_{\alpha,B}^{\,2}\mathscr D_{AC;S}
 &\le
 Cd_Cq_{\alpha,B}^{\,2}
 [1+s_\alpha(A_A+A_B+A_C)]^2
 \min\{s_\alpha A_AA_C,1\}^2I,
 \label{eq:V69-net-AC}\\
 q_{\alpha,C}^{\,2}\mathscr D_{AB;S}
 &\le
 Cd_Cq_{\alpha,C}^{\,2}
 [1+s_\alpha(A_S+A_C)]^2
 \min\{s_\alpha A_AA_B,1\}^2I.
 \label{eq:V69-net-AB}
\end{align}
Thus the unresolved part of the pair-traced square-function gate is
the single positive operator
\(\mathscr D_{A\mid BC;S}\).
\end{theorem}

\begin{proof}
Equation \eqref{eq:V69-net-three} is
\eqref{eq:V69-three-square-conditional}.  The other two estimates are
\eqref{eq:V69-AC-crossing-bound} and
\eqref{eq:V69-aligned-contribution}, with the coefficient from
\eqref{eq:V69-net-three} restored.  Positivity permits their direct
addition and introduces no cancellation hypothesis.
\end{proof}

\begin{center}
\begin{tabular}{|p{0.34\textwidth}|p{0.15\textwidth}|p{0.39\textwidth}|}
\hline
V69 item & Status & Exact disposition\\
\hline
Five-term ternary symbol & Closed at square level &
Exact covariance identity followed by the three-square operator
inequality; no \(25\)-term expansion.\\
\hline
Aligned \(AB\) square & Closed &
Exact scalar formula and pair-Casimir second moment.\\
\hline
Simple \(AC\) crossing square & Closed &
Positive global domination, trace over \(C\), and Schur scalarization.\\
\hline
Composite \(A\mid BC\) square & Open &
Must retain the \(BC\) intermediate while tracing only \(C\).\\
\hline
Pair-traced matrix collision gate & Open &
Depends on a summable bound for the preceding composite square.\\
\hline
Constructive Yang--Mills mass gap & Open &
No full mass-gap theorem is asserted.\\
\hline
\end{tabular}
\end{center}

\begin{openproblem}[V70 mandate: composite-cut partial trace]
\label{op:V70-mandate}
Before further collision estimates, perform the prescribed
five-version comparison with the current Standard Model and
Navier--Stokes notes.  Then start from
\(\mathscr D_{A\mid BC;S}\) itself and:
\begin{enumerate}
\item keep the \(BC\) intermediate, including its multiplicity
indices, through the partial trace over \(C\);
\item derive the exact two-stage partial-trace formula, identifying
the matrix that remains on the \(AB\)-ancestry space;
\item separate the regimes in which the composite \(BC\) mass is
carried principally by \(B\) or by \(C\), using reverse fusion rather
than an output-only maximum;
\item preserve the fixed orientation required by the V58 pair
factors and test the resulting bound on both the dual mesoscopic model
and the alternating-height model.
\end{enumerate}
The target is a positive matrix bound strong enough for
\Cref{cor:V68-pair-core-square}; an entrywise absolute-value estimate is
not an acceptable substitute.
\end{openproblem}

\section{Firewall after V69}
\label{sec:V69-firewall}

\begin{firewall}[What V69 does and does not prove]
\label{fire:V69-final}
V69 proves an exact covariance decomposition of the local ternary
symbol, an operator three-square reduction, the exact aligned-square
formula, and a scalar Schur estimate for the simple crossing square.
It does not bound the composite \(A\mid BC\) square, close the
pair-traced matrix collision gate, sum the full collision forest,
construct the limiting Yang--Mills measure and transfer operator, or
prove a positive spectral gap.  In particular, the Clay Yang--Mills
existence and mass-gap problem remains open in these notes.
\end{firewall}

\section{Append-only record for V69}
\label{sec:V69-record}

V69 was appended to the complete V68 manuscript.  The exact V68
parent had \(50124\) lines, \(1815051\) bytes, and SHA--256 digest
\[
 \texttt{8dbff7f83eb5938414df1c01aebecf9837e2624e78673d0fbf497ec036f2f221}.
\]
The V69 appendix begins only after that parent content; the sole final
\verb|\end{document}| is the one below.

% =============================================================================
% BEGIN APPENDED VERSION V70 MATERIAL
% =============================================================================

\chapter{V70: Companion audit and the composite-cut Racah POVM}
\label{ch:V70}

\section{Purpose and nonduplication boundary}
\label{sec:V70-purpose}

V69 reduced the squared ternary cumulant to three positive binary-cut
squares and bounded the aligned \(AB\) square and the simple \(AC\)
crossing square.  The only remaining square in that decomposition is
\(\mathscr D_{A\mid BC;S}\).  Its apparent obstruction is not the
binary defect itself: in the \(A\otimes(BC)\) bracketing that defect
is diagonal.  The issue is what remains after tracing \(C\) while
retaining \(A\) and \(B\).

This version first performs the prescribed five-version companion
audit.  It then writes the remaining square as an exact
operator-valued conditional law on the \(AB\)-ancestry space.  The
law is a positive operator-valued measure built from complete Racah
matrices.  This identifies exactly where a dimension debit can occur
and proves that normalization alone cannot supply it.

\begin{firewall}[Scope of V70]
\label{fire:V70-scope}
No Standard Model or Navier--Stokes estimate is imported into the
Yang--Mills argument.  Their notes are used only to choose the order
of exact reductions.  The Racah POVM formula below is a
finite-dimensional \(\mathrm{SU}(3)\) identity proved directly in
this appendix.
\end{firewall}

\section{Mandatory V70 companion audit}
\label{sec:V70-audit}

Two consecutive byte reads gave the following stable snapshots:
\begin{center}
\begin{tabular}{|p{0.43\textwidth}|r|r|p{0.27\textwidth}|}
\hline
Companion file & Lines & Bytes & SHA--256\\
\hline
\texttt{Renewal\_SM\_Einstein\_Continuum\_Research\_Notes\_V78.tex}
& \(63389\) & \(2290762\) &
\texttt{0580b36579e4cd3a7221f57dd357f64de07af97b352f91b8d415ddf1edda8dd5}\\
\hline
\texttt{Renewal\_NS\_Stability\_Research\_Notes\_V31.tex}
& \(23052\) & \(887537\) &
\texttt{f1121b62238ad5491bb7cf03635ea9934942edf573ed8faaa9ee79e1d38a6696}\\
\hline
\end{tabular}
\end{center}

The Standard Model--Einstein companion has advanced from the V65
snapshot V75 to V78.  Its V77--V78 calculation proves that a
prematurely minimalized corner quotient removes a nonzero physical
charge; the complete side, cap, corner, and continuum endpoint data
must be retained until the oriented two-form is assembled.  The
Navier--Stokes companion is byte-identical to the V65 snapshot.  Its
terminal V31 calculation proves that flat dyadic marking is exact,
but unmarking restores the critical first moment unless the recent
entry, signed triad work, and parent tail are kept correlated.

\begin{assessment}[The only transferable conclusion]
\label{ass:V70-audit-transfer}
The shared lesson is an order-of-operations rule.  In the present
Yang--Mills problem one must keep the \(BC\) intermediate, final
output, and Racah ancestry correlated until after the \(C\)-partial
trace and the pair response have been assembled.  Neither companion
proves a Yang--Mills inequality, an \(\mathrm{SU}(3)\) Racah estimate,
or a mass-gap statement.
\end{assessment}

\begin{firewall}[Companion separation]
\label{fire:V70-audit-separation}
The companion snapshots, terminology, and methodological comparison
are not hypotheses in any theorem below.  Every displayed
Yang--Mills identity is derived from unitary recoupling, Schur's
lemma, the tensor-dimension identity, or a previously proved
Yang--Mills estimate cited explicitly.
\end{firewall}

\section{The two path bases and their Racah unitary}
\label{sec:V70-path-bases}

Fix irreducibles \(A,B,C,T\).  In the \(AB\)-first bracketing write a
path label as
\[
 i=(S,\alpha,\beta),\qquad
 1\le\alpha\le N_{AB}^{S},\quad
 1\le\beta\le N_{SC}^{T}.
\]
In the \(BC\)-first bracketing write
\[
 j=(V,\eta,\gamma),\qquad
 1\le\eta\le N_{BC}^{V},\quad
 1\le\gamma\le N_{AV}^{T}.
\]
Choose orthonormal intertwiner bases in both bracketings.  Their
change-of-basis matrix
\begin{equation}
 U^T_{S\alpha\beta;\,V\eta\gamma}
 :=
 \left\langle
 L^{AB,T}_{S\alpha\beta},
 L^{BC,T}_{V\eta\gamma}
 \right\rangle_{\operatorname{Hom}_G(H_T,H_A\otimes H_B\otimes H_C)}
 \label{eq:V70-Racah-unitary}
\end{equation}
is unitary, including every Littlewood--Richardson copy.  Thus
\begin{align}
 \sum_{V,\eta,\gamma}
 U^T_{S\alpha\beta;\,V\eta\gamma}
 \overline{U^T_{S'\alpha'\beta';\,V\eta\gamma}}
 &=
 \delta_{SS'}\delta_{\alpha\alpha'}\delta_{\beta\beta'},
 \label{eq:V70-Racah-row-unitarity}\\
 \sum_{S,\alpha,\beta}
 \overline{U^T_{S\alpha\beta;\,V\eta\gamma}}
 U^T_{S\alpha\beta;\,V'\eta'\gamma'}
 &=
 \delta_{VV'}\delta_{\eta\eta'}\delta_{\gamma\gamma'}.
 \label{eq:V70-Racah-column-unitarity}
\end{align}

Put
\begin{equation}
 \delta_{A,V}(T):=
 q_{\alpha,T}-q_{\alpha,A}q_{\alpha,V},
 \qquad
 f_{A,V}(T):=
 \bigl(1+s_\alpha\Xi(T)\bigr)^2
 \delta_{A,V}(T)^2.
 \label{eq:V70-composite-symbol}
\end{equation}
In the \(BC\)-first basis,
\(D_{A\mid BC}^{T}\) is diagonal with eigenvalue
\(\delta_{A,V}(T)\), independently of \(\eta,\gamma\).

\section{Exact composite-cut conditional POVM}
\label{sec:V70-composite-POVM}

For fixed \(S\), define the matrix on
\(\mathbb C^{N_{AB}^{S}}\)
\begin{align}
 \Pi_S(V,T)_{\alpha\alpha'}
 &:=
 \frac{d_T}{d_Sd_C}
 \sum_{\beta=1}^{N_{SC}^{T}}
 \sum_{\eta=1}^{N_{BC}^{V}}
 \sum_{\gamma=1}^{N_{AV}^{T}}
 U^T_{S\alpha\beta;\,V\eta\gamma}
 \overline{
 U^T_{S\alpha'\beta;\,V\eta\gamma}}.
 \label{eq:V70-POVM-atom}
\end{align}
An impossible fusion channel contributes the zero matrix.

\begin{theorem}[The Racah atoms form a normalized POVM]
\label{thm:V70-POVM-normalization}
For every \(S\),
\begin{equation}
 \Pi_S(V,T)\ge0,
 \qquad
 \boxed{\sum_{V,T}\Pi_S(V,T)
 =I_{\mathbb C^{N_{AB}^{S}}}.}
 \label{eq:V70-POVM-normalization}
\end{equation}
\end{theorem}

\begin{proof}
For fixed \(S,V,T,\beta,\eta,\gamma\), the summand in
\eqref{eq:V70-POVM-atom} is the rank-one matrix
\[
 u\,u^*,\qquad
 u_\alpha=
 U^T_{S\alpha\beta;\,V\eta\gamma}.
\]
All coefficients are nonnegative, proving positivity.  Sum first in
\(V,\eta,\gamma\).  Row unitarity
\eqref{eq:V70-Racah-row-unitarity} gives
\[
 \sum_{V,\eta,\gamma}
 U^T_{S\alpha\beta;\,V\eta\gamma}
 \overline{U^T_{S\alpha'\beta;\,V\eta\gamma}}
 =\delta_{\alpha\alpha'}.
\]
There are \(N_{SC}^{T}\) values of \(\beta\).  Hence
\[
 \sum_{V,T}\Pi_S(V,T)
 =
 \left[
 \frac1{d_Sd_C}\sum_Td_TN_{SC}^{T}
 \right]I.
\]
The tensor-dimension identity
\(\sum_TN_{SC}^{T}d_T=d_Sd_C\) proves the claim.
\end{proof}

\begin{theorem}[Exact two-stage partial-trace formula]
\label{thm:V70-composite-partial-trace}
The unresolved V69 square function is exactly
\begin{equation}
 \boxed{
 \mathscr D_{A\mid BC;S}
 =
 d_C\sum_{V,T}f_{A,V}(T)\Pi_S(V,T).}
 \label{eq:V70-composite-partial-trace}
\end{equation}
Equivalently, if
\[
 G_{A\mid BC}:=
 \bigoplus_T
 I_{H_T}\otimes
 \bigl(1+s_\alpha\Xi(T)\bigr)^2
 (D_{A\mid BC}^{T})^2,
\]
then
\begin{equation}
 \operatorname{Tr}_C G_{A\mid BC}
 =
 \bigoplus_S
 I_{H_S}\otimes
 \mathscr D_{A\mid BC;S}.
 \label{eq:V70-global-partial-trace}
\end{equation}
\end{theorem}

\begin{proof}
In the \(BC\)-first basis the multiplicity operator in the
\(T\)-block is
\[
 \operatorname{diag}_{V,\eta,\gamma}
 f_{A,V}(T).
\]
Conjugation by the unitary \eqref{eq:V70-Racah-unitary} shows that its
\((S,\alpha,\beta)\), \((S,\alpha',\beta)\) principal block is
\[
 \sum_{V,\eta,\gamma}
 f_{A,V}(T)
 U^T_{S\alpha\beta;\,V\eta\gamma}
 \overline{U^T_{S\alpha'\beta;\,V\eta\gamma}}.
\]
Insert this expression into the definition
\eqref{eq:V69-composite-square-function}.  The coefficient
\(d_T/d_S\) there is \(d_C\) times the coefficient in
\eqref{eq:V70-POVM-atom}, proving
\eqref{eq:V70-composite-partial-trace}.

Equation \eqref{eq:V70-global-partial-trace} is the same calculation
before taking an \(S\)-block.  More explicitly, apply the exact
matrix partial-trace theorem
\eqref{eq:V67-matrix-partial-trace} to the positive weights
\((1+s_\alpha\Xi(T))^2(D_{A\mid BC}^{T})^2\).
\end{proof}

\begin{corollary}[Exact trace-average law]
\label{cor:V70-trace-average}
One has
\begin{equation}
 \boxed{
 \sum_Sd_S\operatorname{Tr}_{\mathbb C^{N_{AB}^{S}}}
 \mathscr D_{A\mid BC;S}
 =
 \sum_{V,T}
 N_{BC}^{V}N_{AV}^{T}d_T f_{A,V}(T).}
 \label{eq:V70-trace-average}
\end{equation}
Consequently
\[
 \frac1{d_Ad_Bd_C}
 \sum_{V,T}N_{BC}^{V}N_{AV}^{T}d_T f_{A,V}(T)
\]
is the exact dimension-law mean of the squared weighted composite
defect.
\end{corollary}

\begin{proof}
Take traces in \eqref{eq:V70-composite-partial-trace}, multiply by
\(d_S\), and sum \(S\).  After inserting
\eqref{eq:V70-POVM-atom}, column unitarity
\eqref{eq:V70-Racah-column-unitarity} makes the sum over
\(S,\alpha,\beta\) equal to one for each \(V,\eta,\gamma,T\).
The remaining \(\eta,\gamma\) counts are
\(N_{BC}^{V}N_{AV}^{T}\).  Finally,
\[
 \sum_{V,T}N_{BC}^{V}N_{AV}^{T}d_T
 =d_Ad_Bd_C,
\]
the dimension of the full tensor product.
\end{proof}

\section{A valid scalar envelope and its exact loss}
\label{sec:V70-scalar-envelope}

Set \(A_R:=1+C_2(R)\).  Fusion geometry and the V46 binary-defect
estimates imply
\begin{align}
 1+s_\alpha\Xi(T)
 &\le C[1+s_\alpha(A_A+A_V)],
 \label{eq:V70-output-fusion}\\
 |\delta_{A,V}(T)|
 &\le C\min\{s_\alpha A_AA_V,1\},
 \label{eq:V70-binary-cap}\\
 (1+s_\alpha\Xi(T))
 |\delta_{A,V}(T)|
 &\le C.
 \label{eq:V70-weighted-defect-cap}
\end{align}
The last line uses
\eqref{eq:V46-defect-first-moment} together with the unweighted
defect cap.  Define
\begin{equation}
 g_s(A,V):=
 \min\left\{
 1,\,
 [1+s_\alpha(A_A+A_V)]^2
 \min\{s_\alpha A_AA_V,1\}^2
 \right\}.
 \label{eq:V70-BC-envelope}
\end{equation}
Then \(f_{A,V}(T)\le Cg_s(A,V)\).

Let
\begin{equation}
 \nu_{BC}(V,\eta):=
 \frac{d_V}{d_Bd_C}
 \label{eq:V70-BC-dimension-law}
\end{equation}
be the tensor-dimension probability law, including LR copies.

\begin{theorem}[Schur scalarization after the legitimate maximum]
\label{thm:V70-composite-scalar-envelope}
Uniformly in the retained intermediate \(S\),
\begin{equation}
 \boxed{
 \mathscr D_{A\mid BC;S}
 \le
 Cd_C\,
 \mathbb E_{\nu_{BC}}[g_s(A,V)]\,I
 \le Cd_C I.}
 \label{eq:V70-composite-scalar-envelope}
\end{equation}
In the uniformly cold input regime this gives
\begin{align}
 \mathscr D_{A\mid BC;S}
 &\le
 Cd_C[1+s_\alpha(A_A+A_B+A_C)]^2
 \notag\\[-0.2em]
 &\qquad\cdot
 \min\{s_\alpha A_A(A_B+A_C),1\}^2I.
 \label{eq:V70-composite-cold-envelope}
\end{align}
\end{theorem}

\begin{proof}
For a \(BC\) path \((V,\eta)\), sum the complete \(T,\gamma\)
resolution on \(H_A\otimes H_V\).  Since
\(f_{A,V}(T)\le Cg_s(A,V)\), positivity gives on the full
three-factor space
\begin{equation}
 G_{A\mid BC}
 \le
 CI_A\otimes
 \sum_{V,\eta}g_s(A,V)P_{V,\eta}^{BC}.
 \label{eq:V70-global-BC-domination}
\end{equation}
Trace \(C\).  The \(BC\) operator on the right is invariant, so its
partial trace is scalar on the irreducible \(H_B\).  Total trace
determines that scalar:
\begin{align}
 \operatorname{Tr}_C
 \sum_{V,\eta}g_s(A,V)P_{V,\eta}^{BC}
 &=
 \left[
 \frac1{d_B}
 \sum_{V,\eta}d_Vg_s(A,V)
 \right]I_B
 \notag\\
 &=d_C\mathbb E_{\nu_{BC}}[g_s(A,V)]I_B.
 \label{eq:V70-BC-Schur}
\end{align}
Combine \eqref{eq:V70-global-BC-domination} and
\eqref{eq:V70-BC-Schur}, then compress to an \(AB\) intermediate
\(S\), using \eqref{eq:V70-global-partial-trace}.  This proves the
first bound; \(g_s\le1\) proves the second.

For every \(V\subset B\otimes C\), coordinatewise fusion geometry
gives \(A_V\le C(A_B+A_C)\).  Insert this in
\eqref{eq:V70-BC-envelope} and then in the expectation to obtain
\eqref{eq:V70-composite-cold-envelope}.
\end{proof}

\begin{assessment}[Correction of the V69 diagnosis]
\label{ass:V70-V69-correction}
The mismatch noted at the end of V69 does not prevent all
scalarization: after taking the legitimate maximum over the complete
\(A\otimes V\) output resolution, Schur's lemma on \(B\) proves
\eqref{eq:V70-composite-scalar-envelope}.  What the mismatch prevents
is a scalarization that retains the \(T\)-dependent low-output
rarity.  That information lives in the exact POVM
\eqref{eq:V70-composite-partial-trace} and is lost in
\eqref{eq:V70-global-BC-domination}.
\end{assessment}

\section{Insertion into the V68 matrix gate}
\label{sec:V70-gate-insertion}

Define the normalized conditional square
\begin{equation}
 \Phi_{ab}(A,B,C):=
 \max_S
 \left\|
 \sum_{V,T}f_{A,V}(T)\Pi_S(V,T)
 \right\|_{\rm op}.
 \label{eq:V70-Phi-ab}
\end{equation}
Then
\begin{equation}
 \max_S\|\mathscr D_{A\mid BC;S}\|_{\rm op}
 =d_C\Phi_{ab}(A,B,C),
 \qquad
 0\le\Phi_{ab}\le C.
 \label{eq:V70-Phi-identity}
\end{equation}
Combining V69 with V70 gives the explicit estimate
\begin{align}
 \Sigma_{ab}
 &\le
 C d_C\Phi_{ab}
 \notag\\
 &\quad+
 Cd_Cq_{\alpha,B}^{\,2}
 [1+s_\alpha(A_A+A_B+A_C)]^2
 \min\{s_\alpha A_AA_C,1\}^2
 \notag\\
 &\quad+
 Cd_Cq_{\alpha,C}^{\,2}
 \max_S[1+s_\alpha(A_S+A_C)]^2
 \min\{s_\alpha A_AA_B,1\}^2.
 \label{eq:V70-Sigma-decomposition}
\end{align}
The maximum in the last line is restricted to
\(S\subset A\otimes B\), and hence
\(A_S\le C(A_A+A_B)\).

\begin{corollary}[The exact composite contribution to the gate]
\label{cor:V70-composite-gate-term}
The composite square contributes to the \(ab\)-oriented term of
\eqref{eq:V68-square-function-gate} by at most
\begin{equation}
 \boxed{
 CP_{ab}d_C\sqrt{\Phi_{ab}(A,B,C)}.}
 \label{eq:V70-composite-gate-term}
\end{equation}
The complete \(ab\) bound is obtained by adding inside the square
root the two explicit contributions in
\eqref{eq:V70-Sigma-decomposition}; no switch to the \(ac\)
orientation is permitted inside one coordinate or one ancestry
branch.
\end{corollary}

\begin{proof}
The composite part of
\(\sqrt{d_C\Sigma_{ab}}\) is, by
\eqref{eq:V70-Phi-identity},
\[
 \sqrt{d_C\cdot Cd_C\Phi_{ab}}
 =C d_C\sqrt{\Phi_{ab}}.
\]
Multiply by \(P_{ab}\).  The orientation statement is built into
the definitions \eqref{eq:V68-Pab}--\eqref{eq:V68-Sigmaab}; the
minimum over \(ab\) and \(ac\) is taken only after each complete
oriented expression has been formed.
\end{proof}

\section{Sharp diagnostics: what normalization cannot prove}
\label{sec:V70-diagnostics}

\begin{proposition}[POVM normalization has the constant-symbol
extremizer]
\label{prop:V70-POVM-no-gain}
For an abstract nonnegative symbol
\(f_{A,V}(T)\equiv a^2\),
\begin{equation}
 \mathscr D_{A\mid BC;S}=d_Ca^2I,
 \qquad
 \Phi_{ab}=a^2.
 \label{eq:V70-constant-symbol}
\end{equation}
Therefore no argument using only positivity, the pointwise bound
\(f\le C\), and POVM normalization can improve the natural
\(d_C\) size in \eqref{eq:V70-composite-scalar-envelope}.
\end{proposition}

\begin{proof}
Substitute the constant symbol in
\eqref{eq:V70-composite-partial-trace} and use
\eqref{eq:V70-POVM-normalization}.  The assertion about proof
methods follows because the constant family satisfies all three
listed properties and attains the claimed size.
\end{proof}

\begin{assessment}[Dual mesoscopic and alternating-height tests]
\label{ass:V70-regression-tests}
On the dual boundary ray \(C=R_n\) or \(Q_n\),
\(d_C^{1/2}\asymp n\).  The universal estimate
\(\Phi_{ab}\le C\) leaves
\eqref{eq:V70-composite-gate-term} at its natural
\(P_{ab}d_C\) size and supplies none of the inverse-\(n\) companion
debit isolated in V65 and V68.  Thus
\eqref{eq:V70-composite-scalar-envelope} is a correct finite bound
but is not a closure mechanism in the mesoscopic regression.

For the V60 alternating-height array, the formula has a second
advantage: it makes the forbidden orientation switch visible.  One
must estimate either the complete \(ab\) family
\(\{\Pi_S(V,T)\}\) and its \(P_{ab}\), or the complete \(ac\)
analogue and its \(P_{ac}\).  Coordinatewise minima of the two
POVMs do not define either partial trace and cannot be inserted into
\eqref{eq:V68-square-function-gate}.
\end{assessment}

\begin{firewall}[Scope of the regression]
\label{fire:V70-regression-scope}
The constant-symbol example is an abstract sharpness test for the
POVM axioms, not a claim that the actual Yang--Mills symbol is
constant.  Conversely, the fact that the actual symbol is a binary
heat defect does not by itself prove the required conditional
smallness.  A new estimate must use that exact symbol together with
the actual Racah atoms or the upstream pair response.
\end{firewall}

\section{V70 research ledger and next acquisition}
\label{sec:V70-ledger}

\begin{theorem}[Integrated V70 statement]
\label{thm:V70-integrated}
Preserving V1--V69, V70 proves:
\begin{enumerate}[label=\textup{(V70.\arabic*)},leftmargin=5.7em]
\item the stable V70 SM/NS companion audit and its strict
      nontransfer firewall;
\item the normalized positive Racah law
      \eqref{eq:V70-POVM-normalization};
\item the exact two-stage composite-cut partial trace
      \eqref{eq:V70-composite-partial-trace};
\item the exact global trace-average identity
      \eqref{eq:V70-trace-average};
\item the valid Schur scalar envelope
      \eqref{eq:V70-composite-scalar-envelope} and its cold-input
      refinement;
\item the exact remaining conditional factor
      \(\Phi_{ab}\) and its insertion
      \eqref{eq:V70-composite-gate-term}; and
\item the sharp fact that positivity and POVM normalization alone
      cannot pay the required companion dimension.
\end{enumerate}
\end{theorem}

\begin{proof}
Items \textup{(V70.2)--(V70.4)} are
\Cref{thm:V70-POVM-normalization,thm:V70-composite-partial-trace,
cor:V70-trace-average}.  Item \textup{(V70.5)} is
\Cref{thm:V70-composite-scalar-envelope}.  Item
\textup{(V70.6)} is \Cref{cor:V70-composite-gate-term}, and item
\textup{(V70.7)} is \Cref{prop:V70-POVM-no-gain}.
\end{proof}

\begin{openproblem}[V71 mandate: combine the pair tax with the POVM]
\label{op:V71-mandate}
Do not seek a standalone estimate from POVM normalization.  Starting
from the exact oriented quantity
\[
 P_{ab}d_C
 \max_S\left\|
 \sum_{V,T}f_{A,V}(T)\Pi_S(V,T)
 \right\|_{\rm op}^{1/2},
\]
perform the following upstream attack:
\begin{enumerate}
\item insert the actual pair partial-trace weights defining \(P_{ab}\)
      before taking the ancestry maximum;
\item retain the low-\(T\) part of
      \(q_{\alpha,T}-q_{\alpha,A}q_{\alpha,V}\) together with the
      \(BC\) intermediate rather than replacing it by
      \(\sup_T f_{A,V}(T)\);
\item derive the exact marginal atom for the singlet and bounded
      output windows, including all dimension and LR factors;
\item prove either an inverse-\(d_C\) conditional second-moment bound
      for the bad window, or a combined pair--POVM inequality whose
      square root supplies the inverse-\(d_C^{1/2}\) companion;
\item test the result on the dual mesoscopic family and retain one
      fixed orientation on the alternating-height family.
\end{enumerate}
If the conditional atom can concentrate on one ancestry eigendirection,
move further upstream to the untraced row-factorized four-input
operator; do not average the concentration away.
\end{openproblem}

\section{Firewall after V70}
\label{sec:V70-firewall}

\begin{firewall}[What remains open]
\label{fire:V70-final}
V70 exactly identifies the positive conditional law governing the
last V69 square and supplies a correct scalar envelope.  It does not
prove the inverse-dimension conditional estimate required by the
matrix collision gate.  Therefore JREC, simultaneous sideways
closure, puncture aggregation, multi-collision summation, the full
\([3,2]\), \([3,3]\), and \([4]\) sectors, construction of the
cofinal Yang--Mills transfer operator, and a positive continuum
spectral gap all remain open.
\end{firewall}

\section{Append-only record for V70}
\label{sec:V70-record}

V70 was appended to the complete V69 manuscript.  The exact V69
parent had \(50642\) lines, \(1830131\) bytes, and SHA--256 digest
\[
 \texttt{c00aa1ece0621ca26568a2fb7c845787d8231f577bc16536278a6a4c8cc15d69}.
\]
No V69 mathematical content was changed.  The sole final
\verb|\end{document}| is restored below.

% =============================================================================
% END APPENDED VERSION V70 MATERIAL
% =============================================================================
% =============================================================================
% BEGIN APPENDED VERSION V71 MATERIAL
% =============================================================================

\chapter{V71: Exact output marginal, singlet debit, and the mesoscopic
conditional barrier}
\label{ch:V71}

\section{Purpose}
\label{sec:V71-purpose}

V70 represented the remaining composite square as a normalized Racah
POVM.  The first task in V71 is to compute its output marginal before
estimating the binary symbol.  That marginal is exactly scalar, even
though the individual \(BC\)-intermediate atoms are matrices.  It
gives an inverse-\(d_C\) bound for every fixed finite output window
and an exact inverse-\(d_C\) singlet contribution.

The second task is diagnostic.  Fourth Wilson-multiplier damping
controls the remote output tail, but the dual-symmetric
\(k\asymp\sqrt n\) shell has unit cumulative Schur mass.  Thus the
singlet is not the obstruction: the remaining loss is a
thermal-mesoscopic conditional average, which must be combined with
the upstream pair tax.

\begin{firewall}[Scope of V71]
\label{fire:V71-scope}
All results below concern the single composite square
\(\mathscr D_{A\mid BC;S}\) in one fixed \(ab\) orientation.  The
finite-window bounds do not yet prove the complete matrix gate,
because their constants grow with the window and the thermal
mesoscopic window grows with \(s_\alpha^{-1}\).
\end{firewall}

\section{The scalar output marginal}
\label{sec:V71-output-marginal}

For the V70 Racah POVM, define
\begin{equation}
 p_{S,C}(T):=
 \frac{d_TN_{SC}^{T}}{d_Sd_C}.
 \label{eq:V71-output-probability}
\end{equation}
The tensor-dimension identity says
\(\sum_Tp_{S,C}(T)=1\).

\begin{theorem}[Exact output marginal of the Racah POVM]
\label{thm:V71-output-marginal}
For every \(S,T\),
\begin{equation}
 \boxed{
 \sum_V\Pi_S(V,T)
 =
 p_{S,C}(T)
 I_{\mathbb C^{N_{AB}^{S}}}.}
 \label{eq:V71-output-marginal}
\end{equation}
Consequently, for every nonnegative scalar function \(h(T)\),
\begin{equation}
 \boxed{
 d_C\sum_{V,T}h(T)\Pi_S(V,T)
 =
 \left[
 \frac1{d_S}\sum_Td_TN_{SC}^{T}h(T)
 \right]I.}
 \label{eq:V71-output-functional}
\end{equation}
\end{theorem}

\begin{proof}
Insert \eqref{eq:V70-POVM-atom} and sum
\(V,\eta,\gamma\).  For every fixed \(T,S,\alpha,\alpha',\beta\),
row unitarity \eqref{eq:V70-Racah-row-unitarity} gives
\[
 \sum_{V,\eta,\gamma}
 U^T_{S\alpha\beta;\,V\eta\gamma}
 \overline{U^T_{S\alpha'\beta;\,V\eta\gamma}}
 =\delta_{\alpha\alpha'}.
\]
The sum over \(\beta\) contributes \(N_{SC}^{T}\), proving
\eqref{eq:V71-output-marginal}.  Multiply by \(h(T)\), sum \(T\),
and multiply by \(d_C\) to obtain
\eqref{eq:V71-output-functional}.
\end{proof}

Whenever \(p_{S,C}(T)>0\), define the conditional POVM
\begin{equation}
 \widehat\Pi_S(V\mid T):=
 \frac{\Pi_S(V,T)}{p_{S,C}(T)}.
 \label{eq:V71-conditional-POVM}
\end{equation}
Equation \eqref{eq:V71-output-marginal} gives
\(\sum_V\widehat\Pi_S(V\mid T)=I\).  Put
\begin{equation}
 \widehat F_S(T):=
 \sum_V f_{A,V}(T)\widehat\Pi_S(V\mid T).
 \label{eq:V71-conditional-symbol}
\end{equation}

\begin{corollary}[Exact output--conditional factorization]
\label{cor:V71-output-conditional}
The composite square has the exact positive decomposition
\begin{equation}
 \boxed{
 \mathscr D_{A\mid BC;S}
 =
 d_C\sum_Tp_{S,C}(T)\widehat F_S(T).}
 \label{eq:V71-output-conditional}
\end{equation}
Every \(\widehat F_S(T)\) is positive and
\[
 \|\widehat F_S(T)\|_{\rm op}
 \le\sup_{\{V:N_{BC}^{V}N_{AV}^{T}>0\}}f_{A,V}(T).
\]
\end{corollary}

\begin{proof}
Group \eqref{eq:V70-composite-partial-trace} by \(T\), and use
\eqref{eq:V71-conditional-POVM}.  Positivity and the norm bound
follow from the fact that the conditional atoms sum to the identity.
\end{proof}

\section{The singlet atom is exactly dimension-paid}
\label{sec:V71-singlet}

\begin{theorem}[Exact singlet Racah atom]
\label{thm:V71-singlet-atom}
For final output \(T=\mathbf1\),
\begin{equation}
 \boxed{
 \Pi_S(V,\mathbf1)
 =
 \mathbf1_{\{S=C^*,\,V=A^*\}}
 \frac1{d_C^2}
 I_{\mathbb C^{N_{AB}^{C^*}}}.}
 \label{eq:V71-singlet-atom}
\end{equation}
In particular, the singlet part of the composite square is
\begin{equation}
 \boxed{
 \mathscr D_{A\mid BC;S}^{\mathbf1}
 =
 \mathbf1_{\{S=C^*\}}
 \frac{
 (1-q_{\alpha,A}q_{\alpha,A^*})^2
 }{d_C}
 I_{\mathbb C^{N_{AB}^{C^*}}}.}
 \label{eq:V71-singlet-square}
\end{equation}
For the central inversion-symmetric Wilson density,
\(q_{\alpha,A^*}=q_{\alpha,A}\).
\end{theorem}

\begin{proof}
The trivial representation occurs in \(S\otimes C\) if and only if
\(S=C^*\), and then with multiplicity one.  It occurs in
\(A\otimes V\) if and only if \(V=A^*\), again with multiplicity
one.  Thus every singlet atom other than the one displayed is zero.

Set
\[
 m=N_{AB}^{C^*}=N_{BC}^{A^*}.
\]
The equality is the dimension of the common invariant space
\(\operatorname{Hom}_G(\mathbf1,A\otimes B\otimes C)\).
At \(T=\mathbf1\), the Racah matrix is therefore an \(m\)-by-\(m\)
unitary between the \(\alpha\) and \(\eta\) bases; \(\beta\) and
\(\gamma\) have only one value.  Formula
\eqref{eq:V70-POVM-atom} gives
\[
 \Pi_{C^*}(A^*,\mathbf1)_{\alpha\alpha'}
 =
 \frac1{d_C^2}
 \sum_{\eta=1}^{m}
 U^{\mathbf1}_{\alpha;\eta}
 \overline{U^{\mathbf1}_{\alpha';\eta}}
 =
 \frac{\delta_{\alpha\alpha'}}{d_C^2}.
\]
This proves \eqref{eq:V71-singlet-atom}.

Now \(\Xi(\mathbf1)=0\), \(q_{\alpha,\mathbf1}=1\), and
\(\delta_{A,A^*}(\mathbf1)=1-q_{\alpha,A}q_{\alpha,A^*}\).
Multiply \eqref{eq:V71-singlet-atom} by the outer \(d_C\) in
\eqref{eq:V70-composite-partial-trace} to obtain
\eqref{eq:V71-singlet-square}.  Inversion symmetry of the central
one-link law identifies the two conjugate Fourier multipliers.
\end{proof}

\begin{assessment}[The dual singlet is not the missing square root]
\label{ass:V71-singlet-not-obstruction}
The singlet contribution is \(O(d_C^{-1})\).  After the V68
square-function bridge it costs
\[
 \sqrt{d_C\cdot O(d_C^{-1})}=O(1).
\]
Thus the exact dual singlet already supplies the full companion
dimension.  Any remaining \(d_C^{1/2}\) loss must come from a growing
output window, not from the isolated invariant vector.
\end{assessment}

\section{Finite output windows}
\label{sec:V71-finite-windows}

For a set \(\Omega\) of irreducibles, define
\begin{equation}
 \mathscr D_{A\mid BC;S}^{\Omega}
 :=
 d_C\sum_{T\in\Omega}\sum_V
 f_{A,V}(T)\Pi_S(V,T).
 \label{eq:V71-window-square}
\end{equation}

\begin{theorem}[Uniform fixed-window dimension debit]
\label{thm:V71-window-debit}
For every finite output set \(\Omega\),
\begin{align}
 0\le
 \mathscr D_{A\mid BC;S}^{\Omega}
 &\le
 \frac C{d_S}
 \sum_{T\in\Omega}d_TN_{SC}^{T}I
 \label{eq:V71-window-first}\\
 &\le
 \boxed{
 \frac C{d_C}
 \left(\sum_{T\in\Omega}d_T^2\right)I.}
 \label{eq:V71-window-dimension-debit}
\end{align}
The constants are independent of \(A,B,C,S\) and of all LR
multiplicities.
\end{theorem}

\begin{proof}
The weighted binary-defect bound
\eqref{eq:V70-weighted-defect-cap} gives
\(0\le f_{A,V}(T)\le C\).  Sum the POVM atoms in \(V\) and apply
\eqref{eq:V71-output-marginal}; multiplication by \(d_C\) proves
\eqref{eq:V71-window-first}.

If \(N_{SC}^{T}=0\), its summand vanishes.  Otherwise Frobenius
reciprocity gives
\[
 N_{SC}^{T}=N_{T,S^*}^{C^*}.
\]
Dimension counting in \(T\otimes S^*\) yields
\[
 N_{SC}^{T}d_C\le d_Td_S.
\]
Consequently
\[
 \frac{d_TN_{SC}^{T}}{d_S}
 \le\frac{d_T^2}{d_C}.
\]
Sum this inequality over \(T\in\Omega\) to prove
\eqref{eq:V71-window-dimension-debit}.
\end{proof}

Let
\[
 \Omega_L:=\{T:C_2(T)\le L\}.
\]

\begin{corollary}[Casimir-ball debit for \(\mathrm{SU}(3)\)]
\label{cor:V71-Casimir-ball}
For every \(L\ge0\),
\begin{equation}
 \boxed{
 \mathscr D_{A\mid BC;S}^{\Omega_L}
 \le
 \frac{C(1+L)^4}{d_C}I.}
 \label{eq:V71-Casimir-ball}
\end{equation}
\end{corollary}

\begin{proof}
Write an \(\mathrm{SU}(3)\) highest weight as \((p,q)\).
The Weyl formulas give
\[
 C_2(p,q)\asymp p^2+pq+q^2+p+q,
 \qquad
 d_{p,q}=\frac{(p+1)(q+1)(p+q+2)}2.
\]
The ball \(C_2(p,q)\le L\) contains \(O(1+L)\) lattice points, and
on it \(d_{p,q}^2\le C(1+L)^3\).  Hence
\[
 \sum_{T\in\Omega_L}d_T^2\le C(1+L)^4.
\]
Apply \eqref{eq:V71-window-dimension-debit}.
\end{proof}

\section{Fourth-moment control of the remote output tail}
\label{sec:V71-remote-tail}

\begin{lemma}[Sixth-power decay of the squared weighted defect]
\label{lem:V71-weighted-defect-tail}
For all admissible binary channels \(T\subset A\otimes V\),
\begin{equation}
 \boxed{
 f_{A,V}(T)
 \le
 \frac C{[1+s_\alpha\Xi(T)]^6}.}
 \label{eq:V71-weighted-defect-tail}
\end{equation}
\end{lemma}

\begin{proof}
The fourth product multiplier moment
\eqref{eq:V43-balanced-fourth-moment}, with the bounded region
absorbed into the constant, gives
\begin{equation}
 q_{\alpha,R}
 \le\frac C{[1+s_\alpha\Xi(R)]^4}
 \label{eq:V71-fourth-pointwise}
\end{equation}
for every product channel \(R\).  The same order-one through
order-four moment bounds, or simply multiplication of
\eqref{eq:V71-fourth-pointwise} for \(A\) and \(V\), give
\[
 q_{\alpha,A}q_{\alpha,V}
 \le
 \frac C{
 [1+s_\alpha(\Xi(A)+\Xi(V))]^4}.
\]
Fusion geometry implies
\[
 1+s_\alpha\Xi(T)
 \le C[1+s_\alpha(\Xi(A)+\Xi(V))].
\]
Therefore both \(q_{\alpha,T}\) and
\(q_{\alpha,A}q_{\alpha,V}\) are bounded by
\(C[1+s_\alpha\Xi(T)]^{-4}\).  Their difference obeys the same
bound.  Multiplication by \(1+s_\alpha\Xi(T)\) and squaring prove
\eqref{eq:V71-weighted-defect-tail}.
\end{proof}

\begin{theorem}[Low-output debit plus remote-tail bound]
\label{thm:V71-low-high-split}
For every \(L\ge0\),
\begin{equation}
 \boxed{
 \mathscr D_{A\mid BC;S}
 \le
 C\left[
 \frac{(1+L)^4}{d_C}
 +
 \frac{d_C}{(1+s_\alpha L)^6}
 \right]I.}
 \label{eq:V71-low-high-split}
\end{equation}
In the regime \(d_C\ge s_\alpha^{-2}\), optimization at
\(L=s_\alpha^{-3/5}d_C^{1/5}\) gives
\begin{equation}
 \mathscr D_{A\mid BC;S}
 \le
 C s_\alpha^{-12/5}d_C^{-1/5}I.
 \label{eq:V71-optimized-tail}
\end{equation}
For all \(d_C\), this may be combined with the V70 ceiling as
\[
 \mathscr D_{A\mid BC;S}
 \le
 C\min\{d_C,\,
 s_\alpha^{-12/5}d_C^{-1/5}\}I
\]
whenever the second entry is invoked in its stated regime.
\end{theorem}

\begin{proof}
The contribution of \(C_2(T)\le L\) is
\eqref{eq:V71-Casimir-ball}.  On the complement,
\(\Xi(T)\ge c(1+L)\) up to harmless conventions at the trivial
representation.  Apply \eqref{eq:V71-weighted-defect-tail} in
\eqref{eq:V70-composite-partial-trace}, then sum all remaining POVM
atoms and use \eqref{eq:V70-POVM-normalization}.  This gives the
second term in \eqref{eq:V71-low-high-split}.

If \(d_C\ge s_\alpha^{-2}\), then
\(s_\alpha L=s_\alpha^{2/5}d_C^{1/5}\ge1\).
The two terms at the displayed choice of \(L\) are each bounded by
\[
 C s_\alpha^{-12/5}d_C^{-1/5}.
\]
This proves \eqref{eq:V71-optimized-tail}.  The V70 ceiling
\eqref{eq:V70-composite-scalar-envelope} supplies the other entry.
\end{proof}

\begin{assessment}[Why the fourth-moment split is not closure]
\label{ass:V71-tail-not-closure}
At the critical relation \(d_C\asymp s_\alpha^{-2}\), the optimized
right side of \eqref{eq:V71-optimized-tail} is
\(O(s_\alpha^{-2})=O(d_C)\), exactly the natural ceiling.  Thus the
remote tail improves genuinely higher-dimensional regimes, but it
does not pay the thermal mesoscopic boundary where the known
square-root loss occurs.
\end{assessment}

\section{Exact dual-symmetric regression}
\label{sec:V71-dual-regression}

Let
\[
 R_n=(n,0),\qquad Q_n=(0,n),\qquad
 d_n=d_{R_n}=d_{Q_n}=\frac{(n+1)(n+2)}2.
\]
The multiplicity-free decomposition
\[
 R_n\otimes Q_n=\bigoplus_{k=0}^{n}T_k,
 \qquad T_k=(k,k),\qquad d_{T_k}=(k+1)^3
\]
was used in V64--V66.

\begin{proposition}[The POVM marginal reproduces the mesoscopic shell]
\label{prop:V71-dual-marginal}
For \(S=R_n\), \(C=Q_n\),
\begin{equation}
 \sum_V\Pi_{R_n}(V,T_k)
 =
 \frac{(k+1)^3}{d_n^2}I.
 \label{eq:V71-dual-output-atom}
\end{equation}
Hence, using only the uniform symbol cap,
\begin{equation}
 \mathscr D_{A\mid BC;R_n}^{\{T_0,\ldots,T_K\}}
 \le
 \frac C{d_n}\sum_{k=0}^{K}(k+1)^3I
 =
 C D_n(K)I,
 \label{eq:V71-dual-window}
\end{equation}
where \(D_n(K)\) is exactly the V66 cumulative Schur mass
\eqref{eq:V66-dual-cumulative-Schur-mass}.  In particular,
\[
 D_n(\lfloor c\sqrt n\rfloor)\asymp1
\]
for fixed suitable \(c>0\), whereas the isolated \(k=0\) term is
\(d_n^{-1}\).
\end{proposition}

\begin{proof}
Equation \eqref{eq:V71-dual-output-atom} is
\eqref{eq:V71-output-marginal} with
\(N_{R_nQ_n}^{T_k}=1\) and
\(d_S=d_C=d_n\).  Multiply the sum from \(k=0\) to \(K\) by the
outer \(d_C=d_n\) and the uniform bound \(f\le C\) to obtain
\eqref{eq:V71-dual-window}.  Finally,
\[
 \sum_{k=0}^{K}(k+1)^3\asymp(K+1)^4,
 \qquad d_n\asymp n^2,
\]
so \(K\asymp\sqrt n\) gives order one.  The \(K=0\) value is
\(1/d_n\), consistently with
\eqref{eq:V71-singlet-square}.
\end{proof}

\begin{assessment}[The exact remaining conditional quantity]
\label{ass:V71-conditional-barrier}
The dual marginal shows precisely why low-output counting alone
stops one square root short.  On \(k\lesssim\sqrt n\), the scalar
weights \(d_np_{R_n,Q_n}(T_k)\) have total mass of order one.  To
reduce the composite square from order one to order \(d_n^{-1}\),
one would need the conditional matrices
\[
 \widehat F_{R_n}(T_k)
 =
 \sum_V f_{A,V}(T_k)
 \widehat\Pi_{R_n}(V\mid T_k)
\]
to gain order \(d_n^{-1}\) on average.  No theorem through V71 proves
that gain, and the available Wilson bounds permit thermal outputs
with \(s_\alpha C_2(T_k)\asymp1\).  The correct next object is
therefore the product of this conditional average with the actual
pair response \(P_{ab}\), not a further output-marginal estimate.
\end{assessment}

\section{V71 research ledger and next acquisition}
\label{sec:V71-ledger}

\begin{theorem}[Integrated V71 statement]
\label{thm:V71-integrated}
Preserving V1--V70, V71 proves:
\begin{enumerate}[label=\textup{(V71.\arabic*)},leftmargin=5.7em]
\item the exact scalar output marginal
      \eqref{eq:V71-output-marginal} and output--conditional
      factorization \eqref{eq:V71-output-conditional};
\item the exact singlet atom and inverse-\(d_C\) singlet square
      \eqref{eq:V71-singlet-square};
\item the inverse-\(d_C\) fixed-window estimate
      \eqref{eq:V71-window-dimension-debit};
\item the \(\mathrm{SU}(3)\) Casimir-ball estimate
      \eqref{eq:V71-Casimir-ball};
\item the sixth-power remote-output bound obtained from the fourth
      Wilson moment;
\item the optimized low/high estimate
      \eqref{eq:V71-optimized-tail}; and
\item the exact recovery of the V66 unit-mass mesoscopic shell from
      the Racah output marginal.
\end{enumerate}
\end{theorem}

\begin{proof}
The cited results are respectively
\Cref{thm:V71-output-marginal,cor:V71-output-conditional,
thm:V71-singlet-atom,thm:V71-window-debit,
cor:V71-Casimir-ball,lem:V71-weighted-defect-tail,
thm:V71-low-high-split,prop:V71-dual-marginal}.
\end{proof}

\begin{openproblem}[V72 mandate: the thermal conditional pair--POVM
estimate]
\label{op:V72-mandate}
Work on the exact \(ab\)-oriented thermal-window quantity
\begin{equation}
 P_{ab}
 \max_S\left\|
 d_C^2
 \sum_{\{T:s_\alpha C_2(T)\asymp1\}}
 p_{S,C}(T)\widehat F_S(T)
 \right\|_{\rm op}^{1/2},
 \label{eq:V72-target-preview}
\end{equation}
which is the corresponding portion of
\(P_{ab}\sqrt{d_C\Sigma_{ab}}\).  The next calculation must:
\begin{enumerate}
\item insert the \((oa)\) and \((bc)\) pair partial-trace weights
      before the maximum over \(S\);
\item on the dual family, compute the pair response and the
      conditional \(V\)-law on
      \(T_k=(k,k)\), \(k\asymp\sqrt n\);
\item determine whether the inverse-\(n\) debit comes from the pair
      response, from the conditional symbol, or from their correlated
      product;
\item retain one complete orientation on the alternating-height
      family;
\item if no such debit is present, exhibit the precise finite
      representation family disproving the standalone matrix gate
      and return to the untraced row-factorized collision operator.
\end{enumerate}
Do not spend another version sharpening fixed output balls: V71
proves that the growing thermal shell, not the fixed ball, is the
remaining bottleneck.
\end{openproblem}

\section{Firewall after V71}
\label{sec:V71-firewall}

\begin{firewall}[Current claim boundary]
\label{fire:V71-final}
V71 pays the singlet and every fixed finite output window, and it
controls the remote output tail.  It does not control the
thermal-mesoscopic conditional matrix, nor its correlated product
with the two pair banks.  The pair-traced matrix gate, JREC,
simultaneous sideways closure, puncture aggregation, multi-collision
summation, the remaining partition sectors, the cofinal transfer
operator, and the continuum Yang--Mills mass gap remain unproved.
\end{firewall}

\section{Append-only record for V71}
\label{sec:V71-record}

V71 was appended to the complete V70 manuscript.  The exact V70
parent had \(51271\) lines, \(1850297\) bytes, and SHA--256 digest
\[
 \texttt{9b14080c3d0eefc60df153d713d6252b75baf0583c0a93d03262e603912c1b89}.
\]
No V70 mathematical content was changed.  The sole final
\verb|\end{document}| is restored below.

% =============================================================================
% END APPENDED VERSION V71 MATERIAL
% =============================================================================
% =============================================================================
% BEGIN APPENDED VERSION V72 MATERIAL
% =============================================================================

\chapter{V72: Product-state trace taxes and an upstream separable
square-function gate}
\label{ch:V72}

\section{Purpose and upstream correction}
\label{sec:V72-purpose}

V71 isolated a thermal mesoscopic square of natural order one.  The
V68 operator-norm bridge then charges \(d_C^{1/2}\).  Before
attempting an entrywise \(\mathrm{SU}(3)\) Racah estimate, one must
check whether that operator norm is required by the original
row-factorized coefficient problem.

It is not.  The proof of the V67 complete matrix tax tests the
positive carrier only on tensor products of one density matrix per
input row.  Replacing those product tests by the operator norm on the
joint retained \(AB\) carrier admits entangled \(A\)-\(B\) vectors
which need not arise from a row-factorized input.  V72 retains the
exact product-state capacity throughout the polar
Cauchy--Schwarz argument.  This gives a strictly weaker sufficient
gate and keeps Clebsch--Gordan entanglement as a possible source of
the missing dimension debit.

\begin{firewall}[Scope of V72]
\label{fire:V72-scope}
The new gate is sufficient for the same row-factorized trace-class
coefficients, with arbitrary within-row spectators.  It is not a
claim for arbitrary trace-class inputs entangled across distinct
rows.  Nor does V72 yet estimate the product-state capacity of the
actual mesoscopic Yang--Mills square.
\end{firewall}

\section{Positive product-state capacity}
\label{sec:V72-product-capacity}

Let
\[
 H=\bigotimes_{i=1}^{m}H_i,\qquad Q\ge0,
\]
and let \(\varnothing\ne I\subseteq[m]\).  Define
\begin{align}
 \kappa_I^\otimes(Q)
 &:=
 \sup_{\substack{\rho_i\ge0,\ \operatorname{Tr}\rho_i=1\\i\in I}}
 \operatorname{Tr}\left[
 Q\left(
 \bigotimes_{i\in I}\rho_i
 \otimes I_{I^c}
 \right)\right]
 \label{eq:V72-product-capacity}\\
 &=
 \sup_{\substack{\rho_i\ge0,\ \operatorname{Tr}\rho_i=1\\i\in I}}
 \operatorname{Tr}\left[
 (\operatorname{Tr}_{I^c}Q)
 \bigotimes_{i\in I}\rho_i
 \right].
 \notag
\end{align}
The placement of tensor factors in the first line is the canonical
one.  Convexity permits the supremum to be restricted to pure
states.  One always has
\begin{equation}
 \kappa_I^\otimes(Q)
 \le
 \|\operatorname{Tr}_{I^c}Q\|_{\rm op},
 \label{eq:V72-product-below-operator}
\end{equation}
with strict inequality possible.

\begin{lemma}[Exact positive product-state tax]
\label{lem:V72-positive-product-tax}
For arbitrary positive trace-class \(\sigma_i\) on \(H_i\),
\begin{equation}
 \boxed{
 \operatorname{Tr}\left[Q\bigotimes_i\sigma_i\right]
 \le
 \kappa_I^\otimes(Q)
 \prod_i\operatorname{Tr}\sigma_i.}
 \label{eq:V72-positive-product-tax}
\end{equation}
\end{lemma}

\begin{proof}
The assertion is immediate if a trace vanishes.  Otherwise normalize
each \(\sigma_i\) to a density matrix.  On the traced factors,
\(0\le\sigma_{I^c}\le I_{I^c}\); hence
\[
 \operatorname{Tr}[Q(\sigma_I\otimes\sigma_{I^c})]
 \le
 \operatorname{Tr}[Q(\sigma_I\otimes I_{I^c})]
 \le\kappa_I^\otimes(Q).
\]
Restore the removed traces.
\end{proof}

\section{Complete spectator tax without cross-row entanglement}
\label{sec:V72-complete-tax}

Use the V67 setup: isometries
\(V_\omega:G_\omega\to H\) have mutually orthogonal ranges,
\(K_\omega=K_\omega^*\), and
\[
 Q_{\rm abs}=\sum_\omega V_\omega|K_\omega|V_\omega^*.
\]
For \(A_i\in\mathcal S_1(H_i\otimes\mathcal K_i)\), let
\(Y_\omega(\bigotimes_iA_i)\) be the spectator-valued coefficient
defined in \eqref{eq:V67-spectator-coefficient}.

\begin{theorem}[Complete product-state matrix tax]
\label{thm:V72-complete-product-tax}
For every nonempty \(I\subseteq[m]\),
\begin{equation}
 \boxed{
 \sum_\omega
 \left\|Y_\omega\!\left(\bigotimes_iA_i\right)\right\|_1
 \le
 \min\left\{
 \|Q_{\rm abs}\|_{\rm op},
 \kappa_I^\otimes(Q_{\rm abs})
 \right\}
 \prod_i\|A_i\|_1.}
 \label{eq:V72-complete-product-tax}
\end{equation}
The constant is one, and the bound is uniform in every spectator
dimension.
\end{theorem}

\begin{proof}
Repeat the trace-norm duality step in
\Cref{thm:V67-complete-matrix-tax}.  For contractions
\(Z_\omega\) on the total spectator space, it produces an operator
\[
 L_Z=
 \sum_\omega
 (V_\omega K_\omega V_\omega^*)\otimes Z_\omega
\]
such that
\[
 |L_Z|,\ |L_Z^*|
 \le Q_{\rm abs}\otimes I_{\mathcal K}.
\]
Take first \(A_i=|x_i\rangle\langle y_i|\).
The polar Cauchy--Schwarz inequality bounds the resulting matrix
coefficient by the geometric mean of an \(x\)-expectation of
\(|L_Z|\) and a \(y\)-expectation of \(|L_Z^*|\).

The carrier reduction of
\(|\bigotimes_ix_i\rangle\langle\bigotimes_ix_i|\) is
\(\bigotimes_i\sigma_i^x\), where
\(\sigma_i^x\ge0\) and
\(\operatorname{Tr}\sigma_i^x=\|x_i\|^2\).  Apply
\eqref{eq:V72-positive-product-tax} to
\(Q_{\rm abs}\) and these reductions.  Do the same for the \(y_i\).
This gives
\[
 |\operatorname{Tr}[L_Z\bigotimes_iA_i]|
 \le
 \kappa_I^\otimes(Q_{\rm abs})
 \prod_i\|x_i\|\,\|y_i\|.
\]
The ordinary operator-norm estimate gives the other entry of the
minimum.  Expand general \(A_i\) in singular values, sum the
nonnegative coefficients, take the supremum over \(Z_\omega\), and
remove the arbitrarily small duality error.
\end{proof}

\begin{remark}[Relation to V67]
\label{rem:V72-relation-V67}
Equation \eqref{eq:V72-product-below-operator} shows that
\eqref{eq:V72-complete-product-tax} strengthens
\eqref{eq:V67-complete-matrix-tax} on the original row-factorized
domain.  If arbitrary cross-row entanglement is admitted, the
operator norm in V67 is the correct replacement; V72 makes no claim
on that larger domain.
\end{remark}

\section{A product-state square-function bridge}
\label{sec:V72-square-bridge}

On \(H_A\otimes H_B\otimes H_C\), let
\begin{equation}
 \mathbf K^\sim
 :=
 \bigoplus_T
 I_{H_T}\otimes\widetilde K_T,
 \qquad
 \widetilde K_T=(1+s_\alpha\Xi(T))K_T.
 \label{eq:V72-global-weighted-K}
\end{equation}
Define positive operators on \(H_A\otimes H_B\)
\begin{equation}
 M_{ab}^{(1)}:=
 \operatorname{Tr}_C|\mathbf K^\sim|,
 \qquad
 M_{ab}^{(2)}:=
 \operatorname{Tr}_C(\mathbf K^\sim)^2.
 \label{eq:V72-first-second-operators}
\end{equation}
By V67--V68,
\begin{align}
 M_{ab}^{(1)}
 &=
 \bigoplus_S I_{H_S}\otimes\mathscr A_S^\sim,
 \label{eq:V72-first-blocks}\\
 M_{ab}^{(2)}
 &=
 \bigoplus_S I_{H_S}\otimes\mathscr S_S^\sim.
 \label{eq:V72-second-blocks}
\end{align}
Put
\begin{equation}
 \Gamma_{ab}^{\otimes}
 :=
 \sup_{\substack{\rho_A,\rho_B\ge0\\
                  \operatorname{Tr}\rho_A=
                  \operatorname{Tr}\rho_B=1}}
 \operatorname{Tr}\left[
 M_{ab}^{(2)}(\rho_A\otimes\rho_B)
 \right].
 \label{eq:V72-Gamma-product}
\end{equation}

\begin{theorem}[Kadison product-state square bridge]
\label{thm:V72-product-square-bridge}
One has the operator inequality
\begin{equation}
 \boxed{
 (M_{ab}^{(1)})^2
 \le d_C M_{ab}^{(2)}}
 \label{eq:V72-Kadison-operator}
\end{equation}
and the product-state consequence
\begin{equation}
 \boxed{
 \sup_{\rho_A,\rho_B}
 \operatorname{Tr}[M_{ab}^{(1)}(\rho_A\otimes\rho_B)]
 \le
 \sqrt{d_C\Gamma_{ab}^{\otimes}}.}
 \label{eq:V72-product-square-bridge}
\end{equation}
\end{theorem}

\begin{proof}
The normalized partial trace
\[
 \Phi(X):=d_C^{-1}\operatorname{Tr}_C X
\]
is unital and completely positive.  Kadison's inequality applied to
the self-adjoint operator \(|\mathbf K^\sim|\) gives
\[
 \Phi(|\mathbf K^\sim|)^2
 \le
 \Phi((\mathbf K^\sim)^2).
\]
Multiply by \(d_C^2\) to obtain
\eqref{eq:V72-Kadison-operator}.

For a product density \(\rho=\rho_A\otimes\rho_B\),
scalar Cauchy--Schwarz for the state
\(\operatorname{Tr}(\rho\,\cdot)\) gives
\[
 \operatorname{Tr}(\rho M_{ab}^{(1)})^2
 \le
 \operatorname{Tr}(\rho(M_{ab}^{(1)})^2)
 \le
 d_C\operatorname{Tr}(\rho M_{ab}^{(2)}).
\]
Take the supremum over product densities.
\end{proof}

\begin{corollary}[Exact Clebsch--Gordan form of the new target]
\label{cor:V72-CG-target}
For unit vectors \(u\in H_A\), \(v\in H_B\), let
\[
 c_S(u,v):=
 P_S^{AB}(u\otimes v)
 \in H_S\otimes\mathbb C^{N_{AB}^{S}}.
\]
Then
\begin{equation}
 \boxed{
 \Gamma_{ab}^{\otimes}
 =
 \sup_{\|u\|=\|v\|=1}
 \sum_S
 \left\langle
 c_S(u,v),
 (I_{H_S}\otimes\mathscr S_S^\sim)c_S(u,v)
 \right\rangle.}
 \label{eq:V72-CG-target}
\end{equation}
Thus V68's
\(\max_S\|\mathscr S_S^\sim\|_{\rm op}\) is replaced by its exact
Clebsch--Gordan average on product vectors.
\end{corollary}

\begin{proof}
Convexity reduces \eqref{eq:V72-Gamma-product} to pure product
states.  Insert the orthogonal block decomposition
\eqref{eq:V72-second-blocks}.  The summands are exactly the displayed
quadratic forms.
\end{proof}

\section{Insertion of the pair banks}
\label{sec:V72-pair-insertion}

Let the two positive weighted pair-bank operators be
\[
 Q_{oa}:=\sum_\lambda a_\lambda P_{oa,\lambda},
 \qquad
 Q_{bc}:=\sum_\mu b_\mu P_{bc,\mu}.
\]
In the collision-core sector, before the still-suspended
same-label puncture banks, the V67 positive carrier is the tensor
product
\begin{equation}
 Q_{ab}^{\rm abs}
 =
 Q_{oa}\otimes Q_{bc}\otimes|\mathbf K^\sim|.
 \label{eq:V72-three-bank-product}
\end{equation}
The relevant pair traces are scalar:
\begin{align}
 \operatorname{Tr}_{o}Q_{oa}
 &=
 \left(\sum_\lambda a_\lambda x_\lambda^a\right)I_a,
 \notag\\
 \operatorname{Tr}_{c}Q_{bc}
 &=
 \left(\sum_\mu b_\mu y_\mu^b\right)I_b.
 \label{eq:V72-pair-scalar-traces}
\end{align}

\begin{theorem}[Exact product-state pair--core factorization]
\label{thm:V72-pair-core-product}
For the \(ab\)-retained row choice,
\begin{equation}
 \boxed{
 \kappa_{\{a,b\}}^\otimes(Q_{ab}^{\rm abs})
 =
 P_{ab}
 \sup_{\rho_A,\rho_B}
 \operatorname{Tr}[
 M_{ab}^{(1)}(\rho_A\otimes\rho_B)]
 \le
 P_{ab}\sqrt{d_C\Gamma_{ab}^{\otimes}}.}
 \label{eq:V72-pair-core-product}
\end{equation}
All within-row pair and core spectators are allowed.
\end{theorem}

\begin{proof}
Trace every row-bank factor except the retained \(a\) and \(b\)
factors.  Equations \eqref{eq:V72-three-bank-product} and
\eqref{eq:V72-pair-scalar-traces} give
\[
 P_{ab}\,
 I_{\mathrm{pair},a}\otimes I_{\mathrm{pair},b}
 \otimes M_{ab}^{(1)}.
\]
An arbitrary density matrix on row \(a\) may entangle its retained
pair-bank and core factors, but the expectation of the identity on
the pair-bank depends only on its reduced core density.  Every core
density is obtainable by tensoring with a fixed pair-bank pure state.
The same holds on row \(b\).  Therefore the product-state supremum
equals the first expression in
\eqref{eq:V72-pair-core-product}.  Apply
\eqref{eq:V72-product-square-bridge}.
\end{proof}

\begin{corollary}[Product-state sufficient collision gate]
\label{cor:V72-product-collision-gate}
Define \(\Gamma_{ac}^{\otimes}\) by the other fixed row orientation.
The direct matrix JREC conclusion of V67 follows on the
row-factorized coefficient domain if
\begin{equation}
 \boxed{
 \min\left\{
 P_{ab}\sqrt{d_C\Gamma_{ab}^{\otimes}},
 P_{ac}\sqrt{d_B\Gamma_{ac}^{\otimes}},
 \kappa_{\{o,a,b,c\}}^\otimes(Q_{ab}^{\rm abs})
 \right\}
 \le
 Cs_\alpha\sum_{r=b,c}
 \frac{h_{oa}\widehat H_{ar}h_{bc}}
      {\Xi_a\Xi_r}.}
 \label{eq:V72-product-collision-gate}
\end{equation}
This is a weaker sufficient hypothesis than
\eqref{eq:V68-square-function-gate}.
\end{corollary}

\begin{proof}
Apply \Cref{thm:V72-complete-product-tax} to the three displayed
retained-row choices, use
\Cref{thm:V72-pair-core-product} and its \(ac\) analogue, and take
the minimum.  The conclusion is the same JREC trace-norm bound because
\eqref{eq:V72-complete-product-tax} estimates the original
row-factorized coefficient tensor directly.  Finally,
\eqref{eq:V72-product-below-operator} and
\[
 \Gamma_{ab}^{\otimes}
 \le\|M_{ab}^{(2)}\|_{\rm op}
 =\max_S\|\mathscr S_S^\sim\|_{\rm op}
\]
show that every V68-admissible bound is also V72-admissible.
\end{proof}

\begin{assessment}[Where pair--core correlation does and does not
remain]
\label{ass:V72-correlation-location}
After \eqref{eq:V72-three-bank-product}, the pair output labels
\(\lambda,\mu\) have no probabilistic correlation with the core
labels \(S,V,T\); their retained partial traces are the scalars in
\eqref{eq:V72-pair-scalar-traces}.  Reordering those label sums
cannot create a new debit.

What remains correlated is the \emph{row-product constraint}.  The
joint \(AB\) operator may have a large eigenvalue on an entangled
Clebsch--Gordan vector while every admissible test is
\(\rho_A\otimes\rho_B\).  Formula \eqref{eq:V72-CG-target}, not a
pair-label rearrangement, is the genuine upstream object.
\end{assessment}

\section{Sharp examples for the product-state replacement}
\label{sec:V72-examples}

\begin{proposition}[Singlet projector gains one representation
dimension]
\label{prop:V72-singlet-product-norm}
Let \(H_B=H_A^*\), \(d_A=d\), and let
\(P_{\mathbf1}\) be the rank-one invariant projector in
\(H_A\otimes H_A^*\).  Then
\begin{equation}
 \boxed{
 \|P_{\mathbf1}\|_{\rm op}=1,
 \qquad
 \kappa_{\{A,B\}}^\otimes(P_{\mathbf1})=\frac1d.}
 \label{eq:V72-singlet-product-norm}
\end{equation}
\end{proposition}

\begin{proof}
In dual orthonormal bases the invariant unit vector is
\[
 \Omega=d^{-1/2}\sum_{j=1}^{d}e_j\otimes e_j^*.
\]
For unit vectors \(u,v\),
\[
 |\langle\Omega,u\otimes v\rangle|^2
 =
 \frac1d|\langle\overline v,u\rangle|^2
 \le\frac1d.
\]
Equality is attained when \(u=\overline v\).  Convexity shows that
mixed product states cannot give a larger value.  The operator norm
of a nonzero orthogonal projector is one.
\end{proof}

\begin{proposition}[No uniform gain on every irrep block]
\label{prop:V72-highest-product-norm}
Let \(P_{A+B}\) denote the Cartan highest-weight summand in
\(H_A\otimes H_B\).  Then
\begin{equation}
 \kappa_{\{A,B\}}^\otimes(P_{A+B})=1.
 \label{eq:V72-highest-product-norm}
\end{equation}
\end{proposition}

\begin{proof}
The tensor product of unit highest-weight vectors
\(v_A^{\rm hw}\otimes v_B^{\rm hw}\) belongs to the Cartan summand.
Its expectation of \(P_{A+B}\) is one.  The reverse inequality holds
for every projection.
\end{proof}

\begin{assessment}[Meaning of the two sharp tests]
\label{ass:V72-sharp-tests}
Product-state capacity can supply exactly the dimension factor missed
by an operator norm on an entangled singlet, but it supplies no
universal gain on a product-visible highest-weight block.  Therefore
the next theorem must locate the actual V71 thermal mesoscopic
square among these two extremes.  A generic appeal to
\emph{Clebsch--Gordan entanglement} is insufficient.
\end{assessment}

\section{V72 ledger and next acquisition}
\label{sec:V72-ledger}

\begin{theorem}[Integrated V72 statement]
\label{thm:V72-integrated}
Preserving V1--V71, V72 proves:
\begin{enumerate}[label=\textup{(V72.\arabic*)},leftmargin=5.7em]
\item the exact positive product-state capacity and tax
      \eqref{eq:V72-positive-product-tax};
\item the complete spectator trace-norm theorem
      \eqref{eq:V72-complete-product-tax};
\item the Kadison operator inequality
      \eqref{eq:V72-Kadison-operator} and product-state square bridge;
\item the exact Clebsch--Gordan target
      \eqref{eq:V72-CG-target};
\item the pair--core factorization
      \eqref{eq:V72-pair-core-product};
\item the weaker sufficient collision gate
      \eqref{eq:V72-product-collision-gate}; and
\item the sharp singlet-gain and highest-weight-no-gain examples.
\end{enumerate}
\end{theorem}

\begin{proof}
This is the collection of
\Cref{lem:V72-positive-product-tax,
thm:V72-complete-product-tax,thm:V72-product-square-bridge,
cor:V72-CG-target,thm:V72-pair-core-product,
cor:V72-product-collision-gate,prop:V72-singlet-product-norm,
prop:V72-highest-product-norm}.
\end{proof}

\begin{openproblem}[V73 mandate: product visibility of the thermal
mesoscopic square]
\label{op:V73-mandate}
Return to the actual residual collision monomial and identify the
irreducibles \(A,B\) which produce the retained dual-symmetric
intermediate \(S=R_n\) or \(Q_n\).  Then:
\begin{enumerate}
\item evaluate or sharply bound
      \(\sup_{u,v}\|P_S^{AB}(u\otimes v)\|^2\) for that exact
      family, including LR multiplicity;
\item keep the conditional matrices
      \(\widehat F_S(T_k)\) for \(k\asymp\sqrt n\) inside the
      Clebsch--Gordan quadratic form
      \eqref{eq:V72-CG-target};
\item prove the required inverse-\(n\) gain in
      \(P_{ab}\sqrt{d_C\Gamma_{ab}^{\otimes}}\), or exhibit a product
      vector which prevents it;
\item perform the same calculation in one coherent \(ac\)
      orientation before taking the minimum;
\item if a product vector saturates both orientations, return further
      upstream to the exact non-positive connected coefficient rather
      than strengthening the separable norm by assumption.
\end{enumerate}
The next compulsory companion audit remains V75.
\end{openproblem}

\section{Firewall after V72}
\label{sec:V72-firewall}

\begin{firewall}[Current claim boundary]
\label{fire:V72-final}
V72 removes an unnecessary cross-row entangled test from the
sufficient JREC criterion.  It does not estimate the resulting
product-state square on the actual thermal mesoscopic family.  Hence
the product-state collision gate, JREC, simultaneous sideways
closure, puncture aggregation, multi-collision summation, the
remaining partition sectors, the cofinal Yang--Mills construction,
and the continuum mass gap remain open.
\end{firewall}

\section{Append-only record for V72}
\label{sec:V72-record}

V72 was appended to the complete V71 manuscript.  The exact V71
parent had \(51842\) lines, \(1867690\) bytes, and SHA--256 digest
\[
 \texttt{7ccb5a4af32affb3967cacb73f78bef038ec125bf2a025b7010c4c34bbd37e50}.
\]
No V71 mathematical content was changed.  The sole final
\verb|\end{document}| is restored below.

% =============================================================================
% END APPENDED VERSION V72 MATERIAL
% =============================================================================
% =============================================================================
% BEGIN APPENDED VERSION V73 MATERIAL
% =============================================================================

\chapter{V73: Double-Schur product bounds and shellwise
Hilbert--Schmidt capacity}
\label{ch:V73}

\section{Purpose}
\label{sec:V73-purpose}

V72 replaced the joint \(AB\) operator norm by the product-state
capacity dictated by the original row-factorized coefficient domain.
V73 derives an exact representation-theoretic upper bound for that
capacity.  A positive operator commuting with the diagonal
\(\mathrm{SU}(3)\) action has scalar one-row partial traces, so its
product-state capacity is bounded by its total trace divided by the
larger retained representation dimension.

The estimate must be applied shell by shell.  Applying it once to the
complete positive square would replace a maximal ancestry
concentration by a full Hilbert--Schmidt mass and can be worse than
V68.  The shellwise minimum retains both alternatives and turns the
next obstruction into an explicit weighted ternary-multiplicity sum.

\begin{firewall}[Scope of V73]
\label{fire:V73-scope}
V73 proves a new sufficient gate, not the estimate required by that
gate.  No bound on ternary Littlewood--Richardson multiplicities
stronger than dimension counting is assumed.
\end{firewall}

\section{The double-Schur inequality}
\label{sec:V73-double-Schur}

For a positive operator \(M\) on \(H_A\otimes H_B\), define
\begin{equation}
 \|M\|_{A\otimes B,\mathrm{prod}}
 :=
 \sup_{\substack{\rho_A,\rho_B\ge0\\
 \operatorname{Tr}\rho_A=\operatorname{Tr}\rho_B=1}}
 \operatorname{Tr}[M(\rho_A\otimes\rho_B)].
 \label{eq:V73-product-norm}
\end{equation}

\begin{theorem}[Double-Schur product-state bound]
\label{thm:V73-double-Schur}
Suppose \(M\ge0\) commutes with
\(D^A(g)\otimes D^B(g)\) for every \(g\in\mathrm{SU}(3)\).  Then
\begin{equation}
 \boxed{
 \|M\|_{A\otimes B,\mathrm{prod}}
 \le
 \min\left\{
 \|M\|_{\rm op},
 \frac{\operatorname{Tr}M}{d_A},
 \frac{\operatorname{Tr}M}{d_B}
 \right\}
 =
 \min\left\{
 \|M\|_{\rm op},
 \frac{\operatorname{Tr}M}{\max\{d_A,d_B\}}
 \right\}.}
 \label{eq:V73-double-Schur}
\end{equation}
\end{theorem}

\begin{proof}
Invariance of \(M\) implies
\[
 D^B(g)(\operatorname{Tr}_A M)D^B(g)^*
 =\operatorname{Tr}_A M.
\]
Since \(B\) is irreducible, Schur's lemma and total trace give
\[
 \operatorname{Tr}_A M
 =\frac{\operatorname{Tr}M}{d_B}I_B.
\]
For density matrices \(\rho_A,\rho_B\), use
\(0\le\rho_A\le I_A\):
\[
 \operatorname{Tr}[M(\rho_A\otimes\rho_B)]
 \le
 \operatorname{Tr}[M(I_A\otimes\rho_B)]
 =
 \frac{\operatorname{Tr}M}{d_B}.
\]
Interchanging \(A,B\) gives the quotient by \(d_A\).  The operator
norm gives the first entry.  Take the minimum.  Since the smaller of
\(\operatorname{Tr}M/d_A\) and
\(\operatorname{Tr}M/d_B\) has the larger denominator, the two forms
are equal.
\end{proof}

\begin{remark}[Sharpness at both ends]
\label{rem:V73-double-Schur-sharp}
For the dual singlet projector,
\(\operatorname{Tr}M=1\) and \(d_A=d_B\), so
\eqref{eq:V73-double-Schur} is the exact value
\eqref{eq:V72-singlet-product-norm}.  For a Cartan product projector
with a product highest-weight vector, the operator-norm entry one is
sharp, as in \eqref{eq:V72-highest-product-norm}.  Neither entry may
be deleted.
\end{remark}

\begin{corollary}[Positive shell summation]
\label{cor:V73-shell-sum}
If \(M=\sum_{j\in J}M_j\), where every \(M_j\ge0\) is invariant, then
\begin{equation}
 \boxed{
 \|M\|_{A\otimes B,\mathrm{prod}}
 \le
 \sum_{j\in J}
 \min\left\{
 \|M_j\|_{\rm op},
 \frac{\operatorname{Tr}M_j}{\max\{d_A,d_B\}}
 \right\}.}
 \label{eq:V73-shell-sum}
\end{equation}
\end{corollary}

\begin{proof}
For every product density, linearity gives the sum of the
expectations of \(M_j\).  The supremum of a sum of nonnegative
functions is at most the sum of their suprema.  Apply
\eqref{eq:V73-double-Schur} to each summand.
\end{proof}

\section{Output-shell decomposition of the Yang--Mills square}
\label{sec:V73-output-shells}

Let \(\{\Omega_j\}_{j\in J}\) be a partition of the final output
irreducibles, and let \(P_{\Omega_j}\) be the corresponding central
projection in \(H_A\otimes H_B\otimes H_C\).  With
\(\mathbf K^\sim\) from \eqref{eq:V72-global-weighted-K}, put
\begin{equation}
 M_{ab,j}^{(2)}
 :=
 \operatorname{Tr}_C
 \left[
 P_{\Omega_j}(\mathbf K^\sim)^2
 \right].
 \label{eq:V73-shell-square}
\end{equation}
The central projection commutes with \(\mathbf K^\sim\), every
\(M_{ab,j}^{(2)}\) is positive and invariant, and
\[
 M_{ab}^{(2)}=\sum_jM_{ab,j}^{(2)}.
\]
Define
\begin{align}
 \Sigma_{ab,j}^{(2)}
 &:=
 \|M_{ab,j}^{(2)}\|_{\rm op},
 \label{eq:V73-shell-operator-size}\\
 \mathcal H_{ab,j}^{(2)}
 &:=
 \operatorname{Tr}M_{ab,j}^{(2)}.
 \label{eq:V73-shell-HS-size}
\end{align}

\begin{theorem}[Exact shell Hilbert--Schmidt identity]
\label{thm:V73-shell-HS}
For every output shell,
\begin{equation}
 \boxed{
 \mathcal H_{ab,j}^{(2)}
 =
 \sum_{T\in\Omega_j}
 d_T\operatorname{Tr}_{\mathcal M_T}
 (\widetilde K_T)^2.}
 \label{eq:V73-shell-HS}
\end{equation}
Consequently
\begin{equation}
 \boxed{
 \Gamma_{ab}^{\otimes}
 \le
 \sum_j
 \min\left\{
 \Sigma_{ab,j}^{(2)},
 \frac{\mathcal H_{ab,j}^{(2)}}
      {\max\{d_A,d_B\}}
 \right\}.}
 \label{eq:V73-Gamma-shell-gate}
\end{equation}
\end{theorem}

\begin{proof}
Partial trace preserves total trace.  On the final \(T\)-isotypic
space, \((\mathbf K^\sim)^2\) is
\(I_{H_T}\otimes(\widetilde K_T)^2\), whose trace is
\(d_T\operatorname{Tr}_{\mathcal M_T}(\widetilde K_T)^2\).
Sum \(T\in\Omega_j\) to prove
\eqref{eq:V73-shell-HS}.  Equation
\eqref{eq:V73-Gamma-shell-gate} is
\Cref{cor:V73-shell-sum} applied to
\eqref{eq:V73-shell-square}.
\end{proof}

\begin{corollary}[Shellwise product-state collision gate]
\label{cor:V73-shell-collision-gate}
The \(ab\)-oriented entry in the V72 collision gate is bounded by
\begin{equation}
 \boxed{
 P_{ab}\sqrt{
 d_C
 \sum_j
 \min\left\{
 \Sigma_{ab,j}^{(2)},
 \frac{\mathcal H_{ab,j}^{(2)}}
      {\max\{d_A,d_B\}}
 \right\}}.}
 \label{eq:V73-shell-collision-gate}
\end{equation}
The analogous formula holds in the complete \(ac\) orientation.
\end{corollary}

\begin{proof}
Insert \eqref{eq:V73-Gamma-shell-gate} into
\eqref{eq:V72-pair-core-product}.
\end{proof}

\section{A completely explicit multiplicity majorant}
\label{sec:V73-multiplicity}

Let
\begin{equation}
 m_{ABC}(T):=
 \dim\operatorname{Hom}_G
 (H_T,H_A\otimes H_B\otimes H_C)
 =
 \sum_SN_{AB}^{S}N_{SC}^{T}.
 \label{eq:V73-ternary-multiplicity}
\end{equation}
By the other bracketing it also equals
\(\sum_VN_{BC}^{V}N_{AV}^{T}\).

\begin{lemma}[Four dimension-counting bounds]
\label{lem:V73-four-dimension-bounds}
For every \(T\),
\begin{equation}
 \boxed{
 m_{ABC}(T)
 \le
 \min\left\{
 \frac{d_Ad_Bd_C}{d_T},
 \frac{d_Bd_Cd_T}{d_A},
 \frac{d_Ad_Cd_T}{d_B},
 \frac{d_Ad_Bd_T}{d_C}
 \right\}.}
 \label{eq:V73-four-dimension-bounds}
\end{equation}
\end{lemma}

\begin{proof}
The first inequality is ordinary dimension counting in
\(A\otimes B\otimes C\).  Frobenius reciprocity identifies the same
multiplicity with the multiplicity of \(A^*\) in
\(B\otimes C\otimes T^*\); dimension counting there gives the second.
The last two follow by cycling \(A,B,C\).
\end{proof}

\begin{theorem}[Weighted ternary-multiplicity trace bound]
\label{thm:V73-multiplicity-trace}
For every shell \(\Omega_j\),
\begin{equation}
 \boxed{
 \mathcal H_{ab,j}^{(2)}
 \le
 C\sum_{T\in\Omega_j}
 \frac{d_Tm_{ABC}(T)}
      {[1+s_\alpha\Xi(T)]^2}.}
 \label{eq:V73-multiplicity-trace}
\end{equation}
In particular, the right side together with
\eqref{eq:V73-four-dimension-bounds} is a fully explicit sufficient
input for \eqref{eq:V73-shell-collision-gate}.
\end{theorem}

\begin{proof}
The V63 weighted-cut decay
\eqref{eq:V63-cut-weight-decay} is precisely
\[
 \|\widetilde K_T\|_{\rm op}
 \le\frac C{1+s_\alpha\Xi(T)}.
\]
The multiplicity space has dimension \(m_{ABC}(T)\), so positivity
gives
\[
 \operatorname{Tr}_{\mathcal M_T}(\widetilde K_T)^2
 \le
 m_{ABC}(T)\|\widetilde K_T\|_{\rm op}^2
 \le
 \frac{C\,m_{ABC}(T)}
      {[1+s_\alpha\Xi(T)]^2}.
\]
Insert this in \eqref{eq:V73-shell-HS}.
\end{proof}

\begin{assessment}[What the multiplicity bound does not capture]
\label{ass:V73-multiplicity-limit}
Equation \eqref{eq:V73-multiplicity-trace} uses the norm of
\(\widetilde K_T\) on its entire multiplicity space.  It ignores the
three covariance squares of V69 and the conditional \(V\)-law of
V71.  It is therefore a safe shell input, but not a reason to expect
the thermal mesoscopic shell to close.  The shellwise minimum in
\eqref{eq:V73-Gamma-shell-gate} must be retained; replacing every
shell by the trace alternative can turn a low-rank gain into a large
global multiplicity count.
\end{assessment}

\section{Finite-window and singlet consequences}
\label{sec:V73-window-consequences}

\begin{corollary}[Product-state fixed-output certificate]
\label{cor:V73-fixed-output-certificate}
For any finite output set \(\Omega\), its contribution
\(\Gamma_{ab,\Omega}^{\otimes}\) to the product-state square satisfies
\begin{equation}
 \boxed{
 \Gamma_{ab,\Omega}^{\otimes}
 \le
 \min\left\{
 \Sigma_{ab,\Omega}^{(2)},
 \frac C{\max\{d_A,d_B\}}
 \sum_{T\in\Omega}d_Tm_{ABC}(T)
 \right\}.}
 \label{eq:V73-fixed-output-certificate}
\end{equation}
For \(T=\mathbf1\), the trace alternative is
\begin{equation}
 \Gamma_{ab,\mathbf1}^{\otimes}
 \le
 \frac{
 \operatorname{Tr}_{\mathcal M_{\mathbf1}}
 (\widetilde K_{\mathbf1})^2}
 {\max\{d_A,d_B\}}.
 \label{eq:V73-singlet-trace-certificate}
\end{equation}
\end{corollary}

\begin{proof}
Use one shell \(\Omega\) in
\eqref{eq:V73-Gamma-shell-gate} and drop the denominator
\([1+s_\alpha\Xi(T)]^2\ge1\).  For the singlet,
\(d_{\mathbf1}=1\) and
\(\Xi(\mathbf1)=0\), while
\eqref{eq:V73-shell-HS} is exact.
\end{proof}

\begin{remark}[Two independent singlet debits]
\label{rem:V73-two-singlet-debits}
V71 obtained an inverse-\(d_C\) debit by tracing the final singlet
through \(C\).  Equation \eqref{eq:V73-singlet-trace-certificate}
can additionally divide its residual Hilbert--Schmidt mass by
\(\max\{d_A,d_B\}\) when testing product rows.  These are different
mechanisms and should not be multiplied unless both appear in the
same exact positive operator.
\end{remark}

\section{The thermal-shell trace gate}
\label{sec:V73-thermal-gate}

Let
\[
 \Omega_{\rm th}(c_-,c_+)
 :=
 \{T:c_-\le s_\alpha\Xi(T)\le c_+\}
\]
for fixed \(0<c_-<c_+<\infty\).  The multiplier denominator in
\eqref{eq:V73-multiplicity-trace} is then comparable to one.

\begin{proposition}[Exact remaining Hilbert--Schmidt alternative]
\label{prop:V73-thermal-alternative}
The thermal-shell contribution satisfies
\begin{equation}
 \boxed{
 \Gamma_{ab,\rm th}^{\otimes}
 \le
 \min\left\{
 \Sigma_{ab,\rm th}^{(2)},
 \frac1{\max\{d_A,d_B\}}
 \sum_{T\in\Omega_{\rm th}}
 d_T\operatorname{Tr}_{\mathcal M_T}
 (\widetilde K_T)^2
 \right\}.}
 \label{eq:V73-thermal-alternative}
\end{equation}
Consequently the inverse-\(d_C^{1/2}\) companion needed by the V72
square bridge would follow on this shell from the sufficient
Hilbert--Schmidt condition
\begin{equation}
 \sum_{T\in\Omega_{\rm th}}
 d_T\operatorname{Tr}_{\mathcal M_T}
 (\widetilde K_T)^2
 \le
 \frac{\max\{d_A,d_B\}}{d_C}.
 \label{eq:V73-thermal-HS-target}
\end{equation}
\end{proposition}

\begin{proof}
The first formula is the one-shell case of
\eqref{eq:V73-Gamma-shell-gate}, with the exact trace identity
\eqref{eq:V73-shell-HS}.  Condition
\eqref{eq:V73-thermal-HS-target} makes the trace entry at most
\(d_C^{-1}\).  The square bridge then gives
\(\sqrt{d_C\Gamma_{ab,\rm th}^{\otimes}}\le1\).
\end{proof}

\begin{assessment}[The next bottleneck is now quantitative]
\label{ass:V73-thermal-bottleneck}
The sufficient target
\eqref{eq:V73-thermal-HS-target} is not proved by the norm majorant
\eqref{eq:V73-multiplicity-trace}; on the thermal shell that majorant
retains the full ternary multiplicity.  The next attack must compute
the Hilbert--Schmidt mass of the \emph{actual covariance symbol},
using the V69 three-cut decomposition before taking a multiplicity
norm.  If that mass is too large, one must keep the product-vector
quadratic form and cannot close by trace alone.
\end{assessment}

\section{V73 ledger and next acquisition}
\label{sec:V73-ledger}

\begin{theorem}[Integrated V73 statement]
\label{thm:V73-integrated}
Preserving V1--V72, V73 proves:
\begin{enumerate}[label=\textup{(V73.\arabic*)},leftmargin=5.7em]
\item the double-Schur product-state inequality
      \eqref{eq:V73-double-Schur};
\item its positive shell-summation form
      \eqref{eq:V73-shell-sum};
\item the exact output-shell Hilbert--Schmidt identity
      \eqref{eq:V73-shell-HS};
\item the shellwise product collision gate
      \eqref{eq:V73-shell-collision-gate};
\item the four dimension-counting bounds for the exact ternary
      multiplicity;
\item the weighted multiplicity trace majorant
      \eqref{eq:V73-multiplicity-trace}; and
\item the explicit thermal Hilbert--Schmidt target
      \eqref{eq:V73-thermal-HS-target}.
\end{enumerate}
\end{theorem}

\begin{proof}
These are
\Cref{thm:V73-double-Schur,cor:V73-shell-sum,
thm:V73-shell-HS,cor:V73-shell-collision-gate,
lem:V73-four-dimension-bounds,thm:V73-multiplicity-trace,
prop:V73-thermal-alternative}.
\end{proof}

\begin{openproblem}[V74 mandate: covariance Hilbert--Schmidt mass]
\label{op:V74-mandate}
On one fixed orientation and one thermal output shell:
\begin{enumerate}
\item insert the exact V69 identity
      \(K_T=D_{A\mid BC}^T-q_{\alpha,B}D_{AC}^T
      -q_{\alpha,C}D_{AB}^T\) before taking multiplicity trace;
\item compute the three positive Hilbert--Schmidt square sums using
      complete dimension laws, retaining the negative covariance
      cross terms whenever their trace is explicit;
\item compare the result with
      \eqref{eq:V73-thermal-HS-target}, not merely with the uniform
      multiplicity majorant;
\item if the trace target fails, construct a product vector test in
      \eqref{eq:V72-CG-target} to decide whether the failure is real
      or only a Hilbert--Schmidt overcount;
\item keep the \(ac\) orientation separate and postpone the minimum
      until both complete bounds are formed.
\end{enumerate}
The next mandatory SM/NS audit is V75.
\end{openproblem}

\section{Firewall after V73}
\label{sec:V73-firewall}

\begin{firewall}[Current claim boundary]
\label{fire:V73-final}
V73 converts the product-state collision problem into a shellwise
minimum of an operator concentration and an exact
Hilbert--Schmidt trace.  It does not prove the thermal trace target
or the alternative product-vector bound.  Therefore JREC,
simultaneous sideways closure, puncture aggregation, multi-collision
summation, the remaining partition sectors, the cofinal Yang--Mills
construction, and the positive continuum spectral gap remain open.
\end{firewall}

\section{Append-only record for V73}
\label{sec:V73-record}

V73 was appended to the complete V72 manuscript.  The exact V72
parent had \(52418\) lines, \(1885694\) bytes, and SHA--256 digest
\[
 \texttt{d0e9fbf291e3a79d718bc49b6748edc9b78423b39513dbbc10ae0f66b7a7c1f0}.
\]
No V72 mathematical content was changed.  The sole final
\verb|\end{document}| is restored below.

% =============================================================================
% END APPENDED VERSION V73 MATERIAL
% =============================================================================
% =============================================================================
% BEGIN APPENDED VERSION V74 MATERIAL
% =============================================================================

\chapter{V74: Sequential dimension laws for the covariance
Hilbert--Schmidt energy}
\label{ch:V74}

\section{Purpose}
\label{sec:V74-purpose}

V73 reduced the product-state route to an exact shell
Hilbert--Schmidt target.  The multiplicity-norm majorant in that
version ignored the covariance structure of the ternary symbol.
V74 now computes the multiplicity trace of each of the three V69
binary-cut squares.  Because trace is invariant under Racah
conjugation, no entrywise recoupling coefficient appears: each answer
is a scalar expectation under an exact sequential tensor-dimension
law.

\begin{firewall}[Scope of V74]
\label{fire:V74-scope}
The calculation yields an exact positive upper bound for the
Hilbert--Schmidt energy and a necessary quantitative scale for the
trace route.  It does not prove that this scale is attained on the
thermal shell.  In particular, replacing the exact expectations by
their uniform caps does not close the V73 target.
\end{firewall}

\section{Three sequential tensor-dimension laws}
\label{sec:V74-dimension-laws}

Every path label below includes its full LR multiplicity index.
For the \(BC\)-first bracketing, define
\begin{equation}
 \mathbb P_{BC,A}(V,\eta,T,\gamma)
 :=
 \frac{d_T}{d_Ad_Bd_C},
 \quad
 1\le\eta\le N_{BC}^{V},\quad
 1\le\gamma\le N_{AV}^{T}.
 \label{eq:V74-BC-A-law}
\end{equation}
For the \(AC\)-first and \(AB\)-first bracketings define
\begin{align}
 \mathbb P_{AC,B}(U,\mu,T,\nu)
 &:=
 \frac{d_T}{d_Ad_Bd_C},
 \label{eq:V74-AC-B-law}\\
 \mathbb P_{AB,C}(S,\alpha,T,\beta)
 &:=
 \frac{d_T}{d_Ad_Bd_C}.
 \label{eq:V74-AB-C-law}
\end{align}
The indicated labels range over
\[
 U\subset A\otimes C,\quad T\subset U\otimes B,
 \qquad
 S\subset A\otimes B,\quad T\subset S\otimes C.
\]

\begin{lemma}[Normalization of the sequential laws]
\label{lem:V74-law-normalization}
Each of
\eqref{eq:V74-BC-A-law}--\eqref{eq:V74-AB-C-law} has total mass one.
\end{lemma}

\begin{proof}
For the first law,
\[
 \sum_{V,\eta,T,\gamma}d_T
 =
 \sum_VN_{BC}^{V}
 \sum_TN_{AV}^{T}d_T
 =
 \sum_VN_{BC}^{V}d_Ad_V
 =
 d_Ad_Bd_C.
\]
Divide by the denominator in \eqref{eq:V74-BC-A-law}.  The other two
identities follow by permuting the inputs.
\end{proof}

Put
\[
 w(T):=1+s_\alpha\Xi(T),
 \qquad x=q_{\alpha,A},\quad
 y=q_{\alpha,B},\quad z=q_{\alpha,C}.
\]
With the binary defects of V69--V70, define the nonnegative energies
\begin{align}
 \mathcal E_{A\mid BC}
 &:=
 \mathbb E_{BC,A}
 \left[
 w(T)^2
 (q_{\alpha,T}-xq_{\alpha,V})^2
 \right],
 \label{eq:V74-E-composite}\\
 \mathcal E_{AC}
 &:=
 \mathbb E_{AC,B}
 \left[
 w(T)^2
 (q_{\alpha,U}-xz)^2
 \right],
 \label{eq:V74-E-AC}\\
 \mathcal E_{AB}
 &:=
 \mathbb E_{AB,C}
 \left[
 w(T)^2
 (q_{\alpha,S}-xy)^2
 \right].
 \label{eq:V74-E-AB}
\end{align}
For an output set \(\Omega\), add a subscript \(\Omega\) by inserting
\(\mathbf1_{\{T\in\Omega\}}\) in each expectation.

\section{Exact traces of the three positive squares}
\label{sec:V74-three-traces}

\begin{theorem}[Sequential trace identities]
\label{thm:V74-sequential-traces}
On \(H_A\otimes H_B\otimes H_C\),
\begin{align}
 \operatorname{Tr}\left[
 P_\Omega w(\mathbf T)^2
 (D_{A\mid BC})^2
 \right]
 &=
 d_Ad_Bd_C\,\mathcal E_{A\mid BC,\Omega},
 \label{eq:V74-trace-composite}\\
 \operatorname{Tr}\left[
 P_\Omega w(\mathbf T)^2
 (D_{AC})^2
 \right]
 &=
 d_Ad_Bd_C\,\mathcal E_{AC,\Omega},
 \label{eq:V74-trace-AC}\\
 \operatorname{Tr}\left[
 P_\Omega w(\mathbf T)^2
 (D_{AB})^2
 \right]
 &=
 d_Ad_Bd_C\,\mathcal E_{AB,\Omega}.
 \label{eq:V74-trace-AB}
\end{align}
Here \(w(\mathbf T)\) is the scalar \(w(T)\) on the final
\(T\)-isotypic block.
\end{theorem}

\begin{proof}
In the \(BC\)-first bracketing,
\(D_{A\mid BC}\) is diagonal on the path
\((V,\eta,T,\gamma)\) with eigenvalue
\(q_{\alpha,T}-xq_{\alpha,V}\).  The representation factor of that
path has dimension \(d_T\).  Therefore the first trace is
\[
 \sum_{\substack{V,\eta,T,\gamma\\T\in\Omega}}
 d_Tw(T)^2(q_{\alpha,T}-xq_{\alpha,V})^2,
\]
which is \eqref{eq:V74-trace-composite} by
\eqref{eq:V74-BC-A-law}.

In the \(AC\)-first bracketing, \(D_{AC}\) has eigenvalue
\(q_{\alpha,U}-xz\), independently of the subsequent \(T,\nu\).
The same path trace proves \eqref{eq:V74-trace-AC}.  The
\(AB\)-first calculation is identical and proves
\eqref{eq:V74-trace-AB}.  No recoupling coefficient enters because
the trace may be evaluated in the basis which diagonalizes the
chosen square.
\end{proof}

\begin{corollary}[Aligned energy in one-stage form]
\label{cor:V74-aligned-energy}
The \(AB\) energy also has the exact form
\begin{equation}
 \boxed{
 d_Ad_Bd_C\,\mathcal E_{AB}
 =
 \sum_S
 N_{AB}^{S}d_SM_2(S,C)
 (q_{\alpha,S}-xy)^2,}
 \label{eq:V74-aligned-energy}
\end{equation}
where \(M_2(S,C)\) is \eqref{eq:V69-one-sided-output-second}.
\end{corollary}

\begin{proof}
In \eqref{eq:V74-trace-AB}, sum first over
\(T,\beta\) at fixed \(S,\alpha\).  The inner sum is
\[
 \sum_TN_{SC}^{T}d_Tw(T)^2=d_SM_2(S,C).
\]
There are \(N_{AB}^{S}\) values of \(\alpha\).
\end{proof}

\section{Covariance energy of the complete ternary square}
\label{sec:V74-complete-energy}

For an output set \(\Omega\), put
\begin{equation}
 \mathcal H_K(\Omega):=
 \sum_{T\in\Omega}
 d_T\operatorname{Tr}_{\mathcal M_T}
 (\widetilde K_T)^2.
 \label{eq:V74-HK}
\end{equation}
This is exactly the V73 shell Hilbert--Schmidt mass.

\begin{theorem}[Three-law covariance energy bound]
\label{thm:V74-covariance-energy}
For every output set \(\Omega\),
\begin{equation}
 \boxed{
 \mathcal H_K(\Omega)
 \le
 3d_Ad_Bd_C
 \left[
 \mathcal E_{A\mid BC,\Omega}
 +y^2\mathcal E_{AC,\Omega}
 +z^2\mathcal E_{AB,\Omega}
 \right].}
 \label{eq:V74-covariance-energy}
\end{equation}
\end{theorem}

\begin{proof}
The V69 operator inequality
\eqref{eq:V69-three-square-bound}, multiplied by the commuting
positive scalar \(w(T)^2\), gives
\[
 (\widetilde K_T)^2
 \le
 3w(T)^2
 \left[
 (D_{A\mid BC}^{T})^2
 +y^2(D_{AC}^{T})^2
 +z^2(D_{AB}^{T})^2
 \right]
\]
on every final block.  Multiply by \(P_\Omega\), take the full trace,
and apply
\eqref{eq:V74-trace-composite}--\eqref{eq:V74-trace-AB}.
The left side is \eqref{eq:V74-HK}.
\end{proof}

\begin{remark}[What cancellation has already been retained]
\label{rem:V74-cancellation-scope}
The five original cumulant terms were combined exactly into the three
binary covariances before the factor three was used.  Thus
\eqref{eq:V74-covariance-energy} is not the forbidden \(25\)-term
absolute expansion.  It does, however, discard the three cross
traces between the covariance cuts.  Recovering those cross traces is
permitted if the positive bound above is still too large, but their
sign has not been assumed.
\end{remark}

\section{Rigorous scalar bounds for the three energies}
\label{sec:V74-energy-bounds}

\begin{proposition}[Dimension-law energy bounds]
\label{prop:V74-energy-bounds}
With \(g_s\) from \eqref{eq:V70-BC-envelope},
\begin{align}
 \mathcal E_{A\mid BC}
 &\le
 C\mathbb E_{\nu_{BC}}[g_s(A,V)]
 \le C,
 \label{eq:V74-E-composite-bound}\\
 \mathcal E_{AC}
 &\le
 C\mathbb E_{\nu_{AC}}\left[
 [1+s_\alpha(A_B+\xi(U))]^2
 \min\{s_\alpha A_AA_C,1\}^2
 \right],
 \label{eq:V74-E-AC-bound}\\
 \mathcal E_{AB}
 &\le
 C\mathbb E_{\nu_{AB}}\left[
 [1+s_\alpha(A_C+\xi(S))]^2
 \min\{s_\alpha A_AA_B,1\}^2
 \right].
 \label{eq:V74-E-AB-bound}
\end{align}
The laws \(\nu_{BC},\nu_{AC},\nu_{AB}\) are the exact
tensor-dimension laws of the indicated first fusion.
\end{proposition}

\begin{proof}
For the first line, condition
\eqref{eq:V74-BC-A-law} on \(V,\eta\).  The conditional \(T,\gamma\)
weights have total mass one, and
\(f_{A,V}(T)\le Cg_s(A,V)\) by V70.  Average over
\(\nu_{BC}\).

For the second line, fusion in \(U\otimes B\) gives
\[
 w(T)\le C[1+s_\alpha(A_B+\xi(U))].
\]
The binary defect cap gives
\[
 |q_{\alpha,U}-xz|
 \le C\min\{s_\alpha A_AA_C,1\}.
\]
Condition on \(U,\mu\), sum the complete \(T,\nu\) law, and then
average over \(\nu_{AC}\).  The third line is its \(B,C\)
permutation.
\end{proof}

\begin{assessment}[Uniform energy bounds are intentionally
insufficient]
\label{ass:V74-uniform-insufficient}
The first bound in \eqref{eq:V74-E-composite-bound} is at most a
constant, and the other two may also be of constant size on a
thermal input configuration.  Inserting only those caps in
\eqref{eq:V74-covariance-energy} yields the full ambient dimension
\(d_Ad_Bd_C\).  This is a correct bound but is worse than the
conditional operator estimates of V70--V71.  The trace route can
advance only through small \emph{averaged covariance energy}, not
through normalization alone.
\end{assessment}

\section{Exact scale demanded by the double-Schur route}
\label{sec:V74-trace-target}

For one shell \(\Omega\), define
\begin{equation}
 \mathcal E_{\rm cov}(\Omega)
 :=
 \mathcal E_{A\mid BC,\Omega}
 +y^2\mathcal E_{AC,\Omega}
 +z^2\mathcal E_{AB,\Omega}.
 \label{eq:V74-Ecov}
\end{equation}

\begin{corollary}[Dimension-law sufficient condition for the trace
gate]
\label{cor:V74-energy-trace-gate}
The V73 Hilbert--Schmidt target
\eqref{eq:V73-thermal-HS-target} follows from
\begin{equation}
 \boxed{
 \mathcal E_{\rm cov}(\Omega_{\rm th})
 \le
 \frac{c}{
 \min\{d_A,d_B\}\,d_C^2}.}
 \label{eq:V74-energy-trace-target}
\end{equation}
More generally, a shell \(\Omega_j\) may use the trace entry of
\eqref{eq:V73-Gamma-shell-gate} whenever the right side of
\eqref{eq:V74-covariance-energy}, divided by
\(\max\{d_A,d_B\}\), is smaller than its operator entry.
\end{corollary}

\begin{proof}
Since
\[
 \frac{\max\{d_A,d_B\}}{d_Ad_B}
 =\frac1{\min\{d_A,d_B\}},
\]
\eqref{eq:V74-energy-trace-target} and
\eqref{eq:V74-covariance-energy} give
\[
 \mathcal H_K(\Omega_{\rm th})
 \le
 3d_Ad_Bd_C
 \frac{c}{\min\{d_A,d_B\}d_C^2}
 \le
 \frac{\max\{d_A,d_B\}}{d_C}
\]
after choosing the numerical \(c\) to absorb the factor three.  This
is \eqref{eq:V73-thermal-HS-target}.
\end{proof}

\begin{assessment}[Quantitative verdict on the trace alternative]
\label{ass:V74-trace-verdict}
The required averaged energy in
\eqref{eq:V74-energy-trace-target} contains two representation
dimensions beyond a constant probability bound.  Neither V44
Casimir moments nor V46 pointwise defect caps provide those
dimensions.  Therefore the Hilbert--Schmidt alternative has not
closed the thermal shell.  It remains useful on genuinely low-rank
shells, but at the mesoscopic boundary the next attack should return
to the product-vector quadratic form
\eqref{eq:V72-CG-target} unless an exact negative covariance cross
trace can be proved.
\end{assessment}

\section{An exact cross-trace decision card}
\label{sec:V74-cross-card}

Write
\[
 L_1=D_{A\mid BC},\qquad
 L_2=-yD_{AC},\qquad
 L_3=-zD_{AB}.
\]
On an output shell \(\Omega\), define
\begin{equation}
 \mathcal X_{ij}(\Omega)
 :=
 \frac1{d_Ad_Bd_C}
 \operatorname{ReTr}\left[
 P_\Omega w(\mathbf T)^2L_iL_j
 \right].
 \label{eq:V74-cross-energy}
\end{equation}

\begin{proposition}[Exact trace expansion after covariance
compression]
\label{prop:V74-cross-expansion}
One has
\begin{equation}
 \boxed{
 \frac{\mathcal H_K(\Omega)}{d_Ad_Bd_C}
 =
 \mathcal E_{\rm cov}(\Omega)
 +2\sum_{1\le i<j\le3}\mathcal X_{ij}(\Omega).}
 \label{eq:V74-cross-expansion}
\end{equation}
Moreover
\begin{equation}
 |\mathcal X_{ij}(\Omega)|
 \le
 \mathcal E_i(\Omega)^{1/2}
 \mathcal E_j(\Omega)^{1/2},
 \label{eq:V74-cross-CS}
\end{equation}
where the scalar coefficients \(y,z\) are included in
\(\mathcal E_2,\mathcal E_3\).
\end{proposition}

\begin{proof}
Use the exact V69 identity
\(\widetilde K=w(\mathbf T)(L_1+L_2+L_3)\), expand its squared
Hilbert--Schmidt norm, and divide by the total tensor dimension.
This gives \eqref{eq:V74-cross-expansion}.  Hilbert--Schmidt
Cauchy--Schwarz on
\(P_\Omega^{1/2}wL_i\) and
\(P_\Omega^{1/2}wL_j\), noting that the final central projection and
weight commute with every invariant \(L_i\), proves
\eqref{eq:V74-cross-CS}.
\end{proof}

\begin{firewall}[Cross terms are a target, not an assumed
cancellation]
\label{fire:V74-cross-scope}
Equation \eqref{eq:V74-cross-expansion} records the only place where
the factor-three positive bound can improve at trace level.  V74 has
not determined the signs of \(\mathcal X_{ij}\).  Quoting
\eqref{eq:V74-cross-CS} with a negative sign would be invalid.
\end{firewall}

\section{V74 ledger and next acquisition}
\label{sec:V74-ledger}

\begin{theorem}[Integrated V74 statement]
\label{thm:V74-integrated}
Preserving V1--V73, V74 proves:
\begin{enumerate}[label=\textup{(V74.\arabic*)},leftmargin=5.7em]
\item three normalized sequential tensor-dimension laws;
\item the exact traces
      \eqref{eq:V74-trace-composite}--\eqref{eq:V74-trace-AB};
\item the aligned one-stage formula
      \eqref{eq:V74-aligned-energy};
\item the three-law covariance energy bound
      \eqref{eq:V74-covariance-energy};
\item rigorous scalar bounds for each energy;
\item the exact averaged-energy scale
      \eqref{eq:V74-energy-trace-target} required by the V73 trace
      route; and
\item the exact three-cross-term decision card
      \eqref{eq:V74-cross-expansion}.
\end{enumerate}
\end{theorem}

\begin{proof}
These are
\Cref{lem:V74-law-normalization,thm:V74-sequential-traces,
cor:V74-aligned-energy,thm:V74-covariance-energy,
prop:V74-energy-bounds,cor:V74-energy-trace-gate,
prop:V74-cross-expansion}.
\end{proof}

\begin{openproblem}[V75 mandate after the compulsory companion audit]
\label{op:V75-mandate}
First inspect stable current snapshots of the SM and NS notes.  Then,
guided only by transferable proof order:
\begin{enumerate}
\item compute at least one cross trace
      \(\mathcal X_{ij}(\Omega_{\rm th})\) exactly, using a common
      group-integral or Racah trace representation;
\item if no sign-definite cancellation results, abandon the total
      trace route on the thermal shell and return to
      \eqref{eq:V72-CG-target};
\item identify the exact product-visible retained intermediate on the
      simultaneous-sideways monomial, rather than an arbitrary dual
      proxy;
\item preserve one complete \(ab\) or \(ac\) orientation throughout;
\item keep all later collision and partition sectors suspended until
      this single-collision gate is honestly resolved.
\end{enumerate}
\end{openproblem}

\section{Firewall after V74}
\label{sec:V74-firewall}

\begin{firewall}[Current claim boundary]
\label{fire:V74-final}
V74 computes the exact scalar dimension laws governing the three
covariance-square traces and quantifies the missing thermal energy.
It does not prove the required cross cancellation, product-vector
bound, JREC, simultaneous sideways closure, puncture aggregation,
multi-collision summation, the remaining partition sectors, the
cofinal Yang--Mills construction, or a positive continuum spectral
gap.
\end{firewall}

\section{Append-only record for V74}
\label{sec:V74-record}

V74 was appended to the complete V73 manuscript.  The exact V73
parent had \(52923\) lines, \(1901194\) bytes, and SHA--256 digest
\[
 \texttt{6520cfec457aa1a9367e0a19248d7487ed2ac23ba33a4e4385519ac2e1a74e5b}.
\]
No V73 mathematical content was changed.  The sole final
\verb|\end{document}| is restored below.

% =============================================================================
% END APPENDED VERSION V74 MATERIAL
% =============================================================================
% =============================================================================
% BEGIN APPENDED VERSION V75 MATERIAL
% =============================================================================

\chapter{V75: Companion audit and exact central character formulas
for the covariance cross traces}
\label{ch:V75}

\section{Purpose}
\label{sec:V75-purpose}

V74 isolated three undetermined cross traces in the
Hilbert--Schmidt energy.  V75 performs the required five-version
companion audit and then places all three traces in one common
central group-integral calculus.  The resulting formulas are exact
for every finite output shell and include all multiplicities
automatically.

\begin{firewall}[Scope of V75]
\label{fire:V75-scope}
The formulas below compute the cross traces, but do not assign them a
favorable sign.  A common integral representation is not itself a
cancellation theorem.
\end{firewall}

\section{Mandatory V75 companion audit}
\label{sec:V75-audit}

Two consecutive byte reads gave stable identical data:
\begin{center}
\begin{tabular}{|p{0.43\textwidth}|r|r|p{0.27\textwidth}|}
\hline
Companion file & Lines & Bytes & SHA--256\\
\hline
\texttt{Renewal\_SM\_Einstein\_Continuum\_Research\_Notes\_V80.tex}
& \(64341\) & \(2329670\) &
\texttt{16035310193d0c9068f44145ffc82fd6a503425afd89059d413fd4f9b1e7e6c0}\\
\hline
\texttt{Renewal\_NS\_Stability\_Research\_Notes\_V31.tex}
& \(23052\) & \(887537\) &
\texttt{f1121b62238ad5491bb7cf03635ea9934942edf573ed8faaa9ee79e1d38a6696}\\
\hline
\end{tabular}
\end{center}

The SM companion has advanced from V78 to V80.  V79 derives an exact
cap transgression and proves cancellation only after the complete
oriented side--cap--corner form is assembled.  V80 then splits the
exact joint zero mode from the mean-zero sector before using any
\(Q^{-2}\) formula; the zero fibre is not obtained by taking a
singular inverse-frequency limit.  The NS companion is again
byte-identical to its V65 and V70 snapshots.  Its V31 warning remains
that separate positive annular bounds lose the correlation carried
by the signed formation object.

\begin{assessment}[Transferable proof order at V75]
\label{ass:V75-audit-transfer}
For Yang--Mills, the exact \(T=\mathbf1\) fibre and fixed output
windows must remain separate from the thermal shell, as in V71.
Within the thermal shell all covariance cross terms must be placed in
one common representation before their signs are assessed.  This is
only a proof-order transfer; neither companion supplies a compact
group character estimate.
\end{assessment}

\begin{firewall}[Companion separation]
\label{fire:V75-companion-separation}
No Einstein Green-form identity and no Navier--Stokes stopping-time
estimate is used below.  The companion audit changes only the order
of the Yang--Mills calculation.
\end{firewall}

\section{Central realization of all moment operators}
\label{sec:V75-central-realization}

Work first at one compact-group coordinate.  Let
\(\mu_\alpha\) be the central inversion-symmetric Wilson probability
law whose normalized Fourier multipliers satisfy
\begin{equation}
 \int_G D^R(g)\,\mu_\alpha(\dd g)
 =q_{\alpha,R}I_{H_R}.
 \label{eq:V75-Fourier-multiplier}
\end{equation}
Write
\[
 \varphi_R(g):=\frac{\chi_R(g)}{d_R}.
\]
For \(J\subseteq\{A,B,C\}\), let
\begin{equation}
 \rho_J(g):=
 \bigotimes_{R\in J}D^R(g)
 \otimes
 \bigotimes_{R\notin J}I_{H_R},
 \qquad
 Q_J:=\int_G\rho_J(g)\,\mu_\alpha(\dd g),
 \label{eq:V75-QJ}
\end{equation}
with \(Q_\varnothing=I\).  Then
\[
 Q_{\{A,B,C\}}=Q_{ABC},\qquad
 Q_{\{B,C\}}=Q_{BC},\qquad
 Q_{\{A,C\}}=Q_{AC},\qquad
 Q_{\{A,B\}}=Q_{AB}.
\]
These are exactly the total and pair moment operators used in V69.

\begin{lemma}[Moment-operator identification]
\label{lem:V75-moment-identification}
On every irreducible output block of the indicated tensor product,
\(Q_J\) acts by the Wilson multiplier of that output.  In particular,
\begin{equation}
 K=
 Q_{ABC}-xQ_{BC}-yQ_{AC}-zQ_{AB}+2xyzI.
 \label{eq:V75-K-central}
\end{equation}
\end{lemma}

\begin{proof}
Equation \eqref{eq:V75-Fourier-multiplier} and Schur's lemma show that
the central average of a representation acts by its normalized
Fourier multiplier on each irreducible summand.  Tensoring with the
identity on rows outside \(J\) proves the first statement.  The
second is the common-basis cumulant formula
\eqref{eq:V47-common-cumulant}.
\end{proof}

\section{A finite central kernel for an output shell}
\label{sec:V75-shell-kernel}

Let \(\mathcal T(A,B,C)\) be the finite set of irreducibles occurring
in \(A\otimes B\otimes C\).  For an output set \(\Omega\), put
\[
 r_\Omega(T):=
 \mathbf1_{\{T\in\Omega\}}
 [1+s_\alpha\Xi(T)]^2
\]
and define the finite central character polynomial
\begin{equation}
 \kappa_\Omega(u):=
 \sum_{T\in\mathcal T(A,B,C)}
 d_Tr_\Omega(T)\overline{\chi_T(u)}.
 \label{eq:V75-shell-kernel}
\end{equation}

\begin{lemma}[Exact shell functional calculus]
\label{lem:V75-shell-functional-calculus}
The positive final-shell multiplier
\[
 \mathcal R_\Omega:=
 P_\Omega[1+s_\alpha\Xi(\mathbf T)]^2
\]
has the exact representation
\begin{equation}
 \boxed{
 \mathcal R_\Omega
 =
 \int_G
 \kappa_\Omega(u)
 \rho_{\{A,B,C\}}(u)\,\dd u.}
 \label{eq:V75-shell-functional-calculus}
\end{equation}
The kernel \(\kappa_\Omega\) need not be pointwise positive.
\end{lemma}

\begin{proof}
Schur orthogonality gives, for every irreducible \(R\),
\[
 \int_G
 \overline{\chi_T(u)}D^R(u)\,\dd u
 =
 \frac{\delta_{TR}}{d_R}I_{H_R}.
\]
Multiplication by \(d_Tr_\Omega(T)\) and summation in \(T\) show that
the right side of
\eqref{eq:V75-shell-functional-calculus} acts by
\(r_\Omega(R)\) on every \(R\)-isotypic block.  These are exactly the
eigenvalues of \(\mathcal R_\Omega\).
\end{proof}

\section{The common three-variable character integral}
\label{sec:V75-common-integral}

Let
\[
 \tau(X):=\frac{\operatorname{Tr}X}{d_Ad_Bd_C}
\]
be normalized trace.  For \(J,L\subseteq\{A,B,C\}\), define
\begin{align}
 \Psi_\Omega(J,L)
 &:=
 \int_G\kappa_\Omega(u)\,\dd u
 \int_G\mu_\alpha(\dd g)
 \int_G\mu_\alpha(\dd h)
 \notag\\
 &\quad\cdot
 \prod_{R\in\{A,B,C\}}
 \varphi_R\!\left(
 u\,g^{\mathbf1_{\{R\in J\}}}
 h^{\mathbf1_{\{R\in L\}}}
 \right).
 \label{eq:V75-Psi}
\end{align}
Exponent zero means that the corresponding factor is the identity;
the order in the character argument is \(u\), then \(g\), then
\(h\).

\begin{theorem}[Exact common-integral trace formula]
\label{thm:V75-common-integral}
For all \(J,L\subseteq\{A,B,C\}\),
\begin{equation}
 \boxed{
 \tau(\mathcal R_\Omega Q_JQ_L)
 =
 \Psi_\Omega(J,L).}
 \label{eq:V75-common-integral}
\end{equation}
In particular,
\(\Psi_\Omega(J,L)=\Psi_\Omega(L,J)\) is real.
\end{theorem}

\begin{proof}
Insert \eqref{eq:V75-QJ} and
\eqref{eq:V75-shell-functional-calculus}.  Fubini is immediate
because all spaces are finite dimensional and the measures are
finite.  On row \(R\), the product of the three representation
matrices is
\[
 D^R\!\left(
 u\,g^{\mathbf1_{\{R\in J\}}}
 h^{\mathbf1_{\{R\in L\}}}
 \right).
\]
Normalized trace factorizes over \(A,B,C\) and gives the product of
normalized characters in \eqref{eq:V75-Psi}.  This proves
\eqref{eq:V75-common-integral}.  Cyclicity of trace gives symmetry in
\(J,L\); the trace of a product of two self-adjoint operators against
the commuting self-adjoint \(\mathcal R_\Omega\) is real.
\end{proof}

\begin{remark}[Product-group coordinates]
\label{rem:V75-product-coordinates}
For a punctured core with several coordinate groups, apply
\eqref{eq:V75-common-integral} coordinatewise.  The measures,
characters, and shell kernel factor only when the output shell itself
factors.  A thermal shell defined by total mass generally does not
factor and must retain its single central kernel.
\end{remark}

\section{All three covariance cross traces}
\label{sec:V75-cross-traces}

Use the V74 signs
\[
 L_1=Q_{ABC}-xQ_{BC},\qquad
 L_2=-yQ_{AC}+xyzI,\qquad
 L_3=-zQ_{AB}+xyzI.
\]

\begin{theorem}[Exact character formulas for the cross energies]
\label{thm:V75-cross-formulas}
The normalized cross traces
\(\mathcal X_{ij}(\Omega)=
\operatorname{Re}\tau(\mathcal R_\Omega L_iL_j)\)
are
\begin{align}
 \mathcal X_{12}(\Omega)
 &=
 -y\Psi_\Omega(ABC,AC)
 +xyz\Psi_\Omega(ABC,\varnothing)
 \notag\\
 &\quad
 +xy\Psi_\Omega(BC,AC)
 -x^2yz\Psi_\Omega(BC,\varnothing),
 \label{eq:V75-X12}\\
 \mathcal X_{13}(\Omega)
 &=
 -z\Psi_\Omega(ABC,AB)
 +xyz\Psi_\Omega(ABC,\varnothing)
 \notag\\
 &\quad
 +xz\Psi_\Omega(BC,AB)
 -x^2yz\Psi_\Omega(BC,\varnothing),
 \label{eq:V75-X13}\\
 \mathcal X_{23}(\Omega)
 &=
 yz\Psi_\Omega(AC,AB)
 -xy^2z\Psi_\Omega(AC,\varnothing)
 \notag\\
 &\quad
 -xyz^2\Psi_\Omega(AB,\varnothing)
 +x^2y^2z^2\Psi_\Omega(\varnothing,\varnothing).
 \label{eq:V75-X23}
\end{align}
Together with \eqref{eq:V74-cross-expansion}, these formulas are an
exact scalar representation of the full shell Hilbert--Schmidt
energy.
\end{theorem}

\begin{proof}
Expand \(L_1L_2\):
\[
 L_1L_2
 =
 -yQ_{ABC}Q_{AC}
 +xyzQ_{ABC}
 +xyQ_{BC}Q_{AC}
 -x^2yzQ_{BC}.
\]
Apply \eqref{eq:V75-common-integral} term by term to obtain
\eqref{eq:V75-X12}.  The identical expansion of \(L_1L_3\) gives
\eqref{eq:V75-X13}.  Finally,
\[
 L_2L_3
 =
 yzQ_{AC}Q_{AB}
 -xy^2zQ_{AC}
 -xyz^2Q_{AB}
 +x^2y^2z^2I,
\]
which gives \eqref{eq:V75-X23}.
\end{proof}

\begin{corollary}[Exact cross-corrected thermal trace criterion]
\label{cor:V75-cross-corrected-target}
The V73 thermal Hilbert--Schmidt target holds if
\begin{align}
 &\mathcal E_{\rm cov}(\Omega_{\rm th})
 +2\mathcal X_{12}(\Omega_{\rm th})
 +2\mathcal X_{13}(\Omega_{\rm th})
 +2\mathcal X_{23}(\Omega_{\rm th})
 \notag\\
 &\hspace{7em}\le
 \frac1{\min\{d_A,d_B\}d_C^2}.
 \label{eq:V75-cross-corrected-target}
\end{align}
Every term on the left now has an exact finite group-integral
formula.
\end{corollary}

\begin{proof}
Use the equality \eqref{eq:V74-cross-expansion} in place of the
factor-three upper bound in the proof of
\Cref{cor:V74-energy-trace-gate}.
\end{proof}

\section{Sign audit and route decision}
\label{sec:V75-sign-audit}

\begin{proposition}[Centrality gives no cross-term sign]
\label{prop:V75-no-sign}
Self-adjointness, centrality of the Wilson law, and positivity of
\(\mathcal R_\Omega\) do not imply
\(\mathcal X_{ij}(\Omega)\le0\) or
\(\mathcal X_{ij}(\Omega)\ge0\).
\end{proposition}

\begin{proof}
Even in a one-dimensional commuting block, take real numbers
\(\ell_i\) and a positive scalar \(r\).  The cross trace is
\(r\ell_i\ell_j\), which has either sign as the two scalars vary.
All formal properties listed in the statement hold in this special
case.  Therefore a sign theorem requires a specific relation among
the Yang--Mills multipliers beyond centrality and positivity.
\end{proof}

\begin{assessment}[Outcome of the V74 cross-trace test]
\label{ass:V75-cross-outcome}
Equations \eqref{eq:V75-X12}--\eqref{eq:V75-X23} complete the exact
cross-trace calculation requested by V74.  They reveal no
sign-definite cancellation at the structural level.  Proving
\eqref{eq:V75-cross-corrected-target} would require a new
three-variable normalized-character inequality at the very strong
scale \((d_{\min}d_C^2)^{-1}\).  No preceding theorem supplies such
an inequality.

Accordingly the thermal trace route is suspended.  The next attack
returns to the product-vector quantity
\eqref{eq:V72-CG-target}, where a small product overlap can be used
without paying the full multiplicity trace.
\end{assessment}

\begin{firewall}[The exact singlet remains separately closed]
\label{fire:V75-singlet-separation}
The route decision concerns only the growing thermal shell.  It does
not reopen the exact singlet and fixed-window debits proved in V71.
As in the SM V80 zero-mode split, no thermal inverse or shell limit is
used to define the exact \(T=\mathbf1\) fibre.
\end{firewall}

\section{V75 ledger and next acquisition}
\label{sec:V75-ledger}

\begin{theorem}[Integrated V75 statement]
\label{thm:V75-integrated}
Preserving V1--V74, V75 proves:
\begin{enumerate}[label=\textup{(V75.\arabic*)},leftmargin=5.7em]
\item the stable mandatory SM V80/NS V31 audit with a strict
      nontransfer firewall;
\item the central realization \eqref{eq:V75-K-central} of all moment
      operators;
\item the finite shell functional calculus
      \eqref{eq:V75-shell-functional-calculus};
\item the common three-variable character integral
      \eqref{eq:V75-common-integral};
\item the exact cross traces
      \eqref{eq:V75-X12}--\eqref{eq:V75-X23};
\item the cross-corrected thermal trace target
      \eqref{eq:V75-cross-corrected-target}; and
\item the rigorous conclusion that centrality alone supplies no sign,
      so the product-vector route is the next noncircular attack.
\end{enumerate}
\end{theorem}

\begin{proof}
Items \textup{(V75.2)--(V75.6)} are
\Cref{lem:V75-moment-identification,
lem:V75-shell-functional-calculus,thm:V75-common-integral,
thm:V75-cross-formulas,cor:V75-cross-corrected-target}.
Item \textup{(V75.7)} is
\Cref{prop:V75-no-sign,ass:V75-cross-outcome}.
\end{proof}

\begin{openproblem}[V76 mandate: product-visible thermal ancestry]
\label{op:V76-mandate}
Abandon the full thermal Hilbert--Schmidt trace unless a new
character inequality is independently proved.  Instead:
\begin{enumerate}
\item express the thermal part of
      \(\Gamma_{ab}^{\otimes}\) directly as a supremum over product
      vectors and the V70 Racah POVM;
\item retain the distribution
      \(\|P_S^{AB}(u\otimes v)\|^2\) over all \(S\), rather than
      selecting the maximal ancestry block;
\item isolate Cartan product-visible ancestries from entangled
      low-output ancestries and assign the V69 covariance factors to
      the same split;
\item prove a fixed-orientation bound or return the first explicit
      product-vector family which saturates the thermal square;
\item keep puncture banks and later collision sectors suspended.
\end{enumerate}
The next companion audit is V80.
\end{openproblem}

\section{Firewall after V75}
\label{sec:V75-firewall}

\begin{firewall}[Current claim boundary]
\label{fire:V75-final}
V75 computes the covariance cross traces exactly and rules out a
purely formal sign argument.  It does not estimate the thermal
product-vector square, prove JREC, close simultaneous sideways
endpoints, aggregate puncture or multiple-collision sectors, complete
the remaining partitions, construct the cofinal Yang--Mills theory,
or prove a continuum mass gap.
\end{firewall}

\section{Append-only record for V75}
\label{sec:V75-record}

V75 was appended to the complete V74 manuscript.  The exact V74
parent had \(53457\) lines, \(1916687\) bytes, and SHA--256 digest
\[
 \texttt{44d9ce735e75f464ef6739104b25d81a1b6b003295bd826c33e57c0bfa1fbe35}.
\]
No V74 mathematical content was changed.  The sole final
\verb|\end{document}| is restored below.

% =============================================================================
% END APPENDED VERSION V75 MATERIAL
% =============================================================================
% =============================================================================
% BEGIN APPENDED VERSION V76 MATERIAL
% =============================================================================

\chapter{V76: Product-channel disintegration and the exact coherent
dual thermal law}
\label{ch:V76}

\section{Purpose and route}
\label{sec:V76-purpose}

V75 showed that centrality alone does not control the signs of the
three covariance cross traces.  This version therefore returns to the
exact product-vector target \eqref{eq:V72-CG-target}.  It first
disintegrates that target simultaneously over the retained
\(AB\)-channel \(S\) and the V70 Racah outcomes \((V,T)\).  It then
proves a general product-visibility bound for each \(S\).

The decisive test is the multiplicity-free dual family
\[
 (n,0)\otimes(0,n)=\bigoplus_{k=0}^{n}(k,k).
\]
For a concrete coherent product vector, the probability of the
\((k,k)\) channel is computed exactly.  Its mass is concentrated at
\(k\asymp\sqrt n\).  Thus the mesoscopic scale detected in V71 is
not merely an artefact of a full trace: it can be visible to an
admissible product vector.

\begin{firewall}[Scope of V76]
\label{fire:V76-scope}
The coherent calculation below is an exact obstruction to obtaining
a decaying estimate from channel visibility alone.  It is not a
lower bound on the complete signed cumulant square, because the V69
three-square domination may discard cancellations between its three
binary-cut summands.
\end{firewall}

\section{Exact product-vector and Racah disintegration}
\label{sec:V76-disintegration}

For unit vectors \(u\in H_A\), \(v\in H_B\), retain the V72 notation
\[
 c_S(u,v)=P_S^{AB}(u\otimes v)
 \in H_S\otimes\mathbb C^{N_{AB}^{S}}
\]
and put
\begin{equation}
 p_S(u,v):=\|c_S(u,v)\|^2.
 \label{eq:V76-channel-law}
\end{equation}
When \(p_S(u,v)>0\), define the multiplicity density matrix
\begin{equation}
 \sigma_S(u,v):=
 \frac1{p_S(u,v)}
 \operatorname{Tr}_{H_S}
 |c_S(u,v)\rangle\langle c_S(u,v)|.
 \label{eq:V76-multiplicity-state}
\end{equation}
When \(p_S=0\), choose any density matrix on the nonzero
multiplicity space; it will always be multiplied by \(p_S\).

\begin{theorem}[Exact two-level product law]
\label{thm:V76-two-level-law}
For every unit product vector:
\begin{enumerate}[label=\textup{(\roman*)}]
\item \(p_S(u,v)\ge0\) and \(\sum_Sp_S(u,v)=1\);
\item the full V72 square has the exact disintegration
\begin{equation}
 \left\langle u\otimes v,M_{ab}^{(2)}u\otimes v\right\rangle
 =
 \sum_Sp_S(u,v)
 \operatorname{Tr}\!\left[
 \mathscr S_S^\sim\sigma_S(u,v)\right];
 \label{eq:V76-Gamma-disintegration}
\end{equation}
\item the numbers
\begin{equation}
 \lambda_{u,v}(S,V,T):=
 p_S(u,v)
 \operatorname{Tr}\!\left[
 \Pi_S(V,T)\sigma_S(u,v)\right]
 \label{eq:V76-joint-Racah-law}
\end{equation}
form a probability law, with
\(\sum_{V,T}\lambda_{u,v}(S,V,T)=p_S(u,v)\); and
\item the composite-cut positive square satisfies the exact identity
\begin{equation}
 \boxed{
 \left\langle u\otimes v,
 \operatorname{Tr}_C G_{A\mid BC}\,
 u\otimes v\right\rangle
 =
 d_C\sum_{S,V,T}
 \lambda_{u,v}(S,V,T)f_{A,V}(T).}
 \label{eq:V76-composite-product-law}
\end{equation}
\end{enumerate}
\end{theorem}

\begin{proof}
The orthogonal projectors \(P_S^{AB}\), including all multiplicity
copies inside the \(S\)-isotypic projection, resolve the identity on
\(H_A\otimes H_B\).  This proves the first assertion.

Insert \eqref{eq:V72-second-blocks} and decompose \(u\otimes v\) by
the same projectors.  For an operator of the form
\(I_{H_S}\otimes X_S\), the defining property of partial trace gives
\[
 \langle c_S,(I_{H_S}\otimes X_S)c_S\rangle
 =
 p_S\operatorname{Tr}(X_S\sigma_S).
\]
Summation proves \eqref{eq:V76-Gamma-disintegration}.

Every V70 atom \(\Pi_S(V,T)\) is positive, so
\(\lambda_{u,v}\ge0\).  The POVM normalization
\eqref{eq:V70-POVM-normalization} gives
\[
 \sum_{V,T}\lambda_{u,v}(S,V,T)
 =
 p_S\operatorname{Tr}\sigma_S=p_S.
\]
Summing also in \(S\) proves that \(\lambda_{u,v}\) has total mass
one.  Finally insert the exact conditional formula
\eqref{eq:V70-composite-partial-trace} into the preceding
partial-trace identity.  The result is
\eqref{eq:V76-composite-product-law}.
\end{proof}

\begin{definition}[Three fixed-orientation product capacities]
\label{def:V76-three-capacities}
Set
\begin{align}
 \mathfrak G_{A\mid BC}^{\otimes}
 &:=
 \sup_{u,v}
 d_C\sum_{S,V,T}
 \lambda_{u,v}(S,V,T)f_{A,V}(T),
 \label{eq:V76-G-composite}\\
 \mathfrak G_{AC}^{\otimes}
 &:=
 \sup_{u,v}
 \sum_Sp_S(u,v)
 \operatorname{Tr}\!\left[
 \mathscr D_{AC;S}\sigma_S(u,v)\right],
 \label{eq:V76-G-AC}\\
 \mathfrak G_{AB}^{\otimes}
 &:=
 \sup_{u,v}
 \sum_Sp_S(u,v)
 \delta_{AB}(S)^2M_2(S,C).
 \label{eq:V76-G-AB}
\end{align}
All suprema are over unit \(u\in H_A\), \(v\in H_B\).
\end{definition}

\begin{corollary}[Complete \(ab\)-oriented three-square insertion]
\label{cor:V76-three-square-insertion}
The product-vector square obeys
\begin{equation}
 \boxed{
 \Gamma_{ab}^{\otimes}
 \le
 3\left[
 \mathfrak G_{A\mid BC}^{\otimes}
 +y^2\mathfrak G_{AC}^{\otimes}
 +z^2\mathfrak G_{AB}^{\otimes}
 \right].}
 \label{eq:V76-three-square-insertion}
\end{equation}
Moreover,
\begin{align}
 \mathfrak G_{AC}^{\otimes}
 &\le
 Cd_C[1+s_\alpha(A_A+A_B+A_C)]^2
 \min\{s_\alpha A_AA_C,1\}^2,
 \label{eq:V76-G-AC-bound}\\
 \mathfrak G_{AB}^{\otimes}
 &\le
 Cd_C
 \sup_{u,v}\sum_Sp_S(u,v)
 [1+s_\alpha(A_S+A_C)]^2
 \delta_{AB}(S)^2.
 \label{eq:V76-G-AB-bound}
\end{align}
\end{corollary}

\begin{proof}
Evaluate the positive block inequality
\eqref{eq:V69-net-three} in every multiplicity state
\(\sigma_S(u,v)\), multiply by \(p_S(u,v)\), and sum in \(S\).
The three terms are precisely
\eqref{eq:V76-G-composite}--\eqref{eq:V76-G-AB}; for the last one use
the exact aligned formula
\eqref{eq:V69-aligned-square}.  Taking the supremum proves
\eqref{eq:V76-three-square-insertion}.  The scalar crossing estimate
\eqref{eq:V69-AC-crossing-bound} proves
\eqref{eq:V76-G-AC-bound}.  The output second-moment estimate
\eqref{eq:V69-output-second-bound} proves
\eqref{eq:V76-G-AB-bound}.
\end{proof}

\begin{remark}[What has remained correlated]
\label{rem:V76-correlation-retained}
Equation \eqref{eq:V76-composite-product-law} does not maximize first
over \(S\), \(V\), or \(T\).  The product-channel law \(p_S\), its
multiplicity state \(\sigma_S\), the Racah POVM, the final output,
and the covariance symbol all remain in the same probability
measure \(\lambda_{u,v}\).
\end{remark}

\section{A channel-by-channel product-visibility theorem}
\label{sec:V76-visibility}

Define the exact product visibility
\begin{equation}
 \beta_{AB}(S):=
 \sup_{\|u\|=\|v\|=1}
 \|P_S^{AB}(u\otimes v)\|^2.
 \label{eq:V76-beta}
\end{equation}

\begin{theorem}[Dimension cap for every retained channel]
\label{thm:V76-visibility-cap}
For every \(S\subset A\otimes B\),
\begin{equation}
 \boxed{
 \beta_{AB}(S)
 \le
 \min\left\{
 1,\frac{d_SN_{AB}^{S}}{\max\{d_A,d_B\}}
 \right\}.}
 \label{eq:V76-visibility-cap}
\end{equation}
The two alternatives are both necessary: the value is one on the
Cartan highest-weight summand, while it is \(1/d_A\) on the singlet
in \(A\otimes A^*\).
\end{theorem}

\begin{proof}
Apply the double-Schur inequality
\eqref{eq:V73-double-Schur} to the invariant positive operator
\(P_S^{AB}\).  Its operator norm is one and its trace is
\(d_SN_{AB}^{S}\), which gives
\eqref{eq:V76-visibility-cap}.  Sharpness on the two named channels
is exactly
\Cref{prop:V72-highest-product-norm,
prop:V72-singlet-product-norm}.
\end{proof}

For a set \(E\) of \(AB\)-channels, write
\[
 B_{AB}(E):=\sum_{S\in E}\beta_{AB}(S),
 \qquad
 \phi_S:=
 \left\|
 \sum_{V,T}f_{A,V}(T)\Pi_S(V,T)
 \right\|_{\rm op}.
\]

\begin{corollary}[Product-visible ancestry split]
\label{cor:V76-ancestry-split}
For every product vector and every channel set \(E\),
\begin{equation}
 \sum_{S\in E,V,T}\lambda_{u,v}(S,V,T)
 \le
 \min\{1,B_{AB}(E)\}.
 \label{eq:V76-set-mass}
\end{equation}
The contribution of \(E\) to the composite square satisfies
\begin{align}
 &d_C\sum_{S\in E,V,T}
 \lambda_{u,v}(S,V,T)f_{A,V}(T)
 \notag\\
 &\quad\le
 d_C\min\left\{
 \sum_{S\in E}\beta_{AB}(S)\phi_S,\,
 \left(\max_{S\in E}\phi_S\right)
 \min\{1,B_{AB}(E)\}
 \right\}.
 \label{eq:V76-visible-composite-bound}
\end{align}
Thus low-visibility ancestries and product-visible Cartan
ancestries may be separated without changing orientation or losing
the covariance factor.
\end{corollary}

\begin{proof}
The \(S\)-marginal of \(\lambda_{u,v}\) is \(p_S(u,v)\), by
\Cref{thm:V76-two-level-law}.  Definition
\eqref{eq:V76-beta} gives
\(p_S(u,v)\le\beta_{AB}(S)\), while
\(\sum_Sp_S=1\), proving \eqref{eq:V76-set-mass}.

Conditional on \(S\),
\[
 \sum_{V,T}f_{A,V}(T)
 \operatorname{Tr}[\Pi_S(V,T)\sigma_S]
 \le\phi_S.
\]
Hence the left side of
\eqref{eq:V76-visible-composite-bound} is at most both
\[
 d_C\sum_{S\in E}\beta_{AB}(S)\phi_S
\]
and
\[
 d_C\left(\max_{S\in E}\phi_S\right)
 \sum_{S\in E}p_S.
\]
Use \eqref{eq:V76-set-mass} in the second bound and take the minimum.
\end{proof}

\begin{assessment}[The Cartan/low-output division]
\label{ass:V76-Cartan-division}
The Cartan channel \(S=A+B\) has \(\beta_{AB}(S)=1\), so it cannot be
discounted by product entanglement.  It must instead be controlled
by its own conditional covariance factor \(\phi_{A+B}\), normally
through high-output Wilson damping.

At the opposite end, one low-dimensional channel has a small
\(\beta_{AB}(S)\), but a thermal window contains many channels.
The sum \(B_{AB}(E)\), not the visibility of a single member of
\(E\), is the relevant debit.  The next section shows that this
cumulative distinction is sharp at the mesoscopic dual scale.
\end{assessment}

\section{Exact coherent product law for the dual family}
\label{sec:V76-coherent-dual}

Let
\[
 R_n=(n,0)=\operatorname{Sym}^n(\mathbb C^3),
 \qquad
 Q_n=(0,n)=R_n^*,
 \qquad
 d_n=\frac{(n+1)(n+2)}2.
\]
The Littlewood--Richardson rule gives the multiplicity-free
decomposition
\begin{equation}
 R_n\otimes Q_n
 =
 \bigoplus_{k=0}^{n}T_k,
 \qquad T_k=(k,k),
 \qquad d_{T_k}=(k+1)^3.
 \label{eq:V76-dual-decomposition}
\end{equation}
Let \(e\in\mathbb C^3\) be a unit vector,
\(u_n=e^{\odot n}\in R_n\), and set
\[
 \eta_n:=u_n\otimes u_n^*\in R_n\otimes Q_n.
\]
This is a unit product vector.  Denote by \(P_{n,k}\) the projection
onto \(T_k\).

\begin{theorem}[Exact coherent dual channel probabilities]
\label{thm:V76-coherent-law}
For \(0\le k\le n\),
\begin{equation}
 \boxed{
 p_{n,k}:=\|P_{n,k}\eta_n\|^2
 =
 \frac{2(k+1)^3(n!)^2}
 {(n-k)!(n+k+2)!}.}
 \label{eq:V76-coherent-law}
\end{equation}
In particular \(p_{n,k}>0\) and
\(\sum_{k=0}^{n}p_{n,k}=1\).
\end{theorem}

\begin{proof}
Let \(H=S(U(1)\times U(2))\) be the stabilizer of the line
\(\mathbb Ce\).  The vector \(\eta_n\) is \(H\)-fixed.  If
\(t=|\langle e,ge\rangle|^2\), its diagonal matrix coefficient is
\begin{equation}
 \langle\eta_n,
 (R_n(g)\otimes Q_n(g))\eta_n\rangle=t^n.
 \label{eq:V76-coherent-coefficient}
\end{equation}
Under normalized Haar measure on \(\mathrm{SU}(3)\), the first
column is uniform on the unit sphere in \(\mathbb C^3\).  Therefore
\(t\) has density \(2(1-t)\,\dd t\) on \([0,1]\).

The compact rank-one pair \((\mathrm{SU}(3),H)\) has one
\(H\)-fixed line in each \(T_k=(k,k)\).  Its normalized spherical
function is
\begin{equation}
 \varphi_k(t)
 =
 \frac{P_k^{(1,0)}(2t-1)}{k+1},
 \qquad
 \varphi_k(1)=1.
 \label{eq:V76-spherical-Jacobi}
\end{equation}
For completeness, this identification follows without a character
asymptotic: the \(H\)-bi-invariant coefficients appearing in
\(\bigoplus_{m\le n}R_m\otimes Q_m\) are the polynomials in \(t\) of
degree at most \(n\).  Distinct irreducible spherical coefficients
are orthogonal for the weight \(2(1-t)\).  The unique degree-\(k\)
orthogonal polynomial with value one at \(t=1\) is precisely the
normalized Jacobi polynomial in
\eqref{eq:V76-spherical-Jacobi}.

Schur orthogonality for normalized spherical coefficients gives
\[
 \int_G\varphi_k(g)\overline{\varphi_\ell(g)}\,\dd g
 =
 \frac{\delta_{k\ell}}{d_{T_k}}.
\]
Decomposing \(\eta_n\) according to
\eqref{eq:V76-dual-decomposition} and applying this orthogonality to
\eqref{eq:V76-coherent-coefficient} yields
\begin{align}
 p_{n,k}
 &=
 d_{T_k}\int_0^1
 2(1-t)t^n
 \frac{P_k^{(1,0)}(2t-1)}{k+1}\,\dd t.
 \label{eq:V76-probability-integral}
\end{align}
The shifted Rodrigues formula is
\[
 P_k^{(1,0)}(2t-1)
 =
 \frac{(-1)^k}{k!(1-t)}
 \frac{\dd^k}{\dd t^k}
 \left[(1-t)^{k+1}t^k\right].
\]
Integrating by parts \(k\) times in
\eqref{eq:V76-probability-integral}; all boundary terms vanish.
For \(k\le n\), the remaining beta integral gives
\begin{align*}
 &\int_0^1
 2(1-t)t^nP_k^{(1,0)}(2t-1)\,\dd t\\
 &\qquad=
 \frac{2n!}{k!(n-k)!}
 \int_0^1t^n(1-t)^{k+1}\,\dd t\\
 &\qquad=
 \frac{2(k+1)(n!)^2}
 {(n-k)!(n+k+2)!}.
\end{align*}
Since \(d_{T_k}=(k+1)^3\), division by \(k+1\) in
\eqref{eq:V76-probability-integral} proves
\eqref{eq:V76-coherent-law}.  Positivity is immediate.  Total mass
one also follows directly from the orthogonal resolution of
\(\eta_n\).
\end{proof}

\begin{corollary}[Thermal square-root limit law]
\label{cor:V76-thermal-limit}
If \(k/\sqrt n\to t\in[0,\infty)\), then
\begin{equation}
 \sqrt n\,p_{n,k}\longrightarrow
 2t^3e^{-t^2}.
 \label{eq:V76-local-limit}
\end{equation}
The convergence is uniform for \(0\le k\le M\sqrt n\), for each
fixed \(M\).  Consequently, for \(0\le a<b<\infty\),
\begin{equation}
 \boxed{
 \lim_{n\to\infty}
 \sum_{a\sqrt n\le k\le b\sqrt n}p_{n,k}
 =
 \int_a^b2t^3e^{-t^2}\,\dd t.}
 \label{eq:V76-window-limit}
\end{equation}
In particular
\[
 \lim_{n\to\infty}
 \sum_{k\le c\sqrt n}p_{n,k}
 =
 1-(1+c^2)e^{-c^2}.
\]
\end{corollary}

\begin{proof}
Rewrite \eqref{eq:V76-coherent-law} as
\begin{equation}
 p_{n,k}
 =
 \frac{2(k+1)^3}{(n+1)(n+2)}
 \prod_{j=0}^{k-1}
 \frac{n-j}{n+j+3}.
 \label{eq:V76-product-form}
\end{equation}
Uniformly for \(k\le M\sqrt n\), Taylor expansion of the logarithm
gives
\[
 \log\prod_{j=0}^{k-1}\frac{n-j}{n+j+3}
 =
 -\frac{k^2+2k}{n}+O_M(n^{-1/2}).
\]
Multiplying the resulting exponential by the prefactor in
\eqref{eq:V76-product-form} proves the uniform local limit
\eqref{eq:V76-local-limit}.  The window sum is then a Riemann sum on
every bounded interval.  Integration gives
\[
 \int_0^c2t^3e^{-t^2}\,\dd t
 =
 1-(1+c^2)e^{-c^2}.
\]
\end{proof}

\begin{corollary}[Cumulative visibility has no dimension decay]
\label{cor:V76-cumulative-visibility}
For \(K\le n\), the general visibility theorem gives
\begin{equation}
 \sum_{k=0}^{K}p_{n,k}
 \le
 \min\left\{
 1,\frac1{d_n}\sum_{k=0}^{K}(k+1)^3
 \right\}
 =
 \min\left\{
 1,
 \frac{[(K+1)(K+2)/2]^2}{d_n}
 \right\}.
 \label{eq:V76-cumulative-cap}
\end{equation}
For \(K=c\sqrt n\), the untruncated second entry converges to
\(c^4/2\), while the exact coherent mass converges to the positive
quantity \(1-(1+c^2)e^{-c^2}\).  Hence neither the exact product law
nor its channelwise dimension cap supplies any power of \(n^{-1}\)
on a fixed square-root window.
\end{corollary}

\begin{proof}
Here every multiplicity is one and
\(\max\{d_{R_n},d_{Q_n}\}=d_n\).  Apply
\eqref{eq:V76-visibility-cap}, sum, and use
\(\sum_{j=1}^{K+1}j^3=[(K+1)(K+2)/2]^2\).
The limits follow from elementary division and
\Cref{cor:V76-thermal-limit}.
\end{proof}

\section{Consequence for the covariance-square route}
\label{sec:V76-consequence}

In the multiplicity-free dual family, the V70 conditional operator
on each retained \(T_k\)-channel is scalar.  For fixed third input
\(C\), write
\begin{equation}
 \phi_{n,k}^{C}:=
 \sum_{V,T}f_{R_n,V}(T)\Pi_{T_k}(V,T).
 \label{eq:V76-conditional-phi-dual}
\end{equation}
This is a nonnegative number.

\begin{proposition}[Exact coherent test of the composite square]
\label{prop:V76-coherent-composite-test}
For \(A=R_n\), \(B=Q_n\), and the coherent product vector
\(\eta_n\),
\begin{equation}
 \boxed{
 \left\langle\eta_n,
 \operatorname{Tr}_C G_{R_n\mid Q_nC}\,
 \eta_n\right\rangle
 =
 d_C\sum_{k=0}^{n}p_{n,k}\phi_{n,k}^{C}.}
 \label{eq:V76-coherent-composite-test}
\end{equation}
If, for some \(0<a<b<\infty\) and \(\varepsilon>0\),
\[
 \inf_{a\sqrt n\le k\le b\sqrt n}\phi_{n,k}^{C}
 \ge\varepsilon
\]
along an infinite sequence of \(n\), then along that sequence
\begin{align}
 \liminf_{n\to\infty}
 \frac1{d_C}
 \left\langle\eta_n,
 \operatorname{Tr}_C G_{R_n\mid Q_nC}\,
 \eta_n\right\rangle
 \ge
 \varepsilon\int_a^b2t^3e^{-t^2}\,\dd t>0.
 \label{eq:V76-coherent-obstruction}
\end{align}
\end{proposition}

\begin{proof}
Since \(N_{R_nQ_n}^{T_k}=1\), the multiplicity density matrix in
each nonzero channel is the scalar one.  Insert the \(S\)-marginal
\(p_{n,k}\) from \eqref{eq:V76-coherent-law} into
\eqref{eq:V76-composite-product-law}; this proves
\eqref{eq:V76-coherent-composite-test}.  Restrict its nonnegative sum
to the stated window and apply
\eqref{eq:V76-window-limit}.
\end{proof}

\begin{assessment}[The new exact dichotomy]
\label{ass:V76-dichotomy}
The product-vector route has not failed, but generic
Clebsch--Gordan entanglement is no longer an available explanation
for the missing thermal debit.  The coherent product vector sees an
order-one portion of the same square-root channel window.

There are now two honest possibilities.  Either the conditional
covariance numbers \(\phi_{n,k}^{C}\) decay on that window at the
scale required by the collision gate, or the positive
three-square domination \eqref{eq:V76-three-square-insertion} is too
coarse and one must evaluate the signed cumulant \(K\) on the same
coherent vector before squaring.  Passing again to a full trace
would erase the diagnostic and is not a valid third option.
\end{assessment}

\begin{firewall}[What the coherent lower bound does not say]
\label{fire:V76-no-full-lower-bound}
Equation \eqref{eq:V76-coherent-obstruction} is a lower bound only
for the composite positive square \(G_{A\mid BC}\).  Since
\[
 K=D_{A\mid BC}-yD_{AC}-zD_{AB},
\]
the complete \(K^2\) may be smaller through vector-level
cancellation.  V76 therefore does not infer a lower bound for
\(\Gamma_{ab}^{\otimes}\), nor does it declare the Yang--Mills
programme impossible.
\end{firewall}

\section{V76 ledger and next acquisition}
\label{sec:V76-ledger}

\begin{theorem}[Integrated V76 statement]
\label{thm:V76-integrated}
Preserving V1--V75, V76 proves:
\begin{enumerate}[label=\textup{(V76.\arabic*)},leftmargin=5.7em]
\item the exact product/Racah probability law
      \eqref{eq:V76-joint-Racah-law};
\item the composite product expectation
      \eqref{eq:V76-composite-product-law};
\item the fixed-orientation three-square insertion
      \eqref{eq:V76-three-square-insertion};
\item the channel visibility cap
      \eqref{eq:V76-visibility-cap} and ancestry split
      \eqref{eq:V76-visible-composite-bound};
\item the exact coherent dual law
      \eqref{eq:V76-coherent-law};
\item its square-root thermal limit
      \eqref{eq:V76-window-limit}; and
\item the exact coherent diagnostic
      \eqref{eq:V76-coherent-composite-test}.
\end{enumerate}
\end{theorem}

\begin{proof}
These are
\Cref{thm:V76-two-level-law,cor:V76-three-square-insertion,
thm:V76-visibility-cap,cor:V76-ancestry-split,
thm:V76-coherent-law,cor:V76-thermal-limit,
prop:V76-coherent-composite-test}.
\end{proof}

\begin{openproblem}[V77 mandate: covariance on the coherent thermal
window]
\label{op:V77-mandate}
Keep the \(ab\) orientation and the coherent dual product vector
fixed.  Then:
\begin{enumerate}
\item choose the actual third input \(C\) arising in the
      simultaneous-sideways collision monomial, rather than a free
      proxy;
\item compute the conditional Racah scalar
      \(\phi_{n,k}^{C}\) for \(k\asymp\sqrt n\), retaining
      \((V,T)\) until the Wilson covariance is inserted;
\item if its coherent average has the required decay, prove that
      decay uniformly and insert it in
      \eqref{eq:V76-three-square-insertion};
\item if it does not, return upstream from the positive
      three-square bound and compute
      \(\langle\eta_n,\operatorname{Tr}_C K^2\eta_n\rangle\)
      with all three signed binary cuts in one Racah representation;
\item do not reopen puncture banks or later collision sectors before
      this fixed-orientation decision is complete.
\end{enumerate}
The next compulsory companion audit remains V80.
\end{openproblem}

\section{Firewall after V76}
\label{sec:V76-firewall}

\begin{firewall}[Current claim boundary]
\label{fire:V76-final}
V76 proves that an admissible coherent product vector sees the
dual mesoscopic thermal window with order-one probability.  It does
not yet estimate the Wilson covariance on that law, close the
product-state collision gate, prove JREC, close simultaneous
sideways endpoints, aggregate punctures or multiple collisions,
construct the cofinal Yang--Mills theory, or prove a positive
continuum mass gap.
\end{firewall}

\section{Append-only record for V76}
\label{sec:V76-record}

V76 was appended to the complete V75 manuscript.  The exact V75
parent had \(53926\) lines, \(1931673\) bytes, and SHA--256 digest
\[
 \texttt{97cbe592228834c7339bd17c092f81bffe568b937e4b27e4531e3b20eb266272}.
\]
No V75 mathematical content was changed.  The sole final
\verb|\end{document}| is restored below.

% =============================================================================
% END APPENDED VERSION V76 MATERIAL
% =============================================================================
% =============================================================================
% BEGIN APPENDED VERSION V77 MATERIAL
% =============================================================================

\chapter{V77: Ancestry-layer correction and an exact signed
product-vector character calculus}
\label{ch:V77}

\section{Purpose and correction of the regression layer}
\label{sec:V77-purpose}

V76 computed a new and exact fact: a coherent product vector in
\(R_n\otimes Q_n\) distributes its \(AB\)-intermediate over
\((k,k)\) at the square-root scale.  The earlier V66--V71
dual-symmetric obstruction, however, lives one layer later:
\begin{equation}
 A\otimes B\longrightarrow S=(n,0),
 \qquad
 S\otimes C,\quad C=(0,n),
 \longrightarrow T=(k,k).
 \label{eq:V77-actual-dual-layer}
\end{equation}
Thus the V76 probability law is valid but is not the conditional law
in \eqref{eq:V77-actual-dual-layer}.  This version records that
distinction and returns one step upstream.  It derives the exact
product-vector analogue of the V75 common character integral for the
\emph{complete signed} cumulant square.  No positive three-square
domination is used.

\begin{firewall}[Nature of the correction]
\label{fire:V77-layer-correction}
No V76 formula is withdrawn.  Its coherent law concerns the
distribution of \(S\) inside a chosen \(A\otimes B\).  The old
regression fixes \(S\) and \(C\) and leaves the lift
\(A\otimes B\to S\) to be controlled uniformly.  Confusing these two
conditional levels would be an invalid substitution.
\end{firewall}

\section{The lifting fibre of a retained intermediate}
\label{sec:V77-lifting-fibre}

For an irreducible \(S\), define its ordered lifting fibre
\begin{equation}
 \mathfrak L(S):=
 \left\{(A,B,\alpha):
 1\le\alpha\le N_{AB}^{S}\right\}.
 \label{eq:V77-lifting-fibre}
\end{equation}
In the collision core, \(A,B,C\) are the punctured-core
representations of the original rows \(a,b,c\); the \(oa\) and
\(bc\) collision pair banks remain outside this ternary core as in
\eqref{eq:V72-three-bank-product}.

\begin{proposition}[What the dual regression fixes]
\label{prop:V77-dual-data}
The V66--V71 dual regression fixes
\[
 S=R_n,\qquad C=Q_n,\qquad
 T=T_k,\quad k\asymp\sqrt n,
\]
but it does not fix a unique member \((A,B,\alpha)\) of
\(\mathfrak L(R_n)\).  Consequently a product-state estimate for the
actual collision must be uniform over the lifting fibre and must
retain the correlation between:
\begin{enumerate}[label=\textup{(\roman*)}]
\item the product visibility of \(A\otimes B\to R_n\);
\item the three signed covariance cuts in \(K\); and
\item the pair factors \(P_{ab}\) belonging to the same row
      orientation.
\end{enumerate}
\end{proposition}

\begin{proof}
The sequential construction
\eqref{eq:V63-sequential-embedding} first chooses
\(\iota_{S,\alpha}:H_S\to H_A\otimes H_B\), then
\(\jmath_{T,\beta}^{S}:H_T\to H_S\otimes H_C\).
The dual-symmetric specification in
\Cref{ass:V66-mesoscopic-Racah,op:V67} concerns only the second
arrow.  The first arrow and its original row labels remain among the
uniform collision data.  The three listed correlations are exactly
the data in \eqref{eq:V72-CG-target},
\eqref{eq:V57-cut-identity}, and
\eqref{eq:V72-pair-core-product}, respectively.
\end{proof}

\section{A product-vector central kernel}
\label{sec:V77-product-kernel}

Retain the V75 notation \(Q_J\), \(\mathcal R_\Omega\), and
\(\kappa_\Omega\).  Let
\[
 \rho_{u,v;C}:=
 |u\rangle\langle u|\otimes
 |v\rangle\langle v|\otimes
 \frac{I_C}{d_C}.
\]
For a representation \(R\) and unit vector \(\xi\in H_R\), write
\begin{equation}
 m_{R,\xi}(g):=
 \langle\xi,D^R(g)\xi\rangle.
 \label{eq:V77-vector-coefficient}
\end{equation}
For \(J,L\subseteq\{A,B,C\}\), define the ordered scalar
\begin{align}
 \Theta_{\Omega;u,v}(J,L)
 &:=
 \operatorname{Tr}\left[
 \rho_{u,v;C}\mathcal R_\Omega Q_JQ_L\right].
 \label{eq:V77-Theta-operator}
\end{align}

\begin{theorem}[Exact product-vector common integral]
\label{thm:V77-product-common-integral}
For all \(J,L\),
\begin{align}
 \Theta_{\Omega;u,v}(J,L)
 &=
 \int_G\kappa_\Omega(w)\,\dd w
 \int_G\mu_\alpha(\dd g)
 \int_G\mu_\alpha(\dd h)
 \notag\\
 &\quad\cdot
 m_{A,u}\!\left(
 w g^{\mathbf1_{\{A\in J\}}}
 h^{\mathbf1_{\{A\in L\}}}\right)
 \notag\\
 &\quad\cdot
 m_{B,v}\!\left(
 w g^{\mathbf1_{\{B\in J\}}}
 h^{\mathbf1_{\{B\in L\}}}\right)
 \notag\\
 &\quad\cdot
 \varphi_C\!\left(
 w g^{\mathbf1_{\{C\in J\}}}
 h^{\mathbf1_{\{C\in L\}}}\right).
 \label{eq:V77-product-common-integral}
\end{align}
Moreover
\begin{equation}
 \Theta_{\Omega;u,v}(L,J)
 =
 \overline{\Theta_{\Omega;u,v}(J,L)}.
 \label{eq:V77-Theta-Hermitian}
\end{equation}
\end{theorem}

\begin{proof}
Insert the central representations
\eqref{eq:V75-QJ} and
\eqref{eq:V75-shell-functional-calculus} in
\eqref{eq:V77-Theta-operator}.  The \(A\)- and \(B\)-row traces
against the two rank-one densities are the vector coefficients in
\eqref{eq:V77-vector-coefficient}.  The normalized \(C\)-trace is
its normalized character.  Fubini proves
\eqref{eq:V77-product-common-integral}.

All \(Q_J\) and \(\mathcal R_\Omega\) are self-adjoint, and
\(\mathcal R_\Omega\) commutes with every invariant \(Q_J\).  Hence
\begin{align*}
 \overline{\operatorname{Tr}[
 \rho_{u,v;C}\mathcal R_\Omega Q_JQ_L]}
 &=
 \operatorname{Tr}[Q_LQ_J\mathcal R_\Omega\rho_{u,v;C}]\\
 &=
 \operatorname{Tr}[
 \rho_{u,v;C}\mathcal R_\Omega Q_LQ_J],
\end{align*}
by cyclicity, proving \eqref{eq:V77-Theta-Hermitian}.
\end{proof}

\section{The complete signed cumulant square}
\label{sec:V77-signed-square}

Let
\[
 \mathcal J:=
 \{ABC,BC,AC,AB,\varnothing\}
\]
and assign the real coefficients
\begin{equation}
 c_{ABC}=1,\qquad
 c_{BC}=-x,\qquad
 c_{AC}=-y,\qquad
 c_{AB}=-z,\qquad
 c_{\varnothing}=2xyz.
 \label{eq:V77-cumulant-coefficients}
\end{equation}
Then \(K=\sum_{J\in\mathcal J}c_JQ_J\).

\begin{theorem}[Exact signed product-shell energy]
\label{thm:V77-signed-product-energy}
For every output shell \(\Omega\) and unit product vector \(u\otimes
v\),
\begin{align}
 &\frac1{d_C}
 \left\langle u\otimes v,
 \operatorname{Tr}_C[
 \mathcal R_\Omega K^2]\,
 u\otimes v\right\rangle
 \notag\\
 &\qquad=
 \boxed{
 \sum_{J,L\in\mathcal J}
 c_Jc_L\Theta_{\Omega;u,v}(J,L)}
 \label{eq:V77-signed-product-energy}\\
 &\qquad=
 \sum_{J\in\mathcal J}c_J^2
 \Theta_{\Omega;u,v}(J,J)
 +
 2\sum_{\substack{J,L\in\mathcal J\\J<L}}
 c_Jc_L
 \operatorname{Re}\Theta_{\Omega;u,v}(J,L).
 \notag
\end{align}
In particular, with \(\Omega=\mathcal T(A,B,C)\),
\begin{equation}
 \boxed{
 \Gamma_{ab}^{\otimes}
 =
 d_C\sup_{\|u\|=\|v\|=1}
 \sum_{J,L\in\mathcal J}
 c_Jc_L\Theta_{\rm all;u,v}(J,L).}
 \label{eq:V77-Gamma-exact-integral}
\end{equation}
\end{theorem}

\begin{proof}
Expand \(K^2\) using
\eqref{eq:V77-cumulant-coefficients}, evaluate against
\(\rho_{u,v;C}\), and use
\eqref{eq:V77-Theta-operator}.  This proves the first equality.
Equation \eqref{eq:V77-Theta-Hermitian} pairs the two ordered
off-diagonal terms and proves the real expansion.

For the full output family,
\[
 \mathcal R_{\rm all}
 =[1+s_\alpha\Xi(\mathbf T)]^2
\]
and
\(\mathcal R_{\rm all}K^2=(\mathbf K^\sim)^2\).
Convexity reduces the product-density supremum in
\eqref{eq:V72-Gamma-product} to pure product vectors.  The
normalization of \(\rho_{u,v;C}\) accounts for the displayed factor
\(d_C\), proving \eqref{eq:V77-Gamma-exact-integral}.
\end{proof}

\begin{corollary}[No Racah entrywise estimate is required]
\label{cor:V77-no-entrywise-Racah}
The fixed-orientation product square can be studied through the
single scalar integral
\eqref{eq:V77-Gamma-exact-integral}.  All final multiplicities and
all changes among \(AB\), \(AC\), and \(BC\) bracketings are already
contained in the common group variables \((w,g,h)\).  Taking an
absolute value before summing the \(25\) ordered pairs would revert
to the V69 positive-square loss.
\end{corollary}

\begin{proof}
The first assertion is the content of
\Cref{thm:V77-product-common-integral,
thm:V77-signed-product-energy}.  The second follows because
the \(25\) terms are exactly the expansion of the one operator
\(K^2\); termwise absolute values discard their cross signs.
\end{proof}

\section{Visibility and connectedness are correlated}
\label{sec:V77-connectedness}

\begin{proposition}[Exact deletion of a trivial input]
\label{prop:V77-trivial-input-deletion}
If any one of \(A,B,C\) is the trivial representation, then
\begin{equation}
 K=0
 \quad\hbox{on every final output and every multiplicity space}.
 \label{eq:V77-trivial-input-deletion}
\end{equation}
Consequently its contribution to every shell in
\eqref{eq:V77-signed-product-energy} is zero.
\end{proposition}

\begin{proof}
Suppose \(B=\mathbf1\).  Then \(y=1\),
\[
 Q_{ABC}=Q_{AC},\qquad
 Q_{BC}=zI,\qquad
 Q_{AB}=xI.
\]
Substitution in \eqref{eq:V75-K-central} gives
\[
 K=Q_{AC}-xzI-Q_{AC}-xzI+2xzI=0.
\]
The other two cases follow cyclically.  Multiplication by any shell
operator preserves zero.
\end{proof}

\begin{assessment}[Why maximal visibility is not by itself bad]
\label{ass:V77-visibility-connectedness}
The lift \(A=S\), \(B=\mathbf1\) has
\(\beta_{AB}(S)=1\), but its ternary connected coefficient vanishes
by \Cref{prop:V77-trivial-input-deletion}.  Thus maximizing product
visibility and then maximizing the cumulant independently is
provably illegitimate.  The sought estimate is a joint
visibility--connectedness inequality on \(\mathfrak L(S)\).
\end{assessment}

\section{The genuinely product-visible lift of the old dual shell}
\label{sec:V77-Cartan-lift}

Let \(a,b\ge0\), \(a+b=n\), and set
\[
 A=R_a=(a,0),\qquad
 B=R_b=(b,0),\qquad
 C=Q_n=(0,n).
\]
Choose the same unit vector \(e\in\mathbb C^3\) and the coherent
highest-weight vectors
\[
 u_a=e^{\odot a},\qquad v_b=e^{\odot b}.
\]

\begin{proposition}[Cartan lift with unit product visibility]
\label{prop:V77-Cartan-lift}
The product vector \(u_a\otimes v_b\) lies entirely in the Cartan
summand
\[
 R_a\otimes R_b\supset R_{a+b}=R_n.
\]
Hence
\begin{equation}
 p_{R_n}(u_a,v_b)=1.
 \label{eq:V77-Cartan-unit-visibility}
\end{equation}
Coupling this retained intermediate with \(C=Q_n\) produces exactly
the old dual outputs
\[
 R_n\otimes Q_n=\bigoplus_{k=0}^{n}T_k.
\]
If \(a=0\) or \(b=0\), the cumulant is identically zero; if
\(a,b>0\), product visibility supplies no debit.
\end{proposition}

\begin{proof}
The tensor product of highest-weight vectors has weight
\((a+b,0)\), is annihilated by all raising operators, and generates
the unique Cartan summand \(R_{a+b}\).  This proves
\eqref{eq:V77-Cartan-unit-visibility}.  The second decomposition is
\eqref{eq:V76-dual-decomposition}.  The endpoint statement follows
from \Cref{prop:V77-trivial-input-deletion}.
\end{proof}

Put
\[
 \zeta(g):=\langle e,ge\rangle.
\]
Since \(R_m=\operatorname{Sym}^m(\mathbb C^3)\),
\begin{equation}
 m_{R_m,e^{\odot m}}(g)=\zeta(g)^m.
 \label{eq:V77-coherent-matrix-coefficient}
\end{equation}

\begin{corollary}[Scalar integral for every Cartan lift]
\label{cor:V77-Cartan-scalar-integral}
For the lift in \Cref{prop:V77-Cartan-lift},
\(\Theta_{\Omega;u_a,v_b}(J,L)\) equals
\begin{align}
 &\int_G\kappa_\Omega(w)\,\dd w
 \int_G\mu_\alpha(\dd g)
 \int_G\mu_\alpha(\dd h)
 \notag\\
 &\quad\cdot
 \zeta\!\left(
 w g^{\mathbf1_{\{A\in J\}}}
 h^{\mathbf1_{\{A\in L\}}}\right)^a
 \notag\\
 &\quad\cdot
 \zeta\!\left(
 w g^{\mathbf1_{\{B\in J\}}}
 h^{\mathbf1_{\{B\in L\}}}\right)^b
 \notag\\
 &\quad\cdot
 \varphi_{Q_n}\!\left(
 w g^{\mathbf1_{\{C\in J\}}}
 h^{\mathbf1_{\{C\in L\}}}\right).
 \label{eq:V77-Cartan-scalar-integral}
\end{align}
Substitution into \eqref{eq:V77-signed-product-energy} is an exact
scalar integral for the complete cumulant square on the most
product-visible lift of
\eqref{eq:V77-actual-dual-layer}.
\end{corollary}

\begin{proof}
Insert \eqref{eq:V77-coherent-matrix-coefficient} twice in
\eqref{eq:V77-product-common-integral}.  No multiplicity
approximation is made.
\end{proof}

\begin{assessment}[The reduced hard family]
\label{ass:V77-reduced-hard-family}
The lifting ambiguity has now been narrowed to an explicit boundary
family.  Low-visibility lifts may use
\eqref{eq:V76-visible-composite-bound}.  Trivial Cartan endpoints
vanish exactly.  The remaining product-visible test is
\[
 a,b\ge1,\qquad a+b=n,\qquad C=Q_n,
\]
with the complete signed square represented by
\eqref{eq:V77-Cartan-scalar-integral}.  A uniform estimate must keep
the dependence on the split \(a+b=n\); replacing both coherent
powers by their modulus-one cap would erase connectedness again.
\end{assessment}

\section{V77 ledger and next acquisition}
\label{sec:V77-ledger}

\begin{theorem}[Integrated V77 statement]
\label{thm:V77-integrated}
Preserving V1--V76, V77 proves:
\begin{enumerate}[label=\textup{(V77.\arabic*)},leftmargin=5.7em]
\item the precise distinction between the V76 \(AB\)-channel law and
      the V66--V71 \(S\otimes C\) regression;
\item the lifting-fibre formulation
      \eqref{eq:V77-lifting-fibre};
\item the exact product-vector common integral
      \eqref{eq:V77-product-common-integral};
\item the complete signed product-shell energy
      \eqref{eq:V77-signed-product-energy};
\item exact deletion of every trivial input
      \eqref{eq:V77-trivial-input-deletion};
\item a nontrivial Cartan lift with unit product visibility
      \eqref{eq:V77-Cartan-unit-visibility}; and
\item the scalar coherent formula
      \eqref{eq:V77-Cartan-scalar-integral} for the remaining hard
      dual family.
\end{enumerate}
\end{theorem}

\begin{proof}
These are
\Cref{prop:V77-dual-data,thm:V77-product-common-integral,
thm:V77-signed-product-energy,
prop:V77-trivial-input-deletion,prop:V77-Cartan-lift,
cor:V77-Cartan-scalar-integral}.
\end{proof}

\begin{openproblem}[V78 mandate: connected coherent estimate]
\label{op:V78-mandate}
Work on the exact Cartan lift
\[
 R_a\otimes R_b\to R_n,\qquad
 R_n\otimes Q_n\to T_k,\qquad a+b=n.
\]
\begin{enumerate}
\item Split off \(a=0\) or \(b=0\) by exact deletion, not by a
      limiting estimate.
\item For \(a,b\ge1\), keep the \(25\) terms combined as the square in
      \eqref{eq:V77-signed-product-energy}.
\item Use the coherent powers in
      \eqref{eq:V77-Cartan-scalar-integral} to extract a factor
      depending on \(\min\{a,b\}\), uniformly in the split
      \(a+b=n\).
\item Resolve the final shell \(k\asymp\sqrt n\) only after that
      connected factor is present.
\item If the central shell kernel prevents a sign-safe estimate,
      return to the positive operator
      \(\mathcal R_\Omega^{1/2}K\rho_{u,v;C}^{1/2}\) and estimate its
      Hilbert--Schmidt norm directly.
\item Keep pair banks outside until the core estimate has a declared
      dimension scale; then test the same \(ab\) orientation in
      \eqref{eq:V72-pair-core-product}.
\end{enumerate}
The next compulsory companion audit remains V80.
\end{openproblem}

\section{Firewall after V77}
\label{sec:V77-firewall}

\begin{firewall}[Current claim boundary]
\label{fire:V77-final}
V77 corrects the ancestry layer and gives an exact signed scalar
integral for the genuinely product-visible dual Cartan lift.  It does
not yet estimate that integral at the required scale, close the
product-state collision gate, prove JREC, close simultaneous
sideways endpoints, aggregate punctures or multiple collisions,
construct the continuum Yang--Mills theory, or prove a positive mass
gap.
\end{firewall}

\section{Append-only record for V77}
\label{sec:V77-record}

V77 was appended to the complete V76 manuscript.  The exact V76
parent had \(54631\) lines, \(1953118\) bytes, and SHA--256 digest
\[
 \texttt{2b029774608a1e8de587f0440a22505b42cafa630815b982f1b48327e3a26471}.
\]
No V76 mathematical content was changed.  The sole final
\verb|\end{document}| is restored below.

% =============================================================================
% END APPENDED VERSION V77 MATERIAL
% =============================================================================
% =============================================================================
% BEGIN APPENDED VERSION V78 MATERIAL
% =============================================================================

\chapter{V78: Fully centred cumulants and a positive
Hilbert--Schmidt product test}
\label{ch:V78}

\section{Purpose}
\label{sec:V78-purpose}

V77 reduced the genuinely product-visible dual obstruction to the
Cartan lifts
\[
 R_a\otimes R_b\to R_{a+b}=R_n,\qquad C=Q_n.
\]
The scalar kernel \(\kappa_\Omega\) in
\eqref{eq:V77-Cartan-scalar-integral} need not be positive, so taking
its absolute value would lose the signed cumulant structure.  V78
implements the alternative prescribed there: the complete ternary
cumulant is first written as the expectation of three centred
representation matrices, and the shell energy is then the square of
one Hilbert--Schmidt norm.

\begin{firewall}[Scope of V78]
\label{fire:V78-scope}
The unweighted estimate below is uniform and exact at the level of
two Wilson variances.  On the mesoscopic shell, the elementary
spectral maximum still costs \((1+s_\alpha n)^2\).  Thus V78 closes
the near-trivial Cartan endpoints but does not yet estimate the
balanced Cartan shell at the collision scale.
\end{firewall}

\section{The fully centred operator identity}
\label{sec:V78-centred-identity}

For an irreducible \(R\), define
\begin{equation}
 Z_R(g):=D^R(g)-q_{\alpha,R}I_{H_R}.
 \label{eq:V78-centred-representation}
\end{equation}

\begin{theorem}[The five-term cumulant is one centred moment]
\label{thm:V78-centred-cumulant}
On \(H_A\otimes H_B\otimes H_C\),
\begin{equation}
 \boxed{
 K=
 \int_G
 Z_A(g)\otimes Z_B(g)\otimes Z_C(g)\,
 \mu_\alpha(\dd g).}
 \label{eq:V78-centred-cumulant}
\end{equation}
This identity holds before any output or ancestry decomposition.
\end{theorem}

\begin{proof}
Expand the tensor product on the right.  The term containing all
three uncentred matrices is \(Q_{ABC}\).  The three terms containing
one scalar mean are
\[
 -xQ_{BC}-yQ_{AC}-zQ_{AB}.
\]
The three terms containing two scalar means each integrate the
remaining representation to its scalar mean and hence each equal
\(xyzI\).  The fully scalar term is \(-xyzI\).  Their net scalar
coefficient is \(3xyz-xyz=2xyz\).  This is exactly
\eqref{eq:V75-K-central}.
\end{proof}

\begin{corollary}[Exact connectedness in each input]
\label{cor:V78-connectedness}
The operator \(K\) vanishes if any \(Z_A,Z_B,Z_C\) vanishes
identically.  In particular
\Cref{prop:V77-trivial-input-deletion} holds on every output without
separate block computation.
\end{corollary}

\begin{proof}
Immediate from \eqref{eq:V78-centred-cumulant}.
\end{proof}

\section{The product-column Hilbert--Schmidt norm}
\label{sec:V78-HS-column}

For unit \(u\in H_A\), \(v\in H_B\), let
\begin{equation}
 W_{u,v}:H_C\longrightarrow
 H_A\otimes H_B\otimes H_C,
 \qquad
 W_{u,v}\xi=u\otimes v\otimes\xi.
 \label{eq:V78-column-isometry}
\end{equation}
It is an isometry.  Define the normalized unweighted product energy
\begin{equation}
 \mathfrak e_{u,v}:=
 \frac1{d_C}\|KW_{u,v}\|_{\rm HS}^2.
 \label{eq:V78-unweighted-energy}
\end{equation}

\begin{theorem}[Exact positive double-integral formula]
\label{thm:V78-positive-double-integral}
One has
\begin{align}
 \mathfrak e_{u,v}
 &=
 \frac1{d_C}
 \left\langle u\otimes v,
 \operatorname{Tr}_C(K^2)u\otimes v
 \right\rangle
 \label{eq:V78-energy-partial-trace}\\
 &=
 \int_G\!\!\int_G
 \langle Z_A(g)u,Z_A(h)u\rangle
 \langle Z_B(g)v,Z_B(h)v\rangle
 \notag\\
 &\hspace{7em}\cdot
 \frac1{d_C}
 \operatorname{Tr}\!\left[
 Z_C(g)^*Z_C(h)\right]\,
 \mu_\alpha(\dd g)\mu_\alpha(\dd h).
 \label{eq:V78-positive-double-integral}
\end{align}
Although the integrand may be complex, the double integral is the
nonnegative squared norm in
\eqref{eq:V78-unweighted-energy}.
\end{theorem}

\begin{proof}
Since the Wilson law is inversion-symmetric, \(K=K^*\).  Therefore
\[
 \|KW_{u,v}\|_{\rm HS}^2
 =
 \operatorname{Tr}_{H_C}(W_{u,v}^*K^2W_{u,v}),
\]
which is
\eqref{eq:V78-energy-partial-trace}.  Insert
\eqref{eq:V78-centred-cumulant} twice in the Hilbert--Schmidt inner
product.  Tensor-factorization gives the three kernels in
\eqref{eq:V78-positive-double-integral}.  Fubini is valid in finite
dimension.
\end{proof}

\section{An exact variance calculation}
\label{sec:V78-variance}

Put
\begin{equation}
 \upsilon_R:=1-q_{\alpha,R}^2.
 \label{eq:V78-Wilson-variance}
\end{equation}
The multiplier is real because the law is inversion-symmetric.

\begin{lemma}[Vector and Hilbert--Schmidt variances]
\label{lem:V78-exact-variances}
For every irreducible \(R\) and every unit \(\xi\in H_R\),
\begin{align}
 \int_G\|Z_R(g)\xi\|^2\,\mu_\alpha(\dd g)
 &=\upsilon_R,
 \label{eq:V78-vector-variance}\\
 \frac1{d_R}\int_G
 \|Z_R(g)\|_{\rm HS}^2\,\mu_\alpha(\dd g)
 &=\upsilon_R.
 \label{eq:V78-HS-variance}
\end{align}
\end{lemma}

\begin{proof}
Unitarity and
\(\int D^R(g)\,\mu_\alpha(\dd g)=q_{\alpha,R}I\) give
\begin{align*}
 \int\|(D^R(g)-q_{\alpha,R}I)\xi\|^2\,\dd\mu_\alpha(g)
 &=1-2q_{\alpha,R}^2+q_{\alpha,R}^2\\
 &=1-q_{\alpha,R}^2.
\end{align*}
This proves \eqref{eq:V78-vector-variance}.  Sum the same identity
over an orthonormal basis and divide by \(d_R\) to obtain
\eqref{eq:V78-HS-variance}.
\end{proof}

\begin{theorem}[Uniform two-variance product bound]
\label{thm:V78-two-variance-bound}
For all irreducibles and all unit \(u,v\),
\begin{equation}
 \boxed{
 \mathfrak e_{u,v}
 \le
 4\min\left\{
 \upsilon_A\upsilon_B,\,
 \upsilon_A\upsilon_C,\,
 \upsilon_B\upsilon_C
 \right\}.}
 \label{eq:V78-two-variance-bound}
\end{equation}
Consequently
\begin{equation}
 \frac1{d_C}
 \left\|
 \operatorname{Tr}_C K^2
 \right\|_{A\otimes B,\mathrm{prod}}
 \le
 4\min_{\{R,Q\}\subset\{A,B,C\}}
 \upsilon_R\upsilon_Q.
 \label{eq:V78-product-two-variance}
\end{equation}
\end{theorem}

\begin{proof}
Minkowski's integral inequality and
\eqref{eq:V78-centred-cumulant} give
\begin{align*}
 \frac1{\sqrt{d_C}}\|KW_{u,v}\|_{\rm HS}
 &\le
 \int_G
 \|Z_A(g)u\|\,
 \|Z_B(g)v\|\,
 \frac{\|Z_C(g)\|_{\rm HS}}{\sqrt{d_C}}\,
 \mu_\alpha(\dd g).
\end{align*}
Every centred unitary difference has operator norm at most two.
Discard the \(C\)-factor by this bound and apply Cauchy--Schwarz to
the other two.  By \eqref{eq:V78-vector-variance},
\[
 \frac1{\sqrt{d_C}}\|KW_{u,v}\|_{\rm HS}
 \le2\sqrt{\upsilon_A\upsilon_B}.
\]
Instead discard the \(B\)-factor and apply Cauchy--Schwarz using
\eqref{eq:V78-vector-variance} for \(A\) and
\eqref{eq:V78-HS-variance} for \(C\).  This gives
\[
 \frac1{\sqrt{d_C}}\|KW_{u,v}\|_{\rm HS}
 \le2\sqrt{\upsilon_A\upsilon_C}.
\]
The third pairing is identical.  Square all three valid bounds and
take their minimum, proving
\eqref{eq:V78-two-variance-bound}.  Supremizing over product vectors
gives \eqref{eq:V78-product-two-variance}.
\end{proof}

\begin{remark}[Comparison with the final-singlet identity]
\label{rem:V78-singlet-comparison}
V66 obtained a sharper difference of two squares on the
multiplicity-one final singlet.  The estimate
\eqref{eq:V78-two-variance-bound} is weaker there, but it holds on
the complete tensor carrier, with arbitrary multiplicity and before
any final-output projection.
\end{remark}

\section{Positive insertion of an output shell}
\label{sec:V78-shell}

Define
\begin{equation}
 \mathfrak e_{\Omega;u,v}:=
 \frac1{d_C}
 \left\|
 \mathcal R_\Omega^{1/2}KW_{u,v}
 \right\|_{\rm HS}^2.
 \label{eq:V78-shell-HS-energy}
\end{equation}

\begin{theorem}[Sign-safe shell form]
\label{thm:V78-sign-safe-shell}
For every shell \(\Omega\),
\begin{align}
 \mathfrak e_{\Omega;u,v}
 &=
 \frac1{d_C}
 \left\langle u\otimes v,
 \operatorname{Tr}_C[
 \mathcal R_\Omega K^2]\,
 u\otimes v\right\rangle,
 \label{eq:V78-shell-equality}\\
 0\le\mathfrak e_{\Omega;u,v}
 &\le
 r_\Omega^*
 \mathfrak e_{u,v},
 \qquad
 r_\Omega^*:=
 \max_{T\in\Omega}
 [1+s_\alpha\Xi(T)]^2.
 \label{eq:V78-shell-spectral-bound}
\end{align}
Thus
\begin{equation}
 \boxed{
 \mathfrak e_{\Omega;u,v}
 \le
 4r_\Omega^*
 \min_{\{R,Q\}\subset\{A,B,C\}}
 \upsilon_R\upsilon_Q.}
 \label{eq:V78-shell-two-variance}
\end{equation}
\end{theorem}

\begin{proof}
The central output multiplier \(\mathcal R_\Omega\) commutes with
\(K\).  Hence the first squared norm is
\[
 d_C^{-1}\operatorname{Tr}
 [W_{u,v}^*K\mathcal R_\Omega KW_{u,v}]
 =
 d_C^{-1}\operatorname{Tr}
 [W_{u,v}^*\mathcal R_\Omega K^2W_{u,v}],
\]
proving \eqref{eq:V78-shell-equality}.  Positivity is now explicit.
The spectral bound
\(0\le\mathcal R_\Omega\le r_\Omega^*I\) proves
\eqref{eq:V78-shell-spectral-bound}.  Combine with
\eqref{eq:V78-two-variance-bound}.
\end{proof}

\section{Application to the Cartan split}
\label{sec:V78-Cartan-split}

For \(R_m=(m,0)\), the Casimir formula and the one-link upper Wilson
gap give
\begin{equation}
 \upsilon_{R_m}
 \le
 C\min\{s_\alpha(1+m)^2,1\}.
 \label{eq:V78-symmetric-variance}
\end{equation}

\begin{corollary}[Rigorous endpoint factor for the Cartan lift]
\label{cor:V78-Cartan-endpoint}
Let \(a+b=n\), \(A=R_a\), \(B=R_b\), \(C=Q_n\), and use the coherent
product vector of \Cref{prop:V77-Cartan-lift}.  Then
\begin{equation}
 \boxed{
 \mathfrak e_{u_a,v_b}
 \le
 C
 \min\{s_\alpha(1+a)^2,1\}
 \min\{s_\alpha(1+b)^2,1\}.}
 \label{eq:V78-Cartan-unweighted}
\end{equation}
If \(\Omega_{n;A_0,B_0}\) is the one-coordinate dual shell
\[
 A_0\sqrt n\le k\le B_0\sqrt n,
\qquad T=T_k=(k,k),
\]
then
\begin{align}
 \mathfrak e_{\Omega_{n;A_0,B_0};u_a,v_b}
 &\le
 C_{A_0,B_0}(1+s_\alpha n)^2
 \notag\\
 &\quad\cdot
 \min\{s_\alpha(1+a)^2,1\}
 \min\{s_\alpha(1+b)^2,1\}.
 \label{eq:V78-Cartan-shell}
\end{align}
For \(a=0\) or \(b=0\), both energies are exactly zero.
\end{corollary}

\begin{proof}
Use the \(AB\) entry in
\eqref{eq:V78-two-variance-bound} and
\eqref{eq:V78-symmetric-variance}.  On the stated shell,
\(\Xi(T_k)\asymp(k+1)^2\le C_{A_0,B_0}n\); hence
\(r_\Omega^*\le C_{A_0,B_0}(1+s_\alpha n)^2\).
Apply \eqref{eq:V78-shell-two-variance}.  Exact endpoint vanishing is
\Cref{prop:V77-trivial-input-deletion}.
\end{proof}

\begin{assessment}[Edge splits versus the balanced core]
\label{ass:V78-edge-balanced}
The split parameter \(a+b=n\) can no longer disappear from the
estimate.  Formula \eqref{eq:V78-Cartan-unweighted} gives an exact
connectedness gain near either trivial endpoint and zero at the
endpoint itself.  When both
\(s_\alpha a^2\) and \(s_\alpha b^2\) are large, however, both
variances saturate and the bound becomes constant.  In that balanced
region the output projection, not unweighted connectedness, must
supply the missing dimension debit.

The remaining positive object is therefore
\[
 \left\|
 P_{\Omega_{n;A_0,B_0}}
 [1+s_\alpha\Xi(\mathbf T)]
 K W_{u_a,v_b}
 \right\|_{\rm HS}^2,
\]
uniformly for \(a,b\ge1\), \(a+b=n\).  This is narrower than the V75
trace problem and preserves the original product column.
\end{assessment}

\section{V78 ledger and next acquisition}
\label{sec:V78-ledger}

\begin{theorem}[Integrated V78 statement]
\label{thm:V78-integrated}
Preserving V1--V77, V78 proves:
\begin{enumerate}[label=\textup{(V78.\arabic*)},leftmargin=5.7em]
\item the fully centred identity
      \eqref{eq:V78-centred-cumulant};
\item the exact positive double integral
      \eqref{eq:V78-positive-double-integral};
\item the vector and Hilbert--Schmidt variance identities
      \eqref{eq:V78-vector-variance}--%
      \eqref{eq:V78-HS-variance};
\item the uniform two-variance product bound
      \eqref{eq:V78-two-variance-bound};
\item the sign-safe shell norm
      \eqref{eq:V78-shell-HS-energy}; and
\item the rigorous Cartan endpoint estimates
      \eqref{eq:V78-Cartan-unweighted}--%
      \eqref{eq:V78-Cartan-shell}.
\end{enumerate}
\end{theorem}

\begin{proof}
These are
\Cref{thm:V78-centred-cumulant,
thm:V78-positive-double-integral,
lem:V78-exact-variances,thm:V78-two-variance-bound,
thm:V78-sign-safe-shell,cor:V78-Cartan-endpoint}.
\end{proof}

\begin{openproblem}[V79 mandate: projected balanced Cartan column]
\label{op:V79-mandate}
On \(a+b=n\), separate the endpoint regimes already controlled by
\eqref{eq:V78-Cartan-shell}.  In the balanced remainder:
\begin{enumerate}
\item decompose the vector-valued column
      \(KW_{u_a,v_b}\) into the final \(T_k=(k,k)\) outputs;
\item use the common coherent line of \(u_a,v_b\) before applying any
      norm, so the \(a\)- and \(b\)-powers remain correlated;
\item prove an exact Plancherel or spherical-polynomial formula for
      the \(T_k\)-component;
\item compare that component with the V76 density
      \(2t^3e^{-t^2}\), but do not identify the two probability laws;
\item if the projected column has order-one mass on
      \(k\asymp\sqrt n\), insert the external pair factor in the same
      \(ab\) orientation before declaring failure;
\item perform the mandatory companion audit at V80 after this
      calculation.
\end{enumerate}
\end{openproblem}

\section{Firewall after V78}
\label{sec:V78-firewall}

\begin{firewall}[Current claim boundary]
\label{fire:V78-final}
V78 proves a sign-safe product-column formulation and controls
near-trivial Cartan lifts by two exact Wilson variances.  It does not
yet control the projected balanced Cartan column, close the
product-state collision gate, prove JREC, close simultaneous
sideways endpoints, aggregate punctures or multiple collisions,
construct the continuum Yang--Mills theory, or prove a positive mass
gap.
\end{firewall}

\section{Append-only record for V78}
\label{sec:V78-record}

V78 was appended to the complete V77 manuscript.  The exact V77
parent had \(55135\) lines, \(1969081\) bytes, and SHA--256 digest
\[
 \texttt{1f9c638c0ef98138963c04f6679906331bb683a655aae9a8711e27c5f6f52b92}.
\]
No V77 mathematical content was changed.  The sole final
\verb|\end{document}| is restored below.

% =============================================================================
% END APPENDED VERSION V78 MATERIAL
% =============================================================================
% =============================================================================
% BEGIN APPENDED VERSION V79 MATERIAL
% =============================================================================

\chapter{V79: Exact final-output Plancherel reduction for the
balanced Cartan column}
\label{ch:V79}

\section{Purpose}
\label{sec:V79-purpose}

V78 replaced the sign-indefinite shell kernel estimate by the
positive norm
\(\|\mathcal R_\Omega^{1/2}KW_{u,v}\|_{\rm HS}^2\).
V79 decomposes that norm by final irreducible output while retaining
one fixed \(AB\) ancestry.  Schur's lemma removes every carrier
coordinate and leaves one scalar: the squared norm of the complete
signed multiplicity vector \(K_Te_{S,\alpha,\beta}\).

For the Cartan lift
\[
 R_a\otimes R_b\to R_n,\qquad R_n\otimes Q_n\to T_k,
\]
both arrows are multiplicity one.  The complete product-column shell
energy is therefore an exact one-dimensional weighted sum.  This is
the promised Plancherel reduction; it neither replaces \(K_T\) by
its operator norm nor separates its covariance cuts.

\begin{firewall}[Scope of V79]
\label{fire:V79-scope}
The theorem below identifies the exact scalar sequence that must
decay.  It does not prove that decay.  In particular, the dimension
weights have order-one total mass on \(k\asymp\sqrt n\), so output
dimension counting alone cannot close the balanced Cartan case.
\end{firewall}

\section{A fixed-ancestry partial-trace Plancherel formula}
\label{sec:V79-fixed-ancestry}

Let \(X\ge0\) be invariant on
\(H_A\otimes H_B\otimes H_C\), with final-output blocks
\begin{equation}
 X=\bigoplus_T I_{H_T}\otimes X_T,
 \qquad X_T\ge0
 \quad\hbox{on }\mathcal M_T.
 \label{eq:V79-invariant-positive-blocks}
\end{equation}
Use the \(AB\)-first path basis
\[
 e_{S,\alpha,\beta}^{T},
 \qquad
 1\le\alpha\le N_{AB}^{S},\quad
 1\le\beta\le N_{SC}^{T}.
\]

\begin{theorem}[Fixed-ancestry Plancherel identity]
\label{thm:V79-fixed-ancestry-Plancherel}
Suppose a unit product vector lies entirely in one \(AB\) copy:
\[
 u\otimes v=\iota_{S,\alpha}w,
 \qquad \|w\|=1.
\]
Then
\begin{equation}
 \boxed{
 \left\langle u\otimes v,
 \operatorname{Tr}_C X\,u\otimes v\right\rangle
 =
 \frac1{d_S}
 \sum_Td_T
 \sum_{\beta=1}^{N_{SC}^{T}}
 \left\langle
 e_{S,\alpha,\beta}^{T},
 X_Te_{S,\alpha,\beta}^{T}
 \right\rangle.}
 \label{eq:V79-fixed-ancestry-Plancherel}
\end{equation}
The right side is independent of the carrier vector \(w\).
\end{theorem}

\begin{proof}
The matrix partial-trace theorem underlying
\eqref{eq:V67-absolute-cut-trace} applies to any positive
multiplicity block \(X_T\), not only to \(|K_T|\).  It gives
\[
 \operatorname{Tr}_C X
 =
 \bigoplus_S I_{H_S}\otimes
 \left[
 \frac1{d_S}\sum_Td_T
 \sum_{\beta=1}^{N_{SC}^{T}}X_T[S,\beta]
 \right],
\]
where \(X_T[S,\beta]\) is the principal matrix on the first-stage
multiplicity indices \(\alpha\).  Evaluate this operator on
\(\iota_{S,\alpha}w\).  The identity on \(H_S\) removes \(w\), and
the \(\alpha\alpha\) entry of each principal matrix is the displayed
quadratic form.  This proves
\eqref{eq:V79-fixed-ancestry-Plancherel}.
\end{proof}

\begin{corollary}[Signed cumulant output formula]
\label{cor:V79-signed-output-formula}
Under the hypotheses of
\Cref{thm:V79-fixed-ancestry-Plancherel}, for every output set
\(\Omega\),
\begin{align}
 &\left\langle u\otimes v,
 \operatorname{Tr}_C[
 \mathcal R_\Omega K^2]\,
 u\otimes v\right\rangle
 \notag\\
 &\quad=
 \frac1{d_S}
 \sum_{T\in\Omega}d_T
 [1+s_\alpha\Xi(T)]^2
 \sum_{\beta=1}^{N_{SC}^{T}}
 \left\|
 K_Te_{S,\alpha,\beta}^{T}
 \right\|^2.
 \label{eq:V79-signed-output-formula}
\end{align}
\end{corollary}

\begin{proof}
In \eqref{eq:V79-invariant-positive-blocks}, take
\[
 X_T=
 \mathbf1_{\{T\in\Omega\}}
 [1+s_\alpha\Xi(T)]^2K_T^2.
\]
Since \(K_T=K_T^*\), its quadratic form on a path vector is
\(\|K_Te\|^2\).
\end{proof}

\section{The exact Cartan dual sequence}
\label{sec:V79-Cartan-sequence}

Fix \(a,b\ge0\), \(a+b=n\), and use the Cartan coherent product
vector of \Cref{prop:V77-Cartan-lift}.  Let
\[
 e_{a,b;n,k}\in\mathcal M_{T_k}
\]
be the unit path vector corresponding to
\[
 R_a\otimes R_b\longrightarrow R_n,
 \qquad
 R_n\otimes Q_n\longrightarrow T_k.
\]
Both path multiplicities are one.  Define the complete signed
ancestral vector
\begin{equation}
 \boldsymbol\kappa_{a,b;n,k}:=
 K_{T_k}e_{a,b;n,k}.
 \label{eq:V79-ancestral-vector}
\end{equation}

\begin{theorem}[Exact balanced-Cartan shell energy]
\label{thm:V79-Cartan-shell-energy}
For every \(I\subseteq\{0,\dots,n\}\),
\begin{align}
 &\left\langle u_a\otimes v_b,
 \operatorname{Tr}_{Q_n}\left[
 P_{\{T_k:k\in I\}}
 [1+s_\alpha\Xi(\mathbf T)]^2K^2
 \right]
 u_a\otimes v_b\right\rangle
 \notag\\
 &\qquad=
 \boxed{
 \sum_{k\in I}
 \frac{(k+1)^3}{d_n}
 [1+s_\alpha\Xi(T_k)]^2
 \|\boldsymbol\kappa_{a,b;n,k}\|^2.}
 \label{eq:V79-Cartan-shell-energy}
\end{align}
There is no unrecorded multiplicity or representation-formation
factor.
\end{theorem}

\begin{proof}
The product vector lies entirely in the unique Cartan copy \(S=R_n\)
by \Cref{prop:V77-Cartan-lift}.  In
\eqref{eq:V79-signed-output-formula}, use
\[
 d_S=d_n,\qquad
 N_{R_nQ_n}^{T_k}=1,\qquad
 d_{T_k}=(k+1)^3.
\]
The only remaining multiplicity quadratic form is
\(\|K_{T_k}e_{a,b;n,k}\|^2\), which is
\eqref{eq:V79-ancestral-vector}.
\end{proof}

\begin{corollary}[Order-one dimension mass on the thermal window]
\label{cor:V79-thermal-dimension-mass}
For fixed \(0\le A_0<B_0<\infty\),
\begin{equation}
 \boxed{
 \lim_{n\to\infty}
 \sum_{A_0\sqrt n\le k\le B_0\sqrt n}
 \frac{(k+1)^3}{d_n}
 =
 \int_{A_0}^{B_0}2t^3\,\dd t
 =
 \frac{B_0^4-A_0^4}{2}.}
 \label{eq:V79-thermal-dimension-mass}
\end{equation}
Thus the weights in
\eqref{eq:V79-Cartan-shell-energy}, unlike the probabilities
\eqref{eq:V76-coherent-law}, have no exponential factor
\(e^{-t^2}\).  The two laws must not be identified.
\end{corollary}

\begin{proof}
Since \(d_n\sim n^2/2\), uniformly on the stated window
\[
 \frac{(k+1)^3}{d_n}
 =
 \frac1{\sqrt n}
 \left[
 2\left(\frac{k}{\sqrt n}\right)^3+o(1)
 \right].
\]
The sum is a Riemann sum, proving
\eqref{eq:V79-thermal-dimension-mass}.  The V76 probability contains
the additional factorial quotient in
\eqref{eq:V76-coherent-law}; that quotient converges to
\(e^{-t^2}\) after its own normalization.
\end{proof}

\section{Six scalar moments, not twenty-five absolute terms}
\label{sec:V79-six-moments}

On the \(AB\)-first path vector \(e=e_{a,b;n,k}\), the \(AB\) pair
moment is scalar:
\[
 Q_{AB}^{T_k}e=q_{\alpha,R_n}e.
\]
Write
\[
 x=q_{\alpha,R_a},\qquad
 y=q_{\alpha,R_b},\qquad
 z=q_{\alpha,Q_n},\qquad
 s_n=q_{\alpha,R_n},\qquad
 t_k=q_{\alpha,T_k}
\]
and put
\begin{equation}
 \vartheta_{a,b;n,k}:=
 t_k-zs_n+2xyz.
 \label{eq:V79-scalar-part}
\end{equation}

\begin{proposition}[Exact ancestral-vector formula]
\label{prop:V79-ancestral-vector-formula}
The signed vector in \eqref{eq:V79-ancestral-vector} is
\begin{equation}
 \boxed{
 \boldsymbol\kappa_{a,b;n,k}
 =
 \vartheta_{a,b;n,k}e
 -xQ_{BC}^{T_k}e
 -yQ_{AC}^{T_k}e.}
 \label{eq:V79-ancestral-vector-formula}
\end{equation}
Consequently
\begin{align}
 \|\boldsymbol\kappa_{a,b;n,k}\|^2
 &=
 \vartheta_{a,b;n,k}^2
 +x^2\langle e,(Q_{BC}^{T_k})^2e\rangle
 +y^2\langle e,(Q_{AC}^{T_k})^2e\rangle
 \notag\\
 &\quad
 -2x\vartheta_{a,b;n,k}
 \langle e,Q_{BC}^{T_k}e\rangle
 -2y\vartheta_{a,b;n,k}
 \langle e,Q_{AC}^{T_k}e\rangle
 \notag\\
 &\quad
 +2xy\operatorname{Re}
 \langle Q_{BC}^{T_k}e,Q_{AC}^{T_k}e\rangle.
 \label{eq:V79-six-moment-square}
\end{align}
All six quantities occur in one nonnegative square; none has been
estimated separately.
\end{proposition}

\begin{proof}
Restrict the common-basis cumulant
\eqref{eq:V47-common-cumulant} to the final \(T_k\) block.  The total
moment gives \(t_kI\), and the \(AB\) moment gives \(s_nI\) on the
chosen path.  Collect those scalar terms with \(2xyzI\), yielding
\(\vartheta_{a,b;n,k}e\).  The two remaining pair moments give
\eqref{eq:V79-ancestral-vector-formula}.  Take its squared norm and
use self-adjointness of the pair moment operators to obtain
\eqref{eq:V79-six-moment-square}.
\end{proof}

\begin{remark}[Endpoint cancellation is invisible termwise]
\label{rem:V79-endpoint-cancellation}
When \(a=0\), one has \(x=1\), but the individual terms in
\eqref{eq:V79-six-moment-square} need not vanish.  Their assembled
vector \eqref{eq:V79-ancestral-vector-formula} is zero by
\eqref{eq:V77-trivial-input-deletion}.  This is a concrete reason
not to place absolute values on the six scalar moments.
\end{remark}

\section{A sharp necessary thermal decision}
\label{sec:V79-decision}

\begin{proposition}[Balanced-shell obstruction test]
\label{prop:V79-balanced-obstruction}
Fix \(0<A_0<B_0<\infty\).  If there is an
\(\varepsilon>0\) and a sequence
\[
 a_n+b_n=n,\qquad
 \min\{a_n,b_n\}\longrightarrow\infty
\]
such that
\begin{equation}
 \inf_{A_0\sqrt n\le k\le B_0\sqrt n}
 \|\boldsymbol\kappa_{a_n,b_n;n,k}\|
 \ge\varepsilon,
 \label{eq:V79-nondecay-hypothesis}
\end{equation}
then the unweighted thermal part of the Cartan product-column energy
has the positive lower limit
\begin{align}
 \liminf_{n\to\infty}
 &\left\langle u_{a_n}\otimes v_{b_n},
 \operatorname{Tr}_{Q_n}\left[
 P_{\Omega_n}K^2\right]
 u_{a_n}\otimes v_{b_n}\right\rangle
 \notag\\
 &\qquad\ge
 \frac{\varepsilon^2}{2}(B_0^4-A_0^4)>0,
 \label{eq:V79-obstruction-conclusion}
\end{align}
where
\(\Omega_n=\{T_k:A_0\sqrt n\le k\le B_0\sqrt n\}\).
\end{proposition}

\begin{proof}
Drop the weight, restrict the nonnegative sum
\eqref{eq:V79-Cartan-shell-energy} to \(\Omega_n\), insert
\eqref{eq:V79-nondecay-hypothesis}, and apply
\eqref{eq:V79-thermal-dimension-mass}.
\end{proof}

\begin{assessment}[Exact remaining acquisition]
\label{ass:V79-acquisition}
The balanced Cartan question is no longer an unspecified Racah
estimate.  It is the uniform decay of the single vector
\[
 \vartheta e-xQ_{BC}^{T_k}e-yQ_{AC}^{T_k}e
\]
on \(k\asymp\sqrt n\), with \(a+b=n\).  If that vector is
nondecaying, \Cref{prop:V79-balanced-obstruction} gives an actual
product-vector obstruction to a core-only vanishing estimate, and
the external pair bank must be inserted before any conclusion.  If
it decays, its squared norm can be summed directly with the exact
weights in \eqref{eq:V79-Cartan-shell-energy}.

The calculation needed next is therefore a first- and second-moment
Racah calculation for the two pair moment operators on the one
Cartan ancestry vector, not an entrywise estimate of the complete
Racah matrix.
\end{assessment}

\section{V79 ledger and next acquisition}
\label{sec:V79-ledger}

\begin{theorem}[Integrated V79 statement]
\label{thm:V79-integrated}
Preserving V1--V78, V79 proves:
\begin{enumerate}[label=\textup{(V79.\arabic*)},leftmargin=5.7em]
\item the fixed-ancestry Plancherel identity
      \eqref{eq:V79-fixed-ancestry-Plancherel};
\item the exact signed output decomposition
      \eqref{eq:V79-signed-output-formula};
\item the Cartan dual sequence
      \eqref{eq:V79-Cartan-shell-energy};
\item the order-one thermal dimension law
      \eqref{eq:V79-thermal-dimension-mass};
\item the six-moment signed square
      \eqref{eq:V79-six-moment-square}; and
\item the explicit nondecay obstruction
      \eqref{eq:V79-obstruction-conclusion}.
\end{enumerate}
\end{theorem}

\begin{proof}
These are
\Cref{thm:V79-fixed-ancestry-Plancherel,
cor:V79-signed-output-formula,thm:V79-Cartan-shell-energy,
cor:V79-thermal-dimension-mass,
prop:V79-ancestral-vector-formula,
prop:V79-balanced-obstruction}.
\end{proof}

\begin{openproblem}[V80 mandate after the compulsory companion audit]
\label{op:V80-mandate}
First inspect stable current SM and NS snapshots.  Then:
\begin{enumerate}
\item compute
      \(\langle e,Q_{AC}^{T_k}e\rangle\),
      \(\langle e,(Q_{AC}^{T_k})^2e\rangle\), and their \(BC\)
      analogues on the Cartan path;
\item compute the mixed moment
      \(\operatorname{Re}\langle Q_{BC}^{T_k}e,
      Q_{AC}^{T_k}e\rangle\);
\item use a common orthogonal-polynomial or group-integral
      representation so endpoint deletion remains exact;
\item test uniformly over \(a+b=n\), separating an edge strip from
      balanced \(a/n\);
\item if nondecay occurs, insert the actual \(P_{ab}\) pair factor in
      this same orientation before changing routes.
\end{enumerate}
The following compulsory audit after V80 will be V85.
\end{openproblem}

\section{Firewall after V79}
\label{sec:V79-firewall}

\begin{firewall}[Current claim boundary]
\label{fire:V79-final}
V79 reduces the projected balanced Cartan product column to one exact
weighted sequence of signed ancestral-vector norms.  It does not yet
compute the six Racah moments, close the product-state collision
gate, prove JREC, close simultaneous sideways endpoints, aggregate
punctures or multiple collisions, construct the continuum
Yang--Mills theory, or prove a positive mass gap.
\end{firewall}

\section{Append-only record for V79}
\label{sec:V79-record}

V79 was appended to the complete V78 manuscript.  The exact V78
parent had \(55605\) lines, \(1983031\) bytes, and SHA--256 digest
\[
 \texttt{e3599adb142a72d646ed4d43a09d4f8d2abcc0ff5f227a0e985bca7f397f4357}.
\]
No V78 mathematical content was changed.  The sole final
\verb|\end{document}| is restored below.

% =============================================================================
% END APPENDED VERSION V79 MATERIAL
% =============================================================================
% =============================================================================
% BEGIN APPENDED VERSION V80 MATERIAL
% =============================================================================

\chapter{V80: Companion audit and probability laws for the six
Cartan-path moments}
\label{ch:V80}

\section{Purpose}
\label{sec:V80-purpose}

V79 reduced the balanced Cartan product column to the exact ancestral
vector
\[
 \boldsymbol\kappa_{a,b;n,k}
 =
 \vartheta e-xQ_{BC}^{T_k}e-yQ_{AC}^{T_k}e.
\]
V80 performs the compulsory five-version companion audit.  It then
replaces the six scalar moments in
\eqref{eq:V79-six-moment-square} by two honest spectral probability
laws and one correlation coefficient.  This separates the remaining
problem into:
\begin{enumerate}[label=\textup{(\roman*)}]
\item cancellation of the scalar conditional means;
\item matching of the two conditional variances; and
\item alignment of the two centred Racah vectors.
\end{enumerate}
All three remain attached to the same path vector.

\begin{firewall}[Scope of V80]
\label{fire:V80-scope}
V80 computes exact integral and probability representations of the
remaining moments.  It does not estimate the resulting
\(\mathrm{SU}(3)\) path laws uniformly on the balanced thermal
window.  The mass gap is not inferred from the companion audit.
\end{firewall}

\section{Mandatory V80 companion audit}
\label{sec:V80-audit}

Two consecutive byte reads produced identical SHA--256 digests.  The
stable snapshots are:
\begin{center}
\begin{tabular}{|p{0.43\textwidth}|r|r|p{0.27\textwidth}|}
\hline
Companion file & Lines & Bytes & SHA--256\\
\hline
\texttt{Renewal\_SM\_Einstein\_Continuum\_Research\_Notes\_V83.tex}
& \(65953\) & \(2391873\) &
\texttt{5871582e4f49c665c001fc8fc9c69ad6ebeea5abf8f641ac6cff55eb5e3e01a2}\\
\hline
\texttt{Renewal\_NS\_Stability\_Research\_Notes\_V31.tex}
& \(23052\) & \(887537\) &
\texttt{f1121b62238ad5491bb7cf03635ea9934942edf573ed8faaa9ee79e1d38a6696}\\
\hline
\end{tabular}
\end{center}

The SM companion has advanced from V80 to V83.  V81 proves that
canceling a formal inverse multiplier on the physical range can still
leave a sharp residual infrared loss.  V82 returns upstream to an
unsplit Sobolev--Green phase and closes the opposite cocycles before
solving a dressing equation.  V83 proves a family of supercritical
closed reductions, but also shows that an unbounded explicit dressing
section does not decide the quotient norm: the remaining object is
the infimum over the complete closed residual orbit.

The NS file is byte-identical to the V65, V70, and V75 snapshots.  Its
terminal lesson remains unchanged: a flat probability mark is exact,
but removal of the mark restores the critical first moment unless
recent entry, signed nonlinear work, and the parent tail remain
correlated.

\begin{assessment}[Transferable proof order at V80]
\label{ass:V80-audit-transfer}
For Yang--Mills, a core-only nondecay test is analogous to a bad
chosen section, not to the norm of the complete collision capacity.
The signed ancestral vector must first be evaluated on its true
product path.  If it is nondecaying, the same-orientation pair bank
must be assembled before failure of the collision gate can be
claimed.  Conversely, a termwise Racah bound is analogous to splitting
the SM cocycles or the NS formation factors before their exact
correlation has been used.
\end{assessment}

\begin{firewall}[Companion separation]
\label{fire:V80-companion-separation}
No Sobolev trace theorem, gravitational quotient estimate,
Navier--Stokes stopping-time bound, or companion conclusion is used
as a Yang--Mills hypothesis.  The comparison determines only the
order in which the finite-dimensional Yang--Mills objects below are
kept assembled.
\end{firewall}

\section{Spectral probability laws on one Cartan path}
\label{sec:V80-path-laws}

Fix \(a+b=n\), a final output \(T_k=(k,k)\), and the unit Cartan path
vector
\[
 e=e_{a,b;n,k}\in\mathcal M_{T_k}
\]
from V79.  In the \(AC\)-first bracketing, let
\(\mathsf P_{AC}^{T_k}(U)\) denote the orthogonal projection on the
sum of all path copies with \(AC\)-intermediate \(U\).  Define
\(\mathsf P_{BC}^{T_k}(V)\) analogously.  Put
\begin{align}
 \pi_{AC}^{a,b;n,k}(U)
 &:=
 \|\mathsf P_{AC}^{T_k}(U)e\|^2,
 \label{eq:V80-pi-AC}\\
 \pi_{BC}^{a,b;n,k}(V)
 &:=
 \|\mathsf P_{BC}^{T_k}(V)e\|^2.
 \label{eq:V80-pi-BC}
\end{align}

\begin{lemma}[The two Racah spectral laws are probabilities]
\label{lem:V80-path-probabilities}
One has
\begin{equation}
 \pi_{AC}(U)\ge0,\qquad
 \sum_U\pi_{AC}(U)=1,
 \qquad
 \pi_{BC}(V)\ge0,\qquad
 \sum_V\pi_{BC}(V)=1.
 \label{eq:V80-path-probabilities}
\end{equation}
Every Littlewood--Richardson copy is included inside the
corresponding spectral projection.
\end{lemma}

\begin{proof}
For a fixed final \(T_k\), the complete \(AC\)-first path
decomposition is an orthogonal resolution of the multiplicity
identity.  Apply it to the unit vector \(e\).  The \(BC\) statement is
identical.
\end{proof}

Define the first moments
\begin{align}
 m_{AC}&:=\sum_Uq_{\alpha,U}\pi_{AC}(U),
 &
 m_{BC}&:=\sum_Vq_{\alpha,V}\pi_{BC}(V),
 \label{eq:V80-path-means}
\end{align}
the second moments
\begin{align}
 s_{AC}&:=\sum_Uq_{\alpha,U}^2\pi_{AC}(U),
 &
 s_{BC}&:=\sum_Vq_{\alpha,V}^2\pi_{BC}(V),
 \label{eq:V80-path-seconds}
\end{align}
and the variances
\begin{equation}
 \nu_{AC}:=s_{AC}-m_{AC}^2,
 \qquad
 \nu_{BC}:=s_{BC}-m_{BC}^2.
 \label{eq:V80-path-variances}
\end{equation}

\begin{theorem}[Exact first and second Racah moments]
\label{thm:V80-first-second-moments}
The four unmixed moments in
\eqref{eq:V79-six-moment-square} are exactly
\begin{align}
 \langle e,Q_{AC}^{T_k}e\rangle&=m_{AC},
 &
 \langle e,(Q_{AC}^{T_k})^2e\rangle&=s_{AC},
 \label{eq:V80-AC-moments}\\
 \langle e,Q_{BC}^{T_k}e\rangle&=m_{BC},
 &
 \langle e,(Q_{BC}^{T_k})^2e\rangle&=s_{BC}.
 \label{eq:V80-BC-moments}
\end{align}
In particular \(\nu_{AC},\nu_{BC}\ge0\).
\end{theorem}

\begin{proof}
In the \(AC\)-first basis,
\[
 Q_{AC}^{T_k}
 =
 \sum_Uq_{\alpha,U}\mathsf P_{AC}^{T_k}(U).
\]
The spectral projections are orthogonal, so squaring replaces
\(q_{\alpha,U}\) by \(q_{\alpha,U}^2\).  Evaluation on \(e\) gives
\eqref{eq:V80-AC-moments}.  The \(BC\) equations are identical.
Nonnegativity of the variances is the scalar variance identity for
the probability laws in \eqref{eq:V80-path-probabilities}.
\end{proof}

\section{A common path-character integral}
\label{sec:V80-path-characters}

Let
\[
 L_{a,b;n,k}:H_{T_k}\longrightarrow
 H_{R_a}\otimes H_{R_b}\otimes H_{Q_n}
\]
be the normalized sequential isometry for the Cartan path
\[
 R_a\otimes R_b\to R_n,\qquad
 R_n\otimes Q_n\to T_k.
\]
Define the normalized path functions
\begin{align}
 \chi_{AC}^{a,b;n,k}(g)
 &:=
 \frac1{d_{T_k}}
 \operatorname{Tr}\!\left[
 L_{a,b;n,k}^*
 (D^{R_a}(g)\otimes I_{R_b}\otimes D^{Q_n}(g))
 L_{a,b;n,k}\right],
 \label{eq:V80-path-character-AC}\\
 \chi_{BC}^{a,b;n,k}(g)
 &:=
 \frac1{d_{T_k}}
 \operatorname{Tr}\!\left[
 L_{a,b;n,k}^*
 (I_{R_a}\otimes D^{R_b}(g)\otimes D^{Q_n}(g))
 L_{a,b;n,k}\right].
 \label{eq:V80-path-character-BC}
\end{align}
For two group variables set
\begin{align}
 \chi_{\times}^{a,b;n,k}(g,h)
 &:=
 \frac1{d_{T_k}}
 \operatorname{Tr}\!\left[
 L_{a,b;n,k}^*
 (I_{R_a}\otimes D^{R_b}(g)\otimes D^{Q_n}(g))
 \right.
 \notag\\[-0.2em]
 &\hspace{11em}\left.
 \cdot
 (D^{R_a}(h)\otimes I_{R_b}\otimes D^{Q_n}(h))
 L_{a,b;n,k}\right].
 \label{eq:V80-mixed-path-character}
\end{align}

\begin{theorem}[Exact common-integral formulas for all five moments]
\label{thm:V80-path-integrals}
Let \(\mu_\alpha^{*2}\) be convolution of the Wilson law with itself.
Then
\begin{align}
 m_{AC}
 &=\int_G\chi_{AC}^{a,b;n,k}(g)\,\mu_\alpha(\dd g),
 &
 s_{AC}
 &=\int_G\chi_{AC}^{a,b;n,k}(g)\,
             \mu_\alpha^{*2}(\dd g),
 \label{eq:V80-AC-integrals}\\
 m_{BC}
 &=\int_G\chi_{BC}^{a,b;n,k}(g)\,\mu_\alpha(\dd g),
 &
 s_{BC}
 &=\int_G\chi_{BC}^{a,b;n,k}(g)\,
             \mu_\alpha^{*2}(\dd g),
 \label{eq:V80-BC-integrals}\\
 \operatorname{Re}\langle
 Q_{BC}^{T_k}e,Q_{AC}^{T_k}e\rangle
 &=
 \operatorname{Re}
 \int_G\!\!\int_G
 \chi_{\times}^{a,b;n,k}(g,h)\,
 \mu_\alpha(\dd g)\mu_\alpha(\dd h).
 \label{eq:V80-mixed-integral}
\end{align}
\end{theorem}

\begin{proof}
An invariant operator on the final \(T_k\)-isotypic component is
\(I_{H_{T_k}}\) tensored with its multiplicity matrix.  Therefore its
matrix element on the path vector \(e\) is the normalized carrier
trace of its compression by \(L_{a,b;n,k}\).

Apply this observation to
\[
 Q_{AC}=\int_G
 D^{R_a}(g)\otimes I_{R_b}\otimes D^{Q_n}(g)\,
 \mu_\alpha(\dd g)
\]
to obtain the first formula in
\eqref{eq:V80-AC-integrals}.  The product of two such row
representations at \(g,h\) is the same representation at \(gh\);
integrating gives the convolution law and the second formula.
The \(BC\) equations are identical.

Finally,
\[
 \langle Q_{BC}^{T_k}e,Q_{AC}^{T_k}e\rangle
 =
 \langle e,Q_{BC}^{T_k}Q_{AC}^{T_k}e\rangle
\]
because \(Q_{BC}^{T_k}\) is self-adjoint.  Insert the two Wilson
integrals in their displayed order and use the normalized path trace.
This is \eqref{eq:V80-mixed-integral}.
\end{proof}

\begin{remark}[Why these are reduced path characters]
\label{rem:V80-reduced-characters}
The functions \(\chi_{AC}\) and \(\chi_{BC}\) are normalized traces
of a partial group action on one fixed diagonal-\(G\) path.  They are
not ordinary irreducible characters and need not be central.
Replacing them by \(\varphi_{R_a}\varphi_{Q_n}\) would discard the
final \(T_k\) conditioning.
\end{remark}

\section{Mean, variance, and correlation form of the obstruction}
\label{sec:V80-variance-form}

Put
\begin{equation}
 \xi_{AC}:=(Q_{AC}^{T_k}-m_{AC}I)e,
 \qquad
 \xi_{BC}:=(Q_{BC}^{T_k}-m_{BC}I)e.
 \label{eq:V80-centred-Racah-vectors}
\end{equation}
Then
\begin{equation}
 \xi_{AC}\perp e,\qquad
 \xi_{BC}\perp e,\qquad
 \|\xi_{AC}\|^2=\nu_{AC},\qquad
 \|\xi_{BC}\|^2=\nu_{BC}.
 \label{eq:V80-centred-vector-norms}
\end{equation}
When both variances are nonzero, define
\begin{equation}
 \rho_{a,b;n,k}:=
 \frac{\operatorname{Re}\langle\xi_{BC},\xi_{AC}\rangle}
 {\sqrt{\nu_{BC}\nu_{AC}}};
 \label{eq:V80-correlation}
\end{equation}
set \(\rho=0\) when their product is zero.  Cauchy--Schwarz gives
\(|\rho|\le1\).

\begin{theorem}[Exact three-scalar decision card]
\label{thm:V80-three-scalar-card}
Let
\begin{equation}
 \Delta_{a,b;n,k}:=
 \vartheta_{a,b;n,k}-xm_{BC}-ym_{AC}.
 \label{eq:V80-mean-defect}
\end{equation}
Then
\begin{align}
 \boldsymbol\kappa_{a,b;n,k}
 &=
 \Delta_{a,b;n,k}e-x\xi_{BC}-y\xi_{AC},
 \label{eq:V80-vector-decomposition}\\
 \boxed{
 \|\boldsymbol\kappa_{a,b;n,k}\|^2
 &=
 |\Delta_{a,b;n,k}|^2
 +x^2\nu_{BC}+y^2\nu_{AC}
 +2xy\rho_{a,b;n,k}\sqrt{\nu_{BC}\nu_{AC}}.}
 \label{eq:V80-three-scalar-card}
\end{align}
In particular,
\begin{align}
 |\Delta|^2+
 \left(
 |x|\sqrt{\nu_{BC}}-|y|\sqrt{\nu_{AC}}
 \right)^2
 &\le
 \|\boldsymbol\kappa_{a,b;n,k}\|^2
 \label{eq:V80-card-lower}\\
 &\le
 |\Delta|^2+
 \left(
 |x|\sqrt{\nu_{BC}}+|y|\sqrt{\nu_{AC}}
 \right)^2.
 \label{eq:V80-card-upper}
\end{align}
\end{theorem}

\begin{proof}
Substitute
\[
 Q_{AC}^{T_k}e=m_{AC}e+\xi_{AC},
 \qquad
 Q_{BC}^{T_k}e=m_{BC}e+\xi_{BC}
\]
in \eqref{eq:V79-ancestral-vector-formula}.  Orthogonality in
\eqref{eq:V80-centred-vector-norms} separates the scalar component
from the two centred vectors.  Expanding their squared norm and using
\eqref{eq:V80-correlation} proves
\eqref{eq:V80-three-scalar-card}.  The triangle and reverse-triangle
inequalities for \(x\xi_{BC}+y\xi_{AC}\) prove the two bounds.
\end{proof}

\begin{corollary}[Exact endpoint sum rules]
\label{cor:V80-endpoint-sum-rules}
If \(a=0\) or \(b=0\), then for every \(n,k\),
\begin{equation}
 \Delta_{a,b;n,k}=0,
 \qquad
 x\xi_{BC}+y\xi_{AC}=0.
 \label{eq:V80-endpoint-sum-rules}
\end{equation}
Thus endpoint deletion is a simultaneous mean and centred-vector
cancellation, not the vanishing of the individual moment terms.
\end{corollary}

\begin{proof}
At either endpoint the complete ancestral vector is zero by
\Cref{prop:V77-trivial-input-deletion,
rem:V79-endpoint-cancellation}.  Its component parallel to \(e\) and
its component in \(e^\perp\) vanish separately in the orthogonal
decomposition \eqref{eq:V80-vector-decomposition}.
\end{proof}

\begin{assessment}[The balanced thermal acquisition after V80]
\label{ass:V80-balanced-acquisition}
Uniform decay of the Cartan vector on
\(k\asymp\sqrt n\) requires two simultaneous facts:
\[
 \Delta_{a,b;n,k}\longrightarrow0,
 \qquad
 x\xi_{BC}+y\xi_{AC}\longrightarrow0
\]
at the collision scale.  Equality of the two scalar variances is not
enough; the centred vectors must also have the correct relative
orientation.  Conversely, separate upper bounds on
\(\nu_{AC}\) and \(\nu_{BC}\) cannot see a negative correlation
\(\rho\).

The next representation-theoretic task is therefore to compute the
two path laws \(\pi_{AC},\pi_{BC}\) and their joint Racah angle on the
Cartan lift.  Equations
\eqref{eq:V80-AC-integrals}--\eqref{eq:V80-mixed-integral} provide a
common integral in which that computation can be made without
entrywise absolute values.
\end{assessment}

\begin{firewall}[Core obstruction versus complete collision capacity]
\label{fire:V80-core-versus-capacity}
If \eqref{eq:V80-card-lower} proves nondecay for a balanced sequence,
that establishes a core product-column obstruction.  In accordance
with the V80 audit, it does not by itself disprove the complete
collision estimate: the external \(oa\) and \(bc\) pair banks and
their \(P_{ab}\) factor must then be inserted in the same orientation.
\end{firewall}

\section{V80 ledger and next acquisition}
\label{sec:V80-ledger}

\begin{theorem}[Integrated V80 statement]
\label{thm:V80-integrated}
Preserving V1--V79, V80 proves:
\begin{enumerate}[label=\textup{(V80.\arabic*)},leftmargin=5.7em]
\item the stable mandatory SM V83/NS V31 audit and nontransfer
      firewall;
\item the exact path probabilities
      \eqref{eq:V80-path-probabilities};
\item all four unmixed Racah moments
      \eqref{eq:V80-AC-moments}--%
      \eqref{eq:V80-BC-moments};
\item the five common path-integral formulas
      \eqref{eq:V80-AC-integrals}--%
      \eqref{eq:V80-mixed-integral};
\item the mean--variance--correlation identity
      \eqref{eq:V80-three-scalar-card}; and
\item the exact endpoint sum rules
      \eqref{eq:V80-endpoint-sum-rules}.
\end{enumerate}
\end{theorem}

\begin{proof}
These are
\Cref{lem:V80-path-probabilities,
thm:V80-first-second-moments,thm:V80-path-integrals,
thm:V80-three-scalar-card,cor:V80-endpoint-sum-rules}.
\end{proof}

\begin{openproblem}[V81 mandate: explicit Cartan path laws]
\label{op:V81-mandate}
For \(a+b=n\), \(T_k=(k,k)\), and
\(k\asymp\sqrt n\):
\begin{enumerate}
\item identify all \(AC\) and \(BC\) intermediates which contribute
      to the Cartan path and their exact Racah probabilities
      \(\pi_{AC},\pi_{BC}\);
\item exploit the exchange symmetry \(a\leftrightarrow b\) before
      estimating the balanced region;
\item compute the reduced path characters in
      \eqref{eq:V80-path-character-AC}--%
      \eqref{eq:V80-mixed-path-character} as explicit orthogonal
      polynomials or terminating hypergeometric sums;
\item obtain asymptotics of \(\Delta,\nu_{AC},\nu_{BC},\rho\)
      uniformly for \(a/n\) in compact subsets of \((0,1)\);
\item if the lower card shows nondecay, insert \(P_{ab}\) immediately
      in the same orientation; if the upper card decays, sum it with
      the exact V79 dimension weights.
\end{enumerate}
The next compulsory companion audit is V85.
\end{openproblem}

\section{Firewall after V80}
\label{sec:V80-firewall}

\begin{firewall}[Current claim boundary]
\label{fire:V80-final}
V80 performs the mandatory audit and reduces the six balanced-Cartan
moments to two spectral probability laws, two variances, one mean
defect, and one Racah correlation.  It does not yet compute their
thermal asymptotics, close the product-state collision gate, prove
JREC, close simultaneous sideways endpoints, aggregate punctures or
multiple collisions, construct the continuum Yang--Mills theory, or
prove a positive mass gap.
\end{firewall}

\section{Append-only record for V80}
\label{sec:V80-record}

V80 was appended to the complete V79 manuscript.  The exact V79
parent had \(56055\) validation lines, \(1996483\) bytes, and
SHA--256 digest
\[
 \texttt{cdbbccac70508a092fa3915bb1bbcadf341284c58b0655d38fa8210178d43d43}.
\]
No V79 mathematical content was changed.  The sole final
\verb|\end{document}| is restored below.

% =============================================================================
% END APPENDED VERSION V80 MATERIAL
% =============================================================================
% =============================================================================
% BEGIN APPENDED VERSION V81 MATERIAL
% =============================================================================

\chapter{V81: Pieri support and exact exchange-parity cancellation
on the balanced Cartan lift}
\label{ch:V81}

\section{Purpose}
\label{sec:V81-purpose}

V80 expressed the hard Cartan-path vector through two Racah
probability laws and their correlation.  V81 first determines the
exact support of those laws by the \(\mathrm{SU}(3)\) Pieri rule.
It then treats the exactly balanced split \(a=b\).  Interchange of
the two identical symmetric-power factors conjugates \(Q_{AC}\) to
\(Q_{BC}\), while the \(r\)-th \(AB\) ancestry summand has exchange
parity \((-1)^r\).  Consequently every odd-ancestry component of the
two centred pair vectors cancels in the signed cumulant.  The
remaining vector loss is precisely its even-ancestry leakage.

\begin{firewall}[Scope of V81]
\label{fire:V81-scope}
The parity cancellation is exact, but V81 does not estimate the
remaining even leakage or the scalar mean defect on
\(k\asymp\sqrt n\).  It therefore reduces rather than closes the
balanced collision-core gate.
\end{firewall}

\section{Pieri support of the two path laws}
\label{sec:V81-Pieri-support}

Retain \(a+b=n\) and \(T_k=(k,k)\).  The symmetric--dual Pieri rule
gives
\begin{align}
 R_a\otimes Q_n
 &=
 \bigoplus_{j=0}^{a}U_j,
 &
 U_j&:=(a-j,n-j),
 \label{eq:V81-AC-Pieri}\\
 R_b\otimes Q_n
 &=
 \bigoplus_{\ell=0}^{b}V_\ell,
 &
 V_\ell&:=(b-\ell,n-\ell).
 \label{eq:V81-BC-Pieri}
\end{align}
Both decompositions are multiplicity free.

\begin{theorem}[Exact support reaching the final \(T_k\)]
\label{thm:V81-Pieri-support}
An \(AC\)-intermediate \(U_j\) can couple with \(R_b\) to \(T_k\) if
and only if
\begin{equation}
 \max\{0,a-k\}\le j\le\min\{a,n-k\}.
 \label{eq:V81-AC-support}
\end{equation}
Every such coupling has multiplicity one.  Similarly, a
\(BC\)-intermediate \(V_\ell\) contributes if and only if
\begin{equation}
 \max\{0,b-k\}\le \ell\le\min\{b,n-k\},
 \label{eq:V81-BC-support}
\end{equation}
again with multiplicity one.

In particular, if \(k\le\min\{a,b\}\), both supports have exactly
\(k+1\) members:
\[
 j=a-k,\dots,a,\qquad
 \ell=b-k,\dots,b.
\]
\end{theorem}

\begin{proof}
The first two decompositions are the standard contraction form of
the Pieri rule:
\[
 (a,0)\otimes(0,n)
 =
 \bigoplus_{j=0}^{\min\{a,n\}}(a-j,n-j),
\]
and \(a\le n\).  To determine which \(U_j=(p,q)\), with
\[
 p=a-j,\qquad q=n-j,
\]
reaches \(T_k\) after tensoring by \(R_b=(b,0)\), use the
horizontal-strip form of Pieri.  Add nonnegative numbers of boxes
\(r_1,r_2,r_3\) to the three rows, with
\[
 r_1+r_2+r_3=b,\qquad r_2\le p,\qquad r_3\le q.
\]
The resulting Dynkin labels are
\[
 (p+r_1-r_2,\ q+r_2-r_3).
\]
Equating both labels to \(k\) gives
\[
 r_1=k-a+j+r_2,\qquad
 r_3=n-j+r_2-k.
\]
Their sum is \(b+3r_2\), because \(n-a=b\).  Hence the Pieri sum
condition forces
\[
 r_2=0,\qquad
 r_1=k-a+j,\qquad
 r_3=n-j-k.
\]
Nonnegativity is exactly
\[
 j\ge a-k,\qquad j\le n-k,
\]
together with \(0\le j\le a\).  This proves
\eqref{eq:V81-AC-support}.  The solution \((r_1,r_2,r_3)\) is unique,
so its Pieri multiplicity is one.  Interchange \(a,b\) for
\eqref{eq:V81-BC-support}.  If
\(k\le\min\{a,b\}\), the lower and upper endpoints reduce to
\(a-k,a\) and \(b-k,b\), respectively.
\end{proof}

\begin{corollary}[Finite forms of the V80 laws]
\label{cor:V81-finite-path-laws}
On \(k\le\min\{a,b\}\), the V80 means and variances are supported on
only \(k+1\) Wilson multipliers:
\begin{align}
 m_{AC}
 &=
 \sum_{r=0}^{k}
 q_{\alpha,(k-r,b+k-r)}
 \,\pi_{AC}(U_{a-k+r}),
 \label{eq:V81-AC-finite-mean}\\
 m_{BC}
 &=
 \sum_{r=0}^{k}
 q_{\alpha,(k-r,a+k-r)}
 \,\pi_{BC}(V_{b-k+r}).
 \label{eq:V81-BC-finite-mean}
\end{align}
The second moments have the same formulas with the multipliers
squared.
\end{corollary}

\begin{proof}
For \(j=a-k+r\),
\[
 U_j=(a-j,n-j)=(k-r,b+k-r).
\]
Insert this reindexing in \eqref{eq:V80-path-means}.  The \(BC\)
formula is its \(a,b\) interchange.  The statement for second
moments follows from \eqref{eq:V80-path-seconds}.
\end{proof}

\begin{remark}[Thermal support width]
\label{rem:V81-thermal-support-width}
In the balanced regime \(a,b\asymp n\) and
\(k\asymp\sqrt n\), each conditional path law has only
\(O(\sqrt n)\) atoms, not \(O(n)\).  This support count alone is not
a probability debit, because the atoms have total mass one.
\end{remark}

\section{Exchange symmetry of the path laws}
\label{sec:V81-exchange}

Let \(\mathsf F_{AB}\) denote the unitary which interchanges the first
two tensor factors, using the canonical identification when their
representations are equal.

\begin{proposition}[Exchange covariance]
\label{prop:V81-exchange-covariance}
Under interchange \(a\leftrightarrow b\),
\begin{align}
 \pi_{AC}^{a,b;n,k}(U_j)
 &=
 \pi_{BC}^{b,a;n,k}(U_j),
 \label{eq:V81-law-exchange}\\
 m_{AC}^{a,b;n,k}
 &=
 m_{BC}^{b,a;n,k},
 \qquad
 \nu_{AC}^{a,b;n,k}
 =
 \nu_{BC}^{b,a;n,k}.
 \label{eq:V81-moment-exchange}
\end{align}
If \(a=b\), then
\begin{equation}
 m_{AC}=m_{BC}=:m_*,
 \qquad
 \nu_{AC}=\nu_{BC}=:\nu_*.
 \label{eq:V81-balanced-equal-moments}
\end{equation}
\end{proposition}

\begin{proof}
The flip sends the Cartan embedding
\(R_a\otimes R_b\to R_n\) to the Cartan embedding
\(R_b\otimes R_a\to R_n\), with the same normalization.  It fixes the
subsequent \(R_n\otimes Q_n\to T_k\) arrow and interchanges the
\(AC\)- and \(BC\)-first spectral resolutions.  Taking squared
projection norms gives \eqref{eq:V81-law-exchange}.  The moment
identities follow by summing the same Wilson multipliers against the
interchanged laws.  Set \(a=b\) for
\eqref{eq:V81-balanced-equal-moments}.
\end{proof}

\section{Parity of the equal symmetric-power tensor square}
\label{sec:V81-parity}

Set \(n=2m\) and \(a=b=m\).  The \(AB\) decomposition is
\begin{equation}
 R_m\otimes R_m
 =
 \bigoplus_{r=0}^{m}S_r,
 \qquad
 S_r=(2m-2r,r),
 \label{eq:V81-symmetric-square-decomposition}
\end{equation}
with multiplicity one.  Let
\(\mathsf P_r^{T_k}\) be the projection in the final multiplicity
space onto all paths whose \(AB\)-intermediate is \(S_r\).

\begin{lemma}[Exact exchange parity]
\label{lem:V81-exchange-parity}
On the \(S_r\) summand,
\begin{equation}
 \mathsf F_{AB}=(-1)^rI.
 \label{eq:V81-exchange-parity}
\end{equation}
In particular the Cartan path \(e=e_{m,m;2m,k}\), which has \(r=0\),
is fixed by \(\mathsf F_{AB}\).
\end{lemma}

\begin{proof}
The polynomial Pieri decomposition of
\(\operatorname{Sym}^mV\otimes\operatorname{Sym}^mV\) is indexed by
the two-row Young diagrams
\[
 (2m-r,r),\qquad 0\le r\le m,
\]
whose SU(3) Dynkin labels are
\((2m-2r,r)\).  The corresponding Young symmetrizer uses \(r\)
two-box columns joining the two tensor factors.  Interchanging the
factors reverses each of these antisymmetrized column pairs and hence
acts by \((-1)^r\).  The remaining row symmetrizations are fixed.
This proves \eqref{eq:V81-exchange-parity}.  The Cartan summand is
\(r=0\).
\end{proof}

\section{Exact cancellation of all odd ancestry leakage}
\label{sec:V81-odd-cancellation}

In the balanced case, let
\[
 \xi:=\xi_{AC}=(Q_{AC}^{T_k}-m_*I)e.
\]

\begin{lemma}[The second centred vector is the flipped first]
\label{lem:V81-centred-flip}
One has
\begin{equation}
 \xi_{BC}=\mathsf F_{AB}\xi.
 \label{eq:V81-centred-flip}
\end{equation}
Writing
\[
 \xi_r:=\mathsf P_r^{T_k}\xi,
\]
one has \(\xi_0=0\) and
\begin{align}
 \nu_*&=\sum_{r=1}^{m}\|\xi_r\|^2,
 \label{eq:V81-variance-parity-sum}\\
 \operatorname{Re}\langle\xi_{BC},\xi_{AC}\rangle
 &=
 \sum_{r=1}^{m}(-1)^r\|\xi_r\|^2.
 \label{eq:V81-correlation-parity-sum}
\end{align}
\end{lemma}

\begin{proof}
For identical first two representations,
\[
 Q_{BC}^{T_k}
 =
 \mathsf F_{AB}Q_{AC}^{T_k}\mathsf F_{AB}.
\]
Use \(\mathsf F_{AB}e=e\) and the equality of means in
\eqref{eq:V81-balanced-equal-moments} to prove
\eqref{eq:V81-centred-flip}.  The vector \(\xi\) is orthogonal to
\(e\).  The \(r=0\) path from \(R_{2m}\otimes Q_{2m}\) to \(T_k\)
has multiplicity one, so its entire \(r=0\) subspace is spanned by
\(e\); hence \(\xi_0=0\).  Orthogonality of the ancestry resolution
gives \eqref{eq:V81-variance-parity-sum}.  Apply the parity
\eqref{eq:V81-exchange-parity} to each \(\xi_r\) for
\eqref{eq:V81-correlation-parity-sum}.
\end{proof}

Define the even and odd leakage masses
\begin{equation}
 E_{m,k}^{\rm even}:=
 \sum_{\substack{2\le r\le m\\r\ {\rm even}}}\|\xi_r\|^2,
 \qquad
 E_{m,k}^{\rm odd}:=
 \sum_{\substack{1\le r\le m\\r\ {\rm odd}}}\|\xi_r\|^2.
 \label{eq:V81-even-odd-leakage}
\end{equation}

\begin{theorem}[Balanced parity cancellation identity]
\label{thm:V81-balanced-parity-cancellation}
Let
\[
 q_m:=q_{\alpha,R_m},
 \qquad
 \Delta_{m,k}:=
 \vartheta_{m,m;2m,k}-2q_m m_*.
\]
Then
\begin{equation}
 \boxed{
 \|\boldsymbol\kappa_{m,m;2m,k}\|^2
 =
 |\Delta_{m,k}|^2
 +4q_m^2E_{m,k}^{\rm even}.}
 \label{eq:V81-balanced-parity-cancellation}
\end{equation}
Every odd-ancestry component cancels exactly.  If \(\nu_*>0\), the
V80 correlation coefficient is
\begin{equation}
 \rho_{m,m;2m,k}
 =
 \frac{E_{m,k}^{\rm even}-E_{m,k}^{\rm odd}}
 {E_{m,k}^{\rm even}+E_{m,k}^{\rm odd}}
 =
 -1+\frac{2E_{m,k}^{\rm even}}{\nu_*}.
 \label{eq:V81-balanced-correlation}
\end{equation}
\end{theorem}

\begin{proof}
Equation \eqref{eq:V80-vector-decomposition}, with \(x=y=q_m\), is
\[
 \boldsymbol\kappa
 =
 \Delta_{m,k}e-q_m(I+\mathsf F_{AB})\xi.
\]
The scalar vector \(e\) is orthogonal to \(\xi\).  By
\eqref{eq:V81-exchange-parity},
\[
 (I+\mathsf F_{AB})\xi
 =
 2\sum_{\substack{r\ge2\\r\ {\rm even}}}\xi_r.
\]
Orthogonality of the ancestry components proves
\eqref{eq:V81-balanced-parity-cancellation}.  The correlation formula
follows from
\eqref{eq:V81-variance-parity-sum}--%
\eqref{eq:V81-correlation-parity-sum}.
\end{proof}

\begin{corollary}[Exact balanced thermal target]
\label{cor:V81-balanced-thermal-target}
For an output index set \(I\subseteq\{0,\dots,2m\}\), the balanced
Cartan product-column energy is
\begin{align}
 &\sum_{k\in I}
 \frac{(k+1)^3}{d_{2m}}
 [1+s_\alpha\Xi(T_k)]^2
 \notag\\
 &\qquad\cdot
 \left[
 |\Delta_{m,k}|^2
 +4q_m^2E_{m,k}^{\rm even}
 \right].
 \label{eq:V81-balanced-thermal-target}
\end{align}
Thus closure on \(k\asymp\sqrt m\) requires only the scalar mean
defect and the even ancestry leakage; the entire odd leakage is
irrelevant to the signed cumulant.
\end{corollary}

\begin{proof}
Insert
\eqref{eq:V81-balanced-parity-cancellation} into
\eqref{eq:V79-Cartan-shell-energy}.
\end{proof}

\begin{assessment}[What parity has bought]
\label{ass:V81-parity-gain}
The separate variance bound of V80 treats
\(\nu_*=E^{\rm even}+E^{\rm odd}\).  The exact signed square instead
keeps only \(E^{\rm even}\).  Hence a Racah law with most centred mass
on odd \(AB\) ancestries gives a genuine cancellation invisible to
every positive bound on \(Q_{AC}\) and \(Q_{BC}\) separately.

The remaining question is now sharply two-part:
\[
 \Delta_{m,k}
 \quad\hbox{and}\quad
 E_{m,k}^{\rm even},
 \qquad k\asymp\sqrt m.
\]
Computing the total variance without its ancestry parity is no longer
sufficient.
\end{assessment}

\section{V81 ledger and next acquisition}
\label{sec:V81-ledger}

\begin{theorem}[Integrated V81 statement]
\label{thm:V81-integrated}
Preserving V1--V80, V81 proves:
\begin{enumerate}[label=\textup{(V81.\arabic*)},leftmargin=5.7em]
\item the exact Pieri supports
      \eqref{eq:V81-AC-support}--%
      \eqref{eq:V81-BC-support};
\item their \(k+1\)-atom thermal specialization;
\item exchange covariance of both path laws and variances;
\item the ancestry parity
      \eqref{eq:V81-exchange-parity};
\item the centred correlation sum
      \eqref{eq:V81-correlation-parity-sum}; and
\item the exact cancellation identity
      \eqref{eq:V81-balanced-parity-cancellation}.
\end{enumerate}
\end{theorem}

\begin{proof}
These are
\Cref{thm:V81-Pieri-support,cor:V81-finite-path-laws,
prop:V81-exchange-covariance,lem:V81-exchange-parity,
lem:V81-centred-flip,thm:V81-balanced-parity-cancellation}.
\end{proof}

\begin{openproblem}[V82 mandate: scalar defect and even leakage]
\label{op:V82-mandate}
On \(a=b=m\) and \(k\asymp\sqrt m\):
\begin{enumerate}
\item compute the \(k+1\) probabilities in
      \eqref{eq:V81-AC-finite-mean}, using the symmetric-power
      recoupling normalization rather than a guessed hypergeometric
      law;
\item evaluate the scalar defect \(\Delta_{m,k}\) from that law;
\item resolve \(Q_{AC}^{T_k}e\) back into the \(AB\) ancestries
      \(S_r=(2m-2r,r)\) and compute the even mass
      \(E_{m,k}^{\rm even}\);
\item use the parity projection
      \(\frac12(I+\mathsf F_{AB})\) before any absolute value;
\item if the balanced case closes, extend by perturbation to
      \(|a-b|=o(n)\); if it fails, insert the same-orientation pair
      bank before drawing a collision conclusion.
\end{enumerate}
The next compulsory companion audit is V85.
\end{openproblem}

\section{Firewall after V81}
\label{sec:V81-firewall}

\begin{firewall}[Current claim boundary]
\label{fire:V81-final}
V81 proves exact odd-ancestry cancellation in the balanced Cartan
lift.  It does not estimate the scalar defect or even leakage,
control unbalanced lifts, close the product-state collision gate,
prove JREC, close simultaneous sideways endpoints, aggregate
punctures or multiple collisions, construct the continuum
Yang--Mills theory, or prove a positive mass gap.
\end{firewall}

\section{Append-only record for V81}
\label{sec:V81-record}

V81 was appended to the complete V80 manuscript.  The exact V80
parent had \(56584\) validation lines, \(2013230\) bytes, and
SHA--256 digest
\[
 \texttt{c78ca926e5986e5fe72111cd2d0980020fe9ff06b009aaedd3f6b43f138f273c}.
\]
No V80 mathematical content was changed.  The sole final
\verb|\end{document}| is restored below.

% =============================================================================
% END APPENDED VERSION V81 MATERIAL
% =============================================================================
% =============================================================================
% BEGIN APPENDED VERSION V82 MATERIAL
% =============================================================================

\chapter{V82: Pair-Casimir affine cancellation and a curvature-only
bound for the even leakage}
\label{ch:V82}

\section{Purpose}
\label{sec:V82-purpose}

V81 showed that the balanced Cartan cumulant sees only a scalar mean
defect and the even \(AB\)-ancestry part of
\((Q_{AC}-m_*I)e\).  V82 proves that this even vector part vanishes
for every affine function of the \(AC\) pair Casimir.  Therefore it
is controlled by the best affine approximation to the Wilson
multipliers on the \(k+1\) Pieri nodes.  The individual Racah
probabilities disappear from the resulting upper bound.

This is an upstream reduction: the missing estimate is now a
second-divided-difference bound for the Wilson Fourier multiplier
along one explicit representation segment, together with one scalar
compatibility defect.

\begin{firewall}[Scope of V82]
\label{fire:V82-scope}
No uniform second-divided-difference estimate is proved here.  The
existing two-sided gap and tail bounds control the size of each
multiplier but not its affine curvature across a growing family of
representations.  V82 identifies that stronger input explicitly.
\end{firewall}

\section{The pair-Casimir sum on a fixed ancestry vector}
\label{sec:V82-Casimir-sum}

Let \(\mathbf C_{AB}\), \(\mathbf C_{AC}\), and
\(\mathbf C_{BC}\) be the quadratic Casimir operators for the three
diagonal pair actions on \(H_A\otimes H_B\otimes H_C\).  Let
\(\mathbf C_{ABC}\) be the total diagonal Casimir.  Write
\(c_R:=C_2(R)\).

\begin{lemma}[Exact pair-Casimir identity]
\label{lem:V82-pair-Casimir-identity}
On the complete tensor carrier,
\begin{equation}
 \boxed{
 \mathbf C_{AC}+\mathbf C_{BC}
 =
 (c_A+c_B+c_C)I
 +\mathbf C_{ABC}-\mathbf C_{AB}.}
 \label{eq:V82-pair-Casimir-identity}
\end{equation}
Consequently, on a path vector \(e_{S,\alpha,\beta}^{T}\),
\begin{equation}
 (\mathbf C_{AC}+\mathbf C_{BC})e_{S,\alpha,\beta}^{T}
 =
 \Lambda_{A,B,C;S,T}e_{S,\alpha,\beta}^{T},
 \label{eq:V82-pair-Casimir-path}
\end{equation}
where
\begin{equation}
 \Lambda_{A,B,C;S,T}
 :=
 c_A+c_B+c_C+c_T-c_S.
 \label{eq:V82-Lambda-general}
\end{equation}
\end{lemma}

\begin{proof}
Choose an orthonormal Lie-algebra basis and write
\(\Omega_{IJ}=\sum_\ell X_\ell^{(I)}X_\ell^{(J)}\).  Then
\[
 \mathbf C_{AC}=(c_A+c_C)I+2\Omega_{AC},
 \qquad
 \mathbf C_{BC}=(c_B+c_C)I+2\Omega_{BC}.
\]
Also
\[
 \mathbf C_{ABC}
 =
 \mathbf C_{AB}+c_CI
 +2\Omega_{AC}+2\Omega_{BC}.
\]
Eliminating the two cross terms proves
\eqref{eq:V82-pair-Casimir-identity}.  On the stated path,
\(\mathbf C_{ABC}\) and \(\mathbf C_{AB}\) have eigenvalues
\(c_T\) and \(c_S\), respectively, proving
\eqref{eq:V82-pair-Casimir-path}.
\end{proof}

For the balanced Cartan lift
\[
 A=B=R_m,\qquad
 C=Q_{2m},\qquad
 S=R_{2m},\qquad
 T=T_k,
\]
duality gives \(c_C=c_S\), and hence
\begin{equation}
 \Lambda_{m,k}:=
 \Lambda_{R_m,R_m,Q_{2m};R_{2m},T_k}
 =
 2c_{R_m}+c_{T_k}.
 \label{eq:V82-Lambda-balanced}
\end{equation}

\begin{corollary}[Mean pair Casimir]
\label{cor:V82-mean-pair-Casimir}
On the balanced Cartan path vector \(e\),
\begin{equation}
 \langle e,\mathbf C_{AC}e\rangle
 =
 \langle e,\mathbf C_{BC}e\rangle
 =
 \frac{\Lambda_{m,k}}2.
 \label{eq:V82-mean-pair-Casimir}
\end{equation}
\end{corollary}

\begin{proof}
The two expectations are equal by exchange symmetry.  Their sum is
\(\Lambda_{m,k}\) by \eqref{eq:V82-pair-Casimir-path}.
\end{proof}

\section{Affine functions have no even vector leakage}
\label{sec:V82-affine-cancellation}

Let
\[
 \mathsf P_+:=\frac12(I+\mathsf F_{AB}),
 \qquad
 \mathsf P_e:=|e\rangle\langle e|
\]
on the final multiplicity space.  By V81,
\begin{equation}
 E_{m,k}^{\rm even}
 =
 \|(\mathsf P_+-\mathsf P_e)Q_{AC}^{T_k}e\|^2.
 \label{eq:V82-even-projection}
\end{equation}

\begin{theorem}[Exact affine-Casimir annihilation]
\label{thm:V82-affine-annihilation}
For every affine function
\(\ell(c)=\gamma_0+\gamma_1c\),
\begin{equation}
 \boxed{
 (\mathsf P_+-\mathsf P_e)
 \ell(\mathbf C_{AC})e=0.}
 \label{eq:V82-affine-annihilation}
\end{equation}
\end{theorem}

\begin{proof}
The flip fixes \(e\) and conjugates
\(\mathbf C_{AC}\) to \(\mathbf C_{BC}\).  Therefore
\begin{align*}
 \mathsf P_+\ell(\mathbf C_{AC})e
 &=
 \frac12[
 \ell(\mathbf C_{AC})+\ell(\mathbf C_{BC})]e\\
 &=
 \left[
 \gamma_0+
 \frac{\gamma_1}{2}
 (\mathbf C_{AC}+\mathbf C_{BC})
 \right]e\\
 &=
 \ell(\Lambda_{m,k}/2)e,
\end{align*}
using \eqref{eq:V82-pair-Casimir-path}.  This vector lies in the
range of \(\mathsf P_e\), proving
\eqref{eq:V82-affine-annihilation}.
\end{proof}

\section{Best affine error of the Wilson multiplier}
\label{sec:V82-affine-error}

On the balanced thermal support from V81, put
\begin{equation}
 U_r:=(k-r,m+k-r),\qquad 0\le r\le k,
 \label{eq:V82-balanced-Pieri-nodes}
\end{equation}
and define
\begin{equation}
 \epsilon_{m,k}^{\rm aff}
 :=
 \inf_{\gamma_0,\gamma_1\in\mathbb R}
 \max_{0\le r\le k}
 \left|
 q_{\alpha,U_r}
 -\gamma_0-\gamma_1c_{U_r}
 \right|.
 \label{eq:V82-best-affine-error}
\end{equation}

\begin{theorem}[Curvature-only even-leakage bound]
\label{thm:V82-even-leakage-bound}
The exact even leakage in
\eqref{eq:V81-balanced-parity-cancellation} satisfies
\begin{equation}
 \boxed{
 E_{m,k}^{\rm even}
 \le
 (\epsilon_{m,k}^{\rm aff})^2.}
 \label{eq:V82-even-leakage-bound}
\end{equation}
Moreover, for every affine \(\ell\) with uniform residual
\[
 \epsilon(\ell):=
 \max_{0\le r\le k}
 |q_{\alpha,U_r}-\ell(c_{U_r})|,
\]
the path mean satisfies
\begin{equation}
 \left|
 m_*-\ell(\Lambda_{m,k}/2)
 \right|
 \le\epsilon(\ell).
 \label{eq:V82-mean-affine-error}
\end{equation}
\end{theorem}

\begin{proof}
On the spectral support of \(e\) for \(\mathbf C_{AC}\),
\[
 Q_{AC}e
 =
 \sum_{r=0}^{k}
 q_{\alpha,U_r}\mathsf P_{AC}(U_r)e.
\]
Hence
\[
 \|[Q_{AC}-\ell(\mathbf C_{AC})]e\|
 \le\epsilon(\ell).
\]
Apply the contraction
\(\mathsf P_+-\mathsf P_e\).  The affine term vanishes by
\eqref{eq:V82-affine-annihilation}, while
\eqref{eq:V82-even-projection} identifies the remaining norm with
\((E_{m,k}^{\rm even})^{1/2}\).  Minimize over \(\ell\) to prove
\eqref{eq:V82-even-leakage-bound}.

For the mean, take the expectation of
\(Q_{AC}-\ell(\mathbf C_{AC})\) on \(e\).  Its modulus is at most
\(\epsilon(\ell)\), and
\eqref{eq:V82-mean-pair-Casimir} gives
\(\langle e,\ell(\mathbf C_{AC})e\rangle
=\ell(\Lambda_{m,k}/2)\).
\end{proof}

\begin{corollary}[Affine master bound for the signed vector]
\label{cor:V82-affine-master-bound}
For every affine \(\ell\),
\begin{align}
 \|\boldsymbol\kappa_{m,m;2m,k}\|^2
 &\le
 \left(
 \left|
 \vartheta_{m,m;2m,k}
 -2q_{\alpha,R_m}\ell(\Lambda_{m,k}/2)
 \right|
 +2|q_{\alpha,R_m}|\epsilon(\ell)
 \right)^2
 \notag\\
 &\quad+
 4q_{\alpha,R_m}^2\epsilon(\ell)^2.
 \label{eq:V82-affine-master-bound}
\end{align}
In particular, the best bound is obtained by taking the infimum of
the right side over all affine \(\ell\).
\end{corollary}

\begin{proof}
The balanced parity identity
\eqref{eq:V81-balanced-parity-cancellation} gives
\[
 \|\boldsymbol\kappa\|^2
 =
 |\vartheta-2q_{\alpha,R_m}m_*|^2
 +4q_{\alpha,R_m}^2E_{m,k}^{\rm even}.
\]
Use \eqref{eq:V82-mean-affine-error} in the scalar term and
\eqref{eq:V82-even-leakage-bound} in the vector term.
\end{proof}

\begin{assessment}[A precise conditional closure criterion]
\label{ass:V82-conditional-closure}
The balanced thermal core closes if one can choose affine functions
\(\ell_{m,k}\) such that the right side of
\eqref{eq:V82-affine-master-bound}, after multiplication by
\([1+s_\alpha\Xi(T_k)]^2\), is summable against
\((k+1)^3/d_{2m}\) at the collision scale.  This criterion no longer
contains an unknown Racah probability or correlation coefficient.

It has two independent requirements:
\begin{enumerate}[label=\textup{(\roman*)}]
\item the Wilson values on the Pieri segment must be nearly affine
      in pair Casimir; and
\item that affine value at the exact mean Casimir must match the
      scalar coefficient
      \(\vartheta/(2q_{\alpha,R_m})\).
\end{enumerate}
Both are multiplier statements.  Neither follows from a pointwise
two-sided character gap.
\end{assessment}

\section{Explicit Casimir nodes and divided differences}
\label{sec:V82-divided-differences}

With the standard SU(3) normalization
\[
 C_2(p,q)=\frac{p^2+q^2+pq+3p+3q}{3},
\]
the node \(U_r\) in
\eqref{eq:V82-balanced-Pieri-nodes} has
\begin{equation}
 c_r:=c_{U_r}
 =
 (k-r)^2+m(k-r)+2(k-r)+\frac{m^2}{3}+m.
 \label{eq:V82-Casimir-nodes}
\end{equation}
Thus
\begin{align}
 c_r-c_{r+1}
 &=m+2(k-r)+1,
 \label{eq:V82-node-spacing}\\
 c_0-c_k
 &=k(m+k+2).
 \label{eq:V82-node-diameter}
\end{align}

For \(k\ge2\), define the discrete curvature
\begin{equation}
 \mathfrak D_{m,k}^{(2)}
 :=
 \max_{0\le r<s<t\le k}
 \left|
 q_{\alpha,U_r},
 q_{\alpha,U_s},
 q_{\alpha,U_t}
 \right|_{[c_r,c_s,c_t]},
 \label{eq:V82-discrete-curvature}
\end{equation}
where the notation on the right means the absolute second divided
difference of the scalar data \(c\mapsto q\).  Set
\(\mathfrak D_{m,k}^{(2)}=0\) for \(k\le1\).

\begin{proposition}[Secant bound for the affine error]
\label{prop:V82-secant-bound}
One has
\begin{equation}
 \boxed{
 \epsilon_{m,k}^{\rm aff}
 \le
 \frac14
 \mathfrak D_{m,k}^{(2)}
 k^2(m+k+2)^2.}
 \label{eq:V82-secant-bound}
\end{equation}
\end{proposition}

\begin{proof}
For \(k\le1\), any data on at most two distinct nodes are exactly
affine.  For \(k\ge2\), let \(\ell\) be the secant through
\((c_0,q_{\alpha,U_0})\) and
\((c_k,q_{\alpha,U_k})\).  The divided-difference interpolation
identity gives, for every interior node \(c_r\),
\[
 q_{\alpha,U_r}-\ell(c_r)
 =
 [c_0,c_r,c_k]q\,
 (c_r-c_0)(c_r-c_k).
\]
The product of the two distances is at most one quarter of the
squared interval length.  Use
\eqref{eq:V82-node-diameter}, then minimize over affine functions.
\end{proof}

\begin{firewall}[The new multiplier bottleneck]
\label{fire:V82-curvature-bottleneck}
The manuscript proves
\(1-q_{\alpha,R}\asymp\min\{s_\alpha C_2(R),1\}\) and several
weighted pointwise tail bounds.  Those statements do not control
\(\mathfrak D_{m,k}^{(2)}\).  Treating
\eqref{eq:V82-secant-bound} as small without a uniform
divided-difference theorem would be circular.
\end{firewall}

\section{V82 ledger and next acquisition}
\label{sec:V82-ledger}

\begin{theorem}[Integrated V82 statement]
\label{thm:V82-integrated}
Preserving V1--V81, V82 proves:
\begin{enumerate}[label=\textup{(V82.\arabic*)},leftmargin=5.7em]
\item the pair-Casimir identity
      \eqref{eq:V82-pair-Casimir-identity};
\item the exact balanced mean
      \eqref{eq:V82-mean-pair-Casimir};
\item annihilation of affine Casimir functions
      \eqref{eq:V82-affine-annihilation};
\item the curvature-only leakage bound
      \eqref{eq:V82-even-leakage-bound};
\item the affine master estimate
      \eqref{eq:V82-affine-master-bound}; and
\item the explicit divided-difference reduction
      \eqref{eq:V82-secant-bound}.
\end{enumerate}
\end{theorem}

\begin{proof}
These are
\Cref{lem:V82-pair-Casimir-identity,
cor:V82-mean-pair-Casimir,thm:V82-affine-annihilation,
thm:V82-even-leakage-bound,cor:V82-affine-master-bound,
prop:V82-secant-bound}.
\end{proof}

\begin{openproblem}[V83 mandate: Wilson multiplier curvature]
\label{op:V83-mandate}
Return to the one-link central Wilson density and:
\begin{enumerate}
\item prove a representation-uniform second-divided-difference bound
      for \(q_{\alpha,(p,q)}\) along the explicit nodes
      \eqref{eq:V82-Casimir-nodes};
\item distinguish the cold range
      \(s_\alpha c_r\ll1\), the transition range, and the remote tail;
\item in the cold range, upgrade the fixed-representation expansion
      \(q_{\alpha,R}=1-s_\alpha C_2(R)+O_R(s_\alpha^2)\)
      to a remainder uniform on the growing Pieri segment;
\item in the tail, retain the available multiplier decay before
      applying \eqref{eq:V82-secant-bound};
\item estimate the scalar compatibility term in
      \eqref{eq:V82-affine-master-bound} with the same affine
      function;
\item if such curvature is unavailable for the Wilson law, return
      further upstream and compare the central convolution directly
      with the heat semigroup on this restricted representation
      segment.
\end{enumerate}
The next compulsory companion audit is V85.
\end{openproblem}

\section{Firewall after V82}
\label{sec:V82-firewall}

\begin{firewall}[Current claim boundary]
\label{fire:V82-final}
V82 proves that only non-affine Wilson multiplier curvature can feed
the even ancestry leakage in the balanced Cartan lift.  It does not
bound that curvature uniformly, control nearby unbalanced lifts,
close the product-state collision gate, prove JREC, close
simultaneous sideways endpoints, aggregate punctures or multiple
collisions, construct the continuum Yang--Mills theory, or prove a
positive mass gap.
\end{firewall}

\section{Append-only record for V82}
\label{sec:V82-record}

V82 was appended to the complete V81 manuscript.  The exact V81
parent had \(57065\) validation lines, \(2027325\) bytes, and
SHA--256 digest
\[
 \texttt{f6225385ee5b1b5c13dfa245df8593126f7ff3ee82baf7531f1302793034b3a8}.
\]
No V81 mathematical content was changed.  The sole final
\verb|\end{document}| is restored below.

% =============================================================================
% END APPENDED VERSION V82 MATERIAL
% =============================================================================

% =============================================================================
% BEGIN APPENDED VERSION V83 MATERIAL
% =============================================================================

\chapter{V83: Heat-Kernel Benchmark, Transition-Scale Obstruction,
and the Wilson Comparison Gate}
\label{ch:V83}

\section{Purpose and change of direction}
\label{sec:V83-purpose}

V82 reduced the balanced Cartan product column to two multiplier
questions: affine accuracy on a growing Pieri segment and scalar
compatibility at the exact mean pair Casimir.  The last alternative
in \Cref{op:V83-mandate} was to compare the Wilson central law with
the heat semigroup if a direct Wilson curvature theorem could not be
proved from the available estimates.

This note carries out that upstream comparison at the level at which
it is presently rigorous.  The heat-kernel multipliers are exactly
exponential in Casimir, so their curvature is explicit.  Unexpectedly,
the curvature leakage vanishes on the balanced thermal window while
the scalar compatibility defect has a nonzero transition-scale
limit.  Thus the heat benchmark does not close the core-only
balanced lift.  It proves a positive obstruction profile.  A
restricted Wilson-to-heat comparison would transfer that obstruction
to the Wilson law.

This is not an obstruction to the full collision mechanism.  In
particular, it does not include the complete same-orientation
external pair bank identified in the V80 companion audit.  The
correct next move is therefore to enlarge the collision space before
concluding nondecay, rather than to continue trying to make the
core-only scalar defect vanish.

\section{Exact heat-semigroup multipliers}
\label{sec:V83-heat-multipliers}

Let \(G=\mathrm{SU}(3)\), and normalize the nonnegative bi-invariant
Laplacian so that
\[
 -\Delta_G\chi_R=c_R\chi_R,
 \qquad c_R=C_2(R).
\]
Let \(\mathsf h_\tau(g)\,\dd g\) be the heat-kernel probability
measure at time \(\tau\ge0\), and write
\[
 h_{\tau,R}:=e^{-\tau c_R}.
\]

\begin{lemma}[Exact central Fourier scalar]
\label{lem:V83-heat-Fourier-scalar}
For every irreducible unitary representation \(R\),
\begin{equation}
 \frac1{d_R}\int_G\chi_R(g)\mathsf h_\tau(g)\,\dd g
 =
 h_{\tau,R}
 =
 e^{-\tau c_R}.
 \label{eq:V83-heat-Fourier-scalar}
\end{equation}
Equivalently,
\[
 \int_GD^R(g)\mathsf h_\tau(g)\,\dd g
 =
 e^{-\tau c_R}I_{H_R}.
\]
\end{lemma}

\begin{proof}
The heat semigroup acts on every matrix coefficient of \(R\) by
\(e^{-\tau c_R}\), since that Peter--Weyl isotypic space is an
eigenspace of \(-\Delta_G\) with eigenvalue \(c_R\).  Evaluating the
semigroup at the identity gives the matrix formula.  Taking its trace
and dividing by \(d_R\) gives \eqref{eq:V83-heat-Fourier-scalar}.
\end{proof}

\begin{remark}[Fundamental calibration]
\label{rem:V83-fundamental-calibration}
The manuscript defines
\[
 s_\alpha=-\frac{\log q_{\alpha,F}}{c_F}.
\]
Consequently the heat benchmark with \(\tau=s_\alpha\) agrees
exactly with the Wilson multiplier on the fundamental
representation:
\[
 h_{s_\alpha,F}=q_{\alpha,F}.
\]
This calibration alone gives no growing-representation comparison.
\end{remark}

\section{Heat curvature on the balanced Pieri segment}
\label{sec:V83-heat-curvature}

Retain the balanced lift of V81--V82:
\[
 A=B=R_m,\qquad C=Q_{2m},\qquad
 S=R_{2m},\qquad T_k=(k,k).
\]
Set
\begin{align}
 c_m&:=c_{R_m}=\frac{m^2+3m}{3},
 \label{eq:V83-cm}\\
 c_{2m}&:=c_{R_{2m}}=\frac{4m^2+6m}{3},
 \label{eq:V83-c2m}\\
 c_{T_k}&:=c_{(k,k)}=k^2+2k,
 \label{eq:V83-cTk}\\
 \bar c_{m,k}&:=\frac{\Lambda_{m,k}}2
 =c_m+\frac{c_{T_k}}2,
 \label{eq:V83-cbar}\\
 L_{m,k}&:=c_0-c_k=k(m+k+2).
 \label{eq:V83-Lmk}
\end{align}
By \eqref{eq:V82-Casimir-nodes}, the balanced \(AC\) pair-Casimir
nodes fill the interval
\[
 c_k=c_m\le c_r\le c_0=c_m+L_{m,k}.
\]
Moreover \(\bar c_{m,k}\) lies in that interval.

Define the tangent affine function
\begin{equation}
 \ell_{\tau;m,k}(c)
 :=
 e^{-\tau\bar c_{m,k}}
 \bigl[1-\tau(c-\bar c_{m,k})\bigr].
 \label{eq:V83-heat-tangent}
\end{equation}

\begin{proposition}[Uniform heat affine error]
\label{prop:V83-heat-affine-error}
On every balanced Pieri node,
\begin{equation}
 \max_{0\le r\le k}
 \left|
 e^{-\tau c_r}-\ell_{\tau;m,k}(c_r)
 \right|
 \le
 \mathcal E_{\tau;m,k},
 \qquad
 \mathcal E_{\tau;m,k}
 :=
 \frac{\tau^2}{2}e^{-\tau c_m}L_{m,k}^2.
 \label{eq:V83-heat-affine-error}
\end{equation}
Consequently the heat even-ancestry leakage satisfies
\begin{equation}
 E_{m,k}^{\mathrm{even},\mathsf h}
 \le
 \mathcal E_{\tau;m,k}^2.
 \label{eq:V83-heat-even-leakage}
\end{equation}
\end{proposition}

\begin{proof}
For \(f(c)=e^{-\tau c}\), Taylor's formula about
\(\bar c_{m,k}\) gives
\[
 |f(c)-f(\bar c)-f'(\bar c)(c-\bar c)|
 \le
 \frac12
 \sup_{\zeta\in[c_m,c_m+L_{m,k}]}|f''(\zeta)|
 |c-\bar c|^2.
\]
Here \(f''(\zeta)=\tau^2e^{-\tau\zeta}\), the supremum is
\(\tau^2e^{-\tau c_m}\), and both \(c\) and \(\bar c\) belong to an
interval of length \(L_{m,k}\).  This proves
\eqref{eq:V83-heat-affine-error}.  Apply
\eqref{eq:V82-even-leakage-bound} with the heat multipliers and the
affine function \eqref{eq:V83-heat-tangent} to obtain
\eqref{eq:V83-heat-even-leakage}.
\end{proof}

\begin{corollary}[Thermal-window disappearance of heat curvature]
\label{cor:V83-heat-curvature-disappears}
Suppose
\begin{equation}
 \tau_j\downarrow0,\qquad
 m_j\longrightarrow\infty,\qquad
 \tau_jm_j^2\longrightarrow u\in(0,\infty),
 \label{eq:V83-transition-scaling}
\end{equation}
and fix \(0<A_0<B_0<\infty\).  Uniformly for integers
\[
 A_0\sqrt{2m_j}\le k\le B_0\sqrt{2m_j},
\]
one has
\begin{equation}
 \mathcal E_{\tau_j;m_j,k}\longrightarrow0,
 \qquad
 E_{m_j,k}^{\mathrm{even},\mathsf h}\longrightarrow0.
 \label{eq:V83-curvature-vanishes}
\end{equation}
\end{corollary}

\begin{proof}
On the stated window,
\[
 L_{m,k}=k(m+k+2)=O(m^{3/2}).
\]
Therefore
\[
 \mathcal E_{\tau;m,k}
 \le C\tau^2m^3
 =
 C\frac{(\tau m^2)^2}{m},
\]
uniformly on the window.  This tends to zero under
\eqref{eq:V83-transition-scaling}.  The leakage conclusion follows
from \eqref{eq:V83-heat-even-leakage}.
\end{proof}

\section{The exact heat scalar defect}
\label{sec:V83-heat-scalar-defect}

For the heat multipliers, the balanced scalar coefficient
\eqref{eq:V79-scalar-part} is
\begin{equation}
 \vartheta_{\tau;m,k}^{\mathsf h}
 =
 e^{-\tau c_{T_k}}
 -e^{-2\tau c_{2m}}
 +2e^{-\tau(2c_m+c_{2m})}.
 \label{eq:V83-heat-vartheta}
\end{equation}
Since
\(\ell_{\tau;m,k}(\bar c_{m,k})
=e^{-\tau\bar c_{m,k}}\), define
\begin{align}
 \mathcal D_{\tau;m,k}
 &:=
 \vartheta_{\tau;m,k}^{\mathsf h}
 -2e^{-\tau c_m}e^{-\tau\bar c_{m,k}}
 \notag\\
 &=
 e^{-\tau c_{T_k}}
 -e^{-2\tau c_{2m}}
 +2e^{-\tau(2c_m+c_{2m})}
 -2e^{-\tau(2c_m+c_{T_k}/2)}.
 \label{eq:V83-heat-scalar-defect}
\end{align}

\begin{lemma}[Second-order cancellation and two elementary bounds]
\label{lem:V83-scalar-cancellation}
The four coefficients and four exponents in
\eqref{eq:V83-heat-scalar-defect} obey
\begin{align}
 1-1+2-2&=0,
 \label{eq:V83-zero-order-cancellation}\\
 c_{T_k}-2c_{2m}
 +2(2c_m+c_{2m})
 -2(2c_m+c_{T_k}/2)&=0.
 \label{eq:V83-first-order-cancellation}
\end{align}
Hence
\begin{align}
 |\mathcal D_{\tau;m,k}|
 &\le
 \frac{\tau^2}{2}
 \left[
 c_{T_k}^2
 +(2c_{2m})^2
 +2(2c_m+c_{2m})^2
 +2(2c_m+c_{T_k}/2)^2
 \right],
 \label{eq:V83-defect-Taylor-bound}\\
 |\mathcal D_{\tau;m,k}|
 &\le
 e^{-\tau c_{T_k}}
 +e^{-2\tau c_{2m}}
 +2e^{-\tau(2c_m+c_{2m})}
 +2e^{-\tau(2c_m+c_{T_k}/2)}.
 \label{eq:V83-defect-tail-bound}
\end{align}
\end{lemma}

\begin{proof}
The identities
\eqref{eq:V83-zero-order-cancellation} and
\eqref{eq:V83-first-order-cancellation} are direct.  Write
\[
 e^{-x}=1-x+\mathfrak r(x),
 \qquad |\mathfrak r(x)|\le\frac{x^2}{2},
 \qquad x\ge0,
\]
in each term of \eqref{eq:V83-heat-scalar-defect}.  The constant and
linear contributions vanish by the two identities, and the triangle
inequality gives \eqref{eq:V83-defect-Taylor-bound}.  Applying the
triangle inequality directly gives \eqref{eq:V83-defect-tail-bound}.
\end{proof}

\begin{theorem}[Positive transition profile]
\label{thm:V83-positive-transition-profile}
Under \eqref{eq:V83-transition-scaling}, uniformly for
\[
 A_0\sqrt{2m_j}\le k\le B_0\sqrt{2m_j},
\]
one has
\begin{equation}
 \mathcal D_{\tau_j;m_j,k}
 \longrightarrow
 \mathfrak d(u)
 :=
 1-e^{-8u/3}+2e^{-2u}-2e^{-2u/3}.
 \label{eq:V83-transition-profile}
\end{equation}
The profile factors exactly as
\begin{equation}
 \boxed{
 \mathfrak d(u)
 =
 (1-e^{-2u/3})^3(1+e^{-2u/3})>0
 \quad (u>0).}
 \label{eq:V83-transition-profile-factor}
\end{equation}
\end{theorem}

\begin{proof}
The thermal restriction on \(k\) and
\eqref{eq:V83-transition-scaling} imply
\[
 \tau_jc_{T_k}
 =
 O(\tau_jm_j)
 =
 O((\tau_jm_j^2)/m_j)
 \longrightarrow0
\]
uniformly in \(k\).  Equations
\eqref{eq:V83-cm}--\eqref{eq:V83-c2m} give
\[
 \tau_jc_{m_j}\longrightarrow\frac u3,
 \qquad
 \tau_jc_{2m_j}\longrightarrow\frac{4u}{3}.
\]
Substitution in \eqref{eq:V83-heat-scalar-defect} proves
\eqref{eq:V83-transition-profile}.  With \(y=e^{-2u/3}\),
\[
 1-e^{-8u/3}+2e^{-2u}-2e^{-2u/3}
 =
 1-y^4+2y^3-2y
 =
 (1-y)^3(1+y).
\]
For \(u>0\), \(0<y<1\), proving strict positivity.
\end{proof}

\section{From the affine mean to the exact heat collision vector}
\label{sec:V83-exact-heat-vector}

Let \(m_{\tau;m,k}^{\mathsf h}\) denote the exact heat version of the
balanced pair mean \(m_*\) in V80--V82:
\[
 m_{\tau;m,k}^{\mathsf h}
 =
 \sum_{r=0}^ke^{-\tau c_r}\pi_{AC}(U_r).
\]
The Racah probabilities are properties of the representation path;
they do not depend on the chosen central probability measure.

\begin{lemma}[Exact heat mean error]
\label{lem:V83-heat-mean-error}
One has
\begin{equation}
 \left|
 m_{\tau;m,k}^{\mathsf h}
 -e^{-\tau\bar c_{m,k}}
 \right|
 \le
 \mathcal E_{\tau;m,k}.
 \label{eq:V83-heat-mean-error}
\end{equation}
\end{lemma}

\begin{proof}
Use the probability resolution
\(\sum_r\pi_{AC}(U_r)=1\), the exact Casimir mean
\(\sum_rc_r\pi_{AC}(U_r)=\bar c_{m,k}\) from
\eqref{eq:V82-mean-pair-Casimir}, and the affine function
\eqref{eq:V83-heat-tangent}.  Its average is
\[
 \sum_r\ell_{\tau;m,k}(c_r)\pi_{AC}(U_r)
 =
 \ell_{\tau;m,k}(\bar c_{m,k})
 =
 e^{-\tau\bar c_{m,k}}.
\]
Average \eqref{eq:V83-heat-affine-error}.
\end{proof}

Define the exact heat scalar component
\[
 \Delta_{\tau;m,k}^{\mathsf h}
 :=
 \vartheta_{\tau;m,k}^{\mathsf h}
 -2e^{-\tau c_m}m_{\tau;m,k}^{\mathsf h}.
\]

\begin{corollary}[Heat scalar component follows the positive profile]
\label{cor:V83-heat-scalar-profile}
For all \(\tau,m,k\),
\begin{equation}
 \left|
 \Delta_{\tau;m,k}^{\mathsf h}
 -\mathcal D_{\tau;m,k}
 \right|
 \le2\mathcal E_{\tau;m,k}.
 \label{eq:V83-heat-scalar-profile-error}
\end{equation}
Under \eqref{eq:V83-transition-scaling}, uniformly on the thermal
\(k\)-window,
\begin{equation}
 \Delta_{\tau_j;m_j,k}^{\mathsf h}
 \longrightarrow\mathfrak d(u)>0.
 \label{eq:V83-heat-scalar-profile-limit}
\end{equation}
In particular, the exact balanced heat collision vector obeys
\begin{equation}
 \inf_{A_0\sqrt{2m_j}\le k\le B_0\sqrt{2m_j}}
 \|\boldsymbol\kappa_{m_j,m_j;2m_j,k}^{\mathsf h}\|
 \longrightarrow\mathfrak d(u).
 \label{eq:V83-heat-vector-nondecay}
\end{equation}
\end{corollary}

\begin{proof}
The first assertion follows from
\eqref{eq:V83-heat-mean-error} and
\(e^{-\tau c_m}\le1\).  Combine it with
\Cref{cor:V83-heat-curvature-disappears,
thm:V83-positive-transition-profile} for the scalar limit.

The balanced parity identity
\eqref{eq:V81-balanced-parity-cancellation}, applied to the heat
multipliers, gives
\[
 \|\boldsymbol\kappa^{\mathsf h}\|^2
 =
 |\Delta^{\mathsf h}|^2
4e^{-2\tau c_m}E_{m,k}^{\mathrm{even},\mathsf h}.
\]
The first term converges uniformly to \(\mathfrak d(u)^2\), and the
second converges uniformly to zero by
\eqref{eq:V83-curvature-vanishes}.  Taking square roots proves
\eqref{eq:V83-heat-vector-nondecay}.
\end{proof}

\begin{assessment}[The heat model rules out core-only decay]
\label{ass:V83-heat-core-obstruction}
The exact heat semigroup, despite perfect Casimir regularity, has
nondecaying balanced product-column energy at
\(\tau m^2\to u>0\) and \(k\asymp\sqrt m\).  Indeed
\eqref{eq:V79-thermal-dimension-mass} supplies a positive amount of
dimension mass on every fixed thermal \(k\)-window, while
\eqref{eq:V83-heat-vector-nondecay} supplies a positive vector norm.
Thus improving the V82 divided-difference estimate cannot by itself
close this core-only channel.
\end{assessment}

\section{Restricted Wilson-to-heat comparison}
\label{sec:V83-Wilson-heat-comparison}

For fixed \(\alpha,m,k,\tau\), introduce the finite representation
bank
\begin{equation}
 \mathcal B_{m,k}
 :=
 \{R_m,R_{2m},Q_{2m},T_k,U_0,\ldots,U_k\}
 \label{eq:V83-comparison-bank}
\end{equation}
and its multiplier discrepancy
\begin{equation}
 \delta_{\alpha;m,k}(\tau)
 :=
 \max_{R\in\mathcal B_{m,k}}
 \left|q_{\alpha,R}-e^{-\tau c_R}\right|.
 \label{eq:V83-restricted-discrepancy}
\end{equation}

\begin{lemma}[Normalized central multipliers are contractions]
\label{lem:V83-multiplier-contraction}
For every central probability measure and every irreducible unitary
representation,
\begin{equation}
 |q_R|\le1.
 \label{eq:V83-multiplier-contraction}
\end{equation}
\end{lemma}

\begin{proof}
The averaged representation matrix
\(\int D^R(g)\,\dd\mu(g)=q_RI\) is an average of unitary
contractions.  Its operator norm is at most one, proving the claim.
\end{proof}

Let \(\vartheta_{\alpha;m,k}\), \(m_{\alpha;m,k}\), and
\(\Delta_{\alpha;m,k}\) denote the Wilson versions of the balanced
scalar coefficient, pair mean, and scalar collision component.

\begin{proposition}[Finite-bank transfer estimate]
\label{prop:V83-finite-bank-transfer}
With \(\delta=\delta_{\alpha;m,k}(\tau)\),
\begin{align}
 \left|
 \vartheta_{\alpha;m,k}
 -\vartheta_{\tau;m,k}^{\mathsf h}
 \right|
 &\le9\delta,
 \label{eq:V83-vartheta-transfer}\\
 \left|
 m_{\alpha;m,k}
 -m_{\tau;m,k}^{\mathsf h}
 \right|
 &\le\delta,
 \label{eq:V83-mean-transfer}\\
 \left|
 \Delta_{\alpha;m,k}
 -\Delta_{\tau;m,k}^{\mathsf h}
 \right|
 &\le13\delta.
 \label{eq:V83-Delta-transfer}
\end{align}
Moreover, using the heat tangent
\eqref{eq:V83-heat-tangent} in the V82 master estimate gives
\begin{equation}
 \|\boldsymbol\kappa_{m,m;2m,k}^{\alpha}\|^2
 \le
 \left(
 |\mathcal D_{\tau;m,k}|
 +2\mathcal E_{\tau;m,k}
 +13\delta
 \right)^2
4\left(\mathcal E_{\tau;m,k}+\delta\right)^2.
 \label{eq:V83-Wilson-heat-upper-transfer}
\end{equation}
\end{proposition}

\begin{proof}
Write \(q_m=q_{\alpha,R_m}\),
\(q_n=q_{\alpha,R_{2m}}=q_{\alpha,Q_{2m}}\), and use the analogous
notation \(h_m,h_n,h_T\) for the heat multipliers.  Then
\[
 \vartheta_\alpha=q_T-q_n^2+2q_m^2q_n,
 \qquad
 \vartheta_{\mathsf h}=h_T-h_n^2+2h_m^2h_n.
\]
By \eqref{eq:V83-multiplier-contraction},
\[
 |q_n^2-h_n^2|\le2\delta,
 \qquad
 |q_m^2q_n-h_m^2h_n|\le3\delta.
\]
Together with \(|q_T-h_T|\le\delta\), this proves
\eqref{eq:V83-vartheta-transfer}.

The same Racah probabilities occur in both pair means, so
\[
 |m_\alpha-m_{\mathsf h}|
 \le
 \sum_r\pi_{AC}(U_r)
 |q_{\alpha,U_r}-h_{\tau,U_r}|
 \le\delta.
\]
Also \(|m_\alpha|,|m_{\mathsf h}|\le1\), whence
\[
 2|q_mm_\alpha-h_mm_{\mathsf h}|
 \le4\delta.
\]
Combine this with \eqref{eq:V83-vartheta-transfer} to prove
\eqref{eq:V83-Delta-transfer}.

Finally, on every Pieri node the Wilson error relative to the heat
tangent is at most
\(\mathcal E_{\tau;m,k}+\delta\).  The scalar term relative to the
same tangent is bounded by
\[
 |\mathcal D_{\tau;m,k}|+11\delta,
\]
because
\(\ell_{\tau;m,k}(\bar c)=e^{-\tau\bar c}\le1\).
Insert these estimates and \(|q_m|\le1\) into
\eqref{eq:V82-affine-master-bound}.  Combining the additional
\(2(\mathcal E+\delta)\) with the scalar bound gives
\eqref{eq:V83-Wilson-heat-upper-transfer}.
\end{proof}

\begin{theorem}[Conditional transfer of the core obstruction]
\label{thm:V83-conditional-obstruction-transfer}
Let \(\alpha_j\) be such that \(s_{\alpha_j}\downarrow0\), and choose
\(m_j\) so that
\[
 s_{\alpha_j}m_j^2\longrightarrow u\in(0,\infty).
\]
If, for some \(0<A_0<B_0<\infty\),
\begin{equation}
 \sup_{A_0\sqrt{2m_j}\le k\le B_0\sqrt{2m_j}}
 \delta_{\alpha_j;m_j,k}(s_{\alpha_j})
 \longrightarrow0,
 \label{eq:V83-restricted-comparison-hypothesis}
\end{equation}
then
\begin{equation}
 \inf_{A_0\sqrt{2m_j}\le k\le B_0\sqrt{2m_j}}
 \|\boldsymbol\kappa_{m_j,m_j;2m_j,k}^{\alpha_j}\|
 \longrightarrow\mathfrak d(u)>0.
 \label{eq:V83-Wilson-core-nondecay}
\end{equation}
Consequently the core-only balanced Cartan product column satisfies
the nondecay hypothesis \eqref{eq:V79-nondecay-hypothesis}.
\end{theorem}

\begin{proof}
By \eqref{eq:V83-Delta-transfer},
\eqref{eq:V83-restricted-comparison-hypothesis}, and
\eqref{eq:V83-heat-scalar-profile-limit}, the Wilson scalar
component \(\Delta_{\alpha_j;m_j,k}\) converges uniformly to
\(\mathfrak d(u)\).  The exact Wilson parity identity
\eqref{eq:V81-balanced-parity-cancellation} implies
\[
 \|\boldsymbol\kappa^\alpha\|
 \ge|\Delta_\alpha|.
\]
This proves \eqref{eq:V83-Wilson-core-nondecay}; the final statement
is precisely \Cref{prop:V79-balanced-obstruction} with
\(a_j=b_j=m_j\) and \(n_j=2m_j\).
\end{proof}

\section{A stronger measure-level comparison}
\label{sec:V83-measure-comparison}

Let \(\nu_\alpha\) be the one-link Wilson probability measure and
let \(\mathsf h_\tau(g)\,\dd g\) be the heat-kernel measure.  Use the
total-variation norm convention
\[
 \|\sigma\|_{\mathrm{TV}}
 :=
 \sup_{\|f\|_\infty\le1}
 \left|\int_Gf\,\dd\sigma\right|.
\]

\begin{proposition}[Total variation controls every representation]
\label{prop:V83-TV-controls-multipliers}
For every irreducible \(R\),
\begin{equation}
 \left|q_{\alpha,R}-e^{-\tau c_R}\right|
 \le
 \left\|
 \nu_\alpha-\mathsf h_\tau(g)\,\dd g
 \right\|_{\mathrm{TV}}.
 \label{eq:V83-TV-multiplier-bound}
\end{equation}
In particular,
\begin{equation}
 \delta_{\alpha;m,k}(\tau)
 \le
 \left\|
 \nu_\alpha-\mathsf h_\tau(g)\,\dd g
 \right\|_{\mathrm{TV}}.
 \label{eq:V83-TV-restricted-bound}
\end{equation}
\end{proposition}

\begin{proof}
For a unitary representation,
\[
 \left|\frac{\chi_R(g)}{d_R}\right|\le1.
\]
Use this normalized character as a test function in the definition
of total variation, and apply
\eqref{eq:V83-heat-Fourier-scalar}.  Taking the maximum over
\(\mathcal B_{m,k}\) proves \eqref{eq:V83-TV-restricted-bound}.
\end{proof}

\begin{firewall}[What has and has not been compared]
\label{fire:V83-comparison-firewall}
No estimate in the manuscript currently proves
\[
 \left\|
 \nu_\alpha-\mathsf h_{s_\alpha}(g)\,\dd g
 \right\|_{\mathrm{TV}}\longrightarrow0
\]
or even the weaker restricted comparison
\eqref{eq:V83-restricted-comparison-hypothesis}.  Fixed-\(R\)
asymptotics and pointwise two-sided gaps do not imply uniform
convergence for highest weights of order \(s_\alpha^{-1/2}\).
Therefore \Cref{thm:V83-conditional-obstruction-transfer} is
conditional, not a Wilson theorem.
\end{firewall}

\section{Regime diagnosis and the next object}
\label{sec:V83-regime-diagnosis}

The heat calculation separates the three multiplier regimes without
an assumed Wilson curvature estimate:
\begin{enumerate}[label=\textup{(\roman*)}]
\item If \(\tau m^2\to0\), the Taylor cancellation
      \eqref{eq:V83-first-order-cancellation} makes the scalar defect
      higher order, and
      \eqref{eq:V83-heat-affine-error} controls curvature.
\item If \(\tau m^2\to u\in(0,\infty)\), curvature leakage vanishes
      but the scalar defect tends to the strictly positive profile
      \eqref{eq:V83-transition-profile-factor}.
\item If \(\tau m^2\to\infty\), the external multipliers in
      \eqref{eq:V83-defect-tail-bound} decay, except that the final
      \(T_k\) multiplier must be retained whenever
      \(\tau c_{T_k}\) is not large.
\end{enumerate}

\begin{assessment}[Why the collision bank must now be enlarged]
\label{ass:V83-enlarge-bank}
The V82 affine programme can remove only the non-affine ancestry
leakage.  It cannot remove the positive scalar profile
\(\mathfrak d(u)\), because that profile remains when the heat
curvature error tends to zero.  Repeating stronger curvature
estimates would therefore be circular with respect to the actual
obstruction.

The next analysis must insert the external same-orientation pair
bank before the core \(AB\) regression is judged.  Algebraically, it
must identify the extra signed columns that share the scalar
\(R_{2m}\) ancestry and determine whether their Schur complement
cancels \(\mathfrak d(u)\) while preserving positivity of the full
collision Gram operator.
\end{assessment}

\section{V83 ledger and V84 mandate}
\label{sec:V83-ledger}

\begin{theorem}[Integrated V83 statement]
\label{thm:V83-integrated}
Preserving V1--V82, V83 proves:
\begin{enumerate}[label=\textup{(V83.\arabic*)},leftmargin=5.7em]
\item exact heat-semigroup multipliers
      \eqref{eq:V83-heat-Fourier-scalar};
\item the uniform heat affine and even-leakage estimates
      \eqref{eq:V83-heat-affine-error} and
      \eqref{eq:V83-heat-even-leakage};
\item the exact positive transition profile
      \eqref{eq:V83-transition-profile-factor};
\item nondecay of the exact balanced heat collision vector
      \eqref{eq:V83-heat-vector-nondecay};
\item the finite-bank Wilson-to-heat transfer estimates
      \eqref{eq:V83-vartheta-transfer}--\eqref{eq:V83-Delta-transfer};
\item conditional transfer of the core obstruction to Wilson
      multipliers in
      \eqref{eq:V83-Wilson-core-nondecay}; and
\item the measure-level total-variation bound
      \eqref{eq:V83-TV-multiplier-bound}.
\end{enumerate}
\end{theorem}

\begin{proof}
These are
\Cref{lem:V83-heat-Fourier-scalar,
prop:V83-heat-affine-error,
thm:V83-positive-transition-profile,
cor:V83-heat-scalar-profile,
prop:V83-finite-bank-transfer,
thm:V83-conditional-obstruction-transfer,
prop:V83-TV-controls-multipliers}.
\end{proof}

\begin{openproblem}[V84 mandate: complete external pair bank]
\label{op:V84-mandate}
Return to the collision construction upstream of the core-only
regression and:
\begin{enumerate}
\item reconstruct the complete same-orientation external pair bank
      associated with the balanced \(R_m,R_m,Q_{2m}\) lift;
\item write the full signed collision Gram matrix before eliminating
      any external pair column;
\item isolate the common scalar \(R_{2m}\) ancestry and its exact
      Schur complement;
\item test whether the added columns cancel the heat transition
      profile \(\mathfrak d(u)\);
\item preserve the full positive square while taking that Schur
      complement;
\item only after the complete-bank calculation, decide whether a
      Wilson-to-heat estimate such as
      \eqref{eq:V83-restricted-comparison-hypothesis} is helpful or
      obstructive.
\end{enumerate}
The next compulsory SM/NS companion audit remains V85.
\end{openproblem}

\section{Firewall after V83}
\label{sec:V83-firewall}

\begin{firewall}[Current claim boundary]
\label{fire:V83-final}
V83 proves a heat-semigroup obstruction to decay of the core-only
balanced Cartan product column and gives an exact conditional route
for transferring it to Wilson multipliers.  It does not prove the
required Wilson-to-heat comparison, show that the complete collision
bank is nondecaying, rule out cancellation by external pair columns,
prove JREC, close simultaneous sideways endpoints, aggregate
punctures or multiple collisions, construct the continuum
Yang--Mills theory, or prove a positive mass gap.
\end{firewall}

\section{Append-only record for V83}
\label{sec:V83-record}

V83 was appended to the complete V82 manuscript.  The exact V82
parent had \(57538\) validation lines, \(2040892\) bytes, and
SHA--256 digest
\[
 \texttt{c86ed98fa5ac62fc6dd898300ee0a623d696a727c7c608cd6984f4444e53f312}.
\]
No V82 mathematical content was changed.  The sole final
\verb|\end{document}| is restored below.

% =============================================================================
% END APPENDED VERSION V83 MATERIAL
% =============================================================================

% =============================================================================
% BEGIN APPENDED VERSION V84 MATERIAL
% =============================================================================

\chapter{V84: Exact External-Pair Factorization and the
Pair-Response Debit}
\label{ch:V84}

\section{Purpose and correction of the proposed Schur route}
\label{sec:V84-purpose}

V83 proved that the core-only balanced Cartan column has a positive
transition profile in the heat benchmark and conditionally for the
Wilson law.  Its V84 mandate proposed looking for additive external
columns whose Schur complement might cancel that profile.

Returning to the exact V67 and V72 collision carrier shows that this
is not the algebra of the already-defined external \(oa\) and \(bc\)
pair banks.  Those banks are disjoint from the ternary core and enter
the positive carrier as tensor factors.  After their resolved
partial traces, they produce the scalar \(P_{ab}\).  They can make an
oriented collision estimate small by supplying a multiplicative pair
debit, but they cannot cancel the core scalar defect additively.

This note proves that distinction exactly.  It also converts the V83
square obstruction into a lower bound for the absolute-core first
moment used by the collision capacity.  The remaining same-orientation
question is thereby reduced to an explicit estimate of \(P_{ab}\) in
the simultaneous-sideways residual left open by V58--V60.

\section{The complete positive pair--core carrier}
\label{sec:V84-positive-carrier}

Recall the positive weighted pair-bank operators
\begin{equation}
 Q_{oa}=\sum_\lambda a_\lambda P_{oa,\lambda},
 \qquad
 Q_{bc}=\sum_\mu b_\mu P_{bc,\mu},
 \label{eq:V84-pair-operators}
\end{equation}
and the weighted core operator
\[
 \mathbf K^\sim
 =
 \bigoplus_T I_{H_T}\otimes\widetilde K_T,
 \qquad
 \widetilde K_T=(1+s_\alpha\Xi(T))K_T.
\]
Before suspended same-label puncture banks are inserted, the exact
positive carrier is
\begin{equation}
 Q_{ab}^{\mathrm{abs}}
 =
 Q_{oa}\otimes Q_{bc}\otimes|\mathbf K^\sim|.
 \label{eq:V84-exact-positive-carrier}
\end{equation}
This is \eqref{eq:V72-three-bank-product}, recalled here because it
decides the algebraic alternative raised in V83.

Write
\begin{align}
 p_{oa\to a}
 &:=
 \sum_\lambda a_\lambda x_\lambda^a,
 \label{eq:V84-left-pair-response}\\
 p_{bc\to b}
 &:=
 \sum_\mu b_\mu y_\mu^b,
 \label{eq:V84-right-pair-response}\\
 P_{ab}
 &:=
 p_{oa\to a}p_{bc\to b}.
 \label{eq:V84-Pab-factor}
\end{align}
The pair partial traces are
\begin{equation}
 \operatorname{Tr}_oQ_{oa}=p_{oa\to a}I_a,
 \qquad
 \operatorname{Tr}_cQ_{bc}=p_{bc\to b}I_b.
 \label{eq:V84-pair-Schur-traces}
\end{equation}

\begin{lemma}[Dimension formula for a resolved pair copy]
\label{lem:V84-pair-copy-dimension}
Let \(P_{XY;S,\mu}\) be the orthogonal projector onto one
multiplicity-resolved copy of an irreducible \(S\) in \(X\otimes Y\).
Then
\begin{equation}
 \operatorname{Tr}_X P_{XY;S,\mu}
 =
 \frac{d_S}{d_Y}I_Y,
 \qquad
 \operatorname{Tr}_Y P_{XY;S,\mu}
 =
 \frac{d_S}{d_X}I_X.
 \label{eq:V84-pair-copy-partial-trace}
\end{equation}
Consequently, with \(\lambda=(S,\mu)\),
\begin{equation}
 p_{oa\to a}
 =
 \frac1{d_a}\sum_\lambda a_\lambda d_{S_\lambda},
 \qquad
 p_{bc\to b}
 =
 \frac1{d_b}\sum_\mu b_\mu d_{S_\mu}.
 \label{eq:V84-pair-response-dimension-sum}
\end{equation}
Here \(d_a,d_b\) are the dimensions of the relevant irreducible row
carriers on the corresponding pair banks.
\end{lemma}

\begin{proof}
The first partial trace commutes with \(D^Y(g)\), hence is scalar by
Schur's lemma.  Its total trace is the rank \(d_S\) of one resolved
copy, so the scalar is \(d_S/d_Y\).  The second identity is
identical.  Summing with the nonnegative weights in
\eqref{eq:V84-pair-operators} gives
\eqref{eq:V84-pair-response-dimension-sum}.
\end{proof}

\begin{remark}[The response is not a formation probability]
\label{rem:V84-not-formation-probability}
The dimension sum in
\eqref{eq:V84-pair-response-dimension-sum} has one row dimension in
the denominator.  It is the exact one-sided Schur trace appearing in
the coefficient tax.  It is not the tensor-product dimension law
\(d_SN_{XY}^{S}/(d_Xd_Y)\), and replacing it by that probability
would reintroduce the normalization error excluded in V57.
\end{remark}

\section{No additive Schur cancellation}
\label{sec:V84-no-Schur-cancellation}

Define the absolute-core product capacity
\begin{equation}
 \mathfrak C_{ab}^{(1)}
 :=
 \sup_{\substack{\rho_A,\rho_B\ge0\\
 \operatorname{Tr}\rho_A=\operatorname{Tr}\rho_B=1}}
 \operatorname{Tr}\left[
 M_{ab}^{(1)}(\rho_A\otimes\rho_B)
 \right],
 \qquad
 M_{ab}^{(1)}=\operatorname{Tr}_C|\mathbf K^\sim|.
 \label{eq:V84-core-first-capacity}
\end{equation}

\begin{theorem}[Exact external-pair factorization]
\label{thm:V84-external-pair-factorization}
For the \(ab\)-retained row orientation,
\begin{equation}
 \boxed{
 \kappa_{\{a,b\}}^\otimes(Q_{ab}^{\mathrm{abs}})
 =
 P_{ab}\mathfrak C_{ab}^{(1)}.}
 \label{eq:V84-external-pair-factorization}
\end{equation}
In particular, the external pair banks introduce no additive cross
term with the ternary core.
\end{theorem}

\begin{proof}
Take the partial trace of
\eqref{eq:V84-exact-positive-carrier} over every factor except the
retained \(a\) and \(b\) rows.  Tensor-product factorization and
\eqref{eq:V84-pair-Schur-traces} give
\[
 P_{ab}\,
 I_{\mathrm{pair},a}\otimes I_{\mathrm{pair},b}
 \otimes M_{ab}^{(1)}.
\]
The identities on the retained pair carriers make their possible
within-row entanglement irrelevant: the expectation depends only on
the reduced core densities.  Every pair of reduced core densities is
realized by tensoring with fixed pure pair-carrier states.  Taking
the product-state supremum gives
\eqref{eq:V84-external-pair-factorization}.
\end{proof}

\begin{corollary}[Why a Schur complement would change the operator]
\label{cor:V84-no-additive-Schur}
There is no block matrix of the form
\[
 \begin{pmatrix}
 \textup{core}&\textup{pair--core cross}\\
 \textup{cross}^*&\textup{pair}
 \end{pmatrix}
\]
whose Schur complement represents
\eqref{eq:V84-exact-positive-carrier}.  The three factors act on
tensor-product carriers, not on an additive direct-sum decomposition.
Thus no Schur-complement cancellation of the V83 scalar profile is
available from the \(oa\) and \(bc\) banks already present in the
collision definition.
\end{corollary}

\begin{proof}
Equation \eqref{eq:V84-exact-positive-carrier} is a tensor product of
positive operators.  Its retained-row partial trace is the scalar
multiple in \eqref{eq:V84-external-pair-factorization}.  A nonzero
additive pair--core cross block would contradict this exact
factorization.  Introducing one would therefore define a different
carrier rather than refine the existing one.
\end{proof}

\begin{remark}[The same conclusion before the positive carrier]
\label{rem:V84-before-absolute-value}
Before trace-norm duality, the two resolved pair defects and the
ternary cut still occur multiplicatively on disjoint coordinate
banks.  For operators \(D_1,D_2,K\),
\[
 |D_1\otimes D_2\otimes K|
 =
 |D_1|\otimes|D_2|\otimes|K|.
\]
Hence returning only one step before the absolute value does not
create additive cancellation columns.  A genuinely new additive
mechanism would have to come from a regrouping of distinct connected
monomials before the V60 carrier is formed, not from the external
pair factors in \eqref{eq:V84-exact-positive-carrier}.
\end{remark}

\section{A reverse square bridge from the cut cap}
\label{sec:V84-reverse-square}

Recall
\[
 M_{ab}^{(2)}
 :=
 \operatorname{Tr}_C(\mathbf K^\sim)^2,
 \qquad
 \Gamma_{ab}^{\otimes}
 :=
 \sup_{\rho_A,\rho_B}
 \operatorname{Tr}\left[
 M_{ab}^{(2)}(\rho_A\otimes\rho_B)\right].
\]
The V63 weighted-cut estimate gives a numerical constant \(C_K\)
such that
\begin{equation}
 \|\widetilde K_T\|_{\mathrm{op}}
 \le
 \frac{C_K}{1+s_\alpha\Xi(T)}
 \le C_K
 \quad\hbox{for every }T.
 \label{eq:V84-uniform-weighted-cut-cap}
\end{equation}

\begin{theorem}[Reverse square-to-first comparison]
\label{thm:V84-reverse-square-bridge}
One has
\begin{align}
 (\mathbf K^\sim)^2
 &\le C_K|\mathbf K^\sim|,
 \label{eq:V84-square-below-absolute}\\
 M_{ab}^{(2)}
 &\le C_KM_{ab}^{(1)},
 \label{eq:V84-partial-square-below-first}\\
 \boxed{
 \mathfrak C_{ab}^{(1)}
 \ge C_K^{-1}\Gamma_{ab}^{\otimes}.}
 \label{eq:V84-reverse-square-bridge}
\end{align}
\end{theorem}

\begin{proof}
For a self-adjoint matrix \(L\) with
\(\|L\|_{\mathrm{op}}\le C_K\), scalar functional calculus gives
\(L^2\le C_K|L|\).  Apply this on every final block in
\(\mathbf K^\sim\) to obtain
\eqref{eq:V84-square-below-absolute}.  Positivity of the partial
trace gives \eqref{eq:V84-partial-square-below-first}.  Evaluate on
an arbitrary product density and take the supremum to prove
\eqref{eq:V84-reverse-square-bridge}.
\end{proof}

\begin{remark}[Two one-sided bridges]
\label{rem:V84-two-bridges}
V72 proved
\[
 \mathfrak C_{ab}^{(1)}
 \le\sqrt{d_C\Gamma_{ab}^{\otimes}}.
\]
Equation \eqref{eq:V84-reverse-square-bridge} supplies the converse
direction needed for an obstruction test.  Together,
\[
 C_K^{-1}\Gamma_{ab}^{\otimes}
 \le
 \mathfrak C_{ab}^{(1)}
 \le
 \sqrt{d_C\Gamma_{ab}^{\otimes}}.
\]
Neither inequality asserts that the two quantities have the same
dimension scale.
\end{remark}

\section{The V83 transition profile reaches the absolute carrier}
\label{sec:V84-profile-transfer}

Fix \(0<A_0<B_0<\infty\).  For the balanced product vector from
V77--V83 and \(n=2m\), define its thermal square contribution
\begin{equation}
 \mathfrak G_{\alpha;m}^{A_0,B_0}
 :=
 \sum_{A_0\sqrt{2m}\le k\le B_0\sqrt{2m}}
 \frac{(k+1)^3}{d_{2m}}
 [1+s_\alpha\Xi(T_k)]^2
 \|\boldsymbol\kappa_{m,m;2m,k}^{\alpha}\|^2.
 \label{eq:V84-thermal-square-contribution}
\end{equation}
It is one product-vector expectation of a positive output shell in
\(M_{ab}^{(2)}\).

\begin{theorem}[Positive absolute-core transition capacity]
\label{thm:V84-positive-absolute-core}
Assume the hypotheses of
\Cref{thm:V83-conditional-obstruction-transfer}; in particular,
\[
 s_{\alpha_j}m_j^2\longrightarrow u\in(0,\infty)
\]
and the restricted Wilson-to-heat discrepancy tends uniformly to
zero on the thermal \(k\)-window.  Then
\begin{align}
 \liminf_{j\to\infty}
 \mathfrak G_{\alpha_j;m_j}^{A_0,B_0}
 &\ge
 \frac{B_0^4-A_0^4}{2}\mathfrak d(u)^2,
 \label{eq:V84-positive-square-limit}\\
 \liminf_{j\to\infty}
 \Gamma_{ab}^{\otimes}
 &\ge
 \frac{B_0^4-A_0^4}{2}\mathfrak d(u)^2,
 \label{eq:V84-positive-Gamma-limit}\\
 \liminf_{j\to\infty}
 \mathfrak C_{ab}^{(1)}
 &\ge
 \frac{B_0^4-A_0^4}{2C_K}\mathfrak d(u)^2.
 \label{eq:V84-positive-first-limit}
\end{align}
\end{theorem}

\begin{proof}
By \eqref{eq:V83-Wilson-core-nondecay}, the squared vector norm
converges uniformly to \(\mathfrak d(u)^2\) on the thermal window.
The weight \([1+s_\alpha\Xi(T_k)]^2\) is at least one.  The exact
dimension Riemann sum
\eqref{eq:V79-thermal-dimension-mass} therefore gives
\eqref{eq:V84-positive-square-limit}.

The quantity in \eqref{eq:V84-thermal-square-contribution} is the
expectation on one admissible product vector of a positive suboperator
of \(M_{ab}^{(2)}\).  Hence
\(\Gamma_{ab}^{\otimes}\ge
\mathfrak G_{\alpha;m}^{A_0,B_0}\), proving
\eqref{eq:V84-positive-Gamma-limit}.  Apply
\eqref{eq:V84-reverse-square-bridge} for
\eqref{eq:V84-positive-first-limit}.
\end{proof}

\begin{corollary}[The pair factor is the only external debit in this
orientation]
\label{cor:V84-pair-only-debit}
Under the hypotheses of
\Cref{thm:V84-positive-absolute-core}, there is a constant
\(c_{u,A_0,B_0}>0\) such that, for all sufficiently large \(j\),
\begin{equation}
 \boxed{
 \kappa_{\{a,b\}}^\otimes
 (Q_{ab}^{\mathrm{abs}})
 \ge
 c_{u,A_0,B_0}P_{ab}.}
 \label{eq:V84-pair-only-debit}
\end{equation}
\end{corollary}

\begin{proof}
Combine
\eqref{eq:V84-external-pair-factorization} with
\eqref{eq:V84-positive-first-limit}.
\end{proof}

\section{The exact pair-response gate}
\label{sec:V84-pair-response-gate}

For the \(ab\)-path define its marked-tree target
\begin{equation}
 \mathfrak p_{ab}
 :=
 s_\alpha
 \frac{h_{oa}\widehat H_{ab}h_{bc}}
      {\Xi_a\Xi_b}.
 \label{eq:V84-ab-path-target}
\end{equation}

\begin{proposition}[Necessary pair debit for the obstructed
orientation]
\label{prop:V84-necessary-pair-debit}
Along any family satisfying
\Cref{thm:V84-positive-absolute-core}, an estimate
\begin{equation}
 \kappa_{\{a,b\}}^\otimes(Q_{ab}^{\mathrm{abs}})
 \le C\mathfrak p_{ab}
 \label{eq:V84-oriented-target}
\end{equation}
can hold only if
\begin{equation}
 \boxed{
 P_{ab}
 \le
 \frac{C}{c_{u,A_0,B_0}}\mathfrak p_{ab}.}
 \label{eq:V84-necessary-Pab-bound}
\end{equation}
\end{proposition}

\begin{proof}
Compare \eqref{eq:V84-oriented-target} with the lower bound
\eqref{eq:V84-pair-only-debit}.
\end{proof}

\begin{assessment}[What is already closed by V58]
\label{ass:V84-V58-region}
The factors in \(P_{ab}\) are the operator-level one-sided traces of
the same weighted pair defects whose coefficient-level normalized
responses are \(\mathfrak A_{oa}\) and \(\mathfrak A_{bc}\) in V58.
V58 already proves the complete double-pair marked-tree estimate when
the relevant endpoints are thermal or selected-heavy on their
incident collision edges.  Therefore
\eqref{eq:V84-necessary-Pab-bound} is not a new request on that closed
region.  Its unresolved content is precisely the simultaneous-sideways
residual retained in the definition of JREC.
\end{assessment}

\begin{firewall}[Core labels do not determine the pair response]
\label{fire:V84-independent-pair-data}
The balanced core labels
\[
 A=B=R_m,\qquad C=Q_{2m},\qquad S=R_{2m}
\]
belong to the punctured ternary bank.  They do not determine the
irreducible inputs or resolved outputs on the disjoint \(oa\) and
\(bc\) pair banks.  Hence \(P_{ab}\) cannot be computed from \(m\),
\(k\), and the V83 Racah path alone.

For example, if either external pair bank has one trivial input, its
binary covariance defect is identically zero and the corresponding
factor in \eqref{eq:V84-Pab-factor} vanishes.  Other nontrivial pair
data give positive responses.  Thus the V83 core obstruction neither
proves nor disproves the complete collision estimate without the
actual external pair data.
\end{firewall}

\section{The orientation minimum remains essential}
\label{sec:V84-orientation-minimum}

The complete product-state collision tax uses the minimum of the
full-carrier capacity and the two coherent retained-row orientations:
\begin{equation}
 \min\left\{
 \kappa_{\{o,a,b,c\}}^\otimes(Q_{ab}^{\mathrm{abs}}),
 P_{ab}\mathfrak C_{ab}^{(1)},
 P_{ac}\mathfrak C_{ac}^{(1)}
 \right\}.
 \label{eq:V84-complete-orientation-minimum}
\end{equation}

\begin{proposition}[One obstructed orientation is not a complete
counterexample]
\label{prop:V84-one-orientation-not-enough}
The lower bound
\eqref{eq:V84-pair-only-debit} for the \(ab\) entry does not give a
lower bound for \eqref{eq:V84-complete-orientation-minimum}.  To
obstruct the complete product-state collision capacity one would
also need compatible lower bounds for the \(ac\) entry and the
full-carrier entry on the same row-factorized input family.
\end{proposition}

\begin{proof}
For nonnegative numbers \(x,y,z\), a lower bound on \(y\) alone gives
no lower bound on \(\min\{x,y,z\}\).  The three entries in
\eqref{eq:V84-complete-orientation-minimum} are valid taxes for the
same carrier, but their optimizing states and partial traces are
different.  Switching orientations inside a coordinate remains
forbidden; taking their minimum only after each complete orientation
has been evaluated is valid.
\end{proof}

\begin{assessment}[Exact bifurcation after V84]
\label{ass:V84-bifurcation}
The balanced transition family leaves two rigorous routes:
\begin{enumerate}[label=\textup{(\roman*)}]
\item prove the pair debit
      \eqref{eq:V84-necessary-Pab-bound}, or a stronger sufficient
      version, from the actual external pair banks on the
      simultaneous-sideways residual; or
\item if that debit fails, evaluate the complete \(ac\) and
      full-carrier entries on the same family before declaring a
      collision-capacity obstruction.
\end{enumerate}
There is no third route in which the already-factorized external
pair banks additively cancel \(\mathfrak d(u)\).
\end{assessment}

\section{V84 ledger and V85 mandate}
\label{sec:V84-ledger}

\begin{theorem}[Integrated V84 statement]
\label{thm:V84-integrated}
Preserving V1--V83, V84 proves:
\begin{enumerate}[label=\textup{(V84.\arabic*)},leftmargin=5.7em]
\item the resolved-copy dimension formula
      \eqref{eq:V84-pair-copy-partial-trace};
\item exact external-pair factorization
      \eqref{eq:V84-external-pair-factorization};
\item absence of an additive pair--core Schur cancellation;
\item the reverse square bridge
      \eqref{eq:V84-reverse-square-bridge};
\item transfer of the V83 transition profile to a positive
      absolute-core first capacity
      \eqref{eq:V84-positive-first-limit};
\item the necessary external pair debit
      \eqref{eq:V84-necessary-Pab-bound}; and
\item the proof that one obstructed orientation does not decide the
      complete collision minimum.
\end{enumerate}
\end{theorem}

\begin{proof}
These are
\Cref{lem:V84-pair-copy-dimension,
thm:V84-external-pair-factorization,
cor:V84-no-additive-Schur,
thm:V84-reverse-square-bridge,
thm:V84-positive-absolute-core,
prop:V84-necessary-pair-debit,
prop:V84-one-orientation-not-enough}.
\end{proof}

\begin{openproblem}[V85 mandate: audited external-pair response]
\label{op:V85-mandate}
After the compulsory SM/NS companion audit:
\begin{enumerate}
\item return to the exact V52 layer-cake formula for each external
      pair response in
      \eqref{eq:V84-pair-response-dimension-sum};
\item impose the simultaneous-sideways residual conditions left by
      V58--V60, not the already-closed selected-heavy alternatives;
\item keep the two pair responses in the same \(ab\) orientation and
      estimate their product before taking a maximum over core
      ancestry;
\item determine whether the exact level-set dimensions and Wilson
      defect weights prove
      \eqref{eq:V84-necessary-Pab-bound} at the transition scale;
\item if the pair debit fails on an explicit family, carry that same
      family into the \(ac\) and full-carrier entries of
      \eqref{eq:V84-complete-orientation-minimum}.
\end{enumerate}
\end{openproblem}

\section{Firewall after V84}
\label{sec:V84-firewall}

\begin{firewall}[Current claim boundary]
\label{fire:V84-final}
V84 proves that the external collision pair banks multiply rather
than cancel the V83 core profile, and it identifies the exact pair
debit required in the obstructed orientation.  It does not prove
that debit in the simultaneous-sideways residual, analyze the
complete \(ac\) or full-carrier alternatives on the transition
family, prove JREC, close puncture or multi-collision aggregation,
construct the continuum Yang--Mills theory, or prove a positive mass
gap.
\end{firewall}

\section{Append-only record for V84}
\label{sec:V84-record}

V84 was appended to the complete V83 manuscript.  The exact V83
parent had \(58355\) validation lines, \(2064937\) bytes, and
SHA--256 digest
\[
 \texttt{7c1b8e1254ad17a866a243be66539b1793c5402f39b0013f99706a5182e9232f}.
\]
No V83 mathematical content was changed.  The sole final
\verb|\end{document}| is restored below.

% =============================================================================
% END APPENDED VERSION V84 MATERIAL
% =============================================================================

% =============================================================================
% BEGIN APPENDED VERSION V85 MATERIAL
% =============================================================================

\chapter{V85: Companion Audit and the Exact External-Pair
Product Quantile}
\label{ch:V85}

\section{Mandatory V85 companion audit}
\label{sec:V85-audit}

The two companion files were read twice after the concurrent updates.
Their sizes and SHA--256 digests were stable:
\begin{center}
\begin{tabular}{p{0.48\textwidth}r r p{0.25\textwidth}}
\toprule
companion file & physical lines & bytes & SHA--256\\
\midrule
\texttt{Renewal\_SM\_Einstein\_Continuum\_Research\_Notes\_V86.tex}
& \(67643\) & \(2452803\) &
\texttt{ea93b753f2764d169940d2b0af1ad396c8966f9a53a79fd589d0871d1739d35c}\\
\texttt{Renewal\_NS\_Stability\_Research\_Notes\_V31.tex}
& \(23052\) & \(887537\) &
\texttt{f1121b62238ad5491bb7cf03635ea9934942edf573ed8faaa9ee79e1d38a6696}\\
\bottomrule
\end{tabular}
\end{center}
Neither companion file was modified.

The relevant SM advances since the V80 audit are structural.  SM V84
distinguishes an unbounded chosen section from a genuine
orbit-invariant quotient obstruction.  SM V85 then shows that a
generic trace obstruction can disappear on the exact special Hessian
image.  SM V86 removes a divergent corner coordinate by an exact
residual orbit vector, while warning that the resulting metric and
conormal blocks must still be estimated together.

For Yang--Mills, V83 is a genuine obstruction for the selected
core product column, and V84 proves that the already-factorized
external pair banks cannot cancel it additively.  The audit nevertheless
forbids promoting that statement to a complete collision obstruction
before the exact capped product-level capacity and both coherent
orientations have been evaluated.  A separated first-moment carrier is
a valid sufficient route, but it is not the only route to the original
joint capacity.

The NS file is unchanged.  Its terminal V31 theorem says that a flat
probability mark is exact, but removing the mark restores precisely
the missing dyadic first moment.  The corresponding proof-order lesson
here is to retain the product threshold and scalar cap through level
integration.  Separate pair first moments can discard the quantile
information needed at the collision scale.

\begin{firewall}[Companion separation]
\label{fire:V85-audit-separation}
No Einstein quotient estimate, Hessian trace theorem, or
Navier--Stokes formation inequality is imported into Yang--Mills.
Only the order of exact assembly is transferred.  Every estimate
below is proved directly from the existing Yang--Mills pair and core
level sets.
\end{firewall}

\section{One-sided layer functions for the external pairs}
\label{sec:V85-one-sided-layers}

For the \(oa\) pair bank in the \(ab\) orientation, define
\begin{equation}
 D_{oa}^a(t)
 :=
 \sum_{\{\lambda:a_\lambda>t\}}x_\lambda^a,
 \qquad 0\le t\le A_*,
 \label{eq:V85-left-layer}
\end{equation}
and for the \(bc\) bank define
\begin{equation}
 D_{bc}^b(t)
 :=
 \sum_{\{\mu:b_\mu>t\}}y_\mu^b,
 \qquad 0\le t\le A_*.
 \label{eq:V85-right-layer}
\end{equation}
The universal cap \(A_*\) is the one in
\eqref{eq:V52-defect-weight-cap}.  Multiplicity-resolved copies are
summed separately.

\begin{proposition}[Exact one-sided layer identities]
\label{prop:V85-one-sided-layer-identities}
The pair responses in V84 satisfy
\begin{align}
 p_{oa\to a}
 &=
 \int_0^{A_*}D_{oa}^a(t)\,\dd t,
 \label{eq:V85-left-layer-integral}\\
 p_{bc\to b}
 &=
 \int_0^{A_*}D_{bc}^b(t)\,\dd t.
 \label{eq:V85-right-layer-integral}
\end{align}
If \(\lambda=(S,\gamma)\), then
\begin{equation}
 D_{oa}^a(t)
 =
 \frac1{d_a}
 \sum_{\{\lambda:a_\lambda>t\}}d_{S_\lambda},
 \label{eq:V85-left-dimension-layer}
\end{equation}
and analogously on \(bc\).
\end{proposition}

\begin{proof}
For every nonnegative \(a_\lambda\le A_*\),
\[
 a_\lambda
 =
 \int_0^{A_*}\boldsymbol1_{\{a_\lambda>t\}}\,\dd t.
\]
Insert this in
\eqref{eq:V84-left-pair-response} and use Tonelli.  This gives
\eqref{eq:V85-left-layer-integral}; the other identity is identical.
Equation \eqref{eq:V85-left-dimension-layer} follows from the
resolved-copy trace formula
\eqref{eq:V84-pair-copy-partial-trace}.
\end{proof}

\begin{remark}[Relation to the V52 capacity envelope]
\label{rem:V85-relation-V52}
For pair inputs \(X,Y\), V52 uses
\[
 \int_0^{A_*}
 \min\left\{
 \frac{D_{\rm dim}(t)}{d_X},
 \frac{D_{\rm dim}(t)}{d_Y},1
 \right\}\,\dd t.
\]
The one-sided trace in
\eqref{eq:V85-left-layer-integral} retains one specified denominator
and has no scalar cap.  Thus the V52 minimum is no larger than either
one-sided first moment.  It cannot be reversed without information
on the actual level-set dimensions.
\end{remark}

\section{The joint pair-product layer}
\label{sec:V85-product-layer}

Define the same-orientation external-pair product layer
\begin{equation}
 G_{ab}(r)
 :=
 \sum_{\{(\lambda,\mu):a_\lambda b_\mu>r\}}
 x_\lambda^ay_\mu^b,
 \qquad 0\le r\le A_*^2.
 \label{eq:V85-pair-product-layer}
\end{equation}
It is nonincreasing and contains the two external pair banks before
their level integrations have been separated.

\begin{theorem}[Exact product-layer first moment]
\label{thm:V85-product-layer-first-moment}
One has
\begin{equation}
 \boxed{
 \int_0^{A_*^2}G_{ab}(r)\,\dd r
 =
 P_{ab}.}
 \label{eq:V85-product-layer-first-moment}
\end{equation}
Equivalently,
\begin{equation}
 P_{ab}
 =
 \int_0^{A_*}\int_0^{A_*}
 D_{oa}^a(t)D_{bc}^b(u)\,\dd t\,\dd u.
 \label{eq:V85-separated-double-layer}
\end{equation}
\end{theorem}

\begin{proof}
Tonelli and
\[
 a_\lambda b_\mu
 =
 \int_0^{A_*^2}
 \boldsymbol1_{\{a_\lambda b_\mu>r\}}\,\dd r
\]
give
\[
 \int_0^{A_*^2}G_{ab}(r)\,\dd r
 =
 \sum_{\lambda,\mu}
 a_\lambda b_\mu x_\lambda^ay_\mu^b
 =
 p_{oa\to a}p_{bc\to b}.
\]
This is \eqref{eq:V85-product-layer-first-moment}.
Equation \eqref{eq:V85-separated-double-layer} follows from
\eqref{eq:V85-left-layer-integral}--%
\eqref{eq:V85-right-layer-integral}.
\end{proof}

\begin{assessment}[What the separated carrier forgets]
\label{ass:V85-separated-forgets}
The number \(P_{ab}\) is the full first moment of the product layer.
The original JREC capacity does not use that first moment without a
cap.  At every level it first compares the active one-sided Schur mass
with one, and only then integrates.  Therefore two pair families can
have the same \(P_{ab}\) but different collision capacities, depending
on whether their product mass lies below or above the unit Schur
threshold.
\end{assessment}

\section{Insertion of the V83 bad core shell before unmarking}
\label{sec:V85-core-slice}

Use the exact hyperbolic scalarization of V65.  For a fixed retained
core ancestry \(S\),
\begin{equation}
 F_{ab}(t;S)
 =
 \sum_{\{a_\lambda b_\mu k_\nu>t\}}
 x_\lambda^ay_\mu^bz_\nu^C(S).
 \label{eq:V85-Fab-recall}
\end{equation}
Define the oriented capped capacity
\begin{equation}
 \mathcal K_{\mathrm{hyp}}^{ab}
 :=
 \int_0^{W_*}
 \min\{1,\max_SF_{ab}(t;S)\}\,\dd t.
 \label{eq:V85-oriented-hyperbolic-capacity}
\end{equation}
It is one oriented entry, not the full two-orientation minimum in
\eqref{eq:V65-hyperbolic-capacity}.

\begin{theorem}[Core-slice lower reduction to the pair-product
quantile]
\label{thm:V85-core-slice-quantile}
Fix an ancestry \(S_0\) and a set \(I\) of core output labels.  Suppose
\begin{equation}
 k_\nu\ge\kappa_0>0
 \quad(\nu\in I),
 \qquad
 Z_I(S_0):=\sum_{\nu\in I}z_\nu^C(S_0)>0.
 \label{eq:V85-core-slice-hypothesis}
\end{equation}
Then
\begin{equation}
 F_{ab}(t;S_0)
 \ge
 Z_I(S_0)G_{ab}(t/\kappa_0)
 \quad
 (0\le t\le\kappa_0A_*^2),
 \label{eq:V85-Fab-slice-lower}
\end{equation}
and consequently
\begin{equation}
 \boxed{
 \mathcal K_{\mathrm{hyp}}^{ab}
 \ge
 \kappa_0
 \int_0^{A_*^2}
 \min\{1,Z_I(S_0)G_{ab}(r)\}\,\dd r.}
 \label{eq:V85-oriented-quantile-lower}
\end{equation}
\end{theorem}

\begin{proof}
If \(a_\lambda b_\mu>t/\kappa_0\) and \(\nu\in I\), then
\[
 a_\lambda b_\mu k_\nu>t.
\]
All summands in \eqref{eq:V85-Fab-recall} are nonnegative.  Keeping
only these triples gives
\[
 F_{ab}(t;S_0)
 \ge
 \sum_{\{a_\lambda b_\mu>t/\kappa_0\}}
 x_\lambda^ay_\mu^b
 \sum_{\nu\in I}z_\nu^C(S_0),
\]
which is \eqref{eq:V85-Fab-slice-lower}.  Insert this lower bound in
\eqref{eq:V85-oriented-hyperbolic-capacity}, restrict the integral to
\([0,\kappa_0A_*^2]\), and change variables
\(t=\kappa_0r\).
\end{proof}

\begin{corollary}[The balanced transition shell supplies a fixed
core mark]
\label{cor:V85-balanced-core-mark}
Assume the Wilson-to-heat comparison hypothesis
\eqref{eq:V83-restricted-comparison-hypothesis}, let
\[
 S_0=R_{2m_j},\qquad C=Q_{2m_j},
\]
and take
\[
 I_j=
 \{T_k:A_0\sqrt{2m_j}\le k\le B_0\sqrt{2m_j}\}.
\]
For all sufficiently large \(j\), the hypotheses of
\Cref{thm:V85-core-slice-quantile} hold with constants
\(\kappa_0,Z_0>0\) independent of \(j\).  More precisely, one may
take any
\[
 0<\kappa_0<\mathfrak d(u)
\]
for large \(j\), while
\begin{equation}
 Z_{I_j}(R_{2m_j})
 =
 \sum_{k\in I_j}\frac{(k+1)^3}{d_{2m_j}}
 \longrightarrow
 \frac{B_0^4-A_0^4}{2}.
 \label{eq:V85-balanced-core-Schur-mass}
\end{equation}
\end{corollary}

\begin{proof}
For the multiplicity-one path
\(R_{2m}\otimes Q_{2m}\to T_k\), V65 gives
\[
 z_k^C(R_{2m})=\frac{d_{T_k}}{d_{2m}}
 =\frac{(k+1)^3}{d_{2m}}.
\]
Equation \eqref{eq:V85-balanced-core-Schur-mass} is
\eqref{eq:V79-thermal-dimension-mass}.

The scalar hyperbolic core weight satisfies
\[
 k_k
 =
 [1+s_\alpha\Xi(T_k)]\|K_{T_k}\|_{\mathrm{op}}
 \ge
 \|K_{T_k}e_{m,m;2m,k}\|
 =
 \|\boldsymbol\kappa_{m,m;2m,k}^{\alpha}\|.
\]
The last equality is
\eqref{eq:V79-ancestral-vector-formula}.  Uniform nondecay
\eqref{eq:V83-Wilson-core-nondecay} proves the claimed lower bound
for \(k_k\).
\end{proof}

\begin{assessment}[The bad core does not disappear under the exact
product level]
\label{ass:V85-core-does-not-disappear}
V84 showed that the bad core forces a positive absolute-core first
capacity.  The present result is sharper upstream: before separate
pair integration, the same core supplies a fixed mark
\((\kappa_0,Z_0)\) to the pair-product layer.  Hence the oriented
hyperbolic capacity can be small only if the capped product-layer
integral on the right of
\eqref{eq:V85-oriented-quantile-lower} is small.
\end{assessment}

\section{Exact quantile--tail decomposition}
\label{sec:V85-quantile-tail}

For \(Z>0\), put
\begin{align}
 \mathcal Q_{ab}(Z)
 &:=
 \int_0^{A_*^2}\min\{1,ZG_{ab}(r)\}\,\dd r,
 \label{eq:V85-pair-quantile-capacity}\\
 q_{ab}(Z)
 &:=
 \sup\{r\in[0,A_*^2]:ZG_{ab}(r)\ge1\},
 \label{eq:V85-pair-Schur-quantile}
\end{align}
with the supremum of the empty set equal to zero.

\begin{theorem}[Pair-product quantile and tail identity]
\label{thm:V85-quantile-tail-identity}
Up to the value at the single threshold point, which does not affect
the integral,
\begin{equation}
 \boxed{
 \mathcal Q_{ab}(Z)
 =
 q_{ab}(Z)
 +Z\int_{q_{ab}(Z)}^{A_*^2}G_{ab}(r)\,\dd r.}
 \label{eq:V85-quantile-tail-identity}
\end{equation}
In particular,
\begin{align}
 q_{ab}(Z)
 &\le\mathcal Q_{ab}(Z),
 \label{eq:V85-quantile-lower}\\
 \mathcal Q_{ab}(Z)
 &\le
 \min\{A_*^2,ZP_{ab}\}.
 \label{eq:V85-quantile-first-upper}
\end{align}
\end{theorem}

\begin{proof}
The function \(G_{ab}\) is nonincreasing.  Apart from a possible jump
at \(q=q_{ab}(Z)\), one has \(ZG_{ab}(r)\ge1\) for \(r<q\) and
\(ZG_{ab}(r)<1\) for \(r>q\).  Split the integral at \(q\) to obtain
\eqref{eq:V85-quantile-tail-identity}.  The first inequality follows
by discarding the nonnegative tail.  For the second, use
\(\min\{1,ZG\}\le1\) and
\(\min\{1,ZG\}\le ZG\), then apply
\eqref{eq:V85-product-layer-first-moment}.
\end{proof}

\begin{firewall}[A first-moment bound is not a quantile theorem]
\label{fire:V85-first-moment-not-quantile}
Equation \eqref{eq:V85-quantile-first-upper} recovers the V84
separated pair-factor route.  It can be sharp when the product layer
stays below the cap, but it can be much larger than the true capped
capacity when a large first moment is concentrated on a small level
interval.  Conversely, an upper bound on \(P_{ab}\) alone gives no
information sharper than \eqref{eq:V85-quantile-first-upper}.

The new unresolved object is the actual distribution function
\(G_{ab}(r)\), or equivalently its unit Schur quantile and post-quantile
tail in \eqref{eq:V85-quantile-tail-identity}.
\end{firewall}

\section{Why one orientation still cannot decide the full capacity}
\label{sec:V85-full-capacity}

The exact V65 capacity is
\[
 \mathcal K_{\mathrm{hyp}}
 =
 \int_0^{W_*}
 \min\{1,\max_SF_{ab}(t;S),\max_{S'}F_{ac}(t;S')\}\,\dd t.
\]

\begin{proposition}[Required common-level obstruction]
\label{prop:V85-common-level-obstruction}
A lower bound for
\(\mathcal K_{\mathrm{hyp}}^{ab}\) obtained from
\eqref{eq:V85-oriented-quantile-lower} does not lower-bound
\(\mathcal K_{\mathrm{hyp}}\).  Such a lower bound transfers to the
full capacity only on levels where both
\(\max_SF_{ab}(t;S)\) and
\(\max_{S'}F_{ac}(t;S')\) have simultaneous positive lower bounds.
\end{proposition}

\begin{proof}
The integrand of the full capacity is the minimum of the two
orientation functions and one.  A lower bound on only one of the two
functions gives no lower bound on their minimum.  If both bounds hold
on a shared level set, their smaller lower bound does pass through the
minimum.  The shared-level requirement is essential: estimates on
different level subsets cannot be combined pointwise.
\end{proof}

\begin{assessment}[Audited direction after V85]
\label{ass:V85-audited-direction}
The V83 transition profile is now known to survive:
\begin{enumerate}[label=\textup{(\roman*)}]
\item passage from the signed vector to the positive square;
\item passage from the square to the absolute core first moment; and
\item restriction to an exact marked slice of the hyperbolic
      product-level capacity.
\end{enumerate}
What remains undecided is external and two-ended: the pair-product
level distribution \(G_{ab}\), followed by the common-level minimum
with the complete \(ac\) orientation.  This is the analogue of
testing a quotient invariant rather than a chosen section.  No further
core curvature calculation can decide it.
\end{assessment}

\section{V85 ledger and V86 mandate}
\label{sec:V85-ledger}

\begin{theorem}[Integrated V85 statement]
\label{thm:V85-integrated}
Preserving V1--V84, V85 proves:
\begin{enumerate}[label=\textup{(V85.\arabic*)},leftmargin=5.7em]
\item the stable SM V86 and NS V31 companion audit with a strict
      nontransfer firewall;
\item the exact one-sided pair layer identities
      \eqref{eq:V85-left-layer-integral}--%
      \eqref{eq:V85-right-layer-integral};
\item the exact pair-product first moment
      \eqref{eq:V85-product-layer-first-moment};
\item the core-slice reduction
      \eqref{eq:V85-oriented-quantile-lower};
\item insertion of the V83 balanced bad shell with fixed core Schur
      mass \eqref{eq:V85-balanced-core-Schur-mass};
\item the exact product quantile--tail identity
      \eqref{eq:V85-quantile-tail-identity}; and
\item the necessity of simultaneous same-level control of both
      orientations for a full-capacity obstruction.
\end{enumerate}
\end{theorem}

\begin{proof}
These are
\Cref{prop:V85-one-sided-layer-identities,
thm:V85-product-layer-first-moment,
thm:V85-core-slice-quantile,
cor:V85-balanced-core-mark,
thm:V85-quantile-tail-identity,
prop:V85-common-level-obstruction}.
\end{proof}

\begin{openproblem}[V86 mandate: actual pair-product distribution]
\label{op:V86-mandate}
Work with the exact external-pair layer
\eqref{eq:V85-pair-product-layer}:
\begin{enumerate}
\item write \(a_\lambda,b_\mu\) as the actual Wilson binary defect
      weights, including their output-mass factors;
\item impose the simultaneous-sideways inequalities
      \(s_\alpha\Xi_v>M\) and
      \(x_v<\varepsilon\Xi_v\) at the relevant endpoints;
\item derive a level-set dimension estimate for
      \(G_{ab}(r)\), not merely its first moment;
\item insert that estimate into the exact quantile--tail identity
      \eqref{eq:V85-quantile-tail-identity};
\item preserve the shared level \(r\) when constructing the full
      \(ac\) orientation, so the minimum in
      \eqref{eq:V65-hyperbolic-capacity} is taken pointwise;
\item test the result on the V60 alternating-height array and on
      Cartan and dual-singlet pair outputs.
\end{enumerate}
If only marginal first moments are recoverable, return to the
row-factorized four-input coefficient before the two pair labels are
separated.
\end{openproblem}

\section{Firewall after V85}
\label{sec:V85-firewall}

\begin{firewall}[Current claim boundary]
\label{fire:V85-final}
V85 proves the mandatory companion audit and the exact
pair-product quantile through which the V83 bad core enters one
oriented hyperbolic capacity.  It does not estimate the actual Wilson
pair-product level distribution on the simultaneous-sideways
residual, synchronize the \(ab\) and \(ac\) orientations, prove or
disprove the full JREC capacity, close puncture or multi-collision
aggregation, construct the continuum Yang--Mills theory, or prove a
positive mass gap.
\end{firewall}

\section{Append-only record for V85}
\label{sec:V85-record}

V85 was appended to the complete V84 manuscript.  The exact V84
parent had \(58945\) validation lines, \(2084556\) bytes, and
SHA--256 digest
\[
 \texttt{d61f43eb2fe57fdef64fd8a921599b3dc6bd109c2f5f13609e26ce1ad38e8803}.
\]
No V84 mathematical content was changed.  The sole final
\verb|\end{document}| is restored below.

% =============================================================================
% END APPENDED VERSION V85 MATERIAL
% =============================================================================

% =============================================================================
% BEGIN APPENDED VERSION V86 MATERIAL
% =============================================================================

\chapter{V86: Spectator-Inflation No-Go and the Forced
Overlapping-Pair Carrier}
\label{ch:V86}

\section{Purpose}
\label{sec:V86-purpose}

V85 isolated the external pair-product distribution
\(G_{ab}\) seen by the balanced bad core shell.  Its V86 mandate
asked whether the simultaneous-sideways inequalities themselves
force decay of that distribution.

They do not.  Those inequalities compare the collision-edge input
mass with the total mass of an original row.  A representation on a
disjoint punctured-core or spectator coordinate can increase the
total row mass without changing any external collision-pair output,
weight, or Schur layer.  This note proves that spectator-inflation
no-go exactly.

The conclusion does not abandon the programme.  The V59 incidence
partition and V60 puncture alphabet show where the inflated mass must
occur in an actually unclosed monomial: a cut-heavy bank, a retained
singleton mean, an off-cut core bank, or a genuine internal pair
defect.  The first two alternatives are already paid.  In the
remaining branch, a heavy internal pair must be assembled jointly
with the two collision pairs and the ternary cut.  Because that pair
can share a retained row, its product-state capacity need not
factorize into the marginal \(P_{ab}\).

\section{External pair layers are blind to disjoint row mass}
\label{sec:V86-spectator-inflation}

Fix a collision pair bank \(e=uv\), its input representations on the
coordinates of \(e\), all multiplicity-resolved outputs, and the
weighted defect values
\[
 a_{e,\lambda}
 =
 [1+s_\alpha\Xi(S_\lambda)]
 |q_{\alpha,S_\lambda}
   -q_{\alpha,R_e}q_{\alpha,Q_e}|.
\]
Let \(D_e^v(t)\) be its one-sided layer function as in V85.  A
\emph{disjoint spectator extension at \(v\)} tensors the row-\(v\)
carrier with an irreducible representation on a new coordinate not
belonging to \(e\), leaving the pair-bank carrier and operator
unchanged.

\begin{lemma}[Spectator invariance of one-pair data]
\label{lem:V86-spectator-invariance}
Under a disjoint spectator extension at either endpoint:
\begin{enumerate}[label=\textup{(\roman*)}]
\item every pair output \(S_\lambda\), multiplicity, and weighted
      defect \(a_{e,\lambda}\) is unchanged;
\item every one-sided Schur scalar
      \(x_\lambda^v=d_{S_\lambda}/d_v^{\,e}\) is unchanged, where
      \(d_v^{\,e}\) is the dimension of the row-\(v\) carrier on the
      pair bank alone;
\item the complete layer function \(D_e^v(t)\) and its one-sided
      first moment are unchanged.
\end{enumerate}
\label{lem:V86-spectator-invariance-end}
\end{lemma}

\begin{proof}
All quantities in (i) are defined by the tensor product on the fixed
coordinate set \(e\).  A new disjoint coordinate tensors every pair
projector with the identity and does not change the central
multiplier attached to the old output.

For (ii), partial trace on the pair coordinate is still the
resolved-copy identity
\eqref{eq:V84-pair-copy-partial-trace}.  The spectator identity acts
on a different tensor factor and is not included in
\(d_v^{\,e}\).  Thus the scalar is unchanged.  The layer function and
its integral are sums of the unchanged data, proving (iii).
\end{proof}

\begin{theorem}[Spectator-inflation no-go]
\label{thm:V86-spectator-inflation}
Fix nonzero external pair data and numbers \(M>0\),
\(0<\varepsilon<1\).  For every \(s_\alpha>0\), there are disjoint
spectator extensions for which
\begin{equation}
 s_\alpha\Xi_v>M,
 \qquad
 x_v<\varepsilon\Xi_v,
 \label{eq:V86-sideways-by-inflation}
\end{equation}
while the complete one-pair layer function is unchanged.  The
construction can be performed simultaneously at any finite
collection of endpoints, leaving the external pair-product layer
\(G_{ab}\) unchanged.
\end{theorem}

\begin{proof}
Let \(x_v\) be the fixed input mass on the collision pair.  The
Casimirs of the irreducibles \(R_N=(N,0)\) tend to infinity.  Choose
\(N\) so that the mass \(L_N=\Xi(R_N)\) of a new disjoint coordinate
satisfies
\[
 L_N>
 \max\left\{\frac{M}{s_\alpha},\frac{x_v}{\varepsilon}\right\}.
\]
After tensoring this coordinate onto row \(v\),
\[
 \Xi_v\ge L_N,
\]
so both inequalities in \eqref{eq:V86-sideways-by-inflation} hold
after an immaterial strict enlargement of \(N\).  The layer function
is unchanged by
\Cref{lem:V86-spectator-invariance}.  Use disjoint new coordinates at
different endpoints.  Both factors entering
\eqref{eq:V85-pair-product-layer} then remain unchanged, so
\(G_{ab}\) is unchanged.
\end{proof}

\begin{corollary}[No pair-only sideways level estimate]
\label{cor:V86-no-pair-only-estimate}
There is no universal estimate of the form
\begin{equation}
 G_{ab}(r)
 \le
 \Phi\left(
 r,\frac{x_a}{\Xi_a},\frac{x_b}{\Xi_b},
 s_\alpha\Xi_a,s_\alpha\Xi_b
 \right)
 \label{eq:V86-forbidden-pair-estimate}
\end{equation}
whose right side tends to zero solely because
\(x_v/\Xi_v\to0\) and \(s_\alpha\Xi_v\to\infty\), uniformly over the
remaining row data.  The same is true of
\(\mathcal Q_{ab}(Z)\) and \(q_{ab}(Z)\).
\end{corollary}

\begin{proof}
Begin with any pair data for which \(G_{ab}(r_0)>0\) at some
\(r_0>0\).  Apply
\Cref{thm:V86-spectator-inflation} with spectator masses tending to
infinity.  The purported right side tends to zero while the left side
is the fixed positive number \(G_{ab}(r_0)\).  Since
\(\mathcal Q_{ab}(Z)\) and the quantile are functionals of the same
unchanged \(G_{ab}\), the identical argument applies to them.
\end{proof}

\section{A concrete nonzero external pair layer}
\label{sec:V86-concrete-pair}

\begin{proposition}[Fundamental--dual singlet gives a nonzero layer]
\label{prop:V86-fundamental-dual-layer}
Let one external pair input be
\(F\otimes F^*\), and retain its singlet output.  For a nondegenerate
Wilson law with \(0<q_{\alpha,F}<1\), its weighted defect satisfies
\begin{equation}
 a_{\mathbf1}
 \ge
 1-q_{\alpha,F}^2>0,
 \label{eq:V86-fundamental-singlet-defect}
\end{equation}
and either one-sided Schur scalar of that resolved copy is \(1/3\).
If both external pair banks contain this resolved copy, then
\begin{equation}
 G_{ab}(r)\ge\frac19
 \quad
 \text{for }
 0\le r<a_{\mathbf1}^{(oa)}a_{\mathbf1}^{(bc)}.
 \label{eq:V86-fundamental-product-layer}
\end{equation}
\end{proposition}

\begin{proof}
The singlet multiplier is one, and
\(q_{\alpha,F^*}=q_{\alpha,F}\).  Hence its binary defect is
\[
 q_{\alpha,\mathbf1}
 -q_{\alpha,F}q_{\alpha,F^*}
 =
 1-q_{\alpha,F}^2.
\]
The output-mass factor in the weighted defect is at least one, giving
\eqref{eq:V86-fundamental-singlet-defect}.  The resolved singlet has
dimension one and either input has dimension three, so
\eqref{eq:V84-pair-copy-partial-trace} gives the scalar \(1/3\).
The product of the two selected Schur scalars contributes \(1/9\) to
every level below the product of their two weights.
\end{proof}

\begin{remark}[Scope of the example]
\label{rem:V86-example-scope}
At fixed fundamental representations,
\(1-q_{\alpha,F}^2\) tends to zero in the weak-coupling limit.  The
example proves the logical independence of pair layers from sideways
mass ratios; it is not a uniform weak-coupling counterexample to
JREC.
\end{remark}

\section{Where the missing mass actually goes}
\label{sec:V86-missing-mass}

V59 replaced the noncanonical V58 proposal by the exact incidence
identities
\eqref{eq:V59-a-mass-split}--%
\eqref{eq:V59-r-mass-split}.  V60 further proved, for every fixed
puncture monomial,
\[
 Y_i=P_{i,\mu}^{\rm pun}
 +\sum_{j\ne i}D_{ij,\mu}^{\rm pun}.
\]

\begin{theorem}[Forced heavy internal bank in the unclosed
one-sideways branch]
\label{thm:V86-forced-internal-bank}
Fix the small constant \(\varepsilon\) used in
\Cref{thm:V59-one-sideways-reduction}, and consider a fixed local
partition monomial.
\begin{enumerate}[label=\textup{(\roman*)}]
\item If \(v=a\) is the unique sideways endpoint and the branch is
      not already closed, then \(Y_a\ge\varepsilon\Xi_a\).  Either
      \[
      P_{a,\mu}^{\rm pun}\ge\frac\varepsilon2\Xi_a,
      \]
      in which case the retained-mean mechanism closes the monomial,
      or there is \(j\ne a\) such that
      \begin{equation}
      D_{aj,\mu}^{\rm pun}
      \ge\frac\varepsilon6\Xi_a.
      \label{eq:V86-heavy-puncture-pair-a}
      \end{equation}
\item If \(v=r\in\{b,c\}\) is the unique sideways endpoint and the
      branch is not already closed, then
      \[
      D_r^{\rm off}+Y_r\ge\varepsilon\Xi_r.
      \]
      Hence either
      \begin{equation}
      D_r^{\rm off}\ge\frac\varepsilon2\Xi_r,
      \label{eq:V86-heavy-offcut-r}
      \end{equation}
      or \(Y_r\ge\varepsilon\Xi_r/2\).  In the latter case, either a
      puncture singleton mean is heavy and closes the monomial, or
      there is \(j\ne r\) such that
      \begin{equation}
      D_{rj,\mu}^{\rm pun}
      \ge\frac\varepsilon{12}\Xi_r.
      \label{eq:V86-heavy-puncture-pair-r}
      \end{equation}
\end{enumerate}
\end{theorem}

\begin{proof}
The first residual inequality is
\eqref{eq:V59-one-sideways-residual}.  Apply the exact V60 puncture
alphabet and its pigeonhole bounds
\eqref{eq:V60-puncture-mean-heavy}--%
\eqref{eq:V60-puncture-pair-heavy} with \(\eta=\varepsilon\).
The mean-heavy alternative is paid by
\Cref{cor:V60-puncture-mean-closed}.

For \(v=r\), the second inequality in
\eqref{eq:V59-one-sideways-residual} gives the displayed dichotomy
between \(D_r^{\rm off}\) and \(Y_r\).  If the latter is heavy, apply
the same V60 alphabet with \(\eta=\varepsilon/2\).  This gives the
constants in
\eqref{eq:V86-heavy-puncture-pair-r} and the same mean-heavy closure.
\end{proof}

\begin{assessment}[Meaning of spectator inflation inside an actual
monomial]
\label{ass:V86-inflation-meaning}
The abstract spectator used in
\Cref{thm:V86-spectator-inflation} cannot remain analytically
anonymous in the true collision decomposition.  The exact incidence
and puncture partitions force it into one of the alternatives in
\Cref{thm:V86-forced-internal-bank}.  Cut-heavy and singleton-mean
mass is already paid.  The unresolved inflation is therefore visible
as a heavy off-cut core bank or a genuine internal pair defect.
\end{assessment}

\section{The overlapping-pair positive carrier}
\label{sec:V86-overlapping-carrier}

Let \(\mathcal J\) be the finite collection of genuine internal pair
banks selected by a fixed connected monomial after the V60
regrouping, and define
\[
 Q_e:=\sum_{\gamma}c_{e,\gamma}P_{e,\gamma},
 \qquad e\in\mathcal J,
\]
with the same nonnegative weighted-defect convention as in
\eqref{eq:V84-pair-operators}.  The complete positive carrier before
marginal pair factorization is
\begin{equation}
 Q_{ab;\mathcal J}^{\rm abs}
 :=
 Q_{oa}\otimes Q_{bc}
 \otimes
 \left(\bigotimes_{e\in\mathcal J}Q_e\right)
 \otimes|\mathbf K^\sim|.
 \label{eq:V86-overlapping-positive-carrier}
\end{equation}
All tensor factors refer to disjoint coordinate banks.  Edges in
\(\mathcal J\) may nevertheless share original rows with \(oa\),
\(bc\), or each other.

\begin{proposition}[Exact retained-row reduction]
\label{prop:V86-overlapping-retained-reduction}
Fix retained rows \(I\subseteq\{o,a,b,c\}\).  Trace every row outside
\(I\) in
\eqref{eq:V86-overlapping-positive-carrier}.  An internal pair
operator \(Q_e\):
\begin{enumerate}[label=\textup{(\roman*)}]
\item becomes its Schur scalar times the identity if exactly one
      endpoint of \(e\) is retained;
\item becomes its total trace if neither endpoint is retained; and
\item remains the positive operator \(Q_e\) if both endpoints are
      retained.
\end{enumerate}
Consequently the retained capacity factorizes into marginal pair
scalars only when no internal edge of \(\mathcal J\) lies entirely
inside \(I\).
\end{proposition}

\begin{proof}
If exactly one endpoint is retained, the partial trace is scalar by
Schur's lemma, exactly as in
\eqref{eq:V84-pair-Schur-traces}.  If neither is retained, take its
total trace.  If both are retained, no partial trace acts on that
bank, so \(Q_e\) remains.  Tensor-product partial trace gives the
three alternatives simultaneously.  The last assertion follows
because a surviving \(Q_e\) can be correlated with all other
retained factors by a density matrix on the corresponding original
row.
\end{proof}

\begin{corollary}[Why overlapping bilinear taxes cannot be multiplied]
\label{cor:V86-no-overlapping-tax-product}
Suppose an internal pair \(e=ij\) has both endpoints in the retained
row set \(I\).  Then the exact product-state tax is
\begin{equation}
 \kappa_I^\otimes(Q_{ab;\mathcal J}^{\rm abs}),
 \label{eq:V86-overlapping-product-capacity}
\end{equation}
with \(Q_e\) inside the same expectation as the retained core
operator.  Replacing it by a product of its two one-edge response
numbers is not a consequence of
\Cref{lem:V58-tax-composition}, because that lemma applies to the two
row-disjoint pairs \(oa\) and \(bc\).
\end{corollary}

\begin{proof}
The exact positive product-state tax is
\Cref{thm:V72-complete-product-tax} applied to
\eqref{eq:V86-overlapping-positive-carrier}.  By
\Cref{prop:V86-overlapping-retained-reduction}, the internal
operator \(Q_e\) survives the complement trace.  It acts on two
retained row carriers and cannot be replaced by a scalar partial
trace.  The hypothesis of
\Cref{lem:V58-tax-composition} that the two taxed edges use the
row-disjoint partition \((oa)\mid(bc)\) is absent.
\end{proof}

\begin{assessment}[The corrected upstream target]
\label{ass:V86-corrected-target}
The pair-only layer \(G_{ab}\) is the right object only after every
large off-edge row mass has already been paid.  On the remaining
sideways branch, the next object is the product-state capacity
\eqref{eq:V86-overlapping-product-capacity}, with the heavy internal
pair from \Cref{thm:V86-forced-internal-bank} retained in the same
positive expectation as the ternary cut.

This is a stricter assembly requirement than multiplying another V53
number into \(P_{ab}\), but it is narrower than returning to an
arbitrary four-row operator: the coordinate banks, graph type, and
heavy internal edge are all fixed.
\end{assessment}

\section{V86 ledger and V87 mandate}
\label{sec:V86-ledger}

\begin{theorem}[Integrated V86 statement]
\label{thm:V86-integrated}
Preserving V1--V85, V86 proves:
\begin{enumerate}[label=\textup{(V86.\arabic*)},leftmargin=5.7em]
\item spectator invariance of every external pair layer;
\item the spectator-inflation no-go
      \eqref{eq:V86-sideways-by-inflation};
\item impossibility of a pair-only level estimate from sideways mass
      ratios;
\item an explicit nonzero fundamental--dual singlet layer
      \eqref{eq:V86-fundamental-product-layer};
\item the forced heavy puncture-pair or off-cut alternative
      \eqref{eq:V86-heavy-puncture-pair-a}--%
      \eqref{eq:V86-heavy-puncture-pair-r};
\item the exact overlapping-pair carrier
      \eqref{eq:V86-overlapping-positive-carrier}; and
\item the retained-row criterion deciding whether an internal pair
      scalarizes or remains correlated with the core.
\end{enumerate}
\end{theorem}

\begin{proof}
These are
\Cref{lem:V86-spectator-invariance,
thm:V86-spectator-inflation,
cor:V86-no-pair-only-estimate,
prop:V86-fundamental-dual-layer,
thm:V86-forced-internal-bank,
prop:V86-overlapping-retained-reduction,
cor:V86-no-overlapping-tax-product}.
\end{proof}

\begin{openproblem}[V87 mandate: heavy internal-pair capacity]
\label{op:V87-mandate}
Fix one graph type and the heavy internal pair guaranteed by V86.
Then:
\begin{enumerate}
\item choose a retained-row orientation before tracing any bank;
\item write the exact complement trace of
      \eqref{eq:V86-overlapping-positive-carrier}, including every
      Schur scalar and every pair operator that survives;
\item combine the surviving heavy pair operator with
      \((\mathbf K^\sim)^2\) in one product-state square;
\item derive a shellwise double-Schur bound that uses the heavy input
      fraction without applying two taxes to a shared row;
\item test the bound on the balanced Cartan heat profile, the
      fundamental--dual singlet, and the V60 alternating-height
      family;
\item if the heavy off-cut alternative
      \eqref{eq:V86-heavy-offcut-r} occurs instead, resolve its exact
      V47 common/exclusive-pair ancestry before choosing the
      orientation.
\end{enumerate}
The next compulsory companion audit is V90.
\end{openproblem}

\section{Firewall after V86}
\label{sec:V86-firewall}

\begin{firewall}[Current claim boundary]
\label{fire:V86-final}
V86 proves that sideways mass ratios alone cannot control the
external pair-product level and forces the missing mass into an exact
heavy internal-pair or off-cut carrier.  It does not estimate the
resulting overlapping-pair product-state capacity, synchronize both
orientations, prove JREC, close puncture or multi-collision
aggregation, construct the continuum Yang--Mills theory, or prove a
positive mass gap.
\end{firewall}

\section{Append-only record for V86}
\label{sec:V86-record}

V86 was appended to the complete V85 manuscript.  The exact V85
parent had \(59502\) validation lines, \(2102308\) bytes, and
SHA--256 digest
\[
 \texttt{581298126d1652207a57dc11393e920a5e47969a8ffb31f9a1aaac61c279a4b8}.
\]
No V85 mathematical content was changed.  The sole final
\verb|\end{document}| is restored below.

% =============================================================================
% END APPENDED VERSION V86 MATERIAL
% =============================================================================

% =============================================================================
% BEGIN APPENDED VERSION V87 MATERIAL
% =============================================================================

\chapter{V87: Heavy Internal-Pair Layer Tensorized with the
Core Square}
\label{ch:V87}

\section{Purpose}
\label{sec:V87-purpose}

V86 proved that the external pair-product layer cannot see the mass
which makes an endpoint sideways.  In an actually unclosed
one-sideways monomial, the exact incidence and puncture partitions
force a heavy off-cut bank or a heavy internal pair defect.  The
internal pair can share a retained row with the collision carrier, so
its bilinear tax must not be multiplied with an already-used tax on
that row.

This note gives a sign-safe alternative.  Resolve the heavy internal
pair by its level projectors, tensor each level projector with one
positive core output shell, and only then apply the double-Schur
product-state estimate to the full retained row carriers.  The result
is an explicit shell-by-layer capacity.  It uses the heavy internal
pair debit exactly once and retains its correlation with the core
shell through the minimum of an operator entry and a trace entry.

\section{Layer resolution of one heavy internal pair}
\label{sec:V87-internal-layer}

Let \(e=ij\) be one genuine internal pair bank selected as in
\Cref{thm:V86-forced-internal-bank}.  Let \(X,Y\) be its two
product-group input irreducibles.  For every multiplicity-resolved
output copy \(\gamma=(S,\mu)\), let
\[
 c_\gamma
 =
 [1+s_\alpha\Xi(S)]
 |q_{\alpha,S}-q_{\alpha,X}q_{\alpha,Y}|,
 \qquad 0\le c_\gamma\le A_*.
\]
Define
\begin{align}
 Q_e
 &:=
 \sum_\gamma c_\gamma P_{e,\gamma},
 \label{eq:V87-internal-pair-operator}\\
 P_e(t)
 &:=
 \sum_{\{\gamma:c_\gamma>t\}}P_{e,\gamma},
 \label{eq:V87-internal-level-projector}\\
 D_e(t)
 &:=
 \operatorname{Tr}P_e(t)
 =
 \sum_{\{\gamma:c_\gamma>t\}}d_{S_\gamma}.
 \label{eq:V87-internal-level-rank}
\end{align}

\begin{lemma}[Exact operator layer cake]
\label{lem:V87-operator-layer-cake}
In the strong and operator-norm senses,
\begin{equation}
 \boxed{
 Q_e=\int_0^{A_*}P_e(t)\,\dd t.}
 \label{eq:V87-operator-layer-cake}
\end{equation}
\end{lemma}

\begin{proof}
The projectors \(P_{e,\gamma}\) are mutually orthogonal and
\[
 c_\gamma
 =
 \int_0^{A_*}
 \boldsymbol1_{\{c_\gamma>t\}}\,\dd t.
\]
Insert this scalar identity in
\eqref{eq:V87-internal-pair-operator}.  Finite-dimensional
Peter--Weyl blocks permit termwise integration, proving the claim.
\end{proof}

Define the one-pair capped Schur capacity
\begin{equation}
 \mathcal K_e
 :=
 \int_0^{A_*}
 \min\left\{
 1,\frac{D_e(t)}{d_X},\frac{D_e(t)}{d_Y}
 \right\}\,\dd t.
 \label{eq:V87-internal-pair-capacity}
\end{equation}
This is the V52 capacity envelope for the selected internal pair.

\section{A complete pair--shell double-Schur estimate}
\label{sec:V87-pair-shell}

Let \(M_j\ge0\) be one invariant core shell on
\(H_A\otimes H_B\).  In the application,
\[
 M_j
 =
 \operatorname{Tr}_C
 [P_{\Omega_j}(\mathbf K^\sim)^2].
\]
Write
\begin{equation}
 \Sigma_j:=\|M_j\|_{\mathrm{op}},
 \qquad
 \mathcal H_j:=\operatorname{Tr}M_j.
 \label{eq:V87-core-shell-sizes}
\end{equation}
The full retained row carriers are
\[
 H_{\mathsf a}=H_X\otimes H_A,
 \qquad
 H_{\mathsf b}=H_Y\otimes H_B.
\]
They are irreducible for the product of the internal-pair group and
the punctured-core group.

\begin{theorem}[Pair-level tensor core-shell bound]
\label{thm:V87-pair-shell-bound}
For every \(t\) and every positive invariant core shell,
\begin{equation}
 \boxed{
 \left\|
 P_e(t)\otimes M_j
 \right\|_{\mathsf a\otimes\mathsf b,\mathrm{prod}}
 \le
 \min\left\{
 \Sigma_j,
 \frac{D_e(t)\mathcal H_j}
 {\max\{d_Xd_A,d_Yd_B\}}
 \right\}.}
 \label{eq:V87-pair-shell-bound}
\end{equation}
\end{theorem}

\begin{proof}
The tensor product \(P_e(t)\otimes M_j\) is positive and commutes
with the product-group action on
\(H_{\mathsf a}\otimes H_{\mathsf b}\).  Its operator norm is
\[
 \|P_e(t)\|_{\mathrm{op}}\|M_j\|_{\mathrm{op}}
 \le\Sigma_j.
\]
Its trace is \(D_e(t)\mathcal H_j\), while the retained row
dimensions are \(d_Xd_A\) and \(d_Yd_B\).  Apply the double-Schur
product-state bound \eqref{eq:V73-double-Schur}.
\end{proof}

\begin{corollary}[Integrated heavy-pair shell bound]
\label{cor:V87-integrated-pair-shell}
One has
\begin{align}
 \left\|Q_e\otimes M_j\right\|_
 {\mathsf a\otimes\mathsf b,\mathrm{prod}}
 &\le
 \int_0^{A_*}
 \min\left\{
 \Sigma_j,
 \frac{D_e(t)\mathcal H_j}
 {\max\{d_Xd_A,d_Yd_B\}}
 \right\}\,\dd t
 \label{eq:V87-integrated-pair-shell}\\
 &\le
 \mathcal K_e
 \max\left\{
 \Sigma_j,
 \frac{\mathcal H_j}{\min\{d_A,d_B\}}
 \right\}.
 \label{eq:V87-factorized-pair-shell}
\end{align}
\end{corollary}

\begin{proof}
Use the operator layer cake
\eqref{eq:V87-operator-layer-cake}, linearity of expectations, and
the fact that the supremum of an integral of nonnegative functions is
at most the integral of their suprema.  Then apply
\eqref{eq:V87-pair-shell-bound}.

For the second estimate, put
\[
 R_e(t):=
 \frac{D_e(t)}{\max\{d_X,d_Y\}}.
\]
Since
\[
 \max\{d_Xd_A,d_Yd_B\}
 \ge
 \max\{d_X,d_Y\}\min\{d_A,d_B\},
\]
the trace entry in
\eqref{eq:V87-integrated-pair-shell} is at most
\[
 R_e(t)\frac{\mathcal H_j}{\min\{d_A,d_B\}}.
\]
For \(a,b,R\ge0\),
\[
 \min\{a,bR\}
 \le
 \max\{a,b\}\min\{1,R\}.
\]
Integrate this inequality and observe that
\[
 \min\{1,R_e(t)\}
 =
 \min\left\{
 1,\frac{D_e(t)}{d_X},\frac{D_e(t)}{d_Y}
 \right\}.
\]
This proves \eqref{eq:V87-factorized-pair-shell}.
\end{proof}

\begin{assessment}[The un-factorized integral is the sharper result]
\label{ass:V87-unfactorized-sharper}
Equation \eqref{eq:V87-integrated-pair-shell} retains the shared
minimum between the pair level and the core shell.  The simpler
\eqref{eq:V87-factorized-pair-shell} is useful for importing a proved
one-pair capacity, but it can lose the correlation between small
\(D_e(t)\) and a trace-large core shell.  The programme should use the
first display whenever an explicit level distribution is available.
\end{assessment}

\section{Summation over core shells}
\label{sec:V87-shell-sum}

Let
\[
 M^{(2)}=\sum_{j\in J}M_j
\]
be a positive invariant shell decomposition of the core square.

\begin{theorem}[Combined heavy-pair/core-square capacity]
\label{thm:V87-combined-capacity}
The exact product-state square capacity obeys
\begin{align}
 \left\|Q_e\otimes M^{(2)}\right\|_
 {\mathsf a\otimes\mathsf b,\mathrm{prod}}
 &\le
 \sum_{j\in J}\int_0^{A_*}
 \min\left\{
 \Sigma_j,
 \frac{D_e(t)\mathcal H_j}
 {\max\{d_Xd_A,d_Yd_B\}}
 \right\}\,\dd t
 \label{eq:V87-combined-shell-layer}\\
 &\le
 \mathcal K_e
 \sum_{j\in J}
 \max\left\{
 \Sigma_j,
 \frac{\mathcal H_j}{\min\{d_A,d_B\}}
 \right\}.
 \label{eq:V87-combined-factorized}
\end{align}
No two bilinear taxes sharing a row are multiplied.
\end{theorem}

\begin{proof}
For every product density, positivity and linearity give the sum of
the shell expectations.  Its supremum is at most the sum of their
suprema.  Apply
\Cref{cor:V87-integrated-pair-shell} to each shell.
\end{proof}

\begin{remark}[Insertion of the remaining external pairs]
\label{rem:V87-external-scalars}
If the chosen orientation traces exactly one endpoint of each
external pair \(oa\) and \(bc\), their complement traces give the
same scalar factors as in
\eqref{eq:V84-pair-Schur-traces}.  They multiply the right sides of
\eqref{eq:V87-combined-shell-layer}--%
\eqref{eq:V87-combined-factorized}.  An additional internal pair
whose two endpoints are retained must instead remain inside the joint
capacity and can be treated by a further level variable.
\end{remark}

\section{The heavy-edge capacity supplied by V52--V53}
\label{sec:V87-heavy-edge-capacity}

The one-edge theorem in V53 is usually stated as a coefficient tax.
Its proof also yields the following capacity statement, which is the
quantity required in \eqref{eq:V87-combined-factorized}.

\begin{proposition}[Selected-heavy one-pair capacity]
\label{prop:V87-selected-heavy-capacity}
Let the input-\(v\) mass \(x\) on the internal pair satisfy
\[
 x\ge\varepsilon\Xi_v.
\]
With \(h\) and \(N_{\rm eff}=x^2/h\) as in V53,
\begin{equation}
 \boxed{
 \mathcal K_e
 \le
 C_\varepsilon
 \frac{h}{s_\alpha\Xi_v^2}.}
 \label{eq:V87-selected-heavy-capacity}
\end{equation}
\end{proposition}

\begin{proof}
If \(s_\alpha N_{\rm eff}\le M\), then
\[
 \frac{h}{s_\alpha\Xi_v^2}
 =
 \frac{x^2}{s_\alpha N_{\rm eff}\Xi_v^2}
 \ge
 \frac{\varepsilon^2}{M},
\]
while \(\mathcal K_e\le A_*\).  This proves the claim in the first
regime.

Assume \(s_\alpha N_{\rm eff}>M\).  Split the output family into
\(\Xi(S)<\kappa N_{\rm eff}\) and its complement, with the fixed
\(\kappa\) of
\Cref{thm:V53-general-effective-capacity}.  The cold family has
capped one-sided Schur mass at most
\[
 \tau_{\rm eff}(\kappa)
 \le\frac C{s_\alpha N_{\rm eff}}
\]
by \eqref{eq:V53-general-effective-capacity}.  Its contribution to
\(\mathcal K_e\) is at most \(A_*\tau_{\rm eff}\).

On the complementary family, the proof of
\eqref{eq:V53-effective-functional} gives the pointwise weighted
defect bound
\[
 c_\gamma\le\frac C{s_\alpha N_{\rm eff}}.
\]
Therefore no hot-family output occurs above that level; its entire
capped level integral is at most the same quantity.  Thus
\[
 \mathcal K_e
 \le\frac C{s_\alpha N_{\rm eff}}
 =
 \frac{Ch}{s_\alpha x^2}
 \le
 C_\varepsilon\frac h{s_\alpha\Xi_v^2}.
\]
\end{proof}

\begin{corollary}[Heavy internal-pair debit in the joint square]
\label{cor:V87-heavy-pair-joint-debit}
If the internal pair in
\Cref{thm:V87-combined-capacity} carries a fixed fraction of a
sideways endpoint mass, then
\begin{equation}
 \left\|Q_e\otimes M^{(2)}\right\|_{\mathrm{prod}}
 \le
 C_\varepsilon
 \frac{h_e}{s_\alpha\Xi_v^2}
 \sum_{j\in J}
 \max\left\{
 \Sigma_j,
 \frac{\mathcal H_j}{\min\{d_A,d_B\}}
 \right\}.
 \label{eq:V87-heavy-pair-joint-debit}
\end{equation}
\end{corollary}

\begin{proof}
Insert \eqref{eq:V87-selected-heavy-capacity} into
\eqref{eq:V87-combined-factorized}.
\end{proof}

\section{Sharp model tests}
\label{sec:V87-tests}

\begin{proposition}[Internal singlet gains the full combined row
dimension]
\label{prop:V87-internal-singlet-test}
Suppose one internal level projector is a rank-one singlet, so
\(D_e(t)=1\) on its active level interval.  Then
\eqref{eq:V87-pair-shell-bound} gives
\begin{equation}
 \left\|P_{\mathbf1}\otimes M_j\right\|_{\mathrm{prod}}
 \le
 \min\left\{
 \Sigma_j,
 \frac{\mathcal H_j}
 {\max\{d_Xd_A,d_Yd_B\}}
 \right\}.
 \label{eq:V87-singlet-combined-gain}
\end{equation}
Thus the trace branch sees the dimension of the full retained row,
not merely the internal-pair input dimension.
\end{proposition}

\begin{proof}
Set \(D_e(t)=1\) in
\eqref{eq:V87-pair-shell-bound}.
\end{proof}

\begin{proposition}[Cartan product visibility prevents a universal
operator gain]
\label{prop:V87-Cartan-product-test}
Let \(P_e\) be a Cartan highest-weight projector containing a product
highest-weight vector, and let \(M_j\) have a product eigenvector for
its top eigenvalue on the core carriers.  Then
\begin{equation}
 \|P_e\otimes M_j\|_{\mathrm{prod}}=\Sigma_j.
 \label{eq:V87-Cartan-no-gain}
\end{equation}
Hence the operator entry in
\eqref{eq:V87-pair-shell-bound} cannot be removed.
\end{proposition}

\begin{proof}
Tensor the product highest-weight vector in the internal pair with
the product top eigenvector of \(M_j\), regrouping the two factors on
each original row.  This is an admissible product vector across the
two rows.  Its expectation is \(\Sigma_j\).  The reverse inequality
is the operator norm bound.
\end{proof}

\begin{assessment}[Interaction with the V83 bad core]
\label{ass:V87-V83-test}
The V83 balanced heat profile supplies a lower bound on one
product-visible core column.  Proposition
\ref{prop:V87-Cartan-product-test} shows that a simultaneously
product-visible Cartan internal pair need not reduce it through the
operator branch.  A singlet or dimension-small internal output,
however, couples to the trace branch with the enlarged full-row
denominator in
\eqref{eq:V87-singlet-combined-gain}.  The correct V88 calculation is
therefore the actual level distribution \(D_e(t)\) together with the
thermal core shell traces, not an assertion that every heavy internal
pair is entangled.
\end{assessment}

\section{V87 ledger and V88 mandate}
\label{sec:V87-ledger}

\begin{theorem}[Integrated V87 statement]
\label{thm:V87-integrated}
Preserving V1--V86, V87 proves:
\begin{enumerate}[label=\textup{(V87.\arabic*)},leftmargin=5.7em]
\item the exact internal-pair operator layer cake
      \eqref{eq:V87-operator-layer-cake};
\item the pair-level/core-shell double-Schur estimate
      \eqref{eq:V87-pair-shell-bound};
\item the sharper integrated shell-by-layer bound
      \eqref{eq:V87-integrated-pair-shell};
\item the complete shell sum
      \eqref{eq:V87-combined-shell-layer};
\item the selected-heavy internal-pair capacity
      \eqref{eq:V87-selected-heavy-capacity};
\item its insertion as a single nonoverlapping debit
      \eqref{eq:V87-heavy-pair-joint-debit}; and
\item the sharp singlet-gain and Cartan-no-gain tests.
\end{enumerate}
\end{theorem}

\begin{proof}
These are
\Cref{lem:V87-operator-layer-cake,
thm:V87-pair-shell-bound,
cor:V87-integrated-pair-shell,
thm:V87-combined-capacity,
prop:V87-selected-heavy-capacity,
cor:V87-heavy-pair-joint-debit,
prop:V87-internal-singlet-test,
prop:V87-Cartan-product-test}.
\end{proof}

\begin{openproblem}[V88 mandate: thermal shell trace with one heavy
pair]
\label{op:V88-mandate}
Apply \eqref{eq:V87-combined-shell-layer} to the actual V86 heavy
internal bank:
\begin{enumerate}
\item resolve core shells so that
      \(\Sigma_j\) and \(\mathcal H_j\) retain the final thermal
      label \(T_k\);
\item insert the exact ternary multiplicity trace
      \eqref{eq:V73-shell-HS} before using a dimension bound;
\item keep \(D_e(t)\) inside the minimum with that trace;
\item test separately singlet-small, Cartan-visible, and intermediate
      internal pair levels;
\item in the V83 balanced family, distinguish the single bad product
      column from the full multiplicity trace;
\item decide whether
      \eqref{eq:V87-heavy-pair-joint-debit} reaches the marked-tree
      scale or whether only the un-factorized
      \eqref{eq:V87-combined-shell-layer} can do so.
\end{enumerate}
The next compulsory companion audit is V90.
\end{openproblem}

\section{Firewall after V87}
\label{sec:V87-firewall}

\begin{firewall}[Current claim boundary]
\label{fire:V87-final}
V87 proves a rigorous joint heavy-pair/core-shell capacity that avoids
multiplying overlapping row taxes and imports the proved V53 pair
debit once.  It does not estimate the resulting thermal shell sum,
close the off-cut alternative, synchronize both orientations, prove
JREC, close puncture or multi-collision aggregation, construct the
continuum Yang--Mills theory, or prove a positive mass gap.
\end{firewall}

\section{Append-only record for V87}
\label{sec:V87-record}

V87 was appended to the complete V86 manuscript.  The exact V86
parent had \(59971\) validation lines, \(2119966\) bytes, and
SHA--256 digest
\[
 \texttt{8a2ab7f491c63d81ac909787ec5aba6a139bb11c739b629459add8d2dd1605a8}.
\]
No V86 mathematical content was changed.  The sole final
\verb|\end{document}| is restored below.

% =============================================================================
% END APPENDED VERSION V87 MATERIAL
% =============================================================================

% =============================================================================
% BEGIN APPENDED VERSION V88 MATERIAL
% =============================================================================

\chapter{V88: Core Product Branch and the Exact Balanced
Thermal Multiplicity}
\label{ch:V88}

\section{Purpose}
\label{sec:V88-purpose}

V87 tensorized a heavy internal-pair level projector with a positive
core shell and applied double Schur on the full retained rows.  Its
operator entry used the unrestricted core operator norm
\(\Sigma_j\).  That is unnecessarily large: since
\(0\le P_e(t)\le I\), the pair--core expectation is also bounded by
the product-state norm of the core shell alone.

This refinement removes \(\Sigma_j\) from the factorized estimate.
The remaining core quantity is its trace divided by a retained core
dimension.  On the balanced Cartan family, the ternary multiplicity
can be counted exactly by the V81 Pieri support: it is \(k+1\) on
the output \(T_k\).  This yields a sharp statement of the new gap.
The V83 bad product column forces order-one total thermal trace, while
the current uniform bound permits order \(\sqrt m\).  Closing the
heavy-pair route therefore requires a multiplicity-average estimate,
or use of the un-factorized pair-level/core-shell minimum.

\section{Replacing the core operator norm by its product norm}
\label{sec:V88-core-product-branch}

For a positive core shell \(M_j\) on \(H_A\otimes H_B\), define
\begin{equation}
 \Pi_j^{\rm core}
 :=
 \|M_j\|_{A\otimes B,\mathrm{prod}}.
 \label{eq:V88-core-product-norm}
\end{equation}
The V73 double-Schur theorem gives
\begin{equation}
 \Pi_j^{\rm core}
 \le
 \min\left\{
 \Sigma_j,\frac{\mathcal H_j}{\max\{d_A,d_B\}}
 \right\}.
 \label{eq:V88-core-product-Schur}
\end{equation}

\begin{lemma}[Pair projection is dominated by the core product
branch]
\label{lem:V88-pair-below-core-product}
For every internal-pair level projector \(P_e(t)\),
\begin{equation}
 \left\|P_e(t)\otimes M_j\right\|_
 {\mathsf a\otimes\mathsf b,\mathrm{prod}}
 \le
 \Pi_j^{\rm core}.
 \label{eq:V88-pair-below-core-product}
\end{equation}
\end{lemma}

\begin{proof}
Since \(0\le P_e(t)\le I_{X\otimes Y}\),
\[
 0\le P_e(t)\otimes M_j\le I_{X\otimes Y}\otimes M_j.
\]
Evaluate this inequality on a product density
\(\rho_{XA}\otimes\rho_{YB}\).  Tracing \(X\) and \(Y\) from the
right side gives
\[
 \rho_A\otimes\rho_B,
 \qquad
 \rho_A=\operatorname{Tr}_X\rho_{XA},
 \quad
 \rho_B=\operatorname{Tr}_Y\rho_{YB}.
\]
These are density matrices, so the right-side expectation is at most
\(\Pi_j^{\rm core}\).  Take the supremum over the original full-row
product densities.
\end{proof}

\begin{theorem}[Refined pair-level/core-shell minimum]
\label{thm:V88-refined-pair-shell}
The V87 pair--shell estimate improves to
\begin{equation}
 \boxed{
 \left\|P_e(t)\otimes M_j\right\|_
 {\mathsf a\otimes\mathsf b,\mathrm{prod}}
 \le
 \min\left\{
 \Pi_j^{\rm core},
 \frac{D_e(t)\mathcal H_j}
 {\max\{d_Xd_A,d_Yd_B\}}
 \right\}.}
 \label{eq:V88-refined-pair-shell}
\end{equation}
Consequently
\begin{align}
 \left\|Q_e\otimes M_j\right\|_
 {\mathsf a\otimes\mathsf b,\mathrm{prod}}
 &\le
 \int_0^{A_*}
 \min\left\{
 \Pi_j^{\rm core},
 \frac{D_e(t)\mathcal H_j}
 {\max\{d_Xd_A,d_Yd_B\}}
 \right\}\,\dd t
 \label{eq:V88-refined-integral}\\
 &\le
 \boxed{
 \mathcal K_e
 \frac{\mathcal H_j}{\min\{d_A,d_B\}}.}
 \label{eq:V88-refined-factorized}
\end{align}
\end{theorem}

\begin{proof}
The first entry in
\eqref{eq:V88-refined-pair-shell} is
\Cref{lem:V88-pair-below-core-product}; the second is the trace entry
in \eqref{eq:V87-pair-shell-bound}.  Integrate by the operator layer
cake as in
\Cref{cor:V87-integrated-pair-shell}.

Put
\[
 R_e(t)=\frac{D_e(t)}{\max\{d_X,d_Y\}},
 \qquad
 B_j=\frac{\mathcal H_j}{\min\{d_A,d_B\}}.
\]
The full-row dimension inequality used in V87 bounds the trace entry
by \(R_e(t)B_j\), while
\eqref{eq:V88-core-product-Schur} gives
\(\Pi_j^{\rm core}\le B_j\).  Hence
\[
 \min\{\Pi_j^{\rm core},R_e(t)B_j\}
 \le
 B_j\min\{1,R_e(t)\}.
\]
The integral of the last factor is exactly
\(\mathcal K_e\), proving
\eqref{eq:V88-refined-factorized}.
\end{proof}

\begin{corollary}[Refined heavy-pair shell sum]
\label{cor:V88-refined-shell-sum}
For a positive shell decomposition \(M^{(2)}=\sum_jM_j\),
\begin{equation}
 \boxed{
 \|Q_e\otimes M^{(2)}\|_{\mathrm{prod}}
 \le
 \mathcal K_e
 \sum_j\frac{\mathcal H_j}{\min\{d_A,d_B\}}.}
 \label{eq:V88-refined-shell-sum}
\end{equation}
If the internal pair is selected-heavy at endpoint \(v\), then
\begin{equation}
 \|Q_e\otimes M^{(2)}\|_{\mathrm{prod}}
 \le
 C_\varepsilon
 \frac{h_e}{s_\alpha\Xi_v^2}
 \sum_j\frac{\mathcal H_j}{\min\{d_A,d_B\}}.
 \label{eq:V88-selected-heavy-shell-sum}
\end{equation}
\end{corollary}

\begin{proof}
Sum \eqref{eq:V88-refined-factorized} over positive shells.
Then use \eqref{eq:V87-selected-heavy-capacity}.
\end{proof}

\begin{remark}[Why this does not multiply overlapping taxes]
\label{rem:V88-no-overlap}
The factor \(\mathcal K_e\) is used once, through the layer projector
which remains inside the full-row product-state expectation.  The
core trace ratio comes from Schur's lemma on the tensorized full row.
No second trace-norm bilinear map is applied to either endpoint of
\(e\).
\end{remark}

\section{Exact balanced ternary multiplicity}
\label{sec:V88-balanced-multiplicity}

Set
\[
 A=B=R_m,\qquad C=Q_{2m},\qquad T_k=(k,k).
\]

\begin{theorem}[Multiplicity \(k+1\) on the balanced thermal path]
\label{thm:V88-balanced-multiplicity}
For every \(0\le k\le m\),
\begin{equation}
 \boxed{
 m_{R_m,R_m,Q_{2m}}(T_k)=k+1.}
 \label{eq:V88-balanced-multiplicity}
\end{equation}
\end{theorem}

\begin{proof}
Use the \(AC\)-first bracketing.  The Pieri decomposition is
\[
 R_m\otimes Q_{2m}
 =
 \bigoplus_{j=0}^{m}(m-j,2m-j).
\]
By \eqref{eq:V81-AC-support}, an intermediate reaches \(T_k\) after
tensoring by the remaining \(R_m\) if and only if
\[
 m-k\le j\le m.
\]
There are exactly \(k+1\) such intermediates, and each of the two
Pieri arrows has multiplicity one.  Summing the path multiplicities
gives \eqref{eq:V88-balanced-multiplicity}.
\end{proof}

\begin{corollary}[Exact dimension of the final multiplicity space]
\label{cor:V88-final-multiplicity-dimension}
On \(0\le k\le m\),
\begin{equation}
 \dim\mathcal M_{T_k}=k+1,
 \qquad
 \operatorname{rank}
 (I_{H_{T_k}}\otimes I_{\mathcal M_{T_k}})
 =(k+1)^4.
 \label{eq:V88-final-isotypic-rank}
\end{equation}
\end{corollary}

\begin{proof}
The first identity is
\eqref{eq:V88-balanced-multiplicity}.  The SU(3) dimension formula
gives \(d_{T_k}=(k+1)^3\); multiply the two dimensions.
\end{proof}

\section{Balanced thermal trace bounds}
\label{sec:V88-thermal-trace}

For a single output \(T_k\), put
\begin{equation}
 \mathcal H_{m,k}^{(2)}
 :=
 d_{T_k}
 \operatorname{Tr}_{\mathcal M_{T_k}}
 (\widetilde K_{T_k})^2.
 \label{eq:V88-single-shell-trace}
\end{equation}
This is the exact V73 shell Hilbert--Schmidt trace for the singleton
shell \(\{T_k\}\).

\begin{proposition}[Universal upper trace on a balanced output]
\label{prop:V88-balanced-trace-upper}
For \(0\le k\le m\),
\begin{equation}
 \boxed{
 \mathcal H_{m,k}^{(2)}
 \le
 \frac{C(k+1)^4}
 {[1+s_\alpha\Xi(T_k)]^2}.}
 \label{eq:V88-balanced-trace-upper}
\end{equation}
\end{proposition}

\begin{proof}
The V63 weighted-cut bound is
\[
 \|\widetilde K_{T_k}\|_{\mathrm{op}}
 \le
 \frac C{1+s_\alpha\Xi(T_k)}.
\]
The trace of its square on a space of dimension \(k+1\) is at most
\((k+1)\|\widetilde K_{T_k}\|_{\mathrm{op}}^2\).  Multiply by
\(d_{T_k}=(k+1)^3\).
\end{proof}

Let \(d_m=d_{R_m}=(m+1)(m+2)/2\), and define the normalized thermal
trace
\begin{equation}
 \mathfrak T_{\alpha;m}^{A_0,B_0}
 :=
 \sum_{A_0\sqrt{2m}\le k\le B_0\sqrt{2m}}
 \frac{\mathcal H_{m,k}^{(2)}}{d_m}.
 \label{eq:V88-normalized-thermal-trace}
\end{equation}

\begin{theorem}[The exact remaining multiplicity window]
\label{thm:V88-multiplicity-window}
For fixed \(0<A_0<B_0<\infty\),
\begin{equation}
 \boxed{
 \mathfrak T_{\alpha;m}^{A_0,B_0}
 \le C_{A_0,B_0}\sqrt m.}
 \label{eq:V88-thermal-trace-upper}
\end{equation}
Under the hypotheses of
\Cref{thm:V83-conditional-obstruction-transfer},
\begin{equation}
 \boxed{
 \liminf_{j\to\infty}
 \mathfrak T_{\alpha_j;m_j}^{A_0,B_0}
 \ge
 \frac{B_0^4-A_0^4}{2}\mathfrak d(u)^2>0.}
 \label{eq:V88-thermal-trace-lower}
\end{equation}
\end{theorem}

\begin{proof}
On the thermal window \(k=O(\sqrt m)\).  By
\eqref{eq:V88-balanced-trace-upper} and
\(d_m\asymp m^2\), every summand is at most
\[
 C\frac{(k+1)^4}{m^2}\le C_{B_0}.
\]
There are \(O(\sqrt m)\) integers in the window, proving
\eqref{eq:V88-thermal-trace-upper}.

For the lower bound, let \(e=e_{m,m;2m,k}\) be the normalized Cartan
path vector.  Positivity gives
\begin{align*}
 \mathcal H_{m,k}^{(2)}
 &\ge
 d_{T_k}\|\widetilde K_{T_k}e\|^2\\
 &=
 (k+1)^3[1+s_\alpha\Xi(T_k)]^2
 \|\boldsymbol\kappa_{m,m;2m,k}^{\alpha}\|^2.
\end{align*}
Divide by \(d_m\), sum the window, and use
\eqref{eq:V84-positive-square-limit}.
\end{proof}

\begin{assessment}[The gap is an ancestry average]
\label{ass:V88-ancestry-average-gap}
The lower bound
\eqref{eq:V88-thermal-trace-lower} uses one distinguished Cartan
ancestry vector in each \((k+1)\)-dimensional multiplicity space.
The upper bound
\eqref{eq:V88-thermal-trace-upper} permits all \(k+1\) ancestry
directions to carry the V63 operator cap.  Since the output
dimension weight already has order-one total mass on
\(k\asymp\sqrt m\), the unresolved factor \(\sqrt m\) is exactly the
number of ancestry directions per thermal output.

Thus no improvement of the output count can close this trace branch.
One must prove an average square gain across the \(k+1\) ancestry
directions, or keep the internal-pair level projector inside the
minimum in \eqref{eq:V88-refined-integral}.
\end{assessment}

\section{Insertion into the heavy-pair gate}
\label{sec:V88-heavy-pair-insertion}

\begin{corollary}[Balanced heavy-pair thermal gate]
\label{cor:V88-balanced-heavy-pair-gate}
On the balanced thermal shell, the factorized V88 estimate is
\begin{equation}
 \boxed{
 \|Q_e\otimes M_{\mathrm{th}}^{(2)}\|_{\mathrm{prod}}
 \le
 \mathcal K_e
 \mathfrak T_{\alpha;m}^{A_0,B_0}.}
 \label{eq:V88-balanced-heavy-pair-gate}
\end{equation}
The presently proved uniform bounds give only
\begin{equation}
 \|Q_e\otimes M_{\mathrm{th}}^{(2)}\|_{\mathrm{prod}}
 \le
 C\mathcal K_e\sqrt m.
 \label{eq:V88-current-balanced-bound}
\end{equation}
\end{corollary}

\begin{proof}
In the balanced case
\(\min\{d_A,d_B\}=d_m\).  Restrict
\eqref{eq:V88-refined-shell-sum} to the thermal outputs.  Then apply
\eqref{eq:V88-thermal-trace-upper}.
\end{proof}

\begin{firewall}[No inverse square-root pair debit is assumed]
\label{fire:V88-no-assumed-debit}
The selected-heavy estimate
\[
 \mathcal K_e
 \le C_\varepsilon h_e/(s_\alpha\Xi_v^2)
\]
does not universally imply
\(\mathcal K_e\lesssim m^{-1/2}\) on the balanced core family.
The local edge geometry and the core label \(m\) belong to different
coordinate banks.  Declaring that inverse square-root factor without
a joint mass relation would repeat the spectator-inflation error
excluded by V86.
\end{firewall}

\section{V88 ledger and V89 mandate}
\label{sec:V88-ledger}

\begin{theorem}[Integrated V88 statement]
\label{thm:V88-integrated}
Preserving V1--V87, V88 proves:
\begin{enumerate}[label=\textup{(V88.\arabic*)},leftmargin=5.7em]
\item domination by the core product branch
      \eqref{eq:V88-pair-below-core-product};
\item the refined pair-level/core-shell minimum
      \eqref{eq:V88-refined-pair-shell};
\item the factorized heavy-pair trace gate
      \eqref{eq:V88-refined-factorized};
\item the exact balanced ternary multiplicity
      \eqref{eq:V88-balanced-multiplicity};
\item the universal normalized thermal trace bound
      \eqref{eq:V88-thermal-trace-upper};
\item the positive V83 lower trace
      \eqref{eq:V88-thermal-trace-lower}; and
\item localization of the remaining square-root loss to the average
      over \(k+1\) ancestry directions.
\end{enumerate}
\end{theorem}

\begin{proof}
These are
\Cref{lem:V88-pair-below-core-product,
thm:V88-refined-pair-shell,
cor:V88-refined-shell-sum,
thm:V88-balanced-multiplicity,
prop:V88-balanced-trace-upper,
thm:V88-multiplicity-window,
cor:V88-balanced-heavy-pair-gate}.
\end{proof}

\begin{openproblem}[V89 mandate: ancestry-averaged heat trace]
\label{op:V89-mandate}
Return to the exact heat benchmark on the full
\((k+1)\)-dimensional balanced multiplicity space:
\begin{enumerate}
\item express \(\operatorname{Tr}_{\mathcal M_{T_k}}K_{T_k}^2\)
      as a trace of the three pair-Casimir heat operators before
      choosing the Cartan vector;
\item use trace invariance under the three Racah bases to compute all
      diagonal square terms;
\item retain the signed cross traces, using the common group-integral
      representation from V75 rather than termwise absolute values;
\item determine whether the mean square per ancestry direction is
      \(O(k^{-1})\), \(O(1)\), or has a positive heat profile;
\item insert the result into
      \eqref{eq:V88-balanced-heavy-pair-gate};
\item if the heat trace is order \(k\), return to the un-factorized
      pair-level minimum rather than assuming an internal-pair
      inverse-\(k\) debit.
\end{enumerate}
The next compulsory companion audit is V90.
\end{openproblem}

\section{Firewall after V88}
\label{sec:V88-firewall}

\begin{firewall}[Current claim boundary]
\label{fire:V88-final}
V88 proves the exact balanced multiplicity and reduces the
factorized heavy-pair thermal estimate to a normalized core trace
lying rigorously between a positive constant and \(C\sqrt m\).  It
does not determine the ancestry-averaged heat trace, close the
un-factorized pair-level minimum, treat the off-cut alternative,
synchronize both orientations, prove JREC, construct the continuum
Yang--Mills theory, or prove a positive mass gap.
\end{firewall}

\section{Append-only record for V88}
\label{sec:V88-record}

V88 was appended to the complete V87 manuscript.  The exact V87
parent had \(60494\) validation lines, \(2135510\) bytes, and
SHA--256 digest
\[
 \texttt{bb409198a9c20bc54c5f586e40b805e2be3ea5469414c7f56abad9c883fc4914}.
\]
No V87 mathematical content was changed.  The sole final
\verb|\end{document}| is restored below.

% =============================================================================
% END APPENDED VERSION V88 MATERIAL
% =============================================================================

% =============================================================================
% BEGIN APPENDED VERSION V89 MATERIAL
% =============================================================================

\chapter{V89: Full-Ancestry Heat Scalarization and the
Factorized-Trace No-Go}
\label{ch:V89}

\section{Purpose}
\label{sec:V89-purpose}

V88 isolated a possible factor \(\sqrt m\): the balanced output
\(T_k=(k,k)\) has a \((k+1)\)-dimensional ancestry space, and the
available trace bound paid every ancestry direction.  The next
question was whether this was merely a rough use of the operator norm.
This chapter answers that question exactly for the heat benchmark.

The answer is negative.  On the entire multiplicity space, each of the
three pair heat operators is uniformly close to a scalar at
\(\tau m^2\to u\in(0,\infty)\) and \(k\asymp\sqrt m\).  The assembled
connected operator therefore converges in operator norm to
\(\mathfrak d(u)I\), where the V83 transition profile
\(\mathfrak d(u)\) is strictly positive.  Consequently the mean square
per ancestry direction tends to \(\mathfrak d(u)^2\), and the V88
thermal trace is genuinely of order \(\sqrt m\).

This closes the ancestry-average question without computing Racah
coefficients or separating signed cross traces.  It also rules out the
factorized trace route unless an independent inverse-\(\sqrt m\)
internal-pair debit is proved.  By the V86 spectator-inflation no-go,
such a debit cannot be inferred from the core label alone.  The next
step must therefore retain the pair level inside the V88 minimum.

\section{The missing \(AB\)-first Pieri support}
\label{sec:V89-AB-support}

Continue with
\[
 A=B=R_m=(m,0),\qquad
 C=Q_{2m}=(0,2m),\qquad
 T_k=(k,k),\qquad 0\le k\le m.
\]
The \(AB\) Pieri decomposition is
\begin{equation}
 R_m\otimes R_m
 =
 \bigoplus_{r=0}^{m}S_r,
 \qquad
 S_r:=(2m-2r,r).
 \label{eq:V89-AB-Pieri}
\end{equation}

\begin{theorem}[Exact \(AB\)-first thermal support]
\label{thm:V89-AB-support}
For \(0\le k\le m\),
\begin{equation}
 N_{S_r,Q_{2m}}^{T_k}
 =
 \begin{cases}
 1,&0\le r\le k,\\
 0,&k<r\le m.
 \end{cases}
 \label{eq:V89-AB-support}
\end{equation}
Thus the \(AB\)-first paths reaching \(T_k\) are exactly
\(S_0,\ldots,S_k\), each with multiplicity one.
\end{theorem}

\begin{proof}
Frobenius reciprocity and \(Q_{2m}^*=R_{2m}\) give
\[
 N_{S_r,Q_{2m}}^{T_k}
 =
 N_{T_k,R_{2m}}^{S_r}.
\]
Use the \(U(3)\) partition representatives
\[
 T_k\longleftrightarrow(2k,k,0),
 \qquad
 R_{2m}\longleftrightarrow(2m,0,0).
\]
The \(SU(3)\) representation \(S_r=(2m-2r,r)\) has canonical
partition \((2m-r,r,0)\).  To have the same total number
\(3k+2m\) of boxes as the product, its only possible \(U(3)\)
representative is obtained by adding \(k\) determinant columns:
\[
 \lambda_{m,k,r}
 =
 (2m-r+k,r+k,k).
\]
Pieri's rule for multiplication by the one-row partition \((2m)\)
says that \(\lambda_{m,k,r}/(2k,k,0)\) is a horizontal strip if and
only if
\[
 \lambda_1\ge2k\ge\lambda_2\ge k
 \ge\lambda_3\ge0.
\]
For the displayed \(\lambda_{m,k,r}\), the only nonautomatic
condition is \(2k\ge r+k\), which is exactly \(r\le k\).
The remaining first-row condition
\(2m-r+k\ge2k\) follows from \(r\le k\le m\).
Pieri multiplicities are one, proving
\eqref{eq:V89-AB-support}.
\end{proof}

\begin{corollary}[Three multiplicity-one Racah spectra]
\label{cor:V89-three-Racah-spectra}
On the multiplicity space
\(\mathcal M_{T_k}\), of dimension \(k+1\), the three pair
intermediate supports are
\begin{align}
 AB &: S_r=(2m-2r,r), &&0\le r\le k,
 \label{eq:V89-AB-nodes}\\
 AC &: U_r=(k-r,m+k-r), &&0\le r\le k,
 \label{eq:V89-AC-nodes}\\
 BC &: U_r=(k-r,m+k-r), &&0\le r\le k.
 \label{eq:V89-BC-nodes}
\end{align}
Every listed path occurs once.
\end{corollary}

\begin{proof}
The \(AB\) assertion is
\Cref{thm:V89-AB-support}.  The \(AC\) assertion is the exact V81
support used in \Cref{thm:V88-balanced-multiplicity}, after writing
the intermediate index as \(j=m-k+r\).  The \(BC\) assertion follows
by exchanging the equal inputs \(A\) and \(B\).
\end{proof}

\section{Uniform spectral collapse of all three heat pair operators}
\label{sec:V89-spectral-collapse}

Write
\[
 h_{\tau,R}:=e^{-\tau c_R},
 \qquad h_m:=e^{-\tau c_m},
 \qquad h_{2m}:=e^{-\tau c_{2m}},
 \qquad h_T:=e^{-\tau c_{T_k}}.
\]
Let \(Q_{AB}^{\mathsf h,T_k}\),
\(Q_{AC}^{\mathsf h,T_k}\), and
\(Q_{BC}^{\mathsf h,T_k}\) denote the heat pair-moment operators on
\(\mathcal M_{T_k}\).

\begin{lemma}[Exact \(AB\) Casimir segment]
\label{lem:V89-AB-Casimir}
For \(S_r=(2m-2r,r)\),
\begin{align}
 c_{S_r}
 &=
 c_{2m}-r(2m-r+1),
 \label{eq:V89-AB-Casimir}\\
 0\le c_{2m}-c_{S_r}
 &\le
 B_{m,k}:=k(2m-k+1),
 \qquad 0\le r\le k.
 \label{eq:V89-AB-diameter}
\end{align}
\end{lemma}

\begin{proof}
Insert \(p=2m-2r\) and \(q=r\) in
\[
 C_2(p,q)=\frac{p^2+q^2+pq+3p+3q}{3}.
\]
This gives
\[
 c_{S_r}
 =
 \frac{4m^2}{3}-2mr+r^2+2m-r
 =
 c_{2m}-r(2m-r+1).
\]
The function \(r(2m-r+1)\) is increasing on \(0\le r\le m\), so
its maximum on \(0\le r\le k\) is \(B_{m,k}\).
\end{proof}

\begin{theorem}[Three operator-norm scalarization estimates]
\label{thm:V89-three-pair-scalarization}
For every \(\tau\ge0\) and \(0\le k\le m\),
\begin{align}
 \left\|Q_{AB}^{\mathsf h,T_k}-h_{2m}I\right\|_{\mathrm{op}}
 &\le \tau B_{m,k},
 \label{eq:V89-AB-scalarization}\\
 \left\|Q_{AC}^{\mathsf h,T_k}-h_mI\right\|_{\mathrm{op}}
 &\le \tau L_{m,k},
 \label{eq:V89-AC-scalarization}\\
 \left\|Q_{BC}^{\mathsf h,T_k}-h_mI\right\|_{\mathrm{op}}
 &\le \tau L_{m,k},
 \label{eq:V89-BC-scalarization}
\end{align}
where \(L_{m,k}=k(m+k+2)\) is the V83 \(AC\)-node diameter.
\end{theorem}

\begin{proof}
In the \(AB\)-first Racah basis,
\[
 Q_{AB}^{\mathsf h,T_k}
 =
 \sum_{r=0}^{k}e^{-\tau c_{S_r}}P_{AB}(S_r),
\]
where the \(P_{AB}(S_r)\) are mutually orthogonal rank-one
projections summing to the identity.  Since
\[
 |e^{-\tau x}-e^{-\tau y}|\le\tau|x-y|,
 \qquad x,y\ge0,
\]
\Cref{lem:V89-AB-Casimir} proves
\eqref{eq:V89-AB-scalarization}.

In the \(AC\)-first Racah basis, the eigenvalues are
\(e^{-\tau c_{U_r}}\), \(0\le r\le k\).  Equations
\eqref{eq:V82-Casimir-nodes} and
\eqref{eq:V82-node-diameter} show
\[
 c_m=c_{U_k}\le c_{U_r}\le c_m+L_{m,k}.
\]
The same Lipschitz bound proves
\eqref{eq:V89-AC-scalarization}.  Exchange \(A\) and \(B\) for
\eqref{eq:V89-BC-scalarization}.  No comparison of the three Racah
bases is needed because each estimate is basis invariant.
\end{proof}

\begin{corollary}[Uniform collapse on the thermal window]
\label{cor:V89-thermal-pair-collapse}
Suppose
\[
 \tau_jm_j^2\longrightarrow u\in(0,\infty),
 \qquad m_j\longrightarrow\infty,
\]
and fix \(0<A_0<B_0<\infty\).  Uniformly for
\[
 A_0\sqrt{2m_j}\le k\le B_0\sqrt{2m_j},
\]
all three right sides in
\eqref{eq:V89-AB-scalarization}--\eqref{eq:V89-BC-scalarization}
converge to zero.
\end{corollary}

\begin{proof}
On the stated window,
\[
 B_{m,k}=O(m^{3/2}),
 \qquad L_{m,k}=O(m^{3/2}).
\]
Since \(\tau m^2\) remains bounded,
\[
 \tau(B_{m,k}+L_{m,k})=O(m^{-1/2})
\]
uniformly in \(k\).
\end{proof}

\section{The full connected heat operator}
\label{sec:V89-full-heat-operator}

The exact five-term connected operator on \(\mathcal M_{T_k}\) is
\begin{equation}
 K_{\tau;m,k}^{\mathsf h}
 :=
 h_T I
 -h_mQ_{BC}^{\mathsf h,T_k}
 -h_mQ_{AC}^{\mathsf h,T_k}
 -h_{2m}Q_{AB}^{\mathsf h,T_k}
 +2h_m^2h_{2m}I.
 \label{eq:V89-heat-connected-operator}
\end{equation}
Define its endpoint scalar reference
\begin{equation}
 \mathfrak s_{\tau;m,k}
 :=
 h_T-h_{2m}^2-2h_m^2+2h_m^2h_{2m}.
 \label{eq:V89-endpoint-scalar}
\end{equation}
This finite-\(m\) scalar is not the V83 tangent-mean scalar
\(\mathcal D_{\tau;m,k}\); the two have the same transition limit.

\begin{theorem}[Full-ancestry heat scalarization]
\label{thm:V89-full-heat-scalarization}
For every \(\tau\ge0\) and \(0\le k\le m\),
\begin{equation}
 \boxed{
 \left\|
 K_{\tau;m,k}^{\mathsf h}
 -\mathfrak s_{\tau;m,k}I
 \right\|_{\mathrm{op}}
 \le
 \tau\bigl(B_{m,k}+2L_{m,k}\bigr).}
 \label{eq:V89-full-heat-scalarization}
\end{equation}
Under the transition scaling of
\Cref{cor:V89-thermal-pair-collapse},
\begin{equation}
 \boxed{
 \sup_{A_0\sqrt{2m_j}\le k\le B_0\sqrt{2m_j}}
 \left\|
 K_{\tau_j;m_j,k}^{\mathsf h}
 -\mathfrak d(u)I
 \right\|_{\mathrm{op}}
 \longrightarrow0.}
 \label{eq:V89-full-heat-limit}
\end{equation}
\end{theorem}

\begin{proof}
Subtract \eqref{eq:V89-endpoint-scalar} times the identity from
\eqref{eq:V89-heat-connected-operator}.  The result is
\[
 -h_m(Q_{BC}^{\mathsf h,T_k}-h_mI)
 -h_m(Q_{AC}^{\mathsf h,T_k}-h_mI)
 -h_{2m}(Q_{AB}^{\mathsf h,T_k}-h_{2m}I).
\]
All heat multipliers lie in \([0,1]\).  The triangle inequality and
\Cref{thm:V89-three-pair-scalarization} prove
\eqref{eq:V89-full-heat-scalarization}.

On the thermal window,
\[
 \tau_jc_{T_k}\longrightarrow0,\qquad
 \tau_jc_{m_j}\longrightarrow\frac u3,\qquad
 \tau_jc_{2m_j}\longrightarrow\frac{4u}{3}
\]
uniformly in \(k\).  Therefore
\[
 \mathfrak s_{\tau_j;m_j,k}
 \longrightarrow
 1-e^{-8u/3}-2e^{-2u/3}+2e^{-2u}
 =\mathfrak d(u)
\]
uniformly.  The operator error tends to zero by
\Cref{cor:V89-thermal-pair-collapse}, proving
\eqref{eq:V89-full-heat-limit}.
\end{proof}

\begin{corollary}[Uniform positivity on every ancestry direction]
\label{cor:V89-uniform-ancestry-positivity}
Under the hypotheses of
\Cref{thm:V89-full-heat-scalarization}, for all sufficiently large
\(j\), uniformly on the thermal \(k\)-window,
\begin{equation}
 K_{\tau_j;m_j,k}^{\mathsf h}
 \ge\frac{\mathfrak d(u)}2 I,
 \qquad
 \frac1{k+1}
 \operatorname{Tr}_{\mathcal M_{T_k}}
 \left(K_{\tau_j;m_j,k}^{\mathsf h}\right)^2
 \longrightarrow\mathfrak d(u)^2.
 \label{eq:V89-mean-square-limit}
\end{equation}
\end{corollary}

\begin{proof}
The heat connected operator is self-adjoint.  Since
\(\mathfrak d(u)>0\) by
\eqref{eq:V83-transition-profile-factor}, the first assertion follows
from \eqref{eq:V89-full-heat-limit}.  For the second, use
\[
 \|K^2-\mathfrak d(u)^2I\|_{\mathrm{op}}
 \le
 \|K-\mathfrak d(u)I\|_{\mathrm{op}}
 \bigl(\|K\|_{\mathrm{op}}+\mathfrak d(u)\bigr).
\]
The second factor is uniformly bounded, and normalized trace is at
most operator norm in absolute value.
\end{proof}

\section{Exact thermal trace saturation}
\label{sec:V89-trace-saturation}

Define the unweighted heat trace on the thermal window by
\begin{equation}
 \mathfrak T_{\tau;m}^{\mathsf h,0;A_0,B_0}
 :=
 \sum_{A_0\sqrt{2m}\le k\le B_0\sqrt{2m}}
 \frac{d_{T_k}}{d_m}
 \operatorname{Tr}_{\mathcal M_{T_k}}
 \left(K_{\tau;m,k}^{\mathsf h}\right)^2.
 \label{eq:V89-unweighted-heat-trace}
\end{equation}

\begin{theorem}[The V88 square-root trace loss is sharp]
\label{thm:V89-heat-trace-saturation}
Under the transition scaling
\(\tau_jm_j^2\to u\in(0,\infty)\),
\begin{equation}
 \boxed{
 \frac1{\sqrt{m_j}}
 \mathfrak T_{\tau_j;m_j}^{\mathsf h,0;A_0,B_0}
 \longrightarrow
 \frac{8\sqrt2}{5}
 (B_0^5-A_0^5)\mathfrak d(u)^2.}
 \label{eq:V89-heat-trace-saturation}
\end{equation}
In particular, the mean square per ancestry direction is not
\(O(k^{-1})\); it has the strictly positive limit
\(\mathfrak d(u)^2\).
\end{theorem}

\begin{proof}
By \Cref{cor:V89-uniform-ancestry-positivity},
\[
 \operatorname{Tr}_{\mathcal M_{T_k}}
 \left(K_{\tau;m,k}^{\mathsf h}\right)^2
 =
 (k+1)\bigl[\mathfrak d(u)^2+o(1)\bigr]
\]
uniformly on the thermal window.  Since
\(d_{T_k}=(k+1)^3\) and \(d_m=(m+1)(m+2)/2\),
\begin{align*}
 &\frac1{\sqrt m}
 \sum_{A_0\sqrt{2m}\le k\le B_0\sqrt{2m}}
 \frac{(k+1)^4}{d_m}\\
 &\qquad\longrightarrow
 8\sqrt2\int_{A_0}^{B_0}t^4\,\dd t
 =
 \frac{8\sqrt2}{5}(B_0^5-A_0^5).
\end{align*}
Indeed, after writing \(n=2m\), one has
\((k+1)^4/d_m=8(k/\sqrt n)^4+o(1)\) uniformly, and
the mesh in \(t=k/\sqrt n\) is \(n^{-1/2}\).
Multiplication by the uniform mean-square limit proves
\eqref{eq:V89-heat-trace-saturation}.
\end{proof}

\begin{corollary}[Weighted heat trace cannot restore the average gain]
\label{cor:V89-weighted-heat-lower}
If the V88 weighted block is formed with
\[
 \widetilde K_{\tau;m,k}^{\mathsf h}
 =
 [1+\tau\Xi(T_k)]K_{\tau;m,k}^{\mathsf h},
\]
then its normalized thermal trace satisfies
\begin{equation}
 \liminf_{j\to\infty}
 \frac1{\sqrt{m_j}}
 \sum_k
 \frac{d_{T_k}}{d_{m_j}}
 \operatorname{Tr}
 \left(\widetilde K_{\tau_j;m_j,k}^{\mathsf h}\right)^2
 \ge
 \frac{8\sqrt2}{5}
 (B_0^5-A_0^5)\mathfrak d(u)^2.
 \label{eq:V89-weighted-heat-lower}
\end{equation}
The sum is over the same thermal window.
\end{corollary}

\begin{proof}
The scalar weight \(1+\tau\Xi(T_k)\) is at least one.  Hence every
positive trace summand dominates its unweighted counterpart in
\eqref{eq:V89-unweighted-heat-trace}.  Apply
\Cref{thm:V89-heat-trace-saturation}.
\end{proof}

\section{Conditional full-space Wilson transfer}
\label{sec:V89-Wilson-transfer}

The V83 bank controlled the distinguished Cartan vector.  Full-space
transfer also requires the \(AB\)-first nodes.  Define
\begin{align}
 \widehat{\mathcal B}_{m,k}
 &:=
 \mathcal B_{m,k}\cup\{S_0,\ldots,S_k\},
 \label{eq:V89-full-comparison-bank}\\
 \widehat\delta_{\alpha;m,k}(\tau)
 &:=
 \max_{R\in\widehat{\mathcal B}_{m,k}}
 |q_{\alpha,R}-e^{-\tau c_R}|.
 \label{eq:V89-full-discrepancy}
\end{align}

\begin{proposition}[Full multiplicity-block transfer]
\label{prop:V89-full-block-transfer}
Let \(K_{\alpha;m,k}\) be the Wilson connected operator on
\(\mathcal M_{T_k}\), and let
\(\delta=\widehat\delta_{\alpha;m,k}(\tau)\).  Then
\begin{align}
 \|Q_{AB}^{\alpha,T_k}-Q_{AB}^{\mathsf h,T_k}\|_{\mathrm{op}}
 &\le\delta,
 \label{eq:V89-AB-operator-transfer}\\
 \|Q_{AC}^{\alpha,T_k}-Q_{AC}^{\mathsf h,T_k}\|_{\mathrm{op}}
 &\le\delta,
 \label{eq:V89-AC-operator-transfer}\\
 \|Q_{BC}^{\alpha,T_k}-Q_{BC}^{\mathsf h,T_k}\|_{\mathrm{op}}
 &\le\delta,
 \label{eq:V89-BC-operator-transfer}\\
 \boxed{
 \|K_{\alpha;m,k}-K_{\tau;m,k}^{\mathsf h}\|_{\mathrm{op}}
 \le13\delta.}
 \label{eq:V89-K-operator-transfer}
\end{align}
\end{proposition}

\begin{proof}
In each pair Racah basis, the Wilson and heat operators have the same
spectral projections.  Their eigenvalue difference on every support
node is at most \(\delta\).  This proves the first three estimates.

Use the balanced five-term identity
\[
 K_\alpha
 =
 q_{\alpha,T_k}I
 -q_{\alpha,R_m}Q_{BC}^{\alpha,T_k}
 -q_{\alpha,R_m}Q_{AC}^{\alpha,T_k}
 -q_{\alpha,R_{2m}}Q_{AB}^{\alpha,T_k}
 +2q_{\alpha,R_m}^2q_{\alpha,R_{2m}}I.
\]
All scalar multipliers and pair operators are contractions.  The
output term costs \(\delta\), the two \(AC,BC\) terms cost at most
\(2\delta\) each, the \(AB\) term costs at most \(2\delta\), and
\[
 2|q_m^2q_{2m}-h_m^2h_{2m}|\le6\delta.
\]
The total is \(13\delta\).
\end{proof}

\begin{theorem}[Conditional Wilson trace saturation]
\label{thm:V89-conditional-Wilson-saturation}
Let \(s_{\alpha_j}m_j^2\to u\in(0,\infty)\).  Suppose
\begin{equation}
 \sup_{A_0\sqrt{2m_j}\le k\le B_0\sqrt{2m_j}}
 \widehat\delta_{\alpha_j;m_j,k}(s_{\alpha_j})
 \longrightarrow0.
 \label{eq:V89-full-comparison-hypothesis}
\end{equation}
Then, uniformly on the thermal window,
\begin{equation}
 \|K_{\alpha_j;m_j,k}-\mathfrak d(u)I\|_{\mathrm{op}}
 \longrightarrow0.
 \label{eq:V89-Wilson-full-limit}
\end{equation}
Consequently, with
\begin{equation}
 \mathfrak T_{\alpha;m}^{0;A_0,B_0}
 :=
 \sum_k\frac{d_{T_k}}{d_m}
 \operatorname{Tr}_{\mathcal M_{T_k}}K_{\alpha;m,k}^2,
 \label{eq:V89-Wilson-unweighted-trace}
\end{equation}
one has
\begin{equation}
 \boxed{
 \frac1{\sqrt{m_j}}
 \mathfrak T_{\alpha_j;m_j}^{0;A_0,B_0}
 \longrightarrow
 \frac{8\sqrt2}{5}
 (B_0^5-A_0^5)\mathfrak d(u)^2.}
 \label{eq:V89-Wilson-trace-saturation}
\end{equation}
The weighted V88 trace therefore obeys the same expression as a
strictly positive lower limit.
\end{theorem}

\begin{proof}
Combine \eqref{eq:V89-K-operator-transfer},
\eqref{eq:V89-full-comparison-hypothesis}, and
\eqref{eq:V89-full-heat-limit}.  The normalized trace-square
convergence follows exactly as in
\Cref{cor:V89-uniform-ancestry-positivity}.  The Riemann sum in
\Cref{thm:V89-heat-trace-saturation} then gives
\eqref{eq:V89-Wilson-trace-saturation}.  Finally,
\[
 \operatorname{Tr}\widetilde K_{T_k}^2
 =
 [1+s_\alpha\Xi(T_k)]^2\operatorname{Tr}K_{T_k}^2
 \ge \operatorname{Tr}K_{T_k}^2
\]
because the scalar weight is at least one.
\end{proof}

\begin{firewall}[The full-bank comparison remains conditional]
\label{fire:V89-full-comparison}
The new hypothesis
\eqref{eq:V89-full-comparison-hypothesis} is stronger than the V83
restricted comparison because it includes the \(AB\)-first segment
\(S_0,\ldots,S_k\).  Neither comparison is presently proved for the
one-link Wilson measure.  The total-variation estimate
\eqref{eq:V83-TV-multiplier-bound} would imply both, but the required
transition-scale total-variation convergence is also unproved.
\end{firewall}

\section{Consequence for the heavy-pair route}
\label{sec:V89-heavy-pair-consequence}

\begin{theorem}[No ancestry-average gain in the factorized route]
\label{thm:V89-factorized-route-no-go}
For the exact heat benchmark, the core trace factor appearing in
\eqref{eq:V88-balanced-heavy-pair-gate} is bounded below by a positive
constant times \(\sqrt m\) in the transition window.  Under
\eqref{eq:V89-full-comparison-hypothesis}, the same is true for the
Wilson block.  Therefore a uniform bound obtained by the factorized
estimate
\[
 \mathcal K_e\,
 \mathfrak T_{\alpha;m}^{A_0,B_0}
\]
would require an independent estimate
\(\mathcal K_e=O(m^{-1/2})\) on this family.
\end{theorem}

\begin{proof}
The heat assertion is
\eqref{eq:V89-weighted-heat-lower}.  The Wilson assertion is
\Cref{thm:V89-conditional-Wilson-saturation}.  Multiplying a quantity
bounded below by \(c\sqrt m\) by \(\mathcal K_e\) can be uniformly
bounded through this factorized estimate only if
\(\mathcal K_e\sqrt m=O(1)\).
\end{proof}

\begin{assessment}[Why this sends the proof upstream]
\label{ass:V89-return-to-minimum}
The conclusion of
\Cref{thm:V89-factorized-route-no-go} is a statement about what the
factorized \emph{upper bound} would need; it is not a lower bound on
the original pair--core product norm.  One must not reverse an upper
bound.

V86 already showed that a spectator may inflate the core label \(m\)
without changing the internal-pair layer distribution.  Thus the
required inverse-\(\sqrt m\) capacity cannot be deduced from the core
label alone.  The remaining legitimate mechanism is the
un-factorized integral
\[
 \int_0^{A_*}
 \min\left\{
 \Pi_j^{\rm core},
 \frac{D_e(t)\mathcal H_j}
 {\max\{d_Xd_A,d_Yd_B\}}
 \right\}\,\dd t,
\]
summed before separating its pair and core variables.  It can exploit
the threshold at which the pair rank \(D_e(t)\) meets the
ancestry-resolved core product cap; the factorized capacity
\(\mathcal K_e\) discards that threshold.
\end{assessment}

\section{Append-only correction to one V83 display}
\label{sec:V89-V83-display-correction}

In \eqref{eq:V83-Wilson-heat-upper-transfer}, the two displayed
squares are to be joined by a plus sign:
\begin{equation}
 \|\boldsymbol\kappa_{m,m;2m,k}^{\alpha}\|^2
 \le
 \left(
 |\mathcal D_{\tau;m,k}|
 +2\mathcal E_{\tau;m,k}
 +13\delta
 \right)^2
 +4\left(\mathcal E_{\tau;m,k}+\delta\right)^2.
 \label{eq:V89-corrected-V83-transfer}
\end{equation}
This is the sum produced by the scalar and even-leakage terms in the
proof of \Cref{prop:V83-finite-bank-transfer}; no later argument used
the accidental juxtaposition in the earlier display.

\section{V89 ledger and V90 mandate}
\label{sec:V89-ledger}

\begin{theorem}[Integrated V89 statement]
\label{thm:V89-integrated}
Preserving V1--V88, V89 proves:
\begin{enumerate}[label=\textup{(V89.\arabic*)},leftmargin=5.7em]
\item the exact \(AB\)-first support
      \eqref{eq:V89-AB-support};
\item operator-norm scalarization of all three heat pair operators
      \eqref{eq:V89-AB-scalarization}--\eqref{eq:V89-BC-scalarization};
\item full connected heat scalarization
      \eqref{eq:V89-full-heat-limit};
\item the positive mean square on every ancestry direction
      \eqref{eq:V89-mean-square-limit};
\item the exact \(\sqrt m\) heat-trace asymptotic
      \eqref{eq:V89-heat-trace-saturation};
\item the full-bank Wilson operator transfer
      \eqref{eq:V89-K-operator-transfer};
\item conditional Wilson trace saturation
      \eqref{eq:V89-Wilson-trace-saturation}; and
\item the factorized-trace no-go and the resulting return to the
      un-factorized pair-level minimum.
\end{enumerate}
\end{theorem}

\begin{proof}
These are
\Cref{thm:V89-AB-support,
thm:V89-three-pair-scalarization,
thm:V89-full-heat-scalarization,
cor:V89-uniform-ancestry-positivity,
thm:V89-heat-trace-saturation,
prop:V89-full-block-transfer,
thm:V89-conditional-Wilson-saturation,
thm:V89-factorized-route-no-go}.
\end{proof}

\begin{openproblem}[V90 mandate: audited level/core quantile]
\label{op:V90-mandate}
Before the V90 mathematics, perform the compulsory five-version audit
of the current Einstein--Standard-Model and Navier--Stokes notebooks.
Then return to the un-factorized V88 minimum and:
\begin{enumerate}
\item preserve the internal-pair level variable \(t\) while resolving
      the core multiplicity space into \(AB\)-first ancestry paths;
\item derive the exact threshold/quantile formula for
      \(\int\min\{\Pi_j^{\rm core},R_e(t)B_j\}\,\dd t\);
\item determine which threshold parameter is actually tied to the
      V59--V60 heavy-puncture ledger, rather than to the spectator core
      label \(m\);
\item use the V89 scalarization only as a test family: it rules out a
      free ancestry average but does not rule out level/core
      correlation;
\item if the level and ancestry variables again factor, identify the
      exact missing joint estimate and move upstream to the selected
      heavy-path construction rather than postulating it.
\end{enumerate}
\end{openproblem}

\section{Firewall after V89}
\label{sec:V89-firewall}

\begin{firewall}[Current claim boundary]
\label{fire:V89-final}
V89 determines the full balanced heat ancestry trace and proves that
the factor \(\sqrt m\) in the V88 factorized trace bound is sharp.  It
does not prove the full-bank Wilson-to-heat comparison, close the
un-factorized pair-level minimum, control the off-cut alternative,
synchronize both orientations, prove JREC, construct the continuum
Yang--Mills theory, or prove a positive mass gap.
\end{firewall}

\section{Append-only record for V89}
\label{sec:V89-record}

V89 was appended to the complete V88 manuscript.  The exact V88
parent had \(60975\) physical lines, \(2150040\) bytes, and
SHA--256 digest
\[
 \texttt{4b76be6cb95e3c194528c333eb3467adfd35ccf00378213e9d05eb8ba1ab0744}.
\]
No V88 mathematical content was changed.  The sole final
\verb|\end{document}| is restored below.

% =============================================================================
% END APPENDED VERSION V89 MATERIAL
% =============================================================================

% =============================================================================
% BEGIN APPENDED VERSION V90 MATERIAL
% =============================================================================

\chapter{V90: Audited Product-Carrier Capacity and Uniform Closure of
the Balanced Thermal Pair--Core Shell}
\label{ch:V90}

\section{Purpose and compulsory companion audit}
\label{sec:V90-purpose-audit}

V89 proved that the factorized trace route genuinely pays all
\(k+1\) balanced ancestry directions.  That result forces a return to
the first branch of the V88 minimum, namely the core product norm
\(\Pi_j^{\rm core}\).  The compulsory V90 audit supports exactly this
move: both companion programmes distinguish a restricted joint object
from the larger marginal norm obtained after forgetting its graph,
stopping, or common-domain structure.

\begin{audit}[V90 companion audit]
\label{audit:V90-companions}
The current companion files were read without modification.
\begin{enumerate}[label=\textup{(\roman*)},leftmargin=3.2em]
\item The Einstein--Standard-Model notebook was
\[
 \texttt{Renewal\_SM\_Einstein\_Continuum\_Research\_Notes\_V90.tex},
\]
with \(69939\) physical lines, \(2543677\) bytes, and SHA--256
\[
 \texttt{e396b0f4a3679439afc8691878a407c3ed00fba7c0325d13941eb1b8eabfa835}.
\]
Its V87--V90 additions prove that a genuine quotient conormal cannot
be controlled in the stronger canonical completion, replace that
false comparison by a rigged port graph, and require one common
renormalized form domain before passing a Ward cancellation to the
limit.  In particular, finite marginal sector bounds do not prove the
common gate.
\item The Navier--Stokes notebook remained
\[
 \texttt{Renewal\_NS\_Stability\_Research\_Notes\_V31.tex},
\]
with \(23052\) physical lines, \(887537\) bytes, and SHA--256
\[
 \texttt{f1121b62238ad5491bb7cf03635ea9934942edf573ed8faaa9ee79e1d38a6696}.
\]
\end{enumerate}
\end{audit}

The transferable NS conclusion is the exact flat-mark lesson:
removing a probability mark restores its first moment, while the
sharper remaining object is the stopping-time-correlated formation
occupation.  The transferable SM conclusion is likewise structural:
one must retain the common graph or form domain through the
cancellation.  No SM or NS inequality is imported into Yang--Mills.
Here the corresponding retained object is the product norm of the
partially traced \(T_k\)-carrier, before replacing it by the full
multiplicity trace.

\section{The partially traced balanced output carrier}
\label{sec:V90-output-carrier}

Continue with the balanced family
\[
 A=B=R_m,\qquad C=Q_{2m},\qquad T_k=(k,k).
\]
Let \(\mathbf P_{m,k}\) be the full \(T_k\)-isotypic projector in
\(A\otimes B\otimes C\), including all \(k+1\) multiplicity copies,
and define
\begin{equation}
 \mathbf N_{m,k}
 :=
 \operatorname{Tr}_{C}\mathbf P_{m,k}
 \quad\hbox{on }H_A\otimes H_B.
 \label{eq:V90-partial-output-carrier}
\end{equation}

\begin{lemma}[Exact pathwise partial trace]
\label{lem:V90-path-partial-trace}
Let \(\mathbf P_{m,k;r}^{AB}\) be the \(T_k\) path projector through
\(S_r=(2m-2r,r)\), \(0\le r\le k\).  Then
\begin{equation}
 \operatorname{Tr}_{C}\mathbf P_{m,k;r}^{AB}
 =
 \frac{d_{T_k}}{d_{S_r}}P_{AB}(S_r).
 \label{eq:V90-path-partial-trace}
\end{equation}
Consequently
\begin{equation}
 \boxed{
 \mathbf N_{m,k}
 =
 \sum_{r=0}^{k}
 \frac{d_{T_k}}{d_{S_r}}P_{AB}(S_r).}
 \label{eq:V90-output-carrier-spectrum}
\end{equation}
\end{lemma}

\begin{proof}
By \Cref{thm:V89-AB-support}, both arrows in the path
\[
 A\otimes B\longrightarrow S_r,\qquad
 S_r\otimes C\longrightarrow T_k
\]
have multiplicity one.  The partial trace of the path projector
commutes with the irreducible \(S_r\)-action and vanishes on every
other \(AB\)-summand.  Schur's lemma therefore makes it
\(\lambda_rP_{AB}(S_r)\).  The path projector has rank \(d_{T_k}\);
equality of traces gives
\(\lambda_rd_{S_r}=d_{T_k}\), proving
\eqref{eq:V90-path-partial-trace}.  Sum the mutually orthogonal path
projectors.
\end{proof}

\begin{lemma}[Dimension monotonicity on the thermal \(AB\) segment]
\label{lem:V90-S-dimension-monotonicity}
If \(m\ge2\) and \(0\le r\le m/2\), then
\begin{equation}
 d_{S_r}\ge d_{S_0}=d_{2m}.
 \label{eq:V90-S-dimension-monotonicity}
\end{equation}
\end{lemma}

\begin{proof}
The \(SU(3)\) dimension formula gives
\[
 d_{S_r}
 =
 \frac12(2m-2r+1)(r+1)(2m-r+2).
\]
Direct subtraction and factorization give
\begin{align*}
 2(d_{S_r}-d_{S_0})
 &=
 r\left[
 2r^2-(6m+3)r+4m^2-3
 \right].
\end{align*}
The bracket is decreasing on \(0\le r\le m/2\), and at \(r=m/2\)
it equals
\[
 \frac32(m-2)(m+1)\ge0.
\]
This proves the claim.
\end{proof}

\begin{theorem}[Exact balanced product-carrier norm]
\label{thm:V90-balanced-product-carrier}
For \(m\ge2\) and \(0\le k\le m/2\),
\begin{equation}
 \boxed{
 \|\mathbf N_{m,k}\|_{A\otimes B,\mathrm{prod}}
 =
 \frac{d_{T_k}}{d_{2m}}.}
 \label{eq:V90-balanced-product-carrier}
\end{equation}
By contrast,
\begin{equation}
 \operatorname{Tr}\mathbf N_{m,k}
 =
 d_{T_k}(k+1).
 \label{eq:V90-output-carrier-trace}
\end{equation}
Thus the normalized trace contains exactly one additional factor
\(k+1\) which the product-carrier norm does not contain.
\end{theorem}

\begin{proof}
Equations \eqref{eq:V90-output-carrier-spectrum} and
\eqref{eq:V90-S-dimension-monotonicity} give
\[
 \|\mathbf N_{m,k}\|_{\mathrm{op}}
 =
 \frac{d_{T_k}}{d_{2m}}.
\]
The product-state norm is at most the operator norm.  Equality is
attained by the tensor product of highest-weight vectors in
\(R_m\otimes R_m\): this vector belongs entirely to the Cartan
summand \(S_0=R_{2m}\).  This proves
\eqref{eq:V90-balanced-product-carrier}.

For the trace, either sum the traces of
\eqref{eq:V90-path-partial-trace}, each of which is \(d_{T_k}\), or
use preservation of trace under \(\operatorname{Tr}_C\) and the rank
\(d_{T_k}(k+1)\) of the full isotypic projector.
\end{proof}

\section{A general product-carrier bound before multiplicity trace}
\label{sec:V90-general-product-carrier}

The preceding geometry does not require the connected operator to be
diagonal in the \(AB\)-first path basis.  Let \(T\) be any final
output, let \(\mathbf P_T\) be its complete isotypic projector, and
let \(L_T=L_T^*\) be an arbitrary operator on its multiplicity space.
Put
\begin{align}
 \mathbf N_T&:=\operatorname{Tr}_C\mathbf P_T,
 \label{eq:V90-general-output-carrier}\\
 M_T[L]&:=
 \operatorname{Tr}_C\left[
 I_{H_T}\otimes L_T^2
 \right],
 \label{eq:V90-general-core-square}
\end{align}
where the operator in the second line is extended by zero away from
the \(T\)-isotypic block.

\begin{theorem}[Product-carrier domination]
\label{thm:V90-product-carrier-domination}
For every such \(T\) and \(L_T\),
\begin{equation}
 0\le M_T[L]\le
 \|L_T\|_{\mathrm{op}}^2\mathbf N_T,
 \label{eq:V90-product-carrier-order}
\end{equation}
and hence
\begin{equation}
 \boxed{
 \|M_T[L]\|_{A\otimes B,\mathrm{prod}}
 \le
 \|L_T\|_{\mathrm{op}}^2
 \|\mathbf N_T\|_{A\otimes B,\mathrm{prod}}.}
 \label{eq:V90-product-carrier-domination}
\end{equation}
No multiplicity dimension appears on the right.
\end{theorem}

\begin{proof}
On \(H_T\otimes\mathcal M_T\),
\[
 0\le I_{H_T}\otimes L_T^2
 \le
 \|L_T\|_{\mathrm{op}}^2
 I_{H_T}\otimes I_{\mathcal M_T}.
\]
Extend by zero and apply the positive map
\(\operatorname{Tr}_C\).  This proves
\eqref{eq:V90-product-carrier-order}.  Evaluate the order inequality
on product density matrices and take the supremum.
\end{proof}

\begin{corollary}[Balanced singleton core product cap]
\label{cor:V90-balanced-singleton-cap}
Let
\begin{equation}
 M_{m,k}^{(2)}
 :=
 \operatorname{Tr}_{Q_{2m}}\left[
 P_{T_k}(\mathbf K^\sim)^2
 \right],
 \qquad
 \Pi_{m,k}:=
 \|M_{m,k}^{(2)}\|_{R_m\otimes R_m,\mathrm{prod}}.
 \label{eq:V90-balanced-singleton-core}
\end{equation}
For \(k\le m/2\),
\begin{equation}
 \boxed{
 \Pi_{m,k}
 \le
 \|\widetilde K_{T_k}\|_{\mathrm{op}}^2
 \frac{d_{T_k}}{d_{2m}}
 \le
 \frac{C_K^2}{[1+s_\alpha\Xi(T_k)]^2}
 \frac{d_{T_k}}{d_{2m}}.}
 \label{eq:V90-balanced-singleton-cap}
\end{equation}
\end{corollary}

\begin{proof}
Apply \Cref{thm:V90-product-carrier-domination} with
\(L_{T_k}=\widetilde K_{T_k}\), and then use
\eqref{eq:V90-balanced-product-carrier} and the V84 weighted-cut cap
\eqref{eq:V84-uniform-weighted-cut-cap}.
\end{proof}

\begin{assessment}[The trace and product branches count different
objects]
\label{ass:V90-trace-versus-product}
The V88 trace is
\[
 \mathcal H_{m,k}^{(2)}
 =
 d_{T_k}\operatorname{Tr}_{\mathcal M_{T_k}}
 \widetilde K_{T_k}^2
\]
and therefore may pay \(k+1\) active ancestry directions, as V89
proves it does in the heat benchmark.  The product cap
\eqref{eq:V90-balanced-singleton-cap} instead tests the partially
traced output carrier on one \(A|B\) product state.  Its exact
geometric cost is \(d_{T_k}/d_{2m}\), not
\(d_{T_k}(k+1)/d_m\).  Replacing this cap by the trace branch before
summing is precisely the lost \(1/(k+1)\).
\end{assessment}

\section{The exact truncated internal-pair capacity}
\label{sec:V90-truncated-capacity}

For the internal pair \(e=XY\), put
\[
 d_e:=\max\{d_X,d_Y\},
 \qquad
 R_e(t):=\frac{D_e(t)}{d_e}.
\]
For a positive core shell \(M_j\), retain the V88 quantities
\[
 \Pi_j:=\|M_j\|_{\mathrm{prod}},
 \qquad
 B_j:=\frac{\mathcal H_j}{\min\{d_A,d_B\}}.
\]
By \eqref{eq:V88-core-product-Schur}, \(0\le\Pi_j\le B_j\).
When \(B_j>0\), define
\begin{equation}
 \eta_j:=\frac{\Pi_j}{B_j}\in[0,1],
 \label{eq:V90-core-product-ratio}
\end{equation}
and put \(\eta_j=0\) if \(B_j=0\).

\begin{definition}[Truncated one-pair capacity]
\label{def:V90-truncated-pair-capacity}
For \(0\le\eta\le1\), set
\begin{equation}
 \mathcal K_e(\eta)
 :=
 \int_0^{A_*}\min\{\eta,R_e(t)\}\,\dd t.
 \label{eq:V90-truncated-pair-capacity}
\end{equation}
\end{definition}

\begin{theorem}[Un-factorized shell bound in ratio form]
\label{thm:V90-ratio-shell-bound}
Every positive invariant core shell satisfies
\begin{equation}
 \boxed{
 \|Q_e\otimes M_j\|_{\mathrm{prod}}
 \le
 B_j\mathcal K_e(\eta_j).}
 \label{eq:V90-ratio-shell-bound}
\end{equation}
Moreover,
\begin{equation}
 \mathcal K_e(\eta)
 \le
 \min\{\mathcal K_e(1),A_*\eta\}.
 \label{eq:V90-truncated-elementary-bounds}
\end{equation}
\end{theorem}

\begin{proof}
The second entry of the V88 minimum obeys
\[
 \frac{D_e(t)\mathcal H_j}
 {\max\{d_Xd_A,d_Yd_B\}}
 \le R_e(t)B_j.
\]
The first entry is
\(\Pi_j=\eta_jB_j\).  Therefore
\[
 \min\left\{
 \Pi_j,
 \frac{D_e(t)\mathcal H_j}
 {\max\{d_Xd_A,d_Yd_B\}}
 \right\}
 \le
 B_j\min\{\eta_j,R_e(t)\}.
\]
Integrate and use \eqref{eq:V88-refined-integral}.
The two elementary bounds follow pointwise from
\[
 \min\{\eta,R\}\le\min\{1,R\},
 \qquad
 \min\{\eta,R\}\le\eta.
\]
\end{proof}

\begin{theorem}[Exact spectral/Ky--Fan formula]
\label{thm:V90-Ky-Fan-capacity}
List the eigenvalues of \(Q_e\), with multiplicity, as
\[
 c_1\ge c_2\ge\cdots\ge c_N\ge0,
 \qquad N=d_Xd_Y,
\]
and let \(n=\lfloor\eta d_e\rfloor\).  Then
\begin{equation}
 \boxed{
 \mathcal K_e(\eta)
 =
 \frac1{d_e}\left[
 \sum_{\ell=1}^{n}c_\ell
 +(\eta d_e-n)c_{n+1}
 \right],}
 \label{eq:V90-Ky-Fan-capacity}
\end{equation}
where the final term is omitted when
\(\eta d_e=n\) or \(n=N\).  In particular,
\(\mathcal K_e(\eta)\) is increasing and concave in \(\eta\).
\end{theorem}

\begin{proof}
The rank layer is
\[
 D_e(t)=\#\{\ell:c_\ell>t\}.
\]
Layer cake in the second variable gives
\begin{align*}
 \mathcal K_e(\eta)
 &=
 \int_0^{A_*}\int_0^\eta
 \boldsymbol1_{\{r<R_e(t)\}}\,\dd r\,\dd t\\
 &=
 \int_0^\eta
 \left|\{t:D_e(t)>rd_e\}\right|\,\dd r\\
 &=
 \int_0^\eta c_{\lfloor rd_e\rfloor+1}\,\dd r.
\end{align*}
Partition the last integral at the points \(\ell/d_e\) to obtain
\eqref{eq:V90-Ky-Fan-capacity}.  Its integrand is nonincreasing in
\(r\), which proves concavity; nonnegativity proves monotonicity.
\end{proof}

\begin{remark}[What the truncation retains]
\label{rem:V90-truncation-retains}
The factorized V88 capacity is \(\mathcal K_e(1)\).  The exact
un-factorized shell sees only the top
\(\eta_jd_e\) eigenvalue mass of the internal-pair operator.  Thus the
small core product ratio truncates the pair spectrum before the two
variables are separated.  This is the precise level/core analogue of
the restricted quantities retained in both companion programmes.
\end{remark}

\section{Uniform closure of the balanced thermal pair--core shell}
\label{sec:V90-balanced-closure}

Let
\begin{equation}
 M_{\mathrm{th};m}^{(2)}
 :=
 \sum_{A_0\sqrt{2m}\le k\le B_0\sqrt{2m}}
 M_{m,k}^{(2)}.
 \label{eq:V90-balanced-thermal-core}
\end{equation}
For fixed \(A_0,B_0\), the condition \(k\le m/2\) holds throughout
this window once \(m\) is sufficiently large.

\begin{theorem}[Uniform balanced product-cap sum]
\label{thm:V90-balanced-product-sum}
For every Wilson parameter \(\alpha\) and all sufficiently large
\(m\),
\begin{align}
 \sum_{A_0\sqrt{2m}\le k\le B_0\sqrt{2m}}\Pi_{m,k}
 &\le
 C_K^2
 \sum_{A_0\sqrt{2m}\le k\le B_0\sqrt{2m}}
 \frac{d_{T_k}}{d_{2m}},
 \label{eq:V90-product-sum-bound}\\
 \limsup_{m\to\infty}
 \sum_{A_0\sqrt{2m}\le k\le B_0\sqrt{2m}}\Pi_{m,k}
 &\le
 \frac{C_K^2}{2}(B_0^4-A_0^4).
 \label{eq:V90-product-sum-limit}
\end{align}
The constants are independent of the balanced core label \(m\).
\end{theorem}

\begin{proof}
Drop the denominator
\([1+s_\alpha\Xi(T_k)]^2\ge1\) in
\eqref{eq:V90-balanced-singleton-cap} and sum.  The exact Riemann sum
\eqref{eq:V79-thermal-dimension-mass}, with \(n=2m\), is
\[
 \sum_{A_0\sqrt{2m}\le k\le B_0\sqrt{2m}}
 \frac{d_{T_k}}{d_{2m}}
 \longrightarrow
 \frac{B_0^4-A_0^4}{2}.
\]
This proves both assertions.
\end{proof}

\begin{theorem}[Balanced thermal pair--core closure]
\label{thm:V90-balanced-pair-core-closure}
For every internal-pair positive operator \(Q_e\) from V87,
\begin{equation}
 \boxed{
 \left\|
 Q_e\otimes M_{\mathrm{th};m}^{(2)}
 \right\|_{\mathrm{prod}}
 \le
 A_*
 \sum_{A_0\sqrt{2m}\le k\le B_0\sqrt{2m}}\Pi_{m,k}.}
 \label{eq:V90-balanced-pair-core-closure}
\end{equation}
Consequently
\begin{equation}
 \limsup_{m\to\infty}
 \left\|
 Q_e\otimes M_{\mathrm{th};m}^{(2)}
 \right\|_{\mathrm{prod}}
 \le
 \frac{A_*C_K^2}{2}(B_0^4-A_0^4).
 \label{eq:V90-balanced-pair-core-uniform}
\end{equation}
No selected-heavy hypothesis, inverse-\(\sqrt m\) pair debit, or
Wilson-to-heat comparison is used.
\end{theorem}

\begin{proof}
Positivity and the triangle inequality for product expectations give
\[
 \|Q_e\otimes M_{\mathrm{th};m}^{(2)}\|_{\mathrm{prod}}
 \le
 \sum_k\|Q_e\otimes M_{m,k}^{(2)}\|_{\mathrm{prod}}.
\]
Apply \eqref{eq:V90-ratio-shell-bound} and then
\eqref{eq:V90-truncated-elementary-bounds}:
\[
 \|Q_e\otimes M_{m,k}^{(2)}\|_{\mathrm{prod}}
 \le
 B_{m,k}\mathcal K_e(\eta_{m,k})
 \le
 A_*B_{m,k}\eta_{m,k}
 =
 A_*\Pi_{m,k}.
\]
Sum and use \eqref{eq:V90-product-sum-limit}.
\end{proof}

\begin{remark}[This does not reverse the V89 no-go]
\label{rem:V90-no-conflict-V89}
V89 proved that the \emph{trace} factor
\(\sum_k\mathcal H_{m,k}^{(2)}/d_m\) is of order \(\sqrt m\) in the
heat benchmark.  The present theorem never bounds that trace by a
constant.  It retains the smaller product norm
\(\Pi_{m,k}\), whose total output-carrier mass is the order-one
Riemann sum with denominator \(d_{2m}\).  Thus V89 and V90 identify
the two sides of the same sharp distinction.
\end{remark}

\section{Sharp value of the product/trace ratio in the heat model}
\label{sec:V90-sharp-ratio}

For a heat or Wilson balanced block, write
\[
 \epsilon_{m,k}
 :=
 \|K_{m,k}-\mathfrak d(u)I\|_{\mathrm{op}},
 \qquad
 w_{m,k}:=1+s_\alpha\Xi(T_k).
\]
For the heat block take \(s_\alpha=\tau\).  Under the full-bank Wilson
comparison, take \(s_\alpha\) and \(K_{\alpha;m,k}\).

\begin{theorem}[The missing factor is exactly \(k+1\)]
\label{thm:V90-sharp-product-trace-ratio}
Assume \(\epsilon_{m,k}\to0\) uniformly on the transition window.
Let
\[
 \mathcal H_{m,k}^{(2)}=\operatorname{Tr}M_{m,k}^{(2)},
 \qquad
 B_{m,k}=\frac{\mathcal H_{m,k}^{(2)}}{d_m},
 \qquad
 \eta_{m,k}=\frac{\Pi_{m,k}}{B_{m,k}}.
\]
Then
\begin{equation}
 \boxed{
 \sup_{A_0\sqrt{2m}\le k\le B_0\sqrt{2m}}
 \left|
 (k+1)\eta_{m,k}-\frac14
 \right|
 \longrightarrow0.}
 \label{eq:V90-sharp-product-trace-ratio}
\end{equation}
Moreover,
\begin{equation}
 \sum_{A_0\sqrt{2m}\le k\le B_0\sqrt{2m}}
 \Pi_{m,k}
 \longrightarrow
 \frac{B_0^4-A_0^4}{2}\mathfrak d(u)^2.
 \label{eq:V90-sharp-product-sum}
\end{equation}
The assertions hold unconditionally for the heat benchmark and
conditionally for Wilson under
\eqref{eq:V89-full-comparison-hypothesis}.
\end{theorem}

\begin{proof}
Put \(d=\mathfrak d(u)>0\).  For sufficiently large \(m\),
\(\epsilon_{m,k}<d\), and self-adjointness gives on the full
multiplicity space
\[
 (d-\epsilon_{m,k})^2I
 \le K_{m,k}^2
 \le(d+\epsilon_{m,k})^2I.
\]
After multiplication by \(w_{m,k}^2\), extension to the
\(T_k\)-isotypic block, and positive partial trace over \(C\), this
becomes
\[
 w_{m,k}^2(d-\epsilon_{m,k})^2\mathbf N_{m,k}
 \le M_{m,k}^{(2)}
 \le
 w_{m,k}^2(d+\epsilon_{m,k})^2\mathbf N_{m,k}.
\]
Use \eqref{eq:V90-balanced-product-carrier} for product norms and
\eqref{eq:V90-output-carrier-trace} for traces.  One obtains
\begin{align*}
 &w^2(d-\epsilon)^2\frac{d_{T_k}}{d_{2m}}
 \le\Pi_{m,k}\le
 w^2(d+\epsilon)^2\frac{d_{T_k}}{d_{2m}},\\
 &w^2(d-\epsilon)^2d_{T_k}(k+1)
 \le\mathcal H_{m,k}^{(2)}\le
 w^2(d+\epsilon)^2d_{T_k}(k+1).
\end{align*}
Therefore
\begin{equation}
 \frac{d_m}{d_{2m}(k+1)}
 \left(\frac{d-\epsilon}{d+\epsilon}\right)^2
 \le\eta_{m,k}\le
 \frac{d_m}{d_{2m}(k+1)}
 \left(\frac{d+\epsilon}{d-\epsilon}\right)^2.
 \label{eq:V90-ratio-squeeze}
\end{equation}
Since \(d_m/d_{2m}\to1/4\), the uniform squeeze proves
\eqref{eq:V90-sharp-product-trace-ratio}.

For the sum, note that on a one-link output
\(\Xi(T_k)=1+c_{T_k}\).  The transition scaling and
\(k=O(\sqrt m)\) give \(s_\alpha\Xi(T_k)\to0\) uniformly, so
\(w_{m,k}\to1\).  The product-norm squeeze and
\eqref{eq:V79-thermal-dimension-mass} then give
\eqref{eq:V90-sharp-product-sum}.  Heat scalarization is
\eqref{eq:V89-full-heat-limit}; conditional Wilson scalarization is
\eqref{eq:V89-Wilson-full-limit}.
\end{proof}

\begin{corollary}[Sharp truncated-capacity scale]
\label{cor:V90-sharp-truncation-scale}
Under the hypotheses of
\Cref{thm:V90-sharp-product-trace-ratio},
\begin{equation}
 \eta_{m,k}
 =
 \frac{1+o(1)}{4(k+1)}
 \label{eq:V90-eta-scale}
\end{equation}
uniformly on the thermal window.  Hence the un-factorized pair layer
uses only a \(k^{-1}\)-fraction of the normalized top spectral mass:
\begin{equation}
 \mathcal K_e(\eta_{m,k})
 =
 \frac1{d_e}\left[
 \sum_{\ell\le \eta_{m,k}d_e}c_\ell
 +\hbox{one fractional eigenvalue}
 \right].
 \label{eq:V90-thermal-truncated-Ky-Fan}
\end{equation}
\end{corollary}

\begin{proof}
The first statement is
\eqref{eq:V90-sharp-product-trace-ratio}; the second is
\eqref{eq:V90-Ky-Fan-capacity}.
\end{proof}

\section{A representation-geometric gate for general core shells}
\label{sec:V90-general-gate}

For a final output \(T\subset A\otimes B\otimes C\), define
\begin{equation}
 \nu_{A|B;C}(T)
 :=
 \left\|
 \operatorname{Tr}_C\mathbf P_T
 \right\|_{A\otimes B,\mathrm{prod}}.
 \label{eq:V90-geometric-visibility}
\end{equation}
This number depends only on the representation geometry and the
chosen retained-row cut, not on the Wilson multipliers.

\begin{theorem}[General product-carrier shell criterion]
\label{thm:V90-general-shell-criterion}
Let \(\Omega\) be any finite output shell and set
\[
 M_\Omega^{(2)}
 :=
 \sum_{T\in\Omega}
 \operatorname{Tr}_C[
 I_{H_T}\otimes\widetilde K_T^2].
\]
Then every internal-pair operator \(Q_e\) obeys
\begin{equation}
 \boxed{
 \|Q_e\otimes M_\Omega^{(2)}\|_{\mathrm{prod}}
 \le
 A_*
 \sum_{T\in\Omega}
 \|\widetilde K_T\|_{\mathrm{op}}^2
 \nu_{A|B;C}(T).}
 \label{eq:V90-general-shell-criterion}
\end{equation}
Thus summability of the product-carrier visibility, rather than the
full multiplicity trace, is a sufficient gate for the pair--core
shell.
\end{theorem}

\begin{proof}
For each \(T\), apply
\eqref{eq:V90-truncated-elementary-bounds} in
\Cref{thm:V90-ratio-shell-bound} to get
\[
 \|Q_e\otimes M_T^{(2)}\|_{\mathrm{prod}}
 \le A_*\|M_T^{(2)}\|_{\mathrm{prod}}.
\]
Then use \eqref{eq:V90-product-carrier-domination}, sum the positive
shells, and take the product-state supremum.
\end{proof}

\begin{corollary}[Balanced value of the geometric gate]
\label{cor:V90-balanced-geometric-gate}
For the balanced family and all sufficiently large \(m\),
\begin{equation}
 \nu_{R_m|R_m;Q_{2m}}(T_k)
 =
 \frac{d_{T_k}}{d_{2m}},
 \label{eq:V90-balanced-geometric-visibility}
\end{equation}
and its thermal sum is order one:
\begin{equation}
 \sum_{A_0\sqrt{2m}\le k\le B_0\sqrt{2m}}
 \nu_{R_m|R_m;Q_{2m}}(T_k)
 \longrightarrow
 \frac{B_0^4-A_0^4}{2}.
 \label{eq:V90-balanced-geometric-sum}
\end{equation}
\end{corollary}

\begin{proof}
Use \Cref{thm:V90-balanced-product-carrier} and
\eqref{eq:V79-thermal-dimension-mass}.
\end{proof}

\begin{assessment}[Audit-directed change of route]
\label{ass:V90-audit-direction}
The companion audit does not contribute a Yang--Mills estimate.
It does identify the same proof-order failure in three settings:
\begin{enumerate}[label=\textup{(\roman*)}]
\item the SM energy pivot does not control the strong grazing port
      observation, but its closed graph domain retains the Green
      cancellation;
\item separate regulated SM sector values do not imply a common Ward
      identity without a common invariant form domain;
\item the NS flat mark becomes continuation-class after its
      stopping-time/high-parent correlation is discarded.
\end{enumerate}
Accordingly, V90 does not seek a stronger marginal trace estimate
after V89 proved that estimate false.  It retains the partially traced
output carrier and obtains the missing \(1/(k+1)\) from its exact
product geometry.  The next generalization must preserve this carrier
through the actual V59--V60 marked-tree alternatives.
\end{assessment}

\section{V90 ledger and V91 mandate}
\label{sec:V90-ledger}

\begin{theorem}[Integrated V90 statement]
\label{thm:V90-integrated}
Preserving V1--V89, V90 proves:
\begin{enumerate}[label=\textup{(V90.\arabic*)},leftmargin=5.7em]
\item the compulsory read-only SM V90 and NS V31 audit;
\item the exact pathwise partial trace
      \eqref{eq:V90-path-partial-trace};
\item the exact balanced product-carrier norm
      \eqref{eq:V90-balanced-product-carrier};
\item product-carrier domination before multiplicity trace
      \eqref{eq:V90-product-carrier-domination};
\item the balanced singleton core product cap
      \eqref{eq:V90-balanced-singleton-cap};
\item the exact truncated pair capacity and Ky--Fan formula
      \eqref{eq:V90-ratio-shell-bound} and
      \eqref{eq:V90-Ky-Fan-capacity};
\item the unconditional uniform balanced pair--core closure
      \eqref{eq:V90-balanced-pair-core-uniform};
\item the sharp heat/conditional-Wilson ratio
      \eqref{eq:V90-sharp-product-trace-ratio}; and
\item the general representation-geometric shell criterion
      \eqref{eq:V90-general-shell-criterion}.
\end{enumerate}
\end{theorem}

\begin{proof}
These are
\Cref{audit:V90-companions,
lem:V90-path-partial-trace,
thm:V90-balanced-product-carrier,
thm:V90-product-carrier-domination,
cor:V90-balanced-singleton-cap,
thm:V90-ratio-shell-bound,
thm:V90-Ky-Fan-capacity,
thm:V90-balanced-pair-core-closure,
thm:V90-sharp-product-trace-ratio,
thm:V90-general-shell-criterion}.
\end{proof}

\begin{openproblem}[V91 mandate: product-carrier atlas]
\label{op:V91-mandate}
Insert \eqref{eq:V90-general-shell-criterion} into the actual V86
one-sideways heavy-puncture alternatives.
\begin{enumerate}
\item For every retained-row orientation, express
      \(\operatorname{Tr}_C\mathbf P_T\) in a first-stage path basis,
      including path multiplicities.
\item Bound its product visibility
      \(\nu_{A|B;C}(T)\) before summing over \(T\); do not replace it
      by \(\operatorname{Tr}\mathbf P_T\).
\item Identify which Pieri or Littlewood--Richardson sectors admit a
      product-visible maximal eigenvalue as in
      \eqref{eq:V90-balanced-product-carrier}, and which require a
      genuine truncated Ky--Fan estimate.
\item Sum the resulting visibility with the V63 weighted-cut cap on
      the selected heavy-puncture branch.
\item Treat the off-cut-core alternative separately; V90 closes only
      the balanced thermal core family.
\item Repeat the construction in the opposite orientation and retain
      the same level variable when the two orientations overlap.
\item If a general visibility sum fails, build a sharp
      representation packet showing the failure and return upstream
      to the V59 path selection rather than restoring the full trace.
\end{enumerate}
The next compulsory companion audit is V95.
\end{openproblem}

\section{Firewall after V90}
\label{sec:V90-firewall}

\begin{firewall}[Current claim boundary]
\label{fire:V90-final}
V90 closes the balanced thermal pair--core shell uniformly through
the exact product carrier, without a Wilson-to-heat hypothesis or an
inverse-\(\sqrt m\) pair debit.  This is a genuine removal of the V88
square-root loss, but only for the balanced \(SU(3)\) family and one
retained-row orientation.  V90 does not prove the general
product-carrier atlas bound, close the off-cut alternative,
synchronize both orientations, prove JREC, aggregate all punctures
and collisions, construct the continuum Yang--Mills theory, or prove
a positive mass gap.
\end{firewall}

\section{Append-only record for V90}
\label{sec:V90-record}

V90 was appended to the complete V89 manuscript.  The exact V89
parent had \(61721\) physical lines, \(2172915\) bytes, and
SHA--256 digest
\[
 \texttt{fc466f7d2c17cc22f4096aaba01ee5c72351299ae481f79191704776fbdafd08}.
\]
No V89 mathematical content was changed.  The companion audit was
read-only.  The sole final \verb|\end{document}| is restored below.

% =============================================================================
% END APPENDED VERSION V90 MATERIAL
% =============================================================================

% =============================================================================
% BEGIN APPENDED VERSION V91 MATERIAL
% =============================================================================

\chapter{V91: The Cartan-Split Product-Carrier Atlas}
\label{ch:V91}

\section{Purpose}
\label{sec:V91-purpose}

V90 closed the balanced family \(R_m,R_m,Q_{2m}\) by computing the
product norm of the partially traced output carrier.  The equal split
is not essential.  This chapter proves the same carrier identity for
every Cartan split
\[
 A=R_a,\qquad B=R_b,\qquad C=Q_n,\qquad a+b=n,
\]
including highly asymmetric splits.  The ancestry multiplicity may
shrink when \(\min\{a,b\}<k\), but the product-carrier norm remains
exactly \(d_{T_k}/d_n\).  Hence the entire Cartan-split thermal atlas
has the same order-one visibility sum and the same unconditional
pair--core closure.

The proof also records the general first-stage spectrum of
\(\operatorname{Tr}_C\mathbf P_T\), with arbitrary tensor-product
multiplicities.  That formula identifies the next non-Cartan gate as
a maximal ratio \(N_{SC}^T/d_S\) together with product visibility of
the maximizing \(AB\)-isotypic space.

\section{General first-stage spectrum of an output carrier}
\label{sec:V91-general-first-stage}

Let
\[
 H_A\otimes H_B
 =
 \bigoplus_S H_S\otimes\mathcal M_{AB}^S
\]
be the \(AB\)-decomposition, and let \(\mathbf P_T\) be the complete
\(T\)-isotypic projector in \(A\otimes B\otimes C\).  Put
\(\mathbf N_T=\operatorname{Tr}_C\mathbf P_T\).

\begin{theorem}[Exact first-stage carrier spectrum]
\label{thm:V91-general-carrier-spectrum}
On the \(S\)-isotypic part of \(A\otimes B\),
\begin{equation}
 \boxed{
 \mathbf N_T
 =
 \bigoplus_{S\subset A\otimes B}
 \frac{d_TN_{SC}^{T}}{d_S}
 I_{H_S}\otimes I_{\mathcal M_{AB}^{S}}.}
 \label{eq:V91-general-carrier-spectrum}
\end{equation}
Consequently
\begin{align}
 \operatorname{Tr}\mathbf N_T
 &=
 d_T\sum_S N_{AB}^{S}N_{SC}^{T}
 =
 d_Tm_{ABC}(T),
 \label{eq:V91-general-carrier-trace}\\
 \|\mathbf N_T\|_{\mathrm{op}}
 &=
 d_T\max_{\substack{S\subset A\otimes B\\N_{SC}^{T}>0}}
 \frac{N_{SC}^{T}}{d_S}.
 \label{eq:V91-general-carrier-operator}
\end{align}
\end{theorem}

\begin{proof}
The character-projector formula for \(\mathbf P_T\) is a central
average of the representation on
\((A\otimes B)\otimes C\).  In the displayed \(AB\)-decomposition it
acts trivially on every first-stage multiplicity factor
\(\mathcal M_{AB}^S\).  After tracing \(C\), Schur's lemma removes
cross terms between inequivalent \(S\)'s and makes the restriction to
each \(H_S\) a scalar.

Fix one copy of \(S\).  Inside \(S\otimes C\), the complete
\(T\)-isotypic projector has rank \(d_TN_{SC}^{T}\).  If its partial
trace is \(\lambda_{S,T}I_{H_S}\), equality of traces gives
\[
 \lambda_{S,T}d_S=d_TN_{SC}^{T}.
\]
This proves \eqref{eq:V91-general-carrier-spectrum}.  Taking the trace
and the largest block eigenvalue proves
\eqref{eq:V91-general-carrier-trace} and
\eqref{eq:V91-general-carrier-operator}.
\end{proof}

\begin{corollary}[General visibility bound and equality criterion]
\label{cor:V91-general-visibility}
The geometric visibility of V90 satisfies
\begin{equation}
 \boxed{
 \nu_{A|B;C}(T)
 \le
 d_T\max_S\frac{N_{SC}^{T}}{d_S}.}
 \label{eq:V91-general-visibility-bound}
\end{equation}
Equality holds whenever an \(AB\)-product vector lies in an isotypic
space attaining the maximum on the right.
\end{corollary}

\begin{proof}
Product norm is at most operator norm, so use
\eqref{eq:V91-general-carrier-operator}.  If a unit product vector
belongs to a maximizing eigenspace, its expectation attains the
operator norm.
\end{proof}

\begin{remark}[What multiplicity survives]
\label{rem:V91-multiplicity-survival}
The full trace pays
\(\sum_SN_{AB}^{S}N_{SC}^{T}\).  The operator and product-carrier
branches do not pay the first-stage multiplicity \(N_{AB}^{S}\);
they retain only the second-arrow ratio \(N_{SC}^{T}/d_S\).  Thus the
non-Cartan atlas should estimate that ratio and product visibility
before summing multiplicities.
\end{remark}

\section{Exact Cartan-split path support}
\label{sec:V91-Cartan-support}

Fix \(a,b\ge0\), put \(n=a+b\), and write
\begin{equation}
 R_a\otimes R_b
 =
 \bigoplus_{r=0}^{\min\{a,b\}}S_r^{(n)},
 \qquad
 S_r^{(n)}:=(n-2r,r).
 \label{eq:V91-Cartan-AB-Pieri}
\end{equation}

\begin{theorem}[Exact Cartan-split \(AB\)-first support]
\label{thm:V91-Cartan-support}
For \(0\le k\le n\),
\begin{equation}
 N_{S_r^{(n)},Q_n}^{T_k}
 =
 \begin{cases}
 1,&r\le k\ \hbox{and}\ r\le n-k,\\
 0,&\hbox{otherwise}.
 \end{cases}
 \label{eq:V91-Cartan-support}
\end{equation}
In particular, if \(k\le n/2\), the paths reaching \(T_k\) are
exactly
\[
 0\le r\le \min\{a,b,k\},
\]
each with multiplicity one.
\end{theorem}

\begin{proof}
Frobenius reciprocity gives
\[
 N_{S_r^{(n)},Q_n}^{T_k}
 =
 N_{T_k,R_n}^{S_r^{(n)}}.
\]
Represent \(T_k\) by the partition \((2k,k,0)\).  The only partition
representing \(S_r^{(n)}\) with the total size \(3k+n\) is
\[
 \lambda=(n-r+k,r+k,k).
\]
Pieri's horizontal-strip interlacing is
\[
 \lambda_1\ge2k\ge\lambda_2\ge k\ge\lambda_3\ge0.
\]
Its two nonautomatic inequalities are
\[
 n-r\ge k,\qquad k\ge r,
\]
equivalent to \(r\le n-k\) and \(r\le k\).  Pieri multiplicity is
one.  Finally intersect this support with the range
\(0\le r\le\min\{a,b\}\) in
\eqref{eq:V91-Cartan-AB-Pieri}.
\end{proof}

\section{Exact Cartan-split product visibility}
\label{sec:V91-Cartan-visibility}

\begin{lemma}[Dimension monotonicity for \(S_r^{(n)}\)]
\label{lem:V91-Cartan-dimension}
If \(n\ge4\) and \(0\le r\le n/4\), then
\begin{equation}
 d_{S_r^{(n)}}\ge d_{S_0^{(n)}}=d_n.
 \label{eq:V91-Cartan-dimension}
\end{equation}
\end{lemma}

\begin{proof}
The dimension formula gives
\[
 d_{S_r^{(n)}}
 =
 \frac12(n-2r+1)(r+1)(n-r+2).
\]
Direct factorization yields
\begin{equation}
 2(d_{S_r^{(n)}}-d_n)
 =
 r\left[
 2r^2-(3n+3)r+n^2-3
 \right].
 \label{eq:V91-Cartan-dimension-factor}
\end{equation}
The bracket is decreasing on \(0\le r\le n/4\).  At \(r=n/4\)
it is
\[
 \frac38(n-4)(n+2)\ge0.
\]
Equation \eqref{eq:V91-Cartan-dimension-factor} proves the claim.
\end{proof}

Let \(\mathbf P_{a,b;n,k}\) denote the complete \(T_k\)-isotypic
projector in
\[
 R_a\otimes R_b\otimes Q_n,
 \qquad a+b=n,
\]
and put
\[
 \mathbf N_{a,b;n,k}
 :=
 \operatorname{Tr}_{Q_n}\mathbf P_{a,b;n,k}.
\]

\begin{theorem}[Split-independent Cartan product-carrier norm]
\label{thm:V91-Cartan-product-carrier}
If \(n\ge4\) and \(0\le k\le n/4\), then for every split
\(a+b=n\),
\begin{equation}
 \boxed{
 \|\mathbf N_{a,b;n,k}\|_
 {R_a\otimes R_b,\mathrm{prod}}
 =
 \frac{d_{T_k}}{d_n}.}
 \label{eq:V91-Cartan-product-carrier}
\end{equation}
The ancestry multiplicity is
\begin{equation}
 m_{R_a,R_b,Q_n}(T_k)
 =
 1+\min\{a,b,k\},
 \label{eq:V91-Cartan-multiplicity}
\end{equation}
but it does not enter
\eqref{eq:V91-Cartan-product-carrier}.
\end{theorem}

\begin{proof}
By \Cref{thm:V91-Cartan-support}, the carrier spectrum is
\[
 \mathbf N_{a,b;n,k}
 =
 \sum_{r=0}^{\min\{a,b,k\}}
 \frac{d_{T_k}}{d_{S_r^{(n)}}}
 P_{AB}(S_r^{(n)}).
\]
Every active \(r\) is at most \(k\le n/4\), so
\Cref{lem:V91-Cartan-dimension} shows that the largest eigenvalue is
\(d_{T_k}/d_n\), attained at \(r=0\).  The tensor product of the
highest-weight vectors of \(R_a\) and \(R_b\) lies in the Cartan
summand \(S_0^{(n)}=R_n\).  It is an admissible \(A|B\) product
vector and attains that eigenvalue.  This proves
\eqref{eq:V91-Cartan-product-carrier}.

The active path indices are the integers from zero through
\(\min\{a,b,k\}\), each with multiplicity one, proving
\eqref{eq:V91-Cartan-multiplicity}.
\end{proof}

\begin{corollary}[The highly asymmetric edge is harmless]
\label{cor:V91-asymmetric-Cartan}
If \(\min\{a,b\}<k\), the full trace pays only
\(1+\min\{a,b\}\) ancestry paths rather than \(k+1\), while the
product-carrier norm remains \(d_{T_k}/d_n\).  If \(a=0\) or \(b=0\),
there is one path and the same carrier formula holds.
\end{corollary}

\begin{proof}
This is immediate from
\eqref{eq:V91-Cartan-product-carrier} and
\eqref{eq:V91-Cartan-multiplicity}.
\end{proof}

\section{Uniform closure of the entire Cartan-split thermal atlas}
\label{sec:V91-Cartan-closure}

For each split \(a+b=n\), define the singleton core shell
\begin{equation}
 M_{a,b;n,k}^{(2)}
 :=
 \operatorname{Tr}_{Q_n}\left[
 P_{T_k}(\mathbf K^\sim)^2
 \right],
 \qquad
 \Pi_{a,b;n,k}
 :=
 \|M_{a,b;n,k}^{(2)}\|_{R_a\otimes R_b,\mathrm{prod}}.
 \label{eq:V91-Cartan-core-shell}
\end{equation}

\begin{theorem}[Split-uniform singleton cap]
\label{thm:V91-Cartan-singleton-cap}
For fixed \(0<A_0<B_0<\infty\), all sufficiently large \(n\), every
split \(a+b=n\), and
\(A_0\sqrt n\le k\le B_0\sqrt n\),
\begin{equation}
 \boxed{
 \Pi_{a,b;n,k}
 \le
 \frac{C_K^2}{[1+s_\alpha\Xi(T_k)]^2}
 \frac{d_{T_k}}{d_n}.}
 \label{eq:V91-Cartan-singleton-cap}
\end{equation}
\end{theorem}

\begin{proof}
The thermal window lies in \(k\le n/4\) for all sufficiently large
\(n\).  Apply the general product-carrier domination
\eqref{eq:V90-product-carrier-domination}, the exact visibility
\eqref{eq:V91-Cartan-product-carrier}, and the V84 weighted-cut cap
\eqref{eq:V84-uniform-weighted-cut-cap}.
\end{proof}

\begin{theorem}[Cartan-split thermal pair--core closure]
\label{thm:V91-Cartan-pair-core-closure}
Let
\[
 M_{\mathrm{th};a,b,n}^{(2)}
 :=
 \sum_{A_0\sqrt n\le k\le B_0\sqrt n}
 M_{a,b;n,k}^{(2)}.
\]
For every internal-pair positive operator \(Q_e\),
\begin{align}
 \|Q_e\otimes M_{\mathrm{th};a,b,n}^{(2)}\|_{\mathrm{prod}}
 &\le
 A_*\sum_{A_0\sqrt n\le k\le B_0\sqrt n}
 \Pi_{a,b;n,k},
 \label{eq:V91-Cartan-pair-core-first}\\
 \limsup_{n\to\infty}
 \sup_{a+b=n}
 \|Q_e\otimes M_{\mathrm{th};a,b,n}^{(2)}\|_{\mathrm{prod}}
 &\le
 \frac{A_*C_K^2}{2}(B_0^4-A_0^4).
 \label{eq:V91-Cartan-pair-core-uniform}
\end{align}
The estimate is unconditional and uniform over all Cartan splits.
\end{theorem}

\begin{proof}
The first line is the same un-factorized product branch as
\eqref{eq:V90-balanced-pair-core-closure}.  For the second, use
\eqref{eq:V91-Cartan-singleton-cap}, drop the denominator
\([1+s_\alpha\Xi(T_k)]^2\ge1\), and apply the exact Riemann sum
\[
 \sum_{A_0\sqrt n\le k\le B_0\sqrt n}
 \frac{d_{T_k}}{d_n}
 \longrightarrow
 \frac{B_0^4-A_0^4}{2}.
\]
No bound depends on \(a\) or \(b\).
\end{proof}

\begin{assessment}[What V91 removes from the obstruction atlas]
\label{ass:V91-atlas-removal}
Every core column of the V79 Cartan form
\[
 R_a\otimes R_b\otimes Q_{a+b}\longrightarrow T_k,
 \qquad k\asymp\sqrt{a+b},
\]
is now paid uniformly when tensorized with one retained internal pair.
The V83/V89 heat obstruction is therefore not an obstruction to the
un-factorized product route, even away from the equal split.  Future
counterexamples or estimates must leave this Cartan-split atlas and
activate either a non-Cartan first-stage multiplicity, a large
second-arrow ratio \(N_{SC}^{T}/d_S\), or an off-cut core sector.
\end{assessment}

\section{The fusion Markov kernel and a sharper joint-shell gate}
\label{sec:V91-fusion-Markov}

For every \(S,C\), define
\begin{equation}
 \mathsf p_C(T\mid S)
 :=
 \frac{d_TN_{SC}^{T}}{d_Sd_C}.
 \label{eq:V91-fusion-Markov-kernel}
\end{equation}

\begin{lemma}[Fusion dimensions form a probability kernel]
\label{lem:V91-fusion-probability}
For every \(S,C\),
\begin{equation}
 \mathsf p_C(T\mid S)\ge0,
 \qquad
 \sum_T\mathsf p_C(T\mid S)=1.
 \label{eq:V91-fusion-probability}
\end{equation}
Moreover, for a unit \(AB\)-product vector \(v=v_A\otimes v_B\),
\begin{equation}
 \langle v,\mathbf N_Tv\rangle
 =
 d_C\sum_S
 \mathsf p_C(T\mid S)\,
 \pi_v(S),
 \qquad
 \pi_v(S):=\langle v,P_{AB}(S)v\rangle.
 \label{eq:V91-product-fusion-mixture}
\end{equation}
\end{lemma}

\begin{proof}
The tensor-product dimension identity is
\[
 \sum_Td_TN_{SC}^{T}=d_Sd_C,
\]
which proves \eqref{eq:V91-fusion-probability}.  The isotypic
probabilities \(\pi_v(S)\) are nonnegative and sum to one.  Insert
\eqref{eq:V91-fusion-Markov-kernel} in the carrier spectrum
\eqref{eq:V91-general-carrier-spectrum} and evaluate on \(v\) to
obtain \eqref{eq:V91-product-fusion-mixture}.
\end{proof}

\begin{theorem}[Joint-shell fusion-kernel gate]
\label{thm:V91-joint-shell-fusion-gate}
Let \(\Omega\) be an output shell and suppose positive core operators
\(M_T\) satisfy
\begin{equation}
 0\le M_T\le a_T\mathbf N_T,
 \qquad T\in\Omega,
 \label{eq:V91-carrier-majorants}
\end{equation}
for scalars \(a_T\ge0\).  Then
\begin{align}
 \left\|\sum_{T\in\Omega}M_T\right\|_{\mathrm{prod}}
 &\le
 d_C
 \max_{S\subset A\otimes B}
 \sum_{T\in\Omega}
 a_T\mathsf p_C(T\mid S),
 \label{eq:V91-joint-shell-core-gate}\\
 \boxed{
 \left\|Q_e\otimes\sum_{T\in\Omega}M_T\right\|_{\mathrm{prod}}
 &\le
 A_*d_C
 \max_{S\subset A\otimes B}
 \sum_{T\in\Omega}
 a_T\mathsf p_C(T\mid S).}
 \label{eq:V91-joint-shell-pair-gate}
\end{align}
For the connected square one may take
\begin{equation}
 a_T=\|\widetilde K_T\|_{\mathrm{op}}^2
 \le\frac{C_K^2}{[1+s_\alpha\Xi(T)]^2}.
 \label{eq:V91-connected-majorant}
\end{equation}
\end{theorem}

\begin{proof}
Sum \eqref{eq:V91-carrier-majorants}.  By
\eqref{eq:V91-general-carrier-spectrum}, the positive majorant
\(\sum_Ta_T\mathbf N_T\) is block scalar on every \(S\)-isotypic
space, with eigenvalue
\[
 d_C\sum_{T\in\Omega}a_T\mathsf p_C(T\mid S).
\]
Its product norm is at most its operator norm, which is the maximum
of these eigenvalues.  This proves
\eqref{eq:V91-joint-shell-core-gate}.  Since
\(0\le Q_e\le A_*I\), the product-state domination argument of
\Cref{lem:V88-pair-below-core-product} gives the additional factor
\(A_*\), proving \eqref{eq:V91-joint-shell-pair-gate}.
Finally, \eqref{eq:V91-connected-majorant} is
\Cref{thm:V90-product-carrier-domination} together with the V84 cap.
\end{proof}

\begin{remark}[Why the maximum is taken after the output sum]
\label{rem:V91-max-after-sum}
The earlier sufficient criterion
\eqref{eq:V90-general-shell-criterion} sums
\(\sup_v\langle v,\mathbf N_Tv\rangle\) separately for each \(T\).
Equation \eqref{eq:V91-joint-shell-pair-gate} first sums the output
majorants and then takes one product-state supremum.  This preserves
the common first-stage distribution \(\pi_v(S)\) and is never worse
than summing outputwise maxima.  It is the precise common-domain
improvement suggested by the V90 companion audit.
\end{remark}

\begin{corollary}[Recovery of the Cartan-split bound]
\label{cor:V91-Markov-recovers-Cartan}
For the Cartan split and the thermal shell, every contributing
\(S_r^{(n)}\) has \(r\le B_0\sqrt n\), hence
\(d_{S_r^{(n)}}\ge d_n\) for large \(n\).  Therefore
\begin{equation}
 d_n\max_r
 \sum_{A_0\sqrt n\le k\le B_0\sqrt n}
 a_{T_k}\mathsf p_{Q_n}(T_k\mid S_r^{(n)})
 \le
 \sum_{A_0\sqrt n\le k\le B_0\sqrt n}
 a_{T_k}\frac{d_{T_k}}{d_n}.
 \label{eq:V91-Markov-Cartan-bound}
\end{equation}
With \(a_{T_k}\le C_K^2\), this is exactly the order-one bound in
\Cref{thm:V91-Cartan-pair-core-closure}.
\end{corollary}

\begin{proof}
The support theorem gives
\[
 \mathsf p_{Q_n}(T_k\mid S_r^{(n)})
 =
 \frac{d_{T_k}}{d_{S_r^{(n)}}d_n}
\]
when the path exists and zero otherwise.  If \(r>B_0\sqrt n\), no
thermal \(k\) reaches that intermediate.  Otherwise
\Cref{lem:V91-Cartan-dimension} gives
\(d_{S_r^{(n)}}\ge d_n\).  Sum and maximize.
\end{proof}

\section{V91 ledger and V92 mandate}
\label{sec:V91-ledger}

\begin{theorem}[Integrated V91 statement]
\label{thm:V91-integrated}
Preserving V1--V90, V91 proves:
\begin{enumerate}[label=\textup{(V91.\arabic*)},leftmargin=5.7em]
\item the general first-stage carrier spectrum
      \eqref{eq:V91-general-carrier-spectrum};
\item the general visibility bound
      \eqref{eq:V91-general-visibility-bound};
\item the exact Cartan-split support
      \eqref{eq:V91-Cartan-support};
\item the split-independent carrier norm
      \eqref{eq:V91-Cartan-product-carrier};
\item unconditional closure of the entire Cartan-split thermal atlas
      \eqref{eq:V91-Cartan-pair-core-uniform};
\item the normalized fusion probability kernel
      \eqref{eq:V91-fusion-probability}; and
\item the joint-shell fusion-kernel gate
      \eqref{eq:V91-joint-shell-pair-gate}.
\end{enumerate}
\end{theorem}

\begin{proof}
These are
\Cref{thm:V91-general-carrier-spectrum,
cor:V91-general-visibility,
thm:V91-Cartan-support,
thm:V91-Cartan-product-carrier,
thm:V91-Cartan-pair-core-closure,
lem:V91-fusion-probability,
thm:V91-joint-shell-fusion-gate}.
\end{proof}

\begin{openproblem}[V92 mandate: fusion-kernel concentration on the
selected heavy atlas]
\label{op:V92-mandate}
Return to the actual V59--V60 selected heavy-puncture alternatives
and use the joint-shell gate before any outputwise maximum.
\begin{enumerate}
\item Identify the first-stage \(S\) and remaining input \(C\) in
      each retained-row orientation.
\item Express the admissible output shell as an event for the fusion
      probability \(\mathsf p_C(T\mid S)\).
\item Prove a uniform bound for
      \[
      d_C\sup_S\sum_{T\in\Omega}
      \frac{\mathsf p_C(T\mid S)}
      {[1+s_\alpha\Xi(T)]^2}
      \]
      on the selected heavy branch, retaining its level and path
      indicators.
\item Separate Cartan-visible intermediates, now closed by V91, from
      genuinely multiplicity-rich non-Cartan intermediates.
\item Test the remaining supremum on boundary Littlewood--Richardson
      families before claiming concentration.
\item Carry the same output shell into the opposite orientation;
      do not replace the two kernel averages by independent maxima.
\item If concentration fails, determine whether the failure is an
      off-cut alternative already payable by V59--V60.
\end{enumerate}
The next compulsory companion audit remains V95.
\end{openproblem}

\section{Firewall after V91}
\label{sec:V91-firewall}

\begin{firewall}[Current claim boundary]
\label{fire:V91-final}
V91 closes every Cartan-split thermal core column and replaces the
general multiplicity trace by an exact fusion-kernel gate with the
output sum inside the first-stage maximum.  It does not yet prove the
required fusion-kernel concentration on all selected heavy paths,
close the off-cut alternative, synchronize both orientations, prove
JREC, aggregate all punctures and collisions, construct the continuum
Yang--Mills theory, or prove a positive mass gap.
\end{firewall}

\section{Append-only record for V91}
\label{sec:V91-record}

V91 was appended to the complete V90 manuscript.  The exact V90
parent had \(62583\) physical lines, \(2199028\) bytes, and
SHA--256 digest
\[
 \texttt{2517f5a46784bb97757d66a5bb437362f747b91c7af83a96930364161b067355}.
\]
No V90 mathematical content was changed.  The sole final
\verb|\end{document}| is restored below.

% =============================================================================
% END APPENDED VERSION V91 MATERIAL
% =============================================================================

% =============================================================================
% BEGIN APPENDED VERSION V92 MATERIAL
% =============================================================================

\chapter{V92: Failure of Raw Fusion Concentration and the
Product-Visible Fusion Capacity}
\label{ch:V92}

\section{Purpose and placement in the selected-heavy ledger}
\label{sec:V92-purpose}

In one retained-row orientation, \(A\) and \(B\) are the two retained
punctured-core row representations, \(C\) is the core row traced in
the square, \(S\subset A\otimes B\) is the first-stage path, and
\(T\subset S\otimes C\) is the final core output.  The opposite
orientation cyclically changes which core row is traced.  The
internal-pair level operator \(Q_e\) remains tensorized outside this
core carrier, exactly as in V87--V91.

V91 proposed bounding the joint shell by
\[
 d_C\sup_S\sum_{T\in\Omega}
 \frac{\mathsf p_C(T\mid S)}
 {[1+s_\alpha\Xi(T)]^2}.
\]
This chapter proves that this raw first-stage maximum is not uniformly
bounded, even on a one-output boundary Littlewood--Richardson family.
The failure is caused by an entangled \(AB\)-singlet which no product
state can see with unit probability.  The original product carrier
remains bounded because its singlet visibility supplies the missing
inverse dimension.

The correct object is therefore not a maximum over all \(S\), but the
fusion average over first-stage probabilities actually attainable by
one \(A|B\) product state.  An exact threshold estimate reduces that
object to high-fusion sectors weighted by their product visibility.

\section{A sharp singlet regression for the raw maximum}
\label{sec:V92-singlet-regression}

Let \(U\) be an irreducible unitary representation of dimension
\(d_U\).  Take
\[
 A=U,\qquad B=U^*,\qquad C=U,
\]
and use the first-stage singlet \(S=\mathbf1\subset U\otimes U^*\).

\begin{proposition}[Deterministic singlet fusion destroys the raw
maximum]
\label{prop:V92-raw-max-fails}
For the final output \(T=U\),
\begin{equation}
 \mathsf p_U(U\mid\mathbf1)=1.
 \label{eq:V92-deterministic-singlet-fusion}
\end{equation}
Consequently, whenever \(s_\alpha\Xi(U)\le1\),
\begin{equation}
 d_U\sup_{S\subset U\otimes U^*}
 \sum_{T\in\{U\}}
 \frac{\mathsf p_U(T\mid S)}
 {[1+s_\alpha\Xi(T)]^2}
 \ge\frac{d_U}{4}.
 \label{eq:V92-raw-max-divergence}
\end{equation}
For \(U=R_m\), the right side diverges like \(m^2\).
\end{proposition}

\begin{proof}
The tensor product \(\mathbf1\otimes U\) is exactly \(U\), with
multiplicity one.  Hence
\[
 \mathsf p_U(U\mid\mathbf1)
 =
 \frac{d_UN_{\mathbf1,U}^{U}}{d_{\mathbf1}d_U}=1.
\]
The weight denominator is at most \(4\) under the stated thermal
condition.  Finally
\(d_{R_m}=(m+1)(m+2)/2\).
\end{proof}

\begin{firewall}[Scope of the regression]
\label{fire:V92-regression-scope}
The proposition disproves a representation-uniform bound for the raw
maximum in the V92 mandate.  It does not prove failure of the original
product-state capacity.  The remaining input \(U\) may be a thermal
core or spectator while a disjoint puncture coordinate carries the
selected-heavy mass, so the V59 mass algebra alone does not remove
this regression.
\end{firewall}

\section{The exact compensating singlet visibility}
\label{sec:V92-singlet-visibility}

\begin{lemma}[Product norm of the dual singlet]
\label{lem:V92-singlet-product-norm}
Let \(P_{\mathbf1}^{U,U^*}\) be the rank-one singlet projector in
\(U\otimes U^*\).  Then
\begin{equation}
 \boxed{
 \|P_{\mathbf1}^{U,U^*}\|_{U\otimes U^*,\mathrm{prod}}
 =
 \frac1{d_U}.}
 \label{eq:V92-singlet-product-norm}
\end{equation}
\end{lemma}

\begin{proof}
In dual orthonormal bases, a normalized singlet is
\[
 \Omega_U=\frac1{\sqrt{d_U}}
 \sum_{j=1}^{d_U}e_j\otimes e_j^*.
\]
For unit vectors \(x\in U\), \(y\in U^*\),
\[
 |\langle x\otimes y,\Omega_U\rangle|^2
 =
 \frac1{d_U}|y(x)|^2
 \le\frac1{d_U}.
\]
Equality holds when \(y\) is the metric dual of \(x\).
\end{proof}

\begin{proposition}[The actual singlet path carrier is bounded]
\label{prop:V92-singlet-carrier-bounded}
The pathwise partial trace for
\[
 U\otimes U^*\to\mathbf1,\qquad
 \mathbf1\otimes U\to U
\]
is
\begin{equation}
 \operatorname{Tr}_{U}\mathbf P_{U;\mathbf1}
 =
 d_U P_{\mathbf1}^{U,U^*},
 \label{eq:V92-singlet-path-carrier}
\end{equation}
and therefore
\begin{equation}
 \left\|
 \operatorname{Tr}_{U}\mathbf P_{U;\mathbf1}
 \right\|_{\mathrm{prod}}
 =1.
 \label{eq:V92-singlet-carrier-unit}
\end{equation}
Thus the factor \(d_U\) in
\eqref{eq:V92-raw-max-divergence} is cancelled exactly by product
visibility.
\end{proposition}

\begin{proof}
The pathwise Schur trace formula
\eqref{eq:V90-path-partial-trace} applies with
\(d_T=d_U\) and \(d_S=1\), proving
\eqref{eq:V92-singlet-path-carrier}.  Apply
\eqref{eq:V92-singlet-product-norm}.
\end{proof}

\section{The exact product-visible fusion functional}
\label{sec:V92-visible-functional}

For \(S\subset A\otimes B\), define its product visibility
\begin{equation}
 \zeta_{A|B}(S)
 :=
 \|P_{AB}(S)\|_{A\otimes B,\mathrm{prod}}.
 \label{eq:V92-first-stage-visibility}
\end{equation}
For an output shell \(\Omega\) and nonnegative weights \(a_T\), put
\begin{equation}
 F_{\Omega,a}(S)
 :=
 \sum_{T\in\Omega}a_T\mathsf p_C(T\mid S).
 \label{eq:V92-fusion-profile}
\end{equation}

\begin{theorem}[Exact common-product fusion functional]
\label{thm:V92-exact-visible-functional}
One has
\begin{equation}
 \boxed{
 \left\|
 \sum_{T\in\Omega}a_T\mathbf N_T
 \right\|_{A\otimes B,\mathrm{prod}}
 =
 d_C
 \sup_{\substack{\|v_A\|=\|v_B\|=1}}
 \sum_S
 F_{\Omega,a}(S)
 \pi_{v_A\otimes v_B}(S).}
 \label{eq:V92-exact-visible-functional}
\end{equation}
In particular, the first-stage probabilities in the supremum are
generated by one common product vector; they are not arbitrary
probability measures on \(S\).
\end{theorem}

\begin{proof}
The positive operator on the left is block scalar by
\eqref{eq:V91-general-carrier-spectrum}.  Evaluating it on
\(v_A\otimes v_B\) and using
\eqref{eq:V91-product-fusion-mixture} gives the right side.
The product-state supremum may be taken over pure states because the
expectation is affine separately in each density matrix and extreme
points are rank one.
\end{proof}

\begin{definition}[Product-visible fusion capacity]
\label{def:V92-visible-fusion-capacity}
For a nonnegative profile \(F\) on the finite \(AB\) spectrum, set
\begin{align}
 Z_{A|B;F}(t)
 &:=
 \sum_{\{S:F(S)>t\}}\zeta_{A|B}(S),
 \label{eq:V92-visible-level-mass}\\
 \mathcal V_{A|B}(F)
 &:=
 \int_0^{\|F\|_\infty}
 \min\{1,Z_{A|B;F}(t)\}\,\dd t.
 \label{eq:V92-visible-fusion-capacity}
\end{align}
\end{definition}

\begin{theorem}[Visibility-capacity bound]
\label{thm:V92-visibility-capacity-bound}
For every output shell and nonnegative weight bank,
\begin{equation}
 \boxed{
 \left\|
 \sum_{T\in\Omega}a_T\mathbf N_T
 \right\|_{\mathrm{prod}}
 \le
 d_C\mathcal V_{A|B}(F_{\Omega,a}).}
 \label{eq:V92-visibility-capacity-bound}
\end{equation}
Equivalently, for every threshold \(\lambda\ge0\),
\begin{equation}
 \left\|
 \sum_{T\in\Omega}a_T\mathbf N_T
 \right\|_{\mathrm{prod}}
 \le
 d_C\left[
 \lambda+
 \sum_{\{S:F_{\Omega,a}(S)>\lambda\}}
 F_{\Omega,a}(S)\zeta_{A|B}(S)
 \right].
 \label{eq:V92-threshold-visible-bound}
\end{equation}
\end{theorem}

\begin{proof}
For a fixed product vector \(v\), layer cake gives
\[
 \sum_SF(S)\pi_v(S)
 =
 \int_0^{\|F\|_\infty}
 \sum_{\{S:F(S)>t\}}\pi_v(S)\,\dd t.
\]
The inner sum is at most one.  It is also at most
\[
 \sum_{\{S:F(S)>t\}}\zeta_{A|B}(S),
\]
because \(\pi_v(S)\le\zeta_{A|B}(S)\) for every \(S\).
Integrate, take the product supremum, and use
\eqref{eq:V92-exact-visible-functional}.  This proves
\eqref{eq:V92-visibility-capacity-bound}.

For the threshold form, split
\(\sum_SF(S)\pi_v(S)\) into \(F\le\lambda\) and
\(F>\lambda\).  The first part is at most \(\lambda\); bound every
\(\pi_v(S)\) in the second by \(\zeta_{A|B}(S)\).
\end{proof}

\begin{corollary}[Corrected joint pair--core gate]
\label{cor:V92-corrected-pair-core-gate}
Suppose
\[
 0\le M_T\le a_T\mathbf N_T
\]
and let \(Q_e\) be an internal-pair operator.  Then
\begin{equation}
 \boxed{
 \left\|Q_e\otimes\sum_{T\in\Omega}M_T\right\|_{\mathrm{prod}}
 \le
 A_*d_C
 \mathcal V_{A|B}(F_{\Omega,a}).}
 \label{eq:V92-corrected-pair-core-gate}
\end{equation}
For connected squares one may again take
\[
 a_T=\|\widetilde K_T\|_{\mathrm{op}}^2
 \le\frac{C_K^2}{[1+s_\alpha\Xi(T)]^2}.
\]
\end{corollary}

\begin{proof}
Sum the positive carrier majorants, use
\Cref{thm:V92-visibility-capacity-bound}, and apply
\(0\le Q_e\le A_*I\) as in
\Cref{thm:V91-joint-shell-fusion-gate}.
\end{proof}

\begin{assessment}[Disposition of the singlet branch]
\label{ass:V92-singlet-disposition}
On the deterministic singlet path,
\[
 F_{\{U\},1}(\mathbf1)=1,
 \qquad
 \zeta_{U|U^*}(\mathbf1)=d_U^{-1}.
\]
Its visible capacity is \(d_U^{-1}\), and the prefactor \(d_C=d_U\)
in \eqref{eq:V92-corrected-pair-core-gate} gives order one.  This is
the product-visible version of the V87 rank-one singlet gain.
Therefore the failure of the raw maximum is not an off-cut
obstruction; it is already paid by the retained product geometry.
\end{assessment}

\section{V92 ledger and V93 mandate}
\label{sec:V92-ledger}

\begin{theorem}[Integrated V92 statement]
\label{thm:V92-integrated}
Preserving V1--V91, V92 proves:
\begin{enumerate}[label=\textup{(V92.\arabic*)},leftmargin=5.7em]
\item the deterministic singlet regression
      \eqref{eq:V92-raw-max-divergence};
\item the exact dual-singlet product visibility
      \eqref{eq:V92-singlet-product-norm};
\item cancellation of the raw dimension in the actual path carrier
      \eqref{eq:V92-singlet-carrier-unit};
\item the exact common-product fusion functional
      \eqref{eq:V92-exact-visible-functional};
\item the level-set visibility capacity
      \eqref{eq:V92-visible-fusion-capacity};
\item its threshold estimate
      \eqref{eq:V92-threshold-visible-bound}; and
\item the corrected joint pair--core gate
      \eqref{eq:V92-corrected-pair-core-gate}.
\end{enumerate}
\end{theorem}

\begin{proof}
These are
\Cref{prop:V92-raw-max-fails,
lem:V92-singlet-product-norm,
prop:V92-singlet-carrier-bounded,
thm:V92-exact-visible-functional,
thm:V92-visibility-capacity-bound,
cor:V92-corrected-pair-core-gate}.
\end{proof}

\begin{openproblem}[V93 mandate: first-stage visibility tails]
\label{op:V93-mandate}
Apply \eqref{eq:V92-corrected-pair-core-gate} to the complete
same-label regrouping in V60.
\begin{enumerate}
\item Compute or bound
      \(\zeta_{A|B}(S)=\|P_{AB}(S)\|_{\rm prod}\) for the boundary
      Littlewood--Richardson families which maximize the fusion
      profile.
\item Separate rank-one dual singlets, Cartan summands, and
      intermediate multiplicity sectors before summing their
      visibility.
\item Prove a capped level estimate for
      \(Z_{A|B;F}(t)\), with the selected-heavy path and puncture
      indicators retained.
\item Use the same product vector and same final output shell in the
      opposite orientation; independent visibility sums would repeat
      the V60 alternating regression.
\item Insert the resulting capacity into JREC only before the single
      resolvent is used.
\item If the visibility tail is too large, exhibit the responsible
      product-visible LR family and decide whether V59's
      cut-heavy/private-mean alternatives already pay it.
\end{enumerate}
The next compulsory companion audit remains V95.
\end{openproblem}

\section{Firewall after V92}
\label{sec:V92-firewall}

\begin{firewall}[Current claim boundary]
\label{fire:V92-final}
V92 disproves the raw fusion-maximum concentration requested by V91
and replaces it by a rigorous product-visible fusion capacity.  The
singlet regression is paid exactly, but the visibility tails of
general multiplicity-rich LR sectors are not yet bounded.  V92 does
not prove JREC, close both orientations or the off-cut alternative,
aggregate all punctures and collisions, construct the continuum
Yang--Mills theory, or prove a positive mass gap.
\end{firewall}

\section{Append-only record for V92}
\label{sec:V92-record}

V92 was appended to the complete V91 manuscript.  The exact V91
parent had \(63203\) physical lines, \(2217948\) bytes, and
SHA--256 digest
\[
 \texttt{f890332d4e38a8e998d379282cfcba110da7d908e457c6a62cbd6fea70750cda}.
\]
No V91 mathematical content was changed.  The sole final
\verb|\end{document}| is restored below.

% =============================================================================
% END APPENDED VERSION V92 MATERIAL
% =============================================================================

% =============================================================================
% BEGIN APPENDED VERSION V93 MATERIAL
% =============================================================================

\chapter{V93: Visibility Provenance, Trace Collapse, and the
Same-Coordinate No-Go}
\label{ch:V93}

\section{Purpose and provenance correction}
\label{sec:V93-purpose}

V92 introduced
\(\zeta_{A|B}(S)=\|P_{AB}(S)\|_{\rm prod}\).  This is not a new
channel invariant: it is exactly the V76 product visibility
\(\beta_{AB}(S)\) from \eqref{eq:V76-beta}.  Likewise, V76 already
proved the dimension cap
\[
 \beta_{AB}(S)
 \le
 \min\left\{1,
 \frac{d_SN_{AB}^{S}}{\max\{d_A,d_B\}}\right\}
\]
and the setwise ancestry split
\eqref{eq:V76-visible-composite-bound}.  V92's genuinely new content
is the fusion-profile specialization
\eqref{eq:V92-exact-visible-functional} and its continuous level-set
capacity \eqref{eq:V92-visible-fusion-capacity}.

This chapter asks whether substituting the V76 dimension cap into that
new capacity closes the V60 branch.  It does not: after the sums are
evaluated, the result is exactly the V73 multiplicity trace divided by
the larger retained dimension.  Thus no new gain can come from
dimension counting alone.

A second no-go then prevents use of the total selected-heavy mass to
improve visibility.  Product visibility is multiplicative over
disjoint group coordinates, and a product-visible Cartan spectator
can carry arbitrarily large Casimir mass with visibility one.  Any
decay must therefore be tied to the same coordinate which carries the
fusion/Racah profile.

\begin{remark}[Identification of notation]
\label{rem:V93-zeta-beta}
Throughout V93,
\begin{equation}
 \zeta_{A|B}(S)=\beta_{AB}(S).
 \label{eq:V93-zeta-beta}
\end{equation}
All dimension-cap references are to
\Cref{thm:V76-visibility-cap}; they are not claimed as V92 results.
\end{remark}

\section{Exact collapse of the dimension-only visibility estimate}
\label{sec:V93-trace-collapse}

Put
\[
 D_{AB}:=\max\{d_A,d_B\},
 \qquad
 F(S):=\sum_{T\in\Omega}a_T\mathsf p_C(T\mid S),
\]
with \(a_T\ge0\).

\begin{theorem}[Dimension visibility reproduces the multiplicity
trace]
\label{thm:V93-dimension-trace-collapse}
The V76 dimension cap gives
\begin{align}
 d_C\sup_{v_A,v_B}
 \sum_SF(S)\pi_{v_A\otimes v_B}(S)
 &\le
 \frac{d_C}{D_{AB}}
 \sum_Sd_SN_{AB}^{S}F(S)
 \label{eq:V93-dimension-visible-bound}\\
 &=
 \boxed{
 \frac1{D_{AB}}
 \sum_{T\in\Omega}
 a_Td_Tm_{ABC}(T).}
 \label{eq:V93-trace-collapse}
\end{align}
The right side is precisely the shell Hilbert--Schmidt trace majorant
divided by the larger retained representation dimension.
\end{theorem}

\begin{proof}
For every product vector,
\[
 \pi_v(S)\le\beta_{AB}(S)
 \le\frac{d_SN_{AB}^{S}}{D_{AB}}.
\]
Multiply by \(F(S)\), sum, and take the supremum to prove the first
line.  Now insert
\[
 F(S)
 =
 \sum_{T\in\Omega}
 a_T\frac{d_TN_{SC}^{T}}{d_Sd_C}.
\]
Cancellation of \(d_Sd_C\) gives
\[
 \frac1{D_{AB}}
 \sum_{S,T\in\Omega}
 a_Td_TN_{AB}^{S}N_{SC}^{T}.
\]
The associativity identity
\[
 m_{ABC}(T)=\sum_SN_{AB}^{S}N_{SC}^{T}
\]
proves \eqref{eq:V93-trace-collapse}.
\end{proof}

\begin{corollary}[No dimension-only repair of the V89 trace loss]
\label{cor:V93-no-dimension-repair}
For the balanced V89 heat family, the right side of
\eqref{eq:V93-trace-collapse}, with
\(a_{T_k}=\|\widetilde K_{T_k}\|_{\rm op}^2\), permits order
\(\sqrt m\) on the thermal window.  Hence substituting only the V76
dimension cap into the V92 visibility capacity cannot recover V90's
order-one product-carrier bound.
\end{corollary}

\begin{proof}
Here \(D_{AB}=d_m\),
\(m_{ABC}(T_k)=k+1\), and \(d_{T_k}=(k+1)^3\).  Thus
\eqref{eq:V93-trace-collapse} is the V88 sum
\[
 \sum_k
 \frac{d_{T_k}(k+1)}{d_m}
 \|\widetilde K_{T_k}\|_{\rm op}^2.
\]
V89 proves that the corresponding heat trace is genuinely of order
\(\sqrt m\).  V90 improves it only by computing the actual
product-carrier geometry before this dimension substitution.
\end{proof}

\begin{assessment}[The exact missing datum]
\label{ass:V93-missing-datum}
The dimension cap knows only the trace of \(P_{AB}(S)\).  It cannot
distinguish a high-dimensional Cartan channel with visibility one
from a high-dimensional non-Cartan channel whose product variety meets
the isotypic space weakly.  Closing the V92 capacity requires either:
\begin{enumerate}[label=\textup{(\roman*)}]
\item an actual estimate for \(\beta_{AB}(S)\) below its dimension cap
      on the high-fusion level set; or
\item a direct bound on
      \(\sum_SF(S)\pi_v(S)\) which preserves the signed/Racah
      structure generating \(F\).
\end{enumerate}
Replacing either datum by \(d_SN_{AB}^{S}\) returns exactly to V73.
\end{assessment}

\section{Spectator inflation cannot improve product visibility}
\label{sec:V93-spectator-no-go}

Consider two disjoint group-coordinate banks.  Write
\[
 A=A_1\boxtimes A_2,\qquad
 B=B_1\boxtimes B_2,\qquad
 S=S_1\boxtimes S_2.
\]
Under the regrouping
\[
 (A_1\otimes B_1)\boxtimes(A_2\otimes B_2)
 \cong
 (A_1\boxtimes A_2)\otimes(B_1\boxtimes B_2),
\]
the isotypic projector is \(P_{S_1}\boxtimes P_{S_2}\).

\begin{lemma}[Tensor lower stability of visibility]
\label{lem:V93-visibility-tensor-lower}
One always has
\begin{equation}
 \boxed{
 \beta_{A|B}(S_1\boxtimes S_2)
 \ge
 \beta_{A_1|B_1}(S_1)
 \beta_{A_2|B_2}(S_2).}
 \label{eq:V93-visibility-tensor-lower}
\end{equation}
\end{lemma}

\begin{proof}
Choose unit product vectors
\(u_i\otimes v_i\) approaching the two separate suprema.  Then
\[
 (u_1\boxtimes u_2)\otimes(v_1\boxtimes v_2)
\]
is an admissible product vector across the combined \(A|B\) cut, and
its projector expectation factors.  Let the two approximation errors
tend to zero.
\end{proof}

\begin{proposition}[Cartan-spectator inflation no-go]
\label{prop:V93-Cartan-spectator-no-go}
Fix any first-bank channel \(S_1\) with
\(\beta_{A_1|B_1}(S_1)>0\).  On a disjoint \(SU(3)\) coordinate take
\[
 A_2=R_N,\qquad B_2=R_N,\qquad S_2=R_{2N}.
\]
Then
\begin{equation}
 \beta_{A_2|B_2}(S_2)=1,
 \qquad
 \beta_{A|B}(S_1\boxtimes S_2)
 \ge\beta_{A_1|B_1}(S_1),
 \label{eq:V93-Cartan-spectator-visibility}
\end{equation}
while the combined balanced mass satisfies
\begin{equation}
 \Xi(S_1\boxtimes S_2)
 \ge1+C_2(R_{2N})\longrightarrow\infty.
 \label{eq:V93-Cartan-spectator-mass}
\end{equation}
Consequently there is no representation-uniform bound
\[
 \beta_{A|B}(S)\le\varepsilon(s_\alpha\Xi(S))
\]
with \(\varepsilon(x)\to0\), based only on total channel mass.
\end{proposition}

\begin{proof}
The tensor product of highest-weight vectors of
\(R_N\otimes R_N\) lies entirely in the Cartan summand \(R_{2N}\),
so its visibility is one.  Apply
\Cref{lem:V93-visibility-tensor-lower}.  The Casimir formula gives
\[
 C_2(R_{2N})=\frac{4N^2+6N}{3},
\]
which proves the mass divergence.  Choosing \(N\) so that
\(s_\alpha C_2(R_{2N})\to\infty\) contradicts any proposed decaying
bound while leaving the first-bank visibility fixed below.
\end{proof}

\begin{corollary}[Fusion profiles can be spectator-inflated as well]
\label{cor:V93-fusion-spectator-no-go}
Take the remaining spectator input \(C_2=\mathbf1\) and final
spectator output \(T_2=S_2\).  Then
\[
 \mathsf p_{\mathbf1}(S_2\mid S_2)=1.
\]
Tensoring any first-bank fusion event with this deterministic Cartan
event leaves its fusion probability and its nonzero product visibility
unchanged from below, while its total mass diverges.  Hence the
selected-heavy inequality \(s_\alpha\Xi\gg1\), if carried only by a
disjoint spectator coordinate, cannot bound the V92 visible fusion
capacity.
\end{corollary}

\begin{proof}
Fusion with the trivial representation is deterministic.  Fusion
probabilities on product groups factor, and
\eqref{eq:V93-visibility-tensor-lower} supplies the visibility lower
bound.  Use \eqref{eq:V93-Cartan-spectator-mass}.
\end{proof}

\begin{firewall}[Only same-coordinate mass may be used]
\label{fire:V93-same-coordinate}
The V59--V60 ledger distinguishes the collision pair, punctured core,
residual puncture pairs, and spectators before trace norm.  V93 proves
that this distinction is essential for visibility: mass on a
disjoint Cartan spectator cannot be converted into a product-visibility
debit.  A valid V94 estimate must couple the high fusion/Racah level
to the same coordinate bank that supplies the selected-heavy mass, or
use a previously proved private-mean/cut-heavy payment for that bank.
\end{firewall}

\section{A nonduplicative reformulation of the remaining gate}
\label{sec:V93-remaining-gate}

Return to the full V76 law
\(\lambda_{u,v}(S,V,T)\), which retains the multiplicity state and
Racah POVM.  Let \(F_{\rm JREC}(S,V,T)\) denote the complete
nonnegative square majorant obtained only after the two collision
pairs, same-label puncture pairs, and core cut have been assembled.
The remaining fixed-orientation capacity is
\begin{equation}
 \mathfrak V_{\rm JREC}^{A|B}
 :=
 d_C\sup_{u,v}
 \sum_{S,V,T}
 \lambda_{u,v}(S,V,T)
 F_{\rm JREC}(S,V,T).
 \label{eq:V93-full-visible-JREC-gate}
\end{equation}
Unlike the simplified fusion profile in V92, this expression does not
discard the multiplicity density \(\sigma_S(u,v)\) or maximize over
the Racah outcome before the product-channel law is applied.

\begin{assessment}[Exact acquisition after V93]
\label{ass:V93-exact-acquisition}
The new work required for JREC is not another proof of the V76
dimension cap.  It is a same-coordinate level estimate for
\eqref{eq:V93-full-visible-JREC-gate} on the simultaneous-sideways and
puncture-pair-heavy region.  It must show that a high value of the
assembled profile is accompanied by one of:
\begin{enumerate}[label=\textup{(\roman*)}]
\item small actual product-channel mass \(p_S(u,v)\);
\item a same-coordinate selected-heavy pair capacity;
\item a private singleton mean; or
\item a cut-heavy original-row influence.
\end{enumerate}
The last three alternatives already have proved payments in
V52--V60.  Only the first requires new product/Racah geometry.
\end{assessment}

\section{V93 ledger and V94 mandate}
\label{sec:V93-ledger}

\begin{theorem}[Integrated V93 statement]
\label{thm:V93-integrated}
Preserving V1--V92, V93 proves:
\begin{enumerate}[label=\textup{(V93.\arabic*)},leftmargin=5.7em]
\item the identification \(\zeta_{A|B}=\beta_{AB}\) with the prior
      V76 invariant;
\item exact collapse of the dimension-only visibility estimate to the
      V73 multiplicity trace
      \eqref{eq:V93-trace-collapse};
\item impossibility of repairing the V89 square-root loss by that
      cap alone;
\item tensor lower stability of product visibility
      \eqref{eq:V93-visibility-tensor-lower};
\item the Cartan-spectator inflation no-go
      \eqref{eq:V93-Cartan-spectator-visibility}; and
\item the full same-coordinate visible JREC target
      \eqref{eq:V93-full-visible-JREC-gate}.
\end{enumerate}
\end{theorem}

\begin{proof}
These are
\Cref{rem:V93-zeta-beta,
thm:V93-dimension-trace-collapse,
cor:V93-no-dimension-repair,
lem:V93-visibility-tensor-lower,
prop:V93-Cartan-spectator-no-go,
ass:V93-exact-acquisition}.
\end{proof}

\begin{openproblem}[V94 mandate: same-coordinate Racah level split]
\label{op:V94-mandate}
Work inside one completely regrouped V60 monomial and one retained-row
orientation.
\begin{enumerate}
\item Keep the V76 law
      \(\lambda_{u,v}(S,V,T)\) intact through the complete assembled
      JREC square profile.
\item Define its level sets before maximizing over \(S,V,T\).
\item On each high level, apply the exact V59 incidence partition and
      V60 puncture alphabet to allocate the mass on the same
      coordinate bank.
\item Remove private-mean, cut-heavy, and selected-heavy-pair pieces
      using their existing theorems, without charging them again.
\item On the residual piece, prove or refute a capped bound for the
      actual probability
      \(\sum_{(S,V,T)\in E_t}\lambda_{u,v}(S,V,T)\).
\item Test the bound on the dual singlet and Cartan-split families,
      and on a spectator-inflated copy; the latter must be invariant
      after spectators are traced.
\item State the second orientation with the same local level variable,
      preparing the first genuinely joint two-orientation estimate.
\end{enumerate}
V95 is the next compulsory companion audit.
\end{openproblem}

\section{Firewall after V93}
\label{sec:V93-firewall}

\begin{firewall}[Current claim boundary]
\label{fire:V93-final}
V93 proves that dimension-only visibility is exactly the old trace
branch and that total selected-heavy mass cannot improve visibility
through disjoint spectators.  It isolates the same-coordinate
product/Racah probability as the remaining new datum.  It does not
yet bound that probability on every V60 level, prove JREC, close both
orientations or the off-cut alternative, aggregate all punctures and
collisions, construct the continuum Yang--Mills theory, or prove a
positive mass gap.
\end{firewall}

\section{Append-only record for V93}
\label{sec:V93-record}

V93 was appended to the complete V92 manuscript.  The exact V92
parent had \(63610\) physical lines, \(2230591\) bytes, and
SHA--256 digest
\[
 \texttt{69c807606d17e35ddec8e127f9643b4221c00961fa75f181485e701e747e91bc}.
\]
No V92 mathematical content was changed.  The sole final
\verb|\end{document}| is restored below.

% =============================================================================
% END APPENDED VERSION V93 MATERIAL
% =============================================================================

% =============================================================================
% BEGIN APPENDED VERSION V94 MATERIAL
% =============================================================================

\chapter{V94: Spectator Quotient, Residual Racah Levels, and the
Common Two-Orientation Target}
\label{ch:V94}

\section{Purpose}
\label{sec:V94-purpose}

V93 showed that total mass and dimension counting cannot control the
remaining product/Racah probability.  This chapter makes the
same-coordinate restriction exact.  First, arbitrary row states may
be entangled between active and spectator coordinates, but partial
trace still removes identity spectators without changing the
product-state supremum.  A resolved spectator effect can only decrease
the active capacity.  Thus spectator inflation is quotiented before
any level estimate.

Second, the complete assembled nonnegative JREC profile is written by
layer cake against the exact V76 law.  The V59 incidence partition and
V60 puncture alphabet give a disjoint, pre-trace classification of
every level into already-paid private-mean, cut-heavy, and
selected-heavy-pair pieces plus one residual piece.  Only the residual
level occupation is a new estimate.

Finally, the two retained-row orientations are put at the same level
parameter as positive effects on the original row tensor product.
They need not commute, so no fictitious joint Racah probability is
introduced.  The remaining target takes one product-state supremum
after adding the two effects.

\section{Exact removal of spectators}
\label{sec:V94-spectator-quotient}

Write the retained row spaces as
\[
 H_A=H_A^{\rm act}\otimes H_A^{\rm sp},
 \qquad
 H_B=H_B^{\rm act}\otimes H_B^{\rm sp}.
\]

\begin{lemma}[Identity-spectator quotient]
\label{lem:V94-identity-spectator-quotient}
For every positive operator \(M_{\rm act}\) on
\(H_A^{\rm act}\otimes H_B^{\rm act}\),
\begin{equation}
 \boxed{
 \|M_{\rm act}\otimes I_{A^{\rm sp}}\otimes I_{B^{\rm sp}}\|_
 {A\otimes B,\mathrm{prod}}
 =
 \|M_{\rm act}\|_{A^{\rm act}\otimes B^{\rm act},\mathrm{prod}}.}
 \label{eq:V94-identity-spectator-quotient}
\end{equation}
The tensor factors are regrouped by retained row.
\end{lemma}

\begin{proof}
For product density matrices \(\rho_A\otimes\rho_B\), trace out the
spectator factor within each row:
\[
 \rho_A^{\rm act}:=\operatorname{Tr}_{A^{\rm sp}}\rho_A,
 \qquad
 \rho_B^{\rm act}:=\operatorname{Tr}_{B^{\rm sp}}\rho_B.
\]
The expectation of the left-side operator is exactly
\(\operatorname{Tr}[M_{\rm act}
(\rho_A^{\rm act}\otimes\rho_B^{\rm act})]\), giving the upper bound.
Tensor active optimizing states with arbitrary spectator states for
the reverse bound.
\end{proof}

\begin{lemma}[Resolved spectators are contractive]
\label{lem:V94-resolved-spectator-contractive}
If \(0\le E_{\rm sp}\le I\) is any positive spectator effect, then
\begin{equation}
 \|M_{\rm act}\otimes E_{\rm sp}\|_{A\otimes B,\mathrm{prod}}
 \le
 \|M_{\rm act}\|_{A^{\rm act}\otimes B^{\rm act},\mathrm{prod}}.
 \label{eq:V94-resolved-spectator-contractive}
\end{equation}
\end{lemma}

\begin{proof}
The operator order
\[
 0\le M_{\rm act}\otimes E_{\rm sp}
 \le M_{\rm act}\otimes I_{\rm sp}
\]
and \Cref{lem:V94-identity-spectator-quotient} prove the claim.
\end{proof}

\begin{theorem}[Spectator invariance of the active V76 marginal]
\label{thm:V94-V76-spectator-marginal}
Suppose the completely resolved V76/Racah measurement is performed on
an active coordinate bank and all spectator outcomes are summed.
For every product row state \(\rho_A\otimes\rho_B\), its marginal law
is exactly the active V76 law evaluated in the reduced product state
\[
 \rho_A^{\rm act}\otimes\rho_B^{\rm act}.
 \]
Consequently the supremum of every active level probability is
unchanged by adjoining arbitrary spectators.  Conditioning on a
resolved spectator channel can only decrease it.
\end{theorem}

\begin{proof}
Summing the spectator POVM outcomes gives the spectator identity.
Every active event effect therefore has the form
\(E_{\rm act}\otimes I_{\rm sp}\).  Apply
\Cref{lem:V94-identity-spectator-quotient}.  A resolved spectator
channel replaces the identity by an effect between zero and one, so
\Cref{lem:V94-resolved-spectator-contractive} applies.
\end{proof}

\begin{corollary}[Disposition of the V93 inflation test]
\label{cor:V94-inflation-disposition}
After the spectator coordinates are marginalized, the Cartan
spectator of \Cref{prop:V93-Cartan-spectator-no-go} disappears from
the active level law.  It may increase total \(\Xi\), but it neither
increases nor decreases the active level capacity.  Thus only active
same-coordinate mass may enter a V94 debit.
\end{corollary}

\begin{proof}
Use \Cref{thm:V94-V76-spectator-marginal}.  The mass statement is the
content of \Cref{prop:V93-Cartan-spectator-no-go}.
\end{proof}

\section{Exact layer cake for the assembled JREC profile}
\label{sec:V94-JREC-layer-cake}

Fix one completely regrouped V60 monomial and one retained-row
orientation.  Let
\[
 \lambda_{u,v}(\omega),
 \qquad
 \omega=(S,V,T),
\]
be the exact V76 probability law, and let
\(F_{\rm JREC}(\omega)\ge0\) be the complete assembled square profile
after the collision pairs, same-label puncture pairs, and core cut
have been combined.  Put
\begin{align}
 E_t&:=\{\omega:F_{\rm JREC}(\omega)>t\},
 \label{eq:V94-JREC-level-set}\\
 \Lambda_{u,v}(t)&:=
 \sum_{\omega\in E_t}\lambda_{u,v}(\omega).
 \label{eq:V94-JREC-level-probability}
\end{align}

\begin{theorem}[Exact Racah level identity]
\label{thm:V94-Racah-level-identity}
For every product vector,
\begin{equation}
 \boxed{
 d_C\sum_\omega
 \lambda_{u,v}(\omega)F_{\rm JREC}(\omega)
 =
 d_C\int_0^{\|F_{\rm JREC}\|_\infty}
 \Lambda_{u,v}(t)\,\dd t.}
 \label{eq:V94-Racah-level-identity}
\end{equation}
Hence the fixed-orientation product capacity is bounded by
\begin{equation}
 d_C\int_0^{\|F_{\rm JREC}\|_\infty}
 \sup_{u,v}\Lambda_{u,v}(t)\,\dd t.
 \label{eq:V94-Racah-level-capacity}
\end{equation}
\end{theorem}

\begin{proof}
Use
\[
 F_{\rm JREC}(\omega)
 =
 \int_0^{\|F_{\rm JREC}\|_\infty}
 \boldsymbol1_{\{\omega\in E_t\}}\,\dd t
\]
and Tonelli's theorem.  Taking the supremum of the integral is at most
the integral of the suprema, proving the second statement.
\end{proof}

\section{Disjoint removal of already-paid level branches}
\label{sec:V94-paid-level-removal}

For each resolved outcome in \(E_t\), apply the V59 incidence selector
and then the V60 puncture alphabet before taking an absolute value.
This gives a deterministic disjoint partition
\begin{equation}
 E_t
 =
 E_t^{\rm mean}\mathbin{\dot\cup}
 E_t^{\rm cut}\mathbin{\dot\cup}
 E_t^{\rm pair}\mathbin{\dot\cup}
 E_t^{\rm res},
 \label{eq:V94-level-partition}
\end{equation}
where the first three sets respectively satisfy a private-mean
hypothesis, a cut-heavy hypothesis, or a same-coordinate
selected-heavy pair hypothesis.  The residual set contains only the
simultaneous-sideways remainder and the explicit off-cut alternative.

\begin{lemma}[Exact probability decomposition at every level]
\label{lem:V94-level-probability-decomposition}
For every product vector and every \(t\),
\begin{equation}
 \Lambda_{u,v}(t)
 =
 \Lambda_{u,v}^{\rm mean}(t)
 +\Lambda_{u,v}^{\rm cut}(t)
 +\Lambda_{u,v}^{\rm pair}(t)
 +\Lambda_{u,v}^{\rm res}(t),
 \label{eq:V94-level-probability-decomposition}
\end{equation}
with all four terms nonnegative.
\end{lemma}

\begin{proof}
The incidence sets in
\eqref{eq:V59-a-mass-split}--\eqref{eq:V59-r-mass-split} are disjoint
coordinate sets.  On the puncture residual,
\eqref{eq:V60-puncture-mass-split} is likewise a disjoint local
alphabet.  Fix the priority order mean, cut, pair, residual to assign
an outcome satisfying more than one paid condition to its first
condition.  This makes \eqref{eq:V94-level-partition} disjoint and
exhaustive.  Additivity of the positive law \(\lambda_{u,v}\) proves
the identity.
\end{proof}

\begin{theorem}[No-double-charge residual reduction]
\label{thm:V94-residual-reduction}
In a fixed V60 monomial, the level integrals over
\[
 E_t^{\rm mean},\qquad E_t^{\rm cut},\qquad E_t^{\rm pair}
\]
are assigned, respectively, to
\Cref{lem:V59-private-mean-response,
thm:V59-cut-heavy-promotion,
prop:V87-selected-heavy-capacity}
after complete same-label regrouping.  They must then be removed from
the new level ledger.  The only new fixed-orientation quantity is
\begin{equation}
 \boxed{
 \mathcal C_{\rm res}^{A|B}
 :=
 d_C\int_0^{\|F_{\rm JREC}\|_\infty}
 \sup_{u,v}
 \Lambda_{u,v}^{\rm res}(t)\,\dd t.}
 \label{eq:V94-residual-capacity}
\end{equation}
If this quantity has the target V60 path bound, then the assembled
square precursor is closed in that orientation without reusing any
paid debit.
\end{theorem}

\begin{proof}
Insert \eqref{eq:V94-level-probability-decomposition} in
\eqref{eq:V94-Racah-level-identity}.  The private-mean set has the
exact disjoint singleton bank required by
\Cref{lem:V59-private-mean-response}.  The cut set has the fixed
incidence fraction required by
\Cref{thm:V59-cut-heavy-promotion}.  The pair set is first regrouped
by its same label as required in
\Cref{rem:V60-no-monomial-tax}, after which the selected-heavy
capacity is \eqref{eq:V87-selected-heavy-capacity}.  These payments
belong to disjoint outcomes under the priority partition and are each
used once.  What remains is exactly
\eqref{eq:V94-residual-capacity}.
\end{proof}

\begin{firewall}[The residual probability is not yet bounded]
\label{fire:V94-residual-unproved}
Equation \eqref{eq:V94-residual-capacity} is an exact reduction, not a
new estimate.  The V76 dimension cap would return to the trace stock
by \Cref{thm:V93-dimension-trace-collapse}; total selected-heavy mass
is forbidden by \Cref{prop:V93-Cartan-spectator-no-go}.  A bound must
use the actual same-coordinate residual event.
\end{firewall}

\section{Sharp tests of the reduced level capacity}
\label{sec:V94-tests}

\begin{proposition}[Three required tests]
\label{prop:V94-three-tests}
The V94 reduction has the following exact behavior.
\begin{enumerate}[label=\textup{(\roman*)}]
\item On a dual-singlet first-stage path, its level probability is
      bounded by \(1/d_U\), and the traced-row factor \(d_C=d_U\)
      is cancelled.
\item On every Cartan-split thermal family of V91, the complete
      product-carrier level integral is uniformly bounded.
\item Adjoining arbitrary spectators and summing their outcomes leaves
      the active residual capacity unchanged; resolving a spectator
      channel can only decrease it.
\end{enumerate}
\end{proposition}

\begin{proof}
Part \textup{(i)} is
\eqref{eq:V92-singlet-product-norm} together with
\eqref{eq:V92-singlet-carrier-unit}.  Part \textup{(ii)} is
\eqref{eq:V91-Cartan-pair-core-uniform}.  Part \textup{(iii)} is
\Cref{thm:V94-V76-spectator-marginal}.
\end{proof}

\section{One level parameter for two noncommuting orientations}
\label{sec:V94-two-orientation}

Let \(o\in\{ab,ac\}\) denote two retained-row orientations.  For each
orientation, let \(\mathsf E_t^o\) be the positive event effect whose
expectation in an original-row product state is the residual level
probability after all spectators and paid pieces have been removed.
The two effect families act on the same original row tensor product
but need not commute.

\begin{definition}[Common-level two-orientation capacity]
\label{def:V94-two-orientation-capacity}
With the appropriate traced-row dimensions \(d_{C_o}\), define
\begin{equation}
 \boxed{
 \mathcal C_{\rm res}^{(2)}
 :=
 \int_0^{A_{\rm JREC}}
 \sup_{\rho_{\rm rows}^{\rm prod}}
 \operatorname{Tr}\left[
 \rho_{\rm rows}^{\rm prod}
 \left(
 d_{C_{ab}}\mathsf E_t^{ab}
 +d_{C_{ac}}\mathsf E_t^{ac}
 \right)\right]\dd t.}
 \label{eq:V94-two-orientation-capacity}
\end{equation}
Here \(A_{\rm JREC}\) is a common uniform cap for the two assembled
profiles, extending either effect by zero above its own cap.
\end{definition}

\begin{proposition}[Why no joint Racah law is needed]
\label{prop:V94-no-joint-Racah-law}
Each marginal of
\eqref{eq:V94-two-orientation-capacity} is the corresponding residual
capacity \eqref{eq:V94-residual-capacity}.  Moreover,
\begin{equation}
 \mathcal C_{\rm res}^{(2)}
 \le
 \mathcal C_{\rm res}^{ab}
 +\mathcal C_{\rm res}^{ac},
 \label{eq:V94-two-orientation-trivial-bound}
\end{equation}
but equality need not hold.  Any improvement must be proved at the
operator-effect level before taking two independent product suprema.
\end{proposition}

\begin{proof}
Set either effect to zero to recover a marginal.  For the upper bound,
the supremum of a sum is at most the sum of the suprema, followed by
integration.  No commutation or joint measurability is used.
\end{proof}

\begin{assessment}[The first genuinely joint target]
\label{ass:V94-joint-target}
The V60 alternating regression forbids paying two sideways endpoints
from scalar influence algebra.  Equation
\eqref{eq:V94-two-orientation-capacity} is the first target in the
programme which keeps:
\begin{enumerate}[label=\textup{(\roman*)}]
\item the same physical row product state;
\item the same local level parameter;
\item both noncommuting Racah event effects; and
\item the prior paid-branch removals
\end{enumerate}
before the final supremum.  Proving a saving over
\eqref{eq:V94-two-orientation-trivial-bound} is now the precise
simultaneous-sideways acquisition.
\end{assessment}

\section{V94 ledger and V95 audited mandate}
\label{sec:V94-ledger}

\begin{theorem}[Integrated V94 statement]
\label{thm:V94-integrated}
Preserving V1--V93, V94 proves:
\begin{enumerate}[label=\textup{(V94.\arabic*)},leftmargin=5.7em]
\item exact quotienting of identity spectators
      \eqref{eq:V94-identity-spectator-quotient};
\item contraction under resolved spectator effects
      \eqref{eq:V94-resolved-spectator-contractive};
\item invariance of the active V76 marginal under arbitrary
      spectator states;
\item the exact assembled Racah layer cake
      \eqref{eq:V94-Racah-level-identity};
\item the disjoint paid/residual probability decomposition
      \eqref{eq:V94-level-probability-decomposition};
\item the no-double-charge residual capacity
      \eqref{eq:V94-residual-capacity}; and
\item the common-level two-orientation target
      \eqref{eq:V94-two-orientation-capacity}.
\end{enumerate}
\end{theorem}

\begin{proof}
These are
\Cref{lem:V94-identity-spectator-quotient,
lem:V94-resolved-spectator-contractive,
thm:V94-V76-spectator-marginal,
thm:V94-Racah-level-identity,
lem:V94-level-probability-decomposition,
thm:V94-residual-reduction,
def:V94-two-orientation-capacity}.
\end{proof}

\begin{openproblem}[V95 mandate: audit and operator-effect overlap]
\label{op:V95-mandate}
First perform the compulsory read-only audit of the current SM and NS
notebooks.  Then:
\begin{enumerate}
\item express \(\mathsf E_t^{ab}\) and
      \(\mathsf E_t^{ac}\) by their common group-integral kernels;
\item estimate their overlap or anticommutator on original-row
      product states without assuming commutation;
\item preserve the V94 spectator quotient and paid-branch partition;
\item test dual singlet, Cartan split, and alternating-mass families;
\item determine whether a common-level saving closes the
      simultaneous-sideways square precursor or whether the residual
      is exactly an off-cut V59 alternative.
\end{enumerate}
\end{openproblem}

\section{Firewall after V94}
\label{sec:V94-firewall}

\begin{firewall}[Current claim boundary]
\label{fire:V94-final}
V94 removes spectators exactly, isolates the sole unpaid Racah level
occupation, and formulates a common-effect two-orientation capacity
without inventing a joint law.  It does not yet bound the residual
capacity, prove an overlap saving, prove JREC, close the off-cut
alternative, aggregate all punctures and collisions, construct the
continuum Yang--Mills theory, or prove a positive mass gap.
\end{firewall}

\section{Append-only record for V94}
\label{sec:V94-record}

V94 was appended to the complete V93 manuscript.  The exact V93
parent had \(64006\) physical lines, \(2244092\) bytes, and
SHA--256 digest
\[
 \texttt{8160352b3f5e0abb02a6b4cc8e2824d6c06ed31d6f536099749c72617f766352}.
\]
No V93 mathematical content was changed.  The sole final
\verb|\end{document}| is restored below.

% =============================================================================
% END APPENDED VERSION V94 MATERIAL
% =============================================================================
% =============================================================================
% BEGIN APPENDED VERSION V95 MATERIAL
% =============================================================================

\chapter{V95: Audited Racah Event Kernels, Noncommuting Overlap,
and the Levelwise-Supremum Firewall}
\label{ch:V95}

\section{Purpose and compulsory companion audit}
\label{sec:V95-purpose-audit}

V94 isolated two positive residual event effects on the same original-row
carrier.  The first task here is the compulsory five-version read-only audit.
The live companion files were stable under repeated reads:
\begin{center}
\begin{tabular}{|p{0.43\textwidth}|r|r|p{0.27\textwidth}|}
\hline
Companion file & Lines & Bytes & SHA--256\\
\hline
\texttt{Renewal\_SM\_Einstein\_Continuum\_Research\_Notes\_V93.tex}
& \(71431\) & \(2606694\) &
\texttt{3795303205fe58586169dc8dd473157a6eedeaa2bff73e071daaf84fdf1f11d1}\\
\hline
\texttt{Renewal\_NS\_Stability\_Research\_Notes\_V31.tex}
& \(23052\) & \(887537\) &
\texttt{f1121b62238ad5491bb7cf03635ea9934942edf573ed8faaa9ee79e1d38a6696}\\
\hline
\end{tabular}
\end{center}

The SM companion has advanced from the V90 snapshot to V93.  Its current
finite counterterm block is completely enumerated and its variation map is
injective, but explicit left annihilators still detect missing metric and
Higgs directions.  The transferable point is that a finite operator
description is not a range estimate: one must identify the directions which
the proposed columns do not see.

The NS companion is byte-identical to the V90 snapshot.  Its sharp
formation tests again show that probability marking is useful only while the
actual stopping-time, signed source, and parent-tail correlation is retained.
It also insists that persistent stock not be charged again on a recent-entry
branch.

\begin{assessment}[Transfer to the V95 proof order]
\label{ass:V95-audit-transfer}
For Yang--Mills the two orientations must therefore be added as operators
before their common product-state supremum, and the paid V94 branches must
remain deleted.  In addition, the exact integral of the event effects must be
kept separate from the stronger operation which maximizes independently at
every level.  The latter is a sufficient envelope, not the exact residual
square.
\end{assessment}

\begin{firewall}[Companion separation at V95]
\label{fire:V95-companion-separation}
No counterterm theorem from the SM notebook and no stopping-time inequality
from the NS notebook is used as a Yang--Mills estimate.  The audit changes
only the order in which the finite-dimensional Yang--Mills objects below are
assembled.
\end{firewall}

\section{The residual Racah effects are finite group-integral kernels}
\label{sec:V95-event-kernels}

First quotient every disjoint identity spectator by
\Cref{lem:V94-identity-spectator-quotient}.  Let
\[
 \mathcal H_{\rm act}=H_A\otimes H_B\otimes H_C
\]
denote the remaining union of the rows active in the two orientations.  For
\(o=ab\), let
\[
 \mathcal U_{ab}:H_A\otimes H_B
 \longrightarrow
 \bigoplus_S H_S\otimes\mathcal M_{ab}^{S}
\]
be a fixed Clebsch--Gordan unitary.  Lift the V70 Racah atom to the
retained-row carrier by
\begin{equation}
 \widehat\Pi_{S}^{ab}(V,T)
 :=
 \mathcal U_{ab}^{*}
 \bigl(I_{H_S}\otimes\Pi_{S}^{ab}(V,T)\bigr)
 \mathcal U_{ab},
 \label{eq:V95-lifted-Racah-atom}
\end{equation}
and then tensor it with \(I_{H_C}\).  Define
\(\widehat\Pi_{S}^{ac}(V,T)\) cyclically and tensor it with
\(I_{H_B}\).  Impossible atoms are zero.

\begin{lemma}[The lifted atoms are a POVM on the original rows]
\label{lem:V95-lifted-POVM}
For either \(o\in\{ab,ac\}\),
\begin{equation}
 \widehat\Pi_{S}^{o}(V,T)\ge0,
 \qquad
 \sum_{S,V,T}\widehat\Pi_{S}^{o}(V,T)=I_{\mathcal H_{\rm act}}.
 \label{eq:V95-lifted-POVM}
\end{equation}
Moreover
\begin{equation}
 \sum_{V,T}\widehat\Pi_{S}^{o}(V,T)=P_{S}^{o},
 \qquad
 P_{S}^{o}\widehat\Pi_{S}^{o}(V,T)P_{S}^{o}
 =\widehat\Pi_{S}^{o}(V,T).
 \label{eq:V95-atom-support}
\end{equation}
\end{lemma}

\begin{proof}
Positivity and the sum over \(V,T\) follow from
\eqref{eq:V70-POVM-normalization} after conjugation by
\(\mathcal U_o\).  The sum is the identity on the \(S\)-isotypic block,
which proves the first identity in \eqref{eq:V95-atom-support}.  Summing in
\(S\) uses the orthogonal Clebsch--Gordan decomposition and gives
\eqref{eq:V95-lifted-POVM}.  The support identity is immediate from the
block definition.
\end{proof}

If \(E_t^{o,\rm res}\) is the V94 residual set at level \(t\), then its
effect is exactly
\begin{equation}
 \mathsf E_t^o
 =
 \sum_{(S,V,T)\in E_t^{o,\rm res}}
 \widehat\Pi_{S}^{o}(V,T).
 \label{eq:V95-residual-event-effect}
\end{equation}
In particular \(0\le\mathsf E_t^o\le I\), and its expectation in an
original-row product state is the residual V76 probability.

For \(o=ab\), put
\[
 \rho_{ab}(g)=D^A(g)\otimes D^B(g)\otimes I_{H_C},
\]
and define \(\rho_{ac}\) cyclically.  Haar measure is normalized to one.

\begin{theorem}[Exact common two-variable group kernel]
\label{thm:V95-common-event-kernel}
For \(o\in\{ab,ac\}\), define the finite operator-valued kernel
\begin{align}
 \mathsf K_t^o(g,h)
 &:=
 \sum_{(S,V,T)\in E_t^{o,\rm res}}
 d_S^2\overline{\chi_S(g)}\chi_S(h)
 \rho_o(g)\widehat\Pi_S^o(V,T)\rho_o(h)^* .
 \label{eq:V95-event-kernel}
\end{align}
Then
\begin{equation}
 \boxed{
 \mathsf E_t^o
 =
 \int_G\!\!\int_G
 \mathsf K_t^o(g,h)\,\dd g\,\dd h.}
 \label{eq:V95-event-kernel-integral}
\end{equation}
Thus both orientations live in one \(G^2\) calculus on
\(\mathcal H_{\rm act}\), without a joint Racah probability.
\end{theorem}

\begin{proof}
The isotypic projection is
\[
 P_S^o
 =
 d_S\int_G\overline{\chi_S(g)}\rho_o(g)\,\dd g
 =
 d_S\int_G\chi_S(h)\rho_o(h)^*\,\dd h.
\]
Insert these two formulas on the left and right of every atom in
\eqref{eq:V95-atom-support}, sum over the residual event, and use finite
Fubini.  This is exactly \eqref{eq:V95-event-kernel-integral}.
\end{proof}

\begin{corollary}[Exact common-kernel anticommutator]
\label{cor:V95-kernel-anticommutator}
With \(\{X,Y\}=XY+YX\),
\begin{align}
 \{\mathsf E_t^{ab},\mathsf E_t^{ac}\}
 &=
 \int_{G^4}
 \bigl[
 \mathsf K_t^{ab}(g,h)\mathsf K_t^{ac}(k,\ell)
 +
 \mathsf K_t^{ac}(k,\ell)\mathsf K_t^{ab}(g,h)
 \bigr]
 \,\dd g\,\dd h\,\dd k\,\dd\ell .
 \label{eq:V95-kernel-anticommutator}
\end{align}
No commutation is asserted.
\end{corollary}

\begin{proof}
Multiply the two finite integrals in
\eqref{eq:V95-event-kernel-integral} in both orders and use Fubini.
\end{proof}

\section{The exact integrated residual precedes the levelwise supremum}
\label{sec:V95-integrated-residual}

The deterministic residual assignment in V94 need not be independent of
\(t\).  Define the residual weight of an atom without making that
assumption:
\begin{equation}
 F_o^{\rm res}(S,V,T)
 :=
 \int_0^{A_{\rm JREC}}
 \boldsymbol1_{\{(S,V,T)\in E_t^{o,\rm res}\}}\,\dd t.
 \label{eq:V95-residual-atom-weight}
\end{equation}
Put
\begin{equation}
 \mathsf M_o^{\rm res}
 :=
 \int_0^{A_{\rm JREC}}\mathsf E_t^o\,\dd t
 =
 \sum_{S,V,T}
 F_o^{\rm res}(S,V,T)\widehat\Pi_S^o(V,T).
 \label{eq:V95-integrated-residual-effect}
\end{equation}

\begin{theorem}[Exact residual operator and the strong-envelope loss]
\label{thm:V95-exact-integrated-residual}
Let \(a=d_{C_{ab}}\) and \(b=d_{C_{ac}}\).  The actual simultaneous
two-orientation residual product capacity is
\begin{equation}
 \boxed{
 \widetilde{\mathcal C}_{\rm res}^{(2)}
 :=
 \left\|
 a\mathsf M_{ab}^{\rm res}
 +b\mathsf M_{ac}^{\rm res}
 \right\|_{\rm prod}.}
 \label{eq:V95-exact-two-orientation-capacity}
\end{equation}
It satisfies
\begin{equation}
 \widetilde{\mathcal C}_{\rm res}^{(2)}
 \le
 \mathcal C_{\rm res}^{(2)},
 \label{eq:V95-strong-level-envelope}
\end{equation}
where the right side is the V94 levelwise-supremum capacity.
Equality is not automatic.
\end{theorem}

\begin{proof}
For a fixed original-row product density \(\rho\), Tonelli and
\eqref{eq:V95-integrated-residual-effect} give
\[
 \int_0^{A_{\rm JREC}}
 \operatorname{Tr}\rho
 [a\mathsf E_t^{ab}+b\mathsf E_t^{ac}]\,\dd t
 =
 \operatorname{Tr}\rho
 [a\mathsf M_{ab}^{\rm res}+b\mathsf M_{ac}^{\rm res}].
\]
Taking one supremum after the integral gives
\eqref{eq:V95-exact-two-orientation-capacity}.  Moving the supremum inside
the integral gives \eqref{eq:V94-two-orientation-capacity}, proving
\eqref{eq:V95-strong-level-envelope}.
\end{proof}

\begin{firewall}[Do not replace the exact gate by a stronger conjecture]
\label{fire:V95-levelwise-supremum}
It is sufficient to bound \(\mathcal C_{\rm res}^{(2)}\), but failure of
that stronger levelwise envelope would not disprove the required bound on
\(\widetilde{\mathcal C}_{\rm res}^{(2)}\).  The optimizing product state
may vary with \(t\) on the right side of
\eqref{eq:V95-strong-level-envelope}.  This is exactly the kind of
correlation loss flagged by the V95 NS audit.
\end{firewall}

\section{A noncommuting two-effect overlap inequality}
\label{sec:V95-overlap}

For positive operators \(E,F\), set
\[
 \vartheta(E,F):=\|\sqrt E\sqrt F\|_{\rm op}.
\]

\begin{theorem}[Weighted noncommuting overlap bound]
\label{thm:V95-weighted-overlap}
For finite-dimensional positive operators \(E,F\) and \(a,b\ge0\),
\begin{align}
 \|aE+bF\|_{\rm op}
 &\le
 \max\{a\|E\|_{\rm op},b\|F\|_{\rm op}\}
 +\sqrt{ab}\,\vartheta(E,F),
 \label{eq:V95-weighted-overlap-op}\\
 \|aE+bF\|_{\rm prod}
 &\le
 \min\Bigl\{
 a\|E\|_{\rm prod}+b\|F\|_{\rm prod},
 \notag\\
 &\hspace{7em}
 \max\{a\|E\|_{\rm op},b\|F\|_{\rm op}\}
 +\sqrt{ab}\,\vartheta(E,F)
 \Bigr\}.
 \label{eq:V95-weighted-overlap-product}
\end{align}
No commutation of \(E\) and \(F\) is required.
\end{theorem}

\begin{proof}
Define
\[
 X:\mathcal H\oplus\mathcal H\longrightarrow\mathcal H,
 \qquad
 X(\xi,\eta)=\sqrt a\,\sqrt E\,\xi+\sqrt b\,\sqrt F\,\eta.
\]
Then \(XX^*=aE+bF\), whereas
\[
 X^*X=
 \begin{pmatrix}
 aE&\sqrt{ab}\sqrt E\sqrt F\\
 \sqrt{ab}\sqrt F\sqrt E&bF
 \end{pmatrix}.
\]
The norm of the diagonal block is
\(\max\{a\|E\|,b\|F\|\}\), and the norm of the off-diagonal block is
\(\sqrt{ab}\|\sqrt E\sqrt F\|\).  The triangle inequality and
\(\|XX^*\|=\|X^*X\|\) prove
\eqref{eq:V95-weighted-overlap-op}.  Product norm is at most operator
norm.  Its independent marginal bound follows directly from the
definition.  Take the smaller of the two valid estimates.
\end{proof}

\begin{corollary}[Level and integrated overlap envelopes]
\label{cor:V95-level-integrated-overlap}
At every level,
\begin{align}
 &\left\|
 a\mathsf E_t^{ab}+b\mathsf E_t^{ac}
 \right\|_{\rm prod}
 \notag\\
 &\quad\le
 \min\left\{
 a\|\mathsf E_t^{ab}\|_{\rm prod}
 +b\|\mathsf E_t^{ac}\|_{\rm prod},
 \max\{a,b\}
 +\sqrt{ab}\,\vartheta_t
 \right\},
 \label{eq:V95-level-overlap-envelope}\\
 &\vartheta_t:=
 \vartheta(\mathsf E_t^{ab},\mathsf E_t^{ac}).
 \notag
\end{align}
The same theorem applied to
\(\mathsf M_{ab}^{\rm res},\mathsf M_{ac}^{\rm res}\) gives
\begin{align}
 \widetilde{\mathcal C}_{\rm res}^{(2)}
 \le
 \min\Bigl\{&
 a\|\mathsf M_{ab}^{\rm res}\|_{\rm prod}
 +b\|\mathsf M_{ac}^{\rm res}\|_{\rm prod},
 \notag\\
 &
 \max\{a\|\mathsf M_{ab}^{\rm res}\|_{\rm op},
       b\|\mathsf M_{ac}^{\rm res}\|_{\rm op}\}
 +\sqrt{ab}\,
 \vartheta(\mathsf M_{ab}^{\rm res},
            \mathsf M_{ac}^{\rm res})
 \Bigr\}.
 \label{eq:V95-integrated-overlap-envelope}
\end{align}
\end{corollary}

\begin{proof}
Use \(0\le\mathsf E_t^o\le I\) in
\eqref{eq:V95-weighted-overlap-product}.  The integrated assertion is the
same inequality together with
\eqref{eq:V95-exact-two-orientation-capacity}.
\end{proof}

\begin{lemma}[A computable common-kernel overlap certificate]
\label{lem:V95-trace-overlap-certificate}
After the V94 identity-spectator quotient,
\begin{align}
 \vartheta_t^2
 &=
 \|\sqrt{\mathsf E_t^{ac}}\,
       \mathsf E_t^{ab}\,
       \sqrt{\mathsf E_t^{ac}}\|_{\rm op}
 \notag\\
 &\le
 \operatorname{Tr}_{\mathcal H_{\rm act}}
 [\mathsf E_t^{ab}\mathsf E_t^{ac}]
 \label{eq:V95-trace-overlap-certificate}\\
 &=
 \int_{G^4}
 \operatorname{Tr}_{\mathcal H_{\rm act}}
 [\mathsf K_t^{ab}(g,h)\mathsf K_t^{ac}(k,\ell)]
 \,\dd g\,\dd h\,\dd k\,\dd\ell .
 \label{eq:V95-trace-overlap-kernel}
\end{align}
The trace in \eqref{eq:V95-trace-overlap-certificate} is real and
nonnegative even though the integrand in
\eqref{eq:V95-trace-overlap-kernel} need not be pointwise real.
\end{lemma}

\begin{proof}
The sandwiched operator in the first line is positive, so its norm is at
most its trace.  Cyclicity identifies that trace with
\(\operatorname{Tr}[\mathsf E_t^{ab}\mathsf E_t^{ac}]\).
Insert \eqref{eq:V95-event-kernel-integral} twice to obtain
\eqref{eq:V95-trace-overlap-kernel}.
\end{proof}

\begin{remark}[Why the spectator quotient must precede the trace test]
\label{rem:V95-trace-after-quotient}
Tensoring both effects with an unused identity multiplies
\(\operatorname{Tr}(EF)\) by the spectator dimension, although it leaves
all product capacities unchanged.  Hence
\eqref{eq:V95-trace-overlap-certificate} is to be used only on the minimal
active union after \Cref{lem:V94-identity-spectator-quotient}.  A resolved
tensor-factor spectator can only decrease the overlap because
\(\|\sqrt R\sqrt S\|\le1\) for spectator effects \(0\le R,S\le I\).
\end{remark}

\section{Sharp tests and an overlap-only no-go}
\label{sec:V95-tests}

\begin{proposition}[Exact common-row dual-singlet test]
\label{prop:V95-dual-singlet-overlap}
Let \(U\) be a \(d\)-dimensional irreducible, let
\(H_A=U\), and let \(H_B=H_C=U^*\).  Write
\[
 P=P_{\mathbf1}^{AB}\otimes I_C,
 \qquad
 Q=P_{\mathbf1}^{AC}\otimes I_B.
\]
Then
\begin{equation}
 \|PQ\|_{\rm op}=\frac1d,
 \qquad
 \vartheta(P,Q)=\frac1d,
 \qquad
 \boxed{\|dP+dQ\|_{\rm prod}=2.}
 \label{eq:V95-dual-singlet-overlap}
\end{equation}
Thus both traced-row factors are cancelled on product states, although the
two singlet projections do not commute.
\end{proposition}

\begin{proof}
Let
\(\Omega=d^{-1/2}\sum_i e_i\otimes e_i^*\).
On the range of \(Q\), a vector has the form
\(\Omega_{AC}\otimes y_B\).  Direct contraction of the common \(A\)
index gives
\[
 \|P(\Omega_{AC}\otimes y_B)\|=\frac1d\|y_B\|.
\]
Hence \(\|PQ\|=1/d\).  Since \(P,Q\) are projections,
\(\vartheta(P,Q)=\|PQ\|\).

For a product unit vector \(x_A\otimes y_B\otimes z_C\), the Schmidt
coefficients of \(\Omega\) give
\[
 \langle P\rangle\le\frac1d,\qquad
 \langle Q\rangle\le\frac1d.
\]
Equality holds simultaneously for
\(x=e_i,y=e_i^*,z=e_i^*\).  Multiplication by \(d\) proves the last
identity.
\end{proof}

\begin{proposition}[Even zero overlap does not supply inverse dimension]
\label{prop:V95-zero-overlap-no-go}
For every \(d\ge1\) there are two orthogonal rank-one product projections
\(P,Q\) such that
\begin{equation}
 \vartheta(P,Q)=0,
 \qquad
 \|P\|_{\rm prod}=\|Q\|_{\rm prod}=1,
 \qquad
 \|dP+dQ\|_{\rm prod}=d.
 \label{eq:V95-zero-overlap-no-go}
\end{equation}
Consequently an overlap estimate alone cannot cancel a traced-row
dimension.  It can synchronize two orientations only after a
fixed-orientation product-visibility or symbol debit is present.
\end{proposition}

\begin{proof}
Choose two orthogonal vectors from a product basis and let \(P,Q\) be their
rank-one projections.  Their ranges are orthogonal, so the overlap is zero.
Each range contains a product vector, giving product norm one.  The sum has
largest eigenvalue \(d\), attained on either of those product vectors.
\end{proof}

\begin{assessment}[Disposition of the Cartan test]
\label{ass:V95-Cartan-test}
The Cartan highest-weight event has unit product visibility, so
\Cref{prop:V95-zero-overlap-no-go} forbids treating small orientation
overlap as its missing debit.  The balanced thermal Cartan pair-core shell is
instead uniformly bounded by the exact carrier/symbol calculation
\eqref{eq:V91-Cartan-pair-core-uniform}.  Under the V94 priority partition
that contribution belongs to the already-paid pair branch and may not be
charged again as an overlap gain.
\end{assessment}

\begin{proposition}[Alternating scalar mass does not determine effect
overlap]
\label{prop:V95-alternating-overlap-independence}
The scalar data in the V60 alternating-height regression do not determine
\(\vartheta_t\).  At the level of positive-effect axioms, identical
one-orientation ranks, traces, operator norms, and product norms are
compatible both with \(\vartheta_t=1\) and with \(\vartheta_t=0\).
\end{proposition}

\begin{proof}
On a carrier with two orthogonal product vectors, fix the rank-one product
projection \(P\).  In the first realization take
\((\mathsf E_t^{ab},\mathsf E_t^{ac})=(P,P)\), which has overlap one.  In
the second take \((P,Q)\), where \(Q\) is the orthogonal product projection
from \Cref{prop:V95-zero-overlap-no-go}; its overlap is zero.  The two
effects in both realizations have the same rank, trace, operator norm, and
product norm.  The V60 quantities
\(\Xi_a,\Xi_r,\widehat H_{ar}\) contain none of the relative range data
distinguishing these realizations.  This is an abstract regression for the
mass algebra, not a claim that either arbitrary pair is a Yang--Mills Racah
effect.
\end{proof}

\begin{assessment}[What the overlap calculation decides]
\label{ass:V95-overlap-decision}
The common-kernel coefficient
\(\vartheta_t\) is a genuine, computable noncommutative datum.  It is not
controlled by the alternating scalar masses.  More importantly, even its
best possible value zero leaves one full traced-row dimension on a
unit-visible residual event.  Therefore the V95 overlap route cannot by
itself close the simultaneous-sideways precursor.  The attack must go
upstream: first obtain the missing fixed-orientation inverse-dimension
debit from the complete residual symbol, or prove that failure of that
debit forces the off-cut V59 alternative.  Only then can overlap synchronize
the two surviving orientations.
\end{assessment}

\section{V95 ledger and upstream V96 mandate}
\label{sec:V95-ledger}

\begin{theorem}[Integrated V95 statement]
\label{thm:V95-integrated}
Preserving V1--V94, V95 proves:
\begin{enumerate}[label=\textup{(V95.\arabic*)},leftmargin=5.7em]
\item the compulsory SM V93 and NS V31 audit;
\item the lifted original-row Racah POVM
      \eqref{eq:V95-lifted-POVM};
\item the exact common group-integral event kernels
      \eqref{eq:V95-event-kernel-integral} and their noncommuting
      anticommutator \eqref{eq:V95-kernel-anticommutator};
\item the exact integrated residual operator
      \eqref{eq:V95-integrated-residual-effect};
\item the distinction
      \eqref{eq:V95-strong-level-envelope} between the required product
      capacity and the stronger levelwise-supremum envelope;
\item the weighted noncommuting overlap inequality
      \eqref{eq:V95-weighted-overlap-product};
\item the common-kernel trace certificate
      \eqref{eq:V95-trace-overlap-kernel};
\item the exact common-row dual-singlet identity
      \eqref{eq:V95-dual-singlet-overlap}; and
\item the zero-overlap and alternating-mass no-go results
      \eqref{eq:V95-zero-overlap-no-go} and
      \Cref{prop:V95-alternating-overlap-independence}.
\end{enumerate}
\end{theorem}

\begin{proof}
These are
\Cref{ass:V95-audit-transfer,
lem:V95-lifted-POVM,
thm:V95-common-event-kernel,
cor:V95-kernel-anticommutator,
thm:V95-exact-integrated-residual,
thm:V95-weighted-overlap,
lem:V95-trace-overlap-certificate,
prop:V95-dual-singlet-overlap,
prop:V95-zero-overlap-no-go,
prop:V95-alternating-overlap-independence}.
\end{proof}

\begin{openproblem}[V96 mandate: upstream fixed-orientation residual symbol]
\label{op:V96-mandate}
Do not seek closure from the anticommutator alone.  Work upstream of the
V94 levelwise supremum as follows.
\begin{enumerate}
\item Keep the exact integrated operators
      \(\mathsf M_o^{\rm res}\) and one product-state supremum.
\item Reinsert the complete assembled JREC symbol into the V95 lifted atoms
      before replacing it by threshold events.
\item Use the V78 centred group integral for the connected Wilson factor and
      the V59/V60 deterministic residual conditions on the same coordinate
      bank.
\item Prove a fixed-orientation dichotomy: either
      \(d_{C_o}\|\mathsf M_o^{\rm res}\|_{\rm prod}\) has the target path
      bound, or the same outcome carries the heavy off-cut bank in
      \eqref{eq:V59-one-sideways-residual}.
\item Preserve the exact spectator quotient and do not restore a trace over
      deleted identity factors.
\item Test the dichotomy on the dual singlet, the unit-visible Cartan lift,
      the V60 alternating array, and the V89 balanced thermal family.
\item Apply the V95 overlap inequality only after the fixed-orientation
      debit is proved; do not count the V94 paid branches again.
\end{enumerate}
The next compulsory companion audit is V100.
\end{openproblem}

\section{Firewall after V95}
\label{sec:V95-firewall}

\begin{firewall}[Current claim boundary]
\label{fire:V95-final}
V95 gives exact common group kernels for both residual orientations, proves
a noncommuting overlap inequality, identifies the stronger
levelwise-supremum loss, and proves that overlap alone cannot provide the
missing inverse-dimension factor.  It does not yet prove the upstream
fixed-orientation residual-symbol dichotomy, close the off-cut alternative,
prove JREC, aggregate all punctures and collisions, construct the continuum
Yang--Mills theory, or prove a positive mass gap.
\end{firewall}

\section{Append-only record for V95}
\label{sec:V95-record}

V95 was appended to the complete V94 manuscript after the compulsory
companion audit.  The exact V94 parent had \(64464\) physical lines,
\(2260562\) bytes, and SHA--256 digest
\[
 \texttt{b4e48dc6ef69e39aeb804b49832d2d8e7c23d0511100b685dd3ff772784068f9}.
\]
No V94 mathematical content was changed.  The sole final
\verb|\end{document}| is restored below.

% =============================================================================
% END APPENDED VERSION V95 MATERIAL
% =============================================================================
% =============================================================================
% BEGIN APPENDED VERSION V96 MATERIAL
% =============================================================================

\chapter{V96: Canonical Racah Dilation, Orthogonal Off-Cut Splitting,
and the Exact Injective-Image Gate}
\label{ch:V96}

\section{Purpose}
\label{sec:V96-purpose}

V95 proves that orientation overlap cannot create the missing
fixed-orientation inverse-dimension debit.  This chapter therefore goes
upstream of the levelwise supremum.  For one orientation, the complete
residual V76 POVM is placed in its canonical analysis dilation.  The
simultaneous-sideways and off-cut branches then become genuinely orthogonal
output coordinates, even though their compressed effects on the physical
rows need not be projections or commute.

The exact product capacity is next written as an injective tensor norm of
the adjoint analysis image.  This separates two mechanisms without
confusing them: a dual singlet is small because every analysis-image vector
is entangled across the retained rows, whereas a Cartan highest-weight
range contains a product vector and must instead be small through its
assembled symbol.  Finally, the V78 centred cumulant estimate is extended
to arbitrary noncommuting positive filters.  The remaining issue is stated
as an exact image estimate, not as another dimension count.

\section{Canonical dilation of the full residual law}
\label{sec:V96-dilation}

Fix one orientation \(o=A|B\), after applying the V94 identity-spectator
quotient, and put \(H_o=H_A\otimes H_B\).  Let
\[
 \Omega_o=\{(S,V,T):\widehat\Pi_S^o(V,T)\ne0\},
 \qquad
 E_\omega^o:=\widehat\Pi_S^o(V,T).
\]
The identity on the row traced in this orientation is suppressed.  By
\Cref{lem:V95-lifted-POVM},
\[
 E_\omega^o\ge0,\qquad \sum_{\omega\in\Omega_o}E_\omega^o=I_{H_o}.
\]
Define
\[
 \mathcal K_o:=\bigoplus_{\omega\in\Omega_o}H_o
\]
and let \(Q_\omega\) be its coordinate projections.

\begin{theorem}[Canonical analysis dilation of the Racah POVM]
\label{thm:V96-canonical-dilation}
The map
\begin{equation}
 \mathcal V_o:H_o\longrightarrow\mathcal K_o,
 \qquad
 (\mathcal V_o\psi)_\omega=(E_\omega^o)^{1/2}\psi,
 \label{eq:V96-analysis-isometry}
\end{equation}
is an isometry and
\begin{equation}
 \boxed{
 E_\omega^o=\mathcal V_o^*Q_\omega\mathcal V_o.}
 \label{eq:V96-Naimark-compression}
\end{equation}
For arbitrary nonnegative weights \(w_\omega\),
\begin{equation}
 \sum_\omega w_\omega E_\omega^o
 =
 \mathcal V_o^*
 \left(\bigoplus_\omega w_\omega I_{H_o}\right)
 \mathcal V_o.
 \label{eq:V96-weighted-dilation}
\end{equation}
\end{theorem}

\begin{proof}
For \(\psi\in H_o\),
\[
 \|\mathcal V_o\psi\|^2
 =
 \sum_\omega\langle\psi,E_\omega^o\psi\rangle
 =
 \|\psi\|^2.
\]
Thus \(\mathcal V_o^*\mathcal V_o=I\).  Keeping only the
\(\omega\)-coordinate gives
\[
 \langle\phi,\mathcal V_o^*Q_\omega\mathcal V_o\psi\rangle
 =
 \langle(E_\omega^o)^{1/2}\phi,
        (E_\omega^o)^{1/2}\psi\rangle
 =
 \langle\phi,E_\omega^o\psi\rangle,
\]
which proves \eqref{eq:V96-Naimark-compression}.  Sum with weights.
\end{proof}

\begin{remark}[Why the dilation is not a fictitious joint law]
\label{rem:V96-no-joint-law}
The construction is made separately for one physical orientation and is
the canonical dilation of its already-proved V70 POVM.  It does not couple
the \(ab\) and \(ac\) outcome labels and therefore does not assert joint
measurability of the two physical effects.
\end{remark}

\section{A level-augmented analysis operator}
\label{sec:V96-level-analysis}

Let
\[
 \mathcal L_o
 :=
 L^2\!\left(
 [0,A_{\rm JREC}];
 \mathcal K_o
 \right).
\]
For \(X\in\{\mathrm{res},\mathrm{side},\mathrm{off}\}\), define
\begin{equation}
 (\mathcal A_o^X\psi)(t)_\omega
 :=
 \boldsymbol1_{\{\omega\in E_t^{o,X}\}}
 (E_\omega^o)^{1/2}\psi.
 \label{eq:V96-level-analysis-operator}
\end{equation}
Here the V94 residual family is partitioned, before trace norm, by
\begin{equation}
 E_t^{o,\mathrm{res}}
 =
 E_t^{o,\mathrm{side}}
 \mathbin{\dot\cup}
 E_t^{o,\mathrm{off}},
 \label{eq:V96-side-off-partition}
\end{equation}
where the first set is the genuine simultaneous-sideways remainder and the
second contains the explicit heavy off-cut alternative left in
\eqref{eq:V59-one-sideways-residual}.  The priority convention is inherited
from \eqref{eq:V94-level-partition}.

\begin{theorem}[Exact level-analysis factorization]
\label{thm:V96-level-analysis-factorization}
For each \(X\),
\begin{equation}
 (\mathcal A_o^X)^*\mathcal A_o^X
 =
 \int_0^{A_{\rm JREC}}
 \sum_{\omega\in E_t^{o,X}}E_\omega^o\,\dd t
 =:\mathsf M_o^X.
 \label{eq:V96-analysis-Gram}
\end{equation}
In particular
\begin{equation}
 \mathsf M_o^{\rm res}
 =
 \mathsf M_o^{\rm side}+\mathsf M_o^{\rm off},
 \qquad
 \langle\mathcal A_o^{\rm side}\phi,
        \mathcal A_o^{\rm off}\psi\rangle_{\mathcal L_o}=0.
 \label{eq:V96-orthogonal-side-off}
\end{equation}
The orthogonality holds in the dilation even when the two compressed
operators on \(H_o\) are not orthogonal.
\end{theorem}

\begin{proof}
For \(\phi,\psi\in H_o\), Tonelli gives
\begin{align*}
 \langle\mathcal A_o^X\phi,\mathcal A_o^X\psi\rangle
 &=
 \int_0^{A_{\rm JREC}}
 \sum_{\omega\in E_t^{o,X}}
 \langle(E_\omega^o)^{1/2}\phi,
        (E_\omega^o)^{1/2}\psi\rangle\,\dd t\\
 &=
 \left\langle\phi,
 \left[
 \int_0^{A_{\rm JREC}}
 \sum_{\omega\in E_t^{o,X}}E_\omega^o\,\dd t
 \right]\psi\right\rangle.
\end{align*}
This is \eqref{eq:V96-analysis-Gram}.  The disjoint partition
\eqref{eq:V96-side-off-partition} gives the operator sum.  At every
\((t,\omega)\), at most one of the two indicators is nonzero, proving
orthogonality in \(\mathcal L_o\).
\end{proof}

\begin{corollary}[No-double-charge side/off reduction]
\label{cor:V96-side-off-reduction}
For each orientation,
\begin{equation}
 d_{C_o}\|\mathsf M_o^{\rm res}\|_{\rm prod}
 \le
 d_{C_o}\|\mathsf M_o^{\rm side}\|_{\rm prod}
 +
 d_{C_o}\|\mathsf M_o^{\rm off}\|_{\rm prod}.
 \label{eq:V96-side-off-product-bound}
\end{equation}
The first term contains no private-mean, cut-heavy, or selected-heavy-pair
outcome already removed in V94.  The second term is precisely the off-cut
branch and is not silently reassigned to the simultaneous-sideways term.
\end{corollary}

\begin{proof}
Use positivity, \eqref{eq:V96-orthogonal-side-off}, and the triangle
inequality for the product norm.  The branch statement is the defining
priority partition.
\end{proof}

\section{The exact injective-image formula}
\label{sec:V96-injective-image}

For \(\xi\in H_A\otimes H_B\), define its injective vector norm by
\begin{equation}
 \|\xi\|_\varepsilon
 :=
 \sup_{\|u\|=\|v\|=1}
 |\langle\xi,u\otimes v\rangle|.
 \label{eq:V96-injective-vector-norm}
\end{equation}

\begin{theorem}[Product capacity is an exact adjoint-image norm]
\label{thm:V96-exact-image-norm}
Let \(\mathcal A:H_A\otimes H_B\to\mathcal L\) be any linear map and
let \(M=\mathcal A^*\mathcal A\).  Then
\begin{equation}
 \boxed{
 \|M\|_{\rm prod}
 =
 \sup_{\|z\|_{\mathcal L}=1}
 \|\mathcal A^*z\|_\varepsilon^2.}
 \label{eq:V96-exact-image-norm}
\end{equation}
Applied to \(\mathcal A_o^X\), this gives
\begin{equation}
 d_{C_o}\|\mathsf M_o^X\|_{\rm prod}
 =
 d_{C_o}
 \sup_{\|z\|_{\mathcal L_o}=1}
 \|(\mathcal A_o^X)^*z\|_\varepsilon^2.
 \label{eq:V96-exact-JREC-image}
\end{equation}
\end{theorem}

\begin{proof}
Positive product norm is attained after restricting to pure product states,
and hence
\begin{align*}
 \|M\|_{\rm prod}
 &=
 \sup_{\|u\|=\|v\|=1}
 \|\mathcal A(u\otimes v)\|^2\\
 &=
 \sup_{\|u\|=\|v\|=1}
 \sup_{\|z\|=1}
 |\langle z,\mathcal A(u\otimes v)\rangle|^2\\
 &=
 \sup_{\|z\|=1}
 \sup_{\|u\|=\|v\|=1}
 |\langle\mathcal A^*z,u\otimes v\rangle|^2.
\end{align*}
The two suprema commute because they range over a Cartesian product.
The last line is \eqref{eq:V96-exact-image-norm}.  Use
\eqref{eq:V96-analysis-Gram} for the application.
\end{proof}

Let \(P_{\mathcal A}\) be the orthogonal projection onto
\(\operatorname{Ran}\mathcal A^*\), and put
\begin{equation}
 \eta_{\rm ran}(\mathcal A)
 :=
 \sup_{\substack{\xi\in\operatorname{Ran}\mathcal A^*\\
                  \|\xi\|=1}}
 \|\xi\|_\varepsilon^2.
 \label{eq:V96-range-visibility}
\end{equation}

\begin{lemma}[Range visibility and amplitude factorization]
\label{lem:V96-range-amplitude}
One has
\begin{align}
 \eta_{\rm ran}(\mathcal A)
 &=
 \|P_{\mathcal A}\|_{\rm prod},
 \label{eq:V96-range-projector-visibility}\\
 \|\mathcal A^*\mathcal A\|_{\rm prod}
 &\le
 \|\mathcal A\|_{\rm op}^2
 \eta_{\rm ran}(\mathcal A).
 \label{eq:V96-range-amplitude-bound}
\end{align}
\end{lemma}

\begin{proof}
For a product unit vector \(p=u\otimes v\),
\[
 \langle p,P_{\mathcal A}p\rangle
 =
 \|P_{\mathcal A}p\|^2
 =
 \sup_{\substack{\xi\in\operatorname{Ran}\mathcal A^*\\
                  \|\xi\|=1}}
 |\langle\xi,p\rangle|^2.
\]
Taking the product-vector supremum and commuting the two suprema proves
\eqref{eq:V96-range-projector-visibility}.  For every unit \(z\),
\[
 \|\mathcal A^*z\|_\varepsilon^2
 \le
 \eta_{\rm ran}(\mathcal A)\|\mathcal A^*z\|^2
 \le
 \eta_{\rm ran}(\mathcal A)\|\mathcal A\|_{\rm op}^2.
\]
Apply \Cref{thm:V96-exact-image-norm}.
\end{proof}

\begin{corollary}[Recovery and extension of the V76 visibility]
\label{cor:V96-recover-beta}
If \(M=P_S^{AB}\), take \(\mathcal A=P_S^{AB}\).  Then
\[
 \eta_{\rm ran}(\mathcal A)
 =
 \|P_S^{AB}\|_{\rm prod}
 =
 \beta_{AB}(S).
\]
For a weighted, nonprojective aggregate Racah effect, the exact quantity in
\eqref{eq:V96-exact-image-norm} retains the weights and multiplicity
directions inside one adjoint image; it is not a sum of the individual
\(\beta_{AB}(S)\).
\end{corollary}

\begin{proof}
The first statement is \eqref{eq:V76-beta}.  The second follows directly
from the analysis construction
\eqref{eq:V96-level-analysis-operator}.
\end{proof}

\begin{firewall}[The separated range bound can lose symbol--visibility
correlation]
\label{fire:V96-range-factorization-loss}
Equation \eqref{eq:V96-range-amplitude-bound} is a sufficient
factorization, but \eqref{eq:V96-exact-image-norm} is the sharper target.
A large-singular-value direction can have small injective norm while a
product-visible range direction has a small singular value.  Maximizing
amplitude and range visibility separately discards that correlation and may
return to the V93 dimension/trace loss.
\end{firewall}

\section{Identity spectators do not alter the image gate}
\label{sec:V96-spectator-image}

\begin{proposition}[Spectator stability of range visibility]
\label{prop:V96-spectator-range-stability}
Let \(M_{\rm act}\ge0\) act on
\(H_A^{\rm act}\otimes H_B^{\rm act}\), and let
\(P_{\rm act}\) be its support projection.  After adjoining arbitrary
identity spectators and regrouping by rows,
\begin{equation}
 \|P_{\rm act}\otimes I_{\rm sp}\|_{\rm prod}
 =
 \|P_{\rm act}\|_{\rm prod},
 \qquad
 \|M_{\rm act}\otimes I_{\rm sp}\|_{\rm prod}
 =
 \|M_{\rm act}\|_{\rm prod}.
 \label{eq:V96-spectator-image-stability}
\end{equation}
Thus neither the exact image norm nor its range-visibility factor can be
improved or worsened by the V93 Cartan spectator inflation.
\end{proposition}

\begin{proof}
Apply \Cref{lem:V94-identity-spectator-quotient} first to
\(P_{\rm act}\) and then to \(M_{\rm act}\).  The support of
\(M_{\rm act}\otimes I_{\rm sp}\) is
\(P_{\rm act}\otimes I_{\rm sp}\), so
\eqref{eq:V96-range-projector-visibility} identifies the first equality
with range-visibility stability.
\end{proof}

\section{The centred cumulant survives an arbitrary residual filter}
\label{sec:V96-filtered-cumulant}

Retain the V78 notation
\[
 Z(g):=Z_A(g)\otimes Z_B(g)\otimes Z_C(g),
 \qquad
 K=\int_G Z(g)\,\mu_\alpha(\dd g).
\]
For \(\psi\in H_A\otimes H_B\), let
\[
 W_\psi:H_C\longrightarrow H_A\otimes H_B\otimes H_C,
 \qquad
 W_\psi\xi=\psi\otimes\xi.
\]

\begin{theorem}[Sign-safe arbitrary-filter variance bound]
\label{thm:V96-arbitrary-filter}
Let \(0\le R\le rI\) be any positive operator on the three-row carrier;
it need not commute with \(K\), the output projections, or the Racah
effects.  For product \(\psi=u\otimes v\),
\begin{align}
 \mathfrak e_{R;u,v}
 &:=
 \frac1{d_C}
 \|R^{1/2}KW_{u\otimes v}\|_{\rm HS}^2
 \label{eq:V96-filtered-energy}\\
 &=
 \int_G\!\!\int_G
 \frac1{d_C}
 \operatorname{Tr}\!\left[
 W_{u\otimes v}^*
 Z(g)^*RZ(h)
 W_{u\otimes v}\right]
 \mu_\alpha(\dd g)\mu_\alpha(\dd h)
 \label{eq:V96-filtered-double-integral}\\
 &\le
 4r\min\{
 \upsilon_A\upsilon_B,\,
 \upsilon_A\upsilon_C,\,
 \upsilon_B\upsilon_C
 \}.
 \label{eq:V96-filtered-two-variance}
\end{align}
\end{theorem}

\begin{proof}
The first expression is nonnegative by construction.  Insert
\eqref{eq:V78-centred-cumulant} on both sides of \(R\) and use finite
Fubini to obtain \eqref{eq:V96-filtered-double-integral}.  Since
\(0\le R\le rI\),
\[
 \|R^{1/2}KW_{u\otimes v}\|_{\rm HS}^2
 \le
 r\|KW_{u\otimes v}\|_{\rm HS}^2.
\]
Now apply \eqref{eq:V78-two-variance-bound}.
\end{proof}

\begin{corollary}[Filtered cumulant as an injective-image problem]
\label{cor:V96-filtered-image}
Define a linear map
\[
 \mathcal B_R:H_A\otimes H_B
 \longrightarrow
 \mathcal S_2(H_C,H_A\otimes H_B\otimes H_C),
 \qquad
 \mathcal B_R\psi
 =
 d_C^{-1/2}R^{1/2}KW_\psi.
\]
Then
\begin{equation}
 \mathcal B_R^*\mathcal B_R
 =
 \frac1{d_C}\operatorname{Tr}_C(KRK),
 \label{eq:V96-filtered-Gram}
\end{equation}
and
\begin{equation}
 \left\|
 \frac1{d_C}\operatorname{Tr}_C(KRK)
 \right\|_{\rm prod}
 =
 \sup_{\|z\|_{\rm HS}=1}
 \|\mathcal B_R^*z\|_\varepsilon^2
 \le
 4r\min_{\{R_1,R_2\}\subset\{A,B,C\}}
 \upsilon_{R_1}\upsilon_{R_2}.
 \label{eq:V96-filtered-image-bound}
\end{equation}
\end{corollary}

\begin{proof}
The Hilbert--Schmidt inner-product definition of partial trace gives
\eqref{eq:V96-filtered-Gram}.  Apply
\Cref{thm:V96-exact-image-norm,thm:V96-arbitrary-filter}.
\end{proof}

\begin{assessment}[What must be proved before applying the filter theorem]
\label{ass:V96-filter-interface}
Theorem~\ref{thm:V96-arbitrary-filter} shows that a noncommuting residual
selection does not by itself destroy the centred Wilson cancellation.
However, the V93 object \(F_{\rm JREC}\) is a complete nonnegative square
majorant assembled after pair and cut estimates.  It is not automatic that
its side analysis operator factors as a contraction composed with the raw
centred cumulant map \(\mathcal B_R\).  Establishing such a factorization,
or retaining the exact pre-majorant collision operator from which it
follows, is the precise upstream interface still required.
\end{assessment}

\section{Regression tests}
\label{sec:V96-tests}

\begin{proposition}[Dual singlet is a range-entanglement debit]
\label{prop:V96-singlet-image-test}
For the singlet projector \(P_{\mathbf1}^{U,U^*}\) with \(d_U=d\),
\begin{equation}
 \eta_{\rm ran}(P_{\mathbf1})=\frac1d,
 \qquad
 \|P_{\mathbf1}\|_{\rm prod}=\frac1d.
 \label{eq:V96-singlet-range}
\end{equation}
Thus the factor \(d_C=d\) is cancelled entirely by the injective image,
with no appeal to a trace estimate.
\end{proposition}

\begin{proof}
The range is spanned by the normalized maximally entangled vector
\(\Omega=d^{-1/2}\sum_ie_i\otimes e_i^*\), whose largest Schmidt
coefficient squared is \(1/d\).  Use
\Cref{thm:V96-exact-image-norm,lem:V96-range-amplitude}.
\end{proof}

\begin{proposition}[Cartan visibility forces a symbol debit]
\label{prop:V96-Cartan-image-test}
For the Cartan highest-weight projection containing
\(u_a\otimes v_b\) from \Cref{prop:V77-Cartan-lift},
\begin{equation}
 \eta_{\rm ran}=1.
 \label{eq:V96-Cartan-range-one}
\end{equation}
Hence no entanglement or orientation-overlap estimate can make that range
small.  Any uniform bound must retain the singular-value weights of the
complete connected symbol, as in the exact image
\eqref{eq:V96-exact-image-norm}, or use the already-proved V91 pair-core
carrier calculation.
\end{proposition}

\begin{proof}
The range contains the unit product vector \(u_a\otimes v_b\), so its
support projector has product norm one.  Apply
\eqref{eq:V96-range-projector-visibility}.
\end{proof}

\begin{assessment}[Balanced thermal and alternating-mass tests]
\label{ass:V96-balanced-alternating-tests}
On the balanced thermal atlas, the exact image norm reproduces the
product-carrier quantity already bounded in V91; replacing it by
\(\|\mathcal A\|^2\eta_{\rm ran}\) can be strictly weaker and is not used
to recharge that paid branch.  On the V60 alternating array, the scalar
masses specify neither the range of \((\mathcal A_o^{\rm side})^*\) nor
its singular values.  Therefore they cannot decide
\eqref{eq:V96-exact-JREC-image}; this is the image-level version of
\Cref{prop:V95-alternating-overlap-independence}.
\end{assessment}

\section{V96 ledger and V97 acquisition}
\label{sec:V96-ledger}

\begin{theorem}[Integrated V96 statement]
\label{thm:V96-integrated}
Preserving V1--V95, V96 proves:
\begin{enumerate}[label=\textup{(V96.\arabic*)},leftmargin=5.7em]
\item the canonical dilation
      \eqref{eq:V96-Naimark-compression} of the full V76 Racah POVM;
\item the exact level-analysis Gram factorization
      \eqref{eq:V96-analysis-Gram};
\item the orthogonal simultaneous-sideways/off-cut split
      \eqref{eq:V96-orthogonal-side-off};
\item the exact injective-image identity
      \eqref{eq:V96-exact-image-norm};
\item the range-visibility factorization
      \eqref{eq:V96-range-amplitude-bound} and its correlation firewall;
\item spectator stability
      \eqref{eq:V96-spectator-image-stability};
\item the arbitrary noncommuting filter extension
      \eqref{eq:V96-filtered-two-variance} of the V78 centred estimate;
\item exact singlet and Cartan image tests
      \eqref{eq:V96-singlet-range} and
      \eqref{eq:V96-Cartan-range-one}.
\end{enumerate}
\end{theorem}

\begin{proof}
These are
\Cref{thm:V96-canonical-dilation,
thm:V96-level-analysis-factorization,
cor:V96-side-off-reduction,
thm:V96-exact-image-norm,
lem:V96-range-amplitude,
prop:V96-spectator-range-stability,
thm:V96-arbitrary-filter,
prop:V96-singlet-image-test,
prop:V96-Cartan-image-test}.
\end{proof}

\begin{openproblem}[V97 mandate: pre-majorant factorization or failure
certificate]
\label{op:V97-mandate}
Work on one fixed simultaneous-sideways monomial after the V96 off-cut
split.
\begin{enumerate}
\item Return to the exact pre-majorant collision operator before the V69
      three-square domination.
\item Construct its row-to-dilation analysis map explicitly in the V60
      collision, puncture, and final Racah labels.
\item Determine whether the side map factors through the centred cumulant
      map \(\mathcal B_R\) with a contraction carrying the two collision
      pair weights and the punctured-core weight.
\item If it does, apply \eqref{eq:V96-filtered-image-bound} while retaining
      the exact injective image; do not maximize its singular value and
      range visibility separately.
\item If it does not, exhibit a left functional or a concrete channel
      which annihilates the centred image but not the V93 majorant.  Such a
      certificate identifies the precise cancellation lost in V69.
\item Keep the off-cut analysis map separate and do not use the V94 paid
      branches.
\item Regress the resulting map on the dual singlet, Cartan lift, balanced
      thermal family, and alternating-height array.
\end{enumerate}
The next compulsory companion audit remains V100.
\end{openproblem}

\section{Firewall after V96}
\label{sec:V96-firewall}

\begin{firewall}[Current claim boundary]
\label{fire:V96-final}
V96 replaces the abstract fixed-orientation residual probability by an
exact analysis operator, separates the off-cut output orthogonally, and
identifies the required inverse-dimension debit as an injective norm of the
weighted adjoint image.  It also proves that arbitrary residual filters
preserve the V78 centred two-variance bound.  It does not yet factor the
actual pre-majorant JREC collision map through that centred operator, bound
the simultaneous-sideways image, close the off-cut alternative, prove
JREC, construct the continuum Yang--Mills theory, or prove a positive mass
gap.
\end{firewall}

\section{Append-only record for V96}
\label{sec:V96-record}

V96 was appended to the complete V95 manuscript.  The exact V95 parent had
\(65104\) physical lines, \(2282563\) bytes, and SHA--256 digest
\[
 \texttt{84d6d1afff87786765a748937cc0306b4235d29caa93ff005178cedbcafac9ac}.
\]
No V95 mathematical content was changed.  The sole final
\verb|\end{document}| is restored below.

% =============================================================================
% END APPENDED VERSION V96 MATERIAL
% =============================================================================
% =============================================================================
% BEGIN APPENDED VERSION V97 MATERIAL
% =============================================================================

\chapter{V97: The Three-Square Kernel Mismatch and the Exact
Signed-Collision Repair}
\label{ch:V97}

\section{Purpose}
\label{sec:V97-purpose}

V96 asked whether the V93 positive residual majorant factors through the
raw centred cumulant.  V97 decides this question at its earliest source.
The V69 three-square map does not factor through the five-term cumulant on
the full collision family.  There is an exact Yang--Mills channel on which
the cumulant vanishes while two covariance cuts are nonzero.  Thus the
positive majorant contains a disconnected direction which the exact
collision operator deletes.

This is not merely an abstract matrix counterexample.  It occurs for
\(U\otimes\mathbf1\otimes U^*\to\mathbf1\) under every nondegenerate
inversion-symmetric Wilson law.  The repair is forced: trivial-input blocks
must be deleted before any positive cut majorant, and the residual outcome
filters must be inserted between two copies of the complete signed
cumulant.  A multiple-filter analysis map is constructed below and placed
in the exact V96 injective-image gate.

\section{The factorization criterion for a positive square majorant}
\label{sec:V97-factorization-criterion}

On one final \(T\)-multiplicity space, abbreviate
\begin{equation}
 L_1:=D_{A\mid BC}^{T},
 \qquad
 L_2:=-yD_{AC}^{T},
 \qquad
 L_3:=-zD_{AB}^{T},
 \qquad
 K_T=L_1+L_2+L_3
 \label{eq:V97-three-cuts}
\end{equation}
as in \eqref{eq:V69-covariance-decomposition}.  Define the covariance
analysis map
\begin{equation}
 \mathcal A_{\rm cov}^{T}\xi
 :=
 (L_1\xi,L_2\xi,L_3\xi).
 \label{eq:V97-covariance-analysis}
\end{equation}
Then
\[
 (\mathcal A_{\rm cov}^{T})^*\mathcal A_{\rm cov}^{T}
 =
 (D_{A\mid BC}^{T})^2
 +y^2(D_{AC}^{T})^2
 +z^2(D_{AB}^{T})^2.
\]

\begin{theorem}[Exact finite-dimensional factorization criterion]
\label{thm:V97-factorization-criterion}
Let \(K:H\to H'\) and \(\mathcal A:H\to\mathcal L\) be linear maps between
finite-dimensional Hilbert spaces.  The following are equivalent.
\begin{enumerate}[label=\textup{(\roman*)}]
\item There is a linear map \(C:H'\to\mathcal L\) such that
      \(\mathcal A=CK\).
\item \(\ker K\subseteq\ker\mathcal A\).
\item For some finite \(c\ge0\),
      \begin{equation}
       \mathcal A^*\mathcal A\le cK^*K.
       \label{eq:V97-Douglas-inequality}
      \end{equation}
\end{enumerate}
When these statements hold, the least \(c\) is
\begin{equation}
 c_*=
 \sup_{\substack{\xi\perp\ker K\\\xi\ne0}}
 \frac{\|\mathcal A\xi\|^2}{\|K\xi\|^2}.
 \label{eq:V97-optimal-factor-constant}
\end{equation}
\end{theorem}

\begin{proof}
The implication \textup{(i)}\(\Rightarrow\)\textup{(ii)} is immediate.
Assume \textup{(ii)}.  Define \(C_0\) on \(\operatorname{Ran}K\) by
\[
 C_0(K\xi):=\mathcal A\xi.
\]
The kernel inclusion makes this well defined.  It is bounded because all
spaces are finite dimensional.  Extend it by zero on
\((\operatorname{Ran}K)^\perp\), which proves \textup{(i)}.
If \(\mathcal A=CK\), then
\[
 \|\mathcal A\xi\|^2\le\|C\|^2\|K\xi\|^2,
\]
so \textup{(iii)} holds.  Conversely, evaluating
\eqref{eq:V97-Douglas-inequality} on \(\ker K\) proves
\textup{(ii)}.  Finally, the least valid quadratic-form constant is exactly
the Rayleigh quotient in \eqref{eq:V97-optimal-factor-constant}.
\end{proof}

\begin{assessment}[Which inequality V69 supplies]
\label{ass:V97-V69-direction}
The V69 estimate is
\[
 K_T^2\le3(\mathcal A_{\rm cov}^{T})^*
                 \mathcal A_{\rm cov}^{T}.
\]
To recover the centred cancellation from the three-square map one would
need the reverse comparison
\eqref{eq:V97-Douglas-inequality}.  The next theorem shows that its optimal
constant is infinite on the full Yang--Mills channel family.
\end{assessment}

\section{An exact Wilson channel in the lost kernel}
\label{sec:V97-trivial-channel}

\begin{theorem}[Trivial-input failure certificate]
\label{thm:V97-trivial-input-certificate}
Let \(U\) be a nontrivial irreducible with
\[
 0<q_{\alpha,U}=q_{\alpha,U^*}<1.
\]
Take
\[
 A=U,\qquad B=\mathbf1,\qquad C=U^*,
 \qquad T=\mathbf1,
\]
and use the unique fusion path.  Put
\(\delta_U:=1-q_{\alpha,U}^2>0\).  Then
\begin{align}
 D_{A\mid BC}^{T}&=\delta_UI,
 &
 D_{AC}^{T}&=\delta_UI,
 &
 D_{AB}^{T}&=0,
 \label{eq:V97-trivial-cut-values}\\
 K_T&=0,
 &
 (\mathcal A_{\rm cov}^{T})^*
 \mathcal A_{\rm cov}^{T}
 &=2\delta_U^2I.
 \label{eq:V97-kernel-mismatch}
\end{align}
Consequently
\begin{equation}
 \ker K_T\not\subseteq\ker\mathcal A_{\rm cov}^{T},
 \label{eq:V97-kernel-inclusion-fails}
\end{equation}
and there is no finite \(c\) for
\((\mathcal A_{\rm cov}^{T})^*\mathcal A_{\rm cov}^{T}
\le cK_T^2\).
\end{theorem}

\begin{proof}
Here \(x=z=q_{\alpha,U}\) and \(y=q_{\alpha,\mathbf1}=1\).
The \(BC\) intermediate is \(U^*\), the \(AC\) intermediate is the
singlet, and the \(AB\) intermediate is \(U\).  Hence
\[
 q_{\alpha,T}=1,\qquad
 Q_{BC}^{T}=zI,\qquad
 Q_{AC}^{T}=I,\qquad
 Q_{AB}^{T}=xI.
\]
Substitution in \eqref{eq:V69-composite-cut}--%
\eqref{eq:V69-AB-defect} gives
\eqref{eq:V97-trivial-cut-values}.  Since \(y=1\),
\[
 K_T
 =
 D_{A\mid BC}^{T}-D_{AC}^{T}-zD_{AB}^{T}=0.
\]
The squared norm of the covariance analysis map is the sum of the three
squared cuts and equals \(2\delta_U^2I\).  Apply
\Cref{thm:V97-factorization-criterion}.
\end{proof}

\begin{corollary}[The V69 positive majorant is not a connected observable]
\label{cor:V97-majorant-disconnected}
On the channel in \Cref{thm:V97-trivial-input-certificate}, the exact
five-term collision observable and every filtered square \(K_TRK_T\)
vanish, whereas the V69 three-square majorant is strictly positive.
Thus no level set constructed from the three separate covariance squares
can, on the full family, be identified with a level set of the connected
cumulant square.
\end{corollary}

\begin{proof}
The first claim follows from \(K_T=0\).  The second is
\eqref{eq:V97-kernel-mismatch}.  A positive level below
\(2\delta_U^2\) contains the majorant channel but not the zero exact
observable.
\end{proof}

\begin{remark}[Consistency with the earlier exact deletion theorem]
\label{rem:V97-consistency-V77}
Equation \eqref{eq:V97-kernel-mismatch} does not contradict V77.  It
sharpens \Cref{prop:V77-trivial-input-deletion}: the complete cumulant
vanishes exactly when \(B=\mathbf1\), but its two displayed covariance
pieces cancel rather than vanish separately.
\end{remark}

\section{The connected quotient must precede positivity}
\label{sec:V97-connected-quotient}

Let \(\mathscr I_{\rm conn}\) be the set of input triples for which no input
is trivial, and let \(\mathscr I_{\rm disc}\) be its complement.

\begin{proposition}[Exact deletion of the disconnected input bank]
\label{prop:V97-delete-disconnected-bank}
In the direct sum over input labels,
\begin{equation}
 K\big|_{\mathscr I_{\rm disc}}=0.
 \label{eq:V97-disconnected-bank-zero}
\end{equation}
Therefore every exact positive collision form \(KRK\), every exact
collision analysis map, and every exact JREC trace functional receives zero
from \(\mathscr I_{\rm disc}\).  Deleting that bank before taking a
three-square or trace-norm majorant is exact and incurs no multiplicity
loss.
\end{proposition}

\begin{proof}
Apply \Cref{prop:V77-trivial-input-deletion} on every block with a trivial
input.  Direct sums preserve zero, and multiplication on either side by an
arbitrary filter \(R\) preserves zero.
\end{proof}

\begin{firewall}[Connected deletion does not prove reverse factorization]
\label{fire:V97-connected-not-factorized}
Proposition~\ref{prop:V97-delete-disconnected-bank} removes the explicit
certificate above.  It does not prove
\(\ker K_T\subseteq\ker\mathcal A_{\rm cov}^{T}\) on every remaining
nontrivial-input block.  Any reuse of the three-square majorant on the
connected bank must either verify the criterion
\eqref{eq:V97-Douglas-inequality} block by block or retain the exact signed
\(K_T\) instead.
\end{firewall}

\section{The exact multiple-filter signed analysis map}
\label{sec:V97-signed-analysis}

Let \(\{R_j\}_{j\in J}\) be any finite family of positive filters on
\(H_A\otimes H_B\otimes H_C\), let \(w_j\ge0\), and assume
\begin{equation}
 \mathcal R:=\sum_{j\in J}w_jR_j\le rI.
 \label{eq:V97-aggregate-filter}
\end{equation}
No filter is required to commute with \(K\) or with another filter.  Define
\begin{equation}
 \mathcal S_{K,\mathcal R}\psi
 :=
 \left(
 d_C^{-1/2}\sqrt{w_j}\,R_j^{1/2}KW_\psi
 \right)_{j\in J}
 \label{eq:V97-signed-analysis-map}
\end{equation}
as a map from \(H_A\otimes H_B\) to the direct sum of the corresponding
Hilbert--Schmidt spaces.

\begin{theorem}[Exact signed filtered Gram operator]
\label{thm:V97-signed-filtered-Gram}
The map \eqref{eq:V97-signed-analysis-map} satisfies
\begin{align}
 \mathcal S_{K,\mathcal R}^*
 \mathcal S_{K,\mathcal R}
 &=
 \frac1{d_C}\operatorname{Tr}_C(K\mathcal RK),
 \label{eq:V97-signed-filtered-Gram}\\
 \left\|
 \frac1{d_C}\operatorname{Tr}_C(K\mathcal RK)
 \right\|_{\rm prod}
 &=
 \sup_{\|z\|=1}
 \|\mathcal S_{K,\mathcal R}^*z\|_\varepsilon^2
 \label{eq:V97-signed-injective-image}\\
 &\le
 4r\min_{\{R_1,R_2\}\subset\{A,B,C\}}
 \upsilon_{R_1}\upsilon_{R_2}.
 \label{eq:V97-signed-filtered-variance}
\end{align}
Every block with a trivial input is zero before any norm is taken.
\end{theorem}

\begin{proof}
For \(\phi,\psi\in H_A\otimes H_B\), sum the Hilbert--Schmidt inner
products of the \(j\)-components:
\begin{align*}
 \langle\mathcal S_{K,\mathcal R}\phi,
        \mathcal S_{K,\mathcal R}\psi\rangle
 &=
 \frac1{d_C}\sum_jw_j
 \operatorname{Tr}\!\left[
 W_\phi^*KR_jKW_\psi\right]\\
 &=
 \left\langle\phi,
 \frac1{d_C}\operatorname{Tr}_C(K\mathcal RK)\psi
 \right\rangle.
\end{align*}
This proves \eqref{eq:V97-signed-filtered-Gram}.  Apply
\Cref{thm:V96-exact-image-norm} for
\eqref{eq:V97-signed-injective-image}.  Since
\(0\le\mathcal R\le rI\),
\Cref{thm:V96-arbitrary-filter} proves
\eqref{eq:V97-signed-filtered-variance}.  Exact deletion follows from
\eqref{eq:V97-disconnected-bank-zero}.
\end{proof}

\begin{corollary}[A POVM residual filter has the required aggregate cap]
\label{cor:V97-POVM-filter-cap}
If the \(R_j\) form a POVM and \(0\le w_j\le A_*\), then
\[
 0\le\mathcal R=\sum_jw_jR_j\le A_*I.
\]
Hence \eqref{eq:V97-signed-filtered-variance} holds with \(r=A_*\),
without commuting or maximizing the filters separately.
\end{corollary}

\begin{proof}
Use \(w_jR_j\le A_*R_j\) and sum the POVM identity.
\end{proof}

\begin{assessment}[The repaired pre-majorant target]
\label{ass:V97-repaired-target}
The V96 side analysis operator can factor through the exact centred
collision only if its full-carrier residual filters can be chosen so that
its Gram operator is
\[
 d_C^{-1}\operatorname{Tr}_C(K\mathcal R_{\rm side}K)
\]
with the two collision-pair and punctured-core weights still attached.
Equations \eqref{eq:V97-signed-filtered-Gram}--%
\eqref{eq:V97-signed-injective-image} then retain both signed covariance
cancellation and product-row geometry.  Replacing this form by the V69
three-square effects before applying the V96 image norm is invalid on the
full bank by \Cref{thm:V97-trivial-input-certificate}.
\end{assessment}

\section{Regression tests}
\label{sec:V97-tests}

\begin{proposition}[Four exact dispositions]
\label{prop:V97-four-dispositions}
The repaired signed analysis has the following behavior.
\begin{enumerate}[label=\textup{(\roman*)}]
\item On every trivial-input channel it vanishes exactly, including the
      explicit fundamental--trivial--dual singlet certificate.
\item On a dual-singlet retained-row range, its product capacity retains the
      \(1/d_U\) injective-image debit of
      \eqref{eq:V96-singlet-range}.
\item On the Cartan lift, its final-output decomposition is exactly the
      signed ancestral-vector formula
      \eqref{eq:V79-signed-output-formula}; no covariance square is separated.
\item On the V60 alternating array, the scalar masses do not decide whether
      the three cuts cancel in \(K_T\), so the signed map must be evaluated
      before any mass-only conclusion.
\end{enumerate}
\end{proposition}

\begin{proof}
Part \textup{(i)} is
\eqref{eq:V97-disconnected-bank-zero}.  Part \textup{(ii)} follows from
\Cref{prop:V96-singlet-image-test}.  Part \textup{(iii)} is
\Cref{cor:V79-signed-output-formula}.  For \textup{(iv)}, the V60 masses
contain no multiplicity-space matrix coefficients of the three operators
in \eqref{eq:V97-three-cuts}; therefore they cannot determine their vector
sum.
\end{proof}

\begin{assessment}[Balanced thermal warning]
\label{ass:V97-balanced-warning}
The balanced V89 family shows that the complete heat cumulant can remain
order one on a visible ancestry space.  V97 does not claim that restoring
the sign automatically closes that family.  It proves the logically prior
fact that the disconnected directions of the three-square majorant are
spurious and must not be included in the residual acquisition.
\end{assessment}

\section{V97 ledger and V98 acquisition}
\label{sec:V97-ledger}

\begin{theorem}[Integrated V97 statement]
\label{thm:V97-integrated}
Preserving V1--V96, V97 proves:
\begin{enumerate}[label=\textup{(V97.\arabic*)},leftmargin=5.7em]
\item the exact factorization criterion
      \eqref{eq:V97-Douglas-inequality};
\item the fundamental--trivial--dual failure certificate
      \eqref{eq:V97-kernel-mismatch};
\item nonfactorization of the V69 three-square map through the complete
      cumulant on the full collision bank;
\item exact deletion of all disconnected input blocks
      \eqref{eq:V97-disconnected-bank-zero};
\item the exact multiple-filter signed Gram operator
      \eqref{eq:V97-signed-filtered-Gram};
\item its injective-image and two-variance bounds
      \eqref{eq:V97-signed-injective-image}--%
      \eqref{eq:V97-signed-filtered-variance}; and
\item the four regression dispositions in
      \Cref{prop:V97-four-dispositions}.
\end{enumerate}
\end{theorem}

\begin{proof}
These are
\Cref{thm:V97-factorization-criterion,
thm:V97-trivial-input-certificate,
cor:V97-majorant-disconnected,
prop:V97-delete-disconnected-bank,
thm:V97-signed-filtered-Gram,
prop:V97-four-dispositions}.
\end{proof}

\begin{openproblem}[V98 mandate: reconstruct the physical residual filter]
\label{op:V98-mandate}
Stay upstream of the V69 three-square majorant.
\begin{enumerate}
\item On one fixed V60 collision/puncture monomial, lift the final Racah,
      pair-bank, same-label puncture, and core selections to positive
      filters \(R_j\) on the full three-row carrier.
\item Keep the complete \(K_T\) inside every filtered Gram term and verify
      the aggregate cap \eqref{eq:V97-aggregate-filter}.
\item Prove that the resulting signed analysis map controls the original
      row-factorized V60 trace functional before absolute values, or
      identify the precise polar-decomposition term which prevents it.
\item Delete the disconnected input bank first and preserve the V96
      orthogonal side/off-cut split.
\item Insert the external collision-pair operators only after the core
      signed image is defined, using the exact V72 pair--core
      factorization rather than a new pair-label probability.
\item Test the physical filter on the trivial-input certificate, dual
      singlet, Cartan lift, balanced thermal family, and alternating array.
\end{enumerate}
The next compulsory companion audit remains V100.
\end{openproblem}

\section{Firewall after V97}
\label{sec:V97-firewall}

\begin{firewall}[Current claim boundary]
\label{fire:V97-final}
V97 proves that the V69 three-square residual does not factor through the
connected cumulant on the full physical channel family and supplies an
exact Wilson-channel certificate.  It constructs the correct signed
multiple-filter analysis map, but it does not yet identify the complete
physical V60 residual filter, control the original trace norm from that
map, close the simultaneous-sideways or off-cut alternatives, prove JREC,
construct the continuum Yang--Mills theory, or prove a positive mass gap.
\end{firewall}

\section{Append-only record for V97}
\label{sec:V97-record}

V97 was appended to the complete V96 manuscript.  The exact V96 parent had
\(65748\) physical lines, \(2303268\) bytes, and SHA--256 digest
\[
 \texttt{38d90c3f1c034fd1964b4af0d63f90b9ced22c8cf1dfd057b7bd73f04d59de3e}.
\]
No V96 mathematical content was changed.  The sole final
\verb|\end{document}| is restored below.

% =============================================================================
% END APPENDED VERSION V97 MATERIAL
% =============================================================================
% =============================================================================
% BEGIN APPENDED VERSION V98 MATERIAL
% =============================================================================

\chapter{V98: The Physical Residual Absolute Carrier and the
Racah-Instrument Firewall}
\label{ch:V98}

\section{Purpose and correction of the V97 acquisition}
\label{sec:V98-purpose}

V97 proves that the V69 three-square map contains disconnected directions
which the exact cumulant deletes.  Its final mandate asked for a physical
residual filter built from the later Racah effects.  The earlier V67/V70/V72
theorems force a correction to that wording.

The physical trace-norm coefficient is controlled by the blockwise absolute
carrier \(|\widetilde K_T|\) on the complete final-output block.  The V70
Racah atoms arise only after a conditional partial trace and a change of
bracketing.  They form a valid POVM and a useful diagnostic law, but they do
not define a joint orthogonal pinching of the noncommuting \(AB\)- and
\(BC\)-intermediate observables.  Thus they cannot be inserted back into the
raw cumulant as if a physical measurement had occurred.

V98 reconstructs the correct residual object.  The V59/V60 incidence,
puncture, paid, simultaneous-sideways, and off-cut decisions are made on the
orthogonal outer collision blocks before trace norm.  Within every surviving
block the full signed \(K_T\) is retained.  The resulting absolute carrier
is then controlled directly by the V72 product-state tax.  Its exact square
and injective-image gates are derived for every retained-row choice.

\section{Outer collision blocks and the physical signed operator}
\label{sec:V98-outer-blocks}

Fix one completely regrouped V60 collision/puncture monomial.  Let
\(\mathfrak J\) be its finite family of complete orthogonal outer output
labels.  A label \(j\in\mathfrak J\) contains:
\begin{enumerate}[label=\textup{(\roman*)}]
\item the two external collision-pair outputs;
\item all already-regrouped same-label puncture outputs;
\item the punctured-core input triple;
\item the complete final irreducible \(T(j)\); and
\item every genuine orthogonal copy label required by the physical
      compression.
\end{enumerate}
It does not split one final multiplicity space into incompatible \(AB\)- and
\(BC\)-first sharp labels.

Let \(V_j:G_j\to\mathcal H_{\rm rows}\) be the physical isometry of the
complete block, with mutually orthogonal ranges.  Let
\begin{equation}
 L_j=c_j\widetilde K_{T(j)}
 \quad\hbox{on }G_j,
 \qquad
 c_j\ge0,
 \label{eq:V98-physical-block}
\end{equation}
where \(c_j\) contains the two collision-pair and same-label puncture
weights, while
\[
 \widetilde K_T=(1+s_\alpha\Xi(T))K_T
\]
is the complete signed weighted ternary cumulant.  All scalar phases from
trace-norm duality are deliberately excluded from \(L_j\); they do not
change its absolute value.

The pre-trace V59/V60 priority selector gives a disjoint partition
\begin{equation}
 \mathfrak J
 =
 \mathfrak J_{\rm paid}
 \mathbin{\dot\cup}
 \mathfrak J_{\rm side}
 \mathbin{\dot\cup}
 \mathfrak J_{\rm off}
 \mathbin{\dot\cup}
 \mathfrak J_{\rm disc}.
 \label{eq:V98-outer-partition}
\end{equation}
The last set has a trivial punctured-core input; the first contains the
private-mean, cut-heavy, and selected-heavy-pair alternatives; the middle
two are respectively the genuine simultaneous-sideways and heavy off-cut
branches.

For \(X\in\{\mathrm{paid},\mathrm{side},\mathrm{off},\mathrm{disc}\}\),
put
\begin{align}
 \mathbf L_X
 &:=
 \sum_{j\in\mathfrak J_X}V_jL_jV_j^*,
 \label{eq:V98-physical-signed-operator}\\
 \mathbf Q_X
 &:=
 \sum_{j\in\mathfrak J_X}V_j|L_j|V_j^*.
 \label{eq:V98-physical-absolute-carrier}
\end{align}

\begin{theorem}[Exact residual absolute carrier]
\label{thm:V98-exact-absolute-carrier}
For every branch \(X\),
\begin{equation}
 \boxed{
 |\mathbf L_X|=\mathbf Q_X,
 \qquad
 \mathbf L_X^2
 =
 \sum_{j\in\mathfrak J_X}V_jL_j^2V_j^*.}
 \label{eq:V98-absolute-square-blocks}
\end{equation}
Distinct branches in \eqref{eq:V98-outer-partition} have orthogonal carrier
supports.  Moreover,
\begin{equation}
 \mathbf L_{\rm disc}=\mathbf Q_{\rm disc}=0.
 \label{eq:V98-disconnected-outer-zero}
\end{equation}
\end{theorem}

\begin{proof}
The ranges of the \(V_j\) are mutually orthogonal, so
\(\mathbf L_X\) is the direct sum of the self-adjoint blocks \(L_j\).
Functional calculus on a direct sum gives both identities in
\eqref{eq:V98-absolute-square-blocks}.  Disjoint index sets give orthogonal
supports.  On \(\mathfrak J_{\rm disc}\), the ternary cumulant is zero by
\eqref{eq:V97-disconnected-bank-zero}; hence every \(L_j\) and its absolute
value vanish.
\end{proof}

\begin{corollary}[Paid branches are removed before the physical norm]
\label{cor:V98-paid-removal}
The unresolved physical collision operator is the orthogonal direct sum
\[
 \mathbf L_{\rm res}
 =
 \mathbf L_{\rm side}\oplus\mathbf L_{\rm off}.
\]
No private-mean, cut-heavy, selected-heavy-pair, or disconnected block
appears in either residual absolute carrier.
\end{corollary}

\begin{proof}
Use \eqref{eq:V98-outer-partition},
\eqref{eq:V98-disconnected-outer-zero}, and the V94 priority convention.
\end{proof}

\section{The Racah POVM is not a physical joint pinching}
\label{sec:V98-instrument-firewall}

\begin{lemma}[Sharp joint measurability forces commutation]
\label{lem:V98-sharp-joint-commutation}
Let \(\{P_s\}_s\) and \(\{Q_v\}_v\) be two projection-valued resolutions of
the identity on a finite-dimensional Hilbert space.  Suppose positive
operators \(G_{sv}\) satisfy
\begin{equation}
 \sum_vG_{sv}=P_s,
 \qquad
 \sum_sG_{sv}=Q_v.
 \label{eq:V98-sharp-joint-marginals}
\end{equation}
Then
\begin{equation}
 P_sQ_v=Q_vP_s=G_{sv}
 \label{eq:V98-sharp-joint-product}
\end{equation}
for all \(s,v\).  In particular, the two sharp resolutions commute.
\end{lemma}

\begin{proof}
Positivity and the first marginal give \(0\le G_{sv}\le P_s\), hence
\[
 P_sG_{sv}=G_{sv}=G_{sv}P_s.
\]
Likewise \(Q_vG_{sv}=G_{sv}=G_{sv}Q_v\).  If \(s'\ne s\), then
\(P_sG_{s'v}=0\).  Therefore
\[
 P_sQ_v
 =
 P_s\sum_{s'}G_{s'v}=G_{sv}.
\]
The same calculation in the reverse order gives \(Q_vP_s=G_{sv}\).
\end{proof}

\begin{theorem}[Racah-instrument firewall]
\label{thm:V98-Racah-instrument-firewall}
On any multiplicity block where the sharp \(AB\)-intermediate projections
and sharp \(BC\)-intermediate projections do not commute, there is no
positive joint family having both sharp families as marginals.  Hence the
V76 law
\[
 \lambda_{u,v}(S,V,T)
\]
is a valid first-stage-channel/conditional-POVM law, but it is not the law
of a joint physical orthogonal pinching by \(S\) and \(V\).  In particular,
the canonical dilation \eqref{eq:V96-Naimark-compression} cannot be inserted
inside the physical \(K_T\) without specifying and justifying an additional
measurement instrument.
\end{theorem}

\begin{proof}
If a joint positive family with both sharp marginals existed, then
\Cref{lem:V98-sharp-joint-commutation} would force commutation.  Racah
recoupling is nontrivial precisely on blocks where the two intermediate
PVMs do not commute.  The V70 construction avoids contradiction because,
after fixing the \(AB\) channel, \(\Pi_S(V,T)\) is a conditional POVM on
the \(AB\)-multiplicity space; it is not a second sharp marginal on the
original full carrier.  A Naimark dilation realizes that POVM on an enlarged
space, but the dilation is not a physical operation already present in the
collision coefficient.
\end{proof}

\begin{corollary}[Exact status of the V94--V96 event effects]
\label{cor:V98-status-event-effects}
The V94 residual event effects, their V95 group kernels, and their V96
canonical dilation remain correct as finite-dimensional probability and
square-function diagnostics.  They may give sufficient bounds for the
corresponding partially traced positive square.  They are not, without an
additional instrument theorem, an exact decomposition of the physical
absolute carrier \(\mathbf Q_{\rm res}\).
\end{corollary}

\begin{proof}
Their POVM normalization and expectation identities were proved in
\Cref{thm:V76-two-level-law,lem:V95-lifted-POVM,
thm:V96-canonical-dilation}.  The non-identification with a physical joint
pinching is \Cref{thm:V98-Racah-instrument-firewall}.
\end{proof}

\section{Direct product-state control of the residual absolute carrier}
\label{sec:V98-direct-tax}

For a nonempty set \(I\) of retained original rows, use the V72 product
capacity
\[
 \kappa_I^\otimes(\mathbf Q_X)
 =
 \sup_{\rho_I^{\rm prod}}
 \operatorname{Tr}\left[
 \mathbf Q_X(\rho_I^{\rm prod}\otimes I_{I^c})\right].
\]
Let \(Y_j(\bigotimes_iB_i)\) be the physical compressed spectator
coefficient of V67.

\begin{theorem}[Restricted physical matrix tax]
\label{thm:V98-restricted-matrix-tax}
For every branch \(X\) and every retained-row set \(I\),
\begin{equation}
 \boxed{
 \sum_{j\in\mathfrak J_X}
 \left\|Y_j\!\left(\bigotimes_iB_i\right)\right\|_1
 \le
 \min\left\{
 \|\mathbf Q_X\|_{\rm op},
 \kappa_I^\otimes(\mathbf Q_X)
 \right\}
 \prod_i\|B_i\|_1.}
 \label{eq:V98-restricted-matrix-tax}
\end{equation}
The estimate has constant one, allows arbitrary within-row spectators, and
uses the full signed \(K_T\) before taking its blockwise absolute value.
\end{theorem}

\begin{proof}
Apply \Cref{thm:V72-complete-product-tax} to the orthogonal subfamily
\(\{V_j,L_j\}_{j\in\mathfrak J_X}\).  Its absolute carrier is exactly
\(\mathbf Q_X\) by \Cref{thm:V98-exact-absolute-carrier}.
\end{proof}

\begin{corollary}[Exact side/off physical capacity]
\label{cor:V98-side-off-physical-capacity}
Define
\begin{equation}
 \mathcal K_X^{\rm phys}
 :=
 \min\left\{
 \|\mathbf Q_X\|_{\rm op},
 \min_{\varnothing\ne I\subseteq\{o,a,b,c\}}
 \kappa_I^\otimes(\mathbf Q_X)
 \right\}.
 \label{eq:V98-physical-capacity}
\end{equation}
Then the unresolved physical collision contribution is at most
\begin{equation}
 \left(
 \mathcal K_{\rm side}^{\rm phys}
 +\mathcal K_{\rm off}^{\rm phys}
 \right)
 \prod_{i=o,a,b,c}\|B_i\|_1.
 \label{eq:V98-side-off-physical-bound}
\end{equation}
Each block is charged exactly once.
\end{corollary}

\begin{proof}
Apply \eqref{eq:V98-restricted-matrix-tax} separately to the two orthogonal
residual index sets and add the results.
\end{proof}

\section{Exact restricted square bridges}
\label{sec:V98-square-bridges}

For a fixed \(X\) and retained-row set \(I\), put
\begin{align}
 M_{I,X}^{(1)}
 &:=
 \operatorname{Tr}_{I^c}\mathbf Q_X,
 \label{eq:V98-restricted-first}\\
 M_{I,X}^{(2)}
 &:=
 \operatorname{Tr}_{I^c}\mathbf L_X^2,
 \label{eq:V98-restricted-second}\\
 \Gamma_{I,X}^{\otimes}
 &:=
 \|M_{I,X}^{(2)}\|_{I,\rm prod}.
 \label{eq:V98-restricted-Gamma}
\end{align}
Let \(D_{I^c}=\dim H_{I^c}\) for the carrier factors traced at this stage.

\begin{theorem}[Residual Kadison bridge]
\label{thm:V98-residual-Kadison}
For every branch \(X\),
\begin{align}
 (M_{I,X}^{(1)})^2
 &\le
 D_{I^c}M_{I,X}^{(2)},
 \label{eq:V98-residual-Kadison}\\
 \kappa_I^\otimes(\mathbf Q_X)
 &\le
 \sqrt{D_{I^c}\Gamma_{I,X}^{\otimes}}.
 \label{eq:V98-residual-square-bridge}
\end{align}
For \(I\) equal to all original rows, the sharper dimension-free statement
is
\begin{equation}
 \|\mathbf Q_X\|_{\rm rows,prod}
 \le
 \sqrt{\|\mathbf L_X^2\|_{\rm rows,prod}}.
 \label{eq:V98-full-row-square-bridge}
\end{equation}
\end{theorem}

\begin{proof}
The normalized partial trace
\[
 \Phi_I(Z):=D_{I^c}^{-1}\operatorname{Tr}_{I^c}Z
\]
is unital and completely positive.  Kadison's inequality applied to
\(\mathbf Q_X=|\mathbf L_X|\) gives
\[
 \Phi_I(\mathbf Q_X)^2
 \le
 \Phi_I(\mathbf Q_X^2)
 =
 \Phi_I(\mathbf L_X^2).
\]
Multiply by \(D_{I^c}^2\) to prove
\eqref{eq:V98-residual-Kadison}.  Scalar Cauchy--Schwarz in each product
state on the retained rows, followed by the product-state supremum, gives
\eqref{eq:V98-residual-square-bridge}.

If no carrier factor is traced, scalar Cauchy--Schwarz directly gives
\[
 \operatorname{Tr}(\rho|\mathbf L_X|)^2
 \le
 \operatorname{Tr}(\rho\mathbf L_X^2)
\]
for every full-row product state \(\rho\).  Take the supremum.
\end{proof}

\section{The canonical physical analysis image}
\label{sec:V98-physical-image}

For a retained-row set \(I\), identify
\(\mathcal H_{\rm rows}=H_I\otimes H_{I^c}\).  For
\(\psi\in H_I\), let \(W_\psi:H_{I^c}\to\mathcal H_{\rm rows}\) be
\(W_\psi\xi=\psi\otimes\xi\), and define
\begin{equation}
 \mathcal A_{I,X}\psi:=\mathbf L_XW_\psi
 \quad\hbox{in }\mathcal S_2(H_{I^c},\mathcal H_{\rm rows}).
 \label{eq:V98-physical-analysis-map}
\end{equation}

\begin{theorem}[Exact physical injective-image gate]
\label{thm:V98-physical-image-gate}
The physical analysis map satisfies
\begin{align}
 \mathcal A_{I,X}^*\mathcal A_{I,X}
 &=
 M_{I,X}^{(2)},
 \label{eq:V98-physical-analysis-Gram}\\
 \boxed{
 \Gamma_{I,X}^{\otimes}
 =
 \sup_{\|z\|_{\rm HS}=1}
 \|\mathcal A_{I,X}^*z\|_{\varepsilon,I}^2.}
 \label{eq:V98-physical-injective-image}
\end{align}
Here \(\|\cdot\|_{\varepsilon,I}\) is the injective vector norm across the
retained original rows.  No Racah measurement instrument occurs in this
formula.
\end{theorem}

\begin{proof}
For \(\phi,\psi\in H_I\),
\[
 \langle\mathcal A_{I,X}\phi,\mathcal A_{I,X}\psi\rangle_{\rm HS}
 =
 \operatorname{Tr}\left[
 W_\phi^*\mathbf L_X^2W_\psi\right]
 =
 \langle\phi,M_{I,X}^{(2)}\psi\rangle.
\]
This proves \eqref{eq:V98-physical-analysis-Gram}.  Apply
\Cref{thm:V96-exact-image-norm} with the product structure of \(H_I\).
\end{proof}

\begin{corollary}[A rigorous residual collision gate]
\label{cor:V98-rigorous-residual-gate}
The simultaneous-sideways physical collision contribution closes at the
V60 marked-tree scale if, for at least one retained-row choice \(I\),
\begin{equation}
 \min\left\{
 \|\mathbf Q_{\rm side}\|_{\rm op},
 \sqrt{D_{I^c}}
 \sup_{\|z\|_{\rm HS}=1}
 \|\mathcal A_{I,\rm side}^*z\|_{\varepsilon,I}
 \right\}
 \le
 Cs_\alpha\sum_{r=b,c}
 \frac{h_{oa}\widehat H_{ar}h_{bc}}
      {\Xi_a\Xi_r}.
 \label{eq:V98-side-image-gate}
\end{equation}
The off-cut contribution obeys the identical separate criterion with
\(\mathrm{side}\) replaced by \(\mathrm{off}\).
\end{corollary}

\begin{proof}
The squared injective-image supremum is
\(\Gamma_{I,X}^{\otimes}\) by
\eqref{eq:V98-physical-injective-image}.  Insert
\eqref{eq:V98-residual-square-bridge} into
\eqref{eq:V98-restricted-matrix-tax}, then use the output-mass allocation
and final-resolvent argument of
\Cref{thm:V67-matrix-capacity-sufficient}.  The two branches remain
separate by \eqref{eq:V98-side-off-physical-bound}.
\end{proof}

\begin{remark}[When the V72 pair--core factor reappears]
\label{rem:V98-rectangular-factorization}
If a residual outer set is rectangular in the two external pair-output
labels and the core labels, then the exact scalar partial traces
\eqref{eq:V72-pair-scalar-traces} factor out and
\eqref{eq:V98-side-image-gate} reduces to the corresponding
\[
 P_{ab}\sqrt{d_C\Gamma_{ab,X}^{\otimes}}
 \quad\hbox{or}\quad
 P_{ac}\sqrt{d_B\Gamma_{ac,X}^{\otimes}}
\]
entry of the V72 gate.  If the residual selector correlates those labels,
that factorization must not be asserted; the complete
\(\kappa_I^\otimes(\mathbf Q_X)\) in
\eqref{eq:V98-physical-capacity} remains the exact object.
\end{remark}

\section{Regression tests and direction}
\label{sec:V98-tests}

\begin{proposition}[Disposition of the four required regressions]
\label{prop:V98-regressions}
The physical carrier construction has the following exact behavior.
\begin{enumerate}[label=\textup{(\roman*)}]
\item The fundamental--trivial--dual certificate is absent because
      \(\mathbf Q_{\rm disc}=0\).
\item A dual-singlet retained range keeps the \(1/d_U\) product debit of
      \eqref{eq:V96-singlet-range}.
\item A Cartan highest-weight range remains product-visible; all possible
      smallness stays in the exact singular values of
      \(\widetilde K_T\), as required by
      \eqref{eq:V79-signed-output-formula}.
\item The V60 alternating array is assigned to one complete outer
      orientation before the product tax.  No coordinatewise minimum of
      incompatible Racah POVMs is used.
\end{enumerate}
\end{proposition}

\begin{proof}
Part \textup{(i)} is
\eqref{eq:V98-disconnected-outer-zero}.  Part \textup{(ii)} follows by
evaluating the product capacity on the singlet support as in
\Cref{prop:V96-singlet-image-test}.  Part \textup{(iii)} follows because
\(\mathbf L_X\), not its covariance-square majorant, occurs in
\eqref{eq:V98-physical-analysis-map}.  Part \textup{(iv)} is built into
the outer partition \eqref{eq:V98-outer-partition} and the firewall
\Cref{thm:V98-Racah-instrument-firewall}.
\end{proof}

\begin{assessment}[Position after reconstructing the physical object]
\label{ass:V98-position}
The remaining simultaneous-sideways question is no longer an abstract
overlap of two conditional POVMs.  It is the product injective image
\eqref{eq:V98-physical-injective-image} of the signed, branch-restricted,
outer-block collision operator.  The V94--V96 level effects remain useful
upper diagnostics after squaring, but any successful proof must return to
\(\mathbf L_{\rm side}\) before asserting a physical filter or an
orientation switch.
\end{assessment}

\section{V98 ledger and V99 acquisition}
\label{sec:V98-ledger}

\begin{theorem}[Integrated V98 statement]
\label{thm:V98-integrated}
Preserving V1--V97, V98 proves:
\begin{enumerate}[label=\textup{(V98.\arabic*)},leftmargin=5.7em]
\item the exact physical outer-block partition
      \eqref{eq:V98-outer-partition};
\item the residual signed and absolute carriers
      \eqref{eq:V98-physical-signed-operator}--%
      \eqref{eq:V98-physical-absolute-carrier};
\item exact disconnected deletion
      \eqref{eq:V98-disconnected-outer-zero};
\item the sharp-joint-measurability obstruction
      \eqref{eq:V98-sharp-joint-product};
\item the Racah-instrument firewall;
\item the restricted physical matrix tax
      \eqref{eq:V98-restricted-matrix-tax};
\item the residual Kadison bridges
      \eqref{eq:V98-residual-square-bridge}--%
      \eqref{eq:V98-full-row-square-bridge};
\item the exact physical injective-image formula
      \eqref{eq:V98-physical-injective-image}; and
\item the rigorous side/off residual gates
      \eqref{eq:V98-side-image-gate}.
\end{enumerate}
\end{theorem}

\begin{proof}
These are
\Cref{thm:V98-exact-absolute-carrier,
cor:V98-paid-removal,
lem:V98-sharp-joint-commutation,
thm:V98-Racah-instrument-firewall,
thm:V98-restricted-matrix-tax,
thm:V98-residual-Kadison,
thm:V98-physical-image-gate,
cor:V98-rigorous-residual-gate}.
\end{proof}

\begin{openproblem}[V99 mandate: estimate the signed side image]
\label{op:V99-mandate}
Work with \(\mathbf L_{\rm side}\), not with an inserted Racah instrument.
\begin{enumerate}
\item Write \(\mathcal A_{I,\rm side}^*\) explicitly in the complete
      outer labels and the V79 final-output path basis.
\item Delete all trivial-input blocks before summing and retain the three
      signed covariance cuts inside \(K_T\).
\item Determine whether the external pair/same-label weights form a
      rectangular set.  Factor them only if
      \Cref{rem:V98-rectangular-factorization} applies.
\item On nonrectangular residual sets, estimate the complete product
      capacity \(\kappa_I^\otimes(\mathbf Q_{\rm side})\) without replacing
      it by independent label marginals.
\item Test the exact image first on the fundamental--dual singlet and
      Cartan path vectors, then on the V89 balanced thermal ancestry bank.
\item If the balanced bank remains order one, identify whether it lies in
      the already-paid selected-heavy-pair branch or in a genuine
      simultaneous-sideways outer block.
\item Keep \(\mathbf L_{\rm off}\) separate.  V99 need not close it by
      reclassifying it as a side event.
\end{enumerate}
The next compulsory companion audit is V100.
\end{openproblem}

\section{Firewall after V98}
\label{sec:V98-firewall}

\begin{firewall}[Current claim boundary]
\label{fire:V98-final}
V98 reconstructs the exact physical residual absolute carrier, proves that
noncommuting Racah outcomes cannot be treated as a joint sharp pinching, and
reduces each residual branch to a signed physical injective-image estimate.
It does not yet bound the simultaneous-sideways image at the required path
scale, close the off-cut alternative, aggregate all punctures or multiple
collisions, prove the complete \([3,2]\), \([3,3]\), or \([4]\) sectors,
construct the continuum Yang--Mills theory, or prove a positive mass gap.
\end{firewall}

\section{Append-only record for V98}
\label{sec:V98-record}

V98 was appended to the complete V97 manuscript.  The exact V97 parent had
\(66211\) physical lines, \(2320055\) bytes, and SHA--256 digest
\[
 \texttt{84499369b26baeb36672ef2c0f6f299ae9c238d10bb241520857d71a8357bb8d}.
\]
No V97 mathematical content was changed.  The sole final
\verb|\end{document}| is restored below.

% =============================================================================
% END APPENDED VERSION V98 MATERIAL
% =============================================================================
% =============================================================================
% BEGIN APPENDED VERSION V99 MATERIAL
% =============================================================================

\chapter{V99: Cartan Exhaustion of Unit Product Visibility and the
Exact Cartan--Non-Cartan Image Split}
\label{ch:V99}

\section{Purpose}
\label{sec:V99-purpose}

V98 reconstructed the physical residual absolute carrier and reduced its
signed square to the injective image of the physical outer-block operator.
The first obstruction is now geometric: which first-stage
\(AB\)-isotypic spaces can contain an admissible product vector and hence
have product visibility one?

This chapter gives a complete answer.  For a connected complex reductive
group, a nonzero decomposable tensor lying entirely in one isotypic
component of \(V_\lambda\otimes V_\mu\) forces that component to be the
Cartan component \(V_{\lambda+\mu}\).  The proof below is self-contained
apart from the standard Borel fixed-point theorem.  It also handles
isotypic multiplicity, which is essential for the V91 carrier formula.

Consequently every exactly unit-visible first-stage Yang--Mills channel is
Cartan.  The Cartan part of the V98 signed square reduces to one explicit
Plancherel scalar, while the non-Cartan part receives a strict fixed-label
product debit.  An exact symmetric-power example then shows that this
strict debit can decay exponentially with the labels.  Thus the
classification closes a logical gap but does not manufacture the uniform
asymptotic estimate still required.

\section{Decomposable tensors in one isotypic component}
\label{sec:V99-decomposable}

Let \(G_{\mathbb C}\) be a connected complex reductive group, choose a
Borel subgroup \(B\), and let \(V_\lambda,V_\mu\) be irreducible
highest-weight modules.  Write
\[
 V_\lambda\otimes V_\mu
 =
 \bigoplus_\nu V_\nu\otimes\mathcal M_\nu.
\]
The component \(V_{\lambda+\mu}\) has multiplicity one and is called the
Cartan component.

\begin{lemma}[A decomposable highest vector is the Cartan highest vector]
\label{lem:V99-decomposable-highest}
Suppose \(0\ne v\otimes w\in V_\lambda\otimes V_\mu\) is a highest-weight
vector for the diagonal action.  Then \(v\) and \(w\) are highest-weight
vectors in their respective modules, and
\begin{equation}
 \operatorname{wt}(v\otimes w)=\lambda+\mu.
 \label{eq:V99-highest-weight-sum}
\end{equation}
\end{lemma}

\begin{proof}
First let \(H\) lie in a Cartan subalgebra and suppose
\[
 (Hv)\otimes w+v\otimes(Hw)
 =
 \nu(H)v\otimes w.
\]
Pass the first tensor factor to \(V_\lambda/\mathbb Cv\).  Since \(w\ne0\),
the resulting equality implies \(Hv\in\math Cv\).  Interchanging the
factors gives \(Hw\in\math Cw\).  Thus \(v,w\) are weight vectors, say of
weights \(\lambda'\) and \(\mu'\), with \(\nu=\lambda'+\mu'\).

For every positive-root operator \(E\),
\[
 (Ev)\otimes w+v\otimes(Ew)=0.
\]
The same quotient argument gives \(Ev\in\math Cv\) and
\(Ew\in\math Cw\).  A root operator is nilpotent in a finite-dimensional
highest-weight module, so both proportionality constants are zero.
Therefore \(v,w\) are killed by every positive-root operator.  Irreducibility
makes their lines the unique highest-weight lines of
\(V_\lambda,V_\mu\).  Hence \(\lambda'=\lambda\),
\(\mu'=\mu\), and \eqref{eq:V99-highest-weight-sum} follows.
\end{proof}

\begin{theorem}[Cartan exhaustion of decomposable isotypic vectors]
\label{thm:V99-Cartan-exhaustion}
Let \(0\ne v\otimes w\in V_\lambda\otimes V_\mu\) be decomposable.
If it lies entirely in one isotypic component
\(V_\nu\otimes\mathcal M_\nu\), then
\begin{equation}
 \boxed{\nu=\lambda+\mu.}
 \label{eq:V99-Cartan-exhaustion}
\end{equation}
Thus no non-Cartan isotypic component, even one with multiplicity, contains
a nonzero decomposable tensor.
\end{theorem}

\begin{proof}
Let \(p=[v\otimes w]\) in projective space and let
\[
 X:=\overline{G_{\mathbb C}\cdot p}.
\]
The Segre variety of decomposable projective tensors is Zariski closed and
\(G_{\mathbb C}\)-invariant, so \(X\) lies in that variety.  Since the
\(\nu\)-isotypic space is invariant and closed, \(X\) also lies in
\(\mathbb P(V_\nu\otimes\mathcal M_\nu)\).

The Borel fixed-point theorem gives a \(B\)-fixed point
\([z]\in X\).  The line \(\math Cz\) is fixed by the maximal torus and by
the unipotent radical of \(B\); hence \(z\) is a highest-weight vector of
weight \(\nu\) in the \(\nu\)-isotypic component.  Because \(X\) lies in
the Segre variety, \(z=v_0\otimes w_0\) is decomposable.  Apply
\Cref{lem:V99-decomposable-highest} to obtain
\(\nu=\lambda+\mu\).
\end{proof}

\begin{remark}[Provenance and scope]
\label{rem:V99-Baur-provenance}
The irreducible-component form of
\Cref{thm:V99-Cartan-exhaustion} is the decomposable-tensor theorem in
K.~Baur, \emph{Cartan components and decomposable tensors},
Transformation Groups \textbf{8} (2003), 309--319,
\href{https://doi.org/10.1007/s00031-003-1203-2}
{doi:10.1007/s00031-003-1203-2}.  The orbit-closure proof above records
explicitly why the same conclusion holds for a full isotypic component with
multiplicity.
\end{remark}

\section{Exact classification of unit product visibility}
\label{sec:V99-unit-visibility}

Let \(P_\nu^{\lambda,\mu}\) be the full \(\nu\)-isotypic projection and
retain
\[
 \beta_{\lambda,\mu}(\nu)
 :=
 \sup_{\|v\|=\|w\|=1}
 \|P_\nu^{\lambda,\mu}(v\otimes w)\|^2.
\]

\begin{corollary}[Unit visibility if and only if Cartan]
\label{cor:V99-unit-visibility-Cartan}
For every isotypic component,
\begin{equation}
 \boxed{
 \beta_{\lambda,\mu}(\nu)=1
 \quad\Longleftrightarrow\quad
 \nu=\lambda+\mu.}
 \label{eq:V99-unit-visibility-Cartan}
\end{equation}
Every non-Cartan component satisfies
\(\beta_{\lambda,\mu}(\nu)<1\).
\end{corollary}

\begin{proof}
If \(\nu=\lambda+\mu\), the tensor product of normalized highest-weight
vectors lies in the Cartan component, so the visibility is one.
Conversely, the product-vector domain is compact and the projection norm is
continuous.  If its supremum is one, it is attained by a product vector
lying entirely in the \(\nu\)-isotypic space.  Apply
\Cref{thm:V99-Cartan-exhaustion}.  The final assertion follows from the
same compactness.
\end{proof}

Let \(P_{\rm Car}=P_{\lambda+\mu}^{\lambda,\mu}\).

\begin{theorem}[Positive fixed-label Cartan floor]
\label{thm:V99-positive-Cartan-floor}
For fixed \(\lambda,\mu\),
\begin{equation}
 \alpha_{\lambda,\mu}
 :=
 \inf_{\|v\|=\|w\|=1}
 \|P_{\rm Car}(v\otimes w)\|^2
 >0.
 \label{eq:V99-Cartan-floor}
\end{equation}
Equivalently,
\begin{equation}
 \|I-P_{\rm Car}\|_{\rm prod}
 =
 1-\alpha_{\lambda,\mu}<1.
 \label{eq:V99-nonCartan-complement-visibility}
\end{equation}
\end{theorem}

\begin{proof}
Suppose first that a nonzero decomposable tensor \(v\otimes w\) were
orthogonal to the Cartan component.  Its projective orbit closure would lie
in the invariant sum of the non-Cartan isotypic components and in the
Segre variety.  Borel fixed point would again produce a decomposable
highest-weight vector in that non-Cartan sum.  A \(B\)-fixed projective line
is a torus weight line and therefore lies in a single highest-weight
isotypic space.  Lemma~\ref{lem:V99-decomposable-highest} would give weight
\(\lambda+\mu\), contradicting exclusion of the Cartan component.
Therefore the displayed Cartan projection is nonzero on every nonzero
product tensor.

The squared projection norm is a positive continuous function on the
compact product of the two unit spheres, so its infimum is positive.  For a
unit product tensor \(p\),
\[
 \langle p,(I-P_{\rm Car})p\rangle
 =
 1-\|P_{\rm Car}p\|^2.
\]
Taking the supremum proves
\eqref{eq:V99-nonCartan-complement-visibility}.
\end{proof}

\section{The Cartan floor is not uniform in growing labels}
\label{sec:V99-floor-sharpness}

Let \(R_a=\operatorname{Sym}^a(\mathbb C^3)\) and
\(R_b=\operatorname{Sym}^b(\mathbb C^3)\).  Their Cartan component is
\(R_{a+b}\).

\begin{proposition}[Exact orthogonal-monomial Cartan probability]
\label{prop:V99-orthogonal-monomial}
For orthonormal \(e_1,e_2\in\mathbb C^3\),
\begin{equation}
 \left\|
 P_{R_{a+b}}
 \left(e_1^{\odot a}\otimes e_2^{\odot b}\right)
 \right\|^2
 =
 \binom{a+b}{a}^{-1}.
 \label{eq:V99-binomial-Cartan-probability}
\end{equation}
Consequently
\begin{equation}
 0<\alpha_{R_a,R_b}
 \le
 \binom{a+b}{a}^{-1}.
 \label{eq:V99-Cartan-floor-upper}
\end{equation}
For balanced \(a=b=m\), this upper bound is exponentially small in \(m\).
\end{proposition}

\begin{proof}
Embed
\[
 R_a\otimes R_b
 \subset
 (\mathbb C^3)^{\otimes(a+b)}.
\]
The Cartan projection is full symmetrization over the \(a+b\) tensor
positions.  There are
\[
 N=\binom{a+b}{a}
\]
orthonormal words containing \(a\) copies of \(e_1\) and \(b\) copies of
\(e_2\).  Symmetrization of the displayed word is the average of those
\(N\) words, so its squared norm is
\(N\cdot N^{-2}=N^{-1}\).  The floor is at most this value.  For
\(a=b=m\), the central binomial coefficient grows exponentially.
\end{proof}

\begin{firewall}[Strict is not uniform]
\label{fire:V99-strict-not-uniform}
The inequality \(\beta_{\lambda,\mu}(\nu)<1\) for every fixed non-Cartan
component cannot be substituted for a cutoff-uniform debit.
Equation \eqref{eq:V99-Cartan-floor-upper} shows that a product tensor may
place all but an exponentially small fraction of its mass in the
non-Cartan complement as the input labels grow.  A valid continuum estimate
must correlate the actual channel probabilities with the connected symbol
or fusion kernel.
\end{firewall}

\section{Exact Cartan column of the physical signed square}
\label{sec:V99-Cartan-column}

Return to one V98 branch \(X\) and one retained orientation \(A|B\).
Let \(S_{\rm Car}\subset A\otimes B\) be the Cartan component.  It occurs
with first-stage multiplicity one.  In each final \(T\)-multiplicity space,
let
\[
 e_{S_{\rm Car},\beta}^{j,T},
 \qquad
 1\le\beta\le N_{S_{\rm Car},C}^{T},
\]
be the \(AB\)-first path vectors in the physical outer block \(j\).
The block weight \(c_j\) and final weight are included in \(L_j\).

\begin{theorem}[General Cartan-column Plancherel scalar]
\label{thm:V99-general-Cartan-column}
The restriction of the V98 square
\[
 M_{AB,X}^{(2)}
 =
 \operatorname{Tr}_C\mathbf L_X^2
\]
to the Cartan first-stage component is
\begin{equation}
 P_{\rm Car}M_{AB,X}^{(2)}P_{\rm Car}
 =
 \mathfrak g_{\rm Car}^{X}P_{\rm Car},
 \label{eq:V99-Cartan-column-scalar}
\end{equation}
where
\begin{equation}
 \boxed{
 \mathfrak g_{\rm Car}^{X}
 =
 \frac1{d_{S_{\rm Car}}}
 \sum_{\substack{j\in\mathfrak J_X\\T}}
 d_T
 \sum_{\beta=1}^{N_{S_{\rm Car},C}^{T}}
 \left\|
 L_{j,T}e_{S_{\rm Car},\beta}^{j,T}
 \right\|^2.}
 \label{eq:V99-Cartan-Plancherel-scalar}
\end{equation}
Moreover the product norm of this isolated column is exactly
\begin{equation}
 \boxed{
 \|P_{\rm Car}M_{AB,X}^{(2)}P_{\rm Car}\|_{\rm prod}
 =
 \mathfrak g_{\rm Car}^{X}.}
 \label{eq:V99-Cartan-column-product}
\end{equation}
\end{theorem}

\begin{proof}
Apply the matrix partial-trace formula
\eqref{eq:V67-matrix-partial-trace} to the positive block
\(\mathbf L_X^2\).  Since the Cartan first-stage multiplicity is one, every
principal block on that first-stage multiplicity space is scalar.  Its
diagonal entry on the path \(e_{S_{\rm Car},\beta}^{j,T}\) is
\(\|L_{j,T}e_{S_{\rm Car},\beta}^{j,T}\|^2\).  The partial-trace
coefficient is \(d_T/d_{S_{\rm Car}}\), which proves
\eqref{eq:V99-Cartan-column-scalar}--%
\eqref{eq:V99-Cartan-Plancherel-scalar}.

The tensor product of highest-weight vectors lies entirely in
\(S_{\rm Car}\), so it attains the scalar eigenvalue in
\eqref{eq:V99-Cartan-column-scalar}.  Positivity gives the matching upper
bound, proving \eqref{eq:V99-Cartan-column-product}.
\end{proof}

\begin{corollary}[All unit-visible columns are scalar Plancherel gates]
\label{cor:V99-all-unit-columns}
If an \(AB\)-isotypic column has product visibility one, it is the Cartan
column and its exact signed-square product capacity is
\eqref{eq:V99-Cartan-Plancherel-scalar}.  There is no other
unit-visible Racah ancestry requiring a matrix estimate.
\end{corollary}

\begin{proof}
Use \Cref{cor:V99-unit-visibility-Cartan,
thm:V99-general-Cartan-column}.
\end{proof}

\section{A quantitative fixed-label Cartan--non-Cartan split}
\label{sec:V99-capacity-split}

Let
\[
 M=M_{AB,X}^{(2)},\qquad
 M_{\rm nc}:=(I-P_{\rm Car})M(I-P_{\rm Car}),
 \qquad
 \Lambda_{\rm nc}:=\|M_{\rm nc}\|_{\rm op}.
\]
Invariance makes the Cartan and non-Cartan pieces reducing subspaces of
\(M\).

\begin{theorem}[Exact fixed-label product-capacity envelope]
\label{thm:V99-Cartan-nonCartan-envelope}
With
\(\alpha=\alpha_{\lambda,\mu}\) and
\(\mathfrak g=\mathfrak g_{\rm Car}^{X}\),
\begin{equation}
 \boxed{
 \|M\|_{\rm prod}
 \le
 \begin{cases}
 \mathfrak g,&\mathfrak g\ge\Lambda_{\rm nc},\\[0.3em]
 \alpha\mathfrak g+(1-\alpha)\Lambda_{\rm nc},
 &\mathfrak g<\Lambda_{\rm nc}.
 \end{cases}}
 \label{eq:V99-Cartan-nonCartan-envelope}
\end{equation}
In particular,
\begin{equation}
 \|M_{\rm nc}\|_{\rm prod}
 \le
 (1-\alpha)\Lambda_{\rm nc}.
 \label{eq:V99-nonCartan-strict-debit}
\end{equation}
\end{theorem}

\begin{proof}
For a unit product vector \(p\), put
\[
 q=\|P_{\rm Car}p\|^2.
\]
Theorem~\ref{thm:V99-positive-Cartan-floor} gives
\(\alpha\le q\le1\).  Block reduction and
\eqref{eq:V99-Cartan-column-scalar} give
\[
 \langle p,Mp\rangle
 \le
 q\mathfrak g+(1-q)\Lambda_{\rm nc}.
\]
This affine function of \(q\in[\alpha,1]\) is maximized at \(q=1\) when
\(\mathfrak g\ge\Lambda_{\rm nc}\), and at \(q=\alpha\) otherwise.
This proves \eqref{eq:V99-Cartan-nonCartan-envelope}.  Set
\(\mathfrak g=0\) for the isolated non-Cartan operator to obtain
\eqref{eq:V99-nonCartan-strict-debit}.
\end{proof}

\begin{assessment}[What the exact split gains]
\label{ass:V99-split-gain}
The V91 equality criterion is now exhaustive.  An operator-norm maximum
attained on one isotypic product vector can occur only in the Cartan
column, where the full signed square is the explicit scalar
\eqref{eq:V99-Cartan-Plancherel-scalar}.  Every non-Cartan column has a
strict fixed-label debit, and the whole non-Cartan complement has the
strict debit \eqref{eq:V99-nonCartan-strict-debit}.  The remaining
continuum problem is quantitative: \(\alpha_{\lambda,\mu}\) can be
exponentially small, so one must retain the distribution of the
non-Cartan singular values rather than use only their maximum.
\end{assessment}

\section{Regression tests}
\label{sec:V99-tests}

\begin{proposition}[Three exact tests]
\label{prop:V99-three-tests}
\begin{enumerate}[label=\textup{(\roman*)}]
\item For \(F\otimes F^*=\mathbf1\oplus\mathrm{Ad}\), the Cartan component
      is \(\mathrm{Ad}\) and
      \begin{equation}
       \alpha_{F,F^*}=\frac23.
       \label{eq:V99-fundamental-Cartan-floor}
      \end{equation}
\item For \(R_a\otimes R_b\), the highest-weight product has Cartan
      probability one, while the orthogonal-monomial product has the
      probability \eqref{eq:V99-binomial-Cartan-probability}.
\item In the V89 balanced family \(R_m\otimes R_m\), the possible
      exponential smallness of the global Cartan floor does not reopen the
      actual Cartan thermal column, which is already paid by
      \eqref{eq:V91-Cartan-pair-core-uniform}.
\end{enumerate}
\end{proposition}

\begin{proof}
For \textup{(i)}, the singlet probability of a product vector is at most
\(1/3\), with equality on a matched basis vector, by
\eqref{eq:V92-singlet-product-norm}.  Its orthogonal complement is the
Cartan adjoint, giving a minimum probability \(2/3\).
Part \textup{(ii)} combines the highest-weight construction with
\Cref{prop:V99-orthogonal-monomial}.  Part \textup{(iii)} is the V91
un-factorized carrier theorem; no global Cartan-floor estimate is used
there.
\end{proof}

\begin{assessment}[Alternating-mass disposition]
\label{ass:V99-alternating}
The Cartan classification is orientation-local.  On the V60 alternating
array one must first choose the complete \(ab\) or \(ac\) physical carrier
as in V98, then apply the corresponding Cartan split.  Taking the better
Cartan floor coordinate by coordinate would repeat the forbidden
orientation switch and is not permitted.
\end{assessment}

\section{V99 ledger and compulsory V100 acquisition}
\label{sec:V99-ledger}

\begin{theorem}[Integrated V99 statement]
\label{thm:V99-integrated}
Preserving V1--V98, V99 proves:
\begin{enumerate}[label=\textup{(V99.\arabic*)},leftmargin=5.7em]
\item Cartan exhaustion of decomposable isotypic vectors
      \eqref{eq:V99-Cartan-exhaustion};
\item the exact classification
      \eqref{eq:V99-unit-visibility-Cartan} of unit product visibility;
\item the positive fixed-label Cartan floor
      \eqref{eq:V99-Cartan-floor};
\item the exact binomial regression
      \eqref{eq:V99-binomial-Cartan-probability};
\item the general physical Cartan-column scalar
      \eqref{eq:V99-Cartan-Plancherel-scalar}; and
\item the quantitative Cartan--non-Cartan product envelope
      \eqref{eq:V99-Cartan-nonCartan-envelope}.
\end{enumerate}
\end{theorem}

\begin{proof}
These are
\Cref{thm:V99-Cartan-exhaustion,
cor:V99-unit-visibility-Cartan,
thm:V99-positive-Cartan-floor,
prop:V99-orthogonal-monomial,
thm:V99-general-Cartan-column,
thm:V99-Cartan-nonCartan-envelope}.
\end{proof}

\begin{openproblem}[V100 mandate: audit and quantitative non-Cartan image]
\label{op:V100-mandate}
First perform the compulsory read-only audit of the current SM and NS
notebooks.  Then:
\begin{enumerate}
\item keep the V98 physical side carrier and its exact outer orientation;
\item split its first-stage square into the Cartan scalar
      \eqref{eq:V99-Cartan-Plancherel-scalar} and the complete non-Cartan
      operator;
\item use the actual distribution of
      \(\mathscr S_{S,X}\), not the exponentially weak global floor
      \(\alpha_{\lambda,\mu}\);
\item express the non-Cartan product capacity through the common Casimir or
      fusion-kernel law before taking an operator maximum;
\item determine whether a high-visibility non-Cartan sequence must approach
      the Cartan variety and, if so, whether the centred cumulant or an
      internal-pair level supplies the compensating small singular value;
\item test every proposed estimate on the orthogonal-monomial family
      \eqref{eq:V99-binomial-Cartan-probability}, the dual singlet, and the
      V89 balanced thermal family;
\item keep the V98 off-cut carrier separate.
\end{enumerate}
The next compulsory companion audit after V100 is V105.
\end{openproblem}

\section{Firewall after V99}
\label{sec:V99-firewall}

\begin{firewall}[Current claim boundary]
\label{fire:V99-final}
V99 proves that the Cartan component exhausts all unit-visible
first-stage channels and gives an exact signed-square scalar for that
column.  It also proves a strict non-Cartan debit for every fixed input
pair, but the debit is not uniform in growing labels.  It does not yet
control the complete non-Cartan side image, close the general Cartan scalar
outside the V91 atlas, close the off-cut alternative, prove all collision
and partition sectors, construct the continuum Yang--Mills theory, or prove
a positive mass gap.
\end{firewall}

\section{Append-only record for V99}
\label{sec:V99-record}

V99 was appended to the complete V98 manuscript.  The exact V98 parent had
\(66811\) physical lines, \(2341184\) bytes, and SHA--256 digest
\[
 \texttt{c240d2f5ce21c7882f28228d93922097c6f9eda1e35ab841f92477911395f6b3}.
\]
No V98 mathematical content was changed.  The sole final
\verb|\end{document}| is restored below.

% =============================================================================
% END APPENDED VERSION V99 MATERIAL
% =============================================================================
% =============================================================================
% BEGIN APPENDED VERSION V100 MATERIAL
% =============================================================================

\chapter{V100: The Exact Casimir--Singular Law, a Near-Unit
Non-Cartan Sequence, and the Upstream Symbol Gate}
\label{ch:V100}

\section{Purpose and compulsory companion audit}
\label{sec:V100-purpose-audit}

V99 proved that an isotypic component has product visibility one if and
only if it is the Cartan component.  It also proved a strictly positive
Cartan floor for every fixed pair of input labels.  The remaining question
was whether that strict fixed-label statement has a useful uniform
quantitative form.  This chapter gives a negative answer in the simplest
nontrivial multiplicity-free family:
\[
 F\otimes\operatorname{Sym}^{N}\mathbb C^3.
\]
The sole non-Cartan component has product visibility \(N/(N+1)\), and hence
tends to one.  Moreover, the product vectors which realize this limit do
not approach the decomposable Cartan locus.  They become asymptotically
orthogonal to it in projective fidelity.

Consequently a representation-uniform residual estimate cannot come from
the designation non-Cartan alone.  It must use the singular-value
distribution of the actual V98 physical signed cumulant.  The first part
of this chapter therefore writes that distribution exactly.  Before any
operator maximum is taken, it is an exact matrix-valued average against
the V91 fusion kernel.  The common Casimir then records one exact moment
of the same product-state law.  Its least concave majorant gives the sharp
relaxation which can be obtained from that one moment alone.

This is the fifth version after V95, so the current companion notebooks
were inspected before choosing this direction.  The read-only snapshot,
refreshed after the SM notebook advanced during the audit, was
\begin{center}
\begin{tabular}{@{}llll@{}}
programme & endpoint & physical lines and bytes & SHA--256\\
\hline
SM--Einstein &
\texttt{V103} &
\(74538,\ 2734252\) &
\texttt{fb9342b982b80c51b216f312191838baf3959938b3edcb927c3c044c8a9482a4}\\
Navier--Stokes &
\texttt{V31} &
\(23052,\ 887537\) &
\texttt{f1121b62238ad5491bb7cf03635ea9934942edf573ed8faaa9ee79e1d38a6696}
\end{tabular}
\end{center}

SM V103 assembles its conservative principal boundary domain and embeds
the V35 Yang--Mills/ghost glancing sequence into the full direct sum.
Its exact lesson is that positive floors on decoupled sectors do not lift
a singular vector supported in the gauge sector.  NS V31 proves an exact
dyadic first-moment/coarea identity and a heat-formation ladder, but its
positive high-parent estimate returns the continuation-class stock.
It therefore retains the complete signed triad work before separating
positive annular pieces.  In the present problem these facts determine a
test order, not an imported estimate: retain the V98 signed physical block,
form its canonical square, preserve the common product-state distribution,
and only then relax it by a scalar profile.

\begin{firewall}[Scope of the V100 companion audit]
\label{fire:V100-audit-scope}
No SM boundary estimate and no NS heat estimate is used to bound a
Yang--Mills collision coefficient.  The matrix-valued fusion law, the
Casimir moment, and the counterexample below are proved directly in the
finite-dimensional representation problem.  The companion audit only
prevents two invalid moves: repairing a gauge-supported defect with a
decoupled positive sector, and replacing a signed common law by separately
maximized positive pieces before its kernel and ancestry have been tested.
\end{firewall}

\section{The exact matrix-valued fusion law of the physical square}
\label{sec:V100-exact-fusion-law}

Fix one V98 physical branch
\[
 X\in\{\mathrm{side},\mathrm{off}\}
\]
and one complete physical orientation \(A|B\), with the remaining core
factor denoted by \(C\).  All formulas in this chapter are branchwise.
In particular, no side/off sum is inserted inside a product-state
supremum.

Write
\[
 H_A\otimes H_B
 =
 \bigoplus_S H_S\otimes\mathcal M_{AB}^{S},
 \qquad
 \dim\mathcal M_{AB}^{S}=N_{AB}^{S}.
\]
For every physical outer block \(j\) with final irrep \(T\), put
\[
 W_{j,T}:=L_{j,T}^{\,2}\ge0
\]
on its complete \(T\)-multiplicity space.  In the \(AB\)-first path basis,
let
\[
 W_{j,T}[S,\beta]\ge0,
 \qquad 1\le\beta\le N_{SC}^{T},
\]
be the V67 principal matrix acting on \(\mathcal M_{AB}^{S}\).  A path
absent from \(j\) is represented by the zero matrix in the complete
first-stage copy space.

For \(N_{SC}^{T}>0\), define the exact output-conditioned matrix
\begin{equation}
 \mathsf R_{S,T,X}
 :=
 \frac1{N_{SC}^{T}}
 \sum_{\substack{j\in\mathfrak J_X\\T(j)=T}}
 \sum_{\beta=1}^{N_{SC}^{T}}
 W_{j,T}[S,\beta].
 \label{eq:V100-output-conditioned-matrix}
\end{equation}
Terms with \(N_{SC}^{T}=0\) are omitted.  This is a positive matrix on
\(\mathcal M_{AB}^{S}\); it need not be scalar and matrices for different
\(T\)'s need not commute.

\begin{theorem}[Exact matrix-valued fusion representation]
\label{thm:V100-matrix-fusion}
The V98 retained-pair square
\[
 M_{AB,X}^{(2)}=\operatorname{Tr}_{C}\mathbf L_X^2
\]
has the invariant block decomposition
\begin{align}
 M_{AB,X}^{(2)}
 &=
 \bigoplus_S I_{H_S}\otimes\mathsf M_{S,X},
 \label{eq:V100-square-isotypic-blocks}\\
 \mathsf M_{S,X}
 &=
 \frac1{d_S}
 \sum_{\substack{j\in\mathfrak J_X\\T}}
 d_T\sum_{\beta=1}^{N_{SC}^{T}}
 W_{j,T}[S,\beta]
 \label{eq:V100-exact-multiplicity-matrix}\\
 &=
 \boxed{\;
 d_C\sum_T
 \mathsf p_C(T\mid S)\mathsf R_{S,T,X}.
 \;}
 \label{eq:V100-matrix-fusion-average}
\end{align}
Thus the fusion sum is taken before any multiplicity eigenvalue or
operator maximum.
\end{theorem}

\begin{proof}
Apply the exact V67 matrix-valued partial-trace formula
\eqref{eq:V67-matrix-partial-trace} to \(W_{j,T}=L_{j,T}^2\) in each
orthogonal physical outer block, and then sum the blocks.  This proves
\eqref{eq:V100-square-isotypic-blocks} and
\eqref{eq:V100-exact-multiplicity-matrix}.  The V91 fusion kernel satisfies
\[
 d_C\mathsf p_C(T\mid S)
 =
 \frac{d_TN_{SC}^{T}}{d_S}.
\]
Insert definition \eqref{eq:V100-output-conditioned-matrix} to obtain
\eqref{eq:V100-matrix-fusion-average}.  Positivity follows because every
principal block of \(L_{j,T}^2\) is positive.
\end{proof}

\begin{remark}[The fusion law is not an outputwise norm bound]
\label{rem:V100-fusion-not-max}
Replacing every \(\mathsf R_{S,T,X}\) separately by
\(\|\mathsf R_{S,T,X}\|_{\rm op}I\) recovers a sufficient scalar
fusion-kernel majorant of the V91 type.  Equation
\eqref{eq:V100-matrix-fusion-average} is stronger data: it retains
noncommuting multiplicity matrices and lets their sum be diagonalized only
after the common output law has been assembled.
\end{remark}

\section{The actual singular-value distribution}
\label{sec:V100-singular-distribution}

Diagonalize the positive multiplicity matrices after the fusion sum:
\begin{equation}
 \mathsf M_{S,X}
 =
 \sum_{r=1}^{m_{S,X}}
 \sigma_{S,r,X}^{\,2}E_{S,r,X},
 \qquad
 \sigma_{S,r,X}^{\,2}\ge0.
 \label{eq:V100-multiplicity-spectrum}
\end{equation}
Repeated eigenvalues may be grouped into one spectral projection.
The multiset
\begin{equation}
 \mathscr S_{S,X}
 :=
 \{\sigma_{S,r,X}^{\,2}:1\le r\le m_{S,X}\}
 \label{eq:V100-actual-singular-profile}
\end{equation}
is the actual squared singular profile requested in the V99 mandate.
It is a profile of the complete V98 signed physical block, not of the
three separated positive covariance squares rejected in V97.

For a unit product vector \(p=u\otimes v\in H_A\otimes H_B\), define
\begin{equation}
 \pi_p(S,r)
 :=
 \left\|
 (I_{H_S}\otimes E_{S,r,X})P_Sp
 \right\|^2.
 \label{eq:V100-joint-product-law}
\end{equation}

\begin{theorem}[Exact singular-law formula for product capacity]
\label{thm:V100-exact-singular-law}
For every unit product vector \(p\),
\begin{align}
 \pi_p(S,r)&\ge0,
 &
 \sum_{S,r}\pi_p(S,r)&=1,
 \label{eq:V100-law-normalization}\\
 \langle p,M_{AB,X}^{(2)}p\rangle
 &=
 \sum_{S,r}\sigma_{S,r,X}^{\,2}\pi_p(S,r).
 \label{eq:V100-exact-product-expectation}
\end{align}
Consequently
\begin{equation}
 \boxed{\;
 \|M_{AB,X}^{(2)}\|_{\rm prod}
 =
 \sup_{\substack{\|u\|=\|v\|=1}}
 \sum_{S,r}
 \sigma_{S,r,X}^{\,2}\pi_{u\otimes v}(S,r).
 \;}
 \label{eq:V100-exact-product-capacity-law}
\end{equation}
On the Cartan component there is one first-stage copy and
\[
 \sigma_{S_{\rm Car},1,X}^{\,2}
 =
 \mathfrak g_{\rm Car}^{X}.
\]
\end{theorem}

\begin{proof}
The isotypic projections \(P_S\) are orthogonal and sum to the identity.
Within each isotypic block, the spectral projections in
\eqref{eq:V100-multiplicity-spectrum} are orthogonal and sum to the
identity on the multiplicity space.  This proves
\eqref{eq:V100-law-normalization}.  Insert the spectral decomposition into
\eqref{eq:V100-square-isotypic-blocks} to prove
\eqref{eq:V100-exact-product-expectation}, then take the supremum over unit
product vectors.  The Cartan assertion is exactly
\eqref{eq:V99-Cartan-column-scalar}.
\end{proof}

\begin{assessment}[What replaces the global Cartan floor]
\label{ass:V100-replaces-floor}
Equation \eqref{eq:V100-exact-product-capacity-law} retains both pieces of
information which the V99 global floor discards: the complete distribution
of a product vector among the first-stage irreps, and its distribution
among singular directions inside every multiplicity space.  Taking
\(\max_{S,r}\sigma_{S,r,X}^2\) at the outset erases both.  The next section
records one exact statistic of the same law which can sometimes be used
without that loss.
\end{assessment}

\section{The common Casimir is one exact moment of the law}
\label{sec:V100-Casimir-moment}

Work first on one compact connected semisimple gauge coordinate.  Abelian
coordinates have a single tensor-product character and do not create a
non-Cartan decomposition.  Let \(\lambda,\mu\) be the highest weights of
\(A,B\), let \(S_{\rm Car}\) have highest weight \(\lambda+\mu\), and
normalize the quadratic Casimir by
\[
 c_\eta=(\eta,\eta+2\rho).
\]
For every \(S\subset A\otimes B\), define its Cartan Casimir deficit
\begin{equation}
 \delta_S:=c_{\lambda+\mu}-c_S.
 \label{eq:V100-Casimir-deficit}
\end{equation}

\begin{lemma}[Strict Casimir separation from the Cartan component]
\label{lem:V100-strict-Casimir}
One has
\[
 \delta_{S_{\rm Car}}=0,
 \qquad
 \delta_S>0\quad(S\ne S_{\rm Car}).
\]
\end{lemma}

\begin{proof}
Every constituent highest weight \(\nu\) of
\(V_\lambda\otimes V_\mu\) is dominated by \(\lambda+\mu\).  Put
\(\gamma=\lambda+\mu-\nu\), a nonnegative integral combination of positive
roots.  If \(S\) is not the Cartan component, then \(\gamma\ne0\).  Direct
subtraction gives
\[
 c_{\lambda+\mu}-c_\nu
 =
 (\gamma,\lambda+\mu+\nu+2\rho).
\]
The second factor is strictly dominant because of \(2\rho\), and its inner
product with every nonzero vector in the positive-root cone is positive.
Thus the displayed difference is positive.  It vanishes for
\(\nu=\lambda+\mu\).
\end{proof}

Choose an orthonormal Lie-algebra basis and self-adjoint infinitesimal
generators \(Y_a^A,Y_a^B\), so that
\[
 C_{AB}
 =
 C_A\otimes I+I\otimes C_B
 +2\sum_aY_a^A\otimes Y_a^B.
\]
For a unit vector \(u\in H_A\), define its moment vector by
\[
 \mathfrak m_A(u)_a:=\langle u,Y_a^Au\rangle,
\]
and similarly for \(B\).

\begin{theorem}[Casimir moment identity]
\label{thm:V100-Casimir-moment}
For a unit product vector \(p=u\otimes v\), put
\[
 \mathcal D_{\lambda,\mu}(u,v)
 :=
 \langle p,(c_{\lambda+\mu}I-C_{AB})p\rangle.
\]
Then
\begin{align}
 \boxed{\;
 \mathcal D_{\lambda,\mu}(u,v)
 &=
 \sum_{S,r}\delta_S\pi_p(S,r)
 \;}
 \label{eq:V100-Casimir-law-moment}\\
 &=
 2\left[
 (\lambda,\mu)
 -
 \mathfrak m_A(u)\cdot\mathfrak m_B(v)
 \right].
 \label{eq:V100-moment-map-deficit}
\end{align}
If
\[
 q_{\rm Car}(p):=\|P_{\rm Car}p\|^2,
\quad
 \delta_{\min}:=\min_{S\ne S_{\rm Car}}\delta_S,
\quad
 \delta_{\max}:=\max_{S\ne S_{\rm Car}}\delta_S,
\]
then
\begin{equation}
 \boxed{\;
 \delta_{\min}(1-q_{\rm Car}(p))
 \le
 \mathcal D_{\lambda,\mu}(u,v)
 \le
 \delta_{\max}(1-q_{\rm Car}(p)).
 \;}
 \label{eq:V100-Casimir-Cartan-sandwich}
\end{equation}
\end{theorem}

\begin{proof}
The Casimir acts by \(c_S\) on the whole \(S\)-isotypic block and is
independent of its multiplicity coordinate.  The spectral projections in
\eqref{eq:V100-joint-product-law} therefore give
\eqref{eq:V100-Casimir-law-moment}.  On a product vector, the coproduct
formula gives
\[
 \langle p,C_{AB}p\rangle
 =
 c_\lambda+c_\mu
 +2\mathfrak m_A(u)\cdot\mathfrak m_B(v).
\]
Since
\[
 c_{\lambda+\mu}-c_\lambda-c_\mu=2(\lambda,\mu),
\]
this proves \eqref{eq:V100-moment-map-deficit}.  Finally, the total
non-Cartan probability is \(1-q_{\rm Car}(p)\), and every non-Cartan
deficit lies between \(\delta_{\min}\) and \(\delta_{\max}\).  Apply these
pointwise bounds in \eqref{eq:V100-Casimir-law-moment}.
\end{proof}

\begin{remark}[Product groups]
\label{rem:V100-product-groups}
For a product of compact semisimple coordinates, the total Casimir and
the deficit are sums of the coordinate Casimirs and deficits.  The law
\eqref{eq:V100-joint-product-law} refines their common joint isotypic law.
One may use a coordinate Casimir only on the same coordinate as the
physical fusion/singular profile.  A disjoint Cartan spectator cannot
supply decay, by the V93 spectator-inflation theorem.
\end{remark}

\section{The sharp one-moment relaxation}
\label{sec:V100-concave-envelope}

Let
\[
 M_{\rm nc}:=(I-P_{\rm Car})M_{AB,X}^{(2)}(I-P_{\rm Car}),
\qquad
 \delta_*:=\max_S\delta_S.
\]
Form the finite data set
\begin{equation}
 \mathcal P_X
 :=
 \{(0,0)\}
 \cup
 \left\{
 (\delta_S,\sigma_{S,r,X}^{\,2}):
 S\ne S_{\rm Car},\ 1\le r\le m_{S,X}
 \right\}.
 \label{eq:V100-profile-points}
\end{equation}
Let \(\widehat{\mathscr S}_X:[0,\delta_*]\to[0,\infty)\) be the least
concave function which majorizes every point in \(\mathcal P_X\).
Equivalently, its graph is the upper concave hull of these finitely many
profile points.

\begin{theorem}[Casimir--singular concave-envelope gate]
\label{thm:V100-concave-envelope-gate}
For every unit product vector \(p=u\otimes v\),
\begin{equation}
 \boxed{\;
 \langle p,M_{\rm nc}p\rangle
 \le
 \widehat{\mathscr S}_X
 \bigl(\mathcal D_{\lambda,\mu}(u,v)\bigr).
 \;}
 \label{eq:V100-concave-envelope-bound}
\end{equation}
Hence
\begin{equation}
 \|M_{\rm nc}\|_{\rm prod}
 \le
 \sup_{\|u\|=\|v\|=1}
 \widehat{\mathscr S}_X
 \left(
 2[(\lambda,\mu)
 -\mathfrak m_A(u)\cdot\mathfrak m_B(v)]
 \right).
 \label{eq:V100-moment-map-capacity-gate}
\end{equation}
Among relaxations which know only the normalization of a probability law,
the allowed profile points \(\mathcal P_X\), and its first Casimir moment,
the upper bound \(\widehat{\mathscr S}_X\) is sharp.
\end{theorem}

\begin{proof}
Give the Cartan outcome the value zero, as \(M_{\rm nc}\) kills that
block.  For every outcome \((S,r)\), the definition of the concave majorant
gives
\[
 \mathbf1_{\{S\ne S_{\rm Car}\}}\sigma_{S,r,X}^{\,2}
 \le
 \widehat{\mathscr S}_X(\delta_S).
\]
Multiply by \(\pi_p(S,r)\), sum, and apply Jensen's inequality using
\eqref{eq:V100-law-normalization}.  The argument of the concave function
is exactly \(\mathcal D_{\lambda,\mu}(u,v)\) by
\eqref{eq:V100-Casimir-law-moment}.  This proves
\eqref{eq:V100-concave-envelope-bound}; take the product supremum and use
\eqref{eq:V100-moment-map-deficit} for
\eqref{eq:V100-moment-map-capacity-gate}.

For sharpness of the probability-law relaxation, fix a mean \(t\).
The supremum of
\(\sum_i\theta_i y_i\) subject to
\(\theta_i\ge0\), \(\sum_i\theta_i=1\), and
\(\sum_i\theta_i x_i=t\), over points
\((x_i,y_i)\in\mathcal P_X\), is the upper boundary of their convex hull.
That boundary is precisely the least concave majorant at \(t\).
Product-state laws form a possibly proper subset of these probability
laws, so no claim of product attainability is made.
\end{proof}

\begin{corollary}[Power-profile sufficient conditions]
\label{cor:V100-power-profile}
Suppose, for some \(a\ge0\) and \(0<\theta\le1\),
\[
 \sigma_{S,r,X}^{\,2}
 \le a\delta_S^\theta
 \qquad(S\ne S_{\rm Car}).
\]
Then
\begin{equation}
 \langle p,M_{\rm nc}p\rangle
 \le
 a\mathcal D_{\lambda,\mu}(u,v)^\theta
 \label{eq:V100-power-profile-bound}
\end{equation}
for every product vector \(p=u\otimes v\).  In particular, a linear
profile gives the exact first-moment bound
\[
 \langle p,M_{\rm nc}p\rangle
 \le a\mathcal D_{\lambda,\mu}(u,v).
\]
\end{corollary}

\begin{proof}
The function \(t\mapsto at^\theta\) is concave on
\([0,\infty)\), vanishes at zero, and majorizes every profile point.
Apply the proof of \Cref{thm:V100-concave-envelope-gate} directly.
\end{proof}

\begin{firewall}[A Casimir moment is not the full product law]
\label{fire:V100-one-moment}
The sharpness assertion in
\Cref{thm:V100-concave-envelope-gate} concerns the relaxation to arbitrary
probability laws.  It does not say that every law on the upper hull is
realized by a product vector.  Conversely, the Casimir moment alone cannot
force a representation-uniform debit: the next section gives product laws
whose non-Cartan mass tends to one while their Casimir deficit grows.
\end{firewall}

\section{A non-Cartan component with visibility tending to one}
\label{sec:V100-near-unit-example}

Let
\[
 F=\mathbb C^3,\qquad R_N=\operatorname{Sym}^N\mathbb C^3.
\]
For \(N\ge1\), the Pieri decomposition is
\begin{equation}
 F\otimes R_N
 =
 R_{N+1}\oplus V_{(N-1,1)}.
 \label{eq:V100-Pieri-two-block}
\end{equation}
The first summand is Cartan and the second is the sole non-Cartan summand.
Let their projections be \(P_{\rm Car}^{(N)}\) and
\(P_{\rm nc}^{(N)}\).

For a unit \(w\in R_N\), let \(\gamma_w\) be its normalized one-particle
reduced density matrix, so \(\gamma_w\ge0\) and
\(\operatorname{Tr}\gamma_w=1\).

\begin{theorem}[Exact near-unit non-Cartan visibility]
\label{thm:V100-near-unit-nonCartan}
For unit \(u\in F\) and \(w\in R_N\),
\begin{equation}
 \boxed{\;
 \left\|P_{\rm Car}^{(N)}(u\otimes w)\right\|^2
 =
 \frac{1+N\langle u,\gamma_wu\rangle}{N+1}.
 \;}
 \label{eq:V100-Cartan-RDM-formula}
\end{equation}
Consequently
\begin{align}
 \inf_{\|u\|=\|w\|=1}
 \left\|P_{\rm Car}^{(N)}(u\otimes w)\right\|^2
 &=
 \frac1{N+1},
 \label{eq:V100-exact-Cartan-floor-FRN}\\
 \boxed{\;
 \beta_{F,R_N}\!\left(V_{(N-1,1)}\right)
 =
 \frac{N}{N+1}
 \longrightarrow1.
 \;}
 \label{eq:V100-nonCartan-visibility-to-one}
\end{align}
\end{theorem}

\begin{proof}
View \(F\otimes R_N\) inside \(F^{\otimes(N+1)}\), with the last \(N\)
slots already symmetric.  On this subspace the full symmetrizer is
\[
 P_{\rm Car}^{(N)}
 =
 \frac1{N+1}
 \left(I+\sum_{k=1}^{N}\tau_{0k}\right),
\]
where \(\tau_{0k}\) swaps the first slot with the \(k\)-th symmetric slot.
Every swap has expectation
\[
 \langle u\otimes w,\tau_{0k}(u\otimes w)\rangle
 =
 \langle u,\gamma_wu\rangle.
\]
This proves \eqref{eq:V100-Cartan-RDM-formula}.  Positivity of
\(\gamma_w\) gives the lower bound \(1/(N+1)\).  It is attained by
\[
 u=e_2,\qquad w=e_1^{\odot N},
\]
for which \(\gamma_w=|e_1\rangle\langle e_1|\).  Since
\eqref{eq:V100-Pieri-two-block} has exactly two orthogonal summands,
the product norm of \(P_{\rm nc}^{(N)}\) is one minus the minimum Cartan
probability.  This proves both remaining assertions.
\end{proof}

The preceding witness is the \(a=1,b=N\) member of the V99
orthogonal-monomial family, but
\eqref{eq:V100-Cartan-RDM-formula} proves that its binomial value is the
global minimum, not merely one test value.

\begin{theorem}[The near-unit witnesses recede from the Cartan locus]
\label{thm:V100-recede-Cartan}
Let
\[
 p_N:=e_2\otimes e_1^{\odot N}
\]
and let
\[
 \mathcal Z_N
 :=
 \left\{
 [z\otimes z^{\odot N}]:
 z\in\mathbb C^3,\ \|z\|=1
 \right\}
\]
be the projective locus of decomposable tensors lying entirely in the
Cartan summand.  Then
\begin{equation}
 \boxed{\;
 \sup_{\|z\|=1}
 \left|
 \left\langle
 p_N,z\otimes z^{\odot N}
 \right\rangle
 \right|^2
 =
 \frac{N^N}{(N+1)^{N+1}}
 \longrightarrow0.
 \;}
 \label{eq:V100-Cartan-locus-fidelity}
\end{equation}
Thus the Fubini--Study distance from \([p_N]\) to \(\mathcal Z_N\)
tends to \(\pi/2\).  In particular, high non-Cartan visibility does not
force approach to the Cartan variety.
\end{theorem}

\begin{proof}
A decomposable tensor \(u\otimes w\in F\otimes R_N\) lies in
\(R_{N+1}\) precisely when it is symmetric in all \(N+1\) slots.
If its first-slot flattening has rank one, complete symmetry forces every
slot to have the same one-dimensional range.  Hence
\(u\otimes w\) is proportional to \(z\otimes z^{\odot N}\), proving the
description of \(\mathcal Z_N\).

For a unit \(z\), put
\[
 x=|\langle e_1,z\rangle|^2,
 \qquad
 y=|\langle e_2,z\rangle|^2.
\]
Then \(y\le1-x\) and
\[
 \left|
 \left\langle
 p_N,z\otimes z^{\odot N}
 \right\rangle
 \right|^2
 =
 yx^N
 \le(1-x)x^N.
\]
The last function on \([0,1]\) is maximized at
\(x=N/(N+1)\), with value
\(N^N/(N+1)^{N+1}\).  Equality is obtained by a unit vector in
\(\operatorname{span}\{e_1,e_2\}\) with the corresponding squared
coordinates.  The value is asymptotic to \(e^{-1}/(N+1)\), proving the
limit and the projective-distance assertion.
\end{proof}

\begin{corollary}[No representation-uniform non-Cartan debit]
\label{cor:V100-no-uniform-nonCartan-debit}
There is no constant \(\varepsilon>0\) such that
\[
 \|P_{\rm nc}\|_{\rm prod}\le1-\varepsilon
\]
for every non-Cartan isotypic projector in every pair of \(SU(3)\)
irreducibles.  More sharply, if
\[
 M_N=s_N^2P_{\rm nc}^{(N)},
\]
then
\begin{equation}
 \|M_N\|_{\rm prod}
 =
 s_N^2\frac{N}{N+1}.
 \label{eq:V100-scaled-near-unit-capacity}
\end{equation}
Thus a uniform gain on this family must come from the physical singular
coefficient \(s_N\), not from non-Cartan geometry.
\end{corollary}

\begin{proof}
Use \eqref{eq:V100-nonCartan-visibility-to-one}.  Scalar multiplication of
a positive projector multiplies its product norm by the same scalar.
\end{proof}

\section{Casimir value of the near-unit sequence}
\label{sec:V100-near-unit-Casimir}

\begin{proposition}[Large deficit accompanies near-unit visibility]
\label{prop:V100-near-unit-Casimir}
For the two summands in \eqref{eq:V100-Pieri-two-block},
\begin{align}
 \delta_{R_{N+1}}&=0,
 &
 \delta_{V_{(N-1,1)}}&=N+1.
 \label{eq:V100-two-block-deficits}
\end{align}
The witness \(p_N=e_2\otimes e_1^{\odot N}\) satisfies
\begin{equation}
 \mathcal D_{\omega_1,N\omega_1}(p_N)
 =
 N.
 \label{eq:V100-witness-deficit}
\end{equation}
For \(M_N=s_N^2P_{\rm nc}^{(N)}\), the concave-envelope gate
\eqref{eq:V100-concave-envelope-bound} is exact at \(p_N\).
\end{proposition}

\begin{proof}
Using
\[
 c_{(p,q)}
 =
 \frac{p^2+q^2+pq+3p+3q}{3},
\]
one obtains
\[
 c_{(N+1,0)}-c_{(N-1,1)}
 =
 N+1.
\]
The non-Cartan probability of \(p_N\) is \(N/(N+1)\), so
\eqref{eq:V100-Casimir-law-moment} gives
\[
 \mathcal D(p_N)
 =
 (N+1)\frac{N}{N+1}=N.
\]
For \(M_N\), the profile points are
\((0,0)\) and \((N+1,s_N^2)\).  Their least concave majorant is
\[
 \widehat{\mathscr S}_N(t)
 =
 \frac{s_N^2}{N+1}t.
\]
At \(t=N\), its value is
\(s_N^2N/(N+1)=\langle p_N,M_Np_N\rangle\).
\end{proof}

\begin{assessment}[Disposition of the V99 near-Cartan question]
\label{ass:V100-near-Cartan-disposition}
The conditional clause in the V99 mandate has a negative premise:
near-unit non-Cartan visibility need not arise near the Cartan variety.
In the explicit family above it arises from increasingly misaligned
product factors, the Cartan fidelity tends to zero, and the common Casimir
deficit grows linearly.  A local expansion around the Cartan orbit therefore
cannot control the worst non-Cartan product states.
\end{assessment}

\section{Dual-singlet and balanced-thermal regressions}
\label{sec:V100-regressions}

\begin{proposition}[The dual singlet is exactly recovered]
\label{prop:V100-dual-singlet}
For
\[
 F\otimes F^*=\mathrm{Ad}\oplus\mathbf1,
\]
the Cartan component is \(\mathrm{Ad}\), the singlet deficit is \(3\), and
for \(M_{\rm nc}=s^2P_{\mathbf1}\),
\begin{equation}
 \|M_{\rm nc}\|_{\rm prod}
 =
 \frac{s^2}{3}.
 \label{eq:V100-dual-singlet-capacity}
\end{equation}
The Casimir concave envelope gives equality.
\end{proposition}

\begin{proof}
The adjoint Casimir is \(c_{(1,1)}=3\), while the singlet Casimir is zero.
The singlet probability of a product vector is at most \(1/3\), with
equality on a matched vector, by \eqref{eq:V92-singlet-product-norm}.
Thus \eqref{eq:V100-dual-singlet-capacity} follows.  The profile points are
\((0,0)\) and \((3,s^2)\), so the concave majorant is \(s^2t/3\).
At a maximizing product vector the Casimir moment is \(3(1/3)=1\), and the
envelope value is \(s^2/3\).
\end{proof}

\begin{proposition}[Balanced thermal centering is not a singular debit]
\label{prop:V100-balanced-thermal}
In the V89 family \(R_m\otimes R_m\), the first-stage deficits are
\begin{equation}
 \delta_{S_r}=r(2m-r+1).
 \label{eq:V100-V89-deficit-nodes}
\end{equation}
On the transition window
\(\tau m^2\to u\in(0,\infty)\), \(k\asymp\sqrt m\), the bare connected heat
operator satisfies
\[
 K_{\tau;m,k}^{\mathsf h}
 \longrightarrow
 \mathfrak d(u)I
\]
in operator norm, uniformly in the window, with
\(\mathfrak d(u)>0\).  Therefore its squared singular profile tends
uniformly to the nonzero scalar \(\mathfrak d(u)^2\); centering alone does
not supply a compensating small singular value on those ancestry
directions.
\end{proposition}

\begin{proof}
Equation \eqref{eq:V100-V89-deficit-nodes} is the exact V89 identity
\eqref{eq:V89-AB-Casimir}.  The uniform scalar limit and its positivity are
\eqref{eq:V89-full-heat-limit} and
\eqref{eq:V89-mean-square-limit}.  If self-adjoint matrices converge in
operator norm to \(\mathfrak d(u)I\), their squares converge in operator
norm to \(\mathfrak d(u)^2I\), which proves the singular-profile statement.
\end{proof}

\begin{firewall}[The V89 family remains paid, not residual]
\label{fire:V100-V89-paid}
\Cref{prop:V100-balanced-thermal} is a regression test on the bare connected
heat symbol.  It does not reopen the V89 branch.  The complete Cartan
thermal pair--core carrier was already bounded at order one by V91 before
the V98 residual side/off partition.  Nor does the proposition assert that
the additional physical weights in every residual block have a nonzero
limit.
\end{firewall}

\section{The exact upstream gate left by V100}
\label{sec:V100-upstream-gate}

The near-unit family identifies the only possible source of a uniform gain
on its non-Cartan channel.  Let
\[
 S_N:=V_{(N-1,1)}
 \subset F\otimes R_N.
\]
Because this decomposition is multiplicity-free, the \(S_N\)-block of
\(\mathsf M_{S_N,X}\) is a scalar; denote it by
\[
 \mathfrak s_{N,X}^2.
\]
It is obtained from the physical matrix fusion average
\eqref{eq:V100-matrix-fusion-average}, including all final-output and outer
collision weights.

\begin{theorem}[Necessary compensation on the near-unit channel]
\label{thm:V100-necessary-compensation}
For every branch \(X\) on which the two-block family occurs,
\begin{equation}
 \|M_{AB,X}^{(2)}\|_{\rm prod}
 \ge
 \frac{N}{N+1}\mathfrak s_{N,X}^2.
 \label{eq:V100-necessary-symbol-lower}
\end{equation}
Consequently, any proposed representation-uniform estimate
\[
 \|M_{AB,X}^{(2)}\|_{\rm prod}\le\varepsilon_N
\]
with \(\varepsilon_N\to0\) must prove
\[
 \mathfrak s_{N,X}^2
 \le
 \frac{N+1}{N}\varepsilon_N.
\]
No visibility or Cartan-floor argument can supply this conclusion.
\end{theorem}

\begin{proof}
Evaluate \(M_{AB,X}^{(2)}\) on
\(p_N=e_2\otimes e_1^{\odot N}\), discard all nonnegative blocks except
the scalar \(S_N\)-block, and use
\[
 \|P_{S_N}p_N\|^2=\frac{N}{N+1}.
\]
This proves \eqref{eq:V100-necessary-symbol-lower}.  Rearrangement gives
the necessary condition.  The family has visibility tending to one by
\eqref{eq:V100-nonCartan-visibility-to-one}, so a geometric debit bounded
away from one is impossible.
\end{proof}

\begin{assessment}[Audit-adjusted acquisition]
\label{ass:V100-audit-adjusted-acquisition}
The remaining side-branch question has moved upstream from abstract
product visibility to the physical scalar
\(\mathfrak s_{N,\rm side}^2\).  It must be computed from the complete
signed \(L_j\), or bounded after proving a factorization through its exact
analysis image.  A sum of separate positive pair covariances is
inadmissible because V97 already exhibits a lost connected kernel.
Likewise, a positive estimate on a decoupled coordinate cannot repair this
same-coordinate family, consistent with the V100 companion audit.

The off-cut branch remains a different physical carrier.  None of the
side-branch profile points, product laws, or moment-map optimizers may be
chosen coordinatewise and then reused to lower its capacity.
\end{assessment}

\section{V100 ledger and V101 acquisition}
\label{sec:V100-ledger}

\begin{theorem}[Integrated V100 statement]
\label{thm:V100-integrated}
Preserving V1--V99, V100 proves:
\begin{enumerate}[label=\textup{(V100.\arabic*)},leftmargin=6.3em]
\item the compulsory current SM/NS audit and its exact logical scope;
\item the matrix-valued physical fusion law
      \eqref{eq:V100-matrix-fusion-average};
\item the exact singular-distribution formula
      \eqref{eq:V100-exact-product-capacity-law};
\item the common Casimir and moment-map identities
      \eqref{eq:V100-Casimir-law-moment}--%
      \eqref{eq:V100-moment-map-deficit};
\item the sharp one-moment concave-envelope relaxation
      \eqref{eq:V100-concave-envelope-bound};
\item the exact near-unit non-Cartan visibility
      \eqref{eq:V100-nonCartan-visibility-to-one};
\item the quantitative failure of approach to the Cartan locus
      \eqref{eq:V100-Cartan-locus-fidelity};
\item the exact dual-singlet regression and the balanced-thermal
      non-smallness regression; and
\item the necessary physical-symbol compensation
      \eqref{eq:V100-necessary-symbol-lower}.
\end{enumerate}
\end{theorem}

\begin{proof}
The items are respectively
\Cref{fire:V100-audit-scope,
thm:V100-matrix-fusion,
thm:V100-exact-singular-law,
thm:V100-Casimir-moment,
thm:V100-concave-envelope-gate,
thm:V100-near-unit-nonCartan,
thm:V100-recede-Cartan,
prop:V100-dual-singlet,
prop:V100-balanced-thermal,
thm:V100-necessary-compensation}.
\end{proof}

\begin{openproblem}[V101 mandate: compute the near-unit physical symbol]
\label{op:V101-mandate}
Keep the V98 physical side orientation fixed and specialize one legitimate
side block to
\[
 F\otimes R_N
 =
 R_{N+1}\oplus V_{(N-1,1)}.
\]
Then:
\begin{enumerate}
\item write the exact first-stage non-Cartan scalar
      \(\mathfrak s_{N,\rm side}^2\) from
      \eqref{eq:V100-matrix-fusion-average}, without replacing the signed
      cumulant by three positive squares;
\item separate final outputs and physical outer weights only after the
      common scalar has been assembled;
\item determine whether centering, an internal-pair symbol, or the
      \(1+s_\alpha\Xi(T)\) allocation yields decay on this channel;
\item test the result against the lower bound
      \eqref{eq:V100-necessary-symbol-lower};
\item if the scalar cannot decay, identify which earlier paid selector must
      remove the family, rather than declaring it residual;
\item keep the off-cut carrier separate; and
\item retain the exact fusion average and Casimir law for all
      higher-multiplicity channels not covered by this two-block test.
\end{enumerate}
The next compulsory companion audit is V105.
\end{openproblem}

\section{Firewall after V100}
\label{sec:V100-firewall}

\begin{firewall}[Current claim boundary]
\label{fire:V100-current-boundary}
V100 does not prove the residual side estimate
\eqref{eq:V98-side-image-gate}, close the off-cut branch, aggregate all
punctures and collisions, construct the continuum Yang--Mills measure,
prove the Osterwalder--Schrader axioms, construct a Hamiltonian, or prove a
positive spectral gap.  It proves an exact matrix-valued physical
fusion/singular law and shows that non-Cartan geometry alone cannot provide
a uniform debit.  The full Yang--Mills existence and mass gap remain open.
\end{firewall}

\section{Append-only record for V100}
\label{sec:V100-record}

V100 was appended to the complete V99 manuscript.  The exact V99 parent
had \(67361\) physical lines, \(2360927\) bytes, and SHA--256 digest
\[
 \texttt{c40128565676b32defccb8bc13c7186ca99c91d49fee6f36b7d107d9a5a92d71}.
\]
No V99 mathematical content was changed.  The sole final
\verb|\end{document}| is restored below.

% =============================================================================
% END APPENDED VERSION V100 MATERIAL
% =============================================================================
% =============================================================================
% BEGIN APPENDED VERSION V101 MATERIAL
% =============================================================================

\chapter{V101: Multiplicity-One Physical Columns, Exact Crossing Leakage,
and the Residual-Selector Data Contract}
\label{ch:V101}

\section{Purpose: the V100 test must first be physically instantiated}
\label{sec:V101-purpose}

V100 exhibited the representation-theoretic family
\[
 F\otimes R_N=R_{N+1}\oplus S_N,
 \qquad S_N:=V_{(N-1,1)},
\]
for which \(\beta_{F,R_N}(S_N)=N/(N+1)\).  That theorem is an exact
stress test for any proposed representation-uniform product debit.
It does not, by itself, assert that an arbitrary completion of these
first-stage labels belongs to the V98 simultaneous-sideways bank.

The V94 priority selector assigns a complete physical outer outcome to
the private-mean, cut-heavy, selected-heavy-pair, side, or off-cut bank
before the retained-row norm is taken.  Side membership therefore depends
on the full collision/puncture monomial and not only on \(A,B,S_N\).
Attempting to compute a side scalar from the three first-stage labels alone
would silently discard the very selector which prevents double charging.

This chapter moves upstream accordingly.  First it writes the exact
multiplicity-one scalar for any complete physical outer family.  Second it
splits each full signed cumulant column into two orthogonal nonnegative
pieces: crossing leakage away from the selected path and diagonal mismatch
on that path.  Third it computes the aligned heat defect on the V100 family
and proves that it is \(O(N^{-1})\) in the natural heat window.  The exact
column split then shows why this promising scalar fact is not yet a bound
for the full cumulant: the crossing leakage is independent data.

\begin{firewall}[No arbitrary pair is promoted to the side bank]
\label{fire:V101-no-arbitrary-side-promotion}
All statements below are conditional on a complete physical outer family
already assigned to a fixed V98 branch \(X\).  The chapter never declares
the V100 representation pair to be simultaneous-sideways merely because
its non-Cartan visibility is large.  A genuine side sequence must be
extracted from the V59/V60/V94 selector with all core, puncture, output,
and weight labels present.
\end{firewall}

\section{Exact multiplicity-one scalar as a fusion mean}
\label{sec:V101-exact-scalar}

Fix \(X\in\{\mathrm{side},\mathrm{off}\}\), a physical orientation \(A|B\),
and a multiplicity-one intermediate irrep \(S\subset A\otimes B\):
\[
 N_{AB}^{S}=1.
\]
For every active \(j,T,\beta\), let
\[
 e_{S,\beta}^{j,T}
\]
be the corresponding unit \(AB\)-first path vector.  Put
\[
 w_T:=1+s_\alpha\Xi(T).
\]
By V98,
\[
 L_{j,T}=c_jw_TK_T.
\]

For \(N_{SC}^{T}>0\), define the complete conditional column energy
\begin{equation}
 \mathcal E_{S,T,X}
 :=
 \frac1{N_{SC}^{T}}
 \sum_{\substack{j\in\mathfrak J_X\\T(j)=T}}
 \sum_{\beta=1}^{N_{SC}^{T}}
 c_j^2w_T^2
 \left\|K_Te_{S,\beta}^{j,T}\right\|^2.
 \label{eq:V101-conditional-column-energy}
\end{equation}
Absent paths contribute zero.

\begin{theorem}[Exact physical scalar fusion identity]
\label{thm:V101-exact-scalar-fusion}
The multiplicity-one block of
\[
 M_{AB,X}^{(2)}=\operatorname{Tr}_C\mathbf L_X^2
\]
on \(S\) is
\[
 I_{H_S}\otimes \mathfrak s_{S,X}^2,
\]
where
\begin{align}
 \boxed{\;
 \mathfrak s_{S,X}^2
 &=
 \frac1{d_S}
 \sum_{\substack{j\in\mathfrak J_X\\T}}
 d_T\sum_{\beta=1}^{N_{SC}^{T}}
 c_j^2w_T^2
 \left\|K_Te_{S,\beta}^{j,T}\right\|^2
 \;}
 \label{eq:V101-exact-physical-scalar}\\
 &=
 \boxed{\;
 d_C\sum_T
 \mathsf p_C(T\mid S)\mathcal E_{S,T,X}.
 \;}
 \label{eq:V101-scalar-fusion-mean}
\end{align}
In particular,
\begin{equation}
 \mathfrak s_{S,X}^2=0
 \quad\Longleftrightarrow\quad
 c_jw_TK_Te_{S,\beta}^{j,T}=0
 \ \hbox{for every active }j,T,\beta.
 \label{eq:V101-scalar-kernel-criterion}
\end{equation}
\end{theorem}

\begin{proof}
Since \(N_{AB}^{S}=1\), the matrix
\(\mathsf M_{S,X}\) in
\eqref{eq:V100-exact-multiplicity-matrix} is scalar.  Its
\((S,\beta)\) principal entry from \(L_{j,T}^2\) is
\[
 \left\langle
 e_{S,\beta}^{j,T},L_{j,T}^2e_{S,\beta}^{j,T}
 \right\rangle
 =
 c_j^2w_T^2
 \left\|K_Te_{S,\beta}^{j,T}\right\|^2.
\]
Substitution proves \eqref{eq:V101-exact-physical-scalar}.  Group the
terms by \(T\), divide and multiply by \(N_{SC}^{T}\), and use
\[
 \frac{d_TN_{SC}^{T}}{d_S}
 =
 d_C\mathsf p_C(T\mid S)
\]
to prove \eqref{eq:V101-scalar-fusion-mean}.  Every summand in
\eqref{eq:V101-exact-physical-scalar} is nonnegative and \(w_T>0\).
Therefore the sum vanishes exactly when every active weighted column
vanishes, proving \eqref{eq:V101-scalar-kernel-criterion}.
\end{proof}

\begin{corollary}[Output subsets give honest lower tests]
\label{cor:V101-output-subset-lower}
For every output set \(\Omega\),
\begin{equation}
 \mathfrak s_{S,X}^2
 \ge
 d_C\sum_{T\in\Omega}
 \mathsf p_C(T\mid S)\mathcal E_{S,T,X}.
 \label{eq:V101-output-subset-lower}
\end{equation}
If the right side stays bounded below on a sequence of physical blocks,
no product-visibility argument can make that scalar vanish.
\end{corollary}

\begin{proof}
Delete the nonnegative terms with \(T\notin\Omega\) from
\eqref{eq:V101-scalar-fusion-mean}.
\end{proof}

\begin{assessment}[The first-stage labels do not determine the scalar]
\label{ass:V101-first-stage-underdetermined}
The data \(A,B,S\) determine the aligned multiplier
\(q_{\alpha,S}-q_{\alpha,A}q_{\alpha,B}\).  They do not determine
\(\mathfrak s_{S,X}^2\).  Formula
\eqref{eq:V101-exact-physical-scalar} also requires \(C\), the final
outputs \(T\), all physical outer weights \(c_j\), the actual path
embeddings, and the action of the two crossing pair moments in \(K_T\).
It also requires the prior fact \(j\in\mathfrak J_X\).  This is the exact
data contract for a legitimate V100 near-unit stress test.
\end{assessment}

\section{Exact crossing-leakage decomposition of one column}
\label{sec:V101-crossing-leakage}

Recall the exact V69 covariance decomposition
\[
 K_T
 =
 D_{A\mid BC}^{T}
 -q_{\alpha,B}D_{AC}^{T}
 -q_{\alpha,C}D_{AB}^{T}.
\]
On the \(S\)-ancestry space,
\[
 D_{AB}^{T}
 =
 \delta_{AB}^{\alpha}(S)I,
 \qquad
 \delta_{AB}^{\alpha}(S)
 :=
 q_{\alpha,S}
 -q_{\alpha,A}q_{\alpha,B}.
\]
Put
\begin{equation}
 R_T
 :=
 D_{A\mid BC}^{T}
 -q_{\alpha,B}D_{AC}^{T},
 \qquad
 z:=q_{\alpha,C}.
 \label{eq:V101-crossing-operator}
\end{equation}
Thus, on every \(S\)-path,
\begin{equation}
 K_Te=R_Te-z\delta_{AB}^{\alpha}(S)e.
 \label{eq:V101-column-signed-split}
\end{equation}

\begin{lemma}[Orthogonal leakage--mismatch identity]
\label{lem:V101-column-orthogonal-split}
Let \(e=e_{S,\beta}^{j,T}\) be a unit path vector and
\(P_e=|e\rangle\langle e|\).  Then
\begin{equation}
 \boxed{\;
 \|K_Te\|^2
 =
 \|(I-P_e)R_Te\|^2
 +
 \left|
 \langle e,R_Te\rangle
 -z\delta_{AB}^{\alpha}(S)
 \right|^2.
 \;}
 \label{eq:V101-leakage-mismatch-identity}
\end{equation}
\end{lemma}

\begin{proof}
Decompose
\[
 R_Te
 =
 (I-P_e)R_Te+\langle e,R_Te\rangle e.
\]
The two terms are orthogonal.  Subtracting
\(z\delta_{AB}^{\alpha}(S)e\) changes only the second coefficient.
Pythagoras gives \eqref{eq:V101-leakage-mismatch-identity}.
\end{proof}

Define the fusion-weighted crossing leakage and diagonal mismatch by
\begin{align}
 \mathfrak l_{S,X}^2
 &:=
 \frac1{d_S}
 \sum_{\substack{j\in\mathfrak J_X\\T}}
 d_T\sum_{\beta=1}^{N_{SC}^{T}}
 c_j^2w_T^2
 \left\|
 (I-P_{e_{S,\beta}^{j,T}})
 R_Te_{S,\beta}^{j,T}
 \right\|^2,
 \label{eq:V101-total-crossing-leakage}\\
 \mathfrak d_{S,X}^2
 &:=
 \frac1{d_S}
 \sum_{\substack{j\in\mathfrak J_X\\T}}
 d_T\sum_{\beta=1}^{N_{SC}^{T}}
 c_j^2w_T^2
 \left|
 \left\langle
 e_{S,\beta}^{j,T},
 R_Te_{S,\beta}^{j,T}
 \right\rangle
 -z\delta_{AB}^{\alpha}(S)
 \right|^2.
 \label{eq:V101-total-diagonal-mismatch}
\end{align}

\begin{theorem}[Exact two-obstruction formula]
\label{thm:V101-two-obstruction}
For every multiplicity-one first-stage channel,
\begin{equation}
 \boxed{\;
 \mathfrak s_{S,X}^2
 =
 \mathfrak l_{S,X}^2+\mathfrak d_{S,X}^2.
 \;}
 \label{eq:V101-two-obstruction-scalar}
\end{equation}
For a sequence of complete physical blocks,
\[
 \mathfrak s_{S,X}^2\longrightarrow0
\]
if and only if both
\[
 \mathfrak l_{S,X}^2\longrightarrow0,
 \qquad
 \mathfrak d_{S,X}^2\longrightarrow0.
\]
\end{theorem}

\begin{proof}
Multiply \eqref{eq:V101-leakage-mismatch-identity} by
\(d_Tc_j^2w_T^2/d_S\) and sum the active \(j,T,\beta\).  Comparison with
\eqref{eq:V101-exact-physical-scalar} gives
\eqref{eq:V101-two-obstruction-scalar}.  Both terms are nonnegative, so
their sum tends to zero exactly when each tends to zero.
\end{proof}

\begin{corollary}[The aligned pair defect is only one scalar input]
\label{cor:V101-aligned-insufficient}
Smallness of
\(\delta_{AB}^{\alpha}(S)\) alone does not imply smallness of
\(\mathfrak s_{S,X}^2\).  It neither bounds the orthogonal leakage
\(\mathfrak l_{S,X}^2\) nor forces the diagonal quantities
\(\langle e,R_Te\rangle\) to match
\(z\delta_{AB}^{\alpha}(S)\).
\end{corollary}

\begin{proof}
This follows directly from the two independent nonnegative terms in
\eqref{eq:V101-two-obstruction-scalar}.  At the level of self-adjoint
finite-dimensional operators the logical independence is sharp: for
orthonormal \(e,f\), the self-adjoint operator
\[
 R=|f\rangle\langle e|+|e\rangle\langle f|
\]
has
\[
 \|(I-P_e)Re\|=1,
 \qquad
 \langle e,Re\rangle=0,
\]
regardless of how small the subtracted scalar is.  This example is a
factorization test; it is not asserted to be a physical Wilson block.
\end{proof}

\section{Exact aligned heat defect on the near-unit channel}
\label{sec:V101-aligned-heat}

Now specialize only the aligned \(AB\) multiplier to the heat Wilson law
\[
 q_{\tau,R}=e^{-\tau c_R},
 \qquad \tau\ge0,
\]
with
\[
 A=F,\qquad B=R_N,\qquad S=S_N=V_{(N-1,1)}.
\]

\begin{theorem}[Exact heat defect and thermal asymptotic]
\label{thm:V101-exact-aligned-heat}
The aligned defect is
\begin{equation}
 \boxed{\;
 \delta_{AB}^{\tau}(S_N)
 =
 e^{-\tau (N+1)^2/3}
 \left(
 1-e^{-\tau(N+3)/3}
 \right).
 \;}
 \label{eq:V101-exact-aligned-defect}
\end{equation}
If
\[
 \tau_NN^2\longrightarrow u\in(0,\infty),
\]
then
\begin{equation}
 \boxed{\;
 N\delta_{AB}^{\tau_N}(S_N)
 \longrightarrow
 \frac u3e^{-u/3}.
 \;}
 \label{eq:V101-aligned-defect-limit}
\end{equation}
In particular the aligned square is \(O(N^{-2})\) on this heat window.
\end{theorem}

\begin{proof}
The \(SU(3)\) Casimir formula gives
\begin{align*}
 c_F&=\frac43,\\
 c_{R_N}&=\frac{N^2+3N}{3},\\
 c_{S_N}&=c_{(N-1,1)}=\frac{(N+1)^2}{3}.
\end{align*}
Hence
\[
 c_F+c_{R_N}-c_{S_N}
 =
 \frac{N+3}{3}.
\]
Therefore
\begin{align*}
 \delta_{AB}^{\tau}(S_N)
 &=
 e^{-\tau c_{S_N}}
 -e^{-\tau(c_F+c_{R_N})}\\
 &=
 e^{-\tau (N+1)^2/3}
 \left(1-e^{-\tau(N+3)/3}\right),
\end{align*}
which proves \eqref{eq:V101-exact-aligned-defect}.  Under
\(\tau_NN^2\to u\),
\[
 e^{-\tau_N(N+1)^2/3}\longrightarrow e^{-u/3}
\]
and
\[
 N\left(1-e^{-\tau_N(N+3)/3}\right)
 \longrightarrow\frac u3.
\]
Multiplication proves \eqref{eq:V101-aligned-defect-limit}.
\end{proof}

\begin{assessment}[What the \(N^{-1}\) defect does and does not prove]
\label{ass:V101-aligned-status}
The near-unit non-Cartan visibility costs no uniform geometric factor, but
its aligned heat covariance is small at the thermal scale.  This is real
same-coordinate information.  It is nevertheless only the scalar
subtracted in \eqref{eq:V101-column-signed-split}.  A physical gain requires
the two full fusion-weighted conditions in
\eqref{eq:V101-two-obstruction-scalar}: crossing leakage must be small and
the crossing diagonal must match the aligned scalar.  Estimating only the
V69 aligned square would repeat the positive-majorant error isolated in
V97.
\end{assessment}

\section{Why the core and selector cannot be omitted}
\label{sec:V101-core-selector}

\begin{proposition}[The trivial-core completion is disconnected, not side]
\label{prop:V101-trivial-core}
Keep arbitrary nontrivial \(A,B\), but take \(C=\mathbf1\).  Then the
complete ternary cumulant vanishes:
\[
 K=0.
\]
Every scalar \(\mathfrak s_{S,X}^2\) formed from such columns is zero.
In the V98 physical partition this is a disconnected-core outcome and may
not be used as a zero-energy example inside
\(\mathfrak J_{\rm side}\).
\end{proposition}

\begin{proof}
For the trivial representation,
\[
 Z_C(g)
 =
 D^{\mathbf1}(g)-q_{\alpha,\mathbf1}I
 =
 I-I=0.
\]
The exact centred identity
\[
 K=\int_G Z_A(g)\otimes Z_B(g)\otimes Z_C(g)\,
 \mu_\alpha(\dd g)
\]
therefore gives \(K=0\).  Formula
\eqref{eq:V101-exact-physical-scalar} gives the scalar zero.  The V98
outer partition assigns a trivial punctured-core input to
\(\mathfrak J_{\rm disc}\), and
\eqref{eq:V98-disconnected-outer-zero} deletes it before the residual
norm.
\end{proof}

\begin{firewall}[Zero completions do not prove a side estimate]
\label{fire:V101-zero-completion}
The first-stage pair \(A\otimes B\) can be held fixed while a trivial core
makes every ternary column zero.  This proves that first-stage data do not
determine the scalar, but it supplies no favorable side block: the priority
partition deletes that completion as disconnected.  Conversely, choosing
an arbitrary nontrivial core does not prove side membership.  Both the
operator and its physical classification require the full outer data.
\end{firewall}

\begin{theorem}[Exact residual-selector dichotomy for the V100 family]
\label{thm:V101-selector-dichotomy}
Let \(\mathfrak D_N\) denote complete physical data containing
\[
 A=F,\qquad B=R_N,\qquad S=S_N,
\]
including the core \(C_N\), all final outputs and LR paths, collision and
puncture labels, outer weights \(c_j\), and the V94 priority assignment.
Exactly one of the following alternatives holds.
\begin{enumerate}[label=\textup{(\roman*)}]
\item No sequence \(\mathfrak D_N\) is assigned to
      \(\mathfrak J_{\rm side}\).  Then the V100 near-unit family is not a
      side obstruction; every physical occurrence is deleted, paid, or
      assigned to the separate off-cut branch by the existing selector.
\item A side sequence exists.  On every multiplicity-one occurrence its
      exact product-capacity lower test is
      \begin{equation}
       \|M_{AB,\rm side}^{(2)}\|_{\rm prod}
       \ge
       \frac{N}{N+1}
       \left(
       \mathfrak l_{S_N,\rm side}^2
       +
       \mathfrak d_{S_N,\rm side}^2
       \right).
       \label{eq:V101-side-two-obstruction-lower}
      \end{equation}
      Hence decay on this sequence requires both fusion-weighted terms to
      decay.
\end{enumerate}
\end{theorem}

\begin{proof}
The V94 priority partition is disjoint and exhaustive, and V98 refines its
residual bank into orthogonal side and off-cut outer supports.  This gives
the exhaustive alternative between existence and nonexistence of a side
sequence with the stated data.

In alternative \textup{(ii)}, V100 gives
\[
 \|P_{S_N}\|_{\rm prod}=\frac{N}{N+1}.
\]
The \(S_N\)-block is scalar by multiplicity one.  Evaluate the positive
square on a product vector attaining this visibility, discard all other
nonnegative blocks, and use
\eqref{eq:V101-two-obstruction-scalar}.  This proves
\eqref{eq:V101-side-two-obstruction-lower}.  Nonnegativity makes decay of
both summands necessary.
\end{proof}

\section{Regression tests}
\label{sec:V101-regressions}

\begin{proposition}[Three consistency tests]
\label{prop:V101-three-tests}
The V101 formula has the following exact dispositions.
\begin{enumerate}[label=\textup{(\roman*)}]
\item On the V97 trivial-input channel, the two terms in
      \eqref{eq:V101-leakage-mismatch-identity} vanish together, despite
      the positive separated-covariance majorant.
\item On the dual singlet, the first-stage multiplicity is one, so
      \eqref{eq:V101-exact-physical-scalar} gives the scalar \(s^2\) used
      in \eqref{eq:V100-dual-singlet-capacity}; the additional product
      debit \(1/3\) is external visibility, not a change in the column
      energy.
\item On the V89 balanced thermal window, the full connected operator has
      a nonzero scalar limit.  Therefore its crossing and aligned pieces
      must assemble to the nonzero total in
      \eqref{eq:V101-leakage-mismatch-identity}; smallness of one displayed
      covariance cannot be substituted for smallness of \(K_Te\).
\end{enumerate}
\end{proposition}

\begin{proof}
For \textup{(i)}, V97 proves \(K_Te=0\).  The right side of
\eqref{eq:V101-leakage-mismatch-identity} is a sum of nonnegative terms and
therefore both vanish.  V97 simultaneously proves that its three-square
analysis map is nonzero, showing why that majorant is not interchangeable
with the V101 terms.

Part \textup{(ii)} follows from multiplicity one and the exact product
visibility \(\|P_{\mathbf1}\|_{\rm prod}=1/3\).
Part \textup{(iii)} is
\Cref{prop:V100-balanced-thermal} inserted into the exact Pythagorean
identity.
\end{proof}

\section{V101 ledger and V102 acquisition}
\label{sec:V101-ledger}

\begin{theorem}[Integrated V101 statement]
\label{thm:V101-integrated}
Preserving V1--V100, V101 proves:
\begin{enumerate}[label=\textup{(V101.\arabic*)},leftmargin=6.3em]
\item the exact multiplicity-one physical scalar and fusion mean
      \eqref{eq:V101-exact-physical-scalar}--%
      \eqref{eq:V101-scalar-fusion-mean};
\item its exact column-kernel criterion
      \eqref{eq:V101-scalar-kernel-criterion};
\item the orthogonal crossing-leakage/diagonal-mismatch identity
      \eqref{eq:V101-leakage-mismatch-identity};
\item the two-obstruction physical scalar
      \eqref{eq:V101-two-obstruction-scalar};
\item the exact \(F\otimes R_N\) aligned heat defect and its
      \(N^{-1}\) thermal asymptotic
      \eqref{eq:V101-exact-aligned-defect}--%
      \eqref{eq:V101-aligned-defect-limit};
\item the disconnected disposition of the trivial-core completion; and
\item the exact physical selector dichotomy
      \eqref{eq:V101-side-two-obstruction-lower}.
\end{enumerate}
\end{theorem}

\begin{proof}
These are
\Cref{thm:V101-exact-scalar-fusion,
lem:V101-column-orthogonal-split,
thm:V101-two-obstruction,
thm:V101-exact-aligned-heat,
prop:V101-trivial-core,
thm:V101-selector-dichotomy}.
\end{proof}

\begin{openproblem}[V102 mandate: acquire a complete side realization]
\label{op:V102-mandate}
Return to the V59/V60 incidence and puncture alphabets before any
representation asymptotic.  Determine whether a complete physical
simultaneous-sideways sequence can realize
\[
 F\otimes R_N\supset S_N=V_{(N-1,1)}
\]
on its active same-coordinate \(AB\) bank.
\begin{enumerate}
\item Record the complete collision/puncture monomial, original-row masses,
      and priority inequalities.
\item Prove that private-mean, cut-heavy, selected-heavy-pair, disconnected,
      and off-cut alternatives are absent before assigning the sequence to
      \(\mathfrak J_{\rm side}\).
\item If no such sequence exists, identify the forced paid or off-cut
      disposition and remove the test from the side programme.
\item If it exists, retain \(C_N,T,\beta,c_j\) and compute the two
      quantities \(\mathfrak l_{S_N,\rm side}^2\) and
      \(\mathfrak d_{S_N,\rm side}^2\).
\item Use the exact heat defect
      \eqref{eq:V101-aligned-defect-limit} only inside the diagonal term;
      do not use it to bound crossing leakage.
\item Keep the off-cut acquisition independent.
\end{enumerate}
If the printed selector does not contain enough data to decide existence,
state the missing physical datum explicitly rather than inventing a side
block.
\end{openproblem}

\section{Firewall after V101}
\label{sec:V101-firewall}

\begin{firewall}[Current claim boundary]
\label{fire:V101-current-boundary}
V101 does not prove that the V100 near-unit family occurs in the physical
side bank, estimate its crossing leakage, close either V98 residual branch,
prove JREC, aggregate collisions, construct the continuum Yang--Mills
theory, or prove a mass gap.  It proves that the aligned heat defect is
small and identifies the second, independent physical quantity which must
also be controlled.  The full Yang--Mills existence and mass gap remain
open.
\end{firewall}

\section{Append-only record for V101}
\label{sec:V101-record}

V101 was appended to the complete V100 manuscript.  The exact V100 parent
had \(68369\) physical lines, \(2393616\) bytes, and SHA--256 digest
\[
 \texttt{8d781731a7bf6f5a57ed0b20ff001fe9ce148e440740d79811df295efb8c6af6}.
\]
No V100 mathematical content was changed.  The sole final
\verb|\end{document}| is restored below.

% =============================================================================
% END APPENDED VERSION V101 MATERIAL
% =============================================================================
% =============================================================================
% BEGIN APPENDED VERSION V102 MATERIAL
% =============================================================================

\chapter{V102: Sideways Mass Localization and the Exact
Near-Unit Overlapping-Pair Gate}
\label{ch:V102}

\section{Purpose}
\label{sec:V102-purpose}

V101 proved that the near-unit \(F\otimes R_N\) channel can test the
physical side branch only after a complete V94 priority assignment.  This
chapter returns to the V59/V60 mass identities and extracts what every
such assignment must carry.

At a sideways endpoint, the selected collision edge contains only a small
fraction of the row mass.  The exact incidence partition leaves only
assigned cut mass, a private core mean, off-cut mass, or the residual
puncture bank.  The puncture alphabet then splits the last bank into a
singleton mean and genuine internal pairs.  Therefore, after the cut and
mean alternatives are excluded, a side block contains a quantitatively
heavy internal pair unless it belongs to the separate off-cut branch.
This conclusion is algebraic and remains valid when two endpoints are
simultaneously sideways; it does not assume that the resulting overlapping
pair capacities factor.

The second part inserts the V100 near-unit projector into the exact V88
pair-level/core-shell minimum.  The resulting formula shows that the
near-unit visibility and the asymmetric Schur dimensions produce no
automatic decay when the auxiliary pair legs have bounded dimension.
The remaining possible gain is the actual internal-pair level
distribution, retained in the same product-state expectation.

\section{A simultaneous-endpoint mass-localization lemma}
\label{sec:V102-mass-localization}

Let \(\varepsilon_{\rm sw}\) be the fixed V58 sideways threshold, so a
sideways endpoint \(v\) satisfies
\begin{equation}
 s_\alpha\Xi_v>M,
 \qquad
 x_v<\varepsilon_{\rm sw}\Xi_v.
 \label{eq:V102-sideways-threshold}
\end{equation}
Choose \(\eta>0\) so small that
\begin{equation}
 \delta_0:=1-\varepsilon_{\rm sw}-3\eta>0.
 \label{eq:V102-delta-zero}
\end{equation}

\begin{lemma}[Mass localization at the central endpoint]
\label{lem:V102-central-localization}
Suppose \(a\) is sideways.  Then at least one of the following holds:
\begin{enumerate}[label=\textup{(\roman*)}]
\item \(\widehat H_{ab}\ge a_*\eta\Xi_a\) or
      \(\widehat H_{ac}\ge a_*\eta\Xi_a\);
\item \(P_a^{\rm priv}\ge\eta\Xi_a\);
\item \(P_{a,\mu}^{\rm pun}\ge\delta_0\Xi_a/2\);
\item for some \(j\ne a\),
      \begin{equation}
       D_{aj,\mu}^{\rm pun}
       \ge\frac{\delta_0}{6}\Xi_a.
       \label{eq:V102-central-heavy-pair}
      \end{equation}
\end{enumerate}
\end{lemma}

\begin{proof}
Assume the first two alternatives fail.  The exact V59 identity is
\[
 \Xi_a
 =
 x_a+Z_{a\to b}+Z_{a\to c}+P_a^{\rm priv}+Y_a.
\]
By \eqref{eq:V59-cut-dominates-incidence-mass},
\[
 Z_{a\to r}
 \le a_*^{-1}\widehat H_{ar}
 <\eta\Xi_a,
 \qquad r=b,c.
\]
Use \(x_a<\varepsilon_{\rm sw}\Xi_a\) and
\(P_a^{\rm priv}<\eta\Xi_a\) to obtain
\[
 Y_a\ge
 (1-\varepsilon_{\rm sw}-3\eta)\Xi_a
 =
 \delta_0\Xi_a.
\]
The exact puncture alphabet
\eqref{eq:V60-puncture-mass-split} and its pigeonhole alternatives then
give either
\[
 P_{a,\mu}^{\rm pun}\ge\frac12Y_a
\]
or \(D_{aj,\mu}^{\rm pun}\ge Y_a/6\) for some \(j\ne a\).
These imply alternatives \textup{(iii)} and \textup{(iv)}.
\end{proof}

\begin{lemma}[Mass localization at an outer internal endpoint]
\label{lem:V102-outer-localization}
Let \(r\in\{b,c\}\) be sideways and active in the \(ar\) path.  Then at
least one of the following holds:
\begin{enumerate}[label=\textup{(\roman*)}]
\item \(\widehat H_{ar}\ge a_*\eta\Xi_r\);
\item \(P_r^{\rm priv}\ge\eta\Xi_r\);
\item \(D_r^{\rm off}\ge\eta\Xi_r\);
\item \(P_{r,\mu}^{\rm pun}\ge\delta_0\Xi_r/2\);
\item for some \(j\ne r\),
      \begin{equation}
       D_{rj,\mu}^{\rm pun}
       \ge\frac{\delta_0}{6}\Xi_r.
       \label{eq:V102-outer-heavy-pair}
      \end{equation}
\end{enumerate}
\end{lemma}

\begin{proof}
If the first three alternatives fail, use the exact identity
\[
 \Xi_r
 =
 x_r+Z_{r\leftarrow a}+P_r^{\rm priv}
 +D_r^{\rm off}+Y_r.
\]
The cut domination gives
\[
 Z_{r\leftarrow a}
 \le a_*^{-1}\widehat H_{ar}
 <\eta\Xi_r.
\]
Together with \eqref{eq:V102-sideways-threshold} and the failure of
alternatives \textup{(ii)}--\textup{(iii)}, this yields
\[
 Y_r\ge\delta_0\Xi_r.
\]
Apply the same puncture alphabet and pigeonhole bounds as in
\Cref{lem:V102-central-localization}.
\end{proof}

\begin{theorem}[Forced heavy pair or off-cut carrier]
\label{thm:V102-forced-pair-or-offcut}
Consider any fixed physical monomial with one or several sideways
endpoints.  At each such endpoint, after the cut-heavy and private/puncture
mean alternatives are removed by their established payments, one of the
following remains:
\begin{enumerate}[label=\textup{(\alph*)}]
\item a genuine internal puncture pair carries at least
      \(\delta_0\Xi_v/6\) of that endpoint mass; or
\item when \(v\in\{b,c\}\), the separate off-cut bank carries at least
      \(\eta\Xi_v\).
\end{enumerate}
The assertion is simultaneous: applying it at two endpoints uses the same
fixed monomial and does not multiply any pair estimates.
\end{theorem}

\begin{proof}
Apply \Cref{lem:V102-central-localization} or
\Cref{lem:V102-outer-localization} at every sideways endpoint.
Alternatives involving a large assigned cut are paid by V59 cut-heavy
promotion.  Private and puncture singleton means are paid by the V59/V60
mean mechanism.  What remains is precisely the internal-pair or off-cut
alternative.  Each application uses an exact row-mass identity in the
same monomial; no trace norm, product-state supremum, or bilinear-tax
factorization has yet been taken.
\end{proof}

\begin{firewall}[Mass localization is not joint capacity closure]
\label{fire:V102-localization-not-closure}
\Cref{thm:V102-forced-pair-or-offcut} proves that the missing mass is
represented by a physical carrier.  It does not prove that two heavy
internal pairs sharing a row have multiplicative formation taxes.
Such multiplication is forbidden by V86.  Every forced pair must remain
inside the same full-row product-state expectation as the ternary square.
\end{firewall}

\section{The exact near-unit overlapping-pair level gate}
\label{sec:V102-near-unit-pair-gate}

Let an internal-pair level projector \(P_e(t)\) act on auxiliary factors
\(H_X\otimes H_Y\) attached respectively to the two retained rows carrying
\[
 A=F,\qquad B=R_N.
\]
Let
\[
 M_N:=s_N^2P_{S_N}
\]
be a scalar positive shell on the sole non-Cartan core component
\(S_N=V_{(N-1,1)}\).  Denote the pair-level one-sided dimension function
by \(D_e(t)\), as in V87--V88.

The relevant \(SU(3)\) dimensions are
\begin{align}
 d_F&=3,
 &
 d_{R_N}&=\frac{(N+1)(N+2)}2,
 &
 d_{S_N}&=N(N+2).
 \label{eq:V102-near-unit-dimensions}
\end{align}
Moreover,
\begin{align}
 \Pi_N^{\rm core}
 &:=
 \|M_N\|_{\rm prod}
 =
 s_N^2\frac{N}{N+1},
 \label{eq:V102-near-unit-core-product}\\
 \mathcal H_N
 &:=
 \operatorname{Tr}M_N
 =
 s_N^2N(N+2).
 \label{eq:V102-near-unit-core-trace}
\end{align}

\begin{theorem}[Specialized pair-level/core minimum]
\label{thm:V102-near-unit-pair-level}
For every pair level \(t\),
\begin{equation}
 \boxed{\;
 \|P_e(t)\otimes M_N\|_{\rm prod}
 \le
 s_N^2
 \min\left\{
 \frac{N}{N+1},
 \frac{
 D_e(t)N(N+2)}
 {\max\left\{
 3d_X,\,
 \frac{(N+1)(N+2)}2d_Y
 \right\}}
 \right\}.
 \;}
 \label{eq:V102-near-unit-level-minimum}
\end{equation}
If \(Q_e=\int_0^{A_*}P_e(t)\,\dd t\), then
\begin{equation}
 \boxed{\;
 \|Q_e\otimes M_N\|_{\rm prod}
 \le s_N^2\Theta_{e,N},
 \;}
 \label{eq:V102-near-unit-integrated-gate}
\end{equation}
where
\begin{equation}
 \Theta_{e,N}
 :=
 \int_0^{A_*}
 \min\left\{
 \frac{N}{N+1},
 \frac{
 D_e(t)N(N+2)}
 {\max\left\{
 3d_X,\,
 \frac{(N+1)(N+2)}2d_Y
 \right\}}
 \right\}\,\dd t.
 \label{eq:V102-exact-pair-profile}
\end{equation}
The same formulas hold with the two row assignments exchanged after
exchanging \(d_X,3\) with \(d_Y,d_{R_N}\).
\end{theorem}

\begin{proof}
The dimension formula for \(SU(3)\) is
\[
 d_{(p,q)}
 =
 \frac{(p+1)(q+1)(p+q+2)}2.
\]
This gives \eqref{eq:V102-near-unit-dimensions}.
Equation \eqref{eq:V102-near-unit-core-product} is
\eqref{eq:V100-scaled-near-unit-capacity}; the trace of an isotypic
projector is its irrep dimension, proving
\eqref{eq:V102-near-unit-core-trace}.

Insert these two values and the retained dimensions
\(d_A=3\), \(d_B=(N+1)(N+2)/2\) into the exact V88 minimum
\eqref{eq:V88-refined-pair-shell}.  This proves
\eqref{eq:V102-near-unit-level-minimum}.  Integrate the operator layer
cake exactly as in \eqref{eq:V88-refined-integral} to obtain
\eqref{eq:V102-near-unit-integrated-gate}--%
\eqref{eq:V102-exact-pair-profile}.  Exchanging retained rows gives the
last statement.
\end{proof}

\begin{corollary}[No automatic dimension decay for bounded auxiliary legs]
\label{cor:V102-no-dimension-decay}
If \(d_X,d_Y\) remain fixed and positive as \(N\to\infty\), then
\begin{equation}
 \frac{N(N+2)}
 {\max\left\{
 3d_X,\,
 \frac{(N+1)(N+2)}2d_Y
 \right\}}
 \longrightarrow\frac2{d_Y}.
 \label{eq:V102-trace-ratio-limit}
\end{equation}
Thus neither the near-unit projector nor the retained Schur dimensions
provide an \(N\)-decay in
\eqref{eq:V102-near-unit-level-minimum}.  Any decay must come from
\(s_N^2\), the actual pair-level distribution \(D_e(t)\), or growth of a
same-row auxiliary dimension retained in the denominator.
\end{corollary}

\begin{proof}
For fixed positive \(d_X,d_Y\), the second entry in the maximum grows
quadratically and eventually dominates \(3d_X\).  Cancelling \(N+2\)
then gives
\[
 \frac{2N}{(N+1)d_Y}\longrightarrow\frac2{d_Y}.
\]
The first entry of the minimum tends to one.
\end{proof}

\begin{remark}[Why the un-factorized level profile is essential]
\label{rem:V102-unfactorized-essential}
The factorized V88 corollary would give
\[
 \|Q_e\otimes M_N\|_{\rm prod}
 \le
 \mathcal K_e\,
 \frac{s_N^2N(N+2)}3,
\]
because \(\min\{d_F,d_{R_N}\}=3\).  This can lose a quadratic factor.
The exact profile \(\Theta_{e,N}\) retains the competing near-unit cap
\(N/(N+1)\) at every pair level and is therefore the correct acquisition
object for this asymmetric family.
\end{remark}

\begin{assessment}[The isolated near-unit projector is not the residual carrier]
\label{ass:V102-complete-carrier}
\Cref{thm:V102-forced-pair-or-offcut} and
\Cref{thm:V102-near-unit-pair-level} fit together as follows.  If the
near-unit core channel occurs at a genuinely sideways endpoint and the
outcome is not off-cut, its complete physical carrier includes a heavy
internal pair.  The correct capacity is then the joint pair--core
quantity in \eqref{eq:V102-near-unit-integrated-gate}, not
\(\|P_{S_N}\|_{\rm prod}\) by itself.  The visibility limit one remains a
sharp warning, but it does not authorize deletion of the pair level which
the mass partition forces into the same monomial.
\end{assessment}

\section{Assembly conventions and sharp tests}
\label{sec:V102-assembly-tests}

\begin{remark}[No internal-pair double charge]
\label{rem:V102-no-pair-double-charge}
In \eqref{eq:V102-near-unit-integrated-gate}, \(M_N\) denotes the positive
core shell after the complete signed ternary cumulant has been squared and
after all weights disjoint from the chosen internal pair have been
included, but before \(Q_e\) is inserted.  The physical block is
\(Q_e\otimes M_N\) on disjoint coordinate banks regrouped by the same
original rows.  If one instead starts from the V101 scalar
\(\mathfrak s_{S,X}^2\) with the \(Q_e\) weight already inside \(c_j\),
the pair must not be inserted again.
\end{remark}

\begin{proposition}[Three exact dispositions]
\label{prop:V102-three-dispositions}
\begin{enumerate}[label=\textup{(\roman*)}]
\item If the pair effect is the identity, the joint product capacity is
      exactly the near-unit core capacity
      \(s_N^2N/(N+1)\); hence no theorem based only on
      \(0\le Q_e\le I\) can yield uniform decay.
\item If the auxiliary dimension \(d_Y\to\infty\), the trace entry in
      \eqref{eq:V102-near-unit-level-minimum} may gain
      \(d_Y^{-1}\), but this is same-row pair data and cannot be replaced by
      mass on a disjoint Cartan spectator.
\item The off-cut alternative in
      \Cref{thm:V102-forced-pair-or-offcut} is not estimated by
      \(\Theta_{e,N}\).  It remains the independent V98 off-cut carrier.
\end{enumerate}
\end{proposition}

\begin{proof}
For \textup{(i)}, tensoring with the identity does not change a
product-state norm, and use
\eqref{eq:V102-near-unit-core-product}.  Part \textup{(ii)} is the second
entry of \eqref{eq:V102-near-unit-level-minimum}; the same-coordinate
restriction is V94 spectator quotient and V93 spectator inflation.
Part \textup{(iii)} follows from the orthogonal V98 outer partition.
\end{proof}

\begin{theorem}[The exact acquisition fork after mass localization]
\label{thm:V102-acquisition-fork}
For a complete near-unit physical sequence with a sideways endpoint,
the unresolved analysis has the following disjoint order.
\begin{enumerate}[label=\textup{(\alph*)}]
\item If the V94 priority selector assigns the sequence to a paid cut or
      mean branch, remove it before a residual capacity is formed.
\item If it assigns the heavy mass to the off-cut bank, estimate only
      \(\mathbf Q_{\rm off}\).
\item Otherwise retain at least one forced internal pair and estimate the
      joint level profile \(\Theta_{e,N}\) together with the bare signed
      core scalar \(s_N^2\).
\end{enumerate}
In case \textup{(c)}, a sufficient local decay condition is
\begin{equation}
 s_N^2\Theta_{e,N}\longrightarrow0.
 \label{eq:V102-local-decay-condition}
\end{equation}
Neither \(\beta_{F,R_N}(S_N)\to1\) nor
\(\delta_{AB}^{\tau_N}(S_N)=O(N^{-1})\) decides
\eqref{eq:V102-local-decay-condition} without the other factor.
\end{theorem}

\begin{proof}
The branch order is the V94/V98 disjoint priority partition combined with
\Cref{thm:V102-forced-pair-or-offcut}.  In case \textup{(c)},
\eqref{eq:V102-near-unit-integrated-gate} proves the sufficient condition.
The V100 visibility tends to one, while V101 shows that the aligned defect
does not control crossing leakage.  Thus neither scalar statement alone
determines the product \(s_N^2\Theta_{e,N}\).
\end{proof}

\section{V102 ledger and V103 acquisition}
\label{sec:V102-ledger}

\begin{theorem}[Integrated V102 statement]
\label{thm:V102-integrated}
Preserving V1--V101, V102 proves:
\begin{enumerate}[label=\textup{(V102.\arabic*)},leftmargin=6.3em]
\item quantitative central and outer endpoint localization in
      \Cref{lem:V102-central-localization,
      lem:V102-outer-localization};
\item the simultaneous heavy-internal-pair/off-cut dichotomy
      \Cref{thm:V102-forced-pair-or-offcut};
\item the exact near-unit dimensions, product capacity, and trace
      \eqref{eq:V102-near-unit-dimensions}--%
      \eqref{eq:V102-near-unit-core-trace};
\item the specialized un-factorized pair-level/core minimum
      \eqref{eq:V102-near-unit-level-minimum};
\item the integrated pair profile
      \eqref{eq:V102-exact-pair-profile};
\item the bounded-auxiliary-dimension no-gain limit
      \eqref{eq:V102-trace-ratio-limit}; and
\item the disjoint acquisition fork and local decay gate
      \eqref{eq:V102-local-decay-condition}.
\end{enumerate}
\end{theorem}

\begin{proof}
The listed statements are
\Cref{thm:V102-forced-pair-or-offcut,
thm:V102-near-unit-pair-level,
cor:V102-no-dimension-decay,
thm:V102-acquisition-fork}.
\end{proof}

\begin{openproblem}[V103 mandate: compute the forced-pair level profile]
\label{op:V103-mandate}
Remain on one complete physical side monomial and one forced heavy internal
pair \(e\) from \Cref{thm:V102-forced-pair-or-offcut}.
\begin{enumerate}
\item Start from the exact pair defect spectrum and compute its one-sided
      level function \(D_e(t)\) on the selected-heavy branch.
\item Insert that function into \(\Theta_{e,N}\) before applying the
      factorized \(\mathcal K_e\) bound, which loses \(N^2\) on the
      asymmetric near-unit family.
\item Keep the complete signed ternary square in \(s_N^2\); use the V101
      aligned heat asymptotic only after controlling crossing leakage.
\item Determine whether one forced pair makes
      \(s_N^2\Theta_{e,N}\) reach the V60 path scale.
\item If two simultaneous endpoints force two pairs, formulate one common
      two-pair level law on the full rows.  Do not multiply overlapping
      one-pair capacities.
\item Test the estimate on the fixed-dimensional pair regime, the
      fundamental--dual singlet layer, and the V86 spectator-inflation
      regression.
\item Keep the off-cut branch separate.
\end{enumerate}
The next compulsory companion audit remains V105.
\end{openproblem}

\section{Firewall after V102}
\label{sec:V102-firewall}

\begin{firewall}[Current claim boundary]
\label{fire:V102-current-boundary}
V102 does not estimate the physical function \(D_e(t)\), control the V101
crossing leakage, prove \eqref{eq:V102-local-decay-condition} uniformly,
close the side or off-cut branch, prove JREC, aggregate all collisions,
construct continuum Yang--Mills theory, or prove a positive mass gap.
It proves that every unresolved sideways mass has a concrete pair/off-cut
carrier and gives the exact near-unit joint pair-level gate.  The full
Yang--Mills existence and mass gap remain open.
\end{firewall}

\section{Append-only record for V102}
\label{sec:V102-record}

V102 was appended to the complete V101 manuscript.  The exact V101 parent
had \(69024\) physical lines, \(2414352\) bytes, and SHA--256 digest
\[
 \texttt{253c6437a522199c131dddf10b12c219c91e6fa838c335c61ae4fd87d7293713}.
\]
No V101 mathematical content was changed.  The sole final
\verb|\end{document}| is restored below.

% =============================================================================
% END APPENDED VERSION V102 MATERIAL
% =============================================================================
% =============================================================================
% BEGIN APPENDED VERSION V103 MATERIAL
% =============================================================================

\chapter{V103: Fractional Ky--Fan Resolution of the Forced Pair and the
Exact \(s_\alpha N_{\rm eff}/N^2\) Threshold}
\label{ch:V103}

\section{Purpose}
\label{sec:V103-purpose}

V102 reduced a non-off-cut sideways occurrence of the near-unit
\(F\otimes R_N\) core to a joint internal-pair/core capacity.  Its exact
level profile \(\Theta_{e,N}\) should not be estimated from scratch:
V90 already proved that every un-factorized one-pair/core shell is governed
by a truncated pair capacity and gave an exact fractional Ky--Fan formula.

The new step is to compute the truncation parameter for the near-unit
core and then combine the result with the V53 selected-heavy cold/hot
decomposition.  The core product-to-Schur ratio is exactly
\[
 \eta_N=\frac{3}{(N+1)(N+2)}.
\]
Thus the joint capacity probes only the top
\(\eta_Nd_e\) eigenvalue mass of the forced internal-pair operator.  V53
controls the total normalized rank of its cold outputs and the pointwise
size of its hot outputs.  Combining those two facts gives an explicit
upper bound for this top fractional mean and isolates the scale
\[
 s_\alpha N_{\rm eff}\eta_N
 \asymp
 \frac{s_\alpha N_{\rm eff}}{N^2}.
\]

\section{Exact near-unit core ratio}
\label{sec:V103-core-ratio}

Retain the notation
\[
 M_N=s_N^2P_{S_N},
 \qquad
 S_N=V_{(N-1,1)}.
\]
For the V90 ratio form put
\[
 \Pi_N:=\|M_N\|_{\rm prod},
 \qquad
 B_N:=\frac{\operatorname{Tr}M_N}{\min\{d_F,d_{R_N}\}}.
\]

\begin{lemma}[Exact \(N^{-2}\) truncation ratio]
\label{lem:V103-near-unit-ratio}
For \(s_N>0\),
\begin{align}
 \Pi_N
 &=
 s_N^2\frac{N}{N+1},
 \label{eq:V103-Pi-N}\\
 B_N
 &=
 s_N^2\frac{N(N+2)}3,
 \label{eq:V103-B-N}\\
 \boxed{\;
 \eta_N:=\frac{\Pi_N}{B_N}
 =
 \frac{3}{(N+1)(N+2)}.
 \;}
 \label{eq:V103-eta-N}
\end{align}
The same value is assigned by continuity when \(s_N=0\).
\end{lemma}

\begin{proof}
Equations \eqref{eq:V102-near-unit-core-product} and
\eqref{eq:V102-near-unit-core-trace} give the first formula and
\[
 \operatorname{Tr}M_N=s_N^2N(N+2).
\]
Since
\(\min\{d_F,d_{R_N}\}=3\), this proves
\eqref{eq:V103-B-N}.  Division gives
\eqref{eq:V103-eta-N}.
\end{proof}

\section{The exact fractional Ky--Fan pair gate}
\label{sec:V103-fractional-KyFan}

List the eigenvalues of the forced internal-pair operator \(Q_e\), with
multiplicity, as
\[
 c_1\ge c_2\ge\cdots\ge c_L\ge0,
 \qquad L=d_Xd_Y,
\]
and put
\[
 d_e:=\max\{d_X,d_Y\},
 \qquad
 q_N:=\eta_Nd_e,
 \qquad
 n_N:=\lfloor q_N\rfloor.
\]
When \(q_N>0\), define the top fractional Ky--Fan mean
\begin{equation}
 \overline c_{e,N}
 :=
 \frac1{q_N}
 \left[
 \sum_{\ell=1}^{n_N}c_\ell
 +(q_N-n_N)c_{n_N+1}
 \right],
 \label{eq:V103-top-fractional-mean}
\end{equation}
with the usual omission of the last term when \(q_N=n_N\).

\begin{theorem}[Exact fractional Ky--Fan reduction]
\label{thm:V103-fractional-KyFan}
The V90 un-factorized ratio bound specializes to
\begin{align}
 \|Q_e\otimes M_N\|_{\rm prod}
 &\le
 \frac{s_N^2N(N+2)}{3d_e}
 \left[
 \sum_{\ell=1}^{n_N}c_\ell
 +(q_N-n_N)c_{n_N+1}
 \right]
 \label{eq:V103-explicit-KyFan-gate}\\
 &=
 \boxed{\;
 s_N^2\frac{N}{N+1}\,
 \overline c_{e,N}.
 \;}
 \label{eq:V103-normalized-KyFan-gate}
\end{align}
If \(q_N<1\), then
\begin{equation}
 \overline c_{e,N}=c_1=\|Q_e\|_{\rm op},
 \qquad
 \|Q_e\otimes M_N\|_{\rm prod}
 \le
 s_N^2\frac{N}{N+1}\|Q_e\|_{\rm op}.
 \label{eq:V103-subrank-one-gate}
\end{equation}
\end{theorem}

\begin{proof}
The V90 ratio theorem gives
\[
 \|Q_e\otimes M_N\|_{\rm prod}
 \le
 B_N\mathcal K_e(\eta_N).
\]
Its exact Ky--Fan formula is
\[
 \mathcal K_e(\eta_N)
 =
 \frac1{d_e}
 \left[
 \sum_{\ell=1}^{n_N}c_\ell
 +(q_N-n_N)c_{n_N+1}
 \right].
\]
Insert \eqref{eq:V103-B-N} to obtain
\eqref{eq:V103-explicit-KyFan-gate}.  Since
\[
 B_N\eta_N=\Pi_N=s_N^2\frac{N}{N+1},
\]
definition \eqref{eq:V103-top-fractional-mean} gives
\eqref{eq:V103-normalized-KyFan-gate}.  If \(q_N<1\), then \(n_N=0\)
and the fractional mean is \(q_Nc_1/q_N=c_1\).
\end{proof}

\begin{corollary}[Bounded pair dimensions see only the top eigenvalue]
\label{cor:V103-bounded-pair-top}
If \(d_e=o(N^2)\), then \(q_N<1\) for all sufficiently large \(N\).
In particular, bounded auxiliary pair dimensions yield no spectral
averaging: the entire V90 truncation reduces to the top pair coefficient
as in \eqref{eq:V103-subrank-one-gate}.
\end{corollary}

\begin{proof}
Equation \eqref{eq:V103-eta-N} gives
\[
 q_N
 =
 \frac{3d_e}{(N+1)(N+2)}.
\]
This tends to zero when \(d_e=o(N^2)\).
\end{proof}

\begin{assessment}[What must replace the factorized pair capacity]
\label{ass:V103-top-spectrum}
The selected-heavy theorem controls the full capped average
\(\mathcal K_e(1)\).  The near-unit core probes instead the top
\(\eta_N\)-fractional mean.  A small average does not by itself make this
top mean small.  The missing datum is therefore not the total pair
capacity but concentration of its largest weighted defect eigenvalues.
\end{assessment}

\section{Cold rank plus hot coefficient control}
\label{sec:V103-cold-hot}

Assume the forced pair is selected-heavy at an endpoint, with input mass
\(x\ge\varepsilon\Xi_v\), effective cancellation number
\[
 N_{\rm eff}:=\frac{x^2}{h},
\]
and \(s_\alpha N_{\rm eff}>M\).  Use the fixed V53 cold family
\[
 \Omega_{\rm cold}
 =
 \{(S,\mu):\Xi(S)<\kappa N_{\rm eff}\}.
\]
Let \(D_{\rm cold}\) be its total output dimension with multiplicity and
put
\[
 R:=s_\alpha N_{\rm eff}.
\]
The V53 theorem and the V87 hot-output proof give, for a group constant
\(C\),
\begin{align}
 \min\left\{1,\frac{D_{\rm cold}}{d_e}\right\}
 &\le\frac C R,
 \label{eq:V103-cold-rank-fraction}\\
 c_\gamma&\le\frac C R
 \qquad(\gamma\notin\Omega_{\rm cold}).
 \label{eq:V103-hot-pointwise-cap}
\end{align}
All eigenvalues obey \(0\le c_\gamma\le A_*\).

\begin{lemma}[Top fractional mean from a cold/hot split]
\label{lem:V103-fractional-cold-hot}
Let a positive operator have eigenvalues in \([0,A_*]\).  Suppose a
distinguished cold spectral subspace has rank \(D_{\rm cold}\), while every
eigenvalue outside it is at most \(b\).  For \(0<\eta\le1\), its top
\(\eta d_e\)-fractional mean obeys
\begin{equation}
 \boxed{\;
 \overline c(\eta)
 \le
 b+A_*
 \min\left\{
 1,\frac{D_{\rm cold}}{\eta d_e}
 \right\}.
 \;}
 \label{eq:V103-abstract-cold-hot}
\end{equation}
\end{lemma}

\begin{proof}
Put \(q=\eta d_e\).  In the fractional Ky--Fan sum of spectral mass \(q\),
the hot part is at most \(bq\).  The cold part uses spectral mass at most
\(\min\{q,D_{\rm cold}\}\), and every cold eigenvalue is at most \(A_*\).
Thus the numerator is at most
\[
 bq+A_*\min\{q,D_{\rm cold}\}.
\]
Divide by \(q\).
\end{proof}

\begin{theorem}[Exact effective-number threshold bound]
\label{thm:V103-effective-threshold}
For all sufficiently large \(R=s_\alpha N_{\rm eff}\), the forced-pair
fractional mean on the near-unit core satisfies
\begin{equation}
 \boxed{\;
 \overline c_{e,N}
 \le
 \frac C R
 A_*\min\left\{
 1,\frac{C}{\eta_NR}
 \right\},
 \qquad
 \eta_N=\frac3{(N+1)(N+2)}.
 \;}
 \label{eq:V103-fractional-effective-bound}
\end{equation}
Consequently
\begin{equation}
 \boxed{\;
 \|Q_e\otimes M_N\|_{\rm prod}
 \le
 s_N^2\frac{N}{N+1}
 \left[
 \frac C{s_\alpha N_{\rm eff}}
 +A_*\min\left\{
 1,
 \frac{C(N+1)(N+2)}
 {3s_\alpha N_{\rm eff}}
 \right\}
 \right].
 \;}
 \label{eq:V103-near-unit-effective-gate}
\end{equation}
In particular, if
\begin{equation}
 \frac{s_\alpha N_{\rm eff}}{N^2}\longrightarrow\infty
 \label{eq:V103-superquadratic-effective}
\end{equation}
and \(\sup_Ns_N^2<\infty\), then
\begin{equation}
 \|Q_e\otimes M_N\|_{\rm prod}\longrightarrow0.
 \label{eq:V103-superquadratic-decay}
\end{equation}
\end{theorem}

\begin{proof}
Once \(C/R<1\), equation
\eqref{eq:V103-cold-rank-fraction} implies
\[
 \frac{D_{\rm cold}}{d_e}\le\frac C R.
\]
Apply \Cref{lem:V103-fractional-cold-hot} with
\[
 b=\frac C R,
 \qquad
 \eta=\eta_N.
\]
This proves \eqref{eq:V103-fractional-effective-bound}.  Insert it into
\eqref{eq:V103-normalized-KyFan-gate} to obtain
\eqref{eq:V103-near-unit-effective-gate}.
Condition \eqref{eq:V103-superquadratic-effective} implies both
\[
 R\longrightarrow\infty,
 \qquad
 \eta_NR\longrightarrow\infty.
\]
The bracket in \eqref{eq:V103-near-unit-effective-gate} therefore tends
to zero.  Boundedness of \(s_N^2\) proves
\eqref{eq:V103-superquadratic-decay}.
\end{proof}

\begin{corollary}[Quantitative bound with a growing core scalar]
\label{cor:V103-growing-core-scalar}
Without assuming \(s_N^2\) bounded, the same route closes this one-pair
near-unit channel whenever
\begin{equation}
 s_N^2
 \left[
 \frac1{s_\alpha N_{\rm eff}}
 +
 \min\left\{
 1,\frac{N^2}{s_\alpha N_{\rm eff}}
 \right\}
 \right]
 \longrightarrow0.
 \label{eq:V103-combined-sufficient-scale}
\end{equation}
\end{corollary}

\begin{proof}
Absorb the fixed constants and the ratios
\((N+1)(N+2)/N^2\) into
\eqref{eq:V103-near-unit-effective-gate}.
\end{proof}

\begin{firewall}[The threshold is sharp only for the available information]
\label{fire:V103-information-sharp}
If \(\eta_NR\) stays bounded, the V53 data
\eqref{eq:V103-cold-rank-fraction}--%
\eqref{eq:V103-hot-pointwise-cap} alone allow the cold spectral rank to
fill the entire top \(\eta_N\)-fraction with eigenvalue comparable to
\(A_*\).  The fractional mean then need not decay.  This is a sharpness
statement for the cold-rank/hot-cap information, not a construction of a
physical Wilson pair with that extremal spectrum.  Additional ordering,
fusion, or correlation information could still improve the bound.
\end{firewall}

\begin{proposition}[A small full pair average need not control the top fraction]
\label{prop:V103-average-top-no-go}
There are positive spectra for which
\[
 \mathcal K_e(1)\longrightarrow0
\]
while \(\overline c_{e,N}=1\).  Hence the factorized selected-heavy
capacity cannot replace the fractional profile in the near-unit gate.
\end{proposition}

\begin{proof}
Let \(Q_e\) have one eigenvalue equal to one and all other eigenvalues zero,
and take \(d_e\to\infty\).  The V90 Ky--Fan formula gives
\[
 \mathcal K_e(1)=\frac1{d_e}\longrightarrow0.
\]
If \(\eta_Nd_e<1\), its top fractional mean is the largest eigenvalue,
which equals one, by \eqref{eq:V103-subrank-one-gate}.  For example, choose
\(d_e=N\).  Again, this is a spectral information test, not a claim that
every such spectrum is realized by the Wilson pair operator.
\end{proof}

\section{Mass-scale consequences and regressions}
\label{sec:V103-consequences-tests}

\begin{lemma}[Effective number never exceeds selected input mass]
\label{lem:V103-effective-below-mass}
For the selected pair,
\begin{equation}
 N_{\rm eff}\le x.
 \label{eq:V103-effective-below-mass}
\end{equation}
Consequently, if \(x_N=O(N^2)\), then
\[
 \frac{s_\alpha N_{\rm eff}}{N^2}=O(s_\alpha)
\]
and the superquadratic condition
\eqref{eq:V103-superquadratic-effective} cannot hold when
\(s_\alpha\) stays bounded.
\end{lemma}

\begin{proof}
In the V53 notation,
\[
 h=\sum_e r_eq_e,
 \qquad
 x=\sum_er_e,
\]
and every active mass \(q_e\ge1\).  Hence \(h\ge x\), so
\[
 N_{\rm eff}=\frac{x^2}{h}\le x.
\]
The last assertion follows immediately.
\end{proof}

\begin{assessment}[Natural mass matching returns the obstruction]
\label{ass:V103-mass-matched}
If the forced pair carries only order \(N^2\) input mass, comparable to the
Casimir size of \(R_N\), V53 cold-rank/hot-cap information cannot reach the
superquadratic threshold.  One forced-pair average then does not close the
near-unit family through \eqref{eq:V103-near-unit-effective-gate}.  The
remaining possibilities are a direct top-spectrum estimate, correlation
between the pair spectrum and the signed core column, or a common-level
estimate using the pair forced at the second sideways endpoint.
\end{assessment}

\begin{proposition}[Dual-singlet pair regression]
\label{prop:V103-dual-singlet-pair}
For an internal \(F\otimes F^*\) pair restricted to its singlet output,
the pair operator has one nonzero eigenvalue
\begin{equation}
 c_{\mathbf1}
 =
 [1+s_\alpha\Xi(\mathbf1)]
 |1-q_{\alpha,F}^2|.
 \label{eq:V103-dual-singlet-coefficient}
\end{equation}
Whenever \(q_N<1\), its contribution to the V103 gate is bounded by
\begin{equation}
 s_N^2\frac{N}{N+1}c_{\mathbf1}.
 \label{eq:V103-dual-singlet-near-unit}
\end{equation}
For the heat law with \(s_\alpha=\tau\),
\[
 c_{\mathbf1}
 =
 (1+\tau\Xi(\mathbf1))
 (1-e^{-8\tau/3})
 =
 O(\tau)
\]
as \(\tau\downarrow0\).
\end{proposition}

\begin{proof}
The singlet multiplier is one and
\(q_{\tau,F}=e^{-4\tau/3}\).  Definition
\eqref{eq:V87-internal-pair-operator} gives
\eqref{eq:V103-dual-singlet-coefficient}.  A rank-one singlet supplies
the sole top eigenvalue.  Apply
\eqref{eq:V103-subrank-one-gate}.  The heat asymptotic follows from
\(1-e^{-8\tau/3}=O(\tau)\).
\end{proof}

\begin{proposition}[Spectator inflation leaves the fractional problem unchanged]
\label{prop:V103-spectator-invariance}
Adjoining a disjoint Cartan spectator to make an endpoint mass large does
not change the active pair eigenvalues \(c_\ell\), their level ranks, or
the near-unit core ratio \(\eta_N\).  After the V94 spectator quotient it
therefore does not improve \(\overline c_{e,N}\).
\end{proposition}

\begin{proof}
V86 spectator invariance leaves the complete one-pair layer function and
weighted defect spectrum unchanged.  V94 removes the spectator identity
from the active product capacity.  The ratio \(\eta_N\) depends only on
the active \(F\otimes R_N\) core dimensions and product visibility, so it
is unchanged as well.
\end{proof}

\section{V103 ledger and V104 acquisition}
\label{sec:V103-ledger}

\begin{theorem}[Integrated V103 statement]
\label{thm:V103-integrated}
Preserving V1--V102, V103 proves:
\begin{enumerate}[label=\textup{(V103.\arabic*)},leftmargin=6.3em]
\item the exact near-unit core ratio
      \(\eta_N=3/[(N+1)(N+2)]\);
\item the exact fractional Ky--Fan gate
      \eqref{eq:V103-explicit-KyFan-gate}--%
      \eqref{eq:V103-normalized-KyFan-gate};
\item the top-eigenvalue reduction when \(d_e=o(N^2)\);
\item the cold-rank/hot-coefficient fractional bound
      \eqref{eq:V103-fractional-effective-bound};
\item the effective-number gate
      \eqref{eq:V103-near-unit-effective-gate};
\item sufficient decay at the superquadratic threshold
      \(s_\alpha N_{\rm eff}/N^2\to\infty\);
\item the information-level obstruction below that threshold; and
\item the mass-matched, dual-singlet, and spectator regressions.
\end{enumerate}
\end{theorem}

\begin{proof}
The items are
\Cref{lem:V103-near-unit-ratio,
thm:V103-fractional-KyFan,
cor:V103-bounded-pair-top,
thm:V103-effective-threshold,
fire:V103-information-sharp,
lem:V103-effective-below-mass,
prop:V103-dual-singlet-pair,
prop:V103-spectator-invariance}.
\end{proof}

\begin{openproblem}[V104 mandate: common two-pair fractional law]
\label{op:V104-mandate}
For a complete monomial with two simultaneously sideways endpoints, retain
the two forced internal pairs \(e,f\) in one original-row product
expectation.
\begin{enumerate}
\item Resolve \(Q_e\) and \(Q_f\) by two level variables \(t,s\).
\item Tensor both level projectors with one positive shell of the complete
      signed core square before applying any Schur estimate.
\item Compute the exact combined retained-row dimensions and the trace
      \(D_e(t)D_f(s)\mathcal H_N\), allowing the two pair edges to share an
      original row.
\item Take one minimum with the core product cap and integrate over
      \((t,s)\); do not multiply \(\mathcal K_e\mathcal K_f\).
\item Specialize the resulting truncation surface to the
      \(F\otimes R_N\) near-unit core.
\item Test whether the two V53 cold-rank/hot-cap decompositions together
      lower the threshold from \(s_\alpha N_{\rm eff}\gg N^2\).
\item Retain pair/core spectral correlation if the product of separate
      rank bounds returns the continuation-scale stock.
\item Keep off-cut outcomes outside this two-pair law.
\end{enumerate}
The next compulsory companion audit is V105.
\end{openproblem}

\section{Firewall after V103}
\label{sec:V103-firewall}

\begin{firewall}[Current claim boundary]
\label{fire:V103-current-boundary}
V103 does not prove the superquadratic effective-number condition for
physical side blocks, bound the top fractional pair spectrum in the
mass-matched regime, control the complete signed core scalar, close the
side or off-cut branch, prove JREC, construct the continuum theory, or
prove a positive mass gap.  It computes exactly what fraction of the
forced-pair spectrum the near-unit core sees and proves the strongest
consequence of the existing V53 cold/hot information.  The full
Yang--Mills existence and mass gap remain open.
\end{firewall}

\section{Append-only record for V103}
\label{sec:V103-record}

V103 was appended to the complete V102 manuscript.  The exact V102 parent
had \(69530\) physical lines, \(2432402\) bytes, and SHA--256 digest
\[
 \texttt{b4a616578b653a3b8a97c3da94e2a1bacb94b583b2637364f36b2871d4b8303b}.
\]
No V102 mathematical content was changed.  The sole final
\verb|\end{document}| is restored below.

% =============================================================================
% END APPENDED VERSION V103 MATERIAL
% =============================================================================
% =============================================================================
% BEGIN APPENDED VERSION V104 MATERIAL
% =============================================================================

\chapter{V104: Common Two-Pair Fractional Law and the Linear Threshold}
\label{ch:V104}

\section{Purpose and claim boundary}
\label{sec:V104-purpose}

V103 showed that the \(F\otimes R_N\) near-Cartan shell sees only the top
\(\eta_N\)-fraction of one forced-pair spectrum, where
\[
 \eta_N=\frac{3}{(N+1)(N+2)}\asymp N^{-2}.
\]
One selected-heavy pair therefore closes from the V53 cold/hot data only
at an effective scale much larger than \(N^2\).  A complete sideways
monomial, however, has a forced internal pair at each of its two retained
endpoints.  The present note keeps those two pairs inside one product
expectation and takes one Schur minimum only after both level projectors
have been inserted.

The result has two parts.  First, there is an exact two-level trace bound
which remains valid when the pair edges share the same two original
rows.  Second, after the two pair operators are regarded as one positive
auxiliary operator, the exceptional large spectrum is the
\emph{cold--cold} rectangle.  Its rank is a product, not a union.  Under
bounded orientation loss this lowers the conditional symmetric
effective-size threshold from \(N^2\) to \(N\).  This is a genuine
improvement, but it is conditional on the complete physical term
retaining a tensor-separated positive factor
\(Q_e\otimes Q_f\) next to one positive core shell.  No factorization of
two already-maximized capacities is used.

\section{Two retained pair banks on the same original rows}
\label{sec:V104-two-pair-setup}

Let \(G_e\) and \(G_f\) denote the compact gauge groups carried by two
disjoint internal coordinate banks.  Their irreducible row
representations are
\[
 X_e,\ Y_e
 \quad\hbox{and}\quad
 X_f,\ Y_f,
\]
with dimensions \(d_{X_e},d_{Y_e},d_{X_f},d_{Y_f}\).  Let \(A,B\) be
irreducible representations of the core gauge bank.  The two retained
original rows are
\begin{equation}
 \mathcal H_{\mathsf a}
 =
 \mathcal H_{X_e}\otimes\mathcal H_{X_f}
 \otimes\mathcal H_A,
 \qquad
 \mathcal H_{\mathsf b}
 =
 \mathcal H_{Y_e}\otimes\mathcal H_{Y_f}
 \otimes\mathcal H_B.
 \label{eq:V104-retained-rows}
\end{equation}
Thus the pair edges may share an original row, but their gauge-coordinate
banks remain distinct.

Let \(Q_e,Q_f\) be the positive invariant pair operators obtained after
the V87 shell weighting and V94 spectator quotient.  Write their layer
resolutions as
\begin{equation}
 Q_e=\int_0^{A_e}P_e(t)\,\dd t,
 \qquad
 Q_f=\int_0^{A_f}P_f(s)\,\dd s,
 \label{eq:V104-layer-resolutions}
\end{equation}
where
\[
 P_e(t)=\mathbf1_{\{Q_e>t\}},
 \quad
 P_f(s)=\mathbf1_{\{Q_f>s\}},
 \quad
 D_e(t)=\operatorname{rank}P_e(t),
 \quad
 D_f(s)=\operatorname{rank}P_f(s).
\]
The finite upper limits may be replaced by \(+\infty\); the displayed
form records the uniform coefficient caps available in the selected
heavy regime.  Let \(M\ge0\) be an invariant core shell on
\(\mathcal H_A\otimes\mathcal H_B\), and put
\begin{equation}
 \Pi(M):=\|M\|_{\rm prod},
 \qquad
 \mathcal H(M):=\operatorname{Tr}M.
 \label{eq:V104-core-data}
\end{equation}

\begin{lemma}[Product cap survives two auxiliary projections]
\label{lem:V104-two-projection-product-cap}
For every \(t,s\),
\begin{equation}
 \|P_e(t)\otimes P_f(s)\otimes M\|_{\rm prod}
 \le \Pi(M).
 \label{eq:V104-two-projection-product-cap}
\end{equation}
\end{lemma}

\begin{proof}
Since \(0\le P_e(t)\otimes P_f(s)\le I\), it is enough to prove the
claim with the auxiliary identity.  Let \(u\in\mathcal H_{\mathsf a}\)
and \(v\in\mathcal H_{\mathsf b}\) be unit vectors.  Tracing their pure
states over the auxiliary \(e,f\) factors gives density matrices
\(\rho_A,\rho_B\).  Therefore
\[
 \langle u\otimes v,(I\otimes M)u\otimes v\rangle
 =
 \operatorname{Tr}\bigl[M(\rho_A\otimes\rho_B)\bigr].
\]
Resolve each density matrix into convex combinations of pure states.
The last expression is a convex combination of product-vector
expectations of \(M\), and is at most \(\Pi(M)\).
\end{proof}

\begin{lemma}[Double Schur trace cap]
\label{lem:V104-double-Schur}
For every \(t,s\),
\begin{align}
 &\|P_e(t)\otimes P_f(s)\otimes M\|_{\rm prod}
 \notag\\
 &\quad\le
 \frac{D_e(t)D_f(s)\mathcal H(M)}
 {\max\{
 d_{X_e}d_{X_f}d_A,\,
 d_{Y_e}d_{Y_f}d_B
 \}}.
 \label{eq:V104-double-Schur}
\end{align}
\end{lemma}

\begin{proof}
Set
\[
 T_{t,s}:=P_e(t)\otimes P_f(s)\otimes M.
\]
It is positive and invariant under the product of the three compact
gauge groups.  The row representations in
\eqref{eq:V104-retained-rows} are external tensor products of
irreducibles and hence irreducible for that product group.  Schur's
lemma therefore gives
\begin{align*}
 \operatorname{Tr}_{\mathsf b}T_{t,s}
 &=
 \frac{\operatorname{Tr}T_{t,s}}
 {d_{X_e}d_{X_f}d_A}I_{\mathsf a},\\
 \operatorname{Tr}_{\mathsf a}T_{t,s}
 &=
 \frac{\operatorname{Tr}T_{t,s}}
 {d_{Y_e}d_{Y_f}d_B}I_{\mathsf b}.
\end{align*}
For a unit product vector \(u\otimes v\), positivity implies
\[
 \langle u\otimes v,T_{t,s}u\otimes v\rangle
 \le
 \min\left\{
 \langle u,\operatorname{Tr}_{\mathsf b}T_{t,s}u\rangle,\,
 \langle v,\operatorname{Tr}_{\mathsf a}T_{t,s}v\rangle
 \right\}.
\]
Finally,
\[
 \operatorname{Tr}T_{t,s}
 =D_e(t)D_f(s)\mathcal H(M).
\]
Taking the supremum over \(u,v\) proves the formula.
\end{proof}

\begin{theorem}[One-minimum two-level law]
\label{thm:V104-one-minimum-two-level}
Under the preceding hypotheses,
\begin{align}
 \|Q_e\otimes Q_f\otimes M\|_{\rm prod}
 \le
 \int_0^{A_e}\!\!\int_0^{A_f}
 \min\left\{
 \Pi(M),\,
 \frac{D_e(t)D_f(s)\mathcal H(M)}
 {\max\{
 d_{X_e}d_{X_f}d_A,\,
 d_{Y_e}d_{Y_f}d_B
 \}}
 \right\}\dd s\,\dd t.
 \label{eq:V104-one-minimum-two-level}
\end{align}
\end{theorem}

\begin{proof}
Tensor the two identities in
\eqref{eq:V104-layer-resolutions} and use finite-dimensional Bochner
integration:
\[
 Q_e\otimes Q_f\otimes M
 =
 \int_0^{A_e}\!\!\int_0^{A_f}
 P_e(t)\otimes P_f(s)\otimes M\,\dd s\,\dd t.
\]
The product numerical radius is subadditive on positive operators, so it
is at most the integral of the product radii.  At each common level
\((t,s)\), take the minimum of
\Cref{lem:V104-two-projection-product-cap,lem:V104-double-Schur}.
This proves \eqref{eq:V104-one-minimum-two-level}.  In particular, no
separate maximization of the \(e\)- and \(f\)-pair factors has occurred.
\end{proof}

\begin{firewall}[What would be an illegal multiplication]
\label{fire:V104-no-capacity-product}
The theorem does not assert
\[
 \|Q_e\otimes Q_f\otimes M\|_{\rm prod}
 \le
 \mathcal K_e(1)\mathcal K_f(1)\Pi(M).
\]
Such a statement would multiply capacities after optimizing over
possibly incompatible product vectors.  Formula
\eqref{eq:V104-one-minimum-two-level} instead retains the same product
vector, the same core shell, and the same Schur minimum for every
\((t,s)\).
\end{firewall}

\section{Near-unit specialization and orientation loss}
\label{sec:V104-near-unit}

Take the V102--V103 core
\[
 M_N=s_N^2P_{S_N},
 \qquad
 S_N=V_{(N-1,1)}\subset F\otimes R_N.
\]
Recall
\[
 d_F=3,\qquad
 d_{R_N}=\frac{(N+1)(N+2)}2,\qquad
 d_{S_N}=N(N+2),
\]
and
\[
 \Pi_N=s_N^2\frac N{N+1},
 \qquad
 \mathcal H_N=s_N^2N(N+2).
\]

\begin{corollary}[Exact two-level near-unit surface]
\label{cor:V104-near-unit-surface}
For the near-unit core,
\begin{align}
 &\|Q_e\otimes Q_f\otimes M_N\|_{\rm prod}
 \notag\\
 &\quad\le
 s_N^2\int_0^{A_e}\!\!\int_0^{A_f}
 \min\left\{
 \frac N{N+1},\,
 \frac{D_e(t)D_f(s)N(N+2)}
 {\max\{
 3d_{X_e}d_{X_f},\,
 d_{R_N}d_{Y_e}d_{Y_f}
 \}}
 \right\}\dd s\,\dd t.
 \label{eq:V104-near-unit-surface}
\end{align}
\end{corollary}

\begin{proof}
Substitute the displayed near-unit data into
\eqref{eq:V104-one-minimum-two-level}.
\end{proof}

\begin{definition}[Coherent auxiliary-row dimension]
\label{def:V104-coherent-row-dimension}
Put
\begin{align}
 d_e&:=\max\{d_{X_e},d_{Y_e}\},
 &
 d_f&:=\max\{d_{X_f},d_{Y_f}\},
 \notag\\
 d_{ef}&:=
 \max\{d_{X_e}d_{X_f},d_{Y_e}d_{Y_f}\},
 &
 \chi_{ef}&:=\frac{d_ed_f}{d_{ef}}.
 \label{eq:V104-orientation-factor}
\end{align}
The number \(d_{ef}\) is the larger \emph{actual} auxiliary row
dimension.  The factor \(\chi_{ef}\ge1\) measures the loss incurred by
multiplying the two separate larger-row dimensions.
\end{definition}

\begin{lemma}[Orientation comparison]
\label{lem:V104-orientation-comparison}
For arbitrary positive dimensions,
\begin{equation}
 1\le\chi_{ef}
 \le
 \min\{d_e,d_f\}.
 \label{eq:V104-orientation-range}
\end{equation}
If the larger dimension of both pairs lies on the same original row,
then \(\chi_{ef}=1\).  If
\[
 (d_{X_e},d_{Y_e})=(L,1),
 \qquad
 (d_{X_f},d_{Y_f})=(1,L),
\]
then \(d_{ef}=L\) and \(\chi_{ef}=L\).
\end{lemma}

\begin{proof}
 Both products entering \(d_{ef}\) are at most \(d_ed_f\), proving the
 lower bound.  Every representation dimension is at least one.  If the
 two separate maxima occur on the same row, then \(d_{ef}=d_ed_f\).  If
 they occur on opposite rows, one product entering \(d_{ef}\) is at
 least \(d_e\), and the other is at least \(d_f\).  Hence in all cases
\[
 d_{ef}\ge\max\{d_e,d_f\}.
\]
Therefore
\[
 \chi_{ef}=\frac{d_ed_f}{d_{ef}}
 \le\min\{d_e,d_f\}.
\]
The two asserted special cases follow directly from the definition.
\end{proof}

\begin{proposition}[Ratio relaxation with the coherent row]
\label{prop:V104-coherent-ratio-relaxation}
Define
\begin{equation}
 B(M):=\frac{\mathcal H(M)}{\min\{d_A,d_B\}},
 \qquad
 \eta(M):=\frac{\Pi(M)}{B(M)}.
 \label{eq:V104-general-ratio}
\end{equation}
Then
\begin{align}
 \|Q_e\otimes Q_f\otimes M\|_{\rm prod}
 \le
 B(M)\int_0^{A_e}\!\!\int_0^{A_f}
 \min\left\{
 \eta(M),\frac{D_e(t)D_f(s)}{d_{ef}}
 \right\}\dd s\,\dd t.
 \label{eq:V104-coherent-ratio-bound}
\end{align}
Equivalently, if
\[
 R_e(t):=\frac{D_e(t)}{d_e},
 \qquad
 R_f(s):=\frac{D_f(s)}{d_f},
\]
then the second entry of the minimum is exactly
\begin{equation}
 \frac{D_e(t)D_f(s)}{d_{ef}}
 =
 \chi_{ef}R_e(t)R_f(s).
 \label{eq:V104-oriented-rank-product}
\end{equation}
\end{proposition}

\begin{proof}
The denominator in \eqref{eq:V104-one-minimum-two-level} satisfies
\begin{align*}
 &\max\{
 d_{X_e}d_{X_f}d_A,\,
 d_{Y_e}d_{Y_f}d_B
 \}\\
 &\qquad\ge
 \min\{d_A,d_B\}
 \max\{d_{X_e}d_{X_f},d_{Y_e}d_{Y_f}\}\\
 &\qquad=
 \min\{d_A,d_B\}d_{ef}.
\end{align*}
Use this lower bound in the trace entry, factor out \(B(M)\), and invoke
\(\Pi(M)=B(M)\eta(M)\).  Identity
\eqref{eq:V104-oriented-rank-product} is the definition of
\(\chi_{ef}\).
\end{proof}

\begin{corollary}[Near-unit coherent-row gate]
\label{cor:V104-near-unit-coherent-gate}
For \(M_N=s_N^2P_{S_N}\),
\begin{align}
 \|Q_e\otimes Q_f\otimes M_N\|_{\rm prod}
 \le
 B_N\int_0^{A_e}\!\!\int_0^{A_f}
 \min\left\{
 \eta_N,\frac{D_e(t)D_f(s)}{d_{ef}}
 \right\}\dd s\,\dd t,
 \label{eq:V104-near-unit-coherent-gate}
\end{align}
where
\[
 B_N=s_N^2\frac{N(N+2)}3,
 \qquad
 \eta_N=\frac3{(N+1)(N+2)}.
\]
\end{corollary}

\begin{proof}
Here \(\min\{d_F,d_{R_N}\}=3\).  Apply
\Cref{prop:V104-coherent-ratio-relaxation} and the exact V103 core
identities.
\end{proof}

\begin{remark}[The orientation factor is structural]
\label{rem:V104-orientation-structural}
Replacing \(d_{ef}\) by \(d_ed_f\) would silently assume that the larger
side of each pair lies on the same original row.  The alternating
example in \Cref{lem:V104-orientation-comparison} shows that this can
lose a factor \(L\).  Thus row-orientation data must be retained by the
V94 selector; it cannot be reconstructed from the two scalar pair
capacities.
\end{remark}

\section{The combined spectrum and its exact fractional mean}
\label{sec:V104-combined-spectrum}

Regard
\begin{equation}
 Q_{ef}:=Q_e\otimes Q_f
 \label{eq:V104-combined-operator}
\end{equation}
as one positive invariant operator between the auxiliary row
representations
\[
 X_{ef}:=X_e\boxtimes X_f,
 \qquad
 Y_{ef}:=Y_e\boxtimes Y_f.
\]
Their larger dimension is precisely \(d_{ef}\).  If the eigenvalues of
\(Q_e\) and \(Q_f\), including multiplicities, are
\[
 a_1\ge\cdots\ge a_{m_e}\ge0,
 \qquad
 b_1\ge\cdots\ge b_{m_f}\ge0,
\]
then the eigenvalue multiset of \(Q_{ef}\) is
\begin{equation}
 \{a_ib_j:1\le i\le m_e,\ 1\le j\le m_f\}.
 \label{eq:V104-product-spectrum}
\end{equation}
List these products in nonincreasing order as
\[
 c_1^{ef}\ge c_2^{ef}\ge\cdots\ge0.
\]

\begin{definition}[Combined fractional profile]
\label{def:V104-combined-profile}
Let
\[
 D_{ef}(r):=\operatorname{rank}\mathbf1_{\{Q_{ef}>r\}}.
\]
For \(0\le\eta\le1\), define
\begin{equation}
 \mathcal K_{ef}(\eta)
 :=
 \int_0^\infty
 \min\left\{\eta,\frac{D_{ef}(r)}{d_{ef}}\right\}\dd r.
 \label{eq:V104-combined-profile}
\end{equation}
\end{definition}

\begin{theorem}[Combined Ky--Fan identity and core gate]
\label{thm:V104-combined-KyFan}
Let \(q=\eta d_{ef}\) and \(n=\lfloor q\rfloor\).  Then
\begin{equation}
 \boxed{\;
 \mathcal K_{ef}(\eta)
 =
 \frac1{d_{ef}}
 \left[
 \sum_{\ell=1}^{n}c_\ell^{ef}
 +(q-n)c_{n+1}^{ef}
 \right].
 \;}
 \label{eq:V104-combined-KyFan}
\end{equation}
The fractional final term is omitted when \(q=n\).  For an invariant
core shell \(M\) as above,
\begin{equation}
 \boxed{\;
 \|Q_e\otimes Q_f\otimes M\|_{\rm prod}
 \le B(M)\mathcal K_{ef}(\eta(M)).
 \;}
 \label{eq:V104-combined-core-gate}
\end{equation}
\end{theorem}

\begin{proof}
The layer-cake identity gives
\[
 Q_{ef}=\int_0^\infty\mathbf1_{\{Q_{ef}>r\}}\,\dd r.
\]
For each spectral projector, repeat the proof of
\Cref{lem:V104-two-projection-product-cap,lem:V104-double-Schur} with
the single auxiliary pair \((X_{ef},Y_{ef})\).  The product cap is
\(\Pi(M)\), while the relaxed Schur cap is
\[
 B(M)\frac{D_{ef}(r)}{d_{ef}}.
\]
Integration proves \eqref{eq:V104-combined-core-gate}.  For the identity,
write
\[
 D_{ef}(r)=\#\{\ell:c_\ell^{ef}>r\}
\]
and integrate the clipped rank distribution.  Partitioning at the
eigenvalues, exactly as in the V90 one-pair calculation, gives
\eqref{eq:V104-combined-KyFan}.
\end{proof}

\begin{corollary}[Normalized near-unit combined gate]
\label{cor:V104-normalized-near-unit}
Put \(q_N^{ef}:=\eta_Nd_{ef}\) and, for \(q_N^{ef}>0\), define
\begin{equation}
 \overline c_{ef,N}
 :=
 \frac1{q_N^{ef}}
 \left[
 \sum_{\ell=1}^{\lfloor q_N^{ef}\rfloor}c_\ell^{ef}
 +(q_N^{ef}-\lfloor q_N^{ef}\rfloor)
 c_{\lfloor q_N^{ef}\rfloor+1}^{ef}
 \right].
 \label{eq:V104-combined-top-mean}
\end{equation}
Then
\begin{equation}
 \boxed{\;
 \|Q_e\otimes Q_f\otimes M_N\|_{\rm prod}
 \le
 s_N^2\frac N{N+1}\overline c_{ef,N}.
 \;}
 \label{eq:V104-normalized-near-unit}
\end{equation}
\end{corollary}

\begin{proof}
Use \(B_N\eta_N=\Pi_N=s_N^2N/(N+1)\) in
\eqref{eq:V104-combined-core-gate} and
\eqref{eq:V104-combined-KyFan}.
\end{proof}

\section{Two cold/hot decompositions inside the common truncation}
\label{sec:V104-two-cold-hot}

The gain from two forced pairs is not obtained by multiplying their full
average capacities.  It comes from an elementary but different fact:
an eigenvalue product \(a_ib_j\) can be large only when \emph{both}
factors lie in their exceptional cold sets.

\begin{lemma}[Cold--cold exceptional rectangle]
\label{lem:V104-cold-cold-rectangle}
For \(g=e,f\), suppose the spectrum of \(Q_g\) has a distinguished cold
index set \(\mathscr C_g\) of cardinality \(r_g\), and constants
\(A_g,h_g\ge0\) such that
\begin{equation}
 0\le c_i^g\le A_g
 \quad\hbox{for all \(i\)},\qquad
 c_i^g\le h_g
 \quad\hbox{for \(i\notin\mathscr C_g\)}.
 \label{eq:V104-abstract-cold-hot}
\end{equation}
Then, for \(0<\eta\le1\), the top
\(\eta d_{ef}\)-fractional mean of the product spectrum satisfies
\begin{equation}
 \boxed{\;
 \overline c_{ef}(\eta)
 \le
 h_{ef}
 A_eA_f
 \min\left\{
 1,\frac{r_er_f}{\eta d_{ef}}
 \right\},
 \qquad
 h_{ef}:=\max\{h_eA_f,A_eh_f\}.
 \;}
 \label{eq:V104-cold-cold-fractional}
\end{equation}
\end{lemma}

\begin{proof}
If \((i,j)\notin\mathscr C_e\times\mathscr C_f\), then either
\(i\notin\mathscr C_e\) or \(j\notin\mathscr C_f\).  In the first case
\[
 a_ib_j\le h_eA_f,
\]
and in the second
\[
 a_ib_j\le A_eh_f.
\]
Thus every product eigenvalue outside the cold--cold rectangle is at
most \(h_{ef}\), and that rectangle has rank \(r_er_f\).  A top
\(\eta d_{ef}\)-fractional sum contains baseline at most \(h_{ef}\) on
all of its mass and can receive an additional contribution at most
\(A_eA_f\) on the smaller of its own mass and the exceptional rank.
Divide by \(\eta d_{ef}\).  This proves
\eqref{eq:V104-cold-cold-fractional}.
\end{proof}

\begin{theorem}[Two-pair V53 effective-number gate]
\label{thm:V104-two-pair-effective-gate}
Suppose each forced pair satisfies its V53 selected-heavy decomposition
with
\begin{equation}
 \frac{r_e}{d_e}\le\frac{C}{R_e},
 \qquad
 h_e\le\frac{C}{R_e},
 \qquad
 \frac{r_f}{d_f}\le\frac{C}{R_f},
 \qquad
 h_f\le\frac{C}{R_f},
 \label{eq:V104-two-V53-data}
\end{equation}
where
\[
 R_e=s_\alpha N_{{\rm eff},e},
 \qquad
 R_f=s_\alpha N_{{\rm eff},f},
\]
and suppose \(A_e,A_f\) are upper coefficient bounds.  Then
\begin{align}
 \overline c_{ef,N}
 &\le
 C\max\left\{\frac{A_f}{R_e},\frac{A_e}{R_f}\right\}
 \notag\\
 &\quad+
 A_eA_f
 \min\left\{
 1,\frac{C^2\chi_{ef}}{\eta_NR_eR_f}
 \right\}.
 \label{eq:V104-two-pair-fractional-gate}
\end{align}
Consequently,
\begin{align}
 &\|Q_e\otimes Q_f\otimes M_N\|_{\rm prod}
 \notag\\
 &\quad\le
 s_N^2\frac N{N+1}
 \left[
 C\max\left\{\frac{A_f}{R_e},\frac{A_e}{R_f}\right\}
 +
 A_eA_f
 \min\left\{
 1,\frac{C^2\chi_{ef}(N+1)(N+2)}
 {3R_eR_f}
 \right\}
 \right].
 \label{eq:V104-two-pair-core-gate}
\end{align}
\end{theorem}

\begin{proof}
By \eqref{eq:V104-two-V53-data},
\[
 \frac{r_er_f}{d_{ef}}
 =
 \chi_{ef}\frac{r_e}{d_e}\frac{r_f}{d_f}
 \le
 \frac{C^2\chi_{ef}}{R_eR_f}.
\]
Moreover
\[
 h_{ef}
 \le
 C\max\left\{\frac{A_f}{R_e},\frac{A_e}{R_f}\right\}.
\]
Apply \Cref{lem:V104-cold-cold-rectangle} with
\(\eta=\eta_N\) to obtain
\eqref{eq:V104-two-pair-fractional-gate}.  Insert that estimate into
\eqref{eq:V104-normalized-near-unit} and use the exact value of
\(\eta_N\).
\end{proof}

\begin{corollary}[Conditional linear effective-size threshold]
\label{cor:V104-linear-threshold}
Assume
\begin{equation}
 \sup_N(s_N^2+A_e+A_f+\chi_{ef})<\infty
 \label{eq:V104-bounded-prefactors}
\end{equation}
and
\begin{equation}
 R_e\longrightarrow\infty,\qquad
 R_f\longrightarrow\infty,\qquad
 \frac{R_eR_f}{N^2}\longrightarrow\infty.
 \label{eq:V104-asymmetric-threshold}
\end{equation}
Then
\begin{equation}
 \|Q_e\otimes Q_f\otimes M_N\|_{\rm prod}
 \longrightarrow0.
 \label{eq:V104-two-pair-decay}
\end{equation}
In the symmetric regime \(R_e\asymp R_f\asymp R\), it is sufficient
that
\begin{equation}
 \boxed{\qquad \frac{R}{N}\longrightarrow\infty. \qquad}
 \label{eq:V104-linear-threshold}
\end{equation}
\end{corollary}

\begin{proof}
Under \eqref{eq:V104-bounded-prefactors}--%
\eqref{eq:V104-asymmetric-threshold}, both terms in brackets in
\eqref{eq:V104-two-pair-core-gate} tend to zero.  The factor
\(s_N^2N/(N+1)\) is bounded.  This proves
\eqref{eq:V104-two-pair-decay}.  If \(R_e,R_f\asymp R\), condition
\eqref{eq:V104-asymmetric-threshold} is equivalent, up to fixed
constants, to \(R/N\to\infty\).
\end{proof}

\begin{remark}[Asymmetric allocation law]
\label{rem:V104-asymmetric-allocation}
The gain is governed by the geometric mean:
\[
 \sqrt{\frac{R_eR_f}{\chi_{ef}}}\gg N.
\]
Thus one very strong endpoint does not compensate for a bounded other
endpoint, because the hot baseline in
\eqref{eq:V104-two-pair-fractional-gate} also requires
\(\min\{R_e,R_f\}\to\infty\).  This is the exact allocation condition
visible to the present information.
\end{remark}

\begin{corollary}[Translation to selected pair masses]
\label{cor:V104-pair-mass-translation}
For each endpoint the V53 effective number satisfies
\[
 N_{{\rm eff},g}\le x_g,\qquad g=e,f.
\]
Consequently the symmetric sufficient condition
\eqref{eq:V104-linear-threshold} requires, in particular,
\begin{equation}
 \frac{s_\alpha\sqrt{x_ex_f}}{N}\longrightarrow\infty.
 \label{eq:V104-mass-necessary-for-route}
\end{equation}
This is necessary for this V53 route, not sufficient without lower
control of the effective numbers.
\end{corollary}

\begin{proof}
Apply \Cref{lem:V103-effective-below-mass} separately to the two selected
pairs.  Then
\[
 \sqrt{R_eR_f}
 =
 s_\alpha\sqrt{N_{{\rm eff},e}N_{{\rm eff},f}}
 \le s_\alpha\sqrt{x_ex_f}.
\]
\end{proof}

\section{Sharp tests and the remaining physical extraction problem}
\label{sec:V104-tests}

\begin{proposition}[Rank-one sharpness of the common truncation]
\label{prop:V104-rank-one-sharpness}
For every pair of auxiliary row sizes \(d_e,d_f\), there are positive
pair operators with
\[
 \mathcal K_e(1)=\frac1{d_e},
 \qquad
 \mathcal K_f(1)=\frac1{d_f},
\]
but if
\begin{equation}
 \eta_Nd_{ef}<1,
 \label{eq:V104-rank-one-subcritical}
\end{equation}
then
\[
 \overline c_{ef,N}=1
\]
and the combined near-unit gate gives only
\begin{equation}
 \|Q_e\otimes Q_f\otimes M_N\|_{\rm prod}
 \le s_N^2\frac N{N+1}.
 \label{eq:V104-rank-one-core-cap}
\end{equation}
\end{proposition}

\begin{proof}
Take each \(Q_g\) to have one eigenvalue one and all other eigenvalues
zero.  Its V90 full capacity is \(1/d_g\).  The product operator also
has exactly one nonzero eigenvalue, equal to one.  Under
\eqref{eq:V104-rank-one-subcritical}, the fractional Ky--Fan sum in
\eqref{eq:V104-combined-top-mean} takes a fraction
\(\eta_Nd_{ef}\) of that eigenvalue and divides by the same fraction.
Thus its mean is one, and \eqref{eq:V104-rank-one-core-cap} follows from
\eqref{eq:V104-normalized-near-unit}.
\end{proof}

\begin{assessment}[What the rank-one test rules out]
\label{ass:V104-rank-one-meaning}
Even two vanishing full average capacities do not imply decay against
the near-unit core.  If \(d_{ef}=o(N^2)\), the common truncation can
still see only the single top eigenvalue.  The V104 improvement is
therefore not a consequence of small averages; it comes specifically
from quantitative cold-rank bounds for \emph{both} spectra, combined
before truncation.
\end{assessment}

\begin{proposition}[Alternating orientation can erase the linear gain]
\label{prop:V104-alternating-loss}
Suppose \(R_e\asymp R_f\asymp R\), while
\[
 (d_{X_e},d_{Y_e})=(L,1),
 \qquad
 (d_{X_f},d_{Y_f})=(1,L).
\]
Then \(\chi_{ef}=L\), and the exceptional term in
\eqref{eq:V104-two-pair-fractional-gate} is forced small by the present
data only if
\begin{equation}
 \frac{R^2}{LN^2}\longrightarrow\infty.
 \label{eq:V104-alternating-threshold}
\end{equation}
\end{proposition}

\begin{proof}
Insert \(\chi_{ef}=L\) and
\(\eta_N\asymp N^{-2}\) into
\eqref{eq:V104-two-pair-fractional-gate}.
\end{proof}

\begin{firewall}[Tensor separation is not yet supplied by endpoint forcing]
\label{fire:V104-tensor-separation}
The V102 endpoint localization proves that, outside paid cut/mean and
off-cut alternatives, each sideways endpoint contains a heavy internal
pair.  It does not yet prove that the two chosen pairs:
\begin{enumerate}
\item occupy disjoint active gauge-coordinate banks;
\item survive the V94 selector simultaneously;
\item yield the tensor factor \(Q_e\otimes Q_f\) rather than a correlated
      signed column operator; or
\item have bounded orientation loss \(\chi_{ef}\).
\end{enumerate}
Therefore \Cref{cor:V104-linear-threshold} is a rigorous conditional
closure theorem for the extracted two-pair channel, not yet a closure of
every physical sideways monomial.
\end{firewall}

\begin{assessment}[Upstream bottleneck after the spectral calculation]
\label{ass:V104-upstream-bottleneck}
The near-unit spectral obstruction has now been reduced from a
one-pair \(N^2\) threshold to a two-pair \(N\) threshold whenever the
physical selector supplies coherent tensor separation.  The next
question is combinatorial and upstream: can the two endpoint pair
certificates be chosen disjointly and coherently, or can overlap be
charged to an already-paid cut/off-cut stock?  Continuing to sharpen a
single-pair average would be circular until that extraction question is
settled.
\end{assessment}

\section{V104 ledger and V105 acquisition}
\label{sec:V104-ledger}

\begin{theorem}[Integrated V104 statement]
\label{thm:V104-integrated}
Preserving V1--V103, V104 proves:
\begin{enumerate}[label=\textup{(V104.\arabic*)},leftmargin=6.3em]
\item the exact common two-level law
      \eqref{eq:V104-one-minimum-two-level};
\item the exact near-unit truncation surface
      \eqref{eq:V104-near-unit-surface};
\item the coherent row dimension and orientation factor
      \eqref{eq:V104-orientation-factor};
\item the combined Ky--Fan identity and one-core gate
      \eqref{eq:V104-combined-KyFan}--%
      \eqref{eq:V104-combined-core-gate};
\item the cold--cold exceptional-rank theorem
      \eqref{eq:V104-two-pair-fractional-gate};
\item conditional decay when
      \(R_eR_f/(\chi_{ef}N^2)\to\infty\) and both effective scales
      diverge;
\item the coherent symmetric threshold \(R\gg N\); and
\item rank-one, alternating-orientation, and physical-extraction
      firewalls.
\end{enumerate}
\end{theorem}

\begin{proof}
The items are
\Cref{thm:V104-one-minimum-two-level,
cor:V104-near-unit-surface,
def:V104-coherent-row-dimension,
thm:V104-combined-KyFan,
thm:V104-two-pair-effective-gate,
cor:V104-linear-threshold,
prop:V104-rank-one-sharpness,
prop:V104-alternating-loss,
fire:V104-tensor-separation}.
\end{proof}

\begin{openproblem}[V105 mandate: audited two-endpoint extraction]
\label{op:V105-mandate}
Before continuing, perform the compulsory V105 comparison with the
current Standard Model and Navier--Stokes notes.  Then return upstream
to the complete V94 selector and the V102 simultaneous endpoint
certificates.
\begin{enumerate}
\item State the exact incidence conditions under which the two forced
      internal pairs are disjoint active coordinate banks.
\item If they overlap, determine whether the shared edge or shared
      puncture necessarily crosses a paid cut, creates an off-cut
      outcome, or admits a separate diagonal estimate.
\item Track which auxiliary representation lies on each retained row
      and bound \(\chi_{ef}\) from the actual incidence data.
\item Prove simultaneous selector survival before replacing the physical
      column by \(Q_e\otimes Q_f\).
\item If tensor separation fails, retain the joint signed pair/core
      matrix and formulate its common Ky--Fan profile rather than
      multiplying marginals.
\item Keep the V102 off-cut alternative separate throughout.
\end{enumerate}
\end{openproblem}

\section{Firewall after V104}
\label{sec:V104-firewall}

\begin{firewall}[Current claim boundary]
\label{fire:V104-current-boundary}
V104 does not prove that every physical side block contains two
disjoint coherently oriented pair banks, establish the tensor
factorization used by the conditional theorem, force
\(R_eR_f\gg\chi_{ef}N^2\), bound the complete signed core scalar, close
the off-cut branch, prove JREC, construct the continuum Yang--Mills
theory, or prove a positive mass gap.  It proves the exact common
two-pair truncation law and identifies the strongest consequence of two
independent V53 cold/hot decompositions.  The full Yang--Mills existence
and mass gap remain open.
\end{firewall}

\section{Append-only record for V104}
\label{sec:V104-record}

V104 was appended to the complete V103 manuscript.  The exact V103
parent had \(70107\) physical lines, \(2449827\) bytes, and SHA--256
digest
\[
 \texttt{63fa8eab0754baceaed3b97fd38d3aa00f564eb33cb5832061c5b7d41fd40dbe}.
\]
No V103 mathematical content was changed.  The sole final
\verb|\end{document}| is restored below.

% =============================================================================
% END APPENDED VERSION V104 MATERIAL
% =============================================================================
% =============================================================================
% BEGIN APPENDED VERSION V105 MATERIAL
% =============================================================================

\chapter{V105: Companion Audit and the Exact Two-Endpoint
Pair-Overlap Dichotomy}
\label{ch:V105}

\section{Purpose}
\label{sec:V105-purpose}

V104 proved a conditional linear-threshold theorem when the two
sideways-endpoint pair certificates occupy disjoint coordinate banks and
therefore contribute the tensor-separated auxiliary operator
\(Q_e\otimes Q_f\).  V102 did not prove that separation.  The present
version performs the required five-version companion audit and then returns
upstream to the V59/V60 incidence alphabet.

The incidence calculation gives an exact alternative.  From two heavy pair
certificates one can extract two disjoint coordinate banks, each retaining
one quarter of its original certified mass, unless a single pure overlap
type is heavy for both endpoints.  There are only two such overlap types:
the endpoints are the two rows of the same local pair block, or they occupy
the two complementary pair blocks of a local \([2,2]\) partition.  The first
type supplies only one auxiliary pair operator.  In the second type, V56
gives an exact pre-recoupling tensor product of two pair defects and V58
permits their row-factorized trace taxes to compose.  But both factors
transform under the same coordinate gauge group; after regrouping by
retained row, their row representations are generally reducible.  Thus the
external-product Schur scalar used in V104 is not yet available.  This
classification closes the combinatorial ambiguity left in V104 and
isolates the diagonal-row channel resolution still required.

\section{Mandatory V105 companion audit}
\label{sec:V105-audit}

The companion files were read before selecting the V105 acquisition.  Their
stable snapshots were:
\begin{center}
\small
\begin{tabular}{@{}p{0.50\textwidth}r r p{0.25\textwidth}@{}}
\toprule
\textbf{source}&\textbf{lines}&\textbf{bytes}&\textbf{SHA--256}\\
\midrule
\texttt{Renewal\_SM\_Einstein\_Continuum\_Research\_Notes\_V114.tex}
&78315&2877078&
\texttt{a1c1664ef926163d62d74a5637b8a4a65e291a3e0ccfc561c287f6c8c6e0b772}\\
\texttt{Renewal\_NS\_Stability\_Research\_Notes\_V31.tex}
&23052&887537&
\texttt{f1121b62238ad5491bb7cf03635ea9934942edf573ed8faaa9ee79e1d38a6696}\\
\bottomrule
\end{tabular}
\end{center}

The Navier--Stokes file is unchanged from the V100 audit.  Its V31 exact
dyadic mark still proves that separate probability-normalized controls do
not pay the unbounded first moment: the remaining datum is the correlation
between stopping times, signed triad work, and active parent tails.

The Standard--Model--Einstein file advanced from V103 to V114.  Its
V104--V113 sequence retains the complete correlated stable determinant and
proves a no-go theorem for every finite real first-order conservative
conormal graph.  V114 then goes upstream to the boundary incidence relation:
the free normal ghost jet forces a normal gauge direction into the free
conormal sector, where the preceding determinant obstruction reappears.

\begin{assessment}[Transfer used in V105]
\label{ass:V105-audit-transfer}
No SM boundary estimate or NS PDE estimate is imported into Yang--Mills
fusion.  The transferable proof discipline is exact: preserve a correlated
object until the incidence constraints on it are proved.  Accordingly V105
does not assume that two endpoint certificates are independent and does not
multiply their capacities.  It first classifies their coordinate overlap
inside the V60 local partition.
\end{assessment}

\begin{firewall}[Companion separation]
\label{fire:V105-companion-separation}
The SM root-reflection theorem is not an \(SU(3)\) Racah theorem, and the NS
formation mark is not a Wilson-shell estimate.  The audit changes only the
order of acquisition.  Every estimate below is proved from the existing YM
incidence alphabet.
\end{firewall}

\section{Support form of the V60 pair alphabet}
\label{sec:V105-support-alphabet}

Fix one completely resolved V60 local partition monomial \(\mu\).  Let
\(\mathcal F_\mu\) be its residual puncture-coordinate set.  For distinct
replica rows \(i,j\), define
\begin{equation}
 \mathcal P_{ij}
 :=
 \{f\in\mathcal F_\mu:
   \{i,j\}\text{ is a pair block of the local partition at }f\}.
 \label{eq:V105-pair-support}
\end{equation}
Thus \(\mathcal P_{ij}=\mathcal P_{ji}\).  For a row \(v\) and a coordinate
set \(E\subseteq\mathcal F_\mu\), write
\begin{equation}
 m_v(E):=\sum_{f\in E}A_{v,f}.
 \label{eq:V105-row-mass-measure}
\end{equation}
The exact V60 mass is
\begin{equation}
 D_{ij,\mu}^{\rm pun}=m_i(\mathcal P_{ij}).
 \label{eq:V105-D-support-form}
\end{equation}

\begin{lemma}[Local overlap alphabet for two endpoints]
\label{lem:V105-local-overlap-alphabet}
Fix distinct endpoints \(i,k\), and suppose that at a coordinate \(f\)
row \(i\) lies in the pair block \(\{i,j\}\) while row \(k\) lies in the
pair block \(\{k,\ell\}\).  Exactly one of the following holds:
\begin{enumerate}[label=\textup{(\alph*)}]
\item the blocks coincide, necessarily
      \begin{equation}
       \{i,j\}=\{k,\ell\}=\{i,k\};
       \label{eq:V105-same-pair-type}
      \end{equation}
\item the blocks are disjoint, so \(i,j,k,\ell\) are four distinct rows
      and the local partition contains the complementary pair blocks
      \begin{equation}
       \{i,j\}\mathbin{\dot\cup}\{k,\ell\}.
       \label{eq:V105-complementary-pair-type}
      \end{equation}
\end{enumerate}
In particular, two distinct pair blocks at one coordinate cannot share
exactly one row.
\end{lemma}

\begin{proof}
The blocks of a set partition are either identical or disjoint.  If the two
blocks are identical, that block contains both \(i\) and \(k\); since it has
cardinality two, it is \(\{i,k\}\).  If they are distinct, they are disjoint,
which gives four distinct rows and
\eqref{eq:V105-complementary-pair-type}.  These alternatives are exhaustive.
\end{proof}

For fixed pair labels \((i,j)\) and \((k,\ell)\), put
\begin{align}
 E_i&:=\mathcal P_{ij},
 &
 E_k&:=\mathcal P_{k\ell},
 \notag\\
 O&:=E_i\cap E_k,
 &
 U_i&:=E_i\setminus O,
 &
 U_k&:=E_k\setminus O.
 \label{eq:V105-overlap-exclusive-split}
\end{align}
The sets \(U_i,U_k,O\) are pairwise disjoint.  By
\Cref{lem:V105-local-overlap-alphabet}, the overlap has the disjoint
type split
\begin{equation}
 O=O_{\rm same}\mathbin{\dot\cup}O_{\rm comp},
 \label{eq:V105-overlap-type-split}
\end{equation}
where the first set has the common block \(\{i,k\}\) and the second has
two complementary pair blocks.

\begin{theorem}[Quarter-mass disjoint-or-pure-overlap dichotomy]
\label{thm:V105-quarter-dichotomy}
Assume the two endpoint certificates satisfy
\begin{equation}
 m_i(E_i)\ge h_i,
 \qquad
 m_k(E_k)\ge h_k
 \label{eq:V105-two-heavy-certificates}
\end{equation}
for positive \(h_i,h_k\).  Then exactly one of the following conclusions
can be selected:
\begin{enumerate}[label=\textup{(D)}]
\item there are disjoint coordinate sets
      \(G_i\subseteq E_i\), \(G_k\subseteq E_k\) such that
      \begin{equation}
       m_i(G_i)\ge\frac{h_i}{4},
       \qquad
       m_k(G_k)\ge\frac{h_k}{4};
       \label{eq:V105-quarter-disjoint}
      \end{equation}
\end{enumerate}
or
\begin{enumerate}[label=\textup{(O)}]
\item for one pure type
      \(\tau\in\{\mathrm{same},\mathrm{comp}\}\),
      \begin{equation}
       m_i(O_\tau)\ge\frac{h_i}{4},
       \qquad
       m_k(O_\tau)\ge\frac{h_k}{4}.
       \label{eq:V105-quarter-pure-overlap}
      \end{equation}
\end{enumerate}
The alternatives need not be mutually exclusive, but the stated selection
is always possible.
\end{theorem}

\begin{proof}
The exact split \(E_i=U_i\dot\cup O\) implies that either
\[
 m_i(U_i)\ge\frac{h_i}{2}
 \quad\hbox{or}\quad
 m_i(O)\ge\frac{h_i}{2}.
\]
The analogous alternative holds for row \(k\).  If at least one row chooses
its exclusive set and the other chooses either its exclusive set or \(O\),
the two chosen sets are disjoint and already retain half the corresponding
certificate masses.  This gives \textup{(D)}, with room to weaken one half
to one quarter.

It remains only when
\[
 m_i(O)\ge\frac{h_i}{2},
 \qquad
 m_k(O)\ge\frac{h_k}{2}.
\]
Use \eqref{eq:V105-overlap-type-split}.  For each row, at least one of
\(O_{\rm same},O_{\rm comp}\) carries one quarter of its certificate mass.
If the two rows choose different types, those two type sets are disjoint and
give \textup{(D)}.  If they choose the same type \(\tau\), that type obeys
\eqref{eq:V105-quarter-pure-overlap}, giving \textup{(O)}.
\end{proof}

\begin{corollary}[Quantitative V102 endpoint form]
\label{cor:V105-V102-quarter}
On the non-paid, non-off-cut branch of
\Cref{thm:V102-forced-pair-or-offcut}, two sideways endpoints \(i,k\)
admit either disjoint pair banks satisfying
\begin{equation}
 m_i(G_i)\ge\frac{\delta_0}{24}\Xi_i,
 \qquad
 m_k(G_k)\ge\frac{\delta_0}{24}\Xi_k,
 \label{eq:V105-V102-disjoint-mass}
\end{equation}
or one pure same-pair/complementary-pair overlap bank satisfying both
inequalities with \(G_i,G_k\) replaced by that common bank.
\end{corollary}

\begin{proof}
The V102 certificates have
\[
 h_i=\frac{\delta_0}{6}\Xi_i,
 \qquad
 h_k=\frac{\delta_0}{6}\Xi_k.
\]
Apply \Cref{thm:V105-quarter-dichotomy}.
\end{proof}

\section{Disposition of the disjoint branch}
\label{sec:V105-disjoint-branch}

\begin{proposition}[Disjoint supports give genuine coordinate-bank
factorization]
\label{prop:V105-disjoint-factorization}
Suppose conclusion \textup{(D)} of
\Cref{thm:V105-quarter-dichotomy} holds.  Remove all identity spectators
as in V94 and include every remaining weight outside \(G_i\cup G_k\) in
one positive core shell \(M\).  Then the two selected pair banks act on
distinct coordinate tensor factors, and the positive carrier has the exact
form
\begin{equation}
 Q_{G_i}\otimes Q_{G_k}\otimes M,
 \label{eq:V105-disjoint-factorization}
\end{equation}
regrouped by the same two retained original rows.  Hence the V104 common
two-pair law applies, with the selected-heavy mass constants weakened by at
most the fixed factor four.
\end{proposition}

\begin{proof}
The coordinate construction of the V54 puncture formula is a tensor product
over gauge-coordinate indices.  The sets \(G_i,G_k\) are disjoint, and the
puncture coordinates are disjoint from the selected collision/core
coordinates.  Fubini factorization of the finite Haar integrals therefore
places the operators belonging to \(G_i\), \(G_k\), and the remaining core
on three tensor factors.  Regrouping those factors by retained row gives
\eqref{eq:V105-disjoint-factorization}; it does not alter the product-state
cut.

The selected mass in each new bank is at least one quarter of the V102
certificate by \eqref{eq:V105-quarter-disjoint}.  Every V53 cold-rank/hot
estimate depends on that fixed mass fraction only through its numerical
selected-heavy constant, so the conclusion of
\Cref{thm:V104-two-pair-effective-gate} remains valid after changing \(C\)
by a fixed factor.
\end{proof}

\begin{corollary}[Conditional closure of the disjoint coherent branch]
\label{cor:V105-disjoint-coherent-closure}
On the disjoint branch, assume the reduced banks have V53 effective scales
\[
 R_i=s_\alpha N_{{\rm eff},i},
 \qquad
 R_k=s_\alpha N_{{\rm eff},k},
\]
uniform coefficient caps, bounded \(s_N^2\), and bounded orientation factor
\(\chi_{ik}\).  If
\begin{equation}
 R_i\to\infty,\qquad
 R_k\to\infty,\qquad
 \frac{R_iR_k}{N^2}\to\infty,
 \label{eq:V105-disjoint-effective-condition}
\end{equation}
then its \(F\otimes R_N\) near-unit carrier has product capacity tending to
zero.
\end{corollary}

\begin{proof}
Apply \Cref{prop:V105-disjoint-factorization} and then
\Cref{cor:V104-linear-threshold}.  Fixed losses in the selected mass do not
change any limit in \eqref{eq:V105-disjoint-effective-condition}.
\end{proof}

\begin{firewall}[Incidence does not bound representation orientation]
\label{fire:V105-incidence-no-orientation}
The V60 alphabet records which replica rows belong to a pair block.  It does
not compare the dimensions of the resulting input irreducibles on the two
retained original rows.  Therefore disjoint coordinate supports alone do
not imply \(\chi_{ik}=O(1)\).  The alternating-dimension regression
\Cref{prop:V104-alternating-loss} remains live until the actual pair
representations are read from the resolved outcome.
\end{firewall}

\section{The two pure overlap branches}
\label{sec:V105-overlap-branches}

\begin{proposition}[Same-pair overlap is one pair, not two]
\label{prop:V105-same-pair-one-operator}
On \(O_{\rm same}\), the two endpoint certificates refer to the same local
pair block \(\{i,k\}\) at every coordinate.  After same-label regrouping
there is one pair operator \(Q_{ik,O_{\rm same}}\).  Its occurrence at both
endpoints does not produce
\(Q_{ik,O_{\rm same}}\otimes Q_{ik,O_{\rm same}}\).
The valid near-unit bound is the one-pair V103 gate.
\end{proposition}

\begin{proof}
At every coordinate in \(O_{\rm same}\),
\eqref{eq:V105-same-pair-type} identifies the two claimed blocks.  The local
puncture factor contains that block once.  Multiplication over coordinates
and same-label regrouping therefore create one operator on one coordinate
bank.  Tensoring it with a second copy would duplicate the same Haar factor
and the same physical mass.  The V102 no-double-charge convention then
requires the one-pair carrier \(Q_{ik,O_{\rm same}}\otimes M_N\), to which
\Cref{thm:V103-fractional-KyFan} applies.
\end{proof}

\begin{proposition}[Complementary overlap has an exact pre-recoupling
tensor factor]
\label{prop:V105-complementary-joint-operator}
On \(O_{\rm comp}\), every coordinate has two disjoint local pair blocks
\[
 \{i,j\}\mathbin{\dot\cup}\{k,\ell\}.
\]
The exact V56 collision identity gives the local double-defect factor
\begin{equation}
 \mathbf d_{ij,f}\otimes\mathbf d_{k\ell,f}.
 \label{eq:V105-local-double-defect}
\end{equation}
After multiplication over \(O_{\rm comp}\), the two pair carriers remain
tensor separated before final recoupling, as in V58.  Denote the resulting
positive four-row auxiliary effect, after the required common recoupling,
by
\begin{equation}
 J_{\rm comp}.
 \label{eq:V105-Jcomp}
\end{equation}
Nevertheless V104 cannot yet be applied: the two factors in
\eqref{eq:V105-local-double-defect} transform under the same coordinate
copy of the gauge group, so the retained-row representations are diagonal
tensor products and need not be irreducible.
\end{proposition}

\begin{proof}
Equation \eqref{eq:V56-collision-factor} is exactly
\eqref{eq:V105-local-double-defect} after relabeling the four rows.  The two
blocks act on disjoint row-carrier factors, so multiplication over the bank
preserves their tensor form before the common final recoupling.
\Cref{lem:V58-tax-composition} proves that this deterministic factorization,
not probabilistic independence, is sufficient for composing the two trace
taxes on row-factorized coefficients.

On a retained original row, however, the group acts diagonally on the two
auxiliary factors.  Thus a row such as \(X_e\otimes X_f\) is an ordinary
tensor product representation of one group, rather than an external tensor
product for \(G_e\times G_f\).  It is generally reducible.  The irreducibility
step in \Cref{lem:V104-double-Schur} therefore cannot be invoked before its
Clebsch--Gordan channels and multiplicities are resolved.
\end{proof}

\begin{firewall}[Overlap is not automatically cut or off-cut]
\label{fire:V105-overlap-not-paid}
Both overlap types occur inside \(\mathcal F_\mu\), the residual puncture
coordinate set after the V59 cut/core partition.  Their definition does not
place them in \(D_r^{\rm off}\).  Hence a shared-pair or complementary-pair
overlap cannot be charged to the cut or off-cut ledger merely because two
endpoint certificates use it.
\end{firewall}

\begin{proposition}[Exact incidence regression: separation is not forced]
\label{prop:V105-incidence-regression}
The V59/V60 mass identities and sideways inequalities admit algebraic
incidence data in which two endpoints are sideways, all cut, private-mean,
puncture-singleton, and off-cut masses vanish, and the two forced pair
certificates are supported entirely on:
\begin{enumerate}[label=\textup{(\alph*)}]
\item one common same-pair bank; or
\item one common complementary-pair bank.
\end{enumerate}
Thus disjointness cannot be deduced from those identities alone.
\end{proposition}

\begin{proof}
Choose positive row masses \(\Xi_i,\Xi_k\) large enough that
\[
 s_\alpha\Xi_i>M,\qquad s_\alpha\Xi_k>M.
\]
Set the selected collision masses \(x_i=x_k=0\), and set every assigned-cut,
private, singleton, and outer off-cut mass to zero.  Put
\[
 Y_i=\Xi_i,\qquad Y_k=\Xi_k.
\]
For case \textup{(a)}, choose every residual coordinate to have the pair
block \(\{i,k\}\), and assign its row weights so that
\[
 D_{ik,\mu}^{\rm pun}=\Xi_i,\qquad
 D_{ki,\mu}^{\rm pun}=\Xi_k.
\]
For case \textup{(b)}, use four rows and give every residual coordinate the
partition
\(\{i,j\}\dot\cup\{k,\ell\}\), with the analogous row masses.

The V59 mass splits reduce to \(\Xi_v=Y_v\), and the V60 alphabet reduces to
the displayed pair mass.  The sideways inequalities hold because
\(x_v=0<\varepsilon_{\rm sw}\Xi_v\).  All proposed paid stocks are zero.
This constructs incidence data satisfying every cited algebraic identity.
It is a logical regression for those identities, not a claim that arbitrary
weights of this form realize a complete physical Wilson sequence.
\end{proof}

\section{The correlated overlap profile}
\label{sec:V105-correlated-profile}

Every fixed irreducible row-channel compression has an exact one-operator
truncation law.  Let \(J\ge0\) be a positive invariant overlap operator on
\emph{irreducible} auxiliary row representations \(U,V\), and put
\[
 d_J:=\max\{d_U,d_V\},
 \qquad
 E_J(r):=\mathbf1_{\{J>r\}},
 \qquad
 D_J(r):=\operatorname{rank}E_J(r).
\]
Define
\begin{equation}
 \mathcal K_J(\eta)
 :=
 \int_0^\infty
 \min\left\{\eta,\frac{D_J(r)}{d_J}\right\}\dd r.
 \label{eq:V105-joint-profile}
\end{equation}

\begin{theorem}[Exact fixed-channel correlated-overlap/core gate]
\label{thm:V105-correlated-overlap-gate}
Let \(M\ge0\) be an invariant core shell on irreducible rows \(A,B\), with
\[
 B(M)=\frac{\operatorname{Tr}M}{\min\{d_A,d_B\}},
 \qquad
 \eta(M)=\frac{\|M\|_{\rm prod}}{B(M)}.
\]
If the overlap bank is disjoint from the core coordinate bank, then
\begin{equation}
 \boxed{\;
 \|J\otimes M\|_{\rm prod}
 \le B(M)\mathcal K_J(\eta(M)).
 \;}
 \label{eq:V105-correlated-overlap-core-gate}
\end{equation}
For the near-unit core,
\begin{equation}
 \boxed{\;
 \|J\otimes M_N\|_{\rm prod}
 \le
 B_N\mathcal K_J(\eta_N)
 =
 s_N^2\frac N{N+1}\overline c_{J,N},
 \;}
 \label{eq:V105-correlated-near-unit-gate}
\end{equation}
where \(\overline c_{J,N}\) is the top
\(\eta_Nd_J\)-fractional Ky--Fan mean of the eigenvalues of \(J\).
\end{theorem}

\begin{proof}
Use the layer cake
\[
 J=\int_0^\infty E_J(r)\,\dd r.
\]
Since \(0\le E_J(r)\le I\), the product cap of
\(E_J(r)\otimes M\) is at most \(\|M\|_{\rm prod}\), by the density-matrix
argument of \Cref{lem:V104-two-projection-product-cap}.  Schur's lemma on
the two enlarged rows gives the trace cap
\[
 \frac{D_J(r)\operatorname{Tr}M}
 {\max\{d_Ud_A,d_Vd_B\}}
 \le
 B(M)\frac{D_J(r)}{d_J}.
\]
Take one minimum, integrate, and use
\(\|M\|_{\rm prod}=B(M)\eta(M)\).  This proves
\eqref{eq:V105-correlated-overlap-core-gate}.  The near-unit identity uses
the V103 values of \(B_N,\eta_N\) and the Ky--Fan formula
\eqref{eq:V104-combined-KyFan}, whose proof applies to every positive
operator \(J\).
\end{proof}

\begin{corollary}[Exact fixed-channel overlap acquisition condition]
\label{cor:V105-overlap-acquisition}
A sufficient condition for the pure overlap near-unit branch to decay is
\begin{equation}
 s_N^2\overline c_{J,N}\longrightarrow0.
 \label{eq:V105-overlap-decay-condition}
\end{equation}
For \(O_{\rm same}\), this is exactly the V103 one-pair problem.  For
\(O_{\rm comp}\), the statement applies after \(J_{\rm comp}\) has been
compressed to fixed irreducible channels on both retained rows.  Summing
those diagonal-row channels without a multiplicity or dimension loss is
not asserted here.
\end{corollary}

\begin{proof}
Use \eqref{eq:V105-correlated-near-unit-gate} and \(N/(N+1)\le1\).
The branch identifications are
\Cref{prop:V105-same-pair-one-operator,
prop:V105-complementary-joint-operator}.
\end{proof}

\begin{proposition}[Coincident-cold regression]
\label{prop:V105-coincident-cold-regression}
Two marginal cold-rank statements do not by themselves give a product-rank
statement on a shared coordinate bank.  More precisely, let \(d\to\infty\),
let \(P\) be a rank-\(r\) projection on a \(d\)-dimensional auxiliary row
space, and use the same cold projector for both endpoint labels.  Put
\[
 R:=\frac d r.
\]
Then each marginal cold fraction is \(1/R\), while the common cold
intersection still has rank \(r\), not \(r^2/d\).  If
\begin{equation}
 \eta_Nd\le r,
 \label{eq:V105-coincident-cold-condition}
\end{equation}
the top \(\eta_Nd\)-fractional mean of \(P\) equals one.
\end{proposition}

\begin{proof}
Both labeled cold spaces are \(\operatorname{Ran}P\), so their intersection
is that same space and has rank \(r\).  The spectrum of \(P\) consists of
\(r\) ones followed by zeros.  Under
\eqref{eq:V105-coincident-cold-condition}, the fractional Ky--Fan sum uses
only eigenvalues equal to one, and its normalized mean is one.
\end{proof}

\begin{corollary}[The overlap information can retain the quadratic
threshold]
\label{cor:V105-overlap-quadratic-threshold}
Since \(\eta_N\asymp N^{-2}\), the data in
\Cref{prop:V105-coincident-cold-regression} allow
\(\overline c_{J,N}=1\) whenever \(R\lesssim N^2\).  Thus the V104 linear
threshold is not valid on a common overlap bank without an additional
cold-space transversality theorem.
\end{corollary}

\begin{proof}
Condition \eqref{eq:V105-coincident-cold-condition} is
\(\eta_N\le r/d=1/R\).  This holds, up to fixed constants, for
\(R\lesssim N^2\).
\end{proof}

\begin{firewall}[Scope of the coincident-cold test]
\label{fire:V105-coincident-cold-scope}
The projection model proves insufficiency of two marginal rank bounds.  It
does not assert that every projection spectrum is realized by the physical
four-row Wilson operator \(J_{\rm comp}\).  A special Racah
transversality identity could still force product-rank decay; it must be
proved from the exact same-coordinate kernel.
\end{firewall}

\begin{lemma}[Abstract correlated cold-rank gate]
\label{lem:V105-correlated-cold-rank}
Let the eigenvalues of \(J\) be bounded by \(A_J\).  Suppose all but an
exceptional spectral subspace of rank \(r_J\) have eigenvalue at most
\(h_J\).  Then
\begin{equation}
 \overline c_{J,N}
 \le
 h_J+
 A_J\min\left\{
 1,\frac{r_J}{\eta_Nd_J}
 \right\}.
 \label{eq:V105-correlated-cold-rank}
\end{equation}
Consequently the correlated overlap branch decays when \(s_N^2A_J\) is
bounded,
\begin{equation}
 s_N^2h_J\longrightarrow0,
 \qquad
 \frac{r_J}{\eta_Nd_J}\longrightarrow0.
 \label{eq:V105-joint-sufficient-data}
\end{equation}
\end{lemma}

\begin{proof}
A top fractional spectral sum has baseline at most \(h_J\) on its whole
mass.  Only the smaller of \(r_J\) and \(\eta_Nd_J\) can receive the
additional cap \(A_J\).  Divide by \(\eta_Nd_J\) to obtain
\eqref{eq:V105-correlated-cold-rank}.  Insert the result in
\eqref{eq:V105-correlated-near-unit-gate}; the two conditions in
\eqref{eq:V105-joint-sufficient-data} make the right side tend to zero.
\end{proof}

\begin{definition}[Required complementary-pair transversality datum]
\label{def:V105-required-transversality}
For \(J_{\rm comp}\), a sufficient replacement for tensor separation is an
actual same-coordinate estimate
\begin{equation}
 \frac{r_{\rm comp}}{d_{\rm comp}}
 \le
 C_{\rm tr}
 \frac{\chi_{\rm comp}}{R_iR_k},
 \qquad
 h_{\rm comp}
 \le
 C_{\rm tr}
 \max\left\{\frac1{R_i},\frac1{R_k}\right\},
 \label{eq:V105-required-transversality}
\end{equation}
with \(C_{\rm tr}\), \(\chi_{\rm comp}\), and the coefficient cap uniformly
bounded.  Here \(r_{\rm comp}\) and \(h_{\rm comp}\) refer to the spectral
decomposition of the \emph{joint} operator, not to separately optimized
marginals.
\end{definition}

\begin{corollary}[What joint transversality would recover]
\label{cor:V105-transversality-recovers-linear}
If \eqref{eq:V105-required-transversality} holds, the complementary-overlap
near-unit channel satisfies the same sufficient scale as the disjoint
V104 channel:
\[
 R_i,R_k\to\infty,
 \qquad
 \frac{R_iR_k}{N^2}\to\infty.
\]
\end{corollary}

\begin{proof}
Apply \Cref{lem:V105-correlated-cold-rank} and
\(\eta_N\asymp N^{-2}\).
\end{proof}

\section{Complete structural fork after two endpoint certificates}
\label{sec:V105-complete-fork}

\begin{theorem}[Audited two-endpoint extraction theorem]
\label{thm:V105-audited-extraction}
Fix one complete V60 monomial with two sideways endpoints, and apply the V94
priority selector before taking a product-state supremum.  If the outcome is
not already assigned to a cut, mean, selected paid branch, or the separate
off-cut carrier, then its two V102 pair certificates have the following
exhaustive disposition:
\begin{enumerate}[label=\textup{(\roman*)}]
\item \textbf{disjoint banks:} the carrier contains
      \(Q_{G_i}\otimes Q_{G_k}\otimes M\), with both selected masses bounded
      below by \eqref{eq:V105-V102-disjoint-mass}; use the V104 two-pair
      profile;
\item \textbf{same-pair overlap:} one common pair bank is heavy for both
      endpoints; use the V103 one-pair profile once;
\item \textbf{complementary-pair overlap:} one same-coordinate four-row bank
      is heavy for both endpoints; resolve its reducible retained rows and
      use the fixed-channel joint profiles supplied by
      \eqref{eq:V105-correlated-overlap-core-gate}.
\end{enumerate}
No fourth puncture-incidence geometry is possible.
\end{theorem}

\begin{proof}
\Cref{cor:V105-V102-quarter} gives either disjoint banks or a pure overlap
type.  The disjoint factorization is
\Cref{prop:V105-disjoint-factorization}.  The two pure types are exhaustive
by \Cref{lem:V105-local-overlap-alphabet}, and their operator dispositions
are
\Cref{prop:V105-same-pair-one-operator,
prop:V105-complementary-joint-operator}.  The V94 priority hypothesis keeps
all previously paid and off-cut outcomes outside this fork.
\end{proof}

\begin{assessment}[Direction after the audited extraction]
\label{ass:V105-direction}
The combinatorial disjointness question is now closed.  The disjoint branch
has a rigorous V104 conditional estimate; the same-pair branch is provably
only a one-pair problem.  On the complementary branch, V56 and V58 already
supply the exact local double-defect factor and its legitimate
row-factorized trace-tax composition.  The new object is its decomposition
into irreducible channels after the two factors on each retained row are
made into an ordinary diagonal-group tensor product.  V106 should compute
that recoupling and its channelwise cold spectral space.  What remains
forbidden is replacing the channel-resolved product capacity by a product of
two independently maximized marginal capacities, as emphasized by
\Cref{prop:V105-coincident-cold-regression}.
\end{assessment}

\section{V105 ledger and V106 acquisition}
\label{sec:V105-ledger}

\begin{theorem}[Integrated V105 statement]
\label{thm:V105-integrated}
Preserving V1--V104, V105 proves:
\begin{enumerate}[label=\textup{(V105.\arabic*)},leftmargin=6.3em]
\item the digest-recorded SM V114/NS V31 companion audit and its strict
      transfer firewall;
\item the exact local same-pair/complementary-pair overlap alphabet;
\item the quarter-mass disjoint-or-pure-overlap theorem
      \eqref{eq:V105-quarter-disjoint}--%
      \eqref{eq:V105-quarter-pure-overlap};
\item genuine coordinate tensor factorization on the disjoint branch;
\item identification of same-pair overlap as one pair operator;
\item identification of complementary overlap as an exact pre-recoupling
      tensor factor with reducible diagonal-group rows;
\item the exact fixed-channel correlated-overlap Ky--Fan gate
      \eqref{eq:V105-correlated-overlap-core-gate};
\item the coincident-cold regression excluding a marginal-rank product; and
\item the exhaustive audited extraction fork of
      \Cref{thm:V105-audited-extraction}.
\end{enumerate}
\end{theorem}

\begin{proof}
The items are
\Cref{ass:V105-audit-transfer,
fire:V105-companion-separation,
lem:V105-local-overlap-alphabet,
thm:V105-quarter-dichotomy,
prop:V105-disjoint-factorization,
prop:V105-same-pair-one-operator,
prop:V105-complementary-joint-operator,
thm:V105-correlated-overlap-gate,
prop:V105-coincident-cold-regression,
thm:V105-audited-extraction}.
\end{proof}

\begin{openproblem}[V106 mandate: complementary diagonal-row resolution]
\label{op:V106-mandate}
Remain on one fixed \(O_{\rm comp}\) coordinate bank selected by
\Cref{thm:V105-audited-extraction}.
\begin{enumerate}
\item Start from the exact V56 factor
      \(\mathbf d_{ij,f}\otimes\mathbf d_{k\ell,f}\); do not replace it by an
      unspecified four-row moment.
\item Resolve the ordinary tensor product on each retained row into
      diagonal-group irreducibles; record all multiplicity spaces and Racah
      transforms.
\item Prove positivity of each complete squared channel effect and identify
      the exact channel trace and eigenvalue/multiplicity law.
\item Determine whether its large-eigenvalue space obeys the product-rank
      transversality target \eqref{eq:V105-required-transversality}, only a
      one-pair rank bound, or neither.
\item Test the common-pair singlet, complementary dual-singlet, Cartan, and
      alternating-orientation families.
\item If the joint kernel has a persistent top eigenspace, retain its signed
      coupling to the core column before taking the positive square.
\item Keep the V98 off-cut carrier outside this calculation.
\end{enumerate}
The next compulsory companion audit is V110.
\end{openproblem}

\section{Firewall after V105}
\label{sec:V105-firewall}

\begin{firewall}[Current claim boundary]
\label{fire:V105-current-boundary}
V105 does not resolve or sum the complementary diagonal-row channels, prove
the complementary-pair transversality estimate, bound the same-pair top
spectrum at the natural mass scale, control the complete
signed core scalar, close the disjoint branch without its effective-scale
and orientation hypotheses, close the off-cut branch, prove JREC, construct
the continuum theory, or prove a positive mass gap.  It gives a complete
incidence classification of two forced pair certificates and the exact
correlated capacity that remains on each branch.  The full Yang--Mills
existence and mass gap remain open.
\end{firewall}

\section{Append-only record for V105}
\label{sec:V105-record}

V105 was appended to the complete V104 manuscript.  The exact V104 parent
had \(71054\) physical lines, \(2477653\) bytes, and SHA--256 digest
\[
 \texttt{b30d5eab41246f519af4112da1d440941e36be80bde00041bd8ef1ecd8fee281}.
\]
No V104 mathematical content was changed.  The sole final
\verb|\end{document}| is restored below.

% =============================================================================
% END APPENDED VERSION V105 MATERIAL
% =============================================================================
% =============================================================================
% BEGIN APPENDED VERSION V106 MATERIAL
% =============================================================================

\chapter{V106: Reducible-Row Schur Law and Racah Marginal
Concentration}
\label{ch:V106}

\section{Purpose and correction of the overlap target}
\label{sec:V106-purpose}

V105's audited incidence theorem isolates a complementary \([2,2]\)
overlap bank.  The exact local object is not new: V56 already gives
\[
 \mathbf d_{ij,f}\otimes\mathbf d_{k\ell,f},
\]
and V58 proves that the two trace taxes compose on the row-factorized
pre-recoupling coefficient.  The missing issue is specific to the newer
product-capacity route.  Both pair factors belong to the same gauge
coordinate, so their tensor factors on one retained row carry an ordinary
diagonal-group tensor product.  That row is generally reducible.  A final
Racah recoupling or pinching need not preserve the scalar partial marginals
used in the V104 Schur estimate.

This version replaces the unavailable irreducible-row shortcut by an exact
partial-marginal law valid for every positive effect.  Before recoupling,
the two pair-level projectors have scalar marginals and recover the V104
formula even on a common coordinate.  After recoupling, the sole loss is a
pair of explicitly defined marginal concentration factors.  Their exact
isotypic formula shows what a future Racah calculation must bound.

\section{A Schur law without irreducible auxiliary rows}
\label{sec:V106-general-marginal-law}

Let \(\mathcal X,\mathcal Y\) be arbitrary finite-dimensional auxiliary
row spaces.  Let \(A,B\) be irreducible core representations, and let
\(M\ge0\) be an invariant core shell on \(A\otimes B\).  As before, put
\[
 \Pi(M):=\|M\|_{\rm prod},
 \qquad
 \mathcal H(M):=\operatorname{Tr}M.
\]
For an auxiliary effect \(0\le E\le I_{\mathcal X\otimes\mathcal Y}\),
define its two marginal heights
\begin{equation}
 \alpha(E):=
 \|\operatorname{Tr}_{\mathcal Y}E\|_{\rm op},
 \qquad
 \beta(E):=
 \|\operatorname{Tr}_{\mathcal X}E\|_{\rm op}.
 \label{eq:V106-marginal-heights}
\end{equation}

\begin{theorem}[Exact marginal product-capacity bound]
\label{thm:V106-marginal-product-cap}
For every auxiliary effect \(E\),
\begin{equation}
 \boxed{\;
 \|E\otimes M\|_{\rm prod}
 \le
 \min\left\{
 \Pi(M),\,
 \frac{\mathcal H(M)}{d_A}\alpha(E),\,
 \frac{\mathcal H(M)}{d_B}\beta(E)
 \right\}.
 \;}
 \label{eq:V106-marginal-product-cap}
\end{equation}
No irreducibility assumption is made on \(\mathcal X\) or \(\mathcal Y\).
\end{theorem}

\begin{proof}
The operator inequality \(0\le E\otimes M\le I\otimes M\), followed by
the density-matrix argument of
\Cref{lem:V104-two-projection-product-cap}, gives the first entry
\(\Pi(M)\).

For the second entry, partially trace the positive operator over the
complete right row:
\[
 \operatorname{Tr}_{\mathcal Y\otimes B}(E\otimes M)
 =
 \operatorname{Tr}_{\mathcal Y}E
 \otimes
 \operatorname{Tr}_{B}M.
\]
Core invariance and irreducibility of \(A\) give
\[
 \operatorname{Tr}_{B}M
 =
 \frac{\mathcal H(M)}{d_A}I_A.
\]
For any unit product vector \(u\otimes v\), its expectation against a
positive operator is no larger than the expectation of the corresponding
partial trace in \(u\), hence no larger than the operator norm of that
partial trace.  This proves the second entry.  Tracing the left row instead
gives the third.
\end{proof}

\begin{corollary}[Integrated marginal-level law]
\label{cor:V106-integrated-marginal-law}
Suppose a positive auxiliary operator has an effect resolution
\begin{equation}
 J=\int_{\mathcal T}E_\theta\,\dd\nu(\theta),
 \qquad
 0\le E_\theta\le I.
 \label{eq:V106-effect-resolution}
\end{equation}
Then
\begin{align}
 \|J\otimes M\|_{\rm prod}
 \le
 \int_{\mathcal T}
 \min\left\{
 \Pi(M),\,
 \frac{\mathcal H(M)}{d_A}\alpha(E_\theta),\,
 \frac{\mathcal H(M)}{d_B}\beta(E_\theta)
 \right\}\dd\nu(\theta).
 \label{eq:V106-integrated-marginal-law}
\end{align}
\end{corollary}

\begin{proof}
Use \eqref{eq:V106-effect-resolution}, subadditivity of the product
numerical radius on positive operators, and
\Cref{thm:V106-marginal-product-cap} pointwise.
\end{proof}

\begin{remark}[V104 is the scalar-marginal special case]
\label{rem:V106-V104-special-case}
If \(\mathcal X,\mathcal Y\) are irreducible row representations and
\(E\) is invariant, Schur's lemma gives
\[
 \alpha(E)=\frac{\operatorname{Tr}E}{\dim\mathcal X},
 \qquad
 \beta(E)=\frac{\operatorname{Tr}E}{\dim\mathcal Y}.
\]
Equation \eqref{eq:V106-marginal-product-cap} then becomes the V104 double
Schur cap.  Thus the new law strictly extends that argument rather than
replacing it.
\end{remark}

\section{The exact pre-recoupling complementary-pair level}
\label{sec:V106-pre-recoupling}

Let the two pair-level projectors be
\[
 P_e(t)\ \hbox{on }X_e\otimes Y_e,
 \qquad
 P_f(s)\ \hbox{on }X_f\otimes Y_f,
\]
where all four displayed representations are irreducible for the common
coordinate gauge group.  Put
\[
 D_e(t):=\operatorname{rank}P_e(t),
 \qquad
 D_f(s):=\operatorname{rank}P_f(s).
\]
Regroup the pre-recoupling tensor product by retained row:
\begin{equation}
 E^0_{t,s}
 :=
 P_e(t)\otimes P_f(s)
 \quad\hbox{on}\quad
 (X_e\otimes X_f)\otimes(Y_e\otimes Y_f).
 \label{eq:V106-prerecoupling-effect}
\end{equation}

\begin{lemma}[Each pair level has scalar one-row marginals]
\label{lem:V106-pair-scalar-marginals}
For every \(t,s\),
\begin{align}
 \operatorname{Tr}_{Y_e}P_e(t)
 &=
 \frac{D_e(t)}{d_{X_e}}I_{X_e},
 &
 \operatorname{Tr}_{X_e}P_e(t)
 &=
 \frac{D_e(t)}{d_{Y_e}}I_{Y_e},
 \label{eq:V106-e-scalar-marginals}\\
 \operatorname{Tr}_{Y_f}P_f(s)
 &=
 \frac{D_f(s)}{d_{X_f}}I_{X_f},
 &
 \operatorname{Tr}_{X_f}P_f(s)
 &=
 \frac{D_f(s)}{d_{Y_f}}I_{Y_f}.
 \label{eq:V106-f-scalar-marginals}
\end{align}
\end{lemma}

\begin{proof}
Each spectral projector commutes with the diagonal group action on its pair
carrier.  Its partial trace over one leg therefore commutes with the
irreducible representation on the other leg.  Schur's lemma makes the
partial trace scalar.  Its trace equals the rank of the projector, fixing
the displayed scalar.
\end{proof}

\begin{proposition}[Scalar marginals survive common-coordinate row
regrouping]
\label{prop:V106-prerecoupling-scalar}
Although \(X_e\otimes X_f\) and \(Y_e\otimes Y_f\) may be reducible for the
diagonal group,
\begin{align}
 \alpha(E^0_{t,s})
 &=
 \frac{D_e(t)D_f(s)}{d_{X_e}d_{X_f}},
 \label{eq:V106-prerecoupling-alpha}\\
 \beta(E^0_{t,s})
 &=
 \frac{D_e(t)D_f(s)}{d_{Y_e}d_{Y_f}}.
 \label{eq:V106-prerecoupling-beta}
\end{align}
Consequently the V104 two-level trace denominator is exact before final
recoupling, even when the two pair blocks occupy the same gauge coordinate.
\end{proposition}

\begin{proof}
Partial trace of a tensor product factorizes:
\begin{align*}
 \operatorname{Tr}_{Y_e\otimes Y_f}E^0_{t,s}
 &=
 \operatorname{Tr}_{Y_e}P_e(t)
 \otimes
 \operatorname{Tr}_{Y_f}P_f(s),\\
 \operatorname{Tr}_{X_e\otimes X_f}E^0_{t,s}
 &=
 \operatorname{Tr}_{X_e}P_e(t)
 \otimes
 \operatorname{Tr}_{X_f}P_f(s).
\end{align*}
Insert \eqref{eq:V106-e-scalar-marginals}--%
\eqref{eq:V106-f-scalar-marginals}.  The resulting operators are scalar
identities, and their norms are
\eqref{eq:V106-prerecoupling-alpha}--%
\eqref{eq:V106-prerecoupling-beta}.  Substitution in
\eqref{eq:V106-marginal-product-cap} gives the V104 denominator.
\end{proof}

\begin{assessment}[Correction of the V105 overlap description]
\label{ass:V106-V105-correction}
The complementary overlap is not obstructed merely because the two defects
share a coordinate.  V56 supplies their exact tensor factor, and
\Cref{prop:V106-prerecoupling-scalar} shows that reducibility alone does
not spoil its pre-recoupling marginals.  The only possible loss occurs when
the physical Racah recoupling and aggregate output resolution are applied.
The rest of V106 measures that loss exactly.
\end{assessment}

\section{Racah marginal concentration factors}
\label{sec:V106-concentration-factors}

Let \(\Phi\) be the positive, unital, trace-preserving map on four-row
effects generated by a common unitary recoupling followed by an orthogonal
pinching.  This abstractly records the V57 aggregate resolution while
separating the one datum relevant to product capacity.  Define
\begin{equation}
 E^\Phi_{t,s}:=\Phi(E^0_{t,s}).
 \label{eq:V106-recoupled-effect}
\end{equation}
Then
\[
 0\le E^\Phi_{t,s}\le I,
 \qquad
 \operatorname{Tr}E^\Phi_{t,s}=D_e(t)D_f(s).
\]
For nonzero \(D_e(t)D_f(s)\), define
\begin{align}
 \Gamma_{\mathsf a}(t,s)
 &:=
 \frac{d_{X_e}d_{X_f}}{D_e(t)D_f(s)}
 \alpha(E^\Phi_{t,s}),
 \label{eq:V106-Gamma-a}\\
 \Gamma_{\mathsf b}(t,s)
 &:=
 \frac{d_{Y_e}d_{Y_f}}{D_e(t)D_f(s)}
 \beta(E^\Phi_{t,s}).
 \label{eq:V106-Gamma-b}
\end{align}
Set either factor to zero when the corresponding level rank vanishes.

\begin{lemma}[Concentration factors are at least one]
\label{lem:V106-Gamma-lower}
At every nonzero level,
\begin{equation}
 \Gamma_{\mathsf a}(t,s)\ge1,
 \qquad
 \Gamma_{\mathsf b}(t,s)\ge1.
 \label{eq:V106-Gamma-lower}
\end{equation}
Equality holds before recoupling.
\end{lemma}

\begin{proof}
If \(Z\ge0\) acts on a \(d\)-dimensional space, then
\(\|Z\|_{\rm op}\ge\operatorname{Tr}Z/d\).  Apply this to each partial
trace.  Its trace is
\(\operatorname{Tr}E^\Phi_{t,s}=D_e(t)D_f(s)\), giving
\eqref{eq:V106-Gamma-lower}.  Equality before recoupling is
\Cref{prop:V106-prerecoupling-scalar}.
\end{proof}

\begin{theorem}[Exact recoupled two-level law]
\label{thm:V106-recoupled-two-level}
Put
\[
 J_\Phi:=\Phi(Q_e\otimes Q_f)
 =
 \int_0^{A_e}\!\!\int_0^{A_f}
 E^\Phi_{t,s}\,\dd s\,\dd t.
\]
Then
\begin{align}
 \|J_\Phi\otimes M\|_{\rm prod}
 &\le
 \int_0^{A_e}\!\!\int_0^{A_f}
 \min\Biggl\{
 \Pi(M),\,
 \frac{\mathcal H(M)D_e(t)D_f(s)
       \Gamma_{\mathsf a}(t,s)}
      {d_Ad_{X_e}d_{X_f}},
 \notag\\[-0.2em]
 &\hspace{12em}
 \frac{\mathcal H(M)D_e(t)D_f(s)
       \Gamma_{\mathsf b}(t,s)}
      {d_Bd_{Y_e}d_{Y_f}}
 \Biggr\}\dd s\,\dd t.
 \label{eq:V106-recoupled-two-level}
\end{align}
\end{theorem}

\begin{proof}
Linearity and the two pair layer resolutions give the displayed formula
for \(J_\Phi\).  Apply
\Cref{cor:V106-integrated-marginal-law} to the effects
\(E^\Phi_{t,s}\), and substitute
\eqref{eq:V106-Gamma-a}--\eqref{eq:V106-Gamma-b}.
\end{proof}

\begin{corollary}[Near-unit recoupled surface]
\label{cor:V106-near-unit-recoupled}
For \(M_N=s_N^2P_{S_N}\),
\begin{align}
 &\|J_\Phi\otimes M_N\|_{\rm prod}
 \notag\\
 &\quad\le
 s_N^2\int_0^{A_e}\!\!\int_0^{A_f}
 \min\Biggl\{
 \frac N{N+1},\,
 \frac{D_e(t)D_f(s)N(N+2)
       \Gamma_{\mathsf a}(t,s)}
      {3d_{X_e}d_{X_f}},
 \notag\\[-0.2em]
 &\hspace{15em}
 \frac{2ND_e(t)D_f(s)
       \Gamma_{\mathsf b}(t,s)}
      {(N+1)d_{Y_e}d_{Y_f}}
 \Biggr\}\dd s\,\dd t.
 \label{eq:V106-near-unit-recoupled}
\end{align}
Equivalently, relative to the core cap
\(\Pi_N=s_N^2N/(N+1)\), the integrand is
\begin{equation}
 \Pi_N
 \min\left\{
 1,\,
 \frac{D_e(t)D_f(s)\Gamma_{\mathsf a}(t,s)}
      {\eta_Nd_{X_e}d_{X_f}},\,
 \frac{2D_e(t)D_f(s)\Gamma_{\mathsf b}(t,s)}
      {d_{Y_e}d_{Y_f}}
 \right\}.
 \label{eq:V106-near-unit-relative}
\end{equation}
\end{corollary}

\begin{proof}
Use
\[
 \mathcal H_N=s_N^2N(N+2),
 \quad d_A=3,
 \quad d_B=\frac{(N+1)(N+2)}2
\]
in \eqref{eq:V106-recoupled-two-level}.  Division of the two trace entries
by \(\Pi_N\), together with
\(\eta_N=3/[(N+1)(N+2)]\), gives
\eqref{eq:V106-near-unit-relative}.
\end{proof}

\begin{corollary}[Bounded Racah concentration preserves the V104 gate]
\label{cor:V106-bounded-Gamma}
If
\begin{equation}
 \sup_{t,s,N}
 \max\{\Gamma_{\mathsf a}(t,s),\Gamma_{\mathsf b}(t,s)\}
 \le C_\Gamma<\infty,
 \label{eq:V106-bounded-Gamma}
\end{equation}
then \eqref{eq:V106-recoupled-two-level} is bounded by the V104
two-level surface with the two trace entries multiplied by
\(C_\Gamma\).  All V104 decay conclusions which have asymptotic room by a
fixed constant therefore remain valid.
\end{corollary}

\begin{proof}
Replace each concentration factor by \(C_\Gamma\) in
\eqref{eq:V106-recoupled-two-level}.  This changes no exponent or divergent
effective-scale condition.
\end{proof}

\begin{lemma}[Row-local bistochastic recoupling has no concentration loss]
\label{lem:V106-row-local-no-loss}
Suppose
\begin{equation}
 \Phi=\sum_{\rho}p_\rho
 \bigl(\Phi_{\mathsf a,\rho}\otimes
       \Phi_{\mathsf b,\rho}\bigr),
 \qquad
 p_\rho\ge0,\quad\sum_\rho p_\rho=1,
 \label{eq:V106-row-local-channel}
\end{equation}
where every factor map is positive, unital, and trace preserving on its
row.  Then
\begin{equation}
 \Gamma_{\mathsf a}(t,s)
 =
 \Gamma_{\mathsf b}(t,s)
 =1
 \label{eq:V106-row-local-Gamma-one}
\end{equation}
at every nonzero level.
\end{lemma}

\begin{proof}
For a product channel,
\[
 \operatorname{Tr}_{\mathcal Y}
 [(\Phi_{\mathsf a}\otimes\Phi_{\mathsf b})(E^0)]
 =
 \Phi_{\mathsf a}(\operatorname{Tr}_{\mathcal Y}E^0),
\]
because \(\Phi_{\mathsf b}\) is trace preserving.  The right side is the
same scalar identity as before recoupling, because
\(\Phi_{\mathsf a}\) is unital and
\Cref{prop:V106-prerecoupling-scalar} applies.  Convex averaging preserves
that scalar.  The other row is identical.  Insert the two marginal norms
in \eqref{eq:V106-Gamma-a}--\eqref{eq:V106-Gamma-b}.
\end{proof}

\begin{assessment}[The exact recoupling question]
\label{ass:V106-exact-recoupling-question}
The V57/V58 trace-norm contraction is sufficient for the marked-tree trace
ledger, but it does not by itself imply
\eqref{eq:V106-row-local-channel}.  For the product-capacity route one must
decide whether the actual Racah/output map is row-local, or else bound its
two marginal concentration factors.  No new four-row Haar integral is
needed before this question is answered.
\end{assessment}

\section{Exact isotypic formula for marginal concentration}
\label{sec:V106-isotypic-formula}

Assume \(E\) is invariant under the diagonal action of the common
coordinate group.  Decompose the left auxiliary row as
\begin{equation}
 \mathcal X
 \cong
 \bigoplus_{L\in\widehat G}
 \mathcal H_L\otimes\mathbb C^{m_L^{\mathsf a}},
 \qquad
 d_{\mathcal X}
 =
 \sum_Ld_Lm_L^{\mathsf a}.
 \label{eq:V106-left-isotypic}
\end{equation}
Only channels with positive multiplicity are included.  Invariance implies
\begin{equation}
 \operatorname{Tr}_{\mathcal Y}E
 =
 \bigoplus_L I_L\otimes A_L,
 \qquad
 A_L\ge0.
 \label{eq:V106-left-marginal-blocks}
\end{equation}
For \(D:=\operatorname{Tr}E>0\), define the isotypic trace law
\begin{equation}
 p_L:=
 \frac{d_L\operatorname{Tr}A_L}{D}.
 \label{eq:V106-left-sector-law}
\end{equation}

\begin{lemma}[The sector weights form a probability law]
\label{lem:V106-sector-probability}
One has
\begin{equation}
 p_L\ge0,
 \qquad
 \sum_Lp_L=1.
 \label{eq:V106-sector-probability}
\end{equation}
\end{lemma}

\begin{proof}
Positivity follows from \(A_L\ge0\).  Taking the trace in
\eqref{eq:V106-left-marginal-blocks} gives
\[
 D
 =
 \operatorname{Tr}E
 =
 \sum_Ld_L\operatorname{Tr}A_L.
\]
Divide by \(D\).
\end{proof}

\begin{theorem}[Exact isotypic concentration window]
\label{thm:V106-isotypic-concentration}
The left marginal height and concentration factor are
\begin{align}
 \alpha(E)&=\max_L\|A_L\|_{\rm op},
 \label{eq:V106-alpha-isotypic}\\
 \Gamma_{\mathsf a}(E)
 &:=
 \frac{d_{\mathcal X}}D\alpha(E).
 \label{eq:V106-Gamma-isotypic-definition}
\end{align}
They obey
\begin{equation}
 \boxed{\;
 \max\left\{
 1,\,
 d_{\mathcal X}
 \max_L\frac{p_L}{d_Lm_L^{\mathsf a}}
 \right\}
 \le
 \Gamma_{\mathsf a}(E)
 \le
 d_{\mathcal X}\max_L\frac{p_L}{d_L}.
 \;}
 \label{eq:V106-isotypic-window}
\end{equation}
If the row tensor product is multiplicity free, then
\begin{equation}
 \boxed{\;
 \Gamma_{\mathsf a}(E)
 =
 d_{\mathcal X}\max_L\frac{p_L}{d_L}.
 \;}
 \label{eq:V106-multiplicity-free-Gamma}
\end{equation}
The right row has the identical formula with its own irreducibles,
multiplicities, and sector law.
\end{theorem}

\begin{proof}
Equation \eqref{eq:V106-alpha-isotypic} is the norm of the block diagonal
operator \eqref{eq:V106-left-marginal-blocks}.  The average eigenvalue of
that operator is \(D/d_{\mathcal X}\), so
\(\Gamma_{\mathsf a}\ge1\).

On the multiplicity space of dimension \(m_L^{\mathsf a}\),
\[
 \frac{\operatorname{Tr}A_L}{m_L^{\mathsf a}}
 \le
 \|A_L\|_{\rm op}
 \le
 \operatorname{Tr}A_L.
\]
By \eqref{eq:V106-left-sector-law},
\[
 \operatorname{Tr}A_L=\frac{Dp_L}{d_L}.
\]
Take the maximum over \(L\), multiply by \(d_{\mathcal X}/D\), and combine
with the average lower bound.  If every multiplicity equals one, each
\(A_L\) is scalar and both spectral inequalities are equalities, proving
\eqref{eq:V106-multiplicity-free-Gamma}.
\end{proof}

\begin{corollary}[Exact bounded-concentration criterion]
\label{cor:V106-sector-nonconcentration}
A uniform estimate \(\Gamma_{\mathsf a}(E)\le C\) is equivalent to
\begin{equation}
 \|A_L\|_{\rm op}
 \le
 C\frac D{d_{\mathcal X}}
 \quad\hbox{for every \(L\)}.
 \label{eq:V106-block-nonconcentration}
\end{equation}
In the multiplicity-free case it is equivalent to
\begin{equation}
 p_L\le C\frac{d_L}{d_{\mathcal X}}
 \quad\hbox{for every \(L\)}.
 \label{eq:V106-sector-volume-law}
\end{equation}
\end{corollary}

\begin{proof}
The first equivalence is
\eqref{eq:V106-alpha-isotypic}--%
\eqref{eq:V106-Gamma-isotypic-definition}.  The second is the exact formula
\eqref{eq:V106-multiplicity-free-Gamma}.
\end{proof}

\begin{corollary}[\(F\otimes R_N\) Pieri rows have bounded concentration]
\label{cor:V106-Pieri-bounded-Gamma}
Let \(G=SU(3)\), \(F=V_{(1,0)}\), and \(R_N=V_{(N,0)}\).  Suppose the left
auxiliary row is
\[
 \mathcal X_N=F\otimes R_N
 =
 R_{N+1}\oplus S_N,
 \qquad
 S_N=V_{(N-1,1)},
\]
and \(E\) is invariant.  Then, for every \(N\ge1\),
\begin{equation}
 1\le\Gamma_{\mathsf a}(E)\le3.
 \label{eq:V106-Pieri-Gamma-bound}
\end{equation}
\end{corollary}

\begin{proof}
The decomposition is multiplicity free.  Its dimensions are
\begin{align*}
 d_{\mathcal X_N}
 &=3\frac{(N+1)(N+2)}2,\\
 d_{R_{N+1}}
 &=\frac{(N+2)(N+3)}2,\\
 d_{S_N}&=N(N+2).
\end{align*}
Since each sector probability is at most one,
\eqref{eq:V106-multiplicity-free-Gamma} gives
\[
 \Gamma_{\mathsf a}(E)
 \le
 \max\left\{
 \frac{3(N+1)}{N+3},
 \frac{3(N+1)}{2N}
 \right\}
 \le3.
\]
The lower bound is \eqref{eq:V106-Gamma-lower}.
\end{proof}

\begin{remark}[Scope of the Pieri test]
\label{rem:V106-Pieri-scope}
The corollary is unconditional for a resolved auxiliary row whose actual
representation is \(F\otimes R_N\).  The \(F\otimes R_N\) near-unit
\emph{core} does not imply that an arbitrary forced puncture pair has this
auxiliary row representation.  The V94 outcome labels must supply that
identification before \eqref{eq:V106-Pieri-Gamma-bound} is used.
\end{remark}

\section{Sharp recoupling tests}
\label{sec:V106-sharp-tests}

\begin{proposition}[A global unitary can inflate a scalar marginal by the
full row dimension]
\label{prop:V106-global-unitary-inflation}
Let \(\mathcal X=\mathcal Y=\mathbb C^d\), and let
\[
 \Omega_d=\frac1{\sqrt d}\sum_{j=1}^de_j\otimes e_j,
 \qquad
 E^0=|\Omega_d\rangle\langle\Omega_d|.
\]
There is a unitary conjugation \(\Phi\) for which
\begin{equation}
 \alpha(E^0)=\beta(E^0)=\frac1d,
 \qquad
 \alpha(\Phi(E^0))=\beta(\Phi(E^0))=1.
 \label{eq:V106-global-inflation}
\end{equation}
Thus both concentration factors increase from one to \(d\), although
\(\Phi\) is positive, unital, trace preserving, and rank preserving.
\end{proposition}

\begin{proof}
Choose a global unitary \(U\) carrying the unit vector \(\Omega_d\) to the
product vector \(e_1\otimes e_1\), and set
\(\Phi(Z)=UZU^*\).  The two marginals of \(E^0\) are \(I/d\), while the
marginals of
\[
 \Phi(E^0)=|e_1\otimes e_1\rangle\langle e_1\otimes e_1|
\]
are rank-one projections of norm one.  Since \(D=1\) and each row dimension
is \(d\), the definition of \(\Gamma\) gives the last assertion.
\end{proof}

\begin{firewall}[The global-unitary test is informational]
\label{fire:V106-global-unitary-scope}
The unitary in \Cref{prop:V106-global-unitary-inflation} is not asserted to
be an \(SU(3)\) Racah matrix or group-equivariant recoupling.  The example
proves that global unitarity, trace preservation, rank preservation, and
pinching contractivity alone do not control product capacity.  The actual
Racah coefficients or the isotypic sector law must be used.
\end{firewall}

\begin{proposition}[Double dual-singlet pre-recoupling benchmark]
\label{prop:V106-double-singlet-benchmark}
Let each pair level be its one-dimensional singlet projector in
\(U\otimes U^*\) and \(V\otimes V^*\).  Before final recoupling,
\begin{equation}
 \alpha(E^0)=\frac1{d_Ud_V},
 \qquad
 \beta(E^0)=\frac1{d_Ud_V},
 \qquad
 \Gamma_{\mathsf a}=\Gamma_{\mathsf b}=1.
 \label{eq:V106-double-singlet-marginals}
\end{equation}
Any loss on this cold--cold benchmark is therefore entirely a recoupling
concentration loss.
\end{proposition}

\begin{proof}
Each normalized singlet has maximally mixed one-leg marginal.  Tensor the
two marginal identities, or apply
\Cref{prop:V106-prerecoupling-scalar} with
\(D_e=D_f=1\).
\end{proof}

\begin{assessment}[What V106 has reduced the overlap problem to]
\label{ass:V106-reduction}
The complementary overlap no longer requires a conjectural independence
law.  Its pre-recoupling pair levels already have the desired product rank
and scalar row marginals.  To preserve the V104 linear threshold after the
physical output resolution, it is enough to prove uniform boundedness of
\(\Gamma_{\mathsf a},\Gamma_{\mathsf b}\), equivalently the isotypic
nonconcentration law \eqref{eq:V106-block-nonconcentration}.  Fundamental
Pieri rows pass automatically; general rows and multiplicity spaces remain
open.
\end{assessment}

\section{V106 ledger and V107 acquisition}
\label{sec:V106-ledger}

\begin{theorem}[Integrated V106 statement]
\label{thm:V106-integrated}
Preserving V1--V105, V106 proves:
\begin{enumerate}[label=\textup{(V106.\arabic*)},leftmargin=6.3em]
\item the exact auxiliary marginal product-capacity bound
      \eqref{eq:V106-marginal-product-cap};
\item scalar partial marginals for every invariant pair level;
\item recovery of the V104 two-level denominator before recoupling even on
      a same-coordinate complementary pair;
\item the exact post-recoupling concentration factors
      \eqref{eq:V106-Gamma-a}--\eqref{eq:V106-Gamma-b};
\item the recoupled near-unit surface
      \eqref{eq:V106-near-unit-recoupled};
\item absence of loss for row-local bistochastic recouplings;
\item the exact diagonal-group isotypic concentration window
      \eqref{eq:V106-isotypic-window};
\item the uniform bound \(\Gamma\le3\) on \(F\otimes R_N\) Pieri rows; and
\item global-unitary and double-singlet sharp tests.
\end{enumerate}
\end{theorem}

\begin{proof}
The items are
\Cref{thm:V106-marginal-product-cap,
lem:V106-pair-scalar-marginals,
prop:V106-prerecoupling-scalar,
thm:V106-recoupled-two-level,
cor:V106-near-unit-recoupled,
lem:V106-row-local-no-loss,
thm:V106-isotypic-concentration,
cor:V106-Pieri-bounded-Gamma,
prop:V106-global-unitary-inflation,
prop:V106-double-singlet-benchmark}.
\end{proof}

\begin{openproblem}[V107 mandate: physical Racah marginal law]
\label{op:V107-mandate}
Return to the exact V57 aggregate collision projector and the V70/V76 Racah
POVM.
\begin{enumerate}
\item Write the controlled Racah/output channel \(\Phi\) on the
      multiplicity-resolved pair labels \((\lambda,\mu)\); distinguish
      row-local basis changes from genuinely cross-row recoupling.
\item For every left and right diagonal-group irrep, compute the matrices
      \(A_L(t,s)\) in \eqref{eq:V106-left-marginal-blocks} as quadratic
      forms in the Racah coefficients.
\item Use row and column unitarity before any entrywise estimate and decide
      whether \(\Gamma_{\mathsf a},\Gamma_{\mathsf b}\) are uniformly
      bounded.
\item Separate multiplicity-free Pieri rows, already covered by
      \Cref{cor:V106-Pieri-bounded-Gamma}, from genuine multiplicity
      spaces.
\item Test double dual singlets, the \(F\otimes R_N\) near-unit path,
      Cartan outputs, and alternating row orientation.
\item If a concentration factor grows, retain its exact sector label with
      the signed core column before taking the final square.
\item Keep the same-pair and off-cut branches separate.
\end{enumerate}
The next compulsory companion audit remains V110.
\end{openproblem}

\section{Firewall after V106}
\label{sec:V106-firewall}

\begin{firewall}[Current claim boundary]
\label{fire:V106-current-boundary}
V106 does not compute the physical Racah marginal sector law in general,
prove bounded concentration on arbitrary multiplicity spaces, close the
same-pair branch, force the V104 effective-scale hypothesis, control the
complete signed core scalar, close the off-cut branch, prove JREC, construct
the continuum theory, or prove a positive mass gap.  It corrects the
same-coordinate overlap description and reduces the post-recoupling loss to
two exact, representation-theoretic marginal concentrations.  The full
Yang--Mills existence and mass gap remain open.
\end{firewall}

\section{Append-only record for V106}
\label{sec:V106-record}

V106 was appended to the complete corrected V105 manuscript.  The exact
V105 parent had \(71860\) physical lines, \(2509158\) bytes, and SHA--256
digest
\[
 \texttt{baa212f934dd60b93953c691d409a93b9c0e55233f3f5b19d98ec6d1d4347977}.
\]
No V105 mathematical content was changed after this append began.  The sole
final \verb|\end{document}| is restored below.

% =============================================================================
% END APPENDED VERSION V106 MATERIAL
% =============================================================================
% =============================================================================
% BEGIN APPENDED VERSION V107 MATERIAL
% =============================================================================

\chapter{V107: Double-Singlet Entanglement Swapping and Exact
Sector-Uniform Marginals}
\label{ch:V107}

\section{Purpose}
\label{sec:V107-purpose}

V106 reduced the complementary-pair product-capacity problem to the partial
marginals of the physical Racah/output effect.  This chapter solves the
most dangerous cold--cold benchmark exactly: both pair defects resolve to
their multiplicity-one dual singlets.

After regrouping by retained row, the product of the two singlets is the
canonical maximally entangled vector between \(X=U\otimes V\) and \(X^*\).
Diagonal-group row-channel resolution decomposes this vector into its
isotypic canonical pieces.  Every coherent or pinched collection of those
pieces has partial-marginal operator norm exactly \(1/\dim X\), independent
of the number and dimensions of the selected channels.  Consequently the
common product capacity has no Racah channel-count loss.  When
\(U=F\) and \(V=R_N\), coupling this auxiliary effect to the V103 near-unit
core gives a sharp explicit \(O(N^{-2})\) upper bound, up to the two singlet
coefficients and the core scalar.

\section{Canonical reshuffling of two dual singlets}
\label{sec:V107-canonical-reshuffling}

Let \(U,V\) be finite-dimensional unitary representations of a compact
group \(G\), with dimensions \(d_U,d_V\), and put
\begin{equation}
 X:=U\otimes V,
 \qquad
 d_X=d_Ud_V.
 \label{eq:V107-X-definition}
\end{equation}
Choose orthonormal bases and normalized canonical singlets
\begin{align}
 \Omega_U&:=
 \frac1{\sqrt{d_U}}\sum_{a=1}^{d_U}u_a\otimes\overline{u_a}
 \in U\otimes U^*,
 \label{eq:V107-Omega-U}\\
 \Omega_V&:=
 \frac1{\sqrt{d_V}}\sum_{b=1}^{d_V}v_b\otimes\overline{v_b}
 \in V\otimes V^*.
 \label{eq:V107-Omega-V}
\end{align}
Regroup the four factors from
\((U\otimes U^*)\otimes(V\otimes V^*)\) to
\((U\otimes V)\otimes(U^*\otimes V^*)=X\otimes X^*\).

\begin{lemma}[Exact entanglement-swapping identity]
\label{lem:V107-entanglement-swapping}
Under this row regrouping,
\begin{equation}
 \boxed{\;
 \Omega_U\otimes\Omega_V
 =
 \Omega_X
 :=
 \frac1{\sqrt{d_X}}
 \sum_{a,b}
 (u_a\otimes v_b)\otimes
 (\overline{u_a}\otimes\overline{v_b}).
 \;}
 \label{eq:V107-entanglement-swapping}
\end{equation}
In particular, the two-pair cold--cold vector is maximally entangled across
the retained-row cut.
\end{lemma}

\begin{proof}
Expand the left side using
\eqref{eq:V107-Omega-U}--\eqref{eq:V107-Omega-V}, commute the middle
\(U^*\) and \(V\) tensor factors by the canonical row-regrouping unitary,
and use \(d_X=d_Ud_V\).  The result is exactly the right side.
\end{proof}

\begin{corollary}[Flat full-row marginals]
\label{cor:V107-flat-full-marginals}
Let \(E_X:=|\Omega_X\rangle\langle\Omega_X|\).  Then
\begin{equation}
 \operatorname{Tr}_{X^*}E_X=\frac1{d_X}I_X,
 \qquad
 \operatorname{Tr}_{X}E_X=\frac1{d_X}I_{X^*},
 \qquad
 \|E_X\|_{\rm prod}=\frac1{d_X}.
 \label{eq:V107-flat-full-marginals}
\end{equation}
\end{corollary}

\begin{proof}
The first two identities are the standard partial traces of the displayed
Schmidt decomposition in \eqref{eq:V107-entanglement-swapping}.  For unit
vectors \(x\in X,y\in X^*\),
\[
 |\langle x\otimes y,\Omega_X\rangle|^2
 \le\frac1{d_X},
\]
with equality for a conjugate pair of basis vectors.
\end{proof}

\section{Exact diagonal-row isotypic resolution}
\label{sec:V107-isotypic-resolution}

Choose an orthogonal isotypic decomposition
\begin{equation}
 X
 =
 \bigoplus_{L\in\mathcal I}
 X_L,
 \qquad
 X_L\cong H_L\otimes\mathbb C^{m_L},
 \qquad
 n_L:=\dim X_L=d_Lm_L.
 \label{eq:V107-X-isotypic}
\end{equation}
Let \(P_L\) be the projection onto \(X_L\), and let
\(\overline P_L\) be the conjugate projection on \(X^*\).  Define the
unnormalized channel vectors
\begin{equation}
 \psi_L:=(P_L\otimes\overline P_L)\Omega_X.
 \label{eq:V107-psi-L}
\end{equation}

\begin{lemma}[Orthogonal canonical channel decomposition]
\label{lem:V107-canonical-channel-decomposition}
The vectors \(\psi_L\) are mutually orthogonal and
\begin{align}
 \Omega_X&=\sum_{L\in\mathcal I}\psi_L,
 \label{eq:V107-Omega-sector-sum}\\
 \|\psi_L\|^2&=\frac{n_L}{d_X},
 \label{eq:V107-sector-weight}\\
 \operatorname{Tr}_{X^*}|\psi_L\rangle\langle\psi_L|
 &=\frac1{d_X}P_L,
 \label{eq:V107-sector-left-marginal}\\
 \operatorname{Tr}_{X}|\psi_L\rangle\langle\psi_L|
 &=\frac1{d_X}\overline P_L.
 \label{eq:V107-sector-right-marginal}
\end{align}
\end{lemma}

\begin{proof}
Choose an orthonormal basis adapted to
\eqref{eq:V107-X-isotypic}.  The canonical vector is
\[
 \Omega_X
 =
 \frac1{\sqrt{d_X}}
 \sum_{L,\rho}e_{L,\rho}\otimes\overline{e_{L,\rho}},
\]
where \(1\le\rho\le n_L\).  Projection onto one value of \(L\) gives
\(\psi_L\).  Disjoint basis supports prove orthogonality and the sum.
Counting its \(n_L\) equal coefficients proves
\eqref{eq:V107-sector-weight}; direct partial trace proves
\eqref{eq:V107-sector-left-marginal}--%
\eqref{eq:V107-sector-right-marginal}.
\end{proof}

For a nonempty channel set \(\mathcal A\subseteq\mathcal I\), put
\begin{align}
 P_{\mathcal A}&:=\sum_{L\in\mathcal A}P_L,
 &
 n_{\mathcal A}&:=\operatorname{rank}P_{\mathcal A}
 =\sum_{L\in\mathcal A}n_L,
 \label{eq:V107-channel-set-data}\\
 \psi_{\mathcal A}&:=\sum_{L\in\mathcal A}\psi_L,
 &
 E_{\mathcal A}^{\rm coh}
 &:=|\psi_{\mathcal A}\rangle\langle\psi_{\mathcal A}|,
 \label{eq:V107-coherent-effect}\\
 E_{\mathcal A}^{\rm pin}
 &:=\sum_{L\in\mathcal A}
 |\psi_L\rangle\langle\psi_L|.
 \label{eq:V107-pinched-effect}
\end{align}

\begin{theorem}[Sector-uniform marginal and product-capacity law]
\label{thm:V107-sector-uniform-law}
For either
\[
 E_{\mathcal A}\in
 \{E_{\mathcal A}^{\rm coh},E_{\mathcal A}^{\rm pin}\},
\]
one has
\begin{align}
 \operatorname{Tr}E_{\mathcal A}
 &=\frac{n_{\mathcal A}}{d_X},
 \label{eq:V107-subset-trace}\\
 \operatorname{Tr}_{X^*}E_{\mathcal A}
 &=\frac1{d_X}P_{\mathcal A},
 \label{eq:V107-subset-left-marginal}\\
 \operatorname{Tr}_{X}E_{\mathcal A}
 &=\frac1{d_X}\overline P_{\mathcal A},
 \label{eq:V107-subset-right-marginal}\\
 \boxed{\quad
 \|E_{\mathcal A}\|_{\rm prod}
 =\frac1{d_X}.
 \quad}
 \label{eq:V107-subset-product-cap}
\end{align}
Thus coherent selection and isotypic pinching have the same exact
partial-marginal heights and the same product capacity, independent of the
number of selected channels.
\end{theorem}

\begin{proof}
For the coherent effect, \(\psi_{\mathcal A}\) is
\(d_X^{-1/2}\) times the canonical unnormalized vector on the
\(n_{\mathcal A}\)-dimensional subspace \(P_{\mathcal A}X\).  Its squared
norm and two partial traces are therefore
\eqref{eq:V107-subset-trace}--%
\eqref{eq:V107-subset-right-marginal}.  The same formulas for the pinched
effect follow by summing
\eqref{eq:V107-sector-weight}--%
\eqref{eq:V107-sector-right-marginal}.

For unit \(x\in X\), \(y\in X^*\), the coherent expectation is at most
\(1/d_X\) by the Schmidt coefficients of
\(\psi_{\mathcal A}\), and equality holds on a conjugate basis pair inside
\(P_{\mathcal A}X\).  For the pinched effect,
\[
 \langle x\otimes y,E_{\mathcal A}^{\rm pin}x\otimes y\rangle
 =
 \frac1{d_X}\sum_{L\in\mathcal A}
 |\langle P_Lx,\overline P_L\overline y\rangle|^2.
\]
Each summand is at most
\(\|P_Lx\|^2\|\overline P_L\overline y\|^2\), whose sum is at most one
by Cauchy--Schwarz and orthogonality.  Equality again holds when both
vectors occupy one conjugate basis line in one selected sector.
\end{proof}

\begin{corollary}[No isotypic channel-count tax]
\label{cor:V107-no-channel-count}
For every nonempty \(\mathcal A\),
\[
 \left\|
 \sum_{L\in\mathcal A}
 |\psi_L\rangle\langle\psi_L|
 \right\|_{\rm prod}
 =
 \frac1{d_X}.
\]
Estimating the summands separately and adding their individual product
capacities would instead give
\[
 \frac{|\mathcal A|}{d_X}.
\]
The latter factor is an artefact of taking independent product-state
suprema after the channel split.
\end{corollary}

\begin{proof}
The exact common capacity is
\eqref{eq:V107-subset-product-cap}.  Applying the same equation to each
singleton set \(\{L\}\) and summing gives the displayed loose bound.
\end{proof}

\begin{corollary}[Exact sector concentration and its compensating weight]
\label{cor:V107-sector-Gamma}
For either effect in
\Cref{thm:V107-sector-uniform-law}, the V106 left and right concentration
factors are
\begin{equation}
 \Gamma_{\mathsf a}(\mathcal A)
 =
 \Gamma_{\mathsf b}(\mathcal A)
 =
 \frac{d_X}{n_{\mathcal A}}.
 \label{eq:V107-sector-Gamma}
\end{equation}
The product of trace and concentration is nevertheless constant:
\begin{equation}
 \operatorname{Tr}E_{\mathcal A}\,
 \Gamma_{\mathsf a}(\mathcal A)
 =
 \operatorname{Tr}E_{\mathcal A}\,
 \Gamma_{\mathsf b}(\mathcal A)
 =1.
 \label{eq:V107-trace-Gamma-compensation}
\end{equation}
\end{corollary}

\begin{proof}
Use
\(\alpha=\beta=1/d_X\) and
\(\operatorname{Tr}E_{\mathcal A}=n_{\mathcal A}/d_X\) in the
definition \eqref{eq:V106-Gamma-isotypic-definition}; the right row is
identical.  Multiplication gives the second formula.
\end{proof}

\begin{remark}[Bounded \(\Gamma\) is sufficient, not necessary]
\label{rem:V107-Gamma-not-necessary}
A small sector may have \(d_X/n_{\mathcal A}\gg1\), but its exact sector
weight is \(n_{\mathcal A}/d_X\).  The two factors cancel before the
product-state supremum.  Thus the bounded-concentration criterion of V106
is a convenient sufficient condition for general effects; the invariant
quantity is the absolute marginal height.  The double-singlet branch must
retain its dimension-law sector weight rather than normalize each channel
to mass one.
\end{remark}

\begin{proposition}[Conditional normalization creates a fictitious loss]
\label{prop:V107-normalization-trap}
For one sector \(L\), let
\[
 \widehat E_L
 :=
 \frac{d_X}{n_L}|\psi_L\rangle\langle\psi_L|.
\]
Then \(\widehat E_L\) is a rank-one projector and
\[
 \|\widehat E_L\|_{\rm prod}=\frac1{n_L}.
\]
Multiplying by its physical sector probability \(n_L/d_X\) restores
\[
 \frac{n_L}{d_X}\|\widehat E_L\|_{\rm prod}
 =\frac1{d_X}.
 \label{eq:V107-normalization-restoration}
\]
Hence replacing the unnormalized Racah atom by its conditional state and
forgetting its probability can lose the factor \(d_X/n_L\).
\end{proposition}

\begin{proof}
The normalized vector \(\sqrt{d_X/n_L}\,\psi_L\) is maximally entangled on
the \(n_L\)-dimensional isotypic subspace.  Its largest squared Schmidt
coefficient is \(1/n_L\).  The last identity is arithmetic.
\end{proof}

\section{Weighted double-singlet carrier}
\label{sec:V107-weighted-double-singlet}

Suppose the \(e\)-pair operator is supported on the unique singlet in
\(U\otimes U^*\) with eigenvalue \(a\ge0\), and the \(f\)-pair operator is
supported on the unique singlet in \(V\otimes V^*\) with eigenvalue
\(b\ge0\).  Their pre-recoupling tensor product, regrouped by retained row,
is
\begin{equation}
 Q_e\otimes Q_f
 =
 ab\,E_X.
 \label{eq:V107-weighted-double-singlet}
\end{equation}
Resolving a coherent channel set or pinching its isotypic labels replaces
\(E_X\) by the corresponding \(E_{\mathcal A}\).

\begin{theorem}[Exact double-singlet/core marginal gate]
\label{thm:V107-double-singlet-core-gate}
For every nonempty channel set \(\mathcal A\) and every invariant positive
core shell \(M\) on irreducible rows \(A,B\),
\begin{equation}
 \boxed{\;
 \|ab\,E_{\mathcal A}\otimes M\|_{\rm prod}
 \le
 ab\min\left\{
 \Pi(M),\,
 \frac{\mathcal H(M)}{d_Ad_X},\,
 \frac{\mathcal H(M)}{d_Bd_X}
 \right\}.
 \;}
 \label{eq:V107-double-singlet-core-gate}
\end{equation}
The estimate is independent of \(|\mathcal A|\), the individual sector
dimensions, and their multiplicities.
\end{theorem}

\begin{proof}
Apply \Cref{thm:V106-marginal-product-cap} to \(E_{\mathcal A}\) and use
\eqref{eq:V107-subset-left-marginal}--%
\eqref{eq:V107-subset-right-marginal}.  Multiply the result by \(ab\).
\end{proof}

\begin{corollary}[Explicit \(F\otimes R_N\) double-singlet decay]
\label{cor:V107-near-unit-double-singlet}
Take
\[
 U=F,\qquad V=R_N,
\qquad
 d_X=3d_{R_N}
 =\frac{3(N+1)(N+2)}2,
\]
and couple the auxiliary double-singlet effect to the near-unit core
\(M_N=s_N^2P_{S_N}\) on \(F\otimes R_N\).  For every nonempty coherent or
pinched set of diagonal-row channels,
\begin{equation}
 \boxed{\;
 \|ab\,E_{\mathcal A}\otimes M_N\|_{\rm prod}
 \le
 ab\,s_N^2
 \frac{4N}{3(N+1)^2(N+2)}
 \le
 \frac{4ab\,s_N^2}{3N^2}.
 \;}
 \label{eq:V107-near-unit-double-singlet}
\end{equation}
In particular this channel decays when \(ab\,s_N^2=O(1)\).
\end{corollary}

\begin{proof}
The trace of the near-unit core is
\(\mathcal H_N=s_N^2N(N+2)\), and its larger row has
\[
 d_B=d_{R_N}=\frac{(N+1)(N+2)}2.
\]
The smaller trace entry in
\eqref{eq:V107-double-singlet-core-gate} is therefore
\begin{align*}
 ab\frac{\mathcal H_N}{d_Bd_X}
 &=
 ab\frac{s_N^2N(N+2)}{3d_{R_N}^2}\\
 &=
 ab\,s_N^2
 \frac{4N}{3(N+1)^2(N+2)}.
\end{align*}
For \(N\ge1\),
\[
 N^3\le(N+1)^2(N+2),
\]
which gives the second inequality.
\end{proof}

\begin{assessment}[Disposition of the cold--cold singlet regression]
\label{ass:V107-cold-cold-disposition}
The simultaneous cold-singlet test in V56--V58 is harmless for the present
near-unit product-capacity family once its complete dimension-law sector
weight is retained: it gains \(N^{-2}\) from the larger retained-row
marginal.  This does not close arbitrary cold outputs.  It removes the
specific possibility that two multiplicity-one dual singlets alone create
the persistent top eigenspace allowed by the abstract V105 coincident-cold
model.
\end{assessment}

\begin{firewall}[Required physical identification]
\label{fire:V107-physical-identification}
Equation \eqref{eq:V107-near-unit-double-singlet} applies only when the two
forced pair outputs are the actual dual singlets in
\(F\otimes F^*\) and \(R_N\otimes R_N^*\), the four carriers regroup as
\((F\otimes R_N)\mid(F^*\otimes R_N^*)\), and the output operation is a
coherent isotypic restriction or orthogonal isotypic pinching.  A general
Racah POVM atom may act nontrivially inside multiplicity spaces and must be
analyzed separately.
\end{firewall}

\section{V107 ledger and V108 acquisition}
\label{sec:V107-ledger}

\begin{theorem}[Integrated V107 statement]
\label{thm:V107-integrated}
Preserving V1--V106, V107 proves:
\begin{enumerate}[label=\textup{(V107.\arabic*)},leftmargin=6.3em]
\item the exact double-singlet entanglement-swapping identity
      \eqref{eq:V107-entanglement-swapping};
\item the canonical diagonal-row isotypic decomposition
      \eqref{eq:V107-Omega-sector-sum};
\item sector weights \(n_L/d_X\) and exact flat partial marginals;
\item the common coherent/pinched product capacity \(1/d_X\);
\item absence of an isotypic channel-count tax;
\item exact trace--concentration compensation
      \eqref{eq:V107-trace-Gamma-compensation};
\item the conditional-normalization regression
      \eqref{eq:V107-normalization-restoration}; and
\item the \(O(N^{-2})\) near-unit double-singlet bound
      \eqref{eq:V107-near-unit-double-singlet}.
\end{enumerate}
\end{theorem}

\begin{proof}
The items are
\Cref{lem:V107-entanglement-swapping,
lem:V107-canonical-channel-decomposition,
thm:V107-sector-uniform-law,
cor:V107-no-channel-count,
cor:V107-sector-Gamma,
prop:V107-normalization-trap,
thm:V107-double-singlet-core-gate,
cor:V107-near-unit-double-singlet}.
\end{proof}

\begin{openproblem}[V108 mandate: non-singlet complementary pair levels]
\label{op:V108-mandate}
Keep the exact V56 pre-recoupling tensor factor and move beyond the
multiplicity-one double singlet.
\begin{enumerate}
\item For pair output projectors \(P_\lambda^{U,U'}\) and
      \(P_\mu^{V,V'}\), compute the Schmidt spectrum after the retained-row
      reshuffle \((U\otimes V)\mid(U'\otimes V')\).
\item Resolve diagonal-row isotypic channels without normalizing away their
      physical trace weights.
\item Express the two partial marginals as positive quadratic forms in the
      relevant \(6j\) coefficients, summing a common channel set before the
      product-state supremum.
\item Prove a sector-volume or marginal-height bound for
      multiplicity-free Pieri pair outputs.
\item Test highest-weight Cartan pair outputs, one-singlet/one-Cartan
      mixtures, and the V104 alternating-orientation family.
\item If an individual Racah atom concentrates, determine whether summing
      its complete V70/V76 POVM family restores a flat marginal as in
      \Cref{thm:V107-sector-uniform-law}.
\item Keep the same-pair, signed-core, and off-cut acquisitions separate.
\end{enumerate}
The next compulsory companion audit remains V110.
\end{openproblem}

\section{Firewall after V107}
\label{sec:V107-firewall}

\begin{firewall}[Current claim boundary]
\label{fire:V107-current-boundary}
V107 closes only the complementary double-dual-singlet near-unit channel
under the stated physical identification.  It does not control non-singlet
pair levels, arbitrary Racah multiplicities, the same-pair branch, the
complete signed core scalar, the V104 effective-scale hypotheses, the
off-cut branch, JREC, the continuum construction, or a positive mass gap.
The full Yang--Mills existence and mass gap remain open.
\end{firewall}

\section{Append-only record for V107}
\label{sec:V107-record}

V107 was appended to the complete V106 manuscript.  The exact V106 parent
had \(72686\) physical lines, \(2534914\) bytes, and SHA--256 digest
\[
 \texttt{7c7b63bafab44ad42364b7f428a122d199df1c2bbf89983b6eb5600166176e43}.
\]
No V106 mathematical content was changed.  The sole final
\verb|\end{document}| is restored below.

% =============================================================================
% END APPENDED VERSION V107 MATERIAL
% =============================================================================
% =============================================================================
% BEGIN APPENDED VERSION V108 MATERIAL
% =============================================================================

\chapter{V108: Pair-First Racah Subeffects and No Downstream
Marginal Inflation}
\label{ch:V108}

\section{Purpose and correction of the physical recoupling model}
\label{sec:V108-purpose}

V106 introduced marginal concentration factors for an abstract active
channel \(\Phi\) applied to a pair-level effect.  Its global-unitary example
correctly shows that trace and rank alone do not protect product capacity.
The physical V70/V76 construction has more structure: the Racah atoms form
a normalized POVM, and V57 requires the two pair projectors to be inserted
before the final aggregate output resolution.  A downstream outcome is
therefore a positive \emph{subeffect} of the assembled upstream carrier.
It is not obtained by actively conjugating that carrier and discarding the
comparison with its parent.

Positive order immediately controls both partial marginals.  Combining this
fact with the V106 pre-recoupling scalar-marginal theorem shows that no
downstream Racah event or event set can worsen the V104 two-pair Schur
surface in a fixed orientation.  Consequently the complementary
same-coordinate \([2,2]\) branch inherits the V104 conditional linear
effective-size threshold.  Noncommuting retained-row orientations remain a
separate V94 problem.

\section{A general downstream-subeffect theorem}
\label{sec:V108-subeffect-theorem}

Let \(T\ge0\) act on a bipartite row space
\(\mathcal H_{\mathsf a}\otimes\mathcal H_{\mathsf b}\).  Let
\(\{M_\omega\}_{\omega\in\Omega}\) be a finite POVM on its support:
\begin{equation}
 M_\omega\ge0,
 \qquad
 \sum_{\omega\in\Omega}M_\omega=I.
 \label{eq:V108-POVM}
\end{equation}
Define the pair-first outcome effects
\begin{equation}
 F_\omega:=T^{1/2}M_\omega T^{1/2}.
 \label{eq:V108-pair-first-effect}
\end{equation}

\begin{lemma}[Every event is a subeffect of its upstream carrier]
\label{lem:V108-event-subeffect}
For every event \(E\subseteq\Omega\),
\begin{equation}
 F_E:=\sum_{\omega\in E}F_\omega
 \quad\hbox{satisfies}\quad
 0\le F_E\le T,
 \qquad
 \sum_{\omega\in\Omega}F_\omega=T.
 \label{eq:V108-event-subeffect}
\end{equation}
\end{lemma}

\begin{proof}
The POVM identity gives
\[
 0\le M_E:=\sum_{\omega\in E}M_\omega\le I.
\]
Congruence by \(T^{1/2}\) preserves positive order:
\[
 0\le T^{1/2}M_ET^{1/2}\le T.
\]
This is \(F_E\).  Taking \(E=\Omega\) gives the sum identity.
\end{proof}

\begin{theorem}[No downstream product or marginal inflation]
\label{thm:V108-no-downstream-inflation}
For every event \(E\subseteq\Omega\),
\begin{align}
 \|F_E\|_{\rm prod}
 &\le\|T\|_{\rm prod},
 \label{eq:V108-product-monotonicity}\\
 \|\operatorname{Tr}_{\mathsf b}F_E\|_{\rm op}
 &\le
 \|\operatorname{Tr}_{\mathsf b}T\|_{\rm op},
 \label{eq:V108-left-marginal-monotonicity}\\
 \|\operatorname{Tr}_{\mathsf a}F_E\|_{\rm op}
 &\le
 \|\operatorname{Tr}_{\mathsf a}T\|_{\rm op}.
 \label{eq:V108-right-marginal-monotonicity}
\end{align}
\end{theorem}

\begin{proof}
Equation \eqref{eq:V108-event-subeffect} and positivity show that every
product-vector expectation of \(F_E\) is at most the corresponding
expectation of \(T\), proving
\eqref{eq:V108-product-monotonicity}.  Partial trace is a positive map, so
\[
 0\le\operatorname{Tr}_{\mathsf b}F_E
 \le\operatorname{Tr}_{\mathsf b}T.
\]
Operator norm is monotone on positive operators, proving
\eqref{eq:V108-left-marginal-monotonicity}.  The other row is identical.
\end{proof}

Let \(h:\Omega\to[0,A]\) be a nonnegative resolved profile, and put
\begin{align}
 K_h&:=\sum_{\omega}h(\omega)F_\omega,
 \label{eq:V108-weighted-effect}\\
 E_r&:=\{\omega:h(\omega)>r\},
 \qquad 0\le r\le A.
 \label{eq:V108-profile-events}
\end{align}

\begin{corollary}[Exact downstream layer cake]
\label{cor:V108-downstream-layer-cake}
One has
\begin{equation}
 K_h=\int_0^AF_{E_r}\,\dd r,
 \qquad
 0\le F_{E_r}\le T,
 \label{eq:V108-downstream-layer-cake}
\end{equation}
and hence
\begin{equation}
 \|K_h\|_{\rm prod}
 \le
 \int_0^A\|F_{E_r}\|_{\rm prod}\,\dd r
 \le A\|T\|_{\rm prod}.
 \label{eq:V108-weighted-product-cap}
\end{equation}
The sharper bounds obtained from either upstream partial marginal may be
inserted at every \(r\) before integration.
\end{corollary}

\begin{proof}
Use
\[
 h(\omega)=\int_0^A
 \mathbf1_{\{h(\omega)>r\}}\,\dd r
\]
and interchange the finite sum and integral.  Then apply
\Cref{lem:V108-event-subeffect,thm:V108-no-downstream-inflation}.
\end{proof}

\begin{firewall}[Normalization destroys the order comparison]
\label{fire:V108-no-conditional-normalization}
The conditional effect \(F_E/\operatorname{Tr}F_E\) need not be bounded by
\(T\).  Normalizing a rare channel before estimating it discards the exact
probability which compensates its marginal concentration, as proved in
\Cref{prop:V107-normalization-trap}.  Every V108 use retains the
unnormalized POVM effect.
\end{firewall}

\section{Insertion of two pair levels and one core shell}
\label{sec:V108-pair-core-insertion}

At fixed pair levels \(t,s\), let
\[
 E^0_{t,s}=P_e(t)\otimes P_f(s)
\]
be the V106 pre-recoupling effect, regrouped by retained row.  Let
\(M\ge0\) be a positive invariant core shell on irreducible rows \(A,B\),
and define the complete upstream carrier
\begin{equation}
 T_{t,s}:=E^0_{t,s}\otimes M.
 \label{eq:V108-upstream-carrier}
\end{equation}
All tensor factors are regrouped by the same original-row cut.  Apply any
normalized downstream POVM only after this carrier is formed, as in
\eqref{eq:V108-pair-first-effect}.

\begin{theorem}[Every downstream event inherits the two-pair minimum]
\label{thm:V108-event-two-pair-minimum}
For every downstream event \(E\),
\begin{align}
 \|F_{E;t,s}\|_{\rm prod}
 \le
 \min\Biggl\{
 \Pi(M),\,
 \frac{\mathcal H(M)D_e(t)D_f(s)}
 {d_Ad_{X_e}d_{X_f}},\,
 \frac{\mathcal H(M)D_e(t)D_f(s)}
 {d_Bd_{Y_e}d_{Y_f}}
 \Biggr\}.
 \label{eq:V108-event-two-pair-minimum}
\end{align}
The bound is independent of the event's Racah multiplicity, output
dimension, and conditional probability.
\end{theorem}

\begin{proof}
\Cref{lem:V108-event-subeffect} gives
\(0\le F_{E;t,s}\le T_{t,s}\).  Product-capacity monotonicity gives the
first entry through
\Cref{lem:V104-two-projection-product-cap}.  The two partial marginals of
\(T_{t,s}\) factor into those of \(E^0_{t,s}\) and \(M\).  Use
\Cref{prop:V106-prerecoupling-scalar} for the former and core Schur
scalarization for the latter.  Then
\Cref{thm:V108-no-downstream-inflation} gives the two trace entries.
\end{proof}

\begin{corollary}[Near-unit downstream event surface]
\label{cor:V108-near-unit-event}
For \(M=M_N=s_N^2P_{S_N}\),
\begin{align}
 \|F_{E;t,s}\|_{\rm prod}
 \le
 s_N^2
 \min\Biggl\{
 \frac N{N+1},\,
 \frac{D_e(t)D_f(s)N(N+2)}
 {3d_{X_e}d_{X_f}},\,
 \frac{2ND_e(t)D_f(s)}
 {(N+1)d_{Y_e}d_{Y_f}}
 \Biggr\}.
 \label{eq:V108-near-unit-event}
\end{align}
\end{corollary}

\begin{proof}
Insert the V103 near-unit dimensions and trace in
\eqref{eq:V108-event-two-pair-minimum}.
\end{proof}

\section{The V70/V76 Racah effects are normalized downstream effects}
\label{sec:V108-V70-insertion}

Fix an \(AB\)-ancestry channel \(S\).  The V70 atoms
\(\Pi_S(V,T)\) act on the multiplicity space
\(\mathbb C^{N_{AB}^{S}}\) and obey
\[
 \Pi_S(V,T)\ge0,
 \qquad
 \sum_{V,T}\Pi_S(V,T)=I.
\]
Lift them to the full \(S\)-isotypic carrier by
\begin{equation}
 \mathsf M_S(V,T):=
 I_{H_S}\otimes\Pi_S(V,T).
 \label{eq:V108-lifted-Racah-atom}
\end{equation}

\begin{lemma}[Lifted Racah normalization]
\label{lem:V108-lifted-Racah-normalization}
For every fixed \(S\),
\begin{equation}
 \mathsf M_S(V,T)\ge0,
 \qquad
 \sum_{V,T}\mathsf M_S(V,T)=I_{H_S}\otimes I.
 \label{eq:V108-lifted-Racah-normalization}
\end{equation}
Taking the orthogonal direct sum over \(S\) gives a normalized POVM on the
complete \(AB\) carrier.
\end{lemma}

\begin{proof}
Tensor the V70 POVM normalization
\eqref{eq:V70-POVM-normalization} with \(I_{H_S}\).  The orthogonal
\(S\)-isotypic projectors sum to the identity on \(A\otimes B\).
\end{proof}

\begin{proposition}[The V76 probability is the Born law of the lifted POVM]
\label{prop:V108-V76-Born-law}
For a product vector \(u\otimes v\), condition on \(S\) as in V76.  Then
\begin{equation}
 \operatorname{Tr}\left[
 \mathsf M_S(V,T)
 \frac{|c_S(u,v)\rangle\langle c_S(u,v)|}{p_S(u,v)}
 \right]
 =
 \operatorname{Tr}\left[
 \Pi_S(V,T)\sigma_S(u,v)
 \right].
 \label{eq:V108-V76-Born-law}
\end{equation}
Multiplication by \(p_S(u,v)\) gives exactly
\(\lambda_{u,v}(S,V,T)\).
\end{proposition}

\begin{proof}
Use the defining partial-trace identity
\[
 \operatorname{Tr}\bigl[
 (I_{H_S}\otimes X)|c_S\rangle\langle c_S|
 \bigr]
 =
 p_S\operatorname{Tr}(X\sigma_S)
\]
with \(X=\Pi_S(V,T)\).  This is also the proof of the V76 two-level law.
\end{proof}

\begin{theorem}[Pair-first V70/V76 events are upstream subeffects]
\label{thm:V108-Racah-subeffects}
Insert the two pair levels and the core ancestry projection before the
lifted Racah POVM, in the aggregate order required by V57.  Then every
resolved Racah event effect, and every sum of such event effects, has the
form
\begin{equation}
 F_E=T^{1/2}M_ET^{1/2},
 \qquad
 0\le M_E\le I,
 \label{eq:V108-Racah-subeffect-form}
\end{equation}
where \(T\) is the complete upstream pair-level/core carrier.  Consequently
\begin{equation}
 0\le F_E\le T.
 \label{eq:V108-Racah-below-upstream}
\end{equation}
Any further orthogonal final-output refinement preserves this conclusion.
\end{theorem}

\begin{proof}
\Cref{lem:V108-lifted-Racah-normalization} supplies the normalized
downstream POVM.  Pair-first insertion is the Lüders sandwich
\eqref{eq:V108-Racah-subeffect-form}; when the upstream level is a
projection this is simply its compression of the POVM atom.  The V57
aggregate projector performs precisely this resolution before raw
pinching.  Apply \Cref{lem:V108-event-subeffect}.

A further orthogonal output resolution replaces each POVM atom by positive
subatoms whose sum is the original atom.  The combined family is again a
POVM, so the same argument applies.
\end{proof}

\begin{remark}[Racah conjugation is not an active global unitary]
\label{rem:V108-passive-Racah}
The unitary in \eqref{eq:V70-Racah-unitary} changes coordinates between two
intertwiner bases in the same physical tensor product.  It is used to
compute the matrices \(\Pi_S(V,T)\).  It does not actively replace an
upstream effect \(T\) by \(UTU^*\) while forgetting \(T\).  The latter
operation was the deliberately adversarial information test in
\Cref{prop:V106-global-unitary-inflation}; it is not the V70/V76 outcome
law.
\end{remark}

\section{Closure of the downstream concentration issue}
\label{sec:V108-concentration-closure}

\begin{theorem}[No Racah marginal concentration in a fixed orientation]
\label{thm:V108-no-physical-concentration}
For every complementary same-coordinate pair level \(t,s\), every V70/V76
Racah event \(E\), and every subsequent orthogonal output refinement,
\begin{align}
 \|\operatorname{Tr}_{\mathsf b}F_{E;t,s}\|_{\rm op}
 &\le
 \|\operatorname{Tr}_{\mathsf b}T_{t,s}\|_{\rm op},
 \label{eq:V108-no-left-concentration}\\
 \|\operatorname{Tr}_{\mathsf a}F_{E;t,s}\|_{\rm op}
 &\le
 \|\operatorname{Tr}_{\mathsf a}T_{t,s}\|_{\rm op}.
 \label{eq:V108-no-right-concentration}
\end{align}
The right sides are exactly the scalar pre-recoupling marginals computed in
\Cref{prop:V106-prerecoupling-scalar}, tensored with the core Schur
marginals.  Hence the physical fixed-orientation effect has no positive
Racah concentration loss relative to V104.
\end{theorem}

\begin{proof}
\Cref{thm:V108-Racah-subeffects} gives
\(F_{E;t,s}\le T_{t,s}\).  Apply partial trace and operator-norm
monotonicity, exactly as in
\Cref{thm:V108-no-downstream-inflation}.  The upstream marginal identities
are V106 and core Schur scalarization.
\end{proof}

\begin{corollary}[The abstract V106 \(\Gamma\)-loss is absent here]
\label{cor:V108-Gamma-disposition}
When a physical Racah event is measured relative to its upstream pair
level, its \emph{absolute} marginal heights are bounded by their
pre-recoupling values.  A normalized concentration factor may exceed one
only because the event trace has decreased; the exact event probability
compensates that increase.  Therefore no factor
\(\Gamma_{\mathsf a}\) or \(\Gamma_{\mathsf b}\) is inserted into
\eqref{eq:V108-event-two-pair-minimum}.
\end{corollary}

\begin{proof}
The absolute statement is
\eqref{eq:V108-no-left-concentration}--%
\eqref{eq:V108-no-right-concentration}.  Concentration factors divide those
heights by the event's mean eigenvalue.  The division is a conditional
normalization forbidden by
\Cref{fire:V108-no-conditional-normalization}.
\end{proof}

\begin{theorem}[Complementary overlap inherits the V104 two-pair gate]
\label{thm:V108-complementary-inherits-V104}
On the \(O_{\rm comp}\) branch of
\Cref{thm:V105-audited-extraction}, retain the V56 tensor factor, insert
both pair levels before the V57 aggregate output projector, and remain in
one V94 retained-row orientation.  Then the complete positive event effect
is bounded at every \((t,s)\) by the V104 one-minimum surface.  After
integrating the pair levels,
\begin{equation}
 \|K_{\rm comp}\|_{\rm prod}
 \le
 \int_0^{A_e}\!\!\int_0^{A_f}
 \min\left\{
 \Pi(M),\,
 \frac{D_e(t)D_f(s)\mathcal H(M)}
 {\max\{
 d_{X_e}d_{X_f}d_A,\,
 d_{Y_e}d_{Y_f}d_B
 \}}
 \right\}\dd s\,\dd t.
 \label{eq:V108-complementary-V104-gate}
\end{equation}
No independence of the two pair outputs or of the Racah outcome is assumed.
\end{theorem}

\begin{proof}
The V56 local identity and V58 row-factorized assembly give the two pair
level projector \(E^0_{t,s}\).  The downstream Racah/output event is a
subeffect by \Cref{thm:V108-Racah-subeffects}.  Apply
\eqref{eq:V108-event-two-pair-minimum} and combine the two trace entries
into the displayed maximum denominator.  Finally integrate in \(t,s\).
All event atoms are retained under one product-state supremum, so no
probabilistic independence or separate capacity product occurs.
\end{proof}

\begin{corollary}[Conditional linear threshold on the complementary branch]
\label{cor:V108-complementary-linear-threshold}
The V104 cold--cold rank theorem and its conditional decay statement apply
to the complementary overlap branch.  In particular, with bounded
coefficient caps, bounded \(s_N^2\), bounded orientation loss, and endpoint
effective scales \(R_e,R_f\),
\[
 R_e,R_f\to\infty,
 \qquad
 \frac{R_eR_f}{N^2}\to\infty
\]
imply decay of its near-unit product capacity.  In the symmetric case the
sufficient threshold is \(R\gg N\).
\end{corollary}

\begin{proof}
Equation \eqref{eq:V108-complementary-V104-gate} is the V104 common
two-pair gate for this branch.  Apply
\Cref{thm:V104-two-pair-effective-gate,cor:V104-linear-threshold}.
\end{proof}

\begin{firewall}[Two noncommuting orientations still have no joint POVM]
\label{fire:V108-two-orientation}
V108 works in one fixed retained-row orientation.  The \(ab\)- and
\(ac\)-oriented V94 event families act on the same original-row product
space and need not commute.  Subeffect domination for each family
separately does not manufacture a common parent effect for both.  The
two-orientation capacity \eqref{eq:V94-two-orientation-capacity} therefore
remains open.
\end{firewall}

\section{A non-singlet pure-level benchmark}
\label{sec:V108-pure-level}

\begin{proposition}[Schmidt coefficients multiply under pair reshuffling]
\label{prop:V108-Schmidt-product}
Let \(\xi\in U\otimes U'\) and \(\zeta\in V\otimes V'\) be unit vectors
with squared Schmidt coefficient lists
\[
 p_1\ge p_2\ge\cdots\ge0,
 \qquad
 q_1\ge q_2\ge\cdots\ge0.
\]
After regrouping
\[
 (U\otimes U')\otimes(V\otimes V')
 \longrightarrow
 (U\otimes V)\otimes(U'\otimes V'),
\]
the squared Schmidt coefficients of \(\xi\otimes\zeta\) are the products
\begin{equation}
 \{p_iq_j\}_{i,j}.
 \label{eq:V108-Schmidt-products}
\end{equation}
Consequently
\begin{equation}
 \bigl\|
 |\xi\otimes\zeta\rangle\langle\xi\otimes\zeta|
 \bigr\|_{\rm prod}
 =
 p_1q_1.
 \label{eq:V108-rank-one-product-cap}
\end{equation}
\end{proposition}

\begin{proof}
Write Schmidt decompositions
\[
 \xi=\sum_i\sqrt{p_i}\,u_i\otimes u_i',
 \qquad
 \zeta=\sum_j\sqrt{q_j}\,v_j\otimes v_j'.
\]
After regrouping,
\[
 \xi\otimes\zeta
 =
 \sum_{i,j}\sqrt{p_iq_j}\,
 (u_i\otimes v_j)\otimes(u_i'\otimes v_j').
\]
Both displayed tensor families are orthonormal, so this is a Schmidt
decomposition.  The product capacity of a pure-state projector is its
largest squared Schmidt coefficient.
\end{proof}

\begin{corollary}[Dual singlets are the flat endpoint of the pure law]
\label{cor:V108-singlet-flat-endpoint}
For normalized dual singlets,
\[
 p_i=\frac1{d_U},
 \qquad
 q_j=\frac1{d_V}.
\]
Thus \eqref{eq:V108-rank-one-product-cap} gives
\[
 \frac1{d_Ud_V}=\frac1{d_X},
\]
recovering the V107 full-row product capacity.
\end{corollary}

\begin{proof}
Every Schmidt coefficient of a normalized dual singlet is flat.  Substitute
in \eqref{eq:V108-rank-one-product-cap}.
\end{proof}

\begin{assessment}[What remains after V108]
\label{ass:V108-remaining-branches}
For a fixed orientation, the two-pair spectral mechanism is now available
both when the endpoint pair banks use disjoint coordinates and when they
form a same-coordinate complementary \([2,2]\) block.  The double-singlet
subcase has the stronger V107 \(N^{-2}\) estimate.  The remaining puncture
overlap is the same-pair branch, where both endpoint mass certificates refer
to one physical pair operator and no cold-rank product is legitimate.
Separately, the V94 sum of two noncommuting orientations and the V98 off-cut
carrier remain open.
\end{assessment}

\section{V108 ledger and V109 acquisition}
\label{sec:V108-ledger}

\begin{theorem}[Integrated V108 statement]
\label{thm:V108-integrated}
Preserving V1--V107, V108 proves:
\begin{enumerate}[label=\textup{(V108.\arabic*)},leftmargin=6.3em]
\item every pair-first downstream POVM event is a positive subeffect of its
      complete upstream carrier;
\item product capacity and both partial-marginal heights are monotone under
      that event selection;
\item the exact weighted downstream layer cake
      \eqref{eq:V108-downstream-layer-cake};
\item the pair-level/core event minimum
      \eqref{eq:V108-event-two-pair-minimum};
\item normalized lifting of the V70 Racah POVM and recovery of the V76 Born
      law;
\item absence of physical fixed-orientation Racah marginal inflation;
\item inheritance of the V104 gate and conditional linear threshold by the
      complementary same-coordinate branch; and
\item the exact pure-level Schmidt-product law
      \eqref{eq:V108-Schmidt-products}.
\end{enumerate}
\end{theorem}

\begin{proof}
The items are
\Cref{lem:V108-event-subeffect,
thm:V108-no-downstream-inflation,
cor:V108-downstream-layer-cake,
thm:V108-event-two-pair-minimum,
lem:V108-lifted-Racah-normalization,
prop:V108-V76-Born-law,
thm:V108-no-physical-concentration,
thm:V108-complementary-inherits-V104,
cor:V108-complementary-linear-threshold,
prop:V108-Schmidt-product}.
\end{proof}

\begin{openproblem}[V109 mandate: the bi-heavy same-pair branch]
\label{op:V109-mandate}
Return to the \(O_{\rm same}\) branch of
\Cref{thm:V105-audited-extraction}, where one pair bank is heavy for both
sideways endpoints.
\begin{enumerate}
\item Write the two row-mass lower bounds on the same pair support and the
      exact V52/V53 bilinear response before choosing an endpoint
      denominator.
\item Apply the selected-heavy tax from each endpoint separately and take
      their minimum; do not multiply the same physical pair tax by itself.
\item Determine whether the minimum yields a geometric-mean denominator
      \(1/(\Xi_i\Xi_k)\) sufficient for the V57 row-allocated path stock.
\item Translate the bi-heavy condition into the pair effective number and
      the V103 top-\(\eta_N\) spectral mean.
\item Test rank-one singlet, Cartan highest-weight, highly unbalanced row
      masses, and spectator inflation.
\item If the one-pair top spectrum remains persistent, retain its signed
      coupling to the complete core column rather than repeating a marginal
      capacity estimate.
\item Keep the noncommuting-orientation and off-cut branches separate.
\end{enumerate}
The next compulsory companion audit is V110.
\end{openproblem}

\section{Firewall after V108}
\label{sec:V108-firewall}

\begin{firewall}[Current claim boundary]
\label{fire:V108-current-boundary}
V108 closes the downstream marginal-concentration issue only for pair-first
normalized POVM resolutions in one fixed orientation.  It does not close
the bi-heavy same-pair branch, prove the V104 effective-scale condition for
all physical monomials, control the complete signed core scalar, combine the
two noncommuting orientations, close the off-cut branch, prove JREC,
construct the continuum theory, or prove a positive mass gap.  The full
Yang--Mills existence and mass gap remain open.
\end{firewall}

\section{Append-only record for V108}
\label{sec:V108-record}

V108 was appended to the complete V107 manuscript.  The exact V107 parent
had \(73220\) physical lines, \(2552798\) bytes, and SHA--256 digest
\[
 \texttt{0f510459bdf601424a5c7de0c1c85b6658a79cda22d172ecd97d97300ab1fd44}.
\]
No V107 mathematical content was changed.  The sole final
\verb|\end{document}| is restored below.

% =============================================================================
% END APPENDED VERSION V108 MATERIAL
% =============================================================================
% =============================================================================
% BEGIN APPENDED VERSION V109 MATERIAL
% =============================================================================

\chapter{V109: Bi-Heavy Same-Pair Tax and the Geometric
Two-Endpoint Denominator}
\label{ch:V109}

\section{Purpose}
\label{sec:V109-purpose}

V105 showed that the only overlap not carrying two distinct pair factors is
the same-pair branch: both sideways endpoint certificates use one physical
pair bank.  V103 correctly warns that its spectral capacity cannot be
squared.  There is nevertheless additional information at the
coefficient-response level.  The bank is selected-heavy from both
endpoints, so the universal V53 one-edge tax may be oriented in either
direction.  Taking the smaller of the two valid estimates yields a
geometric two-endpoint denominator.  This uses the same tax once and is
therefore compatible with the no-double-charge rule.

When the rest of the assembled monomial supplies the V57
\(s_\alpha\)-weighted core cut, the inverse \(s_\alpha\) in the pair tax
cancels and the exact marked path weight results.  Thus the same-pair
overlap can be removed from the new product-capacity ledger on that
cut-supported branch.  If no such companion cut is present, the V103
one-pair top-spectrum problem remains.

\section{Two orientations of one physical pair tax}
\label{sec:V109-two-oriented-tax}

Let \(L\) be a fixed same-coordinate pair bank joining rows \(i,k\).  In the
V53 notation, write its coordinate masses as
\begin{equation}
 x_i:=\sum_{e\in L}r_e,
 \qquad
 x_k:=\sum_{e\in L}q_e,
 \qquad
 h:=\sum_{e\in L}r_eq_e.
 \label{eq:V109-pair-masses}
\end{equation}
The two oriented effective numbers are
\begin{equation}
 N_i:=\frac{x_i^2}{h},
 \qquad
 N_k:=\frac{x_k^2}{h}.
 \label{eq:V109-two-effective-numbers}
\end{equation}
Let \(\mathfrak F_L(A_i,A_k)\) be the V52/V53 weighted
formation--defect functional on arbitrary spectator-amplified trace-class
row coefficients.

\begin{theorem}[Bi-heavy geometric one-edge tax]
\label{thm:V109-bi-heavy-tax}
Assume
\begin{equation}
 x_i\ge\varepsilon\Xi_i,
 \qquad
 x_k\ge\varepsilon\Xi_k
 \label{eq:V109-bi-heavy}
\end{equation}
for a fixed \(\varepsilon>0\).  Then
\begin{align}
 \mathfrak F_L(A_i,A_k)
 &\le
 C_\varepsilon
 \frac{h}{s_\alpha
       \max\{\Xi_i^2,\Xi_k^2\}}
 \|A_i\|_1\|A_k\|_1
 \label{eq:V109-max-square-tax}\\
 &\le
 C_\varepsilon
 \frac{h}{s_\alpha\Xi_i\Xi_k}
 \|A_i\|_1\|A_k\|_1.
 \label{eq:V109-geometric-tax}
\end{align}
Only one physical pair functional occurs on the left.
\end{theorem}

\begin{proof}
Orient \Cref{cor:V53-one-edge-tax} from endpoint \(i\).  The first
hypothesis in \eqref{eq:V109-bi-heavy} gives
\[
 \mathfrak F_L(A_i,A_k)
 \le
 C_\varepsilon\frac{h}{s_\alpha\Xi_i^2}
 \|A_i\|_1\|A_k\|_1.
\]
The same theorem is symmetric in the two pair inputs.  Orienting it from
endpoint \(k\) and using the second hypothesis gives the analogous bound
with \(\Xi_k^2\).  Take the minimum of the two upper bounds:
\[
 \min\left\{\frac1{\Xi_i^2},\frac1{\Xi_k^2}\right\}
 =
 \frac1{\max\{\Xi_i^2,\Xi_k^2\}}.
\]
Since
\(\max\{\Xi_i^2,\Xi_k^2\}\ge\Xi_i\Xi_k\), the second inequality follows.
No two estimates are multiplied.
\end{proof}

\begin{corollary}[Effective-number form]
\label{cor:V109-effective-geometric-tax}
Under the hypotheses of
\Cref{thm:V109-bi-heavy-tax},
\begin{equation}
 \boxed{\;
 \mathfrak F_L(A_i,A_k)
 \le
 \frac{C_\varepsilon}
 {s_\alpha\sqrt{N_iN_k}}
 \|A_i\|_1\|A_k\|_1.
 \;}
 \label{eq:V109-effective-geometric-tax}
\end{equation}
\end{corollary}

\begin{proof}
Because \(L\) is a subbank of each row,
\[
 x_i\le\Xi_i,
 \qquad
 x_k\le\Xi_k.
\]
Hence
\[
 \frac{h}{\Xi_i\Xi_k}
 \le
 \frac{h}{x_ix_k}
 =
 \frac1{\sqrt{N_iN_k}},
\]
using \eqref{eq:V109-two-effective-numbers}.  Insert this in
\eqref{eq:V109-geometric-tax}.
\end{proof}

\begin{corollary}[The V105 same-pair overlap is uniformly bi-heavy]
\label{cor:V109-V105-bi-heavy}
On the \(O_{\rm same}\) alternative of
\Cref{thm:V105-quarter-dichotomy} arising from two V102 endpoint
certificates, the common pair bank satisfies
\begin{equation}
 x_i\ge\frac{\delta_0}{24}\Xi_i,
 \qquad
 x_k\ge\frac{\delta_0}{24}\Xi_k.
 \label{eq:V109-V105-bi-heavy}
\end{equation}
Therefore \Cref{thm:V109-bi-heavy-tax} applies with a fixed programme
constant \(\varepsilon=\delta_0/24\).
\end{corollary}

\begin{proof}
This is the pure-overlap alternative
\eqref{eq:V105-quarter-pure-overlap}, specialized to
\(\tau=\mathrm{same}\).
\end{proof}

\begin{remark}[Why the minimum is a legitimate two-endpoint gain]
\label{rem:V109-minimum-legitimate}
Both oriented inequalities estimate the same nonnegative scalar
\(\mathfrak F_L(A_i,A_k)\).  Taking their minimum is ordinary logical
intersection.  By contrast, multiplying the two right sides would claim two
independent occurrences of the same formation defect and would double
charge its coordinate mass.  V109 never makes that multiplication.
\end{remark}

\section{Cancellation against the core-cut scale}
\label{sec:V109-cut-cancellation}

\begin{lemma}[Abstract pair--cut composition]
\label{lem:V109-pair-cut-composition}
Suppose the rest of an assembled monomial, retained inside the arbitrary
spectator slots of the V53 functional, contributes a positive core-cut
response bounded by
\begin{equation}
 \mathfrak C\le C s_\alpha W
 \label{eq:V109-core-cut-scale}
\end{equation}
for some nonnegative cut numerator \(W\).  If the same pair bank is
bi-heavy as in \eqref{eq:V109-bi-heavy}, then the assembled pair--cut
response obeys
\begin{equation}
 \boxed{\;
 \mathfrak R_{ik}
 \le
 C_\varepsilon
 \frac{hW}{\Xi_i\Xi_k}
 \prod_v\|A_v\|_1.
 \;}
 \label{eq:V109-pair-cut-path}
\end{equation}
\end{lemma}

\begin{proof}
Keep the core and all other rowwise factors inside the spectator-amplified
coefficient slots until the V53 pair compression has been performed.  Use
\eqref{eq:V109-geometric-tax} for that compression and
\eqref{eq:V109-core-cut-scale} for the retained cut factor.  The factors
\(s_\alpha^{-1}\) and \(s_\alpha\) cancel.  Aggregate raw pinching is a
trace-norm contraction, leaving
\eqref{eq:V109-pair-cut-path}.
\end{proof}

\begin{corollary}[Exact marked-path weight on a matched V57 cut term]
\label{cor:V109-V57-path}
If a term of the V57 covariance identity is retained before trace norm with
the individual cut bound
\[
 \mathfrak C_{ik}\le Cs_\alpha\widehat H_{ik},
\]
then
\begin{equation}
 \mathfrak R_{ik}
 \le
 C_\varepsilon
 \frac{h\widehat H_{ik}}{\Xi_i\Xi_k}
 \prod_v\|A_v\|_1,
 \label{eq:V109-V57-path}
\end{equation}
which is the marked-tree weight of the path segment joining the two
endpoints through that pair bank and cut.
\end{corollary}

\begin{proof}
Apply \Cref{lem:V109-pair-cut-composition} with
\(W=\widehat H_{ik}\).
\end{proof}

\begin{firewall}[The aggregate V57 cut is not silently split]
\label{fire:V109-no-silent-cut-split}
The proved operator cap \eqref{eq:V57-cut-cap} contains
\[
 H_{a\mid bc}=\widehat H_{ab}+\widehat H_{ac}.
\]
Equation \eqref{eq:V57-original-row-allocation} is a majorant identity, not
by itself a pre-trace-norm coefficient allocation.  Therefore
\Cref{cor:V109-V57-path} closes the matched covariance subterm only when its
individual cut numerator has remained explicit.  Applying the aggregate
cap first yields
\[
 C_\varepsilon
 \frac{h(\widehat H_{ab}+\widehat H_{ac})}{\Xi_i\Xi_k},
\]
and the unmatched summand still requires its correct endpoint allocation.
\end{firewall}

\section{What bi-heaviness does not do to one-pair spectrum}
\label{sec:V109-spectral-boundary}

Let the input irreducibles on the two ends of the same pair bank have
dimensions \(d_i,d_k\), and let
\[
 d_L:=\max\{d_i,d_k\}.
\]
Choose \(v_*\in\{i,k\}\) whose input dimension equals \(d_L\), and denote
the corresponding oriented effective number from
\eqref{eq:V109-two-effective-numbers} by \(N_*\).  Put
\[
 R_*:=s_\alpha N_*.
\]

\begin{proposition}[The valid same-pair fractional gate remains one-sided]
\label{prop:V109-same-pair-fractional}
Under the V53 cold-rank/hot-coefficient data oriented from \(v_*\), the
same-pair top fractional mean satisfies the V103 bound
\begin{equation}
 \overline c_{L,N}
 \le
 \frac C{R_*}
 A_*
 \min\left\{
 1,\frac C{\eta_NR_*}
 \right\}.
 \label{eq:V109-same-pair-fractional}
\end{equation}
The second endpoint's selected-heavy hypothesis validates a second
one-edge coefficient estimate, but does not replace \(R_*\) by
\(\sqrt{R_iR_k}\) or \(R_iR_k\) in
\eqref{eq:V109-same-pair-fractional}.
\end{proposition}

\begin{proof}
The V53 oriented cold rank is normalized by the larger input-row dimension
\(d_L\) when the orientation is chosen from \(v_*\).  Apply
\Cref{lem:V103-fractional-cold-hot,thm:V103-effective-threshold} with
\(R=R_*\).  This proves the displayed bound.

There is only one pair operator and one cold spectral subspace.  The second
orientation supplies another upper bound on the same bilinear functional,
which may be intersected as in
\Cref{thm:V109-bi-heavy-tax}.  It supplies no second tensor factor or
second cold subspace, so the cold-rank product argument of V104 is
unavailable.
\end{proof}

\begin{corollary}[The one-pair spectral route retains its quadratic
threshold]
\label{cor:V109-same-pair-quadratic}
With bounded \(A_*,s_N^2\), the information in
\eqref{eq:V109-same-pair-fractional} forces near-unit decay through the
V103 route only under
\begin{equation}
 \frac{s_\alpha N_*}{N^2}\longrightarrow\infty.
 \label{eq:V109-same-pair-quadratic}
\end{equation}
\end{corollary}

\begin{proof}
This is
\Cref{thm:V103-effective-threshold} with \(R_*=s_\alpha N_*\) and
\(\eta_N\asymp N^{-2}\).
\end{proof}

\begin{proposition}[Bi-heavy rank-one regression]
\label{prop:V109-bi-heavy-rank-one}
There are abstract positive pair spectra with one nonzero eigenvalue
\(c>0\) for which both endpoint labels satisfy identical selected-heavy
cold-rank data, while
\begin{equation}
 \eta_Nd_L<1
 \quad\Longrightarrow\quad
 \overline c_{L,N}=c.
 \label{eq:V109-bi-heavy-rank-one}
\end{equation}
Thus bi-heaviness alone does not dilute the top eigenvalue.
\end{proposition}

\begin{proof}
Use the same rank-one spectral projector for both endpoint descriptions.
The top \(\eta_Nd_L\)-fractional sum takes only a fraction of its sole
eigenvalue when \(\eta_Nd_L<1\), and normalization by that same fraction
returns \(c\).  Endpoint mass inequalities concern the coordinate weights
which generated the operator; they do not create another copy of its
spectrum.
\end{proof}

\begin{firewall}[Coefficient closure and spectral closure are distinct]
\label{fire:V109-coefficient-versus-spectrum}
\Cref{thm:V109-bi-heavy-tax} can close a matched pair--cut coefficient
response even when the normalized one-pair top spectrum in
\eqref{eq:V109-bi-heavy-rank-one} does not decay.  Conversely, a small
one-pair top spectrum would not prove the missing pre-trace-norm allocation
of the unmatched V57 cut term.  The V94 priority ledger must choose one
route and remove its paid outcome; it may not charge both descriptions.
\end{firewall}

\section{The exact same-pair acquisition fork}
\label{sec:V109-same-pair-fork}

\begin{theorem}[Bi-heavy same-pair disposition]
\label{thm:V109-same-pair-disposition}
For a same-pair overlap outcome from
\Cref{thm:V105-audited-extraction}, exactly the following proof routes are
available from the established data:
\begin{enumerate}[label=\textup{(\roman*)}]
\item if the companion covariance term retains a matched individual cut
      numerator, the coefficient response has the marked-path bound
      \eqref{eq:V109-V57-path};
\item if only the aggregate cut
      \(\widehat H_{ab}+\widehat H_{ac}\) remains, the geometric pair tax
      is valid but the unmatched summand still needs a pre-trace-norm row
      allocation;
\item if the branch is kept in the near-unit product-capacity ledger, it
      obeys the one-pair fractional gate
      \eqref{eq:V109-same-pair-fractional}, with no two-pair rank product.
\end{enumerate}
\end{theorem}

\begin{proof}
Part \textup{(i)} is
\Cref{cor:V109-V57-path}; part \textup{(ii)} is
\Cref{fire:V109-no-silent-cut-split}; part \textup{(iii)} is
\Cref{prop:V109-same-pair-fractional}.  These routes use the same physical
pair carrier and are alternatives in the no-double-charge ledger.
\end{proof}

\begin{assessment}[Upstream move forced by V109]
\label{ass:V109-upstream-move}
The same-pair mass geometry is no longer the bottleneck: its two endpoint
denominators follow from one V53 tax by taking a minimum.  The unresolved
coefficient issue is upstream in the V57 ternary covariance: the aggregate
operator cap was taken before its \(ab\) and \(ac\) cut numerators became
separate coefficient effects.  V110 should audit the companions and then
retain those three covariance summands through the aggregate raw projector,
testing whether each receives the matching geometric pair denominator.
\end{assessment}

\section{V109 ledger and V110 acquisition}
\label{sec:V109-ledger}

\begin{theorem}[Integrated V109 statement]
\label{thm:V109-integrated}
Preserving V1--V108, V109 proves:
\begin{enumerate}[label=\textup{(V109.\arabic*)},leftmargin=6.3em]
\item the two oriented effective numbers
      \eqref{eq:V109-two-effective-numbers};
\item the bi-heavy maximum-square and geometric one-edge taxes
      \eqref{eq:V109-max-square-tax}--%
      \eqref{eq:V109-geometric-tax};
\item the effective geometric-mean response
      \eqref{eq:V109-effective-geometric-tax};
\item automatic bi-heaviness of the V105 same-pair overlap bank;
\item the abstract pair--cut cancellation
      \eqref{eq:V109-pair-cut-path};
\item the matched-cut marked-path bound
      \eqref{eq:V109-V57-path};
\item the strict firewall against silently splitting the V57 aggregate cut;
\item the valid one-pair fractional gate and its unchanged quadratic
      threshold; and
\item the complete same-pair acquisition fork of
      \Cref{thm:V109-same-pair-disposition}.
\end{enumerate}
\end{theorem}

\begin{proof}
The items are
\Cref{thm:V109-bi-heavy-tax,
cor:V109-effective-geometric-tax,
cor:V109-V105-bi-heavy,
lem:V109-pair-cut-composition,
cor:V109-V57-path,
fire:V109-no-silent-cut-split,
prop:V109-same-pair-fractional,
cor:V109-same-pair-quadratic,
thm:V109-same-pair-disposition}.
\end{proof}

\begin{openproblem}[V110 mandate: audit and covariance-row allocation]
\label{op:V110-mandate}
First perform the compulsory comparison with the current Standard Model and
Navier--Stokes notes.  Then return to the exact three-covariance identity
\eqref{eq:V57-cut-identity}.
\begin{enumerate}
\item Keep its three signed covariance summands separate through the V57
      aggregate collision projector.
\item For a bi-heavy same-pair bank, attach
      \eqref{eq:V109-geometric-tax} before taking the trace norm of each
      covariance summand.
\item Determine which summands carry
      \(\widehat H_{ab}\) and which carry \(\widehat H_{ac}\), without using
      the aggregate majorant first.
\item Prove the two correct row-allocated path weights or exhibit the exact
      cross term that prevents them.
\item If a cross term survives, retain its signed covariance pairing and
      formulate one common effect rather than two positive majorants.
\item Reinsert the resulting disposition into the V94 priority selector and
      remove every paid same-pair outcome once.
\item Keep the two-orientation and off-cut branches separate.
\end{enumerate}
\end{openproblem}

\section{Firewall after V109}
\label{sec:V109-firewall}

\begin{firewall}[Current claim boundary]
\label{fire:V109-current-boundary}
V109 proves the geometric two-endpoint tax for a bi-heavy same pair and
closes its matched individual pair--cut subterm.  It does not prove a
pre-trace-norm split of the aggregate V57 cut, close the unmatched
covariance term, improve the one-pair near-unit spectral threshold, combine
the two noncommuting orientations, close the off-cut branch, prove JREC,
construct the continuum theory, or prove a positive mass gap.  The full
Yang--Mills existence and mass gap remain open.
\end{firewall}

\section{Append-only record for V109}
\label{sec:V109-record}

V109 was appended to the complete V108 manuscript.  The exact V108 parent
had \(73848\) physical lines, \(2574461\) bytes, and SHA--256 digest
\[
 \texttt{6fcffb02d4178e2e8b54fa450ef9a7e6aa5710dab334de6194d51172a0398668}.
\]
No V108 mathematical content was changed.  The sole final
\verb|\end{document}| is restored below.

% =============================================================================
% END APPENDED VERSION V109 MATERIAL
% =============================================================================
% =============================================================================
% BEGIN APPENDED VERSION V110 MATERIAL
% =============================================================================

\chapter{V110: Exact Hybrid Covariance Allocation and the
Positive-Effect Interference Gate}
\label{ch:V110}

\section{Purpose}
\label{sec:V110-purpose}

V109 proved that one bi-heavy physical pair bank admits either endpoint
orientation of the V53 tax, and therefore the geometric denominator
\(1/(\Xi_i\Xi_k)\), without pretending that two independent pair
operators are present.  Its remaining coefficient question was upstream:
the V57 estimate had combined the two cut numerators
\(\widehat H_{ab}\) and \(\widehat H_{ac}\) before the final trace norm.

V110 resolves that amplitude-level question.  The three signed
covariances in \eqref{eq:V57-cut-identity} first combine to the covariance
of \(F_a\) with the \emph{centred product}
\((F_b-\mathbb EF_b)(F_c-\mathbb EF_c)\).  A one-coordinate hybrid
replacement then decomposes its martingale increment into exactly two
signed operators.  One contains the \(b\)-replacement and has the
individual \(ab\) influence numerator; the other contains the
\(c\)-replacement and has the individual \(ac\) numerator.  There is no
third amplitude remainder.

The same calculation also locates the next interface.  The V94 selector
acts on positive level effects, not on signed amplitudes.  Squaring the
sum of the two row-allocated amplitudes produces their off-diagonal
interference.  A single paid diagonal cannot in general be subtracted as a
positive effect.  V110 therefore prints the complete common block effect,
proves the exact rule for reinserting it into the V94 priority ledger, and
does not silently replace it by two independent positive outcomes.

\section{Mandatory V110 companion audit}
\label{sec:V110-audit}

The read-only audit on 5 September 2026 found
\begin{center}
\begin{tabular}{|p{0.50\textwidth}|r|p{0.30\textwidth}|}
\hline
Companion file & Bytes & SHA--256\\
\hline
\texttt{Renewal\_NS\_Stability\_Research\_Notes\_V31.tex}
& \(887537\) &
\texttt{f1121b62238ad5491bb7cf03635ea9934942edf573ed8faaa9ee79e1d38a6696}\\
\hline
\texttt{Renewal\_SM\_Einstein\_Continuum\_Research\_Notes\_V121.tex}
& \(3008430\) &
\texttt{36d4cdb0e332529df7dcb636a0b54a1519cd61cbfee1796b8fcbef19afe67d1f}\\
\hline
\end{tabular}
\end{center}
The files contain respectively \(23053\) and \(82157\) newline-delimited
source records under the audit counter.  The NS file is byte-identical to
the V105 audit.  It retains the warning that active times, signed work,
and parent tails must remain correlated until the restricted formation
quantity has been formed.

The SM programme advanced from V114 to V121.  V115 corrects an earlier
componentwise BRST claim because a vanishing boundary row does not imply
the vanishing of its normal derivative.  V120 keeps the full
harmonic--complement coupling through one Feshbach map and carries the
finite Hodge density with it.  V121 then distinguishes an on-shell
zero-energy incidence identity from the domain of a fixed heat
realization: the exact nonzero-energy defect is
\[
 B_0(\partial_rc)|_\Sigma=\zeta\,\gamma c.
\]
Thus a valid identity was not allowed to cross a change of analytic
domain without a new proof.

\begin{assessment}[V110 adjustment after the audit]
\label{ass:V110-audit-adjustment}
The transferable rule is an interface rule, not an estimate.  In the
present programme the V57 covariance identity lives at the signed
amplitude level, whereas the V94 layer cake lives at the positive-effect
level.  V110 will:
\begin{enumerate}[label=\textup{(\roman*)}]
\item keep the full signed covariance until its exact row split is formed;
\item carry both row pieces through the same downstream filter;
\item form their common positive block effect before invoking the V94
      selector; and
\item declare an outcome paid only when the \emph{whole} common effect,
      including its interference, has received a valid bound.
\end{enumerate}
No SM or NS analytic estimate is imported.
\end{assessment}

\section{Centred form of the three-covariance identity}
\label{sec:V110-centred}

Retain the V57 product probability space and its finitely many independent
coordinate variables.  The values \(F_a,F_b,F_c\) are bounded operators on
distinct input carriers, so their pointwise values commute across the
three carriers.  Put
\begin{equation}
 \overline F_r:=\mathbb EF_r,
 \qquad
 F_r^\circ:=F_r-\overline F_r
 \quad(r=a,b,c).
 \label{eq:V110-centred-factors}
\end{equation}

\begin{lemma}[The three covariances are one centred-product covariance]
\label{lem:V110-centred-covariance}
The exact V57 symbol satisfies
\begin{align}
 \mathbf K_T^\circ
 &=
 \operatorname{Cov}(F_a,F_bF_c)
 -\overline F_b\,\operatorname{Cov}(F_a,F_c)
 -\overline F_c\,\operatorname{Cov}(F_a,F_b)
 \notag\\
 &=
 \operatorname{Cov}(F_a,F_b^\circ F_c^\circ)
 =
 \mathbb E(F_a^\circ F_b^\circ F_c^\circ).
 \label{eq:V110-centred-covariance}
\end{align}
The identity is valid before recoupling and with arbitrary spectator
amplification.
\end{lemma}

\begin{proof}
Because the three carrier factors commute,
\[
 F_b^\circ F_c^\circ
 =
 F_bF_c-\overline F_bF_c-\overline F_cF_b
 +\overline F_b\overline F_c.
\]
Covariance is linear in its second entry and covariance with the last
constant operator is zero.  This proves the first equality.  Since
\(\mathbb EF_a^\circ=0\), the covariance on the second line equals the
displayed centred third moment.  Tensoring with a spectator identity does
not alter any expectation.
\end{proof}

\begin{firewall}[The three signed terms were not bounded separately]
\label{fire:V110-no-premature-three-norms}
Equation \eqref{eq:V110-centred-covariance} is an exact regrouping of all
three terms in \eqref{eq:V57-cut-identity}.  No absolute value or positive
majorant has yet been applied.  In particular, the two singleton-mean
covariances are the subtractions which centre the two factors of the
composite covariance; they are not two additional positive responses.
\end{firewall}

\section{Hybrid replacement and exact original-row allocation}
\label{sec:V110-hybrid}

Order the independent coordinates and let
\(\mathcal F_e\) reveal the first \(e\) of them.  For any integrable
operator-valued function \(H\), write
\begin{equation}
 D_eH:=\mathbb E(H\mid\mathcal F_e)
       -\mathbb E(H\mid\mathcal F_{e-1}).
 \label{eq:V110-Doob-increment}
\end{equation}
Let \(H^{(e)}\) denote replacement of coordinate \(e\) by an independent
copy, with all other coordinates unchanged.

\begin{lemma}[Exact conditional replacement formula]
\label{lem:V110-conditional-replacement}
For every \(H\),
\begin{equation}
 D_eH
 =
 \mathbb E\!\left(H-H^{(e)}\mid\mathcal F_e\right).
 \label{eq:V110-conditional-replacement}
\end{equation}
For \(G:=F_b^\circ F_c^\circ\), define
\begin{align}
 U_{b,e}
 &:=
 \mathbb E\!\left[
  (F_b-F_b^{(e)})F_c^\circ\mid\mathcal F_e
 \right],
 \label{eq:V110-Ub}\\
 U_{c,e}
 &:=
 \mathbb E\!\left[
  (F_b^{(e)}-\overline F_b)
  (F_c-F_c^{(e)})\mid\mathcal F_e
 \right].
 \label{eq:V110-Uc}
\end{align}
Then
\begin{equation}
 \boxed{D_eG=U_{b,e}+U_{c,e}.}
 \label{eq:V110-hybrid-product-increment}
\end{equation}
\end{lemma}

\begin{proof}
Conditional on \(\mathcal F_e\), integration over the independent
replacement and over all future coordinates gives
\[
 \mathbb E(H^{(e)}\mid\mathcal F_e)
 =\mathbb E(H\mid\mathcal F_{e-1}),
\]
which proves \eqref{eq:V110-conditional-replacement}.  Pointwise,
\begin{align*}
 F_b^\circ F_c^\circ
 -(F_b^{(e)}-\overline F_b)(F_c^{(e)}-\overline F_c)
 ={}&
 (F_b-F_b^{(e)})F_c^\circ\\
 &+(F_b^{(e)}-\overline F_b)(F_c-F_c^{(e)}).
\end{align*}
Apply conditional expectation and
\eqref{eq:V110-conditional-replacement}.
\end{proof}

\begin{theorem}[Exact two-row allocation of the V57 covariance]
\label{thm:V110-exact-row-allocation}
Define the two signed row operators
\begin{align}
 \mathbf K_{ab}^{\rhd}
 &:=
 \sum_e\mathbb E\!\left[(D_eF_a)U_{b,e}\right],
 \label{eq:V110-Kab}\\
 \mathbf K_{ac}^{\lhd}
 &:=
 \sum_e\mathbb E\!\left[(D_eF_a)U_{c,e}\right].
 \label{eq:V110-Kac}
\end{align}
Then, before trace norm, raw collision pinching, or final recoupling,
\begin{equation}
 \boxed{
 \mathbf K_T^\circ
 =
 \mathbf K_{ab}^{\rhd}
 +\mathbf K_{ac}^{\lhd}.}
 \label{eq:V110-exact-row-split}
\end{equation}
The first summand contains an \(a\)-increment and a \(b\)-replacement
difference at the same coordinate; the second contains an \(a\)-increment
and a \(c\)-replacement difference at the same coordinate.  Hence the
split is an exact original-row allocation, not merely the scalar identity
\(\widehat H_{ab}+\widehat H_{ac}\).
\end{theorem}

\begin{proof}
The operator martingale covariance identity used in
\Cref{lem:V50-operator-covariance} gives
\[
 \mathbb E(F_a^\circ G)
 =
 \sum_e\mathbb E[(D_eF_a)(D_eG)].
\]
Insert \eqref{eq:V110-hybrid-product-increment} and use linearity.
\Cref{lem:V110-centred-covariance} identifies the left side with
\(\mathbf K_T^\circ\).  The formulas
\eqref{eq:V110-Ub}--\eqref{eq:V110-Uc} print the selected original row in
each term.
\end{proof}

\begin{remark}[The hybrid order is harmless but must be fixed]
\label{rem:V110-hybrid-order}
Using
\[
 XY-X'Y'=X(Y-Y')+(X-X')Y'
\]
instead gives a second exact pair of row operators.  The individual
operators depend on this hybrid convention, but their sum is always
\(\mathbf K_T^\circ\).  V110 fixes
\eqref{eq:V110-Ub}--\eqref{eq:V110-Uc}; it never combines row pieces
constructed from different hybrid conventions.
\end{remark}

\section{Individual influence bounds after the exact split}
\label{sec:V110-individual-bounds}

For a bounded operator-valued function \(H\), put
\begin{align}
 \lambda_e(H)
 &:=
 \left\|
 \mathbb E[(H-H^{(e)})(H-H^{(e)})^*]
 \right\|_{\rm op}^{1/2},
 \label{eq:V110-left-influence}\\
 \rho_e(H)
 &:=
 \left\|
 \mathbb E[(H-H^{(e)})^*(H-H^{(e)})]
 \right\|_{\rm op}^{1/2}.
 \label{eq:V110-right-influence}
\end{align}

\begin{theorem}[Separate operator bounds for the two original rows]
\label{thm:V110-separate-influence}
Suppose, as for the normalized V57 Wilson factors, that
\(\|F_r\|_{\rm op}\le1\) pointwise.  Then
\begin{align}
 \|\mathbf K_{ab}^{\rhd}\|_{\rm op}
 &\le
 2\sum_e\lambda_e(F_a)\rho_e(F_b),
 \label{eq:V110-Kab-influence}\\
 \|\mathbf K_{ac}^{\lhd}\|_{\rm op}
 &\le
 2\sum_e\lambda_e(F_a)\rho_e(F_c).
 \label{eq:V110-Kac-influence}
\end{align}
In the Wilson influence notation already used in the proof of
\Cref{thm:V57-cut-cap}, these imply the individual cut bounds
\begin{align}
 \boxed{
 \|\mathbf K_{ab}^{\rhd}\|_{\rm op}
 \le Cs_\alpha\widehat H_{ab},
 \qquad
 \|\mathbf K_{ac}^{\lhd}\|_{\rm op}
 \le Cs_\alpha\widehat H_{ac}.}
 \label{eq:V110-individual-cut-bounds}
\end{align}
All four inequalities are stable under spectator amplification.
\end{theorem}

\begin{proof}
Conditional Jensen applied to
\eqref{eq:V110-conditional-replacement} gives
\[
 \left\|\mathbb E[(D_eF_a)(D_eF_a)^*]\right\|_{\rm op}^{1/2}
 \le\lambda_e(F_a).
\]
Also \(\|F_r^\circ\|_{\rm op}\le2\).  Since different carrier factors
commute, conditional Jensen gives
\begin{align*}
 \left\|\mathbb E(U_{b,e}^*U_{b,e})\right\|_{\rm op}^{1/2}
 &\le2\rho_e(F_b),\\
 \left\|\mathbb E(U_{c,e}^*U_{c,e})\right\|_{\rm op}^{1/2}
 &\le2\rho_e(F_c).
\end{align*}
For example, before conditional expectation,
\[
 \bigl((F_b-F_b^{(e)})F_c^\circ\bigr)^*
 \bigl((F_b-F_b^{(e)})F_c^\circ\bigr)
 \preccurlyeq
 4(F_b-F_b^{(e)})^*(F_b-F_b^{(e)}).
\]
Apply operator Cauchy--Schwarz to each expectation in
\eqref{eq:V110-Kab} and \eqref{eq:V110-Kac}, then sum over \(e\).
This proves \eqref{eq:V110-Kab-influence} and
\eqref{eq:V110-Kac-influence}.

The local Wilson replacement estimates used in V57 bound the two sums
separately by
\[
 C\sum_e\gamma_{a,e}\gamma_{b,e}
 \le Cs_\alpha\widehat H_{ab},
 \qquad
 C\sum_e\gamma_{a,e}\gamma_{c,e}
 \le Cs_\alpha\widehat H_{ac}.
\]
Their addition was the last step in the original V57 proof; it is not
needed here.  Tensoring by an identity preserves the conditional Jensen
and Cauchy--Schwarz inequalities.
\end{proof}

\begin{corollary}[The exact split survives the aggregate raw projector]
\label{cor:V110-raw-projector}
Let \(\mathscr P_{\rm raw}\) be the single aggregate pinching of
\Cref{lem:V57-aggregate-collision-pinching}, including any fixed
downstream spectator factors.  Then
\begin{equation}
 \mathscr P_{\rm raw}(\mathbf K_T^\circ)
 =
 \mathscr P_{\rm raw}(\mathbf K_{ab}^{\rhd})
 +\mathscr P_{\rm raw}(\mathbf K_{ac}^{\lhd}),
 \label{eq:V110-pinched-row-split}
\end{equation}
and each summand retains its corresponding bound in
\eqref{eq:V110-individual-cut-bounds}.  No intermediate representation
dimension or channel count is introduced.
\end{corollary}

\begin{proof}
The aggregate pinching is linear, so
\eqref{eq:V110-exact-row-split} gives the identity.  A pinching is a
contraction on operator norm and on the trace-class coefficient norm used
in V57.  Apply it only after the two signed row pieces have been labelled.
\end{proof}

\begin{corollary}[Bi-heavy pair payment for either exact row piece]
\label{cor:V110-bi-heavy-row-payment}
Fix \(r\in\{b,c\}\).  Suppose a same-pair bank \(L_{ar}\) is bi-heavy at
its endpoints as in \eqref{eq:V109-bi-heavy}, with pair mass \(h_{ar}\).
Keep \(\mathbf K_{ar}\) in the spectator slot of the V53 functional, where
\(\mathbf K_{ab}:=\mathbf K_{ab}^{\rhd}\) and
\(\mathbf K_{ac}:=\mathbf K_{ac}^{\lhd}\).  Its assembled coefficient
response satisfies
\begin{equation}
 \boxed{
 \mathfrak R_{ar}
 \le
 C_\varepsilon
 \frac{h_{ar}\widehat H_{ar}}
      {\Xi_a\Xi_r}
 \prod_v\|A_v\|_1.}
 \label{eq:V110-row-path-payment}
\end{equation}
Thus both individual V57 row weights are available, each on the amplitude
piece carrying its actual original-row replacement.
\end{corollary}

\begin{proof}
Use the geometric pair tax
\eqref{eq:V109-geometric-tax} before the final trace norm.  For the
retained signed covariance use the corresponding individual estimate in
\eqref{eq:V110-individual-cut-bounds}.  The factors
\(s_\alpha^{-1}\) and \(s_\alpha\) cancel exactly as in
\Cref{lem:V109-pair-cut-composition}.  Finally use
\Cref{cor:V110-raw-projector}.
\end{proof}

\begin{assessment}[Disposition of the V109 amplitude bottleneck]
\label{ass:V110-amplitude-disposition}
At signed-amplitude level the V109 bottleneck is closed:
\(\widehat H_{ab}\) and \(\widehat H_{ac}\) are attached to the exact
operators \(\mathbf K_{ab}^{\rhd}\) and
\(\mathbf K_{ac}^{\lhd}\), respectively, before trace norm.  The composite
covariance creates no third row term.  This statement alone does not yet
authorize deletion of one of those terms from a positive squared JREC
level effect.
\end{assessment}

\section{The positive-effect interference left by the row split}
\label{sec:V110-interference}

Let \(Q\) be any positive contraction on the downstream output space.  It
may be a resolved V70/V76 event effect, a sum of such effects, or a
downstream subeffect as in V108.  Set
\begin{equation}
 X_b:=Q^{1/2}\mathbf K_{ab}^{\rhd},
 \qquad
 X_c:=Q^{1/2}\mathbf K_{ac}^{\lhd}.
 \label{eq:V110-filtered-row-amplitudes}
\end{equation}
The complete filtered effect is
\begin{equation}
 \mathsf E_Q
 :=
 (\mathbf K_T^\circ)^*Q\mathbf K_T^\circ
 =
 (X_b+X_c)^*(X_b+X_c).
 \label{eq:V110-common-effect}
\end{equation}

\begin{theorem}[Exact common block effect]
\label{thm:V110-common-block-effect}
Let
\[
 \mathcal X\xi:=(X_b\xi,X_c\xi)
 \]
map the input into two copies of the filtered output space.  Then
\begin{equation}
 \boxed{
 \mathsf E_Q
 =
 \mathcal X^*
 \begin{pmatrix}I&I\\ I&I\end{pmatrix}
 \mathcal X.}
 \label{eq:V110-block-effect}
\end{equation}
Equivalently,
\begin{align}
 \mathsf E_Q
 ={}&
 X_b^*X_b+X_c^*X_c
 +X_b^*X_c+X_c^*X_b.
 \label{eq:V110-effect-expansion}
\end{align}
The block matrix in \eqref{eq:V110-block-effect} is positive, but the
off-diagonal sum in \eqref{eq:V110-effect-expansion} has no fixed sign.
\end{theorem}

\begin{proof}
Expand the right side of \eqref{eq:V110-block-effect}.  It is
\((X_b+X_c)^*(X_b+X_c)\), which is
\eqref{eq:V110-common-effect}.  The \(2\times2\) scalar matrix of ones has
eigenvalues \(0,2\), so its amplification is positive.  Taking
\(X_b=1,X_c=-1/2\) on a one-dimensional space makes the off-diagonal sum
equal to \(-1\); taking \(X_b=X_c=1\) makes it \(2\).  Hence that sum is
indefinite as a universal expression.
\end{proof}

\begin{proposition}[A single row diagonal cannot be positively deleted]
\label{prop:V110-no-single-diagonal-deletion}
There is no universal operator inequality
\begin{equation}
 \mathsf E_Q-X_b^*X_b\ge0
 \quad\hbox{or}\quad
 \mathsf E_Q-X_c^*X_c\ge0.
 \label{eq:V110-false-positive-deletion}
\end{equation}
Consequently a coefficient payment for only one row amplitude does not by
itself define a positive paid subeffect of the V94 level effect.
\end{proposition}

\begin{proof}
On a one-dimensional space take \(Q=1\), \(X_b=1\), and \(X_c=-1\).
Then \(\mathsf E_Q=0\), while both diagonal effects equal \(1\).  Either
difference in \eqref{eq:V110-false-positive-deletion} is \(-1\).
\end{proof}

\begin{proposition}[Sharp positive domination and statewise interference]
\label{prop:V110-positive-domination}
For every \(\theta>0\),
\begin{equation}
 \mathsf E_Q
 \preccurlyeq
 (1+\theta)X_b^*X_b
 +(1+\theta^{-1})X_c^*X_c.
 \label{eq:V110-Young-effect}
\end{equation}
For every density matrix \(\rho\), in particular every original-row
product density,
\begin{align}
 &\left|
 \operatorname{Tr}\rho
 (X_b^*X_c+X_c^*X_b)
 \right|
 \notag\\
 &\hspace{15mm}\le
 2
 \left(\operatorname{Tr}\rho X_b^*X_b\right)^{1/2}
 \left(\operatorname{Tr}\rho X_c^*X_c\right)^{1/2}.
 \label{eq:V110-statewise-cross-CS}
\end{align}
Both constants are sharp already for scalar amplitudes.
\end{proposition}

\begin{proof}
The positive square
\[
 (\sqrt\theta X_b-\theta^{-1/2}X_c)^*
 (\sqrt\theta X_b-\theta^{-1/2}X_c)
 \]
gives
\[
 X_b^*X_c+X_c^*X_b
 \preccurlyeq
 \theta X_b^*X_b+\theta^{-1}X_c^*X_c.
\]
Insert this in \eqref{eq:V110-effect-expansion}.  For the trace estimate,
apply Hilbert--Schmidt Cauchy--Schwarz to
\(X_b\rho^{1/2}\) and \(X_c\rho^{1/2}\), then add the adjoint pairing.
Equality occurs when the two scalar amplitudes are positive multiples of
one another.
\end{proof}

\begin{corollary}[Exact V94 reinsertion rule]
\label{cor:V110-V94-reinsertion}
At each common V94 level and for each fixed physical outcome, the following
is a disjoint valid priority rule.
\begin{enumerate}[label=\textup{(\roman*)}]
\item If both row diagonals \(X_b^*X_b\) and \(X_c^*X_c\) have valid
      coefficient/capacity payments on that outcome, apply
      \eqref{eq:V110-Young-effect} and move the \emph{whole}
      \(\mathsf E_Q\) to the first applicable paid class once.
\item If exactly one row diagonal has a payment, retain the whole common
      effect \(\mathsf E_Q\), including its off-diagonal entries, in the
      residual class.  Do not spend the one-row certificate elsewhere.
\item If neither diagonal has a payment, leave the outcome in the prior
      V94 residual class.
\end{enumerate}
This refinement neither assumes that the \(ab\) and \(ac\) event effects
commute nor introduces a joint Racah law.
\end{corollary}

\begin{proof}
In case \textup{(i)}, \eqref{eq:V110-Young-effect} bounds the complete
positive effect by the sum of the two paid positive diagonals.  Assigning
the physical outcome only once preserves the disjoint priority partition
\eqref{eq:V94-level-partition}.  In case \textup{(ii)},
\Cref{prop:V110-no-single-diagonal-deletion} forbids interpreting the
unpaid remainder as a positive effect.  Retaining the complete block is
therefore forced.  Case \textup{(iii)} changes nothing.  All three cases
use one outcome and one common level parameter.
\end{proof}

\begin{firewall}[A paid amplitude is not yet a paid level outcome]
\label{fire:V110-amplitude-effect-interface}
The path bound \eqref{eq:V110-row-path-payment} is a rigorous payment for
one signed row amplitude inside the trace-class coefficient calculation.
If the companion row amplitude is not also controlled, it does not permit
the corresponding positive level outcome to be removed from
\(\mathcal C_{\rm res}^{(2)}\).  This is the precise amplitude/effect
interface exposed by the V110 audit; crossing it without
\eqref{eq:V110-common-effect} would repeat the type of domain change
detected in SM V121.
\end{firewall}

\section{The exact residual after one row has a pair tax}
\label{sec:V110-one-row-residual}

For a common level \(t\), write \(Q_t\) for the downstream event filter and
use \eqref{eq:V110-filtered-row-amplitudes} with \(Q=Q_t\).  For an
original-row product density \(\rho\), put
\begin{align}
 p_b(t,\rho)&:=\operatorname{Tr}\rho X_b(t)^*X_b(t),
 \label{eq:V110-pb}\\
 p_c(t,\rho)&:=\operatorname{Tr}\rho X_c(t)^*X_c(t),
 \label{eq:V110-pc}\\
 \iota_{bc}(t,\rho)&:=
 2\operatorname{Re}\operatorname{Tr}\rho X_b(t)^*X_c(t).
 \label{eq:V110-interference}
\end{align}
Then
\begin{equation}
 \operatorname{Tr}\rho\,\mathsf E_{Q_t}
 =
 p_b(t,\rho)+p_c(t,\rho)+\iota_{bc}(t,\rho).
 \label{eq:V110-level-common-expectation}
\end{equation}

\begin{theorem}[One-row payment leaves one common-effect gate]
\label{thm:V110-one-row-gate}
Assume the \(ab\) row has the bi-heavy path payment
\eqref{eq:V110-row-path-payment}.  The only universal consequence for the
positive level effect is
\begin{equation}
 |\iota_{bc}(t,\rho)|
 \le2\sqrt{p_b(t,\rho)p_c(t,\rho)}.
 \label{eq:V110-one-row-cross-bound}
\end{equation}
No uniform \(o(1)\) gain follows from smallness of \(p_b\) alone unless
the scale of \(p_c\), or a signed cancellation in
\(\iota_{bc}\), is also controlled.
\end{theorem}

\begin{proof}
The displayed inequality is
\eqref{eq:V110-statewise-cross-CS}.  To see that no further universal
conclusion is available, take scalar amplitudes
\[
 X_b=\varepsilon,\qquad X_c=1.
\]
Then \(p_b=\varepsilon^2\) is arbitrarily small, but
\(\mathsf E_Q=(1+\varepsilon)^2\) remains of order one.  With
\(X_c=-1\), the interference reverses sign.  Thus neither magnitude nor
sign is determined by the paid row alone.
\end{proof}

\begin{corollary}[Two sufficient closures of the common gate]
\label{cor:V110-two-common-closures}
The V110 common effect is controlled if either of the following is proved
uniformly on the same level family and the same product state.
\begin{enumerate}[label=\textup{(\roman*)}]
\item Both \(p_b\) and \(p_c\) have the appropriate path envelopes; then
      \eqref{eq:V110-Young-effect}, for example with \(\theta=1\), pays
      the whole outcome by twice their sum.
\item The signed overlap admits a direct common-effect estimate which,
      together with the paid diagonal and the existing residual bound for
      the other diagonal, has the target path stock.
\end{enumerate}
Taking separate product-state suprema before checking either condition is
not an admissible substitute.
\end{corollary}

\begin{proof}
Part \textup{(i)} is
\eqref{eq:V110-Young-effect}.  Part \textup{(ii)} is the exact expansion
\eqref{eq:V110-level-common-expectation}.  The same-state requirement is
part of \eqref{eq:V94-two-orientation-capacity}.
\end{proof}

\begin{remark}[Relation to the V97 signed-collision correction]
\label{rem:V110-V97-relation}
V97 proved that the three positive covariance squares do not factor
through the complete signed cumulant.  V110 does not repeat that result.
It constructs a new two-row hybrid decomposition \emph{of the signed
cumulant itself}.  The common block
\eqref{eq:V110-block-effect} is the required V97-compatible way to square
that new decomposition; discarding its off-diagonal blocks would recreate
the separated-square error.
\end{remark}

\section{V110 ledger and V111 acquisition}
\label{sec:V110-ledger}

\begin{theorem}[Integrated V110 statement]
\label{thm:V110-integrated}
Preserving V1--V109, V110 proves:
\begin{enumerate}[label=\textup{(V110.\arabic*)},leftmargin=6.7em]
\item the mandatory SM V121 and NS V31 audit and its interface adjustment;
\item the centred-product form
      \eqref{eq:V110-centred-covariance} of the V57 identity;
\item the conditional hybrid replacement formula
      \eqref{eq:V110-hybrid-product-increment};
\item the exact pre-trace two-row split
      \eqref{eq:V110-exact-row-split};
\item the separate \(ab\) and \(ac\) cut bounds
      \eqref{eq:V110-individual-cut-bounds};
\item survival of the split through aggregate raw pinching;
\item the bi-heavy marked-path payment
      \eqref{eq:V110-row-path-payment} for either matching row piece;
\item the complete common positive block effect
      \eqref{eq:V110-block-effect};
\item impossibility of positive one-diagonal deletion; and
\item the exact three-case reinsertion rule for the V94 priority selector.
\end{enumerate}
\end{theorem}

\begin{proof}
The audit is \Cref{sec:V110-audit}.  The remaining items are
\Cref{lem:V110-centred-covariance,
lem:V110-conditional-replacement,
thm:V110-exact-row-allocation,
thm:V110-separate-influence,
cor:V110-raw-projector,
cor:V110-bi-heavy-row-payment,
thm:V110-common-block-effect,
prop:V110-no-single-diagonal-deletion,
cor:V110-V94-reinsertion}.
\end{proof}

\begin{openproblem}[V111 mandate: common-effect one-row capacity]
\label{op:V111-mandate}
Work on the common effect
\eqref{eq:V110-block-effect}, not on two separately normalized event laws.
\begin{enumerate}
\item Insert the V108 downstream-subeffect representation into the same
      filters \(Q_t\) in both row entries.
\item On the V105 same-pair branch, identify which row diagonal receives
      \eqref{eq:V110-row-path-payment} and retain the other diagonal and
      the signed overlap in one \(2\times2\) block.
\item Test whether the scalar-marginal product-capacity theorem of V106
      controls the unpaid diagonal without conditioning away its sector
      weight.
\item If it does, combine the two diagonal envelopes with
      \eqref{eq:V110-Young-effect} and remove the whole physical outcome
      once under \Cref{cor:V110-V94-reinsertion}.
\item If it does not, compute the precise rank-one or alternating
      regression for the off-diagonal pairing and return upstream to the
      local \(6j\) kernel.
\item Keep the two noncommuting retained-row orientations in the common
      V94 product-state supremum.
\end{enumerate}
\end{openproblem}

\section{Firewall after V110}
\label{sec:V110-firewall}

\begin{firewall}[Current claim boundary]
\label{fire:V110-current-boundary}
V110 closes the pre-trace covariance-row allocation and proves the
geometric pair payment for either matching signed row piece.  It also
proves that one such amplitude payment is not a positive subeffect payment
after squaring.  It does not control the remaining common-effect diagonal
or interference from one same-pair certificate, improve the one-pair
spectral threshold, combine the two noncommuting orientations, close the
off-cut branch, prove JREC, construct the continuum theory, or prove a
positive mass gap.  The full Yang--Mills existence and mass gap remain
open.
\end{firewall}

\section{Append-only record for V110}
\label{sec:V110-record}

V110 was appended to the complete V109 manuscript.  The exact V109 parent
had \(74326\) newline-delimited source records, \(2591145\) bytes, and
SHA--256 digest
\[
 \texttt{7d0311cab648b65124c7d7e99df80d64215c586ca23a679a270f13b342f81055}.
\]
No V109 mathematical content was changed.  The sole final
\verb|\end{document}| is restored below.

% =============================================================================
% END APPENDED VERSION V110 MATERIAL
% =============================================================================
% =============================================================================
% BEGIN APPENDED VERSION V111 MATERIAL
% =============================================================================

\chapter{V111: Upstream Douglas Domination for the Common
Covariance Effect}
\label{ch:V111}

\section{Purpose}
\label{sec:V111-purpose}

V110 constructed the exact signed row split
\[
 \mathbf K_T^\circ
 =
 \mathbf K_{ab}^{\rhd}+\mathbf K_{ac}^{\lhd}
\]
and showed that the relevant positive object is the common effect
\[
 (\mathbf K_T^\circ)^*Q\mathbf K_T^\circ,
\]
including both off-diagonal row pairings.  V111 tests the first proposed
route for that effect: inheritance of the V106--V108 upstream
product-capacity bound.

The test is exact.  A filtered amplitude inherits a positive upstream
carrier \(T\) precisely when its adjoint square is order dominated by
\(T\).  The least domination constant is not its ordinary operator norm;
it is the norm after division by \(T^{1/2}\) on the support of \(T\).
This is the upstream Douglas constant.  It is harmless on a raw projection
block, but it can diverge on a small eigenvalue of an integrated effect.
Consequently V106--V108 apply automatically level by level only after the
raw support identity has been retained.  They cannot be invoked from the
V110 covariance norm bound alone.

\section{The least upstream domination constant}
\label{sec:V111-Douglas}

All spaces in this section are finite-dimensional representation blocks.
Let \(T\ge0\) act on an input space \(H\), let \(P_T\) be its support
projection, and let \(T^{\dagger/2}\) be the Moore--Penrose inverse square
root.  For \(K:H\to G\), define
\begin{equation}
 \Lambda_T(K):=
 \begin{cases}
 \|KT^{\dagger/2}\|_{\rm op}^2,
   &K=KP_T,\\
 +\infty,&K\ne KP_T.
 \end{cases}
 \label{eq:V111-Douglas-constant}
\end{equation}

\begin{theorem}[Exact upstream Douglas constant]
\label{thm:V111-Douglas-constant}
The number in \eqref{eq:V111-Douglas-constant} is the least
\(\lambda\in[0,\infty]\) such that
\begin{equation}
 K^*K\preccurlyeq\lambda T.
 \label{eq:V111-Douglas-order}
\end{equation}
Equivalently, a finite order domination holds exactly when
\begin{equation}
 K=C\,T^{1/2}
 \label{eq:V111-Douglas-factor}
\end{equation}
for some \(C:H\to G\), and the least value is the squared norm of the
least-norm factor \(C=KT^{\dagger/2}\).
\end{theorem}

\begin{proof}
If \(K=KP_T\), then
\[
 K=(KT^{\dagger/2})T^{1/2}.
\]
Therefore
\[
 K^*K
 =
 T^{1/2}(KT^{\dagger/2})^*
 (KT^{\dagger/2})T^{1/2}
 \preccurlyeq
 \Lambda_T(K)T.
\]
Conversely, \eqref{eq:V111-Douglas-order} forces
\(K\xi=0\) on \(\ker T\), hence \(K=KP_T\).  Conjugating the restriction
of \eqref{eq:V111-Douglas-order} to \(\operatorname{Ran}P_T\) by
\(T^{\dagger/2}\) gives
\[
 (KT^{\dagger/2})^*(KT^{\dagger/2})
 \preccurlyeq\lambda P_T.
\]
Thus \(\Lambda_T(K)\le\lambda\).  This also proves the least-factor
statement.  It is the V97 finite-dimensional factorization criterion
applied to the square-root carrier, now with the optimal constant printed.
\end{proof}

\begin{corollary}[Every downstream filter preserves upstream domination]
\label{cor:V111-filtered-domination}
If \(0\le Q\le I_G\), then
\begin{equation}
 0\le K^*QK
 \preccurlyeq K^*K
 \preccurlyeq\Lambda_T(K)T.
 \label{eq:V111-filtered-order}
\end{equation}
Consequently
\begin{align}
 \|K^*QK\|_{\rm prod}
 &\le\Lambda_T(K)\|T\|_{\rm prod},
 \label{eq:V111-filtered-product}\\
 \|\operatorname{Tr}_{\mathsf b}(K^*QK)\|_{\rm op}
 &\le
 \Lambda_T(K)
 \|\operatorname{Tr}_{\mathsf b}T\|_{\rm op},
 \label{eq:V111-filtered-left}\\
 \|\operatorname{Tr}_{\mathsf a}(K^*QK)\|_{\rm op}
 &\le
 \Lambda_T(K)
 \|\operatorname{Tr}_{\mathsf a}T\|_{\rm op}.
 \label{eq:V111-filtered-right}
\end{align}
\end{corollary}

\begin{proof}
The first order inequality follows from
\[
 K^*(I-Q)K\ge0;
\]
the second is \Cref{thm:V111-Douglas-constant}.  Positive order is
preserved by partial trace and by expectation in every product state.
Taking the appropriate suprema proves the three bounds.
\end{proof}

\begin{theorem}[Coherent two-row upstream bound]
\label{thm:V111-coherent-upstream}
For the V110 row amplitudes write
\[
 K_b:=\mathbf K_{ab}^{\rhd},
 \qquad
 K_c:=\mathbf K_{ac}^{\lhd},
 \qquad
 K:=K_b+K_c.
\]
If both \(\Lambda_T(K_b)\) and \(\Lambda_T(K_c)\) are finite, then
\begin{align}
 \Lambda_T(K)
 &\le
 \left(
 \sqrt{\Lambda_T(K_b)}
 +\sqrt{\Lambda_T(K_c)}
 \right)^2,
 \label{eq:V111-coherent-Douglas}\\
 K^*QK
 &\preccurlyeq
 \left(
 \sqrt{\Lambda_T(K_b)}
 +\sqrt{\Lambda_T(K_c)}
 \right)^2T.
 \label{eq:V111-common-upstream-order}
\end{align}
The common signed sum is formed before its norm; no positive covariance
square is declared to be an independent physical event.
\end{theorem}

\begin{proof}
Both amplitudes vanish on \(\ker T\).  On the support of \(T\),
\[
 KT^{\dagger/2}
 =
 K_bT^{\dagger/2}+K_cT^{\dagger/2}.
\]
Use the operator-norm triangle inequality and square.  Then apply
\Cref{cor:V111-filtered-domination} to \(K\).
\end{proof}

\begin{remark}[Exact cancellation is retained when available]
\label{rem:V111-cancellation-retained}
The right side of \eqref{eq:V111-coherent-Douglas} is an upper bound, not
the definition of \(\Lambda_T(K)\).  If
\(K_bT^{\dagger/2}=-K_cT^{\dagger/2}\), then the exact common constant is
zero even though both separated constants are positive.  Thus the
programme should compute \(\Lambda_T(K)\) directly whenever the signed
Racah kernel is accessible.
\end{remark}

\section{Exact insertion of the V106 marginal theorem}
\label{sec:V111-V106-insertion}

\begin{corollary}[Marginal product capacity for the common effect]
\label{cor:V111-common-marginal-capacity}
Suppose the upstream carrier has the V106 form
\[
 T=E\otimes M
\]
on the retained row bipartition.  Put
\[
 \Lambda_{\rm com}:=\Lambda_T(K_b+K_c).
\]
Then every downstream filter \(0\le Q\le I\) satisfies
\begin{equation}
 \boxed{
 \|K^*QK\|_{\rm prod}
 \le
 \Lambda_{\rm com}
 \min\left\{
 \Pi(M),\,
 \frac{\mathcal H(M)}{d_A}\alpha(E),\,
 \frac{\mathcal H(M)}{d_B}\beta(E)
 \right\}.}
 \label{eq:V111-common-V106}
\end{equation}
The same statement holds with the two-row upper bound
\eqref{eq:V111-coherent-Douglas} in place of
\(\Lambda_{\rm com}\).
\end{corollary}

\begin{proof}
Use \eqref{eq:V111-filtered-product} and then
\Cref{thm:V106-marginal-product-cap}.  The alternative constant follows
from \eqref{eq:V111-coherent-Douglas}.
\end{proof}

\begin{corollary}[Downstream events do not create a second Douglas loss]
\label{cor:V111-no-second-Douglas-loss}
If \(Q=Q_t\) is any V108 event filter under the same upstream carrier,
then \eqref{eq:V111-common-V106} holds with the same
\(\Lambda_{\rm com}\) for every \(t\).  Conditional normalization of the
event is neither needed nor permitted.
\end{corollary}

\begin{proof}
Only \(0\le Q_t\le I\) was used in
\eqref{eq:V111-filtered-order}.  The event probability remains inside the
unnormalized effect, exactly as in
\Cref{fire:V108-no-conditional-normalization}.
\end{proof}

\section{Raw projection blocks and the order of level integration}
\label{sec:V111-raw-blocks}

\begin{lemma}[Weighted projection formula]
\label{lem:V111-weighted-projection}
Let
\[
 T_\lambda=\tau_\lambda P_\lambda,
 \qquad
 \tau_\lambda>0,
\]
where \(P_\lambda\) is an orthogonal projection.  If
\(K_\lambda=K_\lambda P_\lambda\), then
\begin{equation}
 \boxed{
 \Lambda_{T_\lambda}(K_\lambda)
 =
 \frac{\|K_\lambda\|_{\rm op}^2}{\tau_\lambda}.}
 \label{eq:V111-weighted-projection}
\end{equation}
Consequently, for \(0\le Q\le I\),
\begin{equation}
 K_\lambda^*QK_\lambda
 \preccurlyeq
 \|K_\lambda\|_{\rm op}^2P_\lambda.
 \label{eq:V111-projection-effect}
\end{equation}
\end{lemma}

\begin{proof}
On the support,
\[
 T_\lambda^{\dagger/2}
 =\tau_\lambda^{-1/2}P_\lambda.
\]
Insert this in \eqref{eq:V111-Douglas-constant}, then use
\eqref{eq:V111-filtered-order}.
\end{proof}

\begin{corollary}[Computable common bound on a V57 raw block]
\label{cor:V111-raw-common-bound}
Let the V57 aggregate pinching retain one raw block \(P_\lambda\), and let
\[
 K_{b,\lambda}:=P_\lambda K_bP_\lambda,
 \qquad
 K_{c,\lambda}:=P_\lambda K_cP_\lambda.
\]
If the local Wilson influence bounds of
\eqref{eq:V110-individual-cut-bounds} hold on that block, then
\begin{align}
 &\Lambda_{\tau_\lambda P_\lambda}
 (K_{b,\lambda}+K_{c,\lambda})
 \notag\\
 &\qquad\le
 \frac{C s_\alpha^2}
      {\tau_\lambda}
 \left(\widehat H_{ab}+\widehat H_{ac}\right)^2,
 \label{eq:V111-raw-Douglas-bound}\\
 &(K_{b,\lambda}+K_{c,\lambda})^*
 Q
 (K_{b,\lambda}+K_{c,\lambda})
 \notag\\
 &\qquad\preccurlyeq
 C s_\alpha^2
 \left(\widehat H_{ab}+\widehat H_{ac}\right)^2P_\lambda.
 \label{eq:V111-raw-common-effect}
\end{align}
The estimate includes the V110 interference and has no channel-count
factor.
\end{corollary}

\begin{proof}
Both pinched amplitudes have right support in \(P_\lambda\).  Apply
\Cref{lem:V111-weighted-projection} and the triangle inequality to their
operator norms.  Use \eqref{eq:V110-individual-cut-bounds}; then apply
\eqref{eq:V111-projection-effect}.
\end{proof}

\begin{firewall}[The projection estimate is not the marked-path tax]
\label{fire:V111-projection-not-path}
Equation \eqref{eq:V111-raw-common-effect} is a rigorous positive
common-effect bound, but it squares the aggregate cut size.  It does not
contain the V109 geometric pair denominator and therefore does not by
itself prove the required marked-tree path weight.  The V109 statement is
a trace-class coefficient tax.  An order domination such as
\eqref{eq:V111-Douglas-order} is stronger and cannot be inferred merely by
renaming that scalar functional.
\end{firewall}

\begin{proposition}[Orthogonal level integration before inversion]
\label{prop:V111-levelwise-before-inversion}
Suppose
\[
 T=\sum_{\lambda}\tau_\lambda P_\lambda,
 \qquad
 P_\lambda P_\mu=0\quad(\lambda\ne\mu),
\]
and the complete downstream construction is pinched in these same blocks,
so that
\[
 K^*QK
 =
 \sum_\lambda K_\lambda^*Q_\lambda K_\lambda,
 \qquad
 K_\lambda=K_\lambda P_\lambda,
 \qquad
 0\le Q_\lambda\le I.
 \label{eq:V111-block-diagonal-filter}
\]
Then each level may be bounded by
\eqref{eq:V111-filtered-product} with its own
\(\Lambda_{\tau_\lambda P_\lambda}(K_\lambda)\), followed by summation.
Using the integrated carrier first would instead insert
\begin{equation}
 \Lambda_T(K)
 =
 \max_{\lambda:\,K_\lambda\ne0}
 \frac{\|K_\lambda\|_{\rm op}^2}{\tau_\lambda},
 \label{eq:V111-integrated-max-loss}
\end{equation}
and can lose the smallest level weight.
\end{proposition}

\begin{proof}
Under \eqref{eq:V111-block-diagonal-filter}, positivity and product-state
expectation are additive across the orthogonal blocks.  Apply
\Cref{lem:V111-weighted-projection} to each term before summing.  On the
other hand,
\[
 T^{\dagger/2}
 =\sum_\lambda\tau_\lambda^{-1/2}P_\lambda,
\]
and orthogonality gives \eqref{eq:V111-integrated-max-loss}.
\end{proof}

\begin{firewall}[A downstream filter must preserve the level pinching]
\label{fire:V111-filter-level-mixing}
The hypothesis \eqref{eq:V111-block-diagonal-filter} is essential.  If
\(Q\) mixes two raw levels, \(K^*QK\) contains cross-level terms and the
levelwise positive sum is false.  In that case one must either retain the
larger common carrier and compute its Douglas constant, or insert the
physical raw pinching before the filter.  This is the positive-effect
analogue of the aggregate-order requirement in V57.
\end{firewall}

\section{Sharp regressions for ordinary norm and support}
\label{sec:V111-regressions}

\begin{proposition}[Small covariance norm does not control the Douglas
constant]
\label{prop:V111-small-eigenvalue-regression}
On \(\mathbb C^2\), let
\[
 T_\delta=
 \begin{pmatrix}\delta^4&0\\0&1\end{pmatrix},
 \qquad
 K_\delta=
 \begin{pmatrix}\delta&0\\0&0\end{pmatrix},
 \qquad
 0<\delta<1.
\]
Then
\begin{equation}
 \|K_\delta\|_{\rm op}=\delta\longrightarrow0,
 \qquad
 \Lambda_{T_\delta}(K_\delta)=\delta^{-2}
 \longrightarrow\infty.
 \label{eq:V111-small-eigenvalue-regression}
\end{equation}
Thus an operator covariance cap cannot replace the upstream-weighted
constant.
\end{proposition}

\begin{proof}
Directly,
\[
 K_\delta T_\delta^{-1/2}
 =
 \begin{pmatrix}\delta^{-1}&0\\0&0\end{pmatrix}.
\]
Square its operator norm and apply
\eqref{eq:V111-Douglas-constant}.
\end{proof}

\begin{proposition}[Support leakage makes domination impossible]
\label{prop:V111-support-regression}
Let \(T=P_1\) and \(K=\varepsilon P_0\) for two nonzero orthogonal
projections \(P_0P_1=0\).  Then
\[
 \|K\|_{\rm op}=\varepsilon,
 \qquad
 \Lambda_T(K)=+\infty
\]
for every \(\varepsilon>0\).
\end{proposition}

\begin{proof}
Here \(P_T=P_1\), while \(K\ne KP_1\).  Equivalently, a vector in
\(\operatorname{Ran}P_0\) has zero \(T\)-quadratic form and positive
\(K^*K\)-quadratic form, so no finite
\eqref{eq:V111-Douglas-order} is possible.
\end{proof}

\begin{assessment}[Exact applicability verdict for V106--V108]
\label{ass:V111-applicability-verdict}
The V106 marginal theorem and V108 downstream monotonicity do apply to the
V110 common effect when all of the following data are retained:
\begin{enumerate}[label=\textup{(\roman*)}]
\item the exact upstream carrier \(T=E\otimes M\);
\item the raw support identity \(K=KP_T\);
\item a finite, scale-compatible value of \(\Lambda_T(K)\); and
\item either a common raw pinching respected by the downstream filter or a
      direct computation on the unpinched common carrier.
\end{enumerate}
V110 supplies neither a spectral floor for an integrated \(T\) nor an
order version of its V109 pair tax.  Therefore the unpaid diagonal cannot
yet be closed merely by citing V106--V108.  The obstruction is now the
explicit number \(\Lambda_T(K)\), not an unspecified failure of
factorization.
\end{assessment}

\section{The next pair-tax/Douglas interface}
\label{sec:V111-next-interface}

For one raw level \(\lambda\), let
\(\mathscr S_L\) denote the same-pair coefficient compression whose
trace-class response is estimated in
\eqref{eq:V109-geometric-tax}.  The exact stronger acquisition required
for positive-effect removal is
\begin{equation}
 \left\|
 \mathscr S_L(K_{ar,\lambda})
 (\tau_\lambda P_\lambda)^{\dagger/2}
 \right\|_{\rm op}^2
 \le
 \mathfrak D_{ar,\lambda},
 \label{eq:V111-pair-Douglas-target}
\end{equation}
where the level sum of
\(\mathfrak D_{ar,\lambda}\|T_\lambda\|_{\rm prod}\) has the appropriate
marked-path stock.  The inverse square root and the sector weight
\(\tau_\lambda\) must remain together.

\begin{lemma}[A trace tax does not imply the Douglas target]
\label{lem:V111-trace-not-Douglas}
No dimension-free implication from
\[
 \|\mathscr S_L(K)\|_1\le A
\]
to a bound for the left side of
\eqref{eq:V111-pair-Douglas-target} is possible without a lower spectral
bound for \(T_\lambda\) on the support of \(\mathscr S_L(K)\).
\end{lemma}

\begin{proof}
Use the matrices in
\Cref{prop:V111-small-eigenvalue-regression}.  Their trace and operator
norms both tend to zero, while the Douglas constant diverges.  Multiplying
by a fixed spectator or embedding in a larger coefficient block does not
change the conclusion.
\end{proof}

\begin{corollary}[Projection-level reduction of the new target]
\label{cor:V111-projection-target}
If \(T_\lambda=\tau_\lambda P_\lambda\) and
\(\mathscr S_L(K_{ar,\lambda})
=\mathscr S_L(K_{ar,\lambda})P_\lambda\), then
\eqref{eq:V111-pair-Douglas-target} is exactly
\begin{equation}
 \|\mathscr S_L(K_{ar,\lambda})\|_{\rm op}^2
 \le
 \tau_\lambda\mathfrak D_{ar,\lambda}.
 \label{eq:V111-projection-Douglas-target}
\end{equation}
This is the levelwise inequality that must be tested against the physical
pair spectrum; summing the V109 trace estimate first is insufficient.
\end{corollary}

\begin{proof}
Apply \Cref{lem:V111-weighted-projection}.
\end{proof}

\section{V111 ledger and V112 acquisition}
\label{sec:V111-ledger}

\begin{theorem}[Integrated V111 statement]
\label{thm:V111-integrated}
Preserving V1--V110, V111 proves:
\begin{enumerate}[label=\textup{(V111.\arabic*)},leftmargin=6.7em]
\item the optimal upstream Douglas constant
      \eqref{eq:V111-Douglas-constant};
\item equivalence between finite support-weighted factorization and the
      order bound \eqref{eq:V111-Douglas-order};
\item inheritance of product capacity and both marginal heights by every
      downstream filtered square;
\item the coherent common-row order estimate
      \eqref{eq:V111-common-upstream-order};
\item exact insertion of the V106 marginal minimum without conditional
      event normalization;
\item the weighted raw-projection formula
      \eqref{eq:V111-weighted-projection};
\item a common-effect bound which retains V110 interference on each raw
      block;
\item the levelwise-before-inversion rule
      \eqref{eq:V111-integrated-max-loss};
\item sharp small-eigenvalue and support-leakage regressions; and
\item the precise pair-tax/Douglas target
      \eqref{eq:V111-pair-Douglas-target}.
\end{enumerate}
\end{theorem}

\begin{proof}
These are
\Cref{thm:V111-Douglas-constant,
cor:V111-filtered-domination,
thm:V111-coherent-upstream,
cor:V111-common-marginal-capacity,
cor:V111-no-second-Douglas-loss,
lem:V111-weighted-projection,
cor:V111-raw-common-bound,
prop:V111-levelwise-before-inversion,
prop:V111-small-eigenvalue-regression,
prop:V111-support-regression,
lem:V111-trace-not-Douglas,
cor:V111-projection-target}.
\end{proof}

\begin{openproblem}[V112 mandate: physical levelwise Douglas law]
\label{op:V112-mandate}
Return to the same-pair spectral resolution before its level integral.
\begin{enumerate}
\item Write the exact physical level weight
      \(T_\lambda=\tau_\lambda P_\lambda\) and the pinched row amplitudes
      \(P_\lambda K_bP_\lambda\),
      \(P_\lambda K_cP_\lambda\).
\item Determine whether the V53 same-pair compression commutes with
      \(P_\lambda\), or print its support leakage block
      \((I-P_\lambda)\mathscr S_L(K)P_\lambda\).
\item If support is retained, estimate the operator norm in
      \eqref{eq:V111-projection-Douglas-target} using the actual Wilson
      multiplier and local \(6j\) coefficient, not the integrated trace
      tax.
\item Test the rank-one spectrum of
      \Cref{prop:V109-bi-heavy-rank-one}.  Decide whether its physical
      level weight cancels the inverse \(\tau_\lambda\) or saturates the
      V111 small-eigenvalue regression.
\item Insert any successful level estimates into the common signed sum
      before applying \eqref{eq:V111-common-V106}.
\item Only if the whole common effect is paid, remove its physical outcome
      once by \Cref{cor:V110-V94-reinsertion}; otherwise return upstream to
      the pair kernel.
\end{enumerate}
\end{openproblem}

\section{Firewall after V111}
\label{sec:V111-firewall}

\begin{firewall}[Current claim boundary]
\label{fire:V111-current-boundary}
V111 gives the necessary and sufficient order criterion for applying the
V106--V108 product-capacity machinery to the complete V110 covariance
effect.  It proves that raw projection support makes the criterion
computable and that ordinary covariance or trace-norm smallness does not.
It does not prove the physical levelwise pair-Douglas estimate, pay the
one-row common effect, improve the same-pair spectral threshold, combine
the two noncommuting orientations, close the off-cut branch, prove JREC,
construct the continuum theory, or prove a positive mass gap.  The full
Yang--Mills existence and mass gap remain open.
\end{firewall}

\section{Append-only record for V111}
\label{sec:V111-record}

V111 was appended to the complete V110 manuscript.  The exact V110 parent
had \(75099\) newline-delimited source records, \(2618860\) bytes, and
SHA--256 digest
\[
 \texttt{6cd86af4f6ef0d9815fa2bafb6cf1379af5192c5ccf8d58b4741858eb785ea94}.
\]
No V110 mathematical content was changed.  The sole final
\verb|\end{document}| is restored below.

% =============================================================================
% END APPENDED VERSION V111 MATERIAL
% =============================================================================
% =============================================================================
% BEGIN APPENDED VERSION V112 MATERIAL
% =============================================================================

\chapter{V112: Positive Square of the Pair Multiplier and
Common-Effect Interference Closure}
\label{ch:V112}

\section{Purpose and correction of the V111 target}
\label{sec:V112-purpose}

V111 asked whether the V109 trace-class pair tax could be upgraded to a
small upstream Douglas constant.  The exact V52 compression map shows that
this is the wrong location for the dimension gain.  On each pair output
the Clebsch--Gordan compression is dual to an orthogonal input projector.
The small factor is the \emph{product capacity of the aggregate pair
projector}, not the ordinary operator norm of the multiplier on its
support.  A singlet projector has operator norm one and product capacity
\(1/d\), so no small Douglas constant should be expected.

V112 keeps this geometry intact.  It forms the positive aggregate
multiplier \(\mathbf A_L\) whose eigenvalues are the V52 absolute
formation--defect weights.  The square of the actual signed multiplier is
order bounded by \(A_*\mathbf A_L\).  Therefore every downstream filtered
square inherits the V53 product-capacity tax with no inverse level weight.
When the same physical pair multiplier is common to the two V110 row
amplitudes, this also controls their interference.  What remains is no
longer a positive-effect problem: it is the endpoint mismatch between that
one pair denominator and the cut numerator of the other row.

\section{The pair tax as a positive multiplier capacity}
\label{sec:V112-positive-pair}

Use the V52 pair input \(H_R\otimes H_Q\).  For each
multiplicity-resolved output \(\lambda=(S,\mu)\), write
\[
 P_\lambda:=\iota_{S,\mu}\iota_{S,\mu}^*.
\]
These projections are mutually orthogonal.  Let \(d_\lambda\) be the
actual signed weighted defect coefficient and let \(a_\lambda=a(S)\) be
the nonnegative V52 majorant, so that
\begin{equation}
 |d_\lambda|\le a_\lambda\le A_*.
 \label{eq:V112-defect-majorization}
\end{equation}
Define
\begin{equation}
 \mathbf D_L:=\sum_\lambda d_\lambda P_\lambda,
 \qquad
 \mathbf A_L:=\sum_\lambda a_\lambda P_\lambda.
 \label{eq:V112-pair-multipliers}
\end{equation}

\begin{lemma}[Exact dual form of the V52 functional on positive inputs]
\label{lem:V112-V52-dual-positive}
For positive trace-class \(A,B\),
\begin{equation}
 \operatorname{Tr}\!\left[
 \mathbf A_L(A\otimes B)
 \right]
 =
 \sum_\lambda
 a_\lambda
 \|\mathcal C_\lambda(A,B)\|_1
 =
 \mathfrak F_L(A,B).
 \label{eq:V112-positive-duality}
\end{equation}
Consequently every V53 bound
\begin{equation}
 \mathfrak F_L(A,B)
 \le\kappa_L\|A\|_1\|B\|_1
 \label{eq:V112-functional-kappa}
\end{equation}
is exactly the product-capacity statement
\begin{equation}
 \boxed{\|\mathbf A_L\|_{\rm prod}\le\kappa_L.}
 \label{eq:V112-A-product-capacity}
\end{equation}
The same equality holds after arbitrary identity spectators are adjoined
and then traced out within the two retained rows.
\end{lemma}

\begin{proof}
For positive \(A,B\), the compressed matrix
\[
 \mathcal C_\lambda(A,B)
 =
 \iota_\lambda^*(A\otimes B)\iota_\lambda
\]
is positive.  Its trace norm is therefore its trace.  Cyclicity gives
\[
 \|\mathcal C_\lambda(A,B)\|_1
 =
 \operatorname{Tr}[P_\lambda(A\otimes B)].
\]
Multiply by \(a_\lambda\), sum, and use
\eqref{eq:V112-pair-multipliers}.  Taking positive trace-one \(A,B\)
proves the capacity statement.  Conversely those are precisely the
positive inputs in the product-capacity supremum.  Spectators are removed
by the partial traces within each retained row, as in
\Cref{lem:V94-identity-spectator-quotient}.
\end{proof}

\begin{theorem}[Positive square of the signed pair multiplier]
\label{thm:V112-positive-square}
For every positive contraction \(0\le Q\le I\) on the pair carrier,
\begin{equation}
 0\le
 \mathbf D_L^*Q\mathbf D_L
 \preccurlyeq
 \mathbf D_L^*\mathbf D_L
 \preccurlyeq
 \mathbf A_L^2
 \preccurlyeq
 A_*\mathbf A_L.
 \label{eq:V112-pair-square-order}
\end{equation}
Hence
\begin{equation}
 \boxed{
 \|\mathbf D_L^*Q\mathbf D_L\|_{\rm prod}
 \le A_*\kappa_L.}
 \label{eq:V112-filtered-square-capacity}
\end{equation}
No lower spectral bound for \(\mathbf A_L\) and no small Douglas constant
are required.
\end{theorem}

\begin{proof}
The first two inequalities follow from \(0\le Q\le I\).  All
\(P_\lambda\) are orthogonal, so
\[
 \mathbf D_L^*\mathbf D_L
 =
 \sum_\lambda|d_\lambda|^2P_\lambda
 \preccurlyeq
 \sum_\lambda a_\lambda^2P_\lambda
 =
 \mathbf A_L^2.
\]
Finally \(0\le a_\lambda\le A_*\) gives
\(\mathbf A_L^2\preccurlyeq A_*\mathbf A_L\).
Take product-state expectations and use
\eqref{eq:V112-A-product-capacity}.
\end{proof}

\begin{corollary}[Bi-heavy geometric capacity survives squaring]
\label{cor:V112-bi-heavy-square}
On a V109 bi-heavy same-pair bank joining endpoints \(i,k\),
\begin{equation}
 \boxed{
 \|\mathbf D_L^*Q\mathbf D_L\|_{\rm prod}
 \le
 C_\varepsilon A_*
 \frac{h}{s_\alpha\Xi_i\Xi_k}.}
 \label{eq:V112-bi-heavy-square}
\end{equation}
\end{corollary}

\begin{proof}
For positive trace-class inputs,
\Cref{thm:V109-bi-heavy-tax} gives
\[
 \kappa_L
 \le C_\varepsilon
 \frac{h}{s_\alpha\Xi_i\Xi_k}.
\]
Insert this in \eqref{eq:V112-filtered-square-capacity}.
\end{proof}

\begin{remark}[Why the V111 small-eigenvalue regression is avoided]
\label{rem:V112-avoid-Douglas-loss}
V111 would compare \(\mathbf D_L^2\) with a weighted effect by dividing
through its small eigenvalues.  Equation
\eqref{eq:V112-pair-square-order} instead compares it with
\(\mathbf A_L\) itself.  The scalar inequality
\(a_\lambda^2\le A_*a_\lambda\) cancels the dangerous inverse weight
before any supremum is taken.
\end{remark}

\begin{proposition}[Dual-singlet sharp check]
\label{prop:V112-singlet-check}
Let \(P_\Omega\) be a normalized dual-singlet projector on
\(\mathbb C^d\otimes\mathbb C^d\), and let
\(\mathbf D_L=d_0P_\Omega\),
\(\mathbf A_L=|d_0|P_\Omega\).  Then
\begin{align}
 \|\mathbf A_L\|_{\rm op}&=|d_0|,
 &
 \|\mathbf A_L\|_{\rm prod}&=\frac{|d_0|}{d},
 \label{eq:V112-singlet-A}\\
 \|\mathbf D_L^*\mathbf D_L\|_{\rm prod}
 &=\frac{|d_0|^2}{d}.
 \label{eq:V112-singlet-square}
\end{align}
Thus the dimension gain survives the positive square even though the
Douglas constant relative to \(P_\Omega\) is \(|d_0|^2\), with no
\(1/d\) factor.
\end{proposition}

\begin{proof}
The largest squared Schmidt coefficient of the dual singlet is \(1/d\),
so \(\|P_\Omega\|_{\rm prod}=1/d\).  The formulas follow by scalar
multiplication.  Equality in
\eqref{eq:V112-pair-square-order} holds with
\(A_*=|d_0|\).
\end{proof}

\section{A common pair multiplier absorbs the V110 interference}
\label{sec:V112-common-pair}

Work before final recoupling in the pair-first order of V108.  The active
same-pair bank and the punctured covariance core occupy distinct
coordinate tensor factors.  Let
\(\mathcal R_b,\mathcal R_c\) denote the two signed core row operators
after the exact V110 hybrid split, with all remaining fixed coefficient
factors retained as spectators.  Thus on a fixed same-pair outcome,
\begin{equation}
 K_b=\mathbf D_L\otimes\mathcal R_b,
 \qquad
 K_c=\mathbf D_L\otimes\mathcal R_c.
 \label{eq:V112-common-pair-factor}
\end{equation}
This is one occurrence of the physical pair multiplier, common to the two
covariance rows; it is not two copies of the pair defect.

\begin{theorem}[Common-pair filtered effect bound]
\label{thm:V112-common-pair-effect}
Assume
\begin{align}
 \|\mathcal R_b\|_{\rm op}
 &\le C\min\{s_\alpha\widehat H_{ab},1\},
 \label{eq:V112-Rb-cap}\\
 \|\mathcal R_c\|_{\rm op}
 &\le C\min\{s_\alpha\widehat H_{ac},1\}.
 \label{eq:V112-Rc-cap}
\end{align}
For every positive downstream contraction \(0\le Q\le I\) on the complete
pair--core output,
\begin{align}
 &(K_b+K_c)^*Q(K_b+K_c)
 \notag\\
 &\quad\preccurlyeq
 C\min\!\left\{
 s_\alpha(\widehat H_{ab}+\widehat H_{ac}),1
 \right\}^{\!2}
 \bigl(\mathbf D_L^*\mathbf D_L\otimes I_{\rm core}\bigr).
 \label{eq:V112-common-effect-order}
\end{align}
Consequently
\begin{align}
 &\|(K_b+K_c)^*Q(K_b+K_c)\|_{\rm prod}
 \notag\\
 &\quad\le
 C A_*\kappa_L
 \min\!\left\{
 s_\alpha(\widehat H_{ab}+\widehat H_{ac}),1
 \right\}^{\!2}.
 \label{eq:V112-common-effect-capacity}
\end{align}
The bound contains the complete V110 off-diagonal interference.
\end{theorem}

\begin{proof}
Equation \eqref{eq:V112-common-pair-factor} gives
\[
 K_b+K_c
 =
 \mathbf D_L\otimes(\mathcal R_b+\mathcal R_c).
\]
Since \(Q\le I\),
\begin{align*}
 (K_b+K_c)^*Q(K_b+K_c)
 &\preccurlyeq
 \mathbf D_L^*\mathbf D_L
 \otimes
 (\mathcal R_b+\mathcal R_c)^*
 (\mathcal R_b+\mathcal R_c)\\
 &\preccurlyeq
 \|\mathcal R_b+\mathcal R_c\|_{\rm op}^2
 (\mathbf D_L^*\mathbf D_L\otimes I).
\end{align*}
For \(x,y\ge0\),
\[
 \min\{x,1\}+\min\{y,1\}
 \le2\min\{x+y,1\}.
\]
Use this with \eqref{eq:V112-Rb-cap}--%
\eqref{eq:V112-Rc-cap} to obtain
\eqref{eq:V112-common-effect-order}.  The identity-spectator quotient
\eqref{eq:V94-identity-spectator-quotient} and
\eqref{eq:V112-filtered-square-capacity}, with \(Q=I\) there, prove
\eqref{eq:V112-common-effect-capacity}.
\end{proof}

\begin{corollary}[Bi-heavy common-effect path bound]
\label{cor:V112-bi-heavy-common-path}
If the common pair bank is bi-heavy at endpoints \(i,k\), then
\begin{align}
 &\|(K_b+K_c)^*Q(K_b+K_c)\|_{\rm prod}
 \notag\\
 &\qquad\le
 C_\varepsilon
 \frac{h(\widehat H_{ab}+\widehat H_{ac})}
      {\Xi_i\Xi_k}.
 \label{eq:V112-bi-heavy-common-path}
\end{align}
\end{corollary}

\begin{proof}
Insert
\[
 \kappa_L
 \le
 C_\varepsilon\frac{h}{s_\alpha\Xi_i\Xi_k}
\]
from \Cref{thm:V109-bi-heavy-tax} into
\eqref{eq:V112-common-effect-capacity}.  For \(z\ge0\),
\(\min\{z,1\}^2\le z\).  Take
\[
 z=s_\alpha(\widehat H_{ab}+\widehat H_{ac})
\]
and absorb the universal \(A_*\).
\end{proof}

\begin{remark}[One pair tax was used, not squared]
\label{rem:V112-one-pair-tax}
The right side of \eqref{eq:V112-bi-heavy-common-path} contains one factor
\(h/(s_\alpha\Xi_i\Xi_k)\).  The two covariance rows were added inside the
same core norm, and the same pair multiplier was squared only after that
addition.  Its positive square was then bounded by
\(A_*\mathbf A_L\).  Thus neither the pair mass nor its product-capacity
tax was charged twice.
\end{remark}

\section{Exact endpoint mismatch and a closed comparable region}
\label{sec:V112-endpoint-mismatch}

Consider first a same-pair bank with endpoints \(a,b\).  Then
\eqref{eq:V112-bi-heavy-common-path} reads
\begin{equation}
 C_\varepsilon h
 \left[
 \frac{\widehat H_{ab}}{\Xi_a\Xi_b}
 +
 \frac{\widehat H_{ac}}{\Xi_a\Xi_b}
 \right].
 \label{eq:V112-ab-common-bound}
\end{equation}
The first term is the matched V57 row weight.  The second has the correct
cut numerator but the pair bank supplies \(\Xi_b\), whereas the
\(ac\)-allocated V57 path requests \(\Xi_c\).

\begin{theorem}[Row-comparable same-pair outcomes are fully paid]
\label{thm:V112-row-comparable}
Suppose the common same-pair bank joins \(a,b\) and
\begin{equation}
 \Xi_c\le C_{\rm row}\Xi_b.
 \label{eq:V112-ab-comparability}
\end{equation}
Then the complete positive common effect obeys
\begin{equation}
 \boxed{
 \|(K_b+K_c)^*Q(K_b+K_c)\|_{\rm prod}
 \le
 C_{\varepsilon,C_{\rm row}}h
 \left[
 \frac{\widehat H_{ab}}{\Xi_a\Xi_b}
 +
 \frac{\widehat H_{ac}}{\Xi_a\Xi_c}
 \right].}
 \label{eq:V112-ab-two-paths}
\end{equation}
The symmetric statement holds for a common pair \(a,c\) under
\(\Xi_b\le C_{\rm row}\Xi_c\).
\end{theorem}

\begin{proof}
The matched term in \eqref{eq:V112-ab-common-bound} is already the first
term of \eqref{eq:V112-ab-two-paths}.  The comparability assumption gives
\[
 \frac1{\Xi_b}
 \le
 \frac{C_{\rm row}}{\Xi_c},
\]
which converts the second term.  Interchange \(b,c\) for the symmetric
case.
\end{proof}

\begin{corollary}[V94 removal on the comparable same-pair region]
\label{cor:V112-V94-removal}
Every V105 same-pair physical outcome satisfying the corresponding row
comparability condition may be assigned to the paid pair class in the V94
priority partition and removed exactly once from the residual level
effect.  Its two row amplitudes and their interference are paid together
by \eqref{eq:V112-ab-two-paths}.
\end{corollary}

\begin{proof}
Theorem~\ref{thm:V112-row-comparable} bounds the whole positive common
effect.  Apply case \textup{(i)} of the complete-effect principle in
\Cref{cor:V110-V94-reinsertion}; the assignment is made by the pre-existing
disjoint V94 priority rule.
\end{proof}

\begin{proposition}[The remaining mismatch factor is exact]
\label{prop:V112-mismatch-factor}
For a pair \(a,b\), converting the unmatched second term in
\eqref{eq:V112-ab-common-bound} to its \(ac\) row weight costs exactly
\begin{equation}
 \mathfrak m_{b\to c}:=\frac{\Xi_c}{\Xi_b}.
 \label{eq:V112-mismatch-factor}
\end{equation}
No uniform bound on this factor follows from bi-heaviness of the \(ab\)
pair bank.
\end{proposition}

\begin{proof}
The ratio of the unmatched denominator to the target denominator is
\[
 \frac{1/(\Xi_a\Xi_b)}{1/(\Xi_a\Xi_c)}
 =\frac{\Xi_c}{\Xi_b}.
\]
Bi-heaviness gives lower bounds for the pair-bank masses relative to
\(\Xi_a\) and \(\Xi_b\), but contains no upper bound for the mass of the
third row \(c\).  Adjoining active mass to row \(c\) outside the \(ab\)
pair bank leaves those two bi-heavy inequalities unchanged and makes
\(\Xi_c/\Xi_b\) arbitrarily large.
\end{proof}

\begin{assessment}[The interference gate is closed; endpoint transport
remains]
\label{ass:V112-interference-disposition}
For a common physical pair multiplier, V112 pays the positive square and
all V110 interference by the same one-edge capacity used before squaring.
The row-comparable part of the same-pair branch is therefore closed and
removed once.  In the residual regime the exact obstruction is the
unbounded third-row ratio \(\Xi_c/\Xi_b\) (or its symmetric counterpart),
not a missing positive-effect construction and not a channel-count loss.
\end{assessment}

\section{V112 ledger and V113 acquisition}
\label{sec:V112-ledger}

\begin{theorem}[Integrated V112 statement]
\label{thm:V112-integrated}
Preserving V1--V111, V112 proves:
\begin{enumerate}[label=\textup{(V112.\arabic*)},leftmargin=6.7em]
\item exact identification of the V52 positive pair multiplier with the
      product-capacity form of its trace functional;
\item the signed-square order bound
      \eqref{eq:V112-pair-square-order};
\item survival of the V53 and V109 one-edge gains under positive squaring
      and every downstream contraction;
\item the dual-singlet sharp check
      \eqref{eq:V112-singlet-square};
\item factorization of both V110 row amplitudes through one common physical
      pair multiplier;
\item the complete common-effect estimate
      \eqref{eq:V112-common-effect-capacity}, including interference;
\item the bi-heavy common path bound
      \eqref{eq:V112-bi-heavy-common-path};
\item full payment and V94 removal of the row-comparable same-pair region;
      and
\item the exact residual endpoint ratio
      \eqref{eq:V112-mismatch-factor}.
\end{enumerate}
\end{theorem}

\begin{proof}
These are
\Cref{lem:V112-V52-dual-positive,
thm:V112-positive-square,
cor:V112-bi-heavy-square,
prop:V112-singlet-check,
thm:V112-common-pair-effect,
cor:V112-bi-heavy-common-path,
thm:V112-row-comparable,
cor:V112-V94-removal,
prop:V112-mismatch-factor}.
\end{proof}

\begin{openproblem}[V113 mandate: third-row incidence transport]
\label{op:V113-mandate}
Return upstream to the V59 incidence partition before any Racah level
supremum.
\begin{enumerate}
\item On the \(ab\) same-pair branch, split the third-row mass
      \(\Xi_c\) into its private, \(ac\)-cut, \(bc\)-cut, same-pair, and
      off-cut coordinate banks using the exact V59/V60 alphabet.
\item Show that \(\Xi_c\gg\Xi_b\) forces a fixed fraction into an already
      paid private-mean, cut-heavy, or distinct-pair class, unless the
      explicit off-cut alternative occurs.
\item Perform this assignment at the physical-outcome level before
      invoking \eqref{eq:V112-bi-heavy-common-path}; do not use total
      spectator mass which V94 later quotients away.
\item If the large third-row mass is active in the same common pair
      carrier, refine the pair-support alphabet rather than treating it as
      an identity spectator.
\item Prove that the residual same-pair outcome has
      \(\Xi_c\lesssim\Xi_b\), and then apply
      \Cref{thm:V112-row-comparable}, or identify the exact surviving
      off-cut monomial.
\item Keep the symmetric \(ac\) pair case and the two noncommuting V94
      orientations separate.
\end{enumerate}
\end{openproblem}

\section{Firewall after V112}
\label{sec:V112-firewall}

\begin{firewall}[Current claim boundary]
\label{fire:V112-current-boundary}
V112 proves that the one-edge pair-capacity gain survives the complete
positive common square and closes all row-comparable same-pair outcomes,
including their interference.  It does not yet route an arbitrarily large
third-row mass through the V59/V60 incidence alphabet, close the resulting
off-cut alternative, improve the one-pair spectral threshold in isolation,
combine the two noncommuting orientations, prove JREC, construct the
continuum theory, or prove a positive mass gap.  The full Yang--Mills
existence and mass gap remain open.
\end{firewall}

\section{Append-only record for V112}
\label{sec:V112-record}

V112 was appended to the complete V111 manuscript.  The exact V111 parent
had \(75721\) newline-delimited source records, \(2638955\) bytes, and
SHA--256 digest
\[
 \texttt{0ac44442ca4dc530ca73231dfd5734707ba0a3a2dc165a3330ef8d7600db79be}.
\]
No V111 mathematical content was changed.  The sole final
\verb|\end{document}| is restored below.

% =============================================================================
% END APPENDED VERSION V112 MATERIAL
% =============================================================================
% =============================================================================
% BEGIN APPENDED VERSION V113 MATERIAL
% =============================================================================

\chapter{V113: Racah-Instrument Correction, Physical Absolute
Pair Factorization, and Third-Row Routing}
\label{ch:V113}

\section{Purpose and mandatory internal correction}
\label{sec:V113-purpose}

V112 proved a valid positive-multiplier theorem: the product-capacity gain
of a V52/V53 pair bank survives the positive square of its signed
multiplier.  Its final removal statement, however, crossed an interface
which the earlier V98 theorem had already forbidden.  The V94 Racah
effects are normalized diagnostic POVM effects; on a noncommuting
multiplicity block they are not the physical joint outer pinching of the
coefficient operator.  Thus a bound on the V94 common effect cannot by
itself move a \emph{physical} outcome out of the V98 residual carrier.

V113 corrects that scope explicitly.  The V110 hybrid covariance split and
all V112 pair-multiplier inequalities remain true.  Physical payment is
instead made on the V98 signed outer block.  When its surviving pair and
core labels form a rectangle, the block is exactly a tensor product of the
signed pair multiplier and the complete signed core amplitude.  Its
absolute value then factorizes, so the V53 product-capacity tax pays the
physical carrier directly, with no Kadison square root and no Racah
instrument.  The nonrectangular case remains in the V98 physical image.

\begin{proposition}[Precise correction of the V112 removal statement]
\label{prop:V113-V112-correction}
The following status holds.
\begin{enumerate}[label=\textup{(\roman*)}]
\item The exact row decomposition
      \eqref{eq:V110-exact-row-split}, its individual influence bounds,
      and the abstract common block effect
      \eqref{eq:V110-block-effect} remain valid.
\item The positive pair-multiplier inequalities
      \eqref{eq:V112-pair-square-order}--%
      \eqref{eq:V112-bi-heavy-square} remain valid.
\item The effect bound
      \eqref{eq:V112-ab-two-paths} is a valid bound for the stated
      factorized diagnostic effect.
\item The physical-removal conclusion in
      \Cref{cor:V112-V94-removal} is withdrawn.  That corollary is valid
      only as a removal inside the V94 diagnostic level ledger.
\item A physical coefficient branch may be removed only from the
      orthogonal V98 outer partition
      \eqref{eq:V98-outer-partition}, after its absolute carrier has the
      required bound.
\end{enumerate}
\end{proposition}

\begin{proof}
Items \textup{(i)--(iii)} are algebraic operator statements and do not
assert that the filter \(Q\) is a physical sharp Racah instrument.
\Cref{thm:V98-Racah-instrument-firewall} proves that such an identification
is unavailable on noncommuting intermediate PVMs.
\Cref{thm:V98-exact-absolute-carrier,cor:V98-paid-removal} give the exact
physical replacement, proving \textup{(iv)--(v)}.
\end{proof}

\begin{firewall}[Diagnostic and physical selectors]
\label{fire:V113-two-selectors}
The V94 priority partition is still a valid partition of its diagnostic
level law and can prove a sufficient positive-square estimate.  The V98
outer partition is the physical pre-trace partition.  V113 never identifies
their atoms.  A branch may use the same deterministic incidence label in
both ledgers, but physical removal is certified only in the V98 ledger.
\end{firewall}

\section{Rectangular physical outer blocks}
\label{sec:V113-rectangular}

Fix one V98 branch and one retained-row orientation.  Let
\(\Lambda_L\) be a family of orthogonal outputs of the same physical pair
bank and let \(\Lambda_C\) be a family of complete punctured-core outputs.
Call the outer family \emph{rectangular} if its labels are exactly
\(\Lambda_L\times\Lambda_C\), with no rule correlating the retained pair
label with the core label.  After the physical isometries are used to
identify the direct sum with a tensor product, suppose the signed block is
\begin{equation}
 \mathbf L_{\rm rect}
 =
 \mathbf D_L\otimes\mathbf R_C,
 \label{eq:V113-rectangular-signed}
\end{equation}
where \(\mathbf D_L\) is the actual signed pair multiplier and
\(\mathbf R_C\) is the \emph{complete} signed V110 covariance core on that
branch.

\begin{lemma}[Absolute value of a rectangular signed block]
\label{lem:V113-absolute-tensor}
For arbitrary finite-dimensional operators \(D,R\),
\begin{equation}
 |D\otimes R|=|D|\otimes|R|.
 \label{eq:V113-absolute-tensor}
\end{equation}
Consequently
\begin{equation}
 \mathbf Q_{\rm rect}
 :=
 |\mathbf L_{\rm rect}|
 =
 |\mathbf D_L|\otimes|\mathbf R_C|.
 \label{eq:V113-rectangular-absolute}
\end{equation}
\end{lemma}

\begin{proof}
One has
\[
 (D\otimes R)^*(D\otimes R)
 =
 (D^*D)\otimes(R^*R).
\]
The positive square root of a tensor product of positive matrices is the
tensor product of their positive square roots, proving the identity.
\end{proof}

\begin{theorem}[Direct physical absolute-carrier tax]
\label{thm:V113-direct-physical-tax}
Let \(\mathbf A_L\) be the V112 positive pair multiplier and suppose
\[
 |\mathbf D_L|\preccurlyeq\mathbf A_L,
 \qquad
 \|\mathbf A_L\|_{\rm prod}\le\kappa_L.
\]
Then
\begin{equation}
 \boxed{
 \|\mathbf Q_{\rm rect}\|_{\rm rows,prod}
 \le
 \kappa_L\|\mathbf R_C\|_{\rm op}.}
 \label{eq:V113-direct-absolute-capacity}
\end{equation}
The estimate is for the physical V98 absolute carrier and contains no
square-root loss.
\end{theorem}

\begin{proof}
By \eqref{eq:V113-rectangular-absolute} and positive order,
\[
 0\le\mathbf Q_{\rm rect}
 \le
 \|\mathbf R_C\|_{\rm op}
 (\mathbf A_L\otimes I_C).
\]
Take expectations in original-row product states.  The core identity is an
arbitrary within-row spectator, so
\Cref{lem:V94-identity-spectator-quotient} gives
\[
 \|\mathbf A_L\otimes I_C\|_{\rm rows,prod}
 =
 \|\mathbf A_L\|_{\rm prod}.
\]
Use the assumed capacity bound.
\end{proof}

\begin{remark}[Why V113 does not use the square bridge]
\label{rem:V113-no-square-root}
Applying \eqref{eq:V98-full-row-square-bridge} to a square estimate of size
\(W\) gives only \(\sqrt W\).  The direct absolute identity
\eqref{eq:V113-rectangular-absolute} instead retains the linear
pair-capacity factor.  This is the correct physical use of the V52/V53
tax on a rectangular outer family.
\end{remark}

\begin{theorem}[Physical bi-heavy common-path estimate]
\label{thm:V113-physical-common-path}
Suppose the rectangular pair bank is bi-heavy at endpoints \(i,k\), with
pair mass \(h\), and the complete signed core satisfies
\begin{equation}
 \|\mathbf R_C\|_{\rm op}
 \le
 C\min\!\left\{
 s_\alpha(\widehat H_{ab}+\widehat H_{ac}),1
 \right\}.
 \label{eq:V113-complete-core-cap}
\end{equation}
Then its physical V98 absolute carrier obeys
\begin{equation}
 \boxed{
 \|\mathbf Q_{\rm rect}\|_{\rm rows,prod}
 \le
 C_\varepsilon
 \frac{h(\widehat H_{ab}+\widehat H_{ac})}
      {\Xi_i\Xi_k}.}
 \label{eq:V113-physical-common-path}
\end{equation}
All interference between the two V110 covariance rows remains inside
\(|\mathbf R_C|\).
\end{theorem}

\begin{proof}
\Cref{thm:V109-bi-heavy-tax,lem:V112-V52-dual-positive} give
\[
 \kappa_L
 \le
 C_\varepsilon
 \frac{h}{s_\alpha\Xi_i\Xi_k}.
\]
Insert this and \eqref{eq:V113-complete-core-cap} in
\eqref{eq:V113-direct-absolute-capacity}.  Since
\(\min\{s_\alpha H,1\}\le s_\alpha H\), the \(s_\alpha\) factors cancel.
\end{proof}

\section{Positive completion of a correlated outer subset}
\label{sec:V113-positive-completion}

Rectangularity is sufficient for the equality
\eqref{eq:V113-rectangular-absolute}, but a physical capacity upper bound
needs less.  Let \(\mathfrak J_X\) be any orthogonal V98 outer subfamily.
For every \(j\in\mathfrak J_X\), let \(\lambda(j)\) be its same-pair
output label and let \(P_j:=V_jV_j^*\).  Assume the outer projections refine
the pair resolution:
\begin{equation}
 \sum_{\substack{j\in\mathfrak J_X\\\lambda(j)=\lambda}}P_j
 \preccurlyeq
 P_\lambda\otimes I_C.
 \label{eq:V113-outer-refines-pair}
\end{equation}

\begin{theorem}[No rectangularity needed under a common pair majorant]
\label{thm:V113-positive-completion}
Suppose every physical signed outer block obeys
\begin{equation}
 |V_jL_jV_j^*|
 \preccurlyeq
 m\,a_{\lambda(j)}P_j
 \label{eq:V113-blockwise-pair-majorant}
\end{equation}
with one \(m\ge0\) independent of \(j\).  Then the complete physical
absolute carrier satisfies
\begin{align}
 \mathbf Q_X
 &\preccurlyeq
 m(\mathbf A_L\otimes I_C),
 \label{eq:V113-completed-order}\\
 \boxed{
 \|\mathbf Q_X\|_{\rm rows,prod}
 \le m\kappa_L.}
 \label{eq:V113-completed-capacity}
\end{align}
The outer label set may correlate pair and core outputs arbitrarily.
\end{theorem}

\begin{proof}
The V98 blocks are orthogonal, so their absolute values add:
\[
 \mathbf Q_X
 =
 \sum_{j\in\mathfrak J_X}|V_jL_jV_j^*|.
\]
Apply \eqref{eq:V113-blockwise-pair-majorant}, group by \(\lambda\), and
use \eqref{eq:V113-outer-refines-pair}:
\begin{align*}
 \mathbf Q_X
 &\preccurlyeq
 m\sum_\lambda a_\lambda
 \sum_{j:\lambda(j)=\lambda}P_j\\
 &\preccurlyeq
 m\sum_\lambda a_\lambda(P_\lambda\otimes I_C)
 =
 m(\mathbf A_L\otimes I_C).
\end{align*}
Take product-state expectations, quotient the identity spectators by
\Cref{lem:V94-identity-spectator-quotient}, and use
\eqref{eq:V112-A-product-capacity}.
\end{proof}

\begin{corollary}[Physical same-pair completion]
\label{cor:V113-physical-same-pair-completion}
Suppose a V98 outer branch contains the same physical pair multiplier
exactly once, so each block can be written
\[
 L_j=d_{\lambda(j)}R_j,
 \qquad
 \|R_j\|_{\rm op}
 \le
 C\min\!\left\{
 s_\alpha(\widehat H_{ab}+\widehat H_{ac}),1
 \right\}.
\]
Then \eqref{eq:V113-blockwise-pair-majorant} holds and the physical
capacity obeys \eqref{eq:V113-physical-common-path}, without a rectangular
label assumption.
\end{corollary}

\begin{proof}
On the range \(P_j\),
\[
 |L_j|
 \preccurlyeq
 |d_{\lambda(j)}|\,\|R_j\|_{\rm op}P_j
 \preccurlyeq
 a_{\lambda(j)}\,\|R_j\|_{\rm op}P_j.
\]
Apply \Cref{thm:V113-positive-completion}, then use the bi-heavy bound for
\(\kappa_L\) and cancel \(s_\alpha\) as in
\Cref{thm:V113-physical-common-path}.
\end{proof}

\begin{firewall}[What positive completion does not separate]
\label{fire:V113-completion-no-row-split}
The enlargement in \eqref{eq:V113-completed-order} forgets the correlation
between pair and core labels only after the physical absolute values have
been formed on orthogonal outer blocks.  It is therefore a valid positive
order bound, not the forbidden assertion that the outer family is a
rectangle.  It still uses one pair denominator for the complete core cap;
it does not turn that denominator into two different row denominators.
\end{firewall}

\section{Physical endpoint comparison}
\label{sec:V113-endpoint-comparison}

\begin{theorem}[Physical row-comparable same-pair closure]
\label{thm:V113-physical-row-comparable}
Suppose the same physical pair bank joins \(a,b\), is bi-heavy there, and
\[
 \Xi_c\le C_{\rm row}\Xi_b.
\]
On every outer branch satisfying
\Cref{cor:V113-physical-same-pair-completion},
\begin{equation}
 \boxed{
 \|\mathbf Q_X\|_{\rm rows,prod}
 \le
 C_{\varepsilon,C_{\rm row}}h
 \left[
 \frac{\widehat H_{ab}}{\Xi_a\Xi_b}
 +
 \frac{\widehat H_{ac}}{\Xi_a\Xi_c}
 \right].}
 \label{eq:V113-physical-two-paths}
\end{equation}
The symmetric statement holds for a same pair \(a,c\).
\end{theorem}

\begin{proof}
\Cref{cor:V113-physical-same-pair-completion} gives
\[
 C_\varepsilon h
 \left[
 \frac{\widehat H_{ab}}{\Xi_a\Xi_b}
 +
 \frac{\widehat H_{ac}}{\Xi_a\Xi_b}
 \right].
\]
Use \(1/\Xi_b\le C_{\rm row}/\Xi_c\) in the second term.  Interchange
\(b,c\) for the symmetric statement.
\end{proof}

\begin{corollary}[Exact physical outer-partition refinement]
\label{cor:V113-physical-removal}
Refine the V98 side/off partition by moving every same-pair outer label
which satisfies the hypotheses of
\Cref{thm:V113-physical-row-comparable} into a new paid subfamily.  This is
an orthogonal disjoint refinement, and the corresponding physical
coefficient response has the two V57 marked-path weights.  Each moved
outer block is charged once.
\end{corollary}

\begin{proof}
The outer labels were orthogonal before trace norm by
\eqref{eq:V98-outer-partition}.  Partitioning their index set preserves
orthogonality.  Apply
\Cref{thm:V98-restricted-matrix-tax} with all original rows retained and
use \eqref{eq:V113-physical-two-paths}.  The output-mass and resolvent
allocation is already part of \(L_j\), so the resulting coefficient
weights are the displayed V57 paths.
\end{proof}

\section{Third-row incidence routing}
\label{sec:V113-third-row}

Take the same-pair bank to be \(ab\); the \(ac\) case is symmetric.
After identity spectators have been quotiented, apply the exact V59/V60
decomposition to row \(c\):
\begin{equation}
 \Xi_c
 =
 x_c+Z_{c\leftarrow a}+P_c^{\rm priv}
 +D_c^{\rm off}
 +P_{c,\mu}^{\rm pun}
 +\sum_{j\ne c}D_{cj,\mu}^{\rm pun}.
 \label{eq:V113-c-mass-alphabet}
\end{equation}
There are three possible labels \(j\ne c\) among the other four input
rows.  Every term is a mass on a disjoint coordinate bank for the fixed
puncture monomial.

\begin{lemma}[Quantitative third-row alternative]
\label{lem:V113-third-row-alternative}
Fix \(0<\delta<1/7\).  At least one of the following holds:
\begin{enumerate}[label=\textup{(\roman*)}]
\item \(x_c\ge\delta\Xi_c\);
\item \(Z_{c\leftarrow a}\ge\delta\Xi_c\);
\item \(P_c^{\rm priv}\ge\delta\Xi_c\);
\item \(P_{c,\mu}^{\rm pun}\ge\delta\Xi_c\);
\item \(D_{cj,\mu}^{\rm pun}\ge\delta\Xi_c\) for some \(j\ne c\);
\item
      \begin{equation}
       D_c^{\rm off}\ge(1-7\delta)\Xi_c.
       \label{eq:V113-heavy-off-cut}
      \end{equation}
\end{enumerate}
\end{lemma}

\begin{proof}
If the first five alternatives fail, the six terms other than
\(D_c^{\rm off}\) in \eqref{eq:V113-c-mass-alphabet} have total less than
\(7\delta\Xi_c\): four singleton displayed terms and three puncture-pair
terms.  Subtract from \(\Xi_c\).
\end{proof}

\begin{proposition}[Disposition of the first four alternatives]
\label{prop:V113-third-row-disposition}
In the exact pre-trace incidence ledger:
\begin{enumerate}[label=\textup{(\roman*)}]
\item alternative \textup{(i)} supplies the incident-pair endpoint
      response of \Cref{lem:V59-old-endpoint-response};
\item alternative \textup{(ii)} implies
      \(\widehat H_{ac}\ge a_*\delta\Xi_c\) and is the cut-heavy class of
      \Cref{thm:V59-cut-heavy-promotion};
\item alternative \textup{(iii)} is the core private-mean class of
      \Cref{lem:V59-private-mean-response};
\item alternative \textup{(iv)} is its puncture-mean analogue from
      \Cref{cor:V60-puncture-mean-closed}.
\end{enumerate}
Alternatives \textup{(v)} and \textup{(vi)} are respectively an exact
same-label puncture-pair branch and the genuine V98 off-cut branch; they
are not silently declared paid here.
\end{proposition}

\begin{proof}
The incident pair mass in \textup{(i)} is a fixed fraction of
\(\Xi_c\), which is the selected-heavy endpoint condition; the thermal
case was already included in
\Cref{lem:V59-old-endpoint-response}.  For \textup{(ii)}, use
\eqref{eq:V59-cut-dominates-incidence-mass}.  The two singleton
statements are exactly the cited V59 and V60 results.  The definitions of
\(D_{cj,\mu}^{\rm pun}\) and \(D_c^{\rm off}\) give the final
classification.
\end{proof}

\begin{corollary}[The noncomparable third row has only two new exits]
\label{cor:V113-noncomparable-exits}
Remove the endpoint-paid, cut-heavy, and private-mean outer blocks in the
physical V98 priority order.  On a remaining \(ab\) same-pair block for
which \(\Xi_c>C_{\rm row}\Xi_b\), either:
\begin{enumerate}[label=\textup{(\roman*)}]
\item a puncture pair \(cj\) carries a fixed fraction of \(\Xi_c\); or
\item the genuine off-cut mass \(D_c^{\rm off}\) carries a fixed fraction
      of \(\Xi_c\).
\end{enumerate}
The fixed fractions depend only on the chosen priority thresholds, not on
volume or representation dimensions.
\end{corollary}

\begin{proof}
Apply \Cref{lem:V113-third-row-alternative}.  The first four alternatives
belong to the removed classes by
\Cref{prop:V113-third-row-disposition}.  What remains is
\textup{(v)} or \textup{(vi)}.  The noncomparability assumption records
why \Cref{thm:V113-physical-row-comparable} was not already used; it is not
needed for the finite pigeonhole constant.
\end{proof}

\begin{firewall}[The puncture pair is not yet an independent second tax]
\label{fire:V113-puncture-pair-not-paid}
In case \textup{(i)} of
\Cref{cor:V113-noncomparable-exits}, the \(cj\) bank may share an original
row with the \(ab\) bank or with another collision factor.  It must first
be regrouped by its exact same label, as required in V60, and compared
with the V105 overlap alphabet.  V113 does not multiply its V53 capacity
by the \(ab\) capacity before that comparison.
\end{firewall}

\begin{assessment}[Position after the physical correction]
\label{ass:V113-position}
The physical same-pair problem now has a disjoint disposition:
\begin{enumerate}[label=\textup{(\roman*)}]
\item row-comparable blocks are paid directly in the V98 absolute carrier;
\item noncomparable blocks with an earlier endpoint, cut, or mean
      certificate return to the established paid class;
\item every other noncomparable block has a heavy same-label puncture pair
      or heavy genuine off-cut mass.
\end{enumerate}
The remaining work is therefore a pair-label comparison followed by the
already isolated off-cut carrier, not the V110 interference and not a
conditional Racah measurement.
\end{assessment}

\section{V113 ledger and V114 acquisition}
\label{sec:V113-ledger}

\begin{theorem}[Integrated corrected V113 statement]
\label{thm:V113-integrated}
Preserving V1--V112 subject to the explicit correction
\Cref{prop:V113-V112-correction}, V113 proves:
\begin{enumerate}[label=\textup{(V113.\arabic*)},leftmargin=6.7em]
\item restoration of the V98 Racah-instrument firewall;
\item exact absolute-value factorization on rectangular physical blocks;
\item the direct physical pair-capacity estimate
      \eqref{eq:V113-direct-absolute-capacity};
\item positive completion of arbitrary correlated outer subsets under one
      common pair majorant;
\item the physical bi-heavy common-path estimate
      \eqref{eq:V113-physical-common-path};
\item full physical payment of row-comparable same-pair outer blocks;
\item the exact third-row mass alphabet
      \eqref{eq:V113-c-mass-alphabet}; and
\item reduction of every remaining noncomparable same-pair block to a
      heavy puncture pair or the genuine off-cut branch.
\end{enumerate}
\end{theorem}

\begin{proof}
These are
\Cref{prop:V113-V112-correction,
lem:V113-absolute-tensor,
thm:V113-direct-physical-tax,
thm:V113-positive-completion,
cor:V113-physical-same-pair-completion,
thm:V113-physical-common-path,
thm:V113-physical-row-comparable,
cor:V113-physical-removal,
lem:V113-third-row-alternative,
cor:V113-noncomparable-exits}.
\end{proof}

\begin{openproblem}[V114 mandate: heavy puncture-pair disposition]
\label{op:V114-mandate}
Work on alternative \textup{(i)} of
\Cref{cor:V113-noncomparable-exits} before the off-cut branch.
\begin{enumerate}
\item Regroup the heavy \(cj\) coordinates by one exact pair label across
      the complete puncture monomial.
\item Compare that pair with the \(ab\) same-pair bank at the physical
      coordinate level.  Use the V105 same/complementary overlap alphabet
      only where its hypotheses apply.
\item If the two banks occupy disjoint coordinate sets, use the V104
      two-pair capacity; if they form a complementary same-coordinate
      pair, use the V108 physical subeffect theorem.
\item Prove that the remaining same-pair alternative is impossible by its
      row labels, or print the exact repeated carrier if a relabeling
      permits it.
\item Remove every closed physical outer block once in the V98 partition.
\item Hand the sole survivor to \(\mathbf L_{\rm off}\) with its exact
      \(D_c^{\rm off}\) certificate; do not merge the two noncommuting
      retained-row orientations.
\end{enumerate}
\end{openproblem}

\section{Firewall after V113}
\label{sec:V113-firewall}

\begin{firewall}[Current claim boundary]
\label{fire:V113-current-boundary}
V113 corrects the V112 physical-removal overclaim, proves a direct
absolute-carrier pair tax on the proper V98 outer blocks, closes the
row-comparable physical same-pair region, and reduces the noncomparable
region to heavy puncture-pair or genuine off-cut carriers.  It does not yet
close those two exits, combine the two noncommuting orientations, prove
JREC, construct the continuum theory, or prove a positive mass gap.  The
full Yang--Mills existence and mass gap remain open.
\end{firewall}

\section{Append-only record for V113}
\label{sec:V113-record}

V113 was appended to the complete V112 manuscript.  The exact V112 parent
had \(76260\) newline-delimited source records, \(2656938\) bytes, and
SHA--256 digest
\[
 \texttt{f59309e07584a01f3b3e2328a92010d3876952113b5a4dcaa76dd1c56765a912}.
\]
No V112 source was altered; its physical-removal statement is corrected
by \Cref{prop:V113-V112-correction}.  The sole final
\verb|\end{document}| is restored below.

% =============================================================================
% END APPENDED VERSION V113 MATERIAL
% =============================================================================
% =============================================================================
% BEGIN APPENDED VERSION V114 MATERIAL
% =============================================================================

\chapter{V114: Heavy Puncture-Pair Disposition and the
No-Return-to-Same-Pair Theorem}
\label{ch:V114}

\section{Purpose}
\label{sec:V114-purpose}

V113 reduced every noncomparable physical same-pair remainder to either a
heavy puncture pair \(cj\) or the genuine off-cut carrier.  V114 disposes
of the first exit at the exact local-partition level.  The standing common
pair is \(ab\), whereas the new bank contains \(c\).  Hence the two pair
labels cannot coincide.  If they meet at a coordinate, the partition
axiom forces their blocks to be disjoint; with four rows this can only be
the complementary pattern \(ab\mid co\).  Otherwise the coordinate banks
can be thinned to disjoint heavy subbanks.

Thus the V105 same-pair alternative cannot recur.  The new branch is
either the V104 disjoint two-pair carrier or the V108 complementary
same-coordinate carrier.  This is a complete combinatorial disposition.
Its near-unit decay still uses the effective-scale and fixed-orientation
hypotheses already printed in those theorems; V114 does not manufacture
them from pair heaviness alone.

\section{Fixed labels and local exclusion}
\label{sec:V114-local-labels}

Let \(L_{ab}\) be the V113 common same-pair coordinate bank.  Let
\(M_{cj}\) be the complete same-label regrouping of the heavy puncture
pair supplied by alternative \textup{(i)} of
\Cref{cor:V113-noncomparable-exits}, where
\[
 j\in\{o,a,b\}.
\]
At every coordinate of either bank, the displayed pair is a block of the
same fixed local replica partition monomial.

\begin{lemma}[A heavy third-row pair cannot be the old same pair]
\label{lem:V114-no-same-label}
The unordered labels satisfy
\begin{equation}
 \{c,j\}\ne\{a,b\}.
 \label{eq:V114-distinct-labels}
\end{equation}
Moreover:
\begin{enumerate}[label=\textup{(\roman*)}]
\item if \(j=a\) or \(j=b\), then
      \(L_{ab}\cap M_{cj}=\varnothing\);
\item if \(j=o\), then every coordinate in
      \(L_{ab}\cap M_{co}\) carries the complementary local partition
      \begin{equation}
       \{a,b\}\mathbin{\dot\cup}\{c,o\}.
       \label{eq:V114-complementary-local}
      \end{equation}
\end{enumerate}
\end{lemma}

\begin{proof}
The four input rows \(o,a,b,c\) are distinct, so a pair containing \(c\)
cannot equal \(\{a,b\}\).  If \(j=a\), the two distinct pair blocks
\(\{a,b\}\) and \(\{c,a\}\) share exactly one row.  They cannot both be
blocks of one set partition.  Hence their coordinate supports are
disjoint.  The case \(j=b\) is identical.  If \(j=o\), the two blocks have
four distinct entries; whenever both occur they are the complementary
blocks in \eqref{eq:V114-complementary-local}.  This is also the local
overlap alphabet of \Cref{lem:V105-local-overlap-alphabet}.
\end{proof}

\begin{corollary}[No second one-pair spectral branch]
\label{cor:V114-no-second-same-pair}
The heavy \(cj\) exit never produces a second copy of the one-pair
same-label carrier of \Cref{prop:V105-same-pair-one-operator}.  It always
supplies a distinct pair operator on any nonempty retained subbank.
\end{corollary}

\begin{proof}
Use \eqref{eq:V114-distinct-labels}.  Same-label regrouping joins
coordinates with the fixed label \(cj\); it does not change that label to
\(ab\).
\end{proof}

\section{Quantitative disjoint-or-complementary extraction}
\label{sec:V114-extraction}

Let \(m_a\) denote the endpoint-\(a\) mass on \(L_{ab}\) and \(m_c\) the
endpoint-\(c\) mass on \(M_{cj}\).  The V109 and V113 certificates give
\begin{equation}
 m_a\ge\varepsilon_0\Xi_a,
 \qquad
 m_c\ge\delta_0\Xi_c
 \label{eq:V114-two-certificates}
\end{equation}
for fixed positive programme constants.

\begin{theorem}[Heavy puncture-pair dichotomy without a same-pair exit]
\label{thm:V114-puncture-pair-dichotomy}
There is a fixed \(c_0>0\) such that one of the following holds.
\begin{enumerate}[label=\textup{(\mathrm D)}]
\item There are disjoint coordinate subbanks
      \(G_{ab}\subseteq L_{ab}\) and
      \(G_{cj}\subseteq M_{cj}\) satisfying
      \begin{equation}
       m_a(G_{ab})\ge c_0\Xi_a,
       \qquad
       m_c(G_{cj})\ge c_0\Xi_c.
       \label{eq:V114-disjoint-heavy}
      \end{equation}
\end{enumerate}
or
\begin{enumerate}[label=\textup{(\mathrm C)}]
\item Necessarily \(j=o\), and a common coordinate bank
      \(O_{\rm comp}\subseteq L_{ab}\cap M_{co}\) has the complementary
      blocks \eqref{eq:V114-complementary-local} and satisfies
      \begin{equation}
       m_a(O_{\rm comp})\ge c_0\Xi_a,
       \qquad
       m_c(O_{\rm comp})\ge c_0\Xi_c.
       \label{eq:V114-complementary-heavy}
      \end{equation}
\end{enumerate}
No same-pair alternative occurs.
\end{theorem}

\begin{proof}
Apply \Cref{thm:V105-quarter-dichotomy} to the endpoint certificates
\eqref{eq:V114-two-certificates}, choosing endpoints \(a\) and \(c\).
Its disjoint conclusion gives \textup{(D)} with
\[
 c_0=\frac14\min\{\varepsilon_0,\delta_0\}.
\]
On its overlap conclusion, the same-pair type would require both local
blocks to be \(\{a,c\}\).  One of the fixed labels is instead
\(\{a,b\}\), so that type is impossible by
\Cref{lem:V114-no-same-label}.  Only complementary overlap remains, and
the four-row label forces \(j=o\).  This proves \textup{(C)} with the same
constant.
\end{proof}

\begin{corollary}[Exact carrier disposition]
\label{cor:V114-carrier-disposition}
After all identity spectators are quotiented:
\begin{enumerate}[label=\textup{(\roman*)}]
\item on \textup{(D)}, the positive pair/core carrier contains the exact
      tensor factor
      \[
       Q_{G_{ab}}\otimes Q_{G_{cj}}\otimes M
      \]
      and \Cref{prop:V105-disjoint-factorization} applies;
\item on \textup{(C)}, the local V56 factor is
      \[
       \mathbf d_{ab,f}\otimes\mathbf d_{co,f}
      \]
      on every retained coordinate, and
      \Cref{prop:V105-complementary-joint-operator,
      thm:V108-complementary-inherits-V104} apply.
\end{enumerate}
\end{corollary}

\begin{proof}
Part \textup{(i)} is exactly the disjoint-coordinate tensor
factorization of V105.  Part \textup{(ii)} follows from the complementary
local partition \eqref{eq:V114-complementary-local}, the V56 double-defect
identity, and the V108 pair-first physical subeffect theorem.
\end{proof}

\begin{firewall}[No multiplication on an overlapping row]
\label{fire:V114-no-overlap-product}
In case \textup{(C)} the two local pair blocks are disjoint as replica
blocks but occur at the same gauge coordinate.  Their trace taxes compose
only through the V56/V58 row-factorized tensor before recoupling, and their
positive capacity is handled by V108 as one upstream effect.  V114 does
not multiply two independently optimized post-recoupling capacities.
\end{firewall}

\section{Coefficient and near-unit consequences}
\label{sec:V114-consequences}

\begin{proposition}[The pair taxes are already legitimate on both
alternatives]
\label{prop:V114-trace-tax-disposition}
For arbitrary row-factorized trace-class input coefficients:
\begin{enumerate}[label=\textup{(\roman*)}]
\item on the disjoint alternative, the two selected-heavy V53 taxes act on
      distinct coordinate tensor factors and compose before final
      recoupling;
\item on the complementary alternative, the two taxes compose by the exact
      V58 row-factorized product theorem.
\end{enumerate}
In neither case is a same-pair tax reused.
\end{proposition}

\begin{proof}
The heavy subbanks in
\eqref{eq:V114-disjoint-heavy}--%
\eqref{eq:V114-complementary-heavy} satisfy the endpoint-fraction
hypotheses of \Cref{cor:V53-one-edge-tax}.  Disjoint coordinate factors
give the first composition.  In the second case use
\Cref{lem:V58-tax-composition,thm:V58-two-pair-tax} on the V56 local
double-defect factor.  The labels are distinct by
\eqref{eq:V114-distinct-labels}.
\end{proof}

\begin{theorem}[Conditional near-unit closure of the heavy-pair exit]
\label{thm:V114-conditional-near-unit}
Let \(R_{ab}\) and \(R_{cj}\) be the two effective pair scales on the
heavy subbanks selected by
\Cref{thm:V114-puncture-pair-dichotomy}.  Assume the coefficient caps and
\(s_N^2\) are bounded, the fixed retained-row orientation loss is bounded,
and
\begin{equation}
 R_{ab},R_{cj}\longrightarrow\infty,
 \qquad
 \frac{R_{ab}R_{cj}}{N^2}\longrightarrow\infty.
 \label{eq:V114-two-pair-scale}
\end{equation}
Then the near-unit product capacity of the heavy puncture-pair exit tends
to zero.  In the symmetric case
\(R_{ab}\asymp R_{cj}\asymp R\), the sufficient scale is
\begin{equation}
 R\gg N.
 \label{eq:V114-linear-threshold}
\end{equation}
\end{theorem}

\begin{proof}
On alternative \textup{(D)}, apply the V104 common two-pair cold/hot
surface to the carrier in
\Cref{cor:V114-carrier-disposition}.  On alternative \textup{(C)}, apply
\Cref{thm:V108-complementary-inherits-V104,
cor:V108-complementary-linear-threshold}.  Their hypotheses and
conclusions are exactly \eqref{eq:V114-two-pair-scale}.  The symmetric
reduction is immediate.
\end{proof}

\begin{corollary}[Double-singlet subcase]
\label{cor:V114-double-singlet}
If the two extracted pair levels are dual singlets, their complete
dimension-law sector weight is retained and the V107
\(O(N^{-2})\) near-unit estimate applies, independently of the number of
selected diagonal-row channels.
\end{corollary}

\begin{proof}
Use
\Cref{thm:V107-sector-uniform-law,
thm:V107-double-singlet-core-gate}; the physical sector weight is not
conditionally normalized away.
\end{proof}

\begin{firewall}[Pair heaviness does not force the effective scale]
\label{fire:V114-scale-not-forced}
The mass bounds \eqref{eq:V114-two-certificates} select two genuine pair
operators but do not imply \eqref{eq:V114-two-pair-scale}.  In particular,
the V103 information-level cold rank can still fill the near-unit top
fraction when an effective scale is too small.  Failure of
\eqref{eq:V114-two-pair-scale} is now a two-pair spectral-scale branch,
not a recurrence of the one-pair overlap.
\end{firewall}

\begin{theorem}[Disposition of the noncomparable same-pair branch]
\label{thm:V114-noncomparable-disposition}
After the V98 paid-branch removals and the V113 row-comparable physical
closure, every remaining noncomparable \(ab\) same-pair block belongs to
exactly one of the following acquisition classes:
\begin{enumerate}[label=\textup{(\roman*)}]
\item a disjoint two-pair carrier satisfying the V104 gate;
\item a complementary same-coordinate two-pair carrier satisfying the
      V108 gate; or
\item the genuine off-cut carrier certified by
      \(D_c^{\rm off}\ge c\Xi_c\).
\end{enumerate}
The first two classes decay under
\eqref{eq:V114-two-pair-scale}; the third remains separate.
\end{theorem}

\begin{proof}
\Cref{cor:V113-noncomparable-exits} gives a heavy puncture pair or heavy
off-cut mass.  The former has the exhaustive two-case disposition of
\Cref{thm:V114-puncture-pair-dichotomy}; the latter is class
\textup{(iii)}.  Orthogonal V98 outer labels make the classes disjoint
after a fixed priority order is chosen.
\end{proof}

\section{V114 ledger and mandatory V115 audit}
\label{sec:V114-ledger}

\begin{theorem}[Integrated V114 statement]
\label{thm:V114-integrated}
Preserving V1--V113, V114 proves:
\begin{enumerate}[label=\textup{(V114.\arabic*)},leftmargin=6.7em]
\item the local label exclusion
      \eqref{eq:V114-distinct-labels};
\item impossibility of same-coordinate overlap for \(ca\) or \(cb\)
      against the \(ab\) pair;
\item the complementary \(ab\mid co\) identity on every possible overlap;
\item the quantitative disjoint-or-complementary extraction
      \eqref{eq:V114-disjoint-heavy}--%
      \eqref{eq:V114-complementary-heavy};
\item exact legitimacy of both coefficient pair taxes;
\item conditional near-unit decay at the two-pair scale
      \eqref{eq:V114-two-pair-scale}; and
\item reduction of every noncomparable same-pair remainder to the two
      established two-pair mechanisms or the genuine off-cut carrier.
\end{enumerate}
\end{theorem}

\begin{proof}
These are
\Cref{lem:V114-no-same-label,
cor:V114-no-second-same-pair,
thm:V114-puncture-pair-dichotomy,
cor:V114-carrier-disposition,
prop:V114-trace-tax-disposition,
thm:V114-conditional-near-unit,
thm:V114-noncomparable-disposition}.
\end{proof}

\begin{openproblem}[V115 mandate: audit and physical off-cut carrier]
\label{op:V115-mandate}
Before new mathematics, audit the current SM and NS notes.  Then work only
on class \textup{(iii)} of
\Cref{thm:V114-noncomparable-disposition}.
\begin{enumerate}
\item Write \(D_c^{\rm off}\) as its exact pair/private incidence blocks in
      the punctured core, preserving the V98 physical outer label.
\item Identify the pair or higher local block carrying a fixed fraction of
      \(D_c^{\rm off}\); do not treat it as an identity spectator.
\item Test whether positive completion
      \Cref{thm:V113-positive-completion} supplies an absolute-carrier tax
      for that block.
\item If its label crosses the retained-row cut, route it by V53/V58; if it
      lies wholly off the cut, retain it with the signed ternary core and
      compute the corresponding V98 injective image.
\item Keep any failure of \eqref{eq:V114-two-pair-scale} as a separate
      conditional spectral branch.
\item Do not merge the two noncommuting retained-row orientations.
\end{enumerate}
\end{openproblem}

\section{Firewall after V114}
\label{sec:V114-firewall}

\begin{firewall}[Current claim boundary]
\label{fire:V114-current-boundary}
V114 proves that the heavy puncture-pair exit can only be disjoint or
complementary and never returns to the one-pair obstruction.  It closes
the corresponding coefficient-tax geometry and gives conditional
near-unit decay at the established two-pair scale.  It does not force that
scale, close the genuine off-cut carrier, combine the two noncommuting
orientations, prove JREC, construct the continuum theory, or prove a
positive mass gap.  The full Yang--Mills existence and mass gap remain
open.
\end{firewall}

\section{Append-only record for V114}
\label{sec:V114-record}

V114 was appended to the complete V113 manuscript.  The exact V113 parent
had \(76845\) newline-delimited source records, \(2678424\) bytes, and
SHA--256 digest
\[
 \texttt{0b4785821a4db5218c55c7571973d403c35b4032fecccf3d97a3f9093803e1d2}.
\]
No V113 mathematical content was changed.  The sole final
\verb|\end{document}| is restored below.

% =============================================================================
% END APPENDED VERSION V114 MATERIAL
% =============================================================================
% =============================================================================
% BEGIN APPENDED VERSION V115 MATERIAL
% =============================================================================

\chapter{V115: Companion Audit and the Exact Off-Cut
Mean--Covariance Peeling Identity}
\label{ch:V115}

\section{Purpose}
\label{sec:V115-purpose}

V113--V114 reduced the noncomparable same-pair branch to established
two-pair carriers, subject to their printed effective-scale hypotheses, or
to the genuine off-cut mass \(D_c^{\rm off}\).  V115 performs the required
five-version companion audit and then goes upstream of the off-cut
absolute carrier.

On a core coordinate counted by \(D_c^{\rm off}\), row \(c\) is active,
row \(a\) is inactive, and the coordinate is not private to \(c\).
Therefore row \(b\) is active: the entire off-cut bank is a \(bc\)-only
bank relative to \(a\).  It is independent of the remaining coordinates.
This gives an exact two-term cumulant identity.  The product of the two
off-cut singleton means multiplies the reduced ternary cumulant, while
the off-cut \(bc\) covariance multiplies one composite binary covariance
joining \(a\) to the reduced \(bc\) product.  No third term remains.

\section{Mandatory V115 companion audit}
\label{sec:V115-audit}

The read-only audit on 5 September 2026 found
\begin{center}
\begin{tabular}{|p{0.50\textwidth}|r|p{0.30\textwidth}|}
\hline
Companion file & Bytes & SHA--256\\
\hline
\texttt{Renewal\_NS\_Stability\_Research\_Notes\_V31.tex}
& \(887537\) &
\texttt{f1121b62238ad5491bb7cf03635ea9934942edf573ed8faaa9ee79e1d38a6696}\\
\hline
\texttt{Renewal\_SM\_Einstein\_Continuum\_Research\_Notes\_V127.tex}
& \(3115797\) &
\texttt{a72e3644fb6d7b545d317c9b92e29b13007b7ff53c8aa3545d325cdc9f1d4969}\\
\hline
\end{tabular}
\end{center}
The files contain respectively \(23053\) and \(85431\)
newline-delimited source records under the audit counter.  NS is
byte-identical to the V110 audit and supplies no new estimate.

SM advanced from V121 to V127.  V122 replaces a moving
parameter-dependent boundary domain by one fixed enlarged collar before
Schur elimination.  V123 performs longitudinal pairing before splitting
the relative determinant into interface channels.  V125 tests the exact
support and inverse reference weight before expanding near a harmonic
threshold.  V126--V127 distinguish a frozen local parity statement from
the complete pseudodifferential symbol and prove vanishing of one
phase-space divergence only after the covariantly regulated integral.

\begin{assessment}[V115 adjustment after the audit]
\label{ass:V115-audit-adjustment}
The common lesson is to retain the enlarged object through the operation
which creates its cancellation.  V115 therefore factors the full off-cut
probability space into the \(bc\)-only bank and its complement, expands the
complete ternary cumulant exactly, and only then identifies its mean and
pair-covariance branches.  It does not estimate the original five moments
separately or identify a local factor with its final integrated
coefficient.  No SM or NS analytic bound is imported.
\end{assessment}

\section{The exact off-cut coordinate bank}
\label{sec:V115-offcut-bank}

Let \(O_{bc}\) be the set of punctured-core coordinates counted by
\(D_c^{\rm off}\) in \eqref{eq:V59-r-mass-split}.  Write
\[
 \Xi_{c,0}:=\sum_{f\in O_{bc}}A_{c,f}=D_c^{\rm off}.
\]

\begin{lemma}[Off-cut means \(bc\)-only]
\label{lem:V115-offcut-bc-only}
For every \(f\in O_{bc}\),
\begin{equation}
 A_{c,f}>0,\qquad A_{b,f}>0,\qquad A_{a,f}=0.
 \label{eq:V115-bc-only}
\end{equation}
Consequently
\begin{equation}
 \widehat H_{bc}
 \ge
 a_*D_c^{\rm off}.
 \label{eq:V115-offcut-pair-heavy}
\end{equation}
\end{lemma}

\begin{proof}
By definition, a coordinate in \(D_c^{\rm off}\) is active at \(c\), is
not in \(E_{c\leftarrow a}\), and is not private to \(c\).  The first
condition gives \(A_{c,f}>0\); the second gives \(A_{a,f}=0\).  The
punctured ternary core has only rows \(a,b,c\).  Since the coordinate is
not private to \(c\), row \(b\) must be active.  The active mass floor
\eqref{eq:V59-active-mass-floor} gives
\[
 A_{b,f}A_{c,f}\ge a_*A_{c,f}.
\]
Sum over \(O_{bc}\).
\end{proof}

\begin{corollary}[Heavy off-cut is a selected \(bc\) pair bank]
\label{cor:V115-heavy-offcut-pair}
On alternative \eqref{eq:V113-heavy-off-cut},
\begin{equation}
 \widehat H_{bc}
 \ge
 a_*(1-7\delta)\Xi_c.
 \label{eq:V115-heavy-bc-cut}
\end{equation}
Thus the off-cut bank is pair-heavy at endpoint \(c\), although it is not
one of the two center-\(a\) cuts in the original V57 allocation.
\end{corollary}

\begin{proof}
Combine \eqref{eq:V115-offcut-pair-heavy} with
\eqref{eq:V113-heavy-off-cut}.
\end{proof}

\section{Independent-bank notation}
\label{sec:V115-bank-notation}

The product probability space factors as
\[
 \Omega=\Omega_0\times\Omega_1,
\]
where \(\Omega_0\) contains exactly the coordinates in \(O_{bc}\).
Because row \(a\) is inactive there, write the three operator-valued
products as
\begin{equation}
 F_a=A,\qquad
 F_b=B_0B_1,\qquad
 F_c=C_0C_1,
 \label{eq:V115-factor-notation}
\end{equation}
where \(B_0,C_0\) depend only on \(\Omega_0\), while
\(A,B_1,C_1\) depend only on \(\Omega_1\).  Factors on distinct replica
carriers commute.  Put
\begin{align}
 b_0&:=\mathbb E_0B_0,
 &
 c_0&:=\mathbb E_0C_0,
 \label{eq:V115-offcut-means}\\
 \Delta_{bc}^{0}
 &:=
 \operatorname{Cov}_0(B_0,C_0)
 =
 \mathbb E_0(B_0C_0)-b_0c_0.
 \label{eq:V115-offcut-covariance}
\end{align}
Expectations and covariances with subscript \(1\) refer to
\(\Omega_1\).

\begin{theorem}[Exact off-cut mean--covariance peeling]
\label{thm:V115-offcut-peeling}
The complete ternary cumulant has the exact factorization
\begin{equation}
 \boxed{
 \kappa_3(A,B_0B_1,C_0C_1)
 =
 b_0c_0\,\kappa_{3,1}(A,B_1,C_1)
 +
 \Delta_{bc}^{0}\,
 \operatorname{Cov}_1(A,B_1C_1).}
 \label{eq:V115-offcut-peeling}
\end{equation}
The identity holds for operator-valued factors on distinct input carriers,
with arbitrary spectator amplification, before trace norm and final
recoupling.
\end{theorem}

\begin{proof}
Set
\[
 M_0:=\mathbb E_0(B_0C_0)=b_0c_0+\Delta_{bc}^{0}.
\]
Independence of \(\Omega_0,\Omega_1\) and the five-term definition of the
ternary cumulant give
\begin{align*}
 \kappa_3(A,B_0B_1,C_0C_1)
 ={}&
 M_0\bigl[
  \mathbb E_1(AB_1C_1)
  -\mathbb E_1A\,\mathbb E_1(B_1C_1)
 \bigr]\\
 &-b_0c_0\bigl[
  \mathbb E_1B_1\,\mathbb E_1(AC_1)
  +\mathbb E_1C_1\,\mathbb E_1(AB_1)\\
 &\hspace{38mm}
  -2\mathbb E_1A\,\mathbb E_1B_1\,\mathbb E_1C_1
 \bigr].
\end{align*}
The first bracket is
\(\operatorname{Cov}_1(A,B_1C_1)\).  The second bracket equals
\[
 \operatorname{Cov}_1(A,B_1C_1)
 -\kappa_{3,1}(A,B_1,C_1)
\]
by the V57 three-covariance identity.  Substitute
\(M_0=b_0c_0+\Delta_{bc}^{0}\) and cancel the two
\(b_0c_0\operatorname{Cov}_1(A,B_1C_1)\) terms.  This yields
\eqref{eq:V115-offcut-peeling}.  All rearrangements preserve carrier
order because the displayed replica factors commute.
\end{proof}

\begin{corollary}[The off-cut bank cannot create connectivity by itself]
\label{cor:V115-offcut-no-connectivity}
If
\[
 \kappa_{3,1}(A,B_1,C_1)=0,
 \qquad
 \operatorname{Cov}_1(A,B_1C_1)=0,
\]
then the complete ternary cumulant vanishes, regardless of the size or
representation content of the \(bc\)-only bank.
\end{corollary}

\begin{proof}
Both right-hand terms in \eqref{eq:V115-offcut-peeling} vanish.
\end{proof}

\begin{corollary}[Two exact limiting checks]
\label{cor:V115-two-checks}
The peeling identity has the following sharp specializations.
\begin{enumerate}[label=\textup{(\roman*)}]
\item If \(\Delta_{bc}^{0}=0\), then
      \[
       \kappa_3(A,B_0B_1,C_0C_1)
       =
       b_0c_0\,\kappa_{3,1}(A,B_1,C_1).
      \]
\item If \(b_0c_0=0\), then
      \[
       \kappa_3(A,B_0B_1,C_0C_1)
       =
       \Delta_{bc}^{0}\operatorname{Cov}_1(A,B_1C_1).
      \]
\end{enumerate}
\end{corollary}

\begin{proof}
Set the indicated factor to zero in
\eqref{eq:V115-offcut-peeling}.
\end{proof}

\section{Separate influence numerators after peeling}
\label{sec:V115-influence}

\begin{proposition}[Reduced connecting bounds]
\label{prop:V115-reduced-bounds}
For normalized Wilson factors,
\begin{align}
 \|\kappa_{3,1}(A,B_1,C_1)\|_{\rm op}
 &\le
 C\min\!\left\{
 s_\alpha(H_{ab}^{1}+H_{ac}^{1}),1
 \right\},
 \label{eq:V115-reduced-ternary-cap}\\
 \|\operatorname{Cov}_1(A,B_1C_1)\|_{\rm op}
 &\le
 Cs_\alpha(H_{ab}^{1}+H_{ac}^{1}),
 \label{eq:V115-reduced-composite-cap}
\end{align}
where the superscript \(1\) restricts the influence sums to
\(\Omega_1\).  In particular, every term in
\eqref{eq:V115-offcut-peeling} retains an explicit connection from row
\(a\) to row \(b\) or \(c\) outside the off-cut bank.
\end{proposition}

\begin{proof}
The first statement is \Cref{thm:V57-cut-cap} on the reduced product
space.  For the second, apply
\Cref{lem:V50-operator-covariance} to \(A\) and \(B_1C_1\).  The
coordinate influence of the latter is bounded by the sum of the \(B_1\)
and \(C_1\) influences, exactly as in the proof of V57.  Restricting the
sum to \(\Omega_1\) gives the two displayed numerators.
\end{proof}

\begin{corollary}[Norm-level off-cut bound]
\label{cor:V115-offcut-norm}
Since normalized singleton means are contractions,
\begin{align}
 &\|\kappa_3(A,B_0B_1,C_0C_1)\|_{\rm op}
 \notag\\
 &\quad\le
 C\min\!\left\{
 s_\alpha(H_{ab}^{1}+H_{ac}^{1}),1
 \right\}
 +
 Cs_\alpha\|\Delta_{bc}^{0}\|_{\rm op}
 (H_{ab}^{1}+H_{ac}^{1}).
 \label{eq:V115-offcut-norm}
\end{align}
\end{corollary}

\begin{proof}
Take the operator norm in
\eqref{eq:V115-offcut-peeling}, use
\(\|b_0\|_{\rm op},\|c_0\|_{\rm op}\le1\), and apply
\Cref{prop:V115-reduced-bounds}.
\end{proof}

\begin{firewall}[The norm bound is not the off-cut coefficient tax]
\label{fire:V115-norm-not-tax}
Equation \eqref{eq:V115-offcut-norm} records the two connecting
influence numerators, but it does not pay the original-row denominator
associated with the heavy \(bc\)-only bank.  The mean factor \(b_0c_0\)
and the covariance factor \(\Delta_{bc}^{0}\) must remain in their
physical pair blocks through the V98 absolute carrier.  Replacing both by
their operator caps would return to the unresolved scalar stock.
\end{firewall}

\section{Physical two-branch off-cut reduction}
\label{sec:V115-physical-reduction}

After complete same-label regrouping of \(O_{bc}\), let
\(\mathbf M_{bc}^{0}\) denote the physical multiplier \(b_0c_0\), and let
\(\mathbf D_{bc}^{0}\) denote the signed pair-covariance multiplier
\(\Delta_{bc}^{0}\).  Let
\(\mathbf R_{\rm ter}^{1}\) and
\(\mathbf R_{\rm bin}^{1}\) be the two reduced operators in
\eqref{eq:V115-offcut-peeling}, with every outer collision and resolvent
weight retained.  On each compatible V98 outer block,
\begin{equation}
 \mathbf L_{\rm off}
 =
 \mathbf M_{bc}^{0}\otimes\mathbf R_{\rm ter}^{1}
 +
 \mathbf D_{bc}^{0}\otimes\mathbf R_{\rm bin}^{1}.
 \label{eq:V115-physical-offcut-split}
\end{equation}
Denote the two signed summands, assembled over the same orthogonal outer
blocks, by \(\mathbf L_{\rm mean}\) and \(\mathbf L_{\rm cov}\), and put
\[
 \mathbf Q_{\rm mean}:=|\mathbf L_{\rm mean}|,
 \qquad
 \mathbf Q_{\rm cov}:=|\mathbf L_{\rm cov}|.
\]

\begin{theorem}[Exact sufficient off-cut gates]
\label{thm:V115-offcut-gates}
The physical off-cut branch has the target marked-tree bound if both
\begin{align}
 \left\|
 \mathbf Q_{\rm mean}
 \right\|_{\rm rows,prod}
 &\le \mathcal P_{\rm off}^{\rm mean},
 \label{eq:V115-offcut-mean-gate}\\
 \left\|
 \mathbf Q_{\rm cov}
 \right\|_{\rm rows,prod}
 &\le \mathcal P_{\rm off}^{\rm cov}
 \label{eq:V115-offcut-cov-gate}
\end{align}
hold after summing the orthogonal physical outer blocks, where
\(\mathcal P_{\rm off}^{\rm mean}+
\mathcal P_{\rm off}^{\rm cov}\) is bounded by the corresponding sum of
quartic marked-tree weights.  No additional mixed term is required in the
trace-class coefficient bound.
\end{theorem}

\begin{proof}
For every fixed input coefficient tensor, linearity of physical
compression and \eqref{eq:V115-physical-offcut-split} give a sum of a mean
output coefficient and a covariance output coefficient in each outer
block.  The trace-norm triangle inequality bounds its norm by the sum of
their norms.  Apply the physical matrix tax
\eqref{eq:V98-restricted-matrix-tax} separately to the two signed outer
families, whose absolute carriers are
\(\mathbf Q_{\rm mean}\) and \(\mathbf Q_{\rm cov}\), and insert
\eqref{eq:V115-offcut-mean-gate}--%
\eqref{eq:V115-offcut-cov-gate}.  Orthogonal outer blocks add without a
channel count.
\end{proof}

\begin{remark}[Why this split is different from the V97 three squares]
\label{rem:V115-not-three-square}
The two terms in \eqref{eq:V115-physical-offcut-split} are obtained from
an exact factorization of independent coordinate banks before any norm.
They are not the three positive covariance squares rejected in V97.
Possible cancellation between the two terms may still improve the bound;
the two-gate estimate is a sufficient coefficient-level triangle
inequality, not an operator-order assertion about
\(|\mathbf L_{\rm off}|\).
\end{remark}

\begin{assessment}[Disposition of the off-cut bottleneck]
\label{ass:V115-offcut-disposition}
The off-cut carrier is no longer an unnamed positive capacity.  It is the
sum of:
\begin{enumerate}[label=\textup{(\roman*)}]
\item a joint \(bc\) mean times a reduced ternary cumulant; and
\item a selected-heavy \(bc\) pair covariance times a reduced composite
      binary covariance.
\end{enumerate}
Both reduced operators carry only the connecting \(ab/ac\) influence
numerators.  The next work is to obtain the \(c\)-endpoint debit from the
mean in the first branch and from the V53 positive pair multiplier in the
second, while retaining the V98 outer blocks.
\end{assessment}

\section{V115 ledger and V116 acquisition}
\label{sec:V115-ledger}

\begin{theorem}[Integrated V115 statement]
\label{thm:V115-integrated}
Preserving V1--V114, V115 proves:
\begin{enumerate}[label=\textup{(V115.\arabic*)},leftmargin=6.7em]
\item the mandatory SM V127 and NS V31 companion audit;
\item identification of \(D_c^{\rm off}\) as a \(bc\)-only coordinate
      bank and the heavy bound \eqref{eq:V115-heavy-bc-cut};
\item the exact independent-bank cumulant identity
      \eqref{eq:V115-offcut-peeling};
\item proof that the off-cut bank cannot create ternary connectivity by
      itself;
\item the reduced connecting bounds
      \eqref{eq:V115-reduced-ternary-cap}--%
      \eqref{eq:V115-reduced-composite-cap}; and
\item reduction of the physical off-cut carrier to the two exact gates
      \eqref{eq:V115-offcut-mean-gate}--%
      \eqref{eq:V115-offcut-cov-gate}.
\end{enumerate}
\end{theorem}

\begin{proof}
These are
\Cref{sec:V115-audit,
lem:V115-offcut-bc-only,
cor:V115-heavy-offcut-pair,
thm:V115-offcut-peeling,
cor:V115-offcut-no-connectivity,
prop:V115-reduced-bounds,
thm:V115-offcut-gates}.
\end{proof}

\begin{openproblem}[V116 mandate: pay both peeled off-cut branches]
\label{op:V116-mandate}
\begin{enumerate}
\item For the mean branch, keep the full \(bc\) joint mean until its
      product/same-label output has been resolved; derive its endpoint
      first-moment debit without replacing it by one.
\item For the covariance branch, construct the positive V52 multiplier
      majorizing \(|\mathbf D_{bc}^{0}|\) and apply
      \Cref{thm:V113-positive-completion} on the physical V98 outer family.
\item Attach the reduced composite covariance before trace norm and retain
      its separate \(ab\) and \(ac\) hybrid replacements from V110.
\item Determine the exact marked-tree allocation when the new \(bc\) pair
      denominator meets the original \(oa\) and \(bc\) collision
      denominators; do not count the \(bc\) edge twice.
\item Test a pure \(bc\) singlet bank and the alternating-height bank.
\item Keep the conditional two-pair scale deficit and the second retained-
      row orientation separate.
\end{enumerate}
\end{openproblem}

\section{Firewall after V115}
\label{sec:V115-firewall}

\begin{firewall}[Current claim boundary]
\label{fire:V115-current-boundary}
V115 performs the required audit, proves the exact off-cut
mean--covariance peeling identity, and reduces the physical off-cut carrier
to two explicit gates.  It does not yet prove either gate at the full
marked-tree scale, force the V114 two-pair effective scale, combine the two
noncommuting orientations, prove JREC, construct the continuum theory, or
prove a positive mass gap.  The full Yang--Mills existence and mass gap
remain open.
\end{firewall}

\section{Append-only record for V115}
\label{sec:V115-record}

V115 was appended to the complete V114 manuscript.  The exact V114 parent
had \(77217\) newline-delimited source records, \(2693135\) bytes, and
SHA--256 digest
\[
 \texttt{0563d9aaeeb46eb717e6fe12e6eff82ee12b4784c12ec9a0ddb3a42eb947bdee}.
\]
No V114 mathematical content was changed.  The sole final
\verb|\end{document}| is restored below.

% =============================================================================
% END APPENDED VERSION V115 MATERIAL
% =============================================================================
% =============================================================================
% BEGIN APPENDED VERSION V116 MATERIAL
% =============================================================================

\chapter{V116: Payment of the Peeled Mean Branch and the
Maximal-Off-Cut Covariance Branch}
\label{ch:V116}

\section{Purpose}
\label{sec:V116-purpose}

V115 factorized the heavy \(bc\)-only off-cut bank exactly:
\[
 \kappa_3(A,B_0B_1,C_0C_1)
 =
 b_0c_0\,\kappa_{3,1}(A,B_1,C_1)
 +
 \Delta_{bc}^{0}\operatorname{Cov}_1(A,B_1C_1).
\]
V116 pays the first term by retaining the unused scalar debit in \(c_0\).
For the second, it applies the physical positive-completion theorem to the
selected-heavy \(bc\) pair multiplier.  The resulting squared
\(c\)-denominator closes the covariance branch whenever row \(c\) is the
largest internal row.  Thus the only off-cut mass ordering left by this
route is
\[
 \Xi_a>\Xi_c\gg\Xi_b.
\]
No pair capacities sharing an original row are multiplied.

\section{A paired-mean endpoint debit}
\label{sec:V116-mean-debit}

Let \(p_c=D_c^{\rm off}\) be the input-\(c\) mass of the \(bc\)-only bank.
In the mean term of \eqref{eq:V115-offcut-peeling}, \(b_0\) and \(c_0\)
are separate scalar product heat multipliers, even though their
coordinates have the same names in the two replica rows.

\begin{lemma}[The \(c\)-mean debit survives the \(b\)-mean]
\label{lem:V116-paired-mean-debit}
If \(s_\alpha p_c>M\), then
\begin{align}
 |b_0c_0|
 &\le
 \frac{C}{(s_\alpha p_c)^2},
 \label{eq:V116-paired-mean-second}\\
 s_\alpha\Xi_{0}^{\rm out}|b_0c_0|
 &\le
 \frac{C}{s_\alpha p_c},
 \label{eq:V116-paired-mean-output}
\end{align}
where \(\Xi_0^{\rm out}\) is the total resolved output mass of both
off-cut mean factors.  If an exact incident pair defect of weight \(h_c\)
is retained with them, their combined normalized response satisfies
\begin{equation}
 \boxed{
 \mathfrak A_{c}^{\rm off,mean}
 \le C\frac{h_c}{p_c}.}
 \label{eq:V116-offcut-mean-endpoint}
\end{equation}
All other modulus-one, first-moment factors may remain as spectators.
\end{lemma}

\begin{proof}
Apply the product second-moment estimate
\eqref{eq:V55-mean-second-debit} to the exact \(c\)-mean bank and use
\(|b_0|\le1\).  For the output estimate, split
\(\Xi_0^{\rm out}\) between the \(c_0\) and \(b_0\) outputs.  The former
uses \eqref{eq:V55-mean-output-second-debit}.  The latter uses the product
first moment \eqref{eq:V46-product-first-moment} for \(b_0\), while
retaining the second debit of \(c_0\):
\[
 s_\alpha\Xi_b^{\rm out}|b_0c_0|
 \le
 C|c_0|
 \le
 \frac{C}{(s_\alpha p_c)^2}
 \le
 \frac{C_M}{s_\alpha p_c}.
\]
This proves the first two displays.  Now repeat the output allocation in
\Cref{lem:V55-unused-debit}: use the defect cap on its modulus term, its
first moment on its own output, and
\eqref{eq:V116-paired-mean-output} on the mean outputs.  The result is
\(Ch_c/p_c\).
\end{proof}

\begin{corollary}[Heavy off-cut mean gives a \(c\)-endpoint response]
\label{cor:V116-mean-endpoint-response}
On \eqref{eq:V113-heavy-off-cut},
\begin{equation}
 \mathfrak A_c^{\rm off,mean}
 \le
 C_\delta\frac{h_c}{\Xi_c}.
 \label{eq:V116-mean-response-Xi}
\end{equation}
If \(s_\alpha\Xi_c\) is bounded, the same conclusion follows from the raw
thermal endpoint estimate.
\end{corollary}

\begin{proof}
Use \(p_c\ge(1-7\delta)\Xi_c\) in
\eqref{eq:V116-offcut-mean-endpoint}.  The bounded thermal case is
\Cref{lem:V59-old-endpoint-response}.
\end{proof}

\begin{theorem}[The peeled off-cut mean branch is paid]
\label{thm:V116-mean-branch-paid}
In the V113 same-pair configuration, suppose the other active internal
endpoints retain their already established endpoint responses.  Then the
physical mean branch
\[
 \mathbf M_{bc}^{0}\otimes\mathbf R_{\rm ter}^{1}
\]
satisfies the summed quartic marked-path bound.
\end{theorem}

\begin{proof}
The off-cut mean bank is disjoint from the collision coordinate and from
the reduced probability space.  Keep \(b_0c_0\), the incident collision
defect, and the reduced ternary factor assembled until the output mass is
allocated.  \Cref{cor:V116-mean-endpoint-response} supplies the missing
endpoint response at \(c\); the hypothesis supplies the remaining
responses.  The proof of \Cref{cor:V59-endpoint-closure} then applies
verbatim to the reduced ternary core.  Aggregate V98 outer pinching is
contractive, so the conclusion is physical and introduces no channel
count.
\end{proof}

\begin{firewall}[This is a mean payment, not a \(bc\) pair tax]
\label{fire:V116-mean-not-pair}
The proof uses the scalar factor \(c_0\) from the first term of the exact
peeling identity.  It does not use the covariance multiplier
\(\Delta_{bc}^{0}\), and therefore leaves that multiplier available only
for the second peeled branch.  The two branches are added by the
coefficient trace-norm triangle of V115.
\end{firewall}

\section{Positive capacity of the peeled covariance branch}
\label{sec:V116-covariance-capacity}

Resolve the \(bc\)-only bank into its multiplicity-resolved pair outputs.
Let \(\mathbf A_{bc}^{0}\) be the V112 positive multiplier whose
eigenvalues majorize the absolute weighted covariance coefficients of
\(\mathbf D_{bc}^{0}\).

\begin{lemma}[Selected-heavy off-cut pair capacity]
\label{lem:V116-offcut-pair-capacity}
There is a fixed \(\eta_\delta>0\) such that
\[
 p_c\ge\eta_\delta\Xi_c.
\]
Consequently
\begin{equation}
 \boxed{
 \|\mathbf A_{bc}^{0}\|_{\rm prod}
 \le
 C_\delta
 \frac{h_{bc}^{0}}{s_\alpha\Xi_c^2},}
 \label{eq:V116-offcut-pair-capacity}
\end{equation}
where \(h_{bc}^{0}\) is the \(bc\) overlap weight on \(O_{bc}\).
\end{lemma}

\begin{proof}
Take \(\eta_\delta=1-7\delta\) in
\eqref{eq:V113-heavy-off-cut}.  The exact \(bc\) bank is therefore
selected-heavy at endpoint \(c\).  Apply
\Cref{prop:V87-selected-heavy-capacity}; equivalently use
\Cref{lem:V112-V52-dual-positive} with
\Cref{cor:V53-one-edge-tax}.
\end{proof}

\begin{theorem}[Physical covariance-branch capacity]
\label{thm:V116-physical-covariance-capacity}
On every V98 outer subfamily on which the off-cut pair multiplier occurs
once, the covariance branch absolute carrier satisfies
\begin{align}
 \left\|\mathbf Q_{\rm cov}\right\|_{\rm rows,prod}
 \le
 C_\delta
 \frac{h_{bc}^{0}
 (H_{ab}^{1}+H_{ac}^{1})}
 {\Xi_c^2}.
 \label{eq:V116-covariance-absolute}
\end{align}
The statement allows arbitrary correlation between the \(bc\) pair output
and the reduced-core output.
\end{theorem}

\begin{proof}
The reduced binary covariance has the cap
\eqref{eq:V115-reduced-composite-cap}.  Thus every physical outer block
has the common pair majorant
\[
 C s_\alpha(H_{ab}^{1}+H_{ac}^{1})
 a_{\lambda(bc)}P_j.
\]
Apply \Cref{thm:V113-positive-completion} with
\(\mathbf A_L=\mathbf A_{bc}^{0}\), then insert
\eqref{eq:V116-offcut-pair-capacity}.  The \(s_\alpha\) factors cancel.
\end{proof}

\section{Maximal-row conversion and exact residual ordering}
\label{sec:V116-maximal-row}

The capacity estimate
\eqref{eq:V116-covariance-absolute} has two powers of the mass of the
selected-heavy endpoint.  If that endpoint is maximal, those two powers
may be assigned to the two endpoints of either reduced-core edge.  This
is only a scalar comparison after the common physical carrier has been
estimated; it is not a product of operator capacities.

\begin{corollary}[Maximal-\(c\) marked-tree conversion]
\label{cor:V116-maximal-c-conversion}
Assume that all three active masses are positive and
\begin{equation}
 \Xi_c\ge\max\{\Xi_a,\Xi_b\}.
 \label{eq:V116-c-maximal}
\end{equation}
Then
\begin{align}
 \left\|\mathbf Q_{\rm cov}\right\|_{\rm rows,prod}
 &\le
 C_\delta h_{bc}^{0}
 \left(
 \frac{H_{ab}^{1}}{\Xi_a\Xi_b}
 +
 \frac{H_{ac}^{1}}{\Xi_a\Xi_c}
 \right).
 \label{eq:V116-covariance-marked}
\end{align}
\end{corollary}

\begin{proof}
The maximality hypothesis gives
\[
 \Xi_c^2\ge\Xi_a\Xi_b,
 \qquad
 \Xi_c^2\ge\Xi_a\Xi_c.
\]
Apply these two inequalities separately to the \(H_{ab}^{1}\) and
\(H_{ac}^{1}\) summands of
\eqref{eq:V116-covariance-absolute}.
\end{proof}

\begin{remark}[Where the \(bc\) numerator is placed]
\label{rem:V116-no-double-bc}
In an adapted routed graph monomial,
\(h_{bc}^{0}\) in \eqref{eq:V116-covariance-marked} is the one selected
or surplus \(bc\)-pair factor represented by the off-cut covariance
multiplier.  It is inserted once at that edge.  Any pre-existing
collision factor is retained at its own graph edge; the off-cut
multiplier is not charged a second time as a scalar covariance cap.
Thus the statement supplies the two displayed endpoint denominators
only after the exact graph monomial identifies
\(h_{bc}^{0}\) as its corresponding pair factor.  It does not identify
unrelated pair weights.
\end{remark}

\begin{theorem}[The covariance branch is paid in the maximal-\(c\)
region]
\label{thm:V116-covariance-cmax-paid}
On a V98 physical outer subfamily satisfying
\eqref{eq:V116-c-maximal}, suppose its adapted routed monomial contains
the off-cut covariance weight \(h_{bc}^{0}\) once and the reduced
connection is routed through \(ab\) or \(ac\) according to the two
summands in \eqref{eq:V115-reduced-composite-cap}.  Then the covariance
branch obeys the corresponding summed quartic marked-tree bound.
\end{theorem}

\begin{proof}
First apply the positive physical completion in
\Cref{thm:V116-physical-covariance-capacity}; this keeps all correlated
outer labels inside one positive carrier.  Then use
\Cref{cor:V116-maximal-c-conversion}.  The first summand pays the
\(a,b\) endpoint denominator and the second pays the \(a,c\) endpoint
denominator.  The graph hypothesis and
\Cref{rem:V116-no-double-bc} place \(h_{bc}^{0}\) exactly once.
Finally sum the two nonnegative routed contributions.  No
coordinatewise orientation minimum and no multiplication of two
separately optimized capacities occurs.
\end{proof}

\begin{proposition}[Exact mass ordering left by the off-cut route]
\label{prop:V116-exact-residual-order}
In the noncomparable V113 branch one already has
\begin{equation}
 \Xi_c>C_{\rm row}\Xi_b.
 \label{eq:V116-c-over-b}
\end{equation}
After \Cref{thm:V116-covariance-cmax-paid}, the only unclassified
ordering for the peeled covariance term is
\begin{equation}
 \boxed{\Xi_a>\Xi_c>C_{\rm row}\Xi_b.}
 \label{eq:V116-residual-order}
\end{equation}
Here the informal notation \(\Xi_c\gg\Xi_b\) used in
\Cref{sec:V116-purpose} means precisely the fixed comparison
\eqref{eq:V116-c-over-b}.
\end{proposition}

\begin{proof}
Partition the positive-mass region into
\(\Xi_c\ge\Xi_a\) and \(\Xi_a>\Xi_c\).  In the first region,
\eqref{eq:V116-c-over-b} also gives \(\Xi_c>\Xi_b\), so
\eqref{eq:V116-c-maximal} holds and
\Cref{thm:V116-covariance-cmax-paid} applies.  The second region,
together with \eqref{eq:V116-c-over-b}, is exactly
\eqref{eq:V116-residual-order}.  The strict/tied boundary has been
assigned to the paid region by the displayed priority convention.
\end{proof}

\section{Sharp regression checks}
\label{sec:V116-regressions}

\begin{proposition}[Pure \(bc\) dual-singlet check]
\label{prop:V116-dual-singlet-check}
Suppose one resolved off-cut pair block is a normalized dual-singlet
projector \(P_\Omega\) on
\(\mathbb C^d\otimes\mathbb C^d\), and its signed coefficient is
\(d_0P_\Omega\).  The positive multiplier used in V116 is then
\(|d_0|P_\Omega\), and
\begin{equation}
 \bigl\||d_0|P_\Omega\bigr\|_{\rm prod}
 =
 \frac{|d_0|}{d}.
 \label{eq:V116-singlet-one-power}
\end{equation}
Thus the V116 completion retains the sharp single inverse dimension;
it neither loses it under positive completion nor replaces it by the
false arbitrary-coefficient weight \(d^{-2}\).
\end{proposition}

\begin{proof}
This is
\eqref{eq:V112-singlet-A}.  The V113 positive-completion theorem
majorizes each physical block by the same positive pair multiplier and
uses its product capacity.  Hence its input in this block is exactly
\(|d_0|/d\).  The distinction from the representation-dimension law is
the dual-singlet normalization counterexample
\Cref{cth:V45-dual-singlet-normalization}.
\end{proof}

\begin{proposition}[Alternating-height check]
\label{prop:V116-alternating-height-check}
Apply the V60 alternating-height construction with \(r=c\), so that
\[
 \Xi_a=\Xi_c=\frac N2(L+1),
 \qquad
 \widehat H_{ac}=NL.
\]
If this configuration lies in the selected-heavy off-cut covariance
branch, the tie belongs to the maximal-\(c\) region and
\begin{equation}
 \frac{H_{ac}^{1}}{\Xi_c^2}
 =
 \frac{H_{ac}^{1}}{\Xi_a\Xi_c}.
 \label{eq:V116-alternating-exact}
\end{equation}
Consequently V116 gives the required two-endpoint mass denominator for
the \(ac\) summand without selecting an orientation separately on the
two coordinate classes.
\end{proposition}

\begin{proof}
The equality of the two total masses is part of
\Cref{prop:V60-alternating-regression}; the priority convention in
\Cref{prop:V116-exact-residual-order} assigns equality to
\(\Xi_c\ge\Xi_a\).  Identity
\eqref{eq:V116-alternating-exact} is immediate.  More importantly, the
proof of \Cref{thm:V116-physical-covariance-capacity} first completes
the whole \(bc\) multiplier and the V98 outer family, and only then
compares the total masses.  It therefore makes no coordinatewise
choice between the two height classes.
\end{proof}

\begin{firewall}[The regression checks do not enlarge the theorem]
\label{fire:V116-regression-scope}
The singlet calculation verifies the normalization of the positive
pair carrier, not an additional dimension-law tail.  The
alternating-height calculation verifies the endpoint denominator only
for configurations already satisfying the selected-heavy off-cut
hypotheses.  Neither check proves the residual ordering
\eqref{eq:V116-residual-order}.
\end{firewall}

\section{V116 ledger and next upstream acquisition}
\label{sec:V116-ledger}

\begin{theorem}[Integrated V116 statement]
\label{thm:V116-integrated}
Preserving V1--V115, V116 proves:
\begin{enumerate}[label=\textup{(V116.\arabic*)},leftmargin=6.7em]
\item survival of the unused \(c\)-mean debit through the paired
      \(bc\) mean, \eqref{eq:V116-offcut-mean-endpoint};
\item conditional marked-path payment of the peeled mean branch;
\item the selected-heavy positive \(bc\)-pair capacity
      \eqref{eq:V116-offcut-pair-capacity};
\item the physical correlated-outer covariance estimate
      \eqref{eq:V116-covariance-absolute};
\item the maximal-\(c\) marked-tree conversion
      \eqref{eq:V116-covariance-marked}, with the \(bc\) factor counted
      once;
\item closure of the covariance branch in the maximal-\(c\) region;
\item the exact remaining mass order
      \eqref{eq:V116-residual-order}; and
\item the dual-singlet and alternating-height checks
      \eqref{eq:V116-singlet-one-power} and
      \eqref{eq:V116-alternating-exact}.
\end{enumerate}
\end{theorem}

\begin{proof}
The items are respectively
\Cref{lem:V116-paired-mean-debit,
thm:V116-mean-branch-paid,
lem:V116-offcut-pair-capacity,
thm:V116-physical-covariance-capacity,
cor:V116-maximal-c-conversion,
thm:V116-covariance-cmax-paid,
prop:V116-exact-residual-order,
prop:V116-dual-singlet-check,
prop:V116-alternating-height-check}.
\end{proof}

\begin{openproblem}[V117 mandate: one shared-row \(ab\)--\(bc\)
carrier]
\label{op:V117-mandate}
Work in the residual regime
\(\Xi_a>\Xi_c>C_{\rm row}\Xi_b\).
\begin{enumerate}
\item Retain the old \(ab\) pair multiplier, the peeled \(bc\) pair
      multiplier, and the reduced signed core inside one physical V98
      outer block before taking trace norm.
\item Resolve both positive pair majorants by a common two-level
      layer cake.  Their coordinate banks are disjoint but their
      original-row incidence shares \(b\).
\item Apply one Schur minimum to the complete carrier, as in V104;
      do not multiply the separately optimized \(ab\) and \(bc\)
      capacities.
\item Record the actual two retained-row representation dimensions
      and the orientation factor.  In particular, do not infer them
      from \(\Xi_a,\Xi_b,\Xi_c\).
\item Determine whether the ordering
      \eqref{eq:V116-residual-order} supplies a fixed orientation or
      whether the alternating orientation remains a separate spectral
      gate.
\item Test the common carrier on a double dual-singlet block and on
      the V60 alternating-height array.
\item If the V104 effective-scale stock is still needed, isolate it as
      an explicit conditional hypothesis rather than deriving it from
      pair heaviness.
\end{enumerate}
\end{openproblem}

\section{Firewall after V116}
\label{sec:V116-firewall}

\begin{firewall}[Current claim boundary]
\label{fire:V116-current-boundary}
V116 pays the peeled mean branch under the stated endpoint-response
hypothesis and pays the peeled covariance branch when \(c\) is a
maximal internal row.  It does not close the shared-row residual
\eqref{eq:V116-residual-order}, force either V104 effective scale,
combine noncommuting retained-row orientations, prove JREC, prove
quartic or all-order MTA, construct the continuum theory, or prove a
positive spectral gap.  The full Yang--Mills existence and mass gap
remain open.
\end{firewall}

\section{Append-only record for V116}
\label{sec:V116-record}

V116 was appended to the complete V115 manuscript.  The exact V115
parent had \(77707\) newline-delimited source records, \(2710319\)
bytes, and SHA--256 digest
\[
 \texttt{e9a32c083df3c8d5a4bb5728360ec2fab19ee31ad5d3ecc30508fe69d59e01e8}.
\]
No V115 mathematical content was changed.  The sole final
\verb|\end{document}| is restored below.

% =============================================================================
% END APPENDED VERSION V116 MATERIAL
% =============================================================================

% =============================================================================
% BEGIN APPENDED VERSION V117 MATERIAL
% =============================================================================

\chapter{V117: Three-Row Shared-Vertex Pair Law and the Exact
Orientation Obstruction}
\label{ch:V117}

\section{Purpose}
\label{sec:V117-purpose}

The residual V116 carrier contains two disjoint coordinate banks but
three original rows:
\[
 e=ab,
 \qquad
 f=bc.
\]
The banks meet at row \(b\).  Although their Hilbert factors tensor by
coordinate, a product vector over the \emph{original} rows may be
entangled between the \(e\)- and \(f\)-factors inside row \(b\).
Therefore neither the product of two separately optimized pair
capacities nor a verbatim invocation of the two-row V104 law is
legitimate.

V117 derives the required three-row law from first principles.  It
keeps one product vector, one positive core shell, and both pair
levels until a single Schur minimum is taken.  The law is then
compressed into an exact combined Ky--Fan profile.  A double-singlet
calculation is sharp, while an explicit orientation comparison proves
that the residual scalar mass order
\(\Xi_a>\Xi_c>C_{\rm row}\Xi_b\) does not by itself select the
coherent row dimension.

\section{The shared-vertex carrier}
\label{sec:V117-shared-carrier}

Let \(G_e,G_f\) be the compact gauge groups carried by the two
disjoint coordinate banks.  Their row representations are
\[
 X_e\ \hbox{on \(a\)},\quad Y_e\ \hbox{on \(b\)},
 \qquad
 X_f\ \hbox{on \(b\)},\quad Y_f\ \hbox{on \(c\)}.
\]
Let \(A,B,C\) be irreducible row representations of the remaining
core gauge bank.  Thus
\begin{align}
 \mathcal H_a
 &=\mathcal H_{X_e}\otimes\mathcal H_A,
 \notag\\
 \mathcal H_b
 &=\mathcal H_{Y_e}\otimes\mathcal H_{X_f}
   \otimes\mathcal H_B,
 \notag\\
 \mathcal H_c
 &=\mathcal H_{Y_f}\otimes\mathcal H_C.
 \label{eq:V117-three-row-spaces}
\end{align}
Write the corresponding dimensions as
\[
 d_{X_e},d_{Y_e},d_{X_f},d_{Y_f},d_A,d_B,d_C.
\]

Let \(Q_e,Q_f\ge0\) be invariant pair multipliers and let \(M\ge0\)
be an invariant core shell on
\(\mathcal H_A\otimes\mathcal H_B\otimes\mathcal H_C\).  Use the layer
resolutions
\begin{align}
 Q_e&=\int_0^{A_e}P_e(t)\,\dd t,
 &
 Q_f&=\int_0^{A_f}P_f(s)\,\dd s,
 \label{eq:V117-layer-resolutions}\\
 D_e(t)&=\operatorname{rank}P_e(t),
 &
 D_f(s)&=\operatorname{rank}P_f(s).
 \notag
\end{align}
For a positive operator \(T\) on the three original rows, define
\begin{equation}
 \|T\|_{\rm 3prod}
 :=
 \sup_{\substack{\|u\|=\|v\|=\|w\|=1}}
 \langle u\otimes v\otimes w,\,
 T(u\otimes v\otimes w)\rangle.
 \label{eq:V117-three-product-norm}
\end{equation}
Set
\begin{equation}
 \Pi_3(M):=\|M\|_{\rm 3prod},
 \qquad
 \mathcal H_3(M):=\operatorname{Tr}M.
 \label{eq:V117-core-data}
\end{equation}

\begin{lemma}[The core product cap survives both shared-row
projections]
\label{lem:V117-core-cap}
For every \(t,s\),
\begin{equation}
 \|P_e(t)\otimes P_f(s)\otimes M\|_{\rm 3prod}
 \le\Pi_3(M).
 \label{eq:V117-core-cap}
\end{equation}
\end{lemma}

\begin{proof}
Since \(0\le P_e(t)\otimes P_f(s)\le I\), it suffices to use the
auxiliary identity.  Let \(u_a\otimes v_b\otimes w_c\) be a unit
product vector over the three original rows.  The vector \(v_b\) may
be entangled among
\(\mathcal H_{Y_e},\mathcal H_{X_f},\mathcal H_B\).  Nevertheless,
after tracing the auxiliary factors from each \emph{separate original
row}, the core state is
\[
 \rho_A\otimes\rho_B\otimes\rho_C.
\]
Each \(\rho\) is a density matrix.  Resolve all three into convex
combinations of pure states.  Then
\[
 \operatorname{Tr}
 [M(\rho_A\otimes\rho_B\otimes\rho_C)]
\]
is a convex combination of three-row product-vector expectations of
\(M\), hence is at most \(\Pi_3(M)\).
\end{proof}

\begin{lemma}[Three-row double-level Schur cap]
\label{lem:V117-three-row-Schur}
For every \(t,s\),
\begin{align}
 &\|P_e(t)\otimes P_f(s)\otimes M\|_{\rm 3prod}
 \notag\\
 &\quad\le
 \frac{D_e(t)D_f(s)\mathcal H_3(M)}
 {\max\{
 d_{X_e}d_A,\,
 d_{Y_e}d_{X_f}d_B,\,
 d_{Y_f}d_C
 \}}.
 \label{eq:V117-three-row-Schur}
\end{align}
\end{lemma}

\begin{proof}
Put
\[
 T_{t,s}:=P_e(t)\otimes P_f(s)\otimes M.
\]
It is positive and invariant under the product of the two auxiliary
gauge groups and the core gauge group.  Each row representation in
\eqref{eq:V117-three-row-spaces} is an external tensor product of
irreducibles, with trivial factors for groups absent from that row,
and is therefore irreducible for the product group acting on it.
Schur's lemma gives
\begin{align*}
 \operatorname{Tr}_{bc}T_{t,s}
 &=
 \frac{\operatorname{Tr}T_{t,s}}{d_{X_e}d_A}I_a,\\
 \operatorname{Tr}_{ac}T_{t,s}
 &=
 \frac{\operatorname{Tr}T_{t,s}}
 {d_{Y_e}d_{X_f}d_B}I_b,\\
 \operatorname{Tr}_{ab}T_{t,s}
 &=
 \frac{\operatorname{Tr}T_{t,s}}{d_{Y_f}d_C}I_c.
\end{align*}
For every unit \(u\otimes v\otimes w\), positivity bounds its
expectation by each of the three one-row marginal expectations.
Take their minimum and use
\[
 \operatorname{Tr}T_{t,s}
 =D_e(t)D_f(s)\mathcal H_3(M).
\]
The minimum of the three reciprocal dimensions is the reciprocal of
their maximum, proving the claim.
\end{proof}

\begin{theorem}[One-minimum shared-vertex layer law]
\label{thm:V117-one-minimum-law}
Under the preceding hypotheses,
\begin{align}
 &\|Q_e\otimes Q_f\otimes M\|_{\rm 3prod}
 \notag\\
 &\quad\le
 \int_0^{A_e}\!\!\int_0^{A_f}
 \min\left\{
 \Pi_3(M),\,
 \frac{D_e(t)D_f(s)\mathcal H_3(M)}
 {\max\{
 d_{X_e}d_A,\,
 d_{Y_e}d_{X_f}d_B,\,
 d_{Y_f}d_C
 \}}
 \right\}\dd s\,\dd t.
 \label{eq:V117-one-minimum-law}
\end{align}
\end{theorem}

\begin{proof}
Tensor the two layer resolutions in
\eqref{eq:V117-layer-resolutions}.  Finite-dimensional Bochner
integration and subadditivity of the positive product numerical radius
give the integral of the levelwise radii.  At each common level
\((t,s)\), take the minimum of
\Cref{lem:V117-core-cap,lem:V117-three-row-Schur}.
\end{proof}

\begin{firewall}[No multiplication of shared-row capacities]
\label{fire:V117-no-capacity-product}
Equation \eqref{eq:V117-one-minimum-law} does not assert
\[
 \|Q_e\otimes Q_f\otimes M\|_{\rm 3prod}
 \le
 \|Q_e\|_{\rm prod}\|Q_f\|_{\rm prod}\Pi_3(M).
\]
The vector in row \(b\) can entangle the two auxiliary bank factors,
so the separately maximizing pair vectors need not be compatible.
The valid theorem keeps the same three-row product vector and takes
one common minimum.
\end{firewall}

\section{Coherent star dimension and the combined spectrum}
\label{sec:V117-star-dimension}

\begin{definition}[Coherent shared-vertex dimension]
\label{def:V117-star-dimension}
Put
\begin{align}
 d_e&:=\max\{d_{X_e},d_{Y_e}\},
 &
 d_f&:=\max\{d_{X_f},d_{Y_f}\},
 \notag\\
 d_\star
 &:=
 \max\{d_{X_e},d_{Y_e}d_{X_f},d_{Y_f}\},
 &
 \chi_\star&:=\frac{d_ed_f}{d_\star}.
 \label{eq:V117-star-dimension}
\end{align}
The middle entry retains the fact that both pair banks meet inside
one original row.
\end{definition}

\begin{lemma}[Exact range of the star orientation loss]
\label{lem:V117-star-orientation-range}
One has
\begin{equation}
 1\le\chi_\star\le\min\{d_e,d_f\}.
 \label{eq:V117-star-orientation-range}
\end{equation}
If the larger side of both pair representations lies on the common
row \(b\), then \(\chi_\star=1\).  For the outward pattern
\[
 (d_{X_e},d_{Y_e})=(L,1),
 \qquad
 (d_{X_f},d_{Y_f})=(1,L),
\]
one has \(d_\star=L\) and \(\chi_\star=L\).
\end{lemma}

\begin{proof}
Every dimension is at least one.  Hence each entry defining
\(d_\star\) is at most \(d_ed_f\), while those entries jointly imply
\[
 d_\star\ge\max\{d_e,d_f\}.
\]
Dividing \(d_ed_f\) by these two inequalities gives
\eqref{eq:V117-star-orientation-range}.  If both maxima lie on \(b\),
then \(d_{Y_e}d_{X_f}=d_ed_f\).  Direct substitution proves the
outward example.
\end{proof}

\begin{proposition}[Relaxed common-ratio law]
\label{prop:V117-ratio-law}
Define
\begin{equation}
 B_3(M):=
 \frac{\mathcal H_3(M)}{\min\{d_A,d_B,d_C\}},
 \qquad
 \eta_3(M):=\frac{\Pi_3(M)}{B_3(M)}.
 \label{eq:V117-ratio-data}
\end{equation}
Then \(0\le\eta_3(M)\le1\), and
\begin{align}
 \|Q_e\otimes Q_f\otimes M\|_{\rm 3prod}
 \le
 B_3(M)
 \int_0^{A_e}\!\!\int_0^{A_f}
 \min\left\{
 \eta_3(M),\frac{D_e(t)D_f(s)}{d_\star}
 \right\}\dd s\,\dd t.
 \label{eq:V117-ratio-law}
\end{align}
\end{proposition}

\begin{proof}
Apply the three-row Schur argument to \(M\) alone.  It gives
\[
 \Pi_3(M)
 \le
 \frac{\mathcal H_3(M)}{\max\{d_A,d_B,d_C\}}
 \le B_3(M),
\]
so \(0\le\eta_3(M)\le1\).  Moreover,
\begin{align*}
 &\max\{
 d_{X_e}d_A,\,
 d_{Y_e}d_{X_f}d_B,\,
 d_{Y_f}d_C
 \}\\
 &\qquad\ge
 d_\star\min\{d_A,d_B,d_C\}.
\end{align*}
Insert this lower bound into
\eqref{eq:V117-one-minimum-law} and factor out \(B_3(M)\).
\end{proof}

Regard \(Q_{ef}:=Q_e\otimes Q_f\) as one positive operator.  If the
eigenvalues of \(Q_e,Q_f\), with multiplicity, are
\[
 a_1\ge a_2\ge\cdots\ge0,
 \qquad
 b_1\ge b_2\ge\cdots\ge0,
\]
then those of \(Q_{ef}\) are the products \(a_ib_j\).  List them as
\[
 c_1^\star\ge c_2^\star\ge\cdots\ge0
\]
and set
\[
 D_\star(r):=\operatorname{rank}\mathbf1_{\{Q_{ef}>r\}}.
\]

\begin{theorem}[Combined star Ky--Fan law]
\label{thm:V117-star-KyFan}
For \(0\le\eta\le1\), define
\begin{equation}
 \mathcal K_\star(\eta)
 :=
 \int_0^\infty
 \min\left\{\eta,\frac{D_\star(r)}{d_\star}\right\}\dd r.
 \label{eq:V117-star-profile}
\end{equation}
If \(q=\eta d_\star\) and \(n=\lfloor q\rfloor\), then
\begin{equation}
 \boxed{
 \mathcal K_\star(\eta)
 =
 \frac1{d_\star}
 \left[
 \sum_{\ell=1}^{n}c_\ell^\star
 +(q-n)c_{n+1}^\star
 \right],}
 \label{eq:V117-star-KyFan}
\end{equation}
with the fractional term omitted when \(q=n\).  Furthermore,
\begin{equation}
 \boxed{
 \|Q_e\otimes Q_f\otimes M\|_{\rm 3prod}
 \le
 B_3(M)\mathcal K_\star(\eta_3(M)).}
 \label{eq:V117-star-core-gate}
\end{equation}
\end{theorem}

\begin{proof}
Use the layer cake
\[
 Q_{ef}
 =
 \int_0^\infty\mathbf1_{\{Q_{ef}>r\}}\,\dd r.
\]
The spectral projector is invariant, has rank \(D_\star(r)\), and
acts on the same three auxiliary row representations.  Repeating
\Cref{lem:V117-core-cap,lem:V117-three-row-Schur} and using the
relaxation in \Cref{prop:V117-ratio-law} bounds its contribution by
\[
 B_3(M)
 \min\left\{\eta_3(M),\frac{D_\star(r)}{d_\star}\right\}.
\]
Integration proves \eqref{eq:V117-star-core-gate}.  Finally, integrate
the clipped rank distribution across the ordered eigenvalues.  The
first \(n\) eigenvalues enter fully and the next enters with fraction
\(q-n\), which gives \eqref{eq:V117-star-KyFan}.
\end{proof}

\section{A sharp common cold--cold criterion}
\label{sec:V117-cold-cold}

The combined law shows exactly which part of the two pair spectra can
remain visible to a small core product cap.  The next lemma is
elementary, but its denominator is the coherent star dimension rather
than a product of independently selected row dimensions.

\begin{lemma}[Shared-vertex cold--cold rectangle]
\label{lem:V117-cold-cold}
Suppose
\[
 0\le a_i\le A_e,
 \qquad
 0\le b_j\le A_f.
\]
Let \(\mathscr C_e,\mathscr C_f\) have cardinalities \(r_e,r_f\), and
assume
\[
 a_i\le h_e\quad(i\notin\mathscr C_e),
 \qquad
 b_j\le h_f\quad(j\notin\mathscr C_f).
\]
Put
\begin{equation}
 H_{ef}:=\max\{h_eA_f,A_eh_f\}.
 \label{eq:V117-hot-product}
\end{equation}
For \(0<\eta\le1\), let \(q=\eta d_\star\) and define the fractional
top mean
\begin{equation}
 \overline c_\star(\eta)
 :=
 \frac1q
 \left[
 \sum_{\ell=1}^{\lfloor q\rfloor}c_\ell^\star
 +(q-\lfloor q\rfloor)
 c_{\lfloor q\rfloor+1}^\star
 \right].
 \label{eq:V117-star-top-mean}
\end{equation}
Then
\begin{equation}
 \boxed{
 \overline c_\star(\eta)
 \le
 H_{ef}
 +
 A_eA_f
 \min\left\{
 1,\frac{r_er_f}{\eta d_\star}
 \right\}.}
 \label{eq:V117-cold-cold-bound}
\end{equation}
\end{lemma}

\begin{proof}
Outside
\(\mathscr C_e\times\mathscr C_f\), at least one factor is hot and
hence \(a_ib_j\le H_{ef}\).  At most \(r_er_f\) product eigenvalues
belong to the cold--cold rectangle, and every eigenvalue is at most
\(A_eA_f\).  Therefore the fractional sum of the largest \(q\)
eigenvalues is at most
\[
 qH_{ef}+A_eA_f\min\{q,r_er_f\}.
\]
Divide by \(q\).
\end{proof}

\begin{corollary}[Exact sufficient spectral gate]
\label{cor:V117-spectral-gate}
If \(\eta_3(M)>0\), then
\begin{align}
 \|Q_e\otimes Q_f\otimes M\|_{\rm 3prod}
 &\le
 \Pi_3(M)
 \left[
 H_{ef}
 +
 A_eA_f
 \min\left\{
 1,
 \frac{r_er_f}{\eta_3(M)d_\star}
 \right\}
 \right].
 \label{eq:V117-spectral-gate}
\end{align}
Consequently the normalized carrier tends to zero whenever
\begin{equation}
 H_{ef}\longrightarrow0,
 \qquad
 \frac{r_er_f}{\eta_3(M)d_\star}\longrightarrow0,
 \label{eq:V117-spectral-sufficient}
\end{equation}
with \(A_e,A_f\) uniformly bounded.
\end{corollary}

\begin{proof}
By \eqref{eq:V117-star-KyFan},
\[
 \mathcal K_\star(\eta)
 =\eta\,\overline c_\star(\eta).
\]
Insert this identity into
\eqref{eq:V117-star-core-gate}, use
\(B_3(M)\eta_3(M)=\Pi_3(M)\), and apply
\Cref{lem:V117-cold-cold}.  The convergence assertion follows
directly.
\end{proof}

\begin{remark}[The effective-scale stock is explicit]
\label{rem:V117-effective-stock}
The second condition in
\eqref{eq:V117-spectral-sufficient} is the precise shared-vertex
cold-rank requirement.  Pair heaviness supplies neither
\(r_er_f=o(\eta_3d_\star)\) nor a bound on \(\chi_\star\).  These
remain spectral/orientation data of the assembled carrier and are not
silently inferred from \(\Xi_a,\Xi_b,\Xi_c\).
\end{remark}

\section{Sharp tests and the orientation decision}
\label{sec:V117-tests}

\begin{proposition}[Two dual singlets sharing the middle row]
\label{prop:V117-double-singlet}
Let \(P_d\) be the normalized dual-singlet projector on
\(\mathbb C^d_a\otimes\mathbb C^d_{b,e}\), and let \(P_m\) be the
normalized dual-singlet projector on
\(\mathbb C^m_{b,f}\otimes\mathbb C^m_c\).  Then
\begin{equation}
 \boxed{
 \|P_d\otimes P_m\|_{\rm 3prod}
 =\frac1{dm}.}
 \label{eq:V117-double-singlet}
\end{equation}
The equality allows an arbitrary, possibly entangled, unit vector on
\(\mathbb C^d_{b,e}\otimes\mathbb C^m_{b,f}\).
\end{proposition}

\begin{proof}
Write
\[
 \Omega_d=d^{-1/2}\sum_i e_i\otimes e_i^*,
 \qquad
 \Omega_m=m^{-1/2}\sum_j f_j\otimes f_j^*.
\]
For fixed unit vectors \(x_a,z_c\), contraction of
\(\Omega_d\otimes\Omega_m\) against \(x_a\) and \(z_c\) produces in
the middle row a vector of norm \(1/\sqrt{dm}\).  Cauchy--Schwarz
therefore gives
\[
 |\langle\Omega_d\otimes\Omega_m,\,
 x_a\otimes y_b\otimes z_c\rangle|^2
 \le\frac1{dm}.
\]
Choose \(y_b\) proportional to that contracted vector to attain
equality.  This proves the formula.
\end{proof}

\begin{assessment}[Why the singlet equality matters]
\label{ass:V117-singlet-meaning}
For the double-singlet block,
\[
 d_\star=\max\{d,dm,m\}=dm.
\]
Thus the middle-row entry in
\eqref{eq:V117-three-row-Schur} is sharp.  The common law correctly
produces the product \(1/(dm)\) even though multiplication of
\emph{arbitrary separately optimized} pair capacities is forbidden.
The gain comes from one coherent three-row computation, not from
combining two unrelated suprema.
\end{assessment}

\begin{proposition}[The residual mass order does not determine the
star orientation]
\label{prop:V117-mass-not-orientation}
At the level of the shared-vertex carrier hypotheses, the same scalar
triple satisfying
\[
 \Xi_a>\Xi_c>C_{\rm row}\Xi_b
\]
is compatible with either of the dimension patterns
\begin{align}
 &(d_{X_e},d_{Y_e};d_{X_f},d_{Y_f})
 =(1,L;L,1),
 \label{eq:V117-middle-pattern}\\
 &(d_{X_e},d_{Y_e};d_{X_f},d_{Y_f})
 =(L,1;1,L).
 \label{eq:V117-outward-pattern}
\end{align}
In the first, \(d_\star=L^2\) and \(\chi_\star=1\); in the second,
\(d_\star=L\) and \(\chi_\star=L\).  Hence no proof using the scalar
mass order alone can replace \(d_\star\) by \(d_ed_f\).
\end{proposition}

\begin{proof}
The scalar mass order is not among the representation-dimension
hypotheses of
\Cref{thm:V117-one-minimum-law}; fix any positive triple having that
order.  The two displayed dimension arrays have
\(d_e=d_f=L\).  Direct substitution into
\eqref{eq:V117-star-dimension} gives respectively
\[
 d_\star=L^2,\ \chi_\star=1,
 \qquad
 d_\star=L,\ \chi_\star=L.
\]
Thus the conclusion is not a logical consequence of the scalar
inequality.  This is an abstract information test: an application to
the Yang--Mills carrier may still prove a dimension--mass relation
from its exact representation labels.
\end{proof}

\begin{assessment}[Alternating-height regression]
\label{ass:V117-alternating-height}
The V60 alternating-height masses likewise do not encode which side
of either pair representation has the larger dimension.  V117 does
remove the forbidden \emph{coordinatewise} choice: all coordinates
enter \(Q_e,Q_f\), then one \(d_\star\) and one common Ky--Fan
truncation are used.  But
\Cref{prop:V117-mass-not-orientation} shows that the scalar
alternation cannot force the favorable middle-row pattern.  The
remaining loss is exactly \(\chi_\star\), not an unspecified
orientation issue.
\end{assessment}

\section{Physical interface with the V116 residual}
\label{sec:V117-physical-interface}

\begin{corollary}[Conditional physical shared-vertex gate]
\label{cor:V117-physical-gate}
Fix a V98 outer label in the residual ordering
\eqref{eq:V116-residual-order}.  Suppose its absolute compressed
carrier is an isometric compression of, or is positively dominated
by,
\begin{equation}
 Q_{ab}\otimes Q_{bc}^{0}\otimes M,
 \label{eq:V117-physical-factorization}
\end{equation}
where the first two factors are the complete positive majorants of
the old same-pair and peeled off-cut pair multipliers on disjoint
coordinate banks, and \(M\) is the positive square of the remaining
signed core.  Then its physical three-row product capacity is bounded
by
\eqref{eq:V117-one-minimum-law}, equivalently by
\eqref{eq:V117-star-core-gate}.  Under
\eqref{eq:V117-spectral-sufficient}, that normalized block tends to
zero.
\end{corollary}

\begin{proof}
Positive domination and isometric compression cannot increase a
product-vector expectation.  Apply
\Cref{thm:V117-one-minimum-law,thm:V117-star-KyFan} to the right-hand
side of \eqref{eq:V117-physical-factorization}.  The final assertion
is \Cref{cor:V117-spectral-gate}.
\end{proof}

\begin{firewall}[The physical factorization is now the upstream gate]
\label{fire:V117-factorization-open}
V117 proves the analytic estimate \emph{after}
\eqref{eq:V117-physical-factorization}.  It does not assert that the
V98 Racah compression, the two signed pair multipliers, and the
reduced core automatically admit that domination.  Establishing the
factorization without replacing correlated outer labels by
independent ones is the next physical task.  If it holds, the separate
spectral conditions \eqref{eq:V117-spectral-sufficient} must still be
verified.
\end{firewall}

\section{V117 ledger and V118 mandate}
\label{sec:V117-ledger}

\begin{theorem}[Integrated V117 statement]
\label{thm:V117-integrated}
Preserving V1--V116, V117 proves:
\begin{enumerate}[label=\textup{(V117.\arabic*)},leftmargin=6.7em]
\item the exact three-original-row carrier
      \eqref{eq:V117-three-row-spaces};
\item survival of the core product cap despite middle-row
      entanglement, \eqref{eq:V117-core-cap};
\item the common three-row Schur cap
      \eqref{eq:V117-three-row-Schur};
\item the one-minimum shared-vertex law
      \eqref{eq:V117-one-minimum-law};
\item the coherent dimension and exact orientation loss
      \eqref{eq:V117-star-dimension}--%
      \eqref{eq:V117-star-orientation-range};
\item the combined star Ky--Fan law
      \eqref{eq:V117-star-KyFan}--%
      \eqref{eq:V117-star-core-gate};
\item the sufficient common cold--cold condition
      \eqref{eq:V117-spectral-sufficient};
\item the sharp shared-middle-row double-singlet identity
      \eqref{eq:V117-double-singlet}; and
\item the proof that the V116 residual mass ordering alone does not
      determine the star orientation.
\end{enumerate}
\end{theorem}

\begin{proof}
The items are
\Cref{sec:V117-shared-carrier,
lem:V117-core-cap,
lem:V117-three-row-Schur,
thm:V117-one-minimum-law,
def:V117-star-dimension,
lem:V117-star-orientation-range,
thm:V117-star-KyFan,
cor:V117-spectral-gate,
prop:V117-double-singlet,
prop:V117-mass-not-orientation}.
\end{proof}

\begin{openproblem}[V118 mandate: identify the physical star carrier]
\label{op:V118-mandate}
\begin{enumerate}
\item Return to the exact V98 compressed coefficient before trace
      norm and write the old \(ab\) and off-cut \(bc\) multipliers on
      their disjoint coordinate tensor factors.
\item Complete the two signed multipliers jointly.  Prove
      \eqref{eq:V117-physical-factorization}, or retain the exact
      positive remainder obstructing it.
\item Keep the V110 \(ab\)- and \(ac\)-hybrid core rows inside one
      positive \(2\times2\) block so their interference is not
      discarded.
\item Determine the actual
      \(X_e,Y_e,X_f,Y_f\) labels selected by the V98 priority
      partition and calculate \(d_\star\) before any dimension bound.
\item Apply the V53 cold/hot data to the \emph{combined} product
      spectrum and test
      \eqref{eq:V117-spectral-sufficient}; do not multiply marginal
      capacities.
\item If the cold-rank ratio fails, exhibit the exact persistent
      representation family and reroute upstream.
\end{enumerate}
\end{openproblem}

\section{Firewall after V117}
\label{sec:V117-firewall}

\begin{firewall}[Current claim boundary]
\label{fire:V117-current-boundary}
V117 supplies the rigorous analytic shared-vertex law and isolates its
exact orientation and cold-rank data.  It does not prove the physical
factorization \eqref{eq:V117-physical-factorization}, the spectral
limits \eqref{eq:V117-spectral-sufficient}, the residual off-cut
closure, JREC, quartic or all-order MTA, continuum Yang--Mills
existence, or a positive spectral gap.  The full Yang--Mills
existence and mass gap remain open.
\end{firewall}

\section{Append-only record for V117}
\label{sec:V117-record}

V117 was appended to the complete V116 manuscript.  The exact V116
parent had \(78194\) newline-delimited source records, \(2728028\)
bytes, and SHA--256 digest
\[
 \texttt{7c1fa8c9fdb101085c90a83e1a063b5c24d3e8b5b52a0491f43b8bb6bf5bb6d5}.
\]
No V116 mathematical content was changed.  The sole final
\verb|\end{document}| is restored below.

% =============================================================================
% END APPENDED VERSION V117 MATERIAL
% =============================================================================

% =============================================================================
% BEGIN APPENDED VERSION V118 MATERIAL
% =============================================================================

\chapter{V118: Joint Physical Positive Completion for the
Shared-Vertex Off-Cut Carrier}
\label{ch:V118}

\section{Purpose and correction of the provisional interface}
\label{sec:V118-purpose}

V117 left
\eqref{eq:V117-physical-factorization} as a physical identification
gate.  Returning to the V98 outer blocks and the V113
positive-completion argument shows that exact rectangularity is not
needed.  Every complete outer label already refines the output
resolutions of both disjoint coordinate banks.  A two-label version of
V113 therefore completes the correlated family directly.

There is one necessary care.  The reduced composite covariance is the
sum of the exact V110 \(ab\)- and \(ac\)-hybrid amplitudes.  The absolute
value of their sum is not bounded in operator order by the sum of their
absolute values.  V118 instead uses linearity of the physical
coefficient compression and the trace-norm triangle, producing two
separate physical absolute carriers.  Each carrier receives the same
joint \(ab\)--\(bc\) positive multiplier, and no pair capacity is
optimized separately.

\section{Two-label positive completion}
\label{sec:V118-double-completion}

Let
\[
 \mathbf A_e=\sum_\lambda a_\lambda P_\lambda^e,
 \qquad
 \mathbf A_f=\sum_\mu b_\mu P_\mu^f
\]
be positive invariant multipliers on disjoint coordinate banks, with
\(a_\lambda,b_\mu\ge0\).  Let
\(\{V_j:G_j\to\mathcal H_{\rm rows}\}_{j\in\mathfrak J}\) have
orthogonal ranges and put \(P_j=V_jV_j^*\).  Suppose each \(j\) carries
two labels \(\lambda(j),\mu(j)\).

\begin{theorem}[Joint completion of two correlated pair labels]
\label{thm:V118-joint-positive-completion}
Assume that the physical outer projections jointly refine the two pair
resolutions:
\begin{equation}
 \sum_{\substack{j\in\mathfrak J\\
 \lambda(j)=\lambda,\ \mu(j)=\mu}}
 P_j
 \preccurlyeq
 P_\lambda^e\otimes P_\mu^f\otimes I_{\rm core}.
 \label{eq:V118-joint-refinement}
\end{equation}
If every signed block \(L_j\) obeys
\begin{equation}
 |V_jL_jV_j^*|
 \preccurlyeq
 m\,a_{\lambda(j)}b_{\mu(j)}P_j
 \label{eq:V118-joint-block-majorant}
\end{equation}
for one \(m\ge0\), then its complete physical absolute carrier
\[
 \mathbf Q_{\mathfrak J}
 :=
 \sum_{j\in\mathfrak J}|V_jL_jV_j^*|
\]
satisfies
\begin{align}
 \mathbf Q_{\mathfrak J}
 &\preccurlyeq
 m(\mathbf A_e\otimes\mathbf A_f\otimes I_{\rm core}),
 \label{eq:V118-joint-completed-order}\\
 \|\mathbf Q_{\mathfrak J}\|_{\rm rows,prod}
 &\le
 m\|\mathbf A_e\otimes\mathbf A_f\|_{\rm 3prod}.
 \label{eq:V118-joint-completed-capacity}
\end{align}
The label set \(\mathfrak J\) need not be rectangular.
\end{theorem}

\begin{proof}
Orthogonality of the outer ranges gives the displayed sum of block
absolute values.  Apply
\eqref{eq:V118-joint-block-majorant}, group by the pair
\((\lambda,\mu)\), and use
\eqref{eq:V118-joint-refinement}:
\begin{align*}
 \mathbf Q_{\mathfrak J}
 &\preccurlyeq
 m\sum_{\lambda,\mu}a_\lambda b_\mu
 \sum_{j:\lambda(j)=\lambda,\,\mu(j)=\mu}P_j\\
 &\preccurlyeq
 m\sum_{\lambda,\mu}a_\lambda b_\mu
 (P_\lambda^e\otimes P_\mu^f\otimes I_{\rm core})\\
 &=
 m(\mathbf A_e\otimes\mathbf A_f\otimes I_{\rm core}).
\end{align*}
Take expectations in product states over the original rows.  The core
identity is an arbitrary within-row spectator and disappears by the
same reduced-density argument as
\Cref{lem:V117-core-cap}.  This proves
\eqref{eq:V118-joint-completed-capacity}.
\end{proof}

\begin{corollary}[Blockwise factorization is sufficient]
\label{cor:V118-block-factorization}
Suppose
\begin{equation}
 L_j=d_{\lambda(j)}^e d_{\mu(j)}^f R_j,
 \qquad
 \|R_j\|_{\rm op}\le m,
 \label{eq:V118-block-factorization}
\end{equation}
and
\[
 |d_\lambda^e|\le a_\lambda,
 \qquad
 |d_\mu^f|\le b_\mu.
\]
Then \eqref{eq:V118-joint-block-majorant} holds.
\end{corollary}

\begin{proof}
On the range of \(P_j\),
\[
 |L_j|
 =
 |d_{\lambda(j)}^e d_{\mu(j)}^f|\,|R_j|
 \preccurlyeq
 m a_{\lambda(j)}b_{\mu(j)}I_{G_j}.
\]
Conjugate by \(V_j\).
\end{proof}

\begin{firewall}[What the joint completion forgets]
\label{fire:V118-completion-scope}
The passage from \eqref{eq:V118-joint-block-majorant} to
\eqref{eq:V118-joint-completed-order} forgets correlations only after
absolute values have been formed on the genuine orthogonal V98 blocks.
It does not assert a joint sharp \(AB\)-/\(BC\)-Racah measurement and
does not turn the original outer label set into a rectangle.
\end{firewall}

\section{Exact physical application to the off-cut covariance}
\label{sec:V118-offcut-application}

Let \(\mathbf D_{ab}\) be the old signed same-pair multiplier in the
V113 noncomparable block and let \(\mathbf D_{bc}^{0}\) be the peeled
off-cut covariance multiplier.  Their coordinate banks are disjoint by
the V115 factorization.  Let
\[
 \mathbf A_{ab}=\sum_\lambda a_\lambda P_\lambda^{ab},
 \qquad
 \mathbf A_{bc}^{0}=\sum_\mu b_\mu P_\mu^{bc}
\]
be their V112 positive majorants.

On the reduced probability space, use the exact V110 hybrid split
\begin{equation}
 \mathbf R_{\rm bin}^{1}
 =
 \mathbf R_{ab}^{1,\rhd}
 +
 \mathbf R_{ac}^{1,\lhd}.
 \label{eq:V118-reduced-hybrid-split}
\end{equation}
It obeys
\begin{equation}
 \|\mathbf R_{ar}^{1}\|_{\rm op}
 \le Cs_\alpha H_{ar}^{1},
 \qquad r\in\{b,c\},
 \label{eq:V118-reduced-row-bounds}
\end{equation}
where
\(\mathbf R_{ab}^{1}:=\mathbf R_{ab}^{1,\rhd}\) and
\(\mathbf R_{ac}^{1}:=\mathbf R_{ac}^{1,\lhd}\).

\begin{lemma}[Each complete outer block refines both pair outputs]
\label{lem:V118-physical-joint-refinement}
After resolving the exact same-label outputs of the old \(ab\) bank and
the off-cut \(bc\) bank, every compatible V98 covariance outer label
\(j\) has a pair \((\lambda(j),\mu(j))\), and
\begin{equation}
 \sum_{j:\lambda(j)=\lambda,\,\mu(j)=\mu}P_j
 \preccurlyeq
 P_\lambda^{ab}\otimes P_\mu^{bc}\otimes I_{\rm rem}.
 \label{eq:V118-physical-joint-refinement}
\end{equation}
\end{lemma}

\begin{proof}
The V98 outer label contains every already-regrouped same-label
puncture output and every genuine orthogonal copy label.  The old
\(ab\) and off-cut \(bc\) banks occupy disjoint coordinate tensor
factors, so their output projections commute and their common
refinement is
\(P_\lambda^{ab}\otimes P_\mu^{bc}\).  Fixing all remaining output and
copy labels gives mutually orthogonal subprojections of that common
pair projection.  Summing any subfamily with the same
\((\lambda,\mu)\) is therefore bounded by the full common projection
tensored with the remaining identity.
\end{proof}

\begin{theorem}[Physical joint completion of each hybrid row]
\label{thm:V118-physical-joint-completion}
For \(r\in\{b,c\}\), let \(\mathbf L_r\) be the physical signed outer
family obtained by inserting
\(\mathbf R_{ar}^{1}\) from
\eqref{eq:V118-reduced-hybrid-split} into the off-cut covariance
branch, and put
\[
 \mathbf Q_r
 :=
 \sum_j|V_jL_{r,j}V_j^*|.
\]
Then
\begin{align}
 \mathbf Q_r
 &\preccurlyeq
 Cs_\alpha H_{ar}^{1}
 (\mathbf A_{ab}\otimes\mathbf A_{bc}^{0}\otimes I_{\rm rem}),
 \label{eq:V118-physical-row-order}\\
 \boxed{
 \|\mathbf Q_r\|_{\rm rows,prod}
 \le
 Cs_\alpha H_{ar}^{1}\,
 \mathcal K_\star^{ab,bc},}
 \label{eq:V118-physical-row-capacity}\\
 \mathcal K_\star^{ab,bc}
 &:=
 \|\mathbf A_{ab}\otimes\mathbf A_{bc}^{0}\|_{\rm 3prod}.
 \label{eq:V118-star-capacity}
\end{align}
All correlations between the two pair outputs and the remaining outer
labels are allowed.
\end{theorem}

\begin{proof}
On a complete outer block, exact same-label regrouping and the
independent-coordinate factorization give
\[
 L_{r,j}
 =
 d_{\lambda(j)}^{ab}d_{\mu(j)}^{bc}R_{r,j}.
\]
The individual hybrid estimate
\eqref{eq:V118-reduced-row-bounds} gives
\(\|R_{r,j}\|_{\rm op}\le Cs_\alpha H_{ar}^{1}\), uniformly in \(j\).
Apply
\Cref{cor:V118-block-factorization} and then
\Cref{thm:V118-joint-positive-completion}, using
\Cref{lem:V118-physical-joint-refinement}.
\end{proof}

\begin{theorem}[Coefficient-level recombination without a false
absolute-value order]
\label{thm:V118-coefficient-recombination}
For arbitrary row-factorized trace-class input coefficients, the
complete peeled covariance contribution is bounded by
\begin{equation}
 \boxed{
 C s_\alpha\mathcal K_\star^{ab,bc}
 (H_{ab}^{1}+H_{ac}^{1})
 \prod_v\|B_v\|_1.}
 \label{eq:V118-complete-covariance-response}
\end{equation}
\end{theorem}

\begin{proof}
Equation \eqref{eq:V118-reduced-hybrid-split} holds before physical
compression, so linearity gives the sum of the \(r=b\) and \(r=c\)
output coefficients in every outer block.  Apply the trace-norm
triangle to those two coefficients.  For each signed family, use the
V98 physical matrix tax with its genuine absolute carrier
\(\mathbf Q_r\), then insert
\eqref{eq:V118-physical-row-capacity} and sum \(r=b,c\).
No operator inequality comparing
\(|\mathbf L_b+\mathbf L_c|\) with
\(|\mathbf L_b|+|\mathbf L_c|\) is used.
\end{proof}

\begin{corollary}[The V117 physical domination gate is achieved
rowwise]
\label{cor:V118-V117-gate-achieved}
For each exact hybrid row, the positive-domination alternative in
\eqref{eq:V117-physical-factorization} holds with
\[
 Q_e=\mathbf A_{ab},
 \qquad
 Q_f=\mathbf A_{bc}^{0},
 \qquad
 M=Cs_\alpha H_{ar}^{1}I_{\rm rem}.
\]
Thus the unresolved analytic quantity is precisely
\(\mathcal K_\star^{ab,bc}\), not rectangularity of the V98 labels.
\end{corollary}

\begin{proof}
This is the operator order
\eqref{eq:V118-physical-row-order}.  Quotient the core identity as a
within-row spectator.
\end{proof}

\section{Exact bounds for the remaining star capacity}
\label{sec:V118-star-capacity}

Let
\[
 A_{ab,*}:=\|\mathbf A_{ab}\|_{\rm op},
 \qquad
 A_{bc,*}:=\|\mathbf A_{bc}^{0}\|_{\rm op},
\]
and write their ordinary pair product capacities as
\[
 \kappa_{ab}:=\|\mathbf A_{ab}\|_{\rm prod},
 \qquad
 \kappa_{bc}^{0}:=\|\mathbf A_{bc}^{0}\|_{\rm prod}.
\]

\begin{lemma}[Valid one-sided shared-row bounds]
\label{lem:V118-one-sided-star}
Even though the two pair banks share original row \(b\),
\begin{equation}
 \boxed{
 \mathcal K_\star^{ab,bc}
 \le
 \min\{
 A_{ab,*}\kappa_{bc}^{0},
 A_{bc,*}\kappa_{ab}
 \}.}
 \label{eq:V118-one-sided-star}
\end{equation}
\end{lemma}

\begin{proof}
Positive order gives
\[
 \mathbf A_{ab}\otimes\mathbf A_{bc}^{0}
 \preccurlyeq
 A_{ab,*}(I_{ab}\otimes\mathbf A_{bc}^{0}).
\]
In a three-original-row product state, tracing the \(ab\)-bank factors
leaves a product of the reduced state on the \(bc\)-part of row \(b\)
and the state on row \(c\).  The first may be mixed, but resolving it
into pure states bounds the expectation by \(\kappa_{bc}^{0}\).
This proves the first bound.  Interchange the two banks for the second
and take their minimum.
\end{proof}

\begin{corollary}[Recovery of the two single-pair estimates]
\label{cor:V118-single-pair-recovery}
If the positive multipliers have uniform operator caps, then
\begin{align}
 \mathcal K_\star^{ab,bc}
 &\le C\kappa_{bc}^{0}
 \le
 C_\delta\frac{h_{bc}^{0}}{s_\alpha\Xi_c^2},
 \label{eq:V118-recover-bc}\\
 \mathcal K_\star^{ab,bc}
 &\le C\kappa_{ab}.
 \label{eq:V118-recover-ab}
\end{align}
The first line recovers the V116 covariance estimate from
\eqref{eq:V118-complete-covariance-response}.  The second retains the
old same-pair orientation.  Neither line multiplies the two marginal
capacities.
\end{corollary}

\begin{proof}
Use \Cref{lem:V118-one-sided-star}; the selected-heavy bound in the
first line is
\eqref{eq:V116-offcut-pair-capacity}.  Insert
\eqref{eq:V118-recover-bc} in
\eqref{eq:V118-complete-covariance-response} to obtain
\eqref{eq:V116-covariance-absolute} at coefficient level.
\end{proof}

Resolve the spectra of \(\mathbf A_{ab}\) and
\(\mathbf A_{bc}^{0}\) as
\[
 a_1\ge a_2\ge\cdots\ge0,
 \qquad
 b_1\ge b_2\ge\cdots\ge0.
\]
Use the actual three-row labels of
\eqref{eq:V117-three-row-spaces}, with \(e=ab\) and \(f=bc\), to
calculate \(d_\star\).

\begin{theorem}[Combined cold-spectrum bound for the physical
multiplier]
\label{thm:V118-physical-cold-spectrum}
Suppose
\[
 a_i\le A_e,\qquad b_j\le A_f,
\]
and outside cold sets of cardinality \(r_e,r_f\),
\[
 a_i\le h_e,\qquad b_j\le h_f.
\]
Then
\begin{equation}
 \boxed{
 \mathcal K_\star^{ab,bc}
 \le
 H_{ef}
 +
 A_eA_f
 \min\left\{1,\frac{r_er_f}{d_\star}\right\},}
 \label{eq:V118-physical-cold-spectrum}
\end{equation}
where \(H_{ef}=\max\{h_eA_f,A_eh_f\}\).
\end{theorem}

\begin{proof}
After quotienting the remaining identity spectators, take the core in
\Cref{cor:V117-spectral-gate} to be one-dimensional:
\[
 M=1,\qquad
 \Pi_3(M)=B_3(M)=\eta_3(M)=1.
\]
Then \eqref{eq:V117-spectral-gate} is exactly the displayed estimate.
\end{proof}

\begin{corollary}[A fully explicit sufficient physical decay
condition]
\label{cor:V118-physical-decay}
If \(A_e,A_f\) are uniformly bounded and
\begin{equation}
 H_{ef}\longrightarrow0,
 \qquad
 \frac{r_er_f}{d_\star}\longrightarrow0,
 \label{eq:V118-physical-decay}
\end{equation}
then
\[
 \mathcal K_\star^{ab,bc}\longrightarrow0.
\]
Consequently the physical covariance coefficient in
\eqref{eq:V118-complete-covariance-response}, normalized by
\(s_\alpha(H_{ab}^{1}+H_{ac}^{1})\), tends to zero.
\end{corollary}

\begin{proof}
Apply \Cref{thm:V118-physical-cold-spectrum} and then
\eqref{eq:V118-complete-covariance-response}.
\end{proof}

\begin{assessment}[What has and has not been paid]
\label{ass:V118-payment-status}
The physical-interface part of the V117 gate is now paid: arbitrary
correlation of complete outer labels is absorbed by
\Cref{thm:V118-joint-positive-completion}, and the V110 row amplitudes
are kept exact until a valid coefficient-level triangle is taken.
The remaining issue is quantitative.  One must derive
\eqref{eq:V118-physical-decay}, or the stronger graph-weighted analogue
needed at finite scale, from the actual V53 spectra of
\(\mathbf A_{ab}\) and \(\mathbf A_{bc}^{0}\).  The separately known
capacities in \eqref{eq:V118-one-sided-star} do not determine the
product-spectrum cold rectangle.
\end{assessment}

\section{Regression checks}
\label{sec:V118-regressions}

\begin{proposition}[Double-singlet equality survives physical
completion]
\label{prop:V118-double-singlet-physical}
If the two positive multipliers consist of single dual-singlet blocks
\[
 \mathbf A_{ab}=\alpha P_d,
 \qquad
 \mathbf A_{bc}^{0}=\beta P_m,
\]
then
\begin{equation}
 \mathcal K_\star^{ab,bc}
 =
 \frac{\alpha\beta}{dm}.
 \label{eq:V118-double-singlet-physical}
\end{equation}
The joint-completion upper bound is sharp when the physical outer
family contains that complete block and its core bound is attained.
\end{proposition}

\begin{proof}
Scalar multiplication of
\eqref{eq:V117-double-singlet} gives the identity.  If the one block
saturates \eqref{eq:V118-joint-block-majorant}, every inequality in the
completion proof is equality on its maximizing product state.
\end{proof}

\begin{proposition}[Alternating orientation remains visible, not
hidden]
\label{prop:V118-alternating-visible}
For the two abstract dimension arrays
\eqref{eq:V117-middle-pattern} and
\eqref{eq:V117-outward-pattern}, the cold-rank term in
\eqref{eq:V118-physical-cold-spectrum} is respectively
\begin{equation}
 \min\left\{1,\frac{r_er_f}{L^2}\right\},
 \qquad
 \min\left\{1,\frac{r_er_f}{L}\right\}.
 \label{eq:V118-orientation-comparison}
\end{equation}
Thus the common physical completion introduces no artificial
coordinatewise switch, but it also does not erase the genuine factor
\(\chi_\star=L\) in the outward pattern.
\end{proposition}

\begin{proof}
Insert \(d_\star=L^2\) and \(d_\star=L\), respectively, from
\Cref{prop:V117-mass-not-orientation} into
\eqref{eq:V118-physical-cold-spectrum}.
\end{proof}

\section{V118 ledger and V119 mandate}
\label{sec:V118-ledger}

\begin{theorem}[Integrated V118 statement]
\label{thm:V118-integrated}
Preserving V1--V117, V118 proves:
\begin{enumerate}[label=\textup{(V118.\arabic*)},leftmargin=6.7em]
\item joint positive completion of two correlated pair labels,
      \eqref{eq:V118-joint-completed-order}--%
      \eqref{eq:V118-joint-completed-capacity};
\item the exact common refinement of the physical \(ab\) and \(bc\)
      output projections,
      \eqref{eq:V118-physical-joint-refinement};
\item the physical positive domination of each V110 hybrid row,
      \eqref{eq:V118-physical-row-order};
\item the single star-capacity response
      \eqref{eq:V118-complete-covariance-response};
\item achievement of the V117 physical domination gate without a
      rectangular outer-label hypothesis;
\item the valid one-sided bounds
      \eqref{eq:V118-one-sided-star};
\item the combined physical cold-spectrum estimate
      \eqref{eq:V118-physical-cold-spectrum}; and
\item the sharp double-singlet and orientation regressions
      \eqref{eq:V118-double-singlet-physical} and
      \eqref{eq:V118-orientation-comparison}.
\end{enumerate}
\end{theorem}

\begin{proof}
The items are
\Cref{thm:V118-joint-positive-completion,
lem:V118-physical-joint-refinement,
thm:V118-physical-joint-completion,
thm:V118-coefficient-recombination,
cor:V118-V117-gate-achieved,
lem:V118-one-sided-star,
thm:V118-physical-cold-spectrum,
prop:V118-double-singlet-physical,
prop:V118-alternating-visible}.
\end{proof}

\begin{openproblem}[V119 mandate: evaluate the physical star spectrum]
\label{op:V119-mandate}
\begin{enumerate}
\item Insert the exact V52/V53 eigenvalue formulas for
      \(\mathbf A_{ab}\) and \(\mathbf A_{bc}^{0}\) into
      \eqref{eq:V118-physical-cold-spectrum}.
\item Define the cold sets with one fixed threshold before taking the
      product spectrum; keep all multiplicities.
\item Compute \(r_e,r_f,d_\star\) from the actual resolved
      \(SU(3)\) representation labels, not from scalar masses.
\item Compare the cold--cold term with the exact marked-tree monomial,
      listing which collision and pair weights are already contained
      in the two multipliers so that none is counted twice.
\item If \eqref{eq:V118-physical-decay} fails, identify an explicit
      persistent family of pair labels and test whether it is Cartan,
      dual-singlet, or a new multiplicity family.
\item Keep the maximal-\(c\) V116 payment separate; V119 concerns only
      the residual order \eqref{eq:V116-residual-order}.
\end{enumerate}
\end{openproblem}

\section{Firewall after V118}
\label{sec:V118-firewall}

\begin{firewall}[Current claim boundary]
\label{fire:V118-current-boundary}
V118 proves the physical joint positive domination and reduces the
residual off-cut covariance coefficient to one explicit shared-vertex
star capacity.  It does not prove the physical spectral decay
\eqref{eq:V118-physical-decay}, compare that capacity with every exact
marked-tree weight, close the residual ordering, prove JREC, prove
quartic or all-order MTA, construct continuum Yang--Mills theory, or
prove a positive spectral gap.  The full Yang--Mills existence and
mass gap remain open.
\end{firewall}

\section{Append-only record for V118}
\label{sec:V118-record}

V118 was appended to the complete V117 manuscript.  The exact V117
parent had \(78948\) newline-delimited source records, \(2750601\)
bytes, and SHA--256 digest
\[
 \texttt{69f178c257cf7a5eea9edbd67104049165f15c661029a1df96d21fc66e6540ae}.
\]
No V117 mathematical content was changed.  The sole final
\verb|\end{document}| is restored below.

% =============================================================================
% END APPENDED VERSION V118 MATERIAL
% =============================================================================

% =============================================================================
% BEGIN APPENDED VERSION V119 MATERIAL
% =============================================================================

\chapter{V119: Residual Off-Cut Asymmetry and the One Remaining
Effective-Scale Gate}
\label{ch:V119}

\section{Purpose}
\label{sec:V119-purpose}

V118 reduced the physical shared-vertex covariance carrier to
\[
 \mathcal K_\star^{ab,bc}
 =
 \|\mathbf A_{ab}\otimes\mathbf A_{bc}^{0}\|_{\rm 3prod}.
\]
A direct use of two generic cold-rank estimates would retain the star
orientation factor \(\chi_\star\).  The residual V116 mass order contains
more information than that generic argument uses.  The off-cut bank
carries a fixed fraction of all row-\(c\) mass, whereas the \emph{entire}
row-\(b\) mass is smaller by \(C_{\rm row}^{-1}\).  Hence a fixed fraction
of the off-cut bank is locally asymmetric in precisely the reverse-fusion
sense of V53.  Its effective-cold output family is empty.

This removes the common cold--cold rectangle and every orientation factor
from the residual branch.  The joint star capacity becomes the old
\(ab\)-pair capacity times the operator hotness of the off-cut \(bc\)
multiplier.  The \(H_{ab}^{1}\) hybrid row is then paid outright.  The
sole remaining \(H_{ac}^{1}\) row is paid if and only if one explicit
effective-scale factor supplies \(\Xi_b/\Xi_c\).  V119 records that factor
without deriving it from pair heaviness.

\section{The off-cut effective number}
\label{sec:V119-effective-number}

On the off-cut bank \(O_{bc}\), put
\begin{equation}
 r_f:=A_{c,f},
 \qquad
 q_f:=A_{b,f},
 \qquad f\in O_{bc},
 \label{eq:V119-local-masses}
\end{equation}
and define
\begin{align}
 x_0&:=\sum_{f\in O_{bc}}r_f=D_c^{\rm off},
 &
 y_0&:=\sum_{f\in O_{bc}}q_f,
 \notag\\
 h_0&:=\sum_{f\in O_{bc}}r_fq_f,
 &
 N_0&:=\frac{x_0^2}{h_0}
 \quad(h_0>0).
 \label{eq:V119-offcut-data}
\end{align}
If \(h_0=0\), then the off-cut covariance multiplier vanishes and the
covariance branch is absent.  We therefore work with \(h_0>0\).

\begin{lemma}[Basic off-cut mass geometry]
\label{lem:V119-basic-geometry}
On the heavy off-cut branch,
\begin{equation}
 x_0\ge\eta_0\Xi_c,
 \qquad
 y_0\le\Xi_b,
 \qquad
 h_0\le x_0y_0,
 \qquad
 N_0\ge\eta_0\frac{\Xi_c}{\Xi_b},
 \label{eq:V119-basic-geometry}
\end{equation}
where \(\eta_0:=1-7\delta>0\).
\end{lemma}

\begin{proof}
The first inequality is \eqref{eq:V113-heavy-off-cut}.  The off-cut
bank is a subset of the row-\(b\) coordinates, giving the second.
Nonnegativity gives
\[
 \sum_fr_fq_f
 \le
 \left(\sum_fr_f\right)\left(\sum_fq_f\right)
 =x_0y_0.
\]
Thus \(N_0=x_0^2/h_0\ge x_0/y_0\), and the first two inequalities give
the final bound.
\end{proof}

\section{The residual order forces reverse-fusion asymmetry}
\label{sec:V119-asymmetry}

Let \(0<\eta_{\rm rf}<1\) be the fixed reverse-fusion constant used in
\eqref{eq:V53-asymmetric-coordinates}, and define
\begin{equation}
 \mathcal A_0
 :=
 \{f\in O_{bc}:q_f\le\eta_{\rm rf}r_f\}.
 \label{eq:V119-asymmetric-set}
\end{equation}

\begin{lemma}[A fixed asymmetric fraction]
\label{lem:V119-asymmetric-fraction}
Choose the programme constant \(C_{\rm row}\) so that
\begin{equation}
 C_{\rm row}\ge\frac{2}{\eta_{\rm rf}\eta_0}.
 \label{eq:V119-row-threshold}
\end{equation}
Then the residual order
\(\Xi_c>C_{\rm row}\Xi_b\) implies
\begin{equation}
 \sum_{f\in\mathcal A_0}r_f
 \ge\frac12x_0.
 \label{eq:V119-asymmetric-half}
\end{equation}
\end{lemma}

\begin{proof}
On the complement of \(\mathcal A_0\),
\(q_f>\eta_{\rm rf}r_f\).  Therefore
\[
 \sum_{f\notin\mathcal A_0}r_f
 <
 \eta_{\rm rf}^{-1}\sum_{f\notin\mathcal A_0}q_f
 \le
 \eta_{\rm rf}^{-1}y_0
 \le
 \frac{\Xi_b}{\eta_{\rm rf}}
 <
 \frac{\Xi_c}{\eta_{\rm rf}C_{\rm row}}
 \le
 \frac{\eta_0\Xi_c}{2}
 \le\frac{x_0}{2}.
\]
Subtract from \(x_0\).
\end{proof}

\begin{theorem}[The off-cut effective-cold family is empty]
\label{thm:V119-cold-family-empty}
There is a fixed \(\kappa_0>0\), depending only on the programme
constants, such that no resolved output \(S\) of the complete off-cut
\(bc\) pair bank satisfies
\begin{equation}
 \Xi(S)<\kappa_0N_0.
 \label{eq:V119-no-cold-output}
\end{equation}
Equivalently, the V53 effective-cold rank of
\(\mathbf A_{bc}^{0}\) is zero.
\end{theorem}

\begin{proof}
Apply \Cref{thm:V53-asymmetric-star} to the off-cut pair bank with
selected endpoint \(c\), selected mass \(x=x_0\), and
\(\beta=1/2\).  Its asymmetric-mass hypothesis is exactly
\eqref{eq:V119-asymmetric-half}.  The proof of that theorem uses
reverse fusion coordinate by coordinate and shows that every admissible
output has mass at least \(c x_0\ge cN_0\).  Choose
\(\kappa_0<c\).
\end{proof}

\begin{corollary}[No dual-singlet cold obstruction]
\label{cor:V119-no-singlet-cold}
The complete off-cut pair output cannot be a dual singlet in the residual
branch.
\end{corollary}

\begin{proof}
A trivial complete output has mass zero, contradicting
\eqref{eq:V119-no-cold-output}.  Equivalently, on every coordinate
producing a trivial output the two inputs are dual and have equal
Casimir mass, whereas
\eqref{eq:V119-asymmetric-half} supplies a nonzero input mass on
coordinates with \(q_f<r_f\).
\end{proof}

\begin{assessment}[Alternating local heights do not restore a cold
rectangle]
\label{ass:V119-alternating}
The proof of \Cref{lem:V119-asymmetric-fraction} uses only the two exact
bank sums \(x_0,y_0\).  It is unchanged if the larger local height
alternates among coordinates.  Reverse fusion is then applied on the
single fixed set \(\mathcal A_0\), not after a coordinatewise choice of
retained-row orientation.  Thus the V60 alternating-height regression
cannot reintroduce a \(bc\) cold set in this residual mass order.
\end{assessment}

\section{Operator hotness of the off-cut pair multiplier}
\label{sec:V119-operator-hotness}

\begin{theorem}[Uniform off-cut operator-hot bound]
\label{thm:V119-offcut-operator-hot}
The positive off-cut pair multiplier obeys
\begin{equation}
 \boxed{
 \|\mathbf A_{bc}^{0}\|_{\rm op}
 \le
 C\vartheta_0,
 \qquad
 \vartheta_0:=
 \min\left\{1,\frac1{s_\alpha N_0}\right\}.}
 \label{eq:V119-offcut-operator-hot}
\end{equation}
\end{theorem}

\begin{proof}
The uniform coefficient cap
\eqref{eq:V52-defect-weight-cap} gives the first entry of the minimum.
If \(s_\alpha N_0>M\), every output satisfies
\(\Xi(S)\ge\kappa_0N_0\) by
\Cref{thm:V119-cold-family-empty}.  The output-hot calculation in the
proof of \Cref{thm:V53-effective-cold-reduction} then gives, for every
eigenvalue \(a(S)\) of the positive multiplier,
\[
 a(S)\le\frac{C}{s_\alpha N_0}.
\]
Take the maximum.  Enlarging \(C\) covers the bounded
\(s_\alpha N_0\) regime.
\end{proof}

\begin{corollary}[Orientation-free collapse of the star capacity]
\label{cor:V119-star-collapse}
The physical star capacity satisfies
\begin{equation}
 \boxed{
 \mathcal K_\star^{ab,bc}
 \le
 C\vartheta_0\,\kappa_{ab}.}
 \label{eq:V119-star-collapse}
\end{equation}
In particular, neither \(d_\star\) nor \(\chi_\star\) occurs.
\end{corollary}

\begin{proof}
Use the second one-sided inequality in
\eqref{eq:V118-one-sided-star} and then
\eqref{eq:V119-offcut-operator-hot}.
\end{proof}

\section{Disposition of the two exact hybrid rows}
\label{sec:V119-row-disposition}

Let \(h_{ab}\) denote the pair mass of the old bi-heavy \(ab\)
multiplier.  Its positive capacity satisfies
\begin{equation}
 \kappa_{ab}
 \le
 C\frac{h_{ab}}{s_\alpha\Xi_a\Xi_b}
 \label{eq:V119-old-pair-capacity}
\end{equation}
by the V109 geometric pair tax and V112 positive majorant.

\begin{theorem}[The \(ab\) hybrid row is paid; the \(ac\) row has one
scalar gate]
\label{thm:V119-two-row-disposition}
The two coefficient responses of
\eqref{eq:V118-reduced-hybrid-split} satisfy
\begin{align}
 \mathcal R_{ab}^{\rm off,cov}
 &\le
 C\frac{h_{ab}H_{ab}^{1}}{\Xi_a\Xi_b}
 \prod_v\|B_v\|_1,
 \label{eq:V119-ab-row-paid}\\
 \mathcal R_{ac}^{\rm off,cov}
 &\le
 C\vartheta_0
 \frac{h_{ab}H_{ac}^{1}}{\Xi_a\Xi_b}
 \prod_v\|B_v\|_1.
 \label{eq:V119-ac-row-gate}
\end{align}
Consequently the \(ac\) row has its marked denominator
\(\Xi_a\Xi_c\) whenever
\begin{equation}
 \boxed{
 \vartheta_0
 \le
 C\frac{\Xi_b}{\Xi_c}.}
 \label{eq:V119-scalar-gate}
\end{equation}
\end{theorem}

\begin{proof}
For each row, combine
\eqref{eq:V118-physical-row-capacity},
\eqref{eq:V119-star-collapse}, and
\eqref{eq:V119-old-pair-capacity}.  The factor \(s_\alpha\) in
\eqref{eq:V118-physical-row-capacity} cancels the one in the old pair
capacity.  This gives
\[
 \mathcal R_{ar}^{\rm off,cov}
 \le
 C\vartheta_0
 \frac{h_{ab}H_{ar}^{1}}{\Xi_a\Xi_b}
 \prod_v\|B_v\|_1.
\]
For \(r=b\), use \(\vartheta_0\le1\), proving
\eqref{eq:V119-ab-row-paid}.  For \(r=c\), insert
\eqref{eq:V119-scalar-gate}.
\end{proof}

\begin{corollary}[Exact effective-scale form of the last gate]
\label{cor:V119-effective-scale-gate}
A sufficient form of \eqref{eq:V119-scalar-gate} is
\begin{equation}
 \boxed{
 s_\alpha N_0
 \ge
 c\frac{\Xi_c}{\Xi_b}.}
 \label{eq:V119-effective-scale-gate}
\end{equation}
Equivalently,
\begin{equation}
 h_0
 \le
 C s_\alpha x_0^2\frac{\Xi_b}{\Xi_c}.
 \label{eq:V119-overlap-scale-gate}
\end{equation}
\end{corollary}

\begin{proof}
Under \eqref{eq:V119-effective-scale-gate}, the second entry in
\(\vartheta_0\) is at most \(C\Xi_b/\Xi_c\); since
\(\Xi_b/\Xi_c<1\), it is the relevant entry after constants are
adjusted.  Substitute \(N_0=x_0^2/h_0\) and rearrange for the equivalent
form.
\end{proof}

\begin{remark}[What basic mass geometry misses]
\label{rem:V119-missing-heat-factor}
\Cref{lem:V119-basic-geometry} gives only
\[
 N_0\ge\eta_0\frac{\Xi_c}{\Xi_b}.
\]
It does not imply
\eqref{eq:V119-effective-scale-gate}, because that condition contains
the additional factor \(s_\alpha\).  Equivalently, the elementary bound
\(h_0\le x_0y_0\) lacks the factor \(s_\alpha\) required by
\eqref{eq:V119-overlap-scale-gate}.
\end{remark}

\section{Sharpness and direction}
\label{sec:V119-sharpness}

\begin{proposition}[A concentrated bank defeats a mass-only proof of
the last gate]
\label{prop:V119-concentrated-regression}
Fix any proposed constant \(C\) in
\eqref{eq:V119-scalar-gate}.  There are nonnegative one-coordinate
bank data satisfying all inequalities in
\Cref{lem:V119-basic-geometry} and
\(\Xi_c>C_{\rm row}\Xi_b\), but for which
\eqref{eq:V119-scalar-gate} fails.
\end{proposition}

\begin{proof}
Choose \(R>\max\{C_{\rm row},2C\}\), take positive \(Y\), and set
\[
 \Xi_b=Y,
 \qquad
 \Xi_c=RY,
 \qquad
 x_0=RY,
 \qquad
 y_0=Y.
\]
Put the two masses on one coordinate.  Then
\[
 h_0=R Y^2=x_0y_0,
 \qquad
 N_0=\frac{x_0^2}{h_0}=R.
\]
Choose \(s_\alpha R<1\).  It follows that
\[
 \vartheta_0=1,
 \qquad
 C\frac{\Xi_b}{\Xi_c}=\frac CR<\frac12.
\]
Thus \eqref{eq:V119-scalar-gate} fails although the residual order and
every elementary mass inequality hold.
\end{proof}

\begin{firewall}[Scope of the concentrated regression]
\label{fire:V119-concentrated-scope}
\Cref{prop:V119-concentrated-regression} is an information-sharp test
for arguments using only the nonnegative bank masses and the universal
multiplier cap.  It does not assert that an actual Yang--Mills Fourier
coefficient saturates \(\vartheta_0=1\), nor does it contradict the
reverse-fusion deletion of cold outputs.  It proves that the missing
factor must come from the exact hot-output heat symbol, a further
coefficient cancellation, or a different endpoint routing.
\end{firewall}

\begin{corollary}[Fixed positive heat scale would close the scalar gate]
\label{cor:V119-fixed-heat}
If \(s_\alpha\ge s_*>0\), then, after increasing the fixed comparison
constant,
\eqref{eq:V119-scalar-gate} follows from
\Cref{lem:V119-basic-geometry}.  Therefore the obstruction is tied to
uniformity as \(s_\alpha\downarrow0\), not to fixed heat time.
\end{corollary}

\begin{proof}
\Cref{lem:V119-basic-geometry} gives
\[
 \frac1{s_\alpha N_0}
 \le
 \frac1{s_*\eta_0}\frac{\Xi_b}{\Xi_c}.
\]
Use the second entry in the definition of \(\vartheta_0\).
\end{proof}

\begin{assessment}[Upstream target after V119]
\label{ass:V119-upstream-target}
The common-carrier, label-correlation, hybrid-interference,
dual-singlet, cold-rank, and alternating-orientation issues have now
been removed from the residual branch.  The exact open factor is the
small-heat comparison
\[
 \min\left\{1,\frac{h_0}{s_\alpha x_0^2}\right\}
 \frac{\Xi_c}{\Xi_b}.
\]
Repeating the V53 capacity estimate cannot improve it: V53 has already
made the effective-cold family empty and the remaining bound is the
pointwise hot-output estimate.  The next proof must therefore go
upstream to the exact heat difference
\[
 (1+s_\alpha\Xi(S))
 \left|
 q_{\alpha,S}
 -
 q_{\alpha,R}q_{\alpha,Q}
 \right|
\]
on a reverse-fusion asymmetric \(bc\) coordinate, or reroute the
\(ac\) hybrid amplitude before the pair multiplier is replaced by its
operator norm.
\end{assessment}

\section{V119 ledger and mandatory V120 audit}
\label{sec:V119-ledger}

\begin{theorem}[Integrated V119 statement]
\label{thm:V119-integrated}
Preserving V1--V118, V119 proves:
\begin{enumerate}[label=\textup{(V119.\arabic*)},leftmargin=6.7em]
\item the exact off-cut effective-number data
      \eqref{eq:V119-offcut-data} and their mass geometry
      \eqref{eq:V119-basic-geometry};
\item a fixed reverse-fusion asymmetric fraction
      \eqref{eq:V119-asymmetric-half};
\item emptiness of the complete off-cut effective-cold output family;
\item exclusion of a dual-singlet cold obstruction;
\item the operator-hot estimate
      \eqref{eq:V119-offcut-operator-hot};
\item orientation-free collapse of the shared-vertex capacity,
      \eqref{eq:V119-star-collapse};
\item unconditional payment of the \(ab\) hybrid row,
      \eqref{eq:V119-ab-row-paid};
\item reduction of the \(ac\) hybrid row to the sole scalar gate
      \eqref{eq:V119-scalar-gate}, with sufficient effective-scale
      form \eqref{eq:V119-effective-scale-gate}; and
\item the concentrated-bank proof that the missing heat factor cannot
      follow from elementary mass geometry alone.
\end{enumerate}
\end{theorem}

\begin{proof}
The items are
\Cref{lem:V119-basic-geometry,
lem:V119-asymmetric-fraction,
thm:V119-cold-family-empty,
cor:V119-no-singlet-cold,
thm:V119-offcut-operator-hot,
cor:V119-star-collapse,
thm:V119-two-row-disposition,
cor:V119-effective-scale-gate,
prop:V119-concentrated-regression}.
\end{proof}

\begin{openproblem}[V120 mandate: audit, then attack the hot-symbol
scale]
\label{op:V120-mandate}
\begin{enumerate}
\item Perform the mandatory read-only audit of the latest SM and NS
      notes in \texttt{latex\_notes}; record byte counts, hashes, and
      only structurally relevant lessons.
\item On the reverse-fusion asymmetric off-cut bank, retain the exact
      signed heat difference
      \(q_{\alpha,S}-q_{\alpha,R}q_{\alpha,Q}\) rather than its
      universal absolute cap.
\item Optimize the output-hot estimate jointly in
      \(s_\alpha\Xi(S)\), \(s_\alpha x_0\), and the local mass ratio
      \(q_f/r_f\).  Determine whether it supplies the missing factor
      \(s_\alpha\) in \eqref{eq:V119-overlap-scale-gate}.
\item If the concentrated one-coordinate family remains order one,
      insert the exact \(ac\) hybrid replacement before taking the
      off-cut operator norm and test for a signed cancellation.
\item Keep the already-paid \(ab\) hybrid row and maximal-\(c\) V116
      region outside the new estimate.
\item State any surviving scale condition explicitly; do not infer it
      from pair heaviness or from the empty cold family.
\end{enumerate}
\end{openproblem}

\section{Firewall after V119}
\label{sec:V119-firewall}

\begin{firewall}[Current claim boundary]
\label{fire:V119-current-boundary}
V119 removes the cold-spectrum and orientation bottlenecks from the
residual off-cut branch, pays its \(ab\) hybrid row, and reduces its
\(ac\) hybrid row to
\eqref{eq:V119-scalar-gate}.  It does not prove that gate uniformly as
\(s_\alpha\downarrow0\), close the residual off-cut covariance branch,
prove JREC, prove quartic or all-order MTA, construct continuum
Yang--Mills theory, or prove a positive spectral gap.  The full
Yang--Mills existence and mass gap remain open.
\end{firewall}

\section{Append-only record for V119}
\label{sec:V119-record}

V119 was appended to the complete V118 manuscript.  The exact V118
parent had \(79559\) newline-delimited source records, \(2770257\)
bytes, and SHA--256 digest
\[
 \texttt{6b8eb793baebb8de046b12025b8bdf71300dd7bbcd4ed6dd8d84c6d1c8bfe8b1}.
\]
No V118 mathematical content was changed.  The sole final
\verb|\end{document}| is restored below.

% =============================================================================
% END APPENDED VERSION V119 MATERIAL
% =============================================================================

% =============================================================================
% BEGIN APPENDED VERSION V120 MATERIAL
% =============================================================================

\chapter{V120: Companion Audit and the Asymmetric Cartan
Hot-Symbol Barrier}
\label{ch:V120}

\section{Purpose}
\label{sec:V120-purpose}

V119 reduced the residual off-cut covariance problem to the scalar
operator-hot gate
\[
 \vartheta_0\lesssim\frac{\Xi_b}{\Xi_c}.
\]
The V120 mandate was to test whether the exact pair heat difference,
rather than its universal cap, supplies that missing ratio.  V120 first
performs the compulsory companion audit.  It then evaluates the strongest
natural benchmark: the exact heat-semigroup multiplier on an asymmetric
\(SU(3)\) Cartan output.

The benchmark does not close the gate.  When the small-to-large input
Casimir ratio is \(\varepsilon^2\), the weighted Cartan defect is of
order \(\varepsilon\).  The Cartan output has unit product visibility, so
no aggregate Schur tax supplies the second power.  This is not a
counterexample for the Wilson multiplier or for the complete signed V98
coefficient.  It proves that reverse fusion plus exponential heat decay
alone cannot establish the V119 scalar gate.  The programme must keep
the exact \(ac\) hybrid amplitude and the Cartan pair output assembled
before absolute value.

\section{Mandatory V120 companion audit}
\label{sec:V120-audit}

The read-only audit on 5 September 2026 found
\begin{center}
\begin{tabular}{|p{0.50\textwidth}|r|p{0.30\textwidth}|}
\hline
Companion file & Bytes & SHA--256\\
\hline
\texttt{Renewal\_NS\_Stability\_Research\_Notes\_V31.tex}
& \(887537\) &
\texttt{f1121b62238ad5491bb7cf03635ea9934942edf573ed8faaa9ee79e1d38a6696}\\
\hline
\texttt{Renewal\_SM\_Einstein\_Continuum\_Research\_Notes\_V135.tex}
& \(3242135\) &
\texttt{1ef43eb313a226750c7505a4d7f3b505972afa06efb4abeafc3d4164e98742e5}\\
\hline
\end{tabular}
\end{center}
The files contain respectively \(23053\) and \(89158\)
newline-delimited source records under the audit counter.  NS is
byte-identical to the V115 audit and supplies no new analytic estimate.

SM advanced from V127 to V135.  V128--V132 keep the full
parameter-dependent Riccati/Krein object through spectral
differentiation, inverse Laplace transform, and angular averaging; a
frozen zero-parameter symbol gives the wrong heat coefficient.  V133--V134
show that a Lorentzian zero-flux graph cannot be inserted into a Euclidean
heat calculation before its Clifford module and domain are identified.
V135 retains determinant multiplicity through the doubled calculation
and delays the boundary subtraction until the gauge, ghost, fermion,
Higgs, and phase sectors have been assembled.

\begin{assessment}[V120 adjustment after the audit]
\label{ass:V120-audit-adjustment}
The transferable rule is an order-of-operations rule.  For the YM
residual, the analogue of the joint spectral parameter is the complete
resolved pair output \(S\) together with the exact signed hybrid
amplitude.  V120 therefore first evaluates
\[
 (1+s_\alpha\Xi(S))
 |q_{\alpha,S}-q_{\alpha,R}q_{\alpha,Q}|
\]
on a fixed physical output.  It does not average the pair capacity,
replace the output by a marginal heat estimate, or declare a
coefficient cancellation before the V98 signed block is assembled.
No SM or NS analytic estimate is imported.
\end{assessment}

\section{Exact heat benchmark on an asymmetric Cartan channel}
\label{sec:V120-heat-Cartan}

For \(k\ge0\), let \(R_k=(k,0)\) be the \(k\)-th symmetric-power
irreducible of \(SU(3)\).  Its Casimir is
\begin{equation}
 c_k:=C_2(R_k)=\frac{k^2+3k}{3}.
 \label{eq:V120-ray-Casimir}
\end{equation}
The Cartan component of \(R_n\otimes R_m\) is \(R_{n+m}\), with
multiplicity one, and
\begin{equation}
 c_{n+m}=c_n+c_m+\frac{2nm}{3}.
 \label{eq:V120-Cartan-cross}
\end{equation}

Fix \(u>0\), fix an integer \(m\ge1\), and put
\begin{equation}
 \tau_n:=\frac{u}{c_n}.
 \label{eq:V120-transition-time}
\end{equation}
For the exact heat-semigroup multipliers
\(h_{\tau,R}=e^{-\tau C_2(R)}\), define the Cartan pair defect and its
weighted size by
\begin{align}
 \delta_{n,m}^{\mathsf h}
 &:=
 e^{-\tau_nc_{n+m}}
 -
 e^{-\tau_n(c_n+c_m)},
 \label{eq:V120-heat-Cartan-defect}\\
 a_{n,m}^{\mathsf h}
 &:=
 (1+\tau_nc_{n+m})
 |\delta_{n,m}^{\mathsf h}|.
 \label{eq:V120-heat-Cartan-weight}
\end{align}

\begin{theorem}[The asymmetric Cartan defect has only a square-root
ratio]
\label{thm:V120-Cartan-square-root}
For fixed \(u>0\) and \(m\ge1\),
\begin{equation}
 \boxed{
 \lim_{n\to\infty}
 n\,a_{n,m}^{\mathsf h}
 =
 2um(1+u)e^{-u}.}
 \label{eq:V120-Cartan-limit}
\end{equation}
In particular,
\begin{align}
 a_{n,m}^{\mathsf h}
 &\asymp_{u,m}
 \sqrt{\frac{c_m}{c_n}},
 \label{eq:V120-square-root-scale}\\
 \frac{a_{n,m}^{\mathsf h}}{c_m/c_n}
 &\longrightarrow+\infty.
 \label{eq:V120-no-full-ratio}
\end{align}
\end{theorem}

\begin{proof}
Use \eqref{eq:V120-Cartan-cross} to factor
\begin{align*}
 |\delta_{n,m}^{\mathsf h}|
 &=
 e^{-\tau_n(c_n+c_m)}
 \left(
 1-e^{-2\tau_nnm/3}
 \right).
\end{align*}
Since \(c_n=(n^2+3n)/3\),
\begin{align*}
 \tau_n(c_n+c_m)&\longrightarrow u,\\
 \tau_nc_{n+m}&\longrightarrow u,\\
 n\left(1-e^{-2unm/(3c_n)}\right)
 &\longrightarrow 2um.
\end{align*}
Multiplying the three limits proves
\eqref{eq:V120-Cartan-limit}.  Also
\[
 \frac{c_m}{c_n}
 \sim\frac{3c_m}{n^2}.
\]
Thus both \(a_{n,m}^{\mathsf h}\) and
\(\sqrt{c_m/c_n}\) are fixed positive multiples of \(n^{-1}\),
whereas \(c_m/c_n\) is a fixed positive multiple of \(n^{-2}\).
This proves the last two statements.
\end{proof}

\begin{corollary}[The V119 scalar gate fails for the heat benchmark]
\label{cor:V120-heat-gate-fails}
Take a one-coordinate off-cut bank with input-\(c\) representation
\(R_n\), input-\(b\) representation \(R_m\), and Cartan output
\(R_{n+m}\).  Then
\begin{equation}
 \frac{\Xi_b}{\Xi_c}
 \asymp_m\frac{c_m}{c_n},
 \label{eq:V120-endpoint-ratio}
\end{equation}
but its exact heat weighted defect violates every uniform estimate
\begin{equation}
 a_{n,m}^{\mathsf h}
 \le C\frac{\Xi_b}{\Xi_c}.
 \label{eq:V120-false-heat-gate}
\end{equation}
\end{corollary}

\begin{proof}
The programme masses are \(1+C_2\), so
\eqref{eq:V120-endpoint-ratio} follows from
\eqref{eq:V120-ray-Casimir} with fixed \(m\) and \(n\to\infty\).
Then \eqref{eq:V120-no-full-ratio} rules out
\eqref{eq:V120-false-heat-gate}.
\end{proof}

\section{Product visibility of the obstructing output}
\label{sec:V120-product-visibility}

\begin{proposition}[No Schur tax repairs the Cartan square-root loss]
\label{prop:V120-Cartan-visible}
Let \(P_{n+m}\) be the Cartan projection in
\(R_n\otimes R_m\).  Then
\begin{equation}
 \|P_{n+m}\|_{\rm prod}=1.
 \label{eq:V120-Cartan-product-one}
\end{equation}
For the positive heat pair multiplier
\[
 \mathbf A_{n,m}^{\mathsf h}
 =
 \sum_{S\subset R_n\otimes R_m}
 a^{\mathsf h}(S)P_S,
\]
one consequently has
\begin{equation}
 \|\mathbf A_{n,m}^{\mathsf h}\|_{\rm prod}
 \ge a_{n,m}^{\mathsf h}.
 \label{eq:V120-Cartan-capacity-lower}
\end{equation}
\end{proposition}

\begin{proof}
The tensor product of normalized highest-weight vectors of \(R_n\)
and \(R_m\) lies entirely in the Cartan component \(R_{n+m}\).  It is
a product vector and attains expectation one in \(P_{n+m}\), proving
\eqref{eq:V120-Cartan-product-one}.  Evaluate the positive multiplier
on that same vector.  Orthogonality of the isotypic projections leaves
only its Cartan eigenvalue
\(a_{n,m}^{\mathsf h}\), proving
\eqref{eq:V120-Cartan-capacity-lower}.
\end{proof}

\begin{theorem}[Heat decay plus aggregate Schur compression is
insufficient]
\label{thm:V120-heat-Schur-insufficient}
No argument whose only inputs on the off-cut pair are:
\begin{enumerate}[label=\textup{(\roman*)}]
\item the exact heat-semigroup multiplier;
\item reverse-fusion output hotness; and
\item aggregate invariant-projector Schur compression
\end{enumerate}
can prove the full ratio in
\eqref{eq:V119-scalar-gate} uniformly over asymmetric \(SU(3)\)
inputs.
\end{theorem}

\begin{proof}
The family in
\Cref{cor:V120-heat-gate-fails} is output-hot: its output mass is
asymptotic to the large input mass.  Its multiplier is the exact heat
multiplier, not an upper majorant.  Yet
\Cref{thm:V120-Cartan-square-root} gives only the square-root ratio.
Finally \Cref{prop:V120-Cartan-visible} shows that aggregate Schur
compression has norm one on the responsible output.  Thus all three
listed inputs hold while the desired full-ratio conclusion fails.
\end{proof}

\begin{firewall}[Benchmark versus the Wilson law]
\label{fire:V120-heat-not-Wilson}
The Wilson multipliers \(q_{\alpha,R}\) are not known uniformly equal to
\(e^{-s_\alpha C_2(R)}\) on the joint scaling
\(n\to\infty\), \(s_\alpha c_n=u\).  Therefore
\Cref{thm:V120-heat-Schur-insufficient} is not a counterexample to the
Wilson scalar gate.  Conversely, fundamental calibration and
fixed-representation small-\(s_\alpha\) expansion do not justify
discarding the heat obstruction.  Any Wilson improvement would require
a new growing-representation theorem and must still be inserted into the
complete signed V98 block.
\end{firewall}

\section{Exact rowwise recombination of the peeled branches}
\label{sec:V120-rowwise-recombination}

Return to the reduced product space of V115 and abbreviate
\[
 A:=F_a,\qquad B:=B_1,\qquad C:=C_1,
 \qquad
 M_0:=b_0c_0,\qquad
 \Delta_0:=\Delta_{bc}^{0}.
\]
Thus
\[
 q_0:=M_0+\Delta_0=\mathbb E_0(B_0C_0)
\]
is the unresolved joint off-cut mean.  Fix the hybrid order of V110 and
define the two plain-product increments
\begin{align}
 V_{b,e}
 &:=
 \mathbb E[(B-B^{(e)})C\mid\mathcal F_e],
 \label{eq:V120-plain-b-increment}\\
 V_{c,e}
 &:=
 \mathbb E[B^{(e)}(C-C^{(e)})\mid\mathcal F_e].
 \label{eq:V120-plain-c-increment}
\end{align}
Put
\begin{align}
 \mathbf C_b
 &:=
 \sum_e\mathbb E[(D_eA)V_{b,e}],
 &
 \mathbf C_c
 &:=
 \sum_e\mathbb E[(D_eA)V_{c,e}].
 \label{eq:V120-plain-row-pieces}
\end{align}

\begin{lemma}[Plain and centred row pieces]
\label{lem:V120-plain-centred}
The composite covariance and the V110 centred row pieces obey
\begin{align}
 \operatorname{Cov}(A,BC)
 &=
 \mathbf C_b+\mathbf C_c,
 \label{eq:V120-plain-row-split}\\
 \mathbf K_{ab}^{\rhd}
 &=
 \mathbf C_b-\overline C\,\operatorname{Cov}(A,B),
 \label{eq:V120-centred-b}\\
 \mathbf K_{ac}^{\lhd}
 &=
 \mathbf C_c-\overline B\,\operatorname{Cov}(A,C).
 \label{eq:V120-centred-c}
\end{align}
\end{lemma}

\begin{proof}
The pointwise hybrid identity
\[
 BC-B^{(e)}C^{(e)}
 =
 (B-B^{(e)})C+B^{(e)}(C-C^{(e)})
\]
and the martingale covariance formula give
\eqref{eq:V120-plain-row-split}.  In the V110 \(b\)-increment replace
\(C^\circ\) by \(C-\overline C\); the term containing
\(\overline C\) sums to
\(\overline C\,\operatorname{Cov}(A,B)\).  In the \(c\)-increment
replace \(B^{(e)}-\overline B\) by
\(B^{(e)}-\overline B\); the removed constant term sums to
\(\overline B\,\operatorname{Cov}(A,C)\).  These are
\eqref{eq:V120-centred-b} and
\eqref{eq:V120-centred-c}.
\end{proof}

\begin{theorem}[Exact rowwise recombination identity]
\label{thm:V120-rowwise-recombination}
The complete V115 peeled amplitude is
\begin{equation}
 M_0\kappa_3(A,B,C)
 +
 \Delta_0\operatorname{Cov}(A,BC)
 =
 \mathbf T_b+\mathbf T_c,
 \label{eq:V120-recombined-total}
\end{equation}
where
\begin{align}
 \boxed{
 \mathbf T_b
 =
 q_0\mathbf C_b
 -
 M_0\overline C\,\operatorname{Cov}(A,B)
 =
 M_0\mathbf K_{ab}^{\rhd}
 +
 \Delta_0\mathbf C_b,}
 \label{eq:V120-recombined-b}\\
 \boxed{
 \mathbf T_c
 =
 q_0\mathbf C_c
 -
 M_0\overline B\,\operatorname{Cov}(A,C)
 =
 M_0\mathbf K_{ac}^{\lhd}
 +
 \Delta_0\mathbf C_c.}
 \label{eq:V120-recombined-c}
\end{align}
All identities hold before trace norm and final recoupling.
\end{theorem}

\begin{proof}
By \Cref{lem:V110-centred-covariance} and
\Cref{thm:V110-exact-row-allocation},
\[
 \kappa_3(A,B,C)
 =
 \mathbf K_{ab}^{\rhd}+\mathbf K_{ac}^{\lhd}.
\]
Use \eqref{eq:V120-plain-row-split} for the composite covariance,
group the \(b\)- and \(c\)-terms, and substitute
\eqref{eq:V120-centred-b}--\eqref{eq:V120-centred-c}.  Finally use
\(q_0=M_0+\Delta_0\).
\end{proof}

\begin{proposition}[A constant reduced \(b\)-spectator leaves the
Cartan defect]
\label{prop:V120-constant-b}
If \(B=\overline B\) is constant on the reduced probability space, then
\begin{equation}
 \mathbf T_b=0,
 \qquad
 \mathbf T_c
 =
 \Delta_0\overline B\,\operatorname{Cov}(A,C).
 \label{eq:V120-constant-b-remainder}
\end{equation}
Thus rowwise recombination alone cannot universally cancel the
off-cut pair defect.
\end{proposition}

\begin{proof}
When \(B\) is constant,
\[
 \mathbf C_b=0,\qquad
 \mathbf C_c=\overline B\,\operatorname{Cov}(A,C),
\qquad
 \mathbf K_{ab}^{\rhd}=\mathbf K_{ac}^{\lhd}=0.
\]
Insert these identities in
\eqref{eq:V120-recombined-b}--\eqref{eq:V120-recombined-c}, and use
\(q_0-M_0=\Delta_0\).
\end{proof}

\begin{firewall}[Recombination is signed, not an absolute-value order]
\label{fire:V120-recombination-signed}
Equations \eqref{eq:V120-recombined-b}--%
\eqref{eq:V120-recombined-c} are exact signed-amplitude identities.
They allow a coefficient-level trace-norm triangle only after each
physical V98 family has been assembled.  They do not imply an operator
order between the absolute value of the sum and the sum of the two
absolute values.
\end{firewall}

\section{The shared-row Cartan compatibility gate}
\label{sec:V120-compatibility}

Let \(P_e\) be a positive output projection of the old \(ab\) pair bank
and let \(P_f\) be a positive output projection of the off-cut \(bc\)
bank.  Define
\begin{equation}
 \zeta(P_e,P_f)
 :=
 \frac{\|P_e\otimes P_f\|_{\rm 3prod}}
      {\|P_e\|_{\rm prod}}
 \quad\hbox{when }\|P_e\|_{\rm prod}>0.
 \label{eq:V120-compatibility}
\end{equation}

\begin{lemma}[Elementary compatibility range]
\label{lem:V120-compatibility-range}
For pair projections on disjoint coordinate banks sharing row \(b\),
\begin{equation}
 0\le\zeta(P_e,P_f)\le1.
 \label{eq:V120-compatibility-range}
\end{equation}
\end{lemma}

\begin{proof}
Positivity gives the lower bound.  Since \(0\le P_f\le I\), the
one-sided reduced-density proof of
\Cref{lem:V118-one-sided-star} gives
\[
 \|P_e\otimes P_f\|_{\rm 3prod}
 \le\|P_e\|_{\rm prod}.
\]
Divide.
\end{proof}

\begin{proposition}[Cartan--Cartan compatibility can be one]
\label{prop:V120-Cartan-Cartan-compatible}
Suppose \(P_e\) is the Cartan projection in
\(U_a\otimes V_b\) and \(P_f\) is the Cartan projection in
\(V'_b\otimes W_c\), with the two \(b\)-representations on disjoint
coordinate factors.  Then
\begin{equation}
 \|P_e\|_{\rm prod}=1,
 \qquad
 \|P_e\otimes P_f\|_{\rm 3prod}=1,
 \qquad
 \boxed{\zeta(P_e,P_f)=1.}
 \label{eq:V120-Cartan-compatibility-one}
\end{equation}
\end{proposition}

\begin{proof}
Choose normalized highest-weight vectors in
\(U,V,V',W\).  Their three-original-row product tensor has the middle
row equal to the tensor product of the highest vectors in \(V,V'\).
It lies simultaneously in both Cartan components, so the expectation
of \(P_e\otimes P_f\) is one.  The individual statement is the same
Cartan highest-vector argument.
\end{proof}

\begin{corollary}[No universal shared-row second square root]
\label{cor:V120-no-universal-second-root}
There is no universal estimate
\[
 \zeta(P_e,P_f)
 \le
 C\sqrt{\frac{\Xi_b}{\Xi_c}}
\]
for an asymmetric off-cut Cartan output based only on positivity,
invariance, and the two pair labels.
\end{corollary}

\begin{proof}
Use \Cref{prop:V120-Cartan-Cartan-compatible} and let
\(\Xi_b/\Xi_c\to0\).
\end{proof}

\begin{assessment}[Exact branch now requiring physical inspection]
\label{ass:V120-physical-target}
The simultaneous presence of
\begin{enumerate}[label=\textup{(\roman*)}]
\item an asymmetric \(bc\) Cartan output with weighted multiplier of
      square-root size;
\item an old \(ab\) output compatible with the same middle-row
      highest vector; and
\item the constant-\(B\) \(ac\) hybrid remainder
      \eqref{eq:V120-constant-b-remainder}
\end{enumerate}
is the first exact persistent configuration not removed by the
positive-capacity programme.  V120 does not assert that one V98 graph
monomial realizes all three conditions.  It proves that each condition
is allowed by the separate estimates, so the next step must inspect
their joint physical incidence.  Another abstract Schur or heat bound
would be circular.
\end{assessment}

\section{V120 ledger and V121 mandate}
\label{sec:V120-ledger}

\begin{theorem}[Integrated V120 statement]
\label{thm:V120-integrated}
Preserving V1--V119, V120 proves:
\begin{enumerate}[label=\textup{(V120.\arabic*)},leftmargin=6.7em]
\item the mandatory SM V135 and NS V31 companion audit;
\item the exact asymmetric Cartan heat limit
      \eqref{eq:V120-Cartan-limit};
\item the square-root rather than full endpoint ratio
      \eqref{eq:V120-square-root-scale}--%
      \eqref{eq:V120-no-full-ratio};
\item unit product visibility of the obstructing Cartan output;
\item insufficiency of heat decay plus aggregate Schur compression;
\item the exact rowwise recombination
      \eqref{eq:V120-recombined-b}--%
      \eqref{eq:V120-recombined-c};
\item the constant-\(B\) surviving remainder
      \eqref{eq:V120-constant-b-remainder}; and
\item the exact shared-row compatibility coefficient and the
      Cartan--Cartan equality
      \eqref{eq:V120-Cartan-compatibility-one}.
\end{enumerate}
\end{theorem}

\begin{proof}
The items are
\Cref{sec:V120-audit,
thm:V120-Cartan-square-root,
cor:V120-heat-gate-fails,
prop:V120-Cartan-visible,
thm:V120-heat-Schur-insufficient,
thm:V120-rowwise-recombination,
prop:V120-constant-b,
lem:V120-compatibility-range,
prop:V120-Cartan-Cartan-compatible}.
\end{proof}

\begin{openproblem}[V121 mandate: physical incidence of the persistent
Cartan configuration]
\label{op:V121-mandate}
\begin{enumerate}
\item Return to one complete V98 residual graph monomial and impose the
      exact local partition labels defining the old \(ab\) pair, the
      off-cut \(bc\) bank, and the reduced \(ac\) hybrid coordinate.
\item Determine whether an \(ab\) Cartan output and the asymmetric
      \(bc\) Cartan output can occur simultaneously with a constant
      reduced \(b\)-spectator in one final irreducible block.
\item If local incidence excludes the triple, prove the exclusion
      before trace norm and remove the corresponding outer labels.
\item If it is allowed, compute the complete signed Wilson coefficient
      on that block, including
      \eqref{eq:V120-recombined-c}, the old pair defect, collision
      factors, and final resolvent.
\item Test the joint scaling
      \(n\to\infty\), \(s_\alpha c_n\to u\), and
      \(\Xi_b/\Xi_c\to0\).  Either obtain the missing second square root
      from the assembled coefficient or record a counterexample to the
      proposed intermediate marked-tree bound.
\item Do not infer the Wilson result from the heat benchmark; prove any
      growing-representation comparison used.
\end{enumerate}
\end{openproblem}

\section{Firewall after V120}
\label{sec:V120-firewall}

\begin{firewall}[Current claim boundary]
\label{fire:V120-current-boundary}
V120 completes the required audit, proves a sharp heat-benchmark
square-root barrier, and isolates the exact joint physical incidence
which must be tested next.  It does not prove that the Wilson multiplier
has the same asymptotic, that the persistent Cartan triple occurs in a
V98 block, that it is cancelled if it occurs, the V119 scalar gate,
residual off-cut closure, JREC, quartic or all-order MTA, continuum
Yang--Mills existence, or a positive spectral gap.  The full
Yang--Mills existence and mass gap remain open.
\end{firewall}

\section{Append-only record for V120}
\label{sec:V120-record}

V120 was appended to the complete V119 manuscript.  The exact V119
parent had \(80093\) newline-delimited source records, \(2787078\)
bytes, and SHA--256 digest
\[
 \texttt{ad5d480dc29a57238a1f2d6d91e316ea2a534579696069f64bc7df4feef4557d}.
\]
No V119 mathematical content was changed.  The sole final
\verb|\end{document}| is restored below.

% =============================================================================
% END APPENDED VERSION V120 MATERIAL
% =============================================================================

% =============================================================================
% BEGIN APPENDED VERSION V121 MATERIAL
% =============================================================================

\chapter{V121: Physical Incidence of the Persistent Cartan Block and
the Exact Wilson Comparison Gate}
\label{ch:V121}

\section{Purpose}
\label{sec:V121-purpose}

V120 isolated three conditions which the earlier estimates did not
jointly decide: an asymmetric \(bc\) Cartan output, an old \(ab\) output
compatible with the same middle-row highest vector, and a constant
reduced \(b\)-spectator leaving only the \(ac\) hybrid row.  V121 returns
to the local partition monomial before trace norm.

The conditions are compatible.  They may be placed on disjoint
coordinate banks, all relevant Cartan Littlewood--Richardson
multiplicities equal one, and a final Cartan output gives a nonzero
physical outer block.  Thus no local-incidence exclusion removes the
V120 heat obstruction.

For the actual Wilson law, the block is decided by one explicit central
Fourier scalar.  V121 derives its exact character integral and proves a
quantitative comparison theorem: an \(o(n^{-1})\) Wilson-to-heat
comparison on the transition ray transfers the V120 square-root barrier
to the Wilson multiplier.  Conversely, the scalar V119 route can close
only if the Wilson correction cancels the heat Cartan cross term at
order \(n^{-1}\).  That growing-representation asymptotic, not another
capacity estimate, is the next acquisition.

\section{A compatible local-partition realization}
\label{sec:V121-incidence}

Take four original rows \(o,a,b,c\) and four mutually disjoint
coordinate banks
\[
 E_{oa},\qquad E_{ab},\qquad E_{bc},\qquad E_{ac}.
\]
At a coordinate in \(E_{ij}\), choose the local replica partition with
\(\{i,j\}\) as its only nontrivial block; all other rows are singleton
or inactive there.  Put the V113 old same-pair multiplier on \(E_{ab}\),
the V115 off-cut bank on \(E_{bc}\), and the reduced \(ac\) hybrid
connection on \(E_{ac}\).  The optional \(E_{oa}\) bank supplies the
remaining external collision edge and plays no role in the argument.

\begin{lemma}[The four local pair prescriptions are compatible]
\label{lem:V121-local-compatibility}
The preceding prescriptions define a valid product of local replica
partition monomials.  In particular:
\begin{enumerate}[label=\textup{(\roman*)}]
\item the \(ab\) and \(bc\) pair blocks never occur at the same
      coordinate;
\item row \(a\) is inactive on the off-cut bank \(E_{bc}\);
\item row \(b\) is inactive on the reduced connection bank \(E_{ac}\);
      and
\item after the off-cut factor is removed, the reduced \(b\)-factor is
      constant on \(E_{ac}\).
\end{enumerate}
\end{lemma}

\begin{proof}
Each coordinate carries only one displayed two-element block, so every
local prescription is a set partition.  The coordinate banks are
disjoint, hence blocks which share a row never compete inside one local
partition.  The inactive-row statements are part of the definitions of
\(E_{bc}\) and \(E_{ac}\).  The reduced \(b\)-product contains no
nontrivial factor on \(E_{ac}\), proving the last assertion.  The old
\(ab\) output remains an outer same-label multiplier and is not part of
the reduced \(B_1\) factor in the V115 peeling notation.
\end{proof}

Let each coordinate group be \(SU(3)\).  Choose irreducible input labels
on each pair and select its Cartan output.  On \(E_{bc}\), take
\[
 R_m=(m,0)\quad\hbox{on \(b\)},\qquad
 R_n=(n,0)\quad\hbox{on \(c\)},\qquad
 R_{n+m}\quad\hbox{as output}.
\]

\begin{theorem}[The persistent Cartan configuration occurs in a
nonzero outer block]
\label{thm:V121-persistent-block-exists}
There is a complete physical outer label for which:
\begin{enumerate}[label=\textup{(\roman*)}]
\item the old \(ab\) pair output is Cartan;
\item the off-cut \(bc\) output is \(R_{n+m}\);
\item their common-row compatibility is one;
\item the reduced \(b\)-spectator is constant;
\item the reduced signed connection is the \(ac\) row; and
\item every selected intermediate and final Cartan multiplicity is one.
\end{enumerate}
Consequently the three conditions in
\Cref{ass:V120-physical-target} are not excluded by local partition
incidence or by Littlewood--Richardson multiplicity.
\end{theorem}

\begin{proof}
Use the disjoint banks of
\Cref{lem:V121-local-compatibility}.  For irreducibles of a connected
compact group, the Cartan component of a tensor product occurs with
multiplicity one and contains the tensor product of highest-weight
vectors.  Select that component on \(E_{ab}\) and \(E_{bc}\), and select
the Cartan component again when the remaining bank outputs are combined
into the final output.  The tensor product of all local highest-weight
vectors is nonzero and belongs to every selected component, so the
complete physical compression is nonzero and every selected Cartan
multiplicity is one.

On row \(b\), the highest vector is the tensor product of its
\(E_{ab}\) and \(E_{bc}\) highest vectors.  It therefore attains both
Cartan projections simultaneously, proving compatibility one exactly
as in \Cref{prop:V120-Cartan-Cartan-compatible}.  The constant reduced
\(b\)-factor and the \(ac\) hybrid assignment are
\Cref{lem:V121-local-compatibility}\textup{(iii)--(iv)}.
\end{proof}

\begin{firewall}[Existence of a representation path is not coefficient
saturation]
\label{fire:V121-incidence-scope}
\Cref{thm:V121-persistent-block-exists} proves that the relevant V98
outer subspace is nonzero and that no Schur or multiplicity debit is
forced there.  It does not prove that arbitrary input coefficients
saturate every remaining Wilson defect, collision factor, reduced
covariance, and final resolvent simultaneously.  Those signed factors
must now be computed on the displayed block.
\end{firewall}

\section{Exact Wilson scalar on the asymmetric ray}
\label{sec:V121-Wilson-scalar}

Let
\[
 \mu_\alpha(\dd U)=q_\alpha(U)\,\dd U
\]
be the normalized central inversion-symmetric one-link Wilson measure
of \eqref{eq:V11-qalpha}.  Centrality and Schur's lemma define
\(q_{\alpha,R}\) by
\begin{equation}
 \int_{SU(3)}D^R(U)\,\mu_\alpha(\dd U)
 =
 q_{\alpha,R}I_{H_R}.
 \label{eq:V121-Wilson-scalar}
\end{equation}

\begin{lemma}[Exact normalized character integral]
\label{lem:V121-character-integral}
For every irreducible \(R\),
\begin{equation}
 \boxed{
 q_{\alpha,R}
 =
 \frac1{d_R}
 \int_{SU(3)}\chi_R(U)\,\mu_\alpha(\dd U).}
 \label{eq:V121-character-integral}
\end{equation}
The scalar is real and satisfies
\(q_{\alpha,R^*}=q_{\alpha,R}\).
\end{lemma}

\begin{proof}
Take the trace of \eqref{eq:V121-Wilson-scalar} and divide by
\(d_R\).  Inversion symmetry sends \(\chi_R(U)\) to
\(\overline{\chi_R(U)}=\chi_{R^*}(U)\), so the integral is real and is
unchanged under \(R\mapsto R^*\).
\end{proof}

For \(R_k=(k,0)\), define the exact Wilson Cartan defect
\begin{align}
 \delta_{\alpha;n,m}^{\rm W}
 &:=
 q_{\alpha,R_{n+m}}
 -
 q_{\alpha,R_n}q_{\alpha,R_m},
 \label{eq:V121-Wilson-Cartan-defect}\\
 a_{\alpha;n,m}^{\rm W}
 &:=
 (1+s_\alpha c_{n+m})
 |\delta_{\alpha;n,m}^{\rm W}|.
 \label{eq:V121-Wilson-Cartan-weight}
\end{align}
The harmless replacement of \(c_k\) by \(1+c_k\) changes only fixed
constants in the endpoint ratio.

\begin{proposition}[Exact decision for the scalar operator-hot route]
\label{prop:V121-scalar-decision}
Suppose the old \(ab\) multiplier is estimated only through the V109
capacity already used in \eqref{eq:V119-old-pair-capacity}, and no
additional label-conditioned \(ab\) debit or signed cancellation is
invoked.  Then the scalar operator-hot route of V119 closes the persistent
Cartan subfamily if
\begin{equation}
 a_{\alpha;n,m}^{\rm W}
 \le
 C_m\frac{1+c_m}{1+c_n}
 \label{eq:V121-required-Wilson-ratio}
\end{equation}
on the transition scaling
\(s_\alpha c_n\asymp1\).  No product-capacity or multiplicity factor
improves the left side on the block of
\Cref{thm:V121-persistent-block-exists}; consequently this route cannot
close if the displayed ratio fails and the old pair capacity is sharp on
the conditioned label.
\end{proposition}

\begin{proof}
On that block the reduced \(b\)-factor is constant, so
\eqref{eq:V120-constant-b-remainder} leaves the off-cut pair difference
multiplying the \(ac\) covariance.  The \(bc\) and \(ab\) Cartan
projections are simultaneously product-visible by
\Cref{thm:V121-persistent-block-exists}.  Thus the missing endpoint
comparison in \eqref{eq:V119-scalar-gate} must be supplied by the
weighted scalar itself under the stated routing hypothesis.  Here
\[
 \frac{\Xi_b}{\Xi_c}
 =
 \frac{1+c_m}{1+c_n}
\]
for the one-coordinate off-cut bank, yielding
\eqref{eq:V121-required-Wilson-ratio}.
\end{proof}

\begin{firewall}[Visibility does not imply capacity saturation]
\label{fire:V121-visibility-not-saturation}
Cartan product visibility one says that the projector itself has no
Schur debit.  It does not say that the old \(ab\) multiplier eigenvalue
on its Cartan label equals the full capacity
\(\kappa_{ab}\).  Therefore
\Cref{prop:V121-scalar-decision} is a decision theorem for the printed
V119 scalar route, not a universal necessary condition for the complete
signed coefficient.
\end{firewall}

\section{A quantitative Wilson-to-heat comparison decision}
\label{sec:V121-comparison}

Put
\begin{equation}
 \epsilon_{\alpha,k}
 :=
 q_{\alpha,R_k}-e^{-s_\alpha c_k}.
 \label{eq:V121-comparison-error}
\end{equation}

\begin{theorem}[An \(o(n^{-1})\) comparison transfers the heat barrier]
\label{thm:V121-comparison-transfers}
Let \(\alpha_j\to\infty\) and \(n_j\to\infty\) satisfy
\begin{equation}
 s_{\alpha_j}c_{n_j}\longrightarrow u\in(0,\infty),
 \label{eq:V121-Wilson-transition}
\end{equation}
and fix \(m\ge1\).  If
\begin{equation}
 n_j\left(
 |\epsilon_{\alpha_j,n_j+m}|
 +|\epsilon_{\alpha_j,n_j}|
 +|\epsilon_{\alpha_j,m}|
 \right)
 \longrightarrow0,
 \label{eq:V121-little-o-comparison}
\end{equation}
then
\begin{equation}
 \boxed{
 n_j a_{\alpha_j;n_j,m}^{\rm W}
 \longrightarrow
 2um(1+u)e^{-u}.}
 \label{eq:V121-Wilson-square-root-limit}
\end{equation}
In particular,
\eqref{eq:V121-required-Wilson-ratio} fails along this sequence.
\end{theorem}

\begin{proof}
Let \(h_{\alpha,k}:=e^{-s_\alpha c_k}\).  Expansion of the exact
difference gives
\begin{align}
 \delta_{\alpha;n,m}^{\rm W}
 =
 \delta_{n,m}^{\mathsf h}
 +\epsilon_{\alpha,n+m}
 -h_{\alpha,m}\epsilon_{\alpha,n}
 -h_{\alpha,n}\epsilon_{\alpha,m}
 -\epsilon_{\alpha,n}\epsilon_{\alpha,m}.
 \label{eq:V121-defect-comparison}
\end{align}
All heat multipliers have modulus at most one, and the Wilson
multipliers are uniformly bounded.  Hence
\eqref{eq:V121-little-o-comparison} implies
\[
 \delta_{\alpha_j;n_j,m}^{\rm W}
 -
 \delta_{n_j,m}^{\mathsf h}
 =
 o(n_j^{-1}).
\]
Under \eqref{eq:V121-Wilson-transition},
\[
 1+s_{\alpha_j}c_{n_j+m}\longrightarrow1+u.
\]
Now apply \eqref{eq:V120-Cartan-limit}.  Since its limiting constant is
strictly positive, taking absolute values preserves the limit.
Finally \((1+c_m)/(1+c_{n_j})\asymp_m n_j^{-2}\), so the required
full-ratio estimate fails.
\end{proof}

\begin{corollary}[Exact cancellation required by a full-ratio Wilson
law]
\label{cor:V121-required-comparison-cancellation}
If instead
\eqref{eq:V121-required-Wilson-ratio} holds under
\eqref{eq:V121-Wilson-transition}, then necessarily
\begin{align}
 &\epsilon_{\alpha,n+m}
 -e^{-s_\alpha c_m}\epsilon_{\alpha,n}
 -e^{-s_\alpha c_n}\epsilon_{\alpha,m}
 -\epsilon_{\alpha,n}\epsilon_{\alpha,m}
 \notag\\
 &\qquad=
 -\delta_{n,m}^{\mathsf h}
 +O_m(n^{-2}).
 \label{eq:V121-required-cancellation}
\end{align}
The right-hand leading term is of order \(n^{-1}\), with nonzero
coefficient fixed by \eqref{eq:V120-Cartan-limit}.
\end{corollary}

\begin{proof}
The full-ratio hypothesis and the bounded weight
\(1+s_\alpha c_{n+m}\asymp1\) give
\(\delta_{\alpha;n,m}^{\rm W}=O_m(n^{-2})\).
Rearrange the exact identity
\eqref{eq:V121-defect-comparison}.  The size and nonzero leading
coefficient of the heat term are
\Cref{thm:V120-Cartan-square-root}.
\end{proof}

\begin{assessment}[Why fixed-representation convergence is
insufficient]
\label{ass:V121-fixed-not-enough}
For fixed \(m\), the known expansion of \(q_{\alpha,R_m}\) controls the
third error in \eqref{eq:V121-little-o-comparison}.  It does not control
\(\epsilon_{\alpha,n}\) or \(\epsilon_{\alpha,n+m}\) when
\(s_\alpha c_n\to u\).  Pointwise density domination from V11 gives a
one-sided Dirichlet-form bound, not the signed \(o(n^{-1})\) Fourier
comparison required here.  Thus neither earlier result decides
\eqref{eq:V121-required-cancellation}.
\end{assessment}

\section{Exact torus integral for the next asymptotic}
\label{sec:V121-torus-integral}

Let \(T\subset SU(3)\) be the diagonal maximal torus, \(W\) its Weyl
group, and \(\Delta\) the Weyl denominator.  Parametrize
\[
 t(\theta_1,\theta_2)
 =
 \operatorname{diag}
 \bigl(
 e^{i\theta_1},
 e^{i\theta_2},
 e^{-i(\theta_1+\theta_2)}
 \bigr).
\]

\begin{proposition}[Exact Weyl-integral representation]
\label{prop:V121-Weyl-integral}
For every \(k\ge0\),
\begin{equation}
 \boxed{
 q_{\alpha,R_k}
 =
 \frac{
 \displaystyle
 \int_{[-\pi,\pi]^2}
 \frac{\chi_{R_k}(t(\theta))}{d_{R_k}}\,
 e^{-\alpha(1-\frac13\operatorname{ReTr}t(\theta))}
 |\Delta(t(\theta))|^2\,\dd\theta
 }{
 \displaystyle
 \int_{[-\pi,\pi]^2}
 e^{-\alpha(1-\frac13\operatorname{ReTr}t(\theta))}
 |\Delta(t(\theta))|^2\,\dd\theta
 }.}
 \label{eq:V121-Weyl-integral}
\end{equation}
The common Weyl normalization cancels in the ratio.
\end{proposition}

\begin{proof}
Insert the density \eqref{eq:V11-qalpha} into
\eqref{eq:V121-character-integral}.  Apply the Weyl integration formula
to the central numerator and denominator.  The factors
\(|W|^{-1}\) and the torus Haar normalization are identical and cancel.
\end{proof}

\begin{remark}[The required uniformity scale]
\label{rem:V121-uniformity-scale}
In \eqref{eq:V121-Wilson-transition},
\(n\asymp s_\alpha^{-1/2}\asymp\alpha^{1/2}\) under fundamental
calibration.  Thus V122 must analyze
\eqref{eq:V121-Weyl-integral} with the character frequency and the
Laplace concentration scale growing together.  A remainder merely
uniform for fixed \(k\), or merely \(o(1)\), cannot decide the
order-\(n^{-1}\) cancellation in
\eqref{eq:V121-required-cancellation}.
\end{remark}

\section{V121 ledger and V122 mandate}
\label{sec:V121-ledger}

\begin{theorem}[Integrated V121 statement]
\label{thm:V121-integrated}
Preserving V1--V120, V121 proves:
\begin{enumerate}[label=\textup{(V121.\arabic*)},leftmargin=6.7em]
\item compatibility of the \(oa,ab,bc,ac\) local pair prescriptions;
\item existence of a nonzero multiplicity-one physical outer block
      satisfying all three V120 persistent Cartan conditions;
\item the exact Wilson character scalar
      \eqref{eq:V121-character-integral};
\item the precisely scoped scalar route condition
      \eqref{eq:V121-required-Wilson-ratio};
\item the comparison identity
      \eqref{eq:V121-defect-comparison};
\item transfer of the heat square-root barrier under the sharp
      \(o(n^{-1})\) comparison
      \eqref{eq:V121-little-o-comparison};
\item the exact order-\(n^{-1}\) cancellation
      \eqref{eq:V121-required-cancellation} required by a full-ratio
      Wilson law; and
\item the Weyl integral
      \eqref{eq:V121-Weyl-integral} for the remaining asymptotic.
\end{enumerate}
\end{theorem}

\begin{proof}
The items are
\Cref{lem:V121-local-compatibility,
thm:V121-persistent-block-exists,
lem:V121-character-integral,
prop:V121-scalar-decision,
thm:V121-comparison-transfers,
cor:V121-required-comparison-cancellation,
prop:V121-Weyl-integral}.
\end{proof}

\begin{openproblem}[V122 mandate: transition-scale Wilson character
asymptotics]
\label{op:V122-mandate}
\begin{enumerate}
\item Rescale the Weyl integral
      \eqref{eq:V121-Weyl-integral} at
      \(\theta=\alpha^{-1/2}x\), retaining the exact Weyl denominator
      and normalized symmetric-power character.
\item Use inversion symmetry to identify which odd local terms vanish,
      but do not assume that this makes the growing-character remainder
      \(O(\alpha^{-1})\).
\item Prove either
      \eqref{eq:V121-little-o-comparison} or the cancellation
      \eqref{eq:V121-required-cancellation} for fixed \(m\) and
      \(n\asymp\sqrt\alpha\).
\item Control the complement of the normal ball uniformly in \(n\);
      the character bound \(|\chi_R|\le d_R\) may be used only after
      normalization.
\item If the heat square-root term survives, return to the complete
      signed physical coefficient rather than claiming the V119 scalar
      gate.  Decide whether the intermediate marked-tree bound must be
      weakened or rerouted.
\item Keep the conclusion separate from the full Yang--Mills mass-gap
      claim.
\end{enumerate}
\end{openproblem}

\section{Firewall after V121}
\label{sec:V121-firewall}

\begin{firewall}[Current claim boundary]
\label{fire:V121-current-boundary}
V121 proves that the persistent Cartan configuration is physically
incident and reduces the existing scalar route to a concrete
transition-scale Wilson character asymptotic.  It does not prove either
side of that asymptotic, show saturation of the old \(ab\) capacity,
disprove the complete signed marked-tree estimate, close the residual
off-cut branch, prove JREC, prove quartic or all-order MTA, construct
continuum Yang--Mills theory, or prove a positive spectral gap.  The
full Yang--Mills existence and mass gap remain open.
\end{firewall}

\section{Append-only record for V121}
\label{sec:V121-record}

V121 was appended to the complete V120 manuscript.  The exact V120
parent had \(80726\) newline-delimited source records, \(2807370\)
bytes, and SHA--256 digest
\[
 \texttt{7a4c6a0e013590bb8b77fbb09b6a9881ffd59bccea69351d1249dcf655484d07}.
\]
No V120 mathematical content was changed.  The sole final
\verb|\end{document}| is restored below.

% =============================================================================
% END APPENDED VERSION V121 MATERIAL
% =============================================================================

% =============================================================================
% BEGIN APPENDED VERSION V122 MATERIAL
% =============================================================================

\chapter{V122: Corrected Wilson Clock and the Direct
Transition-Scale Cartan Obstruction}
\label{ch:V122}

\section{Purpose and upstream correction}
\label{sec:V122-purpose}

V121 reduced the unresolved scalar route to a growing-representation
Wilson character calculation.  Before carrying it out, one must correct
a factor of two in the conversion between the Wilson concentration
\(\alpha\) and the exact fundamental clock \(s_\alpha\).  With the
Wilson density printed in \eqref{eq:V11-qalpha} and
\(C_F=4/3\), the correct relation is
\[
 s_\alpha=\frac3\alpha+O(\alpha^{-2}),
\]
not \(3/(2\alpha)+O(\alpha^{-2})\).

After that correction, V122 solves the part of the V121 asymptotic which
actually decides the Cartan defect.  The normalized character of
\(R_n=\operatorname{Sym}^n(\mathbb C^3)\) has an exact coherent-state
formula.  On \(n/\sqrt\alpha\to\rho\), the rescaled Wilson law becomes a
Gaussian law on traceless Hermitian matrices, and the coherent-state
linear statistic is exactly \(N(0,2)\).  Consequently
\[
 q_{\alpha,R_n}\longrightarrow e^{-\rho^2}
\]
and, more importantly,
\[
 \sqrt\alpha\,
 \bigl(q_{\alpha,R_{n+m}}-q_{\alpha,R_n}\bigr)
 \longrightarrow
 -2m\rho e^{-\rho^2}.
\]
This direct finite-difference limit proves that the exact Wilson Cartan
defect has the same leading square-root obstruction as the V120 heat
benchmark.  It does not require the stronger individual
\(o(n^{-1})\) comparison postulated in V121.

\section{The fundamental clock from Laplace equipartition}
\label{sec:V122-clock}

Put
\[
 V(U):=1-\frac13\operatorname{ReTr}U.
\]
Thus the one-link probability law is
\[
 \nu_\alpha(\dd U)
 =
 Z_\alpha^{-1}e^{-\alpha V(U)}\,\dd U.
\]
The function \(V\) has a unique minimum at the identity and a
nondegenerate Hessian on the eight-dimensional group \(SU(3)\).

\begin{lemma}[Laplace equipartition for the Wilson well]
\label{lem:V122-equipartition}
As \(\alpha\to\infty\),
\begin{equation}
 \mathbb E_{\nu_\alpha}V
 =
 \frac4\alpha+O(\alpha^{-2}).
 \label{eq:V122-equipartition}
\end{equation}
\end{lemma}

\begin{proof}
Choose a fixed exponential neighbourhood of the identity.  On its
complement \(V\ge c>0\), so its contribution to \(Z_\alpha\) and to all
of its \(\alpha\)-derivatives is exponentially small.  In normal
coordinates on the neighbourhood, the phase has a positive-definite
quadratic part and the phase and Haar Jacobian are analytic.  The
ordinary finite-dimensional Laplace expansion, with Gaussian domination
of the Taylor remainder, gives
\[
 Z_\alpha
 =
 \alpha^{-8/2}
 \left(A_0+\frac{A_1}{\alpha}+O(\alpha^{-2})\right),
 \qquad A_0>0.
\]
Differentiation is justified by the same localized domination.  Since
\[
 \mathbb E_{\nu_\alpha}V
 =
 -\frac{\dd}{\dd\alpha}\log Z_\alpha,
\]
logarithmic differentiation gives
\[
 -\frac{\dd}{\dd\alpha}\log Z_\alpha
 =
 \frac{8}{2\alpha}+O(\alpha^{-2}),
\]
which is \eqref{eq:V122-equipartition}.
\end{proof}

\begin{theorem}[Correct fundamental clock]
\label{thm:V122-correct-clock}
For the exact calibration
\[
 s_\alpha=-\frac{\log\eta_F(\alpha)}{C_F},
 \qquad C_F=\frac43,
\]
one has
\begin{equation}
 \boxed{
 \eta_F(\alpha)=1-\frac4\alpha+O(\alpha^{-2}),
 \qquad
 s_\alpha=\frac3\alpha+O(\alpha^{-2}),
 \qquad
 \alpha s_\alpha\longrightarrow3.}
 \label{eq:V122-correct-clock}
\end{equation}
\end{theorem}

\begin{proof}
Inversion invariance makes the fundamental Fourier scalar real.
Therefore its normalized-character formula gives the exact identity
\[
 \eta_F(\alpha)
 =
 \mathbb E_{\nu_\alpha}
 \frac13\operatorname{ReTr}U
 =
 1-\mathbb E_{\nu_\alpha}V(U).
\]
Apply \Cref{lem:V122-equipartition}.  Finally,
\[
 -\log\left(1-\frac4\alpha+O(\alpha^{-2})\right)
 =
 \frac4\alpha+O(\alpha^{-2}),
\]
and division by \(C_F=4/3\) proves the assertion.
\end{proof}

\begin{corollary}[Correct fixed-representation conversion]
\label{cor:V122-fixed-conversion}
For every fixed irreducible \(R\),
\begin{equation}
 q_{\alpha,R}
 =
 1-\frac{3C_2(R)}{\alpha}+O_R(\alpha^{-2}).
 \label{eq:V122-fixed-conversion}
\end{equation}
In particular, for every fixed \(m\),
\begin{equation}
 1-q_{\alpha,R_m}=O_m(\alpha^{-1}).
 \label{eq:V122-fixed-small}
\end{equation}
\end{corollary}

\begin{proof}
The fixed-representation expansion
\[
 q_{\alpha,R}
 =
 1-s_\alpha C_2(R)+O_R(s_\alpha^2)
\]
is the convention-free content of the calibrated V10 character
calculation.  Insert \eqref{eq:V122-correct-clock}.
\end{proof}

\begin{assessment}[Precise supersession of the factor-two formulas]
\label{ass:V122-factor-two-supersession}
Equation \eqref{eq:V122-correct-clock} supersedes
\eqref{eq:V21-temporal-leading-constant} and every later literal use of
\(\alpha s_\alpha\to3/2\).  In particular:
\begin{enumerate}[label=\textup{(\roman*)}]
\item the exact V21 matched identity
      \(\kappa_H\alpha_\beta=s_\alpha\gamma/3\) is unchanged, but its
      asymptotic is \(\gamma/\alpha\), not \(\gamma/(2\alpha)\);
\item the V22 approximate identity has coefficient \(3/\alpha\):
      \[
       \mathcal C_\alpha u
       =
       u-\frac3\alpha\mathcal L_Eu+O_{u,E}(\alpha^{-2});
      \]
\item the corresponding V22 fixed-tilt formulas replace
      \(3/(2\alpha)\) by \(3/\alpha\), and replace the leading
      chi-square coefficient \(3/\alpha\) by \(6/\alpha\);
\item occurrences in V27 and V29 which retain the exact variable
      \(s_\alpha\) remain valid; only their printed conversion from
      \(s_\alpha\) to \(\alpha^{-1}\) changes.
\end{enumerate}
All later arguments using only \(s_\alpha\asymp\alpha^{-1}\), or exact
identities in \(s_\alpha\), are unaffected.  No prior text is silently
altered; this assessment is the controlling correction.
\end{assessment}

\begin{firewall}[Why this correction is compulsory]
\label{fire:V122-clock-not-convention}
The factor in \eqref{eq:V122-correct-clock} is not a choice of
Lie-algebra norm.  It follows from the exact scalar identity
\(\eta_F=1-\mathbb E_{\nu_\alpha}V\) and the dimension-eight Laplace
equipartition law.  Changing a coordinate metric changes the Hessian and
Gaussian volume together and cannot change \(4/\alpha\).
\end{firewall}

\section{The rescaled Wilson law}
\label{sec:V122-rescaled-law}

Let \(\mathfrak h_0\) be the real vector space of traceless Hermitian
\(3\times3\) matrices, equipped with
\[
 \langle H,K\rangle_0:=\operatorname{Tr}(HK),
 \qquad
 \|H\|_0^2=\operatorname{Tr}(H^2).
\]
It has real dimension eight.  Let \(\gamma\) be the centered Gaussian
probability measure on \(\mathfrak h_0\) with density
\begin{equation}
 \gamma(\dd H)
 =
 C_\gamma e^{-\operatorname{Tr}(H^2)/6}\,\dd H.
 \label{eq:V122-Gaussian-law}
\end{equation}

\begin{lemma}[Rescaled Wilson Laplace limit with tails]
\label{lem:V122-rescaled-Wilson}
In the identity exponential chart, write
\[
 U=\exp(iH/\sqrt\alpha).
\]
The pushforward of \(\nu_\alpha\), restricted to any fixed sufficiently
small normal neighbourhood and then rescaled to \(H\), converges to
\(\gamma\).  More precisely, its density converges locally uniformly to
the density in \eqref{eq:V122-Gaussian-law}, is bounded on the rescaled
normal ball by \(Ce^{-c\|H\|_0^2}\), and the probability outside the
unscaled normal neighbourhood is \(O(e^{-c\alpha})\).  Consequently
convergence of expectations holds for every continuous function of
polynomial growth.
\end{lemma}

\begin{proof}
For bounded \(H\), Taylor expansion gives
\begin{align}
 \alpha V\bigl(e^{iH/\sqrt\alpha}\bigr)
 &=
 \frac16\operatorname{Tr}(H^2)
 +O\!\left(\alpha^{-1}\|H\|_0^4\right),
 \label{eq:V122-phase-rescaling}\\
 J(H/\sqrt\alpha)&=1+O(\alpha^{-1}\|H\|_0^2),
 \label{eq:V122-Jacobian-rescaling}
\end{align}
where \(J\) is the exponential-coordinate Haar Jacobian.  The local
quadratic lower bound for \(V\) supplies
\[
 \alpha V(e^{iH/\sqrt\alpha})\ge c\|H\|_0^2
\]
throughout a fixed sufficiently small normal neighbourhood.  The unique
minimum of \(V\) gives a positive lower bound on its complement.  These
facts yield the Gaussian envelope and exponentially small complement.
After the Jacobian factor \(\alpha^{-4}\) is canceled between numerator
and denominator, dominated convergence proves the assertion.
\end{proof}

\section{Exact coherent-state formula and its finite difference}
\label{sec:V122-coherent-state}

Let \(S^5\subset\mathbb C^3\) carry normalized unitary-invariant
probability measure \(\sigma\), and put
\[
 d_n:=\dim\operatorname{Sym}^n(\mathbb C^3)
 =\binom{n+2}{2}.
\]

\begin{lemma}[Exact symmetric-power coherent-state identity]
\label{lem:V122-coherent-identity}
For every \(U\in SU(3)\) and \(n\ge0\),
\begin{equation}
 \boxed{
 \frac{\chi_{R_n}(U)}{d_n}
 =
 \int_{S^5}\langle v,Uv\rangle^n\,\sigma(\dd v).}
 \label{eq:V122-coherent-identity}
\end{equation}
\end{lemma}

\begin{proof}
On \(\operatorname{Sym}^n(\mathbb C^3)\), unitary invariance and Schur's
lemma give
\[
 \int_{S^5}
 |v^{\otimes n}\rangle\langle v^{\otimes n}|
 \,\sigma(\dd v)
 =
 c_n I.
\]
Taking traces shows \(c_n=1/d_n\).  Multiply by the restriction of
\(U^{\otimes n}\) to the symmetric power and take the trace:
\[
 \frac1{d_n}\operatorname{Tr}_{\operatorname{Sym}^n}(U^{\otimes n})
 =
 \int_{S^5}
 \langle v^{\otimes n},U^{\otimes n}v^{\otimes n}\rangle
 \,\sigma(\dd v).
\]
The integrand is \(\langle v,Uv\rangle^n\), proving
\eqref{eq:V122-coherent-identity}.
\end{proof}

For \(\rho\ge0\) and \(H\in\mathfrak h_0\), define
\begin{equation}
 \Phi_\rho(H)
 :=
 \int_{S^5}
 e^{i\rho\langle v,Hv\rangle}\,\sigma(\dd v).
 \label{eq:V122-Phi}
\end{equation}
If \(\lambda_1+\lambda_2+\lambda_3=0\) are the eigenvalues of \(H\),
then the squared coordinate moduli of a \(\sigma\)-distributed vector
have the \(\operatorname{Dirichlet}(1,1,1)\) law.  Hence equivalently
\[
 \Phi_\rho(H)
 =
 2\int_{\substack{p_1,p_2\ge0\\p_1+p_2\le1}}
 e^{i\rho(p_1\lambda_1+p_2\lambda_2+
                 (1-p_1-p_2)\lambda_3)}
 \,\dd p_1\,\dd p_2.
\]

\begin{lemma}[Pointwise coherent transition and discrete derivative]
\label{lem:V122-coherent-transition}
Let \(\alpha_j\to\infty\), \(n_j\to\infty\), and
\[
 \frac{n_j}{\sqrt{\alpha_j}}\longrightarrow\rho\in(0,\infty).
\]
For fixed \(H\in\mathfrak h_0\) and fixed integer \(m\ge1\),
\begin{align}
 \frac{\chi_{R_{n_j}}(e^{iH/\sqrt{\alpha_j}})}{d_{n_j}}
 &\longrightarrow
 \Phi_\rho(H),
 \label{eq:V122-pointwise-character}\\
 \sqrt{\alpha_j}\left[
 \frac{\chi_{R_{n_j+m}}(e^{iH/\sqrt{\alpha_j}})}{d_{n_j+m}}
 -
 \frac{\chi_{R_{n_j}}(e^{iH/\sqrt{\alpha_j}})}{d_{n_j}}
 \right]
 &\longrightarrow
 m\,\partial_\rho\Phi_\rho(H).
 \label{eq:V122-pointwise-difference}
\end{align}
Moreover the absolute value on the left of
\eqref{eq:V122-pointwise-difference} is bounded by
\(m\|H\|_{\rm op}\), uniformly in \(j\).
\end{lemma}

\begin{proof}
Put
\[
 z_j(v):=\langle v,e^{iH/\sqrt{\alpha_j}}v\rangle,
 \qquad
 a(v):=\langle v,Hv\rangle.
\]
Uniformly in \(v\in S^5\),
\[
 z_j(v)
 =
 1+\frac{i\,a(v)}{\sqrt{\alpha_j}}
 +O_H(\alpha_j^{-1}).
\]
Thus
\[
 z_j(v)^{n_j}\longrightarrow e^{i\rho a(v)}
\]
uniformly in \(v\).  The first conclusion follows from
\Cref{lem:V122-coherent-identity}.  For the second, use the exact
factorization
\begin{align*}
 &\frac{\chi_{R_{n_j+m}}(e^{iH/\sqrt{\alpha_j}})}{d_{n_j+m}}
 -
 \frac{\chi_{R_{n_j}}(e^{iH/\sqrt{\alpha_j}})}{d_{n_j}}
 \\
 &\qquad=
 \int_{S^5}z_j(v)^{n_j}\bigl(z_j(v)^m-1\bigr)\,\sigma(\dd v).
\end{align*}
Since \(e^{iH/\sqrt{\alpha_j}}\) is unitary,
\(|z_j(v)|\le1\), and
\[
 \sqrt{\alpha_j}|z_j(v)^m-1|
 \le
 m\sqrt{\alpha_j}|z_j(v)-1|
 \le m\|H\|_{\rm op}.
\]
Also
\[
 \sqrt{\alpha_j}\bigl(z_j(v)^m-1\bigr)
 \longrightarrow im\,a(v).
\]
Dominated convergence gives
\[
 im\int_{S^5}a(v)e^{i\rho a(v)}\,\sigma(\dd v)
 =
 m\,\partial_\rho\Phi_\rho(H).
\]
\end{proof}

\begin{theorem}[Direct Wilson transition and finite-difference limits]
\label{thm:V122-Wilson-transition}
Under the hypotheses of
\Cref{lem:V122-coherent-transition},
\begin{align}
 q_{\alpha_j,R_{n_j}}
 &\longrightarrow
 F(\rho),
 \label{eq:V122-q-limit}\\
 \sqrt{\alpha_j}
 \left(q_{\alpha_j,R_{n_j+m}}-q_{\alpha_j,R_{n_j}}\right)
 &\longrightarrow
 mF'(\rho),
 \label{eq:V122-q-difference-limit}
\end{align}
where
\begin{equation}
 F(\rho)
 :=
 \int_{\mathfrak h_0}\Phi_\rho(H)\,\gamma(\dd H).
 \label{eq:V122-F-definition}
\end{equation}
\end{theorem}

\begin{proof}
On the rescaled normal neighbourhood,
\Cref{lem:V122-coherent-transition} gives pointwise convergence.
For \eqref{eq:V122-q-limit}, normalized characters have modulus at most
one.  For \eqref{eq:V122-q-difference-limit}, the last bound in
\Cref{lem:V122-coherent-transition} is at most \(m\|H\|_0\).
The uniform Gaussian envelope in
\Cref{lem:V122-rescaled-Wilson} is integrable against both majorants.
Outside the normal neighbourhood the normalized characters are bounded
by one, while its probability is \(O(e^{-c\alpha_j})\); even after
multiplication by \(\sqrt{\alpha_j}\), that contribution vanishes.
Dominated convergence proves both limits and permits differentiation of
\eqref{eq:V122-F-definition}.
\end{proof}

\section{The Gaussian coherent statistic is exactly one-dimensional}
\label{sec:V122-Gaussian-collapse}

\begin{lemma}[Exact Gaussian collapse]
\label{lem:V122-Gaussian-collapse}
For every \(\rho\in\mathbb R\),
\begin{equation}
 \boxed{F(\rho)=e^{-\rho^2}.}
 \label{eq:V122-F-exact}
\end{equation}
Consequently
\begin{equation}
 F'(\rho)=-2\rho e^{-\rho^2}.
 \label{eq:V122-F-derivative}
\end{equation}
\end{lemma}

\begin{proof}
Fix \(v\in S^5\) and put
\[
 P_v:=vv^*-\frac13I.
\]
Since \(H\) is traceless,
\[
 \langle v,Hv\rangle
 =
 \operatorname{Tr}(HP_v).
\]
Under \(\gamma\), the coordinates of \(H\) in any
\(\operatorname{Tr}(HK)\)-orthonormal basis of \(\mathfrak h_0\) are
independent centered Gaussians of variance \(3\).  Moreover
\[
 \operatorname{Tr}(P_v^2)
 =
 1-\frac23+\frac13
 =
 \frac23.
\]
Therefore \(\operatorname{Tr}(HP_v)\) is a centered real Gaussian with
variance
\[
 3\operatorname{Tr}(P_v^2)=2,
\]
independently of \(v\).  Its characteristic function is
\[
 \int_{\mathfrak h_0}
 e^{i\rho\langle v,Hv\rangle}\,\gamma(\dd H)
 =
 e^{-\rho^2}.
\]
Fubini's theorem is applicable because the integrand has modulus one.
Integrating over \(v\) proves \eqref{eq:V122-F-exact}; differentiation
gives \eqref{eq:V122-F-derivative}.
\end{proof}

\begin{corollary}[Exact leading Wilson transition law]
\label{cor:V122-exact-transition}
If \(n_j/\sqrt{\alpha_j}\to\rho\in(0,\infty)\), then for every fixed
\(m\ge1\),
\begin{align}
 q_{\alpha_j,R_{n_j}}
 &\longrightarrow e^{-\rho^2},
 \label{eq:V122-exact-q-limit}\\
 \sqrt{\alpha_j}
 \left(q_{\alpha_j,R_{n_j+m}}-q_{\alpha_j,R_{n_j}}\right)
 &\longrightarrow
 -2m\rho e^{-\rho^2}.
 \label{eq:V122-exact-q-difference}
\end{align}
\end{corollary}

\begin{proof}
Combine \Cref{thm:V122-Wilson-transition} with
\Cref{lem:V122-Gaussian-collapse}.
\end{proof}

\begin{assessment}[What was and was not needed from V121]
\label{ass:V122-no-individual-rate}
Equation \eqref{eq:V122-exact-q-limit} identifies the same leading
transition multiplier as the calibrated heat law because
\[
 s_{\alpha_j}c_{n_j}\longrightarrow\rho^2.
\]
It supplies only an \(o(1)\) individual comparison and does not prove
\eqref{eq:V121-little-o-comparison}.  The stronger individual rate is
unnecessary: the exact coherent-state factorization controls the
discrete difference directly at order \(\alpha^{-1/2}\).
\end{assessment}

\section{Failure of the full-ratio Wilson scalar gate}
\label{sec:V122-scalar-obstruction}

\begin{theorem}[Exact Wilson Cartan square-root asymptotic]
\label{thm:V122-Wilson-Cartan-obstruction}
Fix \(m\ge1\), and let
\(\alpha_j,n_j\to\infty\) satisfy
\[
 \frac{n_j}{\sqrt{\alpha_j}}\longrightarrow\rho\in(0,\infty).
\]
For the exact Wilson defect and weighted size from
\eqref{eq:V121-Wilson-Cartan-defect}--%
\eqref{eq:V121-Wilson-Cartan-weight},
\begin{align}
 \sqrt{\alpha_j}\,
 \delta_{\alpha_j;n_j,m}^{\rm W}
 &\longrightarrow
 -2m\rho e^{-\rho^2},
 \label{eq:V122-Wilson-defect-limit}\\
 \boxed{
 \sqrt{\alpha_j}\,
 a_{\alpha_j;n_j,m}^{\rm W}
 \longrightarrow
 2m\rho(1+\rho^2)e^{-\rho^2},}
 \label{eq:V122-Wilson-weighted-alpha-limit}\\
 \boxed{
 n_j a_{\alpha_j;n_j,m}^{\rm W}
 \longrightarrow
 2m\rho^2(1+\rho^2)e^{-\rho^2}.}
 \label{eq:V122-Wilson-weighted-n-limit}
\end{align}
In particular, for every such transition ray and every constant \(C_m\),
\begin{equation}
 a_{\alpha_j;n_j,m}^{\rm W}
 \le
 C_m\frac{1+c_m}{1+c_{n_j}}
 \label{eq:V122-false-full-ratio}
\end{equation}
fails for all sufficiently large \(j\).
\end{theorem}

\begin{proof}
Rewrite the defect exactly as
\begin{align*}
 \delta_{\alpha;n,m}^{\rm W}
 &=
 \bigl(q_{\alpha,R_{n+m}}-q_{\alpha,R_n}\bigr)
 +q_{\alpha,R_n}\bigl(1-q_{\alpha,R_m}\bigr).
\end{align*}
The first term has the limit
\eqref{eq:V122-exact-q-difference}.  Normalized characters have modulus
at most one, and \eqref{eq:V122-fixed-small} gives
\[
 \sqrt\alpha\,
 \left|q_{\alpha,R_n}(1-q_{\alpha,R_m})\right|
 =
 O_m(\alpha^{-1/2})\longrightarrow0.
\]
This proves \eqref{eq:V122-Wilson-defect-limit}.

By \eqref{eq:V122-correct-clock} and
\(c_k=(k^2+3k)/3\),
\[
 s_{\alpha_j}c_{n_j+m}\longrightarrow\rho^2.
\]
The limit in \eqref{eq:V122-Wilson-defect-limit} is nonzero, so taking
absolute values and multiplying by the weight proves
\eqref{eq:V122-Wilson-weighted-alpha-limit}.  Since
\(n_j/\sqrt{\alpha_j}\to\rho\), this also proves
\eqref{eq:V122-Wilson-weighted-n-limit}.

Finally,
\[
 \frac{1+c_m}{1+c_{n_j}}=O_m(n_j^{-2}),
\]
whereas \eqref{eq:V122-Wilson-weighted-n-limit} says that the left side
of \eqref{eq:V122-false-full-ratio} is asymptotic to a strictly positive
constant times \(n_j^{-1}\).  The asserted failure follows.
\end{proof}

\begin{corollary}[The V120 heat constant is the exact Wilson constant]
\label{cor:V122-heat-Wilson-constant}
Set \(u=\rho^2\).  Then the right side of
\eqref{eq:V122-Wilson-weighted-n-limit} is
\[
 2um(1+u)e^{-u},
\]
exactly the V120 heat constant in
\eqref{eq:V120-Cartan-limit}.
\end{corollary}

\begin{proof}
Substitute \(u=\rho^2\).
\end{proof}

\begin{corollary}[Decision of the V121 scalar route]
\label{cor:V122-V121-route-decided}
The sufficient full-ratio estimate
\eqref{eq:V121-required-Wilson-ratio} is false for the exact one-link
Wilson law.  Consequently the operator-hot scalar route described in
\Cref{prop:V121-scalar-decision} cannot close the persistent Cartan
subfamily on any transition ray on which the old \(ab\) capacity is
sharp and no further signed cancellation is used.
\end{corollary}

\begin{proof}
This is \Cref{thm:V122-Wilson-Cartan-obstruction} inserted into the
precisely scoped decision statement of
\Cref{prop:V121-scalar-decision}.
\end{proof}

\begin{firewall}[Scalar obstruction versus the complete coefficient]
\label{fire:V122-scalar-not-full}
The theorem is an exact counterexample to the printed Wilson
full-ratio scalar estimate; it is not yet a counterexample to the
complete signed V98 marked-tree coefficient.  The latter also contains
the old \(ab\) eigenvalue, collision factors, the \(ac\) covariance, and
the final resolvent.  V121 proved nonzero incidence and absence of a
forced Schur debit, but not a uniform lower bound or simultaneous
saturation for those remaining factors.  They must be evaluated before
one may reject or repair the full marked-tree route.
\end{firewall}

\section{V122 ledger and V123 mandate}
\label{sec:V122-ledger}

\begin{theorem}[Integrated V122 statement]
\label{thm:V122-integrated}
Preserving V1--V121 and explicitly superseding only the identified
factor-two conversions, V122 proves:
\begin{enumerate}[label=\textup{(V122.\arabic*)},leftmargin=7.2em]
\item the convention-free Laplace equipartition
      \eqref{eq:V122-equipartition};
\item the corrected clock
      \(s_\alpha=3/\alpha+O(\alpha^{-2})\);
\item the exact coherent-state identity
      \eqref{eq:V122-coherent-identity};
\item the direct transition and discrete-derivative limits
      \eqref{eq:V122-q-limit}--%
      \eqref{eq:V122-q-difference-limit};
\item the exact Gaussian collapse \(F(\rho)=e^{-\rho^2}\);
\item the exact Wilson square-root limits
      \eqref{eq:V122-Wilson-defect-limit}--%
      \eqref{eq:V122-Wilson-weighted-n-limit};
\item failure of the full endpoint-ratio Wilson scalar estimate; and
\item equality of the exact Wilson and heat leading Cartan constants.
\end{enumerate}
\end{theorem}

\begin{proof}
The items are
\Cref{lem:V122-equipartition,
thm:V122-correct-clock,
lem:V122-coherent-identity,
thm:V122-Wilson-transition,
lem:V122-Gaussian-collapse,
thm:V122-Wilson-Cartan-obstruction,
cor:V122-V121-route-decided,
cor:V122-heat-Wilson-constant}.
\end{proof}

\begin{openproblem}[V123 mandate: compute the complete persistent block]
\label{op:V123-mandate}
\begin{enumerate}
\item Instantiate one minimal V98 graph monomial on the disjoint
      \(E_{oa},E_{ab},E_{bc},E_{ac}\) banks of V121.
\item On its multiplicity-one Cartan path, compute before trace norm the
      old \(ab\) multiplier eigenvalue, the exact collision factors, the
      constant-\(B\) \(ac\) covariance, and the final resolvent.
\item Keep the signed coefficient assembled.  Determine whether its
      order-\(n^{-1}\) term from
      \eqref{eq:V122-Wilson-defect-limit} cancels identically.
\item If it does not cancel, quantify the remaining factors on a
      concrete nonzero input family and decide which intermediate MTA
      statement must be weakened or rerouted.
\item If it cancels, prove the cancellation as an algebraic identity and
      propagate the gained debit through the residual order
      \(\Xi_a>\Xi_c>C_{\rm row}\Xi_b\).
\item Use \eqref{eq:V122-correct-clock} in every future conversion
      between \(s_\alpha\) and \(\alpha^{-1}\).
\end{enumerate}
\end{openproblem}

\section{Firewall after V122}
\label{sec:V122-firewall}

\begin{firewall}[Current claim boundary]
\label{fire:V122-current-boundary}
V122 rigorously decides the growing-representation Wilson scalar
asymptotic and disproves the full-ratio scalar gate used by the existing
operator-hot route.  It does not prove simultaneous lower bounds for
the other factors in the complete physical coefficient, disprove every
possible signed marked-tree estimate, close the residual off-cut branch,
prove JREC, prove quartic or all-order MTA, construct continuum
Yang--Mills theory, or prove a positive spectral gap.  The full
Yang--Mills existence and mass gap remain open.
\end{firewall}

\section{Append-only record for V122}
\label{sec:V122-record}

V122 was appended to the complete V121 manuscript.  The exact V121
parent had \(81232\) newline-delimited source records, \(2825495\)
bytes, and SHA--256 digest
\[
 \texttt{3254638fefd7dcb5f53ff9a426f91b44ef276bd688c5e80c2d08c7f958a65221}.
\]
No V121 mathematical content was changed.  The correction in
\Cref{ass:V122-factor-two-supersession} is append-only and controls the
identified earlier normalization statements.  The sole final
\verb|\end{document}| is restored below.

% =============================================================================
% END APPENDED VERSION V122 MATERIAL
% =============================================================================

% =============================================================================
% BEGIN APPENDED VERSION V123 MATERIAL
% =============================================================================

\chapter{V123: Two Asymmetric Cartan Debits and the
Conditioned Persistent-Block Payment}
\label{ch:V123}

\section{Purpose}
\label{sec:V123-purpose}

V122 proves that one asymmetric large--small Cartan pair defect has only
a square-root endpoint debit.  The persistent V121 physical block,
however, contains two distinct signed pair multipliers: the old
\(ab\) multiplier and the peeled off-cut \(bc\) multiplier.  Their
coordinate banks are disjoint, although they share original row \(b\).
V118 keeps both multipliers exactly once before the reduced \(ac\)
hybrid operator.

V123 shows that this is the missing local mechanism on the conditioned
Cartan corridor.  A uniform coherent-state Lipschitz estimate gives
\[
 |d_{\alpha;k,\ell}^{\rm Car}|
 \lesssim
 \sqrt{\frac{1+c_\ell}{1+c_k}}
\]
when the large input \(k\) is at transition scale.  Applying it to both
pair banks and using the common \(b\)-row mass gives
\[
 |d_{ab}^{\rm Car}d_{bc}^{\rm Car}|
 \lesssim
 \frac{\Xi_b}{\Xi_c}
\]
whenever \(\Xi_a\ge\Xi_c\).  The two square roots therefore restore the
full endpoint ratio on every such conditioned double-Cartan block.  The
V122 limits show that this product debit is sharp and is not an
algebraic cancellation.

This pays the explicit persistent Cartan test which defeated the
one-multiplier V119 estimate.  It does not yet aggregate all non-Cartan
outputs or all representation shapes.

\section{A uniform one-link displacement moment}
\label{sec:V123-displacement}

\begin{lemma}[Wilson displacement]
\label{lem:V123-Wilson-displacement}
There are \(C<\infty\) and \(\alpha_0\) such that
\begin{equation}
 \mathbb E_{\nu_\alpha}\|U-I\|_{\rm op}
 \le\frac{C}{\sqrt\alpha}
 \qquad(\alpha\ge\alpha_0).
 \label{eq:V123-Wilson-displacement}
\end{equation}
\end{lemma}

\begin{proof}
On the compact group \(SU(3)\), the continuous quotient
\[
 \frac{\|U-I\|_{\rm op}^2}{V(U)}
\]
extends boundedly across \(U=I\): numerator and denominator have
positive-definite quadratic leading terms there, and \(V\) vanishes
nowhere else.  Thus
\[
 \|U-I\|_{\rm op}^2\le C V(U)
\]
globally.  Jensen's inequality and
\eqref{eq:V122-equipartition} give
\[
 \mathbb E\|U-I\|_{\rm op}
 \le
 C\bigl(\mathbb EV\bigr)^{1/2}
 \le C\alpha^{-1/2}.
\]
\end{proof}

\section{Uniform asymmetric Cartan Lipschitz law}
\label{sec:V123-one-pair}

For integers \(k,\ell\ge0\), define the signed weighted Cartan
coefficient
\begin{equation}
 d_{\alpha;k,\ell}^{\rm Car}
 :=
 \bigl(1+s_\alpha c_{k+\ell}\bigr)
 \left(
 q_{\alpha,R_{k+\ell}}
 -
 q_{\alpha,R_k}q_{\alpha,R_\ell}
 \right).
 \label{eq:V123-Cartan-coefficient}
\end{equation}
This is the V121 one-link weight convention.  If the V52 programme mass
uses \(1+c_{k+\ell}\) instead of \(c_{k+\ell}\), the actual physical
coefficient has the additional factor
\((1+s_\alpha(1+c_{k+\ell}))/(1+s_\alpha c_{k+\ell})\); it lies between
\(1\) and \(1+s_\alpha\) and does not change any estimate or limit below.

\begin{lemma}[Exact coherent Lipschitz difference]
\label{lem:V123-coherent-Lipschitz}
For all \(k,\ell\ge0\),
\begin{equation}
 \left|
 q_{\alpha,R_{k+\ell}}-q_{\alpha,R_k}
 \right|
 \le
 \ell\,
 \mathbb E_{\nu_\alpha}\|U-I\|_{\rm op}.
 \label{eq:V123-q-Lipschitz}
\end{equation}
\end{lemma}

\begin{proof}
For \(z_v(U):=\langle v,Uv\rangle\), the exact coherent-state identity
\eqref{eq:V122-coherent-identity} gives
\begin{align*}
 \frac{\chi_{R_{k+\ell}}(U)}{d_{k+\ell}}
 -
 \frac{\chi_{R_k}(U)}{d_k}
 &=
 \int_{S^5}z_v(U)^k\bigl(z_v(U)^\ell-1\bigr)\,\sigma(\dd v).
\end{align*}
Since \(U\) is unitary, \(|z_v(U)|\le1\), and hence
\[
 |z_v(U)^\ell-1|
 \le
 \ell|z_v(U)-1|
 \le
 \ell\|U-I\|_{\rm op}.
\]
Integrate first over \(v\) and then over the Wilson law.
\end{proof}

\begin{theorem}[Uniform square-root debit on the Cartan corridor]
\label{thm:V123-one-Cartan-debit}
Fix \(K<\infty\).  There are \(C_K<\infty\) and \(\alpha_K\) such that
whenever
\[
 1\le k\le K\sqrt\alpha,
 \qquad
 0\le\ell\le K\sqrt\alpha,
 \qquad
 \alpha\ge\alpha_K,
\]
one has
\begin{equation}
 \boxed{
 |d_{\alpha;k,\ell}^{\rm Car}|
 \le
 C_K\frac{\ell}{\sqrt\alpha}.}
 \label{eq:V123-Cartan-sqrt-alpha}
\end{equation}
If in addition \(k\ge K^{-1}\sqrt\alpha\), then
\begin{equation}
 \boxed{
 |d_{\alpha;k,\ell}^{\rm Car}|
 \le
 C_K
 \sqrt{\frac{1+c_\ell}{1+c_k}}.}
 \label{eq:V123-Cartan-sqrt-ratio}
\end{equation}
\end{theorem}

\begin{proof}
The exact decomposition
\begin{align}
 q_{\alpha,R_{k+\ell}}
 -q_{\alpha,R_k}q_{\alpha,R_\ell}
 &=
 \bigl(q_{\alpha,R_{k+\ell}}-q_{\alpha,R_k}\bigr)
 +q_{\alpha,R_k}\bigl(1-q_{\alpha,R_\ell}\bigr)
 \label{eq:V123-defect-decomposition}
\end{align}
and \(|q_{\alpha,R_k}|\le1\) give, by
\Cref{lem:V123-coherent-Lipschitz,lem:V123-Wilson-displacement},
\[
 |\delta_{\alpha;k,\ell}^{\rm Car}|
 \le
 C\frac{\ell}{\sqrt\alpha}
 +
 \bigl(1-q_{\alpha,R_\ell}\bigr).
\]
The uniform one-link gap
\eqref{eq:V30-one-link-upper-gap} and the corrected clock
\eqref{eq:V122-correct-clock} imply
\[
 0\le1-q_{\alpha,R_\ell}
 \le
 C\min\left\{\frac{c_\ell}{\alpha},1\right\}.
\]
For \(\ell\le K\sqrt\alpha\) and \(\ell\ge1\),
\[
 \min\left\{\frac{c_\ell}{\alpha},1\right\}
 \le C_K\frac{\ell}{\sqrt\alpha}.
\]
When \(\ell=0\), the defect vanishes exactly.  Also
\[
 1+s_\alpha c_{k+\ell}\le C_K
\]
under the two upper bounds on \(k,\ell\).  This proves
\eqref{eq:V123-Cartan-sqrt-alpha}.

If \(k\ge K^{-1}\sqrt\alpha\), then
\[
 \frac{\ell}{\sqrt\alpha}
 \le
 C_K\sqrt{\frac{1+c_\ell}{1+c_k}},
\]
which proves \eqref{eq:V123-Cartan-sqrt-ratio}.
\end{proof}

\begin{remark}[No fixed-small-label hypothesis]
\label{rem:V123-not-fixed-small}
The estimate permits \(\ell\) to grow proportionally to
\(\sqrt\alpha\).  It is therefore uniform across the interpolation
between the fixed-small-input V122 limit and a two-input transition
regime.  When the displayed square-root ratio is not small, the
inequality merely reduces to the universal V52 coefficient cap.
\end{remark}

\section{Two square roots on the shared-\(b\) Cartan block}
\label{sec:V123-two-pairs}

Take two disjoint coordinate banks sharing original row \(b\).  On the
old \(ab\) bank choose
\[
 R_{k_a}\otimes R_{\ell_1}\longrightarrow R_{k_a+\ell_1},
\]
and on the off-cut \(bc\) bank choose
\[
 R_{\ell_2}\otimes R_{k_c}\longrightarrow R_{\ell_2+k_c}.
\]
Both arrows denote the multiplicity-one Cartan component.  Define the
local masses
\begin{align}
 X_a&:=1+c_{k_a},
 &
 X_c&:=1+c_{k_c},
 \label{eq:V123-large-local-masses}\\
 X_b&:=(1+c_{\ell_1})+(1+c_{\ell_2}).
 \label{eq:V123-middle-local-mass}
\end{align}

\begin{theorem}[Two conditioned Cartan debits give the full endpoint
ratio]
\label{thm:V123-two-Cartan-debits}
Fix \(K<\infty\).  Suppose
\begin{equation}
 K^{-1}\sqrt\alpha
 \le k_a,k_c\le K\sqrt\alpha,
 \qquad
 0\le\ell_1,\ell_2\le K\sqrt\alpha,
 \qquad
 X_a\ge X_c.
 \label{eq:V123-double-corridor}
\end{equation}
Then, for all sufficiently large \(\alpha\),
\begin{equation}
 \boxed{
 |d_{\alpha;k_a,\ell_1}^{\rm Car}
   d_{\alpha;k_c,\ell_2}^{\rm Car}|
 \le
 C_K\frac{X_b}{X_c}.}
 \label{eq:V123-two-debit-ratio}
\end{equation}
The same conclusion holds for the literal V52 one-link weights using
\(1+c_{\rm out}\), after changing \(C_K\).
\end{theorem}

\begin{proof}
Apply \eqref{eq:V123-Cartan-sqrt-ratio} twice:
\begin{align*}
 |d_{\alpha;k_a,\ell_1}^{\rm Car}
   d_{\alpha;k_c,\ell_2}^{\rm Car}|
 &\le
 C_K
 \sqrt{
 \frac{(1+c_{\ell_1})(1+c_{\ell_2})}
      {X_aX_c}}.
\end{align*}
The arithmetic--geometric mean inequality gives
\[
 \sqrt{(1+c_{\ell_1})(1+c_{\ell_2})}
 \le\frac{X_b}{2}.
\]
Since \(X_a\ge X_c\), one has
\(\sqrt{X_aX_c}\ge X_c\).  Combining the last two displays proves
\eqref{eq:V123-two-debit-ratio}.  The ratio between the two weight
conventions is bounded uniformly by \(1+s_\alpha\), as observed after
\eqref{eq:V123-Cartan-coefficient}.
\end{proof}

\begin{corollary}[Translation to the V116 residual masses]
\label{cor:V123-global-mass-translation}
Assume the two displayed banks occur in a V116 residual block, their
large labels carry fixed positive fractions of the \(a\)- and \(c\)-row
masses, \(X_a\ge X_c\), and the residual ordering
\[
 \Xi_a>\Xi_c>C_{\rm row}\Xi_b
\]
holds.  On the transition corridor
\eqref{eq:V123-double-corridor},
\begin{equation}
 |d_{ab}^{\rm Car}d_{bc}^{\rm Car}|
 \le
 C_{K,\delta}\frac{\Xi_b}{\Xi_c}.
 \label{eq:V123-global-full-ratio}
\end{equation}
\end{corollary}

\begin{proof}
The middle-row mass \(\Xi_b\) dominates the two local contributions in
\eqref{eq:V123-middle-local-mass}.  The heavy off-cut hypothesis makes
\(\Xi_c\le C_\delta X_c\).  Insert these comparisons and the stated
orientation \(X_a\ge X_c\) in
\eqref{eq:V123-two-debit-ratio}.
\end{proof}

\begin{firewall}[The local-to-global hypotheses are explicit]
\label{fire:V123-mass-translation-scope}
\Cref{cor:V123-global-mass-translation} is not asserted for an arbitrary
outer label merely because its aggregate row masses obey the residual
ordering.  It requires the selected Cartan labels themselves to carry
the stated fixed fractions.  When the \(a\)-mass is spread over other
coordinates or the \(c\)-mass is spread over several off-cut labels, an
aggregation theorem is still required.
\end{firewall}

\section{Sharpness and the absence of pair-level cancellation}
\label{sec:V123-sharpness}

\begin{theorem}[The two-debit order is sharp]
\label{thm:V123-two-debit-sharp}
Fix integers \(r,m\ge1\).  If
\[
 \frac{k_a}{\sqrt\alpha}\longrightarrow\sigma>0,
 \qquad
 \frac{k_c}{\sqrt\alpha}\longrightarrow\rho>0,
\]
then
\begin{align}
 \alpha\,
 d_{\alpha;k_a,r}^{\rm Car}
 d_{\alpha;k_c,m}^{\rm Car}
 \longrightarrow{}&
 4rm\,\sigma\rho
 (1+\sigma^2)(1+\rho^2)
 e^{-(\sigma^2+\rho^2)}
 >0.
 \label{eq:V123-sharp-product-limit}
\end{align}
Thus the two signed Cartan defects do not cancel each other.  Their
product is of order \(\alpha^{-1}\), the same order as
\(X_b/X_c\) for fixed \(r,m\).
\end{theorem}

\begin{proof}
The signed form of \eqref{eq:V122-Wilson-defect-limit}, together with
the limiting output weight, gives
\begin{align*}
 \sqrt\alpha\,d_{\alpha;k_a,r}^{\rm Car}
 &\longrightarrow
 -2r\sigma(1+\sigma^2)e^{-\sigma^2},\\
 \sqrt\alpha\,d_{\alpha;k_c,m}^{\rm Car}
 &\longrightarrow
 -2m\rho(1+\rho^2)e^{-\rho^2}.
\end{align*}
Multiply.  Both leading factors are negative, so the product limit is
the strictly positive number in
\eqref{eq:V123-sharp-product-limit}.  Finally \(X_b=O_{r,m}(1)\) and
\(X_c\asymp_\rho\alpha\).
\end{proof}

\begin{assessment}[What the sharp limit says]
\label{ass:V123-sharp-meaning}
V122 showed that replacing the off-cut defect alone by a full endpoint
ratio is false.  V123 shows why that failure is not yet fatal to the
physical block: the old pair defect supplies the second square root.
The improvement is multiplicative and label-conditioned, not a
cancellation and not an aggregate Schur tax.
\end{assessment}

\section{Insertion into the constant-\(B\) physical block}
\label{sec:V123-physical-insertion}

Let \(j\) be a complete V98 outer block refining both Cartan outputs
above.  Let \(\mathbf R_{ac,j}^{1,\lhd}\) denote the complete reduced
\(ac\) row operator on that block, with its remaining collision,
spectator, recoupling, and final-resolvent factors retained as in the
V118 physical construction.

\begin{theorem}[Exact conditioned persistent-block factorization]
\label{thm:V123-conditioned-block}
On the V120 constant-reduced-\(b\) specialization, the \(ab\) hybrid row
vanishes and the surviving Cartan compression of the physical block is
\begin{equation}
 \boxed{
 L_{ac,j}
 =
 d_{ab,j}^{\rm Car}\,
 d_{bc,j}^{\rm Car}\,
 \mathbf R_{ac,j}^{1,\lhd}.}
 \label{eq:V123-conditioned-block-factor}
\end{equation}
Consequently, under \eqref{eq:V123-double-corridor},
\begin{equation}
 \|L_{ac,j}\|_{\rm op}
 \le
 C_K\frac{X_b}{X_c}
 \|\mathbf R_{ac,j}^{1,\lhd}\|_{\rm op}.
 \label{eq:V123-conditioned-block-bound}
\end{equation}
If the V118 row estimate is inserted, then
\begin{equation}
 \|L_{ac,j}\|_{\rm op}
 \le
 C_K s_\alpha H_{ac}^{1}
 \frac{X_b}{X_c}.
 \label{eq:V123-conditioned-row-payment}
\end{equation}
\end{theorem}

\begin{proof}
The old \(ab\) pair bank and off-cut \(bc\) bank occupy disjoint
coordinate factors, and the complete outer label \(j\) refines both
pair outputs by \Cref{lem:V118-physical-joint-refinement}.  Therefore
the exact blockwise factorization used in the proof of
\Cref{thm:V118-physical-joint-completion} gives one copy of each signed
pair eigenvalue multiplying the reduced row operator.  By
\Cref{prop:V120-constant-b}, the other row is zero; hence no second
signed summand remains on this specialization.  This proves
\eqref{eq:V123-conditioned-block-factor}.

Take absolute values and apply
\eqref{eq:V123-two-debit-ratio} to prove
\eqref{eq:V123-conditioned-block-bound}.  Finally use
\eqref{eq:V118-reduced-row-bounds}.
\end{proof}

\begin{corollary}[The explicit persistent Cartan test is paid]
\label{cor:V123-persistent-Cartan-paid}
Under the hypotheses of
\Cref{cor:V123-global-mass-translation}, the V121 constant-\(B\)
double-Cartan block has the missing \(c\)-endpoint ratio:
\begin{equation}
 \|L_{ac,j}\|_{\rm op}
 \le
 C s_\alpha H_{ac}^{1}
 \frac{\Xi_b}{\Xi_c}.
 \label{eq:V123-persistent-payment}
\end{equation}
Thus the representation family which disproves the one-multiplier
V119 scalar gate is not a counterexample to the complete two-pair
physical coefficient.
\end{corollary}

\begin{proof}
Combine
\eqref{eq:V123-conditioned-row-payment} with
\eqref{eq:V123-global-full-ratio}.
\end{proof}

\begin{firewall}[One block is not the summed carrier]
\label{fire:V123-block-not-aggregate}
Equation \eqref{eq:V123-persistent-payment} is a complete signed
coefficient estimate on each conditioned double-Cartan outer block.
It does not by itself sum all Cartan labels, control non-Cartan outputs,
or show that the labelwise ratios can be replaced by one aggregate row
ratio without multiplicity loss.  Those are the next positive-carrier
tasks.
\end{firewall}

\section{V123 ledger and V124 mandate}
\label{sec:V123-ledger}

\begin{theorem}[Integrated V123 statement]
\label{thm:V123-integrated}
Preserving V1--V122, V123 proves:
\begin{enumerate}[label=\textup{(V123.\arabic*)},leftmargin=7.2em]
\item the uniform Wilson displacement estimate
      \eqref{eq:V123-Wilson-displacement};
\item the exact coherent Lipschitz difference
      \eqref{eq:V123-q-Lipschitz};
\item the uniform one-pair Cartan square-root debit
      \eqref{eq:V123-Cartan-sqrt-ratio};
\item multiplication of the two shared-\(b\) debits into the full
      endpoint ratio \eqref{eq:V123-two-debit-ratio};
\item its scoped translation to the residual row masses;
\item the nonzero sharp product limit
      \eqref{eq:V123-sharp-product-limit};
\item exact insertion into the constant-\(B\) physical block; and
\item payment of the explicit V121 persistent Cartan test.
\end{enumerate}
\end{theorem}

\begin{proof}
The items are
\Cref{lem:V123-Wilson-displacement,
lem:V123-coherent-Lipschitz,
thm:V123-one-Cartan-debit,
thm:V123-two-Cartan-debits,
cor:V123-global-mass-translation,
thm:V123-two-debit-sharp,
thm:V123-conditioned-block,
cor:V123-persistent-Cartan-paid}.
\end{proof}

\begin{openproblem}[V124 mandate: aggregate beyond the Cartan corridor]
\label{op:V124-mandate}
\begin{enumerate}
\item Replace the symmetric-power Cartan labels by arbitrary
      \(SU(3)\) highest weights and arbitrary multiplicity-resolved
      outputs of both pair banks.
\item Prove a joint positive-carrier inequality which retains the
      product of the two labelwise endpoint debits before summing outer
      labels.
\item Use the common middle-row mass to sum the product by an
      arithmetic--geometric mean, without a channel or multiplicity
      count.
\item Separate low, transition, and operator-hot large-output sectors;
      use V123 only on the transition sector where its hypotheses hold.
\item Decide the complementary configurations in which one selected
      label does not carry a fixed fraction of its row mass.  Return
      that dispersed mass to the coordinate-bank ledger rather than
      replacing it by a maximal label.
\item Preserve the exact V110 row interference until the common
      positive carrier has been formed.
\end{enumerate}
\end{openproblem}

\section{Firewall after V123}
\label{sec:V123-firewall}

\begin{firewall}[Current claim boundary]
\label{fire:V123-current-boundary}
V123 pays the explicit conditioned double-Cartan block which survived
V121--V122.  It does not aggregate the full Cartan family, handle all
highest-weight directions and non-Cartan outputs, close the entire
off-cut carrier, prove JREC, prove quartic or all-order MTA, construct
continuum Yang--Mills theory, or prove a positive spectral gap.  The
full Yang--Mills existence and mass gap remain open.
\end{firewall}

\section{Append-only record for V123}
\label{sec:V123-record}

V123 was appended to the complete V122 manuscript.  The exact V122
parent had \(81977\) newline-delimited source records, \(2848799\)
bytes, and SHA--256 digest
\[
 \texttt{38199115bfe3f34f5d81dac8ba4f5b466d34f261c6c0e7e18f73da030490e88d}.
\]
No V122 mathematical content was changed.  The sole final
\verb|\end{document}| is restored below.

% =============================================================================
% END APPENDED VERSION V123 MATERIAL
% =============================================================================

% =============================================================================
% BEGIN APPENDED VERSION V124 MATERIAL
% =============================================================================

\chapter{V124: Arbitrary Highest-Weight Cartan Debits and the
Uniform Double-Cartan Corridor}
\label{ch:V124}

\section{Purpose}
\label{sec:V124-purpose}

V123 paid the conditioned persistent block on symmetric-power Cartan
rays.  V124 removes that ray restriction.  For arbitrary irreducible
\(SU(3)\) highest weights \(\lambda,\mu\), normalized characters admit
an exact coherent-orbit representation, and the highest matrix
coefficient of the Cartan component \(\lambda+\mu\) is the product of
the two highest matrix coefficients.  This gives an exact difference
formula before averaging over the Wilson law.

The resulting estimate is uniform in all highest-weight directions:
when \(\lambda\) is at transition scale and \(\mu\) is no larger than
transition scale,
\[
 |d_{\alpha;\lambda,\mu}^{\rm Car}|
 \lesssim
 \sqrt{\frac{1+C_2(\mu)}{1+C_2(\lambda)}}.
\]
Two such factors on the old \(ab\) and off-cut \(bc\) banks again
multiply to the full middle-row endpoint ratio.  Thus every
double-Cartan highest-weight direction in the fixed transition corridor
is paid.  The unresolved carrier is reduced to non-Cartan outputs,
dispersed large-row mass, and aggregation across those sectors.

\section{Highest-weight notation and coherent orbits}
\label{sec:V124-coherent}

Write an \(SU(3)\) dominant integral weight as
\[
 \lambda=(p,q),
 \qquad p,q\in\mathbb N_0,
 \qquad
 |\lambda|_1:=p+q.
\]
Let \((\pi_\lambda,H_\lambda)\) be the corresponding unitary
irreducible representation, let \(v_\lambda\) be a unit highest-weight
vector, and put
\[
 c_\lambda:=C_2(\lambda)
 =
 \frac{p^2+q^2+pq+3p+3q}{3}.
\]
There are numerical constants \(0<c<C<\infty\) such that
\begin{equation}
 c(1+|\lambda|_1^2)
 \le
 1+c_\lambda
 \le
 C(1+|\lambda|_1^2)
 \label{eq:V124-height-Casimir}
\end{equation}
for every \(\lambda\).

For \(g,U\in SU(3)\), define the coherent matrix coefficient
\begin{equation}
 z_\lambda(g,U)
 :=
 \left\langle
 v_\lambda,
 \pi_\lambda(g^{-1}Ug)v_\lambda
 \right\rangle.
 \label{eq:V124-coherent-coefficient}
\end{equation}

\begin{lemma}[Exact coherent-orbit character identity]
\label{lem:V124-coherent-character}
For every dominant weight \(\lambda\),
\begin{equation}
 \boxed{
 \frac{\chi_\lambda(U)}{d_\lambda}
 =
 \int_{SU(3)}z_\lambda(g,U)\,\dd g.}
 \label{eq:V124-coherent-character}
\end{equation}
\end{lemma}

\begin{proof}
Haar twirling of the rank-one highest-weight projection gives, by
Schur's lemma,
\[
 \int_{SU(3)}
 \pi_\lambda(g)
 |v_\lambda\rangle\langle v_\lambda|
 \pi_\lambda(g)^{-1}\,\dd g
 =
 \frac1{d_\lambda}I_{H_\lambda}.
\]
Multiply by \(\pi_\lambda(U)\), take the trace, and use cyclicity.
\end{proof}

\begin{lemma}[Exact Cartan coherent factorization]
\label{lem:V124-Cartan-coherent-factor}
For dominant weights \(\lambda,\mu\), choose the Cartan copy
\(H_{\lambda+\mu}\subset H_\lambda\otimes H_\mu\) generated by
\(v_\lambda\otimes v_\mu\).  Then
\begin{equation}
 \boxed{
 z_{\lambda+\mu}(g,U)
 =
 z_\lambda(g,U)z_\mu(g,U).}
 \label{eq:V124-Cartan-coherent-factor}
\end{equation}
Consequently,
\begin{equation}
 \frac{\chi_{\lambda+\mu}(U)}{d_{\lambda+\mu}}
 -
 \frac{\chi_\lambda(U)}{d_\lambda}
 =
 \int_{SU(3)}
 z_\lambda(g,U)\bigl(z_\mu(g,U)-1\bigr)\,\dd g.
 \label{eq:V124-character-difference}
\end{equation}
\end{lemma}

\begin{proof}
The tensor \(v_\lambda\otimes v_\mu\) has weight
\(\lambda+\mu\), is killed by every raising operator, and generates the
multiplicity-one Cartan component.  Its matrix coefficient under
\(\pi_\lambda\otimes\pi_\mu\) is the product of the two matrix
coefficients.  Apply
\Cref{lem:V124-coherent-character} to \(\lambda+\mu\) and to
\(\lambda\), using the same Haar variable \(g\), and subtract.
\end{proof}

\begin{lemma}[Representation displacement by highest-weight height]
\label{lem:V124-representation-displacement}
For every dominant weight \(\mu\) and every \(U\in SU(3)\),
\begin{equation}
 \|\pi_\mu(U)-I\|_{\rm op}
 \le
 |\mu|_1\|U-I\|_{\rm op}.
 \label{eq:V124-representation-displacement}
\end{equation}
In particular,
\begin{equation}
 |z_\mu(g,U)-1|
 \le
 |\mu|_1\|U-I\|_{\rm op}.
 \label{eq:V124-coherent-displacement}
\end{equation}
\end{lemma}

\begin{proof}
For \(\mu=(r,s)\), the representation \(H_\mu\) occurs as an invariant
subspace of
\[
 F^{\otimes r}\otimes\overline F^{\otimes s}.
\]
The telescoping identity for a tensor product of unitary matrices gives
\[
 \left\|
 U^{\otimes r}\otimes\overline U^{\otimes s}-I
 \right\|_{\rm op}
 \le
 r\|U-I\|_{\rm op}
 +s\|\overline U-I\|_{\rm op}
 =
 (r+s)\|U-I\|_{\rm op}.
\]
Restriction to an invariant subspace cannot increase the norm.  The
matrix-coefficient bound and conjugation invariance of the operator norm
give \eqref{eq:V124-coherent-displacement}.
\end{proof}

\section{Uniform Cartan pair estimate in every direction}
\label{sec:V124-one-pair}

For dominant weights \(\lambda,\mu\), put
\begin{equation}
 d_{\alpha;\lambda,\mu}^{\rm Car}
 :=
 \bigl(1+s_\alpha c_{\lambda+\mu}\bigr)
 \left(
 q_{\alpha,\lambda+\mu}
 -
 q_{\alpha,\lambda}q_{\alpha,\mu}
 \right).
 \label{eq:V124-general-Cartan-coefficient}
\end{equation}
As in V123, replacing \(c_{\lambda+\mu}\) by the literal one-link
programme mass changes the coefficient by a uniformly harmless factor.

\begin{lemma}[General coherent Cartan Lipschitz law]
\label{lem:V124-general-Lipschitz}
For all dominant \(\lambda,\mu\),
\begin{equation}
 \left|
 q_{\alpha,\lambda+\mu}-q_{\alpha,\lambda}
 \right|
 \le
 |\mu|_1\,
 \mathbb E_{\nu_\alpha}\|U-I\|_{\rm op}.
 \label{eq:V124-general-q-Lipschitz}
\end{equation}
\end{lemma}

\begin{proof}
In \eqref{eq:V124-character-difference},
\(|z_\lambda(g,U)|\le1\).  Apply
\eqref{eq:V124-coherent-displacement}, integrate over \(g\), and then
average over the Wilson law.
\end{proof}

\begin{theorem}[Arbitrary-direction Cartan square-root debit]
\label{thm:V124-general-Cartan-debit}
Fix \(K<\infty\).  There are \(C_K<\infty\) and \(\alpha_K\) such that,
for \(\alpha\ge\alpha_K\),
\begin{equation}
 K^{-1}\sqrt\alpha
 \le|\lambda|_1\le K\sqrt\alpha,
 \qquad
 |\mu|_1\le K\sqrt\alpha
 \label{eq:V124-general-corridor}
\end{equation}
implies
\begin{align}
 |d_{\alpha;\lambda,\mu}^{\rm Car}|
 &\le
 C_K\frac{|\mu|_1}{\sqrt\alpha},
 \label{eq:V124-general-height-debit}\\
 \boxed{
 |d_{\alpha;\lambda,\mu}^{\rm Car}|
 \le
 C_K
 \sqrt{\frac{1+c_\mu}{1+c_\lambda}}.}
 \label{eq:V124-general-Casimir-debit}
\end{align}
\end{theorem}

\begin{proof}
Decompose exactly
\[
 q_{\alpha,\lambda+\mu}
 -q_{\alpha,\lambda}q_{\alpha,\mu}
 =
 (q_{\alpha,\lambda+\mu}-q_{\alpha,\lambda})
 +q_{\alpha,\lambda}(1-q_{\alpha,\mu}).
\]
The first term is bounded by
\Cref{lem:V124-general-Lipschitz,lem:V123-Wilson-displacement}.
For the second, use \(|q_{\alpha,\lambda}|\le1\),
\eqref{eq:V30-one-link-upper-gap}, and
\eqref{eq:V122-correct-clock}:
\[
 0\le1-q_{\alpha,\mu}
 \le
 C\min\left\{\frac{c_\mu}{\alpha},1\right\}
 \le
 C_K\frac{|\mu|_1}{\sqrt\alpha}.
\]
The last inequality is trivial for \(\mu=0\), and otherwise follows
from \eqref{eq:V124-height-Casimir} and
\(|\mu|_1\le K\sqrt\alpha\).
The output weight in
\eqref{eq:V124-general-Cartan-coefficient} is bounded by \(C_K\)
because
\(|\lambda+\mu|_1\le2K\sqrt\alpha\).  This proves
\eqref{eq:V124-general-height-debit}.  Finally
\eqref{eq:V124-height-Casimir} and the lower bound on
\(|\lambda|_1\) give
\[
 \frac{|\mu|_1}{\sqrt\alpha}
 \le
 C_K\sqrt{\frac{1+c_\mu}{1+c_\lambda}},
\]
proving \eqref{eq:V124-general-Casimir-debit}.
\end{proof}

\section{The arbitrary-direction double-Cartan ratio}
\label{sec:V124-double-Cartan}

On the old \(ab\) bank take dominant input weights
\((\lambda_a,\mu_1)\), and on the off-cut \(bc\) bank take
\((\mu_2,\lambda_c)\).  Select the Cartan outputs
\(\lambda_a+\mu_1\) and \(\mu_2+\lambda_c\).  Put
\begin{align}
 X_a&:=1+c_{\lambda_a},
 &
 X_c&:=1+c_{\lambda_c},
 \label{eq:V124-general-large-masses}\\
 X_b&:=(1+c_{\mu_1})+(1+c_{\mu_2}).
 \label{eq:V124-general-middle-mass}
\end{align}

\begin{theorem}[Uniform double-Cartan endpoint ratio]
\label{thm:V124-general-double-Cartan}
Fix \(K<\infty\).  Suppose
\[
 K^{-1}\sqrt\alpha
 \le
 |\lambda_a|_1,|\lambda_c|_1
 \le
 K\sqrt\alpha,
 \qquad
 |\mu_1|_1,|\mu_2|_1\le K\sqrt\alpha,
\]
and \(X_a\ge X_c\).  Then
\begin{equation}
 \boxed{
 |d_{\alpha;\lambda_a,\mu_1}^{\rm Car}
   d_{\alpha;\lambda_c,\mu_2}^{\rm Car}|
 \le
 C_K\frac{X_b}{X_c}.}
 \label{eq:V124-general-double-ratio}
\end{equation}
\end{theorem}

\begin{proof}
Twice apply \eqref{eq:V124-general-Casimir-debit} and then use the
arithmetic--geometric mean:
\begin{align*}
 |d_{\alpha;\lambda_a,\mu_1}^{\rm Car}
   d_{\alpha;\lambda_c,\mu_2}^{\rm Car}|
 &\le
 C_K
 \frac{\sqrt{(1+c_{\mu_1})(1+c_{\mu_2})}}
      {\sqrt{X_aX_c}}\\
 &\le
 C_K
 \frac{(1+c_{\mu_1})+(1+c_{\mu_2})}
      {2X_c}.
\end{align*}
This is \eqref{eq:V124-general-double-ratio}.
\end{proof}

\begin{remark}[Shared-row entanglement causes no loss]
\label{rem:V124-shared-row}
The proof is labelwise and multiplies the two actual scalar pair
coefficients before any product-state capacity is taken.  The
middle-row vector may be entangled between its two coordinate factors.
No factorization of that vector is used; the common-row Cartan
compatibility established in V120 only ensures that this labelwise
product is physically present.
\end{remark}

\section{Aggregation of the complete double-Cartan outer subfamily}
\label{sec:V124-Cartan-aggregation}

Fix the two pair-input labels above and let
\(\mathfrak J_{\rm Car,Car}\) be all complete V98 outer labels which
refine both Cartan pair outputs and satisfy the V120
constant-reduced-\(b\) specialization.  These labels may differ in
their remaining final output and multiplicity data.  Let \(P_j\) be
their mutually orthogonal physical projections.

\begin{theorem}[No-count aggregation of the double-Cartan family]
\label{thm:V124-Cartan-aggregate}
Assume the corridor and orientation hypotheses of
\Cref{thm:V124-general-double-Cartan}.  If the complete reduced
\(ac\)-row operators satisfy
\begin{equation}
 \|\mathbf R_{ac,j}^{1,\lhd}\|_{\rm op}
 \le
 Cs_\alpha H_{ac}^{1}
 \qquad(j\in\mathfrak J_{\rm Car,Car}),
 \label{eq:V124-uniform-reduced-row}
\end{equation}
then their physical absolute carrier obeys
\begin{align}
 \mathbf Q_{\rm Car,Car}
 &:=
 \sum_{j\in\mathfrak J_{\rm Car,Car}}
 |V_jL_{ac,j}V_j^*|
 \notag\\
 &\preccurlyeq
 C_Ks_\alpha H_{ac}^{1}
 \frac{X_b}{X_c}
 \sum_{j\in\mathfrak J_{\rm Car,Car}}P_j
 \notag\\
 &\preccurlyeq
 C_Ks_\alpha H_{ac}^{1}
 \frac{X_b}{X_c}I.
 \label{eq:V124-Cartan-aggregate-order}
\end{align}
Consequently,
\begin{equation}
 \boxed{
 \|\mathbf Q_{\rm Car,Car}\|_{\rm rows,prod}
 \le
 C_Ks_\alpha H_{ac}^{1}\frac{X_b}{X_c}.}
 \label{eq:V124-Cartan-aggregate-capacity}
\end{equation}
There is no final-output or multiplicity count.
\end{theorem}

\begin{proof}
On each complete block, the exact V118 factorization and the vanishing
of the other hybrid row give
\[
 L_{ac,j}
 =
 d_{\alpha;\lambda_a,\mu_1}^{\rm Car}
 d_{\alpha;\lambda_c,\mu_2}^{\rm Car}
 \mathbf R_{ac,j}^{1,\lhd}.
\]
The pair scalars are independent of the later final-output and copy
labels.  Apply
\eqref{eq:V124-general-double-ratio} and
\eqref{eq:V124-uniform-reduced-row} on each orthogonal range \(P_j\).
Absolute values add across the V98 blocks by
\Cref{thm:V98-exact-absolute-carrier}, while
\(\sum_jP_j\preccurlyeq I\).  This proves the operator order.
Taking expectations in original-row product states proves
\eqref{eq:V124-Cartan-aggregate-capacity}.
\end{proof}

\begin{corollary}[Removal of the transition double-Cartan sector]
\label{cor:V124-Cartan-sector-paid}
Whenever the local masses in
\eqref{eq:V124-general-large-masses} carry fixed positive fractions of
the corresponding residual row masses, \(X_a\ge X_c\), and
\(\Xi_a>\Xi_c>C_{\rm row}\Xi_b\), the whole outer subfamily
\(\mathfrak J_{\rm Car,Car}\) satisfies
\begin{equation}
 \|\mathbf Q_{\rm Car,Car}\|_{\rm rows,prod}
 \le
 C s_\alpha H_{ac}^{1}\frac{\Xi_b}{\Xi_c}.
 \label{eq:V124-Cartan-sector-global}
\end{equation}
It may therefore be moved to the paid part of the V98 outer partition.
\end{corollary}

\begin{proof}
As in \Cref{cor:V123-global-mass-translation},
\(X_b\le\Xi_b\) and heavy off-cut incidence gives
\(\Xi_c\le C X_c\).  Insert these comparisons in
\eqref{eq:V124-Cartan-aggregate-capacity}.  The move is a partition of
orthogonal outer labels, so it introduces no count.
\end{proof}

\begin{firewall}[Exactly which Cartan sector is removed]
\label{fire:V124-Cartan-removal-scope}
The removal applies to all highest-weight directions and to every
remaining final/copy label after the two pair outputs have both been
fixed to their Cartan components.  It still assumes that the two large
input labels carry the relevant row masses and lie in a fixed
transition corridor.  It does not include a pair output below the
Cartan highest weight, a large row whose mass is dispersed among many
coordinates, or a sector outside the corridor.
\end{firewall}

\section{V124 ledger and V125 mandate}
\label{sec:V124-ledger}

\begin{theorem}[Integrated V124 statement]
\label{thm:V124-integrated}
Preserving V1--V123, V124 proves:
\begin{enumerate}[label=\textup{(V124.\arabic*)},leftmargin=7.2em]
\item the exact coherent-orbit character identity
      \eqref{eq:V124-coherent-character};
\item exact factorization of the Cartan coherent coefficient
      \eqref{eq:V124-Cartan-coherent-factor};
\item the representation displacement estimate
      \eqref{eq:V124-representation-displacement};
\item the arbitrary-direction square-root debit
      \eqref{eq:V124-general-Casimir-debit};
\item the uniform double-Cartan endpoint ratio
      \eqref{eq:V124-general-double-ratio};
\item positive no-count aggregation over all remaining physical
      final and copy labels; and
\item payment of the complete conditioned transition
      double-Cartan sector.
\end{enumerate}
\end{theorem}

\begin{proof}
The items are
\Cref{lem:V124-coherent-character,
lem:V124-Cartan-coherent-factor,
lem:V124-representation-displacement,
thm:V124-general-Cartan-debit,
thm:V124-general-double-Cartan,
thm:V124-Cartan-aggregate,
cor:V124-Cartan-sector-paid}.
\end{proof}

\begin{openproblem}[V125 mandate: audit and non-Cartan depth]
\label{op:V125-mandate}
\begin{enumerate}
\item Perform the mandatory five-version audit of the current SM and NS
      companion notes before choosing the next estimate.
\item For an output
      \(\nu<\lambda+\mu\), introduce its lowering depth from the Cartan
      face and retain the exact multiplicity-resolved projection.
\item Prove either a square-root endpoint debit for its signed Wilson
      coefficient or a complementary product-capacity debit for its
      projection.  Do not assume unit Cartan visibility.
\item Aggregate depth layers by their invariant projectors before
      taking product-state capacity; do not count Littlewood--Richardson
      copies.
\item Split dispersed large-row mass at the coordinate-bank level and
      return any heavy subbank to the established selected-heavy
      alternatives.
\item Combine the non-Cartan and dispersed sectors with
      \eqref{eq:V124-Cartan-sector-global}.
\end{enumerate}
\end{openproblem}

\section{Firewall after V124}
\label{sec:V124-firewall}

\begin{firewall}[Current claim boundary]
\label{fire:V124-current-boundary}
V124 removes the complete conditioned transition double-Cartan sector
in every \(SU(3)\) highest-weight direction.  It does not control
non-Cartan pair outputs, dispersed large-row mass, or sectors outside
the transition corridor; therefore it does not close the full off-cut
carrier, prove JREC, prove quartic or all-order MTA, construct continuum
Yang--Mills theory, or prove a positive spectral gap.  The full
Yang--Mills existence and mass gap remain open.
\end{firewall}

\section{Append-only record for V124}
\label{sec:V124-record}

V124 was appended to the complete V123 manuscript.  The exact V123
parent had \(82537\) newline-delimited source records, \(2866408\)
bytes, and SHA--256 digest
\[
 \texttt{ccf6d66bb247d715b5148295463657211455be773ca8d61fe51aa66c679fb94c}.
\]
No V123 mathematical content was changed.  The sole final
\verb|\end{document}| is restored below.

% =============================================================================
% END APPENDED VERSION V124 MATERIAL
% =============================================================================

% =============================================================================
% BEGIN APPENDED VERSION V125 MATERIAL
% =============================================================================

\chapter{V125: Companion Audit and the Universal
Output-Compression Debit}
\label{ch:V125}

\section{Purpose}
\label{sec:V125-purpose}

V124 removed every double-Cartan highest-weight direction in the
conditioned transition corridor.  V125 performs the required
five-version companion audit and then removes the Cartan restriction
from the scalar estimate.

The key step is upstream of any Littlewood--Richardson depth analysis.
If \(H_\nu\) is any irreducible summand of
\(H_\lambda\otimes H_\mu\), then the difference
\(q_{\alpha,\nu}-q_{\alpha,\lambda}\) is the compression to \(H_\nu\)
of
\[
 \mathbb E_{\nu_\alpha}
 \left[
 \pi_\lambda(U)\otimes(\pi_\mu(U)-I)
 \right].
\]
Its norm is therefore controlled directly by the displacement of the
small input representation.  No character asymptotic, output
multiplicity count, or Cartan visibility is involved.  On a
single-coordinate transition bank this gives the same square-root
endpoint debit for \emph{every} output.  Applying it to both physical
pair multipliers pays the complete concentrated transition output
family.

\section{Mandatory V125 companion audit}
\label{sec:V125-audit}

The read-only audit on 5 September 2026 found:
\begin{center}
\begin{tabular}{|p{0.50\textwidth}|r|p{0.30\textwidth}|}
\hline
Companion file & Bytes & SHA--256\\
\hline
\texttt{Renewal\_NS\_Stability\_Research\_Notes\_V31.tex}
& \(887537\) &
\texttt{f1121b62238ad5491bb7cf03635ea9934942edf573ed8faaa9ee79e1d38a6696}\\
\hline
\texttt{Renewal\_SM\_Einstein\_Continuum\_Research\_Notes\_V141.tex}
& \(3350838\) &
\texttt{0b8184707a59b9ce9256adb6ac275102d8340b562d4893e4ea04ebade088d45e}\\
\hline
\end{tabular}
\end{center}
The current line counter gives \(23052\) and \(91992\) source lines,
respectively.  The NS file is byte-identical to the V120 audit; the
one-line difference from the earlier newline-delimited-record
convention is only whether a missing final newline creates an
additional record.

SM advanced from V135 to V141.  V136--V137 keep Dirichlet and Robin
domains separate and retain the coupled gauge--Higgs matrix Hessian
until its boundary symbol is repaired.  V138--V139 separate local
anomaly curvature, global holonomy, and boundary determinant-line
trivialization.  Most relevant here, V140 replaces convergence of
separately chosen phase representatives by direct transition ratios
and exact triangle defects, and warns that a control norm equal to the
desired first variation would be circular.  V141 closes the
finite-cutoff faithful-quotient anomaly card but keeps the cofinal
comparison maps as an independent problem.

\begin{assessment}[V125 adjustment after the audit]
\label{ass:V125-audit-adjustment}
The transferable rules are:
\begin{enumerate}[label=\textup{(\roman*)}]
\item compare the two Wilson multipliers by one direct compressed
      operator difference, rather than prove separate asymptotic
      expansions and subtract them;
\item retain the whole tensor representation through the irreducible
      output compression, just as a coupled boundary matrix is retained
      through its domain calculation;
\item aggregate the resulting physical outer projections before taking
      product-state capacity; and
\item do not postulate a new norm whose value is the missing endpoint
      ratio.
\end{enumerate}
No SM or NS analytic estimate is imported.
\end{assessment}

\section{Direct compression on an arbitrary tensor-product output}
\label{sec:V125-compression}

Let \(\lambda,\mu,\nu\) be irreducible \(SU(3)\) labels with
\(H_\nu\) occurring in \(H_\lambda\otimes H_\mu\).  Resolve a copy by an
isometry
\[
 J_{\nu,\zeta}:H_\nu\longrightarrow H_\lambda\otimes H_\mu.
\]

\begin{theorem}[Universal output-compression difference]
\label{thm:V125-output-compression}
For every output copy \((\nu,\zeta)\),
\begin{equation}
 \boxed{
 |q_{\alpha,\nu}-q_{\alpha,\lambda}|
 \le
 \mathbb E_{\nu_\alpha}
 \|\pi_\mu(U)-I\|_{\rm op}
 \le
 |\mu|_1\,
 \mathbb E_{\nu_\alpha}\|U-I\|_{\rm op}.}
 \label{eq:V125-output-compression}
\end{equation}
The bound is independent of \(\nu\), its lowering depth, and its
multiplicity.
\end{theorem}

\begin{proof}
Because the Wilson law is central, its Fourier transform is scalar in
every irreducible representation.  The intertwining property of
\(J_{\nu,\zeta}\) gives
\begin{align*}
 q_{\alpha,\nu}I_{H_\nu}
 &=
 J_{\nu,\zeta}^*
 \mathbb E_{\nu_\alpha}
 [\pi_\lambda(U)\otimes\pi_\mu(U)]
 J_{\nu,\zeta},\\
 q_{\alpha,\lambda}I_{H_\nu}
 &=
 J_{\nu,\zeta}^*
 \mathbb E_{\nu_\alpha}
 [\pi_\lambda(U)\otimes I]
 J_{\nu,\zeta}.
\end{align*}
Subtract.  Since \(J_{\nu,\zeta}\) is an isometry and
\(\pi_\lambda(U)\) is unitary,
\[
 |q_{\alpha,\nu}-q_{\alpha,\lambda}|
 \le
 \mathbb E_{\nu_\alpha}
 \|\pi_\mu(U)-I\|_{\rm op}.
\]
The second inequality is
\eqref{eq:V124-representation-displacement}.
\end{proof}

\begin{remark}[Why multiplicity is harmless]
\label{rem:V125-multiplicity}
The same scalar \(q_{\alpha,\nu}\) acts on every copy of \(\nu\).
The proof applies to each isometric copy separately with constant one.
It never sums copy dimensions and therefore does not require a
multiplicity estimate.
\end{remark}

\section{A square-root debit for every output}
\label{sec:V125-every-output}

For a resolved output copy \((\nu,\zeta)\subset\lambda\otimes\mu\),
write the literal one-link signed V52 coefficient as
\begin{equation}
 d_{\alpha;\lambda,\mu}^{\nu,\zeta}
 :=
 \bigl(1+s_\alpha(1+c_\nu)\bigr)
 \left(
 q_{\alpha,\nu}
 -
 q_{\alpha,\lambda}q_{\alpha,\mu}
 \right).
 \label{eq:V125-physical-pair-coefficient}
\end{equation}
It is independent of \(\zeta\), but the copy label is retained to match
the physical outer resolution.

\begin{lemma}[Output height bound]
\label{lem:V125-output-height}
If \(\nu\subset\lambda\otimes\mu\), then
\begin{equation}
 |\nu|_1\le|\lambda|_1+|\mu|_1,
 \qquad
 1+c_\nu
 \le
 C\left(1+(|\lambda|_1+|\mu|_1)^2\right).
 \label{eq:V125-output-height}
\end{equation}
\end{lemma}

\begin{proof}
Realize \(H_\lambda\) and \(H_\mu\) as Cartan components of tensor
powers of \(F\) and \(\overline F\), as in
\Cref{lem:V124-representation-displacement}.  Their tensor product is a
subrepresentation of a tensor power with
\(|\lambda|_1+|\mu|_1\) factors.  Every irreducible constituent has
dominant height no larger than the number of factors.  The Casimir
bound follows from \eqref{eq:V124-height-Casimir}.
\end{proof}

\begin{theorem}[Uniform all-output square-root debit]
\label{thm:V125-all-output-debit}
Fix \(K<\infty\).  Suppose
\[
 K^{-1}\sqrt\alpha
 \le|\lambda|_1\le K\sqrt\alpha,
 \qquad
 |\mu|_1\le K\sqrt\alpha,
\]
and let \((\nu,\zeta)\) be any output copy in
\(\lambda\otimes\mu\).  Then
\begin{align}
 |d_{\alpha;\lambda,\mu}^{\nu,\zeta}|
 &\le
 C_K\frac{|\mu|_1}{\sqrt\alpha},
 \label{eq:V125-all-output-height}\\
 \boxed{
 |d_{\alpha;\lambda,\mu}^{\nu,\zeta}|
 \le
 C_K
 \sqrt{\frac{1+c_\mu}{1+c_\lambda}}.}
 \label{eq:V125-all-output-ratio}
\end{align}
The constants are uniform in the output, lowering depth, and
multiplicity.
\end{theorem}

\begin{proof}
Use the exact decomposition
\[
 q_{\alpha,\nu}-q_{\alpha,\lambda}q_{\alpha,\mu}
 =
 (q_{\alpha,\nu}-q_{\alpha,\lambda})
 +q_{\alpha,\lambda}(1-q_{\alpha,\mu}).
\]
The first term is bounded by
\Cref{thm:V125-output-compression,
lem:V123-Wilson-displacement}.  The second is bounded by
\eqref{eq:V30-one-link-upper-gap}.  Exactly as in the proof of
\Cref{thm:V124-general-Cartan-debit}, their sum is at most
\[
 C_K\frac{|\mu|_1}{\sqrt\alpha}.
\]
By \Cref{lem:V125-output-height} and
\eqref{eq:V122-correct-clock}, the weight
\(1+s_\alpha(1+c_\nu)\) is bounded by \(C_K\).
This proves \eqref{eq:V125-all-output-height}.  The Casimir form follows
from \eqref{eq:V124-height-Casimir} and the transition lower bound for
\(\lambda\).
\end{proof}

\section{Support-free Hilbert-space compression}
\label{sec:V125-support-free}

The preceding operator-displacement bound is useful on one coordinate.
For a product bank, a Hilbert-space square avoids the sum of coordinate
heights.

\begin{theorem}[Universal square-root character difference]
\label{thm:V125-square-root-character}
Let \(K\) be a compact group, let \(\nu_\alpha\) be a central
inversion-invariant probability measure on \(K\), and suppose its
Fourier scalars \(q_{\alpha,\rho}\) are real and lie in \([0,1]\).
If an irreducible \(\tau\) occurs in
\(\lambda\otimes\mu\), then
\begin{equation}
 \boxed{
 |q_{\alpha,\tau}-q_{\alpha,\lambda}|
 \le
 \sqrt{2(1-q_{\alpha,\mu})}.}
 \label{eq:V125-square-root-character}
\end{equation}
Consequently,
\begin{equation}
 \boxed{
 |q_{\alpha,\tau}
   -q_{\alpha,\lambda}q_{\alpha,\mu}|
 \le
 \sqrt{2(1-q_{\alpha,\mu})}
 +(1-q_{\alpha,\mu})
 \le
 (1+\sqrt2)\sqrt{1-q_{\alpha,\mu}}.}
 \label{eq:V125-square-root-defect}
\end{equation}
\end{theorem}

\begin{proof}
Let \(J:H_\tau\to H_\lambda\otimes H_\mu\) be an isometric
intertwiner and fix a unit vector \(\xi\in H_\tau\).  Put
\(\eta=J\xi\).  The same compressed difference as in
\Cref{thm:V125-output-compression} gives
\begin{align*}
 q_{\alpha,\tau}-q_{\alpha,\lambda}
 &=
 \mathbb E_{\nu_\alpha}
 \left\langle
 \eta,
 \pi_\lambda(U)\otimes(\pi_\mu(U)-I)\eta
 \right\rangle.
\end{align*}
Scalar Cauchy--Schwarz in the probability variable and unitarity of
\(\pi_\lambda(U)\) yield
\begin{align*}
 |q_{\alpha,\tau}-q_{\alpha,\lambda}|^2
 &\le
 \mathbb E_{\nu_\alpha}
 \left\|
 (I\otimes(\pi_\mu(U)-I))\eta
 \right\|^2\\
 &=
 \left\langle
 \eta,
 I\otimes
 \mathbb E_{\nu_\alpha}
 [(\pi_\mu(U)-I)^*(\pi_\mu(U)-I)]
 \eta
 \right\rangle.
\end{align*}
Centrality and inversion invariance give
\[
 \mathbb E\pi_\mu(U)
 =
 \mathbb E\pi_\mu(U)^*
 =
 q_{\alpha,\mu}I.
\]
Therefore the operator inside the last expectation is exactly
\[
 2(1-q_{\alpha,\mu})I,
\]
proving \eqref{eq:V125-square-root-character}.  Add and subtract
\(q_{\alpha,\lambda}\), use
\(|q_{\alpha,\lambda}|\le1\), and note that
\(0\le1-q_{\alpha,\mu}\le1\) to obtain
\eqref{eq:V125-square-root-defect}.
\end{proof}

\begin{remark}[No coordinate count]
\label{rem:V125-no-coordinate-count}
The proof uses the exact averaged square of the complete representation
\(\pi_\mu\).  If \(K\) is a product group and \(\mu\) is a product
irreducible, its right side is still the single scalar
\(1-q_{\alpha,\mu}\).  It is not the sum of operator displacements of
the coordinates.
\end{remark}

\section{Every output on an arbitrary product bank}
\label{sec:V125-product-bank}

Let \(E\) be a finite coordinate bank.  For product irreducibles write
\[
 \boldsymbol\lambda=(\lambda_e)_{e\in E},
 \qquad
 C_E(\boldsymbol\lambda)
 :=
 \sum_{e\in E}C_2(\lambda_e),
\]
and similarly for \(\boldsymbol\mu,\boldsymbol\nu\).  The Wilson
product scalar is
\[
 q_{\alpha,\boldsymbol\lambda}
 =
 \prod_{e\in E}q_{\alpha,\lambda_e}.
\]

\begin{lemma}[Output Casimir under product-bank fusion]
\label{lem:V125-product-output-Casimir}
There is a numerical \(C<\infty\) such that, whenever
\(\boldsymbol\nu\) is a multiplicity-resolved output of
\(\boldsymbol\lambda\otimes\boldsymbol\mu\),
\begin{equation}
 C_E(\boldsymbol\nu)
 \le
 C\left[
 C_E(\boldsymbol\lambda)+C_E(\boldsymbol\mu)
 \right].
 \label{eq:V125-product-output-Casimir}
\end{equation}
\end{lemma}

\begin{proof}
Coordinatewise, the height argument of
\Cref{lem:V125-output-height} and
\eqref{eq:V124-height-Casimir} give
\[
 C_2(\nu_e)
 \le
 C\bigl(C_2(\lambda_e)+C_2(\mu_e)\bigr).
\]
Sum over \(e\).  Trivial input coordinates contribute zero on both
sides and have only the trivial output.
\end{proof}

For an output copy
\((\boldsymbol\nu,\boldsymbol\zeta)\), define
\begin{align}
 d_{\alpha;\boldsymbol\lambda,\boldsymbol\mu}
 ^{\boldsymbol\nu,\boldsymbol\zeta}
 &:=
 \bigl(1+s_\alpha(1+C_E(\boldsymbol\nu))\bigr)
 \left(
 q_{\alpha,\boldsymbol\nu}
 -
 q_{\alpha,\boldsymbol\lambda}
 q_{\alpha,\boldsymbol\mu}
 \right).
 \label{eq:V125-product-pair-coefficient}
\end{align}

\begin{theorem}[Support-free all-output product-bank debit]
\label{thm:V125-product-bank-debit}
Fix \(K<\infty\).  Suppose
\begin{equation}
 K^{-1}
 \le
 s_\alpha C_E(\boldsymbol\lambda)
 \le K,
 \qquad
 s_\alpha C_E(\boldsymbol\mu)\le K.
 \label{eq:V125-product-transition}
\end{equation}
Then every multiplicity-resolved output obeys
\begin{align}
 \left|
 d_{\alpha;\boldsymbol\lambda,\boldsymbol\mu}
 ^{\boldsymbol\nu,\boldsymbol\zeta}
 \right|
 &\le
 C_K
 \min\left\{
 \sqrt{s_\alpha C_E(\boldsymbol\mu)},1
 \right\},
 \label{eq:V125-product-s-debit}\\
 \boxed{
 \left|
 d_{\alpha;\boldsymbol\lambda,\boldsymbol\mu}
 ^{\boldsymbol\nu,\boldsymbol\zeta}
 \right|
 \le
 C_K
 \sqrt{
 \frac{1+C_E(\boldsymbol\mu)}
      {1+C_E(\boldsymbol\lambda)}}.}
 \label{eq:V125-product-ratio-debit}
\end{align}
The constants are independent of the number of coordinates, output
labels, and multiplicities.
\end{theorem}

\begin{proof}
Apply \eqref{eq:V125-square-root-defect} on the product group
\(SU(3)^E\).  The upper half of the V30 product character-gap theorem
\eqref{eq:V30-two-sided-product-gap} gives
\[
 1-q_{\alpha,\boldsymbol\mu}
 \le
 C\min\{s_\alpha C_E(\boldsymbol\mu),1\}.
\]
The output weight is bounded by \(C_K\) under
\eqref{eq:V125-product-transition} and
\eqref{eq:V125-product-output-Casimir}.  This proves
\eqref{eq:V125-product-s-debit}.  The lower transition bound on
\(\boldsymbol\lambda\), together with
\(s_\alpha\asymp\alpha^{-1}\), converts it to
\eqref{eq:V125-product-ratio-debit}; adding one to each total Casimir
handles the trivial small input.
\end{proof}

\section{Two arbitrary-output banks and the full endpoint ratio}
\label{sec:V125-two-banks}

Let the old \(ab\) bank have product inputs
\((\boldsymbol\lambda_a,\boldsymbol\mu_1)\) and arbitrary resolved output
\(\boldsymbol\nu_1\).  Let the disjoint off-cut \(bc\) bank have product
inputs
\((\boldsymbol\mu_2,\boldsymbol\lambda_c)\) and arbitrary resolved output
\(\boldsymbol\nu_2\).  Put
\begin{align}
 A&:=1+C_E(\boldsymbol\lambda_a),
 &
 C&:=1+C_E(\boldsymbol\lambda_c),
 \label{eq:V125-large-bank-masses}\\
 B_1&:=1+C_E(\boldsymbol\mu_1),
 &
 B_2&:=1+C_E(\boldsymbol\mu_2),
 &
 B&:=B_1+B_2.
 \label{eq:V125-middle-bank-masses}
\end{align}

\begin{theorem}[All-output two-bank endpoint ratio below the hot scale]
\label{thm:V125-two-bank-ratio}
Fix \(K<\infty\).  Suppose
\begin{equation}
 s_\alpha A\le K,
 \qquad
 s_\alpha C\le K,
 \qquad
 s_\alpha B\le K.
 \label{eq:V125-bounded-bank-scales}
\end{equation}
Then every pair of multiplicity-resolved outputs satisfies
\begin{equation}
 \boxed{
 \left|
 d_{\alpha;\boldsymbol\lambda_a,\boldsymbol\mu_1}
 ^{\boldsymbol\nu_1,\boldsymbol\zeta_1}
 d_{\alpha;\boldsymbol\lambda_c,\boldsymbol\mu_2}
 ^{\boldsymbol\nu_2,\boldsymbol\zeta_2}
 \right|
 \le
 C_K\frac{B}{C}.}
 \label{eq:V125-two-bank-ratio}
\end{equation}
No Cartan, concentration, support-size, output-depth, or multiplicity
hypothesis is imposed.
\end{theorem}

\begin{proof}
The proof of \Cref{thm:V125-product-bank-debit}, before the transition
lower bound is used, gives
\[
 \left|
 d_{\alpha;\boldsymbol\lambda,\boldsymbol\mu}
 ^{\boldsymbol\nu,\boldsymbol\zeta}
 \right|
 \le
 C_K
 \min\left\{
 \sqrt{s_\alpha C_E(\boldsymbol\mu)},1
 \right\}
\]
whenever the two input heat scales are bounded by \(K\).  Apply this to
both banks and discard the minima:
\begin{align*}
 |d_1d_2|
 &\le
 C_Ks_\alpha
 \sqrt{
 C_E(\boldsymbol\mu_1)
 C_E(\boldsymbol\mu_2)}\\
 &\le
 \frac{C_Ks_\alpha}{2}
 \left[
 C_E(\boldsymbol\mu_1)+
 C_E(\boldsymbol\mu_2)
 \right]
 \le C_Ks_\alpha B.
\end{align*}
Finally \(s_\alpha C\le K\) gives
\[
 s_\alpha B
 =
 (s_\alpha C)\frac{B}{C}
 \le K\frac{B}{C}.
\]
This proves \eqref{eq:V125-two-bank-ratio}.
\end{proof}

\begin{remark}[Two square roots, now without Cartan labels]
\label{rem:V125-two-general-square-roots}
Each factor in the first line of the proof is one square-root character
gap of the complete middle-row product representation on its bank.
Their product is converted by arithmetic--geometric mean into one
linear middle-row mass.  This is exactly the V123 mechanism, but the
universal compression theorem shows it holds simultaneously on every
irreducible output.
\end{remark}

\section{No-count physical aggregation of all pair outputs}
\label{sec:V125-physical-aggregation}

Fix the pair inputs above.  Let \(\mathfrak J_{\le K}\) be the complete
V98 outer subfamily satisfying
\eqref{eq:V125-bounded-bank-scales} after the already paid \(ab\)
hybrid row is removed.  Allow both pair outputs, all their
multiplicities, and every remaining final/copy label to vary.

\begin{theorem}[Complete bounded-scale covariance-family payment]
\label{thm:V125-complete-family}
Suppose the exact reduced \(ac\) row on every block obeys the V118
bound
\[
 \|\mathbf R_{ac,j}^{1,\lhd}\|_{\rm op}
 \le Cs_\alpha H_{ac}^{1}.
\]
Then
\begin{align}
 \mathbf Q_{\le K}
 &:=
 \sum_{j\in\mathfrak J_{\le K}}
 |V_jL_{ac,j}V_j^*|
 \notag\\
 &\preccurlyeq
 C_Ks_\alpha H_{ac}^{1}\frac{B}{C}
 \sum_{j\in\mathfrak J_{\le K}}P_j
 \preccurlyeq
 C_Ks_\alpha H_{ac}^{1}\frac{B}{C}I,
 \label{eq:V125-complete-family-order}\\
 \boxed{
 \|\mathbf Q_{\le K}\|_{\rm rows,prod}
 \le
 C_Ks_\alpha H_{ac}^{1}\frac{B}{C}.}
 \label{eq:V125-complete-family-capacity}
\end{align}
There is no sum over output depths or Littlewood--Richardson copies.
\end{theorem}

\begin{proof}
On every physical block the V118 factorization is
\[
 L_{ac,j}
 =
 d_{ab,j}d_{bc,j}\mathbf R_{ac,j}^{1,\lhd}.
\]
Apply \eqref{eq:V125-two-bank-ratio} and the reduced-row bound.
The \(P_j\) are orthogonal complete outer projections, so absolute
values add and \(\sum_jP_j\preccurlyeq I\) by
\Cref{thm:V98-exact-absolute-carrier}.  This proves the operator order.
Take original-row product-state expectations for the last statement.
\end{proof}

\begin{corollary}[Bounded-scale off-cut covariance closure]
\label{cor:V125-bounded-offcut-closed}
Assume the old \(ab\) and off-cut \(bc\) banks carry fixed positive
fractions of the relevant row masses, the old pair is bi-heavy with
\(h_{ab}\ge\varepsilon\Xi_a\), and
\[
 \Xi_a>\Xi_c>C_{\rm row}\Xi_b,
 \qquad
 s_\alpha\Xi_a\le K.
\]
Then the complete off-cut covariance subfamily with those selected
banks satisfies
\begin{equation}
\|\mathbf Q_{\le K}\|_{\rm rows,prod}
\le
C_Ks_\alpha H_{ac}^{1}\frac{\Xi_b}{\Xi_c}
\le
C_{K,\varepsilon}
\frac{h_{ab}H_{ac}^{1}}{\Xi_a\Xi_c}.
\label{eq:V125-global-bounded-closure}
\end{equation}
Together with the V116 mean-branch payment and the V119 \(ab\)-row
payment, this closes the full off-cut branch on the bounded-\(a\)-heat
sector.
\end{corollary}

\begin{proof}
Selected-bank heaviness gives
\[
 B\le C\Xi_b,
 \qquad
 \Xi_c\le C C,
 \qquad
 A\le C\Xi_a.
\]
The residual ordering controls the two middle-row bank scales by the
\(c\)-row scale, and \(s_\alpha\Xi_a\le K\) supplies
\eqref{eq:V125-bounded-bank-scales} after constants are enlarged.
Insert the comparisons in
\eqref{eq:V125-complete-family-capacity}.  Moreover
\[
 s_\alpha\Xi_b
 \le s_\alpha\Xi_a
 \le K
 \le\frac K\varepsilon\frac{h_{ab}}{\Xi_a},
\]
which proves the second inequality in
\eqref{eq:V125-global-bounded-closure} and places the result at the
V119 marked-tree scale.  The dispositions of the other two branches are
\Cref{thm:V116-mean-branch-paid,
thm:V119-two-row-disposition}.
\end{proof}

\begin{firewall}[What remains after the universal output theorem]
\label{fire:V125-remaining-sector}
V125 removes arbitrary outputs, multiplicities, coordinate support, and
within-bank mass dispersion from the bounded-\(a\)-heat sector.  The
remaining off-cut covariance regime is the scale mismatch
\[
 s_\alpha\Xi_a\longrightarrow\infty,
 \qquad
 \Xi_a>\Xi_c>C_{\rm row}\Xi_b,
\]
where the old \(ab\) output weight can be hot even when the \(c\)-bank
is not.  The V125 proof must not be extended there by declaring its
bounded output weight uniform.
\end{firewall}

\section{V125 ledger and V126 mandate}
\label{sec:V125-ledger}

\begin{theorem}[Integrated V125 statement]
\label{thm:V125-integrated}
Preserving V1--V124, V125 proves:
\begin{enumerate}[label=\textup{(V125.\arabic*)},leftmargin=7.2em]
\item the mandatory SM V141 and NS V31 companion audit;
\item the direct arbitrary-output compression
      \eqref{eq:V125-output-compression};
\item the support-free square-root difference
      \eqref{eq:V125-square-root-character};
\item the all-output one-coordinate debit
      \eqref{eq:V125-all-output-ratio};
\item its support-free product-bank extension
      \eqref{eq:V125-product-ratio-debit};
\item multiplication of two arbitrary-output debits into
      \eqref{eq:V125-two-bank-ratio};
\item no-count aggregation over all pair, final, and copy outputs; and
\item closure of the complete bounded-\(a\)-heat off-cut branch.
\end{enumerate}
\end{theorem}

\begin{proof}
The items are
\Cref{sec:V125-audit,
thm:V125-output-compression,
thm:V125-square-root-character,
thm:V125-all-output-debit,
thm:V125-product-bank-debit,
thm:V125-two-bank-ratio,
thm:V125-complete-family,
cor:V125-bounded-offcut-closed}.
\end{proof}

\begin{openproblem}[V126 mandate: the hot-\(a\) scale mismatch]
\label{op:V126-mandate}
\begin{enumerate}
\item On \(s_\alpha\Xi_a\gg1\), retain the exact old \(ab\) output
      weight together with its Wilson multiplier; do not bound the
      weight and multiplier separately.
\item Apply the V30/V31 hot-output decay to the complete compressed
      difference in
      \Cref{thm:V125-output-compression}, or derive a weighted
      Cauchy--Schwarz version of
      \eqref{eq:V125-square-root-character}.
\item Split according to whether the old \(ab\) output remains hot or
      falls to a cancellation channel under reverse fusion.
\item In the hot-output case, seek a factor
      \((s_\alpha\Xi_a)^{-1/2}\) before multiplying the off-cut
      square-root debit.
\item In the cancellation case, aggregate the low-output invariant
      projector and use its product capacity; do not count outputs.
\item Combine the result with
      \eqref{eq:V125-global-bounded-closure}.
\end{enumerate}
\end{openproblem}

\section{Firewall after V125}
\label{sec:V125-firewall}

\begin{firewall}[Current claim boundary]
\label{fire:V125-current-boundary}
V125 closes the complete off-cut carrier only when the selected
\(a\)-bank has bounded Wilson heat scale.  The hot-\(a\) scale mismatch
remains.  Therefore V125 does not close the entire residual branch,
prove JREC, prove quartic or all-order MTA, construct continuum
Yang--Mills theory, or prove a positive spectral gap.  The full
Yang--Mills existence and mass gap remain open.
\end{firewall}

\section{Append-only record for V125}
\label{sec:V125-record}

V125 was appended to the complete V124 manuscript.  The exact V124
parent had \(83074\) newline-delimited source records, \(2882996\)
bytes, and SHA--256 digest
\[
 \texttt{7ee6eb40aface23c079f57a0c2b26d9b36f4ada2e90f6ef18ba7eece804f0239}.
\]
No V124 mathematical content was changed.  The companion files were
read only.  The sole final \verb|\end{document}| is restored below.

% =============================================================================
% END APPENDED VERSION V125 MATERIAL
% =============================================================================

% =============================================================================
% BEGIN APPENDED VERSION V126 MATERIAL
% =============================================================================

\chapter{V126: Upstream Maximum-Square Repair and Complete
Off-Cut Closure}
\label{ch:V126}

\section{Purpose and correction of course}
\label{sec:V126-purpose}

V125 left the hot-\(a\) scale mismatch
\[
 s_\alpha\Xi_a\longrightarrow\infty,
 \qquad
 \Xi_a>\Xi_c>C_{\rm row}\Xi_b
\]
as the last part of the off-cut covariance branch.  Its proposed V126
route was to prove a new weighted Wilson estimate after splitting the old
\(ab\) output into hot and cancellation channels.

That route is unnecessary.  The upstream V109 theorem proved two bounds
for the same bi-heavy pair functional.  V119 retained only the weaker
geometric consequence
\[
 \frac{h_{ab}}{s_\alpha\Xi_a\Xi_b},
\]
although the preceding line of V109 gives the stronger maximum-square
bound
\[
 \frac{h_{ab}}
 {s_\alpha\max\{\Xi_a^2,\Xi_b^2\}}.
\]
In the residual ordering, \(\Xi_a\) is larger than both \(\Xi_b\) and
\(\Xi_c\).  Thus the already-proved old-pair capacity has denominator
\(\Xi_a^2\), which is stronger than either marked-path denominator needed
for the two hybrid rows.  Keeping that first V109 inequality through the
V112 positive majorant closes the mismatch before any heat-scale split.

This is an upstream ledger repair, not a new asymptotic assertion about
Wilson coefficients.  The V120--V125 computations remain valid on their
stated sectors, but they are no longer needed to pay the off-cut branch.

\section{The maximum-square tax survives positive completion}
\label{sec:V126-max-square-capacity}

Retain the notation of V109, V112, and V119.  Thus \(L\) is one physical
pair bank joining rows \(i,k\),
\[
 x_i\ge\varepsilon\Xi_i,
 \qquad
 x_k\ge\varepsilon\Xi_k,
 \qquad
 h=\sum_{e\in L}r_eq_e,
\]
and
\[
 \mathbf A_L=\sum_\lambda a_\lambda P_\lambda,
 \qquad
 \kappa_L:=\|\mathbf A_L\|_{\rm prod}
\]
is the V112 positive absolute multiplier and its product-state capacity.

\begin{lemma}[Positive maximum-square pair capacity]
\label{lem:V126-positive-max-square}
Under the V109 bi-heavy hypotheses,
\begin{equation}
 \boxed{
 \kappa_L
 \le
 C_\varepsilon
 \frac{h}
 {s_\alpha\max\{\Xi_i^2,\Xi_k^2\}}.}
 \label{eq:V126-positive-max-square}
\end{equation}
The estimate concerns one occurrence of the physical pair multiplier.
\end{lemma}

\begin{proof}
For positive trace-class inputs \(A_i,A_k\), the V112 duality identity
identifies
\[
 \operatorname{Tr}\!\left[
   \mathbf A_L(A_i\otimes A_k)
 \right]
\]
with the nonnegative V52/V53 formation--defect functional
\(\mathfrak F_L(A_i,A_k)\).  The first conclusion of
\Cref{thm:V109-bi-heavy-tax}, namely
\eqref{eq:V109-max-square-tax}, therefore gives
\[
 \operatorname{Tr}\!\left[
   \mathbf A_L(A_i\otimes A_k)
 \right]
 \le
 C_\varepsilon
 \frac{h}
 {s_\alpha\max\{\Xi_i^2,\Xi_k^2\}}
 \|A_i\|_1\|A_k\|_1.
\]
Take the supremum over positive product inputs with unit trace norm.
This is exactly the definition of \(\kappa_L\).  No square, product, or
second copy of the pair tax has been introduced.
\end{proof}

\begin{remark}[The line dropped in V119]
\label{rem:V126-dropped-line}
Equation \eqref{eq:V119-old-pair-capacity} cited V109 but recorded only
\eqref{eq:V109-geometric-tax}.  It is a correct inequality, yet it
discarded \eqref{eq:V109-max-square-tax}.  The loss is harmless in the
matched \(ab\) row but artificial in the noncomparable \(ac\) row.
\Cref{lem:V126-positive-max-square} merely retains the stronger line.
\end{remark}

\section{Shared-row completion with the stronger old-pair tax}
\label{sec:V126-shared-row}

Let \(\mathbf A_{ab}\) be the old bi-heavy pair multiplier and
\(\mathbf A_{bc}^{0}\) the peeled off-cut multiplier of V118.  Recall
\[
 \mathcal K_\star^{ab,bc}
 =
 \|\mathbf A_{ab}\otimes\mathbf A_{bc}^{0}\|_{\rm 3prod}.
\]
The banks share original row \(b\), so their marginal capacities must not
be multiplied.  The following argument uses only the valid one-sided
positive-order estimate.

\begin{lemma}[Maximum-square shared-row collapse]
\label{lem:V126-star-max-square}
Assume the old \(ab\) bank is bi-heavy at \(a,b\), with pair mass
\(h_{ab}\).  Then
\begin{equation}
 \boxed{
 \mathcal K_\star^{ab,bc}
 \le
 C_\varepsilon A_{bc,*}
 \frac{h_{ab}}
 {s_\alpha\max\{\Xi_a^2,\Xi_b^2\}},}
 \qquad
 A_{bc,*}:=\|\mathbf A_{bc}^{0}\|_{\rm op}.
 \label{eq:V126-star-max-square}
\end{equation}
In particular, using the universal V52 weighted-defect cap,
\begin{equation}
 \mathcal K_\star^{ab,bc}
 \le
 C_\varepsilon
 \frac{h_{ab}}
 {s_\alpha\max\{\Xi_a^2,\Xi_b^2\}}.
 \label{eq:V126-star-uniform-max-square}
\end{equation}
\end{lemma}

\begin{proof}
The second one-sided inequality in
\eqref{eq:V118-one-sided-star} is the positive-order statement
\[
 \mathcal K_\star^{ab,bc}
 \le A_{bc,*}\kappa_{ab}.
\]
Apply \Cref{lem:V126-positive-max-square} to \(\kappa_{ab}\).  This
proves \eqref{eq:V126-star-max-square}.  The V52 uniform coefficient cap
\eqref{eq:V52-defect-weight-cap} gives \(A_{bc,*}\le C\), independently
of \(s_\alpha\), the effective number, and every output label.  Substitution
proves \eqref{eq:V126-star-uniform-max-square}.
\end{proof}

\begin{firewall}[No illicit product of capacities]
\label{fire:V126-no-product-capacity}
The proof does not claim
\(\mathcal K_\star^{ab,bc}\le\kappa_{ab}\kappa_{bc}^{0}\).
It bounds the off-cut factor in operator norm and charges product-state
capacity only to the old \(ab\) factor, exactly as permitted by
\Cref{lem:V118-one-sided-star}.  The two factors may remain arbitrarily
correlated on their shared original row.
\end{firewall}

\section{Closure of the residual noncomparable ordering}
\label{sec:V126-residual-closure}

\begin{theorem}[Both residual hybrid rows have their marked denominators]
\label{thm:V126-residual-hybrid-closure}
Assume the exact V118 off-cut covariance setting, the old \(ab\) bank is
bi-heavy at \(a,b\), and
\begin{equation}
 \Xi_a>\Xi_c>C_{\rm row}\Xi_b.
 \label{eq:V126-residual-order}
\end{equation}
Then, for arbitrary row-factorized trace-class inputs, the two physical
hybrid responses obey
\begin{align}
 \mathcal R_{ab}^{\rm off,cov}
 &\le
 C_\varepsilon
 \frac{h_{ab}H_{ab}^{1}}
      {\Xi_a\Xi_b}
 \prod_v\|B_v\|_1,
 \label{eq:V126-ab-row}\\
 \mathcal R_{ac}^{\rm off,cov}
 &\le
 C_\varepsilon
 \frac{h_{ab}H_{ac}^{1}}
      {\Xi_a\Xi_c}
 \prod_v\|B_v\|_1.
 \label{eq:V126-ac-row}
\end{align}
The constants are uniform for all \(s_\alpha>0\).  In particular, no
assumption on \(s_\alpha\Xi_a\) occurs.
\end{theorem}

\begin{proof}
For \(r\in\{b,c\}\), the exact V118 physical row estimate is
\eqref{eq:V118-physical-row-capacity}:
\[
 \mathcal R_{ar}^{\rm off,cov}
 \le
 C s_\alpha H_{ar}^{1}
 \mathcal K_\star^{ab,bc}
 \prod_v\|B_v\|_1.
\]
Insert \eqref{eq:V126-star-uniform-max-square}.  Because
\eqref{eq:V126-residual-order} implies \(\Xi_a>\Xi_b\), one obtains
\begin{equation}
 \mathcal R_{ar}^{\rm off,cov}
 \le
 C_\varepsilon
 \frac{h_{ab}H_{ar}^{1}}{\Xi_a^2}
 \prod_v\|B_v\|_1.
 \label{eq:V126-square-row-bound}
\end{equation}
For \(r=b\), the inequality \(\Xi_b<\Xi_a\) gives
\[
 \frac1{\Xi_a^2}
 \le
 \frac1{\Xi_a\Xi_b}.
\]
For \(r=c\), the inequality \(\Xi_c<\Xi_a\) gives
\[
 \frac1{\Xi_a^2}
 \le
 \frac1{\Xi_a\Xi_c}.
\]
Applying these two comparisons to
\eqref{eq:V126-square-row-bound} proves
\eqref{eq:V126-ab-row}--\eqref{eq:V126-ac-row}.
\end{proof}

\begin{corollary}[The V119 scalar gate is unnecessary]
\label{cor:V126-scalar-gate-bypassed}
On the residual ordering \eqref{eq:V126-residual-order}, neither
\eqref{eq:V119-scalar-gate} nor its effective-scale sufficient condition
\eqref{eq:V119-effective-scale-gate} is needed.  The complete physical
\(ac\) row is paid even when
\[
 \vartheta_0=1
 \quad\hbox{or}\quad
 s_\alpha\Xi_a\longrightarrow\infty.
\]
\end{corollary}

\begin{proof}
The V119 route first weakened the old-pair capacity to the denominator
\(\Xi_a\Xi_b\), then asked the off-cut operator norm to convert
\(\Xi_b\) into \(\Xi_c\).  The proof of
\Cref{thm:V126-residual-hybrid-closure} retains the denominator
\(\Xi_a^2\), which already dominates \(\Xi_a\Xi_c\).  Consequently only
the uniform estimate \(A_{bc,*}\le C\) is used; the value of
\(\vartheta_0\) is irrelevant.
\end{proof}

\begin{proposition}[The concentrated V119 regression is consistent]
\label{prop:V126-regression-consistent}
The mass data in
\Cref{prop:V119-concentrated-regression} do not obstruct
\Cref{thm:V126-residual-hybrid-closure}.  They disprove only a derivation
of \(\vartheta_0\lesssim\Xi_b/\Xi_c\) from mass geometry.
\end{proposition}

\begin{proof}
The regression has \(\vartheta_0=1\), so the V119 scalar gate fails.
Nevertheless the old pair in the residual branch is bi-heavy and
\(\Xi_a>\Xi_c\).  Its maximum-square capacity is bounded by
\eqref{eq:V126-positive-max-square}; multiplying by the V118 row response
gives \eqref{eq:V126-square-row-bound}.  No assertion about
\(\vartheta_0\) follows or is required.
\end{proof}

\section{Global off-cut disposition}
\label{sec:V126-global-disposition}

\begin{theorem}[Complete off-cut covariance branch]
\label{thm:V126-complete-offcut}
Within the exact V94--V119 priority partition and under its standing
programme hypotheses, every physical off-cut covariance outer block has
the required marked-tree payment.  More precisely:
\begin{enumerate}[label=\textup{(\roman*)},leftmargin=4.2em]
\item the row-comparable same-pair region is paid by
      \Cref{thm:V113-physical-row-comparable};
\item the retained-mean region is paid by
      \Cref{thm:V116-mean-branch-paid};
\item the covariance region with maximal row \(c\) is paid by
      \Cref{thm:V116-covariance-cmax-paid}; and
\item in the remaining noncomparable covariance region, the matched \(ab\) and
      mismatched \(ac\) hybrid rows are both paid by
      \Cref{thm:V126-residual-hybrid-closure}.
\end{enumerate}
The physical absolute carriers are aggregated before product-state
capacity is taken, so no pair-output, final-output, or multiplicity count
is introduced.
\end{theorem}

\begin{proof}
The V113/V116 priority split is disjoint at the level of the orthogonal
V98 outer labels.  After the row-comparable, mean, and maximal-\(c\)
families are removed, V119 leaves
exactly the residual ordering
\(\Xi_a>\Xi_c>C_{\rm row}\Xi_b\), with the old \(ab\) bank bi-heavy.
Apply \Cref{thm:V126-residual-hybrid-closure} there.  Its input is the
already-completed V118 physical row carrier, not a blockwise scalar sum;
hence the V98 no-count aggregation survives.  The three alternatives are
exhaustive by the cited priority partition.
\end{proof}

\begin{corollary}[Hot-symbol investigations V120--V125 are auxiliary]
\label{cor:V126-hot-symbol-auxiliary}
The heat and Wilson Cartan asymptotics of V120--V123 and the all-output
compression theorem of V124--V125 remain rigorous statements in their
recorded scope.  They are not premises of
\Cref{thm:V126-complete-offcut}.  In particular, the square-root
transition obstruction found there does not obstruct the physical
off-cut payment, because the latter uses the independent old-pair
maximum-square capacity before taking the shared-row norm.
\end{corollary}

\section{Next upstream acquisition}
\label{sec:V126-next}

Closing one branch does not by itself prove the complete JREC estimate.
The next step must return to the exact V94 JREC layer cake and its V98
physical outer partition, remove the now-paid off-cut family, and identify
which simultaneous-sideways family, if any, remains outside the existing
V107--V116 payments.  This bookkeeping must precede any claim that JREC
or quartic MTA is closed.

\begin{openproblem}[V127 mandate: recompute the JREC residual after
off-cut deletion]
\label{op:V127-mandate}
\begin{enumerate}
\item Reconstruct the exact list of V94 JREC level-set alternatives and
      the V98 orthogonal physical outer families.
\item Mark every side, same-pair, retained-mean, and off-cut subfamily
      already paid by V99--V116 and
      \Cref{thm:V126-complete-offcut}.
\item Prove that the marked families are disjoint or impose an explicit
      priority order; do not add their product capacities as though the
      underlying Racah projections commuted.
\item State the exact surviving JREC carrier, including its orientation,
      outer projections, row masses, and requested marked-tree weight.
\item If no carrier survives, assemble the proved estimates into JREC.
      Otherwise attack only the surviving carrier and go upstream if its
      proposed norm estimate repeats a previously failed gate.
\end{enumerate}
\end{openproblem}

\section{V126 ledger}
\label{sec:V126-ledger}

\begin{theorem}[Integrated V126 statement]
\label{thm:V126-integrated}
Preserving V1--V125, V126 proves:
\begin{enumerate}[label=\textup{(V126.\arabic*)},leftmargin=7.5em]
\item the positive maximum-square old-pair capacity
      \eqref{eq:V126-positive-max-square};
\item its valid shared-row propagation
      \eqref{eq:V126-star-max-square} without multiplying marginal
      capacities;
\item the stronger two-row estimate
      \eqref{eq:V126-square-row-bound};
\item payment of both residual hybrid rows uniformly in heat scale;
\item removal of the V119 scalar gate as an artefact of weakening the
      V109 denominator; and
\item closure of the complete off-cut covariance branch within the
      established physical priority partition.
\end{enumerate}
\end{theorem}

\begin{proof}
The items are
\Cref{lem:V126-positive-max-square,
lem:V126-star-max-square,
thm:V126-residual-hybrid-closure,
cor:V126-scalar-gate-bypassed,
thm:V126-complete-offcut}.
\end{proof}

\begin{firewall}[Claim boundary after V126]
\label{fire:V126-current-boundary}
V126 closes the off-cut covariance branch under the standing V94--V119
programme hypotheses.  It does not yet prove that every JREC level-set
alternative has been paid, prove JREC, prove quartic or all-order MTA,
construct four-dimensional continuum Yang--Mills theory satisfying the
required axioms, or prove a positive mass gap.  The full Yang--Mills
existence and mass-gap problem remains open.
\end{firewall}

\section{Append-only record for V126}
\label{sec:V126-record}

V126 was appended to the complete V125 manuscript.  The exact V125
parent had \(83848\) newline-delimited source records, \(2906323\)
bytes, and SHA--256 digest
\[
 \texttt{b6e14d0813553508e323603d05be1f44199d5e55557aba9708231f607e256c28}.
\]
No V125 mathematical content was changed.  The sole final
\verb|\end{document}| is restored below.

% =============================================================================
% END APPENDED VERSION V126 MATERIAL
% =============================================================================

% =============================================================================
% BEGIN APPENDED VERSION V127 MATERIAL
% =============================================================================

\chapter{V127: Exact JREC Residual Ledger after Off-Cut Deletion}
\label{ch:V127}

\section{Purpose}
\label{sec:V127-purpose}

V126 closes the complete physical off-cut covariance branch by retaining
the maximum-square half of the V109 bi-heavy tax.  The next task is not to
declare JREC closed, but to recompute its residual from the exact V94
level partition and the V98 physical outer partition.

The outcome is a substantial reduction.  After the paid and disconnected
families are removed, V98 originally left
\(\mathfrak J_{\rm side}\dot\cup\mathfrak J_{\rm off}\).
The second family is now paid.  V102 and V105 split the first into
disjoint-pair, same-pair, and complementary-pair alternatives.  The
row-comparable and off-cut descendants of the same-pair alternative are
now paid, while V114 proves that its only other descendants return to the
disjoint or complementary two-pair geometries.  Hence there is no
standalone one-pair or off-cut carrier left in the JREC residual.

The surviving analytic issue is the finite-scale physical estimate for
the two-pair side carriers.  V104 and V108 give exact positive surfaces
and conditional decay, but pair heaviness alone does not force their
effective-scale hypotheses.  The noncommuting orientation issue does not
create an additional obstruction once each fixed orientation has been
paid: ordinary positivity and subadditivity combine their already-bounded
effects without a joint Racah measurement.

\section{Outer families and permissible refinements}
\label{sec:V127-outer-families}

For a fixed completely regrouped V60 collision/puncture monomial, recall
the V98 physical partition
\begin{equation}
 \mathfrak J
 =
 \mathfrak J_{\rm paid}
 \mathbin{\dot\cup}
 \mathfrak J_{\rm side}
 \mathbin{\dot\cup}
 \mathfrak J_{\rm off}
 \mathbin{\dot\cup}
 \mathfrak J_{\rm disc}.
 \label{eq:V127-V98-partition}
\end{equation}
Every \(j\in\mathfrak J\) is a complete physical outer label: it includes
both external pair outputs, all same-label puncture outputs, the complete
final irreducible, and every genuine multiplicity label.  Thus a
partition of its index set is an orthogonal physical partition, whereas a
simultaneous sharp partition by incompatible Racah intermediates is not.

\begin{lemma}[Safe deterministic refinement]
\label{lem:V127-safe-refinement}
Let
\(\mathfrak J_X=\dot\bigcup_{\sigma\in\Sigma}\mathfrak J_{X,\sigma}\)
be any deterministic refinement of one V98 outer family by incidence,
mass-comparison, and coordinate-support predicates already determined by
the complete label.  Put
\[
 \mathbf Q_{X,\sigma}
 :=
 \sum_{j\in\mathfrak J_{X,\sigma}}
 V_j|L_j|V_j^*.
\]
Then
\begin{equation}
 \mathbf Q_X
 =
 \sum_{\sigma\in\Sigma}\mathbf Q_{X,\sigma},
 \qquad
 \mathbf Q_{X,\sigma}\mathbf Q_{X,\sigma'}=0
 \quad(\sigma\ne\sigma').
 \label{eq:V127-safe-refinement}
\end{equation}
Consequently every block is charged once, and no cardinality factor
appears when positive operator bounds are summed.
\end{lemma}

\begin{proof}
The isometries \(V_j\) have mutually orthogonal ranges by
\Cref{thm:V98-exact-absolute-carrier}.  A partition of their index set
therefore partitions an orthogonal direct sum.  Functional calculus has
already been taken on each complete block, so the displayed carriers are
positive and have pairwise orthogonal support.  Summing their definitions
gives \eqref{eq:V127-safe-refinement}.
\end{proof}

\begin{firewall}[The refinement is not a joint Racah pinching]
\label{fire:V127-no-joint-pinching}
\Cref{lem:V127-safe-refinement} applies only to predicates attached to
the complete V98 outer label.  It does not refine a block simultaneously
by noncommuting \(AB\)-first and \(AC\)-first intermediate projections.
That forbidden operation remains excluded by
\Cref{thm:V98-Racah-instrument-firewall}.
\end{firewall}

\section{Exact deletion of the off-cut family}
\label{sec:V127-offcut-deletion}

\begin{theorem}[Physical residual after V126]
\label{thm:V127-after-offcut}
Under the standing hypotheses of the V94--V119 programme, the unresolved
physical JREC carrier may be reduced from
\[
 \mathbf Q_{\rm res}
 =
 \mathbf Q_{\rm side}\oplus\mathbf Q_{\rm off}
\]
to
\begin{equation}
 \boxed{\mathbf Q_{\rm res}^{\,127}=\mathbf Q_{\rm side}.}
 \label{eq:V127-side-only}
\end{equation}
The removal is performed on complete orthogonal outer blocks and does not
alter the V94 private-mean, cut-heavy, or selected-pair accounts.
\end{theorem}

\begin{proof}
The direct-sum identity is
\Cref{cor:V98-paid-removal}.  The disconnected family is already zero by
\eqref{eq:V98-disconnected-outer-zero}.  The complete off-cut family has
its marked-tree payment by
\Cref{thm:V126-complete-offcut}.  Move that entire orthogonal family to
the paid class using
\Cref{lem:V127-safe-refinement}.  No side block is moved, and no paid
debit is reused.  What remains is exactly \(\mathbf Q_{\rm side}\).
\end{proof}

\section{Exhaustive refinement of the side family}
\label{sec:V127-side-refinement}

For each fixed retained-row orientation \(o\), apply the V102 endpoint
localization and the V105 quarter-mass selection after the V94 priority
selector.  Denote by
\[
 \mathfrak J_{{\rm D},o},\qquad
 \mathfrak J_{{\rm S},o},\qquad
 \mathfrak J_{{\rm C},o}
\]
the resulting disjoint-bank, same-pair, and complementary-pair
subfamilies.  A fixed priority order is used when more than one
certificate is available.

\begin{lemma}[Three-family side partition]
\label{lem:V127-three-side-families}
For each fixed orientation \(o\),
\begin{equation}
 \mathfrak J_{{\rm side},o}
 =
 \mathfrak J_{{\rm D},o}
 \mathbin{\dot\cup}
 \mathfrak J_{{\rm S},o}
 \mathbin{\dot\cup}
 \mathfrak J_{{\rm C},o}.
 \label{eq:V127-three-side-partition}
\end{equation}
On \(\mathfrak J_{{\rm D},o}\) the complete carrier has the V105
disjoint-coordinate two-pair factor.  On
\(\mathfrak J_{{\rm C},o}\) it has the V108 pair-first complementary
two-pair subeffect.  On \(\mathfrak J_{{\rm S},o}\) one physical pair
bank is bi-heavy at both certified endpoints.
\end{lemma}

\begin{proof}
After the V94 paid branches and the now-paid off-cut branch have been
removed, \Cref{thm:V102-forced-pair-or-offcut} supplies a genuine pair
certificate at every sideways endpoint.  Apply
\Cref{thm:V105-audited-extraction}.  Its three alternatives are exhaustive,
and a fixed priority makes them disjoint.  Their carrier statements are
\Cref{prop:V105-disjoint-factorization,
prop:V105-same-pair-one-operator,
thm:V108-complementary-inherits-V104}.  The bi-heavy assertion for the
same-pair case is
\eqref{eq:V109-V105-bi-heavy}.
\end{proof}

\begin{theorem}[The same-pair family has no independent residual]
\label{thm:V127-no-same-pair-residual}
Refine \(\mathfrak J_{{\rm S},o}\) by the V113 row comparison and
third-row incidence priority.  Every resulting block is either:
\begin{enumerate}[label=\textup{(\roman*)},leftmargin=4.2em]
\item already paid by the row-comparable physical theorem;
\item assigned to a private-mean, cut-heavy, or endpoint-pair payment;
\item a disjoint two-pair carrier;
\item a complementary same-coordinate two-pair carrier; or
\item a genuine off-cut carrier, now paid by V126.
\end{enumerate}
Therefore, after paid blocks are deleted,
\begin{equation}
 \boxed{
 \mathfrak J_{{\rm S},o}^{\rm unpd}
 \subseteq
 \mathfrak J_{{\rm D},o}^{(2)}
 \mathbin{\dot\cup}
 \mathfrak J_{{\rm C},o}^{(2)}.}
 \label{eq:V127-same-returns-two-pair}
\end{equation}
Here the superscript only records that the two-pair family arose after
the same-pair third-row refinement.
\end{theorem}

\begin{proof}
The row-comparable part is paid by
\Cref{thm:V113-physical-row-comparable}.  On a noncomparable block,
\Cref{lem:V113-third-row-alternative} and
\Cref{prop:V113-third-row-disposition} remove the first four incidence
alternatives.  The heavy puncture-pair exit has the disjoint-or-complementary
classification of
\Cref{thm:V114-puncture-pair-dichotomy}; its exact carriers are listed in
\Cref{cor:V114-carrier-disposition}.  The only other exit is the genuine
off-cut carrier by
\Cref{thm:V114-noncomparable-disposition}, and it is paid by
\Cref{thm:V126-complete-offcut}.  The V114 no-return theorem excludes a
second same-pair exit.  All refinements are made by the priority rule of
\Cref{lem:V127-safe-refinement}.
\end{proof}

\begin{corollary}[Exact form of the fixed-orientation JREC residual]
\label{cor:V127-fixed-orientation-residual}
After all proved payments through V126 are applied, every fixed-orientation
JREC residual block belongs to one of two physical geometries:
\begin{equation}
 \boxed{
 \mathfrak J_{{\rm res},o}^{\,127}
 =
 \mathfrak J_{{\rm D},o}^{\,127}
 \mathbin{\dot\cup}
 \mathfrak J_{{\rm C},o}^{\,127}.}
 \label{eq:V127-two-family-residual}
\end{equation}
The first contains two heavy pair banks on disjoint coordinate sets; the
second contains the exact complementary double-defect factor on a common
coordinate bank.  There is no remaining private-mean, cut-heavy,
selected-pair, disconnected, standalone same-pair, or off-cut family.
\end{corollary}

\begin{proof}
Combine
\Cref{thm:V127-after-offcut,lem:V127-three-side-families,
thm:V127-no-same-pair-residual}, and absorb the two descendants of
\(\mathfrak J_{{\rm S},o}\) into the matching two-pair families.  Their
outer labels remain disjoint because the refinement records their
ancestry.
\end{proof}

\section{What is proved and unproved on the two survivors}
\label{sec:V127-two-pair-status}

\begin{theorem}[Exact available estimates on the surviving families]
\label{thm:V127-available-two-pair-estimates}
The two families in
\eqref{eq:V127-two-family-residual} have the following rigorous
fixed-orientation estimates.
\begin{enumerate}[label=\textup{(\roman*)},leftmargin=4.2em]
\item On \(\mathfrak J_{{\rm D},o}^{\,127}\), the carrier factors as
      \[
       Q_e\otimes Q_f\otimes M
      \]
      and obeys the exact one-minimum two-level law
      \eqref{eq:V104-one-minimum-two-level}.
\item On \(\mathfrak J_{{\rm C},o}^{\,127}\), every physical downstream
      event is a subeffect of the complete upstream pair-level/core
      carrier and obeys
      \eqref{eq:V108-complementary-V104-gate}.
\item In the near-unit test family, both estimates imply decay under
      bounded coefficient and orientation prefactors and
      \begin{equation}
       R_e,R_f\longrightarrow\infty,
       \qquad
       \frac{R_eR_f}{N^2}\longrightarrow\infty.
       \label{eq:V127-known-two-pair-scale}
      \end{equation}
\end{enumerate}
None of the cited theorems proves
\eqref{eq:V127-known-two-pair-scale} from the physical pair-heaviness
certificates.
\end{theorem}

\begin{proof}
Part \textup{(i)} is
\Cref{prop:V105-disjoint-factorization,
thm:V104-one-minimum-two-level}.  Part \textup{(ii)} is
\Cref{thm:V108-Racah-subeffects,
thm:V108-complementary-inherits-V104}.  The conditional scale in
\textup{(iii)} is
\Cref{cor:V104-linear-threshold,
cor:V108-complementary-linear-threshold,
thm:V114-conditional-near-unit}.  The failure of pair heaviness to force
it is stated in
\Cref{fire:V114-scale-not-forced}.
\end{proof}

\begin{firewall}[Asymptotic diagnostic versus the finite-scale target]
\label{fire:V127-asymptotic-not-JREC}
Condition \eqref{eq:V127-known-two-pair-scale} is a sufficient decay
criterion for the \(F\otimes R_N\) near-unit diagnostic.  It is not the
finite-scale marked-tree inequality required by JREC.  Conversely,
failure to derive this sufficient condition from mass data is not a
counterexample to JREC.  The remaining proof must estimate the complete
physical two-pair carrier with its actual Wilson coefficients and graph
weights.
\end{firewall}

\section{The two orientations are no longer an independent analytic gate}
\label{sec:V127-orientations}

Let \(o\in\{ab,ac\}\) denote the two V94 retained-row orientations and
let \(\mathsf M_o^{(2{\rm p})}\ge0\) be the integrated residual effect
after the reduction
\eqref{eq:V127-two-family-residual}.

\begin{lemma}[Fixed-orientation payments assemble without commutation]
\label{lem:V127-orientation-assembly}
Suppose that, for every original-row product density \(\rho\),
\begin{equation}
 \operatorname{Tr}(\rho\mathsf M_{ab}^{(2{\rm p})})
 \le W_{ab},
 \qquad
 \operatorname{Tr}(\rho\mathsf M_{ac}^{(2{\rm p})})
 \le W_{ac}.
 \label{eq:V127-two-orientation-hypotheses}
\end{equation}
Then
\begin{equation}
 \boxed{
 \sup_{\rho_{\rm rows}^{\rm prod}}
 \operatorname{Tr}\!\left[
 \rho_{\rm rows}^{\rm prod}
 \bigl(
 \mathsf M_{ab}^{(2{\rm p})}
 +\mathsf M_{ac}^{(2{\rm p})}
 \bigr)\right]
 \le W_{ab}+W_{ac}.}
 \label{eq:V127-two-orientation-assembly}
\end{equation}
No commutation or joint POVM is required.
\end{lemma}

\begin{proof}
For each fixed product density, linearity of the trace and
\eqref{eq:V127-two-orientation-hypotheses} give the bound by
\(W_{ab}+W_{ac}\).  Take the supremum.  This is the integrated positive
form of \eqref{eq:V94-two-orientation-trivial-bound}.
\end{proof}

\begin{corollary}[No separate overlap-saving mandate]
\label{cor:V127-no-overlap-mandate}
If the two physical fixed-orientation carriers in
\eqref{eq:V127-two-family-residual} satisfy their respective marked-tree
bounds, their noncommuting V94 sum satisfies the sum of those bounds.
Thus a new orientation-overlap saving is not logically necessary for
JREC closure.  The unpaid analytic work is entirely in the two
fixed-orientation two-pair carrier estimates.
\end{corollary}

\begin{proof}
Use the marked-tree envelopes as \(W_{ab},W_{ac}\) in
\Cref{lem:V127-orientation-assembly}.  The V98 physical construction
ensures that each fixed-orientation estimate is taken on the complete
absolute carrier before the product-state supremum.
\end{proof}

\begin{remark}[Why V95 remains correct]
\label{rem:V127-V95-status}
V95 correctly proved that overlap alone cannot create a missing
inverse-dimension debit.  V127 does not contradict that no-go.  It shows
that no overlap improvement is needed once the missing debit is supplied
inside each fixed-orientation physical carrier.
\end{remark}

\section{The exact V128 acquisition}
\label{sec:V127-next}

The remaining two-pair theorem must be finite-scale and physical.  It
must retain the two weighted pair multipliers and the signed row-allocated
core until one complete V98 outer block has been formed.  The V126
maximum-square repair suggests testing a one-sided positive completion,
but capping the second pair multiplier by a universal constant may lose
its required graph numerator.  That loss must be calculated, not ignored.

\begin{openproblem}[V128 mandate: finite-scale two-pair carrier]
\label{op:V128-mandate}
\begin{enumerate}
\item Write the exact marked-tree numerator and denominator requested on
      \(\mathfrak J_{{\rm D},o}^{\,127}\) and
      \(\mathfrak J_{{\rm C},o}^{\,127}\), separately.
\item Keep both physical pair coefficients through the V110 signed
      row-allocation and the V98 absolute carrier.
\item Test both valid one-sided positive completions.  Record explicitly
      whether bounding the other pair in operator norm preserves or loses
      its graph numerator.
\item If one numerator is lost, seek a joint layer-cake estimate for the
      product multiplier before the core shell is replaced by its trace
      or product norm.
\item Use the exact V31--V34 Wilson moment hierarchy on hot output
      sectors and V53 cold rank on the complementary sectors, but do not
      infer an effective-scale lower bound from pair heaviness.
\item Treat the disjoint-coordinate and same-coordinate complementary
      carriers separately; the latter must use the V108 upstream
      subeffect order and may not multiply post-recoupling capacities.
\item Close both fixed orientations before invoking
      \Cref{lem:V127-orientation-assembly}.
\end{enumerate}
\end{openproblem}

\section{V127 ledger}
\label{sec:V127-ledger}

\begin{theorem}[Integrated V127 statement]
\label{thm:V127-integrated}
Preserving V1--V126, V127 proves:
\begin{enumerate}[label=\textup{(V127.\arabic*)},leftmargin=7.5em]
\item safe orthogonal refinement of complete V98 outer labels,
      \eqref{eq:V127-safe-refinement};
\item exact deletion of the complete off-cut family from the JREC
      residual;
\item the disjoint/same/complementary refinement of the side family;
\item return of every unpaid same-pair descendant to a disjoint or
      complementary two-pair family;
\item the exact two-family fixed-orientation residual
      \eqref{eq:V127-two-family-residual};
\item the precise conditional status of the V104/V108 effective-scale
      estimates; and
\item assembly of successful fixed-orientation estimates without
      commutation, a joint POVM, or an additional overlap saving.
\end{enumerate}
\end{theorem}

\begin{proof}
The items are
\Cref{lem:V127-safe-refinement,
thm:V127-after-offcut,
lem:V127-three-side-families,
thm:V127-no-same-pair-residual,
cor:V127-fixed-orientation-residual,
thm:V127-available-two-pair-estimates,
lem:V127-orientation-assembly}.
\end{proof}

\begin{firewall}[Claim boundary after V127]
\label{fire:V127-current-boundary}
V127 proves that the sole JREC residual left by the established programme
is the pair of fixed-orientation, physical two-pair side carriers in
\eqref{eq:V127-two-family-residual}.  It does not prove their uniform
finite-scale marked-tree estimate, force the V104/V108 effective scales,
prove JREC, complete quartic or all-order MTA, construct continuum
four-dimensional Yang--Mills theory, or prove a positive spectral gap.
The full Yang--Mills existence and mass-gap problem remains open.
\end{firewall}

\section{Append-only record for V127}
\label{sec:V127-record}

V127 was appended to the complete V126 manuscript.  The exact V126
parent had \(84270\) newline-delimited source records, \(2921102\)
bytes, and SHA--256 digest
\[
 \texttt{b467342e6e49128653500d9a7b432a9e9a37f85869cf0789394fc483ab33899d}.
\]
No V126 mathematical content was changed.  The sole final
\verb|\end{document}| is restored below.

% =============================================================================
% END APPENDED VERSION V127 MATERIAL
% =============================================================================

% =============================================================================
% BEGIN APPENDED VERSION V128 MATERIAL
% =============================================================================

\chapter{V128: Row-Incidence Correction and Finite-Scale Closure of
the Row-Disjoint Two-Pair Sectors}
\label{ch:V128}

\section{Purpose and upstream correction}
\label{sec:V128-purpose}

V127 reduced the JREC residual to disjoint-coordinate and complementary
two-pair side carriers.  Its V128 mandate asked whether the older V58
coefficient tax already retains both graph numerators at finite scale.
The answer depends on row incidence.

The V58 product theorem applies to a row-disjoint pair partition such as
\((oa)\mid(bc)\).  V86 proved that two pair banks may occupy disjoint
gauge-coordinate sets and nevertheless share an original row; in that
case their positive product-state capacities cannot be multiplied.
Accordingly, the V114 assertion that the two taxes compose on every
disjoint-coordinate alternative was too broad when the heavy third-row
pair is \(ca\) or \(cb\).  It is valid for row-disjoint pair labels and
for the complementary local partition \(ab\mid co\), but not merely from
coordinate-support disjointness.

After making this correction, V58 does give a genuine advance.  It closes
the row-disjoint part of the V127 disjoint family and the complementary
family directly at finite scale, with both pair numerators retained.
The only two-pair residual is therefore a pair of disjoint coordinate
banks whose labels share at least one original row.  This is exactly the
positive carrier isolated in V86 and analyzed conditionally by V104.

\section{Two independent notions of disjointness}
\label{sec:V128-two-disjointness}

Let \(L_e,L_f\) be two selected pair banks, with unordered row labels
\[
 e=\{i,j\},
 \qquad
 f=\{k,\ell\}.
\]
Write \({\rm supp}_{\rm g}(L)\) for the gauge-coordinate support of a bank.

\begin{definition}[Coordinate-disjoint and row-disjoint pairs]
\label{def:V128-disjointness}
The banks are coordinate-disjoint if
\[
 {\rm supp}_{\rm g}(L_e)\cap{\rm supp}_{\rm g}(L_f)=\varnothing.
\]
Their labels are row-disjoint if
\[
 e\cap f=\varnothing.
\]
Coordinate disjointness tensor-separates the gauge-coordinate carriers.
Row disjointness additionally partitions the original input rows between
the two bilinear compressions.
\end{definition}

\begin{lemma}[Coordinate disjointness does not imply row disjointness]
\label{lem:V128-coordinate-not-row}
For the old pair \(e=ab\) and a heavy third-row pair \(f=cj\) with
\(j\in\{o,a,b\}\):
\begin{enumerate}[label=\textup{(\roman*)},leftmargin=4.2em]
\item if \(j=o\), the labels \(ab\) and \(co\) are row-disjoint;
\item if \(j=a\) or \(j=b\), the labels share one original row, even
      though V114 forces their coordinate supports to be disjoint.
\end{enumerate}
\end{lemma}

\begin{proof}
For \(j=o\), the sets \(\{a,b\}\) and \(\{c,o\}\) have empty
intersection.  For \(j=a\), their intersection is \(\{a\}\); for
\(j=b\), it is \(\{b\}\).  Coordinate disjointness in the latter two
cases is
\Cref{lem:V114-no-same-label}.  It changes no row label.
\end{proof}

\begin{proposition}[Correct scope of V58 composition]
\label{prop:V128-V58-scope}
The hypotheses
\[
 {\rm supp}_{\rm g}(L_e)\cap{\rm supp}_{\rm g}(L_f)=\varnothing,
 \qquad
 e\cap f=\varnothing
\]
place the two pair compressions in the row-factorized form of
\Cref{lem:V58-tax-composition}.  Coordinate disjointness without
\(e\cap f=\varnothing\) does not do so.
\end{proposition}

\begin{proof}
When \(e\cap f=\varnothing\), the four endpoints can be relabeled as
\(o,a,b,c\) so that the input coefficient is
\[
 B_o\otimes B_a\otimes B_b\otimes B_c
\]
and the two pair compressions act on the row-disjoint factors
\((oa)\mid(bc)\).  Gauge-coordinate disjointness supplies the tensor
factorization of their local integrals.  These are exactly the structural
hypotheses in
\eqref{eq:V58-product-tax}.

If the pair labels share a row, tracing the complementary rows leaves the
second pair operator inside the same retained-row expectation as the
first and the core.  This is
\Cref{prop:V86-overlapping-retained-reduction}.  The conclusion that the
two bilinear taxes cannot be multiplied from V58 is
\Cref{cor:V86-no-overlapping-tax-product}.
\end{proof}

\begin{firewall}[Supersession of one V114 sentence]
\label{fire:V128-V114-supersession}
The disjoint part of
\Cref{prop:V114-trace-tax-disposition} is henceforth used only when the
two pair labels are also row-disjoint.  Its proof cited disjoint
coordinate factors but did not check the row-disjoint hypothesis of
\Cref{lem:V58-tax-composition}.  The structural dichotomy
\Cref{thm:V114-puncture-pair-dichotomy}, the exact carriers
\Cref{cor:V114-carrier-disposition}, and the V104 conditional positive
estimate remain valid.  Only the overbroad unconditional composition
claim is restricted.
\end{firewall}

\section{A finite-scale row-disjoint two-pair payment}
\label{sec:V128-row-disjoint-payment}

For \(g\in\{e,f\}\), let \(h_g\) be its pair mass, let \(v_g\) be a
selected endpoint, and let
\(\mathfrak A_g\) be the V58 normalized weighted formation response.
Assume
\begin{equation}
 x_{v_g}(L_g)\ge\varepsilon\Xi_{v_g}.
 \label{eq:V128-selected-heavy}
\end{equation}
On a genuine sideways endpoint one also has
\begin{equation}
 s_\alpha\Xi_{v_g}>M.
 \label{eq:V128-sideways-hot}
\end{equation}

\begin{lemma}[Two hot selected taxes give the marked denominator]
\label{lem:V128-two-hot-taxes}
If both pair labels are row-disjoint and satisfy
\eqref{eq:V128-selected-heavy}--\eqref{eq:V128-sideways-hot}, then
\begin{equation}
 \boxed{
 \mathfrak A_e\mathfrak A_f
 \le
 C_{\varepsilon,M}
 \frac{h_eh_f}{\Xi_{v_e}\Xi_{v_f}}.}
 \label{eq:V128-two-hot-marked}
\end{equation}
Both pair numerators occur exactly once.
\end{lemma}

\begin{proof}
For each pair, the selected-heavy one-edge tax
\eqref{eq:V58-hot-edge-tax} gives
\[
 \mathfrak A_g
 \le
 C_\varepsilon
 \frac{h_g}{s_\alpha\Xi_{v_g}^2}.
\]
The row-disjoint hypotheses permit multiplication by
\Cref{lem:V58-tax-composition}.  Hence
\[
 \mathfrak A_e\mathfrak A_f
 \le
 C_\varepsilon
 \frac{h_eh_f}
 {s_\alpha^2\Xi_{v_e}^2\Xi_{v_f}^2}.
\]
Rewrite the last scalar factor as
\[
 \frac1{\Xi_{v_e}\Xi_{v_f}}\,
 \frac1{s_\alpha\Xi_{v_e}}\,
 \frac1{s_\alpha\Xi_{v_f}}.
\]
Both final factors are at most \(M^{-1}\) by
\eqref{eq:V128-sideways-hot}.  This proves
\eqref{eq:V128-two-hot-marked}.
\end{proof}

\begin{theorem}[Finite-scale V58-admissible side payment]
\label{thm:V128-row-disjoint-side-payment}
Fix one complete V98 outer subfamily for which:
\begin{enumerate}[label=\textup{(\roman*)},leftmargin=4.2em]
\item the two selected pair banks are row-disjoint and are retained
      before final recoupling;
\item each is selected-heavy and hot at its certified sideways endpoint;
\item every remaining pair, mean, resolvent, and punctured-core factor is
      retained in the V58 admissible auxiliary response; and
\item the exact row-allocated core response has the V57/V110 bound
      \(C H_o\) after the single resolvent allocation, where \(H_o\) is
      the corresponding marked core numerator.
\end{enumerate}
Then the complete coefficient contribution of that outer subfamily
satisfies
\begin{equation}
 \boxed{
 \mathcal R_{e,f;o}
 \le
 C_{\varepsilon,M}
 \frac{h_eH_oh_f}
      {\Xi_{v_e}\Xi_{v_f}}
 \prod_r\|B_r\|_1.}
 \label{eq:V128-row-disjoint-side-payment}
\end{equation}
The estimate is finite-scale and contains no representation dimension,
effective number, output count, or orientation factor.
\end{theorem}

\begin{proof}
Before final recoupling, use the row-factorized product of the two
multiplicity-resolved pair compressions.  The proof of
\Cref{lem:V58-three-block-response} applies verbatim to a restricted
orthogonal outer subfamily: on low outputs allocate the single resolvent
to the row-allocated core; on pair-hot outputs use the weighted part of
that pair response; on core-hot outputs use the V57 quotient; and on
ordinary spectators use their first moment.  Therefore
\[
 \mathcal R_{e,f;o}
 \le
 C H_o\,\mathfrak A_e\mathfrak A_f
 \prod_r\|B_r\|_1.
\]
Insert
\eqref{eq:V128-two-hot-marked}.  Final recoupling and aggregate orthogonal
pinching are contractions into a trace-class direct sum by
\Cref{lem:V58-tax-composition}; restricting the output index set only
decreases that sum.  Hence no output or multiplicity count is introduced,
and \eqref{eq:V128-row-disjoint-side-payment} follows.
\end{proof}

\begin{remark}[Thermal endpoints]
\label{rem:V128-thermal-extension}
If one certified endpoint is thermal instead of hot, combine
\eqref{eq:V58-thermal-edge-cap} with the hot tax on the other endpoint.
The scalar comparison in the proof of
\Cref{thm:V58-selected-RACC} again gives the marked denominator.  The
surviving V127 side endpoints are hot by definition, so
\Cref{thm:V128-row-disjoint-side-payment} states only the needed case.
\end{remark}

\section{Application to the V127 two-pair families}
\label{sec:V128-application}

Refine the disjoint-coordinate family of
\eqref{eq:V127-two-family-residual} by its fixed pair labels:
\begin{equation}
 \mathfrak J_{{\rm D},o}^{\,127}
 =
 \mathfrak J_{{\rm D,rd},o}
 \mathbin{\dot\cup}
 \mathfrak J_{{\rm D,sh},o},
 \label{eq:V128-D-split}
\end{equation}
where \({\rm rd}\) means row-disjoint and \({\rm sh}\) means that the
labels share at least one original row.  This is a deterministic
refinement of complete outer labels.

\begin{corollary}[Row-disjoint disjoint-coordinate family is paid]
\label{cor:V128-D-row-disjoint-paid}
Every V58-admissible block in
\(\mathfrak J_{{\rm D,rd},o}\) satisfies its finite-scale marked-tree
weight and may be moved to the paid outer family.
\end{corollary}

\begin{proof}
The coordinate tensor factor is
\Cref{prop:V105-disjoint-factorization}; row disjointness supplies the
additional V58 hypothesis by
\Cref{prop:V128-V58-scope}.  Both selected masses are fixed fractions of
their sideways row masses by
\eqref{eq:V105-V102-disjoint-mass}.  The row-allocated core estimate is
\Cref{thm:V110-separate-influence}.  Apply
\Cref{thm:V128-row-disjoint-side-payment}, and use
\Cref{lem:V127-safe-refinement} for physical aggregation.
\end{proof}

\begin{theorem}[Complementary two-pair family is paid at finite scale]
\label{thm:V128-complementary-paid}
Every V58-admissible complementary-pair block in
\(\mathfrak J_{{\rm C},o}^{\,127}\) satisfies its finite-scale marked-tree
weight and may be moved to the paid outer family.
\end{theorem}

\begin{proof}
At each active coordinate the complementary local partition has two
row-disjoint blocks.  The exact V56 collision identity gives the local
double-defect factor, and
\Cref{lem:V58-tax-composition} composes the two row-factorized taxes
before recoupling.  The V108 downstream Racah/output selection is a
positive subeffect of this upstream carrier by
\Cref{thm:V108-Racah-subeffects}, so it cannot increase either partial
marginal or the coefficient sum.  The selected-heavy mass fractions are
\eqref{eq:V105-quarter-pure-overlap}, or
\eqref{eq:V114-complementary-heavy} for a descendant of the old
same-pair family.  Apply
\Cref{thm:V128-row-disjoint-side-payment}.  Aggregate complete V98 blocks
by \Cref{lem:V127-safe-refinement}.
\end{proof}

\begin{firewall}[What is not inferred from V58]
\label{fire:V128-no-shared-row-product}
The proof of
\Cref{cor:V128-D-row-disjoint-paid} does not apply to
\(\mathfrak J_{{\rm D,sh},o}\).  Those pair banks have disjoint
gauge-coordinate supports, but a density matrix on their common original
row may entangle the two carrier factors.  Their exact positive object is
the V86 overlapping product-state capacity, and their available joint
spectral estimate remains V104.  No two marginal V53 taxes are multiplied
on that family.
\end{firewall}

\begin{corollary}[Exact JREC residual after V128]
\label{cor:V128-exact-residual}
After all established payments through V128, the fixed-orientation JREC
residual is
\begin{equation}
 \boxed{
 \mathfrak J_{{\rm res},o}^{\,128}
 =
 \mathfrak J_{{\rm D,sh},o}.}
 \label{eq:V128-shared-row-only}
\end{equation}
It consists of two selected-heavy pair banks on disjoint gauge-coordinate
supports whose labels share at least one original row, retained with the
complete signed row-allocated core and the complete V98 outer label.
\end{corollary}

\begin{proof}
V127 left exactly the D and C families in
\eqref{eq:V127-two-family-residual}.  Split D by
\eqref{eq:V128-D-split}.  Move the row-disjoint part by
\Cref{cor:V128-D-row-disjoint-paid} and move C by
\Cref{thm:V128-complementary-paid}.  The remaining index set is exactly
\(\mathfrak J_{{\rm D,sh},o}\).  Disjointness follows from the
deterministic outer-label refinement.
\end{proof}

\section{The surviving shared-row carrier}
\label{sec:V128-shared-row-carrier}

On a remaining block, let the two positive weighted pair multipliers be
\(\mathbf A_e,\mathbf A_f\), and let \(\mathbf R_{o,j}\) be the exact
signed row-allocated core remainder.  The physical blocks have the form
\[
 L_{o,j}
 =
 d_{e,j}d_{f,j}\mathbf R_{o,j}
\]
only after both pair outputs and the full outer label have been resolved.
Their absolute carrier is aggregated before the product-state supremum.

\begin{proposition}[Two valid one-sided bounds and their numerator loss]
\label{prop:V128-one-sided-numerator-loss}
Let
\[
 A_{e,*}:=\|\mathbf A_e\|_{\rm op},
 \qquad
 A_{f,*}:=\|\mathbf A_f\|_{\rm op},
 \qquad
 \kappa_e:=\|\mathbf A_e\|_{\rm prod},
 \qquad
 \kappa_f:=\|\mathbf A_f\|_{\rm prod}.
\]
For the shared-row joint capacity \(\mathcal K_{e,f}\), positive order
gives
\begin{equation}
 \mathcal K_{e,f}
 \le
 \min\{A_{e,*}\kappa_f,A_{f,*}\kappa_e\}.
 \label{eq:V128-two-one-sided}
\end{equation}
Using only the universal operator caps and one selected-heavy capacity
gives either
\[
 C\frac{h_e}{s_\alpha\Xi_{v_e}^2}
 \quad\hbox{or}\quad
 C\frac{h_f}{s_\alpha\Xi_{v_f}^2}.
\]
Thus this route retains only one of the two pair numerators.  It cannot
by itself prove a target whose numerator is \(h_eh_f\) uniformly when the
other pair mass may tend to zero.
\end{proposition}

\begin{proof}
The proof of
\Cref{lem:V118-one-sided-star} uses only positive order and the
within-row spectator quotient, and therefore applies to any two
disjoint-coordinate pair banks sharing an original row.  This gives
\eqref{eq:V128-two-one-sided}.  The V52 weighted-defect cap gives
\(A_{e,*},A_{f,*}\le C\).  Apply
\Cref{cor:V53-one-edge-tax} to either selected-heavy pair capacity.
The resulting displayed bounds contain respectively \(h_e\) or \(h_f\),
not their product.  Since the standing mass inequalities impose no
positive lower bound on the omitted pair mass, a uniform factor
\(h_eh_f\) cannot be recovered by scalar comparison.
\end{proof}

\begin{assessment}[The noncircular V129 target]
\label{ass:V128-V129-target}
V128 has removed every sector on which V58 legitimately multiplies the
two taxes.  The remaining problem is not to prove again that each
marginal pair capacity is small.  It is to prove a numerator-preserving
joint estimate for
\(\mathcal K_{e,f}\), or to reorganize the graph so that only one of
\(h_e,h_f\) is requested on this carrier.  The exact V104 two-level
surface is the correct positive starting object.  Any proposed bound must
be tested against the V104 rank-one top-fraction regression and the V86
spectator-entanglement mechanism.
\end{assessment}

\begin{openproblem}[V129 mandate: shared-row numerator preservation]
\label{op:V129-mandate}
\begin{enumerate}
\item Fix the shared row and write the two pair operators as acting on
      its two disjoint within-row coordinate factors, with the other
      endpoint rows and the signed core kept explicit.
\item Write the exact marked-tree numerator requested on that graph
      before taking a product-state supremum.
\item Apply a common layer cake to
      \(\mathbf A_e\otimes\mathbf A_f\) and keep one density matrix on the
      shared original row; do not multiply marginal capacities.
\item Determine whether the Wilson formation coefficients give a
      levelwise operator bound proportional to \(h_eh_f\), rather than
      only separate integrated bounds proportional to \(h_e\) and
      \(h_f\).
\item If such a bound fails, use an explicit finite-dimensional shared-row
      regression to decide whether the marked-tree routing must be
      changed.
\item Keep the complete signed V110 core through the estimate and
      aggregate V98 outer labels only after the joint bound is obtained.
\item Once both fixed orientations close, use
      \Cref{lem:V127-orientation-assembly}; do not seek a joint Racah POVM.
\end{enumerate}
\end{openproblem}

\section{V128 ledger}
\label{sec:V128-ledger}

\begin{theorem}[Integrated V128 statement]
\label{thm:V128-integrated}
Preserving V1--V127, except for the explicit scope correction in
\Cref{fire:V128-V114-supersession}, V128 proves:
\begin{enumerate}[label=\textup{(V128.\arabic*)},leftmargin=7.5em]
\item the distinction between coordinate and row disjointness;
\item the exact row-incidence split of the V114 third-row pairs;
\item the finite-scale two-hot-pair denominator
      \eqref{eq:V128-two-hot-marked};
\item finite-scale payment of the row-disjoint disjoint-coordinate and
      complementary two-pair families;
\item reduction of the fixed-orientation JREC residual to the shared-row
      family \eqref{eq:V128-shared-row-only}; and
\item the exact one-sided positive bounds and their loss of one graph
      numerator.
\end{enumerate}
\end{theorem}

\begin{proof}
The items are
\Cref{lem:V128-coordinate-not-row,
prop:V128-V58-scope,
lem:V128-two-hot-taxes,
thm:V128-row-disjoint-side-payment,
cor:V128-D-row-disjoint-paid,
thm:V128-complementary-paid,
cor:V128-exact-residual,
prop:V128-one-sided-numerator-loss}.
\end{proof}

\begin{firewall}[Claim boundary after V128]
\label{fire:V128-current-boundary}
V128 closes only the V58-admissible row-disjoint two-pair side sectors.
It does not prove the numerator-preserving joint capacity of the
shared-row two-pair carrier, close the full fixed-orientation JREC
residual, prove JREC, complete quartic or all-order MTA, construct
continuum four-dimensional Yang--Mills theory, or prove a positive mass
gap.  The full Yang--Mills existence and mass-gap problem remains open.
\end{firewall}

\section{Append-only record for V128}
\label{sec:V128-record}

V128 was appended to the complete V127 manuscript.  The exact V127
parent had \(84741\) newline-delimited source records, \(2939486\)
bytes, and SHA--256 digest
\[
 \texttt{bfa86569828daaf7eaf800e52ddf8560f88e7e25e3fb69aa5a13cfb3ea14b95f}.
\]
No V127 mathematical content was changed.  The sole final
\verb|\end{document}| is restored below.

% =============================================================================
% END APPENDED VERSION V128 MATERIAL
% =============================================================================

% =============================================================================
% BEGIN APPENDED VERSION V129 MATERIAL
% =============================================================================

\chapter{V129: Exact One-Shared-Row Product Law and Closure of the
Fixed-Monomial JREC Residual}
\label{ch:V129}

\section{Purpose}
\label{sec:V129-purpose}

V128 reduced the fixed-orientation residual to two selected-heavy pair
banks on disjoint gauge-coordinate supports whose labels share an
original row.  V86 correctly required those pair operators to remain in
the same positive expectation as the unresolved core.  The V110 operator
bound now permits that core to be replaced by a scalar multiple of the
identity before the pair capacity is taken.

For two pair banks sharing exactly one row, the remaining joint product
capacity factorizes exactly.  The shared row may be arbitrarily entangled
between its two within-row coordinate factors; after the two outer-row
states are fixed, it sees a tensor product of two positive operators, and
the operator norm of that tensor product is the product of their norms.
Thus no independence of the shared-row state is assumed.

Complete same-label regrouping handles the apparent case in which the
two labels share both endpoints: their coordinate subbanks form one
bi-heavy pair bank and return to the V109/V127 same-pair disposition.
Consequently the genuine V128 residual shares exactly one row, and the
new factorization closes it at the finite-scale marked-tree weight.

\section{The exact one-shared-row capacity identity}
\label{sec:V129-capacity-identity}

Let \(H_A,H_C,H_{B_1},H_{B_2}\) be finite-dimensional Hilbert spaces.
The three original rows are
\[
 H_A,\qquad H_B:=H_{B_1}\otimes H_{B_2},\qquad H_C.
\]
Let
\[
 E\in\mathcal B(H_A\otimes H_{B_1})_+,
 \qquad
 F\in\mathcal B(H_{B_2}\otimes H_C)_+.
\]
Define
\begin{align}
 \kappa_{A|B_1}(E)
 &:=
 \sup_{\rho_A,\rho_{B_1}}
 \operatorname{Tr}\!\left[E(\rho_A\otimes\rho_{B_1})\right],
 \label{eq:V129-left-pair-cap}\\
 \kappa_{B_2|C}(F)
 &:=
 \sup_{\rho_{B_2},\rho_C}
 \operatorname{Tr}\!\left[F(\rho_{B_2}\otimes\rho_C)\right],
 \label{eq:V129-right-pair-cap}\\
 \mathcal K_{A|B|C}(E,F)
 &:=
 \sup_{\rho_A,\rho_B,\rho_C}
 \operatorname{Tr}\!\left[
 (E\otimes F)(\rho_A\otimes\rho_B\otimes\rho_C)
 \right],
 \label{eq:V129-three-row-cap}
\end{align}
where every supremum ranges over density matrices and
\(\rho_B\) may be entangled across \(B_1\mid B_2\).

\begin{theorem}[Exact one-shared-row product law]
\label{thm:V129-one-shared-row-product}
For all positive \(E,F\) as above,
\begin{equation}
 \boxed{
 \mathcal K_{A|B|C}(E,F)
 =
 \kappa_{A|B_1}(E)\,
 \kappa_{B_2|C}(F).}
 \label{eq:V129-one-shared-row-product}
\end{equation}
\end{theorem}

\begin{proof}
Fix \(\rho_A,\rho_C\), and define the positive partial expectations
\begin{align*}
 X_{\rho_A}
 &:=
 \operatorname{Tr}_A\!\left[
 (\rho_A^{1/2}\otimes I_{B_1})
 E
 (\rho_A^{1/2}\otimes I_{B_1})
 \right]\in\mathcal B(H_{B_1})_+,\\
 Y_{\rho_C}
 &:=
 \operatorname{Tr}_C\!\left[
 (I_{B_2}\otimes\rho_C^{1/2})
 F
 (I_{B_2}\otimes\rho_C^{1/2})
 \right]\in\mathcal B(H_{B_2})_+.
\end{align*}
For every density matrix \(\rho_B\) on
\(H_{B_1}\otimes H_{B_2}\), cyclicity of the finite-dimensional trace
and the defining property of partial trace give
\[
 \operatorname{Tr}\!\left[
 (E\otimes F)(\rho_A\otimes\rho_B\otimes\rho_C)
 \right]
 =
 \operatorname{Tr}\!\left[
 (X_{\rho_A}\otimes Y_{\rho_C})\rho_B
 \right].
\]
The supremum over all \(\rho_B\), including entangled ones, is
\[
 \|X_{\rho_A}\otimes Y_{\rho_C}\|_{\rm op}
 =
 \|X_{\rho_A}\|_{\rm op}\,
 \|Y_{\rho_C}\|_{\rm op}.
\]
Moreover,
\begin{align*}
 \sup_{\rho_A}\|X_{\rho_A}\|_{\rm op}
 &=
 \sup_{\rho_A,\rho_{B_1}}
 \operatorname{Tr}\!\left[E(\rho_A\otimes\rho_{B_1})\right]
 =
 \kappa_{A|B_1}(E),\\
 \sup_{\rho_C}\|Y_{\rho_C}\|_{\rm op}
 &=
 \kappa_{B_2|C}(F).
\end{align*}
The variables \(\rho_A,\rho_C\) are independent, so taking their
suprema proves the upper bound in
\eqref{eq:V129-one-shared-row-product}.

For the reverse bound, choose outer-row states within an arbitrary
\(\epsilon>0\) of the two suprema and choose top eigenvector states of
\(X_{\rho_A}\) and \(Y_{\rho_C}\).  Their tensor product is an admissible
state on the shared row \(B_1\otimes B_2\) and attains the product within
\(O(\epsilon)\).  Let \(\epsilon\downarrow0\).
\end{proof}

\begin{remark}[Entanglement is allowed but gives no gain here]
\label{rem:V129-entanglement-allowed}
The theorem does not replace the shared-row state by a product state.
It maximizes over every \(\rho_B\).  Positivity and the exact tensor
factor \(X_{\rho_A}\otimes Y_{\rho_C}\) imply that an extremizing product
of top eigenvectors already reaches its operator norm.
\end{remark}

\begin{firewall}[Why this does not contradict V86]
\label{fire:V129-V86-compatible}
V86 keeps the pair operators in the same expectation as a nontrivial
retained core operator.  Such a core can couple \(B_1\) and \(B_2\), so
the shared row need not see a tensor product and
\Cref{thm:V129-one-shared-row-product} cannot yet be applied.  V129 uses
the V110 uniform operator bound to dominate the complete core by a scalar
multiple of the identity first.  The tensor-product structure then
becomes exact.  No claim in V86 is contradicted.
\end{firewall}

\section{Positive completion before applying the product law}
\label{sec:V129-positive-completion}

Let the two disjoint-coordinate pair banks have positive weighted
multipliers
\[
 \mathbf A_e=\sum_\lambda a_\lambda P_\lambda^e,
 \qquad
 \mathbf A_f=\sum_\mu b_\mu P_\mu^f.
\]
Let \(\mathfrak J_X\) be any complete V98 outer subfamily refining both
pair outputs.  On block \(j\), write the signed scalar/core factorization
\[
 L_j=d_{\lambda(j)}^e d_{\mu(j)}^f R_j.
\]

\begin{lemma}[Joint positive completion with a scalar core cap]
\label{lem:V129-joint-completion}
Suppose
\[
 |d_\lambda^e|\le a_\lambda,
 \qquad
 |d_\mu^f|\le b_\mu,
 \qquad
 \|R_j\|_{\rm op}\le m
\]
for all resolved labels.  Then the complete absolute carrier
\[
 \mathbf Q_X
 :=
 \sum_{j\in\mathfrak J_X}|V_jL_jV_j^*|
\]
obeys
\begin{equation}
 \boxed{
 \mathbf Q_X
 \preccurlyeq
 m\bigl(
 \mathbf A_e\otimes\mathbf A_f\otimes I_{\rm rem}
 \bigr).}
 \label{eq:V129-joint-completion}
\end{equation}
There is no output or multiplicity count.
\end{lemma}

\begin{proof}
On the range \(P_j=V_jV_j^*\),
\[
 |V_jL_jV_j^*|
 \preccurlyeq
 m\,a_{\lambda(j)}b_{\mu(j)}P_j.
\]
Group the orthogonal complete outer labels by \((\lambda,\mu)\).  Their
sum is a subprojection of
\[
 P_\lambda^e\otimes P_\mu^f\otimes I_{\rm rem}
\]
by the same common-refinement argument as
\Cref{lem:V118-physical-joint-refinement}.  Therefore
\begin{align*}
 \mathbf Q_X
 &\preccurlyeq
 m\sum_{\lambda,\mu}a_\lambda b_\mu
 \bigl(P_\lambda^e\otimes P_\mu^f\otimes I_{\rm rem}\bigr)\\
 &=
 m\bigl(
 \mathbf A_e\otimes\mathbf A_f\otimes I_{\rm rem}
 \bigr).
\end{align*}
This is \eqref{eq:V129-joint-completion}.
\end{proof}

\begin{corollary}[Exact capacity after one-row sharing]
\label{cor:V129-completed-capacity}
If the pair labels share exactly one original row, then
\begin{equation}
 \boxed{
 \|\mathbf Q_X\|_{\rm rows,prod}
 \le
 m\,\kappa_e\kappa_f,}
 \label{eq:V129-completed-capacity}
\end{equation}
where
\(\kappa_e=\|\mathbf A_e\|_{\rm prod}\) and
\(\kappa_f=\|\mathbf A_f\|_{\rm prod}\) are the two ordinary pair
capacities.
\end{corollary}

\begin{proof}
Quotient the identity spectators within each original row by
\Cref{lem:V94-identity-spectator-quotient}.  After relabeling the three
active rows, the remaining multiplier is precisely the operator in
\Cref{thm:V129-one-shared-row-product}.  Apply
\eqref{eq:V129-one-shared-row-product} to
\eqref{eq:V129-joint-completion}.
\end{proof}

\section{Finite-scale marked-tree payment}
\label{sec:V129-marked-payment}

Let the selected endpoints of \(e,f\) be \(v_e,v_f\), and suppose
\[
 x_{v_e}(L_e)\ge\varepsilon\Xi_{v_e},
 \qquad
 x_{v_f}(L_f)\ge\varepsilon\Xi_{v_f}.
\]
The V53 positive capacities then satisfy
\begin{equation}
 \kappa_e
 \le
 C_\varepsilon\frac{h_e}{s_\alpha\Xi_{v_e}^2},
 \qquad
 \kappa_f
 \le
 C_\varepsilon\frac{h_f}{s_\alpha\Xi_{v_f}^2}.
 \label{eq:V129-two-capacity-taxes}
\end{equation}

\begin{theorem}[One-shared-row two-pair marked payment]
\label{thm:V129-shared-row-marked-payment}
Assume:
\begin{enumerate}[label=\textup{(\roman*)},leftmargin=4.2em]
\item the two pair banks have disjoint gauge-coordinate supports and
      share exactly one original row;
\item both selected endpoints are sideways, so
      \(s_\alpha\Xi_{v_e}>M\) and
      \(s_\alpha\Xi_{v_f}>M\);
\item the complete V98 outer blocks refine both pair outputs; and
\item the exact V110 row-allocated core remainder satisfies
      \begin{equation}
       \|R_j\|_{\rm op}\le Cs_\alpha H_o
       \label{eq:V129-core-cap}
      \end{equation}
      on the chosen orientation.
\end{enumerate}
On the cofinal Wilson regime \(0<s_\alpha\le1\), the complete physical
absolute carrier obeys
\begin{equation}
 \boxed{
 \|\mathbf Q_X\|_{\rm rows,prod}
 \le
 C_{\varepsilon,M}
 \frac{h_eH_oh_f}
      {\Xi_{v_e}\Xi_{v_f}}.}
 \label{eq:V129-shared-row-marked-payment}
\end{equation}
\end{theorem}

\begin{proof}
Apply
\Cref{lem:V129-joint-completion,cor:V129-completed-capacity} with
\(m=Cs_\alpha H_o\), and then insert
\eqref{eq:V129-two-capacity-taxes}:
\[
 \|\mathbf Q_X\|_{\rm rows,prod}
 \le
 C_\varepsilon
 \frac{h_eH_oh_f}
 {s_\alpha\Xi_{v_e}^2\Xi_{v_f}^2}.
\]
Relative to the right side of
\eqref{eq:V129-shared-row-marked-payment}, the remaining scalar ratio is
\[
 \frac1{s_\alpha\Xi_{v_e}\Xi_{v_f}}.
\]
The two hot inequalities imply
\[
 \Xi_{v_e}\Xi_{v_f}>\frac{M^2}{s_\alpha^2},
 \qquad
 \frac1{s_\alpha\Xi_{v_e}\Xi_{v_f}}
 <\frac{s_\alpha}{M^2}\le\frac1{M^2}.
\]
Absorb \(M^{-2}\) into the constant.  Every output and copy label was
aggregated in
\eqref{eq:V129-joint-completion}, so there is no channel count.
\end{proof}

\begin{remark}[Mixed thermal/hot extension]
\label{rem:V129-mixed-extension}
If one selected endpoint is thermal, use its V58 thermal coefficient cap
and the other endpoint's selected-heavy capacity before positive
completion.  The same scalar comparison as
\Cref{thm:V58-selected-RACC} proves the marked denominator.  The V128
residual consists of two sideways endpoints, so the hot--hot form above
is the only one needed here.
\end{remark}

\section{Same-label regrouping removes two-row sharing}
\label{sec:V129-same-label}

\begin{lemma}[Two common endpoints form one same-label bank]
\label{lem:V129-two-row-sharing-regroups}
Suppose two coordinate-disjoint selected subbanks have the same unordered
pair label \(e=f=\{i,k\}\).  Under the mandatory same-label regrouping,
their union is one physical pair bank \(L_{ik}\).  If the first subbank
carries a fixed fraction of \(\Xi_i\) and the second a fixed fraction of
\(\Xi_k\), then \(L_{ik}\) is bi-heavy at both endpoints.
\end{lemma}

\begin{proof}
Same-label regrouping is the tensor product over every gauge coordinate
carrying the fixed pair block \(\{i,k\}\).  Disjoint coordinate subsets
therefore form subfactors of one regrouped bank, not two different pair
labels.  The input-\(i\) mass of the union is at least that of the first
subbank, and the input-\(k\) mass is at least that of the second.  The two
fixed lower bounds are precisely the V109 bi-heavy hypotheses.
\end{proof}

\begin{corollary}[Every genuine V128 residual shares exactly one row]
\label{cor:V129-exactly-one-shared}
After complete same-label regrouping and the same-pair disposition of
\Cref{thm:V127-no-same-pair-residual}, every block in
\(\mathfrak J_{{\rm D,sh},o}\) has two distinct unordered pair labels
with exactly one common endpoint.
\end{corollary}

\begin{proof}
Two distinct two-element subsets cannot share both endpoints.  If both
endpoints coincide, apply
\Cref{lem:V129-two-row-sharing-regroups} and assign the block to the
same-pair family.  V127 disposes of that family into paid blocks or
two-pair descendants with distinct labels.  A block remaining in
\(\mathfrak J_{{\rm D,sh},o}\) shares at least one row by definition,
and therefore shares exactly one.
\end{proof}

\section{Closure of the fixed-monomial residual}
\label{sec:V129-fixed-monomial-closure}

\begin{theorem}[Every V128 residual block is paid]
\label{thm:V129-V128-residual-paid}
Every complete outer block in
\(\mathfrak J_{{\rm res},o}^{\,128}\) satisfies its fixed-orientation
finite-scale marked-tree payment.  Hence
\begin{equation}
 \boxed{\mathfrak J_{{\rm res},o}^{\,129}=\varnothing.}
 \label{eq:V129-fixed-orientation-empty}
\end{equation}
\end{theorem}

\begin{proof}
By
\eqref{eq:V128-shared-row-only}, the V128 residual is
\(\mathfrak J_{{\rm D,sh},o}\).  By
\Cref{cor:V129-exactly-one-shared}, its two pair labels share exactly one
row.  The V102/V105 certificates make both banks selected-heavy at their
sideways endpoints, and the V58 sideways condition makes both endpoints
hot.  The exact V110 hybrid allocation supplies
\eqref{eq:V129-core-cap} separately in the chosen orientation.  Apply
\Cref{thm:V129-shared-row-marked-payment} and aggregate the deterministic
outer subfamily by
\Cref{lem:V127-safe-refinement}.
\end{proof}

\begin{corollary}[Two orientations and the fixed-monomial JREC residual]
\label{cor:V129-two-orientation-closed}
For one completely regrouped V60 monomial, all V94/V98 residual physical
outer families have their marked-tree payments in both retained-row
orientations.  Their noncommuting sum obeys the sum of those payments.
\end{corollary}

\begin{proof}
The off-cut family is paid by
\Cref{thm:V126-complete-offcut}.  V127 reduces the side family to its
two-pair subfamilies.  V128 pays the row-disjoint and complementary
subfamilies, while
\Cref{thm:V129-V128-residual-paid} pays the last one-shared-row family.
Apply
\Cref{lem:V127-orientation-assembly} to the two fixed orientations.
\end{proof}

\begin{firewall}[Local closure is not global JREC]
\label{fire:V129-local-not-global}
\Cref{cor:V129-two-orientation-closed} closes the residual carrier for
one completely regrouped V60 collision/puncture monomial under the
standing finite-dimensional programme estimates.  It does not yet sum
all puncture monomials or multiple collision coordinates with a uniform
graph constant.  Therefore V129 does not claim the global JREC theorem.
\end{firewall}

\section{Mandatory V130 acquisition}
\label{sec:V129-next}

The next version is divisible by five.  It must first audit the latest SM
and NS notes, then return to the exact V54/V56 puncture and
multi-collision decompositions.  The goal is to determine whether the
fixed-monomial bounds of V129 aggregate by orthogonal direct sum,
finite graph-type counting, or a new generating-function estimate.

\begin{openproblem}[V130 mandate: audit and global collision aggregation]
\label{op:V130-mandate}
\begin{enumerate}
\item Perform the compulsory read-only audit of the latest SM and NS
      files in \texttt{latex\_notes}; record sizes, hashes, and only
      structurally relevant lessons.
\item Reconstruct the exact index set of V54 puncture monomials and V56
      collision patterns before any triangle inequality.
\item Identify which sums are orthogonal physical direct sums and which
      are genuine scalar combinatorial sums.
\item Insert
      \Cref{cor:V129-two-orientation-closed} at the complete-monomial
      level; do not reuse a pair or mean debit across two monomials.
\item Prove a uniform aggregation bound at the quartic level, or isolate
      the precise multiplicity/generating-function factor still missing.
\item Only then reassess complete \([3,2]\), \([3,3]\), \([4]\), and
      global quartic MTA.
\end{enumerate}
\end{openproblem}

\section{V129 ledger}
\label{sec:V129-ledger}

\begin{theorem}[Integrated V129 statement]
\label{thm:V129-integrated}
Preserving V1--V128, V129 proves:
\begin{enumerate}[label=\textup{(V129.\arabic*)},leftmargin=7.5em]
\item the exact one-shared-row product-capacity identity
      \eqref{eq:V129-one-shared-row-product};
\item joint positive completion of all correlated outer labels,
      \eqref{eq:V129-joint-completion};
\item the numerator-preserving capacity bound
      \eqref{eq:V129-completed-capacity};
\item the finite-scale marked payment
      \eqref{eq:V129-shared-row-marked-payment};
\item return of the two-common-endpoint case to the same-label bank;
\item emptiness of the fixed-orientation residual
      \eqref{eq:V129-fixed-orientation-empty}; and
\item closure of both noncommuting orientations for one complete V60
      monomial.
\end{enumerate}
\end{theorem}

\begin{proof}
The items are
\Cref{thm:V129-one-shared-row-product,
lem:V129-joint-completion,
cor:V129-completed-capacity,
thm:V129-shared-row-marked-payment,
lem:V129-two-row-sharing-regroups,
thm:V129-V128-residual-paid,
cor:V129-two-orientation-closed}.
\end{proof}

\begin{firewall}[Claim boundary after V129]
\label{fire:V129-current-boundary}
V129 closes the fixed-monomial physical JREC residual under the standing
programme hypotheses.  It does not yet prove uniform aggregation over
all puncture and collision monomials, prove global JREC, complete
quartic or all-order MTA, construct continuum four-dimensional
Yang--Mills theory satisfying the required axioms, or prove a positive
mass gap.  The full Yang--Mills existence and mass-gap problem remains
open.
\end{firewall}

\section{Append-only record for V129}
\label{sec:V129-record}

V129 was appended to the complete V128 manuscript.  The exact V128
parent had \(85241\) newline-delimited source records, \(2958475\)
bytes, and SHA--256 digest
\[
 \texttt{1cbd1484eef0bc4f4d3c2a9475ac880f2b308dd671e2c713be0a452a0b5fe286}.
\]
No V128 mathematical content was changed.  The sole final
\verb|\end{document}| is restored below.

% =============================================================================
% END APPENDED VERSION V129 MATERIAL
% =============================================================================

% =============================================================================
% BEGIN APPENDED VERSION V130 MATERIAL
% =============================================================================

\chapter{V130: Companion audit and an exact finite-alphabet selector for
quartic collision words}
\label{ch:V130}

\section{Purpose and audit-induced change of order}
\label{sec:V130-purpose}

V129 closed both retained-row orientations of the physical JREC residual
for one completely regrouped V60 monomial.  The unresolved step is not
another local capacity estimate.  It is the passage from the exact V54
coordinate and target-map sum to a global quartic operator without taking
absolute values over coordinate choices.  Because this version is divisible
by five, the prescribed SM/NS audit was performed before that passage.

\begin{assessment}[Mandatory V130 companion audit]
\label{ass:V130-audit}
The read-only audit found the following newest companion files:
\begin{center}
\begin{tabular}{llll}
\toprule
programme & version & source records/bytes & SHA--256\\
\midrule
SM--Einstein & V147 & \(94320/3442657\) &
\texttt{0dc115fcc0416b1a272d07b861a319070ae7b4c1d6827a95b4d7d56cc5e4fb6b}\\
Navier--Stokes & V31 & \(23052/887537\) &
\texttt{f1121b62238ad5491bb7cf03635ea9934942edf573ed8faaa9ee79e1d38a6696}\\
\bottomrule
\end{tabular}
\end{center}
The NS file is unchanged from the V125 audit.  The SM file advanced from
V141 to V147.  Its relevant structural lesson is an order-of-operations
lesson: construct exact finite-regulator comparison, composition, or
pushforward maps before asking whether analytic estimates survive them.
The same audit also sharply separates an algebraic directed limit from
analytic convergence, a form domain from its associated operator domain,
and a complete determinant line from phase transport on that line.
\end{assessment}

\begin{firewall}[Companion transfer firewall]
\label{fire:V130-companion-firewall}
No SM or NS analytic inequality is imported into Yang--Mills.  In
particular, the audit supplies no Wilson-shell estimate, representation
capacity bound, JREC estimate, reflection-positive continuum state, or
spectral gap.  It changes only the proof order: the exact global quartic
selector is formed before any norm or monomial estimate is applied.
\end{firewall}

\section{The quartic local alphabet}
\label{sec:V130-alphabet}

Let \(I:=\{o,a,b,c\}\) and let \(\Pi_4\) be the set of all set partitions
of \(I\).  Thus
\begin{equation}
 |\Pi_4|=B_4=15.
 \label{eq:V130-Bell-four}
\end{equation}
Fix a finite coordinate set \(E\).  At coordinate \(e\), absent rows are
adjoined as singleton entries and \(\Pi_e\subseteq\Pi_4\) denotes the
admissible local partitions.  For \(\pi\in\Pi_e\), let
\begin{equation}
 L_{e,\pi}:=
 \bigotimes_{B\in\pi}\boldsymbol\kappa_{B,e}
 \label{eq:V130-local-letter}
\end{equation}
be the complete local cumulant letter, with a singleton block interpreted
as its first moment and an absent singleton as the identity.  All tensor
placements are those of the exact V39 Wilson replica formula.

A quartic coordinate word is a map
\begin{equation}
 w:E\longrightarrow\Pi_4,
 \qquad w(e)\in\Pi_e.
 \label{eq:V130-word}
\end{equation}
Its operator and its used-letter set are
\begin{equation}
 L(w):=\bigotimes_{e\in E}L_{e,w(e)},
 \qquad
 \operatorname{typ}(w):=\{w(e):e\in E\}\subseteq\Pi_4.
 \label{eq:V130-word-operator}
\end{equation}
Zero local letters may be retained: they contribute zero to every identity
below and require no exceptional convention.

For \(S\subseteq\Pi_4\), define the unrestricted type evaluation
\begin{equation}
 Z(S):=
 \bigotimes_{e\in E}
 \left(\sum_{\pi\in S\cap\Pi_e}L_{e,\pi}\right),
 \label{eq:V130-type-evaluation}
\end{equation}
where an empty local sum is zero.  This object is formed before recoupling,
trace norm, or the final resolvent.

\begin{lemma}[Type evaluation is the exact word sublevel sum]
\label{lem:V130-sublevel-sum}
For every \(S\subseteq\Pi_4\),
\begin{equation}
 Z(S)=
 \sum_{\substack{w(e)\in\Pi_e\ (e\in E)\\
                  \operatorname{typ}(w)\subseteq S}}
 L(w).
 \label{eq:V130-sublevel-sum}
\end{equation}
\end{lemma}

\begin{proof}
Distribute the finite tensor product in
\eqref{eq:V130-type-evaluation}.  A choice of one summand at every
coordinate is exactly a word \(w\), and every selected letter lies in
\(S\).  Conversely, each word with used-letter set contained in \(S\)
selects one and only one summand in the expansion.  No commutation between
letters at the same coordinate is used.
\end{proof}

\section{Exact-support inversion and topology projection}
\label{sec:V130-Boolean}

For \(U\subseteq\Pi_4\), put
\begin{equation}
 Z_{=U}:=
 \sum_{\substack{w(e)\in\Pi_e\ (e\in E)\\
                  \operatorname{typ}(w)=U}}L(w).
 \label{eq:V130-exact-support-def}
\end{equation}

\begin{theorem}[Boolean inversion for the complete coordinate sum]
\label{thm:V130-Boolean-inversion}
For every \(U\subseteq\Pi_4\),
\begin{equation}
 \boxed{
 Z_{=U}=
 \sum_{S\subseteq U}(-1)^{|U|-|S|}Z(S).}
 \label{eq:V130-Boolean-inversion}
\end{equation}
Consequently a fixed used-letter sector is a sum of at most \(2^{15}\)
complete tensor-product evaluations, independently of \(|E|\).
\end{theorem}

\begin{proof}
Insert \eqref{eq:V130-sublevel-sum} on the right.  A fixed word \(w\)
with \(V:=\operatorname{typ}(w)\) has coefficient zero unless
\(V\subseteq U\).  If \(V\subseteq U\), its coefficient is
\[
 \sum_{V\subseteq S\subseteq U}(-1)^{|U|-|S|}
 =(1-1)^{|U\setminus V|}.
\]
This equals one when \(V=U\) and zero otherwise, which is precisely
\eqref{eq:V130-exact-support-def}.  The cardinality assertion follows
from \eqref{eq:V130-Bell-four}.
\end{proof}

A topology rule is any collection
\(\mathcal C\subseteq 2^{\Pi_4}\).  It selects a word when its used-letter
set lies in \(\mathcal C\).  Define
\begin{equation}
 Z_{\mathcal C}:=
 \sum_{\operatorname{typ}(w)\in\mathcal C}L(w).
 \label{eq:V130-topology-sector}
\end{equation}

\begin{corollary}[Finite-alphabet topology selector]
\label{cor:V130-topology-selector}
Every such sector has the exact representation
\begin{align}
 Z_{\mathcal C}
 &=\sum_{S\subseteq\Pi_4}c_{\mathcal C}(S)Z(S),
 \label{eq:V130-topology-selector}\\
 c_{\mathcal C}(S)
 &:=\sum_{\substack{U\in\mathcal C\\S\subseteq U}}
       (-1)^{|U|-|S|}.
 \label{eq:V130-selector-coefficients}
\end{align}
Moreover,
\begin{equation}
 \sum_{S\subseteq\Pi_4}|c_{\mathcal C}(S)|
 \le 3^{15}.
 \label{eq:V130-selector-lone}
\end{equation}
The constant is universal at replica order four and has no support,
volume, representation, or coordinate multiplicity dependence.
\end{corollary}

\begin{proof}
Sum \eqref{eq:V130-Boolean-inversion} over \(U\in\mathcal C\) and
interchange two finite sums to obtain
\eqref{eq:V130-topology-selector}--
\eqref{eq:V130-selector-coefficients}.  The triangle inequality at the
level of scalar coefficients gives
\[
 \sum_S|c_{\mathcal C}(S)|
 \le\sum_{U\in\mathcal C}\sum_{S\subseteq U}1
 \le\sum_{U\subseteq\Pi_4}2^{|U|}
 =(1+2)^{15}.
\]
This proves \eqref{eq:V130-selector-lone}.  No operator norm has been
taken.
\end{proof}

\section{Identification of the V54 and V56 indices}
\label{sec:V130-V54-index}

For a word \(w\), let \(\mathcal H(w)\) be the row hypergraph whose
hyperedges are the nonsingleton blocks of all letters \(w(e)\).  The
priority convention of V40 and V54 is determined entirely by which local
partition types occur: a four-row block, two triple blocks with union
\(I\), a triple plus an attaching pair, and a pair tree can all be detected
from \(\operatorname{typ}(w)\).  Coordinate compatibility is already
encoded in each letter, so repeated occurrences carry no further
topological information.

Fix the omitted row \(o\) and \(T=I\setminus\{o\}\).  Let
\(\mathcal C_{32,o}\) consist of the used-letter sets \(U\) for which:
\begin{enumerate}[label=\textup{(\roman*)},leftmargin=3.7em]
\item no block of size at least three containing \(o\) occurs;
\item the restriction to \(T\) is connected and contains a triple block;
\item at least one pair block \(\{o,t\}\), \(t\in T\), occurs; and
\item the earlier V40 priority classes \([4]\) and \([3,3]\) do not occur.
\end{enumerate}

\begin{proposition}[The Boolean selector is exactly the V54 puncture sum]
\label{prop:V130-V54-equivalence}
Before final recoupling,
\begin{equation}
 \boxed{Z_{\mathcal C_{32,o}}=\mathbf K_{T;o}^{[3,2]}.}
 \label{eq:V130-V54-equivalence}
\end{equation}
In particular, the nonempty coordinate set \(A\), target map
\(\boldsymbol a:A\to T\), punctured ternary word, and complementary
moment choices in \eqref{eq:V54-exact-puncture} form a disjoint
reparametrization of the words selected by \(\mathcal C_{32,o}\).
\end{proposition}

\begin{proof}
Given a selected word, let \(A\) be the coordinates at which the block
containing the row-\(o\) entry is a pair, and let \(a_e\) be its unique
other row.  Conditions \textup{(i)} and \textup{(iii)} make this a
nonempty set and a unique allowed target map.  Delete those selected
pair incidences.  The remaining blocks over \(T\) are precisely one word
in the connected, triple-containing punctured operator
\(\mathbf K_{T;A}^{[3]}\); the unused entries at an attachment coordinate
form one local partition in the complete complement moment
\(\mathbf Q_{T\setminus\{a_e\},e}\).  Row-\(o\) singletons outside \(A\)
give the first product in \eqref{eq:V54-exact-puncture}.

Conversely, expand only the complete complement moments and the punctured
ternary operator on the right side of
\eqref{eq:V54-exact-puncture}.  Each choice produces one local partition
at each coordinate and hence one word.  The V54 connectivity and priority
conditions put its used-letter set in \(\mathcal C_{32,o}\).  The two
constructions recover one another, including the coordinate and target
labels.  Therefore their operator sums agree.
\end{proof}

At a coordinate where the selected block is \(\{o,a\}\), the remaining
rows \(b,c\) admit exactly two partitions.  In the present alphabet these
are
\begin{align}
 \pi^{\rm m}_{oa}&:=oa\mid b\mid c,
 &L_{e,\pi^{\rm m}_{oa}}
   &=\mathbf d_{oa,e}\,q_{b,e}q_{c,e},
 \label{eq:V130-mean-letter}\\
 \pi^{\rm dp}_{oa\mid bc}&:=oa\mid bc,
 &L_{e,\pi^{\rm dp}_{oa\mid bc}}
   &=\mathbf d_{oa,e}\otimes\mathbf d_{bc,e}.
 \label{eq:V130-double-pair-letter}
\end{align}

\begin{corollary}[Exact collision-pattern index]
\label{cor:V130-collision-index}
The V56 collision split is exactly the two-letter alternative
\eqref{eq:V130-mean-letter}--\eqref{eq:V130-double-pair-letter} at every
attachment coordinate.  Thus a complete V54 puncture monomial is uniquely
indexed by its target pair label at every coordinate, its mean/double-pair
collision flag, and one punctured ternary word.  On noncollision pair
banks, V56 has already summed all coordinate choices into its three bank
defects and seven attachment monomials.  The unsummed part addressed by
V129 is therefore the genuine collision-word sum, not the seven-term
noncollision polynomial.
\end{corollary}

\begin{proof}
The complement of \(\{o,a\}\) has two elements.  Its partition lattice
has only \(b\mid c\) and \(bc\), giving the displayed letters and the exact
identity \eqref{eq:V56-collision-factor}.  Coordinatewise uniqueness and
the tensor-product construction prove the first two statements.  The last
statement is \Cref{thm:V56-exact-polynomial}; its subset index is the
nonempty subset of the three fixed pair labels, not a subset of an
arbitrarily large coordinate set.
\end{proof}

\section{Which sums are orthogonal}
\label{sec:V130-orthogonality}

\begin{proposition}[Output blocks are orthogonal; partition words are not]
\label{prop:V130-orthogonality-separation}
Let \(\{P_\lambda\}\) be the complete mutually orthogonal outer
Peter--Weyl/CG projectors used in the V94--V129 priority construction.
After one assembled operator has been fixed, the decomposition
\(X=\bigoplus_\lambda P_\lambda XP_\lambda\) satisfies
\begin{equation}
 \|X\|_1=\sum_\lambda\|P_\lambda XP_\lambda\|_1.
 \label{eq:V130-orthogonal-trace-sum}
\end{equation}
Therefore deterministic refinements of complete physical output blocks
may be summed without a channel factor.

By contrast, two different local partitions at the same coordinate act
on the same input/output carrier in general.  Their coordinate-word sums
are algebraic sums, not orthogonal direct sums.  Neither
\eqref{eq:V130-orthogonal-trace-sum} nor the V127 safe-refinement lemma may
be applied to them before the complete word sum has been assembled.
\end{proposition}

\begin{proof}
The first assertion is the trace norm identity for a block-diagonal
trace-class operator on mutually orthogonal subspaces.  It is precisely
the setting of the complete outer projectors.

For the second assertion it suffices to take a one-dimensional carrier.
Two distinct formal partition letters may both act as multiplication by
the scalar \(1\).  Their ranges then coincide, whereas nonzero orthogonal
ranges in one dimension cannot coexist.  Replacing the second scalar by
(-1) also shows why cancellation can be destroyed: the assembled norm is
zero while the sum of the two termwise norms is two.  Local cumulant
letters are signed operator contributions, so no positivity principle
repairs this failure.  Hence word labels alone never create physical
orthogonality.
\end{proof}

\section{What the finite alphabet proves, and the remaining analytic gate}
\label{sec:V130-gate}

\begin{theorem}[Conditional transfer through the finite selector]
\label{thm:V130-conditional-transfer}
Let \(\mathfrak N\) be any subadditive seminorm on the completely
recoupled quartic operator space.  If, for a fixed nonnegative marked-tree
majorant \(M\), every type evaluation obeys
\begin{equation}
 \mathfrak N(Z(S))\le M,
 \qquad S\subseteq\Pi_4,
 \label{eq:V130-evaluation-hypothesis}
\end{equation}
then every topology sector obeys
\begin{equation}
 \mathfrak N(Z_{\mathcal C})\le3^{15}M.
 \label{eq:V130-conditional-transfer}
\end{equation}
\end{theorem}

\begin{proof}
Apply subadditivity to
\eqref{eq:V130-topology-selector}, then use
\eqref{eq:V130-evaluation-hypothesis} and
\eqref{eq:V130-selector-lone}.
\end{proof}

\begin{firewall}[Why V129 does not supply the hypothesis]
\label{fire:V130-V129-not-hypothesis}
The V129 theorem bounds one completely regrouped partition monomial after
its heavy banks and physical outer blocks have been identified.  It does
not prove \eqref{eq:V130-evaluation-hypothesis}: \(Z(S)\) mixes all
coordinate assignments of the allowed letters, and an unrestricted
evaluation contains disconnected zeroth-order pieces.  Those pieces can
be order one even when the required connected marked-tree majorant
vanishes.  Applying the triangle inequality separately to the \(Z(S)\)
would therefore destroy the very Boolean/Moebius cancellation which
removes the disconnected contributions.  The factor \(3^{15}\) is not the
remaining problem; cancellation-compatible control of the complete
finite difference is.
\end{firewall}

Define the exact V130 collision projection by
\begin{equation}
 \mathsf P_{32,o}Z
 :=\sum_{S\subseteq\Pi_4}
 c_{\mathcal C_{32,o}}(S)Z(S)
 =\mathbf K_{T;o}^{[3,2]}.
 \label{eq:V130-collision-projection}
\end{equation}
After the already closed noncollision part is removed, let
\(\mathsf P_{32,o}^{\rm coll}Z\) denote the remaining words containing at
least one double-pair letter \(\pi^{\rm dp}_{oa\mid bc}\).

\begin{openproblem}[The exact global collision-JREC estimate]
\label{op:V130-global-collision-JREC}
Prove directly for the complete signed finite difference, not for its
summands, that
\begin{equation}
 \boxed{
 \mathfrak J_o\!\left(
   \mathsf P_{32,o}^{\rm coll}Z
 \right)
 \le
 C s_\alpha
 \sum_{r\in\{b,c\}}
 \frac{h_{oa}H_{ar}h_{bc}}{\Xi_a\Xi_r}
 \prod_{i\in I}\|B_i\|_1.}
 \label{eq:V130-global-collision-JREC}
\end{equation}
Here every same-label pair bank is first completely regrouped, the
V94--V129 physical priority partition is applied only afterward, and the
final resolvent is charged once.  The constant must be independent of
\(E\), coordinate multiplicities, representation heights, and final
channels.  V129 proves the required estimate on every fixed word fibre;
V130 proves that all fibres are selected by one exact 15-letter finite
difference.  What remains is an assembled norm estimate that preserves
that finite difference.
\end{openproblem}

\begin{assessment}[Regression tests]
\label{ass:V130-regressions}
The construction passes four necessary tests.  Repeating one collision
type at arbitrarily many coordinates does not enlarge the alphabet.
Mean and double-pair alternatives at the same coordinate remain inside
one signed operator sum and are not falsely declared orthogonal.  Cold
singlet outputs are handled only after complete raw output pinching, so no
output Casimir substitutes for an input mass.  Finally, the selector
depends only on the fixed replica order four; its scalar coefficient bound
does not depend on volume or support size.
\end{assessment}

\section{V130 ledger and next upstream acquisition}
\label{sec:V130-ledger}

\begin{theorem}[Integrated V130 statement]
\label{thm:V130-integrated}
Preserving V1--V129, V130 proves:
\begin{enumerate}[label=\textup{(V130.\arabic*)},leftmargin=7.5em]
\item the mandatory SM/NS audit and its strict transfer firewall;
\item the exact 15-letter quartic word algebra
      \eqref{eq:V130-type-evaluation};
\item Boolean inversion by used-letter set
      \eqref{eq:V130-Boolean-inversion};
\item the finite topology selector
      \eqref{eq:V130-topology-selector} with universal coefficient bound
      \eqref{eq:V130-selector-lone};
\item exact equivalence of its \([3,2]\) selector with the V54 puncture
      formula;
\item the precise two-letter collision index underlying V56; and
\item the separation of safe orthogonal output sums from unsafe algebraic
      partition-word sums.
\end{enumerate}
\end{theorem}

\begin{proof}
The assertions are respectively
\Cref{ass:V130-audit,fire:V130-companion-firewall,
lem:V130-sublevel-sum,thm:V130-Boolean-inversion,
cor:V130-topology-selector,prop:V130-V54-equivalence,
cor:V130-collision-index,prop:V130-orthogonality-separation}.
\end{proof}

\begin{openproblem}[V131 mandate: preserve cancellation analytically]
\label{op:V131-mandate}
\begin{enumerate}
\item Rewrite \(\mathsf P_{32,o}^{\rm coll}\) as an iterated finite
      difference of complete local moment products, keeping the row
      Moebius cancellation visible.
\item Seek an interpolation or connected finite-difference formula whose
      derivatives insert actual pair/triple defects.  A useful formula
      must produce three marked links at order four without expanding
      coordinate subsets.
\item Insert the V129 fixed-word capacity only after those defects have
      been selected by the exact derivative/finite-difference operator.
\item If the resulting norm is not stable under the finite difference,
      exhibit a finite-dimensional counterexample and go upstream to the
      choice of connected interpolation rather than adding another local
      capacity estimate.
\item Do not revisit the companion files until V135 unless an earlier
      contradiction specifically requires them.
\end{enumerate}
\end{openproblem}

\begin{firewall}[Claim boundary after V130]
\label{fire:V130-current-boundary}
V130 removes the coordinate-subset multiplicity at the level of exact
quartic operator algebra and identifies the remaining analytic statement
\eqref{eq:V130-global-collision-JREC}.  It does not prove that estimate,
global JREC, complete \([3,2]\), \([3,3]\), \([4]\), quartic or all-order
MTA, a four-dimensional continuum Yang--Mills theory satisfying the
required axioms, or a positive spectral gap.  The full Yang--Mills
existence and mass-gap problem remains open.
\end{firewall}

\section{Append-only record for V130}
\label{sec:V130-record}

V130 was appended to the complete V129 manuscript.  The exact V129 parent
had \(85771\) newline-delimited source records, \(2976511\) bytes, and
SHA--256 digest
\[
 \texttt{0834853cbc5008f43697aec15a567c61cc1a7fa5b4f98d4da6489ff4118e3ef0}.
\]
No V129 mathematical content was changed.  The sole final
\verb|\end{document}| is restored below.

% =============================================================================
% END APPENDED VERSION V130 MATERIAL
% =============================================================================

% =============================================================================
% BEGIN APPENDED VERSION V131 MATERIAL
% =============================================================================

\chapter{V131: A connected forest interpolation for the complete
quartic collision sector}
\label{ch:V131}

\section{Purpose}
\label{sec:V131-purpose}

V130 replaced the coordinate-subset index by an exact finite alphabet,
but its Boolean representation still wrote the desired topology as a
signed sum of unrestricted products.  Taking norms of those products would
destroy connected cancellation.  This continuation goes upstream to the
row-connectivity projection.  It builds a formal edge interpolation of the
complete local cumulant alphabet, proves the associated forest identity,
and obtains an exact 16-tree formula for the whole collision sector.  The
three differentiations insert actual local cumulant links before any norm
is taken.

\begin{firewall}[Scope of V131]
\label{fire:V131-scope}
The interpolation below is an algebraic polynomial device; it is not an
interpolation of probability measures and requires no positivity at
intermediate parameters.  V131 proves the exact connected formula and its
marking properties.  It does not yet prove that the V129 capacity estimate
is stable under the assembled derivative, and therefore does not claim
global collision JREC.
\end{firewall}

\section{Canonical edge encoding of a local partition}
\label{sec:V131-edge-encoding}

Fix once and for all a total order on the four rows \(I=\{o,a,b,c\}\).
Let
\[
 \mathcal E_4:=\bigl\{\{i,j\}:i,j\in I,\ i\ne j\bigr\}
\]
be the six edges of the complete row graph.  For a nonempty block
\(B\subseteq I\), let \(r(B)\) be its least row and set
\begin{equation}
 \epsilon(B):=\bigl\{\{r(B),j\}:j\in B\setminus\{r(B)\}\bigr\}.
 \label{eq:V131-block-star}
\end{equation}
For \(\pi\in\Pi_4\), define
\begin{equation}
 \epsilon(\pi):=
 \mathop{\dot\bigcup}_{B\in\pi}\epsilon(B).
 \label{eq:V131-partition-forest}
\end{equation}
This is a forest whose connected components are exactly the blocks of
\(\pi\).

For \(x=(x_{ij})_{ij\in\mathcal E_4}\in[0,1]^{\mathcal E_4}\), put
\begin{equation}
 x^{\epsilon(\pi)}:=\prod_{ij\in\epsilon(\pi)}x_{ij}.
 \label{eq:V131-edge-monomial}
\end{equation}
Using the V130 local letters, define
\begin{align}
 M_e(x)&:=\sum_{\pi\in\Pi_e}
 L_{e,\pi}x^{\epsilon(\pi)},
 \label{eq:V131-local-polynomial}\\
 \Phi(x)&:=\bigotimes_{e\in E}M_e(x).
 \label{eq:V131-global-polynomial}
\end{align}
All products are finite, so \(\Phi\) is a polynomial with values in the
finite-dimensional operator space of the fixed regulated problem.

For a row partition \(\rho\in\Pi_4\), let
\(\mathbf1_\rho\in\{0,1\}^{\mathcal E_4}\) be defined by
\begin{equation}
 (\mathbf1_\rho)_{ij}:=
 \begin{cases}
 1,&i,j\text{ lie in the same block of }\rho,\\
 0,&\text{otherwise}.
 \end{cases}
 \label{eq:V131-partition-point}
\end{equation}

\begin{lemma}[Row Moebius projection selects connected words]
\label{lem:V131-row-Mobius}
Let \(\mu\) be the Moebius function of the partition lattice.  Then
\begin{equation}
 \boxed{
 \Phi^{\rm conn}:=
 \sum_{\rho\in\Pi_4}\mu(\rho,\boldsymbol1_I)
 \Phi(\mathbf1_\rho)
 =\sum_{\mathcal H(w)\ {\rm connected}}L(w).}
 \label{eq:V131-row-Mobius}
\end{equation}
\end{lemma}

\begin{proof}
Expand \(\Phi\) into coordinate words.  For a fixed word \(w\), the
monomial in \eqref{eq:V131-global-polynomial} is nonzero at
\(\mathbf1_\rho\) exactly when every canonical edge of every local letter
has both endpoints in a block of \(\rho\).  Since the components of
\(\epsilon(\pi)\) are the blocks of \(\pi), this is equivalent to
\(\rho\ge\bar w\), where \(\bar w\) is the row partition into connected
components of \(\mathcal H(w)\).  The coefficient of \(L(w)\) is therefore
\[
 \sum_{\rho\ge\bar w}\mu(\rho,\boldsymbol1_I),
\]
which is one if \(\bar w=\boldsymbol1_I\) and zero otherwise.  This is the
defining cancellation identity for the partition-lattice Moebius
function.
\end{proof}

\section{The finite forest formula}
\label{sec:V131-forest-formula}

For a graph forest \(F\subseteq\mathcal E_4\) and
\(t=(t_f)_{f\in F}\in[0,1]^F\), define the weakening point
\(x^F(t)\in[0,1]^{\mathcal E_4}\) by
\begin{equation}
 x^F_{ij}(t):=
 \begin{cases}
 \displaystyle\min_{f\in\operatorname{path}_F(i,j)}t_f,
       &i,j\text{ are connected in }F,\\
 0,&\text{otherwise}.
 \end{cases}
 \label{eq:V131-forest-weakening}
\end{equation}
For the empty path set \(x^F_{ii}=1\), although diagonal variables never
occur.  Write
\(\partial_F:=\prod_{ij\in F}\partial_{x_{ij}}\).

\begin{theorem}[Finite forest interpolation]
\label{thm:V131-forest-formula}
For every polynomial \(F_0\) on
\([0,1]^{\mathcal E_4}\) with values in a finite-dimensional vector space,
\begin{equation}
 \boxed{
 F_0(\mathbf1)=
 \sum_{F\ {\rm forest\ on}\ I}
 \int_{[0,1]^F}
 (\partial_FF_0)(x^F(t))\,\dd t.}
 \label{eq:V131-forest-formula}
\end{equation}
The empty-forest term is \(F_0(\mathbf0)\).
\end{theorem}

\begin{proof}
It suffices to prove the identity componentwise.  Repeatedly apply the
one-variable fundamental theorem of calculus, beginning at the point where
all edges between distinct current components are zero.  Whenever one
differentiation is made, record the edge whose variable is differentiated
and merge the two components containing its endpoints.  An edge whose
endpoints are already in the same recorded component is not used for a new
merge.  Hence the recorded edges form a forest.

Order the recorded parameters from largest to smallest.  At a stage with
threshold \(u\), precisely the recorded edges with parameter at least \(u\)
have been switched on.  Two rows are then in the same current component
exactly when the unique path between them in the recorded forest has every
parameter at least \(u\).  Consequently the final value assigned to their
edge variable is the least parameter on that path, namely
\eqref{eq:V131-forest-weakening}.  Conversely, a forest together with an
ordering of its parameters uniquely reconstructs this sequence of merges.

The cubes \([0,1]^F\) are, up to measure-zero equal-parameter faces,
partitioned by those orderings.  The fundamental-theorem boundary at which
no merge is made is \(F_0(\mathbf0)\); every intermediate boundary occurs
once with each sign and cancels; the sole terminal boundary is
\(F_0(\mathbf1)\).  Summing the iterated integral over all merge sequences,
then recombining the ordered simplices into the cubes, gives
\eqref{eq:V131-forest-formula}.  Polynomial regularity justifies every
finite differentiation, integration, and rearrangement.
\end{proof}

\begin{corollary}[Connected projection is a sum over spanning trees]
\label{cor:V131-tree-formula}
For every such \(F_0\),
\begin{equation}
 \boxed{
 \sum_{\rho\in\Pi_4}\mu(\rho,\boldsymbol1_I)
 F_0(\mathbf1_\rho)
 =
 \sum_{\tau\in\mathfrak T_4}
 \int_{[0,1]^3}
 (\partial_\tau F_0)(x^\tau(t))\,\dd t.}
 \label{eq:V131-tree-formula}
\end{equation}
There are \(|\mathfrak T_4|=4^{4-2}=16\) summands.
\end{corollary}

\begin{proof}
Apply \Cref{thm:V131-forest-formula} to the variables internal to every
block of a fixed \(\rho\).  A fixed forest \(F\) then occurs exactly for
the row partitions \(\rho\) coarser than its component partition
\(\operatorname{comp}(F)\).  Its total coefficient in the left side is
\[
 \sum_{\rho\ge\operatorname{comp}(F)}
 \mu(\rho,\boldsymbol1_I),
\]
which vanishes unless \(F\) is connected.  A connected forest on four
vertices is a spanning tree and has three edges.  If it is connected, the
coefficient is one.  This proves the formula.  The count is Cayley's formula
for four labeled vertices; it can also be checked directly by choosing the
two middle entries of a length-two Prüfer word.
\end{proof}

\section{Exact topology flags for the collision sector}
\label{sec:V131-flags}

Fix \(o\) and \(T=I\setminus\{o\}\).  Let \(\mathcal A_o\subseteq\Pi_4\)
be the partitions having no block of cardinality at least three which
contains \(o\).  This alphabet enforces the V54 exclusion of the earlier
\([4]\) and \([3,3]\) priority sectors.

For \(\pi\in\mathcal A_o\), define two zero-one flags:
\begin{align}
 \theta_o(\pi)&:=
 \boldsymbol1_{\{T\text{ is a block of }\pi\}},
 \label{eq:V131-triple-flag}\\
 \gamma_o(\pi)&:=
 \boldsymbol1_{\{\pi\text{ consists of two nonsingleton pair blocks}\}}.
 \label{eq:V131-collision-flag}
\end{align}
In \eqref{eq:V131-collision-flag}, one of the two pair blocks necessarily
contains \(o\); thus it is exactly a V56 double-pair collision letter.

For auxiliary variables \(y,z\), set
\begin{align}
 M_{e,o}(x;y,z)
 &:={}
 \sum_{\pi\in\Pi_e\cap\mathcal A_o}
 L_{e,\pi}
 x^{\epsilon(\pi)}
 y^{\theta_o(\pi)}z^{\gamma_o(\pi)},
 \label{eq:V131-flagged-local}\\
 \Phi_o(x;y,z)
 &:={}
 \bigotimes_{e\in E}M_{e,o}(x;y,z).
 \label{eq:V131-flagged-global}
\end{align}
Let \(\Delta_y f:=f(1,z)-f(0,z)\) and
\(\Delta_z f:=f(y,1)-f(y,0)\).

\begin{lemma}[Presence differences select triple and collision letters]
\label{lem:V131-presence-differences}
The operator
\begin{equation}
 \Delta_y\Delta_z\Phi_o(\mathbf1;y,z)
 \label{eq:V131-presence-sector}
\end{equation}
is the exact sum of all words in the alphabet \(\mathcal A_o\) which use
at least one \(T\)-triple letter and at least one double-pair collision
letter.
\end{lemma}

\begin{proof}
After expansion into words, the powers of \(y\) and \(z\) count the
occurrences of the two flagged letter classes.  For every nonnegative
integer \(k\), \(1^k-0^k\) is one when \(k\ge1\) and zero when \(k=0\),
with the polynomial convention \(0^0=1\).  The two differences therefore
retain exactly the stated words, each with coefficient one.
\end{proof}

\begin{proposition}[The flagged sector is the complete V54 collision
sector]
\label{prop:V131-flagged-is-collision}
With the V40/V54 priority convention,
\begin{equation}
 \boxed{
 \mathbf K_{T;o}^{[3,2],\rm coll}
 =\Delta_y\Delta_z\Phi_o(\mathbf1;y,z).}
 \label{eq:V131-flagged-is-collision}
\end{equation}
Every word on the right is row-connected.
\end{proposition}

\begin{proof}
A retained word contains the local triple block \(T\), which connects all
three non-\(o\) rows.  Its double-pair letter contains one pair incident to
\(o\), so the full row hypergraph is connected and has the V54
\([3,2]\) skeleton.  The alphabet \(\mathcal A_o\) excludes any triple or
four-block containing \(o\), hence excludes the earlier \([4]\) and
\([3,3]\) priority types.  The double-pair flag makes the word a collision
word rather than a purely mean-complement attachment.

Conversely, every V54 collision word has a triple \(T\), contains at least
one V56 double-pair letter, and has no cardinality-three or -four block
containing \(o\).  It is therefore retained by the two presence
differences.  This is a wordwise bijection, and summing proves the identity.
\end{proof}

\section{The exact 16-tree collision formula}
\label{sec:V131-collision-tree}

\begin{theorem}[Connected forest representation of the entire collision
sector]
\label{thm:V131-collision-tree}
The full, unexpanded V54 collision operator satisfies
\begin{equation}
 \boxed{
 \mathbf K_{T;o}^{[3,2],\rm coll}
 =
 \sum_{\tau\in\mathfrak T_4}
 \int_{[0,1]^3}
 \Delta_y\Delta_z
 \bigl(\partial_\tau\Phi_o\bigr)
       (x^\tau(t);y,z)\,\dd t.}
 \label{eq:V131-collision-tree}
\end{equation}
There are exactly sixteen tree integrals and four endpoint evaluations of
the two presence flags.  Neither number depends on \(|E|\).
\end{theorem}

\begin{proof}
By \Cref{prop:V131-flagged-is-collision}, every word selected by
\(\Delta_y\Delta_z\) is connected.  Therefore the row Moebius projection
of \Cref{lem:V131-row-Mobius} acts as the identity on that selected sum.
Apply \Cref{cor:V131-tree-formula} to
\(F_0(x)=\Delta_y\Delta_z\Phi_o(x;y,z)\).  The edge derivatives commute
with the two finite differences because all operations act on distinct
polynomial variables.  This gives \eqref{eq:V131-collision-tree}.  The
counts follow from the sixteen labeled spanning trees and the four terms in
\(\Delta_y\Delta_z\).
\end{proof}

\begin{lemma}[Every tree derivative inserts three genuine links]
\label{lem:V131-genuine-links}
Fix \(\tau\in\mathfrak T_4\).  Every nonzero word contribution to
\(\partial_\tau\Phi_o\) has three distinguished incidences, one for each
edge \(ij\in\tau\).  The incidence carrying \(ij\) is a local cumulant
letter whose block contains \(i,j\).  Several tree edges may be carried by
one triple or four-row letter, but no derivative can mark a singleton
factor or an absent incidence.
\end{lemma}

\begin{proof}
The variable \(x_{ij}\) occurs in a local monomial
\iffalse
\(x^{\epsilon(\pi)}\) only when (ij\in\epsilon(\pi)\).  By
\fi
\(x^{\epsilon(\pi)}\) only when \(ij\in\epsilon(\pi)\).  By
\eqref{eq:V131-block-star}, that edge then belongs to the canonical star of
one nonsingleton block \(B\in\pi\), so \(i,j\in B\).  Applying the Leibniz
rule to the product over coordinates distinguishes an occurrence for each
of the three different variables in \(\partial_\tau\).  If two variables
occur in the same local monomial, both derivatives may act there; this is
exactly the legitimate multi-link marking of a triple or four-row
cumulant.  Singleton and absent factors have empty edge monomial and are
annihilated by every edge derivative.
\end{proof}

\begin{corollary}[Compatibility with fixed-word V129 closure]
\label{cor:V131-V129-compatible}
If the integrand of \eqref{eq:V131-collision-tree} is expanded only after
the three differentiations and two presence differences, each resulting
fixed word is a V60 collision monomial with three valid row links and a
complete same-label regrouping.  Its physical V94--V129 residual is paid in
both retained-row orientations by
\Cref{cor:V129-two-orientation-closed}.
\end{corollary}

\begin{proof}
The presence differences impose the V54/V56 collision alphabet by
\Cref{prop:V131-flagged-is-collision}; the derivative marks are genuine by
\Cref{lem:V131-genuine-links}.  Same-label regrouping is tensor grouping of
all coordinates with the same pair block and commutes with scalar
evaluation of \(x^\tau(t)\).  Thus every fixed word lies in the class
treated by V129.  Apply its two-orientation corollary.
\end{proof}

\begin{firewall}[The derivative must remain assembled]
\label{fire:V131-assembled-derivative}
The last corollary is not permission to take a triangle inequality over
the coordinate choices generated by the Leibniz rule.  Formula
\eqref{eq:V131-collision-tree} is useful precisely because
\(\partial_\tau\Phi_o\) can be kept as a derivative of one complete product.
Global JREC now reduces to a norm estimate on sixteen assembled derivative
integrals, rather than an exponential coordinate-word sum.  Such an
estimate has not yet been proved.
\end{firewall}

\section{The new analytic gate}
\label{sec:V131-gate}

For each tree \(\tau\), define
\begin{equation}
 \mathcal D_{o,\tau}(t):=
 \Delta_y\Delta_z
 \bigl(\partial_\tau\Phi_o\bigr)(x^\tau(t);y,z).
 \label{eq:V131-tree-integrand}
\end{equation}

\begin{openproblem}[Assembled tree-derivative JREC]
\label{op:V131-derivative-JREC}
Prove, with one final resolvent and after complete output pinching,
\begin{equation}
 \boxed{
 \int_{[0,1]^3}
 \mathfrak J_o\bigl(\mathcal D_{o,\tau}(t)\bigr)\,\dd t
 \le
 C s_\alpha
 \sum_{r\in\{b,c\}}
 \frac{h_{oa}H_{ar}h_{bc}}{\Xi_a\Xi_r}
 \prod_{i\in I}\|B_i\|_1}
 \label{eq:V131-derivative-JREC}
\end{equation}
uniformly for all sixteen trees.  It is enough that the constant be uniform
in \(\tau\), since their number is fixed.  A valid proof must estimate the
Leibniz derivative through a product-rule or martingale square function
which sums marked coordinates before trace norm; it may then insert the
V129 joint capacity at the selected collision carriers.
\end{openproblem}

\begin{assessment}[Why this is strictly upstream of V130]
\label{ass:V131-upstream}
The V130 Boolean selector had up to \(3^{15}\) scalar variation and exposed
unconnected order-one evaluations.  Formula
\eqref{eq:V131-collision-tree} cancels every disconnected word exactly and
places three derivatives on actual connected links.  The remaining issue
is therefore no longer a topology count or a missing edge: it is the
operator-norm summation of the marked Leibniz derivative.  Returning to a
new single-word capacity estimate would be circular, because V129 already
closed every such word.
\end{assessment}

\section{V131 ledger and next acquisition}
\label{sec:V131-ledger}

\begin{theorem}[Integrated V131 statement]
\label{thm:V131-integrated}
Preserving V1--V130, V131 proves:
\begin{enumerate}[label=\textup{(V131.\arabic*)},leftmargin=7.5em]
\item a canonical six-edge polynomial encoding of every quartic local
      partition;
\item exact row-connectivity projection by partition-lattice Moebius
      inversion;
\item the finite forest interpolation
      \eqref{eq:V131-forest-formula};
\item its exact sixteen-spanning-tree connected specialization
      \eqref{eq:V131-tree-formula};
\item two exact presence differences selecting the complete V54 collision
      sector; and
\item the unexpanded sixteen-tree collision formula
      \eqref{eq:V131-collision-tree}, whose derivatives mark three genuine
      local cumulant links.
\end{enumerate}
\end{theorem}

\begin{proof}
The items are
\Cref{lem:V131-row-Mobius,thm:V131-forest-formula,
cor:V131-tree-formula,lem:V131-presence-differences,
prop:V131-flagged-is-collision,thm:V131-collision-tree,
lem:V131-genuine-links}.
\end{proof}

\begin{openproblem}[V132 mandate: an assembled Leibniz norm]
\label{op:V132-mandate}
\begin{enumerate}
\item Differentiate the product \(\Phi_o\) at the operator level and group
      the three marked occurrences by coincidence pattern: three distinct
      coordinates, exactly two coincident, or all three coincident.
\item Keep each coincidence class as a complete ordered derivative, not a
      sum of absolute monomials.
\item Test whether the existing one-, two-, and three-block local cumulant
      bounds are stable under every weakening factor
      \(x^\tau_{ij}(t)\in[0,1]\).
\item For the three-distinct-coordinate class, seek a noncommutative
      product-rule estimate whose scalar shadow is the exact derivative of
      a product, and insert V129 only at the final positive completion.
\item If uniformity fails in a coincidence class, isolate a concrete
      finite-dimensional counterexample and revise the interpolation
      encoding upstream.
\end{enumerate}
\end{openproblem}

\begin{firewall}[Claim boundary after V131]
\label{fire:V131-current-boundary}
V131 proves an exact, volume-independent, connected 16-tree representation
of the full quartic \([3,2]\) collision operator.  It does not prove the
assembled derivative estimate \eqref{eq:V131-derivative-JREC}, global JREC,
complete \([3,2]\), \([3,3]\), \([4]\), quartic or all-order MTA, continuum
Yang--Mills existence, or a positive spectral gap.  The full Yang--Mills
existence and mass-gap problem remains open.
\end{firewall}

\section{Append-only record for V131}
\label{sec:V131-record}

V131 was appended to the complete V130 manuscript.  The exact V130 parent
had \(86276\) newline-delimited source records, \(2996979\) bytes, and
SHA--256 digest
\[
 \texttt{dccc813e89a7af6810b6d50e4a0cdd09d654a3636813d2960a331445ad60a3d4}.
\]
No V130 mathematical content was changed.  The sole final
\verb|\end{document}| is restored below.

% =============================================================================
% END APPENDED VERSION V131 MATERIAL
% =============================================================================

% =============================================================================
% BEGIN APPENDED VERSION V132 MATERIAL
% =============================================================================

\chapter{V132: A contraction-preserving partition interpolation and an
assembled Leibniz theorem}
\label{ch:V132}

\section{Purpose and regression of the first interpolation}
\label{sec:V132-purpose}

The V131 forest formula is exact.  Its particular local polynomial
\(M_e(x)=\sum_\pi L_{e,\pi}x^{\epsilon(\pi)}\), however, is not known to
be a contraction at the weakened points \(x^\tau(t)\).  Completeness gives
a moment contraction at the partition points
\(\mathbf1_\rho\), but not at an arbitrary point of the six-dimensional
cube.  A bound larger than one on every unmarked coordinate would be
multiplied once per coordinate and would reintroduce exponential volume
growth.

\begin{proposition}[Endpoint control does not control an interpolation]
\label{prop:V132-endpoint-regression}
For every \(R>0\), the scalar polynomial
\begin{equation}
 p_R(x):=1+Rx(1-x)
 \label{eq:V132-endpoint-regression}
\end{equation}
satisfies \(p_R(0)=p_R(1)=1\), while
\(\sup_{0\le x\le1}|p_R(x)|=1+R/4\).
Consequently contraction at the Boolean partition points, by itself,
cannot justify contraction along a forest interpolation.
\end{proposition}

\begin{proof}
The endpoint identities are immediate, and the concave quadratic
\(x(1-x)\) has its maximum \(1/4\) at \(x=1/2\).
\end{proof}

\begin{firewall}[Meaning of the regression]
\label{fire:V132-regression-scope}
\Cref{prop:V132-endpoint-regression} is a logical counterexample, not a
claim that a Wilson local factor realizes the polynomial \(p_R\).
It invalidates only the proposed inference from endpoint contractions to
interior contractions.  The exact V131 identities remain true.  The
remedy below changes the interpolation, not the connected operator.
\end{firewall}

\section{Complete partial moments}
\label{sec:V132-partial-moments}

For a coordinate \(e\) and a row set \(C\subseteq I\), let
\(\mathbf Q_{e,C}\) be the complete joint Wilson moment on the rows in
\(C\), with \(\mathbf Q_{e,\varnothing}=I\).  For a row partition
\(\rho\in\Pi_4\), define
\begin{equation}
 \mathbf Q_{e,\rho}:=
 \bigotimes_{C\in\rho}\mathbf Q_{e,C}.
 \label{eq:V132-partial-moment}
\end{equation}
Inactive row entries are identities.

\begin{lemma}[Partial moments are contractions]
\label{lem:V132-moment-contractions}
For every \(e\) and \(\rho\in\Pi_4\),
\begin{equation}
 \|\mathbf Q_{e,\rho}\|_{\rm op}\le1.
 \label{eq:V132-moment-contraction}
\end{equation}
\end{lemma}

\begin{proof}
Each \(\mathbf Q_{e,C}\) is the expectation of a tensor product of unitary
representation matrices.  Its operator norm is therefore at most the
expectation of the norm, which is one.  Tensor-product norms multiply, so
\eqref{eq:V132-moment-contraction} follows.
\end{proof}

For a graph \(A\subseteq\mathcal E_4\), let \(\rho(A)\) be its connected
component partition.  For
\(x\in[0,1]^{\mathcal E_4}\), define the Bernoulli weight
\begin{equation}
 \beta_A(x):=
 \prod_{g\in A}x_g
 \prod_{g\in\mathcal E_4\setminus A}(1-x_g)
 \label{eq:V132-Bernoulli-weight}
\end{equation}
and the local partition interpolation
\begin{equation}
 \widetilde M_e(x):=
 \sum_{A\subseteq\mathcal E_4}
 \beta_A(x)\mathbf Q_{e,\rho(A)}.
 \label{eq:V132-local-interpolation}
\end{equation}

\begin{theorem}[Exact endpoints and contraction throughout the cube]
\label{thm:V132-contractive-interpolation}
For all \(x\in[0,1]^{\mathcal E_4}\),
\begin{equation}
 \boxed{\|\widetilde M_e(x)\|_{\rm op}\le1.}
 \label{eq:V132-interior-contraction}
\end{equation}
For every row partition \(\rho\),
\begin{equation}
 \boxed{\widetilde M_e(\mathbf1_\rho)=\mathbf Q_{e,\rho}.}
 \label{eq:V132-exact-partition-endpoint}
\end{equation}
\end{theorem}

\begin{proof}
The coefficients \(\beta_A(x)\) are nonnegative and
\[
 \sum_{A\subseteq\mathcal E_4}\beta_A(x)
 =\prod_{g\in\mathcal E_4}\bigl(x_g+(1-x_g)\bigr)=1.
\]
Thus \(\widetilde M_e(x)\) is a convex combination of the contractions
in \Cref{lem:V132-moment-contractions}, proving
\eqref{eq:V132-interior-contraction}.

At \(x=\mathbf1_\rho\), every edge internal to a block of \(\rho\) has
parameter one and every cross-block edge has parameter zero.  Hence exactly
one Bernoulli weight is nonzero: the graph consisting of all internal
edges.  Its component partition is \(\rho\), proving
\eqref{eq:V132-exact-partition-endpoint}.
\end{proof}

\section{A contraction-stable sixteen-tree formula}
\label{sec:V132-global-tree}

Define
\begin{equation}
 \widetilde\Phi(x):=
 \bigotimes_{e\in E}\widetilde M_e(x).
 \label{eq:V132-global-interpolation}
\end{equation}
Let \(\mathbf K_I^{\rm conn}\) denote the complete regulated connected
four-row Wilson operator before the final resolvent.

\begin{theorem}[Full quartic connected operator from contractive
spectators]
\label{thm:V132-contractive-tree-formula}
The complete connected operator has the exact representation
\begin{align}
 \mathbf K_I^{\rm conn}
 &=
 \sum_{\rho\in\Pi_4}
 \mu(\rho,\boldsymbol1_I)\widetilde\Phi(\mathbf1_\rho)
 \label{eq:V132-row-Mobius}\\
 &=
 \boxed{
 \sum_{\tau\in\mathfrak T_4}
 \int_{[0,1]^3}
 (\partial_\tau\widetilde\Phi)(x^\tau(t))\,\dd t.}
 \label{eq:V132-contractive-tree-formula}
\end{align}
At every integration point, each coordinate factor not differentiated is
an operator contraction.
\end{theorem}

\begin{proof}
By \eqref{eq:V132-exact-partition-endpoint},
\[
 \widetilde\Phi(\mathbf1_\rho)
 =
 \bigotimes_{e\in E}\mathbf Q_{e,\rho}.
\]
Coordinate independence identifies this with the product, over the blocks
\(C\in\rho\), of the complete global moments of the row products in \(C\).
The moment--cumulant formula is therefore exactly
\eqref{eq:V132-row-Mobius}.  Apply
\Cref{cor:V131-tree-formula} to \(\widetilde\Phi\) to obtain
\eqref{eq:V132-contractive-tree-formula}.  The last assertion follows
from \eqref{eq:V132-interior-contraction}.
\end{proof}

\begin{assessment}[Why the full connected operator is now preferable]
\label{ass:V132-full-connected-route}
The exact formula \eqref{eq:V132-contractive-tree-formula} contains all
four V40 topology classes at once.  If its sixteen assembled derivatives
obey the normalized MTA estimate, separate puncture-word aggregation is no
longer required: \([4]\), \([3,3]\), \([3,2]\), and \([2,2,2]\) close
together.  This does not discard V129.  Its joint-capacity results remain
available when the differentiated block-merge operators are resolved into
the collision carriers on which endpoint normalization is difficult.
\end{assessment}

\section{Exact local finite differences}
\label{sec:V132-local-differences}

For \(S\subseteq\mathcal E_4\) and
\(A\subseteq\mathcal E_4\setminus S\), define
\begin{align}
 \beta_A^{S}(x)
 &:=
 \prod_{g\in A}x_g
 \prod_{g\in(\mathcal E_4\setminus S)\setminus A}(1-x_g),
 \label{eq:V132-reduced-Bernoulli}\\
 \Delta_S\mathbf Q_e(A)
 &:=
 \sum_{B\subseteq S}(-1)^{|S|-|B|}
 \mathbf Q_{e,\rho(A\cup B)}.
 \label{eq:V132-block-difference}
\end{align}

\begin{lemma}[Derivative as a convex average of block differences]
\label{lem:V132-local-derivative}
For every \(S\subseteq\mathcal E_4\),
\begin{equation}
 \boxed{
 \partial_S\widetilde M_e(x)
 =
 \sum_{A\subseteq\mathcal E_4\setminus S}
 \beta_A^{S}(x)\Delta_S\mathbf Q_e(A).}
 \label{eq:V132-local-derivative}
\end{equation}
The coefficients on the right are nonnegative and sum to one.  Hence
\begin{equation}
 \|\partial_S\widetilde M_e(x)\|_{\rm op}
 \le
 \sup_{A\subseteq\mathcal E_4\setminus S}
 \|\Delta_S\mathbf Q_e(A)\|_{\rm op}
 \le2^{|S|}.
 \label{eq:V132-crude-difference-cap}
\end{equation}
\end{lemma}

\begin{proof}
Separate in \eqref{eq:V132-local-interpolation} the choice of the edges in
\(S\) from all other edges.  Differentiation in one variable \(x_g\)
replaces the two coefficients \(1-x_g,x_g\) by \(-1,1\), respectively.
After all variables in \(S\) have been differentiated, the resulting sign
is \((-1)^{|S|-|B|}\) when precisely the subfamily \(B\subseteq S\) is
present.  The undifferentiated edges retain the coefficient
\(\beta_A^S(x)\), proving \eqref{eq:V132-local-derivative}.  Their sum is
one by the same Bernoulli-product calculation as before.  The last bound
uses \(\|\mathbf Q_{e,\rho}\|_{\rm op}\le1\) in the \(2^{|S|}\)-term
difference.
\end{proof}

\begin{lemma}[A first derivative is an exact complete covariance defect]
\label{lem:V132-first-bridge}
Let \(g=ij\notin A\).  If \(i,j\) already lie in the same component of
\(A\), then
\(\Delta_{\{g\}}\mathbf Q_e(A)=0\).  Otherwise let \(C,D\) be their two
components.  Then
\begin{equation}
 \Delta_{\{g\}}\mathbf Q_e(A)
 =
 \left(
   \mathbf Q_{e,C\cup D}
   -\mathbf Q_{e,C}\otimes\mathbf Q_{e,D}
 \right)
 \otimes
 \bigotimes_{\substack{H\in\rho(A)\\H\ne C,D}}
 \mathbf Q_{e,H}.
 \label{eq:V132-complete-covariance}
\end{equation}
Thus the differentiated edge inserts one complete covariance between the
two row components it merges, with only contraction spectators.
\end{lemma}

\begin{proof}
If \(i,j\) are already connected, then
\(\rho(A\cup\{g\})=\rho(A)\) and the difference is zero.  Otherwise adding
\(g\) merges precisely \(C,D\) and leaves every other component unchanged.
Subtract the two products \eqref{eq:V132-partial-moment}.
\end{proof}

\section{The assembled Leibniz identity}
\label{sec:V132-Leibniz}

For \(S\subseteq\mathcal E_4\), let \(E^S\) denote the set of maps
\(\varphi:S\to E\).  For \(e\in E\), put
\(S_e(\varphi):=\varphi^{-1}(\{e\})\).

\begin{theorem}[Ordered derivative allocation without a support count]
\label{thm:V132-assembled-Leibniz}
For every \(S\subseteq\mathcal E_4\),
\begin{equation}
 \boxed{
 \partial_S\widetilde\Phi(x)
 =
 \sum_{\varphi\in E^S}
 \bigotimes_{e\in E}
 \partial_{S_e(\varphi)}\widetilde M_e(x).}
 \label{eq:V132-assembled-Leibniz}
\end{equation}
Here \(\partial_\varnothing\widetilde M_e=\widetilde M_e\).
Suppose nonnegative weights \(u_{e,g}\) satisfy, for every
\(\varnothing\ne R\subseteq S\),
\begin{equation}
 \|\partial_R\widetilde M_e(x)\|_{\rm op}
 \le\prod_{g\in R}u_{e,g}.
 \label{eq:V132-local-weight-hypothesis}
\end{equation}
Then
\begin{equation}
 \boxed{
 \|\partial_S\widetilde\Phi(x)\|_{\rm op}
 \le
 \prod_{g\in S}\left(\sum_{e\in E}u_{e,g}\right).}
 \label{eq:V132-global-product-rule}
\end{equation}
No additional factor depending on \(|E|\) occurs.
\end{theorem}

\begin{proof}
The variables in \(S\) are distinct.  Apply their derivatives one at a
time to the tensor product \eqref{eq:V132-global-interpolation}.  Choosing
the coordinate factor hit by each derivative is exactly a map
\(\varphi:S\to E\), which proves
\eqref{eq:V132-assembled-Leibniz}; there are no multinomial coefficients.

Take operator norms.  For a coordinate with
\(S_e(\varphi)=\varnothing\), use
\eqref{eq:V132-interior-contraction}; otherwise use
\eqref{eq:V132-local-weight-hypothesis}.  Thus the summand indexed by
\(\varphi\) is at most
\[
 \prod_{g\in S}u_{\varphi(g),g}.
\]
Summing over all maps factorizes:
\[
 \sum_{\varphi:S\to E}\prod_{g\in S}u_{\varphi(g),g}
 =
 \prod_{g\in S}\sum_{e\in E}u_{e,g}.
\]
This is \eqref{eq:V132-global-product-rule}.
\end{proof}

\begin{corollary}[It suffices to bound finite block differences]
\label{cor:V132-differences-suffice}
If the bound
\begin{equation}
 \|\Delta_R\mathbf Q_e(A)\|_{\rm op}
 \le\prod_{g\in R}u_{e,g}
 \label{eq:V132-difference-weight}
\end{equation}
holds for every \(A\subseteq\mathcal E_4\setminus R\), then
\eqref{eq:V132-local-weight-hypothesis} holds with the same weights.
\end{corollary}

\begin{proof}
Use the convex-average identity
\eqref{eq:V132-local-derivative}.
\end{proof}

\section{The exact remaining local gate}
\label{sec:V132-bridge-gate}

For a row product \(F_C\) at coordinate \(e\), use the two-sided influence
of V50 and denote it by \(\gamma_{C,e}\).  The product rule gives
\begin{equation}
 \gamma_{C,e}\le\sum_{i\in C}\gamma_{i,e},
 \qquad
 \gamma_{i,e}\le
 C\min\{s_\alpha^{1/2}A_{i,e}^{1/2},1\}.
 \label{eq:V132-composite-influence}
\end{equation}
By \Cref{lem:V132-first-bridge} and the V50 operator covariance inequality,
the one-edge block difference obeys
\begin{equation}
 \|\Delta_{\{ij\}}\mathbf Q_e(A)\|_{\rm op}
 \le
 C\gamma_{C_i(A),e}\gamma_{C_j(A),e}
 \label{eq:V132-one-bridge-bound}
\end{equation}
when the two components are distinct, and is zero otherwise.

\begin{proof}
Equation \eqref{eq:V132-composite-influence} is the two-sided
noncommutative product rule already used in V50.  Apply the V50 covariance
inequality to the parenthesis in
\eqref{eq:V132-complete-covariance}; every other factor is a contraction.
\end{proof}

For two or three distinct forest edges, the needed statement is the
following.

\begin{openproblem}[Uniform local bridge stability through order three]
\label{op:V132-local-bridge-stability}
For every forest
\(R\subseteq\mathcal E_4\) with \(1\le|R|\le3\), every background
\(A\subseteq\mathcal E_4\setminus R\), and every coordinate \(e\), decompose
\(\Delta_R\mathbf Q_e(A)\) into a fixed number of terms, each assigning to
every bridge in \(R\) an actual pair of original rows separated by that
bridge at the moment it is inserted, and prove
\begin{equation}
 \boxed{
 \|\Delta_R\mathbf Q_e(A)\|_{\rm op}
 \le
 C_4^{|R|}
 \sum_{\psi}
 \prod_{g\in R}
 \gamma_{p_g(\psi),e}\gamma_{q_g(\psi),e}.}
 \label{eq:V132-local-bridge-stability}
\end{equation}
The number of routings \(\psi\) must depend only on four rows.  The same
decomposition must survive final-channel compression and retain the
complete Wilson multipliers needed for the V46 first-moment and V129
capacity debits.
\end{openproblem}

\begin{assessment}[Known endpoints and the new gap]
\label{ass:V132-known-local-cards}
The case \(|R|=1\) is
\eqref{eq:V132-one-bridge-bound}.  When the background is discrete and the
two or three bridges connect three or four singleton rows, the V37
one-link third cumulant and V41 one-link fourth-cumulant tree cards give
the expected powers of the influences.  Those theorems do not, as stated,
cover every non-discrete background \(A\) in
\eqref{eq:V132-block-difference}, nor do they state simultaneous
compatibility with the later resolved first-moment/capacity debits.
Consequently \eqref{eq:V132-local-bridge-stability} is a genuine finite
local gate, not a result already proved upstream.
\end{assessment}

\begin{corollary}[Conditional unnormalized global quartic closure]
\label{cor:V132-conditional-quartic}
If \eqref{eq:V132-local-bridge-stability} holds in a form that supplies
\eqref{eq:V132-difference-weight} after a fixed finite rerouting, then
every tree integrand in
\eqref{eq:V132-contractive-tree-formula} has a volume-uniform operator
bound by a fixed finite sum of products of three global pair influences.
\end{corollary}

\begin{proof}
Use \Cref{cor:V132-differences-suffice} and then
\Cref{thm:V132-assembled-Leibniz} with \(S=E(\tau)\).  Expanding only the
fixed set of original-row routings in
\eqref{eq:V132-local-bridge-stability} turns each factor
\(\sum_e u_{e,g}\) into a global V50 pair influence.  Four rows and
six possible row pairs give only a universal finite rerouting constant.
\end{proof}

\begin{firewall}[Unnormalized closure is not normalized MTA]
\label{fire:V132-not-normalized}
Even the hypothesis of
\Cref{cor:V132-conditional-quartic} would first give the operator numerator.
The BFA denominators, the complete raw output pinching, and the single final
resolvent still have to be inserted without separating the block
differences.  That resolved version is where the V126--V129 collision
capacity results may be used.  V132 therefore does not claim global
quartic MTA.
\end{firewall}

\section{V132 ledger and next acquisition}
\label{sec:V132-ledger}

\begin{theorem}[Integrated V132 statement]
\label{thm:V132-integrated}
Preserving V1--V131, V132 proves:
\begin{enumerate}[label=\textup{(V132.\arabic*)},leftmargin=7.5em]
\item the endpoint-only regression
      \eqref{eq:V132-endpoint-regression};
\item the Bernoulli partition interpolation
      \eqref{eq:V132-local-interpolation}, contractive throughout the cube;
\item the exact full-quartic sixteen-tree formula
      \eqref{eq:V132-contractive-tree-formula};
\item the convex block-difference identity
      \eqref{eq:V132-local-derivative};
\item identification of a first derivative as one complete block covariance;
\item the exact assembled Leibniz identity
      \eqref{eq:V132-assembled-Leibniz}; and
\item the factorized, volume-uniform product rule
      \eqref{eq:V132-global-product-rule}.
\end{enumerate}
\end{theorem}

\begin{proof}
The items are
\Cref{prop:V132-endpoint-regression,
thm:V132-contractive-interpolation,
thm:V132-contractive-tree-formula,
lem:V132-local-derivative,
lem:V132-first-bridge,
thm:V132-assembled-Leibniz}.
\end{proof}

\begin{openproblem}[V133 mandate: close the finite local bridge table]
\label{op:V133-mandate}
\begin{enumerate}
\item Enumerate the component partitions generated by
      \((A,R)\) for forests \(R\) of size two and three, modulo row
      relabeling.
\item Express every nonzero \(\Delta_R\mathbf Q_e(A)\) as a binary,
      ternary, or quartic connected cumulant of complete component
      products, with contraction spectators.
\item Apply the V31, V37, V41, and V50 local operator estimates only after
      that exact identity is fixed.
\item Record which cases retain a complete pair output on which the V129
      joint capacity acts after final-channel resolution.
\item If a background case produces a lower-order term lacking one bridge
      influence, exhibit it explicitly and revise the Bernoulli
      interpolation rather than hiding it in a universal cap.
\end{enumerate}
\end{openproblem}

\begin{firewall}[Claim boundary after V132]
\label{fire:V132-current-boundary}
V132 replaces the analytically unsafe interpolation by an exact
contraction-preserving one and proves that a suitable finite local bridge
bound would aggregate all marked coordinates without a volume count.  It
does not yet prove the two- and three-bridge table, normalized global
quartic MTA, all-order MTA, WUFC/CPM, continuum Yang--Mills existence, or a
positive mass gap.  The full Yang--Mills existence and mass-gap problem
remains open.
\end{firewall}

\section{Append-only record for V132}
\label{sec:V132-record}

V132 was appended to the complete V131 manuscript.  The exact V131 parent
had \(86776\) newline-delimited source records, \(3016655\) bytes, and
SHA--256 digest
\[
 \texttt{38621a99c50c9714beff26185365be5e25ed07afe1d555887e5b2667e37f7b68}.
\]
No V131 mathematical content was changed.  The sole final
\verb|\end{document}| is restored below.

% =============================================================================
% END APPENDED VERSION V132 MATERIAL
% =============================================================================

% =============================================================================
% BEGIN APPENDED VERSION V133 MATERIAL
% =============================================================================

\chapter{V133: Quotient-graph rank loss and rejection of naive local
bridge stability}
\label{ch:V133}

\section{Purpose}
\label{sec:V133-purpose}

V132 reduced its contraction-preserving interpolation to the local
finite differences \(\Delta_R\mathbf Q_e(A)\).  The proposed bridge table
implicitly expected one small covariance factor for every differentiated
edge in \(R\).  This continuation tests that expectation before attempting
a long case enumeration.  The test fails: undifferentiated background
edges can identify original vertices and turn distinct differentiated
edges into parallel edges in the quotient graph.  A mixed finite
difference along \(k\) parallel edges is only one covariance defect, not a
product of \(k\) defects.

\begin{firewall}[What survives V133]
\label{fire:V133-survival}
The exact endpoint identity, cube contraction, full connected forest
formula, local derivative formula, and assembled Leibniz identity of V132
remain correct.  Only the unproved bridge-stability proposal
\eqref{eq:V132-local-bridge-stability} and its conditional use for a
three-link numerator are rejected.
\end{firewall}

\section{The quotient multigraph of a block difference}
\label{sec:V133-quotient}

Fix one coordinate and suppress its index \(e\).  Let
\(A\subseteq\mathcal E_4\setminus R\).  Contract every connected component
of \(A\) to one vertex.  Each edge \(g=ij\in R\) induces an edge
\(\bar g\) between the contracted components containing \(i,j\).  Loops
and parallel edges are retained.  Denote the resulting multigraph by
\begin{equation}
 H(A;R):=R/A.
 \label{eq:V133-quotient-graph}
\end{equation}

\begin{lemma}[The finite difference depends only on the quotient
multigraph]
\label{lem:V133-quotient-dependence}
For \(B\subseteq R\), the row partition
\(\rho(A\cup B)\) is obtained by taking the connected-component partition
of the sub-multigraph of \(H(A;R)\) with edge set \(\bar B\), then expanding
the contracted vertices back to their \(A\)-components.  Consequently
\(\Delta_R\mathbf Q(A)\) depends on \(A\) only through its component
partition and on \(R\) only through the labeled quotient multigraph
\(H(A;R)\).
\end{lemma}

\begin{proof}
Every path in \(A\cup B\) alternates between paths internal to
\(A\)-components and edges of \(B\).  Contracting the former gives a path
in the quotient subgraph.  Conversely every quotient path lifts by
choosing, within each \(A\)-component, paths between the relevant edge
endpoints.  Thus the two component relations coincide.  Insert this
description in \eqref{eq:V132-block-difference}.
\end{proof}

\begin{corollary}[A loop annihilates the mixed difference]
\label{cor:V133-loop-zero}
If one edge of \(R\) becomes a loop in \(H(A;R)\), then
\begin{equation}
 \Delta_R\mathbf Q(A)=0.
 \label{eq:V133-loop-zero}
\end{equation}
\end{corollary}

\begin{proof}
Let \(g\) be the loop.  For every
\(B\subseteq R\setminus\{g\}\),
\(\rho(A\cup B\cup\{g\})=\rho(A\cup B)\).  Pair the two terms of
\eqref{eq:V132-block-difference} which differ only by the presence of
\(g\); their signs are opposite.
\end{proof}

\section{Parallel differentiated edges collapse to one bridge}
\label{sec:V133-parallel}

\begin{theorem}[Exact parallel-edge collapse]
\label{thm:V133-parallel-collapse}
Suppose the quotient \(H(A;R)\) has exactly two vertices \(C,D\), and all
\(k:=|R|\ge1\) differentiated edges are parallel edges joining them.  Put
\begin{align}
 Q_0&:=\mathbf Q_C\otimes\mathbf Q_D,
 &
 Q_1&:=\mathbf Q_{C\cup D}.
 \label{eq:V133-two-component-moments}
\end{align}
Then
\begin{equation}
 \boxed{
 \Delta_R\mathbf Q(A)
 =(-1)^{k+1}(Q_1-Q_0).}
 \label{eq:V133-parallel-collapse}
\end{equation}
Thus its analytic order is that of one complete covariance, independently
of the number \(k\) of parallel differentiated edges.
\end{theorem}

\begin{proof}
For \(B=\varnothing\), the component partition gives \(Q_0\).  Every
nonempty \(B\subseteq R\) joins \(C,D\) and gives \(Q_1\).  Therefore
\begin{align*}
 \Delta_R\mathbf Q(A)
 &=
 (-1)^kQ_0+
 \sum_{j=1}^k\binom{k}{j}(-1)^{k-j}Q_1\\
 &=
 (-1)^kQ_0-(-1)^kQ_1,
\end{align*}
because the full binomial sum is zero.  This is
\eqref{eq:V133-parallel-collapse}.
\end{proof}

\begin{example}[Four-row two-edge collapse]
\label{ex:V133-four-row-parallel}
Take
\begin{equation}
 R=\{12,34\},
 \qquad
 A=\{13,24\}.
 \label{eq:V133-parallel-example-graphs}
\end{equation}
The background components are
\(C=\{1,3\}\) and \(D=\{2,4\}\).  Both differentiated edges join
\(C\) to \(D\), so
\begin{equation}
 \Delta_{\{12,34\}}\mathbf Q(A)
 =
 \mathbf Q_C\otimes\mathbf Q_D-\mathbf Q_I.
 \label{eq:V133-parallel-example}
\end{equation}
This is minus one binary covariance defect.
\end{example}

\begin{proof}
Apply \Cref{thm:V133-parallel-collapse} with \(k=2\).
\end{proof}

\begin{proposition}[The two-bridge product bound is false algebraically]
\label{prop:V133-product-bound-false}
There is no universal \(C\) for which every bounded scalar moment system
in \Cref{ex:V133-four-row-parallel} satisfies
\begin{equation}
 |\Delta_{\{12,34\}}Q(A)|
 \le C\,w_{12}w_{34}
 \label{eq:V133-false-product}
\end{equation}
whenever each \(w_{12},w_{34}\) is required to have the order of the
single covariance joining the two quotient components.
\end{proposition}

\begin{proof}
Let \(0<s\le1\), let \(\xi\) be a symmetric sign, and put
\(X=Y=\sqrt{s}\,\xi\).  Choose bounded row functions whose products on
\(C,D\) are \(X,Y\), respectively; two of the four row functions may be
the constant one.  Then \(Q_C=Q_D=0\) and
\(Q_I=\mathbb E(XY)=s\).  Equation
\eqref{eq:V133-parallel-example} has modulus \(s\).
Each of the two parallel bridges sees the same quotient covariance and
therefore has single-bridge order \(s\).  The proposed right side has
order \(Cs^2\), which cannot dominate \(s\) as \(s\downarrow0\).
\end{proof}

\begin{firewall}[Logical versus Wilson counterexample]
\label{fire:V133-logical-counterexample}
The scalar construction proves that the component-partition finite
difference does not algebraically contain two covariance factors.  It is
not asserted to be a Wilson representation family.  A special Wilson
identity could only repair the estimate by adding new structure not present
in the V132 quotient-graph algebra; no such identity has been proved.
\end{firewall}

\section{Rank, not derivative order, counts forced mergers}
\label{sec:V133-rank}

Let \(q(A)\) be the number of components of \(A\) and define the merger
rank
\begin{equation}
 r(A;R):=q(A)-q(A\cup R).
 \label{eq:V133-merger-rank}
\end{equation}
Always \(0\le r(A;R)\le|R|\).

\begin{lemma}[Equality of rank and derivative count]
\label{lem:V133-rank-equality}
The equality \(r(A;R)=|R|\) holds exactly when the quotient multigraph
\(H(A;R)\) is a forest without loops or parallel edges.  If the quotient
has a loop, a parallel edge, or a cycle, then
\begin{equation}
 r(A;R)<|R|.
 \label{eq:V133-rank-loss}
\end{equation}
\end{lemma}

\begin{proof}
Adding one non-loop edge decreases the component count by one exactly when
its endpoints were previously in different components.  A loop decreases
it by zero.  A loop-free multigraph with \(v\) vertices and \(c\)
components has rank \(v-c\); this equals its number of edges exactly when
it is a simple forest.  Parallel edges form a two-cycle in the multigraph
sense and hence create excess edge count.  Applying this to \(H(A;R)\)
gives the claim.
\end{proof}

\begin{assessment}[Why the V132 local gate fails]
\label{ass:V133-V132-gate-fails}
The desired three-link tree normalization counts \(|R|\) marked links.
But \(\Delta_R\mathbf Q(A)\) can force only \(r(A;R)\) genuine component
mergers.  The difference between these integers is not a constant loss:
it is a loss of a small heat/influence factor.  The parallel example has
\(|R|=2\) and \(r=1\), and
\Cref{prop:V133-product-bound-false} shows that a universal cap cannot
manufacture the missing factor.  Therefore
\eqref{eq:V132-local-bridge-stability} is false as stated.
\end{assessment}

\section{The obstruction is present at forest weakening points}
\label{sec:V133-weakening}

The preceding failure is not confined to a cube point unused by the forest
formula.  It is already visible for three rows.  Write
\begin{align*}
 q_0&:=Q_1Q_2Q_3,&
 q_{12}&:=Q_{12}Q_3,&
 q_{13}&:=Q_{13}Q_2,\\
 q_{23}&:=Q_{23}Q_1,&
 q_{123}&:=Q_{123},
\end{align*}
and abbreviate
\(a=x_{12}\), \(b=x_{13}\), \(c=x_{23}\).  The three-row version of the
Bernoulli interpolation is
\begin{align}
 \widetilde M(a,b,c)
 ={}&
 q_0(1-a)(1-b)(1-c)
 +q_{12}a(1-b)(1-c)\notag\\
 &+q_{13}b(1-a)(1-c)
 +q_{23}c(1-a)(1-b)\notag\\
 &+q_{123}\bigl(ab+ac+bc-2abc\bigr).
 \label{eq:V133-three-row-reliability}
\end{align}

\begin{lemma}[A tree derivative contains lower-rank moments]
\label{lem:V133-three-row-derivative}
For the tree edges \(12,13\),
\begin{align}
 \partial_a\partial_b\widetilde M(a,b,c)
 ={}&
 (1-c)(q_0-q_{12}-q_{13})
 +c q_{23}
 +(1-2c)q_{123}.
 \label{eq:V133-three-row-derivative}
\end{align}
At the forest weakening point \(c=\min\{a,b\}\), the \(q_{23}\) and
lower-partition contributions do not vanish identically.
\end{lemma}

\begin{proof}
Differentiate \eqref{eq:V133-three-row-reliability}.  The connected-graph
probability \(ab+ac+bc-2abc\) has mixed derivative \(1-2c\), giving the
display.  Substitution of \(c=\min\{a,b\}\) proves the last statement.
\end{proof}

\begin{corollary}[Individual V132 tree integrands are not
derivative-faithful]
\label{cor:V133-not-faithful}
Contraction of every unmarked local factor in V132 does not imply that one
tree integrand carries one small influence for each of its tree edges.
The missing lower-rank terms cancel only after the complete forest/Moebius
sum, and taking a norm of each tree integrand separately may destroy that
cancellation.
\end{corollary}

\begin{proof}
\Cref{lem:V133-three-row-derivative} explicitly displays terms which are
not a two-bridge connected cumulant, even at the prescribed weakening
point.  The four-row parallel collapse gives the corresponding mixed
difference mechanism inside the V132 local formula.  Exact equality of the
sum of all forest integrals with the cumulant is compatible with such
cross-tree cancellation, but subadditivity of a norm is not.
\end{proof}

\section{Required replacement: derivative fidelity}
\label{sec:V133-replacement}

\begin{definition}[Derivative-faithful quartic interpolation]
\label{def:V133-derivative-faithful}
An interpolation \(\Psi\) is derivative-faithful at replica order four if:
\begin{enumerate}[label=\textup{(DF\arabic*)},leftmargin=4.5em]
\item its partition endpoints give the exact row moments, so their Moebius
      sum is the complete fourth cumulant;
\item every unmarked coordinate factor along the integration domain is an
      operator contraction; and
\item each differentiated spanning-tree integrand admits, before trace
      norm, an exact finite decomposition in which all three tree links are
      genuine local defects or one already-proved local fourth-cumulant
      tree fibre.  No term may have merger rank below three unless the
      missing factor is present elsewhere in the same assembled term.
\end{enumerate}
\end{definition}

\begin{proposition}[The Bernoulli component interpolation is not
derivative-faithful]
\label{prop:V133-Bernoulli-not-faithful}
The interpolation \(\widetilde\Phi\) of V132 satisfies
\textup{(DF1)} and \textup{(DF2)}, but fails \textup{(DF3)}.
\end{proposition}

\begin{proof}
\textup{(DF1)} is \Cref{thm:V132-contractive-tree-formula}, and
\textup{(DF2)} is \Cref{thm:V132-contractive-interpolation}.
The parallel collapse
\eqref{eq:V133-parallel-example} has two differentiated edges and only one
defect; \Cref{prop:V133-product-bound-false} rules out recovery of the
second factor by a uniform estimate.  The same lower-rank mechanism embeds
in a three-edge derivative by adjoining a third independent merge.
\end{proof}

\begin{assessment}[Upstream direction]
\label{ass:V133-upstream-direction}
The next construction must not interpolate arbitrary graph connectivity by
independent Bernoulli edges.  Two viable algebraic sources remain:
row-resampling difference operators, whose repeated differences are tied
to distinct copies of the base variable, and an exact cumulant recursion
kept assembled until the V41 local fourth-order excess has supplied the
third link.  In either route, derivative order must count actual
fluctuation insertions rather than redundant edges of a quotient graph.
\end{assessment}

\section{V133 ledger and next acquisition}
\label{sec:V133-ledger}

\begin{theorem}[Integrated V133 statement]
\label{thm:V133-integrated}
Preserving V1--V132, V133 proves:
\begin{enumerate}[label=\textup{(V133.\arabic*)},leftmargin=7.5em]
\item exact quotient-multigraph control of every V132 block difference;
\item annihilation by quotient loops;
\item the parallel-edge collapse
      \eqref{eq:V133-parallel-collapse};
\item a scalar small-parameter disproof of a product tax for two parallel
      differentiated bridges;
\item the merger-rank criterion
      \eqref{eq:V133-rank-loss};
\item an explicit lower-rank term inside a forest weakening integrand; and
\item failure of derivative fidelity for the Bernoulli component
      interpolation.
\end{enumerate}
\end{theorem}

\begin{proof}
The items are
\Cref{lem:V133-quotient-dependence,cor:V133-loop-zero,
thm:V133-parallel-collapse,prop:V133-product-bound-false,
lem:V133-rank-equality,lem:V133-three-row-derivative,
prop:V133-Bernoulli-not-faithful}.
\end{proof}

\begin{openproblem}[V134 mandate: a fluctuation-faithful tree identity]
\label{op:V134-mandate}
\begin{enumerate}
\item Start from the exact fourth row-cumulant Moebius formula, not from
      independent graph-edge interpolation.
\item Introduce independent resampled copies of each coordinate variable
      and derive covariance as a first difference identity.
\item Iterate the identity through order four, keeping all cancellation
      terms assembled, and seek a finite sum indexed by the sixteen
      labeled trees in which three distinct difference operators act.
\item At one-coordinate coincidences, insert the complete V37/V41 local
      cumulant operator rather than forcing it into a product of binary
      covariances.
\item Require contraction spectators and test the two-parallel-edge
      regression \eqref{eq:V133-parallel-example} before claiming a
      derivative-faithful formula.
\end{enumerate}
\end{openproblem}

\begin{firewall}[Claim boundary after V133]
\label{fire:V133-current-boundary}
V133 proves that the V132 local bridge-stability route cannot supply three
marked factors term by term.  It does not disprove the exact V132
connected formula, but it prevents its use via separate tree-integrand
norms.  A derivative-faithful replacement, normalized global quartic MTA,
all-order MTA, WUFC/CPM, continuum Yang--Mills existence, and a positive
mass gap remain open.  The full Yang--Mills existence and mass-gap problem
is not proved.
\end{firewall}

\section{Append-only record for V133}
\label{sec:V133-record}

V133 was appended to the complete V132 manuscript.  The exact V132 parent
had \(87300\) newline-delimited source records, \(3035567\) bytes, and
SHA--256 digest
\[
 \texttt{83aec4a33cdb519e1ad72ea858ddd7201f4cfcbb75e4991a45962e972b8dfc6c}.
\]
No V132 mathematical content was changed.  The sole final
\verb|\end{document}| is restored below.

% =============================================================================
% END APPENDED VERSION V133 MATERIAL
% =============================================================================

% =============================================================================
% BEGIN APPENDED VERSION V134 MATERIAL
% =============================================================================

\chapter{V134: Low-activity global quartic aggregation by canonical
minimal skeletons}
\label{ch:V134}

\section{Purpose}
\label{sec:V134-purpose}

V133 showed that a norm may not be taken tree-integrand by
tree-integrand in the Bernoulli interpolation.  The correct upstream
object is the exact V39 connected coordinate-word sum.  In the regime
where the total norm of all nonsingleton local letters is small, that sum
can be handled directly: select one canonical V40 minimal skeleton, keep
its letters marked, and sum every remaining local choice through one
product envelope.  This continuation proves that low-activity theorem for
all four quartic topology classes at once.

\begin{firewall}[Scope of V134]
\label{fire:V134-scope}
The proof below uses absolute values only after the total local activity
has been assumed at most one.  It does not assert that this hypothesis
holds uniformly for arbitrary supports or representation heights.  The
large-activity case requires a cluster/cut decomposition and remains open.
\end{firewall}

\section{Local connected letters and their activities}
\label{sec:V134-letters}

Let \(I=\{1,2,3,4\}\).  At coordinate \(e\), let
\(m_e\) be the all-singleton local letter.  For each pair, triple, and
quadruple of distinct rows, put
\begin{align}
 p_{ij,e}&:=\|\mathbf d_{ij,e}\|_{\rm op},
 \label{eq:V134-pair-activity}\\
 t_{ijk,e}&:=\|\mathbf t_{ijk,e}\|_{\rm op},
 \label{eq:V134-triple-activity}\\
 u_e&:=\|\mathbf u_{1234,e}\|_{\rm op},
 \label{eq:V134-four-activity}
\end{align}
where the bold symbols are the complete local cumulant operators.  An
inactive block has activity zero.  For a double-pair local partition
\(ij\mid k\ell\), submultiplicativity gives the activity
\(p_{ij,e}p_{k\ell,e}\).

Define
\begin{align}
 a_e:={}&
 \sum_{1\le i<j\le4}p_{ij,e}
 +\sum_{1\le i<j<k\le4}t_{ijk,e}
 +u_e\notag\\
 &+
 p_{12,e}p_{34,e}
 +p_{13,e}p_{24,e}
 +p_{14,e}p_{23,e},
 \label{eq:V134-local-total-activity}\\
 \mathcal A&:=\sum_{e\in E}a_e.
 \label{eq:V134-total-activity}
\end{align}
Every singleton moment and hence \(m_e\) is a contraction.

For global activities, set
\begin{align}
 P_{ij}&:=\sum_{e\in E}p_{ij,e},
 &
 T_{ijk}&:=\sum_{e\in E}t_{ijk,e},
 &
 U_4&:=\sum_{e\in E}u_e.
 \label{eq:V134-global-activities}
\end{align}

\section{The quartic skeleton majorant}
\label{sec:V134-majorant}

Write \(\mathfrak T_4\) for the sixteen labeled pair-edge trees.  Let
\(\mathfrak P_{33}\) be the six unordered pairs of distinct three-element
subsets \(B,C\subset I\); every such pair has \(B\cup C=I\).  Define
\begin{align}
 \mathcal S_4:={}&U_4
 +\sum_{\substack{B\subset I,\ |B|=3\\
                  o=I\setminus B}}
 T_B\sum_{i\in B}P_{oi}
 \notag\\
 &+\sum_{\{B,C\}\in\mathfrak P_{33}}T_BT_C
 +\sum_{\tau\in\mathfrak T_4}
   \prod_{ij\in E(\tau)}P_{ij}.
 \label{eq:V134-skeleton-majorant}
\end{align}
Here \(T_B=T_{ijk}\) when \(B=\{i,j,k\}\).

\begin{lemma}[The majorant is exactly the V40 skeleton list]
\label{lem:V134-majorant-list}
The four lines of activity in
\eqref{eq:V134-skeleton-majorant} correspond respectively to:
\([4]\), \([3,2]\), \([3,3]\), and \([2,2,2]\).
Every inclusion-minimal connected family of nonsingleton local blocks
contributes to at least one displayed monomial.  If two disjoint pair
blocks occur at the same coordinate, their joint activity is still
dominated by the relevant pair-tree product.
\end{lemma}

\begin{proof}
The classification is
\Cref{thm:V40-quartic-skeleton-list}.  A four-row block contributes
\(u_e\).  A triple on \(B\) and an attaching pair \(oi\), with
\(o=I\setminus B\), contribute \(t_{B,e}p_{oi,f}\); their blocks overlap
in row \(i\), so coordinate-partition compatibility forces \(e\ne f\).
Dropping that restriction only enlarges \(T_BP_{oi}\).

Two distinct triples contribute \(t_{B,e}t_{C,f}\), again with
different coordinates because they overlap in two rows.  Three pair
blocks form a labeled tree.  If two disjoint tree edges occur in one
coordinate, their local partition letter has norm at most the product of
their two pair activities; summing it together with the third edge is
bounded by the corresponding product of the three \(P_{ij}\).  Pair edges
sharing a row cannot occur as two blocks of one coordinate partition.
\end{proof}

\section{One marked skeleton and one product envelope}
\label{sec:V134-envelope}

Let \(\mathcal W_{\rm conn}\) be the set of coordinate words in the V39
formula whose row hypergraph is connected.  Fix the V40 ordering of all
eligible minimal connected block families and let
\(\operatorname{sk}(w)\) be the first such family in a connected word
\(w\).  This partitions \(\mathcal W_{\rm conn}\) into disjoint skeleton
fibres.

\begin{lemma}[Unmarked local choices have one exponential envelope]
\label{lem:V134-product-envelope}
For any fixed collection of marked local letters, the sum of the operator
norms of all compatible unmarked local choices is at most
\begin{equation}
 \prod_{e\in E}(1+a_e)
 \le e^{\mathcal A}.
 \label{eq:V134-product-envelope}
\end{equation}
This bound remains valid when marked coordinates are omitted from the
product.
\end{lemma}

\begin{proof}
At coordinate \(e\), the all-singleton letter has norm at most one.
The six single-pair, four triple, one quadruple, and three double-pair
partition letters have total norm at most \(a_e\) by
\eqref{eq:V134-local-total-activity}.  Distributing the tensor product and
using submultiplicativity gives the first product.  The scalar inequality
\(1+x\le e^x\) gives the second.  Replacing selected factors by one only
enlarges the same bound.
\end{proof}

\begin{theorem}[Low-activity global quartic aggregation]
\label{thm:V134-low-activity-quartic}
Let \(\mathbf K_I^{\rm conn}\) be the complete four-row connected operator
in the exact V39 replica formula.  Then
\begin{equation}
 \boxed{
 \|\mathbf K_I^{\rm conn}\|_{\rm op}
 \le C_4 e^{\mathcal A}\mathcal S_4.}
 \label{eq:V134-general-quartic-bound}
\end{equation}
In particular, if \(\mathcal A\le1\),
\begin{equation}
 \boxed{
 \|\mathbf K_I^{\rm conn}\|_{\rm op}
 \le C_4e\,\mathcal S_4,}
 \label{eq:V134-low-activity-bound}
\end{equation}
with no dependence on \(|E|\), coordinate-subset multiplicities, or the
number of local partition words.
\end{theorem}

\begin{proof}
Expand the exact connected operator into its coordinate words, but assign
each word to the unique fibre of its canonical first skeleton.  In one
fibre keep every skeleton letter marked.  Its operator norm is at most the
product of the corresponding pair, triple, or quadruple activities.
After those factors are removed, discard the restrictions saying that the
marked skeleton is first and that no earlier skeleton occurs.  This only
enlarges an absolute majorant.  Sum every remaining local choice with
\Cref{lem:V134-product-envelope}.

Now sum the marked skeleton choices.  By
\Cref{lem:V134-majorant-list}, their activity sum is bounded by
\(\mathcal S_4\), up to the fixed overcount caused by orderings of at most
four rows and at most three marked blocks.  Denote that numerical overcount
by \(C_4\).  This proves
\eqref{eq:V134-general-quartic-bound}.  If \(\mathcal A\le1\), replace
\(e^{\mathcal A}\) by \(e\).
\end{proof}

\begin{remark}[Why no coordinate-subset factor occurs]
\label{rem:V134-no-subset-factor}
The proof does sum marked coordinate positions, but it does so through the
global activities \(P_{ij},T_{ijk},U_4\).  Every unmarked position is
summed once by the product \(\prod_e(1+a_e)\).  There is no binomial
coefficient depending on support size, and collision letters
\(ij\mid k\ell\) remain one local partition choice until their two defect
norms have been multiplied.
\end{remark}

\section{Insertion of the established Wilson local cards}
\label{sec:V134-Wilson-cards}

The V50 influence and V37/V41 one-link estimates give, uniformly in
representations,
\begin{align}
 p_{ij,e}
 &\le\gamma_{i,e}\gamma_{j,e},
 &
 \gamma_{i,e}
 &\le C\min\{s_\alpha^{1/2}A_{i,e}^{1/2},1\},
 \label{eq:V134-pair-Wilson}\\
 t_{ijk,e}
 &\le C\min\{s_\alpha^2A_{i,e}A_{j,e}A_{k,e},1\},
 \label{eq:V134-triple-Wilson}
\end{align}
while V41 decomposes \(u_e\) into the sixteen local tree weights
\begin{equation}
 u_e\le
 C\sum_{\tau\in\mathfrak T_4}
 \prod_{ij\in E(\tau)}
 s_\alpha\sqrt{A_{i,e}A_{j,e}},
 \label{eq:V134-four-Wilson}
\end{equation}
with the capped interpretation proved there.

\begin{corollary}[Low-activity unnormalized Wilson tree numerator]
\label{cor:V134-Wilson-low-activity}
Under \(\mathcal A\le1\), the complete connected quartic operator is
bounded by a fixed finite sum of the V40 hypergraph-skeleton weights.
After applying the V50 junction-to-tree routing to each \(T_{ijk}\) and
the V41 tree decomposition to \(U_4\), every summand contains a labeled
spanning-tree set of three Wilson link marks.  All surplus local activity
is absorbed by the universal factor \(e\).
\end{corollary}

\begin{proof}
Insert \eqref{eq:V134-pair-Wilson}--
\eqref{eq:V134-four-Wilson} into
\eqref{eq:V134-low-activity-bound}.  For a triple block, use the V50
operator numerator and junction-to-tree lemma to choose two incident
row links.  In a \([3,2]\) skeleton the attaching pair supplies the third
link.  Two triples have one surplus link beyond a spanning tree; its
capped activity is at most one in the present envelope and may be
discarded after a spanning tree is selected.  A pair-tree and a four-row
block already have the required V40/V41 tree routing.  Only finitely many
row labels and tree choices occur.
\end{proof}

\begin{firewall}[What the low-activity theorem does not normalize]
\label{fire:V134-not-normalized}
\Cref{cor:V134-Wilson-low-activity} is an operator-numerator statement.
It does not by itself supply the repeated global input-mass denominators,
the exact Fourier/BFA coefficient normalization, or the single final
resolvent.  Those structures remain assembled in the earlier V46--V129
cards.  The theorem proves that coordinate-word aggregation creates no
additional multiplicity loss in the low-activity branch; it does not yet
prove normalized quartic MTA there.
\end{firewall}

\section{The precise high-activity residual}
\label{sec:V134-high-activity}

If \(\mathcal A>1\), the exponential envelope in
\eqref{eq:V134-general-quartic-bound} is not uniform.  Taking the universal
fourth-cumulant cap alone also fails when the desired tree majorant is
small.  The correct residual question is therefore a cluster separation
question.

\begin{openproblem}[Quartic high-activity cluster cut]
\label{op:V134-high-activity-cut}
Given \(\mathcal A>1\), find either:
\begin{enumerate}[label=\textup{(\roman*)},leftmargin=3.7em]
\item three row-pair activities containing a labeled spanning tree and
      each bounded below by a fixed numerical threshold; or
\item a nontrivial row partition \(\rho\) such that all activity crossing
      the cut of \(\rho\) is small.
\end{enumerate}
In branch \textup{(i)}, the universal connected-operator cap should be
absorbed by the three saturated tree weights.  In branch \textup{(ii)},
rewrite the fourth cumulant as a complete covariance across the weak cut,
plus lower cumulants whose cancellation is retained, and apply the V50
operator covariance estimate before resolving outputs.
\end{openproblem}

\begin{assessment}[Why total activity alone is insufficient]
\label{ass:V134-total-not-connected}
The inequality \(\mathcal A>1\) may be caused by arbitrarily large
activity internal to one proper row cluster, for example only the pair
type \(12\).  It does not imply that any spanning-tree product is bounded
below.  The next dichotomy must therefore use cut activity or a maximum
spanning tree, not the scalar total \(\mathcal A\).
\end{assessment}

\section{V134 ledger and next acquisition}
\label{sec:V134-ledger}

\begin{theorem}[Integrated V134 statement]
\label{thm:V134-integrated}
Preserving V1--V133, V134 proves:
\begin{enumerate}[label=\textup{(V134.\arabic*)},leftmargin=7.5em]
\item a complete scalar activity alphabet for the fifteen quartic local
      partitions;
\item the exact V40 skeleton majorant
      \eqref{eq:V134-skeleton-majorant};
\item a canonical first-skeleton partition of all connected words;
\item the single product envelope
      \eqref{eq:V134-product-envelope};
\item the global bound
      \eqref{eq:V134-general-quartic-bound}; and
\item uniform coordinate-word aggregation in the branch
      \(\mathcal A\le1\), with all four quartic topology types retained.
\end{enumerate}
\end{theorem}

\begin{proof}
The items are
\Cref{lem:V134-majorant-list,lem:V134-product-envelope,
thm:V134-low-activity-quartic,cor:V134-Wilson-low-activity}.
\end{proof}

\begin{openproblem}[V135 mandate: audit, then the high-activity cut]
\label{op:V135-mandate}
\begin{enumerate}
\item Perform the prescribed read-only audit of the newest SM and NS
      notes and transfer only exact finite-structure lessons.
\item Prove a finite weighted-graph lemma on four row vertices: either a
      spanning tree has three quantitatively strong edges, or some
      nontrivial cut has quantitatively small total crossing activity.
\item Lift the weak-cut alternative to an exact fourth-cumulant covariance
      identity without norming its cancellation terms separately.
\item Use the strong-tree alternative only after checking that the
      universal operator cap and the normalized BFA/resolvent target have
      the same saturation scale.
\item Combine the resulting high-activity theorem with
      \Cref{thm:V134-low-activity-quartic}; do not claim global quartic MTA
      until the normalized output quotient is also closed.
\end{enumerate}
\end{openproblem}

\begin{firewall}[Claim boundary after V134]
\label{fire:V134-current-boundary}
V134 proves uniform global coordinate-word aggregation for the full
quartic connected operator in the low-total-activity regime.  It does not
close the high-activity cluster cut, normalized quartic MTA, all-order MTA,
WUFC/CPM, continuum Yang--Mills existence, or a positive mass gap.  The
full Yang--Mills existence and mass-gap problem remains open.
\end{firewall}

\section{Append-only record for V134}
\label{sec:V134-record}

V134 was appended to the complete V133 manuscript.  The exact V133 parent
had \(87711\) newline-delimited source records, \(3051393\) bytes, and
SHA--256 digest
\[
 \texttt{c7cbbf6894a94fd856df7e71150e31473ef39df9587fd9dd06af7af418c85d5a}.
\]
No V133 mathematical content was changed.  The sole final
\verb|\end{document}| is restored below.

% =============================================================================
% END APPENDED VERSION V134 MATERIAL
% =============================================================================

% =============================================================================
% BEGIN APPENDED VERSION V135 MATERIAL
% =============================================================================

\chapter{V135: Companion audit and a threshold-cluster proof of the
global quartic operator numerator}
\label{ch:V135}

\section{Purpose and mandatory companion audit}
\label{sec:V135-purpose}

V134 closed coordinate-word aggregation when the total local connected
activity is small.  The remaining numerator problem is the complementary
regime, in which internal activity may be arbitrarily large while every
connection across one row cut is weak.  This version first performs the
prescribed companion audit and then resolves that dichotomy by thresholding
the six global row-pair link weights.

\begin{assessment}[Mandatory V135 companion audit]
\label{ass:V135-audit}
The read-only audit found:
\begin{center}
\begin{tabular}{llll}
\toprule
programme & version & source records/bytes & SHA--256\\
\midrule
SM--Einstein & V154 & \(97623/3572486\) &
\texttt{6b77428d5d6952af7bf5a55ebe314fd304b60d27b668d75c82bb5ccffc263383}\\
Navier--Stokes & V31 & \(23052/887537\) &
\texttt{f1121b62238ad5491bb7cf03635ea9934942edf573ed8faaa9ee79e1d38a6696}\\
\bottomrule
\end{tabular}
\end{center}
The NS file is unchanged.  Since the V130 audit, the SM programme advanced
from V147 to V154.  The structurally relevant lessons are: its V148 and
V150--V151 arguments pass complete connected or Ward expressions before
norming their rows; V152--V153 distinguish global boundedness and
conservation from actual removal of a critical concentration defect; and
V154 goes upstream to an exact principal diagonalization but explicitly
does not infer compactness from it.
\end{assessment}

\begin{firewall}[Audit transfer boundary]
\label{fire:V135-audit-boundary}
No SM stress estimate, compactness theorem, or Einstein-frame argument is
imported into Yang--Mills.  The only transfer is proof order: the complete
quartic cut identity is formed before its lower cumulant rows are bounded,
and a universal operator cap is not treated as though it already contained
three marked links.
\end{firewall}

\section{Global link weights and the tree polynomial}
\label{sec:V135-link-weights}

For the four input rows, retain the V50 overlap masses
\begin{equation}
 H_{ij}:=\sum_{e\in X_i\cap X_j}A_{i,e}A_{j,e}
 \label{eq:V135-overlap}
\end{equation}
and define the untruncated link weights
\begin{equation}
 \Lambda_{ij}:=s_\alpha H_{ij}\ge0.
 \label{eq:V135-link-weight}
\end{equation}
All fixed constants in the V31/V37/V41/V50 local estimates are kept
outside \(\Lambda\).  Put
\begin{equation}
 \mathcal T_4(\Lambda):=
 \sum_{\tau\in\mathfrak T_4}
 \prod_{ij\in E(\tau)}\Lambda_{ij}.
 \label{eq:V135-tree-polynomial}
\end{equation}
For disjoint nonempty row sets \(C,D\), abbreviate
\begin{equation}
 \Lambda(C,D):=
 \sum_{\substack{i\in C\\j\in D}}\Lambda_{ij}.
 \label{eq:V135-cut-weight}
\end{equation}

\begin{lemma}[Universal quartic cap]
\label{lem:V135-universal-cap}
For bounded row operator functions \(\|F_i\|_{\rm op}\le1\), their complete
fourth operator cumulant satisfies
\begin{equation}
 \|\kappa_4(F_1,F_2,F_3,F_4)\|_{\rm op}\le C_{\rm cum,4}.
 \label{eq:V135-universal-cap}
\end{equation}
\end{lemma}

\begin{proof}
Use the fifteen-term moment--cumulant formula.  Every block moment is an
operator contraction, products of block moments act on disjoint row
carriers, and the sum of the absolute Moebius coefficients is a numerical
constant depending only on four.
\end{proof}

\section{Exact product--cumulant identities across row clusters}
\label{sec:V135-cluster-identities}

Let \(\rho=\{C,D\}\) be a two-block row partition and write
\(F_C:=\prod_{i\in C}F_i\), with the fixed tensor order.

\begin{lemma}[Exact two-cluster cut identity]
\label{lem:V135-two-cluster-identity}
One has
\begin{equation}
 \operatorname{Cov}(F_C,F_D)
 =
 \sum_{\substack{\pi\in\Pi_4\\
                  \pi\vee\rho=\boldsymbol1_I}}
 \prod_{B\in\pi}\kappa(F_i:i\in B).
 \label{eq:V135-product-cumulant}
\end{equation}
The term \(\pi=\boldsymbol1_I\) is the desired fourth cumulant.  Every
other term contains at least one cumulant block meeting both \(C\) and
\(D\).
\end{lemma}

\begin{proof}
Expand both moments in
\(\mathbb E(F_CF_D)-\mathbb EF_C\,\mathbb EF_D\) into cumulants of the
individual rows.  The first expansion contains every
\(\pi\in\Pi_4\).  The product of the two separate moment expansions
contains exactly the partitions refined by \(\rho\).  Their difference
therefore retains precisely the partitions whose join with \(\rho\) is
the top partition.  Such a partition has a block meeting both sides,
and the top partition itself occurs once.
\end{proof}

\begin{theorem}[Complete quartic weak-cut bound]
\label{thm:V135-weak-cut}
For every nontrivial row cut \(C\mid D\),
\begin{equation}
 \boxed{
 \|\kappa_4(F_1,F_2,F_3,F_4)\|_{\rm op}
 \le C_4\min\{1,\Lambda(C,D)\}.}
 \label{eq:V135-weak-cut}
\end{equation}
The estimate is stable under identity spectator amplification.
\end{theorem}

\begin{proof}
The V50 operator covariance inequality applied to \(F_C,F_D\), together
with its composite influence product rule, gives
\begin{equation}
 \|\operatorname{Cov}(F_C,F_D)\|_{\rm op}
 \le C\Lambda(C,D).
 \label{eq:V135-composite-covariance}
\end{equation}
Use \Cref{lem:V135-two-cluster-identity} to solve for the top cumulant.
There are only fourteen correction partitions.  Choose in each correction
one block \(B\) meeting both sides of the cut.  If \(|B|=2\), V50
covariance bounds it by the corresponding cross link.  If \(|B|=3\),
the V50 ternary operator tree bound contains two links; because \(B\)
meets both sides, every tree on \(B\) has a cross-cut edge, while its other
link is bounded by the universal cap branch.  Hence the three-row
cumulant is at most \(C\Lambda(C,D)\).  All remaining block cumulants and
singleton moments have universal order-\(\le3\) caps.  If a correction has
two crossing pair blocks, their product is at most a constant times the
cross sum after each factor is capped by one.

The finite correction sum and
\eqref{eq:V135-composite-covariance} give
\(\|\kappa_4\|_{\rm op}\le C\Lambda(C,D)\).
Combine this with \Cref{lem:V135-universal-cap} to obtain the minimum.
All cited covariance and cumulant estimates are stable under identity
amplification.
\end{proof}

Now let \(C_1,C_2,C_3\) be a three-block partition of the four rows.  One
block, say \(C_1\), has two rows.  Write \(G_1=F_{C_1}\) and let \(G_2,G_3\)
be the two singleton row functions.

\begin{lemma}[Exact three-cluster reduction]
\label{lem:V135-three-cluster-reduction}
If \(C_1=\{i,j\}\), then
\begin{align}
 \kappa_3(F_iF_j,G_2,G_3)
 ={}&
 \kappa_4(F_i,F_j,G_2,G_3)\notag\\
 &+\mathbb EF_i\,\kappa_3(F_j,G_2,G_3)
  +\mathbb EF_j\,\kappa_3(F_i,G_2,G_3)\notag\\
 &+\operatorname{Cov}(F_i,G_2)\operatorname{Cov}(F_j,G_3)
  +\operatorname{Cov}(F_i,G_3)\operatorname{Cov}(F_j,G_2).
 \label{eq:V135-three-cluster-identity}
\end{align}
\end{lemma}

\begin{proof}
Expand the five terms in the third cumulant on the left into moments of
the four individual rows.  Equivalently, apply the
Leonov--Shiryaev product-cumulant formula and list the partitions which
become connected after \(i,j\) are identified.  They are the top
four-block, the two triple-plus-singleton terms, and the two crossing
pairings displayed above, each with coefficient one.
\end{proof}

\begin{theorem}[Three-cluster tree bound]
\label{thm:V135-three-cluster-bound}
For a three-block row partition,
\begin{equation}
 \boxed{
 \|\kappa_4(F_1,F_2,F_3,F_4)\|_{\rm op}
 \le
 C_4
 \sum_{\sigma\in\mathfrak T(\{C_1,C_2,C_3\})}
 \prod_{CD\in E(\sigma)}\Lambda(C,D).}
 \label{eq:V135-three-cluster-bound}
\end{equation}
\end{theorem}

\begin{proof}
Solve \eqref{eq:V135-three-cluster-identity} for the fourth cumulant.
Apply the V50 three-variable operator tree bound to
\(\kappa_3(G_1,G_2,G_3)\).  The influence of \(G_1=F_iF_j\) is at most the
sum of the two row influences, so its two cluster link weights are bounded
by \(\Lambda(C_1,C_2)\) and \(\Lambda(C_1,C_3)\).  Each of the two
three-row correction cumulants is bounded by the same three-cluster tree
polynomial, with one possible row inside \(C_1\) absent.  The two covariance
products are exactly two-edge three-cluster trees.  Singleton means are
contractions.  A fixed finite sum proves the display.
\end{proof}

\section{Threshold components}
\label{sec:V135-threshold}

Choose a fixed \(0<\delta_0<1\), to be made sufficiently small below, and
let \(G_{\delta_0}\) be the graph on \(I\) with edge \(ij\) present when
\(\Lambda_{ij}\ge\delta_0\).  Let \(k\) be its number of connected
components.

\begin{lemma}[Strong spanning forest inside threshold components]
\label{lem:V135-strong-forest}
Inside each component choose a spanning tree.  Their union is a forest
with \(4-k\) edges, each of weight at least \(\delta_0\).
\end{lemma}

\begin{proof}
Every connected finite graph has a spanning tree.  A tree on a component
with \(n\) vertices has \(n-1\) edges; summing gives \(4-k\).
\end{proof}

\begin{lemma}[All-weak links imply the V134 small-activity branch]
\label{lem:V135-all-weak-low-activity}
There is a group-dependent numerical \(\delta_0>0\) such that, if
\(\Lambda_{ij}<\delta_0\) for every pair, then the total local activity
\(\mathcal A\) of V134 is at most one and
\begin{equation}
 \mathcal S_4\le C_4\mathcal T_4(\Lambda).
 \label{eq:V135-skeleton-to-tree}
\end{equation}
\end{lemma}

\begin{proof}
The V50 pair estimate bounds the sum of pair activities by
\(C\sum_{i<j}\Lambda_{ij}\).  Its junction-to-tree estimate bounds every
summed triple activity by a fixed sum of two-link products
\(C\Lambda_{ij}\Lambda_{ik}\).  V41 bounds the sum of local four-block
activities by a fixed sum of three-link tree products, and a double-pair
activity is bounded by a two-link product.  Hence
\[
 \mathcal A
 \le C_{\rm act}\left(
    \sum_{i<j}\Lambda_{ij}
   +\Bigl(\sum_{i<j}\Lambda_{ij}\Bigr)^2
   +\Bigl(\sum_{i<j}\Lambda_{ij}\Bigr)^3
 \right).
\]
Choose \(\delta_0\) so that the right side is at most one whenever all six
weights are below \(\delta_0\).

The same local routing sends a \([4]\) skeleton to a three-edge tree, a
\([3,2]\) skeleton to two internal triple links plus its attachment, two
triples to a connected four-edge graph from which one spanning tree is
selected, and a pair-only skeleton to its own tree.  Because every
\(\Lambda_{ij}<1\), a surplus edge may be discarded by enlarging the
majorant.  Summing the fixed row choices proves
\eqref{eq:V135-skeleton-to-tree}.
\end{proof}

\begin{theorem}[Global quartic operator tree numerator]
\label{thm:V135-global-quartic-numerator}
Uniformly in the finite support, volume, representations, and
intermediate/final multiplicities before the final resolvent,
\begin{equation}
 \boxed{
 \|\mathbf K_I^{\rm conn}\|_{\rm op}
 \le C_{\rm qtree}\mathcal T_4(\Lambda).}
 \label{eq:V135-global-quartic-numerator}
\end{equation}
The constant depends only on the fixed group and the earlier local
operator constants.
\end{theorem}

\begin{proof}
Use the components of \(G_{\delta_0}\).

If \(k=1\), \Cref{lem:V135-strong-forest} gives a spanning tree all of
whose edges have weight at least \(\delta_0\).  Hence
\(\mathcal T_4(\Lambda)\ge\delta_0^3\), and
\Cref{lem:V135-universal-cap} proves the result.

If \(k=2\), apply \Cref{thm:V135-weak-cut} to the two threshold
components.  Every crossing edge has weight below \(\delta_0\), and
\[
 \|\mathbf K_I^{\rm conn}\|_{\rm op}
 \le C\sum_{\substack{i\in C_1\\j\in C_2}}\Lambda_{ij}.
\]
For each crossing pair \(ij\), adjoin \(ij\) to the two fixed internal
spanning trees.  The result is a labeled spanning tree on \(I\), and its
weight is at least \(\delta_0^{\,2}\Lambda_{ij}\), because the internal
forest has \(4-k=2\) edges.  Sum the at most four crossing choices and
absorb \(\delta_0^{-2}\).

If \(k=3\), use \Cref{thm:V135-three-cluster-bound}.  Expand every quotient
link \(\Lambda(C_r,C_s)\) into its original row pairs.  Adjoin the one
internal threshold edge from
\Cref{lem:V135-strong-forest} to each two-edge quotient tree.  This gives
an original four-row spanning tree of weight at least \(\delta_0\) times
the corresponding quotient monomial.  The expansion has a fixed number
of terms, so \(\delta_0^{-1}\mathcal T_4(\Lambda)\) bounds the result.

If \(k=4\), every pair weight is below \(\delta_0\).
\Cref{lem:V135-all-weak-low-activity} permits
\Cref{thm:V134-low-activity-quartic} and gives
\[
 \|\mathbf K_I^{\rm conn}\|_{\rm op}
 \le C\mathcal S_4
 \le C\mathcal T_4(\Lambda).
\]
The four cases are exhaustive.  All constants are numerical after
\(\delta_0\) is fixed.
\end{proof}

\begin{corollary}[The coordinate-word numerator bottleneck is closed]
\label{cor:V135-numerator-closed}
The exact V39 quartic coordinate-word sum, including all collision
multiplicities and all four V40 topology classes, admits a
support-uniform operator tree numerator.  No coordinate-subset, Bell,
partition-word, or final-channel multiplicity factor is introduced at
this numerator stage.
\end{corollary}

\begin{proof}
The V39 word sum is the complete connected operator
\(\mathbf K_I^{\rm conn}\).  Apply
\Cref{thm:V135-global-quartic-numerator}.  The proof used complete operator
moments/cumulants and the multiplicity-uniform V37/V41/V50 local cards,
not entrywise output sums.
\end{proof}

\begin{firewall}[The numerator is not the normalized MTA theorem]
\label{fire:V135-numerator-not-MTA}
Equation \eqref{eq:V135-global-quartic-numerator} has unnormalized link
weights \(s_\alpha H_{ij}\).  The BFA target contains repeated factors
\(H_{ij}/(\Xi_i\Xi_j)\), exact input coefficient norms, and one final
resolvent.  Those denominators cannot be attached after taking the
operator norm in \eqref{eq:V135-global-quartic-numerator}.  They must be
inserted into the complete raw output and heavy/light partition using the
V46--V129 capacity machinery.  Thus global normalized quartic MTA remains
open.
\end{firewall}

\section{V135 ledger and next acquisition}
\label{sec:V135-ledger}

\begin{theorem}[Integrated V135 statement]
\label{thm:V135-integrated}
Preserving V1--V134, V135 proves:
\begin{enumerate}[label=\textup{(V135.\arabic*)},leftmargin=7.5em]
\item the mandatory SM V154/NS V31 audit and strict transfer firewall;
\item the exact two-cluster product-cumulant identity
      \eqref{eq:V135-product-cumulant};
\item the complete weak-cut bound
      \eqref{eq:V135-weak-cut};
\item the exact three-cluster reduction
      \eqref{eq:V135-three-cluster-identity};
\item the three-cluster tree estimate
      \eqref{eq:V135-three-cluster-bound};
\item the threshold-component four-case theorem; and
\item the global quartic operator numerator
      \eqref{eq:V135-global-quartic-numerator}.
\end{enumerate}
\end{theorem}

\begin{proof}
The items are
\Cref{ass:V135-audit,fire:V135-audit-boundary,
lem:V135-two-cluster-identity,thm:V135-weak-cut,
lem:V135-three-cluster-reduction,thm:V135-three-cluster-bound,
lem:V135-strong-forest,lem:V135-all-weak-low-activity,
thm:V135-global-quartic-numerator}.
\end{proof}

\begin{openproblem}[V136 mandate: normalize the assembled quartic tree]
\label{op:V136-mandate}
\begin{enumerate}
\item Reinsert the exact Fourier coefficient
      \(d_{\rm in}/d_{\boldsymbol U}\), complete raw pinching, and the
      single final resolvent before using
      \eqref{eq:V135-global-quartic-numerator}.
\item Apply the V94/V98 physical priority partition to each of the sixteen
      global tree orientations, not to coordinate words.
\item Use V126--V129 for every collision branch on which two endpoint
      denominators must be recovered jointly.
\item For noncollision and local \([4]\)/\([3,3]\) branches, identify the
      already-proved V41/V42/V50/V56 spectator quotient which supplies the
      repeated input-mass marks.
\item Produce one normalized global quartic MTA inequality, or isolate the
      exact topology/output family not covered by those quotient cards.
\end{enumerate}
\end{openproblem}

\begin{firewall}[Claim boundary after V135]
\label{fire:V135-current-boundary}
V135 closes the full support-uniform quartic operator numerator, including
the high-activity regime and all coordinate-word aggregation.  It does not
yet prove the normalized quartic MTA estimate, all-order MTA, WUFC/CPM,
continuum Yang--Mills existence, or a positive spectral gap.  The full
Yang--Mills existence and mass-gap problem remains open.
\end{firewall}

\section{Append-only record for V135}
\label{sec:V135-record}

V135 was appended to the complete V134 manuscript.  The exact V134 parent
had \(88079\) newline-delimited source records, \(3066299\) bytes, and
SHA--256 digest
\[
 \texttt{b8072a1f7815e3af2bbe2360b7589800c1268b3f15ad5e449fdd3794a0e66b5b}.
\]
No V134 mathematical content was changed.  The sole final
\verb|\end{document}| is restored below.

% =============================================================================
% END APPENDED VERSION V135 MATERIAL
% =============================================================================

\end{document}
